\newcommand{\runningtitle}{Contents}
% Shared course layout; set a running title per unit.
% This is a small hand-maintained TeX layout, not a conversion pipeline.
\RequirePackage{fix-cm}
\documentclass[11pt,twoside]{article}
\usepackage{fontspec}
\setmainfont{Cambria}[
  Renderer=HarfBuzz,
  BoldFont=cambriab.ttf,
  ItalicFont=cambriai.ttf,
  BoldItalicFont=cambriaz.ttf
]
\setmonofont{lmmono10-regular.otf}[BoldFont=lmmonolt10-bold.otf]
\usepackage[paperwidth=165mm,paperheight=241mm,inner=21mm,outer=19mm,
  top=22mm,bottom=20mm,headheight=14pt,headsep=7mm,footskip=10mm]{geometry}
\usepackage{xcolor,array,booktabs,longtable,tabularx}
\usepackage{enumitem,fancyhdr}
\usepackage[unicode,hidelinks]{hyperref}
\definecolor{ink}{HTML}{000000}
\definecolor{muted}{HTML}{4C5558}
\setlength{\parindent}{0pt}
\setlength{\parskip}{3.5pt}
\AtBeginDocument{\fontsize{11}{14.3}\selectfont}
\setlength{\tabcolsep}{4pt}
\renewcommand{\arraystretch}{1.10}
\setlist[itemize]{leftmargin=1.3em,label=\textbullet,itemsep=2pt,topsep=3pt,parsep=0pt}
\setlist[enumerate]{leftmargin=1.8em,itemsep=3pt,topsep=4pt,parsep=0pt}
\widowpenalty=10000
\clubpenalty=10000
\emergencystretch=1em
\raggedbottom
\pagestyle{fancy}
\fancyhf{}
\providecommand{\runningtitle}{Old English}
\fancyhead[LE,RO]{\small\itshape \runningtitle}
\fancyhead[LO,RE]{\small Old English}
\fancyfoot[LE,RO]{\small\thepage}
\renewcommand{\headrulewidth}{0.3pt}
\newcommand{\booktitle}[2]{
  {\fontsize{9}{11}\selectfont\color{muted}#1\par}
  \vspace{3pt}
  {\fontsize{20}{24}\selectfont\color{ink}#2\par}
  \vspace{6pt}
}
\newcommand{\source}[1]{\par{\footnotesize\color{muted}#1\par}}
\newcommand{\readingtext}[1]{
  \par\addvspace{5pt}
  \begingroup\leftskip=1.5em
    \fontsize{12}{16}\selectfont #1\par
  \endgroup\addvspace{5pt}
}
\makeatletter
\renewcommand{\section}{\@startsection{section}{1}{\z@}
  {-2.4ex \@plus -0.6ex \@minus -0.2ex}{0.8ex \@plus 0.2ex}
  {\normalfont\fontsize{12}{15}\selectfont\bfseries\color{ink}}}
\renewcommand{\subsection}{\@startsection{subsection}{2}{\z@}
  {-1.8ex \@plus -0.5ex \@minus -0.2ex}{0.5ex \@plus 0.2ex}
  {\normalfont\fontsize{11}{14}\selectfont\bfseries\color{ink}}}
\makeatother

\usepackage{multicol}
\hypersetup{pdftitle={Old English: Learner’s Book},pdfsubject={A grammar-led course},bookmarksopen=false,bookmarksnumbered=false}
\hyphenpenalty=10000\exhyphenpenalty=10000\emergencystretch=3em
\makeatletter
\renewcommand{\@pnumwidth}{2.3em}\renewcommand{\@tocrmarg}{3em}
\def\toclevel@section{0}\def\toclevel@subsection{1}\def\toclevel@subsubsection{2}
\newcommand{\BookHiddenContents}[3]{\addtocontents{toc}{\protect\BookHiddenLine{#1}{#2}{#3}}\Hy@writebookmark{}{#2}{\@currentHref}{#1}{toc}}
\makeatother
\newcommand{\BookHiddenLine}[3]{}
\newcommand{\BookRecto}{\clearpage\ifodd\value{page}\else\thispagestyle{empty}\null\clearpage\fi}
\newskip\BookSavedSkip
\newcommand{\BookMeta}[1]{\ifvmode\begingroup\BookSavedSkip=\lastskip\ifdim\BookSavedSkip=0pt #1\else\unskip #1\vskip\BookSavedSkip\fi\endgroup\else #1\fi}
\newcommand{\BookLabel}[1]{\BookMeta{\label{#1}}}
\newcommand{\BookPoint}[1]{\phantomsection\label{#1}}
\newcommand{\BookComponent}[2]{\phantomsection\label{component-#1}\addcontentsline{toc}{section}{#2}}
\newcommand{\BookHeading}[4]{\BookMeta{\label{#1}\ifnum#4=1\addcontentsline{toc}{#2}{#3}\else\BookHiddenContents{\csname toclevel@#2\endcsname}{#3}{\thepage}\fi}}
\newcommand{\BookNeedspace}[1]{\par\begingroup\skip0=\lastskip\vskip-\skip0\dimen0=#1\relax\advance\dimen0 by\skip0\vskip0pt plus\dimen0\penalty-100\vskip0pt plus-\dimen0\vskip\dimen0\penalty9999\vskip-\dimen0\vskip\skip0\endgroup}
\begin{document}\raggedright\pagenumbering{roman}
\thispagestyle{empty}\vspace*{32mm}
\booktitle{LEARNER’S BOOK}{Old English}
{\fontsize{16}{20}\selectfont A grammar-led course\par}
\vfill Cook / Lambdin-style course\par November 2026\par\clearpage
\renewcommand{\runningtitle}{About this volume}
\section*{About this volume}
This volume contains the pronunciation guide, eighteen grammar units, the optional sound-history appendix, and both reference glossaries. Answers to the exercises are collected in the separate Answer Key.\par
The table of contents lists the units and their main sections. The grammatical index points to explanations and paradigms, including explicitly limited recognition material. Its references do not extend the scope stated in a lesson.\par
Section numbers within a unit remain local to that unit. Cook’s exercise and item numbers are retained. Contents and index page numbers refer to this volume; references to printed source pages and scans retain their original locators.\par
The course uses Albert S. Cook’s Exercises in Old English and the second edition of A First Book in Old English (1895). Its practical organization follows the grammar-led model of Lambdin’s Gothic course. Source notes identify the evidence for individual readings, forms and explanations.\par
\clearpage\renewcommand{\runningtitle}{Contents}
\setcounter{tocdepth}{2}
\tableofcontents
\BookRecto\pagenumbering{arabic}
\BookRecto
\begingroup
\renewcommand{\runningtitle}{Pronunciation}
\hyphenpenalty=10000
\exhyphenpenalty=10000
\emergencystretch=3em
\newcommand{\reserve}[1]{\par\penalty-100\begingroup\dimen0=\pagegoal\advance\dimen0 by-\pagetotal\ifdim\dimen0<#1\newpage\fi\endgroup}
\newcommand{\itempara}[2]{\par\addvspace{3pt}\begingroup\raggedright\noindent\hangindent=1.8em\hangafter=1\makebox[1.8em][l]{#1.}#2\par\endgroup}
% book-source-body:begin

\addfontfeatures{Ligatures=TeXOff}
\raggedright
\booktitle{\BookComponent{pronunciation-learner}{A guide to Old English pronunciation}GUIDE}{A guide to Old English pronunciation}

% mdparagraph:P001:begin
This guide gives a practical pronunciation for reading the course. English examples refer to a particular sound in an ordinary American pronunciation, not to the whole word. Where English has no equivalent, the instructions describe the position of the tongue and lips. Every Old English word is glossed; its meaning need not be memorized.
% mdparagraph:P001:end

\reserve{4\baselineskip}
% mdparagraph:P002:begin
\section*{1. Letters and length}\BookHeading{nav-pronunciation-learner-0001}{subsection}{1. Letters and length}{1}\BookLabel{bk-pronunciation-learner-156b3d5ac4b3}
% mdparagraph:P002:end

% mdparagraph:P003:begin
\textbf{Æ æ}, called \emph{ash}, is one vowel letter. \textbf{Þ þ}, \emph{thorn}, and \textbf{Ð ð}, \emph{eth}, represent the sounds written \textbf{th} in English \emph{thin} and \emph{this}. Either letter can represent either sound. This course normally uses \textbf{ð}. Cook's hooked \textbf{ę} and \textbf{ǫ} have the sounds of \textbf{e} and \textbf{o}; the hooks record their history, not additional sounds.[1]
% mdparagraph:P003:end

% mdparagraph:P004:begin
A \textbf{macron}, the line above a vowel, marks length: \textbf{ā, ǣ, ē, ī, ō, ū, ȳ}. Hold a long vowel appreciably longer than a short one. Long does not mean the English letter name, and a macron does not mark stress. Thus \textbf{ac}, “but”, has a short vowel; \textbf{āc}, “oak”, has a long one. Both end with the sound of \textbf{k}.
% mdparagraph:P004:end

% mdparagraph:P005:begin
In a long two-letter vowel spelling, the macron stands on the first letter: \textbf{ēa, ēo, īe}. The two letters belong to one syllable. Preserve the printed marks; a following pair of consonants does not turn an unmarked vowel into a long one.
% mdparagraph:P005:end

\reserve{4\baselineskip}
% mdparagraph:P006:begin
\section*{2. Simple vowels}\BookHeading{nav-pronunciation-learner-0002}{subsection}{2. Simple vowels}{1}\BookLabel{bk-pronunciation-learner-c95193d444e8}
% mdparagraph:P006:end

% mdparagraph:P007:begin
Use these short and long values. Some of the practical English approximations differ in quality as well as duration.[2]
% mdparagraph:P007:end

% mdtable:T01:begin
% course-table-style: 2026-10
% course-panel:T01-P01:begin
\par\begingroup
\fontsize{10.5}{13.2}\selectfont
\setlength{\LTleft}{0pt}\setlength{\LTright}{\fill}
\noindent
\begin{tabular}{@{}>{\raggedright\arraybackslash}p{15.099mm}>{\raggedright\arraybackslash}p{52.318mm}>{\raggedright\arraybackslash}p{50.959mm}@{}}
\toprule
\textbf{Letters} & \textbf{Short vowel} & \textbf{Long vowel} \\
\midrule
a, ā & The vowel of \emph{father}, brief: \textbf{ac}, “but”. & Hold the same vowel: \textbf{āc}, “oak”. \\
æ, ǣ & The vowel of \emph{cat}: \textbf{sæt}, “sat”. & Hold that same vowel: \textbf{sǣd}, “seed”. \\
e, ē & The vowel of \emph{bet}: \textbf{męte}, “food”. & Hold the beginning of \emph{day} without its final glide: \textbf{hē}, “he”. \\
i, ī & The vowel of \emph{bit}: \textbf{is}, “is”. & The vowel of \emph{see}, held: \textbf{īs}, “ice”. \\
o, ō & A brief rounded vowel, described below: \textbf{god}, “god”. & Hold the beginning of \emph{go} without its final glide: \textbf{gōd}, “good”. \\
u, ū & The vowel of \emph{foot}: \textbf{full}, “full”. & The vowel of \emph{food}, held: \textbf{fūl}, “vile”. \\
y, ȳ & The tongue position of \emph{bit} with rounded lips: \textbf{wyrt}, “herb; plant”. & The tongue position of \emph{see} with rounded lips, held: \textbf{fȳr}, “fire”. \\
\bottomrule
\end{tabular}
\par\endgroup
% course-panel:T01-P01:end
% mdtable:T01:end

% mdparagraph:P008:begin
For \textbf{o}, the vowel of \emph{caught} is a useful starting point if you distinguish \emph{cot} from \emph{caught}. If they sound alike, begin with the first part of \emph{go}, lower the jaw a little while keeping the lips rounded, and make the vowel brief. Keep \textbf{o} distinct from the unrounded \emph{father} vowel. For \textbf{ō}, keep the jaw and tongue steady: do not finish with an \textbf{oo} movement.
% mdparagraph:P008:end

% mdparagraph:P009:begin
For \textbf{ē}, likewise stop before the glide at the end of \emph{day}. Do not pronounce \textbf{ǣ} like American \emph{care}: it is the \emph{cat} vowel held longer, without an \textbf{r} sound.
% mdparagraph:P009:end

% mdparagraph:P010:begin
In a \textbf{stressed} syllable, short \textbf{a} before \textbf{m} or \textbf{n} takes the more rounded short-\textbf{o} quality just described. Use this in \textbf{mann}, “man; human being”, and \textbf{lang}, “long”. It is not the vowel of ordinary American \emph{man}. This rule does not require the same rounding in every unstressed \textbf{-an} ending.
% mdparagraph:P010:end

% mdparagraph:P011:begin
For \textbf{y} and \textbf{ȳ}, keep the tongue forward while rounding the lips as for \textbf{oo} in \emph{food}. For short \textbf{y}, begin with the tongue position of \emph{bit}; for long \textbf{ȳ}, use \emph{see}. Rounding the lips must not pull the tongue back into an ordinary \textbf{oo} vowel. These sounds have no usual American English equivalent. Here \textbf{y} is a vowel, not the consonant at the beginning of \emph{yes}.
% mdparagraph:P011:end

\reserve{4\baselineskip}
% mdparagraph:P012:begin
\section*{3. Diphthongs}\BookHeading{nav-pronunciation-learner-0003}{subsection}{3. Diphthongs}{1}\BookLabel{bk-pronunciation-learner-1212ca9f676b}
% mdparagraph:P012:end

% mdparagraph:P013:begin
A \emph{diphthong} is a vowel that moves from one position toward another within one syllable. For these spellings, make one smooth movement toward \textbf{oo}. Do not give the second letter a separate beat.[3]
% mdparagraph:P013:end

% mdtable:T02:begin
% course-table-style: 2026-10
% course-panel:T02-P01:begin
\par\begingroup
\fontsize{10.5}{13.2}\selectfont
\setlength{\LTleft}{0pt}\setlength{\LTright}{\fill}
\noindent
\begin{tabular}{@{}>{\raggedright\arraybackslash}p{19.374mm}>{\raggedright\arraybackslash}p{63.45mm}>{\raggedright\arraybackslash}p{35.552mm}@{}}
\toprule
\textbf{Spellings} & \textbf{Movement} & \textbf{Supplied examples} \\
\midrule
ea, ēa & Start with the vowel of \emph{cat}, then round toward \textbf{oo}. & \textbf{earm}, “arm”; \textbf{ēare}, “ear”. \\
eo, ēo & Start with the steady first vowel of \emph{day}, then round toward \textbf{oo}. & \textbf{eorl}, “man”; \textbf{fēoll}, “fell”. \\
ie, īe & Start with the vowel of \emph{see}, then round toward \textbf{oo}. & \textbf{giefan}, “give”; \textbf{hīeran}, “hear”. \\
\bottomrule
\end{tabular}
\par\endgroup
% course-panel:T02-P01:end
% mdtable:T02:end

% mdparagraph:P014:begin
In \textbf{ēa, ēo, īe}, hold the first part longer before making the movement. The \textbf{eo} of \textbf{eorl} is short; the \textbf{ēo} of \textbf{fēoll} is long. The word \textbf{ēare} has two syllables: its long \textbf{ēa}, followed by a light final \textbf{e}. The opening \textbf{g} of \textbf{giefan} has the consonant sound of \textbf{y} in \emph{yes}.
% mdparagraph:P014:end

% mdparagraph:P015:begin
There are word-specific exceptions. In \textbf{geong}, “young”, \textbf{e} does not supply a vowel sound: the initial \textbf{g} is like \textbf{y} in \emph{yes}, followed by short rounded \textbf{o}. Its \textbf{ng} has the sound heard in \emph{finger}, including the \textbf{g}. Do not give this word the ordinary \textbf{eo} glide or copy the vowel of modern \emph{young}.
% mdparagraph:P015:end

\reserve{4\baselineskip}
% mdparagraph:P016:begin
\section*{4. Consonants}\BookHeading{nav-pronunciation-learner-0004}{subsection}{4. Consonants}{1}
% mdparagraph:P016:end

% mdparagraph:P017:begin
Read \textbf{b, d, l, m, n, p, t, w} as at the beginning of \emph{boy, dog, lip, man, no, pen, top, win}. Use a pronounced American \textbf{r}, as in \emph{run}, as a practical approximation. The remaining letters and combinations need the following distinctions.[4]
% mdparagraph:P017:end

\reserve{4\baselineskip}
% mdparagraph:P018:begin
\subsection*{c, g, and combinations}\BookHeading{nav-pronunciation-learner-0005}{subsubsection}{c, g, and combinations}{0}\BookLabel{bk-pronunciation-learner-16b43725bb81}
% mdparagraph:P018:end

% mdparagraph:P019:begin
Work in the following order. A supplied pronunciation for a particular word or form takes precedence over a general pattern.
% mdparagraph:P019:end

% mdparagraph:P020:begin
\textbf{1. Read a combination first.} Do not apply a rule for single \textbf{c} or \textbf{g} to one letter of these combinations.
% mdparagraph:P020:end

% mdtable:T03:begin
% course-table-style: 2026-10
% course-panel:T03-P01:begin
\par\begingroup
\fontsize{10.5}{13.2}\selectfont
\setlength{\LTleft}{0pt}\setlength{\LTright}{\fill}
\noindent
\begin{tabular}{@{}>{\raggedright\arraybackslash}p{15.233mm}>{\raggedright\arraybackslash}p{35.876mm}>{\raggedright\arraybackslash}p{67.267mm}@{}}
\toprule
\textbf{Spelling} & \textbf{Ordinary sound} & \textbf{Model and supplied exception} \\
\midrule
sc & \textbf{sh} as in \emph{ship}, including before a back vowel. & \textbf{scip}, “ship”, has one \textbf{sh} sound. The exception \textbf{frosc}, “frog”, ends with \textbf{s} followed by \textbf{k}. \\
cg; usually gg & The final sound of \emph{bridge}. & \textbf{ecg}, “edge”. The exception \textbf{docga}, “dog”, has a prolonged hard \textbf{g}. \\
ng & As in \emph{finger}, with audible \textbf{g}. & \textbf{lang}, “long”. In the supplied exception \textbf{sengan}, “singe”, \textbf{ng} has the sound at the end of \emph{singe}. \\
\bottomrule
\end{tabular}
\par\endgroup
% course-panel:T03-P01:end
% mdtable:T03:end

% mdparagraph:P021:begin
For the rules below, the \textbf{back-vowel spellings} are \textbf{a/\allowbreak{}ā, o/\allowbreak{}ō, u/\allowbreak{}ū}. The relevant \textbf{front-vowel spellings} are \textbf{e/\allowbreak{}ē, i/\allowbreak{}ī, æ/\allowbreak{}ǣ}; the diphthongs \textbf{ea/\allowbreak{}ēa, eo/\allowbreak{}ēo, ie/\allowbreak{}īe} also begin at the front of the mouth. Although \textbf{y/\allowbreak{}ȳ} are front vowels in articulation, their spelling is not a reliable instruction to give \textbf{c} or \textbf{g} a front sound. Use the supplied pronunciation in such words.
% mdparagraph:P021:end

% mdparagraph:P022:begin
\textbf{2. At the beginning of a word, decide c from what immediately follows.}\par\nobreak
% mdparagraph:P022:end

% mdtable:T04:begin
% course-table-style: 2026-10
% course-panel:T04-P01:begin
\par\begingroup
\fontsize{10.5}{13.2}\selectfont
\setlength{\LTleft}{0pt}\setlength{\LTright}{\fill}
\noindent
\begin{tabular}{@{}>{\raggedright\arraybackslash}p{49.566mm}>{\raggedright\arraybackslash}p{27.057mm}>{\raggedright\arraybackslash}p{41.752mm}@{}}
\toprule
\textbf{Following sound or spelling} & \textbf{Ordinary c sound} & \textbf{Model} \\
\midrule
A consonant, or a/\allowbreak{}ā, o/\allowbreak{}ō, u/\allowbreak{}ū & \textbf{k} as in \emph{kit} & \textbf{cnīf}, “knife”; \textbf{cuman}, “come”. \\
e/\allowbreak{}ē, i/\allowbreak{}ī, æ/\allowbreak{}ǣ, or a front diphthong & \textbf{ch} as in \emph{chin} & \textbf{cealf}, “calf”, begins with \textbf{ch}. \\
\bottomrule
\end{tabular}
\par\endgroup
% course-panel:T04-P01:end
% mdtable:T04:end

% mdparagraph:P023:begin
These are useful defaults, with lexical exceptions. \textbf{Cēlan}, “cool”, begins with \textbf{k} despite \textbf{ē}. Both \textbf{c} letters in \textbf{cyrice}, “church”, are \textbf{ch}: learn its initial sound with the word, rather than deducing it from \textbf{y}. Keep all three syllables, including the final \textbf{e}.
% mdparagraph:P023:end

% mdparagraph:P024:begin
\textbf{3. Inside or at the end of a word, use the neighbours of c.} Before another consonant, the ordinary sound is \textbf{k}. Immediately after \textbf{i/\allowbreak{}ī}, use \textbf{ch} when \textbf{c} ends the word or is followed by a vowel other than \textbf{a/\allowbreak{}ā, o/\allowbreak{}ō, u/\allowbreak{}ū}. The second \textbf{c} of \textbf{cyrice} follows this pattern. A following back vowel instead gives the ordinary \textbf{k} pattern.
% mdparagraph:P024:end

% mdparagraph:P025:begin
Compare \textbf{dīc}, “ditch; dike”, with \textbf{dīcas}, “ditches; dikes”. Final \textbf{c} after \textbf{ī} has \textbf{ch}; before the \textbf{a} of the longer form it has \textbf{k}. An ending can therefore change the consonant sound. In other positions \textbf{k} is the ordinary starting point, but some words need a supplied pronunciation. Do not extend the after-\textbf{i/\allowbreak{}ī} rule to every front vowel.
% mdparagraph:P025:end

% mdparagraph:P026:begin
Learn \textbf{sēcan}, “seek”, as a supplied exception: its \textbf{c} is \textbf{ch} as in \emph{chin}, despite the following \textbf{a}. The back-vowel pattern does not override a word's stated pronunciation.
% mdparagraph:P026:end

% mdparagraph:P027:begin
\textbf{4. For single g, first decide whether it has the consonant sound y in \emph{yes}.}\par\nobreak
% mdparagraph:P027:end

% mdtable:T05:begin
% course-table-style: 2026-10
% course-panel:T05-P01:begin
\par\begingroup
\fontsize{10.5}{13.2}\selectfont
\setlength{\LTleft}{0pt}\setlength{\LTright}{\fill}
\noindent
\begin{tabular}{@{}>{\raggedright\arraybackslash}p{46.0mm}>{\raggedright\arraybackslash}p{20.0mm}>{\raggedright\arraybackslash}p{52.376mm}@{}}
\toprule
\textbf{Position} & \textbf{Ordinary g sound} & \textbf{Model} \\
\midrule
At the beginning, immediately before a front vowel or front diphthong & \textbf{y} as in \emph{yes} & \textbf{giefan}, “give”. The prefix \textbf{ge-} also begins with this sound. \\
After a front vowel, at the end or when no back vowel immediately follows & \textbf{y} as in \emph{yes} & \textbf{dæg}, “day”, ends with \textbf{y}. A following consonant does not by itself require hard \textbf{g}. \\
\bottomrule
\end{tabular}
\par\endgroup
% course-panel:T05-P01:end
% mdtable:T05:end

% mdparagraph:P028:begin
Otherwise use the non-\textbf{y} branch below, unless a word has a supplied exceptional pronunciation. At the beginning, a consonant between \textbf{g} and a later front vowel prevents the first row from applying: look at the immediate neighbour. The \textbf{g} rule after a front vowel is broader than the \textbf{c} rule after \textbf{i/\allowbreak{}ī}; keep them distinct. The earlier caution about \textbf{y/\allowbreak{}ȳ} still applies.
% mdparagraph:P028:end

% mdparagraph:P029:begin
\textbf{5. For non-y g, decide between stopped and continuous g.}\par\nobreak
% mdparagraph:P029:end

% mdtable:T06:begin
% course-table-style: 2026-10
% course-panel:T06-P01:begin
\par\begingroup
\fontsize{10.5}{13.2}\selectfont
\setlength{\LTleft}{0pt}\setlength{\LTright}{\fill}
\noindent
\begin{tabular}{@{}>{\raggedright\arraybackslash}p{47.364mm}>{\raggedright\arraybackslash}p{37.919mm}>{\raggedright\arraybackslash}p{33.092mm}@{}}
\toprule
\textbf{Position} & \textbf{Sound} & \textbf{Model} \\
\midrule
At the beginning; after n; in a supplied hard doubled consonant & Stopped \textbf{g}, as in \emph{go} & \textbf{god}, “god”; \textbf{lang}, “long”; \textbf{docga}, “dog”. \\
Elsewhere & The continuous voiced back sound described below & \textbf{fugol}, “bird”: single \textbf{g} between back vowels. \\
\bottomrule
\end{tabular}
\par\endgroup
% course-panel:T06-P01:end
% mdtable:T06:end

% mdparagraph:P030:begin
For\BookPoint{bk-pronunciation-learner-b426b0db6cb3} the continuous \textbf{g}, put the back of the tongue near the position for hard \textbf{g}, but leave a narrow gap. Let air pass continuously while the throat vibrates: the sound can be held, unlike the stopped \textbf{g} of \emph{go}. Its phonetic sign is \textbf{[ɣ]}. Do not omit it or substitute \textbf{y}. First selecting the \textbf{y} or non-\textbf{y} branch prevents this rule from changing the final sound of \textbf{dæg} into continuous \textbf{g}.
% mdparagraph:P030:end

% mdparagraph:P031:begin
These patterns let you decide regular new examples. Where a word or form has an exceptional pronunciation, use its stated sound; spelling alone does not settle every case. The practice below supplies any required exception. No historical derivation is required.
% mdparagraph:P031:end

\reserve{4\baselineskip}
% mdparagraph:P032:begin
\subsection*{h and sounded clusters}\BookHeading{nav-pronunciation-learner-0006}{subsubsection}{h and sounded clusters}{0}\BookLabel{bk-pronunciation-learner-fce112ac8644}
% mdparagraph:P032:end

% mdparagraph:P033:begin
At the beginning of a syllable, \textbf{h} is like English \textbf{h} in \emph{hat}. Later in a syllable it has a stronger, continuous sound made by narrowing the passage above the tongue.[5]
% mdparagraph:P033:end

% mdparagraph:P034:begin
In \textbf{brōhte}, “brought”, raise the back of the tongue toward the position for \textbf{k}, leaving a gap for breath. Keep the throat unvoiced. This back sound is \textbf{[x]}. Both it and the following \textbf{t} are heard.
% mdparagraph:P034:end

% mdparagraph:P035:begin
In \textbf{niht}, “night”, the preceding \textbf{i} brings the tongue farther forward. Raise the middle of the tongue toward the roof of the mouth and let breath pass without voice. This front sound is \textbf{[ç]}. Keep the \textbf{t} too: the Old English word does not have modern English's silent \textbf{gh}. Front vowels generally favour this front articulation; back vowels favour the back one.
% mdparagraph:P035:end

% mdparagraph:P036:begin
Sound the initials in \textbf{hl, hn, hr, hw, wr, wl, cn}. In \textbf{hlāf}, “loaf”, begin with audible breath before \textbf{l}; in \textbf{hwæl}, “whale”, keep \textbf{h} with \textbf{w}. In \textbf{wrītan}, “write”, pronounce \textbf{w} and \textbf{r}; in \textbf{cnīf}, “knife”, pronounce \textbf{k} and \textbf{n}. Join the consonants without putting a vowel between them. The spelling \textbf{cw} has \textbf{k} followed by \textbf{w}. Read \textbf{x} and \textbf{hs} as \textbf{k} followed by \textbf{s}; \textbf{weax}, “wax”, ends with that pair.
% mdparagraph:P036:end

\reserve{4\baselineskip}
% mdparagraph:P037:begin
\section*{5. f, s, and th: voice}\BookHeading{nav-pronunciation-learner-0007}{subsection}{5. f, s, and th: voice}{1}\BookLabel{bk-pronunciation-learner-74b76bb98627}
% mdparagraph:P037:end

% mdparagraph:P038:begin
Compare \emph{fan} with \emph{van}, \emph{sip} with \emph{zip}, and \emph{thin} with \emph{this}. The initial consonant of the second word has voice: the throat vibrates. Vowels and consonants such as \textbf{l, m, n, r} also carry voice.[6]
% mdparagraph:P038:end

% mdparagraph:P039:begin
The ordinary rule for single \textbf{f, s, þ/\allowbreak{}ð} is:
% mdparagraph:P039:end

\begin{itemize}
% mdparagraph:P040:begin
\item At the beginning or end of a word, or beside a voiceless consonant, use the sounds of \emph{fan, sip, thin}.
% mdparagraph:P040:end

% mdparagraph:P041:begin
\item Inside a word, between voiced sounds, use the sounds of \emph{van, zip, this}.
% mdparagraph:P041:end

\end{itemize}
% mdtable:T07:begin
% course-table-style: 2026-10
% course-panel:T07-P01:begin
\par\begingroup
\fontsize{10.5}{13.2}\selectfont
\setlength{\LTleft}{0pt}\setlength{\LTright}{\fill}
\noindent
\begin{tabular}{@{}>{\raggedright\arraybackslash}p{26.306mm}>{\raggedright\arraybackslash}p{26.306mm}>{\raggedright\arraybackslash}p{65.764mm}@{}}
\toprule
\textbf{Word} & \textbf{Meaning} & \textbf{Sound to use} \\
\midrule
cnīf & knife & Final \textbf{f} as in \emph{fan}. \\
cnīfas & knives & \textbf{f} as in \emph{van}, between vowels; final \textbf{s} as in \emph{sip}. \\
rīsan & rise & \textbf{s} as in \emph{zip}, between vowels. \\
hūs & house & Final \textbf{s} as in \emph{sip}. \\
ōðer & other & \textbf{ð} as in \emph{this}, between vowels. \\
ðā & then & Initial \textbf{ð} as in \emph{thin}. \\
\bottomrule
\end{tabular}
\par\endgroup
% course-panel:T07-P01:end
% mdtable:T07:end

% mdparagraph:P042:begin
Thus \textbf{is}, “is”, ends with the \textbf{s} of \emph{sip}, not the \textbf{z} of its modern English equivalent. The letter \textbf{ð} does not by itself tell you to use the voiced sound.
% mdparagraph:P042:end

% mdparagraph:P043:begin
Stress and the boundaries between word parts can preserve exceptions. In \textbf{hǣlðu}, “health”, and \textbf{forðām}, “therefore”, use the unvoiced \textbf{th} of \emph{thin}. Learn these as supplied word-specific readings; the surrounding voiced sounds do not override them.
% mdparagraph:P043:end

\reserve{4\baselineskip}
% mdparagraph:P044:begin
\section*{6. Duration, syllables, and stress}\BookHeading{nav-pronunciation-learner-0008}{subsection}{6. Duration, syllables, and stress}{1}\BookLabel{bk-pronunciation-learner-a18005f2eb82}
% mdparagraph:P044:end

% mdparagraph:P045:begin
Keep vowel length separate from consonant length. In \textbf{full}, “full”, \textbf{u} is short; \textbf{fūl}, “vile”, has long \textbf{ū}. A doubled consonant at the end of a word normally has ordinary single-consonant duration: the final \textbf{ll} of \textbf{full} does not require extra holding.
% mdparagraph:P045:end

% mdparagraph:P046:begin
Between vowels, a doubled consonant has extra duration, like the \textbf{n} across the join in \emph{pen-knife}. Do not insert a vowel between its two letters.[7]
% mdparagraph:P046:end

% mdparagraph:P047:begin
With a doubled stop, hold the closure longer. In \textbf{habban}, “have”, keep the lips closed for \textbf{bb}, then release into the following vowel. There are two syllables. Do not add another vowel or another syllable to pronounce the doubling.
% mdparagraph:P047:end

% mdparagraph:P048:begin
Written unstressed vowels remain audible and distinguishable. In \textbf{męte}, “food”, there are two syllables, and final \textbf{e} is a light \emph{bet} vowel. In \textbf{fugol}, “bird”, there are also two syllables; its \textbf{o} is heard. Do not reduce every ending to the \textbf{uh} at the beginning of modern \emph{about}.
% mdparagraph:P048:end

% mdparagraph:P049:begin
The main stress normally falls on the first syllable, as in \textbf{habban}, “have”. Prefixes complicate this: \textbf{ge-} is unstressed, \textbf{be-} and \textbf{for-} usually are, and a prefixed verb generally stresses its root. When \textbf{ge-} precedes a one-syllable root, the word has two syllables and the root receives the stress. A supplied division identifies the prefix where needed. A macron indicates duration regardless of which syllable is stressed.
% mdparagraph:P049:end

% mdparagraph:P050:begin
For these careful reading exercises, observe printed length. In connected speech, unstressed function words and prefixes could be shorter than their conventional written macrons suggest. This does not authorize adding, removing, or moving marks in the course text.
% mdparagraph:P050:end

\reserve{4\baselineskip}
% mdparagraph:P051:begin
\section*{7. Practice}\BookHeading{nav-pronunciation-learner-0009}{subsection}{7. Practice}{1}
% mdparagraph:P051:end

% mdparagraph:P052:begin
The following tasks are new teaching material. Meanings are supplied throughout. Apply the guide's sound patterns to the new words; no translation or grammatical analysis is required. Give an English sound cue or the taught mouth instruction, not a phonetic transcription.[9]
% mdparagraph:P052:end

\reserve{4\baselineskip}
% mdparagraph:P053:begin
\subsection*{A. Vowels and length}\BookHeading{nav-pronunciation-learner-0010}{subsubsection}{A. Vowels and length}{0}
% mdparagraph:P053:end

% mdparagraph:P054:begin
Read each word aloud. Keep duration separate from stress, and give each diphthong one syllable.
% mdparagraph:P054:end

% item: PRON-P01a
% mdparagraph:P055:begin
\textbf{A1.} \textbf{for}, “for”; \textbf{fōr}, “went”. What changes in the vowel? Give its sound and duration in each word.
% mdparagraph:P055:end

% item: PRON-P01b
% mdparagraph:P056:begin
\textbf{A2.} \textbf{węnde}, “went”; \textbf{wēnde}, “thought; supposed”. Explain the vowel contrast. Does the hook under \textbf{ę} require a sound different from ordinary short \textbf{e}? Keep the final \textbf{e} audible in both words.
% mdparagraph:P056:end

% item: PRON-P01c
% mdparagraph:P057:begin
\textbf{A3.} \textbf{dǣd}, “deed”; \textbf{wīd}, “wide”; \textbf{hȳdan}, “hide; conceal”. Give the sound of each marked vowel. Explain how the lips distinguish \textbf{ȳ} from \textbf{ī}, then read all three words.
% mdparagraph:P057:end

% item: PRON-P01d
% mdparagraph:P058:begin
\textbf{A4.} \textbf{eall}, “all”; \textbf{scēap}, “sheep”; \textbf{heorte}, “heart”; \textbf{frēond}, “friend”; \textbf{scield}, “shield”; \textbf{gīet}, “yet; still”. For each diphthong, give its starting sound and movement, and say whether the first part is short or long. Count the syllables in each word without giving the diphthong two beats.
% mdparagraph:P058:end

\reserve{4\baselineskip}
% mdparagraph:P059:begin
\subsection*{B. Consonants and unfamiliar sounds}\BookHeading{nav-pronunciation-learner-0011}{subsubsection}{B. Consonants and unfamiliar sounds}{0}
% mdparagraph:P059:end

% mdparagraph:P060:begin
Choose the sound of the indicated letters and explain which position or neighbouring sound selects it. Then read the whole word. Check combinations before single letters.
% mdparagraph:P060:end

% item: PRON-P02a
% mdparagraph:P061:begin
\textbf{B1.} \textbf{c} in \textbf{cēosað}, “choose” (a plural form). Which initial-\textbf{c} pattern applies?
% mdparagraph:P061:end

% item: PRON-P02b
% mdparagraph:P062:begin
\textbf{B2.} \textbf{c} in \textbf{micel}, “large; great”, and \textbf{micles}, “large; great” (a declined form). Do the two forms have the same consonant sound? Explain from the immediate neighbours of \textbf{c}.
% mdparagraph:P062:end

% item: PRON-P02c
% mdparagraph:P063:begin
\textbf{B3.} \textbf{g} in \textbf{gān}, “go; walk”, and \textbf{glæd}, “glad”. Does the \textbf{æ} of the second word make its initial \textbf{g} a \textbf{y} sound? Apply both stages of the \textbf{g} decision.
% mdparagraph:P063:end

% item: PRON-P02d
% mdparagraph:P064:begin
\textbf{B4.} \textbf{g} in \textbf{fæger}, “beautiful; fair”, and \textbf{regn}, “rain”. Decide from the sounds before and after \textbf{g}. Does the following \textbf{n} require stopped \textbf{g} in the second word?
% mdparagraph:P064:end

% item: PRON-P02e
% mdparagraph:P065:begin
\textbf{B5.} \textbf{g} in \textbf{āgan}, “have; possess”. Choose among \textbf{y}, stopped \textbf{g}, and continuous \textbf{g}, and explain both decisions. Show with the tongue and airflow how to make the selected sound.
% mdparagraph:P065:end

% item: PRON-P02f
% mdparagraph:P066:begin
\textbf{B6.} \textbf{cg} in \textbf{brycg}, “bridge”, and \textbf{ng} in \textbf{singan}, “sing”. Which English cue fits each combination? Keep the short rounded vowel of \textbf{brycg} distinct from ordinary English \textbf{i}.
% mdparagraph:P066:end

% item: PRON-P02g
% mdparagraph:P067:begin
\textbf{B7.} \textbf{h} in \textbf{riht}, “right; direct”, and \textbf{dohtor}, “daughter”. What changes in the tongue position, and which vowel causes the change? Keep the following \textbf{t} audible in both words.
% mdparagraph:P067:end

% item: PRON-P02h
% mdparagraph:P068:begin
\textbf{B8.} \textbf{sc} in \textbf{scūfan}, “thrust”, and \textbf{cw} in \textbf{cwæð}, “said”. Decide whether each spelling gives one consonant sound or a sequence. Does the back vowel \textbf{ū} require a separate \textbf{k} in the first word?
% mdparagraph:P068:end

\reserve{4\baselineskip}
% mdparagraph:P069:begin
\subsection*{C. Voice}\BookHeading{nav-pronunciation-learner-0012}{subsubsection}{C. Voice}{0}
% mdparagraph:P069:end

% mdparagraph:P070:begin
State whether the indicated sound is voiced or unvoiced. Give its English sound cue and explain the context or supplied exception.
% mdparagraph:P070:end

% item: PRON-P03a
% mdparagraph:P071:begin
\textbf{C1.} \textbf{f} in \textbf{ǣfre}, “ever; always”. Identify the sounds on both sides; they need not both be vowels.
% mdparagraph:P071:end

% item: PRON-P03b
% mdparagraph:P072:begin
\textbf{C2.} Both \textbf{f} letters in \textbf{fīf}, “five”. Decide each from its position rather than copying the pronunciation of modern English \emph{five}.
% mdparagraph:P072:end

% item: PRON-P03c
% mdparagraph:P073:begin
\textbf{C3.} \textbf{s} and final \textbf{ð} in \textbf{cēosað}, “choose” (a plural form). Does one voicing decision apply to both letters?
% mdparagraph:P073:end

% item: PRON-P03d
% mdparagraph:P074:begin
\textbf{C4.} \textbf{f} and \textbf{ð} in \textbf{hæfð}, “has”. Begin with the word-final consonant, then use its voicing to decide the preceding \textbf{f}.
% mdparagraph:P074:end

% item: PRON-P03e
% mdparagraph:P075:begin
\textbf{C5.} \textbf{ð} in \textbf{cwæð}, “said”. Would replacing its spelling with thorn, \textbf{þ}, change the sound? Explain the role of position and the choice of letter.
% mdparagraph:P075:end

% item: PRON-P03f
% mdparagraph:P076:begin
\textbf{C6.} Compare \textbf{ð} in \textbf{hǣlðu}, “health”, with \textbf{f} in \textbf{ǣfre}, “ever; always”. The unvoiced \textbf{th} in \textbf{hǣlðu} is supplied lexical information from §5. Explain why that fact takes precedence over the ordinary between-voiced-sounds rule, while the ordinary rule applies to \textbf{ǣfre}. This is an exception comparison, not a request to deduce the exception from its spelling.
% mdparagraph:P076:end

\reserve{4\baselineskip}
% mdparagraph:P077:begin
\subsection*{D. Syllables and stress}\BookHeading{nav-pronunciation-learner-0013}{subsubsection}{D. Syllables and stress}{0}
% mdparagraph:P077:end

% mdparagraph:P078:begin
Give the number of syllables and identify the stressed one. Identify any doubled consonant and explain how it is pronounced.
% mdparagraph:P078:end

% item: PRON-P04a
% mdparagraph:P079:begin
\textbf{D1.} \textbf{sittan}, “sit”. Does prolonging \textbf{tt} lengthen \textbf{i}, or add a syllable?
% mdparagraph:P079:end

% item: PRON-P04b
% mdparagraph:P080:begin
\textbf{D2.} \textbf{sunne}, “sun”. Keep the last vowel distinct. Explain the duration of \textbf{nn} and whether it changes the syllable count.
% mdparagraph:P080:end

% item: PRON-P04c
% mdparagraph:P081:begin
\textbf{D3.} \textbf{gewrit}, “letter; written document”; division \textbf{ge-writ}. Apply the prefix rule to stress, and pronounce both \textbf{w} and \textbf{r} without inserting another vowel.
% mdparagraph:P081:end

% item: PRON-P04d
% mdparagraph:P082:begin
\textbf{D4.} \textbf{forlǣtan}, “abandon; forsake”; division \textbf{for-lǣtan}. This is a prefixed verb. Locate its stress from the prefix rule, and explain what the macron tells you separately.
% mdparagraph:P082:end

\reserve{4\baselineskip}
% mdparagraph:P083:begin
\subsection*{E. A sentence adapted from Cook}\BookHeading{nav-pronunciation-learner-0014}{subsubsection}{E. A sentence adapted from Cook}{0}
% mdparagraph:P083:end

% mdparagraph:P084:begin
This sentence adapts the opening of Cook's Creation reading.[8]
% mdparagraph:P084:end

% mdparagraph:P085:begin
\textbf{God gescēop heofonan and eorðan.}\par\nobreak
% mdparagraph:P085:end

% mdparagraph:P086:begin
“God created heaven and earth.”
% mdparagraph:P086:end

% mdtable:T08:begin
% course-table-style: 2026-10
% course-panel:T08-P01:begin
\par\begingroup
\fontsize{10.5}{13.2}\selectfont
\setlength{\LTleft}{0pt}\setlength{\LTright}{\fill}
\noindent
\begin{tabular}{@{}>{\raggedright\arraybackslash}p{56.106mm}>{\raggedright\arraybackslash}p{65.082mm}@{}}
\toprule
\textbf{Word} & \textbf{Supplied meaning} \\
\midrule
God & God \\
gescēop & created \\
heofonan & heaven \\
and & and \\
eorðan & earth \\
\bottomrule
\end{tabular}
\par\endgroup
% course-panel:T08-P01:end
% mdtable:T08:end

% mdparagraph:P087:begin
The verb has the division \textbf{ge-scēop}, with prefix \textbf{ge-}. The sentence's meaning and grammatical forms are supplied; apply the pronunciation rules to read it aloud, then answer these sound questions.
% mdparagraph:P087:end

% item: PRON-P05a
% mdparagraph:P088:begin
\textbf{E1.} Which sounds do \textbf{g} and \textbf{sc} have in \textbf{gescēop}, “created”? Explain the rules you use. Count its syllables and place the stress, keeping \textbf{ēo} one glide.
% mdparagraph:P088:end

% item: PRON-P05b
% mdparagraph:P089:begin
\textbf{E2.} What sound does \textbf{f} have in \textbf{heofonan}, “heaven”? Identify its neighbours and explain your decision. Count the word's syllables.
% mdparagraph:P089:end

% item: PRON-P05c
% mdparagraph:P090:begin
\textbf{E3.} What sound does \textbf{ð} have in \textbf{eorðan}, “earth”? Identify its neighbours and explain why the ordinary rule applies even though one neighbour is a consonant.
% mdparagraph:P090:end

\reserve{4\baselineskip}
% mdparagraph:P091:begin
\section*{Source notes}\BookHeading{nav-pronunciation-learner-0015}{subsection}{Source notes}{1}
% mdparagraph:P091:end

% mdparagraph:P092:begin
\itempara{1}{Cook, \emph{A First Book in Old English}, second edition (1895), §§2–4, pp.5–6 /\allowbreak{} PDF21–22; Fulk, \emph{An Introductory Grammar of Old English}, corrected open-access reprint (2020), §§15–16, pp.10–12 /\allowbreak{} PDF22–24. Fulk's teaching dots are omitted in the selected examples; \textbf{hǣlðu, forðām} replace his thorn with the course's eth. These declared spelling changes do not change their selected sounds.}
% mdparagraph:P092:end

% mdparagraph:P093:begin
\itempara{2}{Fulk §§12,16,19, pp.10–11,13–14 /\allowbreak{} PDF22–23,25–26, permits the traditional short-vowel approximations used here. Held \emph{cat} for \textbf{ǣ} follows Fulk. Cook §5, pp.6–7 /\allowbreak{} PDF22–23, also offers practical approximations, but some of his English cues suit other accents. American cue adjustments are teaching aids.}
% mdparagraph:P093:end

% mdparagraph:P094:begin
\itempara{3}{Fulk §§20–21, p.14 /\allowbreak{} PDF26; \textbf{giefan}, §135a, p.103 /\allowbreak{} PDF115. The rounded diphthongs follow Fulk's declared convention. Cook's §5 table, p.7 /\allowbreak{} PDF23, gives different component vowels; the two descriptions are not interchangeable.}
% mdparagraph:P094:end

% mdparagraph:P095:begin
\itempara{4}{Fulk §§16,18,135a–d, pp.11–13,103–104 /\allowbreak{} PDF23–25,115–116; Cook §§7,10, pp.8–9 /\allowbreak{} PDF24–25. \textbf{Cealf, dīc, dīcas} follow Fulk §135b, with his pronunciation overpoints omitted; Cook's vocabulary, p.283 /\allowbreak{} PDF299, also gives \textbf{dīc}. For \textbf{sēcan}, Cook's vocabulary, p.311 /\allowbreak{} PDF327, supplies the form and meaning as a headword variant; Fulk §135b supports the \textbf{ch} value with its explicitly prefixed infinitive example. The conditions for noninitial \textbf{c} are narrower than those for \textbf{g}. Earlier vowel changes, later endings and analogy explain some exceptions to the practical spelling patterns; the sound-history appendix treats these matters separately.}
% mdparagraph:P095:end

% mdparagraph:P096:begin
\itempara{5}{Cook §§7,10, pp.8–9 /\allowbreak{} PDF24–25; Fulk §§9,16,18, pp.8,11–13 /\allowbreak{} PDF20,23–25. American \textbf{r} is a practical approximation. The \textbf{k+s} reading of \textbf{x/\allowbreak{}hs} follows Cook's practical value; Fulk also records an earlier h-like sound before \textbf{s}. Neither source requires a trill in these practical descriptions.}
% mdparagraph:P096:end

% mdparagraph:P097:begin
\itempara{6}{Fulk §17, p.12 /\allowbreak{} PDF24, supplies the voicing rule and exceptions used here. Cook §8, pp.8–9 /\allowbreak{} PDF24–25, gives a broader syllable-final voicing rule. This guide selects Fulk's account.}
% mdparagraph:P097:end

% mdparagraph:P098:begin
\itempara{7}{Cook §11, pp.9–10 /\allowbreak{} PDF25–26, distinguishes prolonged medial doubling from final double spellings, normally read with single-consonant duration. See also Cook §23, pp.18–19 /\allowbreak{} PDF34–35, and Fulk §§17,19,22–23, pp.12–15 /\allowbreak{} PDF24–27. Printed vowel length and stress are distinct; connected-speech shortening does not require changing the edited spelling.}
% mdparagraph:P098:end

% mdparagraph:P099:begin
\itempara{8}{Adapted from Cook, Reader I, p.123 /\allowbreak{} PDF139, first prose line: remove \textbf{On anginne}, move \textbf{God} before the verb, and omit the three apparatus cues. The retained words are not respelled. The original sentence and notes remain in the source transcription. The full meaning and grammatical forms are supplied here; the sound questions apply the guide's rules. The separate \textbf{ge-scēop} division identifies a prefix and does not respell the reading.}
% mdparagraph:P099:end

% mdparagraph:P100:begin
\itempara{9}{The new practice forms are checked against Fulk §§16–24,135, pp.11–15,103–104 /\allowbreak{} PDF23–27,115–116; Cook §§4–5,53,121, pp.6–7,37,74 /\allowbreak{} PDF22–23,53,90, Reader I p.123 /\allowbreak{} PDF139, and vocabulary pp.283,288–289,296,309–312 /\allowbreak{} PDF299,304–305,312,325–328. \textbf{Hȳdan} is printed in Cook, \emph{Exercises in Old English}, V.18.9, p.9 /\allowbreak{} PDF17; the original exercise is unchanged. “Thought; supposed” modernizes Cook's gloss “weened” for \textbf{wēnde}. Dictionary division hyphens in \textbf{gewrit} and \textbf{forlǣtan} are omitted, as elsewhere in the course; their separate practice divisions identify the prefixes.}
% mdparagraph:P100:end


% book-source-body:end
\par\label{end-pronunciation-learner}
\clearpage\endgroup

\BookRecto
\begingroup
\renewcommand{\runningtitle}{Unit 1. Nouns, the article, and the copula}
\hyphenpenalty=10000
\exhyphenpenalty=10000
\emergencystretch=3em
\newcommand{\reserve}[1]{\par\penalty-100\begingroup\dimen0=\pagegoal\advance\dimen0 by-\pagetotal\ifdim\dimen0<#1\newpage\fi\endgroup}
\newcommand{\itempara}[2]{\par\addvspace{3pt}\begingroup\raggedright\noindent\hangindent=1.8em\hangafter=1\makebox[1.8em][l]{#1.}#2\par\endgroup}
% book-source-body:begin

\addfontfeatures{Ligatures=TeXOff}
\raggedright
\booktitle{\BookComponent{ch01-learner}{Unit 1. Nouns, the article, and the copula}UNIT 1}{Nouns, the article, and the copula}

\reserve{4\baselineskip}
% mdparagraph:P001:begin
\section*{1. Spelling and case}\BookHeading{nav-ch01-learner-0001}{subsection}{1. Spelling and case}{1}\BookLabel{bk-ch01-learner-624dd0eb97a3}
% mdparagraph:P001:end

% teaching: CH01-GRAM-ORTHOGRAPHY CH01-GRAM-GENDER CH01-GRAM-CASES
% mdparagraph:P002:begin
\textbf{Æ æ} is a single vowel letter. \textbf{Þ þ} (thorn) and \textbf{Ð ð} (eth) represent the sounds written \emph{th} in modern English; the choice between the letters does not by itself distinguish the sounds of \emph{thin} and \emph{this}. This unit follows Cook's use of \textbf{ð}. His hooked \textbf{ę} is an editorial form of \emph{e}, retained in such words as \textbf{męte} and \textbf{męnn}. A macron marks a long vowel or diphthong: \textbf{ā, ǣ, ē, ī, ō, ū, ȳ; ēa, ēo, īe}. Preserve the marks printed in each word; do not add them by analogy.[1]
% mdparagraph:P002:end

% mdparagraph:P003:begin
Nouns have masculine, feminine, or neuter gender, independently of the sex or nature of the thing named. Learn the gender with the noun. A noun's declension determines its case endings; its nominative ending alone does not identify either gender or declension.
% mdparagraph:P003:end

% mdtable:T01:begin
% course-table-style: 2026-10
% course-panel:T01-P01:begin
\par\begingroup
\fontsize{10.5}{13.2}\selectfont
\setlength{\LTleft}{0pt}\setlength{\LTright}{\fill}
\noindent
\begin{tabular}{@{}>{\raggedright\arraybackslash}p{27.369mm}>{\raggedright\arraybackslash}p{93.819mm}@{}}
\toprule
\textbf{Case} & \textbf{Uses introduced here} \\
\midrule
Nominative (N) & Subject; a noun identifying the subject after “be.” \\
Genitive (G) & A dependent expressing possession or another relation, often English “of.” \\
Dative (D) & Recipient; the noun after the selected prepositional constructions below. \\
Accusative (A) & Direct object. \\
\bottomrule
\end{tabular}
\par\endgroup
% course-panel:T01-P01:end
% mdtable:T01:end

% mdparagraph:P004:begin
Singular and plural forms may coincide, and one spelling may serve several cases. Its construction and the agreement of accompanying words determine the analysis. The tables give all four cases, even when their forms are identical.[2]
% mdparagraph:P004:end

\reserve{4\baselineskip}
% mdparagraph:P005:begin
\section*{2. The article}\BookHeading{nav-ch01-learner-0002}{subsection}{2. The article}{1}\BookLabel{bk-ch01-learner-f6bfdc1fc647}
% mdparagraph:P005:end

% teaching: CH01-GRAM-ARTICLE
% mdparagraph:P006:begin
The definite article means \textbf{the}. It agrees with its own noun in gender, number, and case. The plural forms are common to all three genders.[3]
% mdparagraph:P006:end

% mdtable:T02:begin
% course-table-style: 2026-10
% course-panel:T02-P01:begin
\par\begingroup
\fontsize{10.5}{13.2}\selectfont
\setlength{\LTleft}{0pt}\setlength{\LTright}{\fill}
\noindent
\begin{tabular}{@{}>{\raggedright\arraybackslash}p{22.0mm}>{\raggedright\arraybackslash}p{22.688mm}>{\raggedright\arraybackslash}p{22.688mm}>{\raggedright\arraybackslash}p{22.688mm}>{\raggedright\arraybackslash}p{22.688mm}@{}}
\toprule
\textbf{Case} & \textbf{Masculine singular} & \textbf{Neuter singular} & \textbf{Feminine singular} & \textbf{All genders plural} \\
\midrule
Nominative & se & ðæt & sēo & ðā \\
Genitive & ðæs & ðæs & ðǣre & ðāra \\
Dative & ðǣm & ðǣm & ðǣre & ðǣm \\
Accusative & ðone & ðæt & ðā & ðā \\
Instrumental: reference & ðȳ, ðon & ðȳ, ðon & as dative & as dative \\
\bottomrule
\end{tabular}
\par\endgroup
% course-panel:T02-P01:end
% mdtable:T02:end

% mdparagraph:P007:begin
Recognize \textbf{ðȳ} and \textbf{ðon} as instrumental singular forms; their uses are not required in composition here. Nouns have no separate instrumental forms in these paradigms. The source also permits \textbf{ðām} for \textbf{ðǣm} and \textbf{ðǣra} for \textbf{ðāra}; the tables and exercises select \textbf{ðǣm} and \textbf{ðāra}.
% mdparagraph:P007:end

\reserve{4\baselineskip}
% mdparagraph:P008:begin
\section*{3. Strong masculine nouns}\BookHeading{nav-ch01-learner-0003}{subsection}{3. Strong masculine nouns}{1}\BookLabel{bk-ch01-learner-246e855423d2}
% mdparagraph:P008:end

% teaching: CH01-PAR-MASC CH01-PAR-MASC-E
% mdparagraph:P009:begin
\textbf{Fisc} illustrates the common consonant pattern. In \textbf{męte}, the nominative \textbf{-e} is replaced by \textbf{-es, -as, -a, -um} where those endings occur; it is not part of a stem to which another full ending is added.[4]
% mdparagraph:P009:end

% mdtable:T03:begin
% course-table-style: 2026-10
% course-panel:T03-P01:begin
\par\begingroup
\fontsize{10.5}{13.2}\selectfont
\setlength{\LTleft}{0pt}\setlength{\LTright}{\fill}
\noindent
\begin{tabular}{@{}>{\raggedright\arraybackslash}p{22.0mm}>{\raggedright\arraybackslash}p{22.688mm}>{\raggedright\arraybackslash}p{22.688mm}>{\raggedright\arraybackslash}p{22.688mm}>{\raggedright\arraybackslash}p{22.688mm}@{}}
\toprule
\textbf{Case} & \textbf{fisc singular} & \textbf{fisc plural} & \textbf{męte singular} & \textbf{męte plural} \\
\midrule
Nominative & fisc & fiscas & męte & mętas \\
Genitive & fisces & fisca & mętes & męta \\
Dative & fisce & fiscum & męte & mętum \\
Accusative & fisc & fiscas & męte & mętas \\
\bottomrule
\end{tabular}
\par\endgroup
% course-panel:T03-P01:end
% mdtable:T03:end

% mdparagraph:P010:begin
The plural of \textbf{męte} is displayed for recognition of the pattern; this unit uses its singular “food.” \textbf{Heofon, gāst, wæstm} follow \textbf{fisc}. Keep both syllables of \textbf{heofon} before endings. The divine name \textbf{God} has the singular \textbf{God, Godes, Gode, God}; no plural of this proper-name sense is required.
% mdparagraph:P010:end

% teaching: CH01-PAR-MASC-AE-A
% mdparagraph:P011:begin
In \textbf{dæg} and \textbf{hwæl}, singular \textbf{æ} alternates with plural \textbf{a}. The change does not apply to \textbf{wæstm}, whose \textbf{æ} stands before a consonant group.
% mdparagraph:P011:end

% mdtable:T04:begin
% course-table-style: 2026-10
% course-panel:T04-P01:begin
\par\begingroup
\fontsize{10.5}{13.2}\selectfont
\setlength{\LTleft}{0pt}\setlength{\LTright}{\fill}
\noindent
\begin{tabular}{@{}>{\raggedright\arraybackslash}p{22.0mm}>{\raggedright\arraybackslash}p{22.688mm}>{\raggedright\arraybackslash}p{22.688mm}>{\raggedright\arraybackslash}p{22.688mm}>{\raggedright\arraybackslash}p{22.688mm}@{}}
\toprule
\textbf{Case} & \textbf{dæg singular} & \textbf{dæg plural} & \textbf{hwæl singular} & \textbf{hwæl plural} \\
\midrule
Nominative & dæg & dagas & hwæl & hwalas \\
Genitive & dæges & daga & hwæles & hwala \\
Dative & dæge & dagum & hwæle & hwalum \\
Accusative & dæg & dagas & hwæl & hwalas \\
\bottomrule
\end{tabular}
\par\endgroup
% course-panel:T04-P01:end
% mdtable:T04:end

% teaching: CH01-PAR-MASC-SYNCOPE
% mdparagraph:P012:begin
\textbf{Fugol} and \textbf{morgen} lose their unstressed middle vowel before endings in the selected forms. This loss, called \emph{syncope}, is lexical: use the diagnostic forms given for each noun rather than deleting vowels from every word of two syllables. Forms retaining the vowel also occur.[4]
% mdparagraph:P012:end

% mdtable:T05:begin
% course-table-style: 2026-10
% course-panel:T05-P01:begin
\par\begingroup
\fontsize{10.5}{13.2}\selectfont
\setlength{\LTleft}{0pt}\setlength{\LTright}{\fill}
\noindent
\begin{tabular}{@{}>{\raggedright\arraybackslash}p{22.0mm}>{\raggedright\arraybackslash}p{22.688mm}>{\raggedright\arraybackslash}p{22.688mm}>{\raggedright\arraybackslash}p{22.688mm}>{\raggedright\arraybackslash}p{22.688mm}@{}}
\toprule
\textbf{Case} & \textbf{fugol singular} & \textbf{fugol plural} & \textbf{morgen singular} & \textbf{morgen plural} \\
\midrule
Nominative & fugol & fuglas & morgen & morgnas \\
Genitive & fugles & fugla & morgnes & morgna \\
Dative & fugle & fuglum & morgne & morgnum \\
Accusative & fugol & fuglas & morgen & morgnas \\
\bottomrule
\end{tabular}
\par\endgroup
% course-panel:T05-P01:end
% mdtable:T05:end

\reserve{4\baselineskip}
% mdparagraph:P013:begin
\section*{4. Strong neuter nouns}\BookHeading{nav-ch01-learner-0004}{subsection}{4. Strong neuter nouns}{1}\BookLabel{bk-ch01-learner-737b70e12605}
% mdparagraph:P013:end

% teaching: CH01-GRAM-WEIGHT CH01-PAR-NEUT-WEIGHT
% mdparagraph:P014:begin
Neuter nominative and accusative are identical in each number. The genitive and dative endings match those of \textbf{fisc}. In the ordinary monosyllabic pattern, the nominative/\allowbreak{}accusative plural has \textbf{-u} after a short stem, but no ending after a heavy stem. A short vowel followed by one consonant gives the short pattern here; a long vowel or long diphthong, or a short vowel followed by two consonants, gives the heavy pattern. A doubled consonant counts as two. Thus \textbf{hof} has \textbf{hofu}, whereas \textbf{sǣd, gærs, cynn} have no added N/\allowbreak{}A plural ending. Vowel length and syllable weight are not the same thing.[5]
% mdparagraph:P014:end

% mdtable:T06:begin
% course-table-style: 2026-10
% course-panel:T06-P01:begin
\par\begingroup
\fontsize{10.5}{13.2}\selectfont
\setlength{\LTleft}{0pt}\setlength{\LTright}{\fill}
\noindent
\begin{tabular}{@{}>{\raggedright\arraybackslash}p{22.0mm}>{\raggedright\arraybackslash}p{22.688mm}>{\raggedright\arraybackslash}p{22.688mm}>{\raggedright\arraybackslash}p{22.688mm}>{\raggedright\arraybackslash}p{22.688mm}@{}}
\toprule
\textbf{Case} & \textbf{hof singular} & \textbf{hof plural} & \textbf{sǣd singular} & \textbf{sǣd plural} \\
\midrule
Nominative & hof & hofu & sǣd & sǣd \\
Genitive & hofes & hofa & sǣdes & sǣda \\
Dative & hofe & hofum & sǣde & sǣdum \\
Accusative & hof & hofu & sǣd & sǣd \\
\bottomrule
\end{tabular}
\par\endgroup
% course-panel:T06-P01:end
% mdtable:T06:end

% mdparagraph:P015:begin
The heavy pattern also covers \textbf{lēoht, līf, gēar, dēor, anginn, lyft, weorc}. In \textbf{anginn}, apply the rule to the stressed stem \textbf{-ginn}, with \textbf{nn}. Plural forms of mass or abstract nouns do not automatically have the same sense as the singular; the vocabulary below marks those whose plurals are only for recognizing the morphology.
% mdparagraph:P015:end

% teaching: CH01-PAR-NEUT-TREO
% mdparagraph:P016:begin
\textbf{Trēo} has \textbf{w} before every ending and the exceptional N/\allowbreak{}A plural \textbf{trēowu}, despite its long vowel. Bare singular \textbf{trēow} also occurs.
% mdparagraph:P016:end

% mdparagraph:P017:begin
\textbf{trēo}\par\nobreak
% mdparagraph:P017:end

% mdtable:T07:begin
% course-table-style: 2026-10
% course-panel:T07-P01:begin
\par\begingroup
\fontsize{10.5}{13.2}\selectfont
\setlength{\LTleft}{0pt}\setlength{\LTright}{\fill}
\noindent
\begin{tabular}{@{}>{\raggedright\arraybackslash}p{39.459mm}>{\raggedright\arraybackslash}p{39.459mm}>{\raggedright\arraybackslash}p{39.459mm}@{}}
\toprule
\textbf{Case} & \textbf{Singular} & \textbf{Plural} \\
\midrule
Nominative & trēo & trēowu \\
Genitive & trēowes & trēowa \\
Dative & trēowe & trēowum \\
Accusative & trēo & trēowu \\
\bottomrule
\end{tabular}
\par\endgroup
% course-panel:T07-P01:end
% mdtable:T07:end

% teaching: CH01-PAR-NEUT-DISYLLABLE
% mdparagraph:P018:begin
Do not apply the monosyllabic \textbf{-u/\allowbreak{}zero} rule mechanically to nouns of two syllables. Learn the following complete patterns. \textbf{Tācen} loses \textbf{e} before endings; \textbf{ǣfen, wæter, nīeten} retain it in the selected forms. The shorter \textbf{wætru} and related forms also occur.[5]
% mdparagraph:P018:end

% mdtable:T08:begin
% course-table-style: 2026-10
% course-panel:T08-P01:begin
\par\begingroup
\fontsize{10.5}{13.2}\selectfont
\setlength{\LTleft}{0pt}\setlength{\LTright}{\fill}
\noindent
\begin{tabular}{@{}>{\raggedright\arraybackslash}p{22.0mm}>{\raggedright\arraybackslash}p{22.688mm}>{\raggedright\arraybackslash}p{22.688mm}>{\raggedright\arraybackslash}p{22.688mm}>{\raggedright\arraybackslash}p{22.688mm}@{}}
\toprule
\textbf{Case} & \textbf{ǣfen singular} & \textbf{ǣfen plural} & \textbf{tācen singular} & \textbf{tācen plural} \\
\midrule
Nominative & ǣfen & ǣfenu & tācen & tācnu \\
Genitive & ǣfenes & ǣfena & tācnes & tācna \\
Dative & ǣfene & ǣfenum & tācne & tācnum \\
Accusative & ǣfen & ǣfenu & tācen & tācnu \\
\bottomrule
\end{tabular}
\par\endgroup
% course-panel:T08-P01:end
% mdtable:T08:end

% mdtable:T09:begin
% course-table-style: 2026-10
% course-panel:T09-P01:begin
\par\begingroup
\fontsize{10.5}{13.2}\selectfont
\setlength{\LTleft}{0pt}\setlength{\LTright}{\fill}
\noindent
\begin{tabular}{@{}>{\raggedright\arraybackslash}p{22.0mm}>{\raggedright\arraybackslash}p{22.688mm}>{\raggedright\arraybackslash}p{22.688mm}>{\raggedright\arraybackslash}p{22.688mm}>{\raggedright\arraybackslash}p{22.688mm}@{}}
\toprule
\textbf{Case} & \textbf{wæter singular} & \textbf{wæter plural} & \textbf{nīeten singular} & \textbf{nīeten plural} \\
\midrule
Nominative & wæter & wæteru & nīeten & nīetenu \\
Genitive & wæteres & wætera & nīetenes & nīetena \\
Dative & wætere & wæterum & nīetene & nīetenum \\
Accusative & wæter & wæteru & nīeten & nīetenu \\
\bottomrule
\end{tabular}
\par\endgroup
% course-panel:T09-P01:end
% mdtable:T09:end

\reserve{4\baselineskip}
% mdparagraph:P019:begin
\section*{5. Strong feminine nouns}\BookHeading{nav-ch01-learner-0005}{subsection}{5. Strong feminine nouns}{1}\BookLabel{bk-ch01-learner-71357d7f7267}
% mdparagraph:P019:end

% teaching: CH01-PAR-FEM-LONG CH01-PAR-FEM-ACC-ZERO
% mdparagraph:P020:begin
The strong feminines required here have no nominative singular ending. The ordinary genitive, dative, and accusative singular end in \textbf{-e}. \textbf{Wyrt} and \textbf{tīd}, however, have an endingless accusative singular in the selected pattern. The selected N/\allowbreak{}A plural is \textbf{-a}; \textbf{-e} is also found.[6]
% mdparagraph:P020:end

% mdtable:T10:begin
% course-table-style: 2026-10
% course-panel:T10-P01:begin
\par\begingroup
\fontsize{10.5}{13.2}\selectfont
\setlength{\LTleft}{0pt}\setlength{\LTright}{\fill}
\noindent
\begin{tabular}{@{}>{\raggedright\arraybackslash}p{22.0mm}>{\raggedright\arraybackslash}p{22.688mm}>{\raggedright\arraybackslash}p{22.688mm}>{\raggedright\arraybackslash}p{22.688mm}>{\raggedright\arraybackslash}p{22.688mm}@{}}
\toprule
\textbf{Case} & \textbf{gesceaft singular} & \textbf{gesceaft plural} & \textbf{wyrt singular} & \textbf{wyrt plural} \\
\midrule
Nominative & gesceaft & gesceafta & wyrt & wyrta \\
Genitive & gesceafte & gesceafta & wyrte & wyrta \\
Dative & gesceafte & gesceaftum & wyrte & wyrtum \\
Accusative & gesceafte & gesceafta & wyrt & wyrta \\
\bottomrule
\end{tabular}
\par\endgroup
% course-panel:T10-P01:end
% mdtable:T10:end

% mdparagraph:P021:begin
\textbf{tīd}\par\nobreak
% mdparagraph:P021:end

% mdtable:T11:begin
% course-table-style: 2026-10
% course-panel:T11-P01:begin
\par\begingroup
\fontsize{10.5}{13.2}\selectfont
\setlength{\LTleft}{0pt}\setlength{\LTright}{\fill}
\noindent
\begin{tabular}{@{}>{\raggedright\arraybackslash}p{39.459mm}>{\raggedright\arraybackslash}p{39.459mm}>{\raggedright\arraybackslash}p{39.459mm}@{}}
\toprule
\textbf{Case} & \textbf{Singular} & \textbf{Plural} \\
\midrule
Nominative & tīd & tīda \\
Genitive & tīde & tīda \\
Dative & tīde & tīdum \\
Accusative & tīd & tīda \\
\bottomrule
\end{tabular}
\par\endgroup
% course-panel:T11-P01:end
% mdtable:T11:end

% teaching: CH01-PAR-FEM-NES CH01-PAR-FEM-UNG
% mdparagraph:P022:begin
Nouns ending in \textbf{-nes} have \textbf{-ness-} before endings: retain the \textbf{e} and double the \textbf{s}. The same pattern governs \textbf{drȳgnes, gelīcnes, fæstnes, nēowolnes, brādnes}. Nouns in \textbf{-ung}, and the assigned \textbf{līehting}, take the endings shown for \textbf{gegaderung}; \textbf{frætwung} follows it. Cook also records singular oblique \textbf{-a} in this group, especially in the dative. The present tables select \textbf{-e}.[6]
% mdparagraph:P022:end

% mdtable:T12:begin
% course-table-style: 2026-10
% course-panel:T12-P01:begin
\par\begingroup
\fontsize{10.5}{13.2}\selectfont
\setlength{\LTleft}{0pt}\setlength{\LTright}{\fill}
\noindent
\begin{tabular}{@{}>{\raggedright\arraybackslash}p{19.0mm}>{\raggedright\arraybackslash}p{23.0mm}>{\raggedright\arraybackslash}p{23.0mm}>{\raggedright\arraybackslash}p{23.0mm}>{\raggedright\arraybackslash}p{24.752mm}@{}}
\toprule
\textbf{Case} & \textbf{andlīcnes singular} & \textbf{andlīcnes plural} & \textbf{gegaderung singular} & \textbf{gegaderung plural} \\
\midrule
Nominative & andlīcnes & andlīcnessa & gegaderung & gegaderunga \\
Genitive & andlīcnesse & andlīcnessa & gegaderunge & gegaderunga \\
Dative & andlīcnesse & andlīcnessum & gegaderunge & gegaderungum \\
Accusative & andlīcnesse & andlīcnessa & gegaderunge & gegaderunga \\
\bottomrule
\end{tabular}
\par\endgroup
% course-panel:T12-P01:end
% mdtable:T12:end

\reserve{4\baselineskip}
% mdparagraph:P023:begin
\section*{6. Weak nouns}\BookHeading{nav-ch01-learner-0006}{subsection}{6. Weak nouns}{1}\BookLabel{bk-ch01-learner-9c5fd7a11f03}
% mdparagraph:P023:end

% teaching: CH01-PAR-WEAK
% mdparagraph:P024:begin
Weak masculine \textbf{steorra} and weak feminine \textbf{eorðe} differ in the nominative singular. Both have \textbf{-an} in the other singular cases and the N/\allowbreak{}A plural, \textbf{-ena} in the genitive plural, and \textbf{-um} in the dative plural. Replace nominative \textbf{-a/\allowbreak{}-e} with these endings.[7]
% mdparagraph:P024:end

% mdtable:T13:begin
% course-table-style: 2026-10
% course-panel:T13-P01:begin
\par\begingroup
\fontsize{10.5}{13.2}\selectfont
\setlength{\LTleft}{0pt}\setlength{\LTright}{\fill}
\noindent
\begin{tabular}{@{}>{\raggedright\arraybackslash}p{22.0mm}>{\raggedright\arraybackslash}p{22.688mm}>{\raggedright\arraybackslash}p{22.688mm}>{\raggedright\arraybackslash}p{22.688mm}>{\raggedright\arraybackslash}p{22.688mm}@{}}
\toprule
\textbf{Case} & \textbf{steorra singular} & \textbf{steorra plural} & \textbf{eorðe singular} & \textbf{eorðe plural} \\
\midrule
Nominative & steorra & steorran & eorðe & eorðan \\
Genitive & steorran & steorrena & eorðan & eorðena \\
Dative & steorran & steorrum & eorðan & eorðum \\
Accusative & steorran & steorran & eorðan & eorðan \\
\bottomrule
\end{tabular}
\par\endgroup
% course-panel:T13-P01:end
% mdtable:T13:end

% mdparagraph:P025:begin
The plural of \textbf{eorðe} is displayed for recognizing its pattern; the exercises use singular “earth.” An \textbf{-an} form alone cannot identify either case or number.
% mdparagraph:P025:end

\reserve{4\baselineskip}
% mdparagraph:P026:begin
\section*{7. Individual paradigms}\BookHeading{nav-ch01-learner-0007}{subsection}{7. Individual paradigms}{1}\BookLabel{bk-ch01-learner-3d6b15e7b6bf}
% mdparagraph:P026:end

% teaching: CH01-PAR-EXCEPTION-MANN CH01-PAR-EXCEPTION-NIHT
% mdparagraph:P027:begin
\textbf{Mann} has \textbf{męnn} in the dative singular and N/\allowbreak{}A plural. In the selected \textbf{niht} paradigm, every singular case and the N/\allowbreak{}A plural are \textbf{niht}; the genitive and dative plurals are \textbf{nihta, nihtum}. Genitive/\allowbreak{}dative singular \textbf{nihte} is a valid alternative, but is not the form selected for recall here.[8]
% mdparagraph:P027:end

% mdtable:T14:begin
% course-table-style: 2026-10
% course-panel:T14-P01:begin
\par\begingroup
\fontsize{10.5}{13.2}\selectfont
\setlength{\LTleft}{0pt}\setlength{\LTright}{\fill}
\noindent
\begin{tabular}{@{}>{\raggedright\arraybackslash}p{22.0mm}>{\raggedright\arraybackslash}p{22.688mm}>{\raggedright\arraybackslash}p{22.688mm}>{\raggedright\arraybackslash}p{22.688mm}>{\raggedright\arraybackslash}p{22.688mm}@{}}
\toprule
\textbf{Case} & \textbf{mann singular} & \textbf{mann plural} & \textbf{niht singular} & \textbf{niht plural} \\
\midrule
Nominative & mann & męnn & niht & niht \\
Genitive & mannes & manna & niht & nihta \\
Dative & męnn & mannum & niht & nihtum \\
Accusative & mann & męnn & niht & niht \\
\bottomrule
\end{tabular}
\par\endgroup
% course-panel:T14-P01:end
% mdtable:T14:end

% teaching: CH01-PAR-EXCEPTION-SAE CH01-REF-SAE
% mdparagraph:P028:begin
\textbf{Sǣ} is masculine in the selected vocabulary. Its nominative plural is \textbf{sǣs}, not a form in \textbf{-as}. The starred genitive plural below is a reconstruction: consult it only as a reference, and do not use it in these exercises.[8]
% mdparagraph:P028:end

% mdparagraph:P029:begin
\textbf{sǣ}\par\nobreak
% mdparagraph:P029:end

% mdtable:T15:begin
% course-table-style: 2026-10
% course-panel:T15-P01:begin
\par\begingroup
\fontsize{10.5}{13.2}\selectfont
\setlength{\LTleft}{0pt}\setlength{\LTright}{\fill}
\noindent
\begin{tabular}{@{}>{\raggedright\arraybackslash}p{39.459mm}>{\raggedright\arraybackslash}p{39.459mm}>{\raggedright\arraybackslash}p{39.459mm}@{}}
\toprule
\textbf{Case} & \textbf{Singular} & \textbf{Plural} \\
\midrule
Nominative & sǣ & sǣs \\
Genitive & sǣs & *sǣwa — reconstructed; not tested \\
Dative & sǣ & sǣum \\
Accusative & sǣ & sǣs \\
\bottomrule
\end{tabular}
\par\endgroup
% course-panel:T15-P01:end
% mdtable:T15:end

% teaching: CH01-PAR-EXCEPTION-THIESTRU
% mdparagraph:P030:begin
\textbf{Ðīestru}, “darkness,” is plural in Old English. Its complete assigned paradigm is plural only.
% mdparagraph:P030:end

% mdparagraph:P031:begin
\textbf{ðīestru}\par\nobreak
% mdparagraph:P031:end

% mdtable:T16:begin
% course-table-style: 2026-10
% course-panel:T16-P01:begin
\par\begingroup
\fontsize{10.5}{13.2}\selectfont
\setlength{\LTleft}{0pt}\setlength{\LTright}{\fill}
\noindent
\begin{tabular}{@{}>{\raggedright\arraybackslash}p{60.594mm}>{\raggedright\arraybackslash}p{60.594mm}@{}}
\toprule
\textbf{Case} & \textbf{Plural} \\
\midrule
Nominative & ðīestru \\
Genitive & ðīestra \\
Dative & ðīestrum \\
Accusative & ðīestru \\
\bottomrule
\end{tabular}
\par\endgroup
% course-panel:T16-P01:end
% mdtable:T16:end

\reserve{4\baselineskip}
% mdparagraph:P032:begin
\section*{8. Subject pronouns and the indicative copula}\BookHeading{nav-ch01-learner-0008}{subsection}{8. Subject pronouns and the indicative copula}{1}
% mdparagraph:P032:end

% teaching: CH01-GRAM-SUBJECT-PRONOUNS CH01-PAR-SUBJECTS CH01-VOC-SUBJECT-ic CH01-VOC-SUBJECT-thu CH01-VOC-SUBJECT-he
% mdparagraph:P033:begin
The first person is the speaker; the second is the person addressed; the third is the person or thing spoken about. Learn these complete \textbf{nominative subject forms} with the verb “be.” The other cases of the personal pronouns follow in Unit 2.[16]
% mdparagraph:P033:end

% mdtable:T17:begin
% course-table-style: 2026-10
% course-panel:T17-P01:begin
\par\begingroup
\fontsize{10.5}{13.2}\selectfont
\setlength{\LTleft}{0pt}\setlength{\LTright}{\fill}
\noindent
\begin{tabular}{@{}>{\raggedright\arraybackslash}p{21.929mm}>{\raggedright\arraybackslash}p{37.279mm}>{\raggedright\arraybackslash}p{25.218mm}>{\raggedright\arraybackslash}p{31.139mm}@{}}
\toprule
\textbf{Person}\BookPoint{bk-ch01-learner-f4792afbd8d6} & \textbf{Singular} & \textbf{Dual: exactly two} & \textbf{Plural} \\
\midrule
First & ic — I & wit — we two & wē — we \\
Second & ðū — you & git — you two & gē — you \\
Third & hē — he; hit — it; hēo — she & — & hīe — they, all genders \\
\bottomrule
\end{tabular}
\par\endgroup
% course-panel:T17-P01:end
% mdtable:T17:end

% teaching: CH01-GRAM-SUBJECT-REFERENCE
% mdparagraph:P034:begin
In\BookPoint{bk-ch01-learner-8f8ad9d5a04c} the third-person singular, \textbf{hē, hit, hēo} normally refer to masculine, neuter, and feminine nouns respectively. Use grammatical gender when referring to a thing: a feminine noun need not name a female. All plural nouns take \textbf{hīe}. The context identifies which noun the pronoun refers to; gender and number help to check that reference.
% mdparagraph:P034:end

% teaching: CH01-GRAM-COPULA
% mdparagraph:P035:begin
These are the complete present \textbf{eom} and preterite \textbf{wæs} series of the indicative “be.” The preterite expresses a past state. Other present forms, negative forms, moods, and nonfinite constructions will be treated separately.[9]
% mdparagraph:P035:end

% mdtable:T18:begin
% course-table-style: 2026-10
% course-panel:T18-P01:begin
\par\begingroup
\fontsize{10.5}{13.2}\selectfont
\setlength{\LTleft}{0pt}\setlength{\LTright}{\fill}
\noindent
\begin{tabular}{@{}>{\raggedright\arraybackslash}p{39.459mm}>{\raggedright\arraybackslash}p{39.459mm}>{\raggedright\arraybackslash}p{39.459mm}@{}}
\toprule
\textbf{Person and number}\BookPoint{bk-ch01-learner-814e0a394df4} & \textbf{Present} & \textbf{Preterite} \\
\midrule
1 sg.: I & eom & wæs \\
2 sg.: you & eart & wǣre \\
3 sg.: he, she, it & is & wæs \\
1–3 pl.: we, you, they & sindon & wǣron \\
\bottomrule
\end{tabular}
\par\endgroup
% course-panel:T18-P01:end
% mdtable:T18:end

% teaching: CH01-GRAM-DUAL-SUBJECTS
% mdparagraph:P036:begin
\textbf{1–3 pl.}\BookPoint{bk-ch01-learner-f73641fc4bb6} gives the shared form for \textbf{wē, gē, hīe}. The dual subjects \textbf{wit} and \textbf{git} also take these ordinary plural verb forms: there is no separate dual form of the copula. Use a dual when specifically saying “we two” or “you two”; the ordinary plurals do not exclude a group of two. Cook also gives \textbf{sind, sint} beside the selected \textbf{sindon}.[16]
% mdparagraph:P036:end

\reserve{4\baselineskip}
% mdparagraph:P037:begin
\section*{9. Basic constructions}\BookHeading{nav-ch01-learner-0009}{subsection}{9. Basic constructions}{1}\BookLabel{bk-ch01-learner-086b075156aa}
% mdparagraph:P037:end

% teaching: CH01-GRAM-SUBJECT-PREDICATE CH01-GRAM-AGREEMENT
% mdparagraph:P038:begin
The subject is nominative. A noun that identifies it after “be” is also nominative; English \textbf{a/\allowbreak{}an} need not receive a separate word in such a predicate. The copula agrees with the subject in person and number.[10]
% mdparagraph:P038:end

% example: CH01-EXAMPLE07
% mdparagraph:P039:begin
\textbf{Ðā hwalas sindon dēor.} — The whales are animals.
% mdparagraph:P039:end

% mdparagraph:P040:begin
Both subject \textbf{hwalas} and predicate \textbf{dēor} are nominative plural; the neuter plural \textbf{dēor} has no added ending.
% mdparagraph:P040:end

% example: CH01-EXAMPLE09
% mdparagraph:P041:begin
\textbf{Ic eom mann.} — I am a human being.
% mdparagraph:P041:end

% example: CH01-EXAMPLE10
% mdparagraph:P042:begin
\textbf{Wit sindon męnn.} — We two are human beings.
% mdparagraph:P042:end

% teaching: CH01-GRAM-COORDINATION
% mdparagraph:P043:begin
\textbf{And} joins words or phrases. When definite noun phrases are joined, each article agrees with its own noun. Gender belongs to that noun; it is not taken from the other noun in the pair.
% mdparagraph:P043:end

% teaching: CH01-GRAM-ON
% mdparagraph:P044:begin
For the location of an entity, use \textbf{on + dative}, “in/\allowbreak{}on.”
% mdparagraph:P044:end

% example: CH01-EXAMPLE01
% mdparagraph:P045:begin
\textbf{Se fisc is on ðǣm wætere.} — The fish is in the water.
% mdparagraph:P045:end

% example: CH01-EXAMPLE02
% mdparagraph:P046:begin
\textbf{Ðā męnn wǣron on ðǣm hofum.} — The men were in the dwellings.
% mdparagraph:P046:end

% teaching: CH01-GRAM-GENITIVE
% mdparagraph:P047:begin
A dependent genitive expresses possession or a relation such as the source, contents, or kind of something. In the phrase pattern used for Cook's Exercise 3, place the genitive first and omit the article of the \textbf{head noun}, while retaining the article belonging to the dependent. The head remains in the case required by its own role. A proper name such as \textbf{God} takes no article here.[11]
% mdparagraph:P047:end

% example: CH01-EXAMPLE03
% mdparagraph:P048:begin
\textbf{ðæs mannes fugol} — the man's bird
% mdparagraph:P048:end

% mdparagraph:P049:begin
Here \textbf{ðæs mannes} is the dependent genitive; \textbf{fugol} is the head, nominative in the isolated naming phrase. The construction does not mean that every English “the” disappears.
% mdparagraph:P049:end

% teaching: CH01-GRAM-NESTED-GENITIVE
% mdparagraph:P050:begin
The\BookPoint{bk-ch01-learner-0b109d29614d} same dependent may stand inside a prepositional phrase. Its genitive does not change the case required for the head.
% mdparagraph:P050:end

% example: CH01-EXAMPLE08
% mdparagraph:P051:begin
\textbf{on ðæs fugles trēowe} — in the bird's tree
% mdparagraph:P051:end

% mdparagraph:P052:begin
\textbf{Trēowe} is dative after locative \textbf{on}; \textbf{ðæs fugles} remains genitive. \textbf{Ðæs} agrees with \textbf{fugles}, not with \textbf{trēowe}.
% mdparagraph:P052:end

% teaching: CH01-GRAM-OBJECT
% mdparagraph:P053:begin
The\BookPoint{bk-ch01-learner-422453d6be7d} direct object is accusative. In the following example, \textbf{geseah} is supplied: “saw,” third-person singular preterite. Its conjugation is not required here.[12]
% mdparagraph:P053:end

% example: CH01-EXAMPLE04
% mdparagraph:P054:begin
\textbf{Se mann geseah ðone fisc.} — The man saw the fish.
% mdparagraph:P054:end

% mdparagraph:P055:begin
\textbf{Mann} and \textbf{fisc} have identical nominative and accusative singular spellings, but their articles distinguish subject \textbf{se} from object \textbf{ðone}.
% mdparagraph:P055:end

% teaching: CH01-GRAM-RECIPIENT CH01-FRAME-GIVING
% mdparagraph:P056:begin
The\BookPoint{bk-ch01-learner-d64391595e92} recipient of a gift is dative; the thing given is accusative. The supplied \textbf{sealde} means “gave,” third-person singular preterite.[12]
% mdparagraph:P056:end

% example: CH01-EXAMPLE05
% mdparagraph:P057:begin
\textbf{God sealde ðǣm fugle męte.} — God gave the bird food.
% mdparagraph:P057:end

% mdparagraph:P058:begin
\textbf{Ðǣm fugle} is the recipient; \textbf{męte} is the direct object. English may express the recipient with “to,” but this Old English construction uses the dative without a preposition.
% mdparagraph:P058:end

% teaching: CH01-GRAM-FRAM
% mdparagraph:P059:begin
Use\BookPoint{bk-ch01-learner-2b2114170c4e} \textbf{fram + dative} for what something is separated from. The supplied \textbf{gedǣlde} means “separated/\allowbreak{}divided,” third-person singular preterite.[12]
% mdparagraph:P059:end

% example: CH01-EXAMPLE06
% mdparagraph:P060:begin
\textbf{Se mann gedǣlde ðæt sǣd fram ðǣm gærse.} — The man separated the seed from the grass.
% mdparagraph:P060:end

% mdparagraph:P061:begin
The thing separated, \textbf{ðæt sǣd}, is accusative; \textbf{ðǣm gærse} is dative after \textbf{fram}. These are particular constructions. They do not establish one case rule for all spatial prepositions.
% mdparagraph:P061:end

\reserve{4\baselineskip}
% mdparagraph:P062:begin
\section*{10. Vocabulary}\BookHeading{nav-ch01-learner-0010}{subsection}{10. Vocabulary}{1}
% mdparagraph:P062:end

% mdparagraph:P063:begin
Each pattern reference gives a section and its model noun: \textbf{§4: hof}, for example, means the \textbf{hof} pattern in §4. References to \textbf{sǣd}, \textbf{trēo}, \textbf{tācen}, or another named neuter select that particular pattern. They are not interchangeable.
% mdparagraph:P063:end

% mdparagraph:P064:begin
Learn each noun with its gender and diagnostic forms. \textbf{G sg} identifies the singular stem; \textbf{N pl} identifies the plural pattern. The full tables and rules above supply the other cases. \textbf{S} marks a noun used only in the singular in composition here: its plural morphology is for recognition, with no additional plural sense assigned. \textbf{P} marks plural-only \textbf{ðīestru}. The singular divine name \textbf{God} has no plural assigned. Abbreviations: \textbf{m.}, masculine; \textbf{f.}, feminine; \textbf{n.}, neuter.[13]
% mdparagraph:P064:end

% mdtable:T19:begin
% course-table-style: 2026-10
% teaching: CH01-VOC-aefen
% teaching: CH01-VOC-morgen
% teaching: CH01-VOC-daeg
% teaching: CH01-VOC-niht
% teaching: CH01-VOC-eorthe
% teaching: CH01-VOC-heofon
% teaching: CH01-VOC-fisc
% teaching: CH01-VOC-fugol
% teaching: CH01-VOC-gaers
% teaching: CH01-VOC-saed
% teaching: CH01-VOC-treo
% teaching: CH01-VOC-wyrt
% teaching: CH01-VOC-drygnes
% teaching: CH01-VOC-waeter
% teaching: CH01-VOC-andlicnes
% teaching: CH01-VOC-gelicnes
% teaching: CH01-VOC-leoht
% teaching: CH01-VOC-thiestru
% teaching: CH01-VOC-tacen
% teaching: CH01-VOC-tid
% teaching: CH01-VOC-gast
% teaching: CH01-VOC-lif
% teaching: CH01-VOC-gesceaft
% teaching: CH01-VOC-mete
% teaching: CH01-VOC-gear
% teaching: CH01-VOC-mann
% teaching: CH01-VOC-deor
% teaching: CH01-VOC-waestm
% teaching: CH01-VOC-faestnes
% teaching: CH01-VOC-neowolnes
% teaching: CH01-VOC-sae
% teaching: CH01-VOC-hwael
% teaching: CH01-VOC-nieten
% teaching: CH01-VOC-anginn
% teaching: CH01-VOC-god
% teaching: CH01-VOC-bradnes
% teaching: CH01-VOC-gegaderung
% teaching: CH01-VOC-lyft
% teaching: CH01-VOC-fraetwung
% teaching: CH01-VOC-weorc
% teaching: CH01-VOC-cynn
% teaching: CH01-VOC-steorra
% teaching: CH01-VOC-liehting
% teaching: CH01-VOC-hof
% course-panel:T19-P01:begin
\par\begingroup
\fontsize{10.5}{13.2}\selectfont
\setlength{\LTleft}{0pt}\setlength{\LTright}{\fill}
\begin{longtable}{@{}>{\raggedright\arraybackslash}p{24.0mm}>{\raggedright\arraybackslash}p{39.0mm}>{\raggedright\arraybackslash}p{28.0mm}>{\raggedright\arraybackslash}p{24.564mm}@{}}
\toprule
\textbf{Noun; gender} & \textbf{Meaning} & \textbf{Diagnostic forms} & \textbf{Pattern; limit} \\
\midrule
\endfirsthead
\toprule
\textbf{Noun; gender} & \textbf{Meaning} & \textbf{Diagnostic forms} & \textbf{Pattern; limit} \\
\midrule
\endhead
\midrule
\multicolumn{4}{r}{\footnotesize Continued on the following page.}\\
\endfoot
\bottomrule
\endlastfoot
 ǣfen (n.) & evening & G sg ǣfenes N pl ǣfenu & §4: ǣfen \\
 morgen (m.) & morning & G sg morgnes N pl morgnas & §3: morgen \\
 dæg (m.) & day & G sg dæges N pl dagas & §3: dæg \\
 niht (f.) & night & G sg niht N pl niht & §7: niht \\
 eorðe (f.) & earth & G sg eorðan N pl eorðan & §6: eorðe; S \\
 heofon (m.) & heaven; sky & G sg heofones N pl heofonas & §3: fisc \\
 fisc (m.) & fish & G sg fisces N pl fiscas & §3: fisc \\
 fugol (m.) & bird & G sg fugles N pl fuglas & §3: fugol \\
 gærs (n.) & grass & G sg gærses N pl gærs & §4: sǣd; S \\
 sǣd (n.) & seed & G sg sǣdes N pl sǣd & §4: sǣd \\
 trēo (n.) & tree & G sg trēowes N pl trēowu & §4: trēo \\
 wyrt (f.) & herb; plant & G sg wyrte N pl wyrta & §5: wyrt \\
 drȳgnes (f.) & dry land & G sg drȳgnesse N pl drȳgnessa & §5: andlīcnes; S \\
 wæter (n.) & water & G sg wæteres N pl wæteru & §4: wæter \\
 andlīcnes (f.) & image & G sg andlīcnesse N pl andlīcnessa & §5: andlīcnes \\
 gelīcnes (f.) & likeness & G sg gelīcnesse N pl gelīcnessa & §5: andlīcnes \\
 lēoht (n.) & light & G sg lēohtes N pl lēoht & §4: sǣd; S \\
 ðīestru (P) & darkness & G sg — N pl ðīestru & §7: ðīestru; P; G pl ðīestra \\
 tācen (n.) & sign & G sg tācnes N pl tācnu & §4: tācen \\
 tīd (f.) & season; time & G sg tīde N pl tīda & §5: tīd \\
 gāst (m.) & spirit & G sg gāstes N pl gāstas & §3: fisc \\
 līf (n.) & life & G sg līfes N pl līf & §4: sǣd; S \\
 gesceaft (f.) & creature; created being or thing & G sg gesceafte N pl gesceafta & §5: gesceaft \\
 męte (m.) & food & G sg mętes N pl mętas & §3: męte; S \\
 gēar (n.) & year & G sg gēares N pl gēar & §4: sǣd \\
 mann (m.) & man; human being & G sg mannes N pl męnn & §7: mann \\
 dēor (n.) & animal; beast & G sg dēores N pl dēor & §4: sǣd \\
 wæstm (m.) & fruit; produce & G sg wæstmes N pl wæstmas & §3: fisc \\
 fæstnes (f.) & firmament & G sg fæstnesse N pl fæstnessa & §5: andlīcnes; S \\
 nēowolnes (f.) & abyss; the deep & G sg nēowolnesse N pl nēowolnessa & §5: andlīcnes; S \\
 sǣ (m.) & sea & G sg sǣs N pl sǣs & §7: sǣ \\
 hwæl (m.) & whale & G sg hwæles N pl hwalas & §3: hwæl \\
 nīeten (n.) & animal; beast; in the plural, cattle & G sg nīetenes N pl nīetenu & §4: nīeten \\
 anginn (n.) & beginning & G sg anginnes N pl anginn & §4: sǣd \\
 God (m.) & God (proper divine name) & G sg Godes N pl — & §3: fisc; divine name, singular \\
 brādnes (f.) & surface; face (of the deep) & G sg brādnesse N pl brādnessa & §5: andlīcnes; S \\
 gegaderung (f.) & gathering; collection; congregation (of waters) & G sg gegaderunge N pl gegaderunga & §5: gegaderung \\
 lyft (n.) & air & G sg lyftes N pl lyft & §4: sǣd; S \\
 frætwung (f.) & array & G sg frætwunge N pl frætwunga & §5: gegaderung; S \\
 weorc (n.) & work & G sg weorces N pl weorc & §4: sǣd \\
 cynn (n.) & kind & G sg cynnes N pl cynn & §4: sǣd \\
 steorra (m.) & star & G sg steorran N pl steorran & §6: steorra \\
 līehting (f.) & illumination & G sg līehtinge N pl līehtinga & §5: gegaderung; S \\
 hof (n.) & dwelling & G sg hofes N pl hofu & §4: hof \\
\end{longtable}
\par\endgroup
% course-panel:T19-P01:end
% mdtable:T19:end

% mdparagraph:P065:begin
Pronunciation: \textbf{cynn} begins with \textbf{k} as in \emph{kit}; \textbf{gærs} begins with stopped \textbf{g} as in \emph{go}, despite its following \textbf{æ}. Learn these supplied sounds with the words.[17]
% mdparagraph:P065:end

% mdparagraph:P066:begin
The other words required for composition are \textbf{and}, “and”; \textbf{on + dative}, “in/\allowbreak{}on” for location; \textbf{fram + dative}, “from” in separation; the article; and the subject pronouns and indicative forms in §8. Verbs printed in frames and the additional help for the reading are supplied, not additional vocabulary to recall independently.
% mdparagraph:P066:end

\reserve{4\baselineskip}
% mdparagraph:P067:begin
\section*{11. Forms and agreement}\BookHeading{nav-ch01-learner-0011}{subsection}{11. Forms and agreement}{1}
% mdparagraph:P067:end

% mdparagraph:P068:begin
The tasks in §§11–12 and §§14–17 are new practice. Use the selected patterns; the key distinguishes valid source variants from the particular forms requested. In the noun tasks, translate only the marked words, not the rest of the English context. Choose the case from the word's function.
% mdparagraph:P068:end

% item: CH01-F01
% mdparagraph:P069:begin
\itempara{1}{Write the bold noun phrase in Old English and label its case: (a) \textbf{The fish} is food; (b) the body of \textbf{the fish}; (c) God gave food to \textbf{the fish}; (d) the bird saw \textbf{the fish}. First make each fish singular, then rewrite all four phrases for several fish. Include the article.}
% mdparagraph:P069:end

% item: CH01-F02
% mdparagraph:P070:begin
\itempara{2}{Write only the bold noun, without an article: (a) several \textbf{dwellings} are on the earth; (b) the man saw several \textbf{dwellings}; (c) the \textbf{seeds} are food; (d) the bird saw the \textbf{seeds}. Identify each case and number. Explain why the two nouns differ in their plural endings although their roles are parallel.}
% mdparagraph:P070:end

% item: CH01-F03
% mdparagraph:P071:begin
\itempara{3}{Give only the bold noun, without an article, and explain its case: (a) the length of several \textbf{days}; (b) the food of several \textbf{whales}; (c) God gave food to one \textbf{man}; (d) the darkness of several \textbf{nights}. Preserve each noun's own stem changes.}
% mdparagraph:P071:end

% item: CH01-F04
% mdparagraph:P072:begin
\itempara{4}{Give only the bold noun, without an article: (a) the fruit of one \textbf{tree}; (b) the bird is in one \textbf{tree}; (c) several \textbf{trees} are in the dwelling; (d) the meaning of one \textbf{sign}; (e) several \textbf{signs} are in the sky; (f) the sound of the \textbf{water}; (g) the \textbf{waters} are in the sea; (h) several \textbf{evenings} have passed; (i) the \textbf{cattle} are in the dwelling. Label the forms. Explain the \textbf{w} in (c), and contrast the loss or retention of the middle stem vowel in (e) and (g).}
% mdparagraph:P072:end

% item: CH01-F05
% mdparagraph:P073:begin
\itempara{5}{Give only the bold noun, without an article: (a) \textbf{the star} is in the sky; (b) the light of \textbf{the star}; (c) the man is looking at \textbf{the star}, as a direct object; (d) the bird is under \textbf{the star}, using a supplied dative frame. First use one star, then several stars. Next give singular \textbf{earth} as (e) subject, (f) possessor in “the surface of the earth,” (g) location after \textbf{on}, and (h) direct object in “God created the earth.” Explain why an isolated \textbf{-an} form cannot settle the analysis. Only the noun forms are requested; the verbs and the English prepositions in (c–d) are not to be translated.}
% mdparagraph:P073:end

% item: CH01-F06
% mdparagraph:P074:begin
\itempara{6}{Give only the bold noun, without an article: (a) the man saw one \textbf{herb}; (b) the seed of one \textbf{herb}; (c) the man marked one \textbf{season} as the direct object; (d) the beginning of one \textbf{season}; (e) the man saw one \textbf{creature}. Explain why the direct objects do not all take the same feminine ending.}
% mdparagraph:P074:end

% item: CH01-F07
% mdparagraph:P075:begin
\itempara{7}{Give only the bold noun, without an article: (a) in the \textbf{firmament}, using locative \textbf{on}; (b) the form of several \textbf{images}; (c) several \textbf{likenesses} are in the dwelling; (d) in several \textbf{gatherings}, using locative \textbf{on}. Identify the case and number, and explain the stem before each ending.}
% mdparagraph:P075:end

% item: CH01-F08
% mdparagraph:P076:begin
\itempara{8}{Reference check: locate \textbf{ðȳ} and \textbf{ðon} in the article table. Identify their case and singular genders. This checks recognition of the reference forms; their use in composition is not required.}
% mdparagraph:P076:end

\reserve{4\baselineskip}
% mdparagraph:P077:begin
\subsection*{The copula}\BookHeading{nav-ch01-learner-0012}{subsubsection}{The copula}{0}
% mdparagraph:P077:end

% mdparagraph:P078:begin
Write complete Old English sentences using a subject pronoun. Preserve the intended person, number, time, and location. Use a dual wherever the prompt explicitly says “two.”
% mdparagraph:P078:end

% item: CH01-COP01
% mdparagraph:P079:begin
\itempara{1}{You are speaking about yourself: “I was in the tree.” Then change only the time: “I am in the tree.”}
% mdparagraph:P079:end

% item: CH01-COP02
% mdparagraph:P080:begin
\itempara{2}{Address one person: “You are in the water.” Then place the same state in the past.}
% mdparagraph:P080:end

% item: CH01-COP03
% mdparagraph:P081:begin
\itempara{3}{Replace the subject of each sentence with its appropriate pronoun and translate: (a) “The spirit was in the air”; (b) “The sign is in the dwelling”; (c) “The image is on the earth.” Which choice depends on grammatical gender rather than sex?}
% mdparagraph:P081:end

% item: CH01-COP04
% mdparagraph:P082:begin
\itempara{4}{You speak for yourself and one companion: “We two are in the sea.” Put this in the past, then restate the past sentence for a group of three that includes you. Explain which pronoun changes and why the verb's number does not.}
% mdparagraph:P082:end

% item: CH01-COP05
% mdparagraph:P083:begin
\itempara{5}{Address two companions: “You two were in the dwelling.” Put this in the present, then restate the present sentence to address a group of three. Explain the choices of pronoun and verb.}
% mdparagraph:P083:end

% item: CH01-COP06
% mdparagraph:P084:begin
\itempara{6}{The whales have already been mentioned. Say “They were in the water,” then put this in the present. Would the subject pronoun change if it referred instead to several feminine \textbf{gesceafta}, “creatures”? Explain.}
% mdparagraph:P084:end

\reserve{4\baselineskip}
% mdparagraph:P085:begin
\section*{12. Case in sentences}\BookHeading{nav-ch01-learner-0013}{subsection}{12. Case in sentences}{1}
% mdparagraph:P085:end

% mdparagraph:P086:begin
Complete each frame with the noun and article phrases needed for the English sentence. The printed verb is supplied in its exact form; choose the noun cases, numbers, and articles. The prepositions in the frames use the constructions of §9.
% mdparagraph:P086:end

\reserve{7\baselineskip}
% item: CH01-CASE01
% mdparagraph:P087:begin
\itempara{1}{The man saw the bird.}
% mdparagraph:P087:end

% mdparagraph:P088:begin
\textbf{\_\_\_\_\_ geseah \_\_\_\_\_.}\par\nobreak
% mdparagraph:P088:end

% mdparagraph:P089:begin
Supplied: \textbf{geseah}, “saw,” third-person singular preterite.
% mdparagraph:P089:end

\reserve{7\baselineskip}
% item: CH01-CASE02
% mdparagraph:P090:begin
\itempara{2}{The bird saw the man.}
% mdparagraph:P090:end

% mdparagraph:P091:begin
\textbf{\_\_\_\_\_ geseah \_\_\_\_\_.}\par\nobreak
% mdparagraph:P091:end

% mdparagraph:P092:begin
Supplied: \textbf{geseah}, “saw,” third-person singular preterite.
% mdparagraph:P092:end

\reserve{7\baselineskip}
% item: CH01-CASE03
% mdparagraph:P093:begin
\itempara{3}{God created the heavens and the earth. Use masculine heofon for 'heavens'.}
% mdparagraph:P093:end

% mdparagraph:P094:begin
\textbf{\_\_\_\_\_ gescēop \_\_\_\_\_ and \_\_\_\_\_.}\par\nobreak
% mdparagraph:P094:end

% mdparagraph:P095:begin
Supplied: \textbf{gescēop}, “created,” third-person singular preterite.
% mdparagraph:P095:end

\reserve{7\baselineskip}
% item: CH01-CASE04
% mdparagraph:P096:begin
\itempara{4}{The man saw the stars.}
% mdparagraph:P096:end

% mdparagraph:P097:begin
\textbf{\_\_\_\_\_ geseah \_\_\_\_\_.}\par\nobreak
% mdparagraph:P097:end

% mdparagraph:P098:begin
Supplied: \textbf{geseah}, “saw,” third-person singular preterite.
% mdparagraph:P098:end

\reserve{7\baselineskip}
% item: CH01-CASE05
% mdparagraph:P099:begin
\itempara{5}{The man saw the fish in the sea.}
% mdparagraph:P099:end

% mdparagraph:P100:begin
\textbf{\_\_\_\_\_ geseah \_\_\_\_\_ on \_\_\_\_\_.}\par\nobreak
% mdparagraph:P100:end

% mdparagraph:P101:begin
Supplied: \textbf{geseah}, “saw,” third-person singular preterite. The fish is in the sea.
% mdparagraph:P101:end

\reserve{7\baselineskip}
% item: CH01-CASE06
% mdparagraph:P102:begin
\itempara{6}{God separated the light from the darkness.}
% mdparagraph:P102:end

% mdparagraph:P103:begin
\textbf{\_\_\_\_\_ gedǣlde \_\_\_\_\_ fram \_\_\_\_\_.}\par\nobreak
% mdparagraph:P103:end

% mdparagraph:P104:begin
Supplied: \textbf{gedǣlde}, “separated/\allowbreak{}divided,” third-person singular preterite.
% mdparagraph:P104:end

\reserve{7\baselineskip}
% item: CH01-CASE07
% mdparagraph:P105:begin
\itempara{7}{The bird saw the seed in the man's dwelling.}
% mdparagraph:P105:end

% mdparagraph:P106:begin
\textbf{\_\_\_\_\_ geseah \_\_\_\_\_ on \_\_\_\_\_.}\par\nobreak
% mdparagraph:P106:end

% mdparagraph:P107:begin
Supplied: \textbf{geseah}, “saw,” third-person singular preterite. The seed is in the dwelling.
% mdparagraph:P107:end

\reserve{7\baselineskip}
% item: CH01-CASE08
% mdparagraph:P108:begin
\itempara{8}{The bird saw the man in the dwelling.}
% mdparagraph:P108:end

% mdparagraph:P109:begin
\textbf{\_\_\_\_\_ geseah \_\_\_\_\_ on \_\_\_\_\_.}\par\nobreak
% mdparagraph:P109:end

% mdparagraph:P110:begin
Supplied: \textbf{geseah}, “saw,” third-person singular preterite. The man is in the dwelling.
% mdparagraph:P110:end

\reserve{9\baselineskip}
% item: CH01-CASE09
% mdparagraph:P111:begin
\itempara{9}{Complete the printed sealde frame to express: God gave food to the man. Identify the recipient and the thing given.}
% mdparagraph:P111:end

% mdparagraph:P112:begin
\textbf{\_\_\_\_\_ sealde \_\_\_\_\_ \_\_\_\_\_.}\par\nobreak
% mdparagraph:P112:end

% mdparagraph:P113:begin
Supplied: \textbf{sealde}, “gave,” third-person singular preterite. Put the recipient before the thing given in this frame.
% mdparagraph:P113:end

\reserve{9\baselineskip}
% item: CH01-CASE10
% mdparagraph:P114:begin
\itempara{10}{Complete the printed sealde frame to express: The man gave seed to the birds. Identify the recipient and the thing given.}
% mdparagraph:P114:end

% mdparagraph:P115:begin
\textbf{\_\_\_\_\_ sealde \_\_\_\_\_ \_\_\_\_\_.}\par\nobreak
% mdparagraph:P115:end

% mdparagraph:P116:begin
Supplied: \textbf{sealde}, “gave,” third-person singular preterite. Put the recipient before the thing given in this frame.
% mdparagraph:P116:end

\reserve{4\baselineskip}
% mdparagraph:P117:begin
\section*{13. Cook's noun-phrase exercises}\BookHeading{nav-ch01-learner-0014}{subsection}{13. Cook's noun-phrase exercises}{1}
% mdparagraph:P117:end

% mdparagraph:P118:begin
Translate the following original prompts using the vocabulary of this unit. In Exercises 1–2, make the isolated phrases nominative and preserve the requested number. In Exercise 3, keep the head nominative; place its genitive dependent first, omitting the head's own article as explained in §9. The dependent retains its appropriate article. Use \textbf{dēor} for “animal/\allowbreak{}beast” and \textbf{nīeten} for plural “cattle”; \textbf{brādnes} means the surface in “face of the deep.”[14]
% mdparagraph:P118:end

\reserve{4\baselineskip}
% mdparagraph:P119:begin
\subsection*{Group I, Exercise 1}\BookHeading{nav-ch01-learner-0015}{subsubsection}{Group I, Exercise 1}{0}
% mdparagraph:P119:end

% item: EX-I-01-01
% mdparagraph:P120:begin
\itempara{1}{The evening and the morning.}
% mdparagraph:P120:end

% item: EX-I-01-02
% mdparagraph:P121:begin
\itempara{2}{The day and the night.}
% mdparagraph:P121:end

% item: EX-I-01-03
% mdparagraph:P122:begin
\itempara{3}{The earth and the heaven.}
% mdparagraph:P122:end

% item: EX-I-01-04
% mdparagraph:P123:begin
\itempara{4}{The fish and the bird.}
% mdparagraph:P123:end

% item: EX-I-01-05
% mdparagraph:P124:begin
\itempara{5}{The grass and the seed.}
% mdparagraph:P124:end

% item: EX-I-01-06
% mdparagraph:P125:begin
\itempara{6}{The tree and the herb.}
% mdparagraph:P125:end

% item: EX-I-01-07
% mdparagraph:P126:begin
\itempara{7}{The dry land and the water.}
% mdparagraph:P126:end

% item: EX-I-01-08
% mdparagraph:P127:begin
\itempara{8}{The image and the likeness.}
% mdparagraph:P127:end

% item: EX-I-01-09
% mdparagraph:P128:begin
\itempara{9}{The light and the darkness (plur.).}
% mdparagraph:P128:end

% item: EX-I-01-10
% mdparagraph:P129:begin
\itempara{10}{The sign and the season.}
% mdparagraph:P129:end

% item: EX-I-01-11
% mdparagraph:P130:begin
\itempara{11}{The spirit and the life.}
% mdparagraph:P130:end

% item: EX-I-01-12
% mdparagraph:P131:begin
\itempara{12}{The creature and the food.}
% mdparagraph:P131:end

% item: EX-I-01-13
% mdparagraph:P132:begin
\itempara{13}{The year and the day.}
% mdparagraph:P132:end

% item: EX-I-01-14
% mdparagraph:P133:begin
\itempara{14}{The man and the animal.}
% mdparagraph:P133:end

% item: EX-I-01-15
% mdparagraph:P134:begin
\itempara{15}{The tree and the fruit.}
% mdparagraph:P134:end

% item: EX-I-01-16
% mdparagraph:P135:begin
\itempara{16}{The firmament and the abyss.}
% mdparagraph:P135:end

\reserve{4\baselineskip}
% mdparagraph:P136:begin
\subsection*{Group I, Exercise 2}\BookHeading{nav-ch01-learner-0016}{subsubsection}{Group I, Exercise 2}{0}
% mdparagraph:P136:end

% item: EX-I-02-01
% mdparagraph:P137:begin
\itempara{1}{The days and the nights.}
% mdparagraph:P137:end

% item: EX-I-02-02
% mdparagraph:P138:begin
\itempara{2}{The fishes and the birds.}
% mdparagraph:P138:end

% item: EX-I-02-03
% mdparagraph:P139:begin
\itempara{3}{The trees and the herbs.}
% mdparagraph:P139:end

% item: EX-I-02-04
% mdparagraph:P140:begin
\itempara{4}{The images and the likenesses.}
% mdparagraph:P140:end

% item: EX-I-02-05
% mdparagraph:P141:begin
\itempara{5}{The signs and the seasons.}
% mdparagraph:P141:end

% item: EX-I-02-06
% mdparagraph:P142:begin
\itempara{6}{The years and the days.}
% mdparagraph:P142:end

% item: EX-I-02-07
% mdparagraph:P143:begin
\itempara{7}{The men and the animals.}
% mdparagraph:P143:end

% item: EX-I-02-08
% mdparagraph:P144:begin
\itempara{8}{The seas and the heavens.}
% mdparagraph:P144:end

% item: EX-I-02-09
% mdparagraph:P145:begin
\itempara{9}{The seeds and the fruits.}
% mdparagraph:P145:end

% item: EX-I-02-10
% mdparagraph:P146:begin
\itempara{10}{The evenings and the mornings.}
% mdparagraph:P146:end

% item: EX-I-02-11
% mdparagraph:P147:begin
\itempara{11}{The spirits and the men.}
% mdparagraph:P147:end

% item: EX-I-02-12
% mdparagraph:P148:begin
\itempara{12}{The whales and the cattle.}
% mdparagraph:P148:end

\reserve{4\baselineskip}
% mdparagraph:P149:begin
\subsection*{Group I, Exercise 3}\BookHeading{nav-ch01-learner-0017}{subsubsection}{Group I, Exercise 3}{0}
% mdparagraph:P149:end

% item: EX-I-03-01
% mdparagraph:P150:begin
\itempara{1}{The beginning of the day.}
% mdparagraph:P150:end

% item: EX-I-03-02
% mdparagraph:P151:begin
\itempara{2}{The spirit of God.}
% mdparagraph:P151:end

% item: EX-I-03-03
% mdparagraph:P152:begin
\itempara{3}{The light of the life.}
% mdparagraph:P152:end

% item: EX-I-03-04
% mdparagraph:P153:begin
\itempara{4}{The life of the bird.}
% mdparagraph:P153:end

% item: EX-I-03-05
% mdparagraph:P154:begin
\itempara{5}{The food of the fish.}
% mdparagraph:P154:end

% item: EX-I-03-06
% mdparagraph:P155:begin
\itempara{6}{The firmament of the heaven.}
% mdparagraph:P155:end

% item: EX-I-03-07
% mdparagraph:P156:begin
\itempara{7}{The beasts of the earth.}
% mdparagraph:P156:end

% item: EX-I-03-08
% mdparagraph:P157:begin
\itempara{8}{The face of the deep.}
% mdparagraph:P157:end

% item: EX-I-03-09
% mdparagraph:P158:begin
\itempara{9}{The congregations of the waters.}
% mdparagraph:P158:end

% item: EX-I-03-10
% mdparagraph:P159:begin
\itempara{10}{The waters of the sea.}
% mdparagraph:P159:end

% item: EX-I-03-11
% mdparagraph:P160:begin
\itempara{11}{The fishes of the sea.}
% mdparagraph:P160:end

% item: EX-I-03-12
% mdparagraph:P161:begin
\itempara{12}{The birds of the air.}
% mdparagraph:P161:end

% item: EX-I-03-13
% mdparagraph:P162:begin
\itempara{13}{The array of the heavens and of the earth.}
% mdparagraph:P162:end

% item: EX-I-03-14
% mdparagraph:P163:begin
\itempara{14}{The works of God.}
% mdparagraph:P163:end

% item: EX-I-03-15
% mdparagraph:P164:begin
\itempara{15}{The kinds of the animals.}
% mdparagraph:P164:end

% item: EX-I-03-16
% mdparagraph:P165:begin
\itempara{16}{The stars of the heaven.}
% mdparagraph:P165:end

% item: EX-I-03-17
% mdparagraph:P166:begin
\itempara{17}{The seeds of the grass.}
% mdparagraph:P166:end

% item: EX-I-03-18
% mdparagraph:P167:begin
\itempara{18}{The illumination of the nights.}
% mdparagraph:P167:end

\reserve{4\baselineskip}
% mdparagraph:P168:begin
\section*{14. Composition}\BookHeading{nav-ch01-learner-0018}{subsection}{14. Composition}{1}
% mdparagraph:P168:end

% mdparagraph:P169:begin
Use the selected forms and constructions without supplied verb frames.
% mdparagraph:P169:end

% item: CH01-COMP01
% mdparagraph:P170:begin
\itempara{1}{The bird is in the tree.}
% mdparagraph:P170:end

% item: CH01-COMP02
% mdparagraph:P171:begin
\itempara{2}{The birds were in the trees.}
% mdparagraph:P171:end

% item: CH01-COMP03
% mdparagraph:P172:begin
\itempara{3}{The seed is the bird's food.}
% mdparagraph:P172:end

% item: CH01-COMP04
% mdparagraph:P173:begin
\itempara{4}{The fish is an animal.}
% mdparagraph:P173:end

% item: CH01-COMP05
% mdparagraph:P174:begin
\itempara{5}{The signs were in the sky.}
% mdparagraph:P174:end

% item: CH01-COMP06
% mdparagraph:P175:begin
\itempara{6}{The herbs are in the man's dwelling.}
% mdparagraph:P175:end

% item: CH01-COMP07
% mdparagraph:P176:begin
\itempara{7}{The light of the night was in the sky.}
% mdparagraph:P176:end

% item: CH01-COMP08
% mdparagraph:P177:begin
\itempara{8}{The animals' food is the fruit of the trees.}
% mdparagraph:P177:end

\reserve{4\baselineskip}
% mdparagraph:P178:begin
\section*{15. Short readings}\BookHeading{nav-ch01-learner-0019}{subsection}{15. Short readings}{1}
% mdparagraph:P178:end

% mdparagraph:P179:begin
These two connected passages are new composition using the vocabulary and grammar of this unit. Translate from the Old English before answering the questions.
% mdparagraph:P179:end

% item: CH01-MINI01
% mdparagraph:P180:begin
\itempara{1}{Where was the man, and where is he now? Identify the reference of each pronoun. Explain the number of the subject and predicate noun in the fourth sentence, and identify whose dwelling is mentioned at the close.}
% mdparagraph:P180:end

% mdparagraph:P181:begin
Se mann wæs on ðǣm sǣ. Hē is on ðǣm hofe. Ðā fiscas sindon on ðǣm hofe. Hīe sindon ðæs mannes męte. Ðæt hof is on ðǣre eorðan. Hit is ðæs mannes hof.
% mdparagraph:P181:end

% item: CH01-MINI02
% mdparagraph:P182:begin
\itempara{2}{Translate the passage and rewrite it in the present. Explain why \textbf{hēo} does not establish the creature's sex. Then rewrite just the first two sentences about several creatures in the same dwelling, in the present; preserve “animals” as the predicate noun. Identify the reference of \textbf{hit} in the original passage and the case of \textbf{gesceafte}.}
% mdparagraph:P182:end

% mdparagraph:P183:begin
Sēo gesceaft wæs on ðǣm hofe. Hēo wæs dēor. Ðæt sǣd wæs on ðǣm hofe. Hit wæs ðǣre gesceafte męte.
% mdparagraph:P183:end

\reserve{4\baselineskip}
% mdparagraph:P184:begin
\section*{16. The first day of Creation}\BookHeading{nav-ch01-learner-0020}{subsection}{16. The first day of Creation}{1}
% mdparagraph:P184:end

% mdparagraph:P185:begin
Cook, \emph{A First Book in Old English}, Reader I, pp.123–124. The following is a continuous extract through the close of the first day. Only note cues and printed line/\allowbreak{}page breaks have been removed.[15]
% mdparagraph:P185:end

% source: FB-I-CREATION-FIRST-DAY-01
% mdparagraph:P186:begin
\readingtext{On anginne gescēop God heofonan and eorðan. Sēo eorðe sōðlice wæs īdel and ǣmtig; and ðīestru wǣron ofer ðǣre nēowolnesse brādnesse; and Godes gāst wæs gefęred ofer wæteru. God cwæð ðā, “Geweorðe lēoht”; and lēoht wearð geworht. God geseah ðā ðæt hit gōd wæs; and hē gedǣlde ðæt lēoht fram ðǣm ðīestrum. And hēt ðæt lēoht Dæg, and ðā ðīestru Niht. Ðā wæs geworden ǣfen and morgen ān dæg.}
% mdparagraph:P186:end

\reserve{4\baselineskip}
% mdparagraph:P187:begin
\subsection*{Supplied help}\BookHeading{nav-ch01-learner-0021}{subsubsection}{Supplied help}{0}
% mdparagraph:P187:end

% teaching: CH01-READ-SUPPORT CH01-GRAM-ORDER
% mdparagraph:P188:begin
Use this help for the questions and translation. It supplies the actual forms and constructions occurring here; their independent production is not required.
% mdparagraph:P188:end

% mdtable:T20:begin
% course-table-style: 2026-10
% course-panel:T20-P01:begin
\par\begingroup
\fontsize{10.5}{13.2}\selectfont
\setlength{\LTleft}{0pt}\setlength{\LTright}{\fill}
\begin{longtable}{@{}>{\raggedright\arraybackslash}p{32.28mm}>{\raggedright\arraybackslash}p{88.908mm}@{}}
\toprule
\textbf{Expression} & \textbf{Local meaning and analysis} \\
\midrule
\endfirsthead
\toprule
\textbf{Expression} & \textbf{Local meaning and analysis} \\
\midrule
\endhead
\midrule
\multicolumn{2}{r}{\footnotesize Continued on the following page.}\\
\endfoot
\bottomrule
\endlastfoot
\textbf{On anginne} & “In the beginning”: a supplied temporal use of \textbf{on}, with dative \textbf{anginne}. \\
\textbf{gescēop} & “Created,” third-person singular preterite. In the opening, the subject \textbf{God} follows the verb; word order alone does not identify the object. \\
\textbf{heofonan} & “Heaven”: accusative \textbf{singular} of weak feminine \textbf{heofone}, a variant of the masculine \textbf{heofon} learned above. It is not masculine \textbf{heofon}'s plural. \\
\textbf{sōðlice; īdel and ǣmtig} & Connective “now/\allowbreak{}for its part”; “empty and void.” The two adjectives are predicate feminine nominative singular with \textbf{eorðe}. \\
\textbf{ofer ðǣre nēowolnesse brādnesse} & \textbf{Ofer} means “over” and here requires a dative head. Apply the nested-genitive construction of §9 to identify the head and its dependent. \\
\textbf{wæs gefęred ofer wæteru} & “Was borne/\allowbreak{}carried over the waters,” contextually “moved over the waters.” \textbf{Wæs gefęred} is a passive phrase; \textbf{gefęred} is a past participle. Here \textbf{ofer} has accusative plural \textbf{wæteru}, unlike the preceding dative construction. \\
\textbf{cwæð ðā} & “Then said”: \textbf{cwæð} is third-person singular preterite; this \textbf{ðā} is the adverb “then,” not the article. \\
\textbf{Geweorðe lēoht} & “Let light come to be.” \textbf{Geweorðe} is a finite third-person singular present subjunctive (also called optative), expressing a wish or command. \\
\textbf{lēoht wearð geworht} & “Light was made/\allowbreak{}came into being.” \textbf{Wearð} is a third-person singular preterite used here with past participle \textbf{geworht} in a passive phrase. \\
\textbf{geseah ðā ðæt hit gōd wæs} & “Then saw that it was good.” \textbf{Geseah} is third-person singular preterite; \textbf{ðā} is “then”; \textbf{ðæt} is the conjunction “that.” \textbf{Gōd} is a neuter nominative singular predicate adjective. The copula stands at the end of the embedded clause. Use the context and §8 for the reference of \textbf{hit}. \\
\textbf{hē gedǣlde} & \textbf{Gedǣlde} is “separated,” third-person singular preterite. Identify the reference of the subject pronoun from the preceding clause. The object and \textbf{fram} phrase use §9. \\
\textbf{hēt ðæt lēoht Dæg, and ðā ðīestru Niht} & “[He] called the light Day, and the darkness Night.” \textbf{Hēt} is third-person singular preterite, with God/\allowbreak{}he understood as subject. \textbf{Ðæt lēoht} and \textbf{ðā ðīestru} are accusative objects; \textbf{Dæg} and \textbf{Niht} are accusative names applied to them. Here \textbf{ðæt} and \textbf{ðā} are articles. \\
\textbf{Ðā wæs geworden ǣfen and morgen ān dæg.} & “Then there came to be evening and morning: one day.” \textbf{Wæs geworden} expresses a resulting state or completed coming-to-be, “had come to be”; \textbf{geworden} is a past participle. \textbf{Ān dæg} is “one day,” not the ordinal “first day.” Treat this closing formula as a whole. \\
\end{longtable}
\par\endgroup
% course-panel:T20-P01:end
% mdtable:T20:end

\reserve{4\baselineskip}
% mdparagraph:P189:begin
\subsection*{Questions and translation}\BookHeading{nav-ch01-learner-0022}{subsubsection}{Questions and translation}{0}
% mdparagraph:P189:end

% mdparagraph:P190:begin
The unfamiliar verb constructions remain supplied. Questions 4 and 7 check that help; question 8 combines it with the noun, pronoun, and copula rules learned in this unit. They do not require independent recall of the subjunctive, passive, or naming construction.
% mdparagraph:P190:end

% item: CH01-READ01
% mdparagraph:P191:begin
\itempara{1}{Using the local support, identify the subject and two objects of the first sentence. Is heofonan plural? Translate the sentence.}
% mdparagraph:P191:end

% item: CH01-READ02
% mdparagraph:P192:begin
\itempara{2}{Describe the earth and the darkness. Explain wǣron, then identify the head and genitive dependent in the phrase after ofer.}
% mdparagraph:P192:end

% item: CH01-READ03
% mdparagraph:P193:begin
\itempara{3}{With the supplied passive phrase, state what God's spirit was doing and over what. Explain why wæteru cannot be the selected dative plural form.}
% mdparagraph:P193:end

% item: CH01-READ04
% mdparagraph:P194:begin
\itempara{4}{Does God's quoted speech report an existing fact or command/\allowbreak{}wish for something to come about? State what follows.}
% mdparagraph:P194:end

% item: CH01-READ05
% mdparagraph:P195:begin
\itempara{5}{In God geseah ðā ðæt hit gōd wæs, use the supplied functions of ðā and ðæt, and identify the reference of hit from the context. Translate the clause.}
% mdparagraph:P195:end

% item: CH01-READ06
% mdparagraph:P196:begin
\itempara{6}{What did God separate, and from what? Identify the direct-object and separation phrases and their cases.}
% mdparagraph:P196:end

% item: CH01-READ07
% mdparagraph:P197:begin
\itempara{7}{Using the naming and closing-formula support, state the names given, the understood subject of hēt, and the number of days in the close.}
% mdparagraph:P197:end

% item: CH01-READ08
% mdparagraph:P198:begin
\itempara{8}{Translate the entire first-day passage from the printed Old English, using the supplied help. Preserve the sequence, command, naming and one-day close.}
% mdparagraph:P198:end

\reserve{4\baselineskip}
% mdparagraph:P199:begin
\section*{17. Review}\BookHeading{nav-ch01-learner-0023}{subsection}{17. Review}{1}
% mdparagraph:P199:end

% item: CH01-REVIEW01
% mdparagraph:P200:begin
\itempara{1}{A learner writes *dægas and *hwælas for the plurals of dæg and hwæl, then changes wæstm to *wastmas. Correct the intended plurals and explain which vowel alternation applies.}
% mdparagraph:P200:end

% item: CH01-REVIEW02
% mdparagraph:P201:begin
\itempara{2}{In \textbf{on ðāra fugla trēowum}, a learner takes \textbf{ðāra} to agree with \textbf{trēowum}. Translate the phrase; identify the head and dependent, giving the case and number of each. Explain the article's actual agreement.}
% mdparagraph:P201:end

% item: CH01-REVIEW03
% mdparagraph:P202:begin
\itempara{3}{List the possible case/\allowbreak{}number values of \textbf{steorran} in the selected paradigm. Which applies in \textbf{ðæs steorran lēoht}, and which in \textbf{Ðā steorran wǣron on ðǣm heofone}? Translate both. Explain what the context adds that the ending alone cannot tell you.}
% mdparagraph:P202:end

% item: CH01-REVIEW04
% mdparagraph:P203:begin
\itempara{4}{Correct the agreement in \textbf{Ðā ðīestru is on ðǣm heofone. Ðæt lēoht wǣron on ðǣm hofe.} Keep the first sentence present and the second past; explain the corrections. Then complete \textbf{Se mann gedǣlde \_\_\_\_\_ fram \_\_\_\_\_.} to mean “The man separated the birds from the animal.” Use \textbf{dēor} for “animal.” The supplied \textbf{gedǣlde} means “separated,” third-person singular preterite. Identify the case and number of both completed phrases.}
% mdparagraph:P203:end

\reserve{4\baselineskip}
% mdparagraph:P204:begin
\section*{Source notes}\BookHeading{nav-ch01-learner-0024}{subsection}{Source notes}{1}
% mdparagraph:P204:end

% mdparagraph:P205:begin
\itempara{1}{Cook, \emph{A First Book in Old English}, second edition (1895), §§2–4, pp.5–6 /\allowbreak{} PDF21–22. The orthography follows the selected edition; this is not a historical account of all spellings or a complete pronunciation lesson.}
% mdparagraph:P205:end

% mdparagraph:P206:begin
\itempara{2}{Cook, §§38–42, pp.26–27 /\allowbreak{} PDF42–43. Tables composed from the checked paradigms are teaching models; they do not claim separate textual attestation of every form.}
% mdparagraph:P206:end

% mdparagraph:P207:begin
\itempara{3}{Cook, §84, p.50 /\allowbreak{} PDF66. The masculine article is \textbf{se}; the source distinguishes emphatic \textbf{sē}. Instrumental syntax is reserved for later treatment.}
% mdparagraph:P207:end

% mdparagraph:P208:begin
\itempara{4}{Cook, §§43–44, pp.28–30 /\allowbreak{} PDF44–46; especially §43.2 for \textbf{dæg/\allowbreak{}hwæl} and §43.4 for syncope and retained vowels.}
% mdparagraph:P208:end

% mdparagraph:P209:begin
\itempara{5}{Cook, §47, pp.32–33 /\allowbreak{} PDF48–49. The tables retain the selected \textbf{trēowu, tācnu, wæteru, nīetenu, ǣfenu} patterns. The neuter \textbf{ǣfen} plural is derived from its checked pattern, not quoted from a newly located passage.}
% mdparagraph:P209:end

% mdparagraph:P210:begin
\itempara{6}{Cook, §51, pp.34–36 /\allowbreak{} PDF50–52; §51.1 and the vocabulary's \textbf{wyrt} reference, p.324 /\allowbreak{} PDF340, support the endingless accusatives. The alternative N/\allowbreak{}A plural \textbf{-e} and singular \textbf{-a} of the \textbf{-ung} group are source variants.}
% mdparagraph:P210:end

% mdparagraph:P211:begin
\itempara{7}{Cook, §53, p.37 /\allowbreak{} PDF53.}
% mdparagraph:P211:end

% mdparagraph:P212:begin
\itempara{8}{Cook, §§46,52, pp.31–32,36 /\allowbreak{} PDF47–48,52; \textbf{mann/\allowbreak{}męnn}, §17, p.13 /\allowbreak{} PDF29; \textbf{niht}, Reader p.125 with note 28 /\allowbreak{} PDF141; \textbf{sǣ}, vocabulary p.324 /\allowbreak{} PDF340. Wright and Wright, \emph{Old English Grammar}, second edition (1914), §388, p.187 /\allowbreak{} PDF207, explicitly marks genitive plural \textbf{*sǣwa} as reconstructed. It is not tested here.}
% mdparagraph:P212:end

% mdparagraph:P213:begin
\itempara{9}{Cook, §138, pp.78–79 /\allowbreak{} PDF94–95. The bounded indicative series is not the complete \textbf{wesan/\allowbreak{}bēon} system.}
% mdparagraph:P213:end

% mdparagraph:P214:begin
\itempara{10}{Cook, §§149–150, p.88 /\allowbreak{} PDF104; §189, p.102 /\allowbreak{} PDF118.}
% mdparagraph:P214:end

% mdparagraph:P215:begin
\itempara{11}{Cook, \emph{Exercises in Old English} (1895), Exercise 3 introductory note, p.2 /\allowbreak{} PDF10; \emph{First Book}, §153, pp.89–90 /\allowbreak{} PDF105–106, especially §153.i.}
% mdparagraph:P215:end

% mdparagraph:P216:begin
\itempara{12}{Cook, \emph{First Book}, §§160,164(a), pp.93–94 /\allowbreak{} PDF109–110; §166, p.95 /\allowbreak{} PDF111; §§168,172, pp.96–97 /\allowbreak{} PDF112–113. \textbf{Sealde} is printed in §117, p.72 /\allowbreak{} PDF88. The other supplied frames occur in the assigned Creation extract. Examples in §9 are new composition. Supplied finite forms do not introduce their full conjugations.}
% mdparagraph:P216:end

% mdparagraph:P217:begin
\itempara{13}{The noun choices follow Cook's English–Old English vocabulary and the checked \textbf{First Book} paradigms. \textbf{Wæstm}, with short \textbf{æ}, corrects Cook's \textbf{wǣstm}; this correction applies also to its inflected forms. The source spelling is preserved in the editorial record. \textbf{Gærs}, \textbf{męte}, and the selected noun genders are retained. \textbf{Dēor} is not restricted to modern English “deer”; \textbf{nīeten} supplies “cattle” in the plural.}
% mdparagraph:P217:end

% mdparagraph:P218:begin
\itempara{14}{All 46 prompts are reproduced from Cook, \emph{Exercises}, Group I, Exercises 1–3, pp.1–2 /\allowbreak{} PDF9–10. Their English wording and original numbering are unchanged. Directions and answers are new teaching material.}
% mdparagraph:P218:end

% mdparagraph:P219:begin
\itempara{15}{Cook, \emph{First Book}, Reader I, pp.123–124 /\allowbreak{} PDF139–140: the five prose lines on p.123 are editorially counted; the first three on p.124 are counted from the printed marginal line number 5. The extract preserves \textbf{heofonan, gefęred, Geweorðe}, and unlengthened \textbf{geseah}. Its apparatus and original lineation remain separately transcribed. All PDF page numbers in these notes are one-based and refer to the selected scans.}
% mdparagraph:P219:end

% mdparagraph:P220:begin
\itempara{16}{Cook, §81, pp.48–49 /\allowbreak{} PDF64–65, gives the personal-pronoun paradigms. The subject inventory selects \textbf{hēo} and \textbf{hīe}, with \textbf{ðū} following this unit's spelling convention. Ringe and Taylor, \emph{The Development of Old English} (2014), §1.1.2, p.6 /\allowbreak{} preview PDF21, explicitly states that the dual pronouns take plural verbs. Fulk's Genesis 19 reading, p.38 /\allowbreak{} PDF50, uses ordinary second- and first-person plurals for the two angels, illustrating that plurals can include two. Only nominative subject use is introduced here.}
% mdparagraph:P220:end

% mdparagraph:P221:begin
\itempara{17}{Fulk, \emph{An Introductory Grammar of Old English} (2020), glossary p.265 /\allowbreak{} PDF277, prints unpointed \textbf{cyn(n)}; §18, p.12 /\allowbreak{} PDF24, gives its \textbf{k} value. The stopped initial of \textbf{gærs} is a source-based inference: Wright and Wright, \emph{Old English Grammar}, second edition (1914), §§314–315, p.156 /\allowbreak{} PDF176, classify \textbf{græs} with guttural \textbf{g} and describe its stopped initial outcome; §§66 note 4 and 345, pp.43,171 /\allowbreak{} PDF63,191, connect \textbf{gærs} with late r-metathesis. Compare Wright §85 note 1, p.47 /\allowbreak{} PDF67, for breaking before this metathesis, and Fulk §72, p.46 /\allowbreak{} PDF58, for palatalization before breaking. The learner needs only the supplied sound, not this historical derivation.}
% mdparagraph:P221:end


% book-source-body:end
\par\label{end-ch01-learner}
\clearpage\endgroup

\BookRecto
\begingroup
\renewcommand{\runningtitle}{Unit 2. Pronouns and weak verbs}
\hyphenpenalty=10000
\exhyphenpenalty=10000
\emergencystretch=3em
\newcommand{\itempara}[2]{\par\addvspace{3pt}\noindent\hangindent=1.8em\hangafter=1\makebox[1.8em][l]{#1.}#2\par}
\newcommand{\reserve}[1]{\par\begingroup\dimen0=\pagegoal\advance\dimen0 by-\pagetotal\ifdim\dimen0<#1\newpage\fi\endgroup}
% book-source-body:begin

\raggedright
\booktitle{\BookComponent{ch02-learner}{Unit 2. Personal pronouns and weak indicative verbs}UNIT 2}{Personal pronouns and weak indicative verbs}

% mdparagraph:P001:begin
This unit completes the personal-pronoun tables begun in Unit 1 and introduces the present and preterite indicative of weak verbs. Learn the finite forms in the tables, together with their stem changes. The six first- and second-person genitives and the negative verb \textbf{nabban} are for recognition. Infinitive and participle citations are available for lookup; their uses in sentences come later.
% mdparagraph:P001:end

\reserve{7\baselineskip}
% mdparagraph:P002:begin
\section*{1. Personal pronouns}\BookHeading{nav-ch02-learner-0001}{subsection}{1. Personal pronouns}{1}\BookLabel{bk-ch02-learner-b542d326d059}
% mdparagraph:P002:end

% mdparagraph:P003:begin
The tables repeat the subject forms learned in Unit 1 and add the other cases. First person identifies the speaker; second person, the addressee; third person, someone or something else. The first and second persons also have a \textbf{dual}, explicitly referring to two. As in Unit 1, a dual subject takes a \textbf{plural verb form}; there is no separate dual verb ending. Use the dual when an exercise requests “we two” or “you two”; ordinary plurals can also include two. [1]
% mdparagraph:P003:end

% teaching: NEW-U02-PAR-ic NEW-U02-PAR-thu NEW-U02-PAR-he NEW-U02-VOC-ic NEW-U02-VOC-thu NEW-U02-VOC-he NEW-U02-GRAM-personal-pronouns NEW-U02-GRAM-personal-genitive-citation NEW-U02-GRAM-dual-agreement
% mdtable:T01:begin
% course-table-style: 2026-10
\par\addvspace{3pt}
% course-caption:T01:0:begin
\textbf{First person}\par\nobreak
% course-caption:T01:0:end
% course-panel:T01-P01:begin
\par\begingroup
\fontsize{10.5}{13.2}\selectfont
\setlength{\LTleft}{0pt}\setlength{\LTright}{\fill}
\noindent
\begin{tabular}{@{}>{\raggedright\arraybackslash}p{22.0mm}>{\raggedright\arraybackslash}p{31.188mm}>{\raggedright\arraybackslash}p{31.188mm}>{\raggedright\arraybackslash}p{31.188mm}@{}}
\toprule
\textbf{Case} & \textbf{Singular: I} & \textbf{Dual: we two} & \textbf{Plural: we} \\
\midrule
Nominative & ic & wit & wē \\
Genitive & mīn & uncer & ūre \\
Dative & mē & unc & ūs \\
Accusative & mē & unc & ūs \\
\bottomrule
\end{tabular}
\par\endgroup
% course-panel:T01-P01:end
\par\addvspace{5pt}
\par\addvspace{3pt}
% course-caption:T01:1:begin
\textbf{Second person}\par\nobreak
% course-caption:T01:1:end
% course-panel:T01-P02:begin
\par\begingroup
\fontsize{10.5}{13.2}\selectfont
\setlength{\LTleft}{0pt}\setlength{\LTright}{\fill}
\noindent
\begin{tabular}{@{}>{\raggedright\arraybackslash}p{22.0mm}>{\raggedright\arraybackslash}p{31.188mm}>{\raggedright\arraybackslash}p{31.188mm}>{\raggedright\arraybackslash}p{31.188mm}@{}}
\toprule
\textbf{Case} & \textbf{Singular: you} & \textbf{Dual: you two} & \textbf{Plural: you} \\
\midrule
Nominative & ðū & git & gē \\
Genitive & ðīn & incer & ēower \\
Dative & ðē & inc & ēow \\
Accusative & ðē & inc & ēow \\
\bottomrule
\end{tabular}
\par\endgroup
% course-panel:T01-P02:end
\par\addvspace{5pt}
\par\addvspace{3pt}
% course-caption:T01:2:begin
\textbf{Third person}\par\nobreak
% course-caption:T01:2:end
% course-panel:T01-P03:begin
\par\begingroup
\fontsize{10.5}{13.2}\selectfont
\setlength{\LTleft}{0pt}\setlength{\LTright}{\fill}
\noindent
\begin{tabular}{@{}>{\raggedright\arraybackslash}p{22.0mm}>{\raggedright\arraybackslash}p{22.688mm}>{\raggedright\arraybackslash}p{22.688mm}>{\raggedright\arraybackslash}p{22.688mm}>{\raggedright\arraybackslash}p{22.688mm}@{}}
\toprule
\textbf{Case} & \textbf{Masculine singular} & \textbf{Neuter singular} & \textbf{Feminine singular} & \textbf{Plural: all genders} \\
\midrule
Nominative & hē & hit & hēo & hīe \\
Genitive & his & his & hiere & hiera \\
Dative & him & him & hiere & him \\
Accusative & hine & hit & hīe & hīe \\
\bottomrule
\end{tabular}
\par\endgroup
% course-panel:T01-P03:end
% mdtable:T01:end
% mdparagraph:P004:begin
Recognize \textbf{mīn, uncer, ūre, ðīn, incer, ēower} as genitives of the persons and numbers shown. Their corresponding possessive adjectives have endings not yet taught. Do not use this table alone to form attributive “my”, “your” or “our” phrases. Review the nominatives; learn the new dative, accusative and third-person genitive cells for the constructions below.
% mdparagraph:P004:end

% teaching: NEW-U02-GRAM-pronoun-reference
% mdparagraph:P005:begin
The\BookPoint{bk-ch02-learner-8f105697c218} nominative is the subject form. The accusative can be a direct object; the dative can name a recipient. Their spelling often coincides: \textbf{mē} and \textbf{ðē} each serve two cases. \textbf{Him} can be singular or plural; \textbf{hīe} can be feminine singular accusative or plural nominative/\allowbreak{}accusative. The sentence and its context settle the analysis.
% mdparagraph:P005:end

% mdparagraph:P006:begin
Third-person\BookPoint{bk-ch02-learner-fc8f17ff0711} pronouns follow the grammatical gender and number of their referents. English “it” therefore need not correspond to \textbf{hit}: a feminine noun can be resumed by \textbf{hēo}, a masculine one by \textbf{hē}. Case still follows the pronoun’s role in its own clause.
% mdparagraph:P006:end

% mdparagraph:P007:begin
The table selects one spelling for each cell. Cook also has \textbf{mec, ðec, ūsic, ēowic} as accusatives, \textbf{ūser} as first-person plural genitive, \textbf{hiene} beside \textbf{hine}, \textbf{hire} beside \textbf{hiere}, and \textbf{hira/\allowbreak{}heora} beside \textbf{hiera}. Other third-person spellings include \textbf{hīo/\allowbreak{}hī} for feminine nominative, \textbf{hī/\allowbreak{}hēo} for feminine accusative and plural nominative/\allowbreak{}accusative. These are reference alternatives, not extra recall requirements. Preserve the quantities in the selected forms: \textbf{hīe}, but \textbf{hine, hiere, hiera}.
% mdparagraph:P007:end

% teaching: NEW-U02-GRAM-third-person-possessors
% mdparagraph:P008:begin
Third-person\BookPoint{bk-ch02-learner-62b2d81d3d34} genitives can be used as invariant possessors: \textbf{his}, “his/\allowbreak{}its”; \textbf{hiere}, “her”; \textbf{hiera}, “their”. Choose the form by the possessor. It does not change to agree with the possessed noun, which retains its own case and number. Thus the singular or plural of a possessed noun does not change \textbf{his}; a dative possessed noun does not turn \textbf{hiera} into \textbf{him}. These genitives precede the noun without an added article in our model phrases. [2]
% mdparagraph:P008:end

\reserve{7\baselineskip}
% mdparagraph:P009:begin
\section*{2. Reading the verb tables}\BookHeading{nav-ch02-learner-0002}{subsection}{2. Reading the verb tables}{1}\BookLabel{bk-ch02-learner-925bbc29948d}
% mdparagraph:P009:end

% teaching: NEW-U02-GRAM-weak-indicative-scope
% mdparagraph:P010:begin
Weak verbs form their preterite with a dental element, \textbf{d} or \textbf{t}, sometimes accompanied by a stem change. The \textbf{indicative} states an action or state; the \textbf{preterite} is the past tense. The three classes below must be learned with the verb: the infinitive ending alone does not identify the class.
% mdparagraph:P010:end

% mdparagraph:P011:begin
Every paradigm gives the complete singular and plural present and preterite indicative. \textbf{1–3 pl.} includes first, second and third plural; the shared plural row also supplies the form for dual subjects. In principal-part citations, \textbf{1/\allowbreak{}3 sg.} means both first and third singular. All twelve person–number–tense combinations remain explicit.
% mdparagraph:P011:end

% mdparagraph:P012:begin
In the selected preterite, the first and third singular coincide. Add \textbf{-st} to that form for the second singular; replace its final \textbf{-e} with \textbf{-on} for all three plural persons. The present singular distinguishes three persons, while all three plural persons coincide. Observe the actual stems and contractions; do not simply attach endings to an infinitive. [3]
% mdparagraph:P012:end

\reserve{7\baselineskip}
% mdparagraph:P013:begin
\section*{3. Class I: long stems and dental changes}\BookHeading{nav-ch02-learner-0003}{subsection}{3. Class I: long stems and dental changes}{1}\BookLabel{bk-ch02-learner-fee1ccca58cc}
% mdparagraph:P013:end

% teaching: NEW-U02-GRAM-weak-I-long
% mdparagraph:P014:begin
The basic present endings are \textbf{-e, -(e)st, -(e)ð, -að}. In the longer stems, loss of the unstressed \textbf{e} in the second and third singular brings consonants together. The resulting forms must be learned with the preterite. Compare all ten patterns in the following tables. [4]
% mdparagraph:P014:end

% teaching: NEW-U02-PAR-hieran NEW-U02-PAR-gefyllan NEW-U02-PAR-cyssan NEW-U02-PAR-settan NEW-U02-PAR-sendan
% mdtable:T02:begin
% course-table-style: 2026-10
% panel:Present indicative
% panel:Preterite indicative
% course-panel:T02-P01:begin
\par\begingroup
\fontsize{10.5}{13.2}\selectfont
\setlength{\LTleft}{0pt}\setlength{\LTright}{\fill}
\noindent
\begin{tabular}{@{}>{\raggedright\arraybackslash}p{20.075mm}>{\raggedright\arraybackslash}p{20.259mm}>{\raggedright\arraybackslash}p{23.925mm}>{\raggedright\arraybackslash}p{22.917mm}>{\raggedright\arraybackslash}p{25.575mm}@{}}
\toprule
\textbf{Person and number} & \textbf{hīeran · Present} & \textbf{hīeran · Preterite} & \textbf{gefyllan · Present} & \textbf{gefyllan · Preterite} \\
\midrule
1 sg. & hīere & hīerde & gefylle & gefylde \\
2 sg. & hīerst & hīerdest & gefylst & gefyldest \\
3 sg. & hīerð & hīerde & gefylð & gefylde \\
1–3 pl. & hīerað & hīerdon & gefyllað & gefyldon \\
\bottomrule
\end{tabular}
\par\endgroup
% course-panel:T02-P01:end
\par\addvspace{5pt}
% course-panel:T02-P02:begin
\par\begingroup
\fontsize{10.5}{13.2}\selectfont
\setlength{\LTleft}{0pt}\setlength{\LTright}{\fill}
\noindent
\begin{tabular}{@{}>{\raggedright\arraybackslash}p{20.873mm}>{\raggedright\arraybackslash}p{21.064mm}>{\raggedright\arraybackslash}p{24.876mm}>{\raggedright\arraybackslash}p{21.064mm}>{\raggedright\arraybackslash}p{24.876mm}@{}}
\toprule
\textbf{Person and number} & \textbf{cyssan · Present} & \textbf{cyssan · Preterite} & \textbf{sęttan · Present} & \textbf{sęttan · Preterite} \\
\midrule
1 sg. & cysse & cyste & sętte & sętte \\
2 sg. & cyst & cystest & sętst & sęttest \\
3 sg. & cyst & cyste & sętt & sętte \\
1–3 pl. & cyssað & cyston & sęttað & sętton \\
\bottomrule
\end{tabular}
\par\endgroup
% course-panel:T02-P02:end
\par\addvspace{5pt}
\par\addvspace{3pt}
% course-caption:T02:2:begin
\textbf{sęndan}\par\nobreak
% course-caption:T02:2:end
% course-panel:T02-P03:begin
\par\begingroup
\fontsize{10.5}{13.2}\selectfont
\setlength{\LTleft}{0pt}\setlength{\LTright}{\fill}
\noindent
\begin{tabular}{@{}>{\raggedright\arraybackslash}p{39.459mm}>{\raggedright\arraybackslash}p{39.459mm}>{\raggedright\arraybackslash}p{39.459mm}@{}}
\toprule
\textbf{Person and number} & \textbf{Present} & \textbf{Preterite} \\
\midrule
1 sg. & sęnde & sęnde \\
2 sg. & sęntst & sęndest \\
3 sg. & sęnt & sęnde \\
1–3 pl. & sęndað & sęndon \\
\bottomrule
\end{tabular}
\par\endgroup
% course-panel:T02-P03:end
% mdtable:T02:end
% mdparagraph:P015:begin
\textbf{Hīeran} has ordinary \textbf{-de}. In \textbf{gefyllan}, doubled \textbf{ll} becomes single before another consonant: \textbf{gefylst, gefylð, gefylde}. The \textbf{ss/\allowbreak{}ff/\allowbreak{}pp} family has \textbf{-te} with simplification, here \textbf{cyssan → cyste}. Its contracted present second and third singular are both \textbf{cyst}. \textbf{Sęttan} retains \textbf{tt} in \textbf{sętte}; \textbf{sęndan} has \textbf{sęnde}, without an added written dental. First singular present and first/\allowbreak{}third singular preterite can consequently look alike.
% mdparagraph:P015:end

% teaching: NEW-U02-PAR-laedan NEW-U02-PAR-iecan NEW-U02-PAR-gerestan NEW-U02-PAR-metan NEW-U02-PAR-gierwan
% mdtable:T03:begin
% course-table-style: 2026-10
% panel:Present indicative
% panel:Preterite indicative
% course-panel:T03-P01:begin
\par\begingroup
\fontsize{10.5}{13.2}\selectfont
\setlength{\LTleft}{0pt}\setlength{\LTright}{\fill}
\noindent
\begin{tabular}{@{}>{\raggedright\arraybackslash}p{21.528mm}>{\raggedright\arraybackslash}p{20.84mm}>{\raggedright\arraybackslash}p{24.772mm}>{\raggedright\arraybackslash}p{20.84mm}>{\raggedright\arraybackslash}p{24.772mm}@{}}
\toprule
\textbf{Person and number} & \textbf{lǣdan · Present} & \textbf{lǣdan · Preterite} & \textbf{īecan · Present} & \textbf{īecan · Preterite} \\
\midrule
1 sg. & lǣde & lǣdde & īece & īecte \\
2 sg. & lǣtst & lǣddest & īecst & īectest \\
3 sg. & lǣtt & lǣdde & īecð & īecte \\
1–3 pl. & lǣdað & lǣddon & īecað & īecton \\
\bottomrule
\end{tabular}
\par\endgroup
% course-panel:T03-P01:end
\par\addvspace{5pt}
% course-panel:T03-P02:begin
\par\begingroup
\fontsize{10.5}{13.2}\selectfont
\setlength{\LTleft}{0pt}\setlength{\LTright}{\fill}
\noindent
\begin{tabular}{@{}>{\raggedright\arraybackslash}p{20.19mm}>{\raggedright\arraybackslash}p{24.062mm}>{\raggedright\arraybackslash}p{25.722mm}>{\raggedright\arraybackslash}p{19.545mm}>{\raggedright\arraybackslash}p{23.233mm}@{}}
\toprule
\textbf{Person and number} & \textbf{geręstan · Present} & \textbf{geręstan · Preterite} & \textbf{mētan · Present} & \textbf{mētan · Preterite} \\
\midrule
1 sg. & geręste & geręste & mēte & mētte \\
2 sg. & geręstest & geręstest & mētst & mēttest \\
3 sg. & geręst & geręste & mētt & mētte \\
1–3 pl. & geręstað & geręston & mētað & mētton \\
\bottomrule
\end{tabular}
\par\endgroup
% course-panel:T03-P02:end
\par\addvspace{5pt}
\par\addvspace{3pt}
% course-caption:T03:2:begin
\textbf{gierwan}\par\nobreak
% course-caption:T03:2:end
% course-panel:T03-P03:begin
\par\begingroup
\fontsize{10.5}{13.2}\selectfont
\setlength{\LTleft}{0pt}\setlength{\LTright}{\fill}
\noindent
\begin{tabular}{@{}>{\raggedright\arraybackslash}p{39.459mm}>{\raggedright\arraybackslash}p{39.459mm}>{\raggedright\arraybackslash}p{39.459mm}@{}}
\toprule
\textbf{Person and number} & \textbf{Present} & \textbf{Preterite} \\
\midrule
1 sg. & gierwe & gierede \\
2 sg. & gierest & gieredest \\
3 sg. & giereð & gierede \\
1–3 pl. & gierwað & gieredon \\
\bottomrule
\end{tabular}
\par\endgroup
% course-panel:T03-P03:end
% mdtable:T03:end
% mdparagraph:P016:begin
Keep\BookPoint{bk-ch02-learner-1310ddcd10cb} these contrasts together:
% mdparagraph:P016:end

\begin{itemize}
% mdparagraph:P017:begin
\item After consonant + \textbf{d}, \textbf{sęndan → sęnde}; after vowel + \textbf{d}, \textbf{lǣdan → lǣdde}.
% mdparagraph:P017:end
% mdparagraph:P018:begin
\item After consonant + \textbf{t}, \textbf{geręstan → geręste}; after vowel + \textbf{t}, \textbf{mētan → mētte}. Vowel length does not itself remove a dental.
% mdparagraph:P018:end
% mdparagraph:P019:begin
\item The \textbf{c/\allowbreak{}p/\allowbreak{}x} family has \textbf{t} in the preterite, here \textbf{īecan → īecte}. The exceptional \textbf{gewyrcean → geworhte} below must be learned separately.
% mdparagraph:P019:end
% mdparagraph:P020:begin
\item Before the present second/\allowbreak{}third singular endings, dental combinations give \textbf{sęntst/\allowbreak{}sęnt}, \textbf{lǣtst/\allowbreak{}lǣtt}, and \textbf{mētst/\allowbreak{}mētt}. Final \textbf{tt} can also be written \textbf{t}: \textbf{sęt, lǣt, mēt} are valid third-singular alternatives.
% mdparagraph:P020:end
% mdparagraph:P021:begin
\item \textbf{Geręstest} is our selected second singular present; contracted \textbf{geręst} is also valid. The selected third singular is \textbf{geręst}. Do not add another \textbf{t} to the preterite \textbf{geręste}.
% mdparagraph:P021:end
% mdparagraph:P022:begin
\item In the \textbf{rw/\allowbreak{}lw} family, \textbf{gierwan → gierede}. The selected present paradigm retains \textbf{w} in the first singular and plural, but uses \textbf{gierest, giereð} in the other singular persons. The sources also allow present forms with \textbf{w} retained or omitted; the selected paradigm is not a rule deleting every present \textbf{w}.
% mdparagraph:P022:end
\end{itemize}

\reserve{7\baselineskip}
% mdparagraph:P023:begin
\section*{4. Class I: irregular and short stems}\BookHeading{nav-ch02-learner-0004}{subsection}{4. Class I: irregular and short stems}{1}\BookLabel{bk-ch02-learner-759d9baec3e2}
% mdparagraph:P023:end

% teaching: NEW-U02-GRAM-weak-I-irregular
% mdparagraph:P024:begin
Learn the past stems of \textbf{sęllan} and \textbf{gewyrcean} with their meanings. Neither is predicted by the regular long-stem pattern. The lexical prefix \textbf{ge-} is retained in the forms of \textbf{gewyrcean}; no general meaning or tense rule for that prefix is implied. [5]
% mdparagraph:P024:end

% teaching: NEW-U02-PAR-sellan NEW-U02-PAR-gewyrcean
% mdtable:T04:begin
% course-table-style: 2026-10
% panel:Present indicative
% panel:Preterite indicative
% course-panel:T04-P01:begin
\par\begingroup
\fontsize{10.5}{13.2}\selectfont
\setlength{\LTleft}{0pt}\setlength{\LTright}{\fill}
\noindent
\begin{tabular}{@{}>{\raggedright\arraybackslash}p{19.443mm}>{\raggedright\arraybackslash}p{19.621mm}>{\raggedright\arraybackslash}p{23.172mm}>{\raggedright\arraybackslash}p{23.971mm}>{\raggedright\arraybackslash}p{26.546mm}@{}}
\toprule
\textbf{Person and number} & \textbf{sęllan · Present} & \textbf{sęllan · Preterite} & \textbf{gewyrcean · Present} & \textbf{gewyrcean · Preterite} \\
\midrule
1 sg. & sęlle & sealde & gewyrce & geworhte \\
2 sg. & sęlst & sealdest & gewyrcst & geworhtest \\
3 sg. & sęlð & sealde & gewyrcð & geworhte \\
1–3 pl. & sęllað & sealdon & gewyrceað & geworhton \\
\bottomrule
\end{tabular}
\par\endgroup
% course-panel:T04-P01:end
% mdtable:T04:end
% mdparagraph:P025:begin
Cook also prints \textbf{sęlest, sęleð} beside \textbf{sęlst, sęlð}, and \textbf{gewyrcan/\allowbreak{}gewyrcað} beside the infinitive \textbf{gewyrcean} and plural \textbf{gewyrceað}. These alternatives do not change the past \textbf{geworhte}.
% mdparagraph:P025:end

% teaching: NEW-U02-GRAM-weak-I-short NEW-U02-GRAM-weak-I-ian
% mdparagraph:P026:begin
The short-stem families have a linking \textbf{e} in the preterite. \textbf{Gefręmman} retains \textbf{mm} in the present first singular and plural, but has single \textbf{m} before \textbf{-est, -eð, -ede}. The consonant + sonorant type \textbf{ęfnan} has \textbf{ęfnede}. Short Class I \textbf{nęrian} retains \textbf{i} in the present first singular and plural, but has \textbf{nęrest, nęreð, nęrede}.
% mdparagraph:P026:end

% teaching: NEW-U02-PAR-gefremman NEW-U02-PAR-efnan NEW-U02-PAR-nerian
% mdtable:T05:begin
% course-table-style: 2026-10
% panel:Present indicative
% panel:Preterite indicative
% course-panel:T05-P01:begin
\par\begingroup
\fontsize{10.5}{13.2}\selectfont
\setlength{\LTleft}{0pt}\setlength{\LTright}{\fill}
\noindent
\begin{tabular}{@{}>{\raggedright\arraybackslash}p{19.551mm}>{\raggedright\arraybackslash}p{24.104mm}>{\raggedright\arraybackslash}p{27.675mm}>{\raggedright\arraybackslash}p{18.926mm}>{\raggedright\arraybackslash}p{22.497mm}@{}}
\toprule
\textbf{Person and number} & \textbf{gefręmman · Present} & \textbf{gefręmman · Preterite} & \textbf{ęfnan · Present} & \textbf{ęfnan · Preterite} \\
\midrule
1 sg. & gefręmme & gefręmede & ęfne & ęfnede \\
2 sg. & gefręmest & gefręmedest & ęfnest & ęfnedest \\
3 sg. & gefręmeð & gefręmede & ęfneð & ęfnede \\
1–3 pl. & gefręmmað & gefręmedon & ęfnað & ęfnedon \\
\bottomrule
\end{tabular}
\par\endgroup
% course-panel:T05-P01:end
\par\addvspace{5pt}
\par\addvspace{3pt}
% course-caption:T05:1:begin
\textbf{nęrian}\par\nobreak
% course-caption:T05:1:end
% course-panel:T05-P02:begin
\par\begingroup
\fontsize{10.5}{13.2}\selectfont
\setlength{\LTleft}{0pt}\setlength{\LTright}{\fill}
\noindent
\begin{tabular}{@{}>{\raggedright\arraybackslash}p{39.459mm}>{\raggedright\arraybackslash}p{39.459mm}>{\raggedright\arraybackslash}p{39.459mm}@{}}
\toprule
\textbf{Person and number} & \textbf{Present} & \textbf{Preterite} \\
\midrule
1 sg. & nęrie & nęrede \\
2 sg. & nęrest & nęredest \\
3 sg. & nęreð & nęrede \\
1–3 pl. & nęriað & nęredon \\
\bottomrule
\end{tabular}
\par\endgroup
% course-panel:T05-P02:end
% mdtable:T05:end
% mdparagraph:P027:begin
For \textbf{ęfnan}, the contracted past \textbf{ęfnde, ęfndest, ęfndon} is also valid. A source exception for later reference is \textbf{lęcgan}, with past \textbf{lęgde/\allowbreak{}lēde}; it is not a recall item here. Cook’s related example \textbf{hyngran} is impersonal: it must not be substituted into an ordinary personal sentence merely because its endings look familiar.
% mdparagraph:P027:end

\reserve{7\baselineskip}
% mdparagraph:P028:begin
\section*{5. Class II}\BookHeading{nav-ch02-learner-0005}{subsection}{5. Class II}{1}\BookLabel{bk-ch02-learner-cef52d0e20f6}
% mdparagraph:P028:end

% teaching: NEW-U02-GRAM-weak-II
% mdparagraph:P029:begin
Class II has present \textbf{-ie, -ast, -að, -iað} and the selected past \textbf{-ode, -odest, -odon}. Contrast \textbf{lufast/\allowbreak{}lufode} with Class I \textbf{nęrest/\allowbreak{}nęrede}. Both infinitives end in \textbf{-ian}, but their classes differ. [6]
% mdparagraph:P029:end

% teaching: NEW-U02-PAR-lufian NEW-U02-PAR-gebletsian NEW-U02-PAR-gehalgian
% mdtable:T06:begin
% course-table-style: 2026-10
% panel:Present indicative
% panel:Preterite indicative
% course-panel:T06-P01:begin
\par\begingroup
\fontsize{10.5}{13.2}\selectfont
\setlength{\LTleft}{0pt}\setlength{\LTright}{\fill}
\noindent
\begin{tabular}{@{}>{\raggedright\arraybackslash}p{18.692mm}>{\raggedright\arraybackslash}p{18.863mm}>{\raggedright\arraybackslash}p{22.277mm}>{\raggedright\arraybackslash}p{24.753mm}>{\raggedright\arraybackslash}p{28.167mm}@{}}
\toprule
\textbf{Person and number} & \textbf{lufian · Present} & \textbf{lufian · Preterite} & \textbf{geblētsian · Present} & \textbf{geblētsian · Preterite} \\
\midrule
1 sg. & lufie & lufode & geblētsie & geblētsode \\
2 sg. & lufast & lufodest & geblētsast & geblētsodest \\
3 sg. & lufað & lufode & geblētsað & geblētsode \\
1–3 pl. & lufiað & lufodon & geblētsiað & geblētsodon \\
\bottomrule
\end{tabular}
\par\endgroup
% course-panel:T06-P01:end
\par\addvspace{5pt}
\par\addvspace{3pt}
% course-caption:T06:1:begin
\textbf{gehālgian}\par\nobreak
% course-caption:T06:1:end
% course-panel:T06-P02:begin
\par\begingroup
\fontsize{10.5}{13.2}\selectfont
\setlength{\LTleft}{0pt}\setlength{\LTright}{\fill}
\noindent
\begin{tabular}{@{}>{\raggedright\arraybackslash}p{39.459mm}>{\raggedright\arraybackslash}p{39.459mm}>{\raggedright\arraybackslash}p{39.459mm}@{}}
\toprule
\textbf{Person and number} & \textbf{Present} & \textbf{Preterite} \\
\midrule
1 sg. & gehālgie & gehālgode \\
2 sg. & gehālgast & gehālgodest \\
3 sg. & gehālgað & gehālgode \\
1–3 pl. & gehālgiað & gehālgodon \\
\bottomrule
\end{tabular}
\par\endgroup
% course-panel:T06-P02:end
% mdtable:T06:end
% mdparagraph:P030:begin
Cook also prints \textbf{lufedon}. His class-wide alternatives include past \textbf{-ade, -ude, -ede} and infinitive \textbf{-ig(e)an}; these are reference information, not instructions to generate every spelling for every verb. The dictionary has \textbf{geblēdsian} beside \textbf{geblētsian}. Learn the selected forms above.
% mdparagraph:P030:end

% example: NEW-U02-EXAMPLE-08
% mdparagraph:P031:begin
\textbf{Hē gehālgode ðone dæg.} He hallowed the day.
% mdparagraph:P031:end

\reserve{7\baselineskip}
% mdparagraph:P032:begin
\section*{6. Class III and the negative companion}\BookHeading{nav-ch02-learner-0006}{subsection}{6. Class III and the negative companion}{1}\BookLabel{bk-ch02-learner-79fe9bfe07cd}
% mdparagraph:P032:end

% mdparagraph:P033:begin
All four surviving Class III verbs are included. Their present stem changes and past stems form separate systems; learn each complete paradigm. \textbf{Habban} means possession here, \textbf{libban} living, \textbf{sęcgan} saying, and \textbf{hycgan} thinking. [7]
% mdparagraph:P033:end

% teaching: NEW-U02-GRAM-weak-III-habban NEW-U02-GRAM-weak-III-libban NEW-U02-GRAM-weak-III-secgan NEW-U02-GRAM-weak-III-hycgan
% teaching: NEW-U02-PAR-habban NEW-U02-PAR-libban NEW-U02-PAR-secgan NEW-U02-PAR-hycgan
% mdtable:T07:begin
% course-table-style: 2026-10
% panel:Present indicative
% panel:Preterite indicative
% course-panel:T07-P01:begin
\par\begingroup
\fontsize{10.5}{13.2}\selectfont
\setlength{\LTleft}{0pt}\setlength{\LTright}{\fill}
\noindent
\begin{tabular}{@{}>{\raggedright\arraybackslash}p{20.873mm}>{\raggedright\arraybackslash}p{21.064mm}>{\raggedright\arraybackslash}p{24.876mm}>{\raggedright\arraybackslash}p{21.064mm}>{\raggedright\arraybackslash}p{24.876mm}@{}}
\toprule
\textbf{Person and number} & \textbf{habban · Present} & \textbf{habban · Preterite} & \textbf{libban · Present} & \textbf{libban · Preterite} \\
\midrule
1 sg. & hæbbe & hæfde & libbe & lifde \\
2 sg. & hæfst & hæfdest & leofast & lifdest \\
3 sg. & hæfð & hæfde & leofað & lifde \\
1–3 pl. & habbað & hæfdon & libbað & lifdon \\
\bottomrule
\end{tabular}
\par\endgroup
% course-panel:T07-P01:end
\par\addvspace{5pt}
% course-panel:T07-P02:begin
\par\begingroup
\fontsize{10.5}{13.2}\selectfont
\setlength{\LTleft}{0pt}\setlength{\LTright}{\fill}
\noindent
\begin{tabular}{@{}>{\raggedright\arraybackslash}p{20.873mm}>{\raggedright\arraybackslash}p{21.064mm}>{\raggedright\arraybackslash}p{24.876mm}>{\raggedright\arraybackslash}p{21.064mm}>{\raggedright\arraybackslash}p{24.876mm}@{}}
\toprule
\textbf{Person and number} & \textbf{sęcgan · Present} & \textbf{sęcgan · Preterite} & \textbf{hycgan · Present} & \textbf{hycgan · Preterite} \\
\midrule
1 sg. & sęcge & sægde & hycge & hogode \\
2 sg. & sægst & sægdest & hygst & hogodest \\
3 sg. & sægð & sægde & hygð & hogode \\
1–3 pl. & sęcgað & sægdon & hycgað & hogodon \\
\bottomrule
\end{tabular}
\par\endgroup
% course-panel:T07-P02:end
% mdtable:T07:end
% mdparagraph:P034:begin
Reference alternatives include \textbf{hafast, hafað, hæbbað} for \textbf{habban}; \textbf{lifian/\allowbreak{}lifiað} beside \textbf{libban/\allowbreak{}libbað}; \textbf{sęcgean/\allowbreak{}sęcgeað}, \textbf{sęgst/\allowbreak{}sagast}, \textbf{sęgð/\allowbreak{}sagað}, and past \textbf{sǣde, sǣdest, sǣdon} beside the selected \textbf{sęcgan} forms. For \textbf{hycgan}, alternatives include \textbf{hycgean/\allowbreak{}hycgeað}, \textbf{hogast/\allowbreak{}hogað}, and past \textbf{hogde, hogdest, hogdon}. They are valid forms, not extra required paradigms.
% mdparagraph:P034:end

% teaching: NEW-U02-GRAM-nabban-reference
% mdparagraph:P035:begin
\textbf{Nabban}, “have not”, is the contracted negative companion of \textbf{habban}. Recognize its complete system; independent production is not required.
% mdparagraph:P035:end

% teaching: NEW-U02-PAR-nabban
% mdtable:T08:begin
% course-table-style: 2026-10
% panel:Present indicative
% panel:Preterite indicative
\par\addvspace{3pt}
% course-caption:T08:0:begin
\textbf{nabban}\par\nobreak
% course-caption:T08:0:end
% course-panel:T08-P01:begin
\par\begingroup
\fontsize{10.5}{13.2}\selectfont
\setlength{\LTleft}{0pt}\setlength{\LTright}{\fill}
\noindent
\begin{tabular}{@{}>{\raggedright\arraybackslash}p{39.459mm}>{\raggedright\arraybackslash}p{39.459mm}>{\raggedright\arraybackslash}p{39.459mm}@{}}
\toprule
\textbf{Person and number} & \textbf{Present} & \textbf{Preterite} \\
\midrule
1 sg. & næbbe & næfde \\
2 sg. & næfst & næfdest \\
3 sg. & næfð & næfde \\
1–3 pl. & nabbað & næfdon \\
\bottomrule
\end{tabular}
\par\endgroup
% course-panel:T08-P01:end
% mdtable:T08:end
% mdparagraph:P036:begin
\textbf{Nafast, nafað} are further source alternatives. The contraction after \textbf{ne} is frequent, not compulsory: uncontracted \textbf{ne hæbbe} is also valid. Recognition of this series does not teach negative forms of the copula.
% mdparagraph:P036:end

\reserve{7\baselineskip}
% mdparagraph:P037:begin
\section*{7. Using the forms}\BookHeading{nav-ch02-learner-0007}{subsection}{7. Using the forms}{1}\BookLabel{bk-ch02-learner-d8738489e241}
% mdparagraph:P037:end

% mdparagraph:P038:begin
The following examples are newly composed teaching material.
% mdparagraph:P038:end

% teaching: NEW-U02-GRAM-active-agreement NEW-U02-GRAM-direct-object
% mdparagraph:P039:begin
The finite verb agrees in person and number with its nominative subject. An accusative object does not control the verb. Use the object constructions specified in the vocabulary; a weak paradigm alone does not tell us which case every verb governs. [8]
% mdparagraph:P039:end

% example: NEW-U02-EXAMPLE-01
% mdparagraph:P040:begin
\textbf{Wē hīerað ðone mann.} We hear the man.
% mdparagraph:P040:end

% example: NEW-U02-EXAMPLE-02
% mdparagraph:P041:begin
\textbf{Se mann hīerð ūs.} The man hears us.
% mdparagraph:P041:end

% mdparagraph:P042:begin
In the first example, \textbf{wē} requires a plural verb. In the second, singular \textbf{se mann} requires \textbf{hīerð}, even though its object \textbf{ūs} is plural. Subject–verb–object order is a useful model here; it is not the only possible Old English order.
% mdparagraph:P042:end

% teaching: NEW-U02-GRAM-giving-recipient
% mdparagraph:P043:begin
With \textbf{sęllan}, the thing given is accusative and the recipient is dative. A personal pronoun takes the same case as the noun it replaces.
% mdparagraph:P043:end

% example: NEW-U02-EXAMPLE-03
% mdparagraph:P044:begin
\textbf{Ðū sęlst ðǣm fugle sǣd.} You give the bird seed.
% mdparagraph:P044:end

% teaching: NEW-U02-GRAM-libban-location
% mdparagraph:P045:begin
\textbf{Libban} is intransitive. A location can be expressed by the already learned \textbf{on + dative}. Possessive \textbf{his, hiere, hiera} remain unchanged before the noun.
% mdparagraph:P045:end

% example: NEW-U02-EXAMPLE-04
% mdparagraph:P046:begin
\textbf{Hiera fuglas libbað on his hofe.} Their birds live in his dwelling.
% mdparagraph:P046:end

% teaching: NEW-U02-GRAM-reflexive-use
% mdparagraph:P047:begin
Ordinary\BookPoint{bk-ch02-learner-8fc55662d2c6} personal pronouns also serve reflexively, referring to the subject. With \textbf{geręstan}, use an accusative referring to the person resting. Natural English often omits “myself”, “himself”, and so on in translating this construction. With \textbf{sęllan}, a self-recipient is dative, just like any other recipient. Reflexive meaning comes from reference and construction; there is no special reflexive form in the table. [9]
% mdparagraph:P047:end

% example: NEW-U02-EXAMPLE-05
% mdparagraph:P048:begin
\textbf{Hēo geręste hīe.} She rested herself.
% mdparagraph:P048:end

% teaching: NEW-U02-VOC-ne NEW-U02-GRAM-simple-negation
% mdparagraph:P049:begin
\textbf{Ne}\BookPoint{bk-ch02-learner-833c9271fac6} means “not”. In these simple clauses put it immediately before the finite verb. English may need “do not” or “did not”; Old English does not add a corresponding auxiliary. The exercises use noncontracting verbs for negative composition. [10]
% mdparagraph:P049:end

% example: NEW-U02-EXAMPLE-06
% mdparagraph:P050:begin
\textbf{Git ne hīerað ðone mann.} You two do not hear the man.
% mdparagraph:P050:end

% teaching: NEW-U02-GRAM-clause-coordination NEW-U02-GRAM-direct-quotation
% mdparagraph:P051:begin
\textbf{And}\BookPoint{bk-ch02-learner-4f6bce65b786} joins clauses or predicates. When two coordinated verbs have the same subject, that subject may remain understood with the second verb. Keep each verb’s object and tense clear.
% mdparagraph:P051:end

% mdparagraph:P052:begin
\textbf{Sęcgan}\BookPoint{bk-ch02-learner-693038b042d5} can introduce direct speech: the words inside quotation marks form their own independent utterance. Pronouns inside the quotation refer to that utterance’s speaker and addressees. The construction below does not introduce indirect speech. [11]
% mdparagraph:P052:end

% example: NEW-U02-EXAMPLE-07
% mdparagraph:P053:begin
\textbf{Hē sægð: “Wē habbað wæter.”} He says: “We have water.”
% mdparagraph:P053:end

\reserve{7\baselineskip}
% mdparagraph:P054:begin
\section*{8. Vocabulary and principal parts}\BookHeading{nav-ch02-learner-0008}{subsection}{8. Vocabulary and principal parts}{1}
% mdparagraph:P054:end

% mdparagraph:P055:begin
The nouns and the article are those of Unit 1. The following verb senses and constructions are used here. The four entries marked \textbf{forms only} are practised as morphology; their complement constructions are not assigned.
% mdparagraph:P055:end

% teaching: NEW-U02-VOC-hieran NEW-U02-VOC-gefyllan NEW-U02-VOC-cyssan NEW-U02-VOC-settan NEW-U02-VOC-sendan NEW-U02-VOC-laedan NEW-U02-VOC-iecan NEW-U02-VOC-gerestan NEW-U02-VOC-metan NEW-U02-VOC-gierwan NEW-U02-VOC-sellan NEW-U02-VOC-gewyrcean NEW-U02-VOC-gefremman NEW-U02-VOC-efnan NEW-U02-VOC-nerian NEW-U02-VOC-lufian NEW-U02-VOC-gebletsian NEW-U02-VOC-gehalgian NEW-U02-VOC-habban NEW-U02-VOC-libban NEW-U02-VOC-secgan NEW-U02-VOC-hycgan NEW-U02-VOC-nabban
% mdtable:T09:begin
\par\begingroup
\fontsize{10.5}{13.2}\selectfont
\setlength{\LTleft}{0pt}\setlength{\LTright}{\fill}
\begin{longtable}{@{}>{\raggedright\arraybackslash}p{25mm}>{\raggedright\arraybackslash}p{11mm}>{\raggedright\arraybackslash}p{81mm}@{}}
\toprule
\textbf{Verb} & \textbf{Class} & \textbf{Meaning and use} \\
\midrule
\endfirsthead
\toprule
\textbf{Verb} & \textbf{Class} & \textbf{Meaning and use} \\
\midrule
\endhead
\midrule
\multicolumn{3}{r}{\footnotesize Continued on the following page.}\\
\endfoot
\bottomrule
\endlastfoot
hīeran & I & hear; accusative person or thing heard \\
gefyllan & I & finish, accomplish a work; accusative work; no substance-filling construction \\
cyssan & I & kiss; accusative person \\
sęttan & I & set, put down, place; accusative thing; no new direction phrase \\
sęndan & I & send; accusative person or thing; recipient and destination constructions deferred \\
lǣdan & I & lead; accusative person; destination constructions deferred \\
īecan & I & increase; forms only \\
geręstan & I & rest, repose; accusative self-referential pronoun \\
mētan & I & meet, find; accusative person or thing; no “find out” clause \\
gierwan & I & prepare; forms only \\
sęllan & I & give; accusative thing and dative recipient \\
gewyrcean & I & make, build; accusative thing \\
gefręmman & I & perform, accomplish a work; accusative work \\
ęfnan & I & perform; forms only \\
nęrian & I & save; accusative person or thing \\
lufian & II & love; accusative person or thing \\
geblētsian & II & bless; accusative person or thing \\
gehālgian & II & hallow, consecrate, make holy; accusative person or thing \\
habban & III & have, possess; accusative thing possessed; no auxiliary perfect \\
libban & III & live; intransitive, with on + dative for location \\
sęcgan & III & say; direct quotation of an independent utterance \\
hycgan & III & think; forms only \\
nabban & III & have not; complete finite series for recognition only \\
\end{longtable}
\par\endgroup

% mdtable:T09:end
% mdparagraph:P056:begin
The prefixed verbs are separate vocabulary entries. Do not infer a new meaning by freely adding or removing \textbf{ge-}. In particular, \textbf{gefręmman} here means “perform/\allowbreak{}accomplish”, and \textbf{gefyllan} “finish/\allowbreak{}accomplish a work”. The headword variant \textbf{hȳran} belongs with \textbf{hīeran}. [12]
% mdparagraph:P056:end

% teaching: NEW-U02-GRAM-citation-boundary
% mdparagraph:P057:begin
The three principal parts are supplied below for lookup. The middle column is the learned first/\allowbreak{}third singular preterite indicative. The infinitive and participle columns identify dictionary forms; they are not independent recall targets and do not teach infinitival, passive, perfect or participial syntax. Parentheses reproduce optional material in the source citation, not a freely applicable prefix rule.
% mdparagraph:P057:end

% mdtable:T10:begin
\par\begingroup
\fontsize{10.5}{13.2}\selectfont
\setlength{\LTleft}{0pt}\setlength{\LTright}{\fill}
\begin{longtable}{@{}>{\raggedright\arraybackslash}p{37mm}>{\raggedright\arraybackslash}p{36mm}>{\raggedright\arraybackslash}p{44mm}@{}}
\toprule
\textbf{Infinitive citation} & \textbf{Preterite 1/\allowbreak{}3 sg.} & \textbf{Past participle citation} \\
\midrule
\endfirsthead
\toprule
\textbf{Infinitive citation} & \textbf{Preterite 1/\allowbreak{}3 sg.} & \textbf{Past participle citation} \\
\midrule
\endhead
\midrule
\multicolumn{3}{r}{\footnotesize Continued on the following page.}\\
\endfoot
\bottomrule
\endlastfoot
hīeran & hīerde & (ge)hīered \\
gefyllan & gefylde & gefylled \\
cyssan & cyste & (ge)cyssed \\
sęttan & sętte & (ge)sętt \\
sęndan & sęnde & (ge)sęnd(ed) \\
lǣdan & lǣdde & (ge)lǣd(ed) \\
īecan & īecte & (ge)īeced \\
geręstan & geręste & geręst \\
mētan & mētte & (ge)mēt(t) \\
gierwan & gierede & (ge)gier(w)ed \\
sęllan & sealde & (ge)seald \\
gewyrcean & geworhte & geworht \\
gefręmman & gefręmede & gefręmed \\
ęfnan & ęfnede & (ge)ęfned \\
nęrian & nęrede & (ge)nęred \\
lufian & lufode & (ge)lufod \\
geblētsian & geblētsode & geblētsod \\
gehālgian & gehālgode & gehālgod \\
habban & hæfde & (ge)hæfd \\
libban & lifde & (ge)lifd \\
sęcgan & sægde & (ge)sægd \\
hycgan & hogode & (ge)hogod \\
nabban & næfde & (ge)næfd \\
\end{longtable}
\par\endgroup

% mdtable:T10:end
% mdparagraph:P058:begin
The alternative past series of \textbf{sęcgan} and \textbf{hycgan} have the participle citations \textbf{(ge)sǣd} and \textbf{(ge)hogd}. These too are reference forms only.
% mdparagraph:P058:end

\reserve{7\baselineskip}
% mdparagraph:P059:begin
\section*{9. Pronoun practice}\BookHeading{nav-ch02-learner-0009}{subsection}{9. Pronoun practice}{1}
% mdparagraph:P059:end

% mdparagraph:P060:begin
Use the selected forms in §1. Identical spellings must still be labelled with their separate cases.
% mdparagraph:P060:end

\par\begin{minipage}{\linewidth}\raggedright
% item: NEW-U02-P01
% mdparagraph:P061:begin
\itempara{1}{Write the nominative, dative and accusative personal pronouns for the first person singular and the second person singular. Give the English meaning of each row.}
% mdparagraph:P061:end

\end{minipage}\par\addvspace{4pt}

\par\begin{minipage}{\linewidth}\raggedright
% item: NEW-U02-P02
% mdparagraph:P062:begin
\itempara{2}{Write the nominative, dative and accusative personal pronouns for “we two” and “you two”. State which number of verb form goes with each nominative.}
% mdparagraph:P062:end

\end{minipage}\par\addvspace{4pt}

\par\begin{minipage}{\linewidth}\raggedright
% item: NEW-U02-P03
% mdparagraph:P063:begin
\itempara{3}{Write the nominative, dative and accusative personal pronouns for “we” and “you” when more than two people are meant.}
% mdparagraph:P063:end

\end{minipage}\par\addvspace{4pt}

\par\begin{minipage}{\linewidth}\raggedright
% item: NEW-U02-P04
% mdparagraph:P064:begin
\itempara{4}{Write the third-person singular nominative, dative and accusative pronouns for masculine, feminine and neuter referents. Keep the three gender rows distinct.}
% mdparagraph:P064:end

\end{minipage}\par\addvspace{4pt}

\par\begin{minipage}{\linewidth}\raggedright
% item: NEW-U02-P05
% mdparagraph:P065:begin
\itempara{5}{Write the third-person plural nominative, dative and accusative personal pronouns. Are separate forms needed for the three genders?}
% mdparagraph:P065:end

\end{minipage}\par\addvspace{4pt}

\par\begin{minipage}{\linewidth}\raggedright
% item: NEW-U02-P06
% mdparagraph:P066:begin
\itempara{6}{Match these genitive forms to person and number: mīn, ðīn, uncer, incer, ūre, ēower. Use “first” or “second” person and “singular”, “dual” or “plural”.}
% mdparagraph:P066:end

\end{minipage}\par\addvspace{4pt}

\reserve{7\baselineskip}
% mdparagraph:P067:begin
\section*{10. Verb forms}\BookHeading{nav-ch02-learner-0010}{subsection}{10. Verb forms}{1}
% mdparagraph:P067:end

% mdparagraph:P068:begin
Except for the given negative forms in exercise 19 and the lookup exercise 20, produce the requested forms from memory. When a complete paradigm is requested, label all three singular and all three plural persons in both tenses; identical plural forms may be grouped.
% mdparagraph:P068:end

\par\begin{minipage}{\linewidth}\raggedright
% item: NEW-U02-F01
% mdparagraph:P069:begin
\itempara{1}{Write the complete present and preterite indicative of hīeran. Group identical plural forms if you label all three persons. Then give “you two heard”.}
% mdparagraph:P069:end

\end{minipage}\par\addvspace{4pt}

\par\begin{minipage}{\linewidth}\raggedright
% item: NEW-U02-F02
% mdparagraph:P070:begin
\itempara{2}{Give “you (singular) finish”, “he finishes”, “I finished” and “we finished”, using gefyllan in the sense “finish; accomplish”. Explain the difference between the present and preterite stem spellings.}
% mdparagraph:P070:end

\end{minipage}\par\addvspace{4pt}

\par\begin{minipage}{\linewidth}\raggedright
% item: NEW-U02-F03
% mdparagraph:P071:begin
\itempara{3}{Give the second and third singular present indicative and the second singular and plural preterite indicative of cyssan.}
% mdparagraph:P071:end

\end{minipage}\par\addvspace{4pt}

\par\begin{minipage}{\linewidth}\raggedright
% item: NEW-U02-F04
% mdparagraph:P072:begin
\itempara{4}{For sęttan and sęndan, give the first singular present, first singular preterite and plural preterite. Label the tense of each form.}
% mdparagraph:P072:end

\end{minipage}\par\addvspace{4pt}

\par\begin{minipage}{\linewidth}\raggedright
% item: NEW-U02-F05
% mdparagraph:P073:begin
\itempara{5}{For lǣdan and mētan, give the second singular present, third singular present and second singular preterite. Keep every printed length mark.}
% mdparagraph:P073:end

\end{minipage}\par\addvspace{4pt}

\par\begin{minipage}{\linewidth}\raggedright
% item: NEW-U02-F06
% mdparagraph:P074:begin
\itempara{6}{For īecan and gierwan, give the first singular present, third singular present and first singular preterite. This is a form exercise, not a sentence translation.}
% mdparagraph:P074:end

\end{minipage}\par\addvspace{4pt}

\par\begin{minipage}{\linewidth}\raggedright
% item: NEW-U02-F07
% mdparagraph:P075:begin
\itempara{7}{Write the complete present and preterite indicative of geręstan. Then give “he rested himself” using its ordinary reflexive construction.}
% mdparagraph:P075:end

\end{minipage}\par\addvspace{4pt}

\par\begin{minipage}{\linewidth}\raggedright
% item: NEW-U02-F08
% mdparagraph:P076:begin
\itempara{8}{Write the complete present and preterite indicative of sęllan. Then give “I give” and “I gave” with the personal pronoun.}
% mdparagraph:P076:end

\end{minipage}\par\addvspace{4pt}

\par\begin{minipage}{\linewidth}\raggedright
% item: NEW-U02-F09
% mdparagraph:P077:begin
\itempara{9}{Give the second and third singular present indicative and all persons of the singular preterite indicative of gewyrcean.}
% mdparagraph:P077:end

\end{minipage}\par\addvspace{4pt}

\par\begin{minipage}{\linewidth}\raggedright
% item: NEW-U02-F10
% mdparagraph:P078:begin
\itempara{10}{For gefręmman, give the first, second and third singular present, the plural present, and the first singular preterite.}
% mdparagraph:P078:end

\end{minipage}\par\addvspace{4pt}

\par\begin{minipage}{\linewidth}\raggedright
% item: NEW-U02-F11
% mdparagraph:P079:begin
\itempara{11}{For ęfnan, give the first singular present, second singular present, first singular preterite and plural preterite.}
% mdparagraph:P079:end

\end{minipage}\par\addvspace{4pt}

\par\begin{minipage}{\linewidth}\raggedright
% item: NEW-U02-F12
% mdparagraph:P080:begin
\itempara{12}{Write the complete present and preterite indicative of nęrian. Then give “we two save” and “you (singular) saved”.}
% mdparagraph:P080:end

\end{minipage}\par\addvspace{4pt}

\par\begin{minipage}{\linewidth}\raggedright
% item: NEW-U02-F13
% mdparagraph:P081:begin
\itempara{13}{Write the complete present and preterite indicative of lufian. Identify the vowel in the present second singular ending and in the preterite stem.}
% mdparagraph:P081:end

\end{minipage}\par\addvspace{4pt}

\par\begin{minipage}{\linewidth}\raggedright
% item: NEW-U02-F14
% mdparagraph:P082:begin
\itempara{14}{For geblētsian and gehālgian, give the third singular present, second singular preterite and plural preterite.}
% mdparagraph:P082:end

\end{minipage}\par\addvspace{4pt}

\par\begin{minipage}{\linewidth}\raggedright
% item: NEW-U02-F15
% mdparagraph:P083:begin
\itempara{15}{Write the complete present and preterite indicative of habban. Include the personal pronouns in the three singular rows.}
% mdparagraph:P083:end

\end{minipage}\par\addvspace{4pt}

\par\begin{minipage}{\linewidth}\raggedright
% item: NEW-U02-F16
% mdparagraph:P084:begin
\itempara{16}{Write the complete present and preterite indicative of libban.}
% mdparagraph:P084:end

\end{minipage}\par\addvspace{4pt}

\par\begin{minipage}{\linewidth}\raggedright
% item: NEW-U02-F17
% mdparagraph:P085:begin
\itempara{17}{Write the complete present and preterite indicative of sęcgan.}
% mdparagraph:P085:end

\end{minipage}\par\addvspace{4pt}

\par\begin{minipage}{\linewidth}\raggedright
% item: NEW-U02-F18
% mdparagraph:P086:begin
\itempara{18}{Write the complete present and preterite indicative of hycgan.}
% mdparagraph:P086:end

\end{minipage}\par\addvspace{4pt}

\par\begin{minipage}{\linewidth}\raggedright
% item: NEW-U02-F19
% mdparagraph:P087:begin
\itempara{19}{For each given form, identify person, number and tense, allowing more than one answer where necessary: næbbe, næfst, næfð, nabbað, næfde, næfdest, næfdon. Give the negative meaning.}
% mdparagraph:P087:end

\end{minipage}\par\addvspace{4pt}

\par\begin{minipage}{\linewidth}\raggedright
% item: NEW-U02-F20
% mdparagraph:P088:begin
\itempara{20}{Use the principal-part entries to identify the lemma and tense of hīerde, sealde, geworhte, gefręmede, lufode and lifde. Which of these forms can be first or third singular in the indicative?}
% mdparagraph:P088:end

\end{minipage}\par\addvspace{4pt}

\reserve{7\baselineskip}
% mdparagraph:P089:begin
\section*{11. Cases, reference and clauses}\BookHeading{nav-ch02-learner-0011}{subsection}{11. Cases, reference and clauses}{1}
% mdparagraph:P089:end

% mdparagraph:P090:begin
Use the meanings and constructions in §8. Give complete sentences where a translation is requested.
% mdparagraph:P090:end

\par\begin{minipage}{\linewidth}\raggedright
% item: NEW-U02-C01
% mdparagraph:P091:begin
\itempara{1}{The referents are a man and several birds. Translate Hē hīerð hīe and Hīe hīerað hine. Identify the case and referent number of every pronoun, and explain how the two statements differ.}
% mdparagraph:P091:end

\end{minipage}\par\addvspace{4pt}

\par\begin{minipage}{\linewidth}\raggedright
% item: NEW-U02-C02
% mdparagraph:P092:begin
\itempara{2}{Translate “I give the man seed”. Then replace only “the man” with the appropriate personal pronoun, keeping the same meaning.}
% mdparagraph:P092:end

\end{minipage}\par\addvspace{4pt}

\par\begin{minipage}{\linewidth}\raggedright
% item: NEW-U02-C03
% mdparagraph:P093:begin
\itempara{3}{Translate Wit habbað męte. Se mann sęlð unc wæter. Identify the case and number of wit and unc.}
% mdparagraph:P093:end

\end{minipage}\par\addvspace{4pt}

\par\begin{minipage}{\linewidth}\raggedright
% item: NEW-U02-C04
% mdparagraph:P094:begin
\itempara{4}{Translate “You two save me. I hear you two.” Use the dual personal pronouns where appropriate.}
% mdparagraph:P094:end

\end{minipage}\par\addvspace{4pt}

\par\begin{minipage}{\linewidth}\raggedright
% item: NEW-U02-C05
% mdparagraph:P095:begin
\itempara{5}{Replace sēo eorðe in the second sentence with the appropriate personal pronoun, and translate both sentences: Sēo eorðe hæfð wæter. Sēo eorðe hæfð gærs.}
% mdparagraph:P095:end

\end{minipage}\par\addvspace{4pt}

\par\begin{minipage}{\linewidth}\raggedright
% item: NEW-U02-C06
% mdparagraph:P096:begin
\itempara{6}{Translate “His dwelling is on the earth. His birds live in the dwelling.” The possessor is a man in both sentences.}
% mdparagraph:P096:end

\end{minipage}\par\addvspace{4pt}

\par\begin{minipage}{\linewidth}\raggedright
% item: NEW-U02-C07
% mdparagraph:P097:begin
\itempara{7}{Translate “She has a dwelling. Birds live in her dwelling.”}
% mdparagraph:P097:end

\end{minipage}\par\addvspace{4pt}

\par\begin{minipage}{\linewidth}\raggedright
% item: NEW-U02-C08
% mdparagraph:P098:begin
\itempara{8}{Translate “They have a dwelling. Their dwelling has birds.”}
% mdparagraph:P098:end

\end{minipage}\par\addvspace{4pt}

\par\begin{minipage}{\linewidth}\raggedright
% item: NEW-U02-C09
% mdparagraph:P099:begin
\itempara{9}{Translate “I rested myself. He rested himself.” Mark the subject and object in each clause.}
% mdparagraph:P099:end

\end{minipage}\par\addvspace{4pt}

\par\begin{minipage}{\linewidth}\raggedright
% item: NEW-U02-C10
% mdparagraph:P100:begin
\itempara{10}{Translate “You two rest yourselves. They rest themselves.”}
% mdparagraph:P100:end

\end{minipage}\par\addvspace{4pt}

\par\begin{minipage}{\linewidth}\raggedright
% item: NEW-U02-C11
% mdparagraph:P101:begin
\itempara{11}{Change Ic hīere ðone fugol and Hē lufað ðā fuglas into negative statements. Translate the results.}
% mdparagraph:P101:end

\end{minipage}\par\addvspace{4pt}

\par\begin{minipage}{\linewidth}\raggedright
% item: NEW-U02-C12
% mdparagraph:P102:begin
\itempara{12}{Translate “You (singular) did not hear him. We did not send the birds.”}
% mdparagraph:P102:end

\end{minipage}\par\addvspace{4pt}

\par\begin{minipage}{\linewidth}\raggedright
% item: NEW-U02-C13
% mdparagraph:P103:begin
\itempara{13}{Translate “I give myself food. I give him food.” Explain which pronoun is reflexive and why both are dative.}
% mdparagraph:P103:end

\end{minipage}\par\addvspace{4pt}

\reserve{7\baselineskip}
% mdparagraph:P104:begin
\section*{12. Translation}\BookHeading{nav-ch02-learner-0012}{subsection}{12. Translation}{1}
% mdparagraph:P104:end

\par\begin{minipage}{\linewidth}\raggedright
% item: NEW-U02-T01
% mdparagraph:P105:begin
\itempara{1}{Translate “The man put down the seed. The men sent it.”}
% mdparagraph:P105:end

\end{minipage}\par\addvspace{4pt}

\par\begin{minipage}{\linewidth}\raggedright
% item: NEW-U02-T02
% mdparagraph:P106:begin
\itempara{2}{Translate “I met you (singular). You met me. We two rested ourselves.”}
% mdparagraph:P106:end

\end{minipage}\par\addvspace{4pt}

\par\begin{minipage}{\linewidth}\raggedright
% item: NEW-U02-T03
% mdparagraph:P107:begin
\itempara{3}{Translate “He saved the man and loved him. He did not send him.”}
% mdparagraph:P107:end

\end{minipage}\par\addvspace{4pt}

\par\begin{minipage}{\linewidth}\raggedright
% item: NEW-U02-T04
% mdparagraph:P108:begin
\itempara{4}{Translate “You (singular) blessed the day. They hallowed it.”}
% mdparagraph:P108:end

\end{minipage}\par\addvspace{4pt}

\par\begin{minipage}{\linewidth}\raggedright
% item: NEW-U02-T05
% mdparagraph:P109:begin
\itempara{5}{Translate Hē geworhte his hof. Hīe gefręmedon hiera weorc. Hē geręste hine on ðǣm hofe.}
% mdparagraph:P109:end

\end{minipage}\par\addvspace{4pt}

\par\begin{minipage}{\linewidth}\raggedright
% item: NEW-U02-T06
% mdparagraph:P110:begin
\itempara{6}{Translate “We have the man’s seed. We give him water. He hears us.”}
% mdparagraph:P110:end

\end{minipage}\par\addvspace{4pt}

\reserve{7\baselineskip}
% mdparagraph:P111:begin
\section*{13. Two connected readings}\BookHeading{nav-ch02-learner-0013}{subsection}{13. Two connected readings}{1}
% mdparagraph:P111:end

% mdparagraph:P112:begin
Both readings are new compositions using the forms and constructions of Units 1–2.
% mdparagraph:P112:end

\reserve{7\baselineskip}
% mdparagraph:P113:begin
\subsection*{Seed and water}\BookHeading{nav-ch02-learner-0014}{subsubsection}{Seed and water}{0}
% mdparagraph:P113:end

% composition: NEW-U02-TEXT-NARRATIVE-01
% mdparagraph:P114:begin
\par\begin{minipage}{\linewidth}\raggedright
\readingtext{Se mann hæfde sǣd and wæter. Ðā męnn hæfdon hof and fuglas. Ðā fuglas lifdon on hiera hofe. Se mann sealde ðǣm mannum sǣd and wæter. Hīe lufodon ðā fuglas and hīerdon hīe. Hīe sealdon him sǣd and wæter. Hīe ne sęndon ðā fuglas. Hīe gefręmedon hiera weorc and geręston hīe. Se mann sægde: “Ic lufie ðā fuglas.”}
\end{minipage}\par
% mdparagraph:P114:end

% mdparagraph:P115:begin
Plural pronoun spellings do not always identify a unique antecedent. Follow the sequence of giving and caring for the birds, and be prepared to explain any other defensible reference.
% mdparagraph:P115:end

\reserve{7\baselineskip}
% mdparagraph:P116:begin
\subsection*{The workers and the birds}\BookHeading{nav-ch02-learner-0015}{subsubsection}{The workers and the birds}{0}
% mdparagraph:P116:end

% mdparagraph:P117:begin
Four workers, A, B, C and D, are together with the birds. Speaker labels identify the addressees. In the last turn, wē includes all four workers and gē addresses B, C and D.
% mdparagraph:P117:end

% composition: NEW-U02-TEXT-DIALOGUE-01
% mdparagraph:P118:begin
\textbf{A to B:} Ic hæbbe sǣd. Ðū hæfst wæter.
% mdparagraph:P118:end

% mdparagraph:P119:begin
\textbf{B to A:} Wit sęllað ðǣm fuglum męte. Hīe hīerað unc.
% mdparagraph:P119:end

% mdparagraph:P120:begin
\textbf{C to A, B:} Git lufiað ðā fuglas. Ic hīere inc.
% mdparagraph:P120:end

% mdparagraph:P121:begin
\textbf{A to B, C, D:} Wē lufiað hīe. Gē hīerað hīe, and hīe hīerað ēow.
% mdparagraph:P121:end

\reserve{7\baselineskip}
% mdparagraph:P122:begin
\subsection*{Reading exercises}\BookHeading{nav-ch02-learner-0016}{subsubsection}{Reading exercises}{0}
% mdparagraph:P122:end

\par\begin{minipage}{\linewidth}\raggedright
% item: NEW-U02-N01
% mdparagraph:P123:begin
\itempara{1}{Translate the connected account “Seed and water” in full. Keep the sequence of giver, recipients and things given clear.}
% mdparagraph:P123:end

\end{minipage}\par\addvspace{4pt}

\par\begin{minipage}{\linewidth}\raggedright
% item: NEW-U02-N02
% mdparagraph:P124:begin
\itempara{2}{In the account, trace the people and birds referred to by each plural personal pronoun. Which references are settled by the form itself, and which depend on the surrounding actions?}
% mdparagraph:P124:end

\end{minipage}\par\addvspace{4pt}

\par\begin{minipage}{\linewidth}\raggedright
% item: NEW-U02-N03
% mdparagraph:P125:begin
\itempara{3}{Translate the connected conversation in full. In each turn, state who is included in ic, ðū, wit, git, wē and gē.}
% mdparagraph:P125:end

\end{minipage}\par\addvspace{4pt}

\par\begin{minipage}{\linewidth}\raggedright
% item: NEW-U02-N04
% mdparagraph:P126:begin
\itempara{4}{In the conversation, identify the case, number and referents of unc, inc and ēow. In each clause containing one of these objects, identify the subject and explain which pronoun controls the finite verb.}
% mdparagraph:P126:end

\end{minipage}\par\addvspace{4pt}

\par\begin{minipage}{\linewidth}\raggedright
% item: NEW-U02-N05
% mdparagraph:P127:begin
\itempara{5}{Retell the first four sentences of “Seed and water” as present actions or states. Preserve the people, things and cases.}
% mdparagraph:P127:end

\end{minipage}\par\addvspace{4pt}

\par\begin{minipage}{\linewidth}\raggedright
% item: NEW-U02-N06
% mdparagraph:P128:begin
\itempara{6}{In Hīe gefręmedon hiera weorc and geręston hīe, identify the subject, the object of each verb, and the referent of hiera. Translate the whole sentence.}
% mdparagraph:P128:end

\end{minipage}\par\addvspace{4pt}

\reserve{7\baselineskip}
% mdparagraph:P129:begin
\section*{14. Source reading: the Sabbath}\BookHeading{nav-ch02-learner-0017}{subsection}{14. Source reading: the Sabbath}{1}
% mdparagraph:P129:end

% mdparagraph:P130:begin
This is the complete closing paragraph of Cook’s “The Creation of the World”, printed p.128 /\allowbreak{} PDF p.144. Its seven prose lines are editorially counted; only the marginal number \textbf{5} is printed. The linguistic text is unchanged; note cues and line breaks are removed. Use the local help freely. The six weak verbs are now learned, but the passive, adjective, temporal, relative, causal and infinitival constructions below are supplied for this passage. [13]
% mdparagraph:P130:end

% source: FB-I-CREATION-SABBATH-01
% mdparagraph:P131:begin
\par\begin{minipage}{\linewidth}\raggedright
\readingtext{Eornostlice ðā wǣron fullfręmede heofonas and eorðe and eall hiera frætwung. And God ðā gefylde on ðone seofoðan dæg his weorc ðe hē geworhte, and hē geręste hine on ðone seofoðan dæg fram eallum ðǣm weorce ðe hē gefręmede. And God geblētsode ðone seofoðan dæg and hine gehālgode, for-ðon-ðe hē on ðone dæg geswāc his weorces ðe hē gescēop tō wyrceanne.}
\end{minipage}\par
% mdparagraph:P131:end

\reserve{7\baselineskip}
% mdparagraph:P132:begin
\subsection*{Local help}\BookHeading{nav-ch02-learner-0018}{subsubsection}{Local help}{0}
% mdparagraph:P132:end

% teaching: NEW-U02-SOURCE-SUPPORT NEW-U02-GRAM-reading-sabbath-adverbials NEW-U02-GRAM-reading-sabbath-passive
% mdparagraph:P133:begin
\textbf{Eornostlice}, “accordingly, thus”; \textbf{ðā}, “then”, an adverb here rather than the identically spelled article.
% mdparagraph:P133:end

% mdparagraph:P134:begin
\textbf{Wǣron fullfręmede}, “were completed/\allowbreak{}finished”. This is a passive expression. \textbf{Wǣron} is familiar, but its use with the plural participle \textbf{fullfręmede} is supplied. \textbf{Fullfręmede} is not another finite past tense merely because it ends in \textbf{-ede}. The coordinated subject follows the verb: the heavens, the earth, and all their array.
% mdparagraph:P134:end

% teaching: NEW-U02-GRAM-reading-sabbath-adjective-eall NEW-U02-GRAM-reading-sabbath-ordinal NEW-U02-GRAM-reading-sabbath-temporal-on-accusative
% mdparagraph:P135:begin
\textbf{Eall hiera frætwung}, “all their array”: \textbf{eall} modifies feminine nominative singular \textbf{frætwung}. In \textbf{eallum ðǣm weorce}, “all the work”, \textbf{eallum} and \textbf{ðǣm} agree with neuter dative singular \textbf{weorce}. The ending \textbf{-um} does not make this noun phrase plural.
% mdparagraph:P135:end

% mdparagraph:P136:begin
\textbf{Ðone seofoðan dæg}, “the seventh day”: masculine accusative singular, with the supplied weak ordinal form \textbf{seofoðan}. \textbf{On ðone seofoðan dæg} means “on the seventh day”; \textbf{on ðone dæg} means “on that day”, the day just named. These are temporal \textbf{accusative} phrases. They express neither location with the dative nor motion; do not change \textbf{dæg} to \textbf{dæge} by extending the earlier locative rule. [14]
% mdparagraph:P136:end

% teaching: NEW-U02-GRAM-reading-sabbath-relative-object
% mdparagraph:P137:begin
The three relative clauses modify the preceding noun for “work”:
% mdparagraph:P137:end

\begin{itemize}
% mdparagraph:P138:begin
\item \textbf{ðe hē geworhte}, “which he made”, modifies \textbf{weorc}.
% mdparagraph:P138:end
% mdparagraph:P139:begin
\item \textbf{ðe hē gefręmede}, “which he performed”, modifies \textbf{weorce}.
% mdparagraph:P139:end
% mdparagraph:P140:begin
\item \textbf{ðe hē gescēop tō wyrceanne}, “which he created to make”, modifies \textbf{weorces}.
% mdparagraph:P140:end
\end{itemize}

% mdparagraph:P141:begin
In each clause, \textbf{ðe} represents the direct object. The preceding noun has its own case in the larger sentence: accusative, dative and genitive respectively. Relative \textbf{ðe} has no macron; it is not personal \textbf{ðē}. English “had made/\allowbreak{}had performed” may convey the earlier past action in context. General relative syntax is not required. [15]
% mdparagraph:P141:end

% teaching: NEW-U02-GRAM-reading-sabbath-rest-fram
% mdparagraph:P142:begin
\textbf{Geręste hine … fram eallum ðǣm weorce}, “rested … from all the work”. \textbf{Fram} takes dative \textbf{weorce}. This particular rest-from-work construction is supplied; the reflexive pronoun and finite weak verb are learned.
% mdparagraph:P142:end

% teaching: NEW-U02-GRAM-reading-sabbath-causal NEW-U02-GRAM-reading-sabbath-cessation-genitive NEW-U02-GRAM-reading-sabbath-created-to-make
% mdparagraph:P143:begin
\textbf{For-ðon-ðe} introduces the reason, “because”. The causal clause continues to the end of the paragraph.
% mdparagraph:P143:end

% mdparagraph:P144:begin
\textbf{Geswāc his weorces}, “ceased from his work”: \textbf{geswāc} is a supplied strong past form, third singular; the verb here governs genitive \textbf{weorces}. Recognizing the noun’s genitive does not by itself explain this government. [16]
% mdparagraph:P144:end

% mdparagraph:P145:begin
\textbf{Gescēop}, “created”, is another supplied third-singular strong past form. The whole ending \textbf{ðe hē gescēop tō wyrceanne} is “which he created to make”. \textbf{Tō wyrceanne} is an inflected infinitive phrase; it is not a second finite verb. The understood maker is God. Cook explicitly retains this “created to make” relation: preserve it rather than substitute a familiar modern version of the Creation story. A close English rendering may be awkward; no single polished paraphrase is required. [17]
% mdparagraph:P145:end

\reserve{7\baselineskip}
% mdparagraph:P146:begin
\subsection*{Reading exercises}\BookHeading{nav-ch02-learner-0019}{subsubsection}{Reading exercises}{0}
% mdparagraph:P146:end

\par\begin{minipage}{\linewidth}\raggedright
% item: NEW-U02-R01
% mdparagraph:P147:begin
\itempara{1}{Translate the complete Sabbath passage using its local support. Preserve each action, referent and clause link.}
% mdparagraph:P147:end

\end{minipage}\par\addvspace{4pt}

\par\begin{minipage}{\linewidth}\raggedright
% item: NEW-U02-R02
% mdparagraph:P148:begin
\itempara{2}{Using the vocabulary and principal-part tables for headword lookup, find the six weak preterite indicative forms in the passage. Give their headwords, class, person and number.}
% mdparagraph:P148:end

\end{minipage}\par\addvspace{4pt}

\par\begin{minipage}{\linewidth}\raggedright
% item: NEW-U02-R03
% mdparagraph:P149:begin
\itempara{3}{Compare the two occurrences of hine: geręste hine and hine gehālgode. For each, identify its antecedent and case, translate it naturally, and state whether its use is reflexive.}
% mdparagraph:P149:end

\end{minipage}\par\addvspace{4pt}

\par\begin{minipage}{\linewidth}\raggedright
% item: NEW-U02-R04
% mdparagraph:P150:begin
\itempara{4}{In hiera frætwung, his weorc and his weorces, identify the possessor and the possessed noun. Explain why the pronoun spelling is unchanged when the head’s number or case changes. Use the supplied construction note for weorces.}
% mdparagraph:P150:end

\end{minipage}\par\addvspace{4pt}

\par\begin{minipage}{\linewidth}\raggedright
% item: NEW-U02-R05
% mdparagraph:P151:begin
\itempara{5}{Using the supplied opening construction, identify the case and number of heofonas, eorðe and frætwung. Which complete phrase supplies the subject of wǣron?}
% mdparagraph:P151:end

\end{minipage}\par\addvspace{4pt}

\par\begin{minipage}{\linewidth}\raggedright
% item: NEW-U02-R06
% mdparagraph:P152:begin
\itempara{6}{List in English the actions assigned to God on the seventh day, using the passage’s temporal support.}
% mdparagraph:P152:end

\end{minipage}\par\addvspace{4pt}

\par\begin{minipage}{\linewidth}\raggedright
% item: NEW-U02-R07
% mdparagraph:P153:begin
\itempara{7}{Using the local clause support, explain the reason the passage gives for hallowing the day and how the final creation phrase describes the work.}
% mdparagraph:P153:end

\end{minipage}\par\addvspace{4pt}

\reserve{7\baselineskip}
% mdparagraph:P154:begin
\section*{15. Review}\BookHeading{nav-ch02-learner-0020}{subsection}{15. Review}{1}
% mdparagraph:P154:end

\par\begin{minipage}{\linewidth}\raggedright
% item: NEW-U02-X01
% mdparagraph:P155:begin
\itempara{1}{Change this passage to the preterite, then make only its final clause negative: Wit habbað sǣd. Wē sęllað ðǣm fuglum męte. Hīe hīerað ūs.}
% mdparagraph:P155:end

\end{minipage}\par\addvspace{4pt}

\par\begin{minipage}{\linewidth}\raggedright
% item: NEW-U02-X02
% mdparagraph:P156:begin
\itempara{2}{Correct the four forms that prevent this passage from meaning “He gives me seed. We two have it. The birds hear us two”: Hine sęlð ic sǣd. Wit hæfst hit. Ðā fuglas hīerað ūre.}
% mdparagraph:P156:end

\end{minipage}\par\addvspace{4pt}

\par\begin{minipage}{\linewidth}\raggedright
% item: NEW-U02-X03
% mdparagraph:P157:begin
\itempara{3}{Write three connected Old English sentences: say that the man had birds, that he gave them seed, and that he rested himself. Use a personal pronoun after introducing the man.}
% mdparagraph:P157:end

\end{minipage}\par\addvspace{4pt}

\clearpage
% mdparagraph:P158:begin
\section*{Source notes}\BookHeading{nav-ch02-learner-0021}{subsection}{Source notes}{1}
% mdparagraph:P158:end

\begingroup\fontsize{10}{12.8}\selectfont
% mdparagraph:P159:begin
[1] Cook, \emph{A First Book in Old English}, second edition (1895), §§81–82, pp.48–49 /\allowbreak{} PDF64–65, supplies the personal forms. Ringe and Taylor, \emph{The Development of Old English} (2014), §1.1.2, p.6 /\allowbreak{} publisher-preview PDF21, explicitly confirms plural verbs with dual subjects. Cook’s spelling is retained.
% mdparagraph:P159:end

% mdparagraph:P160:begin
[2] Cook §83, p.50 /\allowbreak{} PDF66: invariant third-person genitives and the separate inflecting possessive adjectives.
% mdparagraph:P160:end

% mdparagraph:P161:begin
[3] Cook §96, p.56 /\allowbreak{} PDF72, and the complete weak paradigms, §§117–124, pp.72–76 /\allowbreak{} PDF88–92. Some complete paradigms here combine a checked lexical stem with an explicit source pattern; they do not claim that every displayed cell occurs in the reading.
% mdparagraph:P161:end

% mdparagraph:P162:begin
[4] Cook §113, pp.67–69 /\allowbreak{} PDF83–85, with §§23,33–36, pp.18–19,24–25 /\allowbreak{} PDF34–35,40–41. Independent checks: Wright and Wright, \emph{Old English Grammar}, second edition (1914), §§305,476,524,528–529,533, pp.151,248,267,269–270,273 /\allowbreak{} PDF171,268,287,289–290,293. These corroborate the contractions, inherited consonants and selected \textbf{gierwan} forms; no alternative dialectal spelling is silently imported.
% mdparagraph:P162:end

% mdparagraph:P163:begin
[5] Cook §§114–117, pp.69–72 /\allowbreak{} PDF85–88; §18, p.14 /\allowbreak{} PDF30, for optional \textbf{e} after \textbf{c/\allowbreak{}g} before a further vowel. Wright and Wright §532, p.273 /\allowbreak{} PDF293, confirms the \textbf{ęfnan} type and both past patterns. The short-stem exception and impersonal limit follow Cook §§115,171,190.
% mdparagraph:P163:end

% mdparagraph:P164:begin
[6] Cook §§118–119, pp.73–74 /\allowbreak{} PDF89–90.
% mdparagraph:P164:end

% mdparagraph:P165:begin
[7] Cook §§120–124, pp.74–76 /\allowbreak{} PDF90–92; §29, p.23 /\allowbreak{} PDF39, on frequent negative contraction.
% mdparagraph:P165:end

% mdparagraph:P166:begin
[8] Cook §§168,189–191, pp.96,102 /\allowbreak{} PDF112,118; giving recipients §§160,164(a), pp.93–94 /\allowbreak{} PDF109–110. Lexical constructions remain limited to the senses stated in §8.
% mdparagraph:P166:end

% mdparagraph:P167:begin
[9] Cook §§82,184(a,b), pp.49,100–101 /\allowbreak{} PDF65,116–117, on personal pronouns used reflexively.
% mdparagraph:P167:end

% mdparagraph:P168:begin
[10] Cook’s vocabulary, p.306 /\allowbreak{} PDF322, and \emph{The Day of Judgment}, p.136 /\allowbreak{} PDF152, printed line 6, for \textbf{ne} before the finite verb.
% mdparagraph:P168:end

% mdparagraph:P169:begin
[11] Cook, \emph{Trades and Occupations}, p.132 /\allowbreak{} PDF148, speaker headings before printed lines 2 and 13, for \textbf{sęcgan} introducing a direct utterance. The examples and both connected practice readings in this unit are new teaching material, not quotations from Cook.
% mdparagraph:P169:end

% mdparagraph:P170:begin
[12] Cook’s vocabulary, especially pp.290,292,294,296,299,305–306,312 /\allowbreak{} PDF306,308,310,312,315,321–322,328, checks the assigned lexical senses. Dictionary headword hyphens in prefixed verbs are joined for presentation.
% mdparagraph:P170:end

% mdparagraph:P171:begin
[13] Cook, Reader I, p.128 /\allowbreak{} PDF144, entire final paragraph, prose lines 1–7; only line marker 5 is printed. Source apparatus cues 1–13 are omitted. The original quantities, \textbf{ę}, punctuation and \textbf{for-ðon-ðe} are retained. For the opening passive see also §187, p.101 /\allowbreak{} PDF117.
% mdparagraph:P171:end

% mdparagraph:P172:begin
[14] Cook p.128 notes 3,5 and §§80,172.1, pp.48,97 /\allowbreak{} PDF64,113. The temporal accusative is explicitly noted by the source.
% mdparagraph:P172:end

% mdparagraph:P173:begin
[15] Cook §87(d), p.52 /\allowbreak{} PDF68, and §188, p.101 /\allowbreak{} PDF117.
% mdparagraph:P173:end

% mdparagraph:P174:begin
[16] Cook p.128 note 11 and §156(k), p.92 /\allowbreak{} PDF108, explicitly give the genitive after “cease”.
% mdparagraph:P174:end

% mdparagraph:P175:begin
[17] Cook p.128 note 12 gives Latin \textbf{creavit ut faceret} and “created to make”; §200, p.106 /\allowbreak{} PDF122, treats the inflected infinitive, called “gerund” there. The supplied source wording is \textbf{tō wyrceanne}, including its \textbf{e}.
% mdparagraph:P175:end

\endgroup

% book-source-body:end
\par\label{end-ch02-learner}
\clearpage\endgroup

\BookRecto
\begingroup
\renewcommand{\runningtitle}{Unit 3. Adjectives and number}
\hyphenpenalty=10000
\exhyphenpenalty=10000
\emergencystretch=3em
\newcommand{\itempara}[2]{\par\addvspace{3pt}\noindent\hangindent=1.8em\hangafter=1\makebox[1.8em][l]{#1.}#2\par}
\newcommand{\reserve}[1]{\par\penalty-100\begingroup\dimen0=\pagegoal\advance\dimen0 by-\pagetotal\ifdim\dimen0<#1\newpage\fi\endgroup}
% book-source-body:begin

\raggedright
\booktitle{\BookComponent{ch03-learner}{Unit 3. Adjectives, comparison and number}UNIT 3}{Adjectives, comparison and number}

% mdparagraph:P001:begin
This unit adds description, possession, comparison and counting to the noun and indicative-verb systems of Units 1–2. The grammar and vocabulary below supply the new material needed for independent practice. The final source passage has additional forms and constructions; its local help remains available throughout the reading tasks.
% mdparagraph:P001:end

% mdparagraph:P002:begin
In tables, \textbf{m.}, \textbf{n.} and \textbf{f.} mean masculine, neuter and feminine; \textbf{sg.} and \textbf{pl.} mean singular and plural. \textbf{N/\allowbreak{}G/\allowbreak{}D/\allowbreak{}A} mean nominative, genitive, dative and accusative. Use the selected forms when a task requests them. A documented alternative may be valid Old English even when it does not meet that particular instruction. Length marks and the distinction between e and ę are retained.
% mdparagraph:P002:end

\reserve{7\baselineskip}
% mdparagraph:P003:begin
\section*{1. Descriptive vocabulary}\BookHeading{nav-ch03-learner-0001}{subsection}{1. Descriptive vocabulary}{1}
% mdparagraph:P003:end

% teaching: NEW-U03-VOC-god-adjective
% teaching: NEW-U03-VOC-glaed
% teaching: NEW-U03-VOC-micel
% teaching: NEW-U03-VOC-idel
% teaching: NEW-U03-VOC-aemtig
% teaching: NEW-U03-VOC-aeppelbaere
% teaching: NEW-U03-VOC-lytel
% teaching: NEW-U03-VOC-styriendlic
% teaching: NEW-U03-VOC-leof
% teaching: NEW-U03-VOC-eald
% teaching: NEW-U03-VOC-yfel
% teaching: NEW-U03-VOC-gearu
% teaching: NEW-U03-VOC-heah
% mdtable:T01:begin
\par\begingroup
\fontsize{10.5}{13.2}\selectfont
\setlength{\LTleft}{0pt}\setlength{\LTright}{\fill}
\noindent
\begin{tabular}{@{}>{\raggedright\arraybackslash}p{31.04mm}>{\raggedright\arraybackslash}p{36.68mm}>{\raggedright\arraybackslash}p{50.65mm}@{}}
\toprule
\textbf{Adjective} & \textbf{Meaning} & \textbf{Inflection guide} \\
\midrule
gōd & good & long-stem model \\
glæd & glad & short stem; æ/\allowbreak{}a \\
micel & great; large & micl-/\allowbreak{}micel- \\
īdel & empty; void & īdl-/\allowbreak{}īdel- \\
ǣmtig & empty; void & retain ǣmtig- \\
æppelbǣre & fruit-bearing & -e adjective; æppelbǣr- \\
lȳtel & little; small & lȳtl-/\allowbreak{}lȳtel- \\
styriendlic & moving & regular long stem \\
lēof & dear & regular long stem \\
eald & old & regular long stem \\
yfel & bad; evil & selected full yfel- stem \\
gearu & ready & gearw-/\allowbreak{}gearo- \\
hēah & high & special contraction \\
\bottomrule
\end{tabular}
\par\endgroup

% mdtable:T01:end
% mdparagraph:P004:begin
Learn the meaning with the inflection guide. \textbf{Styriendlic} is an adjective meaning “moving”; its use here does not require a participle paradigm. The particular compound \textbf{æppelbǣre} is Cook’s word for “fruit-bearing”. The following tables and stem map give the complete selected forms of every adjective in this list.
% mdparagraph:P004:end

\reserve{7\baselineskip}
% mdparagraph:P005:begin
\section*{2. Agreement and the two declensions}\BookHeading{nav-ch03-learner-0002}{subsection}{2. Agreement and the two declensions}{1}\BookLabel{bk-ch03-learner-dbf98476e10a}
% mdparagraph:P005:end

% teaching: NEW-U03-GRAM-agreement
% teaching: NEW-U03-GRAM-determiner-conditions
% teaching: NEW-U03-GRAM-predicate-adjective
% mdparagraph:P006:begin
An adjective agrees with its own noun in gender, number and case. An attributive adjective belongs inside a noun phrase. A predicate adjective states something about the subject with the copula and is nominative. Agreement is grammatical: a neuter noun takes a neuter adjective regardless of the English meaning.
% mdparagraph:P006:end

% mdparagraph:P007:begin
Ordinary adjectives have a \textbf{strong} and a \textbf{weak} declension. Use the strong form in a bare noun phrase and for an ordinary positive predicate adjective. Use the weak form after the article or a demonstrative. An article inside the subject does not make its predicate adjective weak. Comparatives are always weak; ordinals and superlatives have the rules given later.
% mdparagraph:P007:end

% mdparagraph:P008:begin
The following examples are newly composed teaching material.
% mdparagraph:P008:end

% example: NEW-U03-EXAMPLE-01
% mdparagraph:P009:begin
Gōd męte is on ðǣre eorðan.
% mdparagraph:P009:end

% mdparagraph:P010:begin
Good food is on the earth.
% mdparagraph:P010:end

% mdparagraph:P011:begin
Gōd agrees with the masculine nominative singular męte. The locative phrase retains the dative taught in Unit 1.
% mdparagraph:P011:end

% example: NEW-U03-EXAMPLE-02
% mdparagraph:P012:begin
Se gōda mann hæfð gōdne męte.
% mdparagraph:P012:end

% mdparagraph:P013:begin
The good man has good food.
% mdparagraph:P013:end

% mdparagraph:P014:begin
Gōda is weak after se and agrees with the subject mann. Gōdne is strong masculine accusative singular and agrees with the bare object męte.
% mdparagraph:P014:end

% example: NEW-U03-EXAMPLE-03
% mdparagraph:P015:begin
Ðæt gōde trēo is micel.
% mdparagraph:P015:end

% mdparagraph:P016:begin
The good tree is great.
% mdparagraph:P016:end

% mdparagraph:P017:begin
Gōde is weak inside the determined subject phrase. Micel is a strong neuter nominative singular predicate.
% mdparagraph:P017:end

% mdparagraph:P018:begin
A\BookPoint{bk-ch03-learner-65f86edd6e8d} following ordinary adjective may be strong or weak after a possessive: Cook describes the weak form as occurring sometimes. Do not turn possession into an invariable weak-adjective rule. The exercises avoid requiring a new general rule for an adjective after a limiting noun genitive.
% mdparagraph:P018:end

\clearpage
% mdparagraph:P019:begin
\section*{3. Complete strong and weak models}\BookHeading{nav-ch03-learner-0003}{subsection}{3. Complete strong and weak models}{1}
% mdparagraph:P019:end

% teaching: NEW-U03-GRAM-strong-adjectives
% teaching: NEW-U03-GRAM-long-stems
\reserve{7\baselineskip}
% mdparagraph:P020:begin
\subsection*{Strong gōd}\BookHeading{nav-ch03-learner-0004}{subsubsection}{Strong gōd}{0}\BookLabel{bk-ch03-learner-7fee105ca648}
% mdparagraph:P020:end

% teaching: NEW-U03-PAR-god-adjective-strong
% mdtable:T02:begin
% course-table-style: 2026-10
% course-panel:T02-P01:begin
\par\begingroup
\fontsize{10.5}{13.2}\selectfont
\setlength{\LTleft}{0pt}\setlength{\LTright}{\fill}
\noindent
\begin{tabular}{@{}>{\raggedright\arraybackslash}p{22.0mm}>{\raggedright\arraybackslash}p{31.188mm}>{\raggedright\arraybackslash}p{31.188mm}>{\raggedright\arraybackslash}p{31.188mm}@{}}
\toprule
\textbf{Case} & \textbf{Masculine singular} & \textbf{Neuter singular} & \textbf{Feminine singular} \\
\midrule
Nominative & gōd & gōd & gōd \\
Genitive & gōdes & gōdes & gōdre \\
Dative & gōdum & gōdum & gōdre \\
Accusative & gōdne & gōd & gōde \\
\bottomrule
\end{tabular}
\par\endgroup
% course-panel:T02-P01:end
\par\addvspace{5pt}
% course-panel:T02-P02:begin
\par\begingroup
\fontsize{10.5}{13.2}\selectfont
\setlength{\LTleft}{0pt}\setlength{\LTright}{\fill}
\noindent
\begin{tabular}{@{}>{\raggedright\arraybackslash}p{22.0mm}>{\raggedright\arraybackslash}p{31.188mm}>{\raggedright\arraybackslash}p{31.188mm}>{\raggedright\arraybackslash}p{31.188mm}@{}}
\toprule
\textbf{Case} & \textbf{Masculine plural} & \textbf{Neuter plural} & \textbf{Feminine plural} \\
\midrule
Nominative & gōde & gōd & gōda \\
Genitive & gōdra & gōdra & gōdra \\
Dative & gōdum & gōdum & gōdum \\
Accusative & gōde & gōd & gōda \\
\bottomrule
\end{tabular}
\par\endgroup
% course-panel:T02-P02:end
% mdtable:T02:end
% mdparagraph:P021:begin
Instrumental singular, recognition only: masculine gōde; neuter gōde.
% mdparagraph:P021:end

% mdparagraph:P022:begin
The endingless feminine nominative singular and neuter nominative/\allowbreak{}accusative plural belong to the selected long-stem pattern. Masculine and feminine plurals need their own endings even when the English adjective is unchanged. Neuter nominative and accusative coincide in both numbers; they remain different cases.
% mdparagraph:P022:end

% teaching: NEW-U03-GRAM-weak-adjectives
\reserve{7\baselineskip}
% mdparagraph:P023:begin
\subsection*{Weak gōd}\BookHeading{nav-ch03-learner-0005}{subsubsection}{Weak gōd}{0}\BookLabel{bk-ch03-learner-8c480e64ad30}
% mdparagraph:P023:end

% teaching: NEW-U03-PAR-god-adjective-weak
% mdtable:T03:begin
% course-table-style: 2026-10
% course-panel:T03-P01:begin
\par\begingroup
\fontsize{10.5}{13.2}\selectfont
\setlength{\LTleft}{0pt}\setlength{\LTright}{\fill}
\noindent
\begin{tabular}{@{}>{\raggedright\arraybackslash}p{22.0mm}>{\raggedright\arraybackslash}p{31.188mm}>{\raggedright\arraybackslash}p{31.188mm}>{\raggedright\arraybackslash}p{31.188mm}@{}}
\toprule
\textbf{Case} & \textbf{Masculine singular} & \textbf{Neuter singular} & \textbf{Feminine singular} \\
\midrule
Nominative & gōda & gōde & gōde \\
Genitive & gōdan & gōdan & gōdan \\
Dative & gōdan & gōdan & gōdan \\
Accusative & gōdan & gōde & gōdan \\
\bottomrule
\end{tabular}
\par\endgroup
% course-panel:T03-P01:end
\par\addvspace{5pt}
% course-panel:T03-P02:begin
\par\begingroup
\fontsize{10.5}{13.2}\selectfont
\setlength{\LTleft}{0pt}\setlength{\LTright}{\fill}
\noindent
\begin{tabular}{@{}>{\raggedright\arraybackslash}p{22.0mm}>{\raggedright\arraybackslash}p{31.188mm}>{\raggedright\arraybackslash}p{31.188mm}>{\raggedright\arraybackslash}p{31.188mm}@{}}
\toprule
\textbf{Case} & \textbf{Masculine plural} & \textbf{Neuter plural} & \textbf{Feminine plural} \\
\midrule
Nominative & gōdan & gōdan & gōdan \\
Genitive & gōdra & gōdra & gōdra \\
Dative & gōdum & gōdum & gōdum \\
Accusative & gōdan & gōdan & gōdan \\
\bottomrule
\end{tabular}
\par\endgroup
% course-panel:T03-P02:end
% mdtable:T03:end
% mdparagraph:P024:begin
The weak adjective has -a in masculine nominative singular, -e in feminine nominative singular and neuter nominative/\allowbreak{}accusative singular, and -an in the other singular cells. All genders share the plural forms. Its selected genitive plural is -ra; do not replace it mechanically with a weak noun’s -ena. Cook also admits weak-adjective genitive plurals in -a and -ena. The comparative exceptions below are learned explicitly.
% mdparagraph:P024:end

% teaching: NEW-U03-GRAM-instrumental-reference
% mdparagraph:P025:begin
The distinct strong instrumental singular is recorded for recognition. It has the same spelling as several ordinary -e forms, so a spelling alone does not identify its case. No instrumental phrase or construction is required in independent production here. Weak instrumental forms coincide with the corresponding datives; no separate instrumental use is being introduced.
% mdparagraph:P025:end

\reserve{7\baselineskip}
% mdparagraph:P026:begin
\section*{4. Stems and important exceptions}\BookHeading{nav-ch03-learner-0006}{subsection}{4. Stems and important exceptions}{1}
% mdparagraph:P026:end

% teaching: NEW-U03-GRAM-short-stems
% teaching: NEW-U03-GRAM-ae-a
\reserve{7\baselineskip}
% mdparagraph:P027:begin
\subsection*{Short-stem glæd}\BookHeading{nav-ch03-learner-0007}{subsubsection}{Short-stem glæd}{0}\BookLabel{bk-ch03-learner-fa69c0c29f5c}
% mdparagraph:P027:end

% teaching: NEW-U03-PAR-glaed-strong
% mdtable:T04:begin
% course-table-style: 2026-10
% course-panel:T04-P01:begin
\par\begingroup
\fontsize{10.5}{13.2}\selectfont
\setlength{\LTleft}{0pt}\setlength{\LTright}{\fill}
\noindent
\begin{tabular}{@{}>{\raggedright\arraybackslash}p{22.0mm}>{\raggedright\arraybackslash}p{31.188mm}>{\raggedright\arraybackslash}p{31.188mm}>{\raggedright\arraybackslash}p{31.188mm}@{}}
\toprule
\textbf{Case} & \textbf{Masculine singular} & \textbf{Neuter singular} & \textbf{Feminine singular} \\
\midrule
Nominative & glæd & glæd & gladu \\
Genitive & glades & glades & glædre \\
Dative & gladum & gladum & glædre \\
Accusative & glædne & glæd & glade \\
\bottomrule
\end{tabular}
\par\endgroup
% course-panel:T04-P01:end
\par\addvspace{5pt}
% course-panel:T04-P02:begin
\par\begingroup
\fontsize{10.5}{13.2}\selectfont
\setlength{\LTleft}{0pt}\setlength{\LTright}{\fill}
\noindent
\begin{tabular}{@{}>{\raggedright\arraybackslash}p{22.0mm}>{\raggedright\arraybackslash}p{31.188mm}>{\raggedright\arraybackslash}p{31.188mm}>{\raggedright\arraybackslash}p{31.188mm}@{}}
\toprule
\textbf{Case} & \textbf{Masculine plural} & \textbf{Neuter plural} & \textbf{Feminine plural} \\
\midrule
Nominative & glade & gladu & glada \\
Genitive & glædra & glædra & glædra \\
Dative & gladum & gladum & gladum \\
Accusative & glade & gladu & glada \\
\bottomrule
\end{tabular}
\par\endgroup
% course-panel:T04-P02:end
% mdtable:T04:end
% mdparagraph:P028:begin
Instrumental singular, recognition only: masculine glade; neuter glade.
% mdparagraph:P028:end

% teaching: NEW-U03-PAR-glaed-weak
% mdparagraph:P029:begin
The selected short-stem feminine nominative singular and neuter nominative/\allowbreak{}accusative plural have -u. Glæd uses glad- before the vowel endings and glæd- before -ne, -re and -ra. Its complete weak declension is the weak gōd pattern with glad- before -a, -e, -an and -um, but glæd- before genitive plural -ra. This gives a complete pattern; it is not a rule for changing æ to a in every adjective.
% mdparagraph:P029:end

% teaching: NEW-U03-GRAM-syncope
% teaching: NEW-U03-GRAM-e-stem
% teaching: NEW-U03-GRAM-u-w
\reserve{7\baselineskip}
% mdparagraph:P030:begin
\subsection*{A complete application map}\BookHeading{nav-ch03-learner-0008}{subsubsection}{A complete application map}{0}\BookLabel{bk-ch03-learner-ccc3cc60a047}
% mdparagraph:P030:end

% mdparagraph:P031:begin
Use the strong and weak gōd panels with this map. \textbf{V-stem} is used before endings beginning with a vowel: -es, -um, -e, -a and -an. \textbf{C-stem} is used before -ne, -re and -ra. The strong masculine nominative singular and neuter nominative/\allowbreak{}accusative singular are the listed headword. The last two columns replace the corresponding endingless gōd cells. The strong instrumental singular is V-stem + -e. Thus every strong and weak N/\allowbreak{}G/\allowbreak{}D/\allowbreak{}A cell is determined; do not add an ending twice to a form already listed in the last columns.
% mdparagraph:P031:end

% teaching: NEW-U03-PAR-micel-strong
% teaching: NEW-U03-PAR-micel-weak
% teaching: NEW-U03-PAR-idel-strong
% teaching: NEW-U03-PAR-idel-weak
% teaching: NEW-U03-PAR-aemtig-strong
% teaching: NEW-U03-PAR-aemtig-weak
% teaching: NEW-U03-PAR-aeppelbaere-strong
% teaching: NEW-U03-PAR-aeppelbaere-weak
% teaching: NEW-U03-PAR-lytel-strong
% teaching: NEW-U03-PAR-lytel-weak
% teaching: NEW-U03-PAR-styriendlic-strong
% teaching: NEW-U03-PAR-styriendlic-weak
% teaching: NEW-U03-PAR-leof-strong
% teaching: NEW-U03-PAR-leof-weak
% teaching: NEW-U03-PAR-eald-strong
% teaching: NEW-U03-PAR-eald-weak
% teaching: NEW-U03-PAR-yfel-strong
% teaching: NEW-U03-PAR-yfel-weak
% teaching: NEW-U03-PAR-gearu-strong
% teaching: NEW-U03-PAR-gearu-weak
% mdtable:T05:begin
\par\begingroup
\fontsize{10.5}{13.2}\selectfont
\setlength{\LTleft}{0pt}\setlength{\LTright}{\fill}
\noindent
\begin{tabular}{@{}>{\raggedright\arraybackslash}p{19.73mm}>{\raggedright\arraybackslash}p{21.53mm}>{\raggedright\arraybackslash}p{21.53mm}>{\raggedright\arraybackslash}p{25.38mm}>{\raggedright\arraybackslash}p{24.58mm}@{}}
\toprule
\textbf{Headword} & \textbf{V-stem} & \textbf{C-stem} & \textbf{Strong f. nom. sg.} & \textbf{Strong n. N/\allowbreak{}A pl.} \\
\midrule
gōd & gōd- & gōd- & gōd & gōd \\
glæd & glad- & glæd- & gladu & gladu \\
micel & micl- & micel- & miclu & miclu \\
īdel & īdl- & īdel- & īdel & īdlu \\
ǣmtig & ǣmtig- & ǣmtig- & ǣmtig & ǣmtig \\
æppelbǣre & æppelbǣr- & æppelbǣr- & æppelbǣru & æppelbǣru \\
lȳtel & lȳtl- & lȳtel- & lȳtlu & lȳtlu \\
styriendlic & styriendlic- & styriendlic- & styriendlic & styriendlic \\
lēof & lēof- & lēof- & lēof & lēof \\
eald & eald- & eald- & eald & eald \\
yfel & yfel- & yfel- & yfel & yfel \\
gearu & gearw- & gearo- & gearu & gearu \\
\bottomrule
\end{tabular}
\par\endgroup

% mdtable:T05:end
% mdparagraph:P032:begin
In\BookPoint{bk-ch03-learner-1b4dba13d70a} micel, lȳtel and īdel, the medial vowel is lost before the selected vowel endings but retained before -ne, -re and -ra. The selected īdel feminine nominative singular is endingless, while its neuter plural is īdlu. Retain the full stem in ǣmtig and in the selected yfel pattern. A long disyllable does not license automatic deletion of every medial vowel.
% mdparagraph:P032:end

% mdparagraph:P033:begin
In æppelbǣre the final headword -e is replaced by the endings, rather than treated as part of the stem. Its C-stem ends in r: adding -re or -ra therefore writes rr. Gearu has w before vowel endings but the selected gearo- before consonantal endings; its listed direct forms remain gearu. The map includes the weak declensions of both words.
% mdparagraph:P033:end

% mdparagraph:P034:begin
Selected alternatives worth recognizing: strong feminine nominative micel or micelu beside miclu; lȳtel or lȳtelu beside lȳtlu; īdelu beside īdel. Corresponding neuter plurals may retain or omit -u, and strong feminine N/\allowbreak{}A plurals often have -e beside selected -a. These alternatives do not permit changing the actual spelling of a source passage. In particular, the reading prints miclu.
% mdparagraph:P034:end

% teaching: NEW-U03-GRAM-h-loss
\reserve{7\baselineskip}
% mdparagraph:P035:begin
\subsection*{Hēah}\BookHeading{nav-ch03-learner-0009}{subsubsection}{Hēah}{0}\BookLabel{bk-ch03-learner-480a776eec52}
% mdparagraph:P035:end

% mdparagraph:P036:begin
Hēah loses h and contracts in much of its inflection. Use this complete selected pattern instead of attaching the ordinary endings directly to the headword.
% mdparagraph:P036:end

% teaching: NEW-U03-PAR-heah-strong
\par\begin{minipage}{\linewidth}\raggedright
% mdparagraph:P037:begin
\textbf{Strong hēah.}
% mdparagraph:P037:end

% mdtable:T06:begin
% course-table-style: 2026-10
% course-panel:T06-P01:begin
\par\begingroup
\fontsize{10.5}{13.2}\selectfont
\setlength{\LTleft}{0pt}\setlength{\LTright}{\fill}
\noindent
\begin{tabular}{@{}>{\raggedright\arraybackslash}p{22.0mm}>{\raggedright\arraybackslash}p{31.188mm}>{\raggedright\arraybackslash}p{31.188mm}>{\raggedright\arraybackslash}p{31.188mm}@{}}
\toprule
\textbf{Case} & \textbf{Masculine singular} & \textbf{Neuter singular} & \textbf{Feminine singular} \\
\midrule
Nominative & hēah & hēah & hēa \\
Genitive & hēas & hēas & hēarre \\
Dative & hēaum & hēaum & hēarre \\
Accusative & hēanne & hēah & hēa \\
\bottomrule
\end{tabular}
\par\endgroup
% course-panel:T06-P01:end
\par\addvspace{5pt}
% course-panel:T06-P02:begin
\par\begingroup
\fontsize{10.5}{13.2}\selectfont
\setlength{\LTleft}{0pt}\setlength{\LTright}{\fill}
\noindent
\begin{tabular}{@{}>{\raggedright\arraybackslash}p{22.0mm}>{\raggedright\arraybackslash}p{31.188mm}>{\raggedright\arraybackslash}p{31.188mm}>{\raggedright\arraybackslash}p{31.188mm}@{}}
\toprule
\textbf{Case} & \textbf{Masculine plural} & \textbf{Neuter plural} & \textbf{Feminine plural} \\
\midrule
Nominative & hēa & hēa & hēa \\
Genitive & hēarra & hēarra & hēarra \\
Dative & hēaum & hēaum & hēaum \\
Accusative & hēa & hēa & hēa \\
\bottomrule
\end{tabular}
\par\endgroup
% course-panel:T06-P02:end
% mdtable:T06:end
\end{minipage}\par
% mdparagraph:P038:begin
Instrumental singular, recognition only: masculine hēa; neuter hēa.
% mdparagraph:P038:end

% teaching: NEW-U03-PAR-heah-weak
\par\begin{minipage}{\linewidth}\raggedright
% mdparagraph:P039:begin
\textbf{Weak hēah.}
% mdparagraph:P039:end

% mdtable:T07:begin
% course-table-style: 2026-10
% course-panel:T07-P01:begin
\par\begingroup
\fontsize{10.5}{13.2}\selectfont
\setlength{\LTleft}{0pt}\setlength{\LTright}{\fill}
\noindent
\begin{tabular}{@{}>{\raggedright\arraybackslash}p{22.0mm}>{\raggedright\arraybackslash}p{31.188mm}>{\raggedright\arraybackslash}p{31.188mm}>{\raggedright\arraybackslash}p{31.188mm}@{}}
\toprule
\textbf{Case} & \textbf{Masculine singular} & \textbf{Neuter singular} & \textbf{Feminine singular} \\
\midrule
Nominative & hēa & hēa & hēa \\
Genitive & hēan & hēan & hēan \\
Dative & hēan & hēan & hēan \\
Accusative & hēan & hēa & hēan \\
\bottomrule
\end{tabular}
\par\endgroup
% course-panel:T07-P01:end
\par\addvspace{5pt}
% course-panel:T07-P02:begin
\par\begingroup
\fontsize{10.5}{13.2}\selectfont
\setlength{\LTleft}{0pt}\setlength{\LTright}{\fill}
\noindent
\begin{tabular}{@{}>{\raggedright\arraybackslash}p{22.0mm}>{\raggedright\arraybackslash}p{31.188mm}>{\raggedright\arraybackslash}p{31.188mm}>{\raggedright\arraybackslash}p{31.188mm}@{}}
\toprule
\textbf{Case} & \textbf{Masculine plural} & \textbf{Neuter plural} & \textbf{Feminine plural} \\
\midrule
Nominative & hēan & hēan & hēan \\
Genitive & hēarra & hēarra & hēarra \\
Dative & hēaum & hēaum & hēaum \\
Accusative & hēan & hēan & hēan \\
\bottomrule
\end{tabular}
\par\endgroup
% course-panel:T07-P02:end
% mdtable:T07:end
\end{minipage}\par
% mdparagraph:P040:begin
The selected hēaum has the supported alternative hēam. Hēanne may appear as hēane; hēarre/\allowbreak{}hēarra as hēare/\allowbreak{}hēara. The tables retain their selected doubling and vowel sequence. These forms show why consonant loss, doubling and contraction must be learned with the paradigm.
% mdparagraph:P040:end

% teaching: NEW-U03-VOC-eall
% teaching: NEW-U03-PAR-eall-strong
\reserve{7\baselineskip}
% mdparagraph:P041:begin
\subsection*{Reference: eall}\BookHeading{nav-ch03-learner-0010}{subsubsection}{Reference: eall}{0}
% mdparagraph:P041:end

% mdparagraph:P042:begin
Eall means “all”. Its strong morphology is for recognition here; no independent quantifier construction is required. Apply the complete strong gōd pattern with headword eall, V-stem eall-, C-stem eal-, feminine nominative singular eall and neuter N/\allowbreak{}A plural eall. The instrumental singular is ealle. The single l before -ne/\allowbreak{}-re/\allowbreak{}-ra is deliberate. Do not construct a weak paradigm from this reference.
% mdparagraph:P042:end

\reserve{7\baselineskip}
% mdparagraph:P043:begin
\section*{5. Things and counted lights}\BookHeading{nav-ch03-learner-0011}{subsection}{5. Things and counted lights}{1}\BookLabel{bk-ch03-learner-aad11cd34c34}
% mdparagraph:P043:end

% teaching: NEW-U03-VOC-thing
% teaching: NEW-U03-VOC-leoht
% teaching: NEW-U03-GRAM-thing-noun
% teaching: NEW-U03-GRAM-light-count-sense
% teaching: NEW-U03-PAR-thing
% teaching: NEW-U03-PAR-leoht-count
% mdparagraph:P044:begin
Ðing, neuter, means “thing”. Lēoht, neuter, can now also name an individual light or luminary. This counted sense differs from the mass sense “light” already learned. Both follow the complete heavy-neuter pattern:
% mdparagraph:P044:end

% mdtable:T08:begin
% course-table-style: 2026-10
% course-panel:T08-P01:begin
\par\begingroup
\fontsize{10.5}{13.2}\selectfont
\setlength{\LTleft}{0pt}\setlength{\LTright}{\fill}
\noindent
\begin{tabular}{@{}>{\raggedright\arraybackslash}p{22.0mm}>{\raggedright\arraybackslash}p{22.688mm}>{\raggedright\arraybackslash}p{22.688mm}>{\raggedright\arraybackslash}p{22.688mm}>{\raggedright\arraybackslash}p{22.688mm}@{}}
\toprule
\textbf{Case} & \textbf{ðing singular} & \textbf{ðing plural} & \textbf{lēoht singular} & \textbf{lēoht plural} \\
\midrule
Nominative & ðing & ðing & lēoht & lēoht \\
Genitive & ðinges & ðinga & lēohtes & lēohta \\
Dative & ðinge & ðingum & lēohte & lēohtum \\
Accusative & ðing & ðing & lēoht & lēoht \\
\bottomrule
\end{tabular}
\par\endgroup
% course-panel:T08-P01:end
% mdtable:T08:end
% mdparagraph:P045:begin
Use all four plural cases of counted lēoht productively. The unchanged spelling lēoht in nominative/\allowbreak{}accusative plural does not make it singular. Number is also shown by its determiner, numeral, adjective and clause context.
% mdparagraph:P045:end

\reserve{7\baselineskip}
% mdparagraph:P046:begin
\subsection*{Adjectives — A1–A10}\BookHeading{nav-ch03-learner-0012}{subsubsection}{Adjectives — A1–A10}{0}
% mdparagraph:P046:end

\par\begin{minipage}{\linewidth}\raggedright
% mdparagraph:P047:begin
\textbf{A1.}
% mdparagraph:P047:end

% item: NEW-U03-A01
% mdparagraph:P048:begin
Give the complete strong nominative, genitive, dative and accusative of gōd in all three genders, singular and plural. Label every case even where spellings coincide.
% mdparagraph:P048:end

\end{minipage}\par\addvspace{4pt}

\par\begin{minipage}{\linewidth}\raggedright
% mdparagraph:P049:begin
\textbf{A2.}
% mdparagraph:P049:end

% item: NEW-U03-A02
% mdparagraph:P050:begin
Give the complete strong nominative, genitive, dative and accusative of glæd in all three genders, singular and plural. Mark the forms with a rather than æ.
% mdparagraph:P050:end

\end{minipage}\par\addvspace{4pt}

\par\begin{minipage}{\linewidth}\raggedright
% mdparagraph:P051:begin
\textbf{A3.}
% mdparagraph:P051:end

% item: NEW-U03-A03
% mdparagraph:P052:begin
Give the complete weak nominative, genitive, dative and accusative of gōd in all three genders, singular and plural. Contrast its genitive plural with a weak noun’s genitive plural.
% mdparagraph:P052:end

\end{minipage}\par\addvspace{4pt}

\par\begin{minipage}{\linewidth}\raggedright
% mdparagraph:P053:begin
\textbf{A4.}
% mdparagraph:P053:end

% item: NEW-U03-A04
% mdparagraph:P054:begin
For micel, lȳtel, īdel and ǣmtig, give the strong masculine genitive singular, strong feminine nominative singular, strong neuter nominative plural and weak neuter nominative singular. Preserve the selected stems and vowel marks.
% mdparagraph:P054:end

\end{minipage}\par\addvspace{4pt}

\par\begin{minipage}{\linewidth}\raggedright
% mdparagraph:P055:begin
\textbf{A5.}
% mdparagraph:P055:end

% item: NEW-U03-A05
% mdparagraph:P056:begin
For æppelbǣre, give the strong masculine and feminine nominative singular, masculine accusative singular, neuter nominative plural and genitive plural. Explain why the headword ends in -e. For gearu and hēah, give the strong masculine nominative, genitive, dative and accusative singular and the weak masculine nominative singular.
% mdparagraph:P056:end

\end{minipage}\par\addvspace{4pt}

\par\begin{minipage}{\linewidth}\raggedright
% mdparagraph:P057:begin
\textbf{A6.}
% mdparagraph:P057:end

% item: NEW-U03-A06
% mdparagraph:P058:begin
Give the four cases of ðing in the singular and plural. Give its meaning in the plural phrase ðā ðing.
% mdparagraph:P058:end

\end{minipage}\par\addvspace{4pt}

\par\begin{minipage}{\linewidth}\raggedright
% mdparagraph:P059:begin
\textbf{A7.}
% mdparagraph:P059:end

% item: NEW-U03-A07
% mdparagraph:P060:begin
These phrases are nominative. Add the appropriate article to each and change the adjective as needed: gōd fugol; gōd wyrt; gōd trēo; gōde fuglas. Keep the same case and number.
% mdparagraph:P060:end

\end{minipage}\par\addvspace{4pt}

\par\begin{minipage}{\linewidth}\raggedright
% mdparagraph:P061:begin
\textbf{A8.}
% mdparagraph:P061:end

% item: NEW-U03-A08
% mdparagraph:P062:begin
Translate and explain the different adjective forms: Ðā gōdan fuglas sindon on ðǣm hofe. Ðā fuglas sindon gōde.
% mdparagraph:P062:end

\end{minipage}\par\addvspace{4pt}

\par\begin{minipage}{\linewidth}\raggedright
% mdparagraph:P063:begin
\textbf{A9.}
% mdparagraph:P063:end

% item: NEW-U03-A09
% mdparagraph:P064:begin
Translate “The birds live in the good man’s dwelling.” Put the limiting genitive before the head noun. Identify the case of the dwelling and of the man.
% mdparagraph:P064:end

\end{minipage}\par\addvspace{4pt}

\par\begin{minipage}{\linewidth}\raggedright
% mdparagraph:P065:begin
\textbf{A10.}
% mdparagraph:P065:end

% item: NEW-U03-A10
% mdparagraph:P066:begin
The strong table contains gōde as masculine nominative plural and as masculine or neuter instrumental singular. State each analysis. Compare the corresponding entries of glæd.
% mdparagraph:P066:end

\end{minipage}\par\addvspace{4pt}

\reserve{7\baselineskip}
% mdparagraph:P067:begin
\section*{6. Agreeing possessives}\BookHeading{nav-ch03-learner-0013}{subsection}{6. Agreeing possessives}{1}\BookLabel{bk-ch03-learner-0a3e4174931b}
% mdparagraph:P067:end

% teaching: NEW-U03-GRAM-possessive-agreement
% teaching: NEW-U03-GRAM-possessor-reference
% teaching: NEW-U03-VOC-min-poss
% teaching: NEW-U03-VOC-thin-poss
% teaching: NEW-U03-VOC-uncer-poss
% teaching: NEW-U03-VOC-incer-poss
% teaching: NEW-U03-VOC-ure-poss
% teaching: NEW-U03-VOC-eower-poss
% mdtable:T09:begin
\par\begingroup
\fontsize{10.5}{13.2}\selectfont
\setlength{\LTleft}{0pt}\setlength{\LTright}{\fill}
\noindent
\begin{tabular}{@{}>{\raggedright\arraybackslash}p{47.38mm}>{\raggedright\arraybackslash}p{39.74mm}>{\raggedright\arraybackslash}p{31.26mm}@{}}
\toprule
\textbf{Possessor} & \textbf{Agreeing possessive} & \textbf{Meaning} \\
\midrule
the speaker & mīn & my \\
one addressee & ðīn & your \\
the speaker and one other & uncer & our, of us two \\
two addressees & incer & your, of you two \\
the speaker and others & ūre & our \\
several addressees & ēower & your \\
\bottomrule
\end{tabular}
\par\endgroup

% mdtable:T09:end
% mdparagraph:P068:begin
The stem identifies the possessor; the ending agrees with the possessed noun. A dual possessor can own one thing or several things. A plural possessive does not require its head noun to be plural. These agreeing adjectives are distinct from the personal-genitive citation forms recognized in Unit 2. The third-person possessors his, hiere and hiera remain invariant.
% mdparagraph:P068:end

% example: NEW-U03-EXAMPLE-04
% mdparagraph:P069:begin
Ðū hæfst mīnne fisc.
% mdparagraph:P069:end

% mdparagraph:P070:begin
You have my fish.
% mdparagraph:P070:end

% mdparagraph:P071:begin
The possessor is the speaker. The possessed object is masculine accusative singular, so the selected possessive is mīnne.
% mdparagraph:P071:end

% example: NEW-U03-EXAMPLE-05
% mdparagraph:P072:begin
Hīe hīerað ūre fuglas.
% mdparagraph:P072:end

% mdparagraph:P073:begin
They hear our birds.
% mdparagraph:P073:end

% mdparagraph:P074:begin
The possessors are the speaker and others; the birds are masculine accusative plural. The possessive therefore has the selected form ūre.
% mdparagraph:P074:end

\reserve{7\baselineskip}
% mdparagraph:P075:begin
\subsection*{Mīn and ðīn}\BookHeading{nav-ch03-learner-0014}{subsubsection}{Mīn and ðīn}{0}\BookLabel{bk-ch03-learner-51b68bb635a4}
% mdparagraph:P075:end

% teaching: NEW-U03-PAR-min-poss-strong
% mdtable:T10:begin
% course-table-style: 2026-10
% course-panel:T10-P01:begin
\par\begingroup
\fontsize{10.5}{13.2}\selectfont
\setlength{\LTleft}{0pt}\setlength{\LTright}{\fill}
\noindent
\begin{tabular}{@{}>{\raggedright\arraybackslash}p{22.0mm}>{\raggedright\arraybackslash}p{31.188mm}>{\raggedright\arraybackslash}p{31.188mm}>{\raggedright\arraybackslash}p{31.188mm}@{}}
\toprule
\textbf{Case} & \textbf{Masculine singular} & \textbf{Neuter singular} & \textbf{Feminine singular} \\
\midrule
Nominative & mīn & mīn & mīn \\
Genitive & mīnes & mīnes & mīnre \\
Dative & mīnum & mīnum & mīnre \\
Accusative & mīnne & mīn & mīne \\
\bottomrule
\end{tabular}
\par\endgroup
% course-panel:T10-P01:end
\par\addvspace{5pt}
% course-panel:T10-P02:begin
\par\begingroup
\fontsize{10.5}{13.2}\selectfont
\setlength{\LTleft}{0pt}\setlength{\LTright}{\fill}
\noindent
\begin{tabular}{@{}>{\raggedright\arraybackslash}p{22.0mm}>{\raggedright\arraybackslash}p{31.188mm}>{\raggedright\arraybackslash}p{31.188mm}>{\raggedright\arraybackslash}p{31.188mm}@{}}
\toprule
\textbf{Case} & \textbf{Masculine plural} & \textbf{Neuter plural} & \textbf{Feminine plural} \\
\midrule
Nominative & mīne & mīn & mīna \\
Genitive & mīnra & mīnra & mīnra \\
Dative & mīnum & mīnum & mīnum \\
Accusative & mīne & mīn & mīna \\
\bottomrule
\end{tabular}
\par\endgroup
% course-panel:T10-P02:end
% mdtable:T10:end
% mdparagraph:P076:begin
Instrumental singular, recognition only: masculine mīne; neuter mīne.
% mdparagraph:P076:end

% teaching: NEW-U03-PAR-thin-poss-strong
% mdparagraph:P077:begin
The complete ðīn paradigm is obtained by replacing mīn- with ðīn- in every cell, preserving every ending and doubled n. Its masculine/\allowbreak{}neuter instrumental singular is ðīne. These two long stems do not take the short-stem -u in the selected feminine nominative singular or neuter N/\allowbreak{}A plural.
% mdparagraph:P077:end

\reserve{7\baselineskip}
% mdparagraph:P078:begin
\subsection*{Uncer and incer}\BookHeading{nav-ch03-learner-0015}{subsubsection}{Uncer and incer}{0}\BookLabel{bk-ch03-learner-c24379f75ac8}
% mdparagraph:P078:end

% teaching: NEW-U03-PAR-uncer-poss-strong
% mdtable:T11:begin
% course-table-style: 2026-10
% course-panel:T11-P01:begin
\par\begingroup
\fontsize{10.5}{13.2}\selectfont
\setlength{\LTleft}{0pt}\setlength{\LTright}{\fill}
\noindent
\begin{tabular}{@{}>{\raggedright\arraybackslash}p{22.0mm}>{\raggedright\arraybackslash}p{31.188mm}>{\raggedright\arraybackslash}p{31.188mm}>{\raggedright\arraybackslash}p{31.188mm}@{}}
\toprule
\textbf{Case} & \textbf{Masculine singular} & \textbf{Neuter singular} & \textbf{Feminine singular} \\
\midrule
Nominative & uncer & uncer & uncru \\
Genitive & uncres & uncres & uncerre \\
Dative & uncrum & uncrum & uncerre \\
Accusative & uncerne & uncer & uncre \\
\bottomrule
\end{tabular}
\par\endgroup
% course-panel:T11-P01:end
\par\addvspace{5pt}
% course-panel:T11-P02:begin
\par\begingroup
\fontsize{10.5}{13.2}\selectfont
\setlength{\LTleft}{0pt}\setlength{\LTright}{\fill}
\noindent
\begin{tabular}{@{}>{\raggedright\arraybackslash}p{22.0mm}>{\raggedright\arraybackslash}p{31.188mm}>{\raggedright\arraybackslash}p{31.188mm}>{\raggedright\arraybackslash}p{31.188mm}@{}}
\toprule
\textbf{Case} & \textbf{Masculine plural} & \textbf{Neuter plural} & \textbf{Feminine plural} \\
\midrule
Nominative & uncre & uncru & uncra \\
Genitive & uncerra & uncerra & uncerra \\
Dative & uncrum & uncrum & uncrum \\
Accusative & uncre & uncru & uncra \\
\bottomrule
\end{tabular}
\par\endgroup
% course-panel:T11-P02:end
% mdtable:T11:end
% mdparagraph:P079:begin
Instrumental singular, recognition only: masculine uncre; neuter uncre.
% mdparagraph:P079:end

% teaching: NEW-U03-PAR-incer-poss-strong
% mdparagraph:P080:begin
The complete incer paradigm replaces uncer-/\allowbreak{}uncr- with incer-/\allowbreak{}incr- in the corresponding cells. Retain the medial e before -ne, -re and -ra; omit it before the selected vowel endings. Its instrumental singular is incre. Unsyncopated unceru/\allowbreak{}inceru are supported alternatives to uncru/\allowbreak{}incru, and single-r forms such as uncere/\allowbreak{}uncera or incere/\allowbreak{}incera occur beside the selected doubled-r forms.
% mdparagraph:P080:end

\reserve{7\baselineskip}
% mdparagraph:P081:begin
\subsection*{Ūre and ēower}\BookHeading{nav-ch03-learner-0016}{subsubsection}{Ūre and ēower}{0}
% mdparagraph:P081:end

% teaching: NEW-U03-PAR-ure-poss-strong
% mdtable:T12:begin
% course-table-style: 2026-10
% course-panel:T12-P01:begin
\par\begingroup
\fontsize{10.5}{13.2}\selectfont
\setlength{\LTleft}{0pt}\setlength{\LTright}{\fill}
\noindent
\begin{tabular}{@{}>{\raggedright\arraybackslash}p{22.0mm}>{\raggedright\arraybackslash}p{31.188mm}>{\raggedright\arraybackslash}p{31.188mm}>{\raggedright\arraybackslash}p{31.188mm}@{}}
\toprule
\textbf{Case} & \textbf{Masculine singular} & \textbf{Neuter singular} & \textbf{Feminine singular} \\
\midrule
Nominative & ūre & ūre & ūre \\
Genitive & ūres & ūres & ūrre \\
Dative & ūrum & ūrum & ūrre \\
Accusative & ūrne & ūre & ūre \\
\bottomrule
\end{tabular}
\par\endgroup
% course-panel:T12-P01:end
\par\addvspace{5pt}
% course-panel:T12-P02:begin
\par\begingroup
\fontsize{10.5}{13.2}\selectfont
\setlength{\LTleft}{0pt}\setlength{\LTright}{\fill}
\noindent
\begin{tabular}{@{}>{\raggedright\arraybackslash}p{22.0mm}>{\raggedright\arraybackslash}p{31.188mm}>{\raggedright\arraybackslash}p{31.188mm}>{\raggedright\arraybackslash}p{31.188mm}@{}}
\toprule
\textbf{Case} & \textbf{Masculine plural} & \textbf{Neuter plural} & \textbf{Feminine plural} \\
\midrule
Nominative & ūre & ūru & ūra \\
Genitive & ūrra & ūrra & ūrra \\
Dative & ūrum & ūrum & ūrum \\
Accusative & ūre & ūru & ūra \\
\bottomrule
\end{tabular}
\par\endgroup
% course-panel:T12-P02:end
% mdtable:T12:end
% mdparagraph:P082:begin
Instrumental singular, recognition only: masculine ūre; neuter ūre.
% mdparagraph:P082:end

% teaching: NEW-U03-PAR-eower-poss-strong
% mdtable:T13:begin
% course-table-style: 2026-10
% course-panel:T13-P01:begin
\par\begingroup
\fontsize{10.5}{13.2}\selectfont
\setlength{\LTleft}{0pt}\setlength{\LTright}{\fill}
\noindent
\begin{tabular}{@{}>{\raggedright\arraybackslash}p{22.0mm}>{\raggedright\arraybackslash}p{31.188mm}>{\raggedright\arraybackslash}p{31.188mm}>{\raggedright\arraybackslash}p{31.188mm}@{}}
\toprule
\textbf{Case} & \textbf{Masculine singular} & \textbf{Neuter singular} & \textbf{Feminine singular} \\
\midrule
Nominative & ēower & ēower & ēowru \\
Genitive & ēowres & ēowres & ēowerre \\
Dative & ēowrum & ēowrum & ēowerre \\
Accusative & ēowerne & ēower & ēowre \\
\bottomrule
\end{tabular}
\par\endgroup
% course-panel:T13-P01:end
\par\addvspace{5pt}
% course-panel:T13-P02:begin
\par\begingroup
\fontsize{10.5}{13.2}\selectfont
\setlength{\LTleft}{0pt}\setlength{\LTright}{\fill}
\noindent
\begin{tabular}{@{}>{\raggedright\arraybackslash}p{22.0mm}>{\raggedright\arraybackslash}p{31.188mm}>{\raggedright\arraybackslash}p{31.188mm}>{\raggedright\arraybackslash}p{31.188mm}@{}}
\toprule
\textbf{Case} & \textbf{Masculine plural} & \textbf{Neuter plural} & \textbf{Feminine plural} \\
\midrule
Nominative & ēowre & ēowru & ēowra \\
Genitive & ēowerra & ēowerra & ēowerra \\
Dative & ēowrum & ēowrum & ēowrum \\
Accusative & ēowre & ēowru & ēowra \\
\bottomrule
\end{tabular}
\par\endgroup
% course-panel:T13-P02:end
% mdtable:T13:end
% mdparagraph:P083:begin
Instrumental singular, recognition only: masculine ēowre; neuter ēowre.
% mdparagraph:P083:end

% mdparagraph:P084:begin
Ūre\BookPoint{bk-ch03-learner-ed122a88dd3d} has feminine nominative singular ūre, as well as masculine and neuter nominative singular ūre; its selected neuter plural is ūru. Keep the table’s ūrre and ūrra distinct from the supported alternatives ūre and ūra. Ēower has ēowr- before vowel endings and ēower- before consonantal endings. Its supported unsyncopated -u alternative is ēoweru; single-r ēowere/\allowbreak{}ēowera also occur. These variants are valid alternatives, not reasons to ignore the selected paradigm when it is requested.
% mdparagraph:P084:end

\reserve{7\baselineskip}
% mdparagraph:P085:begin
\subsection*{Possessives — P1–P6}\BookHeading{nav-ch03-learner-0017}{subsubsection}{Possessives — P1–P6}{0}
% mdparagraph:P085:end

\par\begin{minipage}{\linewidth}\raggedright
% mdparagraph:P086:begin
\textbf{P1.}
% mdparagraph:P086:end

% item: NEW-U03-P01
% mdparagraph:P087:begin
Give the complete nominative, genitive, dative and accusative of the agreeing possessive mīn in all three genders, singular and plural. State what determines its ending.
% mdparagraph:P087:end

\end{minipage}\par\addvspace{4pt}

\par\begin{minipage}{\linewidth}\raggedright
% mdparagraph:P088:begin
\textbf{P2.}
% mdparagraph:P088:end

% item: NEW-U03-P02
% mdparagraph:P089:begin
Translate “I hear your bird. You hear my birds.” Address one person.
% mdparagraph:P089:end

\end{minipage}\par\addvspace{4pt}

\par\begin{minipage}{\linewidth}\raggedright
% mdparagraph:P090:begin
\textbf{P3.}
% mdparagraph:P090:end

% item: NEW-U03-P03
% mdparagraph:P091:begin
Translate “We two have a dwelling. Our dwelling is small. You two have a tree. Your tree is great.” Retain the two-person possessors.
% mdparagraph:P091:end

\end{minipage}\par\addvspace{4pt}

\par\begin{minipage}{\linewidth}\raggedright
% mdparagraph:P092:begin
\textbf{P4.}
% mdparagraph:P092:end

% item: NEW-U03-P04
% mdparagraph:P093:begin
Translate “Our birds live in your dwelling.” The possessors are several speakers and several addressees.
% mdparagraph:P093:end

\end{minipage}\par\addvspace{4pt}

\par\begin{minipage}{\linewidth}\raggedright
% mdparagraph:P094:begin
\textbf{P5.}
% mdparagraph:P094:end

% item: NEW-U03-P05
% mdparagraph:P095:begin
Translate “His birds live in my dwelling. Their dwelling is greater than our dwelling.” Use a male singular possessor for “his”.
% mdparagraph:P095:end

\end{minipage}\par\addvspace{4pt}

\par\begin{minipage}{\linewidth}\raggedright
% mdparagraph:P096:begin
\textbf{P6.}
% mdparagraph:P096:end

% item: NEW-U03-P06
% mdparagraph:P097:begin
The phrases mīne fuglas and uncer fugol are clause subjects; ēowrum hofe follows locative on. Identify each possessor’s person and number, then the possessed noun’s gender, number and case. Which information comes from the possessive stem, and which from its ending?
% mdparagraph:P097:end

\end{minipage}\par\addvspace{4pt}

\reserve{7\baselineskip}
% mdparagraph:P098:begin
\section*{7. Comparison}\BookHeading{nav-ch03-learner-0018}{subsection}{7. Comparison}{1}
% mdparagraph:P098:end

% teaching: NEW-U03-GRAM-comparison-regular
% teaching: NEW-U03-GRAM-comparison-umlaut
% teaching: NEW-U03-GRAM-comparison-irregular
% teaching: NEW-U03-VOC-leof-comp
% teaching: NEW-U03-VOC-leof-superl
% teaching: NEW-U03-PAR-leof-comp-weak
% teaching: NEW-U03-PAR-leof-superl-strong
% teaching: NEW-U03-PAR-leof-superl-weak
% teaching: NEW-U03-VOC-eald-comp
% teaching: NEW-U03-VOC-eald-superl
% teaching: NEW-U03-PAR-eald-comp-weak
% teaching: NEW-U03-PAR-eald-superl-strong
% teaching: NEW-U03-PAR-eald-superl-weak
% teaching: NEW-U03-VOC-god-adjective-comp
% teaching: NEW-U03-VOC-god-adjective-superl
% teaching: NEW-U03-PAR-god-adjective-comp-weak
% teaching: NEW-U03-PAR-god-adjective-superl-strong
% teaching: NEW-U03-PAR-god-adjective-superl-weak
% teaching: NEW-U03-VOC-yfel-comp
% teaching: NEW-U03-VOC-yfel-superl
% teaching: NEW-U03-PAR-yfel-comp-weak
% teaching: NEW-U03-PAR-yfel-superl-strong
% teaching: NEW-U03-PAR-yfel-superl-weak
% teaching: NEW-U03-VOC-micel-comp
% teaching: NEW-U03-VOC-micel-superl
% teaching: NEW-U03-PAR-micel-comp-weak
% teaching: NEW-U03-PAR-micel-superl-strong
% teaching: NEW-U03-PAR-micel-superl-weak
% teaching: NEW-U03-VOC-lytel-comp
% teaching: NEW-U03-VOC-lytel-superl
% teaching: NEW-U03-PAR-lytel-comp-weak
% teaching: NEW-U03-PAR-lytel-superl-strong
% teaching: NEW-U03-PAR-lytel-superl-weak
% mdparagraph:P099:begin
The\BookPoint{bk-ch03-learner-3909b65cce85} positive names a quality; the comparative expresses a greater or lesser degree; the superlative expresses the highest or lowest degree. The regular series adds -ra and -ost. Eald has an umlaut series; the four remaining series have forms to learn separately. The comparative column gives masculine nominative singular. The superlative column gives the strong masculine/\allowbreak{}neuter nominative singular.
% mdparagraph:P099:end

% mdtable:T14:begin
\par\begingroup
\fontsize{10.5}{13.2}\selectfont
\setlength{\LTleft}{0pt}\setlength{\LTright}{\fill}
\noindent
\begin{tabular}{@{}>{\raggedright\arraybackslash}p{20.94mm}>{\raggedright\arraybackslash}p{23.03mm}>{\raggedright\arraybackslash}p{23.03mm}>{\raggedright\arraybackslash}p{48.57mm}@{}}
\toprule
\textbf{Positive} & \textbf{Comparative} & \textbf{Superlative} & \textbf{Degree meanings} \\
\midrule
lēof & lēofra & lēofost & dearer; dearest \\
eald & ieldra & ieldest & older; oldest \\
gōd & bętera & bętst & better; best \\
yfel & wiersa & wierrest & worse; worst \\
micel & māra & mǣst & greater/\allowbreak{}larger; greatest/\allowbreak{}largest \\
lȳtel & lǣssa & lǣst & smaller; smallest \\
\bottomrule
\end{tabular}
\par\endgroup

% mdtable:T14:end
% teaching: NEW-U03-GRAM-comparative-weak
\reserve{7\baselineskip}
% mdparagraph:P100:begin
\subsection*{Complete comparative inflection}\BookHeading{nav-ch03-learner-0019}{subsubsection}{Complete comparative inflection}{0}\BookLabel{bk-ch03-learner-c8e0934fdf36}
% mdparagraph:P100:end

% mdparagraph:P101:begin
A comparative is weak, including when it is a predicate with no preceding article. Māra supplies a complete model:
% mdparagraph:P101:end

% mdtable:T15:begin
% course-table-style: 2026-10
% course-panel:T15-P01:begin
\par\begingroup
\fontsize{10.5}{13.2}\selectfont
\setlength{\LTleft}{0pt}\setlength{\LTright}{\fill}
\noindent
\begin{tabular}{@{}>{\raggedright\arraybackslash}p{22.0mm}>{\raggedright\arraybackslash}p{31.188mm}>{\raggedright\arraybackslash}p{31.188mm}>{\raggedright\arraybackslash}p{31.188mm}@{}}
\toprule
\textbf{Case} & \textbf{Masculine singular} & \textbf{Neuter singular} & \textbf{Feminine singular} \\
\midrule
Nominative & māra & māre & māre \\
Genitive & māran & māran & māran \\
Dative & māran & māran & māran \\
Accusative & māran & māre & māran \\
\bottomrule
\end{tabular}
\par\endgroup
% course-panel:T15-P01:end
\par\addvspace{5pt}
% course-panel:T15-P02:begin
\par\begingroup
\fontsize{10.5}{13.2}\selectfont
\setlength{\LTleft}{0pt}\setlength{\LTright}{\fill}
\noindent
\begin{tabular}{@{}>{\raggedright\arraybackslash}p{22.0mm}>{\raggedright\arraybackslash}p{31.188mm}>{\raggedright\arraybackslash}p{31.188mm}>{\raggedright\arraybackslash}p{31.188mm}@{}}
\toprule
\textbf{Case} & \textbf{Masculine plural} & \textbf{Neuter plural} & \textbf{Feminine plural} \\
\midrule
Nominative & māran & māran & māran \\
Genitive & mārra & mārra & mārra \\
Dative & mārum & mārum & mārum \\
Accusative & māran & māran & māran \\
\bottomrule
\end{tabular}
\par\endgroup
% course-panel:T15-P02:end
% mdtable:T15:end
% mdparagraph:P102:begin
For each other comparative, replace its final masculine nominative -a with -e for neuter/\allowbreak{}feminine nominative singular and neuter accusative singular; with -an for the other singular cases and all nominative/\allowbreak{}accusative plurals; and with -um for all dative plurals. Use the following genitive plurals, which are the same in all genders. These explicit forms complete each comparative paradigm.
% mdparagraph:P102:end

% mdtable:T16:begin
\par\begingroup
\fontsize{10.5}{13.2}\selectfont
\setlength{\LTleft}{0pt}\setlength{\LTright}{\fill}
\noindent
\begin{tabular}{@{}>{\raggedright\arraybackslash}p{65.35mm}>{\raggedright\arraybackslash}p{55.84mm}@{}}
\toprule
\textbf{Masculine nom. sg.} & \textbf{Genitive plural} \\
\midrule
lēofra & lēofra \\
ieldra & ieldra \\
bętera & bęterra \\
wiersa & wiersena \\
māra & mārra \\
lǣssa & lǣssena \\
\bottomrule
\end{tabular}
\par\endgroup

% mdtable:T16:end
% mdparagraph:P103:begin
Do not form an extra sequence of r’s by blindly appending -ra to the comparative stem. The selected -ena in wiersena and lǣssena is a licensed weak ending.
% mdparagraph:P103:end

% teaching: NEW-U03-GRAM-superlative
\reserve{7\baselineskip}
% mdparagraph:P104:begin
\subsection*{Superlatives}\BookHeading{nav-ch03-learner-0020}{subsubsection}{Superlatives}{0}\BookLabel{bk-ch03-learner-2b75c3552b41}
% mdparagraph:P104:end

% mdparagraph:P105:begin
Every listed superlative takes the complete strong and weak gōd patterns, replacing gōd- with the printed superlative stem. All six are selected long-stem patterns: strong feminine nominative singular and neuter N/\allowbreak{}A plural are endingless; strong instrumental singular has -e. An ordinary attributive superlative is normally weak; the taught predicative use is strong. The weak masculine nominatives are:
% mdparagraph:P105:end

% mdtable:T17:begin
\par\begingroup
\fontsize{10.5}{13.2}\selectfont
\setlength{\LTleft}{0pt}\setlength{\LTright}{\fill}
\noindent
\begin{tabular}{@{}>{\raggedright\arraybackslash}p{45.74mm}>{\raggedright\arraybackslash}p{34.41mm}>{\raggedright\arraybackslash}p{38.23mm}@{}}
\toprule
\textbf{Strong headword/\allowbreak{}stem} & \textbf{Weak m. nom. sg.} & \textbf{Weak n./\allowbreak{}f. nom. sg.} \\
\midrule
lēofost & lēofosta & lēofoste \\
ieldest & ieldesta & ieldeste \\
bętst & bętsta & bętste \\
wierrest & wierresta & wierreste \\
mǣst & mǣsta & mǣste \\
lǣst & lǣsta & lǣste \\
\bottomrule
\end{tabular}
\par\endgroup

% mdtable:T17:end
% mdparagraph:P106:begin
Use the same stem before all other strong/\allowbreak{}weak endings. This supplies complete inflection, not new partitive or substantival uses of the superlative.
% mdparagraph:P106:end

% teaching: NEW-U03-GRAM-comparison-thonne
% teaching: NEW-U03-VOC-thonne
\reserve{7\baselineskip}
% mdparagraph:P107:begin
\subsection*{A nominative standard with ðonne}\BookHeading{nav-ch03-learner-0021}{subsubsection}{A nominative standard with ðonne}{0}\BookLabel{bk-ch03-learner-91a7506b6176}
% mdparagraph:P107:end

% mdparagraph:P108:begin
Ðonne means “than” in the comparison taught here. When two subjects are compared with the copula, the second subject is nominative. The repeated copula can remain unexpressed. Restrict this pattern to the parallel copular comparison; it is not a rule for every clause containing “than”.
% mdparagraph:P108:end

% example: NEW-U03-EXAMPLE-07
% mdparagraph:P109:begin
Se mann is ieldra ðonne se fugol.
% mdparagraph:P109:end

% mdparagraph:P110:begin
The man is older than the bird.
% mdparagraph:P110:end

% mdparagraph:P111:begin
Mann and fugol are the contrasted nominative subjects. Ieldra is weak masculine nominative singular although it is a predicate.
% mdparagraph:P111:end

\reserve{7\baselineskip}
% mdparagraph:P112:begin
\subsection*{Comparison — C1–C7}\BookHeading{nav-ch03-learner-0022}{subsubsection}{Comparison — C1–C7}{0}
% mdparagraph:P112:end

\par\begin{minipage}{\linewidth}\raggedright
% mdparagraph:P113:begin
\textbf{C1.}
% mdparagraph:P113:end

% item: NEW-U03-C01
% mdparagraph:P114:begin
Give the selected positive, comparative and superlative masculine nominative singular of lēof, eald, gōd, micel, lȳtel and yfel. Use the weak superlative forms.
% mdparagraph:P114:end

\end{minipage}\par\addvspace{4pt}

\par\begin{minipage}{\linewidth}\raggedright
% mdparagraph:P115:begin
\textbf{C2.}
% mdparagraph:P115:end

% item: NEW-U03-C02
% mdparagraph:P116:begin
Give the complete weak nominative, genitive, dative and accusative of māra in all three genders, singular and plural.
% mdparagraph:P116:end

\end{minipage}\par\addvspace{4pt}

\par\begin{minipage}{\linewidth}\raggedright
% mdparagraph:P117:begin
\textbf{C3.}
% mdparagraph:P117:end

% item: NEW-U03-C03
% mdparagraph:P118:begin
Give the strong nominative singular of mǣst in all three genders and its weak nominative singular in all three genders. Then give the genitive and dative plurals in both declensions.
% mdparagraph:P118:end

\end{minipage}\par\addvspace{4pt}

\par\begin{minipage}{\linewidth}\raggedright
% mdparagraph:P119:begin
\textbf{C4.}
% mdparagraph:P119:end

% item: NEW-U03-C04
% mdparagraph:P120:begin
Translate “The great tree is greater than the small tree.” Identify the case of both trees and explain why the comparative has its selected ending.
% mdparagraph:P120:end

\end{minipage}\par\addvspace{4pt}

\par\begin{minipage}{\linewidth}\raggedright
% mdparagraph:P121:begin
\textbf{C5.}
% mdparagraph:P121:end

% item: NEW-U03-C05
% mdparagraph:P122:begin
Translate “The bird is dearer than the fish. The old dwelling is better than the great dwelling.”
% mdparagraph:P122:end

\end{minipage}\par\addvspace{4pt}

\par\begin{minipage}{\linewidth}\raggedright
% mdparagraph:P123:begin
\textbf{C6.}
% mdparagraph:P123:end

% item: NEW-U03-C06
% mdparagraph:P124:begin
Translate these nominative phrases: “the oldest tree”; “the dearest bird”; “the greatest work”; “the best man”; “the smallest star”.
% mdparagraph:P124:end

\end{minipage}\par\addvspace{4pt}

\par\begin{minipage}{\linewidth}\raggedright
% mdparagraph:P125:begin
\textbf{C7.}
% mdparagraph:P125:end

% item: NEW-U03-C07
% mdparagraph:P126:begin
Explain and correct the adjective ending in Se fugol is lǣss. Keep the meaning “The bird is smaller.” Then change the corrected sentence to refer to several birds.
% mdparagraph:P126:end

\end{minipage}\par\addvspace{4pt}

\reserve{7\baselineskip}
% mdparagraph:P127:begin
\section*{8. Adverbs and degree}\BookHeading{nav-ch03-learner-0023}{subsection}{8. Adverbs and degree}{1}\BookLabel{bk-ch03-learner-b99b7da0a77f}
% mdparagraph:P127:end

% teaching: NEW-U03-GRAM-adverb-formation
% teaching: NEW-U03-GRAM-adverb-use
% teaching: NEW-U03-GRAM-adverb-degree
% teaching: NEW-U03-GRAM-adverb-comparison
% teaching: NEW-U03-VOC-wide
% teaching: NEW-U03-VOC-wel
% teaching: NEW-U03-VOC-swithe
% teaching: NEW-U03-VOC-sothlice
% teaching: NEW-U03-VOC-wid-citation
% teaching: NEW-U03-VOC-swith-citation
% teaching: NEW-U03-VOC-soth-citation
% teaching: NEW-U03-PAR-wide
% teaching: NEW-U03-PAR-wel
% teaching: NEW-U03-PAR-swithe
% teaching: NEW-U03-PAR-sothlice
% mdparagraph:P128:begin
An adverb qualifies an action, a statement, an adjective or another adverb. It does not agree in case, number or gender. The selected adjective citations wīd “wide”, swīð “strong” and sōð “true” are supplied bases for this formation; their adjective declensions are not assigned.
% mdparagraph:P128:end

% mdtable:T18:begin
\par\begingroup
\fontsize{10.5}{13.2}\selectfont
\setlength{\LTleft}{0pt}\setlength{\LTright}{\fill}
\noindent
\begin{tabular}{@{}>{\raggedright\arraybackslash}p{24.43mm}>{\raggedright\arraybackslash}p{28.15mm}>{\raggedright\arraybackslash}p{31.0mm}>{\raggedright\arraybackslash}p{31.99mm}@{}}
\toprule
\textbf{Positive adverb} & \textbf{Meaning} & \textbf{Comparative} & \textbf{Superlative} \\
\midrule
wīde & widely; far and wide & wīdor, more widely & wīdost, most widely \\
wel & well & bęt, better & bętst, best \\
swīðe & very & — & — \\
sōðlice & truly; indeed & — & — \\
\bottomrule
\end{tabular}
\par\endgroup

% mdtable:T18:end
% mdparagraph:P129:begin
The\BookPoint{bk-ch03-learner-b9fb825d30c1} first derivation uses -e; sōðlice uses -lice. The regular adverbial comparative/\allowbreak{}superlative endings are -or/\allowbreak{}-ost. Memorize wel/\allowbreak{}bęt/\allowbreak{}bętst separately. An agreeing comparative adjective and the invariant adverb bęt perform different jobs. Swīðe can qualify an adjective. Sōðlice “truly” here is a new meaning distinct from the supplied connective use in the source reading.
% mdparagraph:P129:end

% example: NEW-U03-EXAMPLE-08
% mdparagraph:P130:begin
Hīe geworhton ðā hofu wel.
% mdparagraph:P130:end

% mdparagraph:P131:begin
They built the dwellings well.
% mdparagraph:P131:end

% mdparagraph:P132:begin
Wel qualifies geworhton; the dwellings remain an accusative object with their own morphology.
% mdparagraph:P132:end

\reserve{7\baselineskip}
% mdparagraph:P133:begin
\subsection*{Adverbs — V1–V4}\BookHeading{nav-ch03-learner-0024}{subsubsection}{Adverbs — V1–V4}{0}
% mdparagraph:P133:end

\par\begin{minipage}{\linewidth}\raggedright
% mdparagraph:P134:begin
\textbf{V1.}
% mdparagraph:P134:end

% item: NEW-U03-V01
% mdparagraph:P135:begin
Using the supplied adjective citation wīd, give the positive, comparative and superlative adverbs. Also form the adverb meaning “truly” from the supplied citation sōð. State which of these adverbs changes for gender, case or number.
% mdparagraph:P135:end

\end{minipage}\par\addvspace{4pt}

\par\begin{minipage}{\linewidth}\raggedright
% mdparagraph:P136:begin
\textbf{V2.}
% mdparagraph:P136:end

% item: NEW-U03-V02
% mdparagraph:P137:begin
Translate “The men built the good dwelling well.” Identify the adjective and the adverb, and state what each modifies.
% mdparagraph:P137:end

\end{minipage}\par\addvspace{4pt}

\par\begin{minipage}{\linewidth}\raggedright
% mdparagraph:P138:begin
\textbf{V3.}
% mdparagraph:P138:end

% item: NEW-U03-V03
% mdparagraph:P139:begin
Translate “He built the dwelling better.” Use the adverb meaning “better”. Then replace it with the adverb meaning “best”. Explain why neither word agrees with the dwelling.
% mdparagraph:P139:end

\end{minipage}\par\addvspace{4pt}

\par\begin{minipage}{\linewidth}\raggedright
% mdparagraph:P140:begin
\textbf{V4.}
% mdparagraph:P140:end

% item: NEW-U03-V04
% mdparagraph:P141:begin
Translate Sōðlice, ðæt weorc is swīðe gōd. Compare it with Ðæt weorc is gōd. What do the two adverbs add, and why do they have no adjective agreement endings?
% mdparagraph:P141:end

\end{minipage}\par\addvspace{4pt}

\reserve{7\baselineskip}
% mdparagraph:P142:begin
\section*{9. Cardinal numbers: one through twelve}\BookHeading{nav-ch03-learner-0025}{subsection}{9. Cardinal numbers: one through twelve}{1}\BookLabel{bk-ch03-learner-8a4d537a7515}
% mdparagraph:P142:end

% teaching: NEW-U03-GRAM-cardinal-one
% teaching: NEW-U03-GRAM-cardinal-two-three
% teaching: NEW-U03-GRAM-cardinal-four-twelve
% teaching: NEW-U03-GRAM-cardinal-counting
% teaching: NEW-U03-VOC-an
% teaching: NEW-U03-VOC-twegen
% teaching: NEW-U03-VOC-thrie
% teaching: NEW-U03-VOC-feower
% teaching: NEW-U03-VOC-fif
% teaching: NEW-U03-VOC-siex
% teaching: NEW-U03-VOC-seofon
% teaching: NEW-U03-VOC-eahta
% teaching: NEW-U03-VOC-nigon
% teaching: NEW-U03-VOC-tien
% teaching: NEW-U03-VOC-endlefan
% teaching: NEW-U03-VOC-twelf
% mdparagraph:P143:begin
Ān modifies a singular count noun. Two and three take plural count nouns and inflect for the noun’s gender and case. Four through twelve are used here in their ordinary invariant position before plural count nouns. The noun keeps the case required by its clause; counting does not supply a new noun dual. A dual personal subject still takes a plural verb, as in Unit 2.
% mdparagraph:P143:end

\reserve{7\baselineskip}
% mdparagraph:P144:begin
\subsection*{Ān}\BookHeading{nav-ch03-learner-0026}{subsubsection}{Ān}{0}\BookLabel{bk-ch03-learner-c138fbf9cb87}
% mdparagraph:P144:end

% teaching: NEW-U03-PAR-an-strong
% mdtable:T19:begin
% course-table-style: 2026-10
% course-panel:T19-P01:begin
\par\begingroup
\fontsize{10.5}{13.2}\selectfont
\setlength{\LTleft}{0pt}\setlength{\LTright}{\fill}
\noindent
\begin{tabular}{@{}>{\raggedright\arraybackslash}p{22.0mm}>{\raggedright\arraybackslash}p{31.188mm}>{\raggedright\arraybackslash}p{31.188mm}>{\raggedright\arraybackslash}p{31.188mm}@{}}
\toprule
\textbf{Case} & \textbf{Masculine singular} & \textbf{Neuter singular} & \textbf{Feminine singular} \\
\midrule
Nominative & ān & ān & ān \\
Genitive & ānes & ānes & ānre \\
Dative & ānum & ānum & ānre \\
Accusative & ānne & ān & āne \\
\bottomrule
\end{tabular}
\par\endgroup
% course-panel:T19-P01:end
\par\addvspace{5pt}
% course-panel:T19-P02:begin
\par\begingroup
\fontsize{10.5}{13.2}\selectfont
\setlength{\LTleft}{0pt}\setlength{\LTright}{\fill}
\noindent
\begin{tabular}{@{}>{\raggedright\arraybackslash}p{22.0mm}>{\raggedright\arraybackslash}p{31.188mm}>{\raggedright\arraybackslash}p{31.188mm}>{\raggedright\arraybackslash}p{31.188mm}@{}}
\toprule
\textbf{Case} & \textbf{Masculine plural: reference only} & \textbf{Neuter plural: reference only} & \textbf{Feminine plural: reference only} \\
\midrule
Nominative & āne & ān & āna \\
Genitive & ānra & ānra & ānra \\
Dative & ānum & ānum & ānum \\
Accusative & āne & ān & āna \\
\bottomrule
\end{tabular}
\par\endgroup
% course-panel:T19-P02:end
% mdtable:T19:end
% mdparagraph:P145:begin
Instrumental singular, recognition only: masculine āne; neuter āne.
% mdparagraph:P145:end

% mdparagraph:P146:begin
Only the singular is used productively for cardinal “one”. The plural table is for recognition/\allowbreak{}reference and does not authorize “one” with plural count nouns. The alternative masculine accusative singular ǣnne is valid beside selected ānne. A weak form meaning “alone” and independent numeral uses are not assigned.
% mdparagraph:P146:end

\reserve{7\baselineskip}
% mdparagraph:P147:begin
\subsection*{Two and three}\BookHeading{nav-ch03-learner-0027}{subsubsection}{Two and three}{0}\BookLabel{bk-ch03-learner-02a6416f7b36}
% mdparagraph:P147:end

% teaching: NEW-U03-PAR-twegen
% mdparagraph:P148:begin
\textbf{twēgen}
% mdparagraph:P148:end

% mdtable:T20:begin
% course-table-style: 2026-10
% course-panel:T20-P01:begin
\par\begingroup
\fontsize{10.5}{13.2}\selectfont
\setlength{\LTleft}{0pt}\setlength{\LTright}{\fill}
\noindent
\begin{tabular}{@{}>{\raggedright\arraybackslash}p{22.0mm}>{\raggedright\arraybackslash}p{31.188mm}>{\raggedright\arraybackslash}p{31.188mm}>{\raggedright\arraybackslash}p{31.188mm}@{}}
\toprule
\textbf{Case} & \textbf{Masculine} & \textbf{Neuter} & \textbf{Feminine} \\
\midrule
Nominative & twēgen & twā & twā \\
Genitive & twēgra & twēgra & twēgra \\
Dative & twǣm & twǣm & twǣm \\
Accusative & twēgen & twā & twā \\
\bottomrule
\end{tabular}
\par\endgroup
% course-panel:T20-P01:end
% mdtable:T20:end
% teaching: NEW-U03-PAR-thrie
% mdparagraph:P149:begin
\textbf{ðrīe}
% mdparagraph:P149:end

% mdtable:T21:begin
% course-table-style: 2026-10
% course-panel:T21-P01:begin
\par\begingroup
\fontsize{10.5}{13.2}\selectfont
\setlength{\LTleft}{0pt}\setlength{\LTright}{\fill}
\noindent
\begin{tabular}{@{}>{\raggedright\arraybackslash}p{22.0mm}>{\raggedright\arraybackslash}p{31.188mm}>{\raggedright\arraybackslash}p{31.188mm}>{\raggedright\arraybackslash}p{31.188mm}@{}}
\toprule
\textbf{Case} & \textbf{Masculine} & \textbf{Neuter} & \textbf{Feminine} \\
\midrule
Nominative & ðrīe & ðrēo & ðrēo \\
Genitive & ðrēora & ðrēora & ðrēora \\
Dative & ðrīm & ðrīm & ðrīm \\
Accusative & ðrīe & ðrēo & ðrēo \\
\bottomrule
\end{tabular}
\par\endgroup
% course-panel:T21-P01:end
% mdtable:T21:end
% mdparagraph:P150:begin
All three genders share each numeral’s genitive and dative. Nominative and accusative plural distinguish the selected gender forms. A recipient still takes the dative numeral and dative noun.
% mdparagraph:P150:end

% example: NEW-U03-EXAMPLE-09
% mdparagraph:P151:begin
Wit habbað twā trēowu.
% mdparagraph:P151:end

% mdparagraph:P152:begin
We two have two trees.
% mdparagraph:P152:end

% mdparagraph:P153:begin
Wit identifies the speaker and one other; twā counts the neuter plural trees. These are separate grammatical choices.
% mdparagraph:P153:end

% teaching: NEW-U03-PAR-feower
% teaching: NEW-U03-PAR-fif
% teaching: NEW-U03-PAR-siex
% teaching: NEW-U03-PAR-seofon
% teaching: NEW-U03-PAR-eahta
% teaching: NEW-U03-PAR-nigon
% teaching: NEW-U03-PAR-tien
% teaching: NEW-U03-PAR-endlefan
% teaching: NEW-U03-PAR-twelf
% mdtable:T22:begin
\par\begingroup
\fontsize{10.5}{13.2}\selectfont
\setlength{\LTleft}{0pt}\setlength{\LTright}{\fill}
\noindent
\begin{tabular}{@{}>{\raggedright\arraybackslash}p{54.96mm}>{\raggedright\arraybackslash}p{66.23mm}@{}}
\toprule
\textbf{Number} & \textbf{Selected cardinal} \\
\midrule
4 & fēower \\
5 & fīf \\
6 & siex \\
7 & seofon \\
8 & eahta \\
9 & nigon \\
10 & tīen \\
11 & endlefan \\
12 & twęlf \\
\bottomrule
\end{tabular}
\par\endgroup

% mdtable:T22:end
% mdparagraph:P154:begin
Use\BookPoint{bk-ch03-learner-2aa77dbd105d} these nine forms without case endings in the stated pre-nominal counting construction. Higher numbers, partitive genitives, and the inflection of postposed or substantival numerals are outside the present tasks.
% mdparagraph:P154:end

\reserve{7\baselineskip}
% mdparagraph:P155:begin
\subsection*{Cardinals — N1–N5}\BookHeading{nav-ch03-learner-0028}{subsubsection}{Cardinals — N1–N5}{0}
% mdparagraph:P155:end

\par\begin{minipage}{\linewidth}\raggedright
% mdparagraph:P156:begin
\textbf{N1.}
% mdparagraph:P156:end

% item: NEW-U03-N01
% mdparagraph:P157:begin
Give the nominative, genitive, dative and accusative singular of ān in all three genders. Use the selected cardinal forms.
% mdparagraph:P157:end

\end{minipage}\par\addvspace{4pt}

\par\begin{minipage}{\linewidth}\raggedright
% mdparagraph:P158:begin
\textbf{N2.}
% mdparagraph:P158:end

% item: NEW-U03-N02
% mdparagraph:P159:begin
Give the nominative, genitive, dative and accusative of twēgen and ðrīe in all three genders. State which number their accompanying count nouns have.
% mdparagraph:P159:end

\end{minipage}\par\addvspace{4pt}

\par\begin{minipage}{\linewidth}\raggedright
% mdparagraph:P160:begin
\textbf{N3.}
% mdparagraph:P160:end

% item: NEW-U03-N03
% mdparagraph:P161:begin
Translate these nominative phrases: “two men”; “two herbs”; “two dwellings”; “three birds”; “three nights”; “three trees”.
% mdparagraph:P161:end

\end{minipage}\par\addvspace{4pt}

\par\begin{minipage}{\linewidth}\raggedright
% mdparagraph:P162:begin
\textbf{N4.}
% mdparagraph:P162:end

% item: NEW-U03-N04
% mdparagraph:P163:begin
Translate “We gave seed to two men. They gave water to three birds.”
% mdparagraph:P163:end

\end{minipage}\par\addvspace{4pt}

\par\begin{minipage}{\linewidth}\raggedright
% mdparagraph:P164:begin
\textbf{N5.}
% mdparagraph:P164:end

% item: NEW-U03-N05
% mdparagraph:P165:begin
Translate these nominative phrases: “four fish”; “five trees”; “six stars”; “seven days”; “eight birds”; “nine men”; “ten herbs”; “eleven years”; “twelve signs”. Explain the number of the nouns and the selected treatment of the cardinals.
% mdparagraph:P165:end

\end{minipage}\par\addvspace{4pt}

\reserve{7\baselineskip}
% mdparagraph:P166:begin
\section*{10. Ordinals: first through seventh}\BookHeading{nav-ch03-learner-0029}{subsection}{10. Ordinals: first through seventh}{1}\BookLabel{bk-ch03-learner-482b1a79fb62}
% mdparagraph:P166:end

% teaching: NEW-U03-GRAM-ordinals
% teaching: NEW-U03-GRAM-ordinal-other
% teaching: NEW-U03-VOC-aeresta
% teaching: NEW-U03-VOC-other
% teaching: NEW-U03-VOC-thridda
% teaching: NEW-U03-VOC-feortha
% teaching: NEW-U03-VOC-fifta
% teaching: NEW-U03-VOC-siexta
% teaching: NEW-U03-VOC-seofotha
% mdtable:T23:begin
\par\begingroup
\fontsize{10.5}{13.2}\selectfont
\setlength{\LTleft}{0pt}\setlength{\LTright}{\fill}
\noindent
\begin{tabular}{@{}>{\raggedright\arraybackslash}p{24.69mm}>{\raggedright\arraybackslash}p{42.22mm}>{\raggedright\arraybackslash}p{51.47mm}@{}}
\toprule
\textbf{Order} & \textbf{Selected masculine nom. sg.} & \textbf{Declension/\allowbreak{}stem} \\
\midrule
1 & ǣresta & weak; ǣrest- \\
2 & ōðer & strong, special ōðer pattern \\
3 & ðridda & weak; ðridd- \\
4 & fēorða & weak; fēorð- \\
5 & fīfta & weak; fīft- \\
6 & siexta & weak; siext- \\
7 & seofoða & weak; seofoð- \\
\bottomrule
\end{tabular}
\par\endgroup

% mdtable:T23:end
% mdparagraph:P167:begin
Ordinals express position in an order. First and third through seventh are weak. The second, ōðer, remains strong even after the article. The case still follows the noun’s role; these ordinal forms do not introduce a new temporal use of on.
% mdparagraph:P167:end

\reserve{7\baselineskip}
% mdparagraph:P168:begin
\subsection*{Complete ordinary ordinal pattern}\BookHeading{nav-ch03-learner-0030}{subsubsection}{Complete ordinary ordinal pattern}{0}\BookLabel{bk-ch03-learner-f4e2c8b803ee}
% mdparagraph:P168:end

% teaching: NEW-U03-PAR-thridda-weak
% mdtable:T24:begin
% course-table-style: 2026-10
% course-panel:T24-P01:begin
\par\begingroup
\fontsize{10.5}{13.2}\selectfont
\setlength{\LTleft}{0pt}\setlength{\LTright}{\fill}
\noindent
\begin{tabular}{@{}>{\raggedright\arraybackslash}p{22.0mm}>{\raggedright\arraybackslash}p{31.188mm}>{\raggedright\arraybackslash}p{31.188mm}>{\raggedright\arraybackslash}p{31.188mm}@{}}
\toprule
\textbf{Case} & \textbf{Masculine singular} & \textbf{Neuter singular} & \textbf{Feminine singular} \\
\midrule
Nominative & ðridda & ðridde & ðridde \\
Genitive & ðriddan & ðriddan & ðriddan \\
Dative & ðriddan & ðriddan & ðriddan \\
Accusative & ðriddan & ðridde & ðriddan \\
\bottomrule
\end{tabular}
\par\endgroup
% course-panel:T24-P01:end
\par\addvspace{5pt}
% course-panel:T24-P02:begin
\par\begingroup
\fontsize{10.5}{13.2}\selectfont
\setlength{\LTleft}{0pt}\setlength{\LTright}{\fill}
\noindent
\begin{tabular}{@{}>{\raggedright\arraybackslash}p{22.0mm}>{\raggedright\arraybackslash}p{31.188mm}>{\raggedright\arraybackslash}p{31.188mm}>{\raggedright\arraybackslash}p{31.188mm}@{}}
\toprule
\textbf{Case} & \textbf{Masculine plural} & \textbf{Neuter plural} & \textbf{Feminine plural} \\
\midrule
Nominative & ðriddan & ðriddan & ðriddan \\
Genitive & ðridra & ðridra & ðridra \\
Dative & ðriddum & ðriddum & ðriddum \\
Accusative & ðriddan & ðriddan & ðriddan \\
\bottomrule
\end{tabular}
\par\endgroup
% course-panel:T24-P02:end
% mdtable:T24:end
% teaching: NEW-U03-PAR-aeresta-weak
% teaching: NEW-U03-PAR-feortha-weak
% teaching: NEW-U03-PAR-fifta-weak
% teaching: NEW-U03-PAR-siexta-weak
% teaching: NEW-U03-PAR-seofotha-weak
% mdparagraph:P169:begin
For first and fourth through seventh, remove masculine nominative -a to obtain the printed stem, then apply every weak gōd ending, including genitive plural -ra and dative plural -um. Ðridda has the exceptional genitive plural ðridra, with one d before r; retain dd in its other cells shown above. No strong alternative is substituted merely because an ordinal lacks an article.
% mdparagraph:P169:end

\reserve{7\baselineskip}
% mdparagraph:P170:begin
\subsection*{Ōðer}\BookHeading{nav-ch03-learner-0031}{subsubsection}{Ōðer}{0}\BookLabel{bk-ch03-learner-d9d0ac45cb4c}
% mdparagraph:P170:end

% teaching: NEW-U03-PAR-other-strong
% mdtable:T25:begin
% course-table-style: 2026-10
% course-panel:T25-P01:begin
\par\begingroup
\fontsize{10.5}{13.2}\selectfont
\setlength{\LTleft}{0pt}\setlength{\LTright}{\fill}
\noindent
\begin{tabular}{@{}>{\raggedright\arraybackslash}p{22.0mm}>{\raggedright\arraybackslash}p{31.188mm}>{\raggedright\arraybackslash}p{31.188mm}>{\raggedright\arraybackslash}p{31.188mm}@{}}
\toprule
\textbf{Case} & \textbf{Masculine singular} & \textbf{Neuter singular} & \textbf{Feminine singular} \\
\midrule
Nominative & ōðer & ōðer & ōðru \\
Genitive & ōðres & ōðres & ōðerre \\
Dative & ōðrum & ōðrum & ōðerre \\
Accusative & ōðerne & ōðer & ōðre \\
\bottomrule
\end{tabular}
\par\endgroup
% course-panel:T25-P01:end
\par\addvspace{5pt}
% course-panel:T25-P02:begin
\par\begingroup
\fontsize{10.5}{13.2}\selectfont
\setlength{\LTleft}{0pt}\setlength{\LTright}{\fill}
\noindent
\begin{tabular}{@{}>{\raggedright\arraybackslash}p{22.0mm}>{\raggedright\arraybackslash}p{31.188mm}>{\raggedright\arraybackslash}p{31.188mm}>{\raggedright\arraybackslash}p{31.188mm}@{}}
\toprule
\textbf{Case} & \textbf{Masculine plural} & \textbf{Neuter plural} & \textbf{Feminine plural} \\
\midrule
Nominative & ōðre & ōðru & ōðra \\
Genitive & ōðerra & ōðerra & ōðerra \\
Dative & ōðrum & ōðrum & ōðrum \\
Accusative & ōðre & ōðru & ōðra \\
\bottomrule
\end{tabular}
\par\endgroup
% course-panel:T25-P02:end
% mdtable:T25:end
% mdparagraph:P171:begin
Instrumental singular, recognition only: masculine ōðre; neuter ōðre.
% mdparagraph:P171:end

% mdparagraph:P172:begin
The selected feminine nominative singular is ōðru; ōðer and ōðeru are supported alternatives. Preserve ōðr- before the selected vowel endings, but ōðer- before -ne, -re and -ra. The article does not change this adjective to the weak pattern.
% mdparagraph:P172:end

% example: NEW-U03-EXAMPLE-10
% mdparagraph:P173:begin
Ūre mann is on ðǣm ōðrum hofe.
% mdparagraph:P173:end

% mdparagraph:P174:begin
Our man is in the second dwelling.
% mdparagraph:P174:end

% mdparagraph:P175:begin
Ōðrum is strong neuter dative singular after the article; hofe is dative after locative on.
% mdparagraph:P175:end

\reserve{7\baselineskip}
% mdparagraph:P176:begin
\subsection*{Ordinals — O1–O3}\BookHeading{nav-ch03-learner-0032}{subsubsection}{Ordinals — O1–O3}{0}
% mdparagraph:P176:end

\par\begin{minipage}{\linewidth}\raggedright
% mdparagraph:P177:begin
\textbf{O1.}
% mdparagraph:P177:end

% item: NEW-U03-O01
% mdparagraph:P178:begin
Give the nominative and genitive singular of “the second day”, “the second night” and “the second year”. Explain the declension of ōðer after the article.
% mdparagraph:P178:end

\end{minipage}\par\addvspace{4pt}

\par\begin{minipage}{\linewidth}\raggedright
% mdparagraph:P179:begin
\textbf{O2.}
% mdparagraph:P179:end

% item: NEW-U03-O02
% mdparagraph:P180:begin
Give the complete weak nominative, genitive, dative and accusative of ðridda in all three genders, singular and plural.
% mdparagraph:P180:end

\end{minipage}\par\addvspace{4pt}

\par\begin{minipage}{\linewidth}\raggedright
% mdparagraph:P181:begin
\textbf{O3.}
% mdparagraph:P181:end

% item: NEW-U03-O03
% mdparagraph:P182:begin
Translate “Three birds live in the first dwelling. The old man gave food to the fifth bird.”
% mdparagraph:P182:end

\end{minipage}\par\addvspace{4pt}

\reserve{7\baselineskip}
% mdparagraph:P183:begin
\section*{11. Demonstratives}\BookHeading{nav-ch03-learner-0033}{subsection}{11. Demonstratives}{1}\BookLabel{bk-ch03-learner-c532ec2b751d}
% mdparagraph:P183:end

% teaching: NEW-U03-GRAM-near-demonstrative
% teaching: NEW-U03-GRAM-far-demonstrative
% teaching: NEW-U03-VOC-se
% teaching: NEW-U03-VOC-thes
% mdparagraph:P184:begin
The article forms se, sēo, ðæt can also point to “that” or “those”. Their forms are familiar, but this pointing meaning is now assigned. Ðes, ðēos, ðis means “this”; its plural means “these”. A demonstrative agrees with its head noun. An ordinary adjective following it is weak. Relative uses and general third-person pronoun uses of these words are not introduced.
% mdparagraph:P184:end

\reserve{7\baselineskip}
% mdparagraph:P185:begin
\subsection*{Se: that/\allowbreak{}those; also the familiar article}\BookHeading{nav-ch03-learner-0034}{subsubsection}{Se: that/those; also the familiar article}{0}\BookLabel{bk-ch03-learner-59601dfdb337}
% mdparagraph:P185:end

% teaching: NEW-U03-PAR-se-demonstrative
% mdtable:T26:begin
% course-table-style: 2026-10
% course-panel:T26-P01:begin
\par\begingroup
\fontsize{10.5}{13.2}\selectfont
\setlength{\LTleft}{0pt}\setlength{\LTright}{\fill}
\noindent
\begin{tabular}{@{}>{\raggedright\arraybackslash}p{22.0mm}>{\raggedright\arraybackslash}p{31.188mm}>{\raggedright\arraybackslash}p{31.188mm}>{\raggedright\arraybackslash}p{31.188mm}@{}}
\toprule
\textbf{Case} & \textbf{Masculine singular} & \textbf{Neuter singular} & \textbf{Feminine singular} \\
\midrule
Nominative & se & ðæt & sēo \\
Genitive & ðæs & ðæs & ðǣre \\
Dative & ðǣm & ðǣm & ðǣre \\
Accusative & ðone & ðæt & ðā \\
\bottomrule
\end{tabular}
\par\endgroup
% course-panel:T26-P01:end
\par\addvspace{5pt}
% course-panel:T26-P02:begin
\par\begingroup
\fontsize{10.5}{13.2}\selectfont
\setlength{\LTleft}{0pt}\setlength{\LTright}{\fill}
\noindent
\begin{tabular}{@{}>{\raggedright\arraybackslash}p{22.0mm}>{\raggedright\arraybackslash}p{31.188mm}>{\raggedright\arraybackslash}p{31.188mm}>{\raggedright\arraybackslash}p{31.188mm}@{}}
\toprule
\textbf{Case} & \textbf{Masculine plural} & \textbf{Neuter plural} & \textbf{Feminine plural} \\
\midrule
Nominative & ðā & ðā & ðā \\
Genitive & ðāra & ðāra & ðāra \\
Dative & ðǣm & ðǣm & ðǣm \\
Accusative & ðā & ðā & ðā \\
\bottomrule
\end{tabular}
\par\endgroup
% course-panel:T26-P02:end
% mdtable:T26:end
% mdparagraph:P186:begin
Instrumental singular, recognition only: masculine ðȳ; neuter ðȳ.
% mdparagraph:P186:end

% mdparagraph:P187:begin
The table repeats the existing article morphology. The instrumental alternatives ðȳ and ðon remain recognition forms.
% mdparagraph:P187:end

\reserve{7\baselineskip}
% mdparagraph:P188:begin
\subsection*{Ðes: this/\allowbreak{}these}\BookHeading{nav-ch03-learner-0035}{subsubsection}{Ðes: this/these}{0}\BookLabel{bk-ch03-learner-78aeff8dce51}
% mdparagraph:P188:end

% teaching: NEW-U03-PAR-thes-demonstrative
% mdtable:T27:begin
% course-table-style: 2026-10
% course-panel:T27-P01:begin
\par\begingroup
\fontsize{10.5}{13.2}\selectfont
\setlength{\LTleft}{0pt}\setlength{\LTright}{\fill}
\noindent
\begin{tabular}{@{}>{\raggedright\arraybackslash}p{22.0mm}>{\raggedright\arraybackslash}p{31.188mm}>{\raggedright\arraybackslash}p{31.188mm}>{\raggedright\arraybackslash}p{31.188mm}@{}}
\toprule
\textbf{Case} & \textbf{Masculine singular} & \textbf{Neuter singular} & \textbf{Feminine singular} \\
\midrule
Nominative & ðes & ðis & ðēos \\
Genitive & ðisses & ðisses & ðisse \\
Dative & ðissum & ðissum & ðisse \\
Accusative & ðisne & ðis & ðās \\
\bottomrule
\end{tabular}
\par\endgroup
% course-panel:T27-P01:end
\par\addvspace{5pt}
% course-panel:T27-P02:begin
\par\begingroup
\fontsize{10.5}{13.2}\selectfont
\setlength{\LTleft}{0pt}\setlength{\LTright}{\fill}
\noindent
\begin{tabular}{@{}>{\raggedright\arraybackslash}p{22.0mm}>{\raggedright\arraybackslash}p{31.188mm}>{\raggedright\arraybackslash}p{31.188mm}>{\raggedright\arraybackslash}p{31.188mm}@{}}
\toprule
\textbf{Case} & \textbf{Masculine plural} & \textbf{Neuter plural} & \textbf{Feminine plural} \\
\midrule
Nominative & ðās & ðās & ðās \\
Genitive & ðissa & ðissa & ðissa \\
Dative & ðissum & ðissum & ðissum \\
Accusative & ðās & ðās & ðās \\
\bottomrule
\end{tabular}
\par\endgroup
% course-panel:T27-P02:end
% mdtable:T27:end
% mdparagraph:P189:begin
Instrumental singular, recognition only: masculine ðȳs; neuter ðȳs.
% mdparagraph:P189:end

% mdparagraph:P190:begin
Keep Cook’s short-e spelling ðes. Do not choose a form from the English word alone: a recipient may require a dative singular demonstrative, while an object may require an accusative plural one.
% mdparagraph:P190:end

% example: NEW-U03-EXAMPLE-06
% mdparagraph:P191:begin
Wē lufiað ðis lȳtle dēor.
% mdparagraph:P191:end

% mdparagraph:P192:begin
We love this small animal.
% mdparagraph:P192:end

% mdparagraph:P193:begin
Ðis and lȳtle are neuter accusative singular with the object dēor; the adjective is weak after the demonstrative.
% mdparagraph:P193:end

\reserve{7\baselineskip}
% mdparagraph:P194:begin
\subsection*{Demonstratives — D1–D4}\BookHeading{nav-ch03-learner-0036}{subsubsection}{Demonstratives — D1–D4}{0}
% mdparagraph:P194:end

\par\begin{minipage}{\linewidth}\raggedright
% mdparagraph:P195:begin
\textbf{D1.}
% mdparagraph:P195:end

% item: NEW-U03-D01
% mdparagraph:P196:begin
Give the complete nominative, genitive, dative and accusative of ðes in all three singular genders and in the plural.
% mdparagraph:P196:end

\end{minipage}\par\addvspace{4pt}

\par\begin{minipage}{\linewidth}\raggedright
% mdparagraph:P197:begin
\textbf{D2.}
% mdparagraph:P197:end

% item: NEW-U03-D02
% mdparagraph:P198:begin
Translate “This tree is great. That tree is small. These birds live in that dwelling.”
% mdparagraph:P198:end

\end{minipage}\par\addvspace{4pt}

\par\begin{minipage}{\linewidth}\raggedright
% mdparagraph:P199:begin
\textbf{D3.}
% mdparagraph:P199:end

% item: NEW-U03-D03
% mdparagraph:P200:begin
Translate “I give seed to this man. I hear these birds.”
% mdparagraph:P200:end

\end{minipage}\par\addvspace{4pt}

\par\begin{minipage}{\linewidth}\raggedright
% mdparagraph:P201:begin
\textbf{D4.}
% mdparagraph:P201:end

% item: NEW-U03-D04
% mdparagraph:P202:begin
Translate these nominative phrases: “this good man”; “these great trees”; “good men”. Explain the declension of each adjective.
% mdparagraph:P202:end

\end{minipage}\par\addvspace{4pt}

\reserve{7\baselineskip}
% mdparagraph:P203:begin
\section*{12. Cook, Group I, Exercise 4}\BookHeading{nav-ch03-learner-0037}{subsection}{12. Cook, Group I, Exercise 4}{1}
% mdparagraph:P203:end

% mdparagraph:P204:begin
Translate items 1–17. For a noun-phrase fragment, use a nominative head unless a task explicitly gives another case. Retain the dependent genitive’s article and place that dependent before its head, as in the earlier Cook exercises. Complete sentences require the case appropriate to each clause role. Use the selected forms when practising a paradigm; the key distinguishes those targets from defensible alternatives.
% mdparagraph:P204:end

% mdparagraph:P205:begin
Cook’s introductory note reads: “[NOTE. — For adjectives preceded by a demonstrative, see 55.]” The reference is to his §55. Item 12 retains Cook’s literal reference “(80)”; it points to the ordinal discussion.
% mdparagraph:P205:end

\par\begin{minipage}{\linewidth}\raggedright
% item: EX-I-04-01
% mdparagraph:P206:begin
\itempara{1}{The great whales.}
% mdparagraph:P206:end

\end{minipage}\par\addvspace{4pt}

\par\begin{minipage}{\linewidth}\raggedright
% item: EX-I-04-02
% mdparagraph:P207:begin
\itempara{2}{The earth was empty and void.}
% mdparagraph:P207:end

\end{minipage}\par\addvspace{4pt}

\par\begin{minipage}{\linewidth}\raggedright
% item: EX-I-04-03
% mdparagraph:P208:begin
\itempara{3}{The things were good.}
% mdparagraph:P208:end

\end{minipage}\par\addvspace{4pt}

\par\begin{minipage}{\linewidth}\raggedright
% item: EX-I-04-04
% mdparagraph:P209:begin
\itempara{4}{The fruit-bearing trees.}
% mdparagraph:P209:end

\end{minipage}\par\addvspace{4pt}

\par\begin{minipage}{\linewidth}\raggedright
% item: EX-I-04-05
% mdparagraph:P210:begin
\itempara{5}{The great lights.}
% mdparagraph:P210:end

\end{minipage}\par\addvspace{4pt}

\par\begin{minipage}{\linewidth}\raggedright
% item: EX-I-04-06
% mdparagraph:P211:begin
\itempara{6}{The little star.}
% mdparagraph:P211:end

\end{minipage}\par\addvspace{4pt}

\par\begin{minipage}{\linewidth}\raggedright
% item: EX-I-04-07
% mdparagraph:P212:begin
\itempara{7}{The smaller bird.}
% mdparagraph:P212:end

\end{minipage}\par\addvspace{4pt}

\par\begin{minipage}{\linewidth}\raggedright
% item: EX-I-04-08
% mdparagraph:P213:begin
\itempara{8}{The larger tree.}
% mdparagraph:P213:end

\end{minipage}\par\addvspace{4pt}

\par\begin{minipage}{\linewidth}\raggedright
% item: EX-I-04-09
% mdparagraph:P214:begin
\itempara{9}{The life of the good man.}
% mdparagraph:P214:end

\end{minipage}\par\addvspace{4pt}

\par\begin{minipage}{\linewidth}\raggedright
% item: EX-I-04-10
% mdparagraph:P215:begin
\itempara{10}{The beginning of the seventh night.}
% mdparagraph:P215:end

\end{minipage}\par\addvspace{4pt}

\par\begin{minipage}{\linewidth}\raggedright
% item: EX-I-04-11
% mdparagraph:P216:begin
\itempara{11}{The waters of the great sea.}
% mdparagraph:P216:end

\end{minipage}\par\addvspace{4pt}

\par\begin{minipage}{\linewidth}\raggedright
% item: EX-I-04-12
% mdparagraph:P217:begin
\itempara{12}{The second (80) day and the second night.}
% mdparagraph:P217:end

\end{minipage}\par\addvspace{4pt}

\par\begin{minipage}{\linewidth}\raggedright
% item: EX-I-04-13
% mdparagraph:P218:begin
\itempara{13}{The third day of the second year.}
% mdparagraph:P218:end

\end{minipage}\par\addvspace{4pt}

\par\begin{minipage}{\linewidth}\raggedright
% item: EX-I-04-14
% mdparagraph:P219:begin
\itempara{14}{The food of the fifth man.}
% mdparagraph:P219:end

\end{minipage}\par\addvspace{4pt}

\par\begin{minipage}{\linewidth}\raggedright
% item: EX-I-04-15
% mdparagraph:P220:begin
\itempara{15}{The light of the sixth evening.}
% mdparagraph:P220:end

\end{minipage}\par\addvspace{4pt}

\par\begin{minipage}{\linewidth}\raggedright
% item: EX-I-04-16
% mdparagraph:P221:begin
\itempara{16}{The works of the moving creatures.}
% mdparagraph:P221:end

\end{minipage}\par\addvspace{4pt}

\par\begin{minipage}{\linewidth}\raggedright
% item: EX-I-04-17
% mdparagraph:P222:begin
\itempara{17}{The morning of the first day.}
% mdparagraph:P222:end

% deferred: EX-I-04-18
% mdparagraph:P223:begin
Cook I.4.18, “Increase (sing.) and be multiplied.”, is reserved for Unit 7, where its command and passive construction can be taught. It is not an additional task in this unit.
% mdparagraph:P223:end

\end{minipage}\par\addvspace{4pt}

\reserve{7\baselineskip}
% mdparagraph:P224:begin
\section*{13. Controlled reading: Two dwellings and three birds}\BookHeading{nav-ch03-learner-0038}{subsection}{13. Controlled reading: Two dwellings and three birds}{1}
% mdparagraph:P224:end

% mdparagraph:P225:begin
This connected account is newly composed teaching material. The old man is one of the builders. In his closing speech, ‘our’ includes all the builders.
% mdparagraph:P225:end

% composition: NEW-U03-TEXT-DWELLINGS-01
% mdparagraph:P226:begin
\par\begin{minipage}{\linewidth}\raggedright
\readingtext{Ðā gōdan męnn geworhton twā hofu. Ān hof wæs micel, and ān hof wæs lȳtel. Ðæt micle hof wæs māre ðonne ðæt lȳtle hof. Hīe sealdon ðæt micle hof ðǣm ealdan męnn. Hē hæfde ðrīe fuglas. Hē sealde his fuglum sǣd and wæter. Ðā fuglas lifdon on ðǣm miclan hofe. Se mann sægde: “Mīne fuglas sindon gōde. Ūre weorc is gōd.”}
\end{minipage}\par
% mdparagraph:P226:end

% mdparagraph:P227:begin
All words and constructions in this account have been taught. Track the case of each noun phrase before deciding who builds, gives, possesses or lives somewhere. Inside the direct quotation, the speaker’s first-person possessives refer to him and to the group identified in the orientation.
% mdparagraph:P227:end

\reserve{7\baselineskip}
% mdparagraph:P228:begin
\subsection*{Controlled reading — L1–L5}\BookHeading{nav-ch03-learner-0039}{subsubsection}{Controlled reading — L1–L5}{0}
% mdparagraph:P228:end

\par\begin{minipage}{\linewidth}\raggedright
% mdparagraph:P229:begin
\textbf{L1.}
% mdparagraph:P229:end

% item: NEW-U03-L01
% mdparagraph:P230:begin
Translate the complete new account “Two dwellings and three birds”. Keep the givers, recipient and possessed things distinct.
% mdparagraph:P230:end

\end{minipage}\par\addvspace{4pt}

\par\begin{minipage}{\linewidth}\raggedright
% mdparagraph:P231:begin
\textbf{L2.}
% mdparagraph:P231:end

% item: NEW-U03-L02
% mdparagraph:P232:begin
List the things counted by twā, ān and ðrīe in the account. Identify the gender, number and case of each numeral in its clause.
% mdparagraph:P232:end

\end{minipage}\par\addvspace{4pt}

\par\begin{minipage}{\linewidth}\raggedright
% mdparagraph:P233:begin
\textbf{L3.}
% mdparagraph:P233:end

% item: NEW-U03-L03
% mdparagraph:P234:begin
In the closing quotation, identify the possessor and the possessed noun in Mīne fuglas and Ūre weorc. Explain why the same final -e occurs with different gender and number analyses.
% mdparagraph:P234:end

\end{minipage}\par\addvspace{4pt}

\par\begin{minipage}{\linewidth}\raggedright
% mdparagraph:P235:begin
\textbf{L4.}
% mdparagraph:P235:end

% item: NEW-U03-L04
% mdparagraph:P236:begin
In the account, identify the strong predicate adjectives, the weak attributive adjectives and the comparative. Give the noun or pronoun governing each agreement.
% mdparagraph:P236:end

\end{minipage}\par\addvspace{4pt}

\par\begin{minipage}{\linewidth}\raggedright
% mdparagraph:P237:begin
\textbf{L5.}
% mdparagraph:P237:end

% item: NEW-U03-L05
% mdparagraph:P238:begin
Retell the first four sentences of the account in the present. Preserve all people, things, cases and adjective endings.
% mdparagraph:P238:end

\end{minipage}\par\addvspace{4pt}

\reserve{7\baselineskip}
% mdparagraph:P239:begin
\section*{14. Source reading: the fourth day}\BookHeading{nav-ch03-learner-0040}{subsection}{14. Source reading: the fourth day}{1}
% mdparagraph:P239:end

% mdparagraph:P240:begin
This is the complete fourth-day account from Cook’s first reading. Read the local help before translating. Its wishes, purpose clauses and resultative formulae are supplied; the adjective, numeral, noun and familiar indicative forms can be analysed with the systems learned so far. Do not infer that every form of a familiar verb has been taught.
% mdparagraph:P240:end

% source: FB-I-CREATION-FOURTH-DAY-01
% mdparagraph:P241:begin
\par\begin{minipage}{\linewidth}\raggedright
\readingtext{God cwæð ðā sōðlice, “Bēon nū lēoht on ðǣre heofonan fæstnesse, and tōdǣlen dæg and niht, and bēon tō tācnum, and tō tīdum, and tō dagum, and tō gēarum. And hīe scīnen on ðǣre heofonan fæstnesse, and ālīehten ðā eorðan”; hit wæs ðā swā geworden. And God geworhte twā miclu lēoht: ðæt māre lēoht tō ðæs dæges līehtinge, and ðæt lǣsse lēoht tō ðǣre niht līehtinge; and steorran hē geworhte. And gesętte hīe on ðǣre heofonan, ðæt hīe scinen ofer eorðan, and gīemden ðæs dæges and ðǣre niht, and tōdǣlden lēoht and ðīestru; God geseah ðā ðæt hit gōd wæs. And wæs geworden ǣfen and męrgen se fēorða dæg.}
\end{minipage}\par
% mdparagraph:P241:end

\reserve{7\baselineskip}
% mdparagraph:P242:begin
\subsection*{Local forms and meanings}\BookHeading{nav-ch03-learner-0041}{subsubsection}{Local forms and meanings}{0}
% mdparagraph:P242:end

% teaching: NEW-U03-SUPPORT-cwethan
% teaching: NEW-U03-SUPPORT-tha-adverb
% teaching: NEW-U03-SUPPORT-sothlice
% teaching: NEW-U03-SUPPORT-wesan
% teaching: NEW-U03-SUPPORT-nu
% teaching: NEW-U03-SUPPORT-heofone
% teaching: NEW-U03-SUPPORT-todaelan
% teaching: NEW-U03-SUPPORT-to
% teaching: NEW-U03-SUPPORT-scinan
% teaching: NEW-U03-SUPPORT-aliehtan
% teaching: NEW-U03-SUPPORT-swa
% teaching: NEW-U03-SUPPORT-geweorthan
% teaching: NEW-U03-SUPPORT-gesettan
% teaching: NEW-U03-SUPPORT-thaet-conjunction
% teaching: NEW-U03-SUPPORT-ofer
% teaching: NEW-U03-SUPPORT-gieman
% teaching: NEW-U03-SUPPORT-geseon
% teaching: NEW-U03-SUPPORT-morgen
% mdtable:T28:begin
\par\begingroup
\fontsize{10.5}{13.2}\selectfont
\setlength{\LTleft}{0pt}\setlength{\LTright}{\fill}
\noindent
\begin{tabular}{@{}>{\raggedright\arraybackslash}p{29.38mm}>{\raggedright\arraybackslash}p{91.81mm}@{}}
\toprule
\textbf{Source expression} & \textbf{Local help} \\
\midrule
cwæð & said; supplied third-person singular past of cweðan. The productive sęcgan paradigm does not supply this form. \\
ðā; sōðlice; nū & then; connective sōðlice (Cook compares Latin autem); now. The connective is distinct from the adverb “truly” learned above. \\
Bēon /\allowbreak{} bēon & let there be /\allowbreak{} let them be: plural present optative forms, expressing a wish or command. They are not infinitives or the indicative sindon. \\
heofonan & weak feminine heofone “heaven”: genitive singular before fæstnesse, but dative singular in on ðǣre heofonan. This is not the strong masculine heofon paradigm. \\
tōdǣlen; scīnen; ālīehten & let them separate; let them shine; let them illuminate. Supplied plural present optatives with the lights as subject. \\
tō tācnum /\allowbreak{} tīdum /\allowbreak{} dagum /\allowbreak{} gēarum & for signs /\allowbreak{} seasons /\allowbreak{} days /\allowbreak{} years. Tō governs these datives in a supplied function or purpose construction. \\
swā; geworden & thus, so; come to pass/\allowbreak{}become, a supplied past participle. Translate the complete formula as given below. \\
tō … līehtinge & for the illumination of … ; dative līehtinge contains a dependent genitive day or night. \\
gesętte & he set/\allowbreak{}placed; supplied past form of the prefixed verb. Its subject continues God/\allowbreak{}hē, although no subject is repeated in this clause. Hīe is the accusative object “them”. \\
ðæt hīe scinen … & so that they might shine … . Here ðæt introduces purpose. Preserve the short i in past optative scinen, contrasting present scīnen earlier. \\
ofer eorðan & over the earth. The whole phrase is supplied. Eorðan can be accusative or dative singular in this spelling; no unique choice is required. \\
gīemden; tōdǣlden & might rule over; might separate. Supplied plural past optatives continuing the purpose clause. Gīemden governs the following genitives. \\
geseah ðā ðæt … & then saw that … ; supplied strong past verb and content clause. Here ðæt is a conjunction, not an article. \\
męrgen & source variant of the already learned morgen “morning”; retain its spelling when reading or quoting this passage. \\
\bottomrule
\end{tabular}
\par\endgroup

% mdtable:T28:end
\reserve{7\baselineskip}
% mdparagraph:P243:begin
\subsection*{Whole constructions}\BookHeading{nav-ch03-learner-0042}{subsubsection}{Whole constructions}{0}
% mdparagraph:P243:end

% teaching: NEW-U03-SUPPORT-U03-R-QUOTE
% mdparagraph:P244:begin
\textbf{Opening speech.} God cwæð ðā sōðlice = God then said; sōðlice is connective here (Latin autem).
% mdparagraph:P244:end

% teaching: NEW-U03-SUPPORT-U03-R-WISH
% mdparagraph:P245:begin
\textbf{The wish.} Bēon/\allowbreak{}bēon = let there be/\allowbreak{}let them be; tōdǣlen = let them separate; scīnen = let them shine; ālīehten = let them illuminate. All four verbs express the wish and take the lights as their plural subject.
% mdparagraph:P245:end

% teaching: NEW-U03-SUPPORT-U03-R-HEAVEN-GEN
% mdparagraph:P246:begin
\textbf{Nested heaven phrase.} on ðǣre heofonan fæstnesse = in the firmament of heaven. Heofonan is genitive singular and fæstnesse dative singular; ðǣre agrees with heofonan.
% mdparagraph:P246:end

% teaching: NEW-U03-SUPPORT-U03-R-TO-FUNCTION
% mdparagraph:P247:begin
\textbf{Assigned functions.} tō tācnum/\allowbreak{}tīdum/\allowbreak{}dagum/\allowbreak{}gēarum = for signs/\allowbreak{}seasons/\allowbreak{}days/\allowbreak{}years; tō ... līehtinge = for the illumination of ...; genitive day/\allowbreak{}night is nested inside the dative phrase.
% mdparagraph:P247:end

% teaching: NEW-U03-SUPPORT-U03-R-FULFILMENT
% mdparagraph:P248:begin
\textbf{Fulfilment.} hit wæs ðā swā geworden = it then came to pass thus /\allowbreak{} it was then so; geworden is supplied.
% mdparagraph:P248:end

% teaching: NEW-U03-SUPPORT-U03-R-SET-HEAVEN
% mdparagraph:P249:begin
\textbf{Setting the lights.} And gesętte hīe = and he set them. The subject continues God/\allowbreak{}hē. on ðǣre heofonan = in heaven, with source weak-feminine dative.
% mdparagraph:P249:end

% teaching: NEW-U03-SUPPORT-U03-R-PURPOSE
% mdparagraph:P250:begin
\textbf{Purpose clause.} ðæt hīe scinen ... and gīemden ... and tōdǣlden ... = so that they might shine ... and rule over ... and separate ... . Scinen has short i; the present scīnen occurs earlier in the speech.
% mdparagraph:P250:end

% teaching: NEW-U03-SUPPORT-U03-R-OVER
% mdparagraph:P251:begin
\textbf{Over the earth.} Ofer eorðan means “over the earth”. Either accusative or dative singular is possible here; the form eorðan does not decide between them.
% mdparagraph:P251:end

% teaching: NEW-U03-SUPPORT-U03-R-RULE-GEN
% mdparagraph:P252:begin
\textbf{Governed genitives.} gīemden ðæs dæges and ðǣre niht = might rule over the day and the night. Here the verb governs the two genitives; this is not an of-relationship between nouns.
% mdparagraph:P252:end

% teaching: NEW-U03-SUPPORT-U03-R-SEE-CONTENT
% mdparagraph:P253:begin
\textbf{Seeing that it was good.} God geseah ðā ðæt hit gōd wæs = God then saw that it was good. Ðæt is a conjunction, not the article. This clause pattern is supplied for the reading.
% mdparagraph:P253:end

% teaching: NEW-U03-SUPPORT-U03-R-CLOSE
% mdparagraph:P254:begin
\textbf{Closing day formula.} And wæs geworden ǣfen and męrgen se fēorða dæg = and evening and morning came to pass, the fourth day. This is a supplied formula, including its word order and singular wæs.
% mdparagraph:P254:end

% mdparagraph:P255:begin
The two occurrences of geworhte belong to the familiar indicative paradigm of gewyrcean. The paired greater and lesser lights use the new counted sense of lēoht. The text assigns their functions without naming the sun and moon. The final ordinal identifies the fourth day. The known wæs forms within fulfilment formulae do not make the larger constructions independent production material.
% mdparagraph:P255:end

\reserve{7\baselineskip}
% mdparagraph:P256:begin
\subsection*{Source reading — R1–R6}\BookHeading{nav-ch03-learner-0043}{subsubsection}{Source reading — R1–R6}{0}
% mdparagraph:P256:end

\par\begin{minipage}{\linewidth}\raggedright
% mdparagraph:P257:begin
\textbf{R1.}
% mdparagraph:P257:end

% item: NEW-U03-R01
% mdparagraph:P258:begin
Translate the complete fourth-day Creation account using the local support. Preserve the command, fulfilment, purposes and closing day formula.
% mdparagraph:P258:end

\end{minipage}\par\addvspace{4pt}

\par\begin{minipage}{\linewidth}\raggedright
% mdparagraph:P259:begin
\textbf{R2.}
% mdparagraph:P259:end

% item: NEW-U03-R02
% mdparagraph:P260:begin
In twā miclu lēoht, identify the gender, number and case of all three words. Give their headwords and explain the selected adjective stem and ending.
% mdparagraph:P260:end

\end{minipage}\par\addvspace{4pt}

\par\begin{minipage}{\linewidth}\raggedright
% mdparagraph:P261:begin
\textbf{R3.}
% mdparagraph:P261:end

% item: NEW-U03-R03
% mdparagraph:P262:begin
Compare ðæt māre lēoht and ðæt lǣsse lēoht. Give the positive adjectives and the analysis of each comparative. Using the local purpose support, state the function assigned to each light.
% mdparagraph:P262:end

\end{minipage}\par\addvspace{4pt}

\par\begin{minipage}{\linewidth}\raggedright
% mdparagraph:P263:begin
\textbf{R4.}
% mdparagraph:P263:end

% item: NEW-U03-R04
% mdparagraph:P264:begin
Find the two ðæt forms that introduce the light noun phrases. Contrast them with the ðæt before hīe scinen and before hit gōd wæs, using the supplied clause notes.
% mdparagraph:P264:end

\end{minipage}\par\addvspace{4pt}

\par\begin{minipage}{\linewidth}\raggedright
% mdparagraph:P265:begin
\textbf{R5.}
% mdparagraph:P265:end

% item: NEW-U03-R05
% mdparagraph:P266:begin
Find both occurrences of geworhte and give their person, number and tense. Identify their objects. Explain why the supplied Bēon and scīnen are not forms of the same indicative pattern.
% mdparagraph:P266:end

\end{minipage}\par\addvspace{4pt}

\par\begin{minipage}{\linewidth}\raggedright
% mdparagraph:P267:begin
\textbf{R6.}
% mdparagraph:P267:end

% item: NEW-U03-R06
% mdparagraph:P268:begin
Analyse gōd in ðæt hit gōd wæs and fēorða in se fēorða dæg. Explain their agreement and their different declensions. Use the supplied closing-formula note for the larger clause.
% mdparagraph:P268:end

\end{minipage}\par\addvspace{4pt}

\reserve{7\baselineskip}
% mdparagraph:P269:begin
\section*{15. Review}\BookHeading{nav-ch03-learner-0044}{subsection}{15. Review}{1}
% mdparagraph:P269:end

% mdparagraph:P270:begin
Use the taught forms and ordinary constructions independently. In X2, asterisks mark proposed errors to correct, not reconstructed historical forms. X3 permits more than one defensible composition; keep its stated meaning and referents clear.
% mdparagraph:P270:end

\reserve{7\baselineskip}
% mdparagraph:P271:begin
\subsection*{Review — X1–X3}\BookHeading{nav-ch03-learner-0045}{subsubsection}{Review — X1–X3}{0}
% mdparagraph:P271:end

\par\begin{minipage}{\linewidth}\raggedright
% mdparagraph:P272:begin
\textbf{X1.}
% mdparagraph:P272:end

% item: NEW-U03-X01
% mdparagraph:P273:begin
Translate “This good man gives my birds seed. My birds live in his small dwelling. The dwelling is greater than that tree.”
% mdparagraph:P273:end

\end{minipage}\par\addvspace{4pt}

\par\begin{minipage}{\linewidth}\raggedright
% mdparagraph:P274:begin
\textbf{X2.}
% mdparagraph:P274:end

% item: NEW-U03-X02
% mdparagraph:P275:begin
Correct the four marked adjective or numeral forms so that the sentences mean “The second day was good. My birds live in the great dwelling. The bird is smaller.”: Se *ōðra dæg wæs gōd. *Mīn fuglas libbað on ðǣm *micel hofe. Se fugol is *lǣss.
% mdparagraph:P275:end

\end{minipage}\par\addvspace{4pt}

\par\begin{minipage}{\linewidth}\raggedright
% mdparagraph:P276:begin
\textbf{X3.}
% mdparagraph:P276:end

% item: NEW-U03-X03
% mdparagraph:P277:begin
Write four connected Old English sentences: introduce two dwellings of unequal size, compare them, say that an old man has three birds, and give a direct quotation in which he says that he gives his birds food. Use an agreeing possessive in the quotation.
% mdparagraph:P277:end

\end{minipage}\par\addvspace{4pt}

\clearpage
% mdparagraph:P278:begin
\section*{16. Source and editorial notes}\BookHeading{nav-ch03-learner-0046}{subsection}{16. Source and editorial notes}{1}
% mdparagraph:P278:end

\begingroup\fontsize{10}{12.8}\selectfont
% mdparagraph:P279:begin
The grammar follows Cook, \textbf{A First Book in Old English}, second edition (1895), §§55–70, 77–80 and 82–85, with the ordinary adjective/\allowbreak{}adverb uses in §§179–182. Wright’s grammar corroborates stem variation, possessives and comparison. The bounded nominative comparison is supported by Fulk, §35. The tables distinguish selected teaching forms from supported alternatives. Forms expanded from a checked stem and paradigm need not be separately attested in a printed source.
% mdparagraph:P279:end

% mdparagraph:P280:begin
Cook’s Exercise I.4 is on printed p.2, one-based PDF p.10 of the selected exercise scan. Its original item wording and references are retained. The source reading is \textbf{First Book}, Reader I, printed p.125, editorially counted lines 8–15, and p.126, lines 1–4; one-based PDF pp.141–142. On p.125 the marginal line numbers 10 and 15 are printed; the other line numbers here are editorial counts. The new reading spans the whole fourth day.
% mdparagraph:P280:end

% mdparagraph:P281:begin
The reading removes apparatus cues, joins line and page breaks, and rejoins the line-divided words tācnum and līehtinge. It otherwise preserves the source’s letters, vowel marks and punctuation, including the colon after twā miclu lēoht. The following eighteen entries preserve Cook’s apparatus for this span. Questions in these historical notes are not additional assigned exercises; use the local support above for forms beyond this unit.
% mdparagraph:P281:end

% mdtable:T29:begin
\par\begingroup
\fontsize{10.5}{13.2}\selectfont
\setlength{\LTleft}{0pt}\setlength{\LTright}{\fill}
\begin{longtable}{@{}>{\raggedright\arraybackslash}p{31.07mm}>{\raggedright\arraybackslash}p{90.12mm}@{}}
\toprule
\textbf{Printed page/\allowbreak{}note} & \textbf{Cook’s note} \\
\midrule
\endfirsthead
\toprule
\textbf{Printed page/\allowbreak{}note} & \textbf{Cook’s note} \\
\midrule
\endhead
\midrule
\multicolumn{2}{r}{\footnotesize Continued on the following page.}\\
\endfoot
\bottomrule
\endlastfoot
125/\allowbreak{}15 & Lat. autem. \\
125/\allowbreak{}16 & See 166. \\
125/\allowbreak{}17 & Gen. sing. \\
125/\allowbreak{}18 & See tōdǣlan. \\
125/\allowbreak{}19 & See tācen, and 24. \\
125/\allowbreak{}20 & See tīd, and 24. \\
125/\allowbreak{}21 & See dæg, and 24. \\
125/\allowbreak{}22 & See gēar. \\
125/\allowbreak{}23 & See 193. a. Write the opt. pret. plur. of this verb. \\
125/\allowbreak{}24 & See twēgen. \\
125/\allowbreak{}25 & See micel. \\
125/\allowbreak{}26 & See 66. \\
125/\allowbreak{}27 & What is the relation of the stem-vowel of this word to that of lēoht? \\
125/\allowbreak{}28 & For niht, instead of nieht, see 19. See 153. d. \\
125/\allowbreak{}29 & See gesęttan, and 189, note. \\
126/\allowbreak{}1 & Opt. pret. = Lat. lucerent. What would be the opt. pres.? \\
126/\allowbreak{}2 & See 156. f. \\
126/\allowbreak{}3 & See 78. \\
\end{longtable}
\par\endgroup

% mdtable:T29:end
% mdparagraph:P282:begin
In these notes “opt.” means optative, Cook’s term for the forms commonly called subjunctive. The supplied present/\allowbreak{}past contrast is sufficient for this reading. His historical vowel question does not require an additional sound-change exercise here.
% mdparagraph:P282:end

\endgroup

% book-source-body:end
\par\label{end-ch03-learner}
\clearpage\endgroup

\BookRecto
\begingroup
\renewcommand{\runningtitle}{Unit 4. Strong verbs I–III}
\hyphenpenalty=10000
\exhyphenpenalty=10000
\emergencystretch=3em
\newcommand{\itempara}[2]{\par\addvspace{3pt}\noindent\hangindent=1.8em\hangafter=1\makebox[1.8em][l]{#1.}#2\par}
\newcommand{\reserve}[1]{\par\penalty-100\begingroup\dimen0=\pagegoal\advance\dimen0 by-\pagetotal\ifdim\dimen0<#1\newpage\fi\endgroup}
% book-source-body:begin

\raggedright
\booktitle{\BookComponent{ch04-learner}{Unit 4. Strong verbs I–III}UNIT 4}{Strong verbs I–III}

% mdparagraph:P001:begin
Strong verbs form their preterite chiefly by changing the stem vowel. This unit gives the complete selected present and preterite indicative of Classes I–III, with the stem changes needed to use them. The noun, pronoun, adjective and weak-verb systems of Units 1–3 remain in use.
% mdparagraph:P001:end

% mdparagraph:P002:begin
In the tables, \textbf{1}, \textbf{2} and \textbf{3} denote first, second and third person; \textbf{sg.} and \textbf{pl.} denote singular and plural. A row labelled \textbf{1–3 pl.} gives the same form for all three plural persons. Dual pronoun subjects also take these plural forms. Use the selected forms when requested; the documented alternatives remain valid Old English. Macrons and the distinction between e and ę are retained.
% mdparagraph:P002:end

\reserve{7\baselineskip}
% mdparagraph:P003:begin
\section*{1. Principal parts and finite stems}\BookHeading{nav-ch04-learner-0001}{subsection}{1. Principal parts and finite stems}{1}\BookLabel{bk-ch04-learner-28ffaf280840}
% mdparagraph:P003:end

% teaching: NEW-U04-GRAM-STRONG-WEAK-CONTRAST
% teaching: NEW-U04-GRAM-PRINCIPAL-PARTS-REFERENCE
% teaching: NEW-U04-GRAM-INDICATIVE-SCOPE
% mdparagraph:P004:begin
The weak verbs already learned add a dental element in the preterite. A strong verb instead has a series of stems: the vowel of its present, singular preterite and plural preterite may differ. The \textbf{four principal parts} identify this series. They are conventionally listed as infinitive, singular preterite, plural preterite and past participle.
% mdparagraph:P004:end

% table: U04-PRINCIPAL-PART-POSITIONS
% mdtable:T01:begin
\par\begingroup
\fontsize{10.5}{13.2}\selectfont
\setlength{\LTleft}{0pt}\setlength{\LTright}{\fill}
\noindent
\begin{tabular}{@{}>{\raggedright\arraybackslash}p{33.37mm}>{\raggedright\arraybackslash}p{87.82mm}@{}}
\toprule
\textbf{Position} & \textbf{Function} \\
\midrule
1. Infinitive & Names the verb and identifies the present stem. \\
2. Singular preterite & Identifies the ordinary first/\allowbreak{}third singular preterite stem. \\
3. Plural preterite & Identifies the stem used in all plural persons and in the second singular. \\
4. Past participle & Completes the cited stem series; its syntax is not yet required. \\
\bottomrule
\end{tabular}
\par\endgroup

% mdtable:T01:end
% mdparagraph:P005:begin
Thus drīfan, drāf, drifon, drifen display the Class I sequence ī : ā : i : i. Drāf and drifon are finite indicative forms. The infinitive and participle columns are \textbf{supplied citations}: they may be consulted, and their independent recall or use in infinitive, participial, perfect or passive constructions is not required here. A principal-part entry identifies a stem; the complete tables below identify the person and number.
% mdparagraph:P005:end

% mdparagraph:P006:begin
Present indicative forms state present actions or conditions; preterite indicative forms state past actions or conditions. Every finite table in this unit covers those two tenses and all three persons in singular and plural. The two rare reference verbs in §5 are for recognition only.
% mdparagraph:P006:end

\reserve{7\baselineskip}
% mdparagraph:P007:begin
\section*{2. Endings and stem changes}\BookHeading{nav-ch04-learner-0002}{subsection}{2. Endings and stem changes}{1}\BookLabel{bk-ch04-learner-3d8a96a1ca23}
% mdparagraph:P007:end

% teaching: NEW-U04-GRAM-PRESENT-ENDINGS
% teaching: NEW-U04-GRAM-PRETERITE-ENDINGS
% table: U04-ENDINGS
% mdtable:T02:begin
% course-table-style: 2026-10
% course-panel:T02-P01:begin
\par\begingroup
\fontsize{10.5}{13.2}\selectfont
\setlength{\LTleft}{0pt}\setlength{\LTright}{\fill}
\noindent
\begin{tabular}{@{}>{\raggedright\arraybackslash}p{39.459mm}>{\raggedright\arraybackslash}p{39.459mm}>{\raggedright\arraybackslash}p{39.459mm}@{}}
\toprule
\textbf{Person and number} & \textbf{Present} & \textbf{Preterite} \\
\midrule
1 sg. & -e & singular preterite stem; no added ending \\
2 sg. & -st, with the listed stem changes & plural preterite stem + -e \\
3 sg. & -ð, with the listed stem changes & singular preterite stem; no added ending \\
1–3 pl. & -að & plural preterite stem + -on \\
\bottomrule
\end{tabular}
\par\endgroup
% course-panel:T02-P01:end
% mdtable:T02:end
% mdparagraph:P008:begin
Contracted verbs have the special present forms printed in their tables. The strong \textbf{second singular preterite} uses the plural stem, not the ordinary first/\allowbreak{}third singular stem: drife belongs with drifon, while drāf is the first/\allowbreak{}third singular. Findan has an additional third-singular exception, given in §4. The indicative tables keep every person explicit.
% mdparagraph:P008:end

% teaching: NEW-U04-GRAM-PRESENT-CONSONANTS
% teaching: NEW-U04-GRAM-GRAMMATICAL-CHANGE
% mdparagraph:P009:begin
The second and third singular present often lose the vowel of older -est and -eð endings. The resulting consonant groups change. Dental d or t before ð gives tt, sometimes written t; d before st gives tst. Stem-final s combines with either ending as st. A stem-final ð may disappear before st, and doubled ð may be reduced. Doubled consonants may also simplify before a consonantal ending. The selected forms below show the result for each verb.
% mdparagraph:P009:end

% mdparagraph:P010:begin
Some verbs change a consonant between their principal-part stems: snīðan and sēoðan have d in their plural-based preterites, while cēosan has r. Such alternation belongs to the particular verb and stem position. Likewise, g/\allowbreak{}h alternatives occur in the stated environments after a long vowel or a liquid; they are not interchangeable throughout a paradigm.
% mdparagraph:P010:end

% mdparagraph:P011:begin
Full -est/\allowbreak{}-eð forms also occur. Their vowels may retain an older change or follow another part of the paradigm by analogy. The notes list documented alternatives; the selected tables do not license automatic manufacture of every possible full-ending spelling.
% mdparagraph:P011:end

\reserve{7\baselineskip}
% mdparagraph:P012:begin
\section*{3. Classes I and II}\BookHeading{nav-ch04-learner-0003}{subsection}{3. Classes I and II}{1}
% mdparagraph:P012:end

\reserve{7\baselineskip}
% mdparagraph:P013:begin
\subsection*{Class I}\BookHeading{nav-ch04-learner-0004}{subsubsection}{Class I}{0}\BookLabel{bk-ch04-learner-70278eafbcc8}
% mdparagraph:P013:end

% teaching: NEW-U04-GRAM-CLASS-I
% mdparagraph:P014:begin
The ordinary Class I series is \textbf{ī : ā : i : i}. The long ī remains in the selected present stem; the preterite singular has ā, and the plural-based preterite has short i. The fourth vowel belongs to the supplied participle citation.
% mdparagraph:P014:end

% teaching: NEW-U04-VOC-drifan
% teaching: NEW-U04-CIT-drifan
% teaching: NEW-U04-VOC-ridan
% teaching: NEW-U04-CIT-ridan
% teaching: NEW-U04-VOC-writan
% teaching: NEW-U04-CIT-writan
% teaching: NEW-U04-VOC-risan
% teaching: NEW-U04-CIT-risan
% teaching: NEW-U04-VOC-snithan
% teaching: NEW-U04-CIT-snithan
% teaching: NEW-U04-VOC-wreon
% teaching: NEW-U04-CIT-wreon
% table: U04-PARTS-I
% mdtable:T03:begin
\par\begingroup
\fontsize{10.5}{13.2}\selectfont
\setlength{\LTleft}{0pt}\setlength{\LTright}{\fill}
\noindent
\begin{tabular}{@{}>{\raggedright\arraybackslash}p{23.14mm}>{\raggedright\arraybackslash}p{21.04mm}>{\raggedright\arraybackslash}p{22.72mm}>{\raggedright\arraybackslash}p{22.72mm}>{\raggedright\arraybackslash}p{23.14mm}@{}}
\toprule
\textbf{Infinitive¹} & \textbf{Meaning} & \textbf{Preterite sg.} & \textbf{Preterite pl.} & \textbf{Participle¹} \\
\midrule
drīfan & drive & drāf & drifon & drifen \\
rīdan & ride & rād & ridon & riden \\
wrītan & write & wrāt & writon & writen \\
rīsan & rise & rās & rison & risen \\
snīðan & cut & snāð & snidon & sniden \\
wrēon & cover; clothe & wrāh & wrigon & wrigen \\
\bottomrule
\end{tabular}
\par\endgroup

% mdtable:T03:end
% mdparagraph:P015:begin
¹ The infinitive and participle are supplied citations in every principal-part table. The two middle columns identify finite preterite stems; their person assignments are given below.
% mdparagraph:P015:end

% teaching: NEW-U04-PAR-drifan
% teaching: NEW-U04-PAR-ridan
% teaching: NEW-U04-PAR-writan
% teaching: NEW-U04-PAR-risan
% teaching: NEW-U04-PAR-snithan
% teaching: NEW-U04-PAR-wreon
\reserve{11\baselineskip}
% mdparagraph:P016:begin
\textbf{Present indicative.}
% mdparagraph:P016:end

% table: U04-PRESENT-I
% mdtable:T04:begin
% course-table-style: 2026-10
% course-panel:T04-P01:begin
\par\begingroup
\fontsize{10.5}{13.2}\selectfont
\setlength{\LTleft}{0pt}\setlength{\LTright}{\fill}
\noindent
\begin{tabular}{@{}>{\raggedright\arraybackslash}p{24.232mm}>{\raggedright\arraybackslash}p{22.13mm}>{\raggedright\arraybackslash}p{22.13mm}>{\raggedright\arraybackslash}p{22.13mm}>{\raggedright\arraybackslash}p{22.13mm}@{}}
\toprule
\textbf{Person and number} & \textbf{drīfan} & \textbf{rīdan} & \textbf{wrītan} & \textbf{rīsan} \\
\midrule
1 sg. & drīfe & rīde & wrīte & rīse \\
2 sg. & drīfst & rītst & wrītst & rīst \\
3 sg. & drīfð & rītt & wrītt & rīst \\
1–3 pl. & drīfað & rīdað & wrītað & rīsað \\
\bottomrule
\end{tabular}
\par\endgroup
% course-panel:T04-P01:end
\par\addvspace{5pt}
% course-panel:T04-P02:begin
\par\begingroup
\fontsize{10.5}{13.2}\selectfont
\setlength{\LTleft}{0pt}\setlength{\LTright}{\fill}
\noindent
\begin{tabular}{@{}>{\raggedright\arraybackslash}p{39.459mm}>{\raggedright\arraybackslash}p{39.459mm}>{\raggedright\arraybackslash}p{39.459mm}@{}}
\toprule
\textbf{Person and number} & \textbf{snīðan} & \textbf{wrēon} \\
\midrule
1 sg. & snīðe & wrēo \\
2 sg. & snīst & wrīhst \\
3 sg. & snīð & wrīhð \\
1–3 pl. & snīðað & wrēoð \\
\bottomrule
\end{tabular}
\par\endgroup
% course-panel:T04-P02:end
% mdtable:T04:end
\reserve{11\baselineskip}
% mdparagraph:P017:begin
\textbf{Preterite indicative.}
% mdparagraph:P017:end

% table: U04-PRETERITE-I
% mdtable:T05:begin
% course-table-style: 2026-10
% course-panel:T05-P01:begin
\par\begingroup
\fontsize{10.5}{13.2}\selectfont
\setlength{\LTleft}{0pt}\setlength{\LTright}{\fill}
\noindent
\begin{tabular}{@{}>{\raggedright\arraybackslash}p{24.232mm}>{\raggedright\arraybackslash}p{22.13mm}>{\raggedright\arraybackslash}p{22.13mm}>{\raggedright\arraybackslash}p{22.13mm}>{\raggedright\arraybackslash}p{22.13mm}@{}}
\toprule
\textbf{Person and number} & \textbf{drīfan} & \textbf{rīdan} & \textbf{wrītan} & \textbf{rīsan} \\
\midrule
1 sg. & drāf & rād & wrāt & rās \\
2 sg. & drife & ride & write & rise \\
3 sg. & drāf & rād & wrāt & rās \\
1–3 pl. & drifon & ridon & writon & rison \\
\bottomrule
\end{tabular}
\par\endgroup
% course-panel:T05-P01:end
\par\addvspace{5pt}
% course-panel:T05-P02:begin
\par\begingroup
\fontsize{10.5}{13.2}\selectfont
\setlength{\LTleft}{0pt}\setlength{\LTright}{\fill}
\noindent
\begin{tabular}{@{}>{\raggedright\arraybackslash}p{39.459mm}>{\raggedright\arraybackslash}p{39.459mm}>{\raggedright\arraybackslash}p{39.459mm}@{}}
\toprule
\textbf{Person and number} & \textbf{snīðan} & \textbf{wrēon} \\
\midrule
1 sg. & snāð & wrāh \\
2 sg. & snide & wrige \\
3 sg. & snāð & wrāh \\
1–3 pl. & snidon & wrigon \\
\bottomrule
\end{tabular}
\par\endgroup
% course-panel:T05-P02:end
% mdtable:T05:end
% mdparagraph:P018:begin
In rīdan and wrītan, the selected present second singular ends in -tst and the third in -tt. Third-singular rīt and wrīt are documented alternatives. Rīsan has rīst in both second and third singular; rīsð is another third-singular spelling. Snīðan has selected snīst and snīð, with snīðð also valid in the third singular. Its preterite switches to d in snide and snidon, while the first/\allowbreak{}third singular snāð retains ð.
% mdparagraph:P018:end

% teaching: NEW-U04-GRAM-CONTRACT-I
% mdparagraph:P019:begin
Wrēon\BookPoint{bk-ch04-learner-04534b37414d} has a contracted present: wrēo, wrīhst, wrīhð, wrēoð. Its preterite uses wrāh beside plural-based wrige and wrigon. The h/\allowbreak{}g contrast is part of this pattern. The table selects its Class I series; documented analogical Class II forms are not mixed into that selection. See source note 2.
% mdparagraph:P019:end

\reserve{7\baselineskip}
% mdparagraph:P020:begin
\subsection*{Class II}\BookHeading{nav-ch04-learner-0005}{subsubsection}{Class II}{0}\BookLabel{bk-ch04-learner-d95b114af330}
% mdparagraph:P020:end

% teaching: NEW-U04-GRAM-CLASS-II
% teaching: NEW-U04-GRAM-PRESENT-MUTATION
% mdparagraph:P021:begin
The ordinary Class II series is \textbf{ēo or ū : ēa : u : o}. In the selected second and third singular present, ēo changes to īe and ū to ȳ. The first singular and plural retain the present vowel shown in the infinitive citation. Cook also records ī beside īe; the tables consistently select īe.
% mdparagraph:P021:end

% teaching: NEW-U04-VOC-beodan
% teaching: NEW-U04-CIT-beodan
% teaching: NEW-U04-VOC-bugan
% teaching: NEW-U04-CIT-bugan
% teaching: NEW-U04-VOC-ceosan
% teaching: NEW-U04-CIT-ceosan
% teaching: NEW-U04-VOC-seothan
% teaching: NEW-U04-CIT-seothan
% teaching: NEW-U04-VOC-dreogan
% teaching: NEW-U04-CIT-dreogan
% teaching: NEW-U04-VOC-teon
% teaching: NEW-U04-CIT-teon
% table: U04-PARTS-II
% mdtable:T06:begin
\par\begingroup
\fontsize{10.5}{13.2}\selectfont
\setlength{\LTleft}{0pt}\setlength{\LTright}{\fill}
\noindent
\begin{tabular}{@{}>{\raggedright\arraybackslash}p{23.14mm}>{\raggedright\arraybackslash}p{21.04mm}>{\raggedright\arraybackslash}p{22.72mm}>{\raggedright\arraybackslash}p{22.72mm}>{\raggedright\arraybackslash}p{23.14mm}@{}}
\toprule
\textbf{Infinitive¹} & \textbf{Meaning} & \textbf{Preterite sg.} & \textbf{Preterite pl.} & \textbf{Participle¹} \\
\midrule
bēodan & offer & bēad & budon & boden \\
būgan & bend & bēag & bugon & bogen \\
cēosan & choose & cēas & curon & coren \\
sēoðan & boil & sēað & sudon & soden \\
drēogan & endure & drēag & drugon & drogen \\
tēon & draw; pull & tēah & tugon & togen \\
\bottomrule
\end{tabular}
\par\endgroup

% mdtable:T06:end
% teaching: NEW-U04-PAR-beodan
% teaching: NEW-U04-PAR-bugan
% teaching: NEW-U04-PAR-ceosan
% teaching: NEW-U04-PAR-seothan
% teaching: NEW-U04-PAR-dreogan
% teaching: NEW-U04-PAR-teon
\reserve{11\baselineskip}
% mdparagraph:P022:begin
\textbf{Present indicative.}
% mdparagraph:P022:end

% table: U04-PRESENT-II
% mdtable:T07:begin
% course-table-style: 2026-10
% course-panel:T07-P01:begin
\par\begingroup
\fontsize{10.5}{13.2}\selectfont
\setlength{\LTleft}{0pt}\setlength{\LTright}{\fill}
\noindent
\begin{tabular}{@{}>{\raggedright\arraybackslash}p{24.232mm}>{\raggedright\arraybackslash}p{22.13mm}>{\raggedright\arraybackslash}p{22.13mm}>{\raggedright\arraybackslash}p{22.13mm}>{\raggedright\arraybackslash}p{22.13mm}@{}}
\toprule
\textbf{Person and number} & \textbf{bēodan} & \textbf{būgan} & \textbf{cēosan} & \textbf{sēoðan} \\
\midrule
1 sg. & bēode & būge & cēose & sēoðe \\
2 sg. & bīetst & bȳgst & cīest & sīest \\
3 sg. & bīett & bȳgð & cīest & sīeð \\
1–3 pl. & bēodað & būgað & cēosað & sēoðað \\
\bottomrule
\end{tabular}
\par\endgroup
% course-panel:T07-P01:end
\par\addvspace{5pt}
% course-panel:T07-P02:begin
\par\begingroup
\fontsize{10.5}{13.2}\selectfont
\setlength{\LTleft}{0pt}\setlength{\LTright}{\fill}
\noindent
\begin{tabular}{@{}>{\raggedright\arraybackslash}p{39.459mm}>{\raggedright\arraybackslash}p{39.459mm}>{\raggedright\arraybackslash}p{39.459mm}@{}}
\toprule
\textbf{Person and number} & \textbf{drēogan} & \textbf{tēon} \\
\midrule
1 sg. & drēoge & tēo \\
2 sg. & drīegst & tīehst \\
3 sg. & drīegð & tīehð \\
1–3 pl. & drēogað & tēoð \\
\bottomrule
\end{tabular}
\par\endgroup
% course-panel:T07-P02:end
% mdtable:T07:end
\reserve{11\baselineskip}
% mdparagraph:P023:begin
\textbf{Preterite indicative.}
% mdparagraph:P023:end

% table: U04-PRETERITE-II
% mdtable:T08:begin
% course-table-style: 2026-10
% course-panel:T08-P01:begin
\par\begingroup
\fontsize{10.5}{13.2}\selectfont
\setlength{\LTleft}{0pt}\setlength{\LTright}{\fill}
\noindent
\begin{tabular}{@{}>{\raggedright\arraybackslash}p{24.232mm}>{\raggedright\arraybackslash}p{22.13mm}>{\raggedright\arraybackslash}p{22.13mm}>{\raggedright\arraybackslash}p{22.13mm}>{\raggedright\arraybackslash}p{22.13mm}@{}}
\toprule
\textbf{Person and number} & \textbf{bēodan} & \textbf{būgan} & \textbf{cēosan} & \textbf{sēoðan} \\
\midrule
1 sg. & bēad & bēag & cēas & sēað \\
2 sg. & bude & buge & cure & sude \\
3 sg. & bēad & bēag & cēas & sēað \\
1–3 pl. & budon & bugon & curon & sudon \\
\bottomrule
\end{tabular}
\par\endgroup
% course-panel:T08-P01:end
\par\addvspace{5pt}
% course-panel:T08-P02:begin
\par\begingroup
\fontsize{10.5}{13.2}\selectfont
\setlength{\LTleft}{0pt}\setlength{\LTright}{\fill}
\noindent
\begin{tabular}{@{}>{\raggedright\arraybackslash}p{39.459mm}>{\raggedright\arraybackslash}p{39.459mm}>{\raggedright\arraybackslash}p{39.459mm}@{}}
\toprule
\textbf{Person and number} & \textbf{drēogan} & \textbf{tēon} \\
\midrule
1 sg. & drēag & tēah \\
2 sg. & druge & tuge \\
3 sg. & drēag & tēah \\
1–3 pl. & drugon & tugon \\
\bottomrule
\end{tabular}
\par\endgroup
% course-panel:T08-P02:end
% mdtable:T08:end
% mdparagraph:P024:begin
Bēodan combines the vowel change with dental endings: bīetst and bīett; third-singular bīet is also valid. Būgan has selected bȳgst/\allowbreak{}bȳgð and the documented alternatives bȳhst/\allowbreak{}bȳhð. Its singular preterite bēag may also be written bēah.
% mdparagraph:P024:end

% mdparagraph:P025:begin
Cēosan has cīest in both second and third singular present, cēas in the first/\allowbreak{}third singular preterite, and r in cure and curon. Sēoðan has sīest and sīeð in the present, sēað in the first/\allowbreak{}third singular preterite, and d in sude and sudon. The fourth principal parts retain the corresponding coren and soden stems.
% mdparagraph:P025:end

% mdparagraph:P026:begin
Drēogan selects drīegst/\allowbreak{}drīegð in the present; drīehst/\allowbreak{}drīehð are documented alternatives. Its first/\allowbreak{}third singular preterite may be drēah beside selected drēag. Keep those g/\allowbreak{}h choices distinct from the s/\allowbreak{}r and ð/\allowbreak{}d alternations of the preceding verbs.
% mdparagraph:P026:end

% teaching: NEW-U04-GRAM-CONTRACT-II
% mdparagraph:P027:begin
Tēon\BookPoint{bk-ch04-learner-cde3298d0700} “draw, pull” is the contracted Class II verb here. Its present is tēo, tīehst, tīehð, tēoð; the preterite has tēah beside tuge and tugon. Its similar infinitive ending does not give it the Class I preterite of wrēon. The meaning identifies the assigned verb; no other homographic verb is included.
% mdparagraph:P027:end

\reserve{7\baselineskip}
% mdparagraph:P028:begin
\section*{4. Class III}\BookHeading{nav-ch04-learner-0006}{subsection}{4. Class III}{1}\BookLabel{bk-ch04-learner-d75e819526a2}
% mdparagraph:P028:end

% teaching: NEW-U04-GRAM-CLASS-III-NASAL
% teaching: NEW-U04-GRAM-CLASS-III-EL-IEL
% teaching: NEW-U04-GRAM-CLASS-III-EOR
% teaching: NEW-U04-GRAM-CLASS-III-METATHESIS
% mdparagraph:P029:begin
Class III contains several related stem patterns. Its vowels in the following series are \textbf{short}. The present second/\allowbreak{}third singular may differ from the infinitive vowel, and the consonant changes described in §2 still apply.
% mdparagraph:P029:end

% table: U04-III-SERIES
% mdtable:T09:begin
\par\begingroup
\fontsize{10.5}{13.2}\selectfont
\setlength{\LTleft}{0pt}\setlength{\LTright}{\fill}
\noindent
\begin{tabular}{@{}>{\raggedright\arraybackslash}p{30.49mm}>{\raggedright\arraybackslash}p{34.43mm}>{\raggedright\arraybackslash}p{53.45mm}@{}}
\toprule
\textbf{Family} & \textbf{Four-part vowel series} & \textbf{Assigned verbs} \\
\midrule
Nasal + consonant & i : a : u : u & bindan, winnan, findan, drincan, singan \\
el & e : ea : u : o & helpan \\
iel & ie : ea : u : o & gieldan \\
eor & eo : ea : u : o & weorpan, weorðan, forweorðan \\
Metathesis & e : æ : u : o & berstan, tōberstan \\
\bottomrule
\end{tabular}
\par\endgroup

% mdtable:T09:end
% mdparagraph:P030:begin
Metathesis is a change in the order of sounds. In the berstan family, r and the neighbouring vowel have changed order. Its selected series is berstan, bærst, burston, borsten; it is not obtained by substituting the weorpan vowel in every cell.
% mdparagraph:P030:end

% teaching: NEW-U04-VOC-bindan
% teaching: NEW-U04-CIT-bindan
% teaching: NEW-U04-VOC-helpan
% teaching: NEW-U04-CIT-helpan
% teaching: NEW-U04-VOC-gieldan
% teaching: NEW-U04-CIT-gieldan
% teaching: NEW-U04-VOC-weorpan
% teaching: NEW-U04-CIT-weorpan
% teaching: NEW-U04-VOC-berstan
% teaching: NEW-U04-CIT-berstan
% teaching: NEW-U04-VOC-weorthan
% teaching: NEW-U04-CIT-weorthan
% teaching: NEW-U04-VOC-winnan
% teaching: NEW-U04-CIT-winnan
% teaching: NEW-U04-VOC-findan
% teaching: NEW-U04-CIT-findan
% teaching: NEW-U04-VOC-drincan
% teaching: NEW-U04-CIT-drincan
% teaching: NEW-U04-VOC-singan
% teaching: NEW-U04-CIT-singan
% teaching: NEW-U04-VOC-toberstan
% teaching: NEW-U04-CIT-toberstan
% teaching: NEW-U04-VOC-forweorthan
% teaching: NEW-U04-CIT-forweorthan
% table: U04-PARTS-III
% mdtable:T10:begin
\par\begingroup
\fontsize{10.5}{13.2}\selectfont
\setlength{\LTleft}{0pt}\setlength{\LTright}{\fill}
\noindent
\begin{tabular}{@{}>{\raggedright\arraybackslash}p{21.22mm}>{\raggedright\arraybackslash}p{28.65mm}>{\raggedright\arraybackslash}p{20.83mm}>{\raggedright\arraybackslash}p{20.83mm}>{\raggedright\arraybackslash}p{21.22mm}@{}}
\toprule
\textbf{Infinitive¹} & \textbf{Meaning} & \textbf{Preterite sg.} & \textbf{Preterite pl.} & \textbf{Participle¹} \\
\midrule
bindan & bind; tie & band & bundon & bunden \\
helpan & help & healp & hulpon & holpen \\
gieldan & yield & geald & guldon & golden \\
weorpan & throw & wearp & wurpon & worpen \\
berstan & burst & bærst & burston & borsten \\
weorðan & become & wearð & wurdon & worden \\
winnan & struggle; toil & wann & wunnon & wunnen \\
findan & find & fand & fundon & funden \\
drincan & drink & dranc & druncon & druncen \\
singan & sing & sang & sungon & sungen \\
tōberstan & break up; go to pieces & tōbærst & tōburston & tōborsten \\
forweorðan & perish & forwearð & forwurdon & forworden \\
\bottomrule
\end{tabular}
\par\endgroup

% mdtable:T10:end
% teaching: NEW-U04-PAR-bindan
% teaching: NEW-U04-PAR-helpan
% teaching: NEW-U04-PAR-gieldan
% teaching: NEW-U04-PAR-weorpan
% teaching: NEW-U04-PAR-berstan
% teaching: NEW-U04-PAR-weorthan
% teaching: NEW-U04-PAR-winnan
% teaching: NEW-U04-PAR-findan
% teaching: NEW-U04-PAR-drincan
% teaching: NEW-U04-PAR-singan
% teaching: NEW-U04-PAR-toberstan
% teaching: NEW-U04-PAR-forweorthan
\reserve{17\baselineskip}
% mdparagraph:P031:begin
\textbf{Present indicative.}
% mdparagraph:P031:end

% table: U04-PRESENT-III
% mdtable:T11:begin
% course-table-style: 2026-10
% course-panel:T11-P01:begin
\par\begingroup
\fontsize{10.5}{13.2}\selectfont
\setlength{\LTleft}{0pt}\setlength{\LTright}{\fill}
\noindent
\begin{tabular}{@{}>{\raggedright\arraybackslash}p{24.232mm}>{\raggedright\arraybackslash}p{22.13mm}>{\raggedright\arraybackslash}p{22.13mm}>{\raggedright\arraybackslash}p{22.13mm}>{\raggedright\arraybackslash}p{22.13mm}@{}}
\toprule
\textbf{Person and number} & \textbf{bindan} & \textbf{helpan} & \textbf{gieldan} & \textbf{weorpan} \\
\midrule
1 sg. & binde & helpe & gielde & weorpe \\
2 sg. & bintst & hilpst & gieltst & wierpst \\
3 sg. & bint & hilpð & gielt & wierpð \\
1–3 pl. & bindað & helpað & gieldað & weorpað \\
\bottomrule
\end{tabular}
\par\endgroup
% course-panel:T11-P01:end
\par\addvspace{5pt}
% course-panel:T11-P02:begin
\par\begingroup
\fontsize{10.5}{13.2}\selectfont
\setlength{\LTleft}{0pt}\setlength{\LTright}{\fill}
\noindent
\begin{tabular}{@{}>{\raggedright\arraybackslash}p{24.232mm}>{\raggedright\arraybackslash}p{22.13mm}>{\raggedright\arraybackslash}p{22.13mm}>{\raggedright\arraybackslash}p{22.13mm}>{\raggedright\arraybackslash}p{22.13mm}@{}}
\toprule
\textbf{Person and number} & \textbf{berstan} & \textbf{weorðan} & \textbf{winnan} & \textbf{findan} \\
\midrule
1 sg. & berste & weorðe & winne & finde \\
2 sg. & bierst & wierst & winst & fintst \\
3 sg. & bierst & wierð & winð & fint \\
1–3 pl. & berstað & weorðað & winnað & findað \\
\bottomrule
\end{tabular}
\par\endgroup
% course-panel:T11-P02:end
\par\addvspace{5pt}
% course-panel:T11-P03:begin
\par\begingroup
\fontsize{10.5}{13.2}\selectfont
\setlength{\LTleft}{0pt}\setlength{\LTright}{\fill}
\noindent
\begin{tabular}{@{}>{\raggedright\arraybackslash}p{24.232mm}>{\raggedright\arraybackslash}p{22.13mm}>{\raggedright\arraybackslash}p{22.13mm}>{\raggedright\arraybackslash}p{22.13mm}>{\raggedright\arraybackslash}p{22.13mm}@{}}
\toprule
\textbf{Person and number} & \textbf{drincan} & \textbf{singan} & \textbf{tōberstan} & \textbf{forweorðan} \\
\midrule
1 sg. & drince & singe & tōberste & forweorðe \\
2 sg. & drincst & singst & tōbierst & forwierst \\
3 sg. & drincð & singð & tōbierst & forwierð \\
1–3 pl. & drincað & singað & tōberstað & forweorðað \\
\bottomrule
\end{tabular}
\par\endgroup
% course-panel:T11-P03:end
% mdtable:T11:end
\reserve{17\baselineskip}
% mdparagraph:P032:begin
\textbf{Preterite indicative.}
% mdparagraph:P032:end

% table: U04-PRETERITE-III
% mdtable:T12:begin
% course-table-style: 2026-10
% course-panel:T12-P01:begin
\par\begingroup
\fontsize{10.5}{13.2}\selectfont
\setlength{\LTleft}{0pt}\setlength{\LTright}{\fill}
\noindent
\begin{tabular}{@{}>{\raggedright\arraybackslash}p{24.232mm}>{\raggedright\arraybackslash}p{22.13mm}>{\raggedright\arraybackslash}p{22.13mm}>{\raggedright\arraybackslash}p{22.13mm}>{\raggedright\arraybackslash}p{22.13mm}@{}}
\toprule
\textbf{Person and number} & \textbf{bindan} & \textbf{helpan} & \textbf{gieldan} & \textbf{weorpan} \\
\midrule
1 sg. & band & healp & geald & wearp \\
2 sg. & bunde & hulpe & gulde & wurpe \\
3 sg. & band & healp & geald & wearp \\
1–3 pl. & bundon & hulpon & guldon & wurpon \\
\bottomrule
\end{tabular}
\par\endgroup
% course-panel:T12-P01:end
\par\addvspace{5pt}
% course-panel:T12-P02:begin
\par\begingroup
\fontsize{10.5}{13.2}\selectfont
\setlength{\LTleft}{0pt}\setlength{\LTright}{\fill}
\noindent
\begin{tabular}{@{}>{\raggedright\arraybackslash}p{24.232mm}>{\raggedright\arraybackslash}p{22.13mm}>{\raggedright\arraybackslash}p{22.13mm}>{\raggedright\arraybackslash}p{22.13mm}>{\raggedright\arraybackslash}p{22.13mm}@{}}
\toprule
\textbf{Person and number} & \textbf{berstan} & \textbf{weorðan} & \textbf{winnan} & \textbf{findan} \\
\midrule
1 sg. & bærst & wearð & wann & fand \\
2 sg. & burste & wurde & wunne & funde \\
3 sg. & bærst & wearð & wann & funde \\
1–3 pl. & burston & wurdon & wunnon & fundon \\
\bottomrule
\end{tabular}
\par\endgroup
% course-panel:T12-P02:end
\par\addvspace{5pt}
% course-panel:T12-P03:begin
\par\begingroup
\fontsize{10.5}{13.2}\selectfont
\setlength{\LTleft}{0pt}\setlength{\LTright}{\fill}
\noindent
\begin{tabular}{@{}>{\raggedright\arraybackslash}p{24.232mm}>{\raggedright\arraybackslash}p{22.13mm}>{\raggedright\arraybackslash}p{22.13mm}>{\raggedright\arraybackslash}p{22.13mm}>{\raggedright\arraybackslash}p{22.13mm}@{}}
\toprule
\textbf{Person and number} & \textbf{drincan} & \textbf{singan} & \textbf{tōberstan} & \textbf{forweorðan} \\
\midrule
1 sg. & dranc & sang & tōbærst & forwearð \\
2 sg. & drunce & sunge & tōburste & forwurde \\
3 sg. & dranc & sang & tōbærst & forwearð \\
1–3 pl. & druncon & sungon & tōburston & forwurdon \\
\bottomrule
\end{tabular}
\par\endgroup
% course-panel:T12-P03:end
% mdtable:T12:end
% mdparagraph:P033:begin
Bindan has bintst/\allowbreak{}bint in the selected present. Bindest or bindst may occur beside bintst, and bindeð or bindð beside bint. Winnan reduces nn before the singular consonantal endings, giving winst and winð, but retains nn in winne, winnað and the preterite stems. Drincan and singan follow the nasal vowel series without those dental or doubled-n changes.
% mdparagraph:P033:end

% mdparagraph:P034:begin
Helpan has hilpst/\allowbreak{}hilpð; gieldan has gieltst/\allowbreak{}gielt; weorpan has wierpst/\allowbreak{}wierpð. The vowel and the final consonant group both matter. Berstan has bierst in both second and third singular present. Weorðan has selected wierst, with wierðst also valid, and third-singular wierð. Its first/\allowbreak{}third singular preterite wearð retains ð, while wurde and wurdon have d.
% mdparagraph:P034:end

% teaching: NEW-U04-GRAM-FINDAN-EXCEPTION
% mdparagraph:P035:begin
Findan\BookPoint{bk-ch04-learner-b7272f944b5d} has selected present fintst/\allowbreak{}fint. Its first singular preterite is \textbf{fand}, its second singular \textbf{funde}, and its usual third singular \textbf{funde}. Thus the two singular funde forms are distinguished by person or context, not by spelling. Third-singular fand is valid; first-singular funde is also documented. A request for the selected pattern requires the forms in the table, but these alternatives are not grammatical errors.
% mdparagraph:P035:end

% teaching: NEW-U04-GRAM-PREFIXED-LEXEMES
% teaching: NEW-U04-GRAM-PREFIXED-INTRANSITIVE
% mdparagraph:P036:begin
Tōberstan “break up, go to pieces” follows berstan; forweorðan “perish” follows weorðan. Both are used intransitively in the source reading: the ship breaks up and the companions perish. Their complete forms and whole-word meanings are learned separately. The prefixes do not by themselves supply a general rule for creating other verbs or predicting their meanings. Forwierðst is valid beside selected forwierst.
% mdparagraph:P036:end

\reserve{7\baselineskip}
% mdparagraph:P037:begin
\section*{5. Complete recognition references}\BookHeading{nav-ch04-learner-0007}{subsection}{5. Complete recognition references}{1}\BookLabel{bk-ch04-learner-01b371fdc98a}
% mdparagraph:P037:end

% teaching: NEW-U04-GRAM-RARE-REFERENCE
% mdparagraph:P038:begin
The following two exceptional systems are supplied as complete \textbf{recognition references}. Consult the tables for the reference exercise; independent production and unintroduced object constructions are not required.
% mdparagraph:P038:end

% teaching: NEW-U04-VOC-bregdan
% teaching: NEW-U04-CIT-bregdan
% teaching: NEW-U04-VOC-feolan
% teaching: NEW-U04-CIT-feolan
% table: U04-PARTS-REFERENCE
% mdtable:T13:begin
\par\begingroup
\fontsize{10.5}{13.2}\selectfont
\setlength{\LTleft}{0pt}\setlength{\LTright}{\fill}
\noindent
\begin{tabular}{@{}>{\raggedright\arraybackslash}p{22.61mm}>{\raggedright\arraybackslash}p{23.13mm}>{\raggedright\arraybackslash}p{22.2mm}>{\raggedright\arraybackslash}p{22.2mm}>{\raggedright\arraybackslash}p{22.61mm}@{}}
\toprule
\textbf{Infinitive¹} & \textbf{Meaning} & \textbf{Preterite sg.} & \textbf{Preterite pl.} & \textbf{Participle¹} \\
\midrule
bregdan & brandish; draw & brægd & brugdon & brogden \\
fēolan & penetrate & fealh & fulgon & folen \\
\bottomrule
\end{tabular}
\par\endgroup

% mdtable:T13:end
% teaching: NEW-U04-PAR-bregdan
% teaching: NEW-U04-PAR-feolan
\reserve{7\baselineskip}
% mdparagraph:P039:begin
\textbf{Present indicative.}
% mdparagraph:P039:end

% table: U04-PRESENT-REFERENCE
% mdtable:T14:begin
% course-table-style: 2026-10
% course-panel:T14-P01:begin
\par\begingroup
\fontsize{10.5}{13.2}\selectfont
\setlength{\LTleft}{0pt}\setlength{\LTright}{\fill}
\noindent
\begin{tabular}{@{}>{\raggedright\arraybackslash}p{39.459mm}>{\raggedright\arraybackslash}p{39.459mm}>{\raggedright\arraybackslash}p{39.459mm}@{}}
\toprule
\textbf{Person and number} & \textbf{bregdan} & \textbf{fēolan} \\
\midrule
1 sg. & bregde & fēole \\
2 sg. & brītst & filhst \\
3 sg. & brītt & filhð \\
1–3 pl. & bregdað & fēolað \\
\bottomrule
\end{tabular}
\par\endgroup
% course-panel:T14-P01:end
% mdtable:T14:end
\reserve{7\baselineskip}
% mdparagraph:P040:begin
\textbf{Preterite indicative.}
% mdparagraph:P040:end

% table: U04-PRETERITE-REFERENCE
% mdtable:T15:begin
% course-table-style: 2026-10
% course-panel:T15-P01:begin
\par\begingroup
\fontsize{10.5}{13.2}\selectfont
\setlength{\LTleft}{0pt}\setlength{\LTright}{\fill}
\noindent
\begin{tabular}{@{}>{\raggedright\arraybackslash}p{39.459mm}>{\raggedright\arraybackslash}p{39.459mm}>{\raggedright\arraybackslash}p{39.459mm}@{}}
\toprule
\textbf{Person and number} & \textbf{bregdan} & \textbf{fēolan} \\
\midrule
1 sg. & brægd & fealh \\
2 sg. & brugde & fulge \\
3 sg. & brægd & fealh \\
1–3 pl. & brugdon & fulgon \\
\bottomrule
\end{tabular}
\par\endgroup
% course-panel:T15-P01:end
% mdtable:T15:end
% mdparagraph:P041:begin
Bregdan loses g and lengthens the vowel in selected brītst/\allowbreak{}brītt; brīt is also valid in the third singular. In its preterite, brǣd may occur beside brægd, and brūdon beside brugdon. The participle citation brogden also has the variant brōden. These are reference alternatives, not additional recall targets.
% mdparagraph:P041:end

% mdparagraph:P042:begin
Fēolan has the exceptional selected present filhst/\allowbreak{}filhð beside fēole/\allowbreak{}fēolað. Its preterite uses fealh, fulge and fulgon; fǣlon is a documented plural alternative. The participle citation is Cook’s folen. Wright records folgen and fōlen separately. Preserve these distinctions rather than adding or removing a length mark. This complete reference combines the sources described in note 3.
% mdparagraph:P042:end

\reserve{7\baselineskip}
% mdparagraph:P043:begin
\section*{6. Four nouns}\BookHeading{nav-ch04-learner-0008}{subsection}{6. Four nouns}{1}
% mdparagraph:P043:end

% teaching: NEW-U04-GRAM-NEW-NOUNS
% teaching: NEW-U04-VOC-scip
% teaching: NEW-U04-VOC-wind
% teaching: NEW-U04-VOC-strand
% teaching: NEW-U04-VOC-gefera
% table: U04-NOUN-VOCABULARY
% mdtable:T16:begin
\par\begingroup
\fontsize{10.5}{13.2}\selectfont
\setlength{\LTleft}{0pt}\setlength{\LTright}{\fill}
\noindent
\begin{tabular}{@{}>{\raggedright\arraybackslash}p{25.77mm}>{\raggedright\arraybackslash}p{25.77mm}>{\raggedright\arraybackslash}p{25.77mm}>{\raggedright\arraybackslash}p{38.26mm}@{}}
\toprule
\textbf{Noun} & \textbf{Gender} & \textbf{Meaning} & \textbf{Pattern} \\
\midrule
scip & neuter & ship & short neuter, like hof \\
wind & masculine & wind & strong masculine \\
strand & masculine & shore & strong masculine \\
gefēra & masculine & companion & weak masculine \\
\bottomrule
\end{tabular}
\par\endgroup

% mdtable:T16:end
% teaching: NEW-U04-PAR-scip
% teaching: NEW-U04-PAR-wind
% table: U04-NOUN-SCIP-WIND
% mdtable:T17:begin
% course-table-style: 2026-10
% course-panel:T17-P01:begin
\par\begingroup
\fontsize{10.5}{13.2}\selectfont
\setlength{\LTleft}{0pt}\setlength{\LTright}{\fill}
\noindent
\begin{tabular}{@{}>{\raggedright\arraybackslash}p{22.0mm}>{\raggedright\arraybackslash}p{22.688mm}>{\raggedright\arraybackslash}p{22.688mm}>{\raggedright\arraybackslash}p{22.688mm}>{\raggedright\arraybackslash}p{22.688mm}@{}}
\toprule
\textbf{Case} & \textbf{scip singular} & \textbf{scip plural} & \textbf{wind singular} & \textbf{wind plural} \\
\midrule
Nominative & scip & scipu & wind & windas \\
Genitive & scipes & scipa & windes & winda \\
Dative & scipe & scipum & winde & windum \\
Accusative & scip & scipu & wind & windas \\
\bottomrule
\end{tabular}
\par\endgroup
% course-panel:T17-P01:end
% mdtable:T17:end
% teaching: NEW-U04-PAR-strand
% teaching: NEW-U04-PAR-gefera
% table: U04-NOUN-STRAND-GEFERA
% mdtable:T18:begin
% course-table-style: 2026-10
% course-panel:T18-P01:begin
\par\begingroup
\fontsize{10.5}{13.2}\selectfont
\setlength{\LTleft}{0pt}\setlength{\LTright}{\fill}
\noindent
\begin{tabular}{@{}>{\raggedright\arraybackslash}p{22.0mm}>{\raggedright\arraybackslash}p{22.688mm}>{\raggedright\arraybackslash}p{22.688mm}>{\raggedright\arraybackslash}p{22.688mm}>{\raggedright\arraybackslash}p{22.688mm}@{}}
\toprule
\textbf{Case} & \textbf{strand singular} & \textbf{strand plural} & \textbf{gefēra singular} & \textbf{gefēra plural} \\
\midrule
Nominative & strand & strandas & gefēra & gefēran \\
Genitive & strandes & stranda & gefēran & gefērena \\
Dative & strande & strandum & gefēran & gefērum \\
Accusative & strand & strandas & gefēran & gefēran \\
\bottomrule
\end{tabular}
\par\endgroup
% course-panel:T18-P01:end
% mdtable:T18:end
% mdparagraph:P044:begin
Scip\BookPoint{bk-ch04-learner-c034d18fad71} has a short stem syllable and therefore nominative/\allowbreak{}accusative plural scipu. The initial sc does not make the vowel syllable heavy: both consonants precede the vowel. Wind and strand retain their ordinary strong masculine stems. Gefēra has the weak oblique singular and nominative/\allowbreak{}accusative plural gefēran, but genitive plural gefērena and dative plural gefērum. The same written form gefēran can have several analyses; its clause supplies the choice.
% mdparagraph:P044:end

\reserve{7\baselineskip}
% mdparagraph:P045:begin
\section*{7. Ordinary clauses}\BookHeading{nav-ch04-learner-0009}{subsection}{7. Ordinary clauses}{1}\BookLabel{bk-ch04-learner-0e4e42b78c4b}
% mdparagraph:P045:end

% teaching: NEW-U04-GRAM-DIRECT-OBJECTS
% mdparagraph:P046:begin
The following selected uses take an \textbf{accusative object}: drīfan with the being driven, bindan with the being bound, cēosan with the person or thing chosen, findan with the person or thing found, and drincan with the liquid drunk. The verb’s particular use determines its construction. These object patterns do not follow merely from membership in a strong class.
% mdparagraph:P046:end

% mdparagraph:P047:begin
The eight examples in this section are newly composed teaching material. The other verbs in the tables may be practised as labelled finite forms without assuming additional object or complement constructions.
% mdparagraph:P047:end

% example: NEW-U04-EXAMPLE-01
% mdparagraph:P048:begin
Ic drīfe ðā nīetenu.
% mdparagraph:P048:end

% mdparagraph:P049:begin
I drive the cattle.
% mdparagraph:P049:end

% mdparagraph:P050:begin
Nīetenu is the accusative plural object. The singular first-person subject takes drīfe.
% mdparagraph:P050:end

% example: NEW-U04-EXAMPLE-03
% mdparagraph:P051:begin
Ðū cīest ðæt scip.
% mdparagraph:P051:end

% mdparagraph:P052:begin
You choose the ship.
% mdparagraph:P052:end

% mdparagraph:P053:begin
Scip is the neuter accusative singular object. Cīest is the selected second singular present of cēosan.
% mdparagraph:P053:end

% example: NEW-U04-EXAMPLE-04
% mdparagraph:P054:begin
Se gefēra dranc wæter.
% mdparagraph:P054:end

% mdparagraph:P055:begin
The companion drank water.
% mdparagraph:P055:end

% mdparagraph:P056:begin
The companion is the subject; wæter is the accusative object. Dranc is the third singular preterite.
% mdparagraph:P056:end

% teaching: NEW-U04-GRAM-INTRANSITIVE-ACTIONS
% mdparagraph:P057:begin
Rīdan “ride” and singan “sing” may stand without a nominal object. A familiar adverb may qualify the action. Riding here has its ordinary horseback sense; these clauses introduce no direction or means phrase.
% mdparagraph:P057:end

% example: NEW-U04-EXAMPLE-02
% mdparagraph:P058:begin
Ðū rītst and hē rītt.
% mdparagraph:P058:end

% mdparagraph:P059:begin
You ride, and he rides.
% mdparagraph:P059:end

% mdparagraph:P060:begin
Each singular subject has its own finite verb. Both forms preserve the dental changes of rīdan.
% mdparagraph:P060:end

% example: NEW-U04-EXAMPLE-08
% mdparagraph:P061:begin
Wit sungon wel.
% mdparagraph:P061:end

% mdparagraph:P062:begin
We two sang well.
% mdparagraph:P062:end

% mdparagraph:P063:begin
The dual subject takes the plural preterite sungon. Wel qualifies the singing.
% mdparagraph:P063:end

% teaching: NEW-U04-GRAM-HELPAN-DATIVE
% mdparagraph:P064:begin
The\BookPoint{bk-ch04-learner-97ceb9686099} selected construction of \textbf{helpan} puts the recipient of help in the \textbf{dative}. A person or another being helped can occupy this position. The genitive also occurs with this verb; it is valid Old English but differs from the dative construction requested here. The use with birds in the cumulative practice is an ordinary application of the same helping sense.
% mdparagraph:P064:end

% example: NEW-U04-EXAMPLE-05
% mdparagraph:P065:begin
Hē hilpð ðǣm męnn.
% mdparagraph:P065:end

% mdparagraph:P066:begin
He helps the man.
% mdparagraph:P066:end

% mdparagraph:P067:begin
Męnn is dative singular, marked also by ðǣm. It is the recipient of help, not an accusative object.
% mdparagraph:P067:end

% teaching: NEW-U04-GRAM-BECOME-PREDICATE
% mdparagraph:P068:begin
Ordinary\BookPoint{bk-ch04-learner-bfb873b12018} \textbf{weorðan} means “become”. Its predicate adjective is nominative and agrees with the subject in gender and number, as a predicate with the copula does. This ordinary use is distinct from the passive expressions previously supplied in readings.
% mdparagraph:P068:end

% example: NEW-U04-EXAMPLE-06
% mdparagraph:P069:begin
Se gefēra wearð glæd.
% mdparagraph:P069:end

% mdparagraph:P070:begin
The companion became glad.
% mdparagraph:P070:end

% mdparagraph:P071:begin
Glæd is strong masculine nominative singular, agreeing with gefēra. Wearð is the third singular preterite of weorðan.
% mdparagraph:P071:end

% teaching: NEW-U04-GRAM-SIMPLE-NEGATION
% mdparagraph:P072:begin
Place \textbf{ne} before the noncontracting strong finite forms used here, just as before the ordinary weak forms already learned. The form of the verb otherwise keeps its person, number and tense. Contracted negative copulas and other negative systems are not part of this rule.
% mdparagraph:P072:end

% example: NEW-U04-EXAMPLE-07
% mdparagraph:P073:begin
Ic ne finde ðæt scip.
% mdparagraph:P073:end

% mdparagraph:P074:begin
I do not find the ship.
% mdparagraph:P074:end

% mdparagraph:P075:begin
Ne negates the first singular present finde. The object remains accusative.
% mdparagraph:P075:end

\reserve{7\baselineskip}
% mdparagraph:P076:begin
\section*{8. Practice}\BookHeading{nav-ch04-learner-0010}{subsection}{8. Practice}{1}
% mdparagraph:P076:end

% mdparagraph:P077:begin
All exercises in this unit are new teaching material. Principal-part citations and the two reference paradigms remain available where the prompt calls for them. Other production tasks use the selected forms and the ordinary constructions taught above.
% mdparagraph:P077:end

\reserve{7\baselineskip}
% mdparagraph:P078:begin
\subsection*{Principal parts — P1–P4}\BookHeading{nav-ch04-learner-0011}{subsubsection}{Principal parts — P1–P4}{0}
% mdparagraph:P078:end

\reserve{5\baselineskip}
\par\begin{minipage}{\linewidth}\raggedright
% mdparagraph:P079:begin
\textbf{P1.}
% mdparagraph:P079:end

% item: NEW-U04-P01
% mdparagraph:P080:begin
The printed entry gives drīfan, drāf, drifon, drifen. Label the four principal-part positions. Which two entries are finite indicative forms, and which two are nonfinite citations?
% mdparagraph:P080:end

\end{minipage}\par\addvspace{4pt}

\par\begin{minipage}{\linewidth}\raggedright
% mdparagraph:P081:begin
\textbf{P2.}
% mdparagraph:P081:end

% item: NEW-U04-P02
% mdparagraph:P082:begin
Using the printed principal-part table, identify the strong class of rīdan, bēodan and bindan. For each verb give the second-person singular preterite and state which principal-part stem it uses.
% mdparagraph:P082:end

\end{minipage}\par\addvspace{4pt}

\par\begin{minipage}{\linewidth}\raggedright
% mdparagraph:P083:begin
\textbf{P3.}
% mdparagraph:P083:end

% item: NEW-U04-P03
% mdparagraph:P084:begin
Give every person and number represented by the finite forms drāf and drifon in the selected indicative paradigm. Explain why the supplied citation drifen is not another person of that paradigm.
% mdparagraph:P084:end

\end{minipage}\par\addvspace{4pt}

\par\begin{minipage}{\linewidth}\raggedright
% mdparagraph:P085:begin
\textbf{P4.}
% mdparagraph:P085:end

% item: NEW-U04-P04
% mdparagraph:P086:begin
Compare ðū ride and ðū hīerdest. Give the person, number and tense of each verb, and explain the different formation of their preterites.
% mdparagraph:P086:end

\end{minipage}\par\addvspace{4pt}

\reserve{7\baselineskip}
% mdparagraph:P087:begin
\subsection*{Class I — A1–A6}\BookHeading{nav-ch04-learner-0012}{subsubsection}{Class I — A1–A6}{0}
% mdparagraph:P087:end

\reserve{5\baselineskip}
\par\begin{minipage}{\linewidth}\raggedright
% mdparagraph:P088:begin
\textbf{A1.}
% mdparagraph:P088:end

% item: NEW-U04-A01
% mdparagraph:P089:begin
Give the complete selected present and preterite indicative of drīfan. Label all three persons in the singular and plural even where forms coincide.
% mdparagraph:P089:end

\end{minipage}\par\addvspace{4pt}

\par\begin{minipage}{\linewidth}\raggedright
% mdparagraph:P090:begin
\textbf{A2.}
% mdparagraph:P090:end

% item: NEW-U04-A02
% mdparagraph:P091:begin
For rīdan and wrītan, give the first-, second- and third-person singular present, then the first- and second-person singular preterite and the plural preterite. Explain the selected present dental forms.
% mdparagraph:P091:end

\end{minipage}\par\addvspace{4pt}

\par\begin{minipage}{\linewidth}\raggedright
% mdparagraph:P092:begin
\textbf{A3.}
% mdparagraph:P092:end

% item: NEW-U04-A03
% mdparagraph:P093:begin
For rīsan and snīðan, give the second- and third-person singular present and the second-person singular and plural preterite. Identify the consonant change in snīðan.
% mdparagraph:P093:end

\end{minipage}\par\addvspace{4pt}

\par\begin{minipage}{\linewidth}\raggedright
% mdparagraph:P094:begin
\textbf{A4.}
% mdparagraph:P094:end

% item: NEW-U04-A04
% mdparagraph:P095:begin
Give the complete selected present and preterite indicative of wrēon. Keep its contracted present forms distinct from its preterite stems.
% mdparagraph:P095:end

\end{minipage}\par\addvspace{4pt}

\par\begin{minipage}{\linewidth}\raggedright
% mdparagraph:P096:begin
\textbf{A5.}
% mdparagraph:P096:end

% item: NEW-U04-A05
% mdparagraph:P097:begin
Translate “I rode. You rode. He rode. We rode.” Address one person in the second sentence. Then give the four corresponding present-tense sentences.
% mdparagraph:P097:end

\end{minipage}\par\addvspace{4pt}

\par\begin{minipage}{\linewidth}\raggedright
% mdparagraph:P098:begin
\textbf{A6.}
% mdparagraph:P098:end

% item: NEW-U04-A06
% mdparagraph:P099:begin
With drīfan, give the present forms agreeing with ic, wē, ðū and hē. Then give the corresponding preterite forms. Keep the person and number of each subject unchanged.
% mdparagraph:P099:end

\end{minipage}\par\addvspace{4pt}

\reserve{7\baselineskip}
% mdparagraph:P100:begin
\subsection*{Class II — B1–B6}\BookHeading{nav-ch04-learner-0013}{subsubsection}{Class II — B1–B6}{0}
% mdparagraph:P100:end

\reserve{5\baselineskip}
\par\begin{minipage}{\linewidth}\raggedright
% mdparagraph:P101:begin
\textbf{B1.}
% mdparagraph:P101:end

% item: NEW-U04-B01
% mdparagraph:P102:begin
Give the complete selected present and preterite indicative of bēodan. Identify the present forms whose vowel differs from that of the infinitive citation.
% mdparagraph:P102:end

\end{minipage}\par\addvspace{4pt}

\par\begin{minipage}{\linewidth}\raggedright
% mdparagraph:P103:begin
\textbf{B2.}
% mdparagraph:P103:end

% item: NEW-U04-B02
% mdparagraph:P104:begin
For būgan, give the first-, second- and third-person singular present and the second-person singular and plural preterite. Contrast its present vowel change with bēodan.
% mdparagraph:P104:end

\end{minipage}\par\addvspace{4pt}

\par\begin{minipage}{\linewidth}\raggedright
% mdparagraph:P105:begin
\textbf{B3.}
% mdparagraph:P105:end

% item: NEW-U04-B03
% mdparagraph:P106:begin
For cēosan, give the second- and third-person singular present, the first- and second-person singular preterite and the plural preterite. Explain where the selected s and r occur.
% mdparagraph:P106:end

\end{minipage}\par\addvspace{4pt}

\par\begin{minipage}{\linewidth}\raggedright
% mdparagraph:P107:begin
\textbf{B4.}
% mdparagraph:P107:end

% item: NEW-U04-B04
% mdparagraph:P108:begin
For sēoðan and drēogan, give the second- and third-person singular present, the first- and second-person singular preterite and the plural preterite. Keep the two verbs’ consonant patterns distinct.
% mdparagraph:P108:end

\end{minipage}\par\addvspace{4pt}

\par\begin{minipage}{\linewidth}\raggedright
% mdparagraph:P109:begin
\textbf{B5.}
% mdparagraph:P109:end

% item: NEW-U04-B05
% mdparagraph:P110:begin
Give the complete selected present and preterite indicative of tēon, “draw”. Use the Class II verb assigned in this unit.
% mdparagraph:P110:end

\end{minipage}\par\addvspace{4pt}

\par\begin{minipage}{\linewidth}\raggedright
% mdparagraph:P111:begin
\textbf{B6.}
% mdparagraph:P111:end

% item: NEW-U04-B06
% mdparagraph:P112:begin
Compare the assigned wrēon and tēon tables. Give each verb’s first-person singular present, first-person singular preterite, second-person singular preterite and plural preterite. Explain why the similar infinitive endings do not make them members of the same strong class.
% mdparagraph:P112:end

\end{minipage}\par\addvspace{4pt}

\reserve{7\baselineskip}
% mdparagraph:P113:begin
\subsection*{Class III — C1–C8}\BookHeading{nav-ch04-learner-0014}{subsubsection}{Class III — C1–C8}{0}
% mdparagraph:P113:end

\reserve{5\baselineskip}
\par\begin{minipage}{\linewidth}\raggedright
% mdparagraph:P114:begin
\textbf{C1.}
% mdparagraph:P114:end

% item: NEW-U04-C01
% mdparagraph:P115:begin
Give the complete selected present and preterite indicative of bindan. Explain the selected second- and third-person singular present forms.
% mdparagraph:P115:end

\end{minipage}\par\addvspace{4pt}

\par\begin{minipage}{\linewidth}\raggedright
% mdparagraph:P116:begin
\textbf{C2.}
% mdparagraph:P116:end

% item: NEW-U04-C02
% mdparagraph:P117:begin
For helpan and gieldan, give the second- and third-person singular present, the first- and second-person singular preterite and the plural preterite. Preserve each verb’s selected stem vowels.
% mdparagraph:P117:end

\end{minipage}\par\addvspace{4pt}

\par\begin{minipage}{\linewidth}\raggedright
% mdparagraph:P118:begin
\textbf{C3.}
% mdparagraph:P118:end

% item: NEW-U04-C03
% mdparagraph:P119:begin
For weorpan and berstan, give the second- and third-person singular present, the first- and second-person singular preterite and the plural preterite. Explain the difference between their first preterite stems.
% mdparagraph:P119:end

\end{minipage}\par\addvspace{4pt}

\par\begin{minipage}{\linewidth}\raggedright
% mdparagraph:P120:begin
\textbf{C4.}
% mdparagraph:P120:end

% item: NEW-U04-C04
% mdparagraph:P121:begin
Give the complete selected present and preterite indicative of weorðan. Identify the cells with ð and those with d.
% mdparagraph:P121:end

\end{minipage}\par\addvspace{4pt}

\par\begin{minipage}{\linewidth}\raggedright
% mdparagraph:P122:begin
\textbf{C5.}
% mdparagraph:P122:end

% item: NEW-U04-C05
% mdparagraph:P123:begin
Give the complete selected present and preterite indicative of findan. State the ordinary third-person singular preterite selected here and its documented alternative. Then give the second- and third-person singular present and plural preterite of winnan.
% mdparagraph:P123:end

\end{minipage}\par\addvspace{4pt}

\par\begin{minipage}{\linewidth}\raggedright
% mdparagraph:P124:begin
\textbf{C6.}
% mdparagraph:P124:end

% item: NEW-U04-C06
% mdparagraph:P125:begin
For drincan and singan, give the second- and third-person singular present, the first- and second-person singular preterite and the plural preterite. State which Class III pattern they follow.
% mdparagraph:P125:end

\end{minipage}\par\addvspace{4pt}

\par\begin{minipage}{\linewidth}\raggedright
% mdparagraph:P126:begin
\textbf{C7.}
% mdparagraph:P126:end

% item: NEW-U04-C07
% mdparagraph:P127:begin
Give the complete selected present and preterite indicative of tōberstan and forweorðan. Identify the base pattern used by each, and give the separately learned meaning of each whole verb.
% mdparagraph:P127:end

\end{minipage}\par\addvspace{4pt}

\par\begin{minipage}{\linewidth}\raggedright
% mdparagraph:P128:begin
\textbf{C8.}
% mdparagraph:P128:end

% item: NEW-U04-C08
% mdparagraph:P129:begin
Consult the complete reference tables of bregdan and fēolan. Identify their third-person singular present and plural preterite entries and the special changes noted with them. These are recognition reference forms.
% mdparagraph:P129:end

\end{minipage}\par\addvspace{4pt}

\reserve{7\baselineskip}
% mdparagraph:P130:begin
\subsection*{Nouns — N1–N3}\BookHeading{nav-ch04-learner-0015}{subsubsection}{Nouns — N1–N3}{0}
% mdparagraph:P130:end

\reserve{5\baselineskip}
\par\begin{minipage}{\linewidth}\raggedright
% mdparagraph:P131:begin
\textbf{N1.}
% mdparagraph:P131:end

% item: NEW-U04-N01
% mdparagraph:P132:begin
Give the nominative, genitive, dative and accusative singular and plural of scip and wind. Explain their different nominative and accusative plural endings.
% mdparagraph:P132:end

\end{minipage}\par\addvspace{4pt}

\par\begin{minipage}{\linewidth}\raggedright
% mdparagraph:P133:begin
\textbf{N2.}
% mdparagraph:P133:end

% item: NEW-U04-N02
% mdparagraph:P134:begin
Give the nominative, genitive, dative and accusative singular and plural of strand and gefēra. Explain why gefēran alone does not identify a unique case and number.
% mdparagraph:P134:end

\end{minipage}\par\addvspace{4pt}

\par\begin{minipage}{\linewidth}\raggedright
% mdparagraph:P135:begin
\textbf{N3.}
% mdparagraph:P135:end

% item: NEW-U04-N03
% mdparagraph:P136:begin
Translate “The companions found two ships on the shore. The good man’s ship was small.” In the second sentence put the limiting genitive before the head noun.
% mdparagraph:P136:end

\end{minipage}\par\addvspace{4pt}

\reserve{7\baselineskip}
% mdparagraph:P137:begin
\subsection*{Clauses — D1–D10}\BookHeading{nav-ch04-learner-0016}{subsubsection}{Clauses — D1–D10}{0}
% mdparagraph:P137:end

\reserve{5\baselineskip}
\par\begin{minipage}{\linewidth}\raggedright
% mdparagraph:P138:begin
\textbf{D1.}
% mdparagraph:P138:end

% item: NEW-U04-D01
% mdparagraph:P139:begin
Translate “The companions choose the small dwelling. We choose the great ship.”
% mdparagraph:P139:end

\end{minipage}\par\addvspace{4pt}

\par\begin{minipage}{\linewidth}\raggedright
% mdparagraph:P140:begin
\textbf{D2.}
% mdparagraph:P140:end

% item: NEW-U04-D02
% mdparagraph:P141:begin
Translate “The man found the good man’s ship. He bound the animal.” Put the limiting genitive before its head noun.
% mdparagraph:P141:end

\end{minipage}\par\addvspace{4pt}

\par\begin{minipage}{\linewidth}\raggedright
% mdparagraph:P142:begin
\textbf{D3.}
% mdparagraph:P142:end

% item: NEW-U04-D03
% mdparagraph:P143:begin
Translate “We two drank the water, and you two sang well.” Retain both dual subjects.
% mdparagraph:P143:end

\end{minipage}\par\addvspace{4pt}

\par\begin{minipage}{\linewidth}\raggedright
% mdparagraph:P144:begin
\textbf{D4.}
% mdparagraph:P144:end

% item: NEW-U04-D04
% mdparagraph:P145:begin
Change the subjects to the corresponding plurals and make the verbs agree: Se mann bint ðæt dēor. Se fugol singð wel. Keep the object of the first sentence singular.
% mdparagraph:P145:end

\end{minipage}\par\addvspace{4pt}

\par\begin{minipage}{\linewidth}\raggedright
% mdparagraph:P146:begin
\textbf{D5.}
% mdparagraph:P146:end

% item: NEW-U04-D05
% mdparagraph:P147:begin
Change Ic rīde wīde and finde ān hof to the preterite. Preserve the first-person subject, the object and the coordination.
% mdparagraph:P147:end

\end{minipage}\par\addvspace{4pt}

\par\begin{minipage}{\linewidth}\raggedright
% mdparagraph:P148:begin
\textbf{D6.}
% mdparagraph:P148:end

% item: NEW-U04-D06
% mdparagraph:P149:begin
Here hē and the companion are different people. Change Hē sang wel, and se gefēra hīerde hine to the present. Identify the subject of each verb and the referent of hine.
% mdparagraph:P149:end

\end{minipage}\par\addvspace{4pt}

\par\begin{minipage}{\linewidth}\raggedright
% mdparagraph:P150:begin
\textbf{D7.}
% mdparagraph:P150:end

% item: NEW-U04-D07
% mdparagraph:P151:begin
Translate “The companion became glad. The companions become glad.” Explain the case and number of the predicate adjectives.
% mdparagraph:P151:end

\end{minipage}\par\addvspace{4pt}

\par\begin{minipage}{\linewidth}\raggedright
% mdparagraph:P152:begin
\textbf{D8.}
% mdparagraph:P152:end

% item: NEW-U04-D08
% mdparagraph:P153:begin
Translate “I helped the old man. The old man helped me.” Use the case required by the taught helpan construction.
% mdparagraph:P153:end

\end{minipage}\par\addvspace{4pt}

\par\begin{minipage}{\linewidth}\raggedright
% mdparagraph:P154:begin
\textbf{D9.}
% mdparagraph:P154:end

% item: NEW-U04-D09
% mdparagraph:P155:begin
Translate “The companions did not find the ship. I do not choose the great dwelling.”
% mdparagraph:P155:end

\end{minipage}\par\addvspace{4pt}

\par\begin{minipage}{\linewidth}\raggedright
% mdparagraph:P156:begin
\textbf{D10.}
% mdparagraph:P156:end

% item: NEW-U04-D10
% mdparagraph:P157:begin
Translate “You did not drink the water. You drink water.” Address one person in the first sentence and several people in the second.
% mdparagraph:P157:end

\end{minipage}\par\addvspace{4pt}

\reserve{7\baselineskip}
% mdparagraph:P158:begin
\section*{9. Controlled reading: Two companions choose a dwelling}\BookHeading{nav-ch04-learner-0017}{subsection}{9. Controlled reading: Two companions choose a dwelling}{1}
% mdparagraph:P158:end

% teaching: NEW-U04-GRAM-PAST-NARRATION
% mdparagraph:P159:begin
This account is newly composed. The two companions are the two men of the account; one speaks to the other at the end. Its preterites narrate the successive actions. The quotation shifts to present statements and to the speaker’s first person and addressee’s second person.
% mdparagraph:P159:end

% composition: NEW-U04-TEXT-COMPANIONS-01
% mdparagraph:P160:begin
\par\begin{minipage}{\linewidth}\raggedright
\readingtext{Twēgen gefēran ridon wīde. Hīe fundon twā hofu. Ān hof wæs micel, and ān hof wæs lȳtel. Hīe curon ðæt lȳtle hof. Hīe hæfdon męte and wæter. Hīe druncon ðæt wæter and geręston hīe on ðǣm lȳtlan hofe. Ān mann sægde: “Wit habbað gōd hof. Ic singe, and ðū hīerst mē.”}
\end{minipage}\par
% mdparagraph:P160:end

\reserve{7\baselineskip}
% mdparagraph:P161:begin
\subsection*{Controlled reading — L1–L6}\BookHeading{nav-ch04-learner-0018}{subsubsection}{Controlled reading — L1–L6}{0}
% mdparagraph:P161:end

\reserve{5\baselineskip}
\par\begin{minipage}{\linewidth}\raggedright
% mdparagraph:P162:begin
\textbf{L1.}
% mdparagraph:P162:end

% item: NEW-U04-L01
% mdparagraph:P163:begin
Translate the complete new account “Two companions choose a dwelling”. Preserve the sequence of actions and the speakers in the closing quotation.
% mdparagraph:P163:end

\end{minipage}\par\addvspace{4pt}

\par\begin{minipage}{\linewidth}\raggedright
% mdparagraph:P164:begin
\textbf{L2.}
% mdparagraph:P164:end

% item: NEW-U04-L02
% mdparagraph:P165:begin
List the strong finite verbs in the account. Give each verb’s headword, class, person, number and tense.
% mdparagraph:P165:end

\end{minipage}\par\addvspace{4pt}

\par\begin{minipage}{\linewidth}\raggedright
% mdparagraph:P166:begin
\textbf{L3.}
% mdparagraph:P166:end

% item: NEW-U04-L03
% mdparagraph:P167:begin
For ridon, fundon, curon and druncon in the account, give the corresponding third-person singular preterite and the second-person singular preterite. Preserve the selected lexical exceptions.
% mdparagraph:P167:end

\end{minipage}\par\addvspace{4pt}

\par\begin{minipage}{\linewidth}\raggedright
% mdparagraph:P168:begin
\textbf{L4.}
% mdparagraph:P168:end

% item: NEW-U04-L04
% mdparagraph:P169:begin
Which dwelling do the companions choose? Who sings in the quotation, and whom does the addressee hear? Explain the reference of wit, ic, ðū and mē.
% mdparagraph:P169:end

\end{minipage}\par\addvspace{4pt}

\par\begin{minipage}{\linewidth}\raggedright
% mdparagraph:P170:begin
\textbf{L5.}
% mdparagraph:P170:end

% item: NEW-U04-L05
% mdparagraph:P171:begin
Retell the first four sentences of the account in the present. Preserve the people, things, case relations and adjective endings.
% mdparagraph:P171:end

\end{minipage}\par\addvspace{4pt}

\par\begin{minipage}{\linewidth}\raggedright
% mdparagraph:P172:begin
\textbf{L6.}
% mdparagraph:P172:end

% item: NEW-U04-L06
% mdparagraph:P173:begin
Replace the closing quotation with “We have a good dwelling. We sing, and you hear us.” Let a group of speakers address another group; use plural personal pronouns.
% mdparagraph:P173:end

\end{minipage}\par\addvspace{4pt}

\reserve{7\baselineskip}
% mdparagraph:P174:begin
\section*{10. Source reading: The shipwreck and the survivor}\BookHeading{nav-ch04-learner-0019}{subsection}{10. Source reading: The shipwreck and the survivor}{1}
% mdparagraph:P174:end

% mdparagraph:P175:begin
This original extract from \textbf{Apollonius of Tyre} gives a complete storm-to-shore episode. The two Class III forms tōbærst and forwurdon use the systems learned here. The additional words and constructions are supplied in the local help and may be consulted throughout these tasks.
% mdparagraph:P175:end

% mdparagraph:P176:begin
\textbf{Setting.} Apollonius has embarked for Cyrene. A sudden storm has already disturbed the sea.
% mdparagraph:P176:end

% source: FB-XII-APOLLONIUS-STORM-SHORE-01
% mdparagraph:P177:begin
\par\begin{minipage}{\linewidth}\raggedright
\readingtext{Ðǣr-tō-ēacan cōmon ēastnorðerne windas, and se angrīslica sūðwesterna wind him ongēan stōd, and ðæt scip eall tōbærst on ðisse ęgeslican hrēohnesse. Apollonius gefēran ealle forwurdon tō dēaðe, and Apollonius āna becōm mid sunde tō Pentapolim ðǣm Cyreniscan lande, and ðǣr ūpēode on ðǣm strande.}
\end{minipage}\par
% mdparagraph:P177:end

\reserve{7\baselineskip}
% mdparagraph:P178:begin
\subsection*{Local help}\BookHeading{nav-ch04-learner-0020}{subsubsection}{Local help}{0}
% mdparagraph:P178:end

% mdparagraph:P179:begin
The case analyses below apply to the actual phrases. The four supplied preterites cōmon, stōd, becōm and ūpēode do not yet require complete paradigms or independent recall.
% mdparagraph:P179:end

\reserve{5\baselineskip}
% mdparagraph:P180:begin
\textbf{Ðǣr-tō-ēacan}
% mdparagraph:P180:end

% teaching: NEW-U04-SUPFORM-thaerto-eacan-invariant
% teaching: NEW-U04-SUPSENSE-thaerto-eacan-in-addition
% teaching: NEW-U04-GRAM-SOURCE-CONNECTIVE
% teaching: NEW-U04-SUPPORT-U04-SRC-HELP01
% mdparagraph:P181:begin
“In addition to that; besides.” Further winds join the storm already described in the setting.
% mdparagraph:P181:end

\reserve{5\baselineskip}
% mdparagraph:P182:begin
\textbf{cōmon ēastnorðerne windas}
% mdparagraph:P182:end

% teaching: NEW-U04-SUPFORM-cuman-pret-3-pl-source
% teaching: NEW-U04-SUPFORM-eastnortherne-nom-m-pl-source
% teaching: NEW-U04-SUPSENSE-cuman-come
% teaching: NEW-U04-SUPSENSE-eastnortherne-northeasterly
% teaching: NEW-U04-SUPPORT-U04-SRC-HELP02
% mdparagraph:P183:begin
Cōmon is third-person plural preterite, “came”. Ēastnorðerne means “northeasterly” and is masculine nominative plural, agreeing with windas. The winds are the subject.
% mdparagraph:P183:end

\reserve{5\baselineskip}
% mdparagraph:P184:begin
\textbf{se angrīslica sūðwesterna wind him ongēan stōd}
% mdparagraph:P184:end

% teaching: NEW-U04-SUPFORM-angrislic-nom-m-sg-weak-source
% teaching: NEW-U04-SUPFORM-suthwesterne-nom-m-sg-weak-source
% teaching: NEW-U04-SUPFORM-ongean-postposition-invariant-source
% teaching: NEW-U04-SUPFORM-standan-pret-3-sg-source
% teaching: NEW-U04-SUPSENSE-angrislic-fierce
% teaching: NEW-U04-SUPSENSE-suthwesterne-southwestern
% teaching: NEW-U04-SUPSENSE-ongean-against
% teaching: NEW-U04-SUPSENSE-standan-stand
% teaching: NEW-U04-GRAM-SOURCE-OPPOSITION
% teaching: NEW-U04-SUPPORT-U04-SRC-HELP03
% mdparagraph:P185:begin
“The fierce southwest wind stood against them.” Angrīslica “fierce” and sūðwesterna “southwestern” are weak masculine nominative singular after se. Stōd is third-person singular preterite, “stood”. Ongēan follows its dative pronoun him. The guiding interpretation takes him as the preceding northeast winds; its spelling can also be masculine or neuter singular, so the form alone does not settle its reference.
% mdparagraph:P185:end

\reserve{5\baselineskip}
% mdparagraph:P186:begin
\textbf{ðæt scip eall}
% mdparagraph:P186:end

% teaching: NEW-U04-SUPFORM-eall-nom-n-sg-strong
% teaching: NEW-U04-SUPSENSE-eall-all
% teaching: NEW-U04-GRAM-SOURCE-WHOLE-SHIP
% teaching: NEW-U04-SUPPORT-U04-SRC-HELP04
% mdparagraph:P187:begin
“The whole ship; all the ship.” Eall is strong neuter nominative singular and follows the noun phrase. Its morphology remains the recognition material of Unit 3; this whole-ship expression is supplied here.
% mdparagraph:P187:end

\reserve{5\baselineskip}
% mdparagraph:P188:begin
\textbf{on ðisse ęgeslican hrēohnesse}
% mdparagraph:P188:end

% teaching: NEW-U04-SUPFORM-egeslic-dat-f-sg-weak-source
% teaching: NEW-U04-SUPFORM-hreohnes-dat-sg-source
% teaching: NEW-U04-SUPSENSE-on-storm-circumstance
% teaching: NEW-U04-SUPSENSE-egeslic-terrible
% teaching: NEW-U04-SUPSENSE-hreohnes-tempest
% teaching: NEW-U04-GRAM-SOURCE-IN-STORM
% teaching: NEW-U04-SUPPORT-U04-SRC-HELP05
% mdparagraph:P189:begin
“In this terrible tempest.” Hrēohnesse is feminine dative singular; ęgeslican “terrible” is weak after the demonstrative. Here on describes the surrounding circumstance in which the ship breaks up.
% mdparagraph:P189:end

\reserve{5\baselineskip}
% mdparagraph:P190:begin
\textbf{Apollonius gefēran ealle}
% mdparagraph:P190:end

% teaching: NEW-U04-SUPFORM-apollonius-gen-sg-source
% teaching: NEW-U04-SUPFORM-eall-nom-m-pl-strong
% teaching: NEW-U04-SUPSENSE-apollonius-person
% teaching: NEW-U04-GRAM-SOURCE-PROPER-GENITIVE
% teaching: NEW-U04-GRAM-SOURCE-ALL-COMPANIONS
% teaching: NEW-U04-SUPPORT-U04-SRC-HELP06
% mdparagraph:P191:begin
“All Apollonius’s companions.” The first Apollonius is genitive, despite its ending; Cook explicitly identifies its case. Gefēran is nominative plural. Ealle “all” is strong masculine nominative plural and follows the noun. The all-quantifier phrase is supplied, while the noun and its plural agreement use the paradigms already available.
% mdparagraph:P191:end

\reserve{5\baselineskip}
% mdparagraph:P192:begin
\textbf{forwurdon tō dēaðe}
% mdparagraph:P192:end

% teaching: NEW-U04-SUPFORM-to-preposition-dative-invariant
% teaching: NEW-U04-SUPFORM-death-dat-sg-source
% teaching: NEW-U04-SUPSENSE-to-death-result
% teaching: NEW-U04-SUPSENSE-death-death
% teaching: NEW-U04-GRAM-SOURCE-DEATH-ENDPOINT
% teaching: NEW-U04-SUPPORT-U04-SRC-HELP07
% mdparagraph:P193:begin
“Perished to death”; naturally, “died” or “perished”. Forwurdon is the taught third-person plural preterite. Dēaðe “death” is masculine dative singular. The whole tō phrase gives the endpoint or result and is supplied here.
% mdparagraph:P193:end

\reserve{5\baselineskip}
% mdparagraph:P194:begin
\textbf{Apollonius āna becōm mid sunde}
% mdparagraph:P194:end

% teaching: NEW-U04-SUPFORM-apollonius-nom-sg-source
% teaching: NEW-U04-SUPFORM-an-nom-m-sg-weak-alone-source
% teaching: NEW-U04-SUPFORM-becuman-pret-3-sg-source
% teaching: NEW-U04-SUPFORM-mid-invariant-source
% teaching: NEW-U04-SUPFORM-sund-dat-sg-source
% teaching: NEW-U04-SUPSENSE-an-alone
% teaching: NEW-U04-SUPSENSE-becuman-arrive
% teaching: NEW-U04-SUPSENSE-mid-by-means
% teaching: NEW-U04-SUPSENSE-sund-swimming
% teaching: NEW-U04-GRAM-SOURCE-ALONE
% teaching: NEW-U04-GRAM-SOURCE-BY-SWIMMING
% teaching: NEW-U04-SUPPORT-U04-SRC-HELP08
% mdparagraph:P195:begin
“Apollonius alone arrived by swimming.” This Apollonius is nominative. Āna is weak masculine nominative singular in the sense “alone”, distinct from the ordinary cardinal “one”. Becōm is third-person singular preterite, “came, arrived”. Sunde is the dative singular of a neuter noun meaning “swimming”; mid sunde gives the means. It is not an infinitive or an instrumental case form.
% mdparagraph:P195:end

\reserve{5\baselineskip}
% mdparagraph:P196:begin
\textbf{tō Pentapolim ðǣm Cyreniscan lande}
% mdparagraph:P196:end

% teaching: NEW-U04-SUPFORM-pentapolis-goal-source
% teaching: NEW-U04-SUPFORM-cyrenisc-dat-n-sg-weak-source
% teaching: NEW-U04-SUPFORM-land-dat-sg-source
% teaching: NEW-U04-SUPSENSE-to-destination
% teaching: NEW-U04-SUPSENSE-pentapolis-place
% teaching: NEW-U04-SUPSENSE-cyrenisc-of-cyrene
% teaching: NEW-U04-SUPSENSE-land-country
% teaching: NEW-U04-GRAM-SOURCE-DESTINATION
% teaching: NEW-U04-SUPPORT-U04-SRC-HELP09
% mdparagraph:P197:begin
“To Pentapolis, the Cyrenian land.” The foreign place-name spelling Pentapolim is preserved. The native phrase that explains the destination is neuter dative singular; Cyreniscan is weak and means “Cyrenian, of Cyrene”; lande means “land, country”. This exact destination and apposition are supplied without a new general rule for foreign-name endings.
% mdparagraph:P197:end

\reserve{5\baselineskip}
% mdparagraph:P198:begin
\textbf{ðǣr ūpēode on ðǣm strande}
% mdparagraph:P198:end

% teaching: NEW-U04-SUPFORM-thaer-invariant-source
% teaching: NEW-U04-SUPFORM-upgan-pret-3-sg-source
% teaching: NEW-U04-SUPSENSE-thaer-there
% teaching: NEW-U04-SUPSENSE-upgan-go-up
% teaching: NEW-U04-GRAM-SOURCE-SHORE-ARRIVAL
% teaching: NEW-U04-SUPPORT-U04-SRC-HELP10
% mdparagraph:P199:begin
“There he came up onto the shore.” Ðǣr means “there”; ūpēode is the actual third-person singular preterite, “went/\allowbreak{}came up”. Apollonius remains the understood subject. Strande is dative singular in the printed expression, even though a natural English translation describes movement.
% mdparagraph:P199:end

\reserve{7\baselineskip}
% mdparagraph:P200:begin
\subsection*{Source reading — R1–R5}\BookHeading{nav-ch04-learner-0021}{subsubsection}{Source reading — R1–R5}{0}
% mdparagraph:P200:end

\reserve{5\baselineskip}
\par\begin{minipage}{\linewidth}\raggedright
% mdparagraph:P201:begin
\textbf{R1.}
% mdparagraph:P201:end

% item: NEW-U04-R01
% mdparagraph:P202:begin
Translate the complete source passage “The shipwreck and the survivor”, using the local support. Preserve the contrasting fates of the ship, the companions and Apollonius.
% mdparagraph:P202:end

\end{minipage}\par\addvspace{4pt}

\par\begin{minipage}{\linewidth}\raggedright
% mdparagraph:P203:begin
\textbf{R2.}
% mdparagraph:P203:end

% item: NEW-U04-R02
% mdparagraph:P204:begin
Analyse tōbærst and forwurdon: give each whole verb’s headword, class, person, number and tense, and identify its subject in the passage.
% mdparagraph:P204:end

\end{minipage}\par\addvspace{4pt}

\par\begin{minipage}{\linewidth}\raggedright
% mdparagraph:P205:begin
\textbf{R3.}
% mdparagraph:P205:end

% item: NEW-U04-R03
% mdparagraph:P206:begin
Using the proper-name note, distinguish the first Apollonius from the second. Explain why gefēran takes plural forwurdon rather than a singular verb.
% mdparagraph:P206:end

\end{minipage}\par\addvspace{4pt}

\par\begin{minipage}{\linewidth}\raggedright
% mdparagraph:P207:begin
\textbf{R4.}
% mdparagraph:P207:end

% item: NEW-U04-R04
% mdparagraph:P208:begin
Identify windas, wind, scip, gefēran and strande in their clauses. Give each noun’s case and number; use the supplied phrase notes wherever the construction is new.
% mdparagraph:P208:end

\end{minipage}\par\addvspace{4pt}

\par\begin{minipage}{\linewidth}\raggedright
% mdparagraph:P209:begin
\textbf{R5.}
% mdparagraph:P209:end

% item: NEW-U04-R05
% mdparagraph:P210:begin
How does Apollonius reach land, and where does the passage leave him? Explain the contribution of the supplied āna without treating it as an ordinary cardinal “one”.
% mdparagraph:P210:end

\end{minipage}\par\addvspace{4pt}

\reserve{7\baselineskip}
% mdparagraph:P211:begin
\section*{11. Cumulative review}\BookHeading{nav-ch04-learner-0022}{subsection}{11. Cumulative review}{1}
% mdparagraph:P211:end

\reserve{7\baselineskip}
% mdparagraph:P212:begin
\subsection*{Review — X1–X4}\BookHeading{nav-ch04-learner-0023}{subsubsection}{Review — X1–X4}{0}
% mdparagraph:P212:end

\reserve{5\baselineskip}
\par\begin{minipage}{\linewidth}\raggedright
% mdparagraph:P213:begin
\textbf{X1.}
% mdparagraph:P213:end

% item: NEW-U04-X01
% mdparagraph:P214:begin
Correct the three marked finite forms: Ðū *drāf. Hīe *cēas ðæt scip. Se mann *bindað ðæt dēor. Preserve the given subjects and tenses.
% mdparagraph:P214:end

\end{minipage}\par\addvspace{4pt}

\par\begin{minipage}{\linewidth}\raggedright
% mdparagraph:P215:begin
\textbf{X2.}
% mdparagraph:P215:end

% item: NEW-U04-X02
% mdparagraph:P216:begin
Write four connected Old English sentences about two companions: they ride, find two ships, choose the small ship, and drink water. Narrate all four actions in the preterite and retain the same two people as subjects.
% mdparagraph:P216:end

\end{minipage}\par\addvspace{4pt}

\par\begin{minipage}{\linewidth}\raggedright
% mdparagraph:P217:begin
\textbf{X3.}
% mdparagraph:P217:end

% item: NEW-U04-X03
% mdparagraph:P218:begin
Translate “The old companion helps my birds. My birds sing well. We found his small ship on the shore.” Distinguish the helper, the birds, the finders and the owner.
% mdparagraph:P218:end

\end{minipage}\par\addvspace{4pt}

\par\begin{minipage}{\linewidth}\raggedright
% mdparagraph:P219:begin
\textbf{X4.}
% mdparagraph:P219:end

% item: NEW-U04-X04
% mdparagraph:P220:begin
Analyse Hē funde, Ðū funde and Hīe fundon. Which finite forms distinguish number? What identifies the different persons in the two singular clauses?
% mdparagraph:P220:end

\end{minipage}\par\addvspace{4pt}

\clearpage
% mdparagraph:P221:begin
\section*{12. Source and variant notes}\BookHeading{nav-ch04-learner-0024}{subsection}{12. Source and variant notes}{1}
% mdparagraph:P221:end

\begingroup\fontsize{10}{12.8}\selectfont
\setlength{\parskip}{1.5pt}
% mdparagraph:P222:begin
\itempara{1}{The principal-part, indicative and stem rules follow Cook, \textbf{A First Book in Old English}, §§92–95 and 101–104, printed pp.54–55 and 58–63. The selected tables preserve the edition’s quantities. Full finite cells formed from a checked stem and the ordinary endings are transparent applications of the pattern; the tables do not claim a separate surviving quotation for every person.}
% mdparagraph:P222:end

% mdparagraph:P223:begin
\itempara{2}{Wright’s \textbf{Old English Grammar} corroborates the selected stem series and recorded alternatives. In particular, its Class III discussion permits first- as well as third-person funde. Its discussion of contracted Class I verbs records analogical movement into Class II; the wrēon table here consistently selects Cook’s Class I pattern. Alternatives in the notes are available for recognition and comparison, not additional independent-recall requirements.}
% mdparagraph:P223:end

% mdparagraph:P224:begin
\itempara{3}{The rare fēolan reference combines Cook’s principal parts and ordinary endings with Bosworth–Toller’s filhst/\allowbreak{}filhþ entries. The displayed filhð changes þ to the course’s ð spelling only. The supplement transfers the earlier entry’s passages to fēolan. This is a complete synthesized reference, not a claim that each printed cell is individually attested. Cook’s folen and Wright’s folgen/\allowbreak{}fōlen remain distinct citations.}
% mdparagraph:P224:end

% mdparagraph:P225:begin
\itempara{4}{The ordinary case frames are independently documented: Cook’s direct-object and predicate rules (§§168 and 150), his occurrences of findan, cēosan and singan, and the Bosworth–Toller entries for drīfan, bindan, rīdan, drincan and helpan. The animal and ship/\allowbreak{}dwelling objects in the new examples are ordinary applications of these senses. Attested human dative recipients of help support its ordinary application to birds; the dictionary also records a genitive construction.}
% mdparagraph:P225:end

% mdparagraph:P226:begin
\itempara{5}{The original extract is Cook, \textbf{A First Book in Old English}, Reader XII, \textbf{Apollonius of Tyre}, printed pp.165–166, one-based PDF181–182. It begins within p.165 line 6 and ends at the first period in p.166 line 4. These line counts expand the edition’s printed five-line intervals. The line-crossing compound Ðǣr-tō-ēacan is joined, and note cues are moved below; the Old English spelling and punctuation are preserved. The English setting is editorial context.}
% mdparagraph:P226:end

\reserve{7\baselineskip}
% mdparagraph:P227:begin
\subsection*{Cook’s notes to the extract}\BookHeading{nav-ch04-learner-0025}{subsubsection}{Cook’s notes to the extract}{0}
% mdparagraph:P227:end

% source-note: FB165-n7
% mdparagraph:P228:begin
\textbf{P.165, note 7.} Lat. (verse): Hinc Notus,hinc Boreas,hinc horridus Africus instat.
% mdparagraph:P228:end

% mdparagraph:P229:begin
The Latin verse evokes winds from different quarters; it does not explicitly settle the antecedent of him.
% mdparagraph:P229:end

% source-note: FB166-n1
% mdparagraph:P230:begin
\textbf{P.166, note 1.} The Latin endings of proper nouns are not always a guide to the case (54). Here we have the genitive.
% mdparagraph:P230:end

% source-note: FB166-n2
% mdparagraph:P231:begin
\textbf{P.166, note 2.} See forweorðan.
% mdparagraph:P231:end

% source-note: FB166-n3
% mdparagraph:P232:begin
\textbf{P.166, note 3.} See 79.
% mdparagraph:P232:end

% mdparagraph:P233:begin
Section 79 explains the special “alone” use supplied with āna above.
% mdparagraph:P233:end

\endgroup

% book-source-body:end
\par\label{end-ch04-learner}
\clearpage\endgroup

\BookRecto
\begingroup
\renewcommand{\runningtitle}{Unit 5. Strong verbs IV–VII}
\hyphenpenalty=10000
\exhyphenpenalty=10000
\emergencystretch=3em
\newcommand{\itempara}[2]{\par\addvspace{3pt}\noindent\hangindent=1.8em\hangafter=1\makebox[1.8em][l]{#1.}#2\par}
\newcommand{\reserve}[1]{\par\penalty-100\begingroup\dimen0=\pagegoal\advance\dimen0 by-\pagetotal\ifdim\dimen0<#1\newpage\fi\endgroup}
% book-source-body:begin

\addfontfeatures{Ligatures=TeXOff}
\raggedright
\booktitle{\BookComponent{ch05-learner}{Unit 5. Strong verbs IV–VII}UNIT 5}{Strong verbs IV–VII}

% mdparagraph:P001:begin
This unit completes the principal strong-verb classes with selected present and preterite indicative systems. Four further verbs from Classes I–III provide the forms needed for the complete fifth Cook exercise.
% mdparagraph:P001:end

% mdparagraph:P002:begin
In the tables, \textbf{1}, \textbf{2} and \textbf{3} mean first, second and third person; \textbf{sg.} and \textbf{pl.} mean singular and plural. \textbf{1–3 pl.} gives the same form for all three plural persons. Dual subjects also take plural verbs. Each verb’s paired finite tables cover all six present and all six preterite indicative cells. Preserve the printed quantities and the distinction between e and ę.
% mdparagraph:P002:end

\reserve{7\baselineskip}
% mdparagraph:P003:begin
\section*{1. Principal parts and endings}\BookHeading{nav-ch05-learner-0001}{subsection}{1. Principal parts and endings}{1}\BookLabel{bk-ch05-learner-4385a8053534}
% mdparagraph:P003:end

% mdparagraph:P004:begin
Strong verbs change their stem vowel in the preterite; the weak verbs learned earlier add a dental element. Four principal parts identify the strong stem series. The infinitive names the verb, the singular preterite gives the first/\allowbreak{}third singular stem, the plural preterite gives the plural stem, and the participle completes the cited series.
% mdparagraph:P004:end

% table: U05-ENDINGS
% mdtable:T01:begin
% course-table-style: 2026-10
% course-panel:T01-P01:begin
\par\begingroup
\fontsize{10.5}{13.2}\selectfont
\setlength{\LTleft}{0pt}\setlength{\LTright}{\fill}
\noindent
\begin{tabular}{@{}>{\raggedright\arraybackslash}p{39.459mm}>{\raggedright\arraybackslash}p{39.459mm}>{\raggedright\arraybackslash}p{39.459mm}@{}}
\toprule
\textbf{Person and number} & \textbf{Present} & \textbf{Preterite} \\
\midrule
1 sg. & -e, with the special contracted and -ie forms below & singular preterite stem; no added ending \\
2 sg. & -st or -est, with the listed stem changes & plural preterite stem + -e \\
3 sg. & -ð or -eð, with the listed stem changes & singular preterite stem; no added ending \\
1–3 pl. & -að, with the listed contracted and -iað forms & plural preterite stem + -on \\
\bottomrule
\end{tabular}
\par\endgroup
% course-panel:T01-P01:end
% mdtable:T01:end
% mdparagraph:P005:begin
The \textbf{second singular preterite uses the plural stem}. Thus beran has bær for the first and third singular, but bǣre beside plural bǣron. In other verbs the singular and plural stem vowels coincide; their different endings still identify the cells.
% mdparagraph:P005:end

% mdparagraph:P006:begin
The infinitive and participle columns are \textbf{supplied citations} throughout this unit. Consult them in the principal-part exercises. They do not require independent nonfinite recall or the use of infinitive, participial, perfect or passive constructions. A participle in a four-part entry is not another finite preterite.
% mdparagraph:P006:end

% mdparagraph:P007:begin
The second and third singular present often have a changed vowel and a reduced ending. Short e may appear as i; short a gives æ, short ea gives short ie, long ēa gives long īe, ā gives ǣ, and ō gives ē in the stated paradigms. These changes are lexical patterns, not an instruction to alter every similar-looking vowel. Dental combinations give forms such as -tst and -tt, and doubled stem consonants simplify before consonantal endings. The tables and notes give the resulting forms together.
% mdparagraph:P007:end

% teaching: NEW-U05-GRAM-PREFIXED-LEXEMES
% mdparagraph:P008:begin
\textbf{Gesēon, gescieppan and geswīcan} are learned as separate verbs. Their meanings and forms are stated explicitly. The prefix is not an instruction to create further verbs or to assign a fixed difference in meaning to every prefixed pair. The checked headword \textbf{gesēon} follows Cook’s dictionary ge-sēon, with the dictionary hyphen omitted; the preterite geseah keeps its printed spelling.
% mdparagraph:P008:end

\reserve{7\baselineskip}
% mdparagraph:P009:begin
\section*{2. Class IV}\BookHeading{nav-ch05-learner-0002}{subsection}{2. Class IV}{1}\BookLabel{bk-ch05-learner-8f7dcc852b82}
% mdparagraph:P009:end

% teaching: NEW-U05-GRAM-CLASS-IV
% mdparagraph:P010:begin
The ordinary Class IV series is \textbf{e : æ : ǣ : o}. Beran shows it directly: the singular past has short æ, while the plural-based past has long ǣ. The fourth vowel belongs to the supplied participle.
% mdparagraph:P010:end

% teaching: NEW-U05-VOC-beran
% teaching: NEW-U05-CIT-beran
% teaching: NEW-U05-VOC-niman
% teaching: NEW-U05-CIT-niman
% teaching: NEW-U05-VOC-cuman
% teaching: NEW-U05-CIT-cuman
% table: U05-PARTS-IV
% mdtable:T02:begin
\par\begingroup
\fontsize{10.5}{13.2}\selectfont
\setlength{\LTleft}{0pt}\setlength{\LTright}{\fill}
\noindent
\begin{tabular}{@{}>{\raggedright\arraybackslash}p{23.14mm}>{\raggedright\arraybackslash}p{21.04mm}>{\raggedright\arraybackslash}p{22.72mm}>{\raggedright\arraybackslash}p{22.72mm}>{\raggedright\arraybackslash}p{23.14mm}@{}}
\toprule
\textbf{Infinitive¹} & \textbf{Meaning} & \textbf{Preterite sg.} & \textbf{Preterite pl.} & \textbf{Participle¹} \\
\midrule
beran & carry & bær & bǣron & boren \\
niman & take & nōm & nōmon & numen \\
cuman & come & cōm & cōmon & cumen \\
\bottomrule
\end{tabular}
\par\endgroup

% mdtable:T02:end
% mdparagraph:P011:begin
¹ Infinitive and participle entries are supplied citations. The middle columns identify the finite preterite stems; the complete person assignments follow.
% mdparagraph:P011:end

% teaching: NEW-U05-PAR-beran
% teaching: NEW-U05-PAR-niman
% teaching: NEW-U05-PAR-cuman
\reserve{8\baselineskip}
% mdparagraph:P012:begin
\textbf{Present indicative.}
% mdparagraph:P012:end

% table: U05-PRESENT-IV
% mdtable:T03:begin
% course-table-style: 2026-10
% course-panel:T03-P01:begin
\par\begingroup
\fontsize{10.5}{13.2}\selectfont
\setlength{\LTleft}{0pt}\setlength{\LTright}{\fill}
\noindent
\begin{tabular}{@{}>{\raggedright\arraybackslash}p{30.902mm}>{\raggedright\arraybackslash}p{28.221mm}>{\raggedright\arraybackslash}p{28.221mm}>{\raggedright\arraybackslash}p{28.221mm}@{}}
\toprule
\textbf{Person and number} & \textbf{beran} & \textbf{niman} & \textbf{cuman} \\
\midrule
1 sg. & bere & nime & cume \\
2 sg. & birst & nimst & cymst \\
3 sg. & birð & nimð & cymð \\
1–3 pl. & berað & nimað & cumað \\
\bottomrule
\end{tabular}
\par\endgroup
% course-panel:T03-P01:end
% mdtable:T03:end
\reserve{8\baselineskip}
% mdparagraph:P013:begin
\textbf{Preterite indicative.}
% mdparagraph:P013:end

% table: U05-PRETERITE-IV
% mdtable:T04:begin
% course-table-style: 2026-10
% course-panel:T04-P01:begin
\par\begingroup
\fontsize{10.5}{13.2}\selectfont
\setlength{\LTleft}{0pt}\setlength{\LTright}{\fill}
\noindent
\begin{tabular}{@{}>{\raggedright\arraybackslash}p{30.902mm}>{\raggedright\arraybackslash}p{28.221mm}>{\raggedright\arraybackslash}p{28.221mm}>{\raggedright\arraybackslash}p{28.221mm}@{}}
\toprule
\textbf{Person and number} & \textbf{beran} & \textbf{niman} & \textbf{cuman} \\
\midrule
1 sg. & bær & nōm & cōm \\
2 sg. & bǣre & nōme & cōme \\
3 sg. & bær & nōm & cōm \\
1–3 pl. & bǣron & nōmon & cōmon \\
\bottomrule
\end{tabular}
\par\endgroup
% course-panel:T04-P01:end
% mdtable:T04:end
% teaching: NEW-U05-GRAM-NASAL-IV
% mdparagraph:P014:begin
Learn\BookPoint{bk-ch05-learner-a2089fd14fa6} \textbf{niman} and \textbf{cuman} separately within this class. This course selects Cook’s nōm, nōme, nōmon and cōm, cōme, cōmon. Cuman changes u to y in cymst and cymð; its first singular cume and plural cumað retain u. Niman keeps i in its present forms.
% mdparagraph:P014:end

% mdparagraph:P015:begin
\textbf{Alternatives.} Beran also has bierst and bierð. Wright records nam and nāmon beside the selected nōm and nōmon. Cook gives cwōm and cwōmon beside cōm and cōmon, and the participle citation cymen beside cumen. These alternatives are available for comparison; use the selected table when a task expressly requests it.
% mdparagraph:P015:end

\reserve{7\baselineskip}
% mdparagraph:P016:begin
\section*{3. Class V}\BookHeading{nav-ch05-learner-0003}{subsection}{3. Class V}{1}\BookLabel{bk-ch05-learner-db2ab453e454}
% mdparagraph:P016:end

% teaching: NEW-U05-GRAM-CLASS-V
% mdparagraph:P017:begin
The ordinary Class V series is \textbf{e : æ : ǣ : e}. It differs from Class IV in the participle vowel. Sprecan gives the ordinary series; several other verbs require their own present or preterite changes.
% mdparagraph:P017:end

% teaching: NEW-U05-VOC-sprecan
% teaching: NEW-U05-CIT-sprecan
% teaching: NEW-U05-VOC-cwethan
% teaching: NEW-U05-CIT-cwethan
% teaching: NEW-U05-VOC-giefan
% teaching: NEW-U05-CIT-giefan
% teaching: NEW-U05-VOC-biddan
% teaching: NEW-U05-CIT-biddan
% teaching: NEW-U05-VOC-licgan
% teaching: NEW-U05-CIT-licgan
% teaching: NEW-U05-VOC-sittan
% teaching: NEW-U05-CIT-sittan
% teaching: NEW-U05-VOC-etan
% teaching: NEW-U05-CIT-etan
% teaching: NEW-U05-VOC-seon
% teaching: NEW-U05-CIT-seon
% teaching: NEW-U05-VOC-geseon
% teaching: NEW-U05-CIT-geseon
% table: U05-PARTS-V
% mdtable:T05:begin
\par\begingroup
\fontsize{10.5}{13.2}\selectfont
\setlength{\LTleft}{0pt}\setlength{\LTright}{\fill}
\noindent
\begin{tabular}{@{}>{\raggedright\arraybackslash}p{22.88mm}>{\raggedright\arraybackslash}p{22.05mm}>{\raggedright\arraybackslash}p{22.47mm}>{\raggedright\arraybackslash}p{22.47mm}>{\raggedright\arraybackslash}p{22.88mm}@{}}
\toprule
\textbf{Infinitive¹} & \textbf{Meaning} & \textbf{Preterite sg.} & \textbf{Preterite pl.} & \textbf{Participle¹} \\
\midrule
sprecan & speak & spræc & sprǣcon & sprecen \\
cweðan & say & cwæð & cwǣdon & cweden \\
giefan & give & geaf & gēafon & giefen \\
biddan & ask; request & bæd & bǣdon & beden \\
licgan & lie; be at rest & læg & lǣgon & legen \\
sittan & sit & sæt & sǣton & seten \\
etan & eat & ǣt & ǣton & eten \\
sēon & see & seah & sāwon & sewen \\
gesēon & see & geseah & gesāwon & gesewen \\
\bottomrule
\end{tabular}
\par\endgroup

% mdtable:T05:end
% mdparagraph:P018:begin
¹ Infinitive and participle entries are supplied citations. The middle columns identify the finite preterite stems; the complete person assignments follow.
% mdparagraph:P018:end

% teaching: NEW-U05-PAR-sprecan
% teaching: NEW-U05-PAR-cwethan
% teaching: NEW-U05-PAR-giefan
% teaching: NEW-U05-PAR-biddan
% teaching: NEW-U05-PAR-licgan
% teaching: NEW-U05-PAR-sittan
% teaching: NEW-U05-PAR-etan
% teaching: NEW-U05-PAR-seon
% teaching: NEW-U05-PAR-geseon
\reserve{14\baselineskip}
% mdparagraph:P019:begin
\textbf{Present indicative.}
% mdparagraph:P019:end

% table: U05-PRESENT-V
% mdtable:T06:begin
% course-table-style: 2026-10
% course-panel:T06-P01:begin
\par\begingroup
\fontsize{10.5}{13.2}\selectfont
\setlength{\LTleft}{0pt}\setlength{\LTright}{\fill}
\noindent
\begin{tabular}{@{}>{\raggedright\arraybackslash}p{24.232mm}>{\raggedright\arraybackslash}p{22.13mm}>{\raggedright\arraybackslash}p{22.13mm}>{\raggedright\arraybackslash}p{22.13mm}>{\raggedright\arraybackslash}p{22.13mm}@{}}
\toprule
\textbf{Person and number} & \textbf{sprecan} & \textbf{cweðan} & \textbf{giefan} & \textbf{biddan} \\
\midrule
1 sg. & sprece & cweðe & giefe & bidde \\
2 sg. & spricst & cwist & giefst & bitst \\
3 sg. & spricð & cwið & giefð & bitt \\
1–3 pl. & sprecað & cweðað & giefað & biddað \\
\bottomrule
\end{tabular}
\par\endgroup
% course-panel:T06-P01:end
\par\addvspace{5pt}
% course-panel:T06-P02:begin
\par\begingroup
\fontsize{10.5}{13.2}\selectfont
\setlength{\LTleft}{0pt}\setlength{\LTright}{\fill}
\noindent
\begin{tabular}{@{}>{\raggedright\arraybackslash}p{24.232mm}>{\raggedright\arraybackslash}p{22.13mm}>{\raggedright\arraybackslash}p{22.13mm}>{\raggedright\arraybackslash}p{22.13mm}>{\raggedright\arraybackslash}p{22.13mm}@{}}
\toprule
\textbf{Person and number} & \textbf{licgan} & \textbf{sittan} & \textbf{etan} & \textbf{sēon} \\
\midrule
1 sg. & licge & sitte & ete & sēo \\
2 sg. & ligst & sitst & itst & siehst \\
3 sg. & ligð & sitt & itt & siehð \\
1–3 pl. & licgað & sittað & etað & sēoð \\
\bottomrule
\end{tabular}
\par\endgroup
% course-panel:T06-P02:end
\par\addvspace{5pt}
% course-panel:T06-P03:begin
\par\begingroup
\fontsize{10.5}{13.2}\selectfont
\setlength{\LTleft}{0pt}\setlength{\LTright}{\fill}
\noindent
\begin{tabular}{@{}>{\raggedright\arraybackslash}p{63.342mm}>{\raggedright\arraybackslash}p{57.846mm}@{}}
\toprule
\textbf{Person and number} & \textbf{gesēon} \\
\midrule
1 sg. & gesēo \\
2 sg. & gesiehst \\
3 sg. & gesiehð \\
1–3 pl. & gesēoð \\
\bottomrule
\end{tabular}
\par\endgroup
% course-panel:T06-P03:end
% mdtable:T06:end
\reserve{14\baselineskip}
% mdparagraph:P020:begin
\textbf{Preterite indicative.}
% mdparagraph:P020:end

% table: U05-PRETERITE-V
% mdtable:T07:begin
% course-table-style: 2026-10
% course-panel:T07-P01:begin
\par\begingroup
\fontsize{10.5}{13.2}\selectfont
\setlength{\LTleft}{0pt}\setlength{\LTright}{\fill}
\noindent
\begin{tabular}{@{}>{\raggedright\arraybackslash}p{24.232mm}>{\raggedright\arraybackslash}p{22.13mm}>{\raggedright\arraybackslash}p{22.13mm}>{\raggedright\arraybackslash}p{22.13mm}>{\raggedright\arraybackslash}p{22.13mm}@{}}
\toprule
\textbf{Person and number} & \textbf{sprecan} & \textbf{cweðan} & \textbf{giefan} & \textbf{biddan} \\
\midrule
1 sg. & spræc & cwæð & geaf & bæd \\
2 sg. & sprǣce & cwǣde & gēafe & bǣde \\
3 sg. & spræc & cwæð & geaf & bæd \\
1–3 pl. & sprǣcon & cwǣdon & gēafon & bǣdon \\
\bottomrule
\end{tabular}
\par\endgroup
% course-panel:T07-P01:end
\par\addvspace{5pt}
% course-panel:T07-P02:begin
\par\begingroup
\fontsize{10.5}{13.2}\selectfont
\setlength{\LTleft}{0pt}\setlength{\LTright}{\fill}
\noindent
\begin{tabular}{@{}>{\raggedright\arraybackslash}p{24.232mm}>{\raggedright\arraybackslash}p{22.13mm}>{\raggedright\arraybackslash}p{22.13mm}>{\raggedright\arraybackslash}p{22.13mm}>{\raggedright\arraybackslash}p{22.13mm}@{}}
\toprule
\textbf{Person and number} & \textbf{licgan} & \textbf{sittan} & \textbf{etan} & \textbf{sēon} \\
\midrule
1 sg. & læg & sæt & ǣt & seah \\
2 sg. & lǣge & sǣte & ǣte & sāwe \\
3 sg. & læg & sæt & ǣt & seah \\
1–3 pl. & lǣgon & sǣton & ǣton & sāwon \\
\bottomrule
\end{tabular}
\par\endgroup
% course-panel:T07-P02:end
\par\addvspace{5pt}
% course-panel:T07-P03:begin
\par\begingroup
\fontsize{10.5}{13.2}\selectfont
\setlength{\LTleft}{0pt}\setlength{\LTright}{\fill}
\noindent
\begin{tabular}{@{}>{\raggedright\arraybackslash}p{63.342mm}>{\raggedright\arraybackslash}p{57.846mm}@{}}
\toprule
\textbf{Person and number} & \textbf{gesēon} \\
\midrule
1 sg. & geseah \\
2 sg. & gesāwe \\
3 sg. & geseah \\
1–3 pl. & gesāwon \\
\bottomrule
\end{tabular}
\par\endgroup
% course-panel:T07-P03:end
% mdtable:T07:end
% teaching: NEW-U05-GRAM-CWETHAN-CONSONANTS
% mdparagraph:P021:begin
\textbf{Cweðan}\BookPoint{bk-ch05-learner-7c5f2c14a3b1} has cwist and cwið in the selected second and third singular present. Its past singular cwæð retains ð, while cwǣde and cwǣdon have d; the supplied participle cweden also has d. Keep this consonant change with these stem positions.
% mdparagraph:P021:end

% teaching: NEW-U05-GRAM-PALATAL-V
% mdparagraph:P022:begin
\textbf{Giefan}\BookPoint{bk-ch05-learner-9cc23aad7df2} has the palatal series giefan, geaf, gēafon, giefen. In particular, geaf has short ea, but gēafe and gēafon have long ēa. The present forms retain the selected ie spelling.
% mdparagraph:P022:end

% teaching: NEW-U05-GRAM-J-PRESENT-V
% mdparagraph:P023:begin
\textbf{Biddan, licgan and sittan}\BookPoint{bk-ch05-learner-67a692c44be4} have doubled consonants in the present first singular and plural, but a single stem consonant before the second/\allowbreak{}third singular endings. The resulting forms are bitst/\allowbreak{}bitt, ligst/\allowbreak{}ligð and sitst/\allowbreak{}sitt. Their strong pasts use bæd/\allowbreak{}bǣdon, læg/\allowbreak{}lǣgon and sæt/\allowbreak{}sǣton. These are often called \textbf{j-presents}: the label explains a present-stem pattern, not a weak preterite.
% mdparagraph:P023:end

% teaching: NEW-U05-GRAM-ETAN-EXCEPTION
% mdparagraph:P024:begin
\textbf{Etan} has itst and itt in the present singular. The selected preterite ǣt has long ǣ, as do ǣte and ǣton. Its singular past therefore differs from the ordinary short-æ Class V pattern.
% mdparagraph:P024:end

% teaching: NEW-U05-GRAM-CONTRACT-V
% mdparagraph:P025:begin
\textbf{Sēon}\BookPoint{bk-ch05-learner-852de88e21a6} and \textbf{gesēon} have contracted presents. Their first singular ends in -o and their plural in -oð. The second and third singular have \textbf{short ie}: siehst/\allowbreak{}siehð and gesiehst/\allowbreak{}gesiehð. Their plural-based past stems have āw: sāwe/\allowbreak{}sāwon and gesāwe/\allowbreak{}gesāwon. Preserve the difference between these forms and the singular seah/\allowbreak{}geseah.
% mdparagraph:P025:end

% mdparagraph:P026:begin
\textbf{Alternatives.} Third-singular bit, sit and it occur beside selected bitt, sitt and itt. Cook gives līst and līð beside ligst and ligð, and short æt beside selected ǣt; Wright also gives lāgon beside lǣgon. For sēon, Cook gives the participle segen beside sewen, and Wright records sǣgon and sawen beside sāwon and sewen. The corresponding prefixed citation alternative is gesegen beside gesewen. These spellings do not replace the selected tables or require extra nonfinite recall.
% mdparagraph:P026:end

\reserve{7\baselineskip}
% mdparagraph:P027:begin
\section*{4. Class VI}\BookHeading{nav-ch05-learner-0004}{subsection}{4. Class VI}{1}\BookLabel{bk-ch05-learner-92ffb8f8186c}
% mdparagraph:P027:end

% teaching: NEW-U05-GRAM-CLASS-VI
% mdparagraph:P028:begin
The ordinary Class VI series is \textbf{a : ō : ō : a}. Faran shows it directly. The singular and plural past stems share long ō, and the second singular adds -e to that stem. Present mutation and the special present stems must still be learned separately.
% mdparagraph:P028:end

% teaching: NEW-U05-VOC-faran
% teaching: NEW-U05-CIT-faran
% teaching: NEW-U05-VOC-slean
% teaching: NEW-U05-CIT-slean
% teaching: NEW-U05-VOC-standan
% teaching: NEW-U05-CIT-standan
% teaching: NEW-U05-VOC-hebban
% teaching: NEW-U05-CIT-hebban
% teaching: NEW-U05-VOC-scieppan
% teaching: NEW-U05-CIT-scieppan
% teaching: NEW-U05-VOC-gescieppan
% teaching: NEW-U05-CIT-gescieppan
% teaching: NEW-U05-VOC-swerian
% teaching: NEW-U05-CIT-swerian
% table: U05-PARTS-VI
% mdtable:T08:begin
\par\begingroup
\fontsize{10.5}{13.2}\selectfont
\setlength{\LTleft}{0pt}\setlength{\LTright}{\fill}
\noindent
\begin{tabular}{@{}>{\raggedright\arraybackslash}p{23.14mm}>{\raggedright\arraybackslash}p{21.04mm}>{\raggedright\arraybackslash}p{22.72mm}>{\raggedright\arraybackslash}p{22.72mm}>{\raggedright\arraybackslash}p{23.14mm}@{}}
\toprule
\textbf{Infinitive¹} & \textbf{Meaning} & \textbf{Preterite sg.} & \textbf{Preterite pl.} & \textbf{Participle¹} \\
\midrule
faran & go; travel & fōr & fōron & faren \\
slēan & strike; smite & slōg & slōgon & slægen \\
standan & stand & stōd & stōdon & standen \\
hębban & raise; lift & hōf & hōfon & hafen \\
scieppan & create; shape & scēop & scēopon & sceapen \\
gescieppan & create & gescēop & gescēopon & gesceapen \\
swęrian & swear & swōr & swōron & sworen \\
\bottomrule
\end{tabular}
\par\endgroup

% mdtable:T08:end
% mdparagraph:P029:begin
¹ Infinitive and participle entries are supplied citations. The middle columns identify the finite preterite stems; the complete person assignments follow.
% mdparagraph:P029:end

% teaching: NEW-U05-PAR-faran
% teaching: NEW-U05-PAR-slean
% teaching: NEW-U05-PAR-standan
% teaching: NEW-U05-PAR-hebban
% teaching: NEW-U05-PAR-scieppan
% teaching: NEW-U05-PAR-gescieppan
% teaching: NEW-U05-PAR-swerian
\reserve{12\baselineskip}
% mdparagraph:P030:begin
\textbf{Present indicative.}
% mdparagraph:P030:end

% table: U05-PRESENT-VI
% mdtable:T09:begin
% course-table-style: 2026-10
% course-panel:T09-P01:begin
\par\begingroup
\fontsize{10.5}{13.2}\selectfont
\setlength{\LTleft}{0pt}\setlength{\LTright}{\fill}
\noindent
\begin{tabular}{@{}>{\raggedright\arraybackslash}p{24.232mm}>{\raggedright\arraybackslash}p{22.13mm}>{\raggedright\arraybackslash}p{22.13mm}>{\raggedright\arraybackslash}p{22.13mm}>{\raggedright\arraybackslash}p{22.13mm}@{}}
\toprule
\textbf{Person and number} & \textbf{faran} & \textbf{slēan} & \textbf{standan} & \textbf{hębban} \\
\midrule
1 sg. & fare & slēa & stande & hębbe \\
2 sg. & færst & sliehst & stęntst & hęfst \\
3 sg. & færð & sliehð & stęnt & hęfð \\
1–3 pl. & farað & slēað & standað & hębbað \\
\bottomrule
\end{tabular}
\par\endgroup
% course-panel:T09-P01:end
\par\addvspace{5pt}
% course-panel:T09-P02:begin
\par\begingroup
\fontsize{10.5}{13.2}\selectfont
\setlength{\LTleft}{0pt}\setlength{\LTright}{\fill}
\noindent
\begin{tabular}{@{}>{\raggedright\arraybackslash}p{30.902mm}>{\raggedright\arraybackslash}p{28.221mm}>{\raggedright\arraybackslash}p{28.221mm}>{\raggedright\arraybackslash}p{28.221mm}@{}}
\toprule
\textbf{Person and number} & \textbf{scieppan} & \textbf{gescieppan} & \textbf{swęrian} \\
\midrule
1 sg. & scieppe & gescieppe & swęrie \\
2 sg. & sciepst & gesciepst & swęrest \\
3 sg. & sciepð & gesciepð & swęreð \\
1–3 pl. & scieppað & gescieppað & swęriað \\
\bottomrule
\end{tabular}
\par\endgroup
% course-panel:T09-P02:end
% mdtable:T09:end
\reserve{12\baselineskip}
% mdparagraph:P031:begin
\textbf{Preterite indicative.}
% mdparagraph:P031:end

% table: U05-PRETERITE-VI
% mdtable:T10:begin
% course-table-style: 2026-10
% course-panel:T10-P01:begin
\par\begingroup
\fontsize{10.5}{13.2}\selectfont
\setlength{\LTleft}{0pt}\setlength{\LTright}{\fill}
\noindent
\begin{tabular}{@{}>{\raggedright\arraybackslash}p{24.232mm}>{\raggedright\arraybackslash}p{22.13mm}>{\raggedright\arraybackslash}p{22.13mm}>{\raggedright\arraybackslash}p{22.13mm}>{\raggedright\arraybackslash}p{22.13mm}@{}}
\toprule
\textbf{Person and number} & \textbf{faran} & \textbf{slēan} & \textbf{standan} & \textbf{hębban} \\
\midrule
1 sg. & fōr & slōg & stōd & hōf \\
2 sg. & fōre & slōge & stōde & hōfe \\
3 sg. & fōr & slōg & stōd & hōf \\
1–3 pl. & fōron & slōgon & stōdon & hōfon \\
\bottomrule
\end{tabular}
\par\endgroup
% course-panel:T10-P01:end
\par\addvspace{5pt}
% course-panel:T10-P02:begin
\par\begingroup
\fontsize{10.5}{13.2}\selectfont
\setlength{\LTleft}{0pt}\setlength{\LTright}{\fill}
\noindent
\begin{tabular}{@{}>{\raggedright\arraybackslash}p{30.902mm}>{\raggedright\arraybackslash}p{28.221mm}>{\raggedright\arraybackslash}p{28.221mm}>{\raggedright\arraybackslash}p{28.221mm}@{}}
\toprule
\textbf{Person and number} & \textbf{scieppan} & \textbf{gescieppan} & \textbf{swęrian} \\
\midrule
1 sg. & scēop & gescēop & swōr \\
2 sg. & scēope & gescēope & swōre \\
3 sg. & scēop & gescēop & swōr \\
1–3 pl. & scēopon & gescēopon & swōron \\
\bottomrule
\end{tabular}
\par\endgroup
% course-panel:T10-P02:end
% mdtable:T10:end
% teaching: NEW-U05-GRAM-CONTRACT-VI
% mdparagraph:P032:begin
\textbf{Slēan}\BookPoint{bk-ch05-learner-02721c417980} has a contracted present: slēa, sliehst, sliehð, slēað. The ie in sliehst and sliehð is short. The selected past has g throughout: slōg, slōge, slōgon. The supplied participle is slægen.
% mdparagraph:P032:end

% teaching: NEW-U05-GRAM-STANDAN-NASAL
% mdparagraph:P033:begin
\textbf{Standan} keeps n in the present but loses it in the past stōd/\allowbreak{}stōdon. Its changed present singulars are stęntst and stęnt. This is the ordinary verb of standing; the earlier reading’s opposing-winds construction requires its own supplied explanation.
% mdparagraph:P033:end

% teaching: NEW-U05-GRAM-J-PRESENT-VI
% mdparagraph:P034:begin
\textbf{Hębban}\BookPoint{bk-ch05-learner-69fa4f57c0e3} and \textbf{scieppan} have doubled present consonants in hębbe/\allowbreak{}hębbað and scieppe/\allowbreak{}scieppað. The selected second and third singulars have f in hęfst/\allowbreak{}hęfð and single p in sciepst/\allowbreak{}sciepð. Their strong pasts are hōf/\allowbreak{}hōfon and scēop/\allowbreak{}scēopon, with supplied participles hafen and sceapen. Gescieppan follows its separately given prefixed table.
% mdparagraph:P034:end

% mdparagraph:P035:begin
\textbf{Swęrian}\BookPoint{bk-ch05-learner-aece1b132820} has swęrie, swęrest, swęreð, swęriað, retaining full endings in the second and third singular. Its past is nevertheless strong: swōr, swōre, swōron. An infinitive in -ian does not by itself identify a weak verb. Keep the full -est and -eð in this selected present.
% mdparagraph:P035:end

% mdparagraph:P036:begin
\textbf{Alternatives.} Cook gives scōp/\allowbreak{}scōpon alongside scēop/\allowbreak{}scēopon. The corresponding prefix-preserving alternatives are gescōp/\allowbreak{}gescōpon. Cook also gives the participle slegen beside slægen; Wright records slōh beside singular slōg and the participle slagen. The selected cells remain those in the tables.
% mdparagraph:P036:end

\reserve{7\baselineskip}
% mdparagraph:P037:begin
\section*{5. Class VII}\BookHeading{nav-ch05-learner-0005}{subsection}{5. Class VII}{1}\BookLabel{bk-ch05-learner-55e28af3ed17}
% mdparagraph:P037:end

% mdparagraph:P038:begin
Cook groups these verbs as \textbf{reduplicating verbs}; they are commonly called Class VII. For present purposes, learn two past-vowel groups, ēo and ē. Within each verb the selected preterite stem is the same in singular and plural. The present and participle usually share their vowel, but the complete entries preserve exceptions such as wēpan/\allowbreak{}wōpen and fōn/\allowbreak{}fangen.
% mdparagraph:P038:end

\reserve{7\baselineskip}
% mdparagraph:P039:begin
\subsection*{Preterites with ēo}\BookHeading{nav-ch05-learner-0006}{subsubsection}{Preterites with ēo}{0}
% mdparagraph:P039:end

% teaching: NEW-U05-GRAM-CLASS-VII-EO
% teaching: NEW-U05-VOC-feallan
% teaching: NEW-U05-CIT-feallan
% teaching: NEW-U05-VOC-healdan
% teaching: NEW-U05-CIT-healdan
% teaching: NEW-U05-VOC-weaxan
% teaching: NEW-U05-CIT-weaxan
% teaching: NEW-U05-VOC-cnawan
% teaching: NEW-U05-CIT-cnawan
% teaching: NEW-U05-VOC-heawan
% teaching: NEW-U05-CIT-heawan
% teaching: NEW-U05-VOC-flowan
% teaching: NEW-U05-CIT-flowan
% teaching: NEW-U05-VOC-wepan
% teaching: NEW-U05-CIT-wepan
% table: U05-PARTS-VII-EO
% mdtable:T11:begin
\par\begingroup
\fontsize{10.5}{13.2}\selectfont
\setlength{\LTleft}{0pt}\setlength{\LTright}{\fill}
\noindent
\begin{tabular}{@{}>{\raggedright\arraybackslash}p{23.14mm}>{\raggedright\arraybackslash}p{21.04mm}>{\raggedright\arraybackslash}p{22.72mm}>{\raggedright\arraybackslash}p{22.72mm}>{\raggedright\arraybackslash}p{23.14mm}@{}}
\toprule
\textbf{Infinitive¹} & \textbf{Meaning} & \textbf{Preterite sg.} & \textbf{Preterite pl.} & \textbf{Participle¹} \\
\midrule
feallan & fall & fēoll & fēollon & feallen \\
healdan & hold & hēold & hēoldon & healden \\
weaxan & grow & wēox & wēoxon & weaxen \\
cnāwan & know & cnēow & cnēowon & cnāwen \\
hēawan & hew & hēow & hēowon & hēawen \\
flōwan & flow & flēow & flēowon & flōwen \\
wēpan & weep & wēop & wēopon & wōpen \\
\bottomrule
\end{tabular}
\par\endgroup

% mdtable:T11:end
% mdparagraph:P040:begin
¹ Infinitive and participle entries are supplied citations. The middle columns identify the finite preterite stems; the complete person assignments follow.
% mdparagraph:P040:end

% teaching: NEW-U05-PAR-feallan
% teaching: NEW-U05-PAR-healdan
% teaching: NEW-U05-PAR-weaxan
% teaching: NEW-U05-PAR-cnawan
% teaching: NEW-U05-PAR-heawan
% teaching: NEW-U05-PAR-flowan
% teaching: NEW-U05-PAR-wepan
\reserve{12\baselineskip}
% mdparagraph:P041:begin
\textbf{Present indicative.}
% mdparagraph:P041:end

% table: U05-PRESENT-VII-EO
% mdtable:T12:begin
% course-table-style: 2026-10
% course-panel:T12-P01:begin
\par\begingroup
\fontsize{10.5}{13.2}\selectfont
\setlength{\LTleft}{0pt}\setlength{\LTright}{\fill}
\noindent
\begin{tabular}{@{}>{\raggedright\arraybackslash}p{24.232mm}>{\raggedright\arraybackslash}p{22.13mm}>{\raggedright\arraybackslash}p{22.13mm}>{\raggedright\arraybackslash}p{22.13mm}>{\raggedright\arraybackslash}p{22.13mm}@{}}
\toprule
\textbf{Person and number} & \textbf{feallan} & \textbf{healdan} & \textbf{weaxan} & \textbf{cnāwan} \\
\midrule
1 sg. & fealle & healde & weaxe & cnāwe \\
2 sg. & fielst & hieltst & wiext & cnǣwst \\
3 sg. & fielð & hielt & wiext & cnǣwð \\
1–3 pl. & feallað & healdað & weaxað & cnāwað \\
\bottomrule
\end{tabular}
\par\endgroup
% course-panel:T12-P01:end
\par\addvspace{5pt}
% course-panel:T12-P02:begin
\par\begingroup
\fontsize{10.5}{13.2}\selectfont
\setlength{\LTleft}{0pt}\setlength{\LTright}{\fill}
\noindent
\begin{tabular}{@{}>{\raggedright\arraybackslash}p{30.902mm}>{\raggedright\arraybackslash}p{28.221mm}>{\raggedright\arraybackslash}p{28.221mm}>{\raggedright\arraybackslash}p{28.221mm}@{}}
\toprule
\textbf{Person and number} & \textbf{hēawan} & \textbf{flōwan} & \textbf{wēpan} \\
\midrule
1 sg. & hēawe & flōwe & wēpe \\
2 sg. & hīewst & flēwst & wēpst \\
3 sg. & hīewð & flēwð & wēpð \\
1–3 pl. & hēawað & flōwað & wēpað \\
\bottomrule
\end{tabular}
\par\endgroup
% course-panel:T12-P02:end
% mdtable:T12:end
\reserve{12\baselineskip}
% mdparagraph:P042:begin
\textbf{Preterite indicative.}
% mdparagraph:P042:end

% table: U05-PRETERITE-VII-EO
% mdtable:T13:begin
% course-table-style: 2026-10
% course-panel:T13-P01:begin
\par\begingroup
\fontsize{10.5}{13.2}\selectfont
\setlength{\LTleft}{0pt}\setlength{\LTright}{\fill}
\noindent
\begin{tabular}{@{}>{\raggedright\arraybackslash}p{24.232mm}>{\raggedright\arraybackslash}p{22.13mm}>{\raggedright\arraybackslash}p{22.13mm}>{\raggedright\arraybackslash}p{22.13mm}>{\raggedright\arraybackslash}p{22.13mm}@{}}
\toprule
\textbf{Person and number} & \textbf{feallan} & \textbf{healdan} & \textbf{weaxan} & \textbf{cnāwan} \\
\midrule
1 sg. & fēoll & hēold & wēox & cnēow \\
2 sg. & fēolle & hēolde & wēoxe & cnēowe \\
3 sg. & fēoll & hēold & wēox & cnēow \\
1–3 pl. & fēollon & hēoldon & wēoxon & cnēowon \\
\bottomrule
\end{tabular}
\par\endgroup
% course-panel:T13-P01:end
\par\addvspace{5pt}
% course-panel:T13-P02:begin
\par\begingroup
\fontsize{10.5}{13.2}\selectfont
\setlength{\LTleft}{0pt}\setlength{\LTright}{\fill}
\noindent
\begin{tabular}{@{}>{\raggedright\arraybackslash}p{30.902mm}>{\raggedright\arraybackslash}p{28.221mm}>{\raggedright\arraybackslash}p{28.221mm}>{\raggedright\arraybackslash}p{28.221mm}@{}}
\toprule
\textbf{Person and number} & \textbf{hēawan} & \textbf{flōwan} & \textbf{wēpan} \\
\midrule
1 sg. & hēow & flēow & wēop \\
2 sg. & hēowe & flēowe & wēope \\
3 sg. & hēow & flēow & wēop \\
1–3 pl. & hēowon & flēowon & wēopon \\
\bottomrule
\end{tabular}
\par\endgroup
% course-panel:T13-P02:end
% mdtable:T13:end
% teaching: NEW-U05-GRAM-VII-PRESENT-CHANGES
% mdparagraph:P043:begin
In the present second and third singular, \textbf{feallan, healdan and weaxan} change short ea to short ie. Healdan also has the dental results hieltst and hielt. Stem-final x combines with either singular ending as xt: \textbf{wiext is both second and third singular} of weaxan. Feallan reduces ll before its endings, giving fielst and fielð.
% mdparagraph:P043:end

% mdparagraph:P044:begin
Cnāwan changes ā to ǣ in cnǣwst/\allowbreak{}cnǣwð; hēawan changes long ēa to long īe in hīewst/\allowbreak{}hīewð; flōwan changes ō to ē in flēwst/\allowbreak{}flēwð. Wēpan keeps ē in wēpst/\allowbreak{}wēpð. In the preterite, cnēow/\allowbreak{}cnēowon and hēow/\allowbreak{}hēowon retain their printed consonants and quantities; do not transfer a consonant from another verb’s past.
% mdparagraph:P044:end

\reserve{7\baselineskip}
% mdparagraph:P045:begin
\subsection*{Preterites with ē}\BookHeading{nav-ch05-learner-0007}{subsubsection}{Preterites with ē}{0}
% mdparagraph:P045:end

% teaching: NEW-U05-GRAM-CLASS-VII-E
% teaching: NEW-U05-VOC-laetan
% teaching: NEW-U05-CIT-laetan
% teaching: NEW-U05-VOC-hatan
% teaching: NEW-U05-CIT-hatan
% teaching: NEW-U05-VOC-fon
% teaching: NEW-U05-CIT-fon
% table: U05-PARTS-VII-E
% mdtable:T14:begin
\par\begingroup
\fontsize{10.5}{13.2}\selectfont
\setlength{\LTleft}{0pt}\setlength{\LTright}{\fill}
\noindent
\begin{tabular}{@{}>{\raggedright\arraybackslash}p{23.14mm}>{\raggedright\arraybackslash}p{21.04mm}>{\raggedright\arraybackslash}p{22.72mm}>{\raggedright\arraybackslash}p{22.72mm}>{\raggedright\arraybackslash}p{23.14mm}@{}}
\toprule
\textbf{Infinitive¹} & \textbf{Meaning} & \textbf{Preterite sg.} & \textbf{Preterite pl.} & \textbf{Participle¹} \\
\midrule
lǣtan & let & lēt & lēton & lǣten \\
hātan & call; name & hēt & hēton & hāten \\
fōn & seize; catch & fēng & fēngon & fangen \\
\bottomrule
\end{tabular}
\par\endgroup

% mdtable:T14:end
% mdparagraph:P046:begin
¹ Infinitive and participle entries are supplied citations. The middle columns identify the finite preterite stems; the complete person assignments follow.
% mdparagraph:P046:end

% teaching: NEW-U05-PAR-laetan
% teaching: NEW-U05-PAR-hatan
% teaching: NEW-U05-PAR-fon
\reserve{8\baselineskip}
% mdparagraph:P047:begin
\textbf{Present indicative.}
% mdparagraph:P047:end

% table: U05-PRESENT-VII-E
% mdtable:T15:begin
% course-table-style: 2026-10
% course-panel:T15-P01:begin
\par\begingroup
\fontsize{10.5}{13.2}\selectfont
\setlength{\LTleft}{0pt}\setlength{\LTright}{\fill}
\noindent
\begin{tabular}{@{}>{\raggedright\arraybackslash}p{30.902mm}>{\raggedright\arraybackslash}p{28.221mm}>{\raggedright\arraybackslash}p{28.221mm}>{\raggedright\arraybackslash}p{28.221mm}@{}}
\toprule
\textbf{Person and number} & \textbf{lǣtan} & \textbf{hātan} & \textbf{fōn} \\
\midrule
1 sg. & lǣte & hāte & fō \\
2 sg. & lǣtst & hǣtst & fēhst \\
3 sg. & lǣtt & hǣtt & fēhð \\
1–3 pl. & lǣtað & hātað & fōð \\
\bottomrule
\end{tabular}
\par\endgroup
% course-panel:T15-P01:end
% mdtable:T15:end
\reserve{8\baselineskip}
% mdparagraph:P048:begin
\textbf{Preterite indicative.}
% mdparagraph:P048:end

% table: U05-PRETERITE-VII-E
% mdtable:T16:begin
% course-table-style: 2026-10
% course-panel:T16-P01:begin
\par\begingroup
\fontsize{10.5}{13.2}\selectfont
\setlength{\LTleft}{0pt}\setlength{\LTright}{\fill}
\noindent
\begin{tabular}{@{}>{\raggedright\arraybackslash}p{30.902mm}>{\raggedright\arraybackslash}p{28.221mm}>{\raggedright\arraybackslash}p{28.221mm}>{\raggedright\arraybackslash}p{28.221mm}@{}}
\toprule
\textbf{Person and number} & \textbf{lǣtan} & \textbf{hātan} & \textbf{fōn} \\
\midrule
1 sg. & lēt & hēt & fēng \\
2 sg. & lēte & hēte & fēnge \\
3 sg. & lēt & hēt & fēng \\
1–3 pl. & lēton & hēton & fēngon \\
\bottomrule
\end{tabular}
\par\endgroup
% course-panel:T16-P01:end
% mdtable:T16:end
% mdparagraph:P049:begin
\textbf{Lǣtan}\BookPoint{bk-ch05-learner-748db6b8492f} has lǣtst and lǣtt in the present; \textbf{hātan} changes ā to ǣ in hǣtst and hǣtt. Their preterites are lēt/\allowbreak{}lēton and hēt/\allowbreak{}hēton. Third-singular lǣt and hǣt are documented alternatives. Wright also records the older preterite hēht; it is not an additional recall requirement here.
% mdparagraph:P049:end

% teaching: NEW-U05-GRAM-CONTRACT-VII
% mdparagraph:P050:begin
\textbf{Fōn}\BookPoint{bk-ch05-learner-be6ba8199b8f} has the contracted present fō, fēhst, fēhð, fōð and the past fēng, fēnge, fēngon. The nasal consonant belongs to the past stem. The participle fangen remains a supplied citation.
% mdparagraph:P050:end

% mdparagraph:P051:begin
The hātan table is active indicative. The special passive form hātte, commands with infinitives, and the reading’s construction with an object and a name are separate matters. The naming phrases in §§11–12 remain supplied in full.
% mdparagraph:P051:end

\reserve{7\baselineskip}
% mdparagraph:P052:begin
\section*{6. Four additions to Classes I–III}\BookHeading{nav-ch05-learner-0008}{subsection}{6. Four additions to Classes I–III}{1}
% mdparagraph:P052:end

% teaching: NEW-U05-GRAM-BRIDGES-I-III
% mdparagraph:P053:begin
The following verbs use classes already learned, but their complete lexical systems are new here. Earlier supplied occurrences do not replace learning these tables. Geswīcan and scīnan follow Class I, flēogan Class II, and swimman the nasal pattern of Class III.
% mdparagraph:P053:end

% teaching: NEW-U05-VOC-geswican
% teaching: NEW-U05-CIT-geswican
% teaching: NEW-U05-VOC-scinan
% teaching: NEW-U05-CIT-scinan
% teaching: NEW-U05-VOC-fleogan
% teaching: NEW-U05-CIT-fleogan
% teaching: NEW-U05-VOC-swimman
% teaching: NEW-U05-CIT-swimman
% table: U05-PARTS-I-III
% mdtable:T17:begin
\par\begingroup
\fontsize{10.5}{13.2}\selectfont
\setlength{\LTleft}{0pt}\setlength{\LTright}{\fill}
\noindent
\begin{tabular}{@{}>{\raggedright\arraybackslash}p{23.14mm}>{\raggedright\arraybackslash}p{21.04mm}>{\raggedright\arraybackslash}p{22.72mm}>{\raggedright\arraybackslash}p{22.72mm}>{\raggedright\arraybackslash}p{23.14mm}@{}}
\toprule
\textbf{Infinitive¹} & \textbf{Meaning} & \textbf{Preterite sg.} & \textbf{Preterite pl.} & \textbf{Participle¹} \\
\midrule
geswīcan & cease & geswāc & geswicon & geswicen \\
scīnan & shine & scān & scinon & scinen \\
flēogan & fly & flēag & flugon & flogen \\
swimman & swim & swamm & swummon & swummen \\
\bottomrule
\end{tabular}
\par\endgroup

% mdtable:T17:end
% mdparagraph:P054:begin
¹ Infinitive and participle entries are supplied citations. The middle columns identify the finite preterite stems; the complete person assignments follow.
% mdparagraph:P054:end

% teaching: NEW-U05-PAR-geswican
% teaching: NEW-U05-PAR-scinan
% teaching: NEW-U05-PAR-fleogan
% teaching: NEW-U05-PAR-swimman
\reserve{9\baselineskip}
% mdparagraph:P055:begin
\textbf{Present indicative.}
% mdparagraph:P055:end

% table: U05-PRESENT-I-III
% mdtable:T18:begin
% course-table-style: 2026-10
% course-panel:T18-P01:begin
\par\begingroup
\fontsize{10.5}{13.2}\selectfont
\setlength{\LTleft}{0pt}\setlength{\LTright}{\fill}
\noindent
\begin{tabular}{@{}>{\raggedright\arraybackslash}p{24.232mm}>{\raggedright\arraybackslash}p{22.13mm}>{\raggedright\arraybackslash}p{22.13mm}>{\raggedright\arraybackslash}p{22.13mm}>{\raggedright\arraybackslash}p{22.13mm}@{}}
\toprule
\textbf{Person and number} & \textbf{geswīcan} & \textbf{scīnan} & \textbf{flēogan} & \textbf{swimman} \\
\midrule
1 sg. & geswīce & scīne & flēoge & swimme \\
2 sg. & geswīcst & scīnst & flīehst & swimst \\
3 sg. & geswīcð & scīnð & flīehð & swimð \\
1–3 pl. & geswīcað & scīnað & flēogað & swimmað \\
\bottomrule
\end{tabular}
\par\endgroup
% course-panel:T18-P01:end
% mdtable:T18:end
\reserve{9\baselineskip}
% mdparagraph:P056:begin
\textbf{Preterite indicative.}
% mdparagraph:P056:end

% table: U05-PRETERITE-I-III
% mdtable:T19:begin
% course-table-style: 2026-10
% course-panel:T19-P01:begin
\par\begingroup
\fontsize{10.5}{13.2}\selectfont
\setlength{\LTleft}{0pt}\setlength{\LTright}{\fill}
\noindent
\begin{tabular}{@{}>{\raggedright\arraybackslash}p{24.232mm}>{\raggedright\arraybackslash}p{22.13mm}>{\raggedright\arraybackslash}p{22.13mm}>{\raggedright\arraybackslash}p{22.13mm}>{\raggedright\arraybackslash}p{22.13mm}@{}}
\toprule
\textbf{Person and number} & \textbf{geswīcan} & \textbf{scīnan} & \textbf{flēogan} & \textbf{swimman} \\
\midrule
1 sg. & geswāc & scān & flēag & swamm \\
2 sg. & geswice & scine & fluge & swumme \\
3 sg. & geswāc & scān & flēag & swamm \\
1–3 pl. & geswicon & scinon & flugon & swummon \\
\bottomrule
\end{tabular}
\par\endgroup
% course-panel:T19-P01:end
% mdtable:T19:end
% mdparagraph:P057:begin
Geswīcan and scīnan have long ī in the present, long ā in the first/\allowbreak{}third singular past, and short i in the plural-based past. Flēogan has ēo in the first singular and plural present, long īe in flīehst/\allowbreak{}flīehð, and u in fluge/\allowbreak{}flugon. Swimman has short a in swamm but u in swumme/\allowbreak{}swummon; mm reduces before its present consonantal endings.
% mdparagraph:P057:end

% mdparagraph:P058:begin
\textbf{Alternatives.} Cook records flīegst and flīegð beside selected flīehst and flīehð. The dictionary also attests swimman past alternatives swom and swommon. The tables select the a/\allowbreak{}u series for swimman and the h spellings for flēogan.
% mdparagraph:P058:end

\reserve{7\baselineskip}
% mdparagraph:P059:begin
\section*{7. Ordinary uses}\BookHeading{nav-ch05-learner-0009}{subsection}{7. Ordinary uses}{1}\BookLabel{bk-ch05-learner-95f29519f213}
% mdparagraph:P059:end

% mdparagraph:P060:begin
The glosses in the principal-part tables identify the verbs. The constructions below supply the ordinary uses needed in sentences. For other verbs, the practice asks for labelled person, number and tense forms; a gloss alone does not specify which objects or clauses may follow.
% mdparagraph:P060:end

% teaching: NEW-U05-GRAM-CARRY-TAKE-OBJECTS
% mdparagraph:P061:begin
\textbf{Beran}, “carry,” and \textbf{niman}, “take,” take an accusative person or thing. The subject controls the finite verb. The following eight examples are new teaching material.
% mdparagraph:P061:end

% example: NEW-U05-E01
% mdparagraph:P062:begin
Ðū birst ðone męte.
% mdparagraph:P062:end

% mdparagraph:P063:begin
You (singular) carry the food.
% mdparagraph:P063:end

% example: NEW-U05-E02
% mdparagraph:P064:begin
Hē nimð ðone fisc.
% mdparagraph:P064:end

% mdparagraph:P065:begin
He takes the fish.
% mdparagraph:P065:end

% teaching: NEW-U05-GRAM-GIVING-RECIPIENT
% mdparagraph:P066:begin
With \textbf{giefan}, the thing given is accusative and the recipient is dative. English may express the recipient with “to”; Old English expresses the relation here by the case of the recipient itself.
% mdparagraph:P066:end

% example: NEW-U05-E03
% mdparagraph:P067:begin
Ic giefe ðǣm męnn ðone męte.
% mdparagraph:P067:end

% mdparagraph:P068:begin
I give the man the food.
% mdparagraph:P068:end

% teaching: NEW-U05-GRAM-INTRANSITIVE-POSTURE
% mdparagraph:P069:begin
\textbf{Sittan}, “sit,” \textbf{licgan}, “lie; be at rest,” and \textbf{standan}, “stand,” describe the subject’s posture. A locative \textbf{on + dative} phrase can state where that subject is. No object is needed for these ordinary posture uses.
% mdparagraph:P069:end

% example: NEW-U05-E04
% mdparagraph:P070:begin
Se mann ligð on ðǣm hofe.
% mdparagraph:P070:end

% mdparagraph:P071:begin
The man lies in the dwelling.
% mdparagraph:P071:end

% example: NEW-U05-E05
% mdparagraph:P072:begin
Se mann stęnt on ðǣm strande.
% mdparagraph:P072:end

% mdparagraph:P073:begin
The man stands on the shore.
% mdparagraph:P073:end

% teaching: NEW-U05-GRAM-INTRANSITIVE-MOTION
% teaching: NEW-U05-GRAM-BRIDGE-INTRANSITIVE
% mdparagraph:P074:begin
\textbf{Feallan}, “fall,” \textbf{weaxan}, “grow,” \textbf{flēogan}, “fly,” \textbf{scīnan}, “shine,” and \textbf{swimman}, “swim,” have ordinary intransitive uses. The growing tree, flying bird, shining star or swimming fish is the subject. A locative on-phrase may state where the swimming occurs.
% mdparagraph:P074:end

% example: NEW-U05-E06
% mdparagraph:P075:begin
Ðæt trēo fielð.
% mdparagraph:P075:end

% mdparagraph:P076:begin
The tree falls.
% mdparagraph:P076:end

% teaching: NEW-U05-GRAM-EAT
% teaching: NEW-U05-GRAM-SEE-OBJECT
% teaching: NEW-U05-GRAM-STRIKE-HOLD-OBJECTS
% teaching: NEW-U05-GRAM-U05-CREATE-OBJECT
% mdparagraph:P077:begin
\textbf{Etan}, “eat,” takes an accusative food object; it also occurs with the food understood. \textbf{Sēon} and \textbf{gesēon}, “see,” take an accusative person or thing in the use practised here. \textbf{Slēan}, “strike; smite,” and \textbf{healdan}, “hold,” likewise take an accusative object. “Strike” does not by itself say that the object is killed. \textbf{Scieppan} and \textbf{gescieppan} take the created being or thing as accusative object in the Creation sentences.
% mdparagraph:P077:end

% teaching: NEW-U05-GRAM-CESSATION-GENITIVE
% mdparagraph:P078:begin
\textbf{Geswīcan},\BookPoint{bk-ch05-learner-e8f7cdddfe3f} “cease,” takes the \textbf{genitive} of the work or action from which one ceases. Keep this government when English uses “cease from.” An agreeing possessive takes the case, number and gender of the work noun; the personal genitives his and hiera remain the possessor forms already learned.
% mdparagraph:P078:end

% example: NEW-U05-E07
% mdparagraph:P079:begin
Wē geswīcað ūres weorces.
% mdparagraph:P079:end

% mdparagraph:P080:begin
We cease from our work.
% mdparagraph:P080:end

% teaching: NEW-U05-GRAM-DIRECT-QUOTATION
% mdparagraph:P081:begin
\textbf{Cweðan},\BookPoint{bk-ch05-learner-a3edb3da0846} “say,” can introduce direct quotation. The words inside the quotation form their own clause; a first-person pronoun there denotes the quoted speaker. Content clauses introduced by “that” remain separately supplied in the reading.
% mdparagraph:P081:end

% example: NEW-U05-E08
% mdparagraph:P082:begin
Se mann cwið: “Ic sēo ðone fugol.”
% mdparagraph:P082:end

% mdparagraph:P083:begin
The man says, “I see the bird.”
% mdparagraph:P083:end

% teaching: NEW-U05-GRAM-SIMPLE-NEGATION
% mdparagraph:P084:begin
For these ordinary indicative clauses, put \textbf{ne before the finite verb} to negate the statement. A past negative still uses the appropriate strong preterite; English “did not” does not add a separate Old English auxiliary here. Subjects in the dual take plural verb forms, while dual possessives still identify two possessors and agree with the possessed noun.
% mdparagraph:P084:end

% mdparagraph:P085:begin
No new productive noun pattern is needed. Continue the complete noun, adjective, article, pronoun and possessive paradigms of Units 1–4. In particular, the learned masculine heofon supplies plural “heavens”; the reading’s weak feminine heofone receives separate local help.
% mdparagraph:P085:end

\reserve{7\baselineskip}
% mdparagraph:P086:begin
\section*{8. Form and clause practice}\BookHeading{nav-ch05-learner-0010}{subsection}{8. Form and clause practice}{1}
% mdparagraph:P086:end

% mdparagraph:P087:begin
P, F and C are new teaching exercises. \textbf{Selected} means the form in this unit’s main tables. Consult the printed principal parts where the prompt supplies them; the nonfinite entries are never independent recall targets. In translation tasks, retain each stated person, number, tense and case relation.
% mdparagraph:P087:end

\reserve{16\baselineskip}
% mdparagraph:P088:begin
\subsection*{Principal parts — P1–P4}\BookHeading{nav-ch05-learner-0011}{subsubsection}{Principal parts — P1–P4}{0}
% mdparagraph:P088:end

\reserve{5\baselineskip}
\par\begin{minipage}{\linewidth}\raggedright
% mdparagraph:P089:begin
\textbf{P1.}
% mdparagraph:P089:end

% item: NEW-U05-P01
% mdparagraph:P090:begin
The printed entry gives beran, bær, bǣron, boren. Label the four principal-part positions. Which entries are finite, and which are nonfinite citations?
% mdparagraph:P090:end

\end{minipage}\par\addvspace{4pt}

\par\begin{minipage}{\linewidth}\raggedright
% mdparagraph:P091:begin
\textbf{P2.}
% mdparagraph:P091:end

% item: NEW-U05-P02
% mdparagraph:P092:begin
Compare geaf and sealde, both meaning gave. Identify the strong form and the weak form; explain the difference in preterite formation.
% mdparagraph:P092:end

\end{minipage}\par\addvspace{4pt}

\par\begin{minipage}{\linewidth}\raggedright
% mdparagraph:P093:begin
\textbf{P3.}
% mdparagraph:P093:end

% item: NEW-U05-P03
% mdparagraph:P094:begin
For beran, give the preterite indicative forms for I, you (singular), he, and we. Identify the forms using the preterite plural stem.
% mdparagraph:P094:end

\end{minipage}\par\addvspace{4pt}

\par\begin{minipage}{\linewidth}\raggedright
% mdparagraph:P095:begin
\textbf{P4.}
% mdparagraph:P095:end

% item: NEW-U05-P04
% mdparagraph:P096:begin
The printed entry gives gescieppan, gescēop, gescēopon, gesceapen. Using it, give you (singular) created and we created. Is gesceapen itself a finite preterite?
% mdparagraph:P096:end

\end{minipage}\par\addvspace{4pt}

\reserve{7\baselineskip}
% mdparagraph:P097:begin
\subsection*{Finite forms — F1–F24}\BookHeading{nav-ch05-learner-0012}{subsubsection}{Finite forms — F1–F24}{0}
% mdparagraph:P097:end

\reserve{5\baselineskip}
\par\begin{minipage}{\linewidth}\raggedright
% mdparagraph:P098:begin
\textbf{F1.}
% mdparagraph:P098:end

% item: NEW-U05-F01
% mdparagraph:P099:begin
Give the selected indicative forms of beran: present 2sg; preterite 1sg and 2sg.
% mdparagraph:P099:end

\end{minipage}\par\addvspace{4pt}

\par\begin{minipage}{\linewidth}\raggedright
% mdparagraph:P100:begin
\textbf{F2.}
% mdparagraph:P100:end

% item: NEW-U05-F02
% mdparagraph:P101:begin
Give the selected indicative forms of niman: present 3sg; preterite 3sg and 1pl.
% mdparagraph:P101:end

\end{minipage}\par\addvspace{4pt}

\par\begin{minipage}{\linewidth}\raggedright
% mdparagraph:P102:begin
\textbf{F3.}
% mdparagraph:P102:end

% item: NEW-U05-F03
% mdparagraph:P103:begin
Give the selected indicative forms of cuman: present 1sg, 2sg and 3sg; preterite 3pl.
% mdparagraph:P103:end

\end{minipage}\par\addvspace{4pt}

\par\begin{minipage}{\linewidth}\raggedright
% mdparagraph:P104:begin
\textbf{F4.}
% mdparagraph:P104:end

% item: NEW-U05-F04
% mdparagraph:P105:begin
Give the selected indicative forms of sprecan: present 2sg and 3sg; preterite 2sg.
% mdparagraph:P105:end

\end{minipage}\par\addvspace{4pt}

\par\begin{minipage}{\linewidth}\raggedright
% mdparagraph:P106:begin
\textbf{F5.}
% mdparagraph:P106:end

% item: NEW-U05-F05
% mdparagraph:P107:begin
Give the selected indicative forms of cweðan: present 2sg and 3sg; preterite 3pl.
% mdparagraph:P107:end

\end{minipage}\par\addvspace{4pt}

\par\begin{minipage}{\linewidth}\raggedright
% mdparagraph:P108:begin
\textbf{F6.}
% mdparagraph:P108:end

% item: NEW-U05-F06
% mdparagraph:P109:begin
Give the selected indicative forms of giefan: present 2sg and 3sg; preterite 1sg and 1pl.
% mdparagraph:P109:end

\end{minipage}\par\addvspace{4pt}

\par\begin{minipage}{\linewidth}\raggedright
% mdparagraph:P110:begin
\textbf{F7.}
% mdparagraph:P110:end

% item: NEW-U05-F07
% mdparagraph:P111:begin
Give the selected indicative forms of biddan: present 1sg and 3sg; preterite 2sg.
% mdparagraph:P111:end

\end{minipage}\par\addvspace{4pt}

\par\begin{minipage}{\linewidth}\raggedright
% mdparagraph:P112:begin
\textbf{F8.}
% mdparagraph:P112:end

% item: NEW-U05-F08
% mdparagraph:P113:begin
Give the selected indicative forms of licgan: present 2sg and 3sg; preterite 1pl.
% mdparagraph:P113:end

\end{minipage}\par\addvspace{4pt}

\par\begin{minipage}{\linewidth}\raggedright
% mdparagraph:P114:begin
\textbf{F9.}
% mdparagraph:P114:end

% item: NEW-U05-F09
% mdparagraph:P115:begin
Give the selected indicative forms of sittan: present 3sg; preterite 1sg and 2sg.
% mdparagraph:P115:end

\end{minipage}\par\addvspace{4pt}

\par\begin{minipage}{\linewidth}\raggedright
% mdparagraph:P116:begin
\textbf{F10.}
% mdparagraph:P116:end

% item: NEW-U05-F10
% mdparagraph:P117:begin
Give the selected indicative forms of etan: present 2sg and 3sg; preterite 1sg and 2sg.
% mdparagraph:P117:end

\end{minipage}\par\addvspace{4pt}

\par\begin{minipage}{\linewidth}\raggedright
% mdparagraph:P118:begin
\textbf{F11.}
% mdparagraph:P118:end

% item: NEW-U05-F11
% mdparagraph:P119:begin
Give the selected indicative forms of sēon: present 1sg, 2sg and 3pl.
% mdparagraph:P119:end

\end{minipage}\par\addvspace{4pt}

\par\begin{minipage}{\linewidth}\raggedright
% mdparagraph:P120:begin
\textbf{F12.}
% mdparagraph:P120:end

% item: NEW-U05-F12
% mdparagraph:P121:begin
Give the selected indicative forms of gesēon: present 3sg; preterite 2sg and 2pl.
% mdparagraph:P121:end

\end{minipage}\par\addvspace{4pt}

\par\begin{minipage}{\linewidth}\raggedright
% mdparagraph:P122:begin
\textbf{F13.}
% mdparagraph:P122:end

% item: NEW-U05-F13
% mdparagraph:P123:begin
Give the selected indicative forms of faran: present 2sg and 3sg; preterite 1pl.
% mdparagraph:P123:end

\end{minipage}\par\addvspace{4pt}

\par\begin{minipage}{\linewidth}\raggedright
% mdparagraph:P124:begin
\textbf{F14.}
% mdparagraph:P124:end

% item: NEW-U05-F14
% mdparagraph:P125:begin
Give the selected indicative forms of slēan: present 1sg, 2sg and 2pl; preterite 2sg.
% mdparagraph:P125:end

\end{minipage}\par\addvspace{4pt}

\par\begin{minipage}{\linewidth}\raggedright
% mdparagraph:P126:begin
\textbf{F15.}
% mdparagraph:P126:end

% item: NEW-U05-F15
% mdparagraph:P127:begin
Give the selected indicative forms of standan: present 2sg and 3sg; preterite 3sg.
% mdparagraph:P127:end

\end{minipage}\par\addvspace{4pt}

\par\begin{minipage}{\linewidth}\raggedright
% mdparagraph:P128:begin
\textbf{F16.}
% mdparagraph:P128:end

% item: NEW-U05-F16
% mdparagraph:P129:begin
Give the selected indicative forms of hębban: present 1sg and 3sg; preterite 1pl.
% mdparagraph:P129:end

\end{minipage}\par\addvspace{4pt}

\par\begin{minipage}{\linewidth}\raggedright
% mdparagraph:P130:begin
\textbf{F17.}
% mdparagraph:P130:end

% item: NEW-U05-F17
% mdparagraph:P131:begin
Give the selected indicative forms of scieppan and gescieppan: scieppan: present 1sg and 2sg, preterite 3sg; gescieppan: preterite 3pl.
% mdparagraph:P131:end

\end{minipage}\par\addvspace{4pt}

\par\begin{minipage}{\linewidth}\raggedright
% mdparagraph:P132:begin
\textbf{F18.}
% mdparagraph:P132:end

% item: NEW-U05-F18
% mdparagraph:P133:begin
Give the selected indicative forms of swęrian: present 1sg, 2sg and 3sg; preterite 2sg.
% mdparagraph:P133:end

\end{minipage}\par\addvspace{4pt}

\par\begin{minipage}{\linewidth}\raggedright
% mdparagraph:P134:begin
\textbf{F19.}
% mdparagraph:P134:end

% item: NEW-U05-F19
% mdparagraph:P135:begin
Give the selected indicative forms of feallan, healdan and weaxan: present 3sg of each; preterite 1pl of weaxan.
% mdparagraph:P135:end

\end{minipage}\par\addvspace{4pt}

\par\begin{minipage}{\linewidth}\raggedright
% mdparagraph:P136:begin
\textbf{F20.}
% mdparagraph:P136:end

% item: NEW-U05-F20
% mdparagraph:P137:begin
Give the selected indicative forms of cnāwan and wēpan: present 2sg and preterite 2sg of each.
% mdparagraph:P137:end

\end{minipage}\par\addvspace{4pt}

\par\begin{minipage}{\linewidth}\raggedright
% mdparagraph:P138:begin
\textbf{F21.}
% mdparagraph:P138:end

% item: NEW-U05-F21
% mdparagraph:P139:begin
Give the selected indicative forms of hēawan and flōwan: present 2sg and preterite 2sg of each.
% mdparagraph:P139:end

\end{minipage}\par\addvspace{4pt}

\par\begin{minipage}{\linewidth}\raggedright
% mdparagraph:P140:begin
\textbf{F22.}
% mdparagraph:P140:end

% item: NEW-U05-F22
% mdparagraph:P141:begin
Give the selected indicative forms of lǣtan, hātan and fōn: present 3sg of each; preterite 2sg of fōn.
% mdparagraph:P141:end

\end{minipage}\par\addvspace{4pt}

\par\begin{minipage}{\linewidth}\raggedright
% mdparagraph:P142:begin
\textbf{F23.}
% mdparagraph:P142:end

% item: NEW-U05-F23
% mdparagraph:P143:begin
Give the selected indicative forms of geswīcan and scīnan: present 3sg and preterite 2sg of each.
% mdparagraph:P143:end

\end{minipage}\par\addvspace{4pt}

\par\begin{minipage}{\linewidth}\raggedright
% mdparagraph:P144:begin
\textbf{F24.}
% mdparagraph:P144:end

% item: NEW-U05-F24
% mdparagraph:P145:begin
Give the selected indicative forms of flēogan and swimman: present 3sg and preterite 3pl of each.
% mdparagraph:P145:end

\end{minipage}\par\addvspace{4pt}

\reserve{7\baselineskip}
% mdparagraph:P146:begin
\subsection*{Clauses — C1–C8}\BookHeading{nav-ch05-learner-0013}{subsubsection}{Clauses — C1–C8}{0}
% mdparagraph:P146:end

\reserve{5\baselineskip}
\par\begin{minipage}{\linewidth}\raggedright
% mdparagraph:P147:begin
\textbf{C1.}
% mdparagraph:P147:end

% item: NEW-U05-C01
% mdparagraph:P148:begin
Translate: You (singular) carry the food, and I take the fish.
% mdparagraph:P148:end

\end{minipage}\par\addvspace{4pt}

\par\begin{minipage}{\linewidth}\raggedright
% mdparagraph:P149:begin
\textbf{C2.}
% mdparagraph:P149:end

% item: NEW-U05-C02
% mdparagraph:P150:begin
Translate using giefan: We (two) gave the men our food (belonging to us two).
% mdparagraph:P150:end

\end{minipage}\par\addvspace{4pt}

\par\begin{minipage}{\linewidth}\raggedright
% mdparagraph:P151:begin
\textbf{C3.}
% mdparagraph:P151:end

% item: NEW-U05-C03
% mdparagraph:P152:begin
Translate: The companions sat in the dwelling, and the man stood on the shore.
% mdparagraph:P152:end

\end{minipage}\par\addvspace{4pt}

\par\begin{minipage}{\linewidth}\raggedright
% mdparagraph:P153:begin
\textbf{C4.}
% mdparagraph:P153:end

% item: NEW-U05-C04
% mdparagraph:P154:begin
Translate using etan: I did not eat the food; you (singular) ate it.
% mdparagraph:P154:end

\end{minipage}\par\addvspace{4pt}

\par\begin{minipage}{\linewidth}\raggedright
% mdparagraph:P155:begin
\textbf{C5.}
% mdparagraph:P155:end

% item: NEW-U05-C05
% mdparagraph:P156:begin
Rewrite with hīe as subject and plural objects: Hē geseah ðone fugol and nōm ðone fisc.
% mdparagraph:P156:end

\end{minipage}\par\addvspace{4pt}

\par\begin{minipage}{\linewidth}\raggedright
% mdparagraph:P157:begin
\textbf{C6.}
% mdparagraph:P157:end

% item: NEW-U05-C06
% mdparagraph:P158:begin
Translate using cweðan to introduce direct speech: The man says, “I see the trees.”
% mdparagraph:P158:end

\end{minipage}\par\addvspace{4pt}

\par\begin{minipage}{\linewidth}\raggedright
% mdparagraph:P159:begin
\textbf{C7.}
% mdparagraph:P159:end

% item: NEW-U05-C07
% mdparagraph:P160:begin
Translate using geswīcan: You (plural) cease from your work.
% mdparagraph:P160:end

\end{minipage}\par\addvspace{4pt}

\par\begin{minipage}{\linewidth}\raggedright
% mdparagraph:P161:begin
\textbf{C8.}
% mdparagraph:P161:end

% item: NEW-U05-C08
% mdparagraph:P162:begin
Translate: The tree falls, and the men hold the ship.
% mdparagraph:P162:end

\end{minipage}\par\addvspace{4pt}

\reserve{7\baselineskip}
% mdparagraph:P163:begin
\section*{9. Cook I.5}\BookHeading{nav-ch05-learner-0014}{subsection}{9. Cook I.5}{1}
% mdparagraph:P163:end

% mdparagraph:P164:begin
Translate all sixteen sentences into Old English. These are Cook’s unchanged English prompts from \textbf{Exercises in Old English}, Group I, Exercise 5, printed p.3 (one-based PDF p.11). Use the ordinary constructions taught above and the earlier weak-verb and nominal systems. Original item numbers are retained.
% mdparagraph:P164:end

\reserve{5\baselineskip}
\par\begin{minipage}{\linewidth}\raggedright
% mdparagraph:P165:begin
\textbf{I.5.1.}
% mdparagraph:P165:end

% item: EX-I-05-01
% mdparagraph:P166:begin
God created the heavens and the earth.
% mdparagraph:P166:end

\end{minipage}\par\addvspace{4pt}

\par\begin{minipage}{\linewidth}\raggedright
% mdparagraph:P167:begin
\textbf{I.5.2.}
% mdparagraph:P167:end

% item: EX-I-05-02
% mdparagraph:P168:begin
God made two great lights.
% mdparagraph:P168:end

\end{minipage}\par\addvspace{4pt}

\par\begin{minipage}{\linewidth}\raggedright
% mdparagraph:P169:begin
\textbf{I.5.3.}
% mdparagraph:P169:end

% item: EX-I-05-03
% mdparagraph:P170:begin
God made the firmament.
% mdparagraph:P170:end

\end{minipage}\par\addvspace{4pt}

\par\begin{minipage}{\linewidth}\raggedright
% mdparagraph:P171:begin
\textbf{I.5.4.}
% mdparagraph:P171:end

% item: EX-I-05-04
% mdparagraph:P172:begin
God created the great whales.
% mdparagraph:P172:end

\end{minipage}\par\addvspace{4pt}

\par\begin{minipage}{\linewidth}\raggedright
% mdparagraph:P173:begin
\textbf{I.5.5.}
% mdparagraph:P173:end

% item: EX-I-05-05
% mdparagraph:P174:begin
He made the beasts of the earth.
% mdparagraph:P174:end

\end{minipage}\par\addvspace{4pt}

% teaching: NEW-U05-GRAM-U05-SOURCE-ALL-THOSE-THINGS
% teaching: NEW-U05-SUPFORM-eall-acc-n-pl-strong
% teaching: NEW-U05-SUPSENSE-eall-all
% mdparagraph:P175:begin
\textbf{Supplied phrase for I.5.6:} eall ðā ðing — all those things (accusative plural); use the supplied phrase. Here eall is strong neuter accusative plural, ðā is demonstrative accusative plural, and ðing is neuter accusative plural. The whole phrase remains available; neither eall nor the quantifier construction requires independent production.
% mdparagraph:P175:end

\reserve{5\baselineskip}
\par\begin{minipage}{\linewidth}\raggedright
% mdparagraph:P176:begin
\textbf{I.5.6.}
% mdparagraph:P176:end

% item: EX-I-05-06
% mdparagraph:P177:begin
He saw all those things.
% mdparagraph:P177:end

\end{minipage}\par\addvspace{4pt}

\par\begin{minipage}{\linewidth}\raggedright
% mdparagraph:P178:begin
\textbf{I.5.7.}
% mdparagraph:P178:end

% item: EX-I-05-07
% mdparagraph:P179:begin
He finished his work.
% mdparagraph:P179:end

\end{minipage}\par\addvspace{4pt}

\par\begin{minipage}{\linewidth}\raggedright
% mdparagraph:P180:begin
\textbf{I.5.8.}
% mdparagraph:P180:end

% item: EX-I-05-08
% mdparagraph:P181:begin
He blesses the seventh day.
% mdparagraph:P181:end

\end{minipage}\par\addvspace{4pt}

\par\begin{minipage}{\linewidth}\raggedright
% mdparagraph:P182:begin
\textbf{I.5.9.}
% mdparagraph:P182:end

% item: EX-I-05-09
% mdparagraph:P183:begin
The man ceases from his work.
% mdparagraph:P183:end

\end{minipage}\par\addvspace{4pt}

\par\begin{minipage}{\linewidth}\raggedright
% mdparagraph:P184:begin
\textbf{I.5.10.}
% mdparagraph:P184:end

% item: EX-I-05-10
% mdparagraph:P185:begin
I see the bird in the tree.
% mdparagraph:P185:end

\end{minipage}\par\addvspace{4pt}

\par\begin{minipage}{\linewidth}\raggedright
% mdparagraph:P186:begin
\textbf{I.5.11.}
% mdparagraph:P186:end

% item: EX-I-05-11
% mdparagraph:P187:begin
They saw the fish in the water.
% mdparagraph:P187:end

\end{minipage}\par\addvspace{4pt}

\par\begin{minipage}{\linewidth}\raggedright
% mdparagraph:P188:begin
\textbf{I.5.12.}
% mdparagraph:P188:end

% item: EX-I-05-12
% mdparagraph:P189:begin
The birds of the air fly.
% mdparagraph:P189:end

\end{minipage}\par\addvspace{4pt}

\par\begin{minipage}{\linewidth}\raggedright
% mdparagraph:P190:begin
\textbf{I.5.13.}
% mdparagraph:P190:end

% item: EX-I-05-13
% mdparagraph:P191:begin
The trees grew.
% mdparagraph:P191:end

\end{minipage}\par\addvspace{4pt}

\par\begin{minipage}{\linewidth}\raggedright
% mdparagraph:P192:begin
\textbf{I.5.14.}
% mdparagraph:P192:end

% item: EX-I-05-14
% mdparagraph:P193:begin
The stars shone.
% mdparagraph:P193:end

\end{minipage}\par\addvspace{4pt}

\par\begin{minipage}{\linewidth}\raggedright
% mdparagraph:P194:begin
\textbf{I.5.15.}
% mdparagraph:P194:end

% item: EX-I-05-15
% mdparagraph:P195:begin
The fish swims.
% mdparagraph:P195:end

\end{minipage}\par\addvspace{4pt}

\par\begin{minipage}{\linewidth}\raggedright
% mdparagraph:P196:begin
\textbf{I.5.16.}
% mdparagraph:P196:end

% item: EX-I-05-16
% mdparagraph:P197:begin
The men swam.
% mdparagraph:P197:end

\end{minipage}\par\addvspace{4pt}

\reserve{7\baselineskip}
% mdparagraph:P198:begin
\section*{10. Connected reading: Food for the companions and the birds}\BookHeading{nav-ch05-learner-0015}{subsection}{10. Connected reading: Food for the companions and the birds}{1}
% mdparagraph:P198:end

% mdparagraph:P199:begin
This account is new teaching material. The preterites narrate successive actions. Read the pronouns in context: hīe after the giving clause can naturally refer to the companions, but its form alone does not force that antecedent. Ān mann does not say whether the man belongs to the original pair or their companions.
% mdparagraph:P199:end

% composition: NEW-U05-TEXT-A
% mdparagraph:P200:begin
\par\begin{minipage}{\linewidth}\raggedright
\readingtext{Twēgen męnn bǣron męte. Hiera gefēran sǣton on ðǣm hofe. Ðā męnn gēafon ðǣm gefērum ðone męte, and hīe ǣton. Ān mann ne ǣt. Hē geaf his męte ðǣm fuglum. Ðā fuglas ǣton and flugon. Se mann sæt on ðǣm strande and geseah ðā fiscas. Ðā fiscas swummon on ðǣm wætere.}
\end{minipage}\par
% mdparagraph:P200:end

\reserve{7\baselineskip}
% mdparagraph:P201:begin
\subsection*{Connected reading — T1–T4}\BookHeading{nav-ch05-learner-0016}{subsubsection}{Connected reading — T1–T4}{0}
% mdparagraph:P201:end

\reserve{5\baselineskip}
\par\begin{minipage}{\linewidth}\raggedright
% mdparagraph:P202:begin
\textbf{T1.}
% mdparagraph:P202:end

% item: NEW-U05-T01
% mdparagraph:P203:begin
In the new account, what did one man do with his food, and what did its recipients do next?
% mdparagraph:P203:end

\end{minipage}\par\addvspace{4pt}

\par\begin{minipage}{\linewidth}\raggedright
% mdparagraph:P204:begin
\textbf{T2.}
% mdparagraph:P204:end

% item: NEW-U05-T02
% mdparagraph:P205:begin
In ðǣm gefērum ðone męte, identify the recipient and the thing given, and explain their cases. What subject controls gēafon?
% mdparagraph:P205:end

\end{minipage}\par\addvspace{4pt}

\par\begin{minipage}{\linewidth}\raggedright
% mdparagraph:P206:begin
\textbf{T3.}
% mdparagraph:P206:end

% item: NEW-U05-T03
% mdparagraph:P207:begin
Retell the last two sentences of the new account in the present indicative. Keep the subjects, objects and location phrases.
% mdparagraph:P207:end

\end{minipage}\par\addvspace{4pt}

\par\begin{minipage}{\linewidth}\raggedright
% mdparagraph:P208:begin
\textbf{T4.}
% mdparagraph:P208:end

% item: NEW-U05-T04
% mdparagraph:P209:begin
Translate the complete new account.
% mdparagraph:P209:end

\end{minipage}\par\addvspace{4pt}

\reserve{7\baselineskip}
% mdparagraph:P210:begin
\section*{11. Source reading: The first day revisited}\BookHeading{nav-ch05-learner-0017}{subsection}{11. Source reading: The first day revisited}{1}
% mdparagraph:P210:end

% mdparagraph:P211:begin
This is the same complete first-day passage read in Unit 1, now revisited with the strong finite forms learned in this unit. The original text and assignment are unchanged. Use all the local help below for the constructions that remain supplied. Cook prints the passage in \textbf{A First Book in Old English}, Reader I, pp.123–124 (one-based PDF pp.139–140).
% mdparagraph:P211:end

% source: FB-I-CREATION-FIRST-DAY-01
% mdparagraph:P212:begin
\par\begin{minipage}{\linewidth}\raggedright
\readingtext{On anginne gescēop God heofonan and eorðan. Sēo eorðe sōðlice wæs īdel and ǣmtig; and ðīestru wǣron ofer ðǣre nēowolnesse brādnesse; and Godes gāst wæs gefęred ofer wæteru. God cwæð ðā, “Geweorðe lēoht”; and lēoht wearð geworht. God geseah ðā ðæt hit gōd wæs; and hē gedǣlde ðæt lēoht fram ðǣm ðīestrum. And hēt ðæt lēoht Dæg, and ðā ðīestru Niht. Ðā wæs geworden ǣfen and morgen ān dæg.}
\end{minipage}\par
% mdparagraph:P212:end

\reserve{7\baselineskip}
% mdparagraph:P213:begin
\subsection*{Local help for the first day}\BookHeading{nav-ch05-learner-0018}{subsubsection}{Local help for the first day}{0}
% mdparagraph:P213:end

% mdparagraph:P214:begin
Each explanation applies to the actual printed phrase and remains available during R1–R4. Learned indicative morphology does not make these other constructions independent production targets.
% mdparagraph:P214:end

\reserve{5\baselineskip}
% mdparagraph:P215:begin
\textbf{On anginne}
% mdparagraph:P215:end

% teaching: NEW-U05-SUPSENSE-on-time-supplied
% teaching: NEW-U05-GRAM-READING-ON-ANGINNE
% teaching: NEW-U05-SUPPORT-U05-SRC-D1-ON-BEGINNING
% mdparagraph:P216:begin
On anginne: in the beginning.
% mdparagraph:P216:end

% mdparagraph:P217:begin
This is temporal on with a dative. The whole time phrase is supplied; the independent on-construction already practised states location.
% mdparagraph:P217:end

\reserve{5\baselineskip}
% mdparagraph:P218:begin
\textbf{heofonan}
% mdparagraph:P218:end

% teaching: NEW-U05-SUPFORM-heofone-acc-sg
% teaching: NEW-U05-SUPSENSE-heofone-1
% teaching: NEW-U05-GRAM-READING-HEOFONE
% teaching: NEW-U05-SUPPORT-U05-SRC-D1-HEAVEN
% mdparagraph:P219:begin
heofonan: heaven/\allowbreak{}the sky, accusative singular of weak feminine heofone.
% mdparagraph:P219:end

% mdparagraph:P220:begin
Heofone is a weak feminine noun, and this exact heofonan is accusative singular. It is distinct from the learned masculine heofon and its plural heofonas. No further forms of this feminine noun are required here.
% mdparagraph:P220:end

\reserve{5\baselineskip}
% mdparagraph:P221:begin
\textbf{sōðlice}
% mdparagraph:P221:end

% teaching: NEW-U05-SUPSENSE-sothlice-1
% teaching: NEW-U05-GRAM-U05-SOURCE-CONNECTIVE-SOTHLICE
% teaching: NEW-U05-SUPPORT-U05-SRC-D1-CONNECTIVE
% mdparagraph:P222:begin
sōðlice: now; for its part.
% mdparagraph:P222:end

% mdparagraph:P223:begin
This connective use differs from the earlier ordinary meaning “truly.” Cook’s note gives Latin autem. Keep the contextual meaning supplied here.
% mdparagraph:P223:end

\reserve{5\baselineskip}
% mdparagraph:P224:begin
\textbf{ofer ðǣre nēowolnesse brādnesse;}
% mdparagraph:P224:end

% teaching: NEW-U05-SUPFORM-ofer-invariant
% teaching: NEW-U05-SUPSENSE-ofer-1
% teaching: NEW-U05-GRAM-READING-OFER-DATIVE
% teaching: NEW-U05-SUPPORT-U05-SRC-D1-OFER-DEEP
% mdparagraph:P225:begin
ofer ðǣre nēowolnesse brādnesse: over the surface of the deep.
% mdparagraph:P225:end

% mdparagraph:P226:begin
Ofer governs dative brādnesse. The preceding ðǣre nēowolnesse is feminine genitive singular, dependent on brādnesse. Learn the analysis of this exact phrase without extending it to a complete rule for ofer.
% mdparagraph:P226:end

\reserve{5\baselineskip}
% mdparagraph:P227:begin
\textbf{wæs gefęred ofer wæteru.}
% mdparagraph:P227:end

% teaching: NEW-U05-SUPFORM-geferian-past-participle
% teaching: NEW-U05-SUPSENSE-geferian-1
% teaching: NEW-U05-GRAM-U05-SOURCE-CARRIED-OVER-WATERS
% teaching: NEW-U05-SUPPORT-U05-SRC-D1-CARRIED
% mdparagraph:P228:begin
wæs gefęred ofer wæteru: was borne/\allowbreak{}carried (moved) over the waters.
% mdparagraph:P228:end

% mdparagraph:P229:begin
Wæs is finite; gefęred is the supplied participle in a passive expression. Here ofer has accusative plural wæteru. Both the passive and this local prepositional construction remain supplied.
% mdparagraph:P229:end

\reserve{5\baselineskip}
% mdparagraph:P230:begin
\textbf{ðā /\allowbreak{} Ðā}
% mdparagraph:P230:end

% teaching: NEW-U05-SUPFORM-tha-adverb-invariant
% teaching: NEW-U05-SUPSENSE-tha-adverb-1
% teaching: NEW-U05-GRAM-READING-TEMPORAL-THA
% teaching: NEW-U05-SUPPORT-U05-SRC-D1-THA
% mdparagraph:P231:begin
ðā /\allowbreak{} Ðā: then.
% mdparagraph:P231:end

% mdparagraph:P232:begin
These occurrences are temporal adverbs. They are distinct from the identically spelled plural article or demonstrative.
% mdparagraph:P232:end

\reserve{5\baselineskip}
% mdparagraph:P233:begin
\textbf{“Geweorðe lēoht”;}
% mdparagraph:P233:end

% teaching: NEW-U05-SUPFORM-geweorthan-pres-subj-3-sg
% teaching: NEW-U05-SUPSENSE-geweorthan-1
% teaching: NEW-U05-GRAM-READING-OPTATIVE-COMMAND
% teaching: NEW-U05-SUPPORT-U05-SRC-D1-WISH
% mdparagraph:P234:begin
Geweorðe lēoht: Let light come to be.
% mdparagraph:P234:end

% mdparagraph:P235:begin
Geweorðe is a third-singular present subjunctive used as an optative, expressing a wish or command. The whole expression is supplied; it is outside the indicative tables.
% mdparagraph:P235:end

\reserve{5\baselineskip}
% mdparagraph:P236:begin
\textbf{wearð geworht.}
% mdparagraph:P236:end

% teaching: NEW-U05-SUPFORM-gewyrcean-past-participle
% teaching: NEW-U05-SUPSENSE-weorthan-1
% teaching: NEW-U05-GRAM-READING-PASSIVE-GEWORHT
% teaching: NEW-U05-SUPPORT-U05-SRC-D1-MADE
% mdparagraph:P237:begin
wearð geworht: was made; came into being.
% mdparagraph:P237:end

% mdparagraph:P238:begin
Wearð is the finite preterite. The ordinary meaning of weorðan, “become,” was learned earlier; here wearð serves with the supplied participle geworht in a passive expression. That auxiliary use remains supplied.
% mdparagraph:P238:end

\reserve{5\baselineskip}
% mdparagraph:P239:begin
\textbf{geseah ðā ðæt hit gōd wæs;}
% mdparagraph:P239:end

% teaching: NEW-U05-SUPFORM-thaet-conjunction-invariant
% teaching: NEW-U05-SUPSENSE-thaet-conjunction-1
% teaching: NEW-U05-GRAM-READING-CONTENT-CLAUSE
% teaching: NEW-U05-SUPPORT-U05-SRC-D1-CONTENT
% mdparagraph:P240:begin
God geseah … ðæt hit gōd wæs: God saw … that it was good.
% mdparagraph:P240:end

% mdparagraph:P241:begin
Geseah is now a learned indicative form. The conjunction ðæt introduces the supplied content clause; its final wæs belongs inside that clause. Hit refers to neuter lēoht. This ðæt means “that,” not “the.”
% mdparagraph:P241:end

\reserve{5\baselineskip}
% mdparagraph:P242:begin
\textbf{gedǣlde ðæt lēoht fram ðǣm ðīestrum.}
% mdparagraph:P242:end

% teaching: NEW-U05-SUPFORM-gedaelan-pret-3-sg
% teaching: NEW-U05-SUPSENSE-gedaelan-1
% teaching: NEW-U05-GRAM-U05-SOURCE-GEDAELDE
% teaching: NEW-U05-SUPPORT-U05-SRC-D1-DIVIDE
% mdparagraph:P243:begin
gedǣlde ðæt lēoht fram ðǣm ðīestrum: separated the light from the darkness.
% mdparagraph:P243:end

% mdparagraph:P244:begin
Gedǣlde is the supplied weak third-singular preterite of gedǣlan, “divide; separate.” The separating fram + dative construction is already familiar; this additional verb form remains supplied.
% mdparagraph:P244:end

\reserve{5\baselineskip}
% mdparagraph:P245:begin
\textbf{hēt ðæt lēoht Dæg, and ðā ðīestru Niht.}
% mdparagraph:P245:end

% teaching: NEW-U05-GRAM-READING-NAMING
% teaching: NEW-U05-SUPPORT-U05-SRC-D1-NAMING
% mdparagraph:P246:begin
He called the light Day and the darkness Night.
% mdparagraph:P246:end

% mdparagraph:P247:begin
The object and the name given to it are both accusative: lēoht/\allowbreak{}Dæg and ðīestru/\allowbreak{}Niht. The continuing subject is God; the second naming clause also leaves the verb understood. The finite hēt is learned, while this two-accusative naming construction remains supplied.
% mdparagraph:P247:end

\reserve{5\baselineskip}
% mdparagraph:P248:begin
\textbf{Ðā wæs geworden ǣfen and morgen ān dæg.}
% mdparagraph:P248:end

% teaching: NEW-U05-SUPFORM-geweorthan-past-participle
% teaching: NEW-U05-GRAM-READING-RESULTATIVE-ENDING
% teaching: NEW-U05-SUPPORT-U05-SRC-D1-CLOSE
% mdparagraph:P249:begin
Then evening and morning had come to be—one day.
% mdparagraph:P249:end

% mdparagraph:P250:begin
Wæs geworden is a supplied resultative/\allowbreak{}perfect expression. It describes what has come to be, rather than the passive making in wearð geworht. Preserve ān dæg as “one day.” The formula is given whole; no unique subject/\allowbreak{}predicate division or new agreement rule is required.
% mdparagraph:P250:end

\reserve{7\baselineskip}
% mdparagraph:P251:begin
\subsection*{First-day revisit — R1–R4}\BookHeading{nav-ch05-learner-0019}{subsubsection}{First-day revisit — R1–R4}{0}
% mdparagraph:P251:end

\reserve{5\baselineskip}
\par\begin{minipage}{\linewidth}\raggedright
% mdparagraph:P252:begin
\textbf{R1.}
% mdparagraph:P252:end

% item: NEW-U05-R01
% mdparagraph:P253:begin
In the first-day text, identify the lemma, strong class, tense, person and number of gescēop, cwæð, geseah and hēt. Use the clauses to identify the subject of hēt.
% mdparagraph:P253:end

\end{minipage}\par\addvspace{4pt}

% teaching: NEW-U05-CIT-weorthan
% mdparagraph:P254:begin
\textbf{Supplied citation for R2:} weorðan, “become.” The headword is printed for identification; no infinitival construction or unaided citation recall is required.
% mdparagraph:P254:end

\reserve{5\baselineskip}
\par\begin{minipage}{\linewidth}\raggedright
% mdparagraph:P255:begin
\textbf{R2.}
% mdparagraph:P255:end

% item: NEW-U05-R02
% mdparagraph:P256:begin
Use the first-day notes to distinguish the ordinary meaning of weorðan from its use in wearð geworht. Which word is finite in that supplied phrase?
% mdparagraph:P256:end

\end{minipage}\par\addvspace{4pt}

\par\begin{minipage}{\linewidth}\raggedright
% mdparagraph:P257:begin
\textbf{R3.}
% mdparagraph:P257:end

% item: NEW-U05-R03
% mdparagraph:P258:begin
The first-day text has heofonan; Cook I.5.1 says the heavens. Using the source note, explain the difference in noun, gender and number.
% mdparagraph:P258:end

\end{minipage}\par\addvspace{4pt}

\par\begin{minipage}{\linewidth}\raggedright
% mdparagraph:P259:begin
\textbf{R4.}
% mdparagraph:P259:end

% item: NEW-U05-R04
% mdparagraph:P260:begin
Translate the complete first-day passage, using the printed local help where needed.
% mdparagraph:P260:end

\end{minipage}\par\addvspace{4pt}

\reserve{7\baselineskip}
% mdparagraph:P261:begin
\section*{12. Source reading: The second day}\BookHeading{nav-ch05-learner-0020}{subsection}{12. Source reading: The second day}{1}
% mdparagraph:P261:end

% mdparagraph:P262:begin
Continue with the complete second-day paragraph, Cook’s Reader I, p.124, editorial prose lines 4–10 (one-based PDF p.140). The edition prints marginal line anchors at 5 and 10; the intervening line count is editorial. The new finite knowledge helps with the narrative, while the wishes, relative clauses and closing expressions retain local help.
% mdparagraph:P262:end

% source: FB-I-CREATION-SECOND-DAY-01
% mdparagraph:P263:begin
\par\begin{minipage}{\linewidth}\raggedright
\readingtext{God cwæð ðā ęft, “Geweorðe nū fæstnes tōmiddes ðǣm wæterum, and tōtwǣme ðā wæteru fram ðǣm wæterum.” And God geworhte ðā fæstnesse, and tōtwǣmde ðā wæteru ðe wǣron under ðǣre fæstnesse fram ðǣm ðe wǣron bufan ðǣre fæstnesse; hit wæs ðā swā gedōn. And God hēt ðā fæstnesse Heofonan. And wæs ðā geworden ǣfen and morgen ōðer dæg.}
\end{minipage}\par
% mdparagraph:P263:end

\reserve{7\baselineskip}
% mdparagraph:P264:begin
\subsection*{Local help for the second day}\BookHeading{nav-ch05-learner-0021}{subsubsection}{Local help for the second day}{0}
% mdparagraph:P264:end

% mdparagraph:P265:begin
The following six groups of help remain available during R5–R8. They explain the actual forms and their relations without requiring independent production of the supplied constructions.
% mdparagraph:P265:end

\reserve{5\baselineskip}
% mdparagraph:P266:begin
\textbf{ðā, ęft, nū, swā}
% mdparagraph:P266:end

% teaching: NEW-U05-SUPFORM-eft-invariant
% teaching: NEW-U05-SUPSENSE-eft-again-also
% teaching: NEW-U05-SUPFORM-nu-invariant
% teaching: NEW-U05-SUPSENSE-nu-1
% teaching: NEW-U05-SUPFORM-swa-invariant
% teaching: NEW-U05-SUPSENSE-swa-1
% teaching: NEW-U05-GRAM-U05-SOURCE-SECOND-DAY-ADVERBS
% teaching: NEW-U05-SUPPORT-U05-SRC-D2-ADVERBS
% mdparagraph:P267:begin
ðā: then; ęft: again/\allowbreak{}also; nū: now; swā: so/\allowbreak{}thus.
% mdparagraph:P267:end

% mdparagraph:P268:begin
These exact adverbial uses are supplied. Cook glosses ęft with Latin quoque, “also.” Preserve the printed ę.
% mdparagraph:P268:end

\reserve{5\baselineskip}
% mdparagraph:P269:begin
\textbf{“Geweorðe nū fæstnes tōmiddes ðǣm wæterum, and tōtwǣme ðā wæteru fram ðǣm wæterum.”}
% mdparagraph:P269:end

% teaching: NEW-U05-SUPFORM-tomiddes-invariant
% teaching: NEW-U05-SUPSENSE-tomiddes-midst
% teaching: NEW-U05-SUPFORM-totwaeman-pres-subj-3-sg
% teaching: NEW-U05-SUPSENSE-totwaeman-separate
% teaching: NEW-U05-GRAM-U05-SOURCE-SECOND-DAY-WISH
% teaching: NEW-U05-SUPPORT-U05-SRC-D2-WISH
% mdparagraph:P270:begin
Let a firmament come to be in the midst of the waters, and let it separate the waters from the waters.
% mdparagraph:P270:end

% mdparagraph:P271:begin
Geweorðe and tōtwǣme are coordinated third-singular present optatives; fæstnes is the subject of both. Tōmiddes has dative ðǣm wæterum. Tōtwǣme, “let it separate,” has accusative ðā wæteru and fram + dative. The whole wish and its new preposition remain supplied.
% mdparagraph:P271:end

\reserve{5\baselineskip}
% mdparagraph:P272:begin
\textbf{tōtwǣmde ðā wæteru ðe wǣron under ðǣre fæstnesse fram ðǣm ðe wǣron bufan ðǣre fæstnesse;}
% mdparagraph:P272:end

% teaching: NEW-U05-SUPFORM-totwaeman-pret-3-sg
% teaching: NEW-U05-SUPFORM-the-relative-subject-invariant
% teaching: NEW-U05-SUPSENSE-the-relative-subject-waters
% teaching: NEW-U05-SUPFORM-under-invariant
% teaching: NEW-U05-SUPSENSE-under-below
% teaching: NEW-U05-SUPFORM-bufan-invariant
% teaching: NEW-U05-SUPSENSE-bufan-above
% teaching: NEW-U05-SUPSENSE-se-those-waters-substantival
% teaching: NEW-U05-GRAM-U05-SOURCE-SECOND-DAY-RELATIVES
% teaching: NEW-U05-SUPPORT-U05-SRC-D2-SEPARATING-RELATIVES
% mdparagraph:P273:begin
and separated the waters which were under the firmament from those which were above the firmament.
% mdparagraph:P273:end

% mdparagraph:P274:begin
Tōtwǣmde is the supplied third-singular weak preterite, “separated.” In each relative clause, ðe is the subject of wǣron. The first refers to the waters below the firmament. After fram, dative ðǣm means “those (waters)”; its following relative describes the waters above. Under and bufan both govern dative ðǣre fæstnesse here. These subject-relative clauses and the omitted noun after “those” remain supplied.
% mdparagraph:P274:end

\reserve{5\baselineskip}
% mdparagraph:P275:begin
\textbf{hit wæs ðā swā gedōn.}
% mdparagraph:P275:end

% teaching: NEW-U05-SUPFORM-gedon-past-participle
% teaching: NEW-U05-SUPSENSE-gedon-do-make
% teaching: NEW-U05-GRAM-U05-SOURCE-SECOND-DAY-FULFILLMENT
% teaching: NEW-U05-SUPPORT-U05-SRC-D2-FULFILLMENT
% mdparagraph:P276:begin
hit wæs ðā swā gedōn: it was then done/\allowbreak{}made so.
% mdparagraph:P276:end

% mdparagraph:P277:begin
Wæs is finite; gedōn is the supplied past participle of gedōn, “do; make.” The whole passive/\allowbreak{}resultative expression describes the fulfillment of the preceding action.
% mdparagraph:P277:end

\reserve{5\baselineskip}
% mdparagraph:P278:begin
\textbf{God hēt ðā fæstnesse Heofonan.}
% mdparagraph:P278:end

% teaching: NEW-U05-GRAM-U05-SOURCE-SECOND-DAY-NAMING
% teaching: NEW-U05-SUPPORT-U05-SRC-D2-NAMING
% mdparagraph:P279:begin
God called the firmament Heaven.
% mdparagraph:P279:end

% mdparagraph:P280:begin
Fæstnesse is the accusative singular object; Heofonan is the accusative singular name, from weak feminine heofone. The hēt form is learned, while this naming construction and feminine noun form remain supplied.
% mdparagraph:P280:end

\reserve{5\baselineskip}
% mdparagraph:P281:begin
\textbf{And wæs ðā geworden ǣfen and morgen ōðer dæg.}
% mdparagraph:P281:end

% teaching: NEW-U05-GRAM-U05-SOURCE-SECOND-DAY-CLOSE
% teaching: NEW-U05-SUPPORT-U05-SRC-D2-CLOSE
% mdparagraph:P282:begin
And evening and morning had come to be—a second day.
% mdparagraph:P282:end

% mdparagraph:P283:begin
The resultative closing formula remains supplied. Here ōðer means “second” and is masculine nominative singular, agreeing with dæg. Contrast the first-day formula’s literal “one day.”
% mdparagraph:P283:end

\reserve{7\baselineskip}
% mdparagraph:P284:begin
\subsection*{Second day — R5–R8}\BookHeading{nav-ch05-learner-0022}{subsubsection}{Second day — R5–R8}{0}
% mdparagraph:P284:end

\reserve{5\baselineskip}
\par\begin{minipage}{\linewidth}\raggedright
% mdparagraph:P285:begin
\textbf{R5.}
% mdparagraph:P285:end

% item: NEW-U05-R05
% mdparagraph:P286:begin
In the second-day text, distinguish cwæð and hēt from geworhte and the supplied tōtwǣmde. Identify strong and weak preterites, and identify their common subject.
% mdparagraph:P286:end

\end{minipage}\par\addvspace{4pt}

\par\begin{minipage}{\linewidth}\raggedright
% mdparagraph:P287:begin
\textbf{R6.}
% mdparagraph:P287:end

% item: NEW-U05-R06
% mdparagraph:P288:begin
Using the second-day local help, explain which waters are separated from which. What does ðǣm refer to in fram ðǣm ðe wǣron bufan ðǣre fæstnesse?
% mdparagraph:P288:end

\end{minipage}\par\addvspace{4pt}

\par\begin{minipage}{\linewidth}\raggedright
% mdparagraph:P289:begin
\textbf{R7.}
% mdparagraph:P289:end

% item: NEW-U05-R07
% mdparagraph:P290:begin
What does God call the firmament in the second-day account? Explain why ōðer agrees with dæg in the closing formula, while keeping the formula itself supplied.
% mdparagraph:P290:end

\end{minipage}\par\addvspace{4pt}

\par\begin{minipage}{\linewidth}\raggedright
% mdparagraph:P291:begin
\textbf{R8.}
% mdparagraph:P291:end

% item: NEW-U05-R08
% mdparagraph:P292:begin
Translate the complete second-day passage, using the printed local help where needed.
% mdparagraph:P292:end

\end{minipage}\par\addvspace{4pt}

\clearpage
% mdparagraph:P293:begin
\section*{13. Cumulative review}\BookHeading{nav-ch05-learner-0023}{subsection}{13. Cumulative review}{1}
% mdparagraph:P293:end

% mdparagraph:P294:begin
These three exercises are new teaching material. The printed citations in X2 remain supplied; compare the finite systems without adding nonfinite constructions.
% mdparagraph:P294:end

\reserve{5\baselineskip}
\par\begin{minipage}{\linewidth}\raggedright
% mdparagraph:P295:begin
\textbf{X1.}
% mdparagraph:P295:end

% item: NEW-U05-X01
% mdparagraph:P296:begin
Compare hē geaf and hē sealde. Give the corresponding preterite clauses with ðū, and explain why the endings differ.
% mdparagraph:P296:end

\end{minipage}\par\addvspace{4pt}

% teaching: NEW-U05-CIT-lufian
% mdparagraph:P297:begin
\textbf{Supplied citations for X2:} lufian and (ge)lufod. The parentheses in (ge)lufod preserve the earlier source citation notation; they do not introduce a rule for prefixes, perfect or passive syntax.
% mdparagraph:P297:end

\reserve{5\baselineskip}
\par\begin{minipage}{\linewidth}\raggedright
% mdparagraph:P298:begin
\textbf{X2.}
% mdparagraph:P298:end

% item: NEW-U05-X02
% mdparagraph:P299:begin
The printed entries give swęrian, swōr, swōron, sworen and lufian, lufode, (ge)lufod. Explain why -ian alone does not identify a weak verb. Give the present first-person singular of each.
% mdparagraph:P299:end

\end{minipage}\par\addvspace{4pt}

\par\begin{minipage}{\linewidth}\raggedright
% mdparagraph:P300:begin
\textbf{X3.}
% mdparagraph:P300:end

% item: NEW-U05-X03
% mdparagraph:P301:begin
Within the taught indicative systems, compare gē sēoð and hē sīeð. Identify each verb, class, tense, person and number. These are form-identification examples; add no complement.
% mdparagraph:P301:end

\end{minipage}\par\addvspace{4pt}

\clearpage
% mdparagraph:P302:begin
\section*{14. Sources and original reading notes}\BookHeading{nav-ch05-learner-0024}{subsection}{14. Sources and original reading notes}{1}
% mdparagraph:P302:end

\begingroup\fontsize{10}{12.8}\selectfont
\setlength{\parskip}{1.5pt}
% mdparagraph:P303:begin
The core strong-verb systems follow Cook, \textbf{A First Book in Old English}, §§105–110, printed pp.63–67, with the earlier-class additions from §§101–104. Wright’s \textbf{Old English Grammar} and the checked dictionary entries corroborate the stated alternatives and lexical uses. Complete finite tables apply the ordinary endings to the checked selected stems; they do not claim a separate surviving quotation for every cell.
% mdparagraph:P303:end

% mdparagraph:P304:begin
The supplied participle \textbf{weaxen} is a rule-derived expansion of Cook’s and Wright’s stated Class VII model. It is not a whole-form attestation in the checked rows. Its status remains citation support only.
% mdparagraph:P304:end

% mdparagraph:P305:begin
The lexical constructions use Cook’s case rules and independently checked dictionary evidence for the ordinary verbs. The eight examples and the connected account are new compositions. Their food, companion, bird and fish references are ordinary applications of the stated uses. Diagnostic form practice supplies no additional complement syntax.
% mdparagraph:P305:end

% mdparagraph:P306:begin
The source readings preserve Cook’s Old English spelling, quantity and punctuation. Line breaks are joined and note cues are moved to this section. The first-day passage is an explicit revisit of the Unit 1 text; the second-day paragraph is the new continuation. All relevant original notes follow. References such as “53. 3.” point to Cook’s own grammar or dictionary, not to sections of this unit. The local help above supplies what is needed here.
% mdparagraph:P306:end

\reserve{7\baselineskip}
% mdparagraph:P307:begin
\subsection*{Cook’s notes to the first day}\BookHeading{nav-ch05-learner-0025}{subsubsection}{Cook’s notes to the first day}{0}
% mdparagraph:P307:end

% source-note: FB123-n1
% mdparagraph:P308:begin
\textbf{P.123, note 1.} See gescieppan, and 18.
% mdparagraph:P308:end

% source-note: FB123-n2
% mdparagraph:P309:begin
\textbf{P.123, note 2.} The order is probably determined by the Latin: creavit Deus.
% mdparagraph:P309:end

% source-note: FB123-n3
% mdparagraph:P310:begin
\textbf{P.123, note 3.} 53. 3.
% mdparagraph:P310:end

% source-note: FB123-n4
% mdparagraph:P311:begin
\textbf{P.123, note 4.} See se.
% mdparagraph:P311:end

% source-note: FB123-n5
% mdparagraph:P312:begin
\textbf{P.123, note 5.} Lat. autem.
% mdparagraph:P312:end

% source-note: FB123-n6
% mdparagraph:P313:begin
\textbf{P.123, note 6.} See wesan.
% mdparagraph:P313:end

% source-note: FB123-n7
% mdparagraph:P314:begin
\textbf{P.123, note 7.} Plural, like Lat. tenebrae.
% mdparagraph:P314:end

% source-note: FB123-n8
% mdparagraph:P315:begin
\textbf{P.123, note 8.} Governs brādnesse.
% mdparagraph:P315:end

% source-note: FB123-n9
% mdparagraph:P316:begin
\textbf{P.123, note 9.} Genitive, dependent on brādnesse (153. i).
% mdparagraph:P316:end

% source-note: FB123-n10
% mdparagraph:P317:begin
\textbf{P.123, note 10.} See 166.
% mdparagraph:P317:end

% source-note: FB123-n11
% mdparagraph:P318:begin
\textbf{P.123, note 11.} wæs gefęred = Lat. ferebatur. See gefęrian.
% mdparagraph:P318:end

% source-note: FB123-n12
% mdparagraph:P319:begin
\textbf{P.123, note 12.} See wæter, and 47. 1, 6.
% mdparagraph:P319:end

% source-note: FB123-n13
% mdparagraph:P320:begin
\textbf{P.123, note 13.} See cweðan.
% mdparagraph:P320:end

% source-note: FB123-n14
% mdparagraph:P321:begin
\textbf{P.123, note 14.} See geweorðan, and 193. a.
% mdparagraph:P321:end

% source-note: FB123-n15
% mdparagraph:P322:begin
\textbf{P.123, note 15.} See weorðan.
% mdparagraph:P322:end

% source-note: FB123-n16
% mdparagraph:P323:begin
\textbf{P.123, note 16.} Wearð geworht = facta est. See gewyrcean.
% mdparagraph:P323:end

% source-note: FB123-n17
% mdparagraph:P324:begin
\textbf{P.123, note 17.} See geseon.
% mdparagraph:P324:end

% source-note: FB123-n18
% mdparagraph:P325:begin
\textbf{P.123, note 18.} See hē.
% mdparagraph:P325:end

% source-note: FB124-n1
% mdparagraph:P326:begin
\textbf{P.124, note 1.} See wesan.
% mdparagraph:P326:end

% source-note: FB124-n2
% mdparagraph:P327:begin
\textbf{P.124, note 2.} See gedǣlan.
% mdparagraph:P327:end

% source-note: FB124-n3
% mdparagraph:P328:begin
\textbf{P.124, note 3.} See se.
% mdparagraph:P328:end

% source-note: FB124-n4
% mdparagraph:P329:begin
\textbf{P.124, note 4.} See p.123, note 7.
% mdparagraph:P329:end

% source-note: FB124-n5
% mdparagraph:P330:begin
\textbf{P.124, note 5.} See hātan, and 189, note.
% mdparagraph:P330:end

% source-note: FB124-n6
% mdparagraph:P331:begin
\textbf{P.124, note 6.} Wæs geworden = factum est. See geweorðan.
% mdparagraph:P331:end

% source-note: FB124-n7
% mdparagraph:P332:begin
\textbf{P.124, note 7.} Lat. dies unus.
% mdparagraph:P332:end

\reserve{7\baselineskip}
% mdparagraph:P333:begin
\subsection*{Cook’s notes to the second day}\BookHeading{nav-ch05-learner-0026}{subsubsection}{Cook’s notes to the second day}{0}
% mdparagraph:P333:end

% source-note: FB124-n8
% mdparagraph:P334:begin
\textbf{P.124, note 8.} See cweðan.
% mdparagraph:P334:end

% source-note: FB124-n9
% mdparagraph:P335:begin
\textbf{P.124, note 9.} Lat. quoque.
% mdparagraph:P335:end

% source-note: FB124-n10
% mdparagraph:P336:begin
\textbf{P.124, note 10.} See geweorðan, and 193. a.
% mdparagraph:P336:end

% source-note: FB124-n11
% mdparagraph:P337:begin
\textbf{P.124, note 11.} See wæter, and 47. 1, 6.
% mdparagraph:P337:end

% source-note: FB124-n12
% mdparagraph:P338:begin
\textbf{P.124, note 12.} See tōtwǣman.
% mdparagraph:P338:end

% source-note: FB124-n13
% mdparagraph:P339:begin
\textbf{P.124, note 13.} See 87. d.
% mdparagraph:P339:end

% source-note: FB124-n14
% mdparagraph:P340:begin
\textbf{P.124, note 14.} Past part. of gedōn.
% mdparagraph:P340:end

% source-note: FB124-n15
% mdparagraph:P341:begin
\textbf{P.124, note 15.} See 173.
% mdparagraph:P341:end

% source-note: FB124-n16
% mdparagraph:P342:begin
\textbf{P.124, note 16.} Lat. secundus.
% mdparagraph:P342:end

\endgroup

% book-source-body:end
\par\label{end-ch05-learner}
\clearpage\endgroup

\BookRecto
\begingroup
\renewcommand{\runningtitle}{Unit 6. Irregular verbs}
\hyphenpenalty=10000
\exhyphenpenalty=10000
\emergencystretch=3em
\newcommand{\itempara}[2]{\par\addvspace{3pt}\noindent\hangindent=1.8em\hangafter=1\makebox[1.8em][l]{#1.}#2\par}
\newcommand{\reserve}[1]{\par\penalty-100\begingroup\dimen0=\pagegoal\advance\dimen0 by-\pagetotal\ifdim\dimen0<#1\newpage\fi\endgroup}
% book-source-body:begin

\addfontfeatures{Ligatures=TeXOff}
\raggedright
\booktitle{\BookComponent{ch06-learner}{Unit 6. Frequent irregular verbs, infinitives, and commands}UNIT 6}{Frequent irregular verbs, infinitives, and commands}

% mdparagraph:P001:begin
This unit adds frequent irregular indicative systems, commands addressed to one or more people, and infinitives used after modals or to express purpose.
% mdparagraph:P001:end

% mdparagraph:P002:begin
In the tables, \textbf{1}, \textbf{2}, and \textbf{3} mean first, second, and third person; \textbf{sg.} and \textbf{pl.} mean singular and plural. \textbf{1–3 pl.} gives the same form for all three plural persons. Dual subjects take plural finite verbs. Preserve vowel length and the distinction between e and ę. Examples and exercises are new teaching material unless a source is explicitly identified.
% mdparagraph:P002:end

\reserve{7\baselineskip}
% mdparagraph:P003:begin
\section*{1. Words and meanings}\BookHeading{nav-ch06-learner-0001}{subsection}{1. Words and meanings}{1}
% mdparagraph:P003:end

% mdparagraph:P004:begin
The verb labels in the first column identify systems. Some are conventional dictionary labels rather than usable infinitives; §2 explains the distinction. Learn the meanings given here with their stated constructions.
% mdparagraph:P004:end

% teaching: NEW-U06-VOC-witan
% teaching: NEW-U06-VOC-cunnan
% teaching: NEW-U06-VOC-thurfan
% teaching: NEW-U06-VOC-durran
% teaching: NEW-U06-VOC-sculan
% teaching: NEW-U06-VOC-magan
% teaching: NEW-U06-VOC-motan
% teaching: NEW-U06-VOC-willan
% teaching: NEW-U06-VOC-don
% teaching: NEW-U06-VOC-gan
% teaching: NEW-U06-VOC-thing-matters
% teaching: NEW-U06-VOC-to
% mdtable:T01:begin
\par\begingroup
\fontsize{10.5}{13.2}\selectfont
\setlength{\LTleft}{0pt}\setlength{\LTright}{\fill}
\noindent
\begin{tabular}{@{}>{\raggedright\arraybackslash}p{26.96mm}>{\raggedright\arraybackslash}p{34.04mm}>{\raggedright\arraybackslash}p{57.37mm}@{}}
\toprule
\textbf{Word or verb label} & \textbf{Meaning in this unit} & \textbf{Construction} \\
\midrule
witan & know & accusative fact or matter known \\
cunnan & know how; be able through skill & plain infinitive \\
ðurfan & need to & plain infinitive \\
durran & dare to & plain infinitive \\
sculan & must; be obliged to & plain infinitive \\
mugan, also called magan & can; be able to & plain infinitive \\
mōtan & may; be allowed to & plain infinitive \\
willan & wish; want to & plain infinitive \\
dōn & do; perform & accusative work or thing done \\
gān & go; walk & ordinary intransitive use \\
ðing, neuter & thing; a fact or matter known & the earlier noun forms, with this additional meaning \\
tō & to & before an inflected infinitive \\
\bottomrule
\end{tabular}
\par\endgroup

% mdtable:T01:end
% mdparagraph:P005:begin
The new sense of \textbf{ðing} includes a matter known or discussed. It does not replace the earlier sense of a thing or object. No new noun declension is needed. Infinitival tō is distinguished from locally supplied destination phrases in the source reading.
% mdparagraph:P005:end

\reserve{7\baselineskip}
% mdparagraph:P006:begin
\section*{2. Complete indicative systems}\BookHeading{nav-ch06-learner-0002}{subsection}{2. Complete indicative systems}{1}
% mdparagraph:P006:end

% teaching: NEW-U06-GRAM-PRETERITE-PRESENT-INDICATIVE
\reserve{7\baselineskip}
% mdparagraph:P007:begin
\subsection*{Preterite-present verbs}\BookHeading{nav-ch06-learner-0003}{subsubsection}{Preterite-present verbs}{0}\BookLabel{bk-ch06-learner-f7a979796f68}
% mdparagraph:P007:end

% mdparagraph:P008:begin
These verbs have present forms inherited from an older strong preterite. Thus \textbf{wāt} and \textbf{cann} resemble strong past forms but have present meaning. Their own preterites have a dental element, as weak preterites do. The historical name \textbf{preterite-present} does not make their present tense past in translation.
% mdparagraph:P008:end

% mdparagraph:P009:begin
The first and third singular present coincide. The second singular has its own form, and all three plural persons share another. Learn the singular/\allowbreak{}plural stem changes and the dental preterite together.
% mdparagraph:P009:end

\reserve{12\baselineskip}
% mdparagraph:P010:begin
\textbf{Present indicative.}
% mdparagraph:P010:end

% teaching: NEW-U06-PAR-witan
% teaching: NEW-U06-PAR-cunnan
% teaching: NEW-U06-PAR-thurfan
% teaching: NEW-U06-PAR-durran
% teaching: NEW-U06-PAR-sculan
% teaching: NEW-U06-PAR-magan
% teaching: NEW-U06-PAR-motan
% table: NEW-U06-DISPLAY-PRETERITE-PRESENTS-PRES
% mdtable:T02:begin
% course-table-style: 2026-10
% course-panel:T02-P01:begin
\par\begingroup
\fontsize{10.5}{13.2}\selectfont
\setlength{\LTleft}{0pt}\setlength{\LTright}{\fill}
\noindent
\begin{tabular}{@{}>{\raggedright\arraybackslash}p{24.232mm}>{\raggedright\arraybackslash}p{22.13mm}>{\raggedright\arraybackslash}p{22.13mm}>{\raggedright\arraybackslash}p{22.13mm}>{\raggedright\arraybackslash}p{22.13mm}@{}}
\toprule
\textbf{Person and number} & \textbf{witan} & \textbf{cunnan} & \textbf{ðurfan} & \textbf{durran} \\
\midrule
1 sg. & wāt & cann & ðearf & dearr \\
2 sg. & wāst & canst & ðearft & dearst \\
3 sg. & wāt & cann & ðearf & dearr \\
1–3 pl. & witon & cunnon & ðurfon & durron \\
\bottomrule
\end{tabular}
\par\endgroup
% course-panel:T02-P01:end
\par\addvspace{5pt}
% course-panel:T02-P02:begin
\par\begingroup
\fontsize{10.5}{13.2}\selectfont
\setlength{\LTleft}{0pt}\setlength{\LTright}{\fill}
\noindent
\begin{tabular}{@{}>{\raggedright\arraybackslash}p{30.902mm}>{\raggedright\arraybackslash}p{28.221mm}>{\raggedright\arraybackslash}p{28.221mm}>{\raggedright\arraybackslash}p{28.221mm}@{}}
\toprule
\textbf{Person and number} & \textbf{sculan} & \textbf{mugan} & \textbf{mōtan} \\
\midrule
1 sg. & sceal & mæg & mōt \\
2 sg. & scealt & meaht & mōst \\
3 sg. & sceal & mæg & mōt \\
1–3 pl. & sculon & magon & mōton \\
\bottomrule
\end{tabular}
\par\endgroup
% course-panel:T02-P02:end
% mdtable:T02:end
\reserve{12\baselineskip}
% mdparagraph:P011:begin
\textbf{Preterite indicative.}
% mdparagraph:P011:end

% table: NEW-U06-DISPLAY-PRETERITE-PRESENTS-PRET
% mdtable:T03:begin
% course-table-style: 2026-10
% course-panel:T03-P01:begin
\par\begingroup
\fontsize{10.5}{13.2}\selectfont
\setlength{\LTleft}{0pt}\setlength{\LTright}{\fill}
\noindent
\begin{tabular}{@{}>{\raggedright\arraybackslash}p{24.232mm}>{\raggedright\arraybackslash}p{22.13mm}>{\raggedright\arraybackslash}p{22.13mm}>{\raggedright\arraybackslash}p{22.13mm}>{\raggedright\arraybackslash}p{22.13mm}@{}}
\toprule
\textbf{Person and number} & \textbf{witan} & \textbf{cunnan} & \textbf{ðurfan} & \textbf{durran} \\
\midrule
1 sg. & wiste & cūðe & ðorfte & dorste \\
2 sg. & wistest & cūðest & ðorftest & dorstest \\
3 sg. & wiste & cūðe & ðorfte & dorste \\
1–3 pl. & wiston & cūðon & ðorfton & dorston \\
\bottomrule
\end{tabular}
\par\endgroup
% course-panel:T03-P01:end
\par\addvspace{5pt}
% course-panel:T03-P02:begin
\par\begingroup
\fontsize{10.5}{13.2}\selectfont
\setlength{\LTleft}{0pt}\setlength{\LTright}{\fill}
\noindent
\begin{tabular}{@{}>{\raggedright\arraybackslash}p{30.902mm}>{\raggedright\arraybackslash}p{28.221mm}>{\raggedright\arraybackslash}p{28.221mm}>{\raggedright\arraybackslash}p{28.221mm}@{}}
\toprule
\textbf{Person and number} & \textbf{sculan} & \textbf{mugan} & \textbf{mōtan} \\
\midrule
1 sg. & sceolde & meahte & mōste \\
2 sg. & sceoldest & meahtest & mōstest \\
3 sg. & sceolde & meahte & mōste \\
1–3 pl. & sceoldon & meahton & mōston \\
\bottomrule
\end{tabular}
\par\endgroup
% course-panel:T03-P02:end
% mdtable:T03:end
% mdparagraph:P012:begin
For each verb, the present and preterite tables give all twelve person–number–tense cells. In the preterite, the first and third singular end in \textbf{-e}; the second singular adds \textbf{-st}, and the plural changes \textbf{-e} to \textbf{-on}. Keep each verb’s printed stem. The second singular here does not follow the strong-preterite pattern learned in Units 4–5.
% mdparagraph:P012:end

% mdparagraph:P013:begin
\textbf{Selected forms and alternatives.} This course selects witan, witon, and wiste. Wieton/\allowbreak{}wioton and the wisse-series are documented alternatives. Cook also gives can beside cann and scolde beside sceolde. For mugan, the tables select meaht and meahte; Wright records miht and mehte/\allowbreak{}mihte. Use the selected forms when an exercise asks for them; the alternatives are not errors merely because they differ.
% mdparagraph:P013:end

% teaching: NEW-U06-GRAM-DEFECTIVE-REFERENCE
% mdparagraph:P014:begin
\textbf{Citation labels and missing forms.}\BookPoint{bk-ch06-learner-494df9f87625} Cook’s \textbf{mugan} and the common label \textbf{magan} name the same system. Fulk marks the conventional infinitives *durran, *sculan, *mōtan, and *magan as unattested. Cook marks several labels as doubtful. A label ending in -an is therefore not sufficient evidence for an infinitive that can be put into a sentence. Do not make an inflected infinitive or a command by adding ordinary endings to these labels.
% mdparagraph:P014:end

% mdparagraph:P015:begin
The labels witan, cunnan, and ðurfan are attested uninflected infinitives, but their infinitival uses and complete command systems are not taught here. Learn the finite systems above and the selected infinitive pairs below. A missing cell is different from two genuine cells that happen to share a spelling, such as the first and third singular wāt.
% mdparagraph:P015:end

% mdparagraph:P016:begin
Cook and Wright call certain witan forms commands; Fulk describes such preterite-present command uses as subjunctive substitutes. They do not establish an ordinary imperative pair for every modal. These substitutes and first- or third-person wishes are outside this unit’s command practice.
% mdparagraph:P016:end

% mdparagraph:P017:begin
\textbf{Sources.} Albert S. Cook, \textbf{A First Book in Old English: Grammar, Reader, Notes, and Vocabulary}, second edition (1895), §§125–137, pp.76–78 /\allowbreak{} PDF pp.92–94. R. D. Fulk, \textbf{An Introductory Grammar of Old English with an Anthology of Readings}, corrected 2020 reprint, §123, pp.89–90 /\allowbreak{} PDF pp.101–102, distinguishes unattested labels and command substitutes. Joseph and Elizabeth Mary Wright, \textbf{Old English Grammar}, second edition (1914), gives the corroborating witan and magan variants; see the precise variant references in §9.
% mdparagraph:P017:end

\reserve{7\baselineskip}
% mdparagraph:P018:begin
\subsection*{Dōn, gān, and willan}\BookHeading{nav-ch06-learner-0004}{subsubsection}{Dōn, gān, and willan}{0}\BookLabel{bk-ch06-learner-ee684d5ec1f6}
% mdparagraph:P018:end

% teaching: NEW-U06-GRAM-DON-INDICATIVE
% teaching: NEW-U06-GRAM-GAN-INDICATIVE
% teaching: NEW-U06-GRAM-WILLAN-INDICATIVE
% mdparagraph:P019:begin
Learn each present with its preterite. \textbf{Dōn} changes ō to ē in dēst and dēð, and has the past dyde. \textbf{Gān} changes ā to ǣ in gǣst and gǣð; its past \textbf{ēode} comes from a different stem. \textbf{Willan} has wille in both the first and third singular present, with wilt in the second singular. Do not add an ordinary third-singular -ð to wille.
% mdparagraph:P019:end

\reserve{8\baselineskip}
% mdparagraph:P020:begin
\textbf{Present indicative.}
% mdparagraph:P020:end

% teaching: NEW-U06-PAR-willan
% teaching: NEW-U06-PAR-don
% teaching: NEW-U06-PAR-gan
% table: NEW-U06-DISPLAY-ANOMALOUS-PRES
% mdtable:T04:begin
% course-table-style: 2026-10
% course-panel:T04-P01:begin
\par\begingroup
\fontsize{10.5}{13.2}\selectfont
\setlength{\LTleft}{0pt}\setlength{\LTright}{\fill}
\noindent
\begin{tabular}{@{}>{\raggedright\arraybackslash}p{30.902mm}>{\raggedright\arraybackslash}p{28.221mm}>{\raggedright\arraybackslash}p{28.221mm}>{\raggedright\arraybackslash}p{28.221mm}@{}}
\toprule
\textbf{Person and number} & \textbf{dōn} & \textbf{gān} & \textbf{willan} \\
\midrule
1 sg. & dō & gā & wille \\
2 sg. & dēst & gǣst & wilt \\
3 sg. & dēð & gǣð & wille \\
1–3 pl. & dōð & gāð & willað \\
\bottomrule
\end{tabular}
\par\endgroup
% course-panel:T04-P01:end
% mdtable:T04:end
\reserve{8\baselineskip}
% mdparagraph:P021:begin
\textbf{Preterite indicative.}
% mdparagraph:P021:end

% table: NEW-U06-DISPLAY-ANOMALOUS-PRET
% mdtable:T05:begin
% course-table-style: 2026-10
% course-panel:T05-P01:begin
\par\begingroup
\fontsize{10.5}{13.2}\selectfont
\setlength{\LTleft}{0pt}\setlength{\LTright}{\fill}
\noindent
\begin{tabular}{@{}>{\raggedright\arraybackslash}p{30.902mm}>{\raggedright\arraybackslash}p{28.221mm}>{\raggedright\arraybackslash}p{28.221mm}>{\raggedright\arraybackslash}p{28.221mm}@{}}
\toprule
\textbf{Person and number} & \textbf{dōn} & \textbf{gān} & \textbf{willan} \\
\midrule
1 sg. & dyde & ēode & wolde \\
2 sg. & dydest & ēodest & woldest \\
3 sg. & dyde & ēode & wolde \\
1–3 pl. & dydon & ēodon & woldon \\
\bottomrule
\end{tabular}
\par\endgroup
% course-panel:T05-P01:end
% mdtable:T05:end
% mdparagraph:P022:begin
Cook also gives \textbf{wile} beside wille. The table selects the spelling with double l. The negative contracted companions are not additional paradigms to learn in this unit.
% mdparagraph:P022:end

% mdparagraph:P023:begin
\textbf{Sources.} Cook §§138–141, pp.78–80 /\allowbreak{} PDF pp.94–96. Fulk §134, p.102 /\allowbreak{} PDF p.114, corroborates the complete dōn, gān, and willan indicative systems.
% mdparagraph:P023:end

\reserve{7\baselineskip}
% mdparagraph:P024:begin
\subsection*{The alternative bēo present}\BookHeading{nav-ch06-learner-0005}{subsubsection}{The alternative bēo present}{0}\BookLabel{bk-ch06-learner-666a2007807d}
% mdparagraph:P024:end

% teaching: NEW-U06-GRAM-BEON-INDICATIVE
% mdparagraph:P025:begin
The following alternative present series belongs with the earlier verb “be.” Both present series share the familiar preterite: wæs in the first and third singular, wǣre in the second singular, and wǣron in all plural persons. There is no separate bēon past to construct.
% mdparagraph:P025:end

% teaching: NEW-U06-PAR-wesan-beon
% table: NEW-U06-DISPLAY-BEON
% mdtable:T06:begin
% course-table-style: 2026-10
% course-panel:T06-P01:begin
\par\begingroup
\fontsize{10.5}{13.2}\selectfont
\setlength{\LTleft}{0pt}\setlength{\LTright}{\fill}
\noindent
\begin{tabular}{@{}>{\raggedright\arraybackslash}p{39.459mm}>{\raggedright\arraybackslash}p{39.459mm}>{\raggedright\arraybackslash}p{39.459mm}@{}}
\toprule
\textbf{Person and number} & \textbf{Earlier eom series} & \textbf{Alternative bēo series} \\
\midrule
1 sg. & eom & bēo \\
2 sg. & eart & bist \\
3 sg. & is & bið \\
1–3 pl. & sindon & bēoð \\
\bottomrule
\end{tabular}
\par\endgroup
% course-panel:T06-P01:end
% mdtable:T06:end
% teaching: NEW-U06-GRAM-BEON-USE
% mdparagraph:P026:begin
The bēo series can describe an expected later state. The predicate still follows the agreement already learned. In the following example the ship’s readiness is expected later.
% mdparagraph:P026:end

% example: NEW-U06-EXAMPLE-09
% mdparagraph:P027:begin
Ðæt scip bið gearu.
% mdparagraph:P027:end

% mdparagraph:P028:begin
The ship will be ready, in an explicitly expected later situation.
% mdparagraph:P028:end

% mdparagraph:P029:begin
Bēo forms also occur with habitual or general truths. Context matters: the two present series are not an absolute division into “now” and “future,” and every bið need not be translated “will be.” The exercises here give an explicit expected-later context.
% mdparagraph:P029:end

% mdparagraph:P030:begin
\textbf{Sources.} Cook §138, pp.78–79 /\allowbreak{} PDF pp.94–95; Fulk §52, p.33 /\allowbreak{} PDF p.45.
% mdparagraph:P030:end

\reserve{7\baselineskip}
% mdparagraph:P031:begin
\subsection*{Practice A. Indicative systems and citation labels}\BookHeading{nav-ch06-learner-0006}{subsubsection}{Practice A. Indicative systems and citation labels}{0}
% mdparagraph:P031:end

\reserve{5\baselineskip}
\par\begin{minipage}{\linewidth}\raggedright
% mdparagraph:P032:begin
\textbf{A1.}
% mdparagraph:P032:end

% item: NEW-U06-A01
% mdparagraph:P033:begin
Give all six person-and-number cells of the present indicative and all six of the preterite indicative of witan. Use the selected course forms and label every cell even where spellings coincide.
% mdparagraph:P033:end

\end{minipage}\par\addvspace{4pt}

\par\begin{minipage}{\linewidth}\raggedright
% mdparagraph:P034:begin
\textbf{A2.}
% mdparagraph:P034:end

% item: NEW-U06-A02
% mdparagraph:P035:begin
Give all six person-and-number cells of the present indicative and all six of the preterite indicative of cunnan. Use the selected course forms and label every cell even where spellings coincide.
% mdparagraph:P035:end

\end{minipage}\par\addvspace{4pt}

\par\begin{minipage}{\linewidth}\raggedright
% mdparagraph:P036:begin
\textbf{A3.}
% mdparagraph:P036:end

% item: NEW-U06-A03
% mdparagraph:P037:begin
Give all six person-and-number cells of the present indicative and all six of the preterite indicative of ðurfan. Use the selected course forms and label every cell even where spellings coincide.
% mdparagraph:P037:end

\end{minipage}\par\addvspace{4pt}

\par\begin{minipage}{\linewidth}\raggedright
% mdparagraph:P038:begin
\textbf{A4.}
% mdparagraph:P038:end

% item: NEW-U06-A04
% mdparagraph:P039:begin
Give all six person-and-number cells of the present indicative and all six of the preterite indicative of durran. Use the selected course forms and label every cell even where spellings coincide.
% mdparagraph:P039:end

\end{minipage}\par\addvspace{4pt}

\par\begin{minipage}{\linewidth}\raggedright
% mdparagraph:P040:begin
\textbf{A5.}
% mdparagraph:P040:end

% item: NEW-U06-A05
% mdparagraph:P041:begin
Give all six person-and-number cells of the present indicative and all six of the preterite indicative of sculan. Use the selected course forms and label every cell even where spellings coincide.
% mdparagraph:P041:end

\end{minipage}\par\addvspace{4pt}

\par\begin{minipage}{\linewidth}\raggedright
% mdparagraph:P042:begin
\textbf{A6.}
% mdparagraph:P042:end

% item: NEW-U06-A06
% mdparagraph:P043:begin
Give all six person-and-number cells of the present indicative and all six of the preterite indicative of mugan. Use the selected course forms and label every cell even where spellings coincide.
% mdparagraph:P043:end

\end{minipage}\par\addvspace{4pt}

\par\begin{minipage}{\linewidth}\raggedright
% mdparagraph:P044:begin
\textbf{A7.}
% mdparagraph:P044:end

% item: NEW-U06-A07
% mdparagraph:P045:begin
Give all six person-and-number cells of the present indicative and all six of the preterite indicative of mōtan. Use the selected course forms and label every cell even where spellings coincide.
% mdparagraph:P045:end

\end{minipage}\par\addvspace{4pt}

\par\begin{minipage}{\linewidth}\raggedright
% mdparagraph:P046:begin
\textbf{A8.}
% mdparagraph:P046:end

% item: NEW-U06-A08
% mdparagraph:P047:begin
Give all six person-and-number cells of the present indicative and all six of the preterite indicative of dōn. Use the selected course forms and label every cell even where spellings coincide.
% mdparagraph:P047:end

\end{minipage}\par\addvspace{4pt}

\par\begin{minipage}{\linewidth}\raggedright
% mdparagraph:P048:begin
\textbf{A9.}
% mdparagraph:P048:end

% item: NEW-U06-A09
% mdparagraph:P049:begin
Give all six person-and-number cells of the present indicative and all six of the preterite indicative of gān. Use the selected course forms and label every cell even where spellings coincide.
% mdparagraph:P049:end

\end{minipage}\par\addvspace{4pt}

\par\begin{minipage}{\linewidth}\raggedright
% mdparagraph:P050:begin
\textbf{A10.}
% mdparagraph:P050:end

% item: NEW-U06-A10
% mdparagraph:P051:begin
Give all six person-and-number cells of the present indicative and all six of the preterite indicative of willan. Use the selected course forms and label every cell even where spellings coincide.
% mdparagraph:P051:end

\end{minipage}\par\addvspace{4pt}

\par\begin{minipage}{\linewidth}\raggedright
% mdparagraph:P052:begin
\textbf{A11.}
% mdparagraph:P052:end

% item: NEW-U06-A11
% mdparagraph:P053:begin
Give all six person-and-number cells of the bēo present indicative series. Then give the corresponding earlier eom series. Does bēo supply a separate past-tense paradigm?
% mdparagraph:P053:end

\end{minipage}\par\addvspace{4pt}

\par\begin{minipage}{\linewidth}\raggedright
% mdparagraph:P054:begin
\textbf{A12.}
% mdparagraph:P054:end

% item: NEW-U06-A12
% mdparagraph:P055:begin
Compare wāt and wiste; cann and cūðe; gǣð and ēode; dēð and dyde. Using a third-person singular subject, identify the tense of each. Which forms can also serve a first-person singular subject? Explain why wāt and cann are not translated as past-tense forms.
% mdparagraph:P055:end

\end{minipage}\par\addvspace{4pt}

\par\begin{minipage}{\linewidth}\raggedright
% mdparagraph:P056:begin
\textbf{A13.}
% mdparagraph:P056:end

% item: NEW-U06-A13
% mdparagraph:P057:begin
The reference notes mark some modal infinitive citations as unattested or uncertain. Explain why a dictionary label such as mugan does not authorize you to create an infinitive or a command for every modal verb. Distinguish a missing paradigm cell from two attested cells with the same spelling.
% mdparagraph:P057:end

\end{minipage}\par\addvspace{4pt}

\reserve{7\baselineskip}
% mdparagraph:P058:begin
\section*{3. Infinitives and their uses}\BookHeading{nav-ch06-learner-0007}{subsection}{3. Infinitives and their uses}{1}\BookLabel{bk-ch06-learner-530cdd5762f1}
% mdparagraph:P058:end

% teaching: NEW-U06-GRAM-INFINITIVE-FORMS
% teaching: NEW-U06-INF-MARKER
% mdparagraph:P059:begin
An \textbf{infinitive} names an action without marking person or number. The \textbf{plain infinitive} normally ends in -an or -ian, with contracted forms such as gān and dōn. In earlier units many infinitives were printed only to name verbs. The forms in this unit’s tables are now learned for the stated uses.
% mdparagraph:P059:end

% mdparagraph:P060:begin
The \textbf{inflected infinitive} follows the separate word \textbf{tō}. For the selected -an forms, replace -an with \textbf{-anne}; -ian consequently gives \textbf{-ianne}. Keep the stem consonants: doubled consonants remain doubled in the inflected infinitive. Contracted forms ending in -n add -ne, giving \textbf{-nne}: the complete forms appear in §4. The infinitive is not conjugated for the subject’s person or number.
% mdparagraph:P060:end

% mdparagraph:P061:begin
The tables select \textbf{-anne}. \textbf{-enne} is also a documented ending; it is a valid variant, although an exercise expressly requesting the selected -anne form still asks for that form. Infinitive formation does not supply missing modal forms or teach participles, perfects, or passives.
% mdparagraph:P061:end

\reserve{7\baselineskip}
% mdparagraph:P062:begin
\subsection*{A finite modal with a plain infinitive}\BookHeading{nav-ch06-learner-0008}{subsubsection}{A finite modal with a plain infinitive}{0}\BookLabel{bk-ch06-learner-d0a4c47b0324}
% mdparagraph:P062:end

% teaching: NEW-U06-GRAM-PLAIN-INFINITIVE-COMPLEMENT
% mdparagraph:P063:begin
In the constructions below, the finite verb agrees with the subject and marks present or past. The plain infinitive expresses the action performed by that same subject. It takes neither personal endings nor tō in this construction. Two or more people, including a dual-pronoun subject, take a plural finite verb; the infinitive itself stays unchanged.
% mdparagraph:P063:end

% teaching: NEW-U06-GRAM-MODAL-CUNNAN
% teaching: NEW-U06-GRAM-MODAL-MAGAN
% mdparagraph:P064:begin
\textbf{Cunnan} expresses knowing how through learning or skill. \textbf{Mugan} expresses capability: the subject is able to perform the action. English “can” can represent either, so select the intended sense from the context.
% mdparagraph:P064:end

% example: NEW-U06-EXAMPLE-01
% mdparagraph:P065:begin
Ic cann swimman. Hē mæg ðæt wæter beran.
% mdparagraph:P065:end

% mdparagraph:P066:begin
I know how to swim. He can carry the water.
% mdparagraph:P066:end

% mdparagraph:P067:begin
For the skill statements, learn these additional infinitive pairs. No imperative pair for either verb is required here.
% mdparagraph:P067:end

% teaching: NEW-U06-NONFINITE-singan
% teaching: NEW-U06-NONFINITE-swimman
% table: NEW-U06-DISPLAY-ADDITIONAL-INFINITIVES
% mdtable:T07:begin
\par\begingroup
\fontsize{10.5}{13.2}\selectfont
\setlength{\LTleft}{0pt}\setlength{\LTright}{\fill}
\noindent
\begin{tabular}{@{}>{\raggedright\arraybackslash}p{28.7mm}>{\raggedright\arraybackslash}p{36.45mm}>{\raggedright\arraybackslash}p{53.23mm}@{}}
\toprule
\textbf{Verb} & \textbf{Plain infinitive} & \textbf{Inflected infinitive after tō} \\
\midrule
singan & singan & singanne \\
swimman & swimman & swimmanne \\
\bottomrule
\end{tabular}
\par\endgroup

% mdtable:T07:end
% teaching: NEW-U06-GRAM-MODAL-THURFAN
% teaching: NEW-U06-GRAM-MODAL-DURRAN
% teaching: NEW-U06-GRAM-MODAL-SCULAN
% teaching: NEW-U06-GRAM-MODAL-MOTAN
% mdparagraph:P068:begin
\textbf{Ðurfan}\BookPoint{bk-ch06-learner-1041b72208e0} expresses need to do something; \textbf{durran}, daring to do it. \textbf{Sculan} expresses obligation or necessity in the selected use. \textbf{Mōtan} expresses permission, “be allowed to.” A past form can describe actual past need, daring, obligation, or permission: the morphology alone does not make the statement hypothetical.
% mdparagraph:P068:end

% mdparagraph:P069:begin
In\BookPoint{bk-ch06-learner-95f9e5a43d91} particular, selected \textbf{mōste} can mean “was allowed to”; it need not mean modern English “must.” Other senses occur in the language, but the exercises state the sense required here.
% mdparagraph:P069:end

% example: NEW-U06-EXAMPLE-02
% mdparagraph:P070:begin
Wē sculon ðæt weorc gefręmman.
% mdparagraph:P070:end

% mdparagraph:P071:begin
We must perform the work.
% mdparagraph:P071:end

% teaching: NEW-U06-GRAM-VOLITION-WILLAN
% teaching: NEW-U06-GRAM-FARAN-INTRANSITIVE
% mdparagraph:P072:begin
\textbf{Willan}\BookPoint{bk-ch06-learner-f11fea9bc36d} expresses wishing or wanting. Its preterite \textbf{wolde} can state past desire. \textbf{Faran} is used intransitively for going or travelling; it can name the desired action without an expressed object or destination.
% mdparagraph:P072:end

% example: NEW-U06-EXAMPLE-03
% mdparagraph:P073:begin
Hē wolde faran.
% mdparagraph:P073:end

% mdparagraph:P074:begin
He wanted to travel.
% mdparagraph:P074:end

% mdparagraph:P075:begin
Wanting\BookPoint{bk-ch06-learner-109abaf5847e} to travel does not by itself assert that the journey occurred. Do not translate every wolde as hypothetical “would,” or treat willan as an automatic future-tense auxiliary.
% mdparagraph:P075:end

% teaching: NEW-U06-GRAM-MODAL-NEGATION
% mdparagraph:P076:begin
The familiar negative \textbf{ne} precedes the finite verb. Use this ordinary pattern with the selected durran and mugan clauses. Contracted negative paradigms and negative concord are not needed.
% mdparagraph:P076:end

\reserve{7\baselineskip}
% mdparagraph:P077:begin
\subsection*{Witan with a noun object}\BookHeading{nav-ch06-learner-0009}{subsubsection}{Witan with a noun object}{0}\BookLabel{bk-ch06-learner-9fbb6a52092f}
% mdparagraph:P077:end

% teaching: NEW-U06-GRAM-WITAN-OBJECT
% mdparagraph:P078:begin
\textbf{Witan}, “know,” takes an accusative nominal object for the matter known. The following “things” are facts or matters already discussed. The noun’s form is familiar; its fact/\allowbreak{}matter meaning is new here.
% mdparagraph:P078:end

% example: NEW-U06-EXAMPLE-10
% mdparagraph:P079:begin
Ic wāt ðā ðing.
% mdparagraph:P079:end

% mdparagraph:P080:begin
I know those things, referring to matters already discussed.
% mdparagraph:P080:end

% mdparagraph:P081:begin
This is knowledge of something, whereas cunnan with an infinitive describes knowing how to do something. The object here is a noun phrase; no finite “that” clause is required.
% mdparagraph:P081:end

\reserve{7\baselineskip}
% mdparagraph:P082:begin
\subsection*{Ordinary action and purpose}\BookHeading{nav-ch06-learner-0010}{subsubsection}{Ordinary action and purpose}{0}\BookLabel{bk-ch06-learner-c3a1991273ab}
% mdparagraph:P082:end

% teaching: NEW-U06-GRAM-DON-OBJECT
% teaching: NEW-U06-GRAM-GAN-INTRANSITIVE
% teaching: NEW-U06-GRAM-CUMAN-INTRANSITIVE
% mdparagraph:P083:begin
Use \textbf{dōn}, “do; perform,” with an accusative work or thing done. Use \textbf{gān}, “go,” and \textbf{cuman}, “come,” in ordinary intransitive clauses. These uses do not require a destination phrase and do not establish a general rule for motion prepositions.
% mdparagraph:P083:end

% teaching: NEW-U06-GRAM-PURPOSE-INFINITIVE
% mdparagraph:P084:begin
After \textbf{cuman}, \textbf{tō + inflected infinitive} can state the purpose of coming. The person or people who come are also the understood performers of the infinitive’s action.
% mdparagraph:P084:end

% example: NEW-U06-EXAMPLE-07
% mdparagraph:P085:begin
Hē cōm tō helpanne.
% mdparagraph:P085:end

% mdparagraph:P086:begin
He came to help.
% mdparagraph:P086:end

% mdparagraph:P087:begin
Here he came, and he intended to help. The infinitive names the purpose; it does not assert that the helping was completed. This is the purpose construction practised below. Old English also permits certain bare infinitives after motion verbs, so it would be wrong to say that every grammatical purpose must have tō. When an exercise requests the inflected construction, provide that target.
% mdparagraph:P087:end

% teaching: NEW-U06-GRAM-INFINITIVE-GOVERNMENT
% mdparagraph:P088:begin
An infinitive keeps the object relations of its verb. Carrying, taking, and eating can have accusative objects; \textbf{helpan} takes the person helped in the dative. Giving has an accusative thing and a dative recipient. Adding a modal or changing to an inflected infinitive does not turn those datives into accusatives.
% mdparagraph:P088:end

% mdparagraph:P089:begin
\textbf{Sources.} Cook §§199–200, pp.105–106 /\allowbreak{} PDF pp.121–122; Fulk §58, p.38 /\allowbreak{} PDF p.50. Cook §200 provides cuman with an inflected purpose infinitive; §199.1 also records bare infinitives with motion and purpose. The modal distinctions follow \textbf{Bosworth–Toller, An Anglo-Saxon Dictionary}, main dictionary (1898): CUNNAN II, p.174b; MAGAN III(1), p.665a; ÞURFAN II(1), p.1076b; DURRAN, p.219b; SCULAN II(1), pp.843b–844; MÓTAN I, p.699a–b; WILLAN II(a), p.1226b. Ordinary DÓN is at p.208b, GÁN at p.360b, and FARAN I at p.270b. Witan’s nominal object is evidenced in Cook Reader XII, p.173 /\allowbreak{} PDF p.189; the known-fact sense of þing appears under WITAN I(6)(a), Bosworth–Toller p.1244b. Cuman is glossed in Cook’s vocabulary, p.282 /\allowbreak{} PDF p.298.
% mdparagraph:P089:end

\reserve{7\baselineskip}
% mdparagraph:P090:begin
\section*{4. Commands and complete model tables}\BookHeading{nav-ch06-learner-0011}{subsection}{4. Commands and complete model tables}{1}\BookLabel{bk-ch06-learner-a6c353738baa}
% mdparagraph:P090:end

% teaching: NEW-U06-GRAM-IMPERATIVE-USE
% mdparagraph:P091:begin
The \textbf{imperative} directly addresses the person who is to act. Use the second singular for one addressee and the second plural for two or more. There is no separate dual imperative ending. The subject may be omitted; an expressed subject can precede the verb. The commands used here omit it unless the exercise states otherwise.
% mdparagraph:P091:end

% example: NEW-U06-EXAMPLE-04
% mdparagraph:P092:begin
Nim ðone męte. Nimað ðone męte.
% mdparagraph:P092:end

% mdparagraph:P093:begin
Take the food, one addressee. Take the food, several addressees.
% mdparagraph:P093:end

% mdparagraph:P094:begin
The number of the command identifies the addressees, not the number of its object. A command concerning the same food may address one person or several.
% mdparagraph:P094:end

% mdparagraph:P095:begin
The following tables give each complete command pair beside its complete plain/\allowbreak{}inflected infinitive pair. \textbf{The last column is used after separate tō.} Learn every form printed here. Most plural commands match the present indicative plural; that shared spelling does not make their uses identical.
% mdparagraph:P095:end

\reserve{7\baselineskip}
% mdparagraph:P096:begin
\subsection*{Weak I}\BookHeading{nav-ch06-learner-0012}{subsubsection}{Weak I}{0}\BookLabel{bk-ch06-learner-ba352ea21d51}
% mdparagraph:P096:end

% teaching: NEW-U06-GRAM-WEAK-IMPERATIVE-FORMS
% mdparagraph:P097:begin
The heavy-stem model has no final -e in the singular command: \textbf{hīer}. The short and geminate patterns have \textbf{-e}, with the singular consonant pattern shown in \textbf{sęle, sęte, gefręme, nęre}. In nęre the i of nęrian is absent; it remains in the plural nęriað. \textbf{Gierwan} has \textbf{giere}, without w, beside gierwað.
% mdparagraph:P097:end

% mdparagraph:P098:begin
Learn \textbf{cys} and \textbf{gefyl} with a single final consonant, beside cyssað and gefyllað. Their selected final simplification does not turn them into the -e type. Conversely, sęttan has sęte; its double consonant is not a reason to use the cys pattern. Final consonant simplification is frequent, not a claim that all unsimplified final spellings are impossible.
% mdparagraph:P098:end

% teaching: NEW-U06-IMP-sellan
% teaching: NEW-U06-NONFINITE-sellan
% teaching: NEW-U06-IMP-hieran
% teaching: NEW-U06-NONFINITE-hieran
% teaching: NEW-U06-IMP-gefyllan
% teaching: NEW-U06-NONFINITE-gefyllan
% teaching: NEW-U06-IMP-cyssan
% teaching: NEW-U06-NONFINITE-cyssan
% teaching: NEW-U06-IMP-settan
% teaching: NEW-U06-NONFINITE-settan
% teaching: NEW-U06-IMP-gierwan
% teaching: NEW-U06-NONFINITE-gierwan
% teaching: NEW-U06-IMP-gefremman
% teaching: NEW-U06-NONFINITE-gefremman
% teaching: NEW-U06-IMP-nerian
% teaching: NEW-U06-NONFINITE-nerian
% table: NEW-U06-DISPLAY-WEAK-I
% mdtable:T08:begin
\par\begingroup
\fontsize{10.5}{13.2}\selectfont
\setlength{\LTleft}{0pt}\setlength{\LTright}{\fill}
\noindent
\begin{tabular}{@{}>{\raggedright\arraybackslash}p{18.5mm}>{\raggedright\arraybackslash}p{17.13mm}>{\raggedright\arraybackslash}p{25.07mm}>{\raggedright\arraybackslash}p{21.44mm}>{\raggedright\arraybackslash}p{30.62mm}@{}}
\toprule
\textbf{Verb} & \textbf{Command to one} & \textbf{Command to two or more} & \textbf{Plain infinitive} & \textbf{Inflected infinitive after tō} \\
\midrule
sęllan & sęle & sęllað & sęllan & sęllanne \\
hīeran & hīer & hīerað & hīeran & hīeranne \\
gefyllan & gefyl & gefyllað & gefyllan & gefyllanne \\
cyssan & cys & cyssað & cyssan & cyssanne \\
sęttan & sęte & sęttað & sęttan & sęttanne \\
gierwan & giere & gierwað & gierwan & gierwanne \\
gefręmman & gefręme & gefręmmað & gefręmman & gefręmmanne \\
nęrian & nęre & nęriað & nęrian & nęrianne \\
\bottomrule
\end{tabular}
\par\endgroup

% mdtable:T08:end
\reserve{7\baselineskip}
% mdparagraph:P099:begin
\subsection*{Weak II and III}\BookHeading{nav-ch06-learner-0013}{subsubsection}{Weak II and III}{0}\BookLabel{bk-ch06-learner-8eef50183c9c}
% mdparagraph:P099:end

% mdparagraph:P100:begin
Weak II changes \textbf{-ian} to \textbf{-a} for the singular command and retains \textbf{-iað} in the plural: lufa/\allowbreak{}lufiað. Its inflected infinitive has -ianne.
% mdparagraph:P100:end

% mdparagraph:P101:begin
The four learned weak-III verbs require their own singular forms: \textbf{hafa, leofa, saga, hoga}. Their plural commands use the familiar present-plural stems. Form habbanne from habban; the command stem hafa is different.
% mdparagraph:P101:end

% teaching: NEW-U06-IMP-lufian
% teaching: NEW-U06-NONFINITE-lufian
% teaching: NEW-U06-IMP-habban
% teaching: NEW-U06-NONFINITE-habban
% teaching: NEW-U06-IMP-libban
% teaching: NEW-U06-NONFINITE-libban
% teaching: NEW-U06-IMP-secgan
% teaching: NEW-U06-NONFINITE-secgan
% teaching: NEW-U06-IMP-hycgan
% teaching: NEW-U06-NONFINITE-hycgan
% table: NEW-U06-DISPLAY-WEAK-II-III
% mdtable:T09:begin
\par\begingroup
\fontsize{10.5}{13.2}\selectfont
\setlength{\LTleft}{0pt}\setlength{\LTright}{\fill}
\noindent
\begin{tabular}{@{}>{\raggedright\arraybackslash}p{17.31mm}>{\raggedright\arraybackslash}p{17.57mm}>{\raggedright\arraybackslash}p{23.8mm}>{\raggedright\arraybackslash}p{21.98mm}>{\raggedright\arraybackslash}p{32.1mm}@{}}
\toprule
\textbf{Verb} & \textbf{Command to one} & \textbf{Command to two or more} & \textbf{Plain infinitive} & \textbf{Inflected infinitive after tō} \\
\midrule
lufian & lufa & lufiað & lufian & lufianne \\
habban & hafa & habbað & habban & habbanne \\
libban & leofa & libbað & libban & libbanne \\
sęcgan & saga & sęcgað & sęcgan & sęcganne \\
hycgan & hoga & hycgað & hycgan & hycganne \\
\bottomrule
\end{tabular}
\par\endgroup

% mdtable:T09:end
% mdparagraph:P102:begin
Cook also gives \textbf{sęge} beside saga and \textbf{lifiað} beside libbað, and optional e in \textbf{sęcg(e)að} and \textbf{hycg(e)að}. The tables select saga, libbað, sęcgað, and hycgað. The earlier recognition-only negative companion nabban receives no new command requirement.
% mdparagraph:P102:end

% mdparagraph:P103:begin
\textbf{Sources for the weak forms.} Cook §§112–124, pp.68–76 /\allowbreak{} PDF pp.84–92. Final-geminate simplification is described in §35, pp.24–25 /\allowbreak{} PDF pp.40–41; cys and gefyl apply that rule to the heavy weak-I pattern, rather than being separate continuous-text quotations. Wright §§521–525 and 529–530, pp.266–268 and 270–271 /\allowbreak{} PDF pp.286–288 and 290–291, corroborates the short, geminate, and gierwan patterns.
% mdparagraph:P103:end

\reserve{7\baselineskip}
% mdparagraph:P104:begin
\subsection*{Strong I–III}\BookHeading{nav-ch06-learner-0014}{subsubsection}{Strong I–III}{0}\BookLabel{bk-ch06-learner-634c423f6df0}
% mdparagraph:P104:end

% teaching: NEW-U06-GRAM-STRONG-IMPERATIVE-FORMS
% mdparagraph:P105:begin
The ordinary singular command uses the present stem without the infinitive ending. Do not import the changed vowel of the indicative second or third singular into it. The ordinary plural command matches the present indicative plural.
% mdparagraph:P105:end

% mdparagraph:P106:begin
Contracted verbs need a separate singular command in \textbf{-h}. Class I wrēon gives \textbf{wrīh}, with long ī; Class II tēon gives \textbf{tēoh}, with long ēo. Their plurals and infinitives stay contracted. In the table, drīfan and wrēon belong to Class I; bēodan and tēon to II; bindan, helpan, and drincan to III.
% mdparagraph:P106:end

% teaching: NEW-U06-IMP-drifan
% teaching: NEW-U06-NONFINITE-drifan
% teaching: NEW-U06-IMP-wreon
% teaching: NEW-U06-NONFINITE-wreon
% teaching: NEW-U06-IMP-beodan
% teaching: NEW-U06-NONFINITE-beodan
% teaching: NEW-U06-IMP-teon
% teaching: NEW-U06-NONFINITE-teon
% teaching: NEW-U06-IMP-bindan
% teaching: NEW-U06-NONFINITE-bindan
% teaching: NEW-U06-IMP-helpan
% teaching: NEW-U06-NONFINITE-helpan
% teaching: NEW-U06-IMP-drincan
% teaching: NEW-U06-NONFINITE-drincan
% table: NEW-U06-DISPLAY-STRONG-I-III
% mdtable:T10:begin
\par\begingroup
\fontsize{10.5}{13.2}\selectfont
\setlength{\LTleft}{0pt}\setlength{\LTright}{\fill}
\noindent
\begin{tabular}{@{}>{\raggedright\arraybackslash}p{17.31mm}>{\raggedright\arraybackslash}p{17.57mm}>{\raggedright\arraybackslash}p{23.8mm}>{\raggedright\arraybackslash}p{21.98mm}>{\raggedright\arraybackslash}p{32.1mm}@{}}
\toprule
\textbf{Verb} & \textbf{Command to one} & \textbf{Command to two or more} & \textbf{Plain infinitive} & \textbf{Inflected infinitive after tō} \\
\midrule
drīfan & drīf & drīfað & drīfan & drīfanne \\
wrēon & wrīh & wrēoð & wrēon & wrēonne \\
bēodan & bēod & bēodað & bēodan & bēodanne \\
tēon & tēoh & tēoð & tēon & tēonne \\
bindan & bind & bindað & bindan & bindanne \\
helpan & help & helpað & helpan & helpanne \\
drincan & drinc & drincað & drincan & drincanne \\
\bottomrule
\end{tabular}
\par\endgroup

% mdtable:T10:end
\reserve{7\baselineskip}
% mdparagraph:P107:begin
\subsection*{Strong IV–V}\BookHeading{nav-ch06-learner-0015}{subsubsection}{Strong IV–V}{0}\BookLabel{bk-ch06-learner-eb4f99a5a301}
% mdparagraph:P107:end

% mdparagraph:P108:begin
Beran, niman, and cuman are Class IV. The remaining verbs in this table are Class V. The ordinary singular \textbf{cum} retains u, unlike the changed indicative singular vowel.
% mdparagraph:P108:end

% mdparagraph:P109:begin
The strong j-present verbs have singular commands in \textbf{-e}, like the short weak-I pattern: \textbf{bide, lige, site}. Gemination is absent in these singulars. Licgan has g in lige, beside cg in licgað. Contracted sēon has \textbf{seoh}, with short eo, beside long ēo in sēoð.
% mdparagraph:P109:end

% teaching: NEW-U06-IMP-cuman
% teaching: NEW-U06-NONFINITE-cuman
% teaching: NEW-U06-IMP-beran
% teaching: NEW-U06-NONFINITE-beran
% teaching: NEW-U06-IMP-niman
% teaching: NEW-U06-NONFINITE-niman
% teaching: NEW-U06-IMP-giefan
% teaching: NEW-U06-NONFINITE-giefan
% teaching: NEW-U06-IMP-biddan
% teaching: NEW-U06-NONFINITE-biddan
% teaching: NEW-U06-IMP-licgan
% teaching: NEW-U06-NONFINITE-licgan
% teaching: NEW-U06-IMP-sittan
% teaching: NEW-U06-NONFINITE-sittan
% teaching: NEW-U06-IMP-etan
% teaching: NEW-U06-NONFINITE-etan
% teaching: NEW-U06-IMP-seon
% teaching: NEW-U06-NONFINITE-seon
% table: NEW-U06-DISPLAY-STRONG-IV-V
% mdtable:T11:begin
\par\begingroup
\fontsize{10.5}{13.2}\selectfont
\setlength{\LTleft}{0pt}\setlength{\LTright}{\fill}
\noindent
\begin{tabular}{@{}>{\raggedright\arraybackslash}p{17.31mm}>{\raggedright\arraybackslash}p{17.57mm}>{\raggedright\arraybackslash}p{23.8mm}>{\raggedright\arraybackslash}p{21.98mm}>{\raggedright\arraybackslash}p{32.1mm}@{}}
\toprule
\textbf{Verb} & \textbf{Command to one} & \textbf{Command to two or more} & \textbf{Plain infinitive} & \textbf{Inflected infinitive after tō} \\
\midrule
cuman & cum & cumað & cuman & cumanne \\
beran & ber & berað & beran & beranne \\
niman & nim & nimað & niman & nimanne \\
giefan & gief & giefað & giefan & giefanne \\
biddan & bide & biddað & biddan & biddanne \\
licgan & lige & licgað & licgan & licganne \\
sittan & site & sittað & sittan & sittanne \\
etan & et & etað & etan & etanne \\
sēon & seoh & sēoð & sēon & sēonne \\
\bottomrule
\end{tabular}
\par\endgroup

% mdtable:T11:end
\reserve{7\baselineskip}
% mdparagraph:P110:begin
\subsection*{Strong VI–VII}\BookHeading{nav-ch06-learner-0016}{subsubsection}{Strong VI–VII}{0}\BookLabel{bk-ch06-learner-75a0f72a393d}
% mdparagraph:P110:end

% mdparagraph:P111:begin
Faran, standan, slēan, hębban, scieppan, and swęrian are Class VI. Hātan, healdan, weaxan, lǣtan, and fōn are Class VII.
% mdparagraph:P111:end

% mdparagraph:P112:begin
The j-present singulars are \textbf{hęfe, sciepe, swęre}. Hębban changes bb to f in the singular command. Swęrian has swęre beside swęriað. Contracted slēan gives \textbf{sleah}, with short ea; fōn gives \textbf{fōh}, with long ō. Keep these quantities distinct from those in the corresponding infinitives.
% mdparagraph:P112:end

% teaching: NEW-U06-IMP-hatan
% teaching: NEW-U06-NONFINITE-hatan
% teaching: NEW-U06-IMP-standan
% teaching: NEW-U06-NONFINITE-standan
% teaching: NEW-U06-IMP-faran
% teaching: NEW-U06-NONFINITE-faran
% teaching: NEW-U06-IMP-slean
% teaching: NEW-U06-NONFINITE-slean
% teaching: NEW-U06-IMP-hebban
% teaching: NEW-U06-NONFINITE-hebban
% teaching: NEW-U06-IMP-scieppan
% teaching: NEW-U06-NONFINITE-scieppan
% teaching: NEW-U06-IMP-swerian
% teaching: NEW-U06-NONFINITE-swerian
% teaching: NEW-U06-IMP-healdan
% teaching: NEW-U06-NONFINITE-healdan
% teaching: NEW-U06-IMP-weaxan
% teaching: NEW-U06-NONFINITE-weaxan
% teaching: NEW-U06-IMP-laetan
% teaching: NEW-U06-NONFINITE-laetan
% teaching: NEW-U06-IMP-fon
% teaching: NEW-U06-NONFINITE-fon
% table: NEW-U06-DISPLAY-STRONG-VI-VII
% mdtable:T12:begin
\par\begingroup
\fontsize{10.5}{13.2}\selectfont
\setlength{\LTleft}{0pt}\setlength{\LTright}{\fill}
\noindent
\begin{tabular}{@{}>{\raggedright\arraybackslash}p{17.16mm}>{\raggedright\arraybackslash}p{17.42mm}>{\raggedright\arraybackslash}p{24.54mm}>{\raggedright\arraybackslash}p{21.8mm}>{\raggedright\arraybackslash}p{31.83mm}@{}}
\toprule
\textbf{Verb} & \textbf{Command to one} & \textbf{Command to two or more} & \textbf{Plain infinitive} & \textbf{Inflected infinitive after tō} \\
\midrule
hātan & hāt & hātað & hātan & hātanne \\
standan & stand & standað & standan & standanne \\
faran & far & farað & faran & faranne \\
slēan & sleah & slēað & slēan & slēanne \\
hębban & hęfe & hębbað & hębban & hębbanne \\
scieppan & sciepe & scieppað & scieppan & scieppanne \\
swęrian & swęre & swęriað & swęrian & swęrianne \\
healdan & heald & healdað & healdan & healdanne \\
weaxan & weax & weaxað & weaxan & weaxanne \\
lǣtan & lǣt & lǣtað & lǣtan & lǣtanne \\
fōn & fōh & fōð & fōn & fōnne \\
\bottomrule
\end{tabular}
\par\endgroup

% mdtable:T12:end
% mdparagraph:P113:begin
\textbf{Sources for the strong forms.} Cook §§93–95, pp.54–55 /\allowbreak{} PDF pp.70–71; §§100–101, pp.57–59 /\allowbreak{} PDF pp.73–75; §§106–107, pp.64–65 /\allowbreak{} PDF pp.80–81. Wright §§478 and 480, pp.250–251 /\allowbreak{} PDF pp.270–271, corroborates imperative endings and inflected infinitives. Fulk §127, pp.93–94 /\allowbreak{} PDF pp.105–106, corroborates the j-present singulars and their consonant alternations. Wrīh applies Cook’s stated Class I contracted pattern to wrēon; the model-derived form is not presented as a separate prose quotation.
% mdparagraph:P113:end

\reserve{7\baselineskip}
% mdparagraph:P114:begin
\subsection*{Dōn, gān, wesan, and bēon}\BookHeading{nav-ch06-learner-0017}{subsubsection}{Dōn, gān, wesan, and bēon}{0}\BookLabel{bk-ch06-learner-1a5c7d53adeb}
% mdparagraph:P114:end

% teaching: NEW-U06-GRAM-ANOMALOUS-IMPERATIVE-FORMS
% mdparagraph:P115:begin
Learn \textbf{dō/\allowbreak{}dōð}, \textbf{gā/\allowbreak{}gāð}, \textbf{wes/\allowbreak{}wesað}, and \textbf{bēo/\allowbreak{}bēoð} as complete pairs. Their contracted infinitives keep the contracted stem before -nne. Wesan has the ordinary -anne form.
% mdparagraph:P115:end

% mdparagraph:P116:begin
These second-person commands are distinct from an indicative or a supplied third-person wish that may share a spelling. Neither willan nor the preterite-present verbs receives a regular command pair by analogy with them.
% mdparagraph:P116:end

% teaching: NEW-U06-IMP-wesan
% teaching: NEW-U06-NONFINITE-wesan
% teaching: NEW-U06-IMP-don
% teaching: NEW-U06-NONFINITE-don
% teaching: NEW-U06-IMP-gan
% teaching: NEW-U06-NONFINITE-gan
% teaching: NEW-U06-IMP-wesan-beon
% teaching: NEW-U06-NONFINITE-wesan-beon
% table: NEW-U06-DISPLAY-ANOMALOUS
% mdtable:T13:begin
\par\begingroup
\fontsize{10.5}{13.2}\selectfont
\setlength{\LTleft}{0pt}\setlength{\LTright}{\fill}
\noindent
\begin{tabular}{@{}>{\raggedright\arraybackslash}p{17.31mm}>{\raggedright\arraybackslash}p{17.57mm}>{\raggedright\arraybackslash}p{23.8mm}>{\raggedright\arraybackslash}p{21.98mm}>{\raggedright\arraybackslash}p{32.1mm}@{}}
\toprule
\textbf{Verb} & \textbf{Command to one} & \textbf{Command to two or more} & \textbf{Plain infinitive} & \textbf{Inflected infinitive after tō} \\
\midrule
wesan & wes & wesað & wesan & wesanne \\
dōn & dō & dōð & dōn & dōnne \\
gān & gā & gāð & gān & gānne \\
bēon & bēo & bēoð & bēon & bēonne \\
\bottomrule
\end{tabular}
\par\endgroup

% mdtable:T13:end
% example: NEW-U06-EXAMPLE-08
% mdparagraph:P117:begin
Gā. Dō ðæt weorc wel.
% mdparagraph:P117:end

% mdparagraph:P118:begin
Go. Do the work well, one addressee.
% mdparagraph:P118:end

% mdparagraph:P119:begin
Some verbs remain form models for uses not yet learned. A full command or infinitive row does not by itself teach a new object, a causative complement, or a naming construction. Keep the earlier limits on such models as gierwan and hycgan, and on the complement uses of hātan, lǣtan, bēodan, and biddan.
% mdparagraph:P119:end

% mdparagraph:P120:begin
\textbf{Reference variant.}\BookPoint{bk-ch06-learner-6a76a82f4e05} Cook records a shortened plural \textbf{-e} immediately before a following subject pronoun. Use the full plural forms above when the subject is omitted or precedes the verb. The positional variant is supplied for reference here, not an additional free choice of ending or a first-/\allowbreak{}third-person command system.
% mdparagraph:P120:end

% mdparagraph:P121:begin
\textbf{Sources.} Cook §§138–141, pp.78–80 /\allowbreak{} PDF pp.94–96; Fulk §134, p.102 /\allowbreak{} PDF p.114. The postposed-pronoun variant is Cook §95, p.55 /\allowbreak{} PDF p.71. Inflected -anne/\allowbreak{}-enne is also described by Wright §480, p.251 /\allowbreak{} PDF p.271, and Fulk §58, p.38 /\allowbreak{} PDF p.50.
% mdparagraph:P121:end

\reserve{7\baselineskip}
% mdparagraph:P122:begin
\subsection*{Objects, negative commands, and quoted speech}\BookHeading{nav-ch06-learner-0018}{subsubsection}{Objects, negative commands, and quoted speech}{0}\BookLabel{bk-ch06-learner-c13dd1a2bbd6}
% mdparagraph:P122:end

% mdparagraph:P123:begin
Commands retain the verb’s familiar government. Giving still has a dative recipient and an accusative thing; helpan still takes the dative person helped.
% mdparagraph:P123:end

% example: NEW-U06-EXAMPLE-05
% mdparagraph:P124:begin
Gief mē ðæt wæter. Helpað mē.
% mdparagraph:P124:end

% mdparagraph:P125:begin
Give me the water, one addressee. Help me, several addressees.
% mdparagraph:P125:end

% teaching: NEW-U06-GRAM-IMPERATIVE-NEGATIVE
% mdparagraph:P126:begin
Put \textbf{ne} directly before the imperative for the simple negative command practised here.
% mdparagraph:P126:end

% example: NEW-U06-EXAMPLE-06
% mdparagraph:P127:begin
Ne drinc ðæt wæter.
% mdparagraph:P127:end

% mdparagraph:P128:begin
Do not drink the water, one addressee.
% mdparagraph:P128:end

% teaching: NEW-U06-GRAM-DIRECT-QUOTATION
% mdparagraph:P129:begin
A quotation introduced by \textbf{cweðan} may now contain these direct commands. Within the quotation, “I,” “me,” and “you” belong to the speaker’s and addressee’s situation. The quotation is not an indirect command and does not require a new dependent mood.
% mdparagraph:P129:end

% mdparagraph:P130:begin
\textbf{Sources.} Cook §198, p.105 /\allowbreak{} PDF p.121, gives commands with or without an expressed second-person subject. The negative imperative is directly evidenced in Reader XII, p.183 /\allowbreak{} PDF p.199, editorial line 7, and quoted in §§156h and 159a, pp.92–93 /\allowbreak{} PDF pp.108–109. These citations support the construction; their additional vocabulary is not needed for the exercises.
% mdparagraph:P130:end

\reserve{7\baselineskip}
% mdparagraph:P131:begin
\section*{5. Practice in clauses and forms}\BookHeading{nav-ch06-learner-0019}{subsection}{5. Practice in clauses and forms}{1}
% mdparagraph:P131:end

% mdparagraph:P132:begin
All exercises are new teaching material. Unless an exercise asks for a supplied reference, use the forms and constructions learned in this and earlier units. Preserve the stated modal sense. When a particular course form or construction is requested, distinguish that target from another valid expression.
% mdparagraph:P132:end

\reserve{7\baselineskip}
% mdparagraph:P133:begin
\subsection*{Practice M. Meanings and infinitive complements}\BookHeading{nav-ch06-learner-0020}{subsubsection}{Practice M. Meanings and infinitive complements}{0}
% mdparagraph:P133:end

\reserve{5\baselineskip}
\par\begin{minipage}{\linewidth}\raggedright
% mdparagraph:P134:begin
\textbf{M1.}
% mdparagraph:P134:end

% item: NEW-U06-M01
% mdparagraph:P135:begin
Translate “I know how to swim. We know how to sing.” Use cunnan for learned ability.
% mdparagraph:P135:end

\end{minipage}\par\addvspace{4pt}

\par\begin{minipage}{\linewidth}\raggedright
% mdparagraph:P136:begin
\textbf{M2.}
% mdparagraph:P136:end

% item: NEW-U06-M02
% mdparagraph:P137:begin
Translate “He can carry the water. They can help the man.” Use mugan for ability.
% mdparagraph:P137:end

\end{minipage}\par\addvspace{4pt}

\par\begin{minipage}{\linewidth}\raggedright
% mdparagraph:P138:begin
\textbf{M3.}
% mdparagraph:P138:end

% item: NEW-U06-M03
% mdparagraph:P139:begin
Translate “I need to go. We needed to go.” Use ðurfan.
% mdparagraph:P139:end

\end{minipage}\par\addvspace{4pt}

\par\begin{minipage}{\linewidth}\raggedright
% mdparagraph:P140:begin
\textbf{M4.}
% mdparagraph:P140:end

% item: NEW-U06-M04
% mdparagraph:P141:begin
Translate “You dare to go. You dared to go.” Address one person and use durran.
% mdparagraph:P141:end

\end{minipage}\par\addvspace{4pt}

\par\begin{minipage}{\linewidth}\raggedright
% mdparagraph:P142:begin
\textbf{M5.}
% mdparagraph:P142:end

% item: NEW-U06-M05
% mdparagraph:P143:begin
Translate “We must do the work. They had to do the work.” Use sculan for obligation.
% mdparagraph:P143:end

\end{minipage}\par\addvspace{4pt}

\par\begin{minipage}{\linewidth}\raggedright
% mdparagraph:P144:begin
\textbf{M6.}
% mdparagraph:P144:end

% item: NEW-U06-M06
% mdparagraph:P145:begin
Translate “You are allowed to eat the food. You were allowed to eat the food.” Address one person and use mōtan.
% mdparagraph:P145:end

\end{minipage}\par\addvspace{4pt}

\par\begin{minipage}{\linewidth}\raggedright
% mdparagraph:P146:begin
\textbf{M7.}
% mdparagraph:P146:end

% item: NEW-U06-M07
% mdparagraph:P147:begin
Translate “He wants to give the birds water. They wanted to give the birds water.” Use willan for desire.
% mdparagraph:P147:end

\end{minipage}\par\addvspace{4pt}

\par\begin{minipage}{\linewidth}\raggedright
% mdparagraph:P148:begin
\textbf{M8.}
% mdparagraph:P148:end

% item: NEW-U06-M08
% mdparagraph:P149:begin
Translate “I did not dare to go. We cannot carry the ship.” Use the taught simple ne construction.
% mdparagraph:P149:end

\end{minipage}\par\addvspace{4pt}

\par\begin{minipage}{\linewidth}\raggedright
% mdparagraph:P150:begin
\textbf{M9.}
% mdparagraph:P150:end

% item: NEW-U06-M09
% mdparagraph:P151:begin
Translate “We two know how to swim. You two must help us two.” Keep the two-person references.
% mdparagraph:P151:end

\end{minipage}\par\addvspace{4pt}

\par\begin{minipage}{\linewidth}\raggedright
% mdparagraph:P152:begin
\textbf{M10.}
% mdparagraph:P152:end

% item: NEW-U06-M10
% mdparagraph:P153:begin
Translate “He wanted to go. He went.” Explain why the first sentence alone does not assert that he went.
% mdparagraph:P153:end

\end{minipage}\par\addvspace{4pt}

\par\begin{minipage}{\linewidth}\raggedright
% mdparagraph:P154:begin
\textbf{M11.}
% mdparagraph:P154:end

% item: NEW-U06-M11
% mdparagraph:P155:begin
Translate the printed pairs Ic cann swimman and Ic mæg swimman using the unit’s stated difference between learned ability and present capability. Identify the finite verb and the infinitive in each. Do not claim that one English word always represents only one Old English modal.
% mdparagraph:P155:end

\end{minipage}\par\addvspace{4pt}

\par\begin{minipage}{\linewidth}\raggedright
% mdparagraph:P156:begin
\textbf{M12.}
% mdparagraph:P156:end

% item: NEW-U06-M12
% mdparagraph:P157:begin
Translate “The ship will be ready.” Use the bēo-series indicative to describe the expected later state. Then contrast the earlier eom-series sentence “The ship is ready.” Do not infer that every bēo-series occurrence must be future.
% mdparagraph:P157:end

\end{minipage}\par\addvspace{4pt}

\par\begin{minipage}{\linewidth}\raggedright
% mdparagraph:P158:begin
\textbf{M13.}
% mdparagraph:P158:end

% item: NEW-U06-M13
% mdparagraph:P159:begin
The “things” here are matters already discussed. Translate “I know those things. He knew those things.” Use witan. Explain how this knowledge construction differs from cunnan followed by an infinitive.
% mdparagraph:P159:end

\end{minipage}\par\addvspace{4pt}

\reserve{7\baselineskip}
% mdparagraph:P160:begin
\subsection*{Practice I. Complete command pairs and instructions}\BookHeading{nav-ch06-learner-0021}{subsubsection}{Practice I. Complete command pairs and instructions}{0}
% mdparagraph:P160:end

\reserve{5\baselineskip}
\par\begin{minipage}{\linewidth}\raggedright
% mdparagraph:P161:begin
\textbf{I1.}
% mdparagraph:P161:end

% item: NEW-U06-I01
% mdparagraph:P162:begin
Give the second-person singular and plural imperatives of hīeran, nęrian and gefręmman. Label both numbers.
% mdparagraph:P162:end

\end{minipage}\par\addvspace{4pt}

\par\begin{minipage}{\linewidth}\raggedright
% mdparagraph:P163:begin
\textbf{I2.}
% mdparagraph:P163:end

% item: NEW-U06-I02
% mdparagraph:P164:begin
Give the second-person singular and plural imperatives of cyssan, gefyllan, sęttan, gierwan and sęllan. Label both numbers.
% mdparagraph:P164:end

\end{minipage}\par\addvspace{4pt}

\par\begin{minipage}{\linewidth}\raggedright
% mdparagraph:P165:begin
\textbf{I3.}
% mdparagraph:P165:end

% item: NEW-U06-I03
% mdparagraph:P166:begin
Give the second-person singular and plural imperatives of lufian, habban and libban. Label both numbers.
% mdparagraph:P166:end

\end{minipage}\par\addvspace{4pt}

\par\begin{minipage}{\linewidth}\raggedright
% mdparagraph:P167:begin
\textbf{I4.}
% mdparagraph:P167:end

% item: NEW-U06-I04
% mdparagraph:P168:begin
Give the second-person singular and plural imperatives of sęcgan and hycgan. Label both numbers.
% mdparagraph:P168:end

\end{minipage}\par\addvspace{4pt}

\par\begin{minipage}{\linewidth}\raggedright
% mdparagraph:P169:begin
\textbf{I5.}
% mdparagraph:P169:end

% item: NEW-U06-I05
% mdparagraph:P170:begin
Give the second-person singular and plural imperatives of drīfan, bēodan, bindan, helpan and drincan. Label both numbers.
% mdparagraph:P170:end

\end{minipage}\par\addvspace{4pt}

\par\begin{minipage}{\linewidth}\raggedright
% mdparagraph:P171:begin
\textbf{I6.}
% mdparagraph:P171:end

% item: NEW-U06-I06
% mdparagraph:P172:begin
Give the second-person singular and plural imperatives of wrēon and tēon. Label both numbers.
% mdparagraph:P172:end

\end{minipage}\par\addvspace{4pt}

\par\begin{minipage}{\linewidth}\raggedright
% mdparagraph:P173:begin
\textbf{I7.}
% mdparagraph:P173:end

% item: NEW-U06-I07
% mdparagraph:P174:begin
Give the second-person singular and plural imperatives of beran, niman and cuman. Label both numbers.
% mdparagraph:P174:end

\end{minipage}\par\addvspace{4pt}

\par\begin{minipage}{\linewidth}\raggedright
% mdparagraph:P175:begin
\textbf{I8.}
% mdparagraph:P175:end

% item: NEW-U06-I08
% mdparagraph:P176:begin
Give the second-person singular and plural imperatives of giefan, biddan, sittan, licgan and etan. Label both numbers.
% mdparagraph:P176:end

\end{minipage}\par\addvspace{4pt}

\par\begin{minipage}{\linewidth}\raggedright
% mdparagraph:P177:begin
\textbf{I9.}
% mdparagraph:P177:end

% item: NEW-U06-I09
% mdparagraph:P178:begin
Give the second-person singular and plural imperatives of sēon, slēan and fōn. Label both numbers.
% mdparagraph:P178:end

\end{minipage}\par\addvspace{4pt}

\par\begin{minipage}{\linewidth}\raggedright
% mdparagraph:P179:begin
\textbf{I10.}
% mdparagraph:P179:end

% item: NEW-U06-I10
% mdparagraph:P180:begin
Give the second-person singular and plural imperatives of hębban, scieppan and swęrian. Label both numbers.
% mdparagraph:P180:end

\end{minipage}\par\addvspace{4pt}

\par\begin{minipage}{\linewidth}\raggedright
% mdparagraph:P181:begin
\textbf{I11.}
% mdparagraph:P181:end

% item: NEW-U06-I11
% mdparagraph:P182:begin
Give the second-person singular and plural imperatives of faran, standan and healdan. Label both numbers.
% mdparagraph:P182:end

\end{minipage}\par\addvspace{4pt}

\par\begin{minipage}{\linewidth}\raggedright
% mdparagraph:P183:begin
\textbf{I12.}
% mdparagraph:P183:end

% item: NEW-U06-I12
% mdparagraph:P184:begin
Give the second-person singular and plural imperatives of weaxan, hātan and lǣtan. Label both numbers.
% mdparagraph:P184:end

\end{minipage}\par\addvspace{4pt}

\par\begin{minipage}{\linewidth}\raggedright
% mdparagraph:P185:begin
\textbf{I13.}
% mdparagraph:P185:end

% item: NEW-U06-I13
% mdparagraph:P186:begin
Give the second-person singular and plural imperatives of dōn, gān and both command series of wesan/\allowbreak{}bēon. Label both numbers.
% mdparagraph:P186:end

\end{minipage}\par\addvspace{4pt}

\par\begin{minipage}{\linewidth}\raggedright
% mdparagraph:P187:begin
\textbf{I14.}
% mdparagraph:P187:end

% item: NEW-U06-I14
% mdparagraph:P188:begin
Translate “Take the food and carry it.” First address one person, then several. Identify what “it” refers to and preserve its grammatical gender.
% mdparagraph:P188:end

\end{minipage}\par\addvspace{4pt}

\par\begin{minipage}{\linewidth}\raggedright
% mdparagraph:P189:begin
\textbf{I15.}
% mdparagraph:P189:end

% item: NEW-U06-I15
% mdparagraph:P190:begin
Translate “Help me. Give me the water.” First address one person, then exactly two people. Explain the verb number used for the two addressees.
% mdparagraph:P190:end

\end{minipage}\par\addvspace{4pt}

\par\begin{minipage}{\linewidth}\raggedright
% mdparagraph:P191:begin
\textbf{I16.}
% mdparagraph:P191:end

% item: NEW-U06-I16
% mdparagraph:P192:begin
Translate “Do not drink the water. Sit and eat.” Address several people.
% mdparagraph:P192:end

\end{minipage}\par\addvspace{4pt}

\par\begin{minipage}{\linewidth}\raggedright
% mdparagraph:P193:begin
\textbf{I17.}
% mdparagraph:P193:end

% item: NEW-U06-I17
% mdparagraph:P194:begin
Translate “The man said: ‘Go and do the work well.’” The man addresses one companion.
% mdparagraph:P194:end

\end{minipage}\par\addvspace{4pt}

\reserve{7\baselineskip}
% mdparagraph:P195:begin
\subsection*{Practice N. Infinitives and purpose}\BookHeading{nav-ch06-learner-0022}{subsubsection}{Practice N. Infinitives and purpose}{0}
% mdparagraph:P195:end

\reserve{5\baselineskip}
\par\begin{minipage}{\linewidth}\raggedright
% mdparagraph:P196:begin
\textbf{N1.}
% mdparagraph:P196:end

% item: NEW-U06-N01
% mdparagraph:P197:begin
Give the plain infinitive and the selected inflected infinitive after tō for bindan, beran and sęllan. Use the course’s -anne forms where the table selects them.
% mdparagraph:P197:end

\end{minipage}\par\addvspace{4pt}

\par\begin{minipage}{\linewidth}\raggedright
% mdparagraph:P198:begin
\textbf{N2.}
% mdparagraph:P198:end

% item: NEW-U06-N02
% mdparagraph:P199:begin
Give the plain infinitive and the selected inflected infinitive after tō for hīeran, cyssan, nęrian, gefręmman and gierwan. Use the course’s -anne forms where the table selects them.
% mdparagraph:P199:end

\end{minipage}\par\addvspace{4pt}

\par\begin{minipage}{\linewidth}\raggedright
% mdparagraph:P200:begin
\textbf{N3.}
% mdparagraph:P200:end

% item: NEW-U06-N03
% mdparagraph:P201:begin
Give the plain infinitive and the selected inflected infinitive after tō for lufian, habban, libban, sęcgan and hycgan. Use the course’s -anne forms where the table selects them.
% mdparagraph:P201:end

\end{minipage}\par\addvspace{4pt}

\par\begin{minipage}{\linewidth}\raggedright
% mdparagraph:P202:begin
\textbf{N4.}
% mdparagraph:P202:end

% item: NEW-U06-N04
% mdparagraph:P203:begin
Give the plain infinitive and the selected inflected infinitive after tō for wrēon, tēon, sēon, slēan and fōn. Use the course’s -anne forms where the table selects them.
% mdparagraph:P203:end

\end{minipage}\par\addvspace{4pt}

\par\begin{minipage}{\linewidth}\raggedright
% mdparagraph:P204:begin
\textbf{N5.}
% mdparagraph:P204:end

% item: NEW-U06-N05
% mdparagraph:P205:begin
Give the plain infinitive and the selected inflected infinitive after tō for gān, dōn and both wesan/\allowbreak{}bēon series. Use the course’s -anne forms where the table selects them.
% mdparagraph:P205:end

\end{minipage}\par\addvspace{4pt}

\par\begin{minipage}{\linewidth}\raggedright
% mdparagraph:P206:begin
\textbf{N6.}
% mdparagraph:P206:end

% item: NEW-U06-N06
% mdparagraph:P207:begin
Translate “He wants to eat. He came to eat.” Use the taught plain-infinitive complement in the first sentence and the taught purpose construction in the second.
% mdparagraph:P207:end

\end{minipage}\par\addvspace{4pt}

\par\begin{minipage}{\linewidth}\raggedright
% mdparagraph:P208:begin
\textbf{N7.}
% mdparagraph:P208:end

% item: NEW-U06-N07
% mdparagraph:P209:begin
Translate “The companions came to help the man.” Use the inflected infinitive for their purpose. Identify the person doing the helping and the case of the person helped.
% mdparagraph:P209:end

\end{minipage}\par\addvspace{4pt}

\par\begin{minipage}{\linewidth}\raggedright
% mdparagraph:P210:begin
\textbf{N8.}
% mdparagraph:P210:end

% item: NEW-U06-N08
% mdparagraph:P211:begin
Translate “The man came to give food to the birds.” Use the inflected infinitive for the purpose.
% mdparagraph:P211:end

\end{minipage}\par\addvspace{4pt}

\par\begin{minipage}{\linewidth}\raggedright
% mdparagraph:P212:begin
\textbf{N9.}
% mdparagraph:P212:end

% item: NEW-U06-N09
% mdparagraph:P213:begin
In Hīe sculon ðæt weorc dōn and Hīe cōmon tō dōnne ðæt weorc, identify each finite verb, each infinitive, and the construction that selects the infinitive form. Translate both sentences.
% mdparagraph:P213:end

\end{minipage}\par\addvspace{4pt}

\par\begin{minipage}{\linewidth}\raggedright
% mdparagraph:P214:begin
\textbf{N10.}
% mdparagraph:P214:end

% item: NEW-U06-N10
% mdparagraph:P215:begin
Complete the taught modal construction “We can carry the food” using mugan and beran. Then make the subject “I”. Explain which verb changes for person and number.
% mdparagraph:P215:end

\end{minipage}\par\addvspace{4pt}

\reserve{7\baselineskip}
% mdparagraph:P216:begin
\section*{6. Food and water for the journey}\BookHeading{nav-ch06-learner-0023}{subsection}{6. Food and water for the journey}{1}
% mdparagraph:P216:end

% mdparagraph:P217:begin
The following connected account is new teaching material. Read its statements and instructions together. A command requests an action; the later narrative reports what occurred.
% mdparagraph:P217:end

% composition: NEW-U06-TEXT-A
% mdparagraph:P218:begin
\par\begin{minipage}{\linewidth}\raggedright
\readingtext{Se mann cwæð: ‘Ic wille faran. Mīn scip is gearu. Cumað and helpað mē.’ Twēgen gefēran cōmon tō helpanne. Se mann cwæð: ‘Nimað ðone męte and berað hine. Giefað ðǣm fuglum ðæt wæter.’ Se gefēra cwæð: ‘Ic mæg ðone męte beran. Wē sculon ðæt weorc gefręmman.’ Ðā gefēran bǣron ðone męte and gēafon ðǣm fuglum ðæt wæter. Se mann cwæð: ‘Sittað and etað.’ Ðā gefēran ǣton. Hīe fōron.}
\end{minipage}\par
% mdparagraph:P218:end

\reserve{7\baselineskip}
% mdparagraph:P219:begin
\subsection*{Practice C. The connected account}\BookHeading{nav-ch06-learner-0024}{subsubsection}{Practice C. The connected account}{0}
% mdparagraph:P219:end

\reserve{5\baselineskip}
\par\begin{minipage}{\linewidth}\raggedright
% mdparagraph:P220:begin
\textbf{C1.}
% mdparagraph:P220:end

% item: NEW-U06-C01
% mdparagraph:P221:begin
Read “Food and water for the journey”. Who asks for help, how many companions come, and what is their stated purpose? Answer in English.
% mdparagraph:P221:end

\end{minipage}\par\addvspace{4pt}

\par\begin{minipage}{\linewidth}\raggedright
% mdparagraph:P222:begin
\textbf{C2.}
% mdparagraph:P222:end

% item: NEW-U06-C02
% mdparagraph:P223:begin
List every imperative in the controlled text and give its addressee number. For each, give the singular command form taught in the tables.
% mdparagraph:P223:end

\end{minipage}\par\addvspace{4pt}

\par\begin{minipage}{\linewidth}\raggedright
% mdparagraph:P224:begin
\textbf{C3.}
% mdparagraph:P224:end

% item: NEW-U06-C03
% mdparagraph:P225:begin
In the controlled text, identify the antecedent and case of hine, the recipient of the water, and what the companions subsequently eat. Distinguish what is explicit from what the context suggests.
% mdparagraph:P225:end

\end{minipage}\par\addvspace{4pt}

\par\begin{minipage}{\linewidth}\raggedright
% mdparagraph:P226:begin
\textbf{C4.}
% mdparagraph:P226:end

% item: NEW-U06-C04
% mdparagraph:P227:begin
Identify wille, mæg and sculon in the controlled text by person, number, tense and stated sense. Give the infinitive governed by each.
% mdparagraph:P227:end

\end{minipage}\par\addvspace{4pt}

\par\begin{minipage}{\linewidth}\raggedright
% mdparagraph:P228:begin
\textbf{C5.}
% mdparagraph:P228:end

% item: NEW-U06-C05
% mdparagraph:P229:begin
Translate the whole controlled text. Preserve the distinctions among a wish to travel, ability to carry, obligation, commands and completed past actions.
% mdparagraph:P229:end

\end{minipage}\par\addvspace{4pt}

\par\begin{minipage}{\linewidth}\raggedright
% mdparagraph:P230:begin
\textbf{C6.}
% mdparagraph:P230:end

% item: NEW-U06-C06
% mdparagraph:P231:begin
Rewrite the man’s first quoted commands, Cumað and helpað mē, as commands to one companion. Separately, put the first sentence of the original narrative into the present, keeping its original quotation unchanged.
% mdparagraph:P231:end

\end{minipage}\par\addvspace{4pt}

\reserve{7\baselineskip}
% mdparagraph:P232:begin
\section*{7. The fisherman’s gift and first direction}\BookHeading{nav-ch06-learner-0025}{subsection}{7. The fisherman’s gift and first direction}{1}
% mdparagraph:P232:end

% mdparagraph:P233:begin
\textbf{Editorial context.} Apollonius survived the earlier storm, encountered an old fisherman, and asked for help.
% mdparagraph:P233:end

% mdparagraph:P234:begin
This excerpt continues with the fisherman’s care and gift, ending after his first complete spoken direction. It is from Cook’s Reader XII, \textbf{Apollonius}, printed p.167 /\allowbreak{} one-based PDF p.183, editorial prose lines 3–10. The edition has printed marginal anchors at 5 and 10; the intervening line numbers are editorial. The speech continues in the source. A closing quotation mark is added at this excerpt’s boundary; the Old English wording and spellings are retained.
% mdparagraph:P234:end

% source: FB-XII-APOLLONIUS-FISHERMAN-GIFT-01
% mdparagraph:P235:begin
\par\begin{minipage}{\linewidth}\raggedright
\readingtext{Ðā, sōna swā se fiscere geseah ðæt se geonga mann æt his fōtum læg, hē mid mildheortnesse hine ūpāhōf, and lǣdde hine mid him tō his hūse, and ðā ēstas him beforan lęgde ðe hē him tō bēodanne hæfde. Ðā gīet hē wolde, be his mihte, māran ārfæstnesse him gecȳðan; tōslāt ðā his wǣfels on twā, and sealde Apollonie ðone healfan dǣl, ðus cweðende: “Nim ðæt ic ðē tō sęllanne hæbbe, and gā intō ðǣre ceastre.”}
\end{minipage}\par
% mdparagraph:P235:end

\reserve{7\baselineskip}
% mdparagraph:P236:begin
\subsection*{Local help}\BookHeading{nav-ch06-learner-0026}{subsubsection}{Local help}{0}
% mdparagraph:P236:end

% mdparagraph:P237:begin
These ten groups of help remain available during the reading exercises. Their unfamiliar forms and constructions are supplied for this passage; they are not independent-recall vocabulary or general construction patterns. The learned finite verbs, commands, and infinitive morphology can still be identified within the supported clauses.
% mdparagraph:P237:end

\reserve{5\baselineskip}
% mdparagraph:P238:begin
\textbf{1. Ðā, sōna swā se fiscere geseah ðæt se geonga mann æt his fōtum læg,}
% mdparagraph:P238:end

% support: U06-SRC-SUP-OPENING
% teaching: NEW-U06-SUPFORM-tha-adverb-invariant
% teaching: NEW-U06-SUPSENSE-tha-adverb-1
% teaching: NEW-U06-SUPFORM-sona-invariant-source
% teaching: NEW-U06-SUPSENSE-sona-as-soon-source
% teaching: NEW-U06-SUPFORM-swa-invariant
% teaching: NEW-U06-SUPSENSE-swa-as-correlative-source
% teaching: NEW-U06-SUPFORM-fiscere-nom-sg-source
% teaching: NEW-U06-SUPSENSE-fiscere-fisherman-source
% teaching: NEW-U06-SUPFORM-thaet-conjunction-invariant
% teaching: NEW-U06-SUPSENSE-thaet-conjunction-1
% teaching: NEW-U06-SUPFORM-geong-nom-m-sg-weak-source
% teaching: NEW-U06-SUPSENSE-geong-young-source
% teaching: NEW-U06-SUPFORM-aet-invariant-source
% teaching: NEW-U06-SUPSENSE-aet-at-source
% teaching: NEW-U06-SUPFORM-fot-dat-pl-source
% teaching: NEW-U06-SUPSENSE-fot-foot-source
% teaching: NEW-U06-GRAM-U06-SOURCE-TEMPORAL-CONTENT
% teaching: NEW-U06-GRAM-U06-SOURCE-PREPOSITIONS
% mdparagraph:P239:begin
Then, as soon as the fisherman saw that the young man was lying at his feet.
% mdparagraph:P239:end

% mdparagraph:P240:begin
\textbf{Ðā} means “then”; \textbf{sōna swā}, “as soon as.” Fiscere is nominative singular, “fisherman”; geonga is the weak masculine nominative singular adjective, “young.” The temporal clause contains the whole content clause introduced by ðæt, “that.” Its finite verb læg means “lay; was lying.” In æt his fōtum, “at his feet,” fōtum is dative plural. The feet are the fisherman’s. Both nested clause relations and the prepositional phrase are supplied as a whole.
% mdparagraph:P240:end

\reserve{5\baselineskip}
% mdparagraph:P241:begin
\textbf{2. hē mid mildheortnesse hine ūpāhōf,}
% mdparagraph:P241:end

% support: U06-SRC-SUP-COMPASSION
% teaching: NEW-U06-SUPFORM-mid-invariant-source
% teaching: NEW-U06-SUPSENSE-mid-manner-source
% teaching: NEW-U06-SUPFORM-mildheortnes-dat-sg-source
% teaching: NEW-U06-SUPSENSE-mildheortnes-compassion-source
% teaching: NEW-U06-SUPFORM-up-ahebban-pret-3-sg-source
% teaching: NEW-U06-SUPSENSE-up-ahebban-lift-source
% mdparagraph:P242:begin
He compassionately lifted him up.
% mdparagraph:P242:end

% mdparagraph:P243:begin
\textbf{Mid mildheortnesse}, “with compassion,” is a supplied manner phrase with dative mildheortnesse. \textbf{Ūpāhōf} means “lifted up”: it is the supplied third-singular preterite of a prefixed verb, not the unprefixed hębban form. Hē is the fisherman; hine is Apollonius.
% mdparagraph:P243:end

\reserve{5\baselineskip}
% mdparagraph:P244:begin
\textbf{3. and lǣdde hine mid him tō his hūse,}
% mdparagraph:P244:end

% support: U06-SRC-SUP-HOUSE
% teaching: NEW-U06-SUPSENSE-mid-accompaniment-source
% teaching: NEW-U06-SUPFORM-to-preposition-dative-invariant
% teaching: NEW-U06-SUPSENSE-to-destination
% teaching: NEW-U06-SUPFORM-hus-dat-sg-source
% teaching: NEW-U06-SUPSENSE-hus-house-source
% mdparagraph:P245:begin
And led him with him to his house.
% mdparagraph:P245:end

% mdparagraph:P246:begin
\textbf{Mid him} means “with him,” here accompanying the fisherman. \textbf{Tō his hūse} means “to his house”; hūse is the supplied dative singular of hūs, “house.” In this clause hine is Apollonius, while him and his refer contextually to the fisherman. The accompaniment and destination phrases are supplied locally. Tō before the noun phrase is distinct from tō before an inflected infinitive.
% mdparagraph:P246:end

\reserve{5\baselineskip}
% mdparagraph:P247:begin
\textbf{4. and ðā ēstas him beforan lęgde ðe hē him tō bēodanne hæfde.}
% mdparagraph:P247:end

% support: U06-SRC-SUP-PROVISIONS
% teaching: NEW-U06-SUPFORM-est-acc-m-pl-source
% teaching: NEW-U06-SUPSENSE-est-food-source
% teaching: NEW-U06-SUPFORM-beforan-invariant-source
% teaching: NEW-U06-SUPSENSE-beforan-before-source
% teaching: NEW-U06-SUPFORM-lecgean-pret-3-sg-source
% teaching: NEW-U06-SUPSENSE-lecgean-place-source
% teaching: NEW-U06-SUPFORM-the-relative-invariant
% teaching: NEW-U06-SUPSENSE-the-relative-which
% teaching: NEW-U06-GRAM-U06-SOURCE-AVAILABILITY
% teaching: NEW-U06-GRAM-U06-SOURCE-RELATIVES
% mdparagraph:P248:begin
And set before him the provisions that he had to offer him.
% mdparagraph:P248:end

% mdparagraph:P249:begin
\textbf{Ēstas} is supplied masculine accusative plural, “provisions; food”; \textbf{lęgde} is a supplied third-singular preterite, “placed; set.” \textbf{Him beforan} means “before him,” with the dative person. The following \textbf{ðe} is an object-relative referring to the provisions: “that/\allowbreak{}which he had to offer him.”
% mdparagraph:P249:end

% mdparagraph:P250:begin
Here “had to offer” means \textbf{had available for offering}, not “was obliged to offer.” Tō bēodanne has the learned inflected-infinitive form; bēodan’s “offer” meaning is familiar. The whole availability construction, its relative clause, and its offered-object/\allowbreak{}recipient relation remain supplied. Both occurrences of him refer to Apollonius; the second is not a reflexive “himself.”
% mdparagraph:P250:end

\reserve{5\baselineskip}
% mdparagraph:P251:begin
\textbf{5. Ðā gīet hē wolde, be his mihte, māran ārfæstnesse him gecȳðan;}
% mdparagraph:P251:end

% support: U06-SRC-SUP-KINDNESS
% teaching: NEW-U06-SUPFORM-giet-invariant-source
% teaching: NEW-U06-SUPSENSE-giet-still-source
% teaching: NEW-U06-SUPFORM-be-invariant-source
% teaching: NEW-U06-SUPSENSE-be-according-source
% teaching: NEW-U06-SUPFORM-miht-dat-sg-source
% teaching: NEW-U06-SUPSENSE-miht-power-source
% teaching: NEW-U06-SUPFORM-arfaestnes-acc-sg-source
% teaching: NEW-U06-SUPSENSE-arfaestnes-kindness-source
% teaching: NEW-U06-SUPFORM-gecythan-infinitive-source
% teaching: NEW-U06-SUPSENSE-gecythan-show-source
% mdparagraph:P252:begin
He still wanted, according to his ability, to show him greater kindness.
% mdparagraph:P252:end

% mdparagraph:P253:begin
\textbf{Gīet} means “still; further.” \textbf{Be his mihte}, “according to his ability,” contains dative singular mihte from the noun miht, “ability; power.” It is not a modal verb here. \textbf{Ārfæstnesse} is accusative singular, “kindness”; māran agrees with it as the comparative “greater.” \textbf{Gecȳðan} is a supplied plain infinitive, “show.” In this exact use the thing shown is accusative and its recipient is dative him, Apollonius. The finite wolde expresses the fisherman’s past desire. The prefixed verb and its local government are supplied, not a rule for forming new prefixed verbs.
% mdparagraph:P253:end

\reserve{5\baselineskip}
% mdparagraph:P254:begin
\textbf{6. tōslāt ðā his wǣfels on twā,}
% mdparagraph:P254:end

% support: U06-SRC-SUP-TEARING
% teaching: NEW-U06-SUPFORM-toslitan-pret-3-sg-source
% teaching: NEW-U06-SUPSENSE-toslitan-tear-source
% teaching: NEW-U06-SUPFORM-waefels-acc-sg-source
% teaching: NEW-U06-SUPSENSE-waefels-cloak-source
% teaching: NEW-U06-SUPSENSE-on-divided-into-two-source
% teaching: NEW-U06-SUPSENSE-twegen-divided-in-two-source
% teaching: NEW-U06-GRAM-U06-SOURCE-ELLIPSIS
% mdparagraph:P255:begin
He then tore his cloak in two.
% mdparagraph:P255:end

% mdparagraph:P256:begin
\textbf{Tōslāt} is the supplied third-singular preterite, “tore.” \textbf{Wǣfels} is the accusative singular object, “cloak.” The fisherman remains the subject without a repeated pronoun. Keep \textbf{on twā}, “in two,” as a whole supplied expression. Its usual analysis takes twā as neuter accusative, but that spelling also occurs in feminine forms; it does not identify an unexpressed noun or establish a general rule for motion with on.
% mdparagraph:P256:end

\reserve{5\baselineskip}
% mdparagraph:P257:begin
\textbf{7. and sealde Apollonie ðone healfan dǣl,}
% mdparagraph:P257:end

% support: U06-SRC-SUP-GIFT
% teaching: NEW-U06-SUPFORM-apollonius-dat-sg-source
% teaching: NEW-U06-SUPSENSE-apollonius-person
% teaching: NEW-U06-SUPFORM-healf-acc-m-sg-weak-source
% teaching: NEW-U06-SUPSENSE-healf-half-source
% teaching: NEW-U06-SUPFORM-dael-acc-sg-source
% teaching: NEW-U06-SUPSENSE-dael-part-source
% mdparagraph:P258:begin
And gave Apollonius one half.
% mdparagraph:P258:end

% mdparagraph:P259:begin
\textbf{Apollonie} is the supplied dative singular name, “to Apollonius.” \textbf{Ðone healfan dǣl} means “one half,” literally “the half part.” Dǣl, “part,” is accusative singular; healfan is the weak masculine accusative singular of healf, “half.” This is the familiar giving construction with its recipient and thing given. No second object has to be added.
% mdparagraph:P259:end

\reserve{5\baselineskip}
% mdparagraph:P260:begin
\textbf{8. ðus cweðende:}
% mdparagraph:P260:end

% support: U06-SRC-SUP-SPEECH
% teaching: NEW-U06-SUPFORM-thus-invariant-source
% teaching: NEW-U06-SUPSENSE-thus-thus-source
% teaching: NEW-U06-SUPFORM-cwethan-present-participle-source
% teaching: NEW-U06-GRAM-U06-SOURCE-PARTICIPLE
% mdparagraph:P261:begin
Saying thus.
% mdparagraph:P261:end

% mdparagraph:P262:begin
\textbf{Ðus} means “thus.” \textbf{Cweðende}, “saying,” is a supplied present participle accompanying the fisherman’s action. It introduces his words without serving as a finite past or an imperative. Compare its role with the familiar finite cwæð; no productive participial construction is required here.
% mdparagraph:P262:end

\reserve{5\baselineskip}
% mdparagraph:P263:begin
\textbf{9. “Nim ðæt ic ðē tō sęllanne hæbbe,}
% mdparagraph:P263:end

% support: U06-SRC-SUP-TAKE
% teaching: NEW-U06-SUPSENSE-se-free-relative-what-source
% mdparagraph:P264:begin
Take what I have to give you.
% mdparagraph:P264:end

% mdparagraph:P265:begin
\textbf{Nim} is a singular command. Supply the whole following phrase \textbf{ðæt ic ðē tō sęllanne hæbbe}, “what I have to give you.” Here ðæt introduces a definite free relative, “what/\allowbreak{}that which”; it is not merely the article before an omitted noun in an independently learned rule.
% mdparagraph:P265:end

% mdparagraph:P266:begin
“Have to give” again means \textbf{have available to give}, not obligation. The fisherman is ic; Apollonius is the dative recipient ðē. Sęllanne has the learned inflected-infinitive form. Hæbbe is read here as the familiar first-singular indicative, but its isolated spelling is also shared with subjunctive forms: the letters alone do not uniquely identify its mood. The whole relative and availability construction remains supplied.
% mdparagraph:P266:end

\reserve{5\baselineskip}
% mdparagraph:P267:begin
\textbf{10. and gā intō ðǣre ceastre.”}
% mdparagraph:P267:end

% support: U06-SRC-SUP-GO
% teaching: NEW-U06-SUPFORM-into-invariant-source
% teaching: NEW-U06-SUPSENSE-into-destination-source
% teaching: NEW-U06-SUPFORM-ceaster-dat-sg-source
% teaching: NEW-U06-SUPSENSE-ceaster-city-source
% mdparagraph:P268:begin
And go into the town.
% mdparagraph:P268:end

% mdparagraph:P269:begin
\textbf{Gā} is a singular command. \textbf{Intō ðǣre ceastre} is supplied as “into the town.” Ceastre is dative singular, with matching feminine dative ðǣre. Preserve this actual phrase; the presence of motion does not justify changing it to an accusative or deriving a general rule for every destination. This completes the first spoken direction; do not add the continuation.
% mdparagraph:P269:end

% mdparagraph:P270:begin
\textbf{Source note.} The excerpt is Cook Reader XII, p.167 /\allowbreak{} PDF p.183, editorial lines 3–10. The source has \textbf{læg} without a length mark and \textbf{gīet} with long ī. Line breaks are joined and the attached note cues are moved below. The final closing quotation mark marks the editorial stopping point; it is not printed there in the continuing speech.
% mdparagraph:P270:end

\reserve{7\baselineskip}
% mdparagraph:P271:begin
\subsection*{Cook’s original notes}\BookHeading{nav-ch06-learner-0027}{subsubsection}{Cook’s original notes}{0}
% mdparagraph:P271:end

% mdparagraph:P272:begin
The numbered references below are Cook’s, not sections of this unit. The local help above supplies the needed interpretation.
% mdparagraph:P272:end

% source-note: FB167-n5
% mdparagraph:P273:begin
\textbf{P.167, note 5.} fiscere: See 143.
% mdparagraph:P273:end

% source-note: FB167-n6
% mdparagraph:P274:begin
\textbf{P.167, note 6.} ēstas: Lat. epulas.
% mdparagraph:P274:end

% source-note: FB167-n7
% mdparagraph:P275:begin
\textbf{P.167, note 7.} ārfæstnesse: MS. fæstnesse. Lat. pietati.
% mdparagraph:P275:end

\reserve{7\baselineskip}
% mdparagraph:P276:begin
\subsection*{Practice R. The original excerpt}\BookHeading{nav-ch06-learner-0028}{subsubsection}{Practice R. The original excerpt}{0}
% mdparagraph:P276:end

% mdparagraph:P277:begin
Use all the printed local help. Identification of learned forms does not require independent recall of the supplied clauses.
% mdparagraph:P277:end

\reserve{5\baselineskip}
\par\begin{minipage}{\linewidth}\raggedright
% mdparagraph:P278:begin
\textbf{R1.}
% mdparagraph:P278:end

% item: NEW-U06-R01
% mdparagraph:P279:begin
Read “The fisherman’s gift and first direction”, using its local support. What prompts the fisherman’s action, what does he divide, and what does he give Apollonius?
% mdparagraph:P279:end

\end{minipage}\par\addvspace{4pt}

\par\begin{minipage}{\linewidth}\raggedright
% mdparagraph:P280:begin
\textbf{R2.}
% mdparagraph:P280:end

% item: NEW-U06-R02
% mdparagraph:P281:begin
Find wolde and its plain infinitive in the source. Parse wolde and explain the fisherman’s intention. Use the supplied meaning of gecȳðan.
% mdparagraph:P281:end

\end{minipage}\par\addvspace{4pt}

\par\begin{minipage}{\linewidth}\raggedright
% mdparagraph:P282:begin
\textbf{R3.}
% mdparagraph:P282:end

% item: NEW-U06-R03
% mdparagraph:P283:begin
Identify tō bēodanne and tō sęllanne. Give the familiar plain infinitive behind each. Explain the source’s “have to offer/\allowbreak{}give” phrases using the printed support; do not translate them as modern English “must”.
% mdparagraph:P283:end

\end{minipage}\par\addvspace{4pt}

\par\begin{minipage}{\linewidth}\raggedright
% mdparagraph:P284:begin
\textbf{R4.}
% mdparagraph:P284:end

% item: NEW-U06-R04
% mdparagraph:P285:begin
Find the source’s two singular commands. Parse each and give its plural command form from the taught system.
% mdparagraph:P285:end

\end{minipage}\par\addvspace{4pt}

\par\begin{minipage}{\linewidth}\raggedright
% mdparagraph:P286:begin
\textbf{R5.}
% mdparagraph:P286:end

% item: NEW-U06-R05
% mdparagraph:P287:begin
Use the notes to distinguish the participial cweðende from the finite past cwæð and the command gā. Which form introduces the fisherman’s actual words?
% mdparagraph:P287:end

\end{minipage}\par\addvspace{4pt}

\par\begin{minipage}{\linewidth}\raggedright
% mdparagraph:P288:begin
\textbf{R6.}
% mdparagraph:P288:end

% item: NEW-U06-R06
% mdparagraph:P289:begin
What does the fisherman say he has, and what first direction does he give? Explain the force of the phrase translated “what I have” using the local support.
% mdparagraph:P289:end

\end{minipage}\par\addvspace{4pt}

\par\begin{minipage}{\linewidth}\raggedright
% mdparagraph:P290:begin
\textbf{R7.}
% mdparagraph:P290:end

% item: NEW-U06-R07
% mdparagraph:P291:begin
Identify the source’s tō phrases that contain an inflected infinitive and distinguish them from tō his hūse and intō ðǣre ceastre. State the locally supplied meanings of the destination phrases without deriving a general motion-preposition rule.
% mdparagraph:P291:end

\end{minipage}\par\addvspace{4pt}

\par\begin{minipage}{\linewidth}\raggedright
% mdparagraph:P292:begin
\textbf{R8.}
% mdparagraph:P292:end

% item: NEW-U06-R08
% mdparagraph:P293:begin
Translate the whole 76-word source excerpt with its supports. Stop after the fisherman’s first direction; do not add the continuation of his speech.
% mdparagraph:P293:end

\end{minipage}\par\addvspace{4pt}

\reserve{7\baselineskip}
% mdparagraph:P294:begin
\section*{8. Cumulative practice}\BookHeading{nav-ch06-learner-0029}{subsection}{8. Cumulative practice}{1}
% mdparagraph:P294:end

\reserve{7\baselineskip}
% mdparagraph:P295:begin
\subsection*{Practice X. Combining the constructions}\BookHeading{nav-ch06-learner-0030}{subsubsection}{Practice X. Combining the constructions}{0}
% mdparagraph:P295:end

% mdparagraph:P296:begin
The final task admits more than one coherent answer. Keep to its lexical bank and the constructions taught here; the four sentences should form a connected account.
% mdparagraph:P296:end

\reserve{5\baselineskip}
\par\begin{minipage}{\linewidth}\raggedright
% mdparagraph:P297:begin
\textbf{X1.}
% mdparagraph:P297:end

% item: NEW-U06-X01
% mdparagraph:P298:begin
Translate “We two wanted to give our food to the birds. We two gave it to them.” Retain the two-person possessor and the grammatical gender of “food”.
% mdparagraph:P298:end

\end{minipage}\par\addvspace{4pt}

\par\begin{minipage}{\linewidth}\raggedright
% mdparagraph:P299:begin
\textbf{X2.}
% mdparagraph:P299:end

% item: NEW-U06-X02
% mdparagraph:P300:begin
Translate “The old man cannot carry the great ship. Help him.” Address several companions in the command.
% mdparagraph:P300:end

\end{minipage}\par\addvspace{4pt}

\par\begin{minipage}{\linewidth}\raggedright
% mdparagraph:P301:begin
\textbf{X3.}
% mdparagraph:P301:end

% item: NEW-U06-X03
% mdparagraph:P302:begin
Translate “He came to help us. We wanted to give him food. He sat and ate.” Use the taught elementary purpose and volitional constructions.
% mdparagraph:P302:end

\end{minipage}\par\addvspace{4pt}

\par\begin{minipage}{\linewidth}\raggedright
% mdparagraph:P303:begin
\textbf{X4.}
% mdparagraph:P303:end

% item: NEW-U06-X04
% mdparagraph:P304:begin
Compare the statement Gē gāð with the command Gāð. Explain the shared spelling and the difference in use. Then compare Hē dēð ðæt weorc with Dō ðæt weorc.
% mdparagraph:P304:end

\end{minipage}\par\addvspace{4pt}

\par\begin{minipage}{\linewidth}\raggedright
% mdparagraph:P305:begin
\textbf{X5.}
% mdparagraph:P305:end

% item: NEW-U06-X05
% mdparagraph:P306:begin
Write four connected Old English sentences about a man and two companions: one past arrival with an explicitly stated purpose, one statement of willingness, one direct command addressed to both companions, and one completed past action. Use cuman, helpan, willan and beran, with mann, gefēra and męte and any needed taught articles, numerals or pronouns. Label the three new constructions.
% mdparagraph:P306:end

\end{minipage}\par\addvspace{4pt}

\clearpage
% mdparagraph:P307:begin
\section*{9. Sources and form choices}\BookHeading{nav-ch06-learner-0031}{subsection}{9. Sources and form choices}{1}
% mdparagraph:P307:end

\begingroup\fontsize{10}{12.8}\selectfont
\setlength{\parskip}{1.5pt}
% mdparagraph:P308:begin
The core forms follow Albert S. Cook, \textbf{A First Book in Old English: Grammar, Reader, Notes, and Vocabulary}, second edition, revised and enlarged (Ginn, 1895). The compact tables explicitly identify the three persons covered by each shared plural row and apply the given endings to the selected stems. Such paradigm completion is not a claim that each cell is separately quoted in continuous prose.
% mdparagraph:P308:end

% mdparagraph:P309:begin
Joseph Wright and Elizabeth Mary Wright, \textbf{Old English Grammar}, second edition (Oxford, 1914), corroborates the forms and variants; the precise witan and magan variant locators follow below. R. D. Fulk, \textbf{An Introductory Grammar of Old English with an Anthology of Readings}, author-corrected open-access reprint (2020), supplies the cited distinctions concerning defects, j-presents, the bēon series, and infinitival use. \textbf{Bosworth–Toller, An Anglo-Saxon Dictionary}, main dictionary (1898), supplies the individually cited lexical senses and constructions; page letters a and b identify columns.
% mdparagraph:P309:end

% mdparagraph:P310:begin
All PDF page numbers in this unit are one-based. In the selected Cook scan the PDF page is printed page +16; in Wright, +20; in Fulk, +12. Dictionary facsimiles are cited by printed page and column rather than by an invented PDF position.
% mdparagraph:P310:end

% mdparagraph:P311:begin
The examples and “Food and water for the journey” are new compositions. The original fisherman excerpt preserves its source wording, with only the declared closing quotation mark added. Its supplied nouns, participle, prepositions, and clause constructions remain local help. Earlier participles and other reference-only forms do not become command or infinitive systems through resemblance in spelling.
% mdparagraph:P311:end

% mdparagraph:P312:begin
\textbf{Local reading help.} Compare Cook §§87a and 87d, pp.51–52 /\allowbreak{} PDF pp.67–68, for relatives; §164b, p.94 /\allowbreak{} PDF p.110, for offering and giving; and §§194–200, pp.104–106 /\allowbreak{} PDF pp.120–122, for the verbal forms and uses. The availability and relative phrases in the reading remain supplied as a whole.
% mdparagraph:P312:end

% mdparagraph:P313:begin
The supplied lexical help follows the corresponding entries in Cook’s \textbf{Vocabulary}, whose headword spellings are retained in these references: æt (p.274 /\allowbreak{} PDF p.290); ār-fæstnes (p.276 /\allowbreak{} PDF p.292); be (p.277 /\allowbreak{} PDF p.293); be-foran; bēodan (p.278 /\allowbreak{} PDF p.294); ceaster (p.281 /\allowbreak{} PDF p.297); dǣl; cuman (p.282 /\allowbreak{} PDF p.298); ēst (p.285 /\allowbreak{} PDF p.301); fiscere (p.287 /\allowbreak{} PDF p.303); fōt (p.289 /\allowbreak{} PDF p.305); ge-cȳðan (p.291 /\allowbreak{} PDF p.307); geong (p.293 /\allowbreak{} PDF p.309); gīet (p.296 /\allowbreak{} PDF p.312); healf (p.298 /\allowbreak{} PDF p.314); hūs (p.300 /\allowbreak{} PDF p.316); in-tō (p.301 /\allowbreak{} PDF p.317); lęcgean (p.302 /\allowbreak{} PDF p.318); mid; miht; mild-heortnes (p.305 /\allowbreak{} PDF p.321); sōna (p.313 /\allowbreak{} PDF p.329); tō; tō-slītan (p.316 /\allowbreak{} PDF p.332); ðus (p.319 /\allowbreak{} PDF p.335); ūp-ā-hębban; wǣfels (p.320 /\allowbreak{} PDF p.336).
% mdparagraph:P313:end

% mdparagraph:P314:begin
\textbf{Witan variants.} Wright, §540, printed p.283 /\allowbreak{} PDF p.303.
% mdparagraph:P314:end

% mdparagraph:P315:begin
\textbf{Magan variants.} Wright, §§544–545, printed p.285 /\allowbreak{} PDF p.305.
% mdparagraph:P315:end

\endgroup

% book-source-body:end
\par\label{end-ch06-learner}
\clearpage\endgroup

\BookRecto
\begingroup
\renewcommand{\runningtitle}{Unit 7. Subjunctive and passive}
\hyphenpenalty=10000
\exhyphenpenalty=10000
\emergencystretch=3em
\newcommand{\itempara}[2]{\par\addvspace{3pt}\noindent\hangindent=1.8em\hangafter=1\makebox[1.8em][l]{#1.}#2\par}
\newcommand{\reserve}[1]{\par\penalty-100\begingroup\dimen0=\pagegoal\advance\dimen0 by-\pagetotal\ifdim\dimen0<#1\newpage\fi\endgroup}
% book-source-body:begin

\addfontfeatures{Ligatures=TeXOff}
\raggedright
\booktitle{\BookComponent{ch07-learner}{Unit 7. Subjunctive, participles, and passive}UNIT 7}{Subjunctive, participles, and passive}

% mdparagraph:P001:begin
This unit adds wishes and exhortations, basic conditional clauses, participles, and passive constructions. It gives complete present and preterite subjunctives for the verbs named below, together with the new vocabulary needed for Cook I.6–7.
% mdparagraph:P001:end

% mdparagraph:P002:begin
In verbal tables, \textbf{1}, \textbf{2}, and \textbf{3} mean first, second, and third person; \textbf{sg.} and \textbf{pl.} mean singular and plural. A row headed \textbf{1–3 sg.} or \textbf{1–3 pl.} includes all three persons. Dual subjects take plural finite forms. Noun and adjective tables use nominative, genitive, dative, and accusative; the shorter labels are \textbf{nom}, \textbf{gen}, \textbf{dat}, and \textbf{acc}. Preserve vowel length and the distinction between e and ę. Examples and the account “Binding the animals” are new teaching material; the final Creation paragraph is an original source excerpt.
% mdparagraph:P002:end

\reserve{7\baselineskip}
% mdparagraph:P003:begin
\section*{1. Assertions, wishes, and commands}\BookHeading{nav-ch07-learner-0001}{subsection}{1. Assertions, wishes, and commands}{1}\BookLabel{bk-ch07-learner-bd7d043cab91}
% mdparagraph:P003:end

% teaching: NEW-U07-GRAM-U07-WISH-EXHORTATION
% mdparagraph:P004:begin
The \textbf{indicative} presents an assertion. The \textbf{subjunctive}, called the \textbf{optative} by Cook, can express a wish or an exhortation. In an independent clause its present tense can mean “may …” or “let …”. This English “let” translates the force of the inflected verb; no separate Old English verb meaning “let” is needed.
% mdparagraph:P004:end

% example: NEW-U07-EXAMPLE-01
% mdparagraph:P005:begin
Hīe swimmað on ðǣm sǣ. Swimmen hīe on ðǣm sǣ.
% mdparagraph:P005:end

% mdparagraph:P006:begin
They swim in the sea. Let them swim in the sea.
% mdparagraph:P006:end

% mdparagraph:P007:begin
The first clause asserts an activity; the second calls for it. Swimmað is indicative; swimmen is subjunctive. Both have the plural subject hīe. The sea phrase uses the masculine sǣ paradigm learned earlier.
% mdparagraph:P007:end

% example: NEW-U07-EXAMPLE-02
% mdparagraph:P008:begin
Gāð. Gān wē.
% mdparagraph:P008:end

% mdparagraph:P009:begin
Go. Let us go.
% mdparagraph:P009:end

% mdparagraph:P010:begin
Gāð is an imperative addressed to several people. Gān with wē is a first-person plural subjunctive exhortation: it includes the speaker. The selected subjunctive plural retains -n before its following pronoun.
% mdparagraph:P010:end

% mdparagraph:P011:begin
The familiar imperative gives a direct second-person command. The subjunctive permits first- and third-person wishes as well. The person belongs to the subject and context: its ending distinguishes number, but not first, second, and third person within that number.
% mdparagraph:P011:end

% teaching: NEW-U07-GRAM-U07-MODAL-MOOD-DISTINCTION
% example: NEW-U07-EXAMPLE-03
\reserve{4\baselineskip}
% mdparagraph:P012:begin
Hē mæg swimman. Hē mōt swimman. Swimme hē.
% mdparagraph:P012:end

% mdparagraph:P013:begin
He can swim. He may swim. Let him swim.
% mdparagraph:P013:end

% mdparagraph:P014:begin
Mæg asserts capability; mōt asserts permission. Each governs the plain infinitive swimman. Swimme expresses a wish about his swimming without asserting either capability or permission.
% mdparagraph:P014:end

% mdparagraph:P015:begin
Keep the earlier meanings of cunnan “know how”, ðurfan “need”, durran “dare”, sculan “be obliged”, and willan “wish, want” distinct from inflectional mood. A statement that someone wants something can itself be indicative. Willan is not automatically an English future auxiliary, and wolde is not automatically “would”. Where indicative and subjunctive spellings coincide, the clause setting must decide the intended use.
% mdparagraph:P015:end

\reserve{7\baselineskip}
% mdparagraph:P016:begin
\section*{2. The complete subjunctive systems}\BookHeading{nav-ch07-learner-0002}{subsection}{2. The complete subjunctive systems}{1}\BookLabel{bk-ch07-learner-813bcaedf4f5}
% mdparagraph:P016:end

% teaching: NEW-U07-GRAM-U07-SUBJUNCTIVE-ENDINGS
% teaching: NEW-U07-GRAM-U07-SUBJUNCTIVE-STEMS
% mdparagraph:P017:begin
The ordinary endings are \textbf{-e} throughout the singular and \textbf{-en} throughout the plural, in both present and preterite. Thus four displayed forms for each verb give all twelve person–number–tense combinations. Contracted and irregular forms have their actual spellings in the tables.
% mdparagraph:P017:end

% mdparagraph:P018:begin
For the present, use the selected present stem. Do not carry the indicative second- or third-singular vowel mutation into the whole subjunctive. For a strong preterite, use the \textbf{preterite plural stem}, with its vowel and consonant changes, in both numbers. A weak preterite retains its dental stem. The preterite indicative plural -on and subjunctive plural -en remain distinct.
% mdparagraph:P018:end

% mdparagraph:P019:begin
The named verbs define this unit’s complete subjunctive assignment. They do not assign a subjunctive system to every verb previously encountered in a reading.
% mdparagraph:P019:end

\reserve{7\baselineskip}
% mdparagraph:P020:begin
\subsection*{Weak class I}\BookHeading{nav-ch07-learner-0003}{subsubsection}{Weak class I}{0}\BookLabel{bk-ch07-learner-407c09325356}
% mdparagraph:P020:end

% teaching: NEW-U07-SUBJ-hieran
% teaching: NEW-U07-SUBJ-sellan
% teaching: NEW-U07-SUBJ-gewyrcean
% teaching: NEW-U07-SUBJ-wyrcean
% teaching: NEW-U07-SUBJ-gefremman
% teaching: NEW-U07-SUBJ-gierwan
% teaching: NEW-U07-SUBJ-nerian
% mdparagraph:P021:begin
Keep the irregular past stems with their verbs: sęllan has sealde, and gewyrcean/\allowbreak{}wyrcean have geworhte/\allowbreak{}worhte. Gefręmman retains double mm in the present, but its past has single m before -ede. Gierwan retains w in the selected present and has gierede in the past. Nęrian retains i in -ie/\allowbreak{}-ien; its past is -ede, despite the infinitive’s resemblance to weak class II.
% mdparagraph:P021:end

% table: U07-DISPLAY-SUBJ-1
% mdtable:T01:begin
% course-table-style: 2026-10
% course-panel:T01-P01:begin
\par\begingroup
\fontsize{10.5}{13.2}\selectfont
\setlength{\LTleft}{0pt}\setlength{\LTright}{\fill}
\noindent
\begin{tabular}{@{}>{\raggedright\arraybackslash}p{20.873mm}>{\raggedright\arraybackslash}p{21.064mm}>{\raggedright\arraybackslash}p{24.876mm}>{\raggedright\arraybackslash}p{21.064mm}>{\raggedright\arraybackslash}p{24.876mm}@{}}
\toprule
\textbf{Person and number} & \textbf{hīeran · Present} & \textbf{hīeran · Preterite} & \textbf{sęllan · Present} & \textbf{sęllan · Preterite} \\
\midrule
1–3 sg. & hīere & hīerde & sęlle & sealde \\
1–3 pl. & hīeren & hīerden & sęllen & sealden \\
\bottomrule
\end{tabular}
\par\endgroup
% course-panel:T01-P01:end
\par\addvspace{5pt}
% course-panel:T01-P02:begin
\par\begingroup
\fontsize{10.5}{13.2}\selectfont
\setlength{\LTleft}{0pt}\setlength{\LTright}{\fill}
\noindent
\begin{tabular}{@{}>{\raggedright\arraybackslash}p{19.336mm}>{\raggedright\arraybackslash}p{23.839mm}>{\raggedright\arraybackslash}p{25.429mm}>{\raggedright\arraybackslash}p{20.308mm}>{\raggedright\arraybackslash}p{23.839mm}@{}}
\toprule
\textbf{Person and number} & \textbf{gewyrcean · Present} & \textbf{gewyrcean · Preterite} & \textbf{wyrcean · Present} & \textbf{wyrcean · Preterite} \\
\midrule
1–3 sg. & gewyrce & geworhte & wyrce & worhte \\
1–3 pl. & gewyrcen & geworhten & wyrcen & worhten \\
\bottomrule
\end{tabular}
\par\endgroup
% course-panel:T01-P02:end
\par\addvspace{5pt}
% course-panel:T01-P03:begin
\par\begingroup
\fontsize{10.5}{13.2}\selectfont
\setlength{\LTleft}{0pt}\setlength{\LTright}{\fill}
\noindent
\begin{tabular}{@{}>{\raggedright\arraybackslash}p{19.171mm}>{\raggedright\arraybackslash}p{23.636mm}>{\raggedright\arraybackslash}p{26.175mm}>{\raggedright\arraybackslash}p{20.134mm}>{\raggedright\arraybackslash}p{23.636mm}@{}}
\toprule
\textbf{Person and number} & \textbf{gefręmman · Present} & \textbf{gefręmman · Preterite} & \textbf{gierwan · Present} & \textbf{gierwan · Preterite} \\
\midrule
1–3 sg. & gefręmme & gefręmede & gierwe & gierede \\
1–3 pl. & gefręmmen & gefręmeden & gierwen & giereden \\
\bottomrule
\end{tabular}
\par\endgroup
% course-panel:T01-P03:end
\par\addvspace{5pt}
\par\addvspace{3pt}
% course-caption:T01:3:begin
\textbf{nęrian}\par\nobreak
% course-caption:T01:3:end
% course-panel:T01-P04:begin
\par\begingroup
\fontsize{10.5}{13.2}\selectfont
\setlength{\LTleft}{0pt}\setlength{\LTright}{\fill}
\noindent
\begin{tabular}{@{}>{\raggedright\arraybackslash}p{39.459mm}>{\raggedright\arraybackslash}p{39.459mm}>{\raggedright\arraybackslash}p{39.459mm}@{}}
\toprule
\textbf{Person and number} & \textbf{Present} & \textbf{Preterite} \\
\midrule
1–3 sg. & nęrie & nęrede \\
1–3 pl. & nęrien & nęreden \\
\bottomrule
\end{tabular}
\par\endgroup
% course-panel:T01-P04:end
% mdtable:T01:end
\reserve{7\baselineskip}
% mdparagraph:P022:begin
\subsection*{Weak classes II and III}\BookHeading{nav-ch07-learner-0004}{subsubsection}{Weak classes II and III}{0}\BookLabel{bk-ch07-learner-eb3e06d33ada}
% mdparagraph:P022:end

% teaching: NEW-U07-SUBJ-lufian
% teaching: NEW-U07-SUBJ-gebletsian
% teaching: NEW-U07-SUBJ-gehalgian
% teaching: NEW-U07-SUBJ-habban
% teaching: NEW-U07-SUBJ-libban
% teaching: NEW-U07-SUBJ-secgan
% teaching: NEW-U07-SUBJ-hycgan
% mdparagraph:P023:begin
The first three verbs are class II: present -ie/\allowbreak{}-ien, preterite -ode/\allowbreak{}-oden. The last four are class III and require their own stems. Habban has hæbb- in the present subjunctive; its indicative plural does not supply that vowel. Libban has lifd- in the past; sęcgan has sægd-; hycgan has hogod-.
% mdparagraph:P023:end

% table: U07-DISPLAY-SUBJ-2
% mdtable:T02:begin
% course-table-style: 2026-10
% course-panel:T02-P01:begin
\par\begingroup
\fontsize{10.5}{13.2}\selectfont
\setlength{\LTleft}{0pt}\setlength{\LTright}{\fill}
\noindent
\begin{tabular}{@{}>{\raggedright\arraybackslash}p{18.849mm}>{\raggedright\arraybackslash}p{19.022mm}>{\raggedright\arraybackslash}p{22.464mm}>{\raggedright\arraybackslash}p{24.96mm}>{\raggedright\arraybackslash}p{27.456mm}@{}}
\toprule
\textbf{Person and number} & \textbf{lufian · Present} & \textbf{lufian · Preterite} & \textbf{geblētsian · Present} & \textbf{geblētsian · Preterite} \\
\midrule
1–3 sg. & lufie & lufode & geblētsie & geblētsode \\
1–3 pl. & lufien & lufoden & geblētsien & geblētsoden \\
\bottomrule
\end{tabular}
\par\endgroup
% course-panel:T02-P01:end
\par\addvspace{5pt}
% course-panel:T02-P02:begin
\par\begingroup
\fontsize{10.5}{13.2}\selectfont
\setlength{\LTleft}{0pt}\setlength{\LTright}{\fill}
\noindent
\begin{tabular}{@{}>{\raggedright\arraybackslash}p{19.443mm}>{\raggedright\arraybackslash}p{23.971mm}>{\raggedright\arraybackslash}p{26.546mm}>{\raggedright\arraybackslash}p{19.621mm}>{\raggedright\arraybackslash}p{23.172mm}@{}}
\toprule
\textbf{Person and number} & \textbf{gehālgian · Present} & \textbf{gehālgian · Preterite} & \textbf{habban · Present} & \textbf{habban · Preterite} \\
\midrule
1–3 sg. & gehālgie & gehālgode & hæbbe & hæfde \\
1–3 pl. & gehālgien & gehālgoden & hæbben & hæfden \\
\bottomrule
\end{tabular}
\par\endgroup
% course-panel:T02-P02:end
\par\addvspace{5pt}
% course-panel:T02-P03:begin
\par\begingroup
\fontsize{10.5}{13.2}\selectfont
\setlength{\LTleft}{0pt}\setlength{\LTright}{\fill}
\noindent
\begin{tabular}{@{}>{\raggedright\arraybackslash}p{20.873mm}>{\raggedright\arraybackslash}p{21.064mm}>{\raggedright\arraybackslash}p{24.876mm}>{\raggedright\arraybackslash}p{21.064mm}>{\raggedright\arraybackslash}p{24.876mm}@{}}
\toprule
\textbf{Person and number} & \textbf{libban · Present} & \textbf{libban · Preterite} & \textbf{sęcgan · Present} & \textbf{sęcgan · Preterite} \\
\midrule
1–3 sg. & libbe & lifde & sęcge & sægde \\
1–3 pl. & libben & lifden & sęcgen & sægden \\
\bottomrule
\end{tabular}
\par\endgroup
% course-panel:T02-P03:end
\par\addvspace{5pt}
\par\addvspace{3pt}
% course-caption:T02:3:begin
\textbf{hycgan}\par\nobreak
% course-caption:T02:3:end
% course-panel:T02-P04:begin
\par\begingroup
\fontsize{10.5}{13.2}\selectfont
\setlength{\LTleft}{0pt}\setlength{\LTright}{\fill}
\noindent
\begin{tabular}{@{}>{\raggedright\arraybackslash}p{39.459mm}>{\raggedright\arraybackslash}p{39.459mm}>{\raggedright\arraybackslash}p{39.459mm}@{}}
\toprule
\textbf{Person and number} & \textbf{Present} & \textbf{Preterite} \\
\midrule
1–3 sg. & hycge & hogode \\
1–3 pl. & hycgen & hogoden \\
\bottomrule
\end{tabular}
\par\endgroup
% course-panel:T02-P04:end
% mdtable:T02:end
% mdparagraph:P024:begin
\textbf{Alternatives.} Libban also has present lifie/\allowbreak{}lifien; sęcgan has past sāde/\allowbreak{}sāden. The shorter past of hycgan beside the selected hogode-series is also documented. These alternatives do not replace the selected table forms in a task that expressly asks for them.
% mdparagraph:P024:end

\reserve{7\baselineskip}
% mdparagraph:P025:begin
\subsection*{Strong classes I–III}\BookHeading{nav-ch07-learner-0005}{subsubsection}{Strong classes I–III}{0}\BookLabel{bk-ch07-learner-2bd176407f8c}
% mdparagraph:P025:end

% teaching: NEW-U07-SUBJ-drifan
% teaching: NEW-U07-SUBJ-wreon
% teaching: NEW-U07-SUBJ-beodan
% teaching: NEW-U07-SUBJ-teon
% teaching: NEW-U07-SUBJ-fleogan
% teaching: NEW-U07-SUBJ-bindan
% teaching: NEW-U07-SUBJ-helpan
% teaching: NEW-U07-SUBJ-weorthan
% teaching: NEW-U07-SUBJ-swimman
% mdparagraph:P026:begin
Use the familiar plural-preterite stems, not the first-singular past stems. Bindan therefore has bunde/\allowbreak{}bunden, and helpan has hulpe/\allowbreak{}hulpen. In drīfan the present ī and past i remain distinct. Weorðan changes ð to d in the selected past. The contracted wrēon and tēon have singular wrēo and tēo, without an added -e; their past stems contain g.
% mdparagraph:P026:end

% table: U07-DISPLAY-SUBJ-3
% mdtable:T03:begin
% course-table-style: 2026-10
% course-panel:T03-P01:begin
\par\begingroup
\fontsize{10.5}{13.2}\selectfont
\setlength{\LTleft}{0pt}\setlength{\LTright}{\fill}
\noindent
\begin{tabular}{@{}>{\raggedright\arraybackslash}p{21.195mm}>{\raggedright\arraybackslash}p{21.389mm}>{\raggedright\arraybackslash}p{25.26mm}>{\raggedright\arraybackslash}p{20.518mm}>{\raggedright\arraybackslash}p{24.389mm}@{}}
\toprule
\textbf{Person and number} & \textbf{drīfan · Present} & \textbf{drīfan · Preterite} & \textbf{wrēon · Present} & \textbf{wrēon · Preterite} \\
\midrule
1–3 sg. & drīfe & drife & wrēo & wrige \\
1–3 pl. & drīfen & drifen & wrēon & wrigen \\
\bottomrule
\end{tabular}
\par\endgroup
% course-panel:T03-P01:end
\par\addvspace{5pt}
% course-panel:T03-P02:begin
\par\begingroup
\fontsize{10.5}{13.2}\selectfont
\setlength{\LTleft}{0pt}\setlength{\LTright}{\fill}
\noindent
\begin{tabular}{@{}>{\raggedright\arraybackslash}p{21.528mm}>{\raggedright\arraybackslash}p{21.725mm}>{\raggedright\arraybackslash}p{25.657mm}>{\raggedright\arraybackslash}p{19.955mm}>{\raggedright\arraybackslash}p{23.887mm}@{}}
\toprule
\textbf{Person and number} & \textbf{bēodan · Present} & \textbf{bēodan · Preterite} & \textbf{tēon · Present} & \textbf{tēon · Preterite} \\
\midrule
1–3 sg. & bēode & bude & tēo & tuge \\
1–3 pl. & bēoden & buden & tēon & tugen \\
\bottomrule
\end{tabular}
\par\endgroup
% course-panel:T03-P02:end
\par\addvspace{5pt}
% course-panel:T03-P03:begin
\par\begingroup
\fontsize{10.5}{13.2}\selectfont
\setlength{\LTleft}{0pt}\setlength{\LTright}{\fill}
\noindent
\begin{tabular}{@{}>{\raggedright\arraybackslash}p{20.56mm}>{\raggedright\arraybackslash}p{21.593mm}>{\raggedright\arraybackslash}p{25.348mm}>{\raggedright\arraybackslash}p{20.748mm}>{\raggedright\arraybackslash}p{24.503mm}@{}}
\toprule
\textbf{Person and number} & \textbf{flēogan · Present} & \textbf{flēogan · Preterite} & \textbf{bindan · Present} & \textbf{bindan · Preterite} \\
\midrule
1–3 sg. & flēoge & fluge & binde & bunde \\
1–3 pl. & flēogen & flugen & binden & bunden \\
\bottomrule
\end{tabular}
\par\endgroup
% course-panel:T03-P03:end
\par\addvspace{5pt}
% course-panel:T03-P04:begin
\par\begingroup
\fontsize{10.5}{13.2}\selectfont
\setlength{\LTleft}{0pt}\setlength{\LTright}{\fill}
\noindent
\begin{tabular}{@{}>{\raggedright\arraybackslash}p{20.56mm}>{\raggedright\arraybackslash}p{20.748mm}>{\raggedright\arraybackslash}p{24.503mm}>{\raggedright\arraybackslash}p{21.593mm}>{\raggedright\arraybackslash}p{25.348mm}@{}}
\toprule
\textbf{Person and number} & \textbf{helpan · Present} & \textbf{helpan · Preterite} & \textbf{weorðan · Present} & \textbf{weorðan · Preterite} \\
\midrule
1–3 sg. & helpe & hulpe & weorðe & wurde \\
1–3 pl. & helpen & hulpen & weorðen & wurden \\
\bottomrule
\end{tabular}
\par\endgroup
% course-panel:T03-P04:end
\par\addvspace{5pt}
\par\addvspace{3pt}
% course-caption:T03:4:begin
\textbf{swimman}\par\nobreak
% course-caption:T03:4:end
% course-panel:T03-P05:begin
\par\begingroup
\fontsize{10.5}{13.2}\selectfont
\setlength{\LTleft}{0pt}\setlength{\LTright}{\fill}
\noindent
\begin{tabular}{@{}>{\raggedright\arraybackslash}p{39.459mm}>{\raggedright\arraybackslash}p{39.459mm}>{\raggedright\arraybackslash}p{39.459mm}@{}}
\toprule
\textbf{Person and number} & \textbf{Present} & \textbf{Preterite} \\
\midrule
1–3 sg. & swimme & swumme \\
1–3 pl. & swimmen & swummen \\
\bottomrule
\end{tabular}
\par\endgroup
% course-panel:T03-P05:end
% mdtable:T03:end
% mdparagraph:P027:begin
Tēon here is the class II verb “draw, pull”. The same spelling can name another verb; it is the learned lexical system that determines the past stem.
% mdparagraph:P027:end

\reserve{7\baselineskip}
% mdparagraph:P028:begin
\subsection*{Strong classes IV–VII}\BookHeading{nav-ch07-learner-0006}{subsubsection}{Strong classes IV–VII}{0}\BookLabel{bk-ch07-learner-2c917f3ea037}
% mdparagraph:P028:end

% teaching: NEW-U07-SUBJ-beran
% teaching: NEW-U07-SUBJ-niman
% teaching: NEW-U07-SUBJ-cuman
% teaching: NEW-U07-SUBJ-giefan
% teaching: NEW-U07-SUBJ-cwethan
% teaching: NEW-U07-SUBJ-biddan
% teaching: NEW-U07-SUBJ-licgan
% teaching: NEW-U07-SUBJ-sittan
% teaching: NEW-U07-SUBJ-seon
% teaching: NEW-U07-SUBJ-faran
% teaching: NEW-U07-SUBJ-standan
% teaching: NEW-U07-SUBJ-slean
% teaching: NEW-U07-SUBJ-hebban
% teaching: NEW-U07-SUBJ-scieppan
% teaching: NEW-U07-SUBJ-swerian
% teaching: NEW-U07-SUBJ-weaxan
% teaching: NEW-U07-SUBJ-hatan
% teaching: NEW-U07-SUBJ-fon
% mdparagraph:P029:begin
Retain the selected past-plural vowels of niman and cuman. Cweðan has d in cwǣde/\allowbreak{}cwǣden. The j-present verbs retain their characteristic present consonants: dd, cg, tt, bb, and pp. Swęrian retains i in the present ending. Standan loses its present n in stōde/\allowbreak{}stōden. The contracted sēon, slēan, and fōn require their own present and past forms.
% mdparagraph:P029:end

% table: U07-DISPLAY-SUBJ-4
% mdtable:T04:begin
% course-table-style: 2026-10
% course-panel:T04-P01:begin
\par\begingroup
\fontsize{10.5}{13.2}\selectfont
\setlength{\LTleft}{0pt}\setlength{\LTright}{\fill}
\noindent
\begin{tabular}{@{}>{\raggedright\arraybackslash}p{21.528mm}>{\raggedright\arraybackslash}p{20.84mm}>{\raggedright\arraybackslash}p{24.772mm}>{\raggedright\arraybackslash}p{20.84mm}>{\raggedright\arraybackslash}p{24.772mm}@{}}
\toprule
\textbf{Person and number} & \textbf{beran · Present} & \textbf{beran · Preterite} & \textbf{niman · Present} & \textbf{niman · Preterite} \\
\midrule
1–3 sg. & bere & bǣre & nime & nōme \\
1–3 pl. & beren & bǣren & nimen & nōmen \\
\bottomrule
\end{tabular}
\par\endgroup
% course-panel:T04-P01:end
\par\addvspace{5pt}
% course-panel:T04-P02:begin
\par\begingroup
\fontsize{10.5}{13.2}\selectfont
\setlength{\LTleft}{0pt}\setlength{\LTright}{\fill}
\noindent
\begin{tabular}{@{}>{\raggedright\arraybackslash}p{21.195mm}>{\raggedright\arraybackslash}p{20.518mm}>{\raggedright\arraybackslash}p{24.389mm}>{\raggedright\arraybackslash}p{21.389mm}>{\raggedright\arraybackslash}p{25.26mm}@{}}
\toprule
\textbf{Person and number} & \textbf{cuman · Present} & \textbf{cuman · Preterite} & \textbf{giefan · Present} & \textbf{giefan · Preterite} \\
\midrule
1–3 sg. & cume & cōme & giefe & gēafe \\
1–3 pl. & cumen & cōmen & giefen & gēafen \\
\bottomrule
\end{tabular}
\par\endgroup
% course-panel:T04-P02:end
\par\addvspace{5pt}
% course-panel:T04-P03:begin
\par\begingroup
\fontsize{10.5}{13.2}\selectfont
\setlength{\LTleft}{0pt}\setlength{\LTright}{\fill}
\noindent
\begin{tabular}{@{}>{\raggedright\arraybackslash}p{20.873mm}>{\raggedright\arraybackslash}p{21.064mm}>{\raggedright\arraybackslash}p{24.876mm}>{\raggedright\arraybackslash}p{21.064mm}>{\raggedright\arraybackslash}p{24.876mm}@{}}
\toprule
\textbf{Person and number} & \textbf{cweðan · Present} & \textbf{cweðan · Preterite} & \textbf{biddan · Present} & \textbf{biddan · Preterite} \\
\midrule
1–3 sg. & cweðe & cwǣde & bidde & bǣde \\
1–3 pl. & cweðen & cwǣden & bidden & bǣden \\
\bottomrule
\end{tabular}
\par\endgroup
% course-panel:T04-P03:end
\par\addvspace{5pt}
% course-panel:T04-P04:begin
\par\begingroup
\fontsize{10.5}{13.2}\selectfont
\setlength{\LTleft}{0pt}\setlength{\LTright}{\fill}
\noindent
\begin{tabular}{@{}>{\raggedright\arraybackslash}p{20.873mm}>{\raggedright\arraybackslash}p{21.064mm}>{\raggedright\arraybackslash}p{24.876mm}>{\raggedright\arraybackslash}p{21.064mm}>{\raggedright\arraybackslash}p{24.876mm}@{}}
\toprule
\textbf{Person and number} & \textbf{licgan · Present} & \textbf{licgan · Preterite} & \textbf{sittan · Present} & \textbf{sittan · Preterite} \\
\midrule
1–3 sg. & licge & lǣge & sitte & sǣte \\
1–3 pl. & licgen & lǣgen & sitten & sǣten \\
\bottomrule
\end{tabular}
\par\endgroup
% course-panel:T04-P04:end
\par\addvspace{5pt}
% course-panel:T04-P05:begin
\par\begingroup
\fontsize{10.5}{13.2}\selectfont
\setlength{\LTleft}{0pt}\setlength{\LTright}{\fill}
\noindent
\begin{tabular}{@{}>{\raggedright\arraybackslash}p{21.871mm}>{\raggedright\arraybackslash}p{20.273mm}>{\raggedright\arraybackslash}p{24.268mm}>{\raggedright\arraybackslash}p{21.172mm}>{\raggedright\arraybackslash}p{25.167mm}@{}}
\toprule
\textbf{Person and number} & \textbf{sēon · Present} & \textbf{sēon · Preterite} & \textbf{faran · Present} & \textbf{faran · Preterite} \\
\midrule
1–3 sg. & sēo & sāwe & fare & fōre \\
1–3 pl. & sēon & sāwen & faren & fōren \\
\bottomrule
\end{tabular}
\par\endgroup
% course-panel:T04-P05:end
\par\addvspace{5pt}
% course-panel:T04-P06:begin
\par\begingroup
\fontsize{10.5}{13.2}\selectfont
\setlength{\LTleft}{0pt}\setlength{\LTright}{\fill}
\noindent
\begin{tabular}{@{}>{\raggedright\arraybackslash}p{20.873mm}>{\raggedright\arraybackslash}p{21.921mm}>{\raggedright\arraybackslash}p{25.734mm}>{\raggedright\arraybackslash}p{20.206mm}>{\raggedright\arraybackslash}p{24.018mm}@{}}
\toprule
\textbf{Person and number} & \textbf{standan · Present} & \textbf{standan · Preterite} & \textbf{slēan · Present} & \textbf{slēan · Preterite} \\
\midrule
1–3 sg. & stande & stōde & slēa & slōge \\
1–3 pl. & standen & stōden & slēan & slōgen \\
\bottomrule
\end{tabular}
\par\endgroup
% course-panel:T04-P06:end
\par\addvspace{5pt}
% course-panel:T04-P07:begin
\par\begingroup
\fontsize{10.5}{13.2}\selectfont
\setlength{\LTleft}{0pt}\setlength{\LTright}{\fill}
\noindent
\begin{tabular}{@{}>{\raggedright\arraybackslash}p{20.075mm}>{\raggedright\arraybackslash}p{20.259mm}>{\raggedright\arraybackslash}p{23.925mm}>{\raggedright\arraybackslash}p{22.917mm}>{\raggedright\arraybackslash}p{25.575mm}@{}}
\toprule
\textbf{Person and number} & \textbf{hębban · Present} & \textbf{hębban · Preterite} & \textbf{scieppan · Present} & \textbf{scieppan · Preterite} \\
\midrule
1–3 sg. & hębbe & hōfe & scieppe & scēope \\
1–3 pl. & hębben & hōfen & scieppen & scēopen \\
\bottomrule
\end{tabular}
\par\endgroup
% course-panel:T04-P07:end
\par\addvspace{5pt}
% course-panel:T04-P08:begin
\par\begingroup
\fontsize{10.5}{13.2}\selectfont
\setlength{\LTleft}{0pt}\setlength{\LTright}{\fill}
\noindent
\begin{tabular}{@{}>{\raggedright\arraybackslash}p{20.56mm}>{\raggedright\arraybackslash}p{21.593mm}>{\raggedright\arraybackslash}p{25.348mm}>{\raggedright\arraybackslash}p{20.748mm}>{\raggedright\arraybackslash}p{24.503mm}@{}}
\toprule
\textbf{Person and number} & \textbf{swęrian · Present} & \textbf{swęrian · Preterite} & \textbf{weaxan · Present} & \textbf{weaxan · Preterite} \\
\midrule
1–3 sg. & swęrie & swōre & weaxe & wēoxe \\
1–3 pl. & swęrien & swōren & weaxen & wēoxen \\
\bottomrule
\end{tabular}
\par\endgroup
% course-panel:T04-P08:end
\par\addvspace{5pt}
% course-panel:T04-P09:begin
\par\begingroup
\fontsize{10.5}{13.2}\selectfont
\setlength{\LTleft}{0pt}\setlength{\LTright}{\fill}
\noindent
\begin{tabular}{@{}>{\raggedright\arraybackslash}p{22.106mm}>{\raggedright\arraybackslash}p{21.4mm}>{\raggedright\arraybackslash}p{25.437mm}>{\raggedright\arraybackslash}p{20.188mm}>{\raggedright\arraybackslash}p{23.62mm}@{}}
\toprule
\textbf{Person and number} & \textbf{hātan · Present} & \textbf{hātan · Preterite} & \textbf{fōn · Present} & \textbf{fōn · Preterite} \\
\midrule
1–3 sg. & hāte & hēte & fō & fēnge \\
1–3 pl. & hāten & hēten & fōn & fēngen \\
\bottomrule
\end{tabular}
\par\endgroup
% course-panel:T04-P09:end
% mdtable:T04:end
% mdparagraph:P030:begin
Cuman also has the documented present cyme/\allowbreak{}cymen; the selected forms are cume/\allowbreak{}cumen. Keep scieppan’s selected scēop- past. Slēan retains the ordinary meaning “strike” learned earlier. Hāten is present plural subjunctive in its finite paradigm; the same spelling will occur as a separately identified past participle. Spelling alone cannot establish its function.
% mdparagraph:P030:end

\reserve{7\baselineskip}
% mdparagraph:P031:begin
\subsection*{Preterite-presents and anomalous verbs}\BookHeading{nav-ch07-learner-0007}{subsubsection}{Preterite-presents and anomalous verbs}{0}\BookLabel{bk-ch07-learner-a651adcc0fcb}
% mdparagraph:P031:end

% teaching: NEW-U07-SUBJ-witan
% teaching: NEW-U07-SUBJ-cunnan
% teaching: NEW-U07-SUBJ-thurfan
% teaching: NEW-U07-SUBJ-durran
% teaching: NEW-U07-SUBJ-sculan
% teaching: NEW-U07-SUBJ-magan
% teaching: NEW-U07-SUBJ-motan
% teaching: NEW-U07-SUBJ-willan
% teaching: NEW-U07-SUBJ-don
% teaching: NEW-U07-SUBJ-gan
% teaching: NEW-U07-SUBJ-wesan
% mdparagraph:P032:begin
The preterite-present verbs form their present subjunctives from the selected plural-type or mutated stems, not from a strong-looking indicative singular. Witan has wit-; ðurfan, durran, and sculan have the selected y forms; mugan has mæge. Their dental preterites take the ordinary singular/\allowbreak{}plural subjunctive contrast.
% mdparagraph:P032:end

% mdparagraph:P033:begin
Dōn and gān retain contracted dō/\allowbreak{}dōn and gā/\allowbreak{}gān in the present. Their preterites use dyde and ēode. Willan has wille/\allowbreak{}willen and wolde/\allowbreak{}wolden. Wesan has sīe/\allowbreak{}sīen and wǣre/\allowbreak{}wǣren.
% mdparagraph:P033:end

% table: U07-DISPLAY-SUBJ-5
% mdtable:T05:begin
% course-table-style: 2026-10
% course-panel:T05-P01:begin
\par\begingroup
\fontsize{10.5}{13.2}\selectfont
\setlength{\LTleft}{0pt}\setlength{\LTright}{\fill}
\noindent
\begin{tabular}{@{}>{\raggedright\arraybackslash}p{21.195mm}>{\raggedright\arraybackslash}p{20.518mm}>{\raggedright\arraybackslash}p{24.389mm}>{\raggedright\arraybackslash}p{21.389mm}>{\raggedright\arraybackslash}p{25.26mm}@{}}
\toprule
\textbf{Person and number} & \textbf{witan · Present} & \textbf{witan · Preterite} & \textbf{cunnan · Present} & \textbf{cunnan · Preterite} \\
\midrule
1–3 sg. & wite & wiste & cunne & cūðe \\
1–3 pl. & witen & wisten & cunnen & cūðen \\
\bottomrule
\end{tabular}
\par\endgroup
% course-panel:T05-P01:end
\par\addvspace{5pt}
% course-panel:T05-P02:begin
\par\begingroup
\fontsize{10.5}{13.2}\selectfont
\setlength{\LTleft}{0pt}\setlength{\LTright}{\fill}
\noindent
\begin{tabular}{@{}>{\raggedright\arraybackslash}p{20.873mm}>{\raggedright\arraybackslash}p{21.064mm}>{\raggedright\arraybackslash}p{24.876mm}>{\raggedright\arraybackslash}p{21.064mm}>{\raggedright\arraybackslash}p{24.876mm}@{}}
\toprule
\textbf{Person and number} & \textbf{ðurfan · Present} & \textbf{ðurfan · Preterite} & \textbf{durran · Present} & \textbf{durran · Preterite} \\
\midrule
1–3 sg. & ðyrfe & ðorfte & dyrre & dorste \\
1–3 pl. & ðyrfen & ðorften & dyrren & dorsten \\
\bottomrule
\end{tabular}
\par\endgroup
% course-panel:T05-P02:end
\par\addvspace{5pt}
% course-panel:T05-P03:begin
\par\begingroup
\fontsize{10.5}{13.2}\selectfont
\setlength{\LTleft}{0pt}\setlength{\LTright}{\fill}
\noindent
\begin{tabular}{@{}>{\raggedright\arraybackslash}p{21.195mm}>{\raggedright\arraybackslash}p{21.389mm}>{\raggedright\arraybackslash}p{25.26mm}>{\raggedright\arraybackslash}p{20.518mm}>{\raggedright\arraybackslash}p{24.389mm}@{}}
\toprule
\textbf{Person and number} & \textbf{sculan · Present} & \textbf{sculan · Preterite} & \textbf{mugan · Present} & \textbf{mugan · Preterite} \\
\midrule
1–3 sg. & scyle & sceolde & mæge & meahte \\
1–3 pl. & scylen & sceolden & mægen & meahten \\
\bottomrule
\end{tabular}
\par\endgroup
% course-panel:T05-P03:end
\par\addvspace{5pt}
% course-panel:T05-P04:begin
\par\begingroup
\fontsize{10.5}{13.2}\selectfont
\setlength{\LTleft}{0pt}\setlength{\LTright}{\fill}
\noindent
\begin{tabular}{@{}>{\raggedright\arraybackslash}p{21.195mm}>{\raggedright\arraybackslash}p{20.518mm}>{\raggedright\arraybackslash}p{24.389mm}>{\raggedright\arraybackslash}p{21.389mm}>{\raggedright\arraybackslash}p{25.26mm}@{}}
\toprule
\textbf{Person and number} & \textbf{mōtan · Present} & \textbf{mōtan · Preterite} & \textbf{willan · Present} & \textbf{willan · Preterite} \\
\midrule
1–3 sg. & mōte & mōste & wille & wolde \\
1–3 pl. & mōten & mōsten & willen & wolden \\
\bottomrule
\end{tabular}
\par\endgroup
% course-panel:T05-P04:end
\par\addvspace{5pt}
% course-panel:T05-P05:begin
\par\begingroup
\fontsize{10.5}{13.2}\selectfont
\setlength{\LTleft}{0pt}\setlength{\LTright}{\fill}
\noindent
\begin{tabular}{@{}>{\raggedright\arraybackslash}p{22.716mm}>{\raggedright\arraybackslash}p{20.746mm}>{\raggedright\arraybackslash}p{24.272mm}>{\raggedright\arraybackslash}p{20.746mm}>{\raggedright\arraybackslash}p{24.272mm}@{}}
\toprule
\textbf{Person and number} & \textbf{dōn · Present} & \textbf{dōn · Preterite} & \textbf{gān · Present} & \textbf{gān · Preterite} \\
\midrule
1–3 sg. & dō & dyde & gā & ēode \\
1–3 pl. & dōn & dyden & gān & ēoden \\
\bottomrule
\end{tabular}
\par\endgroup
% course-panel:T05-P05:end
\par\addvspace{5pt}
\par\addvspace{3pt}
% course-caption:T05:5:begin
\textbf{wesan}\par\nobreak
% course-caption:T05:5:end
% course-panel:T05-P06:begin
\par\begingroup
\fontsize{10.5}{13.2}\selectfont
\setlength{\LTleft}{0pt}\setlength{\LTright}{\fill}
\noindent
\begin{tabular}{@{}>{\raggedright\arraybackslash}p{39.459mm}>{\raggedright\arraybackslash}p{39.459mm}>{\raggedright\arraybackslash}p{39.459mm}@{}}
\toprule
\textbf{Person and number} & \textbf{Present} & \textbf{Preterite} \\
\midrule
1–3 sg. & sīe & wǣre \\
1–3 pl. & sīen & wǣren \\
\bottomrule
\end{tabular}
\par\endgroup
% course-panel:T05-P06:end
% mdtable:T05:end
% mdparagraph:P034:begin
\textbf{Alternative present of “be”.} Each row again covers all three persons. The bēo series shares the same preterite wǣre/\allowbreak{}wǣren; it has no separate past paradigm.
% mdparagraph:P034:end

% table: U07-DISPLAY-BEO-SUBJ
% mdtable:T06:begin
% course-table-style: 2026-10
% course-panel:T06-P01:begin
\par\begingroup
\fontsize{10.5}{13.2}\selectfont
\setlength{\LTleft}{0pt}\setlength{\LTright}{\fill}
\noindent
\begin{tabular}{@{}>{\raggedright\arraybackslash}p{63.342mm}>{\raggedright\arraybackslash}p{57.846mm}@{}}
\toprule
\textbf{Person and number} & \textbf{Present} \\
\midrule
1–3 sg. & bēo \\
1–3 pl. & bēon \\
\bottomrule
\end{tabular}
\par\endgroup
% course-panel:T06-P01:end
% mdtable:T06:end
% mdparagraph:P035:begin
\textbf{Alternatives and limits.} Witan has the alternative wisse/\allowbreak{}wissen preterite; cunnan has a cȳðe variant. Beside the selected ðyrfe/\allowbreak{}ðyrfen, dyrre/\allowbreak{}dyrren, and scyle/\allowbreak{}scylen, the corresponding forms with u are documented. The tables continue to select meahte/\allowbreak{}meahten for mugan. Learn the selected series together; a spelling variant is not a new mood.
% mdparagraph:P035:end

% mdparagraph:P036:begin
The labels mugan (also called magan), durran, sculan, and mōtan identify finite systems. Their presence in the table headings does not supply a usable infinitive, imperative, or participle. The new subjunctives do not fill every missing nonfinite or command form. Bēo itself can also be first-singular indicative or second-singular imperative; bēon can be an infinitive. A clause’s subject and purpose matter.
% mdparagraph:P036:end

\reserve{7\baselineskip}
% mdparagraph:P037:begin
\section*{3. Five additional verb systems}\BookHeading{nav-ch07-learner-0008}{subsection}{3. Five additional verb systems}{1}\BookLabel{bk-ch07-learner-4a4ffd61ee9a}
% mdparagraph:P037:end

% teaching: NEW-U07-VOC-wyrcean-yield
% teaching: NEW-U07-VOC-gieman-1
% teaching: NEW-U07-VOC-ricsian-dominion
% teaching: NEW-U07-VOC-gemanigfieldan-multiply
% teaching: NEW-U07-VOC-totwaeman-separate
% teaching: NEW-U07-GRAM-U07-GIEMAN-GENITIVE-COMPLEMENT
% teaching: NEW-U07-GRAM-U07-RICSIAN-INTRANSITIVE
% teaching: NEW-U07-GRAM-U07-WYRCEAN-ACCUSATIVE-FRUIT
% teaching: NEW-U07-GRAM-U07-TOTWAEMAN-ACCUSATIVE-OBJECT-FRAM-DATIVE
% teaching: NEW-U07-GRAM-U07-GEMANIGFIELDAN-PASSIVE-NOMINATIVE-SUBJECT
% mdparagraph:P038:begin
These verbs use familiar weak patterns, but their meanings and complements must be learned individually. The five headwords identify the complete finite systems printed here. Only tōtwǣman adds a new imperative pair; the participle assignment is stated in §5.
% mdparagraph:P038:end

% mdtable:T07:begin
\par\begingroup
\fontsize{10.5}{13.2}\selectfont
\setlength{\LTleft}{0pt}\setlength{\LTright}{\fill}
\noindent
\begin{tabular}{@{}>{\raggedright\arraybackslash}p{26mm}>{\raggedright\arraybackslash}p{23mm}>{\raggedright\arraybackslash}p{69.37mm}@{}}
\toprule
\textbf{Verb} & \textbf{Class} & \textbf{Meaning and construction} \\
\midrule
wyrcean & weak I, irregular & yield fruit; the fruit is accusative \\
gīeman & weak I & rule over; the governed person or thing is genitive \\
rīcsian & weak II & have dominion, bear rule; used intransitively \\
gemanigfieldan & weak I, dental stem & multiply; this unit uses its passive “be multiplied” with the beings as subject \\
tōtwǣman & weak I & divide, separate; the thing divided is accusative, and what it is divided from follows fram in the dative \\
\bottomrule
\end{tabular}
\par\endgroup

% mdtable:T07:end
% mdparagraph:P039:begin
Wyrcean is taught here as its own verb. Knowing prefixed gewyrcean does not by itself establish an unprefixed meaning. The government of gīeman differs from an ordinary accusative object; rīcsian needs no object. Neither construction is an invitation to apply the accusative-object passive rule to it.
% mdparagraph:P039:end

% teaching: NEW-U07-GRAM-U07-FIVE-NEW-INDICATIVES
\reserve{7\baselineskip}
% mdparagraph:P040:begin
\subsection*{Present and preterite indicatives}\BookHeading{nav-ch07-learner-0009}{subsubsection}{Present and preterite indicatives}{0}
% mdparagraph:P040:end

% mdparagraph:P041:begin
The tables give each of the six person–number forms in both tenses. In the weak preterite, first and third singular coincide; second singular adds -st, and the plural has -on. Preserve the lexical stem and the dental changes shown.
% mdparagraph:P041:end

\reserve{5\baselineskip}
% mdparagraph:P042:begin
\textbf{Wyrcean.}
% mdparagraph:P042:end

% teaching: NEW-U07-IND-wyrcean
% mdparagraph:P043:begin
The present has wyrce-, with wyrcst and wyrcð in the second and third singular. The irregular past is worht-, not a regular addition to wyrc-.
% mdparagraph:P043:end

% table: NEW-U07-IND-wyrcean
% mdtable:T08:begin
% course-table-style: 2026-10
% course-panel:T08-P01:begin
\par\begingroup
\fontsize{10.5}{13.2}\selectfont
\setlength{\LTleft}{0pt}\setlength{\LTright}{\fill}
\noindent
\begin{tabular}{@{}>{\raggedright\arraybackslash}p{39.459mm}>{\raggedright\arraybackslash}p{39.459mm}>{\raggedright\arraybackslash}p{39.459mm}@{}}
\toprule
\textbf{Person and number} & \textbf{Present indicative} & \textbf{Preterite indicative} \\
\midrule
1 sg. & wyrce & worhte \\
2 sg. & wyrcst & worhtest \\
3 sg. & wyrcð & worhte \\
1–3 pl. & wyrceað & worhton \\
\bottomrule
\end{tabular}
\par\endgroup
% course-panel:T08-P01:end
% mdtable:T08:end
\reserve{5\baselineskip}
% mdparagraph:P044:begin
\textbf{Gīeman.}
% mdparagraph:P044:end

% teaching: NEW-U07-IND-gieman
% mdparagraph:P045:begin
The long stem has gīemst, gīemð, and plural gīemað; the preterite retains the dental stem gīemd-.
% mdparagraph:P045:end

% table: NEW-U07-IND-gieman
% mdtable:T09:begin
% course-table-style: 2026-10
% course-panel:T09-P01:begin
\par\begingroup
\fontsize{10.5}{13.2}\selectfont
\setlength{\LTleft}{0pt}\setlength{\LTright}{\fill}
\noindent
\begin{tabular}{@{}>{\raggedright\arraybackslash}p{39.459mm}>{\raggedright\arraybackslash}p{39.459mm}>{\raggedright\arraybackslash}p{39.459mm}@{}}
\toprule
\textbf{Person and number} & \textbf{Present indicative} & \textbf{Preterite indicative} \\
\midrule
1 sg. & gīeme & gīemde \\
2 sg. & gīemst & gīemdest \\
3 sg. & gīemð & gīemde \\
1–3 pl. & gīemað & gīemdon \\
\bottomrule
\end{tabular}
\par\endgroup
% course-panel:T09-P01:end
% mdtable:T09:end
\reserve{5\baselineskip}
% mdparagraph:P046:begin
\textbf{Rīcsian.}
% mdparagraph:P046:end

% teaching: NEW-U07-IND-ricsian
% mdparagraph:P047:begin
Class II contrasts rīcsie and plural rīcsiað with rīcsast/\allowbreak{}rīcsað. The preterite has -od-.
% mdparagraph:P047:end

% table: NEW-U07-IND-ricsian
% mdtable:T10:begin
% course-table-style: 2026-10
% course-panel:T10-P01:begin
\par\begingroup
\fontsize{10.5}{13.2}\selectfont
\setlength{\LTleft}{0pt}\setlength{\LTright}{\fill}
\noindent
\begin{tabular}{@{}>{\raggedright\arraybackslash}p{39.459mm}>{\raggedright\arraybackslash}p{39.459mm}>{\raggedright\arraybackslash}p{39.459mm}@{}}
\toprule
\textbf{Person and number} & \textbf{Present indicative} & \textbf{Preterite indicative} \\
\midrule
1 sg. & rīcsie & rīcsode \\
2 sg. & rīcsast & rīcsodest \\
3 sg. & rīcsað & rīcsode \\
1–3 pl. & rīcsiað & rīcsodon \\
\bottomrule
\end{tabular}
\par\endgroup
% course-panel:T10-P01:end
% mdtable:T10:end
\reserve{5\baselineskip}
% mdparagraph:P048:begin
\textbf{Gemanigfieldan.}
% mdparagraph:P048:end

% teaching: NEW-U07-IND-gemanigfieldan
% mdparagraph:P049:begin
The dental stem has the contracted present gemanigfieltst and gemanigfielt. The first-singular present and first-/\allowbreak{}third-singular preterite have the same spelling gemanigfielde; use the sentence to determine tense. Do not add another dental to the printed past stem.
% mdparagraph:P049:end

% table: NEW-U07-IND-gemanigfieldan
% mdtable:T11:begin
% course-table-style: 2026-10
% course-panel:T11-P01:begin
\par\begingroup
\fontsize{10.5}{13.2}\selectfont
\setlength{\LTleft}{0pt}\setlength{\LTright}{\fill}
\noindent
\begin{tabular}{@{}>{\raggedright\arraybackslash}p{39.459mm}>{\raggedright\arraybackslash}p{39.459mm}>{\raggedright\arraybackslash}p{39.459mm}@{}}
\toprule
\textbf{Person and number} & \textbf{Present indicative} & \textbf{Preterite indicative} \\
\midrule
1 sg. & gemanigfielde & gemanigfielde \\
2 sg. & gemanigfieltst & gemanigfieldest \\
3 sg. & gemanigfielt & gemanigfielde \\
1–3 pl. & gemanigfieldað & gemanigfieldon \\
\bottomrule
\end{tabular}
\par\endgroup
% course-panel:T11-P01:end
% mdtable:T11:end
\reserve{5\baselineskip}
% mdparagraph:P050:begin
\textbf{Tōtwǣman.}
% mdparagraph:P050:end

% teaching: NEW-U07-IND-totwaeman
% mdparagraph:P051:begin
The present stem remains tōtwǣm-; the preterite has tōtwǣmd-. The long ǣ remains in both systems.
% mdparagraph:P051:end

% table: NEW-U07-IND-totwaeman
% mdtable:T12:begin
% course-table-style: 2026-10
% course-panel:T12-P01:begin
\par\begingroup
\fontsize{10.5}{13.2}\selectfont
\setlength{\LTleft}{0pt}\setlength{\LTright}{\fill}
\noindent
\begin{tabular}{@{}>{\raggedright\arraybackslash}p{39.459mm}>{\raggedright\arraybackslash}p{39.459mm}>{\raggedright\arraybackslash}p{39.459mm}@{}}
\toprule
\textbf{Person and number} & \textbf{Present indicative} & \textbf{Preterite indicative} \\
\midrule
1 sg. & tōtwǣme & tōtwǣmde \\
2 sg. & tōtwǣmst & tōtwǣmdest \\
3 sg. & tōtwǣmð & tōtwǣmde \\
1–3 pl. & tōtwǣmað & tōtwǣmdon \\
\bottomrule
\end{tabular}
\par\endgroup
% course-panel:T12-P01:end
% mdtable:T12:end
\reserve{7\baselineskip}
% mdparagraph:P052:begin
\subsection*{Their subjunctives}\BookHeading{nav-ch07-learner-0010}{subsubsection}{Their subjunctives}{0}
% mdparagraph:P052:end

% teaching: NEW-U07-SUBJ-gieman
% teaching: NEW-U07-SUBJ-ricsian
% teaching: NEW-U07-SUBJ-totwaeman
% mdparagraph:P053:begin
Wyrcean’s full subjunctive paradigm appears with weak class I in §2. The remaining new verbs follow in these tables. As before, each shared singular or plural row includes all three persons.
% mdparagraph:P053:end

% table: U07-DISPLAY-SUBJ-6
% mdtable:T13:begin
% course-table-style: 2026-10
% course-panel:T13-P01:begin
\par\begingroup
\fontsize{10.5}{13.2}\selectfont
\setlength{\LTleft}{0pt}\setlength{\LTright}{\fill}
\noindent
\begin{tabular}{@{}>{\raggedright\arraybackslash}p{20.56mm}>{\raggedright\arraybackslash}p{20.748mm}>{\raggedright\arraybackslash}p{24.503mm}>{\raggedright\arraybackslash}p{21.593mm}>{\raggedright\arraybackslash}p{25.348mm}@{}}
\toprule
\textbf{Person and number} & \textbf{gīeman · Present} & \textbf{gīeman · Preterite} & \textbf{rīcsian · Present} & \textbf{rīcsian · Preterite} \\
\midrule
1–3 sg. & gīeme & gīemde & rīcsie & rīcsode \\
1–3 pl. & gīemen & gīemden & rīcsien & rīcsoden \\
\bottomrule
\end{tabular}
\par\endgroup
% course-panel:T13-P01:end
\par\addvspace{5pt}
\par\addvspace{3pt}
% course-caption:T13:1:begin
\textbf{tōtwǣman}\par\nobreak
% course-caption:T13:1:end
% course-panel:T13-P02:begin
\par\begingroup
\fontsize{10.5}{13.2}\selectfont
\setlength{\LTleft}{0pt}\setlength{\LTright}{\fill}
\noindent
\begin{tabular}{@{}>{\raggedright\arraybackslash}p{39.459mm}>{\raggedright\arraybackslash}p{39.459mm}>{\raggedright\arraybackslash}p{39.459mm}@{}}
\toprule
\textbf{Person and number} & \textbf{Present} & \textbf{Preterite} \\
\midrule
1–3 sg. & tōtwǣme & tōtwǣmde \\
1–3 pl. & tōtwǣmen & tōtwǣmden \\
\bottomrule
\end{tabular}
\par\endgroup
% course-panel:T13-P02:end
% mdtable:T13:end
% teaching: NEW-U07-SUBJ-gemanigfieldan
% mdparagraph:P054:begin
Gemanigfieldan has identical present and preterite subjunctive spellings. The table is complete in both tenses; the identity of the letters does not erase the distinction.
% mdparagraph:P054:end

% table: U07-DISPLAY-GEMANIGFIELDAN-SUBJ
% mdtable:T14:begin
% course-table-style: 2026-10
% course-panel:T14-P01:begin
\par\begingroup
\fontsize{10.5}{13.2}\selectfont
\setlength{\LTleft}{0pt}\setlength{\LTright}{\fill}
\noindent
\begin{tabular}{@{}>{\raggedright\arraybackslash}p{39.459mm}>{\raggedright\arraybackslash}p{39.459mm}>{\raggedright\arraybackslash}p{39.459mm}@{}}
\toprule
\textbf{Person and number} & \textbf{Present} & \textbf{Preterite} \\
\midrule
1–3 sg. & gemanigfielde & gemanigfielde \\
1–3 pl. & gemanigfielden & gemanigfielden \\
\bottomrule
\end{tabular}
\par\endgroup
% course-panel:T14-P01:end
% mdtable:T14:end
\reserve{7\baselineskip}
% mdparagraph:P055:begin
\subsection*{The imperative of tōtwǣman}\BookHeading{nav-ch07-learner-0011}{subsubsection}{The imperative of tōtwǣman}{0}
% mdparagraph:P055:end

% teaching: NEW-U07-IMP-totwaeman
% teaching: NEW-U07-GRAM-U07-TOTWAEMAN-IMPERATIVE
% mdparagraph:P056:begin
Use the singular to address one person and the plural to address several, including two. Preserve the accusative thing divided and fram with the dative thing separated from.
% mdparagraph:P056:end

% table: NEW-U07-IMP-totwaeman
% mdtable:T15:begin
% course-table-style: 2026-10
% course-panel:T15-P01:begin
\par\begingroup
\fontsize{10.5}{13.2}\selectfont
\setlength{\LTleft}{0pt}\setlength{\LTright}{\fill}
\noindent
\begin{tabular}{@{}>{\raggedright\arraybackslash}p{63.342mm}>{\raggedright\arraybackslash}p{57.846mm}@{}}
\toprule
\textbf{Person and number} & \textbf{Command} \\
\midrule
2 sg. & tōtwǣm \\
2 pl. & tōtwǣmað \\
\bottomrule
\end{tabular}
\par\endgroup
% course-panel:T15-P01:end
% mdtable:T15:end
\reserve{7\baselineskip}
% mdparagraph:P057:begin
\subsection*{Practice A. Forms in use}\BookHeading{nav-ch07-learner-0012}{subsubsection}{Practice A. Forms in use}{0}
% mdparagraph:P057:end

% mdparagraph:P058:begin
Use the selected systems. The four entries in A1–A5 must each be labelled by tense and number; the singular or plural represents all three persons.
% mdparagraph:P058:end

\reserve{5\baselineskip}
\par\begin{minipage}{\linewidth}\raggedright
% mdparagraph:P059:begin
\textbf{A1.}
% mdparagraph:P059:end

% item: NEW-U07-A01
% mdparagraph:P060:begin
Give the present singular, present plural, preterite singular, and preterite plural subjunctive. Each singular/\allowbreak{}plural entry represents all three persons; label every entry. Verbs: hīeran, sęllan.
% mdparagraph:P060:end

\end{minipage}\par\addvspace{4pt}

\par\begin{minipage}{\linewidth}\raggedright
% mdparagraph:P061:begin
\textbf{A2.}
% mdparagraph:P061:end

% item: NEW-U07-A02
% mdparagraph:P062:begin
Give the present singular, present plural, preterite singular, and preterite plural subjunctive. Each singular/\allowbreak{}plural entry represents all three persons; label every entry. Verbs: lufian, libban.
% mdparagraph:P062:end

\end{minipage}\par\addvspace{4pt}

\par\begin{minipage}{\linewidth}\raggedright
% mdparagraph:P063:begin
\textbf{A3.}
% mdparagraph:P063:end

% item: NEW-U07-A03
% mdparagraph:P064:begin
Give the present singular, present plural, preterite singular, and preterite plural subjunctive. Each singular/\allowbreak{}plural entry represents all three persons; label every entry. Verbs: bindan, weaxan.
% mdparagraph:P064:end

\end{minipage}\par\addvspace{4pt}

\par\begin{minipage}{\linewidth}\raggedright
% mdparagraph:P065:begin
\textbf{A4.}
% mdparagraph:P065:end

% item: NEW-U07-A04
% mdparagraph:P066:begin
Give the present singular, present plural, preterite singular, and preterite plural subjunctive. Each singular/\allowbreak{}plural entry represents all three persons; label every entry. Verbs: sēon, biddan.
% mdparagraph:P066:end

\end{minipage}\par\addvspace{4pt}

\par\begin{minipage}{\linewidth}\raggedright
% mdparagraph:P067:begin
\textbf{A5.}
% mdparagraph:P067:end

% item: NEW-U07-A05
% mdparagraph:P068:begin
Give the present singular, present plural, preterite singular, and preterite plural subjunctive. Each singular/\allowbreak{}plural entry represents all three persons; label every entry. Verbs: witan, willan.
% mdparagraph:P068:end

\end{minipage}\par\addvspace{4pt}

\par\begin{minipage}{\linewidth}\raggedright
% mdparagraph:P069:begin
\textbf{A6.}
% mdparagraph:P069:end

% item: NEW-U07-A06
% mdparagraph:P070:begin
Give both selected present subjunctive series of wesan/\allowbreak{}bēon and the shared preterite subjunctive, singular and plural. State which person distinctions the forms express.
% mdparagraph:P070:end

\end{minipage}\par\addvspace{4pt}

\par\begin{minipage}{\linewidth}\raggedright
% mdparagraph:P071:begin
\textbf{A7.}
% mdparagraph:P071:end

% item: NEW-U07-A07
% mdparagraph:P072:begin
For gīeman, rīcsian, gemanigfieldan, tōtwǣman, and wyrcean, give the present indicative third singular and preterite indicative third plural. Then give the present subjunctive plural of each.
% mdparagraph:P072:end

\end{minipage}\par\addvspace{4pt}

\par\begin{minipage}{\linewidth}\raggedright
% mdparagraph:P073:begin
\textbf{A8.}
% mdparagraph:P073:end

% item: NEW-U07-A08
% mdparagraph:P074:begin
Contrast hīe bundon with hīe bunden, and hīe lufodon with hīe lufoden. Parse each finite form and explain what the ending contributes. Do not invent a full English sentence from an isolated mood form.
% mdparagraph:P074:end

\end{minipage}\par\addvspace{4pt}

\reserve{7\baselineskip}
% mdparagraph:P075:begin
\section*{4. Basic conditions with gif}\BookHeading{nav-ch07-learner-0013}{subsection}{4. Basic conditions with gif}{1}\BookLabel{bk-ch07-learner-49a97724b59e}
% mdparagraph:P075:end

% teaching: NEW-U07-VOC-gif
% teaching: NEW-U07-VOC-gif-if
% teaching: NEW-U07-GRAM-U07-BASIC-CONDITIONALS
% mdparagraph:P076:begin
\textbf{Gif} means “if”. It introduces a condition; it does not automatically require the subjunctive. An open condition can be stated in the indicative.
% mdparagraph:P076:end

% example: NEW-U07-EXAMPLE-04
% mdparagraph:P077:begin
Gif hē hæfð męte, hē sęlð ðǣm fuglum męte.
% mdparagraph:P077:end

% mdparagraph:P078:begin
If he has food, he gives the birds food.
% mdparagraph:P078:end

% mdparagraph:P079:begin
This is an open condition: having food is presented as a real possibility, without a specially hypothetical setting. Hæfð and sęlð are indicative. The birds are the dative recipients.
% mdparagraph:P079:end

% mdparagraph:P080:begin
A speaker can instead use a present subjunctive in a contingent or tentative condition. The following clause retains the form appropriate to its own meaning.
% mdparagraph:P080:end

% example: NEW-U07-EXAMPLE-05
% mdparagraph:P081:begin
Gif ðū mē helpe, wē magon ðā dēor bindan.
% mdparagraph:P081:end

% mdparagraph:P082:begin
If you help me, we can bind the animals.
% mdparagraph:P082:end

% mdparagraph:P083:begin
Here the speaker tentatively seeks help. Helpe is second-singular present subjunctive; magon is first-plural present indicative, asserting capability on that condition. Mē remains dative after helpan. The main clause need not become subjunctive simply because the condition is subjunctive.
% mdparagraph:P083:end

% teaching: NEW-U07-GRAM-U07-NONPAST-PRETERITE-SUBJUNCTIVE
% mdparagraph:P084:begin
\textbf{Recognition: a preterite form with nonpast reference.} The following setting and interpretation are supplied: he is not coming, and the situation is explicitly unreal and nonpast.
% mdparagraph:P084:end

% example: NEW-U07-EXAMPLE-06
% mdparagraph:P085:begin
Gif hē cōme, ic him hulpe.
% mdparagraph:P085:end

% mdparagraph:P086:begin
If he came, I would help him.
% mdparagraph:P086:end

% mdparagraph:P087:begin
Cōme\BookPoint{bk-ch07-learner-19b1b69abd23} and hulpe are preterite subjunctives, but this context gives them nonpast reference. The morphological label “preterite” does not itself force past time. Recognize this use in the supplied example; a general system of unreal conditions is not required for independent composition here.
% mdparagraph:P087:end

\reserve{7\baselineskip}
% mdparagraph:P088:begin
\subsection*{Practice M. Mood and meaning}\BookHeading{nav-ch07-learner-0014}{subsubsection}{Practice M. Mood and meaning}{0}
% mdparagraph:P088:end

% mdparagraph:P089:begin
Distinguish what the form establishes from what the stated context establishes. Where two moods share a spelling, identify the intended reading without declaring the spelling unique.
% mdparagraph:P089:end

\reserve{5\baselineskip}
\par\begin{minipage}{\linewidth}\raggedright
% mdparagraph:P090:begin
\textbf{M1.}
% mdparagraph:P090:end

% item: NEW-U07-M01
% mdparagraph:P091:begin
Translate and contrast Hīe swimmað on ðǣm sǣ and Swimmen hīe on ðǣm sǣ. State what is asserted or requested.
% mdparagraph:P091:end

\end{minipage}\par\addvspace{4pt}

\par\begin{minipage}{\linewidth}\raggedright
% mdparagraph:P092:begin
\textbf{M2.}
% mdparagraph:P092:end

% item: NEW-U07-M02
% mdparagraph:P093:begin
Translate Gāð and Gān wē. Identify the persons addressed or included, and explain why they are not the same command form.
% mdparagraph:P093:end

\end{minipage}\par\addvspace{4pt}

\par\begin{minipage}{\linewidth}\raggedright
% mdparagraph:P094:begin
\textbf{M3.}
% mdparagraph:P094:end

% item: NEW-U07-M03
% mdparagraph:P095:begin
Translate Hē mæg swimman, Hē mōt swimman, and Swimme hē. Distinguish capability, permission, and the speaker’s wish.
% mdparagraph:P095:end

\end{minipage}\par\addvspace{4pt}

\par\begin{minipage}{\linewidth}\raggedright
% mdparagraph:P096:begin
\textbf{M4.}
% mdparagraph:P096:end

% item: NEW-U07-M04
% mdparagraph:P097:begin
An open condition describes a real possibility without presenting it as specially hypothetical: Gif hē hæfð męte, hē sęlð ðǣm fuglum męte. Parse and translate the clause pair. Does gif itself require subjunctive morphology?
% mdparagraph:P097:end

\end{minipage}\par\addvspace{4pt}

\par\begin{minipage}{\linewidth}\raggedright
% mdparagraph:P098:begin
\textbf{M5.}
% mdparagraph:P098:end

% item: NEW-U07-M05
% mdparagraph:P099:begin
A speaker tentatively seeks help: Gif ðū mē helpe, wē magon ðā dēor bindan. Parse helpe and magon, translate the sentence, and explain the difference between their roles.
% mdparagraph:P099:end

\end{minipage}\par\addvspace{4pt}

\par\begin{minipage}{\linewidth}\raggedright
% mdparagraph:P100:begin
\textbf{M6.}
% mdparagraph:P100:end

% item: NEW-U07-M06
% mdparagraph:P101:begin
The setting is explicitly unreal and nonpast: he is not coming. The supplied interpretation of Gif hē cōme, ic him hulpe is “If he came, I would help him.” Identify the tense and mood of cōme and hulpe. Explain why their morphological label alone does not force past time.
% mdparagraph:P101:end

\end{minipage}\par\addvspace{4pt}

\par\begin{minipage}{\linewidth}\raggedright
% mdparagraph:P102:begin
\textbf{M7.}
% mdparagraph:P102:end

% item: NEW-U07-M07
% mdparagraph:P103:begin
The narrator asserts a past wish in Hē wolde faran. Parse wolde according to that intended assertion and translate it. State the other mood whose form has the same spelling, and explain why the spelling alone does not require the English word “would.”
% mdparagraph:P103:end

\end{minipage}\par\addvspace{4pt}

\par\begin{minipage}{\linewidth}\raggedright
% mdparagraph:P104:begin
\textbf{M8.}
% mdparagraph:P104:end

% item: NEW-U07-M08
% mdparagraph:P105:begin
Give a third-person singular wish corresponding to Hē leofað, and a first-person plural exhortation corresponding to Wē helpað ðǣm męnn. Preserve the lexical meaning and government.
% mdparagraph:P105:end

\end{minipage}\par\addvspace{4pt}

\reserve{7\baselineskip}
% mdparagraph:P106:begin
\section*{5. Participles}\BookHeading{nav-ch07-learner-0015}{subsection}{5. Participles}{1}\BookLabel{bk-ch07-learner-ae50c7e839f0}
% mdparagraph:P106:end

% teaching: NEW-U07-GRAM-U07-PARTICIPLE-FORMATION
% mdparagraph:P107:begin
A \textbf{participle} is a nonfinite verb form. The present participle in \textbf{-ende} describes an activity; a past participle can describe a resulting condition or enter a passive construction. Neither participle by itself is a finite present or preterite verb.
% mdparagraph:P107:end

% mdparagraph:P108:begin
For the selected present participles, retain the present stem before -ende. Verbs in -ian retain i: nęrian gives nęriende and lufian gives lufiende. Gefręmman retains double mm in its present participle. The strong past participle uses its own principal-part stem and -en, as in bunden, boren, and cweden; it is not made by adding a weak dental ending to the infinitive. The weak past participles have their selected dental forms, including hīered, lufod, and seald.
% mdparagraph:P108:end

% mdparagraph:P109:begin
The following twelve pairs are the formation assignment. Preserve the actual lexical prefix: geblētsod and gehālgod belong to their prefixed verbs; bunden and hāten are the selected unprefixed forms. A dictionary’s optional prefix notation is not a word to copy into a sentence, and it is not a rule for adding ge- to arbitrary forms.
% mdparagraph:P109:end

% teaching: NEW-U07-PART-bindan
% teaching: NEW-U07-PART-hieran
% teaching: NEW-U07-PART-lufian
% teaching: NEW-U07-PART-sellan
% teaching: NEW-U07-PART-hatan
% teaching: NEW-U07-PART-gemanigfieldan
% teaching: NEW-U07-PART-nerian
% teaching: NEW-U07-PART-cwethan
% teaching: NEW-U07-PART-beran
% teaching: NEW-U07-PART-gehalgian
% teaching: NEW-U07-PART-gebletsian
% teaching: NEW-U07-PART-gefremman
% table: U07-DISPLAY-PART-PAIRS
% mdtable:T16:begin
\par\begingroup
\fontsize{10.5}{13.2}\selectfont
\setlength{\LTleft}{0pt}\setlength{\LTright}{\fill}
\noindent
\begin{tabular}{@{}>{\raggedright\arraybackslash}p{37.45mm}>{\raggedright\arraybackslash}p{45.21mm}>{\raggedright\arraybackslash}p{35.71mm}@{}}
\toprule
\textbf{Lemma} & \textbf{Present participle} & \textbf{Past participle} \\
\midrule
bindan & bindende & bunden \\
hīeran & hīerende & hīered \\
lufian & lufiende & lufod \\
sęllan & sęllende & seald \\
hātan & hātende & hāten \\
gemanigfieldan & gemanigfieldende & gemanigfield \\
nęrian & nęriende & nęred \\
cweðan & cweðende & cweden \\
beran & berende & boren \\
gehālgian & gehālgiende & gehālgod \\
geblētsian & geblētsiende & geblētsod \\
gefręmman & gefręmmende & gefręmed \\
\bottomrule
\end{tabular}
\par\endgroup

% mdtable:T16:end
% mdparagraph:P110:begin
The\BookPoint{bk-ch07-learner-3fa3aaeef76b} weak preterite is not a universal past-participle recipe. Hīerde and hīered differ in the placement of e; lufode corresponds to lufod, and sealde to seald. Gefręmed has single m, unlike the present gefręmmende. Learn the actual principal part with its class rather than merely removing the final letter of every past form.
% mdparagraph:P110:end

% teaching: NEW-U07-GRAM-U07-PARTICIPLE-AGREEMENT
\reserve{7\baselineskip}
% mdparagraph:P111:begin
\subsection*{Participles used as adjectives}\BookHeading{nav-ch07-learner-0016}{subsubsection}{Participles used as adjectives}{0}\BookLabel{bk-ch07-learner-5cc5caec500a}
% mdparagraph:P111:end

% mdparagraph:P112:begin
An attributive participle agrees with its noun in \textbf{gender, number, and case}. Use the strong or weak declension under the conditions already learned for adjectives: a definite article calls for the weak attributive form; an otherwise undetermined adjective ordinarily takes the strong form. No new instrumental use is added here.
% mdparagraph:P112:end

% mdparagraph:P113:begin
The five complete models below give all four cases in all three genders and both numbers, in both strong and weak inflection. The present participle follows the grēne-type adjective pattern. It is not an agent noun in -end.
% mdparagraph:P113:end

% example: NEW-U07-EXAMPLE-07
% mdparagraph:P114:begin
Se bindenda mann is gōd.
% mdparagraph:P114:end

% mdparagraph:P115:begin
The binding man is good.
% mdparagraph:P115:end

% mdparagraph:P116:begin
Bindenda is weak masculine nominative singular after se, modifying mann. Gōd is the separate predicate adjective.
% mdparagraph:P116:end

\reserve{7\baselineskip}
% mdparagraph:P117:begin
\subsection*{Bindende: present participle}\BookHeading{nav-ch07-learner-0017}{subsubsection}{Bindende: present participle}{0}\BookLabel{bk-ch07-learner-22bdebc13439}
% mdparagraph:P117:end

% teaching: NEW-U07-PART-ENDE
% mdparagraph:P118:begin
Retain the participial stem bindend- before consonantal endings, and observe the selected -e/\allowbreak{}-u distribution in the strong forms. The weak nominative singular has masculine -a beside feminine and neuter -e.
% mdparagraph:P118:end

\reserve{9\baselineskip}
% mdparagraph:P119:begin
\textbf{Strong singular.}
% mdparagraph:P119:end

% table: NEW-U07-PART-ENDE-1
% mdtable:T17:begin
% course-table-style: 2026-10
% course-panel:T17-P01:begin
\par\begingroup
\fontsize{10.5}{13.2}\selectfont
\setlength{\LTleft}{0pt}\setlength{\LTright}{\fill}
\noindent
\begin{tabular}{@{}>{\raggedright\arraybackslash}p{22.0mm}>{\raggedright\arraybackslash}p{31.188mm}>{\raggedright\arraybackslash}p{31.188mm}>{\raggedright\arraybackslash}p{31.188mm}@{}}
\toprule
\textbf{Case} & \textbf{Masculine} & \textbf{Neuter} & \textbf{Feminine} \\
\midrule
Nominative & bindende & bindende & bindendu \\
Genitive & bindendes & bindendes & bindendre \\
Dative & bindendum & bindendum & bindendre \\
Accusative & bindendne & bindende & bindende \\
\bottomrule
\end{tabular}
\par\endgroup
% course-panel:T17-P01:end
% mdtable:T17:end
\reserve{9\baselineskip}
% mdparagraph:P120:begin
\textbf{Strong plural.}
% mdparagraph:P120:end

% table: NEW-U07-PART-ENDE-2
% mdtable:T18:begin
% course-table-style: 2026-10
% course-panel:T18-P01:begin
\par\begingroup
\fontsize{10.5}{13.2}\selectfont
\setlength{\LTleft}{0pt}\setlength{\LTright}{\fill}
\noindent
\begin{tabular}{@{}>{\raggedright\arraybackslash}p{22.0mm}>{\raggedright\arraybackslash}p{31.188mm}>{\raggedright\arraybackslash}p{31.188mm}>{\raggedright\arraybackslash}p{31.188mm}@{}}
\toprule
\textbf{Case} & \textbf{Masculine} & \textbf{Neuter} & \textbf{Feminine} \\
\midrule
Nominative & bindende & bindendu & bindenda \\
Genitive & bindendra & bindendra & bindendra \\
Dative & bindendum & bindendum & bindendum \\
Accusative & bindende & bindendu & bindenda \\
\bottomrule
\end{tabular}
\par\endgroup
% course-panel:T18-P01:end
% mdtable:T18:end
\reserve{9\baselineskip}
% mdparagraph:P121:begin
\textbf{Weak singular.}
% mdparagraph:P121:end

% table: NEW-U07-PART-ENDE-3
% mdtable:T19:begin
% course-table-style: 2026-10
% course-panel:T19-P01:begin
\par\begingroup
\fontsize{10.5}{13.2}\selectfont
\setlength{\LTleft}{0pt}\setlength{\LTright}{\fill}
\noindent
\begin{tabular}{@{}>{\raggedright\arraybackslash}p{22.0mm}>{\raggedright\arraybackslash}p{31.188mm}>{\raggedright\arraybackslash}p{31.188mm}>{\raggedright\arraybackslash}p{31.188mm}@{}}
\toprule
\textbf{Case} & \textbf{Masculine} & \textbf{Neuter} & \textbf{Feminine} \\
\midrule
Nominative & bindenda & bindende & bindende \\
Genitive & bindendan & bindendan & bindendan \\
Dative & bindendan & bindendan & bindendan \\
Accusative & bindendan & bindende & bindendan \\
\bottomrule
\end{tabular}
\par\endgroup
% course-panel:T19-P01:end
% mdtable:T19:end
\reserve{9\baselineskip}
% mdparagraph:P122:begin
\textbf{Weak plural.}
% mdparagraph:P122:end

% table: NEW-U07-PART-ENDE-4
% mdtable:T20:begin
% course-table-style: 2026-10
% course-panel:T20-P01:begin
\par\begingroup
\fontsize{10.5}{13.2}\selectfont
\setlength{\LTleft}{0pt}\setlength{\LTright}{\fill}
\noindent
\begin{tabular}{@{}>{\raggedright\arraybackslash}p{22.0mm}>{\raggedright\arraybackslash}p{31.188mm}>{\raggedright\arraybackslash}p{31.188mm}>{\raggedright\arraybackslash}p{31.188mm}@{}}
\toprule
\textbf{Case} & \textbf{Masculine} & \textbf{Neuter} & \textbf{Feminine} \\
\midrule
Nominative & bindendan & bindendan & bindendan \\
Genitive & bindendra & bindendra & bindendra \\
Dative & bindendum & bindendum & bindendum \\
Accusative & bindendan & bindendan & bindendan \\
\bottomrule
\end{tabular}
\par\endgroup
% course-panel:T20-P01:end
% mdtable:T20:end
% mdparagraph:P123:begin
\textbf{Variants.} In the strong feminine plural nominative/\allowbreak{}accusative, bindende is also found beside selected bindenda; in the corresponding neuter plural cells, bindende is an alternative to bindendu. Keep the selected forms when a task asks for the model. These variations do not make the participle an uninflected word in every attributive use.
% mdparagraph:P123:end

% teaching: NEW-U07-GRAM-U07-PARTICIPLE-SYNCOPE
\reserve{7\baselineskip}
% mdparagraph:P124:begin
\subsection*{Bunden: strong past participle}\BookHeading{nav-ch07-learner-0018}{subsubsection}{Bunden: strong past participle}{0}\BookLabel{bk-ch07-learner-ea64958329b3}
% mdparagraph:P124:end

% teaching: NEW-U07-PART-BUNDEN
% mdparagraph:P125:begin
The table keeps the full stem bunden- so that the endings remain visible. Its unstressed e can disappear before a vowel-initial ending: this is syncope. Preserve the full stem before -ne, -re, and -ra; bundenne therefore has double n.
% mdparagraph:P125:end

\reserve{9\baselineskip}
% mdparagraph:P126:begin
\textbf{Strong singular.}
% mdparagraph:P126:end

% table: NEW-U07-PART-BUNDEN-1
% mdtable:T21:begin
% course-table-style: 2026-10
% course-panel:T21-P01:begin
\par\begingroup
\fontsize{10.5}{13.2}\selectfont
\setlength{\LTleft}{0pt}\setlength{\LTright}{\fill}
\noindent
\begin{tabular}{@{}>{\raggedright\arraybackslash}p{22.0mm}>{\raggedright\arraybackslash}p{31.188mm}>{\raggedright\arraybackslash}p{31.188mm}>{\raggedright\arraybackslash}p{31.188mm}@{}}
\toprule
\textbf{Case} & \textbf{Masculine} & \textbf{Neuter} & \textbf{Feminine} \\
\midrule
Nominative & bunden & bunden & bundenu \\
Genitive & bundenes & bundenes & bundenre \\
Dative & bundenum & bundenum & bundenre \\
Accusative & bundenne & bunden & bundene \\
\bottomrule
\end{tabular}
\par\endgroup
% course-panel:T21-P01:end
% mdtable:T21:end
\reserve{9\baselineskip}
% mdparagraph:P127:begin
\textbf{Strong plural.}
% mdparagraph:P127:end

% table: NEW-U07-PART-BUNDEN-2
% mdtable:T22:begin
% course-table-style: 2026-10
% course-panel:T22-P01:begin
\par\begingroup
\fontsize{10.5}{13.2}\selectfont
\setlength{\LTleft}{0pt}\setlength{\LTright}{\fill}
\noindent
\begin{tabular}{@{}>{\raggedright\arraybackslash}p{22.0mm}>{\raggedright\arraybackslash}p{31.188mm}>{\raggedright\arraybackslash}p{31.188mm}>{\raggedright\arraybackslash}p{31.188mm}@{}}
\toprule
\textbf{Case} & \textbf{Masculine} & \textbf{Neuter} & \textbf{Feminine} \\
\midrule
Nominative & bundene & bundenu & bundena \\
Genitive & bundenra & bundenra & bundenra \\
Dative & bundenum & bundenum & bundenum \\
Accusative & bundene & bundenu & bundena \\
\bottomrule
\end{tabular}
\par\endgroup
% course-panel:T22-P01:end
% mdtable:T22:end
\reserve{9\baselineskip}
% mdparagraph:P128:begin
\textbf{Weak singular.}
% mdparagraph:P128:end

% table: NEW-U07-PART-BUNDEN-3
% mdtable:T23:begin
% course-table-style: 2026-10
% course-panel:T23-P01:begin
\par\begingroup
\fontsize{10.5}{13.2}\selectfont
\setlength{\LTleft}{0pt}\setlength{\LTright}{\fill}
\noindent
\begin{tabular}{@{}>{\raggedright\arraybackslash}p{22.0mm}>{\raggedright\arraybackslash}p{31.188mm}>{\raggedright\arraybackslash}p{31.188mm}>{\raggedright\arraybackslash}p{31.188mm}@{}}
\toprule
\textbf{Case} & \textbf{Masculine} & \textbf{Neuter} & \textbf{Feminine} \\
\midrule
Nominative & bundena & bundene & bundene \\
Genitive & bundenan & bundenan & bundenan \\
Dative & bundenan & bundenan & bundenan \\
Accusative & bundenan & bundene & bundenan \\
\bottomrule
\end{tabular}
\par\endgroup
% course-panel:T23-P01:end
% mdtable:T23:end
\reserve{9\baselineskip}
% mdparagraph:P129:begin
\textbf{Weak plural.}
% mdparagraph:P129:end

% table: NEW-U07-PART-BUNDEN-4
% mdtable:T24:begin
% course-table-style: 2026-10
% course-panel:T24-P01:begin
\par\begingroup
\fontsize{10.5}{13.2}\selectfont
\setlength{\LTleft}{0pt}\setlength{\LTright}{\fill}
\noindent
\begin{tabular}{@{}>{\raggedright\arraybackslash}p{22.0mm}>{\raggedright\arraybackslash}p{31.188mm}>{\raggedright\arraybackslash}p{31.188mm}>{\raggedright\arraybackslash}p{31.188mm}@{}}
\toprule
\textbf{Case} & \textbf{Masculine} & \textbf{Neuter} & \textbf{Feminine} \\
\midrule
Nominative & bundenan & bundenan & bundenan \\
Genitive & bundenra & bundenra & bundenra \\
Dative & bundenum & bundenum & bundenum \\
Accusative & bundenan & bundenan & bundenan \\
\bottomrule
\end{tabular}
\par\endgroup
% course-panel:T24-P01:end
% mdtable:T24:end
% mdparagraph:P130:begin
\textbf{Syncope and variants.} Bundenes/\allowbreak{}bundnes, bundenum/\allowbreak{}bundnum, bundene/\allowbreak{}bundne, bundena/\allowbreak{}bundna, and bundenan/\allowbreak{}bundnan are the corresponding full and syncopated forms. Apply each pair only to the cells with that full form. Strong feminine plural nominative/\allowbreak{}accusative also permits bundene beside bundena. Strong feminine nominative singular permits bunden beside bundenu; strong neuter plural nominative/\allowbreak{}accusative permits bunden or bundene beside bundenu. The selected -u forms complete the model; they are not claims that every such cell is frequent in continuous texts.
% mdparagraph:P130:end

\reserve{7\baselineskip}
% mdparagraph:P131:begin
\subsection*{Hīered: weak past participle with syncope}\BookHeading{nav-ch07-learner-0019}{subsubsection}{Hīered: weak past participle with syncope}{0}\BookLabel{bk-ch07-learner-49e61c42cd78}
% mdparagraph:P131:end

% teaching: NEW-U07-PART-HIERED
% mdparagraph:P132:begin
The full stem hīered- makes the relation to the dental participle clear. Keep it before -ne, -re, and -ra. Before vowel-initial endings its unstressed e can disappear, producing hīerd-.
% mdparagraph:P132:end

\reserve{9\baselineskip}
% mdparagraph:P133:begin
\textbf{Strong singular.}
% mdparagraph:P133:end

% table: NEW-U07-PART-HIERED-1
% mdtable:T25:begin
% course-table-style: 2026-10
% course-panel:T25-P01:begin
\par\begingroup
\fontsize{10.5}{13.2}\selectfont
\setlength{\LTleft}{0pt}\setlength{\LTright}{\fill}
\noindent
\begin{tabular}{@{}>{\raggedright\arraybackslash}p{22.0mm}>{\raggedright\arraybackslash}p{31.188mm}>{\raggedright\arraybackslash}p{31.188mm}>{\raggedright\arraybackslash}p{31.188mm}@{}}
\toprule
\textbf{Case} & \textbf{Masculine} & \textbf{Neuter} & \textbf{Feminine} \\
\midrule
Nominative & hīered & hīered & hīeredu \\
Genitive & hīeredes & hīeredes & hīeredre \\
Dative & hīeredum & hīeredum & hīeredre \\
Accusative & hīeredne & hīered & hīerede \\
\bottomrule
\end{tabular}
\par\endgroup
% course-panel:T25-P01:end
% mdtable:T25:end
\reserve{9\baselineskip}
% mdparagraph:P134:begin
\textbf{Strong plural.}
% mdparagraph:P134:end

% table: NEW-U07-PART-HIERED-2
% mdtable:T26:begin
% course-table-style: 2026-10
% course-panel:T26-P01:begin
\par\begingroup
\fontsize{10.5}{13.2}\selectfont
\setlength{\LTleft}{0pt}\setlength{\LTright}{\fill}
\noindent
\begin{tabular}{@{}>{\raggedright\arraybackslash}p{22.0mm}>{\raggedright\arraybackslash}p{31.188mm}>{\raggedright\arraybackslash}p{31.188mm}>{\raggedright\arraybackslash}p{31.188mm}@{}}
\toprule
\textbf{Case} & \textbf{Masculine} & \textbf{Neuter} & \textbf{Feminine} \\
\midrule
Nominative & hīerede & hīeredu & hīereda \\
Genitive & hīeredra & hīeredra & hīeredra \\
Dative & hīeredum & hīeredum & hīeredum \\
Accusative & hīerede & hīeredu & hīereda \\
\bottomrule
\end{tabular}
\par\endgroup
% course-panel:T26-P01:end
% mdtable:T26:end
\reserve{9\baselineskip}
% mdparagraph:P135:begin
\textbf{Weak singular.}
% mdparagraph:P135:end

% table: NEW-U07-PART-HIERED-3
% mdtable:T27:begin
% course-table-style: 2026-10
% course-panel:T27-P01:begin
\par\begingroup
\fontsize{10.5}{13.2}\selectfont
\setlength{\LTleft}{0pt}\setlength{\LTright}{\fill}
\noindent
\begin{tabular}{@{}>{\raggedright\arraybackslash}p{22.0mm}>{\raggedright\arraybackslash}p{31.188mm}>{\raggedright\arraybackslash}p{31.188mm}>{\raggedright\arraybackslash}p{31.188mm}@{}}
\toprule
\textbf{Case} & \textbf{Masculine} & \textbf{Neuter} & \textbf{Feminine} \\
\midrule
Nominative & hīereda & hīerede & hīerede \\
Genitive & hīeredan & hīeredan & hīeredan \\
Dative & hīeredan & hīeredan & hīeredan \\
Accusative & hīeredan & hīerede & hīeredan \\
\bottomrule
\end{tabular}
\par\endgroup
% course-panel:T27-P01:end
% mdtable:T27:end
\reserve{9\baselineskip}
% mdparagraph:P136:begin
\textbf{Weak plural.}
% mdparagraph:P136:end

% table: NEW-U07-PART-HIERED-4
% mdtable:T28:begin
% course-table-style: 2026-10
% course-panel:T28-P01:begin
\par\begingroup
\fontsize{10.5}{13.2}\selectfont
\setlength{\LTleft}{0pt}\setlength{\LTright}{\fill}
\noindent
\begin{tabular}{@{}>{\raggedright\arraybackslash}p{22.0mm}>{\raggedright\arraybackslash}p{31.188mm}>{\raggedright\arraybackslash}p{31.188mm}>{\raggedright\arraybackslash}p{31.188mm}@{}}
\toprule
\textbf{Case} & \textbf{Masculine} & \textbf{Neuter} & \textbf{Feminine} \\
\midrule
Nominative & hīeredan & hīeredan & hīeredan \\
Genitive & hīeredra & hīeredra & hīeredra \\
Dative & hīeredum & hīeredum & hīeredum \\
Accusative & hīeredan & hīeredan & hīeredan \\
\bottomrule
\end{tabular}
\par\endgroup
% course-panel:T28-P01:end
% mdtable:T28:end
% mdparagraph:P137:begin
\textbf{Syncope and variants.} Hīeredes/\allowbreak{}hīerdes, hīeredum/\allowbreak{}hīerdum, hīerede/\allowbreak{}hīerde, hīereda/\allowbreak{}hīerda, and hīeredan/\allowbreak{}hīerdan give the corresponding alternatives. Strong feminine plural nominative/\allowbreak{}accusative also permits hīerede beside hīereda. Strong feminine nominative singular permits hīered beside hīeredu; strong neuter plural nominative/\allowbreak{}accusative permits hīered or hīerede beside hīeredu. As in the preceding model, the -u cells are regular paradigm completions with their alternatives stated explicitly.
% mdparagraph:P137:end

\reserve{7\baselineskip}
% mdparagraph:P138:begin
\subsection*{Lufod: class II past participle}\BookHeading{nav-ch07-learner-0020}{subsubsection}{Lufod: class II past participle}{0}\BookLabel{bk-ch07-learner-3ddd770922b3}
% mdparagraph:P138:end

% teaching: NEW-U07-PART-LUFOD
% mdparagraph:P139:begin
Retain o throughout the selected model. Unlike the preceding -u models, selected lufod has no ending in the strong feminine nominative singular or neuter plural nominative/\allowbreak{}accusative. The remaining endings display the familiar strong/\allowbreak{}weak contrasts.
% mdparagraph:P139:end

\reserve{9\baselineskip}
% mdparagraph:P140:begin
\textbf{Strong singular.}
% mdparagraph:P140:end

% table: NEW-U07-PART-LUFOD-1
% mdtable:T29:begin
% course-table-style: 2026-10
% course-panel:T29-P01:begin
\par\begingroup
\fontsize{10.5}{13.2}\selectfont
\setlength{\LTleft}{0pt}\setlength{\LTright}{\fill}
\noindent
\begin{tabular}{@{}>{\raggedright\arraybackslash}p{22.0mm}>{\raggedright\arraybackslash}p{31.188mm}>{\raggedright\arraybackslash}p{31.188mm}>{\raggedright\arraybackslash}p{31.188mm}@{}}
\toprule
\textbf{Case} & \textbf{Masculine} & \textbf{Neuter} & \textbf{Feminine} \\
\midrule
Nominative & lufod & lufod & lufod \\
Genitive & lufodes & lufodes & lufodre \\
Dative & lufodum & lufodum & lufodre \\
Accusative & lufodne & lufod & lufode \\
\bottomrule
\end{tabular}
\par\endgroup
% course-panel:T29-P01:end
% mdtable:T29:end
\reserve{9\baselineskip}
% mdparagraph:P141:begin
\textbf{Strong plural.}
% mdparagraph:P141:end

% table: NEW-U07-PART-LUFOD-2
% mdtable:T30:begin
% course-table-style: 2026-10
% course-panel:T30-P01:begin
\par\begingroup
\fontsize{10.5}{13.2}\selectfont
\setlength{\LTleft}{0pt}\setlength{\LTright}{\fill}
\noindent
\begin{tabular}{@{}>{\raggedright\arraybackslash}p{22.0mm}>{\raggedright\arraybackslash}p{31.188mm}>{\raggedright\arraybackslash}p{31.188mm}>{\raggedright\arraybackslash}p{31.188mm}@{}}
\toprule
\textbf{Case} & \textbf{Masculine} & \textbf{Neuter} & \textbf{Feminine} \\
\midrule
Nominative & lufode & lufod & lufoda \\
Genitive & lufodra & lufodra & lufodra \\
Dative & lufodum & lufodum & lufodum \\
Accusative & lufode & lufod & lufoda \\
\bottomrule
\end{tabular}
\par\endgroup
% course-panel:T30-P01:end
% mdtable:T30:end
\reserve{9\baselineskip}
% mdparagraph:P142:begin
\textbf{Weak singular.}
% mdparagraph:P142:end

% table: NEW-U07-PART-LUFOD-3
% mdtable:T31:begin
% course-table-style: 2026-10
% course-panel:T31-P01:begin
\par\begingroup
\fontsize{10.5}{13.2}\selectfont
\setlength{\LTleft}{0pt}\setlength{\LTright}{\fill}
\noindent
\begin{tabular}{@{}>{\raggedright\arraybackslash}p{22.0mm}>{\raggedright\arraybackslash}p{31.188mm}>{\raggedright\arraybackslash}p{31.188mm}>{\raggedright\arraybackslash}p{31.188mm}@{}}
\toprule
\textbf{Case} & \textbf{Masculine} & \textbf{Neuter} & \textbf{Feminine} \\
\midrule
Nominative & lufoda & lufode & lufode \\
Genitive & lufodan & lufodan & lufodan \\
Dative & lufodan & lufodan & lufodan \\
Accusative & lufodan & lufode & lufodan \\
\bottomrule
\end{tabular}
\par\endgroup
% course-panel:T31-P01:end
% mdtable:T31:end
\reserve{9\baselineskip}
% mdparagraph:P143:begin
\textbf{Weak plural.}
% mdparagraph:P143:end

% table: NEW-U07-PART-LUFOD-4
% mdtable:T32:begin
% course-table-style: 2026-10
% course-panel:T32-P01:begin
\par\begingroup
\fontsize{10.5}{13.2}\selectfont
\setlength{\LTleft}{0pt}\setlength{\LTright}{\fill}
\noindent
\begin{tabular}{@{}>{\raggedright\arraybackslash}p{22.0mm}>{\raggedright\arraybackslash}p{31.188mm}>{\raggedright\arraybackslash}p{31.188mm}>{\raggedright\arraybackslash}p{31.188mm}@{}}
\toprule
\textbf{Case} & \textbf{Masculine} & \textbf{Neuter} & \textbf{Feminine} \\
\midrule
Nominative & lufodan & lufodan & lufodan \\
Genitive & lufodra & lufodra & lufodra \\
Dative & lufodum & lufodum & lufodum \\
Accusative & lufodan & lufodan & lufodan \\
\bottomrule
\end{tabular}
\par\endgroup
% course-panel:T32-P01:end
% mdtable:T32:end
% mdparagraph:P144:begin
\textbf{Variants.} Strong feminine plural nominative/\allowbreak{}accusative lufode is an alternative to lufoda; strong neuter plural nominative/\allowbreak{}accusative lufode is an alternative to lufod. These are particular agreement variants, not a replacement of every final syllable by -e.
% mdparagraph:P144:end

\reserve{7\baselineskip}
% mdparagraph:P145:begin
\subsection*{Seald: monosyllabic dental past participle}\BookHeading{nav-ch07-learner-0021}{subsubsection}{Seald: monosyllabic dental past participle}{0}\BookLabel{bk-ch07-learner-de8dd5a7122c}
% mdparagraph:P145:end

% teaching: NEW-U07-PART-SEALD
% mdparagraph:P146:begin
Seald has no unstressed stem vowel to remove. Add the printed endings directly to seald-. The strong feminine nominative singular and neuter plural nominative/\allowbreak{}accusative are unextended in the selected model.
% mdparagraph:P146:end

\reserve{9\baselineskip}
% mdparagraph:P147:begin
\textbf{Strong singular.}
% mdparagraph:P147:end

% table: NEW-U07-PART-SEALD-1
% mdtable:T33:begin
% course-table-style: 2026-10
% course-panel:T33-P01:begin
\par\begingroup
\fontsize{10.5}{13.2}\selectfont
\setlength{\LTleft}{0pt}\setlength{\LTright}{\fill}
\noindent
\begin{tabular}{@{}>{\raggedright\arraybackslash}p{22.0mm}>{\raggedright\arraybackslash}p{31.188mm}>{\raggedright\arraybackslash}p{31.188mm}>{\raggedright\arraybackslash}p{31.188mm}@{}}
\toprule
\textbf{Case} & \textbf{Masculine} & \textbf{Neuter} & \textbf{Feminine} \\
\midrule
Nominative & seald & seald & seald \\
Genitive & sealdes & sealdes & sealdre \\
Dative & sealdum & sealdum & sealdre \\
Accusative & sealdne & seald & sealde \\
\bottomrule
\end{tabular}
\par\endgroup
% course-panel:T33-P01:end
% mdtable:T33:end
\reserve{9\baselineskip}
% mdparagraph:P148:begin
\textbf{Strong plural.}
% mdparagraph:P148:end

% table: NEW-U07-PART-SEALD-2
% mdtable:T34:begin
% course-table-style: 2026-10
% course-panel:T34-P01:begin
\par\begingroup
\fontsize{10.5}{13.2}\selectfont
\setlength{\LTleft}{0pt}\setlength{\LTright}{\fill}
\noindent
\begin{tabular}{@{}>{\raggedright\arraybackslash}p{22.0mm}>{\raggedright\arraybackslash}p{31.188mm}>{\raggedright\arraybackslash}p{31.188mm}>{\raggedright\arraybackslash}p{31.188mm}@{}}
\toprule
\textbf{Case} & \textbf{Masculine} & \textbf{Neuter} & \textbf{Feminine} \\
\midrule
Nominative & sealde & seald & sealda \\
Genitive & sealdra & sealdra & sealdra \\
Dative & sealdum & sealdum & sealdum \\
Accusative & sealde & seald & sealda \\
\bottomrule
\end{tabular}
\par\endgroup
% course-panel:T34-P01:end
% mdtable:T34:end
\reserve{9\baselineskip}
% mdparagraph:P149:begin
\textbf{Weak singular.}
% mdparagraph:P149:end

% table: NEW-U07-PART-SEALD-3
% mdtable:T35:begin
% course-table-style: 2026-10
% course-panel:T35-P01:begin
\par\begingroup
\fontsize{10.5}{13.2}\selectfont
\setlength{\LTleft}{0pt}\setlength{\LTright}{\fill}
\noindent
\begin{tabular}{@{}>{\raggedright\arraybackslash}p{22.0mm}>{\raggedright\arraybackslash}p{31.188mm}>{\raggedright\arraybackslash}p{31.188mm}>{\raggedright\arraybackslash}p{31.188mm}@{}}
\toprule
\textbf{Case} & \textbf{Masculine} & \textbf{Neuter} & \textbf{Feminine} \\
\midrule
Nominative & sealda & sealde & sealde \\
Genitive & sealdan & sealdan & sealdan \\
Dative & sealdan & sealdan & sealdan \\
Accusative & sealdan & sealde & sealdan \\
\bottomrule
\end{tabular}
\par\endgroup
% course-panel:T35-P01:end
% mdtable:T35:end
\reserve{9\baselineskip}
% mdparagraph:P150:begin
\textbf{Weak plural.}
% mdparagraph:P150:end

% table: NEW-U07-PART-SEALD-4
% mdtable:T36:begin
% course-table-style: 2026-10
% course-panel:T36-P01:begin
\par\begingroup
\fontsize{10.5}{13.2}\selectfont
\setlength{\LTleft}{0pt}\setlength{\LTright}{\fill}
\noindent
\begin{tabular}{@{}>{\raggedright\arraybackslash}p{22.0mm}>{\raggedright\arraybackslash}p{31.188mm}>{\raggedright\arraybackslash}p{31.188mm}>{\raggedright\arraybackslash}p{31.188mm}@{}}
\toprule
\textbf{Case} & \textbf{Masculine} & \textbf{Neuter} & \textbf{Feminine} \\
\midrule
Nominative & sealdan & sealdan & sealdan \\
Genitive & sealdra & sealdra & sealdra \\
Dative & sealdum & sealdum & sealdum \\
Accusative & sealdan & sealdan & sealdan \\
\bottomrule
\end{tabular}
\par\endgroup
% course-panel:T36-P01:end
% mdtable:T36:end
% mdparagraph:P151:begin
\textbf{Variants.} Strong feminine plural nominative/\allowbreak{}accusative sealde is an alternative to sealda; strong neuter plural nominative/\allowbreak{}accusative sealde is an alternative to seald.
% mdparagraph:P151:end

\reserve{7\baselineskip}
% mdparagraph:P152:begin
\section*{6. Elementary passive constructions}\BookHeading{nav-ch07-learner-0022}{subsection}{6. Elementary passive constructions}{1}\BookLabel{bk-ch07-learner-8d010318c225}
% mdparagraph:P152:end

% teaching: NEW-U07-VOC-wesan-passive
% teaching: NEW-U07-VOC-weorthan-1
% teaching: NEW-U07-GRAM-U07-PASSIVE
% mdparagraph:P153:begin
The \textbf{passive} makes the affected person or thing the subject. With the transitive verbs used here, an active accusative object corresponds to a nominative passive subject. Use a finite form of \textbf{wesan/\allowbreak{}bēon} or \textbf{weorðan} with the past participle. The auxiliary agrees with the new subject in person and number and carries the tense and mood. The agent can be omitted.
% mdparagraph:P153:end

% example: NEW-U07-EXAMPLE-09
% mdparagraph:P154:begin
Se mann band ðæt dēor. Ðæt dēor wearð bunden.
% mdparagraph:P154:end

% mdparagraph:P155:begin
The man bound the animal. The animal was bound.
% mdparagraph:P155:end

% mdparagraph:P156:begin
Dēor is accusative in the active clause and nominative in the passive, although its neuter singular spelling is unchanged. Wearð is the singular preterite auxiliary; bunden is the past participle.
% mdparagraph:P156:end

% mdparagraph:P157:begin
Wesan/\allowbreak{}bēon often presents a state or result; weorðan often presents an event or becoming. These are useful tendencies, not an absolute division. Do not assign a fixed modern English tense or aspect to every auxiliary–participle combination. The past participle does not carry the finite tense by itself.
% mdparagraph:P157:end

% teaching: NEW-U07-PRED-PART
% teaching: NEW-U07-GRAM-U07-PREDICATE-PARTICIPLE-VARIANTS
\reserve{7\baselineskip}
% mdparagraph:P158:begin
\subsection*{Verbal predicates and adjective agreement}\BookHeading{nav-ch07-learner-0023}{subsubsection}{Verbal predicates and adjective agreement}{0}\BookLabel{bk-ch07-learner-2de790126d8b}
% mdparagraph:P158:end

% mdparagraph:P159:begin
In a verbal passive the past participle may be \textbf{uninflected}, or it may take \textbf{strong masculine-like endings}. This differs from the agreement rule for an attributive adjective. In particular, a neuter plural subject does not require the attributive neuter-plural -u form in every verbal passive.
% mdparagraph:P159:end

% example: NEW-U07-EXAMPLE-08
% mdparagraph:P160:begin
Ðā dēor sindon bunden. Ðā dēor sindon bundene.
% mdparagraph:P160:end

% mdparagraph:P161:begin
The animals are bound. The animals are bound.
% mdparagraph:P161:end

% mdparagraph:P162:begin
Both are licensed verbal passives: bunden is uninflected and bundene has the masculine-like plural ending. The noun dēor remains neuter plural in both.
% mdparagraph:P162:end

% table: U07-DISPLAY-PREDICATE-PART
% mdtable:T37:begin
\par\begingroup
\fontsize{10.5}{13.2}\selectfont
\setlength{\LTleft}{0pt}\setlength{\LTright}{\fill}
\noindent
\begin{tabular}{@{}>{\raggedright\arraybackslash}p{38.41mm}>{\raggedright\arraybackslash}p{32.92mm}>{\raggedright\arraybackslash}p{47.05mm}@{}}
\toprule
\textbf{Verb} & \textbf{Uninflected} & \textbf{Plural masculine-like} \\
\midrule
bindan & bunden & bundene \\
gemanigfieldan & gemanigfield & gemanigfielde \\
\bottomrule
\end{tabular}
\par\endgroup

% mdtable:T37:end
% mdparagraph:P163:begin
The spellings can overlap across constructions. Attributive neuter plural bundenu has the accepted variants bunden and bundene given in §5. A spelling used in the passive is therefore not automatically forbidden before a noun. Identify the construction first; apply its agreement rule or permitted predicate inflection.
% mdparagraph:P163:end

% mdparagraph:P164:begin
Present participles also have uninflected predicate uses in the sources. The independent adjectival practice here uses the full attributive models. Do not infer an English-style progressive system merely from a form in -ende.
% mdparagraph:P164:end

% teaching: NEW-U07-GRAM-U07-PASSIVE-COMMAND-WISH
\reserve{7\baselineskip}
% mdparagraph:P165:begin
\subsection*{Passive commands and wishes}\BookHeading{nav-ch07-learner-0024}{subsubsection}{Passive commands and wishes}{0}\BookLabel{bk-ch07-learner-162da4968438}
% mdparagraph:P165:end

% mdparagraph:P166:begin
Use the imperative of the auxiliary for a direct command, and its present subjunctive for a wish. The participle remains nonfinite. In the first utterance below the speaker explicitly commands one male addressee; the second is a wish about the birds.
% mdparagraph:P166:end

% example: NEW-U07-EXAMPLE-10
% mdparagraph:P167:begin
Bēo geblētsod. Bēon ðā fuglas gemanigfielde.
% mdparagraph:P167:end

% mdparagraph:P168:begin
Be blessed. Let the birds be multiplied.
% mdparagraph:P168:end

% mdparagraph:P169:begin
Bēo is second-singular imperative in the stated command setting; bēon is third-plural subjunctive with ðā fuglas. Bēo alone could also be a subjunctive spelling, so the command analysis depends on the setting. The source-style plural gemanigfielde is a past participle, not a finite verb.
% mdparagraph:P169:end

% mdparagraph:P170:begin
This unit uses ordinary passives of the selected transitive verbs. It does not add agent phrases, a general perfect system, or a general progressive system. A participle of an intransitive verb may have a different function; the Creation closing formula is supplied locally in §10.
% mdparagraph:P170:end

\reserve{7\baselineskip}
% mdparagraph:P171:begin
\subsection*{Practice P. Participles and voice}\BookHeading{nav-ch07-learner-0025}{subsubsection}{Practice P. Participles and voice}{0}
% mdparagraph:P171:end

% mdparagraph:P172:begin
Use the full adjective models where agreement is requested. For a verbal passive, distinguish required subject–auxiliary agreement from the permitted forms of its predicate participle.
% mdparagraph:P172:end

\reserve{5\baselineskip}
\par\begin{minipage}{\linewidth}\raggedright
% mdparagraph:P173:begin
\textbf{P1.}
% mdparagraph:P173:end

% item: NEW-U07-P01
% mdparagraph:P174:begin
Form the present participle of bindan, nęrian, lufian, and cweðan. Give the uninflected dictionary form and identify the stem retained before -ende.
% mdparagraph:P174:end

\end{minipage}\par\addvspace{4pt}

\par\begin{minipage}{\linewidth}\raggedright
% mdparagraph:P175:begin
\textbf{P2.}
% mdparagraph:P175:end

% item: NEW-U07-P02
% mdparagraph:P176:begin
Give the selected past participle of bindan, beran, hīeran, and gehālgian. Explain why the first two cannot be formed by adding a weak dental ending.
% mdparagraph:P176:end

\end{minipage}\par\addvspace{4pt}

\par\begin{minipage}{\linewidth}\raggedright
% mdparagraph:P177:begin
\textbf{P3.}
% mdparagraph:P177:end

% item: NEW-U07-P03
% mdparagraph:P178:begin
Supply the agreeing form of bindende: (a) a binding man, nominative; (b) the binding man, nominative; (c) with a binding woman, dative feminine singular only; (d) of the binding men, genitive plural only. For (c) and (d), give the participial form only; no new noun or preposition is required.
% mdparagraph:P178:end

\end{minipage}\par\addvspace{4pt}

\par\begin{minipage}{\linewidth}\raggedright
% mdparagraph:P179:begin
\textbf{P4.}
% mdparagraph:P179:end

% item: NEW-U07-P04
% mdparagraph:P180:begin
Supply the agreeing form of bunden: (a) a bound animal, accusative neuter singular; (b) the bound animal, accusative neuter singular; (c) of bound birds, genitive plural; (d) to the bound birds, dative plural. Identify strong or weak inflection in each.
% mdparagraph:P180:end

\end{minipage}\par\addvspace{4pt}

\par\begin{minipage}{\linewidth}\raggedright
% mdparagraph:P181:begin
\textbf{P5.}
% mdparagraph:P181:end

% item: NEW-U07-P05
% mdparagraph:P182:begin
Change Hē bint ðæt dēor into a present passive with wesan. Omit the agent. Identify the subject and the case of dēor before and after the change.
% mdparagraph:P182:end

\end{minipage}\par\addvspace{4pt}

\par\begin{minipage}{\linewidth}\raggedright
% mdparagraph:P183:begin
\textbf{P6.}
% mdparagraph:P183:end

% item: NEW-U07-P06
% mdparagraph:P184:begin
Change Hīe bundon ðā dēor into a past passive with weorðan, describing the binding event. Omit the agent.
% mdparagraph:P184:end

\end{minipage}\par\addvspace{4pt}

\par\begin{minipage}{\linewidth}\raggedright
% mdparagraph:P185:begin
\textbf{P7.}
% mdparagraph:P185:end

% item: NEW-U07-P07
% mdparagraph:P186:begin
Explain why both Ðā dēor sindon bunden and Ðā dēor sindon bundene can serve as verbal passives. Contrast these with an attributive participle directly modifying a neuter plural noun.
% mdparagraph:P186:end

\end{minipage}\par\addvspace{4pt}

\par\begin{minipage}{\linewidth}\raggedright
% mdparagraph:P187:begin
\textbf{P8.}
% mdparagraph:P187:end

% item: NEW-U07-P08
% mdparagraph:P188:begin
The first utterance is explicitly a command to one male addressee: Bēo geblētsod. The second is a wish about the birds: Bēon ðā fuglas gemanigfielde. Parse each auxiliary and participle, explain the addressee/\allowbreak{}subject difference, and translate both. Note any formal ambiguity that the stated context resolves.
% mdparagraph:P188:end

\end{minipage}\par\addvspace{4pt}

\reserve{7\baselineskip}
% mdparagraph:P189:begin
\section*{7. Quantifiers, places, and functions}\BookHeading{nav-ch07-learner-0026}{subsection}{7. Quantifiers, places, and functions}{1}
% mdparagraph:P189:end

% teaching: NEW-U07-VOC-eall-all
% teaching: NEW-U07-VOC-aelc-each
% teaching: NEW-U07-GRAM-U07-QUANTIFIER-AGREEMENT
\reserve{7\baselineskip}
% mdparagraph:P190:begin
\subsection*{Eall and ǣlc}\BookHeading{nav-ch07-learner-0027}{subsubsection}{Eall and ǣlc}{0}\BookLabel{bk-ch07-learner-6ca488fddd27}
% mdparagraph:P190:end

% mdparagraph:P191:begin
\textbf{Eall}, “all; whole”, and \textbf{ǣlc}, “each; every”, agree with the noun they qualify. Learn their complete strong forms. Eall can refer to a whole singular thing or to all the members of a plural group. The ordinary “each/\allowbreak{}every” use of ǣlc in this unit takes a \textbf{singular noun}. Its full plural morphology is given, but no new plural meaning or partitive construction is required in independent translation.
% mdparagraph:P191:end

% teaching: NEW-U07-NOM-eall
\reserve{9\baselineskip}
% mdparagraph:P192:begin
\textbf{Eall: singular.}
% mdparagraph:P192:end

% table: NEW-U07-NOM-eall-1
% mdtable:T38:begin
% course-table-style: 2026-10
% course-panel:T38-P01:begin
\par\begingroup
\fontsize{10.5}{13.2}\selectfont
\setlength{\LTleft}{0pt}\setlength{\LTright}{\fill}
\noindent
\begin{tabular}{@{}>{\raggedright\arraybackslash}p{22.0mm}>{\raggedright\arraybackslash}p{31.188mm}>{\raggedright\arraybackslash}p{31.188mm}>{\raggedright\arraybackslash}p{31.188mm}@{}}
\toprule
\textbf{Case} & \textbf{Masculine} & \textbf{Neuter} & \textbf{Feminine} \\
\midrule
Nominative & eall & eall & eall \\
Genitive & ealles & ealles & ealre \\
Dative & eallum & eallum & ealre \\
Accusative & ealne & eall & ealle \\
\bottomrule
\end{tabular}
\par\endgroup
% course-panel:T38-P01:end
% mdtable:T38:end
\reserve{9\baselineskip}
% mdparagraph:P193:begin
\textbf{Eall: plural.}
% mdparagraph:P193:end

% table: NEW-U07-NOM-eall-2
% mdtable:T39:begin
% course-table-style: 2026-10
% course-panel:T39-P01:begin
\par\begingroup
\fontsize{10.5}{13.2}\selectfont
\setlength{\LTleft}{0pt}\setlength{\LTright}{\fill}
\noindent
\begin{tabular}{@{}>{\raggedright\arraybackslash}p{22.0mm}>{\raggedright\arraybackslash}p{31.188mm}>{\raggedright\arraybackslash}p{31.188mm}>{\raggedright\arraybackslash}p{31.188mm}@{}}
\toprule
\textbf{Case} & \textbf{Masculine} & \textbf{Neuter} & \textbf{Feminine} \\
\midrule
Nominative & ealle & eall & ealla \\
Genitive & ealra & ealra & ealra \\
Dative & eallum & eallum & eallum \\
Accusative & ealle & eall & ealla \\
\bottomrule
\end{tabular}
\par\endgroup
% course-panel:T39-P01:end
% mdtable:T39:end
% mdparagraph:P194:begin
Preserve the single l in ealne, ealre, and ealra, beside eall- in the other printed forms. These are agreement forms, not an invariant English “all”.
% mdparagraph:P194:end

% teaching: NEW-U07-NOM-aelc
\reserve{9\baselineskip}
% mdparagraph:P195:begin
\textbf{Ǣlc: singular.}
% mdparagraph:P195:end

% table: NEW-U07-NOM-aelc-1
% mdtable:T40:begin
% course-table-style: 2026-10
% course-panel:T40-P01:begin
\par\begingroup
\fontsize{10.5}{13.2}\selectfont
\setlength{\LTleft}{0pt}\setlength{\LTright}{\fill}
\noindent
\begin{tabular}{@{}>{\raggedright\arraybackslash}p{22.0mm}>{\raggedright\arraybackslash}p{31.188mm}>{\raggedright\arraybackslash}p{31.188mm}>{\raggedright\arraybackslash}p{31.188mm}@{}}
\toprule
\textbf{Case} & \textbf{Masculine} & \textbf{Neuter} & \textbf{Feminine} \\
\midrule
Nominative & ǣlc & ǣlc & ǣlc \\
Genitive & ǣlces & ǣlces & ǣlcre \\
Dative & ǣlcum & ǣlcum & ǣlcre \\
Accusative & ǣlcne & ǣlc & ǣlce \\
\bottomrule
\end{tabular}
\par\endgroup
% course-panel:T40-P01:end
% mdtable:T40:end
\reserve{9\baselineskip}
% mdparagraph:P196:begin
\textbf{Ǣlc: plural morphology.}
% mdparagraph:P196:end

% table: NEW-U07-NOM-aelc-2
% mdtable:T41:begin
% course-table-style: 2026-10
% course-panel:T41-P01:begin
\par\begingroup
\fontsize{10.5}{13.2}\selectfont
\setlength{\LTleft}{0pt}\setlength{\LTright}{\fill}
\noindent
\begin{tabular}{@{}>{\raggedright\arraybackslash}p{22.0mm}>{\raggedright\arraybackslash}p{31.188mm}>{\raggedright\arraybackslash}p{31.188mm}>{\raggedright\arraybackslash}p{31.188mm}@{}}
\toprule
\textbf{Case} & \textbf{Masculine} & \textbf{Neuter} & \textbf{Feminine} \\
\midrule
Nominative & ǣlce & ǣlc & ǣlca \\
Genitive & ǣlcra & ǣlcra & ǣlcra \\
Dative & ǣlcum & ǣlcum & ǣlcum \\
Accusative & ǣlce & ǣlc & ǣlca \\
\bottomrule
\end{tabular}
\par\endgroup
% course-panel:T41-P01:end
% mdtable:T41:end
% mdparagraph:P197:begin
The vowel is long ǣ throughout. Use the singular case required by the noun’s role when translating “each day” or “every bird”. Displaying the plural endings does not change that ordinary construction. These are strong paradigms; no new weak quantifier system or instrumental use is introduced.
% mdparagraph:P197:end

\reserve{7\baselineskip}
% mdparagraph:P198:begin
\subsection*{Geweald: dominion}\BookHeading{nav-ch07-learner-0028}{subsubsection}{Geweald: dominion}{0}
% mdparagraph:P198:end

% teaching: NEW-U07-NOM-geweald
% teaching: NEW-U07-VOC-geweald-dominion
% mdparagraph:P199:begin
\textbf{Geweald} is a strong neuter noun, “control; rule; dominion”. Use the singular for the ordinary meaning in this unit. Its complete declension includes the plural for recognition and morphological reference; translating independent plural “dominions” is not required.
% mdparagraph:P199:end

% table: NEW-U07-NOM-geweald
% mdtable:T42:begin
% course-table-style: 2026-10
% course-panel:T42-P01:begin
\par\begingroup
\fontsize{10.5}{13.2}\selectfont
\setlength{\LTleft}{0pt}\setlength{\LTright}{\fill}
\noindent
\begin{tabular}{@{}>{\raggedright\arraybackslash}p{39.459mm}>{\raggedright\arraybackslash}p{39.459mm}>{\raggedright\arraybackslash}p{39.459mm}@{}}
\toprule
\textbf{Case} & \textbf{Singular} & \textbf{Plural (morphology reference)} \\
\midrule
Nominative & geweald & geweald \\
Genitive & gewealdes & gewealda \\
Dative & gewealde & gewealdum \\
Accusative & geweald & geweald \\
\bottomrule
\end{tabular}
\par\endgroup
% course-panel:T42-P01:end
% mdtable:T42:end
% mdparagraph:P200:begin
The long neuter stem has the same nominative and accusative spelling in singular and plural. The genitive singular is -es and dative singular -e; genitive plural -a and dative plural -um distinguish the oblique plural forms.
% mdparagraph:P200:end

\reserve{7\baselineskip}
% mdparagraph:P201:begin
\subsection*{Bearing, location, and assigned function}\BookHeading{nav-ch07-learner-0029}{subsubsection}{Bearing, location, and assigned function}{0}
% mdparagraph:P201:end

% teaching: NEW-U07-VOC-beran-produce
% teaching: NEW-U07-GRAM-U07-BERAN-ACCUSATIVE-SEED-FRUIT
% mdparagraph:P202:begin
\textbf{Beran} can mean “bear” in the sense \textbf{produce seed or fruit}, with the seed or fruit as accusative object. This use is additional to the earlier literal “carry”. Wyrcean “yield fruit” likewise takes an accusative fruit object. Keep the short vowel of wæstm, “fruit”.
% mdparagraph:P202:end

% teaching: NEW-U07-VOC-ofer
% teaching: NEW-U07-VOC-ofer-authority
% teaching: NEW-U07-VOC-ofer-1
% teaching: NEW-U07-VOC-bufan
% teaching: NEW-U07-VOC-bufan-above
% teaching: NEW-U07-VOC-under
% teaching: NEW-U07-VOC-under-below
% teaching: NEW-U07-VOC-tomiddes
% teaching: NEW-U07-VOC-tomiddes-midst
% teaching: NEW-U07-VOC-to
% teaching: NEW-U07-VOC-to-function-purpose
% teaching: NEW-U07-GRAM-U07-OFER-ACCUSATIVE
% teaching: NEW-U07-GRAM-U07-SPATIAL-OFER
% teaching: NEW-U07-GRAM-U07-BUFAN-DATIVE
% teaching: NEW-U07-GRAM-U07-UNDER-DATIVE
% teaching: NEW-U07-GRAM-U07-TOMIDDES-DATIVE
% teaching: NEW-U07-GRAM-U07-TO-DATIVE
% teaching: NEW-U07-GRAM-U07-DATIVE-FUNCTION-BENEFICIARY
% mdparagraph:P203:begin
The following prepositions have the stated meanings and cases here. A familiar spelling does not make all its uses interchangeable.
% mdparagraph:P203:end

% table: U07-DISPLAY-FUNCTION-WORDS
% mdtable:T43:begin
\par\begingroup
\fontsize{10.5}{13.2}\selectfont
\setlength{\LTleft}{0pt}\setlength{\LTright}{\fill}
\noindent
\begin{tabular}{@{}>{\raggedright\arraybackslash}p{21.58mm}>{\raggedright\arraybackslash}p{33.02mm}>{\raggedright\arraybackslash}p{63.77mm}@{}}
\toprule
\textbf{Word}\BookPoint{bk-ch07-learner-bae927b6adfb} & \textbf{Meaning here} & \textbf{Construction} \\
\midrule
ofer & over, in authority over & accusative in the authority frame \\
ofer & spatially over, across & accusative or dative, according to the construction \\
bufan & above & dative landmark \\
under & under, below & dative landmark in the location use taught here \\
tōmiddes & in the midst of & dative landmark \\
tō & for, serving as & dative assigned function, such as food or a sign \\
\bottomrule
\end{tabular}
\par\endgroup

% mdtable:T43:end
% mdparagraph:P204:begin
With\BookPoint{bk-ch07-learner-844cd74c2c5d} spatial ofer, the accusative commonly marks direction over or across a landmark, while the dative commonly marks location. This is a tendency, not an exceptionless rule for every use. The source’s ofer eorðan cannot by its spelling alone identify accusative rather than dative: eorðan serves both cases. The authority use is separately taught with the accusative.
% mdparagraph:P204:end

% mdparagraph:P205:begin
Tō\BookPoint{bk-ch07-learner-7c9ec2ca6444} with a dative noun can state what something is \textbf{for} or \textbf{serves as}. A separate dative can name the beneficiary: in ēow tō męte, “to you for food”, ēow is the beneficiary and męte is the assigned function. Neither is a possessive adjective. This tō is different from the earlier tō with an inflected purpose infinitive.
% mdparagraph:P205:end

\reserve{7\baselineskip}
% mdparagraph:P206:begin
\subsection*{Naming}\BookHeading{nav-ch07-learner-0030}{subsubsection}{Naming}{0}\BookLabel{bk-ch07-learner-645dc7063f24}
% mdparagraph:P206:end

% teaching: NEW-U07-GRAM-U07-NAMING-TWO-ACCUSATIVES
% mdparagraph:P207:begin
Hātan “call, name” can take \textbf{two accusatives}: the thing named and the name given to it. In the first sentence below the firmament is the thing named and Heaven the assigned name. The second sentence shows tō with a dative function in a wish.
% mdparagraph:P207:end

% example: NEW-U07-EXAMPLE-11
% mdparagraph:P208:begin
Ic hāte ðā fæstnesse Heofon. Bēo ðæt lēoht tō tācne.
% mdparagraph:P208:end

% mdparagraph:P209:begin
I call the firmament Heaven. Let the light be for a sign.
% mdparagraph:P209:end

% mdparagraph:P210:begin
Ðā fæstnesse and Heofon are both accusative. In the wish, ðæt lēoht is nominative subject and tācne is dative after tō. The “be called” passive uses the learned auxiliary with hāten; a separate hātte system is not needed for this unit.
% mdparagraph:P210:end

\reserve{7\baselineskip}
% mdparagraph:P211:begin
\subsection*{Adverbs and a presentational word}\BookHeading{nav-ch07-learner-0031}{subsubsection}{Adverbs and a presentational word}{0}
% mdparagraph:P211:end

% teaching: NEW-U07-VOC-nu
% teaching: NEW-U07-VOC-nu-1
% teaching: NEW-U07-VOC-swa
% teaching: NEW-U07-VOC-swa-1
% teaching: NEW-U07-VOC-thus
% teaching: NEW-U07-VOC-thus-thus-source
% teaching: NEW-U07-VOC-eac
% teaching: NEW-U07-VOC-eac-also
% teaching: NEW-U07-VOC-swilce
% teaching: NEW-U07-VOC-swilce-likewise
% teaching: NEW-U07-VOC-efne
% teaching: NEW-U07-VOC-efne-behold
% teaching: NEW-U07-GRAM-U07-ADVERBS-INTERJECTION
% mdtable:T44:begin
\par\begingroup
\fontsize{10.5}{13.2}\selectfont
\setlength{\LTleft}{0pt}\setlength{\LTright}{\fill}
\noindent
\begin{tabular}{@{}>{\raggedright\arraybackslash}p{25.49mm}>{\raggedright\arraybackslash}p{95.7mm}@{}}
\toprule
\textbf{Word} & \textbf{Meaning and use} \\
\midrule
nū & now; can open a statement with is before its following nominal subject \\
swā & so, in that way \\
ðus & thus, in this way \\
ēac & also \\
swilce & likewise; the adverb, not a newly taught conjunction \\
efne & behold; introduces the following statement \\
gif & if; introduces the condition described in §4 \\
\bottomrule
\end{tabular}
\par\endgroup

% mdtable:T44:end
% mdparagraph:P212:begin
All these forms are invariant. The source also uses ēac swilce, “also”. The already learned swīðe means “very” before an adjective. Keep the adjective’s own agreement: the intensifier does not alter it.
% mdparagraph:P212:end

% mdparagraph:P213:begin
\textbf{Predicate nouns.}\BookPoint{bk-ch07-learner-73c8a76ced55} The finite verb agrees with its subject. A predicate noun may have its own number: in “Birds are the food of the animal”, the birds are plural but “food” is singular męte. The rule requiring an adjective to agree in number does not turn a singular predicate noun into a plural.
% mdparagraph:P213:end

\reserve{7\baselineskip}
% mdparagraph:P214:begin
\section*{8. Translation practice}\BookHeading{nav-ch07-learner-0032}{subsection}{8. Translation practice}{1}
% mdparagraph:P214:end

\reserve{7\baselineskip}
% mdparagraph:P215:begin
\subsection*{Practice C. Controlled composition}\BookHeading{nav-ch07-learner-0033}{subsubsection}{Practice C. Controlled composition}{0}
% mdparagraph:P215:end

% mdparagraph:P216:begin
Use the stated command, wish, or conditional target. Retain the case required by the verb or preposition. The last item contrasts singular “every” with plural “all”.
% mdparagraph:P216:end

\reserve{5\baselineskip}
\par\begin{minipage}{\linewidth}\raggedright
% mdparagraph:P217:begin
\textbf{C1.}
% mdparagraph:P217:end

% item: NEW-U07-C01
% mdparagraph:P218:begin
Translate: May the birds live.
% mdparagraph:P218:end

\end{minipage}\par\addvspace{4pt}

\par\begin{minipage}{\linewidth}\raggedright
% mdparagraph:P219:begin
\textbf{C2.}
% mdparagraph:P219:end

% item: NEW-U07-C02
% mdparagraph:P220:begin
Translate: Let us help the companions.
% mdparagraph:P220:end

\end{minipage}\par\addvspace{4pt}

\par\begin{minipage}{\linewidth}\raggedright
% mdparagraph:P221:begin
\textbf{C3.}
% mdparagraph:P221:end

% item: NEW-U07-C03
% mdparagraph:P222:begin
Translate, presenting the condition as tentative: If you (sing.) have food, give it to the birds.
% mdparagraph:P222:end

\end{minipage}\par\addvspace{4pt}

\par\begin{minipage}{\linewidth}\raggedright
% mdparagraph:P223:begin
\textbf{C4.}
% mdparagraph:P223:end

% item: NEW-U07-C04
% mdparagraph:P224:begin
Translate as a command to one male addressee: Be blessed.
% mdparagraph:P224:end

\end{minipage}\par\addvspace{4pt}

\par\begin{minipage}{\linewidth}\raggedright
% mdparagraph:P225:begin
\textbf{C5.}
% mdparagraph:P225:end

% item: NEW-U07-C05
% mdparagraph:P226:begin
Translate as a wish: Let the work be accomplished.
% mdparagraph:P226:end

\end{minipage}\par\addvspace{4pt}

\par\begin{minipage}{\linewidth}\raggedright
% mdparagraph:P227:begin
\textbf{C6.}
% mdparagraph:P227:end

% item: NEW-U07-C06
% mdparagraph:P228:begin
Translate: Every bird is above the earth; all trees bear fruit.
% mdparagraph:P228:end

\end{minipage}\par\addvspace{4pt}

\reserve{7\baselineskip}
% mdparagraph:P229:begin
\subsection*{Cook’s exercises}\BookHeading{nav-ch07-learner-0034}{subsubsection}{Cook’s exercises}{0}
% mdparagraph:P229:end

% mdparagraph:P230:begin
Translate into Old English. The original wording and numbering are retained. For English “let” or “may” expressing a wish, use the subjunctive; direct commands use the imperative. Use the taught auxiliary–participle construction for passives and the naming construction for “call”. Keep direct speech as direct speech in I.7.17.
% mdparagraph:P230:end

% mdparagraph:P231:begin
Where English “you” has no number specified, either singular or plural is possible; make the pronouns and related forms consistent. The parenthetical number directions in the original prompts do specify a target. An otherwise valid translation can differ from a model’s chosen wording or number.
% mdparagraph:P231:end

\reserve{16\baselineskip}
% mdparagraph:P232:begin
\subsection*{Cook I.4: the remaining item}\BookHeading{nav-ch07-learner-0035}{subsubsection}{Cook I.4: the remaining item}{0}
% mdparagraph:P232:end

\reserve{5\baselineskip}
\par\begin{minipage}{\linewidth}\raggedright
% mdparagraph:P233:begin
\textbf{Cook I.4.18.}
% mdparagraph:P233:end

% item: EX-I-04-18
% mdparagraph:P234:begin
Increase (sing.) and be multiplied.
% mdparagraph:P234:end

\end{minipage}\par\addvspace{4pt}

\reserve{7\baselineskip}
% mdparagraph:P235:begin
\subsection*{Cook I.6}\BookHeading{nav-ch07-learner-0036}{subsubsection}{Cook I.6}{0}
% mdparagraph:P235:end

\reserve{5\baselineskip}
\par\begin{minipage}{\linewidth}\raggedright
% mdparagraph:P236:begin
\textbf{Cook I.6.1.}
% mdparagraph:P236:end

% item: EX-I-06-01
% mdparagraph:P237:begin
The animal lives.
% mdparagraph:P237:end

\end{minipage}\par\addvspace{4pt}

\par\begin{minipage}{\linewidth}\raggedright
% mdparagraph:P238:begin
\textbf{Cook I.6.2.}
% mdparagraph:P238:end

% item: EX-I-06-02
% mdparagraph:P239:begin
The bird flies.
% mdparagraph:P239:end

\end{minipage}\par\addvspace{4pt}

\par\begin{minipage}{\linewidth}\raggedright
% mdparagraph:P240:begin
\textbf{Cook I.6.3.}
% mdparagraph:P240:end

% item: EX-I-06-03
% mdparagraph:P241:begin
The tree grows.
% mdparagraph:P241:end

\end{minipage}\par\addvspace{4pt}

\par\begin{minipage}{\linewidth}\raggedright
% mdparagraph:P242:begin
\textbf{Cook I.6.4.}
% mdparagraph:P242:end

% item: EX-I-06-04
% mdparagraph:P243:begin
The man rules over the cattle of the earth.
% mdparagraph:P243:end

\end{minipage}\par\addvspace{4pt}

\par\begin{minipage}{\linewidth}\raggedright
% mdparagraph:P244:begin
\textbf{Cook I.6.5.}
% mdparagraph:P244:end

% item: EX-I-06-05
% mdparagraph:P245:begin
Let him be over the birds and (over) the beasts.
% mdparagraph:P245:end

\end{minipage}\par\addvspace{4pt}

\par\begin{minipage}{\linewidth}\raggedright
% mdparagraph:P246:begin
\textbf{Cook I.6.6.}
% mdparagraph:P246:end

% item: EX-I-06-06
% mdparagraph:P247:begin
The stars are above the sea.
% mdparagraph:P247:end

\end{minipage}\par\addvspace{4pt}

\par\begin{minipage}{\linewidth}\raggedright
% mdparagraph:P248:begin
\textbf{Cook I.6.7.}
% mdparagraph:P248:end

% item: EX-I-06-07
% mdparagraph:P249:begin
It is for your food [to you for food].
% mdparagraph:P249:end

\end{minipage}\par\addvspace{4pt}

\par\begin{minipage}{\linewidth}\raggedright
% mdparagraph:P250:begin
\textbf{Cook I.6.8.}
% mdparagraph:P250:end

% item: EX-I-06-08
% mdparagraph:P251:begin
The herb bears seed.
% mdparagraph:P251:end

\end{minipage}\par\addvspace{4pt}

\par\begin{minipage}{\linewidth}\raggedright
% mdparagraph:P252:begin
\textbf{Cook I.6.9.}
% mdparagraph:P252:end

% item: EX-I-06-09
% mdparagraph:P253:begin
The tree yields fruit.
% mdparagraph:P253:end

\end{minipage}\par\addvspace{4pt}

\par\begin{minipage}{\linewidth}\raggedright
% mdparagraph:P254:begin
\textbf{Cook I.6.10.}
% mdparagraph:P254:end

% item: EX-I-06-10
% mdparagraph:P255:begin
The air is very good.
% mdparagraph:P255:end

\end{minipage}\par\addvspace{4pt}

\par\begin{minipage}{\linewidth}\raggedright
% mdparagraph:P256:begin
\textbf{Cook I.6.11.}
% mdparagraph:P256:end

% item: EX-I-06-11
% mdparagraph:P257:begin
The things are very good.
% mdparagraph:P257:end

\end{minipage}\par\addvspace{4pt}

\par\begin{minipage}{\linewidth}\raggedright
% mdparagraph:P258:begin
\textbf{Cook I.6.12.}
% mdparagraph:P258:end

% item: EX-I-06-12
% mdparagraph:P259:begin
The birds were very good.
% mdparagraph:P259:end

\end{minipage}\par\addvspace{4pt}

\par\begin{minipage}{\linewidth}\raggedright
% mdparagraph:P260:begin
\textbf{Cook I.6.13.}
% mdparagraph:P260:end

% item: EX-I-06-13
% mdparagraph:P261:begin
The light of the stars was very good.
% mdparagraph:P261:end

\end{minipage}\par\addvspace{4pt}

\par\begin{minipage}{\linewidth}\raggedright
% mdparagraph:P262:begin
\textbf{Cook I.6.14.}
% mdparagraph:P262:end

% item: EX-I-06-14
% mdparagraph:P263:begin
He gave them all trees.
% mdparagraph:P263:end

\end{minipage}\par\addvspace{4pt}

\par\begin{minipage}{\linewidth}\raggedright
% mdparagraph:P264:begin
\textbf{Cook I.6.15.}
% mdparagraph:P264:end

% item: EX-I-06-15
% mdparagraph:P265:begin
I give you all things upon earth.
% mdparagraph:P265:end

\end{minipage}\par\addvspace{4pt}

\par\begin{minipage}{\linewidth}\raggedright
% mdparagraph:P266:begin
\textbf{Cook I.6.16.}
% mdparagraph:P266:end

% item: EX-I-06-16
% mdparagraph:P267:begin
I call the firmament Heaven.
% mdparagraph:P267:end

\end{minipage}\par\addvspace{4pt}

\reserve{7\baselineskip}
% mdparagraph:P268:begin
\subsection*{Cook I.7}\BookHeading{nav-ch07-learner-0037}{subsubsection}{Cook I.7}{0}
% mdparagraph:P268:end

\reserve{5\baselineskip}
\par\begin{minipage}{\linewidth}\raggedright
% mdparagraph:P269:begin
\textbf{Cook I.7.1.}
% mdparagraph:P269:end

% item: EX-I-07-01
% mdparagraph:P270:begin
He has dominion.
% mdparagraph:P270:end

\end{minipage}\par\addvspace{4pt}

\par\begin{minipage}{\linewidth}\raggedright
% mdparagraph:P271:begin
\textbf{Cook I.7.2.}
% mdparagraph:P271:end

% item: EX-I-07-02
% mdparagraph:P272:begin
He blesses the earth.
% mdparagraph:P272:end

\end{minipage}\par\addvspace{4pt}

\par\begin{minipage}{\linewidth}\raggedright
% mdparagraph:P273:begin
\textbf{Cook I.7.3.}
% mdparagraph:P273:end

% item: EX-I-07-03
% mdparagraph:P274:begin
He says thus.
% mdparagraph:P274:end

\end{minipage}\par\addvspace{4pt}

\par\begin{minipage}{\linewidth}\raggedright
% mdparagraph:P275:begin
\textbf{Cook I.7.4.}
% mdparagraph:P275:end

% item: EX-I-07-04
% mdparagraph:P276:begin
Now is the season.
% mdparagraph:P276:end

\end{minipage}\par\addvspace{4pt}

\par\begin{minipage}{\linewidth}\raggedright
% mdparagraph:P277:begin
\textbf{Cook I.7.5.}
% mdparagraph:P277:end

% item: EX-I-07-05
% mdparagraph:P278:begin
It was so called.
% mdparagraph:P278:end

\end{minipage}\par\addvspace{4pt}

\par\begin{minipage}{\linewidth}\raggedright
% mdparagraph:P279:begin
\textbf{Cook I.7.6.}
% mdparagraph:P279:end

% item: EX-I-07-06
% mdparagraph:P280:begin
Behold, God rules over the earth.
% mdparagraph:P280:end

\end{minipage}\par\addvspace{4pt}

\par\begin{minipage}{\linewidth}\raggedright
% mdparagraph:P281:begin
\textbf{Cook I.7.7.}
% mdparagraph:P281:end

% item: EX-I-07-07
% mdparagraph:P282:begin
The tree is in the midst of the grass.
% mdparagraph:P282:end

\end{minipage}\par\addvspace{4pt}

\par\begin{minipage}{\linewidth}\raggedright
% mdparagraph:P283:begin
\textbf{Cook I.7.8.}
% mdparagraph:P283:end

% item: EX-I-07-08
% mdparagraph:P284:begin
The earth is in the midst of the air.
% mdparagraph:P284:end

\end{minipage}\par\addvspace{4pt}

\par\begin{minipage}{\linewidth}\raggedright
% mdparagraph:P285:begin
\textbf{Cook I.7.9.}
% mdparagraph:P285:end

% item: EX-I-07-09
% mdparagraph:P286:begin
The waters are under the firmament.
% mdparagraph:P286:end

\end{minipage}\par\addvspace{4pt}

\par\begin{minipage}{\linewidth}\raggedright
% mdparagraph:P287:begin
\textbf{Cook I.7.10.}
% mdparagraph:P287:end

% item: EX-I-07-10
% mdparagraph:P288:begin
Birds are the food of the animal.
% mdparagraph:P288:end

\end{minipage}\par\addvspace{4pt}

\par\begin{minipage}{\linewidth}\raggedright
% mdparagraph:P289:begin
\textbf{Cook I.7.11.}
% mdparagraph:P289:end

% item: EX-I-07-11
% mdparagraph:P290:begin
Divide (plur.) the waters from the waters.
% mdparagraph:P290:end

\end{minipage}\par\addvspace{4pt}

\par\begin{minipage}{\linewidth}\raggedright
% mdparagraph:P291:begin
\textbf{Cook I.7.12.}
% mdparagraph:P291:end

% item: EX-I-07-12
% mdparagraph:P292:begin
God's dominion is over all creatures.
% mdparagraph:P292:end

\end{minipage}\par\addvspace{4pt}

\par\begin{minipage}{\linewidth}\raggedright
% mdparagraph:P293:begin
\textbf{Cook I.7.13.}
% mdparagraph:P293:end

% item: EX-I-07-13
% mdparagraph:P294:begin
God hallows each day.
% mdparagraph:P294:end

\end{minipage}\par\addvspace{4pt}

\par\begin{minipage}{\linewidth}\raggedright
% mdparagraph:P295:begin
\textbf{Cook I.7.14.}
% mdparagraph:P295:end

% item: EX-I-07-14
% mdparagraph:P296:begin
Let the birds increase and be multiplied.
% mdparagraph:P296:end

\end{minipage}\par\addvspace{4pt}

\par\begin{minipage}{\linewidth}\raggedright
% mdparagraph:P297:begin
\textbf{Cook I.7.15.}
% mdparagraph:P297:end

% item: EX-I-07-15
% mdparagraph:P298:begin
Let the men swim in the sea.
% mdparagraph:P298:end

\end{minipage}\par\addvspace{4pt}

\par\begin{minipage}{\linewidth}\raggedright
% mdparagraph:P299:begin
\textbf{Cook I.7.16.}
% mdparagraph:P299:end

% item: EX-I-07-16
% mdparagraph:P300:begin
Let the tree bear fruit.
% mdparagraph:P300:end

\end{minipage}\par\addvspace{4pt}

\par\begin{minipage}{\linewidth}\raggedright
% mdparagraph:P301:begin
\textbf{Cook I.7.17.}
% mdparagraph:P301:end

% item: EX-I-07-17
% mdparagraph:P302:begin
He says likewise, Let the light be for a sign.
% mdparagraph:P302:end

\end{minipage}\par\addvspace{4pt}

\reserve{7\baselineskip}
% mdparagraph:P303:begin
\section*{9. Binding the animals}\BookHeading{nav-ch07-learner-0038}{subsection}{9. Binding the animals}{1}
% mdparagraph:P303:end

% teaching: NEW-U07-GRAM-U07-QUOTATION-WISH
% mdparagraph:P304:begin
\textbf{New composition.} The two speeches distinguish a tentative condition, an assertion of ability, willingness to help, and a joint exhortation. Pronouns inside each quotation belong to that speaker’s setting.
% mdparagraph:P304:end

% composition: NEW-U07-TEXT-A
% mdparagraph:P305:begin
\par\begin{minipage}{\linewidth}\raggedright
\readingtext{Twēgen męnn wǣron on ðǣm strande. Hīe hæfdon fīf dēor. Se ǣresta mann cwæð: “Gif ðū mē helpe, wē magon ðā dēor bindan.” Se ōðer mann cwæð: “Ic wille ðē helpan. Binden wē ðā dēor.” Hīe bundon ðā dēor, and se ǣresta mann sealde his gefēran męte. Ðā dēor wǣron bunden, and ðā męnn hæfdon męte.}
\end{minipage}\par
% mdparagraph:P305:end

\reserve{7\baselineskip}
% mdparagraph:P306:begin
\subsection*{Practice D. The connected account}\BookHeading{nav-ch07-learner-0039}{subsubsection}{Practice D. The connected account}{0}
% mdparagraph:P306:end

% mdparagraph:P307:begin
Read the account as a whole before identifying individual forms. Use the sequence of actions and speakers to establish the participants referred to.
% mdparagraph:P307:end

\reserve{5\baselineskip}
\par\begin{minipage}{\linewidth}\raggedright
% mdparagraph:P308:begin
\textbf{D1.}
% mdparagraph:P308:end

% item: NEW-U07-D01
% mdparagraph:P309:begin
Translate the complete new account “Binding the animals.” Keep the speakers and the final result clear.
% mdparagraph:P309:end

\end{minipage}\par\addvspace{4pt}

\par\begin{minipage}{\linewidth}\raggedright
% mdparagraph:P310:begin
\textbf{D2.}
% mdparagraph:P310:end

% item: NEW-U07-D02
% mdparagraph:P311:begin
In the first speech, identify the condition and the claimed ability. In the second, distinguish the speaker’s willingness from the joint exhortation.
% mdparagraph:P311:end

\end{minipage}\par\addvspace{4pt}

\par\begin{minipage}{\linewidth}\raggedright
% mdparagraph:P312:begin
\textbf{D3.}
% mdparagraph:P312:end

% item: NEW-U07-D03
% mdparagraph:P313:begin
Identify the cases of mē, ðē, and his gefēran in the account. Explain the two dative relations, and identify the possessor intended by his from the stated speaker sequence.
% mdparagraph:P313:end

\end{minipage}\par\addvspace{4pt}

\par\begin{minipage}{\linewidth}\raggedright
% mdparagraph:P314:begin
\textbf{D4.}
% mdparagraph:P314:end

% item: NEW-U07-D04
% mdparagraph:P315:begin
Identify the final passive. What earlier event establishes its result, and which participants performed that event? Explain how the narrative supplies this information even though the passive clause has no agent phrase.
% mdparagraph:P315:end

\end{minipage}\par\addvspace{4pt}

\reserve{7\baselineskip}
% mdparagraph:P316:begin
\section*{10. The fifth day of Creation}\BookHeading{nav-ch07-learner-0040}{subsection}{10. The fifth day of Creation}{1}
% mdparagraph:P316:end

% mdparagraph:P317:begin
This complete paragraph is from Cook’s \emph{First Book}, Reader I, printed p.126 /\allowbreak{} one-based PDF p.142, editorial prose lines 5–14. The printed page has marginal line anchors; the consecutive prose-line count is editorial. The paragraph begins with God’s direction to the waters, describes the creatures and blessing, and closes the fifth day.
% mdparagraph:P317:end

% source: FB-I-CREATION-FIFTH-DAY-01
% mdparagraph:P318:begin
\par\begin{minipage}{\linewidth}\raggedright
\readingtext{God cwæð ēac swilce, “Tēon nū ðā wæteru forð swimmendu cynn cucu on līfe, and flēogendu cynn ofer eorðan under ðǣre heofonan fæstnesse.” And God gescēop ðā ðā miclan hwalas, and eall libbendu fisccynn and styriendlicu, ðe ðā wæteru tugon forð on hiera hīwum, and eall flēogendu cynn æfter hiera cynne; God geseah ðā ðæt hit gōd wæs. And blētsode hīe, ðus cweðende, “Weaxað, and bēoð gemanigfielde, and gefyllað ðǣre sǣ wæteru, and ðā fuglas bēon gemanigfielde ofer eorðan.” And ðā wæs geworden ǣfen and męrgen se fīfta dæg.}
\end{minipage}\par
% mdparagraph:P318:end

\reserve{7\baselineskip}
% mdparagraph:P319:begin
\subsection*{Forms for this reading}\BookHeading{nav-ch07-learner-0041}{subsubsection}{Forms for this reading}{0}
% mdparagraph:P319:end

% teaching: NEW-U07-SOURCE-FORMS
% mdparagraph:P320:begin
Recognize the following three \textbf{exact present-participle forms}. Each is strong neuter accusative plural with its noun. Their -u forms are retained from the source. The assignment is recognition of these forms; no independent recall of other participle cells of these three verbs is required.
% mdparagraph:P320:end

% table: U07-DISPLAY-SOURCE-SUPPORT
% mdtable:T45:begin
\par\begingroup
\fontsize{10.5}{13.2}\selectfont
\setlength{\LTleft}{0pt}\setlength{\LTright}{\fill}
\noindent
\begin{tabular}{@{}>{\raggedright\arraybackslash}p{21.01mm}>{\raggedright\arraybackslash}p{19.01mm}>{\raggedright\arraybackslash}p{19.01mm}>{\raggedright\arraybackslash}p{56.54mm}@{}}
\toprule
\textbf{Source form} & \textbf{Verb} & \textbf{Meaning} & \textbf{Agreement in this passage} \\
\midrule
swimmendu & swimman & swimming & neuter accusative plural with cynn \\
flēogendu & flēogan & flying & neuter accusative plural with cynn, at both occurrences \\
libbendu & libban & living & neuter accusative plural with fisccynn \\
\bottomrule
\end{tabular}
\par\endgroup

% mdtable:T45:end
% mdparagraph:P321:begin
The following exact forms are \textbf{supplied help}, available throughout the reading tasks. They do not add independent vocabulary or construction practice. Their phrases are explained in the ten groups below.
% mdparagraph:P321:end

% mdtable:T46:begin
\par\begingroup
\fontsize{10.5}{13.2}\selectfont
\setlength{\LTleft}{0pt}\setlength{\LTright}{\fill}
\noindent
\begin{tabular}{@{}>{\raggedright\arraybackslash}p{27.91mm}>{\raggedright\arraybackslash}p{93.28mm}@{}}
\toprule
\textbf{Source form} & \textbf{Local analysis and meaning} \\
\midrule
forð & adverb “forth”; part of the supplied “bring forth” expression \\
cucu & irregular adjective, neuter accusative plural, “alive” \\
heofonan & genitive singular of weak feminine heofone, “heaven” \\
ðā & temporal adverb “then” where identified below; distinct from the article \\
fisccynn & neuter accusative plural, “fish-kinds” \\
styriendlicu & strong neuter accusative plural, “moving”; this source’s -u variant \\
ðe & relative pronoun, object “which” in the supplied clause \\
hīwum & dative plural of neuter hīw, “forms; kinds” \\
æfter & “according to”, with dative cynne in the supplied phrase \\
ðæt & conjunction “that” introducing the supplied content clause \\
blētsode & third-singular preterite indicative of plain blētsian, “blessed” \\
sǣ & feminine genitive singular, “of the sea”, in ðǣre sǣ wæteru \\
geworden & past participle in the supplied “come to be” closing formula \\
męrgen & nominative singular source spelling of “morning” \\
\bottomrule
\end{tabular}
\par\endgroup

% mdtable:T46:end
% mdparagraph:P322:begin
Gemanigfielde is the learned plural predicate past participle, used twice in the blessing. Its source analysis is nominative plural. Cweðende has the learned present-participle form; the accompanying-action construction in which it occurs still receives local help.
% mdparagraph:P322:end

\reserve{7\baselineskip}
% mdparagraph:P323:begin
\subsection*{Local help}\BookHeading{nav-ch07-learner-0042}{subsubsection}{Local help}{0}
% mdparagraph:P323:end

% mdparagraph:P324:begin
Use all ten groups during the reading exercises. They supply the unfamiliar lexical meanings and larger constructions; the learned finite verbs, participles, and agreement can still be identified inside them.
% mdparagraph:P324:end

\reserve{5\baselineskip}
% mdparagraph:P325:begin
\textbf{1. Tēon … forð /\allowbreak{} tugon forð.}
% mdparagraph:P325:end

% support: U07-SRC-HELP01
% teaching: NEW-U07-VOC-teon-bring-forth-source
% teaching: NEW-U07-VOC-forth-forth-source
% teaching: NEW-U07-GRAM-U07-SOURCE-BRING-FORTH
% mdparagraph:P326:begin
\textbf{Bring forth /\allowbreak{} brought forth.} Tēon is present third-plural subjunctive with ðā wæteru as its subject. Tugon is past third-plural indicative. The split “bring forth” expression and this sense are supplied: the earlier “draw, pull” does not by itself give this meaning.
% mdparagraph:P326:end

\reserve{5\baselineskip}
% mdparagraph:P327:begin
\textbf{2. swimmendu cynn cucu on līfe; flēogendu cynn.}
% mdparagraph:P327:end

% support: U07-SRC-HELP02
% teaching: NEW-U07-VOC-cucu-alive-source
% teaching: NEW-U07-VOC-on-alive-source
% teaching: NEW-U07-GRAM-U07-SOURCE-ALIVE
% mdparagraph:P328:begin
\textbf{Swimming kinds, alive/\allowbreak{}living; flying kinds.} Cynn is neuter accusative plural, the things to be brought forth. The exact participles swimmendu and flēogendu agree with it. Cucu is the supplied irregular neuter accusative plural adjective “alive”; keep cucu on līfe together as the supplied alive expression. Do not infer that every -ende participle must mechanically lose e before -u.
% mdparagraph:P328:end

\reserve{5\baselineskip}
% mdparagraph:P329:begin
\textbf{3. ofer eorðan; under ðǣre heofonan fæstnesse.}
% mdparagraph:P329:end

% support: U07-SRC-HELP03
% teaching: NEW-U07-VOC-heofone-1
% mdparagraph:P330:begin
\textbf{Over the earth; under the firmament of heaven.} Heofonan is the genitive singular of the source’s weak feminine heofone. Fæstnesse is dative after under. The spelling ðǣre can be feminine genitive or dative singular: it may agree with genitive heofonan or with dative fæstnesse. The wording does not require a unique choice between those article attachments.
% mdparagraph:P330:end

% mdparagraph:P331:begin
Eorðan after spatial ofer likewise has the same spelling in accusative and dative. The sense is clear, but the letters alone do not prove a unique case at either occurrence.
% mdparagraph:P331:end

\reserve{5\baselineskip}
% mdparagraph:P332:begin
\textbf{4. ðā ðā miclan hwalas; eall libbendu fisccynn and styriendlicu.}
% mdparagraph:P332:end

% support: U07-SRC-HELP04
% teaching: NEW-U07-VOC-tha-adverb-1
% teaching: NEW-U07-VOC-fisccynn-fish-kinds-source
% mdparagraph:P333:begin
\textbf{Then the great whales; all living and moving fish-kinds.} The first ðā after gescēop means “then”; the second is the article with miclan hwalas. Fisccynn is neuter accusative plural; libbendu and styriendlicu describe it. The source’s styriendlicu is supplied as an actual -u variant, distinct from the selected earlier styriendlic cell. The and joins the descriptions; it does not require a second newly invented noun meaning.
% mdparagraph:P333:end

\reserve{5\baselineskip}
% mdparagraph:P334:begin
\textbf{5. ðe ðā wæteru tugon forð on hiera hīwum.}
% mdparagraph:P334:end

% support: U07-SRC-HELP05
% teaching: NEW-U07-VOC-the-relative-which
% teaching: NEW-U07-VOC-hiw-kind-source
% teaching: NEW-U07-VOC-on-according-form-source
% teaching: NEW-U07-GRAM-U07-SOURCE-RELATIVE
% teaching: NEW-U07-GRAM-U07-SOURCE-CLASSIFICATION
% mdparagraph:P335:begin
\textbf{Which the waters brought forth according to their forms/\allowbreak{}kinds.} Supply the whole relative clause. The relative ðe is its object; ðā wæteru is its subject. Hīwum is dative plural of neuter hīw, “form; kind”. The phrase on hiera hīwum describes the classification of the creatures brought forth. Its local meaning is supplied rather than derived from the ordinary location sense of on.
% mdparagraph:P335:end

\reserve{5\baselineskip}
% mdparagraph:P336:begin
\textbf{6. æfter hiera cynne.}
% mdparagraph:P336:end

% support: U07-SRC-HELP06
% teaching: NEW-U07-VOC-aefter-according-to-source
% mdparagraph:P337:begin
\textbf{According to their kind.} Æfter governs dative cynne here. Supply this classificatory phrase with the flying kinds. It does not establish every use of æfter for independent composition.
% mdparagraph:P337:end

\reserve{5\baselineskip}
% mdparagraph:P338:begin
\textbf{7. God geseah ðā ðæt hit gōd wæs.}
% mdparagraph:P338:end

% support: U07-SRC-HELP07
% teaching: NEW-U07-VOC-thaet-conjunction-1
% teaching: NEW-U07-GRAM-U07-SOURCE-CONTENT
% mdparagraph:P339:begin
\textbf{God then saw that it was good.} The conjunction ðæt introduces the supplied content clause; its verb wæs follows the predicate gōd. Hit refers to the nonhuman situation described. Supply this whole clause relation and word order; recognizing its familiar indicative forms does not require general content-clause production.
% mdparagraph:P339:end

\reserve{5\baselineskip}
% mdparagraph:P340:begin
\textbf{8. And blētsode hīe, ðus cweðende.}
% mdparagraph:P340:end

% support: U07-SRC-HELP08
% teaching: NEW-U07-VOC-bletsian-bless-source
% teaching: NEW-U07-SRC-PART-FORMS
% teaching: NEW-U07-GRAM-U07-SOURCE-ACCOMPANYING-PARTICIPLE
% teaching: NEW-U07-GRAM-U07-SOURCE-SUBJECT-ELLIPSIS
% mdparagraph:P341:begin
\textbf{And [God] blessed them, saying thus.} God continues as the subject without being repeated; hīe is the accusative plural object. Blētsode is supplied as the third-singular preterite of plain blētsian. It is not the prefixed geblētsian of the productive tables. Cweðende refers to God and accompanies the blessing, introducing his words. This attendant participle construction is supplied even though the participle’s form is learned.
% mdparagraph:P341:end

\reserve{5\baselineskip}
% mdparagraph:P342:begin
\textbf{9. Weaxað … bēoð … gefyllað; ðā fuglas bēon gemanigfielde.}
% mdparagraph:P342:end

% support: U07-SRC-HELP09
% teaching: NEW-U07-VOC-gefyllan-fill-source
% mdparagraph:P343:begin
\textbf{Increase, be multiplied, and fill; let the birds be multiplied.} The created beings are first addressed with plural imperatives. The birds are then named as the third-person plural subject of a subjunctive wish. Gemanigfielde is the source’s nominative plural past participle in both clauses. Gefyllað means “fill” here; this lexical sense is supplied in addition to the earlier “finish, accomplish”.
% mdparagraph:P343:end

% mdparagraph:P344:begin
\textbf{Ðǣre sǣ wæteru} means “the waters of the sea”. Wæteru is the head noun and the thing to be filled. Sǣ is feminine genitive singular in this source phrase, and ðǣre agrees with it. This is not a prepositional location, nor is it the masculine genitive sǣs of the earlier selected paradigm. Keep the original gender and form in this supplied phrase.
% mdparagraph:P344:end

\reserve{5\baselineskip}
% mdparagraph:P345:begin
\textbf{10. And ðā wæs geworden ǣfen and męrgen se fīfta dæg.}
% mdparagraph:P345:end

% support: U07-SRC-HELP10
% teaching: NEW-U07-VOC-geweorthan-1
% teaching: NEW-U07-GRAM-U07-SOURCE-RESULTATIVE
% mdparagraph:P346:begin
\textbf{And then evening and morning had come to be, the fifth day.} Supply this resultative closing formula with intransitive geweorðan. It is not the passive of a transitive action: no accusative object has become the subject. Preserve the source spelling męrgen. The formula does not require an independently constructed rule for its agreement or a unique general clause analysis.
% mdparagraph:P346:end

% mdparagraph:P347:begin
\textbf{Source presentation.} The original words, spelling, quantity marks, endings, and punctuation are retained. Printed line wrapping is joined; the line-end divisions swim-/\allowbreak{}mendu and styriend-/\allowbreak{}licu are joined as words. The original note cues are moved to the list below. Curved Unicode quotation marks retain the source’s quotation boundaries; męrgen is not normalized to morgen.
% mdparagraph:P347:end

\reserve{7\baselineskip}
% mdparagraph:P348:begin
\subsection*{Cook’s original notes}\BookHeading{nav-ch07-learner-0043}{subsubsection}{Cook’s original notes}{0}
% mdparagraph:P348:end

% mdparagraph:P349:begin
The note numbers and section references below belong to Cook’s edition. The local help supplies their necessary interpretation without requiring the learner to consult those sections.
% mdparagraph:P349:end

% source-note: FB126-n4
% mdparagraph:P350:begin
\textbf{P.126, note 4 (ēac swilce).} Ēac swilce = etiam.
% mdparagraph:P350:end

% source-note: FB126-n5
% mdparagraph:P351:begin
\textbf{P.126, note 5 (forð).} Producant = tēon . . . forð.
% mdparagraph:P351:end

% source-note: FB126-n6
% mdparagraph:P352:begin
\textbf{P.126, note 6 (cucu).} See cucu.
% mdparagraph:P352:end

% source-note: FB126-n7
% mdparagraph:P353:begin
\textbf{P.126, note 7 (līfe).} See līf.
% mdparagraph:P353:end

% source-note: FB126-n8
% mdparagraph:P354:begin
\textbf{P.126, note 8 (flēogendu).} See flēogan, and 61.
% mdparagraph:P354:end

% source-note: FB126-n9
% mdparagraph:P355:begin
\textbf{P.126, note 9 (ðā (temporal, after gescēop)).} Adverb; see 84. 1.
% mdparagraph:P355:end

% source-note: FB126-n10
% mdparagraph:P356:begin
\textbf{P.126, note 10 (hwalas).} See hwæl.
% mdparagraph:P356:end

% source-note: FB126-n11
% mdparagraph:P357:begin
\textbf{P.126, note 11 (styriendlicu).} Lat. motabilem.
% mdparagraph:P357:end

% source-note: FB126-n12
% mdparagraph:P358:begin
\textbf{P.126, note 12 (ðe).} Acc.
% mdparagraph:P358:end

% source-note: FB126-n13
% mdparagraph:P359:begin
\textbf{P.126, note 13 (ðā (subject of tugon)).} Nom. plur.
% mdparagraph:P359:end

% source-note: FB126-n14
% mdparagraph:P360:begin
\textbf{P.126, note 14 (tugon).} See tēon.
% mdparagraph:P360:end

% source-note: FB126-n15
% mdparagraph:P361:begin
\textbf{P.126, note 15 (forð (after tugon)).} Tugon forð = producerunt.
% mdparagraph:P361:end

% source-note: FB126-n16
% mdparagraph:P362:begin
\textbf{P.126, note 16 (blētsode).} See blētsian, and 33.
% mdparagraph:P362:end

% source-note: FB126-n17
% mdparagraph:P363:begin
\textbf{P.126, note 17 (cweðende).} See cweðan.
% mdparagraph:P363:end

% source-note: FB126-n18
% mdparagraph:P364:begin
\textbf{P.126, note 18 (Weaxað).} See weaxan, and 24.
% mdparagraph:P364:end

% source-note: FB126-n19
% mdparagraph:P365:begin
\textbf{P.126, note 19 (gemanigfielde (first occurrence)).} Past part. in nom. plur.
% mdparagraph:P365:end

% source-note: FB126-n20
% mdparagraph:P366:begin
\textbf{P.126, note 20 (gefyllað).} See gefyllan.
% mdparagraph:P366:end

% source-note: FB126-n21
% mdparagraph:P367:begin
\textbf{P.126, note 21 (bēon).} See 193. a.
% mdparagraph:P367:end

\reserve{7\baselineskip}
% mdparagraph:P368:begin
\subsection*{Practice R. The original paragraph}\BookHeading{nav-ch07-learner-0044}{subsubsection}{Practice R. The original paragraph}{0}
% mdparagraph:P368:end

% mdparagraph:P369:begin
Keep the printed local help available. Recognize the exact three participles and analyse taught forms within their supported clauses. Do not turn a case-syncretic spelling into a unique diagnosis.
% mdparagraph:P369:end

\reserve{5\baselineskip}
\par\begin{minipage}{\linewidth}\raggedright
% mdparagraph:P370:begin
\textbf{R1.}
% mdparagraph:P370:end

% item: NEW-U07-R01
% mdparagraph:P371:begin
Translate the complete fifth-day Creation paragraph with its printed local support.
% mdparagraph:P371:end

\end{minipage}\par\addvspace{4pt}

\par\begin{minipage}{\linewidth}\raggedright
% mdparagraph:P372:begin
\textbf{R2.}
% mdparagraph:P372:end

% item: NEW-U07-R02
% mdparagraph:P373:begin
Parse Tēon, Weaxað, bēoð, and bēon in their actual source clauses. Identify the subjects/\allowbreak{}addressees and explain the commands or wishes.
% mdparagraph:P373:end

\end{minipage}\par\addvspace{4pt}

\par\begin{minipage}{\linewidth}\raggedright
% mdparagraph:P374:begin
\textbf{R3.}
% mdparagraph:P374:end

% item: NEW-U07-R03
% mdparagraph:P375:begin
Find swimmendu, flēogendu, and libbendu in the paragraph. Give the noun each modifies and the case, number, and gender required at each occurrence.
% mdparagraph:P375:end

\end{minipage}\par\addvspace{4pt}

\par\begin{minipage}{\linewidth}\raggedright
% mdparagraph:P376:begin
\textbf{R4.}
% mdparagraph:P376:end

% item: NEW-U07-R04
% mdparagraph:P377:begin
Explain gemanigfielde in the blessing. Identify its two subjects and explain why the source’s plural ending is compatible with the verbal passive.
% mdparagraph:P377:end

\end{minipage}\par\addvspace{4pt}

\par\begin{minipage}{\linewidth}\raggedright
% mdparagraph:P378:begin
\textbf{R5.}
% mdparagraph:P378:end

% item: NEW-U07-R05
% mdparagraph:P379:begin
In ðǣre sǣ wæteru, identify the head noun and the dependent genitive. Translate the phrase and distinguish it from on ðǣm sǣ. Consult the local note on the source’s feminine sea-form.
% mdparagraph:P379:end

\end{minipage}\par\addvspace{4pt}

\par\begin{minipage}{\linewidth}\raggedright
% mdparagraph:P380:begin
\textbf{R6.}
% mdparagraph:P380:end

% item: NEW-U07-R06
% mdparagraph:P381:begin
Identify the limits of the relative clause with ðe using the supplied support. State what the waters brought forth, what God blessed, and which day ends the paragraph.
% mdparagraph:P381:end

\end{minipage}\par\addvspace{4pt}

\reserve{7\baselineskip}
% mdparagraph:P382:begin
\section*{11. Cumulative practice}\BookHeading{nav-ch07-learner-0045}{subsection}{11. Cumulative practice}{1}
% mdparagraph:P382:end

\reserve{7\baselineskip}
% mdparagraph:P383:begin
\subsection*{Practice X. Combining the constructions}\BookHeading{nav-ch07-learner-0046}{subsubsection}{Practice X. Combining the constructions}{0}
% mdparagraph:P383:end

% mdparagraph:P384:begin
Preserve the earlier dative government, recipient constructions, purpose infinitive, and pronoun reference while using this unit’s mood, passive, and naming patterns.
% mdparagraph:P384:end

\reserve{5\baselineskip}
\par\begin{minipage}{\linewidth}\raggedright
% mdparagraph:P385:begin
\textbf{X1.}
% mdparagraph:P385:end

% item: NEW-U07-X01
% mdparagraph:P386:begin
A learner uses an accusative object after gīeman “rule over.” Explain the case required by this unit’s frame, then translate “The man rules over the birds.”
% mdparagraph:P386:end

\end{minipage}\par\addvspace{4pt}

\par\begin{minipage}{\linewidth}\raggedright
% mdparagraph:P387:begin
\textbf{X2.}
% mdparagraph:P387:end

% item: NEW-U07-X02
% mdparagraph:P388:begin
Translate “each tree” as a nominative subject and “all trees” as an accusative object. Explain the different number and agreement choices. Do not invent a plural meaning of ǣlc.
% mdparagraph:P388:end

\end{minipage}\par\addvspace{4pt}

\par\begin{minipage}{\linewidth}\raggedright
% mdparagraph:P389:begin
\textbf{X3.}
% mdparagraph:P389:end

% item: NEW-U07-X03
% mdparagraph:P390:begin
Translate “He can help the man, and he wants to give the birds food.” Identify the case governed by helpan and the two arguments of giefan.
% mdparagraph:P390:end

\end{minipage}\par\addvspace{4pt}

\par\begin{minipage}{\linewidth}\raggedright
% mdparagraph:P391:begin
\textbf{X4.}
% mdparagraph:P391:end

% item: NEW-U07-X04
% mdparagraph:P392:begin
Translate “I call the firmament Heaven,” then change only “I” to “he” and keep the naming relation. Identify both accusatives.
% mdparagraph:P392:end

\end{minipage}\par\addvspace{4pt}

\par\begin{minipage}{\linewidth}\raggedright
% mdparagraph:P393:begin
\textbf{X5.}
% mdparagraph:P393:end

% item: NEW-U07-X05
% mdparagraph:P394:begin
Give a second-plural command and a third-plural wish for “be multiplied.” Explain which word changes and what it tells the reader.
% mdparagraph:P394:end

\end{minipage}\par\addvspace{4pt}

\par\begin{minipage}{\linewidth}\raggedright
% mdparagraph:P395:begin
\textbf{X6.}
% mdparagraph:P395:end

% item: NEW-U07-X06
% mdparagraph:P396:begin
Translate: The companions came to help. The work was accomplished. Use the selected inflected purpose infinitive in the first sentence and a past passive with wesan in the second.
% mdparagraph:P396:end

\end{minipage}\par\addvspace{4pt}

\clearpage
% mdparagraph:P397:begin
\section*{12. Sources and selected forms}\BookHeading{nav-ch07-learner-0047}{subsection}{12. Sources and selected forms}{1}
% mdparagraph:P397:end

\begingroup\fontsize{10}{12.8}\selectfont
\setlength{\parskip}{1.5pt}
% mdparagraph:P398:begin
The principal source is Albert S. Cook, \textbf{A First Book in Old English: Grammar, Reader, Notes, and Vocabulary}, second edition, revised and enlarged (Ginn, 1895). Cook’s \textbf{Exercises in Old English} (Ginn, 1895) supplies I.4.18 and the complete I.6–7 prompts. They retain their original identifiers and wording despite being grouped in this grammatical unit.
% mdparagraph:P398:end

% mdparagraph:P399:begin
The full tables are edited teaching paradigms derived from the checked systems and selected stems. A completed cell is not a claim that the entire table occurs as a quotation in running prose. Joseph Wright and Elizabeth Mary Wright, \textbf{Old English Grammar}, second edition (Oxford, 1914), and R. D. Fulk, \textbf{An Introductory Grammar of Old English with an Anthology of Readings}, author-corrected open-access reprint (2020), corroborate the stated forms, variants, and bounded uses.
% mdparagraph:P399:end

% mdparagraph:P400:begin
All PDF pages are one-based. The selected Cook \emph{First Book} scan is printed page +16; Wright is +20; Fulk is +12. The \emph{Exercises} scan is printed page +8. The original Creation excerpt is Cook Reader I, p.126 /\allowbreak{} PDF p.142, editorial prose lines 5–14, with original notes 4–21. It is not an adaptation of the new account.
% mdparagraph:P400:end

% mdparagraph:P401:begin
\textbf{Subjunctive formation.} Cook §§93–95, pp.54–55 /\allowbreak{} PDF pp.70–71, supplies the basic stem and ending relations. The selected strong systems are in §§102–110, pp.59–67 /\allowbreak{} PDF pp.75–83; weak systems in §§113–124, pp.68–76 /\allowbreak{} PDF pp.84–92; preterite-presents and anomalous systems in §§126–141, pp.76–80 /\allowbreak{} PDF pp.92–96. Fulk §§45 and 88, pp.29 and 59 /\allowbreak{} PDF pp.41 and 71, corroborates the ending system; §123, pp.89–90 /\allowbreak{} PDF pp.101–102, distinguishes defective citation labels.
% mdparagraph:P401:end

% mdparagraph:P402:begin
\textbf{Mood and clauses.} Cook §§192–198, pp.103–105 /\allowbreak{} PDF pp.119–121, and Fulk §§44–45, pp.28–29 /\allowbreak{} PDF pp.40–41, support the elementary uses and the warning that gif does not automatically select subjunctive mood. The nonpast preterite example is a new composition with its unreal setting supplied; wider conditional syntax is deferred.
% mdparagraph:P402:end

% mdparagraph:P403:begin
\textbf{Participles and passive.} Cook §§55–62, pp.38–41 /\allowbreak{} PDF pp.54–57, gives the adjective and participial patterns; §179, pp.99–100 /\allowbreak{} PDF pp.115–116, and §§187–188, pp.101–102 /\allowbreak{} PDF pp.117–118, support the predicate and passive distinctions. Wright §§429–442, pp.211–218 /\allowbreak{} PDF pp.231–238, and Fulk §§54–56, pp.36–37 /\allowbreak{} PDF pp.48–49, corroborate agreement and syncope. The fuller stem spellings in the tables make the endings recoverable; the listed syncopated and plural variants remain valid.
% mdparagraph:P403:end

% mdparagraph:P404:begin
\textbf{Lexical frames.} Cook §156f, p.91 /\allowbreak{} PDF p.107, gives the genitive frame of gīeman; §173, p.97 /\allowbreak{} PDF p.113, gives the two-accusative naming relation. The new noun follows §47, p.32 /\allowbreak{} PDF p.48; ǣlc is identified in §89a, p.52 /\allowbreak{} PDF p.68, with its vocabulary entry on p.274 /\allowbreak{} PDF p.290. Spatial ofer with accusative and dative is explicit in Fulk §36, p.23 /\allowbreak{} PDF p.35; the glossary’s case mark is explained on p.251 /\allowbreak{} PDF p.263 and applied to ofer on p.297 /\allowbreak{} PDF p.309.
% mdparagraph:P404:end

% mdparagraph:P405:begin
The checked Cook vocabulary entries include gīeman (p.296 /\allowbreak{} PDF p.312), rīcsian (p.309 /\allowbreak{} PDF p.325), ge-manig-fieldan (p.293 /\allowbreak{} PDF p.309), tō-twǣman (p.316 /\allowbreak{} PDF p.332), wyrcean (p.324 /\allowbreak{} PDF p.340), and ge-weald (p.295 /\allowbreak{} PDF p.311). Dictionary compound hyphens are omitted in the learner headwords gemanigfieldan, tōtwǣman, and geweald; the source’s running spelling fisccynn is retained. No vowel length is added to gemanigfieldan.
% mdparagraph:P405:end

% mdparagraph:P406:begin
The additional lexical help follows Cook’s entries æfter (p.274 /\allowbreak{} PDF p.290), blētsian (p.279 /\allowbreak{} PDF p.295), cucu and cynn (pp.281–282 /\allowbreak{} PDF pp.297–298), ēac (p.283 /\allowbreak{} PDF p.299), fisc-cynn (p.287 /\allowbreak{} PDF p.303), forð (p.288 /\allowbreak{} PDF p.304), hīw (p.299 /\allowbreak{} PDF p.315), nū and ofer (p.307 /\allowbreak{} PDF p.323), swā (p.314 /\allowbreak{} PDF p.330), swilce and tēon (p.315 /\allowbreak{} PDF p.331), and ðus (p.319 /\allowbreak{} PDF p.335). These references support the particular meanings taught or supplied here, not every dictionary sense.
% mdparagraph:P406:end

% mdparagraph:P407:begin
The chapter’s eleven examples and “Binding the animals” are new compositions. The source’s relative and content clauses, accompanying participle, classification expressions, plain blētsode, and closing formula remain local help. Earlier instrumental and reference-only boundaries remain in force; complete morphology does not by itself teach every construction in which a form can occur.
% mdparagraph:P407:end

\endgroup

% book-source-body:end
\par\label{end-ch07-learner}
\clearpage\endgroup

\BookRecto
\begingroup
\renewcommand{\runningtitle}{Unit 8. Questions and relatives}
\hyphenpenalty=10000
\exhyphenpenalty=10000
\emergencystretch=3em
\newcommand{\itempara}[2]{\par\addvspace{3pt}\noindent\hangindent=1.8em\hangafter=1\makebox[1.8em][l]{#1.}#2\par}
\newcommand{\reserve}[1]{\par\penalty-100\begingroup\dimen0=\pagegoal\advance\dimen0 by-\pagetotal\ifdim\dimen0<#1\newpage\fi\endgroup}
% book-source-body:begin

\addfontfeatures{Ligatures=TeXOff}
\raggedright
\booktitle{\BookComponent{ch08-learner}{Unit 8. Questions and relative clauses}UNIT 8}{Questions and relative clauses}

% mdparagraph:P001:begin
This unit gives direct questions, the remaining selected pronoun and pronominal-adjective paradigms, and ordinary relative clauses. The additional verb systems and lexical constructions prepare the complete Cook I.8 exercise and two connected readings.
% mdparagraph:P001:end

% mdparagraph:P002:begin
Tables distinguish singular and plural, masculine, neuter and feminine, and nominative, genitive, dative and accusative. In verb tables, \textbf{1}, \textbf{2} and \textbf{3} mean first, second and third person. A row labelled \textbf{1–3 sg.} or \textbf{1–3 pl.} includes every person in that number. Dual subjects take plural finite forms. Forms labelled \textbf{reference} are for recognition; their display does not require independent use of a new construction. Preserve vowel length and the distinction between e and ę.
% mdparagraph:P002:end

% mdparagraph:P003:begin
The examples and “The ship and the fish” are new teaching material. The merchant dialogue is an exact source excerpt with separate local help and original notes.
% mdparagraph:P003:end

\reserve{7\baselineskip}
% mdparagraph:P004:begin
\section*{1. Direct questions}\BookHeading{nav-ch08-learner-0001}{subsection}{1. Direct questions}{1}\BookLabel{bk-ch08-learner-da4d22fbf997}
% mdparagraph:P004:end

% teaching: NEW-U08-GRAM-u08-direct-questions
% teaching: NEW-U08-GRAM-u08-clause-order
% mdparagraph:P005:begin
A direct \textbf{yes/\allowbreak{}no question} commonly begins with its finite verb, followed by the subject. The verb retains the person, number and tense required by that subject. Old English does not need an added verb corresponding to the English question auxiliary “do”.
% mdparagraph:P005:end

\reserve{8\baselineskip}
% example: NEW-U08-EXAMPLE-01
% mdparagraph:P006:begin
Hæfð se mann scip?
% mdparagraph:P006:end

% mdparagraph:P007:begin
Does the man have a ship?
% mdparagraph:P007:end

% mdparagraph:P008:begin
A question word or phrase commonly stands early, often before the finite verb. A question word that is itself the subject needs no second subject. These are common orders, not an exceptionless rule requiring one fixed arrangement. A phrase such as “whose ship” is one constituent, even though it contains two words.
% mdparagraph:P008:end

% mdparagraph:P009:begin
More generally, the finite verb commonly stands early in an independent clause and later, often finally, in a dependent clause. The relative clauses below illustrate that tendency. The case endings and the meaning of the clause remain essential when the order varies. This unit uses direct questions; indirect questions and general content complements come later.
% mdparagraph:P009:end

\reserve{7\baselineskip}
% mdparagraph:P010:begin
\subsection*{Who? What?}\BookHeading{nav-ch08-learner-0002}{subsubsection}{Who? What?}{0}\BookLabel{bk-ch08-learner-f4ef5d866f80}
% mdparagraph:P010:end

% teaching: NEW-U08-VOC-hwa
% teaching: NEW-U08-GRAM-u08-hwa-paradigm
% teaching: NEW-U08-GRAM-u08-interrogative-case
% mdparagraph:P011:begin
\textbf{Hwā} asks about a person; \textbf{hwæt} asks about a thing. Their paradigm is singular only. Masculine and feminine share one personal series; there is no separate feminine or plural paradigm. Choose the case from the interrogative’s own function: nominative subject, genitive possessor, accusative direct object, or dative recipient or other governed complement.
% mdparagraph:P011:end

% teaching: NEW-U08-PAR-HWA
\reserve{7\baselineskip}
% mdparagraph:P012:begin
\subsection*{hwā /\allowbreak{} hwæt}\BookHeading{nav-ch08-learner-0003}{subsubsection}{hwā / hwæt}{0}
% mdparagraph:P012:end

% table: NEW-U08-PAR-HWA
% mdtable:T01:begin
\par\begingroup
\fontsize{10.5}{13.2}\selectfont
\setlength{\LTleft}{0pt}\setlength{\LTright}{\fill}
\noindent
\begin{tabular}{@{}>{\raggedright\arraybackslash}p{25.54mm}>{\raggedright\arraybackslash}p{64.36mm}>{\raggedright\arraybackslash}p{28.48mm}@{}}
\toprule
\textbf{Case} & \textbf{Common masculine/\allowbreak{}feminine singular} & \textbf{Neuter singular} \\
\midrule
Nominative & hwā & hwæt \\
Genitive & hwæs & hwæs \\
Dative & hwǣm & hwǣm \\
Accusative & hwone & hwæt \\
\bottomrule
\end{tabular}
\par\endgroup

% mdtable:T01:end
\reserve{8\baselineskip}
% example: NEW-U08-EXAMPLE-02
% mdparagraph:P013:begin
Hwā hæfð ðæt scip? Hwone siehð se mann?
% mdparagraph:P013:end

% mdparagraph:P014:begin
Who has that ship? Whom does the man see?
% mdparagraph:P014:end

% mdparagraph:P015:begin
Hwā is the subject of hæfð. Hwone is the object of siehð; se mann remains the nominative subject. English “whom” corresponds here to an accusative because “see” has an accusative object.
% mdparagraph:P015:end

\reserve{8\baselineskip}
% example: NEW-U08-EXAMPLE-03
% mdparagraph:P016:begin
Hwǣm sęlst ðū ðone fisc? Hwæs scip siehst ðū?
% mdparagraph:P016:end

% mdparagraph:P017:begin
To whom do you give that fish? Whose ship do you see?
% mdparagraph:P017:end

% mdparagraph:P018:begin
Hwǣm is the dative recipient of the fish. Hwæs is genitive because it asks whose ship is seen; scip itself is the accusative object of siehst. The possessor and the possessed noun have different functions within that phrase.
% mdparagraph:P018:end

% mdparagraph:P019:begin
\textbf{Hwǣm} is the selected dative in this unit. Cook also gives \textbf{hwām} as a valid alternative. The distinct instrumental forms below are reference forms; learning them does not introduce general instrumental syntax.
% mdparagraph:P019:end

\reserve{7\baselineskip}
% mdparagraph:P020:begin
\subsection*{hwā reference variants and instrumentals}\BookHeading{nav-ch08-learner-0004}{subsubsection}{hwā reference variants and instrumentals}{0}
% mdparagraph:P020:end

% table: NEW-U08-MATRIX-HWA-REFERENCE
% mdtable:T02:begin
\par\begingroup
\fontsize{10.5}{13.2}\selectfont
\setlength{\LTleft}{0pt}\setlength{\LTright}{\fill}
\noindent
\begin{tabular}{@{}>{\raggedright\arraybackslash}p{87.48mm}>{\raggedright\arraybackslash}p{33.71mm}@{}}
\toprule
\textbf{Function} & \textbf{Form} \\
\midrule
Dative common personal: printed alternative & hwām \\
Dative neuter: printed alternative & hwām \\
Neuter instrumental & hwȳ \\
Neuter instrumental alternative & hwon \\
\bottomrule
\end{tabular}
\par\endgroup

% mdtable:T02:end
\reserve{7\baselineskip}
% mdparagraph:P021:begin
\subsection*{How? Why?}\BookHeading{nav-ch08-learner-0005}{subsubsection}{How? Why?}{0}\BookLabel{bk-ch08-learner-343e3ece044b}
% mdparagraph:P021:end

% teaching: NEW-U08-VOC-hu
% teaching: NEW-U08-VOC-hwy
% teaching: NEW-U08-GRAM-u08-interrogative-adverbs
% mdparagraph:P022:begin
\textbf{Hū} asks about manner; \textbf{hwȳ} asks for a reason. Both are invariant question words in the following use. Lexical hwȳ “why?” can be used without constructing a general instrumental phrase.
% mdparagraph:P022:end

\reserve{7\baselineskip}
% mdparagraph:P023:begin
\subsection*{Question adverbs}\BookHeading{nav-ch08-learner-0006}{subsubsection}{Question adverbs}{0}
% mdparagraph:P023:end

% table: NEW-U08-MATRIX-QUESTION-ADVERBS
% mdtable:T03:begin
\par\begingroup
\fontsize{10.5}{13.2}\selectfont
\setlength{\LTleft}{0pt}\setlength{\LTright}{\fill}
\noindent
\begin{tabular}{@{}>{\raggedright\arraybackslash}p{60.59mm}>{\raggedright\arraybackslash}p{60.59mm}@{}}
\toprule
\textbf{Function} & \textbf{Form} \\
\midrule
How & hū \\
Why & hwȳ \\
\bottomrule
\end{tabular}
\par\endgroup

% mdtable:T03:end
\reserve{8\baselineskip}
% example: NEW-U08-EXAMPLE-04
% mdparagraph:P024:begin
Hwæt hæfð se mann? Hū swimð hē? Hwȳ swimð hē?
% mdparagraph:P024:end

% mdparagraph:P025:begin
What does the man have? How does he swim? Why does he swim?
% mdparagraph:P025:end

% mdparagraph:P026:begin
Hwæt is the accusative thing possessed. Hū and hwȳ ask different questions about the same swimming: how it is done, and why it is done.
% mdparagraph:P026:end

\reserve{7\baselineskip}
% mdparagraph:P027:begin
\section*{2. Interrogative adjectives}\BookHeading{nav-ch08-learner-0007}{subsection}{2. Interrogative adjectives}{1}\BookLabel{bk-ch08-learner-e70ad5862678}
% mdparagraph:P027:end

% teaching: NEW-U08-GRAM-u08-interrogative-adjectives
% mdparagraph:P028:begin
\textbf{Hwilc} means “which?” or “of what kind?”; \textbf{hwæðer}, “which of two?”; \textbf{hūlic}, “of what sort?”. They have strong adjective inflection and agree with their noun or referent in gender, number and case. No weak interrogative-adjective paradigm is needed here.
% mdparagraph:P028:end

% mdparagraph:P029:begin
The strong endings are the familiar adjective endings: masculine and neuter singular genitive -es and dative -um, masculine accusative -ne, feminine genitive and dative -re; plural genitive -ra and dative -um. The tables give the complete forms, including the nominative and accusative distinctions. Feminine plural -a is selected consistently; the previously noted feminine plural alternative -e remains valid.
% mdparagraph:P029:end

% teaching: NEW-U08-VOC-hwilc
% teaching: NEW-U08-PAR-HWILC-STRONG
\reserve{7\baselineskip}
% mdparagraph:P030:begin
\subsection*{hwilc: strong declension, singular}\BookHeading{nav-ch08-learner-0008}{subsubsection}{hwilc: strong declension, singular}{0}
% mdparagraph:P030:end

% table: NEW-U08-PAR-HWILC-STRONG-SG
% mdtable:T04:begin
\par\begingroup
\fontsize{10.5}{13.2}\selectfont
\setlength{\LTleft}{0pt}\setlength{\LTright}{\fill}
\noindent
\begin{tabular}{@{}>{\raggedright\arraybackslash}p{28.89mm}>{\raggedright\arraybackslash}p{28.89mm}>{\raggedright\arraybackslash}p{28.89mm}>{\raggedright\arraybackslash}p{28.89mm}@{}}
\toprule
\textbf{Case} & \textbf{Masculine} & \textbf{Neuter} & \textbf{Feminine} \\
\midrule
Nominative & hwilc & hwilc & hwilc \\
Genitive & hwilces & hwilces & hwilcre \\
Dative & hwilcum & hwilcum & hwilcre \\
Accusative & hwilcne & hwilc & hwilce \\
\bottomrule
\end{tabular}
\par\endgroup

% mdtable:T04:end
\reserve{7\baselineskip}
% mdparagraph:P031:begin
\subsection*{hwilc: strong declension, plural}\BookHeading{nav-ch08-learner-0009}{subsubsection}{hwilc: strong declension, plural}{0}
% mdparagraph:P031:end

% table: NEW-U08-PAR-HWILC-STRONG-PL
% mdtable:T05:begin
\par\begingroup
\fontsize{10.5}{13.2}\selectfont
\setlength{\LTleft}{0pt}\setlength{\LTright}{\fill}
\noindent
\begin{tabular}{@{}>{\raggedright\arraybackslash}p{28.89mm}>{\raggedright\arraybackslash}p{28.89mm}>{\raggedright\arraybackslash}p{28.89mm}>{\raggedright\arraybackslash}p{28.89mm}@{}}
\toprule
\textbf{Case} & \textbf{Masculine} & \textbf{Neuter} & \textbf{Feminine} \\
\midrule
Nominative & hwilce & hwilc & hwilca \\
Genitive & hwilcra & hwilcra & hwilcra \\
Dative & hwilcum & hwilcum & hwilcum \\
Accusative & hwilce & hwilc & hwilca \\
\bottomrule
\end{tabular}
\par\endgroup

% mdtable:T05:end
% mdparagraph:P032:begin
Keep the stem hwilc- before these endings. Hwelc and hwylc are documented spellings; expanded forms such as hwilcere and hwilcera also occur. The tables give the selected forms for recall.
% mdparagraph:P032:end

% teaching: NEW-U08-VOC-hulic
% teaching: NEW-U08-PAR-HULIC-STRONG
\reserve{7\baselineskip}
% mdparagraph:P033:begin
\subsection*{hūlic: strong declension, singular}\BookHeading{nav-ch08-learner-0010}{subsubsection}{hūlic: strong declension, singular}{0}
% mdparagraph:P033:end

% table: NEW-U08-PAR-HULIC-STRONG-SG
% mdtable:T06:begin
\par\begingroup
\fontsize{10.5}{13.2}\selectfont
\setlength{\LTleft}{0pt}\setlength{\LTright}{\fill}
\noindent
\begin{tabular}{@{}>{\raggedright\arraybackslash}p{28.89mm}>{\raggedright\arraybackslash}p{28.89mm}>{\raggedright\arraybackslash}p{28.89mm}>{\raggedright\arraybackslash}p{28.89mm}@{}}
\toprule
\textbf{Case} & \textbf{Masculine} & \textbf{Neuter} & \textbf{Feminine} \\
\midrule
Nominative & hūlic & hūlic & hūlic \\
Genitive & hūlices & hūlices & hūlicre \\
Dative & hūlicum & hūlicum & hūlicre \\
Accusative & hūlicne & hūlic & hūlice \\
\bottomrule
\end{tabular}
\par\endgroup

% mdtable:T06:end
\reserve{7\baselineskip}
% mdparagraph:P034:begin
\subsection*{hūlic: strong declension, plural}\BookHeading{nav-ch08-learner-0011}{subsubsection}{hūlic: strong declension, plural}{0}
% mdparagraph:P034:end

% table: NEW-U08-PAR-HULIC-STRONG-PL
% mdtable:T07:begin
\par\begingroup
\fontsize{10.5}{13.2}\selectfont
\setlength{\LTleft}{0pt}\setlength{\LTright}{\fill}
\noindent
\begin{tabular}{@{}>{\raggedright\arraybackslash}p{28.89mm}>{\raggedright\arraybackslash}p{28.89mm}>{\raggedright\arraybackslash}p{28.89mm}>{\raggedright\arraybackslash}p{28.89mm}@{}}
\toprule
\textbf{Case} & \textbf{Masculine} & \textbf{Neuter} & \textbf{Feminine} \\
\midrule
Nominative & hūlice & hūlic & hūlica \\
Genitive & hūlicra & hūlicra & hūlicra \\
Dative & hūlicum & hūlicum & hūlicum \\
Accusative & hūlice & hūlic & hūlica \\
\bottomrule
\end{tabular}
\par\endgroup

% mdtable:T07:end
% mdparagraph:P035:begin
Hūlic keeps hūlic- throughout. Cook marks the first vowel long; this chapter retains that notation. The separately documented spelling hūlīc does not change the selected table.
% mdparagraph:P035:end

% teaching: NEW-U08-VOC-hwaether
% teaching: NEW-U08-PAR-HWAETHER-STRONG
\reserve{7\baselineskip}
% mdparagraph:P036:begin
\subsection*{hwæðer: strong declension, singular}\BookHeading{nav-ch08-learner-0012}{subsubsection}{hwæðer: strong declension, singular}{0}
% mdparagraph:P036:end

% table: NEW-U08-PAR-HWAETHER-STRONG-SG
% mdtable:T08:begin
\par\begingroup
\fontsize{10.5}{13.2}\selectfont
\setlength{\LTleft}{0pt}\setlength{\LTright}{\fill}
\noindent
\begin{tabular}{@{}>{\raggedright\arraybackslash}p{28.89mm}>{\raggedright\arraybackslash}p{28.89mm}>{\raggedright\arraybackslash}p{28.89mm}>{\raggedright\arraybackslash}p{28.89mm}@{}}
\toprule
\textbf{Case} & \textbf{Masculine} & \textbf{Neuter} & \textbf{Feminine} \\
\midrule
Nominative & hwæðer & hwæðer & hwæðer \\
Genitive & hwæðres & hwæðres & hwæðerre \\
Dative & hwæðrum & hwæðrum & hwæðerre \\
Accusative & hwæðerne & hwæðer & hwæðre \\
\bottomrule
\end{tabular}
\par\endgroup

% mdtable:T08:end
\reserve{7\baselineskip}
% mdparagraph:P037:begin
\subsection*{hwæðer: strong declension, plural}\BookHeading{nav-ch08-learner-0013}{subsubsection}{hwæðer: strong declension, plural}{0}
% mdparagraph:P037:end

% table: NEW-U08-PAR-HWAETHER-STRONG-PL
% mdtable:T09:begin
\par\begingroup
\fontsize{10.5}{13.2}\selectfont
\setlength{\LTleft}{0pt}\setlength{\LTright}{\fill}
\noindent
\begin{tabular}{@{}>{\raggedright\arraybackslash}p{28.89mm}>{\raggedright\arraybackslash}p{28.89mm}>{\raggedright\arraybackslash}p{28.89mm}>{\raggedright\arraybackslash}p{28.89mm}@{}}
\toprule
\textbf{Case} & \textbf{Masculine} & \textbf{Neuter} & \textbf{Feminine} \\
\midrule
Nominative & hwæðre & hwæðer & hwæðra \\
Genitive & hwæðerra & hwæðerra & hwæðerra \\
Dative & hwæðrum & hwæðrum & hwæðrum \\
Accusative & hwæðre & hwæðer & hwæðra \\
\bottomrule
\end{tabular}
\par\endgroup

% mdtable:T09:end
% mdparagraph:P038:begin
Hwæðer loses its unstressed e before the ordinary vowel endings: hwæðres, hwæðrum, hwæðre, hwæðra. It retains e before the consonantal endings in hwæðerne, hwæðerre and hwæðerra. Its feminine nominative singular and neuter nominative/\allowbreak{}accusative plural are hwæðer. Do not copy the ōðru cells of ōðer into this paradigm. Attested alternatives include hwæðeres and feminine hwæðere.
% mdparagraph:P038:end

% mdparagraph:P039:begin
Hwæðer always distinguishes \textbf{two alternatives} in the use taught here. A singular noun can name one of two people or objects; plural morphology can distinguish two explicitly plural groups. It does not turn the choice into one among an unlimited number. The conjunction “whether” is a separate later use.
% mdparagraph:P039:end

% mdparagraph:P040:begin
In the second question below, two men are being compared.
% mdparagraph:P040:end

\reserve{8\baselineskip}
% example: NEW-U08-EXAMPLE-05
% mdparagraph:P041:begin
Hwilc mann hæfð ðæt scip? Hwæðer mann hæfð ðæt scip? Hūlic is ðæt scip?
% mdparagraph:P041:end

% mdparagraph:P042:begin
Which man has that ship? Which of the two men has that ship? What sort of ship is that?
% mdparagraph:P042:end

% mdparagraph:P043:begin
Hwilc and hwæðer modify nominative masculine singular mann. Hūlic is the strong neuter singular predicate referring to scip. The question word’s English meaning does not remove ordinary adjective agreement.
% mdparagraph:P043:end

\reserve{7\baselineskip}
% mdparagraph:P044:begin
\section*{3. Some, any, each, and such}\BookHeading{nav-ch08-learner-0014}{subsection}{3. Some, any, each, and such}{1}\BookLabel{bk-ch08-learner-c37835ca5690}
% mdparagraph:P044:end

% teaching: NEW-U08-GRAM-u08-indefinites
% mdparagraph:P045:begin
The following words also take complete strong adjective inflection. \textbf{Sum} means “a certain” in the singular and “some” in the plural. \textbf{Ǣnig} means “any” in the ordinary direct questions used here. \textbf{Gehwilc} means “each, every”; its ordinary use in this unit is with a singular noun. \textbf{Swilc} means “such”, referring to a kind already understood from the context.
% mdparagraph:P045:end

% teaching: NEW-U08-VOC-sum
% teaching: NEW-U08-PAR-SUM-STRONG
\reserve{7\baselineskip}
% mdparagraph:P046:begin
\subsection*{sum: strong declension, singular}\BookHeading{nav-ch08-learner-0015}{subsubsection}{sum: strong declension, singular}{0}
% mdparagraph:P046:end

% table: NEW-U08-PAR-SUM-STRONG-SG
% mdtable:T10:begin
\par\begingroup
\fontsize{10.5}{13.2}\selectfont
\setlength{\LTleft}{0pt}\setlength{\LTright}{\fill}
\noindent
\begin{tabular}{@{}>{\raggedright\arraybackslash}p{28.89mm}>{\raggedright\arraybackslash}p{28.89mm}>{\raggedright\arraybackslash}p{28.89mm}>{\raggedright\arraybackslash}p{28.89mm}@{}}
\toprule
\textbf{Case} & \textbf{Masculine} & \textbf{Neuter} & \textbf{Feminine} \\
\midrule
Nominative & sum & sum & sumu \\
Genitive & sumes & sumes & sumre \\
Dative & sumum & sumum & sumre \\
Accusative & sumne & sum & sume \\
\bottomrule
\end{tabular}
\par\endgroup

% mdtable:T10:end
\reserve{7\baselineskip}
% mdparagraph:P047:begin
\subsection*{sum: strong declension, plural}\BookHeading{nav-ch08-learner-0016}{subsubsection}{sum: strong declension, plural}{0}
% mdparagraph:P047:end

% table: NEW-U08-PAR-SUM-STRONG-PL
% mdtable:T11:begin
\par\begingroup
\fontsize{10.5}{13.2}\selectfont
\setlength{\LTleft}{0pt}\setlength{\LTright}{\fill}
\noindent
\begin{tabular}{@{}>{\raggedright\arraybackslash}p{28.89mm}>{\raggedright\arraybackslash}p{28.89mm}>{\raggedright\arraybackslash}p{28.89mm}>{\raggedright\arraybackslash}p{28.89mm}@{}}
\toprule
\textbf{Case} & \textbf{Masculine} & \textbf{Neuter} & \textbf{Feminine} \\
\midrule
Nominative & sume & sumu & suma \\
Genitive & sumra & sumra & sumra \\
Dative & sumum & sumum & sumum \\
Accusative & sume & sumu & suma \\
\bottomrule
\end{tabular}
\par\endgroup

% mdtable:T11:end
% mdparagraph:P048:begin
Short-stem sum has -u in feminine nominative singular and neuter nominative/\allowbreak{}accusative plural: sumu. The genitive and dative endings still mark their ordinary functions.
% mdparagraph:P048:end

\reserve{8\baselineskip}
% example: NEW-U08-EXAMPLE-06
% mdparagraph:P049:begin
Sum mann hæfð scip. Sume męnn habbað scipu.
% mdparagraph:P049:end

% mdparagraph:P050:begin
A certain man has a ship. Some men have ships.
% mdparagraph:P050:end

% mdparagraph:P051:begin
Sum mann is singular, with hæfð; sume męnn is plural, with habbað. The object likewise distinguishes scip from plural scipu. Sum supplies “a certain”; the earlier cardinal ān still means “one”.
% mdparagraph:P051:end

% teaching: NEW-U08-VOC-aenig
% teaching: NEW-U08-PAR-AENIG-STRONG
\reserve{7\baselineskip}
% mdparagraph:P052:begin
\subsection*{ǣnig: strong declension, singular}\BookHeading{nav-ch08-learner-0017}{subsubsection}{ǣnig: strong declension, singular}{0}
% mdparagraph:P052:end

% table: NEW-U08-PAR-AENIG-STRONG-SG
% mdtable:T12:begin
\par\begingroup
\fontsize{10.5}{13.2}\selectfont
\setlength{\LTleft}{0pt}\setlength{\LTright}{\fill}
\noindent
\begin{tabular}{@{}>{\raggedright\arraybackslash}p{28.89mm}>{\raggedright\arraybackslash}p{28.89mm}>{\raggedright\arraybackslash}p{28.89mm}>{\raggedright\arraybackslash}p{28.89mm}@{}}
\toprule
\textbf{Case} & \textbf{Masculine} & \textbf{Neuter} & \textbf{Feminine} \\
\midrule
Nominative & ǣnig & ǣnig & ǣnigu \\
Genitive & ǣniges & ǣniges & ǣnigre \\
Dative & ǣnigum & ǣnigum & ǣnigre \\
Accusative & ǣnigne & ǣnig & ǣnige \\
\bottomrule
\end{tabular}
\par\endgroup

% mdtable:T12:end
\reserve{7\baselineskip}
% mdparagraph:P053:begin
\subsection*{ǣnig: strong declension, plural}\BookHeading{nav-ch08-learner-0018}{subsubsection}{ǣnig: strong declension, plural}{0}
% mdparagraph:P053:end

% table: NEW-U08-PAR-AENIG-STRONG-PL
% mdtable:T13:begin
\par\begingroup
\fontsize{10.5}{13.2}\selectfont
\setlength{\LTleft}{0pt}\setlength{\LTright}{\fill}
\noindent
\begin{tabular}{@{}>{\raggedright\arraybackslash}p{28.89mm}>{\raggedright\arraybackslash}p{28.89mm}>{\raggedright\arraybackslash}p{28.89mm}>{\raggedright\arraybackslash}p{28.89mm}@{}}
\toprule
\textbf{Case} & \textbf{Masculine} & \textbf{Neuter} & \textbf{Feminine} \\
\midrule
Nominative & ǣnige & ǣnigu & ǣniga \\
Genitive & ǣnigra & ǣnigra & ǣnigra \\
Dative & ǣnigum & ǣnigum & ǣnigum \\
Accusative & ǣnige & ǣnigu & ǣniga \\
\bottomrule
\end{tabular}
\par\endgroup

% mdtable:T13:end
% mdparagraph:P054:begin
The\BookPoint{bk-ch08-learner-56e2a1cf9d42} selected ǣnig forms retain i: ǣniges, ǣnigum, ǣnigre, ǣnigra. The listed -u cells are ǣnigu; analogical ǣnig without -u is a documented alternative. Syncope in disyllabic adjectives varies, so do not mechanically contract every -ig adjective by analogy with another word. These questions teach a bounded use of “any”; they do not establish its distribution in every negative or conditional construction.
% mdparagraph:P054:end

% teaching: NEW-U08-VOC-gehwilc
% teaching: NEW-U08-PAR-GEHWILC-STRONG
\reserve{7\baselineskip}
% mdparagraph:P055:begin
\subsection*{gehwilc: strong declension, singular}\BookHeading{nav-ch08-learner-0019}{subsubsection}{gehwilc: strong declension, singular}{0}
% mdparagraph:P055:end

% table: NEW-U08-PAR-GEHWILC-STRONG-SG
% mdtable:T14:begin
\par\begingroup
\fontsize{10.5}{13.2}\selectfont
\setlength{\LTleft}{0pt}\setlength{\LTright}{\fill}
\noindent
\begin{tabular}{@{}>{\raggedright\arraybackslash}p{28.89mm}>{\raggedright\arraybackslash}p{28.89mm}>{\raggedright\arraybackslash}p{28.89mm}>{\raggedright\arraybackslash}p{28.89mm}@{}}
\toprule
\textbf{Case} & \textbf{Masculine} & \textbf{Neuter} & \textbf{Feminine} \\
\midrule
Nominative & gehwilc & gehwilc & gehwilc \\
Genitive & gehwilces & gehwilces & gehwilcre \\
Dative & gehwilcum & gehwilcum & gehwilcre \\
Accusative & gehwilcne & gehwilc & gehwilce \\
\bottomrule
\end{tabular}
\par\endgroup

% mdtable:T14:end
\reserve{7\baselineskip}
% mdparagraph:P056:begin
\subsection*{gehwilc: strong declension, plural}\BookHeading{nav-ch08-learner-0020}{subsubsection}{gehwilc: strong declension, plural}{0}
% mdparagraph:P056:end

% table: NEW-U08-PAR-GEHWILC-STRONG-PL
% mdtable:T15:begin
\par\begingroup
\fontsize{10.5}{13.2}\selectfont
\setlength{\LTleft}{0pt}\setlength{\LTright}{\fill}
\noindent
\begin{tabular}{@{}>{\raggedright\arraybackslash}p{28.89mm}>{\raggedright\arraybackslash}p{28.89mm}>{\raggedright\arraybackslash}p{28.89mm}>{\raggedright\arraybackslash}p{28.89mm}@{}}
\toprule
\textbf{Case} & \textbf{Masculine} & \textbf{Neuter} & \textbf{Feminine} \\
\midrule
Nominative & gehwilce & gehwilc & gehwilca \\
Genitive & gehwilcra & gehwilcra & gehwilcra \\
Dative & gehwilcum & gehwilcum & gehwilcum \\
Accusative & gehwilce & gehwilc & gehwilca \\
\bottomrule
\end{tabular}
\par\endgroup

% mdtable:T15:end
% mdparagraph:P057:begin
Gehwilc follows the same strong pattern as hwilc. A singular noun meaning “each man” or “each companion” takes a singular verb. Its complete plural morphology is given for form analysis, but the plural table does not make gehwilc a general substitute for eall “all” in every context. Partitive expressions such as “each of” come later.
% mdparagraph:P057:end

\reserve{8\baselineskip}
% example: NEW-U08-EXAMPLE-07
% mdparagraph:P058:begin
Hæfst ðū ǣnigne fisc? Gehwilc mann hæfð scip.
% mdparagraph:P058:end

% mdparagraph:P059:begin
Do you have any fish? Each man has a ship.
% mdparagraph:P059:end

% mdparagraph:P060:begin
Ǣnigne agrees with singular masculine accusative fisc; the English “any fish” here asks about a fish, not an Old English plural. Gehwilc mann is nominative singular.
% mdparagraph:P060:end

% teaching: NEW-U08-VOC-swilc
% teaching: NEW-U08-PAR-SWILC-STRONG
\reserve{7\baselineskip}
% mdparagraph:P061:begin
\subsection*{swilc: strong declension, singular}\BookHeading{nav-ch08-learner-0021}{subsubsection}{swilc: strong declension, singular}{0}
% mdparagraph:P061:end

% table: NEW-U08-PAR-SWILC-STRONG-SG
% mdtable:T16:begin
\par\begingroup
\fontsize{10.5}{13.2}\selectfont
\setlength{\LTleft}{0pt}\setlength{\LTright}{\fill}
\noindent
\begin{tabular}{@{}>{\raggedright\arraybackslash}p{28.89mm}>{\raggedright\arraybackslash}p{28.89mm}>{\raggedright\arraybackslash}p{28.89mm}>{\raggedright\arraybackslash}p{28.89mm}@{}}
\toprule
\textbf{Case} & \textbf{Masculine} & \textbf{Neuter} & \textbf{Feminine} \\
\midrule
Nominative & swilc & swilc & swilc \\
Genitive & swilces & swilces & swilcre \\
Dative & swilcum & swilcum & swilcre \\
Accusative & swilcne & swilc & swilce \\
\bottomrule
\end{tabular}
\par\endgroup

% mdtable:T16:end
\reserve{7\baselineskip}
% mdparagraph:P062:begin
\subsection*{swilc: strong declension, plural}\BookHeading{nav-ch08-learner-0022}{subsubsection}{swilc: strong declension, plural}{0}
% mdparagraph:P062:end

% table: NEW-U08-PAR-SWILC-STRONG-PL
% mdtable:T17:begin
\par\begingroup
\fontsize{10.5}{13.2}\selectfont
\setlength{\LTleft}{0pt}\setlength{\LTright}{\fill}
\noindent
\begin{tabular}{@{}>{\raggedright\arraybackslash}p{28.89mm}>{\raggedright\arraybackslash}p{28.89mm}>{\raggedright\arraybackslash}p{28.89mm}>{\raggedright\arraybackslash}p{28.89mm}@{}}
\toprule
\textbf{Case} & \textbf{Masculine} & \textbf{Neuter} & \textbf{Feminine} \\
\midrule
Nominative & swilce & swilc & swilca \\
Genitive & swilcra & swilcra & swilcra \\
Dative & swilcum & swilcum & swilcum \\
Accusative & swilce & swilc & swilca \\
\bottomrule
\end{tabular}
\par\endgroup

% mdtable:T17:end
% mdparagraph:P063:begin
Swilc keeps zero ending in feminine nominative singular and neuter nominative/\allowbreak{}accusative plural, as do hwilc, hūlic and gehwilc. In the selected indefinite use, a following ordinary descriptive adjective is strong. Its case, gender and number still agree with the same noun.
% mdparagraph:P063:end

\reserve{8\baselineskip}
% example: NEW-U08-EXAMPLE-10
% mdparagraph:P064:begin
Swilc mann is gōd. Swilce męnn sindon gōde.
% mdparagraph:P064:end

% mdparagraph:P065:begin
Such a man is good. Such men are good.
% mdparagraph:P065:end

% mdparagraph:P066:begin
Here swilce is an inflected masculine plural adjective modifying męnn. The earlier \textbf{swilce “likewise”} is an invariant adverb. Identical spelling does not make the two uses grammatically identical.
% mdparagraph:P066:end

\reserve{7\baselineskip}
% mdparagraph:P067:begin
\section*{4. Other, same, and self}\BookHeading{nav-ch08-learner-0023}{subsection}{4. Other, same, and self}{1}\BookLabel{bk-ch08-learner-a575cb1bd9c9}
% mdparagraph:P067:end

% teaching: NEW-U08-VOC-other
% teaching: NEW-U08-GRAM-u08-other-sense
% mdparagraph:P068:begin
The earlier \textbf{ōðer “second”} also means \textbf{“other, another”}. Context determines the sense. Its strong forms are retained even after an article; this lexical pattern does not acquire a new weak declension.
% mdparagraph:P068:end

% teaching: NEW-U08-PAR-OTHER
\reserve{7\baselineskip}
% mdparagraph:P069:begin
\subsection*{ōðer: other /\allowbreak{} another, singular}\BookHeading{nav-ch08-learner-0024}{subsubsection}{ōðer: other / another, singular}{0}
% mdparagraph:P069:end

% table: NEW-U08-PAR-OTHER-SG
% mdtable:T18:begin
\par\begingroup
\fontsize{10.5}{13.2}\selectfont
\setlength{\LTleft}{0pt}\setlength{\LTright}{\fill}
\noindent
\begin{tabular}{@{}>{\raggedright\arraybackslash}p{28.89mm}>{\raggedright\arraybackslash}p{28.89mm}>{\raggedright\arraybackslash}p{28.89mm}>{\raggedright\arraybackslash}p{28.89mm}@{}}
\toprule
\textbf{Case} & \textbf{Masculine} & \textbf{Neuter} & \textbf{Feminine} \\
\midrule
Nominative & ōðer & ōðer & ōðru \\
Genitive & ōðres & ōðres & ōðerre \\
Dative & ōðrum & ōðrum & ōðerre \\
Accusative & ōðerne & ōðer & ōðre \\
\bottomrule
\end{tabular}
\par\endgroup

% mdtable:T18:end
\reserve{7\baselineskip}
% mdparagraph:P070:begin
\subsection*{ōðer: other /\allowbreak{} another, plural}\BookHeading{nav-ch08-learner-0025}{subsubsection}{ōðer: other / another, plural}{0}
% mdparagraph:P070:end

% table: NEW-U08-PAR-OTHER-PL
% mdtable:T19:begin
\par\begingroup
\fontsize{10.5}{13.2}\selectfont
\setlength{\LTleft}{0pt}\setlength{\LTright}{\fill}
\noindent
\begin{tabular}{@{}>{\raggedright\arraybackslash}p{28.89mm}>{\raggedright\arraybackslash}p{28.89mm}>{\raggedright\arraybackslash}p{28.89mm}>{\raggedright\arraybackslash}p{28.89mm}@{}}
\toprule
\textbf{Case} & \textbf{Masculine} & \textbf{Neuter} & \textbf{Feminine} \\
\midrule
Nominative & ōðre & ōðru & ōðra \\
Genitive & ōðerra & ōðerra & ōðerra \\
Dative & ōðrum & ōðrum & ōðrum \\
Accusative & ōðre & ōðru & ōðra \\
\bottomrule
\end{tabular}
\par\endgroup

% mdtable:T19:end
% mdparagraph:P071:begin
Retain ōðr- before the ordinary vowel endings, but ōðer- in ōðerne, ōðerre and ōðerra. In particular, feminine nominative singular and neuter nominative/\allowbreak{}accusative plural are ōðru. The complete earlier morphology is reused with the new sense.
% mdparagraph:P071:end

% teaching: NEW-U08-VOC-ilca
% teaching: NEW-U08-GRAM-u08-ilca
% mdparagraph:P072:begin
\textbf{Ilca “same” is weak only.}\BookPoint{bk-ch08-learner-788b7927080a} It ordinarily follows the article. Any following ordinary descriptive adjective is also weak in that definite phrase. Ilca has no strong declension.
% mdparagraph:P072:end

% teaching: NEW-U08-PAR-ILCA-WEAK
\reserve{7\baselineskip}
% mdparagraph:P073:begin
\subsection*{ilca: weak declension, singular}\BookHeading{nav-ch08-learner-0026}{subsubsection}{ilca: weak declension, singular}{0}
% mdparagraph:P073:end

% table: NEW-U08-PAR-ILCA-WEAK-SG
% mdtable:T20:begin
\par\begingroup
\fontsize{10.5}{13.2}\selectfont
\setlength{\LTleft}{0pt}\setlength{\LTright}{\fill}
\noindent
\begin{tabular}{@{}>{\raggedright\arraybackslash}p{28.89mm}>{\raggedright\arraybackslash}p{28.89mm}>{\raggedright\arraybackslash}p{28.89mm}>{\raggedright\arraybackslash}p{28.89mm}@{}}
\toprule
\textbf{Case} & \textbf{Masculine} & \textbf{Neuter} & \textbf{Feminine} \\
\midrule
Nominative & ilca & ilce & ilce \\
Genitive & ilcan & ilcan & ilcan \\
Dative & ilcan & ilcan & ilcan \\
Accusative & ilcan & ilce & ilcan \\
\bottomrule
\end{tabular}
\par\endgroup

% mdtable:T20:end
\reserve{7\baselineskip}
% mdparagraph:P074:begin
\subsection*{ilca: weak declension, plural}\BookHeading{nav-ch08-learner-0027}{subsubsection}{ilca: weak declension, plural}{0}
% mdparagraph:P074:end

% table: NEW-U08-PAR-ILCA-WEAK-PL
% mdtable:T21:begin
\par\begingroup
\fontsize{10.5}{13.2}\selectfont
\setlength{\LTleft}{0pt}\setlength{\LTright}{\fill}
\noindent
\begin{tabular}{@{}>{\raggedright\arraybackslash}p{28.89mm}>{\raggedright\arraybackslash}p{28.89mm}>{\raggedright\arraybackslash}p{28.89mm}>{\raggedright\arraybackslash}p{28.89mm}@{}}
\toprule
\textbf{Case} & \textbf{Masculine} & \textbf{Neuter} & \textbf{Feminine} \\
\midrule
Nominative & ilcan & ilcan & ilcan \\
Genitive & ilcra & ilcra & ilcra \\
Dative & ilcum & ilcum & ilcum \\
Accusative & ilcan & ilcan & ilcan \\
\bottomrule
\end{tabular}
\par\endgroup

% mdtable:T21:end
% mdparagraph:P075:begin
The stem ilc- takes masculine nominative singular -a; feminine/\allowbreak{}neuter nominative singular and neuter accusative singular -e; the remaining singular endings are -an. Plural nominative and accusative have -an, genitive -ra and dative -um.
% mdparagraph:P075:end

\reserve{8\baselineskip}
% example: NEW-U08-EXAMPLE-08
% mdparagraph:P076:begin
Se ōðer mann hæfð ōðer scip. Se ilca mann hæfð ðæt scip.
% mdparagraph:P076:end

% mdparagraph:P077:begin
The other man has another ship. The same man has that ship.
% mdparagraph:P077:end

% mdparagraph:P078:begin
The first sentence distinguishes another man and another ship; it need not place them second in a numbered series. The second identifies the same man again. Both ōðer and ilca agree with the masculine nominative singular mann, but each follows its own declensional pattern.
% mdparagraph:P078:end

% teaching: NEW-U08-VOC-self
% teaching: NEW-U08-GRAM-u08-self
% mdparagraph:P079:begin
\textbf{Self}\BookPoint{bk-ch08-learner-d0b92b93014c} emphasizes or reinforces the person or thing referred to. It agrees with that referent in case, number and gender; it is not an invariant ending mechanically attached to an English-style pronoun compound. Both strong and weak forms occur. Strong self is used in the ordinary personal emphasis illustrated below; the weak system is retained separately, with the familiar weak endings.
% mdparagraph:P079:end

% teaching: NEW-U08-PAR-SELF-STRONG
\reserve{7\baselineskip}
% mdparagraph:P080:begin
\subsection*{self: strong declension, singular}\BookHeading{nav-ch08-learner-0028}{subsubsection}{self: strong declension, singular}{0}
% mdparagraph:P080:end

% table: NEW-U08-PAR-SELF-STRONG-SG
% mdtable:T22:begin
\par\begingroup
\fontsize{10.5}{13.2}\selectfont
\setlength{\LTleft}{0pt}\setlength{\LTright}{\fill}
\noindent
\begin{tabular}{@{}>{\raggedright\arraybackslash}p{28.89mm}>{\raggedright\arraybackslash}p{28.89mm}>{\raggedright\arraybackslash}p{28.89mm}>{\raggedright\arraybackslash}p{28.89mm}@{}}
\toprule
\textbf{Case} & \textbf{Masculine} & \textbf{Neuter} & \textbf{Feminine} \\
\midrule
Nominative & self & self & self \\
Genitive & selfes & selfes & selfre \\
Dative & selfum & selfum & selfre \\
Accusative & selfne & self & selfe \\
\bottomrule
\end{tabular}
\par\endgroup

% mdtable:T22:end
\reserve{7\baselineskip}
% mdparagraph:P081:begin
\subsection*{self: strong declension, plural}\BookHeading{nav-ch08-learner-0029}{subsubsection}{self: strong declension, plural}{0}
% mdparagraph:P081:end

% table: NEW-U08-PAR-SELF-STRONG-PL
% mdtable:T23:begin
\par\begingroup
\fontsize{10.5}{13.2}\selectfont
\setlength{\LTleft}{0pt}\setlength{\LTright}{\fill}
\noindent
\begin{tabular}{@{}>{\raggedright\arraybackslash}p{28.89mm}>{\raggedright\arraybackslash}p{28.89mm}>{\raggedright\arraybackslash}p{28.89mm}>{\raggedright\arraybackslash}p{28.89mm}@{}}
\toprule
\textbf{Case} & \textbf{Masculine} & \textbf{Neuter} & \textbf{Feminine} \\
\midrule
Nominative & selfe & self & selfa \\
Genitive & selfra & selfra & selfra \\
Dative & selfum & selfum & selfum \\
Accusative & selfe & self & selfa \\
\bottomrule
\end{tabular}
\par\endgroup

% mdtable:T23:end
% teaching: NEW-U08-PAR-SELF-WEAK
\reserve{7\baselineskip}
% mdparagraph:P082:begin
\subsection*{self: weak declension, singular}\BookHeading{nav-ch08-learner-0030}{subsubsection}{self: weak declension, singular}{0}
% mdparagraph:P082:end

% table: NEW-U08-PAR-SELF-WEAK-SG
% mdtable:T24:begin
\par\begingroup
\fontsize{10.5}{13.2}\selectfont
\setlength{\LTleft}{0pt}\setlength{\LTright}{\fill}
\noindent
\begin{tabular}{@{}>{\raggedright\arraybackslash}p{28.89mm}>{\raggedright\arraybackslash}p{28.89mm}>{\raggedright\arraybackslash}p{28.89mm}>{\raggedright\arraybackslash}p{28.89mm}@{}}
\toprule
\textbf{Case} & \textbf{Masculine} & \textbf{Neuter} & \textbf{Feminine} \\
\midrule
Nominative & selfa & selfe & selfe \\
Genitive & selfan & selfan & selfan \\
Dative & selfan & selfan & selfan \\
Accusative & selfan & selfe & selfan \\
\bottomrule
\end{tabular}
\par\endgroup

% mdtable:T24:end
\reserve{7\baselineskip}
% mdparagraph:P083:begin
\subsection*{self: weak declension, plural}\BookHeading{nav-ch08-learner-0031}{subsubsection}{self: weak declension, plural}{0}
% mdparagraph:P083:end

% table: NEW-U08-PAR-SELF-WEAK-PL
% mdtable:T25:begin
\par\begingroup
\fontsize{10.5}{13.2}\selectfont
\setlength{\LTleft}{0pt}\setlength{\LTright}{\fill}
\noindent
\begin{tabular}{@{}>{\raggedright\arraybackslash}p{28.89mm}>{\raggedright\arraybackslash}p{28.89mm}>{\raggedright\arraybackslash}p{28.89mm}>{\raggedright\arraybackslash}p{28.89mm}@{}}
\toprule
\textbf{Case} & \textbf{Masculine} & \textbf{Neuter} & \textbf{Feminine} \\
\midrule
Nominative & selfan & selfan & selfan \\
Genitive & selfra & selfra & selfra \\
Dative & selfum & selfum & selfum \\
Accusative & selfan & selfan & selfan \\
\bottomrule
\end{tabular}
\par\endgroup

% mdtable:T25:end
\reserve{8\baselineskip}
% example: NEW-U08-EXAMPLE-09
% mdparagraph:P084:begin
Hē self geseah ðæt scip.
% mdparagraph:P084:end

% mdparagraph:P085:begin
He himself saw that ship.
% mdparagraph:P085:end

% mdparagraph:P086:begin
Self\BookPoint{bk-ch08-learner-cdb9f6787322} is strong nominative masculine singular, emphasizing the subject hē. A personal pronoun can already have reflexive reference without self. Adding self reinforces it: an accusative reflexive takes accusative agreement, and a dative reflexive takes dative agreement. For a male accusative referent the strong form is selfne; the first-person form mē can therefore be reinforced by selfne when the speaker is male. Do not choose self’s case from the English translation “myself” alone.
% mdparagraph:P086:end

\reserve{7\baselineskip}
% mdparagraph:P087:begin
\subsection*{Reference forms and limits}\BookHeading{nav-ch08-learner-0032}{subsubsection}{Reference forms and limits}{0}\BookLabel{bk-ch08-learner-14917e578635}
% mdparagraph:P087:end

% teaching: NEW-U08-VOC-gehwa
% teaching: NEW-U08-GRAM-u08-gehwa-reference
% mdparagraph:P088:begin
\textbf{Gehwā /\allowbreak{} gehwæt} is an indefinite ge-compound of the interrogative pattern: “each one, everyone /\allowbreak{} everything”. Recognize the following complete singular paradigm. Its independent indefinite, partitive and distributive uses are reserved for later work; no plural is supplied.
% mdparagraph:P088:end

% teaching: NEW-U08-REF-GEHWA
\reserve{7\baselineskip}
% mdparagraph:P089:begin
\subsection*{gehwā /\allowbreak{} gehwæt: reference forms}\BookHeading{nav-ch08-learner-0033}{subsubsection}{gehwā / gehwæt: reference forms}{0}
% mdparagraph:P089:end

% table: NEW-U08-REF-GEHWA
% mdtable:T26:begin
\par\begingroup
\fontsize{10.5}{13.2}\selectfont
\setlength{\LTleft}{0pt}\setlength{\LTright}{\fill}
\noindent
\begin{tabular}{@{}>{\raggedright\arraybackslash}p{25.54mm}>{\raggedright\arraybackslash}p{64.36mm}>{\raggedright\arraybackslash}p{28.48mm}@{}}
\toprule
\textbf{Case} & \textbf{Common masculine/\allowbreak{}feminine singular} & \textbf{Neuter singular} \\
\midrule
Nominative & gehwā & gehwæt \\
Genitive & gehwæs & gehwæs \\
Dative & gehwǣm & gehwǣm \\
Accusative & gehwone & gehwæt \\
\bottomrule
\end{tabular}
\par\endgroup

% mdtable:T26:end
% mdparagraph:P090:begin
The distinct strong instrumental forms below are also for recognition. Weak instrumental forms coincide with dative forms. These reference tables do not introduce productive instrumental constructions.
% mdparagraph:P090:end

% teaching: NEW-U08-REF-HWILC-INSTR
% teaching: NEW-U08-REF-HULIC-INSTR
% teaching: NEW-U08-REF-HWAETHER-INSTR
% teaching: NEW-U08-REF-SUM-INSTR
% teaching: NEW-U08-REF-AENIG-INSTR
% teaching: NEW-U08-REF-GEHWILC-INSTR
% teaching: NEW-U08-REF-SWILC-INSTR
% teaching: NEW-U08-REF-SELF-INSTR
% teaching: NEW-U08-REF-OTHER-INSTR
\reserve{7\baselineskip}
% mdparagraph:P091:begin
\subsection*{Instrumental reference forms}\BookHeading{nav-ch08-learner-0034}{subsubsection}{Instrumental reference forms}{0}
% mdparagraph:P091:end

% table: NEW-U08-MATRIX-INSTRUMENTAL-REFERENCE
% mdtable:T27:begin
\par\begingroup
\fontsize{10.5}{13.2}\selectfont
\setlength{\LTleft}{0pt}\setlength{\LTright}{\fill}
\noindent
\begin{tabular}{@{}>{\raggedright\arraybackslash}p{34.61mm}>{\raggedright\arraybackslash}p{45.17mm}>{\raggedright\arraybackslash}p{38.59mm}@{}}
\toprule
\textbf{Word} & \textbf{Masculine singular} & \textbf{Neuter singular} \\
\midrule
hwilc & hwilce & hwilce \\
hūlic & hūlice & hūlice \\
hwæðer & hwæðre & hwæðre \\
sum & sume & sume \\
ǣnig & ǣnige & ǣnige \\
gehwilc & gehwilce & gehwilce \\
swilc & swilce & swilce \\
self & selfe & selfe \\
ōðer & ōðre & ōðre \\
\bottomrule
\end{tabular}
\par\endgroup

% mdtable:T27:end
% mdparagraph:P092:begin
The ordinary uses above leave several related systems for Unit 15: indefinite ān “a certain”, negative pronouns such as nān and nǣnig, either/\allowbreak{}neither families, indeclinable “anything/\allowbreak{}nothing” words, -hwega compounds, indefinite hwā “anyone”, and swā-hwā/\allowbreak{}swā-hwilc free-choice clauses. The impersonal use of man “one” is also separate from the count noun mann already learned. The earlier meanings and limits of eall, ǣlc and the possessives remain in force; sīn is not added here.
% mdparagraph:P092:end

\reserve{7\baselineskip}
% mdparagraph:P093:begin
\section*{5. Relative clauses}\BookHeading{nav-ch08-learner-0035}{subsection}{5. Relative clauses}{1}\BookLabel{bk-ch08-learner-287d30ac3868}
% mdparagraph:P093:end

% teaching: NEW-U08-VOC-the-relative
% teaching: NEW-U08-VOC-se
% teaching: NEW-U08-GRAM-u08-relative-the
% teaching: NEW-U08-GRAM-u08-relative-demonstrative
% teaching: NEW-U08-GRAM-u08-relative-internal-case
% mdparagraph:P094:begin
A \textbf{relative clause} identifies or describes an antecedent in another clause. Old English can introduce it with invariant \textbf{ðe}, with an inflected member of \textbf{se/\allowbreak{}sēo/\allowbreak{}ðæt} alone, or with that inflected member followed by \textbf{ðe}. English often uses “who”, “which” or “that”; an English personal relative does not require the Old English interrogative hwā.
% mdparagraph:P094:end

\reserve{7\baselineskip}
% mdparagraph:P095:begin
\subsection*{Invariant relative particle}\BookHeading{nav-ch08-learner-0036}{subsubsection}{Invariant relative particle}{0}\BookLabel{bk-ch08-learner-149b83a2f42c}
% mdparagraph:P095:end

% table: NEW-U08-MATRIX-RELATIVE-PARTICLE
% mdtable:T28:begin
\par\begingroup
\fontsize{10.5}{13.2}\selectfont
\setlength{\LTleft}{0pt}\setlength{\LTright}{\fill}
\noindent
\begin{tabular}{@{}>{\raggedright\arraybackslash}p{92.06mm}>{\raggedright\arraybackslash}p{29.13mm}@{}}
\toprule
\textbf{Function} & \textbf{Form} \\
\midrule
Who /\allowbreak{} which /\allowbreak{} that, with role determined inside the clause & ðe \\
\bottomrule
\end{tabular}
\par\endgroup

% mdtable:T28:end
% mdparagraph:P096:begin
Ðe does not change its spelling for gender, number or case. Determine its role from the relative clause itself.
% mdparagraph:P096:end

\reserve{8\baselineskip}
% example: NEW-U08-EXAMPLE-11
% mdparagraph:P097:begin
Se mann, ðe ðæt scip hæfð, is mīn gefēra.
% mdparagraph:P097:end

% mdparagraph:P098:begin
The man who has that ship is my companion.
% mdparagraph:P098:end

% mdparagraph:P099:begin
Ðe is the subject of hæfð; the ship is its accusative object. The antecedent se mann is the subject of the main clause is mīn gefēra.
% mdparagraph:P099:end

\reserve{8\baselineskip}
% example: NEW-U08-EXAMPLE-12
% mdparagraph:P100:begin
Ic geseah ðæt scip, ðe se mann hæfð.
% mdparagraph:P100:end

% mdparagraph:P101:begin
I saw the ship that the man has.
% mdparagraph:P101:end

% mdparagraph:P102:begin
Here ðe is the object of hæfð, while se mann is the relative clause’s own nominative subject. The antecedent ðæt scip is the object of the main verb geseah. Invariant spelling does not erase these different functions.
% mdparagraph:P102:end

% mdparagraph:P103:begin
The\BookPoint{bk-ch08-learner-748837dad0a8} inflected relative reuses the complete demonstrative forms. Its \textbf{gender and number} normally follow the antecedent. In the ordinary pattern used for composition here, its \textbf{case} follows its function inside the relative clause. The antecedent has the case required by its separate function in the main clause.
% mdparagraph:P103:end

% teaching: NEW-U08-PAR-RELATIVE-SE
\reserve{7\baselineskip}
% mdparagraph:P104:begin
\subsection*{se /\allowbreak{} sēo /\allowbreak{} ðæt: relative forms, singular}\BookHeading{nav-ch08-learner-0037}{subsubsection}{se / sēo / ðæt: relative forms, singular}{0}
% mdparagraph:P104:end

% table: NEW-U08-PAR-RELATIVE-SE-SG
% mdtable:T29:begin
\par\begingroup
\fontsize{10.5}{13.2}\selectfont
\setlength{\LTleft}{0pt}\setlength{\LTright}{\fill}
\noindent
\begin{tabular}{@{}>{\raggedright\arraybackslash}p{28.89mm}>{\raggedright\arraybackslash}p{28.89mm}>{\raggedright\arraybackslash}p{28.89mm}>{\raggedright\arraybackslash}p{28.89mm}@{}}
\toprule
\textbf{Case} & \textbf{Masculine} & \textbf{Neuter} & \textbf{Feminine} \\
\midrule
Nominative & se & ðæt & sēo \\
Genitive & ðæs & ðæs & ðǣre \\
Dative & ðǣm & ðǣm & ðǣre \\
Accusative & ðone & ðæt & ðā \\
\bottomrule
\end{tabular}
\par\endgroup

% mdtable:T29:end
\reserve{7\baselineskip}
% mdparagraph:P105:begin
\subsection*{se /\allowbreak{} sēo /\allowbreak{} ðæt: relative forms, plural}\BookHeading{nav-ch08-learner-0038}{subsubsection}{se / sēo / ðæt: relative forms, plural}{0}
% mdparagraph:P105:end

% table: NEW-U08-PAR-RELATIVE-SE-PL
% mdtable:T30:begin
\par\begingroup
\fontsize{10.5}{13.2}\selectfont
\setlength{\LTleft}{0pt}\setlength{\LTright}{\fill}
\noindent
\begin{tabular}{@{}>{\raggedright\arraybackslash}p{28.89mm}>{\raggedright\arraybackslash}p{28.89mm}>{\raggedright\arraybackslash}p{28.89mm}>{\raggedright\arraybackslash}p{28.89mm}@{}}
\toprule
\textbf{Case} & \textbf{Masculine} & \textbf{Neuter} & \textbf{Feminine} \\
\midrule
Nominative & ðā & ðā & ðā \\
Genitive & ðāra & ðāra & ðāra \\
Dative & ðǣm & ðǣm & ðǣm \\
Accusative & ðā & ðā & ðā \\
\bottomrule
\end{tabular}
\par\endgroup

% mdtable:T30:end
\reserve{8\baselineskip}
% example: NEW-U08-EXAMPLE-13
% mdparagraph:P106:begin
Ic sęlle ðone fisc ðǣm męnn, se ðe mē healp.
% mdparagraph:P106:end

% mdparagraph:P107:begin
I give the fish to the man who helped me.
% mdparagraph:P107:end

% mdparagraph:P108:begin
Ðǣm męnn is dative, the recipient of sęlle. The relative se is nominative, the subject of healp. The person helped is dative mē. The antecedent’s dative does not make the relative subject dative.
% mdparagraph:P108:end

\reserve{8\baselineskip}
% example: NEW-U08-EXAMPLE-14
% mdparagraph:P109:begin
Ic geseah ðone mann, ðǣm ðe ic healp.
% mdparagraph:P109:end

% mdparagraph:P110:begin
I saw the man whom I helped.
% mdparagraph:P110:end

% mdparagraph:P111:begin
Ðone mann is accusative after geseah. Within the relative, ic is subject and helpan governs dative ðǣm. The relative’s gender and number still refer to the masculine singular man.
% mdparagraph:P111:end

\reserve{8\baselineskip}
% example: NEW-U08-EXAMPLE-16
% mdparagraph:P112:begin
Se mann, ðæs scip ic geseah, is gōd.
% mdparagraph:P112:end

% mdparagraph:P113:begin
The man whose ship I saw is good.
% mdparagraph:P113:end

% mdparagraph:P114:begin
Ðæs\BookPoint{bk-ch08-learner-9b001c83b737} is genitive, expressing possession inside the relative: \textbf{whose ship}. Its antecedent mann is masculine singular and nominative in the main clause. Scip is neuter singular but accusative, the object of geseah; ic is the subject. The possessor’s antecedent determines the relative’s gender and number, not the possessed ship. A possessive relative therefore does not force its possessed noun into the genitive.
% mdparagraph:P114:end

% teaching: NEW-U08-GRAM-u08-relative-person-agreement
% mdparagraph:P115:begin
When\BookPoint{bk-ch08-learner-a5e8eb6b52df} a relative subject refers to a personal-pronoun antecedent, its verb can retain that antecedent’s person. In the selected second-person clauses, use the second-person verb form; the English word “who” does not automatically make the Old English verb third person.
% mdparagraph:P115:end

\reserve{8\baselineskip}
% example: NEW-U08-EXAMPLE-15
% mdparagraph:P116:begin
Ðū, ðe mē hulpe, sęle mē ðone fisc.
% mdparagraph:P116:end

% mdparagraph:P117:begin
You, who helped me, give me that fish.
% mdparagraph:P117:end

% mdparagraph:P118:begin
Hulpe is second-person singular preterite, agreeing with the antecedent ðū. Sęle is the singular command addressed to that person. With a plural antecedent the relative verb and the command must both change where their forms require it.
% mdparagraph:P118:end

% teaching: NEW-U08-GRAM-u08-relative-case-attraction
% teaching: NEW-U08-GRAM-u08-relative-personal-reference
% mdparagraph:P119:begin
\textbf{Recognition notes.}\BookPoint{bk-ch08-learner-4171ac2e2f18} An inflected relative sometimes takes its antecedent’s case instead of the case expected from its own role. This is \textbf{case attraction}, an attested exception; it does not replace the ordinary internal-case pattern required in this unit’s compositions. A less common strategy uses ðe together with a personal pronoun whose inflection marks the internal function. Recognize that strategy when explained in a text; independent composition with it comes later. Neither note changes the earlier limits on personal genitives or instrumental syntax.
% mdparagraph:P119:end

\reserve{7\baselineskip}
% mdparagraph:P120:begin
\section*{6. Additional verb systems and lexical constructions}\BookHeading{nav-ch08-learner-0039}{subsection}{6. Additional verb systems and lexical constructions}{1}
% mdparagraph:P120:end

% teaching: NEW-U08-VOC-aeteowian
% teaching: NEW-U08-GRAM-u08-aeteowian-finite
% teaching: NEW-U08-GRAM-u08-appear-intransitive
\reserve{7\baselineskip}
% mdparagraph:P121:begin
\subsection*{Ætēowian: appear}\BookHeading{nav-ch08-learner-0040}{subsubsection}{Ætēowian: appear}{0}
% mdparagraph:P121:end

% mdparagraph:P122:begin
\textbf{Ætēowian} is weak class II and \textbf{intransitive}: the thing appearing is the nominative subject. It is distinct from transitive ætīewan “reveal, display”. The complete present and preterite indicative and subjunctive are given below.
% mdparagraph:P122:end

% mdparagraph:P123:begin
Keep the lexical stem ætēow-. The present indicative has first-singular -ie, second-singular -ast, third-singular -að and plural -iað. The weak preterite has -ode, -odest, -ode and plural -odon. For the subjunctive, add singular -e or plural -en to the present stem ætēowi- or the preterite stem ætēowod-, giving the exact spellings shown.
% mdparagraph:P123:end

% teaching: NEW-U08-PAR-AETEOWIAN-FINITE
\reserve{7\baselineskip}
% mdparagraph:P124:begin
\subsection*{ætēowian: complete indicative}\BookHeading{nav-ch08-learner-0041}{subsubsection}{ætēowian: complete indicative}{0}
% mdparagraph:P124:end

% table: NEW-U08-PAR-AETEOWIAN-FINITE-INDICATIVE
% mdtable:T31:begin
% course-table-style: 2026-10
% course-panel:T31-P01:begin
\par\begingroup
\fontsize{10.5}{13.2}\selectfont
\setlength{\LTleft}{0pt}\setlength{\LTright}{\fill}
\noindent
\begin{tabular}{@{}>{\raggedright\arraybackslash}p{39.459mm}>{\raggedright\arraybackslash}p{39.459mm}>{\raggedright\arraybackslash}p{39.459mm}@{}}
\toprule
\textbf{Person and number} & \textbf{Present indicative} & \textbf{Preterite indicative} \\
\midrule
1 sg. & ætēowie & ætēowode \\
2 sg. & ætēowast & ætēowodest \\
3 sg. & ætēowað & ætēowode \\
1–3 pl. & ætēowiað & ætēowodon \\
\bottomrule
\end{tabular}
\par\endgroup
% course-panel:T31-P01:end
% mdtable:T31:end
\reserve{7\baselineskip}
% mdparagraph:P125:begin
\subsection*{ætēowian: complete subjunctive}\BookHeading{nav-ch08-learner-0042}{subsubsection}{ætēowian: complete subjunctive}{0}
% mdparagraph:P125:end

% table: NEW-U08-PAR-AETEOWIAN-FINITE-SUBJUNCTIVE
% mdtable:T32:begin
% course-table-style: 2026-10
% course-panel:T32-P01:begin
\par\begingroup
\fontsize{10.5}{13.2}\selectfont
\setlength{\LTleft}{0pt}\setlength{\LTright}{\fill}
\noindent
\begin{tabular}{@{}>{\raggedright\arraybackslash}p{39.459mm}>{\raggedright\arraybackslash}p{39.459mm}>{\raggedright\arraybackslash}p{39.459mm}@{}}
\toprule
\textbf{Person and number} & \textbf{Present subjunctive} & \textbf{Preterite subjunctive} \\
\midrule
1–3 sg. & ætēowie & ætēowode \\
1–3 pl. & ætēowien & ætēowoden \\
\bottomrule
\end{tabular}
\par\endgroup
% course-panel:T32-P01:end
% mdtable:T32:end
% mdparagraph:P126:begin
The present subjunctive can express “let … appear” in the wish or direction already learned. No separate lǣtan is needed for that meaning. Cook’s actual form in the Creation text is \textbf{ætēowie}, without g. The documented weak-II i/\allowbreak{}ig(e)/\allowbreak{}g variation does not replace the selected i spellings in these tables.
% mdparagraph:P126:end

\reserve{7\baselineskip}
% mdparagraph:P127:begin
\subsection*{ætēowian reference citations}\BookHeading{nav-ch08-learner-0043}{subsubsection}{ætēowian reference citations}{0}
% mdparagraph:P127:end

% table: NEW-U08-MATRIX-AETEOWIAN-CITATIONS
% mdtable:T33:begin
\par\begingroup
\fontsize{10.5}{13.2}\selectfont
\setlength{\LTleft}{0pt}\setlength{\LTright}{\fill}
\noindent
\begin{tabular}{@{}>{\raggedright\arraybackslash}p{79.47mm}>{\raggedright\arraybackslash}p{41.72mm}@{}}
\toprule
\textbf{Function} & \textbf{Form} \\
\midrule
Dictionary infinitive & ætēowian \\
Past-participle citation & ætēowod \\
\bottomrule
\end{tabular}
\par\endgroup

% mdtable:T33:end
% mdparagraph:P128:begin
The infinitive and past participle in this small table are dictionary reference forms. This unit’s new productive system is finite indicative and subjunctive; it does not add commands or new nonfinite constructions of ætēowian.
% mdparagraph:P128:end

\reserve{7\baselineskip}
% mdparagraph:P129:begin
\subsection*{Scīnan: subjunctive}\BookHeading{nav-ch08-learner-0044}{subsubsection}{Scīnan: subjunctive}{0}
% mdparagraph:P129:end

% teaching: NEW-U08-VOC-scinan
% teaching: NEW-U08-GRAM-u08-scinan-subj
% mdparagraph:P130:begin
The indicative of strong-I \textbf{scīnan “shine”} is familiar. Its complete subjunctive uses present scīn- and the preterite plural stem scin-. The long ī belongs to the present; the past subjunctive has short i, not the scān- of the preterite indicative singular.
% mdparagraph:P130:end

% teaching: NEW-U08-PAR-SCINAN-SUBJ
\reserve{7\baselineskip}
% mdparagraph:P131:begin
\subsection*{scīnan: subjunctive}\BookHeading{nav-ch08-learner-0045}{subsubsection}{scīnan: subjunctive}{0}
% mdparagraph:P131:end

% table: NEW-U08-PAR-SCINAN-SUBJ
% mdtable:T34:begin
% course-table-style: 2026-10
% course-panel:T34-P01:begin
\par\begingroup
\fontsize{10.5}{13.2}\selectfont
\setlength{\LTleft}{0pt}\setlength{\LTright}{\fill}
\noindent
\begin{tabular}{@{}>{\raggedright\arraybackslash}p{39.459mm}>{\raggedright\arraybackslash}p{39.459mm}>{\raggedright\arraybackslash}p{39.459mm}@{}}
\toprule
\textbf{Person and number} & \textbf{Present subjunctive} & \textbf{Preterite subjunctive} \\
\midrule
1–3 sg. & scīne & scine \\
1–3 pl. & scīnen & scinen \\
\bottomrule
\end{tabular}
\par\endgroup
% course-panel:T34-P01:end
% mdtable:T34:end
% mdparagraph:P132:begin
Each singular or plural row covers all three persons, in both tenses. The subject and the clause determine the person; the subjunctive ending itself distinguishes the two numbers.
% mdparagraph:P132:end

\reserve{7\baselineskip}
% mdparagraph:P133:begin
\subsection*{Geblētsian: commands}\BookHeading{nav-ch08-learner-0046}{subsubsection}{Geblētsian: commands}{0}
% mdparagraph:P133:end

% teaching: NEW-U08-VOC-gebletsian
% teaching: NEW-U08-GRAM-u08-gebletsian-imperative
% mdparagraph:P134:begin
Weak-II \textbf{geblētsian “bless”} has singular imperative -a and plural -iað. Retain the prefix ge-. The number of people being blessed is the object’s number; the command form depends on the number of people addressed.
% mdparagraph:P134:end

% teaching: NEW-U08-PAR-GEBLETSIAN-IMP
\reserve{7\baselineskip}
% mdparagraph:P135:begin
\subsection*{geblētsian: commands}\BookHeading{nav-ch08-learner-0047}{subsubsection}{geblētsian: commands}{0}
% mdparagraph:P135:end

% table: NEW-U08-PAR-GEBLETSIAN-IMP
% mdtable:T35:begin
% course-table-style: 2026-10
% course-panel:T35-P01:begin
\par\begingroup
\fontsize{10.5}{13.2}\selectfont
\setlength{\LTleft}{0pt}\setlength{\LTright}{\fill}
\noindent
\begin{tabular}{@{}>{\raggedright\arraybackslash}p{63.342mm}>{\raggedright\arraybackslash}p{57.846mm}@{}}
\toprule
\textbf{Person and number} & \textbf{Imperative} \\
\midrule
2 sg. & geblētsa \\
2 pl. & geblētsiað \\
\bottomrule
\end{tabular}
\par\endgroup
% course-panel:T35-P01:end
% mdtable:T35:end
\reserve{7\baselineskip}
% mdparagraph:P136:begin
\subsection*{Līf and the cognate accusative}\BookHeading{nav-ch08-learner-0048}{subsubsection}{Līf and the cognate accusative}{0}\BookLabel{bk-ch08-learner-f078b5ed9d15}
% mdparagraph:P136:end

% teaching: NEW-U08-VOC-lif
% teaching: NEW-U08-GRAM-u08-lif-plural
% teaching: NEW-U08-GRAM-u08-cognate-accusative
% mdparagraph:P137:begin
Neuter \textbf{līf “life”} can name an individual life, with the ordinary plural “lives”. Its long stem has no ending in nominative or accusative plural. The full declension is repeated here for this use.
% mdparagraph:P137:end

% teaching: NEW-U08-PAR-LIF
\reserve{7\baselineskip}
% mdparagraph:P138:begin
\subsection*{līf: complete declension}\BookHeading{nav-ch08-learner-0049}{subsubsection}{līf: complete declension}{0}
% mdparagraph:P138:end

% table: NEW-U08-PAR-LIF
% mdtable:T36:begin
\par\begingroup
\fontsize{10.5}{13.2}\selectfont
\setlength{\LTleft}{0pt}\setlength{\LTright}{\fill}
\noindent
\begin{tabular}{@{}>{\raggedright\arraybackslash}p{39.46mm}>{\raggedright\arraybackslash}p{39.46mm}>{\raggedright\arraybackslash}p{39.46mm}@{}}
\toprule
\textbf{Case} & \textbf{Singular} & \textbf{Plural} \\
\midrule
Nominative & līf & līf \\
Genitive & līfes & līfa \\
Dative & līfe & līfum \\
Accusative & līf & līf \\
\bottomrule
\end{tabular}
\par\endgroup

% mdtable:T36:end
% mdparagraph:P139:begin
\textbf{Libban} can take the related noun līf as a \textbf{cognate accusative}: “live one’s life”, “live our lives”. The object names the life being lived; it is not a locative phrase. In “our lives”, the possessive is neuter accusative plural \textbf{ūru}, agreeing with plural līf. This is a particular construction of libban with līf, not a rule permitting an arbitrary object after every intransitive verb. Cook I.8.8 refers to his §168.1.
% mdparagraph:P139:end

\reserve{7\baselineskip}
% mdparagraph:P140:begin
\subsection*{Fill; have in one’s control}\BookHeading{nav-ch08-learner-0050}{subsubsection}{Fill; have in one’s control}{0}\BookLabel{bk-ch08-learner-0244beef0db6}
% mdparagraph:P140:end

% teaching: NEW-U08-VOC-gefyllan
% teaching: NEW-U08-GRAM-u08-gefyllan-fill
% mdparagraph:P141:begin
\textbf{Gefyllan}, weak I, now also has the ordinary active sense \textbf{“fill”}, with the thing or place filled in the accusative. Keep this sense distinct from the earlier “finish, accomplish”, as in finishing a work. The already learned finite and imperative forms serve both senses; no new endings are required.
% mdparagraph:P141:end

% teaching: NEW-U08-VOC-on
% teaching: NEW-U08-GRAM-u08-control-on
% mdparagraph:P142:begin
\textbf{Habban + accusative controlled entity + on + dative possessive/\allowbreak{}gewealde} means “have in one’s control”. Geweald is neuter; its dative singular is \textbf{gewealde}. An agreeing possessive takes dative neuter singular: for one addressee, \textbf{ðīnum}. The invariant possessor \textbf{his} keeps its form. This is the established control expression, not a spatial translation of English “under”. The controlled persons or things remain the accusative object of habban.
% mdparagraph:P142:end

\reserve{7\baselineskip}
% mdparagraph:P143:begin
\section*{7. Practice with forms and clauses}\BookHeading{nav-ch08-learner-0051}{subsection}{7. Practice with forms and clauses}{1}
% mdparagraph:P143:end

% mdparagraph:P144:begin
Use the selected table forms when a task requests a paradigm. In composition, preserve the stated person, number, tense and construction. The questions and relatives use the ordinary case patterns explained above.
% mdparagraph:P144:end

\reserve{7\baselineskip}
% mdparagraph:P145:begin
\subsection*{Practice A. Complete forms}\BookHeading{nav-ch08-learner-0052}{subsubsection}{Practice A. Complete forms}{0}
% mdparagraph:P145:end

\reserve{5\baselineskip}
\par\begin{minipage}{\linewidth}\raggedright
% mdparagraph:P146:begin
\textbf{A1.}
% mdparagraph:P146:end

% item: NEW-U08-A01
% mdparagraph:P147:begin
Give the nominative, genitive, dative and accusative singular of hwā and hwæt. Label each form and distinguish the common personal column from the neuter column.
% mdparagraph:P147:end

\end{minipage}\par\addvspace{4pt}

\par\begin{minipage}{\linewidth}\raggedright
% mdparagraph:P148:begin
\textbf{A2.}
% mdparagraph:P148:end

% item: NEW-U08-A02
% mdparagraph:P149:begin
Give the strong singular genitive, dative and accusative of hwilc in all three genders, and its nominative plural in all three genders.
% mdparagraph:P149:end

\end{minipage}\par\addvspace{4pt}

\par\begin{minipage}{\linewidth}\raggedright
% mdparagraph:P150:begin
\textbf{A3.}
% mdparagraph:P150:end

% item: NEW-U08-A03
% mdparagraph:P151:begin
Give the strong masculine accusative singular, feminine dative singular, neuter nominative plural and genitive plural of hwæðer and hūlic. Label every form.
% mdparagraph:P151:end

\end{minipage}\par\addvspace{4pt}

\par\begin{minipage}{\linewidth}\raggedright
% mdparagraph:P152:begin
\textbf{A4.}
% mdparagraph:P152:end

% item: NEW-U08-A04
% mdparagraph:P153:begin
Give the complete masculine singular and neuter plural declension of ilca. Explain why a strong alternative is not the model assigned here.
% mdparagraph:P153:end

\end{minipage}\par\addvspace{4pt}

\par\begin{minipage}{\linewidth}\raggedright
% mdparagraph:P154:begin
\textbf{A5.}
% mdparagraph:P154:end

% item: NEW-U08-A05
% mdparagraph:P155:begin
Give the nominative singular of self in all three genders, first strong and then weak. Also give its strong masculine accusative singular and strong dative plural.
% mdparagraph:P155:end

\end{minipage}\par\addvspace{4pt}

\par\begin{minipage}{\linewidth}\raggedright
% mdparagraph:P156:begin
\textbf{A6.}
% mdparagraph:P156:end

% item: NEW-U08-A06
% mdparagraph:P157:begin
Give the strong masculine accusative singular, feminine genitive singular, neuter nominative plural and dative plural of sum, ǣnig, gehwilc and swilc.
% mdparagraph:P157:end

\end{minipage}\par\addvspace{4pt}

\par\begin{minipage}{\linewidth}\raggedright
% mdparagraph:P158:begin
\textbf{A7.}
% mdparagraph:P158:end

% item: NEW-U08-A07
% mdparagraph:P159:begin
Give the complete singular and plural declension of līf. Explain why the plural in “our lives” is an ordinary contextual plural in this unit.
% mdparagraph:P159:end

\end{minipage}\par\addvspace{4pt}

\par\begin{minipage}{\linewidth}\raggedright
% mdparagraph:P160:begin
\textbf{A8.}
% mdparagraph:P160:end

% item: NEW-U08-A08
% mdparagraph:P161:begin
Give the complete present and preterite indicative of ætēowian, then the singular and plural present and preterite subjunctive. State which persons each subjunctive entry covers.
% mdparagraph:P161:end

\end{minipage}\par\addvspace{4pt}

\par\begin{minipage}{\linewidth}\raggedright
% mdparagraph:P162:begin
\textbf{A9.}
% mdparagraph:P162:end

% item: NEW-U08-A09
% mdparagraph:P163:begin
Give the singular and plural present and preterite subjunctive of scīnan, stating which persons each entry covers. Then give both second-person commands of geblētsian and identify the number of people addressed.
% mdparagraph:P163:end

\end{minipage}\par\addvspace{4pt}

\reserve{7\baselineskip}
% mdparagraph:P164:begin
\subsection*{Practice Q. Questions}\BookHeading{nav-ch08-learner-0053}{subsubsection}{Practice Q. Questions}{0}
% mdparagraph:P164:end

\reserve{5\baselineskip}
\par\begin{minipage}{\linewidth}\raggedright
% mdparagraph:P165:begin
\textbf{Q1.}
% mdparagraph:P165:end

% item: NEW-U08-Q01
% mdparagraph:P166:begin
Turn each statement into a yes/\allowbreak{}no question: Hē hæfð scip. Ðā męnn sindon gōde. Ðū gesāwe ðæt scip. Preserve person, number and tense.
% mdparagraph:P166:end

\end{minipage}\par\addvspace{4pt}

\par\begin{minipage}{\linewidth}\raggedright
% mdparagraph:P167:begin
\textbf{Q2.}
% mdparagraph:P167:end

% item: NEW-U08-Q02
% mdparagraph:P168:begin
Translate into Old English: Who has that ship? Whom does the man see? What does the man have? State the case of each interrogative.
% mdparagraph:P168:end

\end{minipage}\par\addvspace{4pt}

\par\begin{minipage}{\linewidth}\raggedright
% mdparagraph:P169:begin
\textbf{Q3.}
% mdparagraph:P169:end

% item: NEW-U08-Q03
% mdparagraph:P170:begin
Translate into Old English: To whom do you (sing.) give that fish? Whose ship do you (sing.) see? Explain what governs each interrogative case.
% mdparagraph:P170:end

\end{minipage}\par\addvspace{4pt}

\par\begin{minipage}{\linewidth}\raggedright
% mdparagraph:P171:begin
\textbf{Q4.}
% mdparagraph:P171:end

% item: NEW-U08-Q04
% mdparagraph:P172:begin
Two men are being compared. Ask “Which of the two men has that ship?” using attributive hwæðer with singular mann. Then ask “What sort of ship is that?” using hūlic.
% mdparagraph:P172:end

\end{minipage}\par\addvspace{4pt}

\par\begin{minipage}{\linewidth}\raggedright
% mdparagraph:P173:begin
\textbf{Q5.}
% mdparagraph:P173:end

% item: NEW-U08-Q05
% mdparagraph:P174:begin
Translate Hū swimð se mann? and Hwȳ swimð se mann? Explain the difference in the information requested.
% mdparagraph:P174:end

\end{minipage}\par\addvspace{4pt}

\par\begin{minipage}{\linewidth}\raggedright
% mdparagraph:P175:begin
\textbf{Q6.}
% mdparagraph:P175:end

% item: NEW-U08-Q06
% mdparagraph:P176:begin
Ask “Do you (sing.) have any fish?” using ǣnig and singular fisc. A learner adds dōn because English uses do; explain why that addition is unnecessary.
% mdparagraph:P176:end

\end{minipage}\par\addvspace{4pt}

\reserve{7\baselineskip}
% mdparagraph:P177:begin
\subsection*{Practice N. Pronouns and agreement}\BookHeading{nav-ch08-learner-0054}{subsubsection}{Practice N. Pronouns and agreement}{0}
% mdparagraph:P177:end

\reserve{5\baselineskip}
\par\begin{minipage}{\linewidth}\raggedright
% mdparagraph:P178:begin
\textbf{N1.}
% mdparagraph:P178:end

% item: NEW-U08-N01
% mdparagraph:P179:begin
Translate Sum mann hæfð scip and Sume męnn habbað scipu. Identify the subject and object forms and explain the change in number.
% mdparagraph:P179:end

\end{minipage}\par\addvspace{4pt}

\par\begin{minipage}{\linewidth}\raggedright
% mdparagraph:P180:begin
\textbf{N2.}
% mdparagraph:P180:end

% item: NEW-U08-N02
% mdparagraph:P181:begin
Translate: Each man sees a ship. Use gehwilc and singular mann. Then replace “each man” with “each companion” without changing the meaning of the verb.
% mdparagraph:P181:end

\end{minipage}\par\addvspace{4pt}

\par\begin{minipage}{\linewidth}\raggedright
% mdparagraph:P182:begin
\textbf{N3.}
% mdparagraph:P182:end

% item: NEW-U08-N03
% mdparagraph:P183:begin
Translate: The other man has another ship. Use ōðer twice. Explain why “other/\allowbreak{}another” is not automatically “second” here.
% mdparagraph:P183:end

\end{minipage}\par\addvspace{4pt}

\par\begin{minipage}{\linewidth}\raggedright
% mdparagraph:P184:begin
\textbf{N4.}
% mdparagraph:P184:end

% item: NEW-U08-N04
% mdparagraph:P185:begin
Translate: I saw the same man. I give the fish to the same man. Use ilca and give the case of each complete noun phrase.
% mdparagraph:P185:end

\end{minipage}\par\addvspace{4pt}

\par\begin{minipage}{\linewidth}\raggedright
% mdparagraph:P186:begin
\textbf{N5.}
% mdparagraph:P186:end

% item: NEW-U08-N05
% mdparagraph:P187:begin
Translate Hē self geseah ðæt scip. Then translate “I save myself,” spoken by a man, using mē selfne. Distinguish emphasis on the subject from a reflexive object.
% mdparagraph:P187:end

\end{minipage}\par\addvspace{4pt}

\par\begin{minipage}{\linewidth}\raggedright
% mdparagraph:P188:begin
\textbf{N6.}
% mdparagraph:P188:end

% item: NEW-U08-N06
% mdparagraph:P189:begin
Translate Swilce męnn sindon gōde. Explain how the inflected word swilce differs from the invariant adverb swilce “likewise.”
% mdparagraph:P189:end

\end{minipage}\par\addvspace{4pt}

\par\begin{minipage}{\linewidth}\raggedright
% mdparagraph:P190:begin
\textbf{N7.}
% mdparagraph:P190:end

% item: NEW-U08-N07
% mdparagraph:P191:begin
Give a full Old English noun phrase for each: of some men (genitive); to any companion (dative); that same ship (accusative); such little birds (accusative). Label the case and relevant agreement.
% mdparagraph:P191:end

\end{minipage}\par\addvspace{4pt}

\reserve{7\baselineskip}
% mdparagraph:P192:begin
\subsection*{Practice L. Relative clauses}\BookHeading{nav-ch08-learner-0055}{subsubsection}{Practice L. Relative clauses}{0}
% mdparagraph:P192:end

\reserve{5\baselineskip}
\par\begin{minipage}{\linewidth}\raggedright
% mdparagraph:P193:begin
\textbf{L1.}
% mdparagraph:P193:end

% item: NEW-U08-L01
% mdparagraph:P194:begin
Join Se mann is gōd and Se mann hæfð ðæt scip into one sentence with ðe meaning “The man who has that ship is good.” Identify the relative subject.
% mdparagraph:P194:end

\end{minipage}\par\addvspace{4pt}

\par\begin{minipage}{\linewidth}\raggedright
% mdparagraph:P195:begin
\textbf{L2.}
% mdparagraph:P195:end

% item: NEW-U08-L02
% mdparagraph:P196:begin
Join Ic geseah ðæt scip and Se mann hæfð ðæt scip into one sentence with ðe. Retain “I saw” as the main clause.
% mdparagraph:P196:end

\end{minipage}\par\addvspace{4pt}

\par\begin{minipage}{\linewidth}\raggedright
% mdparagraph:P197:begin
\textbf{L3.}
% mdparagraph:P197:end

% item: NEW-U08-L03
% mdparagraph:P198:begin
Combine Ic sęlle ðone fisc ðǣm męnn and Se mann mē healp. Use an inflected demonstrative followed by ðe for the relative. Label the antecedent case and relative case separately.
% mdparagraph:P198:end

\end{minipage}\par\addvspace{4pt}

\par\begin{minipage}{\linewidth}\raggedright
% mdparagraph:P199:begin
\textbf{L4.}
% mdparagraph:P199:end

% item: NEW-U08-L04
% mdparagraph:P200:begin
Combine Ic geseah ðone mann and Ic healp ðǣm męnn. Use an inflected demonstrative followed by ðe. Explain why the two references to the man have different cases.
% mdparagraph:P200:end

\end{minipage}\par\addvspace{4pt}

\par\begin{minipage}{\linewidth}\raggedright
% mdparagraph:P201:begin
\textbf{L5.}
% mdparagraph:P201:end

% item: NEW-U08-L05
% mdparagraph:P202:begin
Translate: You (sing.), who helped me, give me that fish. Use ðe and the preterite of helpan; identify the relative verb’s person and the command’s person.
% mdparagraph:P202:end

\end{minipage}\par\addvspace{4pt}

\par\begin{minipage}{\linewidth}\raggedright
% mdparagraph:P203:begin
\textbf{L6.}
% mdparagraph:P203:end

% item: NEW-U08-L06
% mdparagraph:P204:begin
Change the addressee of your L5 sentence to two or more people. Change every affected pronoun and finite verb, keeping the fish singular.
% mdparagraph:P204:end

\end{minipage}\par\addvspace{4pt}

\par\begin{minipage}{\linewidth}\raggedright
% mdparagraph:P205:begin
\textbf{L7.}
% mdparagraph:P205:end

% item: NEW-U08-L07
% mdparagraph:P206:begin
In Se mann, ðe ðæt scip hæfð, is mīn gefēra and Ic geseah ðæt scip, ðe se mann hæfð, identify each antecedent and each relative role. Does ðe itself change its spelling to show those roles? Then translate Se mann, ðæs scip ic geseah, is gōd and explain the case and function of ðæs and scip.
% mdparagraph:P206:end

\end{minipage}\par\addvspace{4pt}

\par\begin{minipage}{\linewidth}\raggedright
% mdparagraph:P207:begin
\textbf{L8.}
% mdparagraph:P207:end

% item: NEW-U08-L08
% mdparagraph:P208:begin
A learner says that an inflected relative must always have the same case as its antecedent. Correct that rule using L3 or L4. State the ordinary production target and the limited exception noted in the lesson.
% mdparagraph:P208:end

\end{minipage}\par\addvspace{4pt}

\reserve{7\baselineskip}
% mdparagraph:P209:begin
\subsection*{Practice C. Lexical constructions}\BookHeading{nav-ch08-learner-0056}{subsubsection}{Practice C. Lexical constructions}{0}
% mdparagraph:P209:end

\reserve{5\baselineskip}
\par\begin{minipage}{\linewidth}\raggedright
% mdparagraph:P210:begin
\textbf{C1.}
% mdparagraph:P210:end

% item: NEW-U08-C01
% mdparagraph:P211:begin
Translate: They finish the work. They fill the waters. Use gefyllan in both sentences and explain the difference of sense.
% mdparagraph:P211:end

\end{minipage}\par\addvspace{4pt}

\par\begin{minipage}{\linewidth}\raggedright
% mdparagraph:P212:begin
\textbf{C2.}
% mdparagraph:P212:end

% item: NEW-U08-C02
% mdparagraph:P213:begin
Translate: The man has all the animals in his control. Use on with the appropriate form of geweald and an invariant possessor. Then make it a command to one person, “Have all the animals in your control.”
% mdparagraph:P213:end

\end{minipage}\par\addvspace{4pt}

\par\begin{minipage}{\linewidth}\raggedright
% mdparagraph:P214:begin
\textbf{C3.}
% mdparagraph:P214:end

% item: NEW-U08-C03
% mdparagraph:P215:begin
Translate: I live my life. We live our lives. State the case and number of each object and explain why libban can have that object here.
% mdparagraph:P215:end

\end{minipage}\par\addvspace{4pt}

\par\begin{minipage}{\linewidth}\raggedright
% mdparagraph:P216:begin
\textbf{C4.}
% mdparagraph:P216:end

% item: NEW-U08-C04
% mdparagraph:P217:begin
Translate: Dry land appears. Let dry land appear. Use ætēowian and identify the mood of each verb.
% mdparagraph:P217:end

\end{minipage}\par\addvspace{4pt}

\reserve{7\baselineskip}
% mdparagraph:P218:begin
\section*{8. Cook I.8}\BookHeading{nav-ch08-learner-0057}{subsection}{8. Cook I.8}{1}
% mdparagraph:P218:end

% mdparagraph:P219:begin
Translate into Old English. Cook’s complete exercise follows, with its original wording, parentheses and brackets. Where the English leaves an addressee’s number or the number of “fish” unspecified, a consistent defensible choice is possible. Retain an explicitly specified singular or plural.
% mdparagraph:P219:end

\reserve{5\baselineskip}
\par\begin{minipage}{\linewidth}\raggedright
% mdparagraph:P220:begin
\textbf{I.8.1.}
% mdparagraph:P220:end

% item: EX-I-08-01
% mdparagraph:P221:begin
Your works are good.
% mdparagraph:P221:end

\end{minipage}\par\addvspace{4pt}

\par\begin{minipage}{\linewidth}\raggedright
% mdparagraph:P222:begin
\textbf{I.8.2.}
% mdparagraph:P222:end

% item: EX-I-08-02
% mdparagraph:P223:begin
His tree is fruit-bearing.
% mdparagraph:P223:end

\end{minipage}\par\addvspace{4pt}

\par\begin{minipage}{\linewidth}\raggedright
% mdparagraph:P224:begin
\textbf{I.8.3.}
% mdparagraph:P224:end

% item: EX-I-08-03
% mdparagraph:P225:begin
Thou art the man.
% mdparagraph:P225:end

\end{minipage}\par\addvspace{4pt}

\par\begin{minipage}{\linewidth}\raggedright
% mdparagraph:P226:begin
\textbf{I.8.4.}
% mdparagraph:P226:end

% item: EX-I-08-04
% mdparagraph:P227:begin
This year is the fourth.
% mdparagraph:P227:end

\end{minipage}\par\addvspace{4pt}

\par\begin{minipage}{\linewidth}\raggedright
% mdparagraph:P228:begin
\textbf{I.8.5.}
% mdparagraph:P228:end

% item: EX-I-08-05
% mdparagraph:P229:begin
God made these little birds.
% mdparagraph:P229:end

\end{minipage}\par\addvspace{4pt}

\par\begin{minipage}{\linewidth}\raggedright
% mdparagraph:P230:begin
\textbf{I.8.6.}
% mdparagraph:P230:end

% item: EX-I-08-06
% mdparagraph:P231:begin
He created me likewise.
% mdparagraph:P231:end

\end{minipage}\par\addvspace{4pt}

\par\begin{minipage}{\linewidth}\raggedright
% mdparagraph:P232:begin
\textbf{I.8.7.}
% mdparagraph:P232:end

% item: EX-I-08-07
% mdparagraph:P233:begin
Our cattle (plur.) see the grass.
% mdparagraph:P233:end

\end{minipage}\par\addvspace{4pt}

\par\begin{minipage}{\linewidth}\raggedright
% mdparagraph:P234:begin
\textbf{I.8.8.}
% mdparagraph:P234:end

% item: EX-I-08-08
% mdparagraph:P235:begin
We live our lives (168.1).
% mdparagraph:P235:end

\end{minipage}\par\addvspace{4pt}

\par\begin{minipage}{\linewidth}\raggedright
% mdparagraph:P236:begin
\textbf{I.8.9.}
% mdparagraph:P236:end

% item: EX-I-08-09
% mdparagraph:P237:begin
His seeds grow in the earth.
% mdparagraph:P237:end

\end{minipage}\par\addvspace{4pt}

\par\begin{minipage}{\linewidth}\raggedright
% mdparagraph:P238:begin
\textbf{I.8.10.}
% mdparagraph:P238:end

% item: EX-I-08-10
% mdparagraph:P239:begin
The animals rest (themselves).
% mdparagraph:P239:end

\end{minipage}\par\addvspace{4pt}

\par\begin{minipage}{\linewidth}\raggedright
% mdparagraph:P240:begin
\textbf{I.8.11.}
% mdparagraph:P240:end

% item: EX-I-08-11
% mdparagraph:P241:begin
Fill the waters.
% mdparagraph:P241:end

\end{minipage}\par\addvspace{4pt}

\par\begin{minipage}{\linewidth}\raggedright
% mdparagraph:P242:begin
\textbf{I.8.12.}
% mdparagraph:P242:end

% item: EX-I-08-12
% mdparagraph:P243:begin
Let dry land appear.
% mdparagraph:P243:end

\end{minipage}\par\addvspace{4pt}

\par\begin{minipage}{\linewidth}\raggedright
% mdparagraph:P244:begin
\textbf{I.8.13.}
% mdparagraph:P244:end

% item: EX-I-08-13
% mdparagraph:P245:begin
Place (gesęte) the fish in the sea.
% mdparagraph:P245:end

\end{minipage}\par\addvspace{4pt}

% teaching: NEW-U08-COOK13-SUPPORT
% teaching: NEW-U08-VOC-gesettan
% teaching: NEW-U08-GRAM-u08-supplied-placement
% mdparagraph:P246:begin
\textbf{Supplied help for I.8.13.} Gesęte is the exact supplied singular command of weak-I gesęttan “set, place”. Use it as printed. For this placement, the thing placed is accusative and the location uses \textbf{on + dative}. Cook’s comparable placement in heaven has a dative location. This local help supplies the prefixed command and this bounded frame; its unaided recall is not required.
% mdparagraph:P246:end

\reserve{5\baselineskip}
\par\begin{minipage}{\linewidth}\raggedright
% mdparagraph:P247:begin
\textbf{I.8.14.}
% mdparagraph:P247:end

% item: EX-I-08-14
% mdparagraph:P248:begin
Have (sing.) all spirits under [in] your control.
% mdparagraph:P248:end

\end{minipage}\par\addvspace{4pt}

\par\begin{minipage}{\linewidth}\raggedright
% mdparagraph:P249:begin
\textbf{I.8.15.}
% mdparagraph:P249:end

% item: EX-I-08-15
% mdparagraph:P250:begin
Let the stars shine in the heavens.
% mdparagraph:P250:end

\end{minipage}\par\addvspace{4pt}

\par\begin{minipage}{\linewidth}\raggedright
% mdparagraph:P251:begin
\textbf{I.8.16.}
% mdparagraph:P251:end

% item: EX-I-08-16
% mdparagraph:P252:begin
I say indeed, Bless all men.
% mdparagraph:P252:end

\end{minipage}\par\addvspace{4pt}

\par\begin{minipage}{\linewidth}\raggedright
% mdparagraph:P253:begin
\textbf{I.8.17.}
% mdparagraph:P253:end

% item: EX-I-08-17
% mdparagraph:P254:begin
Thou, who madest life, bless me.
% mdparagraph:P254:end

\end{minipage}\par\addvspace{4pt}

\reserve{7\baselineskip}
% mdparagraph:P255:begin
\section*{9. The ship and the fish}\BookHeading{nav-ch08-learner-0058}{subsection}{9. The ship and the fish}{1}
% mdparagraph:P255:end

% mdparagraph:P256:begin
\textbf{New teaching dialogue.} Three men are present. The first speaker asks his companion about the third man.
% mdparagraph:P256:end

% composition: NEW-U08-READING-01
% mdparagraph:P257:begin
\par\begin{minipage}{\linewidth}\raggedright
\readingtext{On ðǣm strande sindon ðrīe męnn and ān scip. Sum mann cwið: “Hwā is se mann, ðe ðæt scip hæfð?” His gefēra cwið: “Se mann is mīn gefēra. Hē hæfð ðā fiscas, ðe wē sēoð.” “Hwǣm sęlst ðū ðone fisc?” “Ic sęlle hine ðǣm męnn, se ðe mē healp.” “Hwilc scip hæfð hē?” “Hē hæfð ðæt lȳtle scip, ðæt ðe ðū gesāwe.”}
\end{minipage}\par
% mdparagraph:P257:end

\reserve{7\baselineskip}
% mdparagraph:P258:begin
\subsection*{Practice D. The connected dialogue}\BookHeading{nav-ch08-learner-0059}{subsubsection}{Practice D. The connected dialogue}{0}
% mdparagraph:P258:end

\reserve{5\baselineskip}
\par\begin{minipage}{\linewidth}\raggedright
% mdparagraph:P259:begin
\textbf{D1.}
% mdparagraph:P259:end

% item: NEW-U08-D01
% mdparagraph:P260:begin
Translate the complete account “The ship and the fish.” Keep the speakers and the three men’s relationships clear.
% mdparagraph:P260:end

\end{minipage}\par\addvspace{4pt}

\par\begin{minipage}{\linewidth}\raggedright
% mdparagraph:P261:begin
\textbf{D2.}
% mdparagraph:P261:end

% item: NEW-U08-D02
% mdparagraph:P262:begin
In the account, list the direct questions. For each interrogative, give its role and the information requested.
% mdparagraph:P262:end

\end{minipage}\par\addvspace{4pt}

\par\begin{minipage}{\linewidth}\raggedright
% mdparagraph:P263:begin
\textbf{D3.}
% mdparagraph:P263:end

% item: NEW-U08-D03
% mdparagraph:P264:begin
In the answer beginning Ic sęlle, explain the roles and cases of hine, ðǣm męnn, se and mē. What identifies the man receiving the fish?
% mdparagraph:P264:end

\end{minipage}\par\addvspace{4pt}

\par\begin{minipage}{\linewidth}\raggedright
% mdparagraph:P265:begin
\textbf{D4.}
% mdparagraph:P265:end

% item: NEW-U08-D04
% mdparagraph:P266:begin
Rewrite the final answer so that the man has two little ships which you (sing.) saw. Preserve its relative construction and change all necessary forms.
% mdparagraph:P266:end

\end{minipage}\par\addvspace{4pt}

\reserve{7\baselineskip}
% mdparagraph:P267:begin
\section*{10. The merchant and his merchandise}\BookHeading{nav-ch08-learner-0060}{subsection}{10. The merchant and his merchandise}{1}
% mdparagraph:P267:end

% mdparagraph:P268:begin
This exact four-turn opening is from Cook’s \emph{First Book}, “Trades and Occupations”, printed pp. 129–130 /\allowbreak{} one-based PDF pp. 145–146. Page 129 has three unnumbered dialogue lines, counted editorially as 1–3; page 130 supplies printed lines 1–7. The original English speaker labels are retained.
% mdparagraph:P268:end

\par\begin{minipage}{\linewidth}\raggedright
% source: FB-TRADES-MERCHANT-OPENING
% mdparagraph:P269:begin
\readingtext{\textbf{Teacher.} Hwæt sægst ðū, mangere?}
% mdparagraph:P269:end

% mdparagraph:P270:begin
\readingtext{\textbf{Merchant.} Ic sęcge ðæt behēfe ic eom ge cyninge and ealdormannum, and weligum, and eallum folce.}
% mdparagraph:P270:end

% mdparagraph:P271:begin
\readingtext{\textbf{Teacher.} And hū?}
% mdparagraph:P271:end

% mdparagraph:P272:begin
\readingtext{\textbf{Merchant.} Ic āstīge mīn scip mid hlæstum mīnum, and rōwe ofer sǣlice dǣlas, and cīepe mīn ðing, and bycge ðing dēorwierðu, ðā on ðissum lande ne bēoð ācęnnede, and ic hit tōgelǣde ēow hider mid miclum plihte ofer sǣ, and hwīlum forlidennesse ic ðolie, mid lyre ealra ðinga mīnra, unēaðe cwic ætberstende.}
% mdparagraph:P272:end

\end{minipage}\par
% mdparagraph:P273:begin
The excerpt ends with the merchant’s escape from shipwreck. The following question, “Hwilc ðing gelǣdst ðū ūs?”, and its goods catalogue are outside this selection.
% mdparagraph:P273:end

\reserve{7\baselineskip}
% mdparagraph:P274:begin
\subsection*{Local help}\BookHeading{nav-ch08-learner-0061}{subsubsection}{Local help}{0}
% mdparagraph:P274:end

% teaching: NEW-U08-SOURCE-HELPS
% mdparagraph:P275:begin
All twelve groups below are supplied for the source-reading tasks. They give the unfamiliar words and constructions in their actual forms. The stated noun gender and verb or adjective class identify the headword; they do not require an unprinted complete paradigm for independent use. The question words and ordinary relative relation are learned material within these supported clauses.
% mdparagraph:P275:end

\reserve{5\baselineskip}
\par\begin{minipage}{\linewidth}\raggedright
% mdparagraph:P276:begin
\textbf{H1. mangere}
% mdparagraph:P276:end

% support: U08-SRC-HELP01
% teaching: NEW-U08-VOC-mangere
% teaching: NEW-U08-GRAM-u08-source-vocative
% mdparagraph:P277:begin
\textbf{Mangere}, masculine strong noun in -e: “merchant”.
% mdparagraph:P277:end

% mdparagraph:P278:begin
mangere, masculine, merchant; the nominative form is used in addressing him.
% mdparagraph:P278:end

\end{minipage}\par\addvspace{4pt}

\par\begin{minipage}{\linewidth}\raggedright
% mdparagraph:P279:begin
\textbf{H2. ðæt behēfe ic eom; And hū?}
% mdparagraph:P279:end

% support: U08-SRC-HELP02
% teaching: NEW-U08-VOC-behefe
% teaching: NEW-U08-VOC-thaet-conjunction
% teaching: NEW-U08-GRAM-u08-source-content-clause
% teaching: NEW-U08-GRAM-u08-source-useful-dative
% teaching: NEW-U08-GRAM-u08-source-elliptical-how
% mdparagraph:P280:begin
\textbf{Behēfe}, strong adjective with nominative -e: “useful”. \textbf{Ðæt} here is the supplied content conjunction.
% mdparagraph:P280:end

% mdparagraph:P281:begin
I say that I am useful. This ðæt introduces a content clause; it is not the neuter demonstrative or a relative. behēfe means useful and takes the dative people to whom the merchant is useful. The next question And hū? abbreviates And how [are you useful]?; its omitted predicate is recovered from this answer. That exact ellipsis is supplied here, without teaching unrestricted question ellipsis.
% mdparagraph:P281:end

\end{minipage}\par\addvspace{4pt}

\par\begin{minipage}{\linewidth}\raggedright
% mdparagraph:P282:begin
\textbf{H3. ge cyninge and ealdormannum, and weligum, and eallum folce}
% mdparagraph:P282:end

% support: U08-SRC-HELP03
% teaching: NEW-U08-VOC-ge-conjunction
% teaching: NEW-U08-VOC-cyning
% teaching: NEW-U08-VOC-ealdormann
% teaching: NEW-U08-VOC-welig
% teaching: NEW-U08-VOC-folc
% teaching: NEW-U08-GRAM-u08-source-ge-and
% teaching: NEW-U08-GRAM-u08-source-substantival-weligum
% mdparagraph:P283:begin
\textbf{Cyning}, masculine a-stem: “king”; \textbf{ealdormann}, masculine compound of mann: “leader, noble”; \textbf{welig}, strong adjective: “wealthy”; \textbf{folc}, neuter a-stem: “people”. \textbf{Ge … and} is the supplied correlative “both … and”.
% mdparagraph:P283:end

% mdparagraph:P284:begin
ge ... and coordinates the recipients: both to the king and to nobles, and to wealthy people, and to all the people. cyninge is dative singular; ealdormannum and weligum are dative plural; folce is dative singular. weligum is an adjective used substantivally.
% mdparagraph:P284:end

\end{minipage}\par\addvspace{4pt}

\par\begin{minipage}{\linewidth}\raggedright
% mdparagraph:P285:begin
\textbf{H4. āstīge mīn scip mid hlæstum mīnum}
% mdparagraph:P285:end

% support: U08-SRC-HELP04
% teaching: NEW-U08-VOC-astigan
% teaching: NEW-U08-VOC-hlaest
% teaching: NEW-U08-VOC-mid
% teaching: NEW-U08-GRAM-u08-source-board-cargo
% mdparagraph:P286:begin
\textbf{Āstīgan}, strong I, here “go aboard, board”; \textbf{hlæst}, neuter a-stem, plural here: “wares, cargo”. \textbf{Mid} means “with” in the locally supplied phrases.
% mdparagraph:P286:end

% mdparagraph:P287:begin
I board my ship with my cargo. āstīge is first singular present; scip is accusative, the thing boarded. mid governs dative hlæstum mīnum here. hlæst is neuter and is plural in this cargo expression.
% mdparagraph:P287:end

\end{minipage}\par\addvspace{4pt}

\par\begin{minipage}{\linewidth}\raggedright
% mdparagraph:P288:begin
\textbf{H5. rōwe ofer sǣlice dǣlas}
% mdparagraph:P288:end

% support: U08-SRC-HELP05
% teaching: NEW-U08-VOC-rowan
% teaching: NEW-U08-VOC-saelic
% teaching: NEW-U08-VOC-dael
% teaching: NEW-U08-GRAM-u08-source-nautical-motion
% mdparagraph:P289:begin
\textbf{Rōwan}, reduplicating strong verb: “row”; \textbf{sǣlic}, strong adjective: “of the sea”; \textbf{dǣl}, masculine a-stem: “part”.
% mdparagraph:P289:end

% mdparagraph:P290:begin
I row over sea regions /\allowbreak{} across stretches of sea. rōwe is first singular present; sǣlice dǣlas is accusative masculine plural after spatial ofer.
% mdparagraph:P290:end

\end{minipage}\par\addvspace{4pt}

\par\begin{minipage}{\linewidth}\raggedright
% mdparagraph:P291:begin
\textbf{H6. cīepe mīn ðing, and bycge ðing dēorwierðu}
% mdparagraph:P291:end

% support: U08-SRC-HELP06
% teaching: NEW-U08-VOC-ciepan
% teaching: NEW-U08-VOC-bycgan
% teaching: NEW-U08-VOC-deorwierth
% teaching: NEW-U08-GRAM-u08-source-buy-sell
% mdparagraph:P292:begin
\textbf{Cīepan}, weak I: “sell”; \textbf{bycgan}, irregular weak I: “buy”; \textbf{dēorwierðe}, strong adjective with nominative -e: “precious, valuable”.
% mdparagraph:P292:end

% mdparagraph:P293:begin
I sell my goods and buy valuable things. cīepe and bycge are first singular presents. Both ðing phrases are accusative plural; mīn and dēorwierðu agree with neuter plural ðing.
% mdparagraph:P293:end

\end{minipage}\par\addvspace{4pt}

\par\begin{minipage}{\linewidth}\raggedright
% mdparagraph:P294:begin
\textbf{H7. ðā on ðissum lande ne bēoð ācęnnede}
% mdparagraph:P294:end

% support: U08-SRC-HELP07
% teaching: NEW-U08-VOC-land
% teaching: NEW-U08-VOC-acennan
% teaching: NEW-U08-GRAM-u08-source-produced-passive
% mdparagraph:P295:begin
\textbf{Land}, neuter a-stem: “land, country”; \textbf{ācęnnan}, weak I: “produce, bring forth”. The entire negative passive \textbf{ne bēoð ācęnnede} is supplied here as “are not produced”.
% mdparagraph:P295:end

% mdparagraph:P296:begin
which are not produced in this land. The target relative ðā is nominative neuter plural, the subject inside its clause; its antecedent ðing is accusative in the main clause. lande is dative singular, neuter land. ācęnnede is a plural predicate past participle, are produced; its exact form is supplied, without assigning the whole verb or a general participial declension.
% mdparagraph:P296:end

\end{minipage}\par\addvspace{4pt}

\par\begin{minipage}{\linewidth}\raggedright
% mdparagraph:P297:begin
\textbf{H8. ic hit tōgelǣde ēow hider}
% mdparagraph:P297:end

% support: U08-SRC-HELP08
% teaching: NEW-U08-VOC-togelaedan
% teaching: NEW-U08-VOC-hider
% teaching: NEW-U08-GRAM-u08-source-bring-goods
% teaching: NEW-U08-GRAM-u08-source-collective-hit
% mdparagraph:P298:begin
\textbf{Tōgelǣdan}, weak I: “bring”; \textbf{hider}, direction adverb: “hither, here”.
% mdparagraph:P298:end

% mdparagraph:P299:begin
I bring it/\allowbreak{}the goods here to you. tōgelǣde is first singular present. ēow is dative plural recipient. hit is the actual neuter singular object used collectively for the merchandise; translate its reference, not as a new unrelated object. hider expresses movement hither.
% mdparagraph:P299:end

\end{minipage}\par\addvspace{4pt}

\par\begin{minipage}{\linewidth}\raggedright
% mdparagraph:P300:begin
\textbf{H9. mid miclum plihte ofer sǣ}
% mdparagraph:P300:end

% support: U08-SRC-HELP09
% teaching: NEW-U08-VOC-pliht
% teaching: NEW-U08-GRAM-u08-source-risk
% mdparagraph:P301:begin
\textbf{Pliht}, masculine a-stem: “peril, risk”. Mid introduces the supplied circumstance of risk here.
% mdparagraph:P301:end

% mdparagraph:P302:begin
at great risk across the sea. mid has a supplied circumstance use here and governs dative masculine singular miclum plihte. In ofer sǣ, sǣ has the same singular form in the accusative and dative. Across the sea motivates an accusative extent interpretation in this context, but the written form does not prove that case. No unique accusative answer is required from this surface alone; the task concerns the meaning across the sea.
% mdparagraph:P302:end

\end{minipage}\par\addvspace{4pt}

\par\begin{minipage}{\linewidth}\raggedright
% mdparagraph:P303:begin
\textbf{H10. hwīlum forlidennesse ic ðolie}
% mdparagraph:P303:end

% support: U08-SRC-HELP10
% teaching: NEW-U08-VOC-hwilum
% teaching: NEW-U08-VOC-forlidennes
% teaching: NEW-U08-VOC-tholian
% teaching: NEW-U08-GRAM-u08-source-shipwreck
% mdparagraph:P304:begin
\textbf{Hwīlum}, adverb: “sometimes”; \textbf{forlidennes}, feminine -nes noun: “shipwreck”; \textbf{ðolian}, weak II: “endure, experience”.
% mdparagraph:P304:end

% mdparagraph:P305:begin
sometimes I suffer shipwreck. hwīlum is a frequency adverb; forlidennesse is accusative feminine singular, the object of first singular present ðolie.
% mdparagraph:P305:end

\end{minipage}\par\addvspace{4pt}

\par\begin{minipage}{\linewidth}\raggedright
% mdparagraph:P306:begin
\textbf{H11. mid lyre ealra ðinga mīnra}
% mdparagraph:P306:end

% support: U08-SRC-HELP11
% teaching: NEW-U08-VOC-lyre
% teaching: NEW-U08-GRAM-u08-source-loss-genitive
% mdparagraph:P307:begin
\textbf{Lyre}, masculine strong noun in -e: “loss”. Mid introduces the supplied circumstance of loss.
% mdparagraph:P307:end

% mdparagraph:P308:begin
with the loss of all my things. lyre is dative masculine singular after mid; ealra ðinga mīnra is a genitive plural dependent of loss.
% mdparagraph:P308:end

\end{minipage}\par\addvspace{4pt}

\par\begin{minipage}{\linewidth}\raggedright
% mdparagraph:P309:begin
\textbf{H12. unēaðe cwic ætberstende}
% mdparagraph:P309:end

% support: U08-SRC-HELP12
% teaching: NEW-U08-VOC-uneathe
% teaching: NEW-U08-VOC-cwic
% teaching: NEW-U08-VOC-aetberstan
% teaching: NEW-U08-GRAM-u08-source-escape-participle
% mdparagraph:P310:begin
\textbf{Unēaðe}, adverb: “with difficulty, barely”; \textbf{cwic}, regular strong adjective: “alive”; \textbf{ætberstan}, strong III: “escape”. Cwic is related to the earlier irregular source form cucu, but retains its own actual spelling and analysis here.
% mdparagraph:P310:end

% mdparagraph:P311:begin
barely escaping alive. unēaðe means with difficulty/\allowbreak{}barely; cwic is nominative masculine singular referring to the merchant; ætberstende is a nominative singular present participle expressing his accompanying circumstance. The construction and exact forms are supplied.
% mdparagraph:P311:end

% mdparagraph:P312:begin
\textbf{Source presentation.} The Old English words, spelling, quantities and punctuation are retained. Printed line wrapping is joined, and the original note cues are moved to the page-qualified list below. The speaker labels remain the source’s English labels. In the notes, typographical emphasis is represented as plain text and line-end word divisions are joined; the continuation of p. 129 note 5 is reunited with its beginning.
% mdparagraph:P312:end

\end{minipage}\par\addvspace{4pt}

\reserve{7\baselineskip}
% mdparagraph:P313:begin
\subsection*{Cook’s original notes}\BookHeading{nav-ch08-learner-0062}{subsubsection}{Cook’s original notes}{0}
% mdparagraph:P313:end

% mdparagraph:P314:begin
These fifteen notes belong to Cook’s edition. Their numbers and grammar references are his, not the section numbers of this course. The local help above supplies the interpretation needed for the reading. All questions within these preserved notes are historical source commentary, not additional course tasks. The long etymological discussion in p. 129 note 5 is preserved as historical source commentary; it is not a new linguistic claim.
% mdparagraph:P314:end

% source-note: U08-ORIG-129-01
% mdparagraph:P315:begin
\textbf{P. 129, note 1 (sægst).} See 123.
% mdparagraph:P315:end

% source-note: U08-ORIG-129-02
% mdparagraph:P316:begin
\textbf{P. 129, note 2 (mangere).} Lat. mercator. Other Old English terms for merchant are cīepa and cīepmann. From a collateral form of the latter, cēapmann, without umlaut, is derived Mod. Eng. chapman. How is chap- related to cheap? See the New English Dictionary (New Eng. Dict.) under these words.
% mdparagraph:P316:end

% source-note: U08-ORIG-129-03
% mdparagraph:P317:begin
\textbf{P. 129, note 3 (behēfe).} Lat. utilis. Cf. the Mod. Eng. noun behoof.
% mdparagraph:P317:end

% source-note: U08-ORIG-129-04
% mdparagraph:P318:begin
\textbf{P. 129, note 4 (ge).} ge . . . and = Lat. et . . . et.
% mdparagraph:P318:end

% source-note: U08-ORIG-129-05
% mdparagraph:P319:begin
\textbf{P. 129, note 5 (cyninge).} Carlyle (Sartor Resartus, Bk. 3, Chap. 7) has the following: “König (King), anciently Könning, means Ken-ning (Cunning), or which is the same thing, Can-ning. Ever must the Sovereign of Mankind be fitly entitled King.” On the other hand Gummere (Germanic Origins, p. 270): “At the head of the family we found, of course, the father; and at the head of the state we naturally look for the king. The word ‘king’ means the child or son of the tribe, its representative or even creation; man of race, man of rank. Gradually the king ceases to be regarded as a creation of his race; his ancestry is pushed back to the gods, and his right is quite above all sanctions of popular choice or approval.” Which of these views is confirmed by etymology?
% mdparagraph:P319:end

% mdparagraph:P320:begin
(This original note continues on p. 130.)
% mdparagraph:P320:end

% source-note: U08-ORIG-129-06
% mdparagraph:P321:begin
\textbf{P. 129, note 6 (ealdormannum).} Lat. ducibus.
% mdparagraph:P321:end

% source-note: U08-ORIG-130-01
% mdparagraph:P322:begin
\textbf{P. 130, note 1 (hlæstum).} Lat. mercibus.
% mdparagraph:P322:end

% source-note: U08-ORIG-130-02
% mdparagraph:P323:begin
\textbf{P. 130, note 2 (rōwe).} Lat. navigo.
% mdparagraph:P323:end

% source-note: U08-ORIG-130-03
% mdparagraph:P324:begin
\textbf{P. 130, note 3 (sǣlice).} Lat. marinas.
% mdparagraph:P324:end

% source-note: U08-ORIG-130-04
% mdparagraph:P325:begin
\textbf{P. 130, note 4 (dǣlas).} Lat. partes.
% mdparagraph:P325:end

% source-note: U08-ORIG-130-05
% mdparagraph:P326:begin
\textbf{P. 130, note 5 (cīepe).} Lat. vendo.
% mdparagraph:P326:end

% source-note: U08-ORIG-130-06
% mdparagraph:P327:begin
\textbf{P. 130, note 6 (ðing/\allowbreak{}dēorwierðu).} Lat. res pretiosas.
% mdparagraph:P327:end

% source-note: U08-ORIG-130-07
% mdparagraph:P328:begin
\textbf{P. 130, note 7 (tōgelǣde).} Lat. adduco.
% mdparagraph:P328:end

% source-note: U08-ORIG-130-08
% mdparagraph:P329:begin
\textbf{P. 130, note 8 (plihte).} Lat. periculo. Mod. Eng. form of pliht?
% mdparagraph:P329:end

% source-note: U08-ORIG-130-09
% mdparagraph:P330:begin
\textbf{P. 130, note 9 (unēaðe/\allowbreak{}cwic/\allowbreak{}ætberstende).} Lat. vix vivus evadens. Note the love for alliteration, even in the Latin.
% mdparagraph:P330:end

\reserve{7\baselineskip}
% mdparagraph:P331:begin
\subsection*{Practice R. The original dialogue}\BookHeading{nav-ch08-learner-0063}{subsubsection}{Practice R. The original dialogue}{0}
% mdparagraph:P331:end

\reserve{5\baselineskip}
\par\begin{minipage}{\linewidth}\raggedright
% mdparagraph:P332:begin
\textbf{R1.}
% mdparagraph:P332:end

% item: NEW-U08-R01
% mdparagraph:P333:begin
Translate the complete Merchant and Teacher excerpt using the local helps. Preserve who speaks each question and answer.
% mdparagraph:P333:end

\end{minipage}\par\addvspace{4pt}

\par\begin{minipage}{\linewidth}\raggedright
% mdparagraph:P334:begin
\textbf{R2.}
% mdparagraph:P334:end

% item: NEW-U08-R02
% mdparagraph:P335:begin
Identify the direct questions in the excerpt and explain what each asks. Use the local help for the merchant’s form of address.
% mdparagraph:P335:end

\end{minipage}\par\addvspace{4pt}

\par\begin{minipage}{\linewidth}\raggedright
% mdparagraph:P336:begin
\textbf{R3.}
% mdparagraph:P336:end

% item: NEW-U08-R03
% mdparagraph:P337:begin
In ðā on ðissum lande ne bēoð ācęnnede, identify the antecedent, the role of ðā and the number of the passive auxiliary. Use the supplied participle meaning.
% mdparagraph:P337:end

\end{minipage}\par\addvspace{4pt}

\par\begin{minipage}{\linewidth}\raggedright
% mdparagraph:P338:begin
\textbf{R4.}
% mdparagraph:P338:end

% item: NEW-U08-R04
% mdparagraph:P339:begin
Using the local helps, explain how the merchant obtains his goods and where he takes them. Distinguish these statements from the dangers he encounters.
% mdparagraph:P339:end

\end{minipage}\par\addvspace{4pt}

\par\begin{minipage}{\linewidth}\raggedright
% mdparagraph:P340:begin
\textbf{R5.}
% mdparagraph:P340:end

% item: NEW-U08-R05
% mdparagraph:P341:begin
Identify two expressions for dangers in the final answer, and explain the contribution of the closing words unēaðe cwic ætberstende. Use the supplied meanings and construction notes.
% mdparagraph:P341:end

\end{minipage}\par\addvspace{4pt}

\reserve{7\baselineskip}
% mdparagraph:P342:begin
\section*{11. Cumulative practice}\BookHeading{nav-ch08-learner-0064}{subsection}{11. Cumulative practice}{1}
% mdparagraph:P342:end

\reserve{7\baselineskip}
% mdparagraph:P343:begin
\subsection*{Practice X. Questions, reference and clause relations}\BookHeading{nav-ch08-learner-0065}{subsubsection}{Practice X. Questions, reference and clause relations}{0}
% mdparagraph:P343:end

\reserve{5\baselineskip}
\par\begin{minipage}{\linewidth}\raggedright
% mdparagraph:P344:begin
\textbf{X1.}
% mdparagraph:P344:end

% item: NEW-U08-X01
% mdparagraph:P345:begin
Translate: Which men have the ships that you (sing.) saw? Use hwilc and ðe. Label the interrogative, antecedent and relative roles.
% mdparagraph:P345:end

\end{minipage}\par\addvspace{4pt}

\par\begin{minipage}{\linewidth}\raggedright
% mdparagraph:P346:begin
\textbf{X2.}
% mdparagraph:P346:end

% item: NEW-U08-X02
% mdparagraph:P347:begin
Translate: I give another fish to the same man who helped me. Use ōðer, ilca and an inflected demonstrative followed by ðe.
% mdparagraph:P347:end

\end{minipage}\par\addvspace{4pt}

\par\begin{minipage}{\linewidth}\raggedright
% mdparagraph:P348:begin
\textbf{X3.}
% mdparagraph:P348:end

% item: NEW-U08-X03
% mdparagraph:P349:begin
Translate: Do you (sing.) see any little birds? Use ǣnig with plural fugol and the appropriate form of lȳtel.
% mdparagraph:P349:end

\end{minipage}\par\addvspace{4pt}

\par\begin{minipage}{\linewidth}\raggedright
% mdparagraph:P350:begin
\textbf{X4.}
% mdparagraph:P350:end

% item: NEW-U08-X04
% mdparagraph:P351:begin
Translate: Each companion himself saw that ship. Use gehwilc and strong self. Then make the whole statement a direct yes/\allowbreak{}no question.
% mdparagraph:P351:end

\end{minipage}\par\addvspace{4pt}

\par\begin{minipage}{\linewidth}\raggedright
% mdparagraph:P352:begin
\textbf{X5.}
% mdparagraph:P352:end

% item: NEW-U08-X05
% mdparagraph:P353:begin
Translate: The man who has the ship says, “Let dry land appear.” Use ðe and retain a direct quotation.
% mdparagraph:P353:end

\end{minipage}\par\addvspace{4pt}

\par\begin{minipage}{\linewidth}\raggedright
% mdparagraph:P354:begin
\textbf{X6.}
% mdparagraph:P354:end

% item: NEW-U08-X06
% mdparagraph:P355:begin
Explain the different functions of ðæt in ðæt scip and ðæt ðe ðū gesāwe in the account. Does their identical spelling make both conjunctions meaning “that”?
% mdparagraph:P355:end

\end{minipage}\par\addvspace{4pt}

\clearpage
% mdparagraph:P356:begin
\section*{Source and editorial notes}\BookHeading{nav-ch08-learner-0066}{subsection}{Source and editorial notes}{1}
% mdparagraph:P356:end

\begingroup\fontsize{10}{12.8}\selectfont
\setlength{\parskip}{1.5pt}
% mdparagraph:P357:begin
The principal sources are Albert S. Cook, \emph{A First Book in Old English: Grammar, Reader, Notes, and Vocabulary}, second edition, revised and enlarged (Ginn, 1895), and \emph{Exercises in Old English} (Ginn, 1895). The complete I.8 prompts are printed on Exercises p. 4 /\allowbreak{} PDF p. 12. The original reader is the opening of “The Merchant and his Merchandise” on First Book pp. 129–130 /\allowbreak{} PDF pp. 145–146. No historical author is claimed for the sixteen teaching examples or the new ship dialogue.
% mdparagraph:P357:end

% mdparagraph:P358:begin
The pronoun systems and relative strategies follow Cook §§86–89, with the strong and weak endings in §§57–60; the supporting checks include Wright §§463–464 and 468–471 and Fulk §§35 and 77. The selected tables are complete paradigms built from checked lexical stems and inflection rules, not claims that every cell is individually attested in this excerpt. The bounded recognition note on case attraction qualifies the ordinary internal-case rule.
% mdparagraph:P358:end

% mdparagraph:P359:begin
The additional systems use Cook’s weak-II account in §§118–119, strong-I scīnan in §102, and the subjunctive endings in §§93–95. Cook §168.1 supplies the cognate accusative of libban. The source appearance form is ætēowie on First Book p. 124 /\allowbreak{} PDF p. 140, with note 20 directing the reader to ætēowian; the dictionary entry is on p. 275 /\allowbreak{} PDF p. 291. The contrasting transitive ætīewan is a different entry. The placement witness is p. 125 /\allowbreak{} PDF p. 141, printed line 15; the have-in-control witness is p. 127 /\allowbreak{} PDF p. 143, printed line 9. These short witnesses support the particular frames taught or supplied here; they do not assign the surrounding Creation passages.
% mdparagraph:P359:end

% mdparagraph:P360:begin
Dictionary compound hyphens are omitted in the displayed headwords ealdormann, āstīgan, dēorwierðe, ācęnnan, tōgelǣdan, forlidennes, unēaðe, ætberstan, ætēowian and gesęttan. Their running source spellings and quantities remain unchanged. Cook’s hūlic is retained beside the separately documented hūlīc spelling. All PDF pages cited here are one-based: the selected First Book scan uses printed page +16 and the Exercises scan +8; Wright uses +20 and Fulk +12.
% mdparagraph:P360:end

\endgroup

% book-source-body:end
\par\label{end-ch08-learner}
\clearpage\endgroup

\BookRecto
\begingroup
\renewcommand{\runningtitle}{Unit 9. Complements and reports}
\hyphenpenalty=10000
\exhyphenpenalty=10000
\emergencystretch=3em
\newcommand{\itempara}[2]{\par\addvspace{3pt}\noindent\hangindent=1.8em\hangafter=1\makebox[1.8em][l]{#1.}#2\par}
\newcommand{\reserve}[1]{\par\penalty-100\begingroup\dimen0=\pagegoal\advance\dimen0 by-\pagetotal\ifdim\dimen0<#1\newpage\fi\endgroup}
% book-source-body:begin

\addfontfeatures{Ligatures=TeXOff}
\raggedright
\booktitle{\BookComponent{ch09-learner}{Unit 9. Complements and reported content}UNIT 9}{Complements and reported content}

% mdparagraph:P001:begin
A clause can state what someone knows, report what someone says, or ask a question within another sentence. An infinitive can also describe an action that someone commands or observes. This unit distinguishes those constructions and adds the vocabulary of occupations used in Cook II.9–12.
% mdparagraph:P001:end

% mdparagraph:P002:begin
The examples and “The fisherman and his companion” are new teaching material. “The Choice of Occupations” is an exact source dialogue with its own help and original notes.
% mdparagraph:P002:end

% mdparagraph:P003:begin
In the tables, 1, 2 and 3 mean first, second and third person. Singular and plural cells are shown even when their spellings coincide. Dual subjects take plural finite forms. Noun and adjective tables give nominative, genitive, dative and accusative. Forms labelled \textbf{reference} are for recognition. Keep the written vowel lengths and the distinction between e and ę.
% mdparagraph:P003:end

\reserve{7\baselineskip}
% mdparagraph:P004:begin
\section*{1. Finite content clauses}\BookHeading{nav-ch09-learner-0001}{subsection}{1. Finite content clauses}{1}\BookLabel{bk-ch09-learner-a7e68e159a4d}
% mdparagraph:P004:end

% teaching: NEW-U09-GRAM-u09-content-that
% teaching: NEW-U09-VOC-thaet-conjunction
% teaching: NEW-U09-VOC-witan
% teaching: NEW-U09-VOC-secgan
% mdparagraph:P005:begin
A \textbf{content clause} states the content of knowing or saying. After the selected uses of \textbf{witan}, “know”, and \textbf{sęcgan}, “say, tell”, it may begin with invariant \textbf{ðæt}, “that”. The clause has its own subject and finite verb.
% mdparagraph:P005:end

% mdparagraph:P006:begin
This conjunction has no case, gender or number. Distinguish it from the inflected \textbf{ðæt} used as a demonstrative or article and from relative \textbf{ðæt}. A relative form has a function inside its clause; the content conjunction introduces the clause without being its subject or object.
% mdparagraph:P006:end

\reserve{8\baselineskip}
% example: NEW-U09-EXAMPLE-01
% mdparagraph:P007:begin
Ic wāt, ðæt se fisc gōd is.
% mdparagraph:P007:end

% mdparagraph:P008:begin
I know that the fish is good.
% mdparagraph:P008:end

% mdparagraph:P009:begin
The speaker presents the fish’s quality as a fact. Se fisc is the subject of is; the whole ðæt clause gives what the speaker knows.
% mdparagraph:P009:end

% teaching: NEW-U09-GRAM-u09-reported-statement
% teaching: NEW-U09-GRAM-u09-report-reference
% mdparagraph:P010:begin
Reported content commonly takes the subjunctive. The narrator presents someone else’s assertion without separately asserting its truth. The subjunctive does not mean that the assertion is false. Indicative content clauses remain possible when the proposition is presented as fact; an exercise that requires a particular mood will state the intended setting.
% mdparagraph:P010:end

\reserve{8\baselineskip}
% example: NEW-U09-EXAMPLE-02
% mdparagraph:P011:begin
Hē sægð, ðæt se fisc gōd sīe.
% mdparagraph:P011:end

% mdparagraph:P012:begin
He says that the fish is good.
% mdparagraph:P012:end

% mdparagraph:P013:begin
Sægð reports the present statement; sīe is present subjunctive. English commonly uses “is” in both this translation and the preceding one.
% mdparagraph:P013:end

\reserve{8\baselineskip}
% example: NEW-U09-EXAMPLE-03
% mdparagraph:P014:begin
Hē sægde, ðæt se fisc gōd wǣre.
% mdparagraph:P014:end

% mdparagraph:P015:begin
He said that the fish was good.
% mdparagraph:P015:end

% mdparagraph:P016:begin
Sægde places the speech in the past. Wǣre is preterite subjunctive: here the fish’s quality belongs to the time of that past speech.
% mdparagraph:P016:end

% mdparagraph:P017:begin
Identify\BookPoint{bk-ch09-learner-dc85c3b339be} who each pronoun denotes when converting direct speech into a report. The original speaker’s “I” may become “he” if another person reports the words; it need not change if the same person is still speaking. Choose person and tense from the actual setting. The present and past examples above give a useful contrast for simultaneous reported content, not an automatic rule that every tense moves backward after a past verb.
% mdparagraph:P017:end

% mdparagraph:P018:begin
The\BookPoint{bk-ch09-learner-8c696664c000} finite verb often comes later, sometimes finally, in a dependent clause. The examples use that common order. Case endings and context remain decisive when a source uses another arrangement.
% mdparagraph:P018:end

\reserve{7\baselineskip}
% mdparagraph:P019:begin
\section*{2. Dependent questions}\BookHeading{nav-ch09-learner-0002}{subsection}{2. Dependent questions}{1}\BookLabel{bk-ch09-learner-6a9c4e61c4a7}
% mdparagraph:P019:end

% teaching: NEW-U09-GRAM-u09-dependent-question
% mdparagraph:P020:begin
A \textbf{dependent question} gives the question whose answer someone knows or requests. It can follow \textbf{witan} or \textbf{sęcgan}. A question word such as \textbf{hwā}, \textbf{hwæt} or \textbf{hwǣm} introduces it.
% mdparagraph:P020:end

% mdparagraph:P021:begin
Choose the question word’s case from its function within the dependent clause: nominative subject, accusative object, dative recipient or other governed complement, and genitive possessor. Its case does not come from the verb of the main clause. A verb such as helpan still requires its dative complement inside a dependent question.
% mdparagraph:P021:end

\reserve{8\baselineskip}
% example: NEW-U09-EXAMPLE-04
% mdparagraph:P022:begin
Ic wāt, hwæt ðū hæbbe.
% mdparagraph:P022:end

% mdparagraph:P023:begin
I know what you have.
% mdparagraph:P023:end

% mdparagraph:P024:begin
Hwæt is the accusative object of hæbbe. Ðū is the subject, so the verb is second-person singular subjunctive even though its spelling matches other singular subjunctive forms.
% mdparagraph:P024:end

% teaching: NEW-U09-GRAM-u09-dependent-whether
% teaching: NEW-U09-GRAM-u09-dative-speaking-recipient
% teaching: NEW-U09-VOC-hwaether-conjunction
% mdparagraph:P025:begin
Invariant \textbf{hwæðer}, “whether”, introduces a dependent yes/\allowbreak{}no question. It differs from the earlier inflected adjective \textbf{hwæðer}, “which of two?”. The conjunction does not agree with a following noun.
% mdparagraph:P025:end

% mdparagraph:P026:begin
With \textbf{sęcgan}, the person addressed is dative. \textbf{Saga mē}, “tell me”, can introduce a question about what happened or whether something is so. This dative addressee differs from the accusative actor in the command construction below.
% mdparagraph:P026:end

\reserve{8\baselineskip}
% example: NEW-U09-EXAMPLE-05
% mdparagraph:P027:begin
Saga mē, hwæðer ðā męnn gearwe sīen.
% mdparagraph:P027:end

% mdparagraph:P028:begin
Tell me whether the men are ready.
% mdparagraph:P028:end

% mdparagraph:P029:begin
Mē names the person to be told. Hwæðer introduces the whole question; ðā męnn is the plural subject of sīen.
% mdparagraph:P029:end

% mdparagraph:P030:begin
Subjunctive is a normal pattern in dependent questions, and is the selected form when the exercises request it. An indicative may be appropriate when the dependent question concerns an asserted fact. The distinction between a dependent question and an ordinary relative clause is one of meaning and construction, not merely the position of a verb.
% mdparagraph:P030:end

\reserve{7\baselineskip}
% mdparagraph:P031:begin
\section*{3. Infinitive complements and object predicates}\BookHeading{nav-ch09-learner-0003}{subsection}{3. Infinitive complements and object predicates}{1}
% mdparagraph:P031:end

\reserve{7\baselineskip}
% mdparagraph:P032:begin
\subsection*{Commands and their actors}\BookHeading{nav-ch09-learner-0004}{subsubsection}{Commands and their actors}{0}\BookLabel{bk-ch09-learner-637a34313b91}
% mdparagraph:P032:end

% teaching: NEW-U09-GRAM-u09-command-infinitive
% teaching: NEW-U09-VOC-hatan
% mdparagraph:P033:begin
\textbf{Hātan}, “command”, takes an \textbf{accusative actor + plain infinitive}. The actor is the person or thing commanded to perform the infinitive’s action. The infinitive has no person ending and no preceding tō. Its own object retains the case required by its verb.
% mdparagraph:P033:end

\reserve{8\baselineskip}
% example: NEW-U09-EXAMPLE-06
% mdparagraph:P034:begin
Se mann hēt hine ðone fisc beran.
% mdparagraph:P034:end

% mdparagraph:P035:begin
The man ordered him to carry the fish.
% mdparagraph:P035:end

% mdparagraph:P036:begin
Se mann is the subject of hēt. Hine is the commanded actor of beran; ðone fisc is the object carried. Two accusatives need not have the same function.
% mdparagraph:P036:end

\reserve{8\baselineskip}
% example: NEW-U09-EXAMPLE-10
% mdparagraph:P037:begin
Ic wille ðone fisc beran.
% mdparagraph:P037:end

% mdparagraph:P038:begin
I want to carry the fish.
% mdparagraph:P038:end

% mdparagraph:P039:begin
Here ic is both the subject of wille and the understood actor of beran. The earlier same-subject modal pattern does not need a separate accusative actor.
% mdparagraph:P039:end

% mdparagraph:P040:begin
The same finite form of hātan can also mean “call, name”. Compare the complement that follows it. \textbf{Cyning} means “king”; its complete noun pattern is given below.
% mdparagraph:P040:end

\reserve{8\baselineskip}
% example: NEW-U09-EXAMPLE-11
% mdparagraph:P041:begin
Hē hēt ðone mann cyning.
% mdparagraph:P041:end

% mdparagraph:P042:begin
He called the man king.
% mdparagraph:P042:end

% mdparagraph:P043:begin
Cyning is an accusative noun predicate identifying ðone mann. This is the naming construction from Unit 7.
% mdparagraph:P043:end

\reserve{8\baselineskip}
% example: NEW-U09-EXAMPLE-12
% mdparagraph:P044:begin
Hē hēt ðone mann gān.
% mdparagraph:P044:end

% mdparagraph:P045:begin
He ordered the man to go.
% mdparagraph:P045:end

% mdparagraph:P046:begin
Gān is a plain infinitive; ðone mann is its commanded actor. The infinitive marks the command construction.
% mdparagraph:P046:end

\reserve{7\baselineskip}
% mdparagraph:P047:begin
\subsection*{Perceived actions}\BookHeading{nav-ch09-learner-0005}{subsubsection}{Perceived actions}{0}\BookLabel{bk-ch09-learner-22cf9e7f4d0c}
% mdparagraph:P047:end

% teaching: NEW-U09-GRAM-u09-perception-infinitive
% teaching: NEW-U09-VOC-geseon
% mdparagraph:P048:begin
\textbf{Gesēon}, “see”, can take an \textbf{accusative actor + plain infinitive} describing the action perceived. The construction says that the observer sees the actor performing that action.
% mdparagraph:P048:end

\reserve{8\baselineskip}
% example: NEW-U09-EXAMPLE-07
% mdparagraph:P049:begin
Ic gesēo ðone mann ðone fisc beran.
% mdparagraph:P049:end

% mdparagraph:P050:begin
I see the man carrying the fish.
% mdparagraph:P050:end

% mdparagraph:P051:begin
Ðone mann is the actor of beran; ðone fisc is its object. English “carrying” translates the Old English infinitive here. It does not make beran a participle.
% mdparagraph:P051:end

\reserve{7\baselineskip}
% mdparagraph:P052:begin
\subsection*{A quality predicated of the object}\BookHeading{nav-ch09-learner-0006}{subsubsection}{A quality predicated of the object}{0}\BookLabel{bk-ch09-learner-33765ec333c9}
% mdparagraph:P052:end

% teaching: NEW-U09-GRAM-u09-object-predicate
% teaching: NEW-U09-VOC-gewyrcean
% mdparagraph:P053:begin
With \textbf{gewyrcean}, “make”, an accusative object may be followed by an adjective predicated of that object. The adjective agrees with the object in case, gender and number. It is strong in this predicate use; an article belonging to the object noun does not make the separate predicate adjective weak.
% mdparagraph:P053:end

\reserve{8\baselineskip}
% example: NEW-U09-EXAMPLE-08
% mdparagraph:P054:begin
God geworhte ðone mann gōdne.
% mdparagraph:P054:end

% mdparagraph:P055:begin
God made the man good.
% mdparagraph:P055:end

% mdparagraph:P056:begin
Gōdne is masculine accusative singular because it describes ðone mann. Its form is not controlled by nominative God.
% mdparagraph:P056:end

\reserve{7\baselineskip}
% mdparagraph:P057:begin
\section*{4. Wishes, preference and related expressions}\BookHeading{nav-ch09-learner-0007}{subsection}{4. Wishes, preference and related expressions}{1}\BookLabel{bk-ch09-learner-84672e8deefc}
% mdparagraph:P057:end

% teaching: NEW-U09-GRAM-u09-unreal-wish
% mdparagraph:P058:begin
The preterite subjunctive can express an unreal wish about the present. The past form need not refer to past time. In such a wish, \textbf{wǣre} means “were” in “would that I were”, and \textbf{hæfde} can mean “had” in “would that I had”. The verb commonly stands first.
% mdparagraph:P058:end

\reserve{8\baselineskip}
% example: NEW-U09-EXAMPLE-09
% mdparagraph:P059:begin
Wǣre ic cyning!
% mdparagraph:P059:end

% mdparagraph:P060:begin
Would that I were king!
% mdparagraph:P060:end

% mdparagraph:P061:begin
The speaker is not king now. This is an unreal present wish, not a statement about a former reign.
% mdparagraph:P061:end

% mdparagraph:P062:begin
Some indicative and subjunctive forms are identical. In particular, first-person singular hæfde can state actual past possession or express an unreal wish. The construction and the stated situation distinguish them; its spelling alone cannot.
% mdparagraph:P062:end

% mdparagraph:P063:begin
The earlier present-subjunctive wishes and exhortations remain available, including first-person plural “let us …” and third-person “let …”. These do not require the English verb “let” to be translated by an additional Old English verb.
% mdparagraph:P063:end

% teaching: NEW-U09-GRAM-u09-volitional-refusal
% teaching: NEW-U09-VOC-willan
% mdparagraph:P064:begin
\textbf{Ne + willan}\BookPoint{bk-ch09-learner-f4a46510a7af} can express unwillingness or refusal: “does not wish to”, “will not”. Use the familiar positive forms of willan with ne before the finite verb. Here English “will not” can concern someone’s willingness; it need not predict the future. No new contracted negative forms are needed in these exercises.
% mdparagraph:P064:end

% teaching: NEW-U09-GRAM-u09-infinitive-subject
% teaching: NEW-U09-GRAM-u09-mid-accompaniment
% teaching: NEW-U09-VOC-leof-comp
% teaching: NEW-U09-VOC-mid
% mdparagraph:P065:begin
The\BookPoint{bk-ch09-learner-57a9fc624ee1} preference expression \textbf{ūs is lēofre wician …}, “it is preferable to us to sojourn …”, has a dative experiencer, \textbf{ūs}, and the neuter comparative \textbf{lēofre}. The infinitive names what is preferable. \textbf{Wician} means “sojourn, lodge”. English “we had rather” can express this present preference; its word “had” does not require a past verb.
% mdparagraph:P065:end

% mdparagraph:P066:begin
\textbf{Mid + dative}\BookPoint{bk-ch09-learner-e3fb5ca9be3f} names an accompanying person or cargo. In the preference construction, compare parallel mid phrases on either side of \textbf{ðonne}, “than”: staying with one person is compared with staying with another. Use the dative after each mid. The accompaniment use also covers going aboard with one’s wares.
% mdparagraph:P066:end

% teaching: NEW-U09-GRAM-u09-free-relative-that
% teaching: NEW-U09-VOC-se
% mdparagraph:P067:begin
In\BookPoint{bk-ch09-learner-5ae9fa426c67} the expression “be what you ought to be”, \textbf{ðæt ðe} means “what, that which”, without an expressed noun antecedent. The expression identifies a role or condition and functions as a predicate with “be”. It is distinct from the statement conjunction ðæt.
% mdparagraph:P067:end

% mdparagraph:P068:begin
The familiar \textbf{sculan + plain infinitive} expresses obligation, including “ought to be”. Keep the person and number of sculan appropriate to its subject. This unit uses the predicative ðæt ðe expression in those be/\allowbreak{}ought-to-be clauses.
% mdparagraph:P068:end

% teaching: NEW-U09-GRAM-u09-partitive-pronoun
% teaching: NEW-U09-VOC-hwilc
% teaching: NEW-U09-VOC-sum
% teaching: NEW-U09-VOC-thu
% mdparagraph:P069:begin
A\BookPoint{bk-ch09-learner-555d1541a4a7} personal genitive can identify the group from which a member is selected: \textbf{hwilc ēower}, “which of you”, and \textbf{sume hiera}, “some of them”. Here \textbf{ēower} is the genitive plural personal pronoun, not an agreeing possessive adjective. The word selecting from the group still has the case required by its own clause.
% mdparagraph:P069:end

% mdparagraph:P070:begin
Cook\BookPoint{bk-ch09-learner-07e15514cef1} also gives \textbf{hīe sume}, “some of them”. In this expression hīe and sume are both nominative plural, standing in apposition. Do not parse hīe as genitive. The appositional expression and the genitive construction are two ways to express the selection; Cook’s reference to §151 asks for the former.
% mdparagraph:P070:end

\reserve{7\baselineskip}
% mdparagraph:P071:begin
\section*{5. Nouns: complete patterns and vocabulary}\BookHeading{nav-ch09-learner-0008}{subsection}{5. Nouns: complete patterns and vocabulary}{1}
% mdparagraph:P071:end

% mdparagraph:P072:begin
Each noun is given with its gender. The tables show complete singular and plural patterns. In the vocabulary tables, \textbf{stem} means the base to which the model’s endings are added. Use the headword itself for nominative singular. In the strong masculine and neuter models below, accusative singular is also unchanged; weak masculine and feminine nouns instead follow their model’s -an. The listed genitive singular and nominative plural verify the stem and principal endings.
% mdparagraph:P072:end

% mdparagraph:P073:begin
Each vocabulary group names its full model. Replace that model’s stem by the listed stem, retaining its endings. Do not remove an internal vowel or change a root vowel unless the diagnostic forms show the change. The three new patterns—sunu, rīce and geðeahtend—are given separately in full.
% mdparagraph:P073:end

% mdparagraph:P074:begin
In meanings marked \textbf{singular}, use the noun in that stated material, abstract or collective sense in the singular. Plural counterparts remain reference forms; they do not automatically mean kinds, portions or separate instances. Meanings marked \textbf{plural} use the plural, with singular counterparts for reference. Unmarked nouns are used in both numbers.
% mdparagraph:P074:end

\reserve{7\baselineskip}
% mdparagraph:P075:begin
\subsection*{Strong masculines}\BookHeading{nav-ch09-learner-0009}{subsubsection}{Strong masculines}{0}
% mdparagraph:P075:end

% mdparagraph:P076:begin
The regular masculine has genitive singular -es and dative -e. Its plural has -as, -a, -um, -as. Cræft keeps æ before its consonant cluster; do not substitute a in its plural.
% mdparagraph:P076:end

% teaching: NEW-U09-NOM-SMITH
\reserve{7\baselineskip}
% mdparagraph:P077:begin
\subsection*{smið}\BookHeading{nav-ch09-learner-0010}{subsubsection}{smið}{0}
% mdparagraph:P077:end

% table: NEW-U09-NOM-SMITH
% mdtable:T01:begin
\par\begingroup
\fontsize{10.5}{13.2}\selectfont
\setlength{\LTleft}{0pt}\setlength{\LTright}{\fill}
\noindent
\begin{tabular}{@{}>{\raggedright\arraybackslash}p{39.46mm}>{\raggedright\arraybackslash}p{39.46mm}>{\raggedright\arraybackslash}p{39.46mm}@{}}
\toprule
\textbf{Case} & \textbf{Singular} & \textbf{Plural} \\
\midrule
Nominative & smið & smiðas \\
Genitive & smiðes & smiða \\
Dative & smiðe & smiðum \\
Accusative & smið & smiðas \\
\bottomrule
\end{tabular}
\par\endgroup

% mdtable:T01:end
% teaching: NEW-U09-NOM-SULHSCEAR
% teaching: NEW-U09-VOC-sulhscear
% teaching: NEW-U09-VOC-smith
% teaching: NEW-U09-NOM-BELG
% teaching: NEW-U09-VOC-belg
% teaching: NEW-U09-NOM-IERTHLING
% teaching: NEW-U09-VOC-ierthling
% teaching: NEW-U09-NOM-MAESSEPREOST
% teaching: NEW-U09-VOC-maessepreost
% teaching: NEW-U09-NOM-CRAEFT
% teaching: NEW-U09-VOC-craeft
% teaching: NEW-U09-NOM-THEOWDOM
% teaching: NEW-U09-VOC-theowdom
% teaching: NEW-U09-NOM-ELP
% teaching: NEW-U09-VOC-elp
% teaching: NEW-U09-NOM-CYNING
% teaching: NEW-U09-VOC-cyning
% teaching: NEW-U09-NOM-MUNUC
% teaching: NEW-U09-VOC-munuc
% teaching: NEW-U09-NOM-DRENC
% teaching: NEW-U09-VOC-drenc
% teaching: NEW-U09-NOM-CEORL
% teaching: NEW-U09-VOC-ceorl
% teaching: NEW-U09-NOM-PLIHT
% teaching: NEW-U09-VOC-pliht
% teaching: NEW-U09-NOM-BEORN
% teaching: NEW-U09-VOC-beorn
% mdparagraph:P078:begin
Vocabulary following \textbf{smið}:
% mdparagraph:P078:end

% mdtable:T02:begin
\par\begingroup
\fontsize{10.5}{13.2}\selectfont
\setlength{\LTleft}{0pt}\setlength{\LTright}{\fill}
\noindent
\begin{tabular}{@{}>{\raggedright\arraybackslash}p{21.3mm}>{\raggedright\arraybackslash}p{20mm}>{\raggedright\arraybackslash}p{22.3mm}>{\raggedright\arraybackslash}p{24.6mm}>{\raggedright\arraybackslash}p{24.55mm}@{}}
\toprule
\textbf{Noun} & \textbf{Gender; meaning} & \textbf{Stem} & \textbf{Genitive singular} & \textbf{Nominative plural} \\
\midrule
sulhscear & m.; plowshare & sulhscear- & sulhsceares & sulhscearas \\
smið & m.; blacksmith & smið- & smiðes & smiðas \\
bęlg & m.; bellows (plural) & bęlg- & bęlges & bęlgas \\
ierðling & m.; farmer; plowman & ierðling- & ierðlinges & ierðlingas \\
mæsseprēost & m.; priest & mæsseprēost- & mæsseprēostes & mæsseprēostas \\
cræft & m.; trade; occupation; craft & cræft- & cræftes & cræftas \\
ðēowdōm & m.; service (singular) & ðēowdōm- & ðēowdōmes & ðēowdōmas \\
elp & m.; elephant & elp- & elpes & elpas \\
cyning & m.; king & cyning- & cyninges & cyningas \\
munuc & m.; monk & munuc- & munuces & munucas \\
dręnc & m.; drink (singular) & dręnc- & dręnces & dręncas \\
ceorl & m.; layman & ceorl- & ceorles & ceorlas \\
pliht & m.; peril; risk & pliht- & plihtes & plihtas \\
beorn & m.; warrior & beorn- & beornes & beornas \\
\bottomrule
\end{tabular}
\par\endgroup

% mdtable:T02:end
\reserve{7\baselineskip}
% mdparagraph:P079:begin
\subsection*{Masculines in -e}\BookHeading{nav-ch09-learner-0011}{subsubsection}{Masculines in -e}{0}
% mdparagraph:P079:end

% mdparagraph:P080:begin
The headword’s final -e remains in nominative and accusative singular. Before another ending use the listed stem: fiscer- + -es gives fisceres. Dative singular ends in -e; the plural has the ordinary masculine endings.
% mdparagraph:P080:end

% teaching: NEW-U09-NOM-FISCERE
\reserve{7\baselineskip}
% mdparagraph:P081:begin
\subsection*{fiscere}\BookHeading{nav-ch09-learner-0012}{subsubsection}{fiscere}{0}
% mdparagraph:P081:end

% table: NEW-U09-NOM-FISCERE
% mdtable:T03:begin
\par\begingroup
\fontsize{10.5}{13.2}\selectfont
\setlength{\LTleft}{0pt}\setlength{\LTright}{\fill}
\noindent
\begin{tabular}{@{}>{\raggedright\arraybackslash}p{39.46mm}>{\raggedright\arraybackslash}p{39.46mm}>{\raggedright\arraybackslash}p{39.46mm}@{}}
\toprule
\textbf{Case} & \textbf{Singular} & \textbf{Plural} \\
\midrule
Nominative & fiscere & fisceras \\
Genitive & fisceres & fiscera \\
Dative & fiscere & fiscerum \\
Accusative & fiscere & fisceras \\
\bottomrule
\end{tabular}
\par\endgroup

% mdtable:T03:end
% teaching: NEW-U09-NOM-MANGERE
% teaching: NEW-U09-VOC-mangere
% teaching: NEW-U09-NOM-ELE
% teaching: NEW-U09-VOC-ele
% teaching: NEW-U09-VOC-fiscere
% teaching: NEW-U09-NOM-SEAMERE
% teaching: NEW-U09-VOC-seamere
% teaching: NEW-U09-NOM-EALDORSCIPE
% teaching: NEW-U09-VOC-ealdorscipe
% teaching: NEW-U09-NOM-LYRE
% teaching: NEW-U09-VOC-lyre
% mdparagraph:P082:begin
Vocabulary following \textbf{fiscere}:
% mdparagraph:P082:end

% mdtable:T04:begin
\par\begingroup
\fontsize{10.5}{13.2}\selectfont
\setlength{\LTleft}{0pt}\setlength{\LTright}{\fill}
\noindent
\begin{tabular}{@{}>{\raggedright\arraybackslash}p{19.5mm}>{\raggedright\arraybackslash}p{29.97mm}>{\raggedright\arraybackslash}p{18.5mm}>{\raggedright\arraybackslash}p{22.39mm}>{\raggedright\arraybackslash}p{22.39mm}@{}}
\toprule
\textbf{Noun} & \textbf{Gender; meaning} & \textbf{Stem} & \textbf{Genitive singular} & \textbf{Nominative plural} \\
\midrule
mangere & m.; merchant & manger- & mangeres & mangeras \\
ęle & m.; oil (singular) & ęl- & ęles & ęlas \\
fiscere & m.; fisherman & fiscer- & fisceres & fisceras \\
sēamere & m.; tailor & sēamer- & sēameres & sēameras \\
ealdorscipe & m.; chief place; primacy (singular) & ealdorscip- & ealdorscipes & ealdorscipas \\
lyre & m.; loss (singular) & lyr- & lyres & lyras \\
\bottomrule
\end{tabular}
\par\endgroup

% mdtable:T04:end
\reserve{7\baselineskip}
% mdparagraph:P083:begin
\subsection*{A masculine with a shortened stem}\BookHeading{nav-ch09-learner-0013}{subsubsection}{A masculine with a shortened stem}{0}
% mdparagraph:P083:end

% mdparagraph:P084:begin
Angel loses its internal e before endings: use angl-. The headword remains angel in nominative and accusative singular.
% mdparagraph:P084:end

% teaching: NEW-U09-NOM-ANGEL
\reserve{7\baselineskip}
% mdparagraph:P085:begin
\subsection*{angel}\BookHeading{nav-ch09-learner-0014}{subsubsection}{angel}{0}
% mdparagraph:P085:end

% table: NEW-U09-NOM-ANGEL
% mdtable:T05:begin
\par\begingroup
\fontsize{10.5}{13.2}\selectfont
\setlength{\LTleft}{0pt}\setlength{\LTright}{\fill}
\noindent
\begin{tabular}{@{}>{\raggedright\arraybackslash}p{39.46mm}>{\raggedright\arraybackslash}p{39.46mm}>{\raggedright\arraybackslash}p{39.46mm}@{}}
\toprule
\textbf{Case} & \textbf{Singular} & \textbf{Plural} \\
\midrule
Nominative & angel & anglas \\
Genitive & angles & angla \\
Dative & angle & anglum \\
Accusative & angel & anglas \\
\bottomrule
\end{tabular}
\par\endgroup

% mdtable:T05:end
% teaching: NEW-U09-VOC-angel
% mdparagraph:P086:begin
Vocabulary following \textbf{angel}:
% mdparagraph:P086:end

% mdtable:T06:begin
\par\begingroup
\fontsize{10.5}{13.2}\selectfont
\setlength{\LTleft}{0pt}\setlength{\LTright}{\fill}
\noindent
\begin{tabular}{@{}>{\raggedright\arraybackslash}p{19.82mm}>{\raggedright\arraybackslash}p{23.18mm}>{\raggedright\arraybackslash}p{19.82mm}>{\raggedright\arraybackslash}p{23.88mm}>{\raggedright\arraybackslash}p{26.06mm}@{}}
\toprule
\textbf{Noun} & \textbf{Gender; meaning} & \textbf{Stem} & \textbf{Genitive singular} & \textbf{Nominative plural} \\
\midrule
angel & m.; fish-hook & angl- & angles & anglas \\
\bottomrule
\end{tabular}
\par\endgroup

% mdtable:T06:end
\reserve{7\baselineskip}
% mdparagraph:P087:begin
\subsection*{Strong neuters with an endingless plural}\BookHeading{nav-ch09-learner-0015}{subsubsection}{Strong neuters with an endingless plural}{0}
% mdparagraph:P087:end

% mdparagraph:P088:begin
These neuters have the usual -es and -e in genitive and dative singular. Nominative and accusative plural have no ending; genitive and dative plural have -a and -um.
% mdparagraph:P088:end

% teaching: NEW-U09-NOM-HUS
\reserve{7\baselineskip}
% mdparagraph:P089:begin
\subsection*{hūs}\BookHeading{nav-ch09-learner-0016}{subsubsection}{hūs}{0}
% mdparagraph:P089:end

% table: NEW-U09-NOM-HUS
% mdtable:T07:begin
\par\begingroup
\fontsize{10.5}{13.2}\selectfont
\setlength{\LTleft}{0pt}\setlength{\LTright}{\fill}
\noindent
\begin{tabular}{@{}>{\raggedright\arraybackslash}p{39.46mm}>{\raggedright\arraybackslash}p{39.46mm}>{\raggedright\arraybackslash}p{39.46mm}@{}}
\toprule
\textbf{Case} & \textbf{Singular} & \textbf{Plural} \\
\midrule
Nominative & hūs & hūs \\
Genitive & hūses & hūsa \\
Dative & hūse & hūsum \\
Accusative & hūs & hūs \\
\bottomrule
\end{tabular}
\par\endgroup

% mdtable:T07:end
% teaching: NEW-U09-VOC-hus
% teaching: NEW-U09-NOM-WIN
% teaching: NEW-U09-VOC-win
% teaching: NEW-U09-NOM-GOLD
% teaching: NEW-U09-VOC-gold
% teaching: NEW-U09-NOM-HORS
% teaching: NEW-U09-VOC-hors
% teaching: NEW-U09-NOM-GODSPELL
% teaching: NEW-U09-VOC-godspell
% teaching: NEW-U09-NOM-GESTREON
% teaching: NEW-U09-VOC-gestreon
% teaching: NEW-U09-NOM-GEDEORF
% teaching: NEW-U09-VOC-gedeorf
% teaching: NEW-U09-NOM-BAN
% teaching: NEW-U09-VOC-ban
% teaching: NEW-U09-NOM-FOLC
% teaching: NEW-U09-VOC-folc
% teaching: NEW-U09-NOM-GEFLIT
% teaching: NEW-U09-VOC-geflit
% teaching: NEW-U09-NOM-WIF
% teaching: NEW-U09-VOC-wif
% teaching: NEW-U09-NOM-HLAEST
% teaching: NEW-U09-VOC-hlaest
% teaching: NEW-U09-NOM-AER-COPPER
% teaching: NEW-U09-VOC-aer-copper
% mdparagraph:P090:begin
Vocabulary following \textbf{hūs}:
% mdparagraph:P090:end

% mdtable:T08:begin
\par\begingroup
\fontsize{10.5}{13.2}\selectfont
\setlength{\LTleft}{0pt}\setlength{\LTright}{\fill}
\noindent
\begin{tabular}{@{}>{\raggedright\arraybackslash}p{16.11mm}>{\raggedright\arraybackslash}p{38.17mm}>{\raggedright\arraybackslash}p{16.11mm}>{\raggedright\arraybackslash}p{21.18mm}>{\raggedright\arraybackslash}p{21.18mm}@{}}
\toprule
\textbf{Noun} & \textbf{Gender; meaning} & \textbf{Stem} & \textbf{Genitive singular} & \textbf{Nominative plural} \\
\midrule
hūs & n.; house & hūs- & hūses & hūs \\
wīn & n.; wine (singular) & wīn- & wīnes & wīn \\
gold & n.; gold (singular) & gold- & goldes & gold \\
hors & n.; horse & hors- & horses & hors \\
godspell & n.; Gospel (singular) & godspell- & godspelles & godspell \\
gestrēon & n.; gain; profit (singular) & gestrēon- & gestrēones & gestrēon \\
gedeorf & n.; labor; toil (singular) & gedeorf- & gedeorfes & gedeorf \\
bān & n.; bone & bān- & bānes & bān \\
folc & n.; people, as a collective (singular) & folc- & folces & folc \\
geflīt & n.; dispute; strife & geflīt- & geflītes & geflīt \\
wīf & n.; wife & wīf- & wīfes & wīf \\
hlæst & n.; wares; merchandise; cargo (plural) & hlæst- & hlæstes & hlæst \\
ǣr & n.; copper (singular) & ǣr- & ǣres & ǣr \\
\bottomrule
\end{tabular}
\par\endgroup

% mdtable:T08:end
\reserve{7\baselineskip}
% mdparagraph:P091:begin
\subsection*{Short-stem neuters}\BookHeading{nav-ch09-learner-0017}{subsubsection}{Short-stem neuters}{0}
% mdparagraph:P091:end

% mdparagraph:P092:begin
The short stem takes -u in nominative and accusative plural. Tin has a short vowel. In this unit the material “tin” is used in the singular; its plural table is a reference model.
% mdparagraph:P092:end

% teaching: NEW-U09-NOM-TIN
\reserve{7\baselineskip}
% mdparagraph:P093:begin
\subsection*{tin}\BookHeading{nav-ch09-learner-0018}{subsubsection}{tin}{0}
% mdparagraph:P093:end

% table: NEW-U09-NOM-TIN
% mdtable:T09:begin
\par\begingroup
\fontsize{10.5}{13.2}\selectfont
\setlength{\LTleft}{0pt}\setlength{\LTright}{\fill}
\noindent
\begin{tabular}{@{}>{\raggedright\arraybackslash}p{39.46mm}>{\raggedright\arraybackslash}p{39.46mm}>{\raggedright\arraybackslash}p{39.46mm}@{}}
\toprule
\textbf{Case} & \textbf{Singular} & \textbf{Plural} \\
\midrule
Nominative & tin & tinu \\
Genitive & tines & tina \\
Dative & tine & tinum \\
Accusative & tin & tinu \\
\bottomrule
\end{tabular}
\par\endgroup

% mdtable:T09:end
% teaching: NEW-U09-VOC-tin
% mdparagraph:P094:begin
Vocabulary following \textbf{tin}:
% mdparagraph:P094:end

% mdtable:T10:begin
\par\begingroup
\fontsize{10.5}{13.2}\selectfont
\setlength{\LTleft}{0pt}\setlength{\LTright}{\fill}
\noindent
\begin{tabular}{@{}>{\raggedright\arraybackslash}p{19.18mm}>{\raggedright\arraybackslash}p{26.08mm}>{\raggedright\arraybackslash}p{19.18mm}>{\raggedright\arraybackslash}p{23.11mm}>{\raggedright\arraybackslash}p{25.22mm}@{}}
\toprule
\textbf{Noun} & \textbf{Gender; meaning} & \textbf{Stem} & \textbf{Genitive singular} & \textbf{Nominative plural} \\
\midrule
tin & n.; tin (singular) & tin- & tines & tinu \\
\bottomrule
\end{tabular}
\par\endgroup

% mdtable:T10:end
\reserve{7\baselineskip}
% mdparagraph:P095:begin
\subsection*{A neuter retaining its internal vowel}\BookHeading{nav-ch09-learner-0019}{subsubsection}{A neuter retaining its internal vowel}{0}
% mdparagraph:P095:end

% mdparagraph:P096:begin
Ðyrel has short y and retains its internal e before endings. Its nominative and accusative plural end in -u.
% mdparagraph:P096:end

% teaching: NEW-U09-NOM-THYREL
\reserve{7\baselineskip}
% mdparagraph:P097:begin
\subsection*{ðyrel}\BookHeading{nav-ch09-learner-0020}{subsubsection}{ðyrel}{0}
% mdparagraph:P097:end

% table: NEW-U09-NOM-THYREL
% mdtable:T11:begin
\par\begingroup
\fontsize{10.5}{13.2}\selectfont
\setlength{\LTleft}{0pt}\setlength{\LTright}{\fill}
\noindent
\begin{tabular}{@{}>{\raggedright\arraybackslash}p{39.46mm}>{\raggedright\arraybackslash}p{39.46mm}>{\raggedright\arraybackslash}p{39.46mm}@{}}
\toprule
\textbf{Case} & \textbf{Singular} & \textbf{Plural} \\
\midrule
Nominative & ðyrel & ðyrelu \\
Genitive & ðyreles & ðyrela \\
Dative & ðyrele & ðyrelum \\
Accusative & ðyrel & ðyrelu \\
\bottomrule
\end{tabular}
\par\endgroup

% mdtable:T11:end
% teaching: NEW-U09-VOC-thyrel
% mdparagraph:P098:begin
Vocabulary following \textbf{ðyrel}:
% mdparagraph:P098:end

% mdtable:T12:begin
\par\begingroup
\fontsize{10.5}{13.2}\selectfont
\setlength{\LTleft}{0pt}\setlength{\LTright}{\fill}
\noindent
\begin{tabular}{@{}>{\raggedright\arraybackslash}p{20.21mm}>{\raggedright\arraybackslash}p{21.42mm}>{\raggedright\arraybackslash}p{20.21mm}>{\raggedright\arraybackslash}p{24.35mm}>{\raggedright\arraybackslash}p{26.57mm}@{}}
\toprule
\textbf{Noun} & \textbf{Gender; meaning} & \textbf{Stem} & \textbf{Genitive singular} & \textbf{Nominative plural} \\
\midrule
ðyrel & n.; hole & ðyrel- & ðyreles & ðyrelu \\
\bottomrule
\end{tabular}
\par\endgroup

% mdtable:T12:end
\reserve{7\baselineskip}
% mdparagraph:P099:begin
\subsection*{A neuter with a shortened stem}\BookHeading{nav-ch09-learner-0021}{subsubsection}{A neuter with a shortened stem}{0}
% mdparagraph:P099:end

% mdparagraph:P100:begin
Fōdor uses fōdr- before endings. Its ordinary sense here is singular “fodder”; the plural counterparts are reference forms.
% mdparagraph:P100:end

% teaching: NEW-U09-NOM-FODOR
\reserve{7\baselineskip}
% mdparagraph:P101:begin
\subsection*{fōdor}\BookHeading{nav-ch09-learner-0022}{subsubsection}{fōdor}{0}
% mdparagraph:P101:end

% table: NEW-U09-NOM-FODOR
% mdtable:T13:begin
\par\begingroup
\fontsize{10.5}{13.2}\selectfont
\setlength{\LTleft}{0pt}\setlength{\LTright}{\fill}
\noindent
\begin{tabular}{@{}>{\raggedright\arraybackslash}p{39.46mm}>{\raggedright\arraybackslash}p{39.46mm}>{\raggedright\arraybackslash}p{39.46mm}@{}}
\toprule
\textbf{Case} & \textbf{Singular} & \textbf{Plural} \\
\midrule
Nominative & fōdor & fōdru \\
Genitive & fōdres & fōdra \\
Dative & fōdre & fōdrum \\
Accusative & fōdor & fōdru \\
\bottomrule
\end{tabular}
\par\endgroup

% mdtable:T13:end
% teaching: NEW-U09-VOC-fodor
% mdparagraph:P102:begin
Vocabulary following \textbf{fōdor}:
% mdparagraph:P102:end

% mdtable:T14:begin
\par\begingroup
\fontsize{10.5}{13.2}\selectfont
\setlength{\LTleft}{0pt}\setlength{\LTright}{\fill}
\noindent
\begin{tabular}{@{}>{\raggedright\arraybackslash}p{18.75mm}>{\raggedright\arraybackslash}p{28.02mm}>{\raggedright\arraybackslash}p{18.75mm}>{\raggedright\arraybackslash}p{22.59mm}>{\raggedright\arraybackslash}p{24.65mm}@{}}
\toprule
\textbf{Noun} & \textbf{Gender; meaning} & \textbf{Stem} & \textbf{Genitive singular} & \textbf{Nominative plural} \\
\midrule
fōdor & n.; fodder (singular) & fōdr- & fōdres & fōdru \\
\bottomrule
\end{tabular}
\par\endgroup

% mdtable:T14:end
\reserve{7\baselineskip}
% mdparagraph:P103:begin
\subsection*{Strong feminines}\BookHeading{nav-ch09-learner-0023}{subsubsection}{Strong feminines}{0}
% mdparagraph:P103:end

% mdparagraph:P104:begin
The nominative singular is the headword. The other singular cases end in -e; the selected plural has -a, -a, -um, -a. Feminine nominative and accusative plural -e is a documented alternative to selected -a. Both ǣl “awl” and slęcg “hammer” are feminine.
% mdparagraph:P104:end

% teaching: NEW-U09-NOM-NAEDL
\reserve{7\baselineskip}
% mdparagraph:P105:begin
\subsection*{nǣdl}\BookHeading{nav-ch09-learner-0024}{subsubsection}{nǣdl}{0}
% mdparagraph:P105:end

% table: NEW-U09-NOM-NAEDL
% mdtable:T15:begin
\par\begingroup
\fontsize{10.5}{13.2}\selectfont
\setlength{\LTleft}{0pt}\setlength{\LTright}{\fill}
\noindent
\begin{tabular}{@{}>{\raggedright\arraybackslash}p{39.46mm}>{\raggedright\arraybackslash}p{39.46mm}>{\raggedright\arraybackslash}p{39.46mm}@{}}
\toprule
\textbf{Case} & \textbf{Singular} & \textbf{Plural} \\
\midrule
Nominative & nǣdl & nǣdla \\
Genitive & nǣdle & nǣdla \\
Dative & nǣdle & nǣdlum \\
Accusative & nǣdle & nǣdla \\
\bottomrule
\end{tabular}
\par\endgroup

% mdtable:T15:end
% teaching: NEW-U09-VOC-naedl
% teaching: NEW-U09-NOM-AEL
% teaching: NEW-U09-VOC-ael
% teaching: NEW-U09-NOM-SLECG
% teaching: NEW-U09-VOC-slecg
% teaching: NEW-U09-NOM-EORTHTILTH
% teaching: NEW-U09-VOC-eorthtilth
% mdparagraph:P106:begin
Vocabulary following \textbf{nǣdl}:
% mdparagraph:P106:end

% mdtable:T16:begin
\par\begingroup
\fontsize{10.5}{13.2}\selectfont
\setlength{\LTleft}{0pt}\setlength{\LTright}{\fill}
\noindent
\begin{tabular}{@{}>{\raggedright\arraybackslash}p{17.76mm}>{\raggedright\arraybackslash}p{31.52mm}>{\raggedright\arraybackslash}p{17.76mm}>{\raggedright\arraybackslash}p{22.37mm}>{\raggedright\arraybackslash}p{23.35mm}@{}}
\toprule
\textbf{Noun} & \textbf{Gender; meaning} & \textbf{Stem} & \textbf{Genitive singular} & \textbf{Nominative plural} \\
\midrule
nǣdl & f.; needle & nǣdl- & nǣdle & nǣdla \\
ǣl & f.; awl & ǣl- & ǣle & ǣla \\
slęcg & f.; hammer & slęcg- & slęcge & slęcga \\
eorðtilð & f.; agriculture (singular) & eorðtilð- & eorðtilðe & eorðtilða \\
\bottomrule
\end{tabular}
\par\endgroup

% mdtable:T16:end
\reserve{7\baselineskip}
% mdparagraph:P107:begin
\subsection*{Feminines in -nes}\BookHeading{nav-ch09-learner-0025}{subsubsection}{Feminines in -nes}{0}
% mdparagraph:P107:end

% mdparagraph:P108:begin
Keep the single s of the headword in nominative singular, but use the stem in -ness- before every ending. These nouns take the same case endings as nǣdl.
% mdparagraph:P108:end

% teaching: NEW-U09-NOM-FORLIDENNES
\reserve{7\baselineskip}
% mdparagraph:P109:begin
\subsection*{forlidennes}\BookHeading{nav-ch09-learner-0026}{subsubsection}{forlidennes}{0}
% mdparagraph:P109:end

% table: NEW-U09-NOM-FORLIDENNES
% mdtable:T17:begin
\par\begingroup
\fontsize{10.5}{13.2}\selectfont
\setlength{\LTleft}{0pt}\setlength{\LTright}{\fill}
\noindent
\begin{tabular}{@{}>{\raggedright\arraybackslash}p{31.99mm}>{\raggedright\arraybackslash}p{41.59mm}>{\raggedright\arraybackslash}p{44.79mm}@{}}
\toprule
\textbf{Case} & \textbf{Singular} & \textbf{Plural} \\
\midrule
Nominative & forlidennes & forlidennessa \\
Genitive & forlidennesse & forlidennessa \\
Dative & forlidennesse & forlidennessum \\
Accusative & forlidennesse & forlidennessa \\
\bottomrule
\end{tabular}
\par\endgroup

% mdtable:T17:end
% teaching: NEW-U09-VOC-forlidennes
% teaching: NEW-U09-NOM-RIHTWISNES
% teaching: NEW-U09-VOC-rihtwisnes
% mdparagraph:P110:begin
Vocabulary following \textbf{forlidennes}:
% mdparagraph:P110:end

% mdtable:T18:begin
\par\begingroup
\fontsize{10.5}{13.2}\selectfont
\setlength{\LTleft}{0pt}\setlength{\LTright}{\fill}
\noindent
\begin{tabular}{@{}>{\raggedright\arraybackslash}p{19mm}>{\raggedright\arraybackslash}p{25.5mm}>{\raggedright\arraybackslash}p{22mm}>{\raggedright\arraybackslash}p{23.125mm}>{\raggedright\arraybackslash}p{23.125mm}@{}}
\toprule
\textbf{Noun} & \textbf{Gender; meaning} & \textbf{Stem} & \textbf{Genitive singular} & \textbf{Nominative plural} \\
\midrule
forlidennes & f.; shipwreck & forlidenness- & forlidennesse & forlidennessa \\
rihtwīsnes & f.; righteousness (singular) & rihtwīsness- & rihtwīsnesse & rihtwīsnessa \\
\bottomrule
\end{tabular}
\par\endgroup

% mdtable:T18:end
\reserve{7\baselineskip}
% mdparagraph:P111:begin
\subsection*{Weak masculines}\BookHeading{nav-ch09-learner-0027}{subsubsection}{Weak masculines}{0}
% mdparagraph:P111:end

% mdparagraph:P112:begin
Weak masculine nominative singular ends in -a. The other singular cases and nominative and accusative plural end in -an; genitive plural has -ena, dative plural -um.
% mdparagraph:P112:end

% teaching: NEW-U09-NOM-FYRSPEARCA
\reserve{7\baselineskip}
% mdparagraph:P113:begin
\subsection*{fȳrspearca}\BookHeading{nav-ch09-learner-0028}{subsubsection}{fȳrspearca}{0}
% mdparagraph:P113:end

% table: NEW-U09-NOM-FYRSPEARCA
% mdtable:T19:begin
\par\begingroup
\fontsize{10.5}{13.2}\selectfont
\setlength{\LTleft}{0pt}\setlength{\LTright}{\fill}
\noindent
\begin{tabular}{@{}>{\raggedright\arraybackslash}p{35.87mm}>{\raggedright\arraybackslash}p{39.46mm}>{\raggedright\arraybackslash}p{43.05mm}@{}}
\toprule
\textbf{Case} & \textbf{Singular} & \textbf{Plural} \\
\midrule
Nominative & fȳrspearca & fȳrspearcan \\
Genitive & fȳrspearcan & fȳrspearcena \\
Dative & fȳrspearcan & fȳrspearcum \\
Accusative & fȳrspearcan & fȳrspearcan \\
\bottomrule
\end{tabular}
\par\endgroup

% mdtable:T19:end
% teaching: NEW-U09-VOC-fyrspearca
% teaching: NEW-U09-NOM-TREOWWYRHTA
% teaching: NEW-U09-VOC-treowwyrhta
% teaching: NEW-U09-NOM-SCEOWYRHTA
% teaching: NEW-U09-VOC-sceowyrhta
% teaching: NEW-U09-NOM-CEMPA
% teaching: NEW-U09-VOC-cempa
% teaching: NEW-U09-NOM-WYRHTA
% teaching: NEW-U09-VOC-wyrhta
\par\begin{minipage}{\linewidth}\raggedright
% mdparagraph:P114:begin
Vocabulary following \textbf{fȳrspearca}:
% mdparagraph:P114:end

% mdtable:T20:begin
\par\begingroup
\fontsize{10.5}{13.2}\selectfont
\setlength{\LTleft}{0pt}\setlength{\LTright}{\fill}
\noindent
\begin{tabular}{@{}>{\raggedright\arraybackslash}p{21.5mm}>{\raggedright\arraybackslash}p{21.57mm}>{\raggedright\arraybackslash}p{20.5mm}>{\raggedright\arraybackslash}p{24.59mm}>{\raggedright\arraybackslash}p{24.59mm}@{}}
\toprule
\textbf{Noun} & \textbf{Gender; meaning} & \textbf{Stem} & \textbf{Genitive singular} & \textbf{Nominative plural} \\
\midrule
fȳrspearca & m.; spark & fȳrspearc- & fȳrspearcan & fȳrspearcan \\
trēowwyrhta & m.; carpenter & trēowwyrht- & trēowwyrhtan & trēowwyrhtan \\
scēowyrhta & m.; shoemaker & scēowyrht- & scēowyrhtan & scēowyrhtan \\
cęmpa & m.; soldier & cęmp- & cęmpan & cęmpan \\
wyrhta & m.; workman; craftsman & wyrht- & wyrhtan & wyrhtan \\
\bottomrule
\end{tabular}
\par\endgroup

% mdtable:T20:end
\end{minipage}\par
\reserve{7\baselineskip}
% mdparagraph:P115:begin
\subsection*{Weak feminines}\BookHeading{nav-ch09-learner-0029}{subsubsection}{Weak feminines}{0}
% mdparagraph:P115:end

% mdparagraph:P116:begin
Weak feminine nominative singular ends in -e. The other seven cells use the same endings as the weak masculine pattern. Smiððe and sīde follow this pattern.
% mdparagraph:P116:end

% teaching: NEW-U09-NOM-SMITHTHE
\reserve{7\baselineskip}
% mdparagraph:P117:begin
\subsection*{smiððe}\BookHeading{nav-ch09-learner-0030}{subsubsection}{smiððe}{0}
% mdparagraph:P117:end

% table: NEW-U09-NOM-SMITHTHE
% mdtable:T21:begin
\par\begingroup
\fontsize{10.5}{13.2}\selectfont
\setlength{\LTleft}{0pt}\setlength{\LTright}{\fill}
\noindent
\begin{tabular}{@{}>{\raggedright\arraybackslash}p{39.46mm}>{\raggedright\arraybackslash}p{39.46mm}>{\raggedright\arraybackslash}p{39.46mm}@{}}
\toprule
\textbf{Case} & \textbf{Singular} & \textbf{Plural} \\
\midrule
Nominative & smiððe & smiððan \\
Genitive & smiððan & smiððena \\
Dative & smiððan & smiððum \\
Accusative & smiððan & smiððan \\
\bottomrule
\end{tabular}
\par\endgroup

% mdtable:T21:end
% teaching: NEW-U09-VOC-smiththe
% teaching: NEW-U09-NOM-SIDE
% teaching: NEW-U09-VOC-side
% mdparagraph:P118:begin
Vocabulary following \textbf{smiððe}:
% mdparagraph:P118:end

% mdtable:T22:begin
\par\begingroup
\fontsize{10.5}{13.2}\selectfont
\setlength{\LTleft}{0pt}\setlength{\LTright}{\fill}
\noindent
\begin{tabular}{@{}>{\raggedright\arraybackslash}p{19.03mm}>{\raggedright\arraybackslash}p{26.74mm}>{\raggedright\arraybackslash}p{19.03mm}>{\raggedright\arraybackslash}p{22.93mm}>{\raggedright\arraybackslash}p{25.02mm}@{}}
\toprule
\textbf{Noun} & \textbf{Gender; meaning} & \textbf{Stem} & \textbf{Genitive singular} & \textbf{Nominative plural} \\
\midrule
smiððe & f.; smithy & smiðð- & smiððan & smiððan \\
sīde & f.; silk (singular) & sīd- & sīdan & sīdan \\
\bottomrule
\end{tabular}
\par\endgroup

% mdtable:T22:end
\reserve{7\baselineskip}
% mdparagraph:P119:begin
\subsection*{A compound of mann}\BookHeading{nav-ch09-learner-0031}{subsubsection}{A compound of mann}{0}
% mdparagraph:P119:end

% mdparagraph:P120:begin
The final member follows mann: dative singular and nominative/\allowbreak{}accusative plural have ealdormęnn. The other forms retain ealdormann-.
% mdparagraph:P120:end

% teaching: NEW-U09-NOM-EALDORMANN
\reserve{7\baselineskip}
% mdparagraph:P121:begin
\subsection*{ealdormann}\BookHeading{nav-ch09-learner-0032}{subsubsection}{ealdormann}{0}
% mdparagraph:P121:end

% table: NEW-U09-NOM-EALDORMANN
% mdtable:T23:begin
\par\begingroup
\fontsize{10.5}{13.2}\selectfont
\setlength{\LTleft}{0pt}\setlength{\LTright}{\fill}
\noindent
\begin{tabular}{@{}>{\raggedright\arraybackslash}p{34.82mm}>{\raggedright\arraybackslash}p{41.78mm}>{\raggedright\arraybackslash}p{41.78mm}@{}}
\toprule
\textbf{Case} & \textbf{Singular} & \textbf{Plural} \\
\midrule
Nominative & ealdormann & ealdormęnn \\
Genitive & ealdormannes & ealdormanna \\
Dative & ealdormęnn & ealdormannum \\
Accusative & ealdormann & ealdormęnn \\
\bottomrule
\end{tabular}
\par\endgroup

% mdtable:T23:end
% teaching: NEW-U09-VOC-ealdormann
% mdparagraph:P122:begin
Vocabulary following \textbf{ealdormann}:
% mdparagraph:P122:end

% mdtable:T24:begin
\par\begingroup
\fontsize{10.5}{13.2}\selectfont
\setlength{\LTleft}{0pt}\setlength{\LTright}{\fill}
\noindent
\begin{tabular}{@{}>{\raggedright\arraybackslash}p{20mm}>{\raggedright\arraybackslash}p{23.1mm}>{\raggedright\arraybackslash}p{21mm}>{\raggedright\arraybackslash}p{25.3mm}>{\raggedright\arraybackslash}p{23.35mm}@{}}
\toprule
\textbf{Noun} & \textbf{Gender; meaning} & \textbf{Stem} & \textbf{Genitive singular} & \textbf{Nominative plural} \\
\midrule
ealdormann & m.; prince; leader; noble & ealdormann- & ealdormannes & ealdormęnn \\
\bottomrule
\end{tabular}
\par\endgroup

% mdtable:T24:end
\reserve{7\baselineskip}
% mdparagraph:P123:begin
\subsection*{Sunu: the u-stem pattern}\BookHeading{nav-ch09-learner-0033}{subsubsection}{Sunu: the u-stem pattern}{0}\BookLabel{bk-ch09-learner-6f36190f5839}
% mdparagraph:P123:end

% teaching: NEW-U09-GRAM-u09-sunu-u-stem
% mdparagraph:P124:begin
Sunu has short u. Nominative and accusative singular end in -u; genitive and dative singular in -a. The plural has -a except in the dative, which has -um. Cook also records -u/\allowbreak{}-o variants outside the genitives and dative plural, and later forms sunes/\allowbreak{}sunas following the ordinary masculine pattern. The table gives the selected pattern.
% mdparagraph:P124:end

% teaching: NEW-U09-NOM-SUNU
\reserve{7\baselineskip}
% mdparagraph:P125:begin
\subsection*{sunu}\BookHeading{nav-ch09-learner-0034}{subsubsection}{sunu}{0}
% mdparagraph:P125:end

% table: NEW-U09-NOM-SUNU
% mdtable:T25:begin
\par\begingroup
\fontsize{10.5}{13.2}\selectfont
\setlength{\LTleft}{0pt}\setlength{\LTright}{\fill}
\noindent
\begin{tabular}{@{}>{\raggedright\arraybackslash}p{39.46mm}>{\raggedright\arraybackslash}p{39.46mm}>{\raggedright\arraybackslash}p{39.46mm}@{}}
\toprule
\textbf{Case} & \textbf{Singular} & \textbf{Plural} \\
\midrule
Nominative & sunu & suna \\
Genitive & suna & suna \\
Dative & suna & sunum \\
Accusative & sunu & suna \\
\bottomrule
\end{tabular}
\par\endgroup

% mdtable:T25:end
% teaching: NEW-U09-VOC-sunu
% mdparagraph:P126:begin
Vocabulary following \textbf{sunu}:
% mdparagraph:P126:end

% mdtable:T26:begin
\par\begingroup
\fontsize{10.5}{13.2}\selectfont
\setlength{\LTleft}{0pt}\setlength{\LTright}{\fill}
\noindent
\begin{tabular}{@{}>{\raggedright\arraybackslash}p{20.21mm}>{\raggedright\arraybackslash}p{21.42mm}>{\raggedright\arraybackslash}p{20.21mm}>{\raggedright\arraybackslash}p{24.35mm}>{\raggedright\arraybackslash}p{26.57mm}@{}}
\toprule
\textbf{Noun} & \textbf{Gender; meaning} & \textbf{Stem} & \textbf{Genitive singular} & \textbf{Nominative plural} \\
\midrule
sunu & m.; son & sun- & suna & suna \\
\bottomrule
\end{tabular}
\par\endgroup

% mdtable:T26:end
\reserve{7\baselineskip}
% mdparagraph:P127:begin
\subsection*{Rīce: a neuter in -e}\BookHeading{nav-ch09-learner-0035}{subsubsection}{Rīce: a neuter in -e}{0}\BookLabel{bk-ch09-learner-47d8491e816c}
% mdparagraph:P127:end

% teaching: NEW-U09-GRAM-u09-rice-neuter-e
% mdparagraph:P128:begin
The final -e remains in the singular where no other ending replaces it. Use rīc- before plural endings: rīcu, rīca, rīcum. Cook also gives forms with inserted i/\allowbreak{}e after c or g; rīciu and rīcia are documented alternatives to selected rīcu and rīca.
% mdparagraph:P128:end

% teaching: NEW-U09-NOM-RICE
\reserve{7\baselineskip}
% mdparagraph:P129:begin
\subsection*{rīce}\BookHeading{nav-ch09-learner-0036}{subsubsection}{rīce}{0}
% mdparagraph:P129:end

% table: NEW-U09-NOM-RICE
% mdtable:T27:begin
\par\begingroup
\fontsize{10.5}{13.2}\selectfont
\setlength{\LTleft}{0pt}\setlength{\LTright}{\fill}
\noindent
\begin{tabular}{@{}>{\raggedright\arraybackslash}p{39.46mm}>{\raggedright\arraybackslash}p{39.46mm}>{\raggedright\arraybackslash}p{39.46mm}@{}}
\toprule
\textbf{Case} & \textbf{Singular} & \textbf{Plural} \\
\midrule
Nominative & rīce & rīcu \\
Genitive & rīces & rīca \\
Dative & rīce & rīcum \\
Accusative & rīce & rīcu \\
\bottomrule
\end{tabular}
\par\endgroup

% mdtable:T27:end
% teaching: NEW-U09-VOC-rice
% mdparagraph:P130:begin
Vocabulary following \textbf{rīce}:
% mdparagraph:P130:end

% mdtable:T28:begin
\par\begingroup
\fontsize{10.5}{13.2}\selectfont
\setlength{\LTleft}{0pt}\setlength{\LTright}{\fill}
\noindent
\begin{tabular}{@{}>{\raggedright\arraybackslash}p{20.21mm}>{\raggedright\arraybackslash}p{21.42mm}>{\raggedright\arraybackslash}p{20.21mm}>{\raggedright\arraybackslash}p{24.35mm}>{\raggedright\arraybackslash}p{26.57mm}@{}}
\toprule
\textbf{Noun} & \textbf{Gender; meaning} & \textbf{Stem} & \textbf{Genitive singular} & \textbf{Nominative plural} \\
\midrule
rīce & n.; kingdom & rīc- & rīces & rīcu \\
\bottomrule
\end{tabular}
\par\endgroup

% mdtable:T28:end
\reserve{7\baselineskip}
% mdparagraph:P131:begin
\subsection*{Geðeahtend: an agent noun in -end}\BookHeading{nav-ch09-learner-0037}{subsubsection}{Geðeahtend: an agent noun in -end}{0}\BookLabel{bk-ch09-learner-a5032b3d8e24}
% mdparagraph:P131:end

% teaching: NEW-U09-GRAM-u09-agent-end
% mdparagraph:P132:begin
Geðeahtend names a person, a counsellor. Its selected nominative and accusative plural have no ending; genitive plural has -ra and dative plural -um. Plural nominative/\allowbreak{}accusative -e and the rarer -as also occur. This noun pattern is distinct from the adjective inflection of a present participle in -ende.
% mdparagraph:P132:end

% teaching: NEW-U09-NOM-GETHEAHTEND
\reserve{7\baselineskip}
% mdparagraph:P133:begin
\subsection*{geðeahtend}\BookHeading{nav-ch09-learner-0038}{subsubsection}{geðeahtend}{0}
% mdparagraph:P133:end

% table: NEW-U09-NOM-GETHEAHTEND
% mdtable:T29:begin
\par\begingroup
\fontsize{10.5}{13.2}\selectfont
\setlength{\LTleft}{0pt}\setlength{\LTright}{\fill}
\noindent
\begin{tabular}{@{}>{\raggedright\arraybackslash}p{34.82mm}>{\raggedright\arraybackslash}p{41.78mm}>{\raggedright\arraybackslash}p{41.78mm}@{}}
\toprule
\textbf{Case} & \textbf{Singular} & \textbf{Plural} \\
\midrule
Nominative & geðeahtend & geðeahtend \\
Genitive & geðeahtendes & geðeahtendra \\
Dative & geðeahtende & geðeahtendum \\
Accusative & geðeahtend & geðeahtend \\
\bottomrule
\end{tabular}
\par\endgroup

% mdtable:T29:end
% teaching: NEW-U09-VOC-getheahtend
% mdparagraph:P134:begin
Vocabulary following \textbf{geðeahtend}:
% mdparagraph:P134:end

% mdtable:T30:begin
\par\begingroup
\fontsize{10.5}{13.2}\selectfont
\setlength{\LTleft}{0pt}\setlength{\LTright}{\fill}
\noindent
\begin{tabular}{@{}>{\raggedright\arraybackslash}p{18.59mm}>{\raggedright\arraybackslash}p{22.77mm}>{\raggedright\arraybackslash}p{20.45mm}>{\raggedright\arraybackslash}p{26.49mm}>{\raggedright\arraybackslash}p{24.45mm}@{}}
\toprule
\textbf{Noun} & \textbf{Gender; meaning} & \textbf{Stem} & \textbf{Genitive singular} & \textbf{Nominative plural} \\
\midrule
geðeahtend & m.; counsellor & geðeahtend- & geðeahtendes & geðeahtend \\
\bottomrule
\end{tabular}
\par\endgroup

% mdtable:T30:end
% mdparagraph:P135:begin
Distinguish \textbf{beorn}, “warrior”, from \textbf{cęmpa}, “soldier”. \textbf{Mæsseprēost} is the selected word for “priest”; \textbf{ceorl} here means “layman”. \textbf{Hlæst} is used as plural “wares”. The selected use of \textbf{bęlg} refers to plural bellows/\allowbreak{}bags; choose verb agreement from the Old English number rather than from the shape of the English word “bellows”.
% mdparagraph:P135:end

\reserve{7\baselineskip}
% mdparagraph:P136:begin
\section*{6. Adjectives and agreement}\BookHeading{nav-ch09-learner-0039}{subsection}{6. Adjectives and agreement}{1}
% mdparagraph:P136:end

% mdparagraph:P137:begin
The complete strong and weak forms follow. Use strong adjectives in ordinary predication or unmarked attribution, weak adjectives after the article/\allowbreak{}demonstrative and in direct address. Comparatives are weak. The tables preserve the selected feminine plural -a; the familiar -e alternative is also valid.
% mdparagraph:P137:end

% mdparagraph:P138:begin
The distinct instrumental is not added to these four-case practice patterns.
% mdparagraph:P138:end

\reserve{7\baselineskip}
% mdparagraph:P139:begin
\subsection*{behēfe — useful}\BookHeading{nav-ch09-learner-0040}{subsubsection}{behēfe — useful}{0}
% mdparagraph:P139:end

% teaching: NEW-U09-VOC-behefe
% mdparagraph:P140:begin
Use behēf- before endings. As an adjective in -e, behēfe retains that e only where another ending does not replace it. The selected strong feminine nominative singular and neuter nominative/\allowbreak{}accusative plural are behēfu.
% mdparagraph:P140:end

% teaching: NEW-U09-ADJ-BEHEFE-STRONG
\reserve{7\baselineskip}
% mdparagraph:P141:begin
\subsection*{behēfe: strong declension, singular}\BookHeading{nav-ch09-learner-0041}{subsubsection}{behēfe: strong declension, singular}{0}
% mdparagraph:P141:end

% table: NEW-U09-ADJ-BEHEFE-STRONG-SG
% mdtable:T31:begin
\par\begingroup
\fontsize{10.5}{13.2}\selectfont
\setlength{\LTleft}{0pt}\setlength{\LTright}{\fill}
\noindent
\begin{tabular}{@{}>{\raggedright\arraybackslash}p{28.89mm}>{\raggedright\arraybackslash}p{28.89mm}>{\raggedright\arraybackslash}p{28.89mm}>{\raggedright\arraybackslash}p{28.89mm}@{}}
\toprule
\textbf{Case} & \textbf{Masculine} & \textbf{Neuter} & \textbf{Feminine} \\
\midrule
Nominative & behēfe & behēfe & behēfu \\
Genitive & behēfes & behēfes & behēfre \\
Dative & behēfum & behēfum & behēfre \\
Accusative & behēfne & behēfe & behēfe \\
\bottomrule
\end{tabular}
\par\endgroup

% mdtable:T31:end
\reserve{7\baselineskip}
% mdparagraph:P142:begin
\subsection*{behēfe: strong declension, plural}\BookHeading{nav-ch09-learner-0042}{subsubsection}{behēfe: strong declension, plural}{0}
% mdparagraph:P142:end

% table: NEW-U09-ADJ-BEHEFE-STRONG-PL
% mdtable:T32:begin
\par\begingroup
\fontsize{10.5}{13.2}\selectfont
\setlength{\LTleft}{0pt}\setlength{\LTright}{\fill}
\noindent
\begin{tabular}{@{}>{\raggedright\arraybackslash}p{28.89mm}>{\raggedright\arraybackslash}p{28.89mm}>{\raggedright\arraybackslash}p{28.89mm}>{\raggedright\arraybackslash}p{28.89mm}@{}}
\toprule
\textbf{Case} & \textbf{Masculine} & \textbf{Neuter} & \textbf{Feminine} \\
\midrule
Nominative & behēfe & behēfu & behēfa \\
Genitive & behēfra & behēfra & behēfra \\
Dative & behēfum & behēfum & behēfum \\
Accusative & behēfe & behēfu & behēfa \\
\bottomrule
\end{tabular}
\par\endgroup

% mdtable:T32:end
% teaching: NEW-U09-ADJ-BEHEFE-WEAK
\reserve{7\baselineskip}
% mdparagraph:P143:begin
\subsection*{behēfe: weak declension, singular}\BookHeading{nav-ch09-learner-0043}{subsubsection}{behēfe: weak declension, singular}{0}
% mdparagraph:P143:end

% table: NEW-U09-ADJ-BEHEFE-WEAK-SG
% mdtable:T33:begin
\par\begingroup
\fontsize{10.5}{13.2}\selectfont
\setlength{\LTleft}{0pt}\setlength{\LTright}{\fill}
\noindent
\begin{tabular}{@{}>{\raggedright\arraybackslash}p{28.89mm}>{\raggedright\arraybackslash}p{28.89mm}>{\raggedright\arraybackslash}p{28.89mm}>{\raggedright\arraybackslash}p{28.89mm}@{}}
\toprule
\textbf{Case} & \textbf{Masculine} & \textbf{Neuter} & \textbf{Feminine} \\
\midrule
Nominative & behēfa & behēfe & behēfe \\
Genitive & behēfan & behēfan & behēfan \\
Dative & behēfan & behēfan & behēfan \\
Accusative & behēfan & behēfe & behēfan \\
\bottomrule
\end{tabular}
\par\endgroup

% mdtable:T33:end
\reserve{7\baselineskip}
% mdparagraph:P144:begin
\subsection*{behēfe: weak declension, plural}\BookHeading{nav-ch09-learner-0044}{subsubsection}{behēfe: weak declension, plural}{0}
% mdparagraph:P144:end

% table: NEW-U09-ADJ-BEHEFE-WEAK-PL
% mdtable:T34:begin
\par\begingroup
\fontsize{10.5}{13.2}\selectfont
\setlength{\LTleft}{0pt}\setlength{\LTright}{\fill}
\noindent
\begin{tabular}{@{}>{\raggedright\arraybackslash}p{28.89mm}>{\raggedright\arraybackslash}p{28.89mm}>{\raggedright\arraybackslash}p{28.89mm}>{\raggedright\arraybackslash}p{28.89mm}@{}}
\toprule
\textbf{Case} & \textbf{Masculine} & \textbf{Neuter} & \textbf{Feminine} \\
\midrule
Nominative & behēfan & behēfan & behēfan \\
Genitive & behēfra & behēfra & behēfra \\
Dative & behēfum & behēfum & behēfum \\
Accusative & behēfan & behēfan & behēfan \\
\bottomrule
\end{tabular}
\par\endgroup

% mdtable:T34:end
\reserve{7\baselineskip}
% mdparagraph:P145:begin
\subsection*{welig — rich; wealthy}\BookHeading{nav-ch09-learner-0045}{subsubsection}{welig — rich; wealthy}{0}
% mdparagraph:P145:end

% teaching: NEW-U09-VOC-welig
% mdparagraph:P146:begin
Keep the stem welig- in these selected forms, including the medial i. Strong feminine nominative singular and neuter nominative/\allowbreak{}accusative plural have -u: weligu.
% mdparagraph:P146:end

% teaching: NEW-U09-ADJ-WELIG-STRONG
\reserve{7\baselineskip}
% mdparagraph:P147:begin
\subsection*{welig: strong declension, singular}\BookHeading{nav-ch09-learner-0046}{subsubsection}{welig: strong declension, singular}{0}
% mdparagraph:P147:end

% table: NEW-U09-ADJ-WELIG-STRONG-SG
% mdtable:T35:begin
\par\begingroup
\fontsize{10.5}{13.2}\selectfont
\setlength{\LTleft}{0pt}\setlength{\LTright}{\fill}
\noindent
\begin{tabular}{@{}>{\raggedright\arraybackslash}p{28.89mm}>{\raggedright\arraybackslash}p{28.89mm}>{\raggedright\arraybackslash}p{28.89mm}>{\raggedright\arraybackslash}p{28.89mm}@{}}
\toprule
\textbf{Case} & \textbf{Masculine} & \textbf{Neuter} & \textbf{Feminine} \\
\midrule
Nominative & welig & welig & weligu \\
Genitive & weliges & weliges & weligre \\
Dative & weligum & weligum & weligre \\
Accusative & weligne & welig & welige \\
\bottomrule
\end{tabular}
\par\endgroup

% mdtable:T35:end
\reserve{7\baselineskip}
% mdparagraph:P148:begin
\subsection*{welig: strong declension, plural}\BookHeading{nav-ch09-learner-0047}{subsubsection}{welig: strong declension, plural}{0}
% mdparagraph:P148:end

% table: NEW-U09-ADJ-WELIG-STRONG-PL
% mdtable:T36:begin
\par\begingroup
\fontsize{10.5}{13.2}\selectfont
\setlength{\LTleft}{0pt}\setlength{\LTright}{\fill}
\noindent
\begin{tabular}{@{}>{\raggedright\arraybackslash}p{28.89mm}>{\raggedright\arraybackslash}p{28.89mm}>{\raggedright\arraybackslash}p{28.89mm}>{\raggedright\arraybackslash}p{28.89mm}@{}}
\toprule
\textbf{Case} & \textbf{Masculine} & \textbf{Neuter} & \textbf{Feminine} \\
\midrule
Nominative & welige & weligu & weliga \\
Genitive & weligra & weligra & weligra \\
Dative & weligum & weligum & weligum \\
Accusative & welige & weligu & weliga \\
\bottomrule
\end{tabular}
\par\endgroup

% mdtable:T36:end
% teaching: NEW-U09-ADJ-WELIG-WEAK
\reserve{7\baselineskip}
% mdparagraph:P149:begin
\subsection*{welig: weak declension, singular}\BookHeading{nav-ch09-learner-0048}{subsubsection}{welig: weak declension, singular}{0}
% mdparagraph:P149:end

% table: NEW-U09-ADJ-WELIG-WEAK-SG
% mdtable:T37:begin
\par\begingroup
\fontsize{10.5}{13.2}\selectfont
\setlength{\LTleft}{0pt}\setlength{\LTright}{\fill}
\noindent
\begin{tabular}{@{}>{\raggedright\arraybackslash}p{28.89mm}>{\raggedright\arraybackslash}p{28.89mm}>{\raggedright\arraybackslash}p{28.89mm}>{\raggedright\arraybackslash}p{28.89mm}@{}}
\toprule
\textbf{Case} & \textbf{Masculine} & \textbf{Neuter} & \textbf{Feminine} \\
\midrule
Nominative & weliga & welige & welige \\
Genitive & weligan & weligan & weligan \\
Dative & weligan & weligan & weligan \\
Accusative & weligan & welige & weligan \\
\bottomrule
\end{tabular}
\par\endgroup

% mdtable:T37:end
\reserve{7\baselineskip}
% mdparagraph:P150:begin
\subsection*{welig: weak declension, plural}\BookHeading{nav-ch09-learner-0049}{subsubsection}{welig: weak declension, plural}{0}
% mdparagraph:P150:end

% table: NEW-U09-ADJ-WELIG-WEAK-PL
% mdtable:T38:begin
\par\begingroup
\fontsize{10.5}{13.2}\selectfont
\setlength{\LTleft}{0pt}\setlength{\LTright}{\fill}
\noindent
\begin{tabular}{@{}>{\raggedright\arraybackslash}p{28.89mm}>{\raggedright\arraybackslash}p{28.89mm}>{\raggedright\arraybackslash}p{28.89mm}>{\raggedright\arraybackslash}p{28.89mm}@{}}
\toprule
\textbf{Case} & \textbf{Masculine} & \textbf{Neuter} & \textbf{Feminine} \\
\midrule
Nominative & weligan & weligan & weligan \\
Genitive & weligra & weligra & weligra \\
Dative & weligum & weligum & weligum \\
Accusative & weligan & weligan & weligan \\
\bottomrule
\end{tabular}
\par\endgroup

% mdtable:T38:end
\reserve{7\baselineskip}
% mdparagraph:P151:begin
\subsection*{wīs — wise}\BookHeading{nav-ch09-learner-0050}{subsubsection}{wīs — wise}{0}
% mdparagraph:P151:end

% teaching: NEW-U09-VOC-wis
% mdparagraph:P152:begin
The heavy stem wīs- has no -u in strong feminine nominative singular or neuter nominative/\allowbreak{}accusative plural.
% mdparagraph:P152:end

% teaching: NEW-U09-ADJ-WIS-STRONG
\reserve{7\baselineskip}
% mdparagraph:P153:begin
\subsection*{wīs: strong declension, singular}\BookHeading{nav-ch09-learner-0051}{subsubsection}{wīs: strong declension, singular}{0}
% mdparagraph:P153:end

% table: NEW-U09-ADJ-WIS-STRONG-SG
% mdtable:T39:begin
\par\begingroup
\fontsize{10.5}{13.2}\selectfont
\setlength{\LTleft}{0pt}\setlength{\LTright}{\fill}
\noindent
\begin{tabular}{@{}>{\raggedright\arraybackslash}p{28.89mm}>{\raggedright\arraybackslash}p{28.89mm}>{\raggedright\arraybackslash}p{28.89mm}>{\raggedright\arraybackslash}p{28.89mm}@{}}
\toprule
\textbf{Case} & \textbf{Masculine} & \textbf{Neuter} & \textbf{Feminine} \\
\midrule
Nominative & wīs & wīs & wīs \\
Genitive & wīses & wīses & wīsre \\
Dative & wīsum & wīsum & wīsre \\
Accusative & wīsne & wīs & wīse \\
\bottomrule
\end{tabular}
\par\endgroup

% mdtable:T39:end
\reserve{7\baselineskip}
% mdparagraph:P154:begin
\subsection*{wīs: strong declension, plural}\BookHeading{nav-ch09-learner-0052}{subsubsection}{wīs: strong declension, plural}{0}
% mdparagraph:P154:end

% table: NEW-U09-ADJ-WIS-STRONG-PL
% mdtable:T40:begin
\par\begingroup
\fontsize{10.5}{13.2}\selectfont
\setlength{\LTleft}{0pt}\setlength{\LTright}{\fill}
\noindent
\begin{tabular}{@{}>{\raggedright\arraybackslash}p{28.89mm}>{\raggedright\arraybackslash}p{28.89mm}>{\raggedright\arraybackslash}p{28.89mm}>{\raggedright\arraybackslash}p{28.89mm}@{}}
\toprule
\textbf{Case} & \textbf{Masculine} & \textbf{Neuter} & \textbf{Feminine} \\
\midrule
Nominative & wīse & wīs & wīsa \\
Genitive & wīsra & wīsra & wīsra \\
Dative & wīsum & wīsum & wīsum \\
Accusative & wīse & wīs & wīsa \\
\bottomrule
\end{tabular}
\par\endgroup

% mdtable:T40:end
% teaching: NEW-U09-ADJ-WIS-WEAK
\reserve{7\baselineskip}
% mdparagraph:P155:begin
\subsection*{wīs: weak declension, singular}\BookHeading{nav-ch09-learner-0053}{subsubsection}{wīs: weak declension, singular}{0}
% mdparagraph:P155:end

% table: NEW-U09-ADJ-WIS-WEAK-SG
% mdtable:T41:begin
\par\begingroup
\fontsize{10.5}{13.2}\selectfont
\setlength{\LTleft}{0pt}\setlength{\LTright}{\fill}
\noindent
\begin{tabular}{@{}>{\raggedright\arraybackslash}p{28.89mm}>{\raggedright\arraybackslash}p{28.89mm}>{\raggedright\arraybackslash}p{28.89mm}>{\raggedright\arraybackslash}p{28.89mm}@{}}
\toprule
\textbf{Case} & \textbf{Masculine} & \textbf{Neuter} & \textbf{Feminine} \\
\midrule
Nominative & wīsa & wīse & wīse \\
Genitive & wīsan & wīsan & wīsan \\
Dative & wīsan & wīsan & wīsan \\
Accusative & wīsan & wīse & wīsan \\
\bottomrule
\end{tabular}
\par\endgroup

% mdtable:T41:end
\reserve{7\baselineskip}
% mdparagraph:P156:begin
\subsection*{wīs: weak declension, plural}\BookHeading{nav-ch09-learner-0054}{subsubsection}{wīs: weak declension, plural}{0}
% mdparagraph:P156:end

% table: NEW-U09-ADJ-WIS-WEAK-PL
% mdtable:T42:begin
\par\begingroup
\fontsize{10.5}{13.2}\selectfont
\setlength{\LTleft}{0pt}\setlength{\LTright}{\fill}
\noindent
\begin{tabular}{@{}>{\raggedright\arraybackslash}p{28.89mm}>{\raggedright\arraybackslash}p{28.89mm}>{\raggedright\arraybackslash}p{28.89mm}>{\raggedright\arraybackslash}p{28.89mm}@{}}
\toprule
\textbf{Case} & \textbf{Masculine} & \textbf{Neuter} & \textbf{Feminine} \\
\midrule
Nominative & wīsan & wīsan & wīsan \\
Genitive & wīsra & wīsra & wīsra \\
Dative & wīsum & wīsum & wīsum \\
Accusative & wīsan & wīsan & wīsan \\
\bottomrule
\end{tabular}
\par\endgroup

% mdtable:T42:end
\reserve{7\baselineskip}
% mdparagraph:P157:begin
\subsection*{manig — many}\BookHeading{nav-ch09-learner-0055}{subsubsection}{manig — many}{0}\BookLabel{bk-ch09-learner-23e51d68091c}
% mdparagraph:P157:end

% teaching: NEW-U09-VOC-manig
% teaching: NEW-U09-GRAM-MANIG
% mdparagraph:P158:begin
Use manig as an ordinary attributive adjective with a plural noun. The singular tables below are reference only; they do not add “many a” or a new substantival use. Keep manig- in the selected forms; neuter nominative/\allowbreak{}accusative plural is manigu.
% mdparagraph:P158:end

% teaching: NEW-U09-ADJ-MANIG-STRONG
\reserve{7\baselineskip}
% mdparagraph:P159:begin
\subsection*{manig: strong declension, singular}\BookHeading{nav-ch09-learner-0056}{subsubsection}{manig: strong declension, singular}{0}
% mdparagraph:P159:end

% table: NEW-U09-ADJ-MANIG-STRONG-SG
% mdtable:T43:begin
\par\begingroup
\fontsize{10.5}{13.2}\selectfont
\setlength{\LTleft}{0pt}\setlength{\LTright}{\fill}
\noindent
\begin{tabular}{@{}>{\raggedright\arraybackslash}p{28.89mm}>{\raggedright\arraybackslash}p{28.89mm}>{\raggedright\arraybackslash}p{28.89mm}>{\raggedright\arraybackslash}p{28.89mm}@{}}
\toprule
\textbf{Case} & \textbf{Masculine} & \textbf{Neuter} & \textbf{Feminine} \\
\midrule
Nominative & manig & manig & manigu \\
Genitive & maniges & maniges & manigre \\
Dative & manigum & manigum & manigre \\
Accusative & manigne & manig & manige \\
\bottomrule
\end{tabular}
\par\endgroup

% mdtable:T43:end
\reserve{7\baselineskip}
% mdparagraph:P160:begin
\subsection*{manig: strong declension, plural}\BookHeading{nav-ch09-learner-0057}{subsubsection}{manig: strong declension, plural}{0}
% mdparagraph:P160:end

% table: NEW-U09-ADJ-MANIG-STRONG-PL
% mdtable:T44:begin
\par\begingroup
\fontsize{10.5}{13.2}\selectfont
\setlength{\LTleft}{0pt}\setlength{\LTright}{\fill}
\noindent
\begin{tabular}{@{}>{\raggedright\arraybackslash}p{28.89mm}>{\raggedright\arraybackslash}p{28.89mm}>{\raggedright\arraybackslash}p{28.89mm}>{\raggedright\arraybackslash}p{28.89mm}@{}}
\toprule
\textbf{Case} & \textbf{Masculine} & \textbf{Neuter} & \textbf{Feminine} \\
\midrule
Nominative & manige & manigu & maniga \\
Genitive & manigra & manigra & manigra \\
Dative & manigum & manigum & manigum \\
Accusative & manige & manigu & maniga \\
\bottomrule
\end{tabular}
\par\endgroup

% mdtable:T44:end
% teaching: NEW-U09-ADJ-MANIG-WEAK
\reserve{7\baselineskip}
% mdparagraph:P161:begin
\subsection*{manig: weak declension, singular}\BookHeading{nav-ch09-learner-0058}{subsubsection}{manig: weak declension, singular}{0}
% mdparagraph:P161:end

% table: NEW-U09-ADJ-MANIG-WEAK-SG
% mdtable:T45:begin
\par\begingroup
\fontsize{10.5}{13.2}\selectfont
\setlength{\LTleft}{0pt}\setlength{\LTright}{\fill}
\noindent
\begin{tabular}{@{}>{\raggedright\arraybackslash}p{28.89mm}>{\raggedright\arraybackslash}p{28.89mm}>{\raggedright\arraybackslash}p{28.89mm}>{\raggedright\arraybackslash}p{28.89mm}@{}}
\toprule
\textbf{Case} & \textbf{Masculine} & \textbf{Neuter} & \textbf{Feminine} \\
\midrule
Nominative & maniga & manige & manige \\
Genitive & manigan & manigan & manigan \\
Dative & manigan & manigan & manigan \\
Accusative & manigan & manige & manigan \\
\bottomrule
\end{tabular}
\par\endgroup

% mdtable:T45:end
\reserve{7\baselineskip}
% mdparagraph:P162:begin
\subsection*{manig: weak declension, plural}\BookHeading{nav-ch09-learner-0059}{subsubsection}{manig: weak declension, plural}{0}
% mdparagraph:P162:end

% table: NEW-U09-ADJ-MANIG-WEAK-PL
% mdtable:T46:begin
\par\begingroup
\fontsize{10.5}{13.2}\selectfont
\setlength{\LTleft}{0pt}\setlength{\LTright}{\fill}
\noindent
\begin{tabular}{@{}>{\raggedright\arraybackslash}p{28.89mm}>{\raggedright\arraybackslash}p{28.89mm}>{\raggedright\arraybackslash}p{28.89mm}>{\raggedright\arraybackslash}p{28.89mm}@{}}
\toprule
\textbf{Case} & \textbf{Masculine} & \textbf{Neuter} & \textbf{Feminine} \\
\midrule
Nominative & manigan & manigan & manigan \\
Genitive & manigra & manigra & manigra \\
Dative & manigum & manigum & manigum \\
Accusative & manigan & manigan & manigan \\
\bottomrule
\end{tabular}
\par\endgroup

% mdtable:T46:end
\reserve{7\baselineskip}
% mdparagraph:P163:begin
\subsection*{wīsra — wiser}\BookHeading{nav-ch09-learner-0060}{subsubsection}{wīsra — wiser}{0}
% mdparagraph:P163:end

% teaching: NEW-U09-VOC-wis-comp
% mdparagraph:P164:begin
Add the comparative r to wīs-, without changing the vowel, and use weak endings. Genitive plural wīsra simplifies doubled r after s; do not build wīsrra mechanically. This is the same type of reduction seen in earlier lēofra.
% mdparagraph:P164:end

% teaching: NEW-U09-ADJ-WIS-COMP-WEAK
\reserve{7\baselineskip}
% mdparagraph:P165:begin
\subsection*{wīsra: weak declension, singular}\BookHeading{nav-ch09-learner-0061}{subsubsection}{wīsra: weak declension, singular}{0}
% mdparagraph:P165:end

% table: NEW-U09-ADJ-WIS-COMP-WEAK-SG
% mdtable:T47:begin
\par\begingroup
\fontsize{10.5}{13.2}\selectfont
\setlength{\LTleft}{0pt}\setlength{\LTright}{\fill}
\noindent
\begin{tabular}{@{}>{\raggedright\arraybackslash}p{28.89mm}>{\raggedright\arraybackslash}p{28.89mm}>{\raggedright\arraybackslash}p{28.89mm}>{\raggedright\arraybackslash}p{28.89mm}@{}}
\toprule
\textbf{Case} & \textbf{Masculine} & \textbf{Neuter} & \textbf{Feminine} \\
\midrule
Nominative & wīsra & wīsre & wīsre \\
Genitive & wīsran & wīsran & wīsran \\
Dative & wīsran & wīsran & wīsran \\
Accusative & wīsran & wīsre & wīsran \\
\bottomrule
\end{tabular}
\par\endgroup

% mdtable:T47:end
\reserve{7\baselineskip}
% mdparagraph:P166:begin
\subsection*{wīsra: weak declension, plural}\BookHeading{nav-ch09-learner-0062}{subsubsection}{wīsra: weak declension, plural}{0}
% mdparagraph:P166:end

% table: NEW-U09-ADJ-WIS-COMP-WEAK-PL
% mdtable:T48:begin
\par\begingroup
\fontsize{10.5}{13.2}\selectfont
\setlength{\LTleft}{0pt}\setlength{\LTright}{\fill}
\noindent
\begin{tabular}{@{}>{\raggedright\arraybackslash}p{28.89mm}>{\raggedright\arraybackslash}p{28.89mm}>{\raggedright\arraybackslash}p{28.89mm}>{\raggedright\arraybackslash}p{28.89mm}@{}}
\toprule
\textbf{Case} & \textbf{Masculine} & \textbf{Neuter} & \textbf{Feminine} \\
\midrule
Nominative & wīsran & wīsran & wīsran \\
Genitive & wīsra & wīsra & wīsra \\
Dative & wīsrum & wīsrum & wīsrum \\
Accusative & wīsran & wīsran & wīsran \\
\bottomrule
\end{tabular}
\par\endgroup

% mdtable:T48:end
% teaching: NEW-U09-GRAM-USEFUL
% mdparagraph:P167:begin
\textbf{Behēfe},\BookPoint{bk-ch09-learner-c78da7bcc201} “useful”, can name the person to whom something is useful in the dative. The adjective itself agrees with the subject; the recipient remains dative regardless of that subject’s gender.
% mdparagraph:P167:end

% teaching: NEW-U09-GRAM-u09-mixed-gender-predicate
% mdparagraph:P168:begin
A\BookPoint{bk-ch09-learner-db7c4490e70c} predicate adjective referring jointly to nouns of different grammatical genders normally takes the neuter plural. Thus the combination of masculine \textbf{ęle} and neuter \textbf{wīn} calls for the selected neuter plural \textbf{gōd}. The zero ending does not make gōd singular in that use.
% mdparagraph:P168:end

% teaching: NEW-U09-GRAM-ADDRESS
% teaching: NEW-U09-VOC-eala
% teaching: NEW-U09-VOC-EALA
% mdparagraph:P169:begin
In\BookPoint{bk-ch09-learner-9466d7df4e42} direct address, use the noun’s nominative form and a weak accompanying adjective. The number still matters: address one person with singular forms, more than one with plural forms. \textbf{Ēalā}, “O”, is an invariant interjection that may introduce an address or an exclamation.
% mdparagraph:P169:end

\reserve{7\baselineskip}
% mdparagraph:P170:begin
\section*{7. Verbs: complete finite forms}\BookHeading{nav-ch09-learner-0063}{subsection}{7. Verbs: complete finite forms}{1}
% mdparagraph:P170:end

% teaching: NEW-U09-GRAM-u09-new-verb-indicatives
% mdparagraph:P171:begin
Each indicative table gives all three persons in both numbers, in the present and preterite. Plural persons share a form, but still refer to different subjects. Each table includes the verb’s stem changes as well as its endings.
% mdparagraph:P171:end

% mdparagraph:P172:begin
The infinitives heading these tables name the verbs. This unit adds their finite forms; the new infinitive used in a sentence is \textbf{wician} in the preference construction. Earlier productive infinitives such as beran, gān and tēon remain available in the constructions already explained.
% mdparagraph:P172:end

\reserve{7\baselineskip}
% mdparagraph:P173:begin
\subsection*{Weak I}\BookHeading{nav-ch09-learner-0064}{subsubsection}{Weak I}{0}
% mdparagraph:P173:end

\reserve{7\baselineskip}
% mdparagraph:P174:begin
\subsection*{cīepan — sell}\BookHeading{nav-ch09-learner-0065}{subsubsection}{cīepan — sell}{0}
% mdparagraph:P174:end

% teaching: NEW-U09-VOC-ciepan
% mdparagraph:P175:begin
Weak I. The object sold is accusative. The preterite has -pte after p: cīepte.
% mdparagraph:P175:end

% teaching: NEW-U09-PAR-ciepan-IND
\reserve{7\baselineskip}
% mdparagraph:P176:begin
\subsection*{cīepan: complete indicative}\BookHeading{nav-ch09-learner-0066}{subsubsection}{cīepan: complete indicative}{0}
% mdparagraph:P176:end

% table: NEW-U09-PAR-ciepan-IND
% mdtable:T49:begin
% course-table-style: 2026-10
% course-panel:T49-P01:begin
\par\begingroup
\fontsize{10.5}{13.2}\selectfont
\setlength{\LTleft}{0pt}\setlength{\LTright}{\fill}
\noindent
\begin{tabular}{@{}>{\raggedright\arraybackslash}p{39.459mm}>{\raggedright\arraybackslash}p{39.459mm}>{\raggedright\arraybackslash}p{39.459mm}@{}}
\toprule
\textbf{Person and number} & \textbf{Present} & \textbf{Preterite} \\
\midrule
1 sg. & cīepe & cīepte \\
2 sg. & cīepst & cīeptest \\
3 sg. & cīepð & cīepte \\
1–3 pl. & cīepað & cīepton \\
\bottomrule
\end{tabular}
\par\endgroup
% course-panel:T49-P01:end
% mdtable:T49:end
\reserve{7\baselineskip}
% mdparagraph:P177:begin
\subsection*{āfēdan — feed; support}\BookHeading{nav-ch09-learner-0067}{subsubsection}{āfēdan — feed; support}{0}
% mdparagraph:P177:end

% teaching: NEW-U09-VOC-afedan
% mdparagraph:P178:begin
Weak I. The person supported is accusative. Final d and the third-singular dental ending yield tt in āfētt. The preterite keeps -dde: āfēdde.
% mdparagraph:P178:end

% teaching: NEW-U09-PAR-afedan-IND
\reserve{7\baselineskip}
% mdparagraph:P179:begin
\subsection*{āfēdan: complete indicative}\BookHeading{nav-ch09-learner-0068}{subsubsection}{āfēdan: complete indicative}{0}
% mdparagraph:P179:end

% table: NEW-U09-PAR-afedan-IND
% mdtable:T50:begin
% course-table-style: 2026-10
% course-panel:T50-P01:begin
\par\begingroup
\fontsize{10.5}{13.2}\selectfont
\setlength{\LTleft}{0pt}\setlength{\LTright}{\fill}
\noindent
\begin{tabular}{@{}>{\raggedright\arraybackslash}p{39.459mm}>{\raggedright\arraybackslash}p{39.459mm}>{\raggedright\arraybackslash}p{39.459mm}@{}}
\toprule
\textbf{Person and number} & \textbf{Present} & \textbf{Preterite} \\
\midrule
1 sg. & āfēde & āfēdde \\
2 sg. & āfētst & āfēddest \\
3 sg. & āfētt & āfēdde \\
1–3 pl. & āfēdað & āfēddon \\
\bottomrule
\end{tabular}
\par\endgroup
% course-panel:T50-P01:end
% mdtable:T50:end
\reserve{7\baselineskip}
% mdparagraph:P180:begin
\subsection*{bycgan — buy}\BookHeading{nav-ch09-learner-0069}{subsubsection}{bycgan — buy}{0}
% mdparagraph:P180:end

% teaching: NEW-U09-VOC-bycgan
% mdparagraph:P181:begin
Weak I with an exceptional preterite. The goods bought are accusative. Keep bycg- in bycge/\allowbreak{}bycgað, but single g in bygest/\allowbreak{}bygeð. The preterite has bohte and bohton; do not form a regular preterite from bycg-.
% mdparagraph:P181:end

% teaching: NEW-U09-PAR-bycgan-IND
\reserve{7\baselineskip}
% mdparagraph:P182:begin
\subsection*{bycgan: complete indicative}\BookHeading{nav-ch09-learner-0070}{subsubsection}{bycgan: complete indicative}{0}
% mdparagraph:P182:end

% table: NEW-U09-PAR-bycgan-IND
% mdtable:T51:begin
% course-table-style: 2026-10
% course-panel:T51-P01:begin
\par\begingroup
\fontsize{10.5}{13.2}\selectfont
\setlength{\LTleft}{0pt}\setlength{\LTright}{\fill}
\noindent
\begin{tabular}{@{}>{\raggedright\arraybackslash}p{39.459mm}>{\raggedright\arraybackslash}p{39.459mm}>{\raggedright\arraybackslash}p{39.459mm}@{}}
\toprule
\textbf{Person and number} & \textbf{Present} & \textbf{Preterite} \\
\midrule
1 sg. & bycge & bohte \\
2 sg. & bygest & bohtest \\
3 sg. & bygeð & bohte \\
1–3 pl. & bycgað & bohton \\
\bottomrule
\end{tabular}
\par\endgroup
% course-panel:T51-P01:end
% mdtable:T51:end
\reserve{7\baselineskip}
% mdparagraph:P183:begin
\subsection*{sēcan — seek}\BookHeading{nav-ch09-learner-0071}{subsubsection}{sēcan — seek}{0}
% mdparagraph:P183:end

% teaching: NEW-U09-VOC-secan
% mdparagraph:P184:begin
Weak I with preterite sōhte/\allowbreak{}sōhton. The thing sought is accusative. The present and preterite stems also supply the complete subjunctive below.
% mdparagraph:P184:end

% teaching: NEW-U09-PAR-secan-IND
\reserve{7\baselineskip}
% mdparagraph:P185:begin
\subsection*{sēcan: complete indicative}\BookHeading{nav-ch09-learner-0072}{subsubsection}{sēcan: complete indicative}{0}
% mdparagraph:P185:end

% table: NEW-U09-PAR-secan-IND
% mdtable:T52:begin
% course-table-style: 2026-10
% course-panel:T52-P01:begin
\par\begingroup
\fontsize{10.5}{13.2}\selectfont
\setlength{\LTleft}{0pt}\setlength{\LTright}{\fill}
\noindent
\begin{tabular}{@{}>{\raggedright\arraybackslash}p{39.459mm}>{\raggedright\arraybackslash}p{39.459mm}>{\raggedright\arraybackslash}p{39.459mm}@{}}
\toprule
\textbf{Person and number} & \textbf{Present} & \textbf{Preterite} \\
\midrule
1 sg. & sēce & sōhte \\
2 sg. & sēcst & sōhtest \\
3 sg. & sēcð & sōhte \\
1–3 pl. & sēcað & sōhton \\
\bottomrule
\end{tabular}
\par\endgroup
% course-panel:T52-P01:end
% mdtable:T52:end
\reserve{7\baselineskip}
% mdparagraph:P186:begin
\subsection*{gestillan — settle; quell}\BookHeading{nav-ch09-learner-0073}{subsubsection}{gestillan — settle; quell}{0}
% mdparagraph:P186:end

% teaching: NEW-U09-VOC-gestillan
% mdparagraph:P187:begin
Weak I. The disturbance or dispute quelled is accusative. The doubled l is reduced before the preterite dental ending: gestilde.
% mdparagraph:P187:end

% teaching: NEW-U09-PAR-gestillan-IND
\reserve{7\baselineskip}
% mdparagraph:P188:begin
\subsection*{gestillan: complete indicative}\BookHeading{nav-ch09-learner-0074}{subsubsection}{gestillan: complete indicative}{0}
% mdparagraph:P188:end

% table: NEW-U09-PAR-gestillan-IND
% mdtable:T53:begin
% course-table-style: 2026-10
% course-panel:T53-P01:begin
\par\begingroup
\fontsize{10.5}{13.2}\selectfont
\setlength{\LTleft}{0pt}\setlength{\LTright}{\fill}
\noindent
\begin{tabular}{@{}>{\raggedright\arraybackslash}p{39.459mm}>{\raggedright\arraybackslash}p{39.459mm}>{\raggedright\arraybackslash}p{39.459mm}@{}}
\toprule
\textbf{Person and number} & \textbf{Present} & \textbf{Preterite} \\
\midrule
1 sg. & gestille & gestilde \\
2 sg. & gestilst & gestildest \\
3 sg. & gestilð & gestilde \\
1–3 pl. & gestillað & gestildon \\
\bottomrule
\end{tabular}
\par\endgroup
% course-panel:T53-P01:end
% mdtable:T53:end
\reserve{7\baselineskip}
% mdparagraph:P189:begin
\subsection*{fręmman — benefit}\BookHeading{nav-ch09-learner-0075}{subsubsection}{fręmman — benefit}{0}\BookLabel{bk-ch09-learner-2a6c766e81e1}
% mdparagraph:P189:end

% teaching: NEW-U09-VOC-fremman
% teaching: NEW-U09-GRAM-u09-benefit-dative
% mdparagraph:P190:begin
Weak I. Keep doubled mm in fręmme/\allowbreak{}fręmmað; the second and third singular have fręmest/\allowbreak{}fręmeð, and the preterite has fręmede. The person benefited is dative.
% mdparagraph:P190:end

% teaching: NEW-U09-PAR-fremman-IND
\reserve{7\baselineskip}
% mdparagraph:P191:begin
\subsection*{fręmman: complete indicative}\BookHeading{nav-ch09-learner-0076}{subsubsection}{fręmman: complete indicative}{0}
% mdparagraph:P191:end

% table: NEW-U09-PAR-fremman-IND
% mdtable:T54:begin
% course-table-style: 2026-10
% course-panel:T54-P01:begin
\par\begingroup
\fontsize{10.5}{13.2}\selectfont
\setlength{\LTleft}{0pt}\setlength{\LTright}{\fill}
\noindent
\begin{tabular}{@{}>{\raggedright\arraybackslash}p{39.459mm}>{\raggedright\arraybackslash}p{39.459mm}>{\raggedright\arraybackslash}p{39.459mm}@{}}
\toprule
\textbf{Person and number} & \textbf{Present} & \textbf{Preterite} \\
\midrule
1 sg. & fręmme & fręmede \\
2 sg. & fręmest & fręmedest \\
3 sg. & fręmeð & fręmede \\
1–3 pl. & fręmmað & fręmedon \\
\bottomrule
\end{tabular}
\par\endgroup
% course-panel:T54-P01:end
% mdtable:T54:end
\reserve{7\baselineskip}
% mdparagraph:P192:begin
\subsection*{ālīehtan — illuminate; give light to}\BookHeading{nav-ch09-learner-0077}{subsubsection}{ālīehtan — illuminate; give light to}{0}
% mdparagraph:P192:end

% teaching: NEW-U09-VOC-aliehtan
% mdparagraph:P193:begin
Weak I. The thing illuminated is accusative. The selected second and third singular forms are uncontracted ālīehtest/\allowbreak{}ālīehteð. The preterite is ālīehte: the dental ending joins the stem’s ht without adding another t. Present first singular and preterite first/\allowbreak{}third singular therefore coincide; present and preterite second singular also coincide.
% mdparagraph:P193:end

% teaching: NEW-U09-PAR-aliehtan-IND
\reserve{7\baselineskip}
% mdparagraph:P194:begin
\subsection*{ālīehtan: complete indicative}\BookHeading{nav-ch09-learner-0078}{subsubsection}{ālīehtan: complete indicative}{0}
% mdparagraph:P194:end

% table: NEW-U09-PAR-aliehtan-IND
% mdtable:T55:begin
% course-table-style: 2026-10
% course-panel:T55-P01:begin
\par\begingroup
\fontsize{10.5}{13.2}\selectfont
\setlength{\LTleft}{0pt}\setlength{\LTright}{\fill}
\noindent
\begin{tabular}{@{}>{\raggedright\arraybackslash}p{39.459mm}>{\raggedright\arraybackslash}p{39.459mm}>{\raggedright\arraybackslash}p{39.459mm}@{}}
\toprule
\textbf{Person and number} & \textbf{Present} & \textbf{Preterite} \\
\midrule
1 sg. & ālīehte & ālīehte \\
2 sg. & ālīehtest & ālīehtest \\
3 sg. & ālīehteð & ālīehte \\
1–3 pl. & ālīehtað & ālīehton \\
\bottomrule
\end{tabular}
\par\endgroup
% course-panel:T55-P01:end
% mdtable:T55:end
\reserve{7\baselineskip}
% mdparagraph:P195:begin
\subsection*{rǣdan — read}\BookHeading{nav-ch09-learner-0079}{subsubsection}{rǣdan — read}{0}
% mdparagraph:P195:end

% teaching: NEW-U09-VOC-raedan
% mdparagraph:P196:begin
The selected preterite is weak: rǣdde/\allowbreak{}rǣddon. Present third singular rǣtt shows the joining of the stem dental and ending. Rǣdan can take an accusative text or stand without an expressed object, “read”. Strong preterites rēd/\allowbreak{}rēdon are documented alternatives, for recognition rather than an additional system to recall here.
% mdparagraph:P196:end

% teaching: NEW-U09-PAR-raedan-IND
\reserve{7\baselineskip}
% mdparagraph:P197:begin
\subsection*{rǣdan: complete indicative}\BookHeading{nav-ch09-learner-0080}{subsubsection}{rǣdan: complete indicative}{0}
% mdparagraph:P197:end

% table: NEW-U09-PAR-raedan-IND
% mdtable:T56:begin
% course-table-style: 2026-10
% course-panel:T56-P01:begin
\par\begingroup
\fontsize{10.5}{13.2}\selectfont
\setlength{\LTleft}{0pt}\setlength{\LTright}{\fill}
\noindent
\begin{tabular}{@{}>{\raggedright\arraybackslash}p{39.459mm}>{\raggedright\arraybackslash}p{39.459mm}>{\raggedright\arraybackslash}p{39.459mm}@{}}
\toprule
\textbf{Person and number} & \textbf{Present} & \textbf{Preterite} \\
\midrule
1 sg. & rǣde & rǣdde \\
2 sg. & rǣtst & rǣddest \\
3 sg. & rǣtt & rǣdde \\
1–3 pl. & rǣdað & rǣddon \\
\bottomrule
\end{tabular}
\par\endgroup
% course-panel:T56-P01:end
% mdtable:T56:end
\reserve{7\baselineskip}
% mdparagraph:P198:begin
\subsection*{Weak II}\BookHeading{nav-ch09-learner-0081}{subsubsection}{Weak II}{0}
% mdparagraph:P198:end

\reserve{7\baselineskip}
% mdparagraph:P199:begin
\subsection*{wician — sojourn; lodge}\BookHeading{nav-ch09-learner-0082}{subsubsection}{wician — sojourn; lodge}{0}
% mdparagraph:P199:end

% teaching: NEW-U09-VOC-wician
% mdparagraph:P200:begin
Weak II. The present has -ie, -ast, -að and plural -iað; the preterite has -ode/\allowbreak{}-odest/\allowbreak{}-odon. Use mid + dative for the person with whom one lodges. The plain infinitive wician is also used in the preference expression.
% mdparagraph:P200:end

% teaching: NEW-U09-PAR-wician-IND
\reserve{7\baselineskip}
% mdparagraph:P201:begin
\subsection*{wician: complete indicative}\BookHeading{nav-ch09-learner-0083}{subsubsection}{wician: complete indicative}{0}
% mdparagraph:P201:end

% table: NEW-U09-PAR-wician-IND
% mdtable:T57:begin
% course-table-style: 2026-10
% course-panel:T57-P01:begin
\par\begingroup
\fontsize{10.5}{13.2}\selectfont
\setlength{\LTleft}{0pt}\setlength{\LTright}{\fill}
\noindent
\begin{tabular}{@{}>{\raggedright\arraybackslash}p{39.459mm}>{\raggedright\arraybackslash}p{39.459mm}>{\raggedright\arraybackslash}p{39.459mm}@{}}
\toprule
\textbf{Person and number} & \textbf{Present} & \textbf{Preterite} \\
\midrule
1 sg. & wicie & wicode \\
2 sg. & wicast & wicodest \\
3 sg. & wicað & wicode \\
1–3 pl. & wiciað & wicodon \\
\bottomrule
\end{tabular}
\par\endgroup
% course-panel:T57-P01:end
% mdtable:T57:end
\reserve{7\baselineskip}
% mdparagraph:P202:begin
\subsection*{ðolian — suffer}\BookHeading{nav-ch09-learner-0084}{subsubsection}{ðolian — suffer}{0}
% mdparagraph:P202:end

% teaching: NEW-U09-VOC-tholian
% mdparagraph:P203:begin
Weak II, with the same endings as wician. The hardship or loss suffered is accusative.
% mdparagraph:P203:end

% teaching: NEW-U09-PAR-tholian-IND
\reserve{7\baselineskip}
% mdparagraph:P204:begin
\subsection*{ðolian: complete indicative}\BookHeading{nav-ch09-learner-0085}{subsubsection}{ðolian: complete indicative}{0}
% mdparagraph:P204:end

% table: NEW-U09-PAR-tholian-IND
% mdtable:T58:begin
% course-table-style: 2026-10
% course-panel:T58-P01:begin
\par\begingroup
\fontsize{10.5}{13.2}\selectfont
\setlength{\LTleft}{0pt}\setlength{\LTright}{\fill}
\noindent
\begin{tabular}{@{}>{\raggedright\arraybackslash}p{39.459mm}>{\raggedright\arraybackslash}p{39.459mm}>{\raggedright\arraybackslash}p{39.459mm}@{}}
\toprule
\textbf{Person and number} & \textbf{Present} & \textbf{Preterite} \\
\midrule
1 sg. & ðolie & ðolode \\
2 sg. & ðolast & ðolodest \\
3 sg. & ðolað & ðolode \\
1–3 pl. & ðoliað & ðolodon \\
\bottomrule
\end{tabular}
\par\endgroup
% course-panel:T58-P01:end
% mdtable:T58:end
\reserve{7\baselineskip}
% mdparagraph:P205:begin
\subsection*{gegaderian — gather}\BookHeading{nav-ch09-learner-0086}{subsubsection}{gegaderian — gather}{0}
% mdparagraph:P205:end

% teaching: NEW-U09-VOC-gegaderian
% mdparagraph:P206:begin
Weak II. The thing or things gathered are accusative. Keep the prefix ge- throughout the selected forms.
% mdparagraph:P206:end

% teaching: NEW-U09-PAR-gegaderian-IND
\reserve{7\baselineskip}
% mdparagraph:P207:begin
\subsection*{gegaderian: complete indicative}\BookHeading{nav-ch09-learner-0087}{subsubsection}{gegaderian: complete indicative}{0}
% mdparagraph:P207:end

% table: NEW-U09-PAR-gegaderian-IND
% mdtable:T59:begin
% course-table-style: 2026-10
% course-panel:T59-P01:begin
\par\begingroup
\fontsize{10.5}{13.2}\selectfont
\setlength{\LTleft}{0pt}\setlength{\LTright}{\fill}
\noindent
\begin{tabular}{@{}>{\raggedright\arraybackslash}p{39.459mm}>{\raggedright\arraybackslash}p{39.459mm}>{\raggedright\arraybackslash}p{39.459mm}@{}}
\toprule
\textbf{Person and number} & \textbf{Present} & \textbf{Preterite} \\
\midrule
1 sg. & gegaderie & gegaderode \\
2 sg. & gegaderast & gegaderodest \\
3 sg. & gegaderað & gegaderode \\
1–3 pl. & gegaderiað & gegaderodon \\
\bottomrule
\end{tabular}
\par\endgroup
% course-panel:T59-P01:end
% mdtable:T59:end
\reserve{7\baselineskip}
% mdparagraph:P208:begin
\subsection*{Strong verbs}\BookHeading{nav-ch09-learner-0088}{subsubsection}{Strong verbs}{0}
% mdparagraph:P208:end

\reserve{7\baselineskip}
% mdparagraph:P209:begin
\subsection*{rōwan — row}\BookHeading{nav-ch09-learner-0089}{subsubsection}{rōwan — row}{0}
% mdparagraph:P209:end

% teaching: NEW-U09-VOC-rowan
% mdparagraph:P210:begin
Strong. The second and third singular present change ō to ē: rēwst/\allowbreak{}rēwð. The preterite has rēow-. The verb can stand without an object; ofer introduces an expressed route across something.
% mdparagraph:P210:end

% teaching: NEW-U09-PAR-rowan-IND
\reserve{7\baselineskip}
% mdparagraph:P211:begin
\subsection*{rōwan: complete indicative}\BookHeading{nav-ch09-learner-0090}{subsubsection}{rōwan: complete indicative}{0}
% mdparagraph:P211:end

% table: NEW-U09-PAR-rowan-IND
% mdtable:T60:begin
% course-table-style: 2026-10
% course-panel:T60-P01:begin
\par\begingroup
\fontsize{10.5}{13.2}\selectfont
\setlength{\LTleft}{0pt}\setlength{\LTright}{\fill}
\noindent
\begin{tabular}{@{}>{\raggedright\arraybackslash}p{39.459mm}>{\raggedright\arraybackslash}p{39.459mm}>{\raggedright\arraybackslash}p{39.459mm}@{}}
\toprule
\textbf{Person and number} & \textbf{Present} & \textbf{Preterite} \\
\midrule
1 sg. & rōwe & rēow \\
2 sg. & rēwst & rēowe \\
3 sg. & rēwð & rēow \\
1–3 pl. & rōwað & rēowon \\
\bottomrule
\end{tabular}
\par\endgroup
% course-panel:T60-P01:end
% mdtable:T60:end
\reserve{7\baselineskip}
% mdparagraph:P212:begin
\subsection*{blāwan — blow}\BookHeading{nav-ch09-learner-0091}{subsubsection}{blāwan — blow}{0}
% mdparagraph:P212:end

% teaching: NEW-U09-VOC-blawan
% mdparagraph:P213:begin
Strong. The second and third singular present have ǣ; the preterite has blēow-. In the bellows sentence the verb stands without an expressed object.
% mdparagraph:P213:end

% teaching: NEW-U09-PAR-blawan-IND
\reserve{7\baselineskip}
% mdparagraph:P214:begin
\subsection*{blāwan: complete indicative}\BookHeading{nav-ch09-learner-0092}{subsubsection}{blāwan: complete indicative}{0}
% mdparagraph:P214:end

% table: NEW-U09-PAR-blawan-IND
% mdtable:T61:begin
% course-table-style: 2026-10
% course-panel:T61-P01:begin
\par\begingroup
\fontsize{10.5}{13.2}\selectfont
\setlength{\LTleft}{0pt}\setlength{\LTright}{\fill}
\noindent
\begin{tabular}{@{}>{\raggedright\arraybackslash}p{39.459mm}>{\raggedright\arraybackslash}p{39.459mm}>{\raggedright\arraybackslash}p{39.459mm}@{}}
\toprule
\textbf{Person and number} & \textbf{Present} & \textbf{Preterite} \\
\midrule
1 sg. & blāwe & blēow \\
2 sg. & blǣwst & blēowe \\
3 sg. & blǣwð & blēow \\
1–3 pl. & blāwað & blēowon \\
\bottomrule
\end{tabular}
\par\endgroup
% course-panel:T61-P01:end
% mdtable:T61:end
\reserve{7\baselineskip}
% mdparagraph:P215:begin
\subsection*{forlǣtan — abandon; leave}\BookHeading{nav-ch09-learner-0093}{subsubsection}{forlǣtan — abandon; leave}{0}
% mdparagraph:P215:end

% teaching: NEW-U09-VOC-forlaetan
% mdparagraph:P216:begin
Strong. The person or thing abandoned is accusative. The preterite stem is forlēt-. Present third singular forlǣtt joins the stem t and dental ending.
% mdparagraph:P216:end

% teaching: NEW-U09-PAR-forlaetan-IND
\reserve{7\baselineskip}
% mdparagraph:P217:begin
\subsection*{forlǣtan: complete indicative}\BookHeading{nav-ch09-learner-0094}{subsubsection}{forlǣtan: complete indicative}{0}
% mdparagraph:P217:end

% table: NEW-U09-PAR-forlaetan-IND
% mdtable:T62:begin
% course-table-style: 2026-10
% course-panel:T62-P01:begin
\par\begingroup
\fontsize{10.5}{13.2}\selectfont
\setlength{\LTleft}{0pt}\setlength{\LTright}{\fill}
\noindent
\begin{tabular}{@{}>{\raggedright\arraybackslash}p{39.459mm}>{\raggedright\arraybackslash}p{39.459mm}>{\raggedright\arraybackslash}p{39.459mm}@{}}
\toprule
\textbf{Person and number} & \textbf{Present} & \textbf{Preterite} \\
\midrule
1 sg. & forlǣte & forlēt \\
2 sg. & forlǣtst & forlēte \\
3 sg. & forlǣtt & forlēt \\
1–3 pl. & forlǣtað & forlēton \\
\bottomrule
\end{tabular}
\par\endgroup
% course-panel:T62-P01:end
% mdtable:T62:end
\reserve{7\baselineskip}
% mdparagraph:P218:begin
\subsection*{āstīgan — go aboard; board}\BookHeading{nav-ch09-learner-0095}{subsubsection}{āstīgan — go aboard; board}{0}
% mdparagraph:P218:end

% teaching: NEW-U09-VOC-astigan
% teaching: NEW-U09-GRAM-BOARD
% mdparagraph:P219:begin
Strong I. Present g becomes h before st/\allowbreak{}ð: āstīhst/\allowbreak{}āstīhð. The selected preterite singular is āstāg, second singular āstige, plural āstigon. Final āstāh is a documented alternative to āstāg. The ship boarded is a direct accusative object.
% mdparagraph:P219:end

% teaching: NEW-U09-PAR-astigan-IND
\reserve{7\baselineskip}
% mdparagraph:P220:begin
\subsection*{āstīgan: complete indicative}\BookHeading{nav-ch09-learner-0096}{subsubsection}{āstīgan: complete indicative}{0}
% mdparagraph:P220:end

% table: NEW-U09-PAR-astigan-IND
% mdtable:T63:begin
% course-table-style: 2026-10
% course-panel:T63-P01:begin
\par\begingroup
\fontsize{10.5}{13.2}\selectfont
\setlength{\LTleft}{0pt}\setlength{\LTright}{\fill}
\noindent
\begin{tabular}{@{}>{\raggedright\arraybackslash}p{39.459mm}>{\raggedright\arraybackslash}p{39.459mm}>{\raggedright\arraybackslash}p{39.459mm}@{}}
\toprule
\textbf{Person and number} & \textbf{Present} & \textbf{Preterite} \\
\midrule
1 sg. & āstīge & āstāg \\
2 sg. & āstīhst & āstige \\
3 sg. & āstīhð & āstāg \\
1–3 pl. & āstīgað & āstigon \\
\bottomrule
\end{tabular}
\par\endgroup
% course-panel:T63-P01:end
% mdtable:T63:end
\reserve{7\baselineskip}
% mdparagraph:P221:begin
\subsection*{bēatan — beat}\BookHeading{nav-ch09-learner-0097}{subsubsection}{bēatan — beat}{0}
% mdparagraph:P221:end

% teaching: NEW-U09-VOC-beatan
% mdparagraph:P222:begin
Strong. The person or thing beaten is accusative. Present second and third singular change ēa to īe; bīett joins the stem dental and ending. The preterite has bēot-.
% mdparagraph:P222:end

% teaching: NEW-U09-PAR-beatan-IND
\reserve{7\baselineskip}
% mdparagraph:P223:begin
\subsection*{bēatan: complete indicative}\BookHeading{nav-ch09-learner-0098}{subsubsection}{bēatan: complete indicative}{0}
% mdparagraph:P223:end

% table: NEW-U09-PAR-beatan-IND
% mdtable:T64:begin
% course-table-style: 2026-10
% course-panel:T64-P01:begin
\par\begingroup
\fontsize{10.5}{13.2}\selectfont
\setlength{\LTleft}{0pt}\setlength{\LTright}{\fill}
\noindent
\begin{tabular}{@{}>{\raggedright\arraybackslash}p{39.459mm}>{\raggedright\arraybackslash}p{39.459mm}>{\raggedright\arraybackslash}p{39.459mm}@{}}
\toprule
\textbf{Person and number} & \textbf{Present} & \textbf{Preterite} \\
\midrule
1 sg. & bēate & bēot \\
2 sg. & bīetst & bēote \\
3 sg. & bīett & bēot \\
1–3 pl. & bēatað & bēoton \\
\bottomrule
\end{tabular}
\par\endgroup
% course-panel:T64-P01:end
% mdtable:T64:end
\reserve{7\baselineskip}
% mdparagraph:P224:begin
\subsection*{The prefixed gān pattern}\BookHeading{nav-ch09-learner-0099}{subsubsection}{The prefixed gān pattern}{0}
% mdparagraph:P224:end

\reserve{7\baselineskip}
% mdparagraph:P225:begin
\subsection*{begān — ply; practise}\BookHeading{nav-ch09-learner-0100}{subsubsection}{begān — ply; practise}{0}
% mdparagraph:P225:end

% teaching: NEW-U09-VOC-began
% mdparagraph:P226:begin
Begān follows the earlier gān pattern with a prefix. The occupation practised is accusative. Keep begǣst/\allowbreak{}begǣð in the present and beēode/\allowbreak{}beēodest/\allowbreak{}beēodon in the preterite; the prefix does not turn the verb into a weak verb.
% mdparagraph:P226:end

% teaching: NEW-U09-PAR-began-IND
\reserve{7\baselineskip}
% mdparagraph:P227:begin
\subsection*{begān: complete indicative}\BookHeading{nav-ch09-learner-0101}{subsubsection}{begān: complete indicative}{0}
% mdparagraph:P227:end

% table: NEW-U09-PAR-began-IND
% mdtable:T65:begin
% course-table-style: 2026-10
% course-panel:T65-P01:begin
\par\begingroup
\fontsize{10.5}{13.2}\selectfont
\setlength{\LTleft}{0pt}\setlength{\LTright}{\fill}
\noindent
\begin{tabular}{@{}>{\raggedright\arraybackslash}p{39.459mm}>{\raggedright\arraybackslash}p{39.459mm}>{\raggedright\arraybackslash}p{39.459mm}@{}}
\toprule
\textbf{Person and number} & \textbf{Present} & \textbf{Preterite} \\
\midrule
1 sg. & begā & beēode \\
2 sg. & begǣst & beēodest \\
3 sg. & begǣð & beēode \\
1–3 pl. & begāð & beēodon \\
\bottomrule
\end{tabular}
\par\endgroup
% course-panel:T65-P01:end
% mdtable:T65:end
\reserve{7\baselineskip}
% mdparagraph:P228:begin
\subsection*{Commands}\BookHeading{nav-ch09-learner-0102}{subsubsection}{Commands}{0}
% mdparagraph:P228:end

% teaching: NEW-U09-GRAM-u09-new-commands
% mdparagraph:P229:begin
The singular command addresses one person; the plural command addresses more than one. The tables give both forms of these six verbs. Their plural commands have the same spelling as the present indicative plurals, but their use in an instruction is different. The singular commands rōw, byge, sēc, begā, gestil and rǣd are shown with their plural counterparts.
% mdparagraph:P229:end

% teaching: NEW-U09-PAR-rowan-IMP
\reserve{7\baselineskip}
% mdparagraph:P230:begin
\subsection*{rōwan: both commands}\BookHeading{nav-ch09-learner-0103}{subsubsection}{rōwan: both commands}{0}
% mdparagraph:P230:end

% table: NEW-U09-PAR-rowan-IMP
% mdtable:T66:begin
% course-table-style: 2026-10
% course-panel:T66-P01:begin
\par\begingroup
\fontsize{10.5}{13.2}\selectfont
\setlength{\LTleft}{0pt}\setlength{\LTright}{\fill}
\noindent
\begin{tabular}{@{}>{\raggedright\arraybackslash}p{63.342mm}>{\raggedright\arraybackslash}p{57.846mm}@{}}
\toprule
\textbf{Person and number} & \textbf{Command} \\
\midrule
2 sg. & rōw \\
2 pl. & rōwað \\
\bottomrule
\end{tabular}
\par\endgroup
% course-panel:T66-P01:end
% mdtable:T66:end
% teaching: NEW-U09-PAR-bycgan-IMP
\reserve{7\baselineskip}
% mdparagraph:P231:begin
\subsection*{bycgan: both commands}\BookHeading{nav-ch09-learner-0104}{subsubsection}{bycgan: both commands}{0}
% mdparagraph:P231:end

% table: NEW-U09-PAR-bycgan-IMP
% mdtable:T67:begin
% course-table-style: 2026-10
% course-panel:T67-P01:begin
\par\begingroup
\fontsize{10.5}{13.2}\selectfont
\setlength{\LTleft}{0pt}\setlength{\LTright}{\fill}
\noindent
\begin{tabular}{@{}>{\raggedright\arraybackslash}p{63.342mm}>{\raggedright\arraybackslash}p{57.846mm}@{}}
\toprule
\textbf{Person and number} & \textbf{Command} \\
\midrule
2 sg. & byge \\
2 pl. & bycgað \\
\bottomrule
\end{tabular}
\par\endgroup
% course-panel:T67-P01:end
% mdtable:T67:end
% teaching: NEW-U09-PAR-secan-IMP
\reserve{7\baselineskip}
% mdparagraph:P232:begin
\subsection*{sēcan: both commands}\BookHeading{nav-ch09-learner-0105}{subsubsection}{sēcan: both commands}{0}
% mdparagraph:P232:end

% table: NEW-U09-PAR-secan-IMP
% mdtable:T68:begin
% course-table-style: 2026-10
% course-panel:T68-P01:begin
\par\begingroup
\fontsize{10.5}{13.2}\selectfont
\setlength{\LTleft}{0pt}\setlength{\LTright}{\fill}
\noindent
\begin{tabular}{@{}>{\raggedright\arraybackslash}p{63.342mm}>{\raggedright\arraybackslash}p{57.846mm}@{}}
\toprule
\textbf{Person and number} & \textbf{Command} \\
\midrule
2 sg. & sēc \\
2 pl. & sēcað \\
\bottomrule
\end{tabular}
\par\endgroup
% course-panel:T68-P01:end
% mdtable:T68:end
% teaching: NEW-U09-PAR-began-IMP
\reserve{7\baselineskip}
% mdparagraph:P233:begin
\subsection*{begān: both commands}\BookHeading{nav-ch09-learner-0106}{subsubsection}{begān: both commands}{0}
% mdparagraph:P233:end

% table: NEW-U09-PAR-began-IMP
% mdtable:T69:begin
% course-table-style: 2026-10
% course-panel:T69-P01:begin
\par\begingroup
\fontsize{10.5}{13.2}\selectfont
\setlength{\LTleft}{0pt}\setlength{\LTright}{\fill}
\noindent
\begin{tabular}{@{}>{\raggedright\arraybackslash}p{63.342mm}>{\raggedright\arraybackslash}p{57.846mm}@{}}
\toprule
\textbf{Person and number} & \textbf{Command} \\
\midrule
2 sg. & begā \\
2 pl. & begāð \\
\bottomrule
\end{tabular}
\par\endgroup
% course-panel:T69-P01:end
% mdtable:T69:end
% teaching: NEW-U09-PAR-gestillan-IMP
\reserve{7\baselineskip}
% mdparagraph:P234:begin
\subsection*{gestillan: both commands}\BookHeading{nav-ch09-learner-0107}{subsubsection}{gestillan: both commands}{0}
% mdparagraph:P234:end

% table: NEW-U09-PAR-gestillan-IMP
% mdtable:T70:begin
% course-table-style: 2026-10
% course-panel:T70-P01:begin
\par\begingroup
\fontsize{10.5}{13.2}\selectfont
\setlength{\LTleft}{0pt}\setlength{\LTright}{\fill}
\noindent
\begin{tabular}{@{}>{\raggedright\arraybackslash}p{63.342mm}>{\raggedright\arraybackslash}p{57.846mm}@{}}
\toprule
\textbf{Person and number} & \textbf{Command} \\
\midrule
2 sg. & gestil \\
2 pl. & gestillað \\
\bottomrule
\end{tabular}
\par\endgroup
% course-panel:T70-P01:end
% mdtable:T70:end
% teaching: NEW-U09-PAR-raedan-IMP
\reserve{7\baselineskip}
% mdparagraph:P235:begin
\subsection*{rǣdan: both commands}\BookHeading{nav-ch09-learner-0108}{subsubsection}{rǣdan: both commands}{0}
% mdparagraph:P235:end

% table: NEW-U09-PAR-raedan-IMP
% mdtable:T71:begin
% course-table-style: 2026-10
% course-panel:T71-P01:begin
\par\begingroup
\fontsize{10.5}{13.2}\selectfont
\setlength{\LTleft}{0pt}\setlength{\LTright}{\fill}
\noindent
\begin{tabular}{@{}>{\raggedright\arraybackslash}p{63.342mm}>{\raggedright\arraybackslash}p{57.846mm}@{}}
\toprule
\textbf{Person and number} & \textbf{Command} \\
\midrule
2 sg. & rǣd \\
2 pl. & rǣdað \\
\bottomrule
\end{tabular}
\par\endgroup
% course-panel:T71-P01:end
% mdtable:T71:end
\reserve{7\baselineskip}
% mdparagraph:P236:begin
\subsection*{Sēcan: subjunctive}\BookHeading{nav-ch09-learner-0109}{subsubsection}{Sēcan: subjunctive}{0}
% mdparagraph:P236:end

% teaching: NEW-U09-GRAM-u09-secan-subjunctive
% mdparagraph:P237:begin
The present subjunctive uses sēc-, the preterite sōht-. Each has singular -e and plural -en for all three persons. Use the present first-person plural in a “let us seek” exhortation. These forms can also be used in the subjunctive constructions explained above.
% mdparagraph:P237:end

% teaching: NEW-U09-PAR-secan-SUBJ
\reserve{7\baselineskip}
% mdparagraph:P238:begin
\subsection*{sēcan: complete subjunctive}\BookHeading{nav-ch09-learner-0110}{subsubsection}{sēcan: complete subjunctive}{0}
% mdparagraph:P238:end

% table: NEW-U09-PAR-secan-SUBJ
% mdtable:T72:begin
% course-table-style: 2026-10
% course-panel:T72-P01:begin
\par\begingroup
\fontsize{10.5}{13.2}\selectfont
\setlength{\LTleft}{0pt}\setlength{\LTright}{\fill}
\noindent
\begin{tabular}{@{}>{\raggedright\arraybackslash}p{39.459mm}>{\raggedright\arraybackslash}p{39.459mm}>{\raggedright\arraybackslash}p{39.459mm}@{}}
\toprule
\textbf{Person and number} & \textbf{Present} & \textbf{Preterite} \\
\midrule
1 sg. & sēce & sōhte \\
2 sg. & sēce & sōhte \\
3 sg. & sēce & sōhte \\
1–3 pl. & sēcen & sōhten \\
\bottomrule
\end{tabular}
\par\endgroup
% course-panel:T72-P01:end
% mdtable:T72:end
\reserve{7\baselineskip}
% mdparagraph:P239:begin
\section*{8. Further uses and small words}\BookHeading{nav-ch09-learner-0111}{subsection}{8. Further uses and small words}{1}\BookLabel{bk-ch09-learner-a6e022a0fd7e}
% mdparagraph:P239:end

% teaching: NEW-U09-GRAM-u09-absolute-striking
% teaching: NEW-U09-VOC-slean
% mdparagraph:P240:begin
\textbf{Slēan}, “strike, smite”, can stand without naming an object when it means “make a stroke”. Its earlier use with an accusative object remains available; the absence of an object here does not require a new conjugation.
% mdparagraph:P240:end

% teaching: NEW-U09-GRAM-u09-bring-forth
% teaching: NEW-U09-VOC-teon
% teaching: NEW-U09-VOC-forth
% mdparagraph:P241:begin
Use \textbf{tēon … forð}, “bring forth, produce”, with an accusative thing produced. \textbf{Forð} is an invariant adverb. The familiar plain infinitive tēon can follow commanding hātan; the present subjunctive can express a wish or exhortation. Keep the actor of the action distinct from what it produces.
% mdparagraph:P241:end

% teaching: NEW-U09-GRAM-u09-temporal-on
% teaching: NEW-U09-VOC-on
% mdparagraph:P242:begin
\textbf{On + accusative} places an event on a numbered day. The article and ordinal agree with the accusative day noun. This time expression differs from ordinary \textbf{on + dative} locating something in a place.
% mdparagraph:P242:end

% teaching: NEW-U09-GRAM-u09-causal-for-tham
% teaching: NEW-U09-VOC-for-tham
% teaching: NEW-U09-VOC-hwilum
% teaching: NEW-U09-VOC-HWILUM
% mdparagraph:P243:begin
\textbf{For-ðām}, “because”, introduces a factual reason; use the indicative in the reason clause. \textbf{Hwīlum}, “sometimes”, is an invariant adverb of frequency. Its ending does not require you to supply a noun case construction.
% mdparagraph:P243:end

% teaching: NEW-U09-GRAM-SERVICE
% mdparagraph:P244:begin
\textbf{Godes ðēowdōm} means “service of God”: God is the one served. A genitive relationship need not be simple ownership. This expression uses the new service sense of ðēowdōm.
% mdparagraph:P244:end

\reserve{7\baselineskip}
% mdparagraph:P245:begin
\section*{9. Practice in forms and constructions}\BookHeading{nav-ch09-learner-0112}{subsection}{9. Practice in forms and constructions}{1}
% mdparagraph:P245:end

% mdparagraph:P246:begin
All exercises labelled with a letter are new teaching material. Preserve vowel length and the selected forms in recall tasks. In translation, preserve the stated meaning and requested construction. When a task specifies a mood, a grammatically possible sentence using another mood may still fail that particular target.
% mdparagraph:P246:end

\reserve{7\baselineskip}
% mdparagraph:P247:begin
\subsection*{Forms and case use}\BookHeading{nav-ch09-learner-0113}{subsubsection}{Forms and case use}{0}
% mdparagraph:P247:end

\reserve{5\baselineskip}
\par\begin{minipage}{\linewidth}\raggedright
% mdparagraph:P248:begin
\textbf{F01.}
% mdparagraph:P248:end

% item: NEW-U09-F01
% mdparagraph:P249:begin
Give the nominative, genitive, dative and accusative singular and plural of sunu. Distinguish the dative singular from the nominative and accusative singular.
% mdparagraph:P249:end

\end{minipage}\par\addvspace{4pt}

\par\begin{minipage}{\linewidth}\raggedright
% mdparagraph:P250:begin
\textbf{F02.}
% mdparagraph:P250:end

% item: NEW-U09-F02
% mdparagraph:P251:begin
Give the complete present and preterite indicative of bycgan. Name all three persons in each number, including cells with identical spellings.
% mdparagraph:P251:end

\end{minipage}\par\addvspace{4pt}

\par\begin{minipage}{\linewidth}\raggedright
% mdparagraph:P252:begin
\textbf{F03.}
% mdparagraph:P252:end

% item: NEW-U09-F03
% mdparagraph:P253:begin
Give the singular and plural commands of rōwan and bycgan. State which forms address one person and which address more than one.
% mdparagraph:P253:end

\end{minipage}\par\addvspace{4pt}

\par\begin{minipage}{\linewidth}\raggedright
% mdparagraph:P254:begin
\textbf{F04.}
% mdparagraph:P254:end

% item: NEW-U09-F04
% mdparagraph:P255:begin
Give the nominative, genitive, dative and accusative singular and plural of rīce and geðeahtend. Identify the characteristic plural endings of each pattern.
% mdparagraph:P255:end

\end{minipage}\par\addvspace{4pt}

\par\begin{minipage}{\linewidth}\raggedright
% mdparagraph:P256:begin
\textbf{F05.}
% mdparagraph:P256:end

% item: NEW-U09-F05
% mdparagraph:P257:begin
Translate the nominative noun phrase ‘the kingdoms of the counsellors’, placing the genitive dependent first and omitting a separate article for the head noun. Then translate ‘I give the fish (sing.) to the counsellors.’ Explain the different cases of ‘counsellors’.
% mdparagraph:P257:end

\end{minipage}\par\addvspace{4pt}

\reserve{7\baselineskip}
% mdparagraph:P258:begin
\subsection*{Content and report}\BookHeading{nav-ch09-learner-0114}{subsubsection}{Content and report}{0}
% mdparagraph:P258:end

\reserve{5\baselineskip}
\par\begin{minipage}{\linewidth}\raggedright
% mdparagraph:P259:begin
\textbf{A01.}
% mdparagraph:P259:end

% item: NEW-U09-A01
% mdparagraph:P260:begin
Classify every ðæt in these sentences as demonstrative/\allowbreak{}article, relative, or content conjunction. For each inflected occurrence, give its case and explain whether it agrees with a noun or functions within the relative clause. (a) Ðæt scip is gōd. (b) Ic wāt, ðæt ðæt scip gōd is. (c) Ðæt scip, ðæt ðū hæfst, is gōd.
% mdparagraph:P260:end

\end{minipage}\par\addvspace{4pt}

\par\begin{minipage}{\linewidth}\raggedright
% mdparagraph:P261:begin
\textbf{A02.}
% mdparagraph:P261:end

% item: NEW-U09-A02
% mdparagraph:P262:begin
Translate: I know that the ships are empty. Present their emptiness as an asserted fact and use the indicative in the content clause.
% mdparagraph:P262:end

\end{minipage}\par\addvspace{4pt}

\par\begin{minipage}{\linewidth}\raggedright
% mdparagraph:P263:begin
\textbf{A03.}
% mdparagraph:P263:end

% item: NEW-U09-A03
% mdparagraph:P264:begin
Se fiscere cwið: “Ic hæbbe twēgen fiscas.” Report his words beginning Se fiscere sægð and using ðæt with the present subjunctive. The fisherman is the possessor in both versions.
% mdparagraph:P264:end

\end{minipage}\par\addvspace{4pt}

\par\begin{minipage}{\linewidth}\raggedright
% mdparagraph:P265:begin
\textbf{A04.}
% mdparagraph:P265:end

% item: NEW-U09-A04
% mdparagraph:P266:begin
Se mann cwæð: “Ic hæbbe twēgen fiscas.” Report the same state at the time of his past speech, beginning Se mann sægde. Use ðæt and the preterite subjunctive; the same man possessed the fish.
% mdparagraph:P266:end

\end{minipage}\par\addvspace{4pt}

\par\begin{minipage}{\linewidth}\raggedright
% mdparagraph:P267:begin
\textbf{A05.}
% mdparagraph:P267:end

% item: NEW-U09-A05
% mdparagraph:P268:begin
A task requests a subjunctive report, but a learner writes Hē sægð, ðæt hē gōd is. Is the sentence necessarily ungrammatical? Give the requested form and explain the distinction.
% mdparagraph:P268:end

\end{minipage}\par\addvspace{4pt}

\reserve{7\baselineskip}
% mdparagraph:P269:begin
\subsection*{Dependent questions}\BookHeading{nav-ch09-learner-0115}{subsubsection}{Dependent questions}{0}
% mdparagraph:P269:end

\reserve{5\baselineskip}
\par\begin{minipage}{\linewidth}\raggedright
% mdparagraph:P270:begin
\textbf{Q01.}
% mdparagraph:P270:end

% item: NEW-U09-Q01
% mdparagraph:P271:begin
Translate both clauses and explain the case of the question word. (a) Ic wāt, hwā ðone fisc hæbbe. (b) Ic wāt, hwǣm ðū ðone fisc sęlle.
% mdparagraph:P271:end

\end{minipage}\par\addvspace{4pt}

\par\begin{minipage}{\linewidth}\raggedright
% mdparagraph:P272:begin
\textbf{Q02.}
% mdparagraph:P272:end

% item: NEW-U09-Q02
% mdparagraph:P273:begin
Turn Hwæt hæfst ðū? into a dependent question after Saga mē. Keep the same addressee and use the present subjunctive.
% mdparagraph:P273:end

\end{minipage}\par\addvspace{4pt}

\par\begin{minipage}{\linewidth}\raggedright
% mdparagraph:P274:begin
\textbf{Q03.}
% mdparagraph:P274:end

% item: NEW-U09-Q03
% mdparagraph:P275:begin
Turn Hæfð se mann twēgen fiscas? into a dependent yes/\allowbreak{}no question after Saga mē. Use hwæðer and the present subjunctive.
% mdparagraph:P275:end

\end{minipage}\par\addvspace{4pt}

\par\begin{minipage}{\linewidth}\raggedright
% mdparagraph:P276:begin
\textbf{Q04.}
% mdparagraph:P276:end

% item: NEW-U09-Q04
% mdparagraph:P277:begin
Translate and distinguish the two uses of hwæðer. In (a), two men are under discussion. (a) Hwæðer mann is ðīn gefēra? (b) Saga mē, hwæðer se mann ðīn gefēra sīe.
% mdparagraph:P277:end

\end{minipage}\par\addvspace{4pt}

\par\begin{minipage}{\linewidth}\raggedright
% mdparagraph:P278:begin
\textbf{Q05.}
% mdparagraph:P278:end

% item: NEW-U09-Q05
% mdparagraph:P279:begin
Translate: I know whom you (sing.) help. Use a dependent question with hwǣm and the present subjunctive.
% mdparagraph:P279:end

\end{minipage}\par\addvspace{4pt}

\reserve{7\baselineskip}
% mdparagraph:P280:begin
\subsection*{Infinitives and object predicates}\BookHeading{nav-ch09-learner-0116}{subsubsection}{Infinitives and object predicates}{0}
% mdparagraph:P280:end

\reserve{5\baselineskip}
\par\begin{minipage}{\linewidth}\raggedright
% mdparagraph:P281:begin
\textbf{I01.}
% mdparagraph:P281:end

% item: NEW-U09-I01
% mdparagraph:P282:begin
Translate both sentences and identify who carries the fish. (a) Ic wille ðone fisc beran. (b) Ic hāte hine ðone fisc beran. Explain why beran has no person ending in either sentence.
% mdparagraph:P282:end

\end{minipage}\par\addvspace{4pt}

\par\begin{minipage}{\linewidth}\raggedright
% mdparagraph:P283:begin
\textbf{I02.}
% mdparagraph:P283:end

% item: NEW-U09-I02
% mdparagraph:P284:begin
The fisherman told two companions, “Berað ðone fisc.” Report this as a past command using hātan, hīe for the companions, and a plain infinitive.
% mdparagraph:P284:end

\end{minipage}\par\addvspace{4pt}

\par\begin{minipage}{\linewidth}\raggedright
% mdparagraph:P285:begin
\textbf{I03.}
% mdparagraph:P285:end

% item: NEW-U09-I03
% mdparagraph:P286:begin
Translate: I see the men carrying the fish (sing.). Use gesēon with an accusative actor and a plain infinitive; label both accusative noun phrases.
% mdparagraph:P286:end

\end{minipage}\par\addvspace{4pt}

\par\begin{minipage}{\linewidth}\raggedright
% mdparagraph:P287:begin
\textbf{I04.}
% mdparagraph:P287:end

% item: NEW-U09-I04
% mdparagraph:P288:begin
Translate both sentences using gewyrcean followed by an object and an agreeing predicate adjective: God made the man good. God made the men good. Explain the adjective endings.
% mdparagraph:P288:end

\end{minipage}\par\addvspace{4pt}

\par\begin{minipage}{\linewidth}\raggedright
% mdparagraph:P289:begin
\textbf{I05.}
% mdparagraph:P289:end

% item: NEW-U09-I05
% mdparagraph:P290:begin
Translate and distinguish the two complements: Hē hēt ðone mann cyning. Hē hēt ðone mann gān. State the case/\allowbreak{}function of cyning and the form/\allowbreak{}function of gān.
% mdparagraph:P290:end

\end{minipage}\par\addvspace{4pt}

\reserve{7\baselineskip}
% mdparagraph:P291:begin
\subsection*{Unreal wishes}\BookHeading{nav-ch09-learner-0117}{subsubsection}{Unreal wishes}{0}
% mdparagraph:P291:end

\reserve{5\baselineskip}
\par\begin{minipage}{\linewidth}\raggedright
% mdparagraph:P292:begin
\textbf{W01.}
% mdparagraph:P292:end

% item: NEW-U09-W01
% mdparagraph:P293:begin
You have no ship. Express the present unreal wish “Would that I had one ship!” using the preterite subjunctive of habban.
% mdparagraph:P293:end

\end{minipage}\par\addvspace{4pt}

\par\begin{minipage}{\linewidth}\raggedright
% mdparagraph:P294:begin
\textbf{W02.}
% mdparagraph:P294:end

% item: NEW-U09-W02
% mdparagraph:P295:begin
Translate these in their stated settings: (a) a report about last year, Ic hæfde ān scip.; (b) a wish by someone who has no ship now, Hæfde ic ān scip! Explain what the shared spelling hæfde can and cannot show.
% mdparagraph:P295:end

\end{minipage}\par\addvspace{4pt}

\reserve{7\baselineskip}
% mdparagraph:P296:begin
\section*{10. The fisherman and his companion}\BookHeading{nav-ch09-learner-0118}{subsection}{10. The fisherman and his companion}{1}
% mdparagraph:P296:end

% mdparagraph:P297:begin
New teaching material. Read the connected exchange, then answer the questions below.
% mdparagraph:P297:end

% composition: NEW-U09-READING-01
% mdparagraph:P298:begin
\par\begin{minipage}{\linewidth}\raggedright
\readingtext{Se fiscere and his gefēra standað on ðǣm strande. Se gefēra cwið: “Saga mē, hwæt ðū hæbbe.” Se fiscere sægð, ðæt hē męte hæbbe. Hē hǣtt his gefēran ðone fisc beran. Se gefēra birð ðone fisc. Se fiscere cwið: “Ic gesēo ðē ðone fisc beran. Nū wāt ic, ðæt ðū mē hilpst.”}
\end{minipage}\par
% mdparagraph:P298:end

\reserve{7\baselineskip}
% mdparagraph:P299:begin
\subsection*{Reading practice}\BookHeading{nav-ch09-learner-0119}{subsubsection}{Reading practice}{0}
% mdparagraph:P299:end

\reserve{5\baselineskip}
\par\begin{minipage}{\linewidth}\raggedright
% mdparagraph:P300:begin
\textbf{R01.}
% mdparagraph:P300:end

% item: NEW-U09-R01
% mdparagraph:P301:begin
In the controlled reading, who requests information, and who is reported to have food? Answer in English and cite the two relevant clauses.
% mdparagraph:P301:end

\end{minipage}\par\addvspace{4pt}

\par\begin{minipage}{\linewidth}\raggedright
% mdparagraph:P302:begin
\textbf{R02.}
% mdparagraph:P302:end

% item: NEW-U09-R02
% mdparagraph:P303:begin
In the controlled reading, compare hwæt ðū hæbbe, ðæt hē męte hæbbe, and ðæt ðū mē hilpst. Identify each clause type and explain the selected moods.
% mdparagraph:P303:end

\end{minipage}\par\addvspace{4pt}

\par\begin{minipage}{\linewidth}\raggedright
% mdparagraph:P304:begin
\textbf{R03.}
% mdparagraph:P304:end

% item: NEW-U09-R03
% mdparagraph:P305:begin
Identify the commanded actor in Hē hǣtt his gefēran ðone fisc beran and the observed actor in Ic gesēo ðē ðone fisc beran. For each sentence, distinguish that actor from the fish and identify the finite verb.
% mdparagraph:P305:end

\end{minipage}\par\addvspace{4pt}

\par\begin{minipage}{\linewidth}\raggedright
% mdparagraph:P306:begin
\textbf{R04.}
% mdparagraph:P306:end

% item: NEW-U09-R04
% mdparagraph:P307:begin
Translate the complete controlled reading, preserving the speakers, the command, the carrying action and the final distinction between report and asserted knowledge.
% mdparagraph:P307:end

\end{minipage}\par\addvspace{4pt}

\reserve{7\baselineskip}
% mdparagraph:P308:begin
\section*{11. The Choice of Occupations}\BookHeading{nav-ch09-learner-0120}{subsection}{11. The Choice of Occupations}{1}
% mdparagraph:P308:end

% mdparagraph:P309:begin
Exact source text from Cook’s \emph{A First Book in Old English}. The Teacher asks the Counsellor to rank occupations. Several expressions are supplied in the local helps for reading; their unfamiliar forms are not required for independent composition.
% mdparagraph:P309:end

\par\begin{minipage}{\linewidth}\raggedright
% source: FB-II-CHOICE-OCCUPATIONS-01
% mdparagraph:P310:begin
\readingtext{\textbf{Teacher.} Hwæt sægst ðū, wīsa? Hwilc cræft ðē is geðūht betweox ðās furðra wesan?}
% mdparagraph:P310:end

% mdparagraph:P311:begin
\readingtext{\textbf{Counsellor.} Ic sęcge ðē, mē is geðūht Godes ðēowdōm betweoh ðās cræftas ealdorscipe healdan, swā-swā hit is gerǣd on godspelle, “Fyrmest sēceað rīce Godes, and rihtwīsnesse his, and ðās ðing eall bēoð tōgēiecte ēow.”}
% mdparagraph:P311:end

% mdparagraph:P312:begin
\readingtext{\textbf{Teacher.} And hwilc ðē is geðūht betweox woruldcræftas healdan ealdordōm?}
% mdparagraph:P312:end

% mdparagraph:P313:begin
\readingtext{\textbf{Counsellor.} Eorðtilð, for-ðām se ierðling ūs ealle fēt.}
% mdparagraph:P313:end

\end{minipage}\par
\reserve{7\baselineskip}
% mdparagraph:P314:begin
\subsection*{Local help}\BookHeading{nav-ch09-learner-0121}{subsubsection}{Local help}{0}
% mdparagraph:P314:end

% teaching: NEW-U09-SOURCE-HELPS
% teaching: NEW-U09-GRAM-u09-source-seceath
% teaching: NEW-U09-VOC-thyncan
% teaching: NEW-U09-VOC-betweox
% teaching: NEW-U09-VOC-furthra
% teaching: NEW-U09-VOC-swa-swa
% teaching: NEW-U09-VOC-fyrmest
% teaching: NEW-U09-VOC-togeiecan
% teaching: NEW-U09-VOC-woruldcraeft
% teaching: NEW-U09-VOC-ealdordom
% teaching: NEW-U09-VOC-fedan
% support: U09-SRC-HELP01
% mdparagraph:P315:begin
\textbf{1. wīsa.} “Wise man; wise one”, in direct address. It is the weak masculine singular of wīs used as a noun. This personal use is given here for reading.
% mdparagraph:P315:end

% support: U09-SRC-HELP02
% mdparagraph:P316:begin
\textbf{2. ðē is geðūht … wesan /\allowbreak{} mē is geðūht … healdan.} “To you it seems … to be /\allowbreak{} to me it seems … to hold.” Read these as complete expressions. The person to whom it seems is dative. Cook identifies furðra as nominative singular; these source expressions do not follow one uniform accusative-actor pattern like commanding hātan.
% mdparagraph:P316:end

% support: U09-SRC-HELP03
% mdparagraph:P317:begin
\textbf{3. betweox /\allowbreak{} betweoh ðās (cræftas).} “Among these (occupations).” The two spellings mean the same thing here. The following phrase is accusative plural. In the first question, cræftas is understood after ðās.
% mdparagraph:P317:end

% support: U09-SRC-HELP04
% mdparagraph:P318:begin
\textbf{4. furðra.} “Foremost; prior.” This is a weak comparative, masculine nominative singular, as Cook’s note 14 states. This one form and sense are given for the passage.
% mdparagraph:P318:end

% support: U09-SRC-HELP05
% mdparagraph:P319:begin
\textbf{5. swā-swā hit is gerǣd on godspelle.} “Just as it is read in the Gospel.” Swā-swā means “just as”; gerǣd is the past participle in this passive expression. Read the whole phrase with this help.
% mdparagraph:P319:end

% support: U09-SRC-HELP06
% mdparagraph:P320:begin
\textbf{6. Fyrmest sēceað.} “First seek”, a plural command. Sēceað is the form printed in the source; recognize it beside the selected plural command sēcað.
% mdparagraph:P320:end

% support: U09-SRC-HELP07
% mdparagraph:P321:begin
\textbf{7. ðās ðing eall bēoð tōgēiecte ēow.} “All these things will be added to you.” Eall agrees with neuter plural ðing. Bēoð is third-person plural here; ēow is the dative recipient. Tōgēiecte is a predicate participle. Keep its printed -e: Cook’s §62 describes this predicate use, so it should not be silently changed to an expected neuter ending.
% mdparagraph:P321:end

% support: U09-SRC-HELP08
% mdparagraph:P322:begin
\textbf{8. woruldcræftas /\allowbreak{} ealdordōm.} “Worldly occupations” (accusative plural) /\allowbreak{} “chief place, primacy” (accusative singular). Cook prints woruldcræftas but records the manuscript words cræftas woruld in note 20.
% mdparagraph:P322:end

% support: U09-SRC-HELP09
% mdparagraph:P323:begin
\textbf{9. ūs ealle fēt.} “Feeds us all.” Ūs is accusative plural; ealle refers to the people fed and is masculine plural. Fēt is present indicative third singular of fēdan. This exact form is given for reading; the prefixed āfēdan table does not supply a complete system for fēdan.
% mdparagraph:P323:end

% support: U09-SRC-HELP10
% mdparagraph:P324:begin
\textbf{10. the three conjectural is insertions.} Cook calls the is in each Teacher question and in swā-swā hit is gerǣd conjectural: they are absent from his manuscripts. The is in mē is geðūht instead belongs to note 15. The printed text here retains Cook’s three insertions.
% mdparagraph:P324:end

% support: U09-SRC-HELP11
% mdparagraph:P325:begin
\textbf{11. Godes ðēowdōm; rihtwīsnesse his.} “God’s service, service of God”; “his righteousness”. The genitives identify the divine possessor or the one served. His follows its noun in the second phrase.
% mdparagraph:P325:end

\reserve{7\baselineskip}
% mdparagraph:P326:begin
\subsection*{Cook’s original notes}\BookHeading{nav-ch09-learner-0122}{subsubsection}{Cook’s original notes}{0}
% mdparagraph:P326:end

% mdparagraph:P327:begin
The numbers are Cook’s original note numbers. “Lat.” identifies his Latin comparison; “MSS.” means manuscripts. His section references point to the grammar in his \emph{First Book}.
% mdparagraph:P327:end

% mdtable:T73:begin
\par\begingroup
\fontsize{10.5}{13.2}\selectfont
\setlength{\LTleft}{0pt}\setlength{\LTright}{\fill}
\noindent
\begin{tabular}{@{}>{\raggedright\arraybackslash}p{22.72mm}>{\raggedright\arraybackslash}p{38.74mm}>{\raggedright\arraybackslash}p{56.92mm}@{}}
\toprule
\textbf{Note} & \textbf{Anchor} & \textbf{Original note} \\
\midrule
11 & cræft & Lat. ars. \\
12 & is (three occurrences) & Conjectural; not in the MSS. \\
13 & geðūht & See ðyncan. \\
14 & furðra & Lat. prior. Nom. sing. \\
15 & mē is geðūht & Lat. videtur. \\
16 & ðēowdōm & See 143 and 149. \\
17 & ealdorscipe; ealdordōm & Lat. primatum. \\
18 & rihtwīsnesse & See 144. \\
19 & tōgēiecte & Lat. adjicientur. See tōgēiecan, and 62. \\
20 & woruldcræftas & Lat. artes seculares. MS. cræftas woruld. \\
21 & Eorðtilð & Lat. agricultura. See 147. \\
22 & ierðling & Lat. arator. \\
23 & fēt & See fēdan. \\
\bottomrule
\end{tabular}
\par\endgroup

% mdtable:T73:end
\reserve{7\baselineskip}
% mdparagraph:P328:begin
\subsection*{Supported source practice}\BookHeading{nav-ch09-learner-0123}{subsubsection}{Supported source practice}{0}
% mdparagraph:P328:end

% mdparagraph:P329:begin
Use the text, vocabulary and local helps. These tasks ask you to understand the supported source; they do not require unaided recall of the unfamiliar source forms.
% mdparagraph:P329:end

\reserve{5\baselineskip}
\par\begin{minipage}{\linewidth}\raggedright
% mdparagraph:P330:begin
\textbf{S01.}
% mdparagraph:P330:end

% item: NEW-U09-S01
% mdparagraph:P331:begin
In the source dialogue, which occupation receives first place overall, and which receives first place among worldly occupations? Explain why the two answers do not contradict one another.
% mdparagraph:P331:end

\end{minipage}\par\addvspace{4pt}

\par\begin{minipage}{\linewidth}\raggedright
% mdparagraph:P332:begin
\textbf{S02.}
% mdparagraph:P332:end

% item: NEW-U09-S02
% mdparagraph:P333:begin
Identify the beginning and end of the Gospel quotation in the source dialogue. What does it command, and what does it promise? Use the printed helps for unfamiliar forms.
% mdparagraph:P333:end

\end{minipage}\par\addvspace{4pt}

\par\begin{minipage}{\linewidth}\raggedright
% mdparagraph:P334:begin
\textbf{S03.}
% mdparagraph:P334:end

% item: NEW-U09-S03
% mdparagraph:P335:begin
In Ic sęcge ðē, mē is geðūht, explain the different roles of ðē and mē using the local help. Who is addressed, and whose judgment follows?
% mdparagraph:P335:end

\end{minipage}\par\addvspace{4pt}

\par\begin{minipage}{\linewidth}\raggedright
% mdparagraph:P336:begin
\textbf{S04.}
% mdparagraph:P336:end

% item: NEW-U09-S04
% mdparagraph:P337:begin
What reason does the Counsellor give for agriculture's place? Explain the reference of ūs ealle in the final turn and identify the subject of fēt.
% mdparagraph:P337:end

\end{minipage}\par\addvspace{4pt}

\par\begin{minipage}{\linewidth}\raggedright
% mdparagraph:P338:begin
\textbf{S05.}
% mdparagraph:P338:end

% item: NEW-U09-S05
% mdparagraph:P339:begin
Translate all four source turns, including the complete Gospel quotation. Preserve the distinction between the overall ranking and the later question about worldly occupations.
% mdparagraph:P339:end

\end{minipage}\par\addvspace{4pt}

\reserve{7\baselineskip}
% mdparagraph:P340:begin
\section*{12. Cook II.9–12}\BookHeading{nav-ch09-learner-0124}{subsection}{12. Cook II.9–12}{1}
% mdparagraph:P340:end

% mdparagraph:P341:begin
Translate the complete English prompts into Old English. Cook’s numbering, parentheses, brackets and grammar references are retained. “Sing.” and “plur.” specify number; “om.” asks you to omit the English word indicated. Section numbers refer to Cook’s grammar. Use the local help printed with the three supported expressions.
% mdparagraph:P341:end

% mdparagraph:P342:begin
In a “let us” or “let …” exhortation, use the appropriate present subjunctive unless another construction is expressly requested. Direct address uses the nominative-shaped noun and a weak accompanying adjective. In the reported-speech item, follow the prompt as printed; a content conjunction has not been silently added.
% mdparagraph:P342:end

\reserve{7\baselineskip}
% mdparagraph:P343:begin
\subsection*{Exercise 9}\BookHeading{nav-ch09-learner-0125}{subsubsection}{Exercise 9}{0}
% mdparagraph:P343:end

\reserve{5\baselineskip}
\par\begin{minipage}{\linewidth}\raggedright
% mdparagraph:P344:begin
\textbf{Cook II.9.1.}
% mdparagraph:P344:end

% item: EX-II-09-01
% mdparagraph:P345:begin
Sparks fly.
% mdparagraph:P345:end

\end{minipage}\par\addvspace{4pt}

\par\begin{minipage}{\linewidth}\raggedright
% mdparagraph:P346:begin
\textbf{Cook II.9.2.}
% mdparagraph:P346:end

% item: EX-II-09-02
% mdparagraph:P347:begin
The house is empty.
% mdparagraph:P347:end

\end{minipage}\par\addvspace{4pt}

\par\begin{minipage}{\linewidth}\raggedright
% mdparagraph:P348:begin
\textbf{Cook II.9.3.}
% mdparagraph:P348:end

% item: EX-II-09-03
% mdparagraph:P349:begin
The ships are empty.
% mdparagraph:P349:end

\end{minipage}\par\addvspace{4pt}

\par\begin{minipage}{\linewidth}\raggedright
% mdparagraph:P350:begin
\textbf{Cook II.9.4.}
% mdparagraph:P350:end

% item: EX-II-09-04
% mdparagraph:P351:begin
The merchant sells needles, awls, fish-hooks, and plowshares.
% mdparagraph:P351:end

\end{minipage}\par\addvspace{4pt}

\par\begin{minipage}{\linewidth}\raggedright
% mdparagraph:P352:begin
\textbf{Cook II.9.5.}
% mdparagraph:P352:end

% item: EX-II-09-05
% mdparagraph:P353:begin
The blacksmith makes plowshares in his smithy.
% mdparagraph:P353:end

\end{minipage}\par\addvspace{4pt}

\par\begin{minipage}{\linewidth}\raggedright
% mdparagraph:P354:begin
\textbf{Cook II.9.6.}
% mdparagraph:P354:end

% item: EX-II-09-06
% mdparagraph:P355:begin
The oil and the wine are good.
% mdparagraph:P355:end

\end{minipage}\par\addvspace{4pt}

\par\begin{minipage}{\linewidth}\raggedright
% mdparagraph:P356:begin
\textbf{Cook II.9.7.}
% mdparagraph:P356:end

% item: EX-II-09-07
% mdparagraph:P357:begin
The fisherman rows over the sea.
% mdparagraph:P357:end

\end{minipage}\par\addvspace{4pt}

\par\begin{minipage}{\linewidth}\raggedright
% mdparagraph:P358:begin
\textbf{Cook II.9.8.}
% mdparagraph:P358:end

% item: EX-II-09-08
% mdparagraph:P359:begin
The carpenter supports his son.
% mdparagraph:P359:end

\end{minipage}\par\addvspace{4pt}

\par\begin{minipage}{\linewidth}\raggedright
% mdparagraph:P360:begin
\textbf{Cook II.9.9.}
% mdparagraph:P360:end

% item: EX-II-09-09
% mdparagraph:P361:begin
The tailor bought the needle.
% mdparagraph:P361:end

\end{minipage}\par\addvspace{4pt}

\par\begin{minipage}{\linewidth}\raggedright
% mdparagraph:P362:begin
\textbf{Cook II.9.10.}
% mdparagraph:P362:end

% item: EX-II-09-10
% mdparagraph:P363:begin
The bellows blow.
% mdparagraph:P363:end

\end{minipage}\par\addvspace{4pt}

\par\begin{minipage}{\linewidth}\raggedright
% mdparagraph:P364:begin
\textbf{Cook II.9.11.}
% mdparagraph:P364:end

% item: EX-II-09-11
% mdparagraph:P365:begin
The hammer smites,
% mdparagraph:P365:end

\end{minipage}\par\addvspace{4pt}

\par\begin{minipage}{\linewidth}\raggedright
% mdparagraph:P366:begin
\textbf{Cook II.9.12.}
% mdparagraph:P366:end

% item: EX-II-09-12
% mdparagraph:P367:begin
Let us seek the kingdom of God.
% mdparagraph:P367:end

\end{minipage}\par\addvspace{4pt}

\par\begin{minipage}{\linewidth}\raggedright
% mdparagraph:P368:begin
\textbf{Cook II.9.13.}
% mdparagraph:P368:end

% item: EX-II-09-13
% mdparagraph:P369:begin
He will not live with the blacksmith.
% mdparagraph:P369:end

\end{minipage}\par\addvspace{4pt}

\par\begin{minipage}{\linewidth}\raggedright
% mdparagraph:P370:begin
\textbf{Cook II.9.14.}
% mdparagraph:P370:end

% item: EX-II-09-14
% mdparagraph:P371:begin
We had rather live [To us it is preferable to sojourn] with the farmer than with the tailor.
% mdparagraph:P371:end

\end{minipage}\par\addvspace{4pt}

\par\begin{minipage}{\linewidth}\raggedright
% mdparagraph:P372:begin
\textbf{Cook II.9.15.}
% mdparagraph:P372:end

% item: EX-II-09-15
% mdparagraph:P373:begin
Gold is better (66) than tin.
% mdparagraph:P373:end

\end{minipage}\par\addvspace{4pt}

\par\begin{minipage}{\linewidth}\raggedright
% mdparagraph:P374:begin
\textbf{Cook II.9.16.}
% mdparagraph:P374:end

% item: EX-II-09-16
% mdparagraph:P375:begin
The carpenter makes a hole.
% mdparagraph:P375:end

\end{minipage}\par\addvspace{4pt}

\reserve{7\baselineskip}
% mdparagraph:P376:begin
\subsection*{Exercise 10}\BookHeading{nav-ch09-learner-0126}{subsubsection}{Exercise 10}{0}
% mdparagraph:P376:end

\reserve{5\baselineskip}
\par\begin{minipage}{\linewidth}\raggedright
% mdparagraph:P377:begin
\textbf{Cook II.10.1.}
% mdparagraph:P377:end

% item: EX-II-10-01
% mdparagraph:P378:begin
The tree is useful to the carpenter.
% mdparagraph:P378:end

\end{minipage}\par\addvspace{4pt}

\par\begin{minipage}{\linewidth}\raggedright
% mdparagraph:P379:begin
\textbf{Cook II.10.2.}
% mdparagraph:P379:end

% item: EX-II-10-02
% mdparagraph:P380:begin
Our horses have fodder.
% mdparagraph:P380:end

\end{minipage}\par\addvspace{4pt}

\par\begin{minipage}{\linewidth}\raggedright
% mdparagraph:P381:begin
\textbf{Cook II.10.3.}
% mdparagraph:P381:end

% item: EX-II-10-03
% mdparagraph:P382:begin
O priest, read the Gospel.
% mdparagraph:P382:end

\end{minipage}\par\addvspace{4pt}

\par\begin{minipage}{\linewidth}\raggedright
% mdparagraph:P383:begin
\textbf{Cook II.10.4.}
% mdparagraph:P383:end

% item: EX-II-10-04
% mdparagraph:P384:begin
Buy (sing.) wine and oil.
% mdparagraph:P384:end

\end{minipage}\par\addvspace{4pt}

\par\begin{minipage}{\linewidth}\raggedright
% mdparagraph:P385:begin
\textbf{Cook II.10.5.}
% mdparagraph:P385:end

% item: EX-II-10-05
% mdparagraph:P386:begin
O good fellows, row over the waters.
% mdparagraph:P386:end

\end{minipage}\par\addvspace{4pt}

\par\begin{minipage}{\linewidth}\raggedright
% mdparagraph:P387:begin
\textbf{Cook II.10.6.}
% mdparagraph:P387:end

% item: EX-II-10-06
% mdparagraph:P388:begin
Her son is rich.
% mdparagraph:P388:end

\end{minipage}\par\addvspace{4pt}

\par\begin{minipage}{\linewidth}\raggedright
% mdparagraph:P389:begin
\textbf{Cook II.10.7.}
% mdparagraph:P389:end

% item: EX-II-10-07
% mdparagraph:P390:begin
Ply (sing.) your trade.
% mdparagraph:P390:end

\end{minipage}\par\addvspace{4pt}

\par\begin{minipage}{\linewidth}\raggedright
% mdparagraph:P391:begin
\textbf{Cook II.10.8.}
% mdparagraph:P391:end

% item: EX-II-10-08
% mdparagraph:P392:begin
God gave us his Gospel.
% mdparagraph:P392:end

\end{minipage}\par\addvspace{4pt}

\par\begin{minipage}{\linewidth}\raggedright
% mdparagraph:P393:begin
\textbf{Cook II.10.9.}
% mdparagraph:P393:end

% item: EX-II-10-09
% mdparagraph:P394:begin
The farmer practises agriculture.
% mdparagraph:P394:end

\end{minipage}\par\addvspace{4pt}

\par\begin{minipage}{\linewidth}\raggedright
% mdparagraph:P395:begin
\textbf{Cook II.10.10.}
% mdparagraph:P395:end

% item: EX-II-10-10
% mdparagraph:P396:begin
The service of God [God’s service] is better (66) than other occupations.
% mdparagraph:P396:end

\end{minipage}\par\addvspace{4pt}

\par\begin{minipage}{\linewidth}\raggedright
% mdparagraph:P397:begin
\textbf{Cook II.10.11.}
% mdparagraph:P397:end

% item: EX-II-10-11
% mdparagraph:P398:begin
Which of you holds (om. the) chief place?
% mdparagraph:P398:end

\end{minipage}\par\addvspace{4pt}

\par\begin{minipage}{\linewidth}\raggedright
% mdparagraph:P399:begin
\textbf{Cook II.10.12.}
% mdparagraph:P399:end

% item: EX-II-10-12
% mdparagraph:P400:begin
My gain is little.
% mdparagraph:P400:end

\end{minipage}\par\addvspace{4pt}

\par\begin{minipage}{\linewidth}\raggedright
% mdparagraph:P401:begin
\textbf{Cook II.10.13.}
% mdparagraph:P401:end

% item: EX-II-10-13
% mdparagraph:P402:begin
My labor is great.
% mdparagraph:P402:end

\end{minipage}\par\addvspace{4pt}

% teaching: NEW-U09-COOK-HELP-IVORY
% mdparagraph:P403:begin
\textbf{Local help for II.10.14:} \textbf{elpes bān} — “ivory”, literally “elephant’s bone”. Elpes is genitive singular; bān is neuter accusative singular as the thing sold. Use this whole expression as given.
% mdparagraph:P403:end

\reserve{5\baselineskip}
\par\begin{minipage}{\linewidth}\raggedright
% mdparagraph:P404:begin
\textbf{Cook II.10.14.}
% mdparagraph:P404:end

% item: EX-II-10-14
% mdparagraph:P405:begin
I sold ivory and silk.
% mdparagraph:P405:end

\end{minipage}\par\addvspace{4pt}

\par\begin{minipage}{\linewidth}\raggedright
% mdparagraph:P406:begin
\textbf{Cook II.10.15.}
% mdparagraph:P406:end

% item: EX-II-10-15
% mdparagraph:P407:begin
O, were I king of this people!
% mdparagraph:P407:end

\end{minipage}\par\addvspace{4pt}

\reserve{7\baselineskip}
% mdparagraph:P408:begin
\subsection*{Exercise 11}\BookHeading{nav-ch09-learner-0127}{subsubsection}{Exercise 11}{0}
% mdparagraph:P408:end

\reserve{5\baselineskip}
\par\begin{minipage}{\linewidth}\raggedright
% mdparagraph:P409:begin
\textbf{Cook II.11.1.}
% mdparagraph:P409:end

% item: EX-II-11-01
% mdparagraph:P410:begin
The monks have food and drink.
% mdparagraph:P410:end

\end{minipage}\par\addvspace{4pt}

\par\begin{minipage}{\linewidth}\raggedright
% mdparagraph:P411:begin
\textbf{Cook II.11.2.}
% mdparagraph:P411:end

% item: EX-II-11-02
% mdparagraph:P412:begin
Settle [quell] your disputes.
% mdparagraph:P412:end

\end{minipage}\par\addvspace{4pt}

\par\begin{minipage}{\linewidth}\raggedright
% mdparagraph:P413:begin
\textbf{Cook II.11.3.}
% mdparagraph:P413:end

% item: EX-II-11-03
% mdparagraph:P414:begin
He suffered shipwreck.
% mdparagraph:P414:end

\end{minipage}\par\addvspace{4pt}

\par\begin{minipage}{\linewidth}\raggedright
% mdparagraph:P415:begin
\textbf{Cook II.11.4.}
% mdparagraph:P415:end

% item: EX-II-11-04
% mdparagraph:P416:begin
He abandoned his wife.
% mdparagraph:P416:end

\end{minipage}\par\addvspace{4pt}

\par\begin{minipage}{\linewidth}\raggedright
% mdparagraph:P417:begin
\textbf{Cook II.11.5.}
% mdparagraph:P417:end

% item: EX-II-11-05
% mdparagraph:P418:begin
Be what [that] you (sing.) ought to be.
% mdparagraph:P418:end

\end{minipage}\par\addvspace{4pt}

\par\begin{minipage}{\linewidth}\raggedright
% mdparagraph:P419:begin
\textbf{Cook II.11.6.}
% mdparagraph:P419:end

% item: EX-II-11-06
% mdparagraph:P420:begin
They do not wish to be what they ought to be.
% mdparagraph:P420:end

\end{minipage}\par\addvspace{4pt}

\par\begin{minipage}{\linewidth}\raggedright
% mdparagraph:P421:begin
\textbf{Cook II.11.7.}
% mdparagraph:P421:end

% item: EX-II-11-07
% mdparagraph:P422:begin
He went aboard the ship with his wares.
% mdparagraph:P422:end

\end{minipage}\par\addvspace{4pt}

\par\begin{minipage}{\linewidth}\raggedright
% mdparagraph:P423:begin
\textbf{Cook II.11.8.}
% mdparagraph:P423:end

% item: EX-II-11-08
% mdparagraph:P424:begin
O good workmen, ply your trades.
% mdparagraph:P424:end

\end{minipage}\par\addvspace{4pt}

\par\begin{minipage}{\linewidth}\raggedright
% mdparagraph:P425:begin
\textbf{Cook II.11.9.}
% mdparagraph:P425:end

% item: EX-II-11-09
% mdparagraph:P426:begin
I tell you (sing.), the shoemaker buys awls, and the fisherman buys (fish-) hooks.
% mdparagraph:P426:end

\end{minipage}\par\addvspace{4pt}

\par\begin{minipage}{\linewidth}\raggedright
% mdparagraph:P427:begin
\textbf{Cook II.11.10.}
% mdparagraph:P427:end

% item: EX-II-11-10
% mdparagraph:P428:begin
Sometimes he reads, because he is wise.
% mdparagraph:P428:end

\end{minipage}\par\addvspace{4pt}

\par\begin{minipage}{\linewidth}\raggedright
% mdparagraph:P429:begin
\textbf{Cook II.11.11.}
% mdparagraph:P429:end

% item: EX-II-11-11
% mdparagraph:P430:begin
The king buys horses.
% mdparagraph:P430:end

\end{minipage}\par\addvspace{4pt}

\par\begin{minipage}{\linewidth}\raggedright
% mdparagraph:P431:begin
\textbf{Cook II.11.12.}
% mdparagraph:P431:end

% item: EX-II-11-12
% mdparagraph:P432:begin
Righteousness is better than gold.
% mdparagraph:P432:end

\end{minipage}\par\addvspace{4pt}

\par\begin{minipage}{\linewidth}\raggedright
% mdparagraph:P433:begin
\textbf{Cook II.11.13.}
% mdparagraph:P433:end

% item: EX-II-11-13
% mdparagraph:P434:begin
They struck me.
% mdparagraph:P434:end

\end{minipage}\par\addvspace{4pt}

\par\begin{minipage}{\linewidth}\raggedright
% mdparagraph:P435:begin
\textbf{Cook II.11.14.}
% mdparagraph:P435:end

% item: EX-II-11-14
% mdparagraph:P436:begin
The farmer beats him.
% mdparagraph:P436:end

\end{minipage}\par\addvspace{4pt}

\par\begin{minipage}{\linewidth}\raggedright
% mdparagraph:P437:begin
\textbf{Cook II.11.15.}
% mdparagraph:P437:end

% item: EX-II-11-15
% mdparagraph:P438:begin
We suffer the loss of copper and tin.
% mdparagraph:P438:end

\end{minipage}\par\addvspace{4pt}

\par\begin{minipage}{\linewidth}\raggedright
% mdparagraph:P439:begin
\textbf{Cook II.11.16.}
% mdparagraph:P439:end

% item: EX-II-11-16
% mdparagraph:P440:begin
The princes of the earth go aboard their ships.
% mdparagraph:P440:end

\end{minipage}\par\addvspace{4pt}

\par\begin{minipage}{\linewidth}\raggedright
% mdparagraph:P441:begin
\textbf{Cook II.11.17.}
% mdparagraph:P441:end

% item: EX-II-11-17
% mdparagraph:P442:begin
The priest benefits (164. e) the laymen.
% mdparagraph:P442:end

\end{minipage}\par\addvspace{4pt}

\reserve{7\baselineskip}
% mdparagraph:P443:begin
\subsection*{Exercise 12}\BookHeading{nav-ch09-learner-0128}{subsubsection}{Exercise 12}{0}
% mdparagraph:P443:end

\reserve{5\baselineskip}
\par\begin{minipage}{\linewidth}\raggedright
% mdparagraph:P444:begin
\textbf{Cook II.12.1.}
% mdparagraph:P444:end

% item: EX-II-12-01
% mdparagraph:P445:begin
The counsellor is wiser (64) than the warrior.
% mdparagraph:P445:end

\end{minipage}\par\addvspace{4pt}

\par\begin{minipage}{\linewidth}\raggedright
% mdparagraph:P446:begin
\textbf{Cook II.12.2.}
% mdparagraph:P446:end

% item: EX-II-12-02
% mdparagraph:P447:begin
The elephant is larger than the horse.
% mdparagraph:P447:end

\end{minipage}\par\addvspace{4pt}

\par\begin{minipage}{\linewidth}\raggedright
% mdparagraph:P448:begin
\textbf{Cook II.12.3.}
% mdparagraph:P448:end

% item: EX-II-12-03
% mdparagraph:P449:begin
God created all (om. the) fishes in (om. the) sea.
% mdparagraph:P449:end

\end{minipage}\par\addvspace{4pt}

\par\begin{minipage}{\linewidth}\raggedright
% mdparagraph:P450:begin
\textbf{Cook II.12.4.}
% mdparagraph:P450:end

% item: EX-II-12-04
% mdparagraph:P451:begin
The cattle seek fodder.
% mdparagraph:P451:end

\end{minipage}\par\addvspace{4pt}

\par\begin{minipage}{\linewidth}\raggedright
% mdparagraph:P452:begin
\textbf{Cook II.12.5.}
% mdparagraph:P452:end

% item: EX-II-12-05
% mdparagraph:P453:begin
God commanded the earth to bring forth grass.
% mdparagraph:P453:end

\end{minipage}\par\addvspace{4pt}

\par\begin{minipage}{\linewidth}\raggedright
% mdparagraph:P454:begin
\textbf{Cook II.12.6.}
% mdparagraph:P454:end

% item: EX-II-12-06
% mdparagraph:P455:begin
On the third day God gathered the waters that were under the heavens.
% mdparagraph:P455:end

\end{minipage}\par\addvspace{4pt}

% teaching: NEW-U09-COOK-HELP-LIVING
% mdparagraph:P456:begin
\textbf{Local help for II.12.7:} \textbf{libbende gesceafta} — “living creatures”, feminine accusative plural. Use this whole phrase as given. Libbende uses the documented feminine -e participial ending; the earlier neuter libbendu would be a different agreement form.
% mdparagraph:P456:end

\reserve{5\baselineskip}
\par\begin{minipage}{\linewidth}\raggedright
% mdparagraph:P457:begin
\textbf{Cook II.12.7.}
% mdparagraph:P457:end

% item: EX-II-12-07
% mdparagraph:P458:begin
Let the earth bring forth living creatures (the verb precedes).
% mdparagraph:P458:end

\end{minipage}\par\addvspace{4pt}

\par\begin{minipage}{\linewidth}\raggedright
% mdparagraph:P459:begin
\textbf{Cook II.12.8.}
% mdparagraph:P459:end

% item: EX-II-12-08
% mdparagraph:P460:begin
God made man good (173).
% mdparagraph:P460:end

\end{minipage}\par\addvspace{4pt}

\par\begin{minipage}{\linewidth}\raggedright
% mdparagraph:P461:begin
\textbf{Cook II.12.9.}
% mdparagraph:P461:end

% item: EX-II-12-09
% mdparagraph:P462:begin
He is (om. the) light that gives light to (ālīehtan) all things.
% mdparagraph:P462:end

\end{minipage}\par\addvspace{4pt}

\par\begin{minipage}{\linewidth}\raggedright
% mdparagraph:P463:begin
\textbf{Cook II.12.10.}
% mdparagraph:P463:end

% item: EX-II-12-10
% mdparagraph:P464:begin
Some of them (151) were rich, and some were wise.
% mdparagraph:P464:end

\end{minipage}\par\addvspace{4pt}

\par\begin{minipage}{\linewidth}\raggedright
% mdparagraph:P465:begin
\textbf{Cook II.12.11.}
% mdparagraph:P465:end

% item: EX-II-12-11
% mdparagraph:P466:begin
The night has many stars.
% mdparagraph:P466:end

\end{minipage}\par\addvspace{4pt}

\par\begin{minipage}{\linewidth}\raggedright
% mdparagraph:P467:begin
\textbf{Cook II.12.12.}
% mdparagraph:P467:end

% item: EX-II-12-12
% mdparagraph:P468:begin
The light shineth in (on) darkness (plur.).
% mdparagraph:P468:end

\end{minipage}\par\addvspace{4pt}

% teaching: NEW-U09-COOK-HELP-365
% teaching: NEW-U09-VOC-hund
% teaching: NEW-U09-VOC-siextig
% mdparagraph:P469:begin
\textbf{Local help for II.12.13:} \textbf{ðrēo hund and fīf and siextig daga} — “three hundred and sixty-five days”. Use this whole phrase as given. Hund is “hundred”, siextig “sixty”; ðrēo counts the neuter hundred-unit, and daga is genitive plural. The phrase is composed from Cook’s numeral rules, not quoted as a sentence from his reader. You are not being asked to recall a complete higher-numeral system.
% mdparagraph:P469:end

\reserve{5\baselineskip}
\par\begin{minipage}{\linewidth}\raggedright
% mdparagraph:P470:begin
\textbf{Cook II.12.13.}
% mdparagraph:P470:end

% item: EX-II-12-13
% mdparagraph:P471:begin
The year has three hundred and sixty-five (78) days.
% mdparagraph:P471:end

\end{minipage}\par\addvspace{4pt}

\par\begin{minipage}{\linewidth}\raggedright
% mdparagraph:P472:begin
\textbf{Cook II.12.14.}
% mdparagraph:P472:end

% item: EX-II-12-14
% mdparagraph:P473:begin
The carpenter has two sons.
% mdparagraph:P473:end

\end{minipage}\par\addvspace{4pt}

\par\begin{minipage}{\linewidth}\raggedright
% mdparagraph:P474:begin
\textbf{Cook II.12.15.}
% mdparagraph:P474:end

% item: EX-II-12-15
% mdparagraph:P475:begin
The soldier has eight houses.
% mdparagraph:P475:end

\end{minipage}\par\addvspace{4pt}

\par\begin{minipage}{\linewidth}\raggedright
% mdparagraph:P476:begin
\textbf{Cook II.12.16.}
% mdparagraph:P476:end

% item: EX-II-12-16
% mdparagraph:P477:begin
The peril was great.
% mdparagraph:P477:end

\end{minipage}\par\addvspace{4pt}

\reserve{7\baselineskip}
% mdparagraph:P478:begin
\section*{13. Cumulative practice}\BookHeading{nav-ch09-learner-0129}{subsection}{13. Cumulative practice}{1}
% mdparagraph:P478:end

\reserve{7\baselineskip}
% mdparagraph:P479:begin
\subsection*{Combining the constructions}\BookHeading{nav-ch09-learner-0130}{subsubsection}{Combining the constructions}{0}
% mdparagraph:P479:end

% mdparagraph:P480:begin
The second task contains deliberate errors to diagnose and correct; its input is not a model of correct Old English.
% mdparagraph:P480:end

\reserve{5\baselineskip}
\par\begin{minipage}{\linewidth}\raggedright
% mdparagraph:P481:begin
\textbf{X01.}
% mdparagraph:P481:end

% item: NEW-U09-X01
% mdparagraph:P482:begin
Translate this connected account: The fisherman commanded the men to carry the fish (sing.). I saw them carry it. I know that they helped him. Use hātan and gesēon with plain infinitives, then an indicative factual content clause.
% mdparagraph:P482:end

\end{minipage}\par\addvspace{4pt}

\par\begin{minipage}{\linewidth}\raggedright
% mdparagraph:P483:begin
\textbf{X02.}
% mdparagraph:P483:end

% item: NEW-U09-X02
% mdparagraph:P484:begin
Repair the case relations in these intended sentences and explain each repair: (a) ‘I see him carrying the fish’, written Ic gesēo hē ðone fisc beran.; (b) ‘God made the man good’, written God geworhte ðone mann gōd. Keep the selected complements.
% mdparagraph:P484:end

\end{minipage}\par\addvspace{4pt}

\clearpage
% mdparagraph:P485:begin
\section*{Sources and editorial notes}\BookHeading{nav-ch09-learner-0131}{subsection}{Sources and editorial notes}{1}
% mdparagraph:P485:end

\begingroup\fontsize{10}{12.8}\selectfont
\setlength{\parskip}{1.5pt}
% mdparagraph:P486:begin
The principal editions are Albert S. Cook, \emph{A First Book in Old English: Grammar, Reader, Notes, and Vocabulary}, second edition, revised and enlarged (Ginn, 1895), and \emph{Exercises in Old English} (Ginn, 1895). The complete II.9–12 prompts are on \emph{Exercises} pp. 4–6 /\allowbreak{} PDF pp. 12–14. Their wording and original grammar references are retained. The twelve examples, fisherman dialogue and lettered tasks are new teaching material.
% mdparagraph:P486:end

% mdparagraph:P487:begin
“The Choice of Occupations” is on \emph{First Book} p. 131 /\allowbreak{} PDF p. 147. The four selected turns occupy prose lines 7–15 counted from the top of the page, with printed marginal numbers 10 and 15 anchoring the count. Cook’s words, spelling and conjectural insertions are retained. Source line breaks and superscript note cues are replaced by paragraphs and the separately printed original notes, whose numbers and anchors identify their locations. The preceding goods catalogue and price exchange are outside this excerpt.
% mdparagraph:P487:end

% mdparagraph:P488:begin
The complement and wish constructions follow Cook §§164–165, 169, 172–173, 193–196 and 199–200, with the partitive and appositional contrast in §§151 and 154b. R. D. Fulk, \emph{An Introductory Grammar of Old English with an Anthology of Readings}, author-corrected open-access reprint (2020), §§44–45 and 58, corroborates the mood and infinitive distinctions; his note on þrēo, p. 47, supports neuter agreement with mixed-gender groups. The preference, trade and service frames have contextual witnesses on Cook pp. 130–133.
% mdparagraph:P488:end

% mdparagraph:P489:begin
The noun and adjective patterns use Cook pp. 27–43 and the individual vocabulary entries. The verbal systems use §§92–119 and 141, with the stem alternations in §§15–19 and 33–37. Complete tables combine these rules with the checked lexical stems; they do not claim that every cell occurs in the reading. Joseph and Elizabeth Mary Wright, \emph{Old English Grammar}, second edition (1914), §§521–530, corroborates the weak-I forms. The selected wīsra genitive plural applies Wright §259.3, p. 130, on simplifying doubled r after another consonant.
% mdparagraph:P489:end

% mdparagraph:P490:begin
Bosworth–Toller, \emph{An Anglo-Saxon Dictionary} (1898), corroborates the specific forms or constructions under āstīgan (p. 55b), bycgan (p. 137a–b), gestillan (p. 446a), sęcgan (p. 855b) and absolute slēan (p. 883a); the \emph{Supplement} (1921), hwā III (pp. 575a–b), supports dependent questions. The main dictionary’s hund entry (p. 566a), together with Cook’s numeral rules on pp. 46–48, supports the locally supplied 365-day expression. That expression and libbende gesceafta are constructed supports, not quotations from the reader.
% mdparagraph:P490:end

% mdparagraph:P491:begin
Compound hyphens in dictionary citations are omitted in the ordinary displayed headwords; running source spelling, including for-ðām and swā-swā, is retained. The comparative wīsra and regular inflected cells are rule-based forms. The edition’s sēceað remains in the source beside selected productive sēcað. PDF locators are one-based: the selected \emph{First Book} scan uses printed page +16, \emph{Exercises} +8, Wright +20 and Fulk +12; dictionary a/\allowbreak{}b locators identify printed columns.
% mdparagraph:P491:end

\endgroup

% book-source-body:end
\par\label{end-ch09-learner}
\clearpage\endgroup

\BookRecto
\begingroup
\renewcommand{\runningtitle}{Unit 10. The case system}
\hyphenpenalty=10000
\exhyphenpenalty=10000
\emergencystretch=3em
\newcommand{\itempara}[2]{\par\addvspace{3pt}\noindent\hangindent=1.8em\hangafter=1\makebox[1.8em][l]{#1.}#2\par}
\newcommand{\reserve}[1]{\par\penalty-100\begingroup\dimen0=\pagegoal\advance\dimen0 by-\pagetotal\ifdim\dimen0<#1\newpage\fi\endgroup}
% book-source-body:begin

\addfontfeatures{Ligatures=TeXOff}
\raggedright
\booktitle{\BookComponent{ch10-learner}{Unit 10. The case system in extended sentences}UNIT 10}{The case system in extended sentences}

% mdparagraph:P001:begin
Case identifies a noun phrase's function in its clause and the relation of one noun to another. A dative recipient may contain a genitive possessor; a verb may itself require a genitive object. The same English word “with” may express accompaniment or means. This unit brings these distinctions together in longer sentences.
% mdparagraph:P001:end

% mdparagraph:P002:begin
The fourteen translated examples, the exercises marked F, A, T, R and S, and “The stranger and the monk” are new teaching material. The three Cook exercise sets retain their original English prompts. “The Day of Judgment” is an exact source passage with supplied help and Cook's original notes.
% mdparagraph:P002:end

% mdparagraph:P003:begin
In the tables, 1, 2 and 3 denote person. A repeated spelling occupies a separate grammatical cell in each position. The ordinary nominative, genitive, dative and accusative forms are for use in composition unless expressly marked \textbf{reference}. Reference forms are for recognition; supplied expressions may be used with their printed help. Keep vowel lengths and the distinction between e and ę.
% mdparagraph:P003:end

\reserve{7\baselineskip}
% mdparagraph:P004:begin
\section*{1. The case of a phrase and the dependent genitive}\BookHeading{nav-ch10-learner-0001}{subsection}{1. The case of a phrase and the dependent genitive}{1}
% mdparagraph:P004:end

% teaching: NEW-U10-GRAM-u10-genitive-relations
% mdparagraph:P005:begin
A\BookPoint{bk-ch10-learner-9d51aaaeaec2} noun phrase receives its case from its role in the clause. Its head noun and agreeing modifiers show that case. A genitive dependent inside the phrase keeps the case required by its relation to the head.
% mdparagraph:P005:end

% mdparagraph:P006:begin
A\BookPoint{bk-ch10-learner-5d1ddbd64d12} dependent genitive may identify a possessor or a kinship relation, the source or belonging of something, or the person associated with an action named by a noun. In a definite possessive expression, this unit normally puts the genitive before the head: \textbf{ðæs cuman sunu}, “the stranger's son”. The article belongs to cuman; do not add a second article before sunu in this selected pattern.
% mdparagraph:P006:end

% mdparagraph:P007:begin
The head can change case without changing its dependent genitive.
% mdparagraph:P007:end

\reserve{8\baselineskip}
% example: NEW-U10-EXAMPLE01
% mdparagraph:P008:begin
Se munuc sealde ðæs cuman suna męte.
% mdparagraph:P008:end

% mdparagraph:P009:begin
The monk gave food to the stranger’s son.
% mdparagraph:P009:end

% mdparagraph:P010:begin
The whole phrase \textbf{ðæs cuman suna} is the dative recipient of sealde. Its head suna is dative singular, while ðæs cuman is genitive singular, “of the stranger”. Męte, the food given, is accusative. The genitive does not make the complete recipient phrase genitive.
% mdparagraph:P010:end

% mdparagraph:P011:begin
With\BookPoint{bk-ch10-learner-606f8f7495ea} a noun naming an action or feeling, a dependent genitive may denote the person performing it, a \textbf{subjective genitive}, or the person or thing toward which it is directed, an \textbf{objective genitive}. Determine the relation from the meaning: “the person's fear” can mean fear felt by that person, while “fear of the person” can mean that someone fears that person. The endings alone do not distinguish these readings. See source note 2 for Cook's objective-genitive example and correction C02.
% mdparagraph:P011:end

% teaching: NEW-U10-GRAM-u10-governed-case-contrast
% teaching: NEW-U10-GRAM-u10-forwiernan-government
% teaching: NEW-U10-GRAM-u10-unnan-government
% teaching: NEW-U10-VOC-FORWIERNAN
% teaching: NEW-U10-VOC-UNNAN
% mdparagraph:P012:begin
A\BookPoint{bk-ch10-learner-4bed4c080848} \textbf{governed genitive} depends directly on a verb. It is not a possessor inside an object phrase. Keep the earlier constructions: \textbf{gīeman + genitive}, “rule over”, and \textbf{geswīcan + genitive}, “cease from” an action or work. An English object does not by itself justify an Old English accusative.
% mdparagraph:P012:end

% mdparagraph:P013:begin
Two further verbs take both a dative person and a genitive thing:
% mdparagraph:P013:end

% mdparagraph:P014:begin
\textbf{Forwiernan},\BookPoint{bk-ch10-learner-3d7f287af6c0} “deny, refuse”: the person denied is dative; the thing refused is genitive.
% mdparagraph:P014:end

% mdparagraph:P015:begin
\textbf{Unnan}, “grant, allow”: the person receiving the grant is dative; the thing granted is genitive.
% mdparagraph:P015:end

% mdparagraph:P016:begin
These case frames remain the same with a statement, a question or a command. The complete verb forms follow in §7.
% mdparagraph:P016:end

\reserve{8\baselineskip}
% example: NEW-U10-EXAMPLE02
% mdparagraph:P017:begin
Ðū forwiernst mē mętes.
% mdparagraph:P017:end

% mdparagraph:P018:begin
You deny me food.
% mdparagraph:P018:end

% mdparagraph:P019:begin
Mē is the dative person; mętes is the genitive of the food refused. Compare sęllan, “give”: its person is also dative, but the thing given is accusative. Similar English word order does not make the constructions identical.
% mdparagraph:P019:end

\reserve{7\baselineskip}
% mdparagraph:P020:begin
\section*{2. Nouns: complete models and lexical diagnostics}\BookHeading{nav-ch10-learner-0002}{subsection}{2. Nouns: complete models and lexical diagnostics}{1}
% mdparagraph:P020:end

% teaching: NEW-U10-LEXICAL-DIAGNOSTICS
% mdparagraph:P021:begin
Each new noun has its gender, meaning, stem and diagnostic forms below. A \textbf{stem} is the part retained when an ending is added. For a noun assigned to a printed model, use the cell with the same case and number, replace the model stem by the listed stem, and retain that cell's ending. A cell with no ending retains no ending in the new noun. The stated full headword is used wherever the rule specifies it.
% mdparagraph:P021:end

% mdparagraph:P022:begin
The tables give all four cases in both numbers. \textbf{Crīst} is a proper name used only in the singular. \textbf{Dōm}, “judgment”, \textbf{rēaf}, “raiment”, and \textbf{woruld}, “world”, are used here in the singular; their complete plural forms are reference forms, without an additional plural meaning to learn. The remaining nouns have both numbers available in their assigned meanings.
% mdparagraph:P022:end

\reserve{7\baselineskip}
% mdparagraph:P023:begin
\subsection*{Masculine models}\BookHeading{nav-ch10-learner-0003}{subsubsection}{Masculine models}{0}
% mdparagraph:P023:end

% mdparagraph:P024:begin
The ordinary consonant masculine uses -es and -e in the singular genitive and dative, and -as, -a, -um, -as in the plural. Its nominative and accusative singular have no ending. \textbf{Dōm} supplies this complete model; its own plural is reference.
% mdparagraph:P024:end

% teaching: NEW-U10-NOM-DOM
% table: NEW-U10-NOM-DOM
% mdtable:T01:begin
\par\begingroup
\fontsize{10.5}{13.2}\selectfont
\setlength{\LTleft}{0pt}\setlength{\LTright}{\fill}
\noindent
\begin{tabular}{@{}>{\raggedright\arraybackslash}p{39.46mm}>{\raggedright\arraybackslash}p{39.46mm}>{\raggedright\arraybackslash}p{39.46mm}@{}}
\toprule
\textbf{Case} & \textbf{Singular} & \textbf{Plural} \\
\midrule
Nominative & dōm & dōmas \\
Genitive & dōmes & dōma \\
Dative & dōme & dōmum \\
Accusative & dōm & dōmas \\
\bottomrule
\end{tabular}
\par\endgroup

% mdtable:T01:end
% mdparagraph:P025:begin
\textbf{Ęngel},\BookPoint{bk-ch10-learner-4c244822ba6e} “angel”, retains its internal e in the nominative and accusative singular. Before an ending it has the shorter stem \textbf{ęngl-}. Use this full-headword/\allowbreak{}oblique-stem distinction also for \textbf{dryhten}, stem \textbf{dryhtn-}, and \textbf{dēofol}, stem \textbf{dēofl-}.
% mdparagraph:P025:end

% teaching: NEW-U10-NOM-ENGEL
% table: NEW-U10-NOM-ENGEL
% mdtable:T02:begin
\par\begingroup
\fontsize{10.5}{13.2}\selectfont
\setlength{\LTleft}{0pt}\setlength{\LTright}{\fill}
\noindent
\begin{tabular}{@{}>{\raggedright\arraybackslash}p{39.46mm}>{\raggedright\arraybackslash}p{39.46mm}>{\raggedright\arraybackslash}p{39.46mm}@{}}
\toprule
\textbf{Case} & \textbf{Singular} & \textbf{Plural} \\
\midrule
Nominative & ęngel & ęnglas \\
Genitive & ęngles & ęngla \\
Dative & ęngle & ęnglum \\
Accusative & ęngel & ęnglas \\
\bottomrule
\end{tabular}
\par\endgroup

% mdtable:T02:end
% mdparagraph:P026:begin
Thus every ending-bearing cell of dryhten and dēofol uses the same ending as the corresponding cell of ęngel, attached to dryhtn- or dēofl-. Their nominative and accusative singular are the full dryhten and dēofol. Retained-vowel variants of dryhten, including genitive \textbf{dryhtenes}, are also valid; the selected table route uses dryhtn- consistently. Do not extend this vowel loss to an unlisted noun.
% mdparagraph:P026:end

% mdparagraph:P027:begin
\textbf{Scēaphierde}, “shepherd”, keeps -e in the nominative and accusative singular. Its stem for endings is \textbf{scēaphierd-}; the genitive singular and the plural show the ordinary masculine endings.
% mdparagraph:P027:end

% teaching: NEW-U10-NOM-SCEAPHIERDE
% table: NEW-U10-NOM-SCEAPHIERDE
% mdtable:T03:begin
\par\begingroup
\fontsize{10.5}{13.2}\selectfont
\setlength{\LTleft}{0pt}\setlength{\LTright}{\fill}
\noindent
\begin{tabular}{@{}>{\raggedright\arraybackslash}p{34.82mm}>{\raggedright\arraybackslash}p{41.78mm}>{\raggedright\arraybackslash}p{41.78mm}@{}}
\toprule
\textbf{Case} & \textbf{Singular} & \textbf{Plural} \\
\midrule
Nominative & scēaphierde & scēaphierdas \\
Genitive & scēaphierdes & scēaphierda \\
Dative & scēaphierde & scēaphierdum \\
Accusative & scēaphierde & scēaphierdas \\
\bottomrule
\end{tabular}
\par\endgroup

% mdtable:T03:end
% mdparagraph:P028:begin
The weak masculine \textbf{cuma}, “stranger, guest”, has -a only in the nominative singular, -an in its other singular cases and in the nominative and accusative plural, -ena in the genitive plural, and -um in the dative plural. \textbf{Nama}, stem \textbf{nam-}, follows every cell of this model.
% mdparagraph:P028:end

% teaching: NEW-U10-NOM-CUMA
% table: NEW-U10-NOM-CUMA
% mdtable:T04:begin
\par\begingroup
\fontsize{10.5}{13.2}\selectfont
\setlength{\LTleft}{0pt}\setlength{\LTright}{\fill}
\noindent
\begin{tabular}{@{}>{\raggedright\arraybackslash}p{39.46mm}>{\raggedright\arraybackslash}p{39.46mm}>{\raggedright\arraybackslash}p{39.46mm}@{}}
\toprule
\textbf{Case} & \textbf{Singular} & \textbf{Plural} \\
\midrule
Nominative & cuma & cuman \\
Genitive & cuman & cumena \\
Dative & cuman & cumum \\
Accusative & cuman & cuman \\
\bottomrule
\end{tabular}
\par\endgroup

% mdtable:T04:end
% mdparagraph:P029:begin
\textbf{Fæder},\BookPoint{bk-ch10-learner-53a33b32add5} “father”, has an exceptional endingless singular. Its plural is fæder- plus -as, -a, -um, -as. Genitive \textbf{fæderes} and dative \textbf{fædere} are accepted variants. The selected singular below has fæder in all four cases; do not derive a changed-vowel dative by analogy with brōðor.
% mdparagraph:P029:end

% teaching: NEW-U10-NOM-FAEDER
% table: NEW-U10-NOM-FAEDER
% mdtable:T05:begin
\par\begingroup
\fontsize{10.5}{13.2}\selectfont
\setlength{\LTleft}{0pt}\setlength{\LTright}{\fill}
\noindent
\begin{tabular}{@{}>{\raggedright\arraybackslash}p{39.46mm}>{\raggedright\arraybackslash}p{39.46mm}>{\raggedright\arraybackslash}p{39.46mm}@{}}
\toprule
\textbf{Case} & \textbf{Singular} & \textbf{Plural} \\
\midrule
Nominative & fæder & fæderas \\
Genitive & fæder & fædera \\
Dative & fæder & fæderum \\
Accusative & fæder & fæderas \\
\bottomrule
\end{tabular}
\par\endgroup

% mdtable:T05:end
\reserve{7\baselineskip}
% mdparagraph:P030:begin
\subsection*{Neuter model}\BookHeading{nav-ch10-learner-0004}{subsubsection}{Neuter model}{0}
% mdparagraph:P030:end

% mdparagraph:P031:begin
The heavy neuter \textbf{scēap}, “sheep”, has identical nominative and accusative forms in both numbers. The genitive and dative singular have -es and -e; the genitive and dative plural have -a and -um. Context and agreement distinguish singular from plural where the noun itself does not.
% mdparagraph:P031:end

% teaching: NEW-U10-NOM-SCEAP
% table: NEW-U10-NOM-SCEAP
% mdtable:T06:begin
\par\begingroup
\fontsize{10.5}{13.2}\selectfont
\setlength{\LTleft}{0pt}\setlength{\LTright}{\fill}
\noindent
\begin{tabular}{@{}>{\raggedright\arraybackslash}p{39.46mm}>{\raggedright\arraybackslash}p{39.46mm}>{\raggedright\arraybackslash}p{39.46mm}@{}}
\toprule
\textbf{Case} & \textbf{Singular} & \textbf{Plural} \\
\midrule
Nominative & scēap & scēap \\
Genitive & scēapes & scēapa \\
Dative & scēape & scēapum \\
Accusative & scēap & scēap \\
\bottomrule
\end{tabular}
\par\endgroup

% mdtable:T06:end
% mdparagraph:P032:begin
Replace scēap- by \textbf{rēaf-}, \textbf{cweartern-}, \textbf{giesthūs-} or \textbf{setl-} in every cell, retaining the same ending, including the zero ending. Cweartern keeps its final -ern throughout: cwearternes, not an automatically shortened stem. Giesthūs keeps the heavy neuter pattern of its final member hūs. No internal vowel is removed from these stems.
% mdparagraph:P032:end

\reserve{7\baselineskip}
% mdparagraph:P033:begin
\subsection*{Feminine models}\BookHeading{nav-ch10-learner-0005}{subsubsection}{Feminine models}{0}
% mdparagraph:P033:end

% mdparagraph:P034:begin
\textbf{Hand},\BookPoint{bk-ch10-learner-df6021e47304} “hand; side”, has an endingless accusative singular in the selected pattern. Its genitive and dative singular have -e; its nominative, genitive and accusative plural have -a, and the dative plural -um. \textbf{Woruld}, stem woruld-, follows these same eight cells.
% mdparagraph:P034:end

% teaching: NEW-U10-NOM-HAND
% table: NEW-U10-NOM-HAND
% mdtable:T07:begin
\par\begingroup
\fontsize{10.5}{13.2}\selectfont
\setlength{\LTleft}{0pt}\setlength{\LTright}{\fill}
\noindent
\begin{tabular}{@{}>{\raggedright\arraybackslash}p{39.46mm}>{\raggedright\arraybackslash}p{39.46mm}>{\raggedright\arraybackslash}p{39.46mm}@{}}
\toprule
\textbf{Case} & \textbf{Singular} & \textbf{Plural} \\
\midrule
Nominative & hand & handa \\
Genitive & hande & handa \\
Dative & hande & handum \\
Accusative & hand & handa \\
\bottomrule
\end{tabular}
\par\endgroup

% mdtable:T07:end
% mdparagraph:P035:begin
For hand, genitive and dative singular \textbf{handa}, accusative singular \textbf{hande} or \textbf{handa}, and nominative/\allowbreak{}accusative plural \textbf{hande} are accepted variants. Woruld has corresponding variants \textbf{worulda} in the genitive/\allowbreak{}dative singular, \textbf{worulde} or \textbf{worulda} in the accusative singular, and \textbf{worulde} in the nominative/\allowbreak{}accusative plural. Learn the selected pattern first; a variant is not an additional case.
% mdparagraph:P035:end

% mdparagraph:P036:begin
\textbf{Ðēod}, “people, nation”, has -e in the accusative singular as well as the genitive and dative. Its plural otherwise has the same endings as hand. \textbf{Wǣd}, “garment”, replaces ðēod- by wǣd- in all eight cells. Both nouns also allow -e in the nominative/\allowbreak{}accusative plural: ðēode, wǣde.
% mdparagraph:P036:end

% teaching: NEW-U10-NOM-THEOD
% table: NEW-U10-NOM-THEOD
% mdtable:T08:begin
\par\begingroup
\fontsize{10.5}{13.2}\selectfont
\setlength{\LTleft}{0pt}\setlength{\LTright}{\fill}
\noindent
\begin{tabular}{@{}>{\raggedright\arraybackslash}p{39.46mm}>{\raggedright\arraybackslash}p{39.46mm}>{\raggedright\arraybackslash}p{39.46mm}@{}}
\toprule
\textbf{Case} & \textbf{Singular} & \textbf{Plural} \\
\midrule
Nominative & ðēod & ðēoda \\
Genitive & ðēode & ðēoda \\
Dative & ðēode & ðēodum \\
Accusative & ðēode & ðēoda \\
\bottomrule
\end{tabular}
\par\endgroup

% mdtable:T08:end
% mdparagraph:P037:begin
The\BookPoint{bk-ch10-learner-55e62ee30134} root noun \textbf{gāt}, “goat”, changes ā to ǣ in the selected dative singular and nominative/\allowbreak{}accusative plural. Its genitive singular has -e; the genitive and dative plural retain ā and add -a and -um. Genitive singular \textbf{gǣt} and dative singular \textbf{gāte} are also valid. The table selects one form per cell; do not make a masculine -as plural or extend ǣ into gātum.
% mdparagraph:P037:end

% teaching: NEW-U10-NOM-GAT
% table: NEW-U10-NOM-GAT
% mdtable:T09:begin
\par\begingroup
\fontsize{10.5}{13.2}\selectfont
\setlength{\LTleft}{0pt}\setlength{\LTright}{\fill}
\noindent
\begin{tabular}{@{}>{\raggedright\arraybackslash}p{39.46mm}>{\raggedright\arraybackslash}p{39.46mm}>{\raggedright\arraybackslash}p{39.46mm}@{}}
\toprule
\textbf{Case} & \textbf{Singular} & \textbf{Plural} \\
\midrule
Nominative & gāt & gǣt \\
Genitive & gāte & gāta \\
Dative & gǣt & gātum \\
Accusative & gāt & gǣt \\
\bottomrule
\end{tabular}
\par\endgroup

% mdtable:T09:end
\reserve{7\baselineskip}
% mdparagraph:P038:begin
\subsection*{Noun vocabulary and model assignments}\BookHeading{nav-ch10-learner-0006}{subsubsection}{Noun vocabulary and model assignments}{0}
% mdparagraph:P038:end

% mdparagraph:P039:begin
In the singular diagnostic column the order is \textbf{genitive /\allowbreak{} dative /\allowbreak{} accusative}. The headword is nominative singular. “Stem; model” specifies the complete table to use above; the nominative plural checks the number pattern. For Crīst use only the singular half of dōm. For dryhten and dēofol use the full headword in nominative/\allowbreak{}accusative singular and the shorter listed stem in every other cell. In every other row, replace the named model's stem directly by the listed stem; the exceptional nouns have their own literal tables.
% mdparagraph:P039:end

% teaching: NEW-U10-VOC-CRIST
% teaching: NEW-U10-NOM-CRIST
% diagnostic: NEW-U10-DIAG-CRIST
% teaching: NEW-U10-VOC-DRYHTEN
% teaching: NEW-U10-NOM-DRYHTEN
% diagnostic: NEW-U10-DIAG-DRYHTEN
% teaching: NEW-U10-VOC-DOM
% diagnostic: NEW-U10-DIAG-DOM
% teaching: NEW-U10-VOC-REAF
% teaching: NEW-U10-NOM-REAF
% diagnostic: NEW-U10-DIAG-REAF
% teaching: NEW-U10-VOC-CUMA
% diagnostic: NEW-U10-DIAG-CUMA
% teaching: NEW-U10-VOC-CWEARTERN
% teaching: NEW-U10-NOM-CWEARTERN
% diagnostic: NEW-U10-DIAG-CWEARTERN
% teaching: NEW-U10-VOC-GIESTHUS
% teaching: NEW-U10-NOM-GIESTHUS
% diagnostic: NEW-U10-DIAG-GIESTHUS
% teaching: NEW-U10-VOC-FAEDER
% diagnostic: NEW-U10-DIAG-FAEDER
% teaching: NEW-U10-VOC-WORULD
% teaching: NEW-U10-NOM-WORULD
% diagnostic: NEW-U10-DIAG-WORULD
% teaching: NEW-U10-VOC-SCEAPHIERDE
% diagnostic: NEW-U10-DIAG-SCEAPHIERDE
% teaching: NEW-U10-VOC-SCEAP
% diagnostic: NEW-U10-DIAG-SCEAP
% teaching: NEW-U10-VOC-SETL
% teaching: NEW-U10-NOM-SETL
% diagnostic: NEW-U10-DIAG-SETL
% teaching: NEW-U10-VOC-HAND
% diagnostic: NEW-U10-DIAG-HAND
% teaching: NEW-U10-VOC-THEOD
% diagnostic: NEW-U10-DIAG-THEOD
% teaching: NEW-U10-VOC-DEOFOL
% teaching: NEW-U10-NOM-DEOFOL
% diagnostic: NEW-U10-DIAG-DEOFOL
% teaching: NEW-U10-VOC-WAED
% teaching: NEW-U10-NOM-WAED
% diagnostic: NEW-U10-DIAG-WAED
% teaching: NEW-U10-VOC-GAT
% diagnostic: NEW-U10-DIAG-GAT
% teaching: NEW-U10-VOC-NAMA
% teaching: NEW-U10-NOM-NAMA
% diagnostic: NEW-U10-DIAG-NAMA
% teaching: NEW-U10-VOC-ENGEL
% diagnostic: NEW-U10-DIAG-ENGEL
% mdtable:T10:begin
\par\begingroup
\fontsize{10.5}{13.2}\selectfont
\setlength{\LTleft}{0pt}\setlength{\LTright}{\fill}
\begin{longtable}{@{}>{\raggedright\arraybackslash}p{37.19mm}>{\raggedright\arraybackslash}p{22.35mm}>{\raggedright\arraybackslash}p{33.52mm}>{\raggedright\arraybackslash}p{22.5mm}@{}}
\toprule
\textbf{Headword and meaning} & \textbf{Stem; model} & \textbf{Singular G /\allowbreak{} D /\allowbreak{} A} & \textbf{Nominative plural} \\
\midrule
\endfirsthead
\toprule
\textbf{Headword and meaning} & \textbf{Stem; model} & \textbf{Singular G /\allowbreak{} D /\allowbreak{} A} & \textbf{Nominative plural} \\
\midrule
\endhead
\midrule
\multicolumn{4}{r}{\footnotesize Continued on the following page.}\\
\endfoot
\bottomrule
\endlastfoot
Crīst (m.), Christ & Crīst; dōm & Crīstes /\allowbreak{} Crīste /\allowbreak{} Crīst & — \\
dryhten (m.), lord; the Lord & dryhtn; ęngel & dryhtnes /\allowbreak{} dryhtne /\allowbreak{} dryhten & dryhtnas \\
dōm (m.), judgment; dōmes dæg, Day of Judgment & dōm; dōm & dōmes /\allowbreak{} dōme /\allowbreak{} dōm & dōmas \\
rēaf (n.), raiment; clothing & rēaf; scēap & rēafes /\allowbreak{} rēafe /\allowbreak{} rēaf & rēaf \\
cuma (m.), stranger; guest & cum; cuma & cuman /\allowbreak{} cuman /\allowbreak{} cuman & cuman \\
cweartern (n.), prison & cweartern; scēap & cwearternes /\allowbreak{} cwearterne /\allowbreak{} cweartern & cweartern \\
giesthūs (n.), inn; guest-house & giesthūs; scēap & giesthūses /\allowbreak{} giesthūse /\allowbreak{} giesthūs & giesthūs \\
fæder (m.), father; also the spiritual relationship & fæder; fæder & fæder /\allowbreak{} fæder /\allowbreak{} fæder & fæderas \\
woruld (f.), world & woruld; hand & worulde /\allowbreak{} worulde /\allowbreak{} woruld & worulda \\
scēaphierde (m.), shepherd & scēaphierd; scēaphierde & scēaphierdes /\allowbreak{} scēaphierde /\allowbreak{} scēaphierde & scēaphierdas \\
scēap (n.), sheep & scēap; scēap & scēapes /\allowbreak{} scēape /\allowbreak{} scēap & scēap \\
setl (n.), seat; throne & setl; scēap & setles /\allowbreak{} setle /\allowbreak{} setl & setl \\
hand (f.), hand; right or left side & hand; hand & hande /\allowbreak{} hande /\allowbreak{} hand & handa \\
ðēod (f.), people; nation & ðēod; ðēod & ðēode /\allowbreak{} ðēode /\allowbreak{} ðēode & ðēoda \\
dēofol (m.), devil; demon & dēofl; ęngel & dēofles /\allowbreak{} dēofle /\allowbreak{} dēofol & dēoflas \\
wǣd (f.), garment & wǣd; ðēod & wǣde /\allowbreak{} wǣde /\allowbreak{} wǣde & wǣda \\
gāt (f.), goat & gāt; gāt & gāte /\allowbreak{} gǣt /\allowbreak{} gāt & gǣt \\
nama (m.), name & nam; cuma & naman /\allowbreak{} naman /\allowbreak{} naman & naman \\
ęngel (m.), angel & ęngl; ęngel & ęngles /\allowbreak{} ęngle /\allowbreak{} ęngel & ęnglas \\
\end{longtable}
\par\endgroup

% mdtable:T10:end
% mdparagraph:P040:begin
The model assignment gives every remaining form. For example, dēofol takes the ending -um of ęnglum after dēofl-, producing \textbf{dēoflum}; the nominative singular remains \textbf{dēofol}. Nama replaces cum- in cumena with nam-, producing \textbf{namena}. Wǣd replaces ðēod- in ðēode with wǣd-, producing \textbf{wǣde} in each corresponding singular oblique cell. Apply the same operation separately to each case and number.
% mdparagraph:P040:end

\reserve{7\baselineskip}
% mdparagraph:P041:begin
\section*{3. Adjectives: agreement, complete patterns and people}\BookHeading{nav-ch10-learner-0007}{subsection}{3. Adjectives: agreement, complete patterns and people}{1}
% mdparagraph:P041:end

% mdparagraph:P042:begin
Adjectives\BookPoint{bk-ch10-learner-f9f1c5d1eb57} agree in gender, number and case with the noun they modify. Use the strong form in the ordinary indefinite and predicate uses already learned; use the weak form with the definite article in a definite attributive phrase. A separately predicated adjective remains strong even when its noun phrase contains an article.
% mdparagraph:P042:end

% mdparagraph:P043:begin
The following tables give complete strong and weak systems. Nominative, genitive, dative and accusative forms are for composition in these uses. The distinct masculine and neuter singular instrumental cells in these new adjective tables are \textbf{reference only}. “As dative” means the form is identical to the dative in the same gender and number. Productive instrumental practice in this unit uses the separately taught gōd and demonstrative forms in §6.
% mdparagraph:P043:end

\reserve{7\baselineskip}
% mdparagraph:P044:begin
\subsection*{The hungrig pattern}\BookHeading{nav-ch10-learner-0008}{subsubsection}{The hungrig pattern}{0}
% mdparagraph:P044:end

% mdparagraph:P045:begin
\textbf{Hungrig}, “hungry”, keeps its internal vowel before every ending in the selected forms. In the strong pattern the feminine nominative singular and neuter nominative/\allowbreak{}accusative plural have -u. The masculine accusative singular has -ne; feminine genitive/\allowbreak{}dative singular -re; genitive plural -ra; dative masculine/\allowbreak{}neuter singular and all plural genders -um. The weak pattern has nominative singular -a in the masculine and -e in the other genders, and -an in the other singular cells except neuter accusative -e. Its plural is -an, -ra, -um, -an.
% mdparagraph:P045:end

% teaching: NEW-U10-ADJ-HUNGRIG-STRONG
% table: NEW-U10-ADJ-HUNGRIG-STRONG
% mdtable:T11:begin
% course-table-style: 2026-10
% course-panel:T11-P01:begin
\par\begingroup
\fontsize{10.5}{13.2}\selectfont
\setlength{\LTleft}{0pt}\setlength{\LTright}{\fill}
\noindent
\begin{tabular}{@{}>{\raggedright\arraybackslash}p{22.0mm}>{\raggedright\arraybackslash}p{31.188mm}>{\raggedright\arraybackslash}p{31.188mm}>{\raggedright\arraybackslash}p{31.188mm}@{}}
\toprule
\textbf{Case} & \textbf{Masculine singular} & \textbf{Neuter singular} & \textbf{Feminine singular} \\
\midrule
Nominative & hungrig & hungrig & hungrigu \\
Genitive & hungriges & hungriges & hungrigre \\
Dative & hungrigum & hungrigum & hungrigre \\
Accusative & hungrigne & hungrig & hungrige \\
Instrumental & hungrige & hungrige & as dative \\
\bottomrule
\end{tabular}
\par\endgroup
% course-panel:T11-P01:end
\par\addvspace{5pt}
% course-panel:T11-P02:begin
\par\begingroup
\fontsize{10.5}{13.2}\selectfont
\setlength{\LTleft}{0pt}\setlength{\LTright}{\fill}
\noindent
\begin{tabular}{@{}>{\raggedright\arraybackslash}p{22.0mm}>{\raggedright\arraybackslash}p{31.188mm}>{\raggedright\arraybackslash}p{31.188mm}>{\raggedright\arraybackslash}p{31.188mm}@{}}
\toprule
\textbf{Case} & \textbf{Masculine plural} & \textbf{Neuter plural} & \textbf{Feminine plural} \\
\midrule
Nominative & hungrige & hungrigu & hungriga \\
Genitive & hungrigra & hungrigra & hungrigra \\
Dative & hungrigum & hungrigum & hungrigum \\
Accusative & hungrige & hungrigu & hungriga \\
Instrumental & as dative & as dative & as dative \\
\bottomrule
\end{tabular}
\par\endgroup
% course-panel:T11-P02:end
% mdtable:T11:end

% teaching: NEW-U10-ADJ-HUNGRIG-WEAK
% table: NEW-U10-ADJ-HUNGRIG-WEAK
% mdtable:T12:begin
% course-table-style: 2026-10
% course-panel:T12-P01:begin
\par\begingroup
\fontsize{10.5}{13.2}\selectfont
\setlength{\LTleft}{0pt}\setlength{\LTright}{\fill}
\noindent
\begin{tabular}{@{}>{\raggedright\arraybackslash}p{22.0mm}>{\raggedright\arraybackslash}p{31.188mm}>{\raggedright\arraybackslash}p{31.188mm}>{\raggedright\arraybackslash}p{31.188mm}@{}}
\toprule
\textbf{Case} & \textbf{Masculine singular} & \textbf{Neuter singular} & \textbf{Feminine singular} \\
\midrule
Nominative & hungriga & hungrige & hungrige \\
Genitive & hungrigan & hungrigan & hungrigan \\
Dative & hungrigan & hungrigan & hungrigan \\
Accusative & hungrigan & hungrige & hungrigan \\
Instrumental & hungrigan & hungrigan & as dative \\
\bottomrule
\end{tabular}
\par\endgroup
% course-panel:T12-P01:end
\par\addvspace{5pt}
% course-panel:T12-P02:begin
\par\begingroup
\fontsize{10.5}{13.2}\selectfont
\setlength{\LTleft}{0pt}\setlength{\LTright}{\fill}
\noindent
\begin{tabular}{@{}>{\raggedright\arraybackslash}p{22.0mm}>{\raggedright\arraybackslash}p{31.188mm}>{\raggedright\arraybackslash}p{31.188mm}>{\raggedright\arraybackslash}p{31.188mm}@{}}
\toprule
\textbf{Case} & \textbf{Masculine plural} & \textbf{Neuter plural} & \textbf{Feminine plural} \\
\midrule
Nominative & hungrigan & hungrigan & hungrigan \\
Genitive & hungrigra & hungrigra & hungrigra \\
Dative & hungrigum & hungrigum & hungrigum \\
Accusative & hungrigan & hungrigan & hungrigan \\
Instrumental & as dative & as dative & as dative \\
\bottomrule
\end{tabular}
\par\endgroup
% course-panel:T12-P02:end
% mdtable:T12:end

\reserve{7\baselineskip}
% mdparagraph:P046:begin
\subsection*{Adjectives with a lexical -e}\BookHeading{nav-ch10-learner-0009}{subsubsection}{Adjectives with a lexical -e}{0}
% mdparagraph:P046:end

% mdparagraph:P047:begin
\textbf{Dēorwierðe},\BookPoint{bk-ch10-learner-85a5a29685fa} “precious, valuable”, has the stem \textbf{dēorwierð-}. Its lexical -e remains in strong masculine/\allowbreak{}neuter nominative singular and neuter accusative singular. Replace that -e before another ending; do not add endings to dēorwierðe-. In the other cells the selected pattern has the endings printed below, including -u in feminine nominative singular and neuter nominative/\allowbreak{}accusative plural.
% mdparagraph:P047:end

% teaching: NEW-U10-ADJ-DEORWIERTH-STRONG
% table: NEW-U10-ADJ-DEORWIERTH-STRONG
% mdtable:T13:begin
% course-table-style: 2026-10
% course-panel:T13-P01:begin
\par\begingroup
\fontsize{10.5}{13.2}\selectfont
\setlength{\LTleft}{0pt}\setlength{\LTright}{\fill}
\noindent
\begin{tabular}{@{}>{\raggedright\arraybackslash}p{22.0mm}>{\raggedright\arraybackslash}p{31.188mm}>{\raggedright\arraybackslash}p{31.188mm}>{\raggedright\arraybackslash}p{31.188mm}@{}}
\toprule
\textbf{Case} & \textbf{Masculine singular} & \textbf{Neuter singular} & \textbf{Feminine singular} \\
\midrule
Nominative & dēorwierðe & dēorwierðe & dēorwierðu \\
Genitive & dēorwierðes & dēorwierðes & dēorwierðre \\
Dative & dēorwierðum & dēorwierðum & dēorwierðre \\
Accusative & dēorwierðne & dēorwierðe & dēorwierðe \\
Instrumental & dēorwierðe & dēorwierðe & as dative \\
\bottomrule
\end{tabular}
\par\endgroup
% course-panel:T13-P01:end
\par\addvspace{5pt}
% course-panel:T13-P02:begin
\par\begingroup
\fontsize{10.5}{13.2}\selectfont
\setlength{\LTleft}{0pt}\setlength{\LTright}{\fill}
\noindent
\begin{tabular}{@{}>{\raggedright\arraybackslash}p{22.0mm}>{\raggedright\arraybackslash}p{31.188mm}>{\raggedright\arraybackslash}p{31.188mm}>{\raggedright\arraybackslash}p{31.188mm}@{}}
\toprule
\textbf{Case} & \textbf{Masculine plural} & \textbf{Neuter plural} & \textbf{Feminine plural} \\
\midrule
Nominative & dēorwierðe & dēorwierðu & dēorwierða \\
Genitive & dēorwierðra & dēorwierðra & dēorwierðra \\
Dative & dēorwierðum & dēorwierðum & dēorwierðum \\
Accusative & dēorwierðe & dēorwierðu & dēorwierða \\
Instrumental & as dative & as dative & as dative \\
\bottomrule
\end{tabular}
\par\endgroup
% course-panel:T13-P02:end
% mdtable:T13:end

% teaching: NEW-U10-ADJ-DEORWIERTH-WEAK
% table: NEW-U10-ADJ-DEORWIERTH-WEAK
% mdtable:T14:begin
% course-table-style: 2026-10
% course-panel:T14-P01:begin
\par\begingroup
\fontsize{10.5}{13.2}\selectfont
\setlength{\LTleft}{0pt}\setlength{\LTright}{\fill}
\noindent
\begin{tabular}{@{}>{\raggedright\arraybackslash}p{22.0mm}>{\raggedright\arraybackslash}p{31.188mm}>{\raggedright\arraybackslash}p{31.188mm}>{\raggedright\arraybackslash}p{31.188mm}@{}}
\toprule
\textbf{Case} & \textbf{Masculine singular} & \textbf{Neuter singular} & \textbf{Feminine singular} \\
\midrule
Nominative & dēorwierða & dēorwierðe & dēorwierðe \\
Genitive & dēorwierðan & dēorwierðan & dēorwierðan \\
Dative & dēorwierðan & dēorwierðan & dēorwierðan \\
Accusative & dēorwierðan & dēorwierðe & dēorwierðan \\
Instrumental & dēorwierðan & dēorwierðan & as dative \\
\bottomrule
\end{tabular}
\par\endgroup
% course-panel:T14-P01:end
\par\addvspace{5pt}
% course-panel:T14-P02:begin
\par\begingroup
\fontsize{10.5}{13.2}\selectfont
\setlength{\LTleft}{0pt}\setlength{\LTright}{\fill}
\noindent
\begin{tabular}{@{}>{\raggedright\arraybackslash}p{22.0mm}>{\raggedright\arraybackslash}p{31.188mm}>{\raggedright\arraybackslash}p{31.188mm}>{\raggedright\arraybackslash}p{31.188mm}@{}}
\toprule
\textbf{Case} & \textbf{Masculine plural} & \textbf{Neuter plural} & \textbf{Feminine plural} \\
\midrule
Nominative & dēorwierðan & dēorwierðan & dēorwierðan \\
Genitive & dēorwierðra & dēorwierðra & dēorwierðra \\
Dative & dēorwierðum & dēorwierðum & dēorwierðum \\
Accusative & dēorwierðan & dēorwierðan & dēorwierðan \\
Instrumental & as dative & as dative & as dative \\
\bottomrule
\end{tabular}
\par\endgroup
% course-panel:T14-P02:end
% mdtable:T14:end

\reserve{7\baselineskip}
% mdparagraph:P048:begin
\subsection*{Long stems and doubled consonants}\BookHeading{nav-ch10-learner-0010}{subsubsection}{Long stems and doubled consonants}{0}
% mdparagraph:P048:end

% mdparagraph:P049:begin
\textbf{Unrihtwīs}, “unrighteous”, has a long final stem syllable. Its strong feminine nominative singular and neuter nominative/\allowbreak{}accusative plural have no -u. Keep the long ī throughout. The other strong and weak endings are shown in full.
% mdparagraph:P049:end

% teaching: NEW-U10-ADJ-UNRIHTWIS-STRONG
% table: NEW-U10-ADJ-UNRIHTWIS-STRONG
% mdtable:T15:begin
% course-table-style: 2026-10
% course-panel:T15-P01:begin
\par\begingroup
\fontsize{10.5}{13.2}\selectfont
\setlength{\LTleft}{0pt}\setlength{\LTright}{\fill}
\noindent
\begin{tabular}{@{}>{\raggedright\arraybackslash}p{22.0mm}>{\raggedright\arraybackslash}p{31.188mm}>{\raggedright\arraybackslash}p{31.188mm}>{\raggedright\arraybackslash}p{31.188mm}@{}}
\toprule
\textbf{Case} & \textbf{Masculine singular} & \textbf{Neuter singular} & \textbf{Feminine singular} \\
\midrule
Nominative & unrihtwīs & unrihtwīs & unrihtwīs \\
Genitive & unrihtwīses & unrihtwīses & unrihtwīsre \\
Dative & unrihtwīsum & unrihtwīsum & unrihtwīsre \\
Accusative & unrihtwīsne & unrihtwīs & unrihtwīse \\
Instrumental & unrihtwīse & unrihtwīse & as dative \\
\bottomrule
\end{tabular}
\par\endgroup
% course-panel:T15-P01:end
\par\addvspace{5pt}
% course-panel:T15-P02:begin
\par\begingroup
\fontsize{10.5}{13.2}\selectfont
\setlength{\LTleft}{0pt}\setlength{\LTright}{\fill}
\noindent
\begin{tabular}{@{}>{\raggedright\arraybackslash}p{22.0mm}>{\raggedright\arraybackslash}p{31.188mm}>{\raggedright\arraybackslash}p{31.188mm}>{\raggedright\arraybackslash}p{31.188mm}@{}}
\toprule
\textbf{Case} & \textbf{Masculine plural} & \textbf{Neuter plural} & \textbf{Feminine plural} \\
\midrule
Nominative & unrihtwīse & unrihtwīs & unrihtwīsa \\
Genitive & unrihtwīsra & unrihtwīsra & unrihtwīsra \\
Dative & unrihtwīsum & unrihtwīsum & unrihtwīsum \\
Accusative & unrihtwīse & unrihtwīs & unrihtwīsa \\
Instrumental & as dative & as dative & as dative \\
\bottomrule
\end{tabular}
\par\endgroup
% course-panel:T15-P02:end
% mdtable:T15:end

% teaching: NEW-U10-ADJ-UNRIHTWIS-WEAK
% table: NEW-U10-ADJ-UNRIHTWIS-WEAK
% mdtable:T16:begin
% course-table-style: 2026-10
% course-panel:T16-P01:begin
\par\begingroup
\fontsize{10.5}{13.2}\selectfont
\setlength{\LTleft}{0pt}\setlength{\LTright}{\fill}
\noindent
\begin{tabular}{@{}>{\raggedright\arraybackslash}p{22.0mm}>{\raggedright\arraybackslash}p{31.188mm}>{\raggedright\arraybackslash}p{31.188mm}>{\raggedright\arraybackslash}p{31.188mm}@{}}
\toprule
\textbf{Case} & \textbf{Masculine singular} & \textbf{Neuter singular} & \textbf{Feminine singular} \\
\midrule
Nominative & unrihtwīsa & unrihtwīse & unrihtwīse \\
Genitive & unrihtwīsan & unrihtwīsan & unrihtwīsan \\
Dative & unrihtwīsan & unrihtwīsan & unrihtwīsan \\
Accusative & unrihtwīsan & unrihtwīse & unrihtwīsan \\
Instrumental & unrihtwīsan & unrihtwīsan & as dative \\
\bottomrule
\end{tabular}
\par\endgroup
% course-panel:T16-P01:end
\par\addvspace{5pt}
% course-panel:T16-P02:begin
\par\begingroup
\fontsize{10.5}{13.2}\selectfont
\setlength{\LTleft}{0pt}\setlength{\LTright}{\fill}
\noindent
\begin{tabular}{@{}>{\raggedright\arraybackslash}p{22.0mm}>{\raggedright\arraybackslash}p{31.188mm}>{\raggedright\arraybackslash}p{31.188mm}>{\raggedright\arraybackslash}p{31.188mm}@{}}
\toprule
\textbf{Case} & \textbf{Masculine plural} & \textbf{Neuter plural} & \textbf{Feminine plural} \\
\midrule
Nominative & unrihtwīsan & unrihtwīsan & unrihtwīsan \\
Genitive & unrihtwīsra & unrihtwīsra & unrihtwīsra \\
Dative & unrihtwīsum & unrihtwīsum & unrihtwīsum \\
Accusative & unrihtwīsan & unrihtwīsan & unrihtwīsan \\
Instrumental & as dative & as dative & as dative \\
\bottomrule
\end{tabular}
\par\endgroup
% course-panel:T16-P02:end
% mdtable:T16:end

% mdparagraph:P050:begin
\textbf{Synfull},\BookPoint{bk-ch10-learner-0c3379bd0a71} “sinful”, likewise has no -u in those strong cells. Preserve the source headword's -ll in bare forms and before vowel endings. Before consonantal -ne, -re and -ra, simplify ll to l: \textbf{synfulne, synfulre, synfulra}. This spelling change affects the genitive plural in both declensions.
% mdparagraph:P050:end

% teaching: NEW-U10-ADJ-SYNFULL-STRONG
% table: NEW-U10-ADJ-SYNFULL-STRONG
% mdtable:T17:begin
% course-table-style: 2026-10
% course-panel:T17-P01:begin
\par\begingroup
\fontsize{10.5}{13.2}\selectfont
\setlength{\LTleft}{0pt}\setlength{\LTright}{\fill}
\noindent
\begin{tabular}{@{}>{\raggedright\arraybackslash}p{22.0mm}>{\raggedright\arraybackslash}p{31.188mm}>{\raggedright\arraybackslash}p{31.188mm}>{\raggedright\arraybackslash}p{31.188mm}@{}}
\toprule
\textbf{Case} & \textbf{Masculine singular} & \textbf{Neuter singular} & \textbf{Feminine singular} \\
\midrule
Nominative & synfull & synfull & synfull \\
Genitive & synfulles & synfulles & synfulre \\
Dative & synfullum & synfullum & synfulre \\
Accusative & synfulne & synfull & synfulle \\
Instrumental & synfulle & synfulle & as dative \\
\bottomrule
\end{tabular}
\par\endgroup
% course-panel:T17-P01:end
\par\addvspace{5pt}
% course-panel:T17-P02:begin
\par\begingroup
\fontsize{10.5}{13.2}\selectfont
\setlength{\LTleft}{0pt}\setlength{\LTright}{\fill}
\noindent
\begin{tabular}{@{}>{\raggedright\arraybackslash}p{22.0mm}>{\raggedright\arraybackslash}p{31.188mm}>{\raggedright\arraybackslash}p{31.188mm}>{\raggedright\arraybackslash}p{31.188mm}@{}}
\toprule
\textbf{Case} & \textbf{Masculine plural} & \textbf{Neuter plural} & \textbf{Feminine plural} \\
\midrule
Nominative & synfulle & synfull & synfulla \\
Genitive & synfulra & synfulra & synfulra \\
Dative & synfullum & synfullum & synfullum \\
Accusative & synfulle & synfull & synfulla \\
Instrumental & as dative & as dative & as dative \\
\bottomrule
\end{tabular}
\par\endgroup
% course-panel:T17-P02:end
% mdtable:T17:end

% teaching: NEW-U10-ADJ-SYNFULL-WEAK
% table: NEW-U10-ADJ-SYNFULL-WEAK
% mdtable:T18:begin
% course-table-style: 2026-10
% course-panel:T18-P01:begin
\par\begingroup
\fontsize{10.5}{13.2}\selectfont
\setlength{\LTleft}{0pt}\setlength{\LTright}{\fill}
\noindent
\begin{tabular}{@{}>{\raggedright\arraybackslash}p{22.0mm}>{\raggedright\arraybackslash}p{31.188mm}>{\raggedright\arraybackslash}p{31.188mm}>{\raggedright\arraybackslash}p{31.188mm}@{}}
\toprule
\textbf{Case} & \textbf{Masculine singular} & \textbf{Neuter singular} & \textbf{Feminine singular} \\
\midrule
Nominative & synfulla & synfulle & synfulle \\
Genitive & synfullan & synfullan & synfullan \\
Dative & synfullan & synfullan & synfullan \\
Accusative & synfullan & synfulle & synfullan \\
Instrumental & synfullan & synfullan & as dative \\
\bottomrule
\end{tabular}
\par\endgroup
% course-panel:T18-P01:end
\par\addvspace{5pt}
% course-panel:T18-P02:begin
\par\begingroup
\fontsize{10.5}{13.2}\selectfont
\setlength{\LTleft}{0pt}\setlength{\LTright}{\fill}
\noindent
\begin{tabular}{@{}>{\raggedright\arraybackslash}p{22.0mm}>{\raggedright\arraybackslash}p{31.188mm}>{\raggedright\arraybackslash}p{31.188mm}>{\raggedright\arraybackslash}p{31.188mm}@{}}
\toprule
\textbf{Case} & \textbf{Masculine plural} & \textbf{Neuter plural} & \textbf{Feminine plural} \\
\midrule
Nominative & synfullan & synfullan & synfullan \\
Genitive & synfulra & synfulra & synfulra \\
Dative & synfullum & synfullum & synfullum \\
Accusative & synfullan & synfullan & synfullan \\
Instrumental & as dative & as dative & as dative \\
\bottomrule
\end{tabular}
\par\endgroup
% course-panel:T18-P02:end
% mdtable:T18:end

\reserve{7\baselineskip}
% mdparagraph:P051:begin
\subsection*{The right side}\BookHeading{nav-ch10-learner-0011}{subsubsection}{The right side}{0}
% mdparagraph:P051:end

% mdparagraph:P052:begin
\textbf{Swīðra}, “right” in the right-hand expression, has weak comparative morphology. Its stem is \textbf{swīðr-}. Use the complete weak table; no positive strong swīð paradigm is introduced here. In the genitive plural, the cluster from swīðr- plus -ra is simplified to \textbf{swīðra}, as in the earlier comparative pattern. The shape is a regular derivation under the cited cluster rule, not an additional strong adjective.
% mdparagraph:P052:end

% teaching: NEW-U10-ADJ-SWITHRA-WEAK
% table: NEW-U10-ADJ-SWITHRA-WEAK
% mdtable:T19:begin
% course-table-style: 2026-10
% course-panel:T19-P01:begin
\par\begingroup
\fontsize{10.5}{13.2}\selectfont
\setlength{\LTleft}{0pt}\setlength{\LTright}{\fill}
\noindent
\begin{tabular}{@{}>{\raggedright\arraybackslash}p{22.0mm}>{\raggedright\arraybackslash}p{31.188mm}>{\raggedright\arraybackslash}p{31.188mm}>{\raggedright\arraybackslash}p{31.188mm}@{}}
\toprule
\textbf{Case} & \textbf{Masculine singular} & \textbf{Neuter singular} & \textbf{Feminine singular} \\
\midrule
Nominative & swīðra & swīðre & swīðre \\
Genitive & swīðran & swīðran & swīðran \\
Dative & swīðran & swīðran & swīðran \\
Accusative & swīðran & swīðre & swīðran \\
Instrumental & swīðran & swīðran & as dative \\
\bottomrule
\end{tabular}
\par\endgroup
% course-panel:T19-P01:end
\par\addvspace{5pt}
% course-panel:T19-P02:begin
\par\begingroup
\fontsize{10.5}{13.2}\selectfont
\setlength{\LTleft}{0pt}\setlength{\LTright}{\fill}
\noindent
\begin{tabular}{@{}>{\raggedright\arraybackslash}p{22.0mm}>{\raggedright\arraybackslash}p{31.188mm}>{\raggedright\arraybackslash}p{31.188mm}>{\raggedright\arraybackslash}p{31.188mm}@{}}
\toprule
\textbf{Case} & \textbf{Masculine plural} & \textbf{Neuter plural} & \textbf{Feminine plural} \\
\midrule
Nominative & swīðran & swīðran & swīðran \\
Genitive & swīðra & swīðra & swīðra \\
Dative & swīðrum & swīðrum & swīðrum \\
Accusative & swīðran & swīðran & swīðran \\
Instrumental & as dative & as dative & as dative \\
\bottomrule
\end{tabular}
\par\endgroup
% course-panel:T19-P02:end
% mdtable:T19:end

\reserve{7\baselineskip}
% mdparagraph:P053:begin
\subsection*{Adjective vocabulary and derivation}\BookHeading{nav-ch10-learner-0012}{subsubsection}{Adjective vocabulary and derivation}{0}
% mdparagraph:P053:end

% mdparagraph:P054:begin
For \textbf{untrum, sellic, cwic} and \textbf{nacod}, replace the exact stem \textbf{hungrig-} by the listed stem in every cell of both hungrig tables. Keep the entire remainder of the model form. Retain the listed internal vowels; do not shorten untrum, sellic or nacod automatically. For \textbf{ēce}, replace \textbf{dēorwierð-} by \textbf{ēc-} in every corresponding cell of both dēorwierðe tables, retaining the final -e wherever the model has it. These operations include all genders, both numbers and both declensions; instrumental cells retain their reference limit.
% mdparagraph:P054:end

% teaching: NEW-U10-VOC-HUNGRIG
% diagnostic: NEW-U10-DIAG-HUNGRIG
% teaching: NEW-U10-VOC-DEORWIERTH
% diagnostic: NEW-U10-DIAG-DEORWIERTH
% teaching: NEW-U10-VOC-UNRIHTWIS
% diagnostic: NEW-U10-DIAG-UNRIHTWIS
% teaching: NEW-U10-VOC-SYNFULL
% diagnostic: NEW-U10-DIAG-SYNFULL
% teaching: NEW-U10-VOC-UNTRUM
% teaching: NEW-U10-ADJ-UNTRUM-STRONG
% teaching: NEW-U10-ADJ-UNTRUM-WEAK
% diagnostic: NEW-U10-DIAG-UNTRUM
% teaching: NEW-U10-VOC-SELLIC
% teaching: NEW-U10-ADJ-SELLIC-STRONG
% teaching: NEW-U10-ADJ-SELLIC-WEAK
% diagnostic: NEW-U10-DIAG-SELLIC
% teaching: NEW-U10-VOC-ECE
% teaching: NEW-U10-ADJ-ECE-STRONG
% teaching: NEW-U10-ADJ-ECE-WEAK
% diagnostic: NEW-U10-DIAG-ECE
% teaching: NEW-U10-VOC-CWIC
% teaching: NEW-U10-ADJ-CWIC-STRONG
% teaching: NEW-U10-ADJ-CWIC-WEAK
% diagnostic: NEW-U10-DIAG-CWIC
% teaching: NEW-U10-VOC-NACOD
% teaching: NEW-U10-ADJ-NACOD-STRONG
% teaching: NEW-U10-ADJ-NACOD-WEAK
% diagnostic: NEW-U10-DIAG-NACOD
% teaching: NEW-U10-VOC-SWITHRA
% diagnostic: NEW-U10-DIAG-SWITHRA
% mdtable:T20:begin
\par\begingroup
\fontsize{10.5}{13.2}\selectfont
\setlength{\LTleft}{0pt}\setlength{\LTright}{\fill}
\noindent
\begin{tabular}{@{}>{\raggedright\arraybackslash}p{46.15mm}>{\raggedright\arraybackslash}p{19mm}>{\raggedright\arraybackslash}p{53.22mm}@{}}
\toprule
\textbf{Headword and meaning} & \textbf{Stem} & \textbf{Complete model and operation} \\
\midrule
hungrig, hungry; hungry people & hungrig & Both full tables above; internal vowel retained. \\
dēorwierðe, precious; valuable & dēorwierð & Both full tables above; lexical -e replaced before other endings. \\
unrihtwīs, unrighteous; unrighteous people & unrihtwīs & Both full tables above; no -u in the long-stem cells. \\
synfull, sinful & synfull & Both full tables above; ll becomes l before -ne, -re, -ra. \\
untrum, sick; sick people & untrum & Replace hungrig- by untrum- in both tables. \\
sellic, strange; remarkable & sellic & Replace hungrig- by sellic- in both tables. \\
ēce, eternal & ēc & Replace dēorwierð- by ēc- in both tables. \\
cwic, alive; living & cwic & Replace hungrig- by cwic- in both tables. \\
nacod, naked; naked people & nacod & Replace hungrig- by nacod- in both tables. \\
swīðra, right, in a right-hand/\allowbreak{}side expression & swīðr & Full weak comparative table only; genitive plural swīðra. \\
\bottomrule
\end{tabular}
\par\endgroup

% mdtable:T20:end
% mdparagraph:P055:begin
For example, hungrigre gives \textbf{untrumre} at the same feminine singular genitive or dative coordinate; weak hungrigan gives \textbf{nacodan} at each corresponding coordinate. Dēorwierðe gives \textbf{ēce} where its lexical -e is retained, while dēorwierðum gives \textbf{ēcum}. Choosing a form still requires the case, gender, number and strong or weak use; matching one ending does not identify a unique cell.
% mdparagraph:P055:end

% teaching: NEW-U10-GRAM-u10-substantival-personal-adjectives
% mdparagraph:P056:begin
The\BookPoint{bk-ch10-learner-f4665e71eb65} article and a weak adjective can name people without an expressed noun: \textbf{the hungry, the sick, the naked, the unrighteous}. Use hungrig, untrum, nacod and unrihtwīs respectively. The whole phrase has the case required by its clause and the number intended. A plural group uses the plural adjective; a phrase describing one expressed person uses singular agreement.
% mdparagraph:P056:end

\reserve{8\baselineskip}
% example: NEW-U10-EXAMPLE14
% mdparagraph:P057:begin
Fēd ðā hungrigan.
% mdparagraph:P057:end

% mdparagraph:P058:begin
Feed the hungry.
% mdparagraph:P058:end

% mdparagraph:P059:begin
Fēd addresses one person. Ðā hungrigan is the accusative plural object, naming the hungry people. The number of the command belongs to its addressee, not to its object. With an expressed masculine singular noun, “the hungry stranger” instead has singular agreement, such as nominative \textbf{se hungriga cuma}.
% mdparagraph:P059:end

\reserve{7\baselineskip}
% mdparagraph:P060:begin
\section*{4. A member of a group and a quantity}\BookHeading{nav-ch10-learner-0013}{subsection}{4. A member of a group and a quantity}{1}
% mdparagraph:P060:end

% teaching: NEW-U10-GRAM-u10-partitive-group
% mdparagraph:P061:begin
A\BookPoint{bk-ch10-learner-2df8a95ea983} word naming a member or some members can take a \textbf{genitive plural whole}: ān, “one”, sum, “some”, and hwilc, “which”, can stand as heads followed by the group from which the selection is made. The head takes its case from the clause; the group remains genitive. Gender and number of the selecting word follow the intended member or members.
% mdparagraph:P061:end

% mdparagraph:P062:begin
Distinguish this from an agreeing word before an expressed noun. In \textbf{ān munuc}, “one monk”, ān modifies nominative singular munuc. In \textbf{ān ðāra munuca}, “one of the monks”, ān names the selected member and ðāra munuca is the genitive plural whole. This cardinal use of ān does not introduce the wider indefinite-article system.
% mdparagraph:P062:end

\reserve{8\baselineskip}
% example: NEW-U10-EXAMPLE04
% mdparagraph:P063:begin
Ān ðāra munuca cōm tō ðǣm hūse.
% mdparagraph:P063:end

% mdparagraph:P064:begin
One of the monks came to the house.
% mdparagraph:P064:end

% mdparagraph:P065:begin
Ān is nominative singular, the subject of cōm; ðāra munuca remains genitive plural. The same relation is available with plural sum, “some of”, and with interrogative hwilc, “which of”. A question can select one member or several; choose the corresponding singular or plural form and verb. These nominal genitive groups extend the earlier selected “of you/\allowbreak{}of them” phrases, without introducing every possible personal-pronoun genitive.
% mdparagraph:P065:end

% teaching: NEW-U10-GRAM-u10-fela-genitive
% teaching: NEW-U10-VOC-FELA
% mdparagraph:P066:begin
\textbf{Fela},\BookPoint{bk-ch10-learner-92c9166ec67f} “many, a large number of”, is invariant in the selected construction and takes a genitive plural. Do not make it agree like an ordinary attributive adjective or a numeral.
% mdparagraph:P066:end

\reserve{8\baselineskip}
% example: NEW-U10-EXAMPLE03
% mdparagraph:P067:begin
Ic hæbbe fela scēapa.
% mdparagraph:P067:end

% mdparagraph:P068:begin
I have many sheep.
% mdparagraph:P068:end

% mdparagraph:P069:begin
Scēapa is genitive plural after fela. In the ordinary counted phrase \textbf{twā scēap}, “two sheep”, the numeral and noun instead share the case required by the clause. The genitive after fela is therefore not the form of scēap required directly by habban. Higher numerals and singular manig are outside the present construction.
% mdparagraph:P069:end

\reserve{7\baselineskip}
% mdparagraph:P070:begin
\section*{5. Datives and experiencers}\BookHeading{nav-ch10-learner-0014}{subsection}{5. Datives and experiencers}{1}\BookLabel{bk-ch10-learner-f12b5c0271f8}
% mdparagraph:P070:end

% teaching: NEW-U10-GRAM-u10-dative-roles
% teaching: NEW-U10-GRAM-u10-don-beneficiary
% teaching: NEW-U10-VOC-DON
% mdparagraph:P071:begin
The dative has several functions. Identify the construction that requires it:
% mdparagraph:P071:end

\begin{itemize}
% mdparagraph:P072:begin
\item A \textbf{recipient} receives something: with sęllan, “give”, the thing is accusative and the recipient dative; tōgelǣdan, “bring”, follows the same case arrangement.
% mdparagraph:P072:end
% mdparagraph:P073:begin
\item A \textbf{beneficiary or affected person} is the person for or to whom a deed is done. In the selected use of dōn, “do”, the deed or “it” is accusative and the person dative.
% mdparagraph:P073:end
% mdparagraph:P074:begin
\item A \textbf{governed person} has the case required by a particular verb or adjective. Continue to use helpan, fręmman and behēfe within their learned meanings and dative frames. Forwiernan and unnan likewise require a dative person, although their thing is genitive.
% mdparagraph:P074:end
% mdparagraph:P075:begin
\item An \textbf{experiencer} is the person to whom something seems a certain way. The construction with ðyncan below uses a dative experiencer.
% mdparagraph:P075:end
\end{itemize}

\reserve{8\baselineskip}
% example: NEW-U10-EXAMPLE05
% mdparagraph:P076:begin
Se fiscere hilpð ðǣm cuman.
% mdparagraph:P076:end

% mdparagraph:P077:begin
The fisherman helps the stranger.
% mdparagraph:P077:end

% mdparagraph:P078:begin
The person helped is dative after hilpð. Its English translation has no preposition.
% mdparagraph:P078:end

\reserve{8\baselineskip}
% example: NEW-U10-EXAMPLE06
% mdparagraph:P079:begin
Sēo gāt is ðǣm scēaphierde behēfu.
% mdparagraph:P079:end

% mdparagraph:P080:begin
The goat is useful to the shepherd.
% mdparagraph:P080:end

% mdparagraph:P081:begin
The goat is the nominative subject. Behēfu agrees with feminine nominative singular gāt; ðǣm scēaphierde is dative, the person to whom the goat is useful. The adjective's agreement and the case of its complement are separate requirements.
% mdparagraph:P081:end

% teaching: NEW-U10-GRAM-u10-thyncan-experiencer
% teaching: NEW-U10-VOC-THYNCAN
% mdparagraph:P082:begin
In\BookPoint{bk-ch10-learner-5ca7ed8d1974} \textbf{ðyncan}, “seem”, the selected construction has a dative experiencer, a singular neuter nominative subject \textbf{hit, ðæt} or \textbf{ðis}, and the strong neuter predicate \textbf{gōd}. The finite verb is third-person singular: present ðyncð, preterite ðūhte. This is the complete present/\allowbreak{}preterite system needed for that defined construction.
% mdparagraph:P082:end

% teaching: NEW-U10-PAR-thyncan-EXPERIENCER-IND
% table: NEW-U10-PAR-thyncan-EXPERIENCER-IND
% mdtable:T21:begin
% course-table-style: 2026-10
% course-panel:T21-P01:begin
\par\begingroup
\fontsize{10.5}{13.2}\selectfont
\setlength{\LTleft}{0pt}\setlength{\LTright}{\fill}
\noindent
\begin{tabular}{@{}>{\raggedright\arraybackslash}p{39.459mm}>{\raggedright\arraybackslash}p{39.459mm}>{\raggedright\arraybackslash}p{39.459mm}@{}}
\toprule
\textbf{Person and number} & \textbf{Present} & \textbf{Preterite} \\
\midrule
3 sg. & ðyncð & ðūhte \\
\bottomrule
\end{tabular}
\par\endgroup
% course-panel:T21-P01:end
% mdtable:T21:end
\reserve{8\baselineskip}
% example: NEW-U10-EXAMPLE07
% mdparagraph:P083:begin
Mē ðyncð ðis gōd.
% mdparagraph:P083:end

% mdparagraph:P084:begin
This seems good to me.
% mdparagraph:P084:end

% mdparagraph:P085:begin
Ðis is the nominative subject: “this” seems good. Mē is the dative experiencer. A plural experiencer does not make this singular subject plural. This use has an expressed nominative subject; it is not the subjectless construction that follows. Do not infer a general content-clause or passive construction from these forms.
% mdparagraph:P085:end

% teaching: NEW-U10-GRAM-u10-hyngran-impersonal
% teaching: NEW-U10-VOC-HYNGRAN
% mdparagraph:P086:begin
\textbf{Hyngran},\BookPoint{bk-ch10-learner-f11ca067063f} “be hungry”, is used here impersonally, with \textbf{no nominative subject}. The experiencer is accusative. The verb remains third-person singular even if the experiencer denotes several people: present hyngreð, preterite hyngrede. These two cells constitute the complete assigned impersonal indicative system; do not construct first- or second-person forms for it.
% mdparagraph:P086:end

% teaching: NEW-U10-PAR-hyngran-IMPERSONAL-IND
% table: NEW-U10-PAR-hyngran-IMPERSONAL-IND
% mdtable:T22:begin
% course-table-style: 2026-10
% course-panel:T22-P01:begin
\par\begingroup
\fontsize{10.5}{13.2}\selectfont
\setlength{\LTleft}{0pt}\setlength{\LTright}{\fill}
\noindent
\begin{tabular}{@{}>{\raggedright\arraybackslash}p{39.459mm}>{\raggedright\arraybackslash}p{39.459mm}>{\raggedright\arraybackslash}p{39.459mm}@{}}
\toprule
\textbf{Person and number} & \textbf{Present} & \textbf{Preterite} \\
\midrule
3 sg. & hyngreð & hyngrede \\
\bottomrule
\end{tabular}
\par\endgroup
% course-panel:T22-P01:end
% mdtable:T22:end
\reserve{8\baselineskip}
% example: NEW-U10-EXAMPLE08
% mdparagraph:P087:begin
Hine hyngrede.
% mdparagraph:P087:end

% mdparagraph:P088:begin
He was hungry.
% mdparagraph:P088:end

% mdparagraph:P089:begin
Hine is accusative, although English requires “he” as its subject. Compare the ordinary copular construction \textbf{Hē wæs hungrig}, “he was hungry”: hē is nominative, wæs agrees with it, and hungrig is a predicate adjective. Both express hunger, but their grammar differs. The accusative experiencer pattern is the one to use when a task expressly requires impersonal hyngran.
% mdparagraph:P089:end

\reserve{7\baselineskip}
% mdparagraph:P090:begin
\section*{6. Accompaniment, means and the instrumental}\BookHeading{nav-ch10-learner-0015}{subsection}{6. Accompaniment, means and the instrumental}{1}
% mdparagraph:P090:end

% teaching: NEW-U10-VOC-MID
% teaching: NEW-U10-GRAM-u10-mid-means
% teaching: NEW-U10-GRAM-u10-fedan-accusative-means
% mdparagraph:P091:begin
\textbf{Mid + dative} can express accompaniment, “with”, or means, “with, by means of”. The preposition governs the dative in both uses. The relation must be read from the situation.
% mdparagraph:P091:end

% mdparagraph:P092:begin
\textbf{Fēdan}, “feed, nourish”, takes the person or animal fed in the accusative. Food used for feeding can be expressed by mid + dative.
% mdparagraph:P092:end

\reserve{8\baselineskip}
% example: NEW-U10-EXAMPLE09
% mdparagraph:P093:begin
Se munuc cōm mid ðǣm cuman.
% mdparagraph:P093:end

% mdparagraph:P094:begin
The monk came with the stranger.
% mdparagraph:P094:end

% mdparagraph:P095:begin
Here mid names the accompanying person.
% mdparagraph:P095:end

\reserve{8\baselineskip}
% example: NEW-U10-EXAMPLE10
% mdparagraph:P096:begin
Se munuc fēdde ðone cuman mid gōdum męte.
% mdparagraph:P096:end

% mdparagraph:P097:begin
The monk fed the stranger with good food.
% mdparagraph:P097:end

% mdparagraph:P098:begin
Here mid names the means of feeding. Ðone cuman is the accusative person fed; gōdum męte is dative after mid. This is the ordinary prose means construction.
% mdparagraph:P098:end

% teaching: NEW-U10-GRAM-u10-bare-instrument
% teaching: NEW-U10-GRAM-u10-instrumental-morphology
% mdparagraph:P099:begin
Means\BookPoint{bk-ch10-learner-accee2a2f1b3} may also be expressed by a \textbf{bare instrumental}, without a preposition. This use is especially characteristic of poetry; mid + dative is the usual prose choice. The present exercises deliberately practise the bare construction in the feeding expression.
% mdparagraph:P099:end

% mdparagraph:P100:begin
The\BookPoint{bk-ch10-learner-c80743ab2698} demonstrative has distinct masculine/\allowbreak{}neuter singular instrumental \textbf{ðȳ}, with variant \textbf{ðon}. Strong gōd has instrumental \textbf{gōde} in those two genders. Compare their dative forms ðǣm and gōdum. Only these distinct instrumental forms are newly required for composition here; the previously displayed instruments of ðes and hwā remain reference.
% mdparagraph:P100:end

% teaching: NEW-U10-PAR-se-SELECTED-INSTR
% table: NEW-U10-PAR-se-SELECTED-INSTR
% mdtable:T23:begin
\par\begingroup
\fontsize{10.5}{13.2}\selectfont
\setlength{\LTleft}{0pt}\setlength{\LTright}{\fill}
\noindent
\begin{tabular}{@{}>{\raggedright\arraybackslash}p{38.0mm}>{\raggedright\arraybackslash}p{34.94mm}>{\raggedright\arraybackslash}p{45.43mm}@{}}
\toprule
\textbf{Gender /\allowbreak{} number} & \textbf{Instrumental} & \textbf{Instrumental variant} \\
\midrule
Masculine singular & ðȳ & ðon \\
Neuter singular & ðȳ & ðon \\
\bottomrule
\end{tabular}
\par\endgroup

% mdtable:T23:end
% teaching: NEW-U10-PAR-god-adjective-SELECTED-INSTR
% table: NEW-U10-PAR-god-adjective-SELECTED-INSTR
% mdtable:T24:begin
\par\begingroup
\fontsize{10.5}{13.2}\selectfont
\setlength{\LTleft}{0pt}\setlength{\LTright}{\fill}
\noindent
\begin{tabular}{@{}>{\raggedright\arraybackslash}p{56.08mm}>{\raggedright\arraybackslash}p{65.11mm}@{}}
\toprule
\textbf{Gender /\allowbreak{} number} & \textbf{Strong instrumental} \\
\midrule
Masculine singular & gōde \\
Neuter singular & gōde \\
\bottomrule
\end{tabular}
\par\endgroup

% mdtable:T24:end
% mdparagraph:P101:begin
Nouns\BookPoint{bk-ch10-learner-458dfac355a0} use the dative-shaped form for this instrumental relation: \textbf{męte} does not acquire a new ending. Feminine singular and all plural modifiers likewise share their forms with the dative. Thus the feminine singular forms are \textbf{ðǣre} and strong \textbf{gōdre}, and the plural forms \textbf{ðǣm} and \textbf{gōdum}. Shared shape does not remove the distinction of function. In a weak adjective, use the corresponding weak dative shape; the new adjectives' separate instrumental cells remain reference.
% mdparagraph:P101:end

\reserve{8\baselineskip}
% example: NEW-U10-EXAMPLE11
% mdparagraph:P102:begin
Se munuc fēdde ðone cuman gōde męte.
% mdparagraph:P102:end

% mdparagraph:P103:begin
The monk fed the stranger with good food.
% mdparagraph:P103:end

% mdparagraph:P104:begin
Gōde is a distinct strong masculine singular instrumental. Męte has the same shape as the dative singular, but the phrase expresses bare means.
% mdparagraph:P104:end

\reserve{8\baselineskip}
% example: NEW-U10-EXAMPLE12
% mdparagraph:P105:begin
Se munuc fēdde ðone cuman ðȳ męte.
% mdparagraph:P105:end

% mdparagraph:P106:begin
The monk fed the stranger with that food.
% mdparagraph:P106:end

% mdparagraph:P107:begin
Ðȳ is the demonstrative instrumental agreeing with masculine męte. Adding mid changes the construction: use the dative \textbf{mid ðǣm męte}. Do not treat mid + a distinct instrumental as the ordinary prose pattern taught here. An instrumental meaning, a distinct morphological form, and the case governed by a preposition are three separate facts.
% mdparagraph:P107:end

\reserve{7\baselineskip}
% mdparagraph:P108:begin
\section*{7. Complete verb systems and their constructions}\BookHeading{nav-ch10-learner-0016}{subsection}{7. Complete verb systems and their constructions}{1}
% mdparagraph:P108:end

% mdparagraph:P109:begin
The following tables give complete present and preterite indicative systems, except where the expressly defined impersonal or experiencer system has already been given. Commands are included only for the verbs with a command table. Unnan has its supported indicative and complete subjunctive system. Infinitives used as dictionary headwords identify verbs; this alone does not assign every infinitive construction. The newly usable plain infinitives are specified in §8.
% mdparagraph:P109:end

% mdparagraph:P110:begin
Plural subjects use the common plural ending in each indicative tense; dual subjects also use the plural finite form. A direct command's number belongs to the addressee. The object's number does not determine it.
% mdparagraph:P110:end

\reserve{7\baselineskip}
% mdparagraph:P111:begin
\subsection*{Weak class I: stem changes and dental preterites}\BookHeading{nav-ch10-learner-0017}{subsubsection}{Weak class I: stem changes and dental preterites}{0}
% mdparagraph:P111:end

% mdparagraph:P112:begin
The ordinary selected present endings are singular -e, -st, -ð and plural -að, with the contractions printed below. The preterite uses a dental stem followed by singular -e, -est, -e and plural -on. A long vowel followed by d gives the dental sequence \textbf{dd} in fēdde, scrȳdde and tōgelǣdde. Their second- and third-singular present forms instead have \textbf{tst} and \textbf{t}. Keep the long vowel in these selected forms.
% mdparagraph:P112:end

% teaching: NEW-U10-VOC-FEDAN
% teaching: NEW-U10-GRAM-u10-fedan-indicative
% teaching: NEW-U10-GRAM-u10-fedan-commands
% mdparagraph:P113:begin
\textbf{Fēdan}, “feed, nourish”, has present stem fēd- and preterite stem fēdd-. It feeds an accusative person or animal; its means expressions were given in §6. The singular command is the bare stem fēd; the plural command equals fēdað.
% mdparagraph:P113:end

\reserve{7\baselineskip}
% mdparagraph:P114:begin
\subsection*{Fēdan: indicative}\BookHeading{nav-ch10-learner-0018}{subsubsection}{Fēdan: indicative}{0}
% mdparagraph:P114:end

% teaching: NEW-U10-PAR-fedan-IND
% table: NEW-U10-PAR-fedan-IND
% mdtable:T25:begin
% course-table-style: 2026-10
% course-panel:T25-P01:begin
\par\begingroup
\fontsize{10.5}{13.2}\selectfont
\setlength{\LTleft}{0pt}\setlength{\LTright}{\fill}
\noindent
\begin{tabular}{@{}>{\raggedright\arraybackslash}p{39.459mm}>{\raggedright\arraybackslash}p{39.459mm}>{\raggedright\arraybackslash}p{39.459mm}@{}}
\toprule
\textbf{Person and number} & \textbf{Present} & \textbf{Preterite} \\
\midrule
1 sg. & fēde & fēdde \\
2 sg. & fētst & fēddest \\
3 sg. & fēt & fēdde \\
1–3 pl. & fēdað & fēddon \\
\bottomrule
\end{tabular}
\par\endgroup
% course-panel:T25-P01:end
% mdtable:T25:end
\reserve{7\baselineskip}
% mdparagraph:P115:begin
\subsection*{Fēdan: commands}\BookHeading{nav-ch10-learner-0019}{subsubsection}{Fēdan: commands}{0}
% mdparagraph:P115:end

% teaching: NEW-U10-PAR-fedan-IMP
% table: NEW-U10-PAR-fedan-IMP
% mdtable:T26:begin
% course-table-style: 2026-10
% course-panel:T26-P01:begin
\par\begingroup
\fontsize{10.5}{13.2}\selectfont
\setlength{\LTleft}{0pt}\setlength{\LTright}{\fill}
\noindent
\begin{tabular}{@{}>{\raggedright\arraybackslash}p{63.342mm}>{\raggedright\arraybackslash}p{57.846mm}@{}}
\toprule
\textbf{Person and number} & \textbf{Command} \\
\midrule
2 sg. & fēd \\
2 pl. & fēdað \\
\bottomrule
\end{tabular}
\par\endgroup
% course-panel:T26-P01:end
% mdtable:T26:end
% teaching: NEW-U10-VOC-SCRYDAN
% teaching: NEW-U10-GRAM-u10-scrydan-indicative
% teaching: NEW-U10-GRAM-u10-scrydan-commands
% teaching: NEW-U10-GRAM-u10-scrydan-accusative
% mdparagraph:P116:begin
\textbf{Scrȳdan}, “clothe, dress”, takes an accusative person. Its present stem is scrȳd-, its preterite stem scrȳdd-. The commands are scrȳd and scrȳdað. The assigned construction names the person clothed; a further garment complement is not required in this unit.
% mdparagraph:P116:end

\reserve{7\baselineskip}
% mdparagraph:P117:begin
\subsection*{Scrȳdan: indicative}\BookHeading{nav-ch10-learner-0020}{subsubsection}{Scrȳdan: indicative}{0}
% mdparagraph:P117:end

% teaching: NEW-U10-PAR-scrydan-IND
% table: NEW-U10-PAR-scrydan-IND
% mdtable:T27:begin
% course-table-style: 2026-10
% course-panel:T27-P01:begin
\par\begingroup
\fontsize{10.5}{13.2}\selectfont
\setlength{\LTleft}{0pt}\setlength{\LTright}{\fill}
\noindent
\begin{tabular}{@{}>{\raggedright\arraybackslash}p{39.459mm}>{\raggedright\arraybackslash}p{39.459mm}>{\raggedright\arraybackslash}p{39.459mm}@{}}
\toprule
\textbf{Person and number} & \textbf{Present} & \textbf{Preterite} \\
\midrule
1 sg. & scrȳde & scrȳdde \\
2 sg. & scrȳtst & scrȳddest \\
3 sg. & scrȳt & scrȳdde \\
1–3 pl. & scrȳdað & scrȳddon \\
\bottomrule
\end{tabular}
\par\endgroup
% course-panel:T27-P01:end
% mdtable:T27:end
\reserve{7\baselineskip}
% mdparagraph:P118:begin
\subsection*{Scrȳdan: commands}\BookHeading{nav-ch10-learner-0021}{subsubsection}{Scrȳdan: commands}{0}
% mdparagraph:P118:end

% teaching: NEW-U10-PAR-scrydan-IMP
% table: NEW-U10-PAR-scrydan-IMP
% mdtable:T28:begin
% course-table-style: 2026-10
% course-panel:T28-P01:begin
\par\begingroup
\fontsize{10.5}{13.2}\selectfont
\setlength{\LTleft}{0pt}\setlength{\LTright}{\fill}
\noindent
\begin{tabular}{@{}>{\raggedright\arraybackslash}p{63.342mm}>{\raggedright\arraybackslash}p{57.846mm}@{}}
\toprule
\textbf{Person and number} & \textbf{Command} \\
\midrule
2 sg. & scrȳd \\
2 pl. & scrȳdað \\
\bottomrule
\end{tabular}
\par\endgroup
% course-panel:T28-P01:end
% mdtable:T28:end
% teaching: NEW-U10-VOC-TOGELAEDAN
% teaching: NEW-U10-GRAM-u10-togelaedan-indicative
% teaching: NEW-U10-GRAM-u10-togelaedan-recipient
% mdparagraph:P119:begin
\textbf{Tōgelǣdan}, “bring”, takes an accusative thing and a dative recipient. Its present stem tōgelǣd- becomes tōgelǣdd- in the preterite. Use this verb when bringing something to someone in the assigned exercises. Its complete indicative is now available; no new command forms are assigned.
% mdparagraph:P119:end

\reserve{7\baselineskip}
% mdparagraph:P120:begin
\subsection*{Tōgelǣdan: indicative}\BookHeading{nav-ch10-learner-0022}{subsubsection}{Tōgelǣdan: indicative}{0}
% mdparagraph:P120:end

% teaching: NEW-U10-PAR-togelaedan-IND
% table: NEW-U10-PAR-togelaedan-IND
% mdtable:T29:begin
% course-table-style: 2026-10
% course-panel:T29-P01:begin
\par\begingroup
\fontsize{10.5}{13.2}\selectfont
\setlength{\LTleft}{0pt}\setlength{\LTright}{\fill}
\noindent
\begin{tabular}{@{}>{\raggedright\arraybackslash}p{39.459mm}>{\raggedright\arraybackslash}p{39.459mm}>{\raggedright\arraybackslash}p{39.459mm}@{}}
\toprule
\textbf{Person and number} & \textbf{Present} & \textbf{Preterite} \\
\midrule
1 sg. & tōgelǣde & tōgelǣdde \\
2 sg. & tōgelǣtst & tōgelǣddest \\
3 sg. & tōgelǣt & tōgelǣdde \\
1–3 pl. & tōgelǣdað & tōgelǣddon \\
\bottomrule
\end{tabular}
\par\endgroup
% course-panel:T29-P01:end
% mdtable:T29:end
% teaching: NEW-U10-GRAM-u10-forwiernan-indicative
% mdparagraph:P121:begin
\textbf{Forwiernan}, “deny, refuse”, has present stem forwiern- and dental preterite stem forwiernd-. Preserve its dative person and genitive thing in every tense and person. The present third singular is forwiernð; the preterite third singular is forwiernde.
% mdparagraph:P121:end

\reserve{7\baselineskip}
% mdparagraph:P122:begin
\subsection*{Forwiernan: indicative}\BookHeading{nav-ch10-learner-0023}{subsubsection}{Forwiernan: indicative}{0}
% mdparagraph:P122:end

% teaching: NEW-U10-PAR-forwiernan-IND
% table: NEW-U10-PAR-forwiernan-IND
% mdtable:T30:begin
% course-table-style: 2026-10
% course-panel:T30-P01:begin
\par\begingroup
\fontsize{10.5}{13.2}\selectfont
\setlength{\LTleft}{0pt}\setlength{\LTright}{\fill}
\noindent
\begin{tabular}{@{}>{\raggedright\arraybackslash}p{39.459mm}>{\raggedright\arraybackslash}p{39.459mm}>{\raggedright\arraybackslash}p{39.459mm}@{}}
\toprule
\textbf{Person and number} & \textbf{Present} & \textbf{Preterite} \\
\midrule
1 sg. & forwierne & forwiernde \\
2 sg. & forwiernst & forwierndest \\
3 sg. & forwiernð & forwiernde \\
1–3 pl. & forwiernað & forwierndon \\
\bottomrule
\end{tabular}
\par\endgroup
% course-panel:T30-P01:end
% mdtable:T30:end
% teaching: NEW-U10-VOC-SCENCAN
% teaching: NEW-U10-GRAM-u10-scencan-indicative
% teaching: NEW-U10-GRAM-u10-scencan-recipient
% mdparagraph:P123:begin
\textbf{Scęncan}, “pour out, give to drink”, takes an accusative drink and a dative recipient. It retains ę and the present stem scęnc-. Its voiceless consonant is followed by \textbf{t} in the preterite stem scęnct-, giving scęncte, not a form in -de.
% mdparagraph:P123:end

\reserve{7\baselineskip}
% mdparagraph:P124:begin
\subsection*{Scęncan: indicative}\BookHeading{nav-ch10-learner-0024}{subsubsection}{Scęncan: indicative}{0}
% mdparagraph:P124:end

% teaching: NEW-U10-PAR-scencan-IND
% table: NEW-U10-PAR-scencan-IND
% mdtable:T31:begin
% course-table-style: 2026-10
% course-panel:T31-P01:begin
\par\begingroup
\fontsize{10.5}{13.2}\selectfont
\setlength{\LTleft}{0pt}\setlength{\LTright}{\fill}
\noindent
\begin{tabular}{@{}>{\raggedright\arraybackslash}p{39.459mm}>{\raggedright\arraybackslash}p{39.459mm}>{\raggedright\arraybackslash}p{39.459mm}@{}}
\toprule
\textbf{Person and number} & \textbf{Present} & \textbf{Preterite} \\
\midrule
1 sg. & scęnce & scęncte \\
2 sg. & scęncst & scęnctest \\
3 sg. & scęncð & scęncte \\
1–3 pl. & scęncað & scęncton \\
\bottomrule
\end{tabular}
\par\endgroup
% course-panel:T31-P01:end
% mdtable:T31:end
% teaching: NEW-U10-VOC-WYRCEAN
% teaching: NEW-U10-GRAM-u10-wyrcean-work
% teaching: NEW-U10-GRAM-u10-wyrcean-commands
% mdparagraph:P125:begin
The familiar \textbf{wyrcean} now also means “work, perform”, with an accusative work or deed, including cognate \textbf{weorc}, “work”. This ordinary sense differs from its earlier yield-fruit use. Its indicative forms are already learned. The two direct commands are given here; retain the different stem spellings wyrc- and wyrce-.
% mdparagraph:P125:end

\reserve{7\baselineskip}
% mdparagraph:P126:begin
\subsection*{Wyrcean: commands}\BookHeading{nav-ch10-learner-0025}{subsubsection}{Wyrcean: commands}{0}
% mdparagraph:P126:end

% teaching: NEW-U10-PAR-wyrcean-IMP
% table: NEW-U10-PAR-wyrcean-IMP
% mdtable:T32:begin
% course-table-style: 2026-10
% course-panel:T32-P01:begin
\par\begingroup
\fontsize{10.5}{13.2}\selectfont
\setlength{\LTleft}{0pt}\setlength{\LTright}{\fill}
\noindent
\begin{tabular}{@{}>{\raggedright\arraybackslash}p{63.342mm}>{\raggedright\arraybackslash}p{57.846mm}@{}}
\toprule
\textbf{Person and number} & \textbf{Command} \\
\midrule
2 sg. & wyrc \\
2 pl. & wyrceað \\
\bottomrule
\end{tabular}
\par\endgroup
% course-panel:T32-P01:end
% mdtable:T32:end
\reserve{7\baselineskip}
% mdparagraph:P127:begin
\subsection*{Weak class II}\BookHeading{nav-ch10-learner-0026}{subsubsection}{Weak class II}{0}
% mdparagraph:P127:end

% mdparagraph:P128:begin
The selected class II endings are present -ie, -ast, -að, plural -iað; preterite -ode, -odest, -ode, plural -odon. The class identifies the inflection, not the case of a verb's complement. Cook also records spellings with -ig(e)- or -g- and preterite -edon; these remain reference variants, not extra forms to reproduce here.
% mdparagraph:P128:end

% teaching: NEW-U10-VOC-ANDSWARIAN
% teaching: NEW-U10-GRAM-u10-andswarian-indicative
% teaching: NEW-U10-GRAM-u10-andswarian-dative
% mdparagraph:P129:begin
\textbf{Andswarian}, “answer, reply”, takes a \textbf{dative addressee}. A direct quotation can give the reply. Use the weak class II stem andswar- in this table; the separate weak class I variant andswerian is not introduced.
% mdparagraph:P129:end

\reserve{7\baselineskip}
% mdparagraph:P130:begin
\subsection*{Andswarian: indicative}\BookHeading{nav-ch10-learner-0027}{subsubsection}{Andswarian: indicative}{0}
% mdparagraph:P130:end

% teaching: NEW-U10-PAR-andswarian-IND
% table: NEW-U10-PAR-andswarian-IND
% mdtable:T33:begin
% course-table-style: 2026-10
% course-panel:T33-P01:begin
\par\begingroup
\fontsize{10.5}{13.2}\selectfont
\setlength{\LTleft}{0pt}\setlength{\LTright}{\fill}
\noindent
\begin{tabular}{@{}>{\raggedright\arraybackslash}p{39.459mm}>{\raggedright\arraybackslash}p{39.459mm}>{\raggedright\arraybackslash}p{39.459mm}@{}}
\toprule
\textbf{Person and number} & \textbf{Present} & \textbf{Preterite} \\
\midrule
1 sg. & andswarie & andswarode \\
2 sg. & andswarast & andswarodest \\
3 sg. & andswarað & andswarode \\
1–3 pl. & andswariað & andswarodon \\
\bottomrule
\end{tabular}
\par\endgroup
% course-panel:T33-P01:end
% mdtable:T33:end
% teaching: NEW-U10-VOC-GENEOSIAN
% teaching: NEW-U10-GRAM-u10-geneosian-indicative
% teaching: NEW-U10-GRAM-u10-geneosian-commands
% teaching: NEW-U10-GRAM-u10-geneosian-accusative
% mdparagraph:P131:begin
\textbf{Genēosian}, “visit”, takes an \textbf{accusative person} in the selected construction. Keep Cook's long ē in genēos-. The singular command is genēosa and the plural genēosiað. Do not transfer the genitive government of the related unprefixed nēosian to this selected prefixed use.
% mdparagraph:P131:end

\reserve{7\baselineskip}
% mdparagraph:P132:begin
\subsection*{Genēosian: indicative}\BookHeading{nav-ch10-learner-0028}{subsubsection}{Genēosian: indicative}{0}
% mdparagraph:P132:end

% teaching: NEW-U10-PAR-geneosian-IND
% table: NEW-U10-PAR-geneosian-IND
% mdtable:T34:begin
% course-table-style: 2026-10
% course-panel:T34-P01:begin
\par\begingroup
\fontsize{10.5}{13.2}\selectfont
\setlength{\LTleft}{0pt}\setlength{\LTright}{\fill}
\noindent
\begin{tabular}{@{}>{\raggedright\arraybackslash}p{39.459mm}>{\raggedright\arraybackslash}p{39.459mm}>{\raggedright\arraybackslash}p{39.459mm}@{}}
\toprule
\textbf{Person and number} & \textbf{Present} & \textbf{Preterite} \\
\midrule
1 sg. & genēosie & genēosode \\
2 sg. & genēosast & genēosodest \\
3 sg. & genēosað & genēosode \\
1–3 pl. & genēosiað & genēosodon \\
\bottomrule
\end{tabular}
\par\endgroup
% course-panel:T34-P01:end
% mdtable:T34:end
\reserve{7\baselineskip}
% mdparagraph:P133:begin
\subsection*{Genēosian: commands}\BookHeading{nav-ch10-learner-0029}{subsubsection}{Genēosian: commands}{0}
% mdparagraph:P133:end

% teaching: NEW-U10-PAR-geneosian-IMP
% table: NEW-U10-PAR-geneosian-IMP
% mdtable:T35:begin
% course-table-style: 2026-10
% course-panel:T35-P01:begin
\par\begingroup
\fontsize{10.5}{13.2}\selectfont
\setlength{\LTleft}{0pt}\setlength{\LTright}{\fill}
\noindent
\begin{tabular}{@{}>{\raggedright\arraybackslash}p{63.342mm}>{\raggedright\arraybackslash}p{57.846mm}@{}}
\toprule
\textbf{Person and number} & \textbf{Command} \\
\midrule
2 sg. & genēosa \\
2 pl. & genēosiað \\
\bottomrule
\end{tabular}
\par\endgroup
% course-panel:T35-P01:end
% mdtable:T35:end
\reserve{7\baselineskip}
% mdparagraph:P134:begin
\subsection*{Strong class I: gewītan}\BookHeading{nav-ch10-learner-0030}{subsubsection}{Strong class I: gewītan}{0}
% mdparagraph:P134:end

% teaching: NEW-U10-VOC-GEWITAN
% teaching: NEW-U10-GRAM-u10-gewitan-indicative
% teaching: NEW-U10-GRAM-u10-gewitan-departure
% mdparagraph:P135:begin
\textbf{Gewītan}, “depart, go away”, can be followed by \textbf{fram + dative}, naming the person or place left. This strong class I verb has long ī in the present, ā in the first and third singular preterite, and short i in the second singular and the plural preterite. The t-stem contracts in gewītst and gewīt; Cook's tt/\allowbreak{}t variation is a reference spelling variation. The forms below select single t.
% mdparagraph:P135:end

% mdparagraph:P136:begin
Distinguish this departure verb from the different verb meaning “know”. An accompanying reflexive dative is attested with departure, but is not required in the selected construction.
% mdparagraph:P136:end

\reserve{7\baselineskip}
% mdparagraph:P137:begin
\subsection*{Gewītan: indicative}\BookHeading{nav-ch10-learner-0031}{subsubsection}{Gewītan: indicative}{0}
% mdparagraph:P137:end

% teaching: NEW-U10-PAR-gewitan-IND
% table: NEW-U10-PAR-gewitan-IND
% mdtable:T36:begin
% course-table-style: 2026-10
% course-panel:T36-P01:begin
\par\begingroup
\fontsize{10.5}{13.2}\selectfont
\setlength{\LTleft}{0pt}\setlength{\LTright}{\fill}
\noindent
\begin{tabular}{@{}>{\raggedright\arraybackslash}p{39.459mm}>{\raggedright\arraybackslash}p{39.459mm}>{\raggedright\arraybackslash}p{39.459mm}@{}}
\toprule
\textbf{Person and number} & \textbf{Present} & \textbf{Preterite} \\
\midrule
1 sg. & gewīte & gewāt \\
2 sg. & gewītst & gewite \\
3 sg. & gewīt & gewāt \\
1–3 pl. & gewītað & gewiton \\
\bottomrule
\end{tabular}
\par\endgroup
% course-panel:T36-P01:end
% mdtable:T36:end
\reserve{7\baselineskip}
% mdparagraph:P138:begin
\subsection*{Unnan: supported forms, wishes and commands}\BookHeading{nav-ch10-learner-0032}{subsubsection}{Unnan: supported forms, wishes and commands}{0}
% mdparagraph:P138:end

% teaching: NEW-U10-GRAM-u10-unnan-indicative
% mdparagraph:P139:begin
\textbf{Unnan}, “grant, allow”, is a preterite-present verb: the present indicative has ann in the first and third singular and unnon throughout the plural. Its dental preterite has ūð-. Remember its dative person and genitive thing.
% mdparagraph:P139:end

% mdparagraph:P140:begin
The\BookPoint{bk-ch10-learner-538269668304} checked sources do not supply a second-singular present indicative form for this table. The dash is an \textbf{evidence gap}, not a form and not a claim that the form could never occur. Do not fill it by analogy.
% mdparagraph:P140:end

\reserve{7\baselineskip}
% mdparagraph:P141:begin
\subsection*{Unnan: indicative}\BookHeading{nav-ch10-learner-0033}{subsubsection}{Unnan: indicative}{0}
% mdparagraph:P141:end

% teaching: NEW-U10-PAR-unnan-IND
% table: NEW-U10-PAR-unnan-IND
% mdtable:T37:begin
% course-table-style: 2026-10
% course-panel:T37-P01:begin
\par\begingroup
\fontsize{10.5}{13.2}\selectfont
\setlength{\LTleft}{0pt}\setlength{\LTright}{\fill}
\noindent
\begin{tabular}{@{}>{\raggedright\arraybackslash}p{39.459mm}>{\raggedright\arraybackslash}p{39.459mm}>{\raggedright\arraybackslash}p{39.459mm}@{}}
\toprule
\textbf{Person and number} & \textbf{Present} & \textbf{Preterite} \\
\midrule
1 sg. & ann & ūðe \\
2 sg. & — & ūðest \\
3 sg. & ann & ūðe \\
1–3 pl. & unnon & ūðon \\
\bottomrule
\end{tabular}
\par\endgroup
% course-panel:T37-P01:end
% mdtable:T37:end
% teaching: NEW-U10-GRAM-u10-unnan-subjunctive-command
% mdparagraph:P142:begin
The present subjunctive is \textbf{unne} in every singular person and \textbf{unnen} in every plural person. The preterite subjunctive is \textbf{ūðe} throughout the singular and \textbf{ūðen} throughout the plural. These forms serve the earlier uses of subjunctive wishes: a present grant or wish uses unne/\allowbreak{}unnen; an unreal present wish can use the preterite subjunctive. Case government remains unchanged.
% mdparagraph:P142:end

\reserve{7\baselineskip}
% mdparagraph:P143:begin
\subsection*{Unnan: subjunctive}\BookHeading{nav-ch10-learner-0034}{subsubsection}{Unnan: subjunctive}{0}
% mdparagraph:P143:end

% teaching: NEW-U10-PAR-unnan-SUBJ
% table: NEW-U10-PAR-unnan-SUBJ
% mdtable:T38:begin
% course-table-style: 2026-10
% course-panel:T38-P01:begin
\par\begingroup
\fontsize{10.5}{13.2}\selectfont
\setlength{\LTleft}{0pt}\setlength{\LTright}{\fill}
\noindent
\begin{tabular}{@{}>{\raggedright\arraybackslash}p{39.459mm}>{\raggedright\arraybackslash}p{39.459mm}>{\raggedright\arraybackslash}p{39.459mm}@{}}
\toprule
\textbf{Person and number} & \textbf{Present} & \textbf{Preterite} \\
\midrule
1 sg. & unne & ūðe \\
2 sg. & unne & ūðe \\
3 sg. & unne & ūðe \\
1–3 pl. & unnen & ūðen \\
\bottomrule
\end{tabular}
\par\endgroup
% course-panel:T38-P01:end
% mdtable:T38:end
% mdparagraph:P144:begin
For\BookPoint{bk-ch10-learner-4160a02bc5fd} a direct addressed command “grant”, use \textbf{unne} to one addressee and \textbf{unnen} to several. These are the second-person subjunctive forms used as command substitutes. Cook labels unne “Imper.”; Fulk describes the subjunctive substitution. Both accounts support the same singular surface. Do not make a formal plural command in -að, or mistake command unne for the missing indicative cell. The preterite plural indicative ūðon and subjunctive ūðen differ in their final vowel.
% mdparagraph:P144:end

\reserve{7\baselineskip}
% mdparagraph:P145:begin
\section*{8. Further constructions for the exercises}\BookHeading{nav-ch10-learner-0035}{subsection}{8. Further constructions for the exercises}{1}
% mdparagraph:P145:end

% teaching: NEW-U10-GRAM-u10-perception-object-predicate
% mdparagraph:P146:begin
\textbf{Gesēon},\BookPoint{bk-ch10-learner-d733552c7139} “see”, can take an accusative object followed by an adjective describing the object's perceived condition. The predicate adjective is strong and agrees with the object in case, gender and number. It does not describe the observer.
% mdparagraph:P146:end

\reserve{8\baselineskip}
% example: NEW-U10-EXAMPLE13
% mdparagraph:P147:begin
Ic geseah ðone cuman untrumne.
% mdparagraph:P147:end

% mdparagraph:P148:begin
I saw the stranger sick.
% mdparagraph:P148:end

% mdparagraph:P149:begin
Untrumne is masculine accusative singular, agreeing with ðone cuman. It is a description of the object, distinct from the earlier accusative actor followed by an infinitive naming an action.
% mdparagraph:P149:end

% teaching: NEW-U10-GRAM-u10-hieran-perception-infinitive
% teaching: NEW-U10-GRAM-u10-promoted-plain-infinitives
% teaching: NEW-U10-VOC-HIERAN
% teaching: NEW-U10-VOC-BLAWAN
% teaching: NEW-U10-VOC-CIEPAN
% teaching: NEW-U10-VOC-BYCGAN
% mdparagraph:P150:begin
\textbf{Hīeran},\BookPoint{bk-ch10-learner-c4f790ec0bcb} “hear”, now takes the same defined action construction as gesēon: an \textbf{accusative perceived actor + plain infinitive}. The infinitive names what the actor is heard doing and has no preceding tō. \textbf{Blāwan}, “blow”, is available as that plain infinitive. If the actor is a noun denoting an instrument or implement, it is still accusative in this construction.
% mdparagraph:P150:end

% mdparagraph:P151:begin
The plain infinitives \textbf{cīepan}, “sell”, and \textbf{bycgan}, “buy”, are also available after volitional \textbf{willan}. The person who wills is the understood subject of the infinitive. Keep each infinitive's ordinary accusative thing sold or bought. This use concerns willingness or intention; it does not establish a mechanical future tense.
% mdparagraph:P151:end

% teaching: NEW-U10-GRAM-u10-present-future
% mdparagraph:P152:begin
An\BookPoint{bk-ch10-learner-5a53bc6a06e5} Old English present can refer to a future event when the context makes the time clear. English “will come”, “will see” or “will separate” can therefore be translated with the ordinary present. \textbf{Willan + infinitive} is appropriate when intention or willingness is actually meant, including the farmer's and workman's stated plans. English “will” alone does not require willan; likewise “shall come” can predict an event without commanding it.
% mdparagraph:P152:end

% mdparagraph:P153:begin
Keep the earlier negative volitional use: ne with willan expresses unwillingness or refusal. Use that meaning when it suits the speaker's refusal to act. Where a prompt admits either prediction or intention, the construction must express the chosen meaning consistently.
% mdparagraph:P153:end

% teaching: NEW-U10-GRAM-u10-to-destination
% teaching: NEW-U10-GRAM-u10-beforan-location
% teaching: NEW-U10-VOC-TO
% teaching: NEW-U10-VOC-BEFORAN
% mdparagraph:P154:begin
\textbf{Tō + dative} names a destination with verbs of coming, going or flying. It is distinct from a purpose/\allowbreak{}function expression and from tō before an inflected infinitive. \textbf{Beforan + dative} means “before, in front of” a person in the standing construction. The exact reading below uses the separately supplied ætforan; do not silently replace its spelling.
% mdparagraph:P154:end

% teaching: NEW-U10-GRAM-u10-on-side-idiom
% teaching: NEW-U10-GRAM-u10-on-name-idiom
% teaching: NEW-U10-VOC-ON
% mdparagraph:P155:begin
In\BookPoint{bk-ch10-learner-902af315de33} the bounded expression \textbf{on … swīðran hand}, “on … right hand/\allowbreak{}side”, the noun \textbf{hand is accusative}. The expression occurs even with static standan, “stand”. Learn this case with the side-expression; motion alone does not explain it. The bare noun hand distinguishes the selected accusative from dative hande or handa. Weak swīðran by itself cannot decide the case, since -an occurs in several weak singular cells. The parallel left-side expression in the reading is supplied locally.
% mdparagraph:P155:end

% mdparagraph:P156:begin
Ordinary \textbf{on + dative} still expresses location, as with a seat, prison or inn. The side-expression does not replace that general pattern.
% mdparagraph:P156:end

% mdparagraph:P157:begin
\textbf{On … naman},\BookPoint{bk-ch10-learner-98bdf06c8402} “in … name”, instead has dative \textbf{naman}. A possessor between on and naman is genitive. As with the recipient phrase in §1, distinguish the case governed by the preposition from the genitive inside its complement.
% mdparagraph:P157:end

% teaching: NEW-U10-GRAM-u10-totwaeman-separation
% teaching: NEW-U10-VOC-TOTWAEMAN
% mdparagraph:P158:begin
The earlier \textbf{tōtwǣman}, “separate”, can now also separate beings: the person, animal or group separated is accusative, and \textbf{fram + dative} names what it is separated from. Use its familiar indicative forms. The source reading has a different verb, tōsceadan, whose exact finite form is supplied there.
% mdparagraph:P158:end

% teaching: NEW-U10-GRAM-u10-hwonne-direct-question
% teaching: NEW-U10-VOC-HWONNE
% mdparagraph:P159:begin
Invariant \textbf{hwonne}, “when?”, introduces a direct question about time. The finite verb expresses the tense and person; use the ordinary interrogative order with hwonne before it and the subject following. English auxiliary “do/\allowbreak{}did” is not an additional Old English verb in these questions. The past form of “see” itself expresses “did see”. This question word does not yet introduce a general system of temporal subordinate clauses.
% mdparagraph:P159:end

% teaching: NEW-U10-GRAM-u10-coordinated-predicate-ellipsis
% mdparagraph:P160:begin
In a coordination with a new expressed subject, a repeated predicate may be left understood from the preceding clause. The following subject still has its own nominative case. This accounts for the two coming subjects in Cook III.15.9 and for the angels' understood coming in the source reading. Supply the repeated action in understanding; this bounded pattern does not permit arbitrary deletion of verbs.
% mdparagraph:P160:end

% teaching: NEW-U10-VOC-MICEL-COMP
% mdparagraph:P161:begin
The familiar comparative \textbf{māra}, neuter \textbf{māre}, can express greater value or importance, “more than”. Use that evaluation sense in the comparison of life and raiment, with the already learned comparative connection ðonne. No new comparative formation is needed.
% mdparagraph:P161:end

\reserve{7\baselineskip}
% mdparagraph:P162:begin
\section*{9. New practice}\BookHeading{nav-ch10-learner-0036}{subsection}{9. New practice}{1}
% mdparagraph:P162:end

\reserve{7\baselineskip}
% mdparagraph:P163:begin
\subsection*{Forms}\BookHeading{nav-ch10-learner-0037}{subsubsection}{Forms}{0}
% mdparagraph:P163:end

\reserve{5\baselineskip}
\par\begin{minipage}{\linewidth}\raggedright
% mdparagraph:P164:begin
\textbf{F01.}
% mdparagraph:P164:end

% item: NEW-U10-F01
% mdparagraph:P165:begin
Give the complete nominative, genitive, dative and accusative singular and plural of fæder. Identify the selected endingless singular cells.
% mdparagraph:P165:end

\end{minipage}\par\addvspace{4pt}

\par\begin{minipage}{\linewidth}\raggedright
% mdparagraph:P166:begin
\textbf{F02.}
% mdparagraph:P166:end

% item: NEW-U10-F02
% mdparagraph:P167:begin
Give the complete nominative, genitive, dative and accusative singular and plural of gāt. Name every cell represented by gǣt in the selected paradigm.
% mdparagraph:P167:end

\end{minipage}\par\addvspace{4pt}

\par\begin{minipage}{\linewidth}\raggedright
% mdparagraph:P168:begin
\textbf{F03.}
% mdparagraph:P168:end

% item: NEW-U10-F03
% mdparagraph:P169:begin
Give all persons of the present and preterite indicative of forwiernan. Keep its present and preterite stems distinct.
% mdparagraph:P169:end

\end{minipage}\par\addvspace{4pt}

\par\begin{minipage}{\linewidth}\raggedright
% mdparagraph:P170:begin
\textbf{F04.}
% mdparagraph:P170:end

% item: NEW-U10-F04
% mdparagraph:P171:begin
Give the supported present and preterite indicative of unnan, marking the unsupported second-singular present cell with a dash. Give its singular and plural subjunctive forms in both tenses and identify the forms used for a grant or wish.
% mdparagraph:P171:end

\end{minipage}\par\addvspace{4pt}

\par\begin{minipage}{\linewidth}\raggedright
% mdparagraph:P172:begin
\textbf{F05.}
% mdparagraph:P172:end

% item: NEW-U10-F05
% mdparagraph:P173:begin
Give the second and third singular and third plural present indicative of fēdan, its first singular and third plural preterite, and both commands.
% mdparagraph:P173:end

\end{minipage}\par\addvspace{4pt}

\par\begin{minipage}{\linewidth}\raggedright
% mdparagraph:P174:begin
\textbf{F06.}
% mdparagraph:P174:end

% item: NEW-U10-F06
% mdparagraph:P175:begin
Give the distinct masculine/\allowbreak{}neuter singular instrumental forms of se and strong gōd. Then give the corresponding dative forms. State how the feminine singular and plural express the instrumental relation.
% mdparagraph:P175:end

\end{minipage}\par\addvspace{4pt}

\reserve{7\baselineskip}
% mdparagraph:P176:begin
\subsection*{Analysis and contrasts}\BookHeading{nav-ch10-learner-0038}{subsubsection}{Analysis and contrasts}{0}
% mdparagraph:P176:end

\reserve{5\baselineskip}
\par\begin{minipage}{\linewidth}\raggedright
% mdparagraph:P177:begin
\textbf{A01.}
% mdparagraph:P177:end

% item: NEW-U10-A01
% mdparagraph:P178:begin
Translate both sentences. Identify the genitive in each and explain what governs it.
% mdparagraph:P178:end

% mdparagraph:P179:begin
\textbf{a.} Ðæs cuman sunu cōm.
% mdparagraph:P179:end

% mdparagraph:P180:begin
\textbf{b.} Ic geswāc ðæs weorces.
% mdparagraph:P180:end

\end{minipage}\par\addvspace{4pt}

\par\begin{minipage}{\linewidth}\raggedright
% mdparagraph:P181:begin
\textbf{A02.}
% mdparagraph:P181:end

% item: NEW-U10-A02
% mdparagraph:P182:begin
Translate the sentence. Identify the recipient, the thing given, and the possessor of the recipient; give their cases.
% mdparagraph:P182:end

% mdparagraph:P183:begin
Se munuc sealde ðæs cuman suna męte.
% mdparagraph:P183:end

\end{minipage}\par\addvspace{4pt}

\par\begin{minipage}{\linewidth}\raggedright
% mdparagraph:P184:begin
\textbf{A03.}
% mdparagraph:P184:end

% item: NEW-U10-A03
% mdparagraph:P185:begin
Translate both sentences and explain the difference in the phrases beginning ān.
% mdparagraph:P185:end

% mdparagraph:P186:begin
\textbf{a.} Ān munuc cōm.
% mdparagraph:P186:end

% mdparagraph:P187:begin
\textbf{b.} Ān ðāra munuca cōm.
% mdparagraph:P187:end

\end{minipage}\par\addvspace{4pt}

\par\begin{minipage}{\linewidth}\raggedright
% mdparagraph:P188:begin
\textbf{A04.}
% mdparagraph:P188:end

% item: NEW-U10-A04
% mdparagraph:P189:begin
Translate both sentences. Explain the different case of scēap and scēapa.
% mdparagraph:P189:end

% mdparagraph:P190:begin
\textbf{a.} Ic hæbbe twā scēap.
% mdparagraph:P190:end

% mdparagraph:P191:begin
\textbf{b.} Ic hæbbe fela scēapa.
% mdparagraph:P191:end

\end{minipage}\par\addvspace{4pt}

\par\begin{minipage}{\linewidth}\raggedright
% mdparagraph:P192:begin
\textbf{A05.}
% mdparagraph:P192:end

% item: NEW-U10-A05
% mdparagraph:P193:begin
Translate both sentences. Explain why the person has the same case but the food has different cases.
% mdparagraph:P193:end

% mdparagraph:P194:begin
\textbf{a.} Hē sealde ðǣm cuman męte.
% mdparagraph:P194:end

% mdparagraph:P195:begin
\textbf{b.} Hē forwiernde ðǣm cuman mętes.
% mdparagraph:P195:end

\end{minipage}\par\addvspace{4pt}

\par\begin{minipage}{\linewidth}\raggedright
% mdparagraph:P196:begin
\textbf{A06.}
% mdparagraph:P196:end

% item: NEW-U10-A06
% mdparagraph:P197:begin
Translate the sentence. Identify who is useful and to whom; explain the forms of gāt, scēaphierde and behēfu.
% mdparagraph:P197:end

% mdparagraph:P198:begin
Sēo gāt is ðǣm scēaphierde behēfu.
% mdparagraph:P198:end

\end{minipage}\par\addvspace{4pt}

\par\begin{minipage}{\linewidth}\raggedright
% mdparagraph:P199:begin
\textbf{A07.}
% mdparagraph:P199:end

% item: NEW-U10-A07
% mdparagraph:P200:begin
Translate the sentence and identify the experiencer. Explain why mē is not a nominative subject.
% mdparagraph:P200:end

% mdparagraph:P201:begin
Mē ðyncð ðis gōd.
% mdparagraph:P201:end

\end{minipage}\par\addvspace{4pt}

\par\begin{minipage}{\linewidth}\raggedright
% mdparagraph:P202:begin
\textbf{A08.}
% mdparagraph:P202:end

% item: NEW-U10-A08
% mdparagraph:P203:begin
Translate both sentences. Compare the case of the pronoun and the agreement of the finite verb.
% mdparagraph:P203:end

% mdparagraph:P204:begin
\textbf{a.} Hine hyngrede.
% mdparagraph:P204:end

% mdparagraph:P205:begin
\textbf{b.} Hē wæs hungrig.
% mdparagraph:P205:end

\end{minipage}\par\addvspace{4pt}

\par\begin{minipage}{\linewidth}\raggedright
% mdparagraph:P206:begin
\textbf{A09.}
% mdparagraph:P206:end

% item: NEW-U10-A09
% mdparagraph:P207:begin
Translate both sentences. Distinguish the relations expressed by mid and state the case it governs in each.
% mdparagraph:P207:end

% mdparagraph:P208:begin
\textbf{a.} Se munuc cōm mid ðǣm cuman.
% mdparagraph:P208:end

% mdparagraph:P209:begin
\textbf{b.} Se munuc fēdde ðone cuman mid gōdum męte.
% mdparagraph:P209:end

\end{minipage}\par\addvspace{4pt}

\par\begin{minipage}{\linewidth}\raggedright
% mdparagraph:P210:begin
\textbf{A10.}
% mdparagraph:P210:end

% item: NEW-U10-A10
% mdparagraph:P211:begin
Translate both sentences. Explain the form of gōde or ðȳ and state whether the noun ending alone distinguishes dative from instrumental.
% mdparagraph:P211:end

% mdparagraph:P212:begin
\textbf{a.} Se munuc fēdde ðone cuman gōde męte.
% mdparagraph:P212:end

% mdparagraph:P213:begin
\textbf{b.} Se munuc fēdde ðone cuman ðȳ męte.
% mdparagraph:P213:end

\end{minipage}\par\addvspace{4pt}

\par\begin{minipage}{\linewidth}\raggedright
% mdparagraph:P214:begin
\textbf{A11.}
% mdparagraph:P214:end

% item: NEW-U10-A11
% mdparagraph:P215:begin
Translate the sentence. State the case, number and gender of cuman and untrumne, and explain the agreement.
% mdparagraph:P215:end

% mdparagraph:P216:begin
Ic geseah ðone cuman untrumne.
% mdparagraph:P216:end

\end{minipage}\par\addvspace{4pt}

\par\begin{minipage}{\linewidth}\raggedright
% mdparagraph:P217:begin
\textbf{A12.}
% mdparagraph:P217:end

% item: NEW-U10-A12
% mdparagraph:P218:begin
Translate both sentences. Explain the number and function of each adjective and the difference in reference.
% mdparagraph:P218:end

% mdparagraph:P219:begin
\textbf{a.} Se munuc scrȳt ðā nacodan.
% mdparagraph:P219:end

% mdparagraph:P220:begin
\textbf{b.} Se munuc scrȳt ðone nacodan cuman.
% mdparagraph:P220:end

\end{minipage}\par\addvspace{4pt}

\reserve{7\baselineskip}
% mdparagraph:P221:begin
\subsection*{Controlled composition}\BookHeading{nav-ch10-learner-0039}{subsubsection}{Controlled composition}{0}
% mdparagraph:P221:end

\reserve{5\baselineskip}
\par\begin{minipage}{\linewidth}\raggedright
% mdparagraph:P222:begin
\textbf{T01.}
% mdparagraph:P222:end

% item: NEW-U10-T01
% mdparagraph:P223:begin
Translate: The monks deny the stranger food. Use forwiernan.
% mdparagraph:P223:end

\end{minipage}\par\addvspace{4pt}

\par\begin{minipage}{\linewidth}\raggedright
% mdparagraph:P224:begin
\textbf{T02.}
% mdparagraph:P224:end

% item: NEW-U10-T02
% mdparagraph:P225:begin
Translate: I grant the stranger food. Use unnan.
% mdparagraph:P225:end

\end{minipage}\par\addvspace{4pt}

\par\begin{minipage}{\linewidth}\raggedright
% mdparagraph:P226:begin
\textbf{T03.}
% mdparagraph:P226:end

% item: NEW-U10-T03
% mdparagraph:P227:begin
Translate both sentences: Some of the monks came with the stranger. Which of the monks came? Use sum and hwilc with a genitive whole.
% mdparagraph:P227:end

\end{minipage}\par\addvspace{4pt}

\par\begin{minipage}{\linewidth}\raggedright
% mdparagraph:P228:begin
\textbf{T04.}
% mdparagraph:P228:end

% item: NEW-U10-T04
% mdparagraph:P229:begin
Translate: The shepherd has many sheep. Use fela.
% mdparagraph:P229:end

\end{minipage}\par\addvspace{4pt}

\par\begin{minipage}{\linewidth}\raggedright
% mdparagraph:P230:begin
\textbf{T05.}
% mdparagraph:P230:end

% item: NEW-U10-T05
% mdparagraph:P231:begin
Translate: The monks feed the strangers with good food. Express the means without a preposition.
% mdparagraph:P231:end

\end{minipage}\par\addvspace{4pt}

\par\begin{minipage}{\linewidth}\raggedright
% mdparagraph:P232:begin
\textbf{T06.}
% mdparagraph:P232:end

% item: NEW-U10-T06
% mdparagraph:P233:begin
Translate: With that food I fed the sick. Use the distinct instrumental of se and an adjective naming the sick people.
% mdparagraph:P233:end

\end{minipage}\par\addvspace{4pt}

\par\begin{minipage}{\linewidth}\raggedright
% mdparagraph:P234:begin
\textbf{T07.}
% mdparagraph:P234:end

% item: NEW-U10-T07
% mdparagraph:P235:begin
Translate: They were hungry. Use the impersonal verb hyngran.
% mdparagraph:P235:end

\end{minipage}\par\addvspace{4pt}

\par\begin{minipage}{\linewidth}\raggedright
% mdparagraph:P236:begin
\textbf{T08.}
% mdparagraph:P236:end

% item: NEW-U10-T08
% mdparagraph:P237:begin
Translate: That seemed good to us. Use the chapter’s ðyncan construction.
% mdparagraph:P237:end

\end{minipage}\par\addvspace{4pt}

\reserve{7\baselineskip}
% mdparagraph:P238:begin
\section*{10. Cook III.13–15}\BookHeading{nav-ch10-learner-0040}{subsection}{10. Cook III.13–15}{1}
% mdparagraph:P238:end

% mdparagraph:P239:begin
Translate into Old English. Preserve the stated number of an addressee and the number of a named object. Where a prompt leaves a second-person number unstated, choose a consistent reading. “Om.” instructs omission of the specified English word; it is not vocabulary to translate. Direct speech stays direct speech. The English “will” and “shall” should be interpreted by the meaning, as explained in §8.
% mdparagraph:P239:end

% mdparagraph:P240:begin
The complete original sets follow. Their grammatical references are preserved. The two local helps supply expressions beyond the independent composition scope; use them for their stated items.
% mdparagraph:P240:end

\reserve{7\baselineskip}
% mdparagraph:P241:begin
\subsection*{III.13}\BookHeading{nav-ch10-learner-0041}{subsubsection}{III.13}{0}
% mdparagraph:P241:end

\reserve{5\baselineskip}
\par\begin{minipage}{\linewidth}\raggedright
% mdparagraph:P242:begin
\textbf{III.13.01.}
% mdparagraph:P242:end

% item: EX-III-13-01
% mdparagraph:P243:begin
Christ is the Son of God [God’s Son].
% mdparagraph:P243:end

\end{minipage}\par\addvspace{4pt}

\par\begin{minipage}{\linewidth}\raggedright
% mdparagraph:P244:begin
\textbf{III.13.02.}
% mdparagraph:P244:end

% item: EX-III-13-02
% mdparagraph:P245:begin
Christ is our Lord.
% mdparagraph:P245:end

\end{minipage}\par\addvspace{4pt}

\par\begin{minipage}{\linewidth}\raggedright
% mdparagraph:P246:begin
\textbf{III.13.03.}
% mdparagraph:P246:end

% item: EX-III-13-03
% mdparagraph:P247:begin
When saw I you (sing.) hungry?
% mdparagraph:P247:end

\end{minipage}\par\addvspace{4pt}

\par\begin{minipage}{\linewidth}\raggedright
% mdparagraph:P248:begin
\textbf{III.13.04.}
% mdparagraph:P248:end

% item: EX-III-13-04
% mdparagraph:P249:begin
Gold is precious.
% mdparagraph:P249:end

\end{minipage}\par\addvspace{4pt}

\par\begin{minipage}{\linewidth}\raggedright
% mdparagraph:P250:begin
\textbf{III.13.05.}
% mdparagraph:P250:end

% item: EX-III-13-05
% mdparagraph:P251:begin
The Day of Judgment will come.
% mdparagraph:P251:end

\end{minipage}\par\addvspace{4pt}

\par\begin{minipage}{\linewidth}\raggedright
% mdparagraph:P252:begin
\textbf{III.13.06.}
% mdparagraph:P252:end

% item: EX-III-13-06
% mdparagraph:P253:begin
The farmer says to the workman, I will sell this house.
% mdparagraph:P253:end

\end{minipage}\par\addvspace{4pt}

\par\begin{minipage}{\linewidth}\raggedright
% mdparagraph:P254:begin
\textbf{III.13.07.}
% mdparagraph:P254:end

% item: EX-III-13-07
% mdparagraph:P255:begin
The workman answers the farmer, I will buy your house.
% mdparagraph:P255:end

\end{minipage}\par\addvspace{4pt}

\par\begin{minipage}{\linewidth}\raggedright
% mdparagraph:P256:begin
\textbf{III.13.08.}
% mdparagraph:P256:end

% item: EX-III-13-08
% mdparagraph:P257:begin
Life is more than raiment.
% mdparagraph:P257:end

\end{minipage}\par\addvspace{4pt}

\par\begin{minipage}{\linewidth}\raggedright
% mdparagraph:P258:begin
\textbf{III.13.09.}
% mdparagraph:P258:end

% item: EX-III-13-09
% mdparagraph:P259:begin
The stranger is in prison.
% mdparagraph:P259:end

\end{minipage}\par\addvspace{4pt}

\par\begin{minipage}{\linewidth}\raggedright
% mdparagraph:P260:begin
\textbf{III.13.10.}
% mdparagraph:P260:end

% item: EX-III-13-10
% mdparagraph:P261:begin
The merchant is in the inn.
% mdparagraph:P261:end

\end{minipage}\par\addvspace{4pt}

\par\begin{minipage}{\linewidth}\raggedright
% mdparagraph:P262:begin
\textbf{III.13.11.}
% mdparagraph:P262:end

% item: EX-III-13-11
% mdparagraph:P263:begin
The father clothes the son.
% mdparagraph:P263:end

\end{minipage}\par\addvspace{4pt}

\par\begin{minipage}{\linewidth}\raggedright
% mdparagraph:P264:begin
\textbf{III.13.12.}
% mdparagraph:P264:end

% item: EX-III-13-12
% mdparagraph:P265:begin
This world is great.
% mdparagraph:P265:end

\end{minipage}\par\addvspace{4pt}

\par\begin{minipage}{\linewidth}\raggedright
% mdparagraph:P266:begin
\textbf{III.13.13.}
% mdparagraph:P266:end

% item: EX-III-13-13
% mdparagraph:P267:begin
The shepherd feeds the sheep (sing.).
% mdparagraph:P267:end

\end{minipage}\par\addvspace{4pt}

\par\begin{minipage}{\linewidth}\raggedright
% mdparagraph:P268:begin
\textbf{III.13.14.}
% mdparagraph:P268:end

% item: EX-III-13-14
% mdparagraph:P269:begin
The merchants are rich.
% mdparagraph:P269:end

\end{minipage}\par\addvspace{4pt}

\par\begin{minipage}{\linewidth}\raggedright
% mdparagraph:P270:begin
\textbf{III.13.15.}
% mdparagraph:P270:end

% item: EX-III-13-15
% mdparagraph:P271:begin
The soldier drinks wine.
% mdparagraph:P271:end

\end{minipage}\par\addvspace{4pt}

\par\begin{minipage}{\linewidth}\raggedright
% mdparagraph:P272:begin
\textbf{III.13.16.}
% mdparagraph:P272:end

% item: EX-III-13-16
% mdparagraph:P273:begin
The monks drank water.
% mdparagraph:P273:end

\end{minipage}\par\addvspace{4pt}

\reserve{7\baselineskip}
% mdparagraph:P274:begin
\subsection*{III.14}\BookHeading{nav-ch10-learner-0042}{subsubsection}{III.14}{0}
% mdparagraph:P274:end

\reserve{5\baselineskip}
\par\begin{minipage}{\linewidth}\raggedright
% mdparagraph:P275:begin
\textbf{III.14.01.}
% mdparagraph:P275:end

% item: EX-III-14-01
% mdparagraph:P276:begin
The king sat on the throne.
% mdparagraph:P276:end

\end{minipage}\par\addvspace{4pt}

\par\begin{minipage}{\linewidth}\raggedright
% mdparagraph:P277:begin
\textbf{III.14.02.}
% mdparagraph:P277:end

% item: EX-III-14-02
% mdparagraph:P278:begin
Stand on my right hand.
% mdparagraph:P278:end

\end{minipage}\par\addvspace{4pt}

\par\begin{minipage}{\linewidth}\raggedright
% mdparagraph:P279:begin
\textbf{III.14.03.}
% mdparagraph:P279:end

% item: EX-III-14-03
% mdparagraph:P280:begin
All nations shall come to the judgment.
% mdparagraph:P280:end

\end{minipage}\par\addvspace{4pt}

\par\begin{minipage}{\linewidth}\raggedright
% mdparagraph:P281:begin
\textbf{III.14.04.}
% mdparagraph:P281:end

% item: EX-III-14-04
% mdparagraph:P282:begin
The devil is (om. the) father of the unrighteous.
% mdparagraph:P282:end

\end{minipage}\par\addvspace{4pt}

\par\begin{minipage}{\linewidth}\raggedright
% mdparagraph:P283:begin
\textbf{III.14.05.}
% mdparagraph:P283:end

% item: EX-III-14-05
% mdparagraph:P284:begin
All men are sinful.
% mdparagraph:P284:end

\end{minipage}\par\addvspace{4pt}

\par\begin{minipage}{\linewidth}\raggedright
% mdparagraph:P285:begin
\textbf{III.14.06.}
% mdparagraph:P285:end

% item: EX-III-14-06
% mdparagraph:P286:begin
The priest is sick.
% mdparagraph:P286:end

\end{minipage}\par\addvspace{4pt}

\par\begin{minipage}{\linewidth}\raggedright
% mdparagraph:P287:begin
\textbf{III.14.07.}
% mdparagraph:P287:end

% item: EX-III-14-07
% mdparagraph:P288:begin
Grant (sing.) me garments.
% mdparagraph:P288:end

\end{minipage}\par\addvspace{4pt}

\par\begin{minipage}{\linewidth}\raggedright
% mdparagraph:P289:begin
\textbf{III.14.08.}
% mdparagraph:P289:end

% item: EX-III-14-08
% mdparagraph:P290:begin
The horse is useful to the king.
% mdparagraph:P290:end

\end{minipage}\par\addvspace{4pt}

\par\begin{minipage}{\linewidth}\raggedright
% mdparagraph:P291:begin
\textbf{III.14.09.}
% mdparagraph:P291:end

% item: EX-III-14-09
% mdparagraph:P292:begin
The loss of all your things was strange.
% mdparagraph:P292:end

\end{minipage}\par\addvspace{4pt}

\par\begin{minipage}{\linewidth}\raggedright
% mdparagraph:P293:begin
\textbf{III.14.10.}
% mdparagraph:P293:end

% item: EX-III-14-10
% mdparagraph:P294:begin
He departed from her.
% mdparagraph:P294:end

\end{minipage}\par\addvspace{4pt}

\par\begin{minipage}{\linewidth}\raggedright
% mdparagraph:P295:begin
\textbf{III.14.11.}
% mdparagraph:P295:end

% item: EX-III-14-11
% mdparagraph:P296:begin
Give (sing.) me to drink.
% mdparagraph:P296:end

\end{minipage}\par\addvspace{4pt}

% teaching: NEW-U10-COOK-HELPS
% teaching: NEW-U10-GRAM-u10-sellan-drink-infinitive
% support: U10-EXHELP-DRINK
% mdparagraph:P297:begin
\textbf{Local help.} For give someone to drink, use sęllan with the person in the dative and plain drincan, without tō. Apply the already known singular command and pronoun forms. This entire construction is supplied for this item.
% mdparagraph:P297:end

\reserve{5\baselineskip}
\par\begin{minipage}{\linewidth}\raggedright
% mdparagraph:P298:begin
\textbf{III.14.12.}
% mdparagraph:P298:end

% item: EX-III-14-12
% mdparagraph:P299:begin
Visit (plur.) the sick.
% mdparagraph:P299:end

\end{minipage}\par\addvspace{4pt}

\par\begin{minipage}{\linewidth}\raggedright
% mdparagraph:P300:begin
\textbf{III.14.13.}
% mdparagraph:P300:end

% item: EX-III-14-13
% mdparagraph:P301:begin
When did you (sing.) see the goat?
% mdparagraph:P301:end

\end{minipage}\par\addvspace{4pt}

\par\begin{minipage}{\linewidth}\raggedright
% mdparagraph:P302:begin
\textbf{III.14.14.}
% mdparagraph:P302:end

% item: EX-III-14-14
% mdparagraph:P303:begin
When was he in the inn?
% mdparagraph:P303:end

\end{minipage}\par\addvspace{4pt}

\par\begin{minipage}{\linewidth}\raggedright
% mdparagraph:P304:begin
\textbf{III.14.15.}
% mdparagraph:P304:end

% item: EX-III-14-15
% mdparagraph:P305:begin
This is eternal life.
% mdparagraph:P305:end

\end{minipage}\par\addvspace{4pt}

\par\begin{minipage}{\linewidth}\raggedright
% mdparagraph:P306:begin
\textbf{III.14.16.}
% mdparagraph:P306:end

% item: EX-III-14-16
% mdparagraph:P307:begin
The prisons were empty.
% mdparagraph:P307:end

\end{minipage}\par\addvspace{4pt}

\par\begin{minipage}{\linewidth}\raggedright
% mdparagraph:P308:begin
\textbf{III.14.17.}
% mdparagraph:P308:end

% item: EX-III-14-17
% mdparagraph:P309:begin
The living God is eternal.
% mdparagraph:P309:end

\end{minipage}\par\addvspace{4pt}

\reserve{7\baselineskip}
% mdparagraph:P310:begin
\subsection*{III.15}\BookHeading{nav-ch10-learner-0043}{subsubsection}{III.15}{0}
% mdparagraph:P310:end

\reserve{5\baselineskip}
\par\begin{minipage}{\linewidth}\raggedright
% mdparagraph:P311:begin
\textbf{III.15.01.}
% mdparagraph:P311:end

% item: EX-III-15-01
% mdparagraph:P312:begin
He stood before the king.
% mdparagraph:P312:end

\end{minipage}\par\addvspace{4pt}

\par\begin{minipage}{\linewidth}\raggedright
% mdparagraph:P313:begin
\textbf{III.15.02.}
% mdparagraph:P313:end

% item: EX-III-15-02
% mdparagraph:P314:begin
The goat eats grass.
% mdparagraph:P314:end

\end{minipage}\par\addvspace{4pt}

\par\begin{minipage}{\linewidth}\raggedright
% mdparagraph:P315:begin
\textbf{III.15.03.}
% mdparagraph:P315:end

% item: EX-III-15-03
% mdparagraph:P316:begin
The angel flies to the earth.
% mdparagraph:P316:end

\end{minipage}\par\addvspace{4pt}

\par\begin{minipage}{\linewidth}\raggedright
% mdparagraph:P317:begin
\textbf{III.15.04.}
% mdparagraph:P317:end

% item: EX-III-15-04
% mdparagraph:P318:begin
You (sing.) deny me food.
% mdparagraph:P318:end

\end{minipage}\par\addvspace{4pt}

\par\begin{minipage}{\linewidth}\raggedright
% mdparagraph:P319:begin
\textbf{III.15.05.}
% mdparagraph:P319:end

% item: EX-III-15-05
% mdparagraph:P320:begin
The monk did it in Christ’s name.
% mdparagraph:P320:end

\end{minipage}\par\addvspace{4pt}

\par\begin{minipage}{\linewidth}\raggedright
% mdparagraph:P321:begin
\textbf{III.15.06.}
% mdparagraph:P321:end

% item: EX-III-15-06
% mdparagraph:P322:begin
He will separate the goats from the sheep.
% mdparagraph:P322:end

\end{minipage}\par\addvspace{4pt}

\par\begin{minipage}{\linewidth}\raggedright
% mdparagraph:P323:begin
\textbf{III.15.07.}
% mdparagraph:P323:end

% item: EX-III-15-07
% mdparagraph:P324:begin
He poured her out wine and water.
% mdparagraph:P324:end

\end{minipage}\par\addvspace{4pt}

\par\begin{minipage}{\linewidth}\raggedright
% mdparagraph:P325:begin
\textbf{III.15.08.}
% mdparagraph:P325:end

% item: EX-III-15-08
% mdparagraph:P326:begin
The Lord will come with his angels.
% mdparagraph:P326:end

\end{minipage}\par\addvspace{4pt}

\par\begin{minipage}{\linewidth}\raggedright
% mdparagraph:P327:begin
\textbf{III.15.09.}
% mdparagraph:P327:end

% item: EX-III-15-09
% mdparagraph:P328:begin
The morning cometh, and also the night.
% mdparagraph:P328:end

\end{minipage}\par\addvspace{4pt}

\par\begin{minipage}{\linewidth}\raggedright
% mdparagraph:P329:begin
\textbf{III.15.10.}
% mdparagraph:P329:end

% item: EX-III-15-10
% mdparagraph:P330:begin
Work the works of God.
% mdparagraph:P330:end

\end{minipage}\par\addvspace{4pt}

\par\begin{minipage}{\linewidth}\raggedright
% mdparagraph:P331:begin
\textbf{III.15.11.}
% mdparagraph:P331:end

% item: EX-III-15-11
% mdparagraph:P332:begin
Clothe (sing.) the naked.
% mdparagraph:P332:end

\end{minipage}\par\addvspace{4pt}

\par\begin{minipage}{\linewidth}\raggedright
% mdparagraph:P333:begin
\textbf{III.15.12.}
% mdparagraph:P333:end

% item: EX-III-15-12
% mdparagraph:P334:begin
I will not sit on the king’s throne.
% mdparagraph:P334:end

\end{minipage}\par\addvspace{4pt}

\par\begin{minipage}{\linewidth}\raggedright
% mdparagraph:P335:begin
\textbf{III.15.13.}
% mdparagraph:P335:end

% item: EX-III-15-13
% mdparagraph:P336:begin
The bird sees the light.
% mdparagraph:P336:end

\end{minipage}\par\addvspace{4pt}

\par\begin{minipage}{\linewidth}\raggedright
% mdparagraph:P337:begin
\textbf{III.15.14.}
% mdparagraph:P337:end

% item: EX-III-15-14
% mdparagraph:P338:begin
The men will see the stars.
% mdparagraph:P338:end

\end{minipage}\par\addvspace{4pt}

\par\begin{minipage}{\linewidth}\raggedright
% mdparagraph:P339:begin
\textbf{III.15.15.}
% mdparagraph:P339:end

% item: EX-III-15-15
% mdparagraph:P340:begin
As long as you (plur.) fed the hungry, you did it to the Lord.
% mdparagraph:P340:end

\end{minipage}\par\addvspace{4pt}

% teaching: NEW-U10-GRAM-u10-swa-lange-swa-supplied
% support: U10-EXHELP-AS-LONG-AS
% mdparagraph:P341:begin
\textbf{Local help.} Use the supplied connection swā lange swā for the opening As long as. Cook’s note17 explains the biblical expression as inasmuch as: the deeds done to the hungry are counted as done to the Lord. Both swā words belong to this whole connection; the earlier so/\allowbreak{}thus sense does not teach it. The two you references are plural; use the taught past forms and case frames.
% mdparagraph:P341:end

\reserve{5\baselineskip}
\par\begin{minipage}{\linewidth}\raggedright
% mdparagraph:P342:begin
\textbf{III.15.16.}
% mdparagraph:P342:end

% item: EX-III-15-16
% mdparagraph:P343:begin
When did you see the stranger sick?
% mdparagraph:P343:end

\end{minipage}\par\addvspace{4pt}

\par\begin{minipage}{\linewidth}\raggedright
% mdparagraph:P344:begin
\textbf{III.15.17.}
% mdparagraph:P344:end

% item: EX-III-15-17
% mdparagraph:P345:begin
The hungry man came to the ships.
% mdparagraph:P345:end

\end{minipage}\par\addvspace{4pt}

\par\begin{minipage}{\linewidth}\raggedright
% mdparagraph:P346:begin
\textbf{III.15.18.}
% mdparagraph:P346:end

% item: EX-III-15-18
% mdparagraph:P347:begin
The soldier brought us the horses’ fodder.
% mdparagraph:P347:end

\end{minipage}\par\addvspace{4pt}

\par\begin{minipage}{\linewidth}\raggedright
% mdparagraph:P348:begin
\textbf{III.15.19.}
% mdparagraph:P348:end

% item: EX-III-15-19
% mdparagraph:P349:begin
The farmer has the blacksmith’s hammer.
% mdparagraph:P349:end

\end{minipage}\par\addvspace{4pt}

\par\begin{minipage}{\linewidth}\raggedright
% mdparagraph:P350:begin
\textbf{III.15.20.}
% mdparagraph:P350:end

% item: EX-III-15-20
% mdparagraph:P351:begin
Who has the tailor’s needles?
% mdparagraph:P351:end

\end{minipage}\par\addvspace{4pt}

\par\begin{minipage}{\linewidth}\raggedright
% mdparagraph:P352:begin
\textbf{III.15.21.}
% mdparagraph:P352:end

% item: EX-III-15-21
% mdparagraph:P353:begin
I hear the bellows (169) blow.
% mdparagraph:P353:end

\end{minipage}\par\addvspace{4pt}

\par\begin{minipage}{\linewidth}\raggedright
% mdparagraph:P354:begin
\textbf{III.15.22.}
% mdparagraph:P354:end

% item: EX-III-15-22
% mdparagraph:P355:begin
Do (om.) you (sing.) see the hammers smite?
% mdparagraph:P355:end

% mdparagraph:P356:begin
In III.15.22, \textbf{(om.) applies to English Do}. Use the finite interrogative construction with “see” and its accusative actor plus infinitive; do not add a translation of auxiliary do. In III.15.21, (169) is Cook's retained reference to his perception construction.
% mdparagraph:P356:end

\end{minipage}\par\addvspace{4pt}

\reserve{7\baselineskip}
% mdparagraph:P357:begin
\section*{11. The stranger and the monk}\BookHeading{nav-ch10-learner-0044}{subsection}{11. The stranger and the monk}{1}
% mdparagraph:P357:end

% mdparagraph:P358:begin
New teaching material. This connected account uses the vocabulary and constructions taught above.
% mdparagraph:P358:end

% composition: NEW-U10-READING-01
% mdparagraph:P359:begin
\par\begin{minipage}{\linewidth}\raggedright
\readingtext{Sum cuma cōm tō ðǣm giesthūse. Hine hyngrede, and hē wæs untrum. Ān ðāra munuca geseah ðone cuman. Se munuc sealde him męte and wæter. Se cuma cwæð: “Ðū ne forwiernst mē mętes.” Se munuc cwæð: “Ic fēde ðē gōde męte.” Se cuma ǣt ðone męte and dranc ðæt wæter.}
\end{minipage}\par
% mdparagraph:P359:end

\reserve{5\baselineskip}
\par\begin{minipage}{\linewidth}\raggedright
% mdparagraph:P360:begin
\textbf{R01.}
% mdparagraph:P360:end

% item: NEW-U10-R01
% mdparagraph:P361:begin
In the new account, identify the referents of hine, him, mē and ðē at each occurrence. State their cases and the words or constructions governing them.
% mdparagraph:P361:end

\end{minipage}\par\addvspace{4pt}

\par\begin{minipage}{\linewidth}\raggedright
% mdparagraph:P362:begin
\textbf{R02.}
% mdparagraph:P362:end

% item: NEW-U10-R02
% mdparagraph:P363:begin
In the new account, explain the case and function of ðāra munuca and mętes. Distinguish the functions of the forms cuma and cuman.
% mdparagraph:P363:end

\end{minipage}\par\addvspace{4pt}

\par\begin{minipage}{\linewidth}\raggedright
% mdparagraph:P364:begin
\textbf{R03.}
% mdparagraph:P364:end

% item: NEW-U10-R03
% mdparagraph:P365:begin
In the monk’s final reply, explain the ending of gōde. Rewrite only that means phrase using mid.
% mdparagraph:P365:end

\end{minipage}\par\addvspace{4pt}

\par\begin{minipage}{\linewidth}\raggedright
% mdparagraph:P366:begin
\textbf{R04.}
% mdparagraph:P366:end

% item: NEW-U10-R04
% mdparagraph:P367:begin
Translate the whole new account. Preserve who gives, who is hungry and who eats or drinks.
% mdparagraph:P367:end

\end{minipage}\par\addvspace{4pt}

\reserve{7\baselineskip}
% mdparagraph:P368:begin
\section*{12. The Day of Judgment}\BookHeading{nav-ch10-learner-0045}{subsection}{12. The Day of Judgment}{1}
% mdparagraph:P368:end

% mdparagraph:P369:begin
Exact source passage. The three complete opening sentences form a scene of arrival, enthronement, gathering and separation. The passage ends before the king's speech. Its original spellings and quantities are retained. The supplied help belongs to this reading; expressions marked there are not additional unaided composition vocabulary.
% mdparagraph:P369:end

\par\begin{minipage}{\linewidth}\raggedright
% source: FB-III-JUDGMENT-TABLEAU-01
% mdparagraph:P370:begin
\readingtext{Witodlice Mannes Bearn cymð on his mægenðrymme, and ealle ęnglas samod mid him tō ðǣm miclan dōme. Ðonne sitt hē on ðǣm setle his mægenðrymnesse, and bēoð gegaderode ætforan him ealle ðēoda, and hē tōscǣt hīe on twā, swā-swā scēaphierde tōscǣt scēap fram gātum. Ðonne gelōgað hē ðā scēap on his swīðran hand, and ðā gǣt on his winstran.}
% mdparagraph:P370:end

\end{minipage}\par
\reserve{7\baselineskip}
% mdparagraph:P371:begin
\subsection*{Local help}\BookHeading{nav-ch10-learner-0046}{subsubsection}{Local help}{0}
% mdparagraph:P371:end

% teaching: NEW-U10-SOURCE-HELPS
% teaching: NEW-U10-VOC-BEARN
% teaching: NEW-U10-VOC-MAEGENTHRYMM
% teaching: NEW-U10-VOC-MAEGENTHRYMNES
% teaching: NEW-U10-VOC-WINSTRE
% teaching: NEW-U10-VOC-WITODLICE
% teaching: NEW-U10-VOC-SAMOD
% teaching: NEW-U10-VOC-TOSCEADAN
% teaching: NEW-U10-VOC-GELOGIAN
% teaching: NEW-U10-VOC-AETFORAN
% teaching: NEW-U10-VOC-THONNE
% teaching: NEW-U10-GRAM-u10-genitive-quality-source
% mdparagraph:P372:begin
The nine groups below supply the difficult words and larger connections. In particular, \textbf{mægenðrymm} and \textbf{mægenðrymnes} both mean “glory, majesty”; \textbf{tōsceadan} means “separate, divide”, and \textbf{gelōgian} “place, set”. Only their actual forms in this passage are supplied. The nominal and adjective tables above do not imply complete paradigms for these source-only words.
% mdparagraph:P372:end

% mdparagraph:P373:begin
The genitive of quality characterizes a noun by an associated quality: the throne “of his majesty” is a majestic throne. This reading supplies that abstract expression. It does not require free composition with abstract genitives of quality.
% mdparagraph:P373:end

% support: U10-HELP01
% mdparagraph:P374:begin
\textbf{1. Witodlice Mannes Bearn.} Witodlice is a connective here, meaning “now; and” (Cook gives Latin autem). Mannes Bearn means the Son of Man: mannes is genitive singular and Bearn neuter nominative singular. The subsequent hē/\allowbreak{}him/\allowbreak{}his refer to this male person.
% mdparagraph:P374:end

% support: U10-HELP02
% mdparagraph:P375:begin
\textbf{2. on his mægenðrymme.} In his glory/\allowbreak{}majesty. Mægenðrymme is dative masculine singular. This abstract expression is supplied as a whole.
% mdparagraph:P375:end

% support: U10-HELP03
% mdparagraph:P376:begin
\textbf{3. and ealle ęnglas samod mid him ….} And all the angels together with him. The source leaves the coming verb understood in this coordinated part; the angels are nominative plural, not objects of with. Mid him separately means with him. Tō ðǣm miclan dōme means to the great judgment.
% mdparagraph:P376:end

% support: U10-HELP04
% mdparagraph:P377:begin
\textbf{4. Ðonne.} Then. Both occurrences are independent sentence adverbs sequencing the scene; they are not the previously taught than or subordinate when.
% mdparagraph:P377:end

% support: U10-HELP05
% mdparagraph:P378:begin
\textbf{5. setle his mægenðrymnesse.} Throne/\allowbreak{}seat of his majesty. Mægenðrymnesse is a feminine genitive singular characterizing the throne; Cook note7 refers to the genitive of quality (§153.f).
% mdparagraph:P378:end

% support: U10-HELP06
% mdparagraph:P379:begin
\textbf{6. bēoð gegaderode ætforan him ealle ðēoda.} All nations will be gathered before him. Bēoð is third-person plural, and ðēoda is the following nominative feminine plural subject. Gegaderode is the exact supplied plural passive participle; ætforan him means before him. Ealle is the feminine plural -e variant here; the ordinary selected target ealla remains valid.
% mdparagraph:P379:end

% support: U10-HELP07
% mdparagraph:P380:begin
\textbf{7. tōscǣt hīe on twā.} Separates them into two groups. Tōscǣt is third-person singular present; hīe is its accusative plural object. On twā is the locally supplied divided-in-two expression; twā is neuter accusative plural here.
% mdparagraph:P380:end

% support: U10-HELP08
% mdparagraph:P381:begin
\textbf{8. swā-swā scēaphierde tōscǣt scēap fram gātum.} Just as a shepherd separates sheep from goats. The whole comparative connection is supplied. Scēaphierde is the nominative subject, scēap the accusative plural object, and gātum dative after fram. The noun scēap has identical singular/\allowbreak{}plural nominative/\allowbreak{}accusative shapes; the source explicitly identifies this occurrence as plural.
% mdparagraph:P381:end

% support: U10-HELP09
% mdparagraph:P382:begin
\textbf{9. gelōgað … on his swīðran hand … on his winstran.} He places the sheep at his right hand and the goats at his left. Gelōgað is an exact supplied third-singular present. The two animal phrases are accusative plural objects. The right/\allowbreak{}left expressions share a placement frame; on takes accusative hand in this bounded right/\allowbreak{}left-side expression (the final hand is understood). The following unassigned source line8 uses the same expression with static standað; this is not a universal motion-versus-rest rule. Swīðran/\allowbreak{}winstran are weak feminine singular forms. Their -an alone does not distinguish accusative from dative; judge the complete clause and the bare hand form.
% mdparagraph:P382:end

% mdparagraph:P383:begin
The familiar \textbf{cymð} and \textbf{sitt} are the selected third-singular forms of cuman and sittan. \textbf{Miclan} is the weak form of micel with its reduced stem. The passage has six overt finite predicates. The angels' coming is understood in the coordination; recover it from cymð rather than treating the angels as the object of mid.
% mdparagraph:P383:end

\reserve{7\baselineskip}
% mdparagraph:P384:begin
\subsection*{Cook's original notes}\BookHeading{nav-ch10-learner-0047}{subsubsection}{Cook's original notes}{0}
% mdparagraph:P384:end

% mdparagraph:P385:begin
The numbers and anchors below preserve Cook's apparatus separately from the passage and the local help. His historical and etymological questions are source notes, not additional exercises assigned here.
% mdparagraph:P385:end

% source-note: U10-FB134-N01
% source-note: U10-FB134-N02
% source-note: U10-FB134-N03
% source-note: U10-FB134-N04
% source-note: U10-FB134-N05
% source-note: U10-FB134-N06
% source-note: U10-FB134-N07
% source-note: U10-FB134-N08
% source-note: U10-FB134-N09
% source-note: U10-FB134-N10
% source-note: U10-FB134-N11
% source-note: U10-FB134-N12
% source-note: U10-FB134-N13
% source-note: U10-FB134-N14
% mdtable:T39:begin
\par\begingroup
\fontsize{10.5}{13.2}\selectfont
\setlength{\LTleft}{0pt}\setlength{\LTright}{\fill}
\noindent
\begin{tabular}{@{}>{\raggedright\arraybackslash}p{8mm}>{\raggedright\arraybackslash}p{31mm}>{\raggedright\arraybackslash}p{79.37mm}@{}}
\toprule
\textbf{Note} & \textbf{Anchor} & \textbf{Original text} \\
\midrule
1 & Witodlice & Lat. autem. \\
2 & cymð & See cuman. \\
3 & ęnglas & See ęngel. What is the history of this word before it entered Old English? \\
4 & miclan & See 55. \\
5 & dōme & In what modern compound does this meaning of dōm persist? \\
6 & sitt & See sittan. \\
7 & mægenðrymnesse & See 153. f. \\
8 & ðēoda & Nom. plur. \\
9 & tōscǣt & See tōsceadan. Account for the vowel ǣ. \\
10 & scēaphierde & In compound words, the vowel of the first syllable is apt to be shortened in Mod. Eng., the more general principle being that shortening is apt to occur before an accumulation of consonants. Besides scēaphierde, shepherd, note e.g. wīsdōm, wisdom. \\
11 & scēap & Plural; account for the form. \\
12 & gātum & See 24. \\
13 & swīðran & See swīð. \\
14 & gǣt & See 52. \\
\bottomrule
\end{tabular}
\par\endgroup

% mdtable:T39:end
\reserve{7\baselineskip}
% mdparagraph:P386:begin
\subsection*{Reading tasks}\BookHeading{nav-ch10-learner-0048}{subsubsection}{Reading tasks}{0}
% mdparagraph:P386:end

% mdparagraph:P387:begin
Use all the printed local help in these tasks.
% mdparagraph:P387:end

\reserve{5\baselineskip}
\par\begin{minipage}{\linewidth}\raggedright
% mdparagraph:P388:begin
\textbf{S01.}
% mdparagraph:P388:end

% item: NEW-U10-S01
% mdparagraph:P389:begin
Using the printed local help, identify who comes, who is gathered, and what separation follows. Give the boundaries of the corresponding clauses.
% mdparagraph:P389:end

\end{minipage}\par\addvspace{4pt}

\par\begin{minipage}{\linewidth}\raggedright
% mdparagraph:P390:begin
\textbf{S02.}
% mdparagraph:P390:end

% item: NEW-U10-S02
% mdparagraph:P391:begin
In the separation and placement clauses, identify the cases and roles of hīe, ðā scēap and ðā gǣt. Distinguish the genders of the animal nouns.
% mdparagraph:P391:end

\end{minipage}\par\addvspace{4pt}

\par\begin{minipage}{\linewidth}\raggedright
% mdparagraph:P392:begin
\textbf{S03.}
% mdparagraph:P392:end

% item: NEW-U10-S03
% mdparagraph:P393:begin
Using the printed help, explain the two occurrences of ðonne and the relation introduced by swā-swā. Distinguish a sequencing adverb from a clause-introducing expression.
% mdparagraph:P393:end

\end{minipage}\par\addvspace{4pt}

\par\begin{minipage}{\linewidth}\raggedright
% mdparagraph:P394:begin
\textbf{S04.}
% mdparagraph:P394:end

% item: NEW-U10-S04
% mdparagraph:P395:begin
Translate the whole selected Judgment passage, using the printed source notes and local help.
% mdparagraph:P395:end

\end{minipage}\par\addvspace{4pt}

\clearpage
% mdparagraph:P396:begin
\section*{Source and editorial notes}\BookHeading{nav-ch10-learner-0049}{subsection}{Source and editorial notes}{1}
% mdparagraph:P396:end

\begingroup\fontsize{10}{12.8}\selectfont
\setlength{\parskip}{1.5pt}
% mdparagraph:P397:begin
\itempara{1}{\textbf{Editions, locators and fidelity.} Cook, \emph{Exercises in Old English} (1895), printed pp. 6–7, one-based PDF pp. 14–15, III.13–15: all 16, 17 and 22 items respectively are retained. Line endings are joined, typographic apostrophes are represented consistently, and the space before the question mark in III.15.20 is closed: source “needles ?”, learner “needles?”. These are declared formatting normalizations, not changes to the wording. “Tailor's”, the bracketed “[God's Son]”, the explicit singular sheep and plural visiting command, the (169) reference, and the omission directions retain their source force. Cook's \emph{A First Book in Old English}, second edition (1895), is abbreviated \textbf{FB} below. Its printed page number + 16 gives the one-based PDF page; the Exercises offset is + 8.}
% mdparagraph:P397:end

% mdparagraph:P398:begin
The reading is FB Reader III, “The Day of Judgment”, printed p. 134 /\allowbreak{} PDF p. 150. Cook's attribution is: “From Ælfric’s Homilies, vol. 2, pp. 106–108; being a paraphrase of Matt. 25:31–46.” The assigned 59-word extract is its first three complete sentences. Editorial line counting gives lines 1–7, anchored by the printed marginal 5; only 5 is a printed line number in this extract. The final assigned words occur partway through editorial line 7; the following royal speech is unassigned. Notes 1–14 are reproduced above without turning their questions into assignments. Note cues and printed line wrapping are kept in the separate apparatus; no word of the passage is modernized. Following material on pp. 134–136 corroborates selected expressions and local helps, without becoming additional assigned reading.
% mdparagraph:P398:end

% mdparagraph:P399:begin
\itempara{2}{\textbf{Case relations and C02.} Dependent genitives: FB §§152–154, pp. 89–90; governed genitives with dative persons: §§155–159, pp. 91–92. Dative constructions and prepositions: §§157–167, pp. 92–95. In §153.d, p. 89 /\allowbreak{} PDF p. 105, Cook prints \textbf{in Frēan ęgesan}, “in fear of the Lord”, but says \textbf{“ęgesan is an obj. gen.”} The dependent objective genitive is \textbf{Frēan}: the Lord is the object of the fear, as Cook's own verbal paraphrase makes clear. Correction \textbf{C02} corrects this analysis and preserves the phrase. The entire citation is supplied here; its words introduce no further production vocabulary. The source quality genitive is referred to §153.f by p. 134 note 7.}
% mdparagraph:P399:end

% mdparagraph:P400:begin
\itempara{3}{\textbf{Noun models.} FB §§43–44, pp. 28–30, strong masculine models; §43.8, p. 30, exceptional fæder; §47, pp. 32–33, strong neuters; §53, p. 37, weak nouns; §§51–52, pp. 35–37, strong feminine variants and root nouns. Hand's glossary entry, p. 297, explicitly refers to §51.1.3: this licenses its endingless accusative as well as its -a oblique variant. Woruld is named among the endingless-accusative examples. Full-headword/\allowbreak{}short-stem behavior is checked against each lexical entry; it is not imposed on compounds automatically.}
% mdparagraph:P400:end

% mdparagraph:P401:begin
\itempara{4}{\textbf{Adjective models and variants.} FB §§55–60, pp. 38–41; personal substantival adjectives, §181, p. 100. Cook's §57.3 allows internal-vowel variation; the chapter selects the explicitly printed retained-vowel stems consistently. The genitive plural swīðra applies the consonant-cluster simplification of §35.c, p. 25, with the earlier comparative model; it is a declared derivation rather than a claim that Cook prints that whole swīðra table. The distinct instruments of the new adjectives are reference forms. Demonstrative instruments: §§83–84, p. 50. Bare means and mid: §§174–178, pp. 97–99. Cook notes the especially poetic use of bare instruments and the Anglian tendency to put a distinct instrument after mid; the latter is not the prose target here.}
% mdparagraph:P401:end

% mdparagraph:P402:begin
\itempara{5}{\textbf{Verbs and clause constructions.} Finite endings: FB §§93–95, pp. 54–55; dental changes: §§33–35, pp. 24–25. Weak I: §§113,117, pp. 68–69,72; weak II: §§118–119, pp. 73–74; strong I: §102, p. 59. Unnan: §129, p. 77 /\allowbreak{} PDF p. 93; subjunctive commands: §193.a, p. 103. Compare Wright, \emph{Old English Grammar} (1914), §542, p. 284, and Fulk, \emph{An Introductory Grammar of Old English} (2020), §123, pp. 89–90: the absent assigned indicative second singular and the subjunctive command analysis are explicit evidence limits. Cook's label “Imper.” for unne is retained alongside Fulk's explanation.}
% mdparagraph:P402:end

% mdparagraph:P403:begin
Impersonal constructions: FB §§189–191, p. 102, with Fulk §66, p. 42. Present with future reference: FB §188, p. 101. Perception: §169, p. 96, and predicative constructions §173, p. 97. Cook §172.1, p. 97, warns against treating motion/\allowbreak{}rest as an invariable predictor of prepositional case. The bounded on … hand idiom is checked against p. 134, including its static standað in the unassigned continuation. Infinitives: §§198–200, pp. 105–106. The give-to-drink help is adapted from p. 136, editorial lines 5–6: the source is a negative past statement, while III.14.11 requests a positive singular command. The swā lange swā help draws on p. 135, editorial lines 12–14 and note 17; Cook explicitly explains the inasmuch-as force. Only these bounded expressions are supplied, not an unrestricted purpose or temporal-clause system.
% mdparagraph:P403:end

% mdparagraph:P404:begin
\itempara{6}{\textbf{Nominal and adjective lexical evidence.} The following FB vocabulary locators establish the headwords, genders and selected meanings used above. Joined course headwords omit the glossary's segmentation hyphens while preserving letters and quantities. Their paradigms combine these lexical diagnostics with the complete grammatical models cited in notes 3–4.}
% mdparagraph:P404:end

% mdtable:T40:begin
\par\begingroup
\fontsize{10.5}{13.2}\selectfont
\setlength{\LTleft}{0pt}\setlength{\LTright}{\fill}
\begin{longtable}{@{}>{\raggedright\arraybackslash}p{28.73mm}>{\raggedright\arraybackslash}p{26.12mm}>{\raggedright\arraybackslash}p{26.12mm}>{\raggedright\arraybackslash}p{34.6mm}@{}}
\toprule
\textbf{Headword} & \textbf{Printed page} & \textbf{PDF page} & \textbf{Vocabulary entry} \\
\midrule
\endfirsthead
\toprule
\textbf{Headword} & \textbf{Printed page} & \textbf{PDF page} & \textbf{Vocabulary entry} \\
\midrule
\endhead
\midrule
\multicolumn{4}{r}{\footnotesize Continued on the following page.}\\
\endfoot
\bottomrule
\endlastfoot
Crīst & 281 & 297 & s.v. Crīst \\
dryhten & 283 & 299 & s.v. dryhten \\
dōm & 283 & 299 & s.v. dōm \\
rēaf & 309 & 325 & s.v. rēaf \\
cuma & 282 & 298 & s.v. cuma \\
cweartern & 282 & 298 & s.v. cweartern \\
giesthūs & 296 & 312 & s.v. giesthūs \\
fæder & 285 & 301 & s.v. fæder \\
woruld & 323 & 339 & s.v. woruld \\
scēaphierde & 310 & 326 & s.v. scēaphierde \\
scēap & 310 & 326 & s.v. scēap \\
setl & 312 & 328 & s.v. setl \\
hand & 297 & 313 & s.v. hand \\
ðēod & 318 & 334 & s.v. ðēod \\
dēofol & 282 & 298 & s.v. dēofol \\
wǣd & 320 & 336 & s.v. wǣd \\
gāt & 290 & 306 & s.v. gāt \\
nama & 306 & 322 & s.v. nama \\
ęngel & 285 & 301 & s.v. ęngel \\
hungrig & 300 & 316 & s.v. hungrig \\
dēorwierðe & 283 & 299 & s.v. dēorwierðe \\
unrihtwīs & 320 & 336 & s.v. unrihtwīs \\
synfull & 315 & 331 & s.v. synfull \\
untrum & 320 & 336 & s.v. untrum \\
sellic & 312 & 328 & s.v. sellic \\
ēce & 284 & 300 & s.v. ēce \\
cwic & 282 & 298 & s.v. cwic \\
nacod & 306 & 322 & s.v. nacod \\
swīðra & 315 & 331 & s.v. swīð \\
\end{longtable}
\par\endgroup

% mdtable:T40:end
% mdparagraph:P405:begin
\itempara{7}{\textbf{Verbal and constructional lexical evidence.} FB vocabulary: andswarian p. 276, beforan p. 278, fēdan and fela p. 286, forwiernan p. 288, genēosian p. 293, gewītan p. 296, hwonne p. 300, scęncan and scrȳdan p. 311, tōgelǣdan p. 316, wyrcean p. 324. Tōgelǣdan's thing-plus-recipient frame also occurs in Reader II, p. 130, editorial line 5. The already learned cīepan, bycgan, blāwan, hīeran, dōn and tōtwǣman keep their established forms; their new uses are explicitly defined in §§5–8.}
% mdparagraph:P405:end

% mdparagraph:P406:begin
Bosworth–Toller, checked dictionary facsimiles, corroborates the selected lexical frames: \textbf{fédan}, main p. 273b and Supplement p. 207b; \textbf{scrýdan}, main p. 843a; \textbf{scencan}, main p. 828b; \textbf{and-swarian}, main p. 41a; \textbf{for-wiernan}, Supplement p. 258a; \textbf{ge-neósian}, main p. 421a; \textbf{hyngrian}, Supplement p. 586a; \textbf{þyncan}, main pp. 1084b–1085a; \textbf{unnan}, main p. 1122b. In ge-neósian, the cited examples with hine establish an accusative visited person; this does not deny the genitive frame of related nēosian. Cook's genēosian quantity remains the chapter's selected spelling. The scrȳdan garment complement is corroborating reference evidence only; the production frame here requires the accusative person.
% mdparagraph:P406:end

% mdparagraph:P407:begin
\itempara{8}{\textbf{Supplied reading vocabulary.} FB vocabulary entries: bearn p. 277, mægenðrymm and mægenðrymnes p. 304, winstre p. 322, witodlice p. 323, samod p. 310, tōsceadan p. 316, gelōgian p. 293, ætforan p. 275. Their exact local senses and forms are determined by p. 134 and its notes, not by an unrestricted glossary translation. Witodlice is connective here (Cook's Latin autem); the two occurrences of ðonne mean “then”. The source feminine plural ealle and supplied plural passive gegaderode are preserved; ordinary selected ealla remains valid elsewhere. The local helps also supply the divided-in-two phrase, the swā-swā comparison, and the elliptical final hand. The new account and all fourteen teaching examples are original teaching material, not adapted quotations.}
% mdparagraph:P407:end

\endgroup

% book-source-body:end
\par\label{end-ch10-learner}
\clearpage\endgroup

\BookRecto
\begingroup
\renewcommand{\runningtitle}{Unit 11. Adverbial clauses}
\hyphenpenalty=10000
\exhyphenpenalty=10000
\emergencystretch=3em
\newcommand{\itempara}[2]{\par\addvspace{3pt}\noindent\hangindent=1.8em\hangafter=1\makebox[1.8em][l]{#1.}#2\par}
\newcommand{\reserve}[1]{\par\penalty-100\begingroup\dimen0=\pagegoal\advance\dimen0 by-\pagetotal\ifdim\dimen0<#1\newpage\fi\endgroup}
% book-source-body:begin

\addfontfeatures{Ligatures=TeXOff}
\raggedright
\booktitle{\BookComponent{ch11-learner}{Unit 11. Adverbial clauses}UNIT 11}{Adverbial clauses}

% mdparagraph:P001:begin
An\BookPoint{bk-ch11-learner-d6160974f921} adverbial clause places an action in time or relates it to a cause, purpose, result, concession or condition. First identify that relation; then consider how the speaker presents it. An indicative commonly asserts an event. A subjunctive can present an intention, an anticipated event, a concession or a supposition. It does not always mean that something is doubtful or false.
% mdparagraph:P001:end

% mdparagraph:P002:begin
The fourteen translated examples, the exercises marked F, A, T, R, B and S, and “The fisherman and his son” are new teaching material. Cook IV.16–17 and V.18 retain their original English prompts. The two later readings are exact passages from Cook’s \emph{First Book}, with separate context, help and original notes.
% mdparagraph:P002:end

% mdparagraph:P003:begin
In the tables, 1, 2 and 3 denote person. Repeated spellings still occupy distinct grammatical cells. Nominative, genitive, dative and accusative forms are for use in composition within their stated meanings. Forms marked \textbf{reference} are for recognition. Supplied expressions may be used with their printed help. Preserve vowel lengths and the distinction between e and ę.
% mdparagraph:P003:end

\reserve{7\baselineskip}
% mdparagraph:P004:begin
\section*{1. Time clauses}\BookHeading{nav-ch11-learner-0001}{subsection}{1. Time clauses}{1}
% mdparagraph:P004:end

% teaching: NEW-U11-GRAM-u11-clause-roles-order
% teaching: NEW-U11-GRAM-u11-past-when-then
% teaching: NEW-U11-VOC-THA-CONJUNCTION
% teaching: NEW-U11-VOC-THA-ADVERB
% mdparagraph:P005:begin
For\BookPoint{bk-ch11-learner-a850a4fba7f6} an event at a particular past time, \textbf{ðā}, “when”, introduces a subordinate clause, normally with the indicative. A second \textbf{ðā}, “then”, may introduce the main clause. The first connects clauses; the second is an adverb inside the main clause. The same spelling also belongs to familiar inflected forms of the article or demonstrative: identify its function from the construction.
% mdparagraph:P005:end

\reserve{8\baselineskip}
% example: NEW-U11-EXAMPLE01
% mdparagraph:P006:begin
Ðā se fiscere cōm, ðā sæt se sunu on ðǣm hūse.
% mdparagraph:P006:end

% mdparagraph:P007:begin
When the fisherman came, the son was sitting in the house.
% mdparagraph:P007:end

% mdparagraph:P008:begin
Cōm and sæt are past indicative forms. The first clause tells when the sitting took place. The second ðā resumes the narrative; English can omit “then”.
% mdparagraph:P008:end

% mdparagraph:P009:begin
A dependent clause often puts its finite verb late. A main clause beginning with an adverb often puts its finite verb before its subject, as in ðā sæt se sunu. These are useful tendencies, not instructions to reverse or move every verb mechanically. Case and agreement still determine the participants.
% mdparagraph:P009:end

% teaching: NEW-U11-GRAM-u11-recurring-when-then
% teaching: NEW-U11-VOC-THONNE-CONJUNCTION
% teaching: NEW-U11-VOC-THONNE
% mdparagraph:P010:begin
Use\BookPoint{bk-ch11-learner-336ad2620b98} \textbf{ðonne}, “when”, for recurring or prospective time. It too can be answered by the adverb \textbf{ðonne}, “then”. In the following example the visit and gift recur. Both finite verbs are present indicative. A future setting does not by itself require a subjunctive or a separate future ending.
% mdparagraph:P010:end

\reserve{8\baselineskip}
% example: NEW-U11-EXAMPLE02
% mdparagraph:P011:begin
Ðonne se munuc cymð, ðonne sęlð se mann him męte.
% mdparagraph:P011:end

% mdparagraph:P012:begin
When the monk comes, the man gives him food.
% mdparagraph:P012:end

% teaching: NEW-U11-GRAM-u11-while
% teaching: NEW-U11-VOC-THA-HWILE-THE
% mdparagraph:P013:begin
\textbf{Ðā-hwīle-ðe},\BookPoint{bk-ch11-learner-8e5721384100} “while”, introduces a period during which another action takes place. In a statement of actual simultaneous events, use the indicative.
% mdparagraph:P013:end

\reserve{8\baselineskip}
% example: NEW-U11-EXAMPLE03
% mdparagraph:P014:begin
Ðā-hwīle-ðe se munuc sang, se sunu dranc wæter.
% mdparagraph:P014:end

% mdparagraph:P015:begin
While the monk sang, the son drank water.
% mdparagraph:P015:end

% mdparagraph:P016:begin
The two actions overlap in time. The conjunction alone does not make either action the cause of the other.
% mdparagraph:P016:end

% teaching: NEW-U11-GRAM-u11-mid-thy-temporal
% teaching: NEW-U11-VOC-MID-THY
% mdparagraph:P017:begin
\textbf{Mid-ðȳ},\BookPoint{bk-ch11-learner-82c1c5623ba7} “when; while”, is another temporal conjunction. Learn the whole connection in this use; do not translate it word by word as a phrase of means.
% mdparagraph:P017:end

\reserve{8\baselineskip}
% example: NEW-U11-EXAMPLE04
% mdparagraph:P018:begin
Mid-ðȳ se sunu ǣt, se fæder sæt on ðǣm hūse.
% mdparagraph:P018:end

% mdparagraph:P019:begin
While the son ate, the father sat in the house.
% mdparagraph:P019:end

% teaching: NEW-U11-GRAM-u11-before-anticipated
% teaching: NEW-U11-VOC-AER-CONJUNCTION
% mdparagraph:P020:begin
\textbf{Ǣr},\BookPoint{bk-ch11-learner-8e868f4d3742} “before”, can introduce an event anticipated from the viewpoint of the main clause. After the following instruction, the monk’s coming is still expected, so \textbf{cume} is present subjunctive. The command itself is plural.
% mdparagraph:P020:end

\reserve{8\baselineskip}
% example: NEW-U11-EXAMPLE05
% mdparagraph:P021:begin
Ne farað, ǣr se munuc cume.
% mdparagraph:P021:end

% mdparagraph:P022:begin
Do not go before the monk comes.
% mdparagraph:P022:end

% teaching: NEW-U11-GRAM-u11-until-anticipated
% teaching: NEW-U11-VOC-OTH-THAET
% teaching: NEW-U11-VOC-HER
% mdparagraph:P023:begin
\textbf{Oð-ðæt},\BookPoint{bk-ch11-learner-81a31350b2d5} “until”, gives the endpoint of an action or state. Use the subjunctive for the selected anticipated endpoint. \textbf{Hēr} means “here”.
% mdparagraph:P023:end

\reserve{8\baselineskip}
% example: NEW-U11-EXAMPLE06
% mdparagraph:P024:begin
Wes hēr, oð-ðæt wē cumen.
% mdparagraph:P024:end

% mdparagraph:P025:begin
Be here until we come.
% mdparagraph:P025:end

% mdparagraph:P026:begin
Wes addresses one person. Cumen is present subjunctive plural: the people expected to come are wē. The English present “come” does not decide the Old English mood.
% mdparagraph:P026:end

% mdparagraph:P027:begin
These\BookPoint{bk-ch11-learner-a018097f3125} are statements about viewpoint, not universal rules equating “before” or “until” with one tense or mood. A narrator can use a subjunctive for an event viewed as still ahead at an earlier point, even if it later happened.
% mdparagraph:P027:end

\reserve{7\baselineskip}
% mdparagraph:P028:begin
\section*{2. Cause, purpose and result}\BookHeading{nav-ch11-learner-0002}{subsection}{2. Cause, purpose and result}{1}
% mdparagraph:P028:end

% teaching: NEW-U11-GRAM-u11-causal-for-thon-the
% teaching: NEW-U11-VOC-FOR-THON-THE
% mdparagraph:P029:begin
\textbf{For-ðon-ðe},\BookPoint{bk-ch11-learner-640cebcb16ea} “because”, introduces a stated reason. Use the indicative for the factual cause illustrated here. The earlier causal for-ðām remains available; the exact reading below also has shorter for-ðon.
% mdparagraph:P029:end

\reserve{8\baselineskip}
% example: NEW-U11-EXAMPLE07
% mdparagraph:P030:begin
Se mann healp ðǣm cuman, for-ðon-ðe hē wæs untrum.
% mdparagraph:P030:end

% mdparagraph:P031:begin
The man helped the stranger because the stranger was sick.
% mdparagraph:P031:end

% mdparagraph:P032:begin
Here hē is understood to mean the stranger. Masculine singular grammar by itself could also fit the man; the intended reference must be established from context. The reason clause explains why the help was given.
% mdparagraph:P032:end

% teaching: NEW-U11-GRAM-u11-purpose
% teaching: NEW-U11-VOC-THAET-CONJUNCTION
% mdparagraph:P033:begin
\textbf{Ðæt + subjunctive}\BookPoint{bk-ch11-learner-0621543f1a64} introduces an intended purpose: “so that; in order that”. Unlike the earlier same-subject purpose infinitive, a finite purpose clause has its own expressed or understood subject. That subject may be the person acting in the main clause or a different person.
% mdparagraph:P033:end

\reserve{8\baselineskip}
% example: NEW-U11-EXAMPLE08
% mdparagraph:P034:begin
Ic cume, ðæt ic ðē helpe.
% mdparagraph:P034:end

% mdparagraph:P035:begin
I come so that I may help you.
% mdparagraph:P035:end

% mdparagraph:P036:begin
The first cume is indicative in the statement “I come”; helpe is subjunctive in the purpose. First-person singular present forms may have the same spelling in both moods. Function and context matter as well as endings.
% mdparagraph:P036:end

\reserve{8\baselineskip}
% example: NEW-U11-EXAMPLE09
% mdparagraph:P037:begin
Se fæder cōm, ðæt hē his suna hulpe.
% mdparagraph:P037:end

% mdparagraph:P038:begin
The father came so that he might help his son.
% mdparagraph:P038:end

% mdparagraph:P039:begin
Hulpe is third singular preterite subjunctive of helpan. The father’s purpose is to help his son; the sentence does not by itself assert that he succeeded. His suna is dative, the person helped, with his as genitive possessor.
% mdparagraph:P039:end

% mdparagraph:P040:begin
The present purpose example uses helpe; the past example uses hulpe. For the already learned subjunctives, present singular -e contrasts with plural -en; the corresponding strong preterite uses the preterite plural stem with singular -e and plural -en. Thus helpe/\allowbreak{}helpen and hulpe/\allowbreak{}hulpen differ in tense and number. The examples here establish these particular present and past settings, without a rule mechanically replacing every subordinate tense.
% mdparagraph:P040:end

% teaching: NEW-U11-GRAM-u11-negative-purpose
% teaching: NEW-U11-VOC-THY-LAES-THE
% mdparagraph:P041:begin
\textbf{Ðȳ-lǣs-ðe + subjunctive},\BookPoint{bk-ch11-learner-5c0d35483e2b} “lest; so that … not”, states an outcome to be prevented. Its negative force belongs to the purpose clause. It does not negate the main instruction.
% mdparagraph:P041:end

\reserve{8\baselineskip}
% example: NEW-U11-EXAMPLE10
% mdparagraph:P042:begin
Heald ðæt dēor, ðȳ-lǣs-ðe hit ðone męte nime.
% mdparagraph:P042:end

% mdparagraph:P043:begin
Hold the animal, lest it take the food.
% mdparagraph:P043:end

% mdparagraph:P044:begin
Nime is third singular present subjunctive. The purpose is that the animal should not take the food. No actual taking is asserted.
% mdparagraph:P044:end

% teaching: NEW-U11-GRAM-u11-factual-result
% mdparagraph:P045:begin
An\BookPoint{bk-ch11-learner-678b89a87f59} asserted result commonly has \textbf{swā … ðæt + indicative}, “so … that”. Swā marks the degree; the ðæt clause says what followed. Compare intended help with an actual consequence.
% mdparagraph:P045:end

\reserve{8\baselineskip}
% example: NEW-U11-EXAMPLE11
% mdparagraph:P046:begin
Se sunu wæs swā glæd, ðæt hē sang.
% mdparagraph:P046:end

% mdparagraph:P047:begin
The son was so glad that he sang.
% mdparagraph:P047:end

% mdparagraph:P048:begin
Sang is past indicative: the son did sing. A purpose clause would present singing as an aim. The earlier content clause, as after witan “know”, instead states what is known. The conjunction ðæt alone does not distinguish these three functions.
% mdparagraph:P048:end

% mdparagraph:P049:begin
Cook also gives a result-like optative curse under his subjunctive examples. That special wish does not make every factual result subjunctive. Read the force of the whole sentence.
% mdparagraph:P049:end

\reserve{7\baselineskip}
% mdparagraph:P050:begin
\section*{3. Concession, condition and comparison}\BookHeading{nav-ch11-learner-0003}{subsection}{3. Concession, condition and comparison}{1}
% mdparagraph:P050:end

% teaching: NEW-U11-GRAM-u11-concession
% teaching: NEW-U11-VOC-THEAH
% mdparagraph:P051:begin
\textbf{Ðēah},\BookPoint{bk-ch11-learner-3669ec8fac8a} “although; even though”, introduces a concession: a circumstance that might have been expected to prevent the main action. The selected construction normally has the subjunctive. The speaker can nevertheless accept the circumstance as a fact.
% mdparagraph:P051:end

\reserve{8\baselineskip}
% example: NEW-U11-EXAMPLE12
% mdparagraph:P052:begin
Ðēah se fiscere sīe eald, hē hilpð his suna.
% mdparagraph:P052:end

% mdparagraph:P053:begin
Although the fisherman is old, he helps his son.
% mdparagraph:P053:end

% mdparagraph:P054:begin
Sīe is third singular present subjunctive. The example concedes the fisherman’s old age; it does not deny it. Hilpð is indicative, asserting the help.
% mdparagraph:P054:end

% teaching: NEW-U11-GRAM-u11-contrary-present-condition
% mdparagraph:P055:begin
The\BookPoint{bk-ch11-learner-32cd9ba34c77} earlier open \textbf{gif}, “if”, construction leaves the condition open. An explicitly \textbf{contrary-present} condition imagines something known not to be so now. In that setting both clauses use preterite subjunctive forms. Their past morphology need not give the sentence past-time reference.
% mdparagraph:P055:end

\reserve{12\baselineskip}
% mdparagraph:P056:begin
Open condition: the speaker does not know whether the monk is here.
% mdparagraph:P056:end

\reserve{8\baselineskip}
% example: NEW-U11-EXAMPLE13
% mdparagraph:P057:begin
Gif se munuc is hēr, hē hilpð ðǣm męnn.
% mdparagraph:P057:end

% mdparagraph:P058:begin
If the monk is here, he helps the man.
% mdparagraph:P058:end

\reserve{12\baselineskip}
% mdparagraph:P059:begin
Contrary-present condition: the monk is absent now.
% mdparagraph:P059:end

\reserve{8\baselineskip}
% example: NEW-U11-EXAMPLE14
% mdparagraph:P060:begin
Gif se munuc wǣre hēr, hē hulpe ðǣm męnn.
% mdparagraph:P060:end

% mdparagraph:P061:begin
If the monk were here, he would help the man.
% mdparagraph:P061:end

% mdparagraph:P062:begin
Wǣre\BookPoint{bk-ch11-learner-126b97da7db2} and hulpe are preterite subjunctives referring to the imagined present. Do not translate them as completed past events merely because of their forms. This contrast does not introduce the full system of contrary-to-fact past conditions.
% mdparagraph:P062:end

% teaching: NEW-U11-GRAM-u11-ac-coordination
% teaching: NEW-U11-VOC-AC
% mdparagraph:P063:begin
\textbf{Ac}, “but”, coordinates contrasting clauses or predicates. It imposes no subordinate-clause mood. Keep the persons and numbers of both subjects even when English repeats “shall” or “will”; the already learned present can refer to the future.
% mdparagraph:P063:end

% teaching: NEW-U11-GRAM-u11-swa-comparative-clause
% teaching: NEW-U11-VOC-SWA
% mdparagraph:P064:begin
A\BookPoint{bk-ch11-learner-af27acafa7ea} finite \textbf{swā} clause can mean “as; in the same way as”, comparing actions. Use the indicative when both actions are asserted. This differs from the degree adverb swā in swā … ðæt.
% mdparagraph:P064:end

% teaching: NEW-U11-GRAM-u11-verbal-comparison-ma
% teaching: NEW-U11-VOC-MA
% mdparagraph:P065:begin
\textbf{Mā},\BookPoint{bk-ch11-learner-d6ee9c5b27b0} “more”, can compare the amount of an action. \textbf{Mā … ðonne} means “more … than”. The second subject remains nominative; a repeated predicate can be understood from the first clause. In a question about what people do, do not turn adverbial mā into an adjective describing the people. Distinguish comparative ðonne, “than”, from temporal ðonne, “when”, and interrogative hwonne, “when?”.
% mdparagraph:P065:end

% teaching: NEW-U11-GRAM-u11-temporal-relative-the
% mdparagraph:P066:begin
After\BookPoint{bk-ch11-learner-c34cc0292c16} a time noun, \textbf{ðe} can introduce a temporal relative: “the day when …”. It attaches the event to that stated day. It is neither a direct question with hwonne nor a conditional with gif. The ordinary subject-relative construction from Unit 8 also remains available: a demonstrative meaning “those” may be followed by ðe, “who”, with a plural verb for the people concerned.
% mdparagraph:P066:end

\reserve{7\baselineskip}
% mdparagraph:P067:begin
\section*{4. Nouns: complete models and stem changes}\BookHeading{nav-ch11-learner-0004}{subsection}{4. Nouns: complete models and stem changes}{1}
% mdparagraph:P067:end

% teaching: NEW-U11-LEXICAL-DIAGNOSTICS
% mdparagraph:P068:begin
For each new noun learn its meaning, gender and stem together. The diagnostic forms below are genitive singular, dative singular and nominative plural. Where a noun follows a printed model, choose the cell with the required case and number, replace the model stem by the listed stem, and keep the ending. An endingless cell stays endingless, except where an explicit full-headword rule says otherwise. Apply the operation to \textbf{all four cases in both numbers}.
% mdparagraph:P068:end

% mdparagraph:P069:begin
Singular and plural are available in the stated meanings unless a limit is given. The plurals of \textbf{lēad, īsern, regn, frōfor, dēað, fȳr} and \textbf{seolfor} are reference forms: no new count meaning such as separate “rains” is required. \textbf{Sunne} and \textbf{mōna} denote the singular sun and moon here; their plurals are reference. \textbf{Hǣlend} is the proper title and has only the four assigned singular cells.
% mdparagraph:P069:end

\reserve{7\baselineskip}
% mdparagraph:P070:begin
\subsection*{Ordinary masculine and neuter endings}\BookHeading{nav-ch11-learner-0005}{subsubsection}{Ordinary masculine and neuter endings}{0}
% mdparagraph:P070:end

% teaching: NEW-U11-VOC-WEALL
% mdparagraph:P071:begin
\textbf{weall}, wall (masculine). Stem \textbf{weall-}. Genitive singular \textbf{wealles}; dative singular \textbf{wealle}; nominative plural \textbf{weallas}. This is the complete ordinary consonant masculine model.
% mdparagraph:P071:end

% mdparagraph:P072:begin
Masculine nominative and accusative singular have no ending. The singular genitive has -es and dative -e; the plural has -as, -a, -um, -as.
% mdparagraph:P072:end

% teaching: NEW-U11-NOM-WEALL
% table: NEW-U11-NOM-WEALL
% mdtable:T01:begin
\par\begingroup
\fontsize{10.5}{13.2}\selectfont
\setlength{\LTleft}{0pt}\setlength{\LTright}{\fill}
\noindent
\begin{tabular}{@{}>{\raggedright\arraybackslash}p{39.46mm}>{\raggedright\arraybackslash}p{39.46mm}>{\raggedright\arraybackslash}p{39.46mm}@{}}
\toprule
\textbf{Case} & \textbf{Singular} & \textbf{Plural} \\
\midrule
Nominative & weall & weallas \\
Genitive & wealles & wealla \\
Dative & wealle & weallum \\
Accusative & weall & weallas \\
\bottomrule
\end{tabular}
\par\endgroup

% mdtable:T01:end
% teaching: NEW-U11-VOC-GIMM
% teaching: NEW-U11-NOM-GIMM
% mdparagraph:P073:begin
\textbf{gimm}, gem; precious stone (masculine). Stem \textbf{gimm-}. Genitive singular \textbf{gimmes}; dative singular \textbf{gimme}; nominative plural \textbf{gimmas}. Use every corresponding cell of \textbf{weall}, replacing weall- by the listed stem.
% mdparagraph:P073:end

% teaching: NEW-U11-VOC-REGN
% teaching: NEW-U11-NOM-REGN
% mdparagraph:P074:begin
\textbf{regn}, rain (masculine). Stem \textbf{regn-}. Genitive singular \textbf{regnes}; dative singular \textbf{regne}; nominative plural \textbf{regnas}. Use every corresponding cell of \textbf{weall}, replacing weall- by the listed stem. The plural is reference only. Cook also gives \textbf{rēn}; regn is the form selected for composition.
% mdparagraph:P074:end

% teaching: NEW-U11-VOC-TORR
% teaching: NEW-U11-NOM-TORR
% mdparagraph:P075:begin
\textbf{torr}, tower (masculine). Stem \textbf{torr-}. Genitive singular \textbf{torres}; dative singular \textbf{torre}; nominative plural \textbf{torras}. Use every corresponding cell of \textbf{weall}, replacing weall- by the listed stem.
% mdparagraph:P075:end

% teaching: NEW-U11-VOC-DEATH
% teaching: NEW-U11-NOM-DEATH
% mdparagraph:P076:begin
\textbf{dēað}, death (masculine). Stem \textbf{dēað-}. Genitive singular \textbf{dēaðes}; dative singular \textbf{dēaðe}; nominative plural \textbf{dēaðas}. Use every corresponding cell of \textbf{weall}, replacing weall- by the listed stem. The plural is reference only.
% mdparagraph:P076:end

% teaching: NEW-U11-VOC-IGLAND
% mdparagraph:P077:begin
\textbf{īgland}, island (neuter). Stem \textbf{īgland-}. Genitive singular \textbf{īglandes}; dative singular \textbf{īglande}; nominative plural \textbf{īgland}. This is the complete long-stem neuter model.
% mdparagraph:P077:end

% mdparagraph:P078:begin
The neuter shares singular -es and -e with the masculine, but nominative and accusative plural have no ending. Plural genitive -a and dative -um remain.
% mdparagraph:P078:end

% teaching: NEW-U11-NOM-IGLAND
% table: NEW-U11-NOM-IGLAND
% mdtable:T02:begin
\par\begingroup
\fontsize{10.5}{13.2}\selectfont
\setlength{\LTleft}{0pt}\setlength{\LTright}{\fill}
\noindent
\begin{tabular}{@{}>{\raggedright\arraybackslash}p{39.46mm}>{\raggedright\arraybackslash}p{39.46mm}>{\raggedright\arraybackslash}p{39.46mm}@{}}
\toprule
\textbf{Case} & \textbf{Singular} & \textbf{Plural} \\
\midrule
Nominative & īgland & īgland \\
Genitive & īglandes & īglanda \\
Dative & īglande & īglandum \\
Accusative & īgland & īgland \\
\bottomrule
\end{tabular}
\par\endgroup

% mdtable:T02:end
% teaching: NEW-U11-VOC-MERESWIN
% teaching: NEW-U11-NOM-MERESWIN
% mdparagraph:P079:begin
\textbf{męreswīn}, dolphin (neuter). Stem \textbf{męreswīn-}. Genitive singular \textbf{męreswīnes}; dative singular \textbf{męreswīne}; nominative plural \textbf{męreswīn}. Use every corresponding cell of \textbf{īgland}, replacing īgland- by the listed stem.
% mdparagraph:P079:end

% teaching: NEW-U11-VOC-LEAD
% teaching: NEW-U11-NOM-LEAD
% mdparagraph:P080:begin
\textbf{lēad}, lead (metal) (neuter). Stem \textbf{lēad-}. Genitive singular \textbf{lēades}; dative singular \textbf{lēade}; nominative plural \textbf{lēad}. Use every corresponding cell of \textbf{īgland}, replacing īgland- by the listed stem. The plural is reference only.
% mdparagraph:P080:end

% teaching: NEW-U11-VOC-ISERN
% teaching: NEW-U11-NOM-ISERN
% mdparagraph:P081:begin
\textbf{īsern}, iron (metal) (neuter). Stem \textbf{īsern-}. Genitive singular \textbf{īsernes}; dative singular \textbf{īserne}; nominative plural \textbf{īsern}. Use every corresponding cell of \textbf{īgland}, replacing īgland- by the listed stem. The plural is reference only. Keep its medial e before endings: the final syllable ends in rn.
% mdparagraph:P081:end

% teaching: NEW-U11-VOC-BEARN
% teaching: NEW-U11-NOM-BEARN
% mdparagraph:P082:begin
\textbf{bearn}, son; child (neuter). Stem \textbf{bearn-}. Genitive singular \textbf{bearnes}; dative singular \textbf{bearne}; nominative plural \textbf{bearn}. Use every corresponding cell of \textbf{īgland}, replacing īgland- by the listed stem. This ordinary “child; son” use does not alter the earlier capitalized title Bearn.
% mdparagraph:P082:end

% teaching: NEW-U11-VOC-FYR
% teaching: NEW-U11-NOM-FYR
% mdparagraph:P083:begin
\textbf{fȳr}, fire (neuter). Stem \textbf{fȳr-}. Genitive singular \textbf{fȳres}; dative singular \textbf{fȳre}; nominative plural \textbf{fȳr}. Use every corresponding cell of \textbf{īgland}, replacing īgland- by the listed stem. The plural is reference only.
% mdparagraph:P083:end

\reserve{7\baselineskip}
% mdparagraph:P084:begin
\subsection*{Vowels that change with number}\BookHeading{nav-ch11-learner-0006}{subsubsection}{Vowels that change with number}{0}
% mdparagraph:P084:end

% teaching: NEW-U11-NOMINAL-VOWELS
% mdparagraph:P085:begin
The\BookPoint{bk-ch11-learner-7e1199ab8347} following vowel changes belong to the named nouns. They cannot be applied to every æ or ea. \textbf{Bæð} keeps æ throughout the singular and has a throughout the plural. \textbf{Geat} keeps ea in the singular and has a in the plural. Their nominative and accusative plurals end in -u.
% mdparagraph:P085:end

% teaching: NEW-U11-VOC-BAETH
% mdparagraph:P086:begin
\textbf{bæð}, bath (neuter). Stem \textbf{bæð-}. Genitive singular \textbf{bæðes}; dative singular \textbf{bæðe}; nominative plural \textbf{baðu}.
% mdparagraph:P086:end

% teaching: NEW-U11-NOM-BAETH
% table: NEW-U11-NOM-BAETH
% mdtable:T03:begin
\par\begingroup
\fontsize{10.5}{13.2}\selectfont
\setlength{\LTleft}{0pt}\setlength{\LTright}{\fill}
\noindent
\begin{tabular}{@{}>{\raggedright\arraybackslash}p{39.46mm}>{\raggedright\arraybackslash}p{39.46mm}>{\raggedright\arraybackslash}p{39.46mm}@{}}
\toprule
\textbf{Case} & \textbf{Singular} & \textbf{Plural} \\
\midrule
Nominative & bæð & baðu \\
Genitive & bæðes & baða \\
Dative & bæðe & baðum \\
Accusative & bæð & baðu \\
\bottomrule
\end{tabular}
\par\endgroup

% mdtable:T03:end
% teaching: NEW-U11-VOC-GEAT
% mdparagraph:P087:begin
\textbf{geat}, gate (neuter). Stem \textbf{geat-}. Genitive singular \textbf{geates}; dative singular \textbf{geate}; nominative plural \textbf{gatu}.
% mdparagraph:P087:end

% teaching: NEW-U11-NOM-GEAT
% table: NEW-U11-NOM-GEAT
% mdtable:T04:begin
\par\begingroup
\fontsize{10.5}{13.2}\selectfont
\setlength{\LTleft}{0pt}\setlength{\LTright}{\fill}
\noindent
\begin{tabular}{@{}>{\raggedright\arraybackslash}p{39.46mm}>{\raggedright\arraybackslash}p{39.46mm}>{\raggedright\arraybackslash}p{39.46mm}@{}}
\toprule
\textbf{Case} & \textbf{Singular} & \textbf{Plural} \\
\midrule
Nominative & geat & gatu \\
Genitive & geates & gata \\
Dative & geate & gatum \\
Accusative & geat & gatu \\
\bottomrule
\end{tabular}
\par\endgroup

% mdtable:T04:end
% mdparagraph:P088:begin
The selected dative plural is \textbf{gatum}. Cook also prints \textbf{geatum} elsewhere; it is an accepted variant, not an extra form required from memory.
% mdparagraph:P088:end

% teaching: NEW-U11-VOC-SITHFAET
% mdparagraph:P089:begin
\textbf{sīðfæt},\BookPoint{bk-ch11-learner-20cd80791fd4} journey; passage (masculine). Stem \textbf{sīðfæt-}. Genitive singular \textbf{sīðfætes}; dative singular \textbf{sīðfæte}; nominative plural \textbf{sīðfatas}. The singular uses sīðfæt- and the plural \textbf{sīðfat-}. This noun is masculine, so its nominative and accusative plural take -as.
% mdparagraph:P089:end

% teaching: NEW-U11-NOM-SITHFAET
% table: NEW-U11-NOM-SITHFAET
% mdtable:T05:begin
\par\begingroup
\fontsize{10.5}{13.2}\selectfont
\setlength{\LTleft}{0pt}\setlength{\LTright}{\fill}
\noindent
\begin{tabular}{@{}>{\raggedright\arraybackslash}p{39.46mm}>{\raggedright\arraybackslash}p{39.46mm}>{\raggedright\arraybackslash}p{39.46mm}@{}}
\toprule
\textbf{Case} & \textbf{Singular} & \textbf{Plural} \\
\midrule
Nominative & sīðfæt & sīðfatas \\
Genitive & sīðfætes & sīðfata \\
Dative & sīðfæte & sīðfatum \\
Accusative & sīðfæt & sīðfatas \\
\bottomrule
\end{tabular}
\par\endgroup

% mdtable:T05:end
\reserve{7\baselineskip}
% mdparagraph:P090:begin
\subsection*{Seasons and h-loss}\BookHeading{nav-ch11-learner-0007}{subsubsection}{Seasons and h-loss}{0}
% mdparagraph:P090:end

% teaching: NEW-U11-NOM-SUMER
% teaching: NEW-U11-VOC-SUMER
% mdparagraph:P091:begin
\textbf{sumer}, summer (masculine). Stem \textbf{sumer-}. Genitive singular \textbf{sumeres}; dative singular \textbf{sumere}; nominative plural \textbf{sumeras}. Keep the medial e before each ending.
% mdparagraph:P091:end

% table: NEW-U11-NOM-SUMER
% mdtable:T06:begin
\par\begingroup
\fontsize{10.5}{13.2}\selectfont
\setlength{\LTleft}{0pt}\setlength{\LTright}{\fill}
\noindent
\begin{tabular}{@{}>{\raggedright\arraybackslash}p{39.46mm}>{\raggedright\arraybackslash}p{39.46mm}>{\raggedright\arraybackslash}p{39.46mm}@{}}
\toprule
\textbf{Case} & \textbf{Singular} & \textbf{Plural} \\
\midrule
Nominative & sumer & sumeras \\
Genitive & sumeres & sumera \\
Dative & sumere & sumerum \\
Accusative & sumer & sumeras \\
\bottomrule
\end{tabular}
\par\endgroup

% mdtable:T06:end
% teaching: NEW-U11-VOC-WINTER
% mdparagraph:P092:begin
\textbf{winter}, winter (masculine). Stem \textbf{wintr-}. Genitive singular \textbf{wintres}; dative singular \textbf{wintre}; nominative plural \textbf{wintras}. The full form winter occurs in nominative and accusative singular; before endings use \textbf{wintr-}.
% mdparagraph:P092:end

% teaching: NEW-U11-NOM-WINTER
% table: NEW-U11-NOM-WINTER
% mdtable:T07:begin
\par\begingroup
\fontsize{10.5}{13.2}\selectfont
\setlength{\LTleft}{0pt}\setlength{\LTright}{\fill}
\noindent
\begin{tabular}{@{}>{\raggedright\arraybackslash}p{39.46mm}>{\raggedright\arraybackslash}p{39.46mm}>{\raggedright\arraybackslash}p{39.46mm}@{}}
\toprule
\textbf{Case} & \textbf{Singular} & \textbf{Plural} \\
\midrule
Nominative & winter & wintras \\
Genitive & wintres & wintra \\
Dative & wintre & wintrum \\
Accusative & winter & wintras \\
\bottomrule
\end{tabular}
\par\endgroup

% mdtable:T07:end
% mdparagraph:P093:begin
Cook records occasional \textbf{sumera} and \textbf{wintra} also as genitive singular, dative singular and nominative plural. They are acceptable alternatives in those cells. Learn the ordinary forms displayed here; the variants need not be recalled independently. The source use of winter for “storm” would be a separate meaning.
% mdparagraph:P093:end

% teaching: NEW-U11-NOM-SEOLH
% teaching: NEW-U11-VOC-SEOLH
% mdparagraph:P094:begin
\textbf{seolh},\BookPoint{bk-ch11-learner-8c7893ec15d2} seal (masculine). Stem \textbf{sēol-}. Genitive singular \textbf{sēoles}; dative singular \textbf{sēole}; nominative plural \textbf{sēolas}. The bare nominative and accusative singular keep short eo and h. Before an ending, h disappears and the diphthong lengthens: \textbf{sēol-}.
% mdparagraph:P094:end

% table: NEW-U11-NOM-SEOLH
% mdtable:T08:begin
\par\begingroup
\fontsize{10.5}{13.2}\selectfont
\setlength{\LTleft}{0pt}\setlength{\LTright}{\fill}
\noindent
\begin{tabular}{@{}>{\raggedright\arraybackslash}p{39.46mm}>{\raggedright\arraybackslash}p{39.46mm}>{\raggedright\arraybackslash}p{39.46mm}@{}}
\toprule
\textbf{Case} & \textbf{Singular} & \textbf{Plural} \\
\midrule
Nominative & seolh & sēolas \\
Genitive & sēoles & sēola \\
Dative & sēole & sēolum \\
Accusative & seolh & sēolas \\
\bottomrule
\end{tabular}
\par\endgroup

% mdtable:T08:end
% mdparagraph:P095:begin
Both parts of the change matter: sēoles has no h and has long ēo. Cook also records a bare form seol; the table retains seolh.
% mdparagraph:P095:end

\reserve{7\baselineskip}
% mdparagraph:P096:begin
\subsection*{Feminines and internal vowel loss}\BookHeading{nav-ch11-learner-0008}{subsubsection}{Feminines and internal vowel loss}{0}
% mdparagraph:P096:end

% teaching: NEW-U11-VOC-STRAET
% mdparagraph:P097:begin
\textbf{strǣt}, street (feminine). Stem \textbf{strǣt-}. Genitive singular \textbf{strǣte}; dative singular \textbf{strǣte}; nominative plural \textbf{strǣta}. This is the complete long-stem feminine model.
% mdparagraph:P097:end

% mdparagraph:P098:begin
The nominative singular is endingless. The other singular cases have -e. Nominative, genitive and accusative plural have -a; dative plural has -um.
% mdparagraph:P098:end

% teaching: NEW-U11-NOM-STRAET
% table: NEW-U11-NOM-STRAET
% mdtable:T09:begin
\par\begingroup
\fontsize{10.5}{13.2}\selectfont
\setlength{\LTleft}{0pt}\setlength{\LTright}{\fill}
\noindent
\begin{tabular}{@{}>{\raggedright\arraybackslash}p{39.46mm}>{\raggedright\arraybackslash}p{39.46mm}>{\raggedright\arraybackslash}p{39.46mm}@{}}
\toprule
\textbf{Case} & \textbf{Singular} & \textbf{Plural} \\
\midrule
Nominative & strǣt & strǣta \\
Genitive & strǣte & strǣta \\
Dative & strǣte & strǣtum \\
Accusative & strǣte & strǣta \\
\bottomrule
\end{tabular}
\par\endgroup

% mdtable:T09:end
% teaching: NEW-U11-VOC-GEBYRD
% teaching: NEW-U11-NOM-GEBYRD
% mdparagraph:P099:begin
\textbf{gebyrd}, birth (feminine). Stem \textbf{gebyrd-}. Genitive singular \textbf{gebyrde}; dative singular \textbf{gebyrde}; nominative plural \textbf{gebyrda}. Replace strǣt- by \textbf{gebyrd-} in every cell of strǣt.
% mdparagraph:P099:end

% teaching: NEW-U11-VOC-MILDHEORTNES
% teaching: NEW-U11-NOM-MILDHEORTNES
% mdparagraph:P100:begin
\textbf{mildheortnes}, mercy; mercies (feminine). Stem \textbf{mildheortness-}. Genitive singular \textbf{mildheortnesse}; dative singular \textbf{mildheortnesse}; nominative plural \textbf{mildheortnessa}. Use the endings of \textbf{strǣt} with \textbf{mildheortness-}, but keep the full headword \textbf{mildheortnes} in the endingless nominative singular. The final s is doubled before each ending. Both “mercy” and plural “mercies” are used here.
% mdparagraph:P100:end

% teaching: NEW-U11-VOC-FROFOR
% mdparagraph:P101:begin
\textbf{frōfor}, comfort; consolation (feminine). Stem \textbf{frōfr-}. Genitive singular \textbf{frōfre}; dative singular \textbf{frōfre}; nominative plural \textbf{frōfra}. The nominative singular is the full \textbf{frōfor}; every ending-bearing cell uses \textbf{frōfr-}. Its plural is reference.
% mdparagraph:P101:end

% teaching: NEW-U11-NOM-FROFOR
% table: NEW-U11-NOM-FROFOR
% mdtable:T10:begin
\par\begingroup
\fontsize{10.5}{13.2}\selectfont
\setlength{\LTleft}{0pt}\setlength{\LTright}{\fill}
\noindent
\begin{tabular}{@{}>{\raggedright\arraybackslash}p{39.46mm}>{\raggedright\arraybackslash}p{39.46mm}>{\raggedright\arraybackslash}p{39.46mm}@{}}
\toprule
\textbf{Case} & \textbf{Singular} & \textbf{Plural} \\
\midrule
Nominative & frōfor & frōfra \\
Genitive & frōfre & frōfra \\
Dative & frōfre & frōfrum \\
Accusative & frōfre & frōfra \\
\bottomrule
\end{tabular}
\par\endgroup

% mdtable:T10:end
% teaching: NEW-U11-VOC-CEASTER
% teaching: NEW-U11-NOM-CEASTER
% mdparagraph:P102:begin
\textbf{ceaster}, town; city (feminine). Stem \textbf{ceastr-}. Genitive singular \textbf{ceastre}; dative singular \textbf{ceastre}; nominative plural \textbf{ceastra}. Use \textbf{frōfor} as the full model: replace its ending-bearing stem frōfr- by \textbf{ceastr-}, and its full nominative singular frōfor by \textbf{ceaster}. Both numbers are available.
% mdparagraph:P102:end

% teaching: NEW-U11-VOC-SEOLFOR
% mdparagraph:P103:begin
\textbf{seolfor}, silver (metal) (neuter). Stem \textbf{seolfr-}. Genitive singular \textbf{seolfres}; dative singular \textbf{seolfre}; nominative plural \textbf{seolfor}. Keep the full headword \textbf{seolfor} in all endingless nominative and accusative cells; use \textbf{seolfr-} before endings. The plural is reference.
% mdparagraph:P103:end

% teaching: NEW-U11-NOM-SEOLFOR
% table: NEW-U11-NOM-SEOLFOR
% mdtable:T11:begin
\par\begingroup
\fontsize{10.5}{13.2}\selectfont
\setlength{\LTleft}{0pt}\setlength{\LTright}{\fill}
\noindent
\begin{tabular}{@{}>{\raggedright\arraybackslash}p{39.46mm}>{\raggedright\arraybackslash}p{39.46mm}>{\raggedright\arraybackslash}p{39.46mm}@{}}
\toprule
\textbf{Case} & \textbf{Singular} & \textbf{Plural} \\
\midrule
Nominative & seolfor & seolfor \\
Genitive & seolfres & seolfra \\
Dative & seolfre & seolfrum \\
Accusative & seolfor & seolfor \\
\bottomrule
\end{tabular}
\par\endgroup

% mdtable:T11:end
% mdparagraph:P104:begin
This loss of an unstressed internal vowel is \textbf{syncope}. Compare frōfr-, ceastr-, seolfr- and wintr-. Sumer and īsern retain their vowels. Learn the lexical distinction rather than deleting every internal vowel before an ending.
% mdparagraph:P104:end

\reserve{7\baselineskip}
% mdparagraph:P105:begin
\subsection*{Weak nouns}\BookHeading{nav-ch11-learner-0009}{subsubsection}{Weak nouns}{0}
% mdparagraph:P105:end

% teaching: NEW-U11-VOC-MEREGREOTE
% mdparagraph:P106:begin
\textbf{męregrēote}, pearl (feminine). Stem \textbf{męregrēot-}. Genitive singular \textbf{męregrēotan}; dative singular \textbf{męregrēotan}; nominative plural \textbf{męregrēotan}. The weak feminine has nominative singular -e, the other singular cases -an, plural nominative and accusative -an, genitive -ena and dative -um.
% mdparagraph:P106:end

% teaching: NEW-U11-NOM-MEREGREOTE
% table: NEW-U11-NOM-MEREGREOTE
% mdtable:T12:begin
\par\begingroup
\fontsize{10.5}{13.2}\selectfont
\setlength{\LTleft}{0pt}\setlength{\LTright}{\fill}
\noindent
\begin{tabular}{@{}>{\raggedright\arraybackslash}p{35.87mm}>{\raggedright\arraybackslash}p{39.46mm}>{\raggedright\arraybackslash}p{43.05mm}@{}}
\toprule
\textbf{Case} & \textbf{Singular} & \textbf{Plural} \\
\midrule
Nominative & męregrēote & męregrēotan \\
Genitive & męregrēotan & męregrēotena \\
Dative & męregrēotan & męregrēotum \\
Accusative & męregrēotan & męregrēotan \\
\bottomrule
\end{tabular}
\par\endgroup

% mdtable:T12:end
% teaching: NEW-U11-VOC-SUNNE
% teaching: NEW-U11-NOM-SUNNE
% mdparagraph:P107:begin
\textbf{sunne}, sun (feminine). Stem \textbf{sunn-}. Genitive singular \textbf{sunnan}; dative singular \textbf{sunnan}; nominative plural \textbf{sunnan}. Use every cell of \textbf{męregrēote}, replacing męregrēot- by \textbf{sunn-}. Thus nominative singular has sunn- + -e; the other singular cases have sunn- + -an. The astronomical singular is for composition; the plural is reference.
% mdparagraph:P107:end

% teaching: NEW-U11-VOC-MONA
% mdparagraph:P108:begin
\textbf{mōna}, moon (masculine). Stem \textbf{mōn-}. Genitive singular \textbf{mōnan}; dative singular \textbf{mōnan}; nominative plural \textbf{mōnan}. The weak masculine differs from the weak feminine in nominative singular -a. Its astronomical singular is for composition; the plural is reference.
% mdparagraph:P108:end

% teaching: NEW-U11-NOM-MONA
% table: NEW-U11-NOM-MONA
% mdtable:T13:begin
\par\begingroup
\fontsize{10.5}{13.2}\selectfont
\setlength{\LTleft}{0pt}\setlength{\LTright}{\fill}
\noindent
\begin{tabular}{@{}>{\raggedright\arraybackslash}p{39.46mm}>{\raggedright\arraybackslash}p{39.46mm}>{\raggedright\arraybackslash}p{39.46mm}@{}}
\toprule
\textbf{Case} & \textbf{Singular} & \textbf{Plural} \\
\midrule
Nominative & mōna & mōnan \\
Genitive & mōnan & mōnena \\
Dative & mōnan & mōnum \\
Accusative & mōnan & mōnan \\
\bottomrule
\end{tabular}
\par\endgroup

% mdtable:T13:end
\reserve{7\baselineskip}
% mdparagraph:P109:begin
\subsection*{Titles and apposition}\BookHeading{nav-ch11-learner-0010}{subsubsection}{Titles and apposition}{0}
% mdparagraph:P109:end

% teaching: NEW-U11-VOC-HAELEND
% teaching: NEW-U11-NOM-HAELEND
% mdparagraph:P110:begin
\textbf{Hǣlend}, the Saviour; Jesus (proper title) (masculine). Stem \textbf{Hǣlend-}. Genitive singular \textbf{Hǣlendes}; dative singular \textbf{Hǣlende}. Use the singular cells of \textbf{weall}, replacing weall- by \textbf{Hǣlend-}. The four title forms are also given together below. There is no assigned proper-title plural.
% mdparagraph:P110:end

% mdtable:T14:begin
\par\begingroup
\fontsize{10.5}{13.2}\selectfont
\setlength{\LTleft}{0pt}\setlength{\LTright}{\fill}
\noindent
\begin{tabular}{@{}>{\raggedright\arraybackslash}p{58.54mm}>{\raggedright\arraybackslash}p{62.64mm}@{}}
\toprule
\textbf{Case} & \textbf{Singular title} \\
\midrule
Nominative & Hǣlend \\
Genitive & Hǣlendes \\
Dative & Hǣlende \\
Accusative & Hǣlend \\
\bottomrule
\end{tabular}
\par\endgroup

% mdtable:T14:end
% teaching: NEW-U11-GRAM-NEW-U11-NOM-APPOSITION
% mdparagraph:P111:begin
Nouns\BookPoint{bk-ch11-learner-ef323531af96} in \textbf{apposition} name the same person or thing and agree in case. In \textbf{ūres Dryhtnes Hǣlendes Crīstes}, “of our Lord Jesus Christ”, the three names or titles are genitive singular; ūres agrees with the masculine singular title. Hǣlend supplies the proper-title equivalent of “Jesus” in this group. In a blessing naming God and the Father, the appositional names remain nominative if they name the subject; a dependent “of …” phrase inside them remains genitive.
% mdparagraph:P111:end

% mdparagraph:P112:begin
Keep the earlier distinction between the case of a whole phrase and a genitive inside it. The composition models normally use the practiced preposed genitive. A supported genitive after its head can also be valid Old English, as the exact reading shows; it is not made ungrammatical by that teaching preference.
% mdparagraph:P112:end

\reserve{7\baselineskip}
% mdparagraph:P113:begin
\section*{5. Adjectives: full models, comparatives and personal uses}\BookHeading{nav-ch11-learner-0011}{subsection}{5. Adjectives: full models, comparatives and personal uses}{1}
% mdparagraph:P113:end

% teaching: NEW-U11-ADJECTIVE-DIAGNOSTICS
% mdparagraph:P114:begin
Use a strong adjective for the ordinary indefinite or predicative construction and a weak adjective in the taught definite construction after the article. Keep the grammatical gender of the noun, including neuter wīf, “wife; woman”. The tables give all four ordinary cases in all three genders and both numbers. Distinct masculine and neuter singular instrumentals are \textbf{reference}; “as dative” points to the corresponding dative cell and introduces no new use.
% mdparagraph:P114:end

% mdparagraph:P115:begin
The familiar strong endings include singular masculine accusative -ne, feminine genitive/\allowbreak{}dative -re, plural genitive -ra and dative -um. Weak nominative singular has masculine -a and feminine/\allowbreak{}neuter -e; neuter accusative singular also has -e. The remaining singular cases and plural nominative/\allowbreak{}accusative have -an; plural genitive has the selected -ra and dative -um. Stem changes must be applied along with these endings.
% mdparagraph:P115:end

% teaching: NEW-U11-VOC-HAT
% mdparagraph:P116:begin
\textbf{hāt}, hot. Stem \textbf{hāt-}. The long stem has no -u in strong feminine nominative singular or neuter nominative/\allowbreak{}accusative plural.
% mdparagraph:P116:end

% teaching: NEW-U11-ADJ-HAT-STRONG
\reserve{18\baselineskip}
% mdparagraph:P117:begin
\textbf{hāt — strong}
% mdparagraph:P117:end

% table: NEW-U11-ADJ-HAT-STRONG
% mdtable:T15:begin
% course-table-style: 2026-10
% course-panel:T15-P01:begin
\par\begingroup
\fontsize{10.5}{13.2}\selectfont
\setlength{\LTleft}{0pt}\setlength{\LTright}{\fill}
\noindent
\begin{tabular}{@{}>{\raggedright\arraybackslash}p{22.0mm}>{\raggedright\arraybackslash}p{31.188mm}>{\raggedright\arraybackslash}p{31.188mm}>{\raggedright\arraybackslash}p{31.188mm}@{}}
\toprule
\textbf{Case} & \textbf{Masculine singular} & \textbf{Neuter singular} & \textbf{Feminine singular} \\
\midrule
Nominative & hāt & hāt & hāt \\
Genitive & hātes & hātes & hātre \\
Dative & hātum & hātum & hātre \\
Accusative & hātne & hāt & hāte \\
Instrumental & hāte & hāte & as dative \\
\bottomrule
\end{tabular}
\par\endgroup
% course-panel:T15-P01:end
\par\addvspace{5pt}
% course-panel:T15-P02:begin
\par\begingroup
\fontsize{10.5}{13.2}\selectfont
\setlength{\LTleft}{0pt}\setlength{\LTright}{\fill}
\noindent
\begin{tabular}{@{}>{\raggedright\arraybackslash}p{22.0mm}>{\raggedright\arraybackslash}p{31.188mm}>{\raggedright\arraybackslash}p{31.188mm}>{\raggedright\arraybackslash}p{31.188mm}@{}}
\toprule
\textbf{Case} & \textbf{Masculine plural} & \textbf{Neuter plural} & \textbf{Feminine plural} \\
\midrule
Nominative & hāte & hāt & hāta \\
Genitive & hātra & hātra & hātra \\
Dative & hātum & hātum & hātum \\
Accusative & hāte & hāt & hāta \\
Instrumental & as dative & as dative & as dative \\
\bottomrule
\end{tabular}
\par\endgroup
% course-panel:T15-P02:end
% mdtable:T15:end

% teaching: NEW-U11-ADJ-HAT-WEAK
\reserve{18\baselineskip}
% mdparagraph:P118:begin
\textbf{hāt — weak}
% mdparagraph:P118:end

% table: NEW-U11-ADJ-HAT-WEAK
% mdtable:T16:begin
% course-table-style: 2026-10
% course-panel:T16-P01:begin
\par\begingroup
\fontsize{10.5}{13.2}\selectfont
\setlength{\LTleft}{0pt}\setlength{\LTright}{\fill}
\noindent
\begin{tabular}{@{}>{\raggedright\arraybackslash}p{22.0mm}>{\raggedright\arraybackslash}p{31.188mm}>{\raggedright\arraybackslash}p{31.188mm}>{\raggedright\arraybackslash}p{31.188mm}@{}}
\toprule
\textbf{Case} & \textbf{Masculine singular} & \textbf{Neuter singular} & \textbf{Feminine singular} \\
\midrule
Nominative & hāta & hāte & hāte \\
Genitive & hātan & hātan & hātan \\
Dative & hātan & hātan & hātan \\
Accusative & hātan & hāte & hātan \\
Instrumental & hātan & hātan & as dative \\
\bottomrule
\end{tabular}
\par\endgroup
% course-panel:T16-P01:end
\par\addvspace{5pt}
% course-panel:T16-P02:begin
\par\begingroup
\fontsize{10.5}{13.2}\selectfont
\setlength{\LTleft}{0pt}\setlength{\LTright}{\fill}
\noindent
\begin{tabular}{@{}>{\raggedright\arraybackslash}p{22.0mm}>{\raggedright\arraybackslash}p{31.188mm}>{\raggedright\arraybackslash}p{31.188mm}>{\raggedright\arraybackslash}p{31.188mm}@{}}
\toprule
\textbf{Case} & \textbf{Masculine plural} & \textbf{Neuter plural} & \textbf{Feminine plural} \\
\midrule
Nominative & hātan & hātan & hātan \\
Genitive & hātra & hātra & hātra \\
Dative & hātum & hātum & hātum \\
Accusative & hātan & hātan & hātan \\
Instrumental & as dative & as dative & as dative \\
\bottomrule
\end{tabular}
\par\endgroup
% course-panel:T16-P02:end
% mdtable:T16:end

% teaching: NEW-U11-VOC-RIHTWIS
% teaching: NEW-U11-ADJ-RIHTWIS-STRONG
% teaching: NEW-U11-ADJ-RIHTWIS-WEAK
% mdparagraph:P119:begin
\textbf{rihtwīs}, righteous. Stem \textbf{rihtwīs-}. For every strong and weak cell, replace hāt- in the corresponding \textbf{hāt} table by the listed stem, retaining the ending or lack of ending. Instrumental reference limits remain the same.
% mdparagraph:P119:end

% teaching: NEW-U11-VOC-DEAD
% teaching: NEW-U11-ADJ-DEAD-STRONG
% teaching: NEW-U11-ADJ-DEAD-WEAK
% mdparagraph:P120:begin
\textbf{dēad}, dead. Stem \textbf{dēad-}. For every strong and weak cell, replace hāt- in the corresponding \textbf{hāt} table by the listed stem, retaining the ending or lack of ending. Instrumental reference limits remain the same.
% mdparagraph:P120:end

% teaching: NEW-U11-VOC-BLAEC
% mdparagraph:P121:begin
\textbf{blæc}, black. Stem \textbf{blæc-}. In the strong system, use \textbf{blac-} before a vowel ending, but retain \textbf{blæc-} before a consonantal ending or no ending. The short stem takes -u in strong feminine nominative singular and neuter nominative/\allowbreak{}accusative plural. The weak system uses \textbf{blac-} before vowel endings and \textbf{blæc-} before genitive plural -ra.
% mdparagraph:P121:end

% teaching: NEW-U11-ADJ-BLAEC-STRONG
\reserve{18\baselineskip}
% mdparagraph:P122:begin
\textbf{blæc — strong}
% mdparagraph:P122:end

% table: NEW-U11-ADJ-BLAEC-STRONG
% mdtable:T17:begin
% course-table-style: 2026-10
% course-panel:T17-P01:begin
\par\begingroup
\fontsize{10.5}{13.2}\selectfont
\setlength{\LTleft}{0pt}\setlength{\LTright}{\fill}
\noindent
\begin{tabular}{@{}>{\raggedright\arraybackslash}p{22.0mm}>{\raggedright\arraybackslash}p{31.188mm}>{\raggedright\arraybackslash}p{31.188mm}>{\raggedright\arraybackslash}p{31.188mm}@{}}
\toprule
\textbf{Case} & \textbf{Masculine singular} & \textbf{Neuter singular} & \textbf{Feminine singular} \\
\midrule
Nominative & blæc & blæc & blacu \\
Genitive & blaces & blaces & blæcre \\
Dative & blacum & blacum & blæcre \\
Accusative & blæcne & blæc & blace \\
Instrumental & blace & blace & as dative \\
\bottomrule
\end{tabular}
\par\endgroup
% course-panel:T17-P01:end
\par\addvspace{5pt}
% course-panel:T17-P02:begin
\par\begingroup
\fontsize{10.5}{13.2}\selectfont
\setlength{\LTleft}{0pt}\setlength{\LTright}{\fill}
\noindent
\begin{tabular}{@{}>{\raggedright\arraybackslash}p{22.0mm}>{\raggedright\arraybackslash}p{31.188mm}>{\raggedright\arraybackslash}p{31.188mm}>{\raggedright\arraybackslash}p{31.188mm}@{}}
\toprule
\textbf{Case} & \textbf{Masculine plural} & \textbf{Neuter plural} & \textbf{Feminine plural} \\
\midrule
Nominative & blace & blacu & blaca \\
Genitive & blæcra & blæcra & blæcra \\
Dative & blacum & blacum & blacum \\
Accusative & blace & blacu & blaca \\
Instrumental & as dative & as dative & as dative \\
\bottomrule
\end{tabular}
\par\endgroup
% course-panel:T17-P02:end
% mdtable:T17:end

% teaching: NEW-U11-ADJ-BLAEC-WEAK
\reserve{18\baselineskip}
% mdparagraph:P123:begin
\textbf{blæc — weak}
% mdparagraph:P123:end

% table: NEW-U11-ADJ-BLAEC-WEAK
% mdtable:T18:begin
% course-table-style: 2026-10
% course-panel:T18-P01:begin
\par\begingroup
\fontsize{10.5}{13.2}\selectfont
\setlength{\LTleft}{0pt}\setlength{\LTright}{\fill}
\noindent
\begin{tabular}{@{}>{\raggedright\arraybackslash}p{22.0mm}>{\raggedright\arraybackslash}p{31.188mm}>{\raggedright\arraybackslash}p{31.188mm}>{\raggedright\arraybackslash}p{31.188mm}@{}}
\toprule
\textbf{Case} & \textbf{Masculine singular} & \textbf{Neuter singular} & \textbf{Feminine singular} \\
\midrule
Nominative & blaca & blace & blace \\
Genitive & blacan & blacan & blacan \\
Dative & blacan & blacan & blacan \\
Accusative & blacan & blace & blacan \\
Instrumental & blacan & blacan & as dative \\
\bottomrule
\end{tabular}
\par\endgroup
% course-panel:T18-P01:end
\par\addvspace{5pt}
% course-panel:T18-P02:begin
\par\begingroup
\fontsize{10.5}{13.2}\selectfont
\setlength{\LTleft}{0pt}\setlength{\LTright}{\fill}
\noindent
\begin{tabular}{@{}>{\raggedright\arraybackslash}p{22.0mm}>{\raggedright\arraybackslash}p{31.188mm}>{\raggedright\arraybackslash}p{31.188mm}>{\raggedright\arraybackslash}p{31.188mm}@{}}
\toprule
\textbf{Case} & \textbf{Masculine plural} & \textbf{Neuter plural} & \textbf{Feminine plural} \\
\midrule
Nominative & blacan & blacan & blacan \\
Genitive & blæcra & blæcra & blæcra \\
Dative & blacum & blacum & blacum \\
Accusative & blacan & blacan & blacan \\
Instrumental & as dative & as dative & as dative \\
\bottomrule
\end{tabular}
\par\endgroup
% course-panel:T18-P02:end
% mdtable:T18:end

% mdparagraph:P124:begin
Thus strong masculine accusative blæcne retains æ before -ne; weak masculine nominative blaca has a before -a. The weak genitive plural \textbf{blæcra} uses \textbf{blæc-} before -ra, the ending shared with the strong declension.
% mdparagraph:P124:end

% teaching: NEW-U11-VOC-FAEGER
% mdparagraph:P125:begin
\textbf{fæger}, beautiful; fair. Stem \textbf{fæger-}. The selected disyllabic model retains its internal e. Strong feminine nominative singular and neuter nominative/\allowbreak{}accusative plural have -u.
% mdparagraph:P125:end

% teaching: NEW-U11-ADJ-FAEGER-STRONG
\reserve{18\baselineskip}
% mdparagraph:P126:begin
\textbf{fæger — strong}
% mdparagraph:P126:end

% table: NEW-U11-ADJ-FAEGER-STRONG
% mdtable:T19:begin
% course-table-style: 2026-10
% course-panel:T19-P01:begin
\par\begingroup
\fontsize{10.5}{13.2}\selectfont
\setlength{\LTleft}{0pt}\setlength{\LTright}{\fill}
\noindent
\begin{tabular}{@{}>{\raggedright\arraybackslash}p{22.0mm}>{\raggedright\arraybackslash}p{31.188mm}>{\raggedright\arraybackslash}p{31.188mm}>{\raggedright\arraybackslash}p{31.188mm}@{}}
\toprule
\textbf{Case} & \textbf{Masculine singular} & \textbf{Neuter singular} & \textbf{Feminine singular} \\
\midrule
Nominative & fæger & fæger & fægeru \\
Genitive & fægeres & fægeres & fægerre \\
Dative & fægerum & fægerum & fægerre \\
Accusative & fægerne & fæger & fægere \\
Instrumental & fægere & fægere & as dative \\
\bottomrule
\end{tabular}
\par\endgroup
% course-panel:T19-P01:end
\par\addvspace{5pt}
% course-panel:T19-P02:begin
\par\begingroup
\fontsize{10.5}{13.2}\selectfont
\setlength{\LTleft}{0pt}\setlength{\LTright}{\fill}
\noindent
\begin{tabular}{@{}>{\raggedright\arraybackslash}p{22.0mm}>{\raggedright\arraybackslash}p{31.188mm}>{\raggedright\arraybackslash}p{31.188mm}>{\raggedright\arraybackslash}p{31.188mm}@{}}
\toprule
\textbf{Case} & \textbf{Masculine plural} & \textbf{Neuter plural} & \textbf{Feminine plural} \\
\midrule
Nominative & fægere & fægeru & fægera \\
Genitive & fægerra & fægerra & fægerra \\
Dative & fægerum & fægerum & fægerum \\
Accusative & fægere & fægeru & fægera \\
Instrumental & as dative & as dative & as dative \\
\bottomrule
\end{tabular}
\par\endgroup
% course-panel:T19-P02:end
% mdtable:T19:end

% teaching: NEW-U11-ADJ-FAEGER-WEAK
\reserve{18\baselineskip}
% mdparagraph:P127:begin
\textbf{fæger — weak}
% mdparagraph:P127:end

% table: NEW-U11-ADJ-FAEGER-WEAK
% mdtable:T20:begin
% course-table-style: 2026-10
% course-panel:T20-P01:begin
\par\begingroup
\fontsize{10.5}{13.2}\selectfont
\setlength{\LTleft}{0pt}\setlength{\LTright}{\fill}
\noindent
\begin{tabular}{@{}>{\raggedright\arraybackslash}p{22.0mm}>{\raggedright\arraybackslash}p{31.188mm}>{\raggedright\arraybackslash}p{31.188mm}>{\raggedright\arraybackslash}p{31.188mm}@{}}
\toprule
\textbf{Case} & \textbf{Masculine singular} & \textbf{Neuter singular} & \textbf{Feminine singular} \\
\midrule
Nominative & fægera & fægere & fægere \\
Genitive & fægeran & fægeran & fægeran \\
Dative & fægeran & fægeran & fægeran \\
Accusative & fægeran & fægere & fægeran \\
Instrumental & fægeran & fægeran & as dative \\
\bottomrule
\end{tabular}
\par\endgroup
% course-panel:T20-P01:end
\par\addvspace{5pt}
% course-panel:T20-P02:begin
\par\begingroup
\fontsize{10.5}{13.2}\selectfont
\setlength{\LTleft}{0pt}\setlength{\LTright}{\fill}
\noindent
\begin{tabular}{@{}>{\raggedright\arraybackslash}p{22.0mm}>{\raggedright\arraybackslash}p{31.188mm}>{\raggedright\arraybackslash}p{31.188mm}>{\raggedright\arraybackslash}p{31.188mm}@{}}
\toprule
\textbf{Case} & \textbf{Masculine plural} & \textbf{Neuter plural} & \textbf{Feminine plural} \\
\midrule
Nominative & fægeran & fægeran & fægeran \\
Genitive & fægerra & fægerra & fægerra \\
Dative & fægerum & fægerum & fægerum \\
Accusative & fægeran & fægeran & fægeran \\
Instrumental & as dative & as dative & as dative \\
\bottomrule
\end{tabular}
\par\endgroup
% course-panel:T20-P02:end
% mdtable:T20:end

% teaching: NEW-U11-VOC-GESYNDIG
% teaching: NEW-U11-ADJ-GESYNDIG-STRONG
% teaching: NEW-U11-ADJ-GESYNDIG-WEAK
% mdparagraph:P128:begin
\textbf{gesyndig}, favorable; fair (of winds). Stem \textbf{gesyndig-}. Replace fæger- by the listed stem in \textbf{every cell of both fæger tables}. Keep the internal vowel, including before -u, -re and -ra; the reference instrumentals follow the same operation.
% mdparagraph:P128:end

% teaching: NEW-U11-VOC-THURSTIG
% teaching: NEW-U11-ADJ-THURSTIG-STRONG
% teaching: NEW-U11-ADJ-THURSTIG-WEAK
% mdparagraph:P129:begin
\textbf{ðurstig}, thirsty. Stem \textbf{ðurstig-}. Replace fæger- by the listed stem in \textbf{every cell of both fæger tables}. Keep the internal vowel, including before -u, -re and -ra; the reference instrumentals follow the same operation.
% mdparagraph:P129:end

% mdparagraph:P130:begin
Cook also allows internal-vowel loss in such disyllabic adjectives, and occasional endingless feminine singular and neuter plural forms. The tables choose the retained-vowel series consistently. The exact \textbf{gesyndge}, “favorable”, in the sea passage is a supplied syncopated form beside the selected \textbf{gesyndige}. Do not silently replace one spelling by the other.
% mdparagraph:P130:end

% teaching: NEW-U11-VOC-SMYLTE
% mdparagraph:P131:begin
\textbf{smylte}, calm; smooth. Stem \textbf{smylt-}. The endingless masculine and neuter singular use the full headword \textbf{smylte}. Elsewhere replace final -e by the required ending on \textbf{smylt-}. For example, feminine nominative singular has -u and dative plural -um.
% mdparagraph:P131:end

% teaching: NEW-U11-ADJ-SMYLTE-STRONG
\reserve{18\baselineskip}
% mdparagraph:P132:begin
\textbf{smylte — strong}
% mdparagraph:P132:end

% table: NEW-U11-ADJ-SMYLTE-STRONG
% mdtable:T21:begin
% course-table-style: 2026-10
% course-panel:T21-P01:begin
\par\begingroup
\fontsize{10.5}{13.2}\selectfont
\setlength{\LTleft}{0pt}\setlength{\LTright}{\fill}
\noindent
\begin{tabular}{@{}>{\raggedright\arraybackslash}p{22.0mm}>{\raggedright\arraybackslash}p{31.188mm}>{\raggedright\arraybackslash}p{31.188mm}>{\raggedright\arraybackslash}p{31.188mm}@{}}
\toprule
\textbf{Case} & \textbf{Masculine singular} & \textbf{Neuter singular} & \textbf{Feminine singular} \\
\midrule
Nominative & smylte & smylte & smyltu \\
Genitive & smyltes & smyltes & smyltre \\
Dative & smyltum & smyltum & smyltre \\
Accusative & smyltne & smylte & smylte \\
Instrumental & smylte & smylte & as dative \\
\bottomrule
\end{tabular}
\par\endgroup
% course-panel:T21-P01:end
\par\addvspace{5pt}
% course-panel:T21-P02:begin
\par\begingroup
\fontsize{10.5}{13.2}\selectfont
\setlength{\LTleft}{0pt}\setlength{\LTright}{\fill}
\noindent
\begin{tabular}{@{}>{\raggedright\arraybackslash}p{22.0mm}>{\raggedright\arraybackslash}p{31.188mm}>{\raggedright\arraybackslash}p{31.188mm}>{\raggedright\arraybackslash}p{31.188mm}@{}}
\toprule
\textbf{Case} & \textbf{Masculine plural} & \textbf{Neuter plural} & \textbf{Feminine plural} \\
\midrule
Nominative & smylte & smyltu & smylta \\
Genitive & smyltra & smyltra & smyltra \\
Dative & smyltum & smyltum & smyltum \\
Accusative & smylte & smyltu & smylta \\
Instrumental & as dative & as dative & as dative \\
\bottomrule
\end{tabular}
\par\endgroup
% course-panel:T21-P02:end
% mdtable:T21:end

% teaching: NEW-U11-ADJ-SMYLTE-WEAK
\reserve{18\baselineskip}
% mdparagraph:P133:begin
\textbf{smylte — weak}
% mdparagraph:P133:end

% table: NEW-U11-ADJ-SMYLTE-WEAK
% mdtable:T22:begin
% course-table-style: 2026-10
% course-panel:T22-P01:begin
\par\begingroup
\fontsize{10.5}{13.2}\selectfont
\setlength{\LTleft}{0pt}\setlength{\LTright}{\fill}
\noindent
\begin{tabular}{@{}>{\raggedright\arraybackslash}p{22.0mm}>{\raggedright\arraybackslash}p{31.188mm}>{\raggedright\arraybackslash}p{31.188mm}>{\raggedright\arraybackslash}p{31.188mm}@{}}
\toprule
\textbf{Case} & \textbf{Masculine singular} & \textbf{Neuter singular} & \textbf{Feminine singular} \\
\midrule
Nominative & smylta & smylte & smylte \\
Genitive & smyltan & smyltan & smyltan \\
Dative & smyltan & smyltan & smyltan \\
Accusative & smyltan & smylte & smyltan \\
Instrumental & smyltan & smyltan & as dative \\
\bottomrule
\end{tabular}
\par\endgroup
% course-panel:T22-P01:end
\par\addvspace{5pt}
% course-panel:T22-P02:begin
\par\begingroup
\fontsize{10.5}{13.2}\selectfont
\setlength{\LTleft}{0pt}\setlength{\LTright}{\fill}
\noindent
\begin{tabular}{@{}>{\raggedright\arraybackslash}p{22.0mm}>{\raggedright\arraybackslash}p{31.188mm}>{\raggedright\arraybackslash}p{31.188mm}>{\raggedright\arraybackslash}p{31.188mm}@{}}
\toprule
\textbf{Case} & \textbf{Masculine plural} & \textbf{Neuter plural} & \textbf{Feminine plural} \\
\midrule
Nominative & smyltan & smyltan & smyltan \\
Genitive & smyltra & smyltra & smyltra \\
Dative & smyltum & smyltum & smyltum \\
Accusative & smyltan & smyltan & smyltan \\
Instrumental & as dative & as dative & as dative \\
\bottomrule
\end{tabular}
\par\endgroup
% course-panel:T22-P02:end
% mdtable:T22:end

% teaching: NEW-U11-VOC-STILLE
% mdparagraph:P134:begin
\textbf{stille}, still; quiet. Stem \textbf{still-}. Use full \textbf{stille} in endingless masculine and neuter singular cells, and \textbf{still-} before endings. Simplify ll to l before consonantal -ne, -re and -ra: \textbf{stilne, stilre, stilra}. Retain ll before vowel endings: \textbf{stille, stilles, stillum}. The weak genitive plural is also stilra.
% mdparagraph:P134:end

% teaching: NEW-U11-ADJ-STILLE-STRONG
\reserve{18\baselineskip}
% mdparagraph:P135:begin
\textbf{stille — strong}
% mdparagraph:P135:end

% table: NEW-U11-ADJ-STILLE-STRONG
% mdtable:T23:begin
% course-table-style: 2026-10
% course-panel:T23-P01:begin
\par\begingroup
\fontsize{10.5}{13.2}\selectfont
\setlength{\LTleft}{0pt}\setlength{\LTright}{\fill}
\noindent
\begin{tabular}{@{}>{\raggedright\arraybackslash}p{22.0mm}>{\raggedright\arraybackslash}p{31.188mm}>{\raggedright\arraybackslash}p{31.188mm}>{\raggedright\arraybackslash}p{31.188mm}@{}}
\toprule
\textbf{Case} & \textbf{Masculine singular} & \textbf{Neuter singular} & \textbf{Feminine singular} \\
\midrule
Nominative & stille & stille & stillu \\
Genitive & stilles & stilles & stilre \\
Dative & stillum & stillum & stilre \\
Accusative & stilne & stille & stille \\
Instrumental & stille & stille & as dative \\
\bottomrule
\end{tabular}
\par\endgroup
% course-panel:T23-P01:end
\par\addvspace{5pt}
% course-panel:T23-P02:begin
\par\begingroup
\fontsize{10.5}{13.2}\selectfont
\setlength{\LTleft}{0pt}\setlength{\LTright}{\fill}
\noindent
\begin{tabular}{@{}>{\raggedright\arraybackslash}p{22.0mm}>{\raggedright\arraybackslash}p{31.188mm}>{\raggedright\arraybackslash}p{31.188mm}>{\raggedright\arraybackslash}p{31.188mm}@{}}
\toprule
\textbf{Case} & \textbf{Masculine plural} & \textbf{Neuter plural} & \textbf{Feminine plural} \\
\midrule
Nominative & stille & stillu & stilla \\
Genitive & stilra & stilra & stilra \\
Dative & stillum & stillum & stillum \\
Accusative & stille & stillu & stilla \\
Instrumental & as dative & as dative & as dative \\
\bottomrule
\end{tabular}
\par\endgroup
% course-panel:T23-P02:end
% mdtable:T23:end

% teaching: NEW-U11-ADJ-STILLE-WEAK
\reserve{18\baselineskip}
% mdparagraph:P136:begin
\textbf{stille — weak}
% mdparagraph:P136:end

% table: NEW-U11-ADJ-STILLE-WEAK
% mdtable:T24:begin
% course-table-style: 2026-10
% course-panel:T24-P01:begin
\par\begingroup
\fontsize{10.5}{13.2}\selectfont
\setlength{\LTleft}{0pt}\setlength{\LTright}{\fill}
\noindent
\begin{tabular}{@{}>{\raggedright\arraybackslash}p{22.0mm}>{\raggedright\arraybackslash}p{31.188mm}>{\raggedright\arraybackslash}p{31.188mm}>{\raggedright\arraybackslash}p{31.188mm}@{}}
\toprule
\textbf{Case} & \textbf{Masculine singular} & \textbf{Neuter singular} & \textbf{Feminine singular} \\
\midrule
Nominative & stilla & stille & stille \\
Genitive & stillan & stillan & stillan \\
Dative & stillan & stillan & stillan \\
Accusative & stillan & stille & stillan \\
Instrumental & stillan & stillan & as dative \\
\bottomrule
\end{tabular}
\par\endgroup
% course-panel:T24-P01:end
\par\addvspace{5pt}
% course-panel:T24-P02:begin
\par\begingroup
\fontsize{10.5}{13.2}\selectfont
\setlength{\LTleft}{0pt}\setlength{\LTright}{\fill}
\noindent
\begin{tabular}{@{}>{\raggedright\arraybackslash}p{22.0mm}>{\raggedright\arraybackslash}p{31.188mm}>{\raggedright\arraybackslash}p{31.188mm}>{\raggedright\arraybackslash}p{31.188mm}@{}}
\toprule
\textbf{Case} & \textbf{Masculine plural} & \textbf{Neuter plural} & \textbf{Feminine plural} \\
\midrule
Nominative & stillan & stillan & stillan \\
Genitive & stilra & stilra & stilra \\
Dative & stillum & stillum & stillum \\
Accusative & stillan & stillan & stillan \\
Instrumental & as dative & as dative & as dative \\
\bottomrule
\end{tabular}
\par\endgroup
% course-panel:T24-P02:end
% mdtable:T24:end

% mdparagraph:P137:begin
The literal stille tables show why substituting still- mechanically into smylte without consonant reduction would fail. Quiet or still is an ordinary adjective predicate; it does not call for a passive participle.
% mdparagraph:P137:end

\reserve{7\baselineskip}
% mdparagraph:P138:begin
\subsection*{Comparatives}\BookHeading{nav-ch11-learner-0012}{subsubsection}{Comparatives}{0}
% mdparagraph:P138:end

% teaching: NEW-U11-VOC-DEORWIERTH-COMP
% mdparagraph:P139:begin
\textbf{dēorwierðra},\BookPoint{bk-ch11-learner-d7ca3a0f46ee} more precious. Stem \textbf{dēorwierðr-}. From dēorwierðe, remove final e and form the comparative stem \textbf{dēorwierðr-}. Comparatives remain weak even as predicates.
% mdparagraph:P139:end

% mdparagraph:P140:begin
Attach the weak endings to the comparative stem. Before genitive plural -ra, the resulting -ðrr- group simplifies to -ðr-: \textbf{dēorwierðra}. Thus masculine nominative singular and genitive plural happen to share a spelling. The instrumental row merely refers to the dative forms.
% mdparagraph:P140:end

% teaching: NEW-U11-ADJ-DEORWIERTH-COMP-WEAK
\reserve{18\baselineskip}
% mdparagraph:P141:begin
\textbf{dēorwierðra — weak}
% mdparagraph:P141:end

% table: NEW-U11-ADJ-DEORWIERTH-COMP-WEAK
% mdtable:T25:begin
% course-table-style: 2026-10
% course-panel:T25-P01:begin
\par\begingroup
\fontsize{10.5}{13.2}\selectfont
\setlength{\LTleft}{0pt}\setlength{\LTright}{\fill}
\noindent
\begin{tabular}{@{}>{\raggedright\arraybackslash}p{22.0mm}>{\raggedright\arraybackslash}p{31.188mm}>{\raggedright\arraybackslash}p{31.188mm}>{\raggedright\arraybackslash}p{31.188mm}@{}}
\toprule
\textbf{Case} & \textbf{Masculine singular} & \textbf{Neuter singular} & \textbf{Feminine singular} \\
\midrule
Nominative & dēorwierðra & dēorwierðre & dēorwierðre \\
Genitive & dēorwierðran & dēorwierðran & dēorwierðran \\
Dative & dēorwierðran & dēorwierðran & dēorwierðran \\
Accusative & dēorwierðran & dēorwierðre & dēorwierðran \\
Instrumental & as dative & as dative & as dative \\
\bottomrule
\end{tabular}
\par\endgroup
% course-panel:T25-P01:end
\par\addvspace{5pt}
% course-panel:T25-P02:begin
\par\begingroup
\fontsize{10.5}{13.2}\selectfont
\setlength{\LTleft}{0pt}\setlength{\LTright}{\fill}
\noindent
\begin{tabular}{@{}>{\raggedright\arraybackslash}p{22.0mm}>{\raggedright\arraybackslash}p{31.188mm}>{\raggedright\arraybackslash}p{31.188mm}>{\raggedright\arraybackslash}p{31.188mm}@{}}
\toprule
\textbf{Case} & \textbf{Masculine plural} & \textbf{Neuter plural} & \textbf{Feminine plural} \\
\midrule
Nominative & dēorwierðran & dēorwierðran & dēorwierðran \\
Genitive & dēorwierðra & dēorwierðra & dēorwierðra \\
Dative & dēorwierðrum & dēorwierðrum & dēorwierðrum \\
Accusative & dēorwierðran & dēorwierðran & dēorwierðran \\
Instrumental & as dative & as dative & as dative \\
\bottomrule
\end{tabular}
\par\endgroup
% course-panel:T25-P02:end
% mdtable:T25:end

% teaching: NEW-U11-VOC-FAEGER-COMP
% teaching: NEW-U11-ADJ-FAEGER-COMP-WEAK
% mdparagraph:P142:begin
\textbf{fægerra}, more beautiful; fairer. Stem \textbf{fægerr-}. Use the entire \textbf{dēorwierðra} table, replacing comparative stem dēorwierðr- by \textbf{fægerr-} and keeping the same weak endings. Before genitive plural -ra, simplify the resulting consonant group to \textbf{fægerra}. The nominative singular is \textbf{fægerra, fægerre, fægerre}; the dative plural is \textbf{fægerrum} in all genders.
% mdparagraph:P142:end

% mdparagraph:P143:begin
A comparative predicate agrees with its own subject: the standard after ðonne has its own gender and case. Gold and silver are neuter; righteousness and night are feminine. Do not make the predicate agree with the noun after “than”.
% mdparagraph:P143:end

\reserve{7\baselineskip}
% mdparagraph:P144:begin
\subsection*{Adjectives used without an expressed noun}\BookHeading{nav-ch11-learner-0013}{subsubsection}{Adjectives used without an expressed noun}{0}
% mdparagraph:P144:end

% teaching: NEW-U11-GRAM-u11-bare-strong-personal-adjectives
% teaching: NEW-U11-VOC-CWIC
% mdparagraph:P145:begin
The\BookPoint{bk-ch11-learner-deff8a6444cf} bare \textbf{strong genitive plurals dēadra} and \textbf{cwicra} can mean “of the dead” and “of the living”, referring to people. No article is expressed in this selected construction. Cwic continues the already learned adjective “alive”; English “living” does not require a new participle or progressive tense.
% mdparagraph:P145:end

% teaching: NEW-U11-GRAM-u11-articled-righteous-persons
% mdparagraph:P146:begin
With an article, \textbf{rihtwīs} and \textbf{unrihtwīs} can name righteous and unrighteous people. Use the appropriate \textbf{weak} adjective ending. The whole personal group receives the case required by its clause or preposition; after the taught ofer in the exercise, use accusative. This extends the earlier personal-adjective construction to these meanings.
% mdparagraph:P146:end

% teaching: NEW-U11-GRAM-u11-eall-substantive
% teaching: NEW-U11-VOC-EALL
% mdparagraph:P147:begin
\textbf{Eall} can stand alone: neuter singular \textbf{eall}, “everything”, or masculine plural \textbf{ealle}, “all people”. State which meaning is intended when English simply says “all”. In “over all”, the preposition governs accusative; the two meanings keep their distinct numbers.
% mdparagraph:P147:end

% teaching: NEW-U11-GRAM-u11-vocative-personal-adjective
% mdparagraph:P148:begin
The\BookPoint{bk-ch11-learner-af84902d6468} earlier superlative of lēof may also stand alone in singular direct address: \textbf{lēofosta}, “dearest”, to a male person; \textbf{lēofoste} to a female person. These are weak forms. The speaker’s choice of addressee determines the gender; an English prompt without that information admits either defensible choice.
% mdparagraph:P148:end

\reserve{7\baselineskip}
% mdparagraph:P149:begin
\section*{6. Verbs and their complements}\BookHeading{nav-ch11-learner-0014}{subsection}{6. Verbs and their complements}{1}
% mdparagraph:P149:end

% mdparagraph:P150:begin
The seven finite tables below give the complete \textbf{present and preterite indicative}, all three persons in both numbers. They do not assign new subjunctive or imperative systems for these verbs. The purpose and other subjunctive exercises use verbs whose full subjunctives have already been learned. Learn each new verb’s case construction together with its forms.
% mdparagraph:P150:end

% teaching: NEW-U11-VOC-SCETHTHAN
% teaching: NEW-U11-GRAM-u11-sceththan-indicative
% teaching: NEW-U11-GRAM-u11-sceththan-dative
% mdparagraph:P151:begin
\textbf{Scęððan}, “hurt; injure”, takes the affected person or thing in the \textbf{dative}. English may make it a direct object, but Old English need not. The selected weak class I system has scęðð- before -e and -að, single ð in \textbf{scęðest, scęðeð}, and an uncontracted dental preterite \textbf{scęð-ede}. Preterite second singular adds -st to the first-singular form; the plural replaces final -e by -on.
% mdparagraph:P151:end

% teaching: NEW-U11-PAR-sceththan-IND
% table: NEW-U11-PAR-sceththan-IND
% mdtable:T26:begin
% course-table-style: 2026-10
% course-panel:T26-P01:begin
\par\begingroup
\fontsize{10.5}{13.2}\selectfont
\setlength{\LTleft}{0pt}\setlength{\LTright}{\fill}
\noindent
\begin{tabular}{@{}>{\raggedright\arraybackslash}p{39.459mm}>{\raggedright\arraybackslash}p{39.459mm}>{\raggedright\arraybackslash}p{39.459mm}@{}}
\toprule
\textbf{Person and number} & \textbf{Present} & \textbf{Preterite} \\
\midrule
1 sg. & scęððe & scęðede \\
2 sg. & scęðest & scęðedest \\
3 sg. & scęðeð & scęðede \\
1–3 pl. & scęððað & scęðedon \\
\bottomrule
\end{tabular}
\par\endgroup
% course-panel:T26-P01:end
% mdtable:T26:end
% mdparagraph:P152:begin
Learn this consistent weak system. Contracted present variants and the separate strong preterite scōd are outside the required series.
% mdparagraph:P152:end

% teaching: NEW-U11-VOC-GEWEORTHIAN
% teaching: NEW-U11-GRAM-u11-geweorthian-indicative
% mdparagraph:P153:begin
\textbf{Geweorðian}, “distinguish; adorn”, is weak class II. It takes the \textbf{accusative} thing adorned or distinguished; \textbf{mid + dative} can name its equipment or adornment. A passive can leave the agent unexpressed. Keep the selected spelling geweorðian.
% mdparagraph:P153:end

% mdparagraph:P154:begin
The class II present has -ie, -ast, -að in the singular and -iað in all three plural persons. The preterite has -ode, -odest, -ode, then -odon in each plural person. The vowel before the dental is o.
% mdparagraph:P154:end

% teaching: NEW-U11-PAR-geweorthian-IND
% table: NEW-U11-PAR-geweorthian-IND
% mdtable:T27:begin
% course-table-style: 2026-10
% course-panel:T27-P01:begin
\par\begingroup
\fontsize{10.5}{13.2}\selectfont
\setlength{\LTleft}{0pt}\setlength{\LTright}{\fill}
\noindent
\begin{tabular}{@{}>{\raggedright\arraybackslash}p{39.459mm}>{\raggedright\arraybackslash}p{39.459mm}>{\raggedright\arraybackslash}p{39.459mm}@{}}
\toprule
\textbf{Person and number} & \textbf{Present} & \textbf{Preterite} \\
\midrule
1 sg. & geweorðie & geweorðode \\
2 sg. & geweorðast & geweorðodest \\
3 sg. & geweorðað & geweorðode \\
1–3 pl. & geweorðiað & geweorðodon \\
\bottomrule
\end{tabular}
\par\endgroup
% course-panel:T27-P01:end
% mdtable:T27:end
% teaching: NEW-U11-VOC-SCEAWIAN
% teaching: NEW-U11-GRAM-u11-sceawian-indicative
% teaching: NEW-U11-GRAM-u11-sceawian-willan
% mdparagraph:P155:begin
\textbf{Scēawian}, “watch; observe”, uses the same class II endings on \textbf{scēaw-}. The person or event watched, including a journey, is \textbf{accusative}. With willan expressing a wish, use plain \textbf{scēawian} as the infinitive.
% mdparagraph:P155:end

% teaching: NEW-U11-PAR-sceawian-IND
% table: NEW-U11-PAR-sceawian-IND
% mdtable:T28:begin
% course-table-style: 2026-10
% course-panel:T28-P01:begin
\par\begingroup
\fontsize{10.5}{13.2}\selectfont
\setlength{\LTleft}{0pt}\setlength{\LTright}{\fill}
\noindent
\begin{tabular}{@{}>{\raggedright\arraybackslash}p{39.459mm}>{\raggedright\arraybackslash}p{39.459mm}>{\raggedright\arraybackslash}p{39.459mm}@{}}
\toprule
\textbf{Person and number} & \textbf{Present} & \textbf{Preterite} \\
\midrule
1 sg. & scēawie & scēawode \\
2 sg. & scēawast & scēawodest \\
3 sg. & scēawað & scēawode \\
1–3 pl. & scēawiað & scēawodon \\
\bottomrule
\end{tabular}
\par\endgroup
% course-panel:T28-P01:end
% mdtable:T28:end
% teaching: NEW-U11-PAR-sceawian-INF
% table: NEW-U11-PAR-sceawian-INF
% mdtable:T29:begin
\par\begingroup
\fontsize{10.5}{13.2}\selectfont
\setlength{\LTleft}{0pt}\setlength{\LTright}{\fill}
\noindent
\begin{tabular}{@{}>{\raggedright\arraybackslash}p{71.72mm}>{\raggedright\arraybackslash}p{49.46mm}@{}}
\toprule
\textbf{Form} & \textbf{Infinitive} \\
\midrule
Plain infinitive & scēawian \\
Inflected infinitive & scēawianne \\
\bottomrule
\end{tabular}
\par\endgroup

% mdtable:T29:end
% mdparagraph:P156:begin
The plain infinitive is for composition after willan. \textbf{Scēawianne} is the inflected infinitive, shown for \textbf{reference}; do not substitute it after willan or attach tō there. The two infinitive shapes do not themselves determine the actor of a finite purpose clause.
% mdparagraph:P156:end

% teaching: NEW-U11-VOC-FEOHTAN
% teaching: NEW-U11-VOC-WITH
% teaching: NEW-U11-GRAM-u11-feohtan-indicative
% teaching: NEW-U11-GRAM-u11-feohtan-opponent
% teaching: NEW-U11-GRAM-u11-with-accusative
% mdparagraph:P157:begin
\textbf{Feohtan}, “fight; contend”, is strong class III. The present has eo except for selected \textbf{fiehtst} and \textbf{fieht} in the second and third singular. The preterite has \textbf{feaht} in first and third singular, but \textbf{fuht-} in second singular and all plural persons: fuhte, fuhton. There is no added dental preterite suffix.
% mdparagraph:P157:end

% teaching: NEW-U11-PAR-feohtan-IND
% table: NEW-U11-PAR-feohtan-IND
% mdtable:T30:begin
% course-table-style: 2026-10
% course-panel:T30-P01:begin
\par\begingroup
\fontsize{10.5}{13.2}\selectfont
\setlength{\LTleft}{0pt}\setlength{\LTright}{\fill}
\noindent
\begin{tabular}{@{}>{\raggedright\arraybackslash}p{39.459mm}>{\raggedright\arraybackslash}p{39.459mm}>{\raggedright\arraybackslash}p{39.459mm}@{}}
\toprule
\textbf{Person and number} & \textbf{Present} & \textbf{Preterite} \\
\midrule
1 sg. & feohte & feaht \\
2 sg. & fiehtst & fuhte \\
3 sg. & fieht & feaht \\
1–3 pl. & feohtað & fuhton \\
\bottomrule
\end{tabular}
\par\endgroup
% course-panel:T30-P01:end
% mdtable:T30:end
% mdparagraph:P158:begin
The opponent follows \textbf{wið + accusative}, “against; with” in the sense of fighting someone. This is opposition, not accompaniment with mid.
% mdparagraph:P158:end

% teaching: NEW-U11-VOC-GENEALAECAN
% teaching: NEW-U11-GRAM-u11-genealaecan-indicative
% teaching: NEW-U11-GRAM-u11-genealaecan-dative
% mdparagraph:P159:begin
\textbf{Genēalǣcan}, “approach; draw near”, takes a \textbf{dative} person, place or event. A time can draw near an event in the same construction. The long stem is \textbf{genēalǣc-}; present second and third singular directly add -st and -ð. In the selected weak preterite the voiceless c is followed by dental t: \textbf{genēalǣc-te}.
% mdparagraph:P159:end

% teaching: NEW-U11-PAR-genealaecan-IND
% table: NEW-U11-PAR-genealaecan-IND
% mdtable:T31:begin
% course-table-style: 2026-10
% course-panel:T31-P01:begin
\par\begingroup
\fontsize{10.5}{13.2}\selectfont
\setlength{\LTleft}{0pt}\setlength{\LTright}{\fill}
\noindent
\begin{tabular}{@{}>{\raggedright\arraybackslash}p{39.459mm}>{\raggedright\arraybackslash}p{39.459mm}>{\raggedright\arraybackslash}p{39.459mm}@{}}
\toprule
\textbf{Person and number} & \textbf{Present} & \textbf{Preterite} \\
\midrule
1 sg. & genēalǣce & genēalǣcte \\
2 sg. & genēalǣcst & genēalǣctest \\
3 sg. & genēalǣcð & genēalǣcte \\
1–3 pl. & genēalǣcað & genēalǣcton \\
\bottomrule
\end{tabular}
\par\endgroup
% course-panel:T31-P01:end
% mdtable:T31:end
% mdparagraph:P160:begin
The second singular preterite has -test and the plural -ton. Attested \textbf{genēalǣhte} is an accepted alternative; the complete selected series follows Cook’s -cte formation and is the one to learn here.
% mdparagraph:P160:end

% teaching: NEW-U11-VOC-HYDAN
% teaching: NEW-U11-GRAM-u11-hydan-indicative
% mdparagraph:P161:begin
\textbf{Hȳdan}, “hide; conceal”, takes an \textbf{accusative} object. Its weak class I dental stem has contracted present \textbf{hȳtst, hȳt}, but hȳd- before -e and -að. The preterite adds the dental suffix to hȳd-, giving \textbf{hȳdde}; keep both d letters. The resulting hidden state may be expressed by the passive construction below.
% mdparagraph:P161:end

% teaching: NEW-U11-PAR-hydan-IND
% table: NEW-U11-PAR-hydan-IND
% mdtable:T32:begin
% course-table-style: 2026-10
% course-panel:T32-P01:begin
\par\begingroup
\fontsize{10.5}{13.2}\selectfont
\setlength{\LTleft}{0pt}\setlength{\LTright}{\fill}
\noindent
\begin{tabular}{@{}>{\raggedright\arraybackslash}p{39.459mm}>{\raggedright\arraybackslash}p{39.459mm}>{\raggedright\arraybackslash}p{39.459mm}@{}}
\toprule
\textbf{Person and number} & \textbf{Present} & \textbf{Preterite} \\
\midrule
1 sg. & hȳde & hȳdde \\
2 sg. & hȳtst & hȳddest \\
3 sg. & hȳt & hȳdde \\
1–3 pl. & hȳdað & hȳddon \\
\bottomrule
\end{tabular}
\par\endgroup
% course-panel:T32-P01:end
% mdtable:T32:end
% teaching: NEW-U11-VOC-BRINGAN
% teaching: NEW-U11-GRAM-u11-bringan-indicative
% teaching: NEW-U11-GRAM-u11-bringan-destination
% mdparagraph:P162:begin
\textbf{Bringan}, “bring”, takes the person or thing brought in the \textbf{accusative}, with \textbf{tō + dative} for its destination, which may itself be a person. Present second singular is \textbf{bringst}, third singular \textbf{bringð}. The selected preterite is weak \textbf{brōhte}, with changed stem \textbf{brōh-} plus dental t. The vowel change does not by itself make the preterite strong.
% mdparagraph:P162:end

% teaching: NEW-U11-PAR-bringan-IND
% table: NEW-U11-PAR-bringan-IND
% mdtable:T33:begin
% course-table-style: 2026-10
% course-panel:T33-P01:begin
\par\begingroup
\fontsize{10.5}{13.2}\selectfont
\setlength{\LTleft}{0pt}\setlength{\LTright}{\fill}
\noindent
\begin{tabular}{@{}>{\raggedright\arraybackslash}p{39.459mm}>{\raggedright\arraybackslash}p{39.459mm}>{\raggedright\arraybackslash}p{39.459mm}@{}}
\toprule
\textbf{Person and number} & \textbf{Present} & \textbf{Preterite} \\
\midrule
1 sg. & bringe & brōhte \\
2 sg. & bringst & brōhtest \\
3 sg. & bringð & brōhte \\
1–3 pl. & bringað & brōhton \\
\bottomrule
\end{tabular}
\par\endgroup
% course-panel:T33-P01:end
% mdtable:T33:end
% mdparagraph:P163:begin
The second singular preterite adds -st to brōhte; all plural persons have brōhton. A separate strong variant is not required.
% mdparagraph:P163:end

\reserve{7\baselineskip}
% mdparagraph:P164:begin
\subsection*{Participles and passive predicates}\BookHeading{nav-ch11-learner-0015}{subsubsection}{Participles and passive predicates}{0}
% mdparagraph:P164:end

% teaching: NEW-U11-GRAM-u11-passive-predicate-bases
% teaching: NEW-U11-GRAM-u11-fon-catch-accusative
% mdparagraph:P165:begin
\textbf{Fōn},\BookPoint{bk-ch11-learner-be45aa398fa7} “catch”, now takes the caught animal, person or thing in the \textbf{accusative}. In the passive, that entity is the subject. The earlier passive with wesan can use a bare past participle, including with a plural subject. Use that construction with the new bases below. Do not invent plural adjective endings from the base alone.
% mdparagraph:P165:end

% mdparagraph:P166:begin
The present-participle ending is \textbf{-nde}: the ordinary linking vowel is e, class II has -iende, and contracted fōn has \textbf{fōnde}. The strong past participles have -en with their own vowel stems: \textbf{fang-en, fund-en}. The weak past participles have the dental ending \textbf{-ed} or class II \textbf{-od}. These are the two base forms of each verb, not full participial adjective declensions.
% mdparagraph:P166:end

% teaching: NEW-U11-GRAM-u11-fon-participle-bases
\reserve{7\baselineskip}
% mdparagraph:P167:begin
\textbf{fōn.}
% mdparagraph:P167:end

% teaching: NEW-U11-PAR-fon-PART
% table: NEW-U11-PAR-fon-PART
% mdtable:T34:begin
\par\begingroup
\fontsize{10.5}{13.2}\selectfont
\setlength{\LTleft}{0pt}\setlength{\LTright}{\fill}
\noindent
\begin{tabular}{@{}>{\raggedright\arraybackslash}p{69.84mm}>{\raggedright\arraybackslash}p{51.35mm}@{}}
\toprule
\textbf{Form} & \textbf{Base} \\
\midrule
Present participle & fōnde \\
Past participle & fangen \\
\bottomrule
\end{tabular}
\par\endgroup

% mdtable:T34:end
% teaching: NEW-U11-GRAM-u11-findan-participle-bases
\reserve{7\baselineskip}
% mdparagraph:P168:begin
\textbf{findan.}
% mdparagraph:P168:end

% teaching: NEW-U11-PAR-findan-PART
% table: NEW-U11-PAR-findan-PART
% mdtable:T35:begin
\par\begingroup
\fontsize{10.5}{13.2}\selectfont
\setlength{\LTleft}{0pt}\setlength{\LTright}{\fill}
\noindent
\begin{tabular}{@{}>{\raggedright\arraybackslash}p{69.84mm}>{\raggedright\arraybackslash}p{51.35mm}@{}}
\toprule
\textbf{Form} & \textbf{Base} \\
\midrule
Present participle & findende \\
Past participle & funden \\
\bottomrule
\end{tabular}
\par\endgroup

% mdtable:T35:end
% teaching: NEW-U11-GRAM-u11-fedan-participle-bases
\reserve{7\baselineskip}
% mdparagraph:P169:begin
\textbf{fēdan.}
% mdparagraph:P169:end

% teaching: NEW-U11-PAR-fedan-PART
% table: NEW-U11-PAR-fedan-PART
% mdtable:T36:begin
\par\begingroup
\fontsize{10.5}{13.2}\selectfont
\setlength{\LTleft}{0pt}\setlength{\LTright}{\fill}
\noindent
\begin{tabular}{@{}>{\raggedright\arraybackslash}p{69.84mm}>{\raggedright\arraybackslash}p{51.35mm}@{}}
\toprule
\textbf{Form} & \textbf{Base} \\
\midrule
Present participle & fēdende \\
Past participle & fēded \\
\bottomrule
\end{tabular}
\par\endgroup

% mdtable:T36:end
% teaching: NEW-U11-GRAM-u11-geweorthian-participle-bases
\reserve{7\baselineskip}
% mdparagraph:P170:begin
\textbf{geweorðian.}
% mdparagraph:P170:end

% teaching: NEW-U11-PAR-geweorthian-PART
% table: NEW-U11-PAR-geweorthian-PART
% mdtable:T37:begin
\par\begingroup
\fontsize{10.5}{13.2}\selectfont
\setlength{\LTleft}{0pt}\setlength{\LTright}{\fill}
\noindent
\begin{tabular}{@{}>{\raggedright\arraybackslash}p{64.38mm}>{\raggedright\arraybackslash}p{56.81mm}@{}}
\toprule
\textbf{Form} & \textbf{Base} \\
\midrule
Present participle & geweorðiende \\
Past participle & geweorðod \\
\bottomrule
\end{tabular}
\par\endgroup

% mdtable:T37:end
% teaching: NEW-U11-GRAM-u11-hydan-participle-bases
\reserve{7\baselineskip}
% mdparagraph:P171:begin
\textbf{hȳdan.}
% mdparagraph:P171:end

% teaching: NEW-U11-PAR-hydan-PART
% table: NEW-U11-PAR-hydan-PART
% mdtable:T38:begin
\par\begingroup
\fontsize{10.5}{13.2}\selectfont
\setlength{\LTleft}{0pt}\setlength{\LTright}{\fill}
\noindent
\begin{tabular}{@{}>{\raggedright\arraybackslash}p{69.84mm}>{\raggedright\arraybackslash}p{51.35mm}@{}}
\toprule
\textbf{Form} & \textbf{Base} \\
\midrule
Present participle & hȳdende \\
Past participle & hȳded \\
\bottomrule
\end{tabular}
\par\endgroup

% mdtable:T38:end
% mdparagraph:P172:begin
The past participles fēded and hȳded also have contracted forms fēd and hȳd. These are acceptable variants; the tables retain fēded and hȳded as the selected bases.
% mdparagraph:P172:end

% mdparagraph:P173:begin
Thus the plural subject still controls a plural auxiliary, while the chosen predicate can remain bare. Present participle forms are learned morphologically here; their display does not create a general progressive tense. Nor do these pairs introduce a new perfect tense or all possible prefixed variants. The exact readings supply their own additional realizations.
% mdparagraph:P173:end

\reserve{7\baselineskip}
% mdparagraph:P174:begin
\subsection*{Departure commands}\BookHeading{nav-ch11-learner-0016}{subsubsection}{Departure commands}{0}
% mdparagraph:P174:end

% teaching: NEW-U11-GRAM-u11-gewitan-imperative
% mdparagraph:P175:begin
The earlier \textbf{gewītan}, “depart; go away”, now has both direct commands. Singular \textbf{gewīt} uses the bare long-vowel stem; plural \textbf{gewītað} adds -að. Retain the earlier departure construction \textbf{fram + dative} for the person or place left.
% mdparagraph:P175:end

% teaching: NEW-U11-PAR-gewitan-IMP
% table: NEW-U11-PAR-gewitan-IMP
% mdtable:T39:begin
% course-table-style: 2026-10
% course-panel:T39-P01:begin
\par\begingroup
\fontsize{10.5}{13.2}\selectfont
\setlength{\LTleft}{0pt}\setlength{\LTright}{\fill}
\noindent
\begin{tabular}{@{}>{\raggedright\arraybackslash}p{63.342mm}>{\raggedright\arraybackslash}p{57.846mm}@{}}
\toprule
\textbf{Person and number} & \textbf{Command} \\
\midrule
2 sg. & gewīt \\
2 pl. & gewītað \\
\bottomrule
\end{tabular}
\par\endgroup
% course-panel:T39-P01:end
% mdtable:T39:end
\reserve{7\baselineskip}
% mdparagraph:P176:begin
\section*{7. Additional meanings and constructions}\BookHeading{nav-ch11-learner-0017}{subsection}{7. Additional meanings and constructions}{1}
% mdparagraph:P176:end

% teaching: NEW-U11-VOC-HIERAN
% teaching: NEW-U11-GRAM-u11-hieran-obey-dative
% mdparagraph:P177:begin
\textbf{Hīeran}, besides “hear”, can mean \textbf{“obey” with a dative} person. Coordinated subjects require a plural verb. Do not infer an accusative from the English object “him”.
% mdparagraph:P177:end

% teaching: NEW-U11-VOC-DON
% teaching: NEW-U11-VOC-ON
% teaching: NEW-U11-GRAM-u11-don-placement
% mdparagraph:P178:begin
\textbf{Dōn}, “put; place”, takes an \textbf{accusative} object. For placement into or onto something, use \textbf{on + accusative} as the goal. This differs from the already taught on + dative for location. In a negative command, retain the taught ne with the appropriate singular or plural command.
% mdparagraph:P178:end

% teaching: NEW-U11-VOC-THURH
% teaching: NEW-U11-GRAM-u11-thurh-accusative
% mdparagraph:P179:begin
\textbf{Þurh + accusative}, “through”, gives the spatial path. Learn the þ spelling for composition. The exact sea reading has \textbf{ðurh}, with ð, and retains it in both its path phrase and the supplied expression ðurh eall, “entirely”.
% mdparagraph:P179:end

% teaching: NEW-U11-VOC-RISAN
% teaching: NEW-U11-VOC-UP
% teaching: NEW-U11-VOC-ONGEAN
% teaching: NEW-U11-GRAM-u11-risan-up-against
% mdparagraph:P180:begin
The familiar \textbf{rīsan}, “rise”, combines with \textbf{ūp}, “up”. For rising against an opponent or adverse power, add \textbf{ongēan + accusative}, “against”. Ūp itself gives no opponent. Keep the already learned unprefixed rīsan forms; this phrase does not require a new verb ārīsan.
% mdparagraph:P180:end

% teaching: NEW-U11-VOC-BUTAN
% teaching: NEW-U11-VOC-WIF
% teaching: NEW-U11-GRAM-u11-butan-dative
% mdparagraph:P181:begin
\textbf{Būtan + dative}, “besides; apart from”, excludes a group from the number just stated. In this use \textbf{wīf} means “woman”, as well as its earlier “wife”; it remains grammatically neuter. Keep the case governed by būtan when the phrase names women and children.
% mdparagraph:P181:end

% teaching: NEW-U11-VOC-BIDDAN
% teaching: NEW-U11-GRAM-u11-biddan-person
% mdparagraph:P182:begin
In the parenthetic \textbf{ic ðē bidde}, “I beseech you; please”, biddan takes the person entreated in the \textbf{accusative}. The same surface ðē can be dative elsewhere, but the verb determines its case here. This bounded request does not introduce every possible complement of biddan.
% mdparagraph:P182:end

% teaching: NEW-U11-VOC-DOM
% teaching: NEW-U11-VOC-OFT
% mdparagraph:P183:begin
\textbf{Dōm} can mean “doom; fate” in addition to the earlier “judgment”. Keep its familiar masculine forms. \textbf{Oft} means “often”; like hēr and mā it is an invariant adverb.
% mdparagraph:P183:end

% teaching: NEW-U11-VOC-NALAES
% teaching: NEW-U11-VOC-HOLUNGA
% teaching: NEW-U11-GRAM-u11-nalaes-negation
% mdparagraph:P184:begin
Learn \textbf{nālæs hōlunga}, “not at all in vain”, as the selected emphatic phrase: nālæs, “not at all”, denies hōlunga, “in vain”. It says that the action was not fruitless. Do not mechanically add another ne. This phrase does not teach the general system of negative concord.
% mdparagraph:P184:end

% mdparagraph:P185:begin
The ordinary senses and forms from earlier units remain in force: helpan governs a dative person; sęllan takes an accusative thing and dative recipient; clothing takes an accusative person; possessive his and hiera name the possessor without agreeing as an adjective. Familiar meanings do not erase the particular new case constructions just learned.
% mdparagraph:P185:end

\reserve{7\baselineskip}
% mdparagraph:P186:begin
\section*{8. Practice: forms}\BookHeading{nav-ch11-learner-0018}{subsection}{8. Practice: forms}{1}
% mdparagraph:P186:end

% mdparagraph:P187:begin
Give the requested forms, retaining their vowel lengths and spelling. Explain the relevant formation rather than merely naming a class. Reference variants are accepted where noted, but are not additional memorization targets.
% mdparagraph:P187:end

\reserve{5\baselineskip}
\par\begin{minipage}{\linewidth}\raggedright
% mdparagraph:P188:begin
\textbf{F01.}
% mdparagraph:P188:end

% item: NEW-U11-F01
% mdparagraph:P189:begin
Give the complete nominative, genitive, dative and accusative singular and plural of bæð and geat. Then give the nominative plural and dative plural of sīðfæt. Explain where each lexical vowel changes.
% mdparagraph:P189:end

\end{minipage}\par\addvspace{4pt}

\par\begin{minipage}{\linewidth}\raggedright
% mdparagraph:P190:begin
\textbf{F02.}
% mdparagraph:P190:end

% item: NEW-U11-F02
% mdparagraph:P191:begin
Give the complete nominative, genitive, dative and accusative singular and plural of seolh. Explain which cells retain h and which have the long-vowel stem.
% mdparagraph:P191:end

\end{minipage}\par\addvspace{4pt}

\par\begin{minipage}{\linewidth}\raggedright
% mdparagraph:P192:begin
\textbf{F03.}
% mdparagraph:P192:end

% item: NEW-U11-F03
% mdparagraph:P193:begin
Give the complete nominative, genitive, dative and accusative singular and plural of sumer and winter, using the selected tables. Explain why the stems before endings differ. The noted -a variants are acceptable but are not required from memory.
% mdparagraph:P193:end

\end{minipage}\par\addvspace{4pt}

\par\begin{minipage}{\linewidth}\raggedright
% mdparagraph:P194:begin
\textbf{F04.}
% mdparagraph:P194:end

% item: NEW-U11-F04
% mdparagraph:P195:begin
Give these forms: blæc, strong masculine accusative singular, strong feminine nominative singular, and weak masculine nominative singular; stille, strong masculine accusative singular; fægerra, nominative singular in each of the three genders; dēorwierðra, dative plural in all genders. Explain the stem and weak-ending choices.
% mdparagraph:P195:end

\end{minipage}\par\addvspace{4pt}

\par\begin{minipage}{\linewidth}\raggedright
% mdparagraph:P196:begin
\textbf{F05.}
% mdparagraph:P196:end

% item: NEW-U11-F05
% mdparagraph:P197:begin
Give the complete present and preterite indicative of scęððan and genēalǣcan: first, second and third person in both numbers for each tense. Explain the selected stem and dental-preterite formation.
% mdparagraph:P197:end

\end{minipage}\par\addvspace{4pt}

\par\begin{minipage}{\linewidth}\raggedright
% mdparagraph:P198:begin
\textbf{F06.}
% mdparagraph:P198:end

% item: NEW-U11-F06
% mdparagraph:P199:begin
For feohtan and bringan give the present indicative second and third singular, the preterite first and second singular, and the preterite first plural. Explain which preterite belongs to a strong vowel series and which contains a dental suffix.
% mdparagraph:P199:end

\end{minipage}\par\addvspace{4pt}

\par\begin{minipage}{\linewidth}\raggedright
% mdparagraph:P200:begin
\textbf{F07.}
% mdparagraph:P200:end

% item: NEW-U11-F07
% mdparagraph:P201:begin
Give the base present and past participles of fōn, findan, fēdan, geweorðian and hȳdan. Identify the present-participle ending and the strong or weak past-participle ending in each pair.
% mdparagraph:P201:end

\end{minipage}\par\addvspace{4pt}

\par\begin{minipage}{\linewidth}\raggedright
% mdparagraph:P202:begin
\textbf{F08.}
% mdparagraph:P202:end

% item: NEW-U11-F08
% mdparagraph:P203:begin
Give the singular and plural commands of gewītan. Then give hȳdan preterite indicative third singular and third plural, and scēawian and geweorðian present indicative third singular and third plural. Explain what changes when one addressee or subject becomes several.
% mdparagraph:P203:end

\end{minipage}\par\addvspace{4pt}

\reserve{7\baselineskip}
% mdparagraph:P204:begin
\section*{9. Practice: clauses and meaning}\BookHeading{nav-ch11-learner-0019}{subsection}{9. Practice: clauses and meaning}{1}
% mdparagraph:P204:end

% mdparagraph:P205:begin
Translate each printed Old English input and answer the questions attached to it. Where two or three inputs are given, treat all of them. The stated situation is part of the task.
% mdparagraph:P205:end

\reserve{5\baselineskip}
\par\begin{minipage}{\linewidth}\raggedright
% mdparagraph:P206:begin
\textbf{A01.}
% mdparagraph:P206:end

% item: NEW-U11-A01
% mdparagraph:P207:begin
Translate. Mark the subordinate and main clauses, identify both finite verbs, and explain the two occurrences of ðā.
% mdparagraph:P207:end

% mdparagraph:P208:begin
Ðā se munuc cōm, ðā sealde se mann him wæter.
% mdparagraph:P208:end

\end{minipage}\par\addvspace{4pt}

\par\begin{minipage}{\linewidth}\raggedright
% mdparagraph:P209:begin
\textbf{A02.}
% mdparagraph:P209:end

% item: NEW-U11-A02
% mdparagraph:P210:begin
The speaker describes what happens on repeated visits. Translate; identify the clause introduced by the first ðonne and the role of the second.
% mdparagraph:P210:end

% mdparagraph:P211:begin
Ðonne gē cumað, ðonne helpe ic ēow.
% mdparagraph:P211:end

\end{minipage}\par\addvspace{4pt}

\par\begin{minipage}{\linewidth}\raggedright
% mdparagraph:P212:begin
\textbf{A03.}
% mdparagraph:P212:end

% item: NEW-U11-A03
% mdparagraph:P213:begin
Translate both statements. Explain the relationship between the actions and whether either statement by itself says that one action caused the other.
% mdparagraph:P213:end

% mdparagraph:P214:begin
(a) Ðā-hwīle-ðe hē ǣt, ic bær wæter.
% mdparagraph:P214:end

% mdparagraph:P215:begin
(b) Mid-ðȳ hē ǣt, ic bær wæter.
% mdparagraph:P215:end

\end{minipage}\par\addvspace{4pt}

\par\begin{minipage}{\linewidth}\raggedright
% mdparagraph:P216:begin
\textbf{A04.}
% mdparagraph:P216:end

% item: NEW-U11-A04
% mdparagraph:P217:begin
Both instructions concern a coming that has not yet occurred. Translate and identify the mood of cume; explain why the English future setting does not require a new Old English future ending.
% mdparagraph:P217:end

% mdparagraph:P218:begin
(a) Ne far, ǣr ic cume.
% mdparagraph:P218:end

% mdparagraph:P219:begin
(b) Wes hēr, oð-ðæt ic cume.
% mdparagraph:P219:end

\end{minipage}\par\addvspace{4pt}

\par\begin{minipage}{\linewidth}\raggedright
% mdparagraph:P220:begin
\textbf{A05.}
% mdparagraph:P220:end

% item: NEW-U11-A05
% mdparagraph:P221:begin
Translate. State which clause gives the reason and which gives the action being explained; identify the finite form in the reason clause.
% mdparagraph:P221:end

% mdparagraph:P222:begin
Ic helpe ðē, for-ðon-ðe ðū eart untrum.
% mdparagraph:P222:end

\end{minipage}\par\addvspace{4pt}

\par\begin{minipage}{\linewidth}\raggedright
% mdparagraph:P223:begin
\textbf{A06.}
% mdparagraph:P223:end

% item: NEW-U11-A06
% mdparagraph:P224:begin
Translate each sentence and classify its ðæt clause as content, intended purpose or asserted result. Parse the finite verb inside that clause.
% mdparagraph:P224:end

% mdparagraph:P225:begin
(a) Ic wāt ðæt se munuc cymð.
% mdparagraph:P225:end

% mdparagraph:P226:begin
(b) Se munuc cymð, ðæt hē mē helpe.
% mdparagraph:P226:end

% mdparagraph:P227:begin
(c) Se munuc is swā gōd, ðæt hē mē hilpð.
% mdparagraph:P227:end

\end{minipage}\par\addvspace{4pt}

\par\begin{minipage}{\linewidth}\raggedright
% mdparagraph:P228:begin
\textbf{A07.}
% mdparagraph:P228:end

% item: NEW-U11-A07
% mdparagraph:P229:begin
Translate. Identify the outcome the man intends to prevent; explain whether the sentence states that the animal actually took the food.
% mdparagraph:P229:end

% mdparagraph:P230:begin
Se mann hēold ðæt dēor, ðȳ-lǣs-ðe hit ðone męte nōme.
% mdparagraph:P230:end

\end{minipage}\par\addvspace{4pt}

\par\begin{minipage}{\linewidth}\raggedright
% mdparagraph:P231:begin
\textbf{A08.}
% mdparagraph:P231:end

% item: NEW-U11-A08
% mdparagraph:P232:begin
Translate and parse sīe. The speaker accepts that the man is sick: can the selected subjunctive still be used, and what relation does the clause express?
% mdparagraph:P232:end

% mdparagraph:P233:begin
Ðēah hē sīe untrum, hē hilpð mē.
% mdparagraph:P233:end

\end{minipage}\par\addvspace{4pt}

\par\begin{minipage}{\linewidth}\raggedright
% mdparagraph:P234:begin
\textbf{A09.}
% mdparagraph:P234:end

% item: NEW-U11-A09
% mdparagraph:P235:begin
In (a), his presence is an open question. In (b), he is absent now. Translate both, identify their finite moods, and state the time referred to by wǣre and hulpe.
% mdparagraph:P235:end

% mdparagraph:P236:begin
(a) Gif hē is hēr, hē hilpð mē.
% mdparagraph:P236:end

% mdparagraph:P237:begin
(b) Gif hē wǣre hēr, hē hulpe mē.
% mdparagraph:P237:end

\end{minipage}\par\addvspace{4pt}

\par\begin{minipage}{\linewidth}\raggedright
% mdparagraph:P238:begin
\textbf{A10.}
% mdparagraph:P238:end

% item: NEW-U11-A10
% mdparagraph:P239:begin
Translate each sentence. Distinguish interrogative hwonne, temporal ðonne, and comparative ðonne; do not infer the function from the spelling alone.
% mdparagraph:P239:end

% mdparagraph:P240:begin
(a) Hwonne cymð se munuc?
% mdparagraph:P240:end

% mdparagraph:P241:begin
(b) Ðonne se munuc cymð, ic helpe him.
% mdparagraph:P241:end

% mdparagraph:P242:begin
(c) Se munuc is ieldra ðonne se fiscere.
% mdparagraph:P242:end

\end{minipage}\par\addvspace{4pt}

\par\begin{minipage}{\linewidth}\raggedright
% mdparagraph:P243:begin
\textbf{A11.}
% mdparagraph:P243:end

% item: NEW-U11-A11
% mdparagraph:P244:begin
Translate. Who is understood to do the helping in each sentence? Explain why a finite purpose clause can make a different subject explicit.
% mdparagraph:P244:end

% mdparagraph:P245:begin
(a) Se fiscere cōm tō helpanne ðǣm męnn.
% mdparagraph:P245:end

% mdparagraph:P246:begin
(b) Se fiscere cōm, ðæt se munuc ðǣm męnn hulpe.
% mdparagraph:P246:end

\end{minipage}\par\addvspace{4pt}

\par\begin{minipage}{\linewidth}\raggedright
% mdparagraph:P247:begin
\textbf{A12.}
% mdparagraph:P247:end

% item: NEW-U11-A12
% mdparagraph:P248:begin
Translate. Identify the antecedent of ðe and explain why this is neither a direct when-question nor a free-standing conditional.
% mdparagraph:P248:end

% mdparagraph:P249:begin
Hit is se dæg ðe hē cōm.
% mdparagraph:P249:end

\end{minipage}\par\addvspace{4pt}

\reserve{7\baselineskip}
% mdparagraph:P250:begin
\section*{10. Controlled composition}\BookHeading{nav-ch11-learner-0020}{subsection}{10. Controlled composition}{1}
% mdparagraph:P250:end

% mdparagraph:P251:begin
Translate into Old English using the stated construction. A grammatical alternative may still miss a requested target; distinguish those two judgments. Preserve the people, number, time, purpose and other relations specified in the English.
% mdparagraph:P251:end

\reserve{5\baselineskip}
\par\begin{minipage}{\linewidth}\raggedright
% mdparagraph:P252:begin
\textbf{T01.}
% mdparagraph:P252:end

% item: NEW-U11-T01
% mdparagraph:P253:begin
Translate using past ðā … ðā: When the father came, the son gave him water.
% mdparagraph:P253:end

\end{minipage}\par\addvspace{4pt}

\par\begin{minipage}{\linewidth}\raggedright
% mdparagraph:P254:begin
\textbf{T02.}
% mdparagraph:P254:end

% item: NEW-U11-T02
% mdparagraph:P255:begin
Translate as a recurring event, using ðonne … ðonne: When the fishermen come, the king gives them food.
% mdparagraph:P255:end

\end{minipage}\par\addvspace{4pt}

\par\begin{minipage}{\linewidth}\raggedright
% mdparagraph:P256:begin
\textbf{T03.}
% mdparagraph:P256:end

% item: NEW-U11-T03
% mdparagraph:P257:begin
Translate with mid-ðȳ: While the fishermen sat in the house, the son carried water.
% mdparagraph:P257:end

\end{minipage}\par\addvspace{4pt}

\par\begin{minipage}{\linewidth}\raggedright
% mdparagraph:P258:begin
\textbf{T04.}
% mdparagraph:P258:end

% item: NEW-U11-T04
% mdparagraph:P259:begin
Address one person. The speaker has not yet arrived. Translate with ǣr: Do not go before I come.
% mdparagraph:P259:end

\end{minipage}\par\addvspace{4pt}

\par\begin{minipage}{\linewidth}\raggedright
% mdparagraph:P260:begin
\textbf{T05.}
% mdparagraph:P260:end

% item: NEW-U11-T05
% mdparagraph:P261:begin
Address several people. The fishermen have not yet arrived. Translate with oð-ðæt: Be in the house until the fishermen come.
% mdparagraph:P261:end

\end{minipage}\par\addvspace{4pt}

\par\begin{minipage}{\linewidth}\raggedright
% mdparagraph:P262:begin
\textbf{T06.}
% mdparagraph:P262:end

% item: NEW-U11-T06
% mdparagraph:P263:begin
Translate with for-ðon-ðe and the adjective ðurstig: The monk gives the man water because the man is thirsty.
% mdparagraph:P263:end

\end{minipage}\par\addvspace{4pt}

\par\begin{minipage}{\linewidth}\raggedright
% mdparagraph:P264:begin
\textbf{T07.}
% mdparagraph:P264:end

% item: NEW-U11-T07
% mdparagraph:P265:begin
Translate with a finite ðæt purpose clause: The fishermen came so that the monk might help them.
% mdparagraph:P265:end

\end{minipage}\par\addvspace{4pt}

\par\begin{minipage}{\linewidth}\raggedright
% mdparagraph:P266:begin
\textbf{T08.}
% mdparagraph:P266:end

% item: NEW-U11-T08
% mdparagraph:P267:begin
Translate the result as an actual past event, using swā … ðæt: The men were so hungry that they ate the food.
% mdparagraph:P267:end

\end{minipage}\par\addvspace{4pt}

\par\begin{minipage}{\linewidth}\raggedright
% mdparagraph:P268:begin
\textbf{T09.}
% mdparagraph:P268:end

% item: NEW-U11-T09
% mdparagraph:P269:begin
Translate with ðȳ-lǣs-ðe: The fishermen held the animals, lest the animals take our food.
% mdparagraph:P269:end

\end{minipage}\par\addvspace{4pt}

\par\begin{minipage}{\linewidth}\raggedright
% mdparagraph:P270:begin
\textbf{T10.}
% mdparagraph:P270:end

% item: NEW-U11-T10
% mdparagraph:P271:begin
The monk was in fact sick. Translate using the selected concessive pattern: Although the monk was sick, he helped the man.
% mdparagraph:P271:end

\end{minipage}\par\addvspace{4pt}

\par\begin{minipage}{\linewidth}\raggedright
% mdparagraph:P272:begin
\textbf{T11.}
% mdparagraph:P272:end

% item: NEW-U11-T11
% mdparagraph:P273:begin
The king is absent now. Translate the contrary-present condition: If the king were here, he would help us.
% mdparagraph:P273:end

\end{minipage}\par\addvspace{4pt}

\reserve{7\baselineskip}
% mdparagraph:P274:begin
\section*{11. Cook IV.16–17 and V.18}\BookHeading{nav-ch11-learner-0021}{subsection}{11. Cook IV.16–17 and V.18}{1}
% mdparagraph:P274:end

% mdparagraph:P275:begin
Translate the three complete exercise sets. Their original wording, grammatical references and parenthetical directions are retained. \textbf{Om.} means “omit”: it applies to the indicated English word, not to the meaning of the sentence. \textbf{P.p.} means “past participle”. Material labelled local help is supplied for use; do not treat it as an additional recall test.
% mdparagraph:P275:end

\reserve{7\baselineskip}
% mdparagraph:P276:begin
\subsection*{IV.16}\BookHeading{nav-ch11-learner-0022}{subsubsection}{IV.16}{0}
% mdparagraph:P276:end

% mdparagraph:P277:begin
For the unspecified “your” command in IV.16.16, state whether you address one person or several and keep the command and possessive consistent. For “all” in IV.16.13, identify whether you mean everything or all people. Use the participle bases from §6 in the passive items.
% mdparagraph:P277:end

\reserve{5\baselineskip}
\par\begin{minipage}{\linewidth}\raggedright
% mdparagraph:P278:begin
\textbf{IV.16.1.}
% mdparagraph:P278:end

% item: EX-IV-16-01
% mdparagraph:P279:begin
Gold is more precious (64) than silver.
% mdparagraph:P279:end

\end{minipage}\par\addvspace{4pt}

\par\begin{minipage}{\linewidth}\raggedright
% mdparagraph:P280:begin
\textbf{IV.16.2.}
% mdparagraph:P280:end

% item: EX-IV-16-02
% mdparagraph:P281:begin
Righteousness is more precious than gold.
% mdparagraph:P281:end

\end{minipage}\par\addvspace{4pt}

\par\begin{minipage}{\linewidth}\raggedright
% mdparagraph:P282:begin
\textbf{IV.16.3.}
% mdparagraph:P282:end

% item: EX-IV-16-03
% mdparagraph:P283:begin
The bath is hot.
% mdparagraph:P283:end

\end{minipage}\par\addvspace{4pt}

\par\begin{minipage}{\linewidth}\raggedright
% mdparagraph:P284:begin
\textbf{IV.16.4.}
% mdparagraph:P284:end

% item: EX-IV-16-04
% mdparagraph:P285:begin
The baths are hot.
% mdparagraph:P285:end

\end{minipage}\par\addvspace{4pt}

\par\begin{minipage}{\linewidth}\raggedright
% mdparagraph:P286:begin
\textbf{IV.16.5.}
% mdparagraph:P286:end

% item: EX-IV-16-05
% mdparagraph:P287:begin
The animals are living.
% mdparagraph:P287:end

\end{minipage}\par\addvspace{4pt}

\par\begin{minipage}{\linewidth}\raggedright
% mdparagraph:P288:begin
\textbf{IV.16.6.}
% mdparagraph:P288:end

% item: EX-IV-16-06
% mdparagraph:P289:begin
Some pearls are black.
% mdparagraph:P289:end

\end{minipage}\par\addvspace{4pt}

\par\begin{minipage}{\linewidth}\raggedright
% mdparagraph:P290:begin
\textbf{IV.16.7.}
% mdparagraph:P290:end

% item: EX-IV-16-07
% mdparagraph:P291:begin
The island is beautiful.
% mdparagraph:P291:end

\end{minipage}\par\addvspace{4pt}

\par\begin{minipage}{\linewidth}\raggedright
% mdparagraph:P292:begin
\textbf{IV.16.8.}
% mdparagraph:P292:end

% item: EX-IV-16-08
% mdparagraph:P293:begin
Dolphins swim through the water.
% mdparagraph:P293:end

\end{minipage}\par\addvspace{4pt}

\par\begin{minipage}{\linewidth}\raggedright
% mdparagraph:P294:begin
\textbf{IV.16.9.}
% mdparagraph:P294:end

% item: EX-IV-16-09
% mdparagraph:P295:begin
Here are fishes often caught.
% mdparagraph:P295:end

\end{minipage}\par\addvspace{4pt}

\par\begin{minipage}{\linewidth}\raggedright
% mdparagraph:P296:begin
\textbf{IV.16.10.}
% mdparagraph:P296:end

% item: EX-IV-16-10
% mdparagraph:P297:begin
The night is more beautiful than the day.
% mdparagraph:P297:end

\end{minipage}\par\addvspace{4pt}

\par\begin{minipage}{\linewidth}\raggedright
% mdparagraph:P298:begin
\textbf{IV.16.11.}
% mdparagraph:P298:end

% item: EX-IV-16-11
% mdparagraph:P299:begin
The city was distinguished with walls and gates.
% mdparagraph:P299:end

\end{minipage}\par\addvspace{4pt}

\par\begin{minipage}{\linewidth}\raggedright
% mdparagraph:P300:begin
\textbf{IV.16.12.}
% mdparagraph:P300:end

% item: EX-IV-16-12
% mdparagraph:P301:begin
Summer comes, and then winter.
% mdparagraph:P301:end

\end{minipage}\par\addvspace{4pt}

\par\begin{minipage}{\linewidth}\raggedright
% mdparagraph:P302:begin
\textbf{IV.16.13.}
% mdparagraph:P302:end

% item: EX-IV-16-13
% mdparagraph:P303:begin
The sun shines over all.
% mdparagraph:P303:end

\end{minipage}\par\addvspace{4pt}

\par\begin{minipage}{\linewidth}\raggedright
% mdparagraph:P304:begin
\textbf{IV.16.14.}
% mdparagraph:P304:end

% item: EX-IV-16-14
% mdparagraph:P305:begin
The old man came to the house.
% mdparagraph:P305:end

\end{minipage}\par\addvspace{4pt}

\par\begin{minipage}{\linewidth}\raggedright
% mdparagraph:P306:begin
\textbf{IV.16.15.}
% mdparagraph:P306:end

% item: EX-IV-16-15
% mdparagraph:P307:begin
Here are also found lead and iron.
% mdparagraph:P307:end

\end{minipage}\par\addvspace{4pt}

\par\begin{minipage}{\linewidth}\raggedright
% mdparagraph:P308:begin
\textbf{IV.16.16.}
% mdparagraph:P308:end

% item: EX-IV-16-16
% mdparagraph:P309:begin
Put not your hand into the fire.
% mdparagraph:P309:end

\end{minipage}\par\addvspace{4pt}

\reserve{7\baselineskip}
% mdparagraph:P310:begin
\subsection*{IV.17}\BookHeading{nav-ch11-learner-0023}{subsubsection}{IV.17}{0}
% mdparagraph:P310:end

% mdparagraph:P311:begin
The original “shall” and “will” can have future reference. Keep the actual person and number of the subject. For the predicative “thirsty” use the adjective ðurstig. The numeral for twelve is familiar; the higher count in item17 is supplied separately.
% mdparagraph:P311:end

\reserve{5\baselineskip}
\par\begin{minipage}{\linewidth}\raggedright
% mdparagraph:P312:begin
\textbf{IV.17.1.}
% mdparagraph:P312:end

% item: EX-IV-17-01
% mdparagraph:P313:begin
The soldier takes the gem.
% mdparagraph:P313:end

\end{minipage}\par\addvspace{4pt}

\par\begin{minipage}{\linewidth}\raggedright
% mdparagraph:P314:begin
\textbf{IV.17.2.}
% mdparagraph:P314:end

% item: EX-IV-17-02
% mdparagraph:P315:begin
The rain injures the tower.
% mdparagraph:P315:end

\end{minipage}\par\addvspace{4pt}

\par\begin{minipage}{\linewidth}\raggedright
% mdparagraph:P316:begin
\textbf{IV.17.3.}
% mdparagraph:P316:end

% item: EX-IV-17-03
% mdparagraph:P317:begin
The whale swims.
% mdparagraph:P317:end

\end{minipage}\par\addvspace{4pt}

\par\begin{minipage}{\linewidth}\raggedright
% mdparagraph:P318:begin
\textbf{IV.17.4.}
% mdparagraph:P318:end

% item: EX-IV-17-04
% mdparagraph:P319:begin
Fishermen caught seals.
% mdparagraph:P319:end

\end{minipage}\par\addvspace{4pt}

\par\begin{minipage}{\linewidth}\raggedright
% mdparagraph:P320:begin
\textbf{IV.17.5.}
% mdparagraph:P320:end

% teaching: NEW-U11-LOCAL-HELPS
% teaching: NEW-U11-GRAM-u11-negative-copula-local
% help: U11-COOK-HELP-17-05
% mdparagraph:P321:begin
\textbf{Local help.} For “are not” in this passive, use the supplied plural auxiliary \textbf{ne sindon}. This expression is given for use here; no full negative paradigm is required.
% mdparagraph:P321:end

% item: EX-IV-17-05
% mdparagraph:P322:begin
The hungry sheep are not fed.
% mdparagraph:P322:end

\end{minipage}\par\addvspace{4pt}

\par\begin{minipage}{\linewidth}\raggedright
% mdparagraph:P323:begin
\textbf{IV.17.6.}
% mdparagraph:P323:end

% item: EX-IV-17-06
% mdparagraph:P324:begin
The man’s wife is beautiful.
% mdparagraph:P324:end

\end{minipage}\par\addvspace{4pt}

\par\begin{minipage}{\linewidth}\raggedright
% mdparagraph:P325:begin
\textbf{IV.17.7.}
% mdparagraph:P325:end

% item: EX-IV-17-07
% mdparagraph:P326:begin
Some men are black.
% mdparagraph:P326:end

\end{minipage}\par\addvspace{4pt}

\par\begin{minipage}{\linewidth}\raggedright
% mdparagraph:P327:begin
\textbf{IV.17.8.}
% mdparagraph:P327:end

% item: EX-IV-17-08
% mdparagraph:P328:begin
The carpenter will clothe his son.
% mdparagraph:P328:end

\end{minipage}\par\addvspace{4pt}

\par\begin{minipage}{\linewidth}\raggedright
% mdparagraph:P329:begin
\textbf{IV.17.9.}
% mdparagraph:P329:end

% item: EX-IV-17-09
% mdparagraph:P330:begin
Depart from me, for I am (om. a) sinful man, (om. O) Lord.
% mdparagraph:P330:end

\end{minipage}\par\addvspace{4pt}

\par\begin{minipage}{\linewidth}\raggedright
% mdparagraph:P331:begin
\textbf{IV.17.10.}
% mdparagraph:P331:end

% item: EX-IV-17-10
% mdparagraph:P332:begin
They shall drink, but ye shall be thirsty.
% mdparagraph:P332:end

\end{minipage}\par\addvspace{4pt}

\par\begin{minipage}{\linewidth}\raggedright
% mdparagraph:P333:begin
\textbf{IV.17.11.}
% mdparagraph:P333:end

% help: U11-COOK-HELP-17-11
% mdparagraph:P334:begin
\textbf{Local help.} The give-to-drink expression remains supplied: \textbf{sęle mē wæter drincan}. Sęle is the singular command “give”; mē is the dative recipient and wæter the accusative thing given. The plain infinitive drincan, without tō, says what the recipient is to do.
% mdparagraph:P334:end

% item: EX-IV-17-11
% mdparagraph:P335:begin
He said to her, Give me, I beseech thee, water to drink; for I am thirsty.
% mdparagraph:P335:end

\end{minipage}\par\addvspace{4pt}

\par\begin{minipage}{\linewidth}\raggedright
% mdparagraph:P336:begin
\textbf{IV.17.12.}
% mdparagraph:P336:end

% help: U11-COOK-HELP-17-12
% mdparagraph:P337:begin
\textbf{Local help.} Use supplied \textbf{nis}, “is not”, the third singular negative form of the copula.
% mdparagraph:P337:end

% item: EX-IV-17-12
% mdparagraph:P338:begin
He is not (om. a) God of (om. the) dead, but of (om. the) living.
% mdparagraph:P338:end

\end{minipage}\par\addvspace{4pt}

\par\begin{minipage}{\linewidth}\raggedright
% mdparagraph:P339:begin
\textbf{IV.17.13.}
% mdparagraph:P339:end

% item: EX-IV-17-13
% mdparagraph:P340:begin
The sun shines on (ofer) the righteous and the unrighteous.
% mdparagraph:P340:end

\end{minipage}\par\addvspace{4pt}

\par\begin{minipage}{\linewidth}\raggedright
% mdparagraph:P341:begin
\textbf{IV.17.14.}
% mdparagraph:P341:end

% item: EX-IV-17-14
% mdparagraph:P342:begin
Ye are (om. the) light of the world.
% mdparagraph:P342:end

\end{minipage}\par\addvspace{4pt}

\par\begin{minipage}{\linewidth}\raggedright
% mdparagraph:P343:begin
\textbf{IV.17.15.}
% mdparagraph:P343:end

% item: EX-IV-17-15
% mdparagraph:P344:begin
The twelve gates were twelve pearls.
% mdparagraph:P344:end

\end{minipage}\par\addvspace{4pt}

\par\begin{minipage}{\linewidth}\raggedright
% mdparagraph:P345:begin
\textbf{IV.17.16.}
% mdparagraph:P345:end

% item: EX-IV-17-16
% mdparagraph:P346:begin
The street of the city was gold.
% mdparagraph:P346:end

\end{minipage}\par\addvspace{4pt}

\par\begin{minipage}{\linewidth}\raggedright
% mdparagraph:P347:begin
\textbf{IV.17.17.}
% mdparagraph:P347:end

% teaching: NEW-U11-GRAM-NEW-U11-SRC-NUM4000
% help: U11-COOK-HELP-17-17
% mdparagraph:P348:begin
\textbf{Local help.} Use the supplied phrase \textbf{fēower ðūsend manna}, “four thousand men”. Manna is genitive plural after this higher numeral. Cook’s §154.c concerns this genitive construction. The phrase is supplied as a whole; a general system of higher numerals is not required.
% mdparagraph:P348:end

% item: EX-IV-17-17
% mdparagraph:P349:begin
Those that ate were four thousand men (154. c), besides women and children.
% mdparagraph:P349:end

\end{minipage}\par\addvspace{4pt}

\reserve{7\baselineskip}
% mdparagraph:P350:begin
\subsection*{V.18}\BookHeading{nav-ch11-learner-0024}{subsubsection}{V.18}{0}
% mdparagraph:P350:end

% mdparagraph:P351:begin
Item1 is a blessing wish with the given past-participle target; keep each title and dependent genitive in its proper case. In item8, the question concerns an actual past event; specify the gender of the singular person addressed as “dearest”. The bracketed “approaches” in item6 explains “draws near”, and (164) refers to the dative construction. The (113) beside hȳdan refers to its weak inflection.
% mdparagraph:P351:end

\reserve{5\baselineskip}
\par\begin{minipage}{\linewidth}\raggedright
% mdparagraph:P352:begin
\textbf{V.18.1.}
% mdparagraph:P352:end

% item: EX-V-18-01
% mdparagraph:P353:begin
Blessed (p.p.) be God, the Father of our Lord Jesus Christ, the Father of mercies, and the God of all comfort.
% mdparagraph:P353:end

\end{minipage}\par\addvspace{4pt}

\par\begin{minipage}{\linewidth}\raggedright
% mdparagraph:P354:begin
\textbf{V.18.2.}
% mdparagraph:P354:end

% item: EX-V-18-02
% mdparagraph:P355:begin
The winds were favorable.
% mdparagraph:P355:end

\end{minipage}\par\addvspace{4pt}

\par\begin{minipage}{\linewidth}\raggedright
% mdparagraph:P356:begin
\textbf{V.18.3.}
% mdparagraph:P356:end

% item: EX-V-18-03
% mdparagraph:P357:begin
He wishes to watch their journey.
% mdparagraph:P357:end

\end{minipage}\par\addvspace{4pt}

\par\begin{minipage}{\linewidth}\raggedright
% mdparagraph:P358:begin
\textbf{V.18.4.}
% mdparagraph:P358:end

% item: EX-V-18-04
% mdparagraph:P359:begin
The wind and the sea obey him.
% mdparagraph:P359:end

\end{minipage}\par\addvspace{4pt}

\par\begin{minipage}{\linewidth}\raggedright
% mdparagraph:P360:begin
\textbf{V.18.5.}
% mdparagraph:P360:end

% item: EX-V-18-05
% mdparagraph:P361:begin
Not at all in vain fought I with death, as he fought with beasts.
% mdparagraph:P361:end

\end{minipage}\par\addvspace{4pt}

\par\begin{minipage}{\linewidth}\raggedright
% mdparagraph:P362:begin
\textbf{V.18.6.}
% mdparagraph:P362:end

% item: EX-V-18-06
% mdparagraph:P363:begin
The time draws near [approaches, 164] (om. the) birth of Christ.
% mdparagraph:P363:end

\end{minipage}\par\addvspace{4pt}

\par\begin{minipage}{\linewidth}\raggedright
% mdparagraph:P364:begin
\textbf{V.18.7.}
% mdparagraph:P364:end

% teaching: NEW-U11-GRAM-u11-birth-passive-supplied
% help: U11-COOK-HELP-18-07
% mdparagraph:P365:begin
\textbf{Local help.} Cook supplies \textbf{ācenned}, “born”, a past participle, with plain e. Use it with the past passive auxiliary for “was born”; the other forms of this verb are not required.
% mdparagraph:P365:end

% item: EX-V-18-07
% mdparagraph:P366:begin
It is the day when (ðe) he was born (ācenned).
% mdparagraph:P366:end

\end{minipage}\par\addvspace{4pt}

\par\begin{minipage}{\linewidth}\raggedright
% mdparagraph:P367:begin
\textbf{V.18.8.}
% mdparagraph:P367:end

% item: EX-V-18-08
% mdparagraph:P368:begin
O, wast thou with me, dearest, then, while I rose up against my doom (dōm)?
% mdparagraph:P368:end

\end{minipage}\par\addvspace{4pt}

\par\begin{minipage}{\linewidth}\raggedright
% mdparagraph:P369:begin
\textbf{V.18.9.}
% mdparagraph:P369:end

% item: EX-V-18-09
% mdparagraph:P370:begin
The moon is hid (hȳdan, 113), the night is still.
% mdparagraph:P370:end

\end{minipage}\par\addvspace{4pt}

\par\begin{minipage}{\linewidth}\raggedright
% mdparagraph:P371:begin
\textbf{V.18.10.}
% mdparagraph:P371:end

% item: EX-V-18-10
% mdparagraph:P372:begin
Thou bringst the fisher to his wife.
% mdparagraph:P372:end

\end{minipage}\par\addvspace{4pt}

\par\begin{minipage}{\linewidth}\raggedright
% mdparagraph:P373:begin
\textbf{V.18.11.}
% mdparagraph:P373:end

% item: EX-V-18-11
% mdparagraph:P374:begin
Calm is the morn(ing).
% mdparagraph:P374:end

\end{minipage}\par\addvspace{4pt}

\par\begin{minipage}{\linewidth}\raggedright
% mdparagraph:P375:begin
\textbf{V.18.12.}
% mdparagraph:P375:end

% item: EX-V-18-12
% mdparagraph:P376:begin
I know that this was life.
% mdparagraph:P376:end

\end{minipage}\par\addvspace{4pt}

\par\begin{minipage}{\linewidth}\raggedright
% mdparagraph:P377:begin
\textbf{V.18.13.}
% mdparagraph:P377:end

% item: EX-V-18-13
% mdparagraph:P378:begin
What do ye more than they?
% mdparagraph:P378:end

\end{minipage}\par\addvspace{4pt}

\reserve{7\baselineskip}
% mdparagraph:P379:begin
\section*{12. The fisherman and his son}\BookHeading{nav-ch11-learner-0025}{subsection}{12. The fisherman and his son}{1}
% mdparagraph:P379:end

% mdparagraph:P380:begin
This connected account is new teaching material. Read it as a sequence: arrival, assistance, speech and the action that follows. Resolve pronouns from both their forms and the context.
% mdparagraph:P380:end

% composition: NEW-U11-READING-01
% mdparagraph:P381:begin
\par\begin{minipage}{\linewidth}\raggedright
\readingtext{Se ealda fiscere cōm tō ðǣm hūse, ðæt hē his suna hulpe. Ðā hē cōm, ðā wæs se sunu hungrig. Se fiscere sealde him męte. Mid-ðȳ se sunu ǣt, se fiscere bær wæter. Ðēah se fiscere wǣre eald, hē healp his suna. Se sunu cwæð: “Wes hēr, oð-ðæt ic cume.” Hē gewāt and cōm and sealde his fæder wæter.}
\end{minipage}\par
% mdparagraph:P381:end

\reserve{5\baselineskip}
\par\begin{minipage}{\linewidth}\raggedright
% mdparagraph:P382:begin
\textbf{R01.}
% mdparagraph:P382:end

% item: NEW-U11-R01
% mdparagraph:P383:begin
Translate the whole new account The fisherman and his son.
% mdparagraph:P383:end

\end{minipage}\par\addvspace{4pt}

\par\begin{minipage}{\linewidth}\raggedright
% mdparagraph:P384:begin
\textbf{R02.}
% mdparagraph:P384:end

% item: NEW-U11-R02
% mdparagraph:P385:begin
In the new account, identify the purpose clause, the while clause and the concession. State what each contributes to the account and parse its finite verb.
% mdparagraph:P385:end

\end{minipage}\par\addvspace{4pt}

\par\begin{minipage}{\linewidth}\raggedright
% mdparagraph:P386:begin
\textbf{R03.}
% mdparagraph:P386:end

% item: NEW-U11-R03
% mdparagraph:P387:begin
In the new account, who gives the first food, who receives it, and who gives the final water? Explain the relevant pronouns and genitive possessors.
% mdparagraph:P387:end

\end{minipage}\par\addvspace{4pt}

\par\begin{minipage}{\linewidth}\raggedright
% mdparagraph:P388:begin
\textbf{R04.}
% mdparagraph:P388:end

% item: NEW-U11-R04
% mdparagraph:P389:begin
In the son’s direct speech, who is ic, when is the coming expected, and what mood is cume? Identify the later action that fulfills the expectation.
% mdparagraph:P389:end

\end{minipage}\par\addvspace{4pt}

\reserve{7\baselineskip}
% mdparagraph:P390:begin
\section*{13. Bede’s description of Britain: metals and jet}\BookHeading{nav-ch11-learner-0026}{subsection}{13. Bede’s description of Britain: metals and jet}{1}
% mdparagraph:P390:end

% mdparagraph:P391:begin
This is an exact extract from Cook, \emph{A First Book in Old English}, printed p.139 /\allowbreak{} PDF p.155, editorial lines1–4, ending at onwęg. Cook’s attribution is “(Eccl. Hist., Bk. I., Chap. I.)”. The passage has 35 words; the dash is punctuation. The line count is editorial, anchored by the printed marginal line5. Original notes are preserved in §15.
% mdparagraph:P391:end

% mdparagraph:P392:begin
\textbf{Context.} The subject is Britain. Bede’s preceding description has discussed the island’s plants, animals, shellfish and springs.
% mdparagraph:P392:end

\par\begin{minipage}{\linewidth}\raggedright
% source: FB-IV-BRITAIN-METALS-01
% mdparagraph:P393:begin
\readingtext{Swilce hit is ēac berende on węcga ōrum — āres and īsernes, lēades and seolfres. Hēr bið ēac gemētt gagates; se stān bið blæc gimm; gif man hine on fȳr dēð, ðonne flēoð ðǣr næddran onwęg.}
% mdparagraph:P393:end

\end{minipage}\par
\reserve{7\baselineskip}
% mdparagraph:P394:begin
\subsection*{Supplied help}\BookHeading{nav-ch11-learner-0027}{subsubsection}{Supplied help}{0}
% mdparagraph:P394:end

% mdparagraph:P395:begin
The following words and constructions may be used with their help throughout the reading tasks. The whole expressions are supported here; they need not be reconstructed from earlier headwords.
% mdparagraph:P395:end

% teaching: NEW-U11-GRAM-NEW-U11-SRC-B01
% help: U11-SRC-HELP-B01
% mdparagraph:P396:begin
\textbf{1. Swilce hit ... seolfres.} The preceding context concerns Britain, the island: hit refers to it. Swilce ... ēac, also. Berende on, rich in /\allowbreak{} productive of; węcga is genitive plural “of metals,” and ōrum dative plural “veins/\allowbreak{}ores” after on. The following genitives name the metals. Āres is Cook’s source spelling for copper here, beside the earlier selected ǣr.
% mdparagraph:P396:end

% teaching: NEW-U11-GRAM-NEW-U11-SRC-B02
% help: U11-SRC-HELP-B02
% mdparagraph:P397:begin
\textbf{2. Hēr bið ... blæc gimm.} Gemētt, found, is the exact past participle of prefixed gemētan; bið gemētt, is found. Gagates is jet, the stone’s Latin loan-name; stān, stone, masculine nominative singular. The second bið links stone with black gem.
% mdparagraph:P397:end

% teaching: NEW-U11-GRAM-NEW-U11-SRC-B03
% help: U11-SRC-HELP-B03
% mdparagraph:P398:begin
\textbf{3. gif man hine on fȳr dēð, ðonne.} Man is the impersonal singular subject “one”; dēð is third singular present indicative. Hine refers to the masculine stone. The newly taught placement frame dōn + object + on with accusative gives “puts it on/\allowbreak{}in a fire”; ðonne correlates with the gif condition.
% mdparagraph:P398:end

% teaching: NEW-U11-GRAM-NEW-U11-SRC-B04
% help: U11-SRC-HELP-B04
% mdparagraph:P399:begin
\textbf{4. flēoð ðǣr næddran onwęg.} Flēoð is third plural present indicative of flēon, flee; it is not the verb fly. Ðǣr, there; næddran, snakes/\allowbreak{}adders, nominative plural of weak feminine næddre; onwęg, away. Both d letters in næddran belong to the printed text.
% mdparagraph:P399:end

\reserve{5\baselineskip}
\par\begin{minipage}{\linewidth}\raggedright
% mdparagraph:P400:begin
\textbf{B01.}
% mdparagraph:P400:end

% item: NEW-U11-B01
% mdparagraph:P401:begin
With the printed helps, translate the complete Bede metals-and-jet passage.
% mdparagraph:P401:end

\end{minipage}\par\addvspace{4pt}

\par\begin{minipage}{\linewidth}\raggedright
% mdparagraph:P402:begin
\textbf{B02.}
% mdparagraph:P402:end

% item: NEW-U11-B02
% mdparagraph:P403:begin
In the Bede passage, identify the condition introduced by gif and its main clause. Parse the two finite verbs using the supplied forms, and explain the role of ðonne.
% mdparagraph:P403:end

\end{minipage}\par\addvspace{4pt}

\par\begin{minipage}{\linewidth}\raggedright
% mdparagraph:P404:begin
\textbf{B03.}
% mdparagraph:P404:end

% item: NEW-U11-B03
% mdparagraph:P405:begin
What does the first hit refer to, and what does hine in the condition refer to? Give the grammatical information that distinguishes them, using the introduction and helps.
% mdparagraph:P405:end

\end{minipage}\par\addvspace{4pt}

\reserve{7\baselineskip}
% mdparagraph:P406:begin
\section*{14. Æthelwald calms the sea}\BookHeading{nav-ch11-learner-0028}{subsection}{14. Æthelwald calms the sea}{1}
% mdparagraph:P406:end

% mdparagraph:P407:begin
This exact connected episode is from Cook’s \emph{First Book}, printed pp.142–143 /\allowbreak{} PDF pp.158–159: editorial line8 on p.142 through line14 on p.143. It begins after standan. in the preceding account and ends with heofonlice forgiefen wæs. The 195 words retain the original two paragraphs. Cook’s original notes, including the cross-reference to an earlier note, follow in §15.
% mdparagraph:P407:end

% mdparagraph:P408:begin
\textbf{Earlier events.} Three visitors are returning by boat from the hermit Æthelwald on Farne. A sudden storm has stopped their progress. After struggling in vain, they look back toward the island. The narrator now tells how they were rescued.
% mdparagraph:P408:end

\par\begin{minipage}{\linewidth}\raggedright
% source: FB-V-SEA-RESCUE-01
% mdparagraph:P409:begin
\readingtext{Ðā wæs æfter langum fæce ðæt wē ūre gesihð feorr ūpāhōfon, ðā gesāwon wē in Farne, ðǣm īglande, Gode ðone lēofan fæder Æðelwald of his dīegelnessum ūtgangende, ðæt hē wolde ūrne sīðfæt scēawian, and gesēon hwæt ūs gelumpe, for-ðon hē gehīerde ðæt gebrec ðǣra storma and ðæs weallendes sǣs.}
% mdparagraph:P409:end

% mdparagraph:P410:begin
\readingtext{Mid-ðȳ hē ðā ūs ēac scēawode, and geseah in gewinne and in ōrmōdnesse gesętte bēon, ðā biegde hē his cnēowu tō Fæder ūres Dryhtnes Hǣlendes Crīstes, and wæs gebiddende for ūre hǣle and for ūrum līfe. And mid-ðȳ hē ðā ðæt gebed gefylde, hē ðā samod ætgædere ge ðone āðundnan sǣ gesmylte ge ðone storm gestilde, tō ðon ðætte ðurh eall sēo rēðnes ðæs stormes wæs blinnende, and gesyndge windas ðurh ðone smyltestan sǣ ūs æt lande gebrōhton. Mid-ðȳ wē ðā ūp cōmon tō lande, and ūre scip ēac swilce fram ðǣm ȳðum ūp ābǣron, ðā sōna se ilca storm ęft hwearf and cōm, se-ðe for ūrum intingan medmicel fæc gestilde, and ealne ðone dæg swīðe micel and strang wæs, ðætte męnn sweotollice ongietan meahton ðætte se medmicla fierst ðǣre stilnesse, ðe ðǣr becōm, tō bēnum ðæs Godes weres for intingan ūre hǣle heofonlice forgiefen wæs.}
% mdparagraph:P410:end

\end{minipage}\par
\reserve{7\baselineskip}
% mdparagraph:P411:begin
\subsection*{Supplied help}\BookHeading{nav-ch11-learner-0029}{subsubsection}{Supplied help}{0}
% mdparagraph:P411:end

% mdparagraph:P412:begin
Follow the shifts from what the travellers see to Æthelwald’s actions and back to the travellers. The complete expressions below are supplied. The sea remains masculine in these phrases, as in the earlier paradigm.
% mdparagraph:P412:end

% teaching: NEW-U11-GRAM-NEW-U11-SRC-S01
% help: U11-SRC-HELP-S01
% mdparagraph:P413:begin
\textbf{1. Ðā wæs æfter langum fæce ðæt.} Read the whole opening as “then, after a long time, it happened that ...”. Æfter governs the dative; langum fæce means a long interval. The opening ðæt belongs to this supplied time-introduction frame.
% mdparagraph:P413:end

% teaching: NEW-U11-GRAM-NEW-U11-SRC-S02
% help: U11-SRC-HELP-S02
% mdparagraph:P414:begin
\textbf{2. wē ūre gesihð feorr ūpāhōfon.} Gesihð is feminine accusative singular, sight/\allowbreak{}gaze; ūre agrees with it. Feorr, far off. Ūpāhōfon, we lifted up, first plural preterite of ūpāhebban. The complete phrase describes looking up and seeing at a distance.
% mdparagraph:P414:end

% teaching: NEW-U11-GRAM-NEW-U11-SRC-S03
% help: U11-SRC-HELP-S03
% mdparagraph:P415:begin
\textbf{3. in Farne ... of his dīegelnessum ūtgangende.} In Farne, on Farne island, with ðǣm īglande as dative apposition. Gode ðone lēofan fæder means the father dear to God: the adjective’s complement is dative. Æðelwald is in apposition to the accusative fæder. Of + dative means out of/\allowbreak{}from; dīegelnessum is the source plural for retreat/\allowbreak{}seclusion. Ūtgangende is the printed participle, coming out; Cook notes that a present participle is used for a Latin past participle here, so “having come out” is also possible.
% mdparagraph:P415:end

% teaching: NEW-U11-GRAM-NEW-U11-SRC-S04
% help: U11-SRC-HELP-S04
% mdparagraph:P416:begin
\textbf{4. ðæt hē wolde ... hwæt ūs gelumpe.} The ðæt clause gives his purpose. Wolde with the infinitives scēawian and gesēon is read here “that he might watch and see” or “in order to watch and see”; do not make private desire the only force. Hwæt ūs gelumpe is a dependent question, what might happen to us. Gelumpe is third singular preterite subjunctive, and ūs is dative.
% mdparagraph:P416:end

% teaching: NEW-U11-GRAM-NEW-U11-SRC-S05
% help: U11-SRC-HELP-S05
% mdparagraph:P417:begin
\textbf{5. for-ðon hē gehīerde ... weallendes sǣs.} For-ðon is the shorter source spelling of the causal conjunction. Gehīerde, heard, is a prefixed past form. Gebrec is the accusative noise/\allowbreak{}crash; ðǣra storma is genitive plural of storms, and ðæs weallendes sǣs genitive singular of the surging sea. Weallendes is an adjectival participle, surging/\allowbreak{}boiling. Ðǣra preserves the source genitive-plural spelling.
% mdparagraph:P417:end

% teaching: NEW-U11-GRAM-NEW-U11-SRC-S06
% help: U11-SRC-HELP-S06
% mdparagraph:P418:begin
\textbf{6. Mid-ðȳ ... in gewinne and in ōrmōdnesse gesętte bēon.} The temporal clause contains both scēawode and geseah. With geseah, understand the earlier ūs as the people seen: he saw us to be situated in distress and despair. In governs the two datives gewinne and ōrmōdnesse. Gesętte is plural past participle here, not the identically spelled finite past; bēon is infinitive. This entire perception complement is supplied.
% mdparagraph:P418:end

% teaching: NEW-U11-GRAM-NEW-U11-SRC-S07
% help: U11-SRC-HELP-S07
% mdparagraph:P419:begin
\textbf{7. ðā biegde hē his cnēowu tō Fæder ... Crīstes.} Biegde, bent, is the exact third singular preterite spelling; cnēowu, knees, accusative plural. Tō Fæder marks the direction of the prayer. Ūres Dryhtnes Hǣlendes Crīstes is a genitive group: of our Lord Jesus Christ, with the proper title Hǣlend and Crīst in apposition.
% mdparagraph:P419:end

% teaching: NEW-U11-GRAM-NEW-U11-SRC-S08
% help: U11-SRC-HELP-S08
% mdparagraph:P420:begin
\textbf{8. wæs gebiddende for ūre hǣle and for ūrum līfe.} Cook143 note1 refers to142 note7: this was + participle wording is read as a simple past here, “prayed.” Do not infer a universal English progressive rule. For ūre hǣle, for our safety/\allowbreak{}rescue, uses accusative; for ūrum līfe, for our life, uses dative. Keep both exact constructions supplied.
% mdparagraph:P420:end

% teaching: NEW-U11-GRAM-NEW-U11-SRC-S09
% help: U11-SRC-HELP-S09
% mdparagraph:P421:begin
\textbf{9. mid-ðȳ ... ðæt gebed gefylde ... samod ætgædere ge.} Gebed is prayer, accusative singular; gefylde is the already taught finished/\allowbreak{}accomplished. Samod ætgædere means together/\allowbreak{}at once. Paired ge ... ge coordinates both ... and; the first ge introduces the sea-calming and the second the storm-stilling. The narrated hē remains Æthelwald, the person who has just completed the prayer.
% mdparagraph:P421:end

% teaching: NEW-U11-GRAM-NEW-U11-SRC-S10
% help: U11-SRC-HELP-S10
% mdparagraph:P422:begin
\textbf{10. ðone āðundnan sǣ gesmylte ge ðone storm gestilde.} Āðundnan, swollen, is a weak participial adjective agreeing with masculine accusative sǣ. Gesmylte, calmed, is third singular preterite; storm is accusative singular. Gestilde retains the taught transitive quell/\allowbreak{}settle sense here.
% mdparagraph:P422:end

% teaching: NEW-U11-GRAM-NEW-U11-SRC-S11
% help: U11-SRC-HELP-S11
% mdparagraph:P423:begin
\textbf{11. tō ðon ðætte ðurh eall sēo rēðnes ... wæs blinnende.} Tō ðon ðætte introduces the result, so that. Ðurh eall is the exact adverbial expression completely/\allowbreak{}entirely. Rēðnes, fury, is the feminine subject; ðæs stormes is genitive. Wæs blinnende is the supplied periphrasis “ceased” here, with present participle blinnende; no new productive progressive tense is introduced.
% mdparagraph:P423:end

% teaching: NEW-U11-GRAM-NEW-U11-SRC-S12
% help: U11-SRC-HELP-S12
% mdparagraph:P424:begin
\textbf{12. gesyndge windas ... æt lande gebrōhton.} Gesyndge is the exact syncopated source form beside taught gesyndige, favorable, nominative plural. Smyltestan is the weak masculine accusative superlative calmest/\allowbreak{}very calm. Ðurh governs ðone ... sǣ; ūs is the accusative people brought. Æt lande means to/\allowbreak{}at land in this exact arrival expression. Gebrōhton is the prefixed third plural preterite brought.
% mdparagraph:P424:end

% teaching: NEW-U11-GRAM-NEW-U11-SRC-S13
% help: U11-SRC-HELP-S13
% mdparagraph:P425:begin
\textbf{13. Mid-ðȳ wē ... fram ðǣm ȳðum ūp ābǣron.} The temporal clause has the two completed actions: we came up onto land and carried our ship up from the waves. Lande is the locally supplied dative of land; ȳðum is dative plural waves after fram. Ābǣron is first plural preterite of prefixed carry. Ūre scip is the accusative object, not a new subject.
% mdparagraph:P425:end

% teaching: NEW-U11-GRAM-NEW-U11-SRC-S14
% help: U11-SRC-HELP-S14
% mdparagraph:P426:begin
\textbf{14. ðā sōna ... se-ðe for ūrum intingan medmicel fæc gestilde.} Sōna, immediately; ęft, again; hwearf, returned. Se-ðe is the exact relative complex referring to masculine storm. For ūrum intingan, for our sake; intingan is dative singular. Medmicel fæc is accusative duration, a short interval. Gestilde is intransitive ceased/\allowbreak{}became still in this relative clause; the earlier transitive sense alone is insufficient. English had ceased may clarify the sequence.
% mdparagraph:P426:end

% teaching: NEW-U11-GRAM-NEW-U11-SRC-S15
% help: U11-SRC-HELP-S15
% mdparagraph:P427:begin
\textbf{15. and ealne ðone dæg ... ðætte męnn sweotollice ongietan meahton ðætte.} Ealne ðone dæg is accusative duration, all that day. Strang, strong, agrees with storm. The first ðætte introduces the result: people could clearly understand. Sweotollice, clearly; ongietan, understand, is the infinitive governed by meahton. The second ðætte introduces what they understand, a content clause.
% mdparagraph:P427:end

% teaching: NEW-U11-GRAM-NEW-U11-SRC-S16
% help: U11-SRC-HELP-S16
% mdparagraph:P428:begin
\textbf{16. se medmicla fierst ðǣre stilnesse, ðe ðǣr becōm.} Fierst is masculine nominative singular, interval/\allowbreak{}period; medmicla is its weak adjective, short. Ðǣre stilnesse, of the calm, is genitive singular. In the relative clause, ðǣr means there and becōm means occurred/\allowbreak{}came about in this context. Relative ðe is the subject of becōm and refers to the interval; it is distinct from the following adverb ðǣr.
% mdparagraph:P428:end

% teaching: NEW-U11-GRAM-NEW-U11-SRC-S17
% help: U11-SRC-HELP-S17
% mdparagraph:P429:begin
\textbf{17. tō bēnum ðæs Godes weres ... heofonlice forgiefen wæs.} Tō bēnum, at/\allowbreak{}in response to the prayers, uses dative plural. Ðæs Godes weres means of the man of God; weres is genitive singular, with Godes depending on it. For intingan ūre hǣle means for the sake of our rescue: hǣle is genitive singular; exact ūre is supplied in this possessive expression. Heofonlice means from heaven; forgiefen means given/\allowbreak{}granted here, expressly not forgiven. The final passive refers to the brief period of calm. Here ūre is the first-person plural genitive used attributively within “our rescue”; this supplied phrase does not establish a new feminine possessive-adjective variant.
% mdparagraph:P429:end

\reserve{5\baselineskip}
\par\begin{minipage}{\linewidth}\raggedright
% mdparagraph:P430:begin
\textbf{S01.}
% mdparagraph:P430:end

% item: NEW-U11-S01
% mdparagraph:P431:begin
With the printed helps, translate the complete rescue-and-resolution passage from Æthelwald Calms the Sea.
% mdparagraph:P431:end

\end{minipage}\par\addvspace{4pt}

\par\begin{minipage}{\linewidth}\raggedright
% mdparagraph:P432:begin
\textbf{S02.}
% mdparagraph:P432:end

% item: NEW-U11-S02
% mdparagraph:P433:begin
Put these events in the order given by the passage: the storm returns; the sailors see Æthelwald; the sailors reach land; Æthelwald prays; the sea becomes calm. State what the sailors have done before the storm returns.
% mdparagraph:P433:end

\end{minipage}\par\addvspace{4pt}

\par\begin{minipage}{\linewidth}\raggedright
% mdparagraph:P434:begin
\textbf{S03.}
% mdparagraph:P434:end

% item: NEW-U11-S03
% mdparagraph:P435:begin
Find the clause that gives Æthelwald’s purpose in coming out and the clause that gives his reason. Explain their different functions; use the local help for wolde and its infinitives.
% mdparagraph:P435:end

\end{minipage}\par\addvspace{4pt}

\par\begin{minipage}{\linewidth}\raggedright
% mdparagraph:P436:begin
\textbf{S04.}
% mdparagraph:P436:end

% item: NEW-U11-S04
% mdparagraph:P437:begin
Find the result introduced by tō ðon ðætte. What happened to the storm, and what brought the sailors to land? Use the help to explain wæs blinnende.
% mdparagraph:P437:end

\end{minipage}\par\addvspace{4pt}

\par\begin{minipage}{\linewidth}\raggedright
% mdparagraph:P438:begin
\textbf{S05.}
% mdparagraph:P438:end

% item: NEW-U11-S05
% mdparagraph:P439:begin
Identify the two mid-ðȳ clauses that introduce Æthelwald’s sight of the sailors and the sailors’ coming ashore. Where does the main action begin in each case?
% mdparagraph:P439:end

\end{minipage}\par\addvspace{4pt}

\par\begin{minipage}{\linewidth}\raggedright
% mdparagraph:P440:begin
\textbf{S06.}
% mdparagraph:P440:end

% item: NEW-U11-S06
% mdparagraph:P441:begin
Explain how wæs gebiddende is to be understood here, using Cook’s note and the help. Why does the present participle here not require an English progressive translation?
% mdparagraph:P441:end

\end{minipage}\par\addvspace{4pt}

\par\begin{minipage}{\linewidth}\raggedright
% mdparagraph:P442:begin
\textbf{S07.}
% mdparagraph:P442:end

% item: NEW-U11-S07
% mdparagraph:P443:begin
At the end, what do people understand about the short period of calm? Say whose prayers and whose rescue are named, and explain the final forgiefen using the supplied note.
% mdparagraph:P443:end

\end{minipage}\par\addvspace{4pt}

\par\begin{minipage}{\linewidth}\raggedright
% mdparagraph:P444:begin
\textbf{S08.}
% mdparagraph:P444:end

% item: NEW-U11-S08
% mdparagraph:P445:begin
Does the passage say that the storm ceased permanently? Explain the return of the storm and how it supports the narrator’s final explanation.
% mdparagraph:P445:end

\end{minipage}\par\addvspace{4pt}

\reserve{8\baselineskip}
% mdparagraph:P446:begin
\section*{15. Cook’s original notes and source information}\BookHeading{nav-ch11-learner-0030}{subsection}{15. Cook’s original notes and source information}{1}
% mdparagraph:P446:end

\begingroup\fontsize{10}{12.8}\selectfont
\setlength{\parskip}{1.5pt}
% mdparagraph:P447:begin
These are Cook’s endnotes, kept separate from the Old English and the new learner help. Numbers belong to their printed pages; they restart on each page. His historical and etymological questions remain source apparatus and are not additional exercises. The anchors identify the words bearing the original cues.
% mdparagraph:P447:end

\reserve{7\baselineskip}
% mdparagraph:P448:begin
\subsection*{Printed page 139}\BookHeading{nav-ch11-learner-0031}{subsubsection}{Printed page 139}{0}
% mdparagraph:P448:end

% source-note: U11-SRC-P139-N01
% mdparagraph:P449:begin
\textbf{1. berende.} Cf. the construction of this word with that above, p. 138, l. 1.
% mdparagraph:P449:end

% source-note: U11-SRC-P139-N02
% mdparagraph:P450:begin
\textbf{2. man.} See 89. e.
% mdparagraph:P450:end

% source-note: U11-SRC-P139-N03
% mdparagraph:P451:begin
\textbf{3. dēð.} See 140.
% mdparagraph:P451:end

% source-note: U11-SRC-P139-N04
% mdparagraph:P452:begin
\textbf{4. næddran.} How did næddre become adder? Cf. OE. nafogār, Mod. Eng. auger. See Skeat, Prin., p. 216.
% mdparagraph:P452:end

% source-note: U11-SRC-P139-N05
% mdparagraph:P453:begin
\textbf{5. onwęg.} There is a parallel form, āwęg, already in OE. The ā- is a contraction of on. Mention other Mod. Eng. words in which the a- represents on.
% mdparagraph:P453:end

\reserve{7\baselineskip}
% mdparagraph:P454:begin
\subsection*{Printed page 142}\BookHeading{nav-ch11-learner-0032}{subsubsection}{Printed page 142}{0}
% mdparagraph:P454:end

% source-note: U11-SRC-P142-N14
% mdparagraph:P455:begin
\textbf{14. feorr.} Translate, from a distance.
% mdparagraph:P455:end

% source-note: U11-SRC-P142-N15
% mdparagraph:P456:begin
\textbf{15. Gode.} Governed by lēofan; = amantissimum Deo. See 165.
% mdparagraph:P456:end

% source-note: U11-SRC-P142-N16
% mdparagraph:P457:begin
\textbf{16. dīegelnessum.} Lat. latibulis.
% mdparagraph:P457:end

% source-note: U11-SRC-P142-N17
% mdparagraph:P458:begin
\textbf{17. ūtgangende.} Translates the Latin past part., egressum.
% mdparagraph:P458:end

% source-note: U11-SRC-P142-N18
% mdparagraph:P459:begin
\textbf{18. ðæt hē wolde.} Translate, that he might, in order to, to. The Latin has the infinitive.
% mdparagraph:P459:end

% source-note: U11-SRC-P142-N19
% mdparagraph:P460:begin
\textbf{19. gelumpe.} See 194. b.
% mdparagraph:P460:end

% source-note: U11-SRC-P142-N20
% mdparagraph:P461:begin
\textbf{20. sǣs.} Lat. fragore procellarum ac ferventis oceani.
% mdparagraph:P461:end

% source-note: U11-SRC-P142-N21
% mdparagraph:P462:begin
\textbf{21. ōrmōdnesse.} Lat. desperatione.
% mdparagraph:P462:end

% source-note: U11-SRC-P142-N22
% mdparagraph:P463:begin
\textbf{22. bēon.} Supply ūs as subject acc. (169).
% mdparagraph:P463:end

\reserve{7\baselineskip}
% mdparagraph:P464:begin
\subsection*{Printed page 143}\BookHeading{nav-ch11-learner-0033}{subsubsection}{Printed page 143}{0}
% mdparagraph:P464:end

% source-note: U11-SRC-P143-N01
% mdparagraph:P465:begin
\textbf{1. wæs gebiddende; wæs blinnende.} See above, p. 142, n. 7.
% mdparagraph:P465:end

% source-note: U11-SRC-P143-N02
% mdparagraph:P466:begin
\textbf{2. gefylde.} Lat. compleret.
% mdparagraph:P466:end

% source-note: U11-SRC-P143-N03
% mdparagraph:P467:begin
\textbf{3. āðundnan.} Lat. tumida.
% mdparagraph:P467:end

% source-note: U11-SRC-P143-N04
% mdparagraph:P468:begin
\textbf{4. tō ðon ðætte.} Lat. adeo ut, nearly = so that.
% mdparagraph:P468:end

% source-note: U11-SRC-P143-N05
% mdparagraph:P469:begin
\textbf{5. ðurh eall.} Lat. per omnia, = entirely.
% mdparagraph:P469:end

% source-note: U11-SRC-P143-N06
% mdparagraph:P470:begin
\textbf{6. gesyndge.} Lat. secunda. What letter (sound) has been lost from the OE. form?
% mdparagraph:P470:end

% source-note: U11-SRC-P143-N07
% mdparagraph:P471:begin
\textbf{7. ēac swilce.} Lat. quoque.
% mdparagraph:P471:end

% source-note: U11-SRC-P143-N08
% mdparagraph:P472:begin
\textbf{8. cōm.} Supply ęft. Latin has only one verb, rediit.
% mdparagraph:P472:end

% source-note: U11-SRC-P143-N09
% mdparagraph:P473:begin
\textbf{9. for ūrum intingan.} Lat. nostri gratia.
% mdparagraph:P473:end

% source-note: U11-SRC-P143-N10
% mdparagraph:P474:begin
\textbf{10. fæc; dæg.} See 170.
% mdparagraph:P474:end

% source-note: U11-SRC-P143-N11
% mdparagraph:P475:begin
\textbf{11. ðone.} Lat. illo. Translate, that.
% mdparagraph:P475:end

% source-note: U11-SRC-P143-N12
% mdparagraph:P476:begin
\textbf{12. ðætte.} From ðæt-ðe (34).
% mdparagraph:P476:end

% source-note: U11-SRC-P143-N13
% mdparagraph:P477:begin
\textbf{13. tō bēnum.} Lat. ad preces. Cf. the Mod. Eng. phrase, ‘bootless bene.’
% mdparagraph:P477:end

% source-note: U11-SRC-P143-N14
% mdparagraph:P478:begin
\textbf{14. Godes.} Dependent on weres.
% mdparagraph:P478:end

% source-note: U11-SRC-P143-N15
% mdparagraph:P479:begin
\textbf{15. weres.} Dependent on bēnum.
% mdparagraph:P479:end

% source-note: U11-SRC-P143-N16
% mdparagraph:P480:begin
\textbf{16. hǣle.} Lat. evasionis.
% mdparagraph:P480:end

% source-note: U11-SRC-P143-N17
% mdparagraph:P481:begin
\textbf{17. heofonlice.} Lat. cælitus, = from heaven.
% mdparagraph:P481:end

% source-note: U11-SRC-P143-N18
% mdparagraph:P482:begin
\textbf{18. forgiefen.} Not forgiven, but given.
% mdparagraph:P482:end

\reserve{7\baselineskip}
% mdparagraph:P483:begin
\subsection*{Referenced note: printed page 142, note 7}\BookHeading{nav-ch11-learner-0034}{subsubsection}{Referenced note: printed page 142, note 7}{0}
% mdparagraph:P483:end

% mdparagraph:P484:begin
Cook p.143 note1 explicitly refers to this earlier explanation. The preceding sentence to which it was first attached is not part of the assigned extract.
% mdparagraph:P484:end

% source-note: U11-SRC-P142-N07
% mdparagraph:P485:begin
Lat. egressi eramus. The pres. part. with the verb is sometimes used in OE. to denote the simple past, as here, and not the progressive.
% mdparagraph:P485:end

\reserve{7\baselineskip}
% mdparagraph:P486:begin
\subsection*{Editions, spelling and grammatical references}\BookHeading{nav-ch11-learner-0035}{subsubsection}{Editions, spelling and grammatical references}{0}
% mdparagraph:P486:end

% mdparagraph:P487:begin
Cook’s \emph{Exercises in Old English} (1895), printed pp.8–9 /\allowbreak{} one-based PDF pp.16–17, supplies all sixteen IV.16 items, all seventeen IV.17 items and all thirteen V.18 items. The original references (64), (154. c), (164) and (113), the omission directions, and the hints dōm and ācenned are retained. The Exercises page offset is +8; for Cook’s \emph{A First Book in Old English}, second edition (1895), the offset is +16. Printed pages and editorial line counts are distinct from PDF pages.
% mdparagraph:P487:end

% mdparagraph:P488:begin
In the exact readings, printed line wraps and line-end divisions are joined and note cues are placed in the endnotes; lexical hyphens, letters, quantities and paragraph divisions remain. Retain āres, næddran, ðǣra, biegde, gesyndge and se-ðe as printed. The first passage describes a belief about jet; translating the statement does not turn it into a modern scientific claim.
% mdparagraph:P488:end

% mdparagraph:P489:begin
The noun patterns and their exceptions follow Cook §§43–54, pp.28–38; adjective patterns §§55–64, pp.38–43. The reduction in comparative genitive plurals continues the earlier rule; compare Wright, \emph{Old English Grammar} (1914), §259, p.130. Selected vocabulary entries establish the genders, stems and meanings. The proper-title construction follows Cook §151, p.88. Personal adjectives follow §§55 and181. The participle bases and indicative systems follow the relevant strong and weak paradigms; complete tables here combine those patterns with the individual lexical stems.
% mdparagraph:P489:end

% mdparagraph:P490:begin
For the clause relations compare Cook §§194–196, pp.103–105, with Fulk, \emph{An Introductory Grammar of Old English} (2020), §45. Constructions with infinitives and perception follow Cook §§169 and198–200; the exact modal, participial and perception expressions in the sea passage remain locally supplied. Cook’s p.143 note18 explicitly requires “give”, not “forgive”, for the final participle.
% mdparagraph:P490:end

\endgroup

% book-source-body:end
\par\label{end-ch11-learner}
\clearpage\endgroup

\BookRecto
\begingroup
\renewcommand{\runningtitle}{Unit 12. Tense, sequence, and voice}
\hyphenpenalty=10000
\exhyphenpenalty=10000
\emergencystretch=3em
\newcommand{\itempara}[2]{\par\addvspace{3pt}\noindent\hangindent=1.8em\hangafter=1\makebox[1.8em][l]{#1.}#2\par}
\newcommand{\reserve}[1]{\par\penalty-100\begingroup\dimen0=\pagegoal\advance\dimen0 by-\pagetotal\ifdim\dimen0<#1\newpage\fi\endgroup}
% book-source-body:begin

\addfontfeatures{Ligatures=TeXOff}
\raggedright
\booktitle{\BookComponent{ch12-learner}{Unit 12. Tense, sequence, and voice}UNIT 12}{Tense, sequence, and voice}

% mdparagraph:P001:begin
Old English has two inflectional tenses, present and preterite. Their interpretation depends on the situation. A present may refer to a future event; a preterite may narrate an event, describe a past state, or refer to something already completed at a past viewpoint. English may consequently need a simple, perfect, pluperfect or future expression where Old English has one finite word.
% mdparagraph:P001:end

% mdparagraph:P002:begin
The twelve translated examples, exercises F, A, C, R and S, and “A house and a return” are new teaching material. Cook VI.19–21 retain all their original English prompts. The final historical reading is an exact connected extract, with separate English context, learner help and Cook’s original notes.
% mdparagraph:P002:end

% mdparagraph:P003:begin
In tables, 1, 2 and 3 denote person. Nominative, genitive, dative and accusative are written out. Repeated spellings occupy distinct grammatical positions. \textbf{Reference} means recognize the stated form or construction; it is not required in independent composition. \textbf{Supplied help} remains available while answering the associated task. Preserve written vowel lengths and the distinction between e and ę.
% mdparagraph:P003:end

\reserve{7\baselineskip}
% mdparagraph:P004:begin
\section*{1. Time in a narrative}\BookHeading{nav-ch12-learner-0001}{subsection}{1. Time in a narrative}{1}
% mdparagraph:P004:end

% teaching: NEW-U12-GRAM-u12-narrative-time
% teaching: NEW-U12-GRAM-u12-preterite-anterior
% mdparagraph:P005:begin
A\BookPoint{bk-ch12-learner-24ebe58b96d5} succession of preterites commonly advances the narrative. Temporal clauses, the meanings of the verbs and the surrounding account establish which event is earlier. A clause does not acquire a different inflection simply because English translates it with “had”.
% mdparagraph:P005:end

\reserve{8\baselineskip}
% example: NEW-U12-EXAMPLE01
% mdparagraph:P006:begin
Se mann cōm, and hē geseah ðæt hūs.
% mdparagraph:P006:end

% mdparagraph:P007:begin
The man came, and he saw the house.
% mdparagraph:P007:end

% mdparagraph:P008:begin
Cōm and geseah are third singular preterite indicatives. In this account the arrival precedes the seeing.
% mdparagraph:P008:end

\reserve{12\baselineskip}
% mdparagraph:P009:begin
Context: The house was finished before the man’s visit.
% mdparagraph:P009:end

\reserve{8\baselineskip}
% example: NEW-U12-EXAMPLE02
% mdparagraph:P010:begin
Se mann geseah ðæt hūs ðe his sunu geworhte.
% mdparagraph:P010:end

% mdparagraph:P011:begin
The man saw the house which his son had built.
% mdparagraph:P011:end

% mdparagraph:P012:begin
Geworhte\BookPoint{bk-ch12-learner-e075379e48d6} is also a simple preterite indicative. Here the stated context places the son’s building before the man’s seeing. English “had built” makes that relation clear without changing the Old English form.
% mdparagraph:P012:end

\reserve{12\baselineskip}
% mdparagraph:P013:begin
Context: The visit is still expected.
% mdparagraph:P013:end

\reserve{8\baselineskip}
% example: NEW-U12-EXAMPLE12
% mdparagraph:P014:begin
Hē cymð.
% mdparagraph:P014:end

% mdparagraph:P015:begin
He will come.
% mdparagraph:P015:end

% mdparagraph:P016:begin
Cymð is present indicative. The context makes the coming future. Do not add willan merely because the English translation contains “will”; willan can retain its meaning “wish; intend”.
% mdparagraph:P016:end

\reserve{7\baselineskip}
% mdparagraph:P017:begin
\section*{2. Habban with a past participle}\BookHeading{nav-ch12-learner-0002}{subsection}{2. Habban with a past participle}{1}
% mdparagraph:P017:end

% teaching: NEW-U12-GRAM-u12-habban-perfect
% teaching: NEW-U12-VOC-HABBAN
% mdparagraph:P018:begin
With\BookPoint{bk-ch12-learner-7f0ea6e7cad6} a transitive past participle and an accusative object, finite \textbf{habban} can present completed action and its result. The auxiliary agrees with the subject. The object is the person or thing affected by the participle. In the following construction the speaker has performed the building, and the completed house is the result.
% mdparagraph:P018:end

\reserve{8\baselineskip}
% example: NEW-U12-EXAMPLE03
% mdparagraph:P019:begin
Nū hæbbe ic ðæt hūs geworht.
% mdparagraph:P019:end

% mdparagraph:P020:begin
Now I have built the house.
% mdparagraph:P020:end

\reserve{12\baselineskip}
% mdparagraph:P021:begin
Context: The building was already complete when the son arrived.
% mdparagraph:P021:end

\reserve{8\baselineskip}
% example: NEW-U12-EXAMPLE04
% mdparagraph:P022:begin
Ðā se sunu cōm, hæfde se mann ðæt hūs geworht.
% mdparagraph:P022:end

% mdparagraph:P023:begin
When the son came, the man had built the house.
% mdparagraph:P023:end

% mdparagraph:P024:begin
Present\BookPoint{bk-ch12-learner-6f6cc62dab0b} hæbbe presents the completed work from the present viewpoint. Past hæfde places the completed work before the stated past arrival. The difference is in the finite auxiliary; geworht remains the selected participle base.
% mdparagraph:P024:end

% teaching: NEW-U12-HABBAN-PERFECT-IND
\reserve{11\baselineskip}
% mdparagraph:P025:begin
\textbf{Habban with an object and past participle: complete present and preterite indicative construction}
% mdparagraph:P025:end

% construction: NEW-U12-HABBAN-PERFECT-IND
% mdtable:T01:begin
% course-table-style: 2026-10
% course-panel:T01-P01:begin
\par\begingroup
\fontsize{10.5}{13.2}\selectfont
\setlength{\LTleft}{0pt}\setlength{\LTright}{\fill}
\noindent
\begin{tabular}{@{}>{\raggedright\arraybackslash}p{39.459mm}>{\raggedright\arraybackslash}p{39.459mm}>{\raggedright\arraybackslash}p{39.459mm}@{}}
\toprule
\textbf{Person and number} & \textbf{Present indicative} & \textbf{Preterite indicative} \\
\midrule
1 sg. & Ic hæbbe ðæt hūs geworht. & Ic hæfde ðæt hūs geworht. \\
2 sg. & Ðū hæfst ðæt hūs geworht. & Ðū hæfdest ðæt hūs geworht. \\
3 sg. & Hē hæfð ðæt hūs geworht. & Hē hæfde ðæt hūs geworht. \\
1 pl. & Wē habbað ðæt hūs geworht. & Wē hæfdon ðæt hūs geworht. \\
2 pl. & Gē habbað ðæt hūs geworht. & Gē hæfdon ðæt hūs geworht. \\
3 pl. & Hīe habbað ðæt hūs geworht. & Hīe hæfdon ðæt hūs geworht. \\
\bottomrule
\end{tabular}
\par\endgroup
% course-panel:T01-P01:end
% mdtable:T01:end
\reserve{11\baselineskip}
% mdparagraph:P026:begin
\textbf{Translations of the twelve positions}
% mdparagraph:P026:end

% mdtable:T02:begin
% course-table-style: 2026-10
% course-panel:T02-P01:begin
\par\begingroup
\fontsize{10.5}{13.2}\selectfont
\setlength{\LTleft}{0pt}\setlength{\LTright}{\fill}
\noindent
\begin{tabular}{@{}>{\raggedright\arraybackslash}p{39.459mm}>{\raggedright\arraybackslash}p{39.459mm}>{\raggedright\arraybackslash}p{39.459mm}@{}}
\toprule
\textbf{Person and number} & \textbf{Present viewpoint} & \textbf{Past viewpoint} \\
\midrule
1 sg. & I have built the house. & I had built the house. \\
2 sg. & You have built the house. & You had built the house. \\
3 sg. & He has built the house. & He had built the house. \\
1 pl. & We have built the house. & We had built the house. \\
2 pl. & You have built the house. & You had built the house. \\
3 pl. & They have built the house. & They had built the house. \\
\bottomrule
\end{tabular}
\par\endgroup
% course-panel:T02-P01:end
% mdtable:T02:end
% mdparagraph:P027:begin
This is a table of a construction in the indicative, not a new tense ending. The object ðæt hūs stays accusative even when the subject changes. The six auxiliary forms in each tense were learned earlier.
% mdparagraph:P027:end

% mdparagraph:P028:begin
Use\BookPoint{bk-ch12-learner-e89b8e1c8555} the bare participle bases \textbf{geworht, gierwed, sętt} in the assigned compositions. Old English also has participles agreeing with the object in this construction. A properly evidenced agreeing form is a valid alternative; the bare forms here do not teach a complete participial adjective declension. Nor does this pattern establish every causative use of habban, or its use with every intransitive verb.
% mdparagraph:P028:end

\reserve{7\baselineskip}
% mdparagraph:P029:begin
\section*{3. Active clauses, passives, and other uses of be}\BookHeading{nav-ch12-learner-0003}{subsection}{3. Active clauses, passives, and other uses of be}{1}
% mdparagraph:P029:end

% teaching: NEW-U12-GRAM-u12-active-passive
% mdparagraph:P030:begin
An\BookPoint{bk-ch12-learner-313b0bb61c54} active transitive clause names an agent as subject and the affected entity as object. In an agentless passive the affected entity becomes nominative subject, and the active agent is no longer named. Use the already learned passive auxiliary with the selected past participle. Do not add a new expression for a passive agent.
% mdparagraph:P030:end

\reserve{8\baselineskip}
% example: NEW-U12-EXAMPLE05
% mdparagraph:P031:begin
Se mann band ðæt hors. Ðæt hors wearð bunden.
% mdparagraph:P031:end

% mdparagraph:P032:begin
The man bound the horse. The horse was bound.
% mdparagraph:P032:end

% mdparagraph:P033:begin
The man is the active subject, and the horse is the active object. In the second clause the horse is the subject. Wearð often presents the event or entrance into a condition; the passive itself does not tell who performed the binding.
% mdparagraph:P033:end

\reserve{12\baselineskip}
% mdparagraph:P034:begin
Context: The context presents the horse in its resulting bound state; wesan does not invariably exclude an event reading.
% mdparagraph:P034:end

\reserve{8\baselineskip}
% example: NEW-U12-EXAMPLE06
% mdparagraph:P035:begin
Ðæt hors wæs bunden.
% mdparagraph:P035:end

% mdparagraph:P036:begin
The horse was bound.
% mdparagraph:P036:end

% mdparagraph:P037:begin
Here\BookPoint{bk-ch12-learner-20f46332400d} wæs presents the resulting condition. This is a contextual contrast: wesan can also occur in accounts of events, and weorðan need not always be translated literally “became”.
% mdparagraph:P037:end

% teaching: NEW-U12-GRAM-u12-wesan-intransitive-perfect
% teaching: NEW-U12-VOC-WESAN
% mdparagraph:P038:begin
For\BookPoint{bk-ch12-learner-a01000ad790b} recognition, \textbf{wesan + gewiten}, the past participle of gewītan, can express a completed departure. The subject is the person who departs, not a patient acted on by an unnamed agent.
% mdparagraph:P038:end

\reserve{12\baselineskip}
% mdparagraph:P039:begin
Context: In a past account, the man’s departure precedes the next event.
% mdparagraph:P039:end

\reserve{8\baselineskip}
% example: NEW-U12-EXAMPLE07
% mdparagraph:P040:begin
Se mann wæs gewiten.
% mdparagraph:P040:end

% mdparagraph:P041:begin
The man had departed.
% mdparagraph:P041:end

% teaching: NEW-U12-GRAM-u12-wesan-present-participle
% mdparagraph:P042:begin
Also\BookPoint{bk-ch12-learner-d03b0ae4da78} recognize \textbf{wesan + sittende}, the present participle of sittan. The context can support a continuing activity or state. English may use a simple past or a progressive; the Old English form does not impose a modern English progressive category.
% mdparagraph:P042:end

\reserve{12\baselineskip}
% mdparagraph:P043:begin
Context: The account describes a period during which the man was seated.
% mdparagraph:P043:end

\reserve{8\baselineskip}
% example: NEW-U12-EXAMPLE08
% mdparagraph:P044:begin
Se mann wæs sittende.
% mdparagraph:P044:end

% mdparagraph:P045:begin
The man sat /\allowbreak{} was sitting.
% mdparagraph:P045:end

% mdparagraph:P046:begin
These two uses are for recognition in the assigned examples and reading. Do not construct an unrestricted intransitive perfect or progressive system from them.
% mdparagraph:P046:end

\reserve{7\baselineskip}
% mdparagraph:P047:begin
\section*{4. Time in reported speech}\BookHeading{nav-ch12-learner-0004}{subsection}{4. Time in reported speech}{1}
% mdparagraph:P047:end

% teaching: NEW-U12-GRAM-u12-reported-sequence
% mdparagraph:P048:begin
A\BookPoint{bk-ch12-learner-f19bc233592b} reported statement has its own temporal viewpoint. Ask whether the reported event is simultaneous with the saying, already completed before it, or intended afterward. The earlier reporting subjunctive remains available. A subjunctive does not mean that the report is false, and a past reporting verb does not require mechanical replacement of every dependent tense.
% mdparagraph:P048:end

\reserve{12\baselineskip}
% mdparagraph:P049:begin
Context: In each report, the presence is simultaneous with the saying.
% mdparagraph:P049:end

\reserve{8\baselineskip}
% example: NEW-U12-EXAMPLE09
% mdparagraph:P050:begin
Hē sægð ðæt hē sīe hēr. Hē sægde ðæt hē wǣre hēr.
% mdparagraph:P050:end

% mdparagraph:P051:begin
He says that he is here. He said that he was here.
% mdparagraph:P051:end

% mdparagraph:P052:begin
Sīe is present subjunctive and wǣre preterite subjunctive. Each places the reported presence at the time of its own saying. The past report does not by itself say that the person is absent now.
% mdparagraph:P052:end

\reserve{12\baselineskip}
% mdparagraph:P053:begin
Context: The building precedes the past saying.
% mdparagraph:P053:end

\reserve{8\baselineskip}
% example: NEW-U12-EXAMPLE10
% mdparagraph:P054:begin
Hē sægde ðæt hē ðæt hūs geworht hæfde.
% mdparagraph:P054:end

% mdparagraph:P055:begin
He said that he had built the house.
% mdparagraph:P055:end

% mdparagraph:P056:begin
In\BookPoint{bk-ch12-learner-234095e10d2a} this assigned report, hæfde is preterite subjunctive singular. It has the same spelling as first and third singular indicative hæfde. The reported-statement construction determines the mood; the spelling alone does not. Geworht with hæfde places the building before the past saying. The complete table in §2 remains labelled indicative.
% mdparagraph:P056:end

\reserve{12\baselineskip}
% mdparagraph:P057:begin
Context: The speaker reports an intention; the sentence does not establish whether the arrival occurred.
% mdparagraph:P057:end

\reserve{8\baselineskip}
% example: NEW-U12-EXAMPLE11
% mdparagraph:P058:begin
Hē sægde ðæt hē wolde cuman.
% mdparagraph:P058:end

% mdparagraph:P059:begin
He said that he intended to come.
% mdparagraph:P059:end

% mdparagraph:P060:begin
Wolde\BookPoint{bk-ch12-learner-339d6a4d03a5} with a plain infinitive reports a wish or intention viewed from the past. The intended coming is later than that intention. The sentence does not assert that the coming took place. This is the use of willan required here, not a general future auxiliary.
% mdparagraph:P060:end

\reserve{7\baselineskip}
% mdparagraph:P061:begin
\section*{5. Verb systems and their uses}\BookHeading{nav-ch12-learner-0005}{subsection}{5. Verb systems and their uses}{1}
% mdparagraph:P061:end

% mdparagraph:P062:begin
The following indicative tables include all three persons in both numbers for each of the two tenses. Their lexical uses are stated beside them. Forms of an additional mood or a nonfinite construction are learned separately, even when a spelling happens to coincide.
% mdparagraph:P062:end

% mdparagraph:P063:begin
In the ordinary strong pattern, present singular endings are -e, -st, -ð and the plural is -að, with the stated stem changes and contractions. Preterite first and third singular share one stem; second singular uses the plural stem plus -e, and all plural persons use -on. Learn the contracted flēon distribution directly.
% mdparagraph:P063:end

% teaching: NEW-U12-GRAM-u12-adrifan-indicative
% teaching: NEW-U12-VOC-ADRIFAN
% mdparagraph:P064:begin
\textbf{ādrīfan: drive out; expel.} Ādrīfan takes an accusative person or group expelled. The earlier fram + dative can identify the place from which they are driven. Present ī contrasts with singular past ā and the short i of second singular and plural past forms.
% mdparagraph:P064:end

% teaching: NEW-U12-PAR-adrifan-IND
\reserve{11\baselineskip}
% mdparagraph:P065:begin
\textbf{ādrīfan; strong verb: present and preterite indicative}
% mdparagraph:P065:end

% table: NEW-U12-PAR-adrifan-IND
% mdtable:T03:begin
% course-table-style: 2026-10
% course-panel:T03-P01:begin
\par\begingroup
\fontsize{10.5}{13.2}\selectfont
\setlength{\LTleft}{0pt}\setlength{\LTright}{\fill}
\noindent
\begin{tabular}{@{}>{\raggedright\arraybackslash}p{39.459mm}>{\raggedright\arraybackslash}p{39.459mm}>{\raggedright\arraybackslash}p{39.459mm}@{}}
\toprule
\textbf{Person and number} & \textbf{Present} & \textbf{Preterite} \\
\midrule
1 sg. & ādrīfe & ādrāf \\
2 sg. & ādrīfst & ādrife \\
3 sg. & ādrīfð & ādrāf \\
1–3 pl. & ādrīfað & ādrifon \\
\bottomrule
\end{tabular}
\par\endgroup
% course-panel:T03-P01:end
% mdtable:T03:end
% teaching: NEW-U12-GRAM-u12-gehatan-indicative
% teaching: NEW-U12-VOC-GEHATAN
% mdparagraph:P066:begin
\textbf{gehātan: promise.} Gehātan takes the accusative thing promised and, when stated, a dative recipient. Learn hā- in the present beside hǣt- in second/\allowbreak{}third singular and hēt- in the preterite. The prefix ge- remains. No additional calling or commanding sense is assigned.
% mdparagraph:P066:end

% teaching: NEW-U12-PAR-gehatan-IND
\reserve{11\baselineskip}
% mdparagraph:P067:begin
\textbf{gehātan; strong verb: present and preterite indicative}
% mdparagraph:P067:end

% table: NEW-U12-PAR-gehatan-IND
% mdtable:T04:begin
% course-table-style: 2026-10
% course-panel:T04-P01:begin
\par\begingroup
\fontsize{10.5}{13.2}\selectfont
\setlength{\LTleft}{0pt}\setlength{\LTright}{\fill}
\noindent
\begin{tabular}{@{}>{\raggedright\arraybackslash}p{39.459mm}>{\raggedright\arraybackslash}p{39.459mm}>{\raggedright\arraybackslash}p{39.459mm}@{}}
\toprule
\textbf{Person and number} & \textbf{Present} & \textbf{Preterite} \\
\midrule
1 sg. & gehāte & gehēt \\
2 sg. & gehǣtst & gehēte \\
3 sg. & gehǣtt & gehēt \\
1–3 pl. & gehātað & gehēton \\
\bottomrule
\end{tabular}
\par\endgroup
% course-panel:T04-P01:end
% mdtable:T04:end
% teaching: NEW-U12-GRAM-u12-fleon-indicative
% teaching: NEW-U12-VOC-FLEON
% mdparagraph:P068:begin
\textbf{flēon: flee; flee from; avoid.} Flēon can stand without an object or take the accusative person or thing avoided. It is distinct from flēogan, fly. Keep flēo-, flīeh-, flēah- and flug- in their printed positions; do not form its contracted paradigm from the infinitive by adding ordinary endings.
% mdparagraph:P068:end

% teaching: NEW-U12-PAR-fleon-IND
\reserve{11\baselineskip}
% mdparagraph:P069:begin
\textbf{flēon; contracted strong verb: present and preterite indicative}
% mdparagraph:P069:end

% table: NEW-U12-PAR-fleon-IND
% mdtable:T05:begin
% course-table-style: 2026-10
% course-panel:T05-P01:begin
\par\begingroup
\fontsize{10.5}{13.2}\selectfont
\setlength{\LTleft}{0pt}\setlength{\LTright}{\fill}
\noindent
\begin{tabular}{@{}>{\raggedright\arraybackslash}p{39.459mm}>{\raggedright\arraybackslash}p{39.459mm}>{\raggedright\arraybackslash}p{39.459mm}@{}}
\toprule
\textbf{Person and number} & \textbf{Present} & \textbf{Preterite} \\
\midrule
1 sg. & flēo & flēah \\
2 sg. & flīehst & fluge \\
3 sg. & flīehð & flēah \\
1–3 pl. & flēoð & flugon \\
\bottomrule
\end{tabular}
\par\endgroup
% course-panel:T05-P01:end
% mdtable:T05:end
% teaching: NEW-U12-GRAM-u12-gemetan-indicative
% teaching: NEW-U12-VOC-GEMETAN
% mdparagraph:P070:begin
Weak I gemētan and grētan retain ē. Their t meets the third-singular dental in -tt, and their preterite has -tte, -ttest, -tton. Ēhtan and cȳðan have their own stated dental combinations.
% mdparagraph:P070:end

% mdparagraph:P071:begin
\textbf{gemētan: find; meet.} Gemētan takes the accusative person or thing found. Distinguish preterite singular gemētte from third singular present gemētt.
% mdparagraph:P071:end

% teaching: NEW-U12-PAR-gemetan-IND
\reserve{11\baselineskip}
% mdparagraph:P072:begin
\textbf{gemētan; weak verb: present and preterite indicative}
% mdparagraph:P072:end

% table: NEW-U12-PAR-gemetan-IND
% mdtable:T06:begin
% course-table-style: 2026-10
% course-panel:T06-P01:begin
\par\begingroup
\fontsize{10.5}{13.2}\selectfont
\setlength{\LTleft}{0pt}\setlength{\LTright}{\fill}
\noindent
\begin{tabular}{@{}>{\raggedright\arraybackslash}p{39.459mm}>{\raggedright\arraybackslash}p{39.459mm}>{\raggedright\arraybackslash}p{39.459mm}@{}}
\toprule
\textbf{Person and number} & \textbf{Present} & \textbf{Preterite} \\
\midrule
1 sg. & gemēte & gemētte \\
2 sg. & gemētst & gemēttest \\
3 sg. & gemētt & gemētte \\
1–3 pl. & gemētað & gemētton \\
\bottomrule
\end{tabular}
\par\endgroup
% course-panel:T06-P01:end
% mdtable:T06:end
% teaching: NEW-U12-GRAM-u12-geopenian-indicative
% teaching: NEW-U12-VOC-GEOPENIAN
% mdparagraph:P073:begin
Weak II geopenian and leornian have present -ie, -ast, -að, -iað and preterite -ode, -odest, -ode, -odon. The plural spelling is shared by all persons.
% mdparagraph:P073:end

% mdparagraph:P074:begin
\textbf{geopenian: open.} Geopenian takes the accusative thing opened, including a mouth or eyes. The table establishes this transitive use.
% mdparagraph:P074:end

% teaching: NEW-U12-PAR-geopenian-IND
\reserve{11\baselineskip}
% mdparagraph:P075:begin
\textbf{geopenian; weak verb: present and preterite indicative}
% mdparagraph:P075:end

% table: NEW-U12-PAR-geopenian-IND
% mdtable:T07:begin
% course-table-style: 2026-10
% course-panel:T07-P01:begin
\par\begingroup
\fontsize{10.5}{13.2}\selectfont
\setlength{\LTleft}{0pt}\setlength{\LTright}{\fill}
\noindent
\begin{tabular}{@{}>{\raggedright\arraybackslash}p{39.459mm}>{\raggedright\arraybackslash}p{39.459mm}>{\raggedright\arraybackslash}p{39.459mm}@{}}
\toprule
\textbf{Person and number} & \textbf{Present} & \textbf{Preterite} \\
\midrule
1 sg. & geopenie & geopenode \\
2 sg. & geopenast & geopenodest \\
3 sg. & geopenað & geopenode \\
1–3 pl. & geopeniað & geopenodon \\
\bottomrule
\end{tabular}
\par\endgroup
% course-panel:T07-P01:end
% mdtable:T07:end
% teaching: NEW-U12-GRAM-u12-gretan-indicative
% teaching: NEW-U12-VOC-GRETAN
% mdparagraph:P076:begin
\textbf{grētan: greet; take leave of.} Grētan takes an accusative person greeted or taken leave of. It also takes the native land as object in the assigned leave-taking expression. This weak verb is distinct from the strong homonym meaning weep.
% mdparagraph:P076:end

% teaching: NEW-U12-PAR-gretan-IND
\reserve{11\baselineskip}
% mdparagraph:P077:begin
\textbf{grētan; weak verb: present and preterite indicative}
% mdparagraph:P077:end

% table: NEW-U12-PAR-gretan-IND
% mdtable:T08:begin
% course-table-style: 2026-10
% course-panel:T08-P01:begin
\par\begingroup
\fontsize{10.5}{13.2}\selectfont
\setlength{\LTleft}{0pt}\setlength{\LTright}{\fill}
\noindent
\begin{tabular}{@{}>{\raggedright\arraybackslash}p{39.459mm}>{\raggedright\arraybackslash}p{39.459mm}>{\raggedright\arraybackslash}p{39.459mm}@{}}
\toprule
\textbf{Person and number} & \textbf{Present} & \textbf{Preterite} \\
\midrule
1 sg. & grēte & grētte \\
2 sg. & grētst & grēttest \\
3 sg. & grētt & grētte \\
1–3 pl. & grētað & grētton \\
\bottomrule
\end{tabular}
\par\endgroup
% course-panel:T08-P01:end
% mdtable:T08:end
% teaching: NEW-U12-GRAM-u12-cythan-indicative
% teaching: NEW-U12-VOC-CYTHAN
% mdparagraph:P078:begin
\textbf{cȳðan: announce; make known.} Cȳðan takes an accusative thing announced or a ðæt content clause; a recipient, when expressed, is dative. The selected present forms are cȳst and cȳð; learn their dental combinations directly. Analogical cȳðst follows the documented retained-dental pattern in the second singular. The selected preterite has cȳðd-; cȳdd- is an attested alternative.
% mdparagraph:P078:end

% teaching: NEW-U12-PAR-cythan-IND
\reserve{11\baselineskip}
% mdparagraph:P079:begin
\textbf{cȳðan; weak verb: present and preterite indicative}
% mdparagraph:P079:end

% table: NEW-U12-PAR-cythan-IND
% mdtable:T09:begin
% course-table-style: 2026-10
% course-panel:T09-P01:begin
\par\begingroup
\fontsize{10.5}{13.2}\selectfont
\setlength{\LTleft}{0pt}\setlength{\LTright}{\fill}
\noindent
\begin{tabular}{@{}>{\raggedright\arraybackslash}p{39.459mm}>{\raggedright\arraybackslash}p{39.459mm}>{\raggedright\arraybackslash}p{39.459mm}@{}}
\toprule
\textbf{Person and number} & \textbf{Present} & \textbf{Preterite} \\
\midrule
1 sg. & cȳðe & cȳðde \\
2 sg. & cȳst & cȳðdest \\
3 sg. & cȳð & cȳðde \\
1–3 pl. & cȳðað & cȳðdon \\
\bottomrule
\end{tabular}
\par\endgroup
% course-panel:T09-P01:end
% mdtable:T09:end
% teaching: NEW-U12-GRAM-u12-ehtan-indicative
% teaching: NEW-U12-VOC-EHTAN
% mdparagraph:P080:begin
\textbf{ēhtan: pursue; persecute.} Use the accusative person pursued in the assigned construction, directly supported by the sheep-and-wolves Gospel passage. Genitive government also occurs and is valid with supported forms. Preserve ēht in third singular present, and ēhte in the preterite.
% mdparagraph:P080:end

% teaching: NEW-U12-PAR-ehtan-IND
\reserve{11\baselineskip}
% mdparagraph:P081:begin
\textbf{ēhtan; weak verb: present and preterite indicative}
% mdparagraph:P081:end

% table: NEW-U12-PAR-ehtan-IND
% mdtable:T10:begin
% course-table-style: 2026-10
% course-panel:T10-P01:begin
\par\begingroup
\fontsize{10.5}{13.2}\selectfont
\setlength{\LTleft}{0pt}\setlength{\LTright}{\fill}
\noindent
\begin{tabular}{@{}>{\raggedright\arraybackslash}p{39.459mm}>{\raggedright\arraybackslash}p{39.459mm}>{\raggedright\arraybackslash}p{39.459mm}@{}}
\toprule
\textbf{Person and number} & \textbf{Present} & \textbf{Preterite} \\
\midrule
1 sg. & ēhte & ēhte \\
2 sg. & ēhtst & ēhtest \\
3 sg. & ēht & ēhte \\
1–3 pl. & ēhtað & ēhton \\
\bottomrule
\end{tabular}
\par\endgroup
% course-panel:T10-P01:end
% mdtable:T10:end
% teaching: NEW-U12-GRAM-u12-agan-indicative
% teaching: NEW-U12-VOC-AGAN
% mdparagraph:P082:begin
\textbf{āgan: own; possess.} Āgan takes an accusative possession. Its preterite-present pattern has āh/\allowbreak{}āhst/\allowbreak{}āh and plural āgon in the present, āhte/\allowbreak{}āhtest/\allowbreak{}āhte and plural āhton in the preterite. This is the positive indicative only. The supplied negative nāh below does not teach a complete negative system.
% mdparagraph:P082:end

% teaching: NEW-U12-PAR-agan-IND
\reserve{11\baselineskip}
% mdparagraph:P083:begin
\textbf{āgan; preterite-present verb: present and preterite indicative}
% mdparagraph:P083:end

% table: NEW-U12-PAR-agan-IND
% mdtable:T11:begin
% course-table-style: 2026-10
% course-panel:T11-P01:begin
\par\begingroup
\fontsize{10.5}{13.2}\selectfont
\setlength{\LTleft}{0pt}\setlength{\LTright}{\fill}
\noindent
\begin{tabular}{@{}>{\raggedright\arraybackslash}p{39.459mm}>{\raggedright\arraybackslash}p{39.459mm}>{\raggedright\arraybackslash}p{39.459mm}@{}}
\toprule
\textbf{Person and number} & \textbf{Present} & \textbf{Preterite} \\
\midrule
1 sg. & āh & āhte \\
2 sg. & āhst & āhtest \\
3 sg. & āh & āhte \\
1–3 pl. & āgon & āhton \\
\bottomrule
\end{tabular}
\par\endgroup
% course-panel:T11-P01:end
% mdtable:T11:end
% teaching: NEW-U12-GRAM-u12-leornian-indicative
% teaching: NEW-U12-VOC-LEORNIAN
% mdparagraph:P084:begin
\textbf{leornian: learn.} An expressed thing learned is accusative. The person from whom one learns is introduced by æt + dative. The phrase learn from me can stand without an expressed object.
% mdparagraph:P084:end

% teaching: NEW-U12-PAR-leornian-IND
\reserve{11\baselineskip}
% mdparagraph:P085:begin
\textbf{leornian; weak verb: present and preterite indicative}
% mdparagraph:P085:end

% table: NEW-U12-PAR-leornian-IND
% mdtable:T12:begin
% course-table-style: 2026-10
% course-panel:T12-P01:begin
\par\begingroup
\fontsize{10.5}{13.2}\selectfont
\setlength{\LTleft}{0pt}\setlength{\LTright}{\fill}
\noindent
\begin{tabular}{@{}>{\raggedright\arraybackslash}p{39.459mm}>{\raggedright\arraybackslash}p{39.459mm}>{\raggedright\arraybackslash}p{39.459mm}@{}}
\toprule
\textbf{Person and number} & \textbf{Present} & \textbf{Preterite} \\
\midrule
1 sg. & leornie & leornode \\
2 sg. & leornast & leornodest \\
3 sg. & leornað & leornode \\
1–3 pl. & leorniað & leornodon \\
\bottomrule
\end{tabular}
\par\endgroup
% course-panel:T12-P01:end
% mdtable:T12:end
\reserve{7\baselineskip}
% mdparagraph:P086:begin
\subsection*{Commands}\BookHeading{nav-ch12-learner-0006}{subsubsection}{Commands}{0}
% mdparagraph:P086:end

% mdparagraph:P087:begin
A singular command addresses one person; a plural command addresses two or more. Both members of each new pair are given. The selected strong long stems and weak I long stems have a bare singular; flēon has flēoh. Weak II uses -a in the singular and -iað in the plural. These commands add no new subjunctive paradigms.
% mdparagraph:P087:end

% teaching: NEW-U12-GRAM-u12-sendan-imperative
% teaching: NEW-U12-PAR-sendan-IMP
\reserve{7\baselineskip}
% mdparagraph:P088:begin
\textbf{sęndan; weak verb: singular and plural commands}
% mdparagraph:P088:end

% table: NEW-U12-PAR-sendan-IMP
% mdtable:T13:begin
% course-table-style: 2026-10
% course-panel:T13-P01:begin
\par\begingroup
\fontsize{10.5}{13.2}\selectfont
\setlength{\LTleft}{0pt}\setlength{\LTright}{\fill}
\noindent
\begin{tabular}{@{}>{\raggedright\arraybackslash}p{63.342mm}>{\raggedright\arraybackslash}p{57.846mm}@{}}
\toprule
\textbf{Person and number} & \textbf{Command} \\
\midrule
2 sg. & sęnd \\
2 pl. & sęndað \\
\bottomrule
\end{tabular}
\par\endgroup
% course-panel:T13-P01:end
% mdtable:T13:end
% teaching: NEW-U12-GRAM-u12-bringan-imperative
% teaching: NEW-U12-PAR-bringan-IMP
\reserve{7\baselineskip}
% mdparagraph:P089:begin
\textbf{bringan; mixed verb: singular and plural commands}
% mdparagraph:P089:end

% table: NEW-U12-PAR-bringan-IMP
% mdtable:T14:begin
% course-table-style: 2026-10
% course-panel:T14-P01:begin
\par\begingroup
\fontsize{10.5}{13.2}\selectfont
\setlength{\LTleft}{0pt}\setlength{\LTright}{\fill}
\noindent
\begin{tabular}{@{}>{\raggedright\arraybackslash}p{63.342mm}>{\raggedright\arraybackslash}p{57.846mm}@{}}
\toprule
\textbf{Person and number} & \textbf{Command} \\
\midrule
2 sg. & bring \\
2 pl. & bringað \\
\bottomrule
\end{tabular}
\par\endgroup
% course-panel:T14-P01:end
% mdtable:T14:end
% teaching: NEW-U12-GRAM-u12-adrifan-imperative
% teaching: NEW-U12-PAR-adrifan-IMP
\reserve{7\baselineskip}
% mdparagraph:P090:begin
\textbf{ādrīfan; strong verb: singular and plural commands}
% mdparagraph:P090:end

% table: NEW-U12-PAR-adrifan-IMP
% mdtable:T15:begin
% course-table-style: 2026-10
% course-panel:T15-P01:begin
\par\begingroup
\fontsize{10.5}{13.2}\selectfont
\setlength{\LTleft}{0pt}\setlength{\LTright}{\fill}
\noindent
\begin{tabular}{@{}>{\raggedright\arraybackslash}p{63.342mm}>{\raggedright\arraybackslash}p{57.846mm}@{}}
\toprule
\textbf{Person and number} & \textbf{Command} \\
\midrule
2 sg. & ādrīf \\
2 pl. & ādrīfað \\
\bottomrule
\end{tabular}
\par\endgroup
% course-panel:T15-P01:end
% mdtable:T15:end
% teaching: NEW-U12-GRAM-u12-fleon-imperative
% teaching: NEW-U12-PAR-fleon-IMP
\reserve{7\baselineskip}
% mdparagraph:P091:begin
\textbf{flēon; contracted strong verb: singular and plural commands}
% mdparagraph:P091:end

% table: NEW-U12-PAR-fleon-IMP
% mdtable:T16:begin
% course-table-style: 2026-10
% course-panel:T16-P01:begin
\par\begingroup
\fontsize{10.5}{13.2}\selectfont
\setlength{\LTleft}{0pt}\setlength{\LTright}{\fill}
\noindent
\begin{tabular}{@{}>{\raggedright\arraybackslash}p{63.342mm}>{\raggedright\arraybackslash}p{57.846mm}@{}}
\toprule
\textbf{Person and number} & \textbf{Command} \\
\midrule
2 sg. & flēoh \\
2 pl. & flēoð \\
\bottomrule
\end{tabular}
\par\endgroup
% course-panel:T16-P01:end
% mdtable:T16:end
% teaching: NEW-U12-GRAM-u12-gemetan-imperative
% teaching: NEW-U12-PAR-gemetan-IMP
\reserve{7\baselineskip}
% mdparagraph:P092:begin
\textbf{gemētan; weak verb: singular and plural commands}
% mdparagraph:P092:end

% table: NEW-U12-PAR-gemetan-IMP
% mdtable:T17:begin
% course-table-style: 2026-10
% course-panel:T17-P01:begin
\par\begingroup
\fontsize{10.5}{13.2}\selectfont
\setlength{\LTleft}{0pt}\setlength{\LTright}{\fill}
\noindent
\begin{tabular}{@{}>{\raggedright\arraybackslash}p{63.342mm}>{\raggedright\arraybackslash}p{57.846mm}@{}}
\toprule
\textbf{Person and number} & \textbf{Command} \\
\midrule
2 sg. & gemēt \\
2 pl. & gemētað \\
\bottomrule
\end{tabular}
\par\endgroup
% course-panel:T17-P01:end
% mdtable:T17:end
% teaching: NEW-U12-GRAM-u12-geopenian-imperative
% teaching: NEW-U12-PAR-geopenian-IMP
\reserve{7\baselineskip}
% mdparagraph:P093:begin
\textbf{geopenian; weak verb: singular and plural commands}
% mdparagraph:P093:end

% table: NEW-U12-PAR-geopenian-IMP
% mdtable:T18:begin
% course-table-style: 2026-10
% course-panel:T18-P01:begin
\par\begingroup
\fontsize{10.5}{13.2}\selectfont
\setlength{\LTleft}{0pt}\setlength{\LTright}{\fill}
\noindent
\begin{tabular}{@{}>{\raggedright\arraybackslash}p{63.342mm}>{\raggedright\arraybackslash}p{57.846mm}@{}}
\toprule
\textbf{Person and number} & \textbf{Command} \\
\midrule
2 sg. & geopena \\
2 pl. & geopeniað \\
\bottomrule
\end{tabular}
\par\endgroup
% course-panel:T18-P01:end
% mdtable:T18:end
% teaching: NEW-U12-GRAM-u12-weorpan-imperative
% teaching: NEW-U12-PAR-weorpan-IMP
\reserve{7\baselineskip}
% mdparagraph:P094:begin
\textbf{weorpan; strong verb: singular and plural commands}
% mdparagraph:P094:end

% table: NEW-U12-PAR-weorpan-IMP
% mdtable:T19:begin
% course-table-style: 2026-10
% course-panel:T19-P01:begin
\par\begingroup
\fontsize{10.5}{13.2}\selectfont
\setlength{\LTleft}{0pt}\setlength{\LTright}{\fill}
\noindent
\begin{tabular}{@{}>{\raggedright\arraybackslash}p{63.342mm}>{\raggedright\arraybackslash}p{57.846mm}@{}}
\toprule
\textbf{Person and number} & \textbf{Command} \\
\midrule
2 sg. & weorp \\
2 pl. & weorpað \\
\bottomrule
\end{tabular}
\par\endgroup
% course-panel:T19-P01:end
% mdtable:T19:end
% teaching: NEW-U12-GRAM-u12-leornian-imperative
% teaching: NEW-U12-PAR-leornian-IMP
\reserve{7\baselineskip}
% mdparagraph:P095:begin
\textbf{leornian; weak verb: singular and plural commands}
% mdparagraph:P095:end

% table: NEW-U12-PAR-leornian-IMP
% mdtable:T20:begin
% course-table-style: 2026-10
% course-panel:T20-P01:begin
\par\begingroup
\fontsize{10.5}{13.2}\selectfont
\setlength{\LTleft}{0pt}\setlength{\LTright}{\fill}
\noindent
\begin{tabular}{@{}>{\raggedright\arraybackslash}p{63.342mm}>{\raggedright\arraybackslash}p{57.846mm}@{}}
\toprule
\textbf{Person and number} & \textbf{Command} \\
\midrule
2 sg. & leorna \\
2 pl. & leorniað \\
\bottomrule
\end{tabular}
\par\endgroup
% course-panel:T20-P01:end
% mdtable:T20:end
% teaching: NEW-U12-GRAM-u12-risan-imperative
% teaching: NEW-U12-PAR-risan-IMP
\reserve{7\baselineskip}
% mdparagraph:P096:begin
\textbf{rīsan; strong verb: singular and plural commands}
% mdparagraph:P096:end

% table: NEW-U12-PAR-risan-IMP
% mdtable:T21:begin
% course-table-style: 2026-10
% course-panel:T21-P01:begin
\par\begingroup
\fontsize{10.5}{13.2}\selectfont
\setlength{\LTleft}{0pt}\setlength{\LTright}{\fill}
\noindent
\begin{tabular}{@{}>{\raggedright\arraybackslash}p{63.342mm}>{\raggedright\arraybackslash}p{57.846mm}@{}}
\toprule
\textbf{Person and number} & \textbf{Command} \\
\midrule
2 sg. & rīs \\
2 pl. & rīsað \\
\bottomrule
\end{tabular}
\par\endgroup
% course-panel:T21-P01:end
% mdtable:T21:end
% mdparagraph:P097:begin
The already learned \textbf{sleah} is the singular command of slēan. For rise use the newly displayed \textbf{rīs}; do not add a new prefix to the selected verb. Earlier \textbf{cum}, \textbf{gā}, \textbf{nim} and \textbf{sęle} remain available in their taught uses.
% mdparagraph:P097:end

\reserve{7\baselineskip}
% mdparagraph:P098:begin
\subsection*{Participle bases}\BookHeading{nav-ch12-learner-0007}{subsubsection}{Participle bases}{0}
% mdparagraph:P098:end

% mdparagraph:P099:begin
A present participle characteristically has -ende; the past bases must be learned with their verbs. The first three pairs below are productive as base forms. Their uses remain limited to the constructions taught here or earlier. A base pair is not a complete adjectival declension.
% mdparagraph:P099:end

% teaching: NEW-U12-GRAM-u12-gewyrcean-participle-bases
% teaching: NEW-U12-PAR-gewyrcean-PART
\reserve{7\baselineskip}
% mdparagraph:P100:begin
\textbf{gewyrcean; weak verb: present and past participle bases}
% mdparagraph:P100:end

% table: NEW-U12-PAR-gewyrcean-PART
% mdtable:T22:begin
\par\begingroup
\fontsize{10.5}{13.2}\selectfont
\setlength{\LTleft}{0pt}\setlength{\LTright}{\fill}
\noindent
\begin{tabular}{@{}>{\raggedright\arraybackslash}p{69.84mm}>{\raggedright\arraybackslash}p{51.35mm}@{}}
\toprule
\textbf{Form} & \textbf{Base} \\
\midrule
Present participle & gewyrcende \\
Past participle & geworht \\
\bottomrule
\end{tabular}
\par\endgroup

% mdtable:T22:end
% teaching: NEW-U12-GRAM-u12-gierwan-participle-bases
% teaching: NEW-U12-PAR-gierwan-PART
\reserve{7\baselineskip}
% mdparagraph:P101:begin
\textbf{gierwan; weak verb: present and past participle bases}
% mdparagraph:P101:end

% table: NEW-U12-PAR-gierwan-PART
% mdtable:T23:begin
\par\begingroup
\fontsize{10.5}{13.2}\selectfont
\setlength{\LTleft}{0pt}\setlength{\LTright}{\fill}
\noindent
\begin{tabular}{@{}>{\raggedright\arraybackslash}p{69.84mm}>{\raggedright\arraybackslash}p{51.35mm}@{}}
\toprule
\textbf{Form} & \textbf{Base} \\
\midrule
Present participle & gierwende \\
Past participle & gierwed \\
\bottomrule
\end{tabular}
\par\endgroup

% mdtable:T23:end
% teaching: NEW-U12-GRAM-u12-settan-participle-bases
% teaching: NEW-U12-PAR-settan-PART
\reserve{7\baselineskip}
% mdparagraph:P102:begin
\textbf{sęttan; weak verb: present and past participle bases}
% mdparagraph:P102:end

% table: NEW-U12-PAR-settan-PART
% mdtable:T24:begin
\par\begingroup
\fontsize{10.5}{13.2}\selectfont
\setlength{\LTleft}{0pt}\setlength{\LTright}{\fill}
\noindent
\begin{tabular}{@{}>{\raggedright\arraybackslash}p{69.84mm}>{\raggedright\arraybackslash}p{51.35mm}@{}}
\toprule
\textbf{Form} & \textbf{Base} \\
\midrule
Present participle & sęttende \\
Past participle & sętt \\
\bottomrule
\end{tabular}
\par\endgroup

% mdtable:T24:end
% teaching: NEW-U12-GRAM-u12-gewitan-participle-bases
% mdparagraph:P103:begin
\textbf{Reference pairs.} Recognize both bases of gewītan and sittan. Only the stated constructions with gewiten and sittende are analyzed here; neither pair is assigned for independent composition.
% mdparagraph:P103:end

% teaching: NEW-U12-PAR-gewitan-PART
\reserve{7\baselineskip}
% mdparagraph:P104:begin
\textbf{gewītan; strong verb: present and past participle bases}
% mdparagraph:P104:end

% table: NEW-U12-PAR-gewitan-PART
% mdtable:T25:begin
\par\begingroup
\fontsize{10.5}{13.2}\selectfont
\setlength{\LTleft}{0pt}\setlength{\LTright}{\fill}
\noindent
\begin{tabular}{@{}>{\raggedright\arraybackslash}p{69.84mm}>{\raggedright\arraybackslash}p{51.35mm}@{}}
\toprule
\textbf{Form} & \textbf{Base} \\
\midrule
Present participle & gewītende \\
Past participle & gewiten \\
\bottomrule
\end{tabular}
\par\endgroup

% mdtable:T25:end
% teaching: NEW-U12-GRAM-u12-sittan-participle-bases
% teaching: NEW-U12-PAR-sittan-PART
\reserve{7\baselineskip}
% mdparagraph:P105:begin
\textbf{sittan; strong verb: present and past participle bases}
% mdparagraph:P105:end

% table: NEW-U12-PAR-sittan-PART
% mdtable:T26:begin
\par\begingroup
\fontsize{10.5}{13.2}\selectfont
\setlength{\LTleft}{0pt}\setlength{\LTright}{\fill}
\noindent
\begin{tabular}{@{}>{\raggedright\arraybackslash}p{69.84mm}>{\raggedright\arraybackslash}p{51.35mm}@{}}
\toprule
\textbf{Form} & \textbf{Base} \\
\midrule
Present participle & sittende \\
Past participle & seten \\
\bottomrule
\end{tabular}
\par\endgroup

% mdtable:T26:end
% mdparagraph:P106:begin
The exact selected bases \textbf{gierwed} and \textbf{sętt} stand beside the earlier citation spellings \textbf{(ge)gier(w)ed} and \textbf{(ge)sętt}. Do not treat the parentheses as letters to write in a sentence. The earlier supplied \textbf{geworht} is now available in the productive constructions of this unit.
% mdparagraph:P106:end

\reserve{7\baselineskip}
% mdparagraph:P107:begin
\subsection*{Plain and inflected infinitives}\BookHeading{nav-ch12-learner-0008}{subsubsection}{Plain and inflected infinitives}{0}
% mdparagraph:P107:end

% mdparagraph:P108:begin
Use the plain infinitive after the assigned modal: ādrīfan after ability expressed by magan, and genēosian after intention expressed by willan. The inflected forms are shown for recognition only; their display does not add new composition with tō.
% mdparagraph:P108:end

% teaching: NEW-U12-GRAM-u12-adrifan-modal-infinitive
% teaching: NEW-U12-PAR-adrifan-INF
\reserve{7\baselineskip}
% mdparagraph:P109:begin
\textbf{ādrīfan; strong verb: plain and inflected infinitives}
% mdparagraph:P109:end

% table: NEW-U12-PAR-adrifan-INF
% mdtable:T27:begin
\par\begingroup
\fontsize{10.5}{13.2}\selectfont
\setlength{\LTleft}{0pt}\setlength{\LTright}{\fill}
\noindent
\begin{tabular}{@{}>{\raggedright\arraybackslash}p{71.72mm}>{\raggedright\arraybackslash}p{49.46mm}@{}}
\toprule
\textbf{Form} & \textbf{Infinitive} \\
\midrule
Plain infinitive & ādrīfan \\
Inflected infinitive & ādrīfanne \\
\bottomrule
\end{tabular}
\par\endgroup

% mdtable:T27:end
% teaching: NEW-U12-GRAM-u12-geneosian-modal-infinitive
% teaching: NEW-U12-PAR-geneosian-INF
\reserve{7\baselineskip}
% mdparagraph:P110:begin
\textbf{genēosian; weak verb: plain and inflected infinitives}
% mdparagraph:P110:end

% table: NEW-U12-PAR-geneosian-INF
% mdtable:T28:begin
\par\begingroup
\fontsize{10.5}{13.2}\selectfont
\setlength{\LTleft}{0pt}\setlength{\LTright}{\fill}
\noindent
\begin{tabular}{@{}>{\raggedright\arraybackslash}p{68.91mm}>{\raggedright\arraybackslash}p{52.28mm}@{}}
\toprule
\textbf{Form} & \textbf{Infinitive} \\
\midrule
Plain infinitive & genēosian \\
Inflected infinitive & genēosianne \\
\bottomrule
\end{tabular}
\par\endgroup

% mdtable:T28:end
\reserve{7\baselineskip}
% mdparagraph:P111:begin
\section*{6. Nouns: full patterns and lexical stems}\BookHeading{nav-ch12-learner-0009}{subsection}{6. Nouns: full patterns and lexical stems}{1}
% mdparagraph:P111:end

% teaching: NEW-U12-NOMINAL-SYSTEMS
% teaching: NEW-U12-LEXICAL-DIAGNOSTICS
% teaching: NEW-U12-WEAK-NOMINALS
% mdparagraph:P112:begin
The complete patterns below give nominative, genitive, dative and accusative in both numbers. Where another noun uses the same pattern, its diagnostic row gives the headword, stem and distinguishing forms. The substitution instruction applies to \textbf{every corresponding case and number}, not merely to the examples in that row. Use the separate headword where expressly stated. Nouns keep their lexical gender.
% mdparagraph:P112:end

% mdparagraph:P113:begin
\textbf{Vowel change and syncope.} Fōt uses fēt in its selected dative singular and nominative/\allowbreak{}accusative plural. Fēond selects fēonde in the dative singular and fīend in the nominative/\allowbreak{}accusative plural; its genitive plural is fēonda. In ēðel, sāwol and wǣpen the medial vowel is absent from the ending-bearing stem: ēðl-, sāwl-, wǣpn-. Hærfest and pęning retain their medial vowel; the ordinary syncope rule requires a short second syllable.
% mdparagraph:P113:end

% mdparagraph:P114:begin
\textbf{Weak nouns.} Andwlita and sǣdtīma use the weak masculine pattern. Fǣmne and heorte use the weak feminine pattern, with fǣmn- and heort- before endings. The abstract suffix in gehīersumnes has -ness- before an ending. Gehīersumnes and ręst are productive in their assigned singular abstract meanings; their plural forms are \textbf{reference only}, without a new count-plural sense.
% mdparagraph:P114:end

% teaching: NEW-U12-VOC-BRYCG
% teaching: NEW-U12-NOM-BRYCG
% diagnostic: NEW-U12-DIAG-BRYCG
% Page review repair: keep the lexical lead with its first display or derivation.
\reserve{8\baselineskip}
% mdparagraph:P115:begin
\textbf{brycg} (feminine): bridge. Diagnostic stem: \textbf{brycg}.
% mdparagraph:P115:end

% mdtable:T29:begin
\par\begingroup
\fontsize{10.5}{13.2}\selectfont
\setlength{\LTleft}{0pt}\setlength{\LTright}{\fill}
\noindent
\begin{tabular}{@{}>{\raggedright\arraybackslash}p{39.24mm}>{\raggedright\arraybackslash}p{36.31mm}>{\raggedright\arraybackslash}p{42.82mm}@{}}
\toprule
\textbf{Genitive singular} & \textbf{Dative singular} & \textbf{Nominative plural} \\
\midrule
brycge & brycge & brycga \\
\bottomrule
\end{tabular}
\par\endgroup

% mdtable:T29:end
\reserve{9\baselineskip}
% mdparagraph:P116:begin
\textbf{brycg; feminine noun}
% mdparagraph:P116:end

% table: NEW-U12-NOM-BRYCG
% mdtable:T30:begin
\par\begingroup
\fontsize{10.5}{13.2}\selectfont
\setlength{\LTleft}{0pt}\setlength{\LTright}{\fill}
\noindent
\begin{tabular}{@{}>{\raggedright\arraybackslash}p{39.46mm}>{\raggedright\arraybackslash}p{39.46mm}>{\raggedright\arraybackslash}p{39.46mm}@{}}
\toprule
\textbf{Case} & \textbf{Singular} & \textbf{Plural} \\
\midrule
Nominative & brycg & brycga \\
Genitive & brycge & brycga \\
Dative & brycge & brycgum \\
Accusative & brycge & brycga \\
\bottomrule
\end{tabular}
\par\endgroup

% mdtable:T30:end
% mdparagraph:P117:begin
Selected nominative and accusative plural brycga have the accepted alternative \textbf{brycge}. The table uses -a consistently.
% mdparagraph:P117:end

% teaching: NEW-U12-VOC-GEHIERSUMNES
% teaching: NEW-U12-NOM-GEHIERSUMNES
% diagnostic: NEW-U12-DIAG-GEHIERSUMNES
% Page review repair: keep the lexical lead with its first display or derivation.
\reserve{8\baselineskip}
% mdparagraph:P118:begin
\textbf{gehīersumnes} (feminine): obedience. Diagnostic stem: \textbf{gehīersumness}.
% mdparagraph:P118:end

% mdtable:T31:begin
\par\begingroup
\fontsize{10.5}{13.2}\selectfont
\setlength{\LTleft}{0pt}\setlength{\LTright}{\fill}
\noindent
\begin{tabular}{@{}>{\raggedright\arraybackslash}p{40.25mm}>{\raggedright\arraybackslash}p{37.89mm}>{\raggedright\arraybackslash}p{40.25mm}@{}}
\toprule
\textbf{Genitive singular} & \textbf{Dative singular} & \textbf{Nominative plural} \\
\midrule
gehīersumnesse & gehīersumnesse & gehīersumnessa \\
\bottomrule
\end{tabular}
\par\endgroup

% mdtable:T31:end
\reserve{9\baselineskip}
% mdparagraph:P119:begin
\textbf{gehīersumnes; feminine noun}
% mdparagraph:P119:end

% table: NEW-U12-NOM-GEHIERSUMNES
% mdtable:T32:begin
\par\begingroup
\fontsize{10.5}{13.2}\selectfont
\setlength{\LTleft}{0pt}\setlength{\LTright}{\fill}
\noindent
\begin{tabular}{@{}>{\raggedright\arraybackslash}p{30.35mm}>{\raggedright\arraybackslash}p{42.49mm}>{\raggedright\arraybackslash}p{45.53mm}@{}}
\toprule
\textbf{Case} & \textbf{Singular} & \textbf{Plural} \\
\midrule
Nominative & gehīersumnes & gehīersumnessa \\
Genitive & gehīersumnesse & gehīersumnessa \\
Dative & gehīersumnesse & gehīersumnessum \\
Accusative & gehīersumnesse & gehīersumnessa \\
\bottomrule
\end{tabular}
\par\endgroup

% mdtable:T32:end
% mdparagraph:P120:begin
The entire plural column above is \textbf{reference only}. In the singular, obedience is an abstract thing promised.
% mdparagraph:P120:end

% teaching: NEW-U12-VOC-ANDWLITA
% teaching: NEW-U12-NOM-ANDWLITA
% diagnostic: NEW-U12-DIAG-ANDWLITA
% Page review repair: keep the lexical lead with its first display or derivation.
\reserve{8\baselineskip}
% mdparagraph:P121:begin
\textbf{andwlita} (masculine): face; countenance. Diagnostic stem: \textbf{andwlit}.
% mdparagraph:P121:end

% mdtable:T33:begin
\par\begingroup
\fontsize{10.5}{13.2}\selectfont
\setlength{\LTleft}{0pt}\setlength{\LTright}{\fill}
\noindent
\begin{tabular}{@{}>{\raggedright\arraybackslash}p{39.83mm}>{\raggedright\arraybackslash}p{36.98mm}>{\raggedright\arraybackslash}p{41.57mm}@{}}
\toprule
\textbf{Genitive singular} & \textbf{Dative singular} & \textbf{Nominative plural} \\
\midrule
andwlitan & andwlitan & andwlitan \\
\bottomrule
\end{tabular}
\par\endgroup

% mdtable:T33:end
\reserve{9\baselineskip}
% mdparagraph:P122:begin
\textbf{andwlita; masculine noun}
% mdparagraph:P122:end

% table: NEW-U12-NOM-ANDWLITA
% mdtable:T34:begin
\par\begingroup
\fontsize{10.5}{13.2}\selectfont
\setlength{\LTleft}{0pt}\setlength{\LTright}{\fill}
\noindent
\begin{tabular}{@{}>{\raggedright\arraybackslash}p{39.46mm}>{\raggedright\arraybackslash}p{39.46mm}>{\raggedright\arraybackslash}p{39.46mm}@{}}
\toprule
\textbf{Case} & \textbf{Singular} & \textbf{Plural} \\
\midrule
Nominative & andwlita & andwlitan \\
Genitive & andwlitan & andwlitena \\
Dative & andwlitan & andwlitum \\
Accusative & andwlitan & andwlitan \\
\bottomrule
\end{tabular}
\par\endgroup

% mdtable:T34:end
% teaching: NEW-U12-VOC-GEWRIT
% teaching: NEW-U12-NOM-GEWRIT
% diagnostic: NEW-U12-DIAG-GEWRIT
% Page review repair: keep the lexical lead with its first display or derivation.
\reserve{8\baselineskip}
% mdparagraph:P123:begin
\textbf{gewrit} (neuter): letter; written document. Diagnostic stem: \textbf{gewrit}.
% mdparagraph:P123:end

% mdtable:T35:begin
\par\begingroup
\fontsize{10.5}{13.2}\selectfont
\setlength{\LTleft}{0pt}\setlength{\LTright}{\fill}
\noindent
\begin{tabular}{@{}>{\raggedright\arraybackslash}p{39.24mm}>{\raggedright\arraybackslash}p{36.31mm}>{\raggedright\arraybackslash}p{42.82mm}@{}}
\toprule
\textbf{Genitive singular} & \textbf{Dative singular} & \textbf{Nominative plural} \\
\midrule
gewrites & gewrite & gewritu \\
\bottomrule
\end{tabular}
\par\endgroup

% mdtable:T35:end
\reserve{9\baselineskip}
% mdparagraph:P124:begin
\textbf{gewrit; neuter noun}
% mdparagraph:P124:end

% table: NEW-U12-NOM-GEWRIT
% mdtable:T36:begin
\par\begingroup
\fontsize{10.5}{13.2}\selectfont
\setlength{\LTleft}{0pt}\setlength{\LTright}{\fill}
\noindent
\begin{tabular}{@{}>{\raggedright\arraybackslash}p{39.46mm}>{\raggedright\arraybackslash}p{39.46mm}>{\raggedright\arraybackslash}p{39.46mm}@{}}
\toprule
\textbf{Case} & \textbf{Singular} & \textbf{Plural} \\
\midrule
Nominative & gewrit & gewritu \\
Genitive & gewrites & gewrita \\
Dative & gewrite & gewritum \\
Accusative & gewrit & gewritu \\
\bottomrule
\end{tabular}
\par\endgroup

% mdtable:T36:end
% teaching: NEW-U12-VOC-HERE
% teaching: NEW-U12-NOM-HERE
% diagnostic: NEW-U12-DIAG-HERE
% Page review repair: keep the lexical lead with its first display or derivation.
\reserve{8\baselineskip}
% mdparagraph:P125:begin
\textbf{hęre} (masculine): army. Diagnostic stem: \textbf{hęr}.
% mdparagraph:P125:end

% mdtable:T37:begin
\par\begingroup
\fontsize{10.5}{13.2}\selectfont
\setlength{\LTleft}{0pt}\setlength{\LTright}{\fill}
\noindent
\begin{tabular}{@{}>{\raggedright\arraybackslash}p{39.24mm}>{\raggedright\arraybackslash}p{36.31mm}>{\raggedright\arraybackslash}p{42.82mm}@{}}
\toprule
\textbf{Genitive singular} & \textbf{Dative singular} & \textbf{Nominative plural} \\
\midrule
hęres & hęre & hęras \\
\bottomrule
\end{tabular}
\par\endgroup

% mdtable:T37:end
\reserve{9\baselineskip}
% mdparagraph:P126:begin
\textbf{hęre; masculine noun}
% mdparagraph:P126:end

% table: NEW-U12-NOM-HERE
% mdtable:T38:begin
\par\begingroup
\fontsize{10.5}{13.2}\selectfont
\setlength{\LTleft}{0pt}\setlength{\LTright}{\fill}
\noindent
\begin{tabular}{@{}>{\raggedright\arraybackslash}p{39.46mm}>{\raggedright\arraybackslash}p{39.46mm}>{\raggedright\arraybackslash}p{39.46mm}@{}}
\toprule
\textbf{Case} & \textbf{Singular} & \textbf{Plural} \\
\midrule
Nominative & hęre & hęras \\
Genitive & hęres & hęra \\
Dative & hęre & hęrum \\
Accusative & hęre & hęras \\
\bottomrule
\end{tabular}
\par\endgroup

% mdtable:T38:end
% mdparagraph:P127:begin
The selected ordinary stem is hęr-. Cook also allows \textbf{hęrg-} or \textbf{hęrig-} before singular endings and \textbf{hęrg-, hęrig-} or \textbf{hęrige-} before plural endings. Thus hęres can have hęrges/\allowbreak{}hęriges; hęras can have hęrgas/\allowbreak{}hęrigas/\allowbreak{}hęrigeas. Apply each stem before the corresponding ending, including genitive and dative endings. These attested alternatives are accepted, not additional recall targets.
% mdparagraph:P127:end

% teaching: NEW-U12-VOC-SCIPHERE
% teaching: NEW-U12-NOM-SCIPHERE
% diagnostic: NEW-U12-DIAG-SCIPHERE
% Page review repair: keep the lexical lead with its first display or derivation.
\reserve{8\baselineskip}
% mdparagraph:P128:begin
\textbf{sciphęre} (masculine): fleet; naval force. Diagnostic stem: \textbf{sciphęr}.
% mdparagraph:P128:end

% mdtable:T39:begin
\par\begingroup
\fontsize{10.5}{13.2}\selectfont
\setlength{\LTleft}{0pt}\setlength{\LTright}{\fill}
\noindent
\begin{tabular}{@{}>{\raggedright\arraybackslash}p{40.42mm}>{\raggedright\arraybackslash}p{35.77mm}>{\raggedright\arraybackslash}p{42.19mm}@{}}
\toprule
\textbf{Genitive singular} & \textbf{Dative singular} & \textbf{Nominative plural} \\
\midrule
sciphęres & sciphęre & sciphęras \\
\bottomrule
\end{tabular}
\par\endgroup

% mdtable:T39:end
% mdparagraph:P129:begin
Use the complete \textbf{hęre; masculine noun} table. In every corresponding cell replace initial \textbf{hęr-} with \textbf{sciphęr-}, retaining the whole printed ending and the same case and number. An empty ending remains empty.
% mdparagraph:P129:end

% mdparagraph:P130:begin
The same alternatives occur in the compound: sciphęrg-/\allowbreak{}sciphęrig- before singular endings and sciphęrg-/\allowbreak{}sciphęrig-/\allowbreak{}sciphęrige- before plural endings. The selected complete model remains sciphęr- + the hęre endings.
% mdparagraph:P130:end

% teaching: NEW-U12-VOC-FEOND
% teaching: NEW-U12-NOM-FEOND
% diagnostic: NEW-U12-DIAG-FEOND
% Page review repair: keep the lexical lead with its first display or derivation.
\reserve{8\baselineskip}
% mdparagraph:P131:begin
\textbf{fēond} (masculine): enemy; foe. Diagnostic stem: \textbf{fēond /\allowbreak{} fīend}.
% mdparagraph:P131:end

% mdtable:T40:begin
\par\begingroup
\fontsize{10.5}{13.2}\selectfont
\setlength{\LTleft}{0pt}\setlength{\LTright}{\fill}
\noindent
\begin{tabular}{@{}>{\raggedright\arraybackslash}p{39.24mm}>{\raggedright\arraybackslash}p{36.31mm}>{\raggedright\arraybackslash}p{42.82mm}@{}}
\toprule
\textbf{Genitive singular} & \textbf{Dative singular} & \textbf{Nominative plural} \\
\midrule
fēondes & fēonde & fīend \\
\bottomrule
\end{tabular}
\par\endgroup

% mdtable:T40:end
\reserve{9\baselineskip}
% mdparagraph:P132:begin
\textbf{fēond; masculine noun}
% mdparagraph:P132:end

% table: NEW-U12-NOM-FEOND
% mdtable:T41:begin
\par\begingroup
\fontsize{10.5}{13.2}\selectfont
\setlength{\LTleft}{0pt}\setlength{\LTright}{\fill}
\noindent
\begin{tabular}{@{}>{\raggedright\arraybackslash}p{39.46mm}>{\raggedright\arraybackslash}p{39.46mm}>{\raggedright\arraybackslash}p{39.46mm}@{}}
\toprule
\textbf{Case} & \textbf{Singular} & \textbf{Plural} \\
\midrule
Nominative & fēond & fīend \\
Genitive & fēondes & fēonda \\
Dative & fēonde & fēondum \\
Accusative & fēond & fīend \\
\bottomrule
\end{tabular}
\par\endgroup

% mdtable:T41:end
% mdparagraph:P133:begin
Accepted alternatives are \textbf{fīend} in the dative singular and \textbf{fēond} or \textbf{fēondas} in the nominative/\allowbreak{}accusative plural. Learn selected genitive plural \textbf{fēonda}; it does not take the -endra of the warrior family.
% mdparagraph:P133:end

% teaching: NEW-U12-VOC-HAERFEST
% teaching: NEW-U12-NOM-HAERFEST
% diagnostic: NEW-U12-DIAG-HAERFEST
% Page review repair: keep the lexical lead with its first display or derivation.
\reserve{8\baselineskip}
% mdparagraph:P134:begin
\textbf{hærfest} (masculine): harvest. Diagnostic stem: \textbf{hærfest}.
% mdparagraph:P134:end

% mdtable:T42:begin
\par\begingroup
\fontsize{10.5}{13.2}\selectfont
\setlength{\LTleft}{0pt}\setlength{\LTright}{\fill}
\noindent
\begin{tabular}{@{}>{\raggedright\arraybackslash}p{40.42mm}>{\raggedright\arraybackslash}p{35.77mm}>{\raggedright\arraybackslash}p{42.19mm}@{}}
\toprule
\textbf{Genitive singular} & \textbf{Dative singular} & \textbf{Nominative plural} \\
\midrule
hærfestes & hærfeste & hærfestas \\
\bottomrule
\end{tabular}
\par\endgroup

% mdtable:T42:end
\reserve{9\baselineskip}
% mdparagraph:P135:begin
\textbf{hærfest; masculine noun}
% mdparagraph:P135:end

% table: NEW-U12-NOM-HAERFEST
% mdtable:T43:begin
\par\begingroup
\fontsize{10.5}{13.2}\selectfont
\setlength{\LTleft}{0pt}\setlength{\LTright}{\fill}
\noindent
\begin{tabular}{@{}>{\raggedright\arraybackslash}p{39.46mm}>{\raggedright\arraybackslash}p{39.46mm}>{\raggedright\arraybackslash}p{39.46mm}@{}}
\toprule
\textbf{Case} & \textbf{Singular} & \textbf{Plural} \\
\midrule
Nominative & hærfest & hærfestas \\
Genitive & hærfestes & hærfesta \\
Dative & hærfeste & hærfestum \\
Accusative & hærfest & hærfestas \\
\bottomrule
\end{tabular}
\par\endgroup

% mdtable:T43:end
% teaching: NEW-U12-VOC-SAEDTIMA
% teaching: NEW-U12-NOM-SAEDTIMA
% diagnostic: NEW-U12-DIAG-SAEDTIMA
% Page review repair: keep the lexical lead with its first display or derivation.
\reserve{8\baselineskip}
% mdparagraph:P136:begin
\textbf{sǣdtīma} (masculine): seedtime; sowing season. Diagnostic stem: \textbf{sǣdtīm}.
% mdparagraph:P136:end

% mdtable:T44:begin
\par\begingroup
\fontsize{10.5}{13.2}\selectfont
\setlength{\LTleft}{0pt}\setlength{\LTright}{\fill}
\noindent
\begin{tabular}{@{}>{\raggedright\arraybackslash}p{39.24mm}>{\raggedright\arraybackslash}p{36.31mm}>{\raggedright\arraybackslash}p{42.82mm}@{}}
\toprule
\textbf{Genitive singular} & \textbf{Dative singular} & \textbf{Nominative plural} \\
\midrule
sǣdtīman & sǣdtīman & sǣdtīman \\
\bottomrule
\end{tabular}
\par\endgroup

% mdtable:T44:end
% mdparagraph:P137:begin
Use the complete \textbf{andwlita; masculine noun} table. In every corresponding cell replace initial \textbf{andwlit-} with \textbf{sǣdtīm-}, retaining the whole printed ending and the same case and number. An empty ending remains empty.
% mdparagraph:P137:end

% teaching: NEW-U12-VOC-WAEPEN
% teaching: NEW-U12-NOM-WAEPEN
% diagnostic: NEW-U12-DIAG-WAEPEN
% Page review repair: keep the lexical lead with its first display or derivation.
\reserve{8\baselineskip}
% mdparagraph:P138:begin
\textbf{wǣpen} (neuter): weapon. Diagnostic stem: \textbf{wǣpn}.
% mdparagraph:P138:end

% mdtable:T45:begin
\par\begingroup
\fontsize{10.5}{13.2}\selectfont
\setlength{\LTleft}{0pt}\setlength{\LTright}{\fill}
\noindent
\begin{tabular}{@{}>{\raggedright\arraybackslash}p{39.24mm}>{\raggedright\arraybackslash}p{36.31mm}>{\raggedright\arraybackslash}p{42.82mm}@{}}
\toprule
\textbf{Genitive singular} & \textbf{Dative singular} & \textbf{Nominative plural} \\
\midrule
wǣpnes & wǣpne & wǣpnu \\
\bottomrule
\end{tabular}
\par\endgroup

% mdtable:T45:end
\reserve{9\baselineskip}
% mdparagraph:P139:begin
\textbf{wǣpen; neuter noun}
% mdparagraph:P139:end

% table: NEW-U12-NOM-WAEPEN
% mdtable:T46:begin
\par\begingroup
\fontsize{10.5}{13.2}\selectfont
\setlength{\LTleft}{0pt}\setlength{\LTright}{\fill}
\noindent
\begin{tabular}{@{}>{\raggedright\arraybackslash}p{39.46mm}>{\raggedright\arraybackslash}p{39.46mm}>{\raggedright\arraybackslash}p{39.46mm}@{}}
\toprule
\textbf{Case} & \textbf{Singular} & \textbf{Plural} \\
\midrule
Nominative & wǣpen & wǣpnu \\
Genitive & wǣpnes & wǣpna \\
Dative & wǣpne & wǣpnum \\
Accusative & wǣpen & wǣpnu \\
\bottomrule
\end{tabular}
\par\endgroup

% mdtable:T46:end
% mdparagraph:P140:begin
The selected nominative/\allowbreak{}accusative plural is \textbf{wǣpnu}. Cook also records \textbf{wǣpen} and \textbf{wǣpn} in these positions. Retain the shortened stem wǣpn- before the other endings.
% mdparagraph:P140:end

% teaching: NEW-U12-VOC-WULF
% teaching: NEW-U12-NOM-WULF
% diagnostic: NEW-U12-DIAG-WULF
% Page review repair: keep the lexical lead with its first display or derivation.
\reserve{8\baselineskip}
% mdparagraph:P141:begin
\textbf{wulf} (masculine): wolf. Diagnostic stem: \textbf{wulf}.
% mdparagraph:P141:end

% mdtable:T47:begin
\par\begingroup
\fontsize{10.5}{13.2}\selectfont
\setlength{\LTleft}{0pt}\setlength{\LTright}{\fill}
\noindent
\begin{tabular}{@{}>{\raggedright\arraybackslash}p{39.24mm}>{\raggedright\arraybackslash}p{36.31mm}>{\raggedright\arraybackslash}p{42.82mm}@{}}
\toprule
\textbf{Genitive singular} & \textbf{Dative singular} & \textbf{Nominative plural} \\
\midrule
wulfes & wulfe & wulfas \\
\bottomrule
\end{tabular}
\par\endgroup

% mdtable:T47:end
% mdparagraph:P142:begin
Use the complete \textbf{hærfest; masculine noun} table. In every corresponding cell replace initial \textbf{hærfest-} with \textbf{wulf-}, retaining the whole printed ending and the same case and number. An empty ending remains empty.
% mdparagraph:P142:end

% teaching: NEW-U12-VOC-GEMAERE
% teaching: NEW-U12-NOM-GEMAERE
% diagnostic: NEW-U12-DIAG-GEMAERE
% Page review repair: keep the lexical lead with its first display or derivation.
\reserve{8\baselineskip}
% mdparagraph:P143:begin
\textbf{gemǣre} (neuter): boundary; border. Diagnostic stem: \textbf{gemǣr}.
% mdparagraph:P143:end

% mdtable:T48:begin
\par\begingroup
\fontsize{10.5}{13.2}\selectfont
\setlength{\LTleft}{0pt}\setlength{\LTright}{\fill}
\noindent
\begin{tabular}{@{}>{\raggedright\arraybackslash}p{39.24mm}>{\raggedright\arraybackslash}p{36.31mm}>{\raggedright\arraybackslash}p{42.82mm}@{}}
\toprule
\textbf{Genitive singular} & \textbf{Dative singular} & \textbf{Nominative plural} \\
\midrule
gemǣres & gemǣre & gemǣru \\
\bottomrule
\end{tabular}
\par\endgroup

% mdtable:T48:end
\reserve{9\baselineskip}
% mdparagraph:P144:begin
\textbf{gemǣre; neuter noun}
% mdparagraph:P144:end

% table: NEW-U12-NOM-GEMAERE
% mdtable:T49:begin
\par\begingroup
\fontsize{10.5}{13.2}\selectfont
\setlength{\LTleft}{0pt}\setlength{\LTright}{\fill}
\noindent
\begin{tabular}{@{}>{\raggedright\arraybackslash}p{39.46mm}>{\raggedright\arraybackslash}p{39.46mm}>{\raggedright\arraybackslash}p{39.46mm}@{}}
\toprule
\textbf{Case} & \textbf{Singular} & \textbf{Plural} \\
\midrule
Nominative & gemǣre & gemǣru \\
Genitive & gemǣres & gemǣra \\
Dative & gemǣre & gemǣrum \\
Accusative & gemǣre & gemǣru \\
\bottomrule
\end{tabular}
\par\endgroup

% mdtable:T49:end
% teaching: NEW-U12-VOC-PENING
% teaching: NEW-U12-NOM-PENING
% diagnostic: NEW-U12-DIAG-PENING
% Page review repair: keep the lexical lead with its first display or derivation.
\reserve{8\baselineskip}
% mdparagraph:P145:begin
\textbf{pęning} (masculine): penny; coin. Diagnostic stem: \textbf{pęning}.
% mdparagraph:P145:end

% mdtable:T50:begin
\par\begingroup
\fontsize{10.5}{13.2}\selectfont
\setlength{\LTleft}{0pt}\setlength{\LTright}{\fill}
\noindent
\begin{tabular}{@{}>{\raggedright\arraybackslash}p{39.24mm}>{\raggedright\arraybackslash}p{36.31mm}>{\raggedright\arraybackslash}p{42.82mm}@{}}
\toprule
\textbf{Genitive singular} & \textbf{Dative singular} & \textbf{Nominative plural} \\
\midrule
pęninges & pęninge & pęningas \\
\bottomrule
\end{tabular}
\par\endgroup

% mdtable:T50:end
% mdparagraph:P146:begin
Use the complete \textbf{hærfest; masculine noun} table. In every corresponding cell replace initial \textbf{hærfest-} with \textbf{pęning-}, retaining the whole printed ending and the same case and number. An empty ending remains empty.
% mdparagraph:P146:end

% teaching: NEW-U12-VOC-MUTH
% teaching: NEW-U12-NOM-MUTH
% diagnostic: NEW-U12-DIAG-MUTH
% Page review repair: keep the lexical lead with its first display or derivation.
\reserve{8\baselineskip}
% mdparagraph:P147:begin
\textbf{mūð} (masculine): mouth. Diagnostic stem: \textbf{mūð}.
% mdparagraph:P147:end

% mdtable:T51:begin
\par\begingroup
\fontsize{10.5}{13.2}\selectfont
\setlength{\LTleft}{0pt}\setlength{\LTright}{\fill}
\noindent
\begin{tabular}{@{}>{\raggedright\arraybackslash}p{39.24mm}>{\raggedright\arraybackslash}p{36.31mm}>{\raggedright\arraybackslash}p{42.82mm}@{}}
\toprule
\textbf{Genitive singular} & \textbf{Dative singular} & \textbf{Nominative plural} \\
\midrule
mūðes & mūðe & mūðas \\
\bottomrule
\end{tabular}
\par\endgroup

% mdtable:T51:end
% mdparagraph:P148:begin
Use the complete \textbf{hærfest; masculine noun} table. In every corresponding cell replace initial \textbf{hærfest-} with \textbf{mūð-}, retaining the whole printed ending and the same case and number. An empty ending remains empty.
% mdparagraph:P148:end

% teaching: NEW-U12-VOC-FAEMNE
% teaching: NEW-U12-NOM-FAEMNE
% diagnostic: NEW-U12-DIAG-FAEMNE
% Page review repair: keep the lexical lead with its first display or derivation.
\reserve{8\baselineskip}
% mdparagraph:P149:begin
\textbf{fǣmne} (feminine): maiden; young woman. Diagnostic stem: \textbf{fǣmn}.
% mdparagraph:P149:end

% mdtable:T52:begin
\par\begingroup
\fontsize{10.5}{13.2}\selectfont
\setlength{\LTleft}{0pt}\setlength{\LTright}{\fill}
\noindent
\begin{tabular}{@{}>{\raggedright\arraybackslash}p{39.24mm}>{\raggedright\arraybackslash}p{36.31mm}>{\raggedright\arraybackslash}p{42.82mm}@{}}
\toprule
\textbf{Genitive singular} & \textbf{Dative singular} & \textbf{Nominative plural} \\
\midrule
fǣmnan & fǣmnan & fǣmnan \\
\bottomrule
\end{tabular}
\par\endgroup

% mdtable:T52:end
\reserve{9\baselineskip}
% mdparagraph:P150:begin
\textbf{fǣmne; feminine noun}
% mdparagraph:P150:end

% table: NEW-U12-NOM-FAEMNE
% mdtable:T53:begin
\par\begingroup
\fontsize{10.5}{13.2}\selectfont
\setlength{\LTleft}{0pt}\setlength{\LTright}{\fill}
\noindent
\begin{tabular}{@{}>{\raggedright\arraybackslash}p{39.46mm}>{\raggedright\arraybackslash}p{39.46mm}>{\raggedright\arraybackslash}p{39.46mm}@{}}
\toprule
\textbf{Case} & \textbf{Singular} & \textbf{Plural} \\
\midrule
Nominative & fǣmne & fǣmnan \\
Genitive & fǣmnan & fǣmnena \\
Dative & fǣmnan & fǣmnum \\
Accusative & fǣmnan & fǣmnan \\
\bottomrule
\end{tabular}
\par\endgroup

% mdtable:T53:end
% teaching: NEW-U12-VOC-STANWEALL
% teaching: NEW-U12-NOM-STANWEALL
% diagnostic: NEW-U12-DIAG-STANWEALL
% Page review repair: keep the lexical lead with its first display or derivation.
\reserve{8\baselineskip}
% mdparagraph:P151:begin
\textbf{stānweall} (masculine): stone wall. Diagnostic stem: \textbf{stānweall}.
% mdparagraph:P151:end

% mdtable:T54:begin
\par\begingroup
\fontsize{10.5}{13.2}\selectfont
\setlength{\LTleft}{0pt}\setlength{\LTright}{\fill}
\noindent
\begin{tabular}{@{}>{\raggedright\arraybackslash}p{40.9mm}>{\raggedright\arraybackslash}p{36.57mm}>{\raggedright\arraybackslash}p{40.9mm}@{}}
\toprule
\textbf{Genitive singular} & \textbf{Dative singular} & \textbf{Nominative plural} \\
\midrule
stānwealles & stānwealle & stānweallas \\
\bottomrule
\end{tabular}
\par\endgroup

% mdtable:T54:end
% mdparagraph:P152:begin
Use the complete \textbf{hærfest; masculine noun} table. In every corresponding cell replace initial \textbf{hærfest-} with \textbf{stānweall-}, retaining the whole printed ending and the same case and number. An empty ending remains empty.
% mdparagraph:P152:end

% teaching: NEW-U12-VOC-LAND
% teaching: NEW-U12-NOM-LAND
% diagnostic: NEW-U12-DIAG-LAND
% Page review repair: keep the lexical lead with its first display or derivation.
\reserve{8\baselineskip}
% mdparagraph:P153:begin
\textbf{land} (neuter): land; country. Diagnostic stem: \textbf{land}.
% mdparagraph:P153:end

% mdtable:T55:begin
\par\begingroup
\fontsize{10.5}{13.2}\selectfont
\setlength{\LTleft}{0pt}\setlength{\LTright}{\fill}
\noindent
\begin{tabular}{@{}>{\raggedright\arraybackslash}p{39.24mm}>{\raggedright\arraybackslash}p{36.31mm}>{\raggedright\arraybackslash}p{42.82mm}@{}}
\toprule
\textbf{Genitive singular} & \textbf{Dative singular} & \textbf{Nominative plural} \\
\midrule
landes & lande & land \\
\bottomrule
\end{tabular}
\par\endgroup

% mdtable:T55:end
\reserve{9\baselineskip}
% mdparagraph:P154:begin
\textbf{land; neuter noun}
% mdparagraph:P154:end

% table: NEW-U12-NOM-LAND
% mdtable:T56:begin
\par\begingroup
\fontsize{10.5}{13.2}\selectfont
\setlength{\LTleft}{0pt}\setlength{\LTright}{\fill}
\noindent
\begin{tabular}{@{}>{\raggedright\arraybackslash}p{39.46mm}>{\raggedright\arraybackslash}p{39.46mm}>{\raggedright\arraybackslash}p{39.46mm}@{}}
\toprule
\textbf{Case} & \textbf{Singular} & \textbf{Plural} \\
\midrule
Nominative & land & land \\
Genitive & landes & landa \\
Dative & lande & landum \\
Accusative & land & land \\
\bottomrule
\end{tabular}
\par\endgroup

% mdtable:T56:end
% teaching: NEW-U12-VOC-FOT
% teaching: NEW-U12-NOM-FOT
% diagnostic: NEW-U12-DIAG-FOT
% Page review repair: keep the lexical lead with its first display or derivation.
\reserve{8\baselineskip}
% mdparagraph:P155:begin
\textbf{fōt} (masculine): foot; also a foot as a measure of length. Diagnostic stem: \textbf{fōt /\allowbreak{} fēt}.
% mdparagraph:P155:end

% mdtable:T57:begin
\par\begingroup
\fontsize{10.5}{13.2}\selectfont
\setlength{\LTleft}{0pt}\setlength{\LTright}{\fill}
\noindent
\begin{tabular}{@{}>{\raggedright\arraybackslash}p{39.24mm}>{\raggedright\arraybackslash}p{36.31mm}>{\raggedright\arraybackslash}p{42.82mm}@{}}
\toprule
\textbf{Genitive singular} & \textbf{Dative singular} & \textbf{Nominative plural} \\
\midrule
fōtes & fēt & fēt \\
\bottomrule
\end{tabular}
\par\endgroup

% mdtable:T57:end
\reserve{9\baselineskip}
% mdparagraph:P156:begin
\textbf{fōt; masculine noun}
% mdparagraph:P156:end

% table: NEW-U12-NOM-FOT
% mdtable:T58:begin
\par\begingroup
\fontsize{10.5}{13.2}\selectfont
\setlength{\LTleft}{0pt}\setlength{\LTright}{\fill}
\noindent
\begin{tabular}{@{}>{\raggedright\arraybackslash}p{39.46mm}>{\raggedright\arraybackslash}p{39.46mm}>{\raggedright\arraybackslash}p{39.46mm}@{}}
\toprule
\textbf{Case} & \textbf{Singular} & \textbf{Plural} \\
\midrule
Nominative & fōt & fēt \\
Genitive & fōtes & fōta \\
Dative & fēt & fōtum \\
Accusative & fōt & fēt \\
\bottomrule
\end{tabular}
\par\endgroup

% mdtable:T58:end
% mdparagraph:P157:begin
Accepted dative singular \textbf{fōte} and occasional nominative/\allowbreak{}accusative plural \textbf{fōtas} occur beside the selected fēt. Use \textbf{fōta} for the genitive plural measure expression taught in §8.
% mdparagraph:P157:end

% teaching: NEW-U12-VOC-ETHEL
% teaching: NEW-U12-NOM-ETHEL
% diagnostic: NEW-U12-DIAG-ETHEL
% Page review repair: keep the lexical lead with its first display or derivation.
\reserve{8\baselineskip}
% mdparagraph:P158:begin
\textbf{ēðel} (masculine): native land; homeland. Diagnostic stem: \textbf{ēðl}.
% mdparagraph:P158:end

% mdtable:T59:begin
\par\begingroup
\fontsize{10.5}{13.2}\selectfont
\setlength{\LTleft}{0pt}\setlength{\LTright}{\fill}
\noindent
\begin{tabular}{@{}>{\raggedright\arraybackslash}p{39.24mm}>{\raggedright\arraybackslash}p{36.31mm}>{\raggedright\arraybackslash}p{42.82mm}@{}}
\toprule
\textbf{Genitive singular} & \textbf{Dative singular} & \textbf{Nominative plural} \\
\midrule
ēðles & ēðle & ēðlas \\
\bottomrule
\end{tabular}
\par\endgroup

% mdtable:T59:end
\reserve{9\baselineskip}
% mdparagraph:P159:begin
\textbf{ēðel; masculine noun}
% mdparagraph:P159:end

% table: NEW-U12-NOM-ETHEL
% mdtable:T60:begin
\par\begingroup
\fontsize{10.5}{13.2}\selectfont
\setlength{\LTleft}{0pt}\setlength{\LTright}{\fill}
\noindent
\begin{tabular}{@{}>{\raggedright\arraybackslash}p{39.46mm}>{\raggedright\arraybackslash}p{39.46mm}>{\raggedright\arraybackslash}p{39.46mm}@{}}
\toprule
\textbf{Case} & \textbf{Singular} & \textbf{Plural} \\
\midrule
Nominative & ēðel & ēðlas \\
Genitive & ēðles & ēðla \\
Dative & ēðle & ēðlum \\
Accusative & ēðel & ēðlas \\
\bottomrule
\end{tabular}
\par\endgroup

% mdtable:T60:end
% teaching: NEW-U12-VOC-HEORTE
% teaching: NEW-U12-NOM-HEORTE
% diagnostic: NEW-U12-DIAG-HEORTE
% Page review repair: keep the lexical lead with its first display or derivation.
\reserve{8\baselineskip}
% mdparagraph:P160:begin
\textbf{heorte} (feminine): heart. Diagnostic stem: \textbf{heort}.
% mdparagraph:P160:end

% mdtable:T61:begin
\par\begingroup
\fontsize{10.5}{13.2}\selectfont
\setlength{\LTleft}{0pt}\setlength{\LTright}{\fill}
\noindent
\begin{tabular}{@{}>{\raggedright\arraybackslash}p{39.24mm}>{\raggedright\arraybackslash}p{36.31mm}>{\raggedright\arraybackslash}p{42.82mm}@{}}
\toprule
\textbf{Genitive singular} & \textbf{Dative singular} & \textbf{Nominative plural} \\
\midrule
heortan & heortan & heortan \\
\bottomrule
\end{tabular}
\par\endgroup

% mdtable:T61:end
% mdparagraph:P161:begin
Use the complete \textbf{fǣmne; feminine noun} table. In every corresponding cell replace initial \textbf{fǣmn-} with \textbf{heort-}, retaining the whole printed ending and the same case and number. An empty ending remains empty.
% mdparagraph:P161:end

% teaching: NEW-U12-VOC-REST
% teaching: NEW-U12-NOM-REST
% diagnostic: NEW-U12-DIAG-REST
% Page review repair: keep the lexical lead with its first display or derivation.
\reserve{8\baselineskip}
% mdparagraph:P162:begin
\textbf{ręst} (feminine): rest; repose. Diagnostic stem: \textbf{ręst}.
% mdparagraph:P162:end

% mdtable:T62:begin
\par\begingroup
\fontsize{10.5}{13.2}\selectfont
\setlength{\LTleft}{0pt}\setlength{\LTright}{\fill}
\noindent
\begin{tabular}{@{}>{\raggedright\arraybackslash}p{39.24mm}>{\raggedright\arraybackslash}p{36.31mm}>{\raggedright\arraybackslash}p{42.82mm}@{}}
\toprule
\textbf{Genitive singular} & \textbf{Dative singular} & \textbf{Nominative plural} \\
\midrule
ręste & ręste & ręsta \\
\bottomrule
\end{tabular}
\par\endgroup

% mdtable:T62:end
% mdparagraph:P163:begin
Use the complete \textbf{brycg; feminine noun} table. In every corresponding cell replace initial \textbf{brycg-} with \textbf{ręst-}, retaining the whole printed ending and the same case and number. An empty ending remains empty.
% mdparagraph:P163:end

% mdparagraph:P164:begin
All derived plural forms of ręst are \textbf{reference only}. The assigned productive sense is singular “rest; repose”, supported by Cook VI.21.9. It is not restricted to the vocabulary’s separate “couch; bed” gloss.
% mdparagraph:P164:end

% teaching: NEW-U12-VOC-SAWOL
% teaching: NEW-U12-NOM-SAWOL
% diagnostic: NEW-U12-DIAG-SAWOL
% Page review repair: keep the lexical lead with its first display or derivation.
\reserve{8\baselineskip}
% mdparagraph:P165:begin
\textbf{sāwol} (feminine): soul. Diagnostic stem: \textbf{sāwl}.
% mdparagraph:P165:end

% mdtable:T63:begin
\par\begingroup
\fontsize{10.5}{13.2}\selectfont
\setlength{\LTleft}{0pt}\setlength{\LTright}{\fill}
\noindent
\begin{tabular}{@{}>{\raggedright\arraybackslash}p{39.24mm}>{\raggedright\arraybackslash}p{36.31mm}>{\raggedright\arraybackslash}p{42.82mm}@{}}
\toprule
\textbf{Genitive singular} & \textbf{Dative singular} & \textbf{Nominative plural} \\
\midrule
sāwle & sāwle & sāwla \\
\bottomrule
\end{tabular}
\par\endgroup

% mdtable:T63:end
\reserve{9\baselineskip}
% mdparagraph:P166:begin
\textbf{sāwol; feminine noun}
% mdparagraph:P166:end

% table: NEW-U12-NOM-SAWOL
% mdtable:T64:begin
\par\begingroup
\fontsize{10.5}{13.2}\selectfont
\setlength{\LTleft}{0pt}\setlength{\LTright}{\fill}
\noindent
\begin{tabular}{@{}>{\raggedright\arraybackslash}p{39.46mm}>{\raggedright\arraybackslash}p{39.46mm}>{\raggedright\arraybackslash}p{39.46mm}@{}}
\toprule
\textbf{Case} & \textbf{Singular} & \textbf{Plural} \\
\midrule
Nominative & sāwol & sāwla \\
Genitive & sāwle & sāwla \\
Dative & sāwle & sāwlum \\
Accusative & sāwle & sāwla \\
\bottomrule
\end{tabular}
\par\endgroup

% mdtable:T64:end
% mdparagraph:P167:begin
Cook also gives the citation form \textbf{sāwl}. The selected headword is sāwol; ending-bearing forms use sāwl-.
% mdparagraph:P167:end

\reserve{7\baselineskip}
% mdparagraph:P168:begin
\section*{7. Adjectives: complete strong and weak patterns}\BookHeading{nav-ch12-learner-0010}{subsection}{7. Adjectives: complete strong and weak patterns}{1}
% mdparagraph:P168:end

% teaching: NEW-U12-ADJECTIVE-DIAGNOSTICS
% mdparagraph:P169:begin
An ordinary positive predicate adjective is strong and nominative, agreeing with its subject. An attributive adjective is weak after the article or demonstrative, strong in the ordinary bare phrase. Agree with the noun, not with the possessor. In direct address the noun has its nominative shape and an accompanying adjective is weak.
% mdparagraph:P169:end

% mdparagraph:P170:begin
Brād, beald, earg, ēaðmōd and mildheort use the complete long strong model and the weak model. Ēadig and the adjective īsern use the selected disyllabic short-model endings. Each model below contains all three genders in both numbers. Distinct instrumental singular cells are \textbf{reference only}.
% mdparagraph:P170:end

% teaching: NEW-U12-VOC-BRAD
% teaching: NEW-U12-ADJ-BRAD-STRONG
% teaching: NEW-U12-ADJ-BRAD-WEAK
% diagnostic: NEW-U12-DIAG-BRAD
% Page review repair: keep the lexical lead with its first display or derivation.
\reserve{24\baselineskip}
% mdparagraph:P171:begin
\textbf{brād} (adjective): broad. Diagnostic stem: \textbf{brād}.
% mdparagraph:P171:end

\reserve{18\baselineskip}
% mdparagraph:P172:begin
\textbf{brād; strong adjective}
% mdparagraph:P172:end

% table: NEW-U12-ADJ-BRAD-STRONG
% mdtable:T65:begin
% course-table-style: 2026-10
% course-panel:T65-P01:begin
\par\begingroup
\fontsize{10.5}{13.2}\selectfont
\setlength{\LTleft}{0pt}\setlength{\LTright}{\fill}
\noindent
\begin{tabular}{@{}>{\raggedright\arraybackslash}p{22.0mm}>{\raggedright\arraybackslash}p{31.188mm}>{\raggedright\arraybackslash}p{31.188mm}>{\raggedright\arraybackslash}p{31.188mm}@{}}
\toprule
\textbf{Case} & \textbf{Masculine singular} & \textbf{Neuter singular} & \textbf{Feminine singular} \\
\midrule
Nominative & brād & brād & brād \\
Genitive & brādes & brādes & brādre \\
Dative & brādum & brādum & brādre \\
Accusative & brādne & brād & brāde \\
Instrumental & brāde & brāde & as dative \\
\bottomrule
\end{tabular}
\par\endgroup
% course-panel:T65-P01:end
\par\addvspace{5pt}
% course-panel:T65-P02:begin
\par\begingroup
\fontsize{10.5}{13.2}\selectfont
\setlength{\LTleft}{0pt}\setlength{\LTright}{\fill}
\noindent
\begin{tabular}{@{}>{\raggedright\arraybackslash}p{22.0mm}>{\raggedright\arraybackslash}p{31.188mm}>{\raggedright\arraybackslash}p{31.188mm}>{\raggedright\arraybackslash}p{31.188mm}@{}}
\toprule
\textbf{Case} & \textbf{Masculine plural} & \textbf{Neuter plural} & \textbf{Feminine plural} \\
\midrule
Nominative & brāde & brād & brāda \\
Genitive & brādra & brādra & brādra \\
Dative & brādum & brādum & brādum \\
Accusative & brāde & brād & brāda \\
Instrumental & as dative & as dative & as dative \\
\bottomrule
\end{tabular}
\par\endgroup
% course-panel:T65-P02:end
% mdtable:T65:end

\reserve{18\baselineskip}
% mdparagraph:P173:begin
\textbf{brād; weak adjective}
% mdparagraph:P173:end

% table: NEW-U12-ADJ-BRAD-WEAK
% mdtable:T66:begin
% course-table-style: 2026-10
% course-panel:T66-P01:begin
\par\begingroup
\fontsize{10.5}{13.2}\selectfont
\setlength{\LTleft}{0pt}\setlength{\LTright}{\fill}
\noindent
\begin{tabular}{@{}>{\raggedright\arraybackslash}p{22.0mm}>{\raggedright\arraybackslash}p{31.188mm}>{\raggedright\arraybackslash}p{31.188mm}>{\raggedright\arraybackslash}p{31.188mm}@{}}
\toprule
\textbf{Case} & \textbf{Masculine singular} & \textbf{Neuter singular} & \textbf{Feminine singular} \\
\midrule
Nominative & brāda & brāde & brāde \\
Genitive & brādan & brādan & brādan \\
Dative & brādan & brādan & brādan \\
Accusative & brādan & brāde & brādan \\
Instrumental & brādan & brādan & as dative \\
\bottomrule
\end{tabular}
\par\endgroup
% course-panel:T66-P01:end
\par\addvspace{5pt}
% course-panel:T66-P02:begin
\par\begingroup
\fontsize{10.5}{13.2}\selectfont
\setlength{\LTleft}{0pt}\setlength{\LTright}{\fill}
\noindent
\begin{tabular}{@{}>{\raggedright\arraybackslash}p{22.0mm}>{\raggedright\arraybackslash}p{31.188mm}>{\raggedright\arraybackslash}p{31.188mm}>{\raggedright\arraybackslash}p{31.188mm}@{}}
\toprule
\textbf{Case} & \textbf{Masculine plural} & \textbf{Neuter plural} & \textbf{Feminine plural} \\
\midrule
Nominative & brādan & brādan & brādan \\
Genitive & brādra & brādra & brādra \\
Dative & brādum & brādum & brādum \\
Accusative & brādan & brādan & brādan \\
Instrumental & as dative & as dative & as dative \\
\bottomrule
\end{tabular}
\par\endgroup
% course-panel:T66-P02:end
% mdtable:T66:end

% teaching: NEW-U12-VOC-BEALD
% teaching: NEW-U12-ADJ-BEALD-STRONG
% teaching: NEW-U12-ADJ-BEALD-WEAK
% diagnostic: NEW-U12-DIAG-BEALD
% Page review repair: keep the lexical lead with its first display or derivation.
\reserve{10\baselineskip}
% mdparagraph:P174:begin
\textbf{beald} (adjective): bold; brave. Diagnostic stem: \textbf{beald}.
% mdparagraph:P174:end

% mdparagraph:P175:begin
Use the complete \textbf{brād; strong adjective} table. In every corresponding cell replace initial \textbf{brād-} with \textbf{beald-}, retaining the whole printed ending and the same case, number, gender and declension. An empty ending remains empty.
% mdparagraph:P175:end

% mdparagraph:P176:begin
Use the complete \textbf{brād; weak adjective} table. In every corresponding cell replace initial \textbf{brād-} with \textbf{beald-}, retaining the whole printed ending and the same case, number, gender and declension. An empty ending remains empty.
% mdparagraph:P176:end

% teaching: NEW-U12-VOC-EARG
% teaching: NEW-U12-ADJ-EARG-STRONG
% teaching: NEW-U12-ADJ-EARG-WEAK
% diagnostic: NEW-U12-DIAG-EARG
% Page review repair: keep the lexical lead with its first display or derivation.
\reserve{10\baselineskip}
% mdparagraph:P177:begin
\textbf{earg} (adjective): cowardly. Diagnostic stem: \textbf{earg}.
% mdparagraph:P177:end

% mdparagraph:P178:begin
Use the complete \textbf{brād; strong adjective} table. In every corresponding cell replace initial \textbf{brād-} with \textbf{earg-}, retaining the whole printed ending and the same case, number, gender and declension. An empty ending remains empty.
% mdparagraph:P178:end

% mdparagraph:P179:begin
Use the complete \textbf{brād; weak adjective} table. In every corresponding cell replace initial \textbf{brād-} with \textbf{earg-}, retaining the whole printed ending and the same case, number, gender and declension. An empty ending remains empty.
% mdparagraph:P179:end

% teaching: NEW-U12-VOC-EADMOD
% teaching: NEW-U12-ADJ-EADMOD-STRONG
% teaching: NEW-U12-ADJ-EADMOD-WEAK
% diagnostic: NEW-U12-DIAG-EADMOD
% Page review repair: keep the lexical lead with its first display or derivation.
\reserve{10\baselineskip}
% mdparagraph:P180:begin
\textbf{ēaðmōd} (adjective): lowly; humble. Diagnostic stem: \textbf{ēaðmōd}.
% mdparagraph:P180:end

% mdparagraph:P181:begin
Use the complete \textbf{brād; strong adjective} table. In every corresponding cell replace initial \textbf{brād-} with \textbf{ēaðmōd-}, retaining the whole printed ending and the same case, number, gender and declension. An empty ending remains empty.
% mdparagraph:P181:end

% mdparagraph:P182:begin
Use the complete \textbf{brād; weak adjective} table. In every corresponding cell replace initial \textbf{brād-} with \textbf{ēaðmōd-}, retaining the whole printed ending and the same case, number, gender and declension. An empty ending remains empty.
% mdparagraph:P182:end

% teaching: NEW-U12-VOC-EADIG
% teaching: NEW-U12-ADJ-EADIG-STRONG
% teaching: NEW-U12-ADJ-EADIG-WEAK
% diagnostic: NEW-U12-DIAG-EADIG
% Page review repair: keep the lexical lead with its first display or derivation.
\reserve{24\baselineskip}
% mdparagraph:P183:begin
\textbf{ēadig} (adjective): happy; blessed. Diagnostic stem: \textbf{ēadig}.
% mdparagraph:P183:end

\reserve{18\baselineskip}
% mdparagraph:P184:begin
\textbf{ēadig; strong adjective}
% mdparagraph:P184:end

% table: NEW-U12-ADJ-EADIG-STRONG
% mdtable:T67:begin
% course-table-style: 2026-10
% course-panel:T67-P01:begin
\par\begingroup
\fontsize{10.5}{13.2}\selectfont
\setlength{\LTleft}{0pt}\setlength{\LTright}{\fill}
\noindent
\begin{tabular}{@{}>{\raggedright\arraybackslash}p{22.0mm}>{\raggedright\arraybackslash}p{31.188mm}>{\raggedright\arraybackslash}p{31.188mm}>{\raggedright\arraybackslash}p{31.188mm}@{}}
\toprule
\textbf{Case} & \textbf{Masculine singular} & \textbf{Neuter singular} & \textbf{Feminine singular} \\
\midrule
Nominative & ēadig & ēadig & ēadigu \\
Genitive & ēadiges & ēadiges & ēadigre \\
Dative & ēadigum & ēadigum & ēadigre \\
Accusative & ēadigne & ēadig & ēadige \\
Instrumental & ēadige & ēadige & as dative \\
\bottomrule
\end{tabular}
\par\endgroup
% course-panel:T67-P01:end
\par\addvspace{5pt}
% course-panel:T67-P02:begin
\par\begingroup
\fontsize{10.5}{13.2}\selectfont
\setlength{\LTleft}{0pt}\setlength{\LTright}{\fill}
\noindent
\begin{tabular}{@{}>{\raggedright\arraybackslash}p{22.0mm}>{\raggedright\arraybackslash}p{31.188mm}>{\raggedright\arraybackslash}p{31.188mm}>{\raggedright\arraybackslash}p{31.188mm}@{}}
\toprule
\textbf{Case} & \textbf{Masculine plural} & \textbf{Neuter plural} & \textbf{Feminine plural} \\
\midrule
Nominative & ēadige & ēadigu & ēadiga \\
Genitive & ēadigra & ēadigra & ēadigra \\
Dative & ēadigum & ēadigum & ēadigum \\
Accusative & ēadige & ēadigu & ēadiga \\
Instrumental & as dative & as dative & as dative \\
\bottomrule
\end{tabular}
\par\endgroup
% course-panel:T67-P02:end
% mdtable:T67:end

\reserve{18\baselineskip}
% mdparagraph:P185:begin
\textbf{ēadig; weak adjective}
% mdparagraph:P185:end

% table: NEW-U12-ADJ-EADIG-WEAK
% mdtable:T68:begin
% course-table-style: 2026-10
% course-panel:T68-P01:begin
\par\begingroup
\fontsize{10.5}{13.2}\selectfont
\setlength{\LTleft}{0pt}\setlength{\LTright}{\fill}
\noindent
\begin{tabular}{@{}>{\raggedright\arraybackslash}p{22.0mm}>{\raggedright\arraybackslash}p{31.188mm}>{\raggedright\arraybackslash}p{31.188mm}>{\raggedright\arraybackslash}p{31.188mm}@{}}
\toprule
\textbf{Case} & \textbf{Masculine singular} & \textbf{Neuter singular} & \textbf{Feminine singular} \\
\midrule
Nominative & ēadiga & ēadige & ēadige \\
Genitive & ēadigan & ēadigan & ēadigan \\
Dative & ēadigan & ēadigan & ēadigan \\
Accusative & ēadigan & ēadige & ēadigan \\
Instrumental & ēadigan & ēadigan & as dative \\
\bottomrule
\end{tabular}
\par\endgroup
% course-panel:T68-P01:end
\par\addvspace{5pt}
% course-panel:T68-P02:begin
\par\begingroup
\fontsize{10.5}{13.2}\selectfont
\setlength{\LTleft}{0pt}\setlength{\LTright}{\fill}
\noindent
\begin{tabular}{@{}>{\raggedright\arraybackslash}p{22.0mm}>{\raggedright\arraybackslash}p{31.188mm}>{\raggedright\arraybackslash}p{31.188mm}>{\raggedright\arraybackslash}p{31.188mm}@{}}
\toprule
\textbf{Case} & \textbf{Masculine plural} & \textbf{Neuter plural} & \textbf{Feminine plural} \\
\midrule
Nominative & ēadigan & ēadigan & ēadigan \\
Genitive & ēadigra & ēadigra & ēadigra \\
Dative & ēadigum & ēadigum & ēadigum \\
Accusative & ēadigan & ēadigan & ēadigan \\
Instrumental & as dative & as dative & as dative \\
\bottomrule
\end{tabular}
\par\endgroup
% course-panel:T68-P02:end
% mdtable:T68:end

% mdparagraph:P186:begin
The selected table retains i throughout. Cook also permits syncope before vowel endings, as in \textbf{ēadges}, and endingless feminine nominative singular and neuter nominative/\allowbreak{}accusative plural alternatives. These evidenced alternatives are accepted; they are not additional required tables.
% mdparagraph:P186:end

% teaching: NEW-U12-VOC-MILDHEORT
% teaching: NEW-U12-ADJ-MILDHEORT-STRONG
% teaching: NEW-U12-ADJ-MILDHEORT-WEAK
% diagnostic: NEW-U12-DIAG-MILDHEORT
% Page review repair: keep the lexical lead with its first display or derivation.
\reserve{10\baselineskip}
% mdparagraph:P187:begin
\textbf{mildheort} (adjective): merciful; compassionate. Diagnostic stem: \textbf{mildheort}.
% mdparagraph:P187:end

% mdparagraph:P188:begin
Use the complete \textbf{brād; strong adjective} table. In every corresponding cell replace initial \textbf{brād-} with \textbf{mildheort-}, retaining the whole printed ending and the same case, number, gender and declension. An empty ending remains empty.
% mdparagraph:P188:end

% mdparagraph:P189:begin
Use the complete \textbf{brād; weak adjective} table. In every corresponding cell replace initial \textbf{brād-} with \textbf{mildheort-}, retaining the whole printed ending and the same case, number, gender and declension. An empty ending remains empty.
% mdparagraph:P189:end

% teaching: NEW-U12-VOC-ISERN-ADJECTIVE
% teaching: NEW-U12-ADJ-ISERN-ADJECTIVE-STRONG
% teaching: NEW-U12-ADJ-ISERN-ADJECTIVE-WEAK
% diagnostic: NEW-U12-DIAG-ISERN-ADJECTIVE
% Page review repair: keep the lexical lead with its first display or derivation.
\reserve{10\baselineskip}
% mdparagraph:P190:begin
\textbf{īsern} (adjective): made of iron. Diagnostic stem: \textbf{īsern}.
% mdparagraph:P190:end

% mdparagraph:P191:begin
Use the complete \textbf{ēadig; strong adjective} table. In every corresponding cell replace initial \textbf{ēadig-} with \textbf{īsern-}, retaining the whole printed ending and the same case, number, gender and declension. An empty ending remains empty.
% mdparagraph:P191:end

% mdparagraph:P192:begin
Use the complete \textbf{ēadig; weak adjective} table. In every corresponding cell replace initial \textbf{ēadig-} with \textbf{īsern-}, retaining the whole printed ending and the same case, number, gender and declension. An empty ending remains empty.
% mdparagraph:P192:end

% mdparagraph:P193:begin
This \textbf{īsern} means “made of iron” and agrees as an adjective. It is distinct from the already learned noun meaning iron. Retain its medial e in the selected forms.
% mdparagraph:P193:end

\reserve{7\baselineskip}
% mdparagraph:P194:begin
\section*{8. Additional words and constructions}\BookHeading{nav-ch12-learner-0011}{subsection}{8. Additional words and constructions}{1}
% mdparagraph:P194:end

\reserve{7\baselineskip}
% mdparagraph:P195:begin
\subsection*{Adverbs}\BookHeading{nav-ch12-learner-0012}{subsubsection}{Adverbs}{0}
% mdparagraph:P195:end

% teaching: NEW-U12-VOC-HAM
% teaching: NEW-U12-VOC-HIDER
% teaching: NEW-U12-VOC-GIET
% teaching: NEW-U12-VOC-TO-DAEG
% teaching: NEW-U12-VOC-NAEFRE
% mdtable:T69:begin
\par\begingroup
\fontsize{10.5}{13.2}\selectfont
\setlength{\LTleft}{0pt}\setlength{\LTright}{\fill}
\noindent
\begin{tabular}{@{}>{\raggedright\arraybackslash}p{28.45mm}>{\raggedright\arraybackslash}p{92.74mm}@{}}
\toprule
\textbf{Word} & \textbf{Meaning and use} \\
\midrule
hām & home; homewards, expressing direction without a preposition \\
hider & hither; toward the speaker \\
gīet & yet; still \\
tō-dæg & today \\
nǣfre & never \\
\bottomrule
\end{tabular}
\par\endgroup

% mdtable:T69:end
\reserve{7\baselineskip}
% mdparagraph:P196:begin
\subsection*{Objects and recipients}\BookHeading{nav-ch12-learner-0013}{subsubsection}{Objects and recipients}{0}
% mdparagraph:P196:end

% teaching: NEW-U12-GRAM-u12-gierwan-accusative
% teaching: NEW-U12-GRAM-u12-writan-accusative
% teaching: NEW-U12-GRAM-u12-sendan-recipient
% teaching: NEW-U12-GRAM-u12-weorpan-accusative
% teaching: NEW-U12-GRAM-u12-gemetan-beneficiary
% mdparagraph:P197:begin
\textbf{Gierwan} takes the accusative thing prepared. The same thing remains accusative with habban + gierwed. \textbf{Wrītan} takes an accusative text or letter written. \textbf{Sęndan} takes the accusative person or thing sent and, if expressed, a dative recipient. \textbf{Weorpan} takes the accusative thing thrown, including a fishhook. Distinguish \textbf{angel}, fishhook, from \textbf{ęngel}, angel.
% mdparagraph:P197:end

% mdparagraph:P198:begin
With \textbf{gemētan}, the thing found is accusative. In the assigned expression about finding rest, the souls benefiting from it are dative. An English “unto” does not require a preposition there. This is the stated beneficiary use, not a rule assigning dative to every object of finding.
% mdparagraph:P198:end

% teaching: NEW-U12-GRAM-u12-geswican-intransitive
% teaching: NEW-U12-GRAM-u12-slean-kill
% teaching: NEW-U12-GRAM-u12-geseon-objectless
% teaching: NEW-U12-VOC-GESWICAN
% teaching: NEW-U12-VOC-SLEAN
% teaching: NEW-U12-VOC-GESEON
% mdparagraph:P199:begin
The earlier \textbf{geswīcan + genitive} expresses ceasing from an activity or rule. Also use geswīcan \textbf{without an expressed object} when seedtime and harvest cease or fail. \textbf{Slēan} can mean kill or slay, with an accusative person; this adds a lethal sense to strike/\allowbreak{}smite. \textbf{Gesēon} can stand without an expressed object in the assigned personification of the sea. Keep the earlier selected masculine gender of \textbf{sǣ} when choosing its pronoun.
% mdparagraph:P199:end

\reserve{7\baselineskip}
% mdparagraph:P200:begin
\subsection*{Prepositions and the meanings of their phrases}\BookHeading{nav-ch12-learner-0014}{subsubsection}{Prepositions and the meanings of their phrases}{0}
% mdparagraph:P200:end

% teaching: NEW-U12-GRAM-u12-on-respect
% teaching: NEW-U12-GRAM-u12-welig-on-dative
% teaching: NEW-U12-GRAM-u12-beforan-motion
% teaching: NEW-U12-GRAM-u12-abutan-accusative
% teaching: NEW-U12-GRAM-u12-gemang-accusative
% teaching: NEW-U12-GRAM-u12-aet-learning-source
% teaching: NEW-U12-GRAM-u12-for-beneficiary
% teaching: NEW-U12-VOC-ON
% teaching: NEW-U12-VOC-WELIG
% teaching: NEW-U12-VOC-BEFORAN
% teaching: NEW-U12-VOC-ABUTAN
% teaching: NEW-U12-VOC-GEMANG
% teaching: NEW-U12-VOC-AET
% teaching: NEW-U12-VOC-FOR
% mdtable:T70:begin
\par\begingroup
\fontsize{10.5}{13.2}\selectfont
\setlength{\LTleft}{0pt}\setlength{\LTright}{\fill}
\noindent
\begin{tabular}{@{}>{\raggedright\arraybackslash}p{34.65mm}>{\raggedright\arraybackslash}p{86.54mm}@{}}
\toprule
\textbf{Construction} & \textbf{Use} \\
\midrule
welig on + dative & rich in a resource or product; this can describe a harvest as well as a person or place \\
ēaðmōd on + dative & lowly in a stated respect; on heortan, in heart \\
sęndan … beforan + dative & send someone before/\allowbreak{}ahead of a person; distinct from a dative recipient \\
ābūtan + accusative & around a person or thing, including building around it \\
sęndan … gemang + accusative & send among a group; use the accusative in the assigned sheep-and-wolves expression \\
leornian æt + dative & learn from a person; the thing learned may remain unexpressed \\
sęllan for + dative & give for/\allowbreak{}on behalf of a person in the assigned payment expression \\
\bottomrule
\end{tabular}
\par\endgroup

% mdtable:T70:end
% mdparagraph:P201:begin
The first two on phrases express respect, not movement. Beforan with the dative extends the earlier static “before” phrase to the person ahead of whom someone is sent. Dative location after gemang also occurs, but the productive motion phrase assigned here uses accusative. The for construction here does not teach every case or causal sense of that preposition.
% mdparagraph:P201:end

\reserve{7\baselineskip}
% mdparagraph:P202:begin
\subsection*{Duration, extent, and possession}\BookHeading{nav-ch12-learner-0015}{subsubsection}{Duration, extent, and possession}{0}
% mdparagraph:P202:end

% teaching: NEW-U12-GRAM-u12-duration-accusative
% teaching: NEW-U12-GRAM-u12-extent-genitive
% teaching: NEW-U12-GRAM-u12-predicate-possession
% mdparagraph:P203:begin
An\BookPoint{bk-ch12-learner-c457eea8440b} accusative quantity of time can express \textbf{duration}: how long an action or state lasts. This differs from identifying the day on which it happens. Higher numerals may govern a genitive noun; the phrase for forty years is supplied at Cook 19.2, so do not force its noun into an accusative shape.
% mdparagraph:P203:end

% mdparagraph:P204:begin
A\BookPoint{bk-ch12-learner-f5ceef68b943} numeral with a \textbf{genitive measure noun} can state breadth or height. The noun fōt has genitive plural \textbf{fōta}. The number of feet measures the extent; the predicate adjective still agrees with the subject being described. The twenty-feet phrase is supplied at Cook 21.3.
% mdparagraph:P204:end

% mdparagraph:P205:begin
Possession\BookPoint{bk-ch12-learner-d4a523783302} can also be predicative. A \textbf{genitive possessor} after be means “is someone’s”; the genitive is not a nominative noun predicate. An \textbf{agreeing possessive adjective} after be agrees with what is possessed. Thus the person/\allowbreak{}number of the possessor selects ðīn, while the gender and number of the possessed dæg or niht determine its grammatical position. The same surface can occupy more than one such position.
% mdparagraph:P205:end

\reserve{7\baselineskip}
% mdparagraph:P206:begin
\subsection*{Comparison}\BookHeading{nav-ch12-learner-0016}{subsubsection}{Comparison}{0}
% mdparagraph:P206:end

% teaching: NEW-U12-GRAM-u12-more-numerous
% teaching: NEW-U12-GRAM-u12-equative-swa-swa
% teaching: NEW-U12-GRAM-u12-swa-predicative-comparison
% teaching: NEW-U12-VOC-MA
% teaching: NEW-U12-VOC-SWA
% mdparagraph:P207:begin
With\BookPoint{bk-ch12-learner-d84bed8d05f0} a plural subject, indeclinable \textbf{mā} can mean “more in number”. A comparison standard after ðonne uses the earlier parallel-subject construction. This quantitative predicate is distinct from the earlier use of mā for a greater amount of doing; it does not require a new comparative adjective māran.
% mdparagraph:P207:end

% mdparagraph:P208:begin
For\BookPoint{bk-ch12-learner-514aac1c9ee6} equality of degree, use \textbf{swā … adjective … swā}. The first swā marks the degree, “as”; the second introduces the standard, “as”. A positive predicate adjective retains strong agreement with the subject. The nominal standard can be nominative, understood as the subject of be. These are two separate words with an adjective between them, not the hyphenated source conjunction swā-swā.
% mdparagraph:P208:end

% mdparagraph:P209:begin
A\BookPoint{bk-ch12-learner-161b6102df9e} single \textbf{swā + noun phrase} can instead compare the persons sent with sheep: “as; like”. The comparing noun phrase refers to the persons who are the object of sending. This is the particular nominal comparison assigned here, distinct from the earlier comparison between finite clauses.
% mdparagraph:P209:end

\reserve{7\baselineskip}
% mdparagraph:P210:begin
\subsection*{Never and never again}\BookHeading{nav-ch12-learner-0017}{subsubsection}{Never and never again}{0}
% mdparagraph:P210:end

% teaching: NEW-U12-GRAM-u12-negative-concord-bounded
% mdparagraph:P211:begin
In the assigned \textbf{nǣfre … ne} construction, nǣfre and the negative finite clause reinforce one negative meaning. Place ne before the finite verb. \textbf{Nǣfre mā} means “never more; never again”. Do not cancel the two negatives as though they made an affirmative. Some verbs have contracted negative forms; nis, nāh and the required nolde remain visibly supplied below. This introduces the stated expressions only, not an entire family of negative pronouns or negative paradigms.
% mdparagraph:P211:end

\reserve{7\baselineskip}
% mdparagraph:P212:begin
\section*{9. Practice}\BookHeading{nav-ch12-learner-0018}{subsection}{9. Practice}{1}
% mdparagraph:P212:end

% mdparagraph:P213:begin
Work through the form exercises before the analyses and compositions. Use the complete printed systems, preserving quantity and stem changes. In analysis, state both the form and its function: a shared spelling alone may not decide mood, case or person. Reference constructions may be analyzed with their explanation in view.
% mdparagraph:P213:end

\reserve{7\baselineskip}
% mdparagraph:P214:begin
\subsection*{F01}\BookHeading{nav-ch12-learner-0019}{subsubsection}{F01}{0}
% mdparagraph:P214:end

\par\begin{minipage}{\linewidth}\raggedright
% item: NEW-U12-F01
% mdparagraph:P215:begin
Give the complete nominative, genitive, dative and accusative singular and plural of fōt. Explain the vowel alternation.
% mdparagraph:P215:end

\end{minipage}\par\addvspace{4pt}

\reserve{7\baselineskip}
% mdparagraph:P216:begin
\subsection*{F02}\BookHeading{nav-ch12-learner-0020}{subsubsection}{F02}{0}
% mdparagraph:P216:end

\par\begin{minipage}{\linewidth}\raggedright
% item: NEW-U12-F02
% mdparagraph:P217:begin
Give the complete nominative, genitive, dative and accusative singular and plural of brycg and hęre, using the selected tables. Explain the differing stems before endings.
% mdparagraph:P217:end

\end{minipage}\par\addvspace{4pt}

\reserve{7\baselineskip}
% mdparagraph:P218:begin
\subsection*{F03}\BookHeading{nav-ch12-learner-0021}{subsubsection}{F03}{0}
% mdparagraph:P218:end

\par\begin{minipage}{\linewidth}\raggedright
% item: NEW-U12-F03
% mdparagraph:P219:begin
Give the complete present and preterite indicative of flēon and grētan: all three persons in both numbers for each tense.
% mdparagraph:P219:end

\end{minipage}\par\addvspace{4pt}

\reserve{7\baselineskip}
% mdparagraph:P220:begin
\subsection*{F04}\BookHeading{nav-ch12-learner-0022}{subsubsection}{F04}{0}
% mdparagraph:P220:end

\par\begin{minipage}{\linewidth}\raggedright
% item: NEW-U12-F04
% mdparagraph:P221:begin
For gemētan, ādrīfan, geopenian and cȳðan, give the present indicative second and third singular, the preterite second singular and the preterite third plural. Explain the stem and ending choices.
% mdparagraph:P221:end

\end{minipage}\par\addvspace{4pt}

\reserve{7\baselineskip}
% mdparagraph:P222:begin
\subsection*{F05}\BookHeading{nav-ch12-learner-0023}{subsubsection}{F05}{0}
% mdparagraph:P222:end

\par\begin{minipage}{\linewidth}\raggedright
% item: NEW-U12-F05
% mdparagraph:P223:begin
Give the singular and plural commands of flēon, ādrīfan, gemētan and geopenian.
% mdparagraph:P223:end

\end{minipage}\par\addvspace{4pt}

\reserve{7\baselineskip}
% mdparagraph:P224:begin
\subsection*{F06}\BookHeading{nav-ch12-learner-0024}{subsubsection}{F06}{0}
% mdparagraph:P224:end

\par\begin{minipage}{\linewidth}\raggedright
% item: NEW-U12-F06
% mdparagraph:P225:begin
Give all six person-and-number forms of the present indicative of habban with ðæt hūs geworht, then all six preterite indicative forms with the same words. Label the finite auxiliary and the unchanged past participle.
% mdparagraph:P225:end

\end{minipage}\par\addvspace{4pt}

\reserve{7\baselineskip}
% mdparagraph:P226:begin
\subsection*{F07}\BookHeading{nav-ch12-learner-0025}{subsubsection}{F07}{0}
% mdparagraph:P226:end

\par\begin{minipage}{\linewidth}\raggedright
% item: NEW-U12-F07
% mdparagraph:P227:begin
Give the base present and past participles of gierwan, sęttan and gewyrcean. Explain which base is used in the habban construction.
% mdparagraph:P227:end

\end{minipage}\par\addvspace{4pt}

\reserve{7\baselineskip}
% mdparagraph:P228:begin
\subsection*{A01}\BookHeading{nav-ch12-learner-0026}{subsubsection}{A01}{0}
% mdparagraph:P228:end

\par\begin{minipage}{\linewidth}\raggedright
% item: NEW-U12-A01
% mdparagraph:P229:begin
Translate and put the two narrated events in order. Identify the finite verb in each clause.
% mdparagraph:P229:end

% mdparagraph:P230:begin
Se mann cōm, and hē geseah ðæt hūs.
% mdparagraph:P230:end

\end{minipage}\par\addvspace{4pt}

\reserve{7\baselineskip}
% mdparagraph:P231:begin
\subsection*{A02}\BookHeading{nav-ch12-learner-0027}{subsubsection}{A02}{0}
% mdparagraph:P231:end

\par\begin{minipage}{\linewidth}\raggedright
% item: NEW-U12-A02
% mdparagraph:P232:begin
The house was finished before the visit. Translate, preserving that order of events. Identify the morphological tense of geworhte and explain why English may use had built.
% mdparagraph:P232:end

% mdparagraph:P233:begin
Se mann geseah ðæt hūs ðe his sunu geworhte.
% mdparagraph:P233:end

\end{minipage}\par\addvspace{4pt}

\reserve{7\baselineskip}
% mdparagraph:P234:begin
\subsection*{A03}\BookHeading{nav-ch12-learner-0028}{subsubsection}{A03}{0}
% mdparagraph:P234:end

\par\begin{minipage}{\linewidth}\raggedright
% item: NEW-U12-A03
% mdparagraph:P235:begin
The speaker is describing completed building work that they performed. Translate both sentences. Identify the builder, the object, the finite auxiliary and the past participle. Explain the difference between the two temporal viewpoints.
% mdparagraph:P235:end

% mdparagraph:P236:begin
Ic hæbbe ðæt hūs geworht.
% mdparagraph:P236:end

% mdparagraph:P237:begin
Ic hæfde ðæt hūs geworht.
% mdparagraph:P237:end

\end{minipage}\par\addvspace{4pt}

\reserve{7\baselineskip}
% mdparagraph:P238:begin
\subsection*{A04}\BookHeading{nav-ch12-learner-0029}{subsubsection}{A04}{0}
% mdparagraph:P238:end

\par\begin{minipage}{\linewidth}\raggedright
% item: NEW-U12-A04
% mdparagraph:P239:begin
In the second sentence, the speaker has just finished building the house. Translate both sentences. What work does habban do in each, and what additional information is introduced by geworht?
% mdparagraph:P239:end

% mdparagraph:P240:begin
Ic hæbbe ðæt hūs.
% mdparagraph:P240:end

% mdparagraph:P241:begin
Ic hæbbe ðæt hūs geworht.
% mdparagraph:P241:end

\end{minipage}\par\addvspace{4pt}

\reserve{7\baselineskip}
% mdparagraph:P242:begin
\subsection*{A05}\BookHeading{nav-ch12-learner-0030}{subsubsection}{A05}{0}
% mdparagraph:P242:end

\par\begin{minipage}{\linewidth}\raggedright
% item: NEW-U12-A05
% mdparagraph:P243:begin
Translate all three sentences. Identify their subjects and the horse’s role. Which sentence names the binder? Explain the usual event/\allowbreak{}state distinction without treating it as an absolute rule.
% mdparagraph:P243:end

% mdparagraph:P244:begin
Se mann band ðæt hors.
% mdparagraph:P244:end

% mdparagraph:P245:begin
Ðæt hors wearð bunden.
% mdparagraph:P245:end

% mdparagraph:P246:begin
Ðæt hors wæs bunden.
% mdparagraph:P246:end

\end{minipage}\par\addvspace{4pt}

\reserve{7\baselineskip}
% mdparagraph:P247:begin
\subsection*{A06}\BookHeading{nav-ch12-learner-0031}{subsubsection}{A06}{0}
% mdparagraph:P247:end

\par\begin{minipage}{\linewidth}\raggedright
% item: NEW-U12-A06
% mdparagraph:P248:begin
Translate. Identify the finite verb and the intransitive past participle. Explain why the man is not the object of an unnamed person’s action.
% mdparagraph:P248:end

% mdparagraph:P249:begin
Se mann wæs gewiten.
% mdparagraph:P249:end

\end{minipage}\par\addvspace{4pt}

\reserve{7\baselineskip}
% mdparagraph:P250:begin
\subsection*{A07}\BookHeading{nav-ch12-learner-0032}{subsubsection}{A07}{0}
% mdparagraph:P250:end

\par\begin{minipage}{\linewidth}\raggedright
% item: NEW-U12-A07
% mdparagraph:P251:begin
Translate, giving a simple-past rendering and a contextually possible English progressive. Identify the two verb forms and explain why the English construction is not imposed by form alone.
% mdparagraph:P251:end

% mdparagraph:P252:begin
Se mann wæs sittende.
% mdparagraph:P252:end

\end{minipage}\par\addvspace{4pt}

\reserve{7\baselineskip}
% mdparagraph:P253:begin
\subsection*{A08}\BookHeading{nav-ch12-learner-0033}{subsubsection}{A08}{0}
% mdparagraph:P253:end

\par\begin{minipage}{\linewidth}\raggedright
% item: NEW-U12-A08
% mdparagraph:P254:begin
Translate both reports. In each, the reported presence is simultaneous with the saying. Identify the main-clause tense and dependent mood; explain which moment the reported presence is related to.
% mdparagraph:P254:end

% mdparagraph:P255:begin
Hē sægð ðæt hē sīe hēr.
% mdparagraph:P255:end

% mdparagraph:P256:begin
Hē sægde ðæt hē wǣre hēr.
% mdparagraph:P256:end

\end{minipage}\par\addvspace{4pt}

\reserve{7\baselineskip}
% mdparagraph:P257:begin
\subsection*{A09}\BookHeading{nav-ch12-learner-0034}{subsubsection}{A09}{0}
% mdparagraph:P257:end

\par\begin{minipage}{\linewidth}\raggedright
% item: NEW-U12-A09
% mdparagraph:P258:begin
Translate. In this report the building precedes the saying. Identify the finite verbs, the participle and the object, and explain how the temporal relation is expressed.
% mdparagraph:P258:end

% mdparagraph:P259:begin
Hē sægde ðæt hē ðæt hūs geworht hæfde.
% mdparagraph:P259:end

\end{minipage}\par\addvspace{4pt}

\reserve{7\baselineskip}
% mdparagraph:P260:begin
\subsection*{A10}\BookHeading{nav-ch12-learner-0035}{subsubsection}{A10}{0}
% mdparagraph:P260:end

\par\begin{minipage}{\linewidth}\raggedright
% item: NEW-U12-A10
% mdparagraph:P261:begin
Translate as a report of intention. Which action is intended, and does the sentence assert that the intended action took place?
% mdparagraph:P261:end

% mdparagraph:P262:begin
Hē sægde ðæt hē wolde cuman.
% mdparagraph:P262:end

\end{minipage}\par\addvspace{4pt}

\reserve{7\baselineskip}
% mdparagraph:P263:begin
\subsection*{A11}\BookHeading{nav-ch12-learner-0036}{subsubsection}{A11}{0}
% mdparagraph:P263:end

\par\begin{minipage}{\linewidth}\raggedright
% item: NEW-U12-A11
% mdparagraph:P264:begin
The visit is still expected. Translate the first sentence in that setting. Translate the second as a statement of desire, and explain why English will does not always call for willan.
% mdparagraph:P264:end

% mdparagraph:P265:begin
Hē cymð.
% mdparagraph:P265:end

% mdparagraph:P266:begin
Hē wille cuman.
% mdparagraph:P266:end

\end{minipage}\par\addvspace{4pt}

\reserve{7\baselineskip}
% mdparagraph:P267:begin
\subsection*{C01}\BookHeading{nav-ch12-learner-0037}{subsubsection}{C01}{0}
% mdparagraph:P267:end

\par\begin{minipage}{\linewidth}\raggedright
% item: NEW-U12-C01
% mdparagraph:P268:begin
Translate: We have built the house. Use present habban with the selected past participle.
% mdparagraph:P268:end

\end{minipage}\par\addvspace{4pt}

\reserve{7\baselineskip}
% mdparagraph:P269:begin
\subsection*{C02}\BookHeading{nav-ch12-learner-0038}{subsubsection}{C02}{0}
% mdparagraph:P269:end

\par\begin{minipage}{\linewidth}\raggedright
% item: NEW-U12-C02
% mdparagraph:P270:begin
Translate: They had built the house. The building is earlier than another past event; use preterite habban with the selected past participle.
% mdparagraph:P270:end

\end{minipage}\par\addvspace{4pt}

\reserve{7\baselineskip}
% mdparagraph:P271:begin
\subsection*{C03}\BookHeading{nav-ch12-learner-0039}{subsubsection}{C03}{0}
% mdparagraph:P271:end

\par\begin{minipage}{\linewidth}\raggedright
% item: NEW-U12-C03
% mdparagraph:P272:begin
Translate: The men bound the horse. Then express the horse’s binding as an agentless passive, using preterite weorðan.
% mdparagraph:P272:end

\end{minipage}\par\addvspace{4pt}

\reserve{7\baselineskip}
% mdparagraph:P273:begin
\subsection*{C04}\BookHeading{nav-ch12-learner-0040}{subsubsection}{C04}{0}
% mdparagraph:P273:end

\par\begin{minipage}{\linewidth}\raggedright
% item: NEW-U12-C04
% mdparagraph:P274:begin
Translate: The monk says that he is here. Use the subjunctive to present his report. Then put the saying in the past, keeping his presence simultaneous with that past saying.
% mdparagraph:P274:end

\end{minipage}\par\addvspace{4pt}

\reserve{7\baselineskip}
% mdparagraph:P275:begin
\subsection*{C05}\BookHeading{nav-ch12-learner-0041}{subsubsection}{C05}{0}
% mdparagraph:P275:end

\par\begin{minipage}{\linewidth}\raggedright
% item: NEW-U12-C05
% mdparagraph:P276:begin
Translate: He said that he had built the house. Use the reporting subjunctive and preterite habban with the selected past participle.
% mdparagraph:P276:end

\end{minipage}\par\addvspace{4pt}

\reserve{7\baselineskip}
% mdparagraph:P277:begin
\subsection*{C06}\BookHeading{nav-ch12-learner-0042}{subsubsection}{C06}{0}
% mdparagraph:P277:end

\par\begin{minipage}{\linewidth}\raggedright
% item: NEW-U12-C06
% mdparagraph:P278:begin
Translate: He said that he intended to come. Use willan with the infinitive; do not assert that he actually came.
% mdparagraph:P278:end

\end{minipage}\par\addvspace{4pt}

\reserve{7\baselineskip}
% mdparagraph:P279:begin
\subsection*{C07}\BookHeading{nav-ch12-learner-0043}{subsubsection}{C07}{0}
% mdparagraph:P279:end

\par\begin{minipage}{\linewidth}\raggedright
% item: NEW-U12-C07
% mdparagraph:P280:begin
Address one person: Prepare the food. Then translate the reply: I have prepared the food. Use gierwan in both.
% mdparagraph:P280:end

\end{minipage}\par\addvspace{4pt}

\reserve{7\baselineskip}
% mdparagraph:P281:begin
\subsection*{C08}\BookHeading{nav-ch12-learner-0044}{subsubsection}{C08}{0}
% mdparagraph:P281:end

\par\begin{minipage}{\linewidth}\raggedright
% item: NEW-U12-C08
% mdparagraph:P282:begin
Translate: The bridge was built. Then a man came. Use an agentless passive with wesan for the completed bridge, followed by a simple preterite for the arrival.
% mdparagraph:P282:end

\end{minipage}\par\addvspace{4pt}

\reserve{7\baselineskip}
% mdparagraph:P283:begin
\section*{10. Cook, Group VI: Exercises 19–21}\BookHeading{nav-ch12-learner-0045}{subsection}{10. Cook, Group VI: Exercises 19–21}{1}
% mdparagraph:P283:end

% teaching: NEW-U12-LOCAL-HELPS
% mdparagraph:P284:begin
Translate all the following into Old English. The English wording, parentheses, omissions and section references are Cook’s. Directions and supplied help are separate. \textbf{Om.} means omit the indicated English word in the Old English construction; it does not authorize dropping the meaning of the sentence. \textbf{Sing.} specifies a singular addressee; \textbf{pret.} requires a preterite. Cook’s numbered references are to his grammar, not the sections of this unit.
% mdparagraph:P284:end

% mdparagraph:P285:begin
Use the taught vocabulary and the stated constructions. A valid alternative may differ from a selected target; when a direction specifies a construction or addressee, satisfy that direction. The fixed expressions and exceptional local forms printed as help remain available during the exercise.
% mdparagraph:P285:end

\reserve{7\baselineskip}
% mdparagraph:P286:begin
\subsection*{Exercise 19 — printed pp. 9–10 /\allowbreak{} PDF pp. 17–18}\BookHeading{nav-ch12-learner-0046}{subsubsection}{Exercise 19 — printed pp. 9–10 / PDF pp. 17–18}{0}
% mdparagraph:P286:end

\reserve{7\baselineskip}
% mdparagraph:P287:begin
\subsection*{19.1}\BookHeading{nav-ch12-learner-0047}{subsubsection}{19.1}{0}
% mdparagraph:P287:end

\par\begin{minipage}{\linewidth}\raggedright
% item: EX-VI-19-01
% mdparagraph:P288:begin
At length the king ceased to bear rule.
% mdparagraph:P288:end

% mdparagraph:P289:begin
Direction: Use geswīcan with a genitive for ceasing from rule.
% mdparagraph:P289:end

\end{minipage}\par\addvspace{4pt}

% teaching: NEW-U12-REUSED-U12-COOK-HELP-19-01
% teaching: NEW-U12-GRAM-U12-COOK-CONSTRUCTION-19-01
% mdparagraph:P290:begin
\textbf{Supplied help.} Use the fixed phrase æt nīehstan, at length. Its ending is supplied; no general superlative system is required.
% mdparagraph:P290:end

\reserve{7\baselineskip}
% mdparagraph:P291:begin
\subsection*{19.2}\BookHeading{nav-ch12-learner-0048}{subsubsection}{19.2}{0}
% mdparagraph:P291:end

\par\begin{minipage}{\linewidth}\raggedright
% item: EX-VI-19-02
% mdparagraph:P292:begin
He had the kingdom forty years.
% mdparagraph:P292:end

% mdparagraph:P293:begin
Direction: Distinguish possessing the kingdom from the new perfect construction.
% mdparagraph:P293:end

\end{minipage}\par\addvspace{4pt}

% teaching: NEW-U12-GRAM-U12-COOK-CONSTRUCTION-19-02
% mdparagraph:P294:begin
\textbf{Supplied help.} Use fēowertig gēara, forty years. This whole duration phrase is supplied; the higher numeral takes the genitive gēara. The duration construction is discussed in this unit.
% mdparagraph:P294:end

\reserve{7\baselineskip}
% mdparagraph:P295:begin
\subsection*{19.3}\BookHeading{nav-ch12-learner-0049}{subsubsection}{19.3}{0}
% mdparagraph:P295:end

\par\begin{minipage}{\linewidth}\raggedright
% item: EX-VI-19-03
% mdparagraph:P296:begin
When were the bridges built?
% mdparagraph:P296:end

% mdparagraph:P297:begin
Direction: Use a past agentless passive, with the bridges as subject.
% mdparagraph:P297:end

\end{minipage}\par\addvspace{4pt}

\reserve{7\baselineskip}
% mdparagraph:P298:begin
\subsection*{19.4}\BookHeading{nav-ch12-learner-0050}{subsubsection}{19.4}{0}
% mdparagraph:P298:end

\par\begin{minipage}{\linewidth}\raggedright
% item: EX-VI-19-04
% mdparagraph:P299:begin
To-day thou shalt be with me.
% mdparagraph:P299:end

% mdparagraph:P300:begin
Direction: Use a taught form of be in its future context.
% mdparagraph:P300:end

\end{minipage}\par\addvspace{4pt}

\reserve{7\baselineskip}
% mdparagraph:P301:begin
\subsection*{19.5}\BookHeading{nav-ch12-learner-0051}{subsubsection}{19.5}{0}
% mdparagraph:P301:end

\par\begin{minipage}{\linewidth}\raggedright
% item: EX-VI-19-05
% mdparagraph:P302:begin
What (om. did) the Romans build (pret.) in Britain?
% mdparagraph:P302:end

% mdparagraph:P303:begin
Direction: Keep the explicit preterite requirement.
% mdparagraph:P303:end

\end{minipage}\par\addvspace{4pt}

% teaching: NEW-U12-GRAM-U12-COOK-CONSTRUCTION-19-05
% mdparagraph:P304:begin
\textbf{Supplied help.} Rōmane, Romans, is the supplied nominative plural; on Breotone, in Britain, supplies the place-name dative. Retain the preterite target.
% mdparagraph:P304:end

\reserve{7\baselineskip}
% mdparagraph:P305:begin
\subsection*{19.6}\BookHeading{nav-ch12-learner-0052}{subsubsection}{19.6}{0}
% mdparagraph:P305:end

\par\begin{minipage}{\linewidth}\raggedright
% item: EX-VI-19-06
% mdparagraph:P306:begin
I send my angel before (beforan, 166) thy face.
% mdparagraph:P306:end

% mdparagraph:P307:begin
Direction: Preserve the original face wording and the specified beforan phrase.
% mdparagraph:P307:end

\end{minipage}\par\addvspace{4pt}

\reserve{7\baselineskip}
% mdparagraph:P308:begin
\subsection*{19.7}\BookHeading{nav-ch12-learner-0053}{subsubsection}{19.7}{0}
% mdparagraph:P308:end

\par\begin{minipage}{\linewidth}\raggedright
% item: EX-VI-19-07
% mdparagraph:P309:begin
Come (sing.), and bring those letters that you wrote.
% mdparagraph:P309:end

% mdparagraph:P310:begin
Direction: Treat you in the relative clause as the single person addressed by the commands.
% mdparagraph:P310:end

\end{minipage}\par\addvspace{4pt}

\reserve{7\baselineskip}
% mdparagraph:P311:begin
\subsection*{19.8}\BookHeading{nav-ch12-learner-0054}{subsubsection}{19.8}{0}
% mdparagraph:P311:end

\par\begin{minipage}{\linewidth}\raggedright
% item: EX-VI-19-08
% mdparagraph:P312:begin
I promise you obedience, if you help (196. d) me.
% mdparagraph:P312:end

% mdparagraph:P313:begin
Direction: Cook §196.d calls for the selected subjunctive in the if-clause. Keep singular or plural you consistent.
% mdparagraph:P313:end

\end{minipage}\par\addvspace{4pt}

\reserve{7\baselineskip}
% mdparagraph:P314:begin
\subsection*{19.9}\BookHeading{nav-ch12-learner-0055}{subsubsection}{19.9}{0}
% mdparagraph:P314:end

\par\begin{minipage}{\linewidth}\raggedright
% item: EX-VI-19-09
% mdparagraph:P315:begin
Send me (om. a) large army.
% mdparagraph:P315:end

% mdparagraph:P316:begin
Direction: Address one person; distinguish the thing sent from its recipient.
% mdparagraph:P316:end

\end{minipage}\par\addvspace{4pt}

\reserve{7\baselineskip}
% mdparagraph:P317:begin
\subsection*{19.10}\BookHeading{nav-ch12-learner-0056}{subsubsection}{19.10}{0}
% mdparagraph:P317:end

\par\begin{minipage}{\linewidth}\raggedright
% item: EX-VI-19-10
% mdparagraph:P318:begin
I fight with my enemies.
% mdparagraph:P318:end

% mdparagraph:P319:begin
Direction: Retain the earlier selected feohtan wið + accusative opponent construction.
% mdparagraph:P319:end

\end{minipage}\par\addvspace{4pt}

\reserve{7\baselineskip}
% mdparagraph:P320:begin
\subsection*{19.11}\BookHeading{nav-ch12-learner-0057}{subsubsection}{19.11}{0}
% mdparagraph:P320:end

\par\begin{minipage}{\linewidth}\raggedright
% item: EX-VI-19-11
% mdparagraph:P321:begin
Expel your foes.
% mdparagraph:P321:end

% mdparagraph:P322:begin
Direction: Address more than one person and make the possessive agree with the foes.
% mdparagraph:P322:end

\end{minipage}\par\addvspace{4pt}

\reserve{7\baselineskip}
% mdparagraph:P323:begin
\subsection*{19.12}\BookHeading{nav-ch12-learner-0058}{subsubsection}{19.12}{0}
% mdparagraph:P323:end

\par\begin{minipage}{\linewidth}\raggedright
% item: EX-VI-19-12
% mdparagraph:P324:begin
The soldiers went home.
% mdparagraph:P324:end

% mdparagraph:P325:begin
Direction: Use the directional adverb home without adding a preposition.
% mdparagraph:P325:end

\end{minipage}\par\addvspace{4pt}

\reserve{7\baselineskip}
% mdparagraph:P326:begin
\subsection*{19.13}\BookHeading{nav-ch12-learner-0059}{subsubsection}{19.13}{0}
% mdparagraph:P326:end

\par\begin{minipage}{\linewidth}\raggedright
% item: EX-VI-19-13
% mdparagraph:P327:begin
He will come with (om. a) fleet.
% mdparagraph:P327:end

% mdparagraph:P328:begin
Direction: Express future reference without treating willan as compulsory.
% mdparagraph:P328:end

\end{minipage}\par\addvspace{4pt}

\reserve{7\baselineskip}
% mdparagraph:P329:begin
\subsection*{19.14}\BookHeading{nav-ch12-learner-0060}{subsubsection}{19.14}{0}
% mdparagraph:P329:end

\par\begin{minipage}{\linewidth}\raggedright
% item: EX-VI-19-14
% mdparagraph:P330:begin
Slay (sing.) all that you find (gemētan).
% mdparagraph:P330:end

% mdparagraph:P331:begin
Direction: Address one person. For this target use all those who with a relative clause.
% mdparagraph:P331:end

\end{minipage}\par\addvspace{4pt}

\reserve{7\baselineskip}
% mdparagraph:P332:begin
\subsection*{19.15}\BookHeading{nav-ch12-learner-0061}{subsubsection}{19.15}{0}
% mdparagraph:P332:end

\par\begin{minipage}{\linewidth}\raggedright
% item: EX-VI-19-15
% mdparagraph:P333:begin
The harvest is rich in fruits.
% mdparagraph:P333:end

% mdparagraph:P334:begin
Direction: Use the rich-in construction taught in §8.
% mdparagraph:P334:end

\end{minipage}\par\addvspace{4pt}

\reserve{7\baselineskip}
% mdparagraph:P335:begin
\subsection*{19.16}\BookHeading{nav-ch12-learner-0062}{subsubsection}{19.16}{0}
% mdparagraph:P335:end

\par\begin{minipage}{\linewidth}\raggedright
% item: EX-VI-19-16
% mdparagraph:P336:begin
Seedtime and harvest never shall fail (geswīcan).
% mdparagraph:P336:end

% mdparagraph:P337:begin
Direction: Use objectless geswīcan, the plural subject, and the taught never construction.
% mdparagraph:P337:end

\end{minipage}\par\addvspace{4pt}

\reserve{7\baselineskip}
% mdparagraph:P338:begin
\subsection*{19.17}\BookHeading{nav-ch12-learner-0063}{subsubsection}{19.17}{0}
% mdparagraph:P338:end

\par\begin{minipage}{\linewidth}\raggedright
% item: EX-VI-19-17
% mdparagraph:P339:begin
The year has four seasons.
% mdparagraph:P339:end

% mdparagraph:P340:begin
Direction: Use ordinary possession, not an auxiliary perfect.
% mdparagraph:P340:end

\end{minipage}\par\addvspace{4pt}

\reserve{7\baselineskip}
% mdparagraph:P341:begin
\subsection*{Exercise 20 — printed p. 10 /\allowbreak{} PDF p. 18}\BookHeading{nav-ch12-learner-0064}{subsubsection}{Exercise 20 — printed p. 10 / PDF p. 18}{0}
% mdparagraph:P341:end

\reserve{7\baselineskip}
% mdparagraph:P342:begin
\subsection*{20.1}\BookHeading{nav-ch12-learner-0065}{subsubsection}{20.1}{0}
% mdparagraph:P342:end

\par\begin{minipage}{\linewidth}\raggedright
% item: EX-VI-20-01
% mdparagraph:P343:begin
The earth is (om. the) Lord’s.
% mdparagraph:P343:end

% mdparagraph:P344:begin
Direction: Use a genitive possessor after be.
% mdparagraph:P344:end

\end{minipage}\par\addvspace{4pt}

\reserve{7\baselineskip}
% mdparagraph:P345:begin
\subsection*{20.2}\BookHeading{nav-ch12-learner-0066}{subsubsection}{20.2}{0}
% mdparagraph:P345:end

\par\begin{minipage}{\linewidth}\raggedright
% item: EX-VI-20-02
% mdparagraph:P346:begin
The sea saw, and it fled.
% mdparagraph:P346:end

% mdparagraph:P347:begin
Direction: Retain the selected masculine gender of sǣ and the objectless seeing sense.
% mdparagraph:P347:end

\end{minipage}\par\addvspace{4pt}

\reserve{7\baselineskip}
% mdparagraph:P348:begin
\subsection*{20.3}\BookHeading{nav-ch12-learner-0067}{subsubsection}{20.3}{0}
% mdparagraph:P348:end

\par\begin{minipage}{\linewidth}\raggedright
% item: EX-VI-20-03
% mdparagraph:P349:begin
He shall flee the iron weapon.
% mdparagraph:P349:end

% mdparagraph:P350:begin
Direction: The thing avoided is accusative; distinguish the adjective iron from the noun.
% mdparagraph:P350:end

\end{minipage}\par\addvspace{4pt}

\reserve{7\baselineskip}
% mdparagraph:P351:begin
\subsection*{20.4}\BookHeading{nav-ch12-learner-0068}{subsubsection}{20.4}{0}
% mdparagraph:P351:end

\par\begin{minipage}{\linewidth}\raggedright
% item: EX-VI-20-04
% mdparagraph:P352:begin
He that (sē ðe) is not (om. the) shepherd, and owns not (127) the sheep, when he sees the wolf, then flees he, and abandons the sheep, and the wolf seizes the sheep.
% mdparagraph:P352:end

% mdparagraph:P353:begin
Direction: Translate the whole relative–temporal–main-clause sequence and preserve its changing subjects.
% mdparagraph:P353:end

\end{minipage}\par\addvspace{4pt}

% teaching: NEW-U12-GRAM-U12-COOK-CONSTRUCTION-20-04
% teaching: NEW-U12-GRAM-u12-negative-agan-local
% mdparagraph:P354:begin
\textbf{Supplied help.} Use nāh, does not own, for the negative possession clause. Cook’s hint sē ðe is retained as printed; the model may use the earlier selected Se ðe, he who. The long-vowel sē is a supplied source spelling, not a newly required paradigm.
% mdparagraph:P354:end

% teaching: NEW-U12-REUSED-U12-COOK-HELP-NIS
% teaching: NEW-U12-GRAM-U12-COOK-CONSTRUCTION-NIS
% mdparagraph:P355:begin
\textbf{Supplied help.} Reuse the supplied contraction nis, is not, for the negative be-clause. This local form does not teach the full negative verb system.
% mdparagraph:P355:end

\reserve{7\baselineskip}
% mdparagraph:P356:begin
\subsection*{20.5}\BookHeading{nav-ch12-learner-0069}{subsubsection}{20.5}{0}
% mdparagraph:P356:end

\par\begin{minipage}{\linewidth}\raggedright
% item: EX-VI-20-05
% mdparagraph:P357:begin
The day is thine, the night also is thine; thou hast prepared the sun and the moon; thou hast set the borders of the earth; thou hast made summer and winter.
% mdparagraph:P357:end

% mdparagraph:P358:begin
Direction: Use agreeing possessive predicates, then habban with the selected past participle for each completed action.
% mdparagraph:P358:end

\end{minipage}\par\addvspace{4pt}

\reserve{7\baselineskip}
% mdparagraph:P359:begin
\subsection*{20.6}\BookHeading{nav-ch12-learner-0070}{subsubsection}{20.6}{0}
% mdparagraph:P359:end

\par\begin{minipage}{\linewidth}\raggedright
% item: EX-VI-20-06
% mdparagraph:P360:begin
Let us build these cities, and make about (ābūtan, 172) them walls, and towers, and gates, while the land is yet (gīet) before us.
% mdparagraph:P360:end

% mdparagraph:P361:begin
Direction: Use first-plural subjunctive exhortations with gewyrcean for both building clauses, the specified ābūtan phrase, and the factual while-clause.
% mdparagraph:P361:end

\end{minipage}\par\addvspace{4pt}

\reserve{7\baselineskip}
% mdparagraph:P362:begin
\subsection*{20.7}\BookHeading{nav-ch12-learner-0071}{subsubsection}{20.7}{0}
% mdparagraph:P362:end

\par\begin{minipage}{\linewidth}\raggedright
% item: EX-VI-20-07
% mdparagraph:P363:begin
These nations are more than I; how can I expel them?
% mdparagraph:P363:end

% mdparagraph:P364:begin
Direction: Use indeclinable quantitative mā and magan with the plain infinitive of expel.
% mdparagraph:P364:end

\end{minipage}\par\addvspace{4pt}

\reserve{7\baselineskip}
% mdparagraph:P365:begin
\subsection*{20.8}\BookHeading{nav-ch12-learner-0072}{subsubsection}{20.8}{0}
% mdparagraph:P365:end

\par\begin{minipage}{\linewidth}\raggedright
% item: EX-VI-20-08
% mdparagraph:P366:begin
Go to the sea, and cast a fishhook, and take the first (ǣrest) fish, and open (geopenian) his mouth (mūð, 43).
% mdparagraph:P366:end

% mdparagraph:P367:begin
Direction: Keep all the commands singular and preserve the fish’s masculine reference.
% mdparagraph:P367:end

\end{minipage}\par\addvspace{4pt}

% teaching: NEW-U12-GRAM-U12-COOK-CONSTRUCTION-20-08
% mdparagraph:P368:begin
\textbf{Supplied help.} Cook prints first (ǣrest). Treat this as a lexical hint: use the previously taught ordinal ǣresta and the needed agreeing form ǣrestan in ðone ǣrestan fisc. The uninflected printed hint is preserved in the prompt.
% mdparagraph:P368:end

\reserve{7\baselineskip}
% mdparagraph:P369:begin
\subsection*{20.9}\BookHeading{nav-ch12-learner-0073}{subsubsection}{20.9}{0}
% mdparagraph:P369:end

\par\begin{minipage}{\linewidth}\raggedright
% item: EX-VI-20-09
% mdparagraph:P370:begin
Thou wilt find a penny; take that, and give (sęllan) for me and for thee.
% mdparagraph:P370:end

% mdparagraph:P371:begin
Direction: Keep the penny masculine; use for + dative in both beneficiary phrases.
% mdparagraph:P371:end

\end{minipage}\par\addvspace{4pt}

\reserve{7\baselineskip}
% mdparagraph:P372:begin
\subsection*{Exercise 21 — printed pp. 10–11 /\allowbreak{} PDF pp. 18–19}\BookHeading{nav-ch12-learner-0074}{subsubsection}{Exercise 21 — printed pp. 10–11 / PDF pp. 18–19}{0}
% mdparagraph:P372:end

\reserve{7\baselineskip}
% mdparagraph:P373:begin
\subsection*{21.1}\BookHeading{nav-ch12-learner-0075}{subsubsection}{21.1}{0}
% mdparagraph:P373:end

\par\begin{minipage}{\linewidth}\raggedright
% item: EX-VI-21-01
% mdparagraph:P374:begin
Then said he, She is not dead.
% mdparagraph:P374:end

% mdparagraph:P375:begin
Direction: Retain the past report and the present state in direct speech.
% mdparagraph:P375:end

% mdparagraph:P376:begin
\textbf{Supplied help.} Reuse the supplied contraction nis, is not, for the negative be-clause. This local form does not teach the full negative verb system.
% mdparagraph:P376:end

\end{minipage}\par\addvspace{4pt}

\reserve{7\baselineskip}
% mdparagraph:P377:begin
\subsection*{21.2}\BookHeading{nav-ch12-learner-0076}{subsubsection}{21.2}{0}
% mdparagraph:P377:end

\par\begin{minipage}{\linewidth}\raggedright
% item: EX-VI-21-02
% mdparagraph:P378:begin
Then took he her hand, and said, Maiden, arise.
% mdparagraph:P378:end

% mdparagraph:P379:begin
Direction: Preserve the feminine possessor and address the maiden with the singular command of rīsan.
% mdparagraph:P379:end

\end{minipage}\par\addvspace{4pt}

\reserve{7\baselineskip}
% mdparagraph:P380:begin
\subsection*{21.3}\BookHeading{nav-ch12-learner-0077}{subsubsection}{21.3}{0}
% mdparagraph:P380:end

\par\begin{minipage}{\linewidth}\raggedright
% item: EX-VI-21-03
% mdparagraph:P381:begin
The stone wall is seven feet broad, and twenty feet high.
% mdparagraph:P381:end

% mdparagraph:P382:begin
Direction: Use both genitive measure phrases and adjectives agreeing with the wall.
% mdparagraph:P382:end

\end{minipage}\par\addvspace{4pt}

% teaching: NEW-U12-GRAM-U12-COOK-CONSTRUCTION-21-03
% mdparagraph:P383:begin
\textbf{Supplied help.} Twēntig means twenty. Use the supplied measure phrase twēntig fōta; the newly taught genitive of extent also governs seofon fōta.
% mdparagraph:P383:end

\reserve{7\baselineskip}
% mdparagraph:P384:begin
\subsection*{21.4}\BookHeading{nav-ch12-learner-0078}{subsubsection}{21.4}{0}
% mdparagraph:P384:end

\par\begin{minipage}{\linewidth}\raggedright
% item: EX-VI-21-04
% mdparagraph:P385:begin
I took leave of her (acc.), and announced (om. to) her that I never more would visit her.
% mdparagraph:P385:end

% mdparagraph:P386:begin
Direction: Use grētan, cȳðan and an intended later genēosian with willan. Use the reporting subjunctive; do not assert that a visit occurred.
% mdparagraph:P386:end

\end{minipage}\par\addvspace{4pt}

% teaching: NEW-U12-GRAM-U12-COOK-CONSTRUCTION-21-04
% mdparagraph:P387:begin
\textbf{Supplied help.} Use the supplied nolde, the contracted negative preterite subjunctive singular of willan, in the reported later intention. Nǣfre mā adds never again; it combines with the negative finite verb. This one supplied contraction does not assign the entire negative paradigm.
% mdparagraph:P387:end

\reserve{7\baselineskip}
% mdparagraph:P388:begin
\subsection*{21.5}\BookHeading{nav-ch12-learner-0079}{subsubsection}{21.5}{0}
% mdparagraph:P388:end

\par\begin{minipage}{\linewidth}\raggedright
% item: EX-VI-21-05
% mdparagraph:P389:begin
The righteous (om. ones) are as bold as (om. a) lion.
% mdparagraph:P389:end

% mdparagraph:P390:begin
Direction: Use the weak adjective with an article for the righteous people, and separated swā … adjective … swā.
% mdparagraph:P390:end

\end{minipage}\par\addvspace{4pt}

% teaching: NEW-U12-GRAM-U12-COOK-CONSTRUCTION-21-05
% mdparagraph:P391:begin
\textbf{Supplied help.} Use the supplied nominative lēo, lion, after the second swā. No full paradigm of this loanword is assigned.
% mdparagraph:P391:end

\reserve{7\baselineskip}
% mdparagraph:P392:begin
\subsection*{21.6}\BookHeading{nav-ch12-learner-0080}{subsubsection}{21.6}{0}
% mdparagraph:P392:end

\par\begin{minipage}{\linewidth}\raggedright
% item: EX-VI-21-06
% mdparagraph:P393:begin
The Britons were cowardly.
% mdparagraph:P393:end

% mdparagraph:P394:begin
Direction: Make the predicate adjective plural.
% mdparagraph:P394:end

\end{minipage}\par\addvspace{4pt}

% teaching: NEW-U12-GRAM-U12-COOK-CONSTRUCTION-21-06
% mdparagraph:P395:begin
\textbf{Supplied help.} Bryttas, Britons, is supplied as a nominative plural.
% mdparagraph:P395:end

\reserve{7\baselineskip}
% mdparagraph:P396:begin
\subsection*{21.7}\BookHeading{nav-ch12-learner-0081}{subsubsection}{21.7}{0}
% mdparagraph:P396:end

\par\begin{minipage}{\linewidth}\raggedright
% item: EX-VI-21-07
% mdparagraph:P397:begin
Come hither, fair maiden.
% mdparagraph:P397:end

% mdparagraph:P398:begin
Direction: Use the weak adjective in singular direct address.
% mdparagraph:P398:end

\end{minipage}\par\addvspace{4pt}

\reserve{7\baselineskip}
% mdparagraph:P399:begin
\subsection*{21.8}\BookHeading{nav-ch12-learner-0082}{subsubsection}{21.8}{0}
% mdparagraph:P399:end

\par\begin{minipage}{\linewidth}\raggedright
% item: EX-VI-21-08
% mdparagraph:P400:begin
Happy is that people whose God is the Lord.
% mdparagraph:P400:end

% mdparagraph:P401:begin
Direction: Distinguish the case of the antecedent people from the genitive function of whose.
% mdparagraph:P401:end

\end{minipage}\par\addvspace{4pt}

\reserve{7\baselineskip}
% mdparagraph:P402:begin
\subsection*{21.9}\BookHeading{nav-ch12-learner-0083}{subsubsection}{21.9}{0}
% mdparagraph:P402:end

\par\begin{minipage}{\linewidth}\raggedright
% item: EX-VI-21-09
% mdparagraph:P403:begin
Learn of (æt) me, for I am merciful and lowly (ēaðmōd) in heart, and ye shall find (gemētan) rest (ręst, 51. b) (om. unto) your souls.
% mdparagraph:P403:end

% mdparagraph:P404:begin
Direction: Address more than one person. Use the supplied original æt hint, on + dative for respect, and dative beneficiaries after finding rest.
% mdparagraph:P404:end

\end{minipage}\par\addvspace{4pt}

\reserve{7\baselineskip}
% mdparagraph:P405:begin
\subsection*{21.10}\BookHeading{nav-ch12-learner-0084}{subsubsection}{21.10}{0}
% mdparagraph:P405:end

\par\begin{minipage}{\linewidth}\raggedright
% item: EX-VI-21-10
% mdparagraph:P406:begin
Now I send you as sheep among (gemang, 172) wolves.
% mdparagraph:P406:end

% mdparagraph:P407:begin
Direction: Use the single-swā nominal comparison and the selected gemang + accusative phrase.
% mdparagraph:P407:end

\end{minipage}\par\addvspace{4pt}

\reserve{7\baselineskip}
% mdparagraph:P408:begin
\subsection*{21.11}\BookHeading{nav-ch12-learner-0085}{subsubsection}{21.11}{0}
% mdparagraph:P408:end

\par\begin{minipage}{\linewidth}\raggedright
% item: EX-VI-21-11
% mdparagraph:P409:begin
When they pursue (ēhtan) you in this city, flee into (on) (om. an)other.
% mdparagraph:P409:end

% mdparagraph:P410:begin
Direction: Use the selected accusative person pursued, dative location, and plural flight command.
% mdparagraph:P410:end

\end{minipage}\par\addvspace{4pt}

% teaching: NEW-U12-GRAM-U12-COOK-CONSTRUCTION-21-11
% teaching: NEW-U12-GRAM-u12-on-elliptic-city
% mdparagraph:P411:begin
\textbf{Supplied help.} Use on ōðre, into another, with feminine accusative singular ōðre and an understood ceastre, city. Cook prints (om. an)other: omit the article an, not the sense another. The city omission is supplied locally.
% mdparagraph:P411:end

\reserve{7\baselineskip}
% mdparagraph:P412:begin
\subsection*{21.12}\BookHeading{nav-ch12-learner-0086}{subsubsection}{21.12}{0}
% mdparagraph:P412:end

\par\begin{minipage}{\linewidth}\raggedright
% item: EX-VI-21-12
% mdparagraph:P413:begin
I take leave of my native land.
% mdparagraph:P413:end

% mdparagraph:P414:begin
Direction: Use present grētan with the accusative native land.
% mdparagraph:P414:end

\end{minipage}\par\addvspace{4pt}

\reserve{7\baselineskip}
% mdparagraph:P415:begin
\section*{11. A house and a return}\BookHeading{nav-ch12-learner-0087}{subsection}{11. A house and a return}{1}
% mdparagraph:P415:end

% mdparagraph:P416:begin
This connected account is new teaching material. Its narrative and direct speech use different viewpoints; the report of an intention does not by itself assert that the intended action happened.
% mdparagraph:P416:end

% composition: NEW-U12-TEXT-A
% mdparagraph:P417:begin
\par\begin{minipage}{\linewidth}\raggedright
\readingtext{Se mann hæfde ān hūs geworht. Ðā cōm his sunu, and se sunu geseah ðæt hūs. Se fæder cwæð: “Nū hæbbe ic ðæt hūs geworht.” Se sunu sægde ðæt hē wolde faran. Ðā se sunu gewāt, ðā sæt se fæder on ðǣm hūse. Ðā cōm se sunu, and hē sealde his fæder męte.}
\end{minipage}\par
% mdparagraph:P417:end

\reserve{7\baselineskip}
% mdparagraph:P418:begin
\subsection*{R01}\BookHeading{nav-ch12-learner-0088}{subsubsection}{R01}{0}
% mdparagraph:P418:end

\par\begin{minipage}{\linewidth}\raggedright
% item: NEW-U12-R01
% mdparagraph:P419:begin
Translate the whole new account, A house and a return. Preserve its changes between narrative, direct speech and reported intention.
% mdparagraph:P419:end

\end{minipage}\par\addvspace{4pt}

\reserve{7\baselineskip}
% mdparagraph:P420:begin
\subsection*{R02}\BookHeading{nav-ch12-learner-0089}{subsubsection}{R02}{0}
% mdparagraph:P420:end

\par\begin{minipage}{\linewidth}\raggedright
% item: NEW-U12-R02
% mdparagraph:P421:begin
Put these events in order: the son’s first arrival; building the house; the son’s departure; the father’s announcement; the son’s return with food. Which verb phrase places one event before the story’s first arrival?
% mdparagraph:P421:end

\end{minipage}\par\addvspace{4pt}

\reserve{7\baselineskip}
% mdparagraph:P422:begin
\subsection*{R03}\BookHeading{nav-ch12-learner-0090}{subsubsection}{R03}{0}
% mdparagraph:P422:end

\par\begin{minipage}{\linewidth}\raggedright
% item: NEW-U12-R03
% mdparagraph:P423:begin
Identify the referents of ic in the father’s speech, hē in the reported intention, and hē in the last sentence. Is the journey an intention only, or does the following narrative show a departure? Cite the relevant words.
% mdparagraph:P423:end

\end{minipage}\par\addvspace{4pt}

\reserve{7\baselineskip}
% mdparagraph:P424:begin
\subsection*{R04}\BookHeading{nav-ch12-learner-0091}{subsubsection}{R04}{0}
% mdparagraph:P424:end

\par\begin{minipage}{\linewidth}\raggedright
% item: NEW-U12-R04
% mdparagraph:P425:begin
In Nū hæbbe ic ðæt hūs geworht, identify the subject and object. Express the house’s completed state in an agentless sentence with present wesan and the same participle.
% mdparagraph:P425:end

\end{minipage}\par\addvspace{4pt}

\reserve{7\baselineskip}
% mdparagraph:P426:begin
\section*{12. First Roman relief and the renewed invasion}\BookHeading{nav-ch12-learner-0092}{subsection}{12. First Roman relief and the renewed invasion}{1}
% mdparagraph:P426:end

% mdparagraph:P427:begin
This exact 155-word extract comes from Cook’s \emph{First Book}, Reader VI, \textbf{The Invasion of Britain by the Picts and Scots}, printed pp. 145–146 /\allowbreak{} PDF pp. 161–162. It begins at editorial prose line 1 of p. 145, includes all seven prose lines there, and continues through foriermdon at the end of the first sentence on editorial prose line 11 of p. 146. These editorial counts differ from Cook’s printed marginal line numbers. The following English context is separate from the Old English.
% mdparagraph:P427:end

% mdparagraph:P428:begin
The preceding passage says that Roman rule in Britain had ended and describes the Britons within the old wall. In this extract the Britons seek military help from Rome after raids by the Picts and Scots. This context is an editorial English summary, not translated words added to the Old English excerpt.
% mdparagraph:P428:end

% teaching: NEW-U12-SOURCE-HELPS
\par\begin{minipage}{\linewidth}\raggedright
% source: FB-VI-PICTS-SCOTS-RELIEF-01
% mdparagraph:P429:begin
\readingtext{Ðā ongunnon twā ðēoda, Pihtas norðan, and Scottas westan, hīe onwinnan, and hiera ǣhta niman and hęrgian; and hīe fela gēara iermdon and hīendon. Ðā, on ðǣre unstilnesse, onsęndon hīe ǣrendwrecan tō Rome mid gewritum and wēpendre bēne; him fultumes bǣdon, and him gehēton ēaðmōde hīernesse and singāle underðēodnesse, gif hīe him gefultumoden ðæt hīe meahten hiera fīend oferwinnan. Ðā onsęndon hīe him micelne hęre tō fultume, and, sōna ðæs-ðe hīe on ðis īgland cōmon, ðā campedon hīe wið hiera fēondum, and him micel wæl ongeslōgon, and of hiera gemǣrum ādrifon and āflīemdon; and lǣrdon ðæt hīe fæsten worhten him tō gebeorge wið hiera fēondum; and swā, mid micle sige, hām fōron.}
% mdparagraph:P429:end

% mdparagraph:P430:begin
\readingtext{Ðā ðæt ðā ongēaton ðā ǣrran gewinnan, ðæt se Rōmanisca hęre wæs onweg gewiten, ðā cōmon hīe sōna mid sciphęre on hiera landgemǣru, and slōgon ealle and cwealdon ðæt hīe gemētton, and swā-swā rīpe ierðe fortrǣdon and fornōmon, and hīe mid ealle foriermdon.}
% mdparagraph:P430:end

\end{minipage}\par
\reserve{7\baselineskip}
% mdparagraph:P431:begin
\subsection*{Reading help}\BookHeading{nav-ch12-learner-0093}{subsubsection}{Reading help}{0}
% mdparagraph:P431:end

% mdparagraph:P432:begin
Use all the following help while translating and analyzing the passage. Extra historical vocabulary and the expressly supplied constructions are not independent-recall material. Earlier forms retain their ordinary uses; new senses are identified where needed. The note numbers cited below refer to Cook’s original apparatus in §13.
% mdparagraph:P432:end

% teaching: NEW-U12-GRAM-NEW-U12-SRC-01
% support: U12-SRC-HELP-01
% mdparagraph:P433:begin
\textbf{1. Ðā ongunnon ... hęrgian.} Ongunnon, began, takes the three infinitives onwinnan, attack, niman, take, and hęrgian, plunder. Pihtas are the Picts, Scottas the Scots; norðan and westan mean from the north and from the west. Hīe refers to the Britons introduced in the separate context; ǣhta means possessions.
% mdparagraph:P433:end

% teaching: NEW-U12-GRAM-NEW-U12-SRC-02
% support: U12-SRC-HELP-02
% mdparagraph:P434:begin
\textbf{2. and hīe fela gēara iermdon and hīendon.} Iermdon and hīendon mean afflicted and oppressed. Fela gēara means for many years, with genitive after fela. Hīe can be the Britons as object with the raiders continuing as subject, or the raiders as subject with a Briton object understood; the events do not change. Do not infer a unique case from the spelling.
% mdparagraph:P434:end

% teaching: NEW-U12-GRAM-NEW-U12-SRC-03
% support: U12-SRC-HELP-03
% mdparagraph:P435:begin
\textbf{3. Ðā, on ðǣre unstilnesse ... bēne.} On ðǣre unstilnesse, amid that unrest. Onsęndon is the exact prefixed sent; the Britons are now the subject. Ǣrendwrecan, messengers; tō Rome, to Rome. Mid gewritum and wēpendre bēne means with letters and a tearful petition. Wēpendre is a present participle, strong dative feminine singular agreeing with bēne.
% mdparagraph:P435:end

% teaching: NEW-U12-GRAM-NEW-U12-SRC-04
% support: U12-SRC-HELP-04
% mdparagraph:P436:begin
\textbf{4. him fultumes bǣdon ... underðēodnesse.} In him fultumes bǣdon, him names the Romans asked and fultumes is the genitive thing requested: help. Hīernes is obedience; singāl is continual; underðēodnes is subjection. Ēaðmōde and singāle are accusative feminine singular. The promise is made by the Britons to the Romans.
% mdparagraph:P436:end

% teaching: NEW-U12-GRAM-NEW-U12-SRC-05
% teaching: NEW-U12-GRAM-u12-source-reported-condition
% support: U12-SRC-HELP-05
% mdparagraph:P437:begin
\textbf{5. gif hīe him gefultumoden ... oferwinnan.} Gefultumoden is a preterite subjunctive: if they, the Romans, would help them, the Britons. The following ðæt clause gives the purpose, so that they might overcome their enemies. Meahten is the known preterite subjunctive plural of magan; oferwinnan means overcome. This whole reported condition/\allowbreak{}purpose sequence is supported locally.
% mdparagraph:P437:end

% teaching: NEW-U12-GRAM-NEW-U12-SRC-06
% support: U12-SRC-HELP-06
% mdparagraph:P438:begin
\textbf{6. Ðā onsęndon ... tō fultume.} Here the subject of onsęndon is the Romans, while him refers to the Britons. Tō fultume means as help/\allowbreak{}to their aid. The Old English says they sent an army in the active voice; Cook compares a passive Latin expression. Preserve the actual Old English voice when analyzing it.
% mdparagraph:P438:end

% teaching: NEW-U12-GRAM-NEW-U12-SRC-07
% teaching: NEW-U12-REUSED-U12-SRC-HELP-07
% support: U12-SRC-HELP-07
% mdparagraph:P439:begin
\textbf{7. and, sōna ðæs-ðe ... cōmon, ðā.} Sōna ðæs-ðe means as soon as. On ðis īgland has the accusative goal, to/\allowbreak{}onto this island. Hīe refers to the Roman force. The following ðā resumes the narrative.
% mdparagraph:P439:end

% teaching: NEW-U12-GRAM-NEW-U12-SRC-08
% teaching: NEW-U12-REUSED-U12-SRC-HELP-08
% support: U12-SRC-HELP-08
% mdparagraph:P440:begin
\textbf{8. campedon ... āflīemdon.} Campedon, fought; wið hiera fēondum, against their enemies, with the source dative after wið. Him micel wæl ongeslōgon means inflicted great slaughter on them. Of hiera gemǣrum means out of their borders, with dative after of; the borders belong to the Britons, Rome’s allies. Āflīemdon means put to flight, reinforcing ādrifon. The enemy object carries on through the coordinated verbs.
% mdparagraph:P440:end

% teaching: NEW-U12-GRAM-NEW-U12-SRC-09
% teaching: NEW-U12-REUSED-U12-SRC-HELP-09
% support: U12-SRC-HELP-09
% mdparagraph:P441:begin
\textbf{9. lǣrdon ðæt ... fēondum.} Lǣrdon means advised/\allowbreak{}instructed. Fæsten is a fortification; worhten is the known unprefixed wyrcean preterite subjunctive, here make/\allowbreak{}build in the advice clause. Him tō gebeorge means for their own protection. The Romans give the advice and the Britons are to build. The source’s wið + dative remains locally supplied.
% mdparagraph:P441:end

% teaching: NEW-U12-GRAM-NEW-U12-SRC-10
% support: U12-SRC-HELP-10
% mdparagraph:P442:begin
\textbf{10. and swā, mid micle sige, hām fōron.} Sige means victory. Mid micle sige is with great victory: the adjective has the instrumental micle beside the dative noun, as Cook’s note directs. Treat this as a supplied source construction. Hām is the directional home, without a preposition.
% mdparagraph:P442:end

% teaching: NEW-U12-GRAM-NEW-U12-SRC-11
% support: U12-SRC-HELP-11
% mdparagraph:P443:begin
\textbf{11. Ðā ðæt ðā ongēaton ðā ǣrran gewinnan, ðæt.} Read the three ðā forms in order as when, then, the. The first ðæt is an anticipatory object, this/\allowbreak{}that, expanded by the later that-clause. Ongēaton means realized; ǣrran gewinnan means former enemies. Gewinnan is a weak masculine plural noun, distinct from the previously supplied neuter gewinn, trouble.
% mdparagraph:P443:end

% teaching: NEW-U12-GRAM-NEW-U12-SRC-12
% support: U12-SRC-HELP-12
% mdparagraph:P444:begin
\textbf{12. se Rōmanisca hęre wæs onweg gewiten.} Rōmanisca is the weak masculine adjective Roman. Wæs onweg gewiten means had gone away: gewiten is the past participle of intransitive gewītan. This is an intransitive completed-state/\allowbreak{}perfect expression, not a passive. The source has plain-e onweg; retain it beside the earlier selected onwęg.
% mdparagraph:P444:end

% teaching: NEW-U12-GRAM-NEW-U12-SRC-13
% teaching: NEW-U12-REUSED-U12-SRC-HELP-13
% support: U12-SRC-HELP-13
% mdparagraph:P445:begin
\textbf{13. ðā cōmon ... landgemǣru.} The returning subject is the Picts and Scots. Sōna means immediately; mid sciphęre, with a fleet. On hiera landgemǣru has accusative plural for movement into their territory; the territory is the Britons’.
% mdparagraph:P445:end

% teaching: NEW-U12-GRAM-NEW-U12-SRC-14
% support: U12-SRC-HELP-14
% mdparagraph:P446:begin
\textbf{14. slōgon ealle and cwealdon ðæt hīe gemētton.} Cwealdon means killed. Slōgon and cwealdon reinforce the attack. Ealle is all, naturally people in context; the following neuter ðæt is a headless/\allowbreak{}collective all that they encountered. A translation all they encountered need not pretend that ðæt is an overt masculine personal relative. Cook’s Latin note reads cædunt.
% mdparagraph:P446:end

% teaching: NEW-U12-GRAM-NEW-U12-SRC-15
% teaching: NEW-U12-REUSED-U12-SRC-HELP-15
% support: U12-SRC-HELP-15
% mdparagraph:P447:begin
\textbf{15. swā-swā rīpe ierðe fortrǣdon and fornōmon.} Swā-swā means just as/\allowbreak{}like. Rīpe ierðe means ripe crops, a collective feminine accusative phrase. Fortrǣdon, trampled down; fornōmon, destroyed/\allowbreak{}laid waste. The Britons are treated as a crop that is trampled; the object continues from the preceding clauses.
% mdparagraph:P447:end

% teaching: NEW-U12-GRAM-NEW-U12-SRC-16
% support: U12-SRC-HELP-16
% mdparagraph:P448:begin
\textbf{16. and hīe mid ealle foriermdon.} Foriermdon means ruined/\allowbreak{}reduced to misery. Hīe is the Britons as object. Mid ealle is the fixed phrase completely/\allowbreak{}utterly; ealle is instrumental here. The complete phrase is supplied.
% mdparagraph:P448:end

\reserve{7\baselineskip}
% mdparagraph:P449:begin
\subsection*{S01}\BookHeading{nav-ch12-learner-0094}{subsubsection}{S01}{0}
% mdparagraph:P449:end

\par\begin{minipage}{\linewidth}\raggedright
% item: NEW-U12-S01
% mdparagraph:P450:begin
Translate the complete Picts and Scots extract, using the printed local helps and original notes. Preserve the renewed attack after the Roman withdrawal.
% mdparagraph:P450:end

\end{minipage}\par\addvspace{4pt}

\reserve{7\baselineskip}
% mdparagraph:P451:begin
\subsection*{S02}\BookHeading{nav-ch12-learner-0095}{subsubsection}{S02}{0}
% mdparagraph:P451:end

\par\begin{minipage}{\linewidth}\raggedright
% item: NEW-U12-S02
% mdparagraph:P452:begin
Put the narrated stages in order: renewed attack; appeal to Rome; first raids; Roman withdrawal; Roman relief. Explain the cause of the renewed attack as the narrator presents it.
% mdparagraph:P452:end

\end{minipage}\par\addvspace{4pt}

\reserve{7\baselineskip}
% mdparagraph:P453:begin
\subsection*{S03}\BookHeading{nav-ch12-learner-0096}{subsubsection}{S03}{0}
% mdparagraph:P453:end

\par\begin{minipage}{\linewidth}\raggedright
% item: NEW-U12-S03
% mdparagraph:P454:begin
In Ðā ðæt ðā ongēaton ðā ǣrran gewinnan, identify the three occurrences of ðā and the role of the first ðæt. Use the printed source help; distinguish grammatical form from clause function.
% mdparagraph:P454:end

\end{minipage}\par\addvspace{4pt}

\reserve{7\baselineskip}
% mdparagraph:P455:begin
\subsection*{S04}\BookHeading{nav-ch12-learner-0097}{subsubsection}{S04}{0}
% mdparagraph:P455:end

\par\begin{minipage}{\linewidth}\raggedright
% item: NEW-U12-S04
% mdparagraph:P456:begin
Who sends messengers to Rome, who sends the army, and who fights the enemies? Explain how the changing subjects are recovered from context, using the original notes where indicated.
% mdparagraph:P456:end

\end{minipage}\par\addvspace{4pt}

\reserve{7\baselineskip}
% mdparagraph:P457:begin
\subsection*{S05}\BookHeading{nav-ch12-learner-0098}{subsubsection}{S05}{0}
% mdparagraph:P457:end

\par\begin{minipage}{\linewidth}\raggedright
% item: NEW-U12-S05
% mdparagraph:P458:begin
Parse wæs onweg gewiten. State which event precedes which other event, and explain why this is not an agentless passive.
% mdparagraph:P458:end

\end{minipage}\par\addvspace{4pt}

\reserve{7\baselineskip}
% mdparagraph:P459:begin
\subsection*{S06}\BookHeading{nav-ch12-learner-0099}{subsubsection}{S06}{0}
% mdparagraph:P459:end

\par\begin{minipage}{\linewidth}\raggedright
% item: NEW-U12-S06
% mdparagraph:P460:begin
Cook’s note on Ðā onsęndon hīe him micelne hęre points out an active/\allowbreak{}passive difference from the Latin source. Identify the Old English subject, recipient and object. Give an English passive rendering that preserves their roles; no Latin or new Old English composition is required.
% mdparagraph:P460:end

\end{minipage}\par\addvspace{4pt}

\reserve{7\baselineskip}
% mdparagraph:P461:begin
\subsection*{S07}\BookHeading{nav-ch12-learner-0100}{subsubsection}{S07}{0}
% mdparagraph:P461:end

\par\begin{minipage}{\linewidth}\raggedright
% item: NEW-U12-S07
% mdparagraph:P462:begin
In lǣrdon ðæt hīe fæsten worhten him tō gebeorge, what action is recommended? Is that action earlier than, simultaneous with, or prospective from the instruction? Does this clause alone state that the work was completed?
% mdparagraph:P462:end

\end{minipage}\par\addvspace{4pt}

\reserve{7\baselineskip}
% mdparagraph:P463:begin
\subsection*{S08}\BookHeading{nav-ch12-learner-0101}{subsubsection}{S08}{0}
% mdparagraph:P463:end

\par\begin{minipage}{\linewidth}\raggedright
% item: NEW-U12-S08
% mdparagraph:P464:begin
Explain the comparison with rīpe ierðe in the final attack. Who are the victims, and how does the comparison contribute to the account? Use the lexical help rather than assuming a new ordinary crop meaning for the people.
% mdparagraph:P464:end

\end{minipage}\par\addvspace{4pt}

\reserve{8\baselineskip}
% mdparagraph:P465:begin
\section*{13. Cook’s original notes}\BookHeading{nav-ch12-learner-0102}{subsection}{13. Cook’s original notes}{1}
% mdparagraph:P465:end

\begingroup\fontsize{10}{12.8}\selectfont
\setlength{\parskip}{1.5pt}
% mdparagraph:P466:begin
The following notes reproduce the apparatus attached to this extract. Their English commentary and questions are Cook’s editorial material, not additional exercises. In particular, p. 145 note 2 preserves nineteenth-century ethnological speculation and terminology as historical apparatus; it is not a present-day account taught by this course. The fragment at the top of p. 145 continuing p. 144 note 10 is outside the extract and is not included. Notes 16–19 on p. 146 likewise belong to later text.
% mdparagraph:P466:end

% mdparagraph:P467:begin
The Old English above remains uninterrupted. This anchor list locates the original note cues; where one note is attached twice, both anchors are given.
% mdparagraph:P467:end

% mdtable:T71:begin
\par\begingroup
\fontsize{10.5}{13.2}\selectfont
\setlength{\LTleft}{0pt}\setlength{\LTright}{\fill}
\begin{longtable}{@{}>{\raggedright\arraybackslash}p{49.35mm}>{\raggedright\arraybackslash}p{71.84mm}@{}}
\toprule
\textbf{Printed page /\allowbreak{} note} & \textbf{Attached word or phrase} \\
\midrule
\endfirsthead
\toprule
\textbf{Printed page /\allowbreak{} note} & \textbf{Attached word or phrase} \\
\midrule
\endhead
\midrule
\multicolumn{2}{r}{\footnotesize Continued on the following page.}\\
\endfoot
\bottomrule
\endlastfoot
145.1 & Ðā \\
145.2 & Pihtas \\
145.3 & ǣrendwrecan \\
145.4 & gewritum \\
145.5 & fultumes \\
145.6 & underðēodnesse \\
146.1 & fīend \\
146.2 & Ðā \\
146.3 & onsęndon \\
146.4 & hīe \\
146.5 & hiera \\
146.6 & āflīemdon \\
146.7 & fæsten \\
146.8 & him \\
146.9 & sige \\
146.10 & fōron \\
146.11 & Ðā \\
146.12 & gewinnan \\
146.13 & slōgon \\
146.13 & cwealdon \\
146.14 & ierðe \\
146.15 & ealle \\
\end{longtable}
\par\endgroup

% mdtable:T71:end
% source-note: U12-SRC-P145-N01
\reserve{5\baselineskip}
% mdparagraph:P468:begin
\textbf{Printed p. 145, note 1.}
% mdparagraph:P468:end

% mdparagraph:P469:begin
“[This account] is pieced together as an abridgment of Gildas, xi.–xvi.; but the turgidity of his style is chastened, and his faulty grammar in several places corrected” (Moberly, pp. 26, 27).
% mdparagraph:P469:end

% source-note: U12-SRC-P145-N02
\reserve{5\baselineskip}
% mdparagraph:P470:begin
\textbf{Printed p. 145, note 2.}
% mdparagraph:P470:end

% mdparagraph:P471:begin
On the Picts the last edition of Chambers’ Encyclopædia remarks: “Four hypotheses have been formed in regard to the language and origin of the Picts. The first, started by Pinkerton and put by Sir Walter Scott into the mouth of the ‘Antiquary,’ is that they were Teutons, speaking a Gothic dialect; the second, maintained by Dr. Skene, is that they were Gaelic-speaking Celts; . . . the third, due to Professor Rhys, is that the Picts were non-Aryans, whose language was overlaid by loans from Welsh and Irish; and the fourth, held by two of the most eminent Celtic scholars of the day, Professor Windisch and Dr. Whitley Stokes, is that they were Celts, but more nearly allied to the Cymry than to the Gael. . . . The conclusion to which we come is that the Picts, whatever traces they show of a non-Aryan racial element, . . . spoke a Celtic language belonging to a branch of Celtic allied to the Cymric, . . . and that this dialect of the Gallo-Cymric stock was a wave of Celtic speech from the continent previous to the Gaulish which held England when Cæsar entered Britain.”
% mdparagraph:P471:end

% source-note: U12-SRC-P145-N03
\reserve{5\baselineskip}
% mdparagraph:P472:begin
\textbf{Printed p. 145, note 3.}
% mdparagraph:P472:end

% mdparagraph:P473:begin
Lat. legatos.
% mdparagraph:P473:end

% source-note: U12-SRC-P145-N04
\reserve{5\baselineskip}
% mdparagraph:P474:begin
\textbf{Printed p. 145, note 4.}
% mdparagraph:P474:end

% mdparagraph:P475:begin
Lat. epistolis.
% mdparagraph:P475:end

% source-note: U12-SRC-P145-N05
\reserve{5\baselineskip}
% mdparagraph:P476:begin
\textbf{Printed p. 145, note 5.}
% mdparagraph:P476:end

% mdparagraph:P477:begin
See 156. b.
% mdparagraph:P477:end

% source-note: U12-SRC-P145-N06
\reserve{5\baselineskip}
% mdparagraph:P478:begin
\textbf{Printed p. 145, note 6.}
% mdparagraph:P478:end

% mdparagraph:P479:begin
This pair of phrases renders subjectionem continuam. What parallel to the use of such synonymous terms may be found in the English Prayer-Book? How is it to be accounted for?
% mdparagraph:P479:end

% source-note: U12-SRC-P146-N01
\reserve{5\baselineskip}
% mdparagraph:P480:begin
\textbf{Printed p. 146, note 1.}
% mdparagraph:P480:end

% mdparagraph:P481:begin
See 46.
% mdparagraph:P481:end

% source-note: U12-SRC-P146-N02
\reserve{5\baselineskip}
% mdparagraph:P482:begin
\textbf{Printed p. 146, note 2.}
% mdparagraph:P482:end

% mdparagraph:P483:begin
This may be anywhere between A.D. 388 and 420. See Moberly, p. 27.
% mdparagraph:P483:end

% source-note: U12-SRC-P146-N03
\reserve{5\baselineskip}
% mdparagraph:P484:begin
\textbf{Printed p. 146, note 3.}
% mdparagraph:P484:end

% mdparagraph:P485:begin
This clause translates Quibus mox legio destinatur armata. Note the use of the active for the passive, which also appears in other sentences of the context.
% mdparagraph:P485:end

% source-note: U12-SRC-P146-N04
\reserve{5\baselineskip}
% mdparagraph:P486:begin
\textbf{Printed p. 146, note 4.}
% mdparagraph:P486:end

% mdparagraph:P487:begin
The legionaries, apparently.
% mdparagraph:P487:end

% source-note: U12-SRC-P146-N05
\reserve{5\baselineskip}
% mdparagraph:P488:begin
\textbf{Printed p. 146, note 5.}
% mdparagraph:P488:end

% mdparagraph:P489:begin
Of the Britons; Lat. sociorum finibus.
% mdparagraph:P489:end

% source-note: U12-SRC-P146-N06
\reserve{5\baselineskip}
% mdparagraph:P490:begin
\textbf{Printed p. 146, note 6.}
% mdparagraph:P490:end

% mdparagraph:P491:begin
Lat. expulit. See above, p. 145, n. 6.
% mdparagraph:P491:end

% source-note: U12-SRC-P146-N07
\reserve{5\baselineskip}
% mdparagraph:P492:begin
\textbf{Printed p. 146, note 7.}
% mdparagraph:P492:end

% mdparagraph:P493:begin
Lat. murum.
% mdparagraph:P493:end

% source-note: U12-SRC-P146-N08
\reserve{5\baselineskip}
% mdparagraph:P494:begin
\textbf{Printed p. 146, note 8.}
% mdparagraph:P494:end

% mdparagraph:P495:begin
See 184. a.
% mdparagraph:P495:end

% source-note: U12-SRC-P146-N09
\reserve{5\baselineskip}
% mdparagraph:P496:begin
\textbf{Printed p. 146, note 9.}
% mdparagraph:P496:end

% mdparagraph:P497:begin
Lat. triumpho. See 175.
% mdparagraph:P497:end

% source-note: U12-SRC-P146-N10
\reserve{5\baselineskip}
% mdparagraph:P498:begin
\textbf{Printed p. 146, note 10.}
% mdparagraph:P498:end

% mdparagraph:P499:begin
A passage of the Latin is here omitted in the translation, describing the construction of the (earthen) wall, between the Friths of Forth and Clyde.
% mdparagraph:P499:end

% source-note: U12-SRC-P146-N11
\reserve{5\baselineskip}
% mdparagraph:P500:begin
\textbf{Printed p. 146, note 11.}
% mdparagraph:P500:end

% mdparagraph:P501:begin
These three ðā’s respectively = when, then, and the.
% mdparagraph:P501:end

% source-note: U12-SRC-P146-N12
\reserve{5\baselineskip}
% mdparagraph:P502:begin
\textbf{Printed p. 146, note 12.}
% mdparagraph:P502:end

% mdparagraph:P503:begin
Lat. inimici.
% mdparagraph:P503:end

% source-note: U12-SRC-P146-N13
\reserve{5\baselineskip}
% mdparagraph:P504:begin
\textbf{Printed p. 146, note 13.}
% mdparagraph:P504:end

% mdparagraph:P505:begin
Lat. cædunt. See above, p. 145, n. 6.
% mdparagraph:P505:end

% source-note: U12-SRC-P146-N14
\reserve{5\baselineskip}
% mdparagraph:P506:begin
\textbf{Printed p. 146, note 14.}
% mdparagraph:P506:end

% mdparagraph:P507:begin
Lat. segetem.
% mdparagraph:P507:end

% source-note: U12-SRC-P146-N15
\reserve{5\baselineskip}
% mdparagraph:P508:begin
\textbf{Printed p. 146, note 15.}
% mdparagraph:P508:end

% mdparagraph:P509:begin
Mid ealle = completely. See 175.
% mdparagraph:P509:end

% mdparagraph:P510:begin
The transcription preserves Cook’s letters, written lengths and lexical spellings. Line wrapping and line-end division have been joined; original note cues are represented by the separate anchor list. \textbf{Rome}, \textbf{onweg}, \textbf{hęre}, \textbf{ādrifon} and \textbf{Rōmanisca} retain their source spellings. Page 146 note 13 reads \textbf{Lat. cædunt}. The supplied modern explanations in §12 are separate from the original notes.
% mdparagraph:P510:end

\reserve{7\baselineskip}
% mdparagraph:P511:begin
\section*{14. Sources for the lesson}\BookHeading{nav-ch12-learner-0103}{subsection}{14. Sources for the lesson}{1}
% mdparagraph:P511:end

% mdparagraph:P512:begin
The exercise prompts are from Albert S. Cook, \emph{Exercises in Old English} (1895), Group VI, Exercises 19–21, printed pp. 9–11; the registered scan uses PDF pp. 17–19. The reading and grammar are from Cook, \emph{A First Book in Old English}, second edition (1895). In that registered scan a printed page corresponds to the one-based PDF page sixteen higher. The reading locators and exclusions are given in §§12–13.
% mdparagraph:P512:end

% mdparagraph:P513:begin
For noun and adjective morphology, compare Cook §§43–60, printed pp. 28–41, and the vocabulary entries for the individual words. Fōt and fēond are treated at §§46–46.3; Wright, \emph{Old English Grammar} (1914), §§416–417, printed p. 202 /\allowbreak{} PDF p. 222, independently supports the fēond class and genitive plural fēonda. Cook §44.2 gives the hęre alternatives; §57.3 gives the optional syncope of ēadig. His vocabulary explicitly distinguishes the noun and adjective īsern.
% mdparagraph:P513:end

% mdparagraph:P514:begin
For the defined verb systems, commands and nonfinite bases, compare Cook §§93–95 and the individual strong and weak classes; §127 for āgan and §139 for willan and its negative forms. The selected dental combinations follow Cook §§33–35 and Wright §476, printed p. 248. Cȳst/\allowbreak{}cȳð and the attested cȳdde alternative are checked with Bosworth–Toller, s.v. cȳðan.
% mdparagraph:P514:end

% mdparagraph:P515:begin
The treatment of time and reported viewpoint uses Cook §§186–188 and §§192–196, with Fulk (2020), §§45–46 and §§48–52. Bosworth–Toller, \emph{An Anglo-Saxon Dictionary}, main entry habban IV–V, printed p. 497a, and the Supplement, habban B, printed p. 493a, provide evidence for both uninflected and object-agreeing participles. The chapter teaches only the selected constructions described above.
% mdparagraph:P515:end

% mdparagraph:P516:begin
Prepositional and case uses draw on Cook §172 and §§160–161, 164, with the individual lexical evidence. Bosworth–Toller’s ēhtan entry includes the same pursuing-in-this-city expression; its gemang entry supports the assigned sheep-and-wolves comparison and accusative motion phrase. Its entries mā and swā support the quantitative and comparative distinctions. The dictionary’s mild-heort entry, printed p. 687b, supports merciful in the learning saying; stán-weall, printed p. 912a, supports stone wall. In teaching lemmas, dictionary division hyphens are joined and the evidenced acute quantity of stán-weall is written with the course’s macron, stānweall. No unmarked vowel length is supplied by guesswork.
% mdparagraph:P516:end

\endgroup

% book-source-body:end
\par\label{end-ch12-learner}
\clearpage\endgroup

\BookRecto
\begingroup
\renewcommand{\runningtitle}{Unit 13. Clause order and reference}
\hyphenpenalty=10000
\exhyphenpenalty=10000
\emergencystretch=3em
\newcommand{\itempara}[2]{\par\addvspace{3pt}\noindent\hangindent=1.8em\hangafter=1\makebox[1.8em][l]{#1.}#2\par}
\newcommand{\reserve}[1]{\par\penalty-100\begingroup\dimen0=\pagegoal\advance\dimen0 by-\pagetotal\ifdim\dimen0<#1\newpage\fi\endgroup}
% book-source-body:begin

\addfontfeatures{Ligatures=TeXOff}
\raggedright
\booktitle{\BookComponent{ch13-learner}{Unit 13. Clause order and reference}UNIT 13}{Clause order and reference}

% mdparagraph:P001:begin
A sentence may begin with its subject, an object, an adverbial, or a dependent clause. Its opening words do not by themselves identify the participants. This unit brings clause structure, agreement and reference together, then applies them to a connected narrative.
% mdparagraph:P001:end

% mdparagraph:P002:begin
Sections 6–8 supply the vocabulary, forms and constructions required by the practice. The four Cook sets are retained in full: 15, 13, 12 and 12 items.
% mdparagraph:P002:end

% mdparagraph:P003:begin
Examples 1–12, Exercises F, A, C, R and S, and The bishop and the monk are new teaching material. The final reading preserves the selected Old English of Cook’s The Passing of Chad; its English orientation and local helps are separate additions. The original notes remain at the end, with their exact attachment points.
% mdparagraph:P003:end

% mdparagraph:P004:begin
\textbf{Scope.} “Production” means use in the senses and constructions taught here or earlier. “Reference” forms are for recognition; their display does not add a new meaning or construction. “Supplied help” stays visible while a task is done and is not a test of unaided recall. In tables, 1, 2 and 3 mean first, second and third person; sg and pl mean singular and plural. Nominative, genitive, dative and accusative are abbreviated N, G, D and A only in diagnostic headings. Preserve the printed distinction between e and ę and the marked vowel lengths.
% mdparagraph:P004:end

\reserve{7\baselineskip}
% mdparagraph:P005:begin
\section*{1. Clause roles and the first constituent}\BookHeading{nav-ch13-learner-0001}{subsection}{1. Clause roles and the first constituent}{1}
% mdparagraph:P005:end

% teaching: NEW-U13-GRAM-U13-CLAUSE-ROLE-PARSING
% teaching: NEW-U13-GRAM-U13-MAIN-FRONTING-ORDER
% mdparagraph:P006:begin
Find the finite verb and its subject before deciding how to translate the order. The nominative marks the subject; an accusative may be a direct object, but its role must also fit the verb. Agreement and the verb’s construction help when endings coincide. A constituent is a whole phrase functioning as a unit: se munuc and ðone sang each contain two words but each is one constituent.
% mdparagraph:P006:end

\reserve{8\baselineskip}
% example: NEW-U13-EXAMPLE01
\reserve{5\baselineskip}
% mdparagraph:P007:begin
\textbf{Example 1.}
% mdparagraph:P007:end

\reserve{5\baselineskip}
% mdparagraph:P008:begin
\textbf{Se munuc hīerde ðone sang.}
% mdparagraph:P008:end

% mdparagraph:P009:begin
The monk heard the song.
% mdparagraph:P009:end

% mdparagraph:P010:begin
Se munuc is nominative; ðone sang is accusative. The finite verb hīerde agrees with the singular subject.
% mdparagraph:P010:end

\reserve{8\baselineskip}
% example: NEW-U13-EXAMPLE02
\reserve{5\baselineskip}
% mdparagraph:P011:begin
\textbf{Example 2.}
% mdparagraph:P011:end

\reserve{5\baselineskip}
% mdparagraph:P012:begin
\textbf{Ðone sang hīerde se munuc.}
% mdparagraph:P012:end

% mdparagraph:P013:begin
The monk heard the song.
% mdparagraph:P013:end

% mdparagraph:P014:begin
The\BookPoint{bk-ch13-learner-54a05581dff2} object phrase comes first. Its accusative ending still makes it the object. The finite verb follows that complete phrase, not the first written word.
% mdparagraph:P014:end

\reserve{8\baselineskip}
% example: NEW-U13-EXAMPLE03
\reserve{5\baselineskip}
% mdparagraph:P015:begin
\textbf{Example 3.}
% mdparagraph:P015:end

\reserve{5\baselineskip}
% mdparagraph:P016:begin
\textbf{Ðā hīerde se munuc ðone sang.}
% mdparagraph:P016:end

% mdparagraph:P017:begin
Then the monk heard the song.
% mdparagraph:P017:end

% mdparagraph:P018:begin
The initial adverb ðā means “then.” Here the finite verb precedes the noun subject. This is a common main-clause order after an initial constituent.
% mdparagraph:P018:end

% mdparagraph:P019:begin
Finite-second\BookPoint{bk-ch13-learner-30eeeceef5e9} order is a tendency, not a requirement. It concerns constituents, not word counts. A coordinating and, ac or oððe does not itself count as the first constituent of the following clause. Fronting gives a constituent a prominent position, but the exact effect depends on context; no change of case follows merely from a change of order.
% mdparagraph:P019:end

\reserve{7\baselineskip}
% mdparagraph:P020:begin
\section*{2. Dependent clauses and limits on order rules}\BookHeading{nav-ch13-learner-0002}{subsection}{2. Dependent clauses and limits on order rules}{1}
% mdparagraph:P020:end

% teaching: NEW-U13-DEPENDENT-ORDER-HELP
% teaching: NEW-U13-GRAM-U13-DEPENDENT-ORDER-TENDENCIES
% teaching: NEW-U13-GRAM-U13-CORRELATIVE-CLAUSE-REFERENCE
% teaching: NEW-U13-GRAM-U13-PRONOUN-POSITION-LIMITS
% mdparagraph:P021:begin
A\BookPoint{bk-ch13-learner-d792ea8f2502} dependent clause has its own finite verb and subject. Its finite verb often comes late, frequently last. The conjunction or relative marker begins the clause; it is not its subject unless the construction makes it so. Identify the boundary of each clause before comparing orders.
% mdparagraph:P021:end

\reserve{8\baselineskip}
% example: NEW-U13-EXAMPLE04
\reserve{5\baselineskip}
% mdparagraph:P022:begin
\textbf{Example 4.}
% mdparagraph:P022:end

\reserve{5\baselineskip}
% mdparagraph:P023:begin
\textbf{Ðā hē ðone sang hīerde, ðā ēode hē tō ðǣm bisceope.}
% mdparagraph:P023:end

% mdparagraph:P024:begin
When he heard the song, he went to the bishop.
% mdparagraph:P024:end

% mdparagraph:P025:begin
The first ðā means “when” and introduces the dependent clause ending in hīerde. The second means “then” and belongs to the main clause. Both occurrences of hē refer to the person already in view; the order alone does not supply his identity.
% mdparagraph:P025:end

\reserve{8\baselineskip}
% example: NEW-U13-EXAMPLE09
\reserve{5\baselineskip}
% mdparagraph:P026:begin
\textbf{Example 9.}
% mdparagraph:P026:end

\reserve{5\baselineskip}
% mdparagraph:P027:begin
\textbf{Se munuc, ðe se bisceop grētte, ēode tō ðǣre cirican.}
% mdparagraph:P027:end

% mdparagraph:P028:begin
The monk whom the bishop greeted went to the church.
% mdparagraph:P028:end

% mdparagraph:P029:begin
Invariant ðe introduces the relative clause and refers back to se munuc. The monk is the subject of ēode, but is the person greeted inside the relative; se bisceop is the relative clause’s subject. Keep the two roles separate.
% mdparagraph:P029:end

\reserve{8\baselineskip}
% example: NEW-U13-EXAMPLE10
\reserve{5\baselineskip}
% mdparagraph:P030:begin
\textbf{Example 10.}
% mdparagraph:P030:end

\reserve{5\baselineskip}
% mdparagraph:P031:begin
\textbf{Ðā se bisceop cōm, ðā sungon ðā brōðor.}
% mdparagraph:P031:end

% mdparagraph:P032:begin
When the bishop came, the brothers sang.
% mdparagraph:P032:end

% mdparagraph:P033:begin
The first ðā introduces “when the bishop came”; the second is the main-clause adverb “then.” The last ðā is the plural demonstrative with brōðor. The plural subject agrees with sungon.
% mdparagraph:P033:end

% mdparagraph:P034:begin
After\BookPoint{bk-ch13-learner-254b66d59004} an initial adverbial a subject pronoun may precede the finite verb. The following short source fragment gives an actual instance. It belongs to the continuation of the Chad passage, outside the continuous extract assigned below.
% mdparagraph:P034:end

% source-example: NEW-U13-SOURCE-ORDER-FRAGMENT
\reserve{5\baselineskip}
% mdparagraph:P035:begin
\textbf{Æfter ðon hē underðēodde}
% mdparagraph:P035:end

% mdparagraph:P036:begin
After that he added
% mdparagraph:P036:end

% teaching: NEW-U13-SOURCE-ORDER-HELP
% support: NEW-U13-SOURCE-ORDER-HELP
% mdparagraph:P037:begin
\textbf{Order help.} This complete fragment is supplied as corroborative evidence outside the assigned continuous reading. Æfter ðon means after that; hē is Chad; underðēodde means added to his speech (Cook note13: Latin subjunxit). Identify its order only; no independent translation or lexical recall is required. The fronted temporal phrase is followed by subject pronoun, then finite verb.
% mdparagraph:P037:end

% teaching: NEW-U13-VOC-SE
% teaching: NEW-U13-VOC-AEFTER
% teaching: NEW-U13-VOC-HE
% teaching: NEW-U13-VOC-UNDERTHEODAN
% teaching: NEW-U13-PAR-SE
% teaching: NEW-U13-PAR-AEFTER
% teaching: NEW-U13-PAR-HE
% teaching: NEW-U13-PAR-UNDERTHEODAN
% teaching: NEW-U13-GRAM-u13-source-after-that
% mdparagraph:P038:begin
The initial phrase Æfter ðon is followed by hē and then underðēodde. Thus even counting constituents does not turn main-clause finite-second order into a universal rule.
% mdparagraph:P038:end

% teaching: NEW-U13-SOURCE-DEPENDENT-ORDER
% mdparagraph:P039:begin
The assigned reading also contains a dependent clause whose finite verb is followed by another constituent:
% mdparagraph:P039:end

% source-example: NEW-U13-SOURCE-DEPENDENT-ORDER
\reserve{5\baselineskip}
% mdparagraph:P040:begin
\textbf{ðæt hē ǣrest gehīerde fram ēastsūðdǣle heofones}
% mdparagraph:P040:end

% support: NEW-U13-DEPENDENT-ORDER-HELP
% mdparagraph:P041:begin
\textbf{Order help.} This is a fragment of the assigned Chad reading: “that he first heard [the sound and song] from the south-eastern part of the sky.” Hē is Owini; gehīerde is the finite preterite “heard”; fram ēastsūðdǣle heofones names the origin of the sound and follows that verb inside the content clause. The objects have been fronted earlier in the full sentence, as Cook note 152.8 explains; they are supplied in brackets only in the English help. Identify the order with this help; no independent recall of the supplied forms is required.
% mdparagraph:P041:end

% teaching: NEW-U13-VOC-EASTSUTHDAEL
% teaching: NEW-U13-VOC-AEREST
% teaching: NEW-U13-SOURCE-HELPS
% teaching: NEW-U13-PAR-AEREST
% mdparagraph:P042:begin
The\BookPoint{bk-ch13-learner-672058608ca5} phrase beginning fram follows gehīerde inside the ðæt clause. “Often late” therefore does not mean “always last.” Cook 25.7 nevertheless asks for a final verb in its specified relative clause. A stated exercise target can select one order without declaring every other order ungrammatical.
% mdparagraph:P042:end

\reserve{7\baselineskip}
% mdparagraph:P043:begin
\section*{3. Pronouns and sustained reference}\BookHeading{nav-ch13-learner-0003}{subsection}{3. Pronouns and sustained reference}{1}
% mdparagraph:P043:end

% teaching: NEW-U13-GRAM-U13-REFERENCE-RESOLUTION-AMBIGUITY
% mdparagraph:P044:begin
A\BookPoint{bk-ch13-learner-1c8bf8495a8f} pronoun agrees with the grammatical gender and number of its antecedent. Its case expresses its role in its own clause. Masculine hē or hine can therefore refer to a masculine thing; English may require “it.” Agreement restricts the possible antecedents but need not select one uniquely.
% mdparagraph:P044:end

\reserve{8\baselineskip}
% example: NEW-U13-EXAMPLE05
\reserve{5\baselineskip}
% mdparagraph:P045:begin
\textbf{Example 5.}
% mdparagraph:P045:end

\reserve{5\baselineskip}
% mdparagraph:P046:begin
\textbf{Se sang wæs swēte. Hē wæs ēac wynsum.}
% mdparagraph:P046:end

% mdparagraph:P047:begin
The song was sweet. It was also pleasant.
% mdparagraph:P047:end

% mdparagraph:P048:begin
Sang is masculine. Hē refers to the song, so English uses “it.” The pronoun’s nominative case belongs to its own clause.
% mdparagraph:P048:end

\reserve{8\baselineskip}
% example: NEW-U13-EXAMPLE06
\reserve{5\baselineskip}
% mdparagraph:P049:begin
\textbf{Example 6.}
% mdparagraph:P049:end

\reserve{5\baselineskip}
% mdparagraph:P050:begin
\textbf{Se bisceop geseah ðone munuc and grētte hine.}
% mdparagraph:P050:end

% mdparagraph:P051:begin
The bishop saw the monk and greeted him.
% mdparagraph:P051:end

% mdparagraph:P052:begin
The second finite verb shares the bishop as its subject. Hine is accusative and refers to the monk, who receives the greeting. The preceding direct object does not become the next subject automatically.
% mdparagraph:P052:end

\reserve{8\baselineskip}
% example: NEW-U13-EXAMPLE07
\reserve{5\baselineskip}
% mdparagraph:P053:begin
\textbf{Example 7.}
% mdparagraph:P053:end

\reserve{5\baselineskip}
% mdparagraph:P054:begin
\textbf{Se bisceop geseah ðone munuc. Hē wæs on ðǣm hūse.}
% mdparagraph:P054:end

% mdparagraph:P055:begin
The bishop saw the monk. He was in the house.
% mdparagraph:P055:end

% mdparagraph:P056:begin
Both\BookPoint{bk-ch13-learner-848f09dddf39} men are masculine singular. Hē can refer to either without a fuller context. Neither the nearest noun nor the previous subject is an infallible rule. Repeating the intended noun phrase removes the ambiguity.
% mdparagraph:P056:end

% mdparagraph:P057:begin
In a longer passage, keep a separate record for each participant and for any thing later named by a pronoun. Recheck the reference when a new subject, apposition or quotation appears. A statement such as hē, se bisceop makes the reference explicit at that point; it need not give every earlier hē the same referent.
% mdparagraph:P057:end

\reserve{7\baselineskip}
% mdparagraph:P058:begin
\section*{4. Coordination and omitted material}\BookHeading{nav-ch13-learner-0004}{subsection}{4. Coordination and omitted material}{1}
% mdparagraph:P058:end

% teaching: NEW-U13-GRAM-U13-COORDINATION-PREDICATE-RECOVERY
% mdparagraph:P059:begin
Two\BookPoint{bk-ch13-learner-4edce9a3cc31} coordinated predicates can share an expressed subject, as in Example 6. Two coordinated clauses may instead have different subjects and omit a repeated predicate. Recover only what the syntax and context support.
% mdparagraph:P059:end

\reserve{8\baselineskip}
% example: NEW-U13-EXAMPLE11
\reserve{5\baselineskip}
% mdparagraph:P060:begin
\textbf{Example 11.}
% mdparagraph:P060:end

\reserve{5\baselineskip}
% mdparagraph:P061:begin
\textbf{Se bisceop wæs on ðǣm hūse, and se munuc on ðǣre cirican.}
% mdparagraph:P061:end

% mdparagraph:P062:begin
The bishop was in the house, and the monk in the church.
% mdparagraph:P062:end

% mdparagraph:P063:begin
The\BookPoint{bk-ch13-learner-9ed625d4233e} second clause has its own subject, se munuc, and omits the repeated copula. Both subjects here are singular. If a restored verb belongs to a plural subject, it must have plural agreement; recovering the lexical predicate does not license copying a singular ending.
% mdparagraph:P063:end

% mdparagraph:P064:begin
Coordinate clauses remain separate for agreement and reference. The same caution applies when one clause contains several verbs: a shared subject is recovered from the construction, not from proximity to the last noun.
% mdparagraph:P064:end

\reserve{7\baselineskip}
% mdparagraph:P065:begin
\section*{5. Speech and long sentences}\BookHeading{nav-ch13-learner-0005}{subsection}{5. Speech and long sentences}{1}
% mdparagraph:P065:end

% teaching: NEW-U13-GRAM-U13-QUOTED-REFERENCE
% teaching: NEW-U13-GRAM-U13-LONG-SENTENCE-ANALYSIS
% teaching: NEW-U13-GRAM-U13-NEGATION-QUESTION-ORDER
% teaching: NEW-U13-GRAM-U13-SUPPORTED-LONG-FRONTING
\reserve{8\baselineskip}
% example: NEW-U13-EXAMPLE08
\reserve{5\baselineskip}
% mdparagraph:P066:begin
\textbf{Example 8.}
% mdparagraph:P066:end

\reserve{5\baselineskip}
% mdparagraph:P067:begin
\textbf{Se bisceop cwæð: “Ic wille mid ēow sprecan.”}
% mdparagraph:P067:end

% mdparagraph:P068:begin
The bishop said, “I want to speak with you.”
% mdparagraph:P068:end

% mdparagraph:P069:begin
Within the quotation, Ic is the bishop and ēow addresses several people. Sprecan with mid and the dative means “speak with”; the plain infinitive follows wille.
% mdparagraph:P069:end

\reserve{8\baselineskip}
% example: NEW-U13-EXAMPLE12
\reserve{5\baselineskip}
% mdparagraph:P070:begin
\textbf{Example 12.}
% mdparagraph:P070:end

\reserve{5\baselineskip}
% mdparagraph:P071:begin
\textbf{Hē sægde ðæt se sang swēte wǣre.}
% mdparagraph:P071:end

% mdparagraph:P072:begin
He said that the song was sweet.
% mdparagraph:P072:end

% mdparagraph:P073:begin
Ðæt introduces reported content. Se sang is its subject and wǣre its finite verb. The report subjunctive does not decide who an otherwise ambiguous pronoun refers to.
% mdparagraph:P073:end

% mdparagraph:P074:begin
Direct\BookPoint{bk-ch13-learner-d1987196da60} speech has its own first and second persons. In an indirect report, a third-person pronoun must be resolved from the reporting context. Questions and negatives keep their previously taught constructions: a yes/\allowbreak{}no question may begin with the finite verb; an interrogative word has a role within its question; ne belongs with the finite verb, and negative concord remains possible. No new negative-copula contraction is required here.
% mdparagraph:P074:end

% mdparagraph:P075:begin
For\BookPoint{bk-ch13-learner-ce6c8914b5b0} a long sentence, first mark each finite verb. Then bracket conjunction and relative clauses, identify each subject, and attach the remaining complements to the verb that governs them. Resolve pronouns only after these relations are clear. Finally restore the temporal sequence and translate in connected English. Coordination, parenthesis and direct speech each need their own boundaries.
% mdparagraph:P075:end

% mdparagraph:P076:begin
One\BookPoint{bk-ch13-learner-52907367739b} supported construction in the Chad reading has objects placed before a reporting clause even though their verb is inside the following content clause. Cook’s note identifies that dependence. Use the printed help to analyse it; the unit does not ask you to create new sentences by freely moving material out of embedded clauses.
% mdparagraph:P076:end

\reserve{7\baselineskip}
% mdparagraph:P077:begin
\section*{6. Nouns and adjectives}\BookHeading{nav-ch13-learner-0006}{subsection}{6. Nouns and adjectives}{1}
% mdparagraph:P077:end

% mdparagraph:P078:begin
The following vocabulary gives the ordinary meanings needed in this unit. The genitive connects one noun with another; prepositions and verbs determine their own governed cases. A complete declension supplies forms, not a permission to substitute any case in any construction.
% mdparagraph:P078:end

% teaching: NEW-U13-VOC-IDELNES
% teaching: NEW-U13-NOM-IDELNES
% teaching: NEW-U13-VOC-MODOR
% teaching: NEW-U13-NOM-MODOR
% teaching: NEW-U13-VOC-SYNN
% teaching: NEW-U13-NOM-SYNN
% teaching: NEW-U13-VOC-DOHTOR
% teaching: NEW-U13-NOM-DOHTOR
% teaching: NEW-U13-VOC-BISCEOP
% teaching: NEW-U13-NOM-BISCEOP
% teaching: NEW-U13-VOC-MYNSTER
% teaching: NEW-U13-NOM-MYNSTER
% teaching: NEW-U13-VOC-LEORNUNG
% teaching: NEW-U13-NOM-LEORNUNG
% teaching: NEW-U13-VOC-CIRICE
% teaching: NEW-U13-NOM-CIRICE
% teaching: NEW-U13-VOC-UNRIHTWISNES
% teaching: NEW-U13-NOM-UNRIHTWISNES
% teaching: NEW-U13-VOC-SANG
% teaching: NEW-U13-NOM-SANG
% teaching: NEW-U13-VOC-LICHAMA
% teaching: NEW-U13-NOM-LICHAMA
% teaching: NEW-U13-VOC-EAGTHYREL
% teaching: NEW-U13-NOM-EAGTHYREL
% teaching: NEW-U13-VOC-DAEL
% teaching: NEW-U13-NOM-DAEL
% teaching: NEW-U13-VOC-ROD
% teaching: NEW-U13-NOM-ROD
% teaching: NEW-U13-VOC-AEX
% teaching: NEW-U13-NOM-AEX
% teaching: NEW-U13-VOC-ADESA
% teaching: NEW-U13-NOM-ADESA
% teaching: NEW-U13-VOC-WYRTTRUMA
% teaching: NEW-U13-NOM-WYRTTRUMA
% teaching: NEW-U13-VOC-UNTRYMNES
% teaching: NEW-U13-NOM-UNTRYMNES
% teaching: NEW-U13-VOC-GEFERSCIPE
% teaching: NEW-U13-NOM-GEFERSCIPE
% teaching: NEW-U13-VOC-GELEAFA
% teaching: NEW-U13-NOM-GELEAFA
% teaching: NEW-U13-VOC-WUNDOR
% teaching: NEW-U13-NOM-WUNDOR
% teaching: NEW-U13-VOC-HELL
% teaching: NEW-U13-NOM-HELL
% teaching: NEW-U13-VOC-HROF
% teaching: NEW-U13-NOM-HROF
% teaching: NEW-U13-VOC-MED
% teaching: NEW-U13-NOM-MED
% teaching: NEW-U13-VOC-AECER
% teaching: NEW-U13-NOM-AECER
% teaching: NEW-U13-VOC-OFEN
% teaching: NEW-U13-NOM-OFEN
% teaching: NEW-U13-VOC-BROTHOR
% teaching: NEW-U13-NOM-BROTHOR
% teaching: NEW-U13-VOC-FORTHFOR
% teaching: NEW-U13-NOM-FORTHFOR
% teaching: NEW-U13-VOC-WEG
% teaching: NEW-U13-NOM-WEG
% teaching: NEW-U13-VOC-FREOND
% teaching: NEW-U13-NOM-FREOND
% teaching: NEW-U13-VOC-BURG
% teaching: NEW-U13-NOM-BURG
% teaching: NEW-U13-VOC-WORD
% teaching: NEW-U13-NOM-WORD
% mdtable:T01:begin
\par\begingroup
\fontsize{10.5}{13.2}\selectfont
\setlength{\LTleft}{0pt}\setlength{\LTright}{\fill}
\begin{longtable}{@{}>{\raggedright\arraybackslash}p{38.39mm}>{\raggedright\arraybackslash}p{31.99mm}>{\raggedright\arraybackslash}p{47.99mm}@{}}
\toprule
\textbf{Noun} & \textbf{Gender} & \textbf{Meaning} \\
\midrule
\endfirsthead
\toprule
\textbf{Noun} & \textbf{Gender} & \textbf{Meaning} \\
\midrule
\endhead
\midrule
\multicolumn{3}{r}{\footnotesize Continued on the following page.}\\
\endfoot
\bottomrule
\endlastfoot
īdelnes & feminine & idleness \\
mōdor & feminine & mother \\
synn & feminine & sin \\
dohtor & feminine & daughter \\
bisceop & masculine & bishop \\
mynster & neuter & monastery \\
leornung & feminine & study; learning \\
cirice & feminine & church \\
unrihtwīsnes & feminine & unrighteousness \\
sang & masculine & song \\
līchama & masculine & body \\
ēagðȳrel & neuter & window \\
dǣl & masculine & part \\
rōd & feminine & cross \\
æx & feminine & ax \\
adesa & masculine & hatchet \\
wyrttruma & masculine & root \\
untrymnes & feminine & infirmity \\
gefērscipe & masculine & companionship \\
gelēafa & masculine & faith \\
wundor & neuter & wonder \\
hęll & feminine & hell \\
hrōf & masculine & roof \\
mēd & feminine & reward \\
æcer & masculine & field \\
ofen & masculine & oven \\
brōðor & masculine & brother \\
forðfōr & feminine & departure \\
weg & masculine & way; path \\
frēond & masculine & friend \\
burg & feminine & town; city \\
word & neuter & word \\
\end{longtable}
\par\endgroup

% mdtable:T01:end
% mdparagraph:P079:begin
The plural forms of īdelnes, leornung, unrihtwīsnes, untrymnes, gefērscipe, gelēafa and hęll are reference forms only. Their assigned abstract senses are singular. All other noun systems below give productive N, G, D and A forms in both numbers, within the meanings just listed. The source spelling ciricean and its sense “oratory” are supplied separately with the reading.
% mdparagraph:P079:end

\reserve{7\baselineskip}
% mdparagraph:P080:begin
\subsection*{6.1. Complete noun models}\BookHeading{nav-ch13-learner-0007}{subsubsection}{6.1. Complete noun models}{0}
% mdparagraph:P080:end

% teaching: NEW-U13-GRAM-U13-NOM-MASC
% teaching: NEW-U13-GRAM-U13-NOM-MASC-E
% teaching: NEW-U13-GRAM-U13-NOM-NEUTER
% teaching: NEW-U13-GRAM-U13-NOM-SYNCOPE
% teaching: NEW-U13-GRAM-U13-NOM-FEMININE
% teaching: NEW-U13-GRAM-U13-NOM-NES-UNG
% teaching: NEW-U13-GRAM-U13-NOM-WEAK
% teaching: NEW-U13-GRAM-U13-NOM-KINSHIP
% teaching: NEW-U13-GRAM-U13-NOM-BURG
% teaching: NEW-U13-GRAM-U13-NOM-FREOND
% mdparagraph:P081:begin
Strong masculines of the bisceop type have singular N/\allowbreak{}A with no ending, G -es and D -e; plural N/\allowbreak{}A -as, G -a and D -um. Retain the vowels in bisceop, æcer and ofen. In gefērscipe, remove citation -e before adding another ending; N/\allowbreak{}A and D singular retain -e.
% mdparagraph:P081:end

% mdparagraph:P082:begin
Strong neuters have G singular -es, D singular -e, G plural -a and D plural -um. Their N/\allowbreak{}A plural is lexical: word and wundor have no added ending; ēagðȳrel and mynster select -u with the shortened stem. Before vowel endings, mynster loses its second e, ēagðȳrel its second e, and wundor its o: mynstr-, ēagðȳrl-, wundr-. Keep the full citation form wherever the tables show it.
% mdparagraph:P082:end

% mdparagraph:P083:begin
Long feminine nouns have singular N with no added ending and G/\allowbreak{}D/\allowbreak{}A -e; selected plural N/\allowbreak{}A/\allowbreak{}G -a and D -um. Plural N/\allowbreak{}A -e is also found. Æx retains æ throughout: no masculine æ/\allowbreak{}a alternation applies. Nouns in -nes double s before an ending. In leornung, the selected oblique singular ending is -e; -a is a documented alternative.
% mdparagraph:P083:end

% mdparagraph:P084:begin
Weak masculine nouns have N singular -a, all other singular cases and N/\allowbreak{}A plural -an, G plural -ena and D plural -um. Weak feminine cirice has N singular -e and follows the same selected endings otherwise. Its productive stem is ciric-.
% mdparagraph:P084:end

% mdparagraph:P085:begin
Kinship nouns require their full tables. Mōdor, dohtor and brōðor have unchanged N/\allowbreak{}G/\allowbreak{}A singular, but D singular mēder, dehter and brēðer. Dohtor has short o. Burg has byrig in selected G/\allowbreak{}D singular and N/\allowbreak{}A plural. Frēond follows the earlier fēond pattern, with frēond- and frīend-; its selected G plural is frēonda.
% mdparagraph:P085:end

\reserve{7\baselineskip}
% mdparagraph:P086:begin
\subsection*{Table 1. īdelnes: feminine noun}\BookHeading{nav-ch13-learner-0008}{subsubsection}{Table 1. īdelnes: feminine noun}{0}
% mdparagraph:P086:end

% table: NEW-U13-NOM-IDELNES
% mdtable:T02:begin
\par\begingroup
\fontsize{10.5}{13.2}\selectfont
\setlength{\LTleft}{0pt}\setlength{\LTright}{\fill}
\noindent
\begin{tabular}{@{}>{\raggedright\arraybackslash}p{39.46mm}>{\raggedright\arraybackslash}p{39.46mm}>{\raggedright\arraybackslash}p{39.46mm}@{}}
\toprule
\textbf{Case} & \textbf{Singular} & \textbf{Plural} \\
\midrule
Nominative & īdelnes & īdelnessa \\
Genitive & īdelnesse & īdelnessa \\
Dative & īdelnesse & īdelnessum \\
Accusative & īdelnesse & īdelnessa \\
\bottomrule
\end{tabular}
\par\endgroup

% mdtable:T02:end
% mdparagraph:P087:begin
Plural: reference only.
% mdparagraph:P087:end

\reserve{7\baselineskip}
% mdparagraph:P088:begin
\subsection*{Table 2. mōdor: feminine noun}\BookHeading{nav-ch13-learner-0009}{subsubsection}{Table 2. mōdor: feminine noun}{0}\BookLabel{bk-ch13-learner-19b96c59b02d}
% mdparagraph:P088:end

% table: NEW-U13-NOM-MODOR
% mdtable:T03:begin
\par\begingroup
\fontsize{10.5}{13.2}\selectfont
\setlength{\LTleft}{0pt}\setlength{\LTright}{\fill}
\noindent
\begin{tabular}{@{}>{\raggedright\arraybackslash}p{39.46mm}>{\raggedright\arraybackslash}p{39.46mm}>{\raggedright\arraybackslash}p{39.46mm}@{}}
\toprule
\textbf{Case} & \textbf{Singular} & \textbf{Plural} \\
\midrule
Nominative & mōdor & mōdru \\
Genitive & mōdor & mōdra \\
Dative & mēder & mōdrum \\
Accusative & mōdor & mōdru \\
\bottomrule
\end{tabular}
\par\endgroup

% mdtable:T03:end
\reserve{7\baselineskip}
% mdparagraph:P089:begin
\subsection*{Table 3. dohtor: feminine noun}\BookHeading{nav-ch13-learner-0010}{subsubsection}{Table 3. dohtor: feminine noun}{0}\BookLabel{bk-ch13-learner-6a05793ee7f4}
% mdparagraph:P089:end

% table: NEW-U13-NOM-DOHTOR
% mdtable:T04:begin
\par\begingroup
\fontsize{10.5}{13.2}\selectfont
\setlength{\LTleft}{0pt}\setlength{\LTright}{\fill}
\noindent
\begin{tabular}{@{}>{\raggedright\arraybackslash}p{39.46mm}>{\raggedright\arraybackslash}p{39.46mm}>{\raggedright\arraybackslash}p{39.46mm}@{}}
\toprule
\textbf{Case} & \textbf{Singular} & \textbf{Plural} \\
\midrule
Nominative & dohtor & dohtor \\
Genitive & dohtor & dohtra \\
Dative & dehter & dohtrum \\
Accusative & dohtor & dohtor \\
\bottomrule
\end{tabular}
\par\endgroup

% mdtable:T04:end
\reserve{7\baselineskip}
% mdparagraph:P090:begin
\subsection*{Table 4. bisceop: masculine noun}\BookHeading{nav-ch13-learner-0011}{subsubsection}{Table 4. bisceop: masculine noun}{0}
% mdparagraph:P090:end

% table: NEW-U13-NOM-BISCEOP
% mdtable:T05:begin
\par\begingroup
\fontsize{10.5}{13.2}\selectfont
\setlength{\LTleft}{0pt}\setlength{\LTright}{\fill}
\noindent
\begin{tabular}{@{}>{\raggedright\arraybackslash}p{39.46mm}>{\raggedright\arraybackslash}p{39.46mm}>{\raggedright\arraybackslash}p{39.46mm}@{}}
\toprule
\textbf{Case} & \textbf{Singular} & \textbf{Plural} \\
\midrule
Nominative & bisceop & bisceopas \\
Genitive & bisceopes & bisceopa \\
Dative & bisceope & bisceopum \\
Accusative & bisceop & bisceopas \\
\bottomrule
\end{tabular}
\par\endgroup

% mdtable:T05:end
\reserve{7\baselineskip}
% mdparagraph:P091:begin
\subsection*{Table 5. leornung: feminine noun}\BookHeading{nav-ch13-learner-0012}{subsubsection}{Table 5. leornung: feminine noun}{0}
% mdparagraph:P091:end

% table: NEW-U13-NOM-LEORNUNG
% mdtable:T06:begin
\par\begingroup
\fontsize{10.5}{13.2}\selectfont
\setlength{\LTleft}{0pt}\setlength{\LTright}{\fill}
\noindent
\begin{tabular}{@{}>{\raggedright\arraybackslash}p{39.46mm}>{\raggedright\arraybackslash}p{39.46mm}>{\raggedright\arraybackslash}p{39.46mm}@{}}
\toprule
\textbf{Case} & \textbf{Singular} & \textbf{Plural} \\
\midrule
Nominative & leornung & leornunga \\
Genitive & leornunge & leornunga \\
Dative & leornunge & leornungum \\
Accusative & leornunge & leornunga \\
\bottomrule
\end{tabular}
\par\endgroup

% mdtable:T06:end
% mdparagraph:P092:begin
Plural: reference only.
% mdparagraph:P092:end

\reserve{7\baselineskip}
% mdparagraph:P093:begin
\subsection*{Table 6. cirice: feminine noun}\BookHeading{nav-ch13-learner-0013}{subsubsection}{Table 6. cirice: feminine noun}{0}
% mdparagraph:P093:end

% table: NEW-U13-NOM-CIRICE
% mdtable:T07:begin
\par\begingroup
\fontsize{10.5}{13.2}\selectfont
\setlength{\LTleft}{0pt}\setlength{\LTright}{\fill}
\noindent
\begin{tabular}{@{}>{\raggedright\arraybackslash}p{39.46mm}>{\raggedright\arraybackslash}p{39.46mm}>{\raggedright\arraybackslash}p{39.46mm}@{}}
\toprule
\textbf{Case} & \textbf{Singular} & \textbf{Plural} \\
\midrule
Nominative & cirice & cirican \\
Genitive & cirican & ciricena \\
Dative & cirican & ciricum \\
Accusative & cirican & cirican \\
\bottomrule
\end{tabular}
\par\endgroup

% mdtable:T07:end
\reserve{7\baselineskip}
% mdparagraph:P094:begin
\subsection*{Table 7. līchama: masculine noun}\BookHeading{nav-ch13-learner-0014}{subsubsection}{Table 7. līchama: masculine noun}{0}
% mdparagraph:P094:end

% table: NEW-U13-NOM-LICHAMA
% mdtable:T08:begin
\par\begingroup
\fontsize{10.5}{13.2}\selectfont
\setlength{\LTleft}{0pt}\setlength{\LTright}{\fill}
\noindent
\begin{tabular}{@{}>{\raggedright\arraybackslash}p{39.46mm}>{\raggedright\arraybackslash}p{39.46mm}>{\raggedright\arraybackslash}p{39.46mm}@{}}
\toprule
\textbf{Case} & \textbf{Singular} & \textbf{Plural} \\
\midrule
Nominative & līchama & līchaman \\
Genitive & līchaman & līchamena \\
Dative & līchaman & līchamum \\
Accusative & līchaman & līchaman \\
\bottomrule
\end{tabular}
\par\endgroup

% mdtable:T08:end
\reserve{7\baselineskip}
% mdparagraph:P095:begin
\subsection*{Table 8. ēagðȳrel: neuter noun}\BookHeading{nav-ch13-learner-0015}{subsubsection}{Table 8. ēagðȳrel: neuter noun}{0}
% mdparagraph:P095:end

% table: NEW-U13-NOM-EAGTHYREL
% mdtable:T09:begin
\par\begingroup
\fontsize{10.5}{13.2}\selectfont
\setlength{\LTleft}{0pt}\setlength{\LTright}{\fill}
\noindent
\begin{tabular}{@{}>{\raggedright\arraybackslash}p{39.46mm}>{\raggedright\arraybackslash}p{39.46mm}>{\raggedright\arraybackslash}p{39.46mm}@{}}
\toprule
\textbf{Case} & \textbf{Singular} & \textbf{Plural} \\
\midrule
Nominative & ēagðȳrel & ēagðȳrlu \\
Genitive & ēagðȳrles & ēagðȳrla \\
Dative & ēagðȳrle & ēagðȳrlum \\
Accusative & ēagðȳrel & ēagðȳrlu \\
\bottomrule
\end{tabular}
\par\endgroup

% mdtable:T09:end
\reserve{7\baselineskip}
% mdparagraph:P096:begin
\subsection*{Table 9. æx: feminine noun}\BookHeading{nav-ch13-learner-0016}{subsubsection}{Table 9. æx: feminine noun}{0}
% mdparagraph:P096:end

% table: NEW-U13-NOM-AEX
% mdtable:T10:begin
\par\begingroup
\fontsize{10.5}{13.2}\selectfont
\setlength{\LTleft}{0pt}\setlength{\LTright}{\fill}
\noindent
\begin{tabular}{@{}>{\raggedright\arraybackslash}p{39.46mm}>{\raggedright\arraybackslash}p{39.46mm}>{\raggedright\arraybackslash}p{39.46mm}@{}}
\toprule
\textbf{Case} & \textbf{Singular} & \textbf{Plural} \\
\midrule
Nominative & æx & æxa \\
Genitive & æxe & æxa \\
Dative & æxe & æxum \\
Accusative & æxe & æxa \\
\bottomrule
\end{tabular}
\par\endgroup

% mdtable:T10:end
\reserve{7\baselineskip}
% mdparagraph:P097:begin
\subsection*{Table 10. gefērscipe: masculine noun}\BookHeading{nav-ch13-learner-0017}{subsubsection}{Table 10. gefērscipe: masculine noun}{0}
% mdparagraph:P097:end

% table: NEW-U13-NOM-GEFERSCIPE
% mdtable:T11:begin
\par\begingroup
\fontsize{10.5}{13.2}\selectfont
\setlength{\LTleft}{0pt}\setlength{\LTright}{\fill}
\noindent
\begin{tabular}{@{}>{\raggedright\arraybackslash}p{36.99mm}>{\raggedright\arraybackslash}p{40.69mm}>{\raggedright\arraybackslash}p{40.69mm}@{}}
\toprule
\textbf{Case} & \textbf{Singular} & \textbf{Plural} \\
\midrule
Nominative & gefērscipe & gefērscipas \\
Genitive & gefērscipes & gefērscipa \\
Dative & gefērscipe & gefērscipum \\
Accusative & gefērscipe & gefērscipas \\
\bottomrule
\end{tabular}
\par\endgroup

% mdtable:T11:end
% mdparagraph:P098:begin
Plural: reference only.
% mdparagraph:P098:end

\reserve{7\baselineskip}
% mdparagraph:P099:begin
\subsection*{Table 11. wundor: neuter noun}\BookHeading{nav-ch13-learner-0018}{subsubsection}{Table 11. wundor: neuter noun}{0}
% mdparagraph:P099:end

% table: NEW-U13-NOM-WUNDOR
% mdtable:T12:begin
\par\begingroup
\fontsize{10.5}{13.2}\selectfont
\setlength{\LTleft}{0pt}\setlength{\LTright}{\fill}
\noindent
\begin{tabular}{@{}>{\raggedright\arraybackslash}p{39.46mm}>{\raggedright\arraybackslash}p{39.46mm}>{\raggedright\arraybackslash}p{39.46mm}@{}}
\toprule
\textbf{Case} & \textbf{Singular} & \textbf{Plural} \\
\midrule
Nominative & wundor & wundor \\
Genitive & wundres & wundra \\
Dative & wundre & wundrum \\
Accusative & wundor & wundor \\
\bottomrule
\end{tabular}
\par\endgroup

% mdtable:T12:end
\reserve{7\baselineskip}
% mdparagraph:P100:begin
\subsection*{Table 12. brōðor: masculine noun}\BookHeading{nav-ch13-learner-0019}{subsubsection}{Table 12. brōðor: masculine noun}{0}\BookLabel{bk-ch13-learner-912457bc89a0}
% mdparagraph:P100:end

% table: NEW-U13-NOM-BROTHOR
% mdtable:T13:begin
\par\begingroup
\fontsize{10.5}{13.2}\selectfont
\setlength{\LTleft}{0pt}\setlength{\LTright}{\fill}
\noindent
\begin{tabular}{@{}>{\raggedright\arraybackslash}p{39.46mm}>{\raggedright\arraybackslash}p{39.46mm}>{\raggedright\arraybackslash}p{39.46mm}@{}}
\toprule
\textbf{Case} & \textbf{Singular} & \textbf{Plural} \\
\midrule
Nominative & brōðor & brōðor \\
Genitive & brōðor & brōðra \\
Dative & brēðer & brōðrum \\
Accusative & brōðor & brōðor \\
\bottomrule
\end{tabular}
\par\endgroup

% mdtable:T13:end
\reserve{7\baselineskip}
% mdparagraph:P101:begin
\subsection*{Table 13. frēond: masculine noun}\BookHeading{nav-ch13-learner-0020}{subsubsection}{Table 13. frēond: masculine noun}{0}\BookLabel{bk-ch13-learner-890b82e6326f}
% mdparagraph:P101:end

% table: NEW-U13-NOM-FREOND
% mdtable:T14:begin
\par\begingroup
\fontsize{10.5}{13.2}\selectfont
\setlength{\LTleft}{0pt}\setlength{\LTright}{\fill}
\noindent
\begin{tabular}{@{}>{\raggedright\arraybackslash}p{39.46mm}>{\raggedright\arraybackslash}p{39.46mm}>{\raggedright\arraybackslash}p{39.46mm}@{}}
\toprule
\textbf{Case} & \textbf{Singular} & \textbf{Plural} \\
\midrule
Nominative & frēond & frīend \\
Genitive & frēondes & frēonda \\
Dative & frēonde & frēondum \\
Accusative & frēond & frīend \\
\bottomrule
\end{tabular}
\par\endgroup

% mdtable:T14:end
\reserve{7\baselineskip}
% mdparagraph:P102:begin
\subsection*{Table 14. burg: feminine noun}\BookHeading{nav-ch13-learner-0021}{subsubsection}{Table 14. burg: feminine noun}{0}\BookLabel{bk-ch13-learner-317b0ddeb02e}
% mdparagraph:P102:end

% table: NEW-U13-NOM-BURG
% mdtable:T15:begin
\par\begingroup
\fontsize{10.5}{13.2}\selectfont
\setlength{\LTleft}{0pt}\setlength{\LTright}{\fill}
\noindent
\begin{tabular}{@{}>{\raggedright\arraybackslash}p{39.46mm}>{\raggedright\arraybackslash}p{39.46mm}>{\raggedright\arraybackslash}p{39.46mm}@{}}
\toprule
\textbf{Case} & \textbf{Singular} & \textbf{Plural} \\
\midrule
Nominative & burg & byrig \\
Genitive & byrig & burga \\
Dative & byrig & burgum \\
Accusative & burg & byrig \\
\bottomrule
\end{tabular}
\par\endgroup

% mdtable:T15:end
\reserve{7\baselineskip}
% mdparagraph:P103:begin
\subsection*{Table 15. word: neuter noun}\BookHeading{nav-ch13-learner-0022}{subsubsection}{Table 15. word: neuter noun}{0}
% mdparagraph:P103:end

% table: NEW-U13-NOM-WORD
% mdtable:T16:begin
\par\begingroup
\fontsize{10.5}{13.2}\selectfont
\setlength{\LTleft}{0pt}\setlength{\LTright}{\fill}
\noindent
\begin{tabular}{@{}>{\raggedright\arraybackslash}p{39.46mm}>{\raggedright\arraybackslash}p{39.46mm}>{\raggedright\arraybackslash}p{39.46mm}@{}}
\toprule
\textbf{Case} & \textbf{Singular} & \textbf{Plural} \\
\midrule
Nominative & word & word \\
Genitive & wordes & worda \\
Dative & worde & wordum \\
Accusative & word & word \\
\bottomrule
\end{tabular}
\par\endgroup

% mdtable:T16:end
\reserve{7\baselineskip}
% mdparagraph:P104:begin
\subsection*{6.2. Diagnostic stems and shared models}\BookHeading{nav-ch13-learner-0023}{subsubsection}{6.2. Diagnostic stems and shared models}{0}
% mdparagraph:P104:end

% mdparagraph:P105:begin
The following rows give every noun’s full diagnostic stem and its G singular, D singular and N plural. For a shared model, use the indicated table in the same case and number, replace its initial stem by the stated stem, and retain the complete ending. A zero ending remains zero. If the recipe specifies the citation form for uninflected N/\allowbreak{}A cells, use the full headword there. This is a complete route through both numbers, not a guess from one oblique form.
% mdparagraph:P105:end

% mdtable:T17:begin
% table-reflow: nominal-diagnostic-panels
\par\begingroup
\fontsize{10.5}{13.2}\selectfont
\setlength{\LTleft}{0pt}\setlength{\LTright}{\fill}
\begin{longtable}{@{}>{\raggedright\arraybackslash}p{48mm}>{\raggedright\arraybackslash}p{48mm}>{\raggedright\arraybackslash}p{22.37mm}@{}}
\toprule
\textbf{Noun} & \textbf{Stem} & \textbf{Model} \\
\midrule
\endfirsthead
\toprule
\textbf{Noun} & \textbf{Stem} & \textbf{Model} \\
\midrule
\endhead
\midrule
\multicolumn{3}{r}{\footnotesize Continued on the following page.}\\
\endfoot
\bottomrule
\endlastfoot
īdelnes & īdelness & Table 1 \\
mōdor & mōdr & Table 2 \\
synn & synn & Table 9 \\
dohtor & dohtr & Table 3 \\
bisceop & bisceop & Table 4 \\
mynster & mynstr & Table 8 \\
leornung & leornung & Table 5 \\
cirice & ciric & Table 6 \\
unrihtwīsnes & unrihtwīsness & Table 1 \\
sang & sang & Table 4 \\
līchama & līcham & Table 7 \\
ēagðȳrel & ēagðȳrl & Table 8 \\
dǣl & dǣl & Table 4 \\
rōd & rōd & Table 9 \\
æx & æx & Table 9 \\
adesa & ades & Table 7 \\
wyrttruma & wyrttrum & Table 7 \\
untrymnes & untrymness & Table 1 \\
gefērscipe & gefērscip & Table 10 \\
gelēafa & gelēaf & Table 7 \\
wundor & wundr & Table 11 \\
hęll & hęll & Table 9 \\
hrōf & hrōf & Table 4 \\
mēd & mēd & Table 9 \\
æcer & æcer & Table 4 \\
ofen & ofen & Table 4 \\
brōðor & brōðr & Table 12 \\
forðfōr & forðfōr & Table 9 \\
weg & weg & Table 4 \\
frēond & frēond /\allowbreak{} frīend & Table 13 \\
burg & burg /\allowbreak{} byrig & Table 14 \\
word & word & Table 15 \\
\end{longtable}
\par\addvspace{5pt}
\setlength{\LTleft}{0pt}\setlength{\LTright}{\fill}
\begin{longtable}{@{}>{\raggedright\arraybackslash}p{31mm}>{\raggedright\arraybackslash}p{29mm}>{\raggedright\arraybackslash}p{29mm}>{\raggedright\arraybackslash}p{26.56mm}@{}}
\toprule
\textbf{Noun} & \textbf{G singular} & \textbf{D singular} & \textbf{N plural} \\
\midrule
\endfirsthead
\toprule
\textbf{Noun} & \textbf{G singular} & \textbf{D singular} & \textbf{N plural} \\
\midrule
\endhead
\midrule
\multicolumn{4}{r}{\footnotesize Continued on the following page.}\\
\endfoot
\bottomrule
\endlastfoot
īdelnes & īdelnesse & īdelnesse & īdelnessa \\
mōdor & mōdor & mēder & mōdru \\
synn & synne & synne & synna \\
dohtor & dohtor & dehter & dohtor \\
bisceop & bisceopes & bisceope & bisceopas \\
mynster & mynstres & mynstre & mynstru \\
leornung & leornunge & leornunge & leornunga \\
cirice & cirican & cirican & cirican \\
unrihtwīsnes & unrihtwīsnesse & unrihtwīsnesse & unrihtwīsnessa \\
sang & sanges & sange & sangas \\
līchama & līchaman & līchaman & līchaman \\
ēagðȳrel & ēagðȳrles & ēagðȳrle & ēagðȳrlu \\
dǣl & dǣles & dǣle & dǣlas \\
rōd & rōde & rōde & rōda \\
æx & æxe & æxe & æxa \\
adesa & adesan & adesan & adesan \\
wyrttruma & wyrttruman & wyrttruman & wyrttruman \\
untrymnes & untrymnesse & untrymnesse & untrymnessa \\
gefērscipe & gefērscipes & gefērscipe & gefērscipas \\
gelēafa & gelēafan & gelēafan & gelēafan \\
wundor & wundres & wundre & wundor \\
hęll & hęlle & hęlle & hęlla \\
hrōf & hrōfes & hrōfe & hrōfas \\
mēd & mēde & mēde & mēda \\
æcer & æceres & æcere & æceras \\
ofen & ofenes & ofene & ofenas \\
brōðor & brōðor & brēðer & brōðor \\
forðfōr & forðfōre & forðfōre & forðfōra \\
weg & weges & wege & wegas \\
frēond & frēondes & frēonde & frīend \\
burg & byrig & byrig & byrig \\
word & wordes & worde & word \\
\end{longtable}
\par\addvspace{5pt}
\par\endgroup
% mdtable:T17:end

% mdtable:T18:begin
\par\begingroup
\fontsize{10.5}{13.2}\selectfont
\setlength{\LTleft}{0pt}\setlength{\LTright}{\fill}
\begin{longtable}{@{}>{\raggedright\arraybackslash}p{25mm}>{\raggedright\arraybackslash}p{19mm}>{\raggedright\arraybackslash}p{74.37mm}@{}}
\toprule
\textbf{Noun} & \textbf{Complete model} & \textbf{Operation in the corresponding cell} \\
\midrule
\endfirsthead
\toprule
\textbf{Noun} & \textbf{Complete model} & \textbf{Operation in the corresponding cell} \\
\midrule
\endhead
\midrule
\multicolumn{3}{r}{\footnotesize Continued on the following page.}\\
\endfoot
\bottomrule
\endlastfoot
synn & Table 9 & Replace æx- by synn-; retain the ending. \\
mynster & Table 8 & In N singular, A singular, use the full citation mynster. Replace ēagðȳrl- by mynstr-; retain the ending. \\
unrihtwīsnes & Table 1 & In N singular, use the full citation unrihtwīsnes. Replace īdelness- by unrihtwīsness-; retain the ending. \\
sang & Table 4 & Replace bisceop- by sang-; retain the ending. \\
dǣl & Table 4 & Replace bisceop- by dǣl-; retain the ending. \\
rōd & Table 9 & Replace æx- by rōd-; retain the ending. \\
adesa & Table 7 & Replace līcham- by ades-; retain the ending. \\
wyrttruma & Table 7 & Replace līcham- by wyrttrum-; retain the ending. \\
untrymnes & Table 1 & In N singular, use the full citation untrymnes. Replace īdelness- by untrymness-; retain the ending. \\
gelēafa & Table 7 & Replace līcham- by gelēaf-; retain the ending. \\
hęll & Table 9 & Replace æx- by hęll-; retain the ending. \\
hrōf & Table 4 & Replace bisceop- by hrōf-; retain the ending. \\
mēd & Table 9 & Replace æx- by mēd-; retain the ending. \\
æcer & Table 4 & Replace bisceop- by æcer-; retain the ending. \\
ofen & Table 4 & Replace bisceop- by ofen-; retain the ending. \\
forðfōr & Table 9 & Replace æx- by forðfōr-; retain the ending. \\
weg & Table 4 & Replace bisceop- by weg-; retain the ending. \\
\end{longtable}
\par\endgroup

% mdtable:T18:end
% mdparagraph:P106:begin
Recognize these documented alternatives without replacing the selected course forms: mōdra for N/\allowbreak{}A plural mōdru; dohtru or dohtra for N/\allowbreak{}A plural dohtor; brōðru for N/\allowbreak{}A plural brōðor; late G singular mēder and dehter; G singular burge beside byrig; D singular frīend and N/\allowbreak{}A plural frēond or frēondas beside the selected frēonde and frīend. Cook’s citation syn(n) includes syn beside selected synn. These alternatives do not change the listed meanings or case relations.
% mdparagraph:P106:end

\reserve{7\baselineskip}
% mdparagraph:P107:begin
\subsection*{6.3. Adjectives and comparison}\BookHeading{nav-ch13-learner-0024}{subsubsection}{6.3. Adjectives and comparison}{0}
% mdparagraph:P107:end

% teaching: NEW-U13-GRAM-U13-ADJ-STRONG
% teaching: NEW-U13-GRAM-U13-ADJ-WEAK
% teaching: NEW-U13-GRAM-U13-ADJ-HALIG
% teaching: NEW-U13-GRAM-U13-ADJ-COMPARISON
% teaching: NEW-U13-VOC-HALIG
% teaching: NEW-U13-ADJ-HALIG
% teaching: NEW-U13-VOC-SWETE
% teaching: NEW-U13-ADJ-SWETE
% teaching: NEW-U13-VOC-LICHAMLIC
% teaching: NEW-U13-ADJ-LICHAMLIC
% teaching: NEW-U13-VOC-HEFIG
% teaching: NEW-U13-ADJ-HEFIG
% teaching: NEW-U13-VOC-WYNSUM
% teaching: NEW-U13-ADJ-WYNSUM
% teaching: NEW-U13-VOC-HWIT
% teaching: NEW-U13-ADJ-HWIT
% teaching: NEW-U13-VOC-WIERTHE
% teaching: NEW-U13-ADJ-WIERTHE
% teaching: NEW-U13-VOC-HEOFONLIC
% teaching: NEW-U13-ADJ-HEOFONLIC
% teaching: NEW-U13-VOC-SWETE-COMP
% teaching: NEW-U13-ADJ-SWETE-COMP
% teaching: NEW-U13-VOC-SWETE-SUPERL
% teaching: NEW-U13-ADJ-SWETE-SUPERL
% mdtable:T19:begin
\par\begingroup
\fontsize{10.5}{13.2}\selectfont
\setlength{\LTleft}{0pt}\setlength{\LTright}{\fill}
\noindent
\begin{tabular}{@{}>{\raggedright\arraybackslash}p{53.62mm}>{\raggedright\arraybackslash}p{67.56mm}@{}}
\toprule
\textbf{Adjective} & \textbf{Meaning} \\
\midrule
hālig & holy \\
swēte & sweet \\
līchamlic & bodily \\
hęfig & grievous \\
wynsum & pleasant; winsome \\
hwīt & white \\
wierðe & worthy \\
heofonlic & heavenly \\
swētra & sweeter \\
swētest & sweetest \\
\bottomrule
\end{tabular}
\par\endgroup

% mdtable:T19:end
% mdparagraph:P108:begin
Use strong and weak adjective forms in the environments already taught. The following complete tables cover N, G, D and A in all three genders and both numbers. Masculine and neuter singular instrumental forms are printed for reference; their display does not add independent productive instrumental syntax. Comparatives have weak inflection only. Superlatives have both strong and weak forms.
% mdparagraph:P108:end

% mdparagraph:P109:begin
An adjective in -e loses that citation ending before another ending: swēte gives swēt-, wierðe gives wierð-. Their strong feminine N singular and neuter N/\allowbreak{}A plural have -u. Long monosyllables and swētest omit that -u. Līchamlic, hęfig, wynsum and heofonlic retain it in those cells.
% mdparagraph:P109:end

% mdparagraph:P110:begin
Hālig needs its own full pattern: feminine N singular and neuter N/\allowbreak{}A plural are hālig. Before vowel endings use hālg-; before -ne, -re and -ra retain hālig-. Hęfig retains its medial vowel throughout the selected pattern. The weak G plural ending is selected -ra; -ena is also documented.
% mdparagraph:P110:end

% mdparagraph:P111:begin
Swēte forms comparative swētra and superlative swētest. In the comparative G plural, swētr- plus -ra simplifies the doubled r after t, giving swētra. Do not mechanically write swētrra. The shared model for hwīt below is morphological only: “white” does not acquire the meaning “sweetest.”
% mdparagraph:P111:end

\reserve{7\baselineskip}
% mdparagraph:P112:begin
\subsection*{Table 16. hālig: strong adjective}\BookHeading{nav-ch13-learner-0025}{subsubsection}{Table 16. hālig: strong adjective}{0}
% mdparagraph:P112:end

% table: NEW-U13-ADJ-HALIG-STRONG
% mdtable:T20:begin
% course-table-style: 2026-10
% course-panel:T20-P01:begin
\par\begingroup
\fontsize{10.5}{13.2}\selectfont
\setlength{\LTleft}{0pt}\setlength{\LTright}{\fill}
\noindent
\begin{tabular}{@{}>{\raggedright\arraybackslash}p{22.0mm}>{\raggedright\arraybackslash}p{31.188mm}>{\raggedright\arraybackslash}p{31.188mm}>{\raggedright\arraybackslash}p{31.188mm}@{}}
\toprule
\textbf{Case} & \textbf{Masculine singular} & \textbf{Neuter singular} & \textbf{Feminine singular} \\
\midrule
Nominative & hālig & hālig & hālig \\
Genitive & hālges & hālges & hāligre \\
Dative & hālgum & hālgum & hāligre \\
Accusative & hāligne & hālig & hālge \\
Instrumental & hālge & hālge & as dative \\
\bottomrule
\end{tabular}
\par\endgroup
% course-panel:T20-P01:end
\par\addvspace{5pt}
% course-panel:T20-P02:begin
\par\begingroup
\fontsize{10.5}{13.2}\selectfont
\setlength{\LTleft}{0pt}\setlength{\LTright}{\fill}
\noindent
\begin{tabular}{@{}>{\raggedright\arraybackslash}p{22.0mm}>{\raggedright\arraybackslash}p{31.188mm}>{\raggedright\arraybackslash}p{31.188mm}>{\raggedright\arraybackslash}p{31.188mm}@{}}
\toprule
\textbf{Case} & \textbf{Masculine plural} & \textbf{Neuter plural} & \textbf{Feminine plural} \\
\midrule
Nominative & hālge & hālig & hālga \\
Genitive & hāligra & hāligra & hāligra \\
Dative & hālgum & hālgum & hālgum \\
Accusative & hālge & hālig & hālga \\
Instrumental & as dative & as dative & as dative \\
\bottomrule
\end{tabular}
\par\endgroup
% course-panel:T20-P02:end
% mdtable:T20:end

\reserve{7\baselineskip}
% mdparagraph:P113:begin
\subsection*{Table 17. hālig: weak adjective}\BookHeading{nav-ch13-learner-0026}{subsubsection}{Table 17. hālig: weak adjective}{0}
% mdparagraph:P113:end

% table: NEW-U13-ADJ-HALIG-WEAK
% mdtable:T21:begin
% course-table-style: 2026-10
% course-panel:T21-P01:begin
\par\begingroup
\fontsize{10.5}{13.2}\selectfont
\setlength{\LTleft}{0pt}\setlength{\LTright}{\fill}
\noindent
\begin{tabular}{@{}>{\raggedright\arraybackslash}p{22.0mm}>{\raggedright\arraybackslash}p{31.188mm}>{\raggedright\arraybackslash}p{31.188mm}>{\raggedright\arraybackslash}p{31.188mm}@{}}
\toprule
\textbf{Case} & \textbf{Masculine singular} & \textbf{Neuter singular} & \textbf{Feminine singular} \\
\midrule
Nominative & hālga & hālge & hālge \\
Genitive & hālgan & hālgan & hālgan \\
Dative & hālgan & hālgan & hālgan \\
Accusative & hālgan & hālge & hālgan \\
\bottomrule
\end{tabular}
\par\endgroup
% course-panel:T21-P01:end
\par\addvspace{5pt}
% course-panel:T21-P02:begin
\par\begingroup
\fontsize{10.5}{13.2}\selectfont
\setlength{\LTleft}{0pt}\setlength{\LTright}{\fill}
\noindent
\begin{tabular}{@{}>{\raggedright\arraybackslash}p{22.0mm}>{\raggedright\arraybackslash}p{31.188mm}>{\raggedright\arraybackslash}p{31.188mm}>{\raggedright\arraybackslash}p{31.188mm}@{}}
\toprule
\textbf{Case} & \textbf{Masculine plural} & \textbf{Neuter plural} & \textbf{Feminine plural} \\
\midrule
Nominative & hālgan & hālgan & hālgan \\
Genitive & hāligra & hāligra & hāligra \\
Dative & hālgum & hālgum & hālgum \\
Accusative & hālgan & hālgan & hālgan \\
\bottomrule
\end{tabular}
\par\endgroup
% course-panel:T21-P02:end
% mdtable:T21:end

\reserve{7\baselineskip}
% mdparagraph:P114:begin
\subsection*{Table 18. swēte: strong adjective}\BookHeading{nav-ch13-learner-0027}{subsubsection}{Table 18. swēte: strong adjective}{0}
% mdparagraph:P114:end

% table: NEW-U13-ADJ-SWETE-STRONG
% mdtable:T22:begin
% course-table-style: 2026-10
% course-panel:T22-P01:begin
\par\begingroup
\fontsize{10.5}{13.2}\selectfont
\setlength{\LTleft}{0pt}\setlength{\LTright}{\fill}
\noindent
\begin{tabular}{@{}>{\raggedright\arraybackslash}p{22.0mm}>{\raggedright\arraybackslash}p{31.188mm}>{\raggedright\arraybackslash}p{31.188mm}>{\raggedright\arraybackslash}p{31.188mm}@{}}
\toprule
\textbf{Case} & \textbf{Masculine singular} & \textbf{Neuter singular} & \textbf{Feminine singular} \\
\midrule
Nominative & swēte & swēte & swētu \\
Genitive & swētes & swētes & swētre \\
Dative & swētum & swētum & swētre \\
Accusative & swētne & swēte & swēte \\
Instrumental & swēte & swēte & as dative \\
\bottomrule
\end{tabular}
\par\endgroup
% course-panel:T22-P01:end
\par\addvspace{5pt}
% course-panel:T22-P02:begin
\par\begingroup
\fontsize{10.5}{13.2}\selectfont
\setlength{\LTleft}{0pt}\setlength{\LTright}{\fill}
\noindent
\begin{tabular}{@{}>{\raggedright\arraybackslash}p{22.0mm}>{\raggedright\arraybackslash}p{31.188mm}>{\raggedright\arraybackslash}p{31.188mm}>{\raggedright\arraybackslash}p{31.188mm}@{}}
\toprule
\textbf{Case} & \textbf{Masculine plural} & \textbf{Neuter plural} & \textbf{Feminine plural} \\
\midrule
Nominative & swēte & swētu & swēta \\
Genitive & swētra & swētra & swētra \\
Dative & swētum & swētum & swētum \\
Accusative & swēte & swētu & swēta \\
Instrumental & as dative & as dative & as dative \\
\bottomrule
\end{tabular}
\par\endgroup
% course-panel:T22-P02:end
% mdtable:T22:end

\reserve{7\baselineskip}
% mdparagraph:P115:begin
\subsection*{Table 19. swēte: weak adjective}\BookHeading{nav-ch13-learner-0028}{subsubsection}{Table 19. swēte: weak adjective}{0}
% mdparagraph:P115:end

% table: NEW-U13-ADJ-SWETE-WEAK
% mdtable:T23:begin
% course-table-style: 2026-10
% course-panel:T23-P01:begin
\par\begingroup
\fontsize{10.5}{13.2}\selectfont
\setlength{\LTleft}{0pt}\setlength{\LTright}{\fill}
\noindent
\begin{tabular}{@{}>{\raggedright\arraybackslash}p{22.0mm}>{\raggedright\arraybackslash}p{31.188mm}>{\raggedright\arraybackslash}p{31.188mm}>{\raggedright\arraybackslash}p{31.188mm}@{}}
\toprule
\textbf{Case} & \textbf{Masculine singular} & \textbf{Neuter singular} & \textbf{Feminine singular} \\
\midrule
Nominative & swēta & swēte & swēte \\
Genitive & swētan & swētan & swētan \\
Dative & swētan & swētan & swētan \\
Accusative & swētan & swēte & swētan \\
\bottomrule
\end{tabular}
\par\endgroup
% course-panel:T23-P01:end
\par\addvspace{5pt}
% course-panel:T23-P02:begin
\par\begingroup
\fontsize{10.5}{13.2}\selectfont
\setlength{\LTleft}{0pt}\setlength{\LTright}{\fill}
\noindent
\begin{tabular}{@{}>{\raggedright\arraybackslash}p{22.0mm}>{\raggedright\arraybackslash}p{31.188mm}>{\raggedright\arraybackslash}p{31.188mm}>{\raggedright\arraybackslash}p{31.188mm}@{}}
\toprule
\textbf{Case} & \textbf{Masculine plural} & \textbf{Neuter plural} & \textbf{Feminine plural} \\
\midrule
Nominative & swētan & swētan & swētan \\
Genitive & swētra & swētra & swētra \\
Dative & swētum & swētum & swētum \\
Accusative & swētan & swētan & swētan \\
\bottomrule
\end{tabular}
\par\endgroup
% course-panel:T23-P02:end
% mdtable:T23:end

\reserve{7\baselineskip}
% mdparagraph:P116:begin
\subsection*{Table 20. līchamlic: strong adjective}\BookHeading{nav-ch13-learner-0029}{subsubsection}{Table 20. līchamlic: strong adjective}{0}
% mdparagraph:P116:end

% table: NEW-U13-ADJ-LICHAMLIC-STRONG
% mdtable:T24:begin
% course-table-style: 2026-10
% course-panel:T24-P01:begin
\par\begingroup
\fontsize{10.5}{13.2}\selectfont
\setlength{\LTleft}{0pt}\setlength{\LTright}{\fill}
\noindent
\begin{tabular}{@{}>{\raggedright\arraybackslash}p{22.0mm}>{\raggedright\arraybackslash}p{31.188mm}>{\raggedright\arraybackslash}p{31.188mm}>{\raggedright\arraybackslash}p{31.188mm}@{}}
\toprule
\textbf{Case} & \textbf{Masculine singular} & \textbf{Neuter singular} & \textbf{Feminine singular} \\
\midrule
Nominative & līchamlic & līchamlic & līchamlicu \\
Genitive & līchamlices & līchamlices & līchamlicre \\
Dative & līchamlicum & līchamlicum & līchamlicre \\
Accusative & līchamlicne & līchamlic & līchamlice \\
Instrumental & līchamlice & līchamlice & as dative \\
\bottomrule
\end{tabular}
\par\endgroup
% course-panel:T24-P01:end
\par\addvspace{5pt}
% course-panel:T24-P02:begin
\par\begingroup
\fontsize{10.5}{13.2}\selectfont
\setlength{\LTleft}{0pt}\setlength{\LTright}{\fill}
\noindent
\begin{tabular}{@{}>{\raggedright\arraybackslash}p{22.0mm}>{\raggedright\arraybackslash}p{31.188mm}>{\raggedright\arraybackslash}p{31.188mm}>{\raggedright\arraybackslash}p{31.188mm}@{}}
\toprule
\textbf{Case} & \textbf{Masculine plural} & \textbf{Neuter plural} & \textbf{Feminine plural} \\
\midrule
Nominative & līchamlice & līchamlicu & līchamlica \\
Genitive & līchamlicra & līchamlicra & līchamlicra \\
Dative & līchamlicum & līchamlicum & līchamlicum \\
Accusative & līchamlice & līchamlicu & līchamlica \\
Instrumental & as dative & as dative & as dative \\
\bottomrule
\end{tabular}
\par\endgroup
% course-panel:T24-P02:end
% mdtable:T24:end

\reserve{7\baselineskip}
% mdparagraph:P117:begin
\subsection*{Table 21. līchamlic: weak adjective}\BookHeading{nav-ch13-learner-0030}{subsubsection}{Table 21. līchamlic: weak adjective}{0}
% mdparagraph:P117:end

% table: NEW-U13-ADJ-LICHAMLIC-WEAK
% mdtable:T25:begin
% course-table-style: 2026-10
% course-panel:T25-P01:begin
\par\begingroup
\fontsize{10.5}{13.2}\selectfont
\setlength{\LTleft}{0pt}\setlength{\LTright}{\fill}
\noindent
\begin{tabular}{@{}>{\raggedright\arraybackslash}p{22.0mm}>{\raggedright\arraybackslash}p{31.188mm}>{\raggedright\arraybackslash}p{31.188mm}>{\raggedright\arraybackslash}p{31.188mm}@{}}
\toprule
\textbf{Case} & \textbf{Masculine singular} & \textbf{Neuter singular} & \textbf{Feminine singular} \\
\midrule
Nominative & līchamlica & līchamlice & līchamlice \\
Genitive & līchamlican & līchamlican & līchamlican \\
Dative & līchamlican & līchamlican & līchamlican \\
Accusative & līchamlican & līchamlice & līchamlican \\
\bottomrule
\end{tabular}
\par\endgroup
% course-panel:T25-P01:end
\par\addvspace{5pt}
% course-panel:T25-P02:begin
\par\begingroup
\fontsize{10.5}{13.2}\selectfont
\setlength{\LTleft}{0pt}\setlength{\LTright}{\fill}
\noindent
\begin{tabular}{@{}>{\raggedright\arraybackslash}p{22.0mm}>{\raggedright\arraybackslash}p{31.188mm}>{\raggedright\arraybackslash}p{31.188mm}>{\raggedright\arraybackslash}p{31.188mm}@{}}
\toprule
\textbf{Case} & \textbf{Masculine plural} & \textbf{Neuter plural} & \textbf{Feminine plural} \\
\midrule
Nominative & līchamlican & līchamlican & līchamlican \\
Genitive & līchamlicra & līchamlicra & līchamlicra \\
Dative & līchamlicum & līchamlicum & līchamlicum \\
Accusative & līchamlican & līchamlican & līchamlican \\
\bottomrule
\end{tabular}
\par\endgroup
% course-panel:T25-P02:end
% mdtable:T25:end

\reserve{7\baselineskip}
% mdparagraph:P118:begin
\subsection*{Table 22. swētra: weak adjective}\BookHeading{nav-ch13-learner-0031}{subsubsection}{Table 22. swētra: weak adjective}{0}
% mdparagraph:P118:end

% table: NEW-U13-ADJ-SWETE-COMP-WEAK
% mdtable:T26:begin
% course-table-style: 2026-10
% course-panel:T26-P01:begin
\par\begingroup
\fontsize{10.5}{13.2}\selectfont
\setlength{\LTleft}{0pt}\setlength{\LTright}{\fill}
\noindent
\begin{tabular}{@{}>{\raggedright\arraybackslash}p{22.0mm}>{\raggedright\arraybackslash}p{31.188mm}>{\raggedright\arraybackslash}p{31.188mm}>{\raggedright\arraybackslash}p{31.188mm}@{}}
\toprule
\textbf{Case} & \textbf{Masculine singular} & \textbf{Neuter singular} & \textbf{Feminine singular} \\
\midrule
Nominative & swētra & swētre & swētre \\
Genitive & swētran & swētran & swētran \\
Dative & swētran & swētran & swētran \\
Accusative & swētran & swētre & swētran \\
\bottomrule
\end{tabular}
\par\endgroup
% course-panel:T26-P01:end
\par\addvspace{5pt}
% course-panel:T26-P02:begin
\par\begingroup
\fontsize{10.5}{13.2}\selectfont
\setlength{\LTleft}{0pt}\setlength{\LTright}{\fill}
\noindent
\begin{tabular}{@{}>{\raggedright\arraybackslash}p{22.0mm}>{\raggedright\arraybackslash}p{31.188mm}>{\raggedright\arraybackslash}p{31.188mm}>{\raggedright\arraybackslash}p{31.188mm}@{}}
\toprule
\textbf{Case} & \textbf{Masculine plural} & \textbf{Neuter plural} & \textbf{Feminine plural} \\
\midrule
Nominative & swētran & swētran & swētran \\
Genitive & swētra & swētra & swētra \\
Dative & swētrum & swētrum & swētrum \\
Accusative & swētran & swētran & swētran \\
\bottomrule
\end{tabular}
\par\endgroup
% course-panel:T26-P02:end
% mdtable:T26:end

\reserve{7\baselineskip}
% mdparagraph:P119:begin
\subsection*{Table 23. swētest: strong adjective}\BookHeading{nav-ch13-learner-0032}{subsubsection}{Table 23. swētest: strong adjective}{0}
% mdparagraph:P119:end

% table: NEW-U13-ADJ-SWETE-SUPERL-STRONG
% mdtable:T27:begin
% course-table-style: 2026-10
% course-panel:T27-P01:begin
\par\begingroup
\fontsize{10.5}{13.2}\selectfont
\setlength{\LTleft}{0pt}\setlength{\LTright}{\fill}
\noindent
\begin{tabular}{@{}>{\raggedright\arraybackslash}p{22.0mm}>{\raggedright\arraybackslash}p{31.188mm}>{\raggedright\arraybackslash}p{31.188mm}>{\raggedright\arraybackslash}p{31.188mm}@{}}
\toprule
\textbf{Case} & \textbf{Masculine singular} & \textbf{Neuter singular} & \textbf{Feminine singular} \\
\midrule
Nominative & swētest & swētest & swētest \\
Genitive & swētestes & swētestes & swētestre \\
Dative & swētestum & swētestum & swētestre \\
Accusative & swētestne & swētest & swēteste \\
Instrumental & swēteste & swēteste & as dative \\
\bottomrule
\end{tabular}
\par\endgroup
% course-panel:T27-P01:end
\par\addvspace{5pt}
% course-panel:T27-P02:begin
\par\begingroup
\fontsize{10.5}{13.2}\selectfont
\setlength{\LTleft}{0pt}\setlength{\LTright}{\fill}
\noindent
\begin{tabular}{@{}>{\raggedright\arraybackslash}p{22.0mm}>{\raggedright\arraybackslash}p{31.188mm}>{\raggedright\arraybackslash}p{31.188mm}>{\raggedright\arraybackslash}p{31.188mm}@{}}
\toprule
\textbf{Case} & \textbf{Masculine plural} & \textbf{Neuter plural} & \textbf{Feminine plural} \\
\midrule
Nominative & swēteste & swētest & swētesta \\
Genitive & swētestra & swētestra & swētestra \\
Dative & swētestum & swētestum & swētestum \\
Accusative & swēteste & swētest & swētesta \\
Instrumental & as dative & as dative & as dative \\
\bottomrule
\end{tabular}
\par\endgroup
% course-panel:T27-P02:end
% mdtable:T27:end

\reserve{7\baselineskip}
% mdparagraph:P120:begin
\subsection*{Table 24. swētest: weak adjective}\BookHeading{nav-ch13-learner-0033}{subsubsection}{Table 24. swētest: weak adjective}{0}
% mdparagraph:P120:end

% table: NEW-U13-ADJ-SWETE-SUPERL-WEAK
% mdtable:T28:begin
% course-table-style: 2026-10
% course-panel:T28-P01:begin
\par\begingroup
\fontsize{10.5}{13.2}\selectfont
\setlength{\LTleft}{0pt}\setlength{\LTright}{\fill}
\noindent
\begin{tabular}{@{}>{\raggedright\arraybackslash}p{22.0mm}>{\raggedright\arraybackslash}p{31.188mm}>{\raggedright\arraybackslash}p{31.188mm}>{\raggedright\arraybackslash}p{31.188mm}@{}}
\toprule
\textbf{Case} & \textbf{Masculine singular} & \textbf{Neuter singular} & \textbf{Feminine singular} \\
\midrule
Nominative & swētesta & swēteste & swēteste \\
Genitive & swētestan & swētestan & swētestan \\
Dative & swētestan & swētestan & swētestan \\
Accusative & swētestan & swēteste & swētestan \\
\bottomrule
\end{tabular}
\par\endgroup
% course-panel:T28-P01:end
\par\addvspace{5pt}
% course-panel:T28-P02:begin
\par\begingroup
\fontsize{10.5}{13.2}\selectfont
\setlength{\LTleft}{0pt}\setlength{\LTright}{\fill}
\noindent
\begin{tabular}{@{}>{\raggedright\arraybackslash}p{22.0mm}>{\raggedright\arraybackslash}p{31.188mm}>{\raggedright\arraybackslash}p{31.188mm}>{\raggedright\arraybackslash}p{31.188mm}@{}}
\toprule
\textbf{Case} & \textbf{Masculine plural} & \textbf{Neuter plural} & \textbf{Feminine plural} \\
\midrule
Nominative & swētestan & swētestan & swētestan \\
Genitive & swētestra & swētestra & swētestra \\
Dative & swētestum & swētestum & swētestum \\
Accusative & swētestan & swētestan & swētestan \\
\bottomrule
\end{tabular}
\par\endgroup
% course-panel:T28-P02:end
% mdtable:T28:end

% mdtable:T29:begin
\par\begingroup
\fontsize{10.5}{13.2}\selectfont
\setlength{\LTleft}{0pt}\setlength{\LTright}{\fill}
\noindent
\begin{tabular}{@{}>{\raggedright\arraybackslash}p{16.83mm}>{\raggedright\arraybackslash}p{16.83mm}>{\raggedright\arraybackslash}p{29.87mm}>{\raggedright\arraybackslash}p{28.19mm}>{\raggedright\arraybackslash}p{21.04mm}@{}}
\toprule
\textbf{Adjective} & \textbf{Stem} & \textbf{Strong feminine N singular} & \textbf{Weak masculine N singular} & \textbf{Complete model(s)} \\
\midrule
hālig & hālig & hālig & hālga & Table 16, Table 17 \\
swēte & swēt & swētu & swēta & Table 18, Table 19 \\
līchamlic & līchamlic & līchamlicu & līchamlica & Table 20, Table 21 \\
hęfig & hęfig & hęfigu & hęfiga & Table 20, Table 21 \\
wynsum & wynsum & wynsumu & wynsuma & Table 20, Table 21 \\
hwīt & hwīt & hwīt & hwīta & Table 23, Table 24 \\
wierðe & wierð & wierðu & wierða & Table 18, Table 19 \\
heofonlic & heofonlic & heofonlicu & heofonlica & Table 20, Table 21 \\
swētra & swētr & No strong comparative & swētra & Table 22 \\
swētest & swētest & swētest & swētesta & Table 23, Table 24 \\
\bottomrule
\end{tabular}
\par\endgroup

% mdtable:T29:end
% mdtable:T30:begin
\par\begingroup
\fontsize{10.5}{13.2}\selectfont
\setlength{\LTleft}{0pt}\setlength{\LTright}{\fill}
\noindent
\begin{tabular}{@{}>{\raggedright\arraybackslash}p{22.31mm}>{\raggedright\arraybackslash}p{27.89mm}>{\raggedright\arraybackslash}p{68.17mm}@{}}
\toprule
\textbf{Adjective} & \textbf{Complete model(s)} & \textbf{Operation in the corresponding cell} \\
\midrule
hęfig & Table 20, Table 21 & Replace līchamlic- by hęfig-; retain the ending. \\
wynsum & Table 20, Table 21 & Replace līchamlic- by wynsum-; retain the ending. \\
hwīt & Table 23, Table 24 & Replace swētest- by hwīt-; retain the ending. \\
wierðe & Table 18, Table 19 & Replace swēt- by wierð-; retain the ending. \\
heofonlic & Table 20, Table 21 & Replace līchamlic- by heofonlic-; retain the ending. \\
\bottomrule
\end{tabular}
\par\endgroup

% mdtable:T30:end
\reserve{7\baselineskip}
% mdparagraph:P121:begin
\section*{7. Verbal systems}\BookHeading{nav-ch13-learner-0034}{subsection}{7. Verbal systems}{1}
% mdparagraph:P121:end

% mdparagraph:P122:begin
The finite tables give all three persons in both numbers. Present and preterite indicative are separate from the subjunctive, imperative, infinitive and participle systems. Learn each listed system as a whole, but use it only with its taught meaning and construction. A citation infinitive used to name a verb does not by itself assign that infinitive for sentence production.
% mdparagraph:P122:end

% teaching: NEW-U13-VOC-FORDILGIAN
% teaching: NEW-U13-PAR-FORDILGIAN
% teaching: NEW-U13-VOC-GEWUNIAN
% teaching: NEW-U13-PAR-GEWUNIAN
% teaching: NEW-U13-VOC-ASCIAN
% teaching: NEW-U13-PAR-ASCIAN
% teaching: NEW-U13-VOC-HALSIAN
% teaching: NEW-U13-PAR-HALSIAN
% teaching: NEW-U13-VOC-WUNIAN
% teaching: NEW-U13-PAR-WUNIAN
% teaching: NEW-U13-VOC-HATIAN
% teaching: NEW-U13-PAR-HATIAN
% teaching: NEW-U13-VOC-ASTRECCEAN
% teaching: NEW-U13-PAR-ASTRECCEAN
% teaching: NEW-U13-VOC-ONLIESAN
% teaching: NEW-U13-PAR-ONLIESAN
% teaching: NEW-U13-VOC-NEALAECAN
% teaching: NEW-U13-PAR-NEALAECAN
% teaching: NEW-U13-VOC-FORCEORFAN
% teaching: NEW-U13-PAR-FORCEORFAN
% teaching: NEW-U13-VOC-ASETTAN
% teaching: NEW-U13-PAR-ASETTAN
% teaching: NEW-U13-VOC-OFERWINNAN
% teaching: NEW-U13-PAR-OFERWINNAN
% teaching: NEW-U13-VOC-GESCRYDAN
% teaching: NEW-U13-PAR-GESCRYDAN
% teaching: NEW-U13-VOC-ONDRAEDAN
% teaching: NEW-U13-PAR-ONDRAEDAN
% teaching: NEW-U13-VOC-FORDON
% teaching: NEW-U13-PAR-FORDON
% teaching: NEW-U13-VOC-INGAN
% teaching: NEW-U13-PAR-INGAN
% teaching: NEW-U13-VOC-ASENDAN
% teaching: NEW-U13-PAR-ASENDAN
% teaching: NEW-U13-VOC-BEHEALDAN
% teaching: NEW-U13-PAR-BEHEALDAN
% teaching: NEW-U13-VOC-BEBEODAN
% teaching: NEW-U13-PAR-BEBEODAN
% teaching: NEW-U13-VOC-EFTHWEORFAN
% teaching: NEW-U13-PAR-EFTHWEORFAN
% teaching: NEW-U13-VOC-LAEDAN
% teaching: NEW-U13-PAR-LAEDAN
% teaching: NEW-U13-VOC-SPRECAN
% teaching: NEW-U13-PAR-SPRECAN
% teaching: NEW-U13-VOC-BEGAN
% teaching: NEW-U13-VOC-ASTIGAN
% teaching: NEW-U13-VOC-SELLAN
% teaching: NEW-U13-VOC-WYRCEAN
% mdtable:T31:begin
\par\begingroup
\fontsize{10.5}{13.2}\selectfont
\setlength{\LTleft}{0pt}\setlength{\LTright}{\fill}
\begin{longtable}{@{}>{\raggedright\arraybackslash}p{44.96mm}>{\raggedright\arraybackslash}p{76.23mm}@{}}
\toprule
\textbf{Verb} & \textbf{Meaning} \\
\midrule
\endfirsthead
\toprule
\textbf{Verb} & \textbf{Meaning} \\
\midrule
\endhead
\midrule
\multicolumn{2}{r}{\footnotesize Continued on the following page.}\\
\endfoot
\bottomrule
\endlastfoot
fordilgian & destroy; annihilate \\
gewunian & be accustomed; be wont \\
āscian & ask a person about something \\
hālsian & entreat; implore \\
wunian & dwell; remain \\
hatian & hate \\
āstręccean & prostrate (oneself) \\
onlīesan & release \\
nēalǣcan & approach; draw near \\
forceorfan & cut down; cut off \\
āsęttan & set; place \\
oferwinnan & overcome; conquer \\
gescrȳdan & clothe \\
ondrǣdan & fear \\
fordōn & destroy \\
ingān & enter; go in \\
āsęndan & send \\
behealdan & behold; look at \\
bebēodan & command; instruct \\
ęfthweorfan & return; turn back \\
lǣdan & lead; bring; take a person \\
sprecan & speak with someone \\
begān & pursue; practise study \\
āstīgan & descend \\
sęllan & sell \\
wyrcean & make; work; perform \\
\end{longtable}
\par\endgroup

% mdtable:T31:end
\reserve{7\baselineskip}
% mdparagraph:P123:begin
\subsection*{7.1. Weak verbs}\BookHeading{nav-ch13-learner-0035}{subsubsection}{7.1. Weak verbs}{0}
% mdparagraph:P123:end

% mdparagraph:P124:begin
The selected weak II pattern has present -ie, -ast, -að, plural -iað; preterite -ode, -odest, -ode, plural -odon. The finite model below supplies every indicative cell for the six verbs in the diagnostic table. Alternative -ige and -ade formations are outside the selected pattern.
% mdparagraph:P124:end

% teaching: NEW-U13-PAR-FORDILGIAN-IND
% teaching: NEW-U13-GRAM-U13-FORDILGIAN-INDICATIVE
\reserve{7\baselineskip}
% mdparagraph:P125:begin
\subsection*{Table 25. fordilgian: present and preterite indicative}\BookHeading{nav-ch13-learner-0036}{subsubsection}{Table 25. fordilgian: present and preterite indicative}{0}
% mdparagraph:P125:end

% table: NEW-U13-PAR-FORDILGIAN-IND
% mdtable:T32:begin
% course-table-style: 2026-10
% course-panel:T32-P01:begin
\par\begingroup
\fontsize{10.5}{13.2}\selectfont
\setlength{\LTleft}{0pt}\setlength{\LTright}{\fill}
\noindent
\begin{tabular}{@{}>{\raggedright\arraybackslash}p{39.459mm}>{\raggedright\arraybackslash}p{39.459mm}>{\raggedright\arraybackslash}p{39.459mm}@{}}
\toprule
\textbf{Person and number} & \textbf{Present} & \textbf{Preterite} \\
\midrule
1 sg. & fordilgie & fordilgode \\
2 sg. & fordilgast & fordilgodest \\
3 sg. & fordilgað & fordilgode \\
1–3 pl. & fordilgiað & fordilgodon \\
\bottomrule
\end{tabular}
\par\endgroup
% course-panel:T32-P01:end
% mdtable:T32:end
% teaching: NEW-U13-PAR-GEWUNIAN-IND
% teaching: NEW-U13-GRAM-U13-GEWUNIAN-INDICATIVE
% teaching: NEW-U13-PAR-ASCIAN-IND
% teaching: NEW-U13-GRAM-U13-ASCIAN-INDICATIVE
% teaching: NEW-U13-PAR-HALSIAN-IND
% teaching: NEW-U13-GRAM-U13-HALSIAN-INDICATIVE
% teaching: NEW-U13-PAR-WUNIAN-IND
% teaching: NEW-U13-GRAM-U13-WUNIAN-INDICATIVE
% teaching: NEW-U13-PAR-HATIAN-IND
% teaching: NEW-U13-GRAM-U13-HATIAN-INDICATIVE
% mdtable:T33:begin
\par\begingroup
\fontsize{10.5}{13.2}\selectfont
\setlength{\LTleft}{0pt}\setlength{\LTright}{\fill}
\noindent
\begin{tabular}{@{}>{\raggedright\arraybackslash}p{21.33mm}>{\raggedright\arraybackslash}p{21.33mm}>{\raggedright\arraybackslash}p{22.08mm}>{\raggedright\arraybackslash}p{25.17mm}>{\raggedright\arraybackslash}p{22.83mm}@{}}
\toprule
\textbf{Verb} & \textbf{Stem} & \textbf{Present 3 sg} & \textbf{Preterite 3 sg} & \textbf{Complete model} \\
\midrule
fordilgian & fordilg & fordilgað & fordilgode & Table 25 \\
gewunian & gewun & gewunað & gewunode & Table 25 \\
āscian & āsc & āscað & āscode & Table 25 \\
hālsian & hāls & hālsað & hālsode & Table 25 \\
wunian & wun & wunað & wunode & Table 25 \\
hatian & hat & hatað & hatode & Table 25 \\
\bottomrule
\end{tabular}
\par\endgroup

% mdtable:T33:end
% mdparagraph:P126:begin
For gewunian, āscian, hālsian, wunian and hatian, replace initial fordilg- in every cell of Table 25 by gewun-, āsc-, hāls-, wun- or hat-, respectively. Retain the rest of the cell exactly. This gives all twelve present and preterite indicative positions for each verb; no mood or nonfinite system is added by this substitution.
% mdparagraph:P126:end

% mdparagraph:P127:begin
Weak I verbs add a dental preterite, but their present stems and consonant outcomes must be learned separately. Āstręccean alternates geminate and ungeminated present consonants; its unmutated past is āstreahte, not a past formed mechanically from stręcc-. Nēalǣcan selects nēalǣcte. The tables retain the distinct present, past and consonant patterns of onlīesan, āsęttan, gescrȳdan and āsęndan.
% mdparagraph:P127:end

% teaching: NEW-U13-PAR-ASTRECCEAN-IND
% teaching: NEW-U13-GRAM-U13-ASTRECCEAN-INDICATIVE
\reserve{7\baselineskip}
% mdparagraph:P128:begin
\subsection*{Table 26. āstręccean: present and preterite indicative}\BookHeading{nav-ch13-learner-0037}{subsubsection}{Table 26. āstręccean: present and preterite indicative}{0}
% mdparagraph:P128:end

% table: NEW-U13-PAR-ASTRECCEAN-IND
% mdtable:T34:begin
% course-table-style: 2026-10
% course-panel:T34-P01:begin
\par\begingroup
\fontsize{10.5}{13.2}\selectfont
\setlength{\LTleft}{0pt}\setlength{\LTright}{\fill}
\noindent
\begin{tabular}{@{}>{\raggedright\arraybackslash}p{39.459mm}>{\raggedright\arraybackslash}p{39.459mm}>{\raggedright\arraybackslash}p{39.459mm}@{}}
\toprule
\textbf{Person and number} & \textbf{Present} & \textbf{Preterite} \\
\midrule
1 sg. & āstręcce & āstreahte \\
2 sg. & āstręcst & āstreahtest \\
3 sg. & āstręcð & āstreahte \\
1–3 pl. & āstręcceað & āstreahton \\
\bottomrule
\end{tabular}
\par\endgroup
% course-panel:T34-P01:end
% mdtable:T34:end
% teaching: NEW-U13-PAR-ONLIESAN-IND
% teaching: NEW-U13-GRAM-U13-ONLIESAN-INDICATIVE
\reserve{7\baselineskip}
% mdparagraph:P129:begin
\subsection*{Table 27. onlīesan: present and preterite indicative}\BookHeading{nav-ch13-learner-0038}{subsubsection}{Table 27. onlīesan: present and preterite indicative}{0}
% mdparagraph:P129:end

% table: NEW-U13-PAR-ONLIESAN-IND
% mdtable:T35:begin
% course-table-style: 2026-10
% course-panel:T35-P01:begin
\par\begingroup
\fontsize{10.5}{13.2}\selectfont
\setlength{\LTleft}{0pt}\setlength{\LTright}{\fill}
\noindent
\begin{tabular}{@{}>{\raggedright\arraybackslash}p{39.459mm}>{\raggedright\arraybackslash}p{39.459mm}>{\raggedright\arraybackslash}p{39.459mm}@{}}
\toprule
\textbf{Person and number} & \textbf{Present} & \textbf{Preterite} \\
\midrule
1 sg. & onlīese & onlīesde \\
2 sg. & onlīest & onlīesdest \\
3 sg. & onlīest & onlīesde \\
1–3 pl. & onlīesað & onlīesdon \\
\bottomrule
\end{tabular}
\par\endgroup
% course-panel:T35-P01:end
% mdtable:T35:end
% teaching: NEW-U13-PAR-NEALAECAN-IND
% teaching: NEW-U13-GRAM-U13-NEALAECAN-INDICATIVE
\reserve{7\baselineskip}
% mdparagraph:P130:begin
\subsection*{Table 28. nēalǣcan: present and preterite indicative}\BookHeading{nav-ch13-learner-0039}{subsubsection}{Table 28. nēalǣcan: present and preterite indicative}{0}
% mdparagraph:P130:end

% table: NEW-U13-PAR-NEALAECAN-IND
% mdtable:T36:begin
% course-table-style: 2026-10
% course-panel:T36-P01:begin
\par\begingroup
\fontsize{10.5}{13.2}\selectfont
\setlength{\LTleft}{0pt}\setlength{\LTright}{\fill}
\noindent
\begin{tabular}{@{}>{\raggedright\arraybackslash}p{39.459mm}>{\raggedright\arraybackslash}p{39.459mm}>{\raggedright\arraybackslash}p{39.459mm}@{}}
\toprule
\textbf{Person and number} & \textbf{Present} & \textbf{Preterite} \\
\midrule
1 sg. & nēalǣce & nēalǣcte \\
2 sg. & nēalǣcst & nēalǣctest \\
3 sg. & nēalǣcð & nēalǣcte \\
1–3 pl. & nēalǣcað & nēalǣcton \\
\bottomrule
\end{tabular}
\par\endgroup
% course-panel:T36-P01:end
% mdtable:T36:end
% teaching: NEW-U13-PAR-ASETTAN-IND
% teaching: NEW-U13-GRAM-U13-ASETTAN-INDICATIVE
\reserve{7\baselineskip}
% mdparagraph:P131:begin
\subsection*{Table 29. āsęttan: present and preterite indicative}\BookHeading{nav-ch13-learner-0040}{subsubsection}{Table 29. āsęttan: present and preterite indicative}{0}
% mdparagraph:P131:end

% table: NEW-U13-PAR-ASETTAN-IND
% mdtable:T37:begin
% course-table-style: 2026-10
% course-panel:T37-P01:begin
\par\begingroup
\fontsize{10.5}{13.2}\selectfont
\setlength{\LTleft}{0pt}\setlength{\LTright}{\fill}
\noindent
\begin{tabular}{@{}>{\raggedright\arraybackslash}p{39.459mm}>{\raggedright\arraybackslash}p{39.459mm}>{\raggedright\arraybackslash}p{39.459mm}@{}}
\toprule
\textbf{Person and number} & \textbf{Present} & \textbf{Preterite} \\
\midrule
1 sg. & āsętte & āsętte \\
2 sg. & āsętst & āsęttest \\
3 sg. & āsętt & āsętte \\
1–3 pl. & āsęttað & āsętton \\
\bottomrule
\end{tabular}
\par\endgroup
% course-panel:T37-P01:end
% mdtable:T37:end
% teaching: NEW-U13-PAR-GESCRYDAN-IND
% teaching: NEW-U13-GRAM-U13-GESCRYDAN-INDICATIVE
\reserve{7\baselineskip}
% mdparagraph:P132:begin
\subsection*{Table 30. gescrȳdan: present and preterite indicative}\BookHeading{nav-ch13-learner-0041}{subsubsection}{Table 30. gescrȳdan: present and preterite indicative}{0}
% mdparagraph:P132:end

% table: NEW-U13-PAR-GESCRYDAN-IND
% mdtable:T38:begin
% course-table-style: 2026-10
% course-panel:T38-P01:begin
\par\begingroup
\fontsize{10.5}{13.2}\selectfont
\setlength{\LTleft}{0pt}\setlength{\LTright}{\fill}
\noindent
\begin{tabular}{@{}>{\raggedright\arraybackslash}p{39.459mm}>{\raggedright\arraybackslash}p{39.459mm}>{\raggedright\arraybackslash}p{39.459mm}@{}}
\toprule
\textbf{Person and number} & \textbf{Present} & \textbf{Preterite} \\
\midrule
1 sg. & gescrȳde & gescrȳdde \\
2 sg. & gescrȳtst & gescrȳddest \\
3 sg. & gescrȳt & gescrȳdde \\
1–3 pl. & gescrȳdað & gescrȳddon \\
\bottomrule
\end{tabular}
\par\endgroup
% course-panel:T38-P01:end
% mdtable:T38:end
% teaching: NEW-U13-PAR-ASENDAN-IND
% teaching: NEW-U13-GRAM-U13-ASENDAN-INDICATIVE
\reserve{7\baselineskip}
% mdparagraph:P133:begin
\subsection*{Table 31. āsęndan: present and preterite indicative}\BookHeading{nav-ch13-learner-0042}{subsubsection}{Table 31. āsęndan: present and preterite indicative}{0}
% mdparagraph:P133:end

% table: NEW-U13-PAR-ASENDAN-IND
% mdtable:T39:begin
% course-table-style: 2026-10
% course-panel:T39-P01:begin
\par\begingroup
\fontsize{10.5}{13.2}\selectfont
\setlength{\LTleft}{0pt}\setlength{\LTright}{\fill}
\noindent
\begin{tabular}{@{}>{\raggedright\arraybackslash}p{39.459mm}>{\raggedright\arraybackslash}p{39.459mm}>{\raggedright\arraybackslash}p{39.459mm}@{}}
\toprule
\textbf{Person and number} & \textbf{Present} & \textbf{Preterite} \\
\midrule
1 sg. & āsęnde & āsęnde \\
2 sg. & āsęntst & āsęndest \\
3 sg. & āsęnt & āsęnde \\
1–3 pl. & āsęndað & āsęndon \\
\bottomrule
\end{tabular}
\par\endgroup
% course-panel:T39-P01:end
% mdtable:T39:end
\reserve{7\baselineskip}
% mdparagraph:P134:begin
\subsection*{7.2. Strong and irregular finite verbs}\BookHeading{nav-ch13-learner-0043}{subsubsection}{7.2. Strong and irregular finite verbs}{0}
% mdparagraph:P134:end

% mdparagraph:P135:begin
Strong verbs distinguish present, preterite singular, preterite plural and past-participle stems. The second singular preterite takes the same vowel grade as the plural. In forceorfan the present second and third singular use forcierf-; the past distinguishes forcearf from forcurf-. Oferwinnan and ęfthweorfan have their own related strong III patterns. Keep every change shown in the tables, including nn plus ð becoming nð and the selected consonant clusters before -st.
% mdparagraph:P135:end

% teaching: NEW-U13-PAR-FORCEORFAN-IND
% teaching: NEW-U13-GRAM-U13-FORCEORFAN-INDICATIVE
\reserve{7\baselineskip}
% mdparagraph:P136:begin
\subsection*{Table 32. forceorfan: present and preterite indicative}\BookHeading{nav-ch13-learner-0044}{subsubsection}{Table 32. forceorfan: present and preterite indicative}{0}
% mdparagraph:P136:end

% table: NEW-U13-PAR-FORCEORFAN-IND
% mdtable:T40:begin
% course-table-style: 2026-10
% course-panel:T40-P01:begin
\par\begingroup
\fontsize{10.5}{13.2}\selectfont
\setlength{\LTleft}{0pt}\setlength{\LTright}{\fill}
\noindent
\begin{tabular}{@{}>{\raggedright\arraybackslash}p{39.459mm}>{\raggedright\arraybackslash}p{39.459mm}>{\raggedright\arraybackslash}p{39.459mm}@{}}
\toprule
\textbf{Person and number} & \textbf{Present} & \textbf{Preterite} \\
\midrule
1 sg. & forceorfe & forcearf \\
2 sg. & forcierfst & forcurfe \\
3 sg. & forcierfð & forcearf \\
1–3 pl. & forceorfað & forcurfon \\
\bottomrule
\end{tabular}
\par\endgroup
% course-panel:T40-P01:end
% mdtable:T40:end
% teaching: NEW-U13-PAR-OFERWINNAN-IND
% teaching: NEW-U13-GRAM-U13-OFERWINNAN-INDICATIVE
\reserve{7\baselineskip}
% mdparagraph:P137:begin
\subsection*{Table 33. oferwinnan: present and preterite indicative}\BookHeading{nav-ch13-learner-0045}{subsubsection}{Table 33. oferwinnan: present and preterite indicative}{0}
% mdparagraph:P137:end

% table: NEW-U13-PAR-OFERWINNAN-IND
% mdtable:T41:begin
% course-table-style: 2026-10
% course-panel:T41-P01:begin
\par\begingroup
\fontsize{10.5}{13.2}\selectfont
\setlength{\LTleft}{0pt}\setlength{\LTright}{\fill}
\noindent
\begin{tabular}{@{}>{\raggedright\arraybackslash}p{39.459mm}>{\raggedright\arraybackslash}p{39.459mm}>{\raggedright\arraybackslash}p{39.459mm}@{}}
\toprule
\textbf{Person and number} & \textbf{Present} & \textbf{Preterite} \\
\midrule
1 sg. & oferwinne & oferwann \\
2 sg. & oferwinst & oferwunne \\
3 sg. & oferwinð & oferwann \\
1–3 pl. & oferwinnað & oferwunnon \\
\bottomrule
\end{tabular}
\par\endgroup
% course-panel:T41-P01:end
% mdtable:T41:end
% teaching: NEW-U13-PAR-EFTHWEORFAN-IND
% teaching: NEW-U13-GRAM-U13-EFTHWEORFAN-INDICATIVE
\reserve{7\baselineskip}
% mdparagraph:P138:begin
\subsection*{Table 34. ęfthweorfan: present and preterite indicative}\BookHeading{nav-ch13-learner-0046}{subsubsection}{Table 34. ęfthweorfan: present and preterite indicative}{0}
% mdparagraph:P138:end

% table: NEW-U13-PAR-EFTHWEORFAN-IND
% mdtable:T42:begin
% course-table-style: 2026-10
% course-panel:T42-P01:begin
\par\begingroup
\fontsize{10.5}{13.2}\selectfont
\setlength{\LTleft}{0pt}\setlength{\LTright}{\fill}
\noindent
\begin{tabular}{@{}>{\raggedright\arraybackslash}p{39.459mm}>{\raggedright\arraybackslash}p{39.459mm}>{\raggedright\arraybackslash}p{39.459mm}@{}}
\toprule
\textbf{Person and number} & \textbf{Present} & \textbf{Preterite} \\
\midrule
1 sg. & ęfthweorfe & ęfthwearf \\
2 sg. & ęfthwierfst & ęfthwurfe \\
3 sg. & ęfthwierfð & ęfthwearf \\
1–3 pl. & ęfthweorfað & ęfthwurfon \\
\bottomrule
\end{tabular}
\par\endgroup
% course-panel:T42-P01:end
% mdtable:T42:end
% teaching: NEW-U13-PRINCIPAL-PARTS-FORCEORFAN
% support: NEW-U13-PRINCIPAL-PARTS-FORCEORFAN
% mdparagraph:P139:begin
\textbf{Principal-part support.} Citation support: forceorfan, forcearf, forcurfon, forcorfen. These four printed parts identify the present, preterite singular, preterite plural and past participle stems. Keep the row visible; only the separately assigned systems and uses are recall targets.
% mdparagraph:P139:end

% mdparagraph:P140:begin
Behealdan and ondrǣdan have the strong VII past stems behēold- and ondrēd-. Bebēodan has the strong II grades bēod-, bēad- and bud-, with bebīetst and bebīett in the present second and third singular. These lexical systems do not make every similarly prefixed verb predictable.
% mdparagraph:P140:end

% teaching: NEW-U13-PAR-BEHEALDAN-IND
% teaching: NEW-U13-GRAM-U13-BEHEALDAN-INDICATIVE
\reserve{7\baselineskip}
% mdparagraph:P141:begin
\subsection*{Table 35. behealdan: present and preterite indicative}\BookHeading{nav-ch13-learner-0047}{subsubsection}{Table 35. behealdan: present and preterite indicative}{0}
% mdparagraph:P141:end

% table: NEW-U13-PAR-BEHEALDAN-IND
% mdtable:T43:begin
% course-table-style: 2026-10
% course-panel:T43-P01:begin
\par\begingroup
\fontsize{10.5}{13.2}\selectfont
\setlength{\LTleft}{0pt}\setlength{\LTright}{\fill}
\noindent
\begin{tabular}{@{}>{\raggedright\arraybackslash}p{39.459mm}>{\raggedright\arraybackslash}p{39.459mm}>{\raggedright\arraybackslash}p{39.459mm}@{}}
\toprule
\textbf{Person and number} & \textbf{Present} & \textbf{Preterite} \\
\midrule
1 sg. & behealde & behēold \\
2 sg. & behieltst & behēolde \\
3 sg. & behielt & behēold \\
1–3 pl. & behealdað & behēoldon \\
\bottomrule
\end{tabular}
\par\endgroup
% course-panel:T43-P01:end
% mdtable:T43:end
% teaching: NEW-U13-PAR-ONDRAEDAN-IND
% teaching: NEW-U13-GRAM-U13-ONDRAEDAN-INDICATIVE
\reserve{7\baselineskip}
% mdparagraph:P142:begin
\subsection*{Table 36. ondrǣdan: present and preterite indicative}\BookHeading{nav-ch13-learner-0048}{subsubsection}{Table 36. ondrǣdan: present and preterite indicative}{0}
% mdparagraph:P142:end

% table: NEW-U13-PAR-ONDRAEDAN-IND
% mdtable:T44:begin
% course-table-style: 2026-10
% course-panel:T44-P01:begin
\par\begingroup
\fontsize{10.5}{13.2}\selectfont
\setlength{\LTleft}{0pt}\setlength{\LTright}{\fill}
\noindent
\begin{tabular}{@{}>{\raggedright\arraybackslash}p{39.459mm}>{\raggedright\arraybackslash}p{39.459mm}>{\raggedright\arraybackslash}p{39.459mm}@{}}
\toprule
\textbf{Person and number} & \textbf{Present} & \textbf{Preterite} \\
\midrule
1 sg. & ondrǣde & ondrēd \\
2 sg. & ondrǣtst & ondrēde \\
3 sg. & ondrǣt & ondrēd \\
1–3 pl. & ondrǣdað & ondrēdon \\
\bottomrule
\end{tabular}
\par\endgroup
% course-panel:T44-P01:end
% mdtable:T44:end
% teaching: NEW-U13-PAR-BEBEODAN-IND
% teaching: NEW-U13-GRAM-U13-BEBEODAN-INDICATIVE
\reserve{7\baselineskip}
% mdparagraph:P143:begin
\subsection*{Table 37. bebēodan: present and preterite indicative}\BookHeading{nav-ch13-learner-0049}{subsubsection}{Table 37. bebēodan: present and preterite indicative}{0}
% mdparagraph:P143:end

% table: NEW-U13-PAR-BEBEODAN-IND
% mdtable:T45:begin
% course-table-style: 2026-10
% course-panel:T45-P01:begin
\par\begingroup
\fontsize{10.5}{13.2}\selectfont
\setlength{\LTleft}{0pt}\setlength{\LTright}{\fill}
\noindent
\begin{tabular}{@{}>{\raggedright\arraybackslash}p{39.459mm}>{\raggedright\arraybackslash}p{39.459mm}>{\raggedright\arraybackslash}p{39.459mm}@{}}
\toprule
\textbf{Person and number} & \textbf{Present} & \textbf{Preterite} \\
\midrule
1 sg. & bebēode & bebēad \\
2 sg. & bebīetst & bebude \\
3 sg. & bebīett & bebēad \\
1–3 pl. & bebēodað & bebudon \\
\bottomrule
\end{tabular}
\par\endgroup
% course-panel:T45-P01:end
% mdtable:T45:end
% mdparagraph:P144:begin
Fordōn and ingān use the corresponding finite patterns of dōn and gān with their attested prefixes. The complete ingān subjunctive follows its indicative table. The new productive use is the present subjunctive in the worthy-content construction below; the other subjunctive cells complete the morphological system without adding new clause uses. The reading’s ingangende is separately supplied and is not generated as a regular participle from ingān.
% mdparagraph:P144:end

% teaching: NEW-U13-PAR-FORDON-IND
% teaching: NEW-U13-GRAM-U13-FORDON-INDICATIVE
\reserve{7\baselineskip}
% mdparagraph:P145:begin
\subsection*{Table 38. fordōn: present and preterite indicative}\BookHeading{nav-ch13-learner-0050}{subsubsection}{Table 38. fordōn: present and preterite indicative}{0}
% mdparagraph:P145:end

% table: NEW-U13-PAR-FORDON-IND
% mdtable:T46:begin
% course-table-style: 2026-10
% course-panel:T46-P01:begin
\par\begingroup
\fontsize{10.5}{13.2}\selectfont
\setlength{\LTleft}{0pt}\setlength{\LTright}{\fill}
\noindent
\begin{tabular}{@{}>{\raggedright\arraybackslash}p{39.459mm}>{\raggedright\arraybackslash}p{39.459mm}>{\raggedright\arraybackslash}p{39.459mm}@{}}
\toprule
\textbf{Person and number} & \textbf{Present} & \textbf{Preterite} \\
\midrule
1 sg. & fordō & fordyde \\
2 sg. & fordēst & fordydest \\
3 sg. & fordēð & fordyde \\
1–3 pl. & fordōð & fordydon \\
\bottomrule
\end{tabular}
\par\endgroup
% course-panel:T46-P01:end
% mdtable:T46:end
% teaching: NEW-U13-PAR-INGAN-IND
% teaching: NEW-U13-GRAM-U13-INGAN-INDICATIVE
\reserve{7\baselineskip}
% mdparagraph:P146:begin
\subsection*{Table 39. ingān: present and preterite indicative}\BookHeading{nav-ch13-learner-0051}{subsubsection}{Table 39. ingān: present and preterite indicative}{0}
% mdparagraph:P146:end

% table: NEW-U13-PAR-INGAN-IND
% mdtable:T47:begin
% course-table-style: 2026-10
% course-panel:T47-P01:begin
\par\begingroup
\fontsize{10.5}{13.2}\selectfont
\setlength{\LTleft}{0pt}\setlength{\LTright}{\fill}
\noindent
\begin{tabular}{@{}>{\raggedright\arraybackslash}p{39.459mm}>{\raggedright\arraybackslash}p{39.459mm}>{\raggedright\arraybackslash}p{39.459mm}@{}}
\toprule
\textbf{Person and number} & \textbf{Present} & \textbf{Preterite} \\
\midrule
1 sg. & ingā & inēode \\
2 sg. & ingǣst & inēodest \\
3 sg. & ingǣð & inēode \\
1–3 pl. & ingāð & inēodon \\
\bottomrule
\end{tabular}
\par\endgroup
% course-panel:T47-P01:end
% mdtable:T47:end
% teaching: NEW-U13-PAR-INGAN-SUBJ
% teaching: NEW-U13-GRAM-U13-INGAN-SUBJUNCTIVE
\reserve{7\baselineskip}
% mdparagraph:P147:begin
\subsection*{Table 40. ingān: present and preterite subjunctive}\BookHeading{nav-ch13-learner-0052}{subsubsection}{Table 40. ingān: present and preterite subjunctive}{0}
% mdparagraph:P147:end

% table: NEW-U13-PAR-INGAN-SUBJ
% mdtable:T48:begin
% course-table-style: 2026-10
% course-panel:T48-P01:begin
\par\begingroup
\fontsize{10.5}{13.2}\selectfont
\setlength{\LTleft}{0pt}\setlength{\LTright}{\fill}
\noindent
\begin{tabular}{@{}>{\raggedright\arraybackslash}p{39.459mm}>{\raggedright\arraybackslash}p{39.459mm}>{\raggedright\arraybackslash}p{39.459mm}@{}}
\toprule
\textbf{Person and number} & \textbf{Present} & \textbf{Preterite} \\
\midrule
1 sg. & ingā & inēode \\
2 sg. & ingā & inēode \\
3 sg. & ingā & inēode \\
1–3 pl. & ingān & inēoden \\
\bottomrule
\end{tabular}
\par\endgroup
% course-panel:T48-P01:end
% mdtable:T48:end
\reserve{7\baselineskip}
% mdparagraph:P148:begin
\subsection*{7.3. Commands, infinitives and participle bases}\BookHeading{nav-ch13-learner-0053}{subsubsection}{7.3. Commands, infinitives and participle bases}{0}
% mdparagraph:P148:end

% teaching: NEW-U13-VOC-WINNAN
% teaching: NEW-U13-PAR-WINNAN
% mdparagraph:P149:begin
The imperative table gives singular and plural commands for all four verbs. Winnan means “toil”; behealdan “behold”; ondrǣdan “fear”; lǣdan “lead or bring a person.” Prohibitions use the already taught ne with the appropriate command. The earlier begān command pair remains begā, begāð, now with the explicit sense “pursue study.”
% mdparagraph:P149:end

% teaching: NEW-U13-PAR-WINNAN-IMP
% teaching: NEW-U13-PAR-BEHEALDAN-IMP
% teaching: NEW-U13-PAR-ONDRAEDAN-IMP
% teaching: NEW-U13-PAR-LAEDAN-IMP
% teaching: NEW-U13-GRAM-U13-WINNAN-IMPERATIVE
% teaching: NEW-U13-GRAM-U13-BEHEALDAN-IMPERATIVE
% teaching: NEW-U13-GRAM-U13-ONDRAEDAN-IMPERATIVE
% teaching: NEW-U13-GRAM-U13-LAEDAN-IMPERATIVE
\reserve{7\baselineskip}
% mdparagraph:P150:begin
\subsection*{Table 41. Complete singular and plural commands}\BookHeading{nav-ch13-learner-0054}{subsubsection}{Table 41. Complete singular and plural commands}{0}
% mdparagraph:P150:end

% table: NEW-U13-DISPLAY-IMP
% mdtable:T49:begin
% course-table-style: 2026-10
% course-panel:T49-P01:begin
\par\begingroup
\fontsize{10.5}{13.2}\selectfont
\setlength{\LTleft}{0pt}\setlength{\LTright}{\fill}
\noindent
\begin{tabular}{@{}>{\raggedright\arraybackslash}p{24.232mm}>{\raggedright\arraybackslash}p{22.13mm}>{\raggedright\arraybackslash}p{22.13mm}>{\raggedright\arraybackslash}p{22.13mm}>{\raggedright\arraybackslash}p{22.13mm}@{}}
\toprule
\textbf{Person and number} & \textbf{winnan} & \textbf{behealdan} & \textbf{ondrǣdan} & \textbf{lǣdan} \\
\midrule
2 sg. & win & beheald & ondrǣd & lǣd \\
2 pl. & winnað & behealdað & ondrǣdað & lǣdað \\
\bottomrule
\end{tabular}
\par\endgroup
% course-panel:T49-P01:end
% mdtable:T49:end
% mdparagraph:P151:begin
The plain infinitive follows taught modal and willan constructions. The inflected infinitive adds -ne after the plain infinitive’s final n, producing endings including -anne, -enne, -ianne and, in fordōnne, -ōnne. Use tō with the inflected form in previously taught functions and in the explicitly supported adjective expression below. The table does not introduce every possible infinitive construction.
% mdparagraph:P151:end

% teaching: NEW-U13-PAR-ASCIAN-INF
% teaching: NEW-U13-PAR-FORDON-INF
% teaching: NEW-U13-PAR-LAEDAN-INF
% teaching: NEW-U13-PAR-SPRECAN-INF
% teaching: NEW-U13-GRAM-U13-ASCIAN-INFINITIVE-PAIR
% teaching: NEW-U13-GRAM-U13-FORDON-INFINITIVE-PAIR
% teaching: NEW-U13-GRAM-U13-LAEDAN-INFINITIVE-PAIR
% teaching: NEW-U13-GRAM-U13-SPRECAN-INFINITIVE-PAIR
\reserve{7\baselineskip}
% mdparagraph:P152:begin
\subsection*{Table 42. Complete plain and inflected infinitives}\BookHeading{nav-ch13-learner-0055}{subsubsection}{Table 42. Complete plain and inflected infinitives}{0}
% mdparagraph:P152:end

% table: NEW-U13-DISPLAY-INF
% mdtable:T50:begin
\par\begingroup
\fontsize{10.5}{13.2}\selectfont
\setlength{\LTleft}{0pt}\setlength{\LTright}{\fill}
\noindent
\begin{tabular}{@{}>{\raggedright\arraybackslash}p{31.82mm}>{\raggedright\arraybackslash}p{40.41mm}>{\raggedright\arraybackslash}p{46.14mm}@{}}
\toprule
\textbf{Verb} & \textbf{Plain infinitive} & \textbf{Inflected infinitive} \\
\midrule
āscian & āscian & āscianne \\
fordōn & fordōn & fordōnne \\
lǣdan & lǣdan & lǣdanne \\
sprecan & sprecan & sprecanne \\
\bottomrule
\end{tabular}
\par\endgroup

% mdtable:T50:end
% mdparagraph:P153:begin
The following table gives complete present and past participle base pairs. Base formation is a production target. In the passive constructions taught here, only the past base is used, uninflected, with wesan; the patient becomes the subject. The table does not introduce full attributive participle declensions or unrestricted uses of every present participle. Gescrȳd has one final d; distinguish the finite present gescrȳt from this past base.
% mdparagraph:P153:end

% teaching: NEW-U13-PAR-FORDILGIAN-PART
% teaching: NEW-U13-PAR-ONLIESAN-PART
% teaching: NEW-U13-PAR-ASETTAN-PART
% teaching: NEW-U13-PAR-GESCRYDAN-PART
% teaching: NEW-U13-PAR-ASENDAN-PART
% teaching: NEW-U13-GRAM-U13-FORDILGIAN-PARTICIPLE-BASES
% teaching: NEW-U13-GRAM-U13-ONLIESAN-PARTICIPLE-BASES
% teaching: NEW-U13-GRAM-U13-ASETTAN-PARTICIPLE-BASES
% teaching: NEW-U13-GRAM-U13-GESCRYDAN-PARTICIPLE-BASES
% teaching: NEW-U13-GRAM-U13-ASENDAN-PARTICIPLE-BASES
\reserve{7\baselineskip}
% mdparagraph:P154:begin
\subsection*{Table 43. Complete present and past participle bases}\BookHeading{nav-ch13-learner-0056}{subsubsection}{Table 43. Complete present and past participle bases}{0}
% mdparagraph:P154:end

% table: NEW-U13-DISPLAY-PART
% mdtable:T51:begin
\par\begingroup
\fontsize{10.5}{13.2}\selectfont
\setlength{\LTleft}{0pt}\setlength{\LTright}{\fill}
\noindent
\begin{tabular}{@{}>{\raggedright\arraybackslash}p{32.04mm}>{\raggedright\arraybackslash}p{47.09mm}>{\raggedright\arraybackslash}p{39.25mm}@{}}
\toprule
\textbf{Verb} & \textbf{Present participle} & \textbf{Past participle} \\
\midrule
fordilgian & fordilgiende & fordilgod \\
onlīesan & onlīesende & onlīesed \\
āsęttan & āsęttende & āsętt \\
gescrȳdan & gescrȳdende & gescrȳd \\
āsęndan & āsęndende & āsęnd \\
\bottomrule
\end{tabular}
\par\endgroup

% mdtable:T51:end
\reserve{7\baselineskip}
% mdparagraph:P155:begin
\section*{8. Ordinary constructions and smaller words}\BookHeading{nav-ch13-learner-0057}{subsection}{8. Ordinary constructions and smaller words}{1}
% mdparagraph:P155:end

% mdparagraph:P156:begin
Learn a verb with the case or clause it governs. English prepositions are not a reliable guide to the Old English construction. The following frames state the uses required by the exercises; an accusative object remains accusative when it is fronted.
% mdparagraph:P156:end

% teaching: NEW-U13-GRAM-U13-FORDILGIAN-ACC
% teaching: NEW-U13-GRAM-U13-GEWUNIAN-INF
% teaching: NEW-U13-GRAM-U13-ASCIAN-PERSON-THING
% teaching: NEW-U13-GRAM-U13-ASTRECCEAN-REFLEXIVE
% teaching: NEW-U13-GRAM-U13-ONLIESAN-ACC-PASSIVE
% teaching: NEW-U13-GRAM-U13-NEALAECAN-TO-DAT
% teaching: NEW-U13-GRAM-U13-HALSIAN-REQUEST
% teaching: NEW-U13-GRAM-U13-FORCEORFAN-ACC
% teaching: NEW-U13-GRAM-U13-ASETTAN-ACC-PASSIVE
% teaching: NEW-U13-GRAM-U13-WUNIAN-LOCATION
% teaching: NEW-U13-GRAM-U13-OFERWINNAN-ACC
% teaching: NEW-U13-GRAM-U13-GESCRYDAN-CLOTHING
% teaching: NEW-U13-GRAM-U13-ONDRAEDAN-ACC
% teaching: NEW-U13-GRAM-U13-FORDON-ACC
% teaching: NEW-U13-GRAM-U13-INGAN-UNDER
% teaching: NEW-U13-GRAM-U13-ASENDAN-PERSON
% teaching: NEW-U13-GRAM-U13-BEHEALDAN-ACC
% teaching: NEW-U13-GRAM-U13-HATIAN-ACC
% teaching: NEW-U13-GRAM-U13-BEBEODAN-CONTENT
% teaching: NEW-U13-GRAM-U13-EFTHWEORFAN-RETURN
% teaching: NEW-U13-GRAM-U13-LAEDAN-PERSON-DIRECTION
% teaching: NEW-U13-GRAM-U13-SPRECAN-CONVERSATION-MID
% teaching: NEW-U13-GRAM-U13-BEGAN-STUDY
% teaching: NEW-U13-GRAM-U13-ASTIGAN-DESCENT
% teaching: NEW-U13-GRAM-U13-SELLAN-SALE
% teaching: NEW-U13-GRAM-U13-WYRCEAN-MAKING
% teaching: NEW-U13-GRAM-U13-INSTRUMENTAL-ROUTE
% teaching: NEW-U13-GRAM-U13-PLEASANT-INFINITIVE
% teaching: NEW-U13-GRAM-U13-GE-PAIRED
% teaching: NEW-U13-GRAM-U13-BEFORAN-POSITIONING-ACC
% teaching: NEW-U13-GRAM-U13-WITH-PROXIMITY-ACC
% teaching: NEW-U13-GRAM-U13-UNDER-DIRECTION-ACC
% teaching: NEW-U13-GRAM-U13-WORTHY-CONTENT
% teaching: NEW-U13-GRAM-U13-HWILC-INDEFINITE
% teaching: NEW-U13-GRAM-U13-DON-WELL-DATIVE
% teaching: NEW-U13-GRAM-U13-SINGAN-SONG
\reserve{7\baselineskip}
% mdparagraph:P157:begin
\subsection*{8.1. New finite verbs}\BookHeading{nav-ch13-learner-0058}{subsubsection}{8.1. New finite verbs}{0}
% mdparagraph:P157:end

% mdtable:T52:begin
\par\begingroup
\fontsize{10.5}{13.2}\selectfont
\setlength{\LTleft}{0pt}\setlength{\LTright}{\fill}
\begin{longtable}{@{}>{\raggedright\arraybackslash}p{21.67mm}>{\raggedright\arraybackslash}p{99.52mm}@{}}
\toprule
\textbf{Word or expression} & \textbf{Construction} \\
\midrule
\endfirsthead
\toprule
\textbf{Word or expression} & \textbf{Construction} \\
\midrule
\endhead
\midrule
\multicolumn{2}{r}{\footnotesize Continued on the following page.}\\
\endfoot
\bottomrule
\endlastfoot
fordilgian & Destroy an accusative thing. The corresponding passive uses the uninflected past participle with wesan. \\
gewunian & Be accustomed to a plain infinitive. The subject of the finite verb also performs the infinitive’s action. \\
āscian & Ask an accusative person about a genitive thing. A supported indirect question can give the content; do not use two accusatives for person and thing. \\
āstręccean & Prostrate oneself: the reflexive pronoun is accusative and refers to the subject. \\
onlīesan & Release an accusative person or thing. In the passive, use the uninflected past base onlīesed with wesan. \\
nēalǣcan & Approach a goal expressed by tō with the dative. \\
hālsian & Entreat an accusative person. A requested action may be expressed by ðæt with the subjunctive. \\
forceorfan & Cut down an accusative tree or other object. \\
āsęttan & Set or place an accusative object. The passive uses the uninflected past base with wesan; the placement phrase has its own case requirements. \\
wunian & Dwell or remain, either without a complement or with an ordinary static location. \\
oferwinnan & Overcome an accusative person or force. \\
gescrȳdan; mid & Clothe an accusative person; mid with the dative names the garment. The passive retains mid with the dative. English “clothed in” does not require Old English on. \\
ondrǣdan & Fear an accusative person or thing. A dative pronoun referring to the subject is possible but need not be added to every sentence. \\
fordōn & Destroy accusative persons or things. Its plain infinitive can follow a taught modal. \\
ingān & Enter or go in. In the supported roof-entry construction, under governs an accusative noun phrase. \\
āsęndan & Send an accusative person or thing. The corresponding passive uses the uninflected past base with wesan. \\
behealdan & Behold or look at an accusative object. \\
hatian & Hate an accusative person. \\
bebēodan & Command a dative addressee. A ðæt clause may express the command with the subjunctive; a relative can identify what is commanded. \\
ęfthweorfan & Return or turn back, without a direct object. A route expression can specify the way. \\
\end{longtable}
\par\endgroup

% mdtable:T52:end
\reserve{7\baselineskip}
% mdparagraph:P158:begin
\subsection*{8.2. Additional senses of familiar verbs}\BookHeading{nav-ch13-learner-0059}{subsubsection}{8.2. Additional senses of familiar verbs}{0}
% mdparagraph:P158:end

% mdtable:T53:begin
\par\begingroup
\fontsize{10.5}{13.2}\selectfont
\setlength{\LTleft}{0pt}\setlength{\LTright}{\fill}
\noindent
\begin{tabular}{@{}>{\raggedright\arraybackslash}p{21.53mm}>{\raggedright\arraybackslash}p{99.66mm}@{}}
\toprule
\textbf{Word or expression} & \textbf{Construction} \\
\midrule
lǣdan & Lead, bring or take an accusative person. Mid with the dative can name a companion, and hider can give the direction “here.” \\
sprecan; mid & Speak with someone: mid with the dative names the interlocutor. \\
begān & Pursue or practise an accusative activity, including leornung “study.” Use the finite and command forms already taught. \\
āstīgan; of & Descend; of with the dative marks the source. This is an additional meaning of the previously encountered verb. \\
sęllan & Sell an accusative possession. This sense is distinct from the earlier meaning “give.” \\
wyrcean & Make an accusative thing. Its earlier “work; perform” sense also permits an accusative deed, including abstract wrongdoing. \\
\bottomrule
\end{tabular}
\par\endgroup

% mdtable:T53:end
\reserve{7\baselineskip}
% mdparagraph:P159:begin
\subsection*{8.3. Other required constructions}\BookHeading{nav-ch13-learner-0060}{subsubsection}{8.3. Other required constructions}{0}
% mdparagraph:P159:end

% mdtable:T54:begin
\par\begingroup
\fontsize{10.5}{13.2}\selectfont
\setlength{\LTleft}{0pt}\setlength{\LTright}{\fill}
\noindent
\begin{tabular}{@{}>{\raggedright\arraybackslash}p{20.97mm}>{\raggedright\arraybackslash}p{100.22mm}@{}}
\toprule
\textbf{Word or expression} & \textbf{Construction} \\
\midrule
ðȳ ilcan wege & “By the same way”: ðȳ is instrumental; ilcan and wege have the already familiar dative shapes in this route expression. Learn the whole expression; it does not add a new ilca declension. \\
wynsum tō hīeranne & “Pleasant to hear”: tō takes the inflected infinitive. The adjective agrees with what is pleasant. This expression supplies the evaluative use needed here. \\
ge … ge & “Both … and”: the paired conjunction joins parallel constituents. Both occurrences have short e. \\
beforan & Before a person, with the accusative in a movement or positioning construction. A static dative is also a valid construal; follow the requested exercise target where one is given. \\
wið & Beside, with the accusative, including static standing beside someone. This accusative does not itself express motion. \\
under & Under, with the accusative for entering beneath a roof. Keep this distinct from the earlier static under with the dative. \\
wierðe & Worthy, with a ðæt clause and present subjunctive for the proposed entry. This bounded construction uses ingān; it does not introduce a general rule for every adjective. \\
hwilc & Any or some, as an indefinite adjective. This use is distinct from the interrogative “which?”; agreement follows the noun qualified. \\
dōn wel & Do well to a person in the dative, without an expressed object naming a deed. \\
singan & Sing an accusative song: sang can be the direct object although it names the action’s own product. \\
\bottomrule
\end{tabular}
\par\endgroup

% mdtable:T54:end
% mdparagraph:P160:begin
In familiar verbs, the earlier inflection remains available, but the meanings just specified must be learned separately. Thus knowing sęllan “give” did not by itself establish “sell,” nor did knowing mid for accompaniment establish its conversational or clothing use.
% mdparagraph:P160:end

% teaching: NEW-U13-VOC-THAER
% teaching: NEW-U13-PAR-THAER
% teaching: NEW-U13-VOC-EFT
% teaching: NEW-U13-PAR-EFT
% teaching: NEW-U13-VOC-OF
% teaching: NEW-U13-PAR-OF
% teaching: NEW-U13-VOC-BUTAN
% teaching: NEW-U13-PAR-BUTAN
% teaching: NEW-U13-VOC-OTH-PREPOSITION
% teaching: NEW-U13-PAR-OTH-PREPOSITION
% teaching: NEW-U13-VOC-THIDER
% teaching: NEW-U13-PAR-THIDER
% teaching: NEW-U13-VOC-AER-ADVERB
% teaching: NEW-U13-PAR-AER-ADVERB
% teaching: NEW-U13-VOC-STYCCEMAELUM
% teaching: NEW-U13-PAR-STYCCEMAELUM
% teaching: NEW-U13-VOC-NEAH-ADVERB
% teaching: NEW-U13-PAR-NEAH-ADVERB
% teaching: NEW-U13-VOC-OTHOTHE
% teaching: NEW-U13-PAR-OTHOTHE
% teaching: NEW-U13-VOC-GE-CONJUNCTION
% teaching: NEW-U13-PAR-GE-CONJUNCTION
% teaching: NEW-U13-VOC-BEFORAN
% teaching: NEW-U13-VOC-WITH
% teaching: NEW-U13-VOC-UNDER
% teaching: NEW-U13-VOC-MID
% teaching: NEW-U13-PAR-MID
% teaching: NEW-U13-VOC-HWILC
% mdtable:T55:begin
\par\begingroup
\fontsize{10.5}{13.2}\selectfont
\setlength{\LTleft}{0pt}\setlength{\LTright}{\fill}
\noindent
\begin{tabular}{@{}>{\raggedright\arraybackslash}p{41.92mm}>{\raggedright\arraybackslash}p{79.27mm}@{}}
\toprule
\textbf{Word} & \textbf{Meaning and use} \\
\midrule
ðǣr & there \\
ęft & again; back \\
of & out of; from, governing dative \\
būtan & without, governing dative \\
oð & up to; until, governing accusative \\
ðider & thither; to that place \\
ǣr & before; earlier \\
ǣrest & first; at first \\
styccemǣlum & little by little \\
nēah & near; at hand \\
oððe & or \\
ge & both ... and (paired ge ... ge) \\
\bottomrule
\end{tabular}
\par\endgroup

% mdtable:T55:end
% mdparagraph:P161:begin
Distinguish ðǣr “there,” a location, from ðider “to that place,” a direction. Here ęft has the explicitly taught sense “again; back”; the reading also supplies its contextual translation “afterwards.” Nēah in the ordinary expression “be near” or “be at hand” is an adverb. Oð with the accusative marks a limit; the reading separately supplies the conjunction oð-ðæt “until.”
% mdparagraph:P161:end

\reserve{7\baselineskip}
% mdparagraph:P162:begin
\section*{9. New practice}\BookHeading{nav-ch13-learner-0061}{subsection}{9. New practice}{1}
% mdparagraph:P162:end

% mdparagraph:P163:begin
For analysis tasks, preserve all possible references when the context does not decide between them. For composition, meet the stated target while keeping the intended participants, agreement and meaning. A valid alternative order and compliance with a particular requested order are separate questions.
% mdparagraph:P163:end

\reserve{7\baselineskip}
% mdparagraph:P164:begin
\subsection*{Forms}\BookHeading{nav-ch13-learner-0062}{subsubsection}{Forms}{0}
% mdparagraph:P164:end

\par\begin{minipage}{\linewidth}\raggedright
% mdparagraph:P165:begin
\subsection*{F01}\BookHeading{nav-ch13-learner-0063}{subsubsection}{F01}{0}
% mdparagraph:P165:end

% item: NEW-U13-F01
% mdparagraph:P166:begin
Give the complete nominative, genitive, dative and accusative singular and plural of mōdor, dohtor and brōðor. Explain the selected oblique and plural stems.
% mdparagraph:P166:end

\end{minipage}\par\addvspace{4pt}

\par\begin{minipage}{\linewidth}\raggedright
% mdparagraph:P167:begin
\subsection*{F02}\BookHeading{nav-ch13-learner-0064}{subsubsection}{F02}{0}
% mdparagraph:P167:end

% item: NEW-U13-F02
% mdparagraph:P168:begin
Give the complete nominative, genitive, dative and accusative singular and plural of æx and ēagðȳrel. Explain the retained stem in æx and the syncopated stem before endings in ēagðȳrel.
% mdparagraph:P168:end

\end{minipage}\par\addvspace{4pt}

\par\begin{minipage}{\linewidth}\raggedright
% mdparagraph:P169:begin
\subsection*{F03}\BookHeading{nav-ch13-learner-0065}{subsubsection}{F03}{0}
% mdparagraph:P169:end

% item: NEW-U13-F03
% mdparagraph:P170:begin
Give the genitive and dative singular and the nominative and genitive plural of cirice, mynster and wyrttruma. State the model used for each.
% mdparagraph:P170:end

\end{minipage}\par\addvspace{4pt}

\par\begin{minipage}{\linewidth}\raggedright
% mdparagraph:P171:begin
\subsection*{F04}\BookHeading{nav-ch13-learner-0066}{subsubsection}{F04}{0}
% mdparagraph:P171:end

% item: NEW-U13-F04
% mdparagraph:P172:begin
For swēte, give all strong nominative, genitive, dative and accusative singular and plural forms in all three genders. Then give the weak masculine nominative singular of its comparative and both strong and weak masculine nominative singular of its superlative. Explain the stems and endings.
% mdparagraph:P172:end

\end{minipage}\par\addvspace{4pt}

\par\begin{minipage}{\linewidth}\raggedright
% mdparagraph:P173:begin
\subsection*{F05}\BookHeading{nav-ch13-learner-0067}{subsubsection}{F05}{0}
% mdparagraph:P173:end

% item: NEW-U13-F05
% mdparagraph:P174:begin
Give the complete present and preterite indicative of forceorfan: all three persons in singular and plural. Using the printed principal parts, explain the present singular changes and the distribution of preterite stems.
% mdparagraph:P174:end

% support: NEW-U13-PRINCIPAL-PARTS-FORCEORFAN
% mdparagraph:P175:begin
\textbf{Supplied help.} Citation support: forceorfan, forcearf, forcurfon, forcorfen. These four printed parts identify the present, preterite singular, preterite plural and past participle stems. Keep the row visible; only the separately assigned systems and uses are recall targets.
% mdparagraph:P175:end

\end{minipage}\par\addvspace{4pt}

\par\begin{minipage}{\linewidth}\raggedright
% mdparagraph:P176:begin
\subsection*{F06}\BookHeading{nav-ch13-learner-0068}{subsubsection}{F06}{0}
% mdparagraph:P176:end

% item: NEW-U13-F06
% mdparagraph:P177:begin
Give the singular and plural commands of winnan, begān and behealdan. For fordilgian and onlīesan, give the base present and past participles.
% mdparagraph:P177:end

\end{minipage}\par\addvspace{4pt}

\reserve{7\baselineskip}
% mdparagraph:P178:begin
\subsection*{Clause analysis and reference}\BookHeading{nav-ch13-learner-0069}{subsubsection}{Clause analysis and reference}{0}
% mdparagraph:P178:end

\par\begin{minipage}{\linewidth}\raggedright
% mdparagraph:P179:begin
\subsection*{A01}\BookHeading{nav-ch13-learner-0070}{subsubsection}{A01}{0}
% mdparagraph:P179:end

% item: NEW-U13-A01
% mdparagraph:P180:begin
Translate both clauses. Identify the subject, direct object and finite verb in each; explain which forms settle the participants.
% mdparagraph:P180:end

% mdparagraph:P181:begin
\textbf{Se bisceop grētte ðone munuc.}
% mdparagraph:P181:end

% mdparagraph:P182:begin
\textbf{Ðone munuc grētte se bisceop.}
% mdparagraph:P182:end

\end{minipage}\par\addvspace{4pt}

\par\begin{minipage}{\linewidth}\raggedright
% mdparagraph:P183:begin
\subsection*{A02}\BookHeading{nav-ch13-learner-0071}{subsubsection}{A02}{0}
% mdparagraph:P183:end

% item: NEW-U13-A02
% mdparagraph:P184:begin
Divide this sentence into clauses; identify each finite verb and its subject. Label every ðā, and translate the whole sentence.
% mdparagraph:P184:end

% mdparagraph:P185:begin
\textbf{Ðā se munuc cōm, ðā grētte se bisceop ðā brōðor.}
% mdparagraph:P185:end

\end{minipage}\par\addvspace{4pt}

\par\begin{minipage}{\linewidth}\raggedright
% mdparagraph:P186:begin
\subsection*{A03}\BookHeading{nav-ch13-learner-0072}{subsubsection}{A03}{0}
% mdparagraph:P186:end

% item: NEW-U13-A03
% mdparagraph:P187:begin
Translate. Identify the main finite verb and the relative-clause finite verb; give the relative antecedent and its role in each clause.
% mdparagraph:P187:end

% mdparagraph:P188:begin
\textbf{Se bisceop, ðe se munuc grētte, sæt on ðǣm hūse.}
% mdparagraph:P188:end

\end{minipage}\par\addvspace{4pt}

\par\begin{minipage}{\linewidth}\raggedright
% mdparagraph:P189:begin
\subsection*{A04}\BookHeading{nav-ch13-learner-0073}{subsubsection}{A04}{0}
% mdparagraph:P189:end

% item: NEW-U13-A04
% mdparagraph:P190:begin
Translate. Identify the person greeted, the person who gave the food, and the recipient. Explain the understood subject and the pronoun cases.
% mdparagraph:P190:end

% mdparagraph:P191:begin
\textbf{Se munuc grētte ðone bisceop and sealde him ðone męte.}
% mdparagraph:P191:end

\end{minipage}\par\addvspace{4pt}

\par\begin{minipage}{\linewidth}\raggedright
% mdparagraph:P192:begin
\subsection*{A05}\BookHeading{nav-ch13-learner-0074}{subsubsection}{A05}{0}
% mdparagraph:P192:end

% item: NEW-U13-A05
% mdparagraph:P193:begin
Translate the second sentence after each context. State the referent of hē in each; explain what morphology establishes and what the context establishes.
% mdparagraph:P193:end

% mdparagraph:P194:begin
\textbf{Se sang wæs swēte. Hē wæs wynsum.}
% mdparagraph:P194:end

% mdparagraph:P195:begin
\textbf{Se bisceop cōm. Hē wæs on ðǣm hūse.}
% mdparagraph:P195:end

\end{minipage}\par\addvspace{4pt}

\par\begin{minipage}{\linewidth}\raggedright
% mdparagraph:P196:begin
\subsection*{A06}\BookHeading{nav-ch13-learner-0075}{subsubsection}{A06}{0}
% mdparagraph:P196:end

% item: NEW-U13-A06
% mdparagraph:P197:begin
Give both possible referents of hē in the second sentence. Is either one selected by gender and number alone? Rewrite that sentence twice, naming each possible subject explicitly.
% mdparagraph:P197:end

% mdparagraph:P198:begin
\textbf{Se bisceop geseah ðone munuc. Hē sæt on ðǣm hūse.}
% mdparagraph:P198:end

\end{minipage}\par\addvspace{4pt}

\par\begin{minipage}{\linewidth}\raggedright
% mdparagraph:P199:begin
\subsection*{A07}\BookHeading{nav-ch13-learner-0076}{subsubsection}{A07}{0}
% mdparagraph:P199:end

% item: NEW-U13-A07
% mdparagraph:P200:begin
The monk addresses the bishop. Translate, and identify each speaker/\allowbreak{}addressee reference in the quotation.
% mdparagraph:P200:end

% mdparagraph:P201:begin
\textbf{Se munuc cwæð: “Ic hīerde ðone sang. Hīerdest ðū hine?”}
% mdparagraph:P201:end

\end{minipage}\par\addvspace{4pt}

\par\begin{minipage}{\linewidth}\raggedright
% mdparagraph:P202:begin
\subsection*{A08}\BookHeading{nav-ch13-learner-0077}{subsubsection}{A08}{0}
% mdparagraph:P202:end

% item: NEW-U13-A08
% mdparagraph:P203:begin
Translate. Identify the content-clause boundary, its finite verb, and the antecedent of hē. The monk is reporting his own presence.
% mdparagraph:P203:end

% mdparagraph:P204:begin
\textbf{Se munuc sægde ðæt hē on ðǣm hūse wǣre.}
% mdparagraph:P204:end

\end{minipage}\par\addvspace{4pt}

\par\begin{minipage}{\linewidth}\raggedright
% mdparagraph:P205:begin
\subsection*{A09}\BookHeading{nav-ch13-learner-0078}{subsubsection}{A09}{0}
% mdparagraph:P205:end

% item: NEW-U13-A09
% mdparagraph:P206:begin
Translate, restoring the omitted repeated predicate in English. Then write the second clause in Old English with that predicate expressed and agreeing with its own subject.
% mdparagraph:P206:end

% mdparagraph:P207:begin
\textbf{Se bisceop wæs on ðǣre cirican, and ðā brōðor on ðǣm hūse.}
% mdparagraph:P207:end

\end{minipage}\par\addvspace{4pt}

\par\begin{minipage}{\linewidth}\raggedright
% mdparagraph:P208:begin
\subsection*{A10}\BookHeading{nav-ch13-learner-0079}{subsubsection}{A10}{0}
% mdparagraph:P208:end

% item: NEW-U13-A10
% mdparagraph:P209:begin
Translate both sentences. Explain which clause each initial phrase belongs to and identify the finite verbs. Are initial words themselves a reliable way to decide who performs an action?
% mdparagraph:P209:end

% mdparagraph:P210:begin
\textbf{On ðǣm hūse sæt se munuc.}
% mdparagraph:P210:end

% mdparagraph:P211:begin
\textbf{Ðā se munuc on ðǣm hūse sæt, ðā cōm se bisceop.}
% mdparagraph:P211:end

\end{minipage}\par\addvspace{4pt}

\par\begin{minipage}{\linewidth}\raggedright
% mdparagraph:P212:begin
\subsection*{A11}\BookHeading{nav-ch13-learner-0080}{subsubsection}{A11}{0}
% mdparagraph:P212:end

% item: NEW-U13-A11
% mdparagraph:P213:begin
Translate and draw the clause structure, using brackets or indentation. Identify the subject of every finite verb and the referent/\allowbreak{}function of ðe and ðā.
% mdparagraph:P213:end

% mdparagraph:P214:begin
\textbf{Ðā se munuc, ðe se bisceop grētte, cōm, ðā sungon ðā brōðor.}
% mdparagraph:P214:end

\end{minipage}\par\addvspace{4pt}

\par\begin{minipage}{\linewidth}\raggedright
% mdparagraph:P215:begin
\subsection*{A12}\BookHeading{nav-ch13-learner-0081}{subsubsection}{A12}{0}
% mdparagraph:P215:end

% item: NEW-U13-A12
% mdparagraph:P216:begin
A learner says: “Every main clause has the finite verb as its second word, and every dependent clause puts it last.” Explain two problems with this statement, referring to the chapter’s checked examples. Then explain why Cook 25.7 can nevertheless require a final verb as a particular exercise target.
% mdparagraph:P216:end

% mdparagraph:P217:begin
The two source fragments and their complete supplied order help are in §2. Keep that help visible for this analysis.
% mdparagraph:P217:end

\end{minipage}\par\addvspace{4pt}

\reserve{7\baselineskip}
% mdparagraph:P218:begin
\subsection*{Composition}\BookHeading{nav-ch13-learner-0082}{subsubsection}{Composition}{0}
% mdparagraph:P218:end

\par\begin{minipage}{\linewidth}\raggedright
% mdparagraph:P219:begin
\subsection*{C01}\BookHeading{nav-ch13-learner-0083}{subsubsection}{C01}{0}
% mdparagraph:P219:end

% item: NEW-U13-C01
% mdparagraph:P220:begin
Translate “The brothers heard the song” twice: first with the subject phrase first, then with the object phrase first. Keep the same participants.
% mdparagraph:P220:end

\end{minipage}\par\addvspace{4pt}

\par\begin{minipage}{\linewidth}\raggedright
% mdparagraph:P221:begin
\subsection*{C02}\BookHeading{nav-ch13-learner-0084}{subsubsection}{C02}{0}
% mdparagraph:P221:end

% item: NEW-U13-C02
% mdparagraph:P222:begin
Translate: Then the bishop greeted the brothers. Begin with ðā and place the finite verb before the nominal subject.
% mdparagraph:P222:end

\end{minipage}\par\addvspace{4pt}

\par\begin{minipage}{\linewidth}\raggedright
% mdparagraph:P223:begin
\subsection*{C03}\BookHeading{nav-ch13-learner-0085}{subsubsection}{C03}{0}
% mdparagraph:P223:end

% item: NEW-U13-C03
% mdparagraph:P224:begin
Translate: When the brothers came, the bishop sang. Use a ðā temporal clause and a main-clause correlative ðā; use a late verb in the dependent clause.
% mdparagraph:P224:end

\end{minipage}\par\addvspace{4pt}

\par\begin{minipage}{\linewidth}\raggedright
% mdparagraph:P225:begin
\subsection*{C04}\BookHeading{nav-ch13-learner-0086}{subsubsection}{C04}{0}
% mdparagraph:P225:end

% item: NEW-U13-C04
% mdparagraph:P226:begin
Translate: The brothers whom the bishop greeted went to the church. Use invariant ðe and keep the relative beside its antecedent.
% mdparagraph:P226:end

\end{minipage}\par\addvspace{4pt}

\par\begin{minipage}{\linewidth}\raggedright
% mdparagraph:P227:begin
\subsection*{C05}\BookHeading{nav-ch13-learner-0087}{subsubsection}{C05}{0}
% mdparagraph:P227:end

% item: NEW-U13-C05
% mdparagraph:P228:begin
Translate as two sentences: The bishop greeted the monk. The monk then went to the church. Repeat the second subject as a noun phrase so that the reference is explicit.
% mdparagraph:P228:end

\end{minipage}\par\addvspace{4pt}

\par\begin{minipage}{\linewidth}\raggedright
% mdparagraph:P229:begin
\subsection*{C06}\BookHeading{nav-ch13-learner-0088}{subsubsection}{C06}{0}
% mdparagraph:P229:end

% item: NEW-U13-C06
% mdparagraph:P230:begin
The monk addresses several brothers. Translate: The monk said, “I want to speak with you.” Preserve the person and number appropriate to direct speech.
% mdparagraph:P230:end

\end{minipage}\par\addvspace{4pt}

\par\begin{minipage}{\linewidth}\raggedright
% mdparagraph:P231:begin
\subsection*{C07}\BookHeading{nav-ch13-learner-0089}{subsubsection}{C07}{0}
% mdparagraph:P231:end

% item: NEW-U13-C07
% mdparagraph:P232:begin
Translate: The monk saw the bishop and greeted him. Express the subject once; both actions have that subject.
% mdparagraph:P232:end

\end{minipage}\par\addvspace{4pt}

\par\begin{minipage}{\linewidth}\raggedright
% mdparagraph:P233:begin
\subsection*{C08}\BookHeading{nav-ch13-learner-0090}{subsubsection}{C08}{0}
% mdparagraph:P233:end

% item: NEW-U13-C08
% mdparagraph:P234:begin
Translate: When the bishop whom the monk greeted came, the brothers sang. Use ðe, temporal ðā and a correlative main-clause ðā. Mark the three clauses after writing the sentence.
% mdparagraph:P234:end

\end{minipage}\par\addvspace{4pt}

\reserve{7\baselineskip}
% mdparagraph:P235:begin
\section*{10. Cook VII.22–25}\BookHeading{nav-ch13-learner-0091}{subsection}{10. Cook VII.22–25}{1}
% mdparagraph:P235:end

% mdparagraph:P236:begin
Translate each English sentence into Old English. The wording, item numbers and original grammatical controls are retained. Parenthesized English and Old English belong to Cook’s prompt. Keep the local helps visible; supplied forms and constructions are not independent recall targets. Original section references are explained in the editorial notes.
% mdparagraph:P236:end

% teaching: NEW-U13-COOK-RELATIVE-CONTROLS
% mdparagraph:P237:begin
Cook’s 87.b requests an inflected demonstrative followed by ðe. Its gender and number follow the antecedent, while its case follows its role inside the relative clause; use this strategy in 24.9, 25.1 and 25.6. Cook’s 87.d requests invariant ðe alone; use it in 25.7. These are the relative strategies taught in Unit 8. A different grammatical strategy can preserve the meaning while missing the stated exercise target.
% mdparagraph:P237:end

\reserve{7\baselineskip}
% mdparagraph:P238:begin
\subsection*{Cook VII.22}\BookHeading{nav-ch13-learner-0092}{subsubsection}{Cook VII.22}{0}
% mdparagraph:P238:end

% mdparagraph:P239:begin
Items 1–15. Use the ordinary senses and verb constructions in §8, including the explicit new senses of familiar words.
% mdparagraph:P239:end

\par\begin{minipage}{\linewidth}\raggedright
% mdparagraph:P240:begin
\subsection*{Cook VII.22.1}\BookHeading{nav-ch13-learner-0093}{subsubsection}{Cook VII.22.1}{0}
% mdparagraph:P240:end

% item: EX-VII-22-01
% mdparagraph:P241:begin
Idleness is the mother of sin.
% mdparagraph:P241:end

\end{minipage}\par\addvspace{4pt}

\par\begin{minipage}{\linewidth}\raggedright
% mdparagraph:P242:begin
\subsection*{Cook VII.22.2}\BookHeading{nav-ch13-learner-0094}{subsubsection}{Cook VII.22.2}{0}
% mdparagraph:P242:end

% item: EX-VII-22-02
% mdparagraph:P243:begin
Etheldred was (om. the) daughter of the king of the East Anglians.
% mdparagraph:P243:end

\end{minipage}\par\addvspace{4pt}

% teaching: U13-COOK-HELP-NAMES
% support: U13-COOK-HELP-NAMES
% mdparagraph:P244:begin
\textbf{Supplied help.} Supply Æðelðrȳð as the nominative personal name and Eastęngla as the genitive plural “of the East Anglians.” No productive proper-name system is assigned.
% mdparagraph:P244:end

% teaching: NEW-U13-VOC-AETHELTHRYTH
% teaching: NEW-U13-VOC-EASTENGLE
% teaching: NEW-U13-LOCAL-HELPS
% teaching: NEW-U13-GRAM-U13-COOK-LOCAL-NAMES
\par\begin{minipage}{\linewidth}\raggedright
% mdparagraph:P245:begin
\subsection*{Cook VII.22.3}\BookHeading{nav-ch13-learner-0095}{subsubsection}{Cook VII.22.3}{0}
% mdparagraph:P245:end

% item: EX-VII-22-03
% mdparagraph:P246:begin
The bishop’s beautiful monastery was destroyed (fordilgian).
% mdparagraph:P246:end

\end{minipage}\par\addvspace{4pt}

\par\begin{minipage}{\linewidth}\raggedright
% mdparagraph:P247:begin
\subsection*{Cook VII.22.4}\BookHeading{nav-ch13-learner-0096}{subsubsection}{Cook VII.22.4}{0}
% mdparagraph:P247:end

% item: EX-VII-22-04
% mdparagraph:P248:begin
Labor [toil, sing.] with your hands.
% mdparagraph:P248:end

\end{minipage}\par\addvspace{4pt}

\par\begin{minipage}{\linewidth}\raggedright
% mdparagraph:P249:begin
\subsection*{Cook VII.22.5}\BookHeading{nav-ch13-learner-0097}{subsubsection}{Cook VII.22.5}{0}
% mdparagraph:P249:end

% item: EX-VII-22-05
% mdparagraph:P250:begin
Pursue (sing.) your study.
% mdparagraph:P250:end

\end{minipage}\par\addvspace{4pt}

\par\begin{minipage}{\linewidth}\raggedright
% mdparagraph:P251:begin
\subsection*{Cook VII.22.6}\BookHeading{nav-ch13-learner-0098}{subsubsection}{Cook VII.22.6}{0}
% mdparagraph:P251:end

% item: EX-VII-22-06
% mdparagraph:P252:begin
He goes to church.
% mdparagraph:P252:end

\end{minipage}\par\addvspace{4pt}

\par\begin{minipage}{\linewidth}\raggedright
% mdparagraph:P253:begin
\subsection*{Cook VII.22.7}\BookHeading{nav-ch13-learner-0099}{subsubsection}{Cook VII.22.7}{0}
% mdparagraph:P253:end

% item: EX-VII-22-07
% mdparagraph:P254:begin
Depart from me, ye that work unrighteousness.
% mdparagraph:P254:end

\end{minipage}\par\addvspace{4pt}

\par\begin{minipage}{\linewidth}\raggedright
% mdparagraph:P255:begin
\subsection*{Cook VII.22.8}\BookHeading{nav-ch13-learner-0100}{subsubsection}{Cook VII.22.8}{0}
% mdparagraph:P255:end

% item: EX-VII-22-08
% mdparagraph:P256:begin
What is your name?
% mdparagraph:P256:end

\end{minipage}\par\addvspace{4pt}

\par\begin{minipage}{\linewidth}\raggedright
% mdparagraph:P257:begin
\subsection*{Cook VII.22.9}\BookHeading{nav-ch13-learner-0101}{subsubsection}{Cook VII.22.9}{0}
% mdparagraph:P257:end

% item: EX-VII-22-09
% mdparagraph:P258:begin
Tell (hātan) our three friends to go thither to him.
% mdparagraph:P258:end

\end{minipage}\par\addvspace{4pt}

\par\begin{minipage}{\linewidth}\raggedright
% mdparagraph:P259:begin
\subsection*{Cook VII.22.10}\BookHeading{nav-ch13-learner-0102}{subsubsection}{Cook VII.22.10}{0}
% mdparagraph:P259:end

% item: EX-VII-22-10
% mdparagraph:P260:begin
The song descends from heaven (plur.) to earth.
% mdparagraph:P260:end

\end{minipage}\par\addvspace{4pt}

\par\begin{minipage}{\linewidth}\raggedright
% mdparagraph:P261:begin
\subsection*{Cook VII.22.11}\BookHeading{nav-ch13-learner-0103}{subsubsection}{Cook VII.22.11}{0}
% mdparagraph:P261:end

% item: EX-VII-22-11
% mdparagraph:P262:begin
I am wont to visit the bishop.
% mdparagraph:P262:end

\end{minipage}\par\addvspace{4pt}

\par\begin{minipage}{\linewidth}\raggedright
% mdparagraph:P263:begin
\subsection*{Cook VII.22.12}\BookHeading{nav-ch13-learner-0104}{subsubsection}{Cook VII.22.12}{0}
% mdparagraph:P263:end

% item: EX-VII-22-12
% mdparagraph:P264:begin
I wish (willan) first (om. to) ask you (sing.) one thing (156).
% mdparagraph:P264:end

\end{minipage}\par\addvspace{4pt}

\par\begin{minipage}{\linewidth}\raggedright
% mdparagraph:P265:begin
\subsection*{Cook VII.22.13}\BookHeading{nav-ch13-learner-0105}{subsubsection}{Cook VII.22.13}{0}
% mdparagraph:P265:end

% item: EX-VII-22-13
% mdparagraph:P266:begin
I wish to take (lǣdan) him with me.
% mdparagraph:P266:end

\end{minipage}\par\addvspace{4pt}

\par\begin{minipage}{\linewidth}\raggedright
% mdparagraph:P267:begin
\subsection*{Cook VII.22.14}\BookHeading{nav-ch13-learner-0106}{subsubsection}{Cook VII.22.14}{0}
% mdparagraph:P267:end

% item: EX-VII-22-14
% mdparagraph:P268:begin
The priest prostrated himself before (beforan, 172) the bishop.
% mdparagraph:P268:end

\end{minipage}\par\addvspace{4pt}

\par\begin{minipage}{\linewidth}\raggedright
% mdparagraph:P269:begin
\subsection*{Cook VII.22.15}\BookHeading{nav-ch13-learner-0107}{subsubsection}{Cook VII.22.15}{0}
% mdparagraph:P269:end

% item: EX-VII-22-15
% mdparagraph:P270:begin
On the third day the holy soul was released from the body.
% mdparagraph:P270:end

\end{minipage}\par\addvspace{4pt}

\reserve{7\baselineskip}
% mdparagraph:P271:begin
\subsection*{Cook VII.23}\BookHeading{nav-ch13-learner-0108}{subsubsection}{Cook VII.23}{0}
% mdparagraph:P271:end

% mdparagraph:P272:begin
Items 1–13. The supplied placement, taking and wonder constructions are local to their items.
% mdparagraph:P272:end

\par\begin{minipage}{\linewidth}\raggedright
% mdparagraph:P273:begin
\subsection*{Cook VII.23.1}\BookHeading{nav-ch13-learner-0109}{subsubsection}{Cook VII.23.1}{0}
% mdparagraph:P273:end

% item: EX-VII-23-01
% mdparagraph:P274:begin
The largest window of the church is two feet (154. c) broad, and four feet high.
% mdparagraph:P274:end

\end{minipage}\par\addvspace{4pt}

\par\begin{minipage}{\linewidth}\raggedright
% mdparagraph:P275:begin
\subsection*{Cook VII.23.2}\BookHeading{nav-ch13-learner-0110}{subsubsection}{Cook VII.23.2}{0}
% mdparagraph:P275:end

% item: EX-VII-23-02
% mdparagraph:P276:begin
Little by little the sweet song drew nigh.
% mdparagraph:P276:end

\end{minipage}\par\addvspace{4pt}

\par\begin{minipage}{\linewidth}\raggedright
% mdparagraph:P277:begin
\subsection*{Cook VII.23.3}\BookHeading{nav-ch13-learner-0111}{subsubsection}{Cook VII.23.3}{0}
% mdparagraph:P277:end

% item: EX-VII-23-03
% mdparagraph:P278:begin
I entreat you to tell [that you tell] me what is (194. b) the angel’s name.
% mdparagraph:P278:end

\end{minipage}\par\addvspace{4pt}

\par\begin{minipage}{\linewidth}\raggedright
% mdparagraph:P279:begin
\subsection*{Cook VII.23.4}\BookHeading{nav-ch13-learner-0112}{subsubsection}{Cook VII.23.4}{0}
% mdparagraph:P279:end

% item: EX-VII-23-04
% mdparagraph:P280:begin
The soldiers took his raiment, and made (wyrcean) four parts, (om. to) every soldier a part; thus did the soldiers.
% mdparagraph:P280:end

\end{minipage}\par\addvspace{4pt}

\par\begin{minipage}{\linewidth}\raggedright
% mdparagraph:P281:begin
\subsection*{Cook VII.23.5}\BookHeading{nav-ch13-learner-0113}{subsubsection}{Cook VII.23.5}{0}
% mdparagraph:P281:end

% item: EX-VII-23-05
% mdparagraph:P282:begin
The mother of Jesus stood by (wið, 172) the cross.
% mdparagraph:P282:end

\end{minipage}\par\addvspace{4pt}

\par\begin{minipage}{\linewidth}\raggedright
% mdparagraph:P283:begin
\subsection*{Cook VII.23.6}\BookHeading{nav-ch13-learner-0114}{subsubsection}{Cook VII.23.6}{0}
% mdparagraph:P283:end

% item: EX-VII-23-06
% mdparagraph:P284:begin
Many read this writing.
% mdparagraph:P284:end

\end{minipage}\par\addvspace{4pt}

\par\begin{minipage}{\linewidth}\raggedright
% mdparagraph:P285:begin
\subsection*{Cook VII.23.7}\BookHeading{nav-ch13-learner-0115}{subsubsection}{Cook VII.23.7}{0}
% mdparagraph:P285:end

% item: EX-VII-23-07
% mdparagraph:P286:begin
With his little hatchet he cut down (forceorfan, 134) the tree.
% mdparagraph:P286:end

\end{minipage}\par\addvspace{4pt}

\par\begin{minipage}{\linewidth}\raggedright
% mdparagraph:P287:begin
\subsection*{Cook VII.23.8}\BookHeading{nav-ch13-learner-0116}{subsubsection}{Cook VII.23.8}{0}
% mdparagraph:P287:end

% item: EX-VII-23-08
% mdparagraph:P288:begin
The ax is placed (āsęttan) into (tō) the root (wyrttruma, 53) of the trees.
% mdparagraph:P288:end

\end{minipage}\par\addvspace{4pt}

% teaching: U13-COOK-HELP-GENITIVE-VARIANT
% support: U13-COOK-HELP-GENITIVE-VARIANT
% mdparagraph:P289:begin
\textbf{Supplied help.} Ðǣra is the supplied genitive plural variant corresponding to ordinary ðāra. It agrees with the plural trees.
% mdparagraph:P289:end

% teaching: NEW-U13-GRAM-U13-COOK-LOCAL-GENITIVE-VARIANT
% teaching: U13-COOK-HELP-ROOT
% support: U13-COOK-HELP-ROOT
% mdparagraph:P290:begin
\textbf{Supplied help.} For the local placement phrase, use tō with dative wyrttruman and a genitive plural trees phrase. The phrase denotes placement at/\allowbreak{}to the root; it does not create a general static tō rule.
% mdparagraph:P290:end

% teaching: NEW-U13-VOC-TO
% teaching: NEW-U13-PAR-TO
% teaching: NEW-U13-GRAM-u13-to-at-root-support
% teaching: NEW-U13-GRAM-U13-COOK-LOCAL-ROOT
\par\begin{minipage}{\linewidth}\raggedright
% mdparagraph:P291:begin
\subsection*{Cook VII.23.9}\BookHeading{nav-ch13-learner-0117}{subsubsection}{Cook VII.23.9}{0}
% mdparagraph:P291:end

% item: EX-VII-23-09
% mdparagraph:P292:begin
He took an ax in his hand.
% mdparagraph:P292:end

\end{minipage}\par\addvspace{4pt}

% teaching: U13-COOK-HELP-TAKE-IN-HAND
% support: U13-COOK-HELP-TAKE-IN-HAND
% mdparagraph:P293:begin
\textbf{Supplied help.} For this taking action, the bounded phrase on his hand means “into his hand,” with accusative hand. This is a declared composition applying Cook’s qualified directional use of on and the ordinary sense of niman; it is not an attested Old English quotation or a general productive niman-plus-goal rule. A static dative “in his hand” is a defensible alternative construal.
% mdparagraph:P293:end

% teaching: NEW-U13-GRAM-U13-COOK-LOCAL-TAKE-IN-HAND
\par\begin{minipage}{\linewidth}\raggedright
% mdparagraph:P294:begin
\subsection*{Cook VII.23.10}\BookHeading{nav-ch13-learner-0118}{subsubsection}{Cook VII.23.10}{0}
% mdparagraph:P294:end

% item: EX-VII-23-10
% mdparagraph:P295:begin
Sell (sęllan) me five hatchets.
% mdparagraph:P295:end

\end{minipage}\par\addvspace{4pt}

\par\begin{minipage}{\linewidth}\raggedright
% mdparagraph:P296:begin
\subsection*{Cook VII.23.11}\BookHeading{nav-ch13-learner-0119}{subsubsection}{Cook VII.23.11}{0}
% mdparagraph:P296:end

% item: EX-VII-23-11
% mdparagraph:P297:begin
His bodily infirmity was very grievous.
% mdparagraph:P297:end

\end{minipage}\par\addvspace{4pt}

\par\begin{minipage}{\linewidth}\raggedright
% mdparagraph:P298:begin
\subsection*{Cook VII.23.12}\BookHeading{nav-ch13-learner-0120}{subsubsection}{Cook VII.23.12}{0}
% mdparagraph:P298:end

% item: EX-VII-23-12
% mdparagraph:P299:begin
Your companionship is very pleasant.
% mdparagraph:P299:end

\end{minipage}\par\addvspace{4pt}

\par\begin{minipage}{\linewidth}\raggedright
% mdparagraph:P300:begin
\subsection*{Cook VII.23.13}\BookHeading{nav-ch13-learner-0121}{subsubsection}{Cook VII.23.13}{0}
% mdparagraph:P300:end

% item: EX-VII-23-13
% mdparagraph:P301:begin
Is it any wonder that [though] the faith of God grew (196. e)?
% mdparagraph:P301:end

\end{minipage}\par\addvspace{4pt}

% teaching: U13-COOK-HELP-WONDER
% support: U13-COOK-HELP-WONDER
% mdparagraph:P302:begin
\textbf{Supplied help.} In the local wonder question, hit anticipates the following circumstance. Wundor is a nominal predicate; ðēah introduces the conceded circumstance with a preterite subjunctive. An admitted fact is compatible with this subjunctive.
% mdparagraph:P302:end

% teaching: NEW-U13-GRAM-U13-COOK-LOCAL-WONDER
\reserve{7\baselineskip}
% mdparagraph:P303:begin
\subsection*{Cook VII.24}\BookHeading{nav-ch13-learner-0122}{subsubsection}{Cook VII.24}{0}
% mdparagraph:P303:end

% mdparagraph:P304:begin
Items 1–12. The relative controls and the supports beside the longer sentences remain part of the task.
% mdparagraph:P304:end

\par\begin{minipage}{\linewidth}\raggedright
% mdparagraph:P305:begin
\subsection*{Cook VII.24.1}\BookHeading{nav-ch13-learner-0123}{subsubsection}{Cook VII.24.1}{0}
% mdparagraph:P305:end

% item: EX-VII-24-01
% mdparagraph:P306:begin
He that (sē ðe) sitteth on the throne shall dwell (wunian) in the midst of them.
% mdparagraph:P306:end

\end{minipage}\par\addvspace{4pt}

% teaching: U13-COOK-HELP-RELATIVE-LONG
% support: U13-COOK-HELP-RELATIVE-LONG
% mdparagraph:P307:begin
\textbf{Supplied help.} Sē is the retained long source spelling of the nominative masculine singular demonstrative/\allowbreak{}relative. The ordinary short se remains a valid taught alternative; the added quantity is not an independent recall target.
% mdparagraph:P307:end

% teaching: NEW-U13-GRAM-U13-COOK-LOCAL-RELATIVE-LONG
\par\begin{minipage}{\linewidth}\raggedright
% mdparagraph:P308:begin
\subsection*{Cook VII.24.2}\BookHeading{nav-ch13-learner-0124}{subsubsection}{Cook VII.24.2}{0}
% mdparagraph:P308:end

% item: EX-VII-24-02
% mdparagraph:P309:begin
Angels are spirits without body.
% mdparagraph:P309:end

\end{minipage}\par\addvspace{4pt}

\par\begin{minipage}{\linewidth}\raggedright
% mdparagraph:P310:begin
\subsection*{Cook VII.24.3}\BookHeading{nav-ch13-learner-0125}{subsubsection}{Cook VII.24.3}{0}
% mdparagraph:P310:end

% item: EX-VII-24-03
% mdparagraph:P311:begin
I saw a star (169) fall from heaven unto [until] the earth.
% mdparagraph:P311:end

\end{minipage}\par\addvspace{4pt}

% teaching: U13-COOK-HELP-FALL
% support: U13-COOK-HELP-FALL
% mdparagraph:P312:begin
\textbf{Supplied help.} Feallan is the supplied plain infinitive “fall.” Use the already taught gesēon + accusative actor + plain infinitive construction.
% mdparagraph:P312:end

% teaching: NEW-U13-GRAM-U13-COOK-LOCAL-FALL
\par\begin{minipage}{\linewidth}\raggedright
% mdparagraph:P313:begin
\subsection*{Cook VII.24.4}\BookHeading{nav-ch13-learner-0126}{subsubsection}{Cook VII.24.4}{0}
% mdparagraph:P313:end

% item: EX-VII-24-04
% mdparagraph:P314:begin
He that overcometh, the same shall be clothed (gescrȳdan) in white (hwīt) raiment.
% mdparagraph:P314:end

% support: U13-COOK-HELP-RELATIVE-LONG
% mdparagraph:P315:begin
\textbf{Supplied help.} Sē is the retained long source spelling of the nominative masculine singular demonstrative/\allowbreak{}relative. The ordinary short se remains a valid taught alternative; the added quantity is not an independent recall target.
% mdparagraph:P315:end

\end{minipage}\par\addvspace{4pt}

\par\begin{minipage}{\linewidth}\raggedright
% mdparagraph:P316:begin
\subsection*{Cook VII.24.5}\BookHeading{nav-ch13-learner-0127}{subsubsection}{Cook VII.24.5}{0}
% mdparagraph:P316:end

% item: EX-VII-24-05
% mdparagraph:P317:begin
These things saith he that is holy.
% mdparagraph:P317:end

% support: U13-COOK-HELP-RELATIVE-LONG
% mdparagraph:P318:begin
\textbf{Supplied help.} Sē is the retained long source spelling of the nominative masculine singular demonstrative/\allowbreak{}relative. The ordinary short se remains a valid taught alternative; the added quantity is not an independent recall target.
% mdparagraph:P318:end

\end{minipage}\par\addvspace{4pt}

\par\begin{minipage}{\linewidth}\raggedright
% mdparagraph:P319:begin
\subsection*{Cook VII.24.6}\BookHeading{nav-ch13-learner-0128}{subsubsection}{Cook VII.24.6}{0}
% mdparagraph:P319:end

% item: EX-VII-24-06
% mdparagraph:P320:begin
Fear him (ðone) that is able to destroy (fordōn) both soul and body in (on, 166) hell.
% mdparagraph:P320:end

\end{minipage}\par\addvspace{4pt}

\par\begin{minipage}{\linewidth}\raggedright
% mdparagraph:P321:begin
\subsection*{Cook VII.24.7}\BookHeading{nav-ch13-learner-0129}{subsubsection}{Cook VII.24.7}{0}
% mdparagraph:P321:end

% item: EX-VII-24-07
% mdparagraph:P322:begin
I am not worthy (wierðe) that thou enter under (172) my roof.
% mdparagraph:P322:end

\end{minipage}\par\addvspace{4pt}

% teaching: U13-COOK-HELP-WORTHY
% support: U13-COOK-HELP-WORTHY
% mdparagraph:P323:begin
\textbf{Supplied help.} Uncontracted ne eom is available for “I am not.” The bounded worthy-content construction uses a present subjunctive; under in the entry phrase takes an accusative roof. No negative-copula contraction is required.
% mdparagraph:P323:end

% teaching: NEW-U13-GRAM-U13-COOK-LOCAL-WORTHY
\par\begin{minipage}{\linewidth}\raggedright
% mdparagraph:P324:begin
\subsection*{Cook VII.24.8}\BookHeading{nav-ch13-learner-0130}{subsubsection}{Cook VII.24.8}{0}
% mdparagraph:P324:end

% item: EX-VII-24-08
% mdparagraph:P325:begin
Great is your reward in heaven (plur.).
% mdparagraph:P325:end

\end{minipage}\par\addvspace{4pt}

\par\begin{minipage}{\linewidth}\raggedright
% mdparagraph:P326:begin
\subsection*{Cook VII.24.9}\BookHeading{nav-ch13-learner-0131}{subsubsection}{Cook VII.24.9}{0}
% mdparagraph:P326:end

% item: EX-VII-24-09
% mdparagraph:P327:begin
If God so clothes (om. the) grass of the field, which (87. b) to-day is, and is (bið) to-morrow sent (āsęndan) into (on, 172) the oven, much (178) more he (om. shall) clothe you, O ye of little faith.
% mdparagraph:P327:end

\end{minipage}\par\addvspace{4pt}

% teaching: U13-COOK-HELP-MUCH-MORE
% support: U13-COOK-HELP-MUCH-MORE
% mdparagraph:P328:begin
\textbf{Supplied help.} Supply micle, instrumental neuter singular used adverbially before mā, for “much more.” Recognition of that cell elsewhere does not establish independent productive recall here.
% mdparagraph:P328:end

% teaching: NEW-U13-GRAM-U13-COOK-LOCAL-MUCH-MORE
% teaching: U13-COOK-HELP-TOMORROW
% support: U13-COOK-HELP-TOMORROW
% mdparagraph:P329:begin
\textbf{Supplied help.} Tō-morgen means “tomorrow” and is supplied as one invariant temporal expression.
% mdparagraph:P329:end

% teaching: NEW-U13-VOC-TO-MORGEN
% teaching: NEW-U13-PAR-TO-MORGEN
% teaching: NEW-U13-GRAM-U13-COOK-LOCAL-TOMORROW
% teaching: U13-COOK-HELP-LITTLE-FAITH
% support: U13-COOK-HELP-LITTLE-FAITH
% mdparagraph:P330:begin
\textbf{Supplied help.} The address “you of little faith” is locally supplied as gē lȳtles gelēafan. Lȳtles gelēafan is the bounded genitive quality expression; no general productive rule is inferred.
% mdparagraph:P330:end

% teaching: NEW-U13-GRAM-U13-COOK-LOCAL-LITTLE-FAITH
\par\begin{minipage}{\linewidth}\raggedright
% mdparagraph:P331:begin
\subsection*{Cook VII.24.10}\BookHeading{nav-ch13-learner-0132}{subsubsection}{Cook VII.24.10}{0}
% mdparagraph:P331:end

% item: EX-VII-24-10
% mdparagraph:P332:begin
Better is the soul than food (męte), and the body than its raiment.
% mdparagraph:P332:end

\end{minipage}\par\addvspace{4pt}

\par\begin{minipage}{\linewidth}\raggedright
% mdparagraph:P333:begin
\subsection*{Cook VII.24.11}\BookHeading{nav-ch13-learner-0133}{subsubsection}{Cook VII.24.11}{0}
% mdparagraph:P333:end

% item: EX-VII-24-11
% mdparagraph:P334:begin
Behold (behealdan) the birds of heaven; your heavenly Father feedeth (fēdan) them.
% mdparagraph:P334:end

\end{minipage}\par\addvspace{4pt}

\par\begin{minipage}{\linewidth}\raggedright
% mdparagraph:P335:begin
\subsection*{Cook VII.24.12}\BookHeading{nav-ch13-learner-0134}{subsubsection}{Cook VII.24.12}{0}
% mdparagraph:P335:end

% item: EX-VII-24-12
% mdparagraph:P336:begin
What shall (sculon) we eat, or what drink?
% mdparagraph:P336:end

\end{minipage}\par\addvspace{4pt}

\reserve{7\baselineskip}
% mdparagraph:P337:begin
\subsection*{Cook VII.25}\BookHeading{nav-ch13-learner-0135}{subsubsection}{Cook VII.25}{0}
% mdparagraph:P337:end

% mdparagraph:P338:begin
Items 1–12. In 25.7, “verb last” applies to the relative translating “that which I command you”: place its verb bebēode last in that relative. The surrounding conditional clause need not be forced into a universal dependent-clause pattern.
% mdparagraph:P338:end

\par\begin{minipage}{\linewidth}\raggedright
% mdparagraph:P339:begin
\subsection*{Cook VII.25.1}\BookHeading{nav-ch13-learner-0136}{subsubsection}{Cook VII.25.1}{0}
% mdparagraph:P339:end

% item: EX-VII-25-01
% mdparagraph:P340:begin
He that (sē ðe) loveth not his brother, whom (87. b) he seeth, how can (mugan) he love God, whom he seeth not?
% mdparagraph:P340:end

% support: U13-COOK-HELP-RELATIVE-LONG
% mdparagraph:P341:begin
\textbf{Supplied help.} Sē is the retained long source spelling of the nominative masculine singular demonstrative/\allowbreak{}relative. The ordinary short se remains a valid taught alternative; the added quantity is not an independent recall target.
% mdparagraph:P341:end

\end{minipage}\par\addvspace{4pt}

\par\begin{minipage}{\linewidth}\raggedright
% mdparagraph:P342:begin
\subsection*{Cook VII.25.2}\BookHeading{nav-ch13-learner-0137}{subsubsection}{Cook VII.25.2}{0}
% mdparagraph:P342:end

% item: EX-VII-25-02
% mdparagraph:P343:begin
The time of my departure is nigh at hand.
% mdparagraph:P343:end

\end{minipage}\par\addvspace{4pt}

\par\begin{minipage}{\linewidth}\raggedright
% mdparagraph:P344:begin
\subsection*{Cook VII.25.3}\BookHeading{nav-ch13-learner-0138}{subsubsection}{Cook VII.25.3}{0}
% mdparagraph:P344:end

% item: EX-VII-25-03
% mdparagraph:P345:begin
He did something (hwæthwega) of this sort, while his companions were going to town (burg).
% mdparagraph:P345:end

\end{minipage}\par\addvspace{4pt}

% teaching: U13-COOK-HELP-INDEFINITE
% support: U13-COOK-HELP-INDEFINITE
% mdparagraph:P346:begin
\textbf{Supplied help.} Hwæthwega means “something” and takes the supplied partitive expression swilces, “of this kind.” The reading’s different exact hwæthwugu spelling remains separately recorded.
% mdparagraph:P346:end

% teaching: NEW-U13-VOC-HWAETHWEGA
% teaching: NEW-U13-PAR-HWAETHWEGA
% teaching: NEW-U13-GRAM-U13-COOK-LOCAL-INDEFINITE
\par\begin{minipage}{\linewidth}\raggedright
% mdparagraph:P347:begin
\subsection*{Cook VII.25.4}\BookHeading{nav-ch13-learner-0139}{subsubsection}{Cook VII.25.4}{0}
% mdparagraph:P347:end

% item: EX-VII-25-04
% mdparagraph:P348:begin
Be with the brethren.
% mdparagraph:P348:end

\end{minipage}\par\addvspace{4pt}

\par\begin{minipage}{\linewidth}\raggedright
% mdparagraph:P349:begin
\subsection*{Cook VII.25.5}\BookHeading{nav-ch13-learner-0140}{subsubsection}{Cook VII.25.5}{0}
% mdparagraph:P349:end

% item: EX-VII-25-05
% mdparagraph:P350:begin
Open the windows of the house.
% mdparagraph:P350:end

\end{minipage}\par\addvspace{4pt}

\par\begin{minipage}{\linewidth}\raggedright
% mdparagraph:P351:begin
\subsection*{Cook VII.25.6}\BookHeading{nav-ch13-learner-0141}{subsubsection}{Cook VII.25.6}{0}
% mdparagraph:P351:end

% item: EX-VII-25-06
% mdparagraph:P352:begin
Love (plur.) your enemies, do (om. to) them well that hate you, that ye (om. may) be (196. f) children of your Father, which (87. b) is in heaven (plur.).
% mdparagraph:P352:end

% support: U13-COOK-HELP-RELATIVE-LONG
% mdparagraph:P353:begin
\textbf{Supplied help.} Sē is the retained long source spelling of the nominative masculine singular demonstrative/\allowbreak{}relative. The ordinary short se remains a valid taught alternative; the added quantity is not an independent recall target.
% mdparagraph:P353:end

\end{minipage}\par\addvspace{4pt}

\par\begin{minipage}{\linewidth}\raggedright
% mdparagraph:P354:begin
\subsection*{Cook VII.25.7}\BookHeading{nav-ch13-learner-0142}{subsubsection}{Cook VII.25.7}{0}
% mdparagraph:P354:end

% item: EX-VII-25-07
% mdparagraph:P355:begin
Ye are my friends, if ye do (ind.) the things which (87. d) I command (bebēodan) you (verb last).
% mdparagraph:P355:end

\end{minipage}\par\addvspace{4pt}

\par\begin{minipage}{\linewidth}\raggedright
% mdparagraph:P356:begin
\subsection*{Cook VII.25.8}\BookHeading{nav-ch13-learner-0143}{subsubsection}{Cook VII.25.8}{0}
% mdparagraph:P356:end

% item: EX-VII-25-08
% mdparagraph:P357:begin
The birds sang the sweetest song, winsome (acc.) to hear.
% mdparagraph:P357:end

\end{minipage}\par\addvspace{4pt}

\par\begin{minipage}{\linewidth}\raggedright
% mdparagraph:P358:begin
\subsection*{Cook VII.25.9}\BookHeading{nav-ch13-learner-0144}{subsubsection}{Cook VII.25.9}{0}
% mdparagraph:P358:end

% item: EX-VII-25-09
% mdparagraph:P359:begin
The roof of the church is very high.
% mdparagraph:P359:end

\end{minipage}\par\addvspace{4pt}

\par\begin{minipage}{\linewidth}\raggedright
% mdparagraph:P360:begin
\subsection*{Cook VII.25.10}\BookHeading{nav-ch13-learner-0145}{subsubsection}{Cook VII.25.10}{0}
% mdparagraph:P360:end

% item: EX-VII-25-10
% mdparagraph:P361:begin
There (ðā) was suddenly with the angel (om. a) multitude [number] of the heavenly host, praising (gen. plur.; cf. 189. 2) God, and thus saying.
% mdparagraph:P361:end

\end{minipage}\par\addvspace{4pt}

% teaching: U13-COOK-HELP-HEAVENLY-HOST
% support: U13-COOK-HELP-HEAVENLY-HOST
% mdparagraph:P362:begin
\textbf{Supplied help.} Unordered local material: cweðendra, “of people saying”; getæl, “number/\allowbreak{}company”; hęriendra, “of people praising”; heofonlices weorodes, “of a heavenly host”; sęmninga, “suddenly.” The genitive plural participles are substantival and supported in this collective-number formula. Identify their functions with this help; do not independently derive a participle paradigm.
% mdparagraph:P362:end

% teaching: NEW-U13-VOC-GETAEL
% teaching: NEW-U13-VOC-WEOROD
% teaching: NEW-U13-VOC-SEMNINGA
% teaching: NEW-U13-VOC-HERIAN
% teaching: NEW-U13-VOC-CWETHAN
% teaching: NEW-U13-PAR-SEMNINGA
% teaching: NEW-U13-PAR-HERIAN
% teaching: NEW-U13-PAR-CWETHAN
% teaching: NEW-U13-GRAM-u13-bounded-formula-support
% teaching: NEW-U13-GRAM-U13-COOK-LOCAL-HEAVENLY-HOST
\par\begin{minipage}{\linewidth}\raggedright
% mdparagraph:P363:begin
\subsection*{Cook VII.25.11}\BookHeading{nav-ch13-learner-0146}{subsubsection}{Cook VII.25.11}{0}
% mdparagraph:P363:end

% item: EX-VII-25-11
% mdparagraph:P364:begin
To God be (193) glory (wuldor) in (om. the) heights, and on earth peace (om. to) men of good will.
% mdparagraph:P364:end

\end{minipage}\par\addvspace{4pt}

% teaching: U13-COOK-HELP-DOXOLOGY
% support: U13-COOK-HELP-DOXOLOGY
% mdparagraph:P365:begin
\textbf{Supplied help.} Unordered local material: Gode, dative recipient; gōdes willan, the bounded quality expression “of good will”; hēanessum, dative plural “heights”; sibb, nominative “peace”; sīe, optative “be”; wuldor, nominative “glory.” Supply the complete doxological construction: dative recipient with nominative glory, and an understood parallel wish of peace for the men described. This formula does not establish unrestricted genitive-quality syntax.
% mdparagraph:P365:end

% teaching: NEW-U13-VOC-HEANES
% teaching: NEW-U13-VOC-WULDOR
% teaching: NEW-U13-VOC-SIBB
% teaching: NEW-U13-VOC-WILLA
% teaching: NEW-U13-GRAM-U13-COOK-LOCAL-DOXOLOGY
\par\begin{minipage}{\linewidth}\raggedright
% mdparagraph:P366:begin
\subsection*{Cook VII.25.12}\BookHeading{nav-ch13-learner-0147}{subsubsection}{Cook VII.25.12}{0}
% mdparagraph:P366:end

% item: EX-VII-25-12
% mdparagraph:P367:begin
He returned (ęfthweorfan) (om. by) the same way that he had come (see ǣr, adv.).
% mdparagraph:P367:end

\end{minipage}\par\addvspace{4pt}

\reserve{7\baselineskip}
% mdparagraph:P368:begin
\section*{11. The bishop and the monk}\BookHeading{nav-ch13-learner-0148}{subsection}{11. The bishop and the monk}{1}
% mdparagraph:P368:end

% mdparagraph:P369:begin
This connected account is new teaching material. It prepares the clause patterns and references of the source reading without claiming to retell the historical episode.
% mdparagraph:P369:end

% passage: NEW-U13-READING
% mdparagraph:P370:begin
\par\begin{minipage}{\linewidth}\raggedright
\readingtext{Se bisceop sæt on ðǣm hūse, and se munuc wæs ðǣr. Ðā se munuc ðone sang hīerde, ðā ēode hē tō ðǣm bisceope. Se munuc cwæð: “Ic hīerde ðone sang, ðe swīðe wynsum wæs. Hīerdest ðū hine?” Se bisceop cwæð: “Gā nū, and lǣd ðā brōðor hider.” Ðā ēode se munuc tō ðǣre cirican, and cōm ęft mid ðǣm brōðrum. Se bisceop geseah hīe and cwæð: “Gē sindon mīne frīend. Ic wille mid ēow sprecan.” Se munuc sæt mid ðǣm brōðrum, and hīe hīerdon ðæs bisceopes word.}
\end{minipage}\par
% mdparagraph:P370:end

\par\begin{minipage}{\linewidth}\raggedright
% mdparagraph:P371:begin
\subsection*{R01}\BookHeading{nav-ch13-learner-0149}{subsubsection}{R01}{0}
% mdparagraph:P371:end

% item: NEW-U13-R01
% mdparagraph:P372:begin
Translate the complete new account, The bishop and the monk. Preserve the inquiry, command and final speech.
% mdparagraph:P372:end

\end{minipage}\par\addvspace{4pt}

\par\begin{minipage}{\linewidth}\raggedright
% mdparagraph:P373:begin
\subsection*{R02}\BookHeading{nav-ch13-learner-0150}{subsubsection}{R02}{0}
% mdparagraph:P373:end

% item: NEW-U13-R02
% mdparagraph:P374:begin
Identify the referents of hē after ēode, hine in the monk’s question, hīe after geseah, and Ic and ēow in the bishop’s final speech. Explain why hine need not refer to a man.
% mdparagraph:P374:end

\end{minipage}\par\addvspace{4pt}

\par\begin{minipage}{\linewidth}\raggedright
% mdparagraph:P375:begin
\subsection*{R03}\BookHeading{nav-ch13-learner-0151}{subsubsection}{R03}{0}
% mdparagraph:P375:end

% item: NEW-U13-R03
% mdparagraph:P376:begin
Find the temporal clause and its main clause, the relative clause, and a coordinated finite verb with an unexpressed repeated subject. Give each finite verb and recover each subject.
% mdparagraph:P376:end

\end{minipage}\par\addvspace{4pt}

\par\begin{minipage}{\linewidth}\raggedright
% mdparagraph:P377:begin
\subsection*{R04}\BookHeading{nav-ch13-learner-0152}{subsubsection}{R04}{0}
% mdparagraph:P377:end

% item: NEW-U13-R04
% mdparagraph:P378:begin
Give the narrative order of hearing the song, the monk’s inquiry, the bishop’s command, the monk’s return and the bishop’s final speech. Who is present at the end, and which words establish this?
% mdparagraph:P378:end

\end{minipage}\par\addvspace{4pt}

\reserve{7\baselineskip}
% mdparagraph:P379:begin
\section*{12. The Passing of Chad: the music and summons}\BookHeading{nav-ch13-learner-0153}{subsection}{12. The Passing of Chad: the music and summons}{1}
% mdparagraph:P379:end

% teaching: NEW-U13-SOURCE-READING
\reserve{7\baselineskip}
% mdparagraph:P380:begin
\subsection*{English context}\BookHeading{nav-ch13-learner-0154}{subsubsection}{English context}{0}
% mdparagraph:P380:end

% mdparagraph:P381:begin
Owini had left a court household to live and work at Chad’s monastery. While the other brothers studied inside, he did needed work outside. Chad was alone in the local oratory; the companions had gone to the community church.
% mdparagraph:P381:end

% mdparagraph:P382:begin
This orientation is new teaching material based on the preceding passage. The two buildings must be distinguished; see original note 152.5. The extract ends with Chad’s summons to the brothers and Owini.
% mdparagraph:P382:end

\reserve{7\baselineskip}
% mdparagraph:P383:begin
\subsection*{Old English text}\BookHeading{nav-ch13-learner-0155}{subsubsection}{Old English text}{0}
% mdparagraph:P383:end

% mdparagraph:P384:begin
Cook, printed pages 152–153, editorial lines 152.3–153.9, ending partway through the last line. The complete original notes and their attachment points follow in §13. Numbers in the local helps below identify explanations, not extra words in the source.
% mdparagraph:P384:end

\par\begin{minipage}{\linewidth}\raggedright
% passage: FB-VII-CHAD-MUSIC-SUMMONS
% mdparagraph:P385:begin
\readingtext{Ðā hē ðā sume dæge hwæthwugu swilces ūte dyde, and his gefēran tō byrig tō ciricean ēodon, swā hīe gelōmlice dydon, and se bisceop, āna in ðǣre ciricean, oððe in bēcrǣdinge oððe in gebedum geornfull wæs, ðā gehīerde hē sęmninga, swā-swā hē ęft æfter ðon sægde, ðā swētestan stefne and ðā fægerestan, singendra and blissiendra, of heofonum oð eorðan āstīgan. Ðā stefne and ðone sang hē cwæð ðæt hē ǣrest gehīerde fram ēastsūðdǣle heofones, ðæt is fram hēanesse ðǣre winterlican sunnan ūpganges; and ðanan tō him styccemǣlum nēalǣcton, oð-ðæt hē becōm tō ðæcean ðǣre ciricean ðe se bisceop in wæs; and, ingangende, eall gefylde, and in ymbhwyrfte ymbsealde. And hē ðā geornlice his mōd āðęnede in ðā ðing ðe hē gehīerde. Ðā gehīerde hē ęft swā-swā healfre tīde fæce, of hrōfe ðǣre ilcan ciricean ūpāstīgan ðone ilcan blissesang, and, ðȳ ilcan wege ðe hē ǣr cōm, ūp oð heofonas mid unasęcgendre swētnesse ęfthweorfan. Ðā wunode hē ðǣr sum fæc tīde, wundriende and wafiende; and mid behygdigum mōde ðōhte and smēade hwæt ðā ðing bēon sceolden. Ðā ontȳnde se bisceop ðæt ēagðȳrel ðǣre ciricean, and mid his handa slōg tācen, swā-swā his gewuna wæs gif hwilc mann ūte wǣre, ðæt hē in tō him ēode. Ðā ēode hē sōna in tō him. Cwæð hē, se bisceop, him tō: “Gang hraðe tō ciricean, and hāt ūre seofon brōðor hider tō mē cuman; and ðū ēac swilce mid wes.”}
% mdparagraph:P385:end

\end{minipage}\par
\reserve{7\baselineskip}
% mdparagraph:P386:begin
\subsection*{Local reading help}\BookHeading{nav-ch13-learner-0156}{subsubsection}{Local reading help}{0}
% mdparagraph:P386:end

% teaching: NEW-U13-GRAM-u13-source-local-idioms
% mdparagraph:P387:begin
All eighteen helps remain available for translation and analysis. They supply exact source forms and bounded constructions; they do not add full paradigms for every word in the extract.
% mdparagraph:P387:end

% teaching: U13-SRC-HELP-01
% support: U13-SRC-HELP-01
% mdparagraph:P388:begin
\textbf{Reading help 1 — sume dæge; hwæthwugu swilces; ūte.} Sume dæge is a time expression, “one day”; sume is the instrumental form. Hwæthwugu means “something”; swilces is neuter genitive, “of this kind.” Ūte means “outside.”
% mdparagraph:P388:end

% teaching: NEW-U13-VOC-HWAETHWUGU
% teaching: NEW-U13-VOC-UTE
% teaching: NEW-U13-PAR-HWAETHWUGU
% teaching: NEW-U13-PAR-UTE
% teaching: U13-SRC-HELP-02
% support: U13-SRC-HELP-02
% mdparagraph:P389:begin
\textbf{Reading help 2 — ciricean.} Ciricean is the source genitive or dative spelling corresponding to productive cirican: genitive when dependent on roof or window, dative after tō or locative in. The same noun can name either the community church or Chad’s oratory; context and Cook’s note distinguish the buildings.
% mdparagraph:P389:end

% teaching: U13-SRC-HELP-03
% support: U13-SRC-HELP-03
% mdparagraph:P390:begin
\textbf{Reading help 3 — gelōmlice; āna; in bēcrǣdinge /\allowbreak{} in gebedum geornfull.} Gelōmlice means “often”; āna, “alone.” Bēcrǣdinge means “reading books”; gebedum is dative plural, “prayers”; geornfull here means “occupied” or “intent.” In governs the dative in these expressions.
% mdparagraph:P390:end

% teaching: NEW-U13-VOC-BECRAEDING
% teaching: NEW-U13-VOC-GEBED
% teaching: NEW-U13-VOC-GEORNFULL
% teaching: NEW-U13-VOC-GELOMLICE
% teaching: NEW-U13-VOC-IN
% teaching: NEW-U13-VOC-AN
% teaching: NEW-U13-PAR-GELOMLICE
% teaching: NEW-U13-PAR-IN
% teaching: NEW-U13-PAR-AN
% teaching: U13-SRC-HELP-04
% support: U13-SRC-HELP-04
% mdparagraph:P391:begin
\textbf{Reading help 4 — gehīerde; sęmninga; swā-swā; ęft æfter ðon.} Gehīerde is prefixed “heard.” Sęmninga means “suddenly”; swā-swā, “just as”; ęft, “again” or “afterwards” as the context requires. Æfter ðon means “after that”; ðon is the source dative demonstrative variant.
% mdparagraph:P391:end

% teaching: NEW-U13-VOC-SWA-SWA
% teaching: NEW-U13-PAR-SWA-SWA
% teaching: U13-SRC-HELP-05
% support: U13-SRC-HELP-05
% mdparagraph:P392:begin
\textbf{Reading help 5 — swētestan /\allowbreak{} fægerestan stefne; singendra and blissiendra; oð; āstīgan.} Stefne is accusative feminine singular, “voice.” Fægerestan means “most beautiful,” agreeing with stefne. Singendra and blissiendra are genitive plural participles used substantivally, “of people singing and rejoicing.” Oð means “as far as.” Āstīgan here means “descend.”
% mdparagraph:P392:end

% teaching: NEW-U13-VOC-STEFN
% teaching: NEW-U13-VOC-FAEGER-SUPERL
% teaching: NEW-U13-VOC-SINGAN
% teaching: NEW-U13-VOC-BLISSIAN
% teaching: NEW-U13-PAR-SINGAN
% teaching: NEW-U13-PAR-BLISSIAN
% teaching: U13-SRC-HELP-06
% support: U13-SRC-HELP-06
% mdparagraph:P393:begin
\textbf{Reading help 6 — Ðā stefne and ðone sang … gehīerde.} The initial accusative objects belong with the later gehīerde inside the reported clause. Cwæð introduces the report. Ǣrest means “first.” Do not turn the fronted accusatives into subjects.
% mdparagraph:P393:end

% teaching: U13-SRC-HELP-07
% support: U13-SRC-HELP-07
% mdparagraph:P394:begin
\textbf{Reading help 7 — ēastsūðdǣle; hēanesse; winterlican sunnan ūpganges; ðanan.} Ēastsūðdǣle means “southeastern part”; hēanesse, “height”; winterlican, “winter,” agrees with genitive sunnan; ūpganges means “rising.” Ðanan means “from there.”
% mdparagraph:P394:end

% teaching: NEW-U13-VOC-WINTERLIC
% teaching: NEW-U13-VOC-UPGANG
% teaching: NEW-U13-VOC-THANAN
% teaching: NEW-U13-PAR-THANAN
% teaching: U13-SRC-HELP-08
% support: U13-SRC-HELP-08
% mdparagraph:P395:begin
\textbf{Reading help 8 — styccemǣlum; nēalǣcton; oð-ðæt hē becōm.} Styccemǣlum means “little by little.” Nēalǣcton is third plural: the voice and song approached. Oð-ðæt means “until.” Becōm means “arrived”; the following singular reference is explained by Cook’s note 9a.
% mdparagraph:P395:end

% teaching: U13-SRC-HELP-09
% support: U13-SRC-HELP-09
% mdparagraph:P396:begin
\textbf{Reading help 9 — ðæcean … ðe … in.} Ðæcean is dative singular, “roof.” The relative ðe is governed by the later in, “in which”; the preposition follows its relative.
% mdparagraph:P396:end

% teaching: NEW-U13-VOC-THAECE
% teaching: U13-SRC-HELP-10
% support: U13-SRC-HELP-10
% mdparagraph:P397:begin
\textbf{Reading help 10 — ingangende; eall; in ymbhwyrfte ymbsealde.} Ingangende is a present participle, “entering.” Eall is accusative neuter singular, “everything.” Gefylde means “filled.” In ymbhwyrfte means “all around”; ymbsealde means “surrounded.” These continue the description of the song.
% mdparagraph:P397:end

% teaching: NEW-U13-VOC-YMBHWYRFT
% teaching: NEW-U13-VOC-INGANGAN
% teaching: NEW-U13-VOC-YMBSELLAN
% teaching: NEW-U13-PAR-INGANGAN
% teaching: NEW-U13-PAR-YMBSELLAN
% teaching: U13-SRC-HELP-11
% support: U13-SRC-HELP-11
% mdparagraph:P398:begin
\textbf{Reading help 11 — geornlice; his mōd āðęnede in ….} Geornlice means “attentively.” Mōd is “mind”; āðęnede is “directed/\allowbreak{}stretched.” The whole attention expression means to direct one’s mind or attention to something; in has the accusative here.
% mdparagraph:P398:end

% teaching: NEW-U13-VOC-MOD
% teaching: NEW-U13-VOC-GEORNLICE
% teaching: NEW-U13-VOC-ATHENIAN
% teaching: NEW-U13-PAR-GEORNLICE
% teaching: NEW-U13-PAR-ATHENIAN
% teaching: U13-SRC-HELP-12
% support: U13-SRC-HELP-12
% mdparagraph:P399:begin
\textbf{Reading help 12 — ęft swā-swā healfre tīde fæce.} Here swā-swā means “about,” not “just as.” Healfre tīde is genitive, “of half an hour”; fæce is dative of the interval after which the event occurs. The hour sense of tīde is supplied separately.
% mdparagraph:P399:end

% teaching: NEW-U13-VOC-HEALF
% teaching: NEW-U13-VOC-FAEC
% teaching: NEW-U13-VOC-TID
% teaching: U13-SRC-HELP-13
% support: U13-SRC-HELP-13
% mdparagraph:P400:begin
\textbf{Reading help 13 — ūpāstīgan; blissesang; ðȳ ilcan wege; unasęcgendre swētnesse; ęfthweorfan.} Ūpāstīgan means “rise up”; blissesang, “song of joy.” Ðȳ ilcan wege is the supported route expression, “by the same way”; ǣr means “previously.” Unasęcgendre swētnesse means “with indescribable sweetness” after mid. Ęfthweorfan means “return.” The two infinitives depend on heard.
% mdparagraph:P400:end

% teaching: NEW-U13-VOC-BLISSESANG
% teaching: NEW-U13-VOC-UNASECGENDE
% teaching: NEW-U13-VOC-SWETNES
% teaching: NEW-U13-VOC-UPASTIGAN
% teaching: NEW-U13-PAR-UPASTIGAN
% teaching: U13-SRC-HELP-14
% support: U13-SRC-HELP-14
% mdparagraph:P401:begin
\textbf{Reading help 14 — wunode; sum fæc tīde; wundriende and wafiende.} Wunode means “remained”; ðǣr, “there.” Sum fæc tīde is an accusative duration expression, “for a while.” Wundriende and wafiende are uninflected participles describing the observer as wondering and amazed.
% mdparagraph:P401:end

% teaching: NEW-U13-VOC-WUNDRIAN
% teaching: NEW-U13-VOC-WAFIAN
% teaching: NEW-U13-PAR-WUNDRIAN
% teaching: NEW-U13-PAR-WAFIAN
% teaching: U13-SRC-HELP-15
% support: U13-SRC-HELP-15
% mdparagraph:P402:begin
\textbf{Reading help 15 — mid behygdigum mōde; ðōhte and smēade; hwæt … bēon sceolden.} Behygdigum mōde means “with an attentive mind.” Ðōhte and smēade mean “thought” and “pondered.” The embedded question asks what the things might be or mean; sceolden is preterite subjunctive plural, with infinitive bēon.
% mdparagraph:P402:end

% teaching: NEW-U13-VOC-BEHYGDIG
% teaching: NEW-U13-VOC-THENCAN
% teaching: NEW-U13-VOC-SMEAGAN
% teaching: NEW-U13-VOC-SCULAN
% teaching: NEW-U13-PAR-THENCAN
% teaching: NEW-U13-PAR-SMEAGAN
% teaching: NEW-U13-PAR-SCULAN
% teaching: U13-SRC-HELP-16
% support: U13-SRC-HELP-16
% mdparagraph:P403:begin
\textbf{Reading help 16 — ontȳnde; ēagðȳrel; handa; slōg tācen; swā-swā; gewuna.} Ontȳnde means “opened.” Handa is the source dative form of hand. Slōg tācen means “made a signal,” literally “struck a sign.” Swā-swā means “just as.” Gewuna means “custom” or “habit.”
% mdparagraph:P403:end

% teaching: NEW-U13-VOC-GEWUNA
% teaching: NEW-U13-VOC-ONTYNAN
% teaching: NEW-U13-VOC-SLEAN
% teaching: NEW-U13-PAR-ONTYNAN
% teaching: NEW-U13-VOC-HAND
% teaching: NEW-U13-PAR-SLEAN
% teaching: U13-SRC-HELP-17
% support: U13-SRC-HELP-17
% mdparagraph:P404:begin
\textbf{Reading help 17 — gif hwilc mann … wǣre; ðæt … ēode; in; sōna.} Hwilc mann here means “anyone.” Wǣre is preterite subjunctive in the habitual condition. The ðæt clause states the content or purpose of the signal. In is the adverb “in”; sōna means “immediately.”
% mdparagraph:P404:end

% teaching: U13-SRC-HELP-18
% support: U13-SRC-HELP-18
% mdparagraph:P405:begin
\textbf{Reading help 18 — Cwæð … him tō; Gang; hāt; ēac swilce mid wes.} The separated him … tō identifies the person addressed. Gang means “go” (singular imperative); hraðe, “quickly.” Hāt is singular “command/\allowbreak{}tell,” with an accusative subject and infinitive. Hider means “hither/\allowbreak{}here.” Ēac swilce means “also”; mid is adverbial “with them”; wes is singular “be.”
% mdparagraph:P405:end

% teaching: NEW-U13-VOC-GAN
% teaching: NEW-U13-VOC-HRATHE
% teaching: NEW-U13-PAR-GAN
% teaching: NEW-U13-PAR-HRATHE
\reserve{7\baselineskip}
% mdparagraph:P406:begin
\subsection*{Source-reading tasks}\BookHeading{nav-ch13-learner-0157}{subsubsection}{Source-reading tasks}{0}
% mdparagraph:P406:end

\par\begin{minipage}{\linewidth}\raggedright
% mdparagraph:P407:begin
\subsection*{S01}\BookHeading{nav-ch13-learner-0158}{subsubsection}{S01}{0}
% mdparagraph:P407:end

% item: NEW-U13-S01
% mdparagraph:P408:begin
Translate the complete music-and-summons extract from The Passing of Chad. Use all printed local helps and original notes.
% mdparagraph:P408:end

\end{minipage}\par\addvspace{4pt}

\par\begin{minipage}{\linewidth}\raggedright
% mdparagraph:P409:begin
\subsection*{S02}\BookHeading{nav-ch13-learner-0159}{subsubsection}{S02}{0}
% mdparagraph:P409:end

% item: NEW-U13-S02
% mdparagraph:P410:begin
Divide the opening sentence through of heofonum oð eorðan āstīgan into its major clauses. Identify the subjects of dyde, ēodon, dydon, wæs, gehīerde and the parenthetical sægde; distinguish the first two occurrences of ðā and the main-clause ðā before gehīerde.
% mdparagraph:P410:end

\end{minipage}\par\addvspace{4pt}

\par\begin{minipage}{\linewidth}\raggedright
% mdparagraph:P411:begin
\subsection*{S03}\BookHeading{nav-ch13-learner-0160}{subsubsection}{S03}{0}
% mdparagraph:P411:end

% item: NEW-U13-S03
% mdparagraph:P412:begin
Follow the sound through the account. Identify what hē refers to in hē becōm and in ðe hē ǣr cōm; describe its path using the printed vocabulary help.
% mdparagraph:P412:end

\end{minipage}\par\addvspace{4pt}

\par\begin{minipage}{\linewidth}\raggedright
% mdparagraph:P413:begin
\subsection*{S04}\BookHeading{nav-ch13-learner-0161}{subsubsection}{S04}{0}
% mdparagraph:P413:end

% item: NEW-U13-S04
% mdparagraph:P414:begin
Identify Owini’s actions and Chad’s actions from the arrival of the sound until the command to fetch the brothers. Where does se bisceop make a reference explicit?
% mdparagraph:P414:end

\end{minipage}\par\addvspace{4pt}

\par\begin{minipage}{\linewidth}\raggedright
% mdparagraph:P415:begin
\subsection*{S05}\BookHeading{nav-ch13-learner-0162}{subsubsection}{S05}{0}
% mdparagraph:P415:end

% item: NEW-U13-S05
% mdparagraph:P416:begin
In ðe se bisceop in wæs, give the antecedent and explain the role of in using the local help. Translate the complete noun phrase and relative clause.
% mdparagraph:P416:end

\end{minipage}\par\addvspace{4pt}

\par\begin{minipage}{\linewidth}\raggedright
% mdparagraph:P417:begin
\subsection*{S06}\BookHeading{nav-ch13-learner-0163}{subsubsection}{S06}{0}
% mdparagraph:P417:end

% item: NEW-U13-S06
% mdparagraph:P418:begin
In Ðā stefne and ðone sang hē cwæð ðæt hē ǣrest gehīerde, identify the two fronted objects and the verb that governs them. Use Cook’s note 8 and the local help. Show the reporting clause and its content clause without assuming that the fronted objects belong to cwæð.
% mdparagraph:P418:end

\end{minipage}\par\addvspace{4pt}

\par\begin{minipage}{\linewidth}\raggedright
% mdparagraph:P419:begin
\subsection*{S07}\BookHeading{nav-ch13-learner-0164}{subsubsection}{S07}{0}
% mdparagraph:P419:end

% item: NEW-U13-S07
% mdparagraph:P420:begin
In Chad’s final speech, identify the speaker, the addressee, and the referents of ūre, mē and ðū. Which brothers are to come, and is Owini included? Cite the wording and the original note distinguishing the churches.
% mdparagraph:P420:end

\end{minipage}\par\addvspace{4pt}

\par\begin{minipage}{\linewidth}\raggedright
% mdparagraph:P421:begin
\subsection*{S08}\BookHeading{nav-ch13-learner-0165}{subsubsection}{S08}{0}
% mdparagraph:P421:end

% item: NEW-U13-S08
% mdparagraph:P422:begin
Choose one main clause and one dependent clause in the extract. Identify their finite verbs, subjects and first constituents. Explain how their actual order relates to the chapter’s tendencies; do not rewrite the source into a compulsory pattern.
% mdparagraph:P422:end

\end{minipage}\par\addvspace{4pt}

\reserve{8\baselineskip}
% mdparagraph:P423:begin
\section*{13. Cook’s original notes and attachments}\BookHeading{nav-ch13-learner-0166}{subsection}{13. Cook’s original notes and attachments}{1}
% mdparagraph:P423:end

\begingroup\fontsize{10}{12.8}\selectfont
\setlength{\parskip}{1.5pt}
% mdparagraph:P424:begin
The following notes preserve Cook’s wording and original numbering. References such as 176 are to sections of Cook’s grammar. His questions about later English words remain source apparatus; they are not additional assessed tasks in this unit. The local helps in §12 supply what is needed to read the extract.
% mdparagraph:P424:end

% mdtable:T56:begin
\par\begingroup
\fontsize{10.5}{13.2}\selectfont
\setlength{\LTleft}{0pt}\setlength{\LTright}{\fill}
\begin{longtable}{@{}>{\raggedright\arraybackslash}p{35.74mm}>{\raggedright\arraybackslash}p{43.41mm}>{\raggedright\arraybackslash}p{39.23mm}@{}}
\toprule
\textbf{Original note} & \textbf{Attached word} & \textbf{Editorial line} \\
\midrule
\endfirsthead
\toprule
\textbf{Original note} & \textbf{Attached word} & \textbf{Editorial line} \\
\midrule
\endhead
\midrule
\multicolumn{3}{r}{\footnotesize Continued on the following page.}\\
\endfoot
\bottomrule
\endlastfoot
152.3 & sume & 152.3 \\
152.3 & dæge & 152.3 \\
152.4 & swilces & 152.3 \\
152.5 & ciricean & 152.5 \\
152.6 & stefne & 152.8 \\
152.7 & blissiendra & 152.8 \\
152.8 & stefne & 152.9 \\
152.8 & sang & 152.10 \\
152.9 & styccemǣlum & 152.12 \\
152.9a & hē & 152.13 \\
152.10 & ðæcean & 152.13 \\
152.11 & ðe & 152.13 \\
152.12 & eall & 152.14 \\
152.13 & mōd & 152.15 \\
152.13 & āðęnede & 152.16 \\
152.14 & swā-swā & 152.17 \\
152.15 & fæce & 152.17 \\
152.16 & blissesang & 152.18 \\
152.17 & wege & 152.19 \\
152.18 & unasęcgendre & 152.19 \\
152.19 & ęfthweorfan & 152.20 \\
153.1 & sum & 153.1 \\
153.1 & fæc & 153.1 \\
153.1 & tīde & 153.1 \\
153.2 & ēagðȳrel & 153.4 \\
153.3 & tō & 153.7 \\
153.4 & ciricean & 153.7 \\
153.5 & hāt & 153.8 \\
153.6 & wes & 153.9 \\
\end{longtable}
\par\endgroup

% mdtable:T56:end
\reserve{7\baselineskip}
% mdparagraph:P425:begin
\subsection*{152.3}\BookHeading{nav-ch13-learner-0167}{subsubsection}{152.3}{0}
% mdparagraph:P425:end

% source-note: FB152-N3
% mdparagraph:P426:begin
See 176.
% mdparagraph:P426:end

\reserve{7\baselineskip}
% mdparagraph:P427:begin
\subsection*{152.4}\BookHeading{nav-ch13-learner-0168}{subsubsection}{152.4}{0}
% mdparagraph:P427:end

% source-note: FB152-N4
% mdparagraph:P428:begin
See 154. b.
% mdparagraph:P428:end

\reserve{7\baselineskip}
% mdparagraph:P429:begin
\subsection*{152.5}\BookHeading{nav-ch13-learner-0169}{subsubsection}{152.5}{0}
% mdparagraph:P429:end

% source-note: FB152-N5
% mdparagraph:P430:begin
Lat. oratorio loci. The monks had gone to the church. Cf. below, p. 153, l. 7.
% mdparagraph:P430:end

\reserve{7\baselineskip}
% mdparagraph:P431:begin
\subsection*{152.6}\BookHeading{nav-ch13-learner-0170}{subsubsection}{152.6}{0}
% mdparagraph:P431:end

% source-note: FB152-N6
% mdparagraph:P432:begin
See 169. For this word see Chaucer, Knight’s Tale 1704.
% mdparagraph:P432:end

\reserve{7\baselineskip}
% mdparagraph:P433:begin
\subsection*{152.7}\BookHeading{nav-ch13-learner-0171}{subsubsection}{152.7}{0}
% mdparagraph:P433:end

% source-note: FB152-N7
% mdparagraph:P434:begin
Lat. vocem suavissimam cantantium atque lætantium. What adjective is concealed in bliss- (see 34)?
% mdparagraph:P434:end

\reserve{7\baselineskip}
% mdparagraph:P435:begin
\subsection*{152.8}\BookHeading{nav-ch13-learner-0172}{subsubsection}{152.8}{0}
% mdparagraph:P435:end

% source-note: FB152-N8
% mdparagraph:P436:begin
Acc. after gehīerde.
% mdparagraph:P436:end

\reserve{7\baselineskip}
% mdparagraph:P437:begin
\subsection*{152.9}\BookHeading{nav-ch13-learner-0173}{subsubsection}{152.9}{0}
% mdparagraph:P437:end

% source-note: FB152-N9
% mdparagraph:P438:begin
See 72.
% mdparagraph:P438:end

\reserve{7\baselineskip}
% mdparagraph:P439:begin
\subsection*{152.9a}\BookHeading{nav-ch13-learner-0174}{subsubsection}{152.9a}{0}
% mdparagraph:P439:end

% source-note: FB152-N9a
% mdparagraph:P440:begin
Se sang.
% mdparagraph:P440:end

\reserve{7\baselineskip}
% mdparagraph:P441:begin
\subsection*{152.10}\BookHeading{nav-ch13-learner-0175}{subsubsection}{152.10}{0}
% mdparagraph:P441:end

% source-note: FB152-N10
% mdparagraph:P442:begin
Lat. tectum, for which l. 17 has hrōf.
% mdparagraph:P442:end

\reserve{7\baselineskip}
% mdparagraph:P443:begin
\subsection*{152.11}\BookHeading{nav-ch13-learner-0176}{subsubsection}{152.11}{0}
% mdparagraph:P443:end

% source-note: FB152-N11
% mdparagraph:P444:begin
Governed by in.
% mdparagraph:P444:end

\reserve{7\baselineskip}
% mdparagraph:P445:begin
\subsection*{152.12}\BookHeading{nav-ch13-learner-0177}{subsubsection}{152.12}{0}
% mdparagraph:P445:end

% source-note: FB152-N12
% mdparagraph:P446:begin
Acc. sing.
% mdparagraph:P446:end

\reserve{7\baselineskip}
% mdparagraph:P447:begin
\subsection*{152.13}\BookHeading{nav-ch13-learner-0178}{subsubsection}{152.13}{0}
% mdparagraph:P447:end

% source-note: FB152-N13
% mdparagraph:P448:begin
Lat. animum intenderet.
% mdparagraph:P448:end

\reserve{7\baselineskip}
% mdparagraph:P449:begin
\subsection*{152.14}\BookHeading{nav-ch13-learner-0179}{subsubsection}{152.14}{0}
% mdparagraph:P449:end

% source-note: FB152-N14
% mdparagraph:P450:begin
Lat. quasi.
% mdparagraph:P450:end

\reserve{7\baselineskip}
% mdparagraph:P451:begin
\subsection*{152.15}\BookHeading{nav-ch13-learner-0180}{subsubsection}{152.15}{0}
% mdparagraph:P451:end

% source-note: FB152-N15
% mdparagraph:P452:begin
See 176.
% mdparagraph:P452:end

\reserve{7\baselineskip}
% mdparagraph:P453:begin
\subsection*{152.16}\BookHeading{nav-ch13-learner-0181}{subsubsection}{152.16}{0}
% mdparagraph:P453:end

% source-note: FB152-N16
% mdparagraph:P454:begin
Lat. lætitiæ canticum. See 169.
% mdparagraph:P454:end

\reserve{7\baselineskip}
% mdparagraph:P455:begin
\subsection*{152.17}\BookHeading{nav-ch13-learner-0182}{subsubsection}{152.17}{0}
% mdparagraph:P455:end

% source-note: FB152-N17
% mdparagraph:P456:begin
See 176. 2.
% mdparagraph:P456:end

\reserve{7\baselineskip}
% mdparagraph:P457:begin
\subsection*{152.18}\BookHeading{nav-ch13-learner-0183}{subsubsection}{152.18}{0}
% mdparagraph:P457:end

% source-note: FB152-N18
% mdparagraph:P458:begin
Lat. ineffabili.
% mdparagraph:P458:end

\reserve{7\baselineskip}
% mdparagraph:P459:begin
\subsection*{152.19}\BookHeading{nav-ch13-learner-0184}{subsubsection}{152.19}{0}
% mdparagraph:P459:end

% source-note: FB152-N19
% mdparagraph:P460:begin
Lat. reverti.
% mdparagraph:P460:end

\reserve{7\baselineskip}
% mdparagraph:P461:begin
\subsection*{153.1}\BookHeading{nav-ch13-learner-0185}{subsubsection}{153.1}{0}
% mdparagraph:P461:end

% source-note: FB153-N1
% mdparagraph:P462:begin
Lat. aliquantulum horæ (170).
% mdparagraph:P462:end

\reserve{7\baselineskip}
% mdparagraph:P463:begin
\subsection*{153.2}\BookHeading{nav-ch13-learner-0186}{subsubsection}{153.2}{0}
% mdparagraph:P463:end

% source-note: FB153-N2
% mdparagraph:P464:begin
In what Mod. Eng. word is a disguised form of -ðȳrel to be found? What is the etymology of window?
% mdparagraph:P464:end

\reserve{7\baselineskip}
% mdparagraph:P465:begin
\subsection*{153.3}\BookHeading{nav-ch13-learner-0187}{subsubsection}{153.3}{0}
% mdparagraph:P465:end

% source-note: FB153-N3
% mdparagraph:P466:begin
See 201. 1.
% mdparagraph:P466:end

\reserve{7\baselineskip}
% mdparagraph:P467:begin
\subsection*{153.4}\BookHeading{nav-ch13-learner-0188}{subsubsection}{153.4}{0}
% mdparagraph:P467:end

% source-note: FB153-N4
% mdparagraph:P468:begin
What is the etymology of church? From what language is it originally derived?
% mdparagraph:P468:end

\reserve{7\baselineskip}
% mdparagraph:P469:begin
\subsection*{153.5}\BookHeading{nav-ch13-learner-0189}{subsubsection}{153.5}{0}
% mdparagraph:P469:end

% source-note: FB153-N5
% mdparagraph:P470:begin
See hātan.
% mdparagraph:P470:end

\reserve{7\baselineskip}
% mdparagraph:P471:begin
\subsection*{153.6}\BookHeading{nav-ch13-learner-0190}{subsubsection}{153.6}{0}
% mdparagraph:P471:end

% source-note: FB153-N6
% mdparagraph:P472:begin
Imper. sing. When followed by the adj. hāl, what Mod. Eng. word does it give rise to?
% mdparagraph:P472:end

\reserve{7\baselineskip}
% mdparagraph:P473:begin
\section*{14. Editorial and source notes}\BookHeading{nav-ch13-learner-0191}{subsection}{14. Editorial and source notes}{1}
% mdparagraph:P473:end

\fontsize{10}{12}\selectfont
% mdparagraph:P474:begin
\textbf{Editions and locators.} The source grammar and reading are Albert S. Cook, A First Book in Old English, second edition (1895). The exercise prompts are Cook’s Exercises in Old English, Group VII, Exercises 22–25, printed pages 11–13 (PDF pages 19–21). The Chad extract is on printed pages 152–153 of the First Book (PDF pages 168–169). Editorial line counts start at the first text line on each printed page; Cook’s marginal five-line numerals remain distinct from this counting convention.
% mdparagraph:P474:end

% mdparagraph:P475:begin
\textbf{The two order fragments.} Æfter ðon hē underðēodde is exact source wording from printed page 153, editorial line 16, outside the continuous extract; Cook’s note 13 glosses the last verb with Latin subjunxit. The dependent-clause fragment is already part of the assigned extract on page 152. Its English brackets supply the objects identified by Cook’s original note 8. The general treatment of order follows the qualified tendencies described by R. D. Fulk, An Introductory Grammar of Old English, §78. Neither source supports an exceptionless finite-second or dependent-final rule.
% mdparagraph:P475:end

% mdparagraph:P476:begin
\textbf{Forms and derivation.} The noun and adjective models follow Cook §§43–63; verbal systems follow the strong, weak and anomalous paradigms on printed pages 54–80 and the lexical entries for the individual verbs. The full kinship systems are also checked against Joseph Wright, Old English Grammar, §415. Wright §259.3 supplies the consonant simplification in swētra; the selected rule-derived genitive plural is not presented as an independently attested quotation. Forceorfan’s selected present forcierf- is a declared Cook-style normalization of dictionary cirf/\allowbreak{}cyrf variants. The movement, proximity and clothing constructions are checked against the corresponding Bosworth–Toller entries, including beforan, wið, under and gescrȳdan.
% mdparagraph:P476:end

% mdparagraph:P477:begin
\textbf{Source typography.} The continuous reading joins source line and page breaks, removes purely typographic word breaks, and represents original note cues through the attachment table. No word has been added to complete its sense. Original quantities, e/\allowbreak{}ę, ð and source variants such as ciricean, hwæthwugu and Gang are retained. Dictionary separators in joined learner headwords, including the compound nouns and adjectives, do not change their letters or vowel quantities. Wright’s thorn is represented by Cook-style eth in the derived kinship paradigms. These are declared presentation choices, not silent changes to the source reading.
% mdparagraph:P477:end

% mdparagraph:P478:begin
\textbf{C03: Cook 23.7.} The original prompt’s reference “forceorfan, 134” is retained. Section 134 of the selected First Book does not supply the forceorfan paradigm required here. Use the complete indicative table and the printed principal-part support in this unit; the original reference is not a prerequisite. No cause is assigned to the discrepancy.
% mdparagraph:P478:end

% mdparagraph:P479:begin
\textbf{Exercise controls and variants.} Cook’s directions choose targets within the language. The long sē and genitive plural ðǣra are supplied where needed; ordinary se and ðāra remain the familiar forms. The English “in white raiment” is rendered through the taught mid-with-dative clothing construction. Static dative interpretations of “before” or “in his hand” can be defensible even where a task selects a different construction. Such alternatives must be judged separately from compliance with the selected target.
% mdparagraph:P479:end

\endgroup

% book-source-body:end
\par\label{end-ch13-learner}
\clearpage\endgroup

\BookRecto
\begingroup
\renewcommand{\runningtitle}{Unit 14. Word formation and lexical families}
\hyphenpenalty=10000
\exhyphenpenalty=10000
\emergencystretch=3em
\newcommand{\itempara}[2]{\par\addvspace{3pt}\noindent\hangindent=1.8em\hangafter=1\makebox[1.8em][l]{#1.}#2\par}
\newcommand{\reserve}[1]{\par\penalty-100\begingroup\dimen0=\pagegoal\advance\dimen0 by-\pagetotal\ifdim\dimen0<#1\newpage\fi\endgroup}
% book-source-body:begin

\addfontfeatures{Ligatures=TeXOff}
\raggedright
\booktitle{\BookComponent{ch14-learner}{Unit 14. Word formation and lexical families}UNIT 14}{Word formation and lexical families}

% mdparagraph:P001:begin
A familiar part can help you understand a word. It does not settle the whole word’s meaning. This unit distinguishes compounds and derived words from their inflected forms, then applies that distinction to the complete Cook exercises and a connected argument about the dangers of high station.
% mdparagraph:P001:end

% mdparagraph:P002:begin
Sections 1–5 supply the forms, meanings and uses needed by the practice. The four translated examples and Exercises F, A, C and S are new teaching material. Cook VIII.26–27 retain their original prompts. The reading preserves the selected Old English; its modern orientation and helps are separate additions.
% mdparagraph:P002:end

% mdparagraph:P003:begin
\textbf{Scope and notation.} Productive forms may be used in the senses and constructions taught here or earlier. A table marked “plural reference” or “singular reference” supplies that number for recognition only. Infinitive and principal-part citations identify a verb; they do not by themselves assign a new syntactic use. Supplied help stays visible during its task and is not tested by unaided recall. In verbal tables, 1, 2 and 3 denote person. Preserve vowel length and the distinction between e and ę.
% mdparagraph:P003:end

\reserve{7\baselineskip}
% mdparagraph:P004:begin
\section*{1. Formation and inflection}\BookHeading{nav-ch14-learner-0001}{subsection}{1. Formation and inflection}{1}
% mdparagraph:P004:end

% teaching: NEW-U14-GRAM-U14-DERIVATION-VERSUS-INFLECTION
% teaching: NEW-U14-GRAM-U14-COMPOUND-ANALYSIS
% mdparagraph:P005:begin
\textbf{Derivation}\BookPoint{bk-ch14-learner-7f05ca79b8a2} forms a lexical word; \textbf{inflection} gives that word a form required in a sentence. Rihtwīs, “righteous”, and rihtwīsnes, “righteousness”, are different words and different parts of speech. The -e of rihtwīsnesse instead marks an inflected form of the noun. An ending can serve more than one case: the clause, agreement and governing word determine which one is present.
% mdparagraph:P005:end

% mdparagraph:P006:begin
A\BookPoint{bk-ch14-learner-79faadb7257c} \textbf{compound} joins lexical components. In the selected noun compounds the final noun is the head: it identifies the kind of thing named and supplies the inflectional pattern. Goldhord combines gold and the supplied citation hord, “treasure; hoard”. The whole noun means “treasure; hoard of gold”; the ending is added after the complete stem goldhord-. A phrase containing an adjective and a noun instead inflects the adjective for agreement as well as the noun for case and number.
% mdparagraph:P006:end

% mdparagraph:P007:begin
The\BookPoint{bk-ch14-learner-719f663802e5} first component can specify material, source or a quality. Heofonfȳre has heofon + fȳr; hēaclifum has hēah + the supplied clif, “cliff”, with hēa- as the combining form. The local helps give these source forms’ meanings and cases. Do not turn that particular hēah/\allowbreak{}hēa- relationship into a rule for dropping h elsewhere. Learn the attested whole-word sense: gōdwębb means “purple cloth”, and wyrtgemang “spice; mixture of spices”. English component substitution will not reliably recover such meanings.\textsuperscript{1}
% mdparagraph:P007:end

% teaching: NEW-U14-GRAM-U14-DERIVED-NOUNS
\reserve{15\baselineskip}
% mdparagraph:P008:begin
A\BookPoint{bk-ch14-learner-70727fa3d47b} suffix follows the family stem. The selected noun suffixes include -ung for an activity, -scipe for a state or relationship, -nes for an abstract quality, and -dōm for a state, position or authority. Agent nouns in -end name a doer. These are aids to analysing attested words, not instructions to invent a noun from any base. The derived noun has its own gender, meaning and declension.
% mdparagraph:P008:end

% mdtable:T01:begin
\par\begingroup
\fontsize{10.5}{13.2}\selectfont
\setlength{\LTleft}{0pt}\setlength{\LTright}{\fill}
\noindent
\begin{tabular}{@{}>{\raggedright\arraybackslash}p{22.05mm}>{\raggedright\arraybackslash}p{52.92mm}>{\raggedright\arraybackslash}p{43.41mm}@{}}
\toprule
\textbf{Family} & \textbf{Parts of speech and learned meanings} & \textbf{Formation} \\
\midrule
leornian → leornung & verb “learn” → feminine noun “learning; study” & leorn- + -ung \\
gefēra → gefērscipe & masculine noun “companion” → masculine noun “companionship” & gefēr- + -scipe; the citation -a is not retained \\
rihtwīs → rihtwīsnes & adjective “righteous” → feminine noun “righteousness” & rihtwīs- + -nes \\
īdel → īdelnes & adjective “empty; void” → feminine noun “idleness” & īdel- + -nes; learn the noun’s established meaning \\
\bottomrule
\end{tabular}
\par\endgroup

% mdtable:T01:end
% mdparagraph:P009:begin
Leornian\BookPoint{bk-ch14-learner-2f64e0f5ae82} is supplied here as a dictionary citation of the earlier verb. Īdelnes demonstrates the limit of a family resemblance: the previously learned noun means “idleness”, not automatically “physical emptiness”. In nouns ending in -nes, s is doubled before endings: -nesse, -nessa, -nessum. That stem spelling belongs to the noun’s inflection.
% mdparagraph:P009:end

% teaching: NEW-U14-GRAM-U14-DERIVED-ADJECTIVES
% teaching: NEW-U14-GRAM-U14-ADVERB-FAMILIES
\reserve{12\baselineskip}
% mdparagraph:P010:begin
The\BookPoint{bk-ch14-learner-92d5683c753b} adjective-forming -lic gives heofonlic, “heavenly”, from heofon, “heaven”. The adjective then takes the ordinary strong or weak endings. In an inflected -lic adjective, separate the suffix from the following agreement ending. The source adjective eorðlican is analysed with its supplied case, number and gender below.
% mdparagraph:P010:end

% mdtable:T02:begin
\par\begingroup
\fontsize{10.5}{13.2}\selectfont
\setlength{\LTleft}{0pt}\setlength{\LTright}{\fill}
\noindent
\begin{tabular}{@{}>{\raggedright\arraybackslash}p{41.26mm}>{\raggedright\arraybackslash}p{38.56mm}>{\raggedright\arraybackslash}p{38.56mm}@{}}
\toprule
\textbf{Citation base} & \textbf{Adverb} & \textbf{Meaning} \\
\midrule
swīð, “strong” & swīðe & very \\
sōð, “true” & sōðlice & truly \\
\bottomrule
\end{tabular}
\par\endgroup

% mdtable:T02:end
% mdparagraph:P011:begin
Swīð\BookPoint{bk-ch14-learner-85c8a677bf69} and sōð are supplied adjective citations, not new full adjective paradigms. The attested adverbs use -e or -lice and do not agree with a noun. A final e is not enough to identify an adverb: it also occurs in many noun and adjective forms. Use the whole word and its clause function.\textsuperscript{2}
% mdparagraph:P011:end

% teaching: NEW-U14-GRAM-U14-PREFIX-LEXICAL-LIMITS
\reserve{10\baselineskip}
% mdparagraph:P012:begin
A\BookPoint{bk-ch14-learner-dce2679d3490} prefix precedes the family stem. In rihtwīsnes /\allowbreak{} unrihtwīsnes, “righteousness /\allowbreak{} unrighteousness”, un- marks opposition. Compare these supplied verb citations with the familiar finite verbs:
% mdparagraph:P012:end

% mdtable:T03:begin
\par\begingroup
\fontsize{10.5}{13.2}\selectfont
\setlength{\LTleft}{0pt}\setlength{\LTright}{\fill}
\noindent
\begin{tabular}{@{}>{\raggedright\arraybackslash}p{20.69mm}>{\raggedright\arraybackslash}p{37.43mm}>{\raggedright\arraybackslash}p{60.26mm}@{}}
\toprule
\textbf{Base} & \textbf{Prefixed word} & \textbf{Learned relation} \\
\midrule
berstan, “burst” & tōberstan, “break up; go to pieces” & tō- has an asunder sense here \\
weorðan, “become” & forweorðan, “perish” & the prefixed whole word has a destructive-result meaning \\
\bottomrule
\end{tabular}
\par\endgroup

% mdtable:T03:end
% mdparagraph:P013:begin
Each\BookPoint{bk-ch14-learner-22f05dfaf32e} whole verb keeps its own meaning, stem changes and government. The prefix for- does not simply equal English “for”, and tō- does not always mean the same thing. Likewise, ge- can belong to a verb’s lexical form or occur in a participle; neither its addition nor its removal is automatic. Use the selected forms below.
% mdparagraph:P013:end

% teaching: NEW-U14-GRAM-U14-NOMINALIZED-EXPRESSIONS
% mdparagraph:P014:begin
An\BookPoint{bk-ch14-learner-ecac515ca5c6} adjective can stand for people without an expressed noun: the earlier ðā rihtwīsan means “the righteous people”. It remains an adjective, with weak inflection after the article and case determined by the clause. The new neuter noun gōd, used in the plural for “goods; property”, has its own lexical entry and noun paradigm. Endings can coincide across these categories; the absence of a following noun does not decide which category a word belongs to.\textsuperscript{3}
% mdparagraph:P014:end

\reserve{7\baselineskip}
% mdparagraph:P015:begin
\section*{2. Nouns}\BookHeading{nav-ch14-learner-0002}{subsection}{2. Nouns}{1}
% mdparagraph:P015:end

% mdparagraph:P016:begin
The following tables give every singular and plural case. Gender and the ordinary meaning are stated with each word. The singular/\allowbreak{}plural reference labels limit lexical use; they do not omit cells. Otherwise both numbers are productive. The strong nouns here use genitive plural -a and dative plural -um; the weak nouns use genitive plural -ena and dative plural -um.\textsuperscript{4}
% mdparagraph:P016:end

\reserve{7\baselineskip}
% mdparagraph:P017:begin
\subsection*{Strong masculine nouns}\BookHeading{nav-ch14-learner-0003}{subsubsection}{Strong masculine nouns}{0}
% mdparagraph:P017:end

% mdparagraph:P018:begin
The nominative and accusative singular have no ending. Add -es for genitive singular and -e for dative singular; nominative and accusative plural take -as. The other plural endings are -a and -um.
% mdparagraph:P018:end

% teaching: NEW-U14-VOC-FLOD
% teaching: NEW-U14-NOM-FLOD
\reserve{11\baselineskip}
% mdparagraph:P019:begin
\textbf{flōd} — masculine; flood; stream of water. Stem \textbf{flōd-}.
% mdparagraph:P019:end

% table: NEW-U14-NOM-FLOD
% mdtable:T04:begin
\par\begingroup
\fontsize{10.5}{13.2}\selectfont
\setlength{\LTleft}{0pt}\setlength{\LTright}{\fill}
\noindent
\begin{tabular}{@{}>{\raggedright\arraybackslash}p{39.46mm}>{\raggedright\arraybackslash}p{39.46mm}>{\raggedright\arraybackslash}p{39.46mm}@{}}
\toprule
\textbf{Case} & \textbf{Singular} & \textbf{Plural} \\
\midrule
Nominative & flōd & flōdas \\
Genitive & flōdes & flōda \\
Dative & flōde & flōdum \\
Accusative & flōd & flōdas \\
\bottomrule
\end{tabular}
\par\endgroup

% mdtable:T04:end
% teaching: NEW-U14-VOC-MUNT
% teaching: NEW-U14-NOM-MUNT
\reserve{11\baselineskip}
% mdparagraph:P020:begin
\textbf{munt} — masculine; mountain. Stem \textbf{munt-}.
% mdparagraph:P020:end

% table: NEW-U14-NOM-MUNT
% mdtable:T05:begin
\par\begingroup
\fontsize{10.5}{13.2}\selectfont
\setlength{\LTleft}{0pt}\setlength{\LTright}{\fill}
\noindent
\begin{tabular}{@{}>{\raggedright\arraybackslash}p{39.46mm}>{\raggedright\arraybackslash}p{39.46mm}>{\raggedright\arraybackslash}p{39.46mm}@{}}
\toprule
\textbf{Case} & \textbf{Singular} & \textbf{Plural} \\
\midrule
Nominative & munt & muntas \\
Genitive & muntes & munta \\
Dative & munte & muntum \\
Accusative & munt & muntas \\
\bottomrule
\end{tabular}
\par\endgroup

% mdtable:T05:end
% teaching: NEW-U14-VOC-HUND-DOG
% teaching: NEW-U14-NOM-HUND-DOG
\reserve{11\baselineskip}
% mdparagraph:P021:begin
\textbf{hund} — masculine; dog. Stem \textbf{hund-}.
% mdparagraph:P021:end

% table: NEW-U14-NOM-HUND-DOG
% mdtable:T06:begin
\par\begingroup
\fontsize{10.5}{13.2}\selectfont
\setlength{\LTleft}{0pt}\setlength{\LTright}{\fill}
\noindent
\begin{tabular}{@{}>{\raggedright\arraybackslash}p{39.46mm}>{\raggedright\arraybackslash}p{39.46mm}>{\raggedright\arraybackslash}p{39.46mm}@{}}
\toprule
\textbf{Case} & \textbf{Singular} & \textbf{Plural} \\
\midrule
Nominative & hund & hundas \\
Genitive & hundes & hunda \\
Dative & hunde & hundum \\
Accusative & hund & hundas \\
\bottomrule
\end{tabular}
\par\endgroup

% mdtable:T06:end
% mdparagraph:P022:begin
This hund is the noun “dog”, distinct from the earlier numeral hund.
% mdparagraph:P022:end

% teaching: NEW-U14-VOC-MIDDANGEARD
% teaching: NEW-U14-NOM-MIDDANGEARD
\reserve{11\baselineskip}
% mdparagraph:P023:begin
\textbf{middangeard} — masculine; world; earth. Stem \textbf{middangeard-}. Plural reference; singular productive.
% mdparagraph:P023:end

% table: NEW-U14-NOM-MIDDANGEARD
% mdtable:T07:begin
\par\begingroup
\fontsize{10.5}{13.2}\selectfont
\setlength{\LTleft}{0pt}\setlength{\LTright}{\fill}
\noindent
\begin{tabular}{@{}>{\raggedright\arraybackslash}p{32.88mm}>{\raggedright\arraybackslash}p{42.75mm}>{\raggedright\arraybackslash}p{42.75mm}@{}}
\toprule
\textbf{Case} & \textbf{Singular} & \textbf{Plural} \\
\midrule
Nominative & middangeard & middangeardas \\
Genitive & middangeardes & middangearda \\
Dative & middangearde & middangeardum \\
Accusative & middangeard & middangeardas \\
\bottomrule
\end{tabular}
\par\endgroup

% mdtable:T07:end
\reserve{7\baselineskip}
% mdparagraph:P024:begin
\subsection*{Strong neuter nouns}\BookHeading{nav-ch14-learner-0004}{subsubsection}{Strong neuter nouns}{0}
% mdparagraph:P024:end

% mdparagraph:P025:begin
In these heavy-stem neuters the nominative and accusative are alike in both numbers and have no ending. Genitive and dative singular take -es and -e; genitive and dative plural take -a and -um. Add the ending to the whole compound stem.
% mdparagraph:P025:end

% teaching: NEW-U14-VOC-DUST
% teaching: NEW-U14-NOM-DUST
\reserve{11\baselineskip}
% mdparagraph:P026:begin
\textbf{dūst} — neuter; dust. Stem \textbf{dūst-}. Plural reference; singular productive.
% mdparagraph:P026:end

% table: NEW-U14-NOM-DUST
% mdtable:T08:begin
\par\begingroup
\fontsize{10.5}{13.2}\selectfont
\setlength{\LTleft}{0pt}\setlength{\LTright}{\fill}
\noindent
\begin{tabular}{@{}>{\raggedright\arraybackslash}p{39.46mm}>{\raggedright\arraybackslash}p{39.46mm}>{\raggedright\arraybackslash}p{39.46mm}@{}}
\toprule
\textbf{Case} & \textbf{Singular} & \textbf{Plural} \\
\midrule
Nominative & dūst & dūst \\
Genitive & dūstes & dūsta \\
Dative & dūste & dūstum \\
Accusative & dūst & dūst \\
\bottomrule
\end{tabular}
\par\endgroup

% mdtable:T08:end
% teaching: NEW-U14-VOC-WYRTGEMANG
% teaching: NEW-U14-NOM-WYRTGEMANG
\reserve{11\baselineskip}
% mdparagraph:P027:begin
\textbf{wyrtgemang} — neuter; spice; mixture of spices. Stem \textbf{wyrtgemang-}.
% mdparagraph:P027:end

% table: NEW-U14-NOM-WYRTGEMANG
% mdtable:T09:begin
\par\begingroup
\fontsize{10.5}{13.2}\selectfont
\setlength{\LTleft}{0pt}\setlength{\LTright}{\fill}
\noindent
\begin{tabular}{@{}>{\raggedright\arraybackslash}p{34.82mm}>{\raggedright\arraybackslash}p{41.78mm}>{\raggedright\arraybackslash}p{41.78mm}@{}}
\toprule
\textbf{Case} & \textbf{Singular} & \textbf{Plural} \\
\midrule
Nominative & wyrtgemang & wyrtgemang \\
Genitive & wyrtgemanges & wyrtgemanga \\
Dative & wyrtgemange & wyrtgemangum \\
Accusative & wyrtgemang & wyrtgemang \\
\bottomrule
\end{tabular}
\par\endgroup

% mdtable:T09:end
% teaching: NEW-U14-VOC-GOLDHORD
% teaching: NEW-U14-NOM-GOLDHORD
\reserve{11\baselineskip}
% mdparagraph:P028:begin
\textbf{goldhord} — neuter; treasure; hoard of gold. Stem \textbf{goldhord-}.
% mdparagraph:P028:end

% table: NEW-U14-NOM-GOLDHORD
% mdtable:T10:begin
\par\begingroup
\fontsize{10.5}{13.2}\selectfont
\setlength{\LTleft}{0pt}\setlength{\LTright}{\fill}
\noindent
\begin{tabular}{@{}>{\raggedright\arraybackslash}p{39.46mm}>{\raggedright\arraybackslash}p{39.46mm}>{\raggedright\arraybackslash}p{39.46mm}@{}}
\toprule
\textbf{Case} & \textbf{Singular} & \textbf{Plural} \\
\midrule
Nominative & goldhord & goldhord \\
Genitive & goldhordes & goldhorda \\
Dative & goldhorde & goldhordum \\
Accusative & goldhord & goldhord \\
\bottomrule
\end{tabular}
\par\endgroup

% mdtable:T10:end
% teaching: NEW-U14-VOC-GODWEBB
% teaching: NEW-U14-NOM-GODWEBB
\reserve{11\baselineskip}
% mdparagraph:P029:begin
\textbf{gōdwębb} — neuter; purple cloth. Stem \textbf{gōdwębb-}. Plural reference; singular productive.
% mdparagraph:P029:end

% table: NEW-U14-NOM-GODWEBB
% mdtable:T11:begin
\par\begingroup
\fontsize{10.5}{13.2}\selectfont
\setlength{\LTleft}{0pt}\setlength{\LTright}{\fill}
\noindent
\begin{tabular}{@{}>{\raggedright\arraybackslash}p{39.46mm}>{\raggedright\arraybackslash}p{39.46mm}>{\raggedright\arraybackslash}p{39.46mm}@{}}
\toprule
\textbf{Case} & \textbf{Singular} & \textbf{Plural} \\
\midrule
Nominative & gōdwębb & gōdwębb \\
Genitive & gōdwębbes & gōdwębba \\
Dative & gōdwębbe & gōdwębbum \\
Accusative & gōdwębb & gōdwębb \\
\bottomrule
\end{tabular}
\par\endgroup

% mdtable:T11:end
% teaching: NEW-U14-VOC-WIC
% teaching: NEW-U14-NOM-WIC
\reserve{11\baselineskip}
% mdparagraph:P030:begin
\textbf{wīc} — neuter; dwelling; abode. Stem \textbf{wīc-}.
% mdparagraph:P030:end

% table: NEW-U14-NOM-WIC
% mdtable:T12:begin
\par\begingroup
\fontsize{10.5}{13.2}\selectfont
\setlength{\LTleft}{0pt}\setlength{\LTright}{\fill}
\noindent
\begin{tabular}{@{}>{\raggedright\arraybackslash}p{39.46mm}>{\raggedright\arraybackslash}p{39.46mm}>{\raggedright\arraybackslash}p{39.46mm}@{}}
\toprule
\textbf{Case} & \textbf{Singular} & \textbf{Plural} \\
\midrule
Nominative & wīc & wīc \\
Genitive & wīces & wīca \\
Dative & wīce & wīcum \\
Accusative & wīc & wīc \\
\bottomrule
\end{tabular}
\par\endgroup

% mdtable:T12:end
% teaching: NEW-U14-VOC-GOD-GOODS
% teaching: NEW-U14-NOM-GOD-GOODS
\reserve{12\baselineskip}
% mdparagraph:P031:begin
\textbf{gōd} — neuter; goods; good things; property (plural). Stem \textbf{gōd-}. Singular reference; plural productive in the goods/\allowbreak{}property sense.
% mdparagraph:P031:end

% table: NEW-U14-NOM-GOD-GOODS
% mdtable:T13:begin
\par\begingroup
\fontsize{10.5}{13.2}\selectfont
\setlength{\LTleft}{0pt}\setlength{\LTright}{\fill}
\noindent
\begin{tabular}{@{}>{\raggedright\arraybackslash}p{39.46mm}>{\raggedright\arraybackslash}p{39.46mm}>{\raggedright\arraybackslash}p{39.46mm}@{}}
\toprule
\textbf{Case} & \textbf{Singular} & \textbf{Plural} \\
\midrule
Nominative & gōd & gōd \\
Genitive & gōdes & gōda \\
Dative & gōde & gōdum \\
Accusative & gōd & gōd \\
\bottomrule
\end{tabular}
\par\endgroup

% mdtable:T13:end
\reserve{7\baselineskip}
% mdparagraph:P032:begin
\subsection*{The reduced stem of wuldor}\BookHeading{nav-ch14-learner-0005}{subsubsection}{The reduced stem of wuldor}{0}
% mdparagraph:P032:end

% mdparagraph:P033:begin
Wuldor\BookPoint{bk-ch14-learner-fdc0479f56b0} loses the unstressed o before an ending: wuldr-. Its reference nominative/\allowbreak{}accusative plural wuldru follows the admitted disyllabic-neuter -u pattern. That plural is a rule-based reference form; it does not assign a new count meaning of “glories”.
% mdparagraph:P033:end

% teaching: NEW-U14-VOC-WULDOR
% teaching: NEW-U14-NOM-WULDOR
\reserve{11\baselineskip}
% mdparagraph:P034:begin
\textbf{wuldor} — neuter; glory. Stem \textbf{wuldr-}. Plural reference; singular productive.
% mdparagraph:P034:end

% table: NEW-U14-NOM-WULDOR
% mdtable:T14:begin
\par\begingroup
\fontsize{10.5}{13.2}\selectfont
\setlength{\LTleft}{0pt}\setlength{\LTright}{\fill}
\noindent
\begin{tabular}{@{}>{\raggedright\arraybackslash}p{39.46mm}>{\raggedright\arraybackslash}p{39.46mm}>{\raggedright\arraybackslash}p{39.46mm}@{}}
\toprule
\textbf{Case} & \textbf{Singular} & \textbf{Plural} \\
\midrule
Nominative & wuldor & wuldru \\
Genitive & wuldres & wuldra \\
Dative & wuldre & wuldrum \\
Accusative & wuldor & wuldru \\
\bottomrule
\end{tabular}
\par\endgroup

% mdtable:T14:end
\reserve{7\baselineskip}
% mdparagraph:P035:begin
\subsection*{Strong feminine nouns}\BookHeading{nav-ch14-learner-0006}{subsubsection}{Strong feminine nouns}{0}
% mdparagraph:P035:end

% mdparagraph:P036:begin
The nominative singular has no ending; genitive, dative and accusative singular take -e. Nominative, genitive and accusative plural take -a, and dative plural -um. In sōðfæstnes the inflectional stem has ss: sōðfæstness-. The tables select genitive plural -a.
% mdparagraph:P036:end

% teaching: NEW-U14-VOC-SOTHFAESTNES
% teaching: NEW-U14-NOM-SOTHFAESTNES
\reserve{11\baselineskip}
% mdparagraph:P037:begin
\textbf{sōðfæstnes} — feminine; truth. Stem \textbf{sōðfæstness-}. Plural reference; singular productive.
% mdparagraph:P037:end

% table: NEW-U14-NOM-SOTHFAESTNES
% mdtable:T15:begin
\par\begingroup
\fontsize{10.5}{13.2}\selectfont
\setlength{\LTleft}{0pt}\setlength{\LTright}{\fill}
\noindent
\begin{tabular}{@{}>{\raggedright\arraybackslash}p{33.82mm}>{\raggedright\arraybackslash}p{40.59mm}>{\raggedright\arraybackslash}p{43.97mm}@{}}
\toprule
\textbf{Case} & \textbf{Singular} & \textbf{Plural} \\
\midrule
Nominative & sōðfæstnes & sōðfæstnessa \\
Genitive & sōðfæstnesse & sōðfæstnessa \\
Dative & sōðfæstnesse & sōðfæstnessum \\
Accusative & sōðfæstnesse & sōðfæstnessa \\
\bottomrule
\end{tabular}
\par\endgroup

% mdtable:T15:end
% teaching: NEW-U14-VOC-SIBB
% teaching: NEW-U14-NOM-SIBB
\reserve{11\baselineskip}
% mdparagraph:P038:begin
\textbf{sibb} — feminine; peace. Stem \textbf{sibb-}. Plural reference; singular productive.
% mdparagraph:P038:end

% table: NEW-U14-NOM-SIBB
% mdtable:T16:begin
\par\begingroup
\fontsize{10.5}{13.2}\selectfont
\setlength{\LTleft}{0pt}\setlength{\LTright}{\fill}
\noindent
\begin{tabular}{@{}>{\raggedright\arraybackslash}p{39.46mm}>{\raggedright\arraybackslash}p{39.46mm}>{\raggedright\arraybackslash}p{39.46mm}@{}}
\toprule
\textbf{Case} & \textbf{Singular} & \textbf{Plural} \\
\midrule
Nominative & sibb & sibba \\
Genitive & sibbe & sibba \\
Dative & sibbe & sibbum \\
Accusative & sibbe & sibba \\
\bottomrule
\end{tabular}
\par\endgroup

% mdtable:T16:end
% teaching: NEW-U14-VOC-WUND
% teaching: NEW-U14-NOM-WUND
\reserve{11\baselineskip}
% mdparagraph:P039:begin
\textbf{wund} — feminine; wound; sore. Stem \textbf{wund-}.
% mdparagraph:P039:end

% table: NEW-U14-NOM-WUND
% mdtable:T17:begin
\par\begingroup
\fontsize{10.5}{13.2}\selectfont
\setlength{\LTleft}{0pt}\setlength{\LTright}{\fill}
\noindent
\begin{tabular}{@{}>{\raggedright\arraybackslash}p{39.46mm}>{\raggedright\arraybackslash}p{39.46mm}>{\raggedright\arraybackslash}p{39.46mm}@{}}
\toprule
\textbf{Case} & \textbf{Singular} & \textbf{Plural} \\
\midrule
Nominative & wund & wunda \\
Genitive & wunde & wunda \\
Dative & wunde & wundum \\
Accusative & wunde & wunda \\
\bottomrule
\end{tabular}
\par\endgroup

% mdtable:T17:end
% mdparagraph:P040:begin
This feminine noun means “wound; sore”; it is distinct from the adjective “wounded”. See correction C04 in the editorial notes.
% mdparagraph:P040:end

\reserve{7\baselineskip}
% mdparagraph:P041:begin
\subsection*{Weak nouns}\BookHeading{nav-ch14-learner-0007}{subsubsection}{Weak nouns}{0}
% mdparagraph:P041:end

% mdparagraph:P042:begin
The weak masculine citation ends in -a; its other singular cases and its nominative/\allowbreak{}accusative plural end in -an. The weak neuter ēage instead keeps -e in both nominative and accusative singular. Its other -an forms coincide with those of the masculine pattern. Thus -an alone cannot decide case or number. All four nouns use genitive plural -ena and dative plural -um; ēage also permits genitive plural ēagna beside selected ēagena.
% mdparagraph:P042:end

% teaching: NEW-U14-VOC-WAEDLA
% teaching: NEW-U14-NOM-WAEDLA
\reserve{11\baselineskip}
% mdparagraph:P043:begin
\textbf{wǣdla} — masculine; poor or destitute man. Stem \textbf{wǣdl-}.
% mdparagraph:P043:end

% table: NEW-U14-NOM-WAEDLA
% mdtable:T18:begin
\par\begingroup
\fontsize{10.5}{13.2}\selectfont
\setlength{\LTleft}{0pt}\setlength{\LTright}{\fill}
\noindent
\begin{tabular}{@{}>{\raggedright\arraybackslash}p{39.46mm}>{\raggedright\arraybackslash}p{39.46mm}>{\raggedright\arraybackslash}p{39.46mm}@{}}
\toprule
\textbf{Case} & \textbf{Singular} & \textbf{Plural} \\
\midrule
Nominative & wǣdla & wǣdlan \\
Genitive & wǣdlan & wǣdlena \\
Dative & wǣdlan & wǣdlum \\
Accusative & wǣdlan & wǣdlan \\
\bottomrule
\end{tabular}
\par\endgroup

% mdtable:T18:end
% teaching: NEW-U14-VOC-WITGA
% teaching: NEW-U14-NOM-WITGA
\reserve{11\baselineskip}
% mdparagraph:P044:begin
\textbf{wītga} — masculine; prophet. Stem \textbf{wītg-}.
% mdparagraph:P044:end

% table: NEW-U14-NOM-WITGA
% mdtable:T19:begin
\par\begingroup
\fontsize{10.5}{13.2}\selectfont
\setlength{\LTleft}{0pt}\setlength{\LTright}{\fill}
\noindent
\begin{tabular}{@{}>{\raggedright\arraybackslash}p{39.46mm}>{\raggedright\arraybackslash}p{39.46mm}>{\raggedright\arraybackslash}p{39.46mm}@{}}
\toprule
\textbf{Case} & \textbf{Singular} & \textbf{Plural} \\
\midrule
Nominative & wītga & wītgan \\
Genitive & wītgan & wītgena \\
Dative & wītgan & wītgum \\
Accusative & wītgan & wītgan \\
\bottomrule
\end{tabular}
\par\endgroup

% mdtable:T19:end
% teaching: NEW-U14-VOC-THEARFA
% teaching: NEW-U14-NOM-THEARFA
\reserve{11\baselineskip}
% mdparagraph:P045:begin
\textbf{ðearfa} — masculine; poor or needy person. Stem \textbf{ðearf-}.
% mdparagraph:P045:end

% table: NEW-U14-NOM-THEARFA
% mdtable:T20:begin
\par\begingroup
\fontsize{10.5}{13.2}\selectfont
\setlength{\LTleft}{0pt}\setlength{\LTright}{\fill}
\noindent
\begin{tabular}{@{}>{\raggedright\arraybackslash}p{39.46mm}>{\raggedright\arraybackslash}p{39.46mm}>{\raggedright\arraybackslash}p{39.46mm}@{}}
\toprule
\textbf{Case} & \textbf{Singular} & \textbf{Plural} \\
\midrule
Nominative & ðearfa & ðearfan \\
Genitive & ðearfan & ðearfena \\
Dative & ðearfan & ðearfum \\
Accusative & ðearfan & ðearfan \\
\bottomrule
\end{tabular}
\par\endgroup

% mdtable:T20:end
% teaching: NEW-U14-VOC-EAGE
% teaching: NEW-U14-NOM-EAGE
\reserve{11\baselineskip}
% mdparagraph:P046:begin
\textbf{ēage} — neuter; eye. Stem \textbf{ēag-}.
% mdparagraph:P046:end

% table: NEW-U14-NOM-EAGE
% mdtable:T21:begin
\par\begingroup
\fontsize{10.5}{13.2}\selectfont
\setlength{\LTleft}{0pt}\setlength{\LTright}{\fill}
\noindent
\begin{tabular}{@{}>{\raggedright\arraybackslash}p{39.46mm}>{\raggedright\arraybackslash}p{39.46mm}>{\raggedright\arraybackslash}p{39.46mm}@{}}
\toprule
\textbf{Case} & \textbf{Singular} & \textbf{Plural} \\
\midrule
Nominative & ēage & ēagan \\
Genitive & ēagan & ēagena \\
Dative & ēagan & ēagum \\
Accusative & ēage & ēagan \\
\bottomrule
\end{tabular}
\par\endgroup

% mdtable:T21:end
\reserve{7\baselineskip}
% mdparagraph:P047:begin
\section*{3. Indicative systems}\BookHeading{nav-ch14-learner-0008}{subsection}{3. Indicative systems}{1}
% mdparagraph:P047:end

% mdparagraph:P048:begin
The nine systems below give all present and preterite indicative persons. A shared \textbf{1–3 pl.} row gives the common form for all three plural persons. The ordinary meanings and government follow in §5. The four principal parts identify the stems; their first and fourth columns here are diagnostic citations. They do not add productive infinitive or participial uses to these nine systems.\textsuperscript{5}
% mdparagraph:P048:end

\reserve{21\baselineskip}
% mdparagraph:P049:begin
The regular selected present endings are singular -e, -st, -ð and plural -að, with the stem changes and contractions stated beside each table. Weak II has its separate present pattern below. Strong preterites have no ending in first/\allowbreak{}third singular, -e with the plural stem in second singular, and -on in the plural. Weak preterites use the dental stem with singular -e, -est, -e and plural -on.
% mdparagraph:P049:end

% mdtable:T22:begin
\par\begingroup
\fontsize{10.5}{13.2}\selectfont
\setlength{\LTleft}{0pt}\setlength{\LTright}{\fill}
\noindent
\begin{tabular}{@{}>{\raggedright\arraybackslash}p{27.02mm}>{\raggedright\arraybackslash}p{28.55mm}>{\raggedright\arraybackslash}p{23.36mm}>{\raggedright\arraybackslash}p{36.63mm}@{}}
\toprule
\textbf{Infinitive citation} & \textbf{Preterite 1/\allowbreak{}3 singular} & \textbf{Preterite plural} & \textbf{Past-participle citation} \\
\midrule
gewęndan & gewęnde & gewęndon & gewęnd \\
gelīefan & gelīefde & gelīefdon & gelīefed \\
forgiefan & forgeaf & forgēafon & forgiefen \\
āhrēosan & āhrēas & āhruron & āhroren \\
gebiddan & gebæd & gebǣdon & gebeden \\
ætīewan & ætīewde & ætīewdon & ætīewed \\
liccian & liccode & liccodon & liccod \\
āhębban & āhōf & āhōfon & āhafen \\
onfōn & onfēng & onfēngon & onfangen \\
\bottomrule
\end{tabular}
\par\endgroup

% mdtable:T22:end
% teaching: NEW-U14-VOC-GEWENDAN
% teaching: NEW-U14-PAR-GEWENDAN
% teaching: NEW-U14-PAR-GEWENDAN-IND
% teaching: NEW-U14-GRAM-U14-GEWENDAN-INDICATIVE
\par\begin{minipage}{\linewidth}\raggedright
% mdparagraph:P050:begin
\subsection*{gewęndan — return; turn}\BookHeading{nav-ch14-learner-0009}{subsubsection}{gewęndan — return; turn}{0}
% mdparagraph:P050:end

% mdparagraph:P051:begin
Weak I. Use present stem gewęnd-, with d devoiced before the second/\allowbreak{}third singular endings: gewęntst, gewęnt. Preterite gewęnde simplifies the doubled dental; its stem is gewęnd- plus the selected past endings. The lexical ge- is retained.
% mdparagraph:P051:end

% table: NEW-U14-PAR-GEWENDAN-IND
% mdtable:T23:begin
% course-table-style: 2026-10
% course-panel:T23-P01:begin
\par\begingroup
\fontsize{10.5}{13.2}\selectfont
\setlength{\LTleft}{0pt}\setlength{\LTright}{\fill}
\noindent
\begin{tabular}{@{}>{\raggedright\arraybackslash}p{39.459mm}>{\raggedright\arraybackslash}p{39.459mm}>{\raggedright\arraybackslash}p{39.459mm}@{}}
\toprule
\textbf{Person and number} & \textbf{Present} & \textbf{Preterite} \\
\midrule
1 sg. & gewęnde & gewęnde \\
2 sg. & gewęntst & gewęndest \\
3 sg. & gewęnt & gewęnde \\
1–3 pl. & gewęndað & gewęndon \\
\bottomrule
\end{tabular}
\par\endgroup
% course-panel:T23-P01:end
% mdtable:T23:end
\end{minipage}\par\addvspace{4pt}
% teaching: NEW-U14-VOC-GELIEFAN
% teaching: NEW-U14-PAR-GELIEFAN
% teaching: NEW-U14-PAR-GELIEFAN-IND
% teaching: NEW-U14-GRAM-U14-GELIEFAN-INDICATIVE
\par\begin{minipage}{\linewidth}\raggedright
% mdparagraph:P052:begin
\subsection*{gelīefan — believe}\BookHeading{nav-ch14-learner-0010}{subsubsection}{gelīefan — believe}{0}
% mdparagraph:P052:end

% mdparagraph:P053:begin
Weak I. Retain long īe throughout. Present endings are -e, -st, -ð and plural -að; the dental past stem is gelīefd-, with -e, -est, -e and plural -on. The infinitive and past-participle citations differ from the finite forms.
% mdparagraph:P053:end

% table: NEW-U14-PAR-GELIEFAN-IND
% mdtable:T24:begin
% course-table-style: 2026-10
% course-panel:T24-P01:begin
\par\begingroup
\fontsize{10.5}{13.2}\selectfont
\setlength{\LTleft}{0pt}\setlength{\LTright}{\fill}
\noindent
\begin{tabular}{@{}>{\raggedright\arraybackslash}p{39.459mm}>{\raggedright\arraybackslash}p{39.459mm}>{\raggedright\arraybackslash}p{39.459mm}@{}}
\toprule
\textbf{Person and number} & \textbf{Present} & \textbf{Preterite} \\
\midrule
1 sg. & gelīefe & gelīefde \\
2 sg. & gelīefst & gelīefdest \\
3 sg. & gelīefð & gelīefde \\
1–3 pl. & gelīefað & gelīefdon \\
\bottomrule
\end{tabular}
\par\endgroup
% course-panel:T24-P01:end
% mdtable:T24:end
\end{minipage}\par\addvspace{4pt}
% teaching: NEW-U14-VOC-FORGIEFAN
% teaching: NEW-U14-PAR-FORGIEFAN
% teaching: NEW-U14-PAR-FORGIEFAN-IND
% teaching: NEW-U14-GRAM-U14-FORGIEFAN-INDICATIVE
\par\begin{minipage}{\linewidth}\raggedright
% mdparagraph:P054:begin
\subsection*{forgiefan — give; grant}\BookHeading{nav-ch14-learner-0011}{subsubsection}{forgiefan — give; grant}{0}
% mdparagraph:P054:end

% mdparagraph:P055:begin
Strong V. Principal stems are forgief-, forgeaf-, forgēaf-, forgief-. The selected present forms retain ie throughout. Past second singular uses the plural stem and -e; first and third singular share forgeaf.
% mdparagraph:P055:end

% table: NEW-U14-PAR-FORGIEFAN-IND
% mdtable:T25:begin
% course-table-style: 2026-10
% course-panel:T25-P01:begin
\par\begingroup
\fontsize{10.5}{13.2}\selectfont
\setlength{\LTleft}{0pt}\setlength{\LTright}{\fill}
\noindent
\begin{tabular}{@{}>{\raggedright\arraybackslash}p{39.459mm}>{\raggedright\arraybackslash}p{39.459mm}>{\raggedright\arraybackslash}p{39.459mm}@{}}
\toprule
\textbf{Person and number} & \textbf{Present} & \textbf{Preterite} \\
\midrule
1 sg. & forgiefe & forgeaf \\
2 sg. & forgiefst & forgēafe \\
3 sg. & forgiefð & forgeaf \\
1–3 pl. & forgiefað & forgēafon \\
\bottomrule
\end{tabular}
\par\endgroup
% course-panel:T25-P01:end
% mdtable:T25:end
\end{minipage}\par\addvspace{4pt}
% teaching: NEW-U14-VOC-AHREOSAN
% teaching: NEW-U14-PAR-AHREOSAN
% teaching: NEW-U14-PAR-AHREOSAN-IND
% teaching: NEW-U14-GRAM-U14-AHREOSAN-INDICATIVE
\par\begin{minipage}{\linewidth}\raggedright
% mdparagraph:P056:begin
\subsection*{āhrēosan — fall; fall upon}\BookHeading{nav-ch14-learner-0012}{subsubsection}{āhrēosan — fall; fall upon}{0}
% mdparagraph:P056:end

% mdparagraph:P057:begin
Strong II. The principal stems are āhrēos-, āhrēas-, āhrur-, āhror-. Present second/\allowbreak{}third singular have long īe and the selected contracted ending -st. Past second singular and all plural persons use āhrur-, including the s/\allowbreak{}r alternation.
% mdparagraph:P057:end

% table: NEW-U14-PAR-AHREOSAN-IND
% mdtable:T26:begin
% course-table-style: 2026-10
% course-panel:T26-P01:begin
\par\begingroup
\fontsize{10.5}{13.2}\selectfont
\setlength{\LTleft}{0pt}\setlength{\LTright}{\fill}
\noindent
\begin{tabular}{@{}>{\raggedright\arraybackslash}p{39.459mm}>{\raggedright\arraybackslash}p{39.459mm}>{\raggedright\arraybackslash}p{39.459mm}@{}}
\toprule
\textbf{Person and number} & \textbf{Present} & \textbf{Preterite} \\
\midrule
1 sg. & āhrēose & āhrēas \\
2 sg. & āhrīest & āhrure \\
3 sg. & āhrīest & āhrēas \\
1–3 pl. & āhrēosað & āhruron \\
\bottomrule
\end{tabular}
\par\endgroup
% course-panel:T26-P01:end
% mdtable:T26:end
\end{minipage}\par\addvspace{4pt}
% teaching: NEW-U14-VOC-GEBIDDAN
% teaching: NEW-U14-PAR-GEBIDDAN
% teaching: NEW-U14-PAR-GEBIDDAN-IND
% teaching: NEW-U14-GRAM-U14-GEBIDDAN-INDICATIVE
\par\begin{minipage}{\linewidth}\raggedright
% mdparagraph:P058:begin
\subsection*{gebiddan — pray}\BookHeading{nav-ch14-learner-0013}{subsubsection}{gebiddan — pray}{0}
% mdparagraph:P058:end

% mdparagraph:P059:begin
Strong V. Principal stems are gebidd-, gebæd-, gebǣd-, gebed-. The present second/\allowbreak{}third singular simplify the consonants to gebitst and gebitt. Past second singular uses the long-vowel plural stem gebǣd- with -e.
% mdparagraph:P059:end

% table: NEW-U14-PAR-GEBIDDAN-IND
% mdtable:T27:begin
% course-table-style: 2026-10
% course-panel:T27-P01:begin
\par\begingroup
\fontsize{10.5}{13.2}\selectfont
\setlength{\LTleft}{0pt}\setlength{\LTright}{\fill}
\noindent
\begin{tabular}{@{}>{\raggedright\arraybackslash}p{39.459mm}>{\raggedright\arraybackslash}p{39.459mm}>{\raggedright\arraybackslash}p{39.459mm}@{}}
\toprule
\textbf{Person and number} & \textbf{Present} & \textbf{Preterite} \\
\midrule
1 sg. & gebidde & gebæd \\
2 sg. & gebitst & gebǣde \\
3 sg. & gebitt & gebæd \\
1–3 pl. & gebiddað & gebǣdon \\
\bottomrule
\end{tabular}
\par\endgroup
% course-panel:T27-P01:end
% mdtable:T27:end
\end{minipage}\par\addvspace{4pt}
% teaching: NEW-U14-VOC-AETIEWAN
% teaching: NEW-U14-PAR-AETIEWAN
% teaching: NEW-U14-PAR-AETIEWAN-IND
% teaching: NEW-U14-GRAM-U14-AETIEWAN-INDICATIVE
\par\begin{minipage}{\linewidth}\raggedright
% mdparagraph:P060:begin
\subsection*{ætīewan — show; display; reveal}\BookHeading{nav-ch14-learner-0014}{subsubsection}{ætīewan — show; display; reveal}{0}
% mdparagraph:P060:end

% mdparagraph:P061:begin
Weak I. Retain w: present ætīew-, dental past ætīewd-. This transitive verb “show” is distinct from earlier weak-II ætēowian, “appear”.
% mdparagraph:P061:end

% table: NEW-U14-PAR-AETIEWAN-IND
% mdtable:T28:begin
% course-table-style: 2026-10
% course-panel:T28-P01:begin
\par\begingroup
\fontsize{10.5}{13.2}\selectfont
\setlength{\LTleft}{0pt}\setlength{\LTright}{\fill}
\noindent
\begin{tabular}{@{}>{\raggedright\arraybackslash}p{39.459mm}>{\raggedright\arraybackslash}p{39.459mm}>{\raggedright\arraybackslash}p{39.459mm}@{}}
\toprule
\textbf{Person and number} & \textbf{Present} & \textbf{Preterite} \\
\midrule
1 sg. & ætīewe & ætīewde \\
2 sg. & ætīewst & ætīewdest \\
3 sg. & ætīewð & ætīewde \\
1–3 pl. & ætīewað & ætīewdon \\
\bottomrule
\end{tabular}
\par\endgroup
% course-panel:T28-P01:end
% mdtable:T28:end
\end{minipage}\par\addvspace{4pt}
% teaching: NEW-U14-VOC-LICCIAN
% teaching: NEW-U14-PAR-LICCIAN
% teaching: NEW-U14-PAR-LICCIAN-IND
% teaching: NEW-U14-GRAM-U14-LICCIAN-INDICATIVE
\par\begin{minipage}{\linewidth}\raggedright
% mdparagraph:P062:begin
\subsection*{liccian — lick}\BookHeading{nav-ch14-learner-0015}{subsubsection}{liccian — lick}{0}
% mdparagraph:P062:end

% mdparagraph:P063:begin
Weak II. Present first singular has -ie, second -ast, third -að and plural -iað. The past stem liccod- takes -e, -est, -e and plural -on. Both c letters remain; the selected past uses -ode, not an unassigned alternative in -ade.
% mdparagraph:P063:end

% table: NEW-U14-PAR-LICCIAN-IND
% mdtable:T29:begin
% course-table-style: 2026-10
% course-panel:T29-P01:begin
\par\begingroup
\fontsize{10.5}{13.2}\selectfont
\setlength{\LTleft}{0pt}\setlength{\LTright}{\fill}
\noindent
\begin{tabular}{@{}>{\raggedright\arraybackslash}p{39.459mm}>{\raggedright\arraybackslash}p{39.459mm}>{\raggedright\arraybackslash}p{39.459mm}@{}}
\toprule
\textbf{Person and number} & \textbf{Present} & \textbf{Preterite} \\
\midrule
1 sg. & liccie & liccode \\
2 sg. & liccast & liccodest \\
3 sg. & liccað & liccode \\
1–3 pl. & licciað & liccodon \\
\bottomrule
\end{tabular}
\par\endgroup
% course-panel:T29-P01:end
% mdtable:T29:end
\end{minipage}\par\addvspace{4pt}
% teaching: NEW-U14-VOC-AHEBBAN
% teaching: NEW-U14-PAR-AHEBBAN
% teaching: NEW-U14-PAR-AHEBBAN-IND
% teaching: NEW-U14-GRAM-U14-AHEBBAN-INDICATIVE
\par\begin{minipage}{\linewidth}\raggedright
% mdparagraph:P064:begin
\subsection*{āhębban — raise; lift}\BookHeading{nav-ch14-learner-0016}{subsubsection}{āhębban — raise; lift}{0}
% mdparagraph:P064:end

% mdparagraph:P065:begin
Strong VI. Present āhębb- appears with bb in first singular and plural, but with f before second/\allowbreak{}third singular endings. The past stem is āhōf-: zero ending in first/\allowbreak{}third singular, -e in second singular and -on in the plural. Retain the prefix ā-.
% mdparagraph:P065:end

% table: NEW-U14-PAR-AHEBBAN-IND
% mdtable:T30:begin
% course-table-style: 2026-10
% course-panel:T30-P01:begin
\par\begingroup
\fontsize{10.5}{13.2}\selectfont
\setlength{\LTleft}{0pt}\setlength{\LTright}{\fill}
\noindent
\begin{tabular}{@{}>{\raggedright\arraybackslash}p{39.459mm}>{\raggedright\arraybackslash}p{39.459mm}>{\raggedright\arraybackslash}p{39.459mm}@{}}
\toprule
\textbf{Person and number} & \textbf{Present} & \textbf{Preterite} \\
\midrule
1 sg. & āhębbe & āhōf \\
2 sg. & āhęfst & āhōfe \\
3 sg. & āhęfð & āhōf \\
1–3 pl. & āhębbað & āhōfon \\
\bottomrule
\end{tabular}
\par\endgroup
% course-panel:T30-P01:end
% mdtable:T30:end
\end{minipage}\par\addvspace{4pt}
% teaching: NEW-U14-VOC-ONFON
% teaching: NEW-U14-PAR-ONFON
% teaching: NEW-U14-PAR-ONFON-IND
% teaching: NEW-U14-GRAM-U14-ONFON-INDICATIVE
\par\begin{minipage}{\linewidth}\raggedright
% mdparagraph:P066:begin
\subsection*{onfōn — receive; accept}\BookHeading{nav-ch14-learner-0017}{subsubsection}{onfōn — receive; accept}{0}
% mdparagraph:P066:end

% mdparagraph:P067:begin
Contracted strong VII. Compare present onfō, onfēhst, onfēhð and plural onfōð. The past stem onfēng- takes no ending in first/\allowbreak{}third singular, -e in second singular and -on in all plural persons. Its government must be learned separately.
% mdparagraph:P067:end

% table: NEW-U14-PAR-ONFON-IND
% mdtable:T31:begin
% course-table-style: 2026-10
% course-panel:T31-P01:begin
\par\begingroup
\fontsize{10.5}{13.2}\selectfont
\setlength{\LTleft}{0pt}\setlength{\LTright}{\fill}
\noindent
\begin{tabular}{@{}>{\raggedright\arraybackslash}p{39.459mm}>{\raggedright\arraybackslash}p{39.459mm}>{\raggedright\arraybackslash}p{39.459mm}@{}}
\toprule
\textbf{Person and number} & \textbf{Present} & \textbf{Preterite} \\
\midrule
1 sg. & onfō & onfēng \\
2 sg. & onfēhst & onfēnge \\
3 sg. & onfēhð & onfēng \\
1–3 pl. & onfōð & onfēngon \\
\bottomrule
\end{tabular}
\par\endgroup
% course-panel:T31-P01:end
% mdtable:T31:end
\end{minipage}\par\addvspace{4pt}
\reserve{7\baselineskip}
% mdparagraph:P068:begin
\section*{4. Subjunctives, participles and commands}\BookHeading{nav-ch14-learner-0018}{subsection}{4. Subjunctives, participles and commands}{1}
% mdparagraph:P068:end

\reserve{7\baselineskip}
% mdparagraph:P069:begin
\subsection*{Five complete subjunctive systems}\BookHeading{nav-ch14-learner-0019}{subsubsection}{Five complete subjunctive systems}{0}
% mdparagraph:P069:end

% mdparagraph:P070:begin
In each tense the subjunctive has one singular form for all three persons, ending in -e, and one plural form, ending in -en. Strong preterite subjunctives use the preterite plural stem; weak ones retain the dental past stem. Syncretism does not remove a distinction of tense or mood: sęnde and gelīefde require their clause context. The already taught purpose and request clauses select subjunctives; a factual ðæt-clause need not do so. The indefinite concessive expression in Cook 27.10 receives separate local help.
% mdparagraph:P070:end

% teaching: NEW-U14-VOC-FORWEORTHAN
% teaching: NEW-U14-PAR-FORWEORTHAN
% teaching: NEW-U14-PAR-FORWEORTHAN-SUBJ
% teaching: NEW-U14-GRAM-U14-FORWEORTHAN-SUBJUNCTIVE
\reserve{13\baselineskip}
% mdparagraph:P071:begin
\textbf{forweorðan.} “Perish”. Present stem forweorð-; past forwurd-. Keep the ð/\allowbreak{}d alternation and the short past u.
% mdparagraph:P071:end

% table: NEW-U14-PAR-FORWEORTHAN-SUBJ
% mdtable:T32:begin
% course-table-style: 2026-10
% course-panel:T32-P01:begin
\par\begingroup
\fontsize{10.5}{13.2}\selectfont
\setlength{\LTleft}{0pt}\setlength{\LTright}{\fill}
\noindent
\begin{tabular}{@{}>{\raggedright\arraybackslash}p{39.459mm}>{\raggedright\arraybackslash}p{39.459mm}>{\raggedright\arraybackslash}p{39.459mm}@{}}
\toprule
\textbf{Person and number} & \textbf{Present} & \textbf{Preterite} \\
\midrule
1 sg. & forweorðe & forwurde \\
2 sg. & forweorðe & forwurde \\
3 sg. & forweorðe & forwurde \\
1–3 pl. & forweorðen & forwurden \\
\bottomrule
\end{tabular}
\par\endgroup
% course-panel:T32-P01:end
% mdtable:T32:end
% teaching: NEW-U14-PAR-GEBIDDAN-SUBJ
% teaching: NEW-U14-GRAM-U14-GEBIDDAN-SUBJUNCTIVE
\reserve{14\baselineskip}
% mdparagraph:P072:begin
\textbf{gebiddan.} “Pray”. Present stem gebidd-; past gebǣd-. The accusative reflexive is retained in the purpose use as well as in factual indicative clauses.
% mdparagraph:P072:end

% table: NEW-U14-PAR-GEBIDDAN-SUBJ
% mdtable:T33:begin
% course-table-style: 2026-10
% course-panel:T33-P01:begin
\par\begingroup
\fontsize{10.5}{13.2}\selectfont
\setlength{\LTleft}{0pt}\setlength{\LTright}{\fill}
\noindent
\begin{tabular}{@{}>{\raggedright\arraybackslash}p{39.459mm}>{\raggedright\arraybackslash}p{39.459mm}>{\raggedright\arraybackslash}p{39.459mm}@{}}
\toprule
\textbf{Person and number} & \textbf{Present} & \textbf{Preterite} \\
\midrule
1 sg. & gebidde & gebǣde \\
2 sg. & gebidde & gebǣde \\
3 sg. & gebidde & gebǣde \\
1–3 pl. & gebidden & gebǣden \\
\bottomrule
\end{tabular}
\par\endgroup
% course-panel:T33-P01:end
% mdtable:T33:end
% teaching: NEW-U14-VOC-SENDAN
% teaching: NEW-U14-PAR-SENDAN
% teaching: NEW-U14-PAR-SENDAN-SUBJ
% teaching: NEW-U14-GRAM-U14-SENDAN-SUBJUNCTIVE
\reserve{15\baselineskip}
% mdparagraph:P073:begin
\textbf{sęndan.} “Send”. The dental past simplifies -ndd- to -nd-, so present and past forms coincide in the selected subjunctive. The earlier recipient frame remains accusative person or thing sent, dative recipient.
% mdparagraph:P073:end

% table: NEW-U14-PAR-SENDAN-SUBJ
% mdtable:T34:begin
% course-table-style: 2026-10
% course-panel:T34-P01:begin
\par\begingroup
\fontsize{10.5}{13.2}\selectfont
\setlength{\LTleft}{0pt}\setlength{\LTright}{\fill}
\noindent
\begin{tabular}{@{}>{\raggedright\arraybackslash}p{39.459mm}>{\raggedright\arraybackslash}p{39.459mm}>{\raggedright\arraybackslash}p{39.459mm}@{}}
\toprule
\textbf{Person and number} & \textbf{Present} & \textbf{Preterite} \\
\midrule
1 sg. & sęnde & sęnde \\
2 sg. & sęnde & sęnde \\
3 sg. & sęnde & sęnde \\
1–3 pl. & sęnden & sęnden \\
\bottomrule
\end{tabular}
\par\endgroup
% course-panel:T34-P01:end
% mdtable:T34:end
% teaching: NEW-U14-PAR-GELIEFAN-SUBJ
% teaching: NEW-U14-GRAM-U14-GELIEFAN-SUBJUNCTIVE
\reserve{13\baselineskip}
% mdparagraph:P074:begin
\textbf{gelīefan.} “Believe”. Present gelīef-; past gelīefd-. Keep the dental and long īe.
% mdparagraph:P074:end

% table: NEW-U14-PAR-GELIEFAN-SUBJ
% mdtable:T35:begin
% course-table-style: 2026-10
% course-panel:T35-P01:begin
\par\begingroup
\fontsize{10.5}{13.2}\selectfont
\setlength{\LTleft}{0pt}\setlength{\LTright}{\fill}
\noindent
\begin{tabular}{@{}>{\raggedright\arraybackslash}p{39.459mm}>{\raggedright\arraybackslash}p{39.459mm}>{\raggedright\arraybackslash}p{39.459mm}@{}}
\toprule
\textbf{Person and number} & \textbf{Present} & \textbf{Preterite} \\
\midrule
1 sg. & gelīefe & gelīefde \\
2 sg. & gelīefe & gelīefde \\
3 sg. & gelīefe & gelīefde \\
1–3 pl. & gelīefen & gelīefden \\
\bottomrule
\end{tabular}
\par\endgroup
% course-panel:T35-P01:end
% mdtable:T35:end
% teaching: NEW-U14-VOC-ARISAN
% teaching: NEW-U14-PAR-ARISAN
% teaching: NEW-U14-PAR-ARISAN-SUBJ
% teaching: NEW-U14-GRAM-U14-ARISAN-SUBJUNCTIVE
\reserve{13\baselineskip}
% mdparagraph:P075:begin
\textbf{ārīsan.} “Arise”. Present ārīs- has long ī; past āris- has short i. The prefix ā- stays long in both tenses.
% mdparagraph:P075:end

% table: NEW-U14-PAR-ARISAN-SUBJ
% mdtable:T36:begin
% course-table-style: 2026-10
% course-panel:T36-P01:begin
\par\begingroup
\fontsize{10.5}{13.2}\selectfont
\setlength{\LTleft}{0pt}\setlength{\LTright}{\fill}
\noindent
\begin{tabular}{@{}>{\raggedright\arraybackslash}p{39.459mm}>{\raggedright\arraybackslash}p{39.459mm}>{\raggedright\arraybackslash}p{39.459mm}@{}}
\toprule
\textbf{Person and number} & \textbf{Present} & \textbf{Preterite} \\
\midrule
1 sg. & ārīse & ārise \\
2 sg. & ārīse & ārise \\
3 sg. & ārīse & ārise \\
1–3 pl. & ārīsen & ārisen \\
\bottomrule
\end{tabular}
\par\endgroup
% course-panel:T36-P01:end
% mdtable:T36:end
\reserve{7\baselineskip}
% mdparagraph:P076:begin
\subsection*{Four participle-base pairs}\BookHeading{nav-ch14-learner-0020}{subsubsection}{Four participle-base pairs}{0}
% mdparagraph:P076:end

% mdparagraph:P077:begin
The present base uses the present stem + -ende. A strong past base uses its fourth stem + -en; a weak past base contains the dental formation. These four pairs teach the bases and the selected past-passive use. They do not establish complete attributive participle declensions. The source’s inflected present-participial phrases are supplied separately.
% mdparagraph:P077:end

% teaching: NEW-U14-VOC-FORCEORFAN
% teaching: NEW-U14-PAR-FORCEORFAN
% teaching: NEW-U14-PAR-FORCEORFAN-PART
% teaching: NEW-U14-GRAM-U14-FORCEORFAN-PARTICIPLE-BASES
\reserve{9\baselineskip}
% mdparagraph:P078:begin
\textbf{forceorfan.} “Cut down; cut off”. Present forceorf-; strong past forcorf- + -en.
% mdparagraph:P078:end

% table: NEW-U14-PAR-FORCEORFAN-PART
% mdtable:T37:begin
\par\begingroup
\fontsize{10.5}{13.2}\selectfont
\setlength{\LTleft}{0pt}\setlength{\LTright}{\fill}
\noindent
\begin{tabular}{@{}>{\raggedright\arraybackslash}p{68.97mm}>{\raggedright\arraybackslash}p{52.22mm}@{}}
\toprule
\textbf{Cell} & \textbf{Form} \\
\midrule
Present participle base & forceorfende \\
Past participle base & forcorfen \\
\bottomrule
\end{tabular}
\par\endgroup

% mdtable:T37:end
% teaching: NEW-U14-VOC-WEORPAN
% teaching: NEW-U14-PAR-WEORPAN
% teaching: NEW-U14-PAR-WEORPAN-PART
% teaching: NEW-U14-GRAM-U14-WEORPAN-PARTICIPLE-BASES
\reserve{10\baselineskip}
% mdparagraph:P079:begin
\textbf{weorpan.} “Throw”. Present weorp-; strong past worp- + -en. Selected worpen is valid without ge-; geworpen is an accepted alternative.
% mdparagraph:P079:end

% table: NEW-U14-PAR-WEORPAN-PART
% mdtable:T38:begin
\par\begingroup
\fontsize{10.5}{13.2}\selectfont
\setlength{\LTleft}{0pt}\setlength{\LTright}{\fill}
\noindent
\begin{tabular}{@{}>{\raggedright\arraybackslash}p{74.31mm}>{\raggedright\arraybackslash}p{46.88mm}@{}}
\toprule
\textbf{Cell} & \textbf{Form} \\
\midrule
Present participle base & weorpende \\
Past participle base & worpen \\
\bottomrule
\end{tabular}
\par\endgroup

% mdtable:T38:end
% teaching: NEW-U14-VOC-GEMETAN
% teaching: NEW-U14-PAR-GEMETAN
% teaching: NEW-U14-PAR-GEMETAN-PART
% teaching: NEW-U14-GRAM-U14-GEMETAN-PARTICIPLE-BASES
\reserve{9\baselineskip}
% mdparagraph:P080:begin
\textbf{gemētan.} “Find; meet”. Present gemēt-; the dental past participle contracts to gemētt. Retain the lexical ge- once.
% mdparagraph:P080:end

% table: NEW-U14-PAR-GEMETAN-PART
% mdtable:T39:begin
\par\begingroup
\fontsize{10.5}{13.2}\selectfont
\setlength{\LTleft}{0pt}\setlength{\LTright}{\fill}
\noindent
\begin{tabular}{@{}>{\raggedright\arraybackslash}p{74.31mm}>{\raggedright\arraybackslash}p{46.88mm}@{}}
\toprule
\textbf{Cell} & \textbf{Form} \\
\midrule
Present participle base & gemētende \\
Past participle base & gemētt \\
\bottomrule
\end{tabular}
\par\endgroup

% mdtable:T39:end
% teaching: NEW-U14-VOC-GEGLENGAN
% teaching: NEW-U14-PAR-GEGLENGAN
% teaching: NEW-U14-PAR-GEGLENGAN-PART
% teaching: NEW-U14-GRAM-U14-GEGLENGAN-PARTICIPLE-BASES
\reserve{9\baselineskip}
% mdparagraph:P081:begin
\textbf{geglęngan.} “Adorn”. Present geglęng-; weak past base geglęnged. Do not add a second ge-.
% mdparagraph:P081:end

% table: NEW-U14-PAR-GEGLENGAN-PART
% mdtable:T40:begin
\par\begingroup
\fontsize{10.5}{13.2}\selectfont
\setlength{\LTleft}{0pt}\setlength{\LTright}{\fill}
\noindent
\begin{tabular}{@{}>{\raggedright\arraybackslash}p{71.54mm}>{\raggedright\arraybackslash}p{49.65mm}@{}}
\toprule
\textbf{Cell} & \textbf{Form} \\
\midrule
Present participle base & geglęngende \\
Past participle base & geglęnged \\
\bottomrule
\end{tabular}
\par\endgroup

% mdtable:T40:end
% mdparagraph:P082:begin
With the earlier passive construction, the person or thing affected is the subject and the finite form of bēon or wesan carries tense. Coordinated past participles may share one auxiliary. An instruction “verb last” in Cook 26.7 or 27.2 selects a final finite auxiliary after the participle group and its complements; it does not merely put a nonfinite participle last.
% mdparagraph:P082:end

\reserve{7\baselineskip}
% mdparagraph:P083:begin
\subsection*{Commands}\BookHeading{nav-ch14-learner-0021}{subsubsection}{Commands}{0}
% mdparagraph:P083:end

% mdparagraph:P084:begin
These weak-II commands have singular -a and plural -iað. The noun or pronoun after the verb takes the case required by its own construction. Preserve ts in gemiltsian and the single l in tilian.
% mdparagraph:P084:end

% teaching: NEW-U14-VOC-GEMILTSIAN
% teaching: NEW-U14-PAR-GEMILTSIAN
% teaching: NEW-U14-PAR-GEMILTSIAN-IMP
% teaching: NEW-U14-GRAM-U14-GEMILTSIAN-IMPERATIVE
\reserve{9\baselineskip}
% mdparagraph:P085:begin
\textbf{gemiltsian} — have compassion; have mercy.
% mdparagraph:P085:end

% table: NEW-U14-PAR-GEMILTSIAN-IMP
% mdtable:T41:begin
% course-table-style: 2026-10
% course-panel:T41-P01:begin
\par\begingroup
\fontsize{10.5}{13.2}\selectfont
\setlength{\LTleft}{0pt}\setlength{\LTright}{\fill}
\noindent
\begin{tabular}{@{}>{\raggedright\arraybackslash}p{63.342mm}>{\raggedright\arraybackslash}p{57.846mm}@{}}
\toprule
\textbf{Person and number} & \textbf{Command} \\
\midrule
2 sg. & gemiltsa \\
2 pl. & gemiltsiað \\
\bottomrule
\end{tabular}
\par\endgroup
% course-panel:T41-P01:end
% mdtable:T41:end
% teaching: NEW-U14-VOC-TILIAN
% teaching: NEW-U14-PAR-TILIAN
% teaching: NEW-U14-PAR-TILIAN-IMP
% teaching: NEW-U14-GRAM-U14-TILIAN-IMPERATIVE
\reserve{9\baselineskip}
% mdparagraph:P086:begin
\textbf{tilian} — obtain; provide.
% mdparagraph:P086:end

% table: NEW-U14-PAR-TILIAN-IMP
% mdtable:T42:begin
% course-table-style: 2026-10
% course-panel:T42-P01:begin
\par\begingroup
\fontsize{10.5}{13.2}\selectfont
\setlength{\LTleft}{0pt}\setlength{\LTright}{\fill}
\noindent
\begin{tabular}{@{}>{\raggedright\arraybackslash}p{63.342mm}>{\raggedright\arraybackslash}p{57.846mm}@{}}
\toprule
\textbf{Person and number} & \textbf{Command} \\
\midrule
2 sg. & tila \\
2 pl. & tiliað \\
\bottomrule
\end{tabular}
\par\endgroup
% course-panel:T42-P01:end
% mdtable:T42:end
\reserve{7\baselineskip}
% mdparagraph:P087:begin
\section*{5. Meanings and ordinary uses}\BookHeading{nav-ch14-learner-0022}{subsection}{5. Meanings and ordinary uses}{1}
% mdparagraph:P087:end

% mdparagraph:P088:begin
A\BookPoint{bk-ch14-learner-ba4182dd4151} verb’s forms and its government are separate facts. The following uses apply with the systems assigned above and earlier. In particular, prefixes do not preserve a base verb’s government automatically.
% mdparagraph:P088:end

% teaching: NEW-U14-VOC-FORLAETAN
% teaching: NEW-U14-VOC-ON
% teaching: NEW-U14-PAR-ON
% teaching: NEW-U14-VOC-MID
% teaching: NEW-U14-PAR-MID
% teaching: NEW-U14-GRAM-U14-GEWENDAN-RETURN
% teaching: NEW-U14-GRAM-U14-AHREOSAN-ON
% teaching: NEW-U14-GRAM-U14-GAN-ON-GOAL
% teaching: NEW-U14-GRAM-U14-BERAN-TRANSPORT
% teaching: NEW-U14-GRAM-U14-ARISAN-SOURCE
% mdparagraph:P089:begin
\textbf{Motion and destination.} Gewęndan, “return”, is used intransitively with tō + dative for the endpoint. Āhrēosan, “fall”, takes on + accusative when falling upon a goal. Gān takes on + accusative for entry or ascent into or onto a place. Beran carries an accusative person or thing to a destination marked tō + dative. Ārīsan is intransitive; of + dative marks the source from which something arises.
% mdparagraph:P089:end

% teaching: NEW-U14-GRAM-U14-GELIEFAN-ON
% mdparagraph:P090:begin
\textbf{Belief.} Gelīefan on + accusative person means “believe in” that person. This is a lexical use of on, not a physical-motion construction. Gelīefan may also stand without an expressed complement, “believe”.
% mdparagraph:P090:end

% teaching: NEW-U14-GRAM-U14-FORLAETAN-TRANSFER
% teaching: NEW-U14-GRAM-U14-FORGIEFAN-GIVING
% teaching: NEW-U14-GRAM-U14-AETIEWAN-SHOW
% teaching: NEW-U14-GRAM-U14-GIEFAN-UNEXPRESSED-THING
% mdparagraph:P091:begin
\textbf{Giving, leaving and showing.} Forgiefan here means “give; grant”, with an accusative thing and a dative recipient. Existing forlǣtan has the same cases in the selected expression “leave peace to someone”; its earlier abandonment meaning alone did not teach this use. Ætīewan shows an accusative thing to a dative person. With giefan the thing given may be left unexpressed when the context supplies it; the recipient remains dative.
% mdparagraph:P091:end

% teaching: NEW-U14-GRAM-U14-GEBIDDAN-REFLEXIVE
% mdparagraph:P092:begin
\textbf{Prayer.}\BookPoint{bk-ch14-learner-61ffa9daf947} Gebiddan uses an accusative personal pronoun referring to its subject. Keep that reflexive in both the purpose clause and the factual when-clause of Cook 26.8. English normally translates the whole construction simply as “pray”. The purpose clause uses subjunctive; the factual past when-clause uses indicative.
% mdparagraph:P092:end

% teaching: NEW-U14-GRAM-U14-LICCIAN-OBJECT
% teaching: NEW-U14-GRAM-U14-AHEBBAN-OBJECT
% mdparagraph:P093:begin
\textbf{Affected objects.} Liccian, “lick”, takes an accusative object. Āhębban, “raise; lift”, also takes an accusative object; ūp may be added as a separate adverb.
% mdparagraph:P093:end

% teaching: NEW-U14-GRAM-U14-WEORPAN-GOAL-PASSIVE
% teaching: NEW-U14-GRAM-U14-GEGLENGAN-MATERIAL
% mdparagraph:P094:begin
\textbf{Passive complements.} Weorpan throws an accusative thing into or onto an on + accusative goal. Its passive retains that goal. Geglęngan adorns a person; the passive makes that person the subject and may add mid + dative for the material or clothing.
% mdparagraph:P094:end

% teaching: NEW-U14-GRAM-U14-ONFON-DATIVE
% teaching: NEW-U14-GRAM-U14-GEMILTSIAN-DATIVE
% mdparagraph:P095:begin
\textbf{Dative government.}\BookPoint{bk-ch14-learner-48f75fc95810} Onfōn, “receive”, takes the dative thing received. This differs from the earlier accusative after fōn, “catch”. Gemiltsian, “have mercy on”, takes the dative person.
% mdparagraph:P095:end

% teaching: NEW-U14-GRAM-U14-TILIAN-GENITIVE-REFLEXIVE
% mdparagraph:P096:begin
\textbf{Providing friends.}\BookPoint{bk-ch14-learner-1ec6cf84d79b} In the selected tilian construction the friends acquired are genitive; the personal beneficiary “for yourselves” is a dative reflexive. Both case requirements must be expressed. Singular and plural commands use the same government.
% mdparagraph:P096:end

% mdparagraph:P097:begin
\textbf{Earlier clause uses.} The bare relative particle ðe can introduce a relative clause without a preceding demonstrative inside that clause. The noun’s role in its main clause need not match the relative’s role. The purpose clauses in Cook 26.2 use subjunctives without an added “may” auxiliary. For Cook 27.9, use the earlier hālsian request construction: the addressee is accusative, and a following ðæt-clause expresses the requested action in the subjunctive. Sęndan retains its earlier accusative object and dative recipient where a recipient is expressed.
% mdparagraph:P097:end

\reserve{8\baselineskip}
% example: NEW-U14-EXAMPLE01
\reserve{5\baselineskip}
% mdparagraph:P098:begin
\textbf{Example 1.}
% mdparagraph:P098:end

\reserve{5\baselineskip}
% mdparagraph:P099:begin
\textbf{Ðæt goldhord is on ðǣm hūse.}
% mdparagraph:P099:end

% mdparagraph:P100:begin
The treasure is in the house.
% mdparagraph:P100:end

% mdparagraph:P101:begin
On + dative gives the location of the treasure. The compound is one neuter noun.
% mdparagraph:P101:end

\reserve{8\baselineskip}
% example: NEW-U14-EXAMPLE02
\reserve{5\baselineskip}
% mdparagraph:P102:begin
\textbf{Example 2.}
% mdparagraph:P102:end

\reserve{5\baselineskip}
% mdparagraph:P103:begin
\textbf{Ic forlǣte ēow sibbe.}
% mdparagraph:P103:end

% mdparagraph:P104:begin
I leave you peace.
% mdparagraph:P104:end

% mdparagraph:P105:begin
Sibbe is the thing left; ēow names the plural recipients.
% mdparagraph:P105:end

\reserve{8\baselineskip}
% example: NEW-U14-EXAMPLE03
\reserve{5\baselineskip}
% mdparagraph:P106:begin
\textbf{Example 3.}
% mdparagraph:P106:end

\reserve{5\baselineskip}
% mdparagraph:P107:begin
\textbf{Hē onfēng ðǣm golde.}
% mdparagraph:P107:end

% mdparagraph:P108:begin
He received the gold.
% mdparagraph:P108:end

% mdparagraph:P109:begin
The thing received is dative after onfōn.
% mdparagraph:P109:end

\reserve{8\baselineskip}
% example: NEW-U14-EXAMPLE04
\reserve{5\baselineskip}
% mdparagraph:P110:begin
\textbf{Example 4.}
% mdparagraph:P110:end

\reserve{5\baselineskip}
% mdparagraph:P111:begin
\textbf{Gemiltsa mē.}
% mdparagraph:P111:end

% mdparagraph:P112:begin
Have mercy on me.
% mdparagraph:P112:end

% mdparagraph:P113:begin
The singular command governs dative mē.
% mdparagraph:P113:end

\reserve{7\baselineskip}
% mdparagraph:P114:begin
\section*{6. Practice in forms and word families}\BookHeading{nav-ch14-learner-0023}{subsection}{6. Practice in forms and word families}{1}
% mdparagraph:P114:end

% mdparagraph:P115:begin
Give the forms or explanations requested. For analysis tasks use the supplied citation meanings and source forms. Their use here does not require you to produce unassigned forms from memory.
% mdparagraph:P115:end

\par\begin{minipage}{\linewidth}\raggedright
% mdparagraph:P116:begin
\subsection*{F01}\BookHeading{nav-ch14-learner-0024}{subsubsection}{F01}{0}
% mdparagraph:P116:end

% item: NEW-U14-F01
% mdparagraph:P117:begin
Give nominative, genitive, dative and accusative, singular and plural, of wǣdla and ēage. Explain the masculine/\allowbreak{}neuter contrast in the singular and which cells ending in -an require context to distinguish them.
% mdparagraph:P117:end

\end{minipage}\par\addvspace{4pt}

\par\begin{minipage}{\linewidth}\raggedright
% mdparagraph:P118:begin
\subsection*{F02}\BookHeading{nav-ch14-learner-0025}{subsubsection}{F02}{0}
% mdparagraph:P118:end

% item: NEW-U14-F02
% mdparagraph:P119:begin
Give the complete nominative, genitive, dative and accusative singular/\allowbreak{}plural of goldhord. Give the dative singular of wyrtgemang and gōdwębb. Explain where the case ending is attached in these compounds.
% mdparagraph:P119:end

\end{minipage}\par\addvspace{4pt}

\par\begin{minipage}{\linewidth}\raggedright
% mdparagraph:P120:begin
\subsection*{F03}\BookHeading{nav-ch14-learner-0026}{subsubsection}{F03}{0}
% mdparagraph:P120:end

% item: NEW-U14-F03
% mdparagraph:P121:begin
Give all six present and all six preterite indicative persons of liccian. Explain the present stems and the preterite stem shown by the paradigm.
% mdparagraph:P121:end

\end{minipage}\par\addvspace{4pt}

\par\begin{minipage}{\linewidth}\raggedright
% mdparagraph:P122:begin
\subsection*{F04}\BookHeading{nav-ch14-learner-0027}{subsubsection}{F04}{0}
% mdparagraph:P122:end

% item: NEW-U14-F04
% mdparagraph:P123:begin
Give these indicative forms: gewęndan present 2 singular; gelīefan present 3 singular; forgiefan present 3 singular; āhrēosan preterite 3 plural; gebiddan present 3 singular; ætīewan preterite 3 plural; āhębban present 3 singular; onfōn preterite 2 singular. Explain the stem choice in the two prefixed strong verbs āhrēosan and onfōn.
% mdparagraph:P123:end

\end{minipage}\par\addvspace{4pt}

\par\begin{minipage}{\linewidth}\raggedright
% mdparagraph:P124:begin
\subsection*{F05}\BookHeading{nav-ch14-learner-0028}{subsubsection}{F05}{0}
% mdparagraph:P124:end

% item: NEW-U14-F05
% mdparagraph:P125:begin
Give the complete present and preterite subjunctive of gebiddan. Then give present subjunctive 3 singular of forweorðan, 2 singular of sęndan, 3 plural of gelīefan and 3 singular of ārīsan.
% mdparagraph:P125:end

\end{minipage}\par\addvspace{4pt}

\par\begin{minipage}{\linewidth}\raggedright
% mdparagraph:P126:begin
\subsection*{F06}\BookHeading{nav-ch14-learner-0029}{subsubsection}{F06}{0}
% mdparagraph:P126:end

% item: NEW-U14-F06
% mdparagraph:P127:begin
Give present and past participle bases of forceorfan, weorpan, gemētan and geglęngan. Give singular and plural commands of gemiltsian and tilian.
% mdparagraph:P127:end

\end{minipage}\par\addvspace{4pt}

\par\begin{minipage}{\linewidth}\raggedright
% mdparagraph:P128:begin
\subsection*{A01}\BookHeading{nav-ch14-learner-0030}{subsubsection}{A01}{0}
% mdparagraph:P128:end

% item: NEW-U14-A01
% mdparagraph:P129:begin
Analyse goldhord and the source forms heofonfȳre and hēaclifum. For each, identify the two components and the final noun, separate any inflectional ending, and give the whole expression’s contextual meaning. Explain why a compound’s exact meaning must still be learned.
% mdparagraph:P129:end

% mdparagraph:P130:begin
\textbf{Inputs, in order:} goldhord · heofonfȳre · hēaclifum
% mdparagraph:P130:end

\end{minipage}\par\addvspace{4pt}

% teaching: U14-DESIGN-COMPOUND-BASES
% support: U14-DESIGN-COMPOUND-BASES
% mdparagraph:P131:begin
\textbf{Supplied help.} For this analysis, hord means treasure/\allowbreak{}hoard and clif means cliff. Heofonfȳre means fire from heaven (here lightning), dative singular; hēaclifum means high cliffs, dative plural. These source forms and component citations are supplied; do not reconstruct further forms from them.
% mdparagraph:P131:end

\par\begin{minipage}{\linewidth}\raggedright
% mdparagraph:P132:begin
\subsection*{A02}\BookHeading{nav-ch14-learner-0031}{subsubsection}{A02}{0}
% mdparagraph:P132:end

% item: NEW-U14-A02
% mdparagraph:P133:begin
Compare leornian/\allowbreak{}leornung, gefēra/\allowbreak{}gefērscipe, rihtwīs/\allowbreak{}rihtwīsnes and īdel/\allowbreak{}īdelnes. Identify the noun-forming suffix in each family, state the parts of speech and give the taught meanings. Explain why īdelnes cannot be translated merely by adding an English suffix to the previously taught meaning of īdel.
% mdparagraph:P133:end

% mdparagraph:P134:begin
\textbf{Inputs, in order:} leornian · leornung · gefēra · gefērscipe · rihtwīs · rihtwīsnes · īdel · īdelnes
% mdparagraph:P134:end

\end{minipage}\par\addvspace{4pt}

% teaching: U14-DESIGN-LEORNIAN-CITATION
% support: U14-DESIGN-LEORNIAN-CITATION
% mdparagraph:P135:begin
\textbf{Supplied help.} Leornian is the supplied infinitive citation of the already taught verb learn. This task analyses its family; it does not require independent infinitive recall or use.
% mdparagraph:P135:end

\par\begin{minipage}{\linewidth}\raggedright
% mdparagraph:P136:begin
\subsection*{A03}\BookHeading{nav-ch14-learner-0032}{subsubsection}{A03}{0}
% mdparagraph:P136:end

% item: NEW-U14-A03
% mdparagraph:P137:begin
Compare heofon/\allowbreak{}heofonlic, swīð/\allowbreak{}swīðe and sōð/\allowbreak{}sōðlice. Identify what forms an adjective and what forms an adverb; give each word’s supplied or taught meaning. Does the final e alone prove that a word is an adverb? Explain.
% mdparagraph:P137:end

% mdparagraph:P138:begin
\textbf{Inputs, in order:} heofon · heofonlic · swīð · swīðe · sōð · sōðlice
% mdparagraph:P138:end

\end{minipage}\par\addvspace{4pt}

% teaching: U14-DESIGN-ADJECTIVE-CITATIONS
% support: U14-DESIGN-ADJECTIVE-CITATIONS
% mdparagraph:P139:begin
\textbf{Supplied help.} Swīð strong and sōð true are supplied citation bases, as in the earlier adverb lesson. No complete adjective systems or independent use of these bases are assigned.
% mdparagraph:P139:end

\par\begin{minipage}{\linewidth}\raggedright
% mdparagraph:P140:begin
\subsection*{A04}\BookHeading{nav-ch14-learner-0033}{subsubsection}{A04}{0}
% mdparagraph:P140:end

% item: NEW-U14-A04
% mdparagraph:P141:begin
Compare rihtwīsnes/\allowbreak{}unrihtwīsnes, berstan/\allowbreak{}tōberstan and weorðan/\allowbreak{}forweorðan. Identify each added prefix and compare the taught whole-word meanings. Explain why the pairs do not license either free coinage or treating every prefixed verb as its unprefixed base.
% mdparagraph:P141:end

% mdparagraph:P142:begin
\textbf{Inputs, in order:} rihtwīsnes · unrihtwīsnes · berstan · tōberstan · weorðan · forweorðan
% mdparagraph:P142:end

\end{minipage}\par\addvspace{4pt}

% teaching: U14-DESIGN-PREFIX-CITATIONS
% support: U14-DESIGN-PREFIX-CITATIONS
% mdparagraph:P143:begin
\textbf{Supplied help.} The infinitive labels berstan, tōberstan, weorðan and forweorðan are supplied here as dictionary citations of already taught finite verbs. Analyse the labels; no independent infinitive recall or nonfinite use is required.
% mdparagraph:P143:end

\par\begin{minipage}{\linewidth}\raggedright
% mdparagraph:P144:begin
\subsection*{A05}\BookHeading{nav-ch14-learner-0034}{subsubsection}{A05}{0}
% mdparagraph:P144:end

% item: NEW-U14-A05
% mdparagraph:P145:begin
In the family rihtwīs → rihtwīsnes → rihtwīsnesse, distinguish derivation from inflection. Parse rihtwīsnesse in Ic lufie rihtwīsnesse. Would the isolated spelling rihtwīsnesse determine its case? Explain.
% mdparagraph:P145:end

% mdparagraph:P146:begin
\textbf{Inputs, in order:} Ic lufie rihtwīsnesse. · rihtwīs · rihtwīsnes
% mdparagraph:P146:end

\end{minipage}\par\addvspace{4pt}

\par\begin{minipage}{\linewidth}\raggedright
% mdparagraph:P147:begin
\subsection*{A06}\BookHeading{nav-ch14-learner-0035}{subsubsection}{A06}{0}
% mdparagraph:P147:end

% item: NEW-U14-A06
% mdparagraph:P148:begin
Translate Hē onfēng gōdum and Ic geseah ðā rihtwīsan. Parse gōdum as the assigned lexical noun and rihtwīsan as a substantival adjective. Explain each word’s case and the adjective’s agreement. Does the absence of an expressed following noun by itself settle a word’s lexical category?
% mdparagraph:P148:end

% mdparagraph:P149:begin
\textbf{Inputs, in order:} Hē onfēng gōdum. · Ic geseah ðā rihtwīsan.
% mdparagraph:P149:end

\end{minipage}\par\addvspace{4pt}

\reserve{7\baselineskip}
% mdparagraph:P150:begin
\section*{7. Short compositions}\BookHeading{nav-ch14-learner-0036}{subsection}{7. Short compositions}{1}
% mdparagraph:P150:end

\par\begin{minipage}{\linewidth}\raggedright
% mdparagraph:P151:begin
\subsection*{C01}\BookHeading{nav-ch14-learner-0037}{subsubsection}{C01}{0}
% mdparagraph:P151:end

% item: NEW-U14-C01
% mdparagraph:P152:begin
Translate: The poor men saw the dogs and the wounds. Use wǣdla for poor man.
% mdparagraph:P152:end

\end{minipage}\par\addvspace{4pt}

\par\begin{minipage}{\linewidth}\raggedright
% mdparagraph:P153:begin
\subsection*{C02}\BookHeading{nav-ch14-learner-0038}{subsubsection}{C02}{0}
% mdparagraph:P153:end

% item: NEW-U14-C02
% mdparagraph:P154:begin
Translate: I leave you [plural] peace, and I give the poor man the treasure. Use forlǣtan, forgiefan, ðearfa and goldhord.
% mdparagraph:P154:end

\end{minipage}\par\addvspace{4pt}

\par\begin{minipage}{\linewidth}\raggedright
% mdparagraph:P155:begin
\subsection*{C03}\BookHeading{nav-ch14-learner-0039}{subsubsection}{C03}{0}
% mdparagraph:P155:end

% item: NEW-U14-C03
% mdparagraph:P156:begin
Translate: The poor men believe in God. Use wǣdla and gelīefan.
% mdparagraph:P156:end

\end{minipage}\par\addvspace{4pt}

\par\begin{minipage}{\linewidth}\raggedright
% mdparagraph:P157:begin
\subsection*{C04}\BookHeading{nav-ch14-learner-0040}{subsubsection}{C04}{0}
% mdparagraph:P157:end

% item: NEW-U14-C04
% mdparagraph:P158:begin
Translate as singular commands: Have mercy on the poor man, and send him the gold. Use gemiltsian, wǣdla and sęndan.
% mdparagraph:P158:end

\end{minipage}\par\addvspace{4pt}

\reserve{7\baselineskip}
% mdparagraph:P159:begin
\section*{8. Cook, Group VIII}\BookHeading{nav-ch14-learner-0041}{subsection}{8. Cook, Group VIII}{1}
% mdparagraph:P159:end

% teaching: NEW-U14-LOCAL-HELPS
% mdparagraph:P160:begin
Translate the complete original prompts. Their parenthetical references are Cook’s grammar section numbers, not sections of this chapter. Modern directions and supplied helps follow separately. “Om.” means that the indicated English word has no separate Old English equivalent. A specified vocabulary or order target must be observed even when another grammatical expression is possible. Keep all help in view; the negative and indefinite expressions below are supplied for these items.
% mdparagraph:P160:end

\reserve{7\baselineskip}
% mdparagraph:P161:begin
\subsection*{Exercise 26}\BookHeading{nav-ch14-learner-0042}{subsubsection}{Exercise 26}{0}
% mdparagraph:P161:end

\par\begin{minipage}{\linewidth}\raggedright
% mdparagraph:P162:begin
\subsection*{Cook VIII.26.1}\BookHeading{nav-ch14-learner-0043}{subsubsection}{Cook VIII.26.1}{0}
% mdparagraph:P162:end

% item: EX-VIII-26-01
% mdparagraph:P163:begin
Thou art dust, and thou returnest (gewendan) to dust.
% mdparagraph:P163:end

\end{minipage}\par\addvspace{4pt}

\par\begin{minipage}{\linewidth}\raggedright
% mdparagraph:P164:begin
\subsection*{Cook VIII.26.2}\BookHeading{nav-ch14-learner-0044}{subsubsection}{Cook VIII.26.2}{0}
% mdparagraph:P164:end

% item: EX-VIII-26-02
% mdparagraph:P165:begin
So much (swīðe) loved God (om. the) world (middangeard), that he gave (sęllan) his only-begotten (āncęnned) Son, that no one (om. may) perish (196. f) that (87. d) believeth in (on, 172) him, but have (196. f) the eternal life.
% mdparagraph:P165:end

\end{minipage}\par\addvspace{4pt}

% teaching: U14-G-COOK-NONE
% support: U14-G-COOK-NONE
% mdparagraph:P166:begin
\textbf{Supplied help.} Nān means “no one” by itself, and “no” in nān mann. Also put ne with the finite verb; the two negatives express one negation.
% mdparagraph:P166:end

% teaching: U14-COOK-LEX-ANCENNED
% support: U14-COOK-LEX-ANCENNED
% mdparagraph:P167:begin
\textbf{Supplied help.} Use āncęnnedan, “only-begotten”: accusative masculine singular weak, agreeing with Sunu after his.
% mdparagraph:P167:end

% mdparagraph:P168:begin
\textbf{Directions.} The references ask for bare relative ðe, on + accusative with belief, and subjunctives in the final purpose. Omit a separate word for “may”.
% mdparagraph:P168:end

\par\begin{minipage}{\linewidth}\raggedright
% mdparagraph:P169:begin
\subsection*{Cook VIII.26.3}\BookHeading{nav-ch14-learner-0045}{subsubsection}{Cook VIII.26.3}{0}
% mdparagraph:P169:end

% item: EX-VIII-26-03
% mdparagraph:P170:begin
Jesus said (cweðan) to him, I am way, and truth, and life; no one (nān) cometh (183) to the Father, except by (ðurh) me.
% mdparagraph:P170:end

% support: U14-G-COOK-NONE
% mdparagraph:P171:begin
\textbf{Supplied help.} Nān means “no one” by itself, and “no” in nān mann. Also put ne with the finite verb; the two negatives express one negation.
% mdparagraph:P171:end

\end{minipage}\par\addvspace{4pt}

% teaching: U14-G-COOK-EXCEPT
% support: U14-G-COOK-EXCEPT
% mdparagraph:P172:begin
\textbf{Supplied help.} Use būtan, “except”, before ðurh + accusative. The supplied spelling ðurh corresponds to earlier þurh. Here būtan introduces the exception; it is not the earlier būtan + dative meaning “without”.
% mdparagraph:P172:end

% teaching: U14-G-COOK-SPEECH-TO
% support: U14-G-COOK-SPEECH-TO
% mdparagraph:P173:begin
\textbf{Supplied help.} Cwæð tō him means “said to him”. The speech addressee is dative after tō. Retain the first person within the direct quotation.
% mdparagraph:P173:end

% teaching: U14-COOK-LEX-WEG
% support: U14-COOK-LEX-WEG
% mdparagraph:P174:begin
\textbf{Supplied help.} Weg here means “way; means of access to the Father”, a figurative sense.
% mdparagraph:P174:end

\par\begin{minipage}{\linewidth}\raggedright
% mdparagraph:P175:begin
\subsection*{Cook VIII.26.4}\BookHeading{nav-ch14-learner-0046}{subsubsection}{Cook VIII.26.4}{0}
% mdparagraph:P175:end

% item: EX-VIII-26-04
% mdparagraph:P176:begin
Peace I leave (om. with) you (164. a); my peace I give (forgiefan) you.
% mdparagraph:P176:end

% mdparagraph:P177:begin
\textbf{Directions.} The recipients are plural and dative. Omit an Old English equivalent of English “with”; use the leave-peace construction in §5.
% mdparagraph:P177:end

\end{minipage}\par\addvspace{4pt}

\par\begin{minipage}{\linewidth}\raggedright
% mdparagraph:P178:begin
\subsection*{Cook VIII.26.5}\BookHeading{nav-ch14-learner-0047}{subsubsection}{Cook VIII.26.5}{0}
% mdparagraph:P178:end

% item: EX-VIII-26-05
% mdparagraph:P179:begin
Then rained it, and there (ðǣr) came floods, and blew winds, and fell (āhrēosan) on that house; and the house fell (feallan).
% mdparagraph:P179:end

\end{minipage}\par\addvspace{4pt}

% teaching: U14-G-COOK-RAIN
% support: U14-G-COOK-RAIN
% mdparagraph:P180:begin
\textbf{Supplied help.} Rīnde is the third singular preterite “rained”. Use impersonal hit, “it”, with this supplied form.
% mdparagraph:P180:end

\par\begin{minipage}{\linewidth}\raggedright
% mdparagraph:P181:begin
\subsection*{Cook VIII.26.6}\BookHeading{nav-ch14-learner-0048}{subsubsection}{Cook VIII.26.6}{0}
% mdparagraph:P181:end

% item: EX-VIII-26-06
% mdparagraph:P182:begin
So every good tree beareth (beran) good fruits (wǣstm), and every evil tree beareth evil fruits.
% mdparagraph:P182:end

% mdparagraph:P183:begin
\textbf{Directions.} C01: retain Cook’s wǣstm in the prompt, but write the earlier selected wæstm, with short æ. See the editorial note.
% mdparagraph:P183:end

\end{minipage}\par\addvspace{4pt}

\par\begin{minipage}{\linewidth}\raggedright
% mdparagraph:P184:begin
\subsection*{Cook VIII.26.7}\BookHeading{nav-ch14-learner-0049}{subsubsection}{Cook VIII.26.7}{0}
% mdparagraph:P184:end

% item: EX-VIII-26-07
% mdparagraph:P185:begin
Every tree that (87. d) beareth not good fruit (sing.) is cut down (forceorfan, 104) and cast into (om. the) fire (verb last).
% mdparagraph:P185:end

% mdparagraph:P186:begin
\textbf{Directions.} Use bare relative ðe. Keep “fruit” singular and omit the article before “fire”. For the stated order target, place the finite passive auxiliary last.
% mdparagraph:P186:end

\end{minipage}\par\addvspace{4pt}

\par\begin{minipage}{\linewidth}\raggedright
% mdparagraph:P187:begin
\subsection*{Cook VIII.26.8}\BookHeading{nav-ch14-learner-0050}{subsubsection}{Cook VIII.26.8}{0}
% mdparagraph:P187:end

% item: EX-VIII-26-08
% mdparagraph:P188:begin
He went into (on, 172) a mountain, that he (om. might) pray (gebiddan, 184. b); and when (ðā) he prayed (184. b), his raiment was white (hwīt), shining.
% mdparagraph:P188:end

\end{minipage}\par\addvspace{4pt}

% teaching: U14-G-COOK-SHINING
% support: U14-G-COOK-SHINING
% mdparagraph:P189:begin
\textbf{Supplied help.} Scīnende, “shining”, is a present participle agreeing with neuter singular rēaf. Keep it in the description of the raiment; no English-style progressive construction is required.
% mdparagraph:P189:end

% mdparagraph:P190:begin
\textbf{Directions.} Reference 184.b requires the accusative reflexive with both forms of gebiddan. The purpose clause and factual when-clause have different moods.
% mdparagraph:P190:end

\par\begin{minipage}{\linewidth}\raggedright
% mdparagraph:P191:begin
\subsection*{Cook VIII.26.9}\BookHeading{nav-ch14-learner-0051}{subsubsection}{Cook VIII.26.9}{0}
% mdparagraph:P191:end

% item: EX-VIII-26-09
% mdparagraph:P192:begin
The mountain, falling, (be)cometh to naught.
% mdparagraph:P192:end

\end{minipage}\par\addvspace{4pt}

% teaching: U14-G-COOK-FALL-NAUGHT
% support: U14-G-COOK-FALL-NAUGHT
% mdparagraph:P193:begin
\textbf{Supplied help.} Hrēosende, “falling”, is nominative masculine singular and describes munt. The fixed expression weorðan tō nāhte means “come to nothing”. Keep the supplied nāhte in this expression.
% mdparagraph:P193:end

\par\begin{minipage}{\linewidth}\raggedright
% mdparagraph:P194:begin
\subsection*{Cook VIII.26.10}\BookHeading{nav-ch14-learner-0052}{subsubsection}{Cook VIII.26.10}{0}
% mdparagraph:P194:end

% item: EX-VIII-26-10
% mdparagraph:P195:begin
The king, in all his glory (wuldor), was not adorned like (swā-swā) one of these.
% mdparagraph:P195:end

\end{minipage}\par\addvspace{4pt}

% teaching: U14-G-COOK-AS
% support: U14-G-COOK-AS
% mdparagraph:P196:begin
\textbf{Supplied help.} Use the connector swā-swā, “like; as”. The following ān + genitive ðissa means “one of these”.
% mdparagraph:P196:end

% teaching: U14-G-COOK-GLORY
% support: U14-G-COOK-GLORY
% mdparagraph:P197:begin
\textbf{Supplied help.} On eallum his wuldre means “in all his glory”. On governs the dative group; eallum agrees with neuter wuldre, while his identifies the possessor. Use ne wæs for “was not”.
% mdparagraph:P197:end

\reserve{7\baselineskip}
% mdparagraph:P198:begin
\subsection*{Exercise 27}\BookHeading{nav-ch14-learner-0053}{subsubsection}{Exercise 27}{0}
% mdparagraph:P198:end

\par\begin{minipage}{\linewidth}\raggedright
% mdparagraph:P199:begin
\subsection*{Cook VIII.27.1}\BookHeading{nav-ch14-learner-0054}{subsubsection}{Cook VIII.27.1}{0}
% mdparagraph:P199:end

% item: EX-VIII-27-01
% mdparagraph:P200:begin
He showed (ætīewan) them his precious things,—the silver, and the gold, and the spices, and all that was found (gemētan) in his treasures.
% mdparagraph:P200:end

% mdparagraph:P201:begin
\textbf{Directions.} The dative recipient differs from the accusative things shown. Strong adjective forms after a possessive remain possible, as taught earlier.
% mdparagraph:P201:end

\end{minipage}\par\addvspace{4pt}

\par\begin{minipage}{\linewidth}\raggedright
% mdparagraph:P202:begin
\subsection*{Cook VIII.27.2}\BookHeading{nav-ch14-learner-0055}{subsubsection}{Cook VIII.27.2}{0}
% mdparagraph:P202:end

% item: EX-VIII-27-02
% mdparagraph:P203:begin
A certain rich man was adorned (geglęngan) with purple (verb last).
% mdparagraph:P203:end

% mdparagraph:P204:begin
\textbf{Directions.} For the stated “verb last” target, put the finite auxiliary after both the material phrase and the participle.
% mdparagraph:P204:end

\end{minipage}\par\addvspace{4pt}

\par\begin{minipage}{\linewidth}\raggedright
% mdparagraph:P205:begin
\subsection*{Cook VIII.27.3}\BookHeading{nav-ch14-learner-0056}{subsubsection}{Cook VIII.27.3}{0}
% mdparagraph:P205:end

% item: EX-VIII-27-03
% mdparagraph:P206:begin
Then (ðā) lay a certain poor man at his gate.
% mdparagraph:P206:end

\end{minipage}\par\addvspace{4pt}

% teaching: U14-G-COOK-AT-GATES
% support: U14-G-COOK-AT-GATES
% mdparagraph:P207:begin
\textbf{Supplied help.} Use æt + dative place: æt his geate, “at his gate”; æt ēowrum gatum, “at your gates”. His names the possessor; ēowrum agrees with plural gatum.
% mdparagraph:P207:end

\par\begin{minipage}{\linewidth}\raggedright
% mdparagraph:P208:begin
\subsection*{Cook VIII.27.4}\BookHeading{nav-ch14-learner-0057}{subsubsection}{Cook VIII.27.4}{0}
% mdparagraph:P208:end

% item: EX-VIII-27-04
% mdparagraph:P209:begin
No man gave (om. unto) him; but dogs (hund, 43) came and licked (liccian, 118) his sores (wund, 43).
% mdparagraph:P209:end

% support: U14-G-COOK-NONE
% mdparagraph:P210:begin
\textbf{Supplied help.} Nān means “no one” by itself, and “no” in nān mann. Also put ne with the finite verb; the two negatives express one negation.
% mdparagraph:P210:end

% mdparagraph:P211:begin
\textbf{Directions.} C04: the noun wund, “wound; sore”, is feminine. Its complete table is in §2; the printed §43 reference is not its feminine declension. The dogs are plural subjects; the wounds belong to the poor man.
% mdparagraph:P211:end

\end{minipage}\par\addvspace{4pt}

\par\begin{minipage}{\linewidth}\raggedright
% mdparagraph:P212:begin
\subsection*{Cook VIII.27.5}\BookHeading{nav-ch14-learner-0058}{subsubsection}{Cook VIII.27.5}{0}
% mdparagraph:P212:end

% item: EX-VIII-27-05
% mdparagraph:P213:begin
Then happened it that the poor man departed (gewītan); and angels bore him to Abraham’s (54) dwelling.
% mdparagraph:P213:end

\end{minipage}\par\addvspace{4pt}

% teaching: U14-G-COOK-HAPPENED
% support: U14-G-COOK-HAPPENED
% mdparagraph:P214:begin
\textbf{Supplied help.} Gelamp means “happened”, singular past. A following ðæt-clause states what happened; optional introductory hit anticipates that clause. The event is factual, so use past indicative in the clause.
% mdparagraph:P214:end

% teaching: U14-COOK-LEX-ABRAHAM
% support: U14-COOK-LEX-ABRAHAM
% mdparagraph:P215:begin
\textbf{Supplied help.} Use Abrahames as the genitive of Abraham. Use unchanged Abraham as the accusative and in direct address.
% mdparagraph:P215:end

\par\begin{minipage}{\linewidth}\raggedright
% mdparagraph:P216:begin
\subsection*{Cook VIII.27.6}\BookHeading{nav-ch14-learner-0059}{subsubsection}{Cook VIII.27.6}{0}
% mdparagraph:P216:end

% item: EX-VIII-27-06
% mdparagraph:P217:begin
Then lifted (āhebban) he his eyes up, and saw Abraham afar off [from afar].
% mdparagraph:P217:end

\end{minipage}\par\addvspace{4pt}

% teaching: U14-G-COOK-AFAR
% support: U14-G-COOK-AFAR
% mdparagraph:P218:begin
\textbf{Supplied help.} Feorran means “from afar”. Its ending already expresses “from”; it needs no additional of.
% mdparagraph:P218:end

% support: U14-COOK-LEX-ABRAHAM
% mdparagraph:P219:begin
\textbf{Supplied help.} Use Abrahames as the genitive of Abraham. Use unchanged Abraham as the accusative and in direct address.
% mdparagraph:P219:end

\par\begin{minipage}{\linewidth}\raggedright
% mdparagraph:P220:begin
\subsection*{Cook VIII.27.7}\BookHeading{nav-ch14-learner-0060}{subsubsection}{Cook VIII.27.7}{0}
% mdparagraph:P220:end

% item: EX-VIII-27-07
% mdparagraph:P221:begin
O father Abraham, have compassion on me, and send Lazarus.
% mdparagraph:P221:end

% support: U14-COOK-LEX-ABRAHAM
% mdparagraph:P222:begin
\textbf{Supplied help.} Use Abrahames as the genitive of Abraham. Use unchanged Abraham as the accusative and in direct address.
% mdparagraph:P222:end

\end{minipage}\par\addvspace{4pt}

% teaching: U14-COOK-LEX-LAZARUS
% support: U14-COOK-LEX-LAZARUS
% mdparagraph:P223:begin
\textbf{Supplied help.} Use Lazarum as accusative singular. Lazaras is accusative plural, “Lazaruses”: people like Lazarus.
% mdparagraph:P223:end

\par\begin{minipage}{\linewidth}\raggedright
% mdparagraph:P224:begin
\subsection*{Cook VIII.27.8}\BookHeading{nav-ch14-learner-0061}{subsubsection}{Cook VIII.27.8}{0}
% mdparagraph:P224:end

% item: EX-VIII-27-08
% mdparagraph:P225:begin
O son, remember that thou didst receive goods in thy life.
% mdparagraph:P225:end

\end{minipage}\par\addvspace{4pt}

% teaching: U14-G-COOK-REMEMBER
% support: U14-G-COOK-REMEMBER
% mdparagraph:P226:begin
\textbf{Supplied help.} Gemun is the singular command “remember”. Its following ðæt-clause recalls an asserted past fact; use indicative onfēnge for “thou didst receive”.
% mdparagraph:P226:end

% teaching: U14-G-COOK-LIFETIME
% support: U14-G-COOK-LIFETIME
% mdparagraph:P227:begin
\textbf{Supplied help.} On ðīnum līfe means “in; during your life”. The whole lifetime phrase is dative singular.
% mdparagraph:P227:end

\par\begin{minipage}{\linewidth}\raggedright
% mdparagraph:P228:begin
\subsection*{Cook VIII.27.9}\BookHeading{nav-ch14-learner-0062}{subsubsection}{Cook VIII.27.9}{0}
% mdparagraph:P228:end

% item: EX-VIII-27-09
% mdparagraph:P229:begin
I beseech thee that thou send him to my father’s house.
% mdparagraph:P229:end

% mdparagraph:P230:begin
\textbf{Directions.} Use the earlier hālsian + ðæt request frame. Distinguish the addressee from the person sent and keep the internal genitive in “my father’s house”.
% mdparagraph:P230:end

\end{minipage}\par\addvspace{4pt}

\par\begin{minipage}{\linewidth}\raggedright
% mdparagraph:P231:begin
\subsection*{Cook VIII.27.10}\BookHeading{nav-ch14-learner-0063}{subsubsection}{Cook VIII.27.10}{0}
% mdparagraph:P231:end

% item: EX-VIII-27-10
% mdparagraph:P232:begin
If they hear not the prophets, they (om. will) not believe, though any one (hwā) arise from death (verb last).
% mdparagraph:P232:end

\end{minipage}\par\addvspace{4pt}

% teaching: U14-G-COOK-INDEFINITE-CONCESSIVE
% support: U14-G-COOK-INDEFINITE-CONCESSIVE
% mdparagraph:P233:begin
\textbf{Supplied help.} Hwā here has the indefinite meaning “anyone”. In the concessive clause use present subjunctive ārīse and obey the requested final-verb order.
% mdparagraph:P233:end

% mdparagraph:P234:begin
\textbf{Directions.} The present belief form can have future reference here; omit an added “will”. The last-verb instruction applies inside the concessive clause.
% mdparagraph:P234:end

\par\begin{minipage}{\linewidth}\raggedright
% mdparagraph:P235:begin
\subsection*{Cook VIII.27.11}\BookHeading{nav-ch14-learner-0064}{subsubsection}{Cook VIII.27.11}{0}
% mdparagraph:P235:end

% item: EX-VIII-27-11
% mdparagraph:P236:begin
Provide for yourselves (184. a) friends (156. n) in (on) God’s poor (ðearfa). Ye have many Lazaruses (Lazaras) now lying at your gates.
% mdparagraph:P236:end

\end{minipage}\par\addvspace{4pt}

% teaching: U14-G-COOK-LYING
% support: U14-G-COOK-LYING
% mdparagraph:P237:begin
\textbf{Supplied help.} Licgende, “lying”, is accusative masculine plural, describing the people denoted by Lazaras. Habbað has these people as its object; the participle describes their situation at the gates.
% mdparagraph:P237:end

% teaching: U14-G-COOK-POOR
% support: U14-G-COOK-POOR
% mdparagraph:P238:begin
\textbf{Supplied help.} In this expression on Godes ðearfum means “in; among God’s poor”. The group after on is dative plural; Godes is its dependent genitive.
% mdparagraph:P238:end

% support: U14-G-COOK-AT-GATES
% mdparagraph:P239:begin
\textbf{Supplied help.} Use æt + dative place: æt his geate, “at his gate”; æt ēowrum gatum, “at your gates”. His names the possessor; ēowrum agrees with plural gatum.
% mdparagraph:P239:end

% support: U14-COOK-LEX-LAZARUS
% mdparagraph:P240:begin
\textbf{Supplied help.} Use Lazarum as accusative singular. Lazaras is accusative plural, “Lazaruses”: people like Lazarus.
% mdparagraph:P240:end

% mdparagraph:P241:begin
\textbf{Directions.} Reference 184.a selects the dative reflexive beneficiary; 156.n selects the genitive object of tilian. The second sentence still addresses more than one person.
% mdparagraph:P241:end

\reserve{7\baselineskip}
% mdparagraph:P242:begin
\section*{9. Reading: The Dangers of Greatness}\BookHeading{nav-ch14-learner-0065}{subsection}{9. Reading: The Dangers of Greatness}{1}
% mdparagraph:P242:end

% teaching: NEW-U14-SOURCE-READING-SUPPORT
% teaching: NEW-U14-SOURCE-HELPS
% teaching: NEW-U14-GRAM-U14-SOURCE-ARGUMENT-STRUCTURE
% teaching: NEW-U14-GRAM-U14-LEXICAL-PARALLELISM
% mdparagraph:P243:begin
The following authentic passage develops a general claim through concrete comparisons and returns to human affairs. Read it as one argument. Repeated words, paired verbs and parallel clauses can reinforce one idea rather than add separate events. Keep track of each subject, its dependent clauses and any pronoun referring back to it.
% mdparagraph:P243:end

% mdparagraph:P244:begin
Cook heads the reading “From Wulfstan’s Homilies, No. 49”, but his introduction questions the attribution. His full qualification is retained in the source notes. The assigned opening ends with tō nāhte; the following argument is not assigned. Page and prose-line divisions and apparatus cues are removed from the reading paragraph below. The linguistic spelling, including source hyphens and vowel marks, is retained; the single page division within geslægene is joined. Note attachments are listed after the reading questions.\textsuperscript{6}
% mdparagraph:P244:end

\par\begin{minipage}{\linewidth}\raggedright
% source: FB-VIII-DANGERS-OPENING-ARGUMENT
% mdparagraph:P245:begin
\readingtext{Ǣghwilc hēah ār, hēr on worulde, bið mid frēcnessum ymbseald; efne swā ðā woruldgeðyngða bēoð māran, swā ðā frēcnessa bēoð swīðran. Swā wē magon, be ðǣm, ðā bȳsena oncnāwan and ongietan. Ðæt trēow, ðonne, ðe wiexð on ðǣm wudubearwe, ðæt hit hlīfað ūp ofer eall ðā ōðru trēowu and brǣt hit, ðonne sęmninga storm gestęnt, and se stranga wind, ðonne bið hit swīðlicor gewǣged and geswęnged ðonne se ōðer wudu. Swā bið ēac gelīce be ðǣm hēaclifum and torrum, ðonne hīe hlīfiað feorr ūp ofer ðā ōðre eorðan, hīe ðonne sęmninga feallan onginnað, and full ðearlice hrēosan tō eorðan. Swilce ēac be ðǣm hēagum muntum and dūnum, ðā-ðe hēah standað ofer ealne middangeard, ðā-hwæðre wīte habbað ðæs ealdordōmes, ðæt hīe bēoð geneahhe mid heofonfȳre geðrēade and geðrǣste, and mid līegum geslægene. Swā ðā hēan mihta hēr on worulde hrēosað, and feallað, and tō lore weorðað, and ðisse worulde welan weorðað tō sorge, and ðās eorðlican wundor weorðað tō nāhte.}
% mdparagraph:P245:end

\end{minipage}\par
\reserve{7\baselineskip}
% mdparagraph:P246:begin
\subsection*{Vocabulary supplied for this reading}\BookHeading{nav-ch14-learner-0066}{subsubsection}{Vocabulary supplied for this reading}{0}
% mdparagraph:P246:end

% mdparagraph:P247:begin
These meanings and source-specific forms remain visible for S01–S07. They do not introduce complete new paradigms. Nouns and adjectives not marked as newly productive in §2 are interpreted here through this help.
% mdparagraph:P247:end

% teaching: U14-READ-LEX-AEGHWILC
% support: U14-READ-LEX-AEGHWILC
% mdparagraph:P248:begin
Ǣghwilc — every; each (nominative feminine singular, modifying ār).
% mdparagraph:P248:end

% teaching: U14-READ-LEX-AR
% support: U14-READ-LEX-AR
% mdparagraph:P249:begin
ār — honour; high office or position (nominative feminine singular; subject).
% mdparagraph:P249:end

% teaching: U14-READ-LEX-FRECNES
% support: U14-READ-LEX-FRECNES
% mdparagraph:P250:begin
frēcnessum — danger; peril (dative feminine plural after mid); frēcnessa — danger; peril (nominative feminine plural; subject).
% mdparagraph:P250:end

% teaching: U14-READ-LEX-WORULDGETHYNGTH
% support: U14-READ-LEX-WORULDGETHYNGTH
% mdparagraph:P251:begin
woruldgeðyngða — worldly honour; worldly dignity (nominative feminine plural; subject).
% mdparagraph:P251:end

% teaching: U14-READ-LEX-BYSEN
% support: U14-READ-LEX-BYSEN
% mdparagraph:P252:begin
bȳsena — example; illustration (accusative feminine plural; object of oncnāwan and ongietan).
% mdparagraph:P252:end

% teaching: U14-READ-LEX-WUDUBEARU
% support: U14-READ-LEX-WUDUBEARU
% mdparagraph:P253:begin
wudubearwe — forest; grove (dative masculine singular after on).
% mdparagraph:P253:end

% teaching: U14-READ-LEX-WUDU
% support: U14-READ-LEX-WUDU
% mdparagraph:P254:begin
wudu — wood; forest, collectively (nominative masculine singular in comparative subject phrase).
% mdparagraph:P254:end

% teaching: U14-READ-LEX-HEACLIF
% support: U14-READ-LEX-HEACLIF
% mdparagraph:P255:begin
hēaclifum — high cliff (dative neuter plural after be; transparent hēah + clif compound).
% mdparagraph:P255:end

% teaching: U14-READ-LEX-DUN
% support: U14-READ-LEX-DUN
% mdparagraph:P256:begin
dūnum — mountain; hill (dative feminine plural after be).
% mdparagraph:P256:end

% teaching: U14-READ-LEX-WITE
% support: U14-READ-LEX-WITE
% mdparagraph:P257:begin
wīte — punishment; penalty (accusative neuter singular; object of habbað).
% mdparagraph:P257:end

% teaching: U14-READ-LEX-HEOFONFYR
% support: U14-READ-LEX-HEOFONFYR
% mdparagraph:P258:begin
heofonfȳre — fire from heaven; heavenly fire (dative neuter singular of instrument after mid).
% mdparagraph:P258:end

% teaching: U14-READ-LEX-LIEG
% support: U14-READ-LEX-LIEG
% mdparagraph:P259:begin
līegum — thunderbolt; levin (lightning) (dative masculine plural of instrument after mid).
% mdparagraph:P259:end

% teaching: U14-READ-LEX-LOR
% support: U14-READ-LEX-LOR
% mdparagraph:P260:begin
lore — destruction; ruin (dative neuter singular in supplied tō lore weorðan).
% mdparagraph:P260:end

% teaching: U14-READ-LEX-WELA
% support: U14-READ-LEX-WELA
% mdparagraph:P261:begin
welan — riches; wealth (nominative masculine plural; subject; modified by genitive ðisse worulde).
% mdparagraph:P261:end

% teaching: U14-READ-LEX-SORG
% support: U14-READ-LEX-SORG
% mdparagraph:P262:begin
sorge — distress; anxiety; trouble; sorrow (dative feminine singular in supplied result tō sorge).
% mdparagraph:P262:end

% teaching: U14-READ-LEX-EORTHLIC
% support: U14-READ-LEX-EORTHLIC
% mdparagraph:P263:begin
eorðlican — earthly (nominative neuter plural weak, agreeing with wundor after ðās).
% mdparagraph:P263:end

% teaching: U14-READ-LEX-TREO
% support: U14-READ-LEX-TREO
% mdparagraph:P264:begin
trēow — tree (nominative neuter singular, source -w variant of taught trēo).
% mdparagraph:P264:end

% teaching: U14-READ-LEX-HEAH
% support: U14-READ-LEX-HEAH
% mdparagraph:P265:begin
hēah — high in rank; exalted (nominative feminine singular exact source form, modifying ār; beside earlier feminine hēa); hēagum — high (dative plural weak after ðǣm, modifying coordinated muntum and dūnum; the source retains g); hēan — high in rank; exalted.
% mdparagraph:P265:end

% teaching: U14-READ-LEX-SWITHRA
% support: U14-READ-LEX-SWITHRA
% mdparagraph:P266:begin
swīðran — stronger; fiercer (degree), distinct from earlier right-side sense.
% mdparagraph:P266:end

% teaching: U14-READ-LEX-STORM
% support: U14-READ-LEX-STORM
% mdparagraph:P267:begin
storm — storm (nominative masculine singular, first subject of gestęnt).
% mdparagraph:P267:end

% teaching: U14-READ-LEX-STRANG
% support: U14-READ-LEX-STRANG
% mdparagraph:P268:begin
stranga — strong, here attributive with wind (nominative masculine singular weak after se).
% mdparagraph:P268:end

% teaching: U14-READ-LEX-OTHER
% support: U14-READ-LEX-OTHER
% mdparagraph:P269:begin
ōðer — remaining; the rest of the (forest/\allowbreak{}earth); ōðre — remaining; the rest of the (forest/\allowbreak{}earth).
% mdparagraph:P269:end

% teaching: U14-READ-LEX-TORR
% support: U14-READ-LEX-TORR
% mdparagraph:P270:begin
torrum — crag, the probable sense here in Cook’s note.
% mdparagraph:P270:end

% teaching: U14-READ-LEX-EALDORDOM
% support: U14-READ-LEX-EALDORDOM
% mdparagraph:P271:begin
ealdordōmes — chief place; primacy; eminence (genitive masculine singular dependent on wīte).
% mdparagraph:P271:end

% teaching: U14-READ-LEX-MIHT
% support: U14-READ-LEX-MIHT
% mdparagraph:P272:begin
mihta — powers; those holding high power or station in the concluding human application (nominative feminine plural; subject of hrēosað, feallað and weorðað).
% mdparagraph:P272:end

\reserve{7\baselineskip}
% mdparagraph:P273:begin
\subsection*{Supplied clause help}\BookHeading{nav-ch14-learner-0067}{subsubsection}{Supplied clause help}{0}
% mdparagraph:P273:end

% teaching: U14-G-SRC-01
% support: U14-G-SRC-01
% mdparagraph:P274:begin
\textbf{Clause help.} Ymbseald means “surrounded; beset”; mid introduces the dangers. Efne emphasizes the statement. Swā … māran, swā … swīðran means “the greater …, the fiercer …”. This supplied proportional comparison differs from equal-degree swā … swā.
% mdparagraph:P274:end

% teaching: U14-G-SRC-02
% support: U14-G-SRC-02
% mdparagraph:P275:begin
\textbf{Clause help.} Be ðǣm is a separate dative phrase, “by; from that”, pointing back to the preceding claim. Ðā bȳsena is the accusative object of both oncnāwan and ongietan, infinitives after magon. The paired verbs reinforce one act of understanding.
% mdparagraph:P275:end

% teaching: U14-G-SRC-03
% support: U14-G-SRC-03
% mdparagraph:P276:begin
\textbf{Clause help.} Relative ðe identifies the tree. Wiexð is the source form “grows”, beside earlier wiext. The result clause ðæt hit hlīfað means “so that it towers”; this hit is nominative. Brǣt hit means “spreads itself”: this hit is accusative and refers to the same tree.
% mdparagraph:P276:end

% teaching: U14-G-SRC-04
% support: U14-G-SRC-04
% mdparagraph:P277:begin
\textbf{Clause help.} Sęmninga means “suddenly”; gestęnt, “assails; comes on”. Storm and se stranga wind are its coordinated subjects, although the source retains singular gestęnt. Gewǣged and geswęnged mean “shaken; buffeted”, passive participles. The main-clause ðonne is “then”; the ðonne following swīðlicor, “more violently”, is “than”. Se ōðer wudu means “the rest of the forest”.
% mdparagraph:P277:end

% teaching: U14-G-SRC-HREOSAN-NOTE
% support: U14-G-SRC-HREOSAN-NOTE
% mdparagraph:P278:begin
\textbf{Clause help.} For the question in original note 157.3, the indicative preterite plural of hrēosan is hruron, in all three persons. The plural stem has u and r in place of present ēo and s. This answer is supplied reference, not an additional exercise.
% mdparagraph:P278:end

% teaching: U14-G-SRC-05
% support: U14-G-SRC-05
% mdparagraph:P279:begin
\textbf{Clause help.} Swā bið ēac gelīce be … means “it is likewise also with; concerning …”. Hīe refers to the cliffs and crags. Hlīfiað means “tower”; feorr ūp, “far up”. Onginnað, “begin”, governs both feallan and hrēosan. Full ðearlice means “very violently”; the paired verbs describe one collapse. Tō eorðan gives its destination.
% mdparagraph:P279:end

% teaching: U14-G-SRC-RELATIVE
% support: U14-G-SRC-RELATIVE
% mdparagraph:P280:begin
\textbf{Clause help.} Ðā-ðe is the source’s joined spelling of the plural relative ðā + ðe, “which”. It is nominative plural, subject of standað, and refers to the mountains named by muntum and dūnum.
% mdparagraph:P280:end

% teaching: U14-G-SRC-06
% support: U14-G-SRC-06
% mdparagraph:P281:begin
\textbf{Clause help.} After Swilce ēac supply “it is likewise with; concerning”. Hēah is adverbial “high”; ðā-hwæðre means “nevertheless”; geneahhe, “often”. The ðæt-clause explains the penalty of pre-eminence. Geðrēade, geðrǣste and geslægene are agreeing plural passive participles, “punished; battered; struck”, with heavenly fire and thunderbolts as the instruments.
% mdparagraph:P281:end

% teaching: U14-G-SRC-07
% support: U14-G-SRC-07
% mdparagraph:P282:begin
\textbf{Clause help.} The high powers hrēosað and feallað, two verbs reinforcing their fall. Tō lore weorðan means “come to destruction”; tō sorge, “to sorrow”; tō nāhte, “to nothing”. These fixed tō-result expressions differ from weorðan followed by a nominative predicate.
% mdparagraph:P282:end

% mdparagraph:P283:begin
For the original question in note 157.6, modern English a down is a hill or upland. The adverb down developed from an expression meaning “from a hill”. This historical information is supplied; no additional etymology exercise is assigned.
% mdparagraph:P283:end

\reserve{7\baselineskip}
% mdparagraph:P284:begin
\subsection*{Reading questions}\BookHeading{nav-ch14-learner-0068}{subsubsection}{Reading questions}{0}
% mdparagraph:P284:end

\par\begin{minipage}{\linewidth}\raggedright
% mdparagraph:P285:begin
\subsection*{S01}\BookHeading{nav-ch14-learner-0069}{subsubsection}{S01}{0}
% mdparagraph:P285:end

% item: NEW-U14-S01
% mdparagraph:P286:begin
Translate the complete assigned Dangers of Greatness passage, using its local helps and notes. Preserve each comparison and the final three outcomes.
% mdparagraph:P286:end

\end{minipage}\par\addvspace{4pt}

\par\begin{minipage}{\linewidth}\raggedright
% mdparagraph:P287:begin
\subsection*{S02}\BookHeading{nav-ch14-learner-0070}{subsubsection}{S02}{0}
% mdparagraph:P287:end

% item: NEW-U14-S02
% mdparagraph:P288:begin
State the opening claim, identify the three examples in their order, and explain how the final sentence applies them to human affairs.
% mdparagraph:P288:end

\end{minipage}\par\addvspace{4pt}

\par\begin{minipage}{\linewidth}\raggedright
% mdparagraph:P289:begin
\subsection*{S03}\BookHeading{nav-ch14-learner-0071}{subsubsection}{S03}{0}
% mdparagraph:P289:end

% item: NEW-U14-S03
% mdparagraph:P290:begin
In the tree sentence, identify the referent of each hit, explain the relative and result clauses, distinguish the temporal ðonne … ðonne from the final comparative ðonne, and explain se ōðer wudu.
% mdparagraph:P290:end

\end{minipage}\par\addvspace{4pt}

\par\begin{minipage}{\linewidth}\raggedright
% mdparagraph:P291:begin
\subsection*{S04}\BookHeading{nav-ch14-learner-0072}{subsubsection}{S04}{0}
% mdparagraph:P291:end

% item: NEW-U14-S04
% mdparagraph:P292:begin
In the cliff sentence, identify what hīe refers to and which verb governs feallan and hrēosan. Explain how the infinitives and the adverbial phrase full ðearlice contribute to the single event.
% mdparagraph:P292:end

\end{minipage}\par\addvspace{4pt}

\par\begin{minipage}{\linewidth}\raggedright
% mdparagraph:P293:begin
\subsection*{S05}\BookHeading{nav-ch14-learner-0073}{subsubsection}{S05}{0}
% mdparagraph:P293:end

% item: NEW-U14-S05
% mdparagraph:P294:begin
Analyse ealdordōmes and eorðlican in the source. Separate the derivational element from the inflection, identify each word’s part of speech and contextual role, and explain the whole expression in English. Use the supplied source analyses; no additional forms are required.
% mdparagraph:P294:end

\end{minipage}\par\addvspace{4pt}

% teaching: U14-DESIGN-SOURCE-DERIVATIONS
% support: U14-DESIGN-SOURCE-DERIVATIONS
% mdparagraph:P295:begin
\textbf{Supplied help.} Ealdor is supplied here as chief/\allowbreak{}king; ealdordōm is supremacy or pre-eminence. The source ealdordōmes is genitive singular, dependent on wīte punishment/\allowbreak{}penalty. Eorðlic is earthly; eorðlican is weak nominative neuter plural agreeing with ðās … wundor. These exact contextual analyses and citation bases are supplied for formation analysis, not assigned as complete new paradigms.
% mdparagraph:P295:end

\par\begin{minipage}{\linewidth}\raggedright
% mdparagraph:P296:begin
\subsection*{S06}\BookHeading{nav-ch14-learner-0074}{subsubsection}{S06}{0}
% mdparagraph:P296:end

% item: NEW-U14-S06
% mdparagraph:P297:begin
In the final sentence, list the three subjects with their stated outcomes. Explain why ðisse worulde is genitive while ðās agrees with wundor.
% mdparagraph:P297:end

\end{minipage}\par\addvspace{4pt}

\par\begin{minipage}{\linewidth}\raggedright
% mdparagraph:P298:begin
\subsection*{S07}\BookHeading{nav-ch14-learner-0075}{subsubsection}{S07}{0}
% mdparagraph:P298:end

% item: NEW-U14-S07
% mdparagraph:P299:begin
Choose two pairs or repetitions in the passage that reinforce its argument. Quote the Old English words and explain their effect in context. Do not assume that paired near-synonyms must describe separate events.
% mdparagraph:P299:end

\end{minipage}\par\addvspace{4pt}

\reserve{8\baselineskip}
% mdparagraph:P300:begin
\section*{10. Original source apparatus}\BookHeading{nav-ch14-learner-0076}{subsection}{10. Original source apparatus}{1}
% mdparagraph:P300:end

\begingroup\fontsize{10}{12.8}\selectfont
\setlength{\parskip}{1.5pt}
% mdparagraph:P301:begin
These are Cook’s attached notes, including his questions and literary parallels. They are historical apparatus, not extra assigned exercises. Modern clause help above supplies the needed grammar. Note labels combine the printed page and original note number; prose-line positions in the attachment table are editorial counts anchored by Cook’s printed marginal numbers.
% mdparagraph:P301:end

\reserve{7\baselineskip}
% mdparagraph:P302:begin
\subsection*{Cook’s introductory notice}\BookHeading{nav-ch14-learner-0077}{subsubsection}{Cook’s introductory notice}{0}
% mdparagraph:P302:end

% mdparagraph:P303:begin
VIII. THE DANGERS OF GREATNESS.
% mdparagraph:P303:end

% mdparagraph:P304:begin
(From Wulfstan’s Homilies, No. 49.)
% mdparagraph:P304:end

% mdparagraph:P305:begin
[Wulfstan — also known by his Latinized name, Lupus — was Bishop of Worcester and Archbishop of York from 1002 to 1023. This homily is one of those attributed to him, but, according to Napier, with insufficient reason, as a portion of it is found in the Blickling Homilies, the manuscript of which bears the date of 971.]
% mdparagraph:P305:end

\reserve{7\baselineskip}
% mdparagraph:P306:begin
\subsection*{Note attachments}\BookHeading{nav-ch14-learner-0078}{subsubsection}{Note attachments}{0}
% mdparagraph:P306:end

% mdtable:T43:begin
\par\begingroup
\fontsize{10.5}{13.2}\selectfont
\setlength{\LTleft}{0pt}\setlength{\LTright}{\fill}
\begin{longtable}{@{}>{\raggedright\arraybackslash}p{24.16mm}>{\raggedright\arraybackslash}p{24.16mm}>{\raggedright\arraybackslash}p{70.06mm}@{}}
\toprule
\textbf{Note} & \textbf{Page.line} & \textbf{Immediately after this source text} \\
\midrule
\endfirsthead
\toprule
\textbf{Note} & \textbf{Page.line} & \textbf{Immediately after this source text} \\
\midrule
\endhead
\midrule
\multicolumn{3}{r}{\footnotesize Continued on the following page.}\\
\endfoot
\bottomrule
\endlastfoot
156.1 & 156.1 & Ǣghwilc hēah ār, hēr on worulde, bið mid frēcnessum \\
156.2 & 156.2 & ymbseald \\
156.3 & 156.2 & ymbseald; efne swā \\
156.4 & 156.4 & bȳsena oncnāwan and ongietan. \\
156.5 & 156.5 & wiexð \\
156.6 & 156.5 & wiexð on ðǣm wudubearwe, ðæt \\
156.6a & 156.6 & ðā ōðru trēowu and brǣt \\
156.7 & 156.6 & ðā ōðru trēowu and brǣt hit, \\
156.8 & 156.6 & ðā ōðru trēowu and brǣt hit, ðonne sęmninga storm \\
156.9 & 156.7 & gestęnt, and se stranga wind, \\
156.10 & 156.7 & gestęnt, and se stranga wind, ðonne \\
156.11 & 156.8 & gewǣged and geswęnged ðonne se ōðer wudu. \\
156.12 & 156.9 & ēac gelīce be ðǣm hēaclifum and torrum, \\
157.1 & 157.1 & hlīfiað feorr ūp ofer ðā ōðre \\
157.2 & 157.2 & feallan onginnað, \\
157.3 & 157.2 & feallan onginnað, and full ðearlice hrēosan \\
157.4 & 157.3 & Swilce \\
157.5 & 157.3 & Swilce ēac be ðǣm hēagum \\
157.6 & 157.3 & Swilce ēac be ðǣm hēagum muntum and dūnum, \\
157.7 & 157.6 & heofonfȳre \\
158.1 & 158.1 & slægene. Swā ðā hēan mihta \\
158.2 & 158.2 & and feallað, \\
158.3 & 158.2 & and feallað, and tō lore weorðað, and ðisse \\
158.3 & 158.2 & and feallað, and tō lore weorðað, and ðisse worulde \\
158.4 & 158.4 & weorðað tō nāhte. \\
\end{longtable}
\par\endgroup

% mdtable:T43:end
\reserve{7\baselineskip}
% mdparagraph:P307:begin
\subsection*{156.1}\BookHeading{nav-ch14-learner-0079}{subsubsection}{156.1}{0}
% mdparagraph:P307:end

% source-note: FB-VIII-DANGERS-156-N1
% mdparagraph:P308:begin
See 144.
% mdparagraph:P308:end

\reserve{7\baselineskip}
% mdparagraph:P309:begin
\subsection*{156.2}\BookHeading{nav-ch14-learner-0080}{subsubsection}{156.2}{0}
% mdparagraph:P309:end

% source-note: FB-VIII-DANGERS-156-N2
% mdparagraph:P310:begin
See 114.
% mdparagraph:P310:end

\reserve{7\baselineskip}
% mdparagraph:P311:begin
\subsection*{156.3}\BookHeading{nav-ch14-learner-0081}{subsubsection}{156.3}{0}
% mdparagraph:P311:end

% source-note: FB-VIII-DANGERS-156-N3
% mdparagraph:P312:begin
Swā . . . māran, swā . . . swīðran = the greater, the fiercer. Note the tendency to antithesis.
% mdparagraph:P312:end

\reserve{7\baselineskip}
% mdparagraph:P313:begin
\subsection*{156.4}\BookHeading{nav-ch14-learner-0082}{subsubsection}{156.4}{0}
% mdparagraph:P313:end

% source-note: FB-VIII-DANGERS-156-N4
% mdparagraph:P314:begin
Observe the redundancy.
% mdparagraph:P314:end

\reserve{7\baselineskip}
% mdparagraph:P315:begin
\subsection*{156.5}\BookHeading{nav-ch14-learner-0083}{subsubsection}{156.5}{0}
% mdparagraph:P315:end

% source-note: FB-VIII-DANGERS-156-N5
% mdparagraph:P316:begin
See weaxan.
% mdparagraph:P316:end

\reserve{7\baselineskip}
% mdparagraph:P317:begin
\subsection*{156.6}\BookHeading{nav-ch14-learner-0084}{subsubsection}{156.6}{0}
% mdparagraph:P317:end

% source-note: FB-VIII-DANGERS-156-N6
% mdparagraph:P318:begin
= so that.
% mdparagraph:P318:end

\reserve{7\baselineskip}
% mdparagraph:P319:begin
\subsection*{156.6a}\BookHeading{nav-ch14-learner-0085}{subsubsection}{156.6a}{0}
% mdparagraph:P319:end

% source-note: FB-VIII-DANGERS-156-N6a
% mdparagraph:P320:begin
See brǣdan (34).
% mdparagraph:P320:end

\reserve{7\baselineskip}
% mdparagraph:P321:begin
\subsection*{156.7}\BookHeading{nav-ch14-learner-0086}{subsubsection}{156.7}{0}
% mdparagraph:P321:end

% source-note: FB-VIII-DANGERS-156-N7
% mdparagraph:P322:begin
See 184. b.
% mdparagraph:P322:end

\reserve{7\baselineskip}
% mdparagraph:P323:begin
\subsection*{156.8}\BookHeading{nav-ch14-learner-0087}{subsubsection}{156.8}{0}
% mdparagraph:P323:end

% source-note: FB-VIII-DANGERS-156-N8
% mdparagraph:P324:begin
Note the alliteration.
% mdparagraph:P324:end

\reserve{7\baselineskip}
% mdparagraph:P325:begin
\subsection*{156.9}\BookHeading{nav-ch14-learner-0088}{subsubsection}{156.9}{0}
% mdparagraph:P325:end

% source-note: FB-VIII-DANGERS-156-N9
% mdparagraph:P326:begin
Second subject of gestęnt.
% mdparagraph:P326:end

\reserve{7\baselineskip}
% mdparagraph:P327:begin
\subsection*{156.10}\BookHeading{nav-ch14-learner-0089}{subsubsection}{156.10}{0}
% mdparagraph:P327:end

% source-note: FB-VIII-DANGERS-156-N10
% mdparagraph:P328:begin
Frequently the second correlative, in such pairs as ðonne . . . ðonne, ðā . . . ðā, need not be translated; it is frequently followed by an inverted order, as here, the verb preceding its subject. See 202.
% mdparagraph:P328:end

\reserve{7\baselineskip}
% mdparagraph:P329:begin
\subsection*{156.11}\BookHeading{nav-ch14-learner-0090}{subsubsection}{156.11}{0}
% mdparagraph:P329:end

% source-note: FB-VIII-DANGERS-156-N11
% mdparagraph:P330:begin
Se ōðer wudu = the rest of the forest, not the other wood.
% mdparagraph:P330:end

\reserve{7\baselineskip}
% mdparagraph:P331:begin
\subsection*{156.12}\BookHeading{nav-ch14-learner-0091}{subsubsection}{156.12}{0}
% mdparagraph:P331:end

% source-note: FB-VIII-DANGERS-156-N12
% mdparagraph:P332:begin
Probably here = crag.
% mdparagraph:P332:end

\reserve{7\baselineskip}
% mdparagraph:P333:begin
\subsection*{157.1}\BookHeading{nav-ch14-learner-0092}{subsubsection}{157.1}{0}
% mdparagraph:P333:end

% source-note: FB-VIII-DANGERS-157-N1
% mdparagraph:P334:begin
See p. 156, n. 12.
% mdparagraph:P334:end

\reserve{7\baselineskip}
% mdparagraph:P335:begin
\subsection*{157.2}\BookHeading{nav-ch14-learner-0093}{subsubsection}{157.2}{0}
% mdparagraph:P335:end

% source-note: FB-VIII-DANGERS-157-N2
% mdparagraph:P336:begin
This resembles the use of gin in Chaucer, almost as an auxiliary tense-sign, like do in Mod. Eng., the latter not being thus used in OE. In Chaucer it usually occurs as the preterit gan, e.g. in the Clerk’s Tale, 392: “til the sonne gan descende.” See Lounsbury’s History of the English Language. An interesting parallel is to be found in New Testament Greek, as, for example, Acts 1. 1: “The former treatise I made, O Theophilus, concerning all that Jesus began (ἤρξατο) both to do and to teach.” According to Thayer, however (Greek-English Lexicon of the New Testament), there is in its employment always a sense of beginning, in its proper meaning.
% mdparagraph:P336:end

\reserve{7\baselineskip}
% mdparagraph:P337:begin
\subsection*{157.3}\BookHeading{nav-ch14-learner-0094}{subsubsection}{157.3}{0}
% mdparagraph:P337:end

% source-note: FB-VIII-DANGERS-157-N3
% mdparagraph:P338:begin
Dependent, like feallan, on onginnað. Give the ind. pret. plur.
% mdparagraph:P338:end

\reserve{7\baselineskip}
% mdparagraph:P339:begin
\subsection*{157.4}\BookHeading{nav-ch14-learner-0095}{subsubsection}{157.4}{0}
% mdparagraph:P339:end

% source-note: FB-VIII-DANGERS-157-N4
% mdparagraph:P340:begin
It would be interesting to know from what literary source these illustrations are ultimately derived. They remind one of Shakespeare (Rich. III., 1. 3. 259–260):— They that stand high have many blasts to shake them; And if they fall, they dash themselves to pieces. Cf. also 3 Hen. VI., 5. 2. 11–15. No doubt many Elizabethan parallels could be found; I have noted in Chapman, Byron’s Conspiracy, Act 3, Scene 1 (p. 232 of Shepherd’s ed.), and Byron’s Tragedy, Act 5, Scene 1 (Ib., p. 272). Perhaps the Elizabethans may have derived them from Seneca; cf. the Chorus in Act 4 of the Hippolytus, vv. 1123–1143; Hercules Furens 201; Œdipus 8–11. Seneca may have caught a suggestion from Sophocles, though the parallel is somewhat remote; see the latter’s Antigone, vv. 712–717, and Horace, C. II. 10.
% mdparagraph:P340:end

\reserve{7\baselineskip}
% mdparagraph:P341:begin
\subsection*{157.5}\BookHeading{nav-ch14-learner-0096}{subsubsection}{157.5}{0}
% mdparagraph:P341:end

% source-note: FB-VIII-DANGERS-157-N5
% mdparagraph:P342:begin
See 58. 1.
% mdparagraph:P342:end

\reserve{7\baselineskip}
% mdparagraph:P343:begin
\subsection*{157.6}\BookHeading{nav-ch14-learner-0097}{subsubsection}{157.6}{0}
% mdparagraph:P343:end

% source-note: FB-VIII-DANGERS-157-N6
% mdparagraph:P344:begin
Redundant. What is dūn in Mod. Eng. (24)? Whence is the adverb down derived?
% mdparagraph:P344:end

\reserve{7\baselineskip}
% mdparagraph:P345:begin
\subsection*{157.7}\BookHeading{nav-ch14-learner-0098}{subsubsection}{157.7}{0}
% mdparagraph:P345:end

% source-note: FB-VIII-DANGERS-157-N7
% mdparagraph:P346:begin
Note the poetical term.
% mdparagraph:P346:end

\reserve{7\baselineskip}
% mdparagraph:P347:begin
\subsection*{158.1}\BookHeading{nav-ch14-learner-0099}{subsubsection}{158.1}{0}
% mdparagraph:P347:end

% source-note: FB-VIII-DANGERS-158-N1
% mdparagraph:P348:begin
This suggests Seneca (Œdipus, Act 1, v. 11):— Imperia sic excelsa Fortunae objacent.
% mdparagraph:P348:end

\reserve{7\baselineskip}
% mdparagraph:P349:begin
\subsection*{158.2}\BookHeading{nav-ch14-learner-0100}{subsubsection}{158.2}{0}
% mdparagraph:P349:end

% source-note: FB-VIII-DANGERS-158-N2
% mdparagraph:P350:begin
Pleonastic.
% mdparagraph:P350:end

\reserve{7\baselineskip}
% mdparagraph:P351:begin
\subsection*{158.3}\BookHeading{nav-ch14-learner-0101}{subsubsection}{158.3}{0}
% mdparagraph:P351:end

% source-note: FB-VIII-DANGERS-158-N3
% mdparagraph:P352:begin
Genitive, dependent on welan.
% mdparagraph:P352:end

\reserve{7\baselineskip}
% mdparagraph:P353:begin
\subsection*{158.4}\BookHeading{nav-ch14-learner-0102}{subsubsection}{158.4}{0}
% mdparagraph:P353:end

% source-note: FB-VIII-DANGERS-158-N4
% mdparagraph:P354:begin
Cf. Mod. Eng. come to naught.
% mdparagraph:P354:end

\reserve{7\baselineskip}
% mdparagraph:P355:begin
\section*{11. Editorial and grammatical notes}\BookHeading{nav-ch14-learner-0103}{subsection}{11. Editorial and grammatical notes}{1}
% mdparagraph:P355:end

% mdparagraph:P356:begin
\textbf{C01 — wæstm.} Cook 26.6 prints wǣstm. The selected corrected form is wæstm, with short æ, as in the earlier noun table. Preserve the original prompt while using the corrected course form. Evidence: Baker’s \emph{Old English Aerobics} glossary, s.v. wæstm; Stenroos (2024), printed 118, Table 1, as checked in the earlier lexical audit.
% mdparagraph:P356:end

% mdparagraph:P357:begin
\textbf{C04 — wund.} Cook 27.4 prints “wund, 43”. The wound/\allowbreak{}sore noun is feminine and takes the selected strong-feminine pattern (§51.b, printed 35 /\allowbreak{} PDF 51), whereas §43 presents masculine nouns. Cook’s glossary adjective wund, “wounded”, is a separate word. Bosworth–Toller’s main volume, printed 1281a, identifies the feminine noun; its Supplement, printed 570b, quotes his wunda in Luke 16.21 and corroborates the accusative plural. The original hint remains printed; no explanation of how the error arose is inferred.
% mdparagraph:P357:end

% mdparagraph:P358:begin
\textbf{C05 — the pointer in note 157.1.} The original says “See p. 156, n. 12.” The parallel explanation “the rest of the forest” is in 156.11, while 156.12 concerns “crag”. This is an apparent reference mismatch; the proposed connection with 156.11 is an inference, not Cook’s stated intention. Independently, ðā ōðre eorðan in the passage means “the rest of the earth”.
% mdparagraph:P358:end

% mdparagraph:P359:begin
\textbf{Quantities and lexical choice.} Cook’s glossary prints dūst and gōd-wębb with long ū and ō; the compound divider is omitted in the course headword gōdwębb. At printed 303 /\allowbreak{} one-based PDF 319, līeg is “thunderbolt; levin”. Accordingly līegum in the reading denotes thunderbolts or lightning. Geneahhe retains the reading’s unmarked spelling; it is not silently replaced by the glossary spelling genēahhe.
% mdparagraph:P359:end

% mdparagraph:P360:begin
\textsuperscript{1}: Cook, \emph{First Book in Old English}, second edition, §§142–147, printed 81–87 (one-based PDF 97–103): selected prefixes, noun and adjective suffixes, composition and the limits of mechanical translation. Lexical meanings follow the selected glossary entries.
% mdparagraph:P360:end

% mdparagraph:P361:begin
\textsuperscript{2}: Cook §70, printed 44 (PDF 60), explicitly compares swīð/\allowbreak{}swīðe and sōð/\allowbreak{}sōðlice. The citation adjectives remain supplied bases.
% mdparagraph:P361:end

% mdparagraph:P362:begin
\textsuperscript{3}: Cook §181, printed 100 (PDF 116), on adjectives as nouns; the personal expression with rihtwīs was taught in Unit 11. Lexical gōd, goods, is separately identified in Cook’s vocabulary, printed 296 (PDF 312), and follows the neuter table.
% mdparagraph:P362:end

% mdparagraph:P363:begin
\textsuperscript{4}: Cook §§43–54, printed 27–38 (PDF 43–54), with the checked glossary entries. The selected reduced wuldr- stem follows §§43.4 and 47.6; reference wuldru follows the disyllabic-neuter route in §47.1. Cook §53.2 gives weak-neuter ēage and genitive-plural ēagena/\allowbreak{}ēagna. These rule-based displays do not claim separate corpus attestations for every reference cell.
% mdparagraph:P363:end

% mdparagraph:P364:begin
\textsuperscript{5}: Cook’s selected weak and strong systems, printed 54–74 (PDF 70–90), with each verb’s glossary entry; present consonant treatment and vowel alternations, printed 24–25 (PDF 40–41). Weak-I gemētt and the selected participle forms follow printed 61–62, 68 and 72. Optional ge- in a participle is discussed at printed 82. Government follows the selected lexical entries and §§156, 164, 172, 184 and 196. The locally supplied speech phrase is supported by Cook printed 153 and 159; the lifetime expression by printed 159.
% mdparagraph:P364:end

% mdparagraph:P365:begin
\textsuperscript{6}: Cook, \emph{First Book in Old English}, second edition, “VIII. The Dangers of Greatness”, printed 156.1–158.4, one-based PDF 172–174. Linguistic text is exact; paragraph joining and removal of source cues are declared layout normalization. The selected edition’s attribution and its qualification are reproduced above. Cook, \emph{Exercises in Old English}, Group VIII, Exercises 26–27, supplies the unchanged English prompts.
% mdparagraph:P365:end

\endgroup

% book-source-body:end
\par\label{end-ch14-learner}
\clearpage\endgroup

\BookRecto
\begingroup
\renewcommand{\runningtitle}{Unit 15. Negation and quantity}
\hyphenpenalty=10000
\exhyphenpenalty=10000
\emergencystretch=3em
\newcommand{\itempara}[2]{\par\addvspace{3pt}\noindent\hangindent=1.8em\hangafter=1\makebox[1.8em][l]{#1.}#2\par}
\newcommand{\reserve}[1]{\par\penalty-100\begingroup\dimen0=\pagegoal\advance\dimen0 by-\pagetotal\ifdim\dimen0<#1\newpage\fi\endgroup}
% book-source-body:begin

\addfontfeatures{Ligatures=TeXOff}
\raggedright
\booktitle{\BookComponent{ch15-learner}{Unit 15. Negation, quantity and comparison in discourse}UNIT 15}{Negation, quantity and comparison in discourse}

% mdparagraph:P001:begin
A negative clause can contain more than one negative word without becoming affirmative. Quantity words distinguish an individual from a group, while comparison can express equality, preference or a proportional change. This unit brings these distinctions into connected arguments and requests.
% mdparagraph:P001:end

% mdparagraph:P002:begin
Sections 1–5 give the vocabulary, complete defined systems and constructions needed by the practice. The eight translated examples and the lettered exercises are new teaching material. Cook IX.28–30 and X, XI.31 retain their complete original prompts. The three readings preserve Cook’s Old English and original notes; modern helps are identified separately.
% mdparagraph:P002:end

% mdparagraph:P003:begin
\textbf{Scope and notation.} Learn the forms for use in the meanings and constructions stated. A reference table supplies forms for recognition, without requiring their independent recall. Supplied expressions remain available during their exercises and readings. “Unassigned” means that no form or use in that position is required here; it does not assert that a historical form was impossible. In verbal tables, 1, 2 and 3 denote person. Preserve vowel length and the distinction between e and ę.
% mdparagraph:P003:end

\reserve{7\baselineskip}
% mdparagraph:P004:begin
\section*{1. Nouns and their meanings}\BookHeading{nav-ch15-learner-0001}{subsection}{1. Nouns and their meanings}{1}
% mdparagraph:P004:end

% mdparagraph:P005:begin
The tables give nominative, genitive, dative and accusative in both numbers. Unless a number is marked as reference or unassigned, both numbers may be used. Endings that coincide still have distinct grammatical functions: the sentence determines the case. The lexical meaning, gender and stem belong to the whole word.\textsuperscript{1}
% mdparagraph:P005:end

\reserve{7\baselineskip}
% mdparagraph:P006:begin
\subsection*{Strong masculine nouns}\BookHeading{nav-ch15-learner-0002}{subsubsection}{Strong masculine nouns}{0}
% mdparagraph:P006:end

% mdparagraph:P007:begin
The nominative and accusative singular have no ending. The other singular endings are genitive -es and dative -e. The plural endings are nominative/\allowbreak{}accusative -as, genitive -a and dative -um.
% mdparagraph:P007:end

% teaching: NEW-U15-VOC-HLAFORD
\reserve{11\baselineskip}
% mdparagraph:P008:begin
\textbf{hlāford} — masculine; lord; master. Stem \textbf{hlāford-}.
% mdparagraph:P008:end

% teaching: NEW-U15-NOM-HLAFORD
% table: NEW-U15-NOM-HLAFORD-DISPLAY
% mdtable:T01:begin
\par\begingroup
\fontsize{10.5}{13.2}\selectfont
\setlength{\LTleft}{0pt}\setlength{\LTright}{\fill}
\noindent
\begin{tabular}{@{}>{\raggedright\arraybackslash}p{39.46mm}>{\raggedright\arraybackslash}p{39.46mm}>{\raggedright\arraybackslash}p{39.46mm}@{}}
\toprule
\textbf{Case} & \textbf{Singular} & \textbf{Plural} \\
\midrule
Nominative & hlāford & hlāfordas \\
Genitive & hlāfordes & hlāforda \\
Dative & hlāforde & hlāfordum \\
Accusative & hlāford & hlāfordas \\
\bottomrule
\end{tabular}
\par\endgroup

% mdtable:T01:end
% teaching: NEW-U15-VOC-WINGEARD
\reserve{11\baselineskip}
% mdparagraph:P009:begin
\textbf{wīngeard} — masculine; vineyard. Stem \textbf{wīngeard-}.
% mdparagraph:P009:end

% teaching: NEW-U15-NOM-WINGEARD
% table: NEW-U15-NOM-WINGEARD-DISPLAY
% mdtable:T02:begin
\par\begingroup
\fontsize{10.5}{13.2}\selectfont
\setlength{\LTleft}{0pt}\setlength{\LTright}{\fill}
\noindent
\begin{tabular}{@{}>{\raggedright\arraybackslash}p{39.46mm}>{\raggedright\arraybackslash}p{39.46mm}>{\raggedright\arraybackslash}p{39.46mm}@{}}
\toprule
\textbf{Case} & \textbf{Singular} & \textbf{Plural} \\
\midrule
Nominative & wīngeard & wīngeardas \\
Genitive & wīngeardes & wīngearda \\
Dative & wīngearde & wīngeardum \\
Accusative & wīngeard & wīngeardas \\
\bottomrule
\end{tabular}
\par\endgroup

% mdtable:T02:end
% teaching: NEW-U15-VOC-WER
\reserve{11\baselineskip}
% mdparagraph:P010:begin
\textbf{wer} — masculine; man. Stem \textbf{wer-}.
% mdparagraph:P010:end

% teaching: NEW-U15-NOM-WER
% table: NEW-U15-NOM-WER-DISPLAY
% mdtable:T03:begin
\par\begingroup
\fontsize{10.5}{13.2}\selectfont
\setlength{\LTleft}{0pt}\setlength{\LTright}{\fill}
\noindent
\begin{tabular}{@{}>{\raggedright\arraybackslash}p{39.46mm}>{\raggedright\arraybackslash}p{39.46mm}>{\raggedright\arraybackslash}p{39.46mm}@{}}
\toprule
\textbf{Case} & \textbf{Singular} & \textbf{Plural} \\
\midrule
Nominative & wer & weras \\
Genitive & weres & wera \\
Dative & were & werum \\
Accusative & wer & weras \\
\bottomrule
\end{tabular}
\par\endgroup

% mdtable:T03:end
% teaching: NEW-U15-VOC-LEORNUNGCNIHT
\reserve{11\baselineskip}
% mdparagraph:P011:begin
\textbf{leornungcniht} — masculine; disciple. Stem \textbf{leornungcniht-}.
% mdparagraph:P011:end

% teaching: NEW-U15-NOM-LEORNUNGCNIHT
% table: NEW-U15-NOM-LEORNUNGCNIHT-DISPLAY
% mdtable:T04:begin
\par\begingroup
\fontsize{10.5}{13.2}\selectfont
\setlength{\LTleft}{0pt}\setlength{\LTright}{\fill}
\noindent
\begin{tabular}{@{}>{\raggedright\arraybackslash}p{29.59mm}>{\raggedright\arraybackslash}p{44.39mm}>{\raggedright\arraybackslash}p{44.39mm}@{}}
\toprule
\textbf{Case} & \textbf{Singular} & \textbf{Plural} \\
\midrule
Nominative & leornungcniht & leornungcnihtas \\
Genitive & leornungcnihtes & leornungcnihta \\
Dative & leornungcnihte & leornungcnihtum \\
Accusative & leornungcniht & leornungcnihtas \\
\bottomrule
\end{tabular}
\par\endgroup

% mdtable:T04:end
% teaching: NEW-U15-VOC-LAREOW
\reserve{11\baselineskip}
% mdparagraph:P012:begin
\textbf{lārēow} — masculine; master; teacher. Stem \textbf{lārēow-}.
% mdparagraph:P012:end

% teaching: NEW-U15-NOM-LAREOW
% table: NEW-U15-NOM-LAREOW-DISPLAY
% mdtable:T05:begin
\par\begingroup
\fontsize{10.5}{13.2}\selectfont
\setlength{\LTleft}{0pt}\setlength{\LTright}{\fill}
\noindent
\begin{tabular}{@{}>{\raggedright\arraybackslash}p{39.46mm}>{\raggedright\arraybackslash}p{39.46mm}>{\raggedright\arraybackslash}p{39.46mm}@{}}
\toprule
\textbf{Case} & \textbf{Singular} & \textbf{Plural} \\
\midrule
Nominative & lārēow & lārēowas \\
Genitive & lārēowes & lārēowa \\
Dative & lārēowe & lārēowum \\
Accusative & lārēow & lārēowas \\
\bottomrule
\end{tabular}
\par\endgroup

% mdtable:T05:end
% teaching: NEW-U15-VOC-HAM-NOUN
\reserve{11\baselineskip}
% mdparagraph:P013:begin
\textbf{hām} — masculine; home; dwelling. Stem \textbf{hām-}. This is the noun “home; dwelling”, distinct from the directional adverb hām.
% mdparagraph:P013:end

% teaching: NEW-U15-NOM-HAM-NOUN
% table: NEW-U15-NOM-HAM-NOUN-DISPLAY
% mdtable:T06:begin
\par\begingroup
\fontsize{10.5}{13.2}\selectfont
\setlength{\LTleft}{0pt}\setlength{\LTright}{\fill}
\noindent
\begin{tabular}{@{}>{\raggedright\arraybackslash}p{39.46mm}>{\raggedright\arraybackslash}p{39.46mm}>{\raggedright\arraybackslash}p{39.46mm}@{}}
\toprule
\textbf{Case} & \textbf{Singular} & \textbf{Plural} \\
\midrule
Nominative & hām & hāmas \\
Genitive & hāmes & hāma \\
Dative & hāme & hāmum \\
Accusative & hām & hāmas \\
\bottomrule
\end{tabular}
\par\endgroup

% mdtable:T06:end
% mdparagraph:P014:begin
Fōtsceamul retains the vowel of its final element -sceamul- before endings. The long vowel in the first element fōt- does not require the loss of that later vowel.
% mdparagraph:P014:end

% teaching: NEW-U15-VOC-FOTSCEAMUL
\reserve{11\baselineskip}
% mdparagraph:P015:begin
\textbf{fōtsceamul} — masculine; footstool. Stem \textbf{fōtsceamul-}.
% mdparagraph:P015:end

% teaching: NEW-U15-NOM-FOTSCEAMUL
% table: NEW-U15-NOM-FOTSCEAMUL-DISPLAY
% mdtable:T07:begin
\par\begingroup
\fontsize{10.5}{13.2}\selectfont
\setlength{\LTleft}{0pt}\setlength{\LTright}{\fill}
\noindent
\begin{tabular}{@{}>{\raggedright\arraybackslash}p{34.82mm}>{\raggedright\arraybackslash}p{41.78mm}>{\raggedright\arraybackslash}p{41.78mm}@{}}
\toprule
\textbf{Case} & \textbf{Singular} & \textbf{Plural} \\
\midrule
Nominative & fōtsceamul & fōtsceamulas \\
Genitive & fōtsceamules & fōtsceamula \\
Dative & fōtsceamule & fōtsceamulum \\
Accusative & fōtsceamul & fōtsceamulas \\
\bottomrule
\end{tabular}
\par\endgroup

% mdtable:T07:end
\reserve{7\baselineskip}
% mdparagraph:P016:begin
\subsection*{Strong neuter nouns}\BookHeading{nav-ch15-learner-0003}{subsubsection}{Strong neuter nouns}{0}
% mdparagraph:P016:end

% mdparagraph:P017:begin
Heavy stems have no ending in nominative and accusative, singular or plural. Genitive and dative singular take -es and -e; genitive and dative plural take -a and -um. Bebod has a short final stem and takes nominative/\allowbreak{}accusative plural -u.
% mdparagraph:P017:end

% teaching: NEW-U15-VOC-HAER
\reserve{11\baselineskip}
% mdparagraph:P018:begin
\textbf{hǣr} — neuter; hair; a hair. Stem \textbf{hǣr-}.
% mdparagraph:P018:end

% teaching: NEW-U15-NOM-HAER
% table: NEW-U15-NOM-HAER-DISPLAY
% mdtable:T08:begin
\par\begingroup
\fontsize{10.5}{13.2}\selectfont
\setlength{\LTleft}{0pt}\setlength{\LTright}{\fill}
\noindent
\begin{tabular}{@{}>{\raggedright\arraybackslash}p{39.46mm}>{\raggedright\arraybackslash}p{39.46mm}>{\raggedright\arraybackslash}p{39.46mm}@{}}
\toprule
\textbf{Case} & \textbf{Singular} & \textbf{Plural} \\
\midrule
Nominative & hǣr & hǣr \\
Genitive & hǣres & hǣra \\
Dative & hǣre & hǣrum \\
Accusative & hǣr & hǣr \\
\bottomrule
\end{tabular}
\par\endgroup

% mdtable:T08:end
% teaching: NEW-U15-VOC-LEOTH
\reserve{11\baselineskip}
% mdparagraph:P019:begin
\textbf{lēoð} — neuter; poem; poetry; verse. Stem \textbf{lēoð-}.
% mdparagraph:P019:end

% teaching: NEW-U15-NOM-LEOTH
% table: NEW-U15-NOM-LEOTH-DISPLAY
% mdtable:T09:begin
\par\begingroup
\fontsize{10.5}{13.2}\selectfont
\setlength{\LTleft}{0pt}\setlength{\LTright}{\fill}
\noindent
\begin{tabular}{@{}>{\raggedright\arraybackslash}p{39.46mm}>{\raggedright\arraybackslash}p{39.46mm}>{\raggedright\arraybackslash}p{39.46mm}@{}}
\toprule
\textbf{Case} & \textbf{Singular} & \textbf{Plural} \\
\midrule
Nominative & lēoð & lēoð \\
Genitive & lēoðes & lēoða \\
Dative & lēoðe & lēoðum \\
Accusative & lēoð & lēoð \\
\bottomrule
\end{tabular}
\par\endgroup

% mdtable:T09:end
% teaching: NEW-U15-VOC-BEBOD
\reserve{11\baselineskip}
% mdparagraph:P020:begin
\textbf{bebod} — neuter; commandment. Stem \textbf{bebod-}.
% mdparagraph:P020:end

% teaching: NEW-U15-NOM-BEBOD
% table: NEW-U15-NOM-BEBOD-DISPLAY
% mdtable:T10:begin
\par\begingroup
\fontsize{10.5}{13.2}\selectfont
\setlength{\LTleft}{0pt}\setlength{\LTright}{\fill}
\noindent
\begin{tabular}{@{}>{\raggedright\arraybackslash}p{39.46mm}>{\raggedright\arraybackslash}p{39.46mm}>{\raggedright\arraybackslash}p{39.46mm}@{}}
\toprule
\textbf{Case} & \textbf{Singular} & \textbf{Plural} \\
\midrule
Nominative & bebod & bebodu \\
Genitive & bebodes & beboda \\
Dative & bebode & bebodum \\
Accusative & bebod & bebodu \\
\bottomrule
\end{tabular}
\par\endgroup

% mdtable:T10:end
% mdparagraph:P021:begin
Hēafod loses the unstressed o before an ending: the ending-bearing stem is hēafd-. Its selected nominative/\allowbreak{}accusative plural is hēafdu. Compare fōtsceamul above: medial vowels are not all deleted by one general rule.
% mdparagraph:P021:end

% teaching: NEW-U15-VOC-HEAFOD
\reserve{11\baselineskip}
% mdparagraph:P022:begin
\textbf{hēafod} — neuter; head. Stem \textbf{hēafd-}.
% mdparagraph:P022:end

% teaching: NEW-U15-NOM-HEAFOD
% table: NEW-U15-NOM-HEAFOD-DISPLAY
% mdtable:T11:begin
\par\begingroup
\fontsize{10.5}{13.2}\selectfont
\setlength{\LTleft}{0pt}\setlength{\LTright}{\fill}
\noindent
\begin{tabular}{@{}>{\raggedright\arraybackslash}p{39.46mm}>{\raggedright\arraybackslash}p{39.46mm}>{\raggedright\arraybackslash}p{39.46mm}@{}}
\toprule
\textbf{Case} & \textbf{Singular} & \textbf{Plural} \\
\midrule
Nominative & hēafod & hēafdu \\
Genitive & hēafdes & hēafda \\
Dative & hēafde & hēafdum \\
Accusative & hēafod & hēafdu \\
\bottomrule
\end{tabular}
\par\endgroup

% mdtable:T11:end
% teaching: NEW-U15-VOC-FLAESC
\reserve{12\baselineskip}
% mdparagraph:P023:begin
\textbf{flǣsc} — neuter; flesh. Stem \textbf{flǣsc-}. \textbf{Plural reference:} the plural forms show the morphology; use the singular for the assigned meaning “flesh”.
% mdparagraph:P023:end

% teaching: NEW-U15-NOM-FLAESC
% table: NEW-U15-NOM-FLAESC-DISPLAY
% mdtable:T12:begin
\par\begingroup
\fontsize{10.5}{13.2}\selectfont
\setlength{\LTleft}{0pt}\setlength{\LTright}{\fill}
\noindent
\begin{tabular}{@{}>{\raggedright\arraybackslash}p{39.46mm}>{\raggedright\arraybackslash}p{39.46mm}>{\raggedright\arraybackslash}p{39.46mm}@{}}
\toprule
\textbf{Case} & \textbf{Singular} & \textbf{Plural} \\
\midrule
Nominative & flǣsc & flǣsc \\
Genitive & flǣsces & flǣsca \\
Dative & flǣsce & flǣscum \\
Accusative & flǣsc & flǣsc \\
\bottomrule
\end{tabular}
\par\endgroup

% mdtable:T12:end
\reserve{7\baselineskip}
% mdparagraph:P024:begin
\subsection*{Strong feminine nouns}\BookHeading{nav-ch15-learner-0004}{subsubsection}{Strong feminine nouns}{0}
% mdparagraph:P024:end

% mdparagraph:P025:begin
These long-stem nouns have no nominative singular ending. The other singular cases take -e. The plural has nominative/\allowbreak{}accusative and genitive -a, dative -um. Nouns in -nes double s before an ending. The -ung and -ing nouns use the same selected endings; oblique singular forms in -a are recognized alternatives.
% mdparagraph:P025:end

% teaching: NEW-U15-VOC-BRYD
\reserve{11\baselineskip}
% mdparagraph:P026:begin
\textbf{brȳd} — feminine; bride. Stem \textbf{brȳd-}.
% mdparagraph:P026:end

% teaching: NEW-U15-NOM-BRYD
% table: NEW-U15-NOM-BRYD-DISPLAY
% mdtable:T13:begin
\par\begingroup
\fontsize{10.5}{13.2}\selectfont
\setlength{\LTleft}{0pt}\setlength{\LTright}{\fill}
\noindent
\begin{tabular}{@{}>{\raggedright\arraybackslash}p{39.46mm}>{\raggedright\arraybackslash}p{39.46mm}>{\raggedright\arraybackslash}p{39.46mm}@{}}
\toprule
\textbf{Case} & \textbf{Singular} & \textbf{Plural} \\
\midrule
Nominative & brȳd & brȳda \\
Genitive & brȳde & brȳda \\
Dative & brȳde & brȳdum \\
Accusative & brȳde & brȳda \\
\bottomrule
\end{tabular}
\par\endgroup

% mdtable:T13:end
% teaching: NEW-U15-VOC-HEFIGNES
\reserve{11\baselineskip}
% mdparagraph:P027:begin
\textbf{hęfignes} — feminine; burden. Stem \textbf{hęfigness-}.
% mdparagraph:P027:end

% teaching: NEW-U15-NOM-HEFIGNES
% table: NEW-U15-NOM-HEFIGNES-DISPLAY
% mdtable:T14:begin
\par\begingroup
\fontsize{10.5}{13.2}\selectfont
\setlength{\LTleft}{0pt}\setlength{\LTright}{\fill}
\noindent
\begin{tabular}{@{}>{\raggedright\arraybackslash}p{38.19mm}>{\raggedright\arraybackslash}p{38.19mm}>{\raggedright\arraybackslash}p{42.0mm}@{}}
\toprule
\textbf{Case} & \textbf{Singular} & \textbf{Plural} \\
\midrule
Nominative & hęfignes & hęfignessa \\
Genitive & hęfignesse & hęfignessa \\
Dative & hęfignesse & hęfignessum \\
Accusative & hęfignesse & hęfignessa \\
\bottomrule
\end{tabular}
\par\endgroup

% mdtable:T14:end
% teaching: NEW-U15-VOC-AWENDEDNES
\reserve{11\baselineskip}
% mdparagraph:P028:begin
\textbf{āwęndednes} — feminine; translation. Stem \textbf{āwęndedness-}.
% mdparagraph:P028:end

% teaching: NEW-U15-NOM-AWENDEDNES
% table: NEW-U15-NOM-AWENDEDNES-DISPLAY
% mdtable:T15:begin
\par\begingroup
\fontsize{10.5}{13.2}\selectfont
\setlength{\LTleft}{0pt}\setlength{\LTright}{\fill}
\noindent
\begin{tabular}{@{}>{\raggedright\arraybackslash}p{33.82mm}>{\raggedright\arraybackslash}p{40.59mm}>{\raggedright\arraybackslash}p{43.97mm}@{}}
\toprule
\textbf{Case} & \textbf{Singular} & \textbf{Plural} \\
\midrule
Nominative & āwęndednes & āwęndednessa \\
Genitive & āwęndednesse & āwęndednessa \\
Dative & āwęndednesse & āwęndednessum \\
Accusative & āwęndednesse & āwęndednessa \\
\bottomrule
\end{tabular}
\par\endgroup

% mdtable:T15:end
% teaching: NEW-U15-VOC-COSTNUNG
\reserve{11\baselineskip}
% mdparagraph:P029:begin
\textbf{costnung} — feminine; temptation. Stem \textbf{costnung-}.
% mdparagraph:P029:end

% teaching: NEW-U15-NOM-COSTNUNG
% table: NEW-U15-NOM-COSTNUNG-DISPLAY
% mdtable:T16:begin
\par\begingroup
\fontsize{10.5}{13.2}\selectfont
\setlength{\LTleft}{0pt}\setlength{\LTright}{\fill}
\noindent
\begin{tabular}{@{}>{\raggedright\arraybackslash}p{39.46mm}>{\raggedright\arraybackslash}p{39.46mm}>{\raggedright\arraybackslash}p{39.46mm}@{}}
\toprule
\textbf{Case} & \textbf{Singular} & \textbf{Plural} \\
\midrule
Nominative & costnung & costnunga \\
Genitive & costnunge & costnunga \\
Dative & costnunge & costnungum \\
Accusative & costnunge & costnunga \\
\bottomrule
\end{tabular}
\par\endgroup

% mdtable:T16:end
% teaching: NEW-U15-VOC-BECRAEDING
\reserve{12\baselineskip}
% mdparagraph:P030:begin
\textbf{bēcrǣding} — feminine; reading; reading books. Stem \textbf{bēcrǣding-}. \textbf{Plural reference:} the meaning practised here is singular “reading”.
% mdparagraph:P030:end

% teaching: NEW-U15-NOM-BECRAEDING
% table: NEW-U15-NOM-BECRAEDING-DISPLAY
% mdtable:T17:begin
\par\begingroup
\fontsize{10.5}{13.2}\selectfont
\setlength{\LTleft}{0pt}\setlength{\LTright}{\fill}
\noindent
\begin{tabular}{@{}>{\raggedright\arraybackslash}p{38.19mm}>{\raggedright\arraybackslash}p{38.19mm}>{\raggedright\arraybackslash}p{42.0mm}@{}}
\toprule
\textbf{Case} & \textbf{Singular} & \textbf{Plural} \\
\midrule
Nominative & bēcrǣding & bēcrǣdinga \\
Genitive & bēcrǣdinge & bēcrǣdinga \\
Dative & bēcrǣdinge & bēcrǣdingum \\
Accusative & bēcrǣdinge & bēcrǣdinga \\
\bottomrule
\end{tabular}
\par\endgroup

% mdtable:T17:end
\reserve{7\baselineskip}
% mdparagraph:P031:begin
\subsection*{Weak masculine nouns}\BookHeading{nav-ch15-learner-0005}{subsubsection}{Weak masculine nouns}{0}
% mdparagraph:P031:end

% mdparagraph:P032:begin
The nominative singular ends in -a. The other singular cases and nominative/\allowbreak{}accusative plural end in -an; genitive plural takes -ena and dative plural -um. In lēoðwyrhta the final element is the noun head, and the ending follows the whole stem.
% mdparagraph:P032:end

% teaching: NEW-U15-VOC-GEREFA
\reserve{11\baselineskip}
% mdparagraph:P033:begin
\textbf{gerēfa} — masculine; steward; reeve. Stem \textbf{gerēf-}.
% mdparagraph:P033:end

% teaching: NEW-U15-NOM-GEREFA
% table: NEW-U15-NOM-GEREFA-DISPLAY
% mdtable:T18:begin
\par\begingroup
\fontsize{10.5}{13.2}\selectfont
\setlength{\LTleft}{0pt}\setlength{\LTright}{\fill}
\noindent
\begin{tabular}{@{}>{\raggedright\arraybackslash}p{39.46mm}>{\raggedright\arraybackslash}p{39.46mm}>{\raggedright\arraybackslash}p{39.46mm}@{}}
\toprule
\textbf{Case} & \textbf{Singular} & \textbf{Plural} \\
\midrule
Nominative & gerēfa & gerēfan \\
Genitive & gerēfan & gerēfena \\
Dative & gerēfan & gerēfum \\
Accusative & gerēfan & gerēfan \\
\bottomrule
\end{tabular}
\par\endgroup

% mdtable:T18:end
% teaching: NEW-U15-VOC-LEOTHWYRHTA
\reserve{11\baselineskip}
% mdparagraph:P034:begin
\textbf{lēoðwyrhta} — masculine; poet. Stem \textbf{lēoðwyrht-}.
% mdparagraph:P034:end

% teaching: NEW-U15-NOM-LEOTHWYRHTA
% table: NEW-U15-NOM-LEOTHWYRHTA-DISPLAY
% mdtable:T19:begin
\par\begingroup
\fontsize{10.5}{13.2}\selectfont
\setlength{\LTleft}{0pt}\setlength{\LTright}{\fill}
\noindent
\begin{tabular}{@{}>{\raggedright\arraybackslash}p{35.87mm}>{\raggedright\arraybackslash}p{39.46mm}>{\raggedright\arraybackslash}p{43.05mm}@{}}
\toprule
\textbf{Case} & \textbf{Singular} & \textbf{Plural} \\
\midrule
Nominative & lēoðwyrhta & lēoðwyrhtan \\
Genitive & lēoðwyrhtan & lēoðwyrhtena \\
Dative & lēoðwyrhtan & lēoðwyrhtum \\
Accusative & lēoðwyrhtan & lēoðwyrhtan \\
\bottomrule
\end{tabular}
\par\endgroup

% mdtable:T19:end
% teaching: NEW-U15-VOC-WILLA
\reserve{11\baselineskip}
% mdparagraph:P035:begin
\textbf{willa} — masculine; will; disposition. Stem \textbf{will-}. \textbf{Plural reference:} the assigned use is singular “will; disposition”.
% mdparagraph:P035:end

% teaching: NEW-U15-NOM-WILLA
% table: NEW-U15-NOM-WILLA-DISPLAY
% mdtable:T20:begin
\par\begingroup
\fontsize{10.5}{13.2}\selectfont
\setlength{\LTleft}{0pt}\setlength{\LTright}{\fill}
\noindent
\begin{tabular}{@{}>{\raggedright\arraybackslash}p{39.46mm}>{\raggedright\arraybackslash}p{39.46mm}>{\raggedright\arraybackslash}p{39.46mm}@{}}
\toprule
\textbf{Case} & \textbf{Singular} & \textbf{Plural} \\
\midrule
Nominative & willa & willan \\
Genitive & willan & willena \\
Dative & willan & willum \\
Accusative & willan & willan \\
\bottomrule
\end{tabular}
\par\endgroup

% mdtable:T20:end
\reserve{7\baselineskip}
% mdparagraph:P036:begin
\subsection*{Root and kinship nouns}\BookHeading{nav-ch15-learner-0006}{subsubsection}{Root and kinship nouns}{0}
% mdparagraph:P036:end

% mdparagraph:P037:begin
Learn\BookPoint{bk-ch15-learner-2acd7f38336d} the vowel changes with their case positions. Bōc has bēc in genitive and dative singular and nominative and accusative plural. Tōð has tēð in dative singular and nominative and accusative plural. Recognize genitive singular bōce, dative singular tōðe and nominative/\allowbreak{}accusative plural tōðas as alternatives to the selected forms below.
% mdparagraph:P037:end

% teaching: NEW-U15-VOC-BOC
\reserve{11\baselineskip}
% mdparagraph:P038:begin
\textbf{bōc} — feminine; book. Stem \textbf{bōc/\allowbreak{}bēc-}.
% mdparagraph:P038:end

% teaching: NEW-U15-NOM-BOC
% table: NEW-U15-NOM-BOC-DISPLAY
% mdtable:T21:begin
\par\begingroup
\fontsize{10.5}{13.2}\selectfont
\setlength{\LTleft}{0pt}\setlength{\LTright}{\fill}
\noindent
\begin{tabular}{@{}>{\raggedright\arraybackslash}p{39.46mm}>{\raggedright\arraybackslash}p{39.46mm}>{\raggedright\arraybackslash}p{39.46mm}@{}}
\toprule
\textbf{Case} & \textbf{Singular} & \textbf{Plural} \\
\midrule
Nominative & bōc & bēc \\
Genitive & bēc & bōca \\
Dative & bēc & bōcum \\
Accusative & bōc & bēc \\
\bottomrule
\end{tabular}
\par\endgroup

% mdtable:T21:end
% teaching: NEW-U15-VOC-TOTH
\reserve{11\baselineskip}
% mdparagraph:P039:begin
\textbf{tōð} — masculine; tooth. Stem \textbf{tōð/\allowbreak{}tēð-}.
% mdparagraph:P039:end

% teaching: NEW-U15-NOM-TOTH
% table: NEW-U15-NOM-TOTH-DISPLAY
% mdtable:T22:begin
\par\begingroup
\fontsize{10.5}{13.2}\selectfont
\setlength{\LTleft}{0pt}\setlength{\LTright}{\fill}
\noindent
\begin{tabular}{@{}>{\raggedright\arraybackslash}p{39.46mm}>{\raggedright\arraybackslash}p{39.46mm}>{\raggedright\arraybackslash}p{39.46mm}@{}}
\toprule
\textbf{Case} & \textbf{Singular} & \textbf{Plural} \\
\midrule
Nominative & tōð & tēð \\
Genitive & tōðes & tōða \\
Dative & tēð & tōðum \\
Accusative & tōð & tēð \\
\bottomrule
\end{tabular}
\par\endgroup

% mdtable:T22:end
% mdparagraph:P040:begin
Sweostor\BookPoint{bk-ch15-learner-bff20a9dbee5} is unchanged throughout the singular and in nominative/\allowbreak{}accusative plural. Only genitive plural sweostra and dative plural sweostrum have different surfaces. Its pattern must be learned separately from the earlier brōðor and mōdor.
% mdparagraph:P040:end

% teaching: NEW-U15-VOC-SWEOSTOR
\reserve{11\baselineskip}
% mdparagraph:P041:begin
\textbf{sweostor} — feminine; sister. Stem \textbf{sweostr-}.
% mdparagraph:P041:end

% teaching: NEW-U15-NOM-SWEOSTOR
% table: NEW-U15-NOM-SWEOSTOR-DISPLAY
% mdtable:T23:begin
\par\begingroup
\fontsize{10.5}{13.2}\selectfont
\setlength{\LTleft}{0pt}\setlength{\LTright}{\fill}
\noindent
\begin{tabular}{@{}>{\raggedright\arraybackslash}p{39.46mm}>{\raggedright\arraybackslash}p{39.46mm}>{\raggedright\arraybackslash}p{39.46mm}@{}}
\toprule
\textbf{Case} & \textbf{Singular} & \textbf{Plural} \\
\midrule
Nominative & sweostor & sweostor \\
Genitive & sweostor & sweostra \\
Dative & sweostor & sweostrum \\
Accusative & sweostor & sweostor \\
\bottomrule
\end{tabular}
\par\endgroup

% mdtable:T23:end
\reserve{7\baselineskip}
% mdparagraph:P042:begin
\subsection*{Abstract nouns in -u}\BookHeading{nav-ch15-learner-0007}{subsubsection}{Abstract nouns in -u}{0}
% mdparagraph:P042:end

% mdparagraph:P043:begin
Hǣtu\BookPoint{bk-ch15-learner-c5bb0898ad26} uses hǣte in the selected oblique singular cases. Hǣlu is unchanged throughout the selected singular. Oblique variants in -u, -o and -e occur, and nominative -o occurs beside -u. Learn the selected singular forms; no plural count meaning is assigned to either word.\textsuperscript{2}
% mdparagraph:P043:end

% teaching: NEW-U15-VOC-HAETU
\reserve{11\baselineskip}
% mdparagraph:P044:begin
\textbf{hǣtu} — feminine; heat. Stem \textbf{hǣt-}. \textbf{Plural unassigned.}
% mdparagraph:P044:end

% teaching: NEW-U15-NOM-HAETU
% table: NEW-U15-NOM-HAETU-DISPLAY
% mdtable:T24:begin
\par\begingroup
\fontsize{10.5}{13.2}\selectfont
\setlength{\LTleft}{0pt}\setlength{\LTright}{\fill}
\noindent
\begin{tabular}{@{}>{\raggedright\arraybackslash}p{39.46mm}>{\raggedright\arraybackslash}p{39.46mm}>{\raggedright\arraybackslash}p{39.46mm}@{}}
\toprule
\textbf{Case} & \textbf{Singular} & \textbf{Plural} \\
\midrule
Nominative & hǣtu & Unassigned \\
Genitive & hǣte & Unassigned \\
Dative & hǣte & Unassigned \\
Accusative & hǣte & Unassigned \\
\bottomrule
\end{tabular}
\par\endgroup

% mdtable:T24:end
% teaching: NEW-U15-VOC-HAELU
\reserve{11\baselineskip}
% mdparagraph:P045:begin
\textbf{hǣlu} — feminine; salvation. Stem \textbf{hǣl-}. \textbf{Plural unassigned.}
% mdparagraph:P045:end

% teaching: NEW-U15-NOM-HAELU
% table: NEW-U15-NOM-HAELU-DISPLAY
% mdtable:T25:begin
\par\begingroup
\fontsize{10.5}{13.2}\selectfont
\setlength{\LTleft}{0pt}\setlength{\LTright}{\fill}
\noindent
\begin{tabular}{@{}>{\raggedright\arraybackslash}p{39.46mm}>{\raggedright\arraybackslash}p{39.46mm}>{\raggedright\arraybackslash}p{39.46mm}@{}}
\toprule
\textbf{Case} & \textbf{Singular} & \textbf{Plural} \\
\midrule
Nominative & hǣlu & Unassigned \\
Genitive & hǣlu & Unassigned \\
Dative & hǣlu & Unassigned \\
Accusative & hǣlu & Unassigned \\
\bottomrule
\end{tabular}
\par\endgroup

% mdtable:T25:end
% mdparagraph:P046:begin
The complete class model strengu follows. It is \textbf{for recognition only in both numbers}. Its spelling is retained from Wright. Besides the selected forms, recognize nominative singular strengo, oblique singular strengu/\allowbreak{}strengo, and nominative/\allowbreak{}accusative plural strenge/\allowbreak{}strengu/\allowbreak{}strengo. The availability of a class pattern does not supply a count meaning for every abstract noun.
% mdparagraph:P046:end

% teaching: NEW-U15-VOC-STRENGU
\reserve{11\baselineskip}
% mdparagraph:P047:begin
\textbf{strengu} — feminine; strength. Stem \textbf{streng-}. \textbf{Singular and plural reference only.}
% mdparagraph:P047:end

% teaching: NEW-U15-NOM-STRENGU
% table: NEW-U15-NOM-STRENGU-DISPLAY
% mdtable:T26:begin
\par\begingroup
\fontsize{10.5}{13.2}\selectfont
\setlength{\LTleft}{0pt}\setlength{\LTright}{\fill}
\noindent
\begin{tabular}{@{}>{\raggedright\arraybackslash}p{39.46mm}>{\raggedright\arraybackslash}p{39.46mm}>{\raggedright\arraybackslash}p{39.46mm}@{}}
\toprule
\textbf{Case} & \textbf{Singular} & \textbf{Plural} \\
\midrule
Nominative & strengu & strenga \\
Genitive & strenge & strenga \\
Dative & strenge & strengum \\
Accusative & strenge & strenga \\
\bottomrule
\end{tabular}
\par\endgroup

% mdtable:T26:end
\reserve{7\baselineskip}
% mdparagraph:P048:begin
\subsection*{People names and an agent noun}\BookHeading{nav-ch15-learner-0008}{subsubsection}{People names and an agent noun}{0}
% mdparagraph:P048:end

% mdparagraph:P049:begin
Engle and Dęne name peoples in the plural. They take nominative/\allowbreak{}accusative -e, genitive -a and dative -um. No singular is assigned. Dęniga and Dęnia are recognized genitive plural alternatives.
% mdparagraph:P049:end

% teaching: NEW-U15-VOC-ENGLE
\reserve{11\baselineskip}
% mdparagraph:P050:begin
\textbf{Engle} — masculine; the English (people). Stem \textbf{Engl-}. \textbf{Singular unassigned.}
% mdparagraph:P050:end

% teaching: NEW-U15-NOM-ENGLE
% table: NEW-U15-NOM-ENGLE-DISPLAY
% mdtable:T27:begin
\par\begingroup
\fontsize{10.5}{13.2}\selectfont
\setlength{\LTleft}{0pt}\setlength{\LTright}{\fill}
\noindent
\begin{tabular}{@{}>{\raggedright\arraybackslash}p{39.46mm}>{\raggedright\arraybackslash}p{39.46mm}>{\raggedright\arraybackslash}p{39.46mm}@{}}
\toprule
\textbf{Case} & \textbf{Singular} & \textbf{Plural} \\
\midrule
Nominative & Unassigned & Engle \\
Genitive & Unassigned & Engla \\
Dative & Unassigned & Englum \\
Accusative & Unassigned & Engle \\
\bottomrule
\end{tabular}
\par\endgroup

% mdtable:T27:end
% teaching: NEW-U15-VOC-DENE
\reserve{11\baselineskip}
% mdparagraph:P051:begin
\textbf{Dęne} — masculine; the Danes. Stem \textbf{Dęn-}. \textbf{Singular unassigned.}
% mdparagraph:P051:end

% teaching: NEW-U15-NOM-DENE
% table: NEW-U15-NOM-DENE-DISPLAY
% mdtable:T28:begin
\par\begingroup
\fontsize{10.5}{13.2}\selectfont
\setlength{\LTleft}{0pt}\setlength{\LTright}{\fill}
\noindent
\begin{tabular}{@{}>{\raggedright\arraybackslash}p{39.46mm}>{\raggedright\arraybackslash}p{39.46mm}>{\raggedright\arraybackslash}p{39.46mm}@{}}
\toprule
\textbf{Case} & \textbf{Singular} & \textbf{Plural} \\
\midrule
Nominative & Unassigned & Dęne \\
Genitive & Unassigned & Dęna \\
Dative & Unassigned & Dęnum \\
Accusative & Unassigned & Dęne \\
\bottomrule
\end{tabular}
\par\endgroup

% mdtable:T28:end
% mdparagraph:P052:begin
The agent noun ālīesend has singular -es and -e as above. Its selected plural has unchanged nominative/\allowbreak{}accusative, genitive -ra and dative -um. Nominative/\allowbreak{}accusative plural forms in -e and -as are recognized alternatives. The plural is morphological reference; the assigned title “Redeemer” is singular.
% mdparagraph:P052:end

% teaching: NEW-U15-VOC-ALIESEND
\reserve{11\baselineskip}
% mdparagraph:P053:begin
\textbf{ālīesend} — masculine; Redeemer; one who delivers. Stem \textbf{ālīesend-}. \textbf{Plural reference.}
% mdparagraph:P053:end

% teaching: NEW-U15-NOM-ALIESEND
% table: NEW-U15-NOM-ALIESEND-DISPLAY
% mdtable:T29:begin
\par\begingroup
\fontsize{10.5}{13.2}\selectfont
\setlength{\LTleft}{0pt}\setlength{\LTright}{\fill}
\noindent
\begin{tabular}{@{}>{\raggedright\arraybackslash}p{39.46mm}>{\raggedright\arraybackslash}p{39.46mm}>{\raggedright\arraybackslash}p{39.46mm}@{}}
\toprule
\textbf{Case} & \textbf{Singular} & \textbf{Plural} \\
\midrule
Nominative & ālīesend & ālīesend \\
Genitive & ālīesendes & ālīesendra \\
Dative & ālīesende & ālīesendum \\
Accusative & ālīesend & ālīesend \\
\bottomrule
\end{tabular}
\par\endgroup

% mdtable:T29:end
\reserve{7\baselineskip}
% mdparagraph:P054:begin
\subsection*{Additional senses of familiar nouns}\BookHeading{nav-ch15-learner-0009}{subsubsection}{Additional senses of familiar nouns}{0}
% mdparagraph:P054:end

% mdparagraph:P055:begin
Keep the earlier declensions. These meanings are needed in the following exercises.
% mdparagraph:P055:end

% teaching: NEW-U15-RECORD-MED-ACC-PL
% teaching: NEW-U15-RECORD-MED-ACC-SG
% teaching: NEW-U15-RECORD-MED-DAT-PL
% teaching: NEW-U15-RECORD-MED-DAT-SG
% teaching: NEW-U15-RECORD-MED-GEN-PL
% teaching: NEW-U15-RECORD-MED-GEN-SG
% teaching: NEW-U15-RECORD-MED-NOM-PL
% teaching: NEW-U15-RECORD-MED-NOM-SG
% teaching: NEW-U15-VOC-MED-SENSE
% mdparagraph:P056:begin
\textbf{mēd} — hire; wages paid for work.
% mdparagraph:P056:end

% teaching: NEW-U15-RECORD-TID-ACC-PL
% teaching: NEW-U15-RECORD-TID-ACC-SG
% teaching: NEW-U15-RECORD-TID-DAT-PL
% teaching: NEW-U15-RECORD-TID-DAT-SG
% teaching: NEW-U15-RECORD-TID-GEN-PL
% teaching: NEW-U15-RECORD-TID-GEN-SG
% teaching: NEW-U15-RECORD-TID-NOM-PL
% teaching: NEW-U15-RECORD-TID-NOM-SG
% teaching: NEW-U15-VOC-TID-SENSE
% mdparagraph:P057:begin
\textbf{tīd} — hour.
% mdparagraph:P057:end

% teaching: NEW-U15-RECORD-STRAET-ACC-PL
% teaching: NEW-U15-RECORD-STRAET-ACC-SG
% teaching: NEW-U15-RECORD-STRAET-DAT-PL
% teaching: NEW-U15-RECORD-STRAET-DAT-SG
% teaching: NEW-U15-RECORD-STRAET-GEN-PL
% teaching: NEW-U15-RECORD-STRAET-GEN-SG
% teaching: NEW-U15-RECORD-STRAET-NOM-PL
% teaching: NEW-U15-RECORD-STRAET-NOM-SG
% teaching: NEW-U15-VOC-STRAET-SENSE
% mdparagraph:P058:begin
\textbf{strǣt} — public place; street as a gathering place.
% mdparagraph:P058:end

% teaching: NEW-U15-VOC-WORD
% mdparagraph:P059:begin
\textbf{word} — speech; utterance.
% mdparagraph:P059:end

\reserve{7\baselineskip}
% mdparagraph:P060:begin
\section*{2. Negation}\BookHeading{nav-ch15-learner-0010}{subsection}{2. Negation}{1}
% mdparagraph:P060:end

\reserve{7\baselineskip}
% mdparagraph:P061:begin
\subsection*{Finite negation and contractions}\BookHeading{nav-ch15-learner-0011}{subsubsection}{Finite negation and contractions}{0}
% mdparagraph:P061:end

% teaching: NEW-U15-GRAM-U15-NEGATIVE-COPULA
% mdparagraph:P062:begin
Place\BookPoint{bk-ch15-learner-435744864c61} ne before the finite verb to deny a proposition. Some common verbs have contracted negative forms. The following complete table covers \textbf{the present indicative eom/\allowbreak{}eart/\allowbreak{}is/\allowbreak{}sindon and the preterite indicative wæs/\allowbreak{}wǣre/\allowbreak{}wǣron only}. It does not replace every form of bēon/\allowbreak{}wesan. Ne eom, ne eart and ne sindon remain two words; nis and the negative preterites are contracted. Separate ne with the ordinary positive form is also valid.\textsuperscript{3}
% mdparagraph:P062:end

% teaching: NEW-U15-GRAM-U15-NEGATIVE-COPULA
% teaching: NEW-U15-PAR-WESAN-NEG-IND
\reserve{13\baselineskip}
\reserve{11\baselineskip}
% mdparagraph:P063:begin
\textbf{Negative be-series: indicative}
% mdparagraph:P063:end

% table: NEW-U15-PAR-WESAN-NEG-IND-DISPLAY
% mdtable:T30:begin
% course-table-style: 2026-10
% course-panel:T30-P01:begin
\par\begingroup
\fontsize{10.5}{13.2}\selectfont
\setlength{\LTleft}{0pt}\setlength{\LTright}{\fill}
\noindent
\begin{tabular}{@{}>{\raggedright\arraybackslash}p{39.459mm}>{\raggedright\arraybackslash}p{39.459mm}>{\raggedright\arraybackslash}p{39.459mm}@{}}
\toprule
\textbf{Person and number} & \textbf{Present} & \textbf{Preterite} \\
\midrule
1 sg. & ne eom & næs \\
2 sg. & ne eart & nǣre \\
3 sg. & nis & næs \\
1–3 pl. & ne sindon & nǣron \\
\bottomrule
\end{tabular}
\par\endgroup
% course-panel:T30-P01:end
% mdtable:T30:end
% teaching: NEW-U15-VOC-NABBAN
% mdparagraph:P064:begin
\textbf{nabban} — have not.
% mdparagraph:P064:end

% teaching: NEW-U15-GRAM-U15-NABBAN-INDICATIVE
% mdparagraph:P065:begin
The\BookPoint{bk-ch15-learner-7af15a67a1b7} complete negative possession series is now available for use. Distinguish first singular næbbe from second and third singular næfst, næfð; the present plural has bb again. A present form denies present possession; a preterite form denies past possession.
% mdparagraph:P065:end

% teaching: NEW-U15-GRAM-U15-NABBAN-INDICATIVE
% teaching: NEW-U15-PAR-NABBAN-IND
\reserve{13\baselineskip}
\reserve{11\baselineskip}
% mdparagraph:P066:begin
\textbf{Nabban: indicative}
% mdparagraph:P066:end

% table: NEW-U15-PAR-NABBAN-IND-DISPLAY
% mdtable:T31:begin
% course-table-style: 2026-10
% course-panel:T31-P01:begin
\par\begingroup
\fontsize{10.5}{13.2}\selectfont
\setlength{\LTleft}{0pt}\setlength{\LTright}{\fill}
\noindent
\begin{tabular}{@{}>{\raggedright\arraybackslash}p{39.459mm}>{\raggedright\arraybackslash}p{39.459mm}>{\raggedright\arraybackslash}p{39.459mm}@{}}
\toprule
\textbf{Person and number} & \textbf{Present} & \textbf{Preterite} \\
\midrule
1 sg. & næbbe & næfde \\
2 sg. & næfst & næfdest \\
3 sg. & næfð & næfde \\
1–3 pl. & nabbað & næfdon \\
\bottomrule
\end{tabular}
\par\endgroup
% course-panel:T31-P01:end
% mdtable:T31:end
% teaching: NEW-U15-VOC-WILLAN
% mdparagraph:P067:begin
\textbf{willan} — be unwilling; refuse (with ne).
% mdparagraph:P067:end

% teaching: NEW-U15-GRAM-U15-NEGATIVE-WILLAN
% mdparagraph:P068:begin
Negative\BookPoint{bk-ch15-learner-ffdec0152920} willan expresses refusal or unwillingness and governs a plain infinitive. The present and preterite distinguish the time of that unwillingness; English “would not” may need a past refusal, not a future form. These are indicative forms. A supplied subjunctive nolde encountered earlier does not thereby acquire every use of this table.
% mdparagraph:P068:end

% teaching: NEW-U15-GRAM-U15-NEGATIVE-WILLAN
% teaching: NEW-U15-PAR-WILLAN-NEG-IND
\reserve{13\baselineskip}
\reserve{11\baselineskip}
% mdparagraph:P069:begin
\textbf{Negative willan: indicative}
% mdparagraph:P069:end

% table: NEW-U15-PAR-WILLAN-NEG-IND-DISPLAY
% mdtable:T32:begin
% course-table-style: 2026-10
% course-panel:T32-P01:begin
\par\begingroup
\fontsize{10.5}{13.2}\selectfont
\setlength{\LTleft}{0pt}\setlength{\LTright}{\fill}
\noindent
\begin{tabular}{@{}>{\raggedright\arraybackslash}p{39.459mm}>{\raggedright\arraybackslash}p{39.459mm}>{\raggedright\arraybackslash}p{39.459mm}@{}}
\toprule
\textbf{Person and number} & \textbf{Present} & \textbf{Preterite} \\
\midrule
1 sg. & nelle & nolde \\
2 sg. & nelt & noldest \\
3 sg. & nele & nolde \\
1–3 pl. & nellað & noldon \\
\bottomrule
\end{tabular}
\par\endgroup
% course-panel:T32-P01:end
% mdtable:T32:end
\reserve{7\baselineskip}
% mdparagraph:P070:begin
\subsection*{Nān and negative concord}\BookHeading{nav-ch15-learner-0012}{subsubsection}{Nān and negative concord}{0}
% mdparagraph:P070:end

% teaching: NEW-U15-VOC-NAN
% mdparagraph:P071:begin
\textbf{nān} — no; no one.
% mdparagraph:P071:end

% mdparagraph:P072:begin
Nān\BookPoint{bk-ch15-learner-1ec05076f1c3} agrees with its noun or stands alone for “no one”. It has the strong long-stem adjective endings in the four cases below. Keep long ā; note masculine accusative singular nānne, feminine plural nāna and neuter nominative/\allowbreak{}accusative plural nān. Only this strong four-case system is required here.
% mdparagraph:P072:end

% teaching: NEW-U15-PAR-NAN-STRONG-NGDA
\reserve{11\baselineskip}
\reserve{16\baselineskip}
% mdparagraph:P073:begin
\textbf{Nān: strong declension}
% mdparagraph:P073:end

% table: NEW-U15-PAR-NAN-STRONG-NGDA-DISPLAY
% mdtable:T33:begin
% course-table-style: 2026-10
% course-panel:T33-P01:begin
\par\begingroup
\fontsize{10.5}{13.2}\selectfont
\setlength{\LTleft}{0pt}\setlength{\LTright}{\fill}
\noindent
\begin{tabular}{@{}>{\raggedright\arraybackslash}p{22.0mm}>{\raggedright\arraybackslash}p{31.188mm}>{\raggedright\arraybackslash}p{31.188mm}>{\raggedright\arraybackslash}p{31.188mm}@{}}
\toprule
\textbf{Case} & \textbf{Masculine singular} & \textbf{Neuter singular} & \textbf{Feminine singular} \\
\midrule
Nominative & nān & nān & nān \\
Genitive & nānes & nānes & nānre \\
Dative & nānum & nānum & nānre \\
Accusative & nānne & nān & nāne \\
\bottomrule
\end{tabular}
\par\endgroup
% course-panel:T33-P01:end
\par\addvspace{5pt}
% course-panel:T33-P02:begin
\par\begingroup
\fontsize{10.5}{13.2}\selectfont
\setlength{\LTleft}{0pt}\setlength{\LTright}{\fill}
\noindent
\begin{tabular}{@{}>{\raggedright\arraybackslash}p{22.0mm}>{\raggedright\arraybackslash}p{31.188mm}>{\raggedright\arraybackslash}p{31.188mm}>{\raggedright\arraybackslash}p{31.188mm}@{}}
\toprule
\textbf{Case} & \textbf{Masculine plural} & \textbf{Neuter plural} & \textbf{Feminine plural} \\
\midrule
Nominative & nāne & nān & nāna \\
Genitive & nānra & nānra & nānra \\
Dative & nānum & nānum & nānum \\
Accusative & nāne & nān & nāna \\
\bottomrule
\end{tabular}
\par\endgroup
% course-panel:T33-P02:end
% mdtable:T33:end

% teaching: NEW-U15-GRAM-U15-NEGATIVE-CONCORD
% mdparagraph:P074:begin
Nān,\BookPoint{bk-ch15-learner-73465f380192} nāwiht and nǣfre can accompany finite ne or a contracted negative verb. They express \textbf{one denial}. This is negative concord: adding a second negative does not turn the clause into an affirmation. Determine what is denied from the roles of the negative expressions, not by counting them.
% mdparagraph:P074:end

% teaching: NEW-U15-VOC-NAWIHT
% mdparagraph:P075:begin
\textbf{nāwiht} — nothing. Use the selected accusative nāwiht as the thing denied: “nothing”. Other cases are not assigned by this form.
% mdparagraph:P075:end

% teaching: NEW-U15-VOC-NA
% mdparagraph:P076:begin
\textbf{nā} — not, negating a contrasted constituent.
% mdparagraph:P076:end

% teaching: NEW-U15-GRAM-U15-CONSTITUENT-NA
% mdparagraph:P077:begin
Nā\BookPoint{bk-ch15-learner-89bd250f0c28} can deny one contrasted participant rather than the whole finite assertion. In ic self and nā ōðer, it excludes “another”; the assertion about “I myself” can remain affirmative. Distinguish this constituent use from finite ne. For a negative command, put ne before the imperative; the familiar affirmative command form remains.
% mdparagraph:P077:end

\reserve{7\baselineskip}
% mdparagraph:P078:begin
\section*{3. Indefinite reference, groups and comparison}\BookHeading{nav-ch15-learner-0013}{subsection}{3. Indefinite reference, groups and comparison}{1}
% mdparagraph:P078:end

\reserve{7\baselineskip}
% mdparagraph:P079:begin
\subsection*{An unnamed person and each member of a group}\BookHeading{nav-ch15-learner-0014}{subsubsection}{An unnamed person and each member of a group}{0}
% mdparagraph:P079:end

% teaching: NEW-U15-VOC-AN
% mdparagraph:P080:begin
\textbf{ān} — a; a certain (introducing an unnamed singular individual). Use the earlier singular strong agreement forms.
% mdparagraph:P080:end

% teaching: NEW-U15-GRAM-U15-AN-INDEFINITE
% mdparagraph:P081:begin
Context\BookPoint{bk-ch15-learner-3439a8c6c2aa} distinguishes cardinal “one” from indefinite “a certain”. Both agree with the same singular noun. An indefinite use does not create a plural cardinal meaning.
% mdparagraph:P081:end

% teaching: NEW-U15-VOC-HWA
% mdparagraph:P082:begin
\textbf{hwā} — anyone. The earlier personal/\allowbreak{}neuter case forms are unchanged.
% mdparagraph:P082:end

% teaching: NEW-U15-GRAM-U15-HWA-INDEFINITE
% mdparagraph:P083:begin
Hwā/\allowbreak{}hwæt\BookPoint{bk-ch15-learner-bfbf45918c8f} can ask “who?/\allowbreak{}what?” or refer indefinitely to “anyone/\allowbreak{}anything” in an appropriate condition or concession. The clause and its setting decide the function. English “any” cannot be replaced mechanically with hwā in every construction. Use the earlier ǣnig when an agreeing “any” is required.
% mdparagraph:P083:end

% teaching: NEW-U15-VOC-GEHWA
% mdparagraph:P084:begin
\textbf{gehwā /\allowbreak{} gehwæt} — each one; everyone /\allowbreak{} everything.
% mdparagraph:P084:end

% teaching: NEW-U15-GRAM-U15-GEHWA-DISTRIBUTIVE
% mdparagraph:P085:begin
Gehwā\BookPoint{bk-ch15-learner-f424e7a8d086} refers to each person; gehwæt to everything. Both are singular, even when they distribute over many members. The genitive can identify the group; the distributive itself takes the case required by its own role. There is one personal column, without separate masculine and feminine forms, and no plural. The selected dative is gehwǣm, and the personal accusative is gehwone.\textsuperscript{4}
% mdparagraph:P085:end

% teaching: NEW-U15-PAR-GEHWA-SINGULAR-NGDA
\reserve{11\baselineskip}
\reserve{9\baselineskip}
% mdparagraph:P086:begin
\textbf{Gehwā/\allowbreak{}gehwæt: singular}
% mdparagraph:P086:end

% table: NEW-U15-PAR-GEHWA-SINGULAR-NGDA-DISPLAY
% mdtable:T34:begin
% course-table-style: 2026-10
% course-panel:T34-P01:begin
\par\begingroup
\fontsize{10.5}{13.2}\selectfont
\setlength{\LTleft}{0pt}\setlength{\LTright}{\fill}
\noindent
\begin{tabular}{@{}>{\raggedright\arraybackslash}p{39.459mm}>{\raggedright\arraybackslash}p{39.459mm}>{\raggedright\arraybackslash}p{39.459mm}@{}}
\toprule
\textbf{Case} & \textbf{Neuter singular} & \textbf{Common singular} \\
\midrule
Nominative & gehwæt & gehwā \\
Genitive & gehwæs & gehwæs \\
Dative & gehwǣm & gehwǣm \\
Accusative & gehwæt & gehwone \\
\bottomrule
\end{tabular}
\par\endgroup
% course-panel:T34-P01:end
% mdtable:T34:end
% teaching: NEW-U15-VOC-MANN
% mdparagraph:P087:begin
\textbf{man} — one; people generally, impersonal subject with singular verb. Only nominative singular man is used in this impersonal function.
% mdparagraph:P087:end

% teaching: NEW-U15-GRAM-U15-MAN-IMPERSONAL
% mdparagraph:P088:begin
Impersonal\BookPoint{bk-ch15-learner-7c10bcfe8c52} man means “one; people in general” and takes a third-person singular verb. Do not give it the oblique or plural forms of the ordinary count noun mann. The count noun still refers to an individual man and retains its earlier declension.
% mdparagraph:P088:end

% teaching: NEW-U15-GRAM-U15-QUANTITY-PARTITIVES
% mdparagraph:P089:begin
Ǣlc\BookPoint{bk-ch15-learner-ca8a7ca630d6} distributes a singular share to each member; eall can gather the members as a plural whole. An explicit plural subject can retain a plural verb while ǣlc his pęning distributes one penny to each person. A genitive after a selection word identifies the whole group; it does not determine the selecting word’s case.
% mdparagraph:P089:end

% teaching: NEW-U15-VOC-IC
% mdparagraph:P090:begin
\textbf{uncer}\BookPoint{bk-ch15-learner-8156c28183ed} — of us two, personal genitive after a selection word. Personal genitive uncer means “of us two”. Its form is fixed by the partitive relation; it is distinct from the earlier agreeing possessive uncer, “our (two people’s)”.
% mdparagraph:P090:end

\reserve{7\baselineskip}
% mdparagraph:P091:begin
\subsection*{Free choice and paired coordination}\BookHeading{nav-ch15-learner-0015}{subsubsection}{Free choice and paired coordination}{0}
% mdparagraph:P091:end

% teaching: NEW-U15-VOC-HWAETHER
% mdparagraph:P092:begin
\textbf{hwæðer} — whichever of two, in swā hwæðer swā.
% mdparagraph:P092:end

% teaching: NEW-U15-VOC-SWA
% mdparagraph:P093:begin
\textbf{swā} — free-choice framing: swā hwā/\allowbreak{}hwæt/\allowbreak{}hwæðer swā; the ... the, proportional relation.
% mdparagraph:P093:end

% teaching: NEW-U15-GRAM-U15-FREE-CHOICE
% mdparagraph:P094:begin
Swā\BookPoint{bk-ch15-learner-81d02f820a47} hwā swā means “whoever”; swā hwæt swā means “whatever”. Swā hwæðer swā selects “whichever of two”. Give the pronoun the case needed within its clause. A pronoun in the main clause can resume the same referent in a different case, such as a subject within the choice clause and a dative recipient in the main clause. A possible or contingent event may take the subjunctive; free choice does not mechanically require one mood. Cook’s exceptional attracted dative in 31.12 is supplied with that item.
% mdparagraph:P094:end

% teaching: NEW-U15-VOC-AEGTHER
% mdparagraph:P095:begin
\textbf{ǣgðer} — both, in ǣgðer ge…ge.
% mdparagraph:P095:end

% teaching: NEW-U15-GRAM-U15-EITHER-BOTH-COORDINATION
% mdparagraph:P096:begin
Ǣgðer\BookPoint{bk-ch15-learner-b36c32144c55} ge … ge means “both … and”. It includes both coordinated constituents. Here ǣgðer is an invariant coordinator; this expression does not require the separate inflectional family meaning “either”.
% mdparagraph:P096:end

\reserve{7\baselineskip}
% mdparagraph:P097:begin
\subsection*{Equality, preference and proportion}\BookHeading{nav-ch15-learner-0016}{subsubsection}{Equality, preference and proportion}{0}
% mdparagraph:P097:end

% teaching: NEW-U15-GRAM-U15-PROPORTIONAL-SWA
% mdparagraph:P098:begin
In\BookPoint{bk-ch15-learner-eacf155021af} equality, swā … swā accompanies a positive adjective: “as … as”. In a proportional comparison, each swā clause contains a comparative: “the more …, the more …”. Each clause retains its own subject, verb and adjective agreement. A proportional relation does not assert equal degree.
% mdparagraph:P098:end

% teaching: NEW-U15-VOC-SWITHE
% mdparagraph:P099:begin
\textbf{swīðe}\BookPoint{bk-ch15-learner-24b5b5f6352f} — rather; more, comparing preference; most; very greatly. The positive swīðe “very; greatly” is familiar. Learn comparative swīðor for preference; superlative swīðost is for recognition.
% mdparagraph:P099:end

% teaching: NEW-U15-PAR-SWITHE-DEGREES
% teaching: NEW-U15-VOC-SWITHE
\reserve{10\baselineskip}
\reserve{8\baselineskip}
% mdparagraph:P100:begin
\textbf{Swīðe: degrees of the adverb}
% mdparagraph:P100:end

% table: NEW-U15-PAR-SWITHE-DEGREES-DISPLAY
% mdtable:T35:begin
\par\begingroup
\fontsize{10.5}{13.2}\selectfont
\setlength{\LTleft}{0pt}\setlength{\LTright}{\fill}
\noindent
\begin{tabular}{@{}>{\raggedright\arraybackslash}p{63.48mm}>{\raggedright\arraybackslash}p{57.71mm}@{}}
\toprule
\textbf{Position} & \textbf{Form} \\
\midrule
Positive & swīðe \\
Comparative & swīðor \\
Superlative & swīðost \\
\bottomrule
\end{tabular}
\par\endgroup

% mdtable:T35:end
% teaching: NEW-U15-GRAM-U15-RATHER-THAN
% mdparagraph:P101:begin
Swīðor\BookPoint{bk-ch15-learner-66e05fdac64f} … ðonne compares preference: “more/\allowbreak{}rather … than”. The compared expressions retain the cases required by their verb. With lufian, both the preferred object and the alternative are accusative; neither becomes nominative merely by following ðonne.
% mdparagraph:P101:end

% teaching: NEW-U15-VOC-MICEL
% mdparagraph:P102:begin
\textbf{micel} — much; an amount. Use the earlier neuter forms substantivally for an amount.
% mdparagraph:P102:end

% teaching: NEW-U15-VOC-MICEL-COMP
% mdparagraph:P103:begin
\textbf{māre}\BookPoint{bk-ch15-learner-6ec683c21516} — more, a greater amount. The neuter substantive māre means “more; a greater amount”.
% mdparagraph:P103:end

% teaching: NEW-U15-RECORD-EALL-U15-DEGREE-ADVERB
% teaching: NEW-U15-VOC-EALL
% mdparagraph:P104:begin
\textbf{eall} — just; exactly, in eall swā. This adverbial use is invariant and distinct from agreeing “all”.
% mdparagraph:P104:end

% teaching: NEW-U15-GRAM-U15-EALL-EQUATIVE
% mdparagraph:P105:begin
Eall\BookPoint{bk-ch15-learner-019bcd600b0c} strengthens an equative in eall swā micel swā, “just as much as”. Micel denotes the amount. When the comparison leaves “give” understood, the second recipient remains dative, just as the first recipient does.
% mdparagraph:P105:end

% mdparagraph:P106:begin
\textbf{Reference limits.} Āwiht “anything” and hwæthwega “something” are supplied as exact objects in A10; they are not models for a new case system. The variants ōht and nōht, and the further family citations nǣnig, āhwæðer, nāhwæðer, ǣghwæðer, āhwā, ǣghwā and swā hwilc swā, require no new recall here. The earlier scope of other indefinite families remains unchanged. Instrumental degree expressions are reference material, and the possessive sīn is not introduced.
% mdparagraph:P106:end

\reserve{7\baselineskip}
% mdparagraph:P107:begin
\section*{4. Additional verbs}\BookHeading{nav-ch15-learner-0017}{subsection}{4. Additional verbs}{1}
% mdparagraph:P107:end

% mdparagraph:P108:begin
The following systems supply the forms required by the exercises. Each indicative table is complete for present and preterite, all three persons, singular and plural. Other moods and nonfinite forms are defined separately. A principal part identifies a stem; its citation alone does not require every possible use.\textsuperscript{5}
% mdparagraph:P108:end

\reserve{7\baselineskip}
% mdparagraph:P109:begin
\subsection*{Strong indicative systems}\BookHeading{nav-ch15-learner-0018}{subsubsection}{Strong indicative systems}{0}
% mdparagraph:P109:end

% teaching: NEW-U15-VOC-FORGIELDAN
% mdparagraph:P110:begin
\textbf{forgieldan} — pay; requite; recompense. Strong III; present forgield-, preterite singular forgeald-, preterite plural forguld-, past participle forgolden.
% mdparagraph:P110:end

% mdparagraph:P111:begin
The second- and third-singular present forms forgieltst and forgielt use the selected t cluster. The preterite second singular uses the plural stem plus -e; first and third singular have the bare singular stem. All present plurals end in -að and preterite plurals in -on. The prefix remains attached throughout.
% mdparagraph:P111:end

% teaching: NEW-U15-PAR-FORGIELDAN-IND
\reserve{13\baselineskip}
\reserve{11\baselineskip}
% mdparagraph:P112:begin
\textbf{Forgieldan: indicative}
% mdparagraph:P112:end

% table: NEW-U15-PAR-FORGIELDAN-IND-DISPLAY
% mdtable:T36:begin
% course-table-style: 2026-10
% course-panel:T36-P01:begin
\par\begingroup
\fontsize{10.5}{13.2}\selectfont
\setlength{\LTleft}{0pt}\setlength{\LTright}{\fill}
\noindent
\begin{tabular}{@{}>{\raggedright\arraybackslash}p{39.459mm}>{\raggedright\arraybackslash}p{39.459mm}>{\raggedright\arraybackslash}p{39.459mm}@{}}
\toprule
\textbf{Person and number} & \textbf{Present} & \textbf{Preterite} \\
\midrule
1 sg. & forgielde & forgeald \\
2 sg. & forgieltst & forgulde \\
3 sg. & forgielt & forgeald \\
1–3 pl. & forgieldað & forguldon \\
\bottomrule
\end{tabular}
\par\endgroup
% course-panel:T36-P01:end
% mdtable:T36:end
% teaching: NEW-U15-VOC-TOWEORPAN
% mdparagraph:P113:begin
\textbf{tōweorpan} — break; overthrow, here break a commandment. Strong III; present tōweorp-, preterite singular tōwearp-, preterite plural tōwurp-, past participle tōworpen.
% mdparagraph:P113:end

% mdparagraph:P114:begin
The second- and third-singular present stem is tōwierp-. The preterite follows the same person/\allowbreak{}number distribution as forgieldan.
% mdparagraph:P114:end

% teaching: NEW-U15-PAR-TOWEORPAN-IND
\reserve{13\baselineskip}
\reserve{11\baselineskip}
% mdparagraph:P115:begin
\textbf{Tōweorpan: indicative}
% mdparagraph:P115:end

% table: NEW-U15-PAR-TOWEORPAN-IND-DISPLAY
% mdtable:T37:begin
% course-table-style: 2026-10
% course-panel:T37-P01:begin
\par\begingroup
\fontsize{10.5}{13.2}\selectfont
\setlength{\LTleft}{0pt}\setlength{\LTright}{\fill}
\noindent
\begin{tabular}{@{}>{\raggedright\arraybackslash}p{39.459mm}>{\raggedright\arraybackslash}p{39.459mm}>{\raggedright\arraybackslash}p{39.459mm}@{}}
\toprule
\textbf{Person and number} & \textbf{Present} & \textbf{Preterite} \\
\midrule
1 sg. & tōweorpe & tōwearp \\
2 sg. & tōwierpst & tōwurpe \\
3 sg. & tōwierpð & tōwearp \\
1–3 pl. & tōweorpað & tōwurpon \\
\bottomrule
\end{tabular}
\par\endgroup
% course-panel:T37-P01:end
% mdtable:T37:end
% teaching: NEW-U15-VOC-ONGIETAN
% mdparagraph:P116:begin
\textbf{ongietan} — understand; perceive. Strong V; present ongiet-, preterite singular ongeat-, preterite plural ongēat-, past participle ongieten.
% mdparagraph:P116:end

% mdparagraph:P117:begin
Second singular ongietst and third singular ongiett retain the selected t-cluster spellings. Third singular ongiet is a recognized alternative. Preterite second singular ongēate takes the plural stem; first and third singular have ongeat. The principal-part infinitive and past participle are reference citations; the complete new system here is indicative.
% mdparagraph:P117:end

% teaching: NEW-U15-PAR-ONGIETAN-IND
\reserve{13\baselineskip}
\reserve{11\baselineskip}
% mdparagraph:P118:begin
\textbf{Ongietan: indicative}
% mdparagraph:P118:end

% table: NEW-U15-PAR-ONGIETAN-IND-DISPLAY
% mdtable:T38:begin
% course-table-style: 2026-10
% course-panel:T38-P01:begin
\par\begingroup
\fontsize{10.5}{13.2}\selectfont
\setlength{\LTleft}{0pt}\setlength{\LTright}{\fill}
\noindent
\begin{tabular}{@{}>{\raggedright\arraybackslash}p{39.459mm}>{\raggedright\arraybackslash}p{39.459mm}>{\raggedright\arraybackslash}p{39.459mm}@{}}
\toprule
\textbf{Person and number} & \textbf{Present} & \textbf{Preterite} \\
\midrule
1 sg. & ongiete & ongeat \\
2 sg. & ongietst & ongēate \\
3 sg. & ongiett & ongeat \\
1–3 pl. & ongietað & ongēaton \\
\bottomrule
\end{tabular}
\par\endgroup
% course-panel:T38-P01:end
% mdtable:T38:end
% teaching: NEW-U15-VOC-GEHATAN
% mdparagraph:P119:begin
\textbf{gehātan} — promise. Strong VII. This indicative is a review of the earlier system.
% mdparagraph:P119:end

% mdparagraph:P120:begin
Gehātan keeps its prefix and the familiar hātan stem changes: ā becomes ǣ in the second- and third-singular present, while the preterite has ē. This review does not add a new prefixed negative or subjunctive system.
% mdparagraph:P120:end

% teaching: NEW-U15-PAR-GEHATAN-IND
\reserve{13\baselineskip}
\reserve{11\baselineskip}
% mdparagraph:P121:begin
\textbf{Gehātan: indicative review}
% mdparagraph:P121:end

% table: NEW-U15-PAR-GEHATAN-IND-DISPLAY
% mdtable:T39:begin
% course-table-style: 2026-10
% course-panel:T39-P01:begin
\par\begingroup
\fontsize{10.5}{13.2}\selectfont
\setlength{\LTleft}{0pt}\setlength{\LTright}{\fill}
\noindent
\begin{tabular}{@{}>{\raggedright\arraybackslash}p{39.459mm}>{\raggedright\arraybackslash}p{39.459mm}>{\raggedright\arraybackslash}p{39.459mm}@{}}
\toprule
\textbf{Person and number} & \textbf{Present} & \textbf{Preterite} \\
\midrule
1 sg. & gehāte & gehēt \\
2 sg. & gehǣtst & gehēte \\
3 sg. & gehǣtt & gehēt \\
1–3 pl. & gehātað & gehēton \\
\bottomrule
\end{tabular}
\par\endgroup
% course-panel:T39-P01:end
% mdtable:T39:end
\reserve{7\baselineskip}
% mdparagraph:P122:begin
\subsection*{Weak indicative systems}\BookHeading{nav-ch15-learner-0019}{subsubsection}{Weak indicative systems}{0}
% mdparagraph:P122:end

% mdparagraph:P123:begin
The long-stem weak I verbs below use present -e, -st, -ð, -að and preterite -de, -dest, -de, -don. The repeated plural ending serves all three persons. Dēman also has the complete subjunctive shown separately.
% mdparagraph:P123:end

% teaching: NEW-U15-VOC-WENAN
% mdparagraph:P124:begin
\textbf{wēnan} — suppose; expect; think. Weak I.
% mdparagraph:P124:end

% teaching: NEW-U15-PAR-WENAN-IND
\reserve{13\baselineskip}
\reserve{11\baselineskip}
% mdparagraph:P125:begin
\textbf{Wēnan: indicative}
% mdparagraph:P125:end

% table: NEW-U15-PAR-WENAN-IND-DISPLAY
% mdtable:T40:begin
% course-table-style: 2026-10
% course-panel:T40-P01:begin
\par\begingroup
\fontsize{10.5}{13.2}\selectfont
\setlength{\LTleft}{0pt}\setlength{\LTright}{\fill}
\noindent
\begin{tabular}{@{}>{\raggedright\arraybackslash}p{39.459mm}>{\raggedright\arraybackslash}p{39.459mm}>{\raggedright\arraybackslash}p{39.459mm}@{}}
\toprule
\textbf{Person and number} & \textbf{Present} & \textbf{Preterite} \\
\midrule
1 sg. & wēne & wēnde \\
2 sg. & wēnst & wēndest \\
3 sg. & wēnð & wēnde \\
1–3 pl. & wēnað & wēndon \\
\bottomrule
\end{tabular}
\par\endgroup
% course-panel:T40-P01:end
% mdtable:T40:end
% teaching: NEW-U15-VOC-FORBAERNAN
% mdparagraph:P126:begin
\textbf{forbærnan} — parch; scorch. Weak I.
% mdparagraph:P126:end

% teaching: NEW-U15-PAR-FORBAERNAN-IND
\reserve{13\baselineskip}
\reserve{11\baselineskip}
% mdparagraph:P127:begin
\textbf{Forbærnan: indicative}
% mdparagraph:P127:end

% table: NEW-U15-PAR-FORBAERNAN-IND-DISPLAY
% mdtable:T41:begin
% course-table-style: 2026-10
% course-panel:T41-P01:begin
\par\begingroup
\fontsize{10.5}{13.2}\selectfont
\setlength{\LTleft}{0pt}\setlength{\LTright}{\fill}
\noindent
\begin{tabular}{@{}>{\raggedright\arraybackslash}p{39.459mm}>{\raggedright\arraybackslash}p{39.459mm}>{\raggedright\arraybackslash}p{39.459mm}@{}}
\toprule
\textbf{Person and number} & \textbf{Present} & \textbf{Preterite} \\
\midrule
1 sg. & forbærne & forbærnde \\
2 sg. & forbærnst & forbærndest \\
3 sg. & forbærnð & forbærnde \\
1–3 pl. & forbærnað & forbærndon \\
\bottomrule
\end{tabular}
\par\endgroup
% course-panel:T41-P01:end
% mdtable:T41:end
% teaching: NEW-U15-VOC-LAERAN
\par\begin{minipage}{\linewidth}\raggedright
% mdparagraph:P128:begin
\textbf{lǣran} — teach; instruct. Weak I.
% mdparagraph:P128:end

% teaching: NEW-U15-PAR-LAERAN-IND
% mdparagraph:P129:begin
\textbf{Lǣran: indicative}
% mdparagraph:P129:end

% table: NEW-U15-PAR-LAERAN-IND-DISPLAY
% mdtable:T42:begin
% course-table-style: 2026-10
% course-panel:T42-P01:begin
\par\begingroup
\fontsize{10.5}{13.2}\selectfont
\setlength{\LTleft}{0pt}\setlength{\LTright}{\fill}
\noindent
\begin{tabular}{@{}>{\raggedright\arraybackslash}p{39.459mm}>{\raggedright\arraybackslash}p{39.459mm}>{\raggedright\arraybackslash}p{39.459mm}@{}}
\toprule
\textbf{Person and number} & \textbf{Present} & \textbf{Preterite} \\
\midrule
1 sg. & lǣre & lǣrde \\
2 sg. & lǣrst & lǣrdest \\
3 sg. & lǣrð & lǣrde \\
1–3 pl. & lǣrað & lǣrdon \\
\bottomrule
\end{tabular}
\par\endgroup
% course-panel:T42-P01:end
% mdtable:T42:end
\end{minipage}\par\addvspace{4pt}
% teaching: NEW-U15-VOC-DEMAN
% mdparagraph:P130:begin
\textbf{dēman} — judge; condemn. Weak I.
% mdparagraph:P130:end

% teaching: NEW-U15-PAR-DEMAN-IND
\reserve{13\baselineskip}
\reserve{11\baselineskip}
% mdparagraph:P131:begin
\textbf{Dēman: indicative}
% mdparagraph:P131:end

% table: NEW-U15-PAR-DEMAN-IND-DISPLAY
% mdtable:T43:begin
% course-table-style: 2026-10
% course-panel:T43-P01:begin
\par\begingroup
\fontsize{10.5}{13.2}\selectfont
\setlength{\LTleft}{0pt}\setlength{\LTright}{\fill}
\noindent
\begin{tabular}{@{}>{\raggedright\arraybackslash}p{39.459mm}>{\raggedright\arraybackslash}p{39.459mm}>{\raggedright\arraybackslash}p{39.459mm}@{}}
\toprule
\textbf{Person and number} & \textbf{Present} & \textbf{Preterite} \\
\midrule
1 sg. & dēme & dēmde \\
2 sg. & dēmst & dēmdest \\
3 sg. & dēmð & dēmde \\
1–3 pl. & dēmað & dēmdon \\
\bottomrule
\end{tabular}
\par\endgroup
% course-panel:T43-P01:end
% mdtable:T43:end
% mdparagraph:P132:begin
The subjunctive singular has one form for all three persons; the plural has -n. Dēme can therefore be indicative first singular or present subjunctive singular; the clause and its function decide. Use the earlier distinctions among assertion, purpose and wish.
% mdparagraph:P132:end

% teaching: NEW-U15-PAR-DEMAN-SUBJ
\reserve{13\baselineskip}
\reserve{11\baselineskip}
% mdparagraph:P133:begin
\textbf{Dēman: subjunctive}
% mdparagraph:P133:end

% table: NEW-U15-PAR-DEMAN-SUBJ-DISPLAY
% mdtable:T44:begin
% course-table-style: 2026-10
% course-panel:T44-P01:begin
\par\begingroup
\fontsize{10.5}{13.2}\selectfont
\setlength{\LTleft}{0pt}\setlength{\LTright}{\fill}
\noindent
\begin{tabular}{@{}>{\raggedright\arraybackslash}p{39.459mm}>{\raggedright\arraybackslash}p{39.459mm}>{\raggedright\arraybackslash}p{39.459mm}@{}}
\toprule
\textbf{Person and number} & \textbf{Present} & \textbf{Preterite} \\
\midrule
1 sg. & dēme & dēmde \\
2 sg. & dēme & dēmde \\
3 sg. & dēme & dēmde \\
1–3 pl. & dēmen & dēmden \\
\bottomrule
\end{tabular}
\par\endgroup
% course-panel:T44-P01:end
% mdtable:T44:end
% teaching: NEW-U15-VOC-HYRIAN
% mdparagraph:P134:begin
\textbf{hȳrian} — hire. Weak II. Keep the long ȳ throughout.
% mdparagraph:P134:end

% mdparagraph:P135:begin
Hȳrian follows lufian: present singular hȳrie, hȳrast, hȳrað, plural hȳriað; the preterite uses hȳrod- and the ordinary weak endings. Distinguish third-singular -að from plural -iað.
% mdparagraph:P135:end

% teaching: NEW-U15-PAR-HYRIAN-IND
\reserve{13\baselineskip}
\reserve{11\baselineskip}
% mdparagraph:P136:begin
\textbf{Hȳrian: indicative}
% mdparagraph:P136:end

% table: NEW-U15-PAR-HYRIAN-IND-DISPLAY
% mdtable:T45:begin
% course-table-style: 2026-10
% course-panel:T45-P01:begin
\par\begingroup
\fontsize{10.5}{13.2}\selectfont
\setlength{\LTleft}{0pt}\setlength{\LTright}{\fill}
\noindent
\begin{tabular}{@{}>{\raggedright\arraybackslash}p{39.459mm}>{\raggedright\arraybackslash}p{39.459mm}>{\raggedright\arraybackslash}p{39.459mm}@{}}
\toprule
\textbf{Person and number} & \textbf{Present} & \textbf{Preterite} \\
\midrule
1 sg. & hȳrie & hȳrode \\
2 sg. & hȳrast & hȳrodest \\
3 sg. & hȳrað & hȳrode \\
1–3 pl. & hȳriað & hȳrodon \\
\bottomrule
\end{tabular}
\par\endgroup
% course-panel:T45-P01:end
% mdtable:T45:end
\reserve{7\baselineskip}
% mdparagraph:P137:begin
\subsection*{Commands}\BookHeading{nav-ch15-learner-0020}{subsubsection}{Commands}{0}
% mdparagraph:P137:end

% mdparagraph:P138:begin
Each table gives the complete second-person singular/\allowbreak{}plural command pair. Forgieldan has the strong singular stem without an ending. The weak I singulars ālīes and ætīew also have no ending; their plurals take -að. Weak II cleopa has singular -a and plural cleopiað.
% mdparagraph:P138:end

% teaching: NEW-U15-PAR-FORGIELDAN-IMP
\reserve{9\baselineskip}
\reserve{7\baselineskip}
% mdparagraph:P139:begin
\textbf{Forgieldan: imperative}
% mdparagraph:P139:end

% table: NEW-U15-PAR-FORGIELDAN-IMP-DISPLAY
% mdtable:T46:begin
% course-table-style: 2026-10
% course-panel:T46-P01:begin
\par\begingroup
\fontsize{10.5}{13.2}\selectfont
\setlength{\LTleft}{0pt}\setlength{\LTright}{\fill}
\noindent
\begin{tabular}{@{}>{\raggedright\arraybackslash}p{63.342mm}>{\raggedright\arraybackslash}p{57.846mm}@{}}
\toprule
\textbf{Person and number} & \textbf{Imperative} \\
\midrule
2 sg. & forgield \\
2 pl. & forgieldað \\
\bottomrule
\end{tabular}
\par\endgroup
% course-panel:T46-P01:end
% mdtable:T46:end
% teaching: NEW-U15-VOC-CLEOPIAN
% mdparagraph:P140:begin
\textbf{cleopian} — call; summon. Weak II; only the command pair is newly required for use.
% mdparagraph:P140:end

% teaching: NEW-U15-PAR-CLEOPIAN-IMP
\reserve{9\baselineskip}
\reserve{7\baselineskip}
% mdparagraph:P141:begin
\textbf{Cleopian: imperative}
% mdparagraph:P141:end

% table: NEW-U15-PAR-CLEOPIAN-IMP-DISPLAY
% mdtable:T47:begin
% course-table-style: 2026-10
% course-panel:T47-P01:begin
\par\begingroup
\fontsize{10.5}{13.2}\selectfont
\setlength{\LTleft}{0pt}\setlength{\LTright}{\fill}
\noindent
\begin{tabular}{@{}>{\raggedright\arraybackslash}p{63.342mm}>{\raggedright\arraybackslash}p{57.846mm}@{}}
\toprule
\textbf{Person and number} & \textbf{Imperative} \\
\midrule
2 sg. & cleopa \\
2 pl. & cleopiað \\
\bottomrule
\end{tabular}
\par\endgroup
% course-panel:T47-P01:end
% mdtable:T47:end
% teaching: NEW-U15-VOC-ALIESAN
% mdparagraph:P142:begin
\textbf{ālīesan} — deliver; release. Weak I; only the command pair is newly required for use.
% mdparagraph:P142:end

% teaching: NEW-U15-PAR-ALIESAN-IMP
\reserve{9\baselineskip}
\reserve{7\baselineskip}
% mdparagraph:P143:begin
\textbf{Ālīesan: imperative}
% mdparagraph:P143:end

% table: NEW-U15-PAR-ALIESAN-IMP-DISPLAY
% mdtable:T48:begin
% course-table-style: 2026-10
% course-panel:T48-P01:begin
\par\begingroup
\fontsize{10.5}{13.2}\selectfont
\setlength{\LTleft}{0pt}\setlength{\LTright}{\fill}
\noindent
\begin{tabular}{@{}>{\raggedright\arraybackslash}p{63.342mm}>{\raggedright\arraybackslash}p{57.846mm}@{}}
\toprule
\textbf{Person and number} & \textbf{Imperative} \\
\midrule
2 sg. & ālīes \\
2 pl. & ālīesað \\
\bottomrule
\end{tabular}
\par\endgroup
% course-panel:T48-P01:end
% mdtable:T48:end
% mdparagraph:P144:begin
\textbf{ætīewan} — show; the earlier verb, now with its complete command pair. Retain short initial æ in ætīew.
% mdparagraph:P144:end

% teaching: NEW-U15-PAR-AETIEWAN-IMP
\reserve{9\baselineskip}
\reserve{7\baselineskip}
% mdparagraph:P145:begin
\textbf{Ætīewan: imperative}
% mdparagraph:P145:end

% table: NEW-U15-PAR-AETIEWAN-IMP-DISPLAY
% mdtable:T49:begin
% course-table-style: 2026-10
% course-panel:T49-P01:begin
\par\begingroup
\fontsize{10.5}{13.2}\selectfont
\setlength{\LTleft}{0pt}\setlength{\LTright}{\fill}
\noindent
\begin{tabular}{@{}>{\raggedright\arraybackslash}p{63.342mm}>{\raggedright\arraybackslash}p{57.846mm}@{}}
\toprule
\textbf{Person and number} & \textbf{Imperative} \\
\midrule
2 sg. & ætīew \\
2 pl. & ætīewað \\
\bottomrule
\end{tabular}
\par\endgroup
% course-panel:T49-P01:end
% mdtable:T49:end
\reserve{7\baselineskip}
% mdparagraph:P146:begin
\subsection*{Infinitives}\BookHeading{nav-ch15-learner-0021}{subsubsection}{Infinitives}{0}
% mdparagraph:P146:end

% mdparagraph:P147:begin
These are complete pairs of plain and inflected infinitives. Use the plain form after the taught modal verbs. The inflected infinitive follows tō in the uses already learned; particular difficult frames are supplied locally. The pair does not introduce a finite system for the verb.
% mdparagraph:P147:end

% teaching: NEW-U15-PAR-FORGIELDAN-INF
\reserve{9\baselineskip}
\reserve{7\baselineskip}
% mdparagraph:P148:begin
\textbf{Forgieldan: infinitives}
% mdparagraph:P148:end

% table: NEW-U15-PAR-FORGIELDAN-INF-DISPLAY
% mdtable:T50:begin
\par\begingroup
\fontsize{10.5}{13.2}\selectfont
\setlength{\LTleft}{0pt}\setlength{\LTright}{\fill}
\noindent
\begin{tabular}{@{}>{\raggedright\arraybackslash}p{66.31mm}>{\raggedright\arraybackslash}p{54.88mm}@{}}
\toprule
\textbf{Position} & \textbf{Form} \\
\midrule
Plain infinitive & forgieldan \\
Inflected infinitive & forgieldanne \\
\bottomrule
\end{tabular}
\par\endgroup

% mdtable:T50:end
% mdparagraph:P149:begin
\textbf{wyrcean} — work; do. The selected pair retains ea before the inflected ending.
% mdparagraph:P149:end

% teaching: NEW-U15-PAR-WYRCEAN-INF
\reserve{9\baselineskip}
\reserve{7\baselineskip}
% mdparagraph:P150:begin
\textbf{Wyrcean: infinitives}
% mdparagraph:P150:end

% table: NEW-U15-PAR-WYRCEAN-INF-DISPLAY
% mdtable:T51:begin
\par\begingroup
\fontsize{10.5}{13.2}\selectfont
\setlength{\LTleft}{0pt}\setlength{\LTright}{\fill}
\noindent
\begin{tabular}{@{}>{\raggedright\arraybackslash}p{78.74mm}>{\raggedright\arraybackslash}p{42.45mm}@{}}
\toprule
\textbf{Position} & \textbf{Form} \\
\midrule
Plain infinitive & wyrcean \\
Inflected infinitive after tō & wyrceanne \\
\bottomrule
\end{tabular}
\par\endgroup

% mdtable:T51:end
% mdparagraph:P151:begin
\textbf{onfōn} — receive; keep the contracted stem in both infinitives.
% mdparagraph:P151:end

% teaching: NEW-U15-PAR-ONFON-INF
\reserve{9\baselineskip}
\reserve{7\baselineskip}
% mdparagraph:P152:begin
\textbf{Onfōn: infinitives}
% mdparagraph:P152:end

% table: NEW-U15-PAR-ONFON-INF-DISPLAY
% mdtable:T52:begin
\par\begingroup
\fontsize{10.5}{13.2}\selectfont
\setlength{\LTleft}{0pt}\setlength{\LTright}{\fill}
\noindent
\begin{tabular}{@{}>{\raggedright\arraybackslash}p{76.63mm}>{\raggedright\arraybackslash}p{44.55mm}@{}}
\toprule
\textbf{Position} & \textbf{Form} \\
\midrule
Infinitive (citation form) & onfōn \\
Inflected infinitive & onfōnne \\
\bottomrule
\end{tabular}
\par\endgroup

% mdtable:T52:end
\par\begin{minipage}{\linewidth}\raggedright
% mdparagraph:P153:begin
\textbf{forlǣtan} — forsake; leave. Keep the long ǣ in both infinitives.
% mdparagraph:P153:end

% teaching: NEW-U15-PAR-FORLAETAN-INF
% mdparagraph:P154:begin
\textbf{Forlǣtan: infinitives}
% mdparagraph:P154:end

% table: NEW-U15-PAR-FORLAETAN-INF-DISPLAY
% mdtable:T53:begin
\par\begingroup
\fontsize{10.5}{13.2}\selectfont
\setlength{\LTleft}{0pt}\setlength{\LTright}{\fill}
\noindent
\begin{tabular}{@{}>{\raggedright\arraybackslash}p{71.72mm}>{\raggedright\arraybackslash}p{49.46mm}@{}}
\toprule
\textbf{Position} & \textbf{Form} \\
\midrule
Plain infinitive & forlǣtan \\
Inflected infinitive & forlǣtanne \\
\bottomrule
\end{tabular}
\par\endgroup

% mdtable:T53:end
\end{minipage}\par\addvspace{4pt}
% teaching: NEW-U15-VOC-BLOSTMIAN
% mdparagraph:P155:begin
\textbf{blōstmian} — blossom; bloom. Weak II; only the infinitive pair is required for use here.
% mdparagraph:P155:end

% teaching: NEW-U15-PAR-BLOSTMIAN-INF
\reserve{9\baselineskip}
\reserve{7\baselineskip}
% mdparagraph:P156:begin
\textbf{Blōstmian: infinitives}
% mdparagraph:P156:end

% table: NEW-U15-PAR-BLOSTMIAN-INF-DISPLAY
% mdtable:T54:begin
\par\begingroup
\fontsize{10.5}{13.2}\selectfont
\setlength{\LTleft}{0pt}\setlength{\LTright}{\fill}
\noindent
\begin{tabular}{@{}>{\raggedright\arraybackslash}p{68.91mm}>{\raggedright\arraybackslash}p{52.28mm}@{}}
\toprule
\textbf{Position} & \textbf{Form} \\
\midrule
Plain infinitive & blōstmian \\
Inflected infinitive & blōstmianne \\
\bottomrule
\end{tabular}
\par\endgroup

% mdtable:T54:end
\reserve{7\baselineskip}
% mdparagraph:P157:begin
\subsection*{Participial bases}\BookHeading{nav-ch15-learner-0022}{subsubsection}{Participial bases}{0}
% mdparagraph:P157:end

% mdparagraph:P158:begin
Each pair gives the present and past participial base, without adding a complete adjective declension. The present participle can express an accompanying activity in the earlier predicative construction; the past participle can enter the earlier passive. Keep full -ed in lǣred and hǣled. Standan is strong VI: its present base standende retains n, and its past base is standen.
% mdparagraph:P158:end

% teaching: NEW-U15-PAR-FORGIELDAN-PART
\reserve{9\baselineskip}
\reserve{7\baselineskip}
% mdparagraph:P159:begin
\textbf{Forgieldan: participial bases}
% mdparagraph:P159:end

% table: NEW-U15-PAR-FORGIELDAN-PART-DISPLAY
% mdtable:T55:begin
\par\begingroup
\fontsize{10.5}{13.2}\selectfont
\setlength{\LTleft}{0pt}\setlength{\LTright}{\fill}
\noindent
\begin{tabular}{@{}>{\raggedright\arraybackslash}p{74.28mm}>{\raggedright\arraybackslash}p{46.91mm}@{}}
\toprule
\textbf{Position} & \textbf{Form} \\
\midrule
Present participle (base form) & forgieldende \\
Past participle (base form) & forgolden \\
\bottomrule
\end{tabular}
\par\endgroup

% mdtable:T55:end
% teaching: NEW-U15-PAR-LAERAN-PART
\reserve{9\baselineskip}
\reserve{7\baselineskip}
% mdparagraph:P160:begin
\textbf{Lǣran: participial bases}
% mdparagraph:P160:end

% table: NEW-U15-PAR-LAERAN-PART-DISPLAY
% mdtable:T56:begin
\par\begingroup
\fontsize{10.5}{13.2}\selectfont
\setlength{\LTleft}{0pt}\setlength{\LTright}{\fill}
\noindent
\begin{tabular}{@{}>{\raggedright\arraybackslash}p{79.4mm}>{\raggedright\arraybackslash}p{41.79mm}@{}}
\toprule
\textbf{Position} & \textbf{Form} \\
\midrule
Present participle (base form) & lǣrende \\
Past participle (base form) & lǣred \\
\bottomrule
\end{tabular}
\par\endgroup

% mdtable:T56:end
% teaching: NEW-U15-VOC-HAELAN
% mdparagraph:P161:begin
\textbf{hǣlan} — save. Weak I; the participial bases are required, not a new finite system.
% mdparagraph:P161:end

% teaching: NEW-U15-PAR-HAELAN-PART
\reserve{9\baselineskip}
\reserve{7\baselineskip}
% mdparagraph:P162:begin
\textbf{Hǣlan: participial bases}
% mdparagraph:P162:end

% table: NEW-U15-PAR-HAELAN-PART-DISPLAY
% mdtable:T57:begin
\par\begingroup
\fontsize{10.5}{13.2}\selectfont
\setlength{\LTleft}{0pt}\setlength{\LTright}{\fill}
\noindent
\begin{tabular}{@{}>{\raggedright\arraybackslash}p{79.4mm}>{\raggedright\arraybackslash}p{41.79mm}@{}}
\toprule
\textbf{Position} & \textbf{Form} \\
\midrule
Present participle (base form) & hǣlende \\
Past participle (base form) & hǣled \\
\bottomrule
\end{tabular}
\par\endgroup

% mdtable:T57:end
% mdparagraph:P163:begin
\textbf{standan} — stand. \textbf{giefan} — give. Their meanings are familiar; the following complete base pairs supply the forms needed here.
% mdparagraph:P163:end

% teaching: NEW-U15-PAR-STANDAN-PART
\reserve{9\baselineskip}
\reserve{7\baselineskip}
% mdparagraph:P164:begin
\textbf{Standan: participial bases}
% mdparagraph:P164:end

% table: NEW-U15-PAR-STANDAN-PART-DISPLAY
% mdtable:T58:begin
\par\begingroup
\fontsize{10.5}{13.2}\selectfont
\setlength{\LTleft}{0pt}\setlength{\LTright}{\fill}
\noindent
\begin{tabular}{@{}>{\raggedright\arraybackslash}p{80.04mm}>{\raggedright\arraybackslash}p{41.15mm}@{}}
\toprule
\textbf{Position} & \textbf{Form} \\
\midrule
Present participle (base form) & standende \\
Past participle (citation form) & standen \\
\bottomrule
\end{tabular}
\par\endgroup

% mdtable:T58:end
% teaching: NEW-U15-PAR-GIEFAN-PART
\reserve{9\baselineskip}
\reserve{7\baselineskip}
% mdparagraph:P165:begin
\textbf{Giefan: participial bases}
% mdparagraph:P165:end

% table: NEW-U15-PAR-GIEFAN-PART-DISPLAY
% mdtable:T59:begin
\par\begingroup
\fontsize{10.5}{13.2}\selectfont
\setlength{\LTleft}{0pt}\setlength{\LTright}{\fill}
\noindent
\begin{tabular}{@{}>{\raggedright\arraybackslash}p{80.04mm}>{\raggedright\arraybackslash}p{41.15mm}@{}}
\toprule
\textbf{Position} & \textbf{Form} \\
\midrule
Present participle (base form) & giefende \\
Past participle (citation form) & giefen \\
\bottomrule
\end{tabular}
\par\endgroup

% mdtable:T59:end
\reserve{32\baselineskip}
% mdparagraph:P166:begin
\subsection*{Principal parts for reference}\BookHeading{nav-ch15-learner-0023}{subsubsection}{Principal parts for reference}{0}
% mdparagraph:P166:end

% mdparagraph:P167:begin
For strong verbs the order is infinitive, preterite singular, preterite plural, past participle. Weak verbs have infinitive, preterite singular and past participle. This list is available for reference. Only forms in a separately defined system above, or already learned earlier, require independent use; the remaining citation forms are supplied.
% mdparagraph:P167:end

% teaching: NEW-U15-PAR-BLOSTMIAN-INF
% teaching: NEW-U15-PAR-DEMAN-IND
% teaching: NEW-U15-PAR-FORBAERNAN-IND
% teaching: NEW-U15-PAR-FORGIELDAN-IND
% teaching: NEW-U15-PAR-FORGIELDAN-INF
% teaching: NEW-U15-PAR-FORGIELDAN-PART
% teaching: NEW-U15-PAR-FORLAETAN-INF
% teaching: NEW-U15-PAR-GEHATAN-IND
% teaching: NEW-U15-PAR-GIEFAN-PART
% teaching: NEW-U15-PAR-HAELAN-PART
% teaching: NEW-U15-PAR-HYRIAN-IND
% teaching: NEW-U15-PAR-LAERAN-IND
% teaching: NEW-U15-PAR-LAERAN-PART
% teaching: NEW-U15-PAR-ONFON-INF
% teaching: NEW-U15-PAR-ONGIETAN-IND
% teaching: NEW-U15-PAR-STANDAN-PART
% teaching: NEW-U15-PAR-TOWEORPAN-IND
% teaching: NEW-U15-PAR-WENAN-IND
% teaching: NEW-U15-PAR-WYRCEAN-INF
% teaching: NEW-U15-PRINCIPAL-PARTS
% mdtable:T60:begin
\par\begingroup
\fontsize{10.5}{13.2}\selectfont
\setlength{\LTleft}{0pt}\setlength{\LTright}{\fill}
\begin{longtable}{@{}>{\raggedright\arraybackslash}p{19.63mm}>{\raggedright\arraybackslash}p{19.63mm}>{\raggedright\arraybackslash}p{25.61mm}>{\raggedright\arraybackslash}p{23.85mm}>{\raggedright\arraybackslash}p{24.04mm}@{}}
\toprule
\textbf{Infinitive} & \textbf{Type} & \textbf{Preterite singular} & \textbf{Preterite plural} & \textbf{Past participle} \\
\midrule
\endfirsthead
\toprule
\textbf{Infinitive} & \textbf{Type} & \textbf{Preterite singular} & \textbf{Preterite plural} & \textbf{Past participle} \\
\midrule
\endhead
\midrule
\multicolumn{5}{r}{\footnotesize Continued on the following page.}\\
\endfoot
\bottomrule
\endlastfoot
dēman & weak & dēmde & — & dēmed \\
forgieldan & strong & forgeald & forguldon & forgolden \\
hȳrian & weak & hȳrode & — & hȳrod \\
wēnan & weak & wēnde & — & wēned \\
tōweorpan & strong & tōwearp & tōwurpon & tōworpen \\
forbærnan & weak & forbærnde & — & forbærned \\
lǣran & weak & lǣrde & — & lǣred \\
ongietan & strong & ongeat & ongēaton & ongieten \\
gehātan & strong & gehēt & gehēton & gehāten \\
cleopian & weak & cleopode & — & cleopod \\
hǣlan & weak & hǣlde & — & hǣled \\
ālīesan & weak & ālīesde & — & ālīesed \\
blōstmian & weak & blōstmode & — & blōstmod \\
ætīewan & weak & ætīewde & — & ætīewed \\
wyrcean & weak & worhte & — & worht \\
onfōn & strong & onfēng & onfēngon & onfangen \\
forlǣtan & strong & forlēt & forlēton & forlǣten \\
standan & strong & stōd & stōdon & standen \\
giefan & strong & geaf & gēafon & giefen \\
\end{longtable}
\par\endgroup

% mdtable:T60:end
\reserve{7\baselineskip}
\par\begin{minipage}{\linewidth}\raggedright
% mdparagraph:P168:begin
\section*{5. Adjectives, meanings and government}\BookHeading{nav-ch15-learner-0024}{subsection}{5. Adjectives, meanings and government}{1}
% mdparagraph:P168:end

% mdparagraph:P169:begin
\subsection*{Complete adjective systems}\BookHeading{nav-ch15-learner-0025}{subsubsection}{Complete adjective systems}{0}
% mdparagraph:P169:end

% mdparagraph:P170:begin
Each adjective has complete strong and weak tables for the four cases, both numbers and all three genders. Use strong forms in the familiar indefinite and predicative constructions, weak forms after the definite determiner as previously learned. No instrumental system is added.\textsuperscript{6}
% mdparagraph:P170:end
\end{minipage}\par\addvspace{4pt}


% mdparagraph:P171:begin
Riht and gelīc follow the ordinary long-stem pattern: strong feminine nominative singular and neuter nominative/\allowbreak{}accusative plural have no ending. The selected strong feminine plural has -a, masculine plural -e. Weak genitive plural has -ra.
% mdparagraph:P171:end

% teaching: NEW-U15-VOC-RIHT
% mdparagraph:P172:begin
\textbf{riht} — right; just.
% mdparagraph:P172:end

% teaching: NEW-U15-PAR-RIHT-STRONG-NGDA
\reserve{11\baselineskip}
\reserve{16\baselineskip}
% mdparagraph:P173:begin
\textbf{Riht: strong}
% mdparagraph:P173:end

% table: NEW-U15-PAR-RIHT-STRONG-NGDA-DISPLAY
% mdtable:T61:begin
% course-table-style: 2026-10
% course-panel:T61-P01:begin
\par\begingroup
\fontsize{10.5}{13.2}\selectfont
\setlength{\LTleft}{0pt}\setlength{\LTright}{\fill}
\noindent
\begin{tabular}{@{}>{\raggedright\arraybackslash}p{22.0mm}>{\raggedright\arraybackslash}p{31.188mm}>{\raggedright\arraybackslash}p{31.188mm}>{\raggedright\arraybackslash}p{31.188mm}@{}}
\toprule
\textbf{Case} & \textbf{Masculine singular} & \textbf{Neuter singular} & \textbf{Feminine singular} \\
\midrule
Nominative & riht & riht & riht \\
Genitive & rihtes & rihtes & rihtre \\
Dative & rihtum & rihtum & rihtre \\
Accusative & rihtne & riht & rihte \\
\bottomrule
\end{tabular}
\par\endgroup
% course-panel:T61-P01:end
\par\addvspace{5pt}
% course-panel:T61-P02:begin
\par\begingroup
\fontsize{10.5}{13.2}\selectfont
\setlength{\LTleft}{0pt}\setlength{\LTright}{\fill}
\noindent
\begin{tabular}{@{}>{\raggedright\arraybackslash}p{22.0mm}>{\raggedright\arraybackslash}p{31.188mm}>{\raggedright\arraybackslash}p{31.188mm}>{\raggedright\arraybackslash}p{31.188mm}@{}}
\toprule
\textbf{Case} & \textbf{Masculine plural} & \textbf{Neuter plural} & \textbf{Feminine plural} \\
\midrule
Nominative & rihte & riht & rihta \\
Genitive & rihtra & rihtra & rihtra \\
Dative & rihtum & rihtum & rihtum \\
Accusative & rihte & riht & rihta \\
\bottomrule
\end{tabular}
\par\endgroup
% course-panel:T61-P02:end
% mdtable:T61:end

% teaching: NEW-U15-PAR-RIHT-WEAK-NGDA
\reserve{11\baselineskip}
\reserve{16\baselineskip}
% mdparagraph:P174:begin
\textbf{Riht: weak}
% mdparagraph:P174:end

% table: NEW-U15-PAR-RIHT-WEAK-NGDA-DISPLAY
% mdtable:T62:begin
% course-table-style: 2026-10
% course-panel:T62-P01:begin
\par\begingroup
\fontsize{10.5}{13.2}\selectfont
\setlength{\LTleft}{0pt}\setlength{\LTright}{\fill}
\noindent
\begin{tabular}{@{}>{\raggedright\arraybackslash}p{22.0mm}>{\raggedright\arraybackslash}p{31.188mm}>{\raggedright\arraybackslash}p{31.188mm}>{\raggedright\arraybackslash}p{31.188mm}@{}}
\toprule
\textbf{Case} & \textbf{Masculine singular} & \textbf{Neuter singular} & \textbf{Feminine singular} \\
\midrule
Nominative & rihta & rihte & rihte \\
Genitive & rihtan & rihtan & rihtan \\
Dative & rihtan & rihtan & rihtan \\
Accusative & rihtan & rihte & rihtan \\
\bottomrule
\end{tabular}
\par\endgroup
% course-panel:T62-P01:end
\par\addvspace{5pt}
% course-panel:T62-P02:begin
\par\begingroup
\fontsize{10.5}{13.2}\selectfont
\setlength{\LTleft}{0pt}\setlength{\LTright}{\fill}
\noindent
\begin{tabular}{@{}>{\raggedright\arraybackslash}p{22.0mm}>{\raggedright\arraybackslash}p{31.188mm}>{\raggedright\arraybackslash}p{31.188mm}>{\raggedright\arraybackslash}p{31.188mm}@{}}
\toprule
\textbf{Case} & \textbf{Masculine plural} & \textbf{Neuter plural} & \textbf{Feminine plural} \\
\midrule
Nominative & rihtan & rihtan & rihtan \\
Genitive & rihtra & rihtra & rihtra \\
Dative & rihtum & rihtum & rihtum \\
Accusative & rihtan & rihtan & rihtan \\
\bottomrule
\end{tabular}
\par\endgroup
% course-panel:T62-P02:end
% mdtable:T62:end

% teaching: NEW-U15-VOC-GELIC
\par\begin{minipage}{\linewidth}\raggedright
% mdparagraph:P175:begin
\textbf{gelīc} — equal; like.
% mdparagraph:P175:end

% teaching: NEW-U15-PAR-GELIC-STRONG-NGDA
% mdparagraph:P176:begin
\textbf{Gelīc: strong}
% mdparagraph:P176:end

% table: NEW-U15-PAR-GELIC-STRONG-NGDA-DISPLAY
% mdtable:T63:begin
% course-table-style: 2026-10
% course-panel:T63-P01:begin
\par\begingroup
\fontsize{10.5}{13.2}\selectfont
\setlength{\LTleft}{0pt}\setlength{\LTright}{\fill}
\noindent
\begin{tabular}{@{}>{\raggedright\arraybackslash}p{22.0mm}>{\raggedright\arraybackslash}p{31.188mm}>{\raggedright\arraybackslash}p{31.188mm}>{\raggedright\arraybackslash}p{31.188mm}@{}}
\toprule
\textbf{Case} & \textbf{Masculine singular} & \textbf{Neuter singular} & \textbf{Feminine singular} \\
\midrule
Nominative & gelīc & gelīc & gelīc \\
Genitive & gelīces & gelīces & gelīcre \\
Dative & gelīcum & gelīcum & gelīcre \\
Accusative & gelīcne & gelīc & gelīce \\
\bottomrule
\end{tabular}
\par\endgroup
% course-panel:T63-P01:end
\par\addvspace{5pt}
% course-panel:T63-P02:begin
\par\begingroup
\fontsize{10.5}{13.2}\selectfont
\setlength{\LTleft}{0pt}\setlength{\LTright}{\fill}
\noindent
\begin{tabular}{@{}>{\raggedright\arraybackslash}p{22.0mm}>{\raggedright\arraybackslash}p{31.188mm}>{\raggedright\arraybackslash}p{31.188mm}>{\raggedright\arraybackslash}p{31.188mm}@{}}
\toprule
\textbf{Case} & \textbf{Masculine plural} & \textbf{Neuter plural} & \textbf{Feminine plural} \\
\midrule
Nominative & gelīce & gelīc & gelīca \\
Genitive & gelīcra & gelīcra & gelīcra \\
Dative & gelīcum & gelīcum & gelīcum \\
Accusative & gelīce & gelīc & gelīca \\
\bottomrule
\end{tabular}
\par\endgroup
% course-panel:T63-P02:end
% mdtable:T63:end
\end{minipage}\par\addvspace{4pt}

% teaching: NEW-U15-PAR-GELIC-WEAK-NGDA
\reserve{11\baselineskip}
\reserve{16\baselineskip}
% mdparagraph:P177:begin
\textbf{Gelīc: weak}
% mdparagraph:P177:end

% table: NEW-U15-PAR-GELIC-WEAK-NGDA-DISPLAY
% mdtable:T64:begin
% course-table-style: 2026-10
% course-panel:T64-P01:begin
\par\begingroup
\fontsize{10.5}{13.2}\selectfont
\setlength{\LTleft}{0pt}\setlength{\LTright}{\fill}
\noindent
\begin{tabular}{@{}>{\raggedright\arraybackslash}p{22.0mm}>{\raggedright\arraybackslash}p{31.188mm}>{\raggedright\arraybackslash}p{31.188mm}>{\raggedright\arraybackslash}p{31.188mm}@{}}
\toprule
\textbf{Case} & \textbf{Masculine singular} & \textbf{Neuter singular} & \textbf{Feminine singular} \\
\midrule
Nominative & gelīca & gelīce & gelīce \\
Genitive & gelīcan & gelīcan & gelīcan \\
Dative & gelīcan & gelīcan & gelīcan \\
Accusative & gelīcan & gelīce & gelīcan \\
\bottomrule
\end{tabular}
\par\endgroup
% course-panel:T64-P01:end
\par\addvspace{5pt}
% course-panel:T64-P02:begin
\par\begingroup
\fontsize{10.5}{13.2}\selectfont
\setlength{\LTleft}{0pt}\setlength{\LTright}{\fill}
\noindent
\begin{tabular}{@{}>{\raggedright\arraybackslash}p{22.0mm}>{\raggedright\arraybackslash}p{31.188mm}>{\raggedright\arraybackslash}p{31.188mm}>{\raggedright\arraybackslash}p{31.188mm}@{}}
\toprule
\textbf{Case} & \textbf{Masculine plural} & \textbf{Neuter plural} & \textbf{Feminine plural} \\
\midrule
Nominative & gelīcan & gelīcan & gelīcan \\
Genitive & gelīcra & gelīcra & gelīcra \\
Dative & gelīcum & gelīcum & gelīcum \\
Accusative & gelīcan & gelīcan & gelīcan \\
\bottomrule
\end{tabular}
\par\endgroup
% course-panel:T64-P02:end
% mdtable:T64:end

% teaching: NEW-U15-VOC-MIHTIG
% mdparagraph:P178:begin
\textbf{mihtig} — mighty; powerful.
% mdparagraph:P178:end

% mdparagraph:P179:begin
The selected forms lose i before a vowel-initial ending: mihtges, mihtgum, mihtge, mihtgan. They retain i before a consonantal ending: mihtigne, mihtigre, mihtigra. Strong feminine nominative singular and neuter nominative/\allowbreak{}accusative plural are mihtig, parallel to the earlier hālig. Unsyncopated forms are recognized variants.
% mdparagraph:P179:end

% teaching: NEW-U15-PAR-MIHTIG-STRONG-NGDA
\reserve{11\baselineskip}
\reserve{16\baselineskip}
% mdparagraph:P180:begin
\textbf{Mihtig: strong}
% mdparagraph:P180:end

% table: NEW-U15-PAR-MIHTIG-STRONG-NGDA-DISPLAY
% mdtable:T65:begin
% course-table-style: 2026-10
% course-panel:T65-P01:begin
\par\begingroup
\fontsize{10.5}{13.2}\selectfont
\setlength{\LTleft}{0pt}\setlength{\LTright}{\fill}
\noindent
\begin{tabular}{@{}>{\raggedright\arraybackslash}p{22.0mm}>{\raggedright\arraybackslash}p{31.188mm}>{\raggedright\arraybackslash}p{31.188mm}>{\raggedright\arraybackslash}p{31.188mm}@{}}
\toprule
\textbf{Case} & \textbf{Masculine singular} & \textbf{Neuter singular} & \textbf{Feminine singular} \\
\midrule
Nominative & mihtig & mihtig & mihtig \\
Genitive & mihtges & mihtges & mihtigre \\
Dative & mihtgum & mihtgum & mihtigre \\
Accusative & mihtigne & mihtig & mihtge \\
\bottomrule
\end{tabular}
\par\endgroup
% course-panel:T65-P01:end
\par\addvspace{5pt}
% course-panel:T65-P02:begin
\par\begingroup
\fontsize{10.5}{13.2}\selectfont
\setlength{\LTleft}{0pt}\setlength{\LTright}{\fill}
\noindent
\begin{tabular}{@{}>{\raggedright\arraybackslash}p{22.0mm}>{\raggedright\arraybackslash}p{31.188mm}>{\raggedright\arraybackslash}p{31.188mm}>{\raggedright\arraybackslash}p{31.188mm}@{}}
\toprule
\textbf{Case} & \textbf{Masculine plural} & \textbf{Neuter plural} & \textbf{Feminine plural} \\
\midrule
Nominative & mihtge & mihtig & mihtga \\
Genitive & mihtigra & mihtigra & mihtigra \\
Dative & mihtgum & mihtgum & mihtgum \\
Accusative & mihtge & mihtig & mihtga \\
\bottomrule
\end{tabular}
\par\endgroup
% course-panel:T65-P02:end
% mdtable:T65:end

% teaching: NEW-U15-PAR-MIHTIG-WEAK-NGDA
\reserve{11\baselineskip}
\reserve{16\baselineskip}
% mdparagraph:P181:begin
\textbf{Mihtig: weak}
% mdparagraph:P181:end

% table: NEW-U15-PAR-MIHTIG-WEAK-NGDA-DISPLAY
% mdtable:T66:begin
% course-table-style: 2026-10
% course-panel:T66-P01:begin
\par\begingroup
\fontsize{10.5}{13.2}\selectfont
\setlength{\LTleft}{0pt}\setlength{\LTright}{\fill}
\noindent
\begin{tabular}{@{}>{\raggedright\arraybackslash}p{22.0mm}>{\raggedright\arraybackslash}p{31.188mm}>{\raggedright\arraybackslash}p{31.188mm}>{\raggedright\arraybackslash}p{31.188mm}@{}}
\toprule
\textbf{Case} & \textbf{Masculine singular} & \textbf{Neuter singular} & \textbf{Feminine singular} \\
\midrule
Nominative & mihtga & mihtge & mihtge \\
Genitive & mihtgan & mihtgan & mihtgan \\
Dative & mihtgan & mihtgan & mihtgan \\
Accusative & mihtgan & mihtge & mihtgan \\
\bottomrule
\end{tabular}
\par\endgroup
% course-panel:T66-P01:end
\par\addvspace{5pt}
% course-panel:T66-P02:begin
\par\begingroup
\fontsize{10.5}{13.2}\selectfont
\setlength{\LTleft}{0pt}\setlength{\LTright}{\fill}
\noindent
\begin{tabular}{@{}>{\raggedright\arraybackslash}p{22.0mm}>{\raggedright\arraybackslash}p{31.188mm}>{\raggedright\arraybackslash}p{31.188mm}>{\raggedright\arraybackslash}p{31.188mm}@{}}
\toprule
\textbf{Case} & \textbf{Masculine plural} & \textbf{Neuter plural} & \textbf{Feminine plural} \\
\midrule
Nominative & mihtgan & mihtgan & mihtgan \\
Genitive & mihtigra & mihtigra & mihtigra \\
Dative & mihtgum & mihtgum & mihtgum \\
Accusative & mihtgan & mihtgan & mihtgan \\
\bottomrule
\end{tabular}
\par\endgroup
% course-panel:T66-P02:end
% mdtable:T66:end

% teaching: NEW-U15-VOC-MAERE
% mdparagraph:P182:begin
\textbf{mǣre} — great; renowned.
% mdparagraph:P182:end

% mdparagraph:P183:begin
Mǣre keeps its final e where there is no added ending and drops it before another ending. Note strong feminine nominative singular and neuter nominative/\allowbreak{}accusative plural mǣru. Stem-final r with -re/\allowbreak{}-ra gives mǣrre/\allowbreak{}mǣrra; weak genitive plural is also mǣrra.
% mdparagraph:P183:end

% teaching: NEW-U15-PAR-MAERE-STRONG-NGDA
\reserve{11\baselineskip}
\reserve{16\baselineskip}
% mdparagraph:P184:begin
\textbf{Mǣre: strong}
% mdparagraph:P184:end

% table: NEW-U15-PAR-MAERE-STRONG-NGDA-DISPLAY
% mdtable:T67:begin
% course-table-style: 2026-10
% course-panel:T67-P01:begin
\par\begingroup
\fontsize{10.5}{13.2}\selectfont
\setlength{\LTleft}{0pt}\setlength{\LTright}{\fill}
\noindent
\begin{tabular}{@{}>{\raggedright\arraybackslash}p{22.0mm}>{\raggedright\arraybackslash}p{31.188mm}>{\raggedright\arraybackslash}p{31.188mm}>{\raggedright\arraybackslash}p{31.188mm}@{}}
\toprule
\textbf{Case} & \textbf{Masculine singular} & \textbf{Neuter singular} & \textbf{Feminine singular} \\
\midrule
Nominative & mǣre & mǣre & mǣru \\
Genitive & mǣres & mǣres & mǣrre \\
Dative & mǣrum & mǣrum & mǣrre \\
Accusative & mǣrne & mǣre & mǣre \\
\bottomrule
\end{tabular}
\par\endgroup
% course-panel:T67-P01:end
\par\addvspace{5pt}
% course-panel:T67-P02:begin
\par\begingroup
\fontsize{10.5}{13.2}\selectfont
\setlength{\LTleft}{0pt}\setlength{\LTright}{\fill}
\noindent
\begin{tabular}{@{}>{\raggedright\arraybackslash}p{22.0mm}>{\raggedright\arraybackslash}p{31.188mm}>{\raggedright\arraybackslash}p{31.188mm}>{\raggedright\arraybackslash}p{31.188mm}@{}}
\toprule
\textbf{Case} & \textbf{Masculine plural} & \textbf{Neuter plural} & \textbf{Feminine plural} \\
\midrule
Nominative & mǣre & mǣru & mǣra \\
Genitive & mǣrra & mǣrra & mǣrra \\
Dative & mǣrum & mǣrum & mǣrum \\
Accusative & mǣre & mǣru & mǣra \\
\bottomrule
\end{tabular}
\par\endgroup
% course-panel:T67-P02:end
% mdtable:T67:end

% teaching: NEW-U15-PAR-MAERE-WEAK-NGDA
\reserve{11\baselineskip}
\reserve{16\baselineskip}
% mdparagraph:P185:begin
\textbf{Mǣre: weak}
% mdparagraph:P185:end

% table: NEW-U15-PAR-MAERE-WEAK-NGDA-DISPLAY
% mdtable:T68:begin
% course-table-style: 2026-10
% course-panel:T68-P01:begin
\par\begingroup
\fontsize{10.5}{13.2}\selectfont
\setlength{\LTleft}{0pt}\setlength{\LTright}{\fill}
\noindent
\begin{tabular}{@{}>{\raggedright\arraybackslash}p{22.0mm}>{\raggedright\arraybackslash}p{31.188mm}>{\raggedright\arraybackslash}p{31.188mm}>{\raggedright\arraybackslash}p{31.188mm}@{}}
\toprule
\textbf{Case} & \textbf{Masculine singular} & \textbf{Neuter singular} & \textbf{Feminine singular} \\
\midrule
Nominative & mǣra & mǣre & mǣre \\
Genitive & mǣran & mǣran & mǣran \\
Dative & mǣran & mǣran & mǣran \\
Accusative & mǣran & mǣre & mǣran \\
\bottomrule
\end{tabular}
\par\endgroup
% course-panel:T68-P01:end
\par\addvspace{5pt}
% course-panel:T68-P02:begin
\par\begingroup
\fontsize{10.5}{13.2}\selectfont
\setlength{\LTleft}{0pt}\setlength{\LTright}{\fill}
\noindent
\begin{tabular}{@{}>{\raggedright\arraybackslash}p{22.0mm}>{\raggedright\arraybackslash}p{31.188mm}>{\raggedright\arraybackslash}p{31.188mm}>{\raggedright\arraybackslash}p{31.188mm}@{}}
\toprule
\textbf{Case} & \textbf{Masculine plural} & \textbf{Neuter plural} & \textbf{Feminine plural} \\
\midrule
Nominative & mǣran & mǣran & mǣran \\
Genitive & mǣrra & mǣrra & mǣrra \\
Dative & mǣrum & mǣrum & mǣrum \\
Accusative & mǣran & mǣran & mǣran \\
\bottomrule
\end{tabular}
\par\endgroup
% course-panel:T68-P02:end
% mdtable:T68:end

% teaching: NEW-U15-VOC-UNGETHANCFULL
% mdparagraph:P186:begin
\textbf{ungeðancfull} — unthankful.
% mdparagraph:P186:end

% mdparagraph:P187:begin
Ungeðancfull retains ll in the bare form and before vowel-initial endings, as in ungeðancfulle and ungeðancfullum. Before consonantal endings the selected forms have one l: ungeðancfulne, ungeðancfulre, ungeðancfulra.
% mdparagraph:P187:end

% teaching: NEW-U15-PAR-UNGETHANCFULL-STRONG-NGDA
\reserve{11\baselineskip}
\reserve{16\baselineskip}
% mdparagraph:P188:begin
\textbf{Ungeðancfull: strong}
% mdparagraph:P188:end

% table: NEW-U15-PAR-UNGETHANCFULL-STRONG-NGDA-DISPLAY
% mdtable:T69:begin
% course-table-style: 2026-10
% course-panel:T69-P01:begin
\par\begingroup
\fontsize{10.5}{13.2}\selectfont
\setlength{\LTleft}{0pt}\setlength{\LTright}{\fill}
\noindent
\begin{tabular}{@{}>{\raggedright\arraybackslash}p{22.0mm}>{\raggedright\arraybackslash}p{31.188mm}>{\raggedright\arraybackslash}p{31.188mm}>{\raggedright\arraybackslash}p{31.188mm}@{}}
\toprule
\textbf{Case} & \textbf{Masculine singular} & \textbf{Neuter singular} & \textbf{Feminine singular} \\
\midrule
Nominative & ungeðancfull & ungeðancfull & ungeðancfull \\
Genitive & ungeðancfulles & ungeðancfulles & ungeðancfulre \\
Dative & ungeðancfullum & ungeðancfullum & ungeðancfulre \\
Accusative & ungeðancfulne & ungeðancfull & ungeðancfulle \\
\bottomrule
\end{tabular}
\par\endgroup
% course-panel:T69-P01:end
\par\addvspace{5pt}
% course-panel:T69-P02:begin
\par\begingroup
\fontsize{10.5}{13.2}\selectfont
\setlength{\LTleft}{0pt}\setlength{\LTright}{\fill}
\noindent
\begin{tabular}{@{}>{\raggedright\arraybackslash}p{22.0mm}>{\raggedright\arraybackslash}p{31.188mm}>{\raggedright\arraybackslash}p{31.188mm}>{\raggedright\arraybackslash}p{31.188mm}@{}}
\toprule
\textbf{Case} & \textbf{Masculine plural} & \textbf{Neuter plural} & \textbf{Feminine plural} \\
\midrule
Nominative & ungeðancfulle & ungeðancfull & ungeðancfulla \\
Genitive & ungeðancfulra & ungeðancfulra & ungeðancfulra \\
Dative & ungeðancfullum & ungeðancfullum & ungeðancfullum \\
Accusative & ungeðancfulle & ungeðancfull & ungeðancfulla \\
\bottomrule
\end{tabular}
\par\endgroup
% course-panel:T69-P02:end
% mdtable:T69:end

% teaching: NEW-U15-PAR-UNGETHANCFULL-WEAK-NGDA
\reserve{11\baselineskip}
\reserve{16\baselineskip}
% mdparagraph:P189:begin
\textbf{Ungeðancfull: weak}
% mdparagraph:P189:end

% table: NEW-U15-PAR-UNGETHANCFULL-WEAK-NGDA-DISPLAY
% mdtable:T70:begin
% course-table-style: 2026-10
% course-panel:T70-P01:begin
\par\begingroup
\fontsize{10.5}{13.2}\selectfont
\setlength{\LTleft}{0pt}\setlength{\LTright}{\fill}
\noindent
\begin{tabular}{@{}>{\raggedright\arraybackslash}p{22.0mm}>{\raggedright\arraybackslash}p{31.188mm}>{\raggedright\arraybackslash}p{31.188mm}>{\raggedright\arraybackslash}p{31.188mm}@{}}
\toprule
\textbf{Case} & \textbf{Masculine singular} & \textbf{Neuter singular} & \textbf{Feminine singular} \\
\midrule
Nominative & ungeðancfulla & ungeðancfulle & ungeðancfulle \\
Genitive & ungeðancfullan & ungeðancfullan & ungeðancfullan \\
Dative & ungeðancfullan & ungeðancfullan & ungeðancfullan \\
Accusative & ungeðancfullan & ungeðancfulle & ungeðancfullan \\
\bottomrule
\end{tabular}
\par\endgroup
% course-panel:T70-P01:end
\par\addvspace{5pt}
% course-panel:T70-P02:begin
\par\begingroup
\fontsize{10.5}{13.2}\selectfont
\setlength{\LTleft}{0pt}\setlength{\LTright}{\fill}
\noindent
\begin{tabular}{@{}>{\raggedright\arraybackslash}p{22.0mm}>{\raggedright\arraybackslash}p{31.188mm}>{\raggedright\arraybackslash}p{31.188mm}>{\raggedright\arraybackslash}p{31.188mm}@{}}
\toprule
\textbf{Case} & \textbf{Masculine plural} & \textbf{Neuter plural} & \textbf{Feminine plural} \\
\midrule
Nominative & ungeðancfullan & ungeðancfullan & ungeðancfullan \\
Genitive & ungeðancfulra & ungeðancfulra & ungeðancfulra \\
Dative & ungeðancfullum & ungeðancfullum & ungeðancfullum \\
Accusative & ungeðancfullan & ungeðancfullan & ungeðancfullan \\
\bottomrule
\end{tabular}
\par\endgroup
% course-panel:T70-P02:end
% mdtable:T70:end

\reserve{7\baselineskip}
% mdparagraph:P190:begin
\subsection*{Additional vocabulary}\BookHeading{nav-ch15-learner-0026}{subsubsection}{Additional vocabulary}{0}
% mdparagraph:P190:end

% mdparagraph:P191:begin
Retain the earlier inflections of the familiar words below. Their stated senses are needed here; the adverbs do not agree with a noun.
% mdparagraph:P191:end

% teaching: NEW-U15-VOC-IDEL
% mdparagraph:P192:begin
\textbf{īdel} — idle; not working.
% mdparagraph:P192:end

% teaching: NEW-U15-VOC-LYTEL-SUPERL
% mdparagraph:P193:begin
\textbf{lǣst} — least in rank or importance.
% mdparagraph:P193:end

% teaching: NEW-U15-VOC-GOD-ADJECTIVE
% mdparagraph:P194:begin
\textbf{gōd} — kind; benevolent.
% mdparagraph:P194:end

% teaching: NEW-U15-VOC-YFEL
% mdparagraph:P195:begin
\textbf{yfel} — evil; evil people.
% mdparagraph:P195:end

% teaching: NEW-U15-VOC-SLEAN
% mdparagraph:P196:begin
\textbf{slēan} — kill.
% mdparagraph:P196:end

% teaching: NEW-U15-RECORD-GELICE-U14-SOURCE-1
% teaching: NEW-U15-VOC-GELICE
% mdparagraph:P197:begin
\textbf{gelīce} — likewise; similarly.
% mdparagraph:P197:end

% teaching: NEW-U15-RECORD-EORNOSTLICE-INVARIANT
% teaching: NEW-U15-VOC-EORNOSTLICE
% mdparagraph:P198:begin
\textbf{eornostlice} — accordingly; thus.
% mdparagraph:P198:end

% teaching: NEW-U15-RECORD-THAERTO-EACAN-INVARIANT
% teaching: NEW-U15-VOC-THAERTO-EACAN
% mdparagraph:P199:begin
\textbf{ðǣr-tō-ēacan} — in addition to that; besides.
% mdparagraph:P199:end

% teaching: NEW-U15-VOC-DON
% mdparagraph:P200:begin
\textbf{dōn} — make; render, with agreeing object predicate.
% mdparagraph:P200:end

% teaching: NEW-U15-RECORD-UT-INVARIANT
% teaching: NEW-U15-VOC-UT
% mdparagraph:P201:begin
\textbf{ūt} — out.
% mdparagraph:P201:end

% teaching: NEW-U15-RECORD-YMBE-INVARIANT
% teaching: NEW-U15-VOC-YMBE
% mdparagraph:P202:begin
\textbf{ymbe} — about, of approximate time, with accusative. Use accusative for the approximate time, including an hour.
% mdparagraph:P202:end

% teaching: NEW-U15-VOC-WITH
% mdparagraph:P203:begin
\textbf{wið} — in exchange or return for, with dative.
% mdparagraph:P203:end

% teaching: NEW-U15-VOC-FOR
% mdparagraph:P204:begin
\textbf{for} — for; in exchange for, with dative.
% mdparagraph:P204:end

% teaching: NEW-U15-VOC-OFER
% mdparagraph:P205:begin
\textbf{ofer} — above in rank, with accusative; toward, in gōd ofer persons, with accusative.
% mdparagraph:P205:end

% teaching: NEW-U15-VOC-THURH
% mdparagraph:P206:begin
\textbf{ðurh} — by, naming an oath standard, with accusative. The spelling ðurh is supplied beside the earlier þurh.
% mdparagraph:P206:end

\reserve{7\baselineskip}
% mdparagraph:P207:begin
\subsection*{Government and clause relations}\BookHeading{nav-ch15-learner-0027}{subsubsection}{Government and clause relations}{0}
% mdparagraph:P207:end

% teaching: NEW-U15-GRAM-U15-FOR-EXCHANGE
% teaching: NEW-U15-GRAM-U15-FORGIELDAN-PAYMENT
% teaching: NEW-U15-GRAM-U15-ONFON-PAYMENT-ACC
% teaching: NEW-U15-GRAM-U15-WID-EXCHANGE
% mdparagraph:P208:begin
\textbf{Payment and exchange.} Forgieldan takes the person paid in the dative and the thing paid in the accusative. Onfōn can take an accusative material payment, including pęning and substantive māre. This payment construction stands alongside the earlier dative use of “receive”. Wið with the dative gives payment or return; for with the dative gives equivalent exchange, as in an eye-for-eye formula. The first noun in that quoted formula is nominative.
% mdparagraph:P208:end

% teaching: NEW-U15-GRAM-U15-GELIC-OBJECT-PREDICATE
% mdparagraph:P209:begin
\textbf{Making equal.}\BookPoint{bk-ch15-learner-249aa2db805e} Dōn takes an accusative object and an agreeing predicate gelīc. The person used as the comparison standard is dative. Distinguish the object’s agreement from the comparison person’s case.
% mdparagraph:P209:end

% teaching: NEW-U15-GRAM-U15-WENAN-CONTENT
% mdparagraph:P210:begin
\textbf{Content after wēnan.} Wēnan can introduce a ðæt clause giving the thought or expectation. In a past report, sceolden, the preterite subjunctive plural of sculan, may express an expected outcome, “that they would receive”, rather than a duty.
% mdparagraph:P210:end

% teaching: NEW-U15-GRAM-U15-OFER-KINDNESS-RANK
% mdparagraph:P211:begin
\textbf{Kindness and rank.} Gōd ofer with the accusative means “kind toward” people. The bare strong accusative plurals ungeðancfulle and yfele can denote “unthankful people” and “evil people” in this use. Ofer with an accusative teacher or master can also mean “above” in rank. Neither use requires physical movement across something.
% mdparagraph:P211:end

% teaching: NEW-U15-GRAM-U15-SWERIAN-THURH
% mdparagraph:P212:begin
\textbf{An oath.} Swęrian with ðurh plus accusative names what one swears by. When a negative command is followed by several ne ðurh phrases, each phrase shares that command; the repeated ne does not create several positive assertions.
% mdparagraph:P212:end

% teaching: NEW-U15-GRAM-U15-ALIESAN-SEPARATION
% teaching: NEW-U15-GRAM-U15-LAERAN-OBJECTS
% mdparagraph:P213:begin
\textbf{Teaching and deliverance.} Lǣran takes an accusative person or the accusative matter taught. Ālīes takes the person delivered in the accusative; fram with the dative names what the person is delivered from. The passive teaching expression in Cook 31.3 and the exact substantive dative yfele, “evil”, in 31.4 remain supplied there.
% mdparagraph:P213:end

\reserve{7\baselineskip}
% mdparagraph:P214:begin
\subsection*{Translated examples — new teaching material}\BookHeading{nav-ch15-learner-0028}{subsubsection}{Translated examples — new teaching material}{0}
% mdparagraph:P214:end

% mdparagraph:P215:begin
These examples distinguish meanings and constructions; they are not exercises.
% mdparagraph:P215:end

% teaching: NEW-U15-EXAMPLE-01
% teaching: NEW-U15-GRAM-U15-NEGATIVE-CONCORD
% teaching: NEW-U15-PAR-NAN-STRONG-NGDA
% teaching: NEW-U15-VOC-NAN
% mdparagraph:P216:begin
\textbf{1.} Nān cyning ne gǣð.
% mdparagraph:P216:end

% mdparagraph:P217:begin
No king goes.
% mdparagraph:P217:end

% mdparagraph:P218:begin
The negative determiner and finite ne express one denial.
% mdparagraph:P218:end

% teaching: NEW-U15-EXAMPLE-02
% teaching: NEW-U15-GRAM-U15-NABBAN-INDICATIVE
% teaching: NEW-U15-PAR-NABBAN-IND
% teaching: NEW-U15-VOC-NABBAN
% mdparagraph:P219:begin
\textbf{2.} Se mann ne hæfð męte; se mann næfð męte.
% mdparagraph:P219:end

% mdparagraph:P220:begin
The man does not have food; the man does not have food.
% mdparagraph:P220:end

% mdparagraph:P221:begin
The two clauses have the same negative meaning. The second uses the selected contracted finite form.
% mdparagraph:P221:end

% teaching: NEW-U15-EXAMPLE-03
% mdparagraph:P222:begin
\textbf{3.} Hwā hæfð ðæt hūs?
% mdparagraph:P222:end

% mdparagraph:P223:begin
Who has that house?
% mdparagraph:P223:end

% mdparagraph:P224:begin
Hwā asks for the person; this is an ordinary direct question.
% mdparagraph:P224:end

% teaching: NEW-U15-EXAMPLE-04
% teaching: NEW-U15-GRAM-U15-HWA-INDEFINITE
% teaching: NEW-U15-VOC-HWA
% mdparagraph:P225:begin
\textbf{4.} Gif hwā gā, hē gesiehð ðæt hūs.
% mdparagraph:P225:end

% mdparagraph:P226:begin
If anyone goes, he will see that house.
% mdparagraph:P226:end

% mdparagraph:P227:begin
Hwā is indefinite in the contingent condition. Gā is third-singular present subjunctive here; the following present states the prospective result.
% mdparagraph:P227:end

% teaching: NEW-U15-EXAMPLE-05
% mdparagraph:P228:begin
\textbf{5.} Ǣlc mann hæfð sunu.
% mdparagraph:P228:end

% mdparagraph:P229:begin
Each man has a son.
% mdparagraph:P229:end

% mdparagraph:P230:begin
Singular mann and hæfð express distribution to each individual.
% mdparagraph:P230:end

% teaching: NEW-U15-EXAMPLE-06
% teaching: NEW-U15-GRAM-U15-QUANTITY-PARTITIVES
% teaching: NEW-U15-NOM-BOC
% teaching: NEW-U15-VOC-BOC
% mdparagraph:P231:begin
\textbf{6.} Se cyning ānum ðissa manna bōc sęlð.
% mdparagraph:P231:end

% mdparagraph:P232:begin
The king gives a book to one of these men.
% mdparagraph:P232:end

% mdparagraph:P233:begin
The selected recipient is dative ānum; the group is genitive ðissa manna. The book is the accusative thing given.
% mdparagraph:P233:end

% teaching: NEW-U15-EXAMPLE-07
% teaching: NEW-U15-GRAM-U15-PROPORTIONAL-SWA
% teaching: NEW-U15-VOC-SWA
% mdparagraph:P234:begin
\textbf{7.} Swā se cyning ieldra is, swā hē wīsra is.
% mdparagraph:P234:end

% mdparagraph:P235:begin
The older the king is, the wiser he is.
% mdparagraph:P235:end

% mdparagraph:P236:begin
The paired clauses correlate degrees. Each comparative agrees with the subject of its own clause.
% mdparagraph:P236:end

% teaching: NEW-U15-EXAMPLE-08
% teaching: NEW-U15-NOM-BOC
% teaching: NEW-U15-VOC-BOC
% mdparagraph:P237:begin
\textbf{8.} Ne sęle ðǣm cyninge ðā bōc.
% mdparagraph:P237:end

% mdparagraph:P238:begin
Do not give the book to the king.
% mdparagraph:P238:end

% mdparagraph:P239:begin
Ne negates a singular command. The king remains the dative recipient and the book the accusative thing.
% mdparagraph:P239:end

\reserve{7\baselineskip}
% mdparagraph:P240:begin
\section*{6. Practice — new teaching material}\BookHeading{nav-ch15-learner-0029}{subsection}{6. Practice — new teaching material}{1}
% mdparagraph:P240:end

% mdparagraph:P241:begin
The form grids identify each requested position even when several answers have the same spelling. Give English explanations for the analysis questions and Old English for composition. In composition, the directions specify the construction being practised. A different grammatical expression may convey the meaning without meeting that particular direction. Keep all supplied help visible.
% mdparagraph:P241:end

\reserve{7\baselineskip}
% mdparagraph:P242:begin
\subsection*{1. Negative propositions and complete indicative systems}\BookHeading{nav-ch15-learner-0030}{subsubsection}{1. Negative propositions and complete indicative systems}{0}
% mdparagraph:P242:end

\reserve{5\baselineskip}
\par\begin{minipage}{\linewidth}\raggedright
% mdparagraph:P243:begin
\textbf{F01}
% mdparagraph:P243:end

% item: NEW-U15-F01
% mdparagraph:P244:begin
Give the complete present and preterite indicative negative be-series for I, you singular, he, we, you plural and they. Use the selected contracted forms where the table gives them; retain separate ne where it does not. Explain why these are twelve logical positions even where forms coincide.
% mdparagraph:P244:end

% mdparagraph:P245:begin
\textbf{Response grid: wesan}
% mdparagraph:P245:end

% output: NEW-U15-F01-O001
% output: NEW-U15-F01-O002
% output: NEW-U15-F01-O003
% output: NEW-U15-F01-O004
% output: NEW-U15-F01-O005
% output: NEW-U15-F01-O006
% output: NEW-U15-F01-O007
% output: NEW-U15-F01-O008
% output: NEW-U15-F01-O009
% output: NEW-U15-F01-O010
% output: NEW-U15-F01-O011
% output: NEW-U15-F01-O012
% mdtable:T71:begin
% course-table-style: 2026-10
% course-panel:T71-P01:begin
\par\begingroup
\fontsize{10.5}{13.2}\selectfont
\setlength{\LTleft}{0pt}\setlength{\LTright}{\fill}
\noindent
\begin{tabular}{@{}>{\raggedright\arraybackslash}p{39.459mm}>{\raggedright\arraybackslash}p{39.459mm}>{\raggedright\arraybackslash}p{39.459mm}@{}}
\toprule
\textbf{Person and number} & \textbf{Present} & \textbf{Preterite} \\
\midrule
1 sg. & (1) \_\_\_\_ & (7) \_\_\_\_ \\
2 sg. & (2) \_\_\_\_ & (8) \_\_\_\_ \\
3 sg. & (3) \_\_\_\_ & (9) \_\_\_\_ \\
1 pl. & (4) \_\_\_\_ & (10) \_\_\_\_ \\
2 pl. & (5) \_\_\_\_ & (11) \_\_\_\_ \\
3 pl. & (6) \_\_\_\_ & (12) \_\_\_\_ \\
\bottomrule
\end{tabular}
\par\endgroup
% course-panel:T71-P01:end
% mdtable:T71:end
% output: NEW-U15-F01-O013
% mdparagraph:P246:begin
\textbf{Then explain:} (13) why twelve positions remain distinct; where ne stays separate.
% mdparagraph:P246:end

\end{minipage}\par\addvspace{4pt}

\par\begin{minipage}{\linewidth}\raggedright
% mdparagraph:P247:begin
\textbf{F02}
% mdparagraph:P247:end

% item: NEW-U15-F02
% mdparagraph:P248:begin
Give the complete present and preterite indicative of nabban for the six person/\allowbreak{}number positions. Interpret the first-person singular and third-person plural in each tense as denial of possession.
% mdparagraph:P248:end

% mdparagraph:P249:begin
\textbf{Response grid: nabban}
% mdparagraph:P249:end

% output: NEW-U15-F02-O001
% output: NEW-U15-F02-O002
% output: NEW-U15-F02-O003
% output: NEW-U15-F02-O004
% output: NEW-U15-F02-O005
% output: NEW-U15-F02-O006
% output: NEW-U15-F02-O007
% output: NEW-U15-F02-O008
% output: NEW-U15-F02-O009
% output: NEW-U15-F02-O010
% output: NEW-U15-F02-O011
% output: NEW-U15-F02-O012
% mdtable:T72:begin
% course-table-style: 2026-10
% course-panel:T72-P01:begin
\par\begingroup
\fontsize{10.5}{13.2}\selectfont
\setlength{\LTleft}{0pt}\setlength{\LTright}{\fill}
\noindent
\begin{tabular}{@{}>{\raggedright\arraybackslash}p{39.459mm}>{\raggedright\arraybackslash}p{39.459mm}>{\raggedright\arraybackslash}p{39.459mm}@{}}
\toprule
\textbf{Person and number} & \textbf{Present} & \textbf{Preterite} \\
\midrule
1 sg. & (1) \_\_\_\_ & (7) \_\_\_\_ \\
2 sg. & (2) \_\_\_\_ & (8) \_\_\_\_ \\
3 sg. & (3) \_\_\_\_ & (9) \_\_\_\_ \\
1 pl. & (4) \_\_\_\_ & (10) \_\_\_\_ \\
2 pl. & (5) \_\_\_\_ & (11) \_\_\_\_ \\
3 pl. & (6) \_\_\_\_ & (12) \_\_\_\_ \\
\bottomrule
\end{tabular}
\par\endgroup
% course-panel:T72-P01:end
% mdtable:T72:end
% output: NEW-U15-F02-O013
% output: NEW-U15-F02-O014
% output: NEW-U15-F02-O015
% output: NEW-U15-F02-O016
% mdparagraph:P250:begin
\textbf{Then explain:} (13) meaning of present first singular; (14) meaning of present third plural; (15) meaning of preterite first singular; (16) meaning of preterite third plural.
% mdparagraph:P250:end

\end{minipage}\par\addvspace{4pt}

\par\begin{minipage}{\linewidth}\raggedright
% mdparagraph:P251:begin
\textbf{F03}
% mdparagraph:P251:end

% item: NEW-U15-F03
% mdparagraph:P252:begin
Give the complete present and preterite indicative of the selected negative willan system. Interpret the first-person singular and second-person singular of each tense as refusal to go.
% mdparagraph:P252:end

% mdparagraph:P253:begin
\textbf{Response grid: willan}
% mdparagraph:P253:end

% output: NEW-U15-F03-O001
% output: NEW-U15-F03-O002
% output: NEW-U15-F03-O003
% output: NEW-U15-F03-O004
% output: NEW-U15-F03-O005
% output: NEW-U15-F03-O006
% output: NEW-U15-F03-O007
% output: NEW-U15-F03-O008
% output: NEW-U15-F03-O009
% output: NEW-U15-F03-O010
% output: NEW-U15-F03-O011
% output: NEW-U15-F03-O012
% mdtable:T73:begin
% course-table-style: 2026-10
% course-panel:T73-P01:begin
\par\begingroup
\fontsize{10.5}{13.2}\selectfont
\setlength{\LTleft}{0pt}\setlength{\LTright}{\fill}
\noindent
\begin{tabular}{@{}>{\raggedright\arraybackslash}p{39.459mm}>{\raggedright\arraybackslash}p{39.459mm}>{\raggedright\arraybackslash}p{39.459mm}@{}}
\toprule
\textbf{Person and number} & \textbf{Present} & \textbf{Preterite} \\
\midrule
1 sg. & (1) \_\_\_\_ & (7) \_\_\_\_ \\
2 sg. & (2) \_\_\_\_ & (8) \_\_\_\_ \\
3 sg. & (3) \_\_\_\_ & (9) \_\_\_\_ \\
1 pl. & (4) \_\_\_\_ & (10) \_\_\_\_ \\
2 pl. & (5) \_\_\_\_ & (11) \_\_\_\_ \\
3 pl. & (6) \_\_\_\_ & (12) \_\_\_\_ \\
\bottomrule
\end{tabular}
\par\endgroup
% course-panel:T73-P01:end
% mdtable:T73:end
% output: NEW-U15-F03-O013
% output: NEW-U15-F03-O014
% output: NEW-U15-F03-O015
% output: NEW-U15-F03-O016
% mdparagraph:P254:begin
\textbf{Then explain:} (13) meaning of present first singular; (14) meaning of present second singular; (15) meaning of preterite first singular; (16) meaning of preterite second singular.
% mdparagraph:P254:end

\end{minipage}\par\addvspace{4pt}

\par\begin{minipage}{\linewidth}\raggedright
% mdparagraph:P255:begin
\textbf{A01}
% mdparagraph:P255:end

% item: NEW-U15-A01
% mdparagraph:P256:begin
Interpret the three clauses. Identify every negative word and state the single proposition denied in each; do not turn the accumulated negatives into an affirmation.
% mdparagraph:P256:end

% mdparagraph:P257:begin
(1) Nān mann ne cymð.
% mdparagraph:P257:end

% mdparagraph:P258:begin
(2) Nǣfre ne cōm se cyning.
% mdparagraph:P258:end

% mdparagraph:P259:begin
(3) Nān mann nāwiht ne hæfð.
% mdparagraph:P259:end

% output: NEW-U15-A01-O001
% output: NEW-U15-A01-O002
% output: NEW-U15-A01-O003
% output: NEW-U15-A01-O004
% output: NEW-U15-A01-O005
% output: NEW-U15-A01-O006
% mdparagraph:P260:begin
\textbf{Responses, in order:} (1) clause 1: negative words; (2) clause 1: proposition denied; (3) clause 2: negative words; (4) clause 2: proposition denied; (5) clause 3: negative words; (6) clause 3: proposition denied.
% mdparagraph:P260:end

\end{minipage}\par\addvspace{4pt}

\par\begin{minipage}{\linewidth}\raggedright
% mdparagraph:P261:begin
\textbf{A02}
% mdparagraph:P261:end

% item: NEW-U15-A02
% mdparagraph:P262:begin
Compare the two clauses. Say whether the verb or one contrasted participant is negated, and identify the contrasting participant.
% mdparagraph:P262:end

% mdparagraph:P263:begin
(1) Ic self gā, and nā ōðer.
% mdparagraph:P263:end

% mdparagraph:P264:begin
(2) Ic ne gā.
% mdparagraph:P264:end

% output: NEW-U15-A02-O001
% output: NEW-U15-A02-O002
% output: NEW-U15-A02-O003
% mdparagraph:P265:begin
\textbf{Responses, in order:} (1) first clause: scope of negation; (2) first clause: contrasting participant; (3) second clause: scope of negation.
% mdparagraph:P265:end

\end{minipage}\par\addvspace{4pt}

\par\begin{minipage}{\linewidth}\raggedright
% mdparagraph:P266:begin
\textbf{C02}
% mdparagraph:P266:end

% item: NEW-U15-C02
% mdparagraph:P267:begin
Translate: I did not have the book, and I did not wish to go. Use the selected contracted negative preterites of have and will.
% mdparagraph:P267:end

% output: NEW-U15-C02-O001
\end{minipage}\par\addvspace{4pt}

\par\begin{minipage}{\linewidth}\raggedright
% mdparagraph:P268:begin
\textbf{C14}
% mdparagraph:P268:end

% item: NEW-U15-C14
% mdparagraph:P269:begin
Translate: He himself goes, and not another. The referent is a man. Use self and constituent nā.
% mdparagraph:P269:end

% output: NEW-U15-C14-O001
\end{minipage}\par\addvspace{4pt}

\reserve{7\baselineskip}
\par\begin{minipage}{\linewidth}\raggedright
% mdparagraph:P270:begin
\subsection*{2. Negative, indefinite and distributive reference}\BookHeading{nav-ch15-learner-0031}{subsubsection}{2. Negative, indefinite and distributive reference}{0}
% mdparagraph:P270:end


% mdparagraph:P271:begin
\textbf{F04}
% mdparagraph:P271:end

% item: NEW-U15-F04
% mdparagraph:P272:begin
Give all nominative, genitive, dative and accusative forms of nān, singular and plural, in masculine, feminine and neuter. Explain why a form such as nānum needs its noun or sentence to identify the particular case/\allowbreak{}number/\allowbreak{}gender.
% mdparagraph:P272:end

% mdparagraph:P273:begin
\textbf{Response grid: nān}
% mdparagraph:P273:end

% output: NEW-U15-F04-O001
% output: NEW-U15-F04-O002
% output: NEW-U15-F04-O003
% output: NEW-U15-F04-O004
% output: NEW-U15-F04-O005
% output: NEW-U15-F04-O006
% output: NEW-U15-F04-O007
% output: NEW-U15-F04-O008
% output: NEW-U15-F04-O009
% output: NEW-U15-F04-O010
% output: NEW-U15-F04-O011
% output: NEW-U15-F04-O012
% output: NEW-U15-F04-O013
% output: NEW-U15-F04-O014
% output: NEW-U15-F04-O015
% output: NEW-U15-F04-O016
% output: NEW-U15-F04-O017
% output: NEW-U15-F04-O018
% output: NEW-U15-F04-O019
% output: NEW-U15-F04-O020
% output: NEW-U15-F04-O021
% output: NEW-U15-F04-O022
% output: NEW-U15-F04-O023
% output: NEW-U15-F04-O024
% mdtable:T74:begin
% course-table-style: 2026-10
% course-panel:T74-P01:begin
\par\begingroup
\fontsize{10.5}{13.2}\selectfont
\setlength{\LTleft}{0pt}\setlength{\LTright}{\fill}
\noindent
\begin{tabular}{@{}>{\raggedright\arraybackslash}p{22.0mm}>{\raggedright\arraybackslash}p{31.188mm}>{\raggedright\arraybackslash}p{31.188mm}>{\raggedright\arraybackslash}p{31.188mm}@{}}
\toprule
\textbf{Case} & \textbf{Masculine singular} & \textbf{Neuter singular} & \textbf{Feminine singular} \\
\midrule
Nominative & (1) \_\_\_\_ & (17) \_\_\_\_ & (9) \_\_\_\_ \\
Genitive & (2) \_\_\_\_ & (18) \_\_\_\_ & (10) \_\_\_\_ \\
Dative & (3) \_\_\_\_ & (19) \_\_\_\_ & (11) \_\_\_\_ \\
Accusative & (4) \_\_\_\_ & (20) \_\_\_\_ & (12) \_\_\_\_ \\
\bottomrule
\end{tabular}
\par\endgroup
% course-panel:T74-P01:end
\par\addvspace{5pt}
% course-panel:T74-P02:begin
\par\begingroup
\fontsize{10.5}{13.2}\selectfont
\setlength{\LTleft}{0pt}\setlength{\LTright}{\fill}
\noindent
\begin{tabular}{@{}>{\raggedright\arraybackslash}p{22.0mm}>{\raggedright\arraybackslash}p{31.188mm}>{\raggedright\arraybackslash}p{31.188mm}>{\raggedright\arraybackslash}p{31.188mm}@{}}
\toprule
\textbf{Case} & \textbf{Masculine plural} & \textbf{Neuter plural} & \textbf{Feminine plural} \\
\midrule
Nominative & (5) \_\_\_\_ & (21) \_\_\_\_ & (13) \_\_\_\_ \\
Genitive & (6) \_\_\_\_ & (22) \_\_\_\_ & (14) \_\_\_\_ \\
Dative & (7) \_\_\_\_ & (23) \_\_\_\_ & (15) \_\_\_\_ \\
Accusative & (8) \_\_\_\_ & (24) \_\_\_\_ & (16) \_\_\_\_ \\
\bottomrule
\end{tabular}
\par\endgroup
% course-panel:T74-P02:end
% mdtable:T74:end

% output: NEW-U15-F04-O025
% mdparagraph:P274:begin
\textbf{Then explain:} (25) how the sentence identifies the form’s case, number and gender.
% mdparagraph:P274:end

\end{minipage}\par\addvspace{4pt}

\par\begin{minipage}{\linewidth}\raggedright
% mdparagraph:P275:begin
\textbf{F05}
% mdparagraph:P275:end

% item: NEW-U15-F05
% mdparagraph:P276:begin
Give all nominative, genitive, dative and accusative singular forms of gehwā/\allowbreak{}gehwæt, distinguishing personal and neuter reference. Explain why no plural column is required for this system.
% mdparagraph:P276:end

% mdparagraph:P277:begin
\textbf{Response grid: gehwā /\allowbreak{} gehwæt}
% mdparagraph:P277:end

% output: NEW-U15-F05-O001
% output: NEW-U15-F05-O002
% output: NEW-U15-F05-O003
% output: NEW-U15-F05-O004
% output: NEW-U15-F05-O005
% output: NEW-U15-F05-O006
% output: NEW-U15-F05-O007
% output: NEW-U15-F05-O008
% mdtable:T75:begin
% course-table-style: 2026-10
% course-panel:T75-P01:begin
\par\begingroup
\fontsize{10.5}{13.2}\selectfont
\setlength{\LTleft}{0pt}\setlength{\LTright}{\fill}
\noindent
\begin{tabular}{@{}>{\raggedright\arraybackslash}p{39.459mm}>{\raggedright\arraybackslash}p{39.459mm}>{\raggedright\arraybackslash}p{39.459mm}@{}}
\toprule
\textbf{Case} & \textbf{Neuter singular} & \textbf{Common singular} \\
\midrule
Nominative & (5) \_\_\_\_ & (1) \_\_\_\_ \\
Genitive & (6) \_\_\_\_ & (2) \_\_\_\_ \\
Dative & (7) \_\_\_\_ & (3) \_\_\_\_ \\
Accusative & (8) \_\_\_\_ & (4) \_\_\_\_ \\
\bottomrule
\end{tabular}
\par\endgroup
% course-panel:T75-P01:end
% mdtable:T75:end
% output: NEW-U15-F05-O009
% mdparagraph:P278:begin
\textbf{Then explain:} (9) reason no plural column is needed.
% mdparagraph:P278:end

\end{minipage}\par\addvspace{4pt}

\par\begin{minipage}{\linewidth}\raggedright
% mdparagraph:P279:begin
\textbf{A03}
% mdparagraph:P279:end

% item: NEW-U15-A03
% mdparagraph:P280:begin
For each clause identify whether hwā is interrogative, indefinite, or part of a free-choice expression; give the clause meaning and identify its subject.
% mdparagraph:P280:end

% mdparagraph:P281:begin
(1) Hwā cymð?
% mdparagraph:P281:end

% mdparagraph:P282:begin
(2) Gif hwā cume, ic him bōc sęlle.
% mdparagraph:P282:end

% mdparagraph:P283:begin
(3) Swā hwā swā cume, ic him bōc sęlle.
% mdparagraph:P283:end

% output: NEW-U15-A03-O001
% output: NEW-U15-A03-O002
% output: NEW-U15-A03-O003
% output: NEW-U15-A03-O004
% output: NEW-U15-A03-O005
% output: NEW-U15-A03-O006
% output: NEW-U15-A03-O007
% output: NEW-U15-A03-O008
% output: NEW-U15-A03-O009
% mdparagraph:P284:begin
\textbf{Responses, in order:} (1) clause 1: function of hwā; (2) clause 1: meaning; (3) clause 1: subject; (4) clause 2: function of hwā; (5) clause 2: meaning; (6) clause 2: subject; (7) clause 3: function of hwā; (8) clause 3: meaning; (9) clause 3: subject.
% mdparagraph:P284:end

\end{minipage}\par\addvspace{4pt}

\par\begin{minipage}{\linewidth}\raggedright
% mdparagraph:P285:begin
\textbf{A04}
% mdparagraph:P285:end

% item: NEW-U15-A04
% mdparagraph:P286:begin
Compare the two statements. Explain each versus all, identify the head of each subject, and account for the finite-verb number.
% mdparagraph:P286:end

% mdparagraph:P287:begin
(1) Ǣlc mann hæfð ānne pęning.
% mdparagraph:P287:end

% mdparagraph:P288:begin
(2) Ealle ðā męnn habbað męte.
% mdparagraph:P288:end

% output: NEW-U15-A04-O001
% output: NEW-U15-A04-O002
% output: NEW-U15-A04-O003
% output: NEW-U15-A04-O004
% mdparagraph:P289:begin
\textbf{Responses, in order:} (1) clause 1: quantifier meaning; (2) clause 1: subject head and verb agreement; (3) clause 2: quantifier meaning; (4) clause 2: subject head and verb agreement.
% mdparagraph:P289:end

\end{minipage}\par\addvspace{4pt}

\par\begin{minipage}{\linewidth}\raggedright
% mdparagraph:P290:begin
\textbf{A06}
% mdparagraph:P290:end

% item: NEW-U15-A06
% mdparagraph:P291:begin
Compare the two uses of man/\allowbreak{}mann. Identify the subject and finite-verb number, and explain why the impersonal use does not introduce a new plural pronoun.
% mdparagraph:P291:end

% mdparagraph:P292:begin
(1) Man sceal gān.
% mdparagraph:P292:end

% mdparagraph:P293:begin
(2) Se mann gǣð.
% mdparagraph:P293:end

% output: NEW-U15-A06-O001
% output: NEW-U15-A06-O002
% output: NEW-U15-A06-O003
% output: NEW-U15-A06-O004
% output: NEW-U15-A06-O005
% mdparagraph:P294:begin
\textbf{Responses, in order:} (1) clause 1: subject and reference; (2) clause 1: verb number; (3) clause 2: subject and reference; (4) clause 2: verb number; (5) reason this does not supply a new plural pronoun.
% mdparagraph:P294:end

\end{minipage}\par\addvspace{4pt}

\par\begin{minipage}{\linewidth}\raggedright
% mdparagraph:P295:begin
\textbf{A10}
% mdparagraph:P295:end

% item: NEW-U15-A10
% mdparagraph:P296:begin
Using the supplied meanings, interpret each clause and say whether it asks a question, denies possession, or asserts possession. Explain how nāwiht and ne work together in the middle clause.
% mdparagraph:P296:end

% mdparagraph:P297:begin
\textbf{Modern help — available throughout this item.}
% mdparagraph:P297:end

\end{minipage}\par\addvspace{4pt}

% teaching: NEW-U15-HELP-A10-INDEFINITES
% teaching: NEW-U15-VOC-AWIHT
% teaching: NEW-U15-VOC-HWAETHWEGA
% teaching: NEW-U15-VOC-NAWIHT
% support: NEW-U15-HELP-A10-INDEFINITES
% mdparagraph:P298:begin
āwiht anything; nāwiht nothing; hwæthwega something. Use these meanings to interpret the three clauses.
% mdparagraph:P298:end

% mdparagraph:P299:begin
(1) Hæfð hē āwiht?
% mdparagraph:P299:end

% mdparagraph:P300:begin
(2) Hē nāwiht ne hæfð.
% mdparagraph:P300:end

% mdparagraph:P301:begin
(3) Hē hæfð hwæthwega.
% mdparagraph:P301:end

% output: NEW-U15-A10-O001
% output: NEW-U15-A10-O002
% output: NEW-U15-A10-O003
% output: NEW-U15-A10-O004
% mdparagraph:P302:begin
\textbf{Responses, in order:} (1) clause 1: meaning and statement type; (2) clause 2: meaning and statement type; (3) clause 3: meaning and statement type; (4) how the two negatives work together.
% mdparagraph:P302:end

\reserve{5\baselineskip}
\par\begin{minipage}{\linewidth}\raggedright
% mdparagraph:P303:begin
\textbf{A11}
% mdparagraph:P303:end

% item: NEW-U15-A11
% mdparagraph:P304:begin
The same words Ān mann cōm occur in two settings: first, a story introduces an unnamed visitor; second, an explicit count contrasts one man with two. Give the sense of ān in each setting, and explain why the nominative agreement is the same.
% mdparagraph:P304:end

% mdparagraph:P305:begin
(1) Ān mann cōm.
% mdparagraph:P305:end

% mdparagraph:P306:begin
(2) Ān mann cōm.
% mdparagraph:P306:end

% output: NEW-U15-A11-O001
% output: NEW-U15-A11-O002
% output: NEW-U15-A11-O003
% mdparagraph:P307:begin
\textbf{Responses, in order:} (1) ān in the first setting; (2) ān in the second setting; (3) reason for the same agreement.
% mdparagraph:P307:end

\end{minipage}\par\addvspace{4pt}

\par\begin{minipage}{\linewidth}\raggedright
% mdparagraph:P308:begin
\textbf{C01}
% mdparagraph:P308:end

% item: NEW-U15-C01
% mdparagraph:P309:begin
Translate: No man has a book; each has food. Use nān and personal gehwā.
% mdparagraph:P309:end

% output: NEW-U15-C01-O001
\end{minipage}\par\addvspace{4pt}

\par\begin{minipage}{\linewidth}\raggedright
% mdparagraph:P310:begin
\textbf{C03}
% mdparagraph:P310:end

% item: NEW-U15-C03
% mdparagraph:P311:begin
Translate: Has the king any book? He has no book. Use ǣnig in the question and nān with finite negation in the reply.
% mdparagraph:P311:end

% output: NEW-U15-C03-O001
\end{minipage}\par\addvspace{4pt}

\par\begin{minipage}{\linewidth}\raggedright
% mdparagraph:P312:begin
\textbf{C04}
% mdparagraph:P312:end

% item: NEW-U15-C04
% mdparagraph:P313:begin
Translate: If anyone has a book, I shall give him a penny. Use indefinite hwā. Treat the condition as a possibility and use the present subjunctive in the if-clause.
% mdparagraph:P313:end

% output: NEW-U15-C04-O001
\end{minipage}\par\addvspace{4pt}

\par\begin{minipage}{\linewidth}\raggedright
% mdparagraph:P314:begin
\textbf{C07}
% mdparagraph:P314:end

% item: NEW-U15-C07
% mdparagraph:P315:begin
Translate: One can go. Use impersonal man with magan.
% mdparagraph:P315:end

% output: NEW-U15-C07-O001
\end{minipage}\par\addvspace{4pt}

\par\begin{minipage}{\linewidth}\raggedright
% mdparagraph:P316:begin
\textbf{C08}
% mdparagraph:P316:end

% item: NEW-U15-C08
% mdparagraph:P317:begin
Translate: Each son has one book; all the sons have food. Use ǣlc for each and eall for all.
% mdparagraph:P317:end

% output: NEW-U15-C08-O001
\end{minipage}\par\addvspace{4pt}

\par\begin{minipage}{\linewidth}\raggedright
% mdparagraph:P318:begin
\textbf{C15}
% mdparagraph:P318:end

% item: NEW-U15-C15
% mdparagraph:P319:begin
Translate: I give nothing to anyone. Use nāwiht for the thing and nān for the recipient. Use finite ne with the two specified negative expressions.
% mdparagraph:P319:end

% output: NEW-U15-C15-O001
\end{minipage}\par\addvspace{4pt}

\par\begin{minipage}{\linewidth}\raggedright
% mdparagraph:P320:begin
\textbf{C16}
% mdparagraph:P320:end

% item: NEW-U15-C16
% mdparagraph:P321:begin
Translate: No man’s son comes. Use nān within the dependent genitive phrase, with finite ne.
% mdparagraph:P321:end

% output: NEW-U15-C16-O001
\end{minipage}\par\addvspace{4pt}

\par\begin{minipage}{\linewidth}\raggedright
% mdparagraph:P322:begin
\textbf{C17}
% mdparagraph:P322:end

% item: NEW-U15-C17
% mdparagraph:P323:begin
Translate: The king has no books. Use contracted negative possession and nān.
% mdparagraph:P323:end

% output: NEW-U15-C17-O001
\end{minipage}\par\addvspace{4pt}

\par\begin{minipage}{\linewidth}\raggedright
% mdparagraph:P324:begin
\textbf{C18}
% mdparagraph:P324:end

% item: NEW-U15-C18
% mdparagraph:P325:begin
Translate: The king sees each person; each person’s son is good. Use personal gehwā for both each-person phrases.
% mdparagraph:P325:end

% output: NEW-U15-C18-O001
\end{minipage}\par\addvspace{4pt}

\par\begin{minipage}{\linewidth}\raggedright
% mdparagraph:P326:begin
\textbf{C19}
% mdparagraph:P326:end

% item: NEW-U15-C19
% mdparagraph:P327:begin
Translate: A certain man came, and he had a book. Use indefinite ān to introduce the unnamed man.
% mdparagraph:P327:end

% output: NEW-U15-C19-O001
\end{minipage}\par\addvspace{4pt}

\reserve{7\baselineskip}
% mdparagraph:P328:begin
\subsection*{3. Free choice, case, parts and groups}\BookHeading{nav-ch15-learner-0032}{subsubsection}{3. Free choice, case, parts and groups}{0}
% mdparagraph:P328:end

\reserve{5\baselineskip}
\par\begin{minipage}{\linewidth}\raggedright
% mdparagraph:P329:begin
\textbf{A05}
% mdparagraph:P329:end

% item: NEW-U15-A05
% mdparagraph:P330:begin
Explain the contribution of ǣgðer ge … ge and identify the two coordinated constituents. State whether the statement chooses just one of them.
% mdparagraph:P330:end

% mdparagraph:P331:begin
(1) Ic lufie ǣgðer ge ðone cyning ge his sunu.
% mdparagraph:P331:end

% output: NEW-U15-A05-O001
% output: NEW-U15-A05-O002
% output: NEW-U15-A05-O003
% output: NEW-U15-A05-O004
% mdparagraph:P332:begin
\textbf{Responses, in order:} (1) meaning of the construction; (2) first coordinated constituent; (3) second coordinated constituent; (4) whether just one is selected.
% mdparagraph:P332:end

\end{minipage}\par\addvspace{4pt}

\par\begin{minipage}{\linewidth}\raggedright
% mdparagraph:P333:begin
\textbf{A08}
% mdparagraph:P333:end

% item: NEW-U15-A08
% mdparagraph:P334:begin
In Swā hwā swā cume, ic him bōc sęlle, identify the case required for the subject that comes and the case required for the recipient of the book. Explain how the explicit pronoun in the main clause keeps the two roles separate.
% mdparagraph:P334:end

% mdparagraph:P335:begin
(1) Swā hwā swā cume, ic him bōc sęlle.
% mdparagraph:P335:end

% output: NEW-U15-A08-O001
% output: NEW-U15-A08-O002
% output: NEW-U15-A08-O003
% mdparagraph:P336:begin
\textbf{Responses, in order:} (1) subject of cume and its case; (2) recipient of the book and its case; (3) function of the explicit main-clause pronoun.
% mdparagraph:P336:end

\end{minipage}\par\addvspace{4pt}

\par\begin{minipage}{\linewidth}\raggedright
% mdparagraph:P337:begin
\textbf{A09}
% mdparagraph:P337:end

% item: NEW-U15-A09
% mdparagraph:P338:begin
Treat ānum uncer as the recipient of a gift and uncer hūs as a subject phrase. Explain the difference between the personal genitive uncer in the first phrase and the agreeing possessive uncer in the second. Identify the role and case of ānum and the number of the group selected by uncer.
% mdparagraph:P338:end

% mdparagraph:P339:begin
(1) ānum uncer
% mdparagraph:P339:end

% mdparagraph:P340:begin
(2) uncer hūs
% mdparagraph:P340:end

% output: NEW-U15-A09-O001
% output: NEW-U15-A09-O002
% output: NEW-U15-A09-O003
% output: NEW-U15-A09-O004
% mdparagraph:P341:begin
\textbf{Responses, in order:} (1) first uncer: function; (2) second uncer: function; (3) ānum: role and case; (4) number of the group.
% mdparagraph:P341:end

\end{minipage}\par\addvspace{4pt}

\par\begin{minipage}{\linewidth}\raggedright
% mdparagraph:P342:begin
\textbf{C05}
% mdparagraph:P342:end

% item: NEW-U15-C05
% mdparagraph:P343:begin
Translate: Whoever goes, I shall give him a penny. Use swā hwā swā and the present subjunctive for the open choice of a possible person.
% mdparagraph:P343:end

% output: NEW-U15-C05-O001
\end{minipage}\par\addvspace{4pt}

\par\begin{minipage}{\linewidth}\raggedright
% mdparagraph:P344:begin
\textbf{C06}
% mdparagraph:P344:end

% item: NEW-U15-C06
% mdparagraph:P345:begin
Translate: Whichever of these two sons comes first, I will give him the book. Use swā hwæðer swā and the present subjunctive for the possible first arrival. Keep the subject of comes nominative and express the recipient again with a dative pronoun in the main clause.
% mdparagraph:P345:end

% output: NEW-U15-C06-O001
\end{minipage}\par\addvspace{4pt}

\par\begin{minipage}{\linewidth}\raggedright
% mdparagraph:P346:begin
\textbf{C09}
% mdparagraph:P346:end

% item: NEW-U15-C09
% mdparagraph:P347:begin
Translate: I see both the king and his son. Use ǣgðer ge … ge.
% mdparagraph:P347:end

% output: NEW-U15-C09-O001
\end{minipage}\par\addvspace{4pt}

\reserve{7\baselineskip}
% mdparagraph:P348:begin
\subsection*{4. Equality, preference and proportion}\BookHeading{nav-ch15-learner-0033}{subsubsection}{4. Equality, preference and proportion}{0}
% mdparagraph:P348:end

\reserve{5\baselineskip}
\par\begin{minipage}{\linewidth}\raggedright
% mdparagraph:P349:begin
\textbf{A07}
% mdparagraph:P349:end

% item: NEW-U15-A07
% mdparagraph:P350:begin
Distinguish the meanings of the three comparisons. Identify the comparison standard or correlated clause in each.
% mdparagraph:P350:end

% mdparagraph:P351:begin
(1) Ðēos bōc is swā gōd swā sēo bōc.
% mdparagraph:P351:end

% mdparagraph:P352:begin
(2) Ic lufie lēoht swīðor ðonne ðīestru.
% mdparagraph:P352:end

% mdparagraph:P353:begin
(3) Swā ðēos bōc māre is, swā hēo bętere is.
% mdparagraph:P353:end

% output: NEW-U15-A07-O001
% output: NEW-U15-A07-O002
% output: NEW-U15-A07-O003
% output: NEW-U15-A07-O004
% output: NEW-U15-A07-O005
% output: NEW-U15-A07-O006
% mdparagraph:P354:begin
\textbf{Responses, in order:} (1) comparison 1: meaning; (2) comparison 1: standard or related clause; (3) comparison 2: meaning; (4) comparison 2: standard or related clause; (5) comparison 3: meaning; (6) comparison 3: standard or related clause.
% mdparagraph:P354:end

\end{minipage}\par\addvspace{4pt}

\par\begin{minipage}{\linewidth}\raggedright
% mdparagraph:P355:begin
\textbf{C10}
% mdparagraph:P355:end

% item: NEW-U15-C10
% mdparagraph:P356:begin
Translate: The larger that house is, the better it is. Use proportional swā … swā and a pronoun referring back to the house.
% mdparagraph:P356:end

% output: NEW-U15-C10-O001
\end{minipage}\par\addvspace{4pt}

\par\begin{minipage}{\linewidth}\raggedright
% mdparagraph:P357:begin
\textbf{C11}
% mdparagraph:P357:end

% item: NEW-U15-C11
% mdparagraph:P358:begin
Translate: This book is as good as that book, but the king’s book is better. Use equal-degree swā … swā in the first clause.
% mdparagraph:P358:end

% output: NEW-U15-C11-O001
\end{minipage}\par\addvspace{4pt}

\par\begin{minipage}{\linewidth}\raggedright
% mdparagraph:P359:begin
\textbf{C12}
% mdparagraph:P359:end

% item: NEW-U15-C12
% mdparagraph:P360:begin
Translate: He loves food more than wine. Use preferential swīðor.
% mdparagraph:P360:end

% output: NEW-U15-C12-O001
\end{minipage}\par\addvspace{4pt}

\reserve{7\baselineskip}
\par\begin{minipage}{\linewidth}\raggedright
% mdparagraph:P361:begin
\subsection*{5. Noun systems and commands}\BookHeading{nav-ch15-learner-0034}{subsubsection}{5. Noun systems and commands}{0}
% mdparagraph:P361:end


% mdparagraph:P362:begin
\textbf{F06}
% mdparagraph:P362:end

% item: NEW-U15-F06
% mdparagraph:P363:begin
Give all nominative, genitive, dative and accusative singular/\allowbreak{}plural forms of bōc and tōð. Identify the forms that change the stem vowel and distinguish a syncretic surface from a uniquely identifiable case.
% mdparagraph:P363:end

% mdparagraph:P364:begin
\textbf{Response grid: bōc}
% mdparagraph:P364:end

% output: NEW-U15-F06-O001
% output: NEW-U15-F06-O002
% output: NEW-U15-F06-O003
% output: NEW-U15-F06-O004
% output: NEW-U15-F06-O005
% output: NEW-U15-F06-O006
% output: NEW-U15-F06-O007
% output: NEW-U15-F06-O008
% mdtable:T76:begin
\par\begingroup
\fontsize{10.5}{13.2}\selectfont
\setlength{\LTleft}{0pt}\setlength{\LTright}{\fill}
\noindent
\begin{tabular}{@{}>{\raggedright\arraybackslash}p{39.46mm}>{\raggedright\arraybackslash}p{39.46mm}>{\raggedright\arraybackslash}p{39.46mm}@{}}
\toprule
\textbf{Case} & \textbf{Singular} & \textbf{Plural} \\
\midrule
Nominative & (1) \_\_\_\_ & (5) \_\_\_\_ \\
Genitive & (2) \_\_\_\_ & (6) \_\_\_\_ \\
Dative & (3) \_\_\_\_ & (7) \_\_\_\_ \\
Accusative & (4) \_\_\_\_ & (8) \_\_\_\_ \\
\bottomrule
\end{tabular}
\par\endgroup

% mdtable:T76:end
% mdparagraph:P365:begin
\textbf{Response grid: tōð}
% mdparagraph:P365:end

% output: NEW-U15-F06-O009
% output: NEW-U15-F06-O010
% output: NEW-U15-F06-O011
% output: NEW-U15-F06-O012
% output: NEW-U15-F06-O013
% output: NEW-U15-F06-O014
% output: NEW-U15-F06-O015
% output: NEW-U15-F06-O016
% mdtable:T77:begin
\par\begingroup
\fontsize{10.5}{13.2}\selectfont
\setlength{\LTleft}{0pt}\setlength{\LTright}{\fill}
\noindent
\begin{tabular}{@{}>{\raggedright\arraybackslash}p{39.46mm}>{\raggedright\arraybackslash}p{39.46mm}>{\raggedright\arraybackslash}p{39.46mm}@{}}
\toprule
\textbf{Case} & \textbf{Singular} & \textbf{Plural} \\
\midrule
Nominative & (9) \_\_\_\_ & (13) \_\_\_\_ \\
Genitive & (10) \_\_\_\_ & (14) \_\_\_\_ \\
Dative & (11) \_\_\_\_ & (15) \_\_\_\_ \\
Accusative & (12) \_\_\_\_ & (16) \_\_\_\_ \\
\bottomrule
\end{tabular}
\par\endgroup

% mdtable:T77:end
% output: NEW-U15-F06-O017
% output: NEW-U15-F06-O018
% mdparagraph:P366:begin
\textbf{Then explain:} (17) positions with vowel changes; (18) forms whose cases need context.
% mdparagraph:P366:end

\end{minipage}\par\addvspace{4pt}

\par\begin{minipage}{\linewidth}\raggedright
% mdparagraph:P367:begin
\textbf{F07}
% mdparagraph:P367:end

% item: NEW-U15-F07
% mdparagraph:P368:begin
Give the complete nominative, genitive, dative and accusative singular/\allowbreak{}plural of sweostor. Compare the number of distinct surfaces with bōc; explain why the frequent surface sweostor still requires a case analysis in a sentence.
% mdparagraph:P368:end

% mdparagraph:P369:begin
\textbf{Response grid: sweostor}
% mdparagraph:P369:end

% output: NEW-U15-F07-O001
% output: NEW-U15-F07-O002
% output: NEW-U15-F07-O003
% output: NEW-U15-F07-O004
% output: NEW-U15-F07-O005
% output: NEW-U15-F07-O006
% output: NEW-U15-F07-O007
% output: NEW-U15-F07-O008
% mdtable:T78:begin
\par\begingroup
\fontsize{10.5}{13.2}\selectfont
\setlength{\LTleft}{0pt}\setlength{\LTright}{\fill}
\noindent
\begin{tabular}{@{}>{\raggedright\arraybackslash}p{39.46mm}>{\raggedright\arraybackslash}p{39.46mm}>{\raggedright\arraybackslash}p{39.46mm}@{}}
\toprule
\textbf{Case} & \textbf{Singular} & \textbf{Plural} \\
\midrule
Nominative & (1) \_\_\_\_ & (5) \_\_\_\_ \\
Genitive & (2) \_\_\_\_ & (6) \_\_\_\_ \\
Dative & (3) \_\_\_\_ & (7) \_\_\_\_ \\
Accusative & (4) \_\_\_\_ & (8) \_\_\_\_ \\
\bottomrule
\end{tabular}
\par\endgroup

% mdtable:T78:end
% output: NEW-U15-F07-O009
% output: NEW-U15-F07-O010
% mdparagraph:P370:begin
\textbf{Then explain:} (9) contrast with bōc; (10) reason a sentence still requires case analysis.
% mdparagraph:P370:end

\end{minipage}\par\addvspace{4pt}

\par\begin{minipage}{\linewidth}\raggedright
% mdparagraph:P371:begin
\textbf{F08}
% mdparagraph:P371:end

% item: NEW-U15-F08
% mdparagraph:P372:begin
Give nominative, genitive, dative and accusative singular of hǣtu and hǣlu. Use the selected forms in the tables. Explain why the plural positions are labelled unassigned and why this neither licenses a count meaning nor claims that no historical plural ever existed.
% mdparagraph:P372:end

% mdparagraph:P373:begin
\textbf{Response grid: hǣtu}
% mdparagraph:P373:end

% output: NEW-U15-F08-O001
% output: NEW-U15-F08-O002
% output: NEW-U15-F08-O003
% output: NEW-U15-F08-O004
% mdtable:T79:begin
\par\begingroup
\fontsize{10.5}{13.2}\selectfont
\setlength{\LTleft}{0pt}\setlength{\LTright}{\fill}
\noindent
\begin{tabular}{@{}>{\raggedright\arraybackslash}p{60.59mm}>{\raggedright\arraybackslash}p{60.59mm}@{}}
\toprule
\textbf{Case} & \textbf{Singular} \\
\midrule
Nominative & (1) \_\_\_\_ \\
Genitive & (2) \_\_\_\_ \\
Dative & (3) \_\_\_\_ \\
Accusative & (4) \_\_\_\_ \\
\bottomrule
\end{tabular}
\par\endgroup

% mdtable:T79:end
% mdparagraph:P374:begin
\textbf{Response grid: hǣlu}
% mdparagraph:P374:end

% output: NEW-U15-F08-O005
% output: NEW-U15-F08-O006
% output: NEW-U15-F08-O007
% output: NEW-U15-F08-O008
% mdtable:T80:begin
\par\begingroup
\fontsize{10.5}{13.2}\selectfont
\setlength{\LTleft}{0pt}\setlength{\LTright}{\fill}
\noindent
\begin{tabular}{@{}>{\raggedright\arraybackslash}p{60.59mm}>{\raggedright\arraybackslash}p{60.59mm}@{}}
\toprule
\textbf{Case} & \textbf{Singular} \\
\midrule
Nominative & (5) \_\_\_\_ \\
Genitive & (6) \_\_\_\_ \\
Dative & (7) \_\_\_\_ \\
Accusative & (8) \_\_\_\_ \\
\bottomrule
\end{tabular}
\par\endgroup

% mdtable:T80:end
% output: NEW-U15-F08-O009
% output: NEW-U15-F08-O010
% mdparagraph:P375:begin
\textbf{Then explain:} (9) selected singular patterns; (10) meaning and limits of the unassigned plural cells.
% mdparagraph:P375:end

\end{minipage}\par\addvspace{4pt}

\par\begin{minipage}{\linewidth}\raggedright
% mdparagraph:P376:begin
\textbf{F09}
% mdparagraph:P376:end

% item: NEW-U15-F09
% mdparagraph:P377:begin
Give the complete nominative, genitive, dative and accusative singular/\allowbreak{}plural of gerēfa and lēoðwyrhta. Explain why the word-formation of lēoðwyrhta does not change where its weak-noun ending belongs.
% mdparagraph:P377:end

% mdparagraph:P378:begin
\textbf{Response grid: gerēfa}
% mdparagraph:P378:end

% output: NEW-U15-F09-O001
% output: NEW-U15-F09-O002
% output: NEW-U15-F09-O003
% output: NEW-U15-F09-O004
% output: NEW-U15-F09-O005
% output: NEW-U15-F09-O006
% output: NEW-U15-F09-O007
% output: NEW-U15-F09-O008
% mdtable:T81:begin
\par\begingroup
\fontsize{10.5}{13.2}\selectfont
\setlength{\LTleft}{0pt}\setlength{\LTright}{\fill}
\noindent
\begin{tabular}{@{}>{\raggedright\arraybackslash}p{39.46mm}>{\raggedright\arraybackslash}p{39.46mm}>{\raggedright\arraybackslash}p{39.46mm}@{}}
\toprule
\textbf{Case} & \textbf{Singular} & \textbf{Plural} \\
\midrule
Nominative & (1) \_\_\_\_ & (5) \_\_\_\_ \\
Genitive & (2) \_\_\_\_ & (6) \_\_\_\_ \\
Dative & (3) \_\_\_\_ & (7) \_\_\_\_ \\
Accusative & (4) \_\_\_\_ & (8) \_\_\_\_ \\
\bottomrule
\end{tabular}
\par\endgroup

% mdtable:T81:end
% mdparagraph:P379:begin
\textbf{Response grid: lēoðwyrhta}
% mdparagraph:P379:end

% output: NEW-U15-F09-O009
% output: NEW-U15-F09-O010
% output: NEW-U15-F09-O011
% output: NEW-U15-F09-O012
% output: NEW-U15-F09-O013
% output: NEW-U15-F09-O014
% output: NEW-U15-F09-O015
% output: NEW-U15-F09-O016
% mdtable:T82:begin
\par\begingroup
\fontsize{10.5}{13.2}\selectfont
\setlength{\LTleft}{0pt}\setlength{\LTright}{\fill}
\noindent
\begin{tabular}{@{}>{\raggedright\arraybackslash}p{39.46mm}>{\raggedright\arraybackslash}p{39.46mm}>{\raggedright\arraybackslash}p{39.46mm}@{}}
\toprule
\textbf{Case} & \textbf{Singular} & \textbf{Plural} \\
\midrule
Nominative & (9) \_\_\_\_ & (13) \_\_\_\_ \\
Genitive & (10) \_\_\_\_ & (14) \_\_\_\_ \\
Dative & (11) \_\_\_\_ & (15) \_\_\_\_ \\
Accusative & (12) \_\_\_\_ & (16) \_\_\_\_ \\
\bottomrule
\end{tabular}
\par\endgroup

% mdtable:T82:end
% output: NEW-U15-F09-O017
% mdparagraph:P380:begin
\textbf{Then explain:} (17) relation of the compound’s structure to its inflection.
% mdparagraph:P380:end

\end{minipage}\par\addvspace{4pt}

\par\begin{minipage}{\linewidth}\raggedright
% mdparagraph:P381:begin
\textbf{A12}
% mdparagraph:P381:end

% item: NEW-U15-A12
% mdparagraph:P382:begin
Use the complete displayed hēafod and fōtsceamul tables. Identify the ending-bearing stem of hēafod, give its selected nominative/\allowbreak{}accusative plural, and contrast the stem retained throughout fōtsceamul. Explain why a general instruction to delete every medial vowel would give wrong forms.
% mdparagraph:P382:end

% mdparagraph:P383:begin
(1) hēafod
% mdparagraph:P383:end

% mdparagraph:P384:begin
(2) fōtsceamul
% mdparagraph:P384:end

% output: NEW-U15-A12-O001
% output: NEW-U15-A12-O002
% output: NEW-U15-A12-O003
% output: NEW-U15-A12-O004
% mdparagraph:P385:begin
\textbf{Responses, in order:} (1) hēafod: stem before endings; (2) hēafod: nominative/\allowbreak{}accusative plural; (3) fōtsceamul: stem before endings; (4) why deleting every medial vowel would fail.
% mdparagraph:P385:end

\end{minipage}\par\addvspace{4pt}

\par\begin{minipage}{\linewidth}\raggedright
% mdparagraph:P386:begin
\textbf{C13}
% mdparagraph:P386:end

% item: NEW-U15-C13
% mdparagraph:P387:begin
Translate: Do not go; give the book to the poor man. Address one person with a negative command followed by an affirmative command.
% mdparagraph:P387:end

% output: NEW-U15-C13-O001
\end{minipage}\par\addvspace{4pt}

\reserve{7\baselineskip}
% mdparagraph:P388:begin
\section*{7. Cook’s exercises}\BookHeading{nav-ch15-learner-0035}{subsection}{7. Cook’s exercises}{1}
% mdparagraph:P388:end

% mdparagraph:P389:begin
The prompts, parentheses, references and translation controls below are Cook’s original material. Translate every item into Old English. Modern help is separate and remains available; it supplies unusual forms or constructions without requiring their unaided recall. Numbers in Cook’s parentheses refer to his grammar or earlier exercises.\textsuperscript{7}
% mdparagraph:P389:end

\reserve{7\baselineskip}
% mdparagraph:P390:begin
\subsection*{Cook IX.28}\BookHeading{nav-ch15-learner-0036}{subsubsection}{Cook IX.28}{0}
% mdparagraph:P390:end

\reserve{5\baselineskip}
\par\begin{minipage}{\linewidth}\raggedright
% mdparagraph:P391:begin
\textbf{Cook IX.28.1}
% mdparagraph:P391:end

% item: EX-IX-28-01
% mdparagraph:P392:begin
Every one that (87. d) forsakes for my name his house, or his brethren, or sisters (52. 3), or father, or mother, or wife, or children, or lands, with a hundredfold (om. it shall) be requited to him (verb last), and he shall have besides eternal life.
% mdparagraph:P392:end

% mdparagraph:P393:begin
\textbf{Modern help — available throughout this item.}
% mdparagraph:P393:end

\end{minipage}\par\addvspace{4pt}

% teaching: U15-COOK-NOM-28-01
% support: U15-COOK-NOM-28-01
% mdparagraph:P394:begin
Keep every alternative, including sisters, and both promised results. Sweostor has its own table; retain the distinct earlier brōðor and mōdor patterns.
% mdparagraph:P394:end

% teaching: NEW-U11-LOCAL-HELPS
% teaching: NEW-U15-LOCAL-HELPS
% teaching: NEW-U15-RECORD-HUNDTEONTIG-U15-SOURCE-DAT-PL
% teaching: NEW-U15-RECORD-HUNDTEONTIG-U15-SOURCE-SENSE
% teaching: U15-H-28-01
% support: U15-H-28-01
% mdparagraph:P395:begin
Use him mid hundtēontigum forgolden bið, “he will be repaid a hundredfold”, with the finite verb last. The initial ǣlc has its own relative clause and is resumed by dative him, the recipient in this supplied passive construction. Mid hundtēontigum means “with a hundredfold”. For mīnum naman means “for the sake of my name”. Keep this whole construction available as help.
% mdparagraph:P395:end

% output: EX-IX-28-01-O001
\reserve{5\baselineskip}
\par\begin{minipage}{\linewidth}\raggedright
% mdparagraph:P396:begin
\textbf{Cook IX.28.2}
% mdparagraph:P396:end

% item: EX-IX-28-02
% mdparagraph:P397:begin
He gave (sęllan) to each a penny in return for his day’s work.
% mdparagraph:P397:end

% output: EX-IX-28-02-O001
\end{minipage}\par\addvspace{4pt}

\par\begin{minipage}{\linewidth}\raggedright
% mdparagraph:P398:begin
\textbf{Cook IX.28.3}
% mdparagraph:P398:end

% item: EX-IX-28-03
% mdparagraph:P399:begin
He sent them into (on, 172) his vineyard.
% mdparagraph:P399:end

% output: EX-IX-28-03-O001
\end{minipage}\par\addvspace{4pt}

\par\begin{minipage}{\linewidth}\raggedright
% mdparagraph:P400:begin
\textbf{Cook IX.28.4}
% mdparagraph:P400:end

% item: EX-IX-28-04
% mdparagraph:P401:begin
He saw others in (om. the) public place standing (169) idle.
% mdparagraph:P401:end

% mdparagraph:P402:begin
\textbf{Modern help — available throughout this item.}
% mdparagraph:P402:end

\end{minipage}\par\addvspace{4pt}

% teaching: U15-COOK-NOM-28-04
% support: U15-COOK-NOM-28-04
% mdparagraph:P403:begin
Strǣt here means “public place”, with the earlier feminine declension. Cook’s reference (169) selects the earlier perception construction: use accusative ōðre with the plain infinitive standan. In 28.7, by contrast, the printed direction explicitly requests a present participle.
% mdparagraph:P403:end

% output: EX-IX-28-04-O001
\reserve{5\baselineskip}
\par\begin{minipage}{\linewidth}\raggedright
% mdparagraph:P404:begin
\textbf{Cook IX.28.5}
% mdparagraph:P404:end

% item: EX-IX-28-05
% mdparagraph:P405:begin
Then said he, Go ye also into my vineyard, and I will give you what (ðæt) is right.
% mdparagraph:P405:end

% output: EX-IX-28-05-O001
\end{minipage}\par\addvspace{4pt}

\par\begin{minipage}{\linewidth}\raggedright
% mdparagraph:P406:begin
\textbf{Cook IX.28.6}
% mdparagraph:P406:end

% item: EX-IX-28-06
% mdparagraph:P407:begin
Once more he went out about (ymbe, 172) the sixth and ninth hour, and did similarly.
% mdparagraph:P407:end

% mdparagraph:P408:begin
\textbf{Modern help — available throughout this item.}
% mdparagraph:P408:end

\end{minipage}\par\addvspace{4pt}

% teaching: U15-COOK-NOM-28-06
% support: U15-COOK-NOM-28-06
% mdparagraph:P409:begin
Tīd means “hour”. Temporal ymbe takes accusative for the approximate time.
% mdparagraph:P409:end

% teaching: NEW-U15-LOCAL-HELPS
% teaching: U15-H-28-06
% support: U15-H-28-06
% mdparagraph:P410:begin
Use supplied nigoðan, the feminine accusative singular weak form of “ninth”, with tīd. Siextan, “sixth”, is already familiar.
% mdparagraph:P410:end

% output: EX-IX-28-06-O001
\reserve{5\baselineskip}
\par\begin{minipage}{\linewidth}\raggedright
% mdparagraph:P411:begin
\textbf{Cook IX.28.7}
% mdparagraph:P411:end

% item: EX-IX-28-07
% mdparagraph:P412:begin
Then about the eleventh hour he went out, and found others standing (pres. part.).
% mdparagraph:P412:end

% mdparagraph:P413:begin
\textbf{Modern help — available throughout this item.}
% mdparagraph:P413:end

\end{minipage}\par\addvspace{4pt}

% teaching: NEW-U15-LOCAL-HELPS
% teaching: U15-H-28-07
% support: U15-H-28-07
% mdparagraph:P414:begin
Use the supplied expression ðā endleftan tīd, “the eleventh hour”, with the feminine accusative singular weak ordinal.
% mdparagraph:P414:end

% output: EX-IX-28-07-O001
\reserve{5\baselineskip}
\par\begin{minipage}{\linewidth}\raggedright
% mdparagraph:P415:begin
\textbf{Cook IX.28.8}
% mdparagraph:P415:end

% item: EX-IX-28-08
% mdparagraph:P416:begin
Then said he, Why stand ye here all day idle?
% mdparagraph:P416:end

% output: EX-IX-28-08-O001
\end{minipage}\par\addvspace{4pt}

\par\begin{minipage}{\linewidth}\raggedright
% mdparagraph:P417:begin
\textbf{Cook IX.28.9}
% mdparagraph:P417:end

% item: EX-IX-28-09
% mdparagraph:P418:begin
Then said they, Because no man (om. hath) hired (hȳrian, 118) us.
% mdparagraph:P418:end

% output: EX-IX-28-09-O001
\end{minipage}\par\addvspace{4pt}

\par\begin{minipage}{\linewidth}\raggedright
% mdparagraph:P419:begin
\textbf{Cook IX.28.10}
% mdparagraph:P419:end

% item: EX-IX-28-10
% mdparagraph:P420:begin
Truly, when even was come [when it was evening become], then said the lord (hlāford) of the vineyard to his steward (gerēfa, 53), Call (cleopian) the workmen, and pay them (164. a) their hire (mēd).
% mdparagraph:P420:end

% mdparagraph:P421:begin
\textbf{Modern help — available throughout this item.}
% mdparagraph:P421:end



% teaching: U15-COOK-NOM-28-10
% support: U15-COOK-NOM-28-10
% mdparagraph:P422:begin
Ǣfen means “evening”; wyrhta, “workman”; gerēfa, “steward”; wīngeard, “vineyard”; mēd here, “hire; wages”.
% mdparagraph:P422:end
\end{minipage}\par\addvspace{4pt}

% teaching: NEW-U14-LOCAL-HELPS
% teaching: NEW-U15-LOCAL-HELPS
% teaching: U15-H-28-10
% support: U15-H-28-10
% mdparagraph:P423:begin
Use ðā hit ǣfen geworden wæs, “when it had become evening”. Geworden is the supplied past participle of geweorðan. With cweðan, tō his gerēfan names the person addressed: tō takes dative.
% mdparagraph:P423:end

% output: EX-IX-28-10-O001
\reserve{5\baselineskip}
\par\begin{minipage}{\linewidth}\raggedright
% mdparagraph:P424:begin
\textbf{Cook IX.28.11}
% mdparagraph:P424:end

% item: EX-IX-28-11
% mdparagraph:P425:begin
Accordingly, when those came that came about the eleventh hour (verb last), then received they each his penny.
% mdparagraph:P425:end

% mdparagraph:P426:begin
\textbf{Modern help — available throughout this item.}
% mdparagraph:P426:end

\end{minipage}\par\addvspace{4pt}

% teaching: NEW-U15-LOCAL-HELPS
% teaching: U15-H-28-07
% support: U15-H-28-07
% mdparagraph:P427:begin
Use the supplied expression ðā endleftan tīd, “the eleventh hour”, with the feminine accusative singular weak ordinal.
% mdparagraph:P427:end

% output: EX-IX-28-11-O001
\reserve{5\baselineskip}
\par\begin{minipage}{\linewidth}\raggedright
% mdparagraph:P428:begin
\textbf{Cook IX.28.12}
% mdparagraph:P428:end

% item: EX-IX-28-12
% mdparagraph:P429:begin
And those, that there (ðǣr) first came, supposed that they should (sculan) receive more (māre).
% mdparagraph:P429:end

% output: EX-IX-28-12-O001
\end{minipage}\par\addvspace{4pt}

\reserve{7\baselineskip}
% mdparagraph:P430:begin
\subsection*{Cook IX.29}\BookHeading{nav-ch15-learner-0037}{subsubsection}{Cook IX.29}{0}
% mdparagraph:P430:end

\reserve{5\baselineskip}
\par\begin{minipage}{\linewidth}\raggedright
% mdparagraph:P431:begin
\textbf{Cook IX.29.1}
% mdparagraph:P431:end

% item: EX-IX-29-01
% mdparagraph:P432:begin
Thou (om. hast) made (dōn) them equal (gelīc) (om. unto) us, which (87. d) (om. have) borne burdens (51. 5) in this day’s heat.
% mdparagraph:P432:end

% mdparagraph:P433:begin
\textbf{Modern help — available throughout this item.}
% mdparagraph:P433:end

\end{minipage}\par\addvspace{4pt}

% teaching: U15-COOK-NOM-29-01
% support: U15-COOK-NOM-29-01
% mdparagraph:P434:begin
Hęfignessa is accusative plural “burdens”; hǣte is the selected dative singular “heat”.
% mdparagraph:P434:end

% teaching: NEW-U15-LOCAL-HELPS
% teaching: U15-H-29-01
% support: U15-H-29-01
% mdparagraph:P435:begin
On ðisses dæges hǣte means “in this day’s heat”: on with dative gives the circumstance; ðisses dæges is a dependent genitive.
% mdparagraph:P435:end

% output: EX-IX-29-01-O001
\reserve{5\baselineskip}
\par\begin{minipage}{\linewidth}\raggedright
% mdparagraph:P436:begin
\textbf{Cook IX.29.2}
% mdparagraph:P436:end

% item: EX-IX-29-02
% mdparagraph:P437:begin
Take that is thine, and go (om. thy way); I will give unto this (om. one) just (eall) as much as (om. unto) thee.
% mdparagraph:P437:end

% output: EX-IX-29-02-O001
\end{minipage}\par\addvspace{4pt}

\par\begin{minipage}{\linewidth}\raggedright
% mdparagraph:P438:begin
\textbf{Cook IX.29.3}
% mdparagraph:P438:end

% item: EX-IX-29-03
% mdparagraph:P439:begin
May I not do what (ðæt) I will?
% mdparagraph:P439:end

% output: EX-IX-29-03-O001
\end{minipage}\par\addvspace{4pt}

\par\begin{minipage}{\linewidth}\raggedright
% mdparagraph:P440:begin
\textbf{Cook IX.29.4}
% mdparagraph:P440:end

% item: EX-IX-29-04
% mdparagraph:P441:begin
Camest thou not to me to work (wyrcean, 200) for (wið) one penny?
% mdparagraph:P441:end

% output: EX-IX-29-04-O001
\end{minipage}\par\addvspace{4pt}

\par\begin{minipage}{\linewidth}\raggedright
% mdparagraph:P442:begin
\textbf{Cook IX.29.5}
% mdparagraph:P442:end

% item: EX-IX-29-05
% mdparagraph:P443:begin
He was (om. a) man (wer) and prophet, mighty in speech and in deed (weorc) before (beforan, 166) God and all (om. the) people (folc).
% mdparagraph:P443:end

% output: EX-IX-29-05-O001
\end{minipage}\par\addvspace{4pt}

\par\begin{minipage}{\linewidth}\raggedright
% mdparagraph:P444:begin
\textbf{Cook IX.29.6}
% mdparagraph:P444:end

% item: EX-IX-29-06
% mdparagraph:P445:begin
That is the judgment, that light (om. is) come into (on, 172) (om. the) world, and men loved darkness (plur.) rather (swīðor) than the light, because their deeds (weorc) were evil.
% mdparagraph:P445:end

% output: EX-IX-29-06-O001
\end{minipage}\par\addvspace{4pt}

\par\begin{minipage}{\linewidth}\raggedright
% mdparagraph:P446:begin
\textbf{Cook IX.29.7}
% mdparagraph:P446:end

% item: EX-IX-29-07
% mdparagraph:P447:begin
God sent not his son into (om. the) world that he (om. should) condemn the world, but that the world (om. may) be saved (hǣlan, 113) by (ðurh) him.
% mdparagraph:P447:end

% mdparagraph:P448:begin
\textbf{Modern help — available throughout this item.}
% mdparagraph:P448:end

\end{minipage}\par\addvspace{4pt}

% teaching: NEW-U15-LOCAL-HELPS
% teaching: U15-H-29-07
% support: U15-H-29-07
% mdparagraph:P449:begin
Use the supplied spelling ðurh beside the familiar þurh. Here ðurh hine means “through/\allowbreak{}by way of him”. Negation rejects the condemning purpose and contrasts it with saving; it does not deny that the Son was sent.
% mdparagraph:P449:end

% output: EX-IX-29-07-O001
\reserve{5\baselineskip}
\par\begin{minipage}{\linewidth}\raggedright
% mdparagraph:P450:begin
\textbf{Cook IX.29.8}
% mdparagraph:P450:end

% item: EX-IX-29-08
% mdparagraph:P451:begin
Your reward (om. shall) be great, and ye shall be (om. the) children of the Highest; because he is kind (gōd) unto (ofer, 172) (om. the) unthankful and evil.
% mdparagraph:P451:end

% mdparagraph:P452:begin
\textbf{Modern help — available throughout this item.}
% mdparagraph:P452:end

\end{minipage}\par\addvspace{4pt}

% teaching: NEW-U15-LOCAL-HELPS
% teaching: U15-H-29-08
% support: U15-H-29-08
% mdparagraph:P453:begin
Use ðæs Hēahstan, “of the Highest”, as a supplied masculine genitive singular divine title. No further superlative forms of hēah are required.
% mdparagraph:P453:end

% output: EX-IX-29-08-O001
\reserve{5\baselineskip}
\par\begin{minipage}{\linewidth}\raggedright
% mdparagraph:P454:begin
\textbf{Cook IX.29.9}
% mdparagraph:P454:end

% item: EX-IX-29-09
% mdparagraph:P455:begin
Give, and (om. it) shall be given you.
% mdparagraph:P455:end

% mdparagraph:P456:begin
\textbf{Modern help — available throughout this item.}
% mdparagraph:P456:end

\end{minipage}\par\addvspace{4pt}

% teaching: NEW-U15-LOCAL-HELPS
% teaching: U15-H-29-09
% support: U15-H-29-09
% mdparagraph:P457:begin
Use the supplied frame ēow bið giefen, “it will be given to you (plural).” Ēow is the dative recipient; no hit is expressed. For one addressee, use ðē bið giefen with a singular command. Keep the frame available while composing; its independent recall is not required.
% mdparagraph:P457:end

% output: EX-IX-29-09-O001
\reserve{5\baselineskip}
\par\begin{minipage}{\linewidth}\raggedright
% mdparagraph:P458:begin
\textbf{Cook IX.29.10}
% mdparagraph:P458:end

% item: EX-IX-29-10
% mdparagraph:P459:begin
The disciple (leornungcniht, 43) is not above (ofer, 172) his master (lārēow).
% mdparagraph:P459:end

% output: EX-IX-29-10-O001
\end{minipage}\par\addvspace{4pt}

\par\begin{minipage}{\linewidth}\raggedright
% mdparagraph:P460:begin
\textbf{Cook IX.29.11}
% mdparagraph:P460:end

% item: EX-IX-29-11
% mdparagraph:P461:begin
Show us thy mercy, and grant us thy salvation.
% mdparagraph:P461:end

% mdparagraph:P462:begin
\textbf{Modern help — available throughout this item.}
% mdparagraph:P462:end

\end{minipage}\par\addvspace{4pt}

% teaching: U15-COOK-NOM-29-11
% support: U15-COOK-NOM-29-11
% mdparagraph:P463:begin
Hǣlu, “salvation”, is unchanged through the selected singular. Unnan governs genitive of the thing granted, so use genitive ðīnre with it even though hǣlu shows no separate genitive ending. Hǣle is a recognized oblique alternative.
% mdparagraph:P463:end

% output: EX-IX-29-11-O001
\reserve{7\baselineskip}
% mdparagraph:P464:begin
\subsection*{Cook IX.30}\BookHeading{nav-ch15-learner-0038}{subsubsection}{Cook IX.30}{0}
% mdparagraph:P464:end

\reserve{5\baselineskip}
\par\begin{minipage}{\linewidth}\raggedright
% mdparagraph:P465:begin
\textbf{Cook IX.30.1}
% mdparagraph:P465:end

% item: EX-IX-30-01
% mdparagraph:P466:begin
Swear (thou) not by (ðurh) heaven, because it is God’s throne; nor by earth, because it is his footstool (fōtsceamul, 43); nor by Jerusalem (Hierusalem), because it is (om. the) city of (om. the) great (mǣre) king.
% mdparagraph:P466:end

% mdparagraph:P467:begin
\textbf{Modern help — available throughout this item.}
% mdparagraph:P467:end

\end{minipage}\par\addvspace{4pt}

% teaching: U15-COOK-NOM-30-01
% support: U15-COOK-NOM-30-01
% mdparagraph:P468:begin
Setl means “seat; throne”. Fōtsceamul retains -sceamul- in its inflected stem.
% mdparagraph:P468:end

% teaching: NEW-U15-LOCAL-HELPS
% teaching: NEW-U15-VOC-THURH
% teaching: U15-H-30-01
% support: U15-H-30-01
% mdparagraph:P469:begin
Use the supplied spelling ðurh before each accusative oath standard. Hierusalem is the supplied place name; its other cases are not required.
% mdparagraph:P469:end

% output: EX-IX-30-01-O001
\reserve{5\baselineskip}
\par\begin{minipage}{\linewidth}\raggedright
% mdparagraph:P470:begin
\textbf{Cook IX.30.2}
% mdparagraph:P470:end

% item: EX-IX-30-02
% mdparagraph:P471:begin
Swear not by thy head, because thou canst not make (wyrcean) one hair (hǣr, 47) white or black.
% mdparagraph:P471:end

% mdparagraph:P472:begin
\textbf{Modern help — available throughout this item.}
% mdparagraph:P472:end

\end{minipage}\par\addvspace{4pt}

% teaching: U15-COOK-NOM-30-02
% support: U15-COOK-NOM-30-02
% mdparagraph:P473:begin
Hēafod uses hēafd- before endings. Hǣr is neuter and has no nominative/\allowbreak{}accusative plural ending. Fōtsceamul, by contrast, retains its medial vowel.
% mdparagraph:P473:end

% teaching: NEW-U15-LOCAL-HELPS
% teaching: NEW-U15-VOC-THURH
% teaching: U15-H-30-01
% support: U15-H-30-01
% mdparagraph:P474:begin
Use the supplied spelling ðurh before each accusative oath standard. Hierusalem is the supplied place name; its other cases are not required.
% mdparagraph:P474:end

% output: EX-IX-30-02-O001
\reserve{5\baselineskip}
\par\begin{minipage}{\linewidth}\raggedright
% mdparagraph:P475:begin
\textbf{Cook IX.30.3}
% mdparagraph:P475:end

% item: EX-IX-30-03
% mdparagraph:P476:begin
Be your (plur.) speech thus, It is, It is not.
% mdparagraph:P476:end

% output: EX-IX-30-03-O001
\end{minipage}\par\addvspace{4pt}

\par\begin{minipage}{\linewidth}\raggedright
% mdparagraph:P477:begin
\textbf{Cook IX.30.4}
% mdparagraph:P477:end

% item: EX-IX-30-04
% mdparagraph:P478:begin
Ye (om. have) heard that (om. it) was said, Eye for (for, 166) eye, and tooth (46) for tooth.
% mdparagraph:P478:end

% output: EX-IX-30-04-O001
\end{minipage}\par\addvspace{4pt}

\par\begin{minipage}{\linewidth}\raggedright
% mdparagraph:P479:begin
\textbf{Cook IX.30.5}
% mdparagraph:P479:end

% item: EX-IX-30-05
% mdparagraph:P480:begin
He that breaks (tōweorpan, 104) one of (of, 166) these least commandments (bebod, 47), and teacheth (lǣran) men so (verb last), he (sē) shall be named least in (om. the) kingdom (rīce) of heaven.
% mdparagraph:P480:end

% mdparagraph:P481:begin
\textbf{Modern help — available throughout this item.}
% mdparagraph:P481:end

\end{minipage}\par\addvspace{4pt}

% teaching: NEW-U15-LOCAL-HELPS
% teaching: NEW-U15-RECORD-OF-U15-SOURCE-PARTITIVE
% teaching: U15-H-30-05
% support: U15-H-30-05
% mdparagraph:P482:begin
The emphatic resumed sē is a supplied long-vowel nominative masculine singular. In ān of ðissum lǣstum bebodum, “one of these least commandments”, of takes dative to identify the group. The same long sē may be used in 30.6.
% mdparagraph:P482:end

% output: EX-IX-30-05-O001
\reserve{5\baselineskip}
\par\begin{minipage}{\linewidth}\raggedright
% mdparagraph:P483:begin
\textbf{Cook IX.30.6}
% mdparagraph:P483:end

% item: EX-IX-30-06
% mdparagraph:P484:begin
He that doeth and teacheth it, he shall be called great (micel).
% mdparagraph:P484:end

% mdparagraph:P485:begin
\textbf{Modern help — available throughout this item.}
% mdparagraph:P485:end

\end{minipage}\par\addvspace{4pt}

% teaching: NEW-U15-LOCAL-HELPS
% teaching: NEW-U15-RECORD-OF-U15-SOURCE-PARTITIVE
% teaching: U15-H-30-05
% support: U15-H-30-05
% mdparagraph:P486:begin
The emphatic resumed sē is a supplied long-vowel nominative masculine singular. In ān of ðissum lǣstum bebodum, “one of these least commandments”, of takes dative to identify the group. The same long sē may be used in 30.6.
% mdparagraph:P486:end

% output: EX-IX-30-06-O001
\reserve{5\baselineskip}
\par\begin{minipage}{\linewidth}\raggedright
% mdparagraph:P487:begin
\textbf{Cook IX.30.7}
% mdparagraph:P487:end

% item: EX-IX-30-07
% mdparagraph:P488:begin
Ye (om. have) heard that (om. it) was said in old times (tīd), Kill (thou) not.
% mdparagraph:P488:end

% mdparagraph:P489:begin
\textbf{Modern help — available throughout this item.}
% mdparagraph:P489:end

\end{minipage}\par\addvspace{4pt}

% teaching: NEW-U15-LOCAL-HELPS
% teaching: NEW-U15-RECORD-ON-U15-SOURCE-DURING
% teaching: U15-H-30-07
% support: U15-H-30-07
% mdparagraph:P490:begin
Use the supplied on ealdum tīdum, “in old times”, with temporal on and dative plural.
% mdparagraph:P490:end

% output: EX-IX-30-07-O001
\reserve{5\baselineskip}
\par\begin{minipage}{\linewidth}\raggedright
% mdparagraph:P491:begin
\textbf{Cook IX.30.8}
% mdparagraph:P491:end

% item: EX-IX-30-08
% mdparagraph:P492:begin
Is it your (sing.) earth?
% mdparagraph:P492:end

% mdparagraph:P493:begin
\textbf{Modern help — available throughout this item.}
% mdparagraph:P493:end

\end{minipage}\par\addvspace{4pt}

% teaching: U15-COOK-NOM-30-08
% support: U15-COOK-NOM-30-08
% mdparagraph:P494:begin
Eorðe, “earth”, is feminine. Its possessive must agree with the noun that it modifies.
% mdparagraph:P494:end

% output: EX-IX-30-08-O001
\reserve{5\baselineskip}
\par\begin{minipage}{\linewidth}\raggedright
% mdparagraph:P495:begin
\textbf{Cook IX.30.9}
% mdparagraph:P495:end

% item: EX-IX-30-09
% mdparagraph:P496:begin
He has from his labor what he can.
% mdparagraph:P496:end

% mdparagraph:P497:begin
\textbf{Modern help — available throughout this item.}
% mdparagraph:P497:end

\end{minipage}\par\addvspace{4pt}

% teaching: U15-COOK-NOM-30-09
% support: U15-COOK-NOM-30-09
% mdparagraph:P498:begin
Gewinne, “labour”, is a supplied noun form in the expression æt his gewinne.
% mdparagraph:P498:end

% teaching: NEW-U15-LOCAL-HELPS
% teaching: U15-H-30-09
% support: U15-H-30-09
% mdparagraph:P499:begin
Use æt his gewinne, “from his labour”. In the supported ðæt hē mæg, “what he can”, an action such as “obtain” is understood. Retain this supplied expression without deriving further forms of gewinn.
% mdparagraph:P499:end

% output: EX-IX-30-09-O001
\reserve{5\baselineskip}
\par\begin{minipage}{\linewidth}\raggedright
% mdparagraph:P500:begin
\textbf{Cook IX.30.10}
% mdparagraph:P500:end

% item: EX-IX-30-10
% mdparagraph:P501:begin
The sun parched his field.
% mdparagraph:P501:end

% output: EX-IX-30-10-O001
\end{minipage}\par\addvspace{4pt}

\par\begin{minipage}{\linewidth}\raggedright
% mdparagraph:P502:begin
\textbf{Cook IX.30.11}
% mdparagraph:P502:end

% item: EX-IX-30-11
% mdparagraph:P503:begin
The roads (weg) should blossom (blōstmian, 118), the roads should bloom, so fair (fæger) (om. a) bride (brȳd, 51. b) shall leave her home (hām).
% mdparagraph:P503:end

% mdparagraph:P504:begin
\textbf{Modern help — available throughout this item.}
% mdparagraph:P504:end

\end{minipage}\par\addvspace{4pt}

% teaching: U15-COOK-NOM-30-11
% support: U15-COOK-NOM-30-11
% mdparagraph:P505:begin
Brȳd is feminine; the selected accusative singular is brȳde. Hām is the noun “home”, here an accusative object, rather than the directional adverb. Use the given infinitive blōstmian for “blossom”.
% mdparagraph:P505:end

% output: EX-IX-30-11-O001
\reserve{5\baselineskip}
\par\begin{minipage}{\linewidth}\raggedright
% mdparagraph:P506:begin
\textbf{Cook IX.30.12}
% mdparagraph:P506:end

% item: EX-IX-30-12
% mdparagraph:P507:begin
If they can, let them live without me.
% mdparagraph:P507:end

% mdparagraph:P508:begin
\textbf{Modern help — available throughout this item.}
% mdparagraph:P508:end

\end{minipage}\par\addvspace{4pt}

% teaching: NEW-U15-LOCAL-HELPS
% teaching: U15-H-30-12
% support: U15-H-30-12
% mdparagraph:P509:begin
Use the supplied permission construction lǣt hīe libban, “let them live”: accusative hīe is the understood subject of libban. Gif hīe magon expresses ability. A third-person optative is also a valid way to express the request.
% mdparagraph:P509:end

% output: EX-IX-30-12-O001
\reserve{5\baselineskip}
\par\begin{minipage}{\linewidth}\raggedright
% mdparagraph:P510:begin
\textbf{Cook IX.30.13}
% mdparagraph:P510:end

% item: EX-IX-30-13
% mdparagraph:P511:begin
He never forsakes those that trust in (on) him (172).
% mdparagraph:P511:end

% output: EX-IX-30-13-O001
\end{minipage}\par\addvspace{4pt}

\reserve{7\baselineskip}
% mdparagraph:P512:begin
\subsection*{Cook X, XI.31}\BookHeading{nav-ch15-learner-0039}{subsubsection}{Cook X, XI.31}{0}
% mdparagraph:P512:end

\reserve{5\baselineskip}
\par\begin{minipage}{\linewidth}\raggedright
% mdparagraph:P513:begin
\textbf{Cook X, XI.31.1}
% mdparagraph:P513:end

% item: EX-X-XI-31-01
% mdparagraph:P514:begin
I know truly that my Redeemer liveth.
% mdparagraph:P514:end

% mdparagraph:P515:begin
\textbf{Modern help — available throughout this item.}
% mdparagraph:P515:end

\end{minipage}\par\addvspace{4pt}

% teaching: U15-COOK-NOM-31-01
% support: U15-COOK-NOM-31-01
% mdparagraph:P516:begin
Use ālīesend as the singular title “Redeemer”. Its displayed plurals remain reference forms.
% mdparagraph:P516:end

% output: EX-X-XI-31-01-O001
\reserve{5\baselineskip}
\par\begin{minipage}{\linewidth}\raggedright
% mdparagraph:P517:begin
\textbf{Cook X, XI.31.2}
% mdparagraph:P517:end

% item: EX-X-XI-31-02
% mdparagraph:P518:begin
In my flesh I shall see God, I (om. my)self, and not (nā) (om. an)other.
% mdparagraph:P518:end

% output: EX-X-XI-31-02-O001
\end{minipage}\par\addvspace{4pt}

\par\begin{minipage}{\linewidth}\raggedright
% mdparagraph:P519:begin
\textbf{Cook X, XI.31.3}
% mdparagraph:P519:end

% item: EX-X-XI-31-03
% mdparagraph:P520:begin
He was taught (187) to do God’s will.
% mdparagraph:P520:end

% mdparagraph:P521:begin
\textbf{Modern help — available throughout this item.}
% mdparagraph:P521:end

\end{minipage}\par\addvspace{4pt}

% teaching: NEW-U15-LOCAL-HELPS
% teaching: U15-H-31-03
% support: U15-H-31-03
% mdparagraph:P522:begin
Use hē wæs lǣred … tō dōnne. The person taught is the passive subject and the understood doer of the infinitive. Godes willan is the infinitive’s accusative object, with dependent genitive Godes. This whole passive-and-infinitive frame is supplied.
% mdparagraph:P522:end

% output: EX-X-XI-31-03-O001
\reserve{5\baselineskip}
\par\begin{minipage}{\linewidth}\raggedright
% mdparagraph:P523:begin
\textbf{Cook X, XI.31.4}
% mdparagraph:P523:end

% item: EX-X-XI-31-04
% mdparagraph:P524:begin
Lead (thou) us not into temptation, but deliver us from evil.
% mdparagraph:P524:end

% mdparagraph:P525:begin
\textbf{Modern help — available throughout this item.}
% mdparagraph:P525:end

\end{minipage}\par\addvspace{4pt}

% teaching: NEW-U15-LOCAL-HELPS
% teaching: U15-H-31-04
% support: U15-H-31-04
% mdparagraph:P526:begin
Use fram yfele, “from evil”. Yfele is a supplied substantive noun in the dative singular, not the ordinary adjective dative yfelum.
% mdparagraph:P526:end

% output: EX-X-XI-31-04-O001
\reserve{5\baselineskip}
\par\begin{minipage}{\linewidth}\raggedright
% mdparagraph:P527:begin
\textbf{Cook X, XI.31.5}
% mdparagraph:P527:end

% item: EX-X-XI-31-05
% mdparagraph:P528:begin
Be thy name hallowed.
% mdparagraph:P528:end

% output: EX-X-XI-31-05-O001
\end{minipage}\par\addvspace{4pt}

\par\begin{minipage}{\linewidth}\raggedright
% mdparagraph:P529:begin
\textbf{Cook X, XI.31.6}
% mdparagraph:P529:end

% item: EX-X-XI-31-06
% mdparagraph:P530:begin
Thy will befall (geweorðan, verb first) in earth as (swā-swā) in heaven(s).
% mdparagraph:P530:end

% mdparagraph:P531:begin
\textbf{Modern help — available throughout this item.}
% mdparagraph:P531:end

\end{minipage}\par\addvspace{4pt}

% teaching: NEW-U09-SOURCE-HELPS
% teaching: NEW-U09-VOC-swa-swa
% teaching: NEW-U10-SOURCE-HELPS
% teaching: NEW-U12-REUSED-U12-SRC-HELP-15
% teaching: NEW-U12-SOURCE-HELPS
% teaching: NEW-U13-PAR-SWA-SWA
% teaching: NEW-U13-VOC-SWA-SWA
% teaching: NEW-U14-LOCAL-HELPS
% teaching: NEW-U15-LOCAL-HELPS
% teaching: U15-H-31-06
% support: U15-H-31-06
% mdparagraph:P532:begin
Begin with supplied Geweorðe, the third-singular present subjunctive of the prefixed verb. Swā-swā means “just as”; retain this supplied hyphenated expression.
% mdparagraph:P532:end

% output: EX-X-XI-31-06-O001
\reserve{5\baselineskip}
\par\begin{minipage}{\linewidth}\raggedright
% mdparagraph:P533:begin
\textbf{Cook X, XI.31.7}
% mdparagraph:P533:end

% item: EX-X-XI-31-07
% mdparagraph:P534:begin
No king of the English was so great as Alfred.
% mdparagraph:P534:end

% mdparagraph:P535:begin
\textbf{Modern help — available throughout this item.}
% mdparagraph:P535:end

\end{minipage}\par\addvspace{4pt}

% teaching: NEW-U15-LOCAL-HELPS
% teaching: NEW-U15-RECORD-AELFRED-U15-SOURCE-NOM-SG
% teaching: U15-H-31-NAMES
% support: U15-H-31-NAMES
% mdparagraph:P536:begin
The names Ælfred, Osburg, Æthelwulf and Asser are supplied; Ælfredes is the genitive for 31.14. Use supplied dative Lǣdene for “Latin”, following Cook’s glossary Lǣden. The Preface reading below retains its printed Lædene without a length mark.
% mdparagraph:P536:end

% output: EX-X-XI-31-07-O001
\reserve{5\baselineskip}
\par\begin{minipage}{\linewidth}\raggedright
% mdparagraph:P537:begin
\textbf{Cook X, XI.31.8}
% mdparagraph:P537:end

% item: EX-X-XI-31-08
% mdparagraph:P538:begin
He fought with the Danes (44. 2); he judged his people; he made translations from Latin; he was poet (lēoðwyrhta).
% mdparagraph:P538:end

% mdparagraph:P539:begin
\textbf{Modern help — available throughout this item.}
% mdparagraph:P539:end

\end{minipage}\par\addvspace{4pt}

% teaching: U15-COOK-NOM-31-08
% support: U15-COOK-NOM-31-08
% mdparagraph:P540:begin
Dęne and Engle name peoples in the plural. Lēoðwyrhta is a weak masculine noun, “poet”. Āwęndednes means “translation” and doubles s before endings. The language-name dative is supplied separately.
% mdparagraph:P540:end

% teaching: NEW-U15-LOCAL-HELPS
% teaching: NEW-U15-RECORD-AELFRED-U15-SOURCE-NOM-SG
% teaching: U15-H-31-NAMES
% support: U15-H-31-NAMES
% mdparagraph:P541:begin
The names Ælfred, Osburg, Æthelwulf and Asser are supplied; Ælfredes is the genitive for 31.14. Use supplied dative Lǣdene for “Latin”, following Cook’s glossary Lǣden. The Preface reading below retains its printed Lædene without a length mark.
% mdparagraph:P541:end

% output: EX-X-XI-31-08-O001
\reserve{5\baselineskip}
\par\begin{minipage}{\linewidth}\raggedright
% mdparagraph:P542:begin
\textbf{Cook X, XI.31.9}
% mdparagraph:P542:end

% item: EX-X-XI-31-09
% mdparagraph:P543:begin
His mother hight Osburg, and his father Æthelwulf.
% mdparagraph:P543:end

% mdparagraph:P544:begin
\textbf{Modern help — available throughout this item.}
% mdparagraph:P544:end

\end{minipage}\par\addvspace{4pt}

% teaching: NEW-U15-LOCAL-HELPS
% teaching: NEW-U15-RECORD-AELFRED-U15-SOURCE-NOM-SG
% teaching: U15-H-31-NAMES
% support: U15-H-31-NAMES
% mdparagraph:P545:begin
The names Ælfred, Osburg, Æthelwulf and Asser are supplied; Ælfredes is the genitive for 31.14. Use supplied dative Lǣdene for “Latin”, following Cook’s glossary Lǣden. The Preface reading below retains its printed Lædene without a length mark.
% mdparagraph:P545:end

% teaching: NEW-U15-LOCAL-HELPS
% teaching: U15-H-31-09
% support: U15-H-31-09
% mdparagraph:P546:begin
Supplied hātte means “is/\allowbreak{}was called”; context supplies past time here. It is a special passive naming form, not the active preterite hēt.
% mdparagraph:P546:end

% output: EX-X-XI-31-09-O001
\reserve{5\baselineskip}
\par\begin{minipage}{\linewidth}\raggedright
% mdparagraph:P547:begin
\textbf{Cook X, XI.31.10}
% mdparagraph:P547:end

% item: EX-X-XI-31-10
% mdparagraph:P548:begin
He lived fifty-two years.
% mdparagraph:P548:end

% mdparagraph:P549:begin
\textbf{Modern help — available throughout this item.}
% mdparagraph:P549:end

\end{minipage}\par\addvspace{4pt}

% teaching: NEW-U15-LOCAL-HELPS
% teaching: U15-H-31-10
% support: U15-H-31-10
% mdparagraph:P550:begin
Use the supplied number phrase twā and fīftig gēara, “fifty-two years”. Fīftig means “fifty”; gēara is genitive plural governed by this numeral expression. No further higher numerals are required.
% mdparagraph:P550:end

% output: EX-X-XI-31-10-O001
\reserve{5\baselineskip}
\par\begin{minipage}{\linewidth}\raggedright
% mdparagraph:P551:begin
\textbf{Cook X, XI.31.11}
% mdparagraph:P551:end

% item: EX-X-XI-31-11
% mdparagraph:P552:begin
He learned reading after he was twelve (155) years (154. c) old.
% mdparagraph:P552:end

% mdparagraph:P553:begin
\textbf{Modern help — available throughout this item.}
% mdparagraph:P553:end

\end{minipage}\par\addvspace{4pt}

% teaching: U15-COOK-NOM-31-11
% support: U15-COOK-NOM-31-11
% mdparagraph:P554:begin
Bēcrǣding means “reading”. Use its singular; the displayed plurals are for reference.
% mdparagraph:P554:end

% teaching: NEW-U15-LOCAL-HELPS
% teaching: U15-H-31-11
% support: U15-H-31-11
% mdparagraph:P555:begin
Use æfter ðǣm ðe … wæs for “after he was …”. Here ðe belongs to the supplied after-conjunction and has no separate subject or object role. In twęlf wintra eald, “twelve years old”, wintra is genitive plural with the numeral and adjective of age. Keep both expressions available as help.
% mdparagraph:P555:end

% output: EX-X-XI-31-11-O001
\reserve{5\baselineskip}
\par\begin{minipage}{\linewidth}\raggedright
% mdparagraph:P556:begin
\textbf{Cook X, XI.31.12}
% mdparagraph:P556:end

% item: EX-X-XI-31-12
% mdparagraph:P557:begin
His mother had a book of English poems, which she promised (om. to) whichever of her two sons (om. should) first learn it.
% mdparagraph:P557:end

% mdparagraph:P558:begin
\textbf{Modern help — available throughout this item.}
% mdparagraph:P558:end

\end{minipage}\par\addvspace{4pt}

% teaching: U15-COOK-NOM-31-12
% support: U15-COOK-NOM-31-12
% mdparagraph:P559:begin
Bōc is accusative singular “book”; bēc belongs to other cases in its table. Lēoða is genitive plural “poems”.
% mdparagraph:P559:end

% teaching: NEW-U15-LOCAL-HELPS
% teaching: U15-H-31-12
% support: U15-H-31-12
% mdparagraph:P560:begin
Use Engliscra lēoða, “of English poems”. The supplied recipient expression is swā hwæðrum hiere twēgra suna swā hīe ǣrest leornode. Hwæðrum is dative because the whole expression is the recipient of the promise; the same person is understood as the learner in the subordinate clause. Leornode is prospective preterite subjunctive, identical in spelling to the indicative; the learning is still to occur. This case attraction is supplied. In the independent free-choice exercises, keep the internal subject nominative and resume it with an explicit dative recipient.
% mdparagraph:P560:end

% output: EX-X-XI-31-12-O001
\reserve{5\baselineskip}
\par\begin{minipage}{\linewidth}\raggedright
% mdparagraph:P561:begin
\textbf{Cook X, XI.31.13}
% mdparagraph:P561:end

% item: EX-X-XI-31-13
% mdparagraph:P562:begin
Alfred asked her, Wilt thou in fact give that book to one of us two (uncer), (om. to) him who first understands (ongietan) it?
% mdparagraph:P562:end

% mdparagraph:P563:begin
\textbf{Modern help — available throughout this item.}
% mdparagraph:P563:end

\end{minipage}\par\addvspace{4pt}

% teaching: NEW-U15-LOCAL-HELPS
% teaching: NEW-U15-RECORD-AELFRED-U15-SOURCE-NOM-SG
% teaching: U15-H-31-NAMES
% support: U15-H-31-NAMES
% mdparagraph:P564:begin
The names Ælfred, Osburg, Æthelwulf and Asser are supplied; Ælfredes is the genitive for 31.14. Use supplied dative Lǣdene for “Latin”, following Cook’s glossary Lǣden. The Preface reading below retains its printed Lædene without a length mark.
% mdparagraph:P564:end

% output: EX-X-XI-31-13-O001
\reserve{5\baselineskip}
\par\begin{minipage}{\linewidth}\raggedright
% mdparagraph:P565:begin
\textbf{Cook X, XI.31.14}
% mdparagraph:P565:end

% item: EX-X-XI-31-14
% mdparagraph:P566:begin
Asser, Alfred’s bishop, wrote a book of (153) the king’s life.
% mdparagraph:P566:end

% mdparagraph:P567:begin
\textbf{Modern help — available throughout this item.}
% mdparagraph:P567:end

\end{minipage}\par\addvspace{4pt}

% teaching: NEW-U15-LOCAL-HELPS
% teaching: NEW-U15-RECORD-AELFRED-U15-SOURCE-NOM-SG
% teaching: U15-H-31-NAMES
% support: U15-H-31-NAMES
% mdparagraph:P568:begin
The names Ælfred, Osburg, Æthelwulf and Asser are supplied; Ælfredes is the genitive for 31.14. Use supplied dative Lǣdene for “Latin”, following Cook’s glossary Lǣden. The Preface reading below retains its printed Lædene without a length mark.
% mdparagraph:P568:end

% output: EX-X-XI-31-14-O001
\reserve{7\baselineskip}
% mdparagraph:P569:begin
\section*{8. Readings}\BookHeading{nav-ch15-learner-0040}{subsection}{8. Readings}{1}
% mdparagraph:P569:end

% mdparagraph:P570:begin
The Old English passages and the numbered original notes are exact source material, with declared typographic normalization. Modern vocabulary and construction helps precede each passage and remain available for all five tasks on it. Original note numbers are prefixed D, P or R here only to distinguish the three sets. Cook’s grammatical references and historical questions are retained as apparatus; answer only the tasks explicitly assigned below.
% mdparagraph:P570:end

% teaching: NEW-U15-GRAM-U15-SOURCE-ARGUMENT
% mdparagraph:P571:begin
Attend to each passage’s movement: a rebuke becomes a promise, an account of translation becomes a request to readers, and an appeal becomes a series of commands and a concluding wish. Distinguish rhetorical questions, assertions, purposes and petitions.
% mdparagraph:P571:end

\reserve{7\baselineskip}
% mdparagraph:P572:begin
\subsection*{D. Duties of the rich toward the poor: the opening argument}\BookHeading{nav-ch15-learner-0041}{subsubsection}{D. Duties of the rich toward the poor: the opening argument}{0}
% mdparagraph:P572:end

% mdparagraph:P573:begin
Cook, First Book, printed p. 159 (PDF p. 175). This is the complete first argument, including its embedded gospel saying. The Savior’s outer speech continues on pp. 160–161; its opening quotation mark is intentionally left unmatched in this selection.
% mdparagraph:P573:end

% mdparagraph:P574:begin
\textbf{Modern vocabulary — supplied for this reading.} Forms are glossed in their local senses; the parses identify the forms encountered, not complete new paradigms.
% mdparagraph:P574:end

% teaching: NEW-U15-RECORD-AEGHWILC-U15-SOURCE-DAT-M-SG
% teaching: NEW-U15-RECORD-AEGHWILC-U15-SOURCE-SENSE
% teaching: U15-READ-LEX-AEGHWILC
% support: U15-READ-LEX-AEGHWILC
% mdparagraph:P575:begin
\textbf{ǣghwilcum} — each; every (strong dative masculine singular).
% mdparagraph:P575:end

% teaching: NEW-U15-RECORD-AFIERRAN-U15-SOURCE-PRES-1-SG
% teaching: NEW-U15-RECORD-AFIERRAN-U15-SOURCE-SENSE
% teaching: U15-READ-LEX-AFIERRAN
% support: U15-READ-LEX-AFIERRAN
% mdparagraph:P576:begin
\textbf{āfierre} — remove; take away (present indicative first singular, future reference).
% mdparagraph:P576:end

% teaching: NEW-U15-RECORD-CWETHAN-U15-SOURCE-GE-PAST-PART
% teaching: NEW-U15-RECORD-CWETHAN-U15-SOURCE-SENSE
% teaching: U15-READ-LEX-CWETHAN
% support: U15-READ-LEX-CWETHAN
% mdparagraph:P577:begin
\textbf{gecweden} — said (uninflected ge-prefixed past participle in passive).
% mdparagraph:P577:end

% teaching: NEW-U15-RECORD-DAED-U15-SOURCE-ACC-SG
% teaching: NEW-U15-RECORD-DAED-U15-SOURCE-SENSE
% teaching: U15-READ-LEX-DAED
% support: U15-READ-LEX-DAED
% mdparagraph:P578:begin
\textbf{dǣd} — deed (endingless accusative feminine singular).
% mdparagraph:P578:end

% teaching: NEW-U15-RECORD-DREAM-U15-SOURCE-ACC-SG
% teaching: NEW-U15-RECORD-DREAM-U15-SOURCE-SENSE
% teaching: U15-READ-LEX-DREAM
% support: U15-READ-LEX-DREAM
% mdparagraph:P579:begin
\textbf{drēam} — joy; bliss (accusative masculine singular).
% mdparagraph:P579:end

% teaching: NEW-U15-RECORD-FAESTHAFOL-U15-SOURCE-NOM-M-SG-STRONG
% teaching: NEW-U15-RECORD-FAESTHAFOL-U15-SOURCE-SENSE
% teaching: U15-READ-LEX-FAESTHAFOL
% support: U15-READ-LEX-FAESTHAFOL
% mdparagraph:P580:begin
\textbf{fæsthafol} — tenacious; withholding, with genitive of what is held (predicate nominative masculine singular).
% mdparagraph:P580:end

% teaching: NEW-U11-LOCAL-HELPS
% teaching: NEW-U15-RECORD-FOR-U15-SOURCE-WHY
% teaching: U15-READ-LEX-FOR
% support: U15-READ-LEX-FOR
% mdparagraph:P581:begin
\textbf{for} — for what reason; why, in for hwȳ/\allowbreak{}hwon. \textbf{for} — for; for the sake of, selected accusative/\allowbreak{}dative expressions supplied locally.
% mdparagraph:P581:end

% teaching: NEW-U15-RECORD-GEDON-U15-SOURCE-PRES-3-SG
% teaching: NEW-U15-RECORD-GEDON-U15-SOURCE-SENSE
% teaching: U15-READ-LEX-GEDON
% support: U15-READ-LEX-GEDON
% mdparagraph:P582:begin
\textbf{gedēð} — do; perform (present indicative third singular).
% mdparagraph:P582:end

% teaching: NEW-U15-RECORD-GESELLAN-U15-SOURCE-PRES-1-SG
% teaching: NEW-U15-RECORD-GESELLAN-U15-SOURCE-SENSE
% teaching: U15-READ-LEX-GESELLAN
% support: U15-READ-LEX-GESELLAN
% mdparagraph:P583:begin
\textbf{gesęlle} — give (present indicative first singular, future reference).
% mdparagraph:P583:end

% teaching: NEW-U15-RECORD-GETHENCEAN-U15-SOURCE-INF
% teaching: NEW-U15-RECORD-GETHENCEAN-U15-SOURCE-SENSE
% teaching: U15-READ-LEX-GETHENCEAN
% support: U15-READ-LEX-GETHENCEAN
% mdparagraph:P584:begin
\textbf{geðęncean} — remember; bear in mind (plain infinitive after noldest).
% mdparagraph:P584:end

% teaching: NEW-U15-RECORD-GIEFU-U15-SOURCE-GEN-SG
% teaching: NEW-U15-RECORD-GIEFU-U15-SOURCE-SENSE
% teaching: U15-READ-LEX-GIEFU
% support: U15-READ-LEX-GIEFU
% mdparagraph:P585:begin
\textbf{giefe} — gift (genitive feminine singular dependent on rēcelēasodest).
% mdparagraph:P585:end

% teaching: NEW-U15-RECORD-HUNDTEONTIG-U15-SOURCE-DAT-PL
% teaching: NEW-U15-RECORD-HUNDTEONTIG-U15-SOURCE-SENSE
% teaching: U15-READ-LEX-HUNDTEONTIG
% support: U15-READ-LEX-HUNDTEONTIG
% mdparagraph:P586:begin
\textbf{hundtēontigum} — a hundredfold; a hundred in return (dative plural after mid in fixed repayment expression).
% mdparagraph:P586:end

% teaching: NEW-U15-RECORD-HWA-U15-SOURCE-WHY
% teaching: U15-READ-LEX-HWA
% support: U15-READ-LEX-HWA
% mdparagraph:P587:begin
\textbf{hwon} — what (reason), instrumental hwon in fixed why expressions (instrumental singular neuter).
% mdparagraph:P587:end

% teaching: NEW-U15-RECORD-OF-U15-SOURCE-PARTITIVE
% teaching: U15-READ-LEX-OF
% support: U15-READ-LEX-OF
% mdparagraph:P588:begin
\textbf{of} — of; out of a set, selecting one of the least.
% mdparagraph:P588:end

% teaching: NEW-U15-RECORD-ON-U15-SOURCE-WORLDLY-LIFE
% teaching: U15-READ-LEX-ON
% support: U15-READ-LEX-ON
% mdparagraph:P589:begin
\textbf{on} — in worldly life, temporal/\allowbreak{}situational expression.
% mdparagraph:P589:end

% teaching: NEW-U15-RECORD-RECELEASIAN-U15-SOURCE-PRET-2-SG
% teaching: NEW-U15-RECORD-RECELEASIAN-U15-SOURCE-SENSE
% teaching: U15-READ-LEX-RECELEASIAN
% support: U15-READ-LEX-RECELEASIAN
% mdparagraph:P590:begin
\textbf{rēcelēasodest} — neglect; despise (preterite indicative second singular; genitive object).
% mdparagraph:P590:end

% teaching: NEW-U15-RECORD-SECGAN-U15-SOURCE-GE-PAST-PART
% teaching: NEW-U15-RECORD-SECGAN-U15-SOURCE-SENSE
% teaching: U15-READ-LEX-SECGAN
% support: U15-READ-LEX-SECGAN
% mdparagraph:P591:begin
\textbf{gesǣd} — said (uninflected ge-prefixed past participle, pleonastic with gecweden).
% mdparagraph:P591:end

% teaching: NEW-U15-RECORD-SELEN-U15-SOURCE-DAT-SG
% teaching: NEW-U15-RECORD-SELEN-U15-SOURCE-SENSE
% teaching: U15-READ-LEX-SELEN
% support: U15-READ-LEX-SELEN
% mdparagraph:P592:begin
\textbf{sęlene} — bounty; bestowal (dative feminine singular after fram).
% mdparagraph:P592:end

% teaching: NEW-U15-RECORD-SIMLE-U15-SOURCE-INVARIANT
% teaching: NEW-U15-RECORD-SIMLE-U15-SOURCE-SENSE
% teaching: U15-READ-LEX-SIMLE
% support: U15-READ-LEX-SIMLE
% mdparagraph:P593:begin
\textbf{simle} — always.
% mdparagraph:P593:end

% teaching: NEW-U15-RECORD-SWA-HWAET-SWA-U15-SOURCE-ACC-N-SG
% teaching: NEW-U15-RECORD-SWA-HWAET-SWA-U15-SOURCE-SENSE
% teaching: U15-READ-LEX-SWA-HWAET-SWA
% support: U15-READ-LEX-SWA-HWAET-SWA
% mdparagraph:P594:begin
\textbf{Swā-hwæt-swā} — whatsoever; whatever (free-choice relative pronoun, accusative neuter object).
% mdparagraph:P594:end

% teaching: NEW-U14-LOCAL-HELPS
% teaching: NEW-U15-RECORD-TO-U15-SOURCE-WHY
% teaching: U15-READ-LEX-TO
% support: U15-READ-LEX-TO
% mdparagraph:P595:begin
\textbf{tō} — to, the person addressed in speech. \textbf{tō} — for what purpose; why, in tō hwon.
% mdparagraph:P595:end

% teaching: NEW-U15-RECORD-WLANC-U15-SOURCE-DAT-M-SG-WEAK
% teaching: NEW-U15-RECORD-WLANC-U15-SOURCE-SENSE
% teaching: U15-READ-LEX-WLANC
% support: U15-READ-LEX-WLANC
% mdparagraph:P596:begin
\textbf{wlancan} — proud; wealthy, here the rich man addressed (substantival weak dative masculine singular).
% mdparagraph:P596:end

% teaching: NEW-U15-RECORD-WORULDLIF-U15-SOURCE-DAT-SG
% teaching: NEW-U15-RECORD-WORULDLIF-U15-SOURCE-SENSE
% teaching: U15-READ-LEX-WORULDLIF
% support: U15-READ-LEX-WORULDLIF
% mdparagraph:P597:begin
\textbf{woruldlīfe} — worldly life (dative neuter singular after on).
% mdparagraph:P597:end

\reserve{5\baselineskip}
% mdparagraph:P598:begin
\textbf{Modern construction help.}
% mdparagraph:P598:end

% teaching: NEW-U15-RECORD-U15-SOURCE-D-ORIENTATION
% teaching: U15-READ-D-ORIENTATION
% support: U15-READ-D-ORIENTATION
% mdparagraph:P599:begin
The Savior addresses a rich man. This selection contains the complete opening argument and its embedded gospel quotation. The Savior’s speech continues on the next page: the unmatched outer quotation mark is present in this bounded source selection.
% mdparagraph:P599:end

% teaching: NEW-U15-RECORD-U15-SOURCE-D-GENITIVE-CONTROLS
% teaching: U15-READ-D-GENITIVE-CONTROLS
% support: U15-READ-D-GENITIVE-CONTROLS
% mdparagraph:P600:begin
For hwȳ and tō hwon ask why. Fæsthafol governs the genitive mīnra gōda, “holding tightly to my goods”; rēcelēasodest takes genitive ðǣre giefe, “neglected the gift.” Relative ðe introduces the following clauses; long ðē is the personal pronoun “you/\allowbreak{}to you.”
% mdparagraph:P600:end

% teaching: NEW-U15-RECORD-U15-SOURCE-D-FUTURE-RESPONSE
% teaching: U15-READ-D-FUTURE-RESPONSE
% support: U15-READ-D-FUTURE-RESPONSE
% mdparagraph:P601:begin
Āfierre, bist, forgielde and gesęlle have present morphology with future reference in this promised response. Noldest is “you would not/\allowbreak{}refused”; geðęncean is its infinitive. Wille expresses the present speaker’s intention even within a remembered proposition: no mechanical backshift is required.
% mdparagraph:P601:end

% teaching: NEW-U15-RECORD-U15-SOURCE-D-GENERIC-AGENT
% teaching: U15-READ-D-GENERIC-AGENT
% support: U15-READ-D-GENERIC-AGENT
% mdparagraph:P602:begin
Ǣghwilcum męnn is the dative recipient; āne gōde dǣd is the accusative object. The subject of gedēð is generic mann, “one/\allowbreak{}a person.” Mid hundtēontigum means “with a hundredfold.” The singular him refers back to the person.
% mdparagraph:P602:end

% teaching: NEW-U15-RECORD-U15-SOURCE-D-EMBEDDED-QUOTATION
% teaching: U15-READ-D-EMBEDDED-QUOTATION
% support: U15-READ-D-EMBEDDED-QUOTATION
% mdparagraph:P603:begin
Gecweden and gesǣd are two past participles with is; Cook calls the pairing pleonastic. Swā-hwæt-swā means “whatever,” the object of sęllað. Ānum of mīnum ðǣm lǣstum means “to one of my least ones”; lǣstum is substantival and dative plural. The repeated gē/\allowbreak{}sęllað address the quoted audience.
% mdparagraph:P603:end

% teaching: NEW-U15-RECORD-U15-SOURCE-D-EXCHANGE
% teaching: U15-READ-D-EXCHANGE
% support: U15-READ-D-EXCHANGE
% mdparagraph:P604:begin
Wið ðǣm means “in return for that,” with dative singular. Drēam here means joy or bliss, not a dream. Ēcne agrees with accusative masculine drēam.
% mdparagraph:P604:end

\reserve{5\baselineskip}
% mdparagraph:P605:begin
\textbf{Cook’s text}
% mdparagraph:P605:end

% mdparagraph:P606:begin
IX. DUTIES OF THE RICH TOWARD THE POOR.
% mdparagraph:P606:end

% mdparagraph:P607:begin
(From the same Homily as the last.)
% mdparagraph:P607:end

% teaching: NEW-U15-READ-FB-IX-DUTIES-FIRST-ARGUMENT
\par\begin{minipage}{\linewidth}\raggedright
% passage: FB-IX-DUTIES-FIRST-ARGUMENT
% mdparagraph:P608:begin
\readingtext{Se Hǣlend cwæð tō ðǣm wlancan\textsuperscript{D1}: “For hwȳ wǣre ðū swā fæsthafol mīnra gōda, ðe ic ðē sealde? Tō hwon\textsuperscript{D2} rēcelēasodest ðū ðǣre giefe, ðe ic ðē geaf? Ic ðē nū āfierre\textsuperscript{D3} fram mīnre sęlene, ðe ic ðē forgeaf; ðonne bist\textsuperscript{D3} ðū wǣdla on woruldlīfe. For hwon\textsuperscript{D2} noldest\textsuperscript{D4} ðū geðęncean ðæt ic wille\textsuperscript{D5} forgieldan ǣghwilcum męnn āne gōde dǣd, ðe for mīnum naman mann gedēð? Mid hundtēontigum ic hit him forgielde,\textsuperscript{D3} swā hit is on mīnum godspelle gecweden and gesǣd,\textsuperscript{D6} ‘Swā-hwæt-swā\textsuperscript{D7} gē sęllað ānum of mīnum ðǣm lǣstum,\textsuperscript{D8} gē hit simle mē sęllað,\textsuperscript{D9} and ic ēow wið\textsuperscript{D10} ðǣm gesęlle\textsuperscript{D3} ēcne drēam\textsuperscript{D11} on heofonum.’}
% mdparagraph:P608:end

\end{minipage}\par
\reserve{5\baselineskip}
% mdparagraph:P609:begin
\textbf{Original notes — Cook.}
% mdparagraph:P609:end

% mdparagraph:P610:begin
\textsuperscript{D1}: From what OE. word is the Mod. Eng. rich derived (see Skeat, Prin., p. 61)? From what OE. meaning is the modern signification derived?
% mdparagraph:P610:end

% mdparagraph:P611:begin
\textsuperscript{D2}: Note Wulfstan’s use of tō hwon, for hwon, in the sense of why. See 88.
% mdparagraph:P611:end

% mdparagraph:P612:begin
\textsuperscript{D3}: See 188.
% mdparagraph:P612:end

% mdparagraph:P613:begin
\textsuperscript{D4}: See 139.
% mdparagraph:P613:end

% mdparagraph:P614:begin
\textsuperscript{D5}: Why should not the preterit be used here?
% mdparagraph:P614:end

% mdparagraph:P615:begin
\textsuperscript{D6}: Pleonastic.
% mdparagraph:P615:end

% mdparagraph:P616:begin
\textsuperscript{D7}: What portion of this is lost, and how is it replaced, in the Mod. Eng. whatsoever?
% mdparagraph:P616:end

% mdparagraph:P617:begin
\textsuperscript{D8}: See 66.
% mdparagraph:P617:end

% mdparagraph:P618:begin
\textsuperscript{D9}: Cf. the form of this sentence with that on p. 135, l. 14.
% mdparagraph:P618:end

% mdparagraph:P619:begin
\textsuperscript{D10}: = in return for. How is this to be reconciled with other senses of wið?
% mdparagraph:P619:end

% mdparagraph:P620:begin
\textsuperscript{D11}: Not dream, but joy, bliss.
% mdparagraph:P620:end

% mdparagraph:P621:begin
\textbf{Tasks D01–D05 — new teaching material.} Keep the modern vocabulary, construction helps and original notes available for every task.
% mdparagraph:P621:end

\reserve{5\baselineskip}
\par\begin{minipage}{\linewidth}\raggedright
% mdparagraph:P622:begin
\textbf{D01}
% mdparagraph:P622:end

% item: NEW-U15-D01
% mdparagraph:P623:begin
State the rich man’s fault, the consequence announced, and the promise made to a giver. Identify who speaks and to whom.
% mdparagraph:P623:end

% output: NEW-U15-D01-O001
% output: NEW-U15-D01-O002
% output: NEW-U15-D01-O003
% output: NEW-U15-D01-O004
% mdparagraph:P624:begin
\textbf{Responses, in order:} (1) fault; (2) consequence; (3) promise; (4) speaker and addressee.
% mdparagraph:P624:end

\end{minipage}\par\addvspace{4pt}

\par\begin{minipage}{\linewidth}\raggedright
% mdparagraph:P625:begin
\textbf{D02}
% mdparagraph:P625:end

% item: NEW-U15-D02
% mdparagraph:P626:begin
In the question beginning For hwon noldest ðū, explain the negative verb and the infinitive it governs. Explain why ic wille inside the following content clause can remain present rather than mechanically copying the question’s past tense.
% mdparagraph:P626:end

% output: NEW-U15-D02-O001
% output: NEW-U15-D02-O002
% mdparagraph:P627:begin
\textbf{Responses, in order:} (1) noldest and its infinitive; (2) tense of wille in this context.
% mdparagraph:P627:end

\end{minipage}\par\addvspace{4pt}

\par\begin{minipage}{\linewidth}\raggedright
% mdparagraph:P628:begin
\textbf{D03}
% mdparagraph:P628:end

% item: NEW-U15-D03
% mdparagraph:P629:begin
In the gospel saying beginning Swā-hwæt-swā, identify what is given, its immediate recipient, the recipient identified by the speaker, and the pronoun that resumes the thing given.
% mdparagraph:P629:end

% output: NEW-U15-D03-O001
% output: NEW-U15-D03-O002
% output: NEW-U15-D03-O003
% output: NEW-U15-D03-O004
% mdparagraph:P630:begin
\textbf{Responses, in order:} (1) thing given; (2) immediate recipient; (3) recipient identified by the speaker; (4) resuming pronoun.
% mdparagraph:P630:end

\end{minipage}\par\addvspace{4pt}

\par\begin{minipage}{\linewidth}\raggedright
% mdparagraph:P631:begin
\textbf{D04}
% mdparagraph:P631:end

% item: NEW-U15-D04
% mdparagraph:P632:begin
Identify the case relations in ānum of mīnum ðǣm lǣstum and wið ðǣm. Use the local helps; explain the difference between choosing a recipient from a group and receiving something in return.
% mdparagraph:P632:end

% output: NEW-U15-D04-O001
% output: NEW-U15-D04-O002
% mdparagraph:P633:begin
\textbf{Responses, in order:} (1) case relations in the first phrase; (2) case relation and meaning in wið ðǣm.
% mdparagraph:P633:end

\end{minipage}\par\addvspace{4pt}

\par\begin{minipage}{\linewidth}\raggedright
% mdparagraph:P634:begin
\textbf{D05}
% mdparagraph:P634:end

% item: NEW-U15-D05
% mdparagraph:P635:begin
Translate the complete assigned opening argument, from Se Hǣlend through heofonum. Preserve the change from rebuke to promise and the embedded gospel quotation.
% mdparagraph:P635:end

% output: NEW-U15-D05-O001
\end{minipage}\par\addvspace{4pt}

\reserve{7\baselineskip}
% mdparagraph:P636:begin
\subsection*{P. Alfred’s preface to Boethius}\BookHeading{nav-ch15-learner-0042}{subsubsection}{P. Alfred’s preface to Boethius}{0}
% mdparagraph:P636:end

% mdparagraph:P637:begin
Cook, First Book, printed p. 162 (PDF p. 178). The complete preface.
% mdparagraph:P637:end

% mdparagraph:P638:begin
\textbf{Modern vocabulary — supplied for this reading.} Forms are glossed in their local senses; the parses identify the forms encountered, not complete new paradigms.
% mdparagraph:P638:end

% teaching: NEW-U15-LOCAL-HELPS
% teaching: NEW-U15-RECORD-AELFRED-U15-SOURCE-NOM-SG
% teaching: U15-READ-LEX-AELFRED
% support: U15-READ-LEX-AELFRED
% mdparagraph:P639:begin
\textbf{Ælfred} — Alfred (nominative masculine singular personal name).
% mdparagraph:P639:end

% teaching: NEW-U15-RECORD-AEMETTA-U15-SOURCE-DAT-SG
% teaching: NEW-U15-RECORD-AEMETTA-U15-SOURCE-SENSE
% teaching: U15-READ-LEX-AEMETTA
% support: U15-READ-LEX-AEMETTA
% mdparagraph:P640:begin
\textbf{ǣmettan} — leisure; available time (dative masculine singular after be).
% mdparagraph:P640:end

% teaching: NEW-U15-RECORD-ANDGIET-U15-SOURCE-ACC-SG
% teaching: NEW-U15-RECORD-ANDGIET-U15-SOURCE-DAT-SG
% teaching: NEW-U15-RECORD-ANDGIET-U15-SOURCE-GEN-SG
% teaching: NEW-U15-RECORD-ANDGIET-U15-SOURCE-SENSE
% teaching: NEW-U15-RECORD-ANDGIET-U15-SOURCE-UNDERSTANDING
% teaching: U15-READ-LEX-ANDGIET
% support: U15-READ-LEX-ANDGIET
% mdparagraph:P641:begin
\textbf{andgiet} — meaning; sense (accusative neuter singular object of sętte). \textbf{andgiete} — meaning; sense (dative neuter singular after of). \textbf{andgietes} — understanding (genitive neuter singular dependent on mǣðe).
% mdparagraph:P641:end

% teaching: NEW-U15-RECORD-ANDGIETFULLICE-U15-SOURCE-SENSE
% teaching: NEW-U15-RECORD-ANDGIETFULLICE-U15-SOURCE-SUPERL
% teaching: U15-READ-LEX-ANDGIETFULLICE
% support: U15-READ-LEX-ANDGIETFULLICE
% mdparagraph:P642:begin
\textbf{andgietfullicost} — most intelligibly (superlative adverb).
% mdparagraph:P642:end

% teaching: NEW-U06-SUPFORM-be-invariant-source
% teaching: NEW-U06-SUPSENSE-be-according-source
% teaching: NEW-U14-LOCAL-HELPS
% teaching: NEW-U15-RECORD-BE-U15-SOURCE-WORD-BY-WORD
% teaching: U15-READ-LEX-BE
% support: U15-READ-LEX-BE
% mdparagraph:P643:begin
\textbf{be} — word by word, literal correspondence in word be worde. \textbf{be} — according to.
% mdparagraph:P643:end

% teaching: NEW-U15-RECORD-BECUMAN-U15-SOURCE-PRET-3-PL
% teaching: NEW-U15-RECORD-BECUMAN-U15-SOURCE-SENSE
% teaching: U15-READ-LEX-BECUMAN
% support: U15-READ-LEX-BECUMAN
% mdparagraph:P644:begin
\textbf{becōmon} — come upon; befall (preterite indicative third plural).
% mdparagraph:P644:end

% teaching: NEW-U15-RECORD-BISGIAN-U15-SOURCE-PRET-3-PL
% teaching: NEW-U15-RECORD-BISGIAN-U15-SOURCE-SENSE
% teaching: U15-READ-LEX-BISGIAN
% support: U15-READ-LEX-BISGIAN
% mdparagraph:P645:begin
\textbf{bisgodon} — occupy; trouble (preterite indicative third plural; syncopated form).
% mdparagraph:P645:end

% teaching: NEW-U15-RECORD-BISGU-U15-SOURCE-NOM-PL
% teaching: NEW-U15-RECORD-BISGU-U15-SOURCE-SENSE
% teaching: U15-READ-LEX-BISGU
% support: U15-READ-LEX-BISGU
% mdparagraph:P646:begin
\textbf{bisga} — occupations; cares (nominative feminine plural subject).
% mdparagraph:P646:end

% teaching: NEW-U15-RECORD-BOCLAEDEN-U15-SOURCE-DAT-SG
% teaching: NEW-U15-RECORD-BOCLAEDEN-U15-SOURCE-SENSE
% teaching: U15-READ-LEX-BOCLAEDEN
% support: U15-READ-LEX-BOCLAEDEN
% mdparagraph:P647:begin
\textbf{Bōclædene} — book-Latin; Latin (dative neuter singular after of).
% mdparagraph:P647:end

% teaching: NEW-U15-RECORD-EARFOTHRIME-U15-SOURCE-NOM-F-PL
% teaching: NEW-U15-RECORD-EARFOTHRIME-U15-SOURCE-SENSE
% teaching: U15-READ-LEX-EARFOTHRIME
% support: U15-READ-LEX-EARFOTHRIME
% mdparagraph:P648:begin
\textbf{earfoðrīme} — hard to count; innumerable (predicate strong nominative feminine plural).
% mdparagraph:P648:end

% teaching: NEW-U15-RECORD-EFT-U15-SOURCE-INVARIANT
% teaching: NEW-U15-RECORD-EFT-U15-SOURCE-SENSE
% teaching: U15-READ-LEX-EFT
% support: U15-READ-LEX-EFT
% mdparagraph:P649:begin
\textbf{eft} — again; afterwards.
% mdparagraph:P649:end

% teaching: NEW-U15-LOCAL-HELPS
% teaching: NEW-U15-RECORD-ENGLISC-U15-SOURCE-DAT-N-SG
% teaching: U15-READ-LEX-ENGLISC
% support: U15-READ-LEX-ENGLISC
% mdparagraph:P650:begin
\textbf{Engliscum} — English (strong dative neuter singular).
% mdparagraph:P650:end

% teaching: NEW-U15-RECORD-ENGLISC-LANGUAGE-U15-SOURCE-ACC-SG
% teaching: NEW-U15-RECORD-ENGLISC-LANGUAGE-U15-SOURCE-SENSE
% teaching: U15-READ-LEX-ENGLISC-LANGUAGE
% support: U15-READ-LEX-ENGLISC-LANGUAGE
% mdparagraph:P651:begin
\textbf{Englisc} — English, the language (accusative language name after on, goal of translation).
% mdparagraph:P651:end

% teaching: NEW-U11-LOCAL-HELPS
% teaching: NEW-U15-RECORD-FOR-U15-SOURCE-CAUSE
% teaching: U15-READ-LEX-FOR
% support: U15-READ-LEX-FOR
% mdparagraph:P652:begin
\textbf{for} — because of, with dative of hindering circumstance. \textbf{for} — for; for the sake of, selected accusative/\allowbreak{}dative expressions supplied locally.
% mdparagraph:P652:end

% teaching: NEW-U15-RECORD-FOR-THAM-THE-U15-SOURCE-INVARIANT
% teaching: NEW-U15-RECORD-FOR-THAM-THE-U15-SOURCE-SENSE
% teaching: U15-READ-LEX-FOR-THAM-THE
% support: U15-READ-LEX-FOR-THAM-THE
% mdparagraph:P653:begin
\textbf{for-ðām-ðe} — because (causal conjunction).
% mdparagraph:P653:end

% teaching: NEW-U05-SUPFORM-gedon-past-participle
% teaching: NEW-U15-RECORD-GEDON-U15-SOURCE-ARRANGEMENT
% teaching: U15-READ-LEX-GEDON
% support: U15-READ-LEX-GEDON
% mdparagraph:P654:begin
\textbf{gedōn} — made; arranged, describing the present form of the translation (past participle).
% mdparagraph:P654:end

% teaching: NEW-U15-RECORD-GELEORNIAN-U15-SOURCE-PAST-PART
% teaching: NEW-U15-RECORD-GELEORNIAN-U15-SOURCE-SENSE
% teaching: U15-READ-LEX-GELEORNIAN
% support: U15-READ-LEX-GELEORNIAN
% mdparagraph:P655:begin
\textbf{geleornod} — learn; study (past participle in hæfde geleornod).
% mdparagraph:P655:end

% teaching: NEW-U15-RECORD-GERECCAN-U15-SOURCE-INF
% teaching: NEW-U15-RECORD-GERECCAN-U15-SOURCE-SENSE
% teaching: U15-READ-LEX-GERECCAN
% support: U15-READ-LEX-GERECCAN
% mdparagraph:P656:begin
\textbf{geręccan} — explain; render (plain infinitive after meahte).
% mdparagraph:P656:end

% teaching: NEW-U14-PAR-GEWENDAN
% teaching: NEW-U15-RECORD-GEWENDAN-U15-SOURCE-TRANSLATE
% teaching: U15-READ-LEX-GEWENDAN
% support: U15-READ-LEX-GEWENDAN
% mdparagraph:P657:begin
\textbf{gewęnd} — translate; turn from one language into another (past participle).
% mdparagraph:P657:end

% teaching: NEW-U15-LOCAL-HELPS
% teaching: NEW-U15-RECORD-LAEDEN-U15-SOURCE-DAT-SG
% teaching: U15-READ-LEX-LAEDEN
% support: U15-READ-LEX-LAEDEN
% mdparagraph:P658:begin
\textbf{Lædene} — Latin, the language (dative neuter singular language name after of).
% mdparagraph:P658:end

% teaching: NEW-U15-RECORD-LYSTAN-U15-SOURCE-PRES-SUBJ-3-SG
% teaching: NEW-U15-RECORD-LYSTAN-U15-SOURCE-SENSE
% teaching: U15-READ-LEX-LYSTAN
% support: U15-READ-LEX-LYSTAN
% mdparagraph:P659:begin
\textbf{lyste} — please; be a desire to (present subjunctive third singular impersonal; understood experiencer accusative from ǣlcne).
% mdparagraph:P659:end

% teaching: NEW-U15-RECORD-MAETH-U15-SOURCE-DAT-SG
% teaching: NEW-U15-RECORD-MAETH-U15-SOURCE-SENSE
% teaching: U15-READ-LEX-MAETH
% support: U15-READ-LEX-MAETH
% mdparagraph:P660:begin
\textbf{mǣðe} — measure; capacity (dative feminine singular after be).
% mdparagraph:P660:end

% teaching: NEW-U15-RECORD-MANIGFEALD-U15-SOURCE-DAT-F-PL-STRONG
% teaching: NEW-U15-RECORD-MANIGFEALD-U15-SOURCE-SENSE
% teaching: U15-READ-LEX-MANIGFEALD
% support: U15-READ-LEX-MANIGFEALD
% mdparagraph:P661:begin
\textbf{manigfealdum} — manifold; many (strong dative feminine plural).
% mdparagraph:P661:end

% teaching: NEW-U15-RECORD-MISLIC-U15-SOURCE-DAT-F-PL-STRONG
% teaching: NEW-U15-RECORD-MISLIC-U15-SOURCE-SENSE
% teaching: U15-READ-LEX-MISLIC
% support: U15-READ-LEX-MISLIC
% mdparagraph:P662:begin
\textbf{mislicum} — various (strong dative feminine plural).
% mdparagraph:P662:end

% teaching: NEW-U13-SOURCE-HELPS
% teaching: U15-READ-LEX-MOD
% support: U15-READ-LEX-MOD
% mdparagraph:P663:begin
\textbf{mōde} — mind; attention (dative neuter singular).
% mdparagraph:P663:end

% teaching: NEW-U15-RECORD-OF-U15-SOURCE-CORRESPONDENCE
% teaching: U15-READ-LEX-OF
% support: U15-READ-LEX-OF
% mdparagraph:P664:begin
\textbf{of} — from; corresponding to, sense for sense.
% mdparagraph:P664:end

% teaching: NEW-U15-RECORD-ON-U15-SOURCE-DURING
% teaching: U15-READ-LEX-ON
% support: U15-READ-LEX-ON
% mdparagraph:P665:begin
\textbf{on} — during, with dative time.
% mdparagraph:P665:end

% teaching: NEW-U15-RECORD-ONGIETAN-U15-SOURCE-PRES-SUBJ-3-SG
% teaching: NEW-U15-VOC-ONGIETAN
% teaching: U15-READ-LEX-ONGIETAN
% support: U15-READ-LEX-ONGIETAN
% mdparagraph:P666:begin
\textbf{ongiete} — understand; perceive (present subjunctive third singular in condition).
% mdparagraph:P666:end

% teaching: NEW-U09-VOC-raedan
% teaching: U15-READ-LEX-RAEDAN
% support: U15-READ-LEX-RAEDAN
% mdparagraph:P667:begin
\textbf{rǣdan} — read (plain infinitive).
% mdparagraph:P667:end

% teaching: NEW-U15-RECORD-RIHTLICE-U15-SOURCE-COMPARATIVE
% teaching: NEW-U15-RECORD-RIHTLICE-U15-SOURCE-SENSE
% teaching: U15-READ-LEX-RIHTLICE
% support: U15-READ-LEX-RIHTLICE
% mdparagraph:P668:begin
\textbf{rihtlicor} — more correctly (comparative adverb).
% mdparagraph:P668:end

% teaching: NEW-U15-RECORD-SPELL-U15-SOURCE-DAT-SG
% teaching: NEW-U15-RECORD-SPELL-U15-SOURCE-SENSE
% teaching: U15-READ-LEX-SPELL
% support: U15-READ-LEX-SPELL
% mdparagraph:P669:begin
\textbf{spelle} — account; discourse (dative neuter singular after tō).
% mdparagraph:P669:end

% teaching: NEW-U09-SOURCE-HELPS
% teaching: NEW-U09-VOC-swa-swa
% teaching: NEW-U10-SOURCE-HELPS
% teaching: NEW-U12-REUSED-U12-SRC-HELP-15
% teaching: NEW-U12-SOURCE-HELPS
% teaching: NEW-U13-PAR-SWA-SWA
% teaching: NEW-U13-VOC-SWA-SWA
% teaching: NEW-U14-LOCAL-HELPS
% teaching: U15-READ-LEX-SWA-SWA
% support: U15-READ-LEX-SWA-SWA
% mdparagraph:P670:begin
\textbf{swā-swā} — just as; like.
% mdparagraph:P670:end

% teaching: NEW-U15-RECORD-SWEOTOLE-U15-SOURCE-SENSE
% teaching: NEW-U15-RECORD-SWEOTOLE-U15-SOURCE-SUPERL
% teaching: U15-READ-LEX-SWEOTOLE
% support: U15-READ-LEX-SWEOTOLE
% mdparagraph:P671:begin
\textbf{sweotolost} — most clearly (superlative adverb).
% mdparagraph:P671:end

% teaching: NEW-U15-RECORD-THAET-THAET-U15-SOURCE-ACC-N-SG
% teaching: NEW-U15-RECORD-THAET-THAET-U15-SOURCE-SENSE
% teaching: U15-READ-LEX-THAET-THAET
% support: U15-READ-LEX-THAET-THAET
% mdparagraph:P672:begin
\textbf{ðæt-ðæt} — that which; what (demonstrative accusative neuter + relative accusative neuter).
% mdparagraph:P672:end

% teaching: NEW-U15-RECORD-THEAH-U15-SOURCE-NEVERTHELESS
% teaching: U15-READ-LEX-THEAH
% support: U15-READ-LEX-THEAH
% mdparagraph:P673:begin
\textbf{ðēah} — nevertheless; yet (independent adverb).
% mdparagraph:P673:end

% teaching: NEW-U15-RECORD-TO-U15-SOURCE-TRANSLATION-RESULT
% teaching: U15-READ-LEX-TO
% support: U15-READ-LEX-TO
% mdparagraph:P674:begin
\textbf{tō} — into; to, result in a new linguistic/\allowbreak{}literary form.
% mdparagraph:P674:end

% teaching: NEW-U15-RECORD-UNDERFON-U15-SOURCE-PAST-PART
% teaching: NEW-U15-RECORD-UNDERFON-U15-SOURCE-SENSE
% teaching: U15-READ-LEX-UNDERFON
% support: U15-READ-LEX-UNDERFON
% mdparagraph:P675:begin
\textbf{underfangen} — receive; take charge of (past participle in anterior construction hæfde underfangen).
% mdparagraph:P675:end

% teaching: NEW-U15-RECORD-WEALHSTOD-U15-SOURCE-NOM-SG
% teaching: NEW-U15-RECORD-WEALHSTOD-U15-SOURCE-SENSE
% teaching: U15-READ-LEX-WEALHSTOD
% support: U15-READ-LEX-WEALHSTOD
% mdparagraph:P676:begin
\textbf{wealhstōd} — interpreter; translator (predicate nominative masculine singular).
% mdparagraph:P676:end

% teaching: NEW-U15-RECORD-WENDAN-U15-SOURCE-PRET-3-SG
% teaching: NEW-U15-RECORD-WENDAN-U15-SOURCE-SENSE
% teaching: U15-READ-LEX-WENDAN
% support: U15-READ-LEX-WENDAN
% mdparagraph:P677:begin
\textbf{węnde} — translate (preterite indicative third singular).
% mdparagraph:P677:end

% teaching: NEW-U15-RECORD-WESAN-U15-SOURCE-SENSE
% teaching: NEW-U15-RECORD-WESAN-U15-SOURCE-SIND-PRES-3-PL
% teaching: U15-READ-LEX-WESAN
% support: U15-READ-LEX-WESAN
% mdparagraph:P678:begin
\textbf{sind} — are (present indicative third plural).
% mdparagraph:P678:end

% teaching: NEW-U15-RECORD-WITAN-BLAME-U15-SOURCE-PRES-SUBJ-3-SG
% teaching: NEW-U15-RECORD-WITAN-BLAME-U15-SOURCE-SENSE
% teaching: U15-READ-LEX-WITAN-BLAME
% support: U15-READ-LEX-WITAN-BLAME
% mdparagraph:P679:begin
\textbf{wīte} — blame; censure (present subjunctive third singular, governed request coordinated with gebidde).
% mdparagraph:P679:end

% teaching: NEW-U15-RECORD-WORULDBISGU-U15-SOURCE-DAT-PL
% teaching: NEW-U15-RECORD-WORULDBISGU-U15-SOURCE-SENSE
% teaching: U15-READ-LEX-WORULDBISGU
% support: U15-READ-LEX-WORULDBISGU
% mdparagraph:P680:begin
\textbf{woruldbisgum} — worldly occupations; cares (dative feminine plural after for).
% mdparagraph:P680:end

\reserve{5\baselineskip}
% mdparagraph:P681:begin
\textbf{Modern construction help.}
% mdparagraph:P681:end

% teaching: NEW-U15-RECORD-U15-SOURCE-P-TRANSLATION
% teaching: U15-READ-P-TRANSLATION
% support: U15-READ-P-TRANSLATION
% mdparagraph:P682:begin
Ælfred is subject. The feminine hīe/\allowbreak{}hēo refers to bōc. The source-language phrases use of plus dative; on Englisc gives the translation’s goal. Word be worde contrasts literal verbal correspondence with andgiet of andgiete, sense for sense. Gedōn describes the work’s resulting arrangement.
% mdparagraph:P682:end

% teaching: NEW-U15-RECORD-U15-SOURCE-P-HINDRANCES
% teaching: U15-READ-P-HINDRANCES
% support: U15-READ-P-HINDRANCES
% mdparagraph:P683:begin
The two superlative adverbs mean “most clearly and most intelligibly”. For with dative introduces the worldly cares that constrained how clearly Alfred could explain it. Ǣgðer ge … ge means “both … and”; on mōde/\allowbreak{}on līchaman identify mind and body. Sind means “are”. The relative after Ðā bisga refers to the cares, which are the subject of becōmon.
% mdparagraph:P683:end

% teaching: NEW-U15-RECORD-U15-SOURCE-P-ANTERIOR-ACTIONS
% teaching: U15-READ-P-ANTERIOR-ACTIONS
% support: U15-READ-P-ANTERIOR-ACTIONS
% mdparagraph:P684:begin
Hæfde underfangen = “had received/\allowbreak{}taken charge of.” The repeated hæfde governs both geleornod and gewęnd: “had studied ... and translated ...”. Ðēah here is “nevertheless”; ðā ... ðā frames “when ... then.” Tō Engliscum spelle and tō lēoðe mark resulting literary form. Hīe/\allowbreak{}hēo still refers to the book.
% mdparagraph:P684:end

% teaching: NEW-U15-RECORD-U15-SOURCE-P-APPEAL
% teaching: U15-READ-P-APPEAL
% support: U15-READ-P-APPEAL
% mdparagraph:P685:begin
Supply hē (Alfred) with bitt and hālsað, as Cook’s note 4 directs. Ǣlcne ðāra ðe ... means “each of those who ...”. Lyste is impersonal: “whom it may please to read this book”; its understood accusative experiencer is the prospective reader. The requested ðæt clause contains subjunctive gebidde, “that he pray for him.”
% mdparagraph:P685:end

% teaching: NEW-U15-RECORD-U15-SOURCE-P-READER-AUTHOR-REFERENCE
% teaching: U15-READ-P-READER-AUTHOR-REFERENCE
% support: U15-READ-P-READER-AUTHOR-REFERENCE
% mdparagraph:P686:begin
Ne wīte continues the request: “and not blame him.” In gif hē hit rihtlicor ongiete ðonne hē meahte, the first hē is the reader and the second is Alfred (Cook’s notes 6–7). Ongiete is subjunctive; wīte is the verb “blame,” not the noun “punishment.”
% mdparagraph:P686:end

% teaching: NEW-U15-RECORD-U15-SOURCE-P-CAPACITY
% teaching: U15-READ-P-CAPACITY
% support: U15-READ-P-CAPACITY
% mdparagraph:P687:begin
Be ... mǣðe and be ... ǣmettan mean according to the measure of one’s understanding and one’s available leisure. Sceal governs sprecan and dōn. Both ðæt ... and ðæt-ðæt ... function as “what/\allowbreak{}that which.” Preserve the source’s hyphenated spelling in the latter.
% mdparagraph:P687:end

\reserve{5\baselineskip}
% mdparagraph:P688:begin
\textbf{Cook’s text}
% mdparagraph:P688:end

% mdparagraph:P689:begin
X. ALFRED’S PREFACE TO BOETHIUS.
% mdparagraph:P689:end

% mdparagraph:P690:begin
(Prefixed to his translation.)
% mdparagraph:P690:end

% teaching: NEW-U15-READ-FB-X-ALFRED-PREFACE
\par\begin{minipage}{\linewidth}\raggedright
% passage: FB-X-ALFRED-PREFACE
% mdparagraph:P691:begin
\readingtext{Ælfred cyning wæs wealhstōd\textsuperscript{P1} ðisse bēc, and hīe of Bōclædene\textsuperscript{P2} on Englisc węnde, swā hēo nū is gedōn. Hwīlum hē sętte word be worde, hwīlum andgiet of andgiete, swā-swā hē hit ðā sweotolost and andgietfullicost geręccan meahte for ðǣm mislicum and manigfealdum woruldbisgum ðe hine oft ǣgðer ge on mōde ge on līchaman bisgodon. Ðā bisga\textsuperscript{P3} ūs sind swīðe earfoðrīme ðe on his dagum on ðā rīcu becōmon ðe hē underfangen hæfde, and ðēah, ðā hē ðās bōc hæfde geleornod, and of Lædene tō Engliscum spelle gewęnd, ðā geworhte hē hīe eft tō lēoðe, swā-swā hēo nū gedōn is. And nū bitt\textsuperscript{P4} and for Godes naman hālsað ǣlcne ðāra ðe ðās bōc rǣdan lyste,\textsuperscript{P5} ðæt hē for hine gebidde, and him ne wīte gif hē\textsuperscript{P6} hit rihtlicor ongiete ðonne hē\textsuperscript{P7} meahte; for-ðām-ðe ǣlc mann sceal be his andgietes mǣðe, and be his ǣmettan, sprecan ðæt hē spricð, and dōn ðæt-ðæt hē dēð.}
% mdparagraph:P691:end

\end{minipage}\par
\reserve{5\baselineskip}
% mdparagraph:P692:begin
\textbf{Original notes — Cook.}
% mdparagraph:P692:end

% mdparagraph:P693:begin
\textsuperscript{P1}: Wealh- signifies foreign (see walnut), and sometimes servant, orig. Celtic, Celt (cf. Wales, Welsh, Cornwall), from Volcæ, the name of a Celtic tribe (Cæsar, Gallic War, Bk. VII.).
% mdparagraph:P693:end

% mdparagraph:P694:begin
\textsuperscript{P2}: Perhaps originally in contrast to the Latin spoken in Britain.
% mdparagraph:P694:end

% mdparagraph:P695:begin
\textsuperscript{P3}: See 51. a.
% mdparagraph:P695:end

% mdparagraph:P696:begin
\textsuperscript{P4}: Supply hē.
% mdparagraph:P696:end

% mdparagraph:P697:begin
\textsuperscript{P5}: See 190.
% mdparagraph:P697:end

% mdparagraph:P698:begin
\textsuperscript{P6}: The reader.
% mdparagraph:P698:end

% mdparagraph:P699:begin
\textsuperscript{P7}: Alfred.
% mdparagraph:P699:end

% mdparagraph:P700:begin
\textbf{Tasks P01–P05 — new teaching material.} Keep the modern vocabulary, construction helps and original notes available for every task.
% mdparagraph:P700:end

\reserve{5\baselineskip}
\par\begin{minipage}{\linewidth}\raggedright
% mdparagraph:P701:begin
\textbf{P01}
% mdparagraph:P701:end

% item: NEW-U15-P01
% mdparagraph:P702:begin
State the two translation methods described, the obstacles Alfred faced, the later change he made to the translation, and his request to readers.
% mdparagraph:P702:end

% output: NEW-U15-P01-O001
% output: NEW-U15-P01-O002
% output: NEW-U15-P01-O003
% output: NEW-U15-P01-O004
% output: NEW-U15-P01-O005
% mdparagraph:P703:begin
\textbf{Responses, in order:} (1) first translation method; (2) second translation method; (3) obstacles; (4) later change; (5) request to readers.
% mdparagraph:P703:end

\end{minipage}\par\addvspace{4pt}

\par\begin{minipage}{\linewidth}\raggedright
% mdparagraph:P704:begin
\textbf{P02}
% mdparagraph:P704:end

% item: NEW-U15-P02
% mdparagraph:P705:begin
In the appeal from ðæt hē for hine gebidde through ðonne hē meahte, identify the referent of each of these five tokens in order: hē, hine, him, hē after gif, hē after ðonne. Explain the resulting contrast between the two comparison subjects.
% mdparagraph:P705:end

% output: NEW-U15-P02-O001
% output: NEW-U15-P02-O002
% output: NEW-U15-P02-O003
% output: NEW-U15-P02-O004
% output: NEW-U15-P02-O005
% output: NEW-U15-P02-O006
% mdparagraph:P706:begin
\textbf{Responses, in order:} (1) hē before gebidde; (2) hine after for; (3) him before ne wīte; (4) hē after gif; (5) hē after ðonne; (6) contrast between the comparison subjects.
% mdparagraph:P706:end

\end{minipage}\par\addvspace{4pt}

\par\begin{minipage}{\linewidth}\raggedright
% mdparagraph:P707:begin
\textbf{P03}
% mdparagraph:P707:end

% item: NEW-U15-P03
% mdparagraph:P708:begin
Explain ǣgðer ge on mōde ge on līchaman. Then trace the subject and finite verb of the sentence beginning Ðā bisga and identify the two relative clauses that follow.
% mdparagraph:P708:end

% output: NEW-U15-P03-O001
% output: NEW-U15-P03-O002
% output: NEW-U15-P03-O003
% output: NEW-U15-P03-O004
% mdparagraph:P709:begin
\textbf{Responses, in order:} (1) meaning and coordination in the first expression; (2) main subject and finite verb; (3) first relative clause; (4) second relative clause.
% mdparagraph:P709:end

\end{minipage}\par\addvspace{4pt}

\par\begin{minipage}{\linewidth}\raggedright
% mdparagraph:P710:begin
\textbf{P04}
% mdparagraph:P710:end

% item: NEW-U15-P04
% mdparagraph:P711:begin
Explain what is being compared by rihtlicor … ðonne and what the second clause leaves understood. State the function of ne in and him ne wīte and the relationship of that clause to the preceding request.
% mdparagraph:P711:end

% output: NEW-U15-P04-O001
% output: NEW-U15-P04-O002
% mdparagraph:P712:begin
\textbf{Responses, in order:} (1) comparison and understood words; (2) function of ne and relation of its clause to the request.
% mdparagraph:P712:end

\end{minipage}\par\addvspace{4pt}

\par\begin{minipage}{\linewidth}\raggedright
% mdparagraph:P713:begin
\textbf{P05}
% mdparagraph:P713:end

% item: NEW-U15-P05
% mdparagraph:P714:begin
Translate Alfred’s complete preface. Preserve the distinction between his account of the translation and his request for a charitable reading.
% mdparagraph:P714:end

% output: NEW-U15-P05-O001
\end{minipage}\par\addvspace{4pt}

\reserve{7\baselineskip}
% mdparagraph:P715:begin
\subsection*{R. A prayer of King Alfred}\BookHeading{nav-ch15-learner-0043}{subsubsection}{R. A prayer of King Alfred}{0}
% mdparagraph:P715:end

% mdparagraph:P716:begin
Cook, First Book, printed p. 163 (PDF p. 179). The complete prayer.
% mdparagraph:P716:end

% mdparagraph:P717:begin
\textbf{Modern vocabulary — supplied for this reading.} Forms are glossed in their local senses; the parses identify the forms encountered, not complete new paradigms.
% mdparagraph:P717:end

% teaching: NEW-U15-RECORD-A-EVER-U15-SOURCE-INVARIANT
% teaching: NEW-U15-RECORD-A-EVER-U15-SOURCE-SENSE
% teaching: U15-READ-LEX-A-EVER
% support: U15-READ-LEX-A-EVER
% mdparagraph:P718:begin
\textbf{ā} — ever; always (adverb repeated for intensification).
% mdparagraph:P718:end

% teaching: NEW-U15-RECORD-AEGHWILC-U15-SOURCE-DAT-M-SG
% teaching: NEW-U15-RECORD-AEGHWILC-U15-SOURCE-SENSE
% teaching: U15-READ-LEX-AEGHWILC
% support: U15-READ-LEX-AEGHWILC
% mdparagraph:P719:begin
\textbf{ǣghwilcum} — each; every (strong dative masculine singular).
% mdparagraph:P719:end

% teaching: NEW-U15-RECORD-AELMIHTIG-U15-SOURCE-SENSE
% teaching: NEW-U15-RECORD-AELMIHTIG-U15-SOURCE-VOC-M-SG-WEAK
% teaching: U15-READ-LEX-AELMIHTIG
% support: U15-READ-LEX-AELMIHTIG
% mdparagraph:P720:begin
\textbf{ælmihtiga} — almighty (weak masculine singular in direct address).
% mdparagraph:P720:end

% teaching: NEW-U15-RECORD-AFIERRAN-U15-SOURCE-IMP-2-SG
% teaching: NEW-U15-RECORD-AFIERRAN-U15-SOURCE-SENSE
% teaching: U15-READ-LEX-AFIERRAN
% support: U15-READ-LEX-AFIERRAN
% mdparagraph:P721:begin
\textbf{āfierr} — remove; take away (imperative second singular).
% mdparagraph:P721:end

% teaching: NEW-U15-RECORD-AMEN-U15-SOURCE-INVARIANT
% teaching: NEW-U15-RECORD-AMEN-U15-SOURCE-SENSE
% teaching: U15-READ-LEX-AMEN
% support: U15-READ-LEX-AMEN
% mdparagraph:P722:begin
\textbf{Amen} — Amen (liturgical closing word).
% mdparagraph:P722:end

% teaching: NEW-U15-RECORD-AWYRCEAN-U15-SOURCE-PRET-1-SG
% teaching: NEW-U15-RECORD-AWYRCEAN-U15-SOURCE-SENSE
% teaching: U15-READ-LEX-AWYRCEAN
% support: U15-READ-LEX-AWYRCEAN
% mdparagraph:P723:begin
\textbf{āworhte} — do; perform, here deserve by one's conduct (preterite indicative first singular with tō ðē).
% mdparagraph:P723:end

% teaching: NEW-U15-RECORD-CLAENE-U15-SOURCE-DAT-M-SG-STRONG
% teaching: NEW-U15-RECORD-CLAENE-U15-SOURCE-DAT-N-SG-STRONG
% teaching: NEW-U15-RECORD-CLAENE-U15-SOURCE-SENSE
% teaching: U15-READ-LEX-CLAENE
% support: U15-READ-LEX-CLAENE
% mdparagraph:P724:begin
\textbf{clǣnum} — pure; clean (strong dative neuter singular). \textbf{clǣnum} — pure; clean (strong dative masculine singular).
% mdparagraph:P724:end

% teaching: NEW-U15-RECORD-EARNUNG-U15-SOURCE-DAT-PL
% teaching: NEW-U15-RECORD-EARNUNG-U15-SOURCE-SENSE
% teaching: U15-READ-LEX-EARNUNG
% support: U15-READ-LEX-EARNUNG
% mdparagraph:P725:begin
\textbf{earnungum} — merits; good services (dative feminine plural coordinated under for).
% mdparagraph:P725:end

% teaching: NEW-U15-RECORD-ENDE-U15-SOURCE-DAT-SG
% teaching: NEW-U15-RECORD-ENDE-U15-SOURCE-SENSE
% teaching: U15-READ-LEX-ENDE
% support: U15-READ-LEX-ENDE
% mdparagraph:P726:begin
\textbf{ęnde} — end (dative masculine singular after būtan).
% mdparagraph:P726:end

% teaching: NEW-U15-RECORD-FOR-U15-SOURCE-APPEAL
% teaching: U15-READ-LEX-FOR
% support: U15-READ-LEX-FOR
% mdparagraph:P727:begin
\textbf{for} — by; for the sake of, naming grounds of appeal in the prayer.
% mdparagraph:P727:end

% teaching: NEW-U15-RECORD-FUL-U15-SOURCE-ACC-F-SG-WEAK
% teaching: NEW-U15-RECORD-FUL-U15-SOURCE-SENSE
% teaching: U15-READ-LEX-FUL
% support: U15-READ-LEX-FUL
% mdparagraph:P728:begin
\textbf{fūlan} — foul (weak accusative feminine singular).
% mdparagraph:P728:end

% teaching: NEW-U15-RECORD-FULTUM-U15-SOURCE-NOM-SG
% teaching: NEW-U15-RECORD-FULTUM-U15-SOURCE-SENSE
% teaching: U15-READ-LEX-FULTUM
% support: U15-READ-LEX-FULTUM
% mdparagraph:P729:begin
\textbf{Fultum} — help; support (predicate nominative masculine singular, address to God).
% mdparagraph:P729:end

% teaching: NEW-U15-RECORD-GALNES-U15-SOURCE-ACC-SG
% teaching: NEW-U15-RECORD-GALNES-U15-SOURCE-SENSE
% teaching: U15-READ-LEX-GALNES
% support: U15-READ-LEX-GALNES
% mdparagraph:P730:begin
\textbf{gālnesse} — lust (accusative feminine singular).
% mdparagraph:P730:end

% teaching: NEW-U15-RECORD-GESCIELDAN-U15-SOURCE-IMP-2-SG
% teaching: NEW-U15-RECORD-GESCIELDAN-U15-SOURCE-SENSE
% teaching: U15-READ-LEX-GESCIELDAN
% support: U15-READ-LEX-GESCIELDAN
% mdparagraph:P731:begin
\textbf{gescield} — shield; protect (imperative second singular).
% mdparagraph:P731:end

% teaching: NEW-U15-RECORD-GESEWENLIC-U15-SOURCE-DAT-M-PL-STRONG
% teaching: NEW-U15-RECORD-GESEWENLIC-U15-SOURCE-SENSE
% teaching: U15-READ-LEX-GESEWENLIC
% support: U15-READ-LEX-GESEWENLIC
% mdparagraph:P732:begin
\textbf{gesewenlicum} — visible (strong dative masculine plural).
% mdparagraph:P732:end

% teaching: NEW-U15-RECORD-GESTATHELIAN-U15-SOURCE-IMP-2-SG
% teaching: NEW-U15-RECORD-GESTATHELIAN-U15-SOURCE-SENSE
% teaching: U15-READ-LEX-GESTATHELIAN
% support: U15-READ-LEX-GESTATHELIAN
% mdparagraph:P733:begin
\textbf{gestaðela} — establish; make steadfast (imperative second singular).
% mdparagraph:P733:end

% teaching: NEW-U15-RECORD-GESTRANGIAN-U15-SOURCE-IMP-2-SG
% teaching: NEW-U15-RECORD-GESTRANGIAN-U15-SOURCE-SENSE
% teaching: U15-READ-LEX-GESTRANGIAN
% support: U15-READ-LEX-GESTRANGIAN
% mdparagraph:P734:begin
\textbf{gestranga} — strengthen (imperative second singular).
% mdparagraph:P734:end

% teaching: NEW-U15-RECORD-GETHANC-U15-SOURCE-DAT-SG
% teaching: NEW-U15-RECORD-GETHANC-U15-SOURCE-SENSE
% teaching: U15-READ-LEX-GETHANC
% support: U15-READ-LEX-GETHANC
% mdparagraph:P735:begin
\textbf{geðance} — thought; mind (dative neuter singular after mid).
% mdparagraph:P735:end

% teaching: NEW-U15-RECORD-GEWISSIAN-U15-SOURCE-IMP-2-SG
% teaching: NEW-U15-RECORD-GEWISSIAN-U15-SOURCE-PRES-SUBJ-2-SG
% teaching: NEW-U15-RECORD-GEWISSIAN-U15-SOURCE-SENSE
% teaching: U15-READ-LEX-GEWISSIAN
% support: U15-READ-LEX-GEWISSIAN
% mdparagraph:P736:begin
\textbf{gewissie} — guide; direct (present subjunctive second singular in requested ðæt clause). \textbf{gewissa} — guide; direct (imperative second singular).
% mdparagraph:P736:end

% teaching: NEW-U15-RECORD-HALIG-U15-SOURCE-GEN-PL-SUBSTANTIVAL
% teaching: NEW-U15-RECORD-HALIG-U15-SOURCE-SENSE
% teaching: U15-READ-LEX-HALIG
% support: U15-READ-LEX-HALIG
% mdparagraph:P737:begin
\textbf{hālgena} — saints; holy ones (substantival genitive plural in -ena).
% mdparagraph:P737:end

% teaching: NEW-U15-RECORD-HE-U15-SOURCE-HIERA-GEN-PL
% teaching: NEW-U15-RECORD-HE-U15-SOURCE-SENSE
% teaching: U15-READ-LEX-HE
% support: U15-READ-LEX-HE
% mdparagraph:P738:begin
\textbf{hīera} — their (third-person genitive plural).
% mdparagraph:P738:end

% teaching: NEW-U15-RECORD-INWEARDLICE-U15-SOURCE-INVARIANT
% teaching: NEW-U15-RECORD-INWEARDLICE-U15-SOURCE-SENSE
% teaching: U15-READ-LEX-INWEARDLICE
% support: U15-READ-LEX-INWEARDLICE
% mdparagraph:P739:begin
\textbf{inweardlīce} — fervently; wholeheartedly.
% mdparagraph:P739:end

% teaching: NEW-U15-RECORD-LOF-U15-SOURCE-NOM-SG
% teaching: NEW-U15-RECORD-LOF-U15-SOURCE-SENSE
% teaching: U15-READ-LEX-LOF
% support: U15-READ-LEX-LOF
% mdparagraph:P740:begin
\textbf{lof} — praise (nominative masculine singular subject).
% mdparagraph:P740:end

% teaching: NEW-U15-RECORD-LUFE-U15-SOURCE-DAT-SG
% teaching: NEW-U15-RECORD-LUFE-U15-SOURCE-SENSE
% teaching: U15-READ-LEX-LUFE
% support: U15-READ-LEX-LUFE
% mdparagraph:P741:begin
\textbf{lufan} — love (weak feminine dative singular after for).
% mdparagraph:P741:end

% teaching: NEW-U15-RECORD-MAEGTHHAD-U15-SOURCE-DAT-SG
% teaching: NEW-U15-RECORD-MAEGTHHAD-U15-SOURCE-SENSE
% teaching: U15-READ-LEX-MAEGTHHAD
% support: U15-READ-LEX-MAEGTHHAD
% mdparagraph:P742:begin
\textbf{mægðhāde} — maidenhood; virginity (dative masculine singular after for).
% mdparagraph:P742:end

% teaching: NEW-U15-RECORD-MARIA-NAME-U15-SOURCE-GEN-CONTEXT
% teaching: NEW-U15-RECORD-MARIA-NAME-U15-SOURCE-SENSE
% teaching: U15-READ-LEX-MARIA-NAME
% support: U15-READ-LEX-MARIA-NAME
% mdparagraph:P743:begin
\textbf{Marian} — Mary (genitive form of a foreign name).
% mdparagraph:P743:end

% teaching: NEW-U15-RECORD-MICHAEL-NAME-U15-SOURCE-GEN-CONTEXT
% teaching: NEW-U15-RECORD-MICHAEL-NAME-U15-SOURCE-SENSE
% teaching: U15-READ-LEX-MICHAEL-NAME
% support: U15-READ-LEX-MICHAEL-NAME
% mdparagraph:P744:begin
\textbf{Michaeles} — Michael (genitive personal name, vernacular -es with Latin Sancti).
% mdparagraph:P744:end

% teaching: NEW-U13-VOC-MID
% teaching: U15-READ-LEX-MID
% support: U15-READ-LEX-MID
% mdparagraph:P745:begin
\textbf{mid} — with; in the manner of the accompanying quality or mental disposition.
% mdparagraph:P745:end

% teaching: NEW-U13-SOURCE-HELPS
% teaching: U15-READ-LEX-MOD
% support: U15-READ-LEX-MOD
% mdparagraph:P746:begin
\textbf{mōd} — mind; attention (accusative neuter singular).
% mdparagraph:P746:end

% teaching: NEW-U15-RECORD-SANCTUS-LATIN-U15-SOURCE-GEN-F-SG
% teaching: NEW-U15-RECORD-SANCTUS-LATIN-U15-SOURCE-GEN-M-SG
% teaching: NEW-U15-RECORD-SANCTUS-LATIN-U15-SOURCE-SENSE
% teaching: U15-READ-LEX-SANCTUS-LATIN
% support: U15-READ-LEX-SANCTUS-LATIN
% mdparagraph:P747:begin
\textbf{Sanctæ} — Saint; holy (Latin genitive feminine singular accompanying Marian). \textbf{Sancti} — Saint; holy (Latin genitive masculine singular).
% mdparagraph:P747:end

% teaching: NEW-U15-RECORD-SCIEPPEND-U15-SOURCE-NOM-SG
% teaching: NEW-U15-RECORD-SCIEPPEND-U15-SOURCE-SENSE
% teaching: U15-READ-LEX-SCIEPPEND
% support: U15-READ-LEX-SCIEPPEND
% mdparagraph:P748:begin
\textbf{Scieppend} — Creator (predicate nominative masculine singular title).
% mdparagraph:P748:end

% teaching: NEW-U15-RECORD-TAECAN-U15-SOURCE-IMP-2-SG
% teaching: NEW-U15-RECORD-TAECAN-U15-SOURCE-SENSE
% teaching: U15-READ-LEX-TAECAN
% support: U15-READ-LEX-TAECAN
% mdparagraph:P749:begin
\textbf{tǣc} — teach; show (imperative second singular).
% mdparagraph:P749:end

% teaching: NEW-U15-RECORD-THEARF-U15-SOURCE-DAT-SG
% teaching: NEW-U15-RECORD-THEARF-U15-SOURCE-SENSE
% teaching: U15-READ-LEX-THEARF
% support: U15-READ-LEX-THEARF
% mdparagraph:P750:begin
\textbf{ðearfe} — benefit; good; need (dative feminine singular after tō, with dependent sāwle).
% mdparagraph:P750:end

% teaching: NEW-U15-RECORD-TO-U15-SOURCE-BENEFIT
% teaching: NEW-U15-RECORD-TO-U15-SOURCE-CONFORMITY
% teaching: NEW-U15-RECORD-TO-U15-SOURCE-DIRECTED-CONDUCT
% teaching: NEW-U15-RECORD-TO-U15-SOURCE-ETERNITY
% teaching: U15-READ-LEX-TO
% support: U15-READ-LEX-TO
% mdparagraph:P751:begin
\textbf{tō} — towards; in relation to the person to whom conduct is directed. \textbf{tō} — to; in accordance with, a will as the goal of guidance. \textbf{tō} — for; to the benefit of. \textbf{tō} — for ever, in fixed tō worulde.
% mdparagraph:P751:end

% teaching: NEW-U15-RECORD-TOFORAN-U15-SOURCE-INVARIANT
% teaching: NEW-U15-RECORD-TOFORAN-U15-SOURCE-SENSE
% teaching: U15-READ-LEX-TOFORAN
% support: U15-READ-LEX-TOFORAN
% mdparagraph:P752:begin
\textbf{tōforan} — before; above all (preposition with dative of comparison or preference).
% mdparagraph:P752:end

% teaching: NEW-U15-RECORD-TOHOPA-U15-SOURCE-NOM-SG
% teaching: NEW-U15-RECORD-TOHOPA-U15-SOURCE-SENSE
% teaching: U15-READ-LEX-TOHOPA
% support: U15-READ-LEX-TOHOPA
% mdparagraph:P753:begin
\textbf{Tōhopa} — hope (predicate nominative masculine singular).
% mdparagraph:P753:end

% teaching: NEW-U15-RECORD-TREOWNES-U15-SOURCE-NOM-SG
% teaching: NEW-U15-RECORD-TREOWNES-U15-SOURCE-SENSE
% teaching: U15-READ-LEX-TREOWNES
% support: U15-READ-LEX-TREOWNES
% mdparagraph:P754:begin
\textbf{Trēownes} — trust; confidence (predicate nominative feminine singular).
% mdparagraph:P754:end

% teaching: NEW-U15-RECORD-UNGESEWENLIC-U15-SOURCE-DAT-M-PL-STRONG
% teaching: NEW-U15-RECORD-UNGESEWENLIC-U15-SOURCE-SENSE
% teaching: U15-READ-LEX-UNGESEWENLIC
% support: U15-READ-LEX-UNGESEWENLIC
% mdparagraph:P755:begin
\textbf{ungesewenlicum} — invisible (strong dative masculine plural).
% mdparagraph:P755:end

% teaching: NEW-U15-RECORD-WEALDEND-U15-SOURCE-SENSE
% teaching: NEW-U15-RECORD-WEALDEND-U15-SOURCE-VOC-SG
% teaching: U15-READ-LEX-WEALDEND
% support: U15-READ-LEX-WEALDEND
% mdparagraph:P756:begin
\textbf{Wealdend} — Ruler (nominative/\allowbreak{}vocative masculine singular title).
% mdparagraph:P756:end

% teaching: NEW-U15-RECORD-WITHERWINNA-U15-SOURCE-DAT-PL
% teaching: NEW-U15-RECORD-WITHERWINNA-U15-SOURCE-SENSE
% teaching: U15-READ-LEX-WITHERWINNA
% support: U15-READ-LEX-WITHERWINNA
% mdparagraph:P757:begin
\textbf{wiðerwinnum} — adversaries (dative masculine plural after wið).
% mdparagraph:P757:end

% teaching: NEW-U15-RECORD-WYRHTA-U15-SOURCE-MAKER
% teaching: U15-READ-LEX-WYRHTA
% support: U15-READ-LEX-WYRHTA
% mdparagraph:P758:begin
\textbf{wyrhta} — Maker; Creator (not the previously taught workman sense) (nominative singular masculine).
% mdparagraph:P758:end

\reserve{5\baselineskip}
% mdparagraph:P759:begin
\textbf{Modern construction help.}
% mdparagraph:P759:end

% teaching: NEW-U15-RECORD-U15-SOURCE-R-ADDRESS
% teaching: U15-READ-R-ADDRESS
% support: U15-READ-R-ADDRESS
% mdparagraph:P760:begin
Dryhten and God are direct address; ælmihtiga is weak singular. Wyrhta and Wealdend are “Maker and Ruler,” with genitive ealra gesceafta. This Maker sense is distinct from the previously taught workman sense.
% mdparagraph:P760:end

% teaching: NEW-U15-RECORD-U15-SOURCE-R-GROUNDS-OF-APPEAL
% teaching: U15-READ-R-GROUNDS-OF-APPEAL
% support: U15-READ-R-GROUNDS-OF-APPEAL
% mdparagraph:P761:begin
The for phrases name grounds of appeal. In ðǣre hālgan rōde tācne, the genitive phrase “of the holy cross” depends on dative tācne. Sanctæ Marian and Sancti Michaeles are supplied names in genitive function with mixed foreign-name inflection; do not derive a new productive Latin paradigm. Gehīersumnesse is obedience. Hālgena means saints; lufan is the weak noun’s dative; earnungum means merits.
% mdparagraph:P761:end

% teaching: NEW-U15-RECORD-U15-SOURCE-R-GUIDANCE
% teaching: U15-READ-R-GUIDANCE
% support: U15-READ-R-GUIDANCE
% mdparagraph:P762:begin
The requested ðæt clause has gewissie, “that you guide.” Bęt ðonne ic āworhte tō ðē means better than my conduct towards you has deserved. The following gewissa is a direct singular command. In mīnre sāwle ðearfe, mīnre sāwle is genitive, dependent on dative ðearfe: “for my soul’s benefit.” Cunne is first-singular subjunctive, “than I myself may know.”
% mdparagraph:P762:end

% teaching: NEW-U15-RECORD-U15-SOURCE-R-COMMANDS
% teaching: U15-READ-R-COMMANDS
% support: U15-READ-R-COMMANDS
% mdparagraph:P763:begin
Gestaðela, gestranga, āfierr and gescield are addressed singular commands. Establish my mind according to your will; strengthen me against the devil’s temptations; remove lust and unrighteousness; protect me against visible and invisible adversaries. The latter adjectives agree with dative plural wiðerwinnum.
% mdparagraph:P763:end

% teaching: NEW-U15-RECORD-U15-SOURCE-R-TEACHING-PURPOSE
% teaching: U15-READ-R-TEACHING-PURPOSE
% support: U15-READ-R-TEACHING-PURPOSE
% mdparagraph:P764:begin
Tǣc takes dative mē. Its following ðīnne willan is the accusative object of wyrceanne (“to do your will”), as Cook’s note 7 specifies. The ðæt clause gives purpose with mæge. Tōforan eallum ðingum means above all things; the two mid phrases give manner: with pure thought and a pure body.
% mdparagraph:P764:end

% teaching: NEW-U15-RECORD-U15-SOURCE-R-TITLES-AND-DOXOLOGY
% teaching: U15-READ-R-TITLES-AND-DOXOLOGY
% support: U15-READ-R-TITLES-AND-DOXOLOGY
% mdparagraph:P765:begin
The final causal clause names God as Creator, Redeemer, help, comfort, trust and hope. Sīe is an optative wish, “may there be”; ðē is dative. Ā is repeated three times, followed by tō worulde būtan ǣghwilcum ęnde, “for ever without any end.”
% mdparagraph:P765:end

\reserve{5\baselineskip}
\par\begin{minipage}{\linewidth}\raggedright
% mdparagraph:P766:begin
\textbf{Cook’s text}
% mdparagraph:P766:end

% mdparagraph:P767:begin
XI. A PRAYER OF KING ALFRED.
% mdparagraph:P767:end

% mdparagraph:P768:begin
(From the end of his translation of Boethius.)
% mdparagraph:P768:end

% teaching: NEW-U15-READ-FB-XI-ALFRED-PRAYER

% passage: FB-XI-ALFRED-PRAYER
% mdparagraph:P769:begin
\readingtext{Dryhten,\textsuperscript{R1} ælmihtiga God,\textsuperscript{R1} Wyrhta and Wealdend ealra gesceafta, ic bidde ðē for ðīnre miclan mildheortnesse, and for ðǣre hālgan rōde tācne,\textsuperscript{R2} and for Sanctæ Marian mægðhāde, and for Sancti Michaeles gehīersumnesse, and for ealra ðīnra hālgena\textsuperscript{R3} lufan and hīera earnungum, ðæt ðū mē gewissie\textsuperscript{R4} bęt ðonne ic āworhte tō ðē; and gewissa mē tō ðīnum willan, and tō mīnre sāwle ðearfe,\textsuperscript{R5} bęt ðonne ic self cunne\textsuperscript{R6}; and gestaðela mīn mōd tō ðīnum willan and tō mīnre sāwle ðearfe; and gestranga mē wið ðæs dēofles\textsuperscript{R3} costnungum; and āfierr fram mē ðā fūlan gālnesse and ǣlce unrihtwīsnesse; and gescield mē wið mīnum wiðerwinnum, gesewenlicum and ungesewenlicum; and tǣc mē ðīnne willan\textsuperscript{R7} tō wyrceanne; ðæt ic mæge\textsuperscript{R8} ðē inweardlīce lufian tōforan eallum ðingum, mid clǣnum geðance and mid clǣnum līchaman. For-ðon-ðe ðū eart mīn Scieppend,\textsuperscript{R9} and mīn Ālīesend, mīn Fultum, mīn Frōfor, mīn Trēownes, and mīn Tōhopa. Sīe ðē lof and wuldor nū and ā ā ā, tō worulde būtan ǣghwilcum ęnde. Amen.}
% mdparagraph:P769:end

\end{minipage}\par
\reserve{5\baselineskip}
% mdparagraph:P770:begin
\textbf{Original notes — Cook.}
% mdparagraph:P770:end

% mdparagraph:P771:begin
\textsuperscript{R1}: See 152.
% mdparagraph:P771:end

% mdparagraph:P772:begin
\textsuperscript{R2}: Governed by for.
% mdparagraph:P772:end

% mdparagraph:P773:begin
\textsuperscript{R3}: See 153. c.
% mdparagraph:P773:end

% mdparagraph:P774:begin
\textsuperscript{R4}: See 194. b.
% mdparagraph:P774:end

% mdparagraph:P775:begin
\textsuperscript{R5}: See 166.
% mdparagraph:P775:end

% mdparagraph:P776:begin
\textsuperscript{R6}: Optative (130).
% mdparagraph:P776:end

% mdparagraph:P777:begin
\textsuperscript{R7}: Object of wyrceanne.
% mdparagraph:P777:end

% mdparagraph:P778:begin
\textsuperscript{R8}: See 196. d.
% mdparagraph:P778:end

% mdparagraph:P779:begin
\textsuperscript{R9}: See 150.
% mdparagraph:P779:end

\par\begin{minipage}{\linewidth}\raggedright
% mdparagraph:P780:begin
\textbf{Tasks R01–R05 — new teaching material.} Keep the modern vocabulary, construction helps and original notes available for every task.
% mdparagraph:P780:end


% mdparagraph:P781:begin
\textbf{R01}
% mdparagraph:P781:end

% item: NEW-U15-R01
% mdparagraph:P782:begin
Identify the opening requested action after ic bidde and each of the six subsequent direct commands. For each give the actual finite form, its mood and its essential meaning.
% mdparagraph:P782:end

% output: NEW-U15-R01-O001
% output: NEW-U15-R01-O002
% output: NEW-U15-R01-O003
% output: NEW-U15-R01-O004
% output: NEW-U15-R01-O005
% output: NEW-U15-R01-O006
% output: NEW-U15-R01-O007
% output: NEW-U15-R01-O008
% output: NEW-U15-R01-O009
% output: NEW-U15-R01-O010
% output: NEW-U15-R01-O011
% output: NEW-U15-R01-O012
% output: NEW-U15-R01-O013
% output: NEW-U15-R01-O014
% output: NEW-U15-R01-O015
% output: NEW-U15-R01-O016
% output: NEW-U15-R01-O017
% output: NEW-U15-R01-O018
% output: NEW-U15-R01-O019
% output: NEW-U15-R01-O020
% output: NEW-U15-R01-O021
% mdtable:T83:begin
\par\begingroup
\fontsize{10.5}{13.2}\selectfont
\setlength{\LTleft}{0pt}\setlength{\LTright}{\fill}
\noindent
\begin{tabular}{@{}>{\raggedright\arraybackslash}p{39.78mm}>{\raggedright\arraybackslash}p{25.26mm}>{\raggedright\arraybackslash}p{25.26mm}>{\raggedright\arraybackslash}p{25.26mm}@{}}
\toprule
\textbf{Requested action} & \textbf{Finite form} & \textbf{Mood} & \textbf{Meaning} \\
\midrule
Opening requested action & (1) \_\_\_\_ & (2) \_\_\_\_ & (3) \_\_\_\_ \\
First direct command & (4) \_\_\_\_ & (5) \_\_\_\_ & (6) \_\_\_\_ \\
Second direct command & (7) \_\_\_\_ & (8) \_\_\_\_ & (9) \_\_\_\_ \\
Third direct command & (10) \_\_\_\_ & (11) \_\_\_\_ & (12) \_\_\_\_ \\
Fourth direct command & (13) \_\_\_\_ & (14) \_\_\_\_ & (15) \_\_\_\_ \\
Fifth direct command & (16) \_\_\_\_ & (17) \_\_\_\_ & (18) \_\_\_\_ \\
Sixth direct command & (19) \_\_\_\_ & (20) \_\_\_\_ & (21) \_\_\_\_ \\
\bottomrule
\end{tabular}
\par\endgroup

% mdtable:T83:end
\end{minipage}\par\addvspace{4pt}

\par\begin{minipage}{\linewidth}\raggedright
% mdparagraph:P783:begin
\textbf{R02}
% mdparagraph:P783:end

% item: NEW-U15-R02
% mdparagraph:P784:begin
Explain the grounds of entreaty introduced by the repeated for phrases. In for ðǣre hālgan rōde tācne, distinguish the noun governed by for from its dependent genitive.
% mdparagraph:P784:end

% output: NEW-U15-R02-O001
% output: NEW-U15-R02-O002
% mdparagraph:P785:begin
\textbf{Responses, in order:} (1) meaning of the for phrases; (2) case relations in ðǣre hālgan rōde tācne.
% mdparagraph:P785:end

\end{minipage}\par\addvspace{4pt}

\par\begin{minipage}{\linewidth}\raggedright
% mdparagraph:P786:begin
\textbf{R03}
% mdparagraph:P786:end

% item: NEW-U15-R03
% mdparagraph:P787:begin
Compare the two bęt ðonne clauses. State the standard in each and explain why the second one uses cunne in this petition.
% mdparagraph:P787:end

% output: NEW-U15-R03-O001
% output: NEW-U15-R03-O002
% output: NEW-U15-R03-O003
% mdparagraph:P788:begin
\textbf{Responses, in order:} (1) first comparison standard; (2) second comparison standard; (3) explanation of cunne.
% mdparagraph:P788:end

\end{minipage}\par\addvspace{4pt}

\par\begin{minipage}{\linewidth}\raggedright
% mdparagraph:P789:begin
\textbf{R04}
% mdparagraph:P789:end

% item: NEW-U15-R04
% mdparagraph:P790:begin
Distinguish the purpose clause ðæt ic mæge … lufian from the concluding Sīe ðē lof and wuldor. Identify the actor/\allowbreak{}addressee, the mood and the function of each.
% mdparagraph:P790:end

% output: NEW-U15-R04-O001
% output: NEW-U15-R04-O002
% output: NEW-U15-R04-O003
% output: NEW-U15-R04-O004
% output: NEW-U15-R04-O005
% output: NEW-U15-R04-O006
% mdparagraph:P791:begin
\textbf{Responses, in order:} (1) first clause: actor/\allowbreak{}addressee; (2) first clause: mood; (3) first clause: function; (4) second clause: actor/\allowbreak{}addressee; (5) second clause: mood; (6) second clause: function.
% mdparagraph:P791:end

\end{minipage}\par\addvspace{4pt}

\par\begin{minipage}{\linewidth}\raggedright
% mdparagraph:P792:begin
\textbf{R05}
% mdparagraph:P792:end

% item: NEW-U15-R05
% mdparagraph:P793:begin
Translate the complete prayer, including the grounds of entreaty, the commands, the purpose and the concluding praise.
% mdparagraph:P793:end

% output: NEW-U15-R05-O001
\end{minipage}\par\addvspace{4pt}

\reserve{8\baselineskip}
% mdparagraph:P794:begin
\section*{9. Sources and selected forms}\BookHeading{nav-ch15-learner-0044}{subsection}{9. Sources and selected forms}{1}
% mdparagraph:P794:end

\begingroup\fontsize{10}{12.8}\selectfont
\setlength{\parskip}{1.5pt}
% mdparagraph:P795:begin
\textsuperscript{1}: Cook, A First Book in Old English, 2nd ed., pp. 28–37, especially §§43–53. Tables give selected complete models; stated alternatives remain valid.
% mdparagraph:P795:end

% mdparagraph:P796:begin
\textsuperscript{2}: Wright, Old English Grammar, p. 184, §§382–383, supplies the complete strengu class and its variants. Bosworth–Toller, p. 500b, treats hǣlu as indeclinable. These facts support the selected singular tables without assigning plural count uses.
% mdparagraph:P796:end

% mdparagraph:P797:begin
\textsuperscript{3}: Cook, First Book, pp. 23, 74, 78–79 and 100. The selected negative indicative forms combine the positive systems with the stated contractions; the multiword entries remain ordinary ne plus a verb.
% mdparagraph:P797:end

% mdparagraph:P798:begin
\textsuperscript{4}: Cook, First Book, pp. 52–53; Wright, Old English Grammar, pp. 239–240. Further indefinite and free-choice families are acknowledged only within the limits stated above.
% mdparagraph:P798:end

% mdparagraph:P799:begin
\textsuperscript{5}: Cook, First Book, pp. 54–74 and the cited lexical entries. The class patterns determine endings and stem distribution; each newly required system is printed in full above.
% mdparagraph:P799:end

% mdparagraph:P800:begin
\textsuperscript{6}: Cook, First Book, pp. 39–41, §§57–59. The selected -ig syncopation and double-consonant spellings do not exclude the stated historical alternatives.
% mdparagraph:P800:end

% mdparagraph:P801:begin
\textsuperscript{7}: Cook, Exercises in Old English (1895), IX.28–30 and X, XI.31, printed pp. 15–18 (PDF pp. 23–26). The preserved source prompts are distinct from the modern directions and helps.
% mdparagraph:P801:end

\endgroup

% book-source-body:end
\par\label{end-ch15-learner}
\clearpage\endgroup

\BookRecto
\begingroup
\renewcommand{\runningtitle}{Unit 16. Narrative and speech}
\hyphenpenalty=10000
\exhyphenpenalty=10000
\emergencystretch=3em
\newcommand{\itempara}[2]{\par\addvspace{3pt}\noindent\hangindent=1.8em\hangafter=1\makebox[1.8em][l]{#1.}#2\par}
\newcommand{\reserve}[1]{\par\penalty-100\begingroup\dimen0=\pagegoal\advance\dimen0 by-\pagetotal\ifdim\dimen0<#1\newpage\fi\endgroup}
% book-source-body:begin

\addfontfeatures{Ligatures=TeXOff}
\raggedright
\booktitle{\BookComponent{ch16-learner}{Unit 16. Extended narrative and speech}UNIT 16}{Extended narrative and speech}

% mdparagraph:P001:begin
A longer narrative combines events, perceptions, reports and quoted speech. The form of each clause establishes its own participants; context connects those participants across the passage. A change of speaker can change the reference of I, you, here and now without changing the people or events described.
% mdparagraph:P001:end

% mdparagraph:P002:begin
The vocabulary and complete systems below prepare the original exercises and Cook’s complete joint Group XII, XIII, Exercises 32–36. The final reading follows Apollonius through a game at the public baths and the invitation that results from it. The four translated examples and the lettered exercises are new material.
% mdparagraph:P002:end

% mdparagraph:P003:begin
\textbf{Scope.} Learn the productive forms for the meanings and constructions stated. A \textbf{reference} form may be consulted; it does not require unaided recall. \textbf{Recognition} forms are identified where they occur. \textbf{Supplied help} remains available during the corresponding task. In verb tables, 1, 2 and 3 denote person; singular you and plural you remain distinct. Preserve vowel length and the distinction between e and ę.
% mdparagraph:P003:end

% teaching: NEW-U16-PRIOR-RECALL
% mdparagraph:P004:begin
The previously learned case functions, strong and weak adjective choice, personal and dual pronouns, direct quotation, finite content clauses, perception with an infinitive, relative clauses, passives and perfect constructions remain in use. A familiar form does not make every meaning or complement of its verb familiar. The new extensions are stated explicitly below.
% mdparagraph:P004:end

\reserve{7\baselineskip}
% mdparagraph:P005:begin
\section*{1. Nouns, adjectives and both}\BookHeading{nav-ch16-learner-0001}{subsection}{1. Nouns, adjectives and both}{1}
% mdparagraph:P005:end

% mdparagraph:P006:begin
Use the nominative for a subject and the accusative for a direct object when the verb takes one. An adnominal genitive expresses possession or another noun-dependent relation. A recipient or a preposition takes its own licensed case. These ordinary uses apply to the nouns below; a special government pattern must be learned separately. Identical endings do not remove the distinction between grammatical positions.\textsuperscript{1}
% mdparagraph:P006:end

\reserve{7\baselineskip}
% mdparagraph:P007:begin
\subsection*{Strong masculine nouns}\BookHeading{nav-ch16-learner-0002}{subsubsection}{Strong masculine nouns}{0}
% mdparagraph:P007:end

% mdparagraph:P008:begin
The nominative and accusative singular have the bare stem; the genitive adds -es and the dative -e. The plural has nominative/\allowbreak{}accusative -as, genitive -a and dative -um.
% mdparagraph:P008:end

% teaching: NEW-U16-VOC-NEAHGEBUR
\par\begin{minipage}{\linewidth}\raggedright
% mdparagraph:P009:begin
\textbf{nēahgebūr} — masculine; neighbour. Stem \textbf{nēahgebūr-}.
% mdparagraph:P009:end

% teaching: NEW-U16-NOM-NEAHGEBUR
% table: NEW-U16-NOM-NEAHGEBUR-DISPLAY
% mdtable:T01:begin
\par\begingroup
\fontsize{10.5}{13.2}\selectfont
\setlength{\LTleft}{0pt}\setlength{\LTright}{\fill}
\noindent
\begin{tabular}{@{}>{\raggedright\arraybackslash}p{36.99mm}>{\raggedright\arraybackslash}p{40.69mm}>{\raggedright\arraybackslash}p{40.69mm}@{}}
\toprule
\textbf{Case} & \textbf{Singular} & \textbf{Plural} \\
\midrule
Nominative & nēahgebūr & nēahgebūras \\
Genitive & nēahgebūres & nēahgebūra \\
Dative & nēahgebūre & nēahgebūrum \\
Accusative & nēahgebūr & nēahgebūras \\
\bottomrule
\end{tabular}
\par\endgroup

% mdtable:T01:end
\end{minipage}\par
% teaching: NEW-U16-VOC-THOTHOR
\par\begin{minipage}{\linewidth}\raggedright
% mdparagraph:P010:begin
\textbf{ðoðor} — masculine; ball. Stem \textbf{ðoðor-}.
% mdparagraph:P010:end

% mdparagraph:P011:begin
The medial o remains before endings. This rare word is learned for the ball exercises and the gymnasium scene.
% mdparagraph:P011:end

% teaching: NEW-U16-NOM-THOTHOR
% table: NEW-U16-NOM-THOTHOR-DISPLAY
% mdtable:T02:begin
\par\begingroup
\fontsize{10.5}{13.2}\selectfont
\setlength{\LTleft}{0pt}\setlength{\LTright}{\fill}
\noindent
\begin{tabular}{@{}>{\raggedright\arraybackslash}p{39.46mm}>{\raggedright\arraybackslash}p{39.46mm}>{\raggedright\arraybackslash}p{39.46mm}@{}}
\toprule
\textbf{Case} & \textbf{Singular} & \textbf{Plural} \\
\midrule
Nominative & ðoðor & ðoðoras \\
Genitive & ðoðores & ðoðora \\
Dative & ðoðore & ðoðorum \\
Accusative & ðoðor & ðoðoras \\
\bottomrule
\end{tabular}
\par\endgroup

% mdtable:T02:end
\end{minipage}\par
% teaching: NEW-U16-VOC-COSS
\par\begin{minipage}{\linewidth}\raggedright
% mdparagraph:P012:begin
\textbf{coss} — masculine; kiss. Stem \textbf{coss-}.
% mdparagraph:P012:end

% teaching: NEW-U16-NOM-COSS
% table: NEW-U16-NOM-COSS-DISPLAY
% mdtable:T03:begin
\par\begingroup
\fontsize{10.5}{13.2}\selectfont
\setlength{\LTleft}{0pt}\setlength{\LTright}{\fill}
\noindent
\begin{tabular}{@{}>{\raggedright\arraybackslash}p{39.46mm}>{\raggedright\arraybackslash}p{39.46mm}>{\raggedright\arraybackslash}p{39.46mm}@{}}
\toprule
\textbf{Case} & \textbf{Singular} & \textbf{Plural} \\
\midrule
Nominative & coss & cossas \\
Genitive & cosses & cossa \\
Dative & cosse & cossum \\
Accusative & coss & cossas \\
\bottomrule
\end{tabular}
\par\endgroup

% mdtable:T03:end
\end{minipage}\par
% teaching: NEW-U16-VOC-DISC
\par\begin{minipage}{\linewidth}\raggedright
% mdparagraph:P013:begin
\textbf{disc} — masculine; dish. Stem \textbf{disc-}.
% mdparagraph:P013:end

% teaching: NEW-U16-NOM-DISC
% table: NEW-U16-NOM-DISC-DISPLAY
% mdtable:T04:begin
\par\begingroup
\fontsize{10.5}{13.2}\selectfont
\setlength{\LTleft}{0pt}\setlength{\LTright}{\fill}
\noindent
\begin{tabular}{@{}>{\raggedright\arraybackslash}p{39.46mm}>{\raggedright\arraybackslash}p{39.46mm}>{\raggedright\arraybackslash}p{39.46mm}@{}}
\toprule
\textbf{Case} & \textbf{Singular} & \textbf{Plural} \\
\midrule
Nominative & disc & discas \\
Genitive & disces & disca \\
Dative & disce & discum \\
Accusative & disc & discas \\
\bottomrule
\end{tabular}
\par\endgroup

% mdtable:T04:end
\end{minipage}\par
% teaching: NEW-U16-VOC-ATH
\par\begin{minipage}{\linewidth}\raggedright
% mdparagraph:P014:begin
\textbf{āð} — masculine; oath. Stem \textbf{āð-}.
% mdparagraph:P014:end

% teaching: NEW-U16-NOM-ATH
% table: NEW-U16-NOM-ATH-DISPLAY
% mdtable:T05:begin
\par\begingroup
\fontsize{10.5}{13.2}\selectfont
\setlength{\LTleft}{0pt}\setlength{\LTright}{\fill}
\noindent
\begin{tabular}{@{}>{\raggedright\arraybackslash}p{39.46mm}>{\raggedright\arraybackslash}p{39.46mm}>{\raggedright\arraybackslash}p{39.46mm}@{}}
\toprule
\textbf{Case} & \textbf{Singular} & \textbf{Plural} \\
\midrule
Nominative & āð & āðas \\
Genitive & āðes & āða \\
Dative & āðe & āðum \\
Accusative & āð & āðas \\
\bottomrule
\end{tabular}
\par\endgroup

% mdtable:T05:end
\end{minipage}\par
% teaching: NEW-U16-VOC-TEAR
\par\begin{minipage}{\linewidth}\raggedright
% mdparagraph:P015:begin
\textbf{tēar} — masculine; tear. Stem \textbf{tēar-}.
% mdparagraph:P015:end

% teaching: NEW-U16-NOM-TEAR
% table: NEW-U16-NOM-TEAR-DISPLAY
% mdtable:T06:begin
\par\begingroup
\fontsize{10.5}{13.2}\selectfont
\setlength{\LTleft}{0pt}\setlength{\LTright}{\fill}
\noindent
\begin{tabular}{@{}>{\raggedright\arraybackslash}p{39.46mm}>{\raggedright\arraybackslash}p{39.46mm}>{\raggedright\arraybackslash}p{39.46mm}@{}}
\toprule
\textbf{Case} & \textbf{Singular} & \textbf{Plural} \\
\midrule
Nominative & tēar & tēaras \\
Genitive & tēares & tēara \\
Dative & tēare & tēarum \\
Accusative & tēar & tēaras \\
\bottomrule
\end{tabular}
\par\endgroup

% mdtable:T06:end
\end{minipage}\par
% teaching: NEW-U16-VOC-MAEGENTHRYMM
\par\begin{minipage}{\linewidth}\raggedright
% mdparagraph:P016:begin
\textbf{mægenðrymm} — masculine; majesty; glory. Stem \textbf{mægenðrymm-}.
% mdparagraph:P016:end

% mdparagraph:P017:begin
Only the singular is learned for this meaning. No count-plural sense is assigned here.
% mdparagraph:P017:end

% teaching: NEW-U16-NOM-MAEGENTHRYMM
% table: NEW-U16-NOM-MAEGENTHRYMM-DISPLAY
% mdtable:T07:begin
\par\begingroup
\fontsize{10.5}{13.2}\selectfont
\setlength{\LTleft}{0pt}\setlength{\LTright}{\fill}
\noindent
\begin{tabular}{@{}>{\raggedright\arraybackslash}p{55.09mm}>{\raggedright\arraybackslash}p{66.1mm}@{}}
\toprule
\textbf{Case} & \textbf{Singular} \\
\midrule
Nominative & mægenðrymm \\
Genitive & mægenðrymmes \\
Dative & mægenðrymme \\
Accusative & mægenðrymm \\
\bottomrule
\end{tabular}
\par\endgroup

% mdtable:T07:end
\end{minipage}\par
\reserve{7\baselineskip}
% mdparagraph:P018:begin
\subsection*{The masculine title in -ere}\BookHeading{nav-ch16-learner-0003}{subsubsection}{The masculine title in -ere}{0}
% mdparagraph:P018:end

% mdparagraph:P019:begin
The singular retains -e in the nominative, accusative and dative; the genitive ends in -es. Here fuluhtere is learned in the singular title “Baptist”.
% mdparagraph:P019:end

% teaching: NEW-U16-VOC-FULUHTERE
\par\begin{minipage}{\linewidth}\raggedright
% mdparagraph:P020:begin
\textbf{fuluhtere} — masculine; baptizer; Baptist as John’s title. Stem \textbf{fuluhter-}.
% mdparagraph:P020:end

% mdparagraph:P021:begin
Only the singular is learned for this meaning. No count-plural sense is assigned here.
% mdparagraph:P021:end

% teaching: NEW-U16-NOM-FULUHTERE
% table: NEW-U16-NOM-FULUHTERE-DISPLAY
% mdtable:T08:begin
\par\begingroup
\fontsize{10.5}{13.2}\selectfont
\setlength{\LTleft}{0pt}\setlength{\LTright}{\fill}
\noindent
\begin{tabular}{@{}>{\raggedright\arraybackslash}p{60.59mm}>{\raggedright\arraybackslash}p{60.59mm}@{}}
\toprule
\textbf{Case} & \textbf{Singular} \\
\midrule
Nominative & fuluhtere \\
Genitive & fuluhteres \\
Dative & fuluhtere \\
Accusative & fuluhtere \\
\bottomrule
\end{tabular}
\par\endgroup

% mdtable:T08:end
\end{minipage}\par
\reserve{7\baselineskip}
% mdparagraph:P022:begin
\subsection*{Weak masculine nouns}\BookHeading{nav-ch16-learner-0004}{subsubsection}{Weak masculine nouns}{0}
% mdparagraph:P022:end

% mdparagraph:P023:begin
The nominative singular ends in -a. The other singular cases and the nominative/\allowbreak{}accusative plural end in -an; the genitive plural has -ena and the dative plural -um.
% mdparagraph:P023:end

% teaching: NEW-U16-VOC-CNAPA
\par\begin{minipage}{\linewidth}\raggedright
% mdparagraph:P024:begin
\textbf{cnapa} — masculine; boy. Stem \textbf{cnap-}.
% mdparagraph:P024:end

% teaching: NEW-U16-NOM-CNAPA
% table: NEW-U16-NOM-CNAPA-DISPLAY
% mdtable:T09:begin
\par\begingroup
\fontsize{10.5}{13.2}\selectfont
\setlength{\LTleft}{0pt}\setlength{\LTright}{\fill}
\noindent
\begin{tabular}{@{}>{\raggedright\arraybackslash}p{39.46mm}>{\raggedright\arraybackslash}p{39.46mm}>{\raggedright\arraybackslash}p{39.46mm}@{}}
\toprule
\textbf{Case} & \textbf{Singular} & \textbf{Plural} \\
\midrule
Nominative & cnapa & cnapan \\
Genitive & cnapan & cnapena \\
Dative & cnapan & cnapum \\
Accusative & cnapan & cnapan \\
\bottomrule
\end{tabular}
\par\endgroup

% mdtable:T09:end
\end{minipage}\par
% teaching: NEW-U16-VOC-GEGIERELA
\par\begin{minipage}{\linewidth}\raggedright
% mdparagraph:P025:begin
\textbf{gegierela} — masculine; garment; clothing. Stem \textbf{gegierel-}.
% mdparagraph:P025:end

% teaching: NEW-U16-NOM-GEGIERELA
% table: NEW-U16-NOM-GEGIERELA-DISPLAY
% mdtable:T10:begin
\par\begingroup
\fontsize{10.5}{13.2}\selectfont
\setlength{\LTleft}{0pt}\setlength{\LTright}{\fill}
\noindent
\begin{tabular}{@{}>{\raggedright\arraybackslash}p{38.19mm}>{\raggedright\arraybackslash}p{38.19mm}>{\raggedright\arraybackslash}p{42.0mm}@{}}
\toprule
\textbf{Case} & \textbf{Singular} & \textbf{Plural} \\
\midrule
Nominative & gegierela & gegierelan \\
Genitive & gegierelan & gegierelena \\
Dative & gegierelan & gegierelum \\
Accusative & gegierelan & gegierelan \\
\bottomrule
\end{tabular}
\par\endgroup

% mdtable:T10:end
\end{minipage}\par
\reserve{7\baselineskip}
% mdparagraph:P026:begin
\subsection*{Strong neuter nouns}\BookHeading{nav-ch16-learner-0005}{subsubsection}{Strong neuter nouns}{0}
% mdparagraph:P026:end

% mdparagraph:P027:begin
Nominative and accusative coincide in each number. The singular genitive and dative use -es and -e; the plural genitive and dative use -a and -um. Learn the nominative/\allowbreak{}accusative plural with the noun. Furlang and nętt have no ending in those plural positions.
% mdparagraph:P027:end

% teaching: NEW-U16-VOC-FURLANG
\par\begin{minipage}{\linewidth}\raggedright
% mdparagraph:P028:begin
\textbf{furlang} — neuter; furlong. Stem \textbf{furlang-}.
% mdparagraph:P028:end

% teaching: NEW-U16-NOM-FURLANG
% table: NEW-U16-NOM-FURLANG-DISPLAY
% mdtable:T11:begin
\par\begingroup
\fontsize{10.5}{13.2}\selectfont
\setlength{\LTleft}{0pt}\setlength{\LTright}{\fill}
\noindent
\begin{tabular}{@{}>{\raggedright\arraybackslash}p{39.46mm}>{\raggedright\arraybackslash}p{39.46mm}>{\raggedright\arraybackslash}p{39.46mm}@{}}
\toprule
\textbf{Case} & \textbf{Singular} & \textbf{Plural} \\
\midrule
Nominative & furlang & furlang \\
Genitive & furlanges & furlanga \\
Dative & furlange & furlangum \\
Accusative & furlang & furlang \\
\bottomrule
\end{tabular}
\par\endgroup

% mdtable:T11:end
\end{minipage}\par
% teaching: NEW-U16-VOC-NETT
\par\begin{minipage}{\linewidth}\raggedright
% mdparagraph:P029:begin
\textbf{nętt} — neuter; net. Stem \textbf{nętt-}.
% mdparagraph:P029:end

% mdparagraph:P030:begin
The zero nominative/\allowbreak{}accusative plural belongs to this noun’s class; it is not inferred merely from the written double consonant.
% mdparagraph:P030:end

% teaching: NEW-U16-NOM-NETT
% table: NEW-U16-NOM-NETT-DISPLAY
% mdtable:T12:begin
\par\begingroup
\fontsize{10.5}{13.2}\selectfont
\setlength{\LTleft}{0pt}\setlength{\LTright}{\fill}
\noindent
\begin{tabular}{@{}>{\raggedright\arraybackslash}p{39.46mm}>{\raggedright\arraybackslash}p{39.46mm}>{\raggedright\arraybackslash}p{39.46mm}@{}}
\toprule
\textbf{Case} & \textbf{Singular} & \textbf{Plural} \\
\midrule
Nominative & nętt & nętt \\
Genitive & nęttes & nętta \\
Dative & nętte & nęttum \\
Accusative & nętt & nętt \\
\bottomrule
\end{tabular}
\par\endgroup

% mdtable:T12:end
\end{minipage}\par
% teaching: NEW-U16-VOC-TEMPEL
\par\begin{minipage}{\linewidth}\raggedright
% mdparagraph:P031:begin
\textbf{tempel} — neuter; temple. Stem \textbf{templ-}.
% mdparagraph:P031:end

% mdparagraph:P032:begin
Use templ- before the displayed vowel endings.
% mdparagraph:P032:end

% teaching: NEW-U16-NOM-TEMPEL
% table: NEW-U16-NOM-TEMPEL-DISPLAY
% mdtable:T13:begin
\par\begingroup
\fontsize{10.5}{13.2}\selectfont
\setlength{\LTleft}{0pt}\setlength{\LTright}{\fill}
\noindent
\begin{tabular}{@{}>{\raggedright\arraybackslash}p{39.46mm}>{\raggedright\arraybackslash}p{39.46mm}>{\raggedright\arraybackslash}p{39.46mm}@{}}
\toprule
\textbf{Case} & \textbf{Singular} & \textbf{Plural} \\
\midrule
Nominative & tempel & templu \\
Genitive & temples & templa \\
Dative & temple & templum \\
Accusative & tempel & templu \\
\bottomrule
\end{tabular}
\par\endgroup

% mdtable:T13:end
\end{minipage}\par
% teaching: NEW-U16-VOC-MAEDEN
\par\begin{minipage}{\linewidth}\raggedright
% mdparagraph:P033:begin
\textbf{mǣden} — neuter; girl; maiden. Stem \textbf{mǣden-}.
% mdparagraph:P033:end

% mdparagraph:P034:begin
Retain mǣden- in the selected forms.
% mdparagraph:P034:end

% teaching: NEW-U16-NOM-MAEDEN
% table: NEW-U16-NOM-MAEDEN-DISPLAY
% mdtable:T14:begin
\par\begingroup
\fontsize{10.5}{13.2}\selectfont
\setlength{\LTleft}{0pt}\setlength{\LTright}{\fill}
\noindent
\begin{tabular}{@{}>{\raggedright\arraybackslash}p{39.46mm}>{\raggedright\arraybackslash}p{39.46mm}>{\raggedright\arraybackslash}p{39.46mm}@{}}
\toprule
\textbf{Case} & \textbf{Singular} & \textbf{Plural} \\
\midrule
Nominative & mǣden & mǣden \\
Genitive & mǣdenes & mǣdena \\
Dative & mǣdene & mǣdenum \\
Accusative & mǣden & mǣden \\
\bottomrule
\end{tabular}
\par\endgroup

% mdtable:T14:end
\end{minipage}\par
% teaching: NEW-U16-VOC-WOLCEN
\par\begin{minipage}{\linewidth}\raggedright
% mdparagraph:P035:begin
\textbf{wolcen} — neuter; cloud. Stem \textbf{wolcn-}.
% mdparagraph:P035:end

% mdparagraph:P036:begin
Use wolcn- before the displayed vowel endings.
% mdparagraph:P036:end

% teaching: NEW-U16-NOM-WOLCEN
% table: NEW-U16-NOM-WOLCEN-DISPLAY
% mdtable:T15:begin
\par\begingroup
\fontsize{10.5}{13.2}\selectfont
\setlength{\LTleft}{0pt}\setlength{\LTright}{\fill}
\noindent
\begin{tabular}{@{}>{\raggedright\arraybackslash}p{39.46mm}>{\raggedright\arraybackslash}p{39.46mm}>{\raggedright\arraybackslash}p{39.46mm}@{}}
\toprule
\textbf{Case} & \textbf{Singular} & \textbf{Plural} \\
\midrule
Nominative & wolcen & wolcnu \\
Genitive & wolcnes & wolcna \\
Dative & wolcne & wolcnum \\
Accusative & wolcen & wolcnu \\
\bottomrule
\end{tabular}
\par\endgroup

% mdtable:T15:end
\end{minipage}\par
\reserve{7\baselineskip}
% mdparagraph:P037:begin
\subsection*{Strong feminine nouns}\BookHeading{nav-ch16-learner-0006}{subsubsection}{Strong feminine nouns}{0}
% mdparagraph:P037:end

% mdparagraph:P038:begin
The selected nouns have no ending in the nominative singular. Their genitive and dative singular have -e. The accusative singular also has -e, except wēn, whose accusative is wēn. The plural, where assigned, has nominative/\allowbreak{}accusative -a, genitive -a and dative -um. Smyltnes doubles its final s before an ending.
% mdparagraph:P038:end

% teaching: NEW-U16-VOC-SMYLTNES
\par\begin{minipage}{\linewidth}\raggedright
% mdparagraph:P039:begin
\textbf{smyltnes} — feminine; calm; calmness. Stem \textbf{smyltness-}.
% mdparagraph:P039:end

% mdparagraph:P040:begin
Only the singular is learned for this meaning. No count-plural sense is assigned here.
% mdparagraph:P040:end

% teaching: NEW-U16-NOM-SMYLTNES
% table: NEW-U16-NOM-SMYLTNES-DISPLAY
% mdtable:T16:begin
\par\begingroup
\fontsize{10.5}{13.2}\selectfont
\setlength{\LTleft}{0pt}\setlength{\LTright}{\fill}
\noindent
\begin{tabular}{@{}>{\raggedright\arraybackslash}p{60.59mm}>{\raggedright\arraybackslash}p{60.59mm}@{}}
\toprule
\textbf{Case} & \textbf{Singular} \\
\midrule
Nominative & smyltnes \\
Genitive & smyltnesse \\
Dative & smyltnesse \\
Accusative & smyltnesse \\
\bottomrule
\end{tabular}
\par\endgroup

% mdtable:T16:end
\end{minipage}\par
% teaching: NEW-U16-VOC-STEFN
\par\begin{minipage}{\linewidth}\raggedright
% mdparagraph:P041:begin
\textbf{stefn} — feminine; voice. Stem \textbf{stefn-}.
% mdparagraph:P041:end

% teaching: NEW-U16-NOM-STEFN
% table: NEW-U16-NOM-STEFN-DISPLAY
% mdtable:T17:begin
\par\begingroup
\fontsize{10.5}{13.2}\selectfont
\setlength{\LTleft}{0pt}\setlength{\LTright}{\fill}
\noindent
\begin{tabular}{@{}>{\raggedright\arraybackslash}p{39.46mm}>{\raggedright\arraybackslash}p{39.46mm}>{\raggedright\arraybackslash}p{39.46mm}@{}}
\toprule
\textbf{Case} & \textbf{Singular} & \textbf{Plural} \\
\midrule
Nominative & stefn & stefna \\
Genitive & stefne & stefna \\
Dative & stefne & stefnum \\
Accusative & stefne & stefna \\
\bottomrule
\end{tabular}
\par\endgroup

% mdtable:T17:end
\end{minipage}\par
% teaching: NEW-U16-VOC-ANSIEN
\par\begin{minipage}{\linewidth}\raggedright
% mdparagraph:P042:begin
\textbf{ansīen} — feminine; appearance. Stem \textbf{ansīen-}.
% mdparagraph:P042:end

% mdparagraph:P043:begin
Only the singular is learned for this meaning. No count-plural sense is assigned here.
% mdparagraph:P043:end

% teaching: NEW-U16-NOM-ANSIEN
% table: NEW-U16-NOM-ANSIEN-DISPLAY
% mdtable:T18:begin
\par\begingroup
\fontsize{10.5}{13.2}\selectfont
\setlength{\LTleft}{0pt}\setlength{\LTright}{\fill}
\noindent
\begin{tabular}{@{}>{\raggedright\arraybackslash}p{60.59mm}>{\raggedright\arraybackslash}p{60.59mm}@{}}
\toprule
\textbf{Case} & \textbf{Singular} \\
\midrule
Nominative & ansīen \\
Genitive & ansīene \\
Dative & ansīene \\
Accusative & ansīene \\
\bottomrule
\end{tabular}
\par\endgroup

% mdtable:T18:end
\end{minipage}\par
% teaching: NEW-U16-VOC-SEALF
\par\begin{minipage}{\linewidth}\raggedright
% mdparagraph:P044:begin
\textbf{sealf} — feminine; ointment. Stem \textbf{sealf-}.
% mdparagraph:P044:end

% mdparagraph:P045:begin
Only the singular is learned for this meaning. No count-plural sense is assigned here.
% mdparagraph:P045:end

% teaching: NEW-U16-NOM-SEALF
% table: NEW-U16-NOM-SEALF-DISPLAY
% mdtable:T19:begin
\par\begingroup
\fontsize{10.5}{13.2}\selectfont
\setlength{\LTleft}{0pt}\setlength{\LTright}{\fill}
\noindent
\begin{tabular}{@{}>{\raggedright\arraybackslash}p{60.59mm}>{\raggedright\arraybackslash}p{60.59mm}@{}}
\toprule
\textbf{Case} & \textbf{Singular} \\
\midrule
Nominative & sealf \\
Genitive & sealfe \\
Dative & sealfe \\
Accusative & sealfe \\
\bottomrule
\end{tabular}
\par\endgroup

% mdtable:T19:end
\end{minipage}\par
% teaching: NEW-U16-VOC-WEN
\par\begin{minipage}{\linewidth}\raggedright
% mdparagraph:P046:begin
\textbf{wēn} — feminine; expectation; chance, especially wēn is ðæt. Stem \textbf{wēn-}.
% mdparagraph:P046:end

% mdparagraph:P047:begin
Only the singular is learned for this meaning. No count-plural sense is assigned here.
% mdparagraph:P047:end

% teaching: NEW-U16-NOM-WEN
% table: NEW-U16-NOM-WEN-DISPLAY
% mdtable:T20:begin
\par\begingroup
\fontsize{10.5}{13.2}\selectfont
\setlength{\LTleft}{0pt}\setlength{\LTright}{\fill}
\noindent
\begin{tabular}{@{}>{\raggedright\arraybackslash}p{60.59mm}>{\raggedright\arraybackslash}p{60.59mm}@{}}
\toprule
\textbf{Case} & \textbf{Singular} \\
\midrule
Nominative & wēn \\
Genitive & wēne \\
Dative & wēne \\
Accusative & wēn \\
\bottomrule
\end{tabular}
\par\endgroup

% mdtable:T20:end
\end{minipage}\par
% teaching: NEW-U16-VOC-SPRAEC
\par\begin{minipage}{\linewidth}\raggedright
% mdparagraph:P048:begin
\textbf{sprǣc} — feminine; saying; speech. Stem \textbf{sprǣc-}.
% mdparagraph:P048:end

% teaching: NEW-U16-NOM-SPRAEC
% table: NEW-U16-NOM-SPRAEC-DISPLAY
% mdtable:T21:begin
\par\begingroup
\fontsize{10.5}{13.2}\selectfont
\setlength{\LTleft}{0pt}\setlength{\LTright}{\fill}
\noindent
\begin{tabular}{@{}>{\raggedright\arraybackslash}p{39.46mm}>{\raggedright\arraybackslash}p{39.46mm}>{\raggedright\arraybackslash}p{39.46mm}@{}}
\toprule
\textbf{Case} & \textbf{Singular} & \textbf{Plural} \\
\midrule
Nominative & sprǣc & sprǣca \\
Genitive & sprǣce & sprǣca \\
Dative & sprǣce & sprǣcum \\
Accusative & sprǣce & sprǣca \\
\bottomrule
\end{tabular}
\par\endgroup

% mdtable:T21:end
\end{minipage}\par
\reserve{7\baselineskip}
% mdparagraph:P049:begin
\subsection*{Adjectives}\BookHeading{nav-ch16-learner-0007}{subsubsection}{Adjectives}{0}
% mdparagraph:P049:end

% mdparagraph:P050:begin
The following tables give the complete strong and weak positive forms. Use the strong or weak series under the conditions already learned, with agreement in case, number and gender. In ordinary predicate use the adjective also agrees with its subject. The strong instrumental singular forms are for recognition only; their appearance here introduces no new instrumental construction. The other displayed positions are productive.\textsuperscript{2}
% mdparagraph:P050:end

% teaching: NEW-U16-VOC-HEALT
\par\begin{minipage}{\linewidth}\raggedright
% mdparagraph:P051:begin
\textbf{healt} — halt; lame.
% mdparagraph:P051:end

% mdparagraph:P052:begin
The heavy stem takes no -u in the strong feminine nominative singular or neuter nominative/\allowbreak{}accusative plural.
% mdparagraph:P052:end

% teaching: NEW-U16-ADJ-HEALT
% table: NEW-U16-ADJ-HEALT-DISPLAY
% mdtable:T22:begin
% course-table-style: 2026-10
% course-panel:T22-P01:begin
\par\begingroup
\fontsize{10.5}{13.2}\selectfont
\setlength{\LTleft}{0pt}\setlength{\LTright}{\fill}
\noindent
\begin{tabular}{@{}>{\raggedright\arraybackslash}p{34.062mm}>{\raggedright\arraybackslash}p{29.341mm}>{\raggedright\arraybackslash}p{25.069mm}>{\raggedright\arraybackslash}p{27.092mm}@{}}
\toprule
\textbf{Series /\allowbreak{} case} & \textbf{Masculine singular} & \textbf{Neuter singular} & \textbf{Feminine singular} \\
\midrule
Strong nominative & healt & healt & healt \\
Strong genitive & healtes & healtes & healtre \\
Strong dative & healtum & healtum & healtre \\
Strong accusative & healtne & healt & healte \\
Strong instrumental & healte & healte & — \\
Weak nominative & healta & healte & healte \\
Weak genitive & healtan & healtan & healtan \\
Weak dative & healtan & healtan & healtan \\
Weak accusative & healtan & healte & healtan \\
\bottomrule
\end{tabular}
\par\endgroup
% course-panel:T22-P01:end
\end{minipage}\par
\par\addvspace{5pt}
% course-panel:T22-P02:begin
\par\begingroup
\fontsize{10.5}{13.2}\selectfont
\setlength{\LTleft}{0pt}\setlength{\LTright}{\fill}
\noindent
\begin{tabular}{@{}>{\raggedright\arraybackslash}p{36.136mm}>{\raggedright\arraybackslash}p{28.98mm}>{\raggedright\arraybackslash}p{23.852mm}>{\raggedright\arraybackslash}p{26.595mm}@{}}
\toprule
\textbf{Series /\allowbreak{} case} & \textbf{Masculine plural} & \textbf{Neuter plural} & \textbf{Feminine plural} \\
\midrule
Strong nominative & healte & healt & healta \\
Strong genitive & healtra & healtra & healtra \\
Strong dative & healtum & healtum & healtum \\
Strong accusative & healte & healt & healta \\
Strong instrumental & — & — & — \\
Weak nominative & healtan & healtan & healtan \\
Weak genitive & healtra & healtra & healtra \\
Weak dative & healtum & healtum & healtum \\
Weak accusative & healtan & healtan & healtan \\
\bottomrule
\end{tabular}
\par\endgroup
% course-panel:T22-P02:end
% mdtable:T22:end

% teaching: NEW-U16-VOC-GELARED
\par\begin{minipage}{\linewidth}\raggedright
% mdparagraph:P053:begin
\textbf{gelǣred} — educated; trained; skilled.
% mdparagraph:P053:end

% mdparagraph:P054:begin
The selected forms keep the full gelǣred- stem before endings. Shorter contracted variants do not replace this table.
% mdparagraph:P054:end

% teaching: NEW-U16-ADJ-GELARED
% table: NEW-U16-ADJ-GELARED-DISPLAY
% mdtable:T23:begin
% course-table-style: 2026-10
% course-panel:T23-P01:begin
\par\begingroup
\fontsize{10.5}{13.2}\selectfont
\setlength{\LTleft}{0pt}\setlength{\LTright}{\fill}
\noindent
\begin{tabular}{@{}>{\raggedright\arraybackslash}p{33.348mm}>{\raggedright\arraybackslash}p{28.726mm}>{\raggedright\arraybackslash}p{25.754mm}>{\raggedright\arraybackslash}p{27.735mm}@{}}
\toprule
\textbf{Series /\allowbreak{} case} & \textbf{Masculine singular} & \textbf{Neuter singular} & \textbf{Feminine singular} \\
\midrule
Strong nominative & gelǣred & gelǣred & gelǣredu \\
Strong genitive & gelǣredes & gelǣredes & gelǣredre \\
Strong dative & gelǣredum & gelǣredum & gelǣredre \\
Strong accusative & gelǣredne & gelǣred & gelǣrede \\
Strong instrumental & gelǣrede & gelǣrede & — \\
Weak nominative & gelǣreda & gelǣrede & gelǣrede \\
Weak genitive & gelǣredan & gelǣredan & gelǣredan \\
Weak dative & gelǣredan & gelǣredan & gelǣredan \\
Weak accusative & gelǣredan & gelǣrede & gelǣredan \\
\bottomrule
\end{tabular}
\par\endgroup
% course-panel:T23-P01:end
\end{minipage}\par
\par\addvspace{5pt}
% course-panel:T23-P02:begin
\par\begingroup
\fontsize{10.5}{13.2}\selectfont
\setlength{\LTleft}{0pt}\setlength{\LTright}{\fill}
\noindent
\begin{tabular}{@{}>{\raggedright\arraybackslash}p{35.157mm}>{\raggedright\arraybackslash}p{28.195mm}>{\raggedright\arraybackslash}p{25.062mm}>{\raggedright\arraybackslash}p{27.151mm}@{}}
\toprule
\textbf{Series /\allowbreak{} case} & \textbf{Masculine plural} & \textbf{Neuter plural} & \textbf{Feminine plural} \\
\midrule
Strong nominative & gelǣrede & gelǣredu & gelǣreda \\
Strong genitive & gelǣredra & gelǣredra & gelǣredra \\
Strong dative & gelǣredum & gelǣredum & gelǣredum \\
Strong accusative & gelǣrede & gelǣredu & gelǣreda \\
Strong instrumental & — & — & — \\
Weak nominative & gelǣredan & gelǣredan & gelǣredan \\
Weak genitive & gelǣredra & gelǣredra & gelǣredra \\
Weak dative & gelǣredum & gelǣredum & gelǣredum \\
Weak accusative & gelǣredan & gelǣredan & gelǣredan \\
\bottomrule
\end{tabular}
\par\endgroup
% course-panel:T23-P02:end
% mdtable:T23:end

% teaching: NEW-U16-VOC-EORTHLIC
\par\begin{minipage}{\linewidth}\raggedright
% mdparagraph:P055:begin
\textbf{eorðlic} — earthly.
% mdparagraph:P055:end

% mdparagraph:P056:begin
The selected forms keep the full eorðlic- stem before endings.
% mdparagraph:P056:end

% teaching: NEW-U16-ADJ-EORTHLIC
% table: NEW-U16-ADJ-EORTHLIC-DISPLAY
% mdtable:T24:begin
% course-table-style: 2026-10
% course-panel:T24-P01:begin
\par\begingroup
\fontsize{10.5}{13.2}\selectfont
\setlength{\LTleft}{0pt}\setlength{\LTright}{\fill}
\noindent
\begin{tabular}{@{}>{\raggedright\arraybackslash}p{33.348mm}>{\raggedright\arraybackslash}p{28.726mm}>{\raggedright\arraybackslash}p{25.754mm}>{\raggedright\arraybackslash}p{27.735mm}@{}}
\toprule
\textbf{Series /\allowbreak{} case} & \textbf{Masculine singular} & \textbf{Neuter singular} & \textbf{Feminine singular} \\
\midrule
Strong nominative & eorðlic & eorðlic & eorðlicu \\
Strong genitive & eorðlices & eorðlices & eorðlicre \\
Strong dative & eorðlicum & eorðlicum & eorðlicre \\
Strong accusative & eorðlicne & eorðlic & eorðlice \\
Strong instrumental & eorðlice & eorðlice & — \\
Weak nominative & eorðlica & eorðlice & eorðlice \\
Weak genitive & eorðlican & eorðlican & eorðlican \\
Weak dative & eorðlican & eorðlican & eorðlican \\
Weak accusative & eorðlican & eorðlice & eorðlican \\
\bottomrule
\end{tabular}
\par\endgroup
% course-panel:T24-P01:end
\end{minipage}\par
\par\addvspace{5pt}
% course-panel:T24-P02:begin
\par\begingroup
\fontsize{10.5}{13.2}\selectfont
\setlength{\LTleft}{0pt}\setlength{\LTright}{\fill}
\noindent
\begin{tabular}{@{}>{\raggedright\arraybackslash}p{35.157mm}>{\raggedright\arraybackslash}p{28.195mm}>{\raggedright\arraybackslash}p{25.062mm}>{\raggedright\arraybackslash}p{27.151mm}@{}}
\toprule
\textbf{Series /\allowbreak{} case} & \textbf{Masculine plural} & \textbf{Neuter plural} & \textbf{Feminine plural} \\
\midrule
Strong nominative & eorðlice & eorðlicu & eorðlica \\
Strong genitive & eorðlicra & eorðlicra & eorðlicra \\
Strong dative & eorðlicum & eorðlicum & eorðlicum \\
Strong accusative & eorðlice & eorðlicu & eorðlica \\
Strong instrumental & — & — & — \\
Weak nominative & eorðlican & eorðlican & eorðlican \\
Weak genitive & eorðlicra & eorðlicra & eorðlicra \\
Weak dative & eorðlicum & eorðlicum & eorðlicum \\
Weak accusative & eorðlican & eorðlican & eorðlican \\
\bottomrule
\end{tabular}
\par\endgroup
% course-panel:T24-P02:end
% mdtable:T24:end

% teaching: NEW-U16-VOC-NEARU
\par\begin{minipage}{\linewidth}\raggedright
% mdparagraph:P057:begin
\textbf{nearu} — narrow.
% mdparagraph:P057:end

% mdparagraph:P058:begin
Learn the alternating stems together: nearu in the strong nominative singular, neuter accusative singular and neuter nominative/\allowbreak{}accusative plural; nearw- before vowel endings; nearo- before consonantal -ne, -re and -ra. The weak forms use nearw- except for genitive plural nearora.
% mdparagraph:P058:end

% teaching: NEW-U16-ADJ-NEARU
% table: NEW-U16-ADJ-NEARU-DISPLAY
% mdtable:T25:begin
% course-table-style: 2026-10
% course-panel:T25-P01:begin
\par\begingroup
\fontsize{10.5}{13.2}\selectfont
\setlength{\LTleft}{0pt}\setlength{\LTright}{\fill}
\noindent
\begin{tabular}{@{}>{\raggedright\arraybackslash}p{34.062mm}>{\raggedright\arraybackslash}p{29.341mm}>{\raggedright\arraybackslash}p{25.069mm}>{\raggedright\arraybackslash}p{27.092mm}@{}}
\toprule
\textbf{Series /\allowbreak{} case} & \textbf{Masculine singular} & \textbf{Neuter singular} & \textbf{Feminine singular} \\
\midrule
Strong nominative & nearu & nearu & nearu \\
Strong genitive & nearwes & nearwes & nearore \\
Strong dative & nearwum & nearwum & nearore \\
Strong accusative & nearone & nearu & nearwe \\
Strong instrumental & nearwe & nearwe & — \\
Weak nominative & nearwa & nearwe & nearwe \\
Weak genitive & nearwan & nearwan & nearwan \\
Weak dative & nearwan & nearwan & nearwan \\
Weak accusative & nearwan & nearwe & nearwan \\
\bottomrule
\end{tabular}
\par\endgroup
% course-panel:T25-P01:end
\end{minipage}\par
\par\addvspace{5pt}
% course-panel:T25-P02:begin
\par\begingroup
\fontsize{10.5}{13.2}\selectfont
\setlength{\LTleft}{0pt}\setlength{\LTright}{\fill}
\noindent
\begin{tabular}{@{}>{\raggedright\arraybackslash}p{36.136mm}>{\raggedright\arraybackslash}p{28.98mm}>{\raggedright\arraybackslash}p{23.852mm}>{\raggedright\arraybackslash}p{26.595mm}@{}}
\toprule
\textbf{Series /\allowbreak{} case} & \textbf{Masculine plural} & \textbf{Neuter plural} & \textbf{Feminine plural} \\
\midrule
Strong nominative & nearwe & nearu & nearwa \\
Strong genitive & nearora & nearora & nearora \\
Strong dative & nearwum & nearwum & nearwum \\
Strong accusative & nearwe & nearu & nearwa \\
Strong instrumental & — & — & — \\
Weak nominative & nearwan & nearwan & nearwan \\
Weak genitive & nearora & nearora & nearora \\
Weak dative & nearwum & nearwum & nearwum \\
Weak accusative & nearwan & nearwan & nearwan \\
\bottomrule
\end{tabular}
\par\endgroup
% course-panel:T25-P02:end
% mdtable:T25:end

\reserve{18\baselineskip}
% mdparagraph:P059:begin
\subsection*{Both}\BookHeading{nav-ch16-learner-0008}{subsubsection}{Both}{0}
% mdparagraph:P059:end

% teaching: NEW-U16-VOC-BEGEN
\par\begin{minipage}{\linewidth}\raggedright
% mdparagraph:P060:begin
\textbf{bēgen}\BookPoint{bk-ch16-learner-02c84df7b74b} — both. Its case and gender follow the noun or pronoun to which it refers. It has a plural declension, although it denotes two: it is not a dual inflection. It can accompany a noun or stand in apposition to a pronoun; \textbf{unc bǣm} means “to us both” and \textbf{inc bǣm}, “to you both”. The dative \textbf{bām} is also valid; neuter \textbf{bū} is an alternative to \textbf{bā}. The table selects one spelling for each position.
% mdparagraph:P060:end

% teaching: NEW-U16-NOM-BEGEN
% table: NEW-U16-NOM-BEGEN-DISPLAY
% mdtable:T26:begin
% course-table-style: 2026-10
% course-panel:T26-P01:begin
\par\begingroup
\fontsize{10.5}{13.2}\selectfont
\setlength{\LTleft}{0pt}\setlength{\LTright}{\fill}
\noindent
\begin{tabular}{@{}>{\raggedright\arraybackslash}p{22.0mm}>{\raggedright\arraybackslash}p{31.188mm}>{\raggedright\arraybackslash}p{31.188mm}>{\raggedright\arraybackslash}p{31.188mm}@{}}
\toprule
\textbf{Case} & \textbf{Masculine plural} & \textbf{Neuter plural} & \textbf{Feminine plural} \\
\midrule
Nominative & bēgen & bā & bā \\
Genitive & bēgra & bēgra & bēgra \\
Dative & bǣm & bǣm & bǣm \\
Accusative & bēgen & bā & bā \\
\bottomrule
\end{tabular}
\par\endgroup
% course-panel:T26-P01:end
% mdtable:T26:end
\end{minipage}\par
% teaching: NEW-U16-VOC-DOM
% mdparagraph:P061:begin
\textbf{dōm} — judgment; a judgment made in judging. Use the previously learned singular forms, including accusative \textbf{dōm}. This ordinary judgment sense is additional to the earlier religious and fate-related uses; no new abstract plural is required.
% mdparagraph:P061:end

\reserve{7\baselineskip}
% mdparagraph:P062:begin
\section*{2. Verbal systems and their basic uses}\BookHeading{nav-ch16-learner-0009}{subsection}{2. Verbal systems and their basic uses}{1}
% mdparagraph:P062:end

% teaching: NEW-U16-GRAMMAR
% mdparagraph:P063:begin
Each indicative table gives all six person-and-number positions in the present and preterite. The basic uses beside it state the objects and complements licensed here. A present can refer to a future event when the context supplies that viewpoint; English will does not by itself require willan. Commands, subjunctives and nonfinite forms are treated separately.\textsuperscript{3}
% mdparagraph:P063:end

% teaching: NEW-U16-PRINCIPAL-PARTS
\reserve{18\baselineskip}
% mdparagraph:P064:begin
\subsection*{Principal parts for reference}\BookHeading{nav-ch16-learner-0010}{subsubsection}{Principal parts for reference}{0}
% mdparagraph:P064:end

% mdparagraph:P065:begin
These dictionary forms show the stems from which the selected systems are formed. They remain available for reference. Productive finite forms are given in the complete tables that follow; productive nonfinite forms are identified in their own tables.
% mdparagraph:P065:end

% mdtable:T27:begin
\par\begingroup
\fontsize{10.5}{13.2}\selectfont
\setlength{\LTleft}{0pt}\setlength{\LTright}{\fill}
\noindent
\begin{tabular}{@{}>{\raggedright\arraybackslash}p{24.84mm}>{\raggedright\arraybackslash}p{36.02mm}>{\raggedright\arraybackslash}p{57.51mm}@{}}
\toprule
\textbf{Verb} & \textbf{Class} & \textbf{Principal parts: reference} \\
\midrule
ǣswician & weak II & ǣswician; ǣswicode; ǣswicod \\
blissian & weak II & blissian; blissode; blissod \\
wuldrian & weak II & wuldrian; wuldrode; wuldrod \\
plegian & weak II & plegian; plegode; plegod \\
behēafdian & weak II & behēafdian; behēafdode; behēafdod \\
cleopian & weak II & cleopian; cleopode; cleopod \\
ārīsan & strong I & ārīsan; ārās; ārison; ārisen \\
āceorfan & strong III & āceorfan; ācearf; ācurfon; ācorfen \\
forlēosan & strong II & forlēosan; forlēas; forluron; forloren \\
iernan & strong III & iernan; arn; urnon; urnen \\
ðwēan & strong VI contracted & ðwēan; ðwōh; ðwōgon; ðwegen \\
bewęndan & weak I & bewęndan; bewęnde; bewęnd \\
smierwan & weak I rw & smierwan; smierede; smiered \\
ālīesan & weak I & ālīesan; ālīesde; ālīesed \\
\bottomrule
\end{tabular}
\par\endgroup

% mdtable:T27:end
\reserve{7\baselineskip}
% mdparagraph:P066:begin
\subsection*{Weak verbs in -ian}\BookHeading{nav-ch16-learner-0011}{subsubsection}{Weak verbs in -ian}{0}
% mdparagraph:P066:end

% teaching: u16-aeswician-frame
\par\begin{minipage}{\linewidth}\raggedright
% mdparagraph:P067:begin
\textbf{ǣswician} — offend; cause someone to stumble. The person affected is accusative. In the exercise, hand or foot is a disjunctive singular subject and takes a singular verb; the meaning is moral offence, not simply physical injury.
% mdparagraph:P067:end

% mdparagraph:P068:begin
Weak II. The selected present endings are -ie, -ast, -að, -iað; the preterite uses -ode, -odest, -odon. Keep ǣswic- before these endings.
% mdparagraph:P068:end

% teaching: NEW-U16-PAR-AESWICIAN-IND
% table: NEW-U16-PAR-AESWICIAN-IND-DISPLAY
% mdtable:T28:begin
% course-table-style: 2026-10
% course-panel:T28-P01:begin
\par\begingroup
\fontsize{10.5}{13.2}\selectfont
\setlength{\LTleft}{0pt}\setlength{\LTright}{\fill}
\noindent
\begin{tabular}{@{}>{\raggedright\arraybackslash}p{39.459mm}>{\raggedright\arraybackslash}p{39.459mm}>{\raggedright\arraybackslash}p{39.459mm}@{}}
\toprule
\textbf{Person and number} & \textbf{Present} & \textbf{Preterite} \\
\midrule
1 sg. & ǣswicie & ǣswicode \\
2 sg. & ǣswicast & ǣswicodest \\
3 sg. & ǣswicað & ǣswicode \\
1–3 pl. & ǣswiciað & ǣswicodon \\
\bottomrule
\end{tabular}
\par\endgroup
% course-panel:T28-P01:end
% mdtable:T28:end
\end{minipage}\par
% teaching: u16-blissian-frame
\par\begin{minipage}{\linewidth}\raggedright
% mdparagraph:P069:begin
\textbf{blissian} — rejoice, intransitive. A companion may be introduced by mid with the dative. Ofer with the accusative introduces what one rejoices over; this use is explained again with comparison below.
% mdparagraph:P069:end

% mdparagraph:P070:begin
Weak II. Add the same selected endings to bliss-.
% mdparagraph:P070:end

% teaching: NEW-U16-PAR-BLISSIAN-IND
% table: NEW-U16-PAR-BLISSIAN-IND-DISPLAY
% mdtable:T29:begin
% course-table-style: 2026-10
% course-panel:T29-P01:begin
\par\begingroup
\fontsize{10.5}{13.2}\selectfont
\setlength{\LTleft}{0pt}\setlength{\LTright}{\fill}
\noindent
\begin{tabular}{@{}>{\raggedright\arraybackslash}p{39.459mm}>{\raggedright\arraybackslash}p{39.459mm}>{\raggedright\arraybackslash}p{39.459mm}@{}}
\toprule
\textbf{Person and number} & \textbf{Present} & \textbf{Preterite} \\
\midrule
1 sg. & blissie & blissode \\
2 sg. & blissast & blissodest \\
3 sg. & blissað & blissode \\
1–3 pl. & blissiað & blissodon \\
\bottomrule
\end{tabular}
\par\endgroup
% course-panel:T29-P01:end
% mdtable:T29:end
\end{minipage}\par
% teaching: u16-wuldrian-frame
\par\begin{minipage}{\linewidth}\raggedright
% mdparagraph:P071:begin
\textbf{wuldrian} — glorify, with an accusative person. A masculine speaker glorifying himself uses the accusative reflexive phrase mē selfne.
% mdparagraph:P071:end

% mdparagraph:P072:begin
Weak II. Add the selected endings to wuldr-.
% mdparagraph:P072:end

% teaching: NEW-U16-PAR-WULDRIAN-IND
% table: NEW-U16-PAR-WULDRIAN-IND-DISPLAY
% mdtable:T30:begin
% course-table-style: 2026-10
% course-panel:T30-P01:begin
\par\begingroup
\fontsize{10.5}{13.2}\selectfont
\setlength{\LTleft}{0pt}\setlength{\LTright}{\fill}
\noindent
\begin{tabular}{@{}>{\raggedright\arraybackslash}p{39.459mm}>{\raggedright\arraybackslash}p{39.459mm}>{\raggedright\arraybackslash}p{39.459mm}@{}}
\toprule
\textbf{Person and number} & \textbf{Present} & \textbf{Preterite} \\
\midrule
1 sg. & wuldrie & wuldrode \\
2 sg. & wuldrast & wuldrodest \\
3 sg. & wuldrað & wuldrode \\
1–3 pl. & wuldriað & wuldrodon \\
\bottomrule
\end{tabular}
\par\endgroup
% course-panel:T30-P01:end
% mdtable:T30:end
\end{minipage}\par
% teaching: u16-plegian-play
\par\begin{minipage}{\linewidth}\raggedright
% mdparagraph:P073:begin
\textbf{plegian} — play, intransitive. It can be modified by time and manner adverbs. The reading’s expressions for fellow players and the ball receive their own local help.
% mdparagraph:P073:end

% mdparagraph:P074:begin
Weak II. The selected stems are plegi- in the plural present and pleg- before -ode/\allowbreak{}-odon. Keep the selected spelling throughout the table.
% mdparagraph:P074:end

% teaching: NEW-U16-PAR-PLEGIAN-IND
% table: NEW-U16-PAR-PLEGIAN-IND-DISPLAY
% mdtable:T31:begin
% course-table-style: 2026-10
% course-panel:T31-P01:begin
\par\begingroup
\fontsize{10.5}{13.2}\selectfont
\setlength{\LTleft}{0pt}\setlength{\LTright}{\fill}
\noindent
\begin{tabular}{@{}>{\raggedright\arraybackslash}p{39.459mm}>{\raggedright\arraybackslash}p{39.459mm}>{\raggedright\arraybackslash}p{39.459mm}@{}}
\toprule
\textbf{Person and number} & \textbf{Present} & \textbf{Preterite} \\
\midrule
1 sg. & plegie & plegode \\
2 sg. & plegast & plegodest \\
3 sg. & plegað & plegode \\
1–3 pl. & plegiað & plegodon \\
\bottomrule
\end{tabular}
\par\endgroup
% course-panel:T31-P01:end
% mdtable:T31:end
\end{minipage}\par
% teaching: u16-beheafdian-frame
\par\begin{minipage}{\linewidth}\raggedright
% mdparagraph:P075:begin
\textbf{behēafdian} — behead, with an accusative person.
% mdparagraph:P075:end

% mdparagraph:P076:begin
Weak II. Retain the prefix and behēafd- stem in the selected forms.
% mdparagraph:P076:end

% teaching: NEW-U16-PAR-BEHEAFDIAN-IND
% table: NEW-U16-PAR-BEHEAFDIAN-IND-DISPLAY
% mdtable:T32:begin
% course-table-style: 2026-10
% course-panel:T32-P01:begin
\par\begingroup
\fontsize{10.5}{13.2}\selectfont
\setlength{\LTleft}{0pt}\setlength{\LTright}{\fill}
\noindent
\begin{tabular}{@{}>{\raggedright\arraybackslash}p{39.459mm}>{\raggedright\arraybackslash}p{39.459mm}>{\raggedright\arraybackslash}p{39.459mm}@{}}
\toprule
\textbf{Person and number} & \textbf{Present} & \textbf{Preterite} \\
\midrule
1 sg. & behēafdie & behēafdode \\
2 sg. & behēafdast & behēafdodest \\
3 sg. & behēafdað & behēafdode \\
1–3 pl. & behēafdiað & behēafdodon \\
\bottomrule
\end{tabular}
\par\endgroup
% course-panel:T32-P01:end
% mdtable:T32:end
\end{minipage}\par
% teaching: u16-cleopian-cry
\par\begin{minipage}{\linewidth}\raggedright
% mdparagraph:P077:begin
\textbf{cleopian} — cry aloud. The use here can introduce quoted words. The command forms are already familiar; the complete indicative is now learned.
% mdparagraph:P077:end

% mdparagraph:P078:begin
Weak II, with the same selected present and preterite endings as the preceding verbs.
% mdparagraph:P078:end

% teaching: NEW-U16-PAR-CLEOPIAN-IND
% table: NEW-U16-PAR-CLEOPIAN-IND-DISPLAY
% mdtable:T33:begin
% course-table-style: 2026-10
% course-panel:T33-P01:begin
\par\begingroup
\fontsize{10.5}{13.2}\selectfont
\setlength{\LTleft}{0pt}\setlength{\LTright}{\fill}
\noindent
\begin{tabular}{@{}>{\raggedright\arraybackslash}p{39.459mm}>{\raggedright\arraybackslash}p{39.459mm}>{\raggedright\arraybackslash}p{39.459mm}@{}}
\toprule
\textbf{Person and number} & \textbf{Present} & \textbf{Preterite} \\
\midrule
1 sg. & cleopie & cleopode \\
2 sg. & cleopast & cleopodest \\
3 sg. & cleopað & cleopode \\
1–3 pl. & cleopiað & cleopodon \\
\bottomrule
\end{tabular}
\par\endgroup
% course-panel:T33-P01:end
% mdtable:T33:end
\end{minipage}\par
\reserve{7\baselineskip}
% mdparagraph:P079:begin
\subsection*{Strong and contracted verbs}\BookHeading{nav-ch16-learner-0012}{subsubsection}{Strong and contracted verbs}{0}
% mdparagraph:P079:end

% mdparagraph:P080:begin
The ordinary strong present has singular endings -e, -st, -ð and plural -að. The preterite first and third singular have no added ending; second singular has -e on the plural stem, and all plural persons have -on. The stem changes and contractions stated below modify these ordinary combinations.
% mdparagraph:P080:end

\par\begin{minipage}{\linewidth}\raggedright
% mdparagraph:P081:begin
\textbf{ārīsan} — arise; get up, intransitive. This indicative system supplements the previously learned subjunctive.
% mdparagraph:P081:end

% mdparagraph:P082:begin
Strong I: ī /\allowbreak{} ā /\allowbreak{} i /\allowbreak{} i. Present -s plus -st or -ð gives the selected -st in both second and third singular. Preterite second singular uses the plural stem.
% mdparagraph:P082:end

% teaching: NEW-U16-PAR-ARISAN-IND
% table: NEW-U16-PAR-ARISAN-IND-DISPLAY
% mdtable:T34:begin
% course-table-style: 2026-10
% course-panel:T34-P01:begin
\par\begingroup
\fontsize{10.5}{13.2}\selectfont
\setlength{\LTleft}{0pt}\setlength{\LTright}{\fill}
\noindent
\begin{tabular}{@{}>{\raggedright\arraybackslash}p{39.459mm}>{\raggedright\arraybackslash}p{39.459mm}>{\raggedright\arraybackslash}p{39.459mm}@{}}
\toprule
\textbf{Person and number} & \textbf{Present} & \textbf{Preterite} \\
\midrule
1 sg. & ārīse & ārās \\
2 sg. & ārīst & ārise \\
3 sg. & ārīst & ārās \\
1–3 pl. & ārīsað & ārison \\
\bottomrule
\end{tabular}
\par\endgroup
% course-panel:T34-P01:end
% mdtable:T34:end
\end{minipage}\par
% teaching: u16-aceorfan-frame
\par\begin{minipage}{\linewidth}\raggedright
% mdparagraph:P083:begin
\textbf{āceorfan} — cut off, with an accusative object. The separate particle of can reinforce “off”; it has no governed noun after it in that use.
% mdparagraph:P083:end

% mdparagraph:P084:begin
Strong III: eo /\allowbreak{} ea /\allowbreak{} u /\allowbreak{} o. The second and third singular present have ie. The second singular preterite has the u-stem of the plural. Retain ā- throughout.
% mdparagraph:P084:end

% teaching: NEW-U16-PAR-ACEORFAN-IND
% table: NEW-U16-PAR-ACEORFAN-IND-DISPLAY
% mdtable:T35:begin
% course-table-style: 2026-10
% course-panel:T35-P01:begin
\par\begingroup
\fontsize{10.5}{13.2}\selectfont
\setlength{\LTleft}{0pt}\setlength{\LTright}{\fill}
\noindent
\begin{tabular}{@{}>{\raggedright\arraybackslash}p{39.459mm}>{\raggedright\arraybackslash}p{39.459mm}>{\raggedright\arraybackslash}p{39.459mm}@{}}
\toprule
\textbf{Person and number} & \textbf{Present} & \textbf{Preterite} \\
\midrule
1 sg. & āceorfe & ācearf \\
2 sg. & ācierfst & ācurfe \\
3 sg. & ācierfð & ācearf \\
1–3 pl. & āceorfað & ācurfon \\
\bottomrule
\end{tabular}
\par\endgroup
% course-panel:T35-P01:end
% mdtable:T35:end
\end{minipage}\par
% teaching: u16-forleosan-frame
\par\begin{minipage}{\linewidth}\raggedright
% mdparagraph:P085:begin
\textbf{forlēosan} — lose, with an accusative thing lost. Its participle can describe that thing as lost.
% mdparagraph:P085:end

% mdparagraph:P086:begin
Strong II: ēo /\allowbreak{} ēa /\allowbreak{} u /\allowbreak{} o. The preterite plural and past participle also change s to r. Both second and third singular present are forlīest; preterite second singular uses forlur-.
% mdparagraph:P086:end

% teaching: NEW-U16-PAR-FORLEOSAN-IND
% table: NEW-U16-PAR-FORLEOSAN-IND-DISPLAY
% mdtable:T36:begin
% course-table-style: 2026-10
% course-panel:T36-P01:begin
\par\begingroup
\fontsize{10.5}{13.2}\selectfont
\setlength{\LTleft}{0pt}\setlength{\LTright}{\fill}
\noindent
\begin{tabular}{@{}>{\raggedright\arraybackslash}p{39.459mm}>{\raggedright\arraybackslash}p{39.459mm}>{\raggedright\arraybackslash}p{39.459mm}@{}}
\toprule
\textbf{Person and number} & \textbf{Present} & \textbf{Preterite} \\
\midrule
1 sg. & forlēose & forlēas \\
2 sg. & forlīest & forlure \\
3 sg. & forlīest & forlēas \\
1–3 pl. & forlēosað & forluron \\
\bottomrule
\end{tabular}
\par\endgroup
% course-panel:T36-P01:end
% mdtable:T36:end
\end{minipage}\par
% teaching: u16-iernan-frame
\par\begin{minipage}{\linewidth}\raggedright
% mdparagraph:P087:begin
\textbf{iernan} — run, intransitive. An accusative actor with its plain infinitive can follow a verb of perception, as in the construction recalled below.
% mdparagraph:P087:end

% mdparagraph:P088:begin
Strong III. The selected stems are iern-, arn, urn- and urnen. Retain the position of r and the present ie spelling; do not replace it with a later irn- or yrn- series.
% mdparagraph:P088:end

% teaching: NEW-U16-PAR-IERNAN-IND
% table: NEW-U16-PAR-IERNAN-IND-DISPLAY
% mdtable:T37:begin
% course-table-style: 2026-10
% course-panel:T37-P01:begin
\par\begingroup
\fontsize{10.5}{13.2}\selectfont
\setlength{\LTleft}{0pt}\setlength{\LTright}{\fill}
\noindent
\begin{tabular}{@{}>{\raggedright\arraybackslash}p{39.459mm}>{\raggedright\arraybackslash}p{39.459mm}>{\raggedright\arraybackslash}p{39.459mm}@{}}
\toprule
\textbf{Person and number} & \textbf{Present} & \textbf{Preterite} \\
\midrule
1 sg. & ierne & arn \\
2 sg. & iernst & urne \\
3 sg. & iernð & arn \\
1–3 pl. & iernað & urnon \\
\bottomrule
\end{tabular}
\par\endgroup
% course-panel:T37-P01:end
% mdtable:T37:end
\end{minipage}\par
% teaching: u16-thwean-frame
\par\begin{minipage}{\linewidth}\raggedright
% mdparagraph:P089:begin
\textbf{ðwēan} — wash, with an accusative thing or body part. A possessor or recipient may retain the case independently required by the clause.
% mdparagraph:P089:end

% mdparagraph:P090:begin
Contracted strong VI, like slēan. The contracted present alternates ēa and ie. The selected preterite singular has h, while the second singular and plural have g; the past participle is ðwegen.
% mdparagraph:P090:end

% teaching: NEW-U16-PAR-THWEAN-IND
% table: NEW-U16-PAR-THWEAN-IND-DISPLAY
% mdtable:T38:begin
% course-table-style: 2026-10
% course-panel:T38-P01:begin
\par\begingroup
\fontsize{10.5}{13.2}\selectfont
\setlength{\LTleft}{0pt}\setlength{\LTright}{\fill}
\noindent
\begin{tabular}{@{}>{\raggedright\arraybackslash}p{39.459mm}>{\raggedright\arraybackslash}p{39.459mm}>{\raggedright\arraybackslash}p{39.459mm}@{}}
\toprule
\textbf{Person and number} & \textbf{Present} & \textbf{Preterite} \\
\midrule
1 sg. & ðwēa & ðwōh \\
2 sg. & ðwiehst & ðwōge \\
3 sg. & ðwiehð & ðwōh \\
1–3 pl. & ðwēað & ðwōgon \\
\bottomrule
\end{tabular}
\par\endgroup
% course-panel:T38-P01:end
% mdtable:T38:end
\end{minipage}\par
\reserve{7\baselineskip}
% mdparagraph:P091:begin
\subsection*{Further weak verbs}\BookHeading{nav-ch16-learner-0013}{subsubsection}{Further weak verbs}{0}
% mdparagraph:P091:end

% mdparagraph:P092:begin
The ordinary weak-I present uses -e, -st, -ð and plural -að, with the contractions shown below. Its dental preterite stem takes singular -e, -est, -e and plural -on. Learn the selected dental stem from the reference principal part and the complete preterite column.
% mdparagraph:P092:end

% teaching: u16-bewendan-frame
\par\begin{minipage}{\linewidth}\raggedright
% mdparagraph:P093:begin
\textbf{bewęndan} — turn. In the subject-directed movement used here, dative him can refer back to the subject: the subject turns toward someone. This particular dative reflexive does not make all reflexives dative.
% mdparagraph:P093:end

% mdparagraph:P094:begin
Weak I with a d-stem, like sęndan. The present contracts to -ntst and -nt. The dental preterite contracts to bewęnde: present first singular and preterite first/\allowbreak{}third singular therefore coincide in spelling.
% mdparagraph:P094:end

% teaching: NEW-U16-PAR-BEWENDAN-IND
% table: NEW-U16-PAR-BEWENDAN-IND-DISPLAY
% mdtable:T39:begin
% course-table-style: 2026-10
% course-panel:T39-P01:begin
\par\begingroup
\fontsize{10.5}{13.2}\selectfont
\setlength{\LTleft}{0pt}\setlength{\LTright}{\fill}
\noindent
\begin{tabular}{@{}>{\raggedright\arraybackslash}p{39.459mm}>{\raggedright\arraybackslash}p{39.459mm}>{\raggedright\arraybackslash}p{39.459mm}@{}}
\toprule
\textbf{Person and number} & \textbf{Present} & \textbf{Preterite} \\
\midrule
1 sg. & bewęnde & bewęnde \\
2 sg. & bewęntst & bewęndest \\
3 sg. & bewęnt & bewęnde \\
1–3 pl. & bewęndað & bewęndon \\
\bottomrule
\end{tabular}
\par\endgroup
% course-panel:T39-P01:end
% mdtable:T39:end
\end{minipage}\par
% teaching: u16-smierwan-frame
\par\begin{minipage}{\linewidth}\raggedright
% mdparagraph:P095:begin
\textbf{smierwan} — anoint. The body part or person anointed is accusative; mid with the dative introduces the substance used.
% mdparagraph:P095:end

% mdparagraph:P096:begin
Weak I, like gierwan. Retain w in present first singular and plural; omit it in second/\allowbreak{}third singular and in the preterite. The selected series has ie, including smierede; the singular command is smiere.
% mdparagraph:P096:end

% teaching: NEW-U16-PAR-SMIERWAN-IND
% table: NEW-U16-PAR-SMIERWAN-IND-DISPLAY
% mdtable:T40:begin
% course-table-style: 2026-10
% course-panel:T40-P01:begin
\par\begingroup
\fontsize{10.5}{13.2}\selectfont
\setlength{\LTleft}{0pt}\setlength{\LTright}{\fill}
\noindent
\begin{tabular}{@{}>{\raggedright\arraybackslash}p{39.459mm}>{\raggedright\arraybackslash}p{39.459mm}>{\raggedright\arraybackslash}p{39.459mm}@{}}
\toprule
\textbf{Person and number} & \textbf{Present} & \textbf{Preterite} \\
\midrule
1 sg. & smierwe & smierede \\
2 sg. & smierest & smieredest \\
3 sg. & smiereð & smierede \\
1–3 pl. & smierwað & smieredon \\
\bottomrule
\end{tabular}
\par\endgroup
% course-panel:T40-P01:end
% mdtable:T40:end
\end{minipage}\par
\par\begin{minipage}{\linewidth}\raggedright
% mdparagraph:P097:begin
\textbf{ālīesan} — release; deliver, with an accusative person or thing. The meaning and commands continue the preceding unit; the indicative is now complete.
% mdparagraph:P097:end

% mdparagraph:P098:begin
Weak I. Continue the selected ālīesan, ālīesde, ālīesed system. In the second and third singular present, contraction after s gives -st. Do not substitute a strong-verb system.
% mdparagraph:P098:end

% teaching: NEW-U16-PAR-ALIESAN-IND
% table: NEW-U16-PAR-ALIESAN-IND-DISPLAY
% mdtable:T41:begin
% course-table-style: 2026-10
% course-panel:T41-P01:begin
\par\begingroup
\fontsize{10.5}{13.2}\selectfont
\setlength{\LTleft}{0pt}\setlength{\LTright}{\fill}
\noindent
\begin{tabular}{@{}>{\raggedright\arraybackslash}p{39.459mm}>{\raggedright\arraybackslash}p{39.459mm}>{\raggedright\arraybackslash}p{39.459mm}@{}}
\toprule
\textbf{Person and number} & \textbf{Present} & \textbf{Preterite} \\
\midrule
1 sg. & ālīese & ālīesde \\
2 sg. & ālīest & ālīesdest \\
3 sg. & ālīest & ālīesde \\
1–3 pl. & ālīesað & ālīesdon \\
\bottomrule
\end{tabular}
\par\endgroup
% course-panel:T41-P01:end
% mdtable:T41:end
\end{minipage}\par
\reserve{18\baselineskip}
% mdparagraph:P099:begin
\subsection*{Commands}\BookHeading{nav-ch16-learner-0014}{subsubsection}{Commands}{0}
% mdparagraph:P099:end

% mdparagraph:P100:begin
The first form addresses one person; the second addresses more than one. Use ne before a negative command. The bare-stem commands preserve their own stem changes; weak-II blissian instead has singular blissa, and smierwan has singular smiere. These are commands, not present indicative forms inferred from a matching spelling.
% mdparagraph:P100:end

% teaching: NEW-U16-PAR-ACEORFAN-IMP
% table: NEW-U16-PAR-ACEORFAN-IMP-DISPLAY
% mdtable:T42:begin
% course-table-style: 2026-10
% course-panel:T42-P01:begin
\par\begingroup
\fontsize{10.5}{13.2}\selectfont
\setlength{\LTleft}{0pt}\setlength{\LTright}{\fill}
\noindent
\begin{tabular}{@{}>{\raggedright\arraybackslash}p{63.342mm}>{\raggedright\arraybackslash}p{57.846mm}@{}}
\toprule
\textbf{Person and number} & \textbf{Imperative} \\
\midrule
2 sg. & āceorf \\
2 pl. & āceorfað \\
\bottomrule
\end{tabular}
\par\endgroup
% course-panel:T42-P01:end
% mdtable:T42:end
% teaching: NEW-U16-PAR-BLISSIAN-IMP
% table: NEW-U16-PAR-BLISSIAN-IMP-DISPLAY
% mdtable:T43:begin
% course-table-style: 2026-10
% course-panel:T43-P01:begin
\par\begingroup
\fontsize{10.5}{13.2}\selectfont
\setlength{\LTleft}{0pt}\setlength{\LTright}{\fill}
\noindent
\begin{tabular}{@{}>{\raggedright\arraybackslash}p{63.342mm}>{\raggedright\arraybackslash}p{57.846mm}@{}}
\toprule
\textbf{Person and number} & \textbf{Imperative} \\
\midrule
2 sg. & blissa \\
2 pl. & blissiað \\
\bottomrule
\end{tabular}
\par\endgroup
% course-panel:T43-P01:end
% mdtable:T43:end
% teaching: NEW-U16-PAR-SMIERWAN-IMP
% table: NEW-U16-PAR-SMIERWAN-IMP-DISPLAY
% mdtable:T44:begin
% course-table-style: 2026-10
% course-panel:T44-P01:begin
\par\begingroup
\fontsize{10.5}{13.2}\selectfont
\setlength{\LTleft}{0pt}\setlength{\LTright}{\fill}
\noindent
\begin{tabular}{@{}>{\raggedright\arraybackslash}p{63.342mm}>{\raggedright\arraybackslash}p{57.846mm}@{}}
\toprule
\textbf{Person and number} & \textbf{Imperative} \\
\midrule
2 sg. & smiere \\
2 pl. & smierwað \\
\bottomrule
\end{tabular}
\par\endgroup
% course-panel:T44-P01:end
% mdtable:T44:end
% teaching: NEW-U16-PAR-ARISAN-IMP
% table: NEW-U16-PAR-ARISAN-IMP-DISPLAY
% mdtable:T45:begin
% course-table-style: 2026-10
% course-panel:T45-P01:begin
\par\begingroup
\fontsize{10.5}{13.2}\selectfont
\setlength{\LTleft}{0pt}\setlength{\LTright}{\fill}
\noindent
\begin{tabular}{@{}>{\raggedright\arraybackslash}p{63.342mm}>{\raggedright\arraybackslash}p{57.846mm}@{}}
\toprule
\textbf{Person and number} & \textbf{Imperative} \\
\midrule
2 sg. & ārīs \\
2 pl. & ārīsað \\
\bottomrule
\end{tabular}
\par\endgroup
% course-panel:T45-P01:end
% mdtable:T45:end
\par\begin{minipage}{\linewidth}\raggedright
% mdparagraph:P101:begin
\textbf{forlǣtan} — leave; let go, with the previously learned objects and sense.
% mdparagraph:P101:end

% teaching: NEW-U16-PAR-FORLAETAN-IMP
% table: NEW-U16-PAR-FORLAETAN-IMP-DISPLAY
% mdtable:T46:begin
% course-table-style: 2026-10
% course-panel:T46-P01:begin
\par\begingroup
\fontsize{10.5}{13.2}\selectfont
\setlength{\LTleft}{0pt}\setlength{\LTright}{\fill}
\noindent
\begin{tabular}{@{}>{\raggedright\arraybackslash}p{63.342mm}>{\raggedright\arraybackslash}p{57.846mm}@{}}
\toprule
\textbf{Person and number} & \textbf{Imperative} \\
\midrule
2 sg. & forlǣt \\
2 pl. & forlǣtað \\
\bottomrule
\end{tabular}
\par\endgroup
% course-panel:T46-P01:end
% mdtable:T46:end
\end{minipage}\par
\par\begin{minipage}{\linewidth}\raggedright
% mdparagraph:P102:begin
\textbf{geswīcan} — cease. It may be used without an expressed complement, including with a wind as subject or in an addressed command. Its genitive complement remains a separate use.
% mdparagraph:P102:end

% teaching: NEW-U16-PAR-GESWICAN-IMP
% table: NEW-U16-PAR-GESWICAN-IMP-DISPLAY
% mdtable:T47:begin
% course-table-style: 2026-10
% course-panel:T47-P01:begin
\par\begingroup
\fontsize{10.5}{13.2}\selectfont
\setlength{\LTleft}{0pt}\setlength{\LTright}{\fill}
\noindent
\begin{tabular}{@{}>{\raggedright\arraybackslash}p{63.342mm}>{\raggedright\arraybackslash}p{57.846mm}@{}}
\toprule
\textbf{Person and number} & \textbf{Imperative} \\
\midrule
2 sg. & geswīc \\
2 pl. & geswīcað \\
\bottomrule
\end{tabular}
\par\endgroup
% course-panel:T47-P01:end
% mdtable:T47:end
\end{minipage}\par
\par\begin{minipage}{\linewidth}\raggedright
% mdparagraph:P103:begin
\textbf{ęfthweorfan} — return; turn back. Retain the prefix in both commands.
% mdparagraph:P103:end

% teaching: NEW-U16-PAR-EFTHWEORFAN-IMP
% table: NEW-U16-PAR-EFTHWEORFAN-IMP-DISPLAY
% mdtable:T48:begin
% course-table-style: 2026-10
% course-panel:T48-P01:begin
\par\begingroup
\fontsize{10.5}{13.2}\selectfont
\setlength{\LTleft}{0pt}\setlength{\LTright}{\fill}
\noindent
\begin{tabular}{@{}>{\raggedright\arraybackslash}p{60.456mm}>{\raggedright\arraybackslash}p{60.732mm}@{}}
\toprule
\textbf{Person and number} & \textbf{Imperative} \\
\midrule
2 sg. & ęfthweorf \\
2 pl. & ęfthweorfað \\
\bottomrule
\end{tabular}
\par\endgroup
% course-panel:T48-P01:end
% mdtable:T48:end
\end{minipage}\par
\par\begin{minipage}{\linewidth}\raggedright
% mdparagraph:P104:begin
\textbf{dēman} — judge. A judgment itself may be expressed as an accusative object, as in the instruction to judge a right judgment.
% mdparagraph:P104:end

% teaching: NEW-U16-PAR-DEMAN-IMP
% table: NEW-U16-PAR-DEMAN-IMP-DISPLAY
% mdtable:T49:begin
% course-table-style: 2026-10
% course-panel:T49-P01:begin
\par\begingroup
\fontsize{10.5}{13.2}\selectfont
\setlength{\LTleft}{0pt}\setlength{\LTright}{\fill}
\noindent
\begin{tabular}{@{}>{\raggedright\arraybackslash}p{63.342mm}>{\raggedright\arraybackslash}p{57.846mm}@{}}
\toprule
\textbf{Person and number} & \textbf{Imperative} \\
\midrule
2 sg. & dēm \\
2 pl. & dēmað \\
\bottomrule
\end{tabular}
\par\endgroup
% course-panel:T49-P01:end
% mdtable:T49:end
\end{minipage}\par
% mdparagraph:P105:begin
\textbf{Immediately before gē}, use the shortened plural command \textbf{dēme gē}. Without that following pronoun the ordinary plural remains \textbf{dēmað}. This positional exception is required in the first clause of Cook 36.1; it does not replace the ordinary pair or introduce a new general command table.\textsuperscript{4}
% mdparagraph:P105:end

\reserve{18\baselineskip}
% mdparagraph:P106:begin
\subsection*{Present subjunctives: expectation and an unspecified provider}\BookHeading{nav-ch16-learner-0015}{subsubsection}{Present subjunctives: expectation and an unspecified provider}{0}
% mdparagraph:P106:end

% mdparagraph:P107:begin
Findan “find” and fēdan “feed” retain their previously learned indicative systems. The complete present subjunctive has -e in every singular person and -en in every plural person. The six positions remain distinct even when the forms coincide.
% mdparagraph:P107:end

% teaching: NEW-U16-PAR-FINDAN-PRES-SUBJ
% table: NEW-U16-PAR-FINDAN-PRES-SUBJ-DISPLAY
% mdtable:T50:begin
% course-table-style: 2026-10
% course-panel:T50-P01:begin
\par\begingroup
\fontsize{10.5}{13.2}\selectfont
\setlength{\LTleft}{0pt}\setlength{\LTright}{\fill}
\noindent
\begin{tabular}{@{}>{\raggedright\arraybackslash}p{63.342mm}>{\raggedright\arraybackslash}p{57.846mm}@{}}
\toprule
\textbf{Person and number} & \textbf{Present} \\
\midrule
1 sg. & finde \\
2 sg. & finde \\
3 sg. & finde \\
1–3 pl. & finden \\
\bottomrule
\end{tabular}
\par\endgroup
% course-panel:T50-P01:end
% mdtable:T50:end
% teaching: NEW-U16-PAR-FEDAN-PRES-SUBJ
% table: NEW-U16-PAR-FEDAN-PRES-SUBJ-DISPLAY
% mdtable:T51:begin
% course-table-style: 2026-10
% course-panel:T51-P01:begin
\par\begingroup
\fontsize{10.5}{13.2}\selectfont
\setlength{\LTleft}{0pt}\setlength{\LTright}{\fill}
\noindent
\begin{tabular}{@{}>{\raggedright\arraybackslash}p{63.342mm}>{\raggedright\arraybackslash}p{57.846mm}@{}}
\toprule
\textbf{Person and number} & \textbf{Present} \\
\midrule
1 sg. & fēde \\
2 sg. & fēde \\
3 sg. & fēde \\
1–3 pl. & fēden \\
\bottomrule
\end{tabular}
\par\endgroup
% course-panel:T51-P01:end
% mdtable:T51:end
% mdparagraph:P108:begin
\textbf{Wēn is ðæt}\BookPoint{bk-ch16-learner-71e2b7958144} with a present subjunctive expresses the bounded expectation “perhaps; there is a prospect that …”. In the hypothetical conditional “if you find no one”, use \textbf{finde}; a relative clause identifying a possible, as-yet-unfound provider uses \textbf{fēde}. These choices concern the status of the situation, not a mechanical conversion of an English future.
% mdparagraph:P108:end

% mdparagraph:P109:begin
Keep the previously selected preterites of findan distinct: first singular \textbf{fand}, third singular \textbf{funde}. A spelling shared with another person or mood must be interpreted from its clause.
% mdparagraph:P109:end

\reserve{7\baselineskip}
% mdparagraph:P110:begin
\subsection*{Participles and infinitives}\BookHeading{nav-ch16-learner-0016}{subsubsection}{Participles and infinitives}{0}
% mdparagraph:P110:end

% mdparagraph:P111:begin
The following pairs give the productive participial bases. These present participles add -ende to the present stem and describe ongoing action. Learn each past-participle stem and prefix as shown. The past participle can enter the previously learned passive or perfect constructions; in a passive, the affected participant is the subject. These bases do not themselves supply every possible adjective inflection.
% mdparagraph:P111:end

\reserve{10\baselineskip}
% mdparagraph:P112:begin
Forlēosan retains the r-stem in its past participle. In the strong neuter accusative singular, \textbf{forloren} can modify a lost thing; this exact attributive position is also learned here.
% mdparagraph:P112:end

% teaching: NEW-U16-PAR-FORLEOSAN-PART
% table: NEW-U16-PAR-FORLEOSAN-PART-DISPLAY
% mdtable:T52:begin
\par\begingroup
\fontsize{10.5}{13.2}\selectfont
\setlength{\LTleft}{0pt}\setlength{\LTright}{\fill}
\noindent
\begin{tabular}{@{}>{\raggedright\arraybackslash}p{69.12mm}>{\raggedright\arraybackslash}p{52.07mm}@{}}
\toprule
\textbf{Position} & \textbf{Form} \\
\midrule
participle, present & forlēosende \\
participle, past & forloren \\
\bottomrule
\end{tabular}
\par\endgroup

% mdtable:T52:end
% teaching: u16-laedan-passive-carry
\reserve{10\baselineskip}
% mdparagraph:P113:begin
Lǣdan here includes “carry; take a person”. A passive makes the person carried its subject. The endpoint in the assigned carrying event uses ōð with the accusative.
% mdparagraph:P113:end

% teaching: NEW-U16-PAR-LAEDAN-PART
% table: NEW-U16-PAR-LAEDAN-PART-DISPLAY
% mdtable:T53:begin
\par\begingroup
\fontsize{10.5}{13.2}\selectfont
\setlength{\LTleft}{0pt}\setlength{\LTright}{\fill}
\noindent
\begin{tabular}{@{}>{\raggedright\arraybackslash}p{71.92mm}>{\raggedright\arraybackslash}p{49.26mm}@{}}
\toprule
\textbf{Position} & \textbf{Form} \\
\midrule
participle, present & lǣdende \\
participle, past & lǣded \\
\bottomrule
\end{tabular}
\par\endgroup

% mdtable:T53:end
% teaching: u16-gefyllan-utterance
\reserve{10\baselineskip}
% mdparagraph:P114:begin
Gefyllan here means “fulfil a spoken prediction”. That prediction is accusative in an active clause and becomes the subject in a passive. This is distinct from physically filling a container.
% mdparagraph:P114:end

% teaching: NEW-U16-PAR-GEFYLLAN-PART
% table: NEW-U16-PAR-GEFYLLAN-PART-DISPLAY
% mdtable:T54:begin
\par\begingroup
\fontsize{10.5}{13.2}\selectfont
\setlength{\LTleft}{0pt}\setlength{\LTright}{\fill}
\noindent
\begin{tabular}{@{}>{\raggedright\arraybackslash}p{71.92mm}>{\raggedright\arraybackslash}p{49.26mm}@{}}
\toprule
\textbf{Position} & \textbf{Form} \\
\midrule
participle, present & gefyllende \\
participle, past & gefylled \\
\bottomrule
\end{tabular}
\par\endgroup

% mdtable:T54:end
% teaching: u16-forgiefan-forgive
\reserve{11\baselineskip}
% mdparagraph:P115:begin
Forgiefan here means “forgive”: the sin or offence is accusative, the person forgiven dative. In a passive the sin becomes nominative and the person remains dative. The uninflected past participle is \textbf{forgiefen} in the assigned predicative use.
% mdparagraph:P115:end

% teaching: NEW-U16-PAR-FORGIEFAN-PART
% table: NEW-U16-PAR-FORGIEFAN-PART-DISPLAY
% mdtable:T55:begin
\par\begingroup
\fontsize{10.5}{13.2}\selectfont
\setlength{\LTleft}{0pt}\setlength{\LTright}{\fill}
\noindent
\begin{tabular}{@{}>{\raggedright\arraybackslash}p{69.12mm}>{\raggedright\arraybackslash}p{52.07mm}@{}}
\toprule
\textbf{Position} & \textbf{Form} \\
\midrule
participle, present & forgiefende \\
participle, past & forgiefen \\
\bottomrule
\end{tabular}
\par\endgroup

% mdtable:T55:end
\reserve{13\baselineskip}
% mdparagraph:P116:begin
The plain infinitive and the inflected infinitive after tō are separate forms. A perception construction takes the plain infinitive and an accusative actor. The previously learned purpose construction takes tō and the inflected form. A new meaning such as “is to come” still requires separate local help; it does not follow merely from knowing the ending.
% mdparagraph:P116:end

% teaching: NEW-U16-PAR-IERNAN-INF
% table: NEW-U16-PAR-IERNAN-INF-DISPLAY
% mdtable:T56:begin
\par\begingroup
\fontsize{10.5}{13.2}\selectfont
\setlength{\LTleft}{0pt}\setlength{\LTright}{\fill}
\noindent
\begin{tabular}{@{}>{\raggedright\arraybackslash}p{68.04mm}>{\raggedright\arraybackslash}p{53.15mm}@{}}
\toprule
\textbf{Position} & \textbf{Form} \\
\midrule
infinitive & iernan \\
infinitive, dat & iernanne \\
\bottomrule
\end{tabular}
\par\endgroup

% mdtable:T56:end
% teaching: NEW-U16-PAR-SENDAN-INF
% table: NEW-U16-PAR-SENDAN-INF-DISPLAY
% mdtable:T57:begin
\par\begingroup
\fontsize{10.5}{13.2}\selectfont
\setlength{\LTleft}{0pt}\setlength{\LTright}{\fill}
\noindent
\begin{tabular}{@{}>{\raggedright\arraybackslash}p{68.04mm}>{\raggedright\arraybackslash}p{53.15mm}@{}}
\toprule
\textbf{Position} & \textbf{Form} \\
\midrule
infinitive & sęndan \\
infinitive, dat & sęndanne \\
\bottomrule
\end{tabular}
\par\endgroup

% mdtable:T57:end
\reserve{7\baselineskip}
% mdparagraph:P117:begin
\section*{3. Clause structure, government and viewpoint}\BookHeading{nav-ch16-learner-0017}{subsection}{3. Clause structure, government and viewpoint}{1}
% mdparagraph:P117:end

\reserve{7\baselineskip}
% mdparagraph:P118:begin
\subsection*{Reference and finite content}\BookHeading{nav-ch16-learner-0018}{subsubsection}{Reference and finite content}{0}
% mdparagraph:P118:end

% mdparagraph:P119:begin
A pronoun does not automatically refer to the nearest noun. Establish the speaker, addressee and continuing participant from the whole passage. A pronoun’s ending can establish its role without uniquely establishing its referent. When a new composition requires an unambiguous reference, an explicit noun phrase can resolve it.
% mdparagraph:P119:end

% teaching: NEW-U16-EXAMPLE-01
\reserve{8\baselineskip}
% example: NEW-U16-EXAMPLE-01
% mdparagraph:P120:begin
Se fæder sægde: ‘Ic eom hēr.’ Se fæder sægde ðæt hē ðǣr wǣre.
% mdparagraph:P120:end

% mdparagraph:P121:begin
The father said, 'I am here.' The father said that he was there.
% mdparagraph:P121:end

% mdparagraph:P122:begin
The\BookPoint{bk-ch16-learner-531548ae175f} quotation uses the father’s I and here. The report uses third-person hē and, under the stated change of reporting place, ðǣr. The subjunctive wǣre belongs to reported content; it does not mean that every English past verb must be replaced by an Old English subjunctive.
% mdparagraph:P122:end

% teaching: NEW-U16-EXAMPLE-02
\reserve{8\baselineskip}
% example: NEW-U16-EXAMPLE-02
% mdparagraph:P123:begin
Hē geseah ðone mann cuman. Hē wiste ðæt se mann cōme.
% mdparagraph:P123:end

% mdparagraph:P124:begin
He saw the man come. He knew that the man came.
% mdparagraph:P124:end

% mdparagraph:P125:begin
With\BookPoint{bk-ch16-learner-88ba24c6485d} perception and a plain infinitive, the actor of the perceived action is accusative. A finite ðæt-clause has its own nominative subject and finite verb. Some finite spellings can represent more than one mood; identify only what the ending and context together warrant.
% mdparagraph:P125:end

% mdparagraph:P126:begin
An earlier speech event does not make every situation in its content earlier still. A reported state can be simultaneous with the saying; habban with a past participle can mark an already completed event; willan with an infinitive can express an intention whose action is still ahead. Interpret the whole verb group and the stated context. Directly quoted now or today retains the quoted speaker’s viewpoint.
% mdparagraph:P126:end

\reserve{7\baselineskip}
% mdparagraph:P127:begin
\subsection*{Knowing, action comparison and causation}\BookHeading{nav-ch16-learner-0019}{subsubsection}{Knowing, action comparison and causation}{0}
% mdparagraph:P127:end

% teaching: u16-cunnan-person
% teaching: u16-witan-person
% mdparagraph:P128:begin
\textbf{Cunnan} can mean “know; be acquainted with” a person and take that person in the accusative. Its earlier use with a plain infinitive expresses knowing how to do something. \textbf{Witan} can introduce factual content or a question, but can also take a personal accusative object. Thus the contrast is useful without making persons and facts an exclusive division between the two verbs.\textsuperscript{5}
% mdparagraph:P128:end

% teaching: NEW-U16-EXAMPLE-03
\reserve{8\baselineskip}
% example: NEW-U16-EXAMPLE-03
% mdparagraph:P129:begin
Gē cunnon ðone cyning; gē cunnon singan.
% mdparagraph:P129:end

% mdparagraph:P130:begin
You know the king; you know how to sing.
% mdparagraph:P130:end

% teaching: u16-don-substitute
% teaching: u16-don-finite-result
% mdparagraph:P131:begin
\textbf{Dōn}\BookPoint{bk-ch16-learner-6a25e97d3147} can repeat an action already expressed. In an action comparison, \textbf{bęt} “better” modifies the action and \textbf{ðonne} introduces the comparison clause. Dōn then carries the tense, person and number of its own subject. It is not an auxiliary inserted to imitate modern English questions or negatives.
% mdparagraph:P131:end

% teaching: NEW-U16-EXAMPLE-04
\reserve{8\baselineskip}
% example: NEW-U16-EXAMPLE-04
% mdparagraph:P132:begin
Hē plegode wel, and hīe dydon swā.
% mdparagraph:P132:end

% mdparagraph:P133:begin
He played well, and they did so.
% mdparagraph:P133:end

% mdparagraph:P134:begin
A\BookPoint{bk-ch16-learner-b481bfab43f6} separate causative pattern is \textbf{dōn + ðæt + finite clause}, “cause that; make it be that”. The subject inside that clause remains nominative. With dual \textbf{git}, use a plural finite verb, such as \textbf{bēoð}; a noun describing both participants is plural. Do not turn the dual subject into an accusative object of dōn in this pattern.
% mdparagraph:P134:end

\reserve{7\baselineskip}
% mdparagraph:P135:begin
\subsection*{Government and lexical meaning}\BookHeading{nav-ch16-learner-0020}{subsubsection}{Government and lexical meaning}{0}
% mdparagraph:P135:end

% mdparagraph:P136:begin
A case after a verb belongs to that verb’s construction. It is not determined by the nearest English preposition.
% mdparagraph:P136:end

% mdparagraph:P137:begin
\textbf{Biddan}\BookPoint{bk-ch16-learner-9b2cc4c38e6b} takes the requested thing in the genitive; the person asked can simultaneously be accusative. Interrogative \textbf{hwæs} asks for that requested thing. \textbf{Gīeman} “take care of” and \textbf{geswīcan} “cease from” retain their learned genitive government. \textbf{Wēnan} “expect” now also takes a genitive thing expected, including a genitive negative or a coordinated exception. Its earlier finite content-clause use does not replace this pattern.
% mdparagraph:P137:end

% teaching: u16-halsian-oath
% mdparagraph:P138:begin
\textbf{Hālsian} “conjure; solemnly charge” takes the person charged in the accusative. \textbf{Ðurh} with an accusative names the oath standard, and a requested \textbf{ðæt}-clause follows. An uncertain embedded identification uses a subjunctive. This person is not an ordinary dative recipient.
% mdparagraph:P138:end

% teaching: u16-healdan-keep
% teaching: u16-weorthan-arise
% mdparagraph:P139:begin
\textbf{Healdan} now includes “keep; observe” a saying, with accusative sprǣc. \textbf{Weorðan} can mean “arise; come about” with an event as its subject. If a calm comes about and no participle is present, this is an intransitive lexical use, not the passive auxiliary.
% mdparagraph:P139:end

% mdparagraph:P140:begin
With\BookPoint{bk-ch16-learner-9a119a3340b1} \textbf{bewęndan}, the dative reflexive can identify the person who turns toward someone. With \textbf{gān}, the genitive route pattern \textbf{ðæs weges ðe …} means “by the way by which …”. Neither construction should be generalized to every motion verb. In Cook 36.10, \textbf{āsęndan} has no expressed patient: understand “sent someone”, followed by the coordinated action.
% mdparagraph:P140:end

% teaching: u16-ofer-rejoicing
% teaching: u16-degree-comparison
% teaching: u16-of-particle
% mdparagraph:P141:begin
With \textbf{blissian}, \textbf{ofer + accusative} names the object of rejoicing. \textbf{Swīðor} here compares its degree, “more”, and \textbf{ðonne} introduces the second object of comparison, with its own ofer phrase. It is not the earlier “rather” meaning. The separate particle \textbf{of}, “off”, has no following governed complement; prepositional \textbf{of} still governs the dative in its own uses.
% mdparagraph:P141:end

\reserve{7\baselineskip}
% mdparagraph:P142:begin
\subsection*{Direction, relation and time words}\BookHeading{nav-ch16-learner-0021}{subsubsection}{Direction, relation and time words}{0}\BookLabel{bk-ch16-learner-5b59859f9b28}
% mdparagraph:P142:end

% teaching: NEW-U16-VOC-HWANON
% teaching: NEW-U16-VOC-HWIDER
% teaching: NEW-U16-VOC-TOGAEDERE
% teaching: NEW-U16-VOC-GEOND
% teaching: NEW-U16-VOC-INNAN
% teaching: NEW-U16-VOC-SITHTHAN
% teaching: NEW-U16-VOC-THRIWA
% teaching: NEW-U16-VOC-SONA
% teaching: NEW-U16-VOC-AEFTER
% teaching: u16-aefter-follow
% teaching: NEW-U16-VOC-INTO
% teaching: NEW-U16-VOC-BE
% teaching: NEW-U16-VOC-THURH
% teaching: NEW-U16-VOC-FOR-THON-THE
% mdtable:T58:begin
\par\begingroup
\fontsize{10.5}{13.2}\selectfont
\setlength{\LTleft}{0pt}\setlength{\LTright}{\fill}
\noindent
\begin{tabular}{@{}>{\raggedright\arraybackslash}p{18.5mm}>{\raggedright\arraybackslash}p{26.31mm}>{\raggedright\arraybackslash}p{73.57mm}@{}}
\toprule
\textbf{Word} & \textbf{Meaning} & \textbf{Use} \\
\midrule
hwanon & whence; from where & Interrogative adverb of origin. \\
hwider & whither; to where & Interrogative adverb of destination. \\
tōgædere & together & Adverb. \\
geond & along; through & Preposition with accusative. \\
innan & within; inside & Adverb in the assigned use. \\
siððan & since; after; afterwards & As a conjunction, introduces a temporal clause; as an adverb, means afterwards. \\
ðrīwa & thrice; three times & Adverb. \\
sōna & immediately; at once & Adverb; a complete as-soon-as construction is separately supplied where required. \\
æfter & after; following; according to & Dative in the three stated uses: time after an event, following a person, or conformity to something. \\
intō & into & Destination with dative in this selected use. \\
be & according to; concerning & Dative in both senses. The reading’s “by the hand” receives separate local help. \\
ðurh & through; by way of a person & Accusative for the mediating person. Oath use is separately identified; source-local means or evidence receives help. \\
for-ðon & therefore; for that reason & Adverbial consequence here. Do not expand it into a because-clause merely from the related longer expression. \\
\bottomrule
\end{tabular}
\par\endgroup

% mdtable:T58:end
% mdparagraph:P143:begin
The following exercises reuse the previously learned masculine \textbf{sǣ}. Preserve the appropriate masculine article or demonstrative with it. \textbf{Bōc} retains its learned feminine plural \textbf{bēc}. Source variants in the reading are supplied locally and do not replace the forms learned for new composition.
% mdparagraph:P143:end

\reserve{7\baselineskip}
% mdparagraph:P144:begin
\section*{4. Form and clause practice}\BookHeading{nav-ch16-learner-0022}{subsection}{4. Form and clause practice}{1}
% mdparagraph:P144:end

% teaching: NEW-U16-DIRECTION-HEADWORDS
% mdparagraph:P145:begin
The dictionary forms printed with the form exercises are given references. Produce the requested inflections from the learned systems. If a requested position has the same form as its given headword, identify that position; it is not a test of recalling an unseen word. Explain the requested contrasts in English.
% mdparagraph:P145:end

\reserve{5\baselineskip}
\par\begin{minipage}{\linewidth}\raggedright
% mdparagraph:P146:begin
\textbf{F01}
% mdparagraph:P146:end

% item: NEW-U16-F01
% mdparagraph:P147:begin
Give the selected forms requested below. For each verb, explain the stem changes linking your three answers. Retain the prefixes and the selected spelling.
% mdparagraph:P147:end

% mdparagraph:P148:begin
(1) āceorfan
% mdparagraph:P148:end

% mdparagraph:P149:begin
(2) forlēosan
% mdparagraph:P149:end

% mdparagraph:P150:begin
(3) ðwēan
% mdparagraph:P150:end

% mdparagraph:P151:begin
(4) iernan
% mdparagraph:P151:end

% mdparagraph:P152:begin
\textbf{Supplied words.} These forms and meanings remain available for this task.
% mdparagraph:P152:end

% mdtable:T59:begin
\par\begingroup
\fontsize{10.5}{13.2}\selectfont
\setlength{\LTleft}{0pt}\setlength{\LTright}{\fill}
\noindent
\begin{tabular}{@{}>{\raggedright\arraybackslash}p{33.87mm}>{\raggedright\arraybackslash}p{33.87mm}>{\raggedright\arraybackslash}p{50.64mm}@{}}
\toprule
\textbf{Form} & \textbf{Meaning} & \textbf{Role} \\
\midrule
āceorfan & cut off & given dictionary form \\
forlēosan & lose & given dictionary form \\
ðwēan & wash & given dictionary form \\
\bottomrule
\end{tabular}
\par\endgroup

% mdtable:T59:end
% output: NEW-U16-F01-O001
% mdparagraph:P153:begin
(1) āceorfan: present indicative, you singular: \_\_\_\_\_\_\_\_\_\_\_\_\_\_\_\_\_\_\_\_
% mdparagraph:P153:end

% output: NEW-U16-F01-O002
% mdparagraph:P154:begin
(2) āceorfan: preterite indicative, you singular: \_\_\_\_\_\_\_\_\_\_\_\_\_\_\_\_\_\_\_\_
% mdparagraph:P154:end

% output: NEW-U16-F01-O003
% mdparagraph:P155:begin
(3) āceorfan: preterite indicative, they: \_\_\_\_\_\_\_\_\_\_\_\_\_\_\_\_\_\_\_\_
% mdparagraph:P155:end

% output: NEW-U16-F01-O004
% mdparagraph:P156:begin
(4) āceorfan: explain the stem relations: \_\_\_\_\_\_\_\_\_\_\_\_\_\_\_\_\_\_\_\_
% mdparagraph:P156:end

% output: NEW-U16-F01-O005
% mdparagraph:P157:begin
(5) forlēosan: present indicative, he: \_\_\_\_\_\_\_\_\_\_\_\_\_\_\_\_\_\_\_\_
% mdparagraph:P157:end

% output: NEW-U16-F01-O006
% mdparagraph:P158:begin
(6) forlēosan: preterite indicative, you singular: \_\_\_\_\_\_\_\_\_\_\_\_\_\_\_\_\_\_\_\_
% mdparagraph:P158:end

% output: NEW-U16-F01-O007
% mdparagraph:P159:begin
(7) forlēosan: preterite indicative, they: \_\_\_\_\_\_\_\_\_\_\_\_\_\_\_\_\_\_\_\_
% mdparagraph:P159:end

% output: NEW-U16-F01-O008
% mdparagraph:P160:begin
(8) forlēosan: explain the stem relations: \_\_\_\_\_\_\_\_\_\_\_\_\_\_\_\_\_\_\_\_
% mdparagraph:P160:end

% output: NEW-U16-F01-O009
% mdparagraph:P161:begin
(9) ðwēan: present indicative, he: \_\_\_\_\_\_\_\_\_\_\_\_\_\_\_\_\_\_\_\_
% mdparagraph:P161:end

% output: NEW-U16-F01-O010
% mdparagraph:P162:begin
(10) ðwēan: preterite indicative, you singular: \_\_\_\_\_\_\_\_\_\_\_\_\_\_\_\_\_\_\_\_
% mdparagraph:P162:end

% output: NEW-U16-F01-O011
% mdparagraph:P163:begin
(11) ðwēan: preterite indicative, they: \_\_\_\_\_\_\_\_\_\_\_\_\_\_\_\_\_\_\_\_
% mdparagraph:P163:end

% output: NEW-U16-F01-O012
% mdparagraph:P164:begin
(12) ðwēan: explain the stem relations: \_\_\_\_\_\_\_\_\_\_\_\_\_\_\_\_\_\_\_\_
% mdparagraph:P164:end

% output: NEW-U16-F01-O013
% mdparagraph:P165:begin
(13) iernan: present indicative, he: \_\_\_\_\_\_\_\_\_\_\_\_\_\_\_\_\_\_\_\_
% mdparagraph:P165:end

% output: NEW-U16-F01-O014
% mdparagraph:P166:begin
(14) iernan: preterite indicative, he: \_\_\_\_\_\_\_\_\_\_\_\_\_\_\_\_\_\_\_\_
% mdparagraph:P166:end

% output: NEW-U16-F01-O015
% mdparagraph:P167:begin
(15) iernan: preterite indicative, they: \_\_\_\_\_\_\_\_\_\_\_\_\_\_\_\_\_\_\_\_
% mdparagraph:P167:end

% output: NEW-U16-F01-O016
% mdparagraph:P168:begin
(16) iernan: explain the stem relations: \_\_\_\_\_\_\_\_\_\_\_\_\_\_\_\_\_\_\_\_
% mdparagraph:P168:end

\end{minipage}\par\addvspace{4pt}

\par\begin{minipage}{\linewidth}\raggedright
% mdparagraph:P169:begin
\textbf{F02}
% mdparagraph:P169:end

% item: NEW-U16-F02
% mdparagraph:P170:begin
Give the selected forms requested below. Then explain the two stated contrasts. These are separate positions even when their spelling coincides.
% mdparagraph:P170:end

% mdparagraph:P171:begin
(1) smierwan
% mdparagraph:P171:end

% mdparagraph:P172:begin
(2) bewęndan
% mdparagraph:P172:end

% mdparagraph:P173:begin
(3) plegian
% mdparagraph:P173:end

% mdparagraph:P174:begin
(4) behēafdian
% mdparagraph:P174:end

% mdparagraph:P175:begin
(5) gefyllan
% mdparagraph:P175:end

% mdparagraph:P176:begin
(6) forgiefan
% mdparagraph:P176:end

% mdparagraph:P177:begin
\textbf{Supplied words.} These forms and meanings remain available for this task.
% mdparagraph:P177:end

% mdtable:T60:begin
\par\begingroup
\fontsize{10.5}{13.2}\selectfont
\setlength{\LTleft}{0pt}\setlength{\LTright}{\fill}
\noindent
\begin{tabular}{@{}>{\raggedright\arraybackslash}p{33.87mm}>{\raggedright\arraybackslash}p{33.87mm}>{\raggedright\arraybackslash}p{50.64mm}@{}}
\toprule
\textbf{Form} & \textbf{Meaning} & \textbf{Role} \\
\midrule
smierwan & anoint & given dictionary form \\
bewęndan & turn & given dictionary form \\
plegian & play & given dictionary form \\
behēafdian & behead & given dictionary form \\
forgiefan & forgive & given dictionary form \\
\bottomrule
\end{tabular}
\par\endgroup

% mdtable:T60:end
% output: NEW-U16-F02-O001
% mdparagraph:P178:begin
(1) smierwan: present indicative, I: \_\_\_\_\_\_\_\_\_\_\_\_\_\_\_\_\_\_\_\_
% mdparagraph:P178:end

% output: NEW-U16-F02-O002
% mdparagraph:P179:begin
(2) smierwan: present indicative, you singular: \_\_\_\_\_\_\_\_\_\_\_\_\_\_\_\_\_\_\_\_
% mdparagraph:P179:end

% output: NEW-U16-F02-O003
% mdparagraph:P180:begin
(3) smierwan: preterite indicative, she: \_\_\_\_\_\_\_\_\_\_\_\_\_\_\_\_\_\_\_\_
% mdparagraph:P180:end

% output: NEW-U16-F02-O004
% mdparagraph:P181:begin
(4) bewęndan: present indicative, I: \_\_\_\_\_\_\_\_\_\_\_\_\_\_\_\_\_\_\_\_
% mdparagraph:P181:end

% output: NEW-U16-F02-O005
% mdparagraph:P182:begin
(5) bewęndan: preterite indicative, he: \_\_\_\_\_\_\_\_\_\_\_\_\_\_\_\_\_\_\_\_
% mdparagraph:P182:end

% output: NEW-U16-F02-O006
% mdparagraph:P183:begin
(6) plegian: preterite indicative, they: \_\_\_\_\_\_\_\_\_\_\_\_\_\_\_\_\_\_\_\_
% mdparagraph:P183:end

% output: NEW-U16-F02-O007
% mdparagraph:P184:begin
(7) behēafdian: preterite indicative, they: \_\_\_\_\_\_\_\_\_\_\_\_\_\_\_\_\_\_\_\_
% mdparagraph:P184:end

% output: NEW-U16-F02-O008
% mdparagraph:P185:begin
(8) gefyllan: uninflected past participle: \_\_\_\_\_\_\_\_\_\_\_\_\_\_\_\_\_\_\_\_
% mdparagraph:P185:end

% output: NEW-U16-F02-O009
% mdparagraph:P186:begin
(9) forgiefan: uninflected past participle: \_\_\_\_\_\_\_\_\_\_\_\_\_\_\_\_\_\_\_\_
% mdparagraph:P186:end

% output: NEW-U16-F02-O010
% mdparagraph:P187:begin
(10) Explain the smierwan stem distribution shown by your answers: \_\_\_\_\_\_\_\_\_\_\_\_\_\_\_\_\_\_\_\_
% mdparagraph:P187:end

% output: NEW-U16-F02-O011
% mdparagraph:P188:begin
(11) Explain what the two bewęndan answers establish about identifying tense from spelling alone: \_\_\_\_\_\_\_\_\_\_\_\_\_\_\_\_\_\_\_\_
% mdparagraph:P188:end

\end{minipage}\par\addvspace{4pt}

\par\begin{minipage}{\linewidth}\raggedright
% mdparagraph:P189:begin
\textbf{F03}
% mdparagraph:P189:end

% item: NEW-U16-F03
% mdparagraph:P190:begin
For each listed verb, give a command to one person and then to more than one. Use the selected command pair taught in this unit.
% mdparagraph:P190:end

% mdparagraph:P191:begin
(1) āceorfan
% mdparagraph:P191:end

% mdparagraph:P192:begin
(2) blissian
% mdparagraph:P192:end

% mdparagraph:P193:begin
(3) smierwan
% mdparagraph:P193:end

% mdparagraph:P194:begin
(4) ārīsan
% mdparagraph:P194:end

% mdparagraph:P195:begin
(5) forlǣtan
% mdparagraph:P195:end

% mdparagraph:P196:begin
(6) geswīcan
% mdparagraph:P196:end

% mdparagraph:P197:begin
(7) ęfthweorfan
% mdparagraph:P197:end

% mdparagraph:P198:begin
\textbf{Supplied words.} These forms and meanings remain available for this task.
% mdparagraph:P198:end

% mdtable:T61:begin
\par\begingroup
\fontsize{10.5}{13.2}\selectfont
\setlength{\LTleft}{0pt}\setlength{\LTright}{\fill}
\noindent
\begin{tabular}{@{}>{\raggedright\arraybackslash}p{34.27mm}>{\raggedright\arraybackslash}p{37.54mm}>{\raggedright\arraybackslash}p{46.57mm}@{}}
\toprule
\textbf{Form} & \textbf{Meaning} & \textbf{Role} \\
\midrule
āceorfan & cut off & given dictionary form \\
blissian & rejoice & given dictionary form \\
smierwan & anoint & given dictionary form \\
ārīsan & arise & given dictionary form \\
forlǣtan & abandon; forsake & given dictionary form \\
geswīcan & cease & given dictionary form \\
ęfthweorfan & return; turn back & given dictionary form \\
\bottomrule
\end{tabular}
\par\endgroup

% mdtable:T61:end
% output: NEW-U16-F03-O001
% mdparagraph:P199:begin
(1) āceorfan: one addressee: \_\_\_\_\_\_\_\_\_\_\_\_\_\_\_\_\_\_\_\_
% mdparagraph:P199:end

% output: NEW-U16-F03-O002
% mdparagraph:P200:begin
(2) āceorfan: more than one addressee: \_\_\_\_\_\_\_\_\_\_\_\_\_\_\_\_\_\_\_\_
% mdparagraph:P200:end

% output: NEW-U16-F03-O003
% mdparagraph:P201:begin
(3) blissian: one addressee: \_\_\_\_\_\_\_\_\_\_\_\_\_\_\_\_\_\_\_\_
% mdparagraph:P201:end

% output: NEW-U16-F03-O004
% mdparagraph:P202:begin
(4) blissian: more than one addressee: \_\_\_\_\_\_\_\_\_\_\_\_\_\_\_\_\_\_\_\_
% mdparagraph:P202:end

% output: NEW-U16-F03-O005
% mdparagraph:P203:begin
(5) smierwan: one addressee: \_\_\_\_\_\_\_\_\_\_\_\_\_\_\_\_\_\_\_\_
% mdparagraph:P203:end

% output: NEW-U16-F03-O006
% mdparagraph:P204:begin
(6) smierwan: more than one addressee: \_\_\_\_\_\_\_\_\_\_\_\_\_\_\_\_\_\_\_\_
% mdparagraph:P204:end

% output: NEW-U16-F03-O007
% mdparagraph:P205:begin
(7) ārīsan: one addressee: \_\_\_\_\_\_\_\_\_\_\_\_\_\_\_\_\_\_\_\_
% mdparagraph:P205:end

% output: NEW-U16-F03-O008
% mdparagraph:P206:begin
(8) ārīsan: more than one addressee: \_\_\_\_\_\_\_\_\_\_\_\_\_\_\_\_\_\_\_\_
% mdparagraph:P206:end

% output: NEW-U16-F03-O009
% mdparagraph:P207:begin
(9) forlǣtan: one addressee: \_\_\_\_\_\_\_\_\_\_\_\_\_\_\_\_\_\_\_\_
% mdparagraph:P207:end

% output: NEW-U16-F03-O010
% mdparagraph:P208:begin
(10) forlǣtan: more than one addressee: \_\_\_\_\_\_\_\_\_\_\_\_\_\_\_\_\_\_\_\_
% mdparagraph:P208:end

% output: NEW-U16-F03-O011
% mdparagraph:P209:begin
(11) geswīcan: one addressee: \_\_\_\_\_\_\_\_\_\_\_\_\_\_\_\_\_\_\_\_
% mdparagraph:P209:end

% output: NEW-U16-F03-O012
% mdparagraph:P210:begin
(12) geswīcan: more than one addressee: \_\_\_\_\_\_\_\_\_\_\_\_\_\_\_\_\_\_\_\_
% mdparagraph:P210:end

% output: NEW-U16-F03-O013
% mdparagraph:P211:begin
(13) ęfthweorfan: one addressee: \_\_\_\_\_\_\_\_\_\_\_\_\_\_\_\_\_\_\_\_
% mdparagraph:P211:end

% output: NEW-U16-F03-O014
% mdparagraph:P212:begin
(14) ęfthweorfan: more than one addressee: \_\_\_\_\_\_\_\_\_\_\_\_\_\_\_\_\_\_\_\_
% mdparagraph:P212:end

\end{minipage}\par\addvspace{4pt}

\par\begin{minipage}{\linewidth}\raggedright
% mdparagraph:P213:begin
\textbf{F04}
% mdparagraph:P213:end

% item: NEW-U16-F04
% mdparagraph:P214:begin
Give the requested present-subjunctive forms and the two infinitive forms. The labels specify the form, not an English tense-conversion rule.
% mdparagraph:P214:end

% mdparagraph:P215:begin
(1) findan
% mdparagraph:P215:end

% mdparagraph:P216:begin
(2) fēdan
% mdparagraph:P216:end

% mdparagraph:P217:begin
(3) iernan
% mdparagraph:P217:end

% mdparagraph:P218:begin
\textbf{Supplied words.} These forms and meanings remain available for this task.
% mdparagraph:P218:end

% mdtable:T62:begin
\par\begingroup
\fontsize{10.5}{13.2}\selectfont
\setlength{\LTleft}{0pt}\setlength{\LTright}{\fill}
\noindent
\begin{tabular}{@{}>{\raggedright\arraybackslash}p{33.87mm}>{\raggedright\arraybackslash}p{33.87mm}>{\raggedright\arraybackslash}p{50.64mm}@{}}
\toprule
\textbf{Form} & \textbf{Meaning} & \textbf{Role} \\
\midrule
findan & find & given dictionary form \\
fēdan & feed; nourish & given dictionary form \\
\bottomrule
\end{tabular}
\par\endgroup

% mdtable:T62:end
% output: NEW-U16-F04-O001
% mdparagraph:P219:begin
(1) findan: present subjunctive, third singular: \_\_\_\_\_\_\_\_\_\_\_\_\_\_\_\_\_\_\_\_
% mdparagraph:P219:end

% output: NEW-U16-F04-O002
% mdparagraph:P220:begin
(2) findan: present subjunctive, third plural: \_\_\_\_\_\_\_\_\_\_\_\_\_\_\_\_\_\_\_\_
% mdparagraph:P220:end

% output: NEW-U16-F04-O003
% mdparagraph:P221:begin
(3) fēdan: present subjunctive, third singular: \_\_\_\_\_\_\_\_\_\_\_\_\_\_\_\_\_\_\_\_
% mdparagraph:P221:end

% output: NEW-U16-F04-O004
% mdparagraph:P222:begin
(4) fēdan: present subjunctive, third plural: \_\_\_\_\_\_\_\_\_\_\_\_\_\_\_\_\_\_\_\_
% mdparagraph:P222:end

% output: NEW-U16-F04-O005
% mdparagraph:P223:begin
(5) iernan: plain infinitive: \_\_\_\_\_\_\_\_\_\_\_\_\_\_\_\_\_\_\_\_
% mdparagraph:P223:end

% output: NEW-U16-F04-O006
% mdparagraph:P224:begin
(6) iernan: inflected infinitive after tō (give the infinitive only): \_\_\_\_\_\_\_\_\_\_\_\_\_\_\_\_\_\_\_\_
% mdparagraph:P224:end

\end{minipage}\par\addvspace{4pt}

\par\begin{minipage}{\linewidth}\raggedright
% mdparagraph:P225:begin
\textbf{F05}
% mdparagraph:P225:end

% item: NEW-U16-F05
% mdparagraph:P226:begin
Give the four requested strong forms of nearu. Then translate the three phrases using bēgen and sunu, without adding an article. Finally explain the stem alternation and the number of both.
% mdparagraph:P226:end

% mdparagraph:P227:begin
(1) nearu
% mdparagraph:P227:end

% mdparagraph:P228:begin
(2) bēgen
% mdparagraph:P228:end

% mdparagraph:P229:begin
(3) sunu
% mdparagraph:P229:end

% output: NEW-U16-F05-O001
% mdparagraph:P230:begin
(1) nearu: masculine nominative singular: \_\_\_\_\_\_\_\_\_\_\_\_\_\_\_\_\_\_\_\_
% mdparagraph:P230:end

% output: NEW-U16-F05-O002
% mdparagraph:P231:begin
(2) nearu: masculine accusative singular: \_\_\_\_\_\_\_\_\_\_\_\_\_\_\_\_\_\_\_\_
% mdparagraph:P231:end

% output: NEW-U16-F05-O003
% mdparagraph:P232:begin
(3) nearu: masculine genitive singular: \_\_\_\_\_\_\_\_\_\_\_\_\_\_\_\_\_\_\_\_
% mdparagraph:P232:end

% output: NEW-U16-F05-O004
% mdparagraph:P233:begin
(4) nearu: masculine dative plural: \_\_\_\_\_\_\_\_\_\_\_\_\_\_\_\_\_\_\_\_
% mdparagraph:P233:end

% output: NEW-U16-F05-O005
% mdparagraph:P234:begin
(5) both sons, accusative: \_\_\_\_\_\_\_\_\_\_\_\_\_\_\_\_\_\_\_\_
% mdparagraph:P234:end

% output: NEW-U16-F05-O006
% mdparagraph:P235:begin
(6) of both sons, genitive: \_\_\_\_\_\_\_\_\_\_\_\_\_\_\_\_\_\_\_\_
% mdparagraph:P235:end

% output: NEW-U16-F05-O007
% mdparagraph:P236:begin
(7) to both sons, dative: \_\_\_\_\_\_\_\_\_\_\_\_\_\_\_\_\_\_\_\_
% mdparagraph:P236:end

% output: NEW-U16-F05-O008
% mdparagraph:P237:begin
(8) Explain the nearu stem alternation shown by your answers: \_\_\_\_\_\_\_\_\_\_\_\_\_\_\_\_\_\_\_\_
% mdparagraph:P237:end

% output: NEW-U16-F05-O009
% mdparagraph:P238:begin
(9) Explain the grammatical number of both in the phrases: \_\_\_\_\_\_\_\_\_\_\_\_\_\_\_\_\_\_\_\_
% mdparagraph:P238:end

\end{minipage}\par\addvspace{4pt}

\par\begin{minipage}{\linewidth}\raggedright
% mdparagraph:P239:begin
\textbf{A01}
% mdparagraph:P239:end

% item: NEW-U16-A01
% mdparagraph:P240:begin
Give the two defensible readings of the final hē. Then rewrite the second sentence twice, replacing only hē with an explicit noun phrase so that each reading is unambiguous. The first sentence identifies the king's son.
% mdparagraph:P240:end

% mdparagraph:P241:begin
(1) Se cyning geseah his sunu. Hē ēode.
% mdparagraph:P241:end

% output: NEW-U16-A01-O001
% mdparagraph:P242:begin
(1) First reading: \_\_\_\_\_\_\_\_\_\_\_\_\_\_\_\_\_\_\_\_
% mdparagraph:P242:end

% output: NEW-U16-A01-O002
% mdparagraph:P243:begin
(2) Second reading: \_\_\_\_\_\_\_\_\_\_\_\_\_\_\_\_\_\_\_\_
% mdparagraph:P243:end

% output: NEW-U16-A01-O003
% mdparagraph:P244:begin
(3) First explicit version: \_\_\_\_\_\_\_\_\_\_\_\_\_\_\_\_\_\_\_\_
% mdparagraph:P244:end

% output: NEW-U16-A01-O004
% mdparagraph:P245:begin
(4) Second explicit version: \_\_\_\_\_\_\_\_\_\_\_\_\_\_\_\_\_\_\_\_
% mdparagraph:P245:end

\end{minipage}\par\addvspace{4pt}

\par\begin{minipage}{\linewidth}\raggedright
% mdparagraph:P246:begin
\textbf{A02}
% mdparagraph:P246:end

% item: NEW-U16-A02
% mdparagraph:P247:begin
Translate both sentences. For each, identify the expression denoting the person who went, its case and its relation to the verb of going. Then compare the two forms of go: which is nonfinite, which is finite, and does the spelling of the finite form distinguish its mood here?
% mdparagraph:P247:end

% mdparagraph:P248:begin
(1) Se cyning geseah ðone sunu gān.
% mdparagraph:P248:end

% mdparagraph:P249:begin
(2) Se cyning wiste ðæt se sunu ēode.
% mdparagraph:P249:end

% output: NEW-U16-A02-O001
% mdparagraph:P250:begin
(1) Sentence 1 translation: \_\_\_\_\_\_\_\_\_\_\_\_\_\_\_\_\_\_\_\_
% mdparagraph:P250:end

% output: NEW-U16-A02-O002
% mdparagraph:P251:begin
(2) Sentence 2 translation: \_\_\_\_\_\_\_\_\_\_\_\_\_\_\_\_\_\_\_\_
% mdparagraph:P251:end

% output: NEW-U16-A02-O003
% mdparagraph:P252:begin
(3) Sentence 1 actor: \_\_\_\_\_\_\_\_\_\_\_\_\_\_\_\_\_\_\_\_
% mdparagraph:P252:end

% output: NEW-U16-A02-O004
% mdparagraph:P253:begin
(4) Sentence 2 actor: \_\_\_\_\_\_\_\_\_\_\_\_\_\_\_\_\_\_\_\_
% mdparagraph:P253:end

% output: NEW-U16-A02-O005
% mdparagraph:P254:begin
(5) Verb-form comparison: \_\_\_\_\_\_\_\_\_\_\_\_\_\_\_\_\_\_\_\_
% mdparagraph:P254:end

\end{minipage}\par\addvspace{4pt}

\par\begin{minipage}{\linewidth}\raggedright
% mdparagraph:P255:begin
\textbf{A03}
% mdparagraph:P255:end

% item: NEW-U16-A03
% mdparagraph:P256:begin
Translate each report. For each, identify and parse every finite verb, identify any participle or infinitive (or state that there is none), and explain what the construction and context show about the time of the reported situation. In sentence 1 the king is there at the time of saying; in sentence 2 the ship-making is complete; in sentence 3 the singing is intended but has not yet occurred. Your explanation must connect these contexts with the actual verb groups and explain why automatic English backshift is insufficient; merely repeating the time hint is not an answer.
% mdparagraph:P256:end

% mdparagraph:P257:begin
(1) Hē sægde ðæt se cyning ðǣr wǣre.
% mdparagraph:P257:end

% mdparagraph:P258:begin
(2) Hē sægde ðæt se cyning ðæt scip geworht hæfde.
% mdparagraph:P258:end

% mdparagraph:P259:begin
(3) Hē sægde ðæt se cyning singan wolde.
% mdparagraph:P259:end

% output: NEW-U16-A03-O001
% mdparagraph:P260:begin
(1) Sentence 1 translation: \_\_\_\_\_\_\_\_\_\_\_\_\_\_\_\_\_\_\_\_
% mdparagraph:P260:end

% output: NEW-U16-A03-O002
% mdparagraph:P261:begin
(2) Sentence 1 finite forms: \_\_\_\_\_\_\_\_\_\_\_\_\_\_\_\_\_\_\_\_
% mdparagraph:P261:end

% output: NEW-U16-A03-O003
% mdparagraph:P262:begin
(3) Sentence 1 nonfinite forms: \_\_\_\_\_\_\_\_\_\_\_\_\_\_\_\_\_\_\_\_
% mdparagraph:P262:end

% output: NEW-U16-A03-O004
% mdparagraph:P263:begin
(4) Sentence 1 construction and time: \_\_\_\_\_\_\_\_\_\_\_\_\_\_\_\_\_\_\_\_
% mdparagraph:P263:end

% output: NEW-U16-A03-O005
% mdparagraph:P264:begin
(5) Sentence 2 translation: \_\_\_\_\_\_\_\_\_\_\_\_\_\_\_\_\_\_\_\_
% mdparagraph:P264:end

% output: NEW-U16-A03-O006
% mdparagraph:P265:begin
(6) Sentence 2 finite forms: \_\_\_\_\_\_\_\_\_\_\_\_\_\_\_\_\_\_\_\_
% mdparagraph:P265:end

% output: NEW-U16-A03-O007
% mdparagraph:P266:begin
(7) Sentence 2 nonfinite form: \_\_\_\_\_\_\_\_\_\_\_\_\_\_\_\_\_\_\_\_
% mdparagraph:P266:end

% output: NEW-U16-A03-O008
% mdparagraph:P267:begin
(8) Sentence 2 construction and time: \_\_\_\_\_\_\_\_\_\_\_\_\_\_\_\_\_\_\_\_
% mdparagraph:P267:end

% output: NEW-U16-A03-O009
% mdparagraph:P268:begin
(9) Sentence 3 translation: \_\_\_\_\_\_\_\_\_\_\_\_\_\_\_\_\_\_\_\_
% mdparagraph:P268:end

% output: NEW-U16-A03-O010
% mdparagraph:P269:begin
(10) Sentence 3 finite forms: \_\_\_\_\_\_\_\_\_\_\_\_\_\_\_\_\_\_\_\_
% mdparagraph:P269:end

% output: NEW-U16-A03-O011
% mdparagraph:P270:begin
(11) Sentence 3 nonfinite form: \_\_\_\_\_\_\_\_\_\_\_\_\_\_\_\_\_\_\_\_
% mdparagraph:P270:end

% output: NEW-U16-A03-O012
% mdparagraph:P271:begin
(12) Sentence 3 construction and time: \_\_\_\_\_\_\_\_\_\_\_\_\_\_\_\_\_\_\_\_
% mdparagraph:P271:end

\end{minipage}\par\addvspace{4pt}

\reserve{7\baselineskip}
% mdparagraph:P272:begin
\section*{5. New compositions}\BookHeading{nav-ch16-learner-0023}{subsection}{5. New compositions}{1}
% mdparagraph:P272:end

% mdparagraph:P273:begin
Translate every clause and identify the requested relations. Keep a grammatical alternative distinct from compliance with an explicitly requested verb or construction.
% mdparagraph:P273:end

\reserve{5\baselineskip}
\par\begin{minipage}{\linewidth}\raggedright
% mdparagraph:P274:begin
\textbf{C01}
% mdparagraph:P274:end

% item: NEW-U16-C01
% mdparagraph:P275:begin
Translate the complete exchange. Use sęcgan for the reporting verb and sęllan for give: The king said to his son, 'Give me the book.' The son gave him the book. Then identify the referents of the first-person and third-person recipient pronouns in your translation.
% mdparagraph:P275:end

% output: NEW-U16-C01-O001
% mdparagraph:P276:begin
(1) Speech and command: \_\_\_\_\_\_\_\_\_\_\_\_\_\_\_\_\_\_\_\_
% mdparagraph:P276:end

% output: NEW-U16-C01-O002
% mdparagraph:P277:begin
(2) Executed action: \_\_\_\_\_\_\_\_\_\_\_\_\_\_\_\_\_\_\_\_
% mdparagraph:P277:end

% output: NEW-U16-C01-O003
% mdparagraph:P278:begin
(3) Recipient reference: \_\_\_\_\_\_\_\_\_\_\_\_\_\_\_\_\_\_\_\_
% mdparagraph:P278:end

\end{minipage}\par\addvspace{4pt}

\par\begin{minipage}{\linewidth}\raggedright
% mdparagraph:P279:begin
\textbf{C02}
% mdparagraph:P279:end

% item: NEW-U16-C02
% mdparagraph:P280:begin
Translate these three statements separately. Use cunnan in the first two and witan in the third; then explain the complement selected in each: I know that man. I know how to sing. I know that he is here.
% mdparagraph:P280:end

% output: NEW-U16-C02-O001
% mdparagraph:P281:begin
(1) Statement 1: \_\_\_\_\_\_\_\_\_\_\_\_\_\_\_\_\_\_\_\_
% mdparagraph:P281:end

% output: NEW-U16-C02-O002
% mdparagraph:P282:begin
(2) Statement 2: \_\_\_\_\_\_\_\_\_\_\_\_\_\_\_\_\_\_\_\_
% mdparagraph:P282:end

% output: NEW-U16-C02-O003
% mdparagraph:P283:begin
(3) Statement 3: \_\_\_\_\_\_\_\_\_\_\_\_\_\_\_\_\_\_\_\_
% mdparagraph:P283:end

% output: NEW-U16-C02-O004
% mdparagraph:P284:begin
(4) Complement comparison: \_\_\_\_\_\_\_\_\_\_\_\_\_\_\_\_\_\_\_\_
% mdparagraph:P284:end

\end{minipage}\par\addvspace{4pt}

\par\begin{minipage}{\linewidth}\raggedright
% mdparagraph:P285:begin
\textbf{C03}
% mdparagraph:P285:end

% item: NEW-U16-C03
% mdparagraph:P286:begin
Translate both clauses, using plegian for the first two expressed actions and dōn for the final repeated action: The men played well; the king played better than they did. Explain what the last verb replaces and why it is plural.
% mdparagraph:P286:end

% output: NEW-U16-C03-O001
% mdparagraph:P287:begin
(1) Whole comparison: \_\_\_\_\_\_\_\_\_\_\_\_\_\_\_\_\_\_\_\_
% mdparagraph:P287:end

% output: NEW-U16-C03-O002
% mdparagraph:P288:begin
(2) Recovered action and agreement: \_\_\_\_\_\_\_\_\_\_\_\_\_\_\_\_\_\_\_\_
% mdparagraph:P288:end

\end{minipage}\par\addvspace{4pt}

\reserve{7\baselineskip}
% mdparagraph:P289:begin
\section*{6. Cook: Group XII, XIII, Exercises 32–36}\BookHeading{nav-ch16-learner-0024}{subsection}{6. Cook: Group XII, XIII, Exercises 32–36}{1}
% mdparagraph:P289:end

% teaching: NEW-U16-COOK-SUPPORT
% mdparagraph:P290:begin
Translate the complete original prompts. Parenthesized words, omissions, number instructions and order directions belong to Cook’s exercises. The modern helps below are additional and remain available; a supplied expression does not require independent recall. Cook’s paragraph numbers refer to his grammar, but the material needed here is stated in this unit.\textsuperscript{6}
% mdparagraph:P290:end

% mdparagraph:P291:begin
\textbf{Source correction for 34.9.} Cook prints “reward (āgyltan, 113)”. The verb meaning “repay; reward” is strong \textbf{āgyldan}, belonging with §104; \textbf{āgylt} is a valid contracted third singular. The complete reward clause is supplied with the item. No additional full conjugation is required. The original prompt is retained.\textsuperscript{7}
% mdparagraph:P291:end

% mdparagraph:P292:begin
\textbf{Names in 36.9–10.} The supplied \textbf{Johannes} is genitive and \textbf{Johannem} accusative. Their endings are directly attested with initial I; initial J is used here to follow Cook’s headword. These name forms remain supplied.\textsuperscript{8}
% mdparagraph:P292:end

\reserve{7\baselineskip}
% mdparagraph:P293:begin
\subsection*{Exercise 32}\BookHeading{nav-ch16-learner-0025}{subsubsection}{Exercise 32}{0}
% mdparagraph:P293:end

\reserve{5\baselineskip}
\par\begin{minipage}{\linewidth}\raggedright
% mdparagraph:P294:begin
\textbf{Cook XII, XIII.32.1}
% mdparagraph:P294:end

% item: EX-XII-XIII-32-01
% mdparagraph:P295:begin
I bid you all farewell (grētan).
% mdparagraph:P295:end

% output: EX-XII-XIII-32-01-O001
\end{minipage}\par\addvspace{4pt}

\par\begin{minipage}{\linewidth}\raggedright
% mdparagraph:P296:begin
\textbf{Cook XII, XIII.32.2}
% mdparagraph:P296:end

% item: EX-XII-XIII-32-02
% mdparagraph:P297:begin
When they had rowed about (swilce) twenty furlongs (furlang, 47), or thirty, then saw they Jesus walking (gān) on (uppan, 166) the sea (verb last).
% mdparagraph:P297:end

% mdparagraph:P298:begin
\textbf{Supplied help:} swilce twēntig furlanga oððe ðrītig; uppan ðǣm sǣ
% mdparagraph:P298:end

% mdparagraph:P299:begin
About twenty furlongs or thirty; upon the sea. The supplied numerical phrase takes genitive furlanga; the prepositional phrase has the dative.
% mdparagraph:P299:end

\end{minipage}\par\addvspace{4pt}

% teaching: NEW-U16-SUP-SWILCE
% teaching: NEW-U16-VOC-SWILCE
% teaching: NEW-U16-SUP-TWENTIG
% teaching: NEW-U16-VOC-TWENTIG
% teaching: NEW-U16-SUP-THRITIG
% teaching: NEW-U16-VOC-THRITIG
% teaching: NEW-U16-SUP-UPPAN
% teaching: NEW-U16-VOC-UPPAN
\reserve{11\baselineskip}
% mdparagraph:P300:begin
\textbf{Supplied words.} These forms and meanings remain available for this task.
% mdparagraph:P300:end

% mdtable:T63:begin
\par\begingroup
\fontsize{10.5}{13.2}\selectfont
\setlength{\LTleft}{0pt}\setlength{\LTright}{\fill}
\noindent
\begin{tabular}{@{}>{\raggedright\arraybackslash}p{31.95mm}>{\raggedright\arraybackslash}p{51.6mm}>{\raggedright\arraybackslash}p{34.83mm}@{}}
\toprule
\textbf{Form} & \textbf{Meaning} & \textbf{Role} \\
\midrule
swilce & about; approximately & invariable \\
twēntig & twenty & as used here \\
ðrītig & thirty & as used here \\
uppan & upon; on top of & invariable \\
\bottomrule
\end{tabular}
\par\endgroup

% mdtable:T63:end
% output: EX-XII-XIII-32-02-O001
\reserve{5\baselineskip}
\par\begin{minipage}{\linewidth}\raggedright
% mdparagraph:P301:begin
\textbf{Cook XII, XIII.32.3}
% mdparagraph:P301:end

% item: EX-XII-XIII-32-03
% mdparagraph:P302:begin
He arose and commanded (bebēodan) the wind (164. b), and said (cweðan) to the sea, Subside.
% mdparagraph:P302:end

% output: EX-XII-XIII-32-03-O001
\end{minipage}\par\addvspace{4pt}

\par\begin{minipage}{\linewidth}\raggedright
% mdparagraph:P303:begin
\textbf{Cook XII, XIII.32.4}
% mdparagraph:P303:end

% item: EX-XII-XIII-32-04
% mdparagraph:P304:begin
The wind ceased (geswīcan), and there was (weorðan) (om. a) great calm.
% mdparagraph:P304:end

% output: EX-XII-XIII-32-04-O001
\end{minipage}\par\addvspace{4pt}

\par\begin{minipage}{\linewidth}\raggedright
% mdparagraph:P305:begin
\textbf{Cook XII, XIII.32.5}
% mdparagraph:P305:end

% item: EX-XII-XIII-32-05
% mdparagraph:P306:begin
Thou hearest his voice, and thou knowest not whence he cometh, nor whither he goeth.
% mdparagraph:P306:end

% output: EX-XII-XIII-32-05-O001
\end{minipage}\par\addvspace{4pt}

\par\begin{minipage}{\linewidth}\raggedright
% mdparagraph:P307:begin
\textbf{Cook XII, XIII.32.6}
% mdparagraph:P307:end

% item: EX-XII-XIII-32-06
% mdparagraph:P308:begin
He saw two brethren casting (sęndan) their net (nętt, 47) into (on, 172) the sea; they were fishers.
% mdparagraph:P308:end

% output: EX-XII-XIII-32-06-O001
\end{minipage}\par\addvspace{4pt}

\par\begin{minipage}{\linewidth}\raggedright
% mdparagraph:P309:begin
\textbf{Cook XII, XIII.32.7}
% mdparagraph:P309:end

% item: EX-XII-XIII-32-07
% mdparagraph:P310:begin
And he said (sęcgean) to them (164. b), Come after (æfter, 166) me, and I (om. will) make (dōn) that ye (dual) are fishers of men (153. d); and they immediately left their nets.
% mdparagraph:P310:end

% output: EX-XII-XIII-32-07-O001
\end{minipage}\par\addvspace{4pt}

\par\begin{minipage}{\linewidth}\raggedright
% mdparagraph:P311:begin
\textbf{Cook XII, XIII.32.8}
% mdparagraph:P311:end

% item: EX-XII-XIII-32-08
% mdparagraph:P312:begin
If thy hand or thy foot offend (ǣswician, 118) thee, cut it (āceorfan, 104) off (of), and cast (om. it) from thee; (om. it) is better (om. for) thee (165) that thou go halt (healt, 58) into life, than (om. that) thou have two hands and two feet, and be cast into everlasting fire.
% mdparagraph:P312:end

% mdparagraph:P313:begin
\textbf{Supplied help:} bętere ðē is ðæt ðū ... gā, ðonne ðū ... hæbbe and ... sīe worpen
% mdparagraph:P313:end

% mdparagraph:P314:begin
It is better for you to go ... than to have ... and be thrown ...
% mdparagraph:P314:end

% mdparagraph:P315:begin
The experiencer is dative ðē; bętere is the familiar neuter comparative predicate. Keep the requested ðæt in the first complement and omit it after ðonne. Earlier hit resumes whichever of the differently gendered body parts is at issue.
% mdparagraph:P315:end

% output: EX-XII-XIII-32-08-O001
\end{minipage}\par\addvspace{4pt}

\par\begin{minipage}{\linewidth}\raggedright
% mdparagraph:P316:begin
\textbf{Cook XII, XIII.32.9}
% mdparagraph:P316:end

% item: EX-XII-XIII-32-09
% mdparagraph:P317:begin
If any (hwilc) man hath (om. a) hundred sheep (154. c), and loseth (forlēosan) one of (of, 166) them [those], (om. doth) he not leave the nine and ninety in the mountains, and goeth and seeketh (sēcean) the one that (87. d) was perishing (forwearð)?
% mdparagraph:P317:end

% mdparagraph:P318:begin
\textbf{Supplied help:} hund scēapa; nigon and hundnigontig; ān of ðǣm
% mdparagraph:P318:end

% mdparagraph:P319:begin
A hundred sheep; ninety-nine; one of those. These numerical and partitive expressions remain supplied.
% mdparagraph:P319:end

\end{minipage}\par\addvspace{4pt}

% teaching: NEW-U16-SUP-HUND
% teaching: NEW-U16-VOC-HUND
% teaching: NEW-U16-SUP-HUNDNIGONTIG
% teaching: NEW-U16-VOC-HUNDNIGONTIG
% teaching: NEW-U16-VOC-OF
\reserve{10\baselineskip}
% mdparagraph:P320:begin
\textbf{Supplied words.} These forms and meanings remain available for this task.
% mdparagraph:P320:end

% mdtable:T64:begin
\par\begingroup
\fontsize{10.5}{13.2}\selectfont
\setlength{\LTleft}{0pt}\setlength{\LTright}{\fill}
\noindent
\begin{tabular}{@{}>{\raggedright\arraybackslash}p{22.44mm}>{\raggedright\arraybackslash}p{77.23mm}>{\raggedright\arraybackslash}p{18.7mm}@{}}
\toprule
\textbf{Form} & \textbf{Meaning} & \textbf{Role} \\
\midrule
hund & hundred, in the supplied hund scēapa phrase & as used here \\
hundnigontig & ninety & as used here \\
of & of; from a portion of a whole, in the supplied phrases ān of ðǣm and of ēowrum ęle & as used here \\
\bottomrule
\end{tabular}
\par\endgroup

% mdtable:T64:end
% output: EX-XII-XIII-32-09-O001
\reserve{5\baselineskip}
\par\begin{minipage}{\linewidth}\raggedright
% mdparagraph:P321:begin
\textbf{Cook XII, XIII.32.10}
% mdparagraph:P321:end

% item: EX-XII-XIII-32-10
% mdparagraph:P322:begin
And if it happeneth that he finds it, truly I say (sęcgean) to you that he rejoiceth more (swīðor) over (ofer, 172) that one than over the nine and ninety.
% mdparagraph:P322:end

% mdparagraph:P323:begin
\textbf{Supplied help:} gif hit gelimpð ðæt hē hit finde
% mdparagraph:P323:end

% mdparagraph:P324:begin
if it happens that he finds it
% mdparagraph:P324:end

% mdparagraph:P325:begin
The first hit is impersonal; the second hit is the sheep. Gelimpð is supplied; finde is the newly taught present subjunctive. No full gelimpan person paradigm is assigned.
% mdparagraph:P325:end

% mdparagraph:P326:begin
\textbf{Supplied help:} hund scēapa; nigon and hundnigontig; ān of ðǣm
% mdparagraph:P326:end

% mdparagraph:P327:begin
A hundred sheep; ninety-nine; one of those. These numerical and partitive expressions remain supplied.
% mdparagraph:P327:end

% mdparagraph:P328:begin
\textbf{Supplied words.} These forms and meanings remain available for this task.
% mdparagraph:P328:end

% mdtable:T65:begin
\par\begingroup
\fontsize{10.5}{13.2}\selectfont
\setlength{\LTleft}{0pt}\setlength{\LTright}{\fill}
\noindent
\begin{tabular}{@{}>{\raggedright\arraybackslash}p{24mm}>{\raggedright\arraybackslash}p{58mm}>{\raggedright\arraybackslash}p{36.38mm}@{}}
\toprule
\textbf{Form} & \textbf{Meaning} & \textbf{Role} \\
\midrule
gelimpð & happen; come about & third person, present, indicative, singular \\
hund & hundred, in the supplied hund scēapa phrase & as used here \\
hundnigontig & ninety & as used here \\
of & of; from a portion of a whole, in the supplied phrases ān of ðǣm and of ēowrum ęle & as used here \\
\bottomrule
\end{tabular}
\par\endgroup

% mdtable:T65:end
% output: EX-XII-XIII-32-10-O001
\end{minipage}\par\addvspace{4pt}

\par\begin{minipage}{\linewidth}\raggedright
% mdparagraph:P329:begin
\textbf{Cook XII, XIII.32.11}
% mdparagraph:P329:end

% item: EX-XII-XIII-32-11
% mdparagraph:P330:begin
When he cometh home (hām), he calleth (cleopian) together his friends and his neighbors (nēahgebūr, 43), and saith (cweðan), Rejoice with me, because I (om. have) found my lost sheep.
% mdparagraph:P330:end

% output: EX-XII-XIII-32-11-O001
\end{minipage}\par\addvspace{4pt}

\reserve{7\baselineskip}
% mdparagraph:P331:begin
\subsection*{Exercise 33}\BookHeading{nav-ch16-learner-0026}{subsubsection}{Exercise 33}{0}
% mdparagraph:P331:end

\reserve{5\baselineskip}
\par\begin{minipage}{\linewidth}\raggedright
% mdparagraph:P332:begin
\textbf{Cook XII, XIII.33.1}
% mdparagraph:P332:end

% item: EX-XII-XIII-33-01
% mdparagraph:P333:begin
And, going (gangan) into (intō, 166) the house, they found the child with Mary (Maria, 53) his mother.
% mdparagraph:P333:end

% mdparagraph:P334:begin
\textbf{Supplied help:} intō ðǣm hūse gangende
% mdparagraph:P334:end

% mdparagraph:P335:begin
going into the house
% mdparagraph:P335:end

% mdparagraph:P336:begin
Gangende accompanies the plural actors hīe. The participial phrase is supplied as a whole; it does not assign a new complete gangan finite system.
% mdparagraph:P336:end

% mdparagraph:P337:begin
\textbf{Supplied help:} mid Marian his mēder
% mdparagraph:P337:end

% mdparagraph:P338:begin
With Mary his mother. Marian is the supplied weak-feminine form after mid here; the earlier genitive use of the name had a different function.
% mdparagraph:P338:end

\end{minipage}\par\addvspace{4pt}

% teaching: NEW-U16-SUP-MARIA-NAME
% teaching: NEW-U16-VOC-MARIA-NAME
\reserve{9\baselineskip}
% mdparagraph:P339:begin
\textbf{Supplied words.} These forms and meanings remain available for this task.
% mdparagraph:P339:end

% mdtable:T66:begin
\par\begingroup
\fontsize{10.5}{13.2}\selectfont
\setlength{\LTleft}{0pt}\setlength{\LTright}{\fill}
\noindent
\begin{tabular}{@{}>{\raggedright\arraybackslash}p{17.92mm}>{\raggedright\arraybackslash}p{51.71mm}>{\raggedright\arraybackslash}p{48.75mm}@{}}
\toprule
\textbf{Form} & \textbf{Meaning} & \textbf{Role} \\
\midrule
gangende & go; enter in the supplied circumstantial phrase & present, participle, nominative, plural, strong \\
Marian & Mary & dative, singular, feminine \\
\bottomrule
\end{tabular}
\par\endgroup

% mdtable:T66:end
% output: EX-XII-XIII-33-01-O001
\reserve{5\baselineskip}
\par\begin{minipage}{\linewidth}\raggedright
% mdparagraph:P340:begin
\textbf{Cook XII, XIII.33.2}
% mdparagraph:P340:end

% item: EX-XII-XIII-33-02
% mdparagraph:P341:begin
Arise, and take (niman) the child and his mother by night, and flee to (tō, 166) (om. the) land of the Egyptians, and be there until (that) I tell thee (verb last).
% mdparagraph:P341:end

% mdparagraph:P342:begin
\textbf{Supplied help:} Egypta land /\allowbreak{} tō Egypta lande /\allowbreak{} of Egyptum
% mdparagraph:P342:end

% mdparagraph:P343:begin
Land of the Egyptians; to their land; from the Egyptians. These foreign-name forms and their place-reference uses remain supplied.
% mdparagraph:P343:end

\end{minipage}\par\addvspace{4pt}

% teaching: NEW-U16-SUP-EGYPTA
% teaching: NEW-U16-VOC-EGYPTA
% teaching: NEW-U16-VOC-ON
\reserve{10\baselineskip}
% mdparagraph:P344:begin
\textbf{Supplied words.} These forms and meanings remain available for this task.
% mdparagraph:P344:end

% mdtable:T67:begin
\par\begingroup
\fontsize{10.5}{13.2}\selectfont
\setlength{\LTleft}{0pt}\setlength{\LTright}{\fill}
\noindent
\begin{tabular}{@{}>{\raggedright\arraybackslash}p{19.8mm}>{\raggedright\arraybackslash}p{64.73mm}>{\raggedright\arraybackslash}p{33.85mm}@{}}
\toprule
\textbf{Form} & \textbf{Meaning} & \textbf{Role} \\
\midrule
Egypta & Egyptians; country-name phrases Egypta land and of Egyptum & genitive, plural, masculine \\
Egyptum & Egyptians; country-name phrases Egypta land and of Egyptum & dative, plural, masculine \\
on & during, with dative time & as used here \\
\bottomrule
\end{tabular}
\par\endgroup

% mdtable:T67:end
% output: EX-XII-XIII-33-02-O001
\reserve{5\baselineskip}
\par\begin{minipage}{\linewidth}\raggedright
% mdparagraph:P345:begin
\textbf{Cook XII, XIII.33.3}
% mdparagraph:P345:end

% item: EX-XII-XIII-33-03
% mdparagraph:P346:begin
He arose then, and took the child and his mother by night, and journeyed into (on, 172) the land of the Egyptians; that it (om. might) be fulfilled (gefyllan) which (87. a) was spoken (cweðan) of (fram, 166) (om. the) Lord by (ðurh, 172) the prophet, From (of, 166) (om. the) Egyptians (om. have) I called (cleopian) my son.
% mdparagraph:P346:end

% mdparagraph:P347:begin
\textbf{Supplied help:} Egypta land /\allowbreak{} tō Egypta lande /\allowbreak{} of Egyptum
% mdparagraph:P347:end

% mdparagraph:P348:begin
Land of the Egyptians; to their land; from the Egyptians. These foreign-name forms and their place-reference uses remain supplied.
% mdparagraph:P348:end

% mdparagraph:P349:begin
\textbf{Supplied words.} These forms and meanings remain available for this task.
% mdparagraph:P349:end

% mdtable:T68:begin
\par\begingroup
\fontsize{10.5}{13.2}\selectfont
\setlength{\LTleft}{0pt}\setlength{\LTright}{\fill}
\noindent
\begin{tabular}{@{}>{\raggedright\arraybackslash}p{19.8mm}>{\raggedright\arraybackslash}p{64.73mm}>{\raggedright\arraybackslash}p{33.85mm}@{}}
\toprule
\textbf{Form} & \textbf{Meaning} & \textbf{Role} \\
\midrule
Egypta & Egyptians; country-name phrases Egypta land and of Egyptum & genitive, plural, masculine \\
Egyptum & Egyptians; country-name phrases Egypta land and of Egyptum & dative, plural, masculine \\
on & during, with dative time & as used here \\
\bottomrule
\end{tabular}
\par\endgroup

% mdtable:T68:end
% output: EX-XII-XIII-33-03-O001
\end{minipage}\par\addvspace{4pt}

\par\begin{minipage}{\linewidth}\raggedright
% mdparagraph:P350:begin
\textbf{Cook XII, XIII.33.4}
% mdparagraph:P350:end

% item: EX-XII-XIII-33-04
% mdparagraph:P351:begin
What can they ask for [request; 156. b]?
% mdparagraph:P351:end

% output: EX-XII-XIII-33-04-O001
\end{minipage}\par\addvspace{4pt}

\par\begin{minipage}{\linewidth}\raggedright
% mdparagraph:P352:begin
\textbf{Cook XII, XIII.33.5}
% mdparagraph:P352:end

% item: EX-XII-XIII-33-05
% mdparagraph:P353:begin
The sun rules over the day (156. f).
% mdparagraph:P353:end

% output: EX-XII-XIII-33-05-O001
\end{minipage}\par\addvspace{4pt}

\par\begin{minipage}{\linewidth}\raggedright
% mdparagraph:P354:begin
\textbf{Cook XII, XIII.33.6}
% mdparagraph:P354:end

% item: EX-XII-XIII-33-06
% mdparagraph:P355:begin
They expected nothing but [except] death itself (156. g).
% mdparagraph:P355:end

% mdparagraph:P356:begin
\textbf{Supplied help:} nāhtes būtan dēaðes selfes
% mdparagraph:P356:end

% mdparagraph:P357:begin
nothing except death itself: the whole exception phrase is supplied with the genitive governed by the expectation construction.
% mdparagraph:P357:end

\end{minipage}\par\addvspace{4pt}

% teaching: NEW-U16-SUP-NAHT
% teaching: NEW-U16-VOC-NAHT
\reserve{9\baselineskip}
% mdparagraph:P358:begin
\textbf{Supplied words.} These forms and meanings remain available for this task.
% mdparagraph:P358:end

% mdtable:T69:begin
\par\begingroup
\fontsize{10.5}{13.2}\selectfont
\setlength{\LTleft}{0pt}\setlength{\LTright}{\fill}
\noindent
\begin{tabular}{@{}>{\raggedright\arraybackslash}p{31.32mm}>{\raggedright\arraybackslash}p{34.92mm}>{\raggedright\arraybackslash}p{52.14mm}@{}}
\toprule
\textbf{Form} & \textbf{Meaning} & \textbf{Role} \\
\midrule
nāhtes & nothing; naught & genitive, singular, neuter \\
būtan & except & as used here \\
\bottomrule
\end{tabular}
\par\endgroup

% mdtable:T69:end
% output: EX-XII-XIII-33-06-O001
\reserve{5\baselineskip}
\par\begin{minipage}{\linewidth}\raggedright
% mdparagraph:P359:begin
\textbf{Cook XII, XIII.33.7}
% mdparagraph:P359:end

% item: EX-XII-XIII-33-07
% mdparagraph:P360:begin
Thou wilt cease from thy work (156. k).
% mdparagraph:P360:end

% output: EX-XII-XIII-33-07-O001
\end{minipage}\par\addvspace{4pt}

\par\begin{minipage}{\linewidth}\raggedright
% mdparagraph:P361:begin
\textbf{Cook XII, XIII.33.8}
% mdparagraph:P361:end

% item: EX-XII-XIII-33-08
% mdparagraph:P362:begin
Perhaps he will find a friend (194. a).
% mdparagraph:P362:end

% output: EX-XII-XIII-33-08-O001
\end{minipage}\par\addvspace{4pt}

\par\begin{minipage}{\linewidth}\raggedright
% mdparagraph:P363:begin
\textbf{Cook XII, XIII.33.9}
% mdparagraph:P363:end

% item: EX-XII-XIII-33-09
% mdparagraph:P364:begin
He said to them, I will (139) go a fishing; they said, And we will go with thee.
% mdparagraph:P364:end

% mdparagraph:P365:begin
\textbf{Supplied help:} on fiscnoð
% mdparagraph:P365:end

% mdparagraph:P366:begin
a-fishing; engaged in fishing. Keep the source spelling fiscnoð, not unmarked replacement fiscnað.
% mdparagraph:P366:end

\end{minipage}\par\addvspace{4pt}

% teaching: NEW-U16-SUP-FISCNOTH
% teaching: NEW-U16-VOC-FISCNOTH
\reserve{9\baselineskip}
% mdparagraph:P367:begin
\textbf{Supplied words.} These forms and meanings remain available for this task.
% mdparagraph:P367:end

% mdtable:T70:begin
\par\begingroup
\fontsize{10.5}{13.2}\selectfont
\setlength{\LTleft}{0pt}\setlength{\LTright}{\fill}
\noindent
\begin{tabular}{@{}>{\raggedright\arraybackslash}p{21.1mm}>{\raggedright\arraybackslash}p{55.07mm}>{\raggedright\arraybackslash}p{42.2mm}@{}}
\toprule
\textbf{Form} & \textbf{Meaning} & \textbf{Role} \\
\midrule
on & in the fishing activity, in supplied on fiscnoð & as used here \\
fiscnoð & fishing & accusative, singular, masculine \\
\bottomrule
\end{tabular}
\par\endgroup

% mdtable:T70:end
% output: EX-XII-XIII-33-09-O001
\reserve{5\baselineskip}
\par\begin{minipage}{\linewidth}\raggedright
% mdparagraph:P368:begin
\textbf{Cook XII, XIII.33.10}
% mdparagraph:P368:end

% item: EX-XII-XIII-33-10
% mdparagraph:P369:begin
They went out (ūt), and went into (om. a) ship, and caught nothing (183) in (on) that night.
% mdparagraph:P369:end

% mdparagraph:P370:begin
\textbf{Supplied words.} These forms and meanings remain available for this task.
% mdparagraph:P370:end

% mdtable:T71:begin
\par\begingroup
\fontsize{10.5}{13.2}\selectfont
\setlength{\LTleft}{0pt}\setlength{\LTright}{\fill}
\noindent
\begin{tabular}{@{}>{\raggedright\arraybackslash}p{34.16mm}>{\raggedright\arraybackslash}p{50.05mm}>{\raggedright\arraybackslash}p{34.16mm}@{}}
\toprule
\textbf{Form} & \textbf{Meaning} & \textbf{Role} \\
\midrule
on & during, with dative time & as used here \\
\bottomrule
\end{tabular}
\par\endgroup

% mdtable:T71:end
% output: EX-XII-XIII-33-10-O001
\end{minipage}\par\addvspace{4pt}

\reserve{7\baselineskip}
% mdparagraph:P371:begin
\subsection*{Exercise 34}\BookHeading{nav-ch16-learner-0027}{subsubsection}{Exercise 34}{0}
% mdparagraph:P371:end

\reserve{5\baselineskip}
\par\begin{minipage}{\linewidth}\raggedright
% mdparagraph:P372:begin
\textbf{Cook XII, XIII.34.1}
% mdparagraph:P372:end

% item: EX-XII-XIII-34-01
% mdparagraph:P373:begin
Whoever you are (193. c), forsake me not.
% mdparagraph:P373:end

% output: EX-XII-XIII-34-01-O001
\end{minipage}\par\addvspace{4pt}

\par\begin{minipage}{\linewidth}\raggedright
% mdparagraph:P374:begin
\textbf{Cook XII, XIII.34.2}
% mdparagraph:P374:end

% item: EX-XII-XIII-34-02
% mdparagraph:P375:begin
If you find nobody (196. d) that will feed you, return hither.
% mdparagraph:P375:end

% output: EX-XII-XIII-34-02-O001
\end{minipage}\par\addvspace{4pt}

\par\begin{minipage}{\linewidth}\raggedright
% mdparagraph:P376:begin
\textbf{Cook XII, XIII.34.3}
% mdparagraph:P376:end

% item: EX-XII-XIII-34-03
% mdparagraph:P377:begin
Let my books be sufficient for both of us (see p. 167, l. 13).
% mdparagraph:P377:end

% mdparagraph:P378:begin
\textbf{Supplied help:} Genyhtsumien mīna bēc unc bǣm.
% mdparagraph:P378:end

% mdparagraph:P379:begin
Let my books suffice for us both.
% mdparagraph:P379:end

% mdparagraph:P380:begin
The verb is plural optative, agreeing with bēc. Unc is dual dative; bǣm is dative plural because both has no separate dual form. This whole sentence remains supplied.
% mdparagraph:P380:end

% mdparagraph:P381:begin
\textbf{Supplied words.} These forms and meanings remain available for this task.
% mdparagraph:P381:end

% mdtable:T72:begin
\par\begingroup
\fontsize{10.5}{13.2}\selectfont
\setlength{\LTleft}{0pt}\setlength{\LTright}{\fill}
\noindent
\begin{tabular}{@{}>{\raggedright\arraybackslash}p{25.97mm}>{\raggedright\arraybackslash}p{37.22mm}>{\raggedright\arraybackslash}p{55.18mm}@{}}
\toprule
\textbf{Form} & \textbf{Meaning} & \textbf{Role} \\
\midrule
genyhtsumien & suffice; be sufficient for & third person, present, subjunctive, plural \\
\bottomrule
\end{tabular}
\par\endgroup

% mdtable:T72:end
% output: EX-XII-XIII-34-03-O001
\end{minipage}\par\addvspace{4pt}

\par\begin{minipage}{\linewidth}\raggedright
% mdparagraph:P382:begin
\textbf{Cook XII, XIII.34.4}
% mdparagraph:P382:end

% item: EX-XII-XIII-34-04
% mdparagraph:P383:begin
Both the boys are well educated.
% mdparagraph:P383:end

% output: EX-XII-XIII-34-04-O001
\end{minipage}\par\addvspace{4pt}

\par\begin{minipage}{\linewidth}\raggedright
% mdparagraph:P384:begin
\textbf{Cook XII, XIII.34.5}
% mdparagraph:P384:end

% item: EX-XII-XIII-34-05
% mdparagraph:P385:begin
They were both righteous before (beforan, 166) God.
% mdparagraph:P385:end

% output: EX-XII-XIII-34-05-O001
\end{minipage}\par\addvspace{4pt}

\par\begin{minipage}{\linewidth}\raggedright
% mdparagraph:P386:begin
\textbf{Cook XII, XIII.34.6}
% mdparagraph:P386:end

% item: EX-XII-XIII-34-06
% mdparagraph:P387:begin
I conjure thee by (ðurh) the living God, that thou tell us if thou be Christ, (om. the) Son of God.
% mdparagraph:P387:end

% mdparagraph:P388:begin
\textbf{Supplied help:} ðurh ðone libbendan God
% mdparagraph:P388:end

% mdparagraph:P389:begin
by the living God
% mdparagraph:P389:end

% mdparagraph:P390:begin
Libbendan is the supplied weak masculine accusative singular participle agreeing with God after ðone. The finite forms of libban do not supply this adjective position.
% mdparagraph:P390:end

% mdparagraph:P391:begin
\textbf{Supplied words.} These forms and meanings remain available for this task.
% mdparagraph:P391:end

% mdtable:T73:begin
\par\begingroup
\fontsize{10.5}{13.2}\selectfont
\setlength{\LTleft}{0pt}\setlength{\LTright}{\fill}
\noindent
\begin{tabular}{@{}>{\raggedright\arraybackslash}p{16.74mm}>{\raggedright\arraybackslash}p{47.8mm}>{\raggedright\arraybackslash}p{53.83mm}@{}}
\toprule
\textbf{Form} & \textbf{Meaning} & \textbf{Role} \\
\midrule
libbendan & living, in the supplied weak participial adjective & present, participle, accusative, singular, masculine, weak \\
\bottomrule
\end{tabular}
\par\endgroup

% mdtable:T73:end
% output: EX-XII-XIII-34-06-O001
\end{minipage}\par\addvspace{4pt}

\par\begin{minipage}{\linewidth}\raggedright
% mdparagraph:P392:begin
\textbf{Cook XII, XIII.34.7}
% mdparagraph:P392:end

% item: EX-XII-XIII-34-07
% mdparagraph:P393:begin
I tell you, After these (om. things), ye shall see the Son of Man sitting on (on, 172) the right hand of (om. the) majesty of God, and coming in (om. the) clouds of heaven.
% mdparagraph:P393:end

% mdparagraph:P394:begin
\textbf{Supplied help:} on ðā swīðran hand Godes mægenðrymmes
% mdparagraph:P394:end

% mdparagraph:P395:begin
at the right hand of God’s majesty
% mdparagraph:P395:end

% mdparagraph:P396:begin
On has an accusative in this fixed position expression; it does not express motion by the seated person. The perceived ongoing actions use plain infinitives sittan and cuman. English -ing does not automatically require an Old English participle.
% mdparagraph:P396:end

% mdparagraph:P397:begin
\textbf{Supplied words.} These forms and meanings remain available for this task.
% mdparagraph:P397:end

% mdtable:T74:begin
\par\begingroup
\fontsize{10.5}{13.2}\selectfont
\setlength{\LTleft}{0pt}\setlength{\LTright}{\fill}
\noindent
\begin{tabular}{@{}>{\raggedright\arraybackslash}p{19.83mm}>{\raggedright\arraybackslash}p{78.72mm}>{\raggedright\arraybackslash}p{19.83mm}@{}}
\toprule
\textbf{Form} & \textbf{Meaning} & \textbf{Role} \\
\midrule
on & at a right/\allowbreak{}left side, with accusative hand in this bounded source expression & as used here \\
\bottomrule
\end{tabular}
\par\endgroup

% mdtable:T74:end
% output: EX-XII-XIII-34-07-O001
\end{minipage}\par\addvspace{4pt}

\par\begin{minipage}{\linewidth}\raggedright
% mdparagraph:P398:begin
\textbf{Cook XII, XIII.34.8}
% mdparagraph:P398:end

% item: EX-XII-XIII-34-08
% mdparagraph:P399:begin
Be (om. it unto) you two (dual, 81) according to (æfter) your faith.
% mdparagraph:P399:end

% output: EX-XII-XIII-34-08-O001
\end{minipage}\par\addvspace{4pt}

\par\begin{minipage}{\linewidth}\raggedright
% mdparagraph:P400:begin
\textbf{Cook XII, XIII.34.9}
% mdparagraph:P400:end

% item: EX-XII-XIII-34-09
% mdparagraph:P401:begin
Truly the Son of Man is to come (200) in (on, 166) his Father’s glory (wuldor), and then (ðonne) (om. shall) reward (āgyltan, 113) every one (164) according to (be) his own work.
% mdparagraph:P401:end

% mdparagraph:P402:begin
\textbf{Supplied help:} is tō cumanne
% mdparagraph:P402:end

% mdparagraph:P403:begin
is to come
% mdparagraph:P403:end

% mdparagraph:P404:begin
This supplied wesan + tō + inflected-infinitive phrase expresses impending coming. Knowing the ordinary purpose infinitive does not by itself establish this use. Retain the complete expression here.
% mdparagraph:P404:end

% mdparagraph:P405:begin
\textbf{Supplied help:} ðonne āgylt hē ǣlcum be his āgnum weorce
% mdparagraph:P405:end

% mdparagraph:P406:begin
then he rewards everyone according to his own work
% mdparagraph:P406:end

% mdparagraph:P407:begin
Āgylt is the contracted third singular of strong āgyldan, and ǣlcum is the dative recipient. The whole clause remains supplied; see the correction above.
% mdparagraph:P407:end

\end{minipage}\par\addvspace{4pt}

% teaching: NEW-U16-SUP-AGEN
% teaching: NEW-U16-VOC-AGEN
\reserve{10\baselineskip}
% mdparagraph:P408:begin
\textbf{Supplied words.} These forms and meanings remain available for this task.
% mdparagraph:P408:end

% mdtable:T75:begin
\par\begingroup
\fontsize{10.5}{13.2}\selectfont
\setlength{\LTleft}{0pt}\setlength{\LTright}{\fill}
\noindent
\begin{tabular}{@{}>{\raggedright\arraybackslash}p{21.06mm}>{\raggedright\arraybackslash}p{43.81mm}>{\raggedright\arraybackslash}p{53.5mm}@{}}
\toprule
\textbf{Form} & \textbf{Meaning} & \textbf{Role} \\
\midrule
āgylt & reward; repay & third person, present, indicative, singular \\
āgnum & own & dative, singular, neuter \\
on & in, in the expression in his glory & as used here \\
\bottomrule
\end{tabular}
\par\endgroup

% mdtable:T75:end
% output: EX-XII-XIII-34-09-O001
\reserve{5\baselineskip}
\par\begin{minipage}{\linewidth}\raggedright
% mdparagraph:P409:begin
\textbf{Cook XII, XIII.34.10}
% mdparagraph:P409:end

% item: EX-XII-XIII-34-10
% mdparagraph:P410:begin
When (ðonne) Christ cometh, then knoweth no man whence he is.
% mdparagraph:P410:end

% output: EX-XII-XIII-34-10-O001
\end{minipage}\par\addvspace{4pt}

\par\begin{minipage}{\linewidth}\raggedright
% mdparagraph:P411:begin
\textbf{Cook XII, XIII.34.11}
% mdparagraph:P411:end

% item: EX-XII-XIII-34-11
% mdparagraph:P412:begin
Jesus cried (cleopian) and taught (lǣran) in the temple, and said, Me ye know (cunnan), and ye know (witan) whence I am; and I came not from myself [me self], but he is true that (87. d) sent me, whom (87. a) ye know not.
% mdparagraph:P412:end

% mdparagraph:P413:begin
\textbf{Supplied help:} sōð
% mdparagraph:P413:end

% mdparagraph:P414:begin
True: supplied nominative masculine singular in predicate use. No complete new adjective system is required.
% mdparagraph:P414:end

\end{minipage}\par\addvspace{4pt}

% teaching: NEW-U16-SUP-SOTH-CITATION
% teaching: NEW-U16-VOC-SOTH-CITATION
\reserve{9\baselineskip}
% mdparagraph:P415:begin
\textbf{Supplied words.} These forms and meanings remain available for this task.
% mdparagraph:P415:end

% mdtable:T76:begin
\par\begingroup
\fontsize{10.5}{13.2}\selectfont
\setlength{\LTleft}{0pt}\setlength{\LTright}{\fill}
\noindent
\begin{tabular}{@{}>{\raggedright\arraybackslash}p{29.59mm}>{\raggedright\arraybackslash}p{29.59mm}>{\raggedright\arraybackslash}p{59.19mm}@{}}
\toprule
\textbf{Form} & \textbf{Meaning} & \textbf{Role} \\
\midrule
sōð & true & nominative, singular, masculine \\
sē & that; those & nominative, singular, masculine \\
\bottomrule
\end{tabular}
\par\endgroup

% mdtable:T76:end
% output: EX-XII-XIII-34-11-O001
\reserve{5\baselineskip}
\par\begin{minipage}{\linewidth}\raggedright
% mdparagraph:P416:begin
\textbf{Cook XII, XIII.34.12}
% mdparagraph:P416:end

% item: EX-XII-XIII-34-12
% mdparagraph:P417:begin
Ye have not known (pret.) him; I know him, and if I say that I know (194) him not (verb last), I (om. shall) be like (om. unto) you (165); but I know him, and I keep (healdan) his saying (sprǣc).
% mdparagraph:P417:end

% output: EX-XII-XIII-34-12-O001
\end{minipage}\par\addvspace{4pt}

\par\begin{minipage}{\linewidth}\raggedright
% mdparagraph:P418:begin
\textbf{Cook XII, XIII.34.13}
% mdparagraph:P418:end

% item: EX-XII-XIII-34-13
% mdparagraph:P419:begin
If I glorify myself, my glory is (183) naught (verb first); my Father (om. it) is, that (87. d) glorifieth me, concerning whom (87. a) ye say (cweðan) that he is (194) your God.
% mdparagraph:P419:end

% mdparagraph:P420:begin
\textbf{Supplied help:} nāwiht
% mdparagraph:P420:end

% mdparagraph:P421:begin
Nothing; naught. Here nāwiht is nominative in the predicate after nis, unlike its earlier accusative object use.
% mdparagraph:P421:end

\end{minipage}\par\addvspace{4pt}

% teaching: NEW-U16-SUP-NAWIHT
% teaching: NEW-U16-VOC-NAWIHT
\reserve{9\baselineskip}
% mdparagraph:P422:begin
\textbf{Supplied words.} These forms and meanings remain available for this task.
% mdparagraph:P422:end

% mdtable:T77:begin
\par\begingroup
\fontsize{10.5}{13.2}\selectfont
\setlength{\LTleft}{0pt}\setlength{\LTright}{\fill}
\noindent
\begin{tabular}{@{}>{\raggedright\arraybackslash}p{28.77mm}>{\raggedright\arraybackslash}p{32.08mm}>{\raggedright\arraybackslash}p{57.53mm}@{}}
\toprule
\textbf{Form} & \textbf{Meaning} & \textbf{Role} \\
\midrule
nāwiht & nothing; naught & nominative, singular, neuter \\
sē & that; those & nominative, singular, masculine \\
\bottomrule
\end{tabular}
\par\endgroup

% mdtable:T77:end
% output: EX-XII-XIII-34-13-O001
\reserve{7\baselineskip}
% mdparagraph:P423:begin
\subsection*{Exercise 35}\BookHeading{nav-ch16-learner-0028}{subsubsection}{Exercise 35}{0}
% mdparagraph:P423:end

\reserve{5\baselineskip}
\par\begin{minipage}{\linewidth}\raggedright
% mdparagraph:P424:begin
\textbf{Cook XII, XIII.35.1}
% mdparagraph:P424:end

% item: EX-XII-XIII-35-01
% mdparagraph:P425:begin
Catch this ball.
% mdparagraph:P425:end

% output: EX-XII-XIII-35-01-O001
\end{minipage}\par\addvspace{4pt}

\par\begin{minipage}{\linewidth}\raggedright
% mdparagraph:P426:begin
\textbf{Cook XII, XIII.35.2}
% mdparagraph:P426:end

% item: EX-XII-XIII-35-02
% mdparagraph:P427:begin
Who caught the black ball?
% mdparagraph:P427:end

% output: EX-XII-XIII-35-02-O001
\end{minipage}\par\addvspace{4pt}

\par\begin{minipage}{\linewidth}\raggedright
% mdparagraph:P428:begin
\textbf{Cook XII, XIII.35.3}
% mdparagraph:P428:end

% item: EX-XII-XIII-35-03
% mdparagraph:P429:begin
Go the way that (87. d) you came.
% mdparagraph:P429:end

% output: EX-XII-XIII-35-03-O001
\end{minipage}\par\addvspace{4pt}

\par\begin{minipage}{\linewidth}\raggedright
% mdparagraph:P430:begin
\textbf{Cook XII, XIII.35.4}
% mdparagraph:P430:end

% item: EX-XII-XIII-35-04
% mdparagraph:P431:begin
I never played better than I did to-day.
% mdparagraph:P431:end

% output: EX-XII-XIII-35-04-O001
\end{minipage}\par\addvspace{4pt}

\par\begin{minipage}{\linewidth}\raggedright
% mdparagraph:P432:begin
\textbf{Cook XII, XIII.35.5}
% mdparagraph:P432:end

% item: EX-XII-XIII-35-05
% mdparagraph:P433:begin
She never sang better than she did now.
% mdparagraph:P433:end

% output: EX-XII-XIII-35-05-O001
\end{minipage}\par\addvspace{4pt}

\par\begin{minipage}{\linewidth}\raggedright
% mdparagraph:P434:begin
\textbf{Cook XII, XIII.35.6}
% mdparagraph:P434:end

% item: EX-XII-XIII-35-06
% mdparagraph:P435:begin
Thrice (ðrīwa) I suffered shipwreck; a night (170) and a day I was in this sea.
% mdparagraph:P435:end

% output: EX-XII-XIII-35-06-O001
\end{minipage}\par\addvspace{4pt}

\par\begin{minipage}{\linewidth}\raggedright
% mdparagraph:P436:begin
\textbf{Cook XII, XIII.35.7}
% mdparagraph:P436:end

% item: EX-XII-XIII-35-07
% mdparagraph:P437:begin
I knew (witan) a certain man that (87. d) was carried (lǣdan) unto (ōð, 172) the third heaven, and there (ðǣr) heard the words that no earthly man may (not) speak (sprecan).
% mdparagraph:P437:end

% mdparagraph:P438:begin
\textbf{Supplied help:} ōð ðone ðriddan heofon
% mdparagraph:P438:end

% mdparagraph:P439:begin
as far as the third heaven; exact preposition plus accusative goal supplied.
% mdparagraph:P439:end

\end{minipage}\par\addvspace{4pt}

% teaching: NEW-U16-SUP-OTH
% teaching: NEW-U16-VOC-OTH
\reserve{9\baselineskip}
% mdparagraph:P440:begin
\textbf{Supplied words.} These forms and meanings remain available for this task.
% mdparagraph:P440:end

% mdtable:T78:begin
\par\begingroup
\fontsize{10.5}{13.2}\selectfont
\setlength{\LTleft}{0pt}\setlength{\LTright}{\fill}
\noindent
\begin{tabular}{@{}>{\raggedright\arraybackslash}p{22.74mm}>{\raggedright\arraybackslash}p{50.15mm}>{\raggedright\arraybackslash}p{45.49mm}@{}}
\toprule
\textbf{Form} & \textbf{Meaning} & \textbf{Role} \\
\midrule
ōð & as far as; up to & invariable \\
sē & who; the one who, referring to the man & nominative, singular, masculine \\
\bottomrule
\end{tabular}
\par\endgroup

% mdtable:T78:end
% output: EX-XII-XIII-35-07-O001
\reserve{5\baselineskip}
\par\begin{minipage}{\linewidth}\raggedright
% mdparagraph:P441:begin
\textbf{Cook XII, XIII.35.8}
% mdparagraph:P441:end

% item: EX-XII-XIII-35-08
% mdparagraph:P442:begin
O, how narrow (nearu, 57. 5) is the gate and the way that (87. d) leads to life (verb last), and very few (om. there) are that find the way.
% mdparagraph:P442:end

% mdparagraph:P443:begin
\textbf{Supplied help:} hū nearu is ðæt geat and se weg ðe tō līfe lǣtt
% mdparagraph:P443:end

% mdparagraph:P444:begin
how narrow is the gate and the way that leads to life
% mdparagraph:P444:end

% mdparagraph:P445:begin
This wording treats gate and way as a combined figurative route and retains a singular predicate. Use this supplied coordination here; it is not a general singular-agreement rule for coordinated subjects.
% mdparagraph:P445:end

% mdparagraph:P446:begin
\textbf{Supplied help:} fēawe
% mdparagraph:P446:end

% mdparagraph:P447:begin
few people; masculine nominative plural substantive supplied, with no general new adjective paradigm.
% mdparagraph:P447:end

\end{minipage}\par\addvspace{4pt}

% teaching: NEW-U16-SUP-FEAWE
% teaching: NEW-U16-VOC-FEAWE
\reserve{8\baselineskip}
% mdparagraph:P448:begin
\textbf{Supplied words.} These forms and meanings remain available for this task.
% mdparagraph:P448:end

% mdtable:T79:begin
\par\begingroup
\fontsize{10.5}{13.2}\selectfont
\setlength{\LTleft}{0pt}\setlength{\LTright}{\fill}
\noindent
\begin{tabular}{@{}>{\raggedright\arraybackslash}p{30.28mm}>{\raggedright\arraybackslash}p{30.28mm}>{\raggedright\arraybackslash}p{57.83mm}@{}}
\toprule
\textbf{Form} & \textbf{Meaning} & \textbf{Role} \\
\midrule
fēawe & few people & nominative, plural, masculine \\
\bottomrule
\end{tabular}
\par\endgroup

% mdtable:T79:end
% output: EX-XII-XIII-35-08-O001
\reserve{5\baselineskip}
\par\begin{minipage}{\linewidth}\raggedright
% mdparagraph:P449:begin
\textbf{Cook XII, XIII.35.9}
% mdparagraph:P449:end

% item: EX-XII-XIII-35-09
% mdparagraph:P450:begin
They come to (tō, 166) you in sheep’s garments (gegierela), but within (innan) they are wolves.
% mdparagraph:P450:end

% mdparagraph:P451:begin
\textbf{Supplied help:} on scēapa gegierelum
% mdparagraph:P451:end

% mdparagraph:P452:begin
in sheep’s clothing; the local clothing use of on is supplied.
% mdparagraph:P452:end

% mdparagraph:P453:begin
\textbf{Supplied words.} These forms and meanings remain available for this task.
% mdparagraph:P453:end

% mdtable:T80:begin
\par\begingroup
\fontsize{10.5}{13.2}\selectfont
\setlength{\LTleft}{0pt}\setlength{\LTright}{\fill}
\noindent
\begin{tabular}{@{}>{\raggedright\arraybackslash}p{23.77mm}>{\raggedright\arraybackslash}p{70.84mm}>{\raggedright\arraybackslash}p{23.77mm}@{}}
\toprule
\textbf{Form} & \textbf{Meaning} & \textbf{Role} \\
\midrule
on & in; wearing clothing, in supplied on scēapa gegierelum & as used here \\
\bottomrule
\end{tabular}
\par\endgroup

% mdtable:T80:end
% output: EX-XII-XIII-35-09-O001
\end{minipage}\par\addvspace{4pt}

\par\begin{minipage}{\linewidth}\raggedright
% mdparagraph:P454:begin
\textbf{Cook XII, XIII.35.10}
% mdparagraph:P454:end

% item: EX-XII-XIII-35-10
% mdparagraph:P455:begin
When he came nigh (nēalǣcan) to the gate (163) of the city, then (ðā) was there (ðǣr) a dead man carried (beran).
% mdparagraph:P455:end

% output: EX-XII-XIII-35-10-O001
\end{minipage}\par\addvspace{4pt}

\par\begin{minipage}{\linewidth}\raggedright
% mdparagraph:P456:begin
\textbf{Cook XII, XIII.35.11}
% mdparagraph:P456:end

% item: EX-XII-XIII-35-11
% mdparagraph:P457:begin
I see a girl running (169) along (172) the street.
% mdparagraph:P457:end

% output: EX-XII-XIII-35-11-O001
\end{minipage}\par\addvspace{4pt}

\par\begin{minipage}{\linewidth}\raggedright
% mdparagraph:P458:begin
\textbf{Cook XII, XIII.35.12}
% mdparagraph:P458:end

% item: EX-XII-XIII-35-12
% mdparagraph:P459:begin
Anoint me with oil.
% mdparagraph:P459:end

% output: EX-XII-XIII-35-12-O001
\end{minipage}\par\addvspace{4pt}

\par\begin{minipage}{\linewidth}\raggedright
% mdparagraph:P460:begin
\textbf{Cook XII, XIII.35.13}
% mdparagraph:P460:end

% item: EX-XII-XIII-35-13
% mdparagraph:P461:begin
Give (sęllan) us of (of, 166) your oil.
% mdparagraph:P461:end

% mdparagraph:P462:begin
\textbf{Supplied help:} of ēowrum ęle
% mdparagraph:P462:end

% mdparagraph:P463:begin
of your oil; source partition/\allowbreak{}substance use of of is supplied.
% mdparagraph:P463:end

% mdparagraph:P464:begin
\textbf{Supplied words.} These forms and meanings remain available for this task.
% mdparagraph:P464:end

% mdtable:T81:begin
\par\begingroup
\fontsize{10.5}{13.2}\selectfont
\setlength{\LTleft}{0pt}\setlength{\LTright}{\fill}
\noindent
\begin{tabular}{@{}>{\raggedright\arraybackslash}p{19.31mm}>{\raggedright\arraybackslash}p{79.75mm}>{\raggedright\arraybackslash}p{19.31mm}@{}}
\toprule
\textbf{Form} & \textbf{Meaning} & \textbf{Role} \\
\midrule
of & of; from a portion of a whole, in the supplied phrases ān of ðǣm and of ēowrum ęle & as used here \\
\bottomrule
\end{tabular}
\par\endgroup

% mdtable:T81:end
% output: EX-XII-XIII-35-13-O001
\end{minipage}\par\addvspace{4pt}

\par\begin{minipage}{\linewidth}\raggedright
% mdparagraph:P465:begin
\textbf{Cook XII, XIII.35.14}
% mdparagraph:P465:end

% item: EX-XII-XIII-35-14
% mdparagraph:P466:begin
If the Son delivers you, ye (om. shall) indeed be free.
% mdparagraph:P466:end

% mdparagraph:P467:begin
\textbf{Supplied help:} frige
% mdparagraph:P467:end

% mdparagraph:P468:begin
free, nominative plural; Cook’s glossary explicitly gives this irregular plural.
% mdparagraph:P468:end

\end{minipage}\par\addvspace{4pt}

% teaching: NEW-U16-SUP-FREO
% teaching: NEW-U16-VOC-FREO
\reserve{8\baselineskip}
% mdparagraph:P469:begin
\textbf{Supplied words.} These forms and meanings remain available for this task.
% mdparagraph:P469:end

% mdtable:T82:begin
\par\begingroup
\fontsize{10.5}{13.2}\selectfont
\setlength{\LTleft}{0pt}\setlength{\LTright}{\fill}
\noindent
\begin{tabular}{@{}>{\raggedright\arraybackslash}p{30.28mm}>{\raggedright\arraybackslash}p{30.28mm}>{\raggedright\arraybackslash}p{57.83mm}@{}}
\toprule
\textbf{Form} & \textbf{Meaning} & \textbf{Role} \\
\midrule
frige & free & nominative, plural, masculine \\
\bottomrule
\end{tabular}
\par\endgroup

% mdtable:T82:end
% output: EX-XII-XIII-35-14-O001
\reserve{7\baselineskip}
% mdparagraph:P470:begin
\subsection*{Exercise 36}\BookHeading{nav-ch16-learner-0029}{subsubsection}{Exercise 36}{0}
% mdparagraph:P470:end

\reserve{5\baselineskip}
\par\begin{minipage}{\linewidth}\raggedright
% mdparagraph:P471:begin
\textbf{Cook XII, XIII.36.1}
% mdparagraph:P471:end

% item: EX-XII-XIII-36-01
% mdparagraph:P472:begin
Judge ye not according to (be) appearance (ansīen, 51. b), but judge right judgment (168. 1).
% mdparagraph:P472:end

% output: EX-XII-XIII-36-01-O001
\end{minipage}\par\addvspace{4pt}

\par\begin{minipage}{\linewidth}\raggedright
% mdparagraph:P473:begin
\textbf{Cook XII, XIII.36.2}
% mdparagraph:P473:end

% item: EX-XII-XIII-36-02
% mdparagraph:P474:begin
Then turned (bewęndan, 184. a) he to the woman, and said, Seest thou this woman? I went into (intō) thy house, thou gavest me not water for (tō) my feet; this (om. one hath) washed (101. VI.) my feet with (mid) her tears.
% mdparagraph:P474:end

% mdparagraph:P475:begin
\textbf{Supplied help:} ðēos
% mdparagraph:P475:end

% mdparagraph:P476:begin
This woman. The feminine pronoun refers to the woman as a person, although wīf is grammatically neuter. Retain this supplied reference in the exchange.
% mdparagraph:P476:end

% mdparagraph:P477:begin
\textbf{Supplied help:} wæter tō mīnum fōtum
% mdparagraph:P477:end

% mdparagraph:P478:begin
water for my feet: the purpose of washing/\allowbreak{}caring for the body part is supplied; tō does not mean motion toward a destination here.
% mdparagraph:P478:end

\end{minipage}\par\addvspace{4pt}

% teaching: NEW-U16-VOC-TO
\reserve{8\baselineskip}
% mdparagraph:P479:begin
\textbf{Supplied words.} These forms and meanings remain available for this task.
% mdparagraph:P479:end

% mdtable:T83:begin
\par\begingroup
\fontsize{10.5}{13.2}\selectfont
\setlength{\LTleft}{0pt}\setlength{\LTright}{\fill}
\noindent
\begin{tabular}{@{}>{\raggedright\arraybackslash}p{22.83mm}>{\raggedright\arraybackslash}p{72.72mm}>{\raggedright\arraybackslash}p{22.83mm}@{}}
\toprule
\textbf{Form} & \textbf{Meaning} & \textbf{Role} \\
\midrule
tō & for the care of a body part, in supplied wæter tō mīnum fōtum & as used here \\
\bottomrule
\end{tabular}
\par\endgroup

% mdtable:T83:end
% output: EX-XII-XIII-36-02-O001
\reserve{5\baselineskip}
\par\begin{minipage}{\linewidth}\raggedright
% mdparagraph:P480:begin
\textbf{Cook XII, XIII.36.3}
% mdparagraph:P480:end

% item: EX-XII-XIII-36-03
% mdparagraph:P481:begin
Thou gavest (verb last) me not (om. a) kiss (coss, 43); this (om. one), since (siððan) I entered, (om. hath) ceased not to kiss [that she kissed] my feet.
% mdparagraph:P481:end

% mdparagraph:P482:begin
\textbf{Supplied help:} ne geswāc ðæt hēo mīne fēt ne cyste
% mdparagraph:P482:end

% mdparagraph:P483:begin
did not stop kissing my feet
% mdparagraph:P483:end

% mdparagraph:P484:begin
Both negatives belong to this attested cessation construction: the second ne does not mean that she withheld kisses. Keep the complete finite clause supplied here.
% mdparagraph:P484:end

% mdparagraph:P485:begin
\textbf{Supplied help:} ðēos
% mdparagraph:P485:end

% mdparagraph:P486:begin
This woman. The feminine pronoun refers to the woman as a person, although wīf is grammatically neuter. Retain this supplied reference in the exchange.
% mdparagraph:P486:end

% output: EX-XII-XIII-36-03-O001
\end{minipage}\par\addvspace{4pt}

\par\begin{minipage}{\linewidth}\raggedright
% mdparagraph:P487:begin
\textbf{Cook XII, XIII.36.4}
% mdparagraph:P487:end

% item: EX-XII-XIII-36-04
% mdparagraph:P488:begin
My head thou with oils (om. didst) not anoint (smierwan, 113); this one anointed my feet with ointment (sealf, 51. b).
% mdparagraph:P488:end

% mdparagraph:P489:begin
\textbf{Supplied help:} ðēos
% mdparagraph:P489:end

% mdparagraph:P490:begin
This woman. The feminine pronoun refers to the woman as a person, although wīf is grammatically neuter. Retain this supplied reference in the exchange.
% mdparagraph:P490:end

% mdparagraph:P491:begin
\textbf{Supplied help:} mid ęlum
% mdparagraph:P491:end

% mdparagraph:P492:begin
With oils. The plural dative ęlum remains supplied.
% mdparagraph:P492:end

\end{minipage}\par\addvspace{4pt}

% teaching: NEW-U16-SUP-ELE
% teaching: NEW-U16-VOC-ELE
\reserve{8\baselineskip}
% mdparagraph:P493:begin
\textbf{Supplied words.} These forms and meanings remain available for this task.
% mdparagraph:P493:end

% mdtable:T84:begin
\par\begingroup
\fontsize{10.5}{13.2}\selectfont
\setlength{\LTleft}{0pt}\setlength{\LTright}{\fill}
\noindent
\begin{tabular}{@{}>{\raggedright\arraybackslash}p{32.7mm}>{\raggedright\arraybackslash}p{32.7mm}>{\raggedright\arraybackslash}p{52.97mm}@{}}
\toprule
\textbf{Form} & \textbf{Meaning} & \textbf{Role} \\
\midrule
ęlum & oil & dative, plural, masculine \\
\bottomrule
\end{tabular}
\par\endgroup

% mdtable:T84:end
% output: EX-XII-XIII-36-04-O001
\reserve{5\baselineskip}
\par\begin{minipage}{\linewidth}\raggedright
% mdparagraph:P494:begin
\textbf{Cook XII, XIII.36.5}
% mdparagraph:P494:end

% item: EX-XII-XIII-36-05
% mdparagraph:P495:begin
Therefore I say (sęcgean) to thee (164. b), Many sins are forgiven (forgiefan) her, because she loved (verb last) me much (swīðe).
% mdparagraph:P495:end

% output: EX-XII-XIII-36-05-O001
\end{minipage}\par\addvspace{4pt}

\par\begin{minipage}{\linewidth}\raggedright
% mdparagraph:P496:begin
\textbf{Cook XII, XIII.36.6}
% mdparagraph:P496:end

% item: EX-XII-XIII-36-06
% mdparagraph:P497:begin
Less (neut. acc. sing.) loveth (om. to) whom (87. b) less forgiven is.
% mdparagraph:P497:end

% mdparagraph:P498:begin
\textbf{Supplied help:} Lǣs lufað ðǣm ðe lǣs forgiefen is.
% mdparagraph:P498:end

% mdparagraph:P499:begin
The one to whom less is forgiven loves less.
% mdparagraph:P499:end

% mdparagraph:P500:begin
The matrix subject is understood from the attracted dative relative. The relative dative belongs to forgiveness; it must not be generalized into a dative subject of ordinary lufian. The prompt attaches neuter accusative singular specifically to its first Less: this is the adverbial degree accusative. The second lǣs is nominative neuter singular, the smaller amount that is forgiven, and is the passive subject. The entire compressed sentence is supplied.
% mdparagraph:P500:end

% mdparagraph:P501:begin
\textbf{Supplied help:} lǣs ... lǣs
% mdparagraph:P501:end

% mdparagraph:P502:begin
First lǣs is the specified adverbial accusative neuter singular; second lǣs is the nominative neuter singular amount forgiven. Both exact cells belong to the supplied clause.
% mdparagraph:P502:end

\end{minipage}\par\addvspace{4pt}

% teaching: NEW-U16-SUP-LAES
% teaching: NEW-U16-VOC-LAES
\reserve{9\baselineskip}
% mdparagraph:P503:begin
\textbf{Supplied words.} These forms and meanings remain available for this task.
% mdparagraph:P503:end

% mdtable:T85:begin
\par\begingroup
\fontsize{10.5}{13.2}\selectfont
\setlength{\LTleft}{0pt}\setlength{\LTright}{\fill}
\noindent
\begin{tabular}{@{}>{\raggedright\arraybackslash}p{27.92mm}>{\raggedright\arraybackslash}p{38.39mm}>{\raggedright\arraybackslash}p{52.07mm}@{}}
\toprule
\textbf{Form} & \textbf{Meaning} & \textbf{Role} \\
\midrule
lǣs & less; a smaller amount & accusative, singular, neuter \\
lǣs & less; a smaller amount & nominative, singular, neuter \\
\bottomrule
\end{tabular}
\par\endgroup

% mdtable:T85:end
% output: EX-XII-XIII-36-06-O001
\reserve{5\baselineskip}
\par\begin{minipage}{\linewidth}\raggedright
% mdparagraph:P504:begin
\textbf{Cook XII, XIII.36.7}
% mdparagraph:P504:end

% item: EX-XII-XIII-36-07
% mdparagraph:P505:begin
It (om. is) not permitted (ālīefan) thee to have her to wife.
% mdparagraph:P505:end

% mdparagraph:P506:begin
\textbf{Supplied help:} Hit ne ālīefð ðē hīe tō wīfe habban.
% mdparagraph:P506:end

% mdparagraph:P507:begin
It is not permitted for you to have her as wife.
% mdparagraph:P507:end

% mdparagraph:P508:begin
Impersonal hit with finite ālīefð, dative permitted person ðē, accusative hīe as object of habban, and tō wīfe as the wife-role phrase. Omit the English is; do not add wesan or turn hīe into a dative object.
% mdparagraph:P508:end

% mdparagraph:P509:begin
\textbf{Supplied help:} tō wīfe
% mdparagraph:P509:end

% mdparagraph:P510:begin
as wife; in the role of wife. This is part of the supplied permission construction, not a destination.
% mdparagraph:P510:end

% mdparagraph:P511:begin
\textbf{Supplied words.} These forms and meanings remain available for this task.
% mdparagraph:P511:end

% mdtable:T86:begin
\par\begingroup
\fontsize{10.5}{13.2}\selectfont
\setlength{\LTleft}{0pt}\setlength{\LTright}{\fill}
\noindent
\begin{tabular}{@{}>{\raggedright\arraybackslash}p{18.8mm}>{\raggedright\arraybackslash}p{51.81mm}>{\raggedright\arraybackslash}p{47.76mm}@{}}
\toprule
\textbf{Form} & \textbf{Meaning} & \textbf{Role} \\
\midrule
ālīefð & permit; allow, in impersonal permission & third person, present, indicative, singular \\
tō & as wife; in the role of wife, in supplied tō wīfe & as used here \\
\bottomrule
\end{tabular}
\par\endgroup

% mdtable:T86:end
% output: EX-XII-XIII-36-07-O001
\end{minipage}\par\addvspace{4pt}

\par\begin{minipage}{\linewidth}\raggedright
% mdparagraph:P512:begin
\textbf{Cook XII, XIII.36.8}
% mdparagraph:P512:end

% item: EX-XII-XIII-36-08
% mdparagraph:P513:begin
Then he promised with (mid) oath (āð, 43) to give (sęllan) her whatsoever she (om. should) ask (biddan) him.
% mdparagraph:P513:end

% mdparagraph:P514:begin
\textbf{Supplied help:} Ðā gehēt hē mid āðe hiere ðæt tō sęllanne, swā hwæs swā hēo hine bǣde.
% mdparagraph:P514:end

% mdparagraph:P515:begin
Then he promised with an oath to give her whatever she asked him for.
% mdparagraph:P515:end

% mdparagraph:P516:begin
Gehēt introduces the promise and the intended-action infinitive. Hiere is the recipient of giving, ðæt the accusative thing given. Within the free-choice request, biddan takes genitive hwæs and accusative hine; bǣde is preterite subjunctive in the open request. The resumptive ðæt lets each clause retain its own case requirement.
% mdparagraph:P516:end

\end{minipage}\par\addvspace{4pt}

% teaching: NEW-U16-VOC-MID
\reserve{8\baselineskip}
% mdparagraph:P517:begin
\textbf{Supplied words.} These forms and meanings remain available for this task.
% mdparagraph:P517:end

% mdtable:T87:begin
\par\begingroup
\fontsize{10.5}{13.2}\selectfont
\setlength{\LTleft}{0pt}\setlength{\LTright}{\fill}
\noindent
\begin{tabular}{@{}>{\raggedright\arraybackslash}p{20.53mm}>{\raggedright\arraybackslash}p{77.31mm}>{\raggedright\arraybackslash}p{20.53mm}@{}}
\toprule
\textbf{Form} & \textbf{Meaning} & \textbf{Role} \\
\midrule
mid & with; in the manner of the accompanying quality or mental disposition & as used here \\
\bottomrule
\end{tabular}
\par\endgroup

% mdtable:T87:end
% output: EX-XII-XIII-36-08-O001
\reserve{5\baselineskip}
\par\begin{minipage}{\linewidth}\raggedright
% mdparagraph:P518:begin
\textbf{Cook XII, XIII.36.9}
% mdparagraph:P518:end

% item: EX-XII-XIII-36-09
% mdparagraph:P519:begin
Then, admonished (myngian) by (fram) her mother, she said, Give me in a dish (disc, 43) the head of John (Johannes) the Baptist (fuluhtere, 44. 1).
% mdparagraph:P519:end

% mdparagraph:P520:begin
\textbf{Supplied help:} fram hiere mēder myngod
% mdparagraph:P520:end

% mdparagraph:P521:begin
admonished by her mother
% mdparagraph:P521:end

% mdparagraph:P522:begin
The past participle is supplied uninflected in its predicative attachment to hēo; fram plus dative marks the agent. This does not introduce a complete myngian system.
% mdparagraph:P522:end

% mdparagraph:P523:begin
\textbf{Supplied help:} Johannes /\allowbreak{} Johannem
% mdparagraph:P523:end

% mdparagraph:P524:begin
John: Johannes is genitive and Johannem accusative. The endings are retained from the attested forms; initial J follows Cook’s headword. See the note on names above.
% mdparagraph:P524:end

\end{minipage}\par\addvspace{4pt}

% teaching: NEW-U16-SUP-JOHANNES-NAME
% teaching: NEW-U16-VOC-JOHANNES-NAME
\reserve{10\baselineskip}
% mdparagraph:P525:begin
\textbf{Supplied words.} These forms and meanings remain available for this task.
% mdparagraph:P525:end

% mdtable:T88:begin
\par\begingroup
\fontsize{10.5}{13.2}\selectfont
\setlength{\LTleft}{0pt}\setlength{\LTright}{\fill}
\noindent
\begin{tabular}{@{}>{\raggedright\arraybackslash}p{26.28mm}>{\raggedright\arraybackslash}p{39.55mm}>{\raggedright\arraybackslash}p{52.55mm}@{}}
\toprule
\textbf{Form} & \textbf{Meaning} & \textbf{Role} \\
\midrule
myngod & admonished; prompted & past, participle \\
Johannes & John & genitive, singular, masculine \\
Johannem & John & accusative, singular, masculine \\
\bottomrule
\end{tabular}
\par\endgroup

% mdtable:T88:end
% output: EX-XII-XIII-36-09-O001
\reserve{5\baselineskip}
\par\begin{minipage}{\linewidth}\raggedright
% mdparagraph:P526:begin
\textbf{Cook XII, XIII.36.10}
% mdparagraph:P526:end

% item: EX-XII-XIII-36-10
% mdparagraph:P527:begin
Then he sent (āsęndan), and beheaded (behēafdian, 118) John in the prison.
% mdparagraph:P527:end

% mdparagraph:P528:begin
\textbf{Supplied help:} Johannes /\allowbreak{} Johannem
% mdparagraph:P528:end

% mdparagraph:P529:begin
John: Johannes is genitive and Johannem accusative. The endings are retained from the attested forms; initial J follows Cook’s headword. See the note on names above.
% mdparagraph:P529:end

% mdparagraph:P530:begin
\textbf{Supplied words.} These forms and meanings remain available for this task.
% mdparagraph:P530:end

% mdtable:T89:begin
\par\begingroup
\fontsize{10.5}{13.2}\selectfont
\setlength{\LTleft}{0pt}\setlength{\LTright}{\fill}
\noindent
\begin{tabular}{@{}>{\raggedright\arraybackslash}p{29.59mm}>{\raggedright\arraybackslash}p{29.59mm}>{\raggedright\arraybackslash}p{59.19mm}@{}}
\toprule
\textbf{Form} & \textbf{Meaning} & \textbf{Role} \\
\midrule
Johannes & John & genitive, singular, masculine \\
Johannem & John & accusative, singular, masculine \\
\bottomrule
\end{tabular}
\par\endgroup

% mdtable:T89:end
% output: EX-XII-XIII-36-10-O001
\end{minipage}\par\addvspace{4pt}

\par\begin{minipage}{\linewidth}\raggedright
% mdparagraph:P531:begin
\textbf{Cook XII, XIII.36.11}
% mdparagraph:P531:end

% item: EX-XII-XIII-36-11
% mdparagraph:P532:begin
And one brought his head in a dish.
% mdparagraph:P532:end

% output: EX-XII-XIII-36-11-O001
\end{minipage}\par\addvspace{4pt}

\reserve{7\baselineskip}
% mdparagraph:P533:begin
\section*{7. The Incidents in the Gymnasium}\BookHeading{nav-ch16-learner-0030}{subsection}{7. The Incidents in the Gymnasium}{1}
% mdparagraph:P533:end

% teaching: NEW-U16-READING-SUPPORT
% mdparagraph:P534:begin
Apollonius has survived a shipwreck. A fisherman has shared clothing with him and shown him the road to the city. This scene begins as he follows that road; it ends with the king’s invitation. The following feast is outside the passage.
% mdparagraph:P534:end

% mdparagraph:P535:begin
The text is Cook’s complete scene at printed pages 168.1–170.16. A locator such as 169.6 means printed page 169, editorially counted line 6, anchored by the edition’s printed line numbers. P169N6 instead means original note 6 on page 169. Original notes and modern help are kept separate; a question in an original note is not an additional exercise.\textsuperscript{9}
% mdparagraph:P535:end

\reserve{7\baselineskip}
% mdparagraph:P536:begin
\subsection*{Supplied vocabulary and constructions}\BookHeading{nav-ch16-learner-0031}{subsubsection}{Supplied vocabulary and constructions}{0}
% mdparagraph:P536:end

% mdparagraph:P537:begin
The following help is local to this reading. Retain the learned productive systems where they apply; the additional forms, spellings, meanings and constructions listed here may be consulted throughout the reading tasks.
% mdparagraph:P537:end

% teaching: NEW-U16-SUP-THES
% teaching: NEW-U16-VOC-THES
% teaching: NEW-U16-SUP-OTH-THAET
% teaching: NEW-U16-VOC-OTH-THAET
% teaching: NEW-U16-SUP-MID-THY
% teaching: NEW-U16-VOC-MID-THY
% teaching: NEW-U16-SUP-FULTUM
% teaching: NEW-U16-VOC-FULTUM
% teaching: NEW-U16-SUP-AN
% teaching: NEW-U16-VOC-AN
% teaching: NEW-U16-SUP-SCIETE
% teaching: NEW-U16-VOC-SCIETE
% teaching: NEW-U16-SUP-GEONG
% teaching: NEW-U16-VOC-GEONG
% teaching: NEW-U16-SUP-PLEGA
% teaching: NEW-U16-VOC-PLEGA
% teaching: NEW-U16-SUP-BAETHSTEDE
% teaching: NEW-U16-VOC-BAETHSTEDE
% teaching: NEW-U16-SUP-CEASTERWARAN
% teaching: NEW-U16-VOC-CEASTERWARAN
% teaching: NEW-U16-SUP-ELTHEODIG
% teaching: NEW-U16-VOC-ELTHEODIG
% teaching: NEW-U16-SUP-THEOW
% teaching: NEW-U16-VOC-THEOW
% teaching: NEW-U16-SUP-AETHELE
% teaching: NEW-U16-VOC-AETHELE
% teaching: NEW-U16-SUP-UNAETHELE
% teaching: NEW-U16-VOC-UNAETHELE
% teaching: NEW-U16-SUP-OPEN
% teaching: NEW-U16-VOC-OPEN
% teaching: NEW-U16-SUP-THA-CONJUNCTION
% teaching: NEW-U16-VOC-THA-CONJUNCTION
% teaching: NEW-U16-SUP-APOLLONIUS
% teaching: NEW-U16-VOC-APOLLONIUS
% teaching: NEW-U16-SUP-HEALF
% teaching: NEW-U16-VOC-HEALF
% teaching: NEW-U16-SUP-SCICCELS
% teaching: NEW-U16-VOC-SCICCELS
% teaching: NEW-U16-SUP-ON-ADVERB
% teaching: NEW-U16-VOC-ON-ADVERB
% teaching: NEW-U16-SUP-THWEAL
% teaching: NEW-U16-VOC-THWEAL
% teaching: NEW-U16-SUP-GELICA
% teaching: NEW-U16-VOC-GELICA
% teaching: NEW-U16-SUP-FLOCC
% teaching: NEW-U16-VOC-FLOCC
% teaching: NEW-U16-SUP-FAERINGA
% teaching: NEW-U16-VOC-FAERINGA
% teaching: NEW-U16-SUP-ARCESTRATES
% teaching: NEW-U16-VOC-ARCESTRATES
% teaching: NEW-U16-SUP-THEOD
% teaching: NEW-U16-VOC-THEOD
% teaching: NEW-U16-SUP-MENIGU
% teaching: NEW-U16-VOC-MENIGU
% teaching: NEW-U16-VOC-WITH
% teaching: NEW-U16-SUP-SWIFT
% teaching: NEW-U16-VOC-SWIFT
% teaching: NEW-U16-SUP-HRAEDNES
% teaching: NEW-U16-VOC-HRAEDNES
% teaching: NEW-U16-SUP-ONGEAN
% teaching: NEW-U16-VOC-ONGEAN
% teaching: NEW-U16-SUP-HRAEDLICE
% teaching: NEW-U16-VOC-HRAEDLICE
% teaching: NEW-U16-SUP-SNELNES
% teaching: NEW-U16-VOC-SNELNES
% teaching: NEW-U16-VOC-THAET-CONJUNCTION
% teaching: NEW-U16-SUP-HEONAN
% teaching: NEW-U16-VOC-HEONAN
% teaching: NEW-U16-SUP-CNIHT
% teaching: NEW-U16-VOC-CNIHT
% teaching: NEW-U16-VOC-THAES-THE
% teaching: NEW-U16-SUP-TOP
% teaching: NEW-U16-VOC-TOP
% teaching: NEW-U16-SUP-SWIFTNES
% teaching: NEW-U16-VOC-SWIFTNES
% teaching: NEW-U16-SUP-WESAN
% teaching: NEW-U16-VOC-WESAN
% teaching: NEW-U16-SUP-IELDU
% teaching: NEW-U16-VOC-IELDU
% teaching: NEW-U16-SUP-GEOGUTH
% teaching: NEW-U16-VOC-GEOGUTH
% teaching: NEW-U16-SUP-CYNESETL
% teaching: NEW-U16-VOC-CYNESETL
% teaching: NEW-U16-SUP-GECWEMLICE
% teaching: NEW-U16-VOC-GECWEMLICE
% teaching: NEW-U16-SUP-GEMAENE
% teaching: NEW-U16-VOC-GEMAENE
% teaching: NEW-U16-SUP-HAELU
% teaching: NEW-U16-VOC-HAELU
% teaching: NEW-U16-SUP-THENUNG
% teaching: NEW-U16-VOC-THENUNG
% teaching: NEW-U16-SUP-GESEON
% teaching: NEW-U16-VOC-GESEON
% teaching: NEW-U16-SUP-HORG
% teaching: NEW-U16-VOC-HORG
% teaching: NEW-U16-SUP-FORLIDEN
% teaching: NEW-U16-VOC-FORLIDEN
% teaching: NEW-U16-SUP-GEREORD
% teaching: NEW-U16-VOC-GEREORD
\reserve{28\baselineskip}
% mdparagraph:P538:begin
\textbf{Reading forms and meanings.} These forms and meanings remain available for this task.
% mdparagraph:P538:end

% mdtable:T90:begin
\par\begingroup
\fontsize{10.5}{13.2}\selectfont
\setlength{\LTleft}{0pt}\setlength{\LTright}{\fill}
\begin{longtable}{@{}>{\raggedright\arraybackslash}p{26mm}>{\raggedright\arraybackslash}p{44.5mm}>{\raggedright\arraybackslash}p{47.88mm}@{}}
\toprule
\textbf{Form} & \textbf{Meaning} & \textbf{Role} \\
\midrule
\endfirsthead
\toprule
\textbf{Form} & \textbf{Meaning} & \textbf{Role} \\
\midrule
\endhead
\midrule
\multicolumn{3}{r}{\footnotesize Continued on the following page.}\\
\endfoot
\bottomrule
\endlastfoot
ðisum & this; these & dative, plural, neuter \\
getǣht & shown; directed & past, participle \\
ōð-ðæt & until & invariable \\
becōm & come; arrive & third singular preterite indicative \\
inēode & went in & third person, preterite, indicative, singular \\
mid-ðȳ-ðe & when; while & invariable \\
ðōhte & thought & as used here \\
meahte & might; could & third singular preterite; indicative or subjunctive in form \\
fultumes & help; support & genitive, singular, masculine \\
ǣnne & one & accusative, singular, masculine \\
sē & that; those & nominative, singular, masculine \\
gesmierwed & anointed & past, participle \\
mid & with a cloth or cloak, accompanying girded or wrapped & as used here \\
scīetan & linen cloth & dative, singular, feminine \\
begyrd & girded; clothed around & past, participle \\
geongra & young & genitive, plural, masculine \\
plegan & play; game & accusative, singular, masculine \\
handa & hand; with right/\allowbreak{}left, side & dative, singular \\
tō & to, complement of belonging & invariable \\
bæðstęde & bathing place; here gymnasium & dative, singular, masculine \\
belimpende & belonging; relating & present, participle \\
Gehīere & hear; listen & second person, imperative, plural \\
ceasterwaran & townspeople; citizens & nominative, plural, masculine \\
ęlðēodige & foreign people; strangers & nominative, plural, masculine \\
frige & free & nominative, plural, masculine \\
ðēowe & unfree; enslaved people & nominative, plural, masculine \\
æðele & noble people & nominative, plural, masculine \\
unæðele & people not noble & nominative, plural, masculine \\
bæðstęde & bathing place; here gymnasium & nominative, singular, masculine \\
open & open & nominative, singular, masculine \\
ðā-ðā & when; factual past time & invariable \\
Apollonius & Apollonius & nominative, singular, masculine \\
gehīerde & heard, third singular preterite indicative of prefixed hear & singular \\
unscrȳdde & undressed; took off & third person, preterite, indicative, singular \\
healfan & half & dative, singular, masculine \\
sciccelse & cloak & dative, singular, masculine \\
on & on, in the relative clause referring to the cloak he wore & invariable \\
ðwēale & bath; bathing place & dative, singular, neuter \\
behēold & observed & third person, preterite, indicative, singular \\
ānra & one & genitive, plural, masculine, strong \\
on & at work; engaged in an activity, in exact source phrases & invariable \\
gelīcan & equal; peer & accusative, singular, masculine \\
findan & find & infinitive \\
flocce & group; company & dative, singular, masculine \\
fǣringa & suddenly & invariable \\
Arcestrates & Arcestrates, king of the region & nominative, singular, masculine \\
ðēode & region; country & genitive, singular, feminine \\
męnige & company; crowd & dative, singular, feminine \\
āgan & began & third person, preterite, indicative, singular \\
plegian & play & infinitive \\
wið & with; against fellow players, with dative in the exact source playing frame & as used here \\
gemęngde & joined; mingled & third person, preterite, indicative, singular \\
swā-swā & just as; like & invariable \\
iernende & run & present, participle, nominative, singular, masculine \\
gelæhte & caught & third person, preterite, indicative, singular \\
mid & with; in the manner of the accompanying quality or mental disposition & as used here \\
swiftre & swift & dative, singular, feminine \\
hrædnesse & speed; quickness & dative, singular, feminine \\
geslægene & struck & past participle, masculine accusative singular; exceptional ending (P169N2) \\
ongēan & back; again & invariable \\
gesęnde & sent back & third person, preterite, indicative, singular \\
plegiendan & play & present, participle, dative, singular, masculine, weak \\
āgēan & back; again & invariable \\
hrædlīce & quickly & invariable \\
slōh & struck & third person, preterite, indicative, singular \\
feallan & fall & infinitive \\
oncnēow & recognized; perceived & third person, preterite, indicative, singular \\
geongan & young & genitive, singular, masculine \\
snelnesse & skill; agility & accusative, singular, feminine \\
ðæt & that; the precise linkage at 169.6 remains uncertain (P169N6) & as used here \\
plegan & play; game & dative, singular, masculine \\
tō & to, the person addressed in speech & invariable \\
heonan & hence; from here & invariable \\
cniht & young man; youth & nominative, singular, masculine \\
ðæs-ðe & as, in the exact parenthesis ðæs-ðe mē ðyncð & as used here \\
gelīca & equal; peer & nominative, singular, masculine \\
hęrede & praised & third singular preterite; indicative reading here, with the same subjunctive spelling possible \\
genēalǣcte & approached & third person, preterite, indicative, singular \\
swang & swung; brandished & third person, preterite, indicative, singular \\
top & ball?, as Cook tentatively interprets it & accusative, singular, masculine \\
swiftnesse & swiftness & dative, singular, feminine \\
wæs & in the impersonal source expression it seemed to the king & third person, preterite, indicative, singular \\
geðūht & seemed & past, participle \\
swilce & as if, introducing the source comparison of apparent rejuvenation & invariable \\
of & from, the state of old age in the supplied rejuvenation comparison & dative \\
ielde & age; old age & dative, singular, feminine \\
tō & to, a resulting state of youth & invariable \\
geoguðe & youth & dative, singular, feminine \\
gewęnd & turned; changed back & past, participle \\
cynesętle & royal seat; throne & dative, singular, neuter \\
gecwēmlīce & pleasantly; agreeably & invariable \\
ðēnode & served; attended & third person, preterite, indicative, singular \\
ūtēode & went out & third person, preterite, indicative, singular \\
be & by, the body part by which a person is led & as used here \\
ðanan & thence; from there & as used here \\
gewęnde & turned; changed back & third person, preterite, indicative, singular \\
āgān & gone away & past, participle \\
wæs & auxiliary with an intransitive past participle & third person, preterite, indicative, singular \\
swęrie & swear & first person, present, indicative, singular \\
gemǣnan & common; shared & accusative, singular, feminine \\
hǣle & safety; rescue & accusative, singular, feminine \\
baðode & bathed & first person, preterite, indicative, singular \\
nāt & do not know & first person, present, indicative, singular \\
ðurh & through; by means of service, or by evidence in the exact source question Ðurh hwæt & as used here \\
geonges & young & genitive, singular, masculine \\
ðēnunge & service; attendance & accusative, singular, feminine \\
beseah & looked; turned his gaze & third person, preterite, indicative, singular \\
gewite & find out; learn & second person, imperative, singular \\
geonga & young & nominative, singular, masculine \\
gehīersumode & served; obeyed & third person, preterite, indicative, singular \\
æfter & after; in search of; about the person sought or inquired after & as used here \\
Apollonio & Apollonius & dative, singular, masculine \\
geseah & see that a stated circumstance holds & third person, preterite, indicative, singular \\
horgum & dirty; filthy & dative, singular, masculine \\
bewǣfed & wrapped; clothed & past, participle \\
węnde & turned; returned & third person, preterite, indicative, singular \\
āscodest & asked about & second person, preterite, indicative, singular \\
forliden & shipwrecked & nominative, singular, masculine \\
forswīge & keep silent about; conceal & third person, present, subjunctive, singular \\
gesweotolað & reveals; makes clear & third person, present, indicative, singular \\
sęge & say; tell & second person, imperative, singular \\
gereorde & meal; banquet & dative, singular, neuter \\
\end{longtable}
\par\endgroup

% mdtable:T90:end
% mdparagraph:P539:begin
\textbf{Further verbal help.} The following notes identify the action or construction in the exact source form.
% mdparagraph:P539:end

% mdparagraph:P540:begin
\textbf{168.02. getǣht — shown; directed.} Past participle in passive getǣht wæs; him is the person to whom the route was indicated.
% mdparagraph:P540:end

% mdparagraph:P541:begin
\textbf{168.03. inēode — went in.} The prefixed form means went in; its finite verb is part of the word.
% mdparagraph:P541:end

% mdparagraph:P542:begin
\textbf{168.03. meahte — might; could.} Inside the dependent question hwone ... biddan meahte; indicative/\allowbreak{}subjunctive surface coincides. The selected interpretation is possibility in the thought.
% mdparagraph:P542:end

% mdparagraph:P543:begin
\textbf{168.05. gesmierwed — anointed.} Exact source prefixed participle and retained w; distinct from selected unprefixed smierwan paradigm. Mid ęle marks the anointing material.
% mdparagraph:P543:end

% mdparagraph:P544:begin
\textbf{168.06. begyrd — girded; clothed around.} Passive/\allowbreak{}predicative participle with mid scīetan, a cloth.
% mdparagraph:P544:end

% mdparagraph:P545:begin
\textbf{168.07. belimpende — belonging; relating.} Tō + dative bæðstęde identifies what the games/\allowbreak{}equipment relate to; preserve P168N6.
% mdparagraph:P545:end

% mdparagraph:P546:begin
\textbf{168.07. cleopode — cried aloud.} Preterite third singular of cleopian in the cry-aloud sense. Micelre stefne expresses the manner or means, “with a loud voice”.
% mdparagraph:P546:end

% mdparagraph:P547:begin
\textbf{168.08. Gehīere — hear; listen.} Exact plural command with shortened -e before following gē, not singular by surface alone. The same form recurs later in the line.
% mdparagraph:P547:end

% mdparagraph:P548:begin
\textbf{168.11. unscrȳdde — undressed; took off.} Hine is the person undressed; ðǣm healfan sciccelse is the removed clothing, supplied case frame.
% mdparagraph:P548:end

% mdparagraph:P549:begin
\textbf{168.13. behēold — observed.} Accusative persons observed; no modern behold quotation formula implied.
% mdparagraph:P549:end

% mdparagraph:P550:begin
\textbf{168.17. inēode — went in.} On ðæt bæð gives the entry destination.
% mdparagraph:P550:end

% mdparagraph:P551:begin
\textbf{168.17. āgan — began.} Strong III preterite of āginnan, not āgan own/\allowbreak{}possess and not āgān depart.
% mdparagraph:P551:end

% mdparagraph:P552:begin
\textbf{168.17. plegian — play.} Infinitive after began; wið + dative companions and mid + dative ball are supported locally.
% mdparagraph:P552:end

% mdparagraph:P553:begin
\textbf{168.18. gemęngde — joined; mingled.} Reflexive hine; on ðæs cyninges plegan continues on 169.1. Plegan is an oblique noun governed by on; ðæs cyninges is its genitive possessor. It is not a finite verb.
% mdparagraph:P553:end

% mdparagraph:P554:begin
\textbf{169.01. iernende — running.} The supplied present participle accompanies Apollonius. The printed division iern-/\allowbreak{}ende is joined in the reading.
% mdparagraph:P554:end

% mdparagraph:P555:begin
\textbf{169.02. gelæhte — caught.} Exact source preterite has short æ as printed; preserve lexical spelling and do not silently import a guessed long form.
% mdparagraph:P555:end

% mdparagraph:P556:begin
\textbf{169.03. geslægene — struck.} The participial phrase refers to the accusative ðone ðoðor. Retain the unusual printed -e, to which Cook draws attention in P169N2.
% mdparagraph:P556:end

% mdparagraph:P557:begin
\textbf{169.03. gesęnde — sent back.} Tō + dative recipient; ongēan is the return direction.
% mdparagraph:P557:end

% mdparagraph:P558:begin
\textbf{169.03. plegiendan — playing.} Weak dative modifier of cyninge after ðǣm; the finite paradigm does not itself teach this participial cell.
% mdparagraph:P558:end

% mdparagraph:P559:begin
\textbf{169.04. slōh — struck.} The supplied h spelling is a valid singular variant of the previously selected slōg.
% mdparagraph:P559:end

% mdparagraph:P560:begin
\textbf{169.05. oncnēow — recognized; perceived.} Accusative snelnesse, followed by a result/\allowbreak{}content sequence whose pronoun referents remain genuinely underdetermined in P169N6.
% mdparagraph:P560:end

% mdparagraph:P561:begin
\textbf{169.09. hęrede — praised.} Hine is the person praised. The selected indicative interpretation fits the known event, but the third-singular preterite subjunctive has the same spelling; the ending alone does not distinguish them.
% mdparagraph:P561:end

% mdparagraph:P562:begin
\textbf{169.10. arn — ran.} Third singular preterite of iernan, as in the learned table.
% mdparagraph:P562:end

% mdparagraph:P563:begin
\textbf{169.10. genēalǣcte — approached.} The prefixed verb with tō + dative goal is supplied; do not rename it as the unprefixed nēalǣcan form.
% mdparagraph:P563:end

% mdparagraph:P564:begin
\textbf{169.11. swang — swung; brandished.} Its direct object is top, tentatively “ball”. The Latin parallel describes rubbing or anointing with ointment; retain the actual Old English action rather than replacing it by that parallel.
% mdparagraph:P564:end

% mdparagraph:P565:begin
\textbf{169.12. geðūht — seemed.} Supply ðǣm cyninge wæs geðūht swilce ... wǣre, it seemed to the king as though ...; dative experiencer and irrealis comparison are explicit.
% mdparagraph:P565:end

% mdparagraph:P566:begin
\textbf{169.13. gewęnd — turned; changed back.} In swilce hē of ielde tō geoguðe gewęnd wǣre: as though he had been turned from age to youth. Do not mechanically label every wesan+participle a motion perfect.
% mdparagraph:P566:end

% mdparagraph:P567:begin
\textbf{169.14. ðēnode — served; attended.} The dative him is the person served. Retain the source’s -ode spelling.
% mdparagraph:P567:end

% mdparagraph:P568:begin
\textbf{169.15. ūtēode — went out.} Prefixed motion verb with of + dative place left.
% mdparagraph:P568:end

% mdparagraph:P569:begin
\textbf{169.16. gewęnde — turned; went away.} Reflexive dative him; ðæs weges is the route genitive, matched by the composition 35.3. Pronoun references need discourse reasoning.
% mdparagraph:P569:end

% mdparagraph:P570:begin
\textbf{170.02. āgān — gone away.} Past participle in āgān wæs, “had gone away”. Its second vowel is long; the second vowel of āgan “began” at 168.17 is short.
% mdparagraph:P570:end

% mdparagraph:P571:begin
\textbf{170.02. swęrie — swear.} First-person quoted oath; ðurh + accusative oath standard; the narrated speaker now owns ic.
% mdparagraph:P571:end

% mdparagraph:P572:begin
\textbf{170.03. baðode — bathed.} Reflexive mē and adverbial comparison with dyde; this bath verb is supplied, substitute dōn is newly productive.
% mdparagraph:P572:end

% mdparagraph:P573:begin
\textbf{170.03. nāt — do not know.} The supplied contracted negative means “I do not know”. Only this local form of nytan is required.
% mdparagraph:P573:end

% mdparagraph:P574:begin
\textbf{170.04. beseah — looked; turned his gaze.} Reflexive hine and tō + dative person toward whom he turned; supply as a local phrase.
% mdparagraph:P574:end

% mdparagraph:P575:begin
\textbf{170.05. gewite — find out; learn.} The preterite-present knowing verb (§126), not strong I gewītan depart; the dependent question has hwæt ... sīe.
% mdparagraph:P575:end

% mdparagraph:P576:begin
\textbf{170.06. gehīersumode — served; obeyed.} Dative mē is the person served; hyphenation gehīer-/\allowbreak{}sumode is print line division only.
% mdparagraph:P576:end

% mdparagraph:P577:begin
\textbf{170.09. bewǣfed — wrapped; clothed.} Mid + dative garment in the passive/\allowbreak{}predicative description.
% mdparagraph:P577:end

% mdparagraph:P578:begin
\textbf{170.09. węnde — turned; returned.} Ongēan + tō destination; distinct from wēnan expected despite the similar modern gloss.
% mdparagraph:P578:end

% mdparagraph:P579:begin
\textbf{170.11. āscodest — asked about.} Æfter links to the relative antecedent; exact source spelling preserves -odest.
% mdparagraph:P579:end

% mdparagraph:P580:begin
\textbf{170.13. forswīge — keep silent about; conceal.} Ðēah concessive clause; hit is the thing withheld. Exact source ending is supplied rather than silently generalized to a weak-II table.
% mdparagraph:P580:end

% mdparagraph:P581:begin
\textbf{170.14. gesweotolað — reveals; makes clear.} Gegierela is the singular subject, hine the person whose condition is revealed.
% mdparagraph:P581:end

% mdparagraph:P582:begin
\textbf{170.15. sęge — say; tell.} Exact source shortened command spelling; the following mixed direct/\allowbreak{}indirect quotation is explicitly supported by P170N15, not presented as a regular report transformation.
% mdparagraph:P582:end

% mdparagraph:P583:begin
\textbf{168.01, 168.02: ðone weg ðe him getǣht wæs.} Relative ðe refers to the road; getǣht wæs is a past passive, with him the person shown the route.
% mdparagraph:P583:end

% mdparagraph:P584:begin
\textbf{168.02: ōð-ðæt hē becōm.} Until he arrived: an endpoint clause; retain the source macron.
% mdparagraph:P584:end

% mdparagraph:P585:begin
\textbf{168.03, 168.04: hwone hē biddan meahte līfes fultumes.} hwone is the person asked, while līfes fultumes names the requested support of life. meahte can be indicative or subjunctive in form; context is a question within his thought, not proof of a uniquely recognizable ending.
% mdparagraph:P585:end

% mdparagraph:P586:begin
\textbf{168.04, 168.05: geseah hē ǣnne nacodne cnapan ... iernan.} Accusative participant with a perception infinitive. Cook’s glossary allows nacod to mean wearing only a tunic here; do not require a claim of total nudity.
% mdparagraph:P586:end

% mdparagraph:P587:begin
\textbf{168.06, 168.07: geongra manna plegan ... tō ðǣm bæðstęde belimpende.} plegan literally denotes play/\allowbreak{}games; carrying equipment for the young men’s games is a contextual explanation, not a new unrestricted dictionary sense. belimpende modifies plegan, as P168N6 says; it is not another finite action of the boy.
% mdparagraph:P587:end

% mdparagraph:P588:begin
\textbf{168.08, 168.09: Gehīere, gē ceasterwaran ....} The call is addressed to the plural audience. Supply the whole source exhortation; do not replace its Gehīere with a silently regularized imperative.
% mdparagraph:P588:end

% mdparagraph:P589:begin
\textbf{168.11, 168.12: hine unscrȳdde ðǣm healfan sciccelse ðe hē on hæfde.} Reflexive undressing and the garment in the dative; the relative clause says he had the cloak on. Preserve the source construction as a whole.
% mdparagraph:P589:end

% mdparagraph:P590:begin
\textbf{168.13: hiera ānra gehwilcne on hiera weorce.} Each single one of them at their work: first hiera participates in the partitive expression; second is possessive. Keep the printed short spelling.
% mdparagraph:P590:end

% mdparagraph:P591:begin
\textbf{169.02, 169.03: mid swiftre hrædnesse geslægene.} P169N2 remarks on the participle’s awkward form. Supply the complete phrase as describing the ball struck with swift speed; retain geslægene rather than silently repairing agreement.
% mdparagraph:P591:end

% mdparagraph:P592:begin
\textbf{169.05, 169.06, 169.07: oncnēow ... ðæt hē wiste ðæt hē næfde his gelīcan.} The nested reference is unclear, as P169N6 states. næfde is negative: the text does not say he had an equal. The Latin comparison concerns the king’s previously unmatched skill; use it as help while retaining the OE uncertainty, not as an invisible replacement text.
% mdparagraph:P592:end

% mdparagraph:P593:begin
\textbf{169.08: ðæs-ðe mē ðyncð.} As it seems to me: exact parenthetic impersonal construction, dative experiencer mē.
% mdparagraph:P593:end

% mdparagraph:P594:begin
\textbf{169.11, 169.12, 169.13: swang ... ðone top ... swilce ... gewęnd wǣre.} Retain the Old English action with top, tentatively a ball. Cook’s Latin parallel describes rubbing or anointing with ointment and its rejuvenating effect; it names no shoulder or neck. Wǣre belongs to the imagined transformation comparison.
% mdparagraph:P594:end

% mdparagraph:P595:begin
\textbf{169.12: ðǣm cyninge wæs geðūht.} It seemed to the king: whole impersonal be-plus-participle construction supplied; do not read it as an ordinary king-as-subject clause.
% mdparagraph:P595:end

% mdparagraph:P596:begin
\textbf{169.14, 169.15, 169.16: him ... hē hine lǣdde be ðǣre handa ... him ... gewęnde.} Use P169N11 and P169N14–P169N16 for the pronoun chain. The later him is reflexive; the by-the-hand phrase specifies how one person leads another.
% mdparagraph:P596:end

% mdparagraph:P597:begin
\textbf{169.16, 169.17: ðæs weges ðe hē ǣr cōm.} The genitive names the route followed, as in the pattern taught above; the relative ðe identifies that route. Retain the past-time reference of the source wording.
% mdparagraph:P597:end

% mdparagraph:P598:begin
\textbf{170.01, 170.02: siððan Apollonius āgān wæs.} After Apollonius had gone away: intransitive completed departure with be, not a passive in which someone sends him away.
% mdparagraph:P598:end

% mdparagraph:P599:begin
\textbf{170.02: ðurh ðā gemǣnan hǣle.} An oath by the common safety or welfare. Hǣle is feminine accusative singular, here in the safety/\allowbreak{}rescue sense. No possessive “our” is printed.
% mdparagraph:P599:end

% mdparagraph:P600:begin
\textbf{170.04, 170.05, 170.06: ðurh hwilces geonges mannes ðēnunge ... hwæt se geonga mann sīe.} The service is the means of the benefit; hwilces ... mannes is its genitive dependent. hwæt asks what sort/\allowbreak{}status of man he is. sīe is the dependent interrogative be form.
% mdparagraph:P600:end

% mdparagraph:P601:begin
\textbf{170.08, 170.09, 170.10, 170.11: ēode æfter Apollonio ... ðe ðū æfter āscodest.} Going after a person, then asking about him: the relative marker is the complement of æfter within the clause; do not force temporal after for either use.
% mdparagraph:P601:end

% mdparagraph:P602:begin
\textbf{170.13: Ðēah hē hit self forswīge.} Concessive clause: even if he himself keeps it unspoken. Supply the source forswīge with its present subjunctive interpretation; do not demand a full new weak-verb paradigm.
% mdparagraph:P602:end

% mdparagraph:P603:begin
\textbf{170.14, 170.15, 170.16: sęge him ðæt ‘se cyning bitt ðē ðæt ðū cume ...’.} P170N15 explicitly diagnoses mixed direct/\allowbreak{}reported construction. First him is the messenger’s addressee; within the quoted message ðē is the person invited (accusative), and ðū/\allowbreak{}cume address him directly. Retain wording and speech levels; no silent repair.
% mdparagraph:P603:end

% mdparagraph:P604:begin
\textbf{The Latin parallels in English.} Original P169N6 compares the difficult knowledge/\allowbreak{}equal clause with a Latin account in which the king has long known his own skill to be without an equal. This parallel helps explain Cook’s concern but does not settle every Old English pronoun or cancel \textbf{næfde}. Original P169N10 quotes a Latin account of rubbing or anointing with an ointment used by wrestlers; it does not name a shoulder or neck. Keep the Old English \textbf{swang ðone top} and the note’s tentative interpretation distinct from that Latin action. No Latin translation is required.
% mdparagraph:P604:end

% mdparagraph:P605:begin
\textbf{Quotation punctuation.} The edition begins a quotation at 170.2 but does not print its expected closing mark before the next narrated action. The final invitation also mixes quoted and reported wording. The text below retains those signs and constructions; they are not silently regularized.
% mdparagraph:P605:end

\reserve{7\baselineskip}
% mdparagraph:P606:begin
\subsection*{Source text}\BookHeading{nav-ch16-learner-0032}{subsubsection}{Source text}{0}
% mdparagraph:P606:end

% passage: FB-XII-APOLLONIUS-GYMNASIUM
% mdparagraph:P607:begin
\readingtext{Æfter ðisum wordum hē ēode on ðone weg ðe him getǣht\textsuperscript{P168N1} wæs, ōð-ðæt hē becōm tō ðǣre ceastre geate, and ðǣr inēode. Mid-ðȳ-ðe hē ðōhte hwone hē biddan meahte līfes fultumes,\textsuperscript{P168N2} ðā geseah hē ǣnne nacodne cnapan geond ðā strǣte iernan. Sē wæs mid\textsuperscript{P168N3} ęle gesmierwed, and mid scīetan begyrd, and bær geongra\textsuperscript{P168N4} manna\textsuperscript{P168N4} plegan\textsuperscript{P168N4} on handa, tō ðǣm bæðstęde\textsuperscript{P168N5} belimpende.\textsuperscript{P168N6} And cleopode\textsuperscript{P168N7} micelre stefne,\textsuperscript{P168N8} and cwæð: “Gehīere,\textsuperscript{P168N9} gē ceasterwaran\textsuperscript{P168N10}! Gehīere, gē ęlðēodige,\textsuperscript{P168N11} frige and ðēowe, æðele and unæðele! Se bæðstęde is open.”}
% mdparagraph:P607:end

% mdparagraph:P608:begin
\readingtext{Ðā-ðā Apollonius ðæt gehīerde, hē hine unscrȳdde ðǣm\textsuperscript{P168N12} healfan sciccelse ðe hē on hæfde, and ēode intō ðǣm ðwēale.\textsuperscript{P168N13} And mid-ðȳ-ðe hē behēold hiera ānra\textsuperscript{P168N14} gehwilcne on hiera weorce, hē sōhte his gelīcan,\textsuperscript{P168N15} ac hē ne meahte hine ðǣr findan on ðǣm flocce. Ðā fǣringa cōm Arcestrates, ealre ðǣre ðēode\textsuperscript{P168N16} cyning,\textsuperscript{P168N17} mid micelre męnige his manna,\textsuperscript{P168N18} and inēode on ðæt bæð. Ðā āgan se cyning plegian wið\textsuperscript{P168N19} his gefērum mid ðoðore.\textsuperscript{P168N20} And Apollonius hine\textsuperscript{P168N21} gemęngde,\textsuperscript{P168N21} swā-swā God wolde, on ðæs cyninges plegan, and, iernende, ðone ðoðor gelæhte,\textsuperscript{P169N1} and, mid swiftre hrædnesse geslægene,\textsuperscript{P169N2} ongēan gesęnde tō ðǣm plegiendan cyninge. Ęft hē āgēan āsęnde; hē hrædlīce slōh, swā hē hine\textsuperscript{P169N3} nǣfre feallan ne lēt. Se cyning ðā oncnēow ðæs geongan snelnesse,\textsuperscript{P169N4} ðæt hē wiste\textsuperscript{P169N5} ðæt hē næfde his gelīcan on ðǣm plegan.\textsuperscript{P169N6} Ðā cwæð hē tō his gefērum: “Gāð ēow heonan; ðes cniht, ðæs-ðe mē ðyncð,\textsuperscript{P169N7} is mīn gelīca.”}
% mdparagraph:P608:end

% mdparagraph:P609:begin
\readingtext{Ðā-ðā Apollonius gehīerde ðæt se cyning hine hęrede, hē arn hrædlīce, and genēalǣcte tō ðǣm cyninge, and mid gelǣredre\textsuperscript{P169N8} handa\textsuperscript{P169N9} hē swang\textsuperscript{P169N10} ðone top mid swā micelre swiftnesse ðæt ðǣm cyninge wæs geðūht swilce hē of ielde tō geoguðe gewęnd wǣre. And, æfter ðǣm, on his cynesętle hē him\textsuperscript{P169N11} gecwēmlīce\textsuperscript{P169N12} ðēnode\textsuperscript{P169N13}; and, ðā-ðā hē ūtēode of ðǣm bæðe, hē\textsuperscript{P169N14} hine\textsuperscript{P169N15} lǣdde be ðǣre handa, and him\textsuperscript{P169N16} ðā siððan ðanan gewęnde, ðæs weges\textsuperscript{P169N17} ðe hē ǣr cōm.}
% mdparagraph:P609:end

% mdparagraph:P610:begin
\readingtext{Ðā cwæð se cyning tō his mannum,\textsuperscript{P170N1} siððan Apollonius āgān\textsuperscript{P170N2} wæs: “Ic swęrie ðurh ðā gemǣnan hǣle\textsuperscript{P170N3} ðæt ic mē nǣfre bęt ne baðode ðonne ic dyde\textsuperscript{P170N4} tō-dæg, nāt ic ðurh\textsuperscript{P170N5} hwilces geonges mannes ðēnunge.\textsuperscript{P170N6} Ðā beseah hē hine tō ānum his manna, and cwæð: “Gā, and gewite hwæt\textsuperscript{P170N7} se geonga mann sīe, ðe mē tō-dæg swā wel gehīersumode.”}
% mdparagraph:P610:end

% mdparagraph:P611:begin
\readingtext{Se mann ðā ēode æfter Apollonio. Mid-ðȳ-ðe hē geseah ðæt hē\textsuperscript{P170N8} wæs mid horgum\textsuperscript{P170N9} sciccelse bewǣfed, ðā węnde hē ongēan tō ðǣm cyninge, and cwæð: “Se geonga mann ðe\textsuperscript{P170N10} ðū æfter āscodest is forliden\textsuperscript{P170N11} mann.\textsuperscript{P170N11}” Ðā cwæð se cyning: “Ðurh\textsuperscript{P170N12} hwæt\textsuperscript{P170N12} wāst\textsuperscript{P170N13} ðū ðæt?” Se mann him andswarode, and cwæð: “Ðēah hē hit self forswīge,\textsuperscript{P170N14} his gegierela hine gesweotolað.” Ðā cwæð se cyning: “Gā hrædlīce, and sęge him ðæt\textsuperscript{P170N15} ‘se cyning bitt ðē ðæt ðū cume\textsuperscript{P170N16} tō his gereorde.’”}
% mdparagraph:P611:end

\par
\reserve{7\baselineskip}
% mdparagraph:P612:begin
\subsection*{Original notes}\BookHeading{nav-ch16-learner-0033}{subsubsection}{Original notes}{0}
% mdparagraph:P612:end

% mdparagraph:P613:begin
\textsuperscript{P168N1}: See 187.
% mdparagraph:P613:end

% mdparagraph:P614:begin
\textsuperscript{P168N2}: MS. fultum. See 159. b.
% mdparagraph:P614:end

% mdparagraph:P615:begin
\textsuperscript{P168N3}: See 174.
% mdparagraph:P615:end

% mdparagraph:P616:begin
\textsuperscript{P168N4}: Lat. lusus juvenales.
% mdparagraph:P616:end

% mdparagraph:P617:begin
\textsuperscript{P168N5}: Lat. gymnasium.
% mdparagraph:P617:end

% mdparagraph:P618:begin
\textsuperscript{P168N6}: Modifies plegan.
% mdparagraph:P618:end

% mdparagraph:P619:begin
\textsuperscript{P168N7}: See 20.
% mdparagraph:P619:end

% mdparagraph:P620:begin
\textsuperscript{P168N8}: See 160. 1. This word is the Chaucerian steven.
% mdparagraph:P620:end

% mdparagraph:P621:begin
\textsuperscript{P168N9}: See 95, note.
% mdparagraph:P621:end

% mdparagraph:P622:begin
\textsuperscript{P168N10}: This is a compound word, formed of a Latin and an English element. Which is Latin, and what is its form in that language?
% mdparagraph:P622:end

% mdparagraph:P623:begin
\textsuperscript{P168N11}: Lat. peregrini. See 152.
% mdparagraph:P623:end

% mdparagraph:P624:begin
\textsuperscript{P168N12}: See 162.
% mdparagraph:P624:end

% mdparagraph:P625:begin
\textsuperscript{P168N13}: Lat. lavacrum.
% mdparagraph:P625:end

% mdparagraph:P626:begin
\textsuperscript{P168N14}: See 154. b.
% mdparagraph:P626:end

% mdparagraph:P627:begin
\textsuperscript{P168N15}: Lat. parem, Eng. peer.
% mdparagraph:P627:end

% mdparagraph:P628:begin
\textsuperscript{P168N16}: Lat. regionis.
% mdparagraph:P628:end

% mdparagraph:P629:begin
\textsuperscript{P168N17}: See 151.
% mdparagraph:P629:end

% mdparagraph:P630:begin
\textsuperscript{P168N18}: Lat. famularum. See 154. a.
% mdparagraph:P630:end

% mdparagraph:P631:begin
\textsuperscript{P168N19}: Why not mid?
% mdparagraph:P631:end

% mdparagraph:P632:begin
\textsuperscript{P168N20}: This curious word is very rare in Old English.
% mdparagraph:P632:end

% mdparagraph:P633:begin
\textsuperscript{P168N21}: Lat. miscuit se. See 184. b.
% mdparagraph:P633:end

% mdparagraph:P634:begin
\textsuperscript{P169N1}: See 114.
% mdparagraph:P634:end

% mdparagraph:P635:begin
\textsuperscript{P169N2}: Lat. subtili velocitate percussam. The OE. participle is a little awkward.
% mdparagraph:P635:end

% mdparagraph:P636:begin
\textsuperscript{P169N3}: The ball.
% mdparagraph:P636:end

% mdparagraph:P637:begin
\textsuperscript{P169N4}: Lat. velocitatem.
% mdparagraph:P637:end

% mdparagraph:P638:begin
\textsuperscript{P169N5}: See 126. What is the latest English quotation that you can find for this word?
% mdparagraph:P638:end

% mdparagraph:P639:begin
\textsuperscript{P169N6}: This clause is not very clear. The Latin has: et quia sciebat se (i.e. Archistrates) in pilæ lusu neminem parem habere, ad suos ait, famuli, recedite; hic enim juvenis, etc.
% mdparagraph:P639:end

% mdparagraph:P640:begin
\textsuperscript{P169N7}: Lat. ut suspicor. See 157. 1 and 164. l. What Mod. Eng. word comes from mē ðyncð?
% mdparagraph:P640:end

% mdparagraph:P641:begin
\textsuperscript{P169N8}: Lat. docta.
% mdparagraph:P641:end

% mdparagraph:P642:begin
\textsuperscript{P169N9}: See 51. 3.
% mdparagraph:P642:end

% mdparagraph:P643:begin
\textsuperscript{P169N10}: Here the English departs from the Latin: ceroma fricavit eum tanta subtilitate, ut de sene juvenem redderet. Top would seem to signify the same as ðoðor.
% mdparagraph:P643:end

% mdparagraph:P644:begin
\textsuperscript{P169N11}: See 164. e.
% mdparagraph:P644:end

% mdparagraph:P645:begin
\textsuperscript{P169N12}: Lat. gratissime.
% mdparagraph:P645:end

% mdparagraph:P646:begin
\textsuperscript{P169N13}: See 28.
% mdparagraph:P646:end

% mdparagraph:P647:begin
\textsuperscript{P169N14}: Apollonius.
% mdparagraph:P647:end

% mdparagraph:P648:begin
\textsuperscript{P169N15}: Archistrates.
% mdparagraph:P648:end

% mdparagraph:P649:begin
\textsuperscript{P169N16}: See 184. a.
% mdparagraph:P649:end

% mdparagraph:P650:begin
\textsuperscript{P169N17}: See 157.
% mdparagraph:P650:end

% mdparagraph:P651:begin
\textsuperscript{P170N1}: Lat. amicos.
% mdparagraph:P651:end

% mdparagraph:P652:begin
\textsuperscript{P170N2}: How is the sense of Mod. Eng. ago related to that of this word?
% mdparagraph:P652:end

% mdparagraph:P653:begin
\textsuperscript{P170N3}: This phrase shows Christian influence.
% mdparagraph:P653:end

% mdparagraph:P654:begin
\textsuperscript{P170N4}: Note this use of dōn to replace a verb of specific meaning.
% mdparagraph:P654:end

% mdparagraph:P655:begin
\textsuperscript{P170N5}: Governs ðēnunge.
% mdparagraph:P655:end

% mdparagraph:P656:begin
\textsuperscript{P170N6}: See 28.
% mdparagraph:P656:end

% mdparagraph:P657:begin
\textsuperscript{P170N7}: How does this, as here used, differ in meaning from hwā?
% mdparagraph:P657:end

% mdparagraph:P658:begin
\textsuperscript{P170N8}: Apollonius.
% mdparagraph:P658:end

% mdparagraph:P659:begin
\textsuperscript{P170N9}: Lat. sordido.
% mdparagraph:P659:end

% mdparagraph:P660:begin
\textsuperscript{P170N10}: Governed by æfter. See 87. c and 201. 1.
% mdparagraph:P660:end

% mdparagraph:P661:begin
\textsuperscript{P170N11}: Lat. naufragus.
% mdparagraph:P661:end

% mdparagraph:P662:begin
\textsuperscript{P170N12}: Lat. unde.
% mdparagraph:P662:end

% mdparagraph:P663:begin
\textsuperscript{P170N13}: See 126.
% mdparagraph:P663:end

% mdparagraph:P664:begin
\textsuperscript{P170N14}: See 196. e.
% mdparagraph:P664:end

% mdparagraph:P665:begin
\textsuperscript{P170N15}: Confusion of two constructions, the direct and the indirect.
% mdparagraph:P665:end

% mdparagraph:P666:begin
\textsuperscript{P170N16}: Lat. ut venias. Translate by the infinitive, as often in such cases.
% mdparagraph:P666:end

\reserve{7\baselineskip}
% mdparagraph:P667:begin
\subsection*{Reading tasks}\BookHeading{nav-ch16-learner-0034}{subsubsection}{Reading tasks}{0}
% mdparagraph:P667:end

% mdparagraph:P668:begin
Use the whole text, the original notes and the supplied modern help for all six tasks. A supported interpretation should preserve the actual negative and clause structure while acknowledging a genuine uncertainty. The line excerpts below are locating aids; interpret the complete clause in the continuous reading.
% mdparagraph:P668:end

\reserve{5\baselineskip}
\par\begin{minipage}{\linewidth}\raggedright
% mdparagraph:P669:begin
\textbf{R01}
% mdparagraph:P669:end

% item: NEW-U16-R01
% mdparagraph:P670:begin
Translate the whole gymnasium scene, from Æfter ðisum wordum through tō his gereorde. Keep the local help and original notes available. Preserve the changes of speaker, the complete action sequence and the places where the edition acknowledges uncertainty.
% mdparagraph:P670:end

% output: NEW-U16-R01-O001
% mdparagraph:P671:begin
(1) Whole translation: \_\_\_\_\_\_\_\_\_\_\_\_\_\_\_\_\_\_\_\_
% mdparagraph:P671:end

\end{minipage}\par\addvspace{4pt}

\par\begin{minipage}{\linewidth}\raggedright
% mdparagraph:P672:begin
\textbf{R02}
% mdparagraph:P672:end

% item: NEW-U16-R02
% mdparagraph:P673:begin
Explain the sequence and grounds of the encounter in five brief answers, citing the relevant words or source lines: what brings Apollonius into the bath; what attracts the king's notice during the game; what the king says about the young man; what the messenger infers and from what evidence; and what the king finally orders.
% mdparagraph:P673:end

% output: NEW-U16-R02-O001
% mdparagraph:P674:begin
(1) Entry into the bath: \_\_\_\_\_\_\_\_\_\_\_\_\_\_\_\_\_\_\_\_
% mdparagraph:P674:end

% output: NEW-U16-R02-O002
% mdparagraph:P675:begin
(2) The game: \_\_\_\_\_\_\_\_\_\_\_\_\_\_\_\_\_\_\_\_
% mdparagraph:P675:end

% output: NEW-U16-R02-O003
% mdparagraph:P676:begin
(3) The king's statement: \_\_\_\_\_\_\_\_\_\_\_\_\_\_\_\_\_\_\_\_
% mdparagraph:P676:end

% output: NEW-U16-R02-O004
% mdparagraph:P677:begin
(4) The messenger's conclusion: \_\_\_\_\_\_\_\_\_\_\_\_\_\_\_\_\_\_\_\_
% mdparagraph:P677:end

% output: NEW-U16-R02-O005
% mdparagraph:P678:begin
(5) Final order: \_\_\_\_\_\_\_\_\_\_\_\_\_\_\_\_\_\_\_\_
% mdparagraph:P678:end

\end{minipage}\par\addvspace{4pt}

\par\begin{minipage}{\linewidth}\raggedright
% mdparagraph:P679:begin
\textbf{R03}
% mdparagraph:P679:end

% item: NEW-U16-R03
% mdparagraph:P680:begin
Identify the referent and grammatical role of each selected pronoun. Use the local context and original notes; distinguish a contextual interpretation from what its ending alone proves.
% mdparagraph:P680:end

% mdparagraph:P681:begin
Selected passages: 169.4 hine after slōh; 169.14 him after hē; 169.15 hē immediately before hine lǣdde; 169.15 hine before lǣdde; 169.16 him before ðā siððan; 170.9 hē after ðæt.
% mdparagraph:P681:end

% mdparagraph:P682:begin
Source wording for the cited locations:
% mdparagraph:P682:end

% mdparagraph:P683:begin
169.04: Ęft hē āgēan āsęnde; hē hrædlīce slōh, swā hē hine
% mdparagraph:P683:end

% mdparagraph:P684:begin
169.14: on his cynesętle hē him gecwēmlīce ðēnode; and,
% mdparagraph:P684:end

% mdparagraph:P685:begin
169.15: ðā-ðā hē ūtēode of ðǣm bæðe, hē hine lǣdde be ðǣre
% mdparagraph:P685:end

% mdparagraph:P686:begin
169.16: handa, and him ðā siððan ðanan gewęnde, ðæs weges
% mdparagraph:P686:end

% mdparagraph:P687:begin
170.09: ðæt hē wæs mid horgum sciccelse bewǣfed, ðā węnde
% mdparagraph:P687:end

% output: NEW-U16-R03-O001
% mdparagraph:P688:begin
(1) 169.4 hine: \_\_\_\_\_\_\_\_\_\_\_\_\_\_\_\_\_\_\_\_
% mdparagraph:P688:end

% output: NEW-U16-R03-O002
% mdparagraph:P689:begin
(2) 169.14 him: \_\_\_\_\_\_\_\_\_\_\_\_\_\_\_\_\_\_\_\_
% mdparagraph:P689:end

% output: NEW-U16-R03-O003
% mdparagraph:P690:begin
(3) 169.15 hē before hine: \_\_\_\_\_\_\_\_\_\_\_\_\_\_\_\_\_\_\_\_
% mdparagraph:P690:end

% output: NEW-U16-R03-O004
% mdparagraph:P691:begin
(4) 169.15 hine: \_\_\_\_\_\_\_\_\_\_\_\_\_\_\_\_\_\_\_\_
% mdparagraph:P691:end

% output: NEW-U16-R03-O005
% mdparagraph:P692:begin
(5) 169.16 him: \_\_\_\_\_\_\_\_\_\_\_\_\_\_\_\_\_\_\_\_
% mdparagraph:P692:end

% output: NEW-U16-R03-O006
% mdparagraph:P693:begin
(6) 170.9 hē: \_\_\_\_\_\_\_\_\_\_\_\_\_\_\_\_\_\_\_\_
% mdparagraph:P693:end

\end{minipage}\par\addvspace{4pt}

\par\begin{minipage}{\linewidth}\raggedright
% mdparagraph:P694:begin
\textbf{R04}
% mdparagraph:P694:end

% item: NEW-U16-R04
% mdparagraph:P695:begin
For each passage give its relation to the surrounding clause, identify its finite verb or finite verb group, and state the form's mood with the help provided. Treat the source-specific constructions within their stated limits.
% mdparagraph:P695:end

% mdparagraph:P696:begin
Selected passages: 168.1–2 ðe him getǣht wæs; 168.2 ōð-ðæt hē becōm; 168.3 hwone hē biddan meahte; 169.9 ðæt se cyning hine hęrede; 170.6 hwæt se geonga mann sīe; 170.13 Ðēah hē hit self forswīge.
% mdparagraph:P696:end

% output: NEW-U16-R04-O001
% mdparagraph:P697:begin
(1) 168.1–2 relation: \_\_\_\_\_\_\_\_\_\_\_\_\_\_\_\_\_\_\_\_
% mdparagraph:P697:end

% output: NEW-U16-R04-O002
% mdparagraph:P698:begin
(2) 168.1–2 finite group: \_\_\_\_\_\_\_\_\_\_\_\_\_\_\_\_\_\_\_\_
% mdparagraph:P698:end

% output: NEW-U16-R04-O003
% mdparagraph:P699:begin
(3) 168.2 relation: \_\_\_\_\_\_\_\_\_\_\_\_\_\_\_\_\_\_\_\_
% mdparagraph:P699:end

% output: NEW-U16-R04-O004
% mdparagraph:P700:begin
(4) 168.2 finite form: \_\_\_\_\_\_\_\_\_\_\_\_\_\_\_\_\_\_\_\_
% mdparagraph:P700:end

% output: NEW-U16-R04-O005
% mdparagraph:P701:begin
(5) 168.3 relation: \_\_\_\_\_\_\_\_\_\_\_\_\_\_\_\_\_\_\_\_
% mdparagraph:P701:end

% output: NEW-U16-R04-O006
% mdparagraph:P702:begin
(6) 168.3 finite group: \_\_\_\_\_\_\_\_\_\_\_\_\_\_\_\_\_\_\_\_
% mdparagraph:P702:end

% output: NEW-U16-R04-O007
% mdparagraph:P703:begin
(7) 169.9 relation: \_\_\_\_\_\_\_\_\_\_\_\_\_\_\_\_\_\_\_\_
% mdparagraph:P703:end

% output: NEW-U16-R04-O008
% mdparagraph:P704:begin
(8) 169.9 finite form: \_\_\_\_\_\_\_\_\_\_\_\_\_\_\_\_\_\_\_\_
% mdparagraph:P704:end

% output: NEW-U16-R04-O009
% mdparagraph:P705:begin
(9) 170.6 relation: \_\_\_\_\_\_\_\_\_\_\_\_\_\_\_\_\_\_\_\_
% mdparagraph:P705:end

% output: NEW-U16-R04-O010
% mdparagraph:P706:begin
(10) 170.6 finite form: \_\_\_\_\_\_\_\_\_\_\_\_\_\_\_\_\_\_\_\_
% mdparagraph:P706:end

% output: NEW-U16-R04-O011
% mdparagraph:P707:begin
(11) 170.13 relation: \_\_\_\_\_\_\_\_\_\_\_\_\_\_\_\_\_\_\_\_
% mdparagraph:P707:end

% output: NEW-U16-R04-O012
% mdparagraph:P708:begin
(12) 170.13 finite form: \_\_\_\_\_\_\_\_\_\_\_\_\_\_\_\_\_\_\_\_
% mdparagraph:P708:end

\end{minipage}\par\addvspace{4pt}

\par\begin{minipage}{\linewidth}\raggedright
% mdparagraph:P709:begin
\textbf{R05}
% mdparagraph:P709:end

% item: NEW-U16-R05
% mdparagraph:P710:begin
Analyse the final invitation at 170.14–16. Identify the outer speaker, his immediate addressee, the person denoted by him, the person denoted by se cyning, and the person denoted by ðē and ðū. Then explain why the original note describes a mixture of direct and indirect construction. Preserve the actual wording while explaining it.
% mdparagraph:P710:end

% mdparagraph:P711:begin
Selected passages: 170.14–16.
% mdparagraph:P711:end

% mdparagraph:P712:begin
Source wording for the cited locations:
% mdparagraph:P712:end

% mdparagraph:P713:begin
170.14: gegierela hine gesweotolað.” Ðā cwæð se cyning: “Gā
% mdparagraph:P713:end

% mdparagraph:P714:begin
170.15: hrædlīce, and sęge him ðæt ‘se cyning bitt ðē ðæt ðū
% mdparagraph:P714:end

% mdparagraph:P715:begin
170.16: cume tō his gereorde.’”
% mdparagraph:P715:end

% output: NEW-U16-R05-O001
% mdparagraph:P716:begin
(1) Outer speaker: \_\_\_\_\_\_\_\_\_\_\_\_\_\_\_\_\_\_\_\_
% mdparagraph:P716:end

% output: NEW-U16-R05-O002
% mdparagraph:P717:begin
(2) Immediate addressee: \_\_\_\_\_\_\_\_\_\_\_\_\_\_\_\_\_\_\_\_
% mdparagraph:P717:end

% output: NEW-U16-R05-O003
% mdparagraph:P718:begin
(3) Him: \_\_\_\_\_\_\_\_\_\_\_\_\_\_\_\_\_\_\_\_
% mdparagraph:P718:end

% output: NEW-U16-R05-O004
% mdparagraph:P719:begin
(4) Se cyning: \_\_\_\_\_\_\_\_\_\_\_\_\_\_\_\_\_\_\_\_
% mdparagraph:P719:end

% output: NEW-U16-R05-O005
% mdparagraph:P720:begin
(5) Ðē and ðū: \_\_\_\_\_\_\_\_\_\_\_\_\_\_\_\_\_\_\_\_
% mdparagraph:P720:end

% output: NEW-U16-R05-O006
% mdparagraph:P721:begin
(6) Construction: \_\_\_\_\_\_\_\_\_\_\_\_\_\_\_\_\_\_\_\_
% mdparagraph:P721:end

\end{minipage}\par\addvspace{4pt}

\par\begin{minipage}{\linewidth}\raggedright
% mdparagraph:P722:begin
\textbf{R06}
% mdparagraph:P722:end

% item: NEW-U16-R06
% mdparagraph:P723:begin
Consult printed page 169, original notes 6 and 10 (not text lines 6 and 10), together with the modern help explaining their Latin parallels in English. For each location, give a defensible translation of the actual Old English and explain the limit the note places on a confident interpretation. No knowledge or translation of Latin is required; do not replace the Old English with the Latin account.
% mdparagraph:P723:end

% mdparagraph:P724:begin
Selected passages: 169.5–7; 169.11–13.
% mdparagraph:P724:end

% mdparagraph:P725:begin
Source wording for the cited locations:
% mdparagraph:P725:end

% mdparagraph:P726:begin
169.05: nǣfre feallan ne lēt. Se cyning ðā oncnēow ðæs geongan
% mdparagraph:P726:end

% mdparagraph:P727:begin
169.06: snelnesse, ðæt hē wiste ðæt hē næfde his gelīcan on
% mdparagraph:P727:end

% mdparagraph:P728:begin
169.07: ðǣm plegan. Ðā cwæð hē tō his gefērum: “Gāð ēow
% mdparagraph:P728:end

% mdparagraph:P729:begin
169.11: mid gelǣredre handa hē swang ðone top mid swā
% mdparagraph:P729:end

% mdparagraph:P730:begin
169.12: micelre swiftnesse ðæt ðǣm cyninge wæs geðūht swilce
% mdparagraph:P730:end

% mdparagraph:P731:begin
169.13: hē of ielde tō geoguðe gewęnd wǣre. And, æfter ðǣm,
% mdparagraph:P731:end

% output: NEW-U16-R06-O001
% mdparagraph:P732:begin
(1) 169.5–7 translation: \_\_\_\_\_\_\_\_\_\_\_\_\_\_\_\_\_\_\_\_
% mdparagraph:P732:end

% output: NEW-U16-R06-O002
% mdparagraph:P733:begin
(2) 169.5–7 limit: \_\_\_\_\_\_\_\_\_\_\_\_\_\_\_\_\_\_\_\_
% mdparagraph:P733:end

% output: NEW-U16-R06-O003
% mdparagraph:P734:begin
(3) 169.11–13 translation: \_\_\_\_\_\_\_\_\_\_\_\_\_\_\_\_\_\_\_\_
% mdparagraph:P734:end

% output: NEW-U16-R06-O004
% mdparagraph:P735:begin
(4) 169.11–13 limit: \_\_\_\_\_\_\_\_\_\_\_\_\_\_\_\_\_\_\_\_
% mdparagraph:P735:end

\end{minipage}\par\addvspace{4pt}

\reserve{8\baselineskip}
% mdparagraph:P736:begin
\section*{Sources and editorial notes}\BookHeading{nav-ch16-learner-0035}{subsection}{Sources and editorial notes}{1}
% mdparagraph:P736:end

\begingroup\fontsize{10}{12.8}\selectfont
\setlength{\parskip}{1.5pt}
% mdparagraph:P737:begin
\textsuperscript{1}: Cook, \emph{A First Book in Old English}, §§43–53, with the noun entries in \emph{Exercises in Old English}, pp. 18–22. For nętt, Wright, \emph{Old English Grammar}, §§355–356 explicitly assigns nett to the cyn(n) type. The limited singular meanings are teaching limits, not claims of historical defectiveness.
% mdparagraph:P737:end

% mdparagraph:P738:begin
\textsuperscript{2}: Cook §§57–62 and 79; Wright §§435–436 for the u/\allowbreak{}o/\allowbreak{}w alternation of nearu. The tables retain the selected complete uncontracted systems where alternatives are possible.
% mdparagraph:P738:end

% mdparagraph:P739:begin
\textsuperscript{3}: Cook §§95–116 and the cited vocabulary entries; Wright pp. 262–263 for contracted slēan/\allowbreak{}ðwēan preterites and pp. 270–274 for the gierwan/\allowbreak{}smierwan pattern. The selected spellings are teaching forms; attested alternatives are not thereby declared ungrammatical.
% mdparagraph:P739:end

% mdparagraph:P740:begin
\textsuperscript{4}: Cook §95, note: the plural command has the shortened form before a following subject pronoun. Here the additional productive position is specifically dēme gē.
% mdparagraph:P740:end

% mdparagraph:P741:begin
\textsuperscript{5}: Bosworth–Toller, s.v. cunnan, distinguishes acquaintance with a person from the infinitival ability construction; personal-object witan is also attested and is expressly requested by Cook 35.7.
% mdparagraph:P741:end

% mdparagraph:P742:begin
\textsuperscript{6}: Cook, \emph{Exercises in Old English} (1895), joint Group XII, XIII, pp. 18–22, Exercises 32–36. All 59 numbered prompts and their controls are retained.
% mdparagraph:P742:end

% mdparagraph:P743:begin
\textsuperscript{7}: The printed hint in 34.9 confuses weak āgyltan with strong āgyldan “repay; reward”. Cook’s strong-verb treatment, Bosworth–Toller s.v. ā-gildan, and the West Saxon Matthew 16.27 support the correction and contracted āgylt. The English prompt itself is unchanged.
% mdparagraph:P743:end

% mdparagraph:P744:begin
\textsuperscript{8}: The West Saxon Gospels, Matthew 14.8 and 14.10, printed p. 114, have genitive Iohannes and accusative Iohannem. Initial I → J is the only declared normalization in the supplied name forms.
% mdparagraph:P744:end

% mdparagraph:P745:begin
\textsuperscript{9}: Cook, \emph{A First Book in Old English}, second edition (1895), pp. 168–170, “The Incidents in the Gymnasium”; one-based PDF pp. 184–186. All 54 original notes and 59 anchors are retained. Modern help interprets the exact reading without replacing it by its Latin parallels. For ceroma, the Latin lexicon meaning is wrestlers’ ointment.
% mdparagraph:P745:end

\endgroup

% book-source-body:end
\par\label{end-ch16-learner}
\clearpage\endgroup

\BookRecto
\begingroup
\renewcommand{\runningtitle}{Unit 17. Exposition and praise}
\hyphenpenalty=10000
\exhyphenpenalty=10000
\emergencystretch=3em
\newcommand{\itempara}[2]{\par\addvspace{3pt}\noindent\hangindent=1.8em\hangafter=1\makebox[1.8em][l]{#1.}#2\par}
\newcommand{\reserve}[1]{\par\penalty-100\begingroup\dimen0=\pagegoal\advance\dimen0 by-\pagetotal\ifdim\dimen0<#1\newpage\fi\endgroup}
% book-source-body:begin

\addfontfeatures{Ligatures=TeXOff}
\raggedright
\booktitle{\BookComponent{ch17-learner}{Unit 17. Connected exposition, praise and figurative language}UNIT 17}{Connected exposition, praise and figurative language}

% mdparagraph:P001:begin
A connected passage uses familiar clauses for different purposes: to report what happened, compare people or actions, praise, request, and interpret an event. Keep the grammar of each clause distinct from the meaning supplied by its setting. A word can retain its ordinary inflection while taking a new contextual sense.
% mdparagraph:P001:end

% mdparagraph:P002:begin
Work through the clause review first, then study the new systems in manageable groups before the form and composition practice. Cook’s complete joint Group XII, XIII, Exercises 37–43 follow. The final reading contains a complete musical performance, criticism, demonstration and audience response. The four translated examples and the lettered tasks are new teaching material.
% mdparagraph:P002:end

% mdparagraph:P003:begin
\textbf{How to use the help.} A dictionary form gives a word and meaning; you may still need to form the required inflection. A complete expression labelled \textbf{Supplied} stays available while you answer and requires no independent recall of its unfamiliar parts. A familiar form does not automatically supply a new meaning. Reference-only and recognition forms are identified with their tables. In verb tables, 1, 2 and 3 denote person; singular and plural address remain distinct.
% mdparagraph:P003:end

\reserve{7\baselineskip}
% mdparagraph:P004:begin
\section*{1. Follow the clauses}\BookHeading{nav-ch17-learner-0001}{subsection}{1. Follow the clauses}{1}
% mdparagraph:P004:end

% mdparagraph:P005:begin
First identify each finite verb, its subject and its complements. Then identify how the clauses are joined. A relative expression can represent a different role inside its clause from the role of its antecedent outside. A content clause reports a proposition. A comparison can repeat an earlier action with dōn. A condition establishes a setting for the following statement or command.
% mdparagraph:P005:end

% teaching: NEW-U17-GRAM-U17-LITERAL-CONTEXTUAL-INTERPRETATION
% mdparagraph:P006:begin
Interpretation\BookPoint{bk-ch17-learner-4a295f1a1bed} comes after this grammatical account. Distinguish an ordinary word meaning, a supplied idiom, and an image whose force depends on context. A complete translation retains the participants and relations as well as the image; an explanation of the image’s point cannot replace the translated sentence.
% mdparagraph:P006:end

% teaching: NEW-U17-EXAMPLE-01-TEACHING
\reserve{5\baselineskip}
% mdparagraph:P007:begin
\textbf{Se cyning sægde ðæt se mann wǣre gōd.}
% mdparagraph:P007:end

% mdparagraph:P008:begin
The king said that the man was good.
% mdparagraph:P008:end

% mdparagraph:P009:begin
Report the man’s condition at the time of the king’s statement; the subjunctive does not itself pronounce the report false.
% mdparagraph:P009:end

% teaching: NEW-U17-EXAMPLE-02-TEACHING
\reserve{5\baselineskip}
% mdparagraph:P010:begin
\textbf{Se cyning sang bęt ðonne se mann dyde.}
% mdparagraph:P010:end

% mdparagraph:P011:begin
The king sang better than the man did.
% mdparagraph:P011:end

% mdparagraph:P012:begin
Dōn repeats singing and has its own subject and tense; bęt modifies the action.
% mdparagraph:P012:end

% teaching: NEW-U17-EXAMPLE-03-TEACHING
\reserve{5\baselineskip}
% mdparagraph:P013:begin
\textbf{Ic geseah ðone mann ðe sang.}
% mdparagraph:P013:end

% mdparagraph:P014:begin
I saw the man who sang.
% mdparagraph:P014:end

% mdparagraph:P015:begin
The man is the object of saw and the subject inside the relative clause; the invariant particle shows no separate case ending.
% mdparagraph:P015:end

% teaching: NEW-U17-EXAMPLE-04-TEACHING
\reserve{5\baselineskip}
% mdparagraph:P016:begin
\textbf{Gif hē cymð, blissiað.}
% mdparagraph:P016:end

% mdparagraph:P017:begin
If he comes, rejoice.
% mdparagraph:P017:end

% mdparagraph:P018:begin
An open condition can use the indicative; the plural command addresses the listeners.
% mdparagraph:P018:end

% mdparagraph:P019:begin
A subjunctive in a report does not by itself say that the report is false. A past form can be simultaneous with a past statement; anteriority needs further evidence. In a condition, establish time and the speaker’s view of its possibility from the whole context. In an equal comparison, swā … swā gives the two sides; in an unequal comparison, identify what the comparative modifies and the standard introduced by ðonne.
% mdparagraph:P019:end

\reserve{8\baselineskip}
% mdparagraph:P020:begin
\subsection*{A01}\BookHeading{nav-ch17-learner-0002}{subsubsection}{A01}{0}
% mdparagraph:P020:end

% item: NEW-U17-A01
% mdparagraph:P021:begin
Analyse the report. The man’s goodness is presented as simultaneous with the king’s past statement.
% mdparagraph:P021:end

% mdparagraph:P022:begin
Se cyning sægde ðæt se mann wǣre gōd.
% mdparagraph:P022:end

% output: NEW-U17-A01-O001
% mdparagraph:P023:begin
\itempara{1}{Identify the two subjects and their finite verbs.}
% mdparagraph:P023:end

% output: NEW-U17-A01-O002
% mdparagraph:P024:begin
\itempara{2}{Classify ðæt and explain the reported time relation.}
% mdparagraph:P024:end

% output: NEW-U17-A01-O003
% mdparagraph:P025:begin
\itempara{3}{Does wǣre require the conclusion that the man was not good? Explain.}
% mdparagraph:P025:end

\reserve{8\baselineskip}
% mdparagraph:P026:begin
\subsection*{A02}\BookHeading{nav-ch17-learner-0003}{subsubsection}{A02}{0}
% mdparagraph:P026:end

% item: NEW-U17-A02
% mdparagraph:P027:begin
Analyse the comparison; then replace the first subject by the men while keeping the king singular.
% mdparagraph:P027:end

% mdparagraph:P028:begin
Se mann sang bęt ðonne se cyning dyde.
% mdparagraph:P028:end

% output: NEW-U17-A02-O001
% mdparagraph:P029:begin
\itempara{1}{Translate the sentence.}
% mdparagraph:P029:end

% output: NEW-U17-A02-O002
% mdparagraph:P030:begin
\itempara{2}{State the action repeated by dyde and explain why it is singular.}
% mdparagraph:P030:end

% output: NEW-U17-A02-O003
% mdparagraph:P031:begin
\itempara{3}{Write the complete revised Old English sentence.}
% mdparagraph:P031:end

\reserve{8\baselineskip}
% mdparagraph:P032:begin
\subsection*{A03}\BookHeading{nav-ch17-learner-0004}{subsubsection}{A03}{0}
% mdparagraph:P032:end

% item: NEW-U17-A03
% mdparagraph:P033:begin
Analyse the relative clause; then replace the man by the men.
% mdparagraph:P033:end

% mdparagraph:P034:begin
Ic geseah ðone mann ðe sang.
% mdparagraph:P034:end

% output: NEW-U17-A03-O001
% mdparagraph:P035:begin
\itempara{1}{Give the outer case and role of ðone mann.}
% mdparagraph:P035:end

% output: NEW-U17-A03-O002
% mdparagraph:P036:begin
\itempara{2}{Give the role represented by ðe inside the relative clause.}
% mdparagraph:P036:end

% output: NEW-U17-A03-O003
% mdparagraph:P037:begin
\itempara{3}{Write the complete revised Old English sentence.}
% mdparagraph:P037:end

\reserve{8\baselineskip}
% mdparagraph:P038:begin
\subsection*{A04}\BookHeading{nav-ch17-learner-0005}{subsubsection}{A04}{0}
% mdparagraph:P038:end

% item: NEW-U17-A04
% mdparagraph:P039:begin
A concerns a possible visit tomorrow. In B the speaker treats his coming now as contrary to fact. Analyse both clauses in each sentence.
% mdparagraph:P039:end

% mdparagraph:P040:begin
\textbf{A.} Gif hē cymð, ic fare.
% mdparagraph:P040:end

% mdparagraph:P041:begin
\textbf{B.} Gif hē cōme, ic fōre.
% mdparagraph:P041:end

% output: NEW-U17-A04-O001
% mdparagraph:P042:begin
\itempara{1}{Translate A and name the mood of each verb.}
% mdparagraph:P042:end

% output: NEW-U17-A04-O002
% mdparagraph:P043:begin
\itempara{2}{Translate B and name the tense and mood of each verb.}
% mdparagraph:P043:end

% output: NEW-U17-A04-O003
% mdparagraph:P044:begin
\itempara{3}{Explain which evidence supplies the time and factual status.}
% mdparagraph:P044:end

\reserve{7\baselineskip}
% mdparagraph:P045:begin
\section*{2. Nouns and adjectives}\BookHeading{nav-ch17-learner-0006}{subsection}{2. Nouns and adjectives}{1}
% mdparagraph:P045:end

% mdparagraph:P046:begin
Choose case from the word’s work in the clause. A subject is nominative, an ordinary direct object accusative, a noun-dependent possessor genitive, and a recipient or governed complement takes its licensed case. A preposition or a particular verb can change the expected choice. Each table below gives all four cases for its stated number and meaning.
% mdparagraph:P046:end

\reserve{7\baselineskip}
% mdparagraph:P047:begin
\subsection*{Strong masculine nouns}\BookHeading{nav-ch17-learner-0007}{subsubsection}{Strong masculine nouns}{0}
% mdparagraph:P047:end

% mdparagraph:P048:begin
The ordinary masculine pattern has a bare nominative/\allowbreak{}accusative singular, genitive -es and dative -e; the plural has nominative/\allowbreak{}accusative -as, genitive -a and dative -um. The -e and agent-noun patterns and any changed stems are explained individually.
% mdparagraph:P048:end

\reserve{7\baselineskip}
% mdparagraph:P049:begin
\subsection*{hyll}\BookHeading{nav-ch17-learner-0008}{subsubsection}{hyll}{0}
% mdparagraph:P049:end

% teaching: NEW-U17-VOC-HYLL-U17-SENSE
\par\begin{minipage}{\linewidth}\raggedright
% mdparagraph:P050:begin
\textbf{hyll} — hill.
% mdparagraph:P050:end

% teaching: NEW-U17-NOM-HYLL
% mdparagraph:P051:begin
Masculine; stem \textbf{hyll-}. Add the endings of the stated pattern to this stem, preserving its vowel length.
% mdparagraph:P051:end

% table: NEW-U17-NOM-HYLL-DISPLAY
% mdtable:T01:begin
\par\begingroup
\fontsize{10.5}{13.2}\selectfont
\setlength{\LTleft}{0pt}\setlength{\LTright}{\fill}
\noindent
\begin{tabular}{@{}>{\raggedright\arraybackslash}p{39.46mm}>{\raggedright\arraybackslash}p{39.46mm}>{\raggedright\arraybackslash}p{39.46mm}@{}}
\toprule
\textbf{Case} & \textbf{Singular} & \textbf{Plural} \\
\midrule
Nominative & hyll & hyllas \\
Genitive & hylles & hylla \\
Dative & hylle & hyllum \\
Accusative & hyll & hyllas \\
\bottomrule
\end{tabular}
\par\endgroup

% mdtable:T01:end
\end{minipage}\par
\reserve{7\baselineskip}
% mdparagraph:P052:begin
\subsection*{cēderbēam}\BookHeading{nav-ch17-learner-0009}{subsubsection}{cēderbēam}{0}
% mdparagraph:P052:end

% teaching: NEW-U17-VOC-CEDERBEAM-U17-SENSE
\par\begin{minipage}{\linewidth}\raggedright
% mdparagraph:P053:begin
\textbf{cēderbēam} — cedar; cedar tree.
% mdparagraph:P053:end

% teaching: NEW-U17-NOM-CEDERBEAM
% mdparagraph:P054:begin
Masculine; stem \textbf{cēderbēam-}. Add the endings of the stated pattern to this stem, preserving its vowel length.
% mdparagraph:P054:end

% table: NEW-U17-NOM-CEDERBEAM-DISPLAY
% mdtable:T02:begin
\par\begingroup
\fontsize{10.5}{13.2}\selectfont
\setlength{\LTleft}{0pt}\setlength{\LTright}{\fill}
\noindent
\begin{tabular}{@{}>{\raggedright\arraybackslash}p{36.99mm}>{\raggedright\arraybackslash}p{40.69mm}>{\raggedright\arraybackslash}p{40.69mm}@{}}
\toprule
\textbf{Case} & \textbf{Singular} & \textbf{Plural} \\
\midrule
Nominative & cēderbēam & cēderbēamas \\
Genitive & cēderbēames & cēderbēama \\
Dative & cēderbēame & cēderbēamum \\
Accusative & cēderbēam & cēderbēamas \\
\bottomrule
\end{tabular}
\par\endgroup

% mdtable:T02:end
\end{minipage}\par
\reserve{7\baselineskip}
% mdparagraph:P055:begin
\subsection*{horn}\BookHeading{nav-ch17-learner-0010}{subsubsection}{horn}{0}
% mdparagraph:P055:end

% teaching: NEW-U17-VOC-HORN-U17-SENSE
\par\begin{minipage}{\linewidth}\raggedright
% mdparagraph:P056:begin
\textbf{horn} — horn; figuratively strength or dignity in 37.14.
% mdparagraph:P056:end

% teaching: NEW-U17-NOM-HORN
% mdparagraph:P057:begin
Masculine; stem \textbf{horn-}. This table selects masculine horn, plural hornas. In Cook 37.14 the horn represents strength or dignity; retain the image when translating.
% mdparagraph:P057:end

% table: NEW-U17-NOM-HORN-DISPLAY
% mdtable:T03:begin
\par\begingroup
\fontsize{10.5}{13.2}\selectfont
\setlength{\LTleft}{0pt}\setlength{\LTright}{\fill}
\noindent
\begin{tabular}{@{}>{\raggedright\arraybackslash}p{39.46mm}>{\raggedright\arraybackslash}p{39.46mm}>{\raggedright\arraybackslash}p{39.46mm}@{}}
\toprule
\textbf{Case} & \textbf{Singular} & \textbf{Plural} \\
\midrule
Nominative & horn & hornas \\
Genitive & hornes & horna \\
Dative & horne & hornum \\
Accusative & horn & hornas \\
\bottomrule
\end{tabular}
\par\endgroup

% mdtable:T03:end
\end{minipage}\par
\reserve{7\baselineskip}
% mdparagraph:P058:begin
\subsection*{lof}\BookHeading{nav-ch17-learner-0011}{subsubsection}{lof}{0}
% mdparagraph:P058:end

% teaching: NEW-U17-VOC-LOF-U15-SOURCE-SENSE
\par\begin{minipage}{\linewidth}\raggedright
% mdparagraph:P059:begin
\textbf{lof} — praise.
% mdparagraph:P059:end

% teaching: NEW-U17-NOM-LOF
% mdparagraph:P060:begin
Masculine; stem \textbf{lof-}. Add the endings of the stated pattern to this stem, preserving its vowel length.
% mdparagraph:P060:end

% mdparagraph:P061:begin
This table gives the complete singular for the stated abstract or mass meaning. The familiar noun class is reused; plural meanings and forms are not assigned here.
% mdparagraph:P061:end

% table: NEW-U17-NOM-LOF-DISPLAY
% mdtable:T04:begin
\par\begingroup
\fontsize{10.5}{13.2}\selectfont
\setlength{\LTleft}{0pt}\setlength{\LTright}{\fill}
\noindent
\begin{tabular}{@{}>{\raggedright\arraybackslash}p{60.59mm}>{\raggedright\arraybackslash}p{60.59mm}@{}}
\toprule
\textbf{Case} & \textbf{Singular} \\
\midrule
Nominative & lof \\
Genitive & lofes \\
Dative & lofe \\
Accusative & lof \\
\bottomrule
\end{tabular}
\par\endgroup

% mdtable:T04:end
\end{minipage}\par
\reserve{7\baselineskip}
% mdparagraph:P062:begin
\subsection*{olfend}\BookHeading{nav-ch17-learner-0012}{subsubsection}{olfend}{0}
% mdparagraph:P062:end

% teaching: NEW-U17-VOC-OLFEND-U17-SENSE
\par\begin{minipage}{\linewidth}\raggedright
% mdparagraph:P063:begin
\textbf{olfend} — camel.
% mdparagraph:P063:end

% teaching: NEW-U17-NOM-OLFEND
% mdparagraph:P064:begin
Masculine; stem \textbf{olfend-}. Retain e in the stem olfend-. The final letters -end do not make this an agent noun of the gyltend pattern.
% mdparagraph:P064:end

% table: NEW-U17-NOM-OLFEND-DISPLAY
% mdtable:T05:begin
\par\begingroup
\fontsize{10.5}{13.2}\selectfont
\setlength{\LTleft}{0pt}\setlength{\LTright}{\fill}
\noindent
\begin{tabular}{@{}>{\raggedright\arraybackslash}p{39.46mm}>{\raggedright\arraybackslash}p{39.46mm}>{\raggedright\arraybackslash}p{39.46mm}@{}}
\toprule
\textbf{Case} & \textbf{Singular} & \textbf{Plural} \\
\midrule
Nominative & olfend & olfendas \\
Genitive & olfendes & olfenda \\
Dative & olfende & olfendum \\
Accusative & olfend & olfendas \\
\bottomrule
\end{tabular}
\par\endgroup

% mdtable:T05:end
\end{minipage}\par
\reserve{7\baselineskip}
% mdparagraph:P065:begin
\subsection*{ðegn}\BookHeading{nav-ch17-learner-0013}{subsubsection}{ðegn}{0}
% mdparagraph:P065:end

% teaching: NEW-U17-VOC-THEGN-U17-SENSE
\par\begin{minipage}{\linewidth}\raggedright
% mdparagraph:P066:begin
\textbf{ðegn} — vassal; retainer; thane.
% mdparagraph:P066:end

% teaching: NEW-U17-NOM-THEGN
% mdparagraph:P067:begin
Masculine; stem \textbf{ðegn-}. Add the endings of the stated pattern to this stem, preserving its vowel length.
% mdparagraph:P067:end

% table: NEW-U17-NOM-THEGN-DISPLAY
% mdtable:T06:begin
\par\begingroup
\fontsize{10.5}{13.2}\selectfont
\setlength{\LTleft}{0pt}\setlength{\LTright}{\fill}
\noindent
\begin{tabular}{@{}>{\raggedright\arraybackslash}p{39.46mm}>{\raggedright\arraybackslash}p{39.46mm}>{\raggedright\arraybackslash}p{39.46mm}@{}}
\toprule
\textbf{Case} & \textbf{Singular} & \textbf{Plural} \\
\midrule
Nominative & ðegn & ðegnas \\
Genitive & ðegnes & ðegna \\
Dative & ðegne & ðegnum \\
Accusative & ðegn & ðegnas \\
\bottomrule
\end{tabular}
\par\endgroup

% mdtable:T06:end
\end{minipage}\par
\reserve{7\baselineskip}
% mdparagraph:P068:begin
\subsection*{wīsdōm}\BookHeading{nav-ch17-learner-0014}{subsubsection}{wīsdōm}{0}
% mdparagraph:P068:end

% teaching: NEW-U17-VOC-WISDOM-U17-SENSE
\par\begin{minipage}{\linewidth}\raggedright
% mdparagraph:P069:begin
\textbf{wīsdōm} — wisdom.
% mdparagraph:P069:end

% teaching: NEW-U17-NOM-WISDOM
% mdparagraph:P070:begin
Masculine; stem \textbf{wīsdōm-}. Retain both lengths in wīsdōm. The final member follows the familiar masculine dōm pattern; only abstract singular wisdom is used here.
% mdparagraph:P070:end

% table: NEW-U17-NOM-WISDOM-DISPLAY
% mdtable:T07:begin
\par\begingroup
\fontsize{10.5}{13.2}\selectfont
\setlength{\LTleft}{0pt}\setlength{\LTright}{\fill}
\noindent
\begin{tabular}{@{}>{\raggedright\arraybackslash}p{60.59mm}>{\raggedright\arraybackslash}p{60.59mm}@{}}
\toprule
\textbf{Case} & \textbf{Singular} \\
\midrule
Nominative & wīsdōm \\
Genitive & wīsdōmes \\
Dative & wīsdōme \\
Accusative & wīsdōm \\
\bottomrule
\end{tabular}
\par\endgroup

% mdtable:T07:end
\end{minipage}\par
\reserve{7\baselineskip}
% mdparagraph:P071:begin
\subsection*{grundweall}\BookHeading{nav-ch17-learner-0015}{subsubsection}{grundweall}{0}
% mdparagraph:P071:end

% teaching: NEW-U17-VOC-GRUNDWEALL-U17-SENSE
\par\begin{minipage}{\linewidth}\raggedright
% mdparagraph:P072:begin
\textbf{grundweall} — foundation.
% mdparagraph:P072:end

% teaching: NEW-U17-NOM-GRUNDWEALL
% mdparagraph:P073:begin
Masculine; stem \textbf{grundweall-}. Add the endings of the stated pattern to this stem, preserving its vowel length.
% mdparagraph:P073:end

% table: NEW-U17-NOM-GRUNDWEALL-DISPLAY
% mdtable:T08:begin
\par\begingroup
\fontsize{10.5}{13.2}\selectfont
\setlength{\LTleft}{0pt}\setlength{\LTright}{\fill}
\noindent
\begin{tabular}{@{}>{\raggedright\arraybackslash}p{34.82mm}>{\raggedright\arraybackslash}p{41.78mm}>{\raggedright\arraybackslash}p{41.78mm}@{}}
\toprule
\textbf{Case} & \textbf{Singular} & \textbf{Plural} \\
\midrule
Nominative & grundweall & grundweallas \\
Genitive & grundwealles & grundwealla \\
Dative & grundwealle & grundweallum \\
Accusative & grundweall & grundweallas \\
\bottomrule
\end{tabular}
\par\endgroup

% mdtable:T08:end
\end{minipage}\par
\reserve{7\baselineskip}
% mdparagraph:P074:begin
\subsection*{cniht}\BookHeading{nav-ch17-learner-0016}{subsubsection}{cniht}{0}
% mdparagraph:P074:end

% teaching: NEW-U17-VOC-CNIHT-U17-SERVANT
\par\begin{minipage}{\linewidth}\raggedright
% mdparagraph:P075:begin
\textbf{cniht} — servant.
% mdparagraph:P075:end

% teaching: NEW-U17-NOM-CNIHT
% mdparagraph:P076:begin
Masculine; stem \textbf{cniht-}. The meaning required here is servant. It is distinct from the earlier occurrence meaning youth; familiarity with the word in that passage did not supply this new meaning.
% mdparagraph:P076:end

% table: NEW-U17-NOM-CNIHT-DISPLAY
% mdtable:T09:begin
\par\begingroup
\fontsize{10.5}{13.2}\selectfont
\setlength{\LTleft}{0pt}\setlength{\LTright}{\fill}
\noindent
\begin{tabular}{@{}>{\raggedright\arraybackslash}p{39.46mm}>{\raggedright\arraybackslash}p{39.46mm}>{\raggedright\arraybackslash}p{39.46mm}@{}}
\toprule
\textbf{Case} & \textbf{Singular} & \textbf{Plural} \\
\midrule
Nominative & cniht & cnihtas \\
Genitive & cnihtes & cnihta \\
Dative & cnihte & cnihtum \\
Accusative & cniht & cnihtas \\
\bottomrule
\end{tabular}
\par\endgroup

% mdtable:T09:end
\end{minipage}\par
\reserve{7\baselineskip}
% mdparagraph:P077:begin
\subsection*{ymbhwyrft}\BookHeading{nav-ch17-learner-0017}{subsubsection}{ymbhwyrft}{0}
% mdparagraph:P077:end

% teaching: NEW-U17-VOC-YMBHWYRFT-U13-SOURCE
\par\begin{minipage}{\linewidth}\raggedright
% mdparagraph:P078:begin
\textbf{ymbhwyrft} — circuit; compass; orbit.
% mdparagraph:P078:end

% teaching: NEW-U17-NOM-YMBHWYRFT
% mdparagraph:P079:begin
Masculine; stem \textbf{ymbhwyrft-}. Add the endings of the stated pattern to this stem, preserving its vowel length.
% mdparagraph:P079:end

% table: NEW-U17-NOM-YMBHWYRFT-DISPLAY
% mdtable:T10:begin
\par\begingroup
\fontsize{10.5}{13.2}\selectfont
\setlength{\LTleft}{0pt}\setlength{\LTright}{\fill}
\noindent
\begin{tabular}{@{}>{\raggedright\arraybackslash}p{36.99mm}>{\raggedright\arraybackslash}p{40.69mm}>{\raggedright\arraybackslash}p{40.69mm}@{}}
\toprule
\textbf{Case} & \textbf{Singular} & \textbf{Plural} \\
\midrule
Nominative & ymbhwyrft & ymbhwyrftas \\
Genitive & ymbhwyrftes & ymbhwyrfta \\
Dative & ymbhwyrfte & ymbhwyrftum \\
Accusative & ymbhwyrft & ymbhwyrftas \\
\bottomrule
\end{tabular}
\par\endgroup

% mdtable:T10:end
\end{minipage}\par
\reserve{7\baselineskip}
% mdparagraph:P080:begin
\subsection*{gylt}\BookHeading{nav-ch17-learner-0018}{subsubsection}{gylt}{0}
% mdparagraph:P080:end

% teaching: NEW-U17-VOC-GYLT-U17-SENSE
\par\begin{minipage}{\linewidth}\raggedright
% mdparagraph:P081:begin
\textbf{gylt} — debt, including the prayer’s figurative moral debt.
% mdparagraph:P081:end

% teaching: NEW-U17-NOM-GYLT
% mdparagraph:P082:begin
Masculine; stem \textbf{gylt-}. Retain debt in translating the prayer. Its moral force comes from the prayer’s context; an ordinary noun ending does not by itself express that interpretation.
% mdparagraph:P082:end

% table: NEW-U17-NOM-GYLT-DISPLAY
% mdtable:T11:begin
\par\begingroup
\fontsize{10.5}{13.2}\selectfont
\setlength{\LTleft}{0pt}\setlength{\LTright}{\fill}
\noindent
\begin{tabular}{@{}>{\raggedright\arraybackslash}p{39.46mm}>{\raggedright\arraybackslash}p{39.46mm}>{\raggedright\arraybackslash}p{39.46mm}@{}}
\toprule
\textbf{Case} & \textbf{Singular} & \textbf{Plural} \\
\midrule
Nominative & gylt & gyltas \\
Genitive & gyltes & gylta \\
Dative & gylte & gyltum \\
Accusative & gylt & gyltas \\
\bottomrule
\end{tabular}
\par\endgroup

% mdtable:T11:end
\end{minipage}\par
\reserve{7\baselineskip}
% mdparagraph:P083:begin
\subsection*{gyltend}\BookHeading{nav-ch17-learner-0019}{subsubsection}{gyltend}{0}
% mdparagraph:P083:end

% teaching: NEW-U17-VOC-GYLTEND-U17-SENSE
\par\begin{minipage}{\linewidth}\raggedright
% mdparagraph:P084:begin
\textbf{gyltend} — debtor.
% mdparagraph:P084:end

% teaching: NEW-U17-NOM-GYLTEND
% mdparagraph:P085:begin
Masculine; stem \textbf{gyltend-}. This masculine agent noun has selected nominative/\allowbreak{}accusative plural gyltende and genitive plural gyltendra. An endingless nominative/\allowbreak{}accusative plural is also recorded. The selected pattern is not the ordinary -as plural.
% mdparagraph:P085:end

% table: NEW-U17-NOM-GYLTEND-DISPLAY
% mdtable:T12:begin
\par\begingroup
\fontsize{10.5}{13.2}\selectfont
\setlength{\LTleft}{0pt}\setlength{\LTright}{\fill}
\noindent
\begin{tabular}{@{}>{\raggedright\arraybackslash}p{39.46mm}>{\raggedright\arraybackslash}p{39.46mm}>{\raggedright\arraybackslash}p{39.46mm}@{}}
\toprule
\textbf{Case} & \textbf{Singular} & \textbf{Plural} \\
\midrule
Nominative & gyltend & gyltende \\
Genitive & gyltendes & gyltendra \\
Dative & gyltende & gyltendum \\
Accusative & gyltend & gyltende \\
\bottomrule
\end{tabular}
\par\endgroup

% mdtable:T12:end
\end{minipage}\par
\reserve{7\baselineskip}
% mdparagraph:P086:begin
\subsection*{finger}\BookHeading{nav-ch17-learner-0020}{subsubsection}{finger}{0}
% mdparagraph:P086:end

% teaching: NEW-U17-VOC-FINGER-U17-SENSE
\par\begin{minipage}{\linewidth}\raggedright
% mdparagraph:P087:begin
\textbf{finger} — finger.
% mdparagraph:P087:end

% teaching: NEW-U17-NOM-FINGER
% mdparagraph:P088:begin
Masculine; stem \textbf{fingr-}. The second-syllable vowel disappears before an ending: finger, but fingres and fingre. Use the oblique stem fingr-.
% mdparagraph:P088:end

% table: NEW-U17-NOM-FINGER-DISPLAY
% mdtable:T13:begin
\par\begingroup
\fontsize{10.5}{13.2}\selectfont
\setlength{\LTleft}{0pt}\setlength{\LTright}{\fill}
\noindent
\begin{tabular}{@{}>{\raggedright\arraybackslash}p{39.46mm}>{\raggedright\arraybackslash}p{39.46mm}>{\raggedright\arraybackslash}p{39.46mm}@{}}
\toprule
\textbf{Case} & \textbf{Singular} & \textbf{Plural} \\
\midrule
Nominative & finger & fingras \\
Genitive & fingres & fingra \\
Dative & fingre & fingrum \\
Accusative & finger & fingras \\
\bottomrule
\end{tabular}
\par\endgroup

% mdtable:T13:end
\end{minipage}\par
\reserve{7\baselineskip}
% mdparagraph:P089:begin
\subsection*{hearpere}\BookHeading{nav-ch17-learner-0021}{subsubsection}{hearpere}{0}
% mdparagraph:P089:end

% teaching: NEW-U17-VOC-HEARPERE-U17-SENSE
\par\begin{minipage}{\linewidth}\raggedright
% mdparagraph:P090:begin
\textbf{hearpere} — harper.
% mdparagraph:P090:end

% teaching: NEW-U17-NOM-HEARPERE
% mdparagraph:P091:begin
Masculine; stem \textbf{hearper-}. The final -e of the headword gives way to endings on hearper-. The genitive plural is hearpera, as in the familiar masculine -ere class.
% mdparagraph:P091:end

% table: NEW-U17-NOM-HEARPERE-DISPLAY
% mdtable:T14:begin
\par\begingroup
\fontsize{10.5}{13.2}\selectfont
\setlength{\LTleft}{0pt}\setlength{\LTright}{\fill}
\noindent
\begin{tabular}{@{}>{\raggedright\arraybackslash}p{39.46mm}>{\raggedright\arraybackslash}p{39.46mm}>{\raggedright\arraybackslash}p{39.46mm}@{}}
\toprule
\textbf{Case} & \textbf{Singular} & \textbf{Plural} \\
\midrule
Nominative & hearpere & hearperas \\
Genitive & hearperes & hearpera \\
Dative & hearpere & hearperum \\
Accusative & hearpere & hearperas \\
\bottomrule
\end{tabular}
\par\endgroup

% mdtable:T14:end
\end{minipage}\par
\reserve{7\baselineskip}
% mdparagraph:P092:begin
\subsection*{weorðscipe}\BookHeading{nav-ch17-learner-0022}{subsubsection}{weorðscipe}{0}
% mdparagraph:P092:end

% teaching: NEW-U17-VOC-WEORTHSCIPE-U17-SENSE
\par\begin{minipage}{\linewidth}\raggedright
% mdparagraph:P093:begin
\textbf{weorðscipe} — honour; dignity.
% mdparagraph:P093:end

% teaching: NEW-U17-NOM-WEORTHSCIPE
% mdparagraph:P094:begin
Masculine; stem \textbf{weorðscip-}. This masculine -e noun retains -e in the nominative, dative and accusative singular. The genitive has -es on weorðscip-. Only singular honour or dignity is assigned here.
% mdparagraph:P094:end

% mdparagraph:P095:begin
This table gives the complete singular for the stated abstract or mass meaning. The familiar noun class is reused; plural meanings and forms are not assigned here.
% mdparagraph:P095:end

% table: NEW-U17-NOM-WEORTHSCIPE-DISPLAY
% mdtable:T15:begin
\par\begingroup
\fontsize{10.5}{13.2}\selectfont
\setlength{\LTleft}{0pt}\setlength{\LTright}{\fill}
\noindent
\begin{tabular}{@{}>{\raggedright\arraybackslash}p{57.71mm}>{\raggedright\arraybackslash}p{63.48mm}@{}}
\toprule
\textbf{Case} & \textbf{Singular} \\
\midrule
Nominative & weorðscipe \\
Genitive & weorðscipes \\
Dative & weorðscipe \\
Accusative & weorðscipe \\
\bottomrule
\end{tabular}
\par\endgroup

% mdtable:T15:end
\end{minipage}\par
\reserve{7\baselineskip}
% mdparagraph:P096:begin
\subsection*{fox}\BookHeading{nav-ch17-learner-0023}{subsubsection}{fox}{0}
% mdparagraph:P096:end

% teaching: NEW-U17-VOC-FOX-U17-SENSE
\par\begin{minipage}{\linewidth}\raggedright
% mdparagraph:P097:begin
\textbf{fox} — fox.
% mdparagraph:P097:end

% teaching: NEW-U17-NOM-FOX
% mdparagraph:P098:begin
Masculine; stem \textbf{fox-}. Add the endings of the stated pattern to this stem, preserving its vowel length.
% mdparagraph:P098:end

% table: NEW-U17-NOM-FOX-DISPLAY
% mdtable:T16:begin
\par\begingroup
\fontsize{10.5}{13.2}\selectfont
\setlength{\LTleft}{0pt}\setlength{\LTright}{\fill}
\noindent
\begin{tabular}{@{}>{\raggedright\arraybackslash}p{39.46mm}>{\raggedright\arraybackslash}p{39.46mm}>{\raggedright\arraybackslash}p{39.46mm}@{}}
\toprule
\textbf{Case} & \textbf{Singular} & \textbf{Plural} \\
\midrule
Nominative & fox & foxas \\
Genitive & foxes & foxa \\
Dative & foxe & foxum \\
Accusative & fox & foxas \\
\bottomrule
\end{tabular}
\par\endgroup

% mdtable:T16:end
\end{minipage}\par
\reserve{7\baselineskip}
% mdparagraph:P099:begin
\subsection*{ðorn}\BookHeading{nav-ch17-learner-0024}{subsubsection}{ðorn}{0}
% mdparagraph:P099:end

% teaching: NEW-U17-VOC-THORN-U17-SENSE
\par\begin{minipage}{\linewidth}\raggedright
% mdparagraph:P100:begin
\textbf{ðorn} — thorn.
% mdparagraph:P100:end

% teaching: NEW-U17-NOM-THORN
% mdparagraph:P101:begin
Masculine; stem \textbf{ðorn-}. Add the endings of the stated pattern to this stem, preserving its vowel length.
% mdparagraph:P101:end

% table: NEW-U17-NOM-THORN-DISPLAY
% mdtable:T17:begin
\par\begingroup
\fontsize{10.5}{13.2}\selectfont
\setlength{\LTleft}{0pt}\setlength{\LTright}{\fill}
\noindent
\begin{tabular}{@{}>{\raggedright\arraybackslash}p{39.46mm}>{\raggedright\arraybackslash}p{39.46mm}>{\raggedright\arraybackslash}p{39.46mm}@{}}
\toprule
\textbf{Case} & \textbf{Singular} & \textbf{Plural} \\
\midrule
Nominative & ðorn & ðornas \\
Genitive & ðornes & ðorna \\
Dative & ðorne & ðornum \\
Accusative & ðorn & ðornas \\
\bottomrule
\end{tabular}
\par\endgroup

% mdtable:T17:end
\end{minipage}\par
\reserve{7\baselineskip}
% mdparagraph:P102:begin
\subsection*{cynehelm}\BookHeading{nav-ch17-learner-0025}{subsubsection}{cynehelm}{0}
% mdparagraph:P102:end

% teaching: NEW-U17-VOC-CYNEHELM-U17-SENSE
\par\begin{minipage}{\linewidth}\raggedright
% mdparagraph:P103:begin
\textbf{cynehelm} — crown.
% mdparagraph:P103:end

% teaching: NEW-U17-NOM-CYNEHELM
% mdparagraph:P104:begin
Masculine; stem \textbf{cynehelm-}. Add the endings of the stated pattern to this stem, preserving its vowel length.
% mdparagraph:P104:end

% table: NEW-U17-NOM-CYNEHELM-DISPLAY
% mdtable:T18:begin
\par\begingroup
\fontsize{10.5}{13.2}\selectfont
\setlength{\LTleft}{0pt}\setlength{\LTright}{\fill}
\noindent
\begin{tabular}{@{}>{\raggedright\arraybackslash}p{39.46mm}>{\raggedright\arraybackslash}p{39.46mm}>{\raggedright\arraybackslash}p{39.46mm}@{}}
\toprule
\textbf{Case} & \textbf{Singular} & \textbf{Plural} \\
\midrule
Nominative & cynehelm & cynehelmas \\
Genitive & cynehelmes & cynehelma \\
Dative & cynehelme & cynehelmum \\
Accusative & cynehelm & cynehelmas \\
\bottomrule
\end{tabular}
\par\endgroup

% mdtable:T18:end
\end{minipage}\par
\reserve{7\baselineskip}
% mdparagraph:P105:begin
\subsection*{gerīmcræft}\BookHeading{nav-ch17-learner-0026}{subsubsection}{gerīmcræft}{0}
% mdparagraph:P105:end

% teaching: NEW-U17-VOC-GERIMCRAEFT-U17-SENSE
\par\begin{minipage}{\linewidth}\raggedright
% mdparagraph:P106:begin
\textbf{gerīmcræft} — arithmetic; computation.
% mdparagraph:P106:end

% teaching: NEW-U17-NOM-GERIMCRAEFT
% mdparagraph:P107:begin
Masculine; stem \textbf{gerīmcræft-}. Add the endings of the stated pattern to this stem, preserving its vowel length.
% mdparagraph:P107:end

% mdparagraph:P108:begin
This table gives the complete singular for the stated abstract or mass meaning. The familiar noun class is reused; plural meanings and forms are not assigned here.
% mdparagraph:P108:end

% table: NEW-U17-NOM-GERIMCRAEFT-DISPLAY
% mdtable:T19:begin
\par\begingroup
\fontsize{10.5}{13.2}\selectfont
\setlength{\LTleft}{0pt}\setlength{\LTright}{\fill}
\noindent
\begin{tabular}{@{}>{\raggedright\arraybackslash}p{55.09mm}>{\raggedright\arraybackslash}p{66.1mm}@{}}
\toprule
\textbf{Case} & \textbf{Singular} \\
\midrule
Nominative & gerīmcræft \\
Genitive & gerīmcræftes \\
Dative & gerīmcræfte \\
Accusative & gerīmcræft \\
\bottomrule
\end{tabular}
\par\endgroup

% mdtable:T19:end
\end{minipage}\par
\reserve{7\baselineskip}
% mdparagraph:P109:begin
\subsection*{stān}\BookHeading{nav-ch17-learner-0027}{subsubsection}{stān}{0}
% mdparagraph:P109:end

% teaching: NEW-U17-VOC-STAN-U11-SOURCE
\par\begin{minipage}{\linewidth}\raggedright
% mdparagraph:P110:begin
\textbf{stān} — stone.
% mdparagraph:P110:end

% teaching: NEW-U17-NOM-STAN
% mdparagraph:P111:begin
Masculine; stem \textbf{stān-}. The full stone paradigm is available for the gifts, building materials and bread-or-stone contrast.
% mdparagraph:P111:end

% table: NEW-U17-NOM-STAN-DISPLAY
% mdtable:T20:begin
\par\begingroup
\fontsize{10.5}{13.2}\selectfont
\setlength{\LTleft}{0pt}\setlength{\LTright}{\fill}
\noindent
\begin{tabular}{@{}>{\raggedright\arraybackslash}p{39.46mm}>{\raggedright\arraybackslash}p{39.46mm}>{\raggedright\arraybackslash}p{39.46mm}@{}}
\toprule
\textbf{Case} & \textbf{Singular} & \textbf{Plural} \\
\midrule
Nominative & stān & stānas \\
Genitive & stānes & stāna \\
Dative & stāne & stānum \\
Accusative & stān & stānas \\
\bottomrule
\end{tabular}
\par\endgroup

% mdtable:T20:end
\end{minipage}\par
\reserve{7\baselineskip}
% mdparagraph:P112:begin
\subsection*{Strong neuter nouns}\BookHeading{nav-ch17-learner-0028}{subsubsection}{Strong neuter nouns}{0}
% mdparagraph:P112:end

% mdparagraph:P113:begin
Nominative and accusative are alike in each number. The singular genitive has -es and dative -e; the plural genitive has -a and dative -um. A light stem normally has nominative/\allowbreak{}accusative plural -u, while a heavy stem normally has no ending. Learn the stated compound and exceptional stems with the table.
% mdparagraph:P113:end

\reserve{7\baselineskip}
% mdparagraph:P114:begin
\subsection*{mægen}\BookHeading{nav-ch17-learner-0029}{subsubsection}{mægen}{0}
% mdparagraph:P114:end

% teaching: NEW-U17-VOC-MAEGEN-U17-POWER
\par\begin{minipage}{\linewidth}\raggedright
% mdparagraph:P115:begin
\textbf{mægen} — power; strength; in 37.2 heavenly powers.
% mdparagraph:P115:end

% teaching: NEW-U17-NOM-MAEGEN
% mdparagraph:P116:begin
Neuter; stem \textbf{mægen-}. Retain the full stem mægen- in mægenes, mægene, mægena and mægenum. The exceptional nominative/\allowbreak{}accusative plural is mægnu. In the praise catalogue the plural denotes heavenly powers; the singular denotes power or strength.
% mdparagraph:P116:end

% table: NEW-U17-NOM-MAEGEN-DISPLAY
% mdtable:T21:begin
\par\begingroup
\fontsize{10.5}{13.2}\selectfont
\setlength{\LTleft}{0pt}\setlength{\LTright}{\fill}
\noindent
\begin{tabular}{@{}>{\raggedright\arraybackslash}p{39.46mm}>{\raggedright\arraybackslash}p{39.46mm}>{\raggedright\arraybackslash}p{39.46mm}@{}}
\toprule
\textbf{Case} & \textbf{Singular} & \textbf{Plural} \\
\midrule
Nominative & mægen & mægnu \\
Genitive & mægenes & mægena \\
Dative & mægene & mægenum \\
Accusative & mægen & mægnu \\
\bottomrule
\end{tabular}
\par\endgroup

% mdtable:T21:end
\end{minipage}\par
\reserve{7\baselineskip}
% mdparagraph:P117:begin
\subsection*{andgiet}\BookHeading{nav-ch17-learner-0030}{subsubsection}{andgiet}{0}
% mdparagraph:P117:end

% teaching: NEW-U17-VOC-ANDGIET-U15-SOURCE-SENSE
% teaching: NEW-U17-VOC-ANDGIET-U15-SOURCE-UNDERSTANDING
\par\begin{minipage}{\linewidth}\raggedright
% mdparagraph:P118:begin
\textbf{andgiet} — meaning; sense; understanding.
% mdparagraph:P118:end

% teaching: NEW-U17-NOM-ANDGIET
% mdparagraph:P119:begin
Neuter; stem \textbf{andgiet-}. Learn the four singular cases for meaning and understanding. These are abstract uses of a familiar neuter class. Plural forms and separate count meanings are outside this unit; the absence of a plural column does not mean that the word has no plural.
% mdparagraph:P119:end

% table: NEW-U17-NOM-ANDGIET-DISPLAY
% mdtable:T22:begin
\par\begingroup
\fontsize{10.5}{13.2}\selectfont
\setlength{\LTleft}{0pt}\setlength{\LTright}{\fill}
\noindent
\begin{tabular}{@{}>{\raggedright\arraybackslash}p{60.59mm}>{\raggedright\arraybackslash}p{60.59mm}@{}}
\toprule
\textbf{Case} & \textbf{Singular} \\
\midrule
Nominative & andgiet \\
Genitive & andgietes \\
Dative & andgiete \\
Accusative & andgiet \\
\bottomrule
\end{tabular}
\par\endgroup

% mdtable:T22:end
\end{minipage}\par
\reserve{7\baselineskip}
% mdparagraph:P120:begin
\subsection*{pund}\BookHeading{nav-ch17-learner-0031}{subsubsection}{pund}{0}
% mdparagraph:P120:end

% teaching: NEW-U17-VOC-PUND-U17-SENSE
\par\begin{minipage}{\linewidth}\raggedright
% mdparagraph:P121:begin
\textbf{pund} — pound; source weight measure.
% mdparagraph:P121:end

% teaching: NEW-U17-NOM-PUND
% mdparagraph:P122:begin
Neuter; stem \textbf{pund-}. Add the endings of the stated pattern to this stem, preserving its vowel length.
% mdparagraph:P122:end

% table: NEW-U17-NOM-PUND-DISPLAY
% mdtable:T23:begin
\par\begingroup
\fontsize{10.5}{13.2}\selectfont
\setlength{\LTleft}{0pt}\setlength{\LTright}{\fill}
\noindent
\begin{tabular}{@{}>{\raggedright\arraybackslash}p{39.46mm}>{\raggedright\arraybackslash}p{39.46mm}>{\raggedright\arraybackslash}p{39.46mm}@{}}
\toprule
\textbf{Case} & \textbf{Singular} & \textbf{Plural} \\
\midrule
Nominative & pund & pund \\
Genitive & pundes & punda \\
Dative & punde & pundum \\
Accusative & pund & pund \\
\bottomrule
\end{tabular}
\par\endgroup

% mdtable:T23:end
\end{minipage}\par
\reserve{7\baselineskip}
% mdparagraph:P123:begin
\subsection*{edlēan}\BookHeading{nav-ch17-learner-0032}{subsubsection}{edlēan}{0}
% mdparagraph:P123:end

% teaching: NEW-U17-VOC-EDLEAN-U17-SENSE
\par\begin{minipage}{\linewidth}\raggedright
% mdparagraph:P124:begin
\textbf{edlēan} — reward; recompense.
% mdparagraph:P124:end

% teaching: NEW-U17-NOM-EDLEAN
% mdparagraph:P125:begin
Neuter; stem \textbf{edlēan-}. Add the endings of the stated pattern to this stem, preserving its vowel length.
% mdparagraph:P125:end

% table: NEW-U17-NOM-EDLEAN-DISPLAY
% mdtable:T24:begin
\par\begingroup
\fontsize{10.5}{13.2}\selectfont
\setlength{\LTleft}{0pt}\setlength{\LTright}{\fill}
\noindent
\begin{tabular}{@{}>{\raggedright\arraybackslash}p{39.46mm}>{\raggedright\arraybackslash}p{39.46mm}>{\raggedright\arraybackslash}p{39.46mm}@{}}
\toprule
\textbf{Case} & \textbf{Singular} & \textbf{Plural} \\
\midrule
Nominative & edlēan & edlēan \\
Genitive & edlēanes & edlēana \\
Dative & edlēane & edlēanum \\
Accusative & edlēan & edlēan \\
\bottomrule
\end{tabular}
\par\endgroup

% mdtable:T24:end
\end{minipage}\par
\reserve{7\baselineskip}
% mdparagraph:P126:begin
\subsection*{ðrymsetl}\BookHeading{nav-ch17-learner-0033}{subsubsection}{ðrymsetl}{0}
% mdparagraph:P126:end

% teaching: NEW-U17-VOC-THRYMSETL-U17-SENSE
\par\begin{minipage}{\linewidth}\raggedright
% mdparagraph:P127:begin
\textbf{ðrymsetl} — throne; throne of glory.
% mdparagraph:P127:end

% teaching: NEW-U17-NOM-THRYMSETL
% mdparagraph:P128:begin
Neuter; stem \textbf{ðrymsetl-}. Retain the source spelling ðrymsetl. The final stem setl is heavy, so the nominative/\allowbreak{}accusative plural has no -u.
% mdparagraph:P128:end

% table: NEW-U17-NOM-THRYMSETL-DISPLAY
% mdtable:T25:begin
\par\begingroup
\fontsize{10.5}{13.2}\selectfont
\setlength{\LTleft}{0pt}\setlength{\LTright}{\fill}
\noindent
\begin{tabular}{@{}>{\raggedright\arraybackslash}p{39.46mm}>{\raggedright\arraybackslash}p{39.46mm}>{\raggedright\arraybackslash}p{39.46mm}@{}}
\toprule
\textbf{Case} & \textbf{Singular} & \textbf{Plural} \\
\midrule
Nominative & ðrymsetl & ðrymsetl \\
Genitive & ðrymsetles & ðrymsetla \\
Dative & ðrymsetle & ðrymsetlum \\
Accusative & ðrymsetl & ðrymsetl \\
\bottomrule
\end{tabular}
\par\endgroup

% mdtable:T25:end
\end{minipage}\par
\reserve{7\baselineskip}
% mdparagraph:P129:begin
\subsection*{palmtwīg}\BookHeading{nav-ch17-learner-0034}{subsubsection}{palmtwīg}{0}
% mdparagraph:P129:end

% teaching: NEW-U17-VOC-PALMTWIG-U17-SENSE
\par\begin{minipage}{\linewidth}\raggedright
% mdparagraph:P130:begin
\textbf{palmtwīg} — palm twig; palm branch.
% mdparagraph:P130:end

% teaching: NEW-U17-NOM-PALMTWIG
% mdparagraph:P131:begin
Neuter; stem \textbf{palmtwīg-}. Long ī makes palmtwīg a heavy neuter. The nominative/\allowbreak{}accusative plural is palmtwīg.
% mdparagraph:P131:end

% table: NEW-U17-NOM-PALMTWIG-DISPLAY
% mdtable:T26:begin
\par\begingroup
\fontsize{10.5}{13.2}\selectfont
\setlength{\LTleft}{0pt}\setlength{\LTright}{\fill}
\noindent
\begin{tabular}{@{}>{\raggedright\arraybackslash}p{39.46mm}>{\raggedright\arraybackslash}p{39.46mm}>{\raggedright\arraybackslash}p{39.46mm}@{}}
\toprule
\textbf{Case} & \textbf{Singular} & \textbf{Plural} \\
\midrule
Nominative & palmtwīg & palmtwīg \\
Genitive & palmtwīges & palmtwīga \\
Dative & palmtwīge & palmtwīgum \\
Accusative & palmtwīg & palmtwīg \\
\bottomrule
\end{tabular}
\par\endgroup

% mdtable:T26:end
\end{minipage}\par
\reserve{7\baselineskip}
% mdparagraph:P132:begin
\subsection*{cnēo}\BookHeading{nav-ch17-learner-0035}{subsubsection}{cnēo}{0}
% mdparagraph:P132:end

% teaching: NEW-U17-VOC-CNEO-U17-KNEE
\par\begin{minipage}{\linewidth}\raggedright
% mdparagraph:P133:begin
\textbf{cnēo} — knee.
% mdparagraph:P133:end

% teaching: NEW-U17-NOM-CNEO
% mdparagraph:P134:begin
Neuter; stem \textbf{cnēow-}. The endingless singular cnēo has cnēow- before endings. Its nominative/\allowbreak{}accusative plural is cnēowu.
% mdparagraph:P134:end

% table: NEW-U17-NOM-CNEO-DISPLAY
% mdtable:T27:begin
\par\begingroup
\fontsize{10.5}{13.2}\selectfont
\setlength{\LTleft}{0pt}\setlength{\LTright}{\fill}
\noindent
\begin{tabular}{@{}>{\raggedright\arraybackslash}p{39.46mm}>{\raggedright\arraybackslash}p{39.46mm}>{\raggedright\arraybackslash}p{39.46mm}@{}}
\toprule
\textbf{Case} & \textbf{Singular} & \textbf{Plural} \\
\midrule
Nominative & cnēo & cnēowu \\
Genitive & cnēowes & cnēowa \\
Dative & cnēowe & cnēowum \\
Accusative & cnēo & cnēowu \\
\bottomrule
\end{tabular}
\par\endgroup

% mdtable:T27:end
\end{minipage}\par
\reserve{7\baselineskip}
% mdparagraph:P135:begin
\subsection*{hol}\BookHeading{nav-ch17-learner-0036}{subsubsection}{hol}{0}
% mdparagraph:P135:end

% teaching: NEW-U17-VOC-HOL-U17-SENSE
\par\begin{minipage}{\linewidth}\raggedright
% mdparagraph:P136:begin
\textbf{hol} — hole.
% mdparagraph:P136:end

% teaching: NEW-U17-NOM-HOL
% mdparagraph:P137:begin
Neuter; stem \textbf{hol-}. The short vowel and single final consonant give the light neuter pattern: nominative/\allowbreak{}accusative plural holu.
% mdparagraph:P137:end

% table: NEW-U17-NOM-HOL-DISPLAY
% mdtable:T28:begin
\par\begingroup
\fontsize{10.5}{13.2}\selectfont
\setlength{\LTleft}{0pt}\setlength{\LTright}{\fill}
\noindent
\begin{tabular}{@{}>{\raggedright\arraybackslash}p{39.46mm}>{\raggedright\arraybackslash}p{39.46mm}>{\raggedright\arraybackslash}p{39.46mm}@{}}
\toprule
\textbf{Case} & \textbf{Singular} & \textbf{Plural} \\
\midrule
Nominative & hol & holu \\
Genitive & holes & hola \\
Dative & hole & holum \\
Accusative & hol & holu \\
\bottomrule
\end{tabular}
\par\endgroup

% mdtable:T28:end
\end{minipage}\par
\reserve{7\baselineskip}
% mdparagraph:P138:begin
\subsection*{nest}\BookHeading{nav-ch17-learner-0037}{subsubsection}{nest}{0}
% mdparagraph:P138:end

% teaching: NEW-U17-VOC-NEST-U17-SENSE
\par\begin{minipage}{\linewidth}\raggedright
% mdparagraph:P139:begin
\textbf{nest} — nest.
% mdparagraph:P139:end

% teaching: NEW-U17-NOM-NEST
% mdparagraph:P140:begin
Neuter; stem \textbf{nest-}. The final consonant cluster makes the stem heavy. The nominative/\allowbreak{}accusative plural is nest.
% mdparagraph:P140:end

% table: NEW-U17-NOM-NEST-DISPLAY
% mdtable:T29:begin
\par\begingroup
\fontsize{10.5}{13.2}\selectfont
\setlength{\LTleft}{0pt}\setlength{\LTright}{\fill}
\noindent
\begin{tabular}{@{}>{\raggedright\arraybackslash}p{39.46mm}>{\raggedright\arraybackslash}p{39.46mm}>{\raggedright\arraybackslash}p{39.46mm}@{}}
\toprule
\textbf{Case} & \textbf{Singular} & \textbf{Plural} \\
\midrule
Nominative & nest & nest \\
Genitive & nestes & nesta \\
Dative & neste & nestum \\
Accusative & nest & nest \\
\bottomrule
\end{tabular}
\par\endgroup

% mdtable:T29:end
\end{minipage}\par
\reserve{7\baselineskip}
% mdparagraph:P141:begin
\subsection*{hrēod}\BookHeading{nav-ch17-learner-0038}{subsubsection}{hrēod}{0}
% mdparagraph:P141:end

% teaching: NEW-U17-VOC-HREOD-U17-SENSE
\par\begin{minipage}{\linewidth}\raggedright
% mdparagraph:P142:begin
\textbf{hrēod} — reed.
% mdparagraph:P142:end

% teaching: NEW-U17-NOM-HREOD
% mdparagraph:P143:begin
Neuter; stem \textbf{hrēod-}. Long ēo makes hrēod a heavy neuter. Its nominative/\allowbreak{}accusative plural is endingless.
% mdparagraph:P143:end

% table: NEW-U17-NOM-HREOD-DISPLAY
% mdtable:T30:begin
\par\begingroup
\fontsize{10.5}{13.2}\selectfont
\setlength{\LTleft}{0pt}\setlength{\LTright}{\fill}
\noindent
\begin{tabular}{@{}>{\raggedright\arraybackslash}p{39.46mm}>{\raggedright\arraybackslash}p{39.46mm}>{\raggedright\arraybackslash}p{39.46mm}@{}}
\toprule
\textbf{Case} & \textbf{Singular} & \textbf{Plural} \\
\midrule
Nominative & hrēod & hrēod \\
Genitive & hrēodes & hrēoda \\
Dative & hrēode & hrēodum \\
Accusative & hrēod & hrēod \\
\bottomrule
\end{tabular}
\par\endgroup

% mdtable:T30:end
\end{minipage}\par
\reserve{7\baselineskip}
% mdparagraph:P144:begin
\subsection*{Strong feminine nouns}\BookHeading{nav-ch17-learner-0039}{subsubsection}{Strong feminine nouns}{0}
% mdparagraph:P144:end

% mdparagraph:P145:begin
The selected pattern has nominative singular without an added ending and genitive, dative and accusative singular -e. Plural nominative, genitive and accusative have -a; dative has -um. The endingless accusative of cwēn and the doubled consonants of the -nes and related stems require particular attention.
% mdparagraph:P145:end

\reserve{7\baselineskip}
% mdparagraph:P146:begin
\subsection*{hēanes}\BookHeading{nav-ch17-learner-0040}{subsubsection}{hēanes}{0}
% mdparagraph:P146:end

% teaching: NEW-U17-VOC-HEANES-HEIGHTS-U13-SOURCE
\par\begin{minipage}{\linewidth}\raggedright
% mdparagraph:P147:begin
\textbf{hēanes} — heights; highest places.
% mdparagraph:P147:end

% teaching: NEW-U17-VOC-HEANES-U13-SOURCE
% mdparagraph:P148:begin
\textbf{hēanes} — height; highest point. Supplied reference only.
% mdparagraph:P148:end

% teaching: NEW-U17-NOM-HEANES
% mdparagraph:P149:begin
Feminine; stem \textbf{hēaness-}. Use the plural for heights or the highest places. The singular column gives reference forms only; singular height is not required from memory here. The stem is hēaness-, with doubled s before every ending and with e retained.
% mdparagraph:P149:end

% table: NEW-U17-NOM-HEANES-DISPLAY
% mdtable:T31:begin
\par\begingroup
\fontsize{10.5}{13.2}\selectfont
\setlength{\LTleft}{0pt}\setlength{\LTright}{\fill}
\noindent
\begin{tabular}{@{}>{\raggedright\arraybackslash}p{39.46mm}>{\raggedright\arraybackslash}p{39.46mm}>{\raggedright\arraybackslash}p{39.46mm}@{}}
\toprule
\textbf{Case} & \textbf{Singular} & \textbf{Plural} \\
\midrule
Nominative & hēanes & hēanessa \\
Genitive & hēanesse & hēanessa \\
Dative & hēanesse & hēanessum \\
Accusative & hēanesse & hēanessa \\
\bottomrule
\end{tabular}
\par\endgroup

% mdtable:T31:end
\end{minipage}\par
\reserve{7\baselineskip}
% mdparagraph:P150:begin
\subsection*{ȳst}\BookHeading{nav-ch17-learner-0041}{subsubsection}{ȳst}{0}
% mdparagraph:P150:end

% teaching: NEW-U17-VOC-YST-U17-SENSE
\par\begin{minipage}{\linewidth}\raggedright
% mdparagraph:P151:begin
\textbf{ȳst} — storm.
% mdparagraph:P151:end

% teaching: NEW-U17-NOM-YST
% mdparagraph:P152:begin
Feminine; stem \textbf{ȳst-}. Add the endings of the stated pattern to this stem, preserving its vowel length.
% mdparagraph:P152:end

% table: NEW-U17-NOM-YST-DISPLAY
% mdtable:T32:begin
\par\begingroup
\fontsize{10.5}{13.2}\selectfont
\setlength{\LTleft}{0pt}\setlength{\LTright}{\fill}
\noindent
\begin{tabular}{@{}>{\raggedright\arraybackslash}p{39.46mm}>{\raggedright\arraybackslash}p{39.46mm}>{\raggedright\arraybackslash}p{39.46mm}@{}}
\toprule
\textbf{Case} & \textbf{Singular} & \textbf{Plural} \\
\midrule
Nominative & ȳst & ȳsta \\
Genitive & ȳste & ȳsta \\
Dative & ȳste & ȳstum \\
Accusative & ȳste & ȳsta \\
\bottomrule
\end{tabular}
\par\endgroup

% mdtable:T32:end
\end{minipage}\par
\reserve{7\baselineskip}
% mdparagraph:P153:begin
\subsection*{cwēn}\BookHeading{nav-ch17-learner-0042}{subsubsection}{cwēn}{0}
% mdparagraph:P153:end

% teaching: NEW-U17-VOC-CWEN-U17-QUEEN
% teaching: NEW-U17-VOC-CWEN-U17-PRINCESS
\par\begin{minipage}{\linewidth}\raggedright
% mdparagraph:P154:begin
\textbf{cwēn} — queen; princess.
% mdparagraph:P154:end

% teaching: NEW-U17-NOM-CWEN
% mdparagraph:P155:begin
Feminine; stem \textbf{cwēn-}. The selected accusative singular is endingless cwēn; the genitive and dative have cwēne. The ordinary-class accusative cwēne is also a defensible variant, but the contrast exercise explicitly asks for cwēn. Queen and princess are separately assigned contextual meanings.
% mdparagraph:P155:end

% table: NEW-U17-NOM-CWEN-DISPLAY
% mdtable:T33:begin
\par\begingroup
\fontsize{10.5}{13.2}\selectfont
\setlength{\LTleft}{0pt}\setlength{\LTright}{\fill}
\noindent
\begin{tabular}{@{}>{\raggedright\arraybackslash}p{39.46mm}>{\raggedright\arraybackslash}p{39.46mm}>{\raggedright\arraybackslash}p{39.46mm}@{}}
\toprule
\textbf{Case} & \textbf{Singular} & \textbf{Plural} \\
\midrule
Nominative & cwēn & cwēna \\
Genitive & cwēne & cwēna \\
Dative & cwēne & cwēnum \\
Accusative & cwēn & cwēna \\
\bottomrule
\end{tabular}
\par\endgroup

% mdtable:T33:end
\end{minipage}\par
\reserve{7\baselineskip}
% mdparagraph:P156:begin
\subsection*{ðēnung}\BookHeading{nav-ch17-learner-0043}{subsubsection}{ðēnung}{0}
% mdparagraph:P156:end

% teaching: NEW-U17-VOC-THENUNG-U16-READING
\par\begin{minipage}{\linewidth}\raggedright
% mdparagraph:P157:begin
\textbf{ðēnung} — service; attendance; service performed.
% mdparagraph:P157:end

% teaching: NEW-U17-NOM-THENUNG
% mdparagraph:P158:begin
Feminine; stem \textbf{ðēnung-}. The selected singular oblique ending is -e. Nouns in -ung can also have -a in the dative, and sometimes in the genitive and accusative; retain such a form when it is given in a source.
% mdparagraph:P158:end

% table: NEW-U17-NOM-THENUNG-DISPLAY
% mdtable:T34:begin
\par\begingroup
\fontsize{10.5}{13.2}\selectfont
\setlength{\LTleft}{0pt}\setlength{\LTright}{\fill}
\noindent
\begin{tabular}{@{}>{\raggedright\arraybackslash}p{39.46mm}>{\raggedright\arraybackslash}p{39.46mm}>{\raggedright\arraybackslash}p{39.46mm}@{}}
\toprule
\textbf{Case} & \textbf{Singular} & \textbf{Plural} \\
\midrule
Nominative & ðēnung & ðēnunga \\
Genitive & ðēnunge & ðēnunga \\
Dative & ðēnunge & ðēnungum \\
Accusative & ðēnunge & ðēnunga \\
\bottomrule
\end{tabular}
\par\endgroup

% mdtable:T34:end
\end{minipage}\par
\reserve{7\baselineskip}
% mdparagraph:P159:begin
\subsection*{getimbrung}\BookHeading{nav-ch17-learner-0044}{subsubsection}{getimbrung}{0}
% mdparagraph:P159:end

% teaching: NEW-U17-VOC-GETIMBRUNG-U17-SENSE
\par\begin{minipage}{\linewidth}\raggedright
% mdparagraph:P160:begin
\textbf{getimbrung} — building; structure built.
% mdparagraph:P160:end

% teaching: NEW-U17-NOM-GETIMBRUNG
% mdparagraph:P161:begin
Feminine; stem \textbf{getimbrung-}. The noun is feminine. A pronoun resuming the building must agree with that grammatical gender.
% mdparagraph:P161:end

% table: NEW-U17-NOM-GETIMBRUNG-DISPLAY
% mdtable:T35:begin
\par\begingroup
\fontsize{10.5}{13.2}\selectfont
\setlength{\LTleft}{0pt}\setlength{\LTright}{\fill}
\noindent
\begin{tabular}{@{}>{\raggedright\arraybackslash}p{35.87mm}>{\raggedright\arraybackslash}p{39.46mm}>{\raggedright\arraybackslash}p{43.05mm}@{}}
\toprule
\textbf{Case} & \textbf{Singular} & \textbf{Plural} \\
\midrule
Nominative & getimbrung & getimbrunga \\
Genitive & getimbrunge & getimbrunga \\
Dative & getimbrunge & getimbrungum \\
Accusative & getimbrunge & getimbrunga \\
\bottomrule
\end{tabular}
\par\endgroup

% mdtable:T35:end
\end{minipage}\par
\reserve{7\baselineskip}
% mdparagraph:P162:begin
\subsection*{ēaðmōdnes}\BookHeading{nav-ch17-learner-0045}{subsubsection}{ēaðmōdnes}{0}
% mdparagraph:P162:end

% teaching: NEW-U17-VOC-EADMODNES-U17-SENSE
\par\begin{minipage}{\linewidth}\raggedright
% mdparagraph:P163:begin
\textbf{ēaðmōdnes} — humility.
% mdparagraph:P163:end

% teaching: NEW-U17-NOM-EADMODNES
% mdparagraph:P164:begin
Feminine; stem \textbf{ēaðmōdness-}. Use ēaðmōdness- before the endings, with doubled s and retained e. Only abstract singular humility is assigned.
% mdparagraph:P164:end

% table: NEW-U17-NOM-EADMODNES-DISPLAY
% mdtable:T36:begin
\par\begingroup
\fontsize{10.5}{13.2}\selectfont
\setlength{\LTleft}{0pt}\setlength{\LTright}{\fill}
\noindent
\begin{tabular}{@{}>{\raggedright\arraybackslash}p{57.71mm}>{\raggedright\arraybackslash}p{63.48mm}@{}}
\toprule
\textbf{Case} & \textbf{Singular} \\
\midrule
Nominative & ēaðmōdnes \\
Genitive & ēaðmōdnesse \\
Dative & ēaðmōdnesse \\
Accusative & ēaðmōdnesse \\
\bottomrule
\end{tabular}
\par\endgroup

% mdtable:T36:end
\end{minipage}\par
\reserve{7\baselineskip}
% mdparagraph:P165:begin
\subsection*{ðīnen}\BookHeading{nav-ch17-learner-0046}{subsubsection}{ðīnen}{0}
% mdparagraph:P165:end

% teaching: NEW-U17-VOC-THINEN-U17-SENSE
\par\begin{minipage}{\linewidth}\raggedright
% mdparagraph:P166:begin
\textbf{ðīnen} — handmaiden.
% mdparagraph:P166:end

% teaching: NEW-U17-NOM-THINEN
% mdparagraph:P167:begin
Feminine; stem \textbf{ðīnenn-}. The oblique stem is ðīnenn-, with doubled n and retained e: ðīnenne. The nominative singular conceals the doubling.
% mdparagraph:P167:end

% table: NEW-U17-NOM-THINEN-DISPLAY
% mdtable:T37:begin
\par\begingroup
\fontsize{10.5}{13.2}\selectfont
\setlength{\LTleft}{0pt}\setlength{\LTright}{\fill}
\noindent
\begin{tabular}{@{}>{\raggedright\arraybackslash}p{39.46mm}>{\raggedright\arraybackslash}p{39.46mm}>{\raggedright\arraybackslash}p{39.46mm}@{}}
\toprule
\textbf{Case} & \textbf{Singular} & \textbf{Plural} \\
\midrule
Nominative & ðīnen & ðīnenna \\
Genitive & ðīnenne & ðīnenna \\
Dative & ðīnenne & ðīnennum \\
Accusative & ðīnenne & ðīnenna \\
\bottomrule
\end{tabular}
\par\endgroup

% mdtable:T37:end
\end{minipage}\par
\reserve{7\baselineskip}
% mdparagraph:P168:begin
\subsection*{cnēoris}\BookHeading{nav-ch17-learner-0047}{subsubsection}{cnēoris}{0}
% mdparagraph:P168:end

% teaching: NEW-U17-VOC-CNEORIS-U17-SENSE
\par\begin{minipage}{\linewidth}\raggedright
% mdparagraph:P169:begin
\textbf{cnēoris} — generation.
% mdparagraph:P169:end

% teaching: NEW-U17-NOM-CNEORIS
% mdparagraph:P170:begin
Feminine; stem \textbf{cnēoriss-}. Use cnēoriss- before endings: cnēorisse, cnēorissa and cnēorissum. The assigned meaning is generation.
% mdparagraph:P170:end

% table: NEW-U17-NOM-CNEORIS-DISPLAY
% mdtable:T38:begin
\par\begingroup
\fontsize{10.5}{13.2}\selectfont
\setlength{\LTleft}{0pt}\setlength{\LTright}{\fill}
\noindent
\begin{tabular}{@{}>{\raggedright\arraybackslash}p{39.46mm}>{\raggedright\arraybackslash}p{39.46mm}>{\raggedright\arraybackslash}p{39.46mm}@{}}
\toprule
\textbf{Case} & \textbf{Singular} & \textbf{Plural} \\
\midrule
Nominative & cnēoris & cnēorissa \\
Genitive & cnēorisse & cnēorissa \\
Dative & cnēorisse & cnēorissum \\
Accusative & cnēorisse & cnēorissa \\
\bottomrule
\end{tabular}
\par\endgroup

% mdtable:T38:end
\end{minipage}\par
\reserve{7\baselineskip}
% mdparagraph:P171:begin
\subsection*{mǣgð}\BookHeading{nav-ch17-learner-0048}{subsubsection}{mǣgð}{0}
% mdparagraph:P171:end

% teaching: NEW-U17-VOC-MAEGTH-TRIBE-U17-SENSE
\par\begin{minipage}{\linewidth}\raggedright
% mdparagraph:P172:begin
\textbf{mǣgð} — tribe; nation.
% mdparagraph:P172:end

% teaching: NEW-U17-NOM-MAEGTH-TRIBE
% mdparagraph:P173:begin
Feminine; stem \textbf{mǣgð-}. Long ǣ distinguishes mǣgð, tribe or nation, from short-vowel mægð, maiden. Keep both the length and this noun’s feminine pattern.
% mdparagraph:P173:end

% table: NEW-U17-NOM-MAEGTH-TRIBE-DISPLAY
% mdtable:T39:begin
\par\begingroup
\fontsize{10.5}{13.2}\selectfont
\setlength{\LTleft}{0pt}\setlength{\LTright}{\fill}
\noindent
\begin{tabular}{@{}>{\raggedright\arraybackslash}p{39.46mm}>{\raggedright\arraybackslash}p{39.46mm}>{\raggedright\arraybackslash}p{39.46mm}@{}}
\toprule
\textbf{Case} & \textbf{Singular} & \textbf{Plural} \\
\midrule
Nominative & mǣgð & mǣgða \\
Genitive & mǣgðe & mǣgða \\
Dative & mǣgðe & mǣgðum \\
Accusative & mǣgðe & mǣgða \\
\bottomrule
\end{tabular}
\par\endgroup

% mdtable:T39:end
\end{minipage}\par
\reserve{7\baselineskip}
% mdparagraph:P174:begin
\subsection*{blētsung}\BookHeading{nav-ch17-learner-0049}{subsubsection}{blētsung}{0}
% mdparagraph:P174:end

% teaching: NEW-U17-VOC-BLETSUNG-U17-SENSE
\par\begin{minipage}{\linewidth}\raggedright
% mdparagraph:P175:begin
\textbf{blētsung} — blessing.
% mdparagraph:P175:end

% teaching: NEW-U17-NOM-BLETSUNG
% mdparagraph:P176:begin
Feminine; stem \textbf{blētsung-}. Add the endings of the stated pattern to this stem, preserving its vowel length.
% mdparagraph:P176:end

% mdparagraph:P177:begin
This table gives the complete singular for the stated abstract or mass meaning. The familiar noun class is reused; plural meanings and forms are not assigned here.
% mdparagraph:P177:end

% table: NEW-U17-NOM-BLETSUNG-DISPLAY
% mdtable:T40:begin
\par\begingroup
\fontsize{10.5}{13.2}\selectfont
\setlength{\LTleft}{0pt}\setlength{\LTright}{\fill}
\noindent
\begin{tabular}{@{}>{\raggedright\arraybackslash}p{60.59mm}>{\raggedright\arraybackslash}p{60.59mm}@{}}
\toprule
\textbf{Case} & \textbf{Singular} \\
\midrule
Nominative & blētsung \\
Genitive & blētsunge \\
Dative & blētsunge \\
Accusative & blētsunge \\
\bottomrule
\end{tabular}
\par\endgroup

% mdtable:T40:end
\end{minipage}\par
\reserve{7\baselineskip}
% mdparagraph:P178:begin
\subsection*{beorhtnes}\BookHeading{nav-ch17-learner-0050}{subsubsection}{beorhtnes}{0}
% mdparagraph:P178:end

% teaching: NEW-U17-VOC-BEORHTNES-U17-SENSE
\par\begin{minipage}{\linewidth}\raggedright
% mdparagraph:P179:begin
\textbf{beorhtnes} — brightness.
% mdparagraph:P179:end

% teaching: NEW-U17-NOM-BEORHTNES
% mdparagraph:P180:begin
Feminine; stem \textbf{beorhtness-}. Add the endings of the stated pattern to this stem, preserving its vowel length.
% mdparagraph:P180:end

% mdparagraph:P181:begin
This table gives the complete singular for the stated abstract or mass meaning. The familiar noun class is reused; plural meanings and forms are not assigned here.
% mdparagraph:P181:end

% table: NEW-U17-NOM-BEORHTNES-DISPLAY
% mdtable:T41:begin
\par\begingroup
\fontsize{10.5}{13.2}\selectfont
\setlength{\LTleft}{0pt}\setlength{\LTright}{\fill}
\noindent
\begin{tabular}{@{}>{\raggedright\arraybackslash}p{57.71mm}>{\raggedright\arraybackslash}p{63.48mm}@{}}
\toprule
\textbf{Case} & \textbf{Singular} \\
\midrule
Nominative & beorhtnes \\
Genitive & beorhtnesse \\
Dative & beorhtnesse \\
Accusative & beorhtnesse \\
\bottomrule
\end{tabular}
\par\endgroup

% mdtable:T41:end
\end{minipage}\par
\reserve{7\baselineskip}
% mdparagraph:P182:begin
\subsection*{ðancung}\BookHeading{nav-ch17-learner-0051}{subsubsection}{ðancung}{0}
% mdparagraph:P182:end

% teaching: NEW-U17-VOC-THANCUNG-U17-SENSE
\par\begin{minipage}{\linewidth}\raggedright
% mdparagraph:P183:begin
\textbf{ðancung} — thanksgiving.
% mdparagraph:P183:end

% teaching: NEW-U17-NOM-THANCUNG
% mdparagraph:P184:begin
Feminine; stem \textbf{ðancung-}. Add the endings of the stated pattern to this stem, preserving its vowel length.
% mdparagraph:P184:end

% mdparagraph:P185:begin
This table gives the complete singular for the stated abstract or mass meaning. The familiar noun class is reused; plural meanings and forms are not assigned here.
% mdparagraph:P185:end

% table: NEW-U17-NOM-THANCUNG-DISPLAY
% mdtable:T42:begin
\par\begingroup
\fontsize{10.5}{13.2}\selectfont
\setlength{\LTleft}{0pt}\setlength{\LTright}{\fill}
\noindent
\begin{tabular}{@{}>{\raggedright\arraybackslash}p{60.59mm}>{\raggedright\arraybackslash}p{60.59mm}@{}}
\toprule
\textbf{Case} & \textbf{Singular} \\
\midrule
Nominative & ðancung \\
Genitive & ðancunge \\
Dative & ðancunge \\
Accusative & ðancunge \\
\bottomrule
\end{tabular}
\par\endgroup

% mdtable:T42:end
\end{minipage}\par
\reserve{7\baselineskip}
% mdparagraph:P186:begin
\subsection*{Weak nouns}\BookHeading{nav-ch17-learner-0052}{subsubsection}{Weak nouns}{0}
% mdparagraph:P186:end

% mdparagraph:P187:begin
Weak masculines normally have nominative singular -a, other singular cases -an, plural nominative/\allowbreak{}accusative -an, genitive -ena and dative -um. Weak feminines have nominative singular -e and otherwise the same endings. The weak neuter ēare keeps -e also in the accusative singular and has a contracted genitive plural.
% mdparagraph:P187:end

\reserve{7\baselineskip}
% mdparagraph:P188:begin
\subsection*{draca}\BookHeading{nav-ch17-learner-0053}{subsubsection}{draca}{0}
% mdparagraph:P188:end

% teaching: NEW-U17-VOC-DRACA-U17-SENSE
\par\begin{minipage}{\linewidth}\raggedright
% mdparagraph:P189:begin
\textbf{draca} — dragon.
% mdparagraph:P189:end

% teaching: NEW-U17-NOM-DRACA
% mdparagraph:P190:begin
Masculine; stem \textbf{drac-}. Add the endings of the stated pattern to this stem, preserving its vowel length.
% mdparagraph:P190:end

% table: NEW-U17-NOM-DRACA-DISPLAY
% mdtable:T43:begin
\par\begingroup
\fontsize{10.5}{13.2}\selectfont
\setlength{\LTleft}{0pt}\setlength{\LTright}{\fill}
\noindent
\begin{tabular}{@{}>{\raggedright\arraybackslash}p{39.46mm}>{\raggedright\arraybackslash}p{39.46mm}>{\raggedright\arraybackslash}p{39.46mm}@{}}
\toprule
\textbf{Case} & \textbf{Singular} & \textbf{Plural} \\
\midrule
Nominative & draca & dracan \\
Genitive & dracan & dracena \\
Dative & dracan & dracum \\
Accusative & dracan & dracan \\
\bottomrule
\end{tabular}
\par\endgroup

% mdtable:T43:end
\end{minipage}\par
\reserve{7\baselineskip}
% mdparagraph:P191:begin
\subsection*{nædre}\BookHeading{nav-ch17-learner-0054}{subsubsection}{nædre}{0}
% mdparagraph:P191:end

% teaching: NEW-U17-VOC-NAEDRE-U17-SENSE
\par\begin{minipage}{\linewidth}\raggedright
% mdparagraph:P192:begin
\textbf{nædre} — serpent.
% mdparagraph:P192:end

% teaching: NEW-U17-NOM-NAEDRE
% mdparagraph:P193:begin
Feminine; stem \textbf{nædr-}. The selected spelling is nædre, with short æ and one d. The source also gives næddre. Add the weak feminine endings to nædr-.
% mdparagraph:P193:end

% table: NEW-U17-NOM-NAEDRE-DISPLAY
% mdtable:T44:begin
\par\begingroup
\fontsize{10.5}{13.2}\selectfont
\setlength{\LTleft}{0pt}\setlength{\LTright}{\fill}
\noindent
\begin{tabular}{@{}>{\raggedright\arraybackslash}p{39.46mm}>{\raggedright\arraybackslash}p{39.46mm}>{\raggedright\arraybackslash}p{39.46mm}@{}}
\toprule
\textbf{Case} & \textbf{Singular} & \textbf{Plural} \\
\midrule
Nominative & nædre & nædran \\
Genitive & nædran & nædrena \\
Dative & nædran & nædrum \\
Accusative & nædran & nædran \\
\bottomrule
\end{tabular}
\par\endgroup

% mdtable:T44:end
\end{minipage}\par
\reserve{7\baselineskip}
% mdparagraph:P194:begin
\subsection*{dēma}\BookHeading{nav-ch17-learner-0055}{subsubsection}{dēma}{0}
% mdparagraph:P194:end

% teaching: NEW-U17-VOC-DEMA-U17-SENSE
\par\begin{minipage}{\linewidth}\raggedright
% mdparagraph:P195:begin
\textbf{dēma} — judge.
% mdparagraph:P195:end

% teaching: NEW-U17-NOM-DEMA
% mdparagraph:P196:begin
Masculine; stem \textbf{dēm-}. Add the endings of the stated pattern to this stem, preserving its vowel length.
% mdparagraph:P196:end

% table: NEW-U17-NOM-DEMA-DISPLAY
% mdtable:T45:begin
\par\begingroup
\fontsize{10.5}{13.2}\selectfont
\setlength{\LTleft}{0pt}\setlength{\LTright}{\fill}
\noindent
\begin{tabular}{@{}>{\raggedright\arraybackslash}p{39.46mm}>{\raggedright\arraybackslash}p{39.46mm}>{\raggedright\arraybackslash}p{39.46mm}@{}}
\toprule
\textbf{Case} & \textbf{Singular} & \textbf{Plural} \\
\midrule
Nominative & dēma & dēman \\
Genitive & dēman & dēmena \\
Dative & dēman & dēmum \\
Accusative & dēman & dēman \\
\bottomrule
\end{tabular}
\par\endgroup

% mdtable:T45:end
\end{minipage}\par
\reserve{7\baselineskip}
% mdparagraph:P197:begin
\subsection*{hālga}\BookHeading{nav-ch17-learner-0056}{subsubsection}{hālga}{0}
% mdparagraph:P197:end

% teaching: NEW-U17-VOC-HALGA-U17-SENSE
\par\begin{minipage}{\linewidth}\raggedright
% mdparagraph:P198:begin
\textbf{hālga} — saint; holy one.
% mdparagraph:P198:end

% teaching: NEW-U17-NOM-HALGA
% mdparagraph:P199:begin
Masculine; stem \textbf{hālg-}. This is a weak masculine noun, saint or holy one. Its genitive plural is hālgena. Keep it distinct from the adjective hālig and its genitive plural hālgra.
% mdparagraph:P199:end

% table: NEW-U17-NOM-HALGA-DISPLAY
% mdtable:T46:begin
\par\begingroup
\fontsize{10.5}{13.2}\selectfont
\setlength{\LTleft}{0pt}\setlength{\LTright}{\fill}
\noindent
\begin{tabular}{@{}>{\raggedright\arraybackslash}p{39.46mm}>{\raggedright\arraybackslash}p{39.46mm}>{\raggedright\arraybackslash}p{39.46mm}@{}}
\toprule
\textbf{Case} & \textbf{Singular} & \textbf{Plural} \\
\midrule
Nominative & hālga & hālgan \\
Genitive & hālgan & hālgena \\
Dative & hālgan & hālgum \\
Accusative & hālgan & hālgan \\
\bottomrule
\end{tabular}
\par\endgroup

% mdtable:T46:end
\end{minipage}\par
\reserve{7\baselineskip}
% mdparagraph:P200:begin
\subsection*{hlīsa}\BookHeading{nav-ch17-learner-0057}{subsubsection}{hlīsa}{0}
% mdparagraph:P200:end

% teaching: NEW-U17-VOC-HLISA-U17-FAME
% teaching: NEW-U17-VOC-HLISA-U17-REPORT
\par\begin{minipage}{\linewidth}\raggedright
% mdparagraph:P201:begin
\textbf{hlīsa} — fame; report; account heard.
% mdparagraph:P201:end

% teaching: NEW-U17-NOM-HLISA
% mdparagraph:P202:begin
Masculine; stem \textbf{hlīs-}. Fame is ordinarily singular in this passage. A report or account can be treated as an individual item. The complete table serves those stated meanings; context determines which is intended.
% mdparagraph:P202:end

% table: NEW-U17-NOM-HLISA-DISPLAY
% mdtable:T47:begin
\par\begingroup
\fontsize{10.5}{13.2}\selectfont
\setlength{\LTleft}{0pt}\setlength{\LTright}{\fill}
\noindent
\begin{tabular}{@{}>{\raggedright\arraybackslash}p{39.46mm}>{\raggedright\arraybackslash}p{39.46mm}>{\raggedright\arraybackslash}p{39.46mm}@{}}
\toprule
\textbf{Case} & \textbf{Singular} & \textbf{Plural} \\
\midrule
Nominative & hlīsa & hlīsan \\
Genitive & hlīsan & hlīsena \\
Dative & hlīsan & hlīsum \\
Accusative & hlīsan & hlīsan \\
\bottomrule
\end{tabular}
\par\endgroup

% mdtable:T47:end
\end{minipage}\par
\reserve{7\baselineskip}
% mdparagraph:P203:begin
\subsection*{ðēowa}\BookHeading{nav-ch17-learner-0058}{subsubsection}{ðēowa}{0}
% mdparagraph:P203:end

% teaching: NEW-U17-VOC-THEOWA-U17-SENSE
\par\begin{minipage}{\linewidth}\raggedright
% mdparagraph:P204:begin
\textbf{ðēowa} — servant.
% mdparagraph:P204:end

% teaching: NEW-U17-NOM-THEOWA
% mdparagraph:P205:begin
Masculine; stem \textbf{ðēow-}. This is weak ðēowa, with nominative/\allowbreak{}accusative plural ðēowan. Do not substitute the plural of the separate strong form ðēow.
% mdparagraph:P205:end

% table: NEW-U17-NOM-THEOWA-DISPLAY
% mdtable:T48:begin
\par\begingroup
\fontsize{10.5}{13.2}\selectfont
\setlength{\LTleft}{0pt}\setlength{\LTright}{\fill}
\noindent
\begin{tabular}{@{}>{\raggedright\arraybackslash}p{39.46mm}>{\raggedright\arraybackslash}p{39.46mm}>{\raggedright\arraybackslash}p{39.46mm}@{}}
\toprule
\textbf{Case} & \textbf{Singular} & \textbf{Plural} \\
\midrule
Nominative & ðēowa & ðēowan \\
Genitive & ðēowan & ðēowena \\
Dative & ðēowan & ðēowum \\
Accusative & ðēowan & ðēowan \\
\bottomrule
\end{tabular}
\par\endgroup

% mdtable:T48:end
\end{minipage}\par
\reserve{7\baselineskip}
% mdparagraph:P206:begin
\subsection*{ēare}\BookHeading{nav-ch17-learner-0059}{subsubsection}{ēare}{0}
% mdparagraph:P206:end

% teaching: NEW-U17-VOC-EARE-U17-SENSE
\par\begin{minipage}{\linewidth}\raggedright
% mdparagraph:P207:begin
\textbf{ēare} — ear.
% mdparagraph:P207:end

% teaching: NEW-U17-NOM-EARE
% mdparagraph:P208:begin
Neuter; stem \textbf{ēar-}. This weak neuter has ēare in both nominative and accusative singular. The other singular cases end in -an. Its genitive plural contracts to ēarna; the dative plural is ēarum.
% mdparagraph:P208:end

% table: NEW-U17-NOM-EARE-DISPLAY
% mdtable:T49:begin
\par\begingroup
\fontsize{10.5}{13.2}\selectfont
\setlength{\LTleft}{0pt}\setlength{\LTright}{\fill}
\noindent
\begin{tabular}{@{}>{\raggedright\arraybackslash}p{39.46mm}>{\raggedright\arraybackslash}p{39.46mm}>{\raggedright\arraybackslash}p{39.46mm}@{}}
\toprule
\textbf{Case} & \textbf{Singular} & \textbf{Plural} \\
\midrule
Nominative & ēare & ēaran \\
Genitive & ēaran & ēarna \\
Dative & ēaran & ēarum \\
Accusative & ēare & ēaran \\
\bottomrule
\end{tabular}
\par\endgroup

% mdtable:T49:end
\end{minipage}\par
\reserve{7\baselineskip}
% mdparagraph:P209:begin
\subsection*{tabele}\BookHeading{nav-ch17-learner-0060}{subsubsection}{tabele}{0}
% mdparagraph:P209:end

% teaching: NEW-U17-VOC-TABELE-U17-SENSE
\par\begin{minipage}{\linewidth}\raggedright
% mdparagraph:P210:begin
\textbf{tabele} — tablet; table of the law.
% mdparagraph:P210:end

% teaching: NEW-U17-NOM-TABELE
% mdparagraph:P211:begin
Feminine; stem \textbf{tabel-}. The weak feminine stem is tabel-. Retain e in the dative plural tabelum.
% mdparagraph:P211:end

% table: NEW-U17-NOM-TABELE-DISPLAY
% mdtable:T50:begin
\par\begingroup
\fontsize{10.5}{13.2}\selectfont
\setlength{\LTleft}{0pt}\setlength{\LTright}{\fill}
\noindent
\begin{tabular}{@{}>{\raggedright\arraybackslash}p{39.46mm}>{\raggedright\arraybackslash}p{39.46mm}>{\raggedright\arraybackslash}p{39.46mm}@{}}
\toprule
\textbf{Case} & \textbf{Singular} & \textbf{Plural} \\
\midrule
Nominative & tabele & tabelan \\
Genitive & tabelan & tabelena \\
Dative & tabelan & tabelum \\
Accusative & tabelan & tabelan \\
\bottomrule
\end{tabular}
\par\endgroup

% mdtable:T50:end
\end{minipage}\par
\reserve{7\baselineskip}
% mdparagraph:P212:begin
\subsection*{hearpe}\BookHeading{nav-ch17-learner-0061}{subsubsection}{hearpe}{0}
% mdparagraph:P212:end

% teaching: NEW-U17-VOC-HEARPE-U17-SENSE
\par\begin{minipage}{\linewidth}\raggedright
% mdparagraph:P213:begin
\textbf{hearpe} — harp.
% mdparagraph:P213:end

% teaching: NEW-U17-NOM-HEARPE
% mdparagraph:P214:begin
Feminine; stem \textbf{hearp-}. Add the endings of the stated pattern to this stem, preserving its vowel length.
% mdparagraph:P214:end

% table: NEW-U17-NOM-HEARPE-DISPLAY
% mdtable:T51:begin
\par\begingroup
\fontsize{10.5}{13.2}\selectfont
\setlength{\LTleft}{0pt}\setlength{\LTright}{\fill}
\noindent
\begin{tabular}{@{}>{\raggedright\arraybackslash}p{39.46mm}>{\raggedright\arraybackslash}p{39.46mm}>{\raggedright\arraybackslash}p{39.46mm}@{}}
\toprule
\textbf{Case} & \textbf{Singular} & \textbf{Plural} \\
\midrule
Nominative & hearpe & hearpan \\
Genitive & hearpan & hearpena \\
Dative & hearpan & hearpum \\
Accusative & hearpan & hearpan \\
\bottomrule
\end{tabular}
\par\endgroup

% mdtable:T51:end
\end{minipage}\par
\reserve{7\baselineskip}
% mdparagraph:P215:begin
\subsection*{blōstma}\BookHeading{nav-ch17-learner-0062}{subsubsection}{blōstma}{0}
% mdparagraph:P215:end

% teaching: NEW-U17-VOC-BLOSTMA-U17-SENSE
\par\begin{minipage}{\linewidth}\raggedright
% mdparagraph:P216:begin
\textbf{blōstma} — blossom.
% mdparagraph:P216:end

% teaching: NEW-U17-NOM-BLOSTMA
% mdparagraph:P217:begin
Masculine; stem \textbf{blōstm-}. Add the endings of the stated pattern to this stem, preserving its vowel length.
% mdparagraph:P217:end

% table: NEW-U17-NOM-BLOSTMA-DISPLAY
% mdtable:T52:begin
\par\begingroup
\fontsize{10.5}{13.2}\selectfont
\setlength{\LTleft}{0pt}\setlength{\LTright}{\fill}
\noindent
\begin{tabular}{@{}>{\raggedright\arraybackslash}p{39.46mm}>{\raggedright\arraybackslash}p{39.46mm}>{\raggedright\arraybackslash}p{39.46mm}@{}}
\toprule
\textbf{Case} & \textbf{Singular} & \textbf{Plural} \\
\midrule
Nominative & blōstma & blōstman \\
Genitive & blōstman & blōstmena \\
Dative & blōstman & blōstmum \\
Accusative & blōstman & blōstman \\
\bottomrule
\end{tabular}
\par\endgroup

% mdtable:T52:end
\end{minipage}\par
\reserve{7\baselineskip}
% mdparagraph:P218:begin
\subsection*{biema}\BookHeading{nav-ch17-learner-0063}{subsubsection}{biema}{0}
% mdparagraph:P218:end

% teaching: NEW-U17-VOC-BIEMA-U17-SENSE
\par\begin{minipage}{\linewidth}\raggedright
% mdparagraph:P219:begin
\textbf{biema} — trumpet; clarion.
% mdparagraph:P219:end

% teaching: NEW-U17-NOM-BIEMA
% mdparagraph:P220:begin
Masculine; stem \textbf{biem-}. Retain the source spelling biema, without adding a length mark. It follows the weak masculine pattern.
% mdparagraph:P220:end

% table: NEW-U17-NOM-BIEMA-DISPLAY
% mdtable:T53:begin
\par\begingroup
\fontsize{10.5}{13.2}\selectfont
\setlength{\LTleft}{0pt}\setlength{\LTright}{\fill}
\noindent
\begin{tabular}{@{}>{\raggedright\arraybackslash}p{39.46mm}>{\raggedright\arraybackslash}p{39.46mm}>{\raggedright\arraybackslash}p{39.46mm}@{}}
\toprule
\textbf{Case} & \textbf{Singular} & \textbf{Plural} \\
\midrule
Nominative & biema & bieman \\
Genitive & bieman & biemena \\
Dative & bieman & biemum \\
Accusative & bieman & bieman \\
\bottomrule
\end{tabular}
\par\endgroup

% mdtable:T53:end
\end{minipage}\par
\reserve{7\baselineskip}
% mdparagraph:P221:begin
\subsection*{hlǣfdige}\BookHeading{nav-ch17-learner-0064}{subsubsection}{hlǣfdige}{0}
% mdparagraph:P221:end

% teaching: NEW-U17-VOC-HLAEFDIGE-U17-SENSE
\par\begin{minipage}{\linewidth}\raggedright
% mdparagraph:P222:begin
\textbf{hlǣfdige} — lady.
% mdparagraph:P222:end

% teaching: NEW-U17-NOM-HLAEFDIGE
% mdparagraph:P223:begin
Feminine; stem \textbf{hlǣfdig-}. Add the endings of the stated pattern to this stem, preserving its vowel length.
% mdparagraph:P223:end

% table: NEW-U17-NOM-HLAEFDIGE-DISPLAY
% mdtable:T54:begin
\par\begingroup
\fontsize{10.5}{13.2}\selectfont
\setlength{\LTleft}{0pt}\setlength{\LTright}{\fill}
\noindent
\begin{tabular}{@{}>{\raggedright\arraybackslash}p{39.46mm}>{\raggedright\arraybackslash}p{39.46mm}>{\raggedright\arraybackslash}p{39.46mm}@{}}
\toprule
\textbf{Case} & \textbf{Singular} & \textbf{Plural} \\
\midrule
Nominative & hlǣfdige & hlǣfdigan \\
Genitive & hlǣfdigan & hlǣfdigena \\
Dative & hlǣfdigan & hlǣfdigum \\
Accusative & hlǣfdigan & hlǣfdigan \\
\bottomrule
\end{tabular}
\par\endgroup

% mdtable:T54:end
\end{minipage}\par
\reserve{7\baselineskip}
% mdparagraph:P224:begin
\subsection*{Strong and weak adjectives}\BookHeading{nav-ch17-learner-0065}{subsubsection}{Strong and weak adjectives}{0}
% mdparagraph:P224:end

% mdparagraph:P225:begin
An adjective agrees with its noun in gender, number and case. Use the familiar strong and weak choice: an article or comparable definite determiner normally calls for weak inflection, while a predicative adjective is normally strong. Direct address also favours weak forms; the praise catalogue uses this pattern even while its global instruction removes the article. The selected personal uses are explained after the tables.
% mdparagraph:P225:end

% mdparagraph:P226:begin
Each panel has masculine, neuter and feminine columns; the singular panel precedes the plural panel. Strong instrumental singular forms in the masculine and neuter columns are for recognition; their display adds no new instrumental construction. A dash is an inapplicable display position. Weak genitive plural -ra is the selected form. These tables print one consistent set of source-licensed forms; a documented alternative is not automatically ungrammatical.
% mdparagraph:P226:end

\reserve{7\baselineskip}
% mdparagraph:P227:begin
\subsection*{wæstmbǣre}\BookHeading{nav-ch17-learner-0066}{subsubsection}{wæstmbǣre}{0}
% mdparagraph:P227:end

% teaching: NEW-U17-VOC-WAESTMBAERE-U17-SENSE
\par\begin{minipage}{\linewidth}\raggedright
% mdparagraph:P228:begin
\textbf{wæstmbǣre} — fruitful.
% mdparagraph:P228:end

% teaching: NEW-U17-ADJ-WAESTMBAERE
% mdparagraph:P229:begin
The headword ends in -e; replace that -e with an overt ending on \textbf{wæstmbǣr-}. The displayed strong forms select feminine nominative singular and neuter nominative/\allowbreak{}accusative plural -u, and feminine plural -a. The source also permits -e in those positions.
% mdparagraph:P229:end

% mdparagraph:P230:begin
\textbf{Strong.}
% mdparagraph:P230:end

% table: NEW-U17-ADJ-WAESTMBAERE-STRONG-DISPLAY
% mdtable:T55:begin
% course-table-style: 2026-10
% course-panel:T55-P01:begin
\par\begingroup
\fontsize{10.5}{13.2}\selectfont
\setlength{\LTleft}{0pt}\setlength{\LTright}{\fill}
\noindent
\begin{tabular}{@{}>{\raggedright\arraybackslash}p{22.0mm}>{\raggedright\arraybackslash}p{31.188mm}>{\raggedright\arraybackslash}p{31.188mm}>{\raggedright\arraybackslash}p{31.188mm}@{}}
\toprule
\textbf{Case} & \textbf{Masculine singular} & \textbf{Neuter singular} & \textbf{Feminine singular} \\
\midrule
Nominative & wæstmbǣre & wæstmbǣre & wæstmbǣru \\
Genitive & wæstmbǣres & wæstmbǣres & wæstmbǣrre \\
Dative & wæstmbǣrum & wæstmbǣrum & wæstmbǣrre \\
Accusative & wæstmbǣrne & wæstmbǣre & wæstmbǣre \\
Instrumental & wæstmbǣre & wæstmbǣre & — \\
\bottomrule
\end{tabular}
\par\endgroup
% course-panel:T55-P01:end
\end{minipage}\par
\par\addvspace{5pt}
% course-panel:T55-P02:begin
\par\begingroup
\fontsize{10.5}{13.2}\selectfont
\setlength{\LTleft}{0pt}\setlength{\LTright}{\fill}
\noindent
\begin{tabular}{@{}>{\raggedright\arraybackslash}p{22.0mm}>{\raggedright\arraybackslash}p{31.188mm}>{\raggedright\arraybackslash}p{31.188mm}>{\raggedright\arraybackslash}p{31.188mm}@{}}
\toprule
\textbf{Case} & \textbf{Masculine plural} & \textbf{Neuter plural} & \textbf{Feminine plural} \\
\midrule
Nominative & wæstmbǣre & wæstmbǣru & wæstmbǣra \\
Genitive & wæstmbǣrra & wæstmbǣrra & wæstmbǣrra \\
Dative & wæstmbǣrum & wæstmbǣrum & wæstmbǣrum \\
Accusative & wæstmbǣre & wæstmbǣru & wæstmbǣra \\
Instrumental & — & — & — \\
\bottomrule
\end{tabular}
\par\endgroup
% course-panel:T55-P02:end
% mdtable:T55:end

\reserve{18\baselineskip}
\reserve{16\baselineskip}
% mdparagraph:P231:begin
\textbf{Weak.}
% mdparagraph:P231:end

% table: NEW-U17-ADJ-WAESTMBAERE-WEAK-DISPLAY
% mdtable:T56:begin
% course-table-style: 2026-10
% course-panel:T56-P01:begin
\par\begingroup
\fontsize{10.5}{13.2}\selectfont
\setlength{\LTleft}{0pt}\setlength{\LTright}{\fill}
\noindent
\begin{tabular}{@{}>{\raggedright\arraybackslash}p{22.0mm}>{\raggedright\arraybackslash}p{31.188mm}>{\raggedright\arraybackslash}p{31.188mm}>{\raggedright\arraybackslash}p{31.188mm}@{}}
\toprule
\textbf{Case} & \textbf{Masculine singular} & \textbf{Neuter singular} & \textbf{Feminine singular} \\
\midrule
Nominative & wæstmbǣra & wæstmbǣre & wæstmbǣre \\
Genitive & wæstmbǣran & wæstmbǣran & wæstmbǣran \\
Dative & wæstmbǣran & wæstmbǣran & wæstmbǣran \\
Accusative & wæstmbǣran & wæstmbǣre & wæstmbǣran \\
\bottomrule
\end{tabular}
\par\endgroup
% course-panel:T56-P01:end
\par\addvspace{5pt}
% course-panel:T56-P02:begin
\par\begingroup
\fontsize{10.5}{13.2}\selectfont
\setlength{\LTleft}{0pt}\setlength{\LTright}{\fill}
\noindent
\begin{tabular}{@{}>{\raggedright\arraybackslash}p{22.0mm}>{\raggedright\arraybackslash}p{31.188mm}>{\raggedright\arraybackslash}p{31.188mm}>{\raggedright\arraybackslash}p{31.188mm}@{}}
\toprule
\textbf{Case} & \textbf{Masculine plural} & \textbf{Neuter plural} & \textbf{Feminine plural} \\
\midrule
Nominative & wæstmbǣran & wæstmbǣran & wæstmbǣran \\
Genitive & wæstmbǣrra & wæstmbǣrra & wæstmbǣrra \\
Dative & wæstmbǣrum & wæstmbǣrum & wæstmbǣrum \\
Accusative & wæstmbǣran & wæstmbǣran & wæstmbǣran \\
\bottomrule
\end{tabular}
\par\endgroup
% course-panel:T56-P02:end
% mdtable:T56:end

\reserve{7\baselineskip}
% mdparagraph:P232:begin
\subsection*{geong}\BookHeading{nav-ch17-learner-0067}{subsubsection}{geong}{0}
% mdparagraph:P232:end

% teaching: NEW-U17-VOC-GEONG-YOUNG-SOURCE
\par\begin{minipage}{\linewidth}\raggedright
% mdparagraph:P233:begin
\textbf{geong} — young; substantivally young people.
% mdparagraph:P233:end

% teaching: NEW-U17-ADJ-GEONG
% mdparagraph:P234:begin
The stem \textbf{geong-} is heavy. Strong feminine nominative singular and neuter nominative/\allowbreak{}accusative plural therefore have no added -u.
% mdparagraph:P234:end

% mdparagraph:P235:begin
\textbf{Strong.}
% mdparagraph:P235:end

% table: NEW-U17-ADJ-GEONG-STRONG-DISPLAY
% mdtable:T57:begin
% course-table-style: 2026-10
% course-panel:T57-P01:begin
\par\begingroup
\fontsize{10.5}{13.2}\selectfont
\setlength{\LTleft}{0pt}\setlength{\LTright}{\fill}
\noindent
\begin{tabular}{@{}>{\raggedright\arraybackslash}p{22.0mm}>{\raggedright\arraybackslash}p{31.188mm}>{\raggedright\arraybackslash}p{31.188mm}>{\raggedright\arraybackslash}p{31.188mm}@{}}
\toprule
\textbf{Case} & \textbf{Masculine singular} & \textbf{Neuter singular} & \textbf{Feminine singular} \\
\midrule
Nominative & geong & geong & geong \\
Genitive & geonges & geonges & geongre \\
Dative & geongum & geongum & geongre \\
Accusative & geongne & geong & geonge \\
Instrumental & geonge & geonge & — \\
\bottomrule
\end{tabular}
\par\endgroup
% course-panel:T57-P01:end
\end{minipage}\par
\par\addvspace{5pt}
% course-panel:T57-P02:begin
\par\begingroup
\fontsize{10.5}{13.2}\selectfont
\setlength{\LTleft}{0pt}\setlength{\LTright}{\fill}
\noindent
\begin{tabular}{@{}>{\raggedright\arraybackslash}p{22.0mm}>{\raggedright\arraybackslash}p{31.188mm}>{\raggedright\arraybackslash}p{31.188mm}>{\raggedright\arraybackslash}p{31.188mm}@{}}
\toprule
\textbf{Case} & \textbf{Masculine plural} & \textbf{Neuter plural} & \textbf{Feminine plural} \\
\midrule
Nominative & geonge & geong & geonga \\
Genitive & geongra & geongra & geongra \\
Dative & geongum & geongum & geongum \\
Accusative & geonge & geong & geonga \\
Instrumental & — & — & — \\
\bottomrule
\end{tabular}
\par\endgroup
% course-panel:T57-P02:end
% mdtable:T57:end

\reserve{18\baselineskip}
\reserve{16\baselineskip}
% mdparagraph:P236:begin
\textbf{Weak.}
% mdparagraph:P236:end

% table: NEW-U17-ADJ-GEONG-WEAK-DISPLAY
% mdtable:T58:begin
% course-table-style: 2026-10
% course-panel:T58-P01:begin
\par\begingroup
\fontsize{10.5}{13.2}\selectfont
\setlength{\LTleft}{0pt}\setlength{\LTright}{\fill}
\noindent
\begin{tabular}{@{}>{\raggedright\arraybackslash}p{22.0mm}>{\raggedright\arraybackslash}p{31.188mm}>{\raggedright\arraybackslash}p{31.188mm}>{\raggedright\arraybackslash}p{31.188mm}@{}}
\toprule
\textbf{Case} & \textbf{Masculine singular} & \textbf{Neuter singular} & \textbf{Feminine singular} \\
\midrule
Nominative & geonga & geonge & geonge \\
Genitive & geongan & geongan & geongan \\
Dative & geongan & geongan & geongan \\
Accusative & geongan & geonge & geongan \\
\bottomrule
\end{tabular}
\par\endgroup
% course-panel:T58-P01:end
\par\addvspace{5pt}
% course-panel:T58-P02:begin
\par\begingroup
\fontsize{10.5}{13.2}\selectfont
\setlength{\LTleft}{0pt}\setlength{\LTright}{\fill}
\noindent
\begin{tabular}{@{}>{\raggedright\arraybackslash}p{22.0mm}>{\raggedright\arraybackslash}p{31.188mm}>{\raggedright\arraybackslash}p{31.188mm}>{\raggedright\arraybackslash}p{31.188mm}@{}}
\toprule
\textbf{Case} & \textbf{Masculine plural} & \textbf{Neuter plural} & \textbf{Feminine plural} \\
\midrule
Nominative & geongan & geongan & geongan \\
Genitive & geongra & geongra & geongra \\
Dative & geongum & geongum & geongum \\
Accusative & geongan & geongan & geongan \\
\bottomrule
\end{tabular}
\par\endgroup
% course-panel:T58-P02:end
% mdtable:T58:end

\reserve{7\baselineskip}
% mdparagraph:P237:begin
\subsection*{manigfeald}\BookHeading{nav-ch17-learner-0068}{subsubsection}{manigfeald}{0}
% mdparagraph:P237:end

% teaching: NEW-U17-VOC-MANIGFEALD-U15-SOURCE-SENSE
\par\begin{minipage}{\linewidth}\raggedright
% mdparagraph:P238:begin
\textbf{manigfeald} — manifold; numerous and varied.
% mdparagraph:P238:end

% teaching: NEW-U17-ADJ-MANIGFEALD
% mdparagraph:P239:begin
The stem \textbf{manigfeald-} is heavy. Strong feminine nominative singular and neuter nominative/\allowbreak{}accusative plural therefore have no added -u.
% mdparagraph:P239:end

% mdparagraph:P240:begin
\textbf{Strong.}
% mdparagraph:P240:end

% table: NEW-U17-ADJ-MANIGFEALD-STRONG-DISPLAY
% mdtable:T59:begin
% course-table-style: 2026-10
% course-panel:T59-P01:begin
\par\begingroup
\fontsize{10.5}{13.2}\selectfont
\setlength{\LTleft}{0pt}\setlength{\LTright}{\fill}
\noindent
\begin{tabular}{@{}>{\raggedright\arraybackslash}p{22.0mm}>{\raggedright\arraybackslash}p{31.188mm}>{\raggedright\arraybackslash}p{31.188mm}>{\raggedright\arraybackslash}p{31.188mm}@{}}
\toprule
\textbf{Case} & \textbf{Masculine singular} & \textbf{Neuter singular} & \textbf{Feminine singular} \\
\midrule
Nominative & manigfeald & manigfeald & manigfeald \\
Genitive & manigfealdes & manigfealdes & manigfealdre \\
Dative & manigfealdum & manigfealdum & manigfealdre \\
Accusative & manigfealdne & manigfeald & manigfealde \\
Instrumental & manigfealde & manigfealde & — \\
\bottomrule
\end{tabular}
\par\endgroup
% course-panel:T59-P01:end
\end{minipage}\par
\par\addvspace{5pt}
% course-panel:T59-P02:begin
\par\begingroup
\fontsize{10.5}{13.2}\selectfont
\setlength{\LTleft}{0pt}\setlength{\LTright}{\fill}
\noindent
\begin{tabular}{@{}>{\raggedright\arraybackslash}p{22.0mm}>{\raggedright\arraybackslash}p{31.188mm}>{\raggedright\arraybackslash}p{31.188mm}>{\raggedright\arraybackslash}p{31.188mm}@{}}
\toprule
\textbf{Case} & \textbf{Masculine plural} & \textbf{Neuter plural} & \textbf{Feminine plural} \\
\midrule
Nominative & manigfealde & manigfeald & manigfealda \\
Genitive & manigfealdra & manigfealdra & manigfealdra \\
Dative & manigfealdum & manigfealdum & manigfealdum \\
Accusative & manigfealde & manigfeald & manigfealda \\
Instrumental & — & — & — \\
\bottomrule
\end{tabular}
\par\endgroup
% course-panel:T59-P02:end
% mdtable:T59:end

\reserve{18\baselineskip}
\reserve{16\baselineskip}
% mdparagraph:P241:begin
\textbf{Weak.}
% mdparagraph:P241:end

% table: NEW-U17-ADJ-MANIGFEALD-WEAK-DISPLAY
% mdtable:T60:begin
% course-table-style: 2026-10
% course-panel:T60-P01:begin
\par\begingroup
\fontsize{10.5}{13.2}\selectfont
\setlength{\LTleft}{0pt}\setlength{\LTright}{\fill}
\noindent
\begin{tabular}{@{}>{\raggedright\arraybackslash}p{22.0mm}>{\raggedright\arraybackslash}p{31.188mm}>{\raggedright\arraybackslash}p{31.188mm}>{\raggedright\arraybackslash}p{31.188mm}@{}}
\toprule
\textbf{Case} & \textbf{Masculine singular} & \textbf{Neuter singular} & \textbf{Feminine singular} \\
\midrule
Nominative & manigfealda & manigfealde & manigfealde \\
Genitive & manigfealdan & manigfealdan & manigfealdan \\
Dative & manigfealdan & manigfealdan & manigfealdan \\
Accusative & manigfealdan & manigfealde & manigfealdan \\
\bottomrule
\end{tabular}
\par\endgroup
% course-panel:T60-P01:end
\par\addvspace{5pt}
% course-panel:T60-P02:begin
\par\begingroup
\fontsize{10.5}{13.2}\selectfont
\setlength{\LTleft}{0pt}\setlength{\LTright}{\fill}
\noindent
\begin{tabular}{@{}>{\raggedright\arraybackslash}p{22.0mm}>{\raggedright\arraybackslash}p{31.188mm}>{\raggedright\arraybackslash}p{31.188mm}>{\raggedright\arraybackslash}p{31.188mm}@{}}
\toprule
\textbf{Case} & \textbf{Masculine plural} & \textbf{Neuter plural} & \textbf{Feminine plural} \\
\midrule
Nominative & manigfealdan & manigfealdan & manigfealdan \\
Genitive & manigfealdra & manigfealdra & manigfealdra \\
Dative & manigfealdum & manigfealdum & manigfealdum \\
Accusative & manigfealdan & manigfealdan & manigfealdan \\
\bottomrule
\end{tabular}
\par\endgroup
% course-panel:T60-P02:end
% mdtable:T60:end

\reserve{7\baselineskip}
% mdparagraph:P242:begin
\subsection*{rīce}\BookHeading{nav-ch17-learner-0069}{subsubsection}{rīce}{0}
% mdparagraph:P242:end

% teaching: NEW-U17-VOC-RICE-ADJECTIVE-U17-SENSE
\par\begin{minipage}{\linewidth}\raggedright
% mdparagraph:P243:begin
\textbf{rīce} — rich; substantivally rich people.
% mdparagraph:P243:end

% teaching: NEW-U17-ADJ-RICE-ADJECTIVE
% mdparagraph:P244:begin
The headword ends in -e; replace that -e with an overt ending on \textbf{rīc-}. The displayed strong forms select feminine nominative singular and neuter nominative/\allowbreak{}accusative plural -u, and feminine plural -a. The source also permits -e in those positions.
% mdparagraph:P244:end

% mdparagraph:P245:begin
\textbf{Strong.}
% mdparagraph:P245:end

% table: NEW-U17-ADJ-RICE-ADJECTIVE-STRONG-DISPLAY
% mdtable:T61:begin
% course-table-style: 2026-10
% course-panel:T61-P01:begin
\par\begingroup
\fontsize{10.5}{13.2}\selectfont
\setlength{\LTleft}{0pt}\setlength{\LTright}{\fill}
\noindent
\begin{tabular}{@{}>{\raggedright\arraybackslash}p{22.0mm}>{\raggedright\arraybackslash}p{31.188mm}>{\raggedright\arraybackslash}p{31.188mm}>{\raggedright\arraybackslash}p{31.188mm}@{}}
\toprule
\textbf{Case} & \textbf{Masculine singular} & \textbf{Neuter singular} & \textbf{Feminine singular} \\
\midrule
Nominative & rīce & rīce & rīcu \\
Genitive & rīces & rīces & rīcre \\
Dative & rīcum & rīcum & rīcre \\
Accusative & rīcne & rīce & rīce \\
Instrumental & rīce & rīce & — \\
\bottomrule
\end{tabular}
\par\endgroup
% course-panel:T61-P01:end
\end{minipage}\par
\par\addvspace{5pt}
% course-panel:T61-P02:begin
\par\begingroup
\fontsize{10.5}{13.2}\selectfont
\setlength{\LTleft}{0pt}\setlength{\LTright}{\fill}
\noindent
\begin{tabular}{@{}>{\raggedright\arraybackslash}p{22.0mm}>{\raggedright\arraybackslash}p{31.188mm}>{\raggedright\arraybackslash}p{31.188mm}>{\raggedright\arraybackslash}p{31.188mm}@{}}
\toprule
\textbf{Case} & \textbf{Masculine plural} & \textbf{Neuter plural} & \textbf{Feminine plural} \\
\midrule
Nominative & rīce & rīcu & rīca \\
Genitive & rīcra & rīcra & rīcra \\
Dative & rīcum & rīcum & rīcum \\
Accusative & rīce & rīcu & rīca \\
Instrumental & — & — & — \\
\bottomrule
\end{tabular}
\par\endgroup
% course-panel:T61-P02:end
% mdtable:T61:end

\reserve{18\baselineskip}
\reserve{16\baselineskip}
% mdparagraph:P246:begin
\textbf{Weak.}
% mdparagraph:P246:end

% table: NEW-U17-ADJ-RICE-ADJECTIVE-WEAK-DISPLAY
% mdtable:T62:begin
% course-table-style: 2026-10
% course-panel:T62-P01:begin
\par\begingroup
\fontsize{10.5}{13.2}\selectfont
\setlength{\LTleft}{0pt}\setlength{\LTright}{\fill}
\noindent
\begin{tabular}{@{}>{\raggedright\arraybackslash}p{22.0mm}>{\raggedright\arraybackslash}p{31.188mm}>{\raggedright\arraybackslash}p{31.188mm}>{\raggedright\arraybackslash}p{31.188mm}@{}}
\toprule
\textbf{Case} & \textbf{Masculine singular} & \textbf{Neuter singular} & \textbf{Feminine singular} \\
\midrule
Nominative & rīca & rīce & rīce \\
Genitive & rīcan & rīcan & rīcan \\
Dative & rīcan & rīcan & rīcan \\
Accusative & rīcan & rīce & rīcan \\
\bottomrule
\end{tabular}
\par\endgroup
% course-panel:T62-P01:end
\par\addvspace{5pt}
% course-panel:T62-P02:begin
\par\begingroup
\fontsize{10.5}{13.2}\selectfont
\setlength{\LTleft}{0pt}\setlength{\LTright}{\fill}
\noindent
\begin{tabular}{@{}>{\raggedright\arraybackslash}p{22.0mm}>{\raggedright\arraybackslash}p{31.188mm}>{\raggedright\arraybackslash}p{31.188mm}>{\raggedright\arraybackslash}p{31.188mm}@{}}
\toprule
\textbf{Case} & \textbf{Masculine plural} & \textbf{Neuter plural} & \textbf{Feminine plural} \\
\midrule
Nominative & rīcan & rīcan & rīcan \\
Genitive & rīcra & rīcra & rīcra \\
Dative & rīcum & rīcum & rīcum \\
Accusative & rīcan & rīcan & rīcan \\
\bottomrule
\end{tabular}
\par\endgroup
% course-panel:T62-P02:end
% mdtable:T62:end

\reserve{7\baselineskip}
% mdparagraph:P247:begin
\subsection*{dæghwāmlīc}\BookHeading{nav-ch17-learner-0070}{subsubsection}{dæghwāmlīc}{0}
% mdparagraph:P247:end

% teaching: NEW-U17-VOC-DAEGHWAMLIC-U17-SENSE
\par\begin{minipage}{\linewidth}\raggedright
% mdparagraph:P248:begin
\textbf{dæghwāmlīc} — daily.
% mdparagraph:P248:end

% teaching: NEW-U17-ADJ-DAEGHWAMLIC
% mdparagraph:P249:begin
The stem \textbf{dæghwāmlīc-} is heavy. Strong feminine nominative singular and neuter nominative/\allowbreak{}accusative plural therefore have no added -u.
% mdparagraph:P249:end

% mdparagraph:P250:begin
\textbf{Strong.}
% mdparagraph:P250:end

% table: NEW-U17-ADJ-DAEGHWAMLIC-STRONG-DISPLAY
% mdtable:T63:begin
% course-table-style: 2026-10
% course-panel:T63-P01:begin
\par\begingroup
\fontsize{10.5}{13.2}\selectfont
\setlength{\LTleft}{0pt}\setlength{\LTright}{\fill}
\noindent
\begin{tabular}{@{}>{\raggedright\arraybackslash}p{22.0mm}>{\raggedright\arraybackslash}p{31.188mm}>{\raggedright\arraybackslash}p{31.188mm}>{\raggedright\arraybackslash}p{31.188mm}@{}}
\toprule
\textbf{Case} & \textbf{Masculine singular} & \textbf{Neuter singular} & \textbf{Feminine singular} \\
\midrule
Nominative & dæghwāmlīc & dæghwāmlīc & dæghwāmlīc \\
Genitive & dæghwāmlīces & dæghwāmlīces & dæghwāmlīcre \\
Dative & dæghwāmlīcum & dæghwāmlīcum & dæghwāmlīcre \\
Accusative & dæghwāmlīcne & dæghwāmlīc & dæghwāmlīce \\
Instrumental & dæghwāmlīce & dæghwāmlīce & — \\
\bottomrule
\end{tabular}
\par\endgroup
% course-panel:T63-P01:end
\end{minipage}\par
\par\addvspace{5pt}
% course-panel:T63-P02:begin
\par\begingroup
\fontsize{10.5}{13.2}\selectfont
\setlength{\LTleft}{0pt}\setlength{\LTright}{\fill}
\noindent
\begin{tabular}{@{}>{\raggedright\arraybackslash}p{22.0mm}>{\raggedright\arraybackslash}p{31.188mm}>{\raggedright\arraybackslash}p{31.188mm}>{\raggedright\arraybackslash}p{31.188mm}@{}}
\toprule
\textbf{Case} & \textbf{Masculine plural} & \textbf{Neuter plural} & \textbf{Feminine plural} \\
\midrule
Nominative & dæghwāmlīce & dæghwāmlīc & dæghwāmlīca \\
Genitive & dæghwāmlīcra & dæghwāmlīcra & dæghwāmlīcra \\
Dative & dæghwāmlīcum & dæghwāmlīcum & dæghwāmlīcum \\
Accusative & dæghwāmlīce & dæghwāmlīc & dæghwāmlīca \\
Instrumental & — & — & — \\
\bottomrule
\end{tabular}
\par\endgroup
% course-panel:T63-P02:end
% mdtable:T63:end

\reserve{18\baselineskip}
\reserve{16\baselineskip}
% mdparagraph:P251:begin
\textbf{Weak.}
% mdparagraph:P251:end

% table: NEW-U17-ADJ-DAEGHWAMLIC-WEAK-DISPLAY
% mdtable:T64:begin
% course-table-style: 2026-10
% course-panel:T64-P01:begin
\par\begingroup
\fontsize{10.5}{13.2}\selectfont
\setlength{\LTleft}{0pt}\setlength{\LTright}{\fill}
\noindent
\begin{tabular}{@{}>{\raggedright\arraybackslash}p{22.0mm}>{\raggedright\arraybackslash}p{31.188mm}>{\raggedright\arraybackslash}p{31.188mm}>{\raggedright\arraybackslash}p{31.188mm}@{}}
\toprule
\textbf{Case} & \textbf{Masculine singular} & \textbf{Neuter singular} & \textbf{Feminine singular} \\
\midrule
Nominative & dæghwāmlīca & dæghwāmlīce & dæghwāmlīce \\
Genitive & dæghwāmlīcan & dæghwāmlīcan & dæghwāmlīcan \\
Dative & dæghwāmlīcan & dæghwāmlīcan & dæghwāmlīcan \\
Accusative & dæghwāmlīcan & dæghwāmlīce & dæghwāmlīcan \\
\bottomrule
\end{tabular}
\par\endgroup
% course-panel:T64-P01:end
\par\addvspace{5pt}
% course-panel:T64-P02:begin
\par\begingroup
\fontsize{10.5}{13.2}\selectfont
\setlength{\LTleft}{0pt}\setlength{\LTright}{\fill}
\noindent
\begin{tabular}{@{}>{\raggedright\arraybackslash}p{22.0mm}>{\raggedright\arraybackslash}p{31.188mm}>{\raggedright\arraybackslash}p{31.188mm}>{\raggedright\arraybackslash}p{31.188mm}@{}}
\toprule
\textbf{Case} & \textbf{Masculine plural} & \textbf{Neuter plural} & \textbf{Feminine plural} \\
\midrule
Nominative & dæghwāmlīcan & dæghwāmlīcan & dæghwāmlīcan \\
Genitive & dæghwāmlīcra & dæghwāmlīcra & dæghwāmlīcra \\
Dative & dæghwāmlīcum & dæghwāmlīcum & dæghwāmlīcum \\
Accusative & dæghwāmlīcan & dæghwāmlīcan & dæghwāmlīcan \\
\bottomrule
\end{tabular}
\par\endgroup
% course-panel:T64-P02:end
% mdtable:T64:end

\reserve{7\baselineskip}
% mdparagraph:P252:begin
\subsection*{hlūd}\BookHeading{nav-ch17-learner-0071}{subsubsection}{hlūd}{0}
% mdparagraph:P252:end

% teaching: NEW-U17-VOC-HLUD-U17-SENSE
\par\begin{minipage}{\linewidth}\raggedright
% mdparagraph:P253:begin
\textbf{hlūd} — loud.
% mdparagraph:P253:end

% teaching: NEW-U17-ADJ-HLUD
% mdparagraph:P254:begin
The stem \textbf{hlūd-} is heavy. Strong feminine nominative singular and neuter nominative/\allowbreak{}accusative plural therefore have no added -u.
% mdparagraph:P254:end

% mdparagraph:P255:begin
\textbf{Strong.}
% mdparagraph:P255:end

% table: NEW-U17-ADJ-HLUD-STRONG-DISPLAY
% mdtable:T65:begin
% course-table-style: 2026-10
% course-panel:T65-P01:begin
\par\begingroup
\fontsize{10.5}{13.2}\selectfont
\setlength{\LTleft}{0pt}\setlength{\LTright}{\fill}
\noindent
\begin{tabular}{@{}>{\raggedright\arraybackslash}p{22.0mm}>{\raggedright\arraybackslash}p{31.188mm}>{\raggedright\arraybackslash}p{31.188mm}>{\raggedright\arraybackslash}p{31.188mm}@{}}
\toprule
\textbf{Case} & \textbf{Masculine singular} & \textbf{Neuter singular} & \textbf{Feminine singular} \\
\midrule
Nominative & hlūd & hlūd & hlūd \\
Genitive & hlūdes & hlūdes & hlūdre \\
Dative & hlūdum & hlūdum & hlūdre \\
Accusative & hlūdne & hlūd & hlūde \\
Instrumental & hlūde & hlūde & — \\
\bottomrule
\end{tabular}
\par\endgroup
% course-panel:T65-P01:end
\end{minipage}\par
\par\addvspace{5pt}
% course-panel:T65-P02:begin
\par\begingroup
\fontsize{10.5}{13.2}\selectfont
\setlength{\LTleft}{0pt}\setlength{\LTright}{\fill}
\noindent
\begin{tabular}{@{}>{\raggedright\arraybackslash}p{22.0mm}>{\raggedright\arraybackslash}p{31.188mm}>{\raggedright\arraybackslash}p{31.188mm}>{\raggedright\arraybackslash}p{31.188mm}@{}}
\toprule
\textbf{Case} & \textbf{Masculine plural} & \textbf{Neuter plural} & \textbf{Feminine plural} \\
\midrule
Nominative & hlūde & hlūd & hlūda \\
Genitive & hlūdra & hlūdra & hlūdra \\
Dative & hlūdum & hlūdum & hlūdum \\
Accusative & hlūde & hlūd & hlūda \\
Instrumental & — & — & — \\
\bottomrule
\end{tabular}
\par\endgroup
% course-panel:T65-P02:end
% mdtable:T65:end

\reserve{18\baselineskip}
\reserve{16\baselineskip}
% mdparagraph:P256:begin
\textbf{Weak.}
% mdparagraph:P256:end

% table: NEW-U17-ADJ-HLUD-WEAK-DISPLAY
% mdtable:T66:begin
% course-table-style: 2026-10
% course-panel:T66-P01:begin
\par\begingroup
\fontsize{10.5}{13.2}\selectfont
\setlength{\LTleft}{0pt}\setlength{\LTright}{\fill}
\noindent
\begin{tabular}{@{}>{\raggedright\arraybackslash}p{22.0mm}>{\raggedright\arraybackslash}p{31.188mm}>{\raggedright\arraybackslash}p{31.188mm}>{\raggedright\arraybackslash}p{31.188mm}@{}}
\toprule
\textbf{Case} & \textbf{Masculine singular} & \textbf{Neuter singular} & \textbf{Feminine singular} \\
\midrule
Nominative & hlūda & hlūde & hlūde \\
Genitive & hlūdan & hlūdan & hlūdan \\
Dative & hlūdan & hlūdan & hlūdan \\
Accusative & hlūdan & hlūde & hlūdan \\
\bottomrule
\end{tabular}
\par\endgroup
% course-panel:T66-P01:end
\par\addvspace{5pt}
% course-panel:T66-P02:begin
\par\begingroup
\fontsize{10.5}{13.2}\selectfont
\setlength{\LTleft}{0pt}\setlength{\LTright}{\fill}
\noindent
\begin{tabular}{@{}>{\raggedright\arraybackslash}p{22.0mm}>{\raggedright\arraybackslash}p{31.188mm}>{\raggedright\arraybackslash}p{31.188mm}>{\raggedright\arraybackslash}p{31.188mm}@{}}
\toprule
\textbf{Case} & \textbf{Masculine plural} & \textbf{Neuter plural} & \textbf{Feminine plural} \\
\midrule
Nominative & hlūdan & hlūdan & hlūdan \\
Genitive & hlūdra & hlūdra & hlūdra \\
Dative & hlūdum & hlūdum & hlūdum \\
Accusative & hlūdan & hlūdan & hlūdan \\
\bottomrule
\end{tabular}
\par\endgroup
% course-panel:T66-P02:end
% mdtable:T66:end

\reserve{7\baselineskip}
% mdparagraph:P257:begin
\subsection*{æðele}\BookHeading{nav-ch17-learner-0072}{subsubsection}{æðele}{0}
% mdparagraph:P257:end

% teaching: NEW-U17-VOC-AETHELE-U17-NOBLE
\par\begin{minipage}{\linewidth}\raggedright
% mdparagraph:P258:begin
\textbf{æðele} — noble; illustrious.
% mdparagraph:P258:end

% teaching: NEW-U17-ADJ-AETHELE
% mdparagraph:P259:begin
The headword ends in -e; replace that -e with an overt ending on \textbf{æðel-}. The displayed strong forms select feminine nominative singular and neuter nominative/\allowbreak{}accusative plural -u, and feminine plural -a. The source also permits -e in those positions.
% mdparagraph:P259:end

% mdparagraph:P260:begin
\textbf{Strong.}
% mdparagraph:P260:end

% table: NEW-U17-ADJ-AETHELE-STRONG-DISPLAY
% mdtable:T67:begin
% course-table-style: 2026-10
% course-panel:T67-P01:begin
\par\begingroup
\fontsize{10.5}{13.2}\selectfont
\setlength{\LTleft}{0pt}\setlength{\LTright}{\fill}
\noindent
\begin{tabular}{@{}>{\raggedright\arraybackslash}p{22.0mm}>{\raggedright\arraybackslash}p{31.188mm}>{\raggedright\arraybackslash}p{31.188mm}>{\raggedright\arraybackslash}p{31.188mm}@{}}
\toprule
\textbf{Case} & \textbf{Masculine singular} & \textbf{Neuter singular} & \textbf{Feminine singular} \\
\midrule
Nominative & æðele & æðele & æðelu \\
Genitive & æðeles & æðeles & æðelre \\
Dative & æðelum & æðelum & æðelre \\
Accusative & æðelne & æðele & æðele \\
Instrumental & æðele & æðele & — \\
\bottomrule
\end{tabular}
\par\endgroup
% course-panel:T67-P01:end
\end{minipage}\par
\par\addvspace{5pt}
% course-panel:T67-P02:begin
\par\begingroup
\fontsize{10.5}{13.2}\selectfont
\setlength{\LTleft}{0pt}\setlength{\LTright}{\fill}
\noindent
\begin{tabular}{@{}>{\raggedright\arraybackslash}p{22.0mm}>{\raggedright\arraybackslash}p{31.188mm}>{\raggedright\arraybackslash}p{31.188mm}>{\raggedright\arraybackslash}p{31.188mm}@{}}
\toprule
\textbf{Case} & \textbf{Masculine plural} & \textbf{Neuter plural} & \textbf{Feminine plural} \\
\midrule
Nominative & æðele & æðelu & æðela \\
Genitive & æðelra & æðelra & æðelra \\
Dative & æðelum & æðelum & æðelum \\
Accusative & æðele & æðelu & æðela \\
Instrumental & — & — & — \\
\bottomrule
\end{tabular}
\par\endgroup
% course-panel:T67-P02:end
% mdtable:T67:end

\reserve{18\baselineskip}
\reserve{16\baselineskip}
% mdparagraph:P261:begin
\textbf{Weak.}
% mdparagraph:P261:end

% table: NEW-U17-ADJ-AETHELE-WEAK-DISPLAY
% mdtable:T68:begin
% course-table-style: 2026-10
% course-panel:T68-P01:begin
\par\begingroup
\fontsize{10.5}{13.2}\selectfont
\setlength{\LTleft}{0pt}\setlength{\LTright}{\fill}
\noindent
\begin{tabular}{@{}>{\raggedright\arraybackslash}p{22.0mm}>{\raggedright\arraybackslash}p{31.188mm}>{\raggedright\arraybackslash}p{31.188mm}>{\raggedright\arraybackslash}p{31.188mm}@{}}
\toprule
\textbf{Case} & \textbf{Masculine singular} & \textbf{Neuter singular} & \textbf{Feminine singular} \\
\midrule
Nominative & æðela & æðele & æðele \\
Genitive & æðelan & æðelan & æðelan \\
Dative & æðelan & æðelan & æðelan \\
Accusative & æðelan & æðele & æðelan \\
\bottomrule
\end{tabular}
\par\endgroup
% course-panel:T68-P01:end
\par\addvspace{5pt}
% course-panel:T68-P02:begin
\par\begingroup
\fontsize{10.5}{13.2}\selectfont
\setlength{\LTleft}{0pt}\setlength{\LTright}{\fill}
\noindent
\begin{tabular}{@{}>{\raggedright\arraybackslash}p{22.0mm}>{\raggedright\arraybackslash}p{31.188mm}>{\raggedright\arraybackslash}p{31.188mm}>{\raggedright\arraybackslash}p{31.188mm}@{}}
\toprule
\textbf{Case} & \textbf{Masculine plural} & \textbf{Neuter plural} & \textbf{Feminine plural} \\
\midrule
Nominative & æðelan & æðelan & æðelan \\
Genitive & æðelra & æðelra & æðelra \\
Dative & æðelum & æðelum & æðelum \\
Accusative & æðelan & æðelan & æðelan \\
\bottomrule
\end{tabular}
\par\endgroup
% course-panel:T68-P02:end
% mdtable:T68:end

\reserve{7\baselineskip}
% mdparagraph:P262:begin
\subsection*{æðelboren}\BookHeading{nav-ch17-learner-0073}{subsubsection}{æðelboren}{0}
% mdparagraph:P262:end

% teaching: NEW-U17-VOC-AETHELBOREN-U17-SENSE
\par\begin{minipage}{\linewidth}\raggedright
% mdparagraph:P263:begin
\textbf{æðelboren} — high-born.
% mdparagraph:P263:end

% teaching: NEW-U17-ADJ-AETHELBOREN
% mdparagraph:P264:begin
Use the retained stem æðelboren- throughout this selected table. Do not remove a vowel merely by analogy with a different word.
% mdparagraph:P264:end

% mdparagraph:P265:begin
\textbf{Strong.}
% mdparagraph:P265:end

% table: NEW-U17-ADJ-AETHELBOREN-STRONG-DISPLAY
% mdtable:T69:begin
% course-table-style: 2026-10
% course-panel:T69-P01:begin
\par\begingroup
\fontsize{10.5}{13.2}\selectfont
\setlength{\LTleft}{0pt}\setlength{\LTright}{\fill}
\noindent
\begin{tabular}{@{}>{\raggedright\arraybackslash}p{22.0mm}>{\raggedright\arraybackslash}p{31.188mm}>{\raggedright\arraybackslash}p{31.188mm}>{\raggedright\arraybackslash}p{31.188mm}@{}}
\toprule
\textbf{Case} & \textbf{Masculine singular} & \textbf{Neuter singular} & \textbf{Feminine singular} \\
\midrule
Nominative & æðelboren & æðelboren & æðelborenu \\
Genitive & æðelborenes & æðelborenes & æðelborenre \\
Dative & æðelborenum & æðelborenum & æðelborenre \\
Accusative & æðelborenne & æðelboren & æðelborene \\
Instrumental & æðelborene & æðelborene & — \\
\bottomrule
\end{tabular}
\par\endgroup
% course-panel:T69-P01:end
\end{minipage}\par
\par\addvspace{5pt}
% course-panel:T69-P02:begin
\par\begingroup
\fontsize{10.5}{13.2}\selectfont
\setlength{\LTleft}{0pt}\setlength{\LTright}{\fill}
\noindent
\begin{tabular}{@{}>{\raggedright\arraybackslash}p{22.0mm}>{\raggedright\arraybackslash}p{31.188mm}>{\raggedright\arraybackslash}p{31.188mm}>{\raggedright\arraybackslash}p{31.188mm}@{}}
\toprule
\textbf{Case} & \textbf{Masculine plural} & \textbf{Neuter plural} & \textbf{Feminine plural} \\
\midrule
Nominative & æðelborene & æðelborenu & æðelborena \\
Genitive & æðelborenra & æðelborenra & æðelborenra \\
Dative & æðelborenum & æðelborenum & æðelborenum \\
Accusative & æðelborene & æðelborenu & æðelborena \\
Instrumental & — & — & — \\
\bottomrule
\end{tabular}
\par\endgroup
% course-panel:T69-P02:end
% mdtable:T69:end

\reserve{18\baselineskip}
\reserve{16\baselineskip}
% mdparagraph:P266:begin
\textbf{Weak.}
% mdparagraph:P266:end

% table: NEW-U17-ADJ-AETHELBOREN-WEAK-DISPLAY
% mdtable:T70:begin
% course-table-style: 2026-10
% course-panel:T70-P01:begin
\par\begingroup
\fontsize{10.5}{13.2}\selectfont
\setlength{\LTleft}{0pt}\setlength{\LTright}{\fill}
\noindent
\begin{tabular}{@{}>{\raggedright\arraybackslash}p{22.0mm}>{\raggedright\arraybackslash}p{31.188mm}>{\raggedright\arraybackslash}p{31.188mm}>{\raggedright\arraybackslash}p{31.188mm}@{}}
\toprule
\textbf{Case} & \textbf{Masculine singular} & \textbf{Neuter singular} & \textbf{Feminine singular} \\
\midrule
Nominative & æðelborena & æðelborene & æðelborene \\
Genitive & æðelborenan & æðelborenan & æðelborenan \\
Dative & æðelborenan & æðelborenan & æðelborenan \\
Accusative & æðelborenan & æðelborene & æðelborenan \\
\bottomrule
\end{tabular}
\par\endgroup
% course-panel:T70-P01:end
\par\addvspace{5pt}
% course-panel:T70-P02:begin
\par\begingroup
\fontsize{10.5}{13.2}\selectfont
\setlength{\LTleft}{0pt}\setlength{\LTright}{\fill}
\noindent
\begin{tabular}{@{}>{\raggedright\arraybackslash}p{22.0mm}>{\raggedright\arraybackslash}p{31.188mm}>{\raggedright\arraybackslash}p{31.188mm}>{\raggedright\arraybackslash}p{31.188mm}@{}}
\toprule
\textbf{Case} & \textbf{Masculine plural} & \textbf{Neuter plural} & \textbf{Feminine plural} \\
\midrule
Nominative & æðelborenan & æðelborenan & æðelborenan \\
Genitive & æðelborenra & æðelborenra & æðelborenra \\
Dative & æðelborenum & æðelborenum & æðelborenum \\
Accusative & æðelborenan & æðelborenan & æðelborenan \\
\bottomrule
\end{tabular}
\par\endgroup
% course-panel:T70-P02:end
% mdtable:T70:end

\reserve{7\baselineskip}
% mdparagraph:P267:begin
\subsection*{Adjectives naming people}\BookHeading{nav-ch17-learner-0074}{subsubsection}{Adjectives naming people}{0}
% mdparagraph:P267:end

% teaching: NEW-U17-GRAM-U17-PERSONAL-ADJECTIVE-GROUPS
% teaching: NEW-U17-VOC-MIHTIG-U17-POWERFUL-PEOPLE
% teaching: NEW-U17-VOC-EADMOD-U17-HUMBLE-PEOPLE
% mdparagraph:P268:begin
The selected adjectives geong, eald, mihtig, ēaðmōd and rīce can name young, old, powerful, humble and rich people. Use weak masculine plural forms in the assigned address and object expressions. Geongan and ealdan address people without an expressed men or article under Cook 37’s instruction. Ðā mihtgan, ðā ēaðmōdan and ðā rīcan are accusative plural objects. These people meanings are distinct from merely describing a noun; this unit assigns these particular plural uses, not an unrestricted noun declension.
% mdparagraph:P268:end

\reserve{7\baselineskip}
% mdparagraph:P269:begin
\section*{3. Verbs and nonfinite forms}\BookHeading{nav-ch17-learner-0075}{subsection}{3. Verbs and nonfinite forms}{1}
% mdparagraph:P269:end

% mdparagraph:P270:begin
Each heading states exactly which verbal system is being learned. A complete indicative table gives all three persons in both numbers for the present and preterite. A command pair gives singular and plural; a subjunctive table gives all persons in its named tense. A two-form nonfinite display gives the stated pair, rather than an unprinted finite conjugation.
% mdparagraph:P270:end

% teaching: NEW-U17-DICTIONARY-REFERENCES
% mdparagraph:P271:begin
\textbf{Dictionary references:} hęrian, praise; ondrǣdan, fear; hyngran, be hungry. These citation forms remain visible. A citation identifies the word; it does not by itself teach a new infinitival construction. The independent forms are the ones printed in the bounded systems below.
% mdparagraph:P271:end

% mdparagraph:P272:begin
Strong verbs change their stem vowels in the preterite; the second singular preterite normally shares the plural stem. Weak verbs use a dental preterite, whose spelling can merge with the last consonant of the stem. Present subjunctives generally have singular -e and plural -en; weak -ian verbs retain i in that pattern. Learn an exceptional stem together with its endings.
% mdparagraph:P272:end

\reserve{7\baselineskip}
% mdparagraph:P273:begin
\subsection*{hęrian: indicative, commands and present subjunctive}\BookHeading{nav-ch17-learner-0076}{subsubsection}{hęrian: indicative, commands and present subjunctive}{0}
% mdparagraph:P273:end

% teaching: NEW-U17-VOC-HERIAN-U17-SELECTED
\par\begin{minipage}{\linewidth}\raggedright
% mdparagraph:P274:begin
\textbf{hęrian} — praise.
% mdparagraph:P274:end

% teaching: u17-herian-indicative12-imperative2-present-subjunctive6
% mdparagraph:P275:begin
This weak verb follows the nęrian type. The present has hęrie, hęrest and hęreð, with -i- in plural hęriað. The preterite stem is hęred-. The command is singular hęre, plural hęriað; the present subjunctive is singular hęrie, plural hęrien. Thus hęriað can be a statement or a command, whereas hęrien expresses a plural wish in the praise exercise.
% mdparagraph:P275:end

% mdparagraph:P276:begin
\textbf{Indicative.}
% mdparagraph:P276:end

% table: u17-herian-indicative12-imperative2-present-subjunctive6-INDICATIVE-DISPLAY
% mdtable:T71:begin
% course-table-style: 2026-10
% course-panel:T71-P01:begin
\par\begingroup
\fontsize{10.5}{13.2}\selectfont
\setlength{\LTleft}{0pt}\setlength{\LTright}{\fill}
\noindent
\begin{tabular}{@{}>{\raggedright\arraybackslash}p{39.459mm}>{\raggedright\arraybackslash}p{39.459mm}>{\raggedright\arraybackslash}p{39.459mm}@{}}
\toprule
\textbf{Person and number} & \textbf{Present} & \textbf{Preterite} \\
\midrule
1 sg. & hęrie & hęrede \\
2 sg. & hęrest & hęredest \\
3 sg. & hęreð & hęrede \\
1–3 pl. & hęriað & hęredon \\
\bottomrule
\end{tabular}
\par\endgroup
% course-panel:T71-P01:end
% mdtable:T71:end
\end{minipage}\par
\reserve{9\baselineskip}
\reserve{7\baselineskip}
% mdparagraph:P277:begin
\textbf{Imperative.}
% mdparagraph:P277:end

% table: u17-herian-indicative12-imperative2-present-subjunctive6-IMPERATIVE-DISPLAY
% mdtable:T72:begin
% course-table-style: 2026-10
% course-panel:T72-P01:begin
\par\begingroup
\fontsize{10.5}{13.2}\selectfont
\setlength{\LTleft}{0pt}\setlength{\LTright}{\fill}
\noindent
\begin{tabular}{@{}>{\raggedright\arraybackslash}p{63.342mm}>{\raggedright\arraybackslash}p{57.846mm}@{}}
\toprule
\textbf{Person and number} & \textbf{Command} \\
\midrule
2 sg. & hęre \\
2 pl. & hęriað \\
\bottomrule
\end{tabular}
\par\endgroup
% course-panel:T72-P01:end
% mdtable:T72:end
\reserve{13\baselineskip}
\reserve{11\baselineskip}
% mdparagraph:P278:begin
\textbf{Subjunctive.}
% mdparagraph:P278:end

% table: u17-herian-indicative12-imperative2-present-subjunctive6-SUBJUNCTIVE-DISPLAY
% mdtable:T73:begin
% course-table-style: 2026-10
% course-panel:T73-P01:begin
\par\begingroup
\fontsize{10.5}{13.2}\selectfont
\setlength{\LTleft}{0pt}\setlength{\LTright}{\fill}
\noindent
\begin{tabular}{@{}>{\raggedright\arraybackslash}p{63.342mm}>{\raggedright\arraybackslash}p{57.846mm}@{}}
\toprule
\textbf{Person and number} & \textbf{Present} \\
\midrule
1 sg. & hęrie \\
2 sg. & hęrie \\
3 sg. & hęrie \\
1–3 pl. & hęrien \\
\bottomrule
\end{tabular}
\par\endgroup
% course-panel:T73-P01:end
% mdtable:T73:end
\reserve{7\baselineskip}
% mdparagraph:P279:begin
\subsection*{gecweðan: present and preterite indicative}\BookHeading{nav-ch17-learner-0077}{subsubsection}{gecweðan: present and preterite indicative}{0}
% mdparagraph:P279:end

% teaching: NEW-U17-VOC-GECWETHAN-U17-SELECTED
\par\begin{minipage}{\linewidth}\raggedright
% mdparagraph:P280:begin
\textbf{gecweðan} — speak; say.
% mdparagraph:P280:end

% teaching: u17-gecwethan-indicative12
% mdparagraph:P281:begin
The present second and third singular have i. In the preterite, first and third singular gecwæð retain æ and ð; second singular gecwǣde and plural gecwǣdon have ǣ and d. Retain the lexical prefix ge-.
% mdparagraph:P281:end

% table: u17-gecwethan-indicative12-INDICATIVE-DISPLAY
% mdtable:T74:begin
% course-table-style: 2026-10
% course-panel:T74-P01:begin
\par\begingroup
\fontsize{10.5}{13.2}\selectfont
\setlength{\LTleft}{0pt}\setlength{\LTright}{\fill}
\noindent
\begin{tabular}{@{}>{\raggedright\arraybackslash}p{39.459mm}>{\raggedright\arraybackslash}p{39.459mm}>{\raggedright\arraybackslash}p{39.459mm}@{}}
\toprule
\textbf{Person and number} & \textbf{Present} & \textbf{Preterite} \\
\midrule
1 sg. & gecweðe & gecwæð \\
2 sg. & gecwist & gecwǣde \\
3 sg. & gecwið & gecwæð \\
1–3 pl. & gecweðað & gecwǣdon \\
\bottomrule
\end{tabular}
\par\endgroup
% course-panel:T74-P01:end
% mdtable:T74:end
\end{minipage}\par
\reserve{7\baselineskip}
% mdparagraph:P282:begin
\subsection*{ūpāhębban: indicative and participles}\BookHeading{nav-ch17-learner-0078}{subsubsection}{ūpāhębban: indicative and participles}{0}
% mdparagraph:P282:end

% teaching: NEW-U17-VOC-UP-AHEBBAN-U17-SELECTED
\par\begin{minipage}{\linewidth}\raggedright
% mdparagraph:P283:begin
\textbf{ūp-ā-hębban} — raise; exalt.
% mdparagraph:P283:end

% teaching: u17-up-ahebban-indicative12-participle-pair2
% mdparagraph:P284:begin
The present stem has doubled bb, but the second and third singular are ūpāhęfst and ūpāhęfð. The preterite has hōf and the past participle hafen, with ūpā- retained. Raising can be physical or, in the explicitly taught exaltation sense, an increase in standing.
% mdparagraph:P284:end

% mdparagraph:P285:begin
\textbf{Indicative.}
% mdparagraph:P285:end

% table: u17-up-ahebban-indicative12-participle-pair2-INDICATIVE-DISPLAY
% mdtable:T75:begin
% course-table-style: 2026-10
% course-panel:T75-P01:begin
\par\begingroup
\fontsize{10.5}{13.2}\selectfont
\setlength{\LTleft}{0pt}\setlength{\LTright}{\fill}
\noindent
\begin{tabular}{@{}>{\raggedright\arraybackslash}p{39.459mm}>{\raggedright\arraybackslash}p{39.459mm}>{\raggedright\arraybackslash}p{39.459mm}@{}}
\toprule
\textbf{Person and number} & \textbf{Present} & \textbf{Preterite} \\
\midrule
1 sg. & ūpāhębbe & ūpāhōf \\
2 sg. & ūpāhęfst & ūpāhōfe \\
3 sg. & ūpāhęfð & ūpāhōf \\
1–3 pl. & ūpāhębbað & ūpāhōfon \\
\bottomrule
\end{tabular}
\par\endgroup
% course-panel:T75-P01:end
% mdtable:T75:end
\end{minipage}\par
\reserve{9\baselineskip}
\reserve{7\baselineskip}
% mdparagraph:P286:begin
\textbf{Participle formation pair.}
% mdparagraph:P286:end

% table: u17-up-ahebban-indicative12-participle-pair2-PARTICIPLES-DISPLAY
% mdtable:T76:begin
\par\begingroup
\fontsize{10.5}{13.2}\selectfont
\setlength{\LTleft}{0pt}\setlength{\LTright}{\fill}
\noindent
\begin{tabular}{@{}>{\raggedright\arraybackslash}p{53.15mm}>{\raggedright\arraybackslash}p{68.04mm}@{}}
\toprule
\textbf{Form} & \textbf{Literal surface} \\
\midrule
present & ūpāhębbende \\
past & ūpāhafen \\
\bottomrule
\end{tabular}
\par\endgroup

% mdtable:T76:end
\reserve{7\baselineskip}
% mdparagraph:P287:begin
\subsection*{oferwyrcean: present and preterite indicative}\BookHeading{nav-ch17-learner-0079}{subsubsection}{oferwyrcean: present and preterite indicative}{0}
% mdparagraph:P287:end

% teaching: NEW-U17-VOC-OFERWYRCEAN-U17-SELECTED
\par\begin{minipage}{\linewidth}\raggedright
% mdparagraph:P288:begin
\textbf{oferwyrcean} — overlay; cover with a worked covering.
% mdparagraph:P288:end

% teaching: u17-oferwyrcean-indicative12
% mdparagraph:P289:begin
The selected present uses -ean and plural -eað. The weak preterite has worht-. Retain ofer- throughout. The thing overlaid is accusative; the covering material follows mid with the dative.
% mdparagraph:P289:end

% table: u17-oferwyrcean-indicative12-INDICATIVE-DISPLAY
% mdtable:T77:begin
% course-table-style: 2026-10
% course-panel:T77-P01:begin
\par\begingroup
\fontsize{10.5}{13.2}\selectfont
\setlength{\LTleft}{0pt}\setlength{\LTright}{\fill}
\noindent
\begin{tabular}{@{}>{\raggedright\arraybackslash}p{39.459mm}>{\raggedright\arraybackslash}p{39.459mm}>{\raggedright\arraybackslash}p{39.459mm}@{}}
\toprule
\textbf{Person and number} & \textbf{Present} & \textbf{Preterite} \\
\midrule
1 sg. & oferwyrce & oferworhte \\
2 sg. & oferwyrcst & oferworhtest \\
3 sg. & oferwyrcð & oferworhte \\
1–3 pl. & oferwyrceað & oferworhton \\
\bottomrule
\end{tabular}
\par\endgroup
% course-panel:T77-P01:end
% mdtable:T77:end
\end{minipage}\par
\reserve{7\baselineskip}
% mdparagraph:P290:begin
\subsection*{gesęttan: present and preterite indicative}\BookHeading{nav-ch17-learner-0080}{subsubsection}{gesęttan: present and preterite indicative}{0}
% mdparagraph:P290:end

% teaching: NEW-U17-VOC-GESETTAN-U17-SELECTED
\par\begin{minipage}{\linewidth}\raggedright
% mdparagraph:P291:begin
\textbf{gesęttan} — set; place; establish.
% mdparagraph:P291:end

% teaching: u17-gesettan-indicative12
% mdparagraph:P292:begin
Retain tt before a vowel. The selected present second singular has -tst and third singular -tt. The weak dental preterite merges with the final t; it does not add a separate -tde sequence.
% mdparagraph:P292:end

% table: u17-gesettan-indicative12-INDICATIVE-DISPLAY
% mdtable:T78:begin
% course-table-style: 2026-10
% course-panel:T78-P01:begin
\par\begingroup
\fontsize{10.5}{13.2}\selectfont
\setlength{\LTleft}{0pt}\setlength{\LTright}{\fill}
\noindent
\begin{tabular}{@{}>{\raggedright\arraybackslash}p{39.459mm}>{\raggedright\arraybackslash}p{39.459mm}>{\raggedright\arraybackslash}p{39.459mm}@{}}
\toprule
\textbf{Person and number} & \textbf{Present} & \textbf{Preterite} \\
\midrule
1 sg. & gesętte & gesętte \\
2 sg. & gesętst & gesęttest \\
3 sg. & gesętt & gesętte \\
1–3 pl. & gesęttað & gesętton \\
\bottomrule
\end{tabular}
\par\endgroup
% course-panel:T78-P01:end
% mdtable:T78:end
\end{minipage}\par
\reserve{7\baselineskip}
% mdparagraph:P293:begin
\subsection*{giernan: present and preterite indicative}\BookHeading{nav-ch17-learner-0081}{subsubsection}{giernan: present and preterite indicative}{0}
% mdparagraph:P293:end

% teaching: NEW-U17-VOC-GIERNAN-U17-SELECTED
\par\begin{minipage}{\linewidth}\raggedright
% mdparagraph:P294:begin
\textbf{giernan} — desire; request; ask for; seek in marriage.
% mdparagraph:P294:end

% teaching: u17-giernan-indicative12
% mdparagraph:P295:begin
This long-stem weak verb has present -st and -ð and a preterite in -de. The thing desired or person sought in marriage is genitive. Æt with the dative names the person from whom something is requested.
% mdparagraph:P295:end

% table: u17-giernan-indicative12-INDICATIVE-DISPLAY
% mdtable:T79:begin
% course-table-style: 2026-10
% course-panel:T79-P01:begin
\par\begingroup
\fontsize{10.5}{13.2}\selectfont
\setlength{\LTleft}{0pt}\setlength{\LTright}{\fill}
\noindent
\begin{tabular}{@{}>{\raggedright\arraybackslash}p{39.459mm}>{\raggedright\arraybackslash}p{39.459mm}>{\raggedright\arraybackslash}p{39.459mm}@{}}
\toprule
\textbf{Person and number} & \textbf{Present} & \textbf{Preterite} \\
\midrule
1 sg. & gierne & giernde \\
2 sg. & giernst & gierndest \\
3 sg. & giernð & giernde \\
1–3 pl. & giernað & gierndon \\
\bottomrule
\end{tabular}
\par\endgroup
% course-panel:T79-P01:end
% mdtable:T79:end
\end{minipage}\par
\reserve{7\baselineskip}
% mdparagraph:P296:begin
\subsection*{mǣrsian: present and preterite indicative}\BookHeading{nav-ch17-learner-0082}{subsubsection}{mǣrsian: present and preterite indicative}{0}
% mdparagraph:P296:end

% teaching: NEW-U17-VOC-MAERSIAN-U17-SELECTED
\par\begin{minipage}{\linewidth}\raggedright
% mdparagraph:P297:begin
\textbf{mǣrsian} — make great; magnify; glorify.
% mdparagraph:P297:end

% teaching: u17-maersian-indicative12
% mdparagraph:P298:begin
The weak -ian pattern has present singular -ie, -ast, -að and plural -iað; its preterite uses -ode, -odest, -odon. The participles have -iende and -od. In the passive, mǣrsod means celebrated or exalted.
% mdparagraph:P298:end

% table: u17-maersian-indicative12-INDICATIVE-DISPLAY
% mdtable:T80:begin
% course-table-style: 2026-10
% course-panel:T80-P01:begin
\par\begingroup
\fontsize{10.5}{13.2}\selectfont
\setlength{\LTleft}{0pt}\setlength{\LTright}{\fill}
\noindent
\begin{tabular}{@{}>{\raggedright\arraybackslash}p{39.459mm}>{\raggedright\arraybackslash}p{39.459mm}>{\raggedright\arraybackslash}p{39.459mm}@{}}
\toprule
\textbf{Person and number} & \textbf{Present} & \textbf{Preterite} \\
\midrule
1 sg. & mǣrsie & mǣrsode \\
2 sg. & mǣrsast & mǣrsodest \\
3 sg. & mǣrsað & mǣrsode \\
1–3 pl. & mǣrsiað & mǣrsodon \\
\bottomrule
\end{tabular}
\par\endgroup
% course-panel:T80-P01:end
% mdtable:T80:end
\end{minipage}\par
\reserve{7\baselineskip}
% page-review repair: U17LF01: participle heading, rules and complete two-form table
\par\begin{minipage}{\linewidth}\raggedright
% mdparagraph:P299:begin
\subsection*{mǣrsian: present and past participles}\BookHeading{nav-ch17-learner-0083}{subsubsection}{mǣrsian: present and past participles}{0}
% mdparagraph:P299:end

% teaching: u17-maersian-participle-pair2
% mdparagraph:P300:begin
The present participle has -iende and the past participle -od. The latter is used in the passive description celebrated or exalted.
% mdparagraph:P300:end

% mdparagraph:P301:begin
This display assigns the complete present/\allowbreak{}past formation pair. It does not add other finite moods or an unrestricted attributive declension; the past form is used in the stated passive frame.
% mdparagraph:P301:end

% table: u17-maersian-participle-pair2-PARTICIPLES-DISPLAY
% mdtable:T81:begin
\par\begingroup
\fontsize{10.5}{13.2}\selectfont
\setlength{\LTleft}{0pt}\setlength{\LTright}{\fill}
\noindent
\begin{tabular}{@{}>{\raggedright\arraybackslash}p{55.85mm}>{\raggedright\arraybackslash}p{65.34mm}@{}}
\toprule
\textbf{Form} & \textbf{Literal surface} \\
\midrule
present & mǣrsiende \\
past & mǣrsod \\
\bottomrule
\end{tabular}
\par\endgroup

% mdtable:T81:end
\end{minipage}\par

\reserve{7\baselineskip}
% mdparagraph:P302:begin
\subsection*{gewilnian: present and preterite indicative}\BookHeading{nav-ch17-learner-0084}{subsubsection}{gewilnian: present and preterite indicative}{0}
% mdparagraph:P302:end

% teaching: NEW-U17-VOC-GEWILNIAN-U17-SELECTED
\par\begin{minipage}{\linewidth}\raggedright
% mdparagraph:P303:begin
\textbf{gewilnian} — desire; wish.
% mdparagraph:P303:end

% teaching: u17-gewilnian-indicative12
% mdparagraph:P304:begin
Use the weak -ian endings: present -ie, -ast, -að, -iað; preterite -ode, -odest, -odon. A following ðæt clause can express what is desired.
% mdparagraph:P304:end

% table: u17-gewilnian-indicative12-INDICATIVE-DISPLAY
% mdtable:T82:begin
% course-table-style: 2026-10
% course-panel:T82-P01:begin
\par\begingroup
\fontsize{10.5}{13.2}\selectfont
\setlength{\LTleft}{0pt}\setlength{\LTright}{\fill}
\noindent
\begin{tabular}{@{}>{\raggedright\arraybackslash}p{39.459mm}>{\raggedright\arraybackslash}p{39.459mm}>{\raggedright\arraybackslash}p{39.459mm}@{}}
\toprule
\textbf{Person and number} & \textbf{Present} & \textbf{Preterite} \\
\midrule
1 sg. & gewilnie & gewilnode \\
2 sg. & gewilnast & gewilnodest \\
3 sg. & gewilnað & gewilnode \\
1–3 pl. & gewilniað & gewilnodon \\
\bottomrule
\end{tabular}
\par\endgroup
% course-panel:T82-P01:end
% mdtable:T82:end
\end{minipage}\par
\reserve{7\baselineskip}
% mdparagraph:P305:begin
\subsection*{ðurhwunian: present and preterite indicative}\BookHeading{nav-ch17-learner-0085}{subsubsection}{ðurhwunian: present and preterite indicative}{0}
% mdparagraph:P305:end

% teaching: NEW-U17-VOC-THURHWUNIAN-U17-SELECTED
\par\begin{minipage}{\linewidth}\raggedright
% mdparagraph:P306:begin
\textbf{ðurhwunian} — remain; endure; last.
% mdparagraph:P306:end

% teaching: u17-thurhwunian-indicative12
% mdparagraph:P307:begin
Use the weak -ian endings: present -ie, -ast, -að, -iað; preterite -ode, -odest, -odon. The meaning is remaining or lasting.
% mdparagraph:P307:end

% table: u17-thurhwunian-indicative12-INDICATIVE-DISPLAY
% mdtable:T83:begin
% course-table-style: 2026-10
% course-panel:T83-P01:begin
\par\begingroup
\fontsize{10.5}{13.2}\selectfont
\setlength{\LTleft}{0pt}\setlength{\LTright}{\fill}
\noindent
\begin{tabular}{@{}>{\raggedright\arraybackslash}p{39.459mm}>{\raggedright\arraybackslash}p{39.459mm}>{\raggedright\arraybackslash}p{39.459mm}@{}}
\toprule
\textbf{Person and number} & \textbf{Present} & \textbf{Preterite} \\
\midrule
1 sg. & ðurhwunie & ðurhwunode \\
2 sg. & ðurhwunast & ðurhwunodest \\
3 sg. & ðurhwunað & ðurhwunode \\
1–3 pl. & ðurhwuniað & ðurhwunodon \\
\bottomrule
\end{tabular}
\par\endgroup
% course-panel:T83-P01:end
% mdtable:T83:end
\end{minipage}\par
\reserve{7\baselineskip}
% mdparagraph:P308:begin
\subsection*{wiðstandan: present and preterite indicative}\BookHeading{nav-ch17-learner-0086}{subsubsection}{wiðstandan: present and preterite indicative}{0}
% mdparagraph:P308:end

% teaching: NEW-U17-VOC-WITHSTANDAN-U17-SELECTED
\par\begin{minipage}{\linewidth}\raggedright
% mdparagraph:P309:begin
\textbf{wiðstandan} — withstand; resist.
% mdparagraph:P309:end

% teaching: u17-withstandan-indicative12
% mdparagraph:P310:begin
The present second and third singular have stęntst and stęnt. The preterite has stōd and stōdon, without the present stem’s n. The thing resisted is dative.
% mdparagraph:P310:end

% table: u17-withstandan-indicative12-INDICATIVE-DISPLAY
% mdtable:T84:begin
% course-table-style: 2026-10
% course-panel:T84-P01:begin
\par\begingroup
\fontsize{10.5}{13.2}\selectfont
\setlength{\LTleft}{0pt}\setlength{\LTright}{\fill}
\noindent
\begin{tabular}{@{}>{\raggedright\arraybackslash}p{39.459mm}>{\raggedright\arraybackslash}p{39.459mm}>{\raggedright\arraybackslash}p{39.459mm}@{}}
\toprule
\textbf{Person and number} & \textbf{Present} & \textbf{Preterite} \\
\midrule
1 sg. & wiðstande & wiðstōd \\
2 sg. & wiðstęntst & wiðstōde \\
3 sg. & wiðstęnt & wiðstōd \\
1–3 pl. & wiðstandað & wiðstōdon \\
\bottomrule
\end{tabular}
\par\endgroup
% course-panel:T84-P01:end
% mdtable:T84:end
\end{minipage}\par
\reserve{7\baselineskip}
% mdparagraph:P311:begin
\subsection*{gearcian: present and preterite indicative}\BookHeading{nav-ch17-learner-0087}{subsubsection}{gearcian: present and preterite indicative}{0}
% mdparagraph:P311:end

% teaching: NEW-U17-VOC-GEARCIAN-U17-SELECTED
\par\begin{minipage}{\linewidth}\raggedright
% mdparagraph:P312:begin
\textbf{gearcian} — prepare.
% mdparagraph:P312:end

% teaching: u17-gearcian-indicative12
% mdparagraph:P313:begin
The verb is unprefixed gearcian. Use the weak -ian endings: present -ie, -ast, -að, -iað; preterite -ode, -odest, -odon. The things prepared are accusative and their beneficiaries dative.
% mdparagraph:P313:end

% table: u17-gearcian-indicative12-INDICATIVE-DISPLAY
% mdtable:T85:begin
% course-table-style: 2026-10
% course-panel:T85-P01:begin
\par\begingroup
\fontsize{10.5}{13.2}\selectfont
\setlength{\LTleft}{0pt}\setlength{\LTright}{\fill}
\noindent
\begin{tabular}{@{}>{\raggedright\arraybackslash}p{39.459mm}>{\raggedright\arraybackslash}p{39.459mm}>{\raggedright\arraybackslash}p{39.459mm}@{}}
\toprule
\textbf{Person and number} & \textbf{Present} & \textbf{Preterite} \\
\midrule
1 sg. & gearcie & gearcode \\
2 sg. & gearcast & gearcodest \\
3 sg. & gearcað & gearcode \\
1–3 pl. & gearciað & gearcodon \\
\bottomrule
\end{tabular}
\par\endgroup
% course-panel:T85-P01:end
% mdtable:T85:end
\end{minipage}\par
\reserve{7\baselineskip}
% mdparagraph:P314:begin
\subsection*{āfyllan: indicative and participles}\BookHeading{nav-ch17-learner-0088}{subsubsection}{āfyllan: indicative and participles}{0}
% mdparagraph:P314:end

% teaching: NEW-U17-VOC-AFYLLAN-U17-SELECTED
\par\begin{minipage}{\linewidth}\raggedright
% mdparagraph:P315:begin
\textbf{āfyllan} — fill.
% mdparagraph:P315:end

% teaching: u17-afyllan-indicative12-participle-pair2
% mdparagraph:P316:begin
Double ll before a vowel, but use single l in āfylst, āfylð and āfylde. The selected past participle is āfylled; contracted āfyld is also recorded. The filled entity is accusative in an active clause and nominative in a passive; mid with the dative names the contents.
% mdparagraph:P316:end

% mdparagraph:P317:begin
\textbf{Indicative.}
% mdparagraph:P317:end

% table: u17-afyllan-indicative12-participle-pair2-INDICATIVE-DISPLAY
% mdtable:T86:begin
% course-table-style: 2026-10
% course-panel:T86-P01:begin
\par\begingroup
\fontsize{10.5}{13.2}\selectfont
\setlength{\LTleft}{0pt}\setlength{\LTright}{\fill}
\noindent
\begin{tabular}{@{}>{\raggedright\arraybackslash}p{39.459mm}>{\raggedright\arraybackslash}p{39.459mm}>{\raggedright\arraybackslash}p{39.459mm}@{}}
\toprule
\textbf{Person and number} & \textbf{Present} & \textbf{Preterite} \\
\midrule
1 sg. & āfylle & āfylde \\
2 sg. & āfylst & āfyldest \\
3 sg. & āfylð & āfylde \\
1–3 pl. & āfyllað & āfyldon \\
\bottomrule
\end{tabular}
\par\endgroup
% course-panel:T86-P01:end
% mdtable:T86:end
\end{minipage}\par
\reserve{9\baselineskip}
\reserve{7\baselineskip}
% mdparagraph:P318:begin
\textbf{Participle formation pair.}
% mdparagraph:P318:end

% table: u17-afyllan-indicative12-participle-pair2-PARTICIPLES-DISPLAY
% mdtable:T87:begin
\par\begingroup
\fontsize{10.5}{13.2}\selectfont
\setlength{\LTleft}{0pt}\setlength{\LTright}{\fill}
\noindent
\begin{tabular}{@{}>{\raggedright\arraybackslash}p{55.85mm}>{\raggedright\arraybackslash}p{65.34mm}@{}}
\toprule
\textbf{Form} & \textbf{Literal surface} \\
\midrule
present & āfyllende \\
past & āfylled \\
\bottomrule
\end{tabular}
\par\endgroup

% mdtable:T87:end
\reserve{7\baselineskip}
% mdparagraph:P319:begin
\subsection*{āweorpan: present and preterite indicative}\BookHeading{nav-ch17-learner-0089}{subsubsection}{āweorpan: present and preterite indicative}{0}
% mdparagraph:P319:end

% teaching: NEW-U17-VOC-AWEORPAN-U17-SELECTED
\par\begin{minipage}{\linewidth}\raggedright
% mdparagraph:P320:begin
\textbf{āweorpan} — cast down; cast away.
% mdparagraph:P320:end

% teaching: u17-aweorpan-indicative12
% mdparagraph:P321:begin
The strong stems alternate eo, ea, u, o. Present second and third singular have ie; the preterite second singular uses the same u-stem as the plural. The displaced person is accusative and of with the dative gives the source.
% mdparagraph:P321:end

% table: u17-aweorpan-indicative12-INDICATIVE-DISPLAY
% mdtable:T88:begin
% course-table-style: 2026-10
% course-panel:T88-P01:begin
\par\begingroup
\fontsize{10.5}{13.2}\selectfont
\setlength{\LTleft}{0pt}\setlength{\LTright}{\fill}
\noindent
\begin{tabular}{@{}>{\raggedright\arraybackslash}p{39.459mm}>{\raggedright\arraybackslash}p{39.459mm}>{\raggedright\arraybackslash}p{39.459mm}@{}}
\toprule
\textbf{Person and number} & \textbf{Present} & \textbf{Preterite} \\
\midrule
1 sg. & āweorpe & āwearp \\
2 sg. & āwierpst & āwurpe \\
3 sg. & āwierpð & āwearp \\
1–3 pl. & āweorpað & āwurpon \\
\bottomrule
\end{tabular}
\par\endgroup
% course-panel:T88-P01:end
% mdtable:T88:end
\end{minipage}\par
\reserve{7\baselineskip}
% mdparagraph:P322:begin
\subsection*{āwrītan: present and preterite indicative}\BookHeading{nav-ch17-learner-0090}{subsubsection}{āwrītan: present and preterite indicative}{0}
% mdparagraph:P322:end

% teaching: NEW-U17-VOC-AWRITAN-U17-SELECTED
\par\begin{minipage}{\linewidth}\raggedright
% mdparagraph:P323:begin
\textbf{āwrītan} — write.
% mdparagraph:P323:end

% teaching: u17-awritan-indicative12
% mdparagraph:P324:begin
The strong vowel sequence is ī, ā, i, i. Present second singular has -tst and the selected third singular -tt; a source variant has single -t. There is no new vowel mutation in these two present forms.
% mdparagraph:P324:end

% table: u17-awritan-indicative12-INDICATIVE-DISPLAY
% mdtable:T89:begin
% course-table-style: 2026-10
% course-panel:T89-P01:begin
\par\begingroup
\fontsize{10.5}{13.2}\selectfont
\setlength{\LTleft}{0pt}\setlength{\LTright}{\fill}
\noindent
\begin{tabular}{@{}>{\raggedright\arraybackslash}p{39.459mm}>{\raggedright\arraybackslash}p{39.459mm}>{\raggedright\arraybackslash}p{39.459mm}@{}}
\toprule
\textbf{Person and number} & \textbf{Present} & \textbf{Preterite} \\
\midrule
1 sg. & āwrīte & āwrāt \\
2 sg. & āwrītst & āwrite \\
3 sg. & āwrītt & āwrāt \\
1–3 pl. & āwrītað & āwriton \\
\bottomrule
\end{tabular}
\par\endgroup
% course-panel:T89-P01:end
% mdtable:T89:end
\end{minipage}\par
\reserve{7\baselineskip}
% mdparagraph:P325:begin
\subsection*{ālūtan: present and preterite indicative}\BookHeading{nav-ch17-learner-0091}{subsubsection}{ālūtan: present and preterite indicative}{0}
% mdparagraph:P325:end

% teaching: NEW-U17-VOC-ALUTAN-U17-SELECTED
\par\begin{minipage}{\linewidth}\raggedright
% mdparagraph:P326:begin
\textbf{ālūtan} — bow; incline oneself.
% mdparagraph:P326:end

% teaching: u17-alutan-indicative12
% mdparagraph:P327:begin
The strong stems alternate ū, ēa, u, o. Present second and third singular change ū to ȳ and have -tst and selected -tt; single -t is also recorded. The actor bows; dative Gode names the one bowed to.
% mdparagraph:P327:end

% table: u17-alutan-indicative12-INDICATIVE-DISPLAY
% mdtable:T90:begin
% course-table-style: 2026-10
% course-panel:T90-P01:begin
\par\begingroup
\fontsize{10.5}{13.2}\selectfont
\setlength{\LTleft}{0pt}\setlength{\LTright}{\fill}
\noindent
\begin{tabular}{@{}>{\raggedright\arraybackslash}p{39.459mm}>{\raggedright\arraybackslash}p{39.459mm}>{\raggedright\arraybackslash}p{39.459mm}@{}}
\toprule
\textbf{Person and number} & \textbf{Present} & \textbf{Preterite} \\
\midrule
1 sg. & ālūte & ālēat \\
2 sg. & ālȳtst & ālute \\
3 sg. & ālȳtt & ālēat \\
1–3 pl. & ālūtað & āluton \\
\bottomrule
\end{tabular}
\par\endgroup
% course-panel:T90-P01:end
% mdtable:T90:end
\end{minipage}\par
\reserve{7\baselineskip}
% mdparagraph:P328:begin
\subsection*{windan: present and preterite indicative}\BookHeading{nav-ch17-learner-0092}{subsubsection}{windan: present and preterite indicative}{0}
% mdparagraph:P328:end

% teaching: NEW-U17-VOC-WINDAN-U17-SELECTED
\par\begin{minipage}{\linewidth}\raggedright
% mdparagraph:P329:begin
\textbf{windan} — wind; plait.
% mdparagraph:P329:end

% teaching: u17-windan-indicative12
% mdparagraph:P330:begin
The strong stems alternate i, a, u, u. Stem d becomes t before the second- and third-singular present endings. The thing plaited is accusative; of with the dative names its material.
% mdparagraph:P330:end

% table: u17-windan-indicative12-INDICATIVE-DISPLAY
% mdtable:T91:begin
% course-table-style: 2026-10
% course-panel:T91-P01:begin
\par\begingroup
\fontsize{10.5}{13.2}\selectfont
\setlength{\LTleft}{0pt}\setlength{\LTright}{\fill}
\noindent
\begin{tabular}{@{}>{\raggedright\arraybackslash}p{39.459mm}>{\raggedright\arraybackslash}p{39.459mm}>{\raggedright\arraybackslash}p{39.459mm}@{}}
\toprule
\textbf{Person and number} & \textbf{Present} & \textbf{Preterite} \\
\midrule
1 sg. & winde & wand \\
2 sg. & wintst & wunde \\
3 sg. & wint & wand \\
1–3 pl. & windað & wundon \\
\bottomrule
\end{tabular}
\par\endgroup
% course-panel:T91-P01:end
% mdtable:T91:end
\end{minipage}\par
\reserve{7\baselineskip}
% mdparagraph:P331:begin
\subsection*{spǣtan: present and preterite indicative}\BookHeading{nav-ch17-learner-0093}{subsubsection}{spǣtan: present and preterite indicative}{0}
% mdparagraph:P331:end

% teaching: NEW-U17-VOC-SPAETAN-U17-SELECTED
\par\begin{minipage}{\linewidth}\raggedright
% mdparagraph:P332:begin
\textbf{spǣtan} — spit.
% mdparagraph:P332:end

% teaching: u17-spaetan-indicative12
% mdparagraph:P333:begin
Retain long ǣ. The t-stem has selected present third singular -tt and weak preterite -tte, -ttest, -tton. In the exercise, on takes the accusative person spat upon.
% mdparagraph:P333:end

% table: u17-spaetan-indicative12-INDICATIVE-DISPLAY
% mdtable:T92:begin
% course-table-style: 2026-10
% course-panel:T92-P01:begin
\par\begingroup
\fontsize{10.5}{13.2}\selectfont
\setlength{\LTleft}{0pt}\setlength{\LTright}{\fill}
\noindent
\begin{tabular}{@{}>{\raggedright\arraybackslash}p{39.459mm}>{\raggedright\arraybackslash}p{39.459mm}>{\raggedright\arraybackslash}p{39.459mm}@{}}
\toprule
\textbf{Person and number} & \textbf{Present} & \textbf{Preterite} \\
\midrule
1 sg. & spǣte & spǣtte \\
2 sg. & spǣtst & spǣttest \\
3 sg. & spǣtt & spǣtte \\
1–3 pl. & spǣtað & spǣtton \\
\bottomrule
\end{tabular}
\par\endgroup
% course-panel:T92-P01:end
% mdtable:T92:end
\end{minipage}\par
\reserve{7\baselineskip}
% mdparagraph:P334:begin
\subsection*{tǣlan: present and preterite indicative}\BookHeading{nav-ch17-learner-0094}{subsubsection}{tǣlan: present and preterite indicative}{0}
% mdparagraph:P334:end

% teaching: NEW-U17-VOC-TAELAN-U17-SELECTED
\par\begin{minipage}{\linewidth}\raggedright
% mdparagraph:P335:begin
\textbf{tǣlan} — blame; censure.
% mdparagraph:P335:end

% teaching: u17-taelan-indicative12
% mdparagraph:P336:begin
This long-stem weak verb has a preterite in -de. The person blamed is accusative. In the reading, tǣlst is second singular present indicative.
% mdparagraph:P336:end

% table: u17-taelan-indicative12-INDICATIVE-DISPLAY
% mdtable:T93:begin
% course-table-style: 2026-10
% course-panel:T93-P01:begin
\par\begingroup
\fontsize{10.5}{13.2}\selectfont
\setlength{\LTleft}{0pt}\setlength{\LTright}{\fill}
\noindent
\begin{tabular}{@{}>{\raggedright\arraybackslash}p{39.459mm}>{\raggedright\arraybackslash}p{39.459mm}>{\raggedright\arraybackslash}p{39.459mm}@{}}
\toprule
\textbf{Person and number} & \textbf{Present} & \textbf{Preterite} \\
\midrule
1 sg. & tǣle & tǣlde \\
2 sg. & tǣlst & tǣldest \\
3 sg. & tǣlð & tǣlde \\
1–3 pl. & tǣlað & tǣldon \\
\bottomrule
\end{tabular}
\par\endgroup
% course-panel:T93-P01:end
% mdtable:T93:end
\end{minipage}\par
\reserve{7\baselineskip}
% mdparagraph:P337:begin
\subsection*{bīegan: present and preterite indicative}\BookHeading{nav-ch17-learner-0095}{subsubsection}{bīegan: present and preterite indicative}{0}
% mdparagraph:P337:end

% teaching: NEW-U17-VOC-BIEGAN-U17-SELECTED
\par\begin{minipage}{\linewidth}\raggedright
% mdparagraph:P338:begin
\textbf{bīegan} — bend.
% mdparagraph:P338:end

% teaching: u17-biegan-indicative12
% mdparagraph:P339:begin
Use the selected long-diphthong stem bīeg- with weak endings. This transitive verb takes the bent knee as an accusative object. Keep it distinct from the intransitive bowing verb.
% mdparagraph:P339:end

% table: u17-biegan-indicative12-INDICATIVE-DISPLAY
% mdtable:T94:begin
% course-table-style: 2026-10
% course-panel:T94-P01:begin
\par\begingroup
\fontsize{10.5}{13.2}\selectfont
\setlength{\LTleft}{0pt}\setlength{\LTright}{\fill}
\noindent
\begin{tabular}{@{}>{\raggedright\arraybackslash}p{39.459mm}>{\raggedright\arraybackslash}p{39.459mm}>{\raggedright\arraybackslash}p{39.459mm}@{}}
\toprule
\textbf{Person and number} & \textbf{Present} & \textbf{Preterite} \\
\midrule
1 sg. & bīege & bīegde \\
2 sg. & bīegst & bīegdest \\
3 sg. & bīegð & bīegde \\
1–3 pl. & bīegað & bīegdon \\
\bottomrule
\end{tabular}
\par\endgroup
% course-panel:T94-P01:end
% mdtable:T94:end
\end{minipage}\par
\reserve{7\baselineskip}
% mdparagraph:P340:begin
\subsection*{gemunan: present and preterite indicative}\BookHeading{nav-ch17-learner-0096}{subsubsection}{gemunan: present and preterite indicative}{0}
% mdparagraph:P340:end

% teaching: NEW-U17-VOC-GEMUNAN-U17-SELECTED
\par\begin{minipage}{\linewidth}\raggedright
% mdparagraph:P341:begin
\textbf{gemunan} — remember.
% mdparagraph:P341:end

% teaching: u17-gemunan-indicative12
% mdparagraph:P342:begin
This preterite-present verb has present first/\allowbreak{}third singular geman, second singular gemanst and plural gemunon. Its dental preterite uses gemund-. The thing remembered is accusative.
% mdparagraph:P342:end

% table: u17-gemunan-indicative12-INDICATIVE-DISPLAY
% mdtable:T95:begin
% course-table-style: 2026-10
% course-panel:T95-P01:begin
\par\begingroup
\fontsize{10.5}{13.2}\selectfont
\setlength{\LTleft}{0pt}\setlength{\LTright}{\fill}
\noindent
\begin{tabular}{@{}>{\raggedright\arraybackslash}p{39.459mm}>{\raggedright\arraybackslash}p{39.459mm}>{\raggedright\arraybackslash}p{39.459mm}@{}}
\toprule
\textbf{Person and number} & \textbf{Present} & \textbf{Preterite} \\
\midrule
1 sg. & geman & gemunde \\
2 sg. & gemanst & gemundest \\
3 sg. & geman & gemunde \\
1–3 pl. & gemunon & gemundon \\
\bottomrule
\end{tabular}
\par\endgroup
% course-panel:T95-P01:end
% mdtable:T95:end
\end{minipage}\par
\reserve{7\baselineskip}
% mdparagraph:P343:begin
\subsection*{ðęncean: present and preterite indicative}\BookHeading{nav-ch17-learner-0097}{subsubsection}{ðęncean: present and preterite indicative}{0}
% mdparagraph:P343:end

% teaching: NEW-U17-VOC-THENCAN-U17-SELECTED
\par\begin{minipage}{\linewidth}\raggedright
% mdparagraph:P344:begin
\textbf{ðęncan} — think.
% mdparagraph:P344:end

% teaching: u17-thencan-indicative12
% mdparagraph:P345:begin
The selected headword variant is ðęncean. Its weak preterite ðōht- changes the vowel and replaces c before the dental ending. In the Sheba exchange, a free-choice expression gives what the queen had thought.
% mdparagraph:P345:end

% table: u17-thencan-indicative12-INDICATIVE-DISPLAY
% mdtable:T96:begin
% course-table-style: 2026-10
% course-panel:T96-P01:begin
\par\begingroup
\fontsize{10.5}{13.2}\selectfont
\setlength{\LTleft}{0pt}\setlength{\LTright}{\fill}
\noindent
\begin{tabular}{@{}>{\raggedright\arraybackslash}p{39.459mm}>{\raggedright\arraybackslash}p{39.459mm}>{\raggedright\arraybackslash}p{39.459mm}@{}}
\toprule
\textbf{Person and number} & \textbf{Present} & \textbf{Preterite} \\
\midrule
1 sg. & ðęnce & ðōhte \\
2 sg. & ðęncst & ðōhtest \\
3 sg. & ðęncð & ðōhte \\
1–3 pl. & ðęnceað & ðōhton \\
\bottomrule
\end{tabular}
\par\endgroup
% course-panel:T96-P01:end
% mdtable:T96:end
\end{minipage}\par
\reserve{7\baselineskip}
% mdparagraph:P346:begin
\subsection*{blētsian: present and past participles}\BookHeading{nav-ch17-learner-0098}{subsubsection}{blētsian: present and past participles}{0}
% mdparagraph:P346:end

% teaching: NEW-U17-VOC-BLETSIAN-U17-BLESS
\par\begin{minipage}{\linewidth}\raggedright
% mdparagraph:P347:begin
\textbf{blētsian} — bless.
% mdparagraph:P347:end

% teaching: u17-bletsian-participle-pair2
% mdparagraph:P348:begin
The present participle has -iende and the past participle -od in this weak -ian pattern.
% mdparagraph:P348:end

% mdparagraph:P349:begin
This display assigns the complete present/\allowbreak{}past formation pair. It does not add other finite moods or an unrestricted attributive declension; the past form is used in the stated passive frame.
% mdparagraph:P349:end

% table: u17-bletsian-participle-pair2-PARTICIPLES-DISPLAY
% mdtable:T97:begin
\par\begingroup
\fontsize{10.5}{13.2}\selectfont
\setlength{\LTleft}{0pt}\setlength{\LTright}{\fill}
\noindent
\begin{tabular}{@{}>{\raggedright\arraybackslash}p{54.47mm}>{\raggedright\arraybackslash}p{66.72mm}@{}}
\toprule
\textbf{Form} & \textbf{Literal surface} \\
\midrule
present & blētsiende \\
past & blētsod \\
\bottomrule
\end{tabular}
\par\endgroup

% mdtable:T97:end
\end{minipage}\par
\reserve{7\baselineskip}
% mdparagraph:P350:begin
\subsection*{ofwundrian: present and past participles}\BookHeading{nav-ch17-learner-0099}{subsubsection}{ofwundrian: present and past participles}{0}
% mdparagraph:P350:end

% teaching: NEW-U17-VOC-OFWUNDRIAN-U17-ASTONISHED
\par\begin{minipage}{\linewidth}\raggedright
% mdparagraph:P351:begin
\textbf{ofwundrian} — be astonished.
% mdparagraph:P351:end

% teaching: u17-ofwundrian-participle-pair2
% mdparagraph:P352:begin
The present participle has -iende and the past participle -od in this weak -ian pattern.
% mdparagraph:P352:end

% mdparagraph:P353:begin
This display assigns the complete present/\allowbreak{}past formation pair. It does not add other finite moods or an unrestricted attributive declension; the past form is used in the stated passive frame.
% mdparagraph:P353:end

% table: u17-ofwundrian-participle-pair2-PARTICIPLES-DISPLAY
% mdtable:T98:begin
\par\begingroup
\fontsize{10.5}{13.2}\selectfont
\setlength{\LTleft}{0pt}\setlength{\LTright}{\fill}
\noindent
\begin{tabular}{@{}>{\raggedright\arraybackslash}p{51.9mm}>{\raggedright\arraybackslash}p{69.29mm}@{}}
\toprule
\textbf{Form} & \textbf{Literal surface} \\
\midrule
present & ofwundriende \\
past & ofwundrod \\
\bottomrule
\end{tabular}
\par\endgroup

% mdtable:T98:end
\end{minipage}\par
\reserve{7\baselineskip}
% mdparagraph:P354:begin
\subsection*{wyrcean: present and past participles}\BookHeading{nav-ch17-learner-0100}{subsubsection}{wyrcean: present and past participles}{0}
% mdparagraph:P354:end

% teaching: u17-wyrcean-participle-pair2
% mdparagraph:P355:begin
The present participle uses the present stem; the past participle is worht, with the same changed stem as the weak preterite.
% mdparagraph:P355:end

\reserve{10\baselineskip}
% mdparagraph:P356:begin
This display assigns the complete present/\allowbreak{}past formation pair. It does not add other finite moods or an unrestricted attributive declension; the past form is used in the stated passive frame.
% mdparagraph:P356:end

% table: u17-wyrcean-participle-pair2-PARTICIPLES-DISPLAY
% mdtable:T99:begin
\par\begingroup
\fontsize{10.5}{13.2}\selectfont
\setlength{\LTleft}{0pt}\setlength{\LTright}{\fill}
\noindent
\begin{tabular}{@{}>{\raggedright\arraybackslash}p{57.3mm}>{\raggedright\arraybackslash}p{63.89mm}@{}}
\toprule
\textbf{Form} & \textbf{Literal surface} \\
\midrule
present & wyrcende \\
past & worht \\
\bottomrule
\end{tabular}
\par\endgroup

% mdtable:T99:end
\reserve{7\baselineskip}
% mdparagraph:P357:begin
\subsection*{gesweotolian: present and past participles}\BookHeading{nav-ch17-learner-0101}{subsubsection}{gesweotolian: present and past participles}{0}
% mdparagraph:P357:end

% teaching: NEW-U17-VOC-GESWEOTOLIAN-U17-MANIFEST
\par\begin{minipage}{\linewidth}\raggedright
% mdparagraph:P358:begin
\textbf{gesweotolian} — make manifest; reveal.
% mdparagraph:P358:end

% teaching: u17-gesweotolian-participle-pair2
% mdparagraph:P359:begin
The present participle has -iende and the past participle -od in this weak -ian pattern.
% mdparagraph:P359:end

% mdparagraph:P360:begin
This display assigns the complete present/\allowbreak{}past formation pair. It does not add other finite moods or an unrestricted attributive declension; the past form is used in the stated passive frame.
% mdparagraph:P360:end

% table: u17-gesweotolian-participle-pair2-PARTICIPLES-DISPLAY
% mdtable:T100:begin
\par\begingroup
\fontsize{10.5}{13.2}\selectfont
\setlength{\LTleft}{0pt}\setlength{\LTright}{\fill}
\noindent
\begin{tabular}{@{}>{\raggedright\arraybackslash}p{49.57mm}>{\raggedright\arraybackslash}p{71.62mm}@{}}
\toprule
\textbf{Form} & \textbf{Literal surface} \\
\midrule
present & gesweotoliende \\
past & gesweotolod \\
\bottomrule
\end{tabular}
\par\endgroup

% mdtable:T100:end
\end{minipage}\par
\reserve{7\baselineskip}
% mdparagraph:P361:begin
\subsection*{ācęnnan: present and past participles}\BookHeading{nav-ch17-learner-0102}{subsubsection}{ācęnnan: present and past participles}{0}
% mdparagraph:P361:end

% teaching: NEW-U17-VOC-ACENNAN-U17-BEGET
\par\begin{minipage}{\linewidth}\raggedright
% mdparagraph:P362:begin
\textbf{ācęnnan} — beget; bring forth.
% mdparagraph:P362:end

% teaching: u17-acennan-participle-pair2
% mdparagraph:P363:begin
The present participle keeps the present stem with nn; the past participle has -ed. The begetting expression below is supplied in full.
% mdparagraph:P363:end

% mdparagraph:P364:begin
This display assigns the complete present/\allowbreak{}past formation pair. It does not add other finite moods or an unrestricted attributive declension; the past form is used in the stated passive frame.
% mdparagraph:P364:end

% table: u17-acennan-participle-pair2-PARTICIPLES-DISPLAY
% mdtable:T101:begin
\par\begingroup
\fontsize{10.5}{13.2}\selectfont
\setlength{\LTleft}{0pt}\setlength{\LTright}{\fill}
\noindent
\begin{tabular}{@{}>{\raggedright\arraybackslash}p{55.85mm}>{\raggedright\arraybackslash}p{65.34mm}@{}}
\toprule
\textbf{Form} & \textbf{Literal surface} \\
\midrule
present & ācęnnende \\
past & ācęnned \\
\bottomrule
\end{tabular}
\par\endgroup

% mdtable:T101:end
\end{minipage}\par
\reserve{7\baselineskip}
% mdparagraph:P365:begin
\subsection*{geopenian: present and past participles}\BookHeading{nav-ch17-learner-0103}{subsubsection}{geopenian: present and past participles}{0}
% mdparagraph:P365:end

% teaching: u17-geopenian-participle-pair2
% mdparagraph:P366:begin
The present participle has -iende and the past participle -od in this weak -ian pattern.
% mdparagraph:P366:end

\reserve{10\baselineskip}
% mdparagraph:P367:begin
This display assigns the complete present/\allowbreak{}past formation pair. It does not add other finite moods or an unrestricted attributive declension; the past form is used in the stated passive frame.
% mdparagraph:P367:end

% table: u17-geopenian-participle-pair2-PARTICIPLES-DISPLAY
% mdtable:T102:begin
\par\begingroup
\fontsize{10.5}{13.2}\selectfont
\setlength{\LTleft}{0pt}\setlength{\LTright}{\fill}
\noindent
\begin{tabular}{@{}>{\raggedright\arraybackslash}p{53.15mm}>{\raggedright\arraybackslash}p{68.04mm}@{}}
\toprule
\textbf{Form} & \textbf{Literal surface} \\
\midrule
present & geopeniende \\
past & geopenod \\
\bottomrule
\end{tabular}
\par\endgroup

% mdtable:T102:end
\reserve{7\baselineskip}
% mdparagraph:P368:begin
\subsection*{cnucian: singular and plural commands}\BookHeading{nav-ch17-learner-0104}{subsubsection}{cnucian: singular and plural commands}{0}
% mdparagraph:P368:end

% teaching: NEW-U17-VOC-CNUCIAN-U17-KNOCK
\par\begin{minipage}{\linewidth}\raggedright
% mdparagraph:P369:begin
\textbf{cnucian} — knock.
% mdparagraph:P369:end

% teaching: u17-cnucian-imperative2
% mdparagraph:P370:begin
The weak -ian command has singular -a and plural -iað. Only this complete command pair is added here.
% mdparagraph:P370:end

% table: u17-cnucian-imperative2-IMPERATIVE-DISPLAY
% mdtable:T103:begin
% course-table-style: 2026-10
% course-panel:T103-P01:begin
\par\begingroup
\fontsize{10.5}{13.2}\selectfont
\setlength{\LTleft}{0pt}\setlength{\LTright}{\fill}
\noindent
\begin{tabular}{@{}>{\raggedright\arraybackslash}p{63.342mm}>{\raggedright\arraybackslash}p{57.846mm}@{}}
\toprule
\textbf{Person and number} & \textbf{Command} \\
\midrule
2 sg. & cnuca \\
2 pl. & cnuciað \\
\bottomrule
\end{tabular}
\par\endgroup
% course-panel:T103-P01:end
% mdtable:T103:end
\end{minipage}\par
\reserve{18\baselineskip}
% mdparagraph:P371:begin
\subsection*{āhębban: singular and plural commands}\BookHeading{nav-ch17-learner-0105}{subsubsection}{āhębban: singular and plural commands}{0}
% mdparagraph:P371:end

% teaching: u17-ahebban-imperative2
% mdparagraph:P372:begin
The strong j-present has singular āhęfe and plural āhębbað. Only this complete command pair is added here.
% mdparagraph:P372:end

% table: u17-ahebban-imperative2-IMPERATIVE-DISPLAY
% mdtable:T104:begin
% course-table-style: 2026-10
% course-panel:T104-P01:begin
\par\begingroup
\fontsize{10.5}{13.2}\selectfont
\setlength{\LTleft}{0pt}\setlength{\LTright}{\fill}
\noindent
\begin{tabular}{@{}>{\raggedright\arraybackslash}p{63.342mm}>{\raggedright\arraybackslash}p{57.846mm}@{}}
\toprule
\textbf{Person and number} & \textbf{Command} \\
\midrule
2 sg. & āhęfe \\
2 pl. & āhębbað \\
\bottomrule
\end{tabular}
\par\endgroup
% course-panel:T104-P01:end
% mdtable:T104:end
\reserve{18\baselineskip}
% mdparagraph:P373:begin
\subsection*{plegian: singular and plural commands}\BookHeading{nav-ch17-learner-0106}{subsubsection}{plegian: singular and plural commands}{0}
% mdparagraph:P373:end

% teaching: u17-plegian-imperative2
% mdparagraph:P374:begin
The weak -ian command has singular -a and plural -iað. The assigned ordinary sense is play, including playing with a ball. Only this complete command pair is added here.
% mdparagraph:P374:end

% table: u17-plegian-imperative2-IMPERATIVE-DISPLAY
% mdtable:T105:begin
% course-table-style: 2026-10
% course-panel:T105-P01:begin
\par\begingroup
\fontsize{10.5}{13.2}\selectfont
\setlength{\LTleft}{0pt}\setlength{\LTright}{\fill}
\noindent
\begin{tabular}{@{}>{\raggedright\arraybackslash}p{63.342mm}>{\raggedright\arraybackslash}p{57.846mm}@{}}
\toprule
\textbf{Person and number} & \textbf{Command} \\
\midrule
2 sg. & plega \\
2 pl. & plegiað \\
\bottomrule
\end{tabular}
\par\endgroup
% course-panel:T105-P01:end
% mdtable:T105:end
\reserve{18\baselineskip}
% mdparagraph:P375:begin
\subsection*{cēosan: singular and plural commands}\BookHeading{nav-ch17-learner-0107}{subsubsection}{cēosan: singular and plural commands}{0}
% mdparagraph:P375:end

% teaching: u17-ceosan-imperative2
% mdparagraph:P376:begin
Use the unmutated present stem in the singular command and add -að for the plural. Only this complete command pair is added here.
% mdparagraph:P376:end

% table: u17-ceosan-imperative2-IMPERATIVE-DISPLAY
% mdtable:T106:begin
% course-table-style: 2026-10
% course-panel:T106-P01:begin
\par\begingroup
\fontsize{10.5}{13.2}\selectfont
\setlength{\LTleft}{0pt}\setlength{\LTright}{\fill}
\noindent
\begin{tabular}{@{}>{\raggedright\arraybackslash}p{63.342mm}>{\raggedright\arraybackslash}p{57.846mm}@{}}
\toprule
\textbf{Person and number} & \textbf{Command} \\
\midrule
2 sg. & cēos \\
2 pl. & cēosað \\
\bottomrule
\end{tabular}
\par\endgroup
% course-panel:T106-P01:end
% mdtable:T106:end
\reserve{7\baselineskip}
% mdparagraph:P377:begin
\subsection*{tōcuman: present subjunctive}\BookHeading{nav-ch17-learner-0108}{subsubsection}{tōcuman: present subjunctive}{0}
% mdparagraph:P377:end

% teaching: NEW-U17-VOC-TOCUMAN-U17-COME
\par\begin{minipage}{\linewidth}\raggedright
% mdparagraph:P378:begin
\textbf{tōcuman} — come; arrive.
% mdparagraph:P378:end

% teaching: u17-tocuman-present-subjunctive6
% mdparagraph:P379:begin
All three singular persons use -e and all three plural persons -en on the present stem, without an indicative second/\allowbreak{}third-singular vowel change. Only this tense of the subjunctive is added here.
% mdparagraph:P379:end

% table: u17-tocuman-present-subjunctive6-SUBJUNCTIVE-DISPLAY
% mdtable:T107:begin
% course-table-style: 2026-10
% course-panel:T107-P01:begin
\par\begingroup
\fontsize{10.5}{13.2}\selectfont
\setlength{\LTleft}{0pt}\setlength{\LTright}{\fill}
\noindent
\begin{tabular}{@{}>{\raggedright\arraybackslash}p{63.342mm}>{\raggedright\arraybackslash}p{57.846mm}@{}}
\toprule
\textbf{Person and number} & \textbf{Present} \\
\midrule
1 sg. & tōcume \\
2 sg. & tōcume \\
3 sg. & tōcume \\
1–3 pl. & tōcumen \\
\bottomrule
\end{tabular}
\par\endgroup
% course-panel:T107-P01:end
% mdtable:T107:end
\end{minipage}\par
\reserve{7\baselineskip}
% mdparagraph:P380:begin
\subsection*{geweorðan: present subjunctive}\BookHeading{nav-ch17-learner-0109}{subsubsection}{geweorðan: present subjunctive}{0}
% mdparagraph:P380:end

% teaching: NEW-U17-VOC-GEWEORTHAN-U17-BEFALL
\par\begin{minipage}{\linewidth}\raggedright
% mdparagraph:P381:begin
\textbf{geweorðan} — come to pass; befall.
% mdparagraph:P381:end

% teaching: u17-geweorthan-present-subjunctive6
% mdparagraph:P382:begin
All three singular persons use -e and all three plural persons -en on the present stem, without an indicative second/\allowbreak{}third-singular vowel change. Only this tense of the subjunctive is added here.
% mdparagraph:P382:end

% table: u17-geweorthan-present-subjunctive6-SUBJUNCTIVE-DISPLAY
% mdtable:T108:begin
% course-table-style: 2026-10
% course-panel:T108-P01:begin
\par\begingroup
\fontsize{10.5}{13.2}\selectfont
\setlength{\LTleft}{0pt}\setlength{\LTright}{\fill}
\noindent
\begin{tabular}{@{}>{\raggedright\arraybackslash}p{63.342mm}>{\raggedright\arraybackslash}p{57.846mm}@{}}
\toprule
\textbf{Person and number} & \textbf{Present} \\
\midrule
1 sg. & geweorðe \\
2 sg. & geweorðe \\
3 sg. & geweorðe \\
1–3 pl. & geweorðen \\
\bottomrule
\end{tabular}
\par\endgroup
% course-panel:T108-P01:end
% mdtable:T108:end
\end{minipage}\par
\reserve{18\baselineskip}
% mdparagraph:P383:begin
\subsection*{leornian: present subjunctive}\BookHeading{nav-ch17-learner-0110}{subsubsection}{leornian: present subjunctive}{0}
% mdparagraph:P383:end

% teaching: u17-leornian-present-subjunctive6
% mdparagraph:P384:begin
All three singular persons use -e and all three plural persons -en on the present stem; there is no indicative second/\allowbreak{}third-singular vowel change. With leornian, retain i in leornie and leornien.
% mdparagraph:P384:end

% table: u17-leornian-present-subjunctive6-SUBJUNCTIVE-DISPLAY
% mdtable:T109:begin
% course-table-style: 2026-10
% course-panel:T109-P01:begin
\par\begingroup
\fontsize{10.5}{13.2}\selectfont
\setlength{\LTleft}{0pt}\setlength{\LTright}{\fill}
\noindent
\begin{tabular}{@{}>{\raggedright\arraybackslash}p{63.342mm}>{\raggedright\arraybackslash}p{57.846mm}@{}}
\toprule
\textbf{Person and number} & \textbf{Present} \\
\midrule
1 sg. & leornie \\
2 sg. & leornie \\
3 sg. & leornie \\
1–3 pl. & leornien \\
\bottomrule
\end{tabular}
\par\endgroup
% course-panel:T109-P01:end
% mdtable:T109:end
\reserve{7\baselineskip}
% mdparagraph:P385:begin
\subsection*{āhyldan: present subjunctive}\BookHeading{nav-ch17-learner-0111}{subsubsection}{āhyldan: present subjunctive}{0}
% mdparagraph:P385:end

% teaching: NEW-U17-VOC-AHYLDAN-U17-LAY
\par\begin{minipage}{\linewidth}\raggedright
% mdparagraph:P386:begin
\textbf{āhyldan} — lay; incline.
% mdparagraph:P386:end

% teaching: u17-ahyldan-present-subjunctive6
% mdparagraph:P387:begin
All three singular persons use -e and all three plural persons -en on the present stem, without an indicative second/\allowbreak{}third-singular vowel change. Only this tense of the subjunctive is added here.
% mdparagraph:P387:end

% table: u17-ahyldan-present-subjunctive6-SUBJUNCTIVE-DISPLAY
% mdtable:T110:begin
% course-table-style: 2026-10
% course-panel:T110-P01:begin
\par\begingroup
\fontsize{10.5}{13.2}\selectfont
\setlength{\LTleft}{0pt}\setlength{\LTright}{\fill}
\noindent
\begin{tabular}{@{}>{\raggedright\arraybackslash}p{63.342mm}>{\raggedright\arraybackslash}p{57.846mm}@{}}
\toprule
\textbf{Person and number} & \textbf{Present} \\
\midrule
1 sg. & āhylde \\
2 sg. & āhylde \\
3 sg. & āhylde \\
1–3 pl. & āhylden \\
\bottomrule
\end{tabular}
\par\endgroup
% course-panel:T110-P01:end
% mdtable:T110:end
\end{minipage}\par
\reserve{7\baselineskip}
% page-review repair: U17LF02: participial agreement introduction with first complete singular panel
\par\begin{minipage}{\linewidth}\raggedright
% mdparagraph:P388:begin
\subsection*{ondrǣdan: present participle: strong and weak agreement}\BookHeading{nav-ch17-learner-0112}{subsubsection}{ondrǣdan: present participle: strong and weak agreement}{0}
% mdparagraph:P388:end

% teaching: u17-ondraedan-present-participle-agreement48
% mdparagraph:P389:begin
Remove final -e from the participle before adding an overt adjective ending. The strong table follows the grēne type: selected feminine nominative singular and neuter nominative/\allowbreak{}accusative plural -u, and feminine plural -a; the corresponding -e variants are also recorded. The weak forms follow the ordinary weak adjective pattern. The new independent use is describing or naming the specified human groups; no complete personal finite hyngran system is assigned.
% mdparagraph:P389:end

% mdparagraph:P390:begin
\textbf{Strong.}
% mdparagraph:P390:end

% table: u17-ondraedan-present-participle-agreement48-STRONG-DISPLAY
% mdtable:T111:begin
% course-table-style: 2026-10
% course-panel:T111-P01:begin
\par\begingroup
\fontsize{10.5}{13.2}\selectfont
\setlength{\LTleft}{0pt}\setlength{\LTright}{\fill}
\noindent
\begin{tabular}{@{}>{\raggedright\arraybackslash}p{22.0mm}>{\raggedright\arraybackslash}p{31.188mm}>{\raggedright\arraybackslash}p{31.188mm}>{\raggedright\arraybackslash}p{31.188mm}@{}}
\toprule
\textbf{Case} & \textbf{Masculine singular} & \textbf{Neuter singular} & \textbf{Feminine singular} \\
\midrule
Nominative & ondrǣdende & ondrǣdende & ondrǣdendu \\
Genitive & ondrǣdendes & ondrǣdendes & ondrǣdendre \\
Dative & ondrǣdendum & ondrǣdendum & ondrǣdendre \\
Accusative & ondrǣdendne & ondrǣdende & ondrǣdende \\
\bottomrule
\end{tabular}
\par\endgroup
% course-panel:T111-P01:end
\end{minipage}\par
\par\addvspace{5pt}
% course-panel:T111-P02:begin
\par\begingroup
\fontsize{10.5}{13.2}\selectfont
\setlength{\LTleft}{0pt}\setlength{\LTright}{\fill}
\noindent
\begin{tabular}{@{}>{\raggedright\arraybackslash}p{22.0mm}>{\raggedright\arraybackslash}p{31.188mm}>{\raggedright\arraybackslash}p{31.188mm}>{\raggedright\arraybackslash}p{31.188mm}@{}}
\toprule
\textbf{Case} & \textbf{Masculine plural} & \textbf{Neuter plural} & \textbf{Feminine plural} \\
\midrule
Nominative & ondrǣdende & ondrǣdendu & ondrǣdenda \\
Genitive & ondrǣdendra & ondrǣdendra & ondrǣdendra \\
Dative & ondrǣdendum & ondrǣdendum & ondrǣdendum \\
Accusative & ondrǣdende & ondrǣdendu & ondrǣdenda \\
\bottomrule
\end{tabular}
\par\endgroup
% course-panel:T111-P02:end
% mdtable:T111:end

\reserve{18\baselineskip}
\reserve{16\baselineskip}
% mdparagraph:P391:begin
\textbf{Weak.}
% mdparagraph:P391:end

% table: u17-ondraedan-present-participle-agreement48-WEAK-DISPLAY
% mdtable:T112:begin
% course-table-style: 2026-10
% course-panel:T112-P01:begin
\par\begingroup
\fontsize{10.5}{13.2}\selectfont
\setlength{\LTleft}{0pt}\setlength{\LTright}{\fill}
\noindent
\begin{tabular}{@{}>{\raggedright\arraybackslash}p{22.0mm}>{\raggedright\arraybackslash}p{31.188mm}>{\raggedright\arraybackslash}p{31.188mm}>{\raggedright\arraybackslash}p{31.188mm}@{}}
\toprule
\textbf{Case} & \textbf{Masculine singular} & \textbf{Neuter singular} & \textbf{Feminine singular} \\
\midrule
Nominative & ondrǣdenda & ondrǣdende & ondrǣdende \\
Genitive & ondrǣdendan & ondrǣdendan & ondrǣdendan \\
Dative & ondrǣdendan & ondrǣdendan & ondrǣdendan \\
Accusative & ondrǣdendan & ondrǣdende & ondrǣdendan \\
\bottomrule
\end{tabular}
\par\endgroup
% course-panel:T112-P01:end
\par\addvspace{5pt}
% course-panel:T112-P02:begin
\par\begingroup
\fontsize{10.5}{13.2}\selectfont
\setlength{\LTleft}{0pt}\setlength{\LTright}{\fill}
\noindent
\begin{tabular}{@{}>{\raggedright\arraybackslash}p{22.0mm}>{\raggedright\arraybackslash}p{31.188mm}>{\raggedright\arraybackslash}p{31.188mm}>{\raggedright\arraybackslash}p{31.188mm}@{}}
\toprule
\textbf{Case} & \textbf{Masculine plural} & \textbf{Neuter plural} & \textbf{Feminine plural} \\
\midrule
Nominative & ondrǣdendan & ondrǣdendan & ondrǣdendan \\
Genitive & ondrǣdendra & ondrǣdendra & ondrǣdendra \\
Dative & ondrǣdendum & ondrǣdendum & ondrǣdendum \\
Accusative & ondrǣdendan & ondrǣdendan & ondrǣdendan \\
\bottomrule
\end{tabular}
\par\endgroup
% course-panel:T112-P02:end
% mdtable:T112:end

\reserve{7\baselineskip}
% mdparagraph:P392:begin
\subsection*{hyngran: present participle: strong and weak agreement}\BookHeading{nav-ch17-learner-0113}{subsubsection}{hyngran: present participle: strong and weak agreement}{0}
% mdparagraph:P392:end

% teaching: NEW-U17-VOC-HYNGRAN-U17-HUNGRY-PARTICIPLE
\par\begin{minipage}{\linewidth}\raggedright
% mdparagraph:P393:begin
\textbf{hyngran} — being hungry; hungering.
% mdparagraph:P393:end

% teaching: u17-hyngran-present-participle-agreement48
% mdparagraph:P394:begin
Remove final -e from the participle before adding an overt adjective ending. The strong table follows the grēne type: selected feminine nominative singular and neuter nominative/\allowbreak{}accusative plural -u, and feminine plural -a; the corresponding -e variants are also recorded. The weak forms follow the ordinary weak adjective pattern. The new independent use is describing or naming the specified human groups; no complete personal finite hyngran system is assigned.
% mdparagraph:P394:end

% mdparagraph:P395:begin
\textbf{Strong.}
% mdparagraph:P395:end

% table: u17-hyngran-present-participle-agreement48-STRONG-DISPLAY
% mdtable:T113:begin
% course-table-style: 2026-10
% course-panel:T113-P01:begin
\par\begingroup
\fontsize{10.5}{13.2}\selectfont
\setlength{\LTleft}{0pt}\setlength{\LTright}{\fill}
\noindent
\begin{tabular}{@{}>{\raggedright\arraybackslash}p{22.0mm}>{\raggedright\arraybackslash}p{31.188mm}>{\raggedright\arraybackslash}p{31.188mm}>{\raggedright\arraybackslash}p{31.188mm}@{}}
\toprule
\textbf{Case} & \textbf{Masculine singular} & \textbf{Neuter singular} & \textbf{Feminine singular} \\
\midrule
Nominative & hyngrende & hyngrende & hyngrendu \\
Genitive & hyngrendes & hyngrendes & hyngrendre \\
Dative & hyngrendum & hyngrendum & hyngrendre \\
Accusative & hyngrendne & hyngrende & hyngrende \\
\bottomrule
\end{tabular}
\par\endgroup
% course-panel:T113-P01:end
\end{minipage}\par
\par\addvspace{5pt}
% course-panel:T113-P02:begin
\par\begingroup
\fontsize{10.5}{13.2}\selectfont
\setlength{\LTleft}{0pt}\setlength{\LTright}{\fill}
\noindent
\begin{tabular}{@{}>{\raggedright\arraybackslash}p{22.0mm}>{\raggedright\arraybackslash}p{31.188mm}>{\raggedright\arraybackslash}p{31.188mm}>{\raggedright\arraybackslash}p{31.188mm}@{}}
\toprule
\textbf{Case} & \textbf{Masculine plural} & \textbf{Neuter plural} & \textbf{Feminine plural} \\
\midrule
Nominative & hyngrende & hyngrendu & hyngrenda \\
Genitive & hyngrendra & hyngrendra & hyngrendra \\
Dative & hyngrendum & hyngrendum & hyngrendum \\
Accusative & hyngrende & hyngrendu & hyngrenda \\
\bottomrule
\end{tabular}
\par\endgroup
% course-panel:T113-P02:end
% mdtable:T113:end

\reserve{18\baselineskip}
\reserve{16\baselineskip}
% mdparagraph:P396:begin
\textbf{Weak.}
% mdparagraph:P396:end

% table: u17-hyngran-present-participle-agreement48-WEAK-DISPLAY
% mdtable:T114:begin
% course-table-style: 2026-10
% course-panel:T114-P01:begin
\par\begingroup
\fontsize{10.5}{13.2}\selectfont
\setlength{\LTleft}{0pt}\setlength{\LTright}{\fill}
\noindent
\begin{tabular}{@{}>{\raggedright\arraybackslash}p{22.0mm}>{\raggedright\arraybackslash}p{31.188mm}>{\raggedright\arraybackslash}p{31.188mm}>{\raggedright\arraybackslash}p{31.188mm}@{}}
\toprule
\textbf{Case} & \textbf{Masculine singular} & \textbf{Neuter singular} & \textbf{Feminine singular} \\
\midrule
Nominative & hyngrenda & hyngrende & hyngrende \\
Genitive & hyngrendan & hyngrendan & hyngrendan \\
Dative & hyngrendan & hyngrendan & hyngrendan \\
Accusative & hyngrendan & hyngrende & hyngrendan \\
\bottomrule
\end{tabular}
\par\endgroup
% course-panel:T114-P01:end
\par\addvspace{5pt}
% course-panel:T114-P02:begin
\par\begingroup
\fontsize{10.5}{13.2}\selectfont
\setlength{\LTleft}{0pt}\setlength{\LTright}{\fill}
\noindent
\begin{tabular}{@{}>{\raggedright\arraybackslash}p{22.0mm}>{\raggedright\arraybackslash}p{31.188mm}>{\raggedright\arraybackslash}p{31.188mm}>{\raggedright\arraybackslash}p{31.188mm}@{}}
\toprule
\textbf{Case} & \textbf{Masculine plural} & \textbf{Neuter plural} & \textbf{Feminine plural} \\
\midrule
Nominative & hyngrendan & hyngrendan & hyngrendan \\
Genitive & hyngrendra & hyngrendra & hyngrendra \\
Dative & hyngrendum & hyngrendum & hyngrendum \\
Accusative & hyngrendan & hyngrendan & hyngrendan \\
\bottomrule
\end{tabular}
\par\endgroup
% course-panel:T114-P02:end
% mdtable:T114:end

\reserve{7\baselineskip}
% mdparagraph:P397:begin
\subsection*{gesēon: preterite subjunctive}\BookHeading{nav-ch17-learner-0114}{subsubsection}{gesēon: preterite subjunctive}{0}
% mdparagraph:P397:end

% teaching: NEW-U17-VOC-GESEON-U17-REGARD
\par\begin{minipage}{\linewidth}\raggedright
% mdparagraph:P398:begin
\textbf{gesēon} — regard; take notice of.
% mdparagraph:P398:end

% teaching: u17-geseon-preterite-subjunctive6
% mdparagraph:P399:begin
Build the past subjunctive on the plural preterite stem gesāw-. All three singular persons have -e and all three plural persons -en. This system supplies the before-clause’s first singular gesāwe.
% mdparagraph:P399:end

% table: u17-geseon-preterite-subjunctive6-SUBJUNCTIVE-DISPLAY
% mdtable:T115:begin
% course-table-style: 2026-10
% course-panel:T115-P01:begin
\par\begingroup
\fontsize{10.5}{13.2}\selectfont
\setlength{\LTleft}{0pt}\setlength{\LTright}{\fill}
\noindent
\begin{tabular}{@{}>{\raggedright\arraybackslash}p{63.342mm}>{\raggedright\arraybackslash}p{57.846mm}@{}}
\toprule
\textbf{Person and number} & \textbf{Preterite} \\
\midrule
1 sg. & gesāwe \\
2 sg. & gesāwe \\
3 sg. & gesāwe \\
1–3 pl. & gesāwen \\
\bottomrule
\end{tabular}
\par\endgroup
% course-panel:T115-P01:end
% mdtable:T115:end
\end{minipage}\par
\reserve{7\baselineskip}
% mdparagraph:P400:begin
\subsection*{sęttan: present subjunctive}\BookHeading{nav-ch17-learner-0115}{subsubsection}{sęttan: present subjunctive}{0}
% mdparagraph:P400:end

% teaching: NEW-U17-VOC-SETTAN-U17-ESTABLISH
\par\begin{minipage}{\linewidth}\raggedright
% mdparagraph:P401:begin
\textbf{sęttan} — appoint; establish; make a decree.
% mdparagraph:P401:end

% teaching: u17-settan-present-subjunctive6
% mdparagraph:P402:begin
All three singular persons use -e and all three plural persons -en on the present stem, without an indicative second/\allowbreak{}third-singular vowel change. Only this tense of the subjunctive is added here.
% mdparagraph:P402:end

% table: u17-settan-present-subjunctive6-SUBJUNCTIVE-DISPLAY
% mdtable:T116:begin
% course-table-style: 2026-10
% course-panel:T116-P01:begin
\par\begingroup
\fontsize{10.5}{13.2}\selectfont
\setlength{\LTleft}{0pt}\setlength{\LTright}{\fill}
\noindent
\begin{tabular}{@{}>{\raggedright\arraybackslash}p{63.342mm}>{\raggedright\arraybackslash}p{57.846mm}@{}}
\toprule
\textbf{Person and number} & \textbf{Present} \\
\midrule
1 sg. & sętte \\
2 sg. & sętte \\
3 sg. & sętte \\
1–3 pl. & sętten \\
\bottomrule
\end{tabular}
\par\endgroup
% course-panel:T116-P01:end
% mdtable:T116:end
\end{minipage}\par
\reserve{18\baselineskip}
% mdparagraph:P403:begin
\subsection*{etan: present subjunctive}\BookHeading{nav-ch17-learner-0116}{subsubsection}{etan: present subjunctive}{0}
% mdparagraph:P403:end

% teaching: u17-etan-present-subjunctive6
% mdparagraph:P404:begin
All three singular persons use -e and all three plural persons -en on the present stem, without an indicative second/\allowbreak{}third-singular vowel change. Only this tense of the subjunctive is added here.
% mdparagraph:P404:end

% table: u17-etan-present-subjunctive6-SUBJUNCTIVE-DISPLAY
% mdtable:T117:begin
% course-table-style: 2026-10
% course-panel:T117-P01:begin
\par\begingroup
\fontsize{10.5}{13.2}\selectfont
\setlength{\LTleft}{0pt}\setlength{\LTright}{\fill}
\noindent
\begin{tabular}{@{}>{\raggedright\arraybackslash}p{63.342mm}>{\raggedright\arraybackslash}p{57.846mm}@{}}
\toprule
\textbf{Person and number} & \textbf{Present} \\
\midrule
1 sg. & ete \\
2 sg. & ete \\
3 sg. & ete \\
1–3 pl. & eten \\
\bottomrule
\end{tabular}
\par\endgroup
% course-panel:T117-P01:end
% mdtable:T117:end
\reserve{18\baselineskip}
% mdparagraph:P405:begin
\subsection*{forgiefan: singular and plural commands}\BookHeading{nav-ch17-learner-0117}{subsubsection}{forgiefan: singular and plural commands}{0}
% mdparagraph:P405:end

% teaching: u17-forgiefan-imperative2
% mdparagraph:P406:begin
Use the unmutated present stem for the singular command and -að for the plural. In the prayer, the meaning is forgive, with the familiar thing/\allowbreak{}person case distinction. Only this complete command pair is added here.
% mdparagraph:P406:end

% table: u17-forgiefan-imperative2-IMPERATIVE-DISPLAY
% mdtable:T118:begin
% course-table-style: 2026-10
% course-panel:T118-P01:begin
\par\begingroup
\fontsize{10.5}{13.2}\selectfont
\setlength{\LTleft}{0pt}\setlength{\LTright}{\fill}
\noindent
\begin{tabular}{@{}>{\raggedright\arraybackslash}p{63.342mm}>{\raggedright\arraybackslash}p{57.846mm}@{}}
\toprule
\textbf{Person and number} & \textbf{Command} \\
\midrule
2 sg. & forgief \\
2 pl. & forgiefað \\
\bottomrule
\end{tabular}
\par\endgroup
% course-panel:T118-P01:end
% mdtable:T118:end
\reserve{18\baselineskip}
% mdparagraph:P407:begin
\subsection*{gesēon: plain and inflected infinitives}\BookHeading{nav-ch17-learner-0118}{subsubsection}{gesēon: plain and inflected infinitives}{0}
% mdparagraph:P407:end

% teaching: u17-geseon-infinitive-pair2
% mdparagraph:P408:begin
The pair distinguishes the plain infinitive from the inflected infinitive. The selected series uses -anne, with -ianne in leornianne and -onne after contracted gesēon. The alternative -enne series also exists. The modal constructions here require the plain infinitive; this table adds no new purpose-clause use.
% mdparagraph:P408:end

% table: u17-geseon-infinitive-pair2-BOUNDED-NONFINITE-DISPLAY
% mdtable:T119:begin
\par\begingroup
\fontsize{10.5}{13.2}\selectfont
\setlength{\LTleft}{0pt}\setlength{\LTright}{\fill}
\noindent
\begin{tabular}{@{}>{\raggedright\arraybackslash}p{76.37mm}>{\raggedright\arraybackslash}p{44.82mm}@{}}
\toprule
\textbf{Form} & \textbf{Literal surface} \\
\midrule
verb form plain infinitive & gesēon \\
verb form inflected infinitive & gesēonne \\
\bottomrule
\end{tabular}
\par\endgroup

% mdtable:T119:end
\reserve{18\baselineskip}
% mdparagraph:P409:begin
\subsection*{gelīefan: plain and inflected infinitives}\BookHeading{nav-ch17-learner-0119}{subsubsection}{gelīefan: plain and inflected infinitives}{0}
% mdparagraph:P409:end

% teaching: u17-geliefan-infinitive-pair2
% mdparagraph:P410:begin
The pair distinguishes the plain infinitive from the inflected infinitive. The selected series uses -anne, with -ianne in leornianne and -onne after contracted gesēon. The alternative -enne series also exists. The modal constructions here require the plain infinitive; this table adds no new purpose-clause use.
% mdparagraph:P410:end

% table: u17-geliefan-infinitive-pair2-BOUNDED-NONFINITE-DISPLAY
% mdtable:T120:begin
\par\begingroup
\fontsize{10.5}{13.2}\selectfont
\setlength{\LTleft}{0pt}\setlength{\LTright}{\fill}
\noindent
\begin{tabular}{@{}>{\raggedright\arraybackslash}p{73.68mm}>{\raggedright\arraybackslash}p{47.51mm}@{}}
\toprule
\textbf{Form} & \textbf{Literal surface} \\
\midrule
verb form plain infinitive & gelīefan \\
verb form inflected infinitive & gelīefanne \\
\bottomrule
\end{tabular}
\par\endgroup

% mdtable:T120:end
\reserve{18\baselineskip}
% mdparagraph:P411:begin
\subsection*{leornian: plain and inflected infinitives}\BookHeading{nav-ch17-learner-0120}{subsubsection}{leornian: plain and inflected infinitives}{0}
% mdparagraph:P411:end

% teaching: u17-leornian-infinitive-pair2
% mdparagraph:P412:begin
The pair distinguishes the plain infinitive from the inflected infinitive. The selected series uses -anne, with -ianne in leornianne and -onne after contracted gesēon. The alternative -enne series also exists. The modal constructions here require the plain infinitive; this table adds no new purpose-clause use.
% mdparagraph:P412:end

% table: u17-leornian-infinitive-pair2-BOUNDED-NONFINITE-DISPLAY
% mdtable:T121:begin
\par\begingroup
\fontsize{10.5}{13.2}\selectfont
\setlength{\LTleft}{0pt}\setlength{\LTright}{\fill}
\noindent
\begin{tabular}{@{}>{\raggedright\arraybackslash}p{73.68mm}>{\raggedright\arraybackslash}p{47.51mm}@{}}
\toprule
\textbf{Form} & \textbf{Literal surface} \\
\midrule
verb form plain infinitive & leornian \\
verb form inflected infinitive & leornianne \\
\bottomrule
\end{tabular}
\par\endgroup

% mdtable:T121:end
\reserve{7\baselineskip}
% mdparagraph:P413:begin
\subsection*{gescieppan: present and past participles}\BookHeading{nav-ch17-learner-0121}{subsubsection}{gescieppan: present and past participles}{0}
% mdparagraph:P413:end

% teaching: u17-gescieppan-participle-pair2
% mdparagraph:P414:begin
Retain the lexical ge-. The present participle uses the present stem, while gesceapen has ea and single p.
% mdparagraph:P414:end

\reserve{10\baselineskip}
% mdparagraph:P415:begin
This display assigns the complete present/\allowbreak{}past formation pair. It does not add other finite moods or an unrestricted attributive declension; the past form is used in the stated passive frame.
% mdparagraph:P415:end

% table: u17-gescieppan-participle-pair2-PARTICIPLES-DISPLAY
% mdtable:T122:begin
\par\begingroup
\fontsize{10.5}{13.2}\selectfont
\setlength{\LTleft}{0pt}\setlength{\LTright}{\fill}
\noindent
\begin{tabular}{@{}>{\raggedright\arraybackslash}p{51.9mm}>{\raggedright\arraybackslash}p{69.29mm}@{}}
\toprule
\textbf{Form} & \textbf{Literal surface} \\
\midrule
present & gescieppende \\
past & gesceapen \\
\bottomrule
\end{tabular}
\par\endgroup

% mdtable:T122:end
\reserve{18\baselineskip}
% mdparagraph:P416:begin
\subsection*{wyrcean: base and plural predicate participle}\BookHeading{nav-ch17-learner-0122}{subsubsection}{wyrcean: base and plural predicate participle}{0}
% mdparagraph:P416:end

% teaching: u17-wyrcean-past-predicate2
% mdparagraph:P417:begin
The base is selected for singular predicates and the -e form for the plural predicate in Cook 37.5. The uninflected base is also permitted with a plural predicate. These are selected alternatives, not a complete attributive declension.
% mdparagraph:P417:end

% table: u17-wyrcean-past-predicate2-DISPLAY
% mdtable:T123:begin
\par\begingroup
\fontsize{10.5}{13.2}\selectfont
\setlength{\LTleft}{0pt}\setlength{\LTright}{\fill}
\noindent
\begin{tabular}{@{}>{\raggedright\arraybackslash}p{87.2mm}>{\raggedright\arraybackslash}p{33.98mm}@{}}
\toprule
\textbf{Predicate form} & \textbf{Literal surface} \\
\midrule
Uninflected predicate base (singular or plural) & worht \\
Agreeing plural predicate & worhte \\
\bottomrule
\end{tabular}
\par\endgroup

% mdtable:T123:end
\reserve{18\baselineskip}
% mdparagraph:P418:begin
\subsection*{gescieppan: base and plural predicate participle}\BookHeading{nav-ch17-learner-0123}{subsubsection}{gescieppan: base and plural predicate participle}{0}
% mdparagraph:P418:end

% teaching: u17-gescieppan-past-predicate2
% mdparagraph:P419:begin
The base is selected for singular predicates and the -e form for the plural predicate in Cook 37.5. The uninflected base is also permitted with a plural predicate. These are selected alternatives, not a complete attributive declension.
% mdparagraph:P419:end

% table: u17-gescieppan-past-predicate2-DISPLAY
% mdtable:T124:begin
\par\begingroup
\fontsize{10.5}{13.2}\selectfont
\setlength{\LTleft}{0pt}\setlength{\LTright}{\fill}
\noindent
\begin{tabular}{@{}>{\raggedright\arraybackslash}p{83.56mm}>{\raggedright\arraybackslash}p{37.63mm}@{}}
\toprule
\textbf{Predicate form} & \textbf{Literal surface} \\
\midrule
Uninflected predicate base (singular or plural) & gesceapen \\
Agreeing plural predicate & gesceapene \\
\bottomrule
\end{tabular}
\par\endgroup

% mdtable:T124:end
\reserve{7\baselineskip}
% mdparagraph:P420:begin
\section*{4. Meanings and complements}\BookHeading{nav-ch17-learner-0124}{subsection}{4. Meanings and complements}{1}
% mdparagraph:P420:end

% mdparagraph:P421:begin
The following meanings extend familiar words. Use their previously learned forms only in the constructions now stated. The new noun and adjective tables above already give their own meanings; the additional distinctions here prevent a familiar spelling from hiding a new lexical task.
% mdparagraph:P421:end

% teaching: NEW-U17-VOC-HEALDAN-U17-MAINTAIN
% mdparagraph:P422:begin
\textbf{healdan} — maintain; sustain.
% mdparagraph:P422:end

% teaching: NEW-U17-VOC-SAED-U17-DESCENDANTS
% mdparagraph:P423:begin
\textbf{sǣd} — seed; offspring; descendants collectively.
% mdparagraph:P423:end

% teaching: NEW-U17-VOC-FOLC-U17-PEOPLES
% mdparagraph:P424:begin
\textbf{folc} — peoples; distinct peoples or national groups.
% mdparagraph:P424:end

% teaching: NEW-U17-VOC-EALD-U17-ELDERLY
% mdparagraph:P425:begin
\textbf{eald} — old people; the elderly.
% mdparagraph:P425:end

% teaching: NEW-U17-VOC-WER-U17-HUSBAND
% mdparagraph:P426:begin
\textbf{wer} — husband.
% mdparagraph:P426:end

% teaching: NEW-U17-VOC-DON-U17-FULFIL
% mdparagraph:P427:begin
\textbf{dōn} — fulfil; carry out a word or command.
% mdparagraph:P427:end

% teaching: NEW-U17-VOC-SECGAN-U17-CALL
% mdparagraph:P428:begin
\textbf{sęcgan} — call; declare to be.
% mdparagraph:P428:end

% teaching: NEW-U17-VOC-LAEDAN-U17-CARRY
% mdparagraph:P429:begin
\textbf{lǣdan} — take; carry along.
% mdparagraph:P429:end

% teaching: NEW-U17-VOC-ONFON-U17-PERSON
% mdparagraph:P430:begin
\textbf{onfōn} — receive; accept a person into care.
% mdparagraph:P430:end

% teaching: NEW-U17-VOC-FORLAETAN-U17-LET-GO
% mdparagraph:P431:begin
\textbf{forlǣtan} — let go; dismiss.
% mdparagraph:P431:end

\reserve{7\baselineskip}
% mdparagraph:P432:begin
\subsection*{Participles naming people}\BookHeading{nav-ch17-learner-0125}{subsubsection}{Participles naming people}{0}
% mdparagraph:P432:end

% teaching: NEW-U17-GRAM-U17-PARTICIPLE-HUMAN-NOMINALIZATION
% mdparagraph:P433:begin
A\BookPoint{bk-ch17-learner-07c1bbec4fe0} present participle can describe people without an expressed noun. Its ending gives the phrase’s role in the surrounding clause, while the verb inside the participle retains its own complements. In hine ondrǣdendum, the dative plural names the beneficiaries of mercy; accusative hine names the person feared. In ðā hyngrendan, the article and weak accusative plural identify hungry people as a direct object. These are nonfinite descriptions, not present-tense assertions. Learn these human uses of ondrǣdende and hyngrende; no personal finite conjugation of hyngran is added.
% mdparagraph:P433:end

\reserve{7\baselineskip}
% mdparagraph:P434:begin
\subsection*{Past participles in a predicate}\BookHeading{nav-ch17-learner-0126}{subsubsection}{Past participles in a predicate}{0}
% mdparagraph:P434:end

% teaching: NEW-U17-GRAM-U17-PREDICATE-PARTICIPLE-PLURAL
% mdparagraph:P435:begin
A\BookPoint{bk-ch17-learner-5930a132dc85} past participle with a form of be can remain uninflected or take the strong masculine agreement endings, regardless of the subject’s gender. The selected plural predicates in Cook 37.5 are worhte and gesceapene. They belong with wǣron; they are not finite weak preterites. The uninflected bases worht and gesceapen are also possible with plural subjects. The two small displays in the preceding section distinguish the base from the selected agreeing plural.
% mdparagraph:P435:end

\reserve{7\baselineskip}
% mdparagraph:P436:begin
\subsection*{Praise: statement, command and wish}\BookHeading{nav-ch17-learner-0127}{subsubsection}{Praise: statement, command and wish}{0}
% mdparagraph:P436:end

% teaching: NEW-U17-GRAM-U17-HERIAN-OBJECT
% mdparagraph:P437:begin
Hęrian\BookPoint{bk-ch17-learner-f0d1bccd5033} takes the praised person or thing in the accusative. A plural command addresses the beings named in the catalogue; a third-person wish has hīe as subject and hęrien as verb. The identical indicative and imperative form hęriað does not make their functions identical.
% mdparagraph:P437:end

\reserve{7\baselineskip}
% mdparagraph:P438:begin
\subsection*{Requesting and seeking in marriage}\BookHeading{nav-ch17-learner-0128}{subsubsection}{Requesting and seeking in marriage}{0}
% mdparagraph:P438:end

% teaching: NEW-U17-GRAM-U17-GIERNAN-REQUEST-WOO
% mdparagraph:P439:begin
Giernan governs a genitive of the requested thing or marriage partner. Æt with the dative gives the person asked. In Cook 39.3, keep the outer accusative role of eall ðæt separate from the role represented by invariant ðe in its relative clause. In Cook 43.14, seeking the maiden in marriage requires her genitive form.
% mdparagraph:P439:end

\reserve{7\baselineskip}
% mdparagraph:P440:begin
\subsection*{Contents and covering}\BookHeading{nav-ch17-learner-0129}{subsubsection}{Contents and covering}{0}
% mdparagraph:P440:end

% teaching: NEW-U17-GRAM-U17-FILLING-MATERIAL
% mdparagraph:P441:begin
Āfyllan takes the filled entity in the accusative and its contents after mid with the dative. In a passive, the filled entity becomes the nominative subject. Oferwyrcean has a similar case arrangement for a different act: the throne is overlaid with gold. Learn each verb’s meaning together with its complements.
% mdparagraph:P441:end

\reserve{7\baselineskip}
% mdparagraph:P442:begin
\subsection*{Receiving a thing or a person}\BookHeading{nav-ch17-learner-0130}{subsubsection}{Receiving a thing or a person}{0}
% mdparagraph:P442:end

% teaching: NEW-U17-GRAM-U17-ONFON-PERSON-DATIVE
% mdparagraph:P443:begin
The selected onfōn construction takes the dative: edlēane for a received reward, and his cnihte Israel for the servant accepted into care. Accusative receiving constructions exist elsewhere; these exercises specifically practise the dative frame.
% mdparagraph:P443:end

\reserve{7\baselineskip}
% mdparagraph:P444:begin
\subsection*{Calling someone blessed}\BookHeading{nav-ch17-learner-0131}{subsubsection}{Calling someone blessed}{0}
% mdparagraph:P444:end

% teaching: NEW-U17-GRAM-U17-SECGAN-PREDICATE
% mdparagraph:P445:begin
Sęcgan\BookPoint{bk-ch17-learner-cbd56e38041e} can take an object and an adjective that agrees with it: the generations call Mary blessed. In Cook 40.3, mē is the first-person singular object, ēadige its feminine accusative singular predicate, and sęcgað agrees with the plural generations. The present verb has future reference in the prediction.
% mdparagraph:P445:end

\reserve{7\baselineskip}
% mdparagraph:P446:begin
\subsection*{New meanings of familiar verbs}\BookHeading{nav-ch17-learner-0132}{subsubsection}{New meanings of familiar verbs}{0}
% mdparagraph:P446:end

% teaching: NEW-U17-GRAM-U17-VERB-LEXICAL-EXTENSIONS
% mdparagraph:P447:begin
Dōn can mean fulfil or carry out a word or command; healdan can mean sustain all things; sęcgan can mean call or declare someone to be something. Learn those meanings explicitly. Wyrcean retains the familiar work or perform sense with the Father’s will and the make or build sense with a foundation.
% mdparagraph:P447:end

\reserve{7\baselineskip}
% mdparagraph:P448:begin
\subsection*{Indefinite relatives and the subjunctive}\BookHeading{nav-ch17-learner-0133}{subsubsection}{Indefinite relatives and the subjunctive}{0}
% mdparagraph:P448:end

% teaching: NEW-U17-GRAM-U17-VIRTUAL-RELATIVE
% mdparagraph:P449:begin
A\BookPoint{bk-ch17-learner-273409eecb31} relative clause can describe an indefinite or possible situation: the one you wish to choose, or a place where one might lay one’s head. Such clauses can take the subjunctive. Wille in Cook 43.12 is second singular present subjunctive; āhylde in Cook 42.7 is third singular present subjunctive. The construction does not make every relative subjunctive, and English may need not be translated by an extra modal.
% mdparagraph:P449:end

\reserve{7\baselineskip}
% mdparagraph:P450:begin
\subsection*{Learning from someone}\BookHeading{nav-ch17-learner-0134}{subsubsection}{Learning from someone}{0}
% mdparagraph:P450:end

% teaching: NEW-U17-GRAM-U17-LEORNIAN-SOURCE
% mdparagraph:P451:begin
Of with the dative can name the person from whom one learns in Cook 43.11. Gewilnian introduces the desired content with ðæt. The embedded ic takes first singular leornie, while plural gē takes habbað. The separate clauses keep their own subjects and agreement.
% mdparagraph:P451:end

\reserve{7\baselineskip}
% mdparagraph:P452:begin
\subsection*{Playing with an object}\BookHeading{nav-ch17-learner-0135}{subsubsection}{Playing with an object}{0}
% mdparagraph:P452:end

% teaching: NEW-U17-GRAM-U17-PLEGIAN-INSTRUMENT
% mdparagraph:P453:begin
Plegian in the ordinary play sense takes mid with the dative instrument. The ball exercise therefore uses the familiar dative ðoðore. The reading’s act or perform sense is supplied separately.
% mdparagraph:P453:end

\reserve{7\baselineskip}
% mdparagraph:P454:begin
\subsection*{Before a past event}\BookHeading{nav-ch17-learner-0136}{subsubsection}{Before a past event}{0}
% mdparagraph:P454:end

% teaching: NEW-U17-GRAM-U17-TEMPORAL-BEFORE
% mdparagraph:P455:begin
In\BookPoint{bk-ch17-learner-41a025a3a7a7} Cook 38.8, nolde expresses a former unwillingness to believe. The seeing is still prospective from that earlier viewpoint: the before-clause uses first singular past subjunctive gesāwe. Its spelling alone does not make it second singular indicative. Retain self and follow the requested final position of the verb.
% mdparagraph:P455:end

\reserve{7\baselineskip}
% mdparagraph:P456:begin
\subsection*{Negative concord}\BookHeading{nav-ch17-learner-0137}{subsubsection}{Negative concord}{0}
% mdparagraph:P456:end

% teaching: NEW-U17-GRAM-U15-NEGATIVE-CONCORD
% mdparagraph:P457:begin
Preverbal ne can combine with nān or another negative word to express one negation. In the coordinated eye, ear and heart sentence, the first mæg can govern the following infinitives as well. Identify the subject of each part; repeated negative words do not make it affirmative. Use the ordinary ne witon for we do not know. The reading’s contracted nāst is supplied separately.
% mdparagraph:P457:end

\reserve{7\baselineskip}
% mdparagraph:P458:begin
\subsection*{Choosing a past tense}\BookHeading{nav-ch17-learner-0138}{subsubsection}{Choosing a past tense}{0}
% mdparagraph:P458:end

% teaching: NEW-U17-GRAM-U17-PAST-TENSE-SELECTION
% mdparagraph:P459:begin
The simple preterite can correspond to English simple past, present perfect or past perfect when the context establishes the time relation. Cook §188 explicitly permits this range. An English has or had therefore does not always require a separate auxiliary in Old English. Follow an explicit omission instruction; elsewhere choose a construction whose exact forms and intended time relation you know.
% mdparagraph:P459:end

\reserve{7\baselineskip}
% mdparagraph:P460:begin
\subsection*{The grounds for a reward}\BookHeading{nav-ch17-learner-0139}{subsubsection}{The grounds for a reward}{0}
% mdparagraph:P460:end

% teaching: NEW-U17-GRAM-U17-GENITIVE-REWARD-CAUSE
% mdparagraph:P461:begin
A\BookPoint{bk-ch17-learner-af5777028dbd} genitive dependent on a reward can express its grounds: his weorces, for his work. In Cook 39.6, edlēane is dative because it is received with onfōn; the work-genitive depends on the reward, and æt Gode gives its source.
% mdparagraph:P461:end

% mdparagraph:P462:begin
For sǣd in the promise to Abraham, the descendants are expressed collectively. The assigned dative sǣde is singular. Folc can denote distinct peoples in the plural catalogue, while the collective people in the dense horn passage remains singular. Wer now also means husband. Keep these contextual meanings separate from the earlier plant seed, people and man meanings.
% mdparagraph:P462:end

\reserve{7\baselineskip}
% mdparagraph:P463:begin
\section*{5. Form contrasts and composition}\BookHeading{nav-ch17-learner-0140}{subsection}{5. Form contrasts and composition}{1}
% mdparagraph:P463:end

% mdparagraph:P464:begin
Complete each requested response. The form contrasts test case, number and clause function in meaningful expressions. Keep supplied dictionary forms available, but choose the required endings from the tables and explanations. In the compositions, preserve every participant and the stated clause relationship.
% mdparagraph:P464:end

\reserve{8\baselineskip}
% mdparagraph:P465:begin
\subsection*{F01}\BookHeading{nav-ch17-learner-0141}{subsubsection}{F01}{0}
% mdparagraph:P465:end

% item: NEW-U17-F01
% mdparagraph:P466:begin
Use cniht in its explicitly taught servant sense and bōc book. Translate the four expressions. In each, choose the noun's case from its role rather than from an English word ending.
% mdparagraph:P466:end

% mdparagraph:P467:begin
cniht
% mdparagraph:P467:end

% mdparagraph:P468:begin
bōc
% mdparagraph:P468:end

% output: NEW-U17-F01-O001
% mdparagraph:P469:begin
\itempara{1}{the servant's book}
% mdparagraph:P469:end

% output: NEW-U17-F01-O002
% mdparagraph:P470:begin
\itempara{2}{the servants' books}
% mdparagraph:P470:end

% output: NEW-U17-F01-O003
% mdparagraph:P471:begin
\itempara{3}{I give the servant the book.}
% mdparagraph:P471:end

% output: NEW-U17-F01-O004
% mdparagraph:P472:begin
\itempara{4}{We give the servants the books.}
% mdparagraph:P472:end

\reserve{8\baselineskip}
% mdparagraph:P473:begin
\subsection*{F02}\BookHeading{nav-ch17-learner-0142}{subsubsection}{F02}{0}
% mdparagraph:P473:end

% item: NEW-U17-F02
% mdparagraph:P474:begin
Use the selected endingless accusative of cwēn queen, while retaining its genitive/\allowbreak{}dative ending. Translate the three expressions, then explain the contrast.
% mdparagraph:P474:end

% mdparagraph:P475:begin
cwēn
% mdparagraph:P475:end

% mdparagraph:P476:begin
bōc
% mdparagraph:P476:end

% output: NEW-U17-F02-O001
% mdparagraph:P477:begin
\itempara{1}{I saw the queen.}
% mdparagraph:P477:end

% output: NEW-U17-F02-O002
% mdparagraph:P478:begin
\itempara{2}{I gave the queen the book.}
% mdparagraph:P478:end

% output: NEW-U17-F02-O003
% mdparagraph:P479:begin
\itempara{3}{the queen's book}
% mdparagraph:P479:end

% output: NEW-U17-F02-O004
% mdparagraph:P480:begin
\itempara{4}{Explain what the three queen forms show about identifying case from an ending alone.}
% mdparagraph:P480:end

\reserve{8\baselineskip}
% mdparagraph:P481:begin
\subsection*{F03}\BookHeading{nav-ch17-learner-0143}{subsubsection}{F03}{0}
% mdparagraph:P481:end

% item: NEW-U17-F03
% mdparagraph:P482:begin
Use fox fox, hol hole and nest nest. Translate each sentence. Keep ordinary locative on with the dative and distinguish the two neuter plural patterns.
% mdparagraph:P482:end

% mdparagraph:P483:begin
fox
% mdparagraph:P483:end

% mdparagraph:P484:begin
hol
% mdparagraph:P484:end

% mdparagraph:P485:begin
nest
% mdparagraph:P485:end

% output: NEW-U17-F03-O001
% mdparagraph:P486:begin
\itempara{1}{The fox is in the hole.}
% mdparagraph:P486:end

% output: NEW-U17-F03-O002
% mdparagraph:P487:begin
\itempara{2}{The foxes are in the holes.}
% mdparagraph:P487:end

% output: NEW-U17-F03-O003
% mdparagraph:P488:begin
\itempara{3}{The nest is in the tree.}
% mdparagraph:P488:end

% output: NEW-U17-F03-O004
% mdparagraph:P489:begin
\itempara{4}{The nests are in the trees.}
% mdparagraph:P489:end

% output: NEW-U17-F03-O005
% mdparagraph:P490:begin
\itempara{5}{The holes are large.}
% mdparagraph:P490:end

\reserve{8\baselineskip}
% mdparagraph:P491:begin
\subsection*{F04}\BookHeading{nav-ch17-learner-0144}{subsubsection}{F04}{0}
% mdparagraph:P491:end

% item: NEW-U17-F04
% mdparagraph:P492:begin
Use hęrian praise. Supply only the finite form required in each stated setting, then explain the contrast. The dictionary form is visible reference help.
% mdparagraph:P492:end

% mdparagraph:P493:begin
hęrian
% mdparagraph:P493:end

% output: NEW-U17-F04-O001
% mdparagraph:P494:begin
\itempara{1}{A narrator states: they praise. Give present indicative third plural.}
% mdparagraph:P494:end

% output: NEW-U17-F04-O002
% mdparagraph:P495:begin
\itempara{2}{A speaker commands several listeners: praise! Give second-plural imperative.}
% mdparagraph:P495:end

% output: NEW-U17-F04-O003
% mdparagraph:P496:begin
\itempara{3}{A speaker wishes: may they praise. Give present subjunctive third plural.}
% mdparagraph:P496:end

% output: NEW-U17-F04-O004
% mdparagraph:P497:begin
\itempara{4}{Which answers are spelled alike, and what distinguishes their functions?}
% mdparagraph:P497:end

\reserve{8\baselineskip}
% mdparagraph:P498:begin
\subsection*{F05}\BookHeading{nav-ch17-learner-0145}{subsubsection}{F05}{0}
% mdparagraph:P498:end

% item: NEW-U17-F05
% mdparagraph:P499:begin
Use the explicitly taught present participles of ondrǣdan fear and hyngran be hungry. Translate the two nominal phrases, then distinguish each phrase's outer case from any object governed inside it.
% mdparagraph:P499:end

% mdparagraph:P500:begin
ondrǣdan
% mdparagraph:P500:end

% mdparagraph:P501:begin
hyngran
% mdparagraph:P501:end

% output: NEW-U17-F05-O001
% mdparagraph:P502:begin
\itempara{1}{to those fearing him: use a dative plural participle without an expressed people noun; him is the one feared}
% mdparagraph:P502:end

% output: NEW-U17-F05-O002
% mdparagraph:P503:begin
\itempara{2}{the hungry: give the weak accusative plural phrase used as a direct object}
% mdparagraph:P503:end

% output: NEW-U17-F05-O003
% mdparagraph:P504:begin
\itempara{3}{State the two grammatical relations in the fear phrase.}
% mdparagraph:P504:end

% output: NEW-U17-F05-O004
% mdparagraph:P505:begin
\itempara{4}{Explain why a present participle can name people here without being a finite verb.}
% mdparagraph:P505:end

\reserve{8\baselineskip}
% mdparagraph:P506:begin
\subsection*{C01}\BookHeading{nav-ch17-learner-0146}{subsubsection}{C01}{0}
% mdparagraph:P506:end

% item: NEW-U17-C01
% mdparagraph:P507:begin
Translate: The king saw the men who sang. Use an invariant relative particle.
% mdparagraph:P507:end

% output: NEW-U17-C01-O001
% mdparagraph:P508:begin
\itempara{1}{Write one complete Old English sentence.}
% mdparagraph:P508:end

\reserve{8\baselineskip}
% mdparagraph:P509:begin
\subsection*{C02}\BookHeading{nav-ch17-learner-0147}{subsubsection}{C02}{0}
% mdparagraph:P509:end

% item: NEW-U17-C02
% mdparagraph:P510:begin
Translate: The king is as good as the men are. Express equal degree with swā … swā.
% mdparagraph:P510:end

% output: NEW-U17-C02-O001
% mdparagraph:P511:begin
\itempara{1}{Write one complete Old English sentence.}
% mdparagraph:P511:end

\reserve{8\baselineskip}
% mdparagraph:P512:begin
\subsection*{C03}\BookHeading{nav-ch17-learner-0148}{subsubsection}{C03}{0}
% mdparagraph:P512:end

% item: NEW-U17-C03
% mdparagraph:P513:begin
A visit is possible tomorrow. Translate: If he comes, give (plural) him the book.
% mdparagraph:P513:end

% output: NEW-U17-C03-O001
% mdparagraph:P514:begin
\itempara{1}{Write one complete Old English sentence.}
% mdparagraph:P514:end

\reserve{8\baselineskip}
% mdparagraph:P515:begin
\subsection*{C04}\BookHeading{nav-ch17-learner-0149}{subsubsection}{C04}{0}
% mdparagraph:P515:end

% item: NEW-U17-C04
% mdparagraph:P516:begin
Translate: The man says that the king is good. Present the content as a report, using the present subjunctive.
% mdparagraph:P516:end

% output: NEW-U17-C04-O001
% mdparagraph:P517:begin
\itempara{1}{Write one complete Old English sentence.}
% mdparagraph:P517:end

\reserve{7\baselineskip}
% mdparagraph:P518:begin
\section*{6. Cook: connected composition}\BookHeading{nav-ch17-learner-0150}{subsection}{6. Cook: connected composition}{1}
% mdparagraph:P518:end

% teaching: NEW-U17-SOURCE-HINTS
% teaching: NEW-U17-VERBAL-CONTEXT-HELP
% mdparagraph:P519:begin
Translate the complete English prompts. Parenthetic Old English is visible lexical or phrase help; a headword may require inflection, while an exact supplied expression is used as given. Parenthetic numbers refer to Cook’s grammar: the relevant constructions are explained here, so the locator is not a substitute for instruction. An omission control removes the specified English word from the Old English wording, while preserving its intended relation. Brackets supply the requested paraphrase. Distinguish a grammatical alternative from compliance with an explicit word, order, number or omission target.
% mdparagraph:P519:end

% mdparagraph:P520:begin
All supplied expressions and meanings remain available during these exercises. Copying an exact given form is a reference response. When the hint gives a headword but the task requires a different inflection, its lexical meaning is still given; the new ending is the part you must choose.
% mdparagraph:P520:end

\reserve{7\baselineskip}
% mdparagraph:P521:begin
\section*{Cook 37. Praise and enumeration}\BookHeading{nav-ch17-learner-0151}{subsection}{Cook 37. Praise and enumeration}{1}
% mdparagraph:P521:end

% mdparagraph:P522:begin
The catalogue alternates addressed beings with statements and wishes about them. Track who is addressed, who acts, and who receives praise. A list can leave a governing verb understood from its setting; keep the listed items complete.
% mdparagraph:P522:end

% mdparagraph:P523:begin
[In this exercise, omit the article and the second personal pronoun throughout.]
% mdparagraph:P523:end

% mdparagraph:P524:begin
This applies to all fourteen items: omit the article and second-person pronoun. It does not remove a third-person subject needed for “let them praise”.
% mdparagraph:P524:end

\par\begin{minipage}{\linewidth}\raggedright
% mdparagraph:P525:begin
\subsection*{Cook 37.1}\BookHeading{nav-ch17-learner-0152}{subsubsection}{Cook 37.1}{0}
% mdparagraph:P525:end

% item: EX-XII-XIII-37-01
% mdparagraph:P526:begin
Praise ye the Lord from (of) the heavens; praise him in the heights.
% mdparagraph:P526:end

% output: EX-XII-XIII-37-01-O001
% mdparagraph:P527:begin
Write one complete Old English translation, keeping the printed controls and the supplied expressions available.
% mdparagraph:P527:end

\end{minipage}\par\addvspace{4pt}

\par\begin{minipage}{\linewidth}\raggedright
% mdparagraph:P528:begin
\subsection*{Cook 37.2}\BookHeading{nav-ch17-learner-0153}{subsubsection}{Cook 37.2}{0}
% mdparagraph:P528:end

% item: EX-XII-XIII-37-02
% mdparagraph:P529:begin
Praise ye him, all his angels; praise him, all his powers (mægen).
% mdparagraph:P529:end

% output: EX-XII-XIII-37-02-O001
% mdparagraph:P530:begin
Write one complete Old English translation, keeping the printed controls and the supplied expressions available.
% mdparagraph:P530:end

\end{minipage}\par\addvspace{4pt}

\par\begin{minipage}{\linewidth}\raggedright
% mdparagraph:P531:begin
\subsection*{Cook 37.3}\BookHeading{nav-ch17-learner-0154}{subsubsection}{Cook 37.3}{0}
% mdparagraph:P531:end

% item: EX-XII-XIII-37-03
% mdparagraph:P532:begin
Praise ye him, sun and moon; praise him, all ye stars and light (lēoht).
% mdparagraph:P532:end

% output: EX-XII-XIII-37-03-O001
% mdparagraph:P533:begin
Write one complete Old English translation, keeping the printed controls and the supplied expressions available.
% mdparagraph:P533:end

\end{minipage}\par\addvspace{4pt}

\par\begin{minipage}{\linewidth}\raggedright
% mdparagraph:P534:begin
\subsection*{Cook 37.4}\BookHeading{nav-ch17-learner-0155}{subsubsection}{Cook 37.4}{0}
% mdparagraph:P534:end

% item: EX-XII-XIII-37-04
% mdparagraph:P535:begin
Praise him, ye heavens of heavens, and ye waters that are above (ofer, 172) the heavens.
% mdparagraph:P535:end

% output: EX-XII-XIII-37-04-O001
% mdparagraph:P536:begin
Write one complete Old English translation, keeping the printed controls and the supplied expressions available.
% mdparagraph:P536:end

\end{minipage}\par\addvspace{4pt}

\par\begin{minipage}{\linewidth}\raggedright
% mdparagraph:P537:begin
\subsection*{Cook 37.5}\BookHeading{nav-ch17-learner-0156}{subsubsection}{Cook 37.5}{0}
% mdparagraph:P537:end

% item: EX-XII-XIII-37-05
% mdparagraph:P538:begin
Let them praise [praise they] the name of the Lord; because he spake (gecweðan), and they were made (wyrcean); he commanded, and they were created (p. 190, l. 19).
% mdparagraph:P538:end

% mdparagraph:P539:begin
\textbf{Construction.} “Praise they” calls for a third-person wish with an expressed subject. Cook’s page 190, line 19 begins the said/\allowbreak{}made example; commanded/\allowbreak{}created follows on line 20. Use the passive auxiliary with the appropriate predicate participles.
% mdparagraph:P539:end

% output: EX-XII-XIII-37-05-O001
% mdparagraph:P540:begin
Write one complete Old English translation, keeping the printed controls and the supplied expressions available.
% mdparagraph:P540:end

\end{minipage}\par\addvspace{4pt}

\par\begin{minipage}{\linewidth}\raggedright
% mdparagraph:P541:begin
\subsection*{Cook 37.6}\BookHeading{nav-ch17-learner-0157}{subsubsection}{Cook 37.6}{0}
% mdparagraph:P541:end

% item: EX-XII-XIII-37-06
% mdparagraph:P542:begin
He hath appointed them for ever and ever; he (om. hath) made (sęttan) a decree (bebod, 47), and (om. it) shall not pass (gewītan).
% mdparagraph:P542:end

\end{minipage}\par\addvspace{4pt}

% teaching: NEW-U17-GRAM-U17-H-37-06
% mdparagraph:P543:begin
\textbf{Construction.} Omit a separate word for hath. In the final coordinated predicate, the decree is the understood subject; God is not the thing that passes away. The supplied present gewīt has future reference.
% mdparagraph:P543:end

% teaching: NEW-U17-VOC-GEWITAN-U17-DECREE-PASS
% mdparagraph:P544:begin
\textbf{Supplied:} hē sętte bebod and ne gewīt. — He established a decree, and it will not pass away. The second predicate has the understood decree as subject. Keep the omitted it unexpressed and preserve the future reference of present gewīt.
% mdparagraph:P544:end

% mdparagraph:P545:begin
\textbf{Supplied:} ā — ever; always.
% mdparagraph:P545:end

% mdparagraph:P546:begin
\textbf{Supplied:} ā — ever; always.
% mdparagraph:P546:end

% mdparagraph:P547:begin
\textbf{Supplied:} ā — ever; always.
% mdparagraph:P547:end

% mdparagraph:P548:begin
\textbf{Supplied:} gewīt — pass away; cease to hold; indicative, present, third person, singular.
% mdparagraph:P548:end

% output: EX-XII-XIII-37-06-O001
% mdparagraph:P549:begin
Write one complete Old English translation, keeping the printed controls and the supplied expressions available.
% mdparagraph:P549:end

\par\begin{minipage}{\linewidth}\raggedright
% mdparagraph:P550:begin
\subsection*{Cook 37.7}\BookHeading{nav-ch17-learner-0158}{subsubsection}{Cook 37.7}{0}
% mdparagraph:P550:end

% item: EX-XII-XIII-37-07
% mdparagraph:P551:begin
Praise the Lord from the earth, ye dragons (draca, 53), and all deeps (nēowolnes).
% mdparagraph:P551:end

% mdparagraph:P552:begin
\textbf{Supplied:} nēowolnessa — abyss; the deep; nominative, plural.
% mdparagraph:P552:end

% output: EX-XII-XIII-37-07-O001
% mdparagraph:P553:begin
Write one complete Old English translation, keeping the printed controls and the supplied expressions available.
% mdparagraph:P553:end

\end{minipage}\par\addvspace{4pt}

\par\begin{minipage}{\linewidth}\raggedright
% mdparagraph:P554:begin
\subsection*{Cook 37.8}\BookHeading{nav-ch17-learner-0159}{subsubsection}{Cook 37.8}{0}
% mdparagraph:P554:end

% item: EX-XII-XIII-37-08
% mdparagraph:P555:begin
Fire and hail (hægel, 43); snow (snāw, 43) and ice (īs, 47); spirits of the storms (ȳst, 51. b) which fulfil (dōn) his word;
% mdparagraph:P555:end

% mdparagraph:P556:begin
\textbf{Construction.} The spirits of the storms fulfil his word; keep this fulfilment sense of dōn distinct from a general act of making.
% mdparagraph:P556:end

\end{minipage}\par\addvspace{4pt}

% teaching: U17-SUP-HAEGEL-NOM-SG
% teaching: U17-SUP-SNAW-NOM-SG
% teaching: U17-SUP-IS-ICE-NOM-SG
% mdparagraph:P557:begin
\textbf{Supplied:} hægel — hail; nominative, masculine, singular.
% mdparagraph:P557:end

% mdparagraph:P558:begin
\textbf{Supplied:} snāw — snow; nominative, masculine, singular.
% mdparagraph:P558:end

% mdparagraph:P559:begin
\textbf{Supplied:} īs — ice; nominative, neuter, singular.
% mdparagraph:P559:end

% output: EX-XII-XIII-37-08-O001
% mdparagraph:P560:begin
Write one complete Old English translation, keeping the printed controls and the supplied expressions available.
% mdparagraph:P560:end

\par\begin{minipage}{\linewidth}\raggedright
% mdparagraph:P561:begin
\subsection*{Cook 37.9}\BookHeading{nav-ch17-learner-0160}{subsubsection}{Cook 37.9}{0}
% mdparagraph:P561:end

% item: EX-XII-XIII-37-09
% mdparagraph:P562:begin
Mountains, and all hills (hyll, 43); fruitful (wæstmbǣre, 59) trees, and all cedars;
% mdparagraph:P562:end

% mdparagraph:P563:begin
\textbf{Construction.} The described trees are addressed directly; use the weak plural adjective form. The quantifier all keeps its own agreement.
% mdparagraph:P563:end

% output: EX-XII-XIII-37-09-O001
% mdparagraph:P564:begin
Write one complete Old English translation, keeping the printed controls and the supplied expressions available.
% mdparagraph:P564:end

\end{minipage}\par\addvspace{4pt}

\par\begin{minipage}{\linewidth}\raggedright
% mdparagraph:P565:begin
\subsection*{Cook 37.10}\BookHeading{nav-ch17-learner-0161}{subsubsection}{Cook 37.10}{0}
% mdparagraph:P565:end

% item: EX-XII-XIII-37-10
% mdparagraph:P566:begin
(Wild) beasts, and all cattle; serpents, and winged (gefiðred) fowl;
% mdparagraph:P566:end

% mdparagraph:P567:begin
\textbf{Construction.} “(Wild) beasts” specifies the intended animal sense. The winged birds are directly addressed and take a weak plural adjective.
% mdparagraph:P567:end

\end{minipage}\par\addvspace{4pt}

% teaching: U17-SUP-GEFITHRED-NOM-PL-WEAK
% mdparagraph:P568:begin
\textbf{Supplied:} gefiðredan — feathered; nominative, masculine, plural, weak.
% mdparagraph:P568:end

% output: EX-XII-XIII-37-10-O001
% mdparagraph:P569:begin
Write one complete Old English translation, keeping the printed controls and the supplied expressions available.
% mdparagraph:P569:end

\par\begin{minipage}{\linewidth}\raggedright
% mdparagraph:P570:begin
\subsection*{Cook 37.11}\BookHeading{nav-ch17-learner-0162}{subsubsection}{Cook 37.11}{0}
% mdparagraph:P570:end

% item: EX-XII-XIII-37-11
% mdparagraph:P571:begin
Kings of the earth, and all peoples; princes, and all judges of the earth;
% mdparagraph:P571:end

% mdparagraph:P572:begin
\textbf{Supplied:} folc — peoples; distinct peoples or national groups; nominative, plural.
% mdparagraph:P572:end

% output: EX-XII-XIII-37-11-O001
% mdparagraph:P573:begin
Write one complete Old English translation, keeping the printed controls and the supplied expressions available.
% mdparagraph:P573:end

\end{minipage}\par\addvspace{4pt}

\par\begin{minipage}{\linewidth}\raggedright
% mdparagraph:P574:begin
\subsection*{Cook 37.12}\BookHeading{nav-ch17-learner-0163}{subsubsection}{Cook 37.12}{0}
% mdparagraph:P574:end

% item: EX-XII-XIII-37-12
% mdparagraph:P575:begin
Young (om. men), and maidens (fǣmne); old (om. men), with children;
% mdparagraph:P575:end

% mdparagraph:P576:begin
\textbf{Construction.} Omit men but retain the people meant by young and old. Their weak masculine plural adjectives carry the direct address.
% mdparagraph:P576:end

% output: EX-XII-XIII-37-12-O001
% mdparagraph:P577:begin
Write one complete Old English translation, keeping the printed controls and the supplied expressions available.
% mdparagraph:P577:end

\end{minipage}\par\addvspace{4pt}

\par\begin{minipage}{\linewidth}\raggedright
% mdparagraph:P578:begin
\subsection*{Cook 37.13}\BookHeading{nav-ch17-learner-0164}{subsubsection}{Cook 37.13}{0}
% mdparagraph:P578:end

% item: EX-XII-XIII-37-13
% mdparagraph:P579:begin
Let them praise the name of the Lord; for the name of him alone is exalted (ūpāhębban), his glory (wuldor) is above (ofer, 172) the earth and heaven.
% mdparagraph:P579:end

% mdparagraph:P580:begin
\textbf{Supplied:} ānes — alone; genitive, masculine, singular, strong.
% mdparagraph:P580:end

% output: EX-XII-XIII-37-13-O001
% mdparagraph:P581:begin
Write one complete Old English translation, keeping the printed controls and the supplied expressions available.
% mdparagraph:P581:end

\end{minipage}\par\addvspace{4pt}

\par\begin{minipage}{\linewidth}\raggedright
% mdparagraph:P582:begin
\subsection*{Cook 37.14}\BookHeading{nav-ch17-learner-0165}{subsubsection}{Cook 37.14}{0}
% mdparagraph:P582:end

% item: EX-XII-XIII-37-14
% mdparagraph:P583:begin
He exalteth the horn (horn, 43) of his people, the praise of all his saints (hālga, 53), of the children (bearn) of Israel (Israel), of the people drawing nigh (om. unto) him.
% mdparagraph:P583:end

\end{minipage}\par\addvspace{4pt}

% teaching: NEW-U17-GRAM-U17-H-37-14
% teaching: U17-SUP-ISRAEL-DAT-SG
% mdparagraph:P584:begin
\textbf{Supplied:} horn his folces, lof ealra his hālgena, bearna Israeles, folces him nēalǣcendes. — The horn of his people, the praise of all his saints, of the children of Israel, of a people approaching him. Horn is the accusative thing exalted and represents strength or dignity; lof is appositional praise. The successive genitives identify the people concerned. Hālgena belongs to the noun hālga; folces is singular collective, nēalǣcendes agrees with it, and him is the dative person approached. Use the complete phrase as supplied.
% mdparagraph:P584:end

% teaching: U17-SUP-ISRAEL-GEN-SG
% mdparagraph:P585:begin
\textbf{Supplied:} Israeles — Israel; genitive, singular.
% mdparagraph:P585:end

% mdparagraph:P586:begin
\textbf{Supplied:} nēalǣcendes — approach; draw near; present participle, genitive, neuter, singular, strong.
% mdparagraph:P586:end

% output: EX-XII-XIII-37-14-O001
% mdparagraph:P587:begin
Write one complete Old English translation, keeping the printed controls and the supplied expressions available.
% mdparagraph:P587:end

\reserve{7\baselineskip}
% mdparagraph:P588:begin
\section*{Cook 38. The queen’s visit}\BookHeading{nav-ch17-learner-0166}{subsection}{Cook 38. The queen’s visit}{1}
% mdparagraph:P588:end

% mdparagraph:P589:begin
Follow the sequence of hearing, arrival, conversation, observation and judgment. Comparison and reporting clauses keep their own subjects and times.
% mdparagraph:P589:end

\par\begin{minipage}{\linewidth}\raggedright
% mdparagraph:P590:begin
\subsection*{Cook 38.1}\BookHeading{nav-ch17-learner-0167}{subsubsection}{Cook 38.1}{0}
% mdparagraph:P590:end

% item: EX-XII-XIII-38-01
% mdparagraph:P591:begin
The king made a great ivory (elpenbǣnen, 57. 3) throne, and overlaid (oferwyrcean) it with (mid) the best gold.
% mdparagraph:P591:end

\end{minipage}\par\addvspace{4pt}

% teaching: U17-SUP-ELPENBAENEN-ACC-N-SG
% mdparagraph:P592:begin
\textbf{Supplied:} elpenbǣnen — ivory; made of ivory; accusative, neuter, singular, strong.
% mdparagraph:P592:end

% output: EX-XII-XIII-38-01-O001
% mdparagraph:P593:begin
Write one complete Old English translation, keeping the printed controls and the supplied expressions available.
% mdparagraph:P593:end

\par\begin{minipage}{\linewidth}\raggedright
% mdparagraph:P594:begin
\subsection*{Cook 38.2}\BookHeading{nav-ch17-learner-0168}{subsubsection}{Cook 38.2}{0}
% mdparagraph:P594:end

% item: EX-XII-XIII-38-02
% mdparagraph:P595:begin
There (om.) was a certain queen in those days in the south, called (hātan) Sheba (Saba).
% mdparagraph:P595:end

\end{minipage}\par\addvspace{4pt}

% teaching: NEW-U17-LOCAL-SENSE-ON-U15-SOURCE-DURING
% teaching: NEW-U17-LOCAL-SENSE-BE-U17-SOUTH
% teaching: NEW-U17-GRAM-U17-SUPPLIED-DAYS
% teaching: NEW-U17-GRAM-U17-SUPPLIED-SOUTH
% mdparagraph:P596:begin
\textbf{Construction.} Omit the English dummy there. The complete time and south expressions below remain available.
% mdparagraph:P596:end

% teaching: U17-SUP-SABA-NOM-SG
% mdparagraph:P597:begin
\textbf{Supplied:} on ðǣm dagum — In those days; a complete dative plural time expression.
% mdparagraph:P597:end

% mdparagraph:P598:begin
\textbf{Supplied:} be sūðan, — In the south. Preserve the short vowel of be and use the complete expression.
% mdparagraph:P598:end

% teaching: U17-SUP-SUTHAN-INVARIANT
% mdparagraph:P599:begin
\textbf{Supplied:} on — in those days; during the stated time.
% mdparagraph:P599:end

% mdparagraph:P600:begin
\textbf{Supplied:} be — in the south, only in supplied be sūðan.
% mdparagraph:P600:end

% mdparagraph:P601:begin
\textbf{Supplied:} sūðan — in/\allowbreak{}from the south.
% mdparagraph:P601:end

% mdparagraph:P602:begin
\textbf{Supplied:} Saba — Sheba; nominative, singular.
% mdparagraph:P602:end

% output: EX-XII-XIII-38-02-O001
% mdparagraph:P603:begin
Write one complete Old English translation, keeping the printed controls and the supplied expressions available.
% mdparagraph:P603:end

\par\begin{minipage}{\linewidth}\raggedright
% mdparagraph:P604:begin
\subsection*{Cook 38.3}\BookHeading{nav-ch17-learner-0169}{subsubsection}{Cook 38.3}{0}
% mdparagraph:P604:end

% item: EX-XII-XIII-38-03
% mdparagraph:P605:begin
Then she heard the fame (hlīsa) of Solomon (Salomones), and came from the southern (sūðerne, 59) boundaries to Solomon within Jerusalem, and her camels bore southern spices [herbs], and precious gems, and a quantity (unrīm) of gold (154. a).
% mdparagraph:P605:end

\end{minipage}\par\addvspace{4pt}

% teaching: NEW-U17-GRAM-U17-SUPPLIED-WITHIN-PLACE
% teaching: NEW-U17-GRAM-U17-SUPPLIED-UNRIM-GOLD
% mdparagraph:P606:begin
\textbf{Construction.} Use the bracketed herbs meaning. Cook §154.a concerns the partitive genitive after a quantity: goldes depends on unrīm. The feminine accusative plural adjective for southern follows its taught inflection; its dictionary form in the hint is not automatically the answer form.
% mdparagraph:P606:end

% teaching: U17-SUP-SALOMON-GEN-SG
% teaching: U17-SUP-SUTHERNE-ACC-F-PL
% teaching: U17-SUP-UNRIM-ACC-SG
% mdparagraph:P607:begin
\textbf{Supplied:} binnan Hierusalem, — Within Jerusalem. The foreign place form and the complete prepositional expression are supplied.
% mdparagraph:P607:end

% teaching: U17-SUP-BINNAN-INVARIANT
% teaching: U17-SUP-HIERUSALEM-DAT-SG
% mdparagraph:P608:begin
\textbf{Supplied:} unrīm goldes. — An immense quantity of gold. Genitive singular goldes depends on the quantity noun unrīm.
% mdparagraph:P608:end

% mdparagraph:P609:begin
\textbf{Supplied:} Salomones — Solomon; genitive, singular.
% mdparagraph:P609:end

% teaching: U17-SUP-SUTHERNE-DAT-PL
% mdparagraph:P610:begin
\textbf{Supplied:} sūðernum — southern; dative, neuter, plural, strong.
% mdparagraph:P610:end

% teaching: U17-SUP-SALOMON-DAT-SG
% mdparagraph:P611:begin
\textbf{Supplied:} Salomone — Solomon; dative, singular.
% mdparagraph:P611:end

% mdparagraph:P612:begin
\textbf{Supplied:} binnan — within.
% mdparagraph:P612:end

% mdparagraph:P613:begin
\textbf{Supplied:} Hierusalem — Jerusalem; dative, singular.
% mdparagraph:P613:end

% mdparagraph:P614:begin
\textbf{Supplied:} sūðerne — southern; accusative, feminine, plural, strong.
% mdparagraph:P614:end

% mdparagraph:P615:begin
\textbf{Supplied:} unrīm — an uncounted quantity; abundance; accusative, neuter, singular.
% mdparagraph:P615:end

% output: EX-XII-XIII-38-03-O001
% mdparagraph:P616:begin
Write one complete Old English translation, keeping the printed controls and the supplied expressions available.
% mdparagraph:P616:end

\par\begin{minipage}{\linewidth}\raggedright
% mdparagraph:P617:begin
\subsection*{Cook 38.4}\BookHeading{nav-ch17-learner-0170}{subsubsection}{Cook 38.4}{0}
% mdparagraph:P617:end

% item: EX-XII-XIII-38-04
% mdparagraph:P618:begin
The queen then (ðā) had speech with (wið) Solomon, and said (om. to) him whatsoever she thought in her heart (verb last).
% mdparagraph:P618:end

\end{minipage}\par\addvspace{4pt}

% teaching: NEW-U17-GRAM-U17-H-38-04
% teaching: NEW-U17-LOCAL-SENSE-WITH-U17-CONVERSATION
% mdparagraph:P619:begin
\textbf{Construction.} Omit a separate to with the following object and put the finite verb last in the specified dependent clause. The supplied conversation frame gives wið its local meaning.
% mdparagraph:P619:end

% mdparagraph:P620:begin
\textbf{Supplied:} Sēo cwēn hæfde ðā sprǣce wið Salomon — The queen then had speech with Solomon. Wið identifies the conversational partner. This supplied arrangement selects accusative Salomon; dative Salomone is also supported for this conversation frame.
% mdparagraph:P620:end

% teaching: U17-SUP-SALOMON-ACC-SG
% mdparagraph:P621:begin
\textbf{Supplied:} wið — with, the conversational counterpart in the supplied speech frame.
% mdparagraph:P621:end

% mdparagraph:P622:begin
\textbf{Supplied:} Salomon — Solomon; accusative, singular.
% mdparagraph:P622:end

% output: EX-XII-XIII-38-04-O001
% mdparagraph:P623:begin
Write one complete Old English translation, keeping the printed controls and the supplied expressions available.
% mdparagraph:P623:end

\par\begin{minipage}{\linewidth}\raggedright
% mdparagraph:P624:begin
\subsection*{Cook 38.5}\BookHeading{nav-ch17-learner-0171}{subsubsection}{Cook 38.5}{0}
% mdparagraph:P624:end

% item: EX-XII-XIII-38-05
% mdparagraph:P625:begin
Solomon then taught her, and told (sęcgean) her the meaning of all the words that (87. d) she asked (āscian) him.
% mdparagraph:P625:end

% mdparagraph:P626:begin
\textbf{Construction.} Cook §87.d points to invariant relative ðe. Work out its clause role from the verb and antecedent; its shape does not show a case ending.
% mdparagraph:P626:end

\end{minipage}\par\addvspace{4pt}

% teaching: U17-SUP-SALOMON-NOM-SG
% mdparagraph:P627:begin
\textbf{Supplied:} Salomon — Solomon; nominative, singular.
% mdparagraph:P627:end

% output: EX-XII-XIII-38-05-O001
% mdparagraph:P628:begin
Write one complete Old English translation, keeping the printed controls and the supplied expressions available.
% mdparagraph:P628:end

\par\begin{minipage}{\linewidth}\raggedright
% mdparagraph:P629:begin
\subsection*{Cook 38.6}\BookHeading{nav-ch17-learner-0172}{subsubsection}{Cook 38.6}{0}
% mdparagraph:P629:end

% item: EX-XII-XIII-38-06
% mdparagraph:P630:begin
Then the queen saw the wisdom of Solomon, and the great (mǣre) temple that (87. d) he (om. had) built, and the king’s manifold services.
% mdparagraph:P630:end

% mdparagraph:P631:begin
\textbf{Construction.} Use invariant relative ðe and retain the complete list of what the queen saw. Omit the separately marked had.
% mdparagraph:P631:end

% mdparagraph:P632:begin
\textbf{Supplied:} Salomones — Solomon; genitive, singular.
% mdparagraph:P632:end

% output: EX-XII-XIII-38-06-O001
% mdparagraph:P633:begin
Write one complete Old English translation, keeping the printed controls and the supplied expressions available.
% mdparagraph:P633:end

\end{minipage}\par\addvspace{4pt}

\par\begin{minipage}{\linewidth}\raggedright
% mdparagraph:P634:begin
\subsection*{Cook 38.7}\BookHeading{nav-ch17-learner-0173}{subsubsection}{Cook 38.7}{0}
% mdparagraph:P634:end

% item: EX-XII-XIII-38-07
% mdparagraph:P635:begin
She was so (tō ðan) greatly astonished (ofwundrian, 118), that she had no spirit more (furðor).
% mdparagraph:P635:end

\end{minipage}\par\addvspace{4pt}

% teaching: NEW-U17-LOCAL-SENSE-TO-U17-DEGREE-RESULT
% teaching: NEW-U17-LOCAL-REFERENCE-SE-U17-DEGREE-REFERENCE-THAN
% teaching: NEW-U17-GRAM-U17-SUPPLIED-DEGREE-RESULT
% teaching: U17-SUP-FURTHOR-INVARIANT
% mdparagraph:P636:begin
\textbf{Supplied:} tō ðan swīðe ofwundrod, — So greatly astonished. Tō ðan gives the degree and the following ðæt clause its result. Keep the exact supplied ðan.
% mdparagraph:P636:end

% mdparagraph:P637:begin
\textbf{Supplied:} tō — to that degree, within supplied tō ðan … ðæt.
% mdparagraph:P637:end

% mdparagraph:P638:begin
\textbf{Supplied:} ðan — that degree, in tō ðan.
% mdparagraph:P638:end

% mdparagraph:P639:begin
\textbf{Supplied:} furðor — further; any more.
% mdparagraph:P639:end

% output: EX-XII-XIII-38-07-O001
% mdparagraph:P640:begin
Write one complete Old English translation, keeping the printed controls and the supplied expressions available.
% mdparagraph:P640:end

\par\begin{minipage}{\linewidth}\raggedright
% mdparagraph:P641:begin
\subsection*{Cook 38.8}\BookHeading{nav-ch17-learner-0174}{subsubsection}{Cook 38.8}{0}
% mdparagraph:P641:end

% item: EX-XII-XIII-38-08
% mdparagraph:P642:begin
She said then to (tō) the king, True is the word that (87. d) I heard in my country concerning thee and thy wisdom (wīsdōm); but I would not believe before I (om. my)self saw (196. b) it (verb last).
% mdparagraph:P642:end

% mdparagraph:P643:begin
\textbf{Construction.} Keep self while omitting my. Cook §196.b supports the past subjunctive in the before-clause, viewed from the earlier unwillingness. Put that finite verb last.
% mdparagraph:P643:end

% mdparagraph:P644:begin
\textbf{Supplied:} tō — to, the person addressed in speech.
% mdparagraph:P644:end

\end{minipage}\par\addvspace{4pt}

% teaching: U17-SUP-SOTH-CITATION-NOM-N-SG
% mdparagraph:P645:begin
\textbf{Supplied:} Sōð — true; nominative, neuter, singular, strong.
% mdparagraph:P645:end

% output: EX-XII-XIII-38-08-O001
% mdparagraph:P646:begin
Write one complete Old English translation, keeping the printed controls and the supplied expressions available.
% mdparagraph:P646:end

\par\begin{minipage}{\linewidth}\raggedright
% mdparagraph:P647:begin
\subsection*{Cook 38.9}\BookHeading{nav-ch17-learner-0175}{subsubsection}{Cook 38.9}{0}
% mdparagraph:P647:end

% item: EX-XII-XIII-38-09
% mdparagraph:P648:begin
Greater (micel, 66) is thy wisdom and thy work than the report (hlīsa) was (wǣre) that I heard.
% mdparagraph:P648:end

\end{minipage}\par\addvspace{4pt}

% teaching: NEW-U17-GRAM-U17-H-38-09
% mdparagraph:P649:begin
\textbf{Construction.} The printed wǣre is supplied exactly. It belongs to the report in the comparison; it does not agree with wisdom and works. The whole expression below construes wisdom and its work together.
% mdparagraph:P649:end

% mdparagraph:P650:begin
\textbf{Supplied:} Māra is ðīn wīsdōm and ðīn weorc ðonne se hlīsa wǣre ðe ic hīerde. — Your wisdom and your work are greater than the report I heard. In this supplied arrangement, singular Māra is construes wisdom and its work together. The exact wǣre is third singular past subjunctive in the comparison, referring to the report. Keep that given form available; a defensible plural arrangement of the main clause can also be grammatical.
% mdparagraph:P650:end

% output: EX-XII-XIII-38-09-O001
% mdparagraph:P651:begin
Write one complete Old English translation, keeping the printed controls and the supplied expressions available.
% mdparagraph:P651:end

\par\begin{minipage}{\linewidth}\raggedright
% mdparagraph:P652:begin
\subsection*{Cook 38.10}\BookHeading{nav-ch17-learner-0176}{subsubsection}{Cook 38.10}{0}
% mdparagraph:P652:end

% item: EX-XII-XIII-38-10
% mdparagraph:P653:begin
Blessed (ēadig) are thy vassals and thy servants (ðēowa, 53), that always (simle) stand before (ætforan) thee, and hear thy wisdom.
% mdparagraph:P653:end

% mdparagraph:P654:begin
\textbf{Supplied:} simle — always.
% mdparagraph:P654:end

% mdparagraph:P655:begin
\textbf{Supplied:} ætforan — before; in front of.
% mdparagraph:P655:end

% output: EX-XII-XIII-38-10-O001
% mdparagraph:P656:begin
Write one complete Old English translation, keeping the printed controls and the supplied expressions available.
% mdparagraph:P656:end

\end{minipage}\par\addvspace{4pt}

\reserve{7\baselineskip}
% mdparagraph:P657:begin
\section*{Cook 39. Praise, gifts and a change of subject}\BookHeading{nav-ch17-learner-0177}{subsection}{Cook 39. Praise, gifts and a change of subject}{1}
% mdparagraph:P657:end

% mdparagraph:P658:begin
The first four items continue praise and reciprocal gifts. Items 5–7 turn to building and to what eye, ear and heart can apprehend; they are a change of material, not a continuation of the queen’s visit.
% mdparagraph:P658:end

\par\begin{minipage}{\linewidth}\raggedright
% mdparagraph:P659:begin
\subsection*{Cook 39.1}\BookHeading{nav-ch17-learner-0178}{subsubsection}{Cook 39.1}{0}
% mdparagraph:P659:end

% item: EX-XII-XIII-39-01
% mdparagraph:P660:begin
Blessed (blētsian) be the almighty God, who (87. d) selected thee, and placed (gesęttan) thee over (ofer, 172) the kingdom of the Israelites (Israela), that thou (om. mightest) establish (sęttan) judgments and righteousness.
% mdparagraph:P660:end

% mdparagraph:P661:begin
\textbf{Construction.} Use the relative particle for “which”. In the that-clause, omit a separate mightest and express the intended purpose with the subjunctive.
% mdparagraph:P661:end

\end{minipage}\par\addvspace{4pt}

% teaching: U17-SUP-ISRAEL-GEN-PL
% teaching: U17-SUP-AELMIHTIG-NOM-M-SG-WEAK
% mdparagraph:P662:begin
\textbf{Supplied:} ælmihtiga — almighty; nominative, masculine, singular, weak.
% mdparagraph:P662:end

% mdparagraph:P663:begin
\textbf{Supplied:} Israela — Israelites, in the supplied genitive plural; genitive, plural.
% mdparagraph:P663:end

% mdparagraph:P664:begin
\textbf{Supplied:} dōmas — judgment; a judgment made in judging; accusative, plural.
% mdparagraph:P664:end

% output: EX-XII-XIII-39-01-O001
% mdparagraph:P665:begin
Write one complete Old English translation, keeping the printed controls and the supplied expressions available.
% mdparagraph:P665:end

\par\begin{minipage}{\linewidth}\raggedright
% mdparagraph:P666:begin
\subsection*{Cook 39.2}\BookHeading{nav-ch17-learner-0179}{subsubsection}{Cook 39.2}{0}
% mdparagraph:P666:end

% item: EX-XII-XIII-39-02
% mdparagraph:P667:begin
She gave (forgiefan) the king (om. a) hundred and twenty pounds (154. c) of gold (154. a), and (om. a) multitude of precious spices (wyrt) and of precious stones.
% mdparagraph:P667:end

\end{minipage}\par\addvspace{4pt}

% teaching: NEW-U17-GRAM-U17-SUPPLIED-120-GOLD
% teaching: NEW-U17-GRAM-U17-SUPPLIED-MENIGU-QUANTITY
% mdparagraph:P668:begin
\textbf{Construction.} Cook §154.c concerns a genitive after a higher numeral; §154.a concerns quantity plus genitive. Both complete quantity expressions are supplied below. Omit the marked indefinite articles.
% mdparagraph:P668:end

% mdparagraph:P669:begin
\textbf{Supplied:} hund and twēntig punda goldes — One hundred and twenty pounds of gold. Punda is genitive plural after the higher numeral; goldes is genitive singular naming the material. Use the complete quantity phrase.
% mdparagraph:P669:end

% mdparagraph:P670:begin
\textbf{Supplied:} męnigeo dēorwierðra wyrta and dēorwierðra stāna. — A great quantity of costly herbs and precious stones. The quantity noun is accusative singular, followed by both groups in the genitive plural. The complete phrase is supplied.
% mdparagraph:P670:end

% teaching: U17-SUP-MENIGU-ACC-SG-VARIANT
% mdparagraph:P671:begin
\textbf{Supplied:} hund — hundred, within the complete supplied quantity of gold.
% mdparagraph:P671:end

% mdparagraph:P672:begin
\textbf{Supplied:} twēntig — twenty.
% mdparagraph:P672:end

% mdparagraph:P673:begin
\textbf{Supplied:} męnigeo — company; crowd; large quantity; accusative, feminine, singular.
% mdparagraph:P673:end

% output: EX-XII-XIII-39-02-O001
% mdparagraph:P674:begin
Write one complete Old English translation, keeping the printed controls and the supplied expressions available.
% mdparagraph:P674:end

\par\begin{minipage}{\linewidth}\raggedright
% mdparagraph:P675:begin
\subsection*{Cook 39.3}\BookHeading{nav-ch17-learner-0180}{subsubsection}{Cook 39.3}{0}
% mdparagraph:P675:end

% item: EX-XII-XIII-39-03
% mdparagraph:P676:begin
Solomon also gave the queen whatsoever she desired (giernan) of (æt) him; and she returned (gewęndan) again (ongēan) to her country with her retainers.
% mdparagraph:P676:end

% mdparagraph:P677:begin
\textbf{Construction.} The requested thing is genitive with giernan; æt with the dative marks the person asked. Keep the main object’s role separate from that of the relative clause.
% mdparagraph:P677:end

\end{minipage}\par\addvspace{4pt}

% teaching: U17-SUP-ONGEAN-BACK
% mdparagraph:P678:begin
\textbf{Supplied:} Salomon — Solomon; nominative, singular.
% mdparagraph:P678:end

% mdparagraph:P679:begin
\textbf{Supplied:} ongēan — back; again.
% mdparagraph:P679:end

% output: EX-XII-XIII-39-03-O001
% mdparagraph:P680:begin
Write one complete Old English translation, keeping the printed controls and the supplied expressions available.
% mdparagraph:P680:end

\par\begin{minipage}{\linewidth}\raggedright
% mdparagraph:P681:begin
\subsection*{Cook 39.4}\BookHeading{nav-ch17-learner-0181}{subsubsection}{Cook 39.4}{0}
% mdparagraph:P681:end

% item: EX-XII-XIII-39-04
% mdparagraph:P682:begin
Solomon then was magnified (mǣrsian, 118) above (ofer, 166) all earthly kings; and all nations (ðēod) desired (gewilnian) to see [that they might see, 194. b] him, and hear his wisdom.
% mdparagraph:P682:end

\end{minipage}\par\addvspace{4pt}

% teaching: NEW-U17-GRAM-U17-H-39-04
% teaching: NEW-U17-LOCAL-SENSE-OFER-U17-RANK-DATIVE
% mdparagraph:P683:begin
\textbf{Construction.} The bracket replaces the infinitive “to see” with the finite clause “that they might see”, expressing what all nations desired. This desired-content clause takes the subjunctive (Cook §194.b). The rank comparison after ofer is dative as Cook §166 directs.
% mdparagraph:P683:end

% mdparagraph:P684:begin
\textbf{Supplied:} Salomon wæs ðā mǣrsod ofer eallum eorðlicum cyningum; — Solomon was then celebrated above all earthly kings. Ofer here expresses rank and takes the explicitly requested dative complement. Keep that local case choice separate from the accusative rank construction encountered earlier.
% mdparagraph:P684:end

% mdparagraph:P685:begin
\textbf{Supplied:} Salomon — Solomon; nominative, singular.
% mdparagraph:P685:end

% mdparagraph:P686:begin
\textbf{Supplied:} ofer — above in rank, with the explicitly supplied dative complement.
% mdparagraph:P686:end

% output: EX-XII-XIII-39-04-O001
% mdparagraph:P687:begin
Write one complete Old English translation, keeping the printed controls and the supplied expressions available.
% mdparagraph:P687:end

\par\begin{minipage}{\linewidth}\raggedright
% mdparagraph:P688:begin
\subsection*{Cook 39.5}\BookHeading{nav-ch17-learner-0182}{subsubsection}{Cook 39.5}{0}
% mdparagraph:P688:end

% item: EX-XII-XIII-39-05
% mdparagraph:P689:begin
Whosoever buildeth upon (ofer, 166) this foundation (grundweall, 43) gold, or silver, or precious stones, or trees, straw (strēaw, 47) or chaff (ceaf, 47), every one’s (ānes gehwilces) work shall be manifested.
% mdparagraph:P689:end

\end{minipage}\par\addvspace{4pt}

% teaching: U17-SUP-STREAW-ACC-SG
% teaching: U17-SUP-CEAF-ACC-SG
% mdparagraph:P690:begin
\textbf{Supplied:} strēaw — straw; accusative, neuter, singular.
% mdparagraph:P690:end

% mdparagraph:P691:begin
\textbf{Supplied:} ceaf — chaff; accusative, neuter, singular.
% mdparagraph:P691:end

% output: EX-XII-XIII-39-05-O001
% mdparagraph:P692:begin
Write one complete Old English translation, keeping the printed controls and the supplied expressions available.
% mdparagraph:P692:end

\par\begin{minipage}{\linewidth}\raggedright
% mdparagraph:P693:begin
\subsection*{Cook 39.6}\BookHeading{nav-ch17-learner-0183}{subsubsection}{Cook 39.6}{0}
% mdparagraph:P693:end

% item: EX-XII-XIII-39-06
% mdparagraph:P694:begin
If the building (getimbrung, 51. 3) continues (ðurhwunian, 118), and withstands the fire (164. m), then shall the workman receive (om. a) reward (edlēan, 47) from (æt) God for his work (153. e).
% mdparagraph:P694:end

% mdparagraph:P695:begin
\textbf{Construction.} Cook §164.m calls for the dative thing resisted. The received reward is also dative in the selected onfōn frame. Under §153.e the work is a genitive dependent on reward, expressing its grounds.
% mdparagraph:P695:end

% output: EX-XII-XIII-39-06-O001
% mdparagraph:P696:begin
Write one complete Old English translation, keeping the printed controls and the supplied expressions available.
% mdparagraph:P696:end

\end{minipage}\par\addvspace{4pt}

\par\begin{minipage}{\linewidth}\raggedright
% mdparagraph:P697:begin
\subsection*{Cook 39.7}\BookHeading{nav-ch17-learner-0184}{subsubsection}{Cook 39.7}{0}
% mdparagraph:P697:end

% item: EX-XII-XIII-39-07
% mdparagraph:P698:begin
No (183) eye can in this life see, nor any (183) ear (53. 2) hear, nor any (183) man’s heart conceive (āsmēan, 113), the things that God prepares (gearcian, 118) (om. for) those that love him.
% mdparagraph:P698:end

\end{minipage}\par\addvspace{4pt}

% teaching: NEW-U17-GRAM-U17-H-39-07
% mdparagraph:P699:begin
\textbf{Construction.} Cook §183 calls for negative concord. Keep the subjects eye, ear and heart distinct; the first mæg can govern all three following infinitives. Ēare is singular weak neuter (§53.2). The beneficiaries are dative without a separate for.
% mdparagraph:P699:end

% teaching: NEW-U17-VOC-ASMEAN-U17-SOURCE-SENSE
% mdparagraph:P700:begin
\textbf{Supplied:} ne nānes mannes heorte āsmēan — Nor can any man’s heart conceive: the earlier mæg also governs this last infinitive. The heart is the final subject, and the things mentioned are its shared object. The exact printed āsmēan and this frame are supplied.
% mdparagraph:P700:end

% mdparagraph:P701:begin
\textbf{Supplied:} on — in worldly life, temporal/\allowbreak{}situational expression.
% mdparagraph:P701:end

% mdparagraph:P702:begin
\textbf{Supplied:} āsmēan — conceive; plain infinitive under the shared mæg, with the heart as subject.
% mdparagraph:P702:end

% output: EX-XII-XIII-39-07-O001
% mdparagraph:P703:begin
Write one complete Old English translation, keeping the printed controls and the supplied expressions available.
% mdparagraph:P703:end

\reserve{7\baselineskip}
% mdparagraph:P704:begin
\section*{Cook 40. Blessing and reversal}\BookHeading{nav-ch17-learner-0185}{subsection}{Cook 40. Blessing and reversal}{1}
% mdparagraph:P704:end

% mdparagraph:P705:begin
The prayer moves between praise, remembered acts, descriptions of people and a promise to descendants. Keep the beneficiaries, objects and predicates distinct.
% mdparagraph:P705:end

\par\begin{minipage}{\linewidth}\raggedright
% mdparagraph:P706:begin
\subsection*{Cook 40.1}\BookHeading{nav-ch17-learner-0186}{subsubsection}{Cook 40.1}{0}
% mdparagraph:P706:end

% item: EX-XII-XIII-40-01
% mdparagraph:P707:begin
The angel saluted Mary, and said that she was (194. b) filled (āfyllan, 113) with God’s grace [gift], and that God was with her (preposition last); and she was blessed (blētsian) among women.
% mdparagraph:P707:end

% mdparagraph:P708:begin
\textbf{Construction.} Cook §194.b calls for the subjunctive in the reported content. Use the bracketed gift meaning and put the specified preposition after its complement; its government is unchanged.
% mdparagraph:P708:end

\end{minipage}\par\addvspace{4pt}

% teaching: U17-SUP-MARIA-ACC-SG
% mdparagraph:P709:begin
\textbf{Supplied:} Marian — Mary; accusative, singular.
% mdparagraph:P709:end

% teaching: U17-SUP-GIEFU-DAT-SG-GRACE
% mdparagraph:P710:begin
\textbf{Supplied:} giefe — gift of grace; divine favour; dative, feminine, singular.
% mdparagraph:P710:end

% teaching: U17-SUP-BETWEOX-INVARIANT
% mdparagraph:P711:begin
\textbf{Supplied:} betwēox wīfum — Among women. Wīfum is dative plural after betwēox in this expression. Keep the whole phrase available; its government is supplied here.
% mdparagraph:P711:end

% output: EX-XII-XIII-40-01-O001
% mdparagraph:P712:begin
Write one complete Old English translation, keeping the printed controls and the supplied expressions available.
% mdparagraph:P712:end

\par\begin{minipage}{\linewidth}\raggedright
% mdparagraph:P713:begin
\subsection*{Cook 40.2}\BookHeading{nav-ch17-learner-0187}{subsubsection}{Cook 40.2}{0}
% mdparagraph:P713:end

% item: EX-XII-XIII-40-02
% mdparagraph:P714:begin
My soul (om. doth) magnify (mǣrsian) the Lord, and my spirit (om. hath) rejoiced in God my Saviour.
% mdparagraph:P714:end

\end{minipage}\par\addvspace{4pt}

% teaching: NEW-U17-GRAM-U17-H-40-02
% teaching: NEW-U17-LOCAL-SENSE-ON-U17-REJOICING
% mdparagraph:P715:begin
\textbf{Construction.} Omit separate hath and doth. The supplied rejoice-in expression gives on its local relation to the source of joy.
% mdparagraph:P715:end

% mdparagraph:P716:begin
\textbf{Supplied:} mīn gāst blissode on Gode mīnum Hǣlende. — My spirit rejoiced in God my Saviour. On with the dative names the source or object of joy in this supplied expression.
% mdparagraph:P716:end

% mdparagraph:P717:begin
\textbf{Supplied:} on — in, the source or object of joy in the supplied rejoice-in frame.
% mdparagraph:P717:end

% output: EX-XII-XIII-40-02-O001
% mdparagraph:P718:begin
Write one complete Old English translation, keeping the printed controls and the supplied expressions available.
% mdparagraph:P718:end

\par\begin{minipage}{\linewidth}\raggedright
% mdparagraph:P719:begin
\subsection*{Cook 40.3}\BookHeading{nav-ch17-learner-0188}{subsubsection}{Cook 40.3}{0}
% mdparagraph:P719:end

% item: EX-XII-XIII-40-03
% mdparagraph:P720:begin
For he regarded (gesēon) the humility of his handmaiden (ðīnen, 51. 5); verily (sōðlice), hence forth all generations (cnēoris) shall call (sęcgean) me blessed (ēadig).
% mdparagraph:P720:end

% mdparagraph:P721:begin
\textbf{Construction.} The plural generations are the subject of call. The first-person object has its own agreeing feminine predicate adjective.
% mdparagraph:P721:end

\end{minipage}\par\addvspace{4pt}

% teaching: U17-SUP-HEONAN-FORTH-INVARIANT
% mdparagraph:P722:begin
\textbf{Supplied:} heonanforð — henceforth.
% mdparagraph:P722:end

% output: EX-XII-XIII-40-03-O001
% mdparagraph:P723:begin
Write one complete Old English translation, keeping the printed controls and the supplied expressions available.
% mdparagraph:P723:end

\par\begin{minipage}{\linewidth}\raggedright
% mdparagraph:P724:begin
\subsection*{Cook 40.4}\BookHeading{nav-ch17-learner-0189}{subsubsection}{Cook 40.4}{0}
% mdparagraph:P724:end

% item: EX-XII-XIII-40-04
% mdparagraph:P725:begin
For he that (sē ðe) is mighty hath done (om. to) me great things; and holy is his name.
% mdparagraph:P725:end

% output: EX-XII-XIII-40-04-O001
% mdparagraph:P726:begin
Write one complete Old English translation, keeping the printed controls and the supplied expressions available.
% mdparagraph:P726:end

\end{minipage}\par\addvspace{4pt}

\par\begin{minipage}{\linewidth}\raggedright
% mdparagraph:P727:begin
\subsection*{Cook 40.5}\BookHeading{nav-ch17-learner-0190}{subsubsection}{Cook 40.5}{0}
% mdparagraph:P727:end

% item: EX-XII-XIII-40-05
% mdparagraph:P728:begin
And his mercy is (om. on those) fearing (161) him from (of) generation to (on) generation.
% mdparagraph:P728:end

\end{minipage}\par\addvspace{4pt}

% teaching: NEW-U17-GRAM-U17-H-40-05
% teaching: NEW-U17-LOCAL-SENSE-OF-U17-GENERATION-RANGE
% teaching: NEW-U17-LOCAL-SENSE-ON-U17-GENERATION-RANGE
% mdparagraph:P729:begin
\textbf{Construction.} Cook §161 marks the dative beneficiaries. Omit on those, but retain the people through the participle. Their fearing still takes accusative him. The generation-to-generation expression is supplied as a whole.
% mdparagraph:P729:end

% mdparagraph:P730:begin
\textbf{Supplied:} of cnēorisse on cnēorisse. — From generation to generation. Both cnēorisse forms are dative singular; of marks the starting generation and on the succeeding extent. Use the whole supplied temporal expression.
% mdparagraph:P730:end

% mdparagraph:P731:begin
\textbf{Supplied:} of — from, the starting generation in the supplied temporal range.
% mdparagraph:P731:end

% mdparagraph:P732:begin
\textbf{Supplied:} on — to, the succeeding generation in the supplied temporal range.
% mdparagraph:P732:end

% output: EX-XII-XIII-40-05-O001
% mdparagraph:P733:begin
Write one complete Old English translation, keeping the printed controls and the supplied expressions available.
% mdparagraph:P733:end

\par\begin{minipage}{\linewidth}\raggedright
% mdparagraph:P734:begin
\subsection*{Cook 40.6}\BookHeading{nav-ch17-learner-0191}{subsubsection}{Cook 40.6}{0}
% mdparagraph:P734:end

% item: EX-XII-XIII-40-06
% mdparagraph:P735:begin
He hath put down (āweorpan) the powerful from (of) (om. their) seat, and lifted up (ūpāhębban) the lowly.
% mdparagraph:P735:end

% mdparagraph:P736:begin
\textbf{Construction.} The mighty and the humble are weak masculine plural adjectives naming people. Omit the marked their without deleting the seat.
% mdparagraph:P736:end

% output: EX-XII-XIII-40-06-O001
% mdparagraph:P737:begin
Write one complete Old English translation, keeping the printed controls and the supplied expressions available.
% mdparagraph:P737:end

\end{minipage}\par\addvspace{4pt}

\par\begin{minipage}{\linewidth}\raggedright
% mdparagraph:P738:begin
\subsection*{Cook 40.7}\BookHeading{nav-ch17-learner-0192}{subsubsection}{Cook 40.7}{0}
% mdparagraph:P738:end

% item: EX-XII-XIII-40-07
% mdparagraph:P739:begin
He hath filled the hungering (om. ones) with good (om. things), and the rich (rīce) he hath let go empty.
% mdparagraph:P739:end

% mdparagraph:P740:begin
\textbf{Construction.} The hungry and the rich name people without added ones. Good things uses the familiar plural noun for goods; omit a separate things.
% mdparagraph:P740:end

% output: EX-XII-XIII-40-07-O001
% mdparagraph:P741:begin
Write one complete Old English translation, keeping the printed controls and the supplied expressions available.
% mdparagraph:P741:end

\end{minipage}\par\addvspace{4pt}

\par\begin{minipage}{\linewidth}\raggedright
% mdparagraph:P742:begin
\subsection*{Cook 40.8}\BookHeading{nav-ch17-learner-0193}{subsubsection}{Cook 40.8}{0}
% mdparagraph:P742:end

% item: EX-XII-XIII-40-08
% mdparagraph:P743:begin
He hath accepted his servant (cniht) Israel, and remembered his mercy.
% mdparagraph:P743:end

% mdparagraph:P744:begin
\textbf{Supplied:} Israel — Israel; dative, singular.
% mdparagraph:P744:end

% output: EX-XII-XIII-40-08-O001
% mdparagraph:P745:begin
Write one complete Old English translation, keeping the printed controls and the supplied expressions available.
% mdparagraph:P745:end

\end{minipage}\par\addvspace{4pt}

\par\begin{minipage}{\linewidth}\raggedright
% mdparagraph:P746:begin
\subsection*{Cook 40.9}\BookHeading{nav-ch17-learner-0194}{subsubsection}{Cook 40.9}{0}
% mdparagraph:P746:end

% item: EX-XII-XIII-40-09
% mdparagraph:P747:begin
As he spake (sprecan) to (tō) our fathers, Abraham and his seed, for ever and ever.
% mdparagraph:P747:end

% mdparagraph:P748:begin
\textbf{Supplied:} tō — to, the person addressed in speech.
% mdparagraph:P748:end

\end{minipage}\par\addvspace{4pt}

% teaching: U17-SUP-ABRAHAM-DAT-SG
% mdparagraph:P749:begin
\textbf{Supplied:} Abrahame — Abraham; dative, singular.
% mdparagraph:P749:end

% mdparagraph:P750:begin
\textbf{Supplied:} ā — ever; always.
% mdparagraph:P750:end

% mdparagraph:P751:begin
\textbf{Supplied:} ā — ever; always.
% mdparagraph:P751:end

% mdparagraph:P752:begin
\textbf{Supplied:} ā — ever; always.
% mdparagraph:P752:end

% output: EX-XII-XIII-40-09-O001
% mdparagraph:P753:begin
Write one complete Old English translation, keeping the printed controls and the supplied expressions available.
% mdparagraph:P753:end

\par\begin{minipage}{\linewidth}\raggedright
% mdparagraph:P754:begin
\subsection*{Cook 40.10}\BookHeading{nav-ch17-learner-0195}{subsubsection}{Cook 40.10}{0}
% mdparagraph:P754:end

% item: EX-XII-XIII-40-10
% mdparagraph:P755:begin
Mary dwelt with her about (swilce) three months, and returned (gewęndan) then to her own house.
% mdparagraph:P755:end

\end{minipage}\par\addvspace{4pt}

% teaching: NEW-U17-GRAM-U17-SUPPLIED-THREE-MONTHS
% teaching: U17-SUP-SWILCE-APPROXIMATE
% mdparagraph:P756:begin
\textbf{Supplied:} swilce ðrīe mōnaðas — About three months. Swilce gives approximation and mōnaðas is accusative plural of duration; use this complete supplied phrase.
% mdparagraph:P756:end

% teaching: U17-SUP-MONATH-ACC-PL
% teaching: U17-SUP-MARIA-NOM-SG
% mdparagraph:P757:begin
\textbf{Supplied:} Maria — Mary; nominative, singular.
% mdparagraph:P757:end

% mdparagraph:P758:begin
\textbf{Supplied:} swilce — about; approximately.
% mdparagraph:P758:end

% mdparagraph:P759:begin
\textbf{Supplied:} mōnaðas — months; accusative, masculine, plural.
% mdparagraph:P759:end

% mdparagraph:P760:begin
\textbf{Supplied:} āgnum — own; dative, neuter, singular, strong.
% mdparagraph:P760:end

% output: EX-XII-XIII-40-10-O001
% mdparagraph:P761:begin
Write one complete Old English translation, keeping the printed controls and the supplied expressions available.
% mdparagraph:P761:end

\reserve{7\baselineskip}
% mdparagraph:P762:begin
\section*{Cook 41. Petitions and explanation}\BookHeading{nav-ch17-learner-0196}{subsection}{Cook 41. Petitions and explanation}{1}
% mdparagraph:P762:end

% mdparagraph:P763:begin
Identify each addressee, the requested act and its recipients. The asking and receiving constructions have their own case patterns.
% mdparagraph:P763:end

\par\begin{minipage}{\linewidth}\raggedright
% mdparagraph:P764:begin
\subsection*{Cook 41.1}\BookHeading{nav-ch17-learner-0197}{subsubsection}{Cook 41.1}{0}
% mdparagraph:P764:end

% item: EX-XII-XIII-41-01
% mdparagraph:P765:begin
God was ever, and ever was his wisdom begotten of (of) him, and the wisdom is his Son, just as mighty as the Father.
% mdparagraph:P765:end

\end{minipage}\par\addvspace{4pt}

% teaching: NEW-U17-GRAM-U17-H-41-01
% teaching: NEW-U17-LOCAL-SENSE-OF-U17-BEGETTING
% mdparagraph:P766:begin
\textbf{Supplied:} ǣfre wæs his wīsdōm of him ācęnned; — His wisdom was always begotten of him. The complete theological relation is supplied: of him gives the origin in this expression. It does not establish a general passive-agent use of of.
% mdparagraph:P766:end

% teaching: U17-SUP-AEFRE-INVARIANT
% mdparagraph:P767:begin
\textbf{Supplied:} ǣfre — ever; always.
% mdparagraph:P767:end

% mdparagraph:P768:begin
\textbf{Supplied:} ǣfre — ever; always.
% mdparagraph:P768:end

% mdparagraph:P769:begin
\textbf{Supplied:} of — of, the origin of wisdom in the supplied begetting statement.
% mdparagraph:P769:end

% output: EX-XII-XIII-41-01-O001
% mdparagraph:P770:begin
Write one complete Old English translation, keeping the printed controls and the supplied expressions available.
% mdparagraph:P770:end

\par\begin{minipage}{\linewidth}\raggedright
% mdparagraph:P771:begin
\subsection*{Cook 41.2}\BookHeading{nav-ch17-learner-0198}{subsubsection}{Cook 41.2}{0}
% mdparagraph:P771:end

% item: EX-XII-XIII-41-02
% mdparagraph:P772:begin
The Spirit of God filleth all the circuit of the world, and he maintaineth all things.
% mdparagraph:P772:end

% output: EX-XII-XIII-41-02-O001
% mdparagraph:P773:begin
Write one complete Old English translation, keeping the printed controls and the supplied expressions available.
% mdparagraph:P773:end

\end{minipage}\par\addvspace{4pt}

\par\begin{minipage}{\linewidth}\raggedright
% mdparagraph:P774:begin
\subsection*{Cook 41.3}\BookHeading{nav-ch17-learner-0199}{subsubsection}{Cook 41.3}{0}
% mdparagraph:P774:end

% item: EX-XII-XIII-41-03
% mdparagraph:P775:begin
He that doeth (wyrcean) my Father’s will, he is my brother, and my mother, and my sister.
% mdparagraph:P775:end

% output: EX-XII-XIII-41-03-O001
% mdparagraph:P776:begin
Write one complete Old English translation, keeping the printed controls and the supplied expressions available.
% mdparagraph:P776:end

\end{minipage}\par\addvspace{4pt}

\par\begin{minipage}{\linewidth}\raggedright
% mdparagraph:P777:begin
\subsection*{Cook 41.4}\BookHeading{nav-ch17-learner-0200}{subsubsection}{Cook 41.4}{0}
% mdparagraph:P777:end

% item: EX-XII-XIII-41-04
% mdparagraph:P778:begin
We brought (bringan) no thing to this world, nor can we take (lǣdan) any (183) thing hence with us.
% mdparagraph:P778:end

% mdparagraph:P779:begin
\textbf{Construction.} Cook §183 calls for negative concord. Retain the specified negative relation through both parts of the sentence.
% mdparagraph:P779:end

\end{minipage}\par\addvspace{4pt}

% teaching: U17-SUP-HEONAN-INVARIANT-VARIANT
% mdparagraph:P780:begin
\textbf{Supplied:} heonon — hence; from here.
% mdparagraph:P780:end

% output: EX-XII-XIII-41-04-O001
% mdparagraph:P781:begin
Write one complete Old English translation, keeping the printed controls and the supplied expressions available.
% mdparagraph:P781:end

\par\begin{minipage}{\linewidth}\raggedright
% mdparagraph:P782:begin
\subsection*{Cook 41.5}\BookHeading{nav-ch17-learner-0201}{subsubsection}{Cook 41.5}{0}
% mdparagraph:P782:end

% item: EX-XII-XIII-41-05
% mdparagraph:P783:begin
Ask (biddan), and (om. it shall) be given you; seek, and ye shall find; knock (cnucian, 118), and it shall be opened (om. to) you.
% mdparagraph:P783:end

\end{minipage}\par\addvspace{4pt}

% teaching: NEW-U17-GRAM-U17-H-41-05
% mdparagraph:P784:begin
\textbf{Construction.} Use plural commands. The supplied impersonal passive expressions take dative addressees; do not insert the omitted it shall.
% mdparagraph:P784:end

% mdparagraph:P785:begin
\textbf{Supplied:} ēow bið seald; — It will be given to you. The addressees are dative plural; singular bið has future reference in this supplied impersonal passive.
% mdparagraph:P785:end

% mdparagraph:P786:begin
\textbf{Supplied:} ēow bið geopenod. — It will be opened to you. The addressees are dative plural; singular bið has future reference in this supplied impersonal passive.
% mdparagraph:P786:end

% output: EX-XII-XIII-41-05-O001
% mdparagraph:P787:begin
Write one complete Old English translation, keeping the printed controls and the supplied expressions available.
% mdparagraph:P787:end

\par\begin{minipage}{\linewidth}\raggedright
% mdparagraph:P788:begin
\subsection*{Cook 41.6}\BookHeading{nav-ch17-learner-0202}{subsubsection}{Cook 41.6}{0}
% mdparagraph:P788:end

% item: EX-XII-XIII-41-06
% mdparagraph:P789:begin
Our Father, thou which (87. d) art in heaven(s), be thy name hallowed.
% mdparagraph:P789:end

% mdparagraph:P790:begin
\textbf{Construction.} Use the indicated plural heavens. In “thou who art”, invariant ðe leaves the verb second singular with the addressed thou.
% mdparagraph:P790:end

% output: EX-XII-XIII-41-06-O001
% mdparagraph:P791:begin
Write one complete Old English translation, keeping the printed controls and the supplied expressions available.
% mdparagraph:P791:end

\end{minipage}\par\addvspace{4pt}

\par\begin{minipage}{\linewidth}\raggedright
% mdparagraph:P792:begin
\subsection*{Cook 41.7}\BookHeading{nav-ch17-learner-0203}{subsubsection}{Cook 41.7}{0}
% mdparagraph:P792:end

% item: EX-XII-XIII-41-07
% mdparagraph:P793:begin
Thy kingdom (rīce) come (tōcuman; verb first); thy will befall (geweorðan) on earth as (swā-swā) in heaven(s) (verb first).
% mdparagraph:P793:end

% mdparagraph:P794:begin
\textbf{Construction.} Put the finite verbs first as directed. The wishes are expressed by present subjunctives. The exact swā-swā comparison remains visible reference help.
% mdparagraph:P794:end

% mdparagraph:P795:begin
\textbf{Supplied:} swā-swā — just as; like.
% mdparagraph:P795:end

% output: EX-XII-XIII-41-07-O001
% mdparagraph:P796:begin
Write one complete Old English translation, keeping the printed controls and the supplied expressions available.
% mdparagraph:P796:end

\end{minipage}\par\addvspace{4pt}

\par\begin{minipage}{\linewidth}\raggedright
% mdparagraph:P797:begin
\subsection*{Cook 41.8}\BookHeading{nav-ch17-learner-0204}{subsubsection}{Cook 41.8}{0}
% mdparagraph:P797:end

% item: EX-XII-XIII-41-08
% mdparagraph:P798:begin
Give (sęllan) us to-day our daily (dæghwāmlīc, 55) bread; and forgive (forgiefan) us our debts (gylt, 43), as we forgive (164) our debtors (gyltend, 43. 6).
% mdparagraph:P798:end

% mdparagraph:P799:begin
\textbf{Construction.} Use a singular command to the addressee. Debts are the thing forgiven; the debtors are dative people forgiven (§164). Preserve the requested forgiveness relation.
% mdparagraph:P799:end

\end{minipage}\par\addvspace{4pt}

% teaching: U17-SUP-TODAEG-INVARIANT
% mdparagraph:P800:begin
\textbf{Supplied:} tōdæg — today.
% mdparagraph:P800:end

% teaching: U17-SUP-HLAF-ACC-SG
% mdparagraph:P801:begin
\textbf{Supplied:} hlāf — bread; accusative, masculine, singular.
% mdparagraph:P801:end

% output: EX-XII-XIII-41-08-O001
% mdparagraph:P802:begin
Write one complete Old English translation, keeping the printed controls and the supplied expressions available.
% mdparagraph:P802:end

\par\begin{minipage}{\linewidth}\raggedright
% mdparagraph:P803:begin
\subsection*{Cook 41.9}\BookHeading{nav-ch17-learner-0205}{subsubsection}{Cook 41.9}{0}
% mdparagraph:P803:end

% item: EX-XII-XIII-41-09
% mdparagraph:P804:begin
And lead thou us not into (on, 172) temptation (costnung), but deliver us from evil.
% mdparagraph:P804:end

\end{minipage}\par\addvspace{4pt}

% teaching: U17-SUP-YFEL-DAT-SG-EVIL
% mdparagraph:P805:begin
\textbf{Supplied:} yfele — evil; dative, neuter, singular.
% mdparagraph:P805:end

% output: EX-XII-XIII-41-09-O001
% mdparagraph:P806:begin
Write one complete Old English translation, keeping the printed controls and the supplied expressions available.
% mdparagraph:P806:end

\par\begin{minipage}{\linewidth}\raggedright
% mdparagraph:P807:begin
\subsection*{Cook 41.10}\BookHeading{nav-ch17-learner-0206}{subsubsection}{Cook 41.10}{0}
% mdparagraph:P807:end

% item: EX-XII-XIII-41-10
% mdparagraph:P808:begin
Be it so.
% mdparagraph:P808:end

% output: EX-XII-XIII-41-10-O001
% mdparagraph:P809:begin
Write one complete Old English translation, keeping the printed controls and the supplied expressions available.
% mdparagraph:P809:end

\end{minipage}\par\addvspace{4pt}

\par\begin{minipage}{\linewidth}\raggedright
% mdparagraph:P810:begin
\subsection*{Cook 41.11}\BookHeading{nav-ch17-learner-0207}{subsubsection}{Cook 41.11}{0}
% mdparagraph:P810:end

% item: EX-XII-XIII-41-11
% mdparagraph:P811:begin
What (hwilc) father will give to his child (164. a) a stone, if it asks (156. b; 159. b) bread of him?
% mdparagraph:P811:end

% mdparagraph:P812:begin
\textbf{Construction.} The recipient of give is dative (§164.a). Ask has an accusative person asked and a genitive thing requested (§156.b; §159.b); the child is the subject of that clause.
% mdparagraph:P812:end

\end{minipage}\par\addvspace{4pt}

% teaching: U17-SUP-HLAF-GEN-SG
% mdparagraph:P813:begin
\textbf{Supplied:} hlāfes — bread; genitive, masculine, singular.
% mdparagraph:P813:end

% output: EX-XII-XIII-41-11-O001
% mdparagraph:P814:begin
Write one complete Old English translation, keeping the printed controls and the supplied expressions available.
% mdparagraph:P814:end

\par\begin{minipage}{\linewidth}\raggedright
% mdparagraph:P815:begin
\subsection*{Cook 41.12}\BookHeading{nav-ch17-learner-0208}{subsubsection}{Cook 41.12}{0}
% mdparagraph:P815:end

% item: EX-XII-XIII-41-12
% mdparagraph:P816:begin
God wrote (āwrītan) the old law (ǣ) with his finger (finger, 43) on the stone tables (tabele, 53).
% mdparagraph:P816:end

\end{minipage}\par\addvspace{4pt}

% teaching: U17-SUP-AE-LAW-ACC-SG
% mdparagraph:P817:begin
\textbf{Supplied:} ǣ — law; accusative, feminine, singular.
% mdparagraph:P817:end

% teaching: U17-SUP-STANEN-DAT-F-PL
% mdparagraph:P818:begin
\textbf{Supplied:} stānenum — made of stone; dative, feminine, plural, weak.
% mdparagraph:P818:end

% output: EX-XII-XIII-41-12-O001
% mdparagraph:P819:begin
Write one complete Old English translation, keeping the printed controls and the supplied expressions available.
% mdparagraph:P819:end

\reserve{7\baselineskip}
% mdparagraph:P820:begin
\section*{Cook 42. Vision and figurative description}\BookHeading{nav-ch17-learner-0209}{subsection}{Cook 42. Vision and figurative description}{1}
% mdparagraph:P820:end

% mdparagraph:P821:begin
The longer vision sentences contain nested descriptions and source-specific number expressions. Their complete supplied phrases remain available. Separate the finite clauses from participles that describe perceived people.
% mdparagraph:P821:end

\par\begin{minipage}{\linewidth}\raggedright
% mdparagraph:P822:begin
\subsection*{Cook 42.1}\BookHeading{nav-ch17-learner-0210}{subsubsection}{Cook 42.1}{0}
% mdparagraph:P822:end

% item: EX-XII-XIII-42-01
% mdparagraph:P823:begin
The young man fell in love with the girl (p. 175, l. 20; cf. mīn lufu befielð on ðē, ‘my love falls on thee,’ ‘I fall in love with thee’).
% mdparagraph:P823:end

\end{minipage}\par\addvspace{4pt}

% teaching: NEW-U17-GRAM-U17-H-42-01
% teaching: NEW-U17-LOCAL-SENSE-ON-U17-LOVE-GOAL
% teaching: NEW-U17-LOCAL-REFERENCE-BEFEALLAN-U17-QUOTE-PRES-3-SG-REFERENCE
% mdparagraph:P824:begin
\textbf{Construction.} The exact quotation supplies the idiom: love is the subject, and on introduces the girl toward whom it is directed. The supplied past expression retains this structure.
% mdparagraph:P824:end

% mdparagraph:P825:begin
\textbf{Supplied:} Ðæs geongan mannes lufu befēoll on ðæt mǣden. — The young man fell in love with the girl: literally, the young man’s love fell upon the girl. Love is the nominative subject, the young man its genitive possessor, and the girl the accusative complement of on. Use the whole supplied past expression; the quoted present befielð in the prompt remains reference help.
% mdparagraph:P825:end

% teaching: U17-SUP-LUFU-NOM-SG
% teaching: NEW-U17-VOC-BEFEALLAN-U17-SOURCE-SENSE
% mdparagraph:P826:begin
\textbf{Supplied:} lufu — love; nominative, feminine, singular.
% mdparagraph:P826:end

% mdparagraph:P827:begin
\textbf{Supplied:} befēoll — fell upon, with love as the subject and the girl as its goal; indicative, preterite, third person, singular.
% mdparagraph:P827:end

% mdparagraph:P828:begin
\textbf{Supplied:} on — on/\allowbreak{}upon, the target of love in the supplied fall-in-love expression.
% mdparagraph:P828:end

% output: EX-XII-XIII-42-01-O001
% mdparagraph:P829:begin
Write one complete Old English translation, keeping the printed controls and the supplied expressions available.
% mdparagraph:P829:end

\par\begin{minipage}{\linewidth}\raggedright
% mdparagraph:P830:begin
\subsection*{Cook 42.2}\BookHeading{nav-ch17-learner-0211}{subsubsection}{Cook 42.2}{0}
% mdparagraph:P830:end

% item: EX-XII-XIII-42-02
% mdparagraph:P831:begin
I heard the voice of harpers harping with their harps; and they sung a new song before the throne, and before the four beasts; and no man could learn that song except the hundred and forty and four thousand.
% mdparagraph:P831:end

\end{minipage}\par\addvspace{4pt}

% teaching: NEW-U17-GRAM-U17-H-42-02-PART
% teaching: NEW-U17-GRAM-U17-H-42-02-NUMBER
% teaching: NEW-U17-LOCAL-REFERENCE-BUTAN-EXCEPT-U17-BUTON-REFERENCE
% teaching: NEW-U17-LOCAL-REFERENCE-HUNDTEONTIG-U17-REFERENCE
% teaching: NEW-U17-LOCAL-REFERENCE-THUSEND-U17-EXCLUSION-FINAL-O
% teaching: NEW-U17-GRAM-U17-SUPPLIED-144000-EXCLUSION
% mdparagraph:P832:begin
\textbf{Construction.} The voice is singular; the harpers are plural. Both the nested harper phrase and the entire exclusion of the 144,000 remain supplied. Retain nīwne sang and the later ðone sang with their visible support.
% mdparagraph:P832:end

% mdparagraph:P833:begin
\textbf{Supplied:} ðāra hearpera stefne mid hiera hearpum hearpiendra; — The voice of the harpers harping with their harps. Stefne is the accusative voice heard. Hearpera and hearpiendra are genitive plural; mid governs dative plural hearpum. The following they refers to the harpers. This complete nested description is supplied.
% mdparagraph:P833:end

% teaching: NEW-U17-VOC-HEARPIAN-U17-SOURCE-SENSE
% mdparagraph:P834:begin
\textbf{Supplied:} būton ðæt hundtēontig and fēowertig and fēower ðūsendo. — Except the one hundred and forty-four thousand. This entire exclusion phrase is supplied, including collective ðæt and final-o ðūsendo. It is not a dative chain. The number’s internal arrangement and unfamiliar forms require no independent recall.
% mdparagraph:P834:end

% mdparagraph:P835:begin
\textbf{Supplied:} hearpiendra — harping; describing the harpers; present participle, genitive, masculine, plural, strong.
% mdparagraph:P835:end

% teaching: U17-SUP-NIWE-ACC-M-SG
% mdparagraph:P836:begin
\textbf{Supplied:} nīwne — new; accusative, masculine, singular, strong.
% mdparagraph:P836:end

% mdparagraph:P837:begin
\textbf{Supplied:} būton — except, within the whole supplied exclusion phrase.
% mdparagraph:P837:end

% mdparagraph:P838:begin
\textbf{Supplied:} hundtēontig — hundred, within the whole supplied 144,000 expression.
% mdparagraph:P838:end

% mdparagraph:P839:begin
\textbf{Supplied:} fēowertig — forty.
% mdparagraph:P839:end

% mdparagraph:P840:begin
\textbf{Supplied:} ðūsendo — thousand, here within the whole supplied 144,000 expression.
% mdparagraph:P840:end

% output: EX-XII-XIII-42-02-O001
% mdparagraph:P841:begin
Write one complete Old English translation, keeping the printed controls and the supplied expressions available.
% mdparagraph:P841:end

\par\begin{minipage}{\linewidth}\raggedright
% mdparagraph:P842:begin
\subsection*{Cook 42.3}\BookHeading{nav-ch17-learner-0212}{subsubsection}{Cook 42.3}{0}
% mdparagraph:P842:end

% item: EX-XII-XIII-42-03
% mdparagraph:P843:begin
I saw so great (om. a) company (męnigu), as no man can number (gerīman), of all nations and of every tribe (mǣgð), standing before the throne (ðrymsetl) of God, all clothed (gescrȳdan) with white garments (gegierela), and holding palm-twigs (palmtwīg, 47) in their hands.
% mdparagraph:P843:end

\end{minipage}\par\addvspace{4pt}

% teaching: NEW-U17-GRAM-U17-H-42-03
% mdparagraph:P844:begin
\textbf{Construction.} The supplied expression keeps the grammatical singular company and the following descriptions of its many members. The participles describe perceived people; they do not become separate progressive finite verbs.
% mdparagraph:P844:end

% teaching: U17-SUP-MENIGU-ACC-SG
% teaching: NEW-U17-VOC-GERIMAN-U17-SOURCE-SENSE
% mdparagraph:P845:begin
\textbf{Supplied:} swā micle męnigu swā nān mann gerīman ne mæg, of eallum ðēodum and of ǣlcre mǣgðe, standende beforan Godes ðrymsetle, ealle mid hwītum gegierelum gescrȳdde and palmtwīg on hiera handum healdende. — So great a company as no man can number, from all nations and every tribe, standing before God’s throne, all clothed in white garments and holding palm branches in their hands. Męnigu is grammatically feminine singular, while the later plural expressions describe its members. Standende and healdende are nonfinite descriptions, and gescrȳdde is a plural predicate past participle. Keep the whole supplied expression; this particular shift from collective to members is not a general agreement rule.
% mdparagraph:P845:end

% mdparagraph:P846:begin
\textbf{Supplied:} męnigu — company; crowd; large quantity; accusative, feminine, singular.
% mdparagraph:P846:end

% mdparagraph:P847:begin
\textbf{Supplied:} gerīman — number; plain infinitive under mæg.
% mdparagraph:P847:end

% mdparagraph:P848:begin
\textbf{Supplied:} gescrȳdde — clothed; plural predicate past participle describing the company’s members.
% mdparagraph:P848:end

% mdparagraph:P849:begin
\textbf{Supplied:} healdende — holding; present participle describing the company’s members.
% mdparagraph:P849:end

% output: EX-XII-XIII-42-03-O001
% mdparagraph:P850:begin
Write one complete Old English translation, keeping the printed controls and the supplied expressions available.
% mdparagraph:P850:end

\par\begin{minipage}{\linewidth}\raggedright
% mdparagraph:P851:begin
\subsection*{Cook 42.4}\BookHeading{nav-ch17-learner-0213}{subsubsection}{Cook 42.4}{0}
% mdparagraph:P851:end

% item: EX-XII-XIII-42-04
% mdparagraph:P852:begin
And they sang with (om. a) loud (hlūd) voice, Be salvation (om. to) our God which sitteth upon (ofer, 166) his throne.
% mdparagraph:P852:end

% output: EX-XII-XIII-42-04-O001
% mdparagraph:P853:begin
Write one complete Old English translation, keeping the printed controls and the supplied expressions available.
% mdparagraph:P853:end

\end{minipage}\par\addvspace{4pt}

\par\begin{minipage}{\linewidth}\raggedright
% mdparagraph:P854:begin
\subsection*{Cook 42.5}\BookHeading{nav-ch17-learner-0214}{subsubsection}{Cook 42.5}{0}
% mdparagraph:P854:end

% item: EX-XII-XIII-42-05
% mdparagraph:P855:begin
And all (om. the) angels stood round about [in circuit of] his throne, and bowed (ālūtan, 103) to God, thus saying,
% mdparagraph:P855:end

\end{minipage}\par\addvspace{4pt}

% teaching: NEW-U17-GRAM-U17-SUPPLIED-ACCOMPANYING-SPEECH
% mdparagraph:P856:begin
\textbf{Construction.} Omit the marked article and retain the supplied “in circuit of” arrangement. Possessive his is not an article.
% mdparagraph:P856:end

% mdparagraph:P857:begin
\textbf{Supplied:} ðus cweðende — Thus saying. The present participle accompanies the angels in 42.5 or the soldiers in 42.8 and introduces their words; it is not another finite verb. Keep this whole accompanying-speech expression available.
% mdparagraph:P857:end

% output: EX-XII-XIII-42-05-O001
% mdparagraph:P858:begin
Write one complete Old English translation, keeping the printed controls and the supplied expressions available.
% mdparagraph:P858:end

\par\begin{minipage}{\linewidth}\raggedright
% mdparagraph:P859:begin
\subsection*{Cook 42.6}\BookHeading{nav-ch17-learner-0215}{subsubsection}{Cook 42.6}{0}
% mdparagraph:P859:end

% item: EX-XII-XIII-42-06
% mdparagraph:P860:begin
Be (om. to) our God blessing and brightness, wisdom and thanksgiving (ðancung), dignity and strength, for ever and ever.
% mdparagraph:P860:end

\end{minipage}\par\addvspace{4pt}

% teaching: NEW-U17-GRAM-U17-H-42-06
% mdparagraph:P861:begin
\textbf{Construction.} The entire singular attribution of blessing and praise is supplied. It is a local formula, not a rule that every coordinated subject takes a singular verb.
% mdparagraph:P861:end

% mdparagraph:P862:begin
\textbf{Supplied:} Ūrum Gode sīe blētsung and beorhtnes, wīsdōm and ðancung, weorðscipe and stręngðu, ā ā ā. — To our God be blessing and brightness, wisdom and thanksgiving, honour and strength, for ever. The recipient is dative and sīe is the singular optative. The whole coordinated formula, including ā ā ā, is supplied.
% mdparagraph:P862:end

% teaching: U17-SUP-STRENGTHU-NOM-SG
% mdparagraph:P863:begin
\textbf{Supplied:} stręngðu — strength; nominative, feminine, singular.
% mdparagraph:P863:end

% mdparagraph:P864:begin
\textbf{Supplied:} ā — ever; always.
% mdparagraph:P864:end

% mdparagraph:P865:begin
\textbf{Supplied:} ā — ever; always.
% mdparagraph:P865:end

% mdparagraph:P866:begin
\textbf{Supplied:} ā — ever; always.
% mdparagraph:P866:end

% output: EX-XII-XIII-42-06-O001
% mdparagraph:P867:begin
Write one complete Old English translation, keeping the printed controls and the supplied expressions available.
% mdparagraph:P867:end

\par\begin{minipage}{\linewidth}\raggedright
% mdparagraph:P868:begin
\subsection*{Cook 42.7}\BookHeading{nav-ch17-learner-0216}{subsubsection}{Cook 42.7}{0}
% mdparagraph:P868:end

% item: EX-XII-XIII-42-07
% mdparagraph:P869:begin
Foxes (fox, 43) have holes (hol, 47), and the birds of heaven nests (nest, 47); but (sōðlice) the Son of Man hath not where (hwǣr) he (om. may) lay (āhyldan, 113) his head.
% mdparagraph:P869:end

% mdparagraph:P870:begin
\textbf{Construction.} Omit a separate may. The subjunctive in the indefinite relative clause expresses a place where one might lay one’s head.
% mdparagraph:P870:end

\end{minipage}\par\addvspace{4pt}

% teaching: U17-SUP-HWAER-INVARIANT
% mdparagraph:P871:begin
\textbf{Supplied:} hwǣr — where.
% mdparagraph:P871:end

% output: EX-XII-XIII-42-07-O001
% mdparagraph:P872:begin
Write one complete Old English translation, keeping the printed controls and the supplied expressions available.
% mdparagraph:P872:end

\par\begin{minipage}{\linewidth}\raggedright
% mdparagraph:P873:begin
\subsection*{Cook 42.8}\BookHeading{nav-ch17-learner-0217}{subsubsection}{Cook 42.8}{0}
% mdparagraph:P873:end

% item: EX-XII-XIII-42-08
% mdparagraph:P874:begin
The soldiers platted (windan, 104) (om. a) crown of (of) thorns (ðorn, 43), and placed on (172) his head, and a reed (hrēod, 47) in his right (om. hand), and bowed their knee before him, thus saying, Hale be thou, King of the Jews (Iudea).
% mdparagraph:P874:end

% mdparagraph:P875:begin
\textbf{Construction.} Cook §172 specifies accusative complements in the marked on phrases, including placement on the head and spitting upon a person. Do not reduce all these meanings to one physical-motion rule. Omit the marked hand.
% mdparagraph:P875:end

\end{minipage}\par\addvspace{4pt}

% teaching: U17-SUP-IUDEA-GEN-PL
% teaching: U17-SUP-HAL-NOM-M-SG
% mdparagraph:P876:begin
\textbf{Supplied:} Hāl — hale; well, in a greeting; nominative, masculine, singular, strong.
% mdparagraph:P876:end

% mdparagraph:P877:begin
\textbf{Supplied:} Iudea — Jews, in the supplied genitive plural; genitive, plural.
% mdparagraph:P877:end

% mdparagraph:P878:begin
\textbf{Supplied:} ðus cweðende — Thus saying. The present participle accompanies the angels in 42.5 or the soldiers in 42.8 and introduces their words; it is not another finite verb. Keep this whole accompanying-speech expression available.
% mdparagraph:P878:end

% output: EX-XII-XIII-42-08-O001
% mdparagraph:P879:begin
Write one complete Old English translation, keeping the printed controls and the supplied expressions available.
% mdparagraph:P879:end

\par\begin{minipage}{\linewidth}\raggedright
% mdparagraph:P880:begin
\subsection*{Cook 42.9}\BookHeading{nav-ch17-learner-0218}{subsubsection}{Cook 42.9}{0}
% mdparagraph:P880:end

% item: EX-XII-XIII-42-09
% mdparagraph:P881:begin
And they spit (spǣtan, 113) on (on, 172) him, and took (om. a) reed and smote (bēatan) his head.
% mdparagraph:P881:end

% mdparagraph:P882:begin
\textbf{Construction.} Keep the actions coordinated and identify their subject before choosing each past form.
% mdparagraph:P882:end

% output: EX-XII-XIII-42-09-O001
% mdparagraph:P883:begin
Write one complete Old English translation, keeping the printed controls and the supplied expressions available.
% mdparagraph:P883:end

\end{minipage}\par\addvspace{4pt}

\reserve{7\baselineskip}
% mdparagraph:P884:begin
\section*{Cook 43. Mixed composition}\BookHeading{nav-ch17-learner-0219}{subsection}{Cook 43. Mixed composition}{1}
% mdparagraph:P884:end

% mdparagraph:P885:begin
Apply the same distinctions to comparisons, requests, reports and descriptions. The change of subject does not change the need to identify each clause’s participants and complements.
% mdparagraph:P885:end

\par\begin{minipage}{\linewidth}\raggedright
% mdparagraph:P886:begin
\subsection*{Cook 43.1}\BookHeading{nav-ch17-learner-0220}{subsubsection}{Cook 43.1}{0}
% mdparagraph:P886:end

% item: EX-XII-XIII-43-01
% mdparagraph:P887:begin
All [each] flesh is grass, and the glory (wuldor) of the flesh is like [as if] the blossom of herbs.
% mdparagraph:P887:end

\end{minipage}\par\addvspace{4pt}

% teaching: NEW-U17-GRAM-U17-H-43-01
% mdparagraph:P888:begin
\textbf{Construction.} The bracketed each and as if are part of the target. Preserve both images, then distinguish a copular identification from an elliptical comparison.
% mdparagraph:P888:end

% mdparagraph:P889:begin
\textbf{Supplied:} Ǣlc flǣsc is gærs, and ðæs flǣsces wuldor is swilce ðāra wyrta blōstma. — Every flesh is grass, and the glory of the flesh is as the blossom of plants. Retain the flesh/\allowbreak{}grass identification and the separate glory/\allowbreak{}blossom comparison. Swilce introduces an elliptical comparison here; no separate finite counterfactual verb is required.
% mdparagraph:P889:end

% teaching: U17-SUP-SWILCE-AS-IF
% mdparagraph:P890:begin
\textbf{Supplied:} swilce — like; as if.
% mdparagraph:P890:end

% output: EX-XII-XIII-43-01-O001
% mdparagraph:P891:begin
Write one complete Old English translation, keeping the printed controls and the supplied expressions available.
% mdparagraph:P891:end

\par\begin{minipage}{\linewidth}\raggedright
% mdparagraph:P892:begin
\subsection*{Cook 43.2}\BookHeading{nav-ch17-learner-0221}{subsubsection}{Cook 43.2}{0}
% mdparagraph:P892:end

% item: EX-XII-XIII-43-02
% mdparagraph:P893:begin
No man heard (183) his voice in (om. the) streets.
% mdparagraph:P893:end

% mdparagraph:P894:begin
\textbf{Construction.} Use negative concord (§183) and omit the marked article.
% mdparagraph:P894:end

% output: EX-XII-XIII-43-02-O001
% mdparagraph:P895:begin
Write one complete Old English translation, keeping the printed controls and the supplied expressions available.
% mdparagraph:P895:end

\end{minipage}\par\addvspace{4pt}

\par\begin{minipage}{\linewidth}\raggedright
% mdparagraph:P896:begin
\subsection*{Cook 43.3}\BookHeading{nav-ch17-learner-0222}{subsubsection}{Cook 43.3}{0}
% mdparagraph:P896:end

% item: EX-XII-XIII-43-03
% mdparagraph:P897:begin
Cry (cleopian), and cease thou not, lift up (āhębban) thy voice like (swā-swā) (om. a) trumpet.
% mdparagraph:P897:end

\end{minipage}\par\addvspace{4pt}

% teaching: NEW-U17-GRAM-U17-H-43-03
% mdparagraph:P898:begin
\textbf{Construction.} Use the command to the stated addressee. The supplied swā-swā phrase is an elliptical comparison; it does not make an actual trumpet the object being lifted.
% mdparagraph:P898:end

% mdparagraph:P899:begin
\textbf{Supplied:} swā-swā bieman. — As a trumpet: the accusative comparison imagines lifting a trumpet. It supplies a standard for the voice being raised, rather than an actual second object of the command. Retain the exact given swā-swā.
% mdparagraph:P899:end

% mdparagraph:P900:begin
\textbf{Supplied:} swā-swā — just as; like.
% mdparagraph:P900:end

% output: EX-XII-XIII-43-03-O001
% mdparagraph:P901:begin
Write one complete Old English translation, keeping the printed controls and the supplied expressions available.
% mdparagraph:P901:end

\par\begin{minipage}{\linewidth}\raggedright
% mdparagraph:P902:begin
\subsection*{Cook 43.4}\BookHeading{nav-ch17-learner-0223}{subsubsection}{Cook 43.4}{0}
% mdparagraph:P902:end

% item: EX-XII-XIII-43-04
% mdparagraph:P903:begin
Lady, whither goest thou?
% mdparagraph:P903:end

% output: EX-XII-XIII-43-04-O001
% mdparagraph:P904:begin
Write one complete Old English translation, keeping the printed controls and the supplied expressions available.
% mdparagraph:P904:end

\end{minipage}\par\addvspace{4pt}

\par\begin{minipage}{\linewidth}\raggedright
% mdparagraph:P905:begin
\subsection*{Cook 43.5}\BookHeading{nav-ch17-learner-0224}{subsubsection}{Cook 43.5}{0}
% mdparagraph:P905:end

% item: EX-XII-XIII-43-05
% mdparagraph:P906:begin
Thou doest ill (yfele).
% mdparagraph:P906:end

\end{minipage}\par\addvspace{4pt}

% teaching: U17-SUP-YFELE-ADVERB-INVARIANT
% mdparagraph:P907:begin
\textbf{Supplied:} yfele — badly; wrongly.
% mdparagraph:P907:end

% output: EX-XII-XIII-43-05-O001
% mdparagraph:P908:begin
Write one complete Old English translation, keeping the printed controls and the supplied expressions available.
% mdparagraph:P908:end

\par\begin{minipage}{\linewidth}\raggedright
% mdparagraph:P909:begin
\subsection*{Cook 43.6}\BookHeading{nav-ch17-learner-0225}{subsubsection}{Cook 43.6}{0}
% mdparagraph:P909:end

% item: EX-XII-XIII-43-06
% mdparagraph:P910:begin
He did ill, because he censured the princess, the king’s daughter.
% mdparagraph:P910:end

% mdparagraph:P911:begin
\textbf{Supplied:} yfele — badly; wrongly.
% mdparagraph:P911:end

% output: EX-XII-XIII-43-06-O001
% mdparagraph:P912:begin
Write one complete Old English translation, keeping the printed controls and the supplied expressions available.
% mdparagraph:P912:end

\end{minipage}\par\addvspace{4pt}

\par\begin{minipage}{\linewidth}\raggedright
% mdparagraph:P913:begin
\subsection*{Cook 43.7}\BookHeading{nav-ch17-learner-0226}{subsubsection}{Cook 43.7}{0}
% mdparagraph:P913:end

% item: EX-XII-XIII-43-07
% mdparagraph:P914:begin
They never learned arithmetic well.
% mdparagraph:P914:end

% output: EX-XII-XIII-43-07-O001
% mdparagraph:P915:begin
Write one complete Old English translation, keeping the printed controls and the supplied expressions available.
% mdparagraph:P915:end

\end{minipage}\par\addvspace{4pt}

\par\begin{minipage}{\linewidth}\raggedright
% mdparagraph:P916:begin
\subsection*{Cook 43.8}\BookHeading{nav-ch17-learner-0227}{subsubsection}{Cook 43.8}{0}
% mdparagraph:P916:end

% item: EX-XII-XIII-43-08
% mdparagraph:P917:begin
You (plur.) will soon know what you do not yet know.
% mdparagraph:P917:end

% mdparagraph:P918:begin
\textbf{Construction.} The addressees are plural. Express their not knowing with ordinary ne and the appropriate finite form of witan.
% mdparagraph:P918:end

% output: EX-XII-XIII-43-08-O001
% mdparagraph:P919:begin
Write one complete Old English translation, keeping the printed controls and the supplied expressions available.
% mdparagraph:P919:end

\end{minipage}\par\addvspace{4pt}

\par\begin{minipage}{\linewidth}\raggedright
% mdparagraph:P920:begin
\subsection*{Cook 43.9}\BookHeading{nav-ch17-learner-0228}{subsubsection}{Cook 43.9}{0}
% mdparagraph:P920:end

% item: EX-XII-XIII-43-09
% mdparagraph:P921:begin
Play (with) ball!
% mdparagraph:P921:end

% mdparagraph:P922:begin
\textbf{Construction.} The parenthetic with calls for mid with the dative instrument. Use the familiar ball paradigm.
% mdparagraph:P922:end

% output: EX-XII-XIII-43-09-O001
% mdparagraph:P923:begin
Write one complete Old English translation, keeping the printed controls and the supplied expressions available.
% mdparagraph:P923:end

\end{minipage}\par\addvspace{4pt}

\par\begin{minipage}{\linewidth}\raggedright
% mdparagraph:P924:begin
\subsection*{Cook 43.10}\BookHeading{nav-ch17-learner-0229}{subsubsection}{Cook 43.10}{0}
% mdparagraph:P924:end

% item: EX-XII-XIII-43-10
% mdparagraph:P925:begin
I eat nothing, although other men eat and are merry.
% mdparagraph:P925:end

\end{minipage}\par\addvspace{4pt}

% teaching: NEW-U17-GRAM-U17-H-43-10
% teaching: NEW-U17-VOC-THE-U17-CONCESSIVE-MARKER
% mdparagraph:P926:begin
\textbf{Construction.} The supplied although-clause has plural present subjunctives. Its meaning is “eat and are merry”, a concession rather than a wish.
% mdparagraph:P926:end

% mdparagraph:P927:begin
\textbf{Supplied:} ðēah ðe ōðre męnn eten and blīðe sīen. — Although other people eat and are merry. Eten and sīen are third plural present subjunctives in the concessive clause. Ðe belongs to the connective ðēah ðe; it is not a relative who or which. The complete clause remains supplied.
% mdparagraph:P927:end

% teaching: U17-SUP-BLITHE-NOM-M-PL
% mdparagraph:P928:begin
\textbf{Supplied:} blīðe — merry; glad; nominative, masculine, plural, strong.
% mdparagraph:P928:end

% output: EX-XII-XIII-43-10-O001
% mdparagraph:P929:begin
Write one complete Old English translation, keeping the printed controls and the supplied expressions available.
% mdparagraph:P929:end

\par\begin{minipage}{\linewidth}\raggedright
% mdparagraph:P930:begin
\subsection*{Cook 43.11}\BookHeading{nav-ch17-learner-0230}{subsubsection}{Cook 43.11}{0}
% mdparagraph:P930:end

% item: EX-XII-XIII-43-11
% mdparagraph:P931:begin
I desire to learn [that I may learn] from you the knowledge that you have.
% mdparagraph:P931:end

% mdparagraph:P932:begin
\textbf{Construction.} The bracket calls for a finite desired-content clause. Its first-person subject and the following plural you keep separate agreement; of with the dative names the source of learning.
% mdparagraph:P932:end

% output: EX-XII-XIII-43-11-O001
% mdparagraph:P933:begin
Write one complete Old English translation, keeping the printed controls and the supplied expressions available.
% mdparagraph:P933:end

\end{minipage}\par\addvspace{4pt}

\par\begin{minipage}{\linewidth}\raggedright
% mdparagraph:P934:begin
\subsection*{Cook 43.12}\BookHeading{nav-ch17-learner-0231}{subsubsection}{Cook 43.12}{0}
% mdparagraph:P934:end

% item: EX-XII-XIII-43-12
% mdparagraph:P935:begin
Choose which of them you (sing.) will.
% mdparagraph:P935:end

% mdparagraph:P936:begin
\textbf{Construction.} Address one person. The relative clause is indefinite: the one you wish to choose. Use the appropriate subjunctive.
% mdparagraph:P936:end

% output: EX-XII-XIII-43-12-O001
% mdparagraph:P937:begin
Write one complete Old English translation, keeping the printed controls and the supplied expressions available.
% mdparagraph:P937:end

\end{minipage}\par\addvspace{4pt}

\par\begin{minipage}{\linewidth}\raggedright
% mdparagraph:P938:begin
\subsection*{Cook 43.13}\BookHeading{nav-ch17-learner-0232}{subsubsection}{Cook 43.13}{0}
% mdparagraph:P938:end

% item: EX-XII-XIII-43-13
% mdparagraph:P939:begin
Within a few hours, two educated and high-born men came toward him.
% mdparagraph:P939:end

\end{minipage}\par\addvspace{4pt}

% teaching: NEW-U17-GRAM-U17-SUPPLIED-WITHIN-TIME
% mdparagraph:P940:begin
\textbf{Supplied:} Binnan fēawum tīdum — Within a few hours. The adjective and time noun are dative plural; keep the whole expression available.
% mdparagraph:P940:end

% teaching: U17-SUP-FEAWE-DAT-F-PL
% mdparagraph:P941:begin
\textbf{Supplied:} Binnan — within.
% mdparagraph:P941:end

% mdparagraph:P942:begin
\textbf{Supplied:} fēawum — few; dative, feminine, plural, strong.
% mdparagraph:P942:end

% teaching: U17-SUP-TOGEANES-INVARIANT
% mdparagraph:P943:begin
\textbf{Supplied:} him tōgēanes — Toward him; to meet him. Him is the dative complement, with tōgēanes placed after it. Keep the whole expression available.
% mdparagraph:P943:end

% output: EX-XII-XIII-43-13-O001
% mdparagraph:P944:begin
Write one complete Old English translation, keeping the printed controls and the supplied expressions available.
% mdparagraph:P944:end

\par\begin{minipage}{\linewidth}\raggedright
% mdparagraph:P945:begin
\subsection*{Cook 43.14}\BookHeading{nav-ch17-learner-0233}{subsubsection}{Cook 43.14}{0}
% mdparagraph:P945:end

% item: EX-XII-XIII-43-14
% mdparagraph:P946:begin
I woo the illustrious maiden.
% mdparagraph:P946:end

% output: EX-XII-XIII-43-14-O001
% mdparagraph:P947:begin
Write one complete Old English translation, keeping the printed controls and the supplied expressions available.
% mdparagraph:P947:end

\end{minipage}\par\addvspace{4pt}

\par\begin{minipage}{\linewidth}\raggedright
% mdparagraph:P948:begin
\subsection*{Cook 43.15}\BookHeading{nav-ch17-learner-0234}{subsubsection}{Cook 43.15}{0}
% mdparagraph:P948:end

% item: EX-XII-XIII-43-15
% mdparagraph:P949:begin
We all say that no one here is (194. b) so acceptable to God as he.
% mdparagraph:P949:end

% mdparagraph:P950:begin
\textbf{Construction.} The denied content takes the selected present subjunctive under Cook §194.b. Keep the local pleasing expression supplied.
% mdparagraph:P950:end

\end{minipage}\par\addvspace{4pt}

% teaching: U17-SUP-GECWEME-NOM-M-SG
% mdparagraph:P951:begin
\textbf{Supplied:} Gode swā gecwēme — So acceptable to God. Gode is dative, naming the person to whom acceptable; gecwēme is the strong masculine nominative singular predicate adjective. The degree word swā stands between them. Keep the whole expression available.
% mdparagraph:P951:end

% output: EX-XII-XIII-43-15-O001
% mdparagraph:P952:begin
Write one complete Old English translation, keeping the printed controls and the supplied expressions available.
% mdparagraph:P952:end

\par\begin{minipage}{\linewidth}\raggedright
% mdparagraph:P953:begin
\subsection*{Cook 43.16}\BookHeading{nav-ch17-learner-0235}{subsubsection}{Cook 43.16}{0}
% mdparagraph:P953:end

% item: EX-XII-XIII-43-16
% mdparagraph:P954:begin
At first he thought that he should be slain [that one should slay him].
% mdparagraph:P954:end

\end{minipage}\par\addvspace{4pt}

% teaching: NEW-U17-GRAM-U17-H-43-16
% mdparagraph:P955:begin
\textbf{Construction.} The bracket supplies an active paraphrase: an indefinite person might slay him. Preserve the distinction between that grammatical subject and the man who held the expectation.
% mdparagraph:P955:end

% mdparagraph:P956:begin
\textbf{Supplied:} ðæt hine man slōge. — That someone should slay him. Man is an indefinite singular subject, and hine is the man who had the expectation. Past subjunctive slōge expresses the contemplated act; it does not assert that he died.
% mdparagraph:P956:end

% output: EX-XII-XIII-43-16-O001
% mdparagraph:P957:begin
Write one complete Old English translation, keeping the printed controls and the supplied expressions available.
% mdparagraph:P957:end

\par\begin{minipage}{\linewidth}\raggedright
% mdparagraph:P958:begin
\subsection*{Cook 43.17}\BookHeading{nav-ch17-learner-0236}{subsubsection}{Cook 43.17}{0}
% mdparagraph:P958:end

% item: EX-XII-XIII-43-17
% mdparagraph:P959:begin
We all fall at (tō) your (sing.) feet.
% mdparagraph:P959:end

% mdparagraph:P960:begin
\textbf{Construction.} The possessor addressed by your is singular. We is the subject of fall and requires a plural indicative verb; the sentence is a statement, not a command.
% mdparagraph:P960:end

% output: EX-XII-XIII-43-17-O001
% mdparagraph:P961:begin
Write one complete Old English translation, keeping the printed controls and the supplied expressions available.
% mdparagraph:P961:end

\end{minipage}\par\addvspace{4pt}

\par\begin{minipage}{\linewidth}\raggedright
% mdparagraph:P962:begin
\subsection*{Cook 43.18}\BookHeading{nav-ch17-learner-0237}{subsubsection}{Cook 43.18}{0}
% mdparagraph:P962:end

% item: EX-XII-XIII-43-18
% mdparagraph:P963:begin
I (om. do) not believe that you are (194. b) my wife.
% mdparagraph:P963:end

% mdparagraph:P964:begin
\textbf{Construction.} Omit the English supporting do. The reported content takes second singular present subjunctive with thou under Cook §194.b.
% mdparagraph:P964:end

% output: EX-XII-XIII-43-18-O001
% mdparagraph:P965:begin
Write one complete Old English translation, keeping the printed controls and the supplied expressions available.
% mdparagraph:P965:end

\end{minipage}\par\addvspace{4pt}

\par\begin{minipage}{\linewidth}\raggedright
% mdparagraph:P966:begin
\subsection*{Cook 43.19}\BookHeading{nav-ch17-learner-0238}{subsubsection}{Cook 43.19}{0}
% mdparagraph:P966:end

% item: EX-XII-XIII-43-19
% mdparagraph:P967:begin
Where is my husband?
% mdparagraph:P967:end

% mdparagraph:P968:begin
\textbf{Supplied:} Hwǣr — where.
% mdparagraph:P968:end

% output: EX-XII-XIII-43-19-O001
% mdparagraph:P969:begin
Write one complete Old English translation, keeping the printed controls and the supplied expressions available.
% mdparagraph:P969:end

\end{minipage}\par\addvspace{4pt}

\par\begin{minipage}{\linewidth}\raggedright
% mdparagraph:P970:begin
\subsection*{Cook 43.20}\BookHeading{nav-ch17-learner-0239}{subsubsection}{Cook 43.20}{0}
% mdparagraph:P970:end

% item: EX-XII-XIII-43-20
% mdparagraph:P971:begin
I lived eleven years in the land of the Egyptians.
% mdparagraph:P971:end

\end{minipage}\par\addvspace{4pt}

% teaching: U17-SUP-ENDLEOFAN-INVARIANT
% mdparagraph:P972:begin
\textbf{Supplied:} endleofan — eleven.
% mdparagraph:P972:end

% teaching: U17-SUP-EGYPTA-GEN-PL
% mdparagraph:P973:begin
\textbf{Supplied:} Ægypta — the Egyptians; genitive, plural.
% mdparagraph:P973:end

% output: EX-XII-XIII-43-20-O001
% mdparagraph:P974:begin
Write one complete Old English translation, keeping the printed controls and the supplied expressions available.
% mdparagraph:P974:end

\reserve{7\baselineskip}
% mdparagraph:P975:begin
\section*{7. Interpreting the images}\BookHeading{nav-ch17-learner-0240}{subsection}{7. Interpreting the images}{1}
% mdparagraph:P975:end

% mdparagraph:P976:begin
Keep the wording of the completed translations beside the following questions. Explain the effect of an image after identifying its literal grammatical structure. Do not replace an entire translated clause with a general statement of its message.
% mdparagraph:P976:end

\reserve{8\baselineskip}
% mdparagraph:P977:begin
\subsection*{A05}\BookHeading{nav-ch17-learner-0241}{subsubsection}{A05}{0}
% mdparagraph:P977:end

% item: NEW-U17-A05
% mdparagraph:P978:begin
After translating Cook 43.1, analyse both images. Keep your complete translation beside the English prompt and the supplied vocabulary.
% mdparagraph:P978:end

% output: NEW-U17-A05-O001
% mdparagraph:P979:begin
\itempara{1}{Identify the two things joined by the first copular image and explain its contextual force.}
% mdparagraph:P979:end

% output: NEW-U17-A05-O002
% mdparagraph:P980:begin
\itempara{2}{Identify the second comparison and distinguish its structure from the first.}
% mdparagraph:P980:end

% output: NEW-U17-A05-O003
% mdparagraph:P981:begin
\itempara{3}{Explain why Human life is brief alone is not a complete translation of 43.1.}
% mdparagraph:P981:end

\reserve{8\baselineskip}
% mdparagraph:P982:begin
\subsection*{A06}\BookHeading{nav-ch17-learner-0242}{subsubsection}{A06}{0}
% mdparagraph:P982:end

% item: NEW-U17-A06
% mdparagraph:P983:begin
Using Cook 37.14 and 40.9 and the printed lexical help, explain how context narrows the meanings of horn and seed. Retain the actual images in translation.
% mdparagraph:P983:end

% output: NEW-U17-A06-O001
% mdparagraph:P984:begin
\itempara{1}{For 37.14, distinguish the horn image from a report about a literal animal.}
% mdparagraph:P984:end

% output: NEW-U17-A06-O002
% mdparagraph:P985:begin
\itempara{2}{For 40.9, explain what seed denotes and identify whose seed it is.}
% mdparagraph:P985:end

% output: NEW-U17-A06-O003
% mdparagraph:P986:begin
\itempara{3}{State what the familiar inflection alone cannot establish in these two examples.}
% mdparagraph:P986:end

\reserve{7\baselineskip}
% mdparagraph:P987:begin
\section*{8. A musical performance}\BookHeading{nav-ch17-learner-0243}{subsection}{8. A musical performance}{1}
% mdparagraph:P987:end

% mdparagraph:P988:begin
After a shipwreck, Apollonius has reached the court of King Arcestrates. At the feast the king’s daughter asks about his misfortunes. The king then tells her to fetch her harp and ease the young man’s sorrow.
% mdparagraph:P988:end

% mdparagraph:P989:begin
The reading is from Cook’s First Book, Reader XII, printed pages 174–175, lines 174.7–175.17 (scan pages 190–191). It begins when the daughter goes out and ends with the audience’s approval. The line references count the physical prose lines, anchored by the original marginal numbers. The text preserves the source’s spelling; line-end word divisions are joined and the original note numbers are identified by page. All help remains available during the reading tasks.
% mdparagraph:P989:end

% source: FB-XII-APOLLONIUS-MUSIC-PERFORMANCE
% mdparagraph:P990:begin
\readingtext{Ðā ēode hēo ūt,\textsuperscript{P174N9} and hēt fęccean hiere hearpan. And sōna swā hēo hearpian ongan, hēo mid wynsumum sange gemęngde ðǣre hearpan swēg. Ðā ongunnon ealle ðā męnn hīe hęrian on hiere swēgcræfte; and Apollonius āna\textsuperscript{P174N10} swīgode. Ðā cwæð se cyning: “Apolloni, nū ðū dēst\textsuperscript{P174N11} yfele, for-ðǣm-ðe ealle męnn hęriað mīne dohtor on hiere swēgcræfte,\textsuperscript{P174N12} and ðū āna hīe, swīgende,\textsuperscript{P174N12a} tǣlst.\textsuperscript{P174N13}” Apollonius cwæð: “Ēalā, ðū gōda cyning, gif ðū mē gelīefst,\textsuperscript{P174N14} ic sęcge ðæt ic ongiete ðæt sōðlice ðīn dohtor gefēoll\textsuperscript{P174N15} on swēgcræft, ac hēo næfð hine nā wel geleornod; ac hāt mē\textsuperscript{P174N16} nū sęllan ðā hearpan, ðonne wāst\textsuperscript{P174N17} ðū nū ðæt ðū gīet nāst.\textsuperscript{P174N17}” Arcestrates se cyning cwæð: “Apolloni, ic oncnāwe sōðlice ðæt ðū eart\textsuperscript{P175N1} on eallum ðingum wel gelǣred.” Ðā hēt se cyning sęllan Apollonie ðā hearpan. Apollonius ðā ūtēode, and hine scrȳdde, and sętte ænne cynehelm upon his hēafod, and nōm ðā hearpan on his hand, and inēode, and swā stōd ðæt se cyning and ealle ðā ymbsittendan wēndon ðæt hē nǣre Apollonius, ac ðæt hē wǣre Apollines,\textsuperscript{P175N2} ðāra hǣðenra god. Ðā wearð stilnes and swīge\textsuperscript{P175N3} geworden innan ðǣre healle. And Apollonius his hearpenægl genōm, and hē ðā hearpestręngas mid cræfte āstyrian ongan, and ðǣre hearpan swēg mid wynsumum sange gemęngde.\textsuperscript{P175N4} And se cyning self, and ealle ðe ðǣr andwearde wǣron, micelre stefne cleopedon and hine hęredon. Æfter ðisum forlēt\textsuperscript{P175N5} Apollonius ðā hearpan, and\textsuperscript{P175N6} plegode, and fela fægerra ðinga\textsuperscript{P175N7} ðǣr forðtēah,\textsuperscript{P175N8} ðe ðǣm folce ungecnāwen wæs and ungewunelic. And him\textsuperscript{P175N9} eallum ðearle līcode ǣlc ðāra ðinga ðe hē forðtēah.}
% mdparagraph:P990:end

\par
\reserve{7\baselineskip}
% mdparagraph:P991:begin
\subsection*{Finding the cited lines}\BookHeading{nav-ch17-learner-0244}{subsubsection}{Finding the cited lines}{0}
% mdparagraph:P991:end

% mdparagraph:P992:begin
The passage is set continuously above. These short openings locate the physical source lines used by the questions; words divided at a line end are printed complete here.
% mdparagraph:P992:end

\begin{itemize}
% mdparagraph:P993:begin
\item 174.7 — Ðā ēode hēo ūt
% mdparagraph:P993:end
% mdparagraph:P994:begin
\item 174.8 — sōna swā hēo hearpian
% mdparagraph:P994:end
% mdparagraph:P995:begin
\item 174.9 — gemęngde ðǣre hearpan swēg.
% mdparagraph:P995:end
% mdparagraph:P996:begin
\item 174.10 — męnn hīe hęrian on
% mdparagraph:P996:end
% mdparagraph:P997:begin
\item 174.11 — āna swīgode. Ðā cwæð
% mdparagraph:P997:end
% mdparagraph:P998:begin
\item 174.12 — dēst yfele, for-ðǣm-ðe ealle
% mdparagraph:P998:end
% mdparagraph:P999:begin
\item 174.13 — on hiere swēgcræfte, and
% mdparagraph:P999:end
% mdparagraph:P1000:begin
\item 174.14 — Apollonius cwæð: “Ēalā, ðū
% mdparagraph:P1000:end
% mdparagraph:P1001:begin
\item 174.15 — gelīefst, ic sęcge ðæt
% mdparagraph:P1001:end
% mdparagraph:P1002:begin
\item 174.16 — gefēoll on swēgcræft, ac
% mdparagraph:P1002:end
% mdparagraph:P1003:begin
\item 174.17 — ac hāt mē nū
% mdparagraph:P1003:end
% mdparagraph:P1004:begin
\item 174.18 — ðū gīet nāst.” Arcestrates
% mdparagraph:P1004:end
% mdparagraph:P1005:begin
\item 175.1 — ic oncnāwe sōðlice ðæt
% mdparagraph:P1005:end
% mdparagraph:P1006:begin
\item 175.2 — gelǣred.”
% mdparagraph:P1006:end
% mdparagraph:P1007:begin
\item 175.3 — Ðā hēt se cyning
% mdparagraph:P1007:end
% mdparagraph:P1008:begin
\item 175.4 — Apollonius ðā ūtēode, and
% mdparagraph:P1008:end
% mdparagraph:P1009:begin
\item 175.5 — cynehelm upon his hēafod
% mdparagraph:P1009:end
% mdparagraph:P1010:begin
\item 175.6 — hand, and inēode, and
% mdparagraph:P1010:end
% mdparagraph:P1011:begin
\item 175.7 — ðā ymbsittendan wēndon ðæt
% mdparagraph:P1011:end
% mdparagraph:P1012:begin
\item 175.8 — hē wǣre Apollines, ðāra
% mdparagraph:P1012:end
% mdparagraph:P1013:begin
\item 175.9 — and swīge geworden innan
% mdparagraph:P1013:end
% mdparagraph:P1014:begin
\item 175.10 — his hearpenægl genōm, and
% mdparagraph:P1014:end
% mdparagraph:P1015:begin
\item 175.11 — cræfte āstyrian ongan, and
% mdparagraph:P1015:end
% mdparagraph:P1016:begin
\item 175.12 — wynsumum sange gemęngde. And
% mdparagraph:P1016:end
% mdparagraph:P1017:begin
\item 175.13 — ðe ðǣr andwearde wǣron
% mdparagraph:P1017:end
% mdparagraph:P1018:begin
\item 175.14 — hine hęredon. Æfter ðisum
% mdparagraph:P1018:end
% mdparagraph:P1019:begin
\item 175.15 — and plegode, and fela
% mdparagraph:P1019:end
% mdparagraph:P1020:begin
\item 175.16 — ðǣm folce ungecnāwen wæs
% mdparagraph:P1020:end
% mdparagraph:P1021:begin
\item 175.17 — eallum ðearle līcode ǣlc
% mdparagraph:P1021:end
\end{itemize}

\reserve{7\baselineskip}
% mdparagraph:P1022:begin
\subsection*{Modern help for the reading}\BookHeading{nav-ch17-learner-0245}{subsubsection}{Modern help for the reading}{0}
% mdparagraph:P1022:end

% mdparagraph:P1023:begin
The following explanations supply the unfamiliar constructions and contextual meanings. Consult them throughout the reading and its questions. The original notes are retained separately at the end; their Latin parallels require no Latin exercise.
% mdparagraph:P1023:end

\reserve{7\baselineskip}
% mdparagraph:P1024:begin
\subsection*{174.7: the harp is fetched}\BookHeading{nav-ch17-learner-0246}{subsubsection}{174.7: the harp is fetched}{0}
% mdparagraph:P1024:end

% teaching: NEW-U17-GRAM-U17-RSUP-01
% teaching: NEW-U17-VOC-FECCEAN-U17-READING
% mdparagraph:P1025:begin
The daughter goes out and has her harp fetched. Hēt governs fęccean, fetch, with the performer of the fetching left unexpressed; hiere gives the harp’s possessor. Do not make the daughter herself the necessary performer of fetch.
% mdparagraph:P1025:end

\reserve{7\baselineskip}
% mdparagraph:P1026:begin
\subsection*{174.8–10: playing, singing and praise}\BookHeading{nav-ch17-learner-0247}{subsubsection}{174.8–10: playing, singing and praise}{0}
% mdparagraph:P1026:end

% teaching: NEW-U17-GRAM-U17-RSUP-02
% teaching: U17-SUP-SONA-TEMPORAL-CORRELATIVE
% teaching: U17-SUP-SWA-TEMPORAL-CORRELATIVE
% teaching: NEW-U17-VOC-HEARPIAN-U17-READING
% teaching: NEW-U17-VOC-ONGINNAN-U17-READING
% teaching: NEW-U17-VOC-GEMENGAN-U17-READING
% teaching: U17-SUP-SWEG-ACC-SG
% teaching: U17-SUP-SWEGCRAEFT-DAT-SG
% teaching: NEW-U17-GRAM-U17-READ-SONA-SWA
% mdparagraph:P1027:begin
Sōna swā means as soon as and introduces the complete time clause. Hearpian is the infinitive after began. Gemęngde means mixed: the harp’s sound is accusative and the song follows mid with the dative. On hiere swēgcræfte means with respect to her musical skill. The men’s praise concerns that skill.
% mdparagraph:P1027:end

\reserve{7\baselineskip}
% mdparagraph:P1028:begin
\subsection*{174.11–13: silence and blame}\BookHeading{nav-ch17-learner-0248}{subsubsection}{174.11–13: silence and blame}{0}
% mdparagraph:P1028:end

% teaching: NEW-U17-GRAM-U17-RSUP-03
% teaching: U17-SUP-APOLLONIUS-READING-NOM
% teaching: U17-SUP-AN-ALONE
% teaching: NEW-U17-VOC-SWIGIAN-U17-READING
% teaching: U17-SUP-APOLLONIUS-READING-VOC
% teaching: U17-SUP-FOR-THAM-THE-SOURCE-INVARIANT
% mdparagraph:P1029:begin
Āna means alone. Swīgode means remained silent; the source participle swīgende is explained by Cook’s note 12a as swīgiende, being silent. It accompanies the addressed you in the accusation. Dēst yfele means you do wrong: dēst is second singular present indicative of dōn, and yfele is an adverb. Tǣlst means you blame or censure. Keep the source spelling of the participle.
% mdparagraph:P1029:end

\reserve{7\baselineskip}
% mdparagraph:P1030:begin
\subsection*{174.14–16: permission and competence}\BookHeading{nav-ch17-learner-0249}{subsubsection}{174.14–16: permission and competence}{0}
% mdparagraph:P1030:end

% teaching: NEW-U17-GRAM-U17-RSUP-04
% teaching: NEW-U17-VOC-GELIEFAN-PERMIT-U17-SOURCE-SENSE
% teaching: NEW-U17-VOC-GEFEALLAN-U17-READING
% teaching: U17-SUP-ON-CHANCED-DOMAIN
% teaching: U17-SUP-SWEGCRAEFT-ACC-SG
% teaching: NEW-U17-VOC-GELEORNIAN-U17-READING
% teaching: NEW-U17-GRAM-U17-READ-ON-MUSIC
% mdparagraph:P1031:begin
Gif ðū mē gelīefst means if you permit me to speak. Gelīefst is second singular present indicative; mē is the dative person granted permission. Speaking is understood from the following ic sęcge, I will say; no speaking-infinitive occurs in the printed clause. This is the permission verb, distinct from believe. Ic ongiete introduces what Apollonius perceives. Gefēoll on swēgcræft means has chanced upon musical skill: on has an accusative complement in this supplied idiom. Masculine hine resumes swēgcræft, and næfð … nā wel geleornod says she has not learned or mastered it well.
% mdparagraph:P1031:end

\reserve{7\baselineskip}
% mdparagraph:P1032:begin
\subsection*{174.17–18: giving and knowing}\BookHeading{nav-ch17-learner-0250}{subsubsection}{174.17–18: giving and knowing}{0}
% mdparagraph:P1032:end

% teaching: NEW-U17-GRAM-U17-RSUP-05
% teaching: U17-SUP-THONNE-THEN-CONSEQUENCE
% teaching: NEW-U17-VOC-NYTAN-U17-READING
% mdparagraph:P1033:begin
Hāt mē nū sęllan ðā hearpan asks for the harp to be given to the speaker. Mē is the dative recipient of giving, as original note 16 directs; it is not the person commanded to give. Ðonne means then in the consequence. Wāst is you know and nāst is the contracted you do not know; the latter is supplied here without a new conjugation. The present wāst refers to the knowledge that will follow. Ðæt ðū gīet nāst means what you do not yet know.
% mdparagraph:P1033:end

\reserve{7\baselineskip}
% mdparagraph:P1034:begin
\subsection*{175.1–2: the king’s acknowledgment}\BookHeading{nav-ch17-learner-0251}{subsubsection}{175.1–2: the king’s acknowledgment}{0}
% mdparagraph:P1034:end

% teaching: NEW-U17-GRAM-U17-RSUP-06
% teaching: NEW-U17-VOC-ONCNAWAN-U17-READING
% mdparagraph:P1035:begin
Oncnāwe means acknowledge or recognize. The king’s content clause asserts that the addressee is well educated, using indicative eart. Original note 1 points to the certainty exception; the presence of ðæt does not mechanically demand a subjunctive. On eallum ðingum means in all matters, and gelǣred is the familiar adjective learned or educated.
% mdparagraph:P1035:end

\reserve{7\baselineskip}
% mdparagraph:P1036:begin
\subsection*{175.3–6: preparing the demonstration}\BookHeading{nav-ch17-learner-0252}{subsubsection}{175.3–6: preparing the demonstration}{0}
% mdparagraph:P1036:end

% teaching: NEW-U17-GRAM-U17-RSUP-07
% teaching: U17-SUP-APOLLONIUS-READING-DAT
% teaching: NEW-U17-VOC-UTGAN-U17-READING
% teaching: NEW-U17-VOC-SCRYDAN-U17-READING
% teaching: U17-SUP-AN-SOURCE-ACC-M-SG
% teaching: U17-SUP-UPON-INVARIANT
% mdparagraph:P1037:begin
The king has the harp given to dative Apollonie. Apollonius then goes out (ūtēode), dresses himself (hine scrȳdde), puts a crown upon his head, takes the harp into his hand and enters (inēode). The supplied source form ænne is masculine accusative singular; retain its unmarked vowel. Upon his hēafod and on his hand describe the destinations of placement and taking.
% mdparagraph:P1037:end

\reserve{7\baselineskip}
% mdparagraph:P1038:begin
\subsection*{175.6–8: appearance and supposition}\BookHeading{nav-ch17-learner-0253}{subsubsection}{175.6–8: appearance and supposition}{0}
% mdparagraph:P1038:end

% teaching: NEW-U17-GRAM-U17-RSUP-08
% teaching: NEW-U17-VOC-YMBSITTAN-U17-READING
% teaching: NEW-U17-VOC-WESAN-U17-READING
% teaching: U17-SUP-APOLLINES-READING-NOM
% teaching: U17-SUP-HAETHEN-GEN-PL
% teaching: U17-SUP-GOD-COMMON-NOM-SG
% mdparagraph:P1039:begin
Swā … ðæt introduces a result: he stood in such a way that the king and those sitting around supposed something about him. Ðā ymbsittendan is a weak plural description naming those people. The coordinated subject takes plural wēndon. The two content clauses have past subjunctives nǣre and wǣre; they present the spectators’ supposition. Apollines is the source’s name-form for Apollo, the god of the heathens. Preserve it without confusing him with Apollonius.
% mdparagraph:P1039:end

\reserve{7\baselineskip}
% mdparagraph:P1040:begin
\subsection*{175.8–9: silence}\BookHeading{nav-ch17-learner-0254}{subsubsection}{175.8–9: silence}{0}
% mdparagraph:P1040:end

% teaching: NEW-U17-GRAM-U17-RSUP-09
% teaching: NEW-U17-VOC-WEORTHAN-U17-READING
% teaching: U17-SUP-STILNES-NOM-SG
% teaching: U17-SUP-SWIGE-NOM-SG
% teaching: NEW-U17-VOC-GEWEORTHAN-U17-READING
% teaching: U17-SUP-INNAN-PREPOSITION
% teaching: U17-SUP-HEALL-DAT-SG
% mdparagraph:P1041:begin
Wearð … geworden presents quietness and silence coming about as a whole result; it is not a transitive passive. Retain singular wearð in this supplied expression. Stilnes and swīge are feminine nouns, and innan ðǣre healle means within the hall, with a dative complement.
% mdparagraph:P1041:end

\reserve{7\baselineskip}
% mdparagraph:P1042:begin
\subsection*{175.10–12: the performance}\BookHeading{nav-ch17-learner-0255}{subsubsection}{175.10–12: the performance}{0}
% mdparagraph:P1042:end

% teaching: NEW-U17-GRAM-U17-RSUP-10
% teaching: U17-SUP-HEARPENAEGL-ACC-SG
% teaching: NEW-U17-VOC-GENIMAN-U17-READING
% teaching: U17-SUP-HEARPESTRENG-ACC-PL
% teaching: U17-SUP-CRAEFT-DAT-SG-SKILL
% teaching: NEW-U17-VOC-CRAEFT-U17-SKILL
% teaching: NEW-U17-VOC-ASTYRIAN-U17-READING
% mdparagraph:P1043:begin
Genōm means took; the object is his hearpenægl, his plectrum or string-plucking implement. Āstyrian ongan means began to set in motion, with the harp strings as accusative object. Mid cræfte means skilfully, with skill. Gemęngde again combines harp sound with song. These supplied senses explain the performance without creating new finite systems for all its verbs.
% mdparagraph:P1043:end

\reserve{7\baselineskip}
% mdparagraph:P1044:begin
\subsection*{175.12–14: the response}\BookHeading{nav-ch17-learner-0256}{subsubsection}{175.12–14: the response}{0}
% mdparagraph:P1044:end

% teaching: NEW-U17-GRAM-U17-RSUP-11
% teaching: U17-SUP-ANDWEARD-NOM-M-PL
% teaching: U17-SUP-MICEL-DAT-F-SG-LOUD
% mdparagraph:P1045:begin
The king himself and all those present are the plural actors of cleopedon, cried out, and hęredon, praised. The invariant relative ðe represents the people who were present. Micelre stefne means with a loud voice; it is a dative expression of manner without a printed mid. Hine is the person praised.
% mdparagraph:P1045:end

\reserve{7\baselineskip}
% mdparagraph:P1046:begin
\subsection*{175.14–17: further accomplishments}\BookHeading{nav-ch17-learner-0257}{subsubsection}{175.14–17: further accomplishments}{0}
% mdparagraph:P1046:end

% teaching: NEW-U17-GRAM-U17-RSUP-12
% teaching: U17-SUP-THES-DAT-N-SG-SOURCE
% teaching: NEW-U17-VOC-FORLAETAN-U17-READING
% teaching: NEW-U17-VOC-PLEGIAN-U17-READING
% teaching: NEW-U17-VOC-FORTHTEON-U17-READING
% teaching: NEW-U17-VOC-UNGECNAWEN-U17-READING
% teaching: U17-SUP-UNGEWUNELIC-NOM-N-SG
% teaching: U17-SUP-THEARLE-INVARIANT
% teaching: NEW-U17-VOC-LICIAN-U17-READING
% mdparagraph:P1047:begin
Forlēt means set aside the harp; plegode means acted or performed; forðtēah means presented or performed the accomplishments. These performance meanings remain supplied. Fela governs genitive plural fægerra ðinga. The following relative construes what was unfamiliar collectively and has singular ungecnāwen wæs and ungewunelic; preserve this source construction without making it a general rule for agreement. In the last sentence, ǣlc is singular subject of līcode, pleased, and him eallum gives the dative people pleased. Ðāra ðinga is the partitive genitive after each; ðearle means greatly.
% mdparagraph:P1047:end

\reserve{7\baselineskip}
% mdparagraph:P1048:begin
\subsection*{Exact word help}\BookHeading{nav-ch17-learner-0258}{subsubsection}{Exact word help}{0}
% mdparagraph:P1048:end

% mdparagraph:P1049:begin
These exact source forms can be consulted. The noun and verb endings below identify their particular occurrence; they do not require an additional complete paradigm from memory.
% mdparagraph:P1049:end

% mdparagraph:P1050:begin
\textbf{fęccean} — fetch; plain infinitive after hēt.
% mdparagraph:P1050:end

% mdparagraph:P1051:begin
\textbf{sōna} — as soon as, with swā.
% mdparagraph:P1051:end

% mdparagraph:P1052:begin
\textbf{swā} — as, in sōna swā.
% mdparagraph:P1052:end

% mdparagraph:P1053:begin
\textbf{hearpian} — play the harp; plain infinitive after began.
% mdparagraph:P1053:end

% mdparagraph:P1054:begin
\textbf{ongan} — begin, with the explicitly supplied playing/\allowbreak{}praising action; indicative, preterite, third person, singular.
% mdparagraph:P1054:end

% mdparagraph:P1055:begin
\textbf{gemęngde} — mix sound with song; indicative, preterite, third person, singular.
% mdparagraph:P1055:end

% mdparagraph:P1056:begin
\textbf{swēg} — music; sound; accusative, masculine, singular.
% mdparagraph:P1056:end

% mdparagraph:P1057:begin
\textbf{ongunnon} — begin, with the explicitly supplied playing/\allowbreak{}praising action; indicative, preterite, third person, plural.
% mdparagraph:P1057:end

% mdparagraph:P1058:begin
\textbf{hęrian} — praise; plain infinitive after began.
% mdparagraph:P1058:end

% mdparagraph:P1059:begin
\textbf{swēgcræfte} — musical skill; dative, masculine, singular.
% mdparagraph:P1059:end

% mdparagraph:P1060:begin
\textbf{Apollonius} — Apollonius, the visitor; nominative, masculine, singular.
% mdparagraph:P1060:end

% mdparagraph:P1061:begin
\textbf{āna} — alone; nominative, masculine, singular, weak.
% mdparagraph:P1061:end

% mdparagraph:P1062:begin
\textbf{swīgode} — be silent; indicative, preterite, third person, singular.
% mdparagraph:P1062:end

% mdparagraph:P1063:begin
\textbf{Apolloni} — Apollonius, the visitor; vocative, masculine, singular.
% mdparagraph:P1063:end

% mdparagraph:P1064:begin
\textbf{yfele} — badly; wrongly.
% mdparagraph:P1064:end

% mdparagraph:P1065:begin
\textbf{for-ðǣm-ðe} — because.
% mdparagraph:P1065:end

% mdparagraph:P1066:begin
\textbf{swīgende} — being silent; the supplied present participle.
% mdparagraph:P1066:end

% mdparagraph:P1067:begin
\textbf{gelīefst} — permit; allow; indicative, present, second person, singular.
% mdparagraph:P1067:end

% mdparagraph:P1068:begin
\textbf{gefēoll} — chance upon; indicative, preterite, third person, singular.
% mdparagraph:P1068:end

% mdparagraph:P1069:begin
\textbf{on} — on/\allowbreak{}into, in the chanced-upon idiom.
% mdparagraph:P1069:end

% mdparagraph:P1070:begin
\textbf{swēgcræft} — musical skill; accusative, masculine, singular.
% mdparagraph:P1070:end

% mdparagraph:P1071:begin
\textbf{geleornod} — learned; mastered; past participle.
% mdparagraph:P1071:end

% mdparagraph:P1072:begin
\textbf{ðonne} — then; at that point.
% mdparagraph:P1072:end

% mdparagraph:P1073:begin
\textbf{nāst} — not know; indicative, present, second person, singular.
% mdparagraph:P1073:end

% teaching: U17-SUP-ARCESTRATES-READING-NOM
% mdparagraph:P1074:begin
\textbf{Arcestrates} — Arcestrates, the king; nominative, masculine, singular.
% mdparagraph:P1074:end

% mdparagraph:P1075:begin
\textbf{oncnāwe} — recognize; acknowledge; indicative, present, first person, singular.
% mdparagraph:P1075:end

% mdparagraph:P1076:begin
\textbf{Apollonie} — Apollonius, the visitor; dative, masculine, singular.
% mdparagraph:P1076:end

% mdparagraph:P1077:begin
\textbf{ūtēode} — go out; indicative, preterite, third person, singular.
% mdparagraph:P1077:end

% mdparagraph:P1078:begin
\textbf{scrȳdde} — dress oneself; indicative, preterite, third person, singular.
% mdparagraph:P1078:end

% mdparagraph:P1079:begin
\textbf{ænne} — a; one; accusative, masculine, singular, strong.
% mdparagraph:P1079:end

% mdparagraph:P1080:begin
\textbf{upon} — on; onto.
% mdparagraph:P1080:end

% mdparagraph:P1081:begin
\textbf{ymbsittendan} — people sitting around; present participle, nominative, masculine, plural, weak.
% mdparagraph:P1081:end

% mdparagraph:P1082:begin
\textbf{nǣre} — was not, in the reported supposition; subjunctive, preterite, third person, singular.
% mdparagraph:P1082:end

% mdparagraph:P1083:begin
\textbf{Apollines} — Apollo, the god; not Apollonius; nominative, masculine, singular.
% mdparagraph:P1083:end

% mdparagraph:P1084:begin
\textbf{hǣðenra} — heathen people; genitive, masculine, plural, weak.
% mdparagraph:P1084:end

% mdparagraph:P1085:begin
\textbf{god} — a god; deity; nominative, masculine, singular.
% mdparagraph:P1085:end

% mdparagraph:P1086:begin
\textbf{wearð} — came about, with geworden in the supplied result expression; indicative, preterite, third person, singular.
% mdparagraph:P1086:end

% mdparagraph:P1087:begin
\textbf{stilnes} — stillness; quietness; nominative, feminine, singular.
% mdparagraph:P1087:end

% mdparagraph:P1088:begin
\textbf{swīge} — silence; nominative, feminine, singular.
% mdparagraph:P1088:end

% mdparagraph:P1089:begin
\textbf{geworden} — come about; past participle in the supplied result expression.
% mdparagraph:P1089:end

% mdparagraph:P1090:begin
\textbf{innan} — within; inside.
% mdparagraph:P1090:end

% mdparagraph:P1091:begin
\textbf{healle} — hall; dative, feminine, singular.
% mdparagraph:P1091:end

% mdparagraph:P1092:begin
\textbf{hearpenægl} — plectrum; harp pick; accusative, masculine, singular.
% mdparagraph:P1092:end

% mdparagraph:P1093:begin
\textbf{genōm} — take; indicative, preterite, third person, singular.
% mdparagraph:P1093:end

% mdparagraph:P1094:begin
\textbf{hearpestręngas} — harp strings; accusative, masculine, plural.
% mdparagraph:P1094:end

% mdparagraph:P1095:begin
\textbf{mid} — with; in the manner of the accompanying quality or mental disposition.
% mdparagraph:P1095:end

% mdparagraph:P1096:begin
\textbf{cræfte} — skill; accomplishment; dative, masculine, singular.
% mdparagraph:P1096:end

% mdparagraph:P1097:begin
\textbf{āstyrian} — move; set in motion; plain infinitive after began.
% mdparagraph:P1097:end

% mdparagraph:P1098:begin
\textbf{ongan} — begin, with the explicitly supplied playing/\allowbreak{}praising action; indicative, preterite, third person, singular.
% mdparagraph:P1098:end

% mdparagraph:P1099:begin
\textbf{ðe} — that; which; who.
% mdparagraph:P1099:end

% mdparagraph:P1100:begin
\textbf{andwearde} — present; nominative, masculine, plural, strong.
% mdparagraph:P1100:end

% mdparagraph:P1101:begin
\textbf{micelre} — great; contextually loud; dative, feminine, singular, strong.
% mdparagraph:P1101:end

% mdparagraph:P1102:begin
\textbf{cleopedon} — cry aloud; indicative, preterite, third person, plural.
% mdparagraph:P1102:end

% mdparagraph:P1103:begin
\textbf{ðisum} — this; this event; dative, neuter, singular.
% mdparagraph:P1103:end

% mdparagraph:P1104:begin
\textbf{forlēt} — set aside; indicative, preterite, third person, singular.
% mdparagraph:P1104:end

% mdparagraph:P1105:begin
\textbf{plegode} — act; perform; indicative, preterite, third person, singular.
% mdparagraph:P1105:end

% mdparagraph:P1106:begin
\textbf{forðtēah} — present; perform; indicative, preterite, third person, singular.
% mdparagraph:P1106:end

% mdparagraph:P1107:begin
\textbf{ungecnāwen} — unknown; singular predicate in the supplied collective expression.
% mdparagraph:P1107:end

% mdparagraph:P1108:begin
\textbf{ungewunelic} — unusual; unfamiliar; nominative, neuter, singular, strong.
% mdparagraph:P1108:end

% mdparagraph:P1109:begin
\textbf{ðearle} — greatly; exceedingly.
% mdparagraph:P1109:end

% mdparagraph:P1110:begin
\textbf{līcode} — please; indicative, preterite, third person, singular.
% mdparagraph:P1110:end

% mdparagraph:P1111:begin
\textbf{forðtēah} — present; perform; indicative, preterite, third person, singular.
% mdparagraph:P1111:end

% mdparagraph:P1112:begin
\textbf{Names:} Apollonius is the visitor; Apolloni addresses him, and Apollonie is the dative form for the recipient. Arcestrates is the king. Apollines is Apollo, the god. These exact name-forms are supplied.
% mdparagraph:P1112:end

% mdparagraph:P1113:begin
\textbf{Source distinctions:} for-ðǣm-ðe means because, with the source’s internal hyphens and ǣ. Āna means alone, rather than cardinal one. The lowercase common noun god names a deity here. Hǣðenra names the heathens in the genitive plural after ðāra. Mid cræfte describes manner, whereas mid wynsumum sange supplies the accompanying song. Æfter ðisum means after this, with the source’s single s.
% mdparagraph:P1113:end

\reserve{7\baselineskip}
% mdparagraph:P1114:begin
\subsection*{Reading tasks}\BookHeading{nav-ch17-learner-0259}{subsubsection}{Reading tasks}{0}
% mdparagraph:P1114:end

% mdparagraph:P1115:begin
Keep the exact passage, the original notes and all modern help open. These questions ask you to relate the supplied forms to one another and to the episode’s argument; they do not test independent recall of source-only forms.
% mdparagraph:P1115:end

\reserve{8\baselineskip}
% mdparagraph:P1116:begin
\subsection*{R01}\BookHeading{nav-ch17-learner-0260}{subsubsection}{R01}{0}
% mdparagraph:P1116:end

% item: NEW-U17-R01
% mdparagraph:P1117:begin
Translate the complete assigned reading in the three stated divisions, keeping all local help available.
% mdparagraph:P1117:end

% output: NEW-U17-R01-O001
% mdparagraph:P1118:begin
\itempara{1}{Translate 174.7–174.13, through tǣlst.}
% mdparagraph:P1118:end

% output: NEW-U17-R01-O002
% mdparagraph:P1119:begin
\itempara{2}{Translate 174.14–175.2, through wel gelǣred.}
% mdparagraph:P1119:end

% output: NEW-U17-R01-O003
% mdparagraph:P1120:begin
\itempara{3}{Translate 175.3–175.17, through the final forðtēah.}
% mdparagraph:P1120:end

\reserve{8\baselineskip}
% mdparagraph:P1121:begin
\subsection*{R02}\BookHeading{nav-ch17-learner-0261}{subsubsection}{R02}{0}
% mdparagraph:P1121:end

% item: NEW-U17-R02
% mdparagraph:P1122:begin
Identify each referent and the grammatical evidence. Use the supplied source notes; say when context is needed.
% mdparagraph:P1122:end

% output: NEW-U17-R02-O001
% mdparagraph:P1123:begin
\itempara{1}{Give the referent of hēo in 174.7.}
% mdparagraph:P1123:end

% output: NEW-U17-R02-O002
% mdparagraph:P1124:begin
\itempara{2}{Give the referent of hine in 174.16 and its role.}
% mdparagraph:P1124:end

% output: NEW-U17-R02-O003
% mdparagraph:P1125:begin
\itempara{3}{Give the referent of hē in the reported clauses 175.7–8.}
% mdparagraph:P1125:end

% output: NEW-U17-R02-O004
% mdparagraph:P1126:begin
\itempara{4}{Give the referent and case of him eallum in 175.16–17.}
% mdparagraph:P1126:end

\reserve{8\baselineskip}
% mdparagraph:P1127:begin
\subsection*{R03}\BookHeading{nav-ch17-learner-0262}{subsubsection}{R03}{0}
% mdparagraph:P1127:end

% item: NEW-U17-R03
% mdparagraph:P1128:begin
Identify the clause relationship in each cited expression and distinguish it from a similar construction elsewhere.
% mdparagraph:P1128:end

% output: NEW-U17-R03-O001
% mdparagraph:P1129:begin
\itempara{1}{Classify sōna swā 174.8 and explain its time relation.}
% mdparagraph:P1129:end

% output: NEW-U17-R03-O002
% mdparagraph:P1130:begin
\itempara{2}{Classify gif 174.14–15 in its context.}
% mdparagraph:P1130:end

% output: NEW-U17-R03-O003
% mdparagraph:P1131:begin
\itempara{3}{Classify swā … ðæt 175.6–7.}
% mdparagraph:P1131:end

% output: NEW-U17-R03-O004
% mdparagraph:P1132:begin
\itempara{4}{Explain why the two ðæt clauses 175.7–8 do not assert that Apollonius actually is Apollo.}
% mdparagraph:P1132:end

\reserve{8\baselineskip}
% mdparagraph:P1133:begin
\subsection*{R04}\BookHeading{nav-ch17-learner-0263}{subsubsection}{R04}{0}
% mdparagraph:P1133:end

% item: NEW-U17-R04
% mdparagraph:P1134:begin
Analyse the supplied participial and predicate expressions without changing the source spelling.
% mdparagraph:P1134:end

% output: NEW-U17-R04-O001
% mdparagraph:P1135:begin
\itempara{1}{Explain swīgende 174.13 using original note 12a.}
% mdparagraph:P1135:end

% output: NEW-U17-R04-O002
% mdparagraph:P1136:begin
\itempara{2}{Explain the referent of ðā ymbsittendan 175.7.}
% mdparagraph:P1136:end

% output: NEW-U17-R04-O003
% mdparagraph:P1137:begin
\itempara{3}{Account for singular ungecnāwen wæs and ungewunelic 175.16.}
% mdparagraph:P1137:end

% output: NEW-U17-R04-O004
% mdparagraph:P1138:begin
\itempara{4}{Identify the subject of līcode 175.17 and distinguish it from him eallum.}
% mdparagraph:P1138:end

\reserve{8\baselineskip}
% mdparagraph:P1139:begin
\subsection*{R05}\BookHeading{nav-ch17-learner-0264}{subsubsection}{R05}{0}
% mdparagraph:P1139:end

% item: NEW-U17-R05
% mdparagraph:P1140:begin
Separate literal wording from the contextual performance sense in each expression, using the local help and original notes.
% mdparagraph:P1140:end

% output: NEW-U17-R05-O001
% mdparagraph:P1141:begin
\itempara{1}{Explain gefēoll on swēgcræft 174.16.}
% mdparagraph:P1141:end

% output: NEW-U17-R05-O002
% mdparagraph:P1142:begin
\itempara{2}{Explain plegode and forðtēah 175.15.}
% mdparagraph:P1142:end

% output: NEW-U17-R05-O003
% mdparagraph:P1143:begin
\itempara{3}{Explain the function of Apollines 175.8.}
% mdparagraph:P1143:end

\reserve{8\baselineskip}
% mdparagraph:P1144:begin
\subsection*{R06}\BookHeading{nav-ch17-learner-0265}{subsubsection}{R06}{0}
% mdparagraph:P1144:end

% item: NEW-U17-R06
% mdparagraph:P1145:begin
Explain the argument of the episode from the actual Old English.
% mdparagraph:P1145:end

% output: NEW-U17-R06-O001
% mdparagraph:P1146:begin
\itempara{1}{State the difference between the king’s reading of silence and Apollonius’s explanation.}
% mdparagraph:P1146:end

% output: NEW-U17-R06-O002
% mdparagraph:P1147:begin
\itempara{2}{Identify two kinds of evidence by which Apollonius answers the criticism of his conduct.}
% mdparagraph:P1147:end

% output: NEW-U17-R06-O003
% mdparagraph:P1148:begin
\itempara{3}{Explain how the final sentence combines novelty with approval.}
% mdparagraph:P1148:end

\reserve{8\baselineskip}
% mdparagraph:P1149:begin
\section*{Source and editorial notes}\BookHeading{nav-ch17-learner-0266}{subsection}{Source and editorial notes}{1}
% mdparagraph:P1149:end

\begingroup\fontsize{10}{12.8}\selectfont
\setlength{\parskip}{1.5pt}
% mdparagraph:P1150:begin
The English exercises are Cook, Exercises in Old English, joint Group XII, XIII, Exercises 37–43, printed pages 22–28 (scan pages 30–36). All original prompts and controls are retained. The Old English reading is Cook, A First Book in Old English, second edition, Reader XII, printed pages 174–175. Original reading note numbers restart on page 175; P174 and P175 below identify their page. The modern explanations and the lettered exercises are separate teaching material.
% mdparagraph:P1150:end

% mdparagraph:P1151:begin
The supplied exclusion phrase for the 144,000 follows the parallel recorded by Bosworth–Toller under þūsend, printed page 1082. Acute length marks are presented as macrons and a divided word is joined; final-o ðūsendo is retained. The supplied nīwne follows the long nīwe adjective’s wilde pattern described by Wright, pages 49, 213–214. In the although-expression, the dictionary’s ðēah-ðe is spaced as ðēah ðe. These limited presentation choices do not change the original English prompts or the reading text.
% mdparagraph:P1151:end

% mdparagraph:P1152:begin
\textsuperscript{P174N9}: Not in the Latin.
% mdparagraph:P1152:end

% mdparagraph:P1153:begin
\textsuperscript{P174N10}: See 79.
% mdparagraph:P1153:end

% mdparagraph:P1154:begin
\textsuperscript{P174N11}: See 140.
% mdparagraph:P1154:end

% mdparagraph:P1155:begin
\textsuperscript{P174N12}: Lat. arte musica.
% mdparagraph:P1155:end

% mdparagraph:P1156:begin
\textsuperscript{P174N12a}: For swīgiende.
% mdparagraph:P1156:end

% mdparagraph:P1157:begin
\textsuperscript{P174N13}: Lat. vituperas.
% mdparagraph:P1157:end

% mdparagraph:P1158:begin
\textsuperscript{P174N14}: See 196. d.
% mdparagraph:P1158:end

% mdparagraph:P1159:begin
\textsuperscript{P174N15}: Lat. incidit. Translate, has chanced.
% mdparagraph:P1159:end

% mdparagraph:P1160:begin
\textsuperscript{P174N16}: See 164. a.
% mdparagraph:P1160:end

% mdparagraph:P1161:begin
\textsuperscript{P174N17}: See 126.
% mdparagraph:P1161:end

% mdparagraph:P1162:begin
\textsuperscript{P175N1}: See 194, note.
% mdparagraph:P1162:end

% mdparagraph:P1163:begin
\textsuperscript{P175N2}: Apollo.
% mdparagraph:P1163:end

% mdparagraph:P1164:begin
\textsuperscript{P175N3}: We are reminded of Æn. II. 1, Conticuere omnes.
% mdparagraph:P1164:end

% mdparagraph:P1165:begin
\textsuperscript{P175N4}: To this sentence there corresponds in the Latin: — arripuit plectrum, animumque accommodat arti; cum chordis miscetur vox cantu modulata.
% mdparagraph:P1165:end

% mdparagraph:P1166:begin
\textsuperscript{P175N5}: Lat. deponens.
% mdparagraph:P1166:end

% mdparagraph:P1167:begin
\textsuperscript{P175N6}: The rest of this sentence paraphrases: induit statum comicum et inauditas actiones expressit, deinde tragicum.
% mdparagraph:P1167:end

% mdparagraph:P1168:begin
\textsuperscript{P175N7}: See 154. a, b.
% mdparagraph:P1168:end

% mdparagraph:P1169:begin
\textsuperscript{P175N8}: Lat. expressit.
% mdparagraph:P1169:end

% mdparagraph:P1170:begin
\textsuperscript{P175N9}: See 164. k.
% mdparagraph:P1170:end

% mdparagraph:P1171:begin
The short Latin glosses identify musical art (174 note 12), censure (174 note 13), chancing upon something (174 note 15), setting aside (175 note 5), and representing or performing (175 note 8). The allusion in 175 note 3 recalls everyone falling silent. The verse in note 4 describes taking the plectrum and mingling song with the strings; note 6 names comic and tragic performance. They illuminate the Old English scene without replacing its actual wording.
% mdparagraph:P1171:end

\endgroup

% book-source-body:end
\par\label{end-ch17-learner}
\clearpage\endgroup

\BookRecto
\begingroup
\renewcommand{\runningtitle}{Unit 18. Prose consolidation}
\hyphenpenalty=10000
\exhyphenpenalty=10000
\emergencystretch=3em
\newcommand{\itempara}[2]{\par\addvspace{3pt}\noindent\hangindent=1.8em\hangafter=1\makebox[1.8em][l]{#1.}#2\par}
\newcommand{\reserve}[1]{\par\penalty-100\begingroup\dimen0=\pagegoal\advance\dimen0 by-\pagetotal\ifdim\dimen0<#1\newpage\fi\endgroup}
% book-source-body:begin

\addfontfeatures{Ligatures=TeXOff}
\raggedright
\booktitle{\BookComponent{ch18-learner}{Unit 18. Independent prose consolidation}UNIT 18}{Independent prose consolidation}

% mdparagraph:P001:begin
A connected narrative changes speakers, refers back across clauses, and presents some actions as wishes, instructions or reports. Identify each clause’s participants and time relation before joining it to the whole.
% mdparagraph:P001:end

% mdparagraph:P002:begin
The final numbered Cook exercises form a sustained Joseph composition sequence. The original Old English reading continues the banquet scene with the princess’s gifts and the farewells. These are separate texts. The two translated examples and lettered tasks are new teaching material.
% mdparagraph:P002:end

% teaching: NEW-U18-DICTIONARY-REFERENCES
% teaching: NEW-U18-VERBAL-CONTEXT-HELP
% mdparagraph:P003:begin
\textbf{Forms and help.} A dictionary form identifies a word; it may still need inflection. The complete tables teach their named systems and stated meanings. Diagnostic stems are references to those systems, not assignments of unlisted moods. An expression labelled \textbf{Supplied} remains visible during the task: use its unfamiliar forms with the explanation given. Its independent recall is not required. In verb tables, 1, 2 and 3 denote person.
% mdparagraph:P003:end

\reserve{7\baselineskip}
% mdparagraph:P004:begin
\section*{1. Clauses, reference and reporting}\BookHeading{nav-ch18-learner-0001}{subsection}{1. Clauses, reference and reporting}{1}
% mdparagraph:P004:end

% mdparagraph:P005:begin
A noun phrase takes its case from its own clause. A relative expression can have a different role from its antecedent. The particle ðe has no visible case ending; identify its verb and complements to establish its role.
% mdparagraph:P005:end

% teaching: NEW-U18-EXAMPLE-01-TEACHING
\reserve{5\baselineskip}
% mdparagraph:P006:begin
\textbf{Se mann ðe se cyning geseah, sang.}
% mdparagraph:P006:end

% mdparagraph:P007:begin
The man whom the king saw sang.
% mdparagraph:P007:end

% mdparagraph:P008:begin
The man is subject of sang, while the relative represents the object of geseah; the king is the relative clause subject.
% mdparagraph:P008:end

% teaching: NEW-U18-EXAMPLE-02-TEACHING
\reserve{5\baselineskip}
% mdparagraph:P009:begin
\textbf{Se cyning hēt ðone mann ðǣm cnihte ðā bōc sęllan.}
% mdparagraph:P009:end

% mdparagraph:P010:begin
The king ordered the man to give the servant the book.
% mdparagraph:P010:end

% mdparagraph:P011:begin
The accusative man performs the infinitive; the servant is its dative recipient and the book its accusative object.
% mdparagraph:P011:end

% mdparagraph:P012:begin
A third-person pronoun follows its referent’s grammatical gender, including that of an inanimate noun. Agreement can exclude a referent without selecting a unique person. There is no general nearest-noun rule.
% mdparagraph:P012:end

% mdparagraph:P013:begin
In direct speech, ic identifies the current speaker. A report revises that reference to fit its reporting clause. A preterite subjunctive under a past report does not itself prove the report false. Establish earlier, simultaneous or prospective time from the whole statement.
% mdparagraph:P013:end

% mdparagraph:P014:begin
With hātan and an infinitive, the accusative actor is the person ordered to act. A dative recipient of that action is a separate complement.
% mdparagraph:P014:end

\reserve{7\baselineskip}
% mdparagraph:P015:begin
\section*{2. Nouns and adjectives}\BookHeading{nav-ch18-learner-0002}{subsection}{2. Nouns and adjectives}{1}
% mdparagraph:P015:end

% mdparagraph:P016:begin
Choose case from the phrase’s function: nominative subject, accusative direct object, genitive possessor or licensed complement, dative recipient or licensed complement. A preposition or particular verb can determine the case.
% mdparagraph:P016:end

% mdparagraph:P017:begin
Strong masculine nouns normally have genitive singular -es, dative -e, nominative/\allowbreak{}accusative plural -as, genitive plural -a and dative plural -um. Strong neuters share -es, -e, -a and -um; nominative and accusative agree, with plural -u or no ending according to the stem. The selected feminine oblique singular ends in -e. Stem changes and weak assa are given individually.
% mdparagraph:P017:end

% mdparagraph:P018:begin
A singular-only table gives all four cases of the stated mass or abstract meaning, using a familiar class. It does not claim that no plural is possible in other senses.
% mdparagraph:P018:end

\reserve{7\baselineskip}
% mdparagraph:P019:begin
\subsection*{sacc}\BookHeading{nav-ch18-learner-0003}{subsubsection}{sacc}{0}
% mdparagraph:P019:end

% teaching: NEW-U18-VOC-SACC-U18-SENSE
\par\begin{minipage}{\linewidth}\raggedright
% mdparagraph:P020:begin
\textbf{sacc} — sack.
% mdparagraph:P020:end

% teaching: NEW-U18-NOM-SACC
% mdparagraph:P021:begin
Masculine; stem sacc-. Use the ordinary strong masculine endings.
% mdparagraph:P021:end

% table: NEW-U18-NOM-SACC-DISPLAY
% mdtable:T01:begin
\par\begingroup
\fontsize{10.5}{13.2}\selectfont
\setlength{\LTleft}{0pt}\setlength{\LTright}{\fill}
\noindent
\begin{tabular}{@{}>{\raggedright\arraybackslash}p{39.46mm}>{\raggedright\arraybackslash}p{39.46mm}>{\raggedright\arraybackslash}p{39.46mm}@{}}
\toprule
\textbf{Case} & \textbf{Singular} & \textbf{Plural} \\
\midrule
Nominative & sacc & saccas \\
Genitive & sacces & sacca \\
Dative & sacce & saccum \\
Accusative & sacc & saccas \\
\bottomrule
\end{tabular}
\par\endgroup

% mdtable:T01:end
\end{minipage}\par
\reserve{7\baselineskip}
% mdparagraph:P022:begin
\subsection*{assa}\BookHeading{nav-ch18-learner-0004}{subsubsection}{assa}{0}
% mdparagraph:P022:end

% teaching: NEW-U18-VOC-ASSA-U18-SENSE
\par\begin{minipage}{\linewidth}\raggedright
% mdparagraph:P023:begin
\textbf{assa} — ass; donkey.
% mdparagraph:P023:end

% teaching: NEW-U18-NOM-ASSA
% mdparagraph:P024:begin
Masculine; weak stem ass-. All oblique singular cases and nominative/\allowbreak{}accusative plural are assan; genitive plural is assena and dative plural assum.
% mdparagraph:P024:end

% table: NEW-U18-NOM-ASSA-DISPLAY
% mdtable:T02:begin
\par\begingroup
\fontsize{10.5}{13.2}\selectfont
\setlength{\LTleft}{0pt}\setlength{\LTright}{\fill}
\noindent
\begin{tabular}{@{}>{\raggedright\arraybackslash}p{39.46mm}>{\raggedright\arraybackslash}p{39.46mm}>{\raggedright\arraybackslash}p{39.46mm}@{}}
\toprule
\textbf{Case} & \textbf{Singular} & \textbf{Plural} \\
\midrule
Nominative & assa & assan \\
Genitive & assan & assena \\
Dative & assan & assum \\
Accusative & assan & assan \\
\bottomrule
\end{tabular}
\par\endgroup

% mdtable:T02:end
\end{minipage}\par
\reserve{7\baselineskip}
% mdparagraph:P025:begin
\subsection*{hunig}\BookHeading{nav-ch18-learner-0005}{subsubsection}{hunig}{0}
% mdparagraph:P025:end

% teaching: NEW-U18-VOC-HUNIG-U18-SENSE
\par\begin{minipage}{\linewidth}\raggedright
% mdparagraph:P026:begin
\textbf{hunig} — honey.
% mdparagraph:P026:end

% teaching: NEW-U18-NOM-HUNIG
% mdparagraph:P027:begin
Neuter; stem hunig-. Retain the medial vowel before endings. This table gives the complete singular of the mass meaning honey.
% mdparagraph:P027:end

% table: NEW-U18-NOM-HUNIG-DISPLAY
% mdtable:T03:begin
\par\begingroup
\fontsize{10.5}{13.2}\selectfont
\setlength{\LTleft}{0pt}\setlength{\LTright}{\fill}
\noindent
\begin{tabular}{@{}>{\raggedright\arraybackslash}p{60.59mm}>{\raggedright\arraybackslash}p{60.59mm}@{}}
\toprule
\textbf{Case} & \textbf{Singular} \\
\midrule
Nominative & hunig \\
Genitive & huniges \\
Dative & hunige \\
Accusative & hunig \\
\bottomrule
\end{tabular}
\par\endgroup

% mdtable:T03:end
\end{minipage}\par
\reserve{7\baselineskip}
% mdparagraph:P028:begin
\subsection*{æcern}\BookHeading{nav-ch18-learner-0006}{subsubsection}{æcern}{0}
% mdparagraph:P028:end

% teaching: NEW-U18-VOC-AECERN-U18-SENSE
\par\begin{minipage}{\linewidth}\raggedright
% mdparagraph:P029:begin
\textbf{æcern} — acorn.
% mdparagraph:P029:end

% teaching: NEW-U18-NOM-AECERN
% mdparagraph:P030:begin
Neuter; stem æcern-. The selected nominative/\allowbreak{}accusative plural is æcernu. Retain the medial vowel in the other cases: final rn makes the second syllable heavy. Preserve unmarked æ.
% mdparagraph:P030:end

% table: NEW-U18-NOM-AECERN-DISPLAY
% mdtable:T04:begin
\par\begingroup
\fontsize{10.5}{13.2}\selectfont
\setlength{\LTleft}{0pt}\setlength{\LTright}{\fill}
\noindent
\begin{tabular}{@{}>{\raggedright\arraybackslash}p{39.46mm}>{\raggedright\arraybackslash}p{39.46mm}>{\raggedright\arraybackslash}p{39.46mm}@{}}
\toprule
\textbf{Case} & \textbf{Singular} & \textbf{Plural} \\
\midrule
Nominative & æcern & æcernu \\
Genitive & æcernes & æcerna \\
Dative & æcerne & æcernum \\
Accusative & æcern & æcernu \\
\bottomrule
\end{tabular}
\par\endgroup

% mdtable:T04:end
\end{minipage}\par
\reserve{7\baselineskip}
% mdparagraph:P031:begin
\subsection*{hwǣte}\BookHeading{nav-ch18-learner-0007}{subsubsection}{hwǣte}{0}
% mdparagraph:P031:end

% teaching: NEW-U18-VOC-HWAETE-U18-SENSE
\par\begin{minipage}{\linewidth}\raggedright
% mdparagraph:P032:begin
\textbf{hwǣte} — wheat.
% mdparagraph:P032:end

% teaching: NEW-U18-NOM-HWAETE
% mdparagraph:P033:begin
Masculine; stem hwǣt-. Nominative and accusative singular retain final -e. This table gives the complete singular of wheat.
% mdparagraph:P033:end

% table: NEW-U18-NOM-HWAETE-DISPLAY
% mdtable:T05:begin
\par\begingroup
\fontsize{10.5}{13.2}\selectfont
\setlength{\LTleft}{0pt}\setlength{\LTright}{\fill}
\noindent
\begin{tabular}{@{}>{\raggedright\arraybackslash}p{60.59mm}>{\raggedright\arraybackslash}p{60.59mm}@{}}
\toprule
\textbf{Case} & \textbf{Singular} \\
\midrule
Nominative & hwǣte \\
Genitive & hwǣtes \\
Dative & hwǣte \\
Accusative & hwǣte \\
\bottomrule
\end{tabular}
\par\endgroup

% mdtable:T05:end
\end{minipage}\par
\reserve{7\baselineskip}
% mdparagraph:P034:begin
\subsection*{læfel}\BookHeading{nav-ch18-learner-0008}{subsubsection}{læfel}{0}
% mdparagraph:P034:end

% teaching: NEW-U18-VOC-LAEFEL-U18-SENSE
\par\begin{minipage}{\linewidth}\raggedright
% mdparagraph:P035:begin
\textbf{læfel} — cup.
% mdparagraph:P035:end

% teaching: NEW-U18-NOM-LAEFEL
% mdparagraph:P036:begin
Masculine; stem læfel-. Retain medial e before endings because the first syllable is short. Preserve unmarked æ.
% mdparagraph:P036:end

% table: NEW-U18-NOM-LAEFEL-DISPLAY
% mdtable:T06:begin
\par\begingroup
\fontsize{10.5}{13.2}\selectfont
\setlength{\LTleft}{0pt}\setlength{\LTright}{\fill}
\noindent
\begin{tabular}{@{}>{\raggedright\arraybackslash}p{39.46mm}>{\raggedright\arraybackslash}p{39.46mm}>{\raggedright\arraybackslash}p{39.46mm}@{}}
\toprule
\textbf{Case} & \textbf{Singular} & \textbf{Plural} \\
\midrule
Nominative & læfel & læfelas \\
Genitive & læfeles & læfela \\
Dative & læfele & læfelum \\
Accusative & læfel & læfelas \\
\bottomrule
\end{tabular}
\par\endgroup

% mdtable:T06:end
\end{minipage}\par
\reserve{7\baselineskip}
% mdparagraph:P037:begin
\subsection*{scēawere}\BookHeading{nav-ch18-learner-0009}{subsubsection}{scēawere}{0}
% mdparagraph:P037:end

% teaching: NEW-U18-VOC-SCEAWERE-U18-SENSE
\par\begin{minipage}{\linewidth}\raggedright
% mdparagraph:P038:begin
\textbf{scēawere} — spy.
% mdparagraph:P038:end

% teaching: NEW-U18-NOM-SCEAWERE
% mdparagraph:P039:begin
Masculine; stem scēawer-. Replace the dictionary form’s final e with the required ending; nominative/\allowbreak{}accusative plural is scēaweras.
% mdparagraph:P039:end

% table: NEW-U18-NOM-SCEAWERE-DISPLAY
% mdtable:T07:begin
\par\begingroup
\fontsize{10.5}{13.2}\selectfont
\setlength{\LTleft}{0pt}\setlength{\LTright}{\fill}
\noindent
\begin{tabular}{@{}>{\raggedright\arraybackslash}p{39.46mm}>{\raggedright\arraybackslash}p{39.46mm}>{\raggedright\arraybackslash}p{39.46mm}@{}}
\toprule
\textbf{Case} & \textbf{Singular} & \textbf{Plural} \\
\midrule
Nominative & scēawere & scēaweras \\
Genitive & scēaweres & scēawera \\
Dative & scēawere & scēawerum \\
Accusative & scēawere & scēaweras \\
\bottomrule
\end{tabular}
\par\endgroup

% mdtable:T07:end
\end{minipage}\par
\reserve{7\baselineskip}
% mdparagraph:P040:begin
\subsection*{ęlcung}\BookHeading{nav-ch18-learner-0010}{subsubsection}{ęlcung}{0}
% mdparagraph:P040:end

% teaching: NEW-U18-VOC-ELCUNG-U18-SENSE
\par\begin{minipage}{\linewidth}\raggedright
% mdparagraph:P041:begin
\textbf{ęlcung} — delay.
% mdparagraph:P041:end

% teaching: NEW-U18-NOM-ELCUNG
% mdparagraph:P042:begin
Feminine; stem ęlcung-. This table selects -e in the oblique singular. The -ung class also admits -a variants. Only abstract singular delay is required.
% mdparagraph:P042:end

% table: NEW-U18-NOM-ELCUNG-DISPLAY
% mdtable:T08:begin
\par\begingroup
\fontsize{10.5}{13.2}\selectfont
\setlength{\LTleft}{0pt}\setlength{\LTright}{\fill}
\noindent
\begin{tabular}{@{}>{\raggedright\arraybackslash}p{60.59mm}>{\raggedright\arraybackslash}p{60.59mm}@{}}
\toprule
\textbf{Case} & \textbf{Singular} \\
\midrule
Nominative & ęlcung \\
Genitive & ęlcunge \\
Dative & ęlcunge \\
Accusative & ęlcunge \\
\bottomrule
\end{tabular}
\par\endgroup

% mdtable:T08:end
\end{minipage}\par
\reserve{7\baselineskip}
% mdparagraph:P043:begin
\subsection*{gebēorscipe}\BookHeading{nav-ch18-learner-0011}{subsubsection}{gebēorscipe}{0}
% mdparagraph:P043:end

% teaching: NEW-U18-VOC-GEBEORSCIPE-U18-SENSE
\par\begin{minipage}{\linewidth}\raggedright
% mdparagraph:P044:begin
\textbf{gebēorscipe} — feast; banquet.
% mdparagraph:P044:end

% teaching: NEW-U18-NOM-GEBEORSCIPE
% mdparagraph:P045:begin
Masculine; stem gebēorscip-. Nominative and accusative singular retain -e; the plural follows the ordinary strong masculine pattern.
% mdparagraph:P045:end

% table: NEW-U18-NOM-GEBEORSCIPE-DISPLAY
% mdtable:T09:begin
\par\begingroup
\fontsize{10.5}{13.2}\selectfont
\setlength{\LTleft}{0pt}\setlength{\LTright}{\fill}
\noindent
\begin{tabular}{@{}>{\raggedright\arraybackslash}p{34.82mm}>{\raggedright\arraybackslash}p{41.78mm}>{\raggedright\arraybackslash}p{41.78mm}@{}}
\toprule
\textbf{Case} & \textbf{Singular} & \textbf{Plural} \\
\midrule
Nominative & gebēorscipe & gebēorscipas \\
Genitive & gebēorscipes & gebēorscipa \\
Dative & gebēorscipe & gebēorscipum \\
Accusative & gebēorscipe & gebēorscipas \\
\bottomrule
\end{tabular}
\par\endgroup

% mdtable:T09:end
\end{minipage}\par
\reserve{7\baselineskip}
% mdparagraph:P046:begin
\subsection*{sārnes}\BookHeading{nav-ch18-learner-0012}{subsubsection}{sārnes}{0}
% mdparagraph:P046:end

% teaching: NEW-U18-VOC-SARNES-U18-SENSE
\par\begin{minipage}{\linewidth}\raggedright
% mdparagraph:P047:begin
\textbf{sārnes} — grief.
% mdparagraph:P047:end

% teaching: NEW-U18-NOM-SARNES
% mdparagraph:P048:begin
Feminine; oblique stem sārness-. Double final s and retain the preceding e: sārnesse. Only abstract singular grief is required.
% mdparagraph:P048:end

% table: NEW-U18-NOM-SARNES-DISPLAY
% mdtable:T10:begin
\par\begingroup
\fontsize{10.5}{13.2}\selectfont
\setlength{\LTleft}{0pt}\setlength{\LTright}{\fill}
\noindent
\begin{tabular}{@{}>{\raggedright\arraybackslash}p{60.59mm}>{\raggedright\arraybackslash}p{60.59mm}@{}}
\toprule
\textbf{Case} & \textbf{Singular} \\
\midrule
Nominative & sārnes \\
Genitive & sārnesse \\
Dative & sārnesse \\
Accusative & sārnesse \\
\bottomrule
\end{tabular}
\par\endgroup

% mdtable:T10:end
\end{minipage}\par
\reserve{7\baselineskip}
% mdparagraph:P049:begin
\subsection*{ealdor}\BookHeading{nav-ch18-learner-0013}{subsubsection}{ealdor}{0}
% mdparagraph:P049:end

% teaching: NEW-U18-VOC-EALDOR-CHIEF-U14-COMPONENT-SOURCE
\par\begin{minipage}{\linewidth}\raggedright
% mdparagraph:P050:begin
\textbf{ealdor} — chief; ruler.
% mdparagraph:P050:end

% teaching: NEW-U18-NOM-EALDOR-CHIEF
% mdparagraph:P051:begin
Masculine; ealdor has the oblique stem ealdr-. Lose the unstressed vowel before an ending: ealdres, ealdre, ealdras.
% mdparagraph:P051:end

% table: NEW-U18-NOM-EALDOR-CHIEF-DISPLAY
% mdtable:T11:begin
\par\begingroup
\fontsize{10.5}{13.2}\selectfont
\setlength{\LTleft}{0pt}\setlength{\LTright}{\fill}
\noindent
\begin{tabular}{@{}>{\raggedright\arraybackslash}p{39.46mm}>{\raggedright\arraybackslash}p{39.46mm}>{\raggedright\arraybackslash}p{39.46mm}@{}}
\toprule
\textbf{Case} & \textbf{Singular} & \textbf{Plural} \\
\midrule
Nominative & ealdor & ealdras \\
Genitive & ealdres & ealdra \\
Dative & ealdre & ealdrum \\
Accusative & ealdor & ealdras \\
\bottomrule
\end{tabular}
\par\endgroup

% mdtable:T11:end
\end{minipage}\par
\reserve{7\baselineskip}
% mdparagraph:P052:begin
\subsection*{ðearf}\BookHeading{nav-ch18-learner-0014}{subsubsection}{ðearf}{0}
% mdparagraph:P052:end

% teaching: NEW-U18-VOC-THEARF-U15-SOURCE-SENSE
\par\begin{minipage}{\linewidth}\raggedright
% mdparagraph:P053:begin
\textbf{ðearf} — benefit; good.
% mdparagraph:P053:end

% teaching: NEW-U18-NOM-THEARF
% mdparagraph:P054:begin
Feminine; stem ðearf-. This unit uses the abstract singular benefit or good; the oblique singular is ðearfe.
% mdparagraph:P054:end

% table: NEW-U18-NOM-THEARF-DISPLAY
% mdtable:T12:begin
\par\begingroup
\fontsize{10.5}{13.2}\selectfont
\setlength{\LTleft}{0pt}\setlength{\LTright}{\fill}
\noindent
\begin{tabular}{@{}>{\raggedright\arraybackslash}p{60.59mm}>{\raggedright\arraybackslash}p{60.59mm}@{}}
\toprule
\textbf{Case} & \textbf{Singular} \\
\midrule
Nominative & ðearf \\
Genitive & ðearfe \\
Dative & ðearfe \\
Accusative & ðearfe \\
\bottomrule
\end{tabular}
\par\endgroup

% mdtable:T12:end
\end{minipage}\par
\reserve{7\baselineskip}
% mdparagraph:P055:begin
\subsection*{ðēow}\BookHeading{nav-ch18-learner-0015}{subsubsection}{ðēow}{0}
% mdparagraph:P055:end

% teaching: NEW-U18-VOC-THEOW-SERVANT-U18-SENSE
\par\begin{minipage}{\linewidth}\raggedright
% mdparagraph:P056:begin
\textbf{ðēow} — servant; enslaved person.
% mdparagraph:P056:end

% teaching: NEW-U18-NOM-THEOW-SERVANT
% mdparagraph:P057:begin
Masculine; strong stem ðēow-. Distinguish this noun from the earlier weak ðēowa. Its selected genitive plural is ðēowa, while the weak noun has ðēowena. The complete phrase in the reading has local help.
% mdparagraph:P057:end

% table: NEW-U18-NOM-THEOW-SERVANT-DISPLAY
% mdtable:T13:begin
\par\begingroup
\fontsize{10.5}{13.2}\selectfont
\setlength{\LTleft}{0pt}\setlength{\LTright}{\fill}
\noindent
\begin{tabular}{@{}>{\raggedright\arraybackslash}p{39.46mm}>{\raggedright\arraybackslash}p{39.46mm}>{\raggedright\arraybackslash}p{39.46mm}@{}}
\toprule
\textbf{Case} & \textbf{Singular} & \textbf{Plural} \\
\midrule
Nominative & ðēow & ðēowas \\
Genitive & ðēowes & ðēowa \\
Dative & ðēowe & ðēowum \\
Accusative & ðēow & ðēowas \\
\bottomrule
\end{tabular}
\par\endgroup

% mdtable:T13:end
\end{minipage}\par
\reserve{7\baselineskip}
% mdparagraph:P058:begin
\subsection*{hunger}\BookHeading{nav-ch18-learner-0016}{subsubsection}{hunger}{0}
% mdparagraph:P058:end

% teaching: NEW-U18-VOC-HUNGER-U18-SENSE
\par\begin{minipage}{\linewidth}\raggedright
% mdparagraph:P059:begin
\textbf{hunger} — famine; hunger.
% mdparagraph:P059:end

% teaching: NEW-U18-NOM-HUNGER
% mdparagraph:P060:begin
Masculine; hunger has the oblique stem hungr-. Retain hunger without an ending; use hungres and hungre before endings. This table gives the mass singular famine or hunger.
% mdparagraph:P060:end

% table: NEW-U18-NOM-HUNGER-DISPLAY
% mdtable:T14:begin
\par\begingroup
\fontsize{10.5}{13.2}\selectfont
\setlength{\LTleft}{0pt}\setlength{\LTright}{\fill}
\noindent
\begin{tabular}{@{}>{\raggedright\arraybackslash}p{60.59mm}>{\raggedright\arraybackslash}p{60.59mm}@{}}
\toprule
\textbf{Case} & \textbf{Singular} \\
\midrule
Nominative & hunger \\
Genitive & hungres \\
Dative & hungre \\
Accusative & hunger \\
\bottomrule
\end{tabular}
\par\endgroup

% mdtable:T14:end
\end{minipage}\par
\reserve{7\baselineskip}
% mdparagraph:P061:begin
\subsection*{sīð}\BookHeading{nav-ch18-learner-0017}{subsubsection}{sīð}{0}
% mdparagraph:P061:end

% teaching: NEW-U18-VOC-SITH-U18-SENSE
\par\begin{minipage}{\linewidth}\raggedright
% mdparagraph:P062:begin
\textbf{sīð} — time; occasion.
% mdparagraph:P062:end

% teaching: NEW-U18-NOM-SITH
% mdparagraph:P063:begin
Masculine; stem sīð-. These are the four ordinary cases. The second-time expression in 44.7 is instrumental and is supplied separately; its noun is not dative merely because both end in -e.
% mdparagraph:P063:end

% table: NEW-U18-NOM-SITH-DISPLAY
% mdtable:T15:begin
\par\begingroup
\fontsize{10.5}{13.2}\selectfont
\setlength{\LTleft}{0pt}\setlength{\LTright}{\fill}
\noindent
\begin{tabular}{@{}>{\raggedright\arraybackslash}p{39.46mm}>{\raggedright\arraybackslash}p{39.46mm}>{\raggedright\arraybackslash}p{39.46mm}@{}}
\toprule
\textbf{Case} & \textbf{Singular} & \textbf{Plural} \\
\midrule
Nominative & sīð & sīðas \\
Genitive & sīðes & sīða \\
Dative & sīðe & sīðum \\
Accusative & sīð & sīðas \\
\bottomrule
\end{tabular}
\par\endgroup

% mdtable:T15:end
\end{minipage}\par
\reserve{7\baselineskip}
% mdparagraph:P064:begin
\subsection*{lāc}\BookHeading{nav-ch18-learner-0018}{subsubsection}{lāc}{0}
% mdparagraph:P064:end

% teaching: NEW-U18-VOC-LAC-U18-SENSE
\par\begin{minipage}{\linewidth}\raggedright
% mdparagraph:P065:begin
\textbf{lāc} — gift; present.
% mdparagraph:P065:end

% teaching: NEW-U18-NOM-LAC
% mdparagraph:P066:begin
Neuter; stem lāc-. The long vowel makes the stem heavy, so nominative/\allowbreak{}accusative plural has no ending.
% mdparagraph:P066:end

% table: NEW-U18-NOM-LAC-DISPLAY
% mdtable:T16:begin
\par\begingroup
\fontsize{10.5}{13.2}\selectfont
\setlength{\LTleft}{0pt}\setlength{\LTright}{\fill}
\noindent
\begin{tabular}{@{}>{\raggedright\arraybackslash}p{39.46mm}>{\raggedright\arraybackslash}p{39.46mm}>{\raggedright\arraybackslash}p{39.46mm}@{}}
\toprule
\textbf{Case} & \textbf{Singular} & \textbf{Plural} \\
\midrule
Nominative & lāc & lāc \\
Genitive & lāces & lāca \\
Dative & lāce & lācum \\
Accusative & lāc & lāc \\
\bottomrule
\end{tabular}
\par\endgroup

% mdtable:T16:end
\end{minipage}\par
\reserve{7\baselineskip}
% mdparagraph:P067:begin
\subsection*{fæt}\BookHeading{nav-ch18-learner-0019}{subsubsection}{fæt}{0}
% mdparagraph:P067:end

% teaching: NEW-U18-VOC-FAET-U18-SENSE
\par\begin{minipage}{\linewidth}\raggedright
% mdparagraph:P068:begin
\textbf{fæt} — vessel; container.
% mdparagraph:P068:end

% teaching: NEW-U18-NOM-FAET
% mdparagraph:P069:begin
Neuter; retain æ throughout the singular and use a throughout the plural: fatu, fata, fatum. The short stem takes -u in nominative/\allowbreak{}accusative plural.
% mdparagraph:P069:end

% table: NEW-U18-NOM-FAET-DISPLAY
% mdtable:T17:begin
\par\begingroup
\fontsize{10.5}{13.2}\selectfont
\setlength{\LTleft}{0pt}\setlength{\LTright}{\fill}
\noindent
\begin{tabular}{@{}>{\raggedright\arraybackslash}p{39.46mm}>{\raggedright\arraybackslash}p{39.46mm}>{\raggedright\arraybackslash}p{39.46mm}@{}}
\toprule
\textbf{Case} & \textbf{Singular} & \textbf{Plural} \\
\midrule
Nominative & fæt & fatu \\
Genitive & fætes & fata \\
Dative & fæte & fatum \\
Accusative & fæt & fatu \\
\bottomrule
\end{tabular}
\par\endgroup

% mdtable:T17:end
\end{minipage}\par
\reserve{7\baselineskip}
% mdparagraph:P070:begin
\subsection*{weorð}\BookHeading{nav-ch18-learner-0020}{subsubsection}{weorð}{0}
% mdparagraph:P070:end

% teaching: NEW-U18-VOC-WEORTH-PRICE-U18-SENSE
\par\begin{minipage}{\linewidth}\raggedright
% mdparagraph:P071:begin
\textbf{weorð} — price; value.
% mdparagraph:P071:end

% teaching: NEW-U18-NOM-WEORTH-PRICE
% mdparagraph:P072:begin
Neuter; stem weorð-. The table gives the complete singular of price or value.
% mdparagraph:P072:end

% table: NEW-U18-NOM-WEORTH-PRICE-DISPLAY
% mdtable:T18:begin
\par\begingroup
\fontsize{10.5}{13.2}\selectfont
\setlength{\LTleft}{0pt}\setlength{\LTright}{\fill}
\noindent
\begin{tabular}{@{}>{\raggedright\arraybackslash}p{60.59mm}>{\raggedright\arraybackslash}p{60.59mm}@{}}
\toprule
\textbf{Case} & \textbf{Singular} \\
\midrule
Nominative & weorð \\
Genitive & weorðes \\
Dative & weorðe \\
Accusative & weorð \\
\bottomrule
\end{tabular}
\par\endgroup

% mdtable:T18:end
\end{minipage}\par
\reserve{7\baselineskip}
% mdparagraph:P073:begin
\subsection*{feoh}\BookHeading{nav-ch18-learner-0021}{subsubsection}{feoh}{0}
% mdparagraph:P073:end

% teaching: NEW-U18-VOC-FEOH-U18-SENSE
\par\begin{minipage}{\linewidth}\raggedright
% mdparagraph:P074:begin
\textbf{feoh} — money.
% mdparagraph:P074:end

% teaching: NEW-U18-NOM-FEOH
% mdparagraph:P075:begin
Neuter; money. Final h disappears before endings: genitive fēos and selected dative fēo. The source prints feoh without a macron before h; preserve that spelling. An absent printed macron does not assert that the historical vowel was short.
% mdparagraph:P075:end

% table: NEW-U18-NOM-FEOH-DISPLAY
% mdtable:T19:begin
\par\begingroup
\fontsize{10.5}{13.2}\selectfont
\setlength{\LTleft}{0pt}\setlength{\LTright}{\fill}
\noindent
\begin{tabular}{@{}>{\raggedright\arraybackslash}p{60.59mm}>{\raggedright\arraybackslash}p{60.59mm}@{}}
\toprule
\textbf{Case} & \textbf{Singular} \\
\midrule
Nominative & feoh \\
Genitive & fēos \\
Dative & fēo \\
Accusative & feoh \\
\bottomrule
\end{tabular}
\par\endgroup

% mdtable:T19:end
\end{minipage}\par
\reserve{7\baselineskip}
% mdparagraph:P076:begin
\subsection*{Strong and weak adjectives}\BookHeading{nav-ch18-learner-0022}{subsubsection}{Strong and weak adjectives}{0}
% mdparagraph:P076:end

% mdparagraph:P077:begin
Adjectives agree in gender, number and case. Use the strong pattern in its familiar undetermined uses and the weak pattern after the determiners that require it. The tables give every ordinary case in all three genders and both numbers. Strong masculine and neuter instrumental singular cells are recognition only; they introduce no new instrumental production.
% mdparagraph:P077:end

\reserve{7\baselineskip}
% mdparagraph:P078:begin
\subsection*{scyldig}\BookHeading{nav-ch18-learner-0023}{subsubsection}{scyldig}{0}
% mdparagraph:P078:end

% teaching: NEW-U18-VOC-SCYLDIG-U18-SENSE
\par\begin{minipage}{\linewidth}\raggedright
% mdparagraph:P079:begin
\textbf{scyldig} — guilty; culpable.
% mdparagraph:P079:end

% teaching: NEW-U18-ADJ-SCYLDIG
% mdparagraph:P080:begin
Use the uncontracted stem scyldig-. Strong feminine nominative singular and neuter nominative/\allowbreak{}accusative plural retain -u; contraction is a documented alternative.
% mdparagraph:P080:end

% mdparagraph:P081:begin
\textbf{scyldig — strong}
% mdparagraph:P081:end

% table: NEW-U18-ADJ-SCYLDIG-STRONG-DISPLAY
% mdtable:T20:begin
% course-table-style: 2026-10
% course-panel:T20-P01:begin
\par\begingroup
\fontsize{10.5}{13.2}\selectfont
\setlength{\LTleft}{0pt}\setlength{\LTright}{\fill}
\noindent
\begin{tabular}{@{}>{\raggedright\arraybackslash}p{22.0mm}>{\raggedright\arraybackslash}p{31.188mm}>{\raggedright\arraybackslash}p{31.188mm}>{\raggedright\arraybackslash}p{31.188mm}@{}}
\toprule
\textbf{Case} & \textbf{Masculine singular} & \textbf{Neuter singular} & \textbf{Feminine singular} \\
\midrule
Nominative & scyldig & scyldig & scyldigu \\
Genitive & scyldiges & scyldiges & scyldigre \\
Dative & scyldigum & scyldigum & scyldigre \\
Accusative & scyldigne & scyldig & scyldige \\
Instrumental & scyldige & scyldige & — \\
\bottomrule
\end{tabular}
\par\endgroup
% course-panel:T20-P01:end
\end{minipage}\par
\par\addvspace{5pt}
% course-panel:T20-P02:begin
\par\begingroup
\fontsize{10.5}{13.2}\selectfont
\setlength{\LTleft}{0pt}\setlength{\LTright}{\fill}
\noindent
\begin{tabular}{@{}>{\raggedright\arraybackslash}p{22.0mm}>{\raggedright\arraybackslash}p{31.188mm}>{\raggedright\arraybackslash}p{31.188mm}>{\raggedright\arraybackslash}p{31.188mm}@{}}
\toprule
\textbf{Case} & \textbf{Masculine plural} & \textbf{Neuter plural} & \textbf{Feminine plural} \\
\midrule
Nominative & scyldige & scyldigu & scyldiga \\
Genitive & scyldigra & scyldigra & scyldigra \\
Dative & scyldigum & scyldigum & scyldigum \\
Accusative & scyldige & scyldigu & scyldiga \\
Instrumental & — & — & — \\
\bottomrule
\end{tabular}
\par\endgroup
% course-panel:T20-P02:end
% mdtable:T20:end

\reserve{18\baselineskip}
\reserve{16\baselineskip}
% mdparagraph:P082:begin
\textbf{scyldig — weak}
% mdparagraph:P082:end

% table: NEW-U18-ADJ-SCYLDIG-WEAK-DISPLAY
% mdtable:T21:begin
% course-table-style: 2026-10
% course-panel:T21-P01:begin
\par\begingroup
\fontsize{10.5}{13.2}\selectfont
\setlength{\LTleft}{0pt}\setlength{\LTright}{\fill}
\noindent
\begin{tabular}{@{}>{\raggedright\arraybackslash}p{22.0mm}>{\raggedright\arraybackslash}p{31.188mm}>{\raggedright\arraybackslash}p{31.188mm}>{\raggedright\arraybackslash}p{31.188mm}@{}}
\toprule
\textbf{Case} & \textbf{Masculine singular} & \textbf{Neuter singular} & \textbf{Feminine singular} \\
\midrule
Nominative & scyldiga & scyldige & scyldige \\
Genitive & scyldigan & scyldigan & scyldigan \\
Dative & scyldigan & scyldigan & scyldigan \\
Accusative & scyldigan & scyldige & scyldigan \\
\bottomrule
\end{tabular}
\par\endgroup
% course-panel:T21-P01:end
\par\addvspace{5pt}
% course-panel:T21-P02:begin
\par\begingroup
\fontsize{10.5}{13.2}\selectfont
\setlength{\LTleft}{0pt}\setlength{\LTright}{\fill}
\noindent
\begin{tabular}{@{}>{\raggedright\arraybackslash}p{22.0mm}>{\raggedright\arraybackslash}p{31.188mm}>{\raggedright\arraybackslash}p{31.188mm}>{\raggedright\arraybackslash}p{31.188mm}@{}}
\toprule
\textbf{Case} & \textbf{Masculine plural} & \textbf{Neuter plural} & \textbf{Feminine plural} \\
\midrule
Nominative & scyldigan & scyldigan & scyldigan \\
Genitive & scyldigra & scyldigra & scyldigra \\
Dative & scyldigum & scyldigum & scyldigum \\
Accusative & scyldigan & scyldigan & scyldigan \\
\bottomrule
\end{tabular}
\par\endgroup
% course-panel:T21-P02:end
% mdtable:T21:end

\reserve{7\baselineskip}
% mdparagraph:P083:begin
\subsection*{twīfeald}\BookHeading{nav-ch18-learner-0024}{subsubsection}{twīfeald}{0}
% mdparagraph:P083:end

% teaching: NEW-U18-VOC-TWIFEALD-U18-SENSE
\par\begin{minipage}{\linewidth}\raggedright
% mdparagraph:P084:begin
\textbf{twīfeald} — double; twofold.
% mdparagraph:P084:end

% teaching: NEW-U18-ADJ-TWIFEALD
% mdparagraph:P085:begin
The heavy stem twīfeald- takes no -u in strong feminine nominative singular or neuter nominative/\allowbreak{}accusative plural.
% mdparagraph:P085:end

% mdparagraph:P086:begin
\textbf{twīfeald — strong}
% mdparagraph:P086:end

% table: NEW-U18-ADJ-TWIFEALD-STRONG-DISPLAY
% mdtable:T22:begin
% course-table-style: 2026-10
% course-panel:T22-P01:begin
\par\begingroup
\fontsize{10.5}{13.2}\selectfont
\setlength{\LTleft}{0pt}\setlength{\LTright}{\fill}
\noindent
\begin{tabular}{@{}>{\raggedright\arraybackslash}p{22.0mm}>{\raggedright\arraybackslash}p{31.188mm}>{\raggedright\arraybackslash}p{31.188mm}>{\raggedright\arraybackslash}p{31.188mm}@{}}
\toprule
\textbf{Case} & \textbf{Masculine singular} & \textbf{Neuter singular} & \textbf{Feminine singular} \\
\midrule
Nominative & twīfeald & twīfeald & twīfeald \\
Genitive & twīfealdes & twīfealdes & twīfealdre \\
Dative & twīfealdum & twīfealdum & twīfealdre \\
Accusative & twīfealdne & twīfeald & twīfealde \\
Instrumental & twīfealde & twīfealde & — \\
\bottomrule
\end{tabular}
\par\endgroup
% course-panel:T22-P01:end
\end{minipage}\par
\par\addvspace{5pt}
% course-panel:T22-P02:begin
\par\begingroup
\fontsize{10.5}{13.2}\selectfont
\setlength{\LTleft}{0pt}\setlength{\LTright}{\fill}
\noindent
\begin{tabular}{@{}>{\raggedright\arraybackslash}p{22.0mm}>{\raggedright\arraybackslash}p{31.188mm}>{\raggedright\arraybackslash}p{31.188mm}>{\raggedright\arraybackslash}p{31.188mm}@{}}
\toprule
\textbf{Case} & \textbf{Masculine plural} & \textbf{Neuter plural} & \textbf{Feminine plural} \\
\midrule
Nominative & twīfealde & twīfeald & twīfealda \\
Genitive & twīfealdra & twīfealdra & twīfealdra \\
Dative & twīfealdum & twīfealdum & twīfealdum \\
Accusative & twīfealde & twīfeald & twīfealda \\
Instrumental & — & — & — \\
\bottomrule
\end{tabular}
\par\endgroup
% course-panel:T22-P02:end
% mdtable:T22:end

\reserve{18\baselineskip}
\reserve{16\baselineskip}
% mdparagraph:P087:begin
\textbf{twīfeald — weak}
% mdparagraph:P087:end

% table: NEW-U18-ADJ-TWIFEALD-WEAK-DISPLAY
% mdtable:T23:begin
% course-table-style: 2026-10
% course-panel:T23-P01:begin
\par\begingroup
\fontsize{10.5}{13.2}\selectfont
\setlength{\LTleft}{0pt}\setlength{\LTright}{\fill}
\noindent
\begin{tabular}{@{}>{\raggedright\arraybackslash}p{22.0mm}>{\raggedright\arraybackslash}p{31.188mm}>{\raggedright\arraybackslash}p{31.188mm}>{\raggedright\arraybackslash}p{31.188mm}@{}}
\toprule
\textbf{Case} & \textbf{Masculine singular} & \textbf{Neuter singular} & \textbf{Feminine singular} \\
\midrule
Nominative & twīfealda & twīfealde & twīfealde \\
Genitive & twīfealdan & twīfealdan & twīfealdan \\
Dative & twīfealdan & twīfealdan & twīfealdan \\
Accusative & twīfealdan & twīfealde & twīfealdan \\
\bottomrule
\end{tabular}
\par\endgroup
% course-panel:T23-P01:end
\par\addvspace{5pt}
% course-panel:T23-P02:begin
\par\begingroup
\fontsize{10.5}{13.2}\selectfont
\setlength{\LTleft}{0pt}\setlength{\LTright}{\fill}
\noindent
\begin{tabular}{@{}>{\raggedright\arraybackslash}p{22.0mm}>{\raggedright\arraybackslash}p{31.188mm}>{\raggedright\arraybackslash}p{31.188mm}>{\raggedright\arraybackslash}p{31.188mm}@{}}
\toprule
\textbf{Case} & \textbf{Masculine plural} & \textbf{Neuter plural} & \textbf{Feminine plural} \\
\midrule
Nominative & twīfealdan & twīfealdan & twīfealdan \\
Genitive & twīfealdra & twīfealdra & twīfealdra \\
Dative & twīfealdum & twīfealdum & twīfealdum \\
Accusative & twīfealdan & twīfealdan & twīfealdan \\
\bottomrule
\end{tabular}
\par\endgroup
% course-panel:T23-P02:end
% mdtable:T23:end

\reserve{7\baselineskip}
% mdparagraph:P088:begin
\subsection*{gesund}\BookHeading{nav-ch18-learner-0025}{subsubsection}{gesund}{0}
% mdparagraph:P088:end

% teaching: NEW-U18-VOC-GESUND-U18-SENSE
\par\begin{minipage}{\linewidth}\raggedright
% mdparagraph:P089:begin
\textbf{gesund} — well; in good health.
% mdparagraph:P089:end

% teaching: NEW-U18-ADJ-GESUND
% mdparagraph:P090:begin
The heavy stem gesund- takes no -u in strong feminine nominative singular or neuter nominative/\allowbreak{}accusative plural.
% mdparagraph:P090:end

% mdparagraph:P091:begin
\textbf{gesund — strong}
% mdparagraph:P091:end

% table: NEW-U18-ADJ-GESUND-STRONG-DISPLAY
% mdtable:T24:begin
% course-table-style: 2026-10
% course-panel:T24-P01:begin
\par\begingroup
\fontsize{10.5}{13.2}\selectfont
\setlength{\LTleft}{0pt}\setlength{\LTright}{\fill}
\noindent
\begin{tabular}{@{}>{\raggedright\arraybackslash}p{22.0mm}>{\raggedright\arraybackslash}p{31.188mm}>{\raggedright\arraybackslash}p{31.188mm}>{\raggedright\arraybackslash}p{31.188mm}@{}}
\toprule
\textbf{Case} & \textbf{Masculine singular} & \textbf{Neuter singular} & \textbf{Feminine singular} \\
\midrule
Nominative & gesund & gesund & gesund \\
Genitive & gesundes & gesundes & gesundre \\
Dative & gesundum & gesundum & gesundre \\
Accusative & gesundne & gesund & gesunde \\
Instrumental & gesunde & gesunde & — \\
\bottomrule
\end{tabular}
\par\endgroup
% course-panel:T24-P01:end
\end{minipage}\par
\par\addvspace{5pt}
% course-panel:T24-P02:begin
\par\begingroup
\fontsize{10.5}{13.2}\selectfont
\setlength{\LTleft}{0pt}\setlength{\LTright}{\fill}
\noindent
\begin{tabular}{@{}>{\raggedright\arraybackslash}p{22.0mm}>{\raggedright\arraybackslash}p{31.188mm}>{\raggedright\arraybackslash}p{31.188mm}>{\raggedright\arraybackslash}p{31.188mm}@{}}
\toprule
\textbf{Case} & \textbf{Masculine plural} & \textbf{Neuter plural} & \textbf{Feminine plural} \\
\midrule
Nominative & gesunde & gesund & gesunda \\
Genitive & gesundra & gesundra & gesundra \\
Dative & gesundum & gesundum & gesundum \\
Accusative & gesunde & gesund & gesunda \\
Instrumental & — & — & — \\
\bottomrule
\end{tabular}
\par\endgroup
% course-panel:T24-P02:end
% mdtable:T24:end

\reserve{18\baselineskip}
\reserve{16\baselineskip}
% mdparagraph:P092:begin
\textbf{gesund — weak}
% mdparagraph:P092:end

% table: NEW-U18-ADJ-GESUND-WEAK-DISPLAY
% mdtable:T25:begin
% course-table-style: 2026-10
% course-panel:T25-P01:begin
\par\begingroup
\fontsize{10.5}{13.2}\selectfont
\setlength{\LTleft}{0pt}\setlength{\LTright}{\fill}
\noindent
\begin{tabular}{@{}>{\raggedright\arraybackslash}p{22.0mm}>{\raggedright\arraybackslash}p{31.188mm}>{\raggedright\arraybackslash}p{31.188mm}>{\raggedright\arraybackslash}p{31.188mm}@{}}
\toprule
\textbf{Case} & \textbf{Masculine singular} & \textbf{Neuter singular} & \textbf{Feminine singular} \\
\midrule
Nominative & gesunda & gesunde & gesunde \\
Genitive & gesundan & gesundan & gesundan \\
Dative & gesundan & gesundan & gesundan \\
Accusative & gesundan & gesunde & gesundan \\
\bottomrule
\end{tabular}
\par\endgroup
% course-panel:T25-P01:end
\par\addvspace{5pt}
% course-panel:T25-P02:begin
\par\begingroup
\fontsize{10.5}{13.2}\selectfont
\setlength{\LTleft}{0pt}\setlength{\LTright}{\fill}
\noindent
\begin{tabular}{@{}>{\raggedright\arraybackslash}p{22.0mm}>{\raggedright\arraybackslash}p{31.188mm}>{\raggedright\arraybackslash}p{31.188mm}>{\raggedright\arraybackslash}p{31.188mm}@{}}
\toprule
\textbf{Case} & \textbf{Masculine plural} & \textbf{Neuter plural} & \textbf{Feminine plural} \\
\midrule
Nominative & gesundan & gesundan & gesundan \\
Genitive & gesundra & gesundra & gesundra \\
Dative & gesundum & gesundum & gesundum \\
Accusative & gesundan & gesundan & gesundan \\
\bottomrule
\end{tabular}
\par\endgroup
% course-panel:T25-P02:end
% mdtable:T25:end

\reserve{7\baselineskip}
% mdparagraph:P093:begin
\subsection*{hāl}\BookHeading{nav-ch18-learner-0026}{subsubsection}{hāl}{0}
% mdparagraph:P093:end

% teaching: NEW-U18-VOC-HAL-U18-HEALTH
\par\begin{minipage}{\linewidth}\raggedright
% mdparagraph:P094:begin
\textbf{hāl} — hale; healthy.
% mdparagraph:P094:end

% teaching: NEW-U18-ADJ-HAL
% mdparagraph:P095:begin
The heavy stem hāl- takes no -u in strong feminine nominative singular or neuter nominative/\allowbreak{}accusative plural. The meaning here is health.
% mdparagraph:P095:end

% mdparagraph:P096:begin
\textbf{hāl — strong}
% mdparagraph:P096:end

% table: NEW-U18-ADJ-HAL-STRONG-DISPLAY
% mdtable:T26:begin
% course-table-style: 2026-10
% course-panel:T26-P01:begin
\par\begingroup
\fontsize{10.5}{13.2}\selectfont
\setlength{\LTleft}{0pt}\setlength{\LTright}{\fill}
\noindent
\begin{tabular}{@{}>{\raggedright\arraybackslash}p{22.0mm}>{\raggedright\arraybackslash}p{31.188mm}>{\raggedright\arraybackslash}p{31.188mm}>{\raggedright\arraybackslash}p{31.188mm}@{}}
\toprule
\textbf{Case} & \textbf{Masculine singular} & \textbf{Neuter singular} & \textbf{Feminine singular} \\
\midrule
Nominative & Hāl & hāl & hāl \\
Genitive & hāles & hāles & hālre \\
Dative & hālum & hālum & hālre \\
Accusative & hālne & hāl & hāle \\
Instrumental & hāle & hāle & — \\
\bottomrule
\end{tabular}
\par\endgroup
% course-panel:T26-P01:end
\end{minipage}\par
\par\addvspace{5pt}
% course-panel:T26-P02:begin
\par\begingroup
\fontsize{10.5}{13.2}\selectfont
\setlength{\LTleft}{0pt}\setlength{\LTright}{\fill}
\noindent
\begin{tabular}{@{}>{\raggedright\arraybackslash}p{22.0mm}>{\raggedright\arraybackslash}p{31.188mm}>{\raggedright\arraybackslash}p{31.188mm}>{\raggedright\arraybackslash}p{31.188mm}@{}}
\toprule
\textbf{Case} & \textbf{Masculine plural} & \textbf{Neuter plural} & \textbf{Feminine plural} \\
\midrule
Nominative & hāle & hāl & hāla \\
Genitive & hālra & hālra & hālra \\
Dative & hālum & hālum & hālum \\
Accusative & hāle & hāl & hāla \\
Instrumental & — & — & — \\
\bottomrule
\end{tabular}
\par\endgroup
% course-panel:T26-P02:end
% mdtable:T26:end

\reserve{18\baselineskip}
\reserve{16\baselineskip}
% mdparagraph:P097:begin
\textbf{hāl — weak}
% mdparagraph:P097:end

% table: NEW-U18-ADJ-HAL-WEAK-DISPLAY
% mdtable:T27:begin
% course-table-style: 2026-10
% course-panel:T27-P01:begin
\par\begingroup
\fontsize{10.5}{13.2}\selectfont
\setlength{\LTleft}{0pt}\setlength{\LTright}{\fill}
\noindent
\begin{tabular}{@{}>{\raggedright\arraybackslash}p{22.0mm}>{\raggedright\arraybackslash}p{31.188mm}>{\raggedright\arraybackslash}p{31.188mm}>{\raggedright\arraybackslash}p{31.188mm}@{}}
\toprule
\textbf{Case} & \textbf{Masculine singular} & \textbf{Neuter singular} & \textbf{Feminine singular} \\
\midrule
Nominative & hāla & hāle & hāle \\
Genitive & hālan & hālan & hālan \\
Dative & hālan & hālan & hālan \\
Accusative & hālan & hāle & hālan \\
\bottomrule
\end{tabular}
\par\endgroup
% course-panel:T27-P01:end
\par\addvspace{5pt}
% course-panel:T27-P02:begin
\par\begingroup
\fontsize{10.5}{13.2}\selectfont
\setlength{\LTleft}{0pt}\setlength{\LTright}{\fill}
\noindent
\begin{tabular}{@{}>{\raggedright\arraybackslash}p{22.0mm}>{\raggedright\arraybackslash}p{31.188mm}>{\raggedright\arraybackslash}p{31.188mm}>{\raggedright\arraybackslash}p{31.188mm}@{}}
\toprule
\textbf{Case} & \textbf{Masculine plural} & \textbf{Neuter plural} & \textbf{Feminine plural} \\
\midrule
Nominative & hālan & hālan & hālan \\
Genitive & hālra & hālra & hālra \\
Dative & hālum & hālum & hālum \\
Accusative & hālan & hālan & hālan \\
\bottomrule
\end{tabular}
\par\endgroup
% course-panel:T27-P02:end
% mdtable:T27:end

\reserve{7\baselineskip}
% mdparagraph:P098:begin
\subsection*{āgen}\BookHeading{nav-ch18-learner-0027}{subsubsection}{āgen}{0}
% mdparagraph:P098:end

% teaching: NEW-U18-VOC-AGEN-U16-SENSE
\par\begin{minipage}{\linewidth}\raggedright
% mdparagraph:P099:begin
\textbf{āgen} — own.
% mdparagraph:P099:end

% teaching: NEW-U18-ADJ-AGEN
% mdparagraph:P100:begin
Use āgn- before vowel-initial endings but āgen- before consonantal endings: compare āgnes, āgnum, āgna, āgnu with āgenne, āgenre, āgenra.
% mdparagraph:P100:end

% mdparagraph:P101:begin
\textbf{āgen — strong}
% mdparagraph:P101:end

% table: NEW-U18-ADJ-AGEN-STRONG-DISPLAY
% mdtable:T28:begin
% course-table-style: 2026-10
% course-panel:T28-P01:begin
\par\begingroup
\fontsize{10.5}{13.2}\selectfont
\setlength{\LTleft}{0pt}\setlength{\LTright}{\fill}
\noindent
\begin{tabular}{@{}>{\raggedright\arraybackslash}p{22.0mm}>{\raggedright\arraybackslash}p{31.188mm}>{\raggedright\arraybackslash}p{31.188mm}>{\raggedright\arraybackslash}p{31.188mm}@{}}
\toprule
\textbf{Case} & \textbf{Masculine singular} & \textbf{Neuter singular} & \textbf{Feminine singular} \\
\midrule
Nominative & āgen & āgen & āgnu \\
Genitive & āgnes & āgnes & āgenre \\
Dative & āgnum & āgnum & āgenre \\
Accusative & āgenne & āgen & āgne \\
Instrumental & āgne & āgne & — \\
\bottomrule
\end{tabular}
\par\endgroup
% course-panel:T28-P01:end
\end{minipage}\par
\par\addvspace{5pt}
% course-panel:T28-P02:begin
\par\begingroup
\fontsize{10.5}{13.2}\selectfont
\setlength{\LTleft}{0pt}\setlength{\LTright}{\fill}
\noindent
\begin{tabular}{@{}>{\raggedright\arraybackslash}p{22.0mm}>{\raggedright\arraybackslash}p{31.188mm}>{\raggedright\arraybackslash}p{31.188mm}>{\raggedright\arraybackslash}p{31.188mm}@{}}
\toprule
\textbf{Case} & \textbf{Masculine plural} & \textbf{Neuter plural} & \textbf{Feminine plural} \\
\midrule
Nominative & āgne & āgnu & āgna \\
Genitive & āgenra & āgenra & āgenra \\
Dative & āgnum & āgnum & āgnum \\
Accusative & āgne & āgnu & āgna \\
Instrumental & — & — & — \\
\bottomrule
\end{tabular}
\par\endgroup
% course-panel:T28-P02:end
% mdtable:T28:end

\reserve{18\baselineskip}
\reserve{16\baselineskip}
% mdparagraph:P102:begin
\textbf{āgen — weak}
% mdparagraph:P102:end

% table: NEW-U18-ADJ-AGEN-WEAK-DISPLAY
% mdtable:T29:begin
% course-table-style: 2026-10
% course-panel:T29-P01:begin
\par\begingroup
\fontsize{10.5}{13.2}\selectfont
\setlength{\LTleft}{0pt}\setlength{\LTright}{\fill}
\noindent
\begin{tabular}{@{}>{\raggedright\arraybackslash}p{22.0mm}>{\raggedright\arraybackslash}p{31.188mm}>{\raggedright\arraybackslash}p{31.188mm}>{\raggedright\arraybackslash}p{31.188mm}@{}}
\toprule
\textbf{Case} & \textbf{Masculine singular} & \textbf{Neuter singular} & \textbf{Feminine singular} \\
\midrule
Nominative & āgna & āgne & āgne \\
Genitive & āgnan & āgnan & āgnan \\
Dative & āgnan & āgnan & āgnan \\
Accusative & āgnan & āgne & āgnan \\
\bottomrule
\end{tabular}
\par\endgroup
% course-panel:T29-P01:end
\par\addvspace{5pt}
% course-panel:T29-P02:begin
\par\begingroup
\fontsize{10.5}{13.2}\selectfont
\setlength{\LTleft}{0pt}\setlength{\LTright}{\fill}
\noindent
\begin{tabular}{@{}>{\raggedright\arraybackslash}p{22.0mm}>{\raggedright\arraybackslash}p{31.188mm}>{\raggedright\arraybackslash}p{31.188mm}>{\raggedright\arraybackslash}p{31.188mm}@{}}
\toprule
\textbf{Case} & \textbf{Masculine plural} & \textbf{Neuter plural} & \textbf{Feminine plural} \\
\midrule
Nominative & āgnan & āgnan & āgnan \\
Genitive & āgenra & āgenra & āgenra \\
Dative & āgnum & āgnum & āgnum \\
Accusative & āgnan & āgnan & āgnan \\
\bottomrule
\end{tabular}
\par\endgroup
% course-panel:T29-P02:end
% mdtable:T29:end

\reserve{7\baselineskip}
% mdparagraph:P103:begin
\subsection*{clǣne}\BookHeading{nav-ch18-learner-0028}{subsubsection}{clǣne}{0}
% mdparagraph:P103:end

% teaching: NEW-U18-VOC-CLAENE-U18-BLAMELESS
\par\begin{minipage}{\linewidth}\raggedright
% mdparagraph:P104:begin
\textbf{clǣne} — clean; pure; blameless.
% mdparagraph:P104:end

% teaching: NEW-U18-ADJ-CLAENE
% mdparagraph:P105:begin
Use clǣn- before endings. Retain dictionary -e where no other ending is present. In the accusation, the sense is blameless or free from guilt.
% mdparagraph:P105:end

% mdparagraph:P106:begin
\textbf{clǣne — strong}
% mdparagraph:P106:end

% table: NEW-U18-ADJ-CLAENE-STRONG-DISPLAY
% mdtable:T30:begin
% course-table-style: 2026-10
% course-panel:T30-P01:begin
\par\begingroup
\fontsize{10.5}{13.2}\selectfont
\setlength{\LTleft}{0pt}\setlength{\LTright}{\fill}
\noindent
\begin{tabular}{@{}>{\raggedright\arraybackslash}p{22.0mm}>{\raggedright\arraybackslash}p{31.188mm}>{\raggedright\arraybackslash}p{31.188mm}>{\raggedright\arraybackslash}p{31.188mm}@{}}
\toprule
\textbf{Case} & \textbf{Masculine singular} & \textbf{Neuter singular} & \textbf{Feminine singular} \\
\midrule
Nominative & clǣne & clǣne & clǣnu \\
Genitive & clǣnes & clǣnes & clǣnre \\
Dative & clǣnum & clǣnum & clǣnre \\
Accusative & clǣnne & clǣne & clǣne \\
Instrumental & clǣne & clǣne & — \\
\bottomrule
\end{tabular}
\par\endgroup
% course-panel:T30-P01:end
\end{minipage}\par
\par\addvspace{5pt}
% course-panel:T30-P02:begin
\par\begingroup
\fontsize{10.5}{13.2}\selectfont
\setlength{\LTleft}{0pt}\setlength{\LTright}{\fill}
\noindent
\begin{tabular}{@{}>{\raggedright\arraybackslash}p{22.0mm}>{\raggedright\arraybackslash}p{31.188mm}>{\raggedright\arraybackslash}p{31.188mm}>{\raggedright\arraybackslash}p{31.188mm}@{}}
\toprule
\textbf{Case} & \textbf{Masculine plural} & \textbf{Neuter plural} & \textbf{Feminine plural} \\
\midrule
Nominative & clǣne & clǣnu & clǣna \\
Genitive & clǣnra & clǣnra & clǣnra \\
Dative & clǣnum & clǣnum & clǣnum \\
Accusative & clǣne & clǣnu & clǣna \\
Instrumental & — & — & — \\
\bottomrule
\end{tabular}
\par\endgroup
% course-panel:T30-P02:end
% mdtable:T30:end

\reserve{18\baselineskip}
\reserve{16\baselineskip}
% mdparagraph:P107:begin
\textbf{clǣne — weak}
% mdparagraph:P107:end

% table: NEW-U18-ADJ-CLAENE-WEAK-DISPLAY
% mdtable:T31:begin
% course-table-style: 2026-10
% course-panel:T31-P01:begin
\par\begingroup
\fontsize{10.5}{13.2}\selectfont
\setlength{\LTleft}{0pt}\setlength{\LTright}{\fill}
\noindent
\begin{tabular}{@{}>{\raggedright\arraybackslash}p{22.0mm}>{\raggedright\arraybackslash}p{31.188mm}>{\raggedright\arraybackslash}p{31.188mm}>{\raggedright\arraybackslash}p{31.188mm}@{}}
\toprule
\textbf{Case} & \textbf{Masculine singular} & \textbf{Neuter singular} & \textbf{Feminine singular} \\
\midrule
Nominative & clǣna & clǣne & clǣne \\
Genitive & clǣnan & clǣnan & clǣnan \\
Dative & clǣnan & clǣnan & clǣnan \\
Accusative & clǣnan & clǣne & clǣnan \\
\bottomrule
\end{tabular}
\par\endgroup
% course-panel:T31-P01:end
\par\addvspace{5pt}
% course-panel:T31-P02:begin
\par\begingroup
\fontsize{10.5}{13.2}\selectfont
\setlength{\LTleft}{0pt}\setlength{\LTright}{\fill}
\noindent
\begin{tabular}{@{}>{\raggedright\arraybackslash}p{22.0mm}>{\raggedright\arraybackslash}p{31.188mm}>{\raggedright\arraybackslash}p{31.188mm}>{\raggedright\arraybackslash}p{31.188mm}@{}}
\toprule
\textbf{Case} & \textbf{Masculine plural} & \textbf{Neuter plural} & \textbf{Feminine plural} \\
\midrule
Nominative & clǣnan & clǣnan & clǣnan \\
Genitive & clǣnra & clǣnra & clǣnra \\
Dative & clǣnum & clǣnum & clǣnum \\
Accusative & clǣnan & clǣnan & clǣnan \\
\bottomrule
\end{tabular}
\par\endgroup
% course-panel:T31-P02:end
% mdtable:T31:end

\reserve{7\baselineskip}
% mdparagraph:P108:begin
\subsection*{ęlðēodig}\BookHeading{nav-ch18-learner-0029}{subsubsection}{ęlðēodig}{0}
% mdparagraph:P108:end

% teaching: NEW-U18-VOC-ELTHEODIG-U18-FOREIGN
\par\begin{minipage}{\linewidth}\raggedright
% mdparagraph:P109:begin
\textbf{ęlðēodig} — foreign.
% mdparagraph:P109:end

% teaching: NEW-U18-ADJ-ELTHEODIG
% mdparagraph:P110:begin
Use the uncontracted stem ęlðēodig-. Strong feminine nominative singular and neuter nominative/\allowbreak{}accusative plural retain -u.
% mdparagraph:P110:end

% mdparagraph:P111:begin
\textbf{ęlðēodig — strong}
% mdparagraph:P111:end

% table: NEW-U18-ADJ-ELTHEODIG-STRONG-DISPLAY
% mdtable:T32:begin
% course-table-style: 2026-10
% course-panel:T32-P01:begin
\par\begingroup
\fontsize{10.5}{13.2}\selectfont
\setlength{\LTleft}{0pt}\setlength{\LTright}{\fill}
\noindent
\begin{tabular}{@{}>{\raggedright\arraybackslash}p{22.0mm}>{\raggedright\arraybackslash}p{31.188mm}>{\raggedright\arraybackslash}p{31.188mm}>{\raggedright\arraybackslash}p{31.188mm}@{}}
\toprule
\textbf{Case} & \textbf{Masculine singular} & \textbf{Neuter singular} & \textbf{Feminine singular} \\
\midrule
Nominative & ęlðēodig & ęlðēodig & ęlðēodigu \\
Genitive & ęlðēodiges & ęlðēodiges & ęlðēodigre \\
Dative & ęlðēodigum & ęlðēodigum & ęlðēodigre \\
Accusative & ęlðēodigne & ęlðēodig & ęlðēodige \\
Instrumental & ęlðēodige & ęlðēodige & — \\
\bottomrule
\end{tabular}
\par\endgroup
% course-panel:T32-P01:end
\end{minipage}\par
\par\addvspace{5pt}
% course-panel:T32-P02:begin
\par\begingroup
\fontsize{10.5}{13.2}\selectfont
\setlength{\LTleft}{0pt}\setlength{\LTright}{\fill}
\noindent
\begin{tabular}{@{}>{\raggedright\arraybackslash}p{22.0mm}>{\raggedright\arraybackslash}p{31.188mm}>{\raggedright\arraybackslash}p{31.188mm}>{\raggedright\arraybackslash}p{31.188mm}@{}}
\toprule
\textbf{Case} & \textbf{Masculine plural} & \textbf{Neuter plural} & \textbf{Feminine plural} \\
\midrule
Nominative & ęlðēodige & ęlðēodigu & ęlðēodiga \\
Genitive & ęlðēodigra & ęlðēodigra & ęlðēodigra \\
Dative & ęlðēodigum & ęlðēodigum & ęlðēodigum \\
Accusative & ęlðēodige & ęlðēodigu & ęlðēodiga \\
Instrumental & — & — & — \\
\bottomrule
\end{tabular}
\par\endgroup
% course-panel:T32-P02:end
% mdtable:T32:end

\reserve{18\baselineskip}
\reserve{16\baselineskip}
% mdparagraph:P112:begin
\textbf{ęlðēodig — weak}
% mdparagraph:P112:end

% table: NEW-U18-ADJ-ELTHEODIG-WEAK-DISPLAY
% mdtable:T33:begin
% course-table-style: 2026-10
% course-panel:T33-P01:begin
\par\begingroup
\fontsize{10.5}{13.2}\selectfont
\setlength{\LTleft}{0pt}\setlength{\LTright}{\fill}
\noindent
\begin{tabular}{@{}>{\raggedright\arraybackslash}p{22.0mm}>{\raggedright\arraybackslash}p{31.188mm}>{\raggedright\arraybackslash}p{31.188mm}>{\raggedright\arraybackslash}p{31.188mm}@{}}
\toprule
\textbf{Case} & \textbf{Masculine singular} & \textbf{Neuter singular} & \textbf{Feminine singular} \\
\midrule
Nominative & ęlðēodiga & ęlðēodige & ęlðēodige \\
Genitive & ęlðēodigan & ęlðēodigan & ęlðēodigan \\
Dative & ęlðēodigan & ęlðēodigan & ęlðēodigan \\
Accusative & ęlðēodigan & ęlðēodige & ęlðēodigan \\
\bottomrule
\end{tabular}
\par\endgroup
% course-panel:T33-P01:end
\par\addvspace{5pt}
% course-panel:T33-P02:begin
\par\begingroup
\fontsize{10.5}{13.2}\selectfont
\setlength{\LTleft}{0pt}\setlength{\LTright}{\fill}
\noindent
\begin{tabular}{@{}>{\raggedright\arraybackslash}p{22.0mm}>{\raggedright\arraybackslash}p{31.188mm}>{\raggedright\arraybackslash}p{31.188mm}>{\raggedright\arraybackslash}p{31.188mm}@{}}
\toprule
\textbf{Case} & \textbf{Masculine plural} & \textbf{Neuter plural} & \textbf{Feminine plural} \\
\midrule
Nominative & ęlðēodigan & ęlðēodigan & ęlðēodigan \\
Genitive & ęlðēodigra & ęlðēodigra & ęlðēodigra \\
Dative & ęlðēodigum & ęlðēodigum & ęlðēodigum \\
Accusative & ęlðēodigan & ęlðēodigan & ęlðēodigan \\
\bottomrule
\end{tabular}
\par\endgroup
% course-panel:T33-P02:end
% mdtable:T33:end

\reserve{7\baselineskip}
% mdparagraph:P113:begin
\section*{3. Verbal systems}\BookHeading{nav-ch18-learner-0030}{subsection}{3. Verbal systems}{1}
% mdparagraph:P113:end

% mdparagraph:P114:begin
Each heading states the complete system added or reviewed. Indicative tables distinguish present and preterite. Subjunctive tables give all persons of the named tense; identical forms do not remove person distinctions. Commands distinguish singular and plural addressees. Infinitive and participle pairs add no unlisted finite forms.
% mdparagraph:P114:end

% mdparagraph:P115:begin
A modal takes a plain infinitive. The inflected infinitive requires its stated construction, commonly with tō. A base past participle may occur in the specified passive or result predicate without a new adjective declension. Meaning and complementation remain separate from the form.
% mdparagraph:P115:end

\reserve{7\baselineskip}
% mdparagraph:P116:begin
\subsection*{beðurfan: present and preterite indicative}\BookHeading{nav-ch18-learner-0031}{subsubsection}{beðurfan: present and preterite indicative}{0}
% mdparagraph:P116:end

% teaching: NEW-U18-VOC-BETHURFAN-U18-SELECTED
\par\begin{minipage}{\linewidth}\raggedright
% mdparagraph:P117:begin
\textbf{beðurfan} — need; require.
% mdparagraph:P117:end

% teaching: u18-bethurfan-indicative12
% mdparagraph:P118:begin
This preterite-present has present singular beðearf, beðearft, beðearf and plural beðurfon. The dental preterite uses beðorfte- in the singular and beðorft- in the plural. The complete indicative is assigned; the headword alone does not introduce an infinitive construction or subjunctive.
% mdparagraph:P118:end

% mdparagraph:P119:begin
\textbf{Diagnostic forms:} beðearf (present indicative first singular); beðurfon (present indicative plural); beðorfte (preterite indicative first singular); beðorfton (preterite indicative plural).
% mdparagraph:P119:end

% table: u18-bethurfan-indicative12-INDICATIVE-DISPLAY
% mdtable:T34:begin
% course-table-style: 2026-10
% course-panel:T34-P01:begin
\par\begingroup
\fontsize{10.5}{13.2}\selectfont
\setlength{\LTleft}{0pt}\setlength{\LTright}{\fill}
\noindent
\begin{tabular}{@{}>{\raggedright\arraybackslash}p{39.459mm}>{\raggedright\arraybackslash}p{39.459mm}>{\raggedright\arraybackslash}p{39.459mm}@{}}
\toprule
\textbf{Person and number} & \textbf{Present} & \textbf{Preterite} \\
\midrule
1 sg. & beðearf & beðorfte \\
2 sg. & beðearft & beðorftest \\
3 sg. & beðearf & beðorfte \\
1–3 pl. & beðurfon & beðorfton \\
\bottomrule
\end{tabular}
\par\endgroup
% course-panel:T34-P01:end
% mdtable:T34:end
\end{minipage}\par
\reserve{7\baselineskip}
% mdparagraph:P120:begin
\subsection*{oncnāwan: present and preterite indicative}\BookHeading{nav-ch18-learner-0032}{subsubsection}{oncnāwan: present and preterite indicative}{0}
% mdparagraph:P120:end

% teaching: NEW-U18-VOC-ONCNAWAN-U18-SELECTED
\par\begin{minipage}{\linewidth}\raggedright
% mdparagraph:P121:begin
\textbf{oncnāwan} — recognize; know by recognition.
% mdparagraph:P121:end

% teaching: u18-oncnawan-indicative12
% mdparagraph:P122:begin
The present stem oncnāw- changes ā to ǣ in the second and third singular. Preterite oncnēow- adds -e in the second singular; first and third singular have no ending. The recognized person is accusative.
% mdparagraph:P122:end

% mdparagraph:P123:begin
\textbf{Diagnostic forms:} oncnāwe (present indicative first singular); oncnāwað (present indicative plural); oncnēow (preterite indicative first singular); oncnēowon (preterite indicative plural).
% mdparagraph:P123:end

% table: u18-oncnawan-indicative12-INDICATIVE-DISPLAY
% mdtable:T35:begin
% course-table-style: 2026-10
% course-panel:T35-P01:begin
\par\begingroup
\fontsize{10.5}{13.2}\selectfont
\setlength{\LTleft}{0pt}\setlength{\LTright}{\fill}
\noindent
\begin{tabular}{@{}>{\raggedright\arraybackslash}p{39.459mm}>{\raggedright\arraybackslash}p{39.459mm}>{\raggedright\arraybackslash}p{39.459mm}@{}}
\toprule
\textbf{Person and number} & \textbf{Present} & \textbf{Preterite} \\
\midrule
1 sg. & oncnāwe & oncnēow \\
2 sg. & oncnǣwst & oncnēowe \\
3 sg. & oncnǣwð & oncnēow \\
1–3 pl. & oncnāwað & oncnēowon \\
\bottomrule
\end{tabular}
\par\endgroup
% course-panel:T35-P01:end
% mdtable:T35:end
\end{minipage}\par
\reserve{7\baselineskip}
% mdparagraph:P124:begin
\subsection*{undōn: present and preterite indicative}\BookHeading{nav-ch18-learner-0033}{subsubsection}{undōn: present and preterite indicative}{0}
% mdparagraph:P124:end

% teaching: NEW-U18-VOC-UNDON-U18-SELECTED
\par\begin{minipage}{\linewidth}\raggedright
% mdparagraph:P125:begin
\textbf{undōn} — undo; open.
% mdparagraph:P125:end

% teaching: u18-undon-indicative12
% mdparagraph:P126:begin
Retain un- with the irregular dōn stems: undō, undēst, undēð, undōð; preterite undyde, undydest, undydon. The accusative object is the tied container opened.
% mdparagraph:P126:end

% mdparagraph:P127:begin
\textbf{Diagnostic forms:} undō (present indicative first singular); undōð (present indicative plural); undyde (preterite indicative first singular); undydon (preterite indicative plural).
% mdparagraph:P127:end

% table: u18-undon-indicative12-INDICATIVE-DISPLAY
% mdtable:T36:begin
% course-table-style: 2026-10
% course-panel:T36-P01:begin
\par\begingroup
\fontsize{10.5}{13.2}\selectfont
\setlength{\LTleft}{0pt}\setlength{\LTright}{\fill}
\noindent
\begin{tabular}{@{}>{\raggedright\arraybackslash}p{39.459mm}>{\raggedright\arraybackslash}p{39.459mm}>{\raggedright\arraybackslash}p{39.459mm}@{}}
\toprule
\textbf{Person and number} & \textbf{Present} & \textbf{Preterite} \\
\midrule
1 sg. & undō & undyde \\
2 sg. & undēst & undydest \\
3 sg. & undēð & undyde \\
1–3 pl. & undōð & undydon \\
\bottomrule
\end{tabular}
\par\endgroup
% course-panel:T36-P01:end
% mdtable:T36:end
\end{minipage}\par
\reserve{7\baselineskip}
% mdparagraph:P128:begin
\subsection*{lęcgean: present and preterite indicative}\BookHeading{nav-ch18-learner-0034}{subsubsection}{lęcgean: present and preterite indicative}{0}
% mdparagraph:P128:end

% teaching: NEW-U18-VOC-LECGEAN-U18-SELECTED
\par\begin{minipage}{\linewidth}\raggedright
% mdparagraph:P129:begin
\textbf{lęcgean} — lay; place.
% mdparagraph:P129:end

% teaching: u18-lecgean-indicative12
% mdparagraph:P130:begin
The short weak verb has lęcge, lęgest, lęgeð and plural lęcgeað. The doubled consonant is absent in the second and third singular. Present cg alternates with preterite g: lęgde, lęgdest, lęgdon. These tables select the uncontracted past; lēde/\allowbreak{}lēdon are documented variants. The reading’s imperative lęcgeað is separately supplied.
% mdparagraph:P130:end

% mdparagraph:P131:begin
\textbf{Diagnostic forms:} lęcge (present indicative first singular); lęcgeað (present indicative plural); lęgde (preterite indicative first singular); lęgdon (preterite indicative plural).
% mdparagraph:P131:end

% table: u18-lecgean-indicative12-INDICATIVE-DISPLAY
% mdtable:T37:begin
% course-table-style: 2026-10
% course-panel:T37-P01:begin
\par\begingroup
\fontsize{10.5}{13.2}\selectfont
\setlength{\LTleft}{0pt}\setlength{\LTright}{\fill}
\noindent
\begin{tabular}{@{}>{\raggedright\arraybackslash}p{39.459mm}>{\raggedright\arraybackslash}p{39.459mm}>{\raggedright\arraybackslash}p{39.459mm}@{}}
\toprule
\textbf{Person and number} & \textbf{Present} & \textbf{Preterite} \\
\midrule
1 sg. & lęcge & lęgde \\
2 sg. & lęgest & lęgdest \\
3 sg. & lęgeð & lęgde \\
1–3 pl. & lęcgeað & lęgdon \\
\bottomrule
\end{tabular}
\par\endgroup
% course-panel:T37-P01:end
% mdtable:T37:end
\end{minipage}\par
\reserve{7\baselineskip}
% mdparagraph:P132:begin
\subsection*{foresęcgean: present and preterite indicative}\BookHeading{nav-ch18-learner-0035}{subsubsection}{foresęcgean: present and preterite indicative}{0}
% mdparagraph:P132:end

% teaching: NEW-U18-VOC-FORESECGEAN-U18-SELECTED
\par\begin{minipage}{\linewidth}\raggedright
% mdparagraph:P133:begin
\textbf{foresęcgean} — mention before; tell of beforehand.
% mdparagraph:P133:end

% teaching: u18-foresecgean-indicative12
% mdparagraph:P134:begin
Add fore- to the irregular sęcgan pattern. Select plural foresęcgeað, second/\allowbreak{}third singular foresægst and foresægð, and preterite foresægde, foresægdest, foresægdon. Retain unmarked æ before gd; contracted sāde is a separate alternative. Here the verb means mention earlier, not predict the future.
% mdparagraph:P134:end

% mdparagraph:P135:begin
\textbf{Diagnostic forms:} foresęcge (present indicative first singular); foresęcgeað (present indicative plural); foresægde (preterite indicative first singular); foresægdon (preterite indicative plural).
% mdparagraph:P135:end

% table: u18-foresecgean-indicative12-INDICATIVE-DISPLAY
% mdtable:T38:begin
% course-table-style: 2026-10
% course-panel:T38-P01:begin
\par\begingroup
\fontsize{10.5}{13.2}\selectfont
\setlength{\LTleft}{0pt}\setlength{\LTright}{\fill}
\noindent
\begin{tabular}{@{}>{\raggedright\arraybackslash}p{39.459mm}>{\raggedright\arraybackslash}p{39.459mm}>{\raggedright\arraybackslash}p{39.459mm}@{}}
\toprule
\textbf{Person and number} & \textbf{Present} & \textbf{Preterite} \\
\midrule
1 sg. & foresęcge & foresægde \\
2 sg. & foresægst & foresægdest \\
3 sg. & foresægð & foresægde \\
1–3 pl. & foresęcgeað & foresægdon \\
\bottomrule
\end{tabular}
\par\endgroup
% course-panel:T38-P01:end
% mdtable:T38:end
\end{minipage}\par
\reserve{7\baselineskip}
% mdparagraph:P136:begin
\subsection*{forstelan: present and preterite indicative}\BookHeading{nav-ch18-learner-0036}{subsubsection}{forstelan: present and preterite indicative}{0}
% mdparagraph:P136:end

% teaching: NEW-U18-VOC-FORSTELAN-U18-SELECTED
\par\begin{minipage}{\linewidth}\raggedright
% mdparagraph:P137:begin
\textbf{forstelan} — steal away; steal.
% mdparagraph:P137:end

% teaching: u18-forstelan-indicative12
% mdparagraph:P138:begin
Strong class IV alternates e, æ, ǣ, o. Present second and third singular have i: forstilst, forstilð. Preterite first/\allowbreak{}third singular is forstæl, second singular forstǣle, and plural forstǣlon. Keep lexical for-; unprefixed stelan is listed separately.
% mdparagraph:P138:end

% mdparagraph:P139:begin
\textbf{Diagnostic forms:} forstele (present indicative first singular); forstelað (present indicative plural); forstæl (preterite indicative first singular); forstǣlon (preterite indicative plural).
% mdparagraph:P139:end

% table: u18-forstelan-indicative12-INDICATIVE-DISPLAY
% mdtable:T39:begin
% course-table-style: 2026-10
% course-panel:T39-P01:begin
\par\begingroup
\fontsize{10.5}{13.2}\selectfont
\setlength{\LTleft}{0pt}\setlength{\LTright}{\fill}
\noindent
\begin{tabular}{@{}>{\raggedright\arraybackslash}p{39.459mm}>{\raggedright\arraybackslash}p{39.459mm}>{\raggedright\arraybackslash}p{39.459mm}@{}}
\toprule
\textbf{Person and number} & \textbf{Present} & \textbf{Preterite} \\
\midrule
1 sg. & forstele & forstæl \\
2 sg. & forstilst & forstǣle \\
3 sg. & forstilð & forstæl \\
1–3 pl. & forstelað & forstǣlon \\
\bottomrule
\end{tabular}
\par\endgroup
% course-panel:T39-P01:end
% mdtable:T39:end
\end{minipage}\par
\reserve{7\baselineskip}
% mdparagraph:P140:begin
\subsection*{cēapian: indicative and infinitives}\BookHeading{nav-ch18-learner-0037}{subsubsection}{cēapian: indicative and infinitives}{0}
% mdparagraph:P140:end

% teaching: NEW-U18-VOC-CEAPIAN-U18-SELECTED
% mdparagraph:P141:begin
\textbf{cēapian} — trade; traffic; buy through trade.
% mdparagraph:P141:end

% teaching: u18-ceapian-indicative12-infinitive2
% mdparagraph:P142:begin
Weak class II has present -ie, -ast, -að, -iað and preterite -ode, -odest, -odon. Retain ēa. The complete infinitive pair is cēapian/\allowbreak{}cēapianne; the special buying-purpose construction remains supplied.
% mdparagraph:P142:end

% mdparagraph:P143:begin
\textbf{Diagnostic forms:} cēapie (present indicative first singular); cēapiað (present indicative plural); cēapode (preterite indicative first singular); cēapodon (preterite indicative plural); cēapian (plain\_infinitive); cēapianne (inflected\_infinitive).
% mdparagraph:P143:end

\reserve{13\baselineskip}
\reserve{11\baselineskip}
% mdparagraph:P144:begin
\textbf{cēapian — indicative}
% mdparagraph:P144:end

% table: u18-ceapian-indicative12-infinitive2-INDICATIVE-DISPLAY
% mdtable:T40:begin
% course-table-style: 2026-10
% course-panel:T40-P01:begin
\par\begingroup
\fontsize{10.5}{13.2}\selectfont
\setlength{\LTleft}{0pt}\setlength{\LTright}{\fill}
\noindent
\begin{tabular}{@{}>{\raggedright\arraybackslash}p{39.459mm}>{\raggedright\arraybackslash}p{39.459mm}>{\raggedright\arraybackslash}p{39.459mm}@{}}
\toprule
\textbf{Person and number} & \textbf{Present} & \textbf{Preterite} \\
\midrule
1 sg. & cēapie & cēapode \\
2 sg. & cēapast & cēapodest \\
3 sg. & cēapað & cēapode \\
1–3 pl. & cēapiað & cēapodon \\
\bottomrule
\end{tabular}
\par\endgroup
% course-panel:T40-P01:end
% mdtable:T40:end
\reserve{9\baselineskip}
\reserve{7\baselineskip}
% mdparagraph:P145:begin
\textbf{cēapian — additional bounded forms}
% mdparagraph:P145:end

% table: u18-ceapian-indicative12-infinitive2-BOUNDED-NONFINITE-DISPLAY
% mdtable:T41:begin
\par\begingroup
\fontsize{10.5}{13.2}\selectfont
\setlength{\LTleft}{0pt}\setlength{\LTright}{\fill}
\noindent
\begin{tabular}{@{}>{\raggedright\arraybackslash}p{71.72mm}>{\raggedright\arraybackslash}p{49.46mm}@{}}
\toprule
\textbf{Type} & \textbf{Form} \\
\midrule
Plain infinitive & cēapian \\
Inflected infinitive & cēapianne \\
\bottomrule
\end{tabular}
\par\endgroup

% mdtable:T41:end
\reserve{7\baselineskip}
% mdparagraph:P146:begin
\subsection*{gedōn: present subjunctive}\BookHeading{nav-ch18-learner-0038}{subsubsection}{gedōn: present subjunctive}{0}
% mdparagraph:P146:end

% teaching: NEW-U18-VOC-GEDON-U18-SELECTED
\par\begin{minipage}{\linewidth}\raggedright
% mdparagraph:P147:begin
\textbf{gedōn} — make; cause to be.
% mdparagraph:P147:end

% teaching: u18-gedon-present-subjunctive6
% mdparagraph:P148:begin
Present subjunctive is gedō in all singular persons and gedōn in all plural persons. The causative construction takes an accusative person and an agreeing adjective.
% mdparagraph:P148:end

% mdparagraph:P149:begin
\textbf{Diagnostic forms:} gedō (present subjunctive sg); gedōn (present subjunctive pl).
% mdparagraph:P149:end

% table: u18-gedon-present-subjunctive6-SUBJUNCTIVE-DISPLAY
% mdtable:T42:begin
% course-table-style: 2026-10
% course-panel:T42-P01:begin
\par\begingroup
\fontsize{10.5}{13.2}\selectfont
\setlength{\LTleft}{0pt}\setlength{\LTright}{\fill}
\noindent
\begin{tabular}{@{}>{\raggedright\arraybackslash}p{63.342mm}>{\raggedright\arraybackslash}p{57.846mm}@{}}
\toprule
\textbf{Person and number} & \textbf{Present} \\
\midrule
1 sg. & gedō \\
2 sg. & gedō \\
3 sg. & gedō \\
1–3 pl. & gedōn \\
\bottomrule
\end{tabular}
\par\endgroup
% course-panel:T42-P01:end
% mdtable:T42:end
\end{minipage}\par
\reserve{7\baselineskip}
% mdparagraph:P150:begin
\subsection*{āgiefan: present subjunctive}\BookHeading{nav-ch18-learner-0039}{subsubsection}{āgiefan: present subjunctive}{0}
% mdparagraph:P150:end

% teaching: NEW-U18-VOC-AGIEFAN-U18-SELECTED
\par\begin{minipage}{\linewidth}\raggedright
% mdparagraph:P151:begin
\textbf{āgiefan} — give back; restore.
% mdparagraph:P151:end

% teaching: u18-agiefan-present-subjunctive6
% mdparagraph:P152:begin
Retain ā- with unchanged gief-. Present subjunctive singular ends in -e and plural in -en. The returned person is accusative and the recipient dative.
% mdparagraph:P152:end

% mdparagraph:P153:begin
\textbf{Diagnostic forms:} āgiefe (present subjunctive sg); āgiefen (present subjunctive pl).
% mdparagraph:P153:end

% table: u18-agiefan-present-subjunctive6-SUBJUNCTIVE-DISPLAY
% mdtable:T43:begin
% course-table-style: 2026-10
% course-panel:T43-P01:begin
\par\begingroup
\fontsize{10.5}{13.2}\selectfont
\setlength{\LTleft}{0pt}\setlength{\LTright}{\fill}
\noindent
\begin{tabular}{@{}>{\raggedright\arraybackslash}p{63.342mm}>{\raggedright\arraybackslash}p{57.846mm}@{}}
\toprule
\textbf{Person and number} & \textbf{Present} \\
\midrule
1 sg. & āgiefe \\
2 sg. & āgiefe \\
3 sg. & āgiefe \\
1–3 pl. & āgiefen \\
\bottomrule
\end{tabular}
\par\endgroup
% course-panel:T43-P01:end
% mdtable:T43:end
\end{minipage}\par
\reserve{7\baselineskip}
% mdparagraph:P154:begin
\subsection*{hlystan: present subjunctive}\BookHeading{nav-ch18-learner-0040}{subsubsection}{hlystan: present subjunctive}{0}
% mdparagraph:P154:end

% teaching: NEW-U18-VOC-HLYSTAN-U18-SELECTED
\par\begin{minipage}{\linewidth}\raggedright
% mdparagraph:P155:begin
\textbf{hlystan} — listen to; heed.
% mdparagraph:P155:end

% teaching: u18-hlystan-present-subjunctive6
% mdparagraph:P156:begin
Long weak hlyst- takes present subjunctive -e/\allowbreak{}-en. Speech or content heeded is genitive.
% mdparagraph:P156:end

% mdparagraph:P157:begin
\textbf{Diagnostic forms:} hlyste (present subjunctive sg); hlysten (present subjunctive pl).
% mdparagraph:P157:end

% table: u18-hlystan-present-subjunctive6-SUBJUNCTIVE-DISPLAY
% mdtable:T44:begin
% course-table-style: 2026-10
% course-panel:T44-P01:begin
\par\begingroup
\fontsize{10.5}{13.2}\selectfont
\setlength{\LTleft}{0pt}\setlength{\LTright}{\fill}
\noindent
\begin{tabular}{@{}>{\raggedright\arraybackslash}p{63.342mm}>{\raggedright\arraybackslash}p{57.846mm}@{}}
\toprule
\textbf{Person and number} & \textbf{Present} \\
\midrule
1 sg. & hlyste \\
2 sg. & hlyste \\
3 sg. & hlyste \\
1–3 pl. & hlysten \\
\bottomrule
\end{tabular}
\par\endgroup
% course-panel:T44-P01:end
% mdtable:T44:end
\end{minipage}\par
\reserve{7\baselineskip}
% mdparagraph:P158:begin
\subsection*{ðēowian: present subjunctive}\BookHeading{nav-ch18-learner-0041}{subsubsection}{ðēowian: present subjunctive}{0}
% mdparagraph:P158:end

% teaching: NEW-U18-VOC-THEOWIAN-U18-SELECTED
\par\begin{minipage}{\linewidth}\raggedright
% mdparagraph:P159:begin
\textbf{ðēowian} — enslave; reduce to servitude.
% mdparagraph:P159:end

% teaching: u18-theowian-present-subjunctive6
% mdparagraph:P160:begin
Weak class II present subjunctive has -ie in the singular and -ien in the plural. The selected sense is enslave, with an accusative person. The source’s -ige variant is the same ending.
% mdparagraph:P160:end

% mdparagraph:P161:begin
\textbf{Diagnostic forms:} ðēowie (present subjunctive sg); ðēowien (present subjunctive pl).
% mdparagraph:P161:end

% table: u18-theowian-present-subjunctive6-SUBJUNCTIVE-DISPLAY
% mdtable:T45:begin
% course-table-style: 2026-10
% course-panel:T45-P01:begin
\par\begingroup
\fontsize{10.5}{13.2}\selectfont
\setlength{\LTleft}{0pt}\setlength{\LTright}{\fill}
\noindent
\begin{tabular}{@{}>{\raggedright\arraybackslash}p{63.342mm}>{\raggedright\arraybackslash}p{57.846mm}@{}}
\toprule
\textbf{Person and number} & \textbf{Present} \\
\midrule
1 sg. & ðēowie \\
2 sg. & ðēowie \\
3 sg. & ðēowie \\
1–3 pl. & ðēowien \\
\bottomrule
\end{tabular}
\par\endgroup
% course-panel:T45-P01:end
% mdtable:T45:end
\end{minipage}\par
\reserve{7\baselineskip}
% mdparagraph:P162:begin
\subsection*{berēafian: present subjunctive}\BookHeading{nav-ch18-learner-0042}{subsubsection}{berēafian: present subjunctive}{0}
% mdparagraph:P162:end

% teaching: NEW-U18-VOC-BEREAFIAN-U18-SELECTED
\par\begin{minipage}{\linewidth}\raggedright
% mdparagraph:P163:begin
\textbf{berēafian} — despoil; deprive of.
% mdparagraph:P163:end

% teaching: u18-bereafian-present-subjunctive6
% mdparagraph:P164:begin
Weak class II gives present subjunctive -ie/\allowbreak{}-ien. The deprivation complement is explained below.
% mdparagraph:P164:end

% mdparagraph:P165:begin
\textbf{Diagnostic forms:} berēafie (present subjunctive sg); berēafien (present subjunctive pl).
% mdparagraph:P165:end

% table: u18-bereafian-present-subjunctive6-SUBJUNCTIVE-DISPLAY
% mdtable:T46:begin
% course-table-style: 2026-10
% course-panel:T46-P01:begin
\par\begingroup
\fontsize{10.5}{13.2}\selectfont
\setlength{\LTleft}{0pt}\setlength{\LTright}{\fill}
\noindent
\begin{tabular}{@{}>{\raggedright\arraybackslash}p{63.342mm}>{\raggedright\arraybackslash}p{57.846mm}@{}}
\toprule
\textbf{Person and number} & \textbf{Present} \\
\midrule
1 sg. & berēafie \\
2 sg. & berēafie \\
3 sg. & berēafie \\
1–3 pl. & berēafien \\
\bottomrule
\end{tabular}
\par\endgroup
% course-panel:T46-P01:end
% mdtable:T46:end
\end{minipage}\par
\reserve{7\baselineskip}
% mdparagraph:P166:begin
\subsection*{gemiltsian: present subjunctive}\BookHeading{nav-ch18-learner-0043}{subsubsection}{gemiltsian: present subjunctive}{0}
% mdparagraph:P166:end

% teaching: u18-gemiltsian-present-subjunctive6
% mdparagraph:P167:begin
This extends the familiar weak class II verb to present subjunctive -ie/\allowbreak{}-ien. The person receiving mercy remains dative.
% mdparagraph:P167:end

\reserve{13\baselineskip}
% mdparagraph:P168:begin
\textbf{Diagnostic forms:} gemiltsie (present subjunctive sg); gemiltsien (present subjunctive pl).
% mdparagraph:P168:end

% table: u18-gemiltsian-present-subjunctive6-SUBJUNCTIVE-DISPLAY
% mdtable:T47:begin
% course-table-style: 2026-10
% course-panel:T47-P01:begin
\par\begingroup
\fontsize{10.5}{13.2}\selectfont
\setlength{\LTleft}{0pt}\setlength{\LTright}{\fill}
\noindent
\begin{tabular}{@{}>{\raggedright\arraybackslash}p{63.342mm}>{\raggedright\arraybackslash}p{57.846mm}@{}}
\toprule
\textbf{Person and number} & \textbf{Present} \\
\midrule
1 sg. & gemiltsie \\
2 sg. & gemiltsie \\
3 sg. & gemiltsie \\
1–3 pl. & gemiltsien \\
\bottomrule
\end{tabular}
\par\endgroup
% course-panel:T47-P01:end
% mdtable:T47:end
\reserve{7\baselineskip}
% mdparagraph:P169:begin
\subsection*{gesēon: present subjunctive}\BookHeading{nav-ch18-learner-0044}{subsubsection}{gesēon: present subjunctive}{0}
% mdparagraph:P169:end

% teaching: u18-geseon-present-subjunctive6
% mdparagraph:P170:begin
The contracted present subjunctive is gesēo in the singular and gesēon in the plural. Only this mood and tense are added here.
% mdparagraph:P170:end

\reserve{13\baselineskip}
% mdparagraph:P171:begin
\textbf{Diagnostic forms:} gesēo (present subjunctive sg); gesēon (present subjunctive pl).
% mdparagraph:P171:end

% table: u18-geseon-present-subjunctive6-SUBJUNCTIVE-DISPLAY
% mdtable:T48:begin
% course-table-style: 2026-10
% course-panel:T48-P01:begin
\par\begingroup
\fontsize{10.5}{13.2}\selectfont
\setlength{\LTleft}{0pt}\setlength{\LTright}{\fill}
\noindent
\begin{tabular}{@{}>{\raggedright\arraybackslash}p{63.342mm}>{\raggedright\arraybackslash}p{57.846mm}@{}}
\toprule
\textbf{Person and number} & \textbf{Present} \\
\midrule
1 sg. & gesēo \\
2 sg. & gesēo \\
3 sg. & gesēo \\
1–3 pl. & gesēon \\
\bottomrule
\end{tabular}
\par\endgroup
% course-panel:T48-P01:end
% mdtable:T48:end
\reserve{7\baselineskip}
% mdparagraph:P172:begin
\subsection*{wunian: present subjunctive}\BookHeading{nav-ch18-learner-0045}{subsubsection}{wunian: present subjunctive}{0}
% mdparagraph:P172:end

% teaching: u18-wunian-present-subjunctive6
% mdparagraph:P173:begin
The familiar weak class II verb adds present subjunctive -ie/\allowbreak{}-ien. Dwelling in a place uses on with dative.
% mdparagraph:P173:end

\reserve{13\baselineskip}
% mdparagraph:P174:begin
\textbf{Diagnostic forms:} wunie (present subjunctive sg); wunien (present subjunctive pl).
% mdparagraph:P174:end

% table: u18-wunian-present-subjunctive6-SUBJUNCTIVE-DISPLAY
% mdtable:T49:begin
% course-table-style: 2026-10
% course-panel:T49-P01:begin
\par\begingroup
\fontsize{10.5}{13.2}\selectfont
\setlength{\LTleft}{0pt}\setlength{\LTright}{\fill}
\noindent
\begin{tabular}{@{}>{\raggedright\arraybackslash}p{63.342mm}>{\raggedright\arraybackslash}p{57.846mm}@{}}
\toprule
\textbf{Person and number} & \textbf{Present} \\
\midrule
1 sg. & wunie \\
2 sg. & wunie \\
3 sg. & wunie \\
1–3 pl. & wunien \\
\bottomrule
\end{tabular}
\par\endgroup
% course-panel:T49-P01:end
% mdtable:T49:end
\reserve{7\baselineskip}
% mdparagraph:P175:begin
\subsection*{bringan: present and preterite subjunctive}\BookHeading{nav-ch18-learner-0046}{subsubsection}{bringan: present and preterite subjunctive}{0}
% mdparagraph:P175:end

% teaching: u18-bringan-subjunctive12
% mdparagraph:P176:begin
Present subjunctive is bringe/\allowbreak{}bringen; preterite subjunctive brōhte/\allowbreak{}brōhten. The irregular dental past changes g to h before the dental group. Each tense has one form for all three persons within a number.
% mdparagraph:P176:end

\reserve{14\baselineskip}
% mdparagraph:P177:begin
\textbf{Diagnostic forms:} bringe (present subjunctive sg); bringen (present subjunctive pl); brōhte (preterite subjunctive sg); brōhten (preterite subjunctive pl).
% mdparagraph:P177:end

% table: u18-bringan-subjunctive12-SUBJUNCTIVE-DISPLAY
% mdtable:T50:begin
% course-table-style: 2026-10
% course-panel:T50-P01:begin
\par\begingroup
\fontsize{10.5}{13.2}\selectfont
\setlength{\LTleft}{0pt}\setlength{\LTright}{\fill}
\noindent
\begin{tabular}{@{}>{\raggedright\arraybackslash}p{39.459mm}>{\raggedright\arraybackslash}p{39.459mm}>{\raggedright\arraybackslash}p{39.459mm}@{}}
\toprule
\textbf{Person and number} & \textbf{Present} & \textbf{Preterite} \\
\midrule
1 sg. & bringe & brōhte \\
2 sg. & bringe & brōhte \\
3 sg. & bringe & brōhte \\
1–3 pl. & bringen & brōhten \\
\bottomrule
\end{tabular}
\par\endgroup
% course-panel:T50-P01:end
% mdtable:T50:end
\reserve{7\baselineskip}
% mdparagraph:P178:begin
\subsection*{lǣdan: preterite subjunctive}\BookHeading{nav-ch18-learner-0047}{subsubsection}{lǣdan: preterite subjunctive}{0}
% mdparagraph:P178:end

% teaching: u18-laedan-preterite-subjunctive6
% mdparagraph:P179:begin
The weak dental past has dd: lǣdde/\allowbreak{}lǣdden. These are the preterite subjunctive singular and plural in all persons, used under a past report of desire.
% mdparagraph:P179:end

\reserve{13\baselineskip}
% mdparagraph:P180:begin
\textbf{Diagnostic forms:} lǣdde (preterite subjunctive sg); lǣdden (preterite subjunctive pl).
% mdparagraph:P180:end

% table: u18-laedan-preterite-subjunctive6-SUBJUNCTIVE-DISPLAY
% mdtable:T51:begin
% course-table-style: 2026-10
% course-panel:T51-P01:begin
\par\begingroup
\fontsize{10.5}{13.2}\selectfont
\setlength{\LTleft}{0pt}\setlength{\LTright}{\fill}
\noindent
\begin{tabular}{@{}>{\raggedright\arraybackslash}p{63.342mm}>{\raggedright\arraybackslash}p{57.846mm}@{}}
\toprule
\textbf{Person and number} & \textbf{Preterite} \\
\midrule
1 sg. & lǣdde \\
2 sg. & lǣdde \\
3 sg. & lǣdde \\
1–3 pl. & lǣdden \\
\bottomrule
\end{tabular}
\par\endgroup
% course-panel:T51-P01:end
% mdtable:T51:end
\reserve{7\baselineskip}
% mdparagraph:P181:begin
\subsection*{stelan: preterite subjunctive}\BookHeading{nav-ch18-learner-0048}{subsubsection}{stelan: preterite subjunctive}{0}
% mdparagraph:P181:end

% teaching: NEW-U18-VOC-STELAN-U18-SELECTED
\par\begin{minipage}{\linewidth}\raggedright
% mdparagraph:P182:begin
\textbf{stelan} — steal.
% mdparagraph:P182:end

% teaching: u18-stelan-preterite-subjunctive6
% mdparagraph:P183:begin
The preterite subjunctive uses the long-vowel plural stem stǣl- in both numbers: stǣle/\allowbreak{}stǣlen. Its use in the accusation concerns earlier theft.
% mdparagraph:P183:end

% mdparagraph:P184:begin
\textbf{Diagnostic forms:} stǣle (preterite subjunctive sg); stǣlen (preterite subjunctive pl).
% mdparagraph:P184:end

% table: u18-stelan-preterite-subjunctive6-SUBJUNCTIVE-DISPLAY
% mdtable:T52:begin
% course-table-style: 2026-10
% course-panel:T52-P01:begin
\par\begingroup
\fontsize{10.5}{13.2}\selectfont
\setlength{\LTleft}{0pt}\setlength{\LTright}{\fill}
\noindent
\begin{tabular}{@{}>{\raggedright\arraybackslash}p{63.342mm}>{\raggedright\arraybackslash}p{57.846mm}@{}}
\toprule
\textbf{Person and number} & \textbf{Preterite} \\
\midrule
1 sg. & stǣle \\
2 sg. & stǣle \\
3 sg. & stǣle \\
1–3 pl. & stǣlen \\
\bottomrule
\end{tabular}
\par\endgroup
% course-panel:T52-P01:end
% mdtable:T52:end
\end{minipage}\par
\reserve{7\baselineskip}
% mdparagraph:P185:begin
\subsection*{gearcian: singular and plural commands}\BookHeading{nav-ch18-learner-0049}{subsubsection}{gearcian: singular and plural commands}{0}
% mdparagraph:P185:end

% teaching: u18-gearcian-imperative2
% mdparagraph:P186:begin
Weak class II commands have singular -a and plural -iað; these extend the earlier indicative system.
% mdparagraph:P186:end

\reserve{9\baselineskip}
% mdparagraph:P187:begin
\textbf{Diagnostic forms:} gearca (imperative sg); gearciað (imperative pl).
% mdparagraph:P187:end

% table: u18-gearcian-imperative2-IMPERATIVE-DISPLAY
% mdtable:T53:begin
% course-table-style: 2026-10
% course-panel:T53-P01:begin
\par\begingroup
\fontsize{10.5}{13.2}\selectfont
\setlength{\LTleft}{0pt}\setlength{\LTright}{\fill}
\noindent
\begin{tabular}{@{}>{\raggedright\arraybackslash}p{63.342mm}>{\raggedright\arraybackslash}p{57.846mm}@{}}
\toprule
\textbf{Person and number} & \textbf{Command} \\
\midrule
2 sg. & gearca \\
2 pl. & gearciað \\
\bottomrule
\end{tabular}
\par\endgroup
% course-panel:T53-P01:end
% mdtable:T53:end
\reserve{7\baselineskip}
% mdparagraph:P188:begin
\subsection*{āscian: singular and plural commands}\BookHeading{nav-ch18-learner-0050}{subsubsection}{āscian: singular and plural commands}{0}
% mdparagraph:P188:end

% teaching: u18-ascian-imperative2
% mdparagraph:P189:begin
Weak class II commands have singular -a and plural -iað. Retain initial ā.
% mdparagraph:P189:end

\reserve{9\baselineskip}
% mdparagraph:P190:begin
\textbf{Diagnostic forms:} āsca (imperative sg); āsciað (imperative pl).
% mdparagraph:P190:end

% table: u18-ascian-imperative2-IMPERATIVE-DISPLAY
% mdtable:T54:begin
% course-table-style: 2026-10
% course-panel:T54-P01:begin
\par\begingroup
\fontsize{10.5}{13.2}\selectfont
\setlength{\LTleft}{0pt}\setlength{\LTright}{\fill}
\noindent
\begin{tabular}{@{}>{\raggedright\arraybackslash}p{63.342mm}>{\raggedright\arraybackslash}p{57.846mm}@{}}
\toprule
\textbf{Person and number} & \textbf{Command} \\
\midrule
2 sg. & āsca \\
2 pl. & āsciað \\
\bottomrule
\end{tabular}
\par\endgroup
% course-panel:T54-P01:end
% mdtable:T54:end
\reserve{7\baselineskip}
% mdparagraph:P191:begin
\subsection*{diernan: plain and inflected infinitives}\BookHeading{nav-ch18-learner-0051}{subsubsection}{diernan: plain and inflected infinitives}{0}
% mdparagraph:P191:end

% teaching: NEW-U18-VOC-DIERNAN-U18-SELECTED
\par\begin{minipage}{\linewidth}\raggedright
% mdparagraph:P192:begin
\textbf{diernan} — conceal; hide.
% mdparagraph:P192:end

% teaching: u18-diernan-infinitive2
% mdparagraph:P193:begin
The pair gives plain and inflected infinitives. Select -anne here; -enne is a documented alternative. A modal takes the plain form; the inflected form still requires a licensed construction.
% mdparagraph:P193:end

% mdparagraph:P194:begin
\textbf{Diagnostic forms:} diernan (plain\_infinitive); diernanne (inflected\_infinitive).
% mdparagraph:P194:end

% table: u18-diernan-infinitive2-BOUNDED-NONFINITE-DISPLAY
% mdtable:T55:begin
\par\begingroup
\fontsize{10.5}{13.2}\selectfont
\setlength{\LTleft}{0pt}\setlength{\LTright}{\fill}
\noindent
\begin{tabular}{@{}>{\raggedright\arraybackslash}p{71.72mm}>{\raggedright\arraybackslash}p{49.46mm}@{}}
\toprule
\textbf{Type} & \textbf{Form} \\
\midrule
Plain infinitive & diernan \\
Inflected infinitive & diernanne \\
\bottomrule
\end{tabular}
\par\endgroup

% mdtable:T55:end
\end{minipage}\par
\reserve{7\baselineskip}
% mdparagraph:P195:begin
\subsection*{ęrian: plain and inflected infinitives}\BookHeading{nav-ch18-learner-0052}{subsubsection}{ęrian: plain and inflected infinitives}{0}
% mdparagraph:P195:end

% teaching: NEW-U18-VOC-ERIAN-U18-SELECTED
\par\begin{minipage}{\linewidth}\raggedright
% mdparagraph:P196:begin
\textbf{ęrian} — plow.
% mdparagraph:P196:end

% teaching: u18-erian-infinitive2
% mdparagraph:P197:begin
The pair gives plain and inflected infinitives, with selected -ianne. The modal construction takes the plain form. Retain printed ę.
% mdparagraph:P197:end

% mdparagraph:P198:begin
\textbf{Diagnostic forms:} ęrian (plain\_infinitive); ęrianne (inflected\_infinitive).
% mdparagraph:P198:end

% table: u18-erian-infinitive2-BOUNDED-NONFINITE-DISPLAY
% mdtable:T56:begin
\par\begingroup
\fontsize{10.5}{13.2}\selectfont
\setlength{\LTleft}{0pt}\setlength{\LTright}{\fill}
\noindent
\begin{tabular}{@{}>{\raggedright\arraybackslash}p{71.72mm}>{\raggedright\arraybackslash}p{49.46mm}@{}}
\toprule
\textbf{Type} & \textbf{Form} \\
\midrule
Plain infinitive & ęrian \\
Inflected infinitive & ęrianne \\
\bottomrule
\end{tabular}
\par\endgroup

% mdtable:T56:end
\end{minipage}\par
\reserve{7\baselineskip}
% mdparagraph:P199:begin
\subsection*{rīpan: plain and inflected infinitives}\BookHeading{nav-ch18-learner-0053}{subsubsection}{rīpan: plain and inflected infinitives}{0}
% mdparagraph:P199:end

% teaching: NEW-U18-VOC-RIPAN-U18-SELECTED
\par\begin{minipage}{\linewidth}\raggedright
% mdparagraph:P200:begin
\textbf{rīpan} — reap.
% mdparagraph:P200:end

% teaching: u18-ripan-infinitive2
% mdparagraph:P201:begin
The pair gives plain and inflected infinitives, with selected -anne. The modal construction takes the plain form. Retain ī.
% mdparagraph:P201:end

% mdparagraph:P202:begin
\textbf{Diagnostic forms:} rīpan (plain\_infinitive); rīpanne (inflected\_infinitive).
% mdparagraph:P202:end

% table: u18-ripan-infinitive2-BOUNDED-NONFINITE-DISPLAY
% mdtable:T57:begin
\par\begingroup
\fontsize{10.5}{13.2}\selectfont
\setlength{\LTleft}{0pt}\setlength{\LTright}{\fill}
\noindent
\begin{tabular}{@{}>{\raggedright\arraybackslash}p{71.72mm}>{\raggedright\arraybackslash}p{49.46mm}@{}}
\toprule
\textbf{Type} & \textbf{Form} \\
\midrule
Plain infinitive & rīpan \\
Inflected infinitive & rīpanne \\
\bottomrule
\end{tabular}
\par\endgroup

% mdtable:T57:end
\end{minipage}\par
\reserve{7\baselineskip}
% mdparagraph:P203:begin
\subsection*{giernan: plain and inflected infinitives}\BookHeading{nav-ch18-learner-0054}{subsubsection}{giernan: plain and inflected infinitives}{0}
% mdparagraph:P203:end

% teaching: u18-giernan-infinitive2
% mdparagraph:P204:begin
The pair extends the earlier finite system to plain and inflected infinitives, with selected -anne. The genitive and desired-content complements are explained separately.
% mdparagraph:P204:end

\reserve{9\baselineskip}
% mdparagraph:P205:begin
\textbf{Diagnostic forms:} giernan (plain\_infinitive); giernanne (inflected\_infinitive).
% mdparagraph:P205:end

% table: u18-giernan-infinitive2-BOUNDED-NONFINITE-DISPLAY
% mdtable:T58:begin
\par\begingroup
\fontsize{10.5}{13.2}\selectfont
\setlength{\LTleft}{0pt}\setlength{\LTright}{\fill}
\noindent
\begin{tabular}{@{}>{\raggedright\arraybackslash}p{71.72mm}>{\raggedright\arraybackslash}p{49.46mm}@{}}
\toprule
\textbf{Type} & \textbf{Form} \\
\midrule
Plain infinitive & giernan \\
Inflected infinitive & giernanne \\
\bottomrule
\end{tabular}
\par\endgroup

% mdtable:T58:end
\reserve{7\baselineskip}
% mdparagraph:P206:begin
\subsection*{āfǣran: present and past base participles}\BookHeading{nav-ch18-learner-0055}{subsubsection}{āfǣran: present and past base participles}{0}
% mdparagraph:P206:end

% teaching: NEW-U18-VOC-AFAERAN-U18-SELECTED
\par\begin{minipage}{\linewidth}\raggedright
% mdparagraph:P207:begin
\textbf{āfǣran} — frighten.
% mdparagraph:P207:end

% teaching: u18-afaeran-participle2
% mdparagraph:P208:begin
The long weak stem gives present āfǣrende and uncontracted past āfǣred. The base past participle can remain uninflected in the supplied predicate with plural subject.
% mdparagraph:P208:end

% mdparagraph:P209:begin
\textbf{Diagnostic forms:} āfǣrende (present participle); āfǣred (past participle).
% mdparagraph:P209:end

% table: u18-afaeran-participle2-PARTICIPLES-DISPLAY
% mdtable:T59:begin
\par\begingroup
\fontsize{10.5}{13.2}\selectfont
\setlength{\LTleft}{0pt}\setlength{\LTright}{\fill}
\noindent
\begin{tabular}{@{}>{\raggedright\arraybackslash}p{60.59mm}>{\raggedright\arraybackslash}p{60.59mm}@{}}
\toprule
\textbf{Type} & \textbf{Form} \\
\midrule
present & āfǣrende \\
past & āfǣred \\
\bottomrule
\end{tabular}
\par\endgroup

% mdtable:T59:end
\end{minipage}\par
\reserve{7\baselineskip}
% mdparagraph:P210:begin
\subsection*{bēodan: present and past base participles}\BookHeading{nav-ch18-learner-0056}{subsubsection}{bēodan: present and past base participles}{0}
% mdparagraph:P210:end

% teaching: NEW-U18-VOC-BEODAN-U18-COMMAND
\par\begin{minipage}{\linewidth}\raggedright
% mdparagraph:P211:begin
\textbf{bēodan} — command; bid.
% mdparagraph:P211:end

% teaching: u18-beodan-participle2
% mdparagraph:P212:begin
Strong class II gives present bēodende and past boden. The selected new sense is command, distinct from offering. The impersonal command passive remains supplied as a whole.
% mdparagraph:P212:end

% mdparagraph:P213:begin
\textbf{Diagnostic forms:} bēodende (present participle); boden (past participle).
% mdparagraph:P213:end

% table: u18-beodan-participle2-PARTICIPLES-DISPLAY
% mdtable:T60:begin
\par\begingroup
\fontsize{10.5}{13.2}\selectfont
\setlength{\LTleft}{0pt}\setlength{\LTright}{\fill}
\noindent
\begin{tabular}{@{}>{\raggedright\arraybackslash}p{60.59mm}>{\raggedright\arraybackslash}p{60.59mm}@{}}
\toprule
\textbf{Type} & \textbf{Form} \\
\midrule
present & bēodende \\
past & boden \\
\bottomrule
\end{tabular}
\par\endgroup

% mdtable:T60:end
\end{minipage}\par
\reserve{7\baselineskip}
% mdparagraph:P214:begin
\subsection*{berēafian: present and past base participles}\BookHeading{nav-ch18-learner-0057}{subsubsection}{berēafian: present and past base participles}{0}
% mdparagraph:P214:end

% teaching: u18-bereafian-participle2
% mdparagraph:P215:begin
Weak class II gives present -iende and past -od. Retain ēa. These are base participles; the deprivation complement is taught separately.
% mdparagraph:P215:end

\reserve{9\baselineskip}
% mdparagraph:P216:begin
\textbf{Diagnostic forms:} berēafiende (present participle); berēafod (past participle).
% mdparagraph:P216:end

% table: u18-bereafian-participle2-PARTICIPLES-DISPLAY
% mdtable:T61:begin
\par\begingroup
\fontsize{10.5}{13.2}\selectfont
\setlength{\LTleft}{0pt}\setlength{\LTright}{\fill}
\noindent
\begin{tabular}{@{}>{\raggedright\arraybackslash}p{57.71mm}>{\raggedright\arraybackslash}p{63.48mm}@{}}
\toprule
\textbf{Type} & \textbf{Form} \\
\midrule
present & berēafiende \\
past & berēafod \\
\bottomrule
\end{tabular}
\par\endgroup

% mdtable:T61:end
\reserve{7\baselineskip}
% mdparagraph:P217:begin
\subsection*{ofslēan: present and past base participles}\BookHeading{nav-ch18-learner-0058}{subsubsection}{ofslēan: present and past base participles}{0}
% mdparagraph:P217:end

% teaching: NEW-U18-VOC-OFSLEAN-U18-SELECTED
\par\begin{minipage}{\linewidth}\raggedright
% mdparagraph:P218:begin
\textbf{ofslēan} — slay; kill.
% mdparagraph:P218:end

% teaching: u18-ofslean-participle2
% mdparagraph:P219:begin
The contracted verb gives present ofslēande and past ofslægen; ofslagen is a documented alternative. Retain of-. These participles do not introduce a finite system.
% mdparagraph:P219:end

% mdparagraph:P220:begin
\textbf{Diagnostic forms:} ofslēande (present participle); ofslægen (past participle).
% mdparagraph:P220:end

% table: u18-ofslean-participle2-PARTICIPLES-DISPLAY
% mdtable:T62:begin
\par\begingroup
\fontsize{10.5}{13.2}\selectfont
\setlength{\LTleft}{0pt}\setlength{\LTright}{\fill}
\noindent
\begin{tabular}{@{}>{\raggedright\arraybackslash}p{60.59mm}>{\raggedright\arraybackslash}p{60.59mm}@{}}
\toprule
\textbf{Type} & \textbf{Form} \\
\midrule
present & ofslēande \\
past & ofslægen \\
\bottomrule
\end{tabular}
\par\endgroup

% mdtable:T62:end
\end{minipage}\par
\reserve{7\baselineskip}
% mdparagraph:P221:begin
\subsection*{ofðyncean: impersonal third-singular indicative/\allowbreak{}subjunctive, present and preterite}\BookHeading{nav-ch18-learner-0059}{subsubsection}{ofðyncean: impersonal third-singular indicative/subjunctive, present and preterite}{0}
% mdparagraph:P221:end

% teaching: NEW-U18-VOC-OFTHYNCEAN-U18-SELECTED
% mdparagraph:P222:begin
\textbf{ofðyncean} — grieve; displease (impersonal).
% mdparagraph:P222:end

% teaching: u18-ofthyncean-impersonal4
% mdparagraph:P223:begin
This\BookPoint{bk-ch18-learner-8cf9d181123a} complete impersonal system is limited to third singular: present indicative ofðyncð, present subjunctive ofðynce, and past ofðūhte in both moods. The experiencer is dative and does not trigger personal or plural agreement. A cause/\allowbreak{}content clause gives what grieves that person.
% mdparagraph:P223:end

% mdparagraph:P224:begin
\textbf{Diagnostic forms:} ofðyncð (present indicative impersonal third singular); ofðynce (present subjunctive sg); ofðūhte (preterite indicative impersonal third singular); ofðūhte (preterite subjunctive sg).
% mdparagraph:P224:end

\reserve{8\baselineskip}
\reserve{6\baselineskip}
% mdparagraph:P225:begin
\textbf{ofðyncean — indicative}
% mdparagraph:P225:end

% table: u18-ofthyncean-impersonal4-INDICATIVE-DISPLAY
% mdtable:T63:begin
% course-table-style: 2026-10
% course-panel:T63-P01:begin
\par\begingroup
\fontsize{10.5}{13.2}\selectfont
\setlength{\LTleft}{0pt}\setlength{\LTright}{\fill}
\noindent
\begin{tabular}{@{}>{\raggedright\arraybackslash}p{39.459mm}>{\raggedright\arraybackslash}p{39.459mm}>{\raggedright\arraybackslash}p{39.459mm}@{}}
\toprule
\textbf{Person and number} & \textbf{Present} & \textbf{Preterite} \\
\midrule
3 sg. & ofðyncð & ofðūhte \\
\bottomrule
\end{tabular}
\par\endgroup
% course-panel:T63-P01:end
% mdtable:T63:end
\reserve{8\baselineskip}
\reserve{6\baselineskip}
% mdparagraph:P226:begin
\textbf{ofðyncean — subjunctive}
% mdparagraph:P226:end

% table: u18-ofthyncean-impersonal4-SUBJUNCTIVE-DISPLAY
% mdtable:T64:begin
% course-table-style: 2026-10
% course-panel:T64-P01:begin
\par\begingroup
\fontsize{10.5}{13.2}\selectfont
\setlength{\LTleft}{0pt}\setlength{\LTright}{\fill}
\noindent
\begin{tabular}{@{}>{\raggedright\arraybackslash}p{39.459mm}>{\raggedright\arraybackslash}p{39.459mm}>{\raggedright\arraybackslash}p{39.459mm}@{}}
\toprule
\textbf{Person and number} & \textbf{Present} & \textbf{Preterite} \\
\midrule
3 sg. & ofðynce & ofðūhte \\
\bottomrule
\end{tabular}
\par\endgroup
% course-panel:T64-P01:end
% mdtable:T64:end
\reserve{7\baselineskip}
% mdparagraph:P227:begin
\section*{4. Meanings and complements}\BookHeading{nav-ch18-learner-0060}{subsection}{4. Meanings and complements}{1}
% mdparagraph:P227:end

% teaching: NEW-U18-VOC-GLAED-U18-GRACIOUS
% mdparagraph:P228:begin
\textbf{glæd} — gracious; favourably disposed toward someone with wið.
% mdparagraph:P228:end

\reserve{7\baselineskip}
% mdparagraph:P229:begin
\subsection*{Need a thing; relative complement}\BookHeading{nav-ch18-learner-0061}{subsubsection}{Need a thing; relative complement}{0}
% mdparagraph:P229:end

% teaching: NEW-U18-GRAM-U18-BETHURFAN-GENITIVE-RELATIVE
% mdparagraph:P230:begin
Beðurfan\BookPoint{bk-ch18-learner-49add5c35eec} takes a genitive of the thing needed in the selected construction: hwǣtes, of wheat. In ðā ðing ðe wē beðurfon, the things are accusative outside the relative because they are bought. Inside it, invariant ðe represents the genitive complement required by beðurfan. Its unchanged spelling does not transfer the outer accusative into the relative. Other government is attested; the genitive is the target here.
% mdparagraph:P230:end

\reserve{7\baselineskip}
% mdparagraph:P231:begin
\subsection*{Recognize an accusative person}\BookHeading{nav-ch18-learner-0062}{subsubsection}{Recognize an accusative person}{0}
% mdparagraph:P231:end

% teaching: NEW-U18-GRAM-U18-ONCNAWAN-OBJECT
% mdparagraph:P232:begin
Oncnāwan in the selected recognition sense takes an accusative recognized person. The nominative recognizer controls the verb; a change of recognizer from ic to wē changes person/\allowbreak{}number independently of the object. Present cnāw-/\allowbreak{}cnǣw- and past cnēow- stay distinct.
% mdparagraph:P232:end

\reserve{7\baselineskip}
% mdparagraph:P233:begin
\subsection*{Open an accusative container}\BookHeading{nav-ch18-learner-0063}{subsubsection}{Open an accusative container}{0}
% mdparagraph:P233:end

% teaching: NEW-U18-GRAM-U18-UNDON-OBJECT
% mdparagraph:P234:begin
Undōn means undo/\allowbreak{}open the tied sacks in 45.8 and takes the accusative container. Present dēð and preterite dyde are different stems; plural preterite undydon follows the plural opener, not the number of sacks.
% mdparagraph:P234:end

\reserve{7\baselineskip}
% mdparagraph:P235:begin
\subsection*{Swear oaths and report their content}\BookHeading{nav-ch18-learner-0064}{subsubsection}{Swear oaths and report their content}{0}
% mdparagraph:P235:end

% teaching: NEW-U18-GRAM-U18-OATH-RECIPIENT-CONTENT
% mdparagraph:P236:begin
Swęrian/\allowbreak{}swerian takes accusative āðas as the sworn utterance. A dative recipient mīnum fæder 48.1 needs no separate tō; ætforan ūs 44.2 names the people in whose presence it was sworn. The reported content is a ðæt clause: mōsten gives past reported permission, brōhte a promised bringing viewed prospectively from the former oath. Keep the recipient distinct from those present at the oath.
% mdparagraph:P236:end

\reserve{7\baselineskip}
% mdparagraph:P237:begin
\subsection*{Desire that a state of affairs occur}\BookHeading{nav-ch18-learner-0065}{subsubsection}{Desire that a state of affairs occur}{0}
% mdparagraph:P237:end

% teaching: NEW-U18-GRAM-U18-GIERNAN-PROPOSITIONAL-DESIRE
% mdparagraph:P238:begin
After\BookPoint{bk-ch18-learner-0bf6031e8afb} giernan, genitive ðæs points forward to a ðæt clause explaining the desired action. The demonstrative is inflected; the following ðæt is a conjunction. In the past report, lǣdden is first-plural preterite subjunctive, describing the action then desired.
% mdparagraph:P238:end

\reserve{7\baselineskip}
% mdparagraph:P239:begin
\subsection*{Make a person be gracious}\BookHeading{nav-ch18-learner-0066}{subsubsection}{Make a person be gracious}{0}
% mdparagraph:P239:end

% teaching: NEW-U18-GRAM-U18-GEDON-OBJECT-PREDICATE
% mdparagraph:P240:begin
Gedōn has accusative person hine plus agreeing accusative masculine singular glædne. In 45.2 gedō is third-singular present subjunctive expressing the wish; it is not an imperative addressed to the named Lord. Wið ēow supplies the graciousness target, as the original hint expressly requires.
% mdparagraph:P240:end

\reserve{7\baselineskip}
% mdparagraph:P241:begin
\subsection*{Give back a person}\BookHeading{nav-ch18-learner-0067}{subsubsection}{Give back a person}{0}
% mdparagraph:P241:end

% teaching: NEW-U18-GRAM-U18-AGIEFAN-RESTORATION
% mdparagraph:P242:begin
Āgiefan selects the returned person as accusative and the recipient as dative. Ēow and ēowerne brōðor have different roles. Present subjunctive āgiefe is in the final purpose clause, not merely an unexplained alternative ending.
% mdparagraph:P242:end

\reserve{7\baselineskip}
% mdparagraph:P243:begin
\subsection*{Despoil someone of possessions}\BookHeading{nav-ch18-learner-0068}{subsubsection}{Despoil someone of possessions}{0}
% mdparagraph:P243:end

% teaching: NEW-U18-GRAM-U18-BEREAFIAN-DEPRIVATION
% mdparagraph:P244:begin
Berēafian\BookPoint{bk-ch18-learner-83e42b478fb6} takes an accusative person and genitive possessions in the selected active construction. In the passive, the deprived person becomes nominative while the possessions remain genitive. In 45.6, assena and ðinga are genitive plurals; selfe agrees with accusative ūs in the following clause. Dative and instrumental alternatives are attested elsewhere, but the local target is genitive.
% mdparagraph:P244:end

\reserve{7\baselineskip}
% mdparagraph:P245:begin
\subsection*{Listen to speech}\BookHeading{nav-ch18-learner-0069}{subsubsection}{Listen to speech}{0}
% mdparagraph:P245:end

% teaching: NEW-U18-GRAM-U18-HLYSTAN-GENITIVE
% mdparagraph:P246:begin
Hlystan governs genitive speech/\allowbreak{}content in 45.7, as the explicit 156 hint directs. Thus ūrre sprǣce is genitive feminine singular, distinct from accusative ūre sprǣce. The people beseech the steward (ðē accusative); ðū is the subject of subordinate hlyste.
% mdparagraph:P246:end

\reserve{7\baselineskip}
% mdparagraph:P247:begin
\subsection*{Enslave an accusative person}\BookHeading{nav-ch18-learner-0070}{subsubsection}{Enslave an accusative person}{0}
% mdparagraph:P247:end

% teaching: NEW-U18-GRAM-U18-THEOWIAN-ENSLAVE
% mdparagraph:P248:begin
The selected ðēowian sense is enslave, with an accusative person. The more common serve sense with different government remains separate. The selected ðēowie normalizes the documented -ige/\allowbreak{}-ie variation under Cook 119 and does not change tense or mood.
% mdparagraph:P248:end

\reserve{7\baselineskip}
% mdparagraph:P249:begin
\subsection*{Command a dative person}\BookHeading{nav-ch18-learner-0071}{subsubsection}{Command a dative person}{0}
% mdparagraph:P249:end

% teaching: NEW-U18-GRAM-U18-BEODAN-COMMAND
% mdparagraph:P250:begin
Bēodan\BookPoint{bk-ch18-learner-a0509728a1ca} takes a dative addressee in 46.7 and 48.8, as 164 expressly directs. It can accompany a direct command or introduce ðæt plus subjunctive commanded action. 47.11 intentionally retains present bīetst (commandest); the lookalike bīetst of bēatan is a different verb here.
% mdparagraph:P250:end

\reserve{7\baselineskip}
% mdparagraph:P251:begin
\subsection*{Lay out things}\BookHeading{nav-ch18-learner-0072}{subsubsection}{Lay out things}{0}
% mdparagraph:P251:end

% teaching: NEW-U18-GRAM-U18-LECGEAN-PLACEMENT
% mdparagraph:P252:begin
Lęcgean takes accusative things laid, with forð marking their presentation in 46.3. Its source reading imperative lęcgeað is separately supplied; learning indicative plural lęcgeað does not by itself assign the command use.
% mdparagraph:P252:end

\reserve{7\baselineskip}
% mdparagraph:P253:begin
\subsection*{Mention a person earlier to someone}\BookHeading{nav-ch18-learner-0073}{subsubsection}{Mention a person earlier to someone}{0}
% mdparagraph:P253:end

% teaching: NEW-U18-GRAM-U18-FORESECGEAN-PAST-MENTION
% mdparagraph:P254:begin
Foresęcgean here means mention previously, not foretell the future. The invariant ðe is the mentioned person (accusative function), and mē the dative hearer. The English of is explicitly omitted.
% mdparagraph:P254:end

\reserve{7\baselineskip}
% mdparagraph:P255:begin
\subsection*{Be grieved by a stated circumstance}\BookHeading{nav-ch18-learner-0074}{subsubsection}{Be grieved by a stated circumstance}{0}
% mdparagraph:P255:end

% teaching: NEW-U18-GRAM-U18-OFTHYNCEAN-IMPERSONAL
% mdparagraph:P256:begin
Ofðyncean remains third singular with a dative experiencer and a cause/\allowbreak{}content ðæt clause. Present subjunctive ofðynce means let it not grieve in the negative wish. English it needs no separate Old English pronoun here. A plural experiencer does not make the verb plural.
% mdparagraph:P256:end

\reserve{7\baselineskip}
% mdparagraph:P257:begin
\subsection*{Fear with genitive object and reflexive dative}\BookHeading{nav-ch18-learner-0075}{subsubsection}{Fear with genitive object and reflexive dative}{0}
% mdparagraph:P257:end

% teaching: NEW-U18-GRAM-U18-ONDRAEDAN-GENITIVE
% mdparagraph:P258:begin
Ondrǣdan\BookPoint{bk-ch18-learner-a19c79a380b1} can take a reflexive dative and a genitive feared thing. In 48.6, ēow names the people themselves and nānes ðinges the feared entity. In 46.1 the reflexive dative occurs without an expressed feared thing. The unchanged spelling of ēow does not make it accusative. This construction supplements the earlier accusative frame.
% mdparagraph:P258:end

\reserve{7\baselineskip}
% mdparagraph:P259:begin
\subsection*{Requite good with evil}\BookHeading{nav-ch18-learner-0076}{subsubsection}{Requite good with evil}{0}
% mdparagraph:P259:end

% teaching: NEW-U18-GRAM-U18-FORGIELDAN-GOOD-EVIL
% mdparagraph:P260:begin
In 46.9 accusative neuter gōd is the good being repaid; mid yfele is the evil given in return. The recipient frame previously used for paying is not by itself this good-for-evil substitution frame. Wolden is third-plural past subjunctive inside the reported why-question; forgieldan is its plain infinitive.
% mdparagraph:P260:end

\reserve{7\baselineskip}
% mdparagraph:P261:begin
\subsection*{Prospective when with subjunctive}\BookHeading{nav-ch18-learner-0077}{subsubsection}{Prospective when with subjunctive}{0}
% mdparagraph:P261:end

% teaching: NEW-U18-GRAM-U18-PROSPECTIVE-WHEN
% mdparagraph:P262:begin
46.9\BookPoint{bk-ch18-learner-0a7d2c4cd048} when you have caught them concerns a future condition for the command to ask. The selected ðonne ðū hīe fō uses second-singular present subjunctive, following 196 b. There is no required present-perfect auxiliary. Original English says you, and the singular steward continues as addressee of Ārīs, gā and āsca.
% mdparagraph:P262:end

\reserve{7\baselineskip}
% mdparagraph:P263:begin
\subsection*{Subjunctive within an optative relative}\BookHeading{nav-ch18-learner-0078}{subsubsection}{Subjunctive within an optative relative}{0}
% mdparagraph:P263:end

% teaching: NEW-U18-GRAM-U18-RELATIVE-SUBJUNCTIVE-ATTRACTION
% mdparagraph:P264:begin
47.3\BookPoint{bk-ch18-learner-bbf20820beab} retains 195: in sīe hē ofslægen mid ðǣm ðe ðū hit finde, finde is second-singular present subjunctive within the relative governed by the hypothetical/\allowbreak{}optative context. Mid takes a dative associate; ðe is invariant. This is not a rule that every relative clause uses subjunctive.
% mdparagraph:P264:end

\reserve{7\baselineskip}
% mdparagraph:P265:begin
\subsection*{Narrative preterite with anterior meaning}\BookHeading{nav-ch18-learner-0079}{subsubsection}{Narrative preterite with anterior meaning}{0}
% mdparagraph:P265:end

% teaching: NEW-U18-GRAM-U18-PRETERITE-ANTERIOR-REVIEW
% mdparagraph:P266:begin
Cook\BookPoint{bk-ch18-learner-8fe39e530b47} 188 permits the simple preterite for English simple, perfect and pluperfect past. In 45.8 cēapodon and ǣr sealdon give the anterior actions; 46.9 fōron ... sumne dǣl ðæs weges conveys had gone without adding a habban perfect. Each clause remains present in the translation.
% mdparagraph:P266:end

\reserve{7\baselineskip}
% mdparagraph:P267:begin
\subsection*{Conceal oneself}\BookHeading{nav-ch18-learner-0080}{subsubsection}{Conceal oneself}{0}
% mdparagraph:P267:end

% teaching: NEW-U18-GRAM-U18-DIERNAN-REFLEXIVE
% mdparagraph:P268:begin
Diernan\BookPoint{bk-ch18-learner-dbac285fef20} takes the concealed person as accusative; hine in 48.4 refers back to nominative Iosep. Meahte is past ability and governs the plain infinitive. Longer is an adverb modifying concealment, not an adjective agreeing with Joseph.
% mdparagraph:P268:end

% mdparagraph:P269:begin
The earlier preterite singular funde of findan remains valid; the alternative fand does not invalidate it. Retain the selected spelling unless using an independently justified alternative.
% mdparagraph:P269:end

\reserve{7\baselineskip}
% mdparagraph:P270:begin
\section*{5. Cumulative analysis and controlled composition}\BookHeading{nav-ch18-learner-0081}{subsection}{5. Cumulative analysis and controlled composition}{1}
% mdparagraph:P270:end

% mdparagraph:P271:begin
Use the stated earlier or newly taught forms. Keep any explicitly supplied reference visible. Explain roles from the actual sentence; do not supply an unstated referent or event.
% mdparagraph:P271:end

\reserve{8\baselineskip}
% mdparagraph:P272:begin
\subsection*{A01}\BookHeading{nav-ch18-learner-0082}{subsubsection}{A01}{0}
% mdparagraph:P272:end

% item: NEW-U18-A01
% mdparagraph:P273:begin
Analyse the two sentences. No further context identifies the reader.
% mdparagraph:P273:end

% mdparagraph:P274:begin
Se cyning sealde ðǣm męnn ðā bōc. Hē hīe rǣdde.
% mdparagraph:P274:end

% output: NEW-U18-A01-O001
% mdparagraph:P275:begin
\itempara{1}{Identify the giver, recipient and thing given.}
% mdparagraph:P275:end

% output: NEW-U18-A01-O002
% mdparagraph:P276:begin
\itempara{2}{Give the referent and case of hīe.}
% mdparagraph:P276:end

% output: NEW-U18-A01-O003
% mdparagraph:P277:begin
\itempara{3}{Which named people can hē refer to? Explain what the forms establish.}
% mdparagraph:P277:end

% output: NEW-U18-A01-O004
% mdparagraph:P278:begin
\itempara{4}{Rewrite only the second sentence so that the recipient is explicitly the reader.}
% mdparagraph:P278:end

\reserve{8\baselineskip}
% mdparagraph:P279:begin
\subsection*{A02}\BookHeading{nav-ch18-learner-0083}{subsubsection}{A02}{0}
% mdparagraph:P279:end

% item: NEW-U18-A02
% mdparagraph:P280:begin
Convert this direct statement into a report. Retain the king as the person who saw the man; use sęcgan with a ðæt clause and a preterite subjunctive in that clause.
% mdparagraph:P280:end

% mdparagraph:P281:begin
Se cyning cwæð: “Ic geseah ðone mann.”
% mdparagraph:P281:end

% mdparagraph:P282:begin
\textbf{Reference:} sęcgan.
% mdparagraph:P282:end

% mdparagraph:P283:begin
\textbf{Reference:} ðæt.
% mdparagraph:P283:end

% output: NEW-U18-A02-O001
% mdparagraph:P284:begin
\itempara{1}{Identify the referent of ic and the time of the seeing relative to the quoted statement.}
% mdparagraph:P284:end

% output: NEW-U18-A02-O002
% mdparagraph:P285:begin
\itempara{2}{Write the complete reported statement.}
% mdparagraph:P285:end

% output: NEW-U18-A02-O003
% mdparagraph:P286:begin
\itempara{3}{Explain whether the chosen subjunctive proves that the king’s statement is false.}
% mdparagraph:P286:end

\reserve{8\baselineskip}
% mdparagraph:P287:begin
\subsection*{A03}\BookHeading{nav-ch18-learner-0084}{subsubsection}{A03}{0}
% mdparagraph:P287:end

% item: NEW-U18-A03
% mdparagraph:P288:begin
Analyse the sentence, then replace the men by the man and keep the king singular.
% mdparagraph:P288:end

% mdparagraph:P289:begin
Ðā męnn ðe ðone cyning gesāwon, sungon.
% mdparagraph:P289:end

% output: NEW-U18-A03-O001
% mdparagraph:P290:begin
\itempara{1}{Translate the complete sentence.}
% mdparagraph:P290:end

% output: NEW-U18-A03-O002
% mdparagraph:P291:begin
\itempara{2}{Identify the role represented by ðe and explain the agreement of both finite verbs.}
% mdparagraph:P291:end

% output: NEW-U18-A03-O003
% mdparagraph:P292:begin
\itempara{3}{Write the complete singular version.}
% mdparagraph:P292:end

\reserve{8\baselineskip}
% mdparagraph:P293:begin
\subsection*{C01}\BookHeading{nav-ch18-learner-0085}{subsubsection}{C01}{0}
% mdparagraph:P293:end

% item: NEW-U18-C01
% mdparagraph:P294:begin
Translate: The man whom the king saw gave the servant the book. Use ðe and sęllan.
% mdparagraph:P294:end

% mdparagraph:P295:begin
\textbf{Reference:} ðe.
% mdparagraph:P295:end

% mdparagraph:P296:begin
\textbf{Reference:} sęllan.
% mdparagraph:P296:end

% output: NEW-U18-C01-O001
% mdparagraph:P297:begin
\itempara{1}{Write one complete Old English sentence.}
% mdparagraph:P297:end

\reserve{8\baselineskip}
% mdparagraph:P298:begin
\subsection*{C02}\BookHeading{nav-ch18-learner-0086}{subsubsection}{C02}{0}
% mdparagraph:P298:end

% item: NEW-U18-C02
% mdparagraph:P299:begin
Translate: The king ordered the men to give the servant the books. Use hātan and a plain infinitive of sęllan.
% mdparagraph:P299:end

% mdparagraph:P300:begin
\textbf{Reference:} hātan.
% mdparagraph:P300:end

% mdparagraph:P301:begin
\textbf{Reference:} sęllan.
% mdparagraph:P301:end

% output: NEW-U18-C02-O001
% mdparagraph:P302:begin
\itempara{1}{Write one complete Old English sentence.}
% mdparagraph:P302:end

\reserve{8\baselineskip}
% mdparagraph:P303:begin
\subsection*{F01}\BookHeading{nav-ch18-learner-0087}{subsubsection}{F01}{0}
% mdparagraph:P303:end

% item: NEW-U18-F01
% mdparagraph:P304:begin
Translate each sentence, then explain the selected ending contrast. Use the newly taught words for ass, sack and wheat.
% mdparagraph:P304:end

% output: NEW-U18-F01-O001
% mdparagraph:P305:begin
\itempara{1}{The ass has the sack.}
% mdparagraph:P305:end

% output: NEW-U18-F01-O002
% mdparagraph:P306:begin
\itempara{2}{The asses have the sacks.}
% mdparagraph:P306:end

% output: NEW-U18-F01-O003
% mdparagraph:P307:begin
\itempara{3}{The men give wheat to the asses.}
% mdparagraph:P307:end

% output: NEW-U18-F01-O004
% mdparagraph:P308:begin
\itempara{4}{The wheat is in the sacks.}
% mdparagraph:P308:end

% output: NEW-U18-F01-O005
% mdparagraph:P309:begin
\itempara{5}{Compare the plural noun endings in the first three translations and the location phrase.}
% mdparagraph:P309:end

\reserve{8\baselineskip}
% mdparagraph:P310:begin
\subsection*{F02}\BookHeading{nav-ch18-learner-0088}{subsubsection}{F02}{0}
% mdparagraph:P310:end

% item: NEW-U18-F02
% mdparagraph:P311:begin
Translate the two sentences using the newly taught verb meaning need. Then explain the governed noun.
% mdparagraph:P311:end

% output: NEW-U18-F02-O001
% mdparagraph:P312:begin
\itempara{1}{The man needs wheat.}
% mdparagraph:P312:end

% output: NEW-U18-F02-O002
% mdparagraph:P313:begin
\itempara{2}{The men needed wheat.}
% mdparagraph:P313:end

% output: NEW-U18-F02-O003
% mdparagraph:P314:begin
\itempara{3}{Identify the tense and number of each finite verb, and explain the case of wheat in both sentences.}
% mdparagraph:P314:end

\reserve{8\baselineskip}
% mdparagraph:P315:begin
\subsection*{F03}\BookHeading{nav-ch18-learner-0089}{subsubsection}{F03}{0}
% mdparagraph:P315:end

% item: NEW-U18-F03
% mdparagraph:P316:begin
Translate the four sentences using the newly taught verbs meaning recognize and undo/\allowbreak{}open. Then explain the roles and number changes.
% mdparagraph:P316:end

% output: NEW-U18-F03-O001
% mdparagraph:P317:begin
\itempara{1}{I recognize the king.}
% mdparagraph:P317:end

% output: NEW-U18-F03-O002
% mdparagraph:P318:begin
\itempara{2}{We recognized the king.}
% mdparagraph:P318:end

% output: NEW-U18-F03-O003
% mdparagraph:P319:begin
\itempara{3}{The man opens the sack.}
% mdparagraph:P319:end

% output: NEW-U18-F03-O004
% mdparagraph:P320:begin
\itempara{4}{The men opened the sacks.}
% mdparagraph:P320:end

% output: NEW-U18-F03-O005
% mdparagraph:P321:begin
\itempara{5}{Identify the direct objects and explain which words change when the subject, object or tense changes.}
% mdparagraph:P321:end

\reserve{7\baselineskip}
% mdparagraph:P322:begin
\section*{6. Cook: the Joseph sequence}\BookHeading{nav-ch18-learner-0090}{subsection}{6. Cook: the Joseph sequence}{1}
% mdparagraph:P322:end

% teaching: NEW-U18-SOURCE-HINTS
% mdparagraph:P323:begin
Translate every numbered prompt into Old English. Parentheses and brackets are Cook’s original controls: a lexical hint selects a word, om. requests omission of a separate word, and an order or number instruction applies where marked. The numbered grammar references are explained beside their prompts. English wording is retained exactly, including commandest in 47.11. Modern hints and supplied expressions remain available during composition. Preserve every participant, clause and time relation.
% mdparagraph:P323:end

\reserve{7\baselineskip}
% mdparagraph:P324:begin
\section*{Cook 44. The return is proposed}\BookHeading{nav-ch18-learner-0091}{subsection}{Cook 44. The return is proposed}{1}
% mdparagraph:P324:end

\par\begin{minipage}{\linewidth}\raggedright
% mdparagraph:P325:begin
\subsection*{Cook 44.1}\BookHeading{nav-ch18-learner-0092}{subsubsection}{Cook 44.1}{0}
% mdparagraph:P325:end

% item: EX-XII-XIII-44-01
% mdparagraph:P326:begin
Jacob (Iacob) said to his sons, Go and buy us a certain amount of food (męte).
% mdparagraph:P326:end

% mdparagraph:P327:begin
\textbf{Hint:} Iacob.
% mdparagraph:P327:end

% mdparagraph:P328:begin
\textbf{Hint:} męte.
% mdparagraph:P328:end

\end{minipage}\par\addvspace{4pt}

% teaching: U18-SUP-IACOB-NOM
% mdparagraph:P329:begin
\textbf{Supplied:} Iacob — Jacob; nominative, masculine, singular.
% mdparagraph:P329:end

% output: EX-XII-XIII-44-01-O001
\par\begin{minipage}{\linewidth}\raggedright
% mdparagraph:P330:begin
\subsection*{Cook 44.2}\BookHeading{nav-ch18-learner-0093}{subsubsection}{Cook 44.2}{0}
% mdparagraph:P330:end

% item: EX-XII-XIII-44-02
% mdparagraph:P331:begin
Then said Judah (Iudas), The chief of the land swore oaths before us that we should not (mōtan) come before him, except we bring our youngest brother with us.
% mdparagraph:P331:end

% mdparagraph:P332:begin
\textbf{Hint:} Iudas.
% mdparagraph:P332:end

% mdparagraph:P333:begin
\textbf{Hint:} mōtan.
% mdparagraph:P333:end

\end{minipage}\par\addvspace{4pt}

% teaching: U18-SUP-IUDAS-NOM
% mdparagraph:P334:begin
\textbf{Supplied:} Iudas — Judah; nominative, masculine, singular.
% mdparagraph:P334:end

% mdparagraph:P335:begin
\textbf{Supplied:} ætforan — before; in front of.
% mdparagraph:P335:end

% mdparagraph:P336:begin
\textbf{Supplied:} ætforan — before; in front of.
% mdparagraph:P336:end

% teaching: U18-SUP-CONDITIONAL
% mdparagraph:P337:begin
\textbf{Supplied:} būtan — unless; except if.
% mdparagraph:P337:end

% teaching: U18-SUP-GINGSTAN-ACC-M-SG
% mdparagraph:P338:begin
\textbf{Supplied:} gingstan — youngest; accusative, masculine, singular.
% mdparagraph:P338:end

% teaching: NEW-U18-SUP-AETFORAN-DATIVE
% mdparagraph:P339:begin
\textbf{Supplied construction:} \textbf{ætforan ūs} — Before us. Ūs is dative plural after ætforan.
% mdparagraph:P339:end

% mdparagraph:P340:begin
\textbf{Supplied construction:} \textbf{cuman ætforan him} — Come before him. Him is dative singular, referring to the chief. Ætforan retains dative government here even with cuman; copy the whole supplied phrase.
% mdparagraph:P340:end

% output: EX-XII-XIII-44-02-O001
\par\begin{minipage}{\linewidth}\raggedright
% mdparagraph:P341:begin
\subsection*{Cook 44.3}\BookHeading{nav-ch18-learner-0094}{subsubsection}{Cook 44.3}{0}
% mdparagraph:P341:end

% item: EX-XII-XIII-44-03
% mdparagraph:P342:begin
If thou wilt send him with us (verb last), we (om. will) go together and buy the things that (87. d) we need (beðurfan, 131).
% mdparagraph:P342:end

% mdparagraph:P343:begin
\textbf{Hint:} beðurfan.
% mdparagraph:P343:end

% mdparagraph:P344:begin
\textbf{Supplied:} ætgædere — together; at once.
% mdparagraph:P344:end

% mdparagraph:P345:begin
\textbf{(verb last)} Put the finite verb last in this specified clause. This local composition target is not a universal Old English order rule.
% mdparagraph:P345:end

% mdparagraph:P346:begin
\textbf{(om. will)} Omit a separate Old English word for the marked English expression will. Preserve its grammatical relation through the surrounding selected construction.
% mdparagraph:P346:end

% mdparagraph:P347:begin
\textbf{(87. d)} Invariant relative ðe does not display case endings. The head ðā ðing is object of buy, while within the relative ðe supplies the complement required by beðurfan; its spelling does not transfer the head’s accusative to that internal role.
% mdparagraph:P347:end

% mdparagraph:P348:begin
\textbf{(beðurfan, 131)} Beðurfan is a preterite-present: its present singular uses beðearf, plural beðurfon; weak dental past uses beðorfte/\allowbreak{}beðorfton. The complete twelve-cell indicative system is displayed, with the selected genitive frame separately taught.
% mdparagraph:P348:end

% output: EX-XII-XIII-44-03-O001
\end{minipage}\par\addvspace{4pt}

\par\begin{minipage}{\linewidth}\raggedright
% mdparagraph:P349:begin
\subsection*{Cook 44.4}\BookHeading{nav-ch18-learner-0095}{subsubsection}{Cook 44.4}{0}
% mdparagraph:P349:end

% item: EX-XII-XIII-44-04
% mdparagraph:P350:begin
Ye dealt ill (earmlice) with (wið) me, when ye said to him that ye had more (77) brethren (154).
% mdparagraph:P350:end

% mdparagraph:P351:begin
\textbf{Hint:} earmlice.
% mdparagraph:P351:end

% mdparagraph:P352:begin
\textbf{Hint:} wið.
% mdparagraph:P352:end

\end{minipage}\par\addvspace{4pt}

% teaching: NEW-U18-GRAM-U18-H-44-04
% teaching: NEW-U18-VOC-DON-U18-DEAL-WITH
% teaching: U18-SUP-EARMLICE-INVARIANT
% teaching: U18-SUP-TOWARD
% mdparagraph:P353:begin
\textbf{Supplied:} Gē earmlice wið mē dydon, — Complete interpersonal clause: you behaved/\allowbreak{}dealt badly toward me. Wið names the affected counterpart; it is not a physical battle assertion. Dydon agrees with plural gē.
% mdparagraph:P353:end

% mdparagraph:P354:begin
\textbf{Supplied:} earmlice — cruelly; miserably.
% mdparagraph:P354:end

% mdparagraph:P355:begin
\textbf{Supplied:} wið — toward; in relation to a person.
% mdparagraph:P355:end

% mdparagraph:P356:begin
\textbf{Supplied:} dydon — deal; behave toward someone; indicative, preterite, second person, plural.
% mdparagraph:P356:end

% mdparagraph:P357:begin
\textbf{(77)} Mā expresses greater quantity, more; it does not agree as a comparative adjective with brothers. Its complement is the genitive plural brōðra.
% mdparagraph:P357:end

% mdparagraph:P358:begin
\textbf{(154)} Partitive genitive after a quantity: mā brōðra means more brothers. The brothers are the set whose quantity is compared.
% mdparagraph:P358:end

% output: EX-XII-XIII-44-04-O001
\par\begin{minipage}{\linewidth}\raggedright
% mdparagraph:P359:begin
\subsection*{Cook 44.5}\BookHeading{nav-ch18-learner-0096}{subsubsection}{Cook 44.5}{0}
% mdparagraph:P359:end

% item: EX-XII-XIII-44-05
% mdparagraph:P360:begin
The man asked (āscian) us all whether our father was living [lived] and whether we had brethren, and we answered him; we knew not that he would (willan) desire (giernan) that (156. a), that we should bring (lǣdan, 194. b) our brother thither.
% mdparagraph:P360:end

% mdparagraph:P361:begin
\textbf{Hint:} āscian.
% mdparagraph:P361:end

% mdparagraph:P362:begin
\textbf{Hint:} willan.
% mdparagraph:P362:end

% mdparagraph:P363:begin
\textbf{Hint:} giernan.
% mdparagraph:P363:end

% mdparagraph:P364:begin
\textbf{Hint:} lǣdan.
% mdparagraph:P364:end

% mdparagraph:P365:begin
\textbf{[lived]} Use the bracketed English as the explicit target wording, preserving the complete sentence and every participant: lived.
% mdparagraph:P365:end

% mdparagraph:P366:begin
\textbf{(156. a)} Giernan takes genitive ðæs here; the following ðæt clause explains the thing desired. Retain both words: ðæs is the genitive demonstrative, ðæt the content conjunction.
% mdparagraph:P366:end

% mdparagraph:P367:begin
\textbf{(lǣdan, 194. b)} The desired bringing is expressed by past subjunctive lǣdden under the former desire. The brothers are plural subject; the youngest brother is accusative object.
% mdparagraph:P367:end

% output: EX-XII-XIII-44-05-O001
\end{minipage}\par\addvspace{4pt}

\par\begin{minipage}{\linewidth}\raggedright
% mdparagraph:P368:begin
\subsection*{Cook 44.6}\BookHeading{nav-ch18-learner-0097}{subsubsection}{Cook 44.6}{0}
% mdparagraph:P368:end

% item: EX-XII-XIII-44-06
% mdparagraph:P369:begin
Judah said, Send the boy with me, that we may (mugan) buy the things that we need, that we perish (forweorðan) not; except I bring him back home, may I be (193. a) guilty (scyldig).
% mdparagraph:P369:end

% mdparagraph:P370:begin
\textbf{Hint:} mugan.
% mdparagraph:P370:end

% mdparagraph:P371:begin
\textbf{Hint:} forweorðan.
% mdparagraph:P371:end

% mdparagraph:P372:begin
\textbf{Hint:} scyldig.
% mdparagraph:P372:end

% mdparagraph:P373:begin
\textbf{Supplied:} Iudas — Judah; nominative, masculine, singular.
% mdparagraph:P373:end

% mdparagraph:P374:begin
\textbf{Supplied:} būtan — unless; except if.
% mdparagraph:P374:end

% mdparagraph:P375:begin
\textbf{(193. a)} Sīe ic scyldig is an independent present-subjunctive wish or self-imprecation: let me be guilty. It is not a present indicative statement of guilt.
% mdparagraph:P375:end

% output: EX-XII-XIII-44-06-O001
\end{minipage}\par\addvspace{4pt}

\par\begin{minipage}{\linewidth}\raggedright
% mdparagraph:P376:begin
\subsection*{Cook 44.7}\BookHeading{nav-ch18-learner-0098}{subsubsection}{Cook 44.7}{0}
% mdparagraph:P376:end

% item: EX-XII-XIII-44-07
% mdparagraph:P377:begin
Were it not for this delay [if this delay were not], we had [were, 188] come (om. a) second time (sīð, 176).
% mdparagraph:P377:end

% mdparagraph:P378:begin
\textbf{Hint:} sīð.
% mdparagraph:P378:end

\end{minipage}\par\addvspace{4pt}

% teaching: NEW-U18-GRAM-U18-H-44-07
% teaching: NEW-U18-VOC-WESAN-U18-ANTERIOR-COUNTERFACTUAL
% teaching: U18-SUP-INS-SG
% mdparagraph:P379:begin
\textbf{Supplied:} Gif ðēos ęlcung nǣre, wē wǣren ōðre sīðe cumen. — Complete contrary condition and anterior result: if this delay were not, we would already have come a second time. Nǣre is ne+wǣre, third-singular past subjunctive; wǣren is first-plural past subjunctive. Cumen is a past participle with intransitive wesan, not the identically spelled present-subjunctive plural. Ōðre sīðe is adverbial instrumental singular, expressly 176; the sīð hint names the lemma and does not suppress its ending. The entire combined construction and contracted form are supplied; no general counterfactual-perfect mastery is inferred.
% mdparagraph:P379:end

% teaching: NEW-U18-GRAM-U18-HELP-SECOND-TIME
% mdparagraph:P380:begin
\textbf{Supplied:} ōðre sīðe — A second time: both visible instrumental singular forms in the complete hypothetical clause of 44.7. The source headword sīð is not an instruction to leave the noun uninflected.
% mdparagraph:P380:end

% mdparagraph:P381:begin
\textbf{Supplied:} nǣre — be: present and past states/\allowbreak{}identification; subjunctive, preterite, third person, singular.
% mdparagraph:P381:end

% mdparagraph:P382:begin
\textbf{Supplied:} wǣren — be as auxiliary of counterfactual anterior coming; subjunctive, preterite, first person, plural.
% mdparagraph:P382:end

% mdparagraph:P383:begin
\textbf{Supplied:} ōðre — second; instrumental, masculine, singular, strong.
% mdparagraph:P383:end

% mdparagraph:P384:begin
\textbf{Supplied:} sīðe — time; occasion; instrumental, masculine, singular.
% mdparagraph:P384:end

% mdparagraph:P385:begin
\textbf{Supplied:} cumen — come; past, participle.
% mdparagraph:P385:end

% mdparagraph:P386:begin
\textbf{[if this delay were not]} Use the bracketed English as the explicit target wording, preserving the complete sentence and every participant: if this delay were not.
% mdparagraph:P386:end

% mdparagraph:P387:begin
\textbf{[were, 188]} Use the bracketed English as the explicit target wording, preserving the complete sentence and every participant: were, 188. The source deliberately says were, pointing to intransitive wesan plus past participle cumen. In the complete supplied contrary condition, past subjunctive wǣren makes the coming anterior to the imagined result; do not replace it with unassigned habban plus an intransitive participle.
% mdparagraph:P387:end

% mdparagraph:P388:begin
\textbf{(om. a)} Omit a separate Old English word for the marked English expression a. Preserve its grammatical relation through the surrounding selected construction.
% mdparagraph:P388:end

% mdparagraph:P389:begin
\textbf{(sīð, 176)} Section 176 is adverbial instrumental. A second time is ōðre sīðe: ōðre instrumental singular, sīðe the noun’s dative/\allowbreak{}instrumental singular. The sīð hint is a lemma, not a command to suppress the noun ending.
% mdparagraph:P389:end

% output: EX-XII-XIII-44-07-O001
\par\begin{minipage}{\linewidth}\raggedright
% mdparagraph:P390:begin
\subsection*{Cook 44.8}\BookHeading{nav-ch18-learner-0099}{subsubsection}{Cook 44.8}{0}
% mdparagraph:P390:end

% item: EX-XII-XIII-44-08
% mdparagraph:P391:begin
Then said Jacob, If ye needs (nīede, 176) must do so, do as ye will; take of (of) your best fruits in vessels, and carry (bringan) presents (om. to) the man, honey (hunig, 47) and acorns (æcern, 47).
% mdparagraph:P391:end

% mdparagraph:P392:begin
\textbf{Hint:} nīede.
% mdparagraph:P392:end

% mdparagraph:P393:begin
\textbf{Hint:} of.
% mdparagraph:P393:end

% mdparagraph:P394:begin
\textbf{Hint:} bringan.
% mdparagraph:P394:end

% mdparagraph:P395:begin
\textbf{Hint:} hunig.
% mdparagraph:P395:end

% mdparagraph:P396:begin
\textbf{Hint:} æcern.
% mdparagraph:P396:end

% mdparagraph:P397:begin
\textbf{Supplied:} Iacob — Jacob; nominative, masculine, singular.
% mdparagraph:P397:end

\end{minipage}\par\addvspace{4pt}

% teaching: U18-SUP-NIEDE-INVARIANT
% mdparagraph:P398:begin
\textbf{Supplied:} nīede — necessarily; of necessity.
% mdparagraph:P398:end

% mdparagraph:P399:begin
\textbf{(nīede, 176)} Nīede is the supplied adverbial instrumental expressing necessity, necessarily/\allowbreak{}by necessity. With sculon it conveys must needs; do not treat it as a finite form of need.
% mdparagraph:P399:end

% mdparagraph:P400:begin
\textbf{(om. to)} Omit a separate Old English word for the marked English expression to. Preserve its grammatical relation through the surrounding selected construction.
% mdparagraph:P400:end

% mdparagraph:P401:begin
\textbf{(hunig, 47)} The number refers to the noun’s declension. Use the complete hunig table in section 2 and its stated stem rule.
% mdparagraph:P401:end

% mdparagraph:P402:begin
\textbf{(æcern, 47)} The number refers to the noun’s declension. Use the complete æcern table in section 2 and its stated stem rule.
% mdparagraph:P402:end

% output: EX-XII-XIII-44-08-O001
\reserve{7\baselineskip}
% mdparagraph:P403:begin
\section*{Cook 45. The brothers return to Egypt}\BookHeading{nav-ch18-learner-0100}{subsection}{Cook 45. The brothers return to Egypt}{1}
% mdparagraph:P403:end

\par\begin{minipage}{\linewidth}\raggedright
% mdparagraph:P404:begin
\subsection*{Cook 45.1}\BookHeading{nav-ch18-learner-0101}{subsubsection}{Cook 45.1}{0}
% mdparagraph:P404:end

% item: EX-XII-XIII-45-01
% mdparagraph:P405:begin
Take your (plur.) brother and go to the man.
% mdparagraph:P405:end

% mdparagraph:P406:begin
\textbf{(plur.)} Address more than one person and select plural personal and finite forms.
% mdparagraph:P406:end

% output: EX-XII-XIII-45-01-O001
\end{minipage}\par\addvspace{4pt}

\par\begin{minipage}{\linewidth}\raggedright
% mdparagraph:P407:begin
\subsection*{Cook 45.2}\BookHeading{nav-ch18-learner-0102}{subsubsection}{Cook 45.2}{0}
% mdparagraph:P407:end

% item: EX-XII-XIII-45-02
% mdparagraph:P408:begin
And my Lord make (gedōn) him gracious (glæd) toward (wið) you, that he may give (āgiefan, 196. f) you your brother, whom he hath with him; I am now bereft [despoiled] of my children (156).
% mdparagraph:P408:end

% mdparagraph:P409:begin
\textbf{Hint:} gedōn.
% mdparagraph:P409:end

% mdparagraph:P410:begin
\textbf{Hint:} glæd.
% mdparagraph:P410:end

% mdparagraph:P411:begin
\textbf{Hint:} wið.
% mdparagraph:P411:end

% mdparagraph:P412:begin
\textbf{Hint:} āgiefan.
% mdparagraph:P412:end

% mdparagraph:P413:begin
\textbf{Supplied:} wið — toward; in relation to a person.
% mdparagraph:P413:end

% mdparagraph:P414:begin
\textbf{(āgiefan, 196. f)} Āgiefe is third-singular present subjunctive in the purpose clause, so that he may restore the brother to you. Ēow is dative recipient and ēowerne brōðor accusative returned person.
% mdparagraph:P414:end

% mdparagraph:P415:begin
\textbf{[despoiled]} Use the bracketed English as the explicit target wording, preserving the complete sentence and every participant: despoiled.
% mdparagraph:P415:end

% mdparagraph:P416:begin
\textbf{(156)} The genitive denotes what someone has been deprived of: mīnra bearna remains genitive after passive eom berēafod. In the active frame the deprived person is accusative; the genitive possessions do not become the passive subject.
% mdparagraph:P416:end

% output: EX-XII-XIII-45-02-O001
\end{minipage}\par\addvspace{4pt}

\par\begin{minipage}{\linewidth}\raggedright
% mdparagraph:P417:begin
\subsection*{Cook 45.3}\BookHeading{nav-ch18-learner-0103}{subsubsection}{Cook 45.3}{0}
% mdparagraph:P417:end

% item: EX-XII-XIII-45-03
% mdparagraph:P418:begin
Then they took the presents and double (twīfeald) money (feoh, 47), and went to the land of the Egyptians to Joseph (Iosep).
% mdparagraph:P418:end

% mdparagraph:P419:begin
\textbf{Hint:} twīfeald.
% mdparagraph:P419:end

% mdparagraph:P420:begin
\textbf{Hint:} feoh.
% mdparagraph:P420:end

% mdparagraph:P421:begin
\textbf{Hint:} Iosep.
% mdparagraph:P421:end

\end{minipage}\par\addvspace{4pt}

% teaching: U18-SUP-EGYPTA-GEN-PL
% mdparagraph:P422:begin
\textbf{Supplied:} Egypta — Egyptians; genitive, masculine, plural.
% mdparagraph:P422:end

% teaching: U18-SUP-IOSEP-DAT
% mdparagraph:P423:begin
\textbf{Supplied:} Iosepe — Joseph; dative, masculine, singular.
% mdparagraph:P423:end

% mdparagraph:P424:begin
\textbf{(feoh, 47)} The number refers to the noun’s declension. Use the complete feoh table in section 2 and its stated stem rule.
% mdparagraph:P424:end

% output: EX-XII-XIII-45-03-O001
\par\begin{minipage}{\linewidth}\raggedright
% mdparagraph:P425:begin
\subsection*{Cook 45.4}\BookHeading{nav-ch18-learner-0104}{subsubsection}{Cook 45.4}{0}
% mdparagraph:P425:end

% item: EX-XII-XIII-45-04
% mdparagraph:P426:begin
When he saw them, then said he to (tō) his steward (gerēfa, 53), Lead in these men, and prepare our feast, that they may (mugan) eat with me at noon (tō middes dæges, 158; 166. 1).
% mdparagraph:P426:end

% mdparagraph:P427:begin
\textbf{Hint:} tō.
% mdparagraph:P427:end

% mdparagraph:P428:begin
\textbf{Hint:} gerēfa.
% mdparagraph:P428:end

% mdparagraph:P429:begin
\textbf{Hint:} mugan.
% mdparagraph:P429:end

% mdparagraph:P430:begin
\textbf{Hint:} tō middes dæges.
% mdparagraph:P430:end

% mdparagraph:P431:begin
\textbf{Supplied:} in — in; inward.
% mdparagraph:P431:end

\end{minipage}\par\addvspace{4pt}

% teaching: U18-SUP-MIDDES-INVARIANT
% mdparagraph:P432:begin
\textbf{Supplied:} middes — midday, in tō middes dæges.
% mdparagraph:P432:end

% mdparagraph:P433:begin
\textbf{(gerēfa, 53)} Gerēfa follows the earlier weak masculine pattern: nominative singular -a; other singular cases and nominative/\allowbreak{}accusative plural -an; genitive plural -ena; dative plural -um.
% mdparagraph:P433:end

% mdparagraph:P434:begin
\textbf{(tō middes dæges, 158; 166. 1)} Tō middes dæges is the source-supplied time phrase at midday. Here tō governs a genitive; middes agrees with genitive dæges. Section 166.1 supplies the parallel principle of adjective agreement within the prepositional phrase, not a direction to change this printed genitive to dative.
% mdparagraph:P434:end

% output: EX-XII-XIII-45-04-O001
\par\begin{minipage}{\linewidth}\raggedright
% mdparagraph:P435:begin
\subsection*{Cook 45.5}\BookHeading{nav-ch18-learner-0105}{subsubsection}{Cook 45.5}{0}
% mdparagraph:P435:end

% item: EX-XII-XIII-45-05
% mdparagraph:P436:begin
And he did as he was commanded [was commanded him; cf. p. 177, l. 10].
% mdparagraph:P436:end

\end{minipage}\par\addvspace{4pt}

% teaching: NEW-U18-GRAM-U18-H-45-05
% mdparagraph:P437:begin
\textbf{Supplied:} hē dyde swā him boden wæs. — Impersonal command passive: he did as was commanded him. Him is dative addressee; boden wæs is singular without promoting him to nominative. Cook 177.10 swā hiere beboden wæs is the explicitly cited parallel; replacing be-bēodan with independently attested command bēodan is a declared lexical adaptation.
% mdparagraph:P437:end

% mdparagraph:P438:begin
\textbf{[was commanded him; cf. p. 177, l. 10]} Use the bracketed English as the explicit target wording, preserving the complete sentence and every participant: was commanded him; cf. p. 177, l. 10. Cook 177.10 has swā hiere beboden wæs. The supplied expression uses bēodan: swā him boden wæs, as was commanded him. Him is a dative addressee in an impersonal passive, not a nominative subject.
% mdparagraph:P438:end

% output: EX-XII-XIII-45-05-O001
\par\begin{minipage}{\linewidth}\raggedright
% mdparagraph:P439:begin
\subsection*{Cook 45.6}\BookHeading{nav-ch18-learner-0106}{subsubsection}{Cook 45.6}{0}
% mdparagraph:P439:end

% item: EX-XII-XIII-45-06
% mdparagraph:P440:begin
When they were there (ðǣr), then they became greatly frightened, and said to one another (betwēonan him), Because of (for) the money that we found in our sacks (sacc, 43), we (om. have) been led (62) in hither, that he may despoil us of our asses, and of our things, and enslave us (our)selves.
% mdparagraph:P440:end

% mdparagraph:P441:begin
\textbf{Hint:} ðǣr.
% mdparagraph:P441:end

% mdparagraph:P442:begin
\textbf{Hint:} betwēonan him.
% mdparagraph:P442:end

% mdparagraph:P443:begin
\textbf{Hint:} for.
% mdparagraph:P443:end

% mdparagraph:P444:begin
\textbf{Hint:} sacc.
% mdparagraph:P444:end

\end{minipage}\par\addvspace{4pt}

% teaching: NEW-U18-GRAM-U18-H-45-06-PASSIVE
% mdparagraph:P445:begin
\textbf{Supplied:} wē sindon hider in lǣded, — We have been led in hither. Present wesan plus the past participle is the selected passive/\allowbreak{}result expression; uninflected lǣded with plural wē is expressly allowed by 62. Inflected plural lǣdde is not needed to satisfy the prompt and is not silently treated as taught.
% mdparagraph:P445:end

% teaching: NEW-U18-GRAM-U18-HELP-FOR-CAUSE
% teaching: U18-SUP-CAUSE
% mdparagraph:P446:begin
\textbf{Supplied:} For ðǣm fēo — Because of that money /\allowbreak{} on account of that money. Ðǣm and fēo are dative neuter singular.
% mdparagraph:P446:end

% teaching: U18-SUP-BETWEONAN-INVARIANT
% mdparagraph:P447:begin
\textbf{Supplied:} betwēonan — between; toward one another.
% mdparagraph:P447:end

% mdparagraph:P448:begin
\textbf{Supplied:} For — because of; on account of; governs dative.
% mdparagraph:P448:end

% mdparagraph:P449:begin
\textbf{Supplied:} in — in; inward.
% mdparagraph:P449:end

% mdparagraph:P450:begin
\textbf{(sacc, 43)} The number refers to the noun’s declension. Use the complete sacc table in section 2 and its stated stem rule.
% mdparagraph:P450:end

% mdparagraph:P451:begin
\textbf{(om. have)} Omit a separate Old English word for the marked English expression have. Preserve its grammatical relation through the surrounding selected construction.
% mdparagraph:P451:end

% mdparagraph:P452:begin
\textbf{(62)} A past participle used predicatively may be uninflected. Thus plural wē can have base lǣded in wē sindon hider in lǣded. No new plural participial ending is required by the prompt.
% mdparagraph:P452:end

% mdparagraph:P453:begin
\textbf{(our)} The parenthesis (our) belongs to the English expression (our)selves, reinforcing the preceding us. In the supplied clause, ūs selfe means ourselves: selfe agrees with accusative plural ūs. This parenthesis does not introduce a possessive modifying the things.
% mdparagraph:P453:end

% output: EX-XII-XIII-45-06-O001
\par\begin{minipage}{\linewidth}\raggedright
% mdparagraph:P454:begin
\subsection*{Cook 45.7}\BookHeading{nav-ch18-learner-0107}{subsubsection}{Cook 45.7}{0}
% mdparagraph:P454:end

% item: EX-XII-XIII-45-07
% mdparagraph:P455:begin
Then said they to (tō) the steward, O sir, we pray (biddan) thee that thou hear (hlystan, 156) our speech; we were here a while ago and bought wheat (hwǣte, 44) for ourselves (161; 184. a).
% mdparagraph:P455:end

% mdparagraph:P456:begin
\textbf{Hint:} tō.
% mdparagraph:P456:end

% mdparagraph:P457:begin
\textbf{Hint:} biddan.
% mdparagraph:P457:end

% mdparagraph:P458:begin
\textbf{Hint:} hlystan.
% mdparagraph:P458:end

% mdparagraph:P459:begin
\textbf{Hint:} hwǣte.
% mdparagraph:P459:end

% mdparagraph:P460:begin
\textbf{(hlystan, 156)} Hlystan governs genitive speech/\allowbreak{}content, so ūrre sprǣce is genitive feminine singular. Hlyste is second-singular present subjunctive after the plea.
% mdparagraph:P460:end

% mdparagraph:P461:begin
\textbf{(hwǣte, 44)} The number refers to the noun’s declension. Use the complete hwǣte table in section 2 and its stated stem rule.
% mdparagraph:P461:end

% mdparagraph:P462:begin
\textbf{(161; 184. a)} Ūs is the benefactive/\allowbreak{}reflexive dative: we bought wheat for ourselves. It refers to the buying subject, while wheat is the accusative thing purchased. Retain the phrase; the prompt gives no omission instruction.
% mdparagraph:P462:end

\end{minipage}\par\addvspace{4pt}

% teaching: NEW-U18-SUP-LEOF-ADDRESS
% mdparagraph:P463:begin
\textbf{Supplied construction:} \textbf{Lēof} — Sir. Lēof is a substantival use of the strong masculine nominative singular adjective in direct address. This respectful address sense is supplied here; the previously taught meaning dear alone does not establish it.
% mdparagraph:P463:end

% teaching: NEW-U18-SUP-HWILUM-AER
% mdparagraph:P464:begin
\textbf{Supplied construction:} \textbf{hwīlum ǣr} — A while ago; formerly. Use hwīlum ǣr as a complete earlier-time expression. Hwīlum here is part of this supplied expression, rather than the frequency meaning sometimes previously learned.
% mdparagraph:P464:end

% output: EX-XII-XIII-45-07-O001
\par\begin{minipage}{\linewidth}\raggedright
% mdparagraph:P465:begin
\subsection*{Cook 45.8}\BookHeading{nav-ch18-learner-0108}{subsubsection}{Cook 45.8}{0}
% mdparagraph:P465:end

% item: EX-XII-XIII-45-08
% mdparagraph:P466:begin
When we had trafficked (cēapian) and were on the way home (hāmweard wǣron), then we undid (undōn) our sacks; then we found the money therein which we had given (see ǣr, adv.).
% mdparagraph:P466:end

% mdparagraph:P467:begin
\textbf{Hint:} cēapian.
% mdparagraph:P467:end

% mdparagraph:P468:begin
\textbf{Hint:} hāmweard wǣron.
% mdparagraph:P468:end

% mdparagraph:P469:begin
\textbf{Hint:} undōn.
% mdparagraph:P469:end

% mdparagraph:P470:begin
\textbf{Hint:} ǣr.
% mdparagraph:P470:end

\end{minipage}\par\addvspace{4pt}

% teaching: U18-SUP-HAMWEARD-INVARIANT
% mdparagraph:P471:begin
\textbf{Supplied:} hāmweard — homeward.
% mdparagraph:P471:end

% teaching: U18-SUP-THAERON-INVARIANT
% mdparagraph:P472:begin
\textbf{Supplied:} ðǣron — therein; in it.
% mdparagraph:P472:end

% output: EX-XII-XIII-45-08-O001
\par\begin{minipage}{\linewidth}\raggedright
% mdparagraph:P473:begin
\subsection*{Cook 45.9}\BookHeading{nav-ch18-learner-0109}{subsubsection}{Cook 45.9}{0}
% mdparagraph:P473:end

% item: EX-XII-XIII-45-09
% mdparagraph:P474:begin
Now we have brought it back, and also other silver to traffic with [with to traffic]; and we know not who put (dōn) it into (on, 172) our sacks.
% mdparagraph:P474:end

% mdparagraph:P475:begin
\textbf{Hint:} dōn.
% mdparagraph:P475:end

% mdparagraph:P476:begin
\textbf{Hint:} on.
% mdparagraph:P476:end

\end{minipage}\par\addvspace{4pt}

% teaching: NEW-U18-GRAM-U18-H-45-09
% teaching: U18-SUP-INSTRUMENT-STRANDED
% mdparagraph:P477:begin
\textbf{Supplied:} and ēac ōðer seolfor mid tō cēapianne; — The earlier we brought governs the coordinated silver-object; mid marks silver as the instrument of trading and is stranded before the purposive inflected infinitive. This does not require a new finite verb or make the silver a human trading partner. Supply the whole compressed phrase.
% mdparagraph:P477:end

% mdparagraph:P478:begin
\textbf{Supplied:} mid — with; by means of.
% mdparagraph:P478:end

% mdparagraph:P479:begin
\textbf{[with to traffic]} Use the bracketed English as the explicit target wording, preserving the complete sentence and every participant: with to traffic.
% mdparagraph:P479:end

% mdparagraph:P480:begin
\textbf{(on, 172)} On plus accusative marks putting something into a container. Ūre saccas is accusative plural after on in the clause about who put the money there; static location would instead take a dative.
% mdparagraph:P480:end

% output: EX-XII-XIII-45-09-O001
\reserve{7\baselineskip}
% mdparagraph:P481:begin
\section*{Cook 46. The welcome and the cup}\BookHeading{nav-ch18-learner-0110}{subsection}{Cook 46. The welcome and the cup}{1}
% mdparagraph:P481:end

\par\begin{minipage}{\linewidth}\raggedright
% mdparagraph:P482:begin
\subsection*{Cook 46.1}\BookHeading{nav-ch18-learner-0111}{subsubsection}{Cook 46.1}{0}
% mdparagraph:P482:end

% item: EX-XII-XIII-46-01
% mdparagraph:P483:begin
Then said the steward, Peace be with you, fear (ondrǣdan, 184) ye not; your God and your father’s God (om. hath) given you treasure in your sacks.
% mdparagraph:P483:end

% mdparagraph:P484:begin
\textbf{Hint:} ondrǣdan.
% mdparagraph:P484:end

\end{minipage}\par\addvspace{4pt}

% teaching: NEW-U18-GRAM-U18-H-46-01
% mdparagraph:P485:begin
\textbf{Supplied:} ēower God and ēowres fæder God sealde ēow goldhord — The one deity is described as your God and your father’s God, explaining singular sealde. Ēowres agrees with the genitive father; second ēow is dative recipient. The expression does not assert two gods.
% mdparagraph:P485:end

% mdparagraph:P486:begin
\textbf{(ondrǣdan, 184)} Ondrǣdað ēow uses a reflexive dative: do not fear. The plural addressees are subject, and ēow refers back to them. It is not the object feared; no feared object is expressed in this item.
% mdparagraph:P486:end

% mdparagraph:P487:begin
\textbf{(om. hath)} Omit a separate Old English word for the marked English expression hath. Preserve its grammatical relation through the surrounding selected construction.
% mdparagraph:P487:end

% output: EX-XII-XIII-46-01-O001
\par\begin{minipage}{\linewidth}\raggedright
% mdparagraph:P488:begin
\subsection*{Cook 46.2}\BookHeading{nav-ch18-learner-0112}{subsubsection}{Cook 46.2}{0}
% mdparagraph:P488:end

% item: EX-XII-XIII-46-02
% mdparagraph:P489:begin
And they washed (101) their feet, and he gave their asses provender [fodder].
% mdparagraph:P489:end

% mdparagraph:P490:begin
\textbf{(101)} Ðwēan is a contracted strong verb. Select the earlier past plural ðwōgon for they washed; the headword’s contracted vowel does not supply the past stem by itself.
% mdparagraph:P490:end

% mdparagraph:P491:begin
\textbf{[fodder]} Use the bracketed English as the explicit target wording, preserving the complete sentence and every participant: fodder.
% mdparagraph:P491:end

% output: EX-XII-XIII-46-02-O001
\end{minipage}\par\addvspace{4pt}

\par\begin{minipage}{\linewidth}\raggedright
% mdparagraph:P492:begin
\subsection*{Cook 46.3}\BookHeading{nav-ch18-learner-0113}{subsubsection}{Cook 46.3}{0}
% mdparagraph:P492:end

% item: EX-XII-XIII-46-03
% mdparagraph:P493:begin
They set (lęcgean) forth their presents, against that (ongēan ðætte) Joseph (om. should) enter.
% mdparagraph:P493:end

% mdparagraph:P494:begin
\textbf{Hint:} lęcgean.
% mdparagraph:P494:end

% mdparagraph:P495:begin
\textbf{Hint:} ongēan ðætte.
% mdparagraph:P495:end

\end{minipage}\par\addvspace{4pt}

% teaching: NEW-U18-GRAM-U18-H-46-03
% teaching: U18-SUP-IOSEP-NOM
% teaching: U18-SUP-TEMPORAL-CONJUNCTION-COMPONENT
% mdparagraph:P496:begin
\textbf{Supplied:} ongēan ðætte Iosep inēode. — Against/\allowbreak{}in readiness for Joseph’s entering. The gifts are laid out in preparation for this prospective arrival; inēode is third-singular past subjunctive, an already productive ingān cell. The whole source conjunction is supplied, not equated with spatial against.
% mdparagraph:P496:end

% mdparagraph:P497:begin
\textbf{Supplied:} ongēan — by the time that.
% mdparagraph:P497:end

% mdparagraph:P498:begin
\textbf{Supplied:} ðætte — that (content conjunction).
% mdparagraph:P498:end

% mdparagraph:P499:begin
\textbf{Supplied:} Iosep — Joseph; nominative, masculine, singular.
% mdparagraph:P499:end

% mdparagraph:P500:begin
\textbf{(om. should)} Omit a separate Old English word for the marked English expression should. Preserve its grammatical relation through the surrounding selected construction.
% mdparagraph:P500:end

% output: EX-XII-XIII-46-03-O001
\par\begin{minipage}{\linewidth}\raggedright
% mdparagraph:P501:begin
\subsection*{Cook 46.4}\BookHeading{nav-ch18-learner-0114}{subsubsection}{Cook 46.4}{0}
% mdparagraph:P501:end

% item: EX-XII-XIII-46-04
% mdparagraph:P502:begin
Then Joseph recognized them, and asked them whether their father was hale, or whether he was living [lived].
% mdparagraph:P502:end

% mdparagraph:P503:begin
\textbf{Supplied:} Iosep — Joseph; nominative, masculine, singular.
% mdparagraph:P503:end

% mdparagraph:P504:begin
\textbf{[lived]} Use the bracketed English as the explicit target wording, preserving the complete sentence and every participant: lived.
% mdparagraph:P504:end

% output: EX-XII-XIII-46-04-O001
\end{minipage}\par\addvspace{4pt}

\par\begin{minipage}{\linewidth}\raggedright
% mdparagraph:P505:begin
\subsection*{Cook 46.5}\BookHeading{nav-ch18-learner-0115}{subsubsection}{Cook 46.5}{0}
% mdparagraph:P505:end

% item: EX-XII-XIII-46-05
% mdparagraph:P506:begin
Then said they, Thy servant (ðēow) our father is well (gesund); he is yet alive [liveth yet].
% mdparagraph:P506:end

% mdparagraph:P507:begin
\textbf{Hint:} ðēow.
% mdparagraph:P507:end

% mdparagraph:P508:begin
\textbf{Hint:} gesund.
% mdparagraph:P508:end

% mdparagraph:P509:begin
\textbf{[liveth yet]} Use the bracketed English as the explicit target wording, preserving the complete sentence and every participant: liveth yet.
% mdparagraph:P509:end

% output: EX-XII-XIII-46-05-O001
\end{minipage}\par\addvspace{4pt}

\par\begin{minipage}{\linewidth}\raggedright
% mdparagraph:P510:begin
\subsection*{Cook 46.6}\BookHeading{nav-ch18-learner-0116}{subsubsection}{Cook 46.6}{0}
% mdparagraph:P510:end

% item: EX-XII-XIII-46-06
% mdparagraph:P511:begin
When Joseph saw his brother Benjamin, then said he, Is this (neut.) the lad, (om. of) whom (87. d) ye told me before (foresęcgean)? And he said, God have compassion on thee, my son.
% mdparagraph:P511:end

% mdparagraph:P512:begin
\textbf{Hint:} foresęcgean.
% mdparagraph:P512:end

% mdparagraph:P513:begin
\textbf{Supplied:} Iosep — Joseph; nominative, masculine, singular.
% mdparagraph:P513:end

\end{minipage}\par\addvspace{4pt}

% teaching: U18-SUP-BENIAMIN-ACC
% mdparagraph:P514:begin
\textbf{Supplied:} Beniamin — Benjamin; accusative, masculine, singular.
% mdparagraph:P514:end

% mdparagraph:P515:begin
\textbf{(neut.)} Use a neuter demonstrative for identifying this, even though the identified boy is a person; the complete identifying clause remains visible in its support.
% mdparagraph:P515:end

% mdparagraph:P516:begin
\textbf{(om. of)} Omit a separate Old English word for the marked English expression of. Preserve its grammatical relation through the surrounding selected construction.
% mdparagraph:P516:end

% mdparagraph:P517:begin
\textbf{(87. d)} The relative pronoun ðe is invariant. It represents the boy as the person previously mentioned; mē is the dative person told. Omit English of as explicitly instructed instead of adding an unsupported preposition.
% mdparagraph:P517:end

% output: EX-XII-XIII-46-06-O001
\par\begin{minipage}{\linewidth}\raggedright
% mdparagraph:P518:begin
\subsection*{Cook 46.7}\BookHeading{nav-ch18-learner-0117}{subsubsection}{Cook 46.7}{0}
% mdparagraph:P518:end

% item: EX-XII-XIII-46-07
% mdparagraph:P519:begin
Then Joseph commanded (bēodan, 164) his steward, and said, Fill their sacks with wheat, and put the money of each of them into (on, 172) his own sack.
% mdparagraph:P519:end

% mdparagraph:P520:begin
\textbf{Hint:} bēodan.
% mdparagraph:P520:end

% mdparagraph:P521:begin
\textbf{Hint:} on.
% mdparagraph:P521:end

% mdparagraph:P522:begin
\textbf{Supplied:} Iosep — Joseph; nominative, masculine, singular.
% mdparagraph:P522:end

% mdparagraph:P523:begin
\textbf{(bēodan, 164)} The person commanded by bēodan is dative: his gerēfan. The direct imperative that follows gives the commanded act. This is the newly assigned command sense, distinct from the earlier offer sense.
% mdparagraph:P523:end

% mdparagraph:P524:begin
\textbf{(on, 172)} On plus accusative gives the destination into his own sack: on his āgenne sacc. The adjective and noun agree as accusative masculine singular.
% mdparagraph:P524:end

\end{minipage}\par\addvspace{4pt}

% teaching: NEW-U18-SUP-GEFYLLAN-MATERIAL
% mdparagraph:P525:begin
\textbf{Supplied construction:} \textbf{Gefyl hiera saccas mid hwǣte} — Fill their sacks with wheat. Gefyl is a singular command; hiera saccas is the accusative plural container, and mid hwǣte names the material with dative singular hwǣte. This full material-complement frame is supplied here.
% mdparagraph:P525:end

% output: EX-XII-XIII-46-07-O001
\par\begin{minipage}{\linewidth}\raggedright
% mdparagraph:P526:begin
\subsection*{Cook 46.8}\BookHeading{nav-ch18-learner-0118}{subsubsection}{Cook 46.8}{0}
% mdparagraph:P526:end

% item: EX-XII-XIII-46-08
% mdparagraph:P527:begin
And take my silver cup (læfel, 43), and the price (weorð, 47) of the wheat, which he gave thee, and put into the sack of the youngest; and he did so.
% mdparagraph:P527:end

% mdparagraph:P528:begin
\textbf{Hint:} læfel.
% mdparagraph:P528:end

% mdparagraph:P529:begin
\textbf{Hint:} weorð.
% mdparagraph:P529:end

\end{minipage}\par\addvspace{4pt}

% teaching: NEW-U18-GRAM-U18-H-46-08
% teaching: U18-SUP-GINGSTAN-GEN-M-SG
% teaching: U18-SUP-SEOLFRENNE-ACC-M-SG
% mdparagraph:P530:begin
\textbf{Supplied:} nim mīnne seolfrenne læfel, and ðæs hwǣtes weorð ðe hē ðē sealde, and dō on ðæs gingstan sacc; — Supplied seolfrenne is the strong accusative masculine singular adjective agreeing with the cup; ðæs hwǣtes specifies the wheat price, and the price is what he gave thee. The coordinated put shares the cup-and-price objects; no unprinted singular it is required. The youngest is the substantival weak adjective identifying the sack’s possessor. The exact silver adjective is supplied; its remaining forms are not assigned here.
% mdparagraph:P530:end

% teaching: NEW-U18-GRAM-U18-HELP-YOUNGEST
% mdparagraph:P531:begin
\textbf{Supplied:} ðæs gingstan — The youngest and eldest refer to individual brothers. Supply the masculine singular human referents and each case, rather than inferring substantive person uses from familiar adjective headwords.
% mdparagraph:P531:end

% teaching: NEW-U18-GRAM-U18-HELP-SILVER-CUP
% mdparagraph:P532:begin
\textbf{Supplied:} mīnne seolfrenne læfel, — My silver cup; the entire strong accusative adjective is supplied. Earlier proposed læfel seolfres is withdrawn because a direct material-genitive frame was not established by the checked evidence.
% mdparagraph:P532:end

% mdparagraph:P533:begin
\textbf{Supplied:} seolfrenne — silver; made of silver; accusative, masculine, singular.
% mdparagraph:P533:end

% mdparagraph:P534:begin
\textbf{Supplied:} gingstan — youngest; genitive, masculine, singular.
% mdparagraph:P534:end

% mdparagraph:P535:begin
\textbf{(læfel, 43)} The number refers to the noun’s declension. Use the complete læfel table in section 2 and its stated stem rule.
% mdparagraph:P535:end

% mdparagraph:P536:begin
\textbf{(weorð, 47)} The number refers to the noun’s declension. Use the complete weorð table in section 2 and its stated stem rule.
% mdparagraph:P536:end

% output: EX-XII-XIII-46-08-O001
\par\begin{minipage}{\linewidth}\raggedright
% mdparagraph:P537:begin
\subsection*{Cook 46.9}\BookHeading{nav-ch18-learner-0119}{subsubsection}{Cook 46.9}{0}
% mdparagraph:P537:end

% item: EX-XII-XIII-46-09
% mdparagraph:P538:begin
In (om. the) morning, as (ðā) they journeyed, and they were without (om. the) city, and had gone (faran) some part of the way, then said Joseph to (tō) his steward, Arise and go after these men, and when you have (196. b) overtaken [caught] them, then ask them why they would (willan) requite good with (mid) evil.
% mdparagraph:P538:end

% mdparagraph:P539:begin
\textbf{Hint:} ðā.
% mdparagraph:P539:end

% mdparagraph:P540:begin
\textbf{Hint:} faran.
% mdparagraph:P540:end

% mdparagraph:P541:begin
\textbf{Hint:} tō.
% mdparagraph:P541:end

% mdparagraph:P542:begin
\textbf{Hint:} willan.
% mdparagraph:P542:end

% mdparagraph:P543:begin
\textbf{Hint:} mid.
% mdparagraph:P543:end

\end{minipage}\par\addvspace{4pt}

% teaching: NEW-U18-GRAM-U18-HELP-OUTSIDE-CITY
% teaching: U18-SUP-OUTSIDE
% mdparagraph:P544:begin
\textbf{Supplied:} būtan byrig — Outside the city. Byrig is dative singular. The preposition here gives a location, rather than the separate without/\allowbreak{}absence or unless uses.
% mdparagraph:P544:end

% teaching: U18-SUP-MORNING
% mdparagraph:P545:begin
\textbf{Supplied:} On — in; at a time.
% mdparagraph:P545:end

% mdparagraph:P546:begin
\textbf{Supplied:} būtan — outside; governs dative.
% mdparagraph:P546:end

% mdparagraph:P547:begin
\textbf{Supplied:} Iosep — Joseph; nominative, masculine, singular.
% mdparagraph:P547:end

% teaching: U18-SUP-ABSTRACT-GOOD-ACC-SG
% mdparagraph:P548:begin
\textbf{Supplied:} gōd — good; benefit; kindness; accusative, neuter, singular.
% mdparagraph:P548:end

% mdparagraph:P549:begin
\textbf{Supplied:} yfele — evil; dative, neuter, singular.
% mdparagraph:P549:end

% mdparagraph:P550:begin
\textbf{(om. the)} Omit a separate Old English word for the marked English expression the. Preserve its grammatical relation through the surrounding selected construction.
% mdparagraph:P550:end

% mdparagraph:P551:begin
\textbf{(om. the)} Omit a separate Old English word for the marked English expression the. Preserve its grammatical relation through the surrounding selected construction.
% mdparagraph:P551:end

% mdparagraph:P552:begin
\textbf{(196. b)} The catching is prospective from the command: when you have caught them, then ask. Cook 196 b permits present subjunctive fō in that future temporal condition; the source’s you is singular here because Joseph addresses one steward. An English perfect does not force an OE perfect auxiliary.
% mdparagraph:P552:end

% mdparagraph:P553:begin
\textbf{[caught]} Use the bracketed English as the explicit target wording, preserving the complete sentence and every participant: caught.
% mdparagraph:P553:end

% teaching: NEW-U18-SUP-ON-MORGEN
% mdparagraph:P554:begin
\textbf{Supplied construction:} \textbf{On morgen} — In the morning. Morgen is accusative singular in this fixed temporal expression. Keep the whole phrase available; it does not introduce a general rule for every time expression with on.
% mdparagraph:P554:end

% output: EX-XII-XIII-46-09-O001
\reserve{7\baselineskip}
% mdparagraph:P555:begin
\section*{Cook 47. The accusation and the plea}\BookHeading{nav-ch18-learner-0120}{subsection}{Cook 47. The accusation and the plea}{1}
% mdparagraph:P555:end

\par\begin{minipage}{\linewidth}\raggedright
% mdparagraph:P556:begin
\subsection*{Cook 47.1}\BookHeading{nav-ch18-learner-0121}{subsubsection}{Cook 47.1}{0}
% mdparagraph:P556:end

% item: EX-XII-XIII-47-01
% mdparagraph:P557:begin
The cup, which ye stole (forstelan, 105), was very precious to my lord (hlāford).
% mdparagraph:P557:end

% mdparagraph:P558:begin
\textbf{Hint:} forstelan.
% mdparagraph:P558:end

% mdparagraph:P559:begin
\textbf{Hint:} hlāford.
% mdparagraph:P559:end

% mdparagraph:P560:begin
\textbf{(forstelan, 105)} Forstelan belongs to strongIV: present e/\allowbreak{}i, past singular æ, past plural ǣ. Second-plural forstǣlon therefore differs from third-singular forstæl in 47.8. The complete twelve-cell indicative system is displayed.
% mdparagraph:P560:end

% output: EX-XII-XIII-47-01-O001
\end{minipage}\par\addvspace{4pt}

\par\begin{minipage}{\linewidth}\raggedright
% mdparagraph:P561:begin
\subsection*{Cook 47.2}\BookHeading{nav-ch18-learner-0122}{subsubsection}{Cook 47.2}{0}
% mdparagraph:P561:end

% item: EX-XII-XIII-47-02
% mdparagraph:P562:begin
He did as was commanded him, and when they were overtaken [caught], they said, The money that we found in our sacks we brought to thee from our country; thinkest (wēnan) thou that we (om. have) stolen (stelan) thy lord’s gold or his silver?
% mdparagraph:P562:end

% mdparagraph:P563:begin
\textbf{Hint:} wēnan.
% mdparagraph:P563:end

% mdparagraph:P564:begin
\textbf{Hint:} stelan.
% mdparagraph:P564:end

% mdparagraph:P565:begin
\textbf{Supplied:} Hē dyde swā him boden wæs, — Impersonal command passive: he did as was commanded him. Him is dative addressee; boden wæs is singular without promoting him to nominative. Cook 177.10 swā hiere beboden wæs is the explicitly cited parallel; replacing be-bēodan with independently attested command bēodan is a declared lexical adaptation.
% mdparagraph:P565:end

% mdparagraph:P566:begin
\textbf{[caught]} Use the bracketed English as the explicit target wording, preserving the complete sentence and every participant: caught.
% mdparagraph:P566:end

% mdparagraph:P567:begin
\textbf{(om. have)} Omit a separate Old English word for the marked English expression have. Preserve its grammatical relation through the surrounding selected construction.
% mdparagraph:P567:end

% output: EX-XII-XIII-47-02-O001
\end{minipage}\par\addvspace{4pt}

\par\begin{minipage}{\linewidth}\raggedright
% mdparagraph:P568:begin
\subsection*{Cook 47.3}\BookHeading{nav-ch18-learner-0123}{subsubsection}{Cook 47.3}{0}
% mdparagraph:P568:end

% item: EX-XII-XIII-47-03
% mdparagraph:P569:begin
Look for [seek out] it; and let him be slain (ofslēan) with whom thou findest (195) it.
% mdparagraph:P569:end

% mdparagraph:P570:begin
\textbf{Hint:} ofslēan.
% mdparagraph:P570:end

% mdparagraph:P571:begin
\textbf{[seek out]} Use the bracketed English as the explicit target wording, preserving the complete sentence and every participant: seek out.
% mdparagraph:P571:end

% mdparagraph:P572:begin
\textbf{(195)} Finde is second-singular present subjunctive in the relative under the optative/\allowbreak{}hypothetical sīe hē ofslægen. Section 195 describes this subjunctive attraction; it does not require subjunctive in every relative clause.
% mdparagraph:P572:end

% output: EX-XII-XIII-47-03-O001
\end{minipage}\par\addvspace{4pt}

\par\begin{minipage}{\linewidth}\raggedright
% mdparagraph:P573:begin
\subsection*{Cook 47.4}\BookHeading{nav-ch18-learner-0124}{subsubsection}{Cook 47.4}{0}
% mdparagraph:P573:end

% item: EX-XII-XIII-47-04
% mdparagraph:P574:begin
Then said he, With whomsoever I find it, let him be my servant, and let the others be blameless (clǣne).
% mdparagraph:P574:end

% mdparagraph:P575:begin
\textbf{Hint:} clǣne.
% mdparagraph:P575:end

\end{minipage}\par\addvspace{4pt}

% teaching: NEW-U18-GRAM-U18-H-47-04
% mdparagraph:P576:begin
\textbf{Supplied:} Mid swā hwǣm swā ic hit finde, sīe hē mīn ðēow, and sīen ðā ōðre clǣne. — The free-choice dative person is governed by mid within the indefinite condition; hē resumes that person as subject of the singular wish. The other people are plural subject of sīen. Clǣne means clear of blame here. Preserve the scope contrast between the one offender and the other people.
% mdparagraph:P576:end

% output: EX-XII-XIII-47-04-O001
\par\begin{minipage}{\linewidth}\raggedright
% mdparagraph:P577:begin
\subsection*{Cook 47.5}\BookHeading{nav-ch18-learner-0125}{subsubsection}{Cook 47.5}{0}
% mdparagraph:P577:end

% item: EX-XII-XIII-47-05
% mdparagraph:P578:begin
And he sought from (fram) the eldest unto [until] the youngest; and he found the cup in Benjamin’s sack; and they turned back to the city.
% mdparagraph:P578:end

% mdparagraph:P579:begin
\textbf{Hint:} fram.
% mdparagraph:P579:end

\end{minipage}\par\addvspace{4pt}

% teaching: U18-SUP-OLDEST-DAT-M-SG-SUBSTANTIVE
% mdparagraph:P580:begin
\textbf{Supplied:} ðǣm ieldestan — The youngest and eldest refer to individual brothers. Supply the masculine singular human referents and each case, rather than inferring substantive person uses from familiar adjective headwords.
% mdparagraph:P580:end

% mdparagraph:P581:begin
\textbf{Supplied:} ðone gingstan; — The youngest and eldest refer to individual brothers. Supply the masculine singular human referents and each case, rather than inferring substantive person uses from familiar adjective headwords.
% mdparagraph:P581:end

% mdparagraph:P582:begin
\textbf{Supplied:} ieldestan — oldest person; eldest brother; dative, masculine, singular.
% mdparagraph:P582:end

% mdparagraph:P583:begin
\textbf{Supplied:} gingstan — youngest; accusative, masculine, singular.
% mdparagraph:P583:end

% teaching: U18-SUP-BENIAMIN-GEN
% mdparagraph:P584:begin
\textbf{Supplied:} Beniamines — Benjamin; genitive, masculine, singular.
% mdparagraph:P584:end

% mdparagraph:P585:begin
\textbf{[until]} Use the bracketed English as the explicit target wording, preserving the complete sentence and every participant: until.
% mdparagraph:P585:end

% output: EX-XII-XIII-47-05-O001
\par\begin{minipage}{\linewidth}\raggedright
% mdparagraph:P586:begin
\subsection*{Cook 47.6}\BookHeading{nav-ch18-learner-0126}{subsubsection}{Cook 47.6}{0}
% mdparagraph:P586:end

% item: EX-XII-XIII-47-06
% mdparagraph:P587:begin
And Judah went in first (fyrmest) with his brethren to Joseph, and they fell all together before Joseph.
% mdparagraph:P587:end

% mdparagraph:P588:begin
\textbf{Hint:} fyrmest.
% mdparagraph:P588:end

% mdparagraph:P589:begin
\textbf{Supplied:} Iudas — Judah; nominative, masculine, singular.
% mdparagraph:P589:end

\end{minipage}\par\addvspace{4pt}

% teaching: U18-SUP-FYRMEST-INVARIANT
% mdparagraph:P590:begin
\textbf{Supplied:} fyrmest — first of all; at the outset.
% mdparagraph:P590:end

% mdparagraph:P591:begin
\textbf{Supplied:} Iosepe — Joseph; dative, masculine, singular.
% mdparagraph:P591:end

% mdparagraph:P592:begin
\textbf{Supplied:} ætgædere — together; at once.
% mdparagraph:P592:end

% mdparagraph:P593:begin
\textbf{Supplied:} Iosepe — Joseph; dative, masculine, singular.
% mdparagraph:P593:end

% output: EX-XII-XIII-47-06-O001
\par\begin{minipage}{\linewidth}\raggedright
% mdparagraph:P594:begin
\subsection*{Cook 47.7}\BookHeading{nav-ch18-learner-0127}{subsubsection}{Cook 47.7}{0}
% mdparagraph:P594:end

% item: EX-XII-XIII-47-07
% mdparagraph:P595:begin
And he said to (tō) them, Why would ye do so (verb last)?
% mdparagraph:P595:end

% mdparagraph:P596:begin
\textbf{Hint:} tō.
% mdparagraph:P596:end

% mdparagraph:P597:begin
\textbf{(verb last)} Put the finite verb last in this specified clause. This local composition target is not a universal Old English order rule.
% mdparagraph:P597:end

% output: EX-XII-XIII-47-07-O001
\end{minipage}\par\addvspace{4pt}

\par\begin{minipage}{\linewidth}\raggedright
% mdparagraph:P598:begin
\subsection*{Cook 47.8}\BookHeading{nav-ch18-learner-0128}{subsubsection}{Cook 47.8}{0}
% mdparagraph:P598:end

% item: EX-XII-XIII-47-08
% mdparagraph:P599:begin
Let him be my servant who stole the cup, and go ye free to your father.
% mdparagraph:P599:end

\end{minipage}\par\addvspace{4pt}

% teaching: U18-SUP-FRIGE-NOM-M-PL
% mdparagraph:P600:begin
\textbf{Supplied:} frige — free; nominative, masculine, plural.
% mdparagraph:P600:end

% output: EX-XII-XIII-47-08-O001
\par\begin{minipage}{\linewidth}\raggedright
% mdparagraph:P601:begin
\subsection*{Cook 47.9}\BookHeading{nav-ch18-learner-0129}{subsubsection}{Cook 47.9}{0}
% mdparagraph:P601:end

% item: EX-XII-XIII-47-09
% mdparagraph:P602:begin
Then said Judah, Thou askedst us before whether we had father or brother.
% mdparagraph:P602:end

% mdparagraph:P603:begin
\textbf{Supplied:} Iudas — Judah; nominative, masculine, singular.
% mdparagraph:P603:end

% output: EX-XII-XIII-47-09-O001
\end{minipage}\par\addvspace{4pt}

\par\begin{minipage}{\linewidth}\raggedright
% mdparagraph:P604:begin
\subsection*{Cook 47.10}\BookHeading{nav-ch18-learner-0130}{subsubsection}{Cook 47.10}{0}
% mdparagraph:P604:end

% item: EX-XII-XIII-47-10
% mdparagraph:P605:begin
And we said, We have (om. an) old father, and he hath with him our youngest brother, and he loveth him alone (ænne) above (ofer) us all.
% mdparagraph:P605:end

% mdparagraph:P606:begin
\textbf{Hint:} ænne.
% mdparagraph:P606:end

% mdparagraph:P607:begin
\textbf{Hint:} ofer.
% mdparagraph:P607:end

\end{minipage}\par\addvspace{4pt}

% teaching: NEW-U18-GRAM-U18-HELP-OFER-PREFERENCE
% teaching: U18-SUP-PREFERENCE
% mdparagraph:P608:begin
\textbf{Supplied:} ofer ūs ealle. — Above us all /\allowbreak{} more than all of us, describing the father’s preference. Ūs and ealle are accusative plural.
% mdparagraph:P608:end

% mdparagraph:P609:begin
\textbf{Supplied:} gingstan — youngest; accusative, masculine, singular.
% mdparagraph:P609:end

% teaching: U18-SUP-AENNE-ACC-M-SG-SOURCE
% mdparagraph:P610:begin
\textbf{Supplied:} ænne — one; single; accusative, masculine, singular.
% mdparagraph:P610:end

% mdparagraph:P611:begin
\textbf{Supplied:} ofer — above; more than, in preference; governs accusative.
% mdparagraph:P611:end

% mdparagraph:P612:begin
\textbf{(om. an)} Omit a separate Old English word for the marked English expression an. Preserve its grammatical relation through the surrounding selected construction.
% mdparagraph:P612:end

% output: EX-XII-XIII-47-10-O001
\par\begin{minipage}{\linewidth}\raggedright
% mdparagraph:P613:begin
\subsection*{Cook 47.11}\BookHeading{nav-ch18-learner-0131}{subsubsection}{Cook 47.11}{0}
% mdparagraph:P613:end

% item: EX-XII-XIII-47-11
% mdparagraph:P614:begin
And thou commandest us that we (om. should) bring him to thee, that thou (om. mightest) see him, and know by (be) that that we were not spies (scēawere, 44).
% mdparagraph:P614:end

% mdparagraph:P615:begin
\textbf{Hint:} be.
% mdparagraph:P615:end

% mdparagraph:P616:begin
\textbf{Hint:} scēawere.
% mdparagraph:P616:end

\end{minipage}\par\addvspace{4pt}

% teaching: NEW-U18-GRAM-U18-H-47-11
% teaching: U18-SUP-INFERENCE
% mdparagraph:P617:begin
\textbf{Supplied:} ðū bīetst ūs ðæt wē hine ðē bringen, ðæt ðū hine gesēo, and wite be ðǣm ðæt wē ne wǣren scēaweras. — Retain actual present commandest: bīetst is present second singular, not an editorial backshift. The commanded bringing has first-plural present subjunctive; seeing and knowing are coordinated second-singular present subjunctives under purpose. The final reported status preserves original were as first-plural past subjunctive. Be ðǣm means by that evidence, and its following ðæt introduces the content known. The complete tense and proof chain is supplied.
% mdparagraph:P617:end

% mdparagraph:P618:begin
\textbf{Supplied:} be — by/\allowbreak{}from that; on that evidence.
% mdparagraph:P618:end

% mdparagraph:P619:begin
\textbf{(om. should)} Omit a separate Old English word for the marked English expression should. Preserve its grammatical relation through the surrounding selected construction.
% mdparagraph:P619:end

% mdparagraph:P620:begin
\textbf{(om. mightest)} Omit a separate Old English word for the marked English expression mightest. Preserve its grammatical relation through the surrounding selected construction.
% mdparagraph:P620:end

% mdparagraph:P621:begin
\textbf{(scēawere, 44)} The number refers to the noun’s declension. Use the complete scēawere table in section 2 and its stated stem rule.
% mdparagraph:P621:end

% output: EX-XII-XIII-47-11-O001
\reserve{7\baselineskip}
% mdparagraph:P622:begin
\section*{Cook 48. Recognition and invitation}\BookHeading{nav-ch18-learner-0132}{subsection}{Cook 48. Recognition and invitation}{1}
% mdparagraph:P622:end

\par\begin{minipage}{\linewidth}\raggedright
% mdparagraph:P623:begin
\subsection*{Cook 48.1}\BookHeading{nav-ch18-learner-0133}{subsubsection}{Cook 48.1}{0}
% mdparagraph:P623:end

% item: EX-XII-XIII-48-01
% mdparagraph:P624:begin
And I swore oaths (om. to) my father that I (om. would) bring him home again to (tō) him, and said, Except I bring him again back to thee, let me be ever guilty (scyldig) toward (wið) thee.
% mdparagraph:P624:end

% mdparagraph:P625:begin
\textbf{Hint:} tō.
% mdparagraph:P625:end

% mdparagraph:P626:begin
\textbf{Hint:} scyldig.
% mdparagraph:P626:end

% mdparagraph:P627:begin
\textbf{Hint:} wið.
% mdparagraph:P627:end

% mdparagraph:P628:begin
\textbf{Supplied:} Būtan — unless; except if.
% mdparagraph:P628:end

% mdparagraph:P629:begin
\textbf{Supplied:} ā — ever; always.
% mdparagraph:P629:end

% mdparagraph:P630:begin
\textbf{Supplied:} wið — toward; in relation to a person.
% mdparagraph:P630:end

% mdparagraph:P631:begin
\textbf{(om. to)} Omit a separate Old English word for the marked English expression to. Preserve its grammatical relation through the surrounding selected construction.
% mdparagraph:P631:end

% mdparagraph:P632:begin
\textbf{(om. would)} Omit a separate Old English word for the marked English expression would. Preserve its grammatical relation through the surrounding selected construction.
% mdparagraph:P632:end

% output: EX-XII-XIII-48-01-O001
\end{minipage}\par\addvspace{4pt}

\par\begin{minipage}{\linewidth}\raggedright
% mdparagraph:P633:begin
\subsection*{Cook 48.2}\BookHeading{nav-ch18-learner-0134}{subsubsection}{Cook 48.2}{0}
% mdparagraph:P633:end

% item: EX-XII-XIII-48-02
% mdparagraph:P634:begin
My lord, let (lǣtan) the lad go home with his brethren, and (om. let) me be thy servant instead of (for, 172) him.
% mdparagraph:P634:end

% mdparagraph:P635:begin
\textbf{Hint:} lǣtan.
% mdparagraph:P635:end

% mdparagraph:P636:begin
\textbf{Hint:} for.
% mdparagraph:P636:end

\end{minipage}\par\addvspace{4pt}

% teaching: NEW-U18-GRAM-U18-H-48-02
% mdparagraph:P637:begin
\textbf{Supplied:} lǣt ðone cnapan mid his brōðrum hām gān, — Complete permission frame: singular lǣt addresses Joseph; accusative ðone cnapan is the actor of plain infinitive gān. The boy goes home with his brothers. The accusative actor is not the recipient of a finite command and gān has no person ending. This local frame does not infer unrestricted causative/\allowbreak{}permissive infinitives from the familiar verb alone.
% mdparagraph:P637:end

% teaching: U18-SUP-SUBSTITUTION
% mdparagraph:P638:begin
\textbf{Supplied:} for — instead of; in place of.
% mdparagraph:P638:end

% mdparagraph:P639:begin
\textbf{(om. let)} Omit a separate Old English word for the marked English expression let. Preserve its grammatical relation through the surrounding selected construction.
% mdparagraph:P639:end

% mdparagraph:P640:begin
\textbf{(for, 172)} For plus accusative marks substitution: for hine means instead of him/\allowbreak{}in his place. It is distinct from the causal for + dative used of the money in 45.6.
% mdparagraph:P640:end

% output: EX-XII-XIII-48-02-O001
\par\begin{minipage}{\linewidth}\raggedright
% mdparagraph:P641:begin
\subsection*{Cook 48.3}\BookHeading{nav-ch18-learner-0135}{subsubsection}{Cook 48.3}{0}
% mdparagraph:P641:end

% item: EX-XII-XIII-48-03
% mdparagraph:P642:begin
I dare not go home without the lad, lest I see my father’s grief.
% mdparagraph:P642:end

\end{minipage}\par\addvspace{4pt}

% teaching: NEW-U18-GRAM-U18-H-48-03
% mdparagraph:P643:begin
\textbf{Supplied:} Ic ne dearr hām gān būtan ðǣm cnapan, ðȳ-lǣs-ðe ic mīnes fæder sārnesse gesēo. — Negative present dare plus plain go has the same speaker subject. Būtan with dative denotes absence of the boy, not the unless conjunction used elsewhere. The following negative-purpose clause lest I see uses first-singular present subjunctive gesēo; it does not assert that the grief was seen. Father is genitive possessor within accusative grief.
% mdparagraph:P643:end

% output: EX-XII-XIII-48-03-O001
\par\begin{minipage}{\linewidth}\raggedright
% mdparagraph:P644:begin
\subsection*{Cook 48.4}\BookHeading{nav-ch18-learner-0136}{subsubsection}{Cook 48.4}{0}
% mdparagraph:P644:end

% item: EX-XII-XIII-48-04
% mdparagraph:P645:begin
Then could (mugan) not Joseph conceal (diernan, 113) himself longer, but he drove all the others out, that no foreign man (om. should) be among them.
% mdparagraph:P645:end

% mdparagraph:P646:begin
\textbf{Hint:} mugan.
% mdparagraph:P646:end

% mdparagraph:P647:begin
\textbf{Hint:} diernan.
% mdparagraph:P647:end

% mdparagraph:P648:begin
\textbf{Supplied:} Iosep — Joseph; nominative, masculine, singular.
% mdparagraph:P648:end

\end{minipage}\par\addvspace{4pt}

% teaching: U18-SUP-LENG-INVARIANT
% mdparagraph:P649:begin
\textbf{Supplied:} lęng — longer.
% mdparagraph:P649:end

% mdparagraph:P650:begin
\textbf{Supplied:} betweox — among; between.
% mdparagraph:P650:end

% mdparagraph:P651:begin
\textbf{(diernan, 113)} Diernan is a long weak class I verb; the assigned system here is the complete plain/\allowbreak{}inflected infinitive pair diernan/\allowbreak{}diernanne, not a silently expanded finite paradigm. The plain infinitive follows past ability meahte, and accusative hine is reflexive.
% mdparagraph:P651:end

% mdparagraph:P652:begin
\textbf{(om. should)} Omit a separate Old English word for the marked English expression should. Preserve its grammatical relation through the surrounding selected construction.
% mdparagraph:P652:end

% output: EX-XII-XIII-48-04-O001
\par\begin{minipage}{\linewidth}\raggedright
% mdparagraph:P653:begin
\subsection*{Cook 48.5}\BookHeading{nav-ch18-learner-0137}{subsubsection}{Cook 48.5}{0}
% mdparagraph:P653:end

% item: EX-XII-XIII-48-05
% mdparagraph:P654:begin
And he wept and cried (clēopian) with a loud (hlūd) voice (160. 1); and he said to (tō) his brethren, I am Joseph; (om. doth) our father live yet (nū gīet)?
% mdparagraph:P654:end

% mdparagraph:P655:begin
\textbf{Hint:} clēopian.
% mdparagraph:P655:end

% mdparagraph:P656:begin
\textbf{Hint:} hlūd.
% mdparagraph:P656:end

% mdparagraph:P657:begin
\textbf{Hint:} tō.
% mdparagraph:P657:end

% mdparagraph:P658:begin
\textbf{Hint:} nū gīet.
% mdparagraph:P658:end

\end{minipage}\par\addvspace{4pt}

% teaching: NEW-U18-GRAM-U18-H-48-05
% mdparagraph:P659:begin
\textbf{Supplied:} clēopode hlūdre stefne; — Exact 48.5 hint has marked clēopian. Supply the derived weak class II past third singular clēopode and the dative/\allowbreak{}instrumental voice phrase meaning cried with a loud voice. This marked source variant is supplied locally; retain the earlier spelling in its own paradigm.
% mdparagraph:P659:end

% mdparagraph:P660:begin
\textbf{Supplied:} clēopode — cry aloud; indicative, preterite, third person, singular.
% mdparagraph:P660:end

% mdparagraph:P661:begin
\textbf{Supplied:} Iosep — Joseph; nominative, masculine, singular.
% mdparagraph:P661:end

% mdparagraph:P662:begin
\textbf{(160. 1)} Hlūdre stefne is a dative/\allowbreak{}instrumental manner phrase, with a loud voice; no extra mid is required. The feminine adjective agrees with singular stefne.
% mdparagraph:P662:end

% mdparagraph:P663:begin
\textbf{(om. doth)} Omit a separate Old English word for the marked English expression doth. Preserve its grammatical relation through the surrounding selected construction.
% mdparagraph:P663:end

% output: EX-XII-XIII-48-05-O001
\par\begin{minipage}{\linewidth}\raggedright
% mdparagraph:P664:begin
\subsection*{Cook 48.6}\BookHeading{nav-ch18-learner-0138}{subsubsection}{Cook 48.6}{0}
% mdparagraph:P664:end

% item: EX-XII-XIII-48-06
% mdparagraph:P665:begin
Fear ye (184) nothing, nor (om. let) it grieve (190) you that ye sold (sęllan) me into this kingdom; indeed (om. it was) for your benefit God sent me into the land of the Egyptians.
% mdparagraph:P665:end

% mdparagraph:P666:begin
\textbf{Hint:} sęllan.
% mdparagraph:P666:end

\end{minipage}\par\addvspace{4pt}

% teaching: NEW-U18-GRAM-U18-HELP-FOR-BENEFIT
% teaching: U18-SUP-BENEFIT
% mdparagraph:P667:begin
\textbf{Supplied:} for ēowerre ðearfe — For your good /\allowbreak{} for your benefit. Ēowerre and ðearfe are dative feminine singular. Use the whole phrase visibly.
% mdparagraph:P667:end

% teaching: U18-SUP-SELLING-GOAL
% mdparagraph:P668:begin
\textbf{Supplied:} on — into; destination of sale.
% mdparagraph:P668:end

% mdparagraph:P669:begin
\textbf{Supplied:} for — for; to the benefit of; governs dative.
% mdparagraph:P669:end

% mdparagraph:P670:begin
\textbf{Supplied:} Egypta — Egyptians; genitive, masculine, plural.
% mdparagraph:P670:end

% mdparagraph:P671:begin
\textbf{(184)} Reflexive dative ēow belongs with fear; nānes ðinges is the separately licensed genitive feared object. The negation is concordant: do not fear anything, not an affirmative double-negative reading.
% mdparagraph:P671:end

% mdparagraph:P672:begin
\textbf{(om. let)} Omit a separate Old English word for the marked English expression let. Preserve its grammatical relation through the surrounding selected construction.
% mdparagraph:P672:end

% mdparagraph:P673:begin
\textbf{(190)} Ofðynce is third-singular impersonal subjunctive, let it not grieve/\allowbreak{}displease you. Dative ēow gives the experiencers; the following ðæt clause gives what troubles them. Plural experiencers do not cause plural verbal agreement.
% mdparagraph:P673:end

% mdparagraph:P674:begin
\textbf{(om. it was)} Omit a separate Old English word for the marked English expression it was. Preserve its grammatical relation through the surrounding selected construction.
% mdparagraph:P674:end

% output: EX-XII-XIII-48-06-O001
\par\begin{minipage}{\linewidth}\raggedright
% mdparagraph:P675:begin
\subsection*{Cook 48.7}\BookHeading{nav-ch18-learner-0139}{subsubsection}{Cook 48.7}{0}
% mdparagraph:P675:end

% item: EX-XII-XIII-48-07
% mdparagraph:P676:begin
Now two years has been [was] famine over (172) all (om. the) earth, and (om. there) yet shall (om. be) five, in which one can (mugan) neither (89. a) plow (ęrian, 116) nor reap (rīpan, 102).
% mdparagraph:P676:end

% mdparagraph:P677:begin
\textbf{Hint:} mugan.
% mdparagraph:P677:end

% mdparagraph:P678:begin
\textbf{Hint:} ęrian.
% mdparagraph:P678:end

% mdparagraph:P679:begin
\textbf{Hint:} rīpan.
% mdparagraph:P679:end

\end{minipage}\par\addvspace{4pt}

% teaching: NEW-U18-GRAM-U18-H-48-07
% teaching: NEW-U18-VOC-SCULAN-U18-FUTURE-ELLIPSIS
% teaching: U18-SUP-NAHWAETHER-INVARIANT
% mdparagraph:P680:begin
\textbf{Supplied:} Nū twā gēar hunger wæs ofer ealle eorðan, and gīet sculon fīf, on ðǣm mann nāhwæðer ne ęrian ne rīpan mæg. — Two years is accusative duration; singular hunger governs wæs. Ofer takes accusative extent ealle eorðan, as 172 requires. Fīf stands for five further years; the source explicitly omits there and be, so sculon has an understood existential complement, not a missing translation item. On ðǣm refers to those years. Indefinite mann is singular subject of mæg; nāhwæðer plus ne...ne excludes both infinitive activities. Cook 89 a supplies the full pronoun; All ellipsis, future use and negative coordination are supplied together.
% mdparagraph:P680:end

% mdparagraph:P681:begin
\textbf{Supplied:} sculon — shall; be yet to come (with an omitted be); indicative, present, third person, plural.
% mdparagraph:P681:end

% mdparagraph:P682:begin
\textbf{Supplied:} nāhwæðer — neither.
% mdparagraph:P682:end

% mdparagraph:P683:begin
\textbf{[was]} Use the bracketed English as the explicit target wording, preserving the complete sentence and every participant: was.
% mdparagraph:P683:end

% mdparagraph:P684:begin
\textbf{(172)} Ofer with accusative ealle eorðan marks extent over all the earth. The noun’s identical weak ending does not make it a dative; the full phrase has the accusative relation required by 172.
% mdparagraph:P684:end

% mdparagraph:P685:begin
\textbf{(om. the)} Omit a separate Old English word for the marked English expression the. Preserve its grammatical relation through the surrounding selected construction.
% mdparagraph:P685:end

% mdparagraph:P686:begin
\textbf{(om. there)} Omit a separate Old English word for the marked English expression there. Preserve its grammatical relation through the surrounding selected construction.
% mdparagraph:P686:end

% mdparagraph:P687:begin
\textbf{(om. be)} Omit a separate Old English word for the marked English expression be. Preserve its grammatical relation through the surrounding selected construction.
% mdparagraph:P687:end

% mdparagraph:P688:begin
\textbf{(89. a)} Nāhwæðer means neither. In the supplied whole clause, nāhwæðer with paired ne ... ne excludes plowing and reaping under singular mæg; mann is the singular indefinite subject. The lexical hint 89 a alone does not teach the coordination, so the complete negative frame is visible.
% mdparagraph:P688:end

% mdparagraph:P689:begin
\textbf{(ęrian, 116)} The source spells ęrian, with ogonek, and assigns short weak class I 116. The complete assigned nonfinite pair is ęrian/\allowbreak{}ęrianne; only plain ęrian occurs after mæg. No new full finite weak class I system is inferred.
% mdparagraph:P689:end

% mdparagraph:P690:begin
\textbf{(rīpan, 102)} Rīpan is strongI 102. The complete assigned infinitive pair rīpan/\allowbreak{}rīpanne preserves the source’s long ī; only plain rīpan occurs after mæg. A class reference does not itself assign all finite forms.
% mdparagraph:P690:end

% output: EX-XII-XIII-48-07-O001
\par\begin{minipage}{\linewidth}\raggedright
% mdparagraph:P691:begin
\subsection*{Cook 48.8}\BookHeading{nav-ch18-learner-0140}{subsubsection}{Cook 48.8}{0}
% mdparagraph:P691:end

% item: EX-XII-XIII-48-08
% mdparagraph:P692:begin
Go quickly to my father and tell him that God has set me as a (tō) lord (om. to) all (om. the) Egyptians; bid (bēodan) him (164) that he come to me and dwell in this land.
% mdparagraph:P692:end

% mdparagraph:P693:begin
\textbf{Hint:} tō.
% mdparagraph:P693:end

% mdparagraph:P694:begin
\textbf{Hint:} bēodan.
% mdparagraph:P694:end

\end{minipage}\par\addvspace{4pt}

% teaching: NEW-U18-GRAM-U18-H-48-08
% teaching: NEW-U18-VOC-GESETTAN-U18-APPOINT-LORD
% teaching: U18-SUP-EGYPTUM-DAT-PL
% mdparagraph:P695:begin
\textbf{Supplied:} God mē tō hlāforde eallum Egyptum gesętte; — God appointed me as lord to all Egyptians. Mē is the accusative appointee, tō hlāforde the role complement, and eallum Egyptum the dative people governed; use the bare dative for the people, following the instruction to omit a separate equivalent of English to. The exact name phrase is supplied.
% mdparagraph:P695:end

% mdparagraph:P696:begin
\textbf{Supplied:} hrædlīce — quickly.
% mdparagraph:P696:end

% mdparagraph:P697:begin
\textbf{Supplied:} Egyptum — Egyptians; dative, masculine, plural.
% mdparagraph:P697:end

% mdparagraph:P698:begin
\textbf{Supplied:} gesętte — appoint as lord; indicative, preterite, third person, singular.
% mdparagraph:P698:end

% mdparagraph:P699:begin
\textbf{(om. to)} Omit a separate Old English word for the marked English expression to. Preserve its grammatical relation through the surrounding selected construction.
% mdparagraph:P699:end

% mdparagraph:P700:begin
\textbf{(om. the)} Omit a separate Old English word for the marked English expression the. Preserve its grammatical relation through the surrounding selected construction.
% mdparagraph:P700:end

% mdparagraph:P701:begin
\textbf{(164)} Him is the dative addressee of plural command bēodað; the father is subject of cume and wunie in the commanded ðæt clause. Keep the commanding brothers distinct from the father who is to come and dwell.
% mdparagraph:P701:end

% output: EX-XII-XIII-48-08-O001
\par\begin{minipage}{\linewidth}\raggedright
% mdparagraph:P702:begin
\subsection*{Cook 48.9}\BookHeading{nav-ch18-learner-0141}{subsubsection}{Cook 48.9}{0}
% mdparagraph:P702:end

% item: EX-XII-XIII-48-09
% mdparagraph:P703:begin
Do thus, that ye perish not.
% mdparagraph:P703:end

% output: EX-XII-XIII-48-09-O001
\end{minipage}\par\addvspace{4pt}

\par\begin{minipage}{\linewidth}\raggedright
% mdparagraph:P704:begin
\subsection*{Cook 48.10}\BookHeading{nav-ch18-learner-0142}{subsubsection}{Cook 48.10}{0}
% mdparagraph:P704:end

% item: EX-XII-XIII-48-10
% mdparagraph:P705:begin
Now ye see how it is with me, and ye hear what I say (sprecan) to (tō) you.
% mdparagraph:P705:end

% mdparagraph:P706:begin
\textbf{Hint:} sprecan.
% mdparagraph:P706:end

% mdparagraph:P707:begin
\textbf{Hint:} tō.
% mdparagraph:P707:end

\end{minipage}\par\addvspace{4pt}

% teaching: NEW-U18-GRAM-U18-H-48-10
% teaching: NEW-U18-VOC-GESEON-U18-STATE-CONTENT
% mdparagraph:P708:begin
\textbf{Supplied:} Nū gē gesēoð hū hit mid mē is, and gē hīerað hwæt ic tō ēow sprece. — Hū introduces how things stand with the speaker; this is a state-of-affairs expression, not physical accompaniment. Hwæt is what he says, heard by plural gē. Sprece retains the first-singular speaker, and tō ēow the plural addressee.
% mdparagraph:P708:end

% mdparagraph:P709:begin
\textbf{Supplied:} gesēoð — see how a state of affairs stands; indicative, present, second person, plural.
% mdparagraph:P709:end

% output: EX-XII-XIII-48-10-O001
\reserve{7\baselineskip}
% mdparagraph:P710:begin
\section*{7. The gifts and the end of the banquet}\BookHeading{nav-ch18-learner-0143}{subsection}{7. The gifts and the end of the banquet}{1}
% mdparagraph:P710:end

% mdparagraph:P711:begin
The king’s daughter has heard and watched Apollonius. The continuation moves through her gifts, their delivery and the farewells. The exact passage is Cook 175.18–177.2, ending before the following teacher scene. Read it with the original notes and modern help below.
% mdparagraph:P711:end

% mdparagraph:P712:begin
Sōðlice, mid-ðȳ-ðe ðæs cyninges dohtor geseah ðæt Apollonius on eallum gōdum cræftum swā wel wæs getogen,\textsuperscript{P175N10} ðā gefēoll hiere mōd on his lufe. Ðā, æfter ðæs bēorscipes geęndunge, cwæð ðæt mǣden tō ðǣm cyninge: “Lēofa\textsuperscript{P175N11} fæder, ðū līefdest mē, lȳtle\textsuperscript{P176N1} ǣr,\textsuperscript{P176N1} ðæt\textsuperscript{P176N2} ic\textsuperscript{P176N2} mōste\textsuperscript{P176N2} giefan Apollonio swā-hwæt-swā ic wolde of ðīnum goldhorde.” Arcestrates se cyning cwæð tō hiere: “Gief him swā-hwæt-swā ðū wille.\textsuperscript{P176N3}” Hēo ðā swīðe blīðe ūtēode,\textsuperscript{P176N4} and cwæð: “Lārēow\textsuperscript{P176N5} Apolloni, ic giefe ðē, be mīnes fæder lēafe, twā hund punda\textsuperscript{P176N6} goldes,\textsuperscript{P176N7} and fēower hund punda\textsuperscript{P176N6} gewihte\textsuperscript{P176N8} seolfres,\textsuperscript{P176N7} and ðone mǣstan dǣl dēorwurðes\textsuperscript{P176N9} rēafes, and twēntig ðēowa\textsuperscript{P176N10} manna.\textsuperscript{P176N10}” And hēo ðā ðus cwæð tō ðǣm ðēowum mannum: “Berað ðās ðing mid ēow, ðe ic behēt Apollonio mīnum lārēowe, and lęcgeað innan būre\textsuperscript{P176N11} beforan mīnum frēondum.” Ðis wearð ðā ðus gedōn, æfter ðǣre cwēne hǣse\textsuperscript{P176N12}; and ealle ðā męnn hiere giefa hęredon ðe\textsuperscript{P176N13} hīe gesāwon. Ðā sōðlice geęndode se gebēorscipe, and ðā męnn ealle ārison,\textsuperscript{P176N14} and grētton ðone cyning and ðā cwēne, and bǣdon hīe gesunde bēon,\textsuperscript{P176N15} and hām gewęndon. Ēac swilce Apollonius cwæð: “Ðū gōda cyning, and earmra\textsuperscript{P176N16} gemiltsiend, and ðū cwēn, lāre\textsuperscript{P176N16} lufiend, bēon gē gesunde.\textsuperscript{P176N17}” Hē beseah ēac tō ðǣm ðēowum mannum, ðe ðæt mǣden him forgiefen hæfde,\textsuperscript{P176N18} and him cwæð tō: “Nimað ðās ðing mid ēow, ðe mē sēo cwēn forgeaf, and gān\textsuperscript{P177N1} wē sēcean ūre giesthūs, ðæt wē mægen ūs\textsuperscript{P177N2} geręstan.”
% mdparagraph:P712:end

\reserve{18\baselineskip}
% mdparagraph:P713:begin
\subsection*{Finding the cited lines}\BookHeading{nav-ch18-learner-0144}{subsubsection}{Finding the cited lines}{0}
% mdparagraph:P713:end

% mdparagraph:P714:begin
These locators retain the source’s physical line divisions, including divided words. The running text joins those words.
% mdparagraph:P714:end

% mdtable:T65:begin
\par\begingroup
\fontsize{10.5}{13.2}\selectfont
\setlength{\LTleft}{0pt}\setlength{\LTright}{\fill}
\begin{longtable}{@{}>{\raggedright\arraybackslash}p{35.03mm}>{\raggedright\arraybackslash}p{86.16mm}@{}}
\toprule
\textbf{Page and line} & \textbf{Beginning of the printed line} \\
\midrule
\endfirsthead
\toprule
\textbf{Page and line} & \textbf{Beginning of the printed line} \\
\midrule
\endhead
\midrule
\multicolumn{2}{r}{\footnotesize Continued on the following page.}\\
\endfoot
\bottomrule
\endlastfoot
175.18 & Sōðlice, mid-ðȳ-ðe ðæs cyninges dohtor … \\
175.19 & lonius on eallum gōdum cræftum … \\
175.20 & ðā gefēoll hiere mōd on … \\
175.21 & geęndunge, cwæð ðæt mǣden tō … \\
176.1 & fæder, ðū līefdest mē, lȳtle … \\
176.2 & Apollonio swā-hwæt-swā ic wolde of … \\
176.3 & Arcestrates se cyning cwæð tō … \\
176.4 & hwæt-swā ðū wille.” Hēo ðā … \\
176.5 & cwæð: “Lārēow Apolloni, ic giefe … \\
176.6 & lēafe, twā hund punda goldes, … \\
176.7 & gewihte seolfres, and ðone mǣstan … \\
176.8 & rēafes, and twēntig ðēowa manna.” … \\
176.9 & cwæð tō ðǣm ðēowum mannum: … \\
176.10 & ēow, ðe ic behēt Apollonio … \\
176.11 & innan būre beforan mīnum frēondum.” … \\
176.12 & ðus gedōn, æfter ðǣre cwēne … \\
176.13 & hiere giefa hęredon ðe hīe … \\
176.14 & ode se gebēorscipe, and ðā … \\
176.15 & grētton ðone cyning and ðā … \\
176.16 & gesunde bēon, and hām gewęndon. … \\
176.17 & lonius cwæð: “Ðū gōda cyning, … \\
176.18 & and ðū cwēn, lāre lufiend, … \\
176.19 & beseah ēac tō ðǣm ðēowum … \\
176.20 & forgiefen hæfde, and him cwæð … \\
177.1 & ēow, ðe mē sēo cwēn … \\
177.2 & giesthūs, ðæt wē mægen ūs … \\
\end{longtable}
\par\endgroup

% mdtable:T65:end
\reserve{7\baselineskip}
% mdparagraph:P715:begin
\subsection*{Modern help for the reading}\BookHeading{nav-ch18-learner-0145}{subsubsection}{Modern help for the reading}{0}
% mdparagraph:P715:end

% mdparagraph:P716:begin
The explanations supply unfamiliar senses and whole constructions. Where an original note already identifies a referent, the later question asks for an explanation with that aid visible. Source spelling is retained.
% mdparagraph:P716:end

\reserve{7\baselineskip}
% mdparagraph:P717:begin
\subsection*{175.18–20: seeing and falling in love}\BookHeading{nav-ch18-learner-0146}{subsubsection}{175.18–20: seeing and falling in love}{0}
% mdparagraph:P717:end

% teaching: NEW-U18-GRAM-U18-RSUP-01
% mdparagraph:P718:begin
The daughter sees that a fact holds; indicative wæs belongs to the content clause. Wæs getogen means was educated or accomplished in good skills. In the following idiom, singular mōd is the subject of gefēoll, hiere its possessor, and on his lufe the state into which her mind falls: she falls in love with him.
% mdparagraph:P718:end

\reserve{7\baselineskip}
% mdparagraph:P719:begin
\subsection*{175.20–176.2: the earlier permission}\BookHeading{nav-ch18-learner-0147}{subsubsection}{175.20–176.2: the earlier permission}{0}
% mdparagraph:P719:end

% teaching: NEW-U18-GRAM-U18-RSUP-02
% mdparagraph:P720:begin
Līefdest means permitted, with mē as dative person permitted. Ðæt ic mōste giefan is naturally translated by permitted me to give. The free-choice object swā-hwæt-swā means whatever; wolde is first-singular preterite subjunctive and leaves giving understood. Of with dative identifies the treasury from which she may take the gift.
% mdparagraph:P720:end

\reserve{7\baselineskip}
% mdparagraph:P721:begin
\subsection*{176.3–4: renewed permission}\BookHeading{nav-ch18-learner-0148}{subsubsection}{176.3–4: renewed permission}{0}
% mdparagraph:P721:end

% teaching: NEW-U18-GRAM-U18-RSUP-03
% mdparagraph:P722:begin
Gief is a singular command to the daughter, with him as recipient. Wille is second-singular present subjunctive within whatever you wish. Its spelling also occurs in other persons and moods; this context fixes the use.
% mdparagraph:P722:end

\reserve{7\baselineskip}
% mdparagraph:P723:begin
\subsection*{176.4: going out}\BookHeading{nav-ch18-learner-0149}{subsubsection}{176.4: going out}{0}
% mdparagraph:P723:end

% teaching: NEW-U18-GRAM-U18-RSUP-04
% mdparagraph:P724:begin
Hēo … ūtēode means she went out. Swīðe blīðe means very gladly. The exact prefixed preterite is supplied.
% mdparagraph:P724:end

\reserve{7\baselineskip}
% mdparagraph:P725:begin
\subsection*{176.5–8: the gifts}\BookHeading{nav-ch18-learner-0150}{subsubsection}{176.5–8: the gifts}{0}
% mdparagraph:P725:end

% teaching: NEW-U18-GRAM-U18-RSUP-05
% mdparagraph:P726:begin
Giefe is first-singular present indicative; ðē is the recipient. Be mīnes fæder lēafe means with my father’s permission. The following quantities all name what she gives. Their complete weight and servant expressions are supplied below.
% mdparagraph:P726:end

\reserve{7\baselineskip}
% mdparagraph:P727:begin
\subsection*{176.9–11: the instructions}\BookHeading{nav-ch18-learner-0151}{subsubsection}{176.9–11: the instructions}{0}
% mdparagraph:P727:end

% teaching: NEW-U18-GRAM-U18-RSUP-06
% mdparagraph:P728:begin
Berað and lęcgeað are plural commands to the enslaved men: carry the gifts along and place them in the chamber before her friends. In the promise relative, behēt is first-singular preterite, ðe represents the promised things, and Apollonio mīnum lārēowe is the dative recipient. The source lęcgeað is imperative despite matching the newly learned indicative plural in spelling.
% mdparagraph:P728:end

\reserve{7\baselineskip}
% mdparagraph:P729:begin
\subsection*{176.12: the command is carried out}\BookHeading{nav-ch18-learner-0152}{subsubsection}{176.12: the command is carried out}{0}
% mdparagraph:P729:end

% teaching: NEW-U18-GRAM-U18-RSUP-07
% mdparagraph:P730:begin
Ðis wearð … gedōn means this was done. Gedōn is a past participle in the singular passive event, not a present subjunctive plural of the same spelling. The action follows the princess’s command.
% mdparagraph:P730:end

\reserve{7\baselineskip}
% mdparagraph:P731:begin
\subsection*{176.12–13: praise and a separated relative}\BookHeading{nav-ch18-learner-0153}{subsubsection}{176.12–13: praise and a separated relative}{0}
% mdparagraph:P731:end

% teaching: NEW-U18-GRAM-U18-RSUP-08
% mdparagraph:P732:begin
All the men who saw the gifts praised her gifts. Original note 13 explicitly identifies męnn as antecedent of ðe. The relative is subject of its clause; hīe is accusative plural, the gifts. The relative stands after the main clause’s object, separated from its own antecedent. Both finite verbs have plural subjects.
% mdparagraph:P732:end

\reserve{7\baselineskip}
% mdparagraph:P733:begin
\subsection*{176.13–16: the guests leave}\BookHeading{nav-ch18-learner-0154}{subsubsection}{176.13–16: the guests leave}{0}
% mdparagraph:P733:end

% teaching: NEW-U18-GRAM-U18-RSUP-09
% mdparagraph:P734:begin
Geęndode is intransitive: the banquet ended. The men rose, took leave of king and princess, wished them well and returned home. In bǣdon hīe gesunde bēon, hīe is the accusative actor of bēon and gesunde its agreeing accusative plural predicate. This complete farewell wish is supplied.
% mdparagraph:P734:end

\reserve{7\baselineskip}
% mdparagraph:P735:begin
\subsection*{176.16–18: Apollonius’s farewell}\BookHeading{nav-ch18-learner-0155}{subsubsection}{176.16–18: Apollonius’s farewell}{0}
% mdparagraph:P735:end

% teaching: NEW-U18-GRAM-U18-RSUP-10
% mdparagraph:P736:begin
The king and princess are addressed separately and then resumed together by plural gē. Bēon is second-plural present subjunctive, may you be; gesunde is nominative plural agreeing with gē. The agent nouns and their genitive complements are supplied below.
% mdparagraph:P736:end

\reserve{7\baselineskip}
% mdparagraph:P737:begin
\subsection*{176.18–20: the servants}\BookHeading{nav-ch18-learner-0156}{subsubsection}{176.18–20: the servants}{0}
% mdparagraph:P737:end

% teaching: NEW-U18-GRAM-U18-RSUP-11
% mdparagraph:P738:begin
Beseah tō means looked toward, with a dative target. In the relative, the maiden is subject, the enslaved men the object represented by ðe, and singular him the recipient Apollonius. Forgiefen hæfde is active past perfect. The later him is plural and refers to the servants; tō follows the verb but still belongs to that dative addressee.
% mdparagraph:P738:end

\reserve{7\baselineskip}
% mdparagraph:P739:begin
\subsection*{177.1–2: the proposed departure}\BookHeading{nav-ch18-learner-0157}{subsubsection}{177.1–2: the proposed departure}{0}
% mdparagraph:P739:end

% teaching: NEW-U18-GRAM-U18-RSUP-12
% mdparagraph:P740:begin
Nimað is a plural command to take the gifts along. Gān wē is a first-plural subjunctive exhortation, let us go; supplied sēcean states the purpose of going. In the gift relative, the princess is subject and mē recipient. Final ðæt expresses purpose, with plural mægen and accusative reflexive ūs before geręstan. The words propose seeking the guesthouse and resting; they do not narrate a completed arrival.
% mdparagraph:P740:end

\reserve{7\baselineskip}
% mdparagraph:P741:begin
\subsection*{Supplied expressions and exact word help}\BookHeading{nav-ch18-learner-0158}{subsubsection}{Supplied expressions and exact word help}{0}
% mdparagraph:P741:end

% teaching: NEW-U18-GRAM-U18-HELP-LITTLE-EARLIER
% mdparagraph:P742:begin
\textbf{lȳtle ǣr} — A little earlier: lȳtle is instrumental of degree/\allowbreak{}difference with the comparative adverb ǣr. Cook’s original note points to 178. It is not a feminine accusative object.
% mdparagraph:P742:end

% teaching: NEW-U18-GRAM-U18-HELP-QUANTITIES
% mdparagraph:P743:begin
\textbf{twā hund punda goldes and fēower hund punda gewihte seolfres} — Two hundred pounds of gold and four hundred pounds by weight of silver. Hund means hundred. Punda is genitive plural after the hundreds; goldes and seolfres are material genitives. Gewihte is a dative expression of measure: by weight.
% mdparagraph:P743:end

% teaching: NEW-U18-GRAM-U18-HELP-PRECIOUS-CLOTHING
% mdparagraph:P744:begin
\textbf{ðone mǣstan dǣl dēorwurðes rēafes} — The greatest part of costly clothing. Mǣstan is weak accusative masculine singular with dǣl; dēorwurðes agrees with the genitive neuter singular rēafes.
% mdparagraph:P744:end

% teaching: NEW-U18-GRAM-U18-HELP-SERVANTS
% mdparagraph:P745:begin
\textbf{twēntig ðēowa manna /\allowbreak{} tō ðǣm ðēowum mannum} — Twenty enslaved men; to the enslaved men. Preserve ðēowa without r. The selected local explanation treats ðēowa and manna as appositional genitive plurals, and ðēowum and mannum as appositional datives. The phrase’s meaning and case are supplied without requiring independent historical classification.
% mdparagraph:P745:end

% teaching: NEW-U18-GRAM-U18-HELP-DELAYED-RELATIVE
% mdparagraph:P746:begin
\textbf{ealle ðā męnn hiere giefa hęredon ðe hīe gesāwon} — All the men who saw them praised her gifts. The relative ðe is subject and reaches back to męnn; hīe is the accusative plural gifts. The relative clause is separated from its antecedent. Keep the people and gifts distinct.
% mdparagraph:P746:end

% teaching: NEW-U18-GRAM-U18-HELP-FAREWELL-ACC
% mdparagraph:P747:begin
\textbf{bǣdon hīe gesunde bēon} — Wished them to be well: hīe refers to the king and princess and is the accusative actor of bēon; gesunde is accusative plural predicate. The complete construction is supplied.
% mdparagraph:P747:end

% teaching: NEW-U18-GRAM-U18-HELP-FAREWELL-NOM
% mdparagraph:P748:begin
\textbf{bēon gē gesunde} — May you be well: gē is nominative plural and gesunde nominative plural predicate. This is a direct wish with a plural subjunctive, distinct from the preceding accusative-infinitive construction despite identical gesunde.
% mdparagraph:P748:end

% teaching: NEW-U18-GRAM-U18-HELP-AGENT-ADDRESSES
% mdparagraph:P749:begin
\textbf{earmra gemiltsiend /\allowbreak{} lāre lufiend} — Pitier of the poor /\allowbreak{} lover of learning. Supply both agent nouns and their object genitives. These phrases address the king and princess respectively; historical or Latin parallels in the original notes are explained in English.
% mdparagraph:P749:end

% teaching: NEW-U18-GRAM-U18-HELP-SERVANT-RELATIVE
% mdparagraph:P750:begin
\textbf{ðēowum mannum ðe ðæt mǣden him forgiefen hæfde} — The enslaved men whom the girl had given to him. Relative ðe is accusative plural; mǣden is the subject; him is Apollonius.
% mdparagraph:P750:end

% teaching: NEW-U18-GRAM-U18-HELP-POSTPOSED-TO
% mdparagraph:P751:begin
\textbf{him cwæð tō} — Said to them: him is dative plural referring to the servants; tō follows the verb while governing him.
% mdparagraph:P751:end

% teaching: NEW-U18-GRAM-U18-HELP-FINAL-PLAN
% mdparagraph:P752:begin
\textbf{gān wē sēcean ūre giesthūs, ðæt wē mægen ūs geręstan} — Let us go seek our guesthouse, so that we may rest. Gān wē is an exhortation, ðæt marks purpose, and ūs is accusative reflexive with geręstan.
% mdparagraph:P752:end

% teaching: U18-SUP-APOLLONIUS-NOM
% mdparagraph:P753:begin
\textbf{Apollonius} — Apollonius; nominative, masculine, singular.
% mdparagraph:P753:end

% teaching: U18-SUP-APOLLONIUS-DAT
% mdparagraph:P754:begin
\textbf{Apollonio} — Apollonius; dative, masculine, singular.
% mdparagraph:P754:end

% teaching: U18-SUP-APOLLONIUS-VOC
% mdparagraph:P755:begin
\textbf{Apolloni} — Apollonius; direct address, masculine, singular.
% mdparagraph:P755:end

% teaching: U18-SUP-ARCESTRATES-NOM
% mdparagraph:P756:begin
\textbf{Arcestrates} — Arcestrates, the king; nominative, masculine, singular.
% mdparagraph:P756:end

% teaching: U18-SUP-LUFE-ACC-SG
% mdparagraph:P757:begin
\textbf{lufe} — love; accusative, feminine, singular.
% mdparagraph:P757:end

% teaching: U18-SUP-BEORSCIPES-GEN-SG
% mdparagraph:P758:begin
\textbf{bēorscipes} — feast; banquet; genitive, masculine, singular.
% mdparagraph:P758:end

% teaching: U18-SUP-GEENDUNGE-DAT-SG
% mdparagraph:P759:begin
\textbf{geęndunge} — end; close; dative, feminine, singular.
% mdparagraph:P759:end

% teaching: U18-SUP-LEAFE-DAT-SG
% mdparagraph:P760:begin
\textbf{lēafe} — permission; dative, feminine, singular.
% mdparagraph:P760:end

% teaching: U18-SUP-GEWIHTE-DAT-SG
% mdparagraph:P761:begin
\textbf{gewihte} — weight; dative, neuter, singular.
% mdparagraph:P761:end

% teaching: U18-SUP-DEORWURTHES-GEN-N-SG
% mdparagraph:P762:begin
\textbf{dēorwurðes} — precious; costly; genitive, neuter, singular.
% mdparagraph:P762:end

% teaching: U18-SUP-REAFES-GEN-SG
% mdparagraph:P763:begin
\textbf{rēafes} — clothing; garments; genitive, neuter, singular.
% mdparagraph:P763:end

% teaching: U18-SUP-BURE-DAT-SG
% mdparagraph:P764:begin
\textbf{būre} — chamber; room; dative, masculine, singular.
% mdparagraph:P764:end

% teaching: U18-SUP-HAESE-DAT-SG
% mdparagraph:P765:begin
\textbf{hǣse} — command; order; dative, feminine, singular.
% mdparagraph:P765:end

% teaching: U18-SUP-GIEFA-ACC-PL
% mdparagraph:P766:begin
\textbf{giefa} — gifts; accusative, feminine, plural.
% mdparagraph:P766:end

% teaching: U18-SUP-CWENE-ACC-SG-VARIANT
% mdparagraph:P767:begin
\textbf{cwēne} — princess; accusative, feminine, singular.
% mdparagraph:P767:end

% teaching: U18-SUP-GEMILTSIEND-ADDRESS
% mdparagraph:P768:begin
\textbf{gemiltsiend} — one who shows mercy; pitier; nominative, masculine, singular.
% mdparagraph:P768:end

% teaching: U18-SUP-LARE-GEN-SG
% mdparagraph:P769:begin
\textbf{lāre} — learning; study; genitive, feminine, singular.
% mdparagraph:P769:end

% teaching: U18-SUP-LUFIEND-ADDRESS
% mdparagraph:P770:begin
\textbf{lufiend} — lover; nominative, masculine, singular.
% mdparagraph:P770:end

% teaching: U18-SUP-BLITHE-JOYOUSLY
% mdparagraph:P771:begin
\textbf{blīðe} — joyfully; gladly.
% mdparagraph:P771:end

% teaching: U18-SUP-THEOWA-GEN-PL-PHRASE
% mdparagraph:P772:begin
\textbf{ðēowa} — enslaved men; servants; genitive, masculine, plural.
% mdparagraph:P772:end

% teaching: U18-SUP-THEOWUM-DAT-PL-PHRASE
% mdparagraph:P773:begin
\textbf{ðēowum} — enslaved men; servants; dative, masculine, plural.
% mdparagraph:P773:end

% teaching: U18-SUP-MOD-NOM-SG
% mdparagraph:P774:begin
\textbf{mōd} — heart; mind; nominative, neuter, singular.
% mdparagraph:P774:end

% teaching: U18-SUP-EARMRA-GEN-PL-SUBSTANTIVE
% mdparagraph:P775:begin
\textbf{earmra} — the poor; unfortunate people; genitive, masculine, plural.
% mdparagraph:P775:end

% teaching: U18-SUP-READING-ADDRESSEE
% mdparagraph:P776:begin
\textbf{tō} — to a speech addressee; governs dative.
% mdparagraph:P776:end

% teaching: U18-SUP-READING-DIRECTED-LOOK
% mdparagraph:P777:begin
\textbf{tō} — toward; target of a look; governs dative.
% mdparagraph:P777:end

% teaching: U18-SUP-READING-PERMISSION
% mdparagraph:P778:begin
\textbf{be} — by; with permission; governs dative.
% mdparagraph:P778:end

% teaching: U18-SUP-READING-LOVE-GOAL
% mdparagraph:P779:begin
\textbf{on} — into love for; governs accusative.
% mdparagraph:P779:end

% teaching: U18-SUP-READING-CONNECTIVE
% mdparagraph:P780:begin
\textbf{sōðlice} — now; indeed; narrative connective.
% mdparagraph:P780:end

% teaching: U18-SUP-READING-ACCOMPLISHMENTS
% mdparagraph:P781:begin
\textbf{cræftum} — skills; accomplishments; dative, masculine, plural.
% mdparagraph:P781:end

% mdparagraph:P782:begin
\textbf{mid-ðȳ-ðe} — when; while.
% mdparagraph:P782:end

% mdparagraph:P783:begin
\textbf{geseah} — see that a stated circumstance holds; indicative, preterite, third person, singular.
% mdparagraph:P783:end

% teaching: NEW-U18-VOC-GETEON-U18-READING
% mdparagraph:P784:begin
\textbf{getogen} — train; educate; past, participle.
% mdparagraph:P784:end

% teaching: NEW-U18-VOC-GEFEALLAN-U18-READING
% mdparagraph:P785:begin
\textbf{gefēoll} — fall, figuratively of the heart falling in love; indicative, preterite, third person, singular.
% mdparagraph:P785:end

% mdparagraph:P786:begin
\textbf{æfter} — after, temporal with dative.
% mdparagraph:P786:end

% teaching: NEW-U18-VOC-LIEFAN-PERMIT-U18-READING
% mdparagraph:P787:begin
\textbf{līefdest} — allow; permit; indicative, preterite, second person, singular.
% mdparagraph:P787:end

% mdparagraph:P788:begin
\textbf{swā-hwæt-swā} — whatsoever; whatever.
% mdparagraph:P788:end

% teaching: NEW-U18-VOC-UTGAN-U18-READING
% mdparagraph:P789:begin
\textbf{ūtēode} — go out; indicative, preterite, third person, singular.
% mdparagraph:P789:end

% mdparagraph:P790:begin
\textbf{hund} — hundred, in the supplied gold and silver weight expressions.
% mdparagraph:P790:end

% mdparagraph:P791:begin
\textbf{twēntig} — twenty.
% mdparagraph:P791:end

% teaching: NEW-U18-VOC-BEHATAN-U18-READING
% mdparagraph:P792:begin
\textbf{behēt} — promise; indicative, preterite, first person, singular.
% mdparagraph:P792:end

% mdparagraph:P793:begin
\textbf{lęcgeað} — lay; place; imperative, second person, plural.
% mdparagraph:P793:end

% mdparagraph:P794:begin
\textbf{innan} — within; inside.
% mdparagraph:P794:end

% teaching: NEW-U18-VOC-GEDON-U18-READING
% mdparagraph:P795:begin
\textbf{gedōn} — do; perform; past, participle.
% mdparagraph:P795:end

% mdparagraph:P796:begin
\textbf{ðe} — that; which; who; nominative, masculine, plural.
% mdparagraph:P796:end

% teaching: NEW-U18-VOC-GEENDIAN-U18-READING
% mdparagraph:P797:begin
\textbf{geęndode} — end; come to an end; indicative, preterite, third person, singular.
% mdparagraph:P797:end

% teaching: NEW-U18-VOC-BIDDAN-U18-FAREWELL
% mdparagraph:P798:begin
\textbf{bǣdon} — wish someone well, in the supplied farewell formula; indicative, preterite, third person, plural.
% mdparagraph:P798:end

% teaching: NEW-U18-VOC-BESEON-U18-READING
% mdparagraph:P799:begin
\textbf{beseah} — look toward; indicative, preterite, third person, singular.
% mdparagraph:P799:end

% mdparagraph:P800:begin
\textbf{sēcean} — seek; plain infinitive.
% mdparagraph:P800:end

% mdparagraph:P801:begin
\textbf{geręstan} — rest; repose; plain infinitive.
% mdparagraph:P801:end

\reserve{7\baselineskip}
% mdparagraph:P802:begin
\subsection*{Reading tasks}\BookHeading{nav-ch18-learner-0159}{subsubsection}{Reading tasks}{0}
% mdparagraph:P802:end

% mdparagraph:P803:begin
Keep the complete text, original notes and modern help available. Translate the whole assigned passage; the analysis questions explain the supplied relations without testing recall of unfamiliar source forms.
% mdparagraph:P803:end

\reserve{8\baselineskip}
% mdparagraph:P804:begin
\subsection*{R01}\BookHeading{nav-ch18-learner-0160}{subsubsection}{R01}{0}
% mdparagraph:P804:end

% item: NEW-U18-R01
% mdparagraph:P805:begin
Translate the complete assigned passage in the following three consecutive divisions. All local help and original notes remain available.
% mdparagraph:P805:end

% output: NEW-U18-R01-O001
% mdparagraph:P806:begin
\itempara{1}{Translate 175.18–176.2, through goldhorde.}
% mdparagraph:P806:end

% output: NEW-U18-R01-O002
% mdparagraph:P807:begin
\itempara{2}{Translate 176.3–176.13, through gesāwon.}
% mdparagraph:P807:end

% output: NEW-U18-R01-O003
% mdparagraph:P808:begin
\itempara{3}{Translate from Ðā sōðlice geęndode in 176.13 through 177.2 geręstan.}
% mdparagraph:P808:end

\reserve{8\baselineskip}
% mdparagraph:P809:begin
\subsection*{R02}\BookHeading{nav-ch18-learner-0161}{subsubsection}{R02}{0}
% mdparagraph:P809:end

% item: NEW-U18-R02
% mdparagraph:P810:begin
For each cited form, name the referent and explain its grammatical role. Keep the whole passage available.
% mdparagraph:P810:end

% output: NEW-U18-R02-O001
% mdparagraph:P811:begin
\itempara{1}{Explain hiere and his in 175.20.}
% mdparagraph:P811:end

% output: NEW-U18-R02-O002
% mdparagraph:P812:begin
\itempara{2}{Identify the speaker and addressee of ic and ðē in 176.5.}
% mdparagraph:P812:end

% output: NEW-U18-R02-O003
% mdparagraph:P813:begin
\itempara{3}{Distinguish him in 176.19 from him in 176.20.}
% mdparagraph:P813:end

\reserve{8\baselineskip}
% mdparagraph:P814:begin
\subsection*{R03}\BookHeading{nav-ch18-learner-0162}{subsubsection}{R03}{0}
% mdparagraph:P814:end

% item: NEW-U18-R03
% mdparagraph:P815:begin
Parse the following selected source expressions and explain their functions. Supplied lexical and quantity explanations remain visible.
% mdparagraph:P815:end

% output: NEW-U18-R03-O001
% mdparagraph:P816:begin
\itempara{1}{Compare the grammatical subjects of geseah 175.18 and gefēoll 175.20.}
% mdparagraph:P816:end

% output: NEW-U18-R03-O002
% mdparagraph:P817:begin
\itempara{2}{Explain the relationships expressed by punda and goldes in 176.6.}
% mdparagraph:P817:end

% output: NEW-U18-R03-O003
% mdparagraph:P818:begin
\itempara{3}{Use original note 13 to explain ðe hīe gesāwon in 176.13 and the role of giefa.}
% mdparagraph:P818:end

% output: NEW-U18-R03-O004
% mdparagraph:P819:begin
\itempara{4}{Compare the case, number and agreement of gesunde in 176.16 and 176.18.}
% mdparagraph:P819:end

\reserve{8\baselineskip}
% mdparagraph:P820:begin
\subsection*{R04}\BookHeading{nav-ch18-learner-0163}{subsubsection}{R04}{0}
% mdparagraph:P820:end

% item: NEW-U18-R04
% mdparagraph:P821:begin
Describe the clause relations in the following spans using the complete source and its local support.
% mdparagraph:P821:end

% output: NEW-U18-R04-O001
% mdparagraph:P822:begin
\itempara{1}{Analyse the opening sentence 175.18–20: identify the temporal frame, the seeing-content clause and the main event.}
% mdparagraph:P822:end

% output: NEW-U18-R04-O002
% mdparagraph:P823:begin
\itempara{2}{Compare ðæt ic mōste 176.1 with swā-hwæt-swā 176.2 and 176.3–4.}
% mdparagraph:P823:end

% output: NEW-U18-R04-O003
% mdparagraph:P824:begin
\itempara{3}{Analyse the final ðæt clause 177.2 and explain how it relates to gān wē sēcean.}
% mdparagraph:P824:end

\reserve{8\baselineskip}
% mdparagraph:P825:begin
\subsection*{R05}\BookHeading{nav-ch18-learner-0164}{subsubsection}{R05}{0}
% mdparagraph:P825:end

% item: NEW-U18-R05
% mdparagraph:P826:begin
Explain the passage’s sequence from the actual statements, distinguishing completed actions from proposals.
% mdparagraph:P826:end

% output: NEW-U18-R05-O001
% mdparagraph:P827:begin
\itempara{1}{What authority does the princess give for her gifts, and who selects them?}
% mdparagraph:P827:end

% output: NEW-U18-R05-O002
% mdparagraph:P828:begin
\itempara{2}{What does the final speech establish about Apollonius’s next action, and what remains unnarrated at this boundary?}
% mdparagraph:P828:end

\reserve{8\baselineskip}
% mdparagraph:P829:begin
\section*{Source and editorial notes}\BookHeading{nav-ch18-learner-0165}{subsection}{Source and editorial notes}{1}
% mdparagraph:P829:end

\begingroup\fontsize{10}{12.8}\selectfont
\setlength{\parskip}{1.5pt}
% mdparagraph:P830:begin
The English prompts are Cook, Exercises in Old English, joint Group XII, XIII, Exercises 44–48, printed pages 29–33 (PDF pages 37–41). The original prose is Cook, A First Book in Old English, second edition, Reader XII, printed 175.18–177.2 (PDF 191–193). Modern explanations and lettered tasks are separate teaching material. Note numbers restart by printed page; P175, P176 and P177 identify the page.
% mdparagraph:P830:end

% mdparagraph:P831:begin
Preserve the selected spellings, including unmarked æcern, læfel and feoh, reading lęcgeað, and ðēowa without added r. The strong servant noun and the acorn plural have separately checked dictionary and manuscript-edition support. Derived paradigm cells apply the given stem and ending rules; their display does not assert that every cell occurs in the passage.
% mdparagraph:P831:end

% mdparagraph:P832:begin
The note on getogen raises an etymological question about modern wanton; here getogen means educated or accomplished. The hǣse note points to modern behest; hǣs here means command. These historical prompts are not additional exercises. The Beowulf comparison describes the company rising and one man greeting another and wishing him health; it illustrates farewell without assigning a poetry reading.
% mdparagraph:P832:end

% mdparagraph:P833:begin
The Latin notes mean a little earlier (176 note 1), teacher or master (note 5), enslaved people (note 10), dining room (note 11), saying farewell (note 15), and farewell or be well (note 17). Note 2 recommends an English infinitive for the finite Old English permission clause. Note 4 observes that the going-out detail is absent from the Latin. Note 9 records a manuscript adjective variant; the learner text retains Cook’s edited dēorwurðes. The other numbered grammatical references are explained in the modern reading help.
% mdparagraph:P833:end

% mdparagraph:P834:begin
\textsuperscript{P175N10}: See getēon. What relation has getogen to Mod. Eng. wanton?
% mdparagraph:P834:end

% mdparagraph:P835:begin
\textsuperscript{P175N11}: See 55.
% mdparagraph:P835:end

% mdparagraph:P836:begin
\textsuperscript{P176N1}: Lat. paulo ante. See 178.
% mdparagraph:P836:end

% mdparagraph:P837:begin
\textsuperscript{P176N2}: Translate by the infinitive sign, to. The OE. follows the Latin.
% mdparagraph:P837:end

% mdparagraph:P838:begin
\textsuperscript{P176N3}: See 197.
% mdparagraph:P838:end

% mdparagraph:P839:begin
\textsuperscript{P176N4}: Not in Latin.
% mdparagraph:P839:end

% mdparagraph:P840:begin
\textsuperscript{P176N5}: Lat. magister.
% mdparagraph:P840:end

% mdparagraph:P841:begin
\textsuperscript{P176N6}: See 154. c.
% mdparagraph:P841:end

% mdparagraph:P842:begin
\textsuperscript{P176N7}: See 153. f.
% mdparagraph:P842:end

% mdparagraph:P843:begin
\textsuperscript{P176N8}: See 174.
% mdparagraph:P843:end

% mdparagraph:P844:begin
\textsuperscript{P176N9}: MS. deorwurðan.
% mdparagraph:P844:end

% mdparagraph:P845:begin
\textsuperscript{P176N10}: Lat. servos.
% mdparagraph:P845:end

% mdparagraph:P846:begin
\textsuperscript{P176N11}: Lat. triclinio.
% mdparagraph:P846:end

% mdparagraph:P847:begin
\textsuperscript{P176N12}: See the derivation of Mod. Eng. behest.
% mdparagraph:P847:end

% mdparagraph:P848:begin
\textsuperscript{P176N13}: Refers to męnn.
% mdparagraph:P848:end

% mdparagraph:P849:begin
\textsuperscript{P176N14}: So in Beowulf (653–655): “Werod eall ārās; grētte ðā . . . guma ōðerne, . . . and him hǣl ābēad.”
% mdparagraph:P849:end

% mdparagraph:P850:begin
\textsuperscript{P176N15}: Lat. vale dicentes.
% mdparagraph:P850:end

% mdparagraph:P851:begin
\textsuperscript{P176N16}: See 153. d.
% mdparagraph:P851:end

% mdparagraph:P852:begin
\textsuperscript{P176N17}: Lat. valete.
% mdparagraph:P852:end

% mdparagraph:P853:begin
\textsuperscript{P176N18}: See 188.
% mdparagraph:P853:end

% mdparagraph:P854:begin
\textsuperscript{P177N1}: See 193. a.
% mdparagraph:P854:end

% mdparagraph:P855:begin
\textsuperscript{P177N2}: See 184. b.
% mdparagraph:P855:end

\endgroup

% book-source-body:end
\par\label{end-ch18-learner}
\clearpage\endgroup

\BookRecto
\begingroup
\renewcommand{\runningtitle}{Sound history}
\hyphenpenalty=10000\exhyphenpenalty=10000
\emergencystretch=3em
\newcommand{\reserve}[1]{\par\penalty-100\begingroup\dimen0=\pagegoal\advance\dimen0 by-\pagetotal\ifdim\dimen0<#1\newpage\fi\endgroup}
% book-source-body:begin

\addfontfeatures{Ligatures=TeXOff}
\raggedright
% appendix-block:B001:title:begin
\booktitle{\BookComponent{sound-history-learner}{Why Old English forms vary: a short sound history}APPENDIX}{Why Old English forms vary: a short sound history}
% appendix-block:B001:end
% appendix-block:B002:paragraph:begin
The\BookPoint{bk-sound-history-learner-5a5f7fd73e0c} course gives the forms needed to read and write Old English. This optional appendix explains some of their history. A change in a neighbouring sound, the loss of an ending, or the reshaping of a word after a familiar pattern can leave an alternation inside a paradigm. Breaking and umlaut are particularly useful explanations; several other processes help distinguish their results.\par
% appendix-block:B002:end
% appendix-block:B003:paragraph:begin
The examples retain the course’s spellings. Consult the pronunciation guide for the reading convention used in this course; the discussion here concerns how sounds changed over time.\par
% appendix-block:B003:end
% section: SH-S01
\reserve{5\baselineskip}
% appendix-block:B004:heading:begin
\section*{1. Reading a historical explanation}\BookHeading{nav-sound-history-learner-0001}{subsection}{1. Reading a historical explanation}{1}
% appendix-block:B004:end
% appendix-block:B005:paragraph:begin
A \emph{front vowel} is formed with the tongue relatively far forward in the mouth; a \emph{back vowel} has the tongue farther back. \emph{Rounding} means rounding the lips. A vowel can move forward without losing its lip rounding: this is relevant to the relation between \textbf{u} and \textbf{y}. \emph{Raising} means moving the tongue higher. These terms describe positions, not loudness or pitch.\par
% appendix-block:B005:end
% appendix-block:B006:paragraph:begin
A \emph{diphthong} changes position within one syllable. Old English distinguishes short and long diphthongs, just as it distinguishes short and long simple vowels. Two written vowel letters therefore do not by themselves indicate length. The macron marks the long forms: compare the spellings \textbf{ea} and \textbf{ēa}.\par
% appendix-block:B006:end
% appendix-block:B007:paragraph:begin
In\BookPoint{bk-sound-history-learner-fd30ba546e22} historical explanations, \textbf{j} represents the consonant sound at the beginning of English \emph{yes}. A \emph{reconstructed form} is an earlier form inferred from evidence rather than preserved as an Old English example. Here a preceding asterisk marks a reconstruction. This local convention differs from Unit 3, where asterisks mark proposed errors in an exercise.\par
% appendix-block:B007:end
% appendix-block:B008:paragraph:begin
An arrow means “developed into”; it may summarize several changes rather than display every intermediate stage. Forms printed beside one another in a paradigm are coexisting grammatical forms. A singular does not become its plural through the historical passage of time.\par
% appendix-block:B008:end
% appendix-block:B009:paragraph:begin
Most of the explanations concern the prehistory of the West Saxon forms used in the course. Later changes and other dialects are labelled where they matter. A sound’s surroundings at the time of a change may differ from the letters visible in the surviving word.\par
% appendix-block:B009:end
% section: SH-S02
% examples: SH-E01
\reserve{5\baselineskip}
% appendix-block:B010:heading:begin
\section*{2. Fronting and restoration}\BookHeading{nav-sound-history-learner-0002}{subsection}{2. Fronting and restoration}{1}
% appendix-block:B010:end
% appendix-block:B011:paragraph:begin
In\BookPoint{bk-sound-history-learner-af4fa27ce6c1} stressed syllables, an early change brought short \textbf{a} forward to \textbf{æ}. This \emph{fronting}, also called \emph{Anglo-Frisian brightening}, did not apply in the same way before nasal consonants, \textbf{m} and \textbf{n}. The nasal environment had already given the vowel a different history.[1]\par
% appendix-block:B011:end
% appendix-block:B012:paragraph:begin
Later, some of the newly fronted \textbf{æ} returned to \textbf{a}. This \emph{restoration} occurred before an immediately following \textbf{w}, and when a back vowel in the next syllable followed an open first syllable. An \emph{open syllable} ends in a vowel. In the relevant pattern, a single consonant between vowels begins the following syllable; a cluster can close the first syllable and prevent restoration.\par
% appendix-block:B012:end
% appendix-block:B013:table:begin
\begingroup
\fontsize{10.5}{13.2}\selectfont
\setlength{\LTleft}{0pt}
\setlength{\LTright}{0pt plus 1fil}
\begin{longtable}{@{}>{\raggedright\arraybackslash}p{\dimexpr0.2500\linewidth-1.0000\tabcolsep\relax}>{\raggedright\arraybackslash}p{\dimexpr0.3500\linewidth-1.4000\tabcolsep\relax}>{\raggedright\arraybackslash}p{\dimexpr0.4000\linewidth-1.6000\tabcolsep\relax}@{}}
\toprule
\textbf{Form} & \textbf{Meaning and identity} & \textbf{What the form illustrates} \\
\midrule
\endfirsthead
\toprule
\textbf{Form} & \textbf{Meaning and identity} & \textbf{What the form illustrates} \\
\midrule
\endhead
\bottomrule
\endfoot
\bottomrule
\endlastfoot
\textbf{dæg} & day; nominative singular & The front vowel remains. \\
\textbf{dagas} & days; nominative plural & A following back vowel allowed restoration in the open first syllable. \\
\textbf{cræftum} & with skills; dative plural & The consonant cluster closes the first syllable, so the back vowel of the ending does not produce the same result. \\
\end{longtable}
\endgroup
% appendix-block:B013:end
% appendix-block:B014:paragraph:begin
Thus the familiar singular/plural alternation has a historical explanation. Its present distribution is learned with the noun: changing every \textbf{æ} to \textbf{a} before an ending would erase the distinction illustrated by \textbf{cræftum}.\par
% appendix-block:B014:end
% section: SH-S03
% examples: SH-E02
\reserve{5\baselineskip}
% appendix-block:B015:heading:begin
\section*{3. Breaking}\BookHeading{nav-sound-history-learner-0003}{subsection}{3. Breaking}{1}
% appendix-block:B015:end
% appendix-block:B016:paragraph:begin
\emph{Breaking}\BookPoint{bk-sound-history-learner-bf4f0271a73a} is the development of a diphthong from a stressed front vowel before particular consonants. A second vowel element arose as the pronunciation moved toward the following consonant. The immediate short-vowel inputs were \textbf{æ}, \textbf{e} and \textbf{i}. Earlier \textbf{a} had already been fronted to \textbf{æ} where this was relevant.[2]\par
% appendix-block:B016:end
% appendix-block:B017:paragraph:begin
The conditions differ according to the vowel. In the following table, “another consonant” includes a second \textbf{r}, but excludes \textbf{j} after \textbf{r}.\par
% appendix-block:B017:end
% appendix-block:B018:table:begin
\begingroup
\fontsize{10.5}{13.2}\selectfont
\setlength{\LTleft}{0pt}
\setlength{\LTright}{0pt plus 1fil}
\begin{longtable}{@{}>{\raggedright\arraybackslash}p{\dimexpr0.2500\linewidth-1.0000\tabcolsep\relax}>{\raggedright\arraybackslash}p{\dimexpr0.3500\linewidth-1.4000\tabcolsep\relax}>{\raggedright\arraybackslash}p{\dimexpr0.4000\linewidth-1.6000\tabcolsep\relax}@{}}
\toprule
\textbf{Earlier short stressed vowel} & \textbf{Result} & \textbf{Principal environments} \\
\midrule
\endfirsthead
\toprule
\textbf{Earlier short stressed vowel} & \textbf{Result} & \textbf{Principal environments} \\
\midrule
\endhead
\bottomrule
\endfoot
\bottomrule
\endlastfoot
\textbf{æ} & \textbf{ea} & Before \textbf{h}; before \textbf{r} plus another consonant; before \textbf{l} plus another consonant. \\
\textbf{e} & \textbf{eo} & Before \textbf{h}; before \textbf{r} plus another consonant. Before \textbf{l} plus a consonant, the regular case is \textbf{lh}. \\
\textbf{i} & \textbf{io} & Before \textbf{h} or \textbf{r} plus another consonant. Breaking before \textbf{l} is not securely established. \\
\end{longtable}
\endgroup
% appendix-block:B018:end
% appendix-block:B019:paragraph:begin
At the time of breaking, \textbf{h} represented a friction sound made farther back in the mouth, rather than the weak initial sound familiar from English \emph{house}. The exclusion of \textbf{rj} and the restricted treatment before \textbf{l} are part of the rule. “A front vowel before a consonant cluster” is too broad a description.\par
% appendix-block:B019:end
% appendix-block:B020:paragraph:begin
The contrast between \textbf{helpan}, “help”, and \textbf{healp}, “helped” (preterite singular), makes the vowel-specific condition visible. The reconstructed earlier preterite had \textbf{æ}:\par
% appendix-block:B020:end
% appendix-block:B021:paragraph:begin
\textbf{*hælp} → \textbf{healp}, “helped”.\par
% appendix-block:B021:end
% appendix-block:B022:paragraph:begin
That vowel broke before \textbf{lp}. The \textbf{e} in \textbf{helpan} did not break before the same cluster: regular breaking of \textbf{e} before \textbf{l} requires the more restricted environment just given. The two forms began with different stem vowels; breaking affected them differently. The adjective \textbf{eald}, “old”, likewise shows the outcome of \textbf{æ} before \textbf{ld}.\par
% appendix-block:B022:end
% appendix-block:B023:paragraph:begin
Breaking a short vowel did not make it long. The \textbf{ea} in \textbf{healp} and \textbf{eald} is short. Their following consonant clusters are a separate matter, relevant to syllable weight (§7). Some long front vowels also underwent breaking before \textbf{h}; the table above is a summary of short-vowel breaking, not its whole history.\par
% appendix-block:B023:end
% appendix-block:B024:paragraph:begin
West Saxon subsequently merged older \textbf{io} with \textbf{eo}. The same written \textbf{eo} can consequently have more than one earlier vowel behind it. Nor does every \textbf{ea} or \textbf{eo} come from breaking: §5 describes other sources of diphthong spellings.\par
% appendix-block:B024:end
% section: SH-S04
\reserve{5\baselineskip}
% appendix-block:B025:heading:begin
\section*{4. Umlaut: the influence of a following syllable}\BookHeading{nav-sound-history-learner-0004}{subsection}{4. Umlaut: the influence of a following syllable}{1}
% appendix-block:B025:end
% appendix-block:B026:paragraph:begin
\emph{Umlaut},\BookPoint{bk-sound-history-learner-b884d4929d35} or \emph{mutation}, is a change in one vowel under the influence of another sound in a following syllable. The most useful kind here is \emph{i-umlaut}, also called \emph{front mutation}. An older \textbf{i} or \textbf{j}, usually in the next syllable, caused the stressed vowel to be pronounced farther forward or higher. The speaker anticipated part of the following sound’s tongue position.[3]\par
% appendix-block:B026:end
% appendix-block:B027:paragraph:begin
The conditioning \textbf{i} or \textbf{j} often disappeared later, or was replaced by another vowel. Its effect on the root could remain. An Old English form may therefore have a mutated vowel even though its surviving ending contains neither \textbf{i} nor \textbf{j}.\par
% appendix-block:B027:end
% appendix-block:B028:paragraph:begin
Some characteristic results are:\par
% appendix-block:B028:end
% appendix-block:B029:table:begin
\begingroup
\fontsize{10.5}{13.2}\selectfont
\setlength{\LTleft}{0pt}
\setlength{\LTright}{0pt plus 1fil}
\begin{longtable}{@{}>{\raggedright\arraybackslash}p{\dimexpr0.2500\linewidth-1.0000\tabcolsep\relax}>{\raggedright\arraybackslash}p{\dimexpr0.3500\linewidth-1.4000\tabcolsep\relax}>{\raggedright\arraybackslash}p{\dimexpr0.4000\linewidth-1.6000\tabcolsep\relax}@{}}
\toprule
\textbf{Earlier vowel} & \textbf{Front-mutation result} & \textbf{Qualification} \\
\midrule
\endfirsthead
\toprule
\textbf{Earlier vowel} & \textbf{Front-mutation result} & \textbf{Qualification} \\
\midrule
\endhead
\bottomrule
\endfoot
\bottomrule
\endlastfoot
\textbf{u}, \textbf{ū} & \textbf{y}, \textbf{ȳ} & The tongue moves forward while the lips remain rounded. \\
\textbf{o}, \textbf{ō} & \textbf{e}, \textbf{ē} & Front vowels replace the back rounded vowels. \\
\textbf{æ} & \textbf{e} & The already front vowel is raised. \\
\textbf{ā} & \textbf{ǣ} & This compact correspondence excludes the separate nasal-vowel history. \\
\textbf{ea}, \textbf{ēa} & \textbf{ie}, \textbf{īe} & West Saxon diphthong correspondences; later spellings can differ. \\
\end{longtable}
\endgroup
% appendix-block:B029:end
% appendix-block:B030:paragraph:begin
This is a historical guide to selected correspondences. It is not an instruction to replace a vowel whenever a later written ending happens to contain \textbf{i}. The relevant environment must have existed when the change operated.\par
% appendix-block:B030:end
% examples: SH-E03
\reserve{5\baselineskip}
% appendix-block:B031:subheading:begin
\subsection*{Nouns}\BookHeading{nav-sound-history-learner-0005}{subsubsection}{Nouns}{0}
% appendix-block:B031:end
% appendix-block:B032:paragraph:begin
\textbf{Fōt},\BookPoint{bk-sound-history-learner-a00d2110e9a4} “foot”, has \textbf{fēt}, “feet”, in the nominative and accusative plural. The same \textbf{fēt} is also the selected dative singular, where it still means “foot”. The long \textbf{ē} reflects a formerly front-mutation environment. Other forms of the noun retain \textbf{ō}.\par
% appendix-block:B032:end
% appendix-block:B033:paragraph:begin
This explains an alternation already learned in the paradigm. Umlaut is not a general Old English plural ending: it occurs in particular inherited formations, and later analogy can also affect its distribution.\par
% appendix-block:B033:end
% examples: SH-E04
\reserve{5\baselineskip}
% appendix-block:B034:subheading:begin
\subsection*{Related words}\BookHeading{nav-sound-history-learner-0006}{subsubsection}{Related words}{0}
% appendix-block:B034:end
% appendix-block:B035:paragraph:begin
The\BookPoint{bk-sound-history-learner-64f61747fcac} noun \textbf{dōm}, “judgment”, and the weak verb \textbf{dēman}, “judge”, belong to the same word family. The earlier formation of the verb included \textbf{j}. Fulk represents that earlier verb as \textbf{*dōm(i)jan}, “judge”:\par
% appendix-block:B035:end
% appendix-block:B036:paragraph:begin
\textbf{*dōm(i)jan} → \textbf{dēman}, “judge”.\par
% appendix-block:B036:end
% appendix-block:B037:paragraph:begin
The arrow summarizes front mutation and the subsequent changes that removed the conditioning sound. The noun and the verb did not have identical endings, so their root vowels need not have developed identically. The surviving \textbf{-an} of \textbf{dēman} conceals the earlier conditioning sound.\par
% appendix-block:B037:end
% examples: SH-E05
\reserve{5\baselineskip}
% appendix-block:B038:subheading:begin
\subsection*{Comparison}\BookHeading{nav-sound-history-learner-0007}{subsubsection}{Comparison}{0}
% appendix-block:B038:end
% appendix-block:B039:table:begin
\begingroup
\fontsize{10.5}{13.2}\selectfont
\setlength{\LTleft}{0pt}
\setlength{\LTright}{0pt plus 1fil}
\begin{longtable}{@{}>{\raggedright\arraybackslash}p{\dimexpr0.2500\linewidth-1.0000\tabcolsep\relax}>{\raggedright\arraybackslash}p{\dimexpr0.3500\linewidth-1.4000\tabcolsep\relax}>{\raggedright\arraybackslash}p{\dimexpr0.4000\linewidth-1.6000\tabcolsep\relax}@{}}
\toprule
\textbf{Positive} & \textbf{Comparative} & \textbf{Superlative} \\
\midrule
\endfirsthead
\toprule
\textbf{Positive} & \textbf{Comparative} & \textbf{Superlative} \\
\midrule
\endhead
\bottomrule
\endfoot
\bottomrule
\endlastfoot
\textbf{eald}, old & \textbf{ieldra}, older & \textbf{ieldest}, oldest \\
\end{longtable}
\endgroup
% appendix-block:B039:end
% appendix-block:B040:paragraph:begin
The comparison forms had an older front-mutation environment absent from the positive. Their \textbf{ie} belongs with the \textbf{ea} of \textbf{eald}. This family also shows that processes can operate in succession: breaking supplied one diphthong, and later mutation could change it. The comparative shown is masculine nominative singular; the table compares stems rather than replacing the complete adjective paradigms.\par
% appendix-block:B040:end
% examples: SH-E06
\reserve{5\baselineskip}
% appendix-block:B041:subheading:begin
\subsection*{Present singulars, and an older change}\BookHeading{nav-sound-history-learner-0008}{subsubsection}{Present singulars, and an older change}{0}
% appendix-block:B041:end
% appendix-block:B042:paragraph:begin
Older second- and third-person singular present endings contained \textbf{i}. This helps explain why a changed vowel can occur in those cells while the infinitive, first singular and plural retain another vowel. Compare \textbf{hātan}, “call; name”, with third singular \textbf{hǣtt}, “calls; names”. Here \textbf{ā} beside \textbf{ǣ} illustrates front mutation; the consonant ending has its own development.\par
% appendix-block:B042:end
% appendix-block:B043:paragraph:begin
The\BookPoint{bk-sound-history-learner-02dd4c54579e} course also places \textbf{helpan}, “help”, beside \textbf{hilpð}, “helps”. Its historical explanation needs a distinction. The raising of \textbf{e} to \textbf{i} before an older \textbf{i} or \textbf{j} took place earlier than the later front-mutation changes just described. Cook and Fulk both recognize this distinction. A practical description of “present mutation” can group the visible alternations together without making them one event in time.[4]\par
% appendix-block:B043:end
% appendix-block:B044:paragraph:begin
The familiar \textbf{ēo} beside \textbf{īe} in some present stems likewise involves an earlier raised vowel and further development. A one-step equation between the written Old English vowels leaves out that history. Keep the learned form in its actual person, tense and mood; a historical explanation does not extend an indicative singular alternation to an entire subjunctive paradigm.\par
% appendix-block:B044:end
% section: SH-S05
\reserve{5\baselineskip}
% appendix-block:B045:heading:begin
\section*{5. Other origins and later treatment of diphthongs}\BookHeading{nav-sound-history-learner-0009}{subsection}{5. Other origins and later treatment of diphthongs}{1}
% appendix-block:B045:end
% examples: SH-E07
\reserve{5\baselineskip}
% appendix-block:B046:subheading:begin
\subsection*{After an initial palatal consonant}\BookHeading{nav-sound-history-learner-0010}{subsubsection}{After an initial palatal consonant}{0}
% appendix-block:B046:end
% appendix-block:B047:paragraph:begin
A\BookPoint{bk-sound-history-learner-43e30165bf50} \emph{palatal consonant} is pronounced with the tongue toward the hard palate, the front part of the roof of the mouth. In West Saxon, an initial palatal consonant could affect the stressed vowel following it. This \emph{palatal diphthongization} changed \textbf{æ}, \textbf{ǣ} to \textbf{ea}, \textbf{ēa}, and \textbf{e}, \textbf{ē} to \textbf{ie}, \textbf{īe}, in the relevant environments.[5]\par
% appendix-block:B047:end
% appendix-block:B048:paragraph:begin
The direction of influence differs from breaking: the vowel follows the conditioning consonant. The change also differs from the consonant itself becoming palatal, discussed in §6. It followed breaking; a vowel already altered by breaking did not simply undergo the same conversion again.\par
% appendix-block:B048:end
% appendix-block:B049:paragraph:begin
\textbf{Geat}, “gate”, illustrates the vowel effect after an initial palatal. Its plural \textbf{gatu}, “gates”, retains a back vowel. The earlier distinction between front and restored back vowels (§2) is therefore part of the explanation; the written letter \textbf{g} alone cannot predict the vowel.\par
% appendix-block:B049:end
% appendix-block:B050:paragraph:begin
The familiar verb provides several related results:\par
% appendix-block:B050:end
% appendix-block:B051:table:begin
\begingroup
\fontsize{10.5}{13.2}\selectfont
\setlength{\LTleft}{0pt}
\setlength{\LTright}{0pt plus 1fil}
\begin{longtable}{@{}>{\raggedright\arraybackslash}p{\dimexpr0.2500\linewidth-1.5000\tabcolsep\relax}>{\raggedright\arraybackslash}p{\dimexpr0.2500\linewidth-1.5000\tabcolsep\relax}>{\raggedright\arraybackslash}p{\dimexpr0.2500\linewidth-1.5000\tabcolsep\relax}>{\raggedright\arraybackslash}p{\dimexpr0.2500\linewidth-1.5000\tabcolsep\relax}@{}}
\toprule
\textbf{Infinitive} & \textbf{Preterite singular} & \textbf{Preterite plural} & \textbf{Past participle} \\
\midrule
\endfirsthead
\toprule
\textbf{Infinitive} & \textbf{Preterite singular} & \textbf{Preterite plural} & \textbf{Past participle} \\
\midrule
\endhead
\bottomrule
\endfoot
\bottomrule
\endlastfoot
\textbf{giefan}, give & \textbf{geaf}, gave & \textbf{gēafon}, gave & \textbf{giefen}, given \\
\end{longtable}
\endgroup
% appendix-block:B051:end
% appendix-block:B052:paragraph:begin
These are four principal parts used in the same language, not successive stages of a change. In particular, \textbf{geaf} has short \textbf{ea}, whereas \textbf{gēafon} has long \textbf{ēa}. Palatal diphthongization preserves that distinction of length. The pattern described here is centred on West Saxon; other Old English dialects need not show the same results.\par
% appendix-block:B052:end
% examples: SH-E08
\reserve{5\baselineskip}
% appendix-block:B053:subheading:begin
\subsection*{Back mutation}\BookHeading{nav-sound-history-learner-0011}{subsubsection}{Back mutation}{0}
% appendix-block:B053:end
% appendix-block:B054:paragraph:begin
\emph{Back mutation}\BookPoint{bk-sound-history-learner-bcd0bead7661} is another influence from a following syllable. A back vowel could cause a short stressed front vowel to acquire a second element: earlier \textbf{i} and \textbf{e} could give \textbf{io} and \textbf{eo}. With the West Saxon merger of \textbf{io} and \textbf{eo}, the usual visible result is \textbf{eo}. \textbf{Heofon}, “sky; heaven”, is an example.[6]\par
% appendix-block:B054:end
% appendix-block:B055:paragraph:begin
West Saxon restricts the change chiefly to a single intervening \emph{labial}, a consonant made with the lips, or a \emph{liquid}, \textbf{l} or \textbf{r}. The relevant labials are \textbf{p}, \textbf{f}, \textbf{m} and \textbf{w}. More than one intervening consonant normally prevents this West Saxon change. The following vowels are also unequal in effect: \textbf{a} affects preceding \textbf{i}, but does not produce this mutation of preceding \textbf{e} in West Saxon.\par
% appendix-block:B055:end
% appendix-block:B056:paragraph:begin
The\BookPoint{bk-sound-history-learner-1ac916116017} distribution is wider in other dialects. Even within West Saxon, speakers could restore an unmutated vowel by following related forms. Such reshaping is called \emph{analogy}. The source spelling \textbf{hefon}, also “sky; heaven”, has a simple vowel; for dictionary lookup, see \textbf{heofon}. A historical tendency can thus coexist with forms that have been altered after it operated.\par
% appendix-block:B056:end
% examples: SH-E15
\reserve{5\baselineskip}
% appendix-block:B057:subheading:begin
\subsection*{Smoothing}\BookHeading{nav-sound-history-learner-0012}{subsubsection}{Smoothing}{0}
% appendix-block:B057:end
% appendix-block:B058:paragraph:begin
\emph{Smoothing}\BookPoint{bk-sound-history-learner-ffea67e44b42} replaces a diphthong with a single vowel. Late West Saxon smoothing affects \textbf{ea}, \textbf{ēa} in certain environments, including before \textbf{c}, \textbf{g}, \textbf{h}, or after an initial palatal. It is inconsistently represented in writing. A labelled example is \textbf{hēah}, “high”, beside the Late West Saxon spelling \textbf{hēh}, also “high”. For \textbf{hēh}, see \textbf{hēah} in the glossary.[7]\par
% appendix-block:B058:end
% appendix-block:B059:paragraph:begin
\emph{Anglian smoothing} has a different and wider scope. It affects more diphthongs before back consonants, sometimes with \textbf{r} or \textbf{l} intervening. The two developments should retain their dialect and period labels even where a resulting spelling looks alike. Neither supplies a universal instruction to simplify Old English diphthongs.\par
% appendix-block:B059:end
% section: SH-S06
\reserve{5\baselineskip}
% appendix-block:B060:heading:begin
\section*{6. Consonants and the order of changes}\BookHeading{nav-sound-history-learner-0013}{subsection}{6. Consonants and the order of changes}{1}
% appendix-block:B060:end
% examples: SH-E09
\reserve{5\baselineskip}
% appendix-block:B061:subheading:begin
\subsection*{Palatalization}\BookHeading{nav-sound-history-learner-0014}{subsubsection}{Palatalization}{0}
% appendix-block:B061:end
% appendix-block:B062:paragraph:begin
\emph{Palatalization}\BookPoint{bk-sound-history-learner-5d52be86de50} changes a consonant’s articulation toward the hard palate. In appropriate historical environments, consonants written \textbf{c} or \textbf{g} acquired the sounds represented by English \textbf{ch} or consonantal \textbf{y}. The pronunciation guide gives their practical values and exceptions. Their history explains why a following front-vowel letter is useful evidence without being a complete spelling rule.[8]\par
% appendix-block:B062:end
% appendix-block:B063:paragraph:begin
Consider \textbf{cēlan}, “cool”, whose \textbf{c} is pronounced like English \textbf{k}. Its earlier form \textbf{*kōljan}, “cool”, had a back vowel next to the initial consonant. The later fronting produced \textbf{ē} after the consonant had remained nonpalatal. Reading the surviving \textbf{ē} as proof that the consonant must have become palatal reverses the order of events.\par
% appendix-block:B063:end
% appendix-block:B064:paragraph:begin
Different consonants and positions had different conditions, and some details remain disputed. The useful principle is chronological: a sound change responded to the sounds present at its own date, including sounds later lost or altered.\par
% appendix-block:B064:end
% examples: SH-E10
\reserve{5\baselineskip}
% appendix-block:B065:subheading:begin
\subsection*{Gemination}\BookHeading{nav-sound-history-learner-0015}{subsubsection}{Gemination}{0}
% appendix-block:B065:end
% appendix-block:B066:paragraph:begin
\emph{Gemination}\BookPoint{bk-sound-history-learner-8fefc2cf70de} is consonant lengthening. A following \textbf{j} commonly lengthened the preceding consonant after a light syllable (see §7); \textbf{r} is an important exception. The \textbf{j} could later disappear while its effect remained.[9]\par
% appendix-block:B066:end
% appendix-block:B067:paragraph:begin
Old English \textbf{biddan}, “ask”, has a doubled consonant. Compare Gothic \textbf{bidjan}, also “ask”, which visibly retains \textbf{j}. Gothic is a related Germanic language: its form is a \emph{cognate}, not an Old English spelling or the parent of the Old English word. The comparison makes the historical conditioning sound easier to recognize. Other processes can also produce doubled consonants, so a double letter alone does not prove an earlier \textbf{j}.\par
% appendix-block:B067:end
% section: SH-S07
\reserve{5\baselineskip}
% appendix-block:B068:heading:begin
\section*{7. Vowel loss, weight and contraction}\BookHeading{nav-sound-history-learner-0016}{subsection}{7. Vowel loss, weight and contraction}{1}
% appendix-block:B068:end
% examples: SH-E11
\reserve{5\baselineskip}
% appendix-block:B069:subheading:begin
\subsection*{Final and internal vowels}\BookHeading{nav-sound-history-learner-0017}{subsubsection}{Final and internal vowels}{0}
% appendix-block:B069:end
% appendix-block:B070:paragraph:begin
\emph{Apocope}\BookPoint{bk-sound-history-learner-f2eedb8e42ac} is loss at the end of a word; \emph{syncope} is loss inside it. Old English lost many unstressed vowels, but the results depended on their position and on syllable weight. Vowel length and weight are distinct: a short vowel can occur in a heavy stem.[10]\par
% appendix-block:B070:end
% appendix-block:B071:paragraph:begin
For\BookPoint{bk-sound-history-learner-771e0549912c} the ordinary noun pattern illustrated here, a short vowel or short diphthong followed by one consonant gives a light stem. A long vowel or long diphthong gives a heavy stem; so does a short vowel followed by two consonants. A doubled consonant counts as two. These differences helped determine whether a final short high vowel survived.\par
% appendix-block:B071:end
% appendix-block:B072:table:begin
\begingroup
\fontsize{10.5}{13.2}\selectfont
\setlength{\LTleft}{0pt}
\setlength{\LTright}{0pt plus 1fil}
\begin{longtable}{@{}>{\raggedright\arraybackslash}p{\dimexpr0.3500\linewidth-1.4000\tabcolsep\relax}>{\raggedright\arraybackslash}p{\dimexpr0.3500\linewidth-1.4000\tabcolsep\relax}>{\raggedright\arraybackslash}p{\dimexpr0.3000\linewidth-1.2000\tabcolsep\relax}@{}}
\toprule
\textbf{Nominative/accusative singular} & \textbf{Nominative/accusative plural} & \textbf{Relevant contrast} \\
\midrule
\endfirsthead
\toprule
\textbf{Nominative/accusative singular} & \textbf{Nominative/accusative plural} & \textbf{Relevant contrast} \\
\midrule
\endhead
\bottomrule
\endfoot
\bottomrule
\endlastfoot
\textbf{hof}, dwelling & \textbf{hofu}, dwellings & Light stem: plural \textbf{-u} remains. \\
\textbf{word}, word & \textbf{word}, words & Short \textbf{o}, but a heavy stem because of \textbf{rd}. \\
\textbf{wīf}, wife & \textbf{wīf}, wives & Heavy stem with long \textbf{ī}. \\
\end{longtable}
\endgroup
% appendix-block:B072:end
% appendix-block:B073:paragraph:begin
Historical high-vowel loss helps explain this distribution; it does not remove all final vowels from every heavy word. Later analogy and the histories of individual noun classes also matter. Use the actual paradigm for a particular noun.\par
% appendix-block:B073:end
% appendix-block:B074:paragraph:begin
For\BookPoint{bk-sound-history-learner-00c58def41fa} internal loss, compare \textbf{hēafod}, “head”, with genitive \textbf{hēafdes}, “head’s”. The unstressed middle vowel is absent before the ending. The source also records full \textbf{hēafodes}, “head’s”, where analogy has restored the fuller stem. For \textbf{hēafodes}, see \textbf{hēafod} in the glossary. The coexistence of the two genitives shows why a broad syncope rule cannot replace lexical and period-specific evidence.\par
% appendix-block:B074:end
% examples: SH-E12
\reserve{5\baselineskip}
% appendix-block:B075:subheading:begin
\subsection*{Loss of a consonant and contraction}\BookHeading{nav-sound-history-learner-0018}{subsubsection}{Loss of a consonant and contraction}{0}
% appendix-block:B075:end
% appendix-block:B076:paragraph:begin
When\BookPoint{bk-sound-history-learner-fb537a53e749} an intervocalic \textbf{h} disappeared, it brought vowels together. \emph{Contraction} could then remove the second vowel and leave a long vowel or diphthong. Thus the reconstructed earlier infinitive \textbf{*seohan}, “see”, developed into \textbf{sēon}, “see”.[11]\par
% appendix-block:B076:end
% appendix-block:B077:paragraph:begin
The\BookPoint{bk-sound-history-learner-d2a70dff307f} imperative singular \textbf{seoh}, “see!”, retains \textbf{h} in a different environment. Its \textbf{eo} is short; the infinitive’s \textbf{ēo} is long. Contraction therefore explains an important part of the pattern without accounting by itself for every form of a contracted verb. Loss of a consonant can also cause \emph{compensatory lengthening}, an increase in the preceding vowel’s duration; the exact environment must be checked separately.\par
% appendix-block:B077:end
% section: SH-S08
\reserve{5\baselineskip}
% appendix-block:B078:heading:begin
\section*{8. Ablaut, inherited alternations and analogy}\BookHeading{nav-sound-history-learner-0019}{subsection}{8. Ablaut, inherited alternations and analogy}{1}
% appendix-block:B078:end
% examples: SH-E13
% appendix-block:B079:paragraph:begin
\emph{Ablaut}\BookPoint{bk-sound-history-learner-7a97cb9f6faa} is an inherited alternation of stem vowels. Strong-verb principal parts preserve such alternations, subsequently reshaped by other sound changes. They are not formed by applying umlaut to an infinitive when a past meaning is needed.[12]\par
% appendix-block:B079:end
% appendix-block:B080:table:begin
\begingroup
\fontsize{10.5}{13.2}\selectfont
\setlength{\LTleft}{0pt}
\setlength{\LTright}{0pt plus 1fil}
\begin{longtable}{@{}>{\raggedright\arraybackslash}p{\dimexpr0.2500\linewidth-1.5000\tabcolsep\relax}>{\raggedright\arraybackslash}p{\dimexpr0.2500\linewidth-1.5000\tabcolsep\relax}>{\raggedright\arraybackslash}p{\dimexpr0.2500\linewidth-1.5000\tabcolsep\relax}>{\raggedright\arraybackslash}p{\dimexpr0.2500\linewidth-1.5000\tabcolsep\relax}@{}}
\toprule
\textbf{Infinitive} & \textbf{Preterite singular} & \textbf{Preterite plural} & \textbf{Past participle} \\
\midrule
\endfirsthead
\toprule
\textbf{Infinitive} & \textbf{Preterite singular} & \textbf{Preterite plural} & \textbf{Past participle} \\
\midrule
\endhead
\bottomrule
\endfoot
\bottomrule
\endlastfoot
\textbf{drīfan}, drive & \textbf{drāf}, drove & \textbf{drifon}, drove & \textbf{drifen}, driven \\
\textbf{helpan}, help & \textbf{healp}, helped & \textbf{hulpon}, helped & \textbf{holpen}, helped \\
\end{longtable}
\endgroup
% appendix-block:B080:end
% appendix-block:B081:paragraph:begin
The first series has \textbf{ī}, \textbf{ā}, \textbf{i}, \textbf{i}. The second shows the combined effects of inherited alternation and later change: the earlier different vowel of the singular preterite underwent breaking, giving \textbf{healp} (§3). Its \textbf{ea} is therefore not an untouched original ablaut vowel. The present singular \textbf{hilpð}, “helps”, belongs to the separate raising history explained in §4.\par
% appendix-block:B081:end
% appendix-block:B082:paragraph:begin
One word family can preserve several layers of history. A different strong class need not have the same series, and a weak verb can also have a changed stem vowel while retaining its dental preterite. The principal parts identify the forms actually needed.\par
% appendix-block:B082:end
% examples: SH-E14
% appendix-block:B083:paragraph:begin
Some\BookPoint{bk-sound-history-learner-880e35b3cd92} inherited alternations involve consonants. Compare \textbf{cweðan}, “say”, and its past participle \textbf{cweden}, “said”. The \textbf{ð} beside \textbf{d} belongs to a history of consonant alternation. \emph{Verner’s law} describes an early accent-conditioned voicing process whose results were further altered by later sound changes. It operated before the ordinary Germanic stress pattern became established.[13]\par
% appendix-block:B083:end
% appendix-block:B084:paragraph:begin
This inherited alternation is distinct from the contextual voicing of fricatives in ordinary Old English pronunciation. It neither makes \textbf{ð} and \textbf{d} interchangeable spellings nor allows the present position of stress to generate all the older forms. Analogy can preserve, extend or remove an alternation, just as it can restore a lost vowel. Historical explanation and the recorded paradigm are therefore complementary kinds of evidence.\par
% appendix-block:B084:end
% section: SH-S09
\reserve{5\baselineskip}
% appendix-block:B085:heading:begin
\section*{9. Reference summary and source notes}\BookHeading{nav-sound-history-learner-0020}{subsection}{9. Reference summary and source notes}{1}
% appendix-block:B085:end
% appendix-block:B086:table:begin
\begingroup
\fontsize{10.5}{13.2}\selectfont
\setlength{\LTleft}{0pt}
\setlength{\LTright}{0pt plus 1fil}
\begin{longtable}{@{}>{\raggedright\arraybackslash}p{\dimexpr0.3400\linewidth-0.6800\tabcolsep\relax}>{\raggedright\arraybackslash}p{\dimexpr0.6600\linewidth-1.3200\tabcolsep\relax}@{}}
\toprule
\textbf{Process} & \textbf{Historical relation to look for} \\
\midrule
\endfirsthead
\toprule
\textbf{Process} & \textbf{Historical relation to look for} \\
\midrule
\endhead
\bottomrule
\endfoot
\bottomrule
\endlastfoot
Fronting and restoration & A stressed front vowel and a restored back vowel, conditioned by an earlier following sound and syllable structure. \\
Breaking & A stressed front vowel acquiring a second element before particular consonants; the conditions differ by vowel. \\
Front mutation & A vowel anticipating older \textbf{i} or \textbf{j}, often lost from the surviving form. \\
Palatal diphthongization & A stressed vowel affected by an initial palatal consonant, especially in West Saxon. \\
Back mutation & A stressed front vowel affected by a following back vowel, with dialect-specific intervening-consonant restrictions. \\
Smoothing & A diphthong becoming a simple vowel in a specified period and dialect. \\
Gemination & A consonant becoming longer; an older \textbf{j} is one possible cause. \\
Vowel loss and contraction & Endings or internal vowels disappearing, or adjacent vowels merging after consonant loss. \\
Ablaut & Inherited stem-vowel alternation, subsequently affected by other changes and analogy. \\
\end{longtable}
\endgroup
% appendix-block:B086:end
% appendix-block:B087:paragraph:begin
The account follows Cook’s selected 1895 second edition of \emph{A First Book in Old English} and R. D. Fulk’s \emph{An Introductory Grammar of Old English with an Anthology of Readings}, in its author-corrected 2020 open-access reprint of the 2014 edition. Fulk’s editorial dots over palatal consonants are omitted in the course-style examples; the pronunciation distinctions are explained in words. The starred earlier forms and Gothic comparison remain explicitly separate from the Old English examples. PDF references below count pages from the beginning of each file, starting at one.\par
% appendix-block:B087:end
% appendix-block:B088:paragraph:begin
1. Fronting and restoration: Fulk §29, printed pp. 18–19 / PDF 30–31; Appendix A IV.26, printed p. 111 / PDF 123. Cook §43.2, printed p. 28 / PDF 44, gives the familiar noun alternation. The back-vowel restoration condition does not extend through arbitrary consonant clusters.\par
% appendix-block:B088:end
% appendix-block:B089:paragraph:begin
2. Breaking: Fulk §62, printed pp. 41–42 / PDF 53–54; §91, printed p. 61 / PDF 73; Appendix A IV.28, printed p. 111 / PDF 123. Cook §21, printed pp. 16–17 / PDF 32–33, explicitly qualifies his shorthand by identifying \textbf{æ} as the immediate input where he first writes \textbf{a}. The vowel-specific restrictions before \textbf{l} follow Fulk. His Appendix A V.45, printed p. 113 / PDF 125, explicitly excludes breaking in unstressed syllables.\par
% appendix-block:B089:end
% appendix-block:B090:paragraph:begin
3. Front mutation: Cook §§15–17, printed pp. 12–13 / PDF 28–29, and Appendix V, printed pp. 268–269 / PDF 284–285; Fulk §74, printed pp. 48–49 / PDF 60–61, §79, printed p. 52 / PDF 64, §82, printed p. 53 / PDF 65, and §109, printed p. 78 / PDF 90. The comparison forms and the selected \textbf{hǣtt} are those of Units 3 and 5.\par
% appendix-block:B090:end
% appendix-block:B091:paragraph:begin
4. Older raising: Cook §17, printed p. 13 / PDF 29, and §104, printed pp. 61–62 / PDF 77–78; Fulk §§73–76, printed pp. 48–49 / PDF 60–61, and Appendix A IV.18, printed p. 110 / PDF 122. Both authors qualify the practical grouping of \textbf{e} beside \textbf{i} with later umlaut.\par
% appendix-block:B091:end
% appendix-block:B092:paragraph:begin
5. Palatal diphthongization: Cook §18, printed pp. 13–14 / PDF 29–30; Fulk §71, printed pp. 45–46 / PDF 57–58, and Appendix B I.5, printed p. 119 / PDF 131. The complete principal-part series and the noun forms agree with Units 5 and 11.\par
% appendix-block:B092:end
% appendix-block:B093:paragraph:begin
6. Back mutation: Fulk §119, printed p. 85 / PDF 97, and Appendix B I.3, printed p. 118 / PDF 130; compare Cook §20, printed pp. 15–16 / PDF 31–32. Fulk supplies the West Saxon restrictions and discusses analogical restoration of a simple vowel. The two spellings compared here occur at separate points in his paragraph.\par
% appendix-block:B093:end
% appendix-block:B094:paragraph:begin
7. Smoothing: Fulk §111d, printed p. 79 / PDF 91; Appendix A IV.35, printed p. 112 / PDF 124; Appendix B III.12, printed p. 119 / PDF 131. Late West Saxon and Anglian smoothing have separate conditions. Cook §19, printed pp. 14–15 / PDF 30–31, is useful for other vowel-spelling correspondences, but is not the authority for this smoothing account.\par
% appendix-block:B094:end
% appendix-block:B095:paragraph:begin
8. Palatal consonants and chronology: Fulk §135a–d, printed pp. 103–104 / PDF 115–116; §71, printed p. 45 / PDF 57. The pronunciation guide, §4, supplies the practical hard-consonant reading of \textbf{cēlan}. Fulk notes disagreement over some of the detailed palatalization conditions.\par
% appendix-block:B095:end
% appendix-block:B096:paragraph:begin
9. Gemination: Fulk §67, printed p. 44 / PDF 56, and Appendix A III.4–6, printed p. 109 / PDF 121; Cook §§33–36, printed pp. 24–25 / PDF 40–41. Cook Appendix V also cautions against treating Gothic as the parent of the other Germanic languages.\par
% appendix-block:B096:end
% appendix-block:B097:paragraph:begin
10. Vowel loss and weight: Cook §4, printed p. 6 / PDF 22; §23, printed pp. 18–19 / PDF 34–35; §43.4, printed pp. 28–29 / PDF 44–45; §§46–47, printed pp. 31–33 / PDF 47–49. Fulk §§32–33, printed p. 21 / PDF 33, and §137a–b, printed p. 104 / PDF 116, supply the historical account and the full-stem variant.\par
% appendix-block:B097:end
% appendix-block:B098:paragraph:begin
11. Loss of \textbf{h} and contraction: Fulk §§131–132, printed pp. 98–99 / PDF 110–111; compare Cook §§27–29, printed pp. 21–23 / PDF 37–39. The imperative is also supplied in Unit 6. Compensatory lengthening and contraction have distinct environments.\par
% appendix-block:B098:end
% appendix-block:B099:paragraph:begin
12. Ablaut: Cook §22, printed pp. 17–18 / PDF 33–34, and §104, printed pp. 61–62 / PDF 77–78; Fulk §85, printed p. 57 / PDF 69, and §91, printed p. 61 / PDF 73. Fulk Appendix A VI offers a fuller account of the inherited vowel system.\par
% appendix-block:B099:end
% appendix-block:B100:paragraph:begin
13. Consonant alternation: Cook §37, printed p. 25 / PDF 41; Fulk §87, printed p. 58 / PDF 70, and Appendix A II, printed pp. 108–109 / PDF 120–121. These locators support a bounded explanation of the historical alternation, rather than a new rule for freely exchanging consonants.\par
% appendix-block:B100:end

% book-source-body:end
\par\label{end-sound-history-learner}
\clearpage\endgroup

\BookRecto
\begingroup
\renewcommand{\runningtitle}{Old English–English}
\raggedcolumns
\newcommand{\needspace}[1]{\par\begingroup\skip0=\lastskip\vskip-\skip0\dimen0=#1\relax\advance\dimen0 by\skip0\vskip0pt plus\dimen0\penalty-100\vskip0pt plus-\dimen0\vskip\dimen0\penalty9999\vskip-\dimen0\vskip\skip0\endgroup}
\setlength{\columnsep}{6mm}
\hyphenpenalty=10000\exhyphenpenalty=10000
\emergencystretch=2em
% book-source-body:begin

\addfontfeatures{Ligatures=TeXOff}
\raggedright
% book-navigation: standalone bookmark replaced by volume hierarchy
% glossary-block:B00001:begin
\booktitle{\BookComponent{glossary-old-english}{Old English–English glossary}GLOSSARY}{Old English–English glossary}
% glossary-block:B00001:end
\fontsize{10.5}{13.2}\selectfont
% glossary-block:B00002:begin
This glossary covers the Old English in all eighteen units and their keys, including tables, readings, alternatives and source notes, and in the pronunciation guide and its key. It is a reference, not a new memorization list. The exercise instructions still determine what you must know and what is supplied.\par
% glossary-block:B00002:end
% glossary-block:B00003:begin
Look under the dictionary headword, normally the infinitive for a verb. Strong verbs give the preterite first singular, preterite plural and past participle; weak verbs give the preterite and participle. Distinctive presents and irregular verbs have their forms labelled. Selected difficult forms, including vowel-changing third-person singulars, have \textbf{see} entries. Routine inflections belong to the unit paradigms.\par
% glossary-block:B00003:end
% glossary-block:B00004:begin
Nouns give gender, genitive singular and nominative plural where established. Adjectives and adverbs give selected comparison forms or useful stem clues. Words appearing only in a source note may have a limited citation rather than a complete set of principal parts. A missing cell supplies no form by implication. Meanings and construction notes retain the limits of the cited use.\par
% glossary-block:B00004:end
% glossary-block:B00005:begin
The order is a, æ, b … t, þ/ð, u … z. Length marks and hooks do not affect alphabetization; they remain essential parts of the printed spelling. Thorn and eth share a sorting position. Words beginning ge- are under g. Ordinary forms are listed in lowercase; proper-name capitalization is retained. Sentence-initial capitals do not create separate entries.\par
% glossary-block:B00005:end
% glossary-block:B00006:begin
Abbreviations: n. noun; v. verb; adj. adjective; adv. adverb; prep. preposition; conj. conjunction; pron. pronoun; interj. interjection; m., f., n. masculine, feminine, neuter; gen. genitive; nom. nominative; pres. present; pret. preterite; ptcp. participle; sg. singular; pl. plural; comp. comparative; superl. superlative. Numbers 1 and 3 identify persons. Parenthesized (ge) preserves the source's optional prefix notation; it does not license adding ge- to other words.\par
% glossary-block:B00006:end
% glossary-block:B00007:begin
Case and construction notes qualify the meaning beside them. A supplied expression is not a general word-for-word translation rule. Source spellings corrected in the course, deliberately wrong examples and reconstructed forms are identified separately at the end.\par
% glossary-block:B00007:end
% glossary-block:B00008:begin
The definitions and spelling follow the course's checked Cook sources, with the recorded Fulk, Wright and dictionary support where used. Principal parts include checked reference formations from the stated grammatical rules as well as exact source citations. This glossary is newly arranged reference material. The companion English–Old English glossary indexes exercise meanings and constructions.\par
% glossary-block:B00008:end
% book-navigation: standalone bookmark replaced by volume hierarchy
\clearpage
\fontsize{10.5}{13.2}\selectfont
\setlength{\parskip}{3.8pt}
\begin{multicols*}{2}
% glossary-block:B00009:begin
\needspace{5\baselineskip}
\section*{Entries}\BookHeading{nav-glossary-old-english-0001}{subsection}{Entries}{1}
% glossary-block:B00009:end
% book-navigation: standalone bookmark replaced by volume hierarchy
% glossary-block:B00010:begin
\needspace{5\baselineskip}
\subsection*{A}\BookHeading{nav-glossary-old-english-0002}{subsubsection}{A}{0}
% glossary-block:B00010:end
% glossary-block:B00011:begin
\textbf{ā} (adv.). ever; always.\par
% glossary-block:B00011:end
% glossary-block:B00012:begin
\needspace{3\baselineskip}
\textbf{ābēodan} (v.). announce; bid a greeting (hǣl ābēodan). hǣl ābēodan: bid someone hail / give someone greeting.\par
% glossary-block:B00012:end
% glossary-block:B00013:begin
\needspace{3\baselineskip}
\textbf{āberan} (v.). carry. Supplied carried form; retain the ā-prefix.\par
% glossary-block:B00013:end
% glossary-block:B00014:begin
\needspace{3\baselineskip}
\textbf{Abraham} (proper name). Abraham. Biblical patriarch; supplied name forms.\par
% glossary-block:B00014:end
% glossary-block:B00015:begin
\needspace{3\baselineskip}
\textbf{ābūtan} (prep.). around; about. With accusative.\par
% glossary-block:B00015:end
% glossary-block:B00016:begin
\needspace{3\baselineskip}
\textbf{āc} (n.). oak. Pronunciation example.\par
% glossary-block:B00016:end
% glossary-block:B00017:begin
\needspace{3\baselineskip}
\textbf{ac} (conj.). but.\par
% glossary-block:B00017:end
% glossary-block:B00018:begin
\needspace{3\baselineskip}
\textbf{ācęnnan} (v.); past ptcp. ācenned, ācęnned. produce; bring forth. beget; bring forth.\par
% glossary-block:B00018:end
% glossary-block:B00019:begin
\needspace{3\baselineskip}
\textbf{āceorfan} (v. strong); pret. 1 sg. ācearf; pret. pl. ācurfon; past ptcp. ācorfen; pres. 3 sg. ācierfð. cut off. Accusative thing cut; of in the assigned expression is a separate off particle.\par
% glossary-block:B00019:end
% glossary-block:B00020:begin
\needspace{3\baselineskip}
\textbf{ācierfð} — pres. ind. 3 sg.: see \textbf{āceorfan}.\par
% glossary-block:B00020:end
% glossary-block:B00021:begin
\needspace{3\baselineskip}
\textbf{adesa} (n. m.); gen. sg. adesan; nom. pl. adesan. hatchet.\par
% glossary-block:B00021:end
% glossary-block:B00022:begin
\needspace{3\baselineskip}
\textbf{ādrīfan} (v. strong); pret. 1 sg. ādrāf; pret. pl. ādrifon; past ptcp. ādrifen. drive out; expel. Accusative person expelled; fram + dative may mark the point of separation.\par
% glossary-block:B00022:end
% glossary-block:B00023:begin
\needspace{3\baselineskip}
\textbf{āfǣran} (v.); past ptcp. āfǣred. frighten.\par
% glossary-block:B00023:end
% glossary-block:B00024:begin
\needspace{3\baselineskip}
\textbf{āfēdan} (v. weak); pret. 1 sg. āfēdde; past ptcp. āfēded. support; nourish. Accusative person supported.\par
% glossary-block:B00024:end
% glossary-block:B00025:begin
\needspace{3\baselineskip}
\textbf{āfierran} (v.). remove; take away.\par
% glossary-block:B00025:end
% glossary-block:B00026:begin
\needspace{3\baselineskip}
\textbf{āflīeman} (v.). put to flight; expel.\par
% glossary-block:B00026:end
% glossary-block:B00027:begin
\needspace{3\baselineskip}
\textbf{āfyllan} (v. weak); pret. 1 sg. āfylde; past ptcp. āfylled. fill.\par
% glossary-block:B00027:end
% glossary-block:B00028:begin
\needspace{3\baselineskip}
\textbf{āgan} (v. preterite-present); pres. 1/3 sg. āh; pres. pl. āgon; pret. 1 sg. āhte; adjectival ptcp. āgen. have; possess. Accusative possession; positive present and preterite indicative in the taught scope. Āgen is adjectival: “own”.\par
% glossary-block:B00028:end
% glossary-block:B00029:begin
\needspace{3\baselineskip}
\textbf{āgān} (v.); past ptcp. āgān. go away. Supplied gone-away form; distinct from āgan possess.\par
% glossary-block:B00029:end
% glossary-block:B00030:begin
\needspace{3\baselineskip}
\textbf{āgēan} — citation/spelling variant: see \textbf{ongēan}.\par
% glossary-block:B00030:end
% glossary-block:B00031:begin
\needspace{3\baselineskip}
\textbf{āgen} (adj.). own.\par
% glossary-block:B00031:end
% glossary-block:B00032:begin
\needspace{3\baselineskip}
\textbf{āgiefan} (v.). give back; restore.\par
% glossary-block:B00032:end
% glossary-block:B00033:begin
\needspace{3\baselineskip}
\textbf{āginnan} (v.). begin. Supplied source form began.\par
% glossary-block:B00033:end
% glossary-block:B00034:begin
\needspace{3\baselineskip}
\textbf{āgyldan} (v.); pres. 3 sg. āgylt. reward; repay.\par
% glossary-block:B00034:end
% glossary-block:B00035:begin
\needspace{3\baselineskip}
\textbf{āh} — pres. ind. 1/3 sg.: see \textbf{āgan}.\par
% glossary-block:B00035:end
% glossary-block:B00036:begin
\needspace{3\baselineskip}
\textbf{āhębban} (v. strong); pret. 1 sg. āhōf; pret. pl. āhōfon; past ptcp. āhafen; pres. 3 sg. āhęfð. raise; lift.\par
% glossary-block:B00036:end
% glossary-block:B00037:begin
\needspace{3\baselineskip}
\textbf{āhęfð} — pres. ind. 3 sg.: see \textbf{āhębban}.\par
% glossary-block:B00037:end
% glossary-block:B00038:begin
\needspace{3\baselineskip}
\textbf{āhrēosan} (v. strong); pret. 1 sg. āhrēas; pret. pl. āhruron; past ptcp. āhroren; pres. 3 sg. āhrīest. fall; fall upon.\par
% glossary-block:B00038:end
% glossary-block:B00039:begin
\needspace{3\baselineskip}
\textbf{āhrīest} — pres. ind. 3 sg.: see \textbf{āhrēosan}.\par
% glossary-block:B00039:end
% glossary-block:B00040:begin
\needspace{3\baselineskip}
\textbf{āhte} — pret. sg.: see \textbf{āgan}.\par
% glossary-block:B00040:end
% glossary-block:B00041:begin
\needspace{3\baselineskip}
\textbf{āhwā} (indefinite pronoun). anyone. Reference in a source or explanatory note.\par
% glossary-block:B00041:end
% glossary-block:B00042:begin
\needspace{3\baselineskip}
\textbf{āhwæðer} (indefinite pronoun/adjective). either; one of two. Reference in a source or explanatory note.\par
% glossary-block:B00042:end
% glossary-block:B00043:begin
\needspace{3\baselineskip}
\textbf{āhyldan} (v.). lay; incline.\par
% glossary-block:B00043:end
% glossary-block:B00044:begin
\needspace{3\baselineskip}
\textbf{ālīefan} (v.); pres. 3 sg. ālīefð. permit; allow. Impersonal permission frame, with dative permitted person; consult the complete supplied expression.\par
% glossary-block:B00044:end
% glossary-block:B00045:begin
\needspace{3\baselineskip}
\textbf{ālīehtan} (v. weak); pret. 1 sg. ālīehte; past ptcp. ālīeht. illuminate; give light to. Accusative person or thing illuminated.\par
% glossary-block:B00045:end
% glossary-block:B00046:begin
\needspace{3\baselineskip}
\textbf{ālīesan} (v. weak); pret. 1 sg. ālīesde; past ptcp. ālīesed. deliver; release. Accusative person delivered or released; fram + dative marks separation.\par
% glossary-block:B00046:end
% glossary-block:B00047:begin
\needspace{3\baselineskip}
\textbf{ālīesend} (n. m.); gen. sg. ālīesendes; nom. pl. ālīesend. Redeemer; one who delivers.\par
% glossary-block:B00047:end
% glossary-block:B00048:begin
\needspace{3\baselineskip}
\textbf{ālūtan} (v. strong); pret. 1 sg. ālēat; pret. pl. āluton; past ptcp. āloten; pres. 3 sg. ālȳtt. bow; incline oneself. Reflexive meaning in the assigned use.\par
% glossary-block:B00048:end
% glossary-block:B00049:begin
\needspace{3\baselineskip}
\textbf{ālȳtt} — pres. ind. 3 sg.: see \textbf{ālūtan}.\par
% glossary-block:B00049:end
% glossary-block:B00050:begin
\needspace{3\baselineskip}
\textbf{Amen} (interj.). Amen.\par
% glossary-block:B00050:end
% glossary-block:B00051:begin
\needspace{3\baselineskip}
\textbf{ān} (numeral; adj.). one. alone. a certain. Introduces an unnamed singular individual; not an obligatory indefinite article. one; single.\par
% glossary-block:B00051:end
% glossary-block:B00052:begin
\needspace{3\baselineskip}
\textbf{āncęnned} (adj.). only-begotten.\par
% glossary-block:B00052:end
% glossary-block:B00053:begin
\needspace{3\baselineskip}
\textbf{and} (conj.). and. Simple coordination.\par
% glossary-block:B00053:end
% glossary-block:B00054:begin
\needspace{3\baselineskip}
\textbf{andgiet} (n. n.); gen. sg. andgietes. meaning; sense. understanding.\par
% glossary-block:B00054:end
% glossary-block:B00055:begin
\needspace{3\baselineskip}
\textbf{andgietfullice} (adv.); superl. andgietfullicost. In superlative andgietfullicost: most intelligibly.\par
% glossary-block:B00055:end
% glossary-block:B00056:begin
\needspace{3\baselineskip}
\textbf{andgietfullicost} — superlative adverb; exact source spelling: see \textbf{andgietfullice}.\par
% glossary-block:B00056:end
% glossary-block:B00057:begin
\needspace{3\baselineskip}
\textbf{andlīcnes} (n. f.); gen. sg. andlīcnesse; nom. pl. andlīcnessa. image.\par
% glossary-block:B00057:end
% glossary-block:B00058:begin
\needspace{3\baselineskip}
\textbf{andswarian} (v. weak); pret. 1 sg. andswarode; past ptcp. andswarod. answer; reply. Dative addressee; the reply may be given as direct speech.\par
% glossary-block:B00058:end
% glossary-block:B00059:begin
\needspace{3\baselineskip}
\textbf{andswerian} — citation/spelling variant: see \textbf{andswarian}.\par
% glossary-block:B00059:end
% glossary-block:B00060:begin
\needspace{3\baselineskip}
\textbf{andweard} (adj.). present.\par
% glossary-block:B00060:end
% glossary-block:B00061:begin
\needspace{3\baselineskip}
\textbf{andwlita} (n. m.); gen. sg. andwlitan; nom. pl. andwlitan. face; countenance.\par
% glossary-block:B00061:end
% glossary-block:B00062:begin
\needspace{3\baselineskip}
\textbf{angel} (n. m.); gen. sg. angles; nom. pl. anglas. fish-hook.\par
% glossary-block:B00062:end
% glossary-block:B00063:begin
\needspace{3\baselineskip}
\textbf{anginn} (n. n.); gen. sg. anginnes; nom. pl. anginn. beginning.\par
% glossary-block:B00063:end
% glossary-block:B00064:begin
\needspace{3\baselineskip}
\textbf{angrīslic} (adj.). fierce; raging.\par
% glossary-block:B00064:end
% glossary-block:B00065:begin
\needspace{3\baselineskip}
\textbf{ann} — pres. ind. 1/3 sg.: see \textbf{unnan}.\par
% glossary-block:B00065:end
% glossary-block:B00066:begin
\needspace{3\baselineskip}
\textbf{ansīen} (n. f.); gen. sg. ansīene. appearance.\par
% glossary-block:B00066:end
% glossary-block:B00067:begin
\needspace{3\baselineskip}
\textbf{Apollines} (proper name; source nom. sg.). Apollo. The god, distinct from Apollonius. Apollo, distinct from Apollonius.\par
% glossary-block:B00067:end
% glossary-block:B00068:begin
\needspace{3\baselineskip}
\textbf{Apollonius} (proper name). Apollonius. Supplied dative person in the assigned source expression. Apollonius.\par
% glossary-block:B00068:end
% glossary-block:B00069:begin
\needspace{3\baselineskip}
\textbf{ār} (n. f.). honour; high position. Contextual rank or office, not a new general meaning of every ār form.\par
% glossary-block:B00069:end
% glossary-block:B00070:begin
\needspace{3\baselineskip}
\textbf{Arcestrates} (proper name). Arcestrates. King of the region in the story; supplied name forms. Arcestrates. The king; supplied name forms.\par
% glossary-block:B00070:end
% glossary-block:B00071:begin
\needspace{3\baselineskip}
\textbf{ār-fæstnes} (n. f.). kindness.\par
% glossary-block:B00071:end
% glossary-block:B00072:begin
\needspace{3\baselineskip}
\textbf{ārīsan} (v. strong); pret. 1 sg. ārās; pret. pl. ārison; past ptcp. ārisen. arise. Intransitive arise; of + dative source in the bounded rise-from-death expression.\par
% glossary-block:B00072:end
% glossary-block:B00073:begin
\needspace{3\baselineskip}
\textbf{arn} — pret. sg.: see \textbf{iernan}.\par
% glossary-block:B00073:end
% glossary-block:B00074:begin
\needspace{3\baselineskip}
\textbf{āscian} (v. weak); pret. 1 sg. āscode; past ptcp. āscod. ask about. Accusative person asked, genitive matter asked about; an indirect question may express the matter. ask about. Supplied source form asked about.\par
% glossary-block:B00074:end
% glossary-block:B00075:begin
\needspace{3\baselineskip}
\textbf{āsęndan} (v. weak); pret. 1 sg. āsęnde; past ptcp. āsęnd. send.\par
% glossary-block:B00075:end
% glossary-block:B00076:begin
\needspace{3\baselineskip}
\textbf{āsęttan} (v. weak); pret. 1 sg. āsętte; past ptcp. āsętt. set; place.\par
% glossary-block:B00076:end
% glossary-block:B00077:begin
\needspace{3\baselineskip}
\textbf{āsmēan} (v.). conceive. Exact supplied infinitive āsmēan in the heart-conceiving expression; not an independently assigned full paradigm.\par
% glossary-block:B00077:end
% glossary-block:B00078:begin
\needspace{3\baselineskip}
\textbf{assa} (n. m.); gen. sg. assan; nom. pl. assan. ass; donkey.\par
% glossary-block:B00078:end
% glossary-block:B00079:begin
\needspace{3\baselineskip}
\textbf{Asser} (proper name). Asser.\par
% glossary-block:B00079:end
% glossary-block:B00080:begin
\needspace{3\baselineskip}
\textbf{āstīgan} (v. strong); pret. 1 sg. āstāg; pret. pl. āstigon; past ptcp. āstigen. go aboard; board. Accusative ship boarded in the assigned frame. descend. Intransitive descend; of + dative departure point.\par
% glossary-block:B00080:end
% glossary-block:B00081:begin
\needspace{3\baselineskip}
\textbf{āstreahte} — pret. sg.: see \textbf{āstręccean}.\par
% glossary-block:B00081:end
% glossary-block:B00082:begin
\needspace{3\baselineskip}
\textbf{āstręccean} (v. weak); pret. 1 sg. āstreahte; past ptcp. āstreaht. prostrate oneself. Reflexive accusative; the before phrase retains its own case/construal.\par
% glossary-block:B00082:end
% glossary-block:B00083:begin
\needspace{3\baselineskip}
\textbf{astyrian} (v.). move; set in motion. Supplied source sense.\par
% glossary-block:B00083:end
% glossary-block:B00084:begin
\needspace{3\baselineskip}
\textbf{āð} (n. m.); gen. sg. āðes; nom. pl. āðas. oath.\par
% glossary-block:B00084:end
% glossary-block:B00085:begin
\needspace{3\baselineskip}
\textbf{āðęnian} (v.). direct; stretch toward. Directing the mind toward what is heard; supplied figurative phrase.\par
% glossary-block:B00085:end
% glossary-block:B00086:begin
\needspace{3\baselineskip}
\textbf{āðindan} (v.). swollen. Supplied weak participial adjective describing the sea.\par
% glossary-block:B00086:end
% glossary-block:B00087:begin
\needspace{3\baselineskip}
\textbf{āwęg} — citation/spelling variant: see \textbf{onwęg}.\par
% glossary-block:B00087:end
% glossary-block:B00088:begin
\needspace{3\baselineskip}
\textbf{āwęndednes} (n. f.); gen. sg. āwęndednesse; nom. pl. āwęndednessa. translation.\par
% glossary-block:B00088:end
% glossary-block:B00089:begin
\needspace{3\baselineskip}
\textbf{āweorpan} (v. strong); pret. 1 sg. āwearp; pret. pl. āwurpon; past ptcp. āworpen; pres. 3 sg. āwierpð. cast down; cast away.\par
% glossary-block:B00089:end
% glossary-block:B00090:begin
\needspace{3\baselineskip}
\textbf{āwierpð} — pres. ind. 3 sg.: see \textbf{āweorpan}.\par
% glossary-block:B00090:end
% glossary-block:B00091:begin
\needspace{3\baselineskip}
\textbf{āwiht} (indefinite substantive). anything.\par
% glossary-block:B00091:end
% glossary-block:B00092:begin
\needspace{3\baselineskip}
\textbf{āwrītan} (v. strong); pret. 1 sg. āwrāt; pret. pl. āwriton; past ptcp. āwriten. write.\par
% glossary-block:B00092:end
% glossary-block:B00093:begin
\needspace{3\baselineskip}
\textbf{āwyrcean} (v.). do; deserve. In the source, deserving through one's conduct; retain the complete local frame.\par
% glossary-block:B00093:end
% book-navigation: standalone bookmark replaced by volume hierarchy
% glossary-block:B00094:begin
\needspace{5\baselineskip}
\subsection*{Æ}\BookHeading{nav-glossary-old-english-0003}{subsubsection}{Æ}{0}
% glossary-block:B00094:end
% glossary-block:B00095:begin
\textbf{ǣ} (n. f.). law.\par
% glossary-block:B00095:end
% glossary-block:B00096:begin
\needspace{3\baselineskip}
\textbf{æcer} (n. m.); gen. sg. æceres; nom. pl. æceras. field.\par
% glossary-block:B00096:end
% glossary-block:B00097:begin
\needspace{3\baselineskip}
\textbf{æcern} (n. n.); gen. sg. æcernes; nom. pl. æcernu. acorn.\par
% glossary-block:B00097:end
% glossary-block:B00098:begin
\needspace{3\baselineskip}
\textbf{ǣfen} (n. n.); gen. sg. ǣfenes; nom. pl. ǣfenu. evening.\par
% glossary-block:B00098:end
% glossary-block:B00099:begin
\needspace{3\baselineskip}
\textbf{ǣfre} (adv.). ever; always.\par
% glossary-block:B00099:end
% glossary-block:B00100:begin
\needspace{3\baselineskip}
\textbf{æfter} (prep.; conj. in fixed expressions). according to. With dative. after. Temporal preposition with dative. after. In the supplied finite-clause conjunction æfter ðǣm ðe. after; in search of. Person sought or inquired after; keep this lexical frame distinct from temporal after. after; following. Dative person followed, including motion after that person.\par
% glossary-block:B00100:end
% glossary-block:B00101:begin
\needspace{3\baselineskip}
\textbf{ǣghwā} (indefinite pronoun). everyone; each one. Reference in a source or explanatory note.\par
% glossary-block:B00101:end
% glossary-block:B00102:begin
\needspace{3\baselineskip}
\textbf{ǣghwæðer} (indefinite pronoun/adjective). each of two; both. Reference in a source or explanatory note.\par
% glossary-block:B00102:end
% glossary-block:B00103:begin
\needspace{3\baselineskip}
\textbf{ǣghwilc} (indefinite adjective). every; each. each; every.\par
% glossary-block:B00103:end
% glossary-block:B00104:begin
\needspace{3\baselineskip}
\textbf{ǣgðer} (coordinating adverb). both. In ǣgðer ge … ge.\par
% glossary-block:B00104:end
% glossary-block:B00105:begin
\needspace{3\baselineskip}
\textbf{ǣht} (n. f.). possessions; goods.\par
% glossary-block:B00105:end
% glossary-block:B00106:begin
\needspace{3\baselineskip}
\textbf{ǣl} (n. f.); gen. sg. ǣle; nom. pl. ǣla. awl.\par
% glossary-block:B00106:end
% glossary-block:B00107:begin
\needspace{3\baselineskip}
\textbf{ǣlc} (adj.). each; every. With singular noun; substantival each follows the case required by its clause.\par
% glossary-block:B00107:end
% glossary-block:B00108:begin
\needspace{3\baselineskip}
\textbf{Ælfred} (proper name); gen. sg. Ælfredes. Alfred.\par
% glossary-block:B00108:end
% glossary-block:B00109:begin
\needspace{3\baselineskip}
\textbf{ælmihtig} (adj.). almighty.\par
% glossary-block:B00109:end
% glossary-block:B00110:begin
\needspace{3\baselineskip}
\textbf{ǣmetta} (n. m.). leisure; available time.\par
% glossary-block:B00110:end
% glossary-block:B00111:begin
\needspace{3\baselineskip}
\textbf{ǣmtig} (adj.). empty; void.\par
% glossary-block:B00111:end
% glossary-block:B00112:begin
\needspace{3\baselineskip}
\textbf{ǣnig} (indefinite adjective). any.\par
% glossary-block:B00112:end
% glossary-block:B00113:begin
\needspace{3\baselineskip}
\textbf{ǣnne} — acc. masc. sg., variant: see \textbf{ān}.\par
% glossary-block:B00113:end
% glossary-block:B00114:begin
\needspace{3\baselineskip}
\textbf{ænne} — acc. masc. sg., variant: see \textbf{ān}.\par
% glossary-block:B00114:end
% glossary-block:B00115:begin
\needspace{3\baselineskip}
\textbf{æppelbǣre} (adj.). fruit-bearing.\par
% glossary-block:B00115:end
% glossary-block:B00116:begin
\needspace{3\baselineskip}
\textbf{ǣr} (adv.). before; earlier.\par
% glossary-block:B00116:end
% glossary-block:B00117:begin
\needspace{3\baselineskip}
\textbf{ǣr} (conj.). before.\par
% glossary-block:B00117:end
% glossary-block:B00118:begin
\needspace{3\baselineskip}
\textbf{ǣr} (n. n.); gen. sg. ǣres; nom. pl. ǣr. copper.\par
% glossary-block:B00118:end
% glossary-block:B00119:begin
\needspace{3\baselineskip}
\textbf{ǣrendwreca} (n. m.). messengers; envoys.\par
% glossary-block:B00119:end
% glossary-block:B00120:begin
\needspace{3\baselineskip}
\textbf{ǣrest} (adv.). first; at first.\par
% glossary-block:B00120:end
% glossary-block:B00121:begin
\needspace{3\baselineskip}
\textbf{ǣresta} (ordinal numeral). first.\par
% glossary-block:B00121:end
% glossary-block:B00122:begin
\needspace{3\baselineskip}
\textbf{ǣrra} (adj.). former; earlier.\par
% glossary-block:B00122:end
% glossary-block:B00123:begin
\needspace{3\baselineskip}
\textbf{ǣswician} (v. weak); pret. 1 sg. ǣswicode; past ptcp. ǣswicod. offend; cause to stumble. Accusative person; ethical/metaphorical offence in the assigned use.\par
% glossary-block:B00123:end
% glossary-block:B00124:begin
\needspace{3\baselineskip}
\textbf{æt} (prep.). at. from. æt + dative person with learn. from; out of. Proceeds of work, in the assigned expression.\par
% glossary-block:B00124:end
% glossary-block:B00125:begin
\needspace{3\baselineskip}
\textbf{ætberstan} (v.). escape.\par
% glossary-block:B00125:end
% glossary-block:B00126:begin
\needspace{3\baselineskip}
\textbf{ætēowian} (v. weak); pret. 1 sg. ætēowode; past ptcp. ætēowod. appear.\par
% glossary-block:B00126:end
% glossary-block:B00127:begin
\needspace{3\baselineskip}
\textbf{ætforan} (prep.). before; in front of. Dative landmark/person in the supplied frame; retain ætforan rather than replacing it with beforan.\par
% glossary-block:B00127:end
% glossary-block:B00128:begin
\needspace{3\baselineskip}
\textbf{ætgædere} (adv.). together; at once. Paired with samod in the supplied source.\par
% glossary-block:B00128:end
% glossary-block:B00129:begin
\needspace{3\baselineskip}
\textbf{Æthelwulf} (proper name). Æthelwulf.\par
% glossary-block:B00129:end
% glossary-block:B00130:begin
\needspace{3\baselineskip}
\textbf{ætīewan} (v. weak); pret. 1 sg. ætīewde; past ptcp. ætīewed. show; display; reveal.\par
% glossary-block:B00130:end
% glossary-block:B00131:begin
\needspace{3\baselineskip}
\textbf{æðelboren} (adj.). high-born.\par
% glossary-block:B00131:end
% glossary-block:B00132:begin
\needspace{3\baselineskip}
\textbf{æðele} (adj.). noble people. Substantival adjective in the source. noble; illustrious.\par
% glossary-block:B00132:end
% glossary-block:B00133:begin
\needspace{3\baselineskip}
\textbf{Æðelðrȳð} (proper name). Etheldred; Æthelthryth. Proper name.\par
% glossary-block:B00133:end
% glossary-block:B00134:begin
\needspace{3\baselineskip}
\textbf{Æðelwald} (proper name). Æthelwald. Proper name in apposition to father.\par
% glossary-block:B00134:end
% glossary-block:B00135:begin
\needspace{3\baselineskip}
\textbf{æx} (n. f.); gen. sg. æxe; nom. pl. æxa. ax.\par
% glossary-block:B00135:end
% book-navigation: standalone bookmark replaced by volume hierarchy
% glossary-block:B00136:begin
\needspace{5\baselineskip}
\subsection*{B}\BookHeading{nav-glossary-old-english-0004}{subsubsection}{B}{0}
% glossary-block:B00136:end
% glossary-block:B00137:begin
\textbf{bā} — nom./acc. fem./neut. pl.: see \textbf{bēgen}.\par
% glossary-block:B00137:end
% glossary-block:B00138:begin
\needspace{3\baselineskip}
\textbf{bām} — dat. pl.: see \textbf{bēgen}.\par
% glossary-block:B00138:end
% glossary-block:B00139:begin
\needspace{3\baselineskip}
\textbf{bān} (n. n.); gen. sg. bānes; nom. pl. bān. bone; ivory. Ivory is elpes bān, literally elephant's bone.\par
% glossary-block:B00139:end
% glossary-block:B00140:begin
\needspace{3\baselineskip}
\textbf{baðian} (v.); pret. 1 sg. baðode. bathe. Supplied source form bathed.\par
% glossary-block:B00140:end
% glossary-block:B00141:begin
\needspace{3\baselineskip}
\textbf{bǣm} — dat. pl., variant: see \textbf{bēgen}.\par
% glossary-block:B00141:end
% glossary-block:B00142:begin
\needspace{3\baselineskip}
\textbf{bæð} (n. n.); gen. sg. bæðes; nom. pl. baðu. bath.\par
% glossary-block:B00142:end
% glossary-block:B00143:begin
\needspace{3\baselineskip}
\textbf{bæðstęde} (n. m.). bathing place; gymnasium. Gymnasium is the contextual interpretation in the supplied reading.\par
% glossary-block:B00143:end
% glossary-block:B00144:begin
\needspace{3\baselineskip}
\textbf{be} (prep.). according to. With dative. concerning; about. With dative. by that; from that. Refers to the preceding proposition; supplied expression. word by word. In word be worde; literal correspondence. by. Body part by which a person is led; supplied grasping phrase. in the south. Only in supplied be sūðan. by that; on that evidence. Evidence/inference frame, not an unrestricted source preposition. by permission. Supplied permission expression.\par
% glossary-block:B00144:end
% glossary-block:B00145:begin
\needspace{3\baselineskip}
\textbf{beald} (adj.). bold; brave.\par
% glossary-block:B00145:end
% glossary-block:B00146:begin
\needspace{3\baselineskip}
\textbf{bearn} (n. n.); gen. sg. bearnes; nom. pl. bearn. son; child.\par
% glossary-block:B00146:end
% glossary-block:B00147:begin
\needspace{3\baselineskip}
\textbf{bēatan} (v. strong); pret. 1 sg. bēot; pret. pl. bēoton; past ptcp. bēaten; pres. 3 sg. bīett. beat; strike. Accusative person or thing struck.\par
% glossary-block:B00147:end
% glossary-block:B00148:begin
\needspace{3\baselineskip}
\textbf{bebēodan} (v. strong); pret. 1 sg. bebēad; pret. pl. bebudon; past ptcp. beboden; pres. 3 sg. bebīett. command; instruct. Dative addressee and a ðæt clause with subjunctive for the commanded action.\par
% glossary-block:B00148:end
% glossary-block:B00149:begin
\needspace{3\baselineskip}
\textbf{bebīett} — pres. ind. 3 sg.: see \textbf{bebēodan}.\par
% glossary-block:B00149:end
% glossary-block:B00150:begin
\needspace{3\baselineskip}
\textbf{bebod} (n. n.); gen. sg. bebodes; nom. pl. bebodu. commandment.\par
% glossary-block:B00150:end
% glossary-block:B00151:begin
\needspace{3\baselineskip}
\textbf{bēc} — gen./dat. sg.; nom./acc. pl.: see \textbf{bōc}.\par
% glossary-block:B00151:end
% glossary-block:B00152:begin
\needspace{3\baselineskip}
\textbf{bēcrǣding} (n. f.); gen. sg. bēcrǣdinge; nom. pl. bēcrǣdinga. reading books.\par
% glossary-block:B00152:end
% glossary-block:B00153:begin
\needspace{3\baselineskip}
\textbf{becuman} (v.). come; arrive. occur; come about. Of the calm in the source. come upon; befall. Retain the source's affected participant.\par
% glossary-block:B00153:end
% glossary-block:B00154:begin
\needspace{3\baselineskip}
\textbf{beēode} — pret. sg.: see \textbf{begān}.\par
% glossary-block:B00154:end
% glossary-block:B00155:begin
\needspace{3\baselineskip}
\textbf{befeallan} (v.); pres. 3 sg. befielð. fall in love. Supplied idiom makes lufu love the subject, its possessor the lover, and on + accusative the beloved.\par
% glossary-block:B00155:end
% glossary-block:B00156:begin
\needspace{3\baselineskip}
\textbf{befielð} — pres. ind. 3 sg.: see \textbf{befeallan}.\par
% glossary-block:B00156:end
% glossary-block:B00157:begin
\needspace{3\baselineskip}
\textbf{be-foran} (prep.). before; in front of. Dative landmark in the static frame. before; ahead of. Dative person, including going or sending ahead of that person. before; in front of. Accusative landmark in the positioning construal; static dative can be a different valid construal.\par
% glossary-block:B00157:end
% glossary-block:B00158:begin
\needspace{3\baselineskip}
\textbf{begān} (v. irregular); pres. 1 sg. begā; pres. 3 sg. begǣð; pres. pl. begāð; pret. 1 sg. beēode; pret. pl. beēodon; past ptcp. begān. ply; practise. Accusative trade or occupation. pursue; practise. Study is the thing pursued.\par
% glossary-block:B00158:end
% glossary-block:B00159:begin
\needspace{3\baselineskip}
\textbf{begǣð} — pres. ind. 3 sg.: see \textbf{begān}.\par
% glossary-block:B00159:end
% glossary-block:B00160:begin
\needspace{3\baselineskip}
\textbf{bēgen} (numeral). both.\par
% glossary-block:B00160:end
% glossary-block:B00161:begin
\needspace{3\baselineskip}
\textbf{begyrdan} (v.); past ptcp. begyrd. gird; clothe. Supplied participial clothing expression.\par
% glossary-block:B00161:end
% glossary-block:B00162:begin
\needspace{3\baselineskip}
\textbf{behātan} (v.); pret. 1 sg. behēt. promise.\par
% glossary-block:B00162:end
% glossary-block:B00163:begin
\needspace{3\baselineskip}
\textbf{behēafdian} (v. weak); pret. 1 sg. behēafdode; past ptcp. behēafdod. behead. Accusative person beheaded.\par
% glossary-block:B00163:end
% glossary-block:B00164:begin
\needspace{3\baselineskip}
\textbf{behealdan} (v. strong); pret. 1 sg. behēold; pret. pl. behēoldon; past ptcp. behealden; pres. 3 sg. behielt. behold; look at. Accusative object seen. observe. Supplied source form observed.\par
% glossary-block:B00164:end
% glossary-block:B00165:begin
\needspace{3\baselineskip}
\textbf{behēfe} (adj.). useful. Dative person to whom something is useful.\par
% glossary-block:B00165:end
% glossary-block:B00166:begin
\needspace{3\baselineskip}
\textbf{behielt} — pres. ind. 3 sg.: see \textbf{behealdan}.\par
% glossary-block:B00166:end
% glossary-block:B00167:begin
\needspace{3\baselineskip}
\textbf{behygdig} (adj.). attentive; thoughtful.\par
% glossary-block:B00167:end
% glossary-block:B00168:begin
\needspace{3\baselineskip}
\textbf{bęlg} (n. m.); gen. sg. bęlges; nom. pl. bęlgas. bellows. Selected plural instrument use.\par
% glossary-block:B00168:end
% glossary-block:B00169:begin
\needspace{3\baselineskip}
\textbf{belimpan} (v.). belong; relate. Supplied participial construction.\par
% glossary-block:B00169:end
% glossary-block:B00170:begin
\needspace{3\baselineskip}
\textbf{bēn} (n. f.). prayer; petition. Source dative plural after tō. petition; appeal.\par
% glossary-block:B00170:end
% glossary-block:B00171:begin
\needspace{3\baselineskip}
\textbf{Beniamin} (proper name); gen. sg. Beniamines. Benjamin.\par
% glossary-block:B00171:end
% glossary-block:B00172:begin
\needspace{3\baselineskip}
\textbf{bēodan} (v. strong); pret. 1 sg. bēad; pret. pl. budon; past ptcp. boden; pres. 3 sg. bīett. offer. command; bid.\par
% glossary-block:B00172:end
% glossary-block:B00173:begin
\needspace{3\baselineskip}
\textbf{bēon} — alternative infinitive: see \textbf{wesan}.\par
% glossary-block:B00173:end
% glossary-block:B00174:begin
\needspace{3\baselineskip}
\textbf{beorhtnes} (n. f.); gen. sg. beorhtnesse. brightness.\par
% glossary-block:B00174:end
% glossary-block:B00175:begin
\needspace{3\baselineskip}
\textbf{beorn} (n. m.); gen. sg. beornes; nom. pl. beornas. warrior.\par
% glossary-block:B00175:end
% glossary-block:B00176:begin
\needspace{3\baselineskip}
\textbf{bēorscipe} (n. m.); gen. sg. bēorscipes. feast; banquet.\par
% glossary-block:B00176:end
% glossary-block:B00177:begin
\needspace{3\baselineskip}
\textbf{beran} (v. strong); pret. 1 sg. bær; pret. pl. bǣron; past ptcp. boren; pres. 3 sg. birð. carry. Accusative person or thing carried. bear; produce. Accusative produce; distinct from physical carrying. rich in; productive of. Supplied berende on + dative ore expression; participial adjective.\par
% glossary-block:B00177:end
% glossary-block:B00178:begin
\needspace{3\baselineskip}
\textbf{berēafian} (v.); past ptcp. berēafod. despoil; deprive of. Accusative person deprived, genitive thing lost; passive retains that genitive.\par
% glossary-block:B00178:end
% glossary-block:B00179:begin
\needspace{3\baselineskip}
\textbf{berstan} (v. strong); pret. 1 sg. bærst; pret. pl. burston; past ptcp. borsten; pres. 3 sg. bierst. burst.\par
% glossary-block:B00179:end
% glossary-block:B00180:begin
\needspace{3\baselineskip}
\textbf{besēon} (v.). look; turn one's gaze. Supplied source form and frame. look toward. Target of the look follows the supplied phrase.\par
% glossary-block:B00180:end
% glossary-block:B00181:begin
\needspace{3\baselineskip}
\textbf{bęt} — comparative adverb: see \textbf{wel}.\par
% glossary-block:B00181:end
% glossary-block:B00182:begin
\needspace{3\baselineskip}
\textbf{bętera} (adj.). better. Comparative of gōd.\par
% glossary-block:B00182:end
% glossary-block:B00183:begin
\needspace{3\baselineskip}
\textbf{bętst} (adj.). best. Superlative of gōd. Also superlative adverb: see \textbf{wel}.\par
% glossary-block:B00183:end
% glossary-block:B00184:begin
\needspace{3\baselineskip}
\textbf{betwēonan} (adverb or particle). between; to one another. Supplied reciprocal betwēonan him in 45.6.\par
% glossary-block:B00184:end
% glossary-block:B00185:begin
\needspace{3\baselineskip}
\textbf{betweox} (prep.). among. Accusative in the exact supplied phrases; not a universal case rule. among; between.\par
% glossary-block:B00185:end
% glossary-block:B00186:begin
\needspace{3\baselineskip}
\textbf{beðearf} — pres. ind. 1/3 sg.: see \textbf{beðurfan}.\par
% glossary-block:B00186:end
% glossary-block:B00187:begin
\needspace{3\baselineskip}
\textbf{beðorfte} — pret. sg.: see \textbf{beðurfan}.\par
% glossary-block:B00187:end
% glossary-block:B00188:begin
\needspace{3\baselineskip}
\textbf{beðurfan} (v. preterite-present); pres. 1/3 sg. beðearf; pres. pl. beðurfon; pret. 1 sg. beðorfte. need; require. Genitive thing needed. No past-participle citation is given in the checked source paradigm.\par
% glossary-block:B00188:end
% glossary-block:B00189:begin
\needspace{3\baselineskip}
\textbf{bewǣfan} (v.); past ptcp. bewǣfed. wrap; clothe. Supplied participial clothing expression.\par
% glossary-block:B00189:end
% glossary-block:B00190:begin
\needspace{3\baselineskip}
\textbf{bewęndan} (v. weak); pret. 1 sg. bewęnde; past ptcp. bewęnd. turn. The assigned reflexive frame uses dative him and tō + dative goal.\par
% glossary-block:B00190:end
% glossary-block:B00191:begin
\needspace{3\baselineskip}
\textbf{biddan} (v. strong); pret. 1 sg. bæd; pret. pl. bǣdon; past ptcp. beden. ask; request. Genitive thing requested; an accusative person may be added. beseech; entreat. Accusative person; the requested action can be a ðæt clause. wish well; bid farewell. Supplied farewell formula.\par
% glossary-block:B00191:end
% glossary-block:B00192:begin
\needspace{3\baselineskip}
\textbf{bīegan} (v. weak); pret. 1 sg. bīegde; past ptcp. bīeged. bend. Supplied past bent; retain the source's unmarked ie in its actual form. bend.\par
% glossary-block:B00192:end
% glossary-block:B00193:begin
\needspace{3\baselineskip}
\textbf{biema} (n. m.); gen. sg. bieman; nom. pl. bieman. trumpet; clarion.\par
% glossary-block:B00193:end
% glossary-block:B00194:begin
\needspace{3\baselineskip}
\textbf{bierst} — pres. ind. 2 sg., variant: see \textbf{beran}; pres. ind. 3 sg.: see \textbf{berstan}.\par
% glossary-block:B00194:end
% glossary-block:B00195:begin
\needspace{3\baselineskip}
\textbf{bierð} — pres. ind. 3 sg., variant: see \textbf{beran}.\par
% glossary-block:B00195:end
% glossary-block:B00196:begin
\needspace{3\baselineskip}
\textbf{bīet} — pres. ind. 3 sg., variant: see \textbf{bēodan}.\par
% glossary-block:B00196:end
% glossary-block:B00197:begin
\needspace{3\baselineskip}
\textbf{bīett} — pres. ind. 3 sg.: see \textbf{bēatan}; pres. ind. 3 sg.: see \textbf{bēodan}.\par
% glossary-block:B00197:end
% glossary-block:B00198:begin
\needspace{3\baselineskip}
\textbf{bindan} (v. strong); pret. 1 sg. band; pret. pl. bundon; past ptcp. bunden. bind; tie. Accusative person or thing bound.\par
% glossary-block:B00198:end
% glossary-block:B00199:begin
\needspace{3\baselineskip}
\textbf{binnan} (prep.). within.\par
% glossary-block:B00199:end
% glossary-block:B00200:begin
\needspace{3\baselineskip}
\textbf{birð} — pres. ind. 3 sg.: see \textbf{beran}.\par
% glossary-block:B00200:end
% glossary-block:B00201:begin
\needspace{3\baselineskip}
\textbf{bisceop} (n. m.); gen. sg. bisceopes; nom. pl. bisceopas. bishop.\par
% glossary-block:B00201:end
% glossary-block:B00202:begin
\needspace{3\baselineskip}
\textbf{bisgian} (v.). occupy; trouble.\par
% glossary-block:B00202:end
% glossary-block:B00203:begin
\needspace{3\baselineskip}
\textbf{bisgu} (n. f.). occupations; cares.\par
% glossary-block:B00203:end
% glossary-block:B00204:begin
\needspace{3\baselineskip}
\textbf{bist} — pres. ind. 2 sg.: see \textbf{wesan}.\par
% glossary-block:B00204:end
% glossary-block:B00205:begin
\needspace{3\baselineskip}
\textbf{bið} — pres. ind. 3 sg.: see \textbf{wesan}.\par
% glossary-block:B00205:end
% glossary-block:B00206:begin
\needspace{3\baselineskip}
\textbf{blāwan} (v. strong); pret. 1 sg. blēow; pret. pl. blēowon; past ptcp. blāwen; pres. 3 sg. blǣwð. blow. Intransitive in the assigned bellows sentence.\par
% glossary-block:B00206:end
% glossary-block:B00207:begin
\needspace{3\baselineskip}
\textbf{blæc} (adj.). black.\par
% glossary-block:B00207:end
% glossary-block:B00208:begin
\needspace{3\baselineskip}
\textbf{blǣwð} — pres. ind. 3 sg.: see \textbf{blāwan}.\par
% glossary-block:B00208:end
% glossary-block:B00209:begin
\needspace{3\baselineskip}
\textbf{blētsian} (v.); past ptcp. blētsod. bless. Supplied third-singular preterite; God is the understood subject. bless.\par
% glossary-block:B00209:end
% glossary-block:B00210:begin
\needspace{3\baselineskip}
\textbf{blētsung} (n. f.); gen. sg. blētsunge. blessing.\par
% glossary-block:B00210:end
% glossary-block:B00211:begin
\needspace{3\baselineskip}
\textbf{blinnan} (v.). cease. Supplied wæs blinnende periphrasis.\par
% glossary-block:B00211:end
% glossary-block:B00212:begin
\needspace{3\baselineskip}
\textbf{blissesang} (n. m.). song of joy.\par
% glossary-block:B00212:end
% glossary-block:B00213:begin
\needspace{3\baselineskip}
\textbf{blissian} (v. weak); pret. 1 sg. blissode; past ptcp. blissod. rejoice. Source genitive plural present-participle use, of people rejoicing. rejoice. Intransitive; mid + dative companion or ofer + accusative cause of joy in the assigned frames.\par
% glossary-block:B00213:end
% glossary-block:B00214:begin
\needspace{3\baselineskip}
\textbf{blīðe} (adj.). merry; glad.\par
% glossary-block:B00214:end
% glossary-block:B00215:begin
\needspace{3\baselineskip}
\textbf{blīðe} (adv.). joyfully; gladly.\par
% glossary-block:B00215:end
% glossary-block:B00216:begin
\needspace{3\baselineskip}
\textbf{blōstma} (n. m.); gen. sg. blōstman; nom. pl. blōstman. blossom.\par
% glossary-block:B00216:end
% glossary-block:B00217:begin
\needspace{3\baselineskip}
\textbf{blōstmian} (v. weak); pret. 1 sg. blōstmode; past ptcp. blōstmod. blossom; bloom.\par
% glossary-block:B00217:end
% glossary-block:B00218:begin
\needspace{3\baselineskip}
\textbf{bōc} (n. f.); gen. sg. bēc; nom. pl. bēc. book. Genitive singular variant bōce.\par
% glossary-block:B00218:end
% glossary-block:B00219:begin
\needspace{3\baselineskip}
\textbf{Bōclæden} (n. n.). book-Latin; Latin.\par
% glossary-block:B00219:end
% glossary-block:B00220:begin
\needspace{3\baselineskip}
\textbf{bohte} — pret. sg.: see \textbf{bycgan}.\par
% glossary-block:B00220:end
% glossary-block:B00221:begin
\needspace{3\baselineskip}
\textbf{brād} (adj.). broad.\par
% glossary-block:B00221:end
% glossary-block:B00222:begin
\needspace{3\baselineskip}
\textbf{brādnes} (n. f.); gen. sg. brādnesse; nom. pl. brādnessa. surface; face. Of the deep in the assigned expression.\par
% glossary-block:B00222:end
% glossary-block:B00223:begin
\needspace{3\baselineskip}
\textbf{brǣd} — pret. sg., contracted variant: see \textbf{bregdan}.\par
% glossary-block:B00223:end
% glossary-block:B00224:begin
\needspace{3\baselineskip}
\textbf{brǣdan} (v.); pres. 3 sg. brǣt. spread; extend. Supplied source form spreads.\par
% glossary-block:B00224:end
% glossary-block:B00225:begin
\needspace{3\baselineskip}
\textbf{bregdan} (v. strong); pret. 1 sg. brægd; pret. pl. brugdon; past ptcp. brogden; pres. 3 sg. brītt. brandish; draw.\par
% glossary-block:B00225:end
% glossary-block:B00226:begin
\needspace{3\baselineskip}
\textbf{Breoton} (proper name). Britain.\par
% glossary-block:B00226:end
% glossary-block:B00227:begin
\needspace{3\baselineskip}
\textbf{brēðer} — dat. sg.: see \textbf{brōðor}.\par
% glossary-block:B00227:end
% glossary-block:B00228:begin
\needspace{3\baselineskip}
\textbf{bringan} (v. weak); pret. 1 sg. brōhte; past ptcp. (ge)brōht. bring. Accusative thing or person brought; tō + dative destination/person, or the assigned bare dative recipient.\par
% glossary-block:B00228:end
% glossary-block:B00229:begin
\needspace{3\baselineskip}
\textbf{brīt} — pres. ind. 3 sg., variant: see \textbf{bregdan}.\par
% glossary-block:B00229:end
% glossary-block:B00230:begin
\needspace{3\baselineskip}
\textbf{brītt} — pres. ind. 3 sg.: see \textbf{bregdan}.\par
% glossary-block:B00230:end
% glossary-block:B00231:begin
\needspace{3\baselineskip}
\textbf{brōden} — past ptcp., contracted variant: see \textbf{bregdan}.\par
% glossary-block:B00231:end
% glossary-block:B00232:begin
\needspace{3\baselineskip}
\textbf{brōhte} — pret. sg.: see \textbf{bringan}.\par
% glossary-block:B00232:end
% glossary-block:B00233:begin
\needspace{3\baselineskip}
\textbf{brōðor} (n. m.); gen. sg. brōðor; nom. pl. brōðor. brother. Dative singular brēðer.\par
% glossary-block:B00233:end
% glossary-block:B00234:begin
\needspace{3\baselineskip}
\textbf{brycg} (n. f.); gen. sg. brycge; nom. pl. brycga. bridge.\par
% glossary-block:B00234:end
% glossary-block:B00235:begin
\needspace{3\baselineskip}
\textbf{brȳd} (n. f.); gen. sg. brȳde; nom. pl. brȳda. bride.\par
% glossary-block:B00235:end
% glossary-block:B00236:begin
\needspace{3\baselineskip}
\textbf{Bryttas} (ethnonym). Britons.\par
% glossary-block:B00236:end
% glossary-block:B00237:begin
\needspace{3\baselineskip}
\textbf{bū} — nom./acc. neut. pl.: see \textbf{bēgen}.\par
% glossary-block:B00237:end
% glossary-block:B00238:begin
\needspace{3\baselineskip}
\textbf{bufan} (prep.). above. With dative landmark in the taught spatial frame.\par
% glossary-block:B00238:end
% glossary-block:B00239:begin
\needspace{3\baselineskip}
\textbf{būgan} (v. strong); pret. 1 sg. bēag; pret. pl. bugon; past ptcp. bogen; pres. 3 sg. bȳgð. bend.\par
% glossary-block:B00239:end
% glossary-block:B00240:begin
\needspace{3\baselineskip}
\textbf{būr} (n. m.). chamber; room.\par
% glossary-block:B00240:end
% glossary-block:B00241:begin
\needspace{3\baselineskip}
\textbf{burg} (n. f.); gen. sg. byrig; nom. pl. byrig. town; city. Genitive singular variant burge.\par
% glossary-block:B00241:end
% glossary-block:B00242:begin
\needspace{3\baselineskip}
\textbf{būtan} (prep.). besides; apart from. Dative excluded person or thing. without. With dative. outside.\par
% glossary-block:B00242:end
% glossary-block:B00243:begin
\needspace{3\baselineskip}
\textbf{būtan} (conj.). except.\par
% glossary-block:B00243:end
% glossary-block:B00244:begin
\needspace{3\baselineskip}
\textbf{būtan} (conj.). unless; except if.\par
% glossary-block:B00244:end
% glossary-block:B00245:begin
\needspace{3\baselineskip}
\textbf{bycgan} (v. weak); pret. 1 sg. bohte; past ptcp. (ge)boht. buy. Accusative goods bought.\par
% glossary-block:B00245:end
% glossary-block:B00246:begin
\needspace{3\baselineskip}
\textbf{bȳgð} — pres. ind. 3 sg.: see \textbf{būgan}.\par
% glossary-block:B00246:end
% glossary-block:B00247:begin
\needspace{3\baselineskip}
\textbf{bȳhð} — pres. ind. 3 sg., variant: see \textbf{būgan}.\par
% glossary-block:B00247:end
% glossary-block:B00248:begin
\needspace{3\baselineskip}
\textbf{byrig} — gen./dat. sg.; nom./acc. pl.: see \textbf{burg}.\par
% glossary-block:B00248:end
% glossary-block:B00249:begin
\needspace{3\baselineskip}
\textbf{bȳsen} (n. f.). example; illustration.\par
% glossary-block:B00249:end
% book-navigation: standalone bookmark replaced by volume hierarchy
% glossary-block:B00250:begin
\needspace{5\baselineskip}
\subsection*{C}\BookHeading{nav-glossary-old-english-0005}{subsubsection}{C}{0}
% glossary-block:B00250:end
% glossary-block:B00251:begin
\textbf{campian} (v.). fight.\par
% glossary-block:B00251:end
% glossary-block:B00252:begin
\needspace{3\baselineskip}
\textbf{cann} — pres. ind. 1/3 sg.: see \textbf{cunnan}.\par
% glossary-block:B00252:end
% glossary-block:B00253:begin
\needspace{3\baselineskip}
\textbf{ceaf} (n. n.). chaff.\par
% glossary-block:B00253:end
% glossary-block:B00254:begin
\needspace{3\baselineskip}
\textbf{cealf} (n.). calf. Only this cited form is supplied.\par
% glossary-block:B00254:end
% glossary-block:B00255:begin
\needspace{3\baselineskip}
\textbf{cēapian} (v. weak); pret. 1 sg. cēapode; past ptcp. cēapod. trade; traffic. Traffic means trade; retain the supplied trading frame.\par
% glossary-block:B00255:end
% glossary-block:B00256:begin
\needspace{3\baselineskip}
\textbf{cēapmann} — citation/spelling variant: see \textbf{cīepmann}.\par
% glossary-block:B00256:end
% glossary-block:B00257:begin
\needspace{3\baselineskip}
\textbf{ceaster} (n. f.); gen. sg. ceastre; nom. pl. ceastra. town; city.\par
% glossary-block:B00257:end
% glossary-block:B00258:begin
\needspace{3\baselineskip}
\textbf{ceasterwaran} (n. m.); nom. pl. ceasterwaran. townspeople; citizens.\par
% glossary-block:B00258:end
% glossary-block:B00259:begin
\needspace{3\baselineskip}
\textbf{cēderbēam} (n. m.); gen. sg. cēderbēames; nom. pl. cēderbēamas. cedar; cedar tree.\par
% glossary-block:B00259:end
% glossary-block:B00260:begin
\needspace{3\baselineskip}
\textbf{cēlan} (v.). cool. Pronunciation example.\par
% glossary-block:B00260:end
% glossary-block:B00261:begin
\needspace{3\baselineskip}
\textbf{cęmpa} (n. m.); gen. sg. cęmpan; nom. pl. cęmpan. soldier.\par
% glossary-block:B00261:end
% glossary-block:B00262:begin
\needspace{3\baselineskip}
\textbf{ceorl} (n. m.); gen. sg. ceorles; nom. pl. ceorlas. layman.\par
% glossary-block:B00262:end
% glossary-block:B00263:begin
\needspace{3\baselineskip}
\textbf{cēosan} (v. strong); pret. 1 sg. cēas; pret. pl. curon; past ptcp. coren; pres. 3 sg. cīest. choose. Accusative person or thing chosen.\par
% glossary-block:B00263:end
% glossary-block:B00264:begin
\needspace{3\baselineskip}
\textbf{cīepa} (n.). merchant. Reference in a source or explanatory note.\par
% glossary-block:B00264:end
% glossary-block:B00265:begin
\needspace{3\baselineskip}
\textbf{cīepan} (v. weak); pret. 1 sg. cīepte; past ptcp. cīeped. sell. Accusative goods sold.\par
% glossary-block:B00265:end
% glossary-block:B00266:begin
\needspace{3\baselineskip}
\textbf{cīepmann} (n.). merchant. Reference in a source or explanatory note.\par
% glossary-block:B00266:end
% glossary-block:B00267:begin
\needspace{3\baselineskip}
\textbf{cīest} — pres. ind. 3 sg.: see \textbf{cēosan}.\par
% glossary-block:B00267:end
% glossary-block:B00268:begin
\needspace{3\baselineskip}
\textbf{cirice} (n. f.); gen. sg. cirican; nom. pl. cirican. church. oratory. Chad's separate place of prayer; distinct from the ordinary church sense.\par
% glossary-block:B00268:end
% glossary-block:B00269:begin
\needspace{3\baselineskip}
\textbf{clǣne} (adj.). pure; clean. clean; pure; blameless.\par
% glossary-block:B00269:end
% glossary-block:B00270:begin
\needspace{3\baselineskip}
\textbf{cleopian} (v. weak); pret. 1 sg. cleopode; past ptcp. cleopod. call; summon. Accusative person summoned; distinguish intransitive cry aloud. cry aloud. Intransitive, with the words cried given as direct speech.\par
% glossary-block:B00270:end
% glossary-block:B00271:begin
\needspace{3\baselineskip}
\textbf{clif} (n. n.). cliff.\par
% glossary-block:B00271:end
% glossary-block:B00272:begin
\needspace{3\baselineskip}
\textbf{cnapa} (n. m.); gen. sg. cnapan; nom. pl. cnapan. boy.\par
% glossary-block:B00272:end
% glossary-block:B00273:begin
\needspace{3\baselineskip}
\textbf{cnāwan} (v. strong); pret. 1 sg. cnēow; pret. pl. cnēowon; past ptcp. cnāwen; pres. 3 sg. cnǣwð. know.\par
% glossary-block:B00273:end
% glossary-block:B00274:begin
\needspace{3\baselineskip}
\textbf{cnǣwð} — pres. ind. 3 sg.: see \textbf{cnāwan}.\par
% glossary-block:B00274:end
% glossary-block:B00275:begin
\needspace{3\baselineskip}
\textbf{cnēo} (n. n.); gen. sg. cnēowes; nom. pl. cnēowu. knee. Supplied accusative plural knees. knee.\par
% glossary-block:B00275:end
% glossary-block:B00276:begin
\needspace{3\baselineskip}
\textbf{cnēoris} (n. f.); gen. sg. cnēorisse; nom. pl. cnēorissa. generation.\par
% glossary-block:B00276:end
% glossary-block:B00277:begin
\needspace{3\baselineskip}
\textbf{cnīf} (n.). knife. Pronunciation example.\par
% glossary-block:B00277:end
% glossary-block:B00278:begin
\needspace{3\baselineskip}
\textbf{cniht} (n. m.); gen. sg. cnihtes; nom. pl. cnihtas. young man; youth. servant.\par
% glossary-block:B00278:end
% glossary-block:B00279:begin
\needspace{3\baselineskip}
\textbf{cnucian} (v.). knock.\par
% glossary-block:B00279:end
% glossary-block:B00280:begin
\needspace{3\baselineskip}
\textbf{coss} (n. m.); gen. sg. cosses; nom. pl. cossas. kiss.\par
% glossary-block:B00280:end
% glossary-block:B00281:begin
\needspace{3\baselineskip}
\textbf{costnung} (n. f.); gen. sg. costnunge; nom. pl. costnunga. temptation.\par
% glossary-block:B00281:end
% glossary-block:B00282:begin
\needspace{3\baselineskip}
\textbf{cræft} (n. m.); gen. sg. cræftes; nom. pl. cræftas. trade; occupation; craft. skill; accomplishment.\par
% glossary-block:B00282:end
% glossary-block:B00283:begin
\needspace{3\baselineskip}
\textbf{Crīst} (n. m.); gen. sg. Crīstes. Christ.\par
% glossary-block:B00283:end
% glossary-block:B00284:begin
\needspace{3\baselineskip}
\textbf{cucu} (adj.). alive; living. Supplied accusative neuter plural adjective.\par
% glossary-block:B00284:end
% glossary-block:B00285:begin
\needspace{3\baselineskip}
\textbf{cuma} (n. m.); gen. sg. cuman; nom. pl. cuman. stranger; guest.\par
% glossary-block:B00285:end
% glossary-block:B00286:begin
\needspace{3\baselineskip}
\textbf{cuman} (v. strong); pret. 1 sg. cōm; pret. pl. cōmon; past ptcp. cumen; pres. 3 sg. cymð. come.\par
% glossary-block:B00286:end
% glossary-block:B00287:begin
\needspace{3\baselineskip}
\textbf{cunnan} (v. preterite-present); pres. 1/3 sg. cann; pres. pl. cunnon; pret. 1 sg. cūðe; past ptcp. (ge)cunnen. know how; be able. Learned skill; plain infinitive complement. know; be acquainted with. Accusative person known; distinct from learned ability.\par
% glossary-block:B00287:end
% glossary-block:B00288:begin
\needspace{3\baselineskip}
\textbf{cūðe} — pret. sg.: see \textbf{cunnan}.\par
% glossary-block:B00288:end
% glossary-block:B00289:begin
\needspace{3\baselineskip}
\textbf{cweartern} (n. n.); gen. sg. cwearternes; nom. pl. cweartern. prison.\par
% glossary-block:B00289:end
% glossary-block:B00290:begin
\needspace{3\baselineskip}
\textbf{cwellan} (v.). kill.\par
% glossary-block:B00290:end
% glossary-block:B00291:begin
\needspace{3\baselineskip}
\textbf{cwēn} (n. f.); gen. sg. cwēne; nom. pl. cwēna. queen. princess.\par
% glossary-block:B00291:end
% glossary-block:B00292:begin
\needspace{3\baselineskip}
\textbf{cweðan} (v. strong); pret. 1 sg. cwæð; pret. pl. cwǣdon; past ptcp. cweden; pres. 3 sg. cwið. say. say. Source genitive plural present participle, of people saying. say. Supplied preterite said.\par
% glossary-block:B00292:end
% glossary-block:B00293:begin
\needspace{3\baselineskip}
\textbf{cwic} (adj.). alive; living. living people. Substantival adjective.\par
% glossary-block:B00293:end
% glossary-block:B00294:begin
\needspace{3\baselineskip}
\textbf{cwið} — pres. ind. 3 sg.: see \textbf{cweðan}.\par
% glossary-block:B00294:end
% glossary-block:B00295:begin
\needspace{3\baselineskip}
\textbf{cwōm} — pret. sg., variant: see \textbf{cuman}.\par
% glossary-block:B00295:end
% glossary-block:B00296:begin
\needspace{3\baselineskip}
\textbf{cymð} — pres. ind. 3 sg.: see \textbf{cuman}.\par
% glossary-block:B00296:end
% glossary-block:B00297:begin
\needspace{3\baselineskip}
\textbf{cynehelm} (n. m.); gen. sg. cynehelmes; nom. pl. cynehelmas. crown.\par
% glossary-block:B00297:end
% glossary-block:B00298:begin
\needspace{3\baselineskip}
\textbf{cynesętl} (n. n.). royal seat; throne.\par
% glossary-block:B00298:end
% glossary-block:B00299:begin
\needspace{3\baselineskip}
\textbf{cyning} (n. m.); gen. sg. cyninges; nom. pl. cyningas. king.\par
% glossary-block:B00299:end
% glossary-block:B00300:begin
\needspace{3\baselineskip}
\textbf{cynn} (n. n.); gen. sg. cynnes; nom. pl. cynn. kind.\par
% glossary-block:B00300:end
% glossary-block:B00301:begin
\needspace{3\baselineskip}
\textbf{Cyrenisc} (adj.). Cyrenian; of Cyrene. Geographical adjective.\par
% glossary-block:B00301:end
% glossary-block:B00302:begin
\needspace{3\baselineskip}
\textbf{cyrice} — citation/spelling variant: see \textbf{cirice}.\par
% glossary-block:B00302:end
% glossary-block:B00303:begin
\needspace{3\baselineskip}
\textbf{cyssan} (v. weak); pret. 1 sg. cyste; past ptcp. (ge)cyssed. kiss.\par
% glossary-block:B00303:end
% glossary-block:B00304:begin
\needspace{3\baselineskip}
\textbf{cȳðan} (v. weak); pret. 1 sg. cȳðde; past ptcp. cȳðed. announce; make known. Dative person informed; accusative thing announced or a ðæt content clause.\par
% glossary-block:B00304:end
% book-navigation: standalone bookmark replaced by volume hierarchy
% glossary-block:B00305:begin
\needspace{5\baselineskip}
\subsection*{D}\BookHeading{nav-glossary-old-english-0006}{subsubsection}{D}{0}
% glossary-block:B00305:end
% glossary-block:B00306:begin
\textbf{dǣd} (n. f.). deed.\par
% glossary-block:B00306:end
% glossary-block:B00307:begin
\needspace{3\baselineskip}
\textbf{dæg} (n. m.); gen. sg. dæges; nom. pl. dagas. day.\par
% glossary-block:B00307:end
% glossary-block:B00308:begin
\needspace{3\baselineskip}
\textbf{dæghwāmlīc} (adj.). daily.\par
% glossary-block:B00308:end
% glossary-block:B00309:begin
\needspace{3\baselineskip}
\textbf{dǣl} (n. m.); gen. sg. dǣles; nom. pl. dǣlas. part.\par
% glossary-block:B00309:end
% glossary-block:B00310:begin
\needspace{3\baselineskip}
\textbf{dēad} (adj.). dead. dead people. Substantival adjective.\par
% glossary-block:B00310:end
% glossary-block:B00311:begin
\needspace{3\baselineskip}
\textbf{dearr} — pres. ind. 1/3 sg.: see \textbf{durran}.\par
% glossary-block:B00311:end
% glossary-block:B00312:begin
\needspace{3\baselineskip}
\textbf{dēað} (n. m.); gen. sg. dēaðes; nom. pl. dēaðas. death.\par
% glossary-block:B00312:end
% glossary-block:B00313:begin
\needspace{3\baselineskip}
\textbf{dehter} — dat. sg.: see \textbf{dohtor}.\par
% glossary-block:B00313:end
% glossary-block:B00314:begin
\needspace{3\baselineskip}
\textbf{dēma} (n. m.); gen. sg. dēman; nom. pl. dēman. judge.\par
% glossary-block:B00314:end
% glossary-block:B00315:begin
\needspace{3\baselineskip}
\textbf{dēman} (v. weak); pret. 1 sg. dēmde; past ptcp. dēmed. judge. condemn. Accusative person condemned.\par
% glossary-block:B00315:end
% glossary-block:B00316:begin
\needspace{3\baselineskip}
\textbf{Dęne} (n. m.); nom. pl. Dęne. Danes. Plural people-name.\par
% glossary-block:B00316:end
% glossary-block:B00317:begin
\needspace{3\baselineskip}
\textbf{dēofol} (n. m.); gen. sg. dēofles; nom. pl. dēoflas. devil; demon.\par
% glossary-block:B00317:end
% glossary-block:B00318:begin
\needspace{3\baselineskip}
\textbf{dēor} (n. n.); gen. sg. dēores; nom. pl. dēor. animal; beast.\par
% glossary-block:B00318:end
% glossary-block:B00319:begin
\needspace{3\baselineskip}
\textbf{dēorwierðe} (adj.); comp. dēorwierðra. precious; valuable.\par
% glossary-block:B00319:end
% glossary-block:B00320:begin
\needspace{3\baselineskip}
\textbf{dēorwierðra} (adj.). more precious. Comparative adjective; distinct from an adverb modifying an action. Comparative of dēorwierðe.\par
% glossary-block:B00320:end
% glossary-block:B00321:begin
\needspace{3\baselineskip}
\textbf{dēorwurð} (adj.). precious; costly.\par
% glossary-block:B00321:end
% glossary-block:B00322:begin
\needspace{3\baselineskip}
\textbf{dēð} — pres. ind. 3 sg.: see \textbf{dōn}.\par
% glossary-block:B00322:end
% glossary-block:B00323:begin
\needspace{3\baselineskip}
\textbf{dīc} (n. m.); gen. sg. dīces; nom. pl. dīcas. ditch; dike.\par
% glossary-block:B00323:end
% glossary-block:B00324:begin
\needspace{3\baselineskip}
\textbf{dīegelnes} (n. f.). retreat; seclusion. Supplied plural after of.\par
% glossary-block:B00324:end
% glossary-block:B00325:begin
\needspace{3\baselineskip}
\textbf{diernan} (v.). conceal; hide.\par
% glossary-block:B00325:end
% glossary-block:B00326:begin
\needspace{3\baselineskip}
\textbf{disc} (n. m.); gen. sg. disces; nom. pl. discas. dish.\par
% glossary-block:B00326:end
% glossary-block:B00327:begin
\needspace{3\baselineskip}
\textbf{docga} (n.). dog. Pronunciation example.\par
% glossary-block:B00327:end
% glossary-block:B00328:begin
\needspace{3\baselineskip}
\textbf{dohtor} (n. f.); gen. sg. dohtor; nom. pl. dohtor. daughter. Dative singular dehter, also a late genitive variant.\par
% glossary-block:B00328:end
% glossary-block:B00329:begin
\needspace{3\baselineskip}
\textbf{dōm} (n. m.); gen. sg. dōmes; nom. pl. dōmas. judgment. dōmes dæg means Day of Judgment. doom; fate. judgment. A judgment made in judging, distinct from doom/fate.\par
% glossary-block:B00329:end
% glossary-block:B00330:begin
\needspace{3\baselineskip}
\textbf{dōn} (v. irregular); pres. 1 sg. dō; pres. 3 sg. dēð; pres. pl. dōð; pret. 1 sg. dyde; pret. pl. dydon; past ptcp. (ge)dōn. do; perform. do for; do to. Dative beneficiary/affected person; thing done is accusative when expressed. put; place. Accusative object; on + accusative goal in the assigned frame. make; render. Accusative object with agreeing object predicate. do. Replaces an action already mentioned; distinct from English dummy do. cause; bring about. With a finite ðæt clause; its actor is nominative within that clause. fulfil; carry out. A word or command is fulfilled. deal with; behave toward. Assigned wið person frame; manner is separately expressed.\par
% glossary-block:B00330:end
% glossary-block:B00331:begin
\needspace{3\baselineskip}
\textbf{dorste} — pret. sg.: see \textbf{durran}.\par
% glossary-block:B00331:end
% glossary-block:B00332:begin
\needspace{3\baselineskip}
\textbf{draca} (n. m.); gen. sg. dracan; nom. pl. dracan. dragon.\par
% glossary-block:B00332:end
% glossary-block:B00333:begin
\needspace{3\baselineskip}
\textbf{drēam} (n. m.). joy; bliss.\par
% glossary-block:B00333:end
% glossary-block:B00334:begin
\needspace{3\baselineskip}
\textbf{dręnc} (n. m.); gen. sg. dręnces; nom. pl. dręncas. drink.\par
% glossary-block:B00334:end
% glossary-block:B00335:begin
\needspace{3\baselineskip}
\textbf{drēogan} (v. strong); pret. 1 sg. drēag; pret. pl. drugon; past ptcp. drogen; pres. 3 sg. drīegð. endure.\par
% glossary-block:B00335:end
% glossary-block:B00336:begin
\needspace{3\baselineskip}
\textbf{drīegð} — pres. ind. 3 sg.: see \textbf{drēogan}.\par
% glossary-block:B00336:end
% glossary-block:B00337:begin
\needspace{3\baselineskip}
\textbf{drīehð} — pres. ind. 3 sg., variant: see \textbf{drēogan}.\par
% glossary-block:B00337:end
% glossary-block:B00338:begin
\needspace{3\baselineskip}
\textbf{drīfan} (v. strong); pret. 1 sg. drāf; pret. pl. drifon; past ptcp. drifen. drive. Accusative person or thing driven.\par
% glossary-block:B00338:end
% glossary-block:B00339:begin
\needspace{3\baselineskip}
\textbf{drincan} (v. strong); pret. 1 sg. dranc; pret. pl. druncon; past ptcp. druncen. drink. Accusative liquid drunk.\par
% glossary-block:B00339:end
% glossary-block:B00340:begin
\needspace{3\baselineskip}
\textbf{drȳgnes} (n. f.); gen. sg. drȳgnesse; nom. pl. drȳgnessa. dry land.\par
% glossary-block:B00340:end
% glossary-block:B00341:begin
\needspace{3\baselineskip}
\textbf{dryhten} (n. m.); gen. sg. dryhtnes; nom. pl. dryhtnas. lord; the Lord.\par
% glossary-block:B00341:end
% glossary-block:B00342:begin
\needspace{3\baselineskip}
\textbf{dūn} (n. f.). mountain; hill.\par
% glossary-block:B00342:end
% glossary-block:B00343:begin
\needspace{3\baselineskip}
\textbf{durran} (v. preterite-present); pres. 1/3 sg. dearr; pres. pl. durron; pret. 1 sg. dorste. dare to. Conventional citation label; no usable infinitive is supplied. See the finite forms in Unit 6. No past-participle citation is given in the checked source paradigm.\par
% glossary-block:B00343:end
% glossary-block:B00344:begin
\needspace{3\baselineskip}
\textbf{dūst} (n. n.); gen. sg. dūstes; nom. pl. dūst. dust.\par
% glossary-block:B00344:end
% glossary-block:B00345:begin
\needspace{3\baselineskip}
\textbf{dyde} — pret. sg.: see \textbf{dōn}.\par
% glossary-block:B00345:end
% book-navigation: standalone bookmark replaced by volume hierarchy
% glossary-block:B00346:begin
\needspace{5\baselineskip}
\subsection*{E}\BookHeading{nav-glossary-old-english-0007}{subsubsection}{E}{0}
% glossary-block:B00346:end
% glossary-block:B00347:begin
\textbf{ēac} (adv.). also. Also in the combination ēac swilce.\par
% glossary-block:B00347:end
% glossary-block:B00348:begin
\needspace{3\baselineskip}
\textbf{ēadig} (adj.). happy; blessed.\par
% glossary-block:B00348:end
% glossary-block:B00349:begin
\needspace{3\baselineskip}
\textbf{ēage} (n. n.); gen. sg. ēagan; nom. pl. ēagan. eye.\par
% glossary-block:B00349:end
% glossary-block:B00350:begin
\needspace{3\baselineskip}
\textbf{ēagðȳrel} (n. n.); gen. sg. ēagðȳrles; nom. pl. ēagðȳrlu. window.\par
% glossary-block:B00350:end
% glossary-block:B00351:begin
\needspace{3\baselineskip}
\textbf{eahta} (cardinal numeral). eight. Cardinal 8.\par
% glossary-block:B00351:end
% glossary-block:B00352:begin
\needspace{3\baselineskip}
\textbf{ēalā} (interj.). O!. Interjection in address, exclamation or wish.\par
% glossary-block:B00352:end
% glossary-block:B00353:begin
\needspace{3\baselineskip}
\textbf{eald} (adj.); comp. ieldra; superl. ieldest. old. old people; the elderly. Substantival adjective.\par
% glossary-block:B00353:end
% glossary-block:B00354:begin
\needspace{3\baselineskip}
\textbf{ealdor} (n. m.); gen. sg. ealdres; nom. pl. ealdras. chief; king. This noun is distinct from the homonymous noun meaning life.\par
% glossary-block:B00354:end
% glossary-block:B00355:begin
\needspace{3\baselineskip}
\textbf{ealdordōm} (n. m.). chief place; primacy.\par
% glossary-block:B00355:end
% glossary-block:B00356:begin
\needspace{3\baselineskip}
\textbf{ealdormann} (n. m.); gen. sg. ealdormannes; nom. pl. ealdormęnn. prince; leader; noble.\par
% glossary-block:B00356:end
% glossary-block:B00357:begin
\needspace{3\baselineskip}
\textbf{ealdorscipe} (n. m.); gen. sg. ealdorscipes; nom. pl. ealdorscipas. chief place; primacy.\par
% glossary-block:B00357:end
% glossary-block:B00358:begin
\needspace{3\baselineskip}
\textbf{eall} (adj.; substantive; adv.). all; whole. all people; everything. Substantival: plural for people, neuter singular for everything. completely. Only in the fixed phrase mid ealle. just; exactly. Degree emphasis in eall swā.\par
% glossary-block:B00358:end
% glossary-block:B00359:begin
\needspace{3\baselineskip}
\textbf{ēare} (n. n.); gen. sg. ēaran; nom. pl. ēaran. ear.\par
% glossary-block:B00359:end
% glossary-block:B00360:begin
\needspace{3\baselineskip}
\textbf{earfoðrīme} (adj.). innumerable; hard to count. Assigned source adjective.\par
% glossary-block:B00360:end
% glossary-block:B00361:begin
\needspace{3\baselineskip}
\textbf{earg} (adj.). cowardly.\par
% glossary-block:B00361:end
% glossary-block:B00362:begin
\needspace{3\baselineskip}
\textbf{earm} (n.). arm. Pronunciation example.\par
% glossary-block:B00362:end
% glossary-block:B00363:begin
\needspace{3\baselineskip}
\textbf{earm} (adj.). the poor; the unfortunate. Substantival adjective.\par
% glossary-block:B00363:end
% glossary-block:B00364:begin
\needspace{3\baselineskip}
\textbf{earmlice} (adverb or particle). cruelly; miserably.\par
% glossary-block:B00364:end
% glossary-block:B00365:begin
\needspace{3\baselineskip}
\textbf{earnung} (n. f.). merit; good service. Source plural merits/services.\par
% glossary-block:B00365:end
% glossary-block:B00366:begin
\needspace{3\baselineskip}
\textbf{eart} — pres. ind. 2 sg.: see \textbf{wesan}.\par
% glossary-block:B00366:end
% glossary-block:B00367:begin
\needspace{3\baselineskip}
\textbf{Eastęngle} (proper name). East Anglians.\par
% glossary-block:B00367:end
% glossary-block:B00368:begin
\needspace{3\baselineskip}
\textbf{ēastnorðerne} (adj.). northeasterly.\par
% glossary-block:B00368:end
% glossary-block:B00369:begin
\needspace{3\baselineskip}
\textbf{ēastsūðdǣl} (n. m.). southeastern part.\par
% glossary-block:B00369:end
% glossary-block:B00370:begin
\needspace{3\baselineskip}
\textbf{ēaðmōd} (adj.). lowly; humble. humble people; the lowly. Substantival adjective.\par
% glossary-block:B00370:end
% glossary-block:B00371:begin
\needspace{3\baselineskip}
\textbf{ēaðmōdnes} (n. f.); gen. sg. ēaðmōdnesse. humility.\par
% glossary-block:B00371:end
% glossary-block:B00372:begin
\needspace{3\baselineskip}
\textbf{ēce} (adj.). eternal.\par
% glossary-block:B00372:end
% glossary-block:B00373:begin
\needspace{3\baselineskip}
\textbf{ecg} (n.). edge. Pronunciation example.\par
% glossary-block:B00373:end
% glossary-block:B00374:begin
\needspace{3\baselineskip}
\textbf{edlēan} (n. n.); gen. sg. edlēanes; nom. pl. edlēan. reward; recompense.\par
% glossary-block:B00374:end
% glossary-block:B00375:begin
\needspace{3\baselineskip}
\textbf{ęfnan} (v. weak); pret. 1 sg. ęfnede; past ptcp. (ge)ęfned. perform.\par
% glossary-block:B00375:end
% glossary-block:B00376:begin
\needspace{3\baselineskip}
\textbf{efne} (interj.). behold!. Interjection, distinct from the verb behealdan look at. just; precisely; even.\par
% glossary-block:B00376:end
% glossary-block:B00377:begin
\needspace{3\baselineskip}
\textbf{eft, ęft} (adv.). again; also. The source note also permits the additive sense. again; back. again; afterwards.\par
% glossary-block:B00377:end
% glossary-block:B00378:begin
\needspace{3\baselineskip}
\textbf{ęfthweorfan} (v. strong); pret. 1 sg. ęfthwearf; pret. pl. ęfthwurfon; past ptcp. ęfthworfen; pres. 3 sg. ęfthwierfð. return; turn back.\par
% glossary-block:B00378:end
% glossary-block:B00379:begin
\needspace{3\baselineskip}
\textbf{ęfthwierfð} — pres. ind. 3 sg.: see \textbf{ęfthweorfan}.\par
% glossary-block:B00379:end
% glossary-block:B00380:begin
\needspace{3\baselineskip}
\textbf{ęgesa} (n. m.). fear; terror. Reference in a source or explanatory note.\par
% glossary-block:B00380:end
% glossary-block:B00381:begin
\needspace{3\baselineskip}
\textbf{ęgeslic} (adj.). dreadful; terrible.\par
% glossary-block:B00381:end
% glossary-block:B00382:begin
\needspace{3\baselineskip}
\textbf{Egypta} (proper name). Egyptians. Egypta land denotes Egypt, literally land of the Egyptians; supplied name forms. Supplied name forms: Egypta / Ægypta, genitive plural; Egyptum, dative plural. Egypta land is the land of the Egyptians.\par
% glossary-block:B00382:end
% glossary-block:B00383:begin
\needspace{3\baselineskip}
\textbf{Egypte} (ethnonym). Egyptians.\par
% glossary-block:B00383:end
% glossary-block:B00384:begin
\needspace{3\baselineskip}
\textbf{ēhtan} (v. weak); pret. 1 sg. ēhte; past ptcp. (ge)ēht. pursue; persecute. Accusative person in the selected production frame; a separately supported genitive alternative is not the selected target.\par
% glossary-block:B00384:end
% glossary-block:B00385:begin
\needspace{3\baselineskip}
\textbf{ęlcung} (n. f.); gen. sg. ęlcunge. delay.\par
% glossary-block:B00385:end
% glossary-block:B00386:begin
\needspace{3\baselineskip}
\textbf{ęle} (n. m.); gen. sg. ęles; nom. pl. ęlas. oil.\par
% glossary-block:B00386:end
% glossary-block:B00387:begin
\needspace{3\baselineskip}
\textbf{elp} (n. m.); gen. sg. elpes; nom. pl. elpas. elephant.\par
% glossary-block:B00387:end
% glossary-block:B00388:begin
\needspace{3\baselineskip}
\textbf{elpenbǣnen} (adj.). ivory; made of ivory. Adjective of material.\par
% glossary-block:B00388:end
% glossary-block:B00389:begin
\needspace{3\baselineskip}
\textbf{ęlðēodig} (adj.). foreigners; strangers. Substantival adjective. foreign.\par
% glossary-block:B00389:end
% glossary-block:B00390:begin
\needspace{3\baselineskip}
\textbf{ęnde} (n. m.). end.\par
% glossary-block:B00390:end
% glossary-block:B00391:begin
\needspace{3\baselineskip}
\textbf{endlefan} (cardinal numeral). eleven. Cardinal 11.\par
% glossary-block:B00391:end
% glossary-block:B00392:begin
\needspace{3\baselineskip}
\textbf{endlefta} (ordinal). eleventh.\par
% glossary-block:B00392:end
% glossary-block:B00393:begin
\needspace{3\baselineskip}
\textbf{endleofan} (numeral). eleven.\par
% glossary-block:B00393:end
% glossary-block:B00394:begin
\needspace{3\baselineskip}
\textbf{ęngel} (n. m.); gen. sg. ęngles; nom. pl. ęnglas. angel.\par
% glossary-block:B00394:end
% glossary-block:B00395:begin
\needspace{3\baselineskip}
\textbf{Engle} (n. m.); nom. pl. Engle. the English. People, distinct from the language/adjective.\par
% glossary-block:B00395:end
% glossary-block:B00396:begin
\needspace{3\baselineskip}
\textbf{Englisc} (adj.). English.\par
% glossary-block:B00396:end
% glossary-block:B00397:begin
\needspace{3\baselineskip}
\textbf{Englisc} (adjective used substantivally). English, the language. Language noun.\par
% glossary-block:B00397:end
% glossary-block:B00398:begin
\needspace{3\baselineskip}
\textbf{ēode} — pret. sg.: see \textbf{gān}.\par
% glossary-block:B00398:end
% glossary-block:B00399:begin
\needspace{3\baselineskip}
\textbf{eom} — pres. ind. 1 sg.: see \textbf{wesan}.\par
% glossary-block:B00399:end
% glossary-block:B00400:begin
\needspace{3\baselineskip}
\textbf{eorl} (n.). man. Pronunciation example.\par
% glossary-block:B00400:end
% glossary-block:B00401:begin
\needspace{3\baselineskip}
\textbf{eornostlice} (adv.). accordingly; thus.\par
% glossary-block:B00401:end
% glossary-block:B00402:begin
\needspace{3\baselineskip}
\textbf{eorðe} (n. f.); gen. sg. eorðan; nom. pl. eorðan. earth.\par
% glossary-block:B00402:end
% glossary-block:B00403:begin
\needspace{3\baselineskip}
\textbf{eorðlic} (adj.). earthly.\par
% glossary-block:B00403:end
% glossary-block:B00404:begin
\needspace{3\baselineskip}
\textbf{eorðtilð} (n. f.); gen. sg. eorðtilðe; nom. pl. eorðtilða. agriculture.\par
% glossary-block:B00404:end
% glossary-block:B00405:begin
\needspace{3\baselineskip}
\textbf{ēow} — dat. pl.; acc. pl.: see \textbf{ðū}.\par
% glossary-block:B00405:end
% glossary-block:B00406:begin
\needspace{3\baselineskip}
\textbf{ēower} (possessive adjective). your; yours. Plural possessors; inflection agrees with the possessed noun. Also gen. pl.: see \textbf{ðū}.\par
% glossary-block:B00406:end
% glossary-block:B00407:begin
\needspace{3\baselineskip}
\textbf{ēowic} — acc. pl., variant: see \textbf{ðū}.\par
% glossary-block:B00407:end
% glossary-block:B00408:begin
\needspace{3\baselineskip}
\textbf{ęrian} (v.). plow.\par
% glossary-block:B00408:end
% glossary-block:B00409:begin
\needspace{3\baselineskip}
\textbf{ēst} (n. m.). provisions; food.\par
% glossary-block:B00409:end
% glossary-block:B00410:begin
\needspace{3\baselineskip}
\textbf{etan} (v. strong); pret. 1 sg. ǣt; pret. pl. ǣton; past ptcp. eten; pres. 3 sg. itt. eat.\par
% glossary-block:B00410:end
% glossary-block:B00411:begin
\needspace{3\baselineskip}
\textbf{ēðel} (n. m.); gen. sg. ēðles; nom. pl. ēðlas. native land; homeland.\par
% glossary-block:B00411:end
% book-navigation: standalone bookmark replaced by volume hierarchy
% glossary-block:B00412:begin
\needspace{5\baselineskip}
\subsection*{F}\BookHeading{nav-glossary-old-english-0008}{subsubsection}{F}{0}
% glossary-block:B00412:end
% glossary-block:B00413:begin
\textbf{faran} (v. strong); pret. 1 sg. fōr; pret. pl. fōron; past ptcp. faren; pres. 3 sg. færð. go; travel.\par
% glossary-block:B00413:end
% glossary-block:B00414:begin
\needspace{3\baselineskip}
\textbf{Farne} (proper name). Farne. Island; supplied proper-name form after in.\par
% glossary-block:B00414:end
% glossary-block:B00415:begin
\needspace{3\baselineskip}
\textbf{fæc} (n. n.). interval; period of time; while. Noun in the supplied time expression.\par
% glossary-block:B00415:end
% glossary-block:B00416:begin
\needspace{3\baselineskip}
\textbf{fæder} (n. m.); gen. sg. fæder; nom. pl. fæderas. father. Includes the spiritual father relationship in the assigned prompts.\par
% glossary-block:B00416:end
% glossary-block:B00417:begin
\needspace{3\baselineskip}
\textbf{fæger} (adj.); comp. fægerra; supplied superl. acc. f. sg. fægerestan. beautiful; fair.\par
% glossary-block:B00417:end
% glossary-block:B00418:begin
\needspace{3\baselineskip}
\textbf{fægerest} (adj.). most beautiful. Superlative of fæger; source form fægerestan.\par
% glossary-block:B00418:end
% glossary-block:B00419:begin
\needspace{3\baselineskip}
\textbf{fægerra} (adj.). more beautiful; fairer. Comparative of fæger.\par
% glossary-block:B00419:end
% glossary-block:B00420:begin
\needspace{3\baselineskip}
\textbf{fǣmne} (n. f.); gen. sg. fǣmnan; nom. pl. fǣmnan. maiden; young woman.\par
% glossary-block:B00420:end
% glossary-block:B00421:begin
\needspace{3\baselineskip}
\textbf{fǣringa} (adv.). suddenly.\par
% glossary-block:B00421:end
% glossary-block:B00422:begin
\needspace{3\baselineskip}
\textbf{færð} — pres. ind. 3 sg.: see \textbf{faran}.\par
% glossary-block:B00422:end
% glossary-block:B00423:begin
\needspace{3\baselineskip}
\textbf{fæsten} (n. n.). fortification; stronghold.\par
% glossary-block:B00423:end
% glossary-block:B00424:begin
\needspace{3\baselineskip}
\textbf{fæsthafol} (adj.). tenacious; withholding. Genitive of what is held back in the source.\par
% glossary-block:B00424:end
% glossary-block:B00425:begin
\needspace{3\baselineskip}
\textbf{fæstnes} (n. f.); gen. sg. fæstnesse; nom. pl. fæstnessa. firmament.\par
% glossary-block:B00425:end
% glossary-block:B00426:begin
\needspace{3\baselineskip}
\textbf{fæt} (n. n.); gen. sg. fætes; nom. pl. fatu. vessel; container.\par
% glossary-block:B00426:end
% glossary-block:B00427:begin
\needspace{3\baselineskip}
\textbf{feallan} (v. strong); pret. 1 sg. fēoll; pret. pl. fēollon; past ptcp. feallen; pres. 3 sg. fielð. fall.\par
% glossary-block:B00427:end
% glossary-block:B00428:begin
\needspace{3\baselineskip}
\textbf{fēawe} (adj.). few people. Substantival use in the assigned expression. few.\par
% glossary-block:B00428:end
% glossary-block:B00429:begin
\needspace{3\baselineskip}
\textbf{feccean} (v.). fetch.\par
% glossary-block:B00429:end
% glossary-block:B00430:begin
\needspace{3\baselineskip}
\textbf{fēdan} (v. weak); pret. 1 sg. fēdde; past ptcp. fēded. feed; nourish. Accusative being fed; food may be mid + dative or a bare instrumental in the assigned expressions.\par
% glossary-block:B00430:end
% glossary-block:B00431:begin
\needspace{3\baselineskip}
\textbf{fēhð} — pres. ind. 3 sg.: see \textbf{fōn}.\par
% glossary-block:B00431:end
% glossary-block:B00432:begin
\needspace{3\baselineskip}
\textbf{fela} (adj.). many; a large number of. Genitive of the things counted in the taught frame.\par
% glossary-block:B00432:end
% glossary-block:B00433:begin
\needspace{3\baselineskip}
\textbf{felgan} — citation/spelling variant: see \textbf{fēolan}.\par
% glossary-block:B00433:end
% glossary-block:B00434:begin
\needspace{3\baselineskip}
\textbf{feoh} (n. n.); gen. sg. fēos. money.\par
% glossary-block:B00434:end
% glossary-block:B00435:begin
\needspace{3\baselineskip}
\textbf{feohtan} (v. strong); pret. 1 sg. feaht; pret. pl. fuhton; past ptcp. fohten; pres. 3 sg. fieht. fight; contend. Intransitive; wið + accusative opponent.\par
% glossary-block:B00435:end
% glossary-block:B00436:begin
\needspace{3\baselineskip}
\textbf{fēolan} (v. strong); pret. 1 sg. fealh; pret. pl. fulgon; past ptcp. folen; pres. 3 sg. filhð. penetrate. Preserve the supplied/taught lexical frame; not a synonym for every kind of entering.\par
% glossary-block:B00436:end
% glossary-block:B00437:begin
\needspace{3\baselineskip}
\textbf{fēond} (n. m.); gen. sg. fēondes; nom. pl. fīend. enemy; foe.\par
% glossary-block:B00437:end
% glossary-block:B00438:begin
\needspace{3\baselineskip}
\textbf{feorr} (adv.). far; far off.\par
% glossary-block:B00438:end
% glossary-block:B00439:begin
\needspace{3\baselineskip}
\textbf{feorran} (adv.). from afar.\par
% glossary-block:B00439:end
% glossary-block:B00440:begin
\needspace{3\baselineskip}
\textbf{fēorða} (ordinal numeral). fourth.\par
% glossary-block:B00440:end
% glossary-block:B00441:begin
\needspace{3\baselineskip}
\textbf{fēower} (cardinal numeral). four. Cardinal 4.\par
% glossary-block:B00441:end
% glossary-block:B00442:begin
\needspace{3\baselineskip}
\textbf{fēowertig} (numeral). forty.\par
% glossary-block:B00442:end
% glossary-block:B00443:begin
\needspace{3\baselineskip}
\textbf{fēran} (v.). go; journey. Reference in a source or explanatory note.\par
% glossary-block:B00443:end
% glossary-block:B00444:begin
\needspace{3\baselineskip}
\textbf{fēt} — pres. ind. 3 sg.: see \textbf{fēdan}; dat. sg.; nom./acc. pl.: see \textbf{fōt}.\par
% glossary-block:B00444:end
% glossary-block:B00445:begin
\needspace{3\baselineskip}
\textbf{fieht} — pres. ind. 3 sg.: see \textbf{feohtan}.\par
% glossary-block:B00445:end
% glossary-block:B00446:begin
\needspace{3\baselineskip}
\textbf{fielð} — pres. ind. 3 sg.: see \textbf{feallan}.\par
% glossary-block:B00446:end
% glossary-block:B00447:begin
\needspace{3\baselineskip}
\textbf{fielþ} — pres. ind. 3 sg., variant: see \textbf{feallan}.\par
% glossary-block:B00447:end
% glossary-block:B00448:begin
\needspace{3\baselineskip}
\textbf{fīend} — nom./acc. pl., variant: see \textbf{fēond}.\par
% glossary-block:B00448:end
% glossary-block:B00449:begin
\needspace{3\baselineskip}
\textbf{fierst} (n. m.). period; interval. Masculine noun in the source time expression.\par
% glossary-block:B00449:end
% glossary-block:B00450:begin
\needspace{3\baselineskip}
\textbf{fīf} (cardinal numeral). five. Cardinal 5.\par
% glossary-block:B00450:end
% glossary-block:B00451:begin
\needspace{3\baselineskip}
\textbf{fīfta} (ordinal numeral). fifth.\par
% glossary-block:B00451:end
% glossary-block:B00452:begin
\needspace{3\baselineskip}
\textbf{fīftig} (numeral). fifty.\par
% glossary-block:B00452:end
% glossary-block:B00453:begin
\needspace{3\baselineskip}
\textbf{filhð} — pres. ind. 3 sg.: see \textbf{fēolan}.\par
% glossary-block:B00453:end
% glossary-block:B00454:begin
\needspace{3\baselineskip}
\textbf{filhþ} — pres. ind. 3 sg., variant: see \textbf{fēolan}.\par
% glossary-block:B00454:end
% glossary-block:B00455:begin
\needspace{3\baselineskip}
\textbf{findan} (v. strong); pret. 1 sg. fand; pret. pl. fundon; past ptcp. funden. find. Accusative person or thing found.\par
% glossary-block:B00455:end
% glossary-block:B00456:begin
\needspace{3\baselineskip}
\textbf{finger} (n. m.); gen. sg. fingres; nom. pl. fingras. finger.\par
% glossary-block:B00456:end
% glossary-block:B00457:begin
\needspace{3\baselineskip}
\textbf{fisc} (n. m.); gen. sg. fisces; nom. pl. fiscas. fish.\par
% glossary-block:B00457:end
% glossary-block:B00458:begin
\needspace{3\baselineskip}
\textbf{fisccynn} (n. n.). kinds of fish. Supplied accusative neuter plural.\par
% glossary-block:B00458:end
% glossary-block:B00459:begin
\needspace{3\baselineskip}
\textbf{fiscere} (n. m.); gen. sg. fisceres; nom. pl. fisceras. fisherman.\par
% glossary-block:B00459:end
% glossary-block:B00460:begin
\needspace{3\baselineskip}
\textbf{fiscnoð} (n. m.). fishing.\par
% glossary-block:B00460:end
% glossary-block:B00461:begin
\needspace{3\baselineskip}
\textbf{flǣsc} (n. n.); gen. sg. flǣsces; nom. pl. flǣsc. flesh.\par
% glossary-block:B00461:end
% glossary-block:B00462:begin
\needspace{3\baselineskip}
\textbf{flēogan} (v. strong); pret. 1 sg. flēag; pret. pl. flugon; past ptcp. flogen; pres. 3 sg. flīehð. fly.\par
% glossary-block:B00462:end
% glossary-block:B00463:begin
\needspace{3\baselineskip}
\textbf{flēon} (v. strong); pret. 1 sg. flēah; pret. pl. flugon; past ptcp. flogen; pres. 3 sg. flīehð. flee. Supplied third-plural present; distinct from flēogan fly. flee; flee from; avoid. Intransitive or with accusative person/thing avoided.\par
% glossary-block:B00463:end
% glossary-block:B00464:begin
\needspace{3\baselineskip}
\textbf{flēwð} — pres. ind. 3 sg.: see \textbf{flōwan}.\par
% glossary-block:B00464:end
% glossary-block:B00465:begin
\needspace{3\baselineskip}
\textbf{flīegð} — pres. ind. 3 sg., variant: see \textbf{flēogan}.\par
% glossary-block:B00465:end
% glossary-block:B00466:begin
\needspace{3\baselineskip}
\textbf{flīehð} — pres. ind. 3 sg.: see \textbf{flēogan}; pres. ind. 3 sg.: see \textbf{flēon}.\par
% glossary-block:B00466:end
% glossary-block:B00467:begin
\needspace{3\baselineskip}
\textbf{flocc} (n. m.). group; company.\par
% glossary-block:B00467:end
% glossary-block:B00468:begin
\needspace{3\baselineskip}
\textbf{flōd} (n. m.); gen. sg. flōdes; nom. pl. flōdas. flood; stream of water.\par
% glossary-block:B00468:end
% glossary-block:B00469:begin
\needspace{3\baselineskip}
\textbf{flōwan} (v. strong); pret. 1 sg. flēow; pret. pl. flēowon; past ptcp. flōwen; pres. 3 sg. flēwð. flow.\par
% glossary-block:B00469:end
% glossary-block:B00470:begin
\needspace{3\baselineskip}
\textbf{fōdor} (n. n.); gen. sg. fōdres; nom. pl. fōdru. fodder.\par
% glossary-block:B00470:end
% glossary-block:B00471:begin
\needspace{3\baselineskip}
\textbf{folc} (n. n.); gen. sg. folces; nom. pl. folc. people. Collective singular. peoples. Distinct peoples or national groups; plural sense.\par
% glossary-block:B00471:end
% glossary-block:B00472:begin
\needspace{3\baselineskip}
\textbf{fōn} (v. strong); pret. 1 sg. fēng; pret. pl. fēngon; past ptcp. fangen; pres. 3 sg. fēhð. seize; catch. Accusative patient in the selected frame.\par
% glossary-block:B00472:end
% glossary-block:B00473:begin
\needspace{3\baselineskip}
\textbf{for} (prep.). for; for the sake of. Case follows each complete supplied accusative/dative expression. for; on behalf of. Dative beneficiary. why; for what reason. In for hwȳ or for hwon. because of. Dative hindering circumstance. by; for the sake of. Grounds of entreaty in the prayer; keep the supplied phrase. in exchange for. With dative. instead of; in place of. Accusative in the supplied substitution frame. because of; on account of. for; to the benefit of.\par
% glossary-block:B00473:end
% glossary-block:B00474:begin
\needspace{3\baselineskip}
\textbf{forbærnan} (v. weak); pret. 1 sg. forbærnde; past ptcp. forbærned. parch; scorch.\par
% glossary-block:B00474:end
% glossary-block:B00475:begin
\needspace{3\baselineskip}
\textbf{forceorfan} (v. strong); pret. 1 sg. forcearf; pret. pl. forcurfon; past ptcp. forcorfen; pres. 3 sg. forcierfð. cut down; cut off. Accusative tree or thing cut.\par
% glossary-block:B00475:end
% glossary-block:B00476:begin
\needspace{3\baselineskip}
\textbf{forcierfð} — pres. ind. 3 sg.: see \textbf{forceorfan}.\par
% glossary-block:B00476:end
% glossary-block:B00477:begin
\needspace{3\baselineskip}
\textbf{fordēð} — pres. ind. 3 sg.: see \textbf{fordōn}.\par
% glossary-block:B00477:end
% glossary-block:B00478:begin
\needspace{3\baselineskip}
\textbf{fordilgian} (v. weak); pret. 1 sg. fordilgode; past ptcp. fordilgod. destroy; annihilate.\par
% glossary-block:B00478:end
% glossary-block:B00479:begin
\needspace{3\baselineskip}
\textbf{fordōn} (v. irregular); pres. 1 sg. fordō; pres. 3 sg. fordēð; pres. pl. fordōð; pret. 1 sg. fordyde; pret. pl. fordydon; past ptcp. fordōn. destroy.\par
% glossary-block:B00479:end
% glossary-block:B00480:begin
\needspace{3\baselineskip}
\textbf{fordyde} — pret. sg.: see \textbf{fordōn}.\par
% glossary-block:B00480:end
% glossary-block:B00481:begin
\needspace{3\baselineskip}
\textbf{foresægde} — pret. sg.: see \textbf{foresęcgean}.\par
% glossary-block:B00481:end
% glossary-block:B00482:begin
\needspace{3\baselineskip}
\textbf{foresægð} — pres. ind. 3 sg.: see \textbf{foresęcgean}.\par
% glossary-block:B00482:end
% glossary-block:B00483:begin
\needspace{3\baselineskip}
\textbf{foresęcgean} (v. weak); pret. 1 sg. foresægde; past ptcp. foresægd; pres. 3 sg. foresægð. mention before; tell of beforehand.\par
% glossary-block:B00483:end
% glossary-block:B00484:begin
\needspace{3\baselineskip}
\textbf{forgiefan} (v. strong); pret. 1 sg. forgeaf; pret. pl. forgēafon; past ptcp. forgiefen. give; grant. Supplied granted sense at Cook 143.18; not forgive here. give; grant. Accusative thing given and dative recipient; the thing may be understood from context. forgive. Accusative sin, dative person forgiven; in a passive the sin becomes nominative and the person remains dative.\par
% glossary-block:B00484:end
% glossary-block:B00485:begin
\needspace{3\baselineskip}
\textbf{forgieldan} (v. strong); pret. 1 sg. forgeald; pret. pl. forguldon; past ptcp. forgolden. pay; repay; requite. Dative recipient, accusative thing repaid in the selected frame.\par
% glossary-block:B00485:end
% glossary-block:B00486:begin
\needspace{3\baselineskip}
\textbf{forierman} (v.). ruin; reduce to misery.\par
% glossary-block:B00486:end
% glossary-block:B00487:begin
\needspace{3\baselineskip}
\textbf{forlǣtan} (v. strong); pret. 1 sg. forlēt; pret. pl. forlēton; past ptcp. forlǣten. abandon; forsake. Accusative person or thing left. leave to. Thing left and recipient retain their taught roles. let go; dismiss. set aside. Supplied source sense.\par
% glossary-block:B00487:end
% glossary-block:B00488:begin
\needspace{3\baselineskip}
\textbf{forlēosan} (v. strong); pret. 1 sg. forlēas; pret. pl. forluron; past ptcp. forloren; pres. 3 sg. forlīest. lose. Accusative thing lost.\par
% glossary-block:B00488:end
% glossary-block:B00489:begin
\needspace{3\baselineskip}
\textbf{forliden} (adj.). shipwrecked.\par
% glossary-block:B00489:end
% glossary-block:B00490:begin
\needspace{3\baselineskip}
\textbf{forlidennes} (n. f.); gen. sg. forlidennesse; nom. pl. forlidennessa. shipwreck.\par
% glossary-block:B00490:end
% glossary-block:B00491:begin
\needspace{3\baselineskip}
\textbf{forlīest} — pres. ind. 3 sg.: see \textbf{forlēosan}.\par
% glossary-block:B00491:end
% glossary-block:B00492:begin
\needspace{3\baselineskip}
\textbf{forniman} (v.). destroy; lay waste.\par
% glossary-block:B00492:end
% glossary-block:B00493:begin
\needspace{3\baselineskip}
\textbf{forstelan} (v. strong); pret. 1 sg. forstæl; pret. pl. forstǣlon; past ptcp. forstolen; pres. 3 sg. forstilð. steal away; steal.\par
% glossary-block:B00493:end
% glossary-block:B00494:begin
\needspace{3\baselineskip}
\textbf{forstilð} — pres. ind. 3 sg.: see \textbf{forstelan}.\par
% glossary-block:B00494:end
% glossary-block:B00495:begin
\needspace{3\baselineskip}
\textbf{forswīgian} (v.). keep silent about; conceal.\par
% glossary-block:B00495:end
% glossary-block:B00496:begin
\needspace{3\baselineskip}
\textbf{fortrēdan} (v.). trample down.\par
% glossary-block:B00496:end
% glossary-block:B00497:begin
\needspace{3\baselineskip}
\textbf{forð} (adv.). forth. Adverbial particle with the bringing/producing verb.\par
% glossary-block:B00497:end
% glossary-block:B00498:begin
\needspace{3\baselineskip}
\textbf{for-ðām} (conj.). because.\par
% glossary-block:B00498:end
% glossary-block:B00499:begin
\needspace{3\baselineskip}
\textbf{forðām} (adv.). therefore. Pronunciation example.\par
% glossary-block:B00499:end
% glossary-block:B00500:begin
\needspace{3\baselineskip}
\textbf{for-ðām-ðe} (conj.). because.\par
% glossary-block:B00500:end
% glossary-block:B00501:begin
\needspace{3\baselineskip}
\textbf{forðfōr} (n. f.); gen. sg. forðfōre; nom. pl. forðfōra. departure.\par
% glossary-block:B00501:end
% glossary-block:B00502:begin
\needspace{3\baselineskip}
\textbf{for-ðon-ðe} (conj.; adv. in for-ðon). because. Causal conjunction. therefore; for that reason. Adverbial consequence, distinct from because.\par
% glossary-block:B00502:end
% glossary-block:B00503:begin
\needspace{3\baselineskip}
\textbf{forðtēon} (v.). present; perform. Source performance sense; retain the actual form's spelling.\par
% glossary-block:B00503:end
% glossary-block:B00504:begin
\needspace{3\baselineskip}
\textbf{forweorðan} (v. strong); pret. 1 sg. forwearð; pret. pl. forwurdon; past ptcp. forworden; pres. 3 sg. forwierð. perish.\par
% glossary-block:B00504:end
% glossary-block:B00505:begin
\needspace{3\baselineskip}
\textbf{forwiernan} (v. weak); pret. 1 sg. forwiernde; past ptcp. forwierned. deny; refuse. Dative person denied, genitive thing denied.\par
% glossary-block:B00505:end
% glossary-block:B00506:begin
\needspace{3\baselineskip}
\textbf{forwierð} — pres. ind. 3 sg.: see \textbf{forweorðan}.\par
% glossary-block:B00506:end
% glossary-block:B00507:begin
\needspace{3\baselineskip}
\textbf{fōt} (n. m.); gen. sg. fōtes; nom. pl. fēt. foot. Supplied body-part form feet; distinct from the length-measure sense. foot. Unit of length, distinct from the body-part sense.\par
% glossary-block:B00507:end
% glossary-block:B00508:begin
\needspace{3\baselineskip}
\textbf{fōtsceamul} (n. m.); gen. sg. fōtsceamules; nom. pl. fōtsceamulas. footstool.\par
% glossary-block:B00508:end
% glossary-block:B00509:begin
\needspace{3\baselineskip}
\textbf{fox} (n. m.); gen. sg. foxes; nom. pl. foxas. fox.\par
% glossary-block:B00509:end
% glossary-block:B00510:begin
\needspace{3\baselineskip}
\textbf{fram} (prep.). from. Separation, with dative in the taught frame. from. Only the supplied rest-from-work construction.\par
% glossary-block:B00510:end
% glossary-block:B00511:begin
\needspace{3\baselineskip}
\textbf{frætwung} (n. f.); gen. sg. frætwunge; nom. pl. frætwunga. array.\par
% glossary-block:B00511:end
% glossary-block:B00512:begin
\needspace{3\baselineskip}
\textbf{frēa} (n. m.). lord. Reference in a source or explanatory note.\par
% glossary-block:B00512:end
% glossary-block:B00513:begin
\needspace{3\baselineskip}
\textbf{frēcnes} (n. f.). danger; peril.\par
% glossary-block:B00513:end
% glossary-block:B00514:begin
\needspace{3\baselineskip}
\textbf{fręmman} (v. weak); pret. 1 sg. fręmede; past ptcp. (ge)fręmed. benefit. Dative beneficiary.\par
% glossary-block:B00514:end
% glossary-block:B00515:begin
\needspace{3\baselineskip}
\textbf{frēo} (adj.). free.\par
% glossary-block:B00515:end
% glossary-block:B00516:begin
\needspace{3\baselineskip}
\textbf{frēond} (n. m.); gen. sg. frēondes; nom. pl. frīend. friend.\par
% glossary-block:B00516:end
% glossary-block:B00517:begin
\needspace{3\baselineskip}
\textbf{frīend} — nom./acc. pl., variant: see \textbf{frēond}.\par
% glossary-block:B00517:end
% glossary-block:B00518:begin
\needspace{3\baselineskip}
\textbf{frōfor} (n. f.); gen. sg. frōfre; nom. pl. frōfra. comfort; consolation.\par
% glossary-block:B00518:end
% glossary-block:B00519:begin
\needspace{3\baselineskip}
\textbf{frosc} (n.). frog. Pronunciation example.\par
% glossary-block:B00519:end
% glossary-block:B00520:begin
\needspace{3\baselineskip}
\textbf{fugol} (n. m.); gen. sg. fugles; nom. pl. fuglas. bird.\par
% glossary-block:B00520:end
% glossary-block:B00521:begin
\needspace{3\baselineskip}
\textbf{fūl} (adj.). foul; vile.\par
% glossary-block:B00521:end
% glossary-block:B00522:begin
\needspace{3\baselineskip}
\textbf{full} (adj.). full. Pronunciation example.\par
% glossary-block:B00522:end
% glossary-block:B00523:begin
\needspace{3\baselineskip}
\textbf{full} (adv.). very.\par
% glossary-block:B00523:end
% glossary-block:B00524:begin
\needspace{3\baselineskip}
\textbf{fullfręmman} (v.). finish; complete. Supplied passive use.\par
% glossary-block:B00524:end
% glossary-block:B00525:begin
\needspace{3\baselineskip}
\textbf{fultum} (n. m.); gen. sg. fultumes. help; aid. help; support.\par
% glossary-block:B00525:end
% glossary-block:B00526:begin
\needspace{3\baselineskip}
\textbf{fuluhtere} (n. m.); gen. sg. fuluhteres. baptizer; Baptist. Baptist as John's title in the assigned passage.\par
% glossary-block:B00526:end
% glossary-block:B00527:begin
\needspace{3\baselineskip}
\textbf{furlang} (n. n.); gen. sg. furlanges; nom. pl. furlang. furlong.\par
% glossary-block:B00527:end
% glossary-block:B00528:begin
\needspace{3\baselineskip}
\textbf{furðor} (adv.). further; any more. Adverb in the supplied no-spirit-more expression.\par
% glossary-block:B00528:end
% glossary-block:B00529:begin
\needspace{3\baselineskip}
\textbf{furðra} (adj.). foremost; prior.\par
% glossary-block:B00529:end
% glossary-block:B00530:begin
\needspace{3\baselineskip}
\textbf{fȳr} (n. n.); gen. sg. fȳres; nom. pl. fȳr. fire.\par
% glossary-block:B00530:end
% glossary-block:B00531:begin
\needspace{3\baselineskip}
\textbf{fyrmest} (adv.). first. first of all; at the outset.\par
% glossary-block:B00531:end
% glossary-block:B00532:begin
\needspace{3\baselineskip}
\textbf{fȳrspearca} (n. m.); gen. sg. fȳrspearcan; nom. pl. fȳrspearcan. spark.\par
% glossary-block:B00532:end
% book-navigation: standalone bookmark replaced by volume hierarchy
% glossary-block:B00533:begin
\needspace{5\baselineskip}
\subsection*{G}\BookHeading{nav-glossary-old-english-0009}{subsubsection}{G}{0}
% glossary-block:B00533:end
% glossary-block:B00534:begin
\textbf{gagates} (n.). jet (Latin loan-name). The stone; Latin loan-name in the source.\par
% glossary-block:B00534:end
% glossary-block:B00535:begin
\needspace{3\baselineskip}
\textbf{gālnes} (n. f.). lust. Assigned source sense.\par
% glossary-block:B00535:end
% glossary-block:B00536:begin
\needspace{3\baselineskip}
\textbf{gān} (v. irregular); pres. 1 sg. gā; pres. 3 sg. gǣð; pres. pl. gāð; pret. 1 sg. ēode; pret. pl. ēodon; past ptcp. (ge)gān. go; walk. go. Supplied imperative gang.\par
% glossary-block:B00536:end
% glossary-block:B00537:begin
\needspace{3\baselineskip}
\textbf{gangan} (v.). go; enter. Supplied circumstantial participial phrase.\par
% glossary-block:B00537:end
% glossary-block:B00538:begin
\needspace{3\baselineskip}
\textbf{gāst} (n. m.); gen. sg. gāstes; nom. pl. gāstas. spirit.\par
% glossary-block:B00538:end
% glossary-block:B00539:begin
\needspace{3\baselineskip}
\textbf{gāt} (n. f.); gen. sg. gāte; nom. pl. gǣt. goat.\par
% glossary-block:B00539:end
% glossary-block:B00540:begin
\needspace{3\baselineskip}
\textbf{gærs} (n. n.); gen. sg. gærses; nom. pl. gærs. grass.\par
% glossary-block:B00540:end
% glossary-block:B00541:begin
\needspace{3\baselineskip}
\textbf{gǣt} — dat. sg.; nom./acc. pl.; also gen. sg. variant: see \textbf{gāt}.\par
% glossary-block:B00541:end
% glossary-block:B00542:begin
\needspace{3\baselineskip}
\textbf{gǣð} — pres. ind. 3 sg.: see \textbf{gān}.\par
% glossary-block:B00542:end
% glossary-block:B00543:begin
\needspace{3\baselineskip}
\textbf{ge} (conj.). both ... and. both … and. Paired ge … ge.\par
% glossary-block:B00543:end
% glossary-block:B00544:begin
\needspace{3\baselineskip}
\textbf{gē} — nom. pl.: see \textbf{ðū}.\par
% glossary-block:B00544:end
% glossary-block:B00545:begin
\needspace{3\baselineskip}
\textbf{gēar} (n. n.); gen. sg. gēares; nom. pl. gēar. year.\par
% glossary-block:B00545:end
% glossary-block:B00546:begin
\needspace{3\baselineskip}
\textbf{gearcian} (v. weak); pret. 1 sg. gearcode; past ptcp. gearcod. prepare.\par
% glossary-block:B00546:end
% glossary-block:B00547:begin
\needspace{3\baselineskip}
\textbf{gearu} (adj.). ready.\par
% glossary-block:B00547:end
% glossary-block:B00548:begin
\needspace{3\baselineskip}
\textbf{geat} (n. n.); gen. sg. geates; nom. pl. gatu. gate.\par
% glossary-block:B00548:end
% glossary-block:B00549:begin
\needspace{3\baselineskip}
\textbf{gebed} (n. n.). prayer.\par
% glossary-block:B00549:end
% glossary-block:B00550:begin
\needspace{3\baselineskip}
\textbf{gebeorg} (n. n.). protection; defence.\par
% glossary-block:B00550:end
% glossary-block:B00551:begin
\needspace{3\baselineskip}
\textbf{gebēorscipe} (n. m.); gen. sg. gebēorscipes; nom. pl. gebēorscipas. feast; banquet.\par
% glossary-block:B00551:end
% glossary-block:B00552:begin
\needspace{3\baselineskip}
\textbf{gebiddan} (v. strong); pret. 1 sg. gebæd; pret. pl. gebǣdon; past ptcp. gebeden. pray. Supplied wæs gebiddende, rendered prayed in this source. pray. With an accusative reflexive pronoun in the assigned prayer frame.\par
% glossary-block:B00552:end
% glossary-block:B00553:begin
\needspace{3\baselineskip}
\textbf{geblēdsian} — citation/spelling variant: see \textbf{geblētsian}.\par
% glossary-block:B00553:end
% glossary-block:B00554:begin
\needspace{3\baselineskip}
\textbf{geblētsian} (v. weak); pret. 1 sg. geblētsode; past ptcp. geblētsod. bless.\par
% glossary-block:B00554:end
% glossary-block:B00555:begin
\needspace{3\baselineskip}
\textbf{gebrec} (n. n.). crash; noise; din. Supplied accusative singular.\par
% glossary-block:B00555:end
% glossary-block:B00556:begin
\needspace{3\baselineskip}
\textbf{gebringan} (v.). bring. Supplied preterite with plural winds as subject; retain the prefix.\par
% glossary-block:B00556:end
% glossary-block:B00557:begin
\needspace{3\baselineskip}
\textbf{gebrōhton} — pret. ind. pl.; only supplied form: see \textbf{gebringan}.\par
% glossary-block:B00557:end
% glossary-block:B00558:begin
\needspace{3\baselineskip}
\textbf{gebyrd} (n. f.); gen. sg. gebyrde; nom. pl. gebyrda. birth.\par
% glossary-block:B00558:end
% glossary-block:B00559:begin
\needspace{3\baselineskip}
\textbf{gecwēme} (adj.). pleasing; acceptable.\par
% glossary-block:B00559:end
% glossary-block:B00560:begin
\needspace{3\baselineskip}
\textbf{gecwēmlīce} (adv.). pleasantly; agreeably.\par
% glossary-block:B00560:end
% glossary-block:B00561:begin
\needspace{3\baselineskip}
\textbf{gecweðan} (v. strong); pret. 1 sg. gecwæð; pret. pl. gecwǣdon; past ptcp. gecweden; pres. 3 sg. gecwið. speak; say.\par
% glossary-block:B00561:end
% glossary-block:B00562:begin
\needspace{3\baselineskip}
\textbf{gecwið} — pres. ind. 3 sg.: see \textbf{gecweðan}.\par
% glossary-block:B00562:end
% glossary-block:B00563:begin
\needspace{3\baselineskip}
\textbf{ge-cȳðan} (v.). show.\par
% glossary-block:B00563:end
% glossary-block:B00564:begin
\needspace{3\baselineskip}
\textbf{gedǣlan} (v.). divide; separate. Preserve the assigned separation frame.\par
% glossary-block:B00564:end
% glossary-block:B00565:begin
\needspace{3\baselineskip}
\textbf{gedeorf} (n. n.); gen. sg. gedeorfes; nom. pl. gedeorf. labor; toil.\par
% glossary-block:B00565:end
% glossary-block:B00566:begin
\needspace{3\baselineskip}
\textbf{gedēð} — pres. ind. 3 sg.: see \textbf{gedōn}.\par
% glossary-block:B00566:end
% glossary-block:B00567:begin
\needspace{3\baselineskip}
\textbf{gedōn} (v.); past ptcp. gedōn. do; make. Supplied passive fulfilment formula. do; perform. make; arrange. The present form/arrangement of the translation in the source. make; cause to be.\par
% glossary-block:B00567:end
% glossary-block:B00568:begin
\needspace{3\baselineskip}
\textbf{geęndian} (v.). end; come to an end.\par
% glossary-block:B00568:end
% glossary-block:B00569:begin
\needspace{3\baselineskip}
\textbf{geęndung} (n. f.). end; close.\par
% glossary-block:B00569:end
% glossary-block:B00570:begin
\needspace{3\baselineskip}
\textbf{gefeallan} (v.). chance upon. fall in love. Figurative falling of the heart; supplied idiom differs in prefix from befeallan in 42.1.\par
% glossary-block:B00570:end
% glossary-block:B00571:begin
\needspace{3\baselineskip}
\textbf{gefēra} (n. m.); gen. sg. gefēran; nom. pl. gefēran. companion; fellow.\par
% glossary-block:B00571:end
% glossary-block:B00572:begin
\needspace{3\baselineskip}
\textbf{gefęrian} (v.); past ptcp. gefęred. carry; bear. Supplied passive movement; was carried/borne.\par
% glossary-block:B00572:end
% glossary-block:B00573:begin
\needspace{3\baselineskip}
\textbf{gefērscipe} (n. m.); gen. sg. gefērscipes; nom. pl. gefērscipas. companionship.\par
% glossary-block:B00573:end
% glossary-block:B00574:begin
\needspace{3\baselineskip}
\textbf{gefiðred} (adj.). feathered.\par
% glossary-block:B00574:end
% glossary-block:B00575:begin
\needspace{3\baselineskip}
\textbf{geflīt} (n. n.); gen. sg. geflītes; nom. pl. geflīt. dispute; strife.\par
% glossary-block:B00575:end
% glossary-block:B00576:begin
\needspace{3\baselineskip}
\textbf{gefręmman} (v. weak); pret. 1 sg. gefręmede; past ptcp. gefręmed. perform; accomplish. A work is the thing accomplished.\par
% glossary-block:B00576:end
% glossary-block:B00577:begin
\needspace{3\baselineskip}
\textbf{gefultumian} (v.). help; assist.\par
% glossary-block:B00577:end
% glossary-block:B00578:begin
\needspace{3\baselineskip}
\textbf{gefyllan} (v. weak); pret. 1 sg. gefylde; past ptcp. gefylled. finish; accomplish. A work is the thing completed; distinct from fill. fill. fulfil. Accusative prediction fulfilled; passive subjunctive in the assigned purpose construction.\par
% glossary-block:B00578:end
% glossary-block:B00579:begin
\needspace{3\baselineskip}
\textbf{gegaderian} (v. weak); pret. 1 sg. gegaderode; past ptcp. gegaderod. gather. Accusative people or things gathered.\par
% glossary-block:B00579:end
% glossary-block:B00580:begin
\needspace{3\baselineskip}
\textbf{gegaderung} (n. f.); gen. sg. gegaderunge; nom. pl. gegaderunga. gathering; collection; congregation. Congregations of waters in the source.\par
% glossary-block:B00580:end
% glossary-block:B00581:begin
\needspace{3\baselineskip}
\textbf{gegierela} (n. m.); gen. sg. gegierelan; nom. pl. gegierelan. garment; clothing.\par
% glossary-block:B00581:end
% glossary-block:B00582:begin
\needspace{3\baselineskip}
\textbf{geglęngan} (v.); past ptcp. geglęnged. adorn. The supplied passive adornment expression uses mid + dative material.\par
% glossary-block:B00582:end
% glossary-block:B00583:begin
\needspace{3\baselineskip}
\textbf{gehālgian} (v. weak); pret. 1 sg. gehālgode; past ptcp. gehālgod. hallow; consecrate.\par
% glossary-block:B00583:end
% glossary-block:B00584:begin
\needspace{3\baselineskip}
\textbf{gehātan} (v. strong); pret. 1 sg. gehēt; pret. pl. gehēton; past ptcp. gehāten; pres. 3 sg. gehǣtt. promise. Accusative thing promised; a dative recipient may be added.\par
% glossary-block:B00584:end
% glossary-block:B00585:begin
\needspace{3\baselineskip}
\textbf{gehǣtt} — pres. ind. 3 sg.: see \textbf{gehātan}.\par
% glossary-block:B00585:end
% glossary-block:B00586:begin
\needspace{3\baselineskip}
\textbf{gehīeran} (v.). hear. Supplied prefixed third-singular preterite. hear; listen.\par
% glossary-block:B00586:end
% glossary-block:B00587:begin
\needspace{3\baselineskip}
\textbf{gehīersumian} (v.). serve; obey. Supplied source form and government.\par
% glossary-block:B00587:end
% glossary-block:B00588:begin
\needspace{3\baselineskip}
\textbf{gehīersumnes} (n. f.); gen. sg. gehīersumnesse; nom. pl. gehīersumnessa. obedience.\par
% glossary-block:B00588:end
% glossary-block:B00589:begin
\needspace{3\baselineskip}
\textbf{gehwā / gehwæt} (pron.). each one; everyone; everything. Personal gehwā and neuter gehwæt; the early reference status does not by itself establish production.\par
% glossary-block:B00589:end
% glossary-block:B00590:begin
\needspace{3\baselineskip}
\textbf{gehwǣm} — dat. m./f. sg.; dat. n. sg.: see \textbf{gehwā / gehwæt}.\par
% glossary-block:B00590:end
% glossary-block:B00591:begin
\needspace{3\baselineskip}
\textbf{gehwæs} — gen. m./f. sg.; gen. n. sg.: see \textbf{gehwā / gehwæt}.\par
% glossary-block:B00591:end
% glossary-block:B00592:begin
\needspace{3\baselineskip}
\textbf{gehwæt} — nom. n. sg.; acc. n. sg.: see \textbf{gehwā / gehwæt}.\par
% glossary-block:B00592:end
% glossary-block:B00593:begin
\needspace{3\baselineskip}
\textbf{gehwilc} (indefinite adjective). each; every.\par
% glossary-block:B00593:end
% glossary-block:B00594:begin
\needspace{3\baselineskip}
\textbf{gehwone} — acc. m./f. sg.: see \textbf{gehwā / gehwæt}.\par
% glossary-block:B00594:end
% glossary-block:B00595:begin
\needspace{3\baselineskip}
\textbf{gelǣccan} (v.). catch. Supplied caught form.\par
% glossary-block:B00595:end
% glossary-block:B00596:begin
\needspace{3\baselineskip}
\textbf{gelǣdan} (v.). bring; carry. Reference in a source or explanatory note.\par
% glossary-block:B00596:end
% glossary-block:B00597:begin
\needspace{3\baselineskip}
\textbf{gelæhte} — pret. sg.; supplied source form: see \textbf{gelǣccan}.\par
% glossary-block:B00597:end
% glossary-block:B00598:begin
\needspace{3\baselineskip}
\textbf{gelǣran} (v.). teach; instruct. Reference in a source or explanatory note.\par
% glossary-block:B00598:end
% glossary-block:B00599:begin
\needspace{3\baselineskip}
\textbf{gelǣred} (adj.). educated; trained; skilled.\par
% glossary-block:B00599:end
% glossary-block:B00600:begin
\needspace{3\baselineskip}
\textbf{gelēafa} (n. m.); gen. sg. gelēafan; nom. pl. gelēafan. faith.\par
% glossary-block:B00600:end
% glossary-block:B00601:begin
\needspace{3\baselineskip}
\textbf{geleornian} (v.); past ptcp. geleornod. learn; study. learn; master. Learn a skill or performance in the source.\par
% glossary-block:B00601:end
% glossary-block:B00602:begin
\needspace{3\baselineskip}
\textbf{gelīc} (adj.). equal; like. Dative standard of comparison.\par
% glossary-block:B00602:end
% glossary-block:B00603:begin
\needspace{3\baselineskip}
\textbf{gelīca} (n. m.). equal; peer.\par
% glossary-block:B00603:end
% glossary-block:B00604:begin
\needspace{3\baselineskip}
\textbf{gelīce} (adv.). likewise; similarly.\par
% glossary-block:B00604:end
% glossary-block:B00605:begin
\needspace{3\baselineskip}
\textbf{gelīcnes} (n. f.); gen. sg. gelīcnesse; nom. pl. gelīcnessa. likeness.\par
% glossary-block:B00605:end
% glossary-block:B00606:begin
\needspace{3\baselineskip}
\textbf{gelīefan} (v. weak); pret. 1 sg. gelīefde; past ptcp. gelīefed. believe.\par
% glossary-block:B00606:end
% glossary-block:B00607:begin
\needspace{3\baselineskip}
\textbf{gelīefst} (v.; supplied pres. ind. 2 sg.). permit; allow. You permit; you allow. In gif ðū mē gelīefst, mē is the dative person granted permission; speaking is understood.\par
% glossary-block:B00607:end
% glossary-block:B00608:begin
\needspace{3\baselineskip}
\textbf{gelimpan} (v.); pres. 3 sg. gelimpð. happen; befall. Supplied preterite subjunctive in might befall us. happen; come about. Supplied preterite source form. happen; come about.\par
% glossary-block:B00608:end
% glossary-block:B00609:begin
\needspace{3\baselineskip}
\textbf{gelōgian} (v.); pres. 3 sg. gelōgað. place; set.\par
% glossary-block:B00609:end
% glossary-block:B00610:begin
\needspace{3\baselineskip}
\textbf{gelōmlice} (adv.). frequently; regularly.\par
% glossary-block:B00610:end
% glossary-block:B00611:begin
\needspace{3\baselineskip}
\textbf{geman} — pres. ind. 1/3 sg.: see \textbf{gemunan}.\par
% glossary-block:B00611:end
% glossary-block:B00612:begin
\needspace{3\baselineskip}
\textbf{gemang} (prep.). among. Accusative group in the selected sending-among frame.\par
% glossary-block:B00612:end
% glossary-block:B00613:begin
\needspace{3\baselineskip}
\textbf{gemanigfieldan} (v. weak); pret. 1 sg. gemanigfielde; past ptcp. gemanigfield. be multiplied. Assigned passive with agreeing gemanigfield/gemanigfielde; no independent active causative frame is assigned.\par
% glossary-block:B00613:end
% glossary-block:B00614:begin
\needspace{3\baselineskip}
\textbf{gemǣne} (adj.). common; shared.\par
% glossary-block:B00614:end
% glossary-block:B00615:begin
\needspace{3\baselineskip}
\textbf{gemǣre} (n. n.); gen. sg. gemǣres; nom. pl. gemǣru. boundary; border.\par
% glossary-block:B00615:end
% glossary-block:B00616:begin
\needspace{3\baselineskip}
\textbf{gemęngan} (v.). join; mingle. Supplied source form. mix. Mix sound with song in the supplied performance expression.\par
% glossary-block:B00616:end
% glossary-block:B00617:begin
\needspace{3\baselineskip}
\textbf{gemētan} (v. weak); pret. 1 sg. gemētte; past ptcp. gemētt. find. Supplied past participle gemētt in a passive. find; meet. Accusative person or thing found/met.\par
% glossary-block:B00617:end
% glossary-block:B00618:begin
\needspace{3\baselineskip}
\textbf{gemiltsian} (v.). pity; have compassion on. Dative person.\par
% glossary-block:B00618:end
% glossary-block:B00619:begin
\needspace{3\baselineskip}
\textbf{gemiltsiend} (n. m.). one who shows mercy; pitier.\par
% glossary-block:B00619:end
% glossary-block:B00620:begin
\needspace{3\baselineskip}
\textbf{gemunan} (v. preterite-present); pres. 1/3 sg. geman; pres. pl. gemunon; pret. 1 sg. gemunde; past ptcp. gemunen. remember.\par
% glossary-block:B00620:end
% glossary-block:B00621:begin
\needspace{3\baselineskip}
\textbf{genēahhe} (adv.). often; frequently.\par
% glossary-block:B00621:end
% glossary-block:B00622:begin
\needspace{3\baselineskip}
\textbf{genēalǣcan} (v. weak); pret. 1 sg. genēalǣcte; past ptcp. genēalǣced. approach; draw near. Dative person, place or event approached. approach. Supplied source form approached.\par
% glossary-block:B00622:end
% glossary-block:B00623:begin
\needspace{3\baselineskip}
\textbf{genēosian} (v. weak); pret. 1 sg. genēosode; past ptcp. genēosod. visit. Accusative person in the selected prefixed-verb frame.\par
% glossary-block:B00623:end
% glossary-block:B00624:begin
\needspace{3\baselineskip}
\textbf{geniman} (v.). take.\par
% glossary-block:B00624:end
% glossary-block:B00625:begin
\needspace{3\baselineskip}
\textbf{genyhtsumian} (v.). suffice; be sufficient for. Complete supplied construction with dative beneficiary.\par
% glossary-block:B00625:end
% glossary-block:B00626:begin
\needspace{3\baselineskip}
\textbf{geoguð} (n. f.). youth.\par
% glossary-block:B00626:end
% glossary-block:B00627:begin
\needspace{3\baselineskip}
\textbf{geond} (prep.). along; through. With accusative.\par
% glossary-block:B00627:end
% glossary-block:B00628:begin
\needspace{3\baselineskip}
\textbf{geong} (adj.); supplied superl. acc./gen. m. sg. gingstan. young. In supplied superlative gingstan: youngest.\par
% glossary-block:B00628:end
% glossary-block:B00629:begin
\needspace{3\baselineskip}
\textbf{geopenian} (v. weak); pret. 1 sg. geopenode; past ptcp. geopenod. open. Accusative thing opened, including mouth or eyes.\par
% glossary-block:B00629:end
% glossary-block:B00630:begin
\needspace{3\baselineskip}
\textbf{geornfull} (adj.). occupied; intent. Source mental-attention context.\par
% glossary-block:B00630:end
% glossary-block:B00631:begin
\needspace{3\baselineskip}
\textbf{geornlice} (adv.). earnestly; attentively.\par
% glossary-block:B00631:end
% glossary-block:B00632:begin
\needspace{3\baselineskip}
\textbf{geręccan} (v.). explain; render.\par
% glossary-block:B00632:end
% glossary-block:B00633:begin
\needspace{3\baselineskip}
\textbf{gerēfa} (n. m.); gen. sg. gerēfan; nom. pl. gerēfan. steward; reeve.\par
% glossary-block:B00633:end
% glossary-block:B00634:begin
\needspace{3\baselineskip}
\textbf{gereord} (n. n.). meal; banquet.\par
% glossary-block:B00634:end
% glossary-block:B00635:begin
\needspace{3\baselineskip}
\textbf{geręstan} (v. weak); pret. 1 sg. geręste; past ptcp. geręst. rest; repose.\par
% glossary-block:B00635:end
% glossary-block:B00636:begin
\needspace{3\baselineskip}
\textbf{gerīman} (v.). number; count. Exact supplied infinitive gerīman; not an independently assigned full paradigm.\par
% glossary-block:B00636:end
% glossary-block:B00637:begin
\needspace{3\baselineskip}
\textbf{gerīmcræft} (n. m.); gen. sg. gerīmcræftes. arithmetic; computation.\par
% glossary-block:B00637:end
% glossary-block:B00638:begin
\needspace{3\baselineskip}
\textbf{gesǣd} — past ptcp., variant: see \textbf{sęcgan}.\par
% glossary-block:B00638:end
% glossary-block:B00639:begin
\needspace{3\baselineskip}
\textbf{gesceaft} (n. f.); gen. sg. gesceafte; nom. pl. gesceafta. creature; created being or thing.\par
% glossary-block:B00639:end
% glossary-block:B00640:begin
\needspace{3\baselineskip}
\textbf{gescieldan} (v.). shield; protect.\par
% glossary-block:B00640:end
% glossary-block:B00641:begin
\needspace{3\baselineskip}
\textbf{gescieppan} (v. strong); pret. 1 sg. gescēop; pret. pl. gescēopon; past ptcp. gesceapen. create.\par
% glossary-block:B00641:end
% glossary-block:B00642:begin
\needspace{3\baselineskip}
\textbf{gescrȳdan} (v. weak); pret. 1 sg. gescrȳdde; past ptcp. gescrȳd. clothe.\par
% glossary-block:B00642:end
% glossary-block:B00643:begin
\needspace{3\baselineskip}
\textbf{gesęllan} (v.). give.\par
% glossary-block:B00643:end
% glossary-block:B00644:begin
\needspace{3\baselineskip}
\textbf{gesęndan} (v.). send back.\par
% glossary-block:B00644:end
% glossary-block:B00645:begin
\needspace{3\baselineskip}
\textbf{gesēon} (v. strong); pret. 1 sg. geseah; pret. pl. gesāwon; past ptcp. gesewen; pres. 3 sg. gesiehð. see. see. Accusative perceived actor and plain infinitive. see. Absolute use, with no expressed object. see. With a content clause stating the circumstance seen. regard; take notice of. see. With the assigned clause describing how a state of affairs stands.\par
% glossary-block:B00645:end
% glossary-block:B00646:begin
\needspace{3\baselineskip}
\textbf{gesęttan} (v. weak); pret. 1 sg. gesętte; past ptcp. gesętt. set; place. Supplied gesętte. placed; situated. Supplied participial complement of understood us; not the identical-looking finite preterite. set; place; establish. appoint as lord. Supplied frame with person appointed, tō + dative role, and dative people governed.\par
% glossary-block:B00646:end
% glossary-block:B00647:begin
\needspace{3\baselineskip}
\textbf{gesewenlic} (adj.). visible.\par
% glossary-block:B00647:end
% glossary-block:B00648:begin
\needspace{3\baselineskip}
\textbf{gesiehð} — pres. ind. 3 sg.: see \textbf{gesēon}.\par
% glossary-block:B00648:end
% glossary-block:B00649:begin
\needspace{3\baselineskip}
\textbf{gesihð} (n. f.). sight; gaze. Supplied accusative singular.\par
% glossary-block:B00649:end
% glossary-block:B00650:begin
\needspace{3\baselineskip}
\textbf{geslēan} (v.). strike; smite. Supplied past/participial source form. strike. Supplied source form struck.\par
% glossary-block:B00650:end
% glossary-block:B00651:begin
\needspace{3\baselineskip}
\textbf{gesmierwan} (v.); past ptcp. gesmierwed. anoint. Supplied source form anointed.\par
% glossary-block:B00651:end
% glossary-block:B00652:begin
\needspace{3\baselineskip}
\textbf{gesmyltan} (v.). become calm. Supplied preterite of the sea becoming calm.\par
% glossary-block:B00652:end
% glossary-block:B00653:begin
\needspace{3\baselineskip}
\textbf{gestandan} (v.); pres. 3 sg. gestęnt. assail; come upon. Supplied source sense.\par
% glossary-block:B00653:end
% glossary-block:B00654:begin
\needspace{3\baselineskip}
\textbf{gestaðelian} (v.). establish; make steadfast.\par
% glossary-block:B00654:end
% glossary-block:B00655:begin
\needspace{3\baselineskip}
\textbf{gestęnt} — pres. ind. 3 sg.: see \textbf{gestandan}.\par
% glossary-block:B00655:end
% glossary-block:B00656:begin
\needspace{3\baselineskip}
\textbf{gestillan} (v. weak); pret. 1 sg. gestilde; past ptcp. gestilled. quell; settle. Accusative disturbance settled; distinguish intransitive become still. cease; become still. Intransitive in the source storm relative.\par
% glossary-block:B00656:end
% glossary-block:B00657:begin
\needspace{3\baselineskip}
\textbf{gestrangian} (v.). strengthen.\par
% glossary-block:B00657:end
% glossary-block:B00658:begin
\needspace{3\baselineskip}
\textbf{gestrēon} (n. n.); gen. sg. gestrēones; nom. pl. gestrēon. gain; profit.\par
% glossary-block:B00658:end
% glossary-block:B00659:begin
\needspace{3\baselineskip}
\textbf{gesund} (adj.). well; in good health.\par
% glossary-block:B00659:end
% glossary-block:B00660:begin
\needspace{3\baselineskip}
\textbf{geswęngan} (v.); past ptcp. geswęnged. toss; shake. Supplied participial source use.\par
% glossary-block:B00660:end
% glossary-block:B00661:begin
\needspace{3\baselineskip}
\textbf{gesweotolian} (v.); pres. 3 sg. gesweotolað; past ptcp. gesweotolod. reveal; make clear. Supplied source form. make manifest; reveal.\par
% glossary-block:B00661:end
% glossary-block:B00662:begin
\needspace{3\baselineskip}
\textbf{geswīcan} (v. strong); pret. 1 sg. geswāc; pret. pl. geswicon; past ptcp. geswicen. cease. Genitive thing ceased from. fail; cease. Intransitive without a genitive complement in the selected frame.\par
% glossary-block:B00662:end
% glossary-block:B00663:begin
\needspace{3\baselineskip}
\textbf{gesyndig} (adj.). favourable; fair. Of winds.\par
% glossary-block:B00663:end
% glossary-block:B00664:begin
\needspace{3\baselineskip}
\textbf{getǣcan} (v.); past ptcp. getǣht. show; direct. Supplied source form.\par
% glossary-block:B00664:end
% glossary-block:B00665:begin
\needspace{3\baselineskip}
\textbf{getǣht} — past ptcp.; supplied source form: see \textbf{getǣcan}.\par
% glossary-block:B00665:end
% glossary-block:B00666:begin
\needspace{3\baselineskip}
\textbf{getæl} (n. n.). number; company.\par
% glossary-block:B00666:end
% glossary-block:B00667:begin
\needspace{3\baselineskip}
\textbf{getēon} (v.); past ptcp. getogen. train; educate.\par
% glossary-block:B00667:end
% glossary-block:B00668:begin
\needspace{3\baselineskip}
\textbf{getimbrung} (n. f.); gen. sg. getimbrunge; nom. pl. getimbrunga. building; structure built.\par
% glossary-block:B00668:end
% glossary-block:B00669:begin
\needspace{3\baselineskip}
\textbf{geðanc} (n. n.). thought; mind.\par
% glossary-block:B00669:end
% glossary-block:B00670:begin
\needspace{3\baselineskip}
\textbf{geðeahtend} (n. m.); gen. sg. geðeahtendes; nom. pl. geðeahtend. counsellor.\par
% glossary-block:B00670:end
% glossary-block:B00671:begin
\needspace{3\baselineskip}
\textbf{geðęncean} (v.). remember; bear in mind.\par
% glossary-block:B00671:end
% glossary-block:B00672:begin
\needspace{3\baselineskip}
\textbf{geðrǣstan} (v.). afflict; beat. Supplied participial source use.\par
% glossary-block:B00672:end
% glossary-block:B00673:begin
\needspace{3\baselineskip}
\textbf{geðūht} — past ptcp.: see \textbf{ðyncan}.\par
% glossary-block:B00673:end
% glossary-block:B00674:begin
\needspace{3\baselineskip}
\textbf{gewǣgan} (v.); past ptcp. gewǣged. buffet; shake. Supplied participial source use.\par
% glossary-block:B00674:end
% glossary-block:B00675:begin
\needspace{3\baselineskip}
\textbf{geweald} (n. n.); gen. sg. gewealdes; nom. pl. geweald. control; rule; dominion. Ordinary singular; on gewealde is the taught control phrase.\par
% glossary-block:B00675:end
% glossary-block:B00676:begin
\needspace{3\baselineskip}
\textbf{gewęndan} (v. weak); pret. 1 sg. gewęnde; past ptcp. gewęnd. return; turn. translate; turn from one language into another. turn; change back. Supplied source form.\par
% glossary-block:B00676:end
% glossary-block:B00677:begin
\needspace{3\baselineskip}
\textbf{geweorðan} (v.); past ptcp. geworden. come to be; become. come to pass; befall. become; come about.\par
% glossary-block:B00677:end
% glossary-block:B00678:begin
\needspace{3\baselineskip}
\textbf{geweorðian} (v. weak); pret. 1 sg. geweorðode; past ptcp. geweorðod. distinguish; adorn. Accusative person adorned; mid + dative equipment or adornment.\par
% glossary-block:B00678:end
% glossary-block:B00679:begin
\needspace{3\baselineskip}
\textbf{gewiht} (n. n.). weight.\par
% glossary-block:B00679:end
% glossary-block:B00680:begin
\needspace{3\baselineskip}
\textbf{gewilnian} (v. weak); pret. 1 sg. gewilnode; past ptcp. gewilnod. desire; wish.\par
% glossary-block:B00680:end
% glossary-block:B00681:begin
\needspace{3\baselineskip}
\textbf{gewinn} (n. n.). distress; struggle. labour; work. Noun, distinct from labour as a verb.\par
% glossary-block:B00681:end
% glossary-block:B00682:begin
\needspace{3\baselineskip}
\textbf{gewinna} (n. m.). enemy; opponent. Supplied plural source forms.\par
% glossary-block:B00682:end
% glossary-block:B00683:begin
\needspace{3\baselineskip}
\textbf{gewissian} (v.). guide; direct.\par
% glossary-block:B00683:end
% glossary-block:B00684:begin
\needspace{3\baselineskip}
\textbf{gewītan} (v. strong); pret. 1 sg. gewāt; pret. pl. gewiton; past ptcp. gewiten. depart; go away. Intransitive; fram + dative departure point is optional. The assigned frame requires no reflexive. pass away; cease to hold. Of a decree.\par
% glossary-block:B00684:end
% glossary-block:B00685:begin
\needspace{3\baselineskip}
\textbf{gewitan} (v.). find out; learn. Knowledge verb, distinct from gewītan depart.\par
% glossary-block:B00685:end
% glossary-block:B00686:begin
\needspace{3\baselineskip}
\textbf{geworhte} — pret. sg.: see \textbf{gewyrcean}.\par
% glossary-block:B00686:end
% glossary-block:B00687:begin
\needspace{3\baselineskip}
\textbf{gewrit} (n. n.); gen. sg. gewrites; nom. pl. gewritu. letter; written document.\par
% glossary-block:B00687:end
% glossary-block:B00688:begin
\needspace{3\baselineskip}
\textbf{gewuna} (n. m.). custom; habit.\par
% glossary-block:B00688:end
% glossary-block:B00689:begin
\needspace{3\baselineskip}
\textbf{gewunian} (v. weak); pret. 1 sg. gewunode; past ptcp. gewunod. be accustomed; be wont.\par
% glossary-block:B00689:end
% glossary-block:B00690:begin
\needspace{3\baselineskip}
\textbf{gewyrcean} (v. weak); pret. 1 sg. geworhte; past ptcp. geworht. make; build. make; render. Accusative object with an agreeing predicate noun/adjective.\par
% glossary-block:B00690:end
% glossary-block:B00691:begin
\needspace{3\baselineskip}
\textbf{giefan} (v. strong); pret. 1 sg. geaf; pret. pl. gēafon; past ptcp. giefen. give.\par
% glossary-block:B00691:end
% glossary-block:B00692:begin
\needspace{3\baselineskip}
\textbf{giefu} (n. f.). gift. grace; divine favour. gift. Supplied plural gifts.\par
% glossary-block:B00692:end
% glossary-block:B00693:begin
\needspace{3\baselineskip}
\textbf{gieldan} (v. strong); pret. 1 sg. geald; pret. pl. guldon; past ptcp. golden. yield.\par
% glossary-block:B00693:end
% glossary-block:B00694:begin
\needspace{3\baselineskip}
\textbf{gīeman} (v. weak); pret. 1 sg. gīemde; past ptcp. gīemed. rule over. Genitive of what is ruled; no additional over preposition.\par
% glossary-block:B00694:end
% glossary-block:B00695:begin
\needspace{3\baselineskip}
\textbf{giernan} (v. weak); pret. 1 sg. giernde; past ptcp. gierned. desire; request; woo. Genitive desired thing/person in the selected frames, including the maiden wooed.\par
% glossary-block:B00695:end
% glossary-block:B00696:begin
\needspace{3\baselineskip}
\textbf{gierwan} (v. weak); pret. 1 sg. gierede; past ptcp. (ge)gier(w)ed. prepare. Accusative thing prepared.\par
% glossary-block:B00696:end
% glossary-block:B00697:begin
\needspace{3\baselineskip}
\textbf{giesthūs} (n. n.); gen. sg. giesthūses; nom. pl. giesthūs. inn; guest-house.\par
% glossary-block:B00697:end
% glossary-block:B00698:begin
\needspace{3\baselineskip}
\textbf{gīet} (adv.). still; further.\par
% glossary-block:B00698:end
% glossary-block:B00699:begin
\needspace{3\baselineskip}
\textbf{gif} (conj.). if. Mood is determined by the open, tentative or counterfactual condition, not by gif alone.\par
% glossary-block:B00699:end
% glossary-block:B00700:begin
\needspace{3\baselineskip}
\textbf{gimm} (n. m.); gen. sg. gimmes; nom. pl. gimmas. gem; precious stone.\par
% glossary-block:B00700:end
% glossary-block:B00701:begin
\needspace{3\baselineskip}
\textbf{gingstan} — inflected superlative; supplied source form: see \textbf{geong}.\par
% glossary-block:B00701:end
% glossary-block:B00702:begin
\needspace{3\baselineskip}
\textbf{git} — nom. dual: see \textbf{ðū}.\par
% glossary-block:B00702:end
% glossary-block:B00703:begin
\needspace{3\baselineskip}
\textbf{glæd} (adj.). glad. gracious; favourably disposed. With wið + person in the supplied gracious-toward frame.\par
% glossary-block:B00703:end
% glossary-block:B00704:begin
\needspace{3\baselineskip}
\textbf{God} (n. m.); gen. sg. Godes. God. Proper divine name; no productive plural sense is assigned. god; deity. Common deity sense, distinct from the proper divine name.\par
% glossary-block:B00704:end
% glossary-block:B00705:begin
\needspace{3\baselineskip}
\textbf{gōd} (adj.); comp. bętera; superl. bętst. good. kind; benevolent.\par
% glossary-block:B00705:end
% glossary-block:B00706:begin
\needspace{3\baselineskip}
\textbf{gōd} (n. n.); gen. sg. gōdes; nom. pl. gōd. goods; property. Plural lexical noun. good; benefit; kindness. Nominal good in the repayment contrast.\par
% glossary-block:B00706:end
% glossary-block:B00707:begin
\needspace{3\baselineskip}
\textbf{godspell} (n. n.); gen. sg. godspelles; nom. pl. godspell. Gospel.\par
% glossary-block:B00707:end
% glossary-block:B00708:begin
\needspace{3\baselineskip}
\textbf{gōdwębb} (n. n.); gen. sg. gōdwębbes; nom. pl. gōdwębb. purple cloth.\par
% glossary-block:B00708:end
% glossary-block:B00709:begin
\needspace{3\baselineskip}
\textbf{gold} (n. n.); gen. sg. goldes; nom. pl. gold. gold.\par
% glossary-block:B00709:end
% glossary-block:B00710:begin
\needspace{3\baselineskip}
\textbf{goldhord} (n. n.); gen. sg. goldhordes; nom. pl. goldhord. treasure; hoard of gold.\par
% glossary-block:B00710:end
% glossary-block:B00711:begin
\needspace{3\baselineskip}
\textbf{grēne} (adj.). green. Reference in a source or explanatory note.\par
% glossary-block:B00711:end
% glossary-block:B00712:begin
\needspace{3\baselineskip}
\textbf{grētan} (v. weak); pret. 1 sg. grētte; past ptcp. grētt. greet; take leave of. Accusative person/place; greeting and farewell depend on context.\par
% glossary-block:B00712:end
% glossary-block:B00713:begin
\needspace{3\baselineskip}
\textbf{grōwan} (v.). grow. Reference in a source or explanatory note.\par
% glossary-block:B00713:end
% glossary-block:B00714:begin
\needspace{3\baselineskip}
\textbf{grundweall} (n. m.); gen. sg. grundwealles; nom. pl. grundweallas. foundation.\par
% glossary-block:B00714:end
% glossary-block:B00715:begin
\needspace{3\baselineskip}
\textbf{guma} (n. m.). man. poetic/source quotation.\par
% glossary-block:B00715:end
% glossary-block:B00716:begin
\needspace{3\baselineskip}
\textbf{gylt} (n. m.); gen. sg. gyltes; nom. pl. gyltas. debt. Includes the prayer's moral debt.\par
% glossary-block:B00716:end
% glossary-block:B00717:begin
\needspace{3\baselineskip}
\textbf{gyltend} (n. m.); gen. sg. gyltendes; nom. pl. gyltende. debtor.\par
% glossary-block:B00717:end
% book-navigation: standalone bookmark replaced by volume hierarchy
% glossary-block:B00718:begin
\needspace{5\baselineskip}
\subsection*{H}\BookHeading{nav-glossary-old-english-0010}{subsubsection}{H}{0}
% glossary-block:B00718:end
% glossary-block:B00719:begin
\textbf{habban} (v. weak); pret. 1 sg. hæfde; past ptcp. (ge)hæfd; pres. 3 sg. hæfð. have; possess. Accusative possession; distinguish the perfect auxiliary. have. Auxiliary with transitive past participle for the taught completed-action/result construction.\par
% glossary-block:B00719:end
% glossary-block:B00720:begin
\needspace{3\baselineskip}
\textbf{hāl} (adj.). hale; well. Supplied greeting. hale; healthy.\par
% glossary-block:B00720:end
% glossary-block:B00721:begin
\needspace{3\baselineskip}
\textbf{hālga} (n. m.); gen. sg. hālgan; nom. pl. hālgan. saint; holy one.\par
% glossary-block:B00721:end
% glossary-block:B00722:begin
\needspace{3\baselineskip}
\textbf{hālig} (adj.). holy. holy ones; saints. Substantival adjective in the source. Before vowel endings hālg-; retain hālig- before -ne, -re and -ra.\par
% glossary-block:B00722:end
% glossary-block:B00723:begin
\needspace{3\baselineskip}
\textbf{hālsian} (v. weak); pret. 1 sg. hālsode; past ptcp. hālsod. entreat; implore. conjure; solemnly charge. Accusative person; ðurh + accusative oath standard; a requested ðæt clause may follow.\par
% glossary-block:B00723:end
% glossary-block:B00724:begin
\needspace{3\baselineskip}
\textbf{hām} (adv.). home; homeward. Directional adverb; distinct from the noun home.\par
% glossary-block:B00724:end
% glossary-block:B00725:begin
\needspace{3\baselineskip}
\textbf{hām} (n. m.); gen. sg. hāmes; nom. pl. hāmas. home; dwelling.\par
% glossary-block:B00725:end
% glossary-block:B00726:begin
\needspace{3\baselineskip}
\textbf{hāmweard} (adverb or particle). homeward.\par
% glossary-block:B00726:end
% glossary-block:B00727:begin
\needspace{3\baselineskip}
\textbf{hand} (n. f.); gen. sg. hande; nom. pl. handa. hand; side. Side in right-/left-hand expressions.\par
% glossary-block:B00727:end
% glossary-block:B00728:begin
\needspace{3\baselineskip}
\textbf{hāt} (adj.). hot.\par
% glossary-block:B00728:end
% glossary-block:B00729:begin
\needspace{3\baselineskip}
\textbf{hātan} (v. strong); pret. 1 sg. hēt; pret. pl. hēton; past ptcp. hāten; pres. 3 sg. hǣtt. call; name. command; order. Accusative actor + plain infinitive for the commanded action; distinct from call/name. be called; be named. Special hātte form in the assigned naming expression.\par
% glossary-block:B00729:end
% glossary-block:B00730:begin
\needspace{3\baselineskip}
\textbf{hatian} (v. weak); pret. 1 sg. hatode; past ptcp. hatod. hate.\par
% glossary-block:B00730:end
% glossary-block:B00731:begin
\needspace{3\baselineskip}
\textbf{hæfde} — pret. sg.: see \textbf{habban}.\par
% glossary-block:B00731:end
% glossary-block:B00732:begin
\needspace{3\baselineskip}
\textbf{hæfð} — pres. ind. 3 sg.: see \textbf{habban}.\par
% glossary-block:B00732:end
% glossary-block:B00733:begin
\needspace{3\baselineskip}
\textbf{hægel} (n. m.). hail.\par
% glossary-block:B00733:end
% glossary-block:B00734:begin
\needspace{3\baselineskip}
\textbf{hǣl} (n. n.). hail; greeting (in hǣl ābēodan, bid someone hail). greeting formula; him is the dative recipient.\par
% glossary-block:B00734:end
% glossary-block:B00735:begin
\needspace{3\baselineskip}
\textbf{hǣlan} (v. weak); pret. 1 sg. hǣlde; past ptcp. hǣled. save.\par
% glossary-block:B00735:end
% glossary-block:B00736:begin
\needspace{3\baselineskip}
\textbf{Hǣlend} (n. m.); gen. sg. Hǣlendes. Saviour; Jesus. Proper title.\par
% glossary-block:B00736:end
% glossary-block:B00737:begin
\needspace{3\baselineskip}
\textbf{hǣlðu} (n.). health. Pronunciation example.\par
% glossary-block:B00737:end
% glossary-block:B00738:begin
\needspace{3\baselineskip}
\textbf{hǣlu} (n. f.); gen. sg. hǣlu. safety; rescue. salvation.\par
% glossary-block:B00738:end
% glossary-block:B00739:begin
\needspace{3\baselineskip}
\textbf{hǣr} (n. n.); gen. sg. hǣres; nom. pl. hǣr. hair; a hair.\par
% glossary-block:B00739:end
% glossary-block:B00740:begin
\needspace{3\baselineskip}
\textbf{hærfest} (n. m.); gen. sg. hærfestes; nom. pl. hærfestas. harvest.\par
% glossary-block:B00740:end
% glossary-block:B00741:begin
\needspace{3\baselineskip}
\textbf{hǣs} (n. f.). command; order.\par
% glossary-block:B00741:end
% glossary-block:B00742:begin
\needspace{3\baselineskip}
\textbf{hǣt} — pres. ind. 3 sg., variant: see \textbf{hātan}.\par
% glossary-block:B00742:end
% glossary-block:B00743:begin
\needspace{3\baselineskip}
\textbf{hǣtt} — pres. ind. 3 sg.: see \textbf{hātan}.\par
% glossary-block:B00743:end
% glossary-block:B00744:begin
\needspace{3\baselineskip}
\textbf{hǣtu} (n. f.); gen. sg. hǣte. heat.\par
% glossary-block:B00744:end
% glossary-block:B00745:begin
\needspace{3\baselineskip}
\textbf{hǣðen} (adj.). heathen people.\par
% glossary-block:B00745:end
% glossary-block:B00746:begin
\needspace{3\baselineskip}
\textbf{hē} (personal pron.). he; she; it; they. Personal pronoun; gender/number and governed case determine the selected form and reference. his; her; its; their. Invariant genitive possessors, unlike agreeing possessive adjectives. their. Source genitive possessor.\par
% glossary-block:B00746:end
% glossary-block:B00747:begin
\needspace{3\baselineskip}
\textbf{hēaclif} (n. n.). high cliff.\par
% glossary-block:B00747:end
% glossary-block:B00748:begin
\needspace{3\baselineskip}
\textbf{hēafod} (n. n.); gen. sg. hēafdes; nom. pl. hēafdu. head.\par
% glossary-block:B00748:end
% glossary-block:B00749:begin
\needspace{3\baselineskip}
\textbf{hēah} (adj.); supplied superl. gen. m. sg. Hēahstan. high. high in rank; exalted. Rank/status sense in the supplied source. the Highest. Divine title. Contracted forms include hēan and hēam; see the Unit 3 pattern.\par
% glossary-block:B00749:end
% glossary-block:B00750:begin
\needspace{3\baselineskip}
\textbf{hēah} (adv.). high. Adverbial height, distinct from an agreeing adjective.\par
% glossary-block:B00750:end
% glossary-block:B00751:begin
\needspace{3\baselineskip}
\textbf{hēahstan} — inflected superlative; supplied source form: see \textbf{hēah}.\par
% glossary-block:B00751:end
% glossary-block:B00752:begin
\needspace{3\baselineskip}
\textbf{healdan} (v. strong); pret. 1 sg. hēold; pret. pl. hēoldon; past ptcp. healden; pres. 3 sg. hielt. hold. keep; observe. Accusative saying or command observed. maintain; sustain.\par
% glossary-block:B00752:end
% glossary-block:B00753:begin
\needspace{3\baselineskip}
\textbf{healf} (adj.). half.\par
% glossary-block:B00753:end
% glossary-block:B00754:begin
\needspace{3\baselineskip}
\textbf{heall} (n. f.). hall.\par
% glossary-block:B00754:end
% glossary-block:B00755:begin
\needspace{3\baselineskip}
\textbf{healt} (adj.). lame; halt. Adjective describing impaired walking; not the command halt/stop.\par
% glossary-block:B00755:end
% glossary-block:B00756:begin
\needspace{3\baselineskip}
\textbf{hēanes} (n. f.); gen. sg. hēanesse; nom. pl. hēanessa. height; highest point. heights; highest places.\par
% glossary-block:B00756:end
% glossary-block:B00757:begin
\needspace{3\baselineskip}
\textbf{hearpe} (n. f.); gen. sg. hearpan; nom. pl. hearpan. harp.\par
% glossary-block:B00757:end
% glossary-block:B00758:begin
\needspace{3\baselineskip}
\textbf{hearpenægl} (n. m.). plectrum; harp pick.\par
% glossary-block:B00758:end
% glossary-block:B00759:begin
\needspace{3\baselineskip}
\textbf{hearpere} (n. m.); gen. sg. hearperes; nom. pl. hearperas. harper.\par
% glossary-block:B00759:end
% glossary-block:B00760:begin
\needspace{3\baselineskip}
\textbf{hearpestręng} (n. m.). harp strings.\par
% glossary-block:B00760:end
% glossary-block:B00761:begin
\needspace{3\baselineskip}
\textbf{hearpian} (v.). harp. Supplied genitive harper/harping construction, with instrument; do not reconstruct it from the infinitive alone. play harp.\par
% glossary-block:B00761:end
% glossary-block:B00762:begin
\needspace{3\baselineskip}
\textbf{hēawan} (v. strong); pret. 1 sg. hēow; pret. pl. hēowon; past ptcp. hēawen; pres. 3 sg. hīewð. hew.\par
% glossary-block:B00762:end
% glossary-block:B00763:begin
\needspace{3\baselineskip}
\textbf{hębban} (v. strong); pret. 1 sg. hōf; pret. pl. hōfon; past ptcp. hafen; pres. 3 sg. hęfð. raise; lift.\par
% glossary-block:B00763:end
% glossary-block:B00764:begin
\needspace{3\baselineskip}
\textbf{hęfig} (adj.). grievous.\par
% glossary-block:B00764:end
% glossary-block:B00765:begin
\needspace{3\baselineskip}
\textbf{hęfignes} (n. f.); gen. sg. hęfignesse; nom. pl. hęfignessa. burden.\par
% glossary-block:B00765:end
% glossary-block:B00766:begin
\needspace{3\baselineskip}
\textbf{hęfð} — pres. ind. 3 sg.: see \textbf{hębban}.\par
% glossary-block:B00766:end
% glossary-block:B00767:begin
\needspace{3\baselineskip}
\textbf{hēht} — pret. sg., older variant: see \textbf{hātan}.\par
% glossary-block:B00767:end
% glossary-block:B00768:begin
\needspace{3\baselineskip}
\textbf{hęll} (n. f.); gen. sg. hęlle; nom. pl. hęlla. hell.\par
% glossary-block:B00768:end
% glossary-block:B00769:begin
\needspace{3\baselineskip}
\textbf{helpan} (v. strong); pret. 1 sg. healp; pret. pl. hulpon; past ptcp. holpen; pres. 3 sg. hilpð. help. Dative person helped.\par
% glossary-block:B00769:end
% glossary-block:B00770:begin
\needspace{3\baselineskip}
\textbf{hēo} — nom. f. sg.; also acc. f. sg. and nom./acc. pl., variant: see \textbf{hē}.\par
% glossary-block:B00770:end
% glossary-block:B00771:begin
\needspace{3\baselineskip}
\textbf{heofon} (n. m.); gen. sg. heofones; nom. pl. heofonas. heaven; sky.\par
% glossary-block:B00771:end
% glossary-block:B00772:begin
\needspace{3\baselineskip}
\textbf{heofone} (n. f.). heaven; sky. Weak feminine lexical alternative in the source.\par
% glossary-block:B00772:end
% glossary-block:B00773:begin
\needspace{3\baselineskip}
\textbf{heofonfȳr} (n. n.). fire from heaven; heavenly fire.\par
% glossary-block:B00773:end
% glossary-block:B00774:begin
\needspace{3\baselineskip}
\textbf{heofonlic} (adj.). heavenly.\par
% glossary-block:B00774:end
% glossary-block:B00775:begin
\needspace{3\baselineskip}
\textbf{heofonlice} (adv.). from heaven; heavenly. Source expression at Cook 143.17.\par
% glossary-block:B00775:end
% glossary-block:B00776:begin
\needspace{3\baselineskip}
\textbf{heonan} (adv.). hence; from here.\par
% glossary-block:B00776:end
% glossary-block:B00777:begin
\needspace{3\baselineskip}
\textbf{heonanforð} (adv.). henceforth.\par
% glossary-block:B00777:end
% glossary-block:B00778:begin
\needspace{3\baselineskip}
\textbf{heora} — gen. pl., variant: see \textbf{hē}.\par
% glossary-block:B00778:end
% glossary-block:B00779:begin
\needspace{3\baselineskip}
\textbf{heorte} (n. f.); gen. sg. heortan; nom. pl. heortan. heart.\par
% glossary-block:B00779:end
% glossary-block:B00780:begin
\needspace{3\baselineskip}
\textbf{hēr} (adv.). here.\par
% glossary-block:B00780:end
% glossary-block:B00781:begin
\needspace{3\baselineskip}
\textbf{hęre} (n. m.); gen. sg. hęres; nom. pl. hęras. army.\par
% glossary-block:B00781:end
% glossary-block:B00782:begin
\needspace{3\baselineskip}
\textbf{hęrgian} (v.). plunder; ravage.\par
% glossary-block:B00782:end
% glossary-block:B00783:begin
\needspace{3\baselineskip}
\textbf{hęrian} (v. weak); pret. 1 sg. hęrede; past ptcp. hęred. praise. Genitive plural present-participle use, of people praising. praise. Supplied source form praised. praise.\par
% glossary-block:B00783:end
% glossary-block:B00784:begin
\needspace{3\baselineskip}
\textbf{hī} — nom./acc. f. sg.; nom./acc. pl., variant: see \textbf{hē}.\par
% glossary-block:B00784:end
% glossary-block:B00785:begin
\needspace{3\baselineskip}
\textbf{hider} (adv.). hither; here. Motion toward the speaker, distinct from location hēr.\par
% glossary-block:B00785:end
% glossary-block:B00786:begin
\needspace{3\baselineskip}
\textbf{hīe} — acc. f sg.; nom. pl.; acc. pl.: see \textbf{hē}.\par
% glossary-block:B00786:end
% glossary-block:B00787:begin
\needspace{3\baselineskip}
\textbf{hielt} — pres. ind. 3 sg.: see \textbf{healdan}.\par
% glossary-block:B00787:end
% glossary-block:B00788:begin
\needspace{3\baselineskip}
\textbf{hīenan} (v.). oppress; humble.\par
% glossary-block:B00788:end
% glossary-block:B00789:begin
\needspace{3\baselineskip}
\textbf{hiene} — acc. m. sg., variant: see \textbf{hē}.\par
% glossary-block:B00789:end
% glossary-block:B00790:begin
\needspace{3\baselineskip}
\textbf{hīera} — gen. pl., source spelling: see \textbf{hē}.\par
% glossary-block:B00790:end
% glossary-block:B00791:begin
\needspace{3\baselineskip}
\textbf{hiera} — gen. pl.: see \textbf{hē}.\par
% glossary-block:B00791:end
% glossary-block:B00792:begin
\needspace{3\baselineskip}
\textbf{hīeran} (v. weak); pret. 1 sg. hīerde; past ptcp. (ge)hīered. hear. hear an action. Accusative perceived actor plus plain infinitive in the taught perception construction. obey. Dative person obeyed.\par
% glossary-block:B00792:end
% glossary-block:B00793:begin
\needspace{3\baselineskip}
\textbf{hiere} — gen. f sg.; dat. f sg.: see \textbf{hē}.\par
% glossary-block:B00793:end
% glossary-block:B00794:begin
\needspace{3\baselineskip}
\textbf{hīernes} (n. f.). obedience.\par
% glossary-block:B00794:end
% glossary-block:B00795:begin
\needspace{3\baselineskip}
\textbf{Hierusalem} (proper name). Hierusalem.\par
% glossary-block:B00795:end
% glossary-block:B00796:begin
\needspace{3\baselineskip}
\textbf{hīewð} — pres. ind. 3 sg.: see \textbf{hēawan}.\par
% glossary-block:B00796:end
% glossary-block:B00797:begin
\needspace{3\baselineskip}
\textbf{hilpð} — pres. ind. 3 sg.: see \textbf{helpan}.\par
% glossary-block:B00797:end
% glossary-block:B00798:begin
\needspace{3\baselineskip}
\textbf{him} — dat. m sg.; dat. n sg.; dat. pl.: see \textbf{hē}.\par
% glossary-block:B00798:end
% glossary-block:B00799:begin
\needspace{3\baselineskip}
\textbf{hine} — acc. m sg.: see \textbf{hē}.\par
% glossary-block:B00799:end
% glossary-block:B00800:begin
\needspace{3\baselineskip}
\textbf{hīo} — nom. f. sg., variant: see \textbf{hē}.\par
% glossary-block:B00800:end
% glossary-block:B00801:begin
\needspace{3\baselineskip}
\textbf{hira} — gen. pl., variant: see \textbf{hē}.\par
% glossary-block:B00801:end
% glossary-block:B00802:begin
\needspace{3\baselineskip}
\textbf{hire} — gen./dat. f. sg., variant: see \textbf{hē}.\par
% glossary-block:B00802:end
% glossary-block:B00803:begin
\needspace{3\baselineskip}
\textbf{his} — gen. m sg.; gen. n sg.: see \textbf{hē}.\par
% glossary-block:B00803:end
% glossary-block:B00804:begin
\needspace{3\baselineskip}
\textbf{hit} — nom. n sg.; acc. n sg.: see \textbf{hē}.\par
% glossary-block:B00804:end
% glossary-block:B00805:begin
\needspace{3\baselineskip}
\textbf{hīw} (n. n.). form; kind. Supplied dative neuter plural.\par
% glossary-block:B00805:end
% glossary-block:B00806:begin
\needspace{3\baselineskip}
\textbf{hlāf} (n. m.); gen. sg. hlāfes. bread; loaf.\par
% glossary-block:B00806:end
% glossary-block:B00807:begin
\needspace{3\baselineskip}
\textbf{hlāford} (n. m.); gen. sg. hlāfordes; nom. pl. hlāfordas. lord; master.\par
% glossary-block:B00807:end
% glossary-block:B00808:begin
\needspace{3\baselineskip}
\textbf{hlǣfdige} (n. f.); gen. sg. hlǣfdigan; nom. pl. hlǣfdigan. lady.\par
% glossary-block:B00808:end
% glossary-block:B00809:begin
\needspace{3\baselineskip}
\textbf{hlæst} (n. n.); gen. sg. hlæstes; nom. pl. hlæst. wares; merchandise; cargo. Selected plural source use.\par
% glossary-block:B00809:end
% glossary-block:B00810:begin
\needspace{3\baselineskip}
\textbf{hlīfian} (v.); pres. 3 sg. hlīfað. tower. Verb meaning stand/tower high in the source; distinct from noun tower.\par
% glossary-block:B00810:end
% glossary-block:B00811:begin
\needspace{3\baselineskip}
\textbf{hlīsa} (n. m.); gen. sg. hlīsan; nom. pl. hlīsan. fame. report; account heard.\par
% glossary-block:B00811:end
% glossary-block:B00812:begin
\needspace{3\baselineskip}
\textbf{hlūd} (adj.). loud.\par
% glossary-block:B00812:end
% glossary-block:B00813:begin
\needspace{3\baselineskip}
\textbf{hlystan} (v.). listen to; heed. Genitive person or thing listened to; hear is the contextual prompt wording.\par
% glossary-block:B00813:end
% glossary-block:B00814:begin
\needspace{3\baselineskip}
\textbf{hof} (n. n.); gen. sg. hofes; nom. pl. hofu. dwelling.\par
% glossary-block:B00814:end
% glossary-block:B00815:begin
\needspace{3\baselineskip}
\textbf{hogað} — pres. ind. 3 sg., variant: see \textbf{hycgan}.\par
% glossary-block:B00815:end
% glossary-block:B00816:begin
\needspace{3\baselineskip}
\textbf{hogde} — pret. sg., variant: see \textbf{hycgan}.\par
% glossary-block:B00816:end
% glossary-block:B00817:begin
\needspace{3\baselineskip}
\textbf{hogode} — pret. sg.: see \textbf{hycgan}.\par
% glossary-block:B00817:end
% glossary-block:B00818:begin
\needspace{3\baselineskip}
\textbf{hol} (n. n.); gen. sg. holes; nom. pl. holu. hole.\par
% glossary-block:B00818:end
% glossary-block:B00819:begin
\needspace{3\baselineskip}
\textbf{hōlunga} (adv.). in vain.\par
% glossary-block:B00819:end
% glossary-block:B00820:begin
\needspace{3\baselineskip}
\textbf{hord} (n. n/m.). hoard; treasure.\par
% glossary-block:B00820:end
% glossary-block:B00821:begin
\needspace{3\baselineskip}
\textbf{horg} (adj.). dirty; filthy.\par
% glossary-block:B00821:end
% glossary-block:B00822:begin
\needspace{3\baselineskip}
\textbf{horn} (n. m.); gen. sg. hornes; nom. pl. hornas. horn; strength; dignity. Strength/dignity is the figurative force in 37.14, not a replacement for the literal source horn.\par
% glossary-block:B00822:end
% glossary-block:B00823:begin
\needspace{3\baselineskip}
\textbf{hors} (n. n.); gen. sg. horses; nom. pl. hors. horse.\par
% glossary-block:B00823:end
% glossary-block:B00824:begin
\needspace{3\baselineskip}
\textbf{hraðe} (adv.). quickly.\par
% glossary-block:B00824:end
% glossary-block:B00825:begin
\needspace{3\baselineskip}
\textbf{hrædlīce} (adv.). quickly.\par
% glossary-block:B00825:end
% glossary-block:B00826:begin
\needspace{3\baselineskip}
\textbf{hrædnes} (n. f.). speed; quickness.\par
% glossary-block:B00826:end
% glossary-block:B00827:begin
\needspace{3\baselineskip}
\textbf{hrēod} (n. n.); gen. sg. hrēodes; nom. pl. hrēod. reed.\par
% glossary-block:B00827:end
% glossary-block:B00828:begin
\needspace{3\baselineskip}
\textbf{hrēohnes} (n.). tempest; storm.\par
% glossary-block:B00828:end
% glossary-block:B00829:begin
\needspace{3\baselineskip}
\textbf{hrēosan} (v.); pret. pl. hruron. fall. Supplied present participle falling.\par
% glossary-block:B00829:end
% glossary-block:B00830:begin
\needspace{3\baselineskip}
\textbf{hrōf} (n. m.); gen. sg. hrōfes; nom. pl. hrōfas. roof.\par
% glossary-block:B00830:end
% glossary-block:B00831:begin
\needspace{3\baselineskip}
\textbf{hū} (adv.). how?.\par
% glossary-block:B00831:end
% glossary-block:B00832:begin
\needspace{3\baselineskip}
\textbf{hūlic} (interrogative adjective). of what sort?.\par
% glossary-block:B00832:end
% glossary-block:B00833:begin
\needspace{3\baselineskip}
\textbf{hund} (cardinal). hundred. Only the complete supplied hundred-sheep expression at its first appearance; later numerical frames remain locally specified.\par
% glossary-block:B00833:end
% glossary-block:B00834:begin
\needspace{3\baselineskip}
\textbf{hund} (n. m.); gen. sg. hundes; nom. pl. hundas. dog.\par
% glossary-block:B00834:end
% glossary-block:B00835:begin
\needspace{3\baselineskip}
\textbf{hundnigontig} (numeral). ninety.\par
% glossary-block:B00835:end
% glossary-block:B00836:begin
\needspace{3\baselineskip}
\textbf{hundtēontig} (numeral). a hundredfold. A hundred in return, in the supplied recompense expression. hundred. Within the complete supplied 144,000 expression; do not infer an independent high-number system.\par
% glossary-block:B00836:end
% glossary-block:B00837:begin
\needspace{3\baselineskip}
\textbf{hunger} (n. m.); gen. sg. hungres. famine; hunger.\par
% glossary-block:B00837:end
% glossary-block:B00838:begin
\needspace{3\baselineskip}
\textbf{hungrig} (adj.). hungry. hungry people. Substantival adjective.\par
% glossary-block:B00838:end
% glossary-block:B00839:begin
\needspace{3\baselineskip}
\textbf{hunig} (n. n.); gen. sg. huniges. honey.\par
% glossary-block:B00839:end
% glossary-block:B00840:begin
\needspace{3\baselineskip}
\textbf{hūs} (n. n.); gen. sg. hūses; nom. pl. hūs. house.\par
% glossary-block:B00840:end
% glossary-block:B00841:begin
\needspace{3\baselineskip}
\textbf{hwā, hwæt} (interrogative/indefinite pron.). who? what?. Case follows role/government; distinguish personal and neuter reference. anyone. what reason. Instrumental hwon in the supplied why expressions.\par
% glossary-block:B00841:end
% glossary-block:B00842:begin
\needspace{3\baselineskip}
\textbf{hwām} — dat. m./f. sg.; dat. n. sg.: see \textbf{hwā, hwæt}.\par
% glossary-block:B00842:end
% glossary-block:B00843:begin
\needspace{3\baselineskip}
\textbf{hwanon} (adv.). whence; from where.\par
% glossary-block:B00843:end
% glossary-block:B00844:begin
\needspace{3\baselineskip}
\textbf{hwæl} (n. m.); gen. sg. hwæles; nom. pl. hwalas. whale.\par
% glossary-block:B00844:end
% glossary-block:B00845:begin
\needspace{3\baselineskip}
\textbf{hwǣm} — dat. m./f. sg.; dat. n. sg.: see \textbf{hwā, hwæt}.\par
% glossary-block:B00845:end
% glossary-block:B00846:begin
\needspace{3\baselineskip}
\textbf{hwǣr} (adv.). where.\par
% glossary-block:B00846:end
% glossary-block:B00847:begin
\needspace{3\baselineskip}
\textbf{hwæs} — gen. m./f. sg.; gen. n. sg.: see \textbf{hwā, hwæt}.\par
% glossary-block:B00847:end
% glossary-block:B00848:begin
\needspace{3\baselineskip}
\textbf{hwæt} — nom. n. sg.; acc. n. sg.: see \textbf{hwā, hwæt}.\par
% glossary-block:B00848:end
% glossary-block:B00849:begin
\needspace{3\baselineskip}
\textbf{hwǣte} (n. m.); gen. sg. hwǣtes. wheat.\par
% glossary-block:B00849:end
% glossary-block:B00850:begin
\needspace{3\baselineskip}
\textbf{hwæthwega} (indefinite pronoun/substantive). something. Supplied local indefinite form; no unrestricted new paradigm.\par
% glossary-block:B00850:end
% glossary-block:B00851:begin
\needspace{3\baselineskip}
\textbf{hwæthwugu} (indefinite). something. Supplied hwæthwugu swilces.\par
% glossary-block:B00851:end
% glossary-block:B00852:begin
\needspace{3\baselineskip}
\textbf{hwæðer} (interrogative pronoun/adjective). which of two?. whichever of two. In swā hwæðer swā; the supplied attraction example has a special case arrangement.\par
% glossary-block:B00852:end
% glossary-block:B00853:begin
\needspace{3\baselineskip}
\textbf{hwæðer} (conj.). whether.\par
% glossary-block:B00853:end
% glossary-block:B00854:begin
\needspace{3\baselineskip}
\textbf{hwelc} — citation/spelling variant: see \textbf{hwilc}.\par
% glossary-block:B00854:end
% glossary-block:B00855:begin
\needspace{3\baselineskip}
\textbf{hweorfan} (v.). return; turn back. Supplied source preterite.\par
% glossary-block:B00855:end
% glossary-block:B00856:begin
\needspace{3\baselineskip}
\textbf{hwider} (adv.). whither; to where.\par
% glossary-block:B00856:end
% glossary-block:B00857:begin
\needspace{3\baselineskip}
\textbf{hwilc} (interrogative/indefinite adj. or pron.). which? of what kind?. which. Selection from a group expressed in genitive. any; some. Indefinite adjective, distinct from an interrogative.\par
% glossary-block:B00857:end
% glossary-block:B00858:begin
\needspace{3\baselineskip}
\textbf{hwīlum} (adv.). sometimes. a while ago. Only in supplied hwīlum ǣr; earlier hwīlum sometimes is not enough.\par
% glossary-block:B00858:end
% glossary-block:B00859:begin
\needspace{3\baselineskip}
\textbf{hwīt} (adj.). white.\par
% glossary-block:B00859:end
% glossary-block:B00860:begin
\needspace{3\baselineskip}
\textbf{hwon} — instr. n. sg.: see \textbf{hwā, hwæt}.\par
% glossary-block:B00860:end
% glossary-block:B00861:begin
\needspace{3\baselineskip}
\textbf{hwone} — acc. m./f. sg.: see \textbf{hwā, hwæt}.\par
% glossary-block:B00861:end
% glossary-block:B00862:begin
\needspace{3\baselineskip}
\textbf{hwonne} (adv.). when?. Direct interrogative.\par
% glossary-block:B00862:end
% glossary-block:B00863:begin
\needspace{3\baselineskip}
\textbf{hwȳ} (adv.). why?. Also instr. n. sg.: see \textbf{hwā, hwæt}.\par
% glossary-block:B00863:end
% glossary-block:B00864:begin
\needspace{3\baselineskip}
\textbf{hwylc} — citation/spelling variant: see \textbf{hwilc}.\par
% glossary-block:B00864:end
% glossary-block:B00865:begin
\needspace{3\baselineskip}
\textbf{hycgan} (v. weak); pret. 1 sg. hogode; past ptcp. (ge)hogod; pres. 3 sg. hygð. think.\par
% glossary-block:B00865:end
% glossary-block:B00866:begin
\needspace{3\baselineskip}
\textbf{hȳdan} (v. weak); pret. 1 sg. hȳdde; past ptcp. hȳded. hide; conceal. Accusative thing hidden; wesan forms the assigned passive.\par
% glossary-block:B00866:end
% glossary-block:B00867:begin
\needspace{3\baselineskip}
\textbf{hygð} — pres. ind. 3 sg.: see \textbf{hycgan}.\par
% glossary-block:B00867:end
% glossary-block:B00868:begin
\needspace{3\baselineskip}
\textbf{hyll} (n. m.); gen. sg. hylles; nom. pl. hyllas. hill.\par
% glossary-block:B00868:end
% glossary-block:B00869:begin
\needspace{3\baselineskip}
\textbf{hyngran} (v. impersonal); pres. 3 sg. hyngreð; pret. 3 sg. hyngrede; past ptcp. (ge)hyngred. be hungry. Impersonal construction with accusative experiencer. hungry; hungering. Present participle used substantivally in the assigned frame. Impersonal use: the course supplies third-singular finite forms.\par
% glossary-block:B00869:end
% glossary-block:B00870:begin
\needspace{3\baselineskip}
\textbf{hȳran} — citation/spelling variant: see \textbf{hīeran}.\par
% glossary-block:B00870:end
% glossary-block:B00871:begin
\needspace{3\baselineskip}
\textbf{hȳrian} (v. weak); pret. 1 sg. hȳrode; past ptcp. hȳrod. hire.\par
% glossary-block:B00871:end
% book-navigation: standalone bookmark replaced by volume hierarchy
% glossary-block:B00872:begin
\needspace{5\baselineskip}
\subsection*{I}\BookHeading{nav-glossary-old-english-0011}{subsubsection}{I}{0}
% glossary-block:B00872:end
% glossary-block:B00873:begin
\textbf{Iacob} (proper name). Jacob.\par
% glossary-block:B00873:end
% glossary-block:B00874:begin
\needspace{3\baselineskip}
\textbf{ic} (personal pron.). I; me; we; us. Singular, dual and plural personal forms; select governed case. of me; of us. Personal genitive citation forms identifying person/number; this record alone does not teach an attributive possessive use. our. Personal genitive plural used attributively in the exact ūre hǣle expression. of us two. Personal genitive dual after a selection word.\par
% glossary-block:B00874:end
% glossary-block:B00875:begin
\needspace{3\baselineskip}
\textbf{īdel} (adj.). empty; void. idle. Distinct from empty/void.\par
% glossary-block:B00875:end
% glossary-block:B00876:begin
\needspace{3\baselineskip}
\textbf{īdelnes} (n. f.); gen. sg. īdelnesse; nom. pl. īdelnessa. idleness.\par
% glossary-block:B00876:end
% glossary-block:B00877:begin
\needspace{3\baselineskip}
\textbf{īecan} (v. weak); pret. 1 sg. īecte; past ptcp. (ge)īeced. increase.\par
% glossary-block:B00877:end
% glossary-block:B00878:begin
\needspace{3\baselineskip}
\textbf{ieldest} (adj.). oldest. eldest; oldest. Substantival person/eldest brother in this source context. Superlative of eald.\par
% glossary-block:B00878:end
% glossary-block:B00879:begin
\needspace{3\baselineskip}
\textbf{ieldra} (adj.). older. Comparative of eald.\par
% glossary-block:B00879:end
% glossary-block:B00880:begin
\needspace{3\baselineskip}
\textbf{ieldu} (n. f.). age; old age.\par
% glossary-block:B00880:end
% glossary-block:B00881:begin
\needspace{3\baselineskip}
\textbf{ierman} (v.). afflict.\par
% glossary-block:B00881:end
% glossary-block:B00882:begin
\needspace{3\baselineskip}
\textbf{iernan} (v. strong); pret. 1 sg. arn; pret. pl. urnon; past ptcp. urnen. run. Intransitive; geond + accusative path. A perception construction may make its actor accusative.\par
% glossary-block:B00882:end
% glossary-block:B00883:begin
\needspace{3\baselineskip}
\textbf{ierð} (n. f.). crops; produce of the earth.\par
% glossary-block:B00883:end
% glossary-block:B00884:begin
\needspace{3\baselineskip}
\textbf{ierðling} (n. m.); gen. sg. ierðlinges; nom. pl. ierðlingas. farmer; plowman.\par
% glossary-block:B00884:end
% glossary-block:B00885:begin
\needspace{3\baselineskip}
\textbf{īgland} (n. n.); gen. sg. īglandes; nom. pl. īgland. island.\par
% glossary-block:B00885:end
% glossary-block:B00886:begin
\needspace{3\baselineskip}
\textbf{ilca} (demonstrative adjective). same.\par
% glossary-block:B00886:end
% glossary-block:B00887:begin
\needspace{3\baselineskip}
\textbf{in} (prep.). in; on. With dative in the exact locations/circumstances. in; inward. Adverbial direction. toward; onto. Supplied phrase directing the mind toward these things.\par
% glossary-block:B00887:end
% glossary-block:B00888:begin
\needspace{3\baselineskip}
\textbf{inc} — dat. dual; acc. dual: see \textbf{ðū}.\par
% glossary-block:B00888:end
% glossary-block:B00889:begin
\needspace{3\baselineskip}
\textbf{incer} (possessive adjective). your; yours. Two possessors; agreeing possessive. Also gen. dual: see \textbf{ðū}.\par
% glossary-block:B00889:end
% glossary-block:B00890:begin
\needspace{3\baselineskip}
\textbf{inēode} — pret. sg.: see \textbf{ingān}.\par
% glossary-block:B00890:end
% glossary-block:B00891:begin
\needspace{3\baselineskip}
\textbf{ingān} (v. irregular); pres. 1 sg. ingā; pres. 3 sg. ingǣð; pres. pl. ingāð; pret. 1 sg. inēode; pret. pl. inēodon; past ptcp. ingān. enter; go in. enter; go in. Supplied past source form.\par
% glossary-block:B00891:end
% glossary-block:B00892:begin
\needspace{3\baselineskip}
\textbf{ingangan} (v.). enter; go in. Supplied present-participle phrase.\par
% glossary-block:B00892:end
% glossary-block:B00893:begin
\needspace{3\baselineskip}
\textbf{ingǣð} — pres. ind. 3 sg.: see \textbf{ingān}.\par
% glossary-block:B00893:end
% glossary-block:B00894:begin
\needspace{3\baselineskip}
\textbf{innan} (adv.). within; inside.\par
% glossary-block:B00894:end
% glossary-block:B00895:begin
\needspace{3\baselineskip}
\textbf{intinga} (n. m.). cause; sake (in for … intingan). Within supplied for intingan expressions.\par
% glossary-block:B00895:end
% glossary-block:B00896:begin
\needspace{3\baselineskip}
\textbf{in-tō} (prep.). into; destination expressed with dative after intō.\par
% glossary-block:B00896:end
% glossary-block:B00897:begin
\needspace{3\baselineskip}
\textbf{inweardlīce} (adv.). fervently; wholeheartedly.\par
% glossary-block:B00897:end
% glossary-block:B00898:begin
\needspace{3\baselineskip}
\textbf{Iosep} (proper name). Joseph.\par
% glossary-block:B00898:end
% glossary-block:B00899:begin
\needspace{3\baselineskip}
\textbf{īs} (n. n.). ice.\par
% glossary-block:B00899:end
% glossary-block:B00900:begin
\needspace{3\baselineskip}
\textbf{is} — pres. ind. 3 sg.: see \textbf{wesan}.\par
% glossary-block:B00900:end
% glossary-block:B00901:begin
\needspace{3\baselineskip}
\textbf{īsern} (n. n.); gen. sg. īsernes; nom. pl. īsern. iron. Metal noun.\par
% glossary-block:B00901:end
% glossary-block:B00902:begin
\needspace{3\baselineskip}
\textbf{īsern} (adj.). iron; made of iron. Adjective of material.\par
% glossary-block:B00902:end
% glossary-block:B00903:begin
\needspace{3\baselineskip}
\textbf{Israel} (proper name); gen. sg. Israeles. Israel.\par
% glossary-block:B00903:end
% glossary-block:B00904:begin
\needspace{3\baselineskip}
\textbf{Israela} (proper name). Israelites. Supplied genitive plural name form. Israela is the supplied genitive plural name-form, of the Israelites.\par
% glossary-block:B00904:end
% glossary-block:B00905:begin
\needspace{3\baselineskip}
\textbf{it} — pres. ind. 3 sg., variant: see \textbf{etan}.\par
% glossary-block:B00905:end
% glossary-block:B00906:begin
\needspace{3\baselineskip}
\textbf{itt} — pres. ind. 3 sg.: see \textbf{etan}.\par
% glossary-block:B00906:end
% glossary-block:B00907:begin
\needspace{3\baselineskip}
\textbf{Iudas} (proper name). Judah.\par
% glossary-block:B00907:end
% glossary-block:B00908:begin
\needspace{3\baselineskip}
\textbf{Iudea} (proper name). Jews. Supplied genitive plural name form. Iudea is the supplied genitive plural name-form, of the Jews.\par
% glossary-block:B00908:end
% book-navigation: standalone bookmark replaced by volume hierarchy
% glossary-block:B00909:begin
\needspace{5\baselineskip}
\subsection*{J}\BookHeading{nav-glossary-old-english-0012}{subsubsection}{J}{0}
% glossary-block:B00909:end
% glossary-block:B00910:begin
\textbf{Johannes} (proper name); gen. sg. Johannes. John.\par
% glossary-block:B00910:end
% book-navigation: standalone bookmark replaced by volume hierarchy
% glossary-block:B00911:begin
\needspace{5\baselineskip}
\subsection*{L}\BookHeading{nav-glossary-old-english-0013}{subsubsection}{L}{0}
% glossary-block:B00911:end
% glossary-block:B00912:begin
\textbf{lāc} (n. n.); gen. sg. lāces; nom. pl. lāc. gift; present.\par
% glossary-block:B00912:end
% glossary-block:B00913:begin
\needspace{3\baselineskip}
\textbf{land} (n. n.); gen. sg. landes; nom. pl. land. land; country.\par
% glossary-block:B00913:end
% glossary-block:B00914:begin
\needspace{3\baselineskip}
\textbf{landgemǣre} (n. n.). territory; borders.\par
% glossary-block:B00914:end
% glossary-block:B00915:begin
\needspace{3\baselineskip}
\textbf{lang} (adj.). long. Adjective agreeing with the time interval fæc.\par
% glossary-block:B00915:end
% glossary-block:B00916:begin
\needspace{3\baselineskip}
\textbf{lange} (adv.). long. In supplied swā lange swā; its biblical inasmuch-as force is explained locally.\par
% glossary-block:B00916:end
% glossary-block:B00917:begin
\needspace{3\baselineskip}
\textbf{lār} (n. f.); gen. sg. lāre. learning; study.\par
% glossary-block:B00917:end
% glossary-block:B00918:begin
\needspace{3\baselineskip}
\textbf{lārēow} (n. m.); gen. sg. lārēowes; nom. pl. lārēowas. master; teacher.\par
% glossary-block:B00918:end
% glossary-block:B00919:begin
\needspace{3\baselineskip}
\textbf{Lazarus} (proper name). Lazarus. The parable's poor man; supplied plural denotes people like him.\par
% glossary-block:B00919:end
% glossary-block:B00920:begin
\needspace{3\baselineskip}
\textbf{lǣdan} (v. weak); pret. 1 sg. lǣdde; past ptcp. (ge)lǣd(ed). lead. lead; bring; take. A person is the object in the assigned frame. carry; take. A person is conveyed in this supplied frame. take; carry along.\par
% glossary-block:B00920:end
% glossary-block:B00921:begin
\needspace{3\baselineskip}
\textbf{Læden} (n. n.). Latin. The language.\par
% glossary-block:B00921:end
% glossary-block:B00922:begin
\needspace{3\baselineskip}
\textbf{læfel} (n. m.); gen. sg. læfeles; nom. pl. læfelas. cup.\par
% glossary-block:B00922:end
% glossary-block:B00923:begin
\needspace{3\baselineskip}
\textbf{lǣran} (v. weak); pret. 1 sg. lǣrde; past ptcp. lǣred. advise; instruct. teach; instruct.\par
% glossary-block:B00923:end
% glossary-block:B00924:begin
\needspace{3\baselineskip}
\textbf{lǣs} (degree word). less; a smaller amount.\par
% glossary-block:B00924:end
% glossary-block:B00925:begin
\needspace{3\baselineskip}
\textbf{lǣssa} (adj.). smaller. Comparative of lȳtel.\par
% glossary-block:B00925:end
% glossary-block:B00926:begin
\needspace{3\baselineskip}
\textbf{lǣst} (adj.). smallest. least. In rank or importance. Superlative of lȳtel.\par
% glossary-block:B00926:end
% glossary-block:B00927:begin
\needspace{3\baselineskip}
\textbf{lǣtan} (v. strong); pret. 1 sg. lēt; pret. pl. lēton; past ptcp. lǣten. let.\par
% glossary-block:B00927:end
% glossary-block:B00928:begin
\needspace{3\baselineskip}
\textbf{lēad} (n. n.); gen. sg. lēades; nom. pl. lēad. lead. Metal noun, not the verb lead.\par
% glossary-block:B00928:end
% glossary-block:B00929:begin
\needspace{3\baselineskip}
\textbf{lēaf} (n. f.). permission.\par
% glossary-block:B00929:end
% glossary-block:B00930:begin
\needspace{3\baselineskip}
\textbf{lęcgean} (v. weak); pret. 1 sg. lęgde; past ptcp. (ge)lęgd. lay; place. Supplied source form. lay; place.\par
% glossary-block:B00930:end
% glossary-block:B00931:begin
\needspace{3\baselineskip}
\textbf{lēde} — pret. sg., contracted variant: see \textbf{lęcgean}.\par
% glossary-block:B00931:end
% glossary-block:B00932:begin
\needspace{3\baselineskip}
\textbf{lęng} (adverb or particle). longer.\par
% glossary-block:B00932:end
% glossary-block:B00933:begin
\needspace{3\baselineskip}
\textbf{lēo} (n.). lion (Latin loan-form).\par
% glossary-block:B00933:end
% glossary-block:B00934:begin
\needspace{3\baselineskip}
\textbf{lēof} (adj.); comp. lēofra; superl. lēofost. dear. sir. Supplied respectful address lēof.\par
% glossary-block:B00934:end
% glossary-block:B00935:begin
\needspace{3\baselineskip}
\textbf{leofað} — pres. ind. 3 sg.: see \textbf{libban}.\par
% glossary-block:B00935:end
% glossary-block:B00936:begin
\needspace{3\baselineskip}
\textbf{lēofost} (adj.). dearest. Superlative of lēof.\par
% glossary-block:B00936:end
% glossary-block:B00937:begin
\needspace{3\baselineskip}
\textbf{lēofra} (adj.). dearer. preferable; dearer. Dative experiencer + is lēofre + infinitive in the taught preference frame. Comparative of lēof.\par
% glossary-block:B00937:end
% glossary-block:B00938:begin
\needspace{3\baselineskip}
\textbf{lēoht} (n. n.); gen. sg. lēohtes; nom. pl. lēoht. light. light; luminary. Individual countable light, distinct from mass light.\par
% glossary-block:B00938:end
% glossary-block:B00939:begin
\needspace{3\baselineskip}
\textbf{leornian} (v. weak); pret. 1 sg. leornode; past ptcp. leornod. learn. Accusative thing learned; æt + dative person/source.\par
% glossary-block:B00939:end
% glossary-block:B00940:begin
\needspace{3\baselineskip}
\textbf{leornung} (n. f.); gen. sg. leornunge; nom. pl. leornunga. study; learning.\par
% glossary-block:B00940:end
% glossary-block:B00941:begin
\needspace{3\baselineskip}
\textbf{leornungcniht} (n. m.); gen. sg. leornungcnihtes; nom. pl. leornungcnihtas. disciple.\par
% glossary-block:B00941:end
% glossary-block:B00942:begin
\needspace{3\baselineskip}
\textbf{lēoð} (n. n.); gen. sg. lēoðes; nom. pl. lēoð. poem; poetry; verse.\par
% glossary-block:B00942:end
% glossary-block:B00943:begin
\needspace{3\baselineskip}
\textbf{lēoðwyrhta} (n. m.); gen. sg. lēoðwyrhtan; nom. pl. lēoðwyrhtan. poet.\par
% glossary-block:B00943:end
% glossary-block:B00944:begin
\needspace{3\baselineskip}
\textbf{libban} (v. weak); pret. 1 sg. lifde; past ptcp. (ge)lifd; pres. 3 sg. leofað. live. living. Supplied weak participial adjective.\par
% glossary-block:B00944:end
% glossary-block:B00945:begin
\needspace{3\baselineskip}
\textbf{liccian} (v. weak); pret. 1 sg. liccode; past ptcp. liccod. lick.\par
% glossary-block:B00945:end
% glossary-block:B00946:begin
\needspace{3\baselineskip}
\textbf{licgan} (v. strong); pret. 1 sg. læg; pret. pl. lǣgon; past ptcp. legen. lie; be at rest. lie; be lying. Supplied participle describing the possessed persons, not a habban perfect.\par
% glossary-block:B00946:end
% glossary-block:B00947:begin
\needspace{3\baselineskip}
\textbf{līchama} (n. m.); gen. sg. līchaman; nom. pl. līchaman. body.\par
% glossary-block:B00947:end
% glossary-block:B00948:begin
\needspace{3\baselineskip}
\textbf{līchamlic} (adj.). bodily.\par
% glossary-block:B00948:end
% glossary-block:B00949:begin
\needspace{3\baselineskip}
\textbf{lician} (v.). please. The thing pleasing is the subject; the persons pleased are dative in the supplied reading.\par
% glossary-block:B00949:end
% glossary-block:B00950:begin
\needspace{3\baselineskip}
\textbf{līefan} (v.). allow; permit.\par
% glossary-block:B00950:end
% glossary-block:B00951:begin
\needspace{3\baselineskip}
\textbf{līeg} (n. m.). thunderbolt; lightning. Source levin sense.\par
% glossary-block:B00951:end
% glossary-block:B00952:begin
\needspace{3\baselineskip}
\textbf{līehting} (n. f.); gen. sg. līehtinge; nom. pl. līehtinga. illumination.\par
% glossary-block:B00952:end
% glossary-block:B00953:begin
\needspace{3\baselineskip}
\textbf{līf} (n. n.); gen. sg. līfes; nom. pl. līf. life. life. An individual's lived life, including cognate accusative with libban.\par
% glossary-block:B00953:end
% glossary-block:B00954:begin
\needspace{3\baselineskip}
\textbf{lifde} — pret. sg.: see \textbf{libban}.\par
% glossary-block:B00954:end
% glossary-block:B00955:begin
\needspace{3\baselineskip}
\textbf{lifian} — citation/spelling variant: see \textbf{libban}.\par
% glossary-block:B00955:end
% glossary-block:B00956:begin
\needspace{3\baselineskip}
\textbf{līð} — pres. ind. 3 sg., variant: see \textbf{licgan}.\par
% glossary-block:B00956:end
% glossary-block:B00957:begin
\needspace{3\baselineskip}
\textbf{lof} (n. m.); gen. sg. lofes. praise.\par
% glossary-block:B00957:end
% glossary-block:B00958:begin
\needspace{3\baselineskip}
\textbf{lor} (n. n.). destruction; ruin.\par
% glossary-block:B00958:end
% glossary-block:B00959:begin
\needspace{3\baselineskip}
\textbf{lufe} (n. f.). love.\par
% glossary-block:B00959:end
% glossary-block:B00960:begin
\needspace{3\baselineskip}
\textbf{lufian} (v. weak); pret. 1 sg. lufode; past ptcp. (ge)lufod. love.\par
% glossary-block:B00960:end
% glossary-block:B00961:begin
\needspace{3\baselineskip}
\textbf{lufiend} (n. m.). lover.\par
% glossary-block:B00961:end
% glossary-block:B00962:begin
\needspace{3\baselineskip}
\textbf{lufu} (n. f.). love.\par
% glossary-block:B00962:end
% glossary-block:B00963:begin
\needspace{3\baselineskip}
\textbf{lyft} (n. n.); gen. sg. lyftes; nom. pl. lyft. air.\par
% glossary-block:B00963:end
% glossary-block:B00964:begin
\needspace{3\baselineskip}
\textbf{lyre} (n. m.); gen. sg. lyres; nom. pl. lyras. loss.\par
% glossary-block:B00964:end
% glossary-block:B00965:begin
\needspace{3\baselineskip}
\textbf{lystan} (v.). please; be desired. Impersonal source expression; keep the complete supplied frame.\par
% glossary-block:B00965:end
% glossary-block:B00966:begin
\needspace{3\baselineskip}
\textbf{lȳtel} (adj.); comp. lǣssa; superl. lǣst. little; small. Stem lȳtl- before vowel endings; lȳtel- before -ne, -re and -ra.\par
% glossary-block:B00966:end
% book-navigation: standalone bookmark replaced by volume hierarchy
% glossary-block:B00967:begin
\needspace{5\baselineskip}
\subsection*{M}\BookHeading{nav-glossary-old-english-0014}{subsubsection}{M}{0}
% glossary-block:B00967:end
% glossary-block:B00968:begin
\textbf{mā} (adv.). more. Adverb of amount. more; again. In the negative phrase never more/never again. more. More in number, distinct from degree or greater value.\par
% glossary-block:B00968:end
% glossary-block:B00969:begin
\needspace{3\baselineskip}
\textbf{magan} — citation/spelling variant: see \textbf{mugan}.\par
% glossary-block:B00969:end
% glossary-block:B00970:begin
\needspace{3\baselineskip}
\textbf{man} — citation/spelling variant: see \textbf{mann}.\par
% glossary-block:B00970:end
% glossary-block:B00971:begin
\needspace{3\baselineskip}
\textbf{mangere} (n. m.); gen. sg. mangeres; nom. pl. mangeras. merchant.\par
% glossary-block:B00971:end
% glossary-block:B00972:begin
\needspace{3\baselineskip}
\textbf{manig} (adj.). many.\par
% glossary-block:B00972:end
% glossary-block:B00973:begin
\needspace{3\baselineskip}
\textbf{manigfeald} (adj.). manifold; many.\par
% glossary-block:B00973:end
% glossary-block:B00974:begin
\needspace{3\baselineskip}
\textbf{mann} (n.; impersonal pronoun m.); gen. sg. mannes; nom. pl. męnn. man; human being. one; people in general. Impersonal man with singular finite verb.\par
% glossary-block:B00974:end
% glossary-block:B00975:begin
\needspace{3\baselineskip}
\textbf{māra} (adj.). greater; larger. Comparative adjective. greater; more important. Value/importance in the life/raiment comparison. more; a greater amount. Nominal/attributive quantity, distinct from more in number. Comparative of micel.\par
% glossary-block:B00975:end
% glossary-block:B00976:begin
\needspace{3\baselineskip}
\textbf{Maria} (proper name). Mary.\par
% glossary-block:B00976:end
% glossary-block:B00977:begin
\needspace{3\baselineskip}
\textbf{mǣden} (n. n.); gen. sg. mǣdenes; nom. pl. mǣden. girl; maiden.\par
% glossary-block:B00977:end
% glossary-block:B00978:begin
\needspace{3\baselineskip}
\textbf{mæg} — pres. ind. 1/3 sg.: see \textbf{mugan}.\par
% glossary-block:B00978:end
% glossary-block:B00979:begin
\needspace{3\baselineskip}
\textbf{mægen} (n. n.); gen. sg. mægenes; nom. pl. mægnu. power; strength. Heavenly powers in 37.2 is the contextual plural use.\par
% glossary-block:B00979:end
% glossary-block:B00980:begin
\needspace{3\baselineskip}
\textbf{mægenðrymm} (n. m.); gen. sg. mægenðrymmes. majesty; glory.\par
% glossary-block:B00980:end
% glossary-block:B00981:begin
\needspace{3\baselineskip}
\textbf{mægenðrymnes} (n. f.); gen. sg. mægenðrymnesse. glory; majesty.\par
% glossary-block:B00981:end
% glossary-block:B00982:begin
\needspace{3\baselineskip}
\textbf{mægð} (n. f.). maiden; girl. Reference in a source or explanatory note.\par
% glossary-block:B00982:end
% glossary-block:B00983:begin
\needspace{3\baselineskip}
\textbf{mǣgð} (n. f.); gen. sg. mǣgðe; nom. pl. mǣgða. tribe; nation.\par
% glossary-block:B00983:end
% glossary-block:B00984:begin
\needspace{3\baselineskip}
\textbf{mægðhād} (n. m.). maidenhood; virginity.\par
% glossary-block:B00984:end
% glossary-block:B00985:begin
\needspace{3\baselineskip}
\textbf{mǣre} (adj.). great; renowned. Distinct from simple physical size where the task requires the renowned sense.\par
% glossary-block:B00985:end
% glossary-block:B00986:begin
\needspace{3\baselineskip}
\textbf{mǣrsian} (v. weak); pret. 1 sg. mǣrsode; past ptcp. mǣrsod. make great; magnify; glorify.\par
% glossary-block:B00986:end
% glossary-block:B00987:begin
\needspace{3\baselineskip}
\textbf{mæsseprēost} (n. m.); gen. sg. mæsseprēostes; nom. pl. mæsseprēostas. priest.\par
% glossary-block:B00987:end
% glossary-block:B00988:begin
\needspace{3\baselineskip}
\textbf{mǣst} (adj.). greatest; largest. Superlative adjective. Superlative of micel.\par
% glossary-block:B00988:end
% glossary-block:B00989:begin
\needspace{3\baselineskip}
\textbf{mǣð} (n. f.). measure; capacity.\par
% glossary-block:B00989:end
% glossary-block:B00990:begin
\needspace{3\baselineskip}
\textbf{mē} — dat. sg.; acc. sg.: see \textbf{ic}.\par
% glossary-block:B00990:end
% glossary-block:B00991:begin
\needspace{3\baselineskip}
\textbf{meahte} — pret. sg.: see \textbf{mugan}.\par
% glossary-block:B00991:end
% glossary-block:B00992:begin
\needspace{3\baselineskip}
\textbf{mec} — acc. sg., variant: see \textbf{ic}.\par
% glossary-block:B00992:end
% glossary-block:B00993:begin
\needspace{3\baselineskip}
\textbf{mēd} (n. f.); gen. sg. mēde; nom. pl. mēda. reward. hire; wages paid for work.\par
% glossary-block:B00993:end
% glossary-block:B00994:begin
\needspace{3\baselineskip}
\textbf{mēder} — dat. sg.: see \textbf{mōdor}.\par
% glossary-block:B00994:end
% glossary-block:B00995:begin
\needspace{3\baselineskip}
\textbf{medmicel} (adj.). short; brief.\par
% glossary-block:B00995:end
% glossary-block:B00996:begin
\needspace{3\baselineskip}
\textbf{mehte} — pret. sg., variant: see \textbf{mugan}.\par
% glossary-block:B00996:end
% glossary-block:B00997:begin
\needspace{3\baselineskip}
\textbf{męnigu} (n. f.). company; crowd.\par
% glossary-block:B00997:end
% glossary-block:B00998:begin
\needspace{3\baselineskip}
\textbf{męnn} — dat. sg.; nom./acc. pl.: see \textbf{mann}.\par
% glossary-block:B00998:end
% glossary-block:B00999:begin
\needspace{3\baselineskip}
\textbf{męregrēote} (n. f.); gen. sg. męregrēotan; nom. pl. męregrēotan. pearl.\par
% glossary-block:B00999:end
% glossary-block:B01000:begin
\needspace{3\baselineskip}
\textbf{męreswīn} (n. n.); gen. sg. męreswīnes; nom. pl. męreswīn. dolphin.\par
% glossary-block:B01000:end
% glossary-block:B01001:begin
\needspace{3\baselineskip}
\textbf{męrgen} — citation/spelling variant: see \textbf{morgen}.\par
% glossary-block:B01001:end
% glossary-block:B01002:begin
\needspace{3\baselineskip}
\textbf{mētan} (v. weak); pret. 1 sg. mētte; past ptcp. (ge)mēt(t). meet; find.\par
% glossary-block:B01002:end
% glossary-block:B01003:begin
\needspace{3\baselineskip}
\textbf{męte} (n. m.); gen. sg. mętes; nom. pl. mętas. food.\par
% glossary-block:B01003:end
% glossary-block:B01004:begin
\needspace{3\baselineskip}
\textbf{micel} (adj.); comp. māra; superl. mǣst. great; large. much; an amount. Quantity sense. great; loud. Loud is the contextual force with a voice. Stem micl- before vowel endings; micel- before -ne, -re and -ra.\par
% glossary-block:B01004:end
% glossary-block:B01005:begin
\needspace{3\baselineskip}
\textbf{Michael} (proper name). Michael.\par
% glossary-block:B01005:end
% glossary-block:B01006:begin
\needspace{3\baselineskip}
\textbf{mid} (prep.; adv. in fixed expressions). with; by means of. Dative instrument or means. with. Manner in the supplied compassion phrase. with; accompanying. Dative person or cargo in the ordinary accompaniment frame. with. Circumstances of risk or loss in the supplied phrase. with. Manner/circumstance of victory. completely. Only in fixed mid ealle. with. Dative garment in the assigned clothing frame. with. Dative conversational partner. with; along with. Supplied mid wes, be with them. with. Accompanying quality or mental disposition in the source. with. Material or clothing in the supplied adornment phrase. with. Exact cloth/cloak complements of begyrd and bewǣfed. with; by means of.\par
% glossary-block:B01006:end
% glossary-block:B01007:begin
\needspace{3\baselineskip}
\textbf{middangeard} (n. m.); gen. sg. middangeardes; nom. pl. middangeardas. world; earth.\par
% glossary-block:B01007:end
% glossary-block:B01008:begin
\needspace{3\baselineskip}
\textbf{middes} (adverb or particle). at noon. Whole supplied tō middes dæges, with genitives.\par
% glossary-block:B01008:end
% glossary-block:B01009:begin
\needspace{3\baselineskip}
\textbf{mid-ðȳ} (conj.). when; while.\par
% glossary-block:B01009:end
% glossary-block:B01010:begin
\needspace{3\baselineskip}
\textbf{miht} (n. f.). ability; power. powers; powerful people. Those holding high station in the source's concluding human application.\par
% glossary-block:B01010:end
% glossary-block:B01011:begin
\needspace{3\baselineskip}
\textbf{mihtig} (adj.). mighty; powerful. the powerful; the mighty. Substantival adjective. Stem mihtg- before vowel endings and mihtig- before consonantal endings; unsyncopated forms also occur.\par
% glossary-block:B01011:end
% glossary-block:B01012:begin
\needspace{3\baselineskip}
\textbf{mildheort} (adj.). merciful; compassionate.\par
% glossary-block:B01012:end
% glossary-block:B01013:begin
\needspace{3\baselineskip}
\textbf{mildheortnes} (n. f.); gen. sg. mildheortnesse; nom. pl. mildheortnessa. compassion. mercy; mercies.\par
% glossary-block:B01013:end
% glossary-block:B01014:begin
\needspace{3\baselineskip}
\textbf{mīn} (possessive adjective). my; mine. Also gen. sg.: see \textbf{ic}.\par
% glossary-block:B01014:end
% glossary-block:B01015:begin
\needspace{3\baselineskip}
\textbf{mislic} (adj.). various.\par
% glossary-block:B01015:end
% glossary-block:B01016:begin
\needspace{3\baselineskip}
\textbf{mōd} (n. n.). mind; attention. heart; mind.\par
% glossary-block:B01016:end
% glossary-block:B01017:begin
\needspace{3\baselineskip}
\textbf{mōdor} (n. f.); gen. sg. mōdor; nom. pl. mōdru. mother. Plural variant mōdra; dative singular mēder, also a late genitive variant.\par
% glossary-block:B01017:end
% glossary-block:B01018:begin
\needspace{3\baselineskip}
\textbf{mōna} (n. m.); gen. sg. mōnan; nom. pl. mōnan. moon.\par
% glossary-block:B01018:end
% glossary-block:B01019:begin
\needspace{3\baselineskip}
\textbf{mōnað} (n. m.). month. Source plural months.\par
% glossary-block:B01019:end
% glossary-block:B01020:begin
\needspace{3\baselineskip}
\textbf{morgen} (n. m.); gen. sg. morgnes; nom. pl. morgnas. morning.\par
% glossary-block:B01020:end
% glossary-block:B01021:begin
\needspace{3\baselineskip}
\textbf{mōste} — pret. sg.: see \textbf{mōtan}.\par
% glossary-block:B01021:end
% glossary-block:B01022:begin
\needspace{3\baselineskip}
\textbf{mōtan} (v. preterite-present); pres. 1/3 sg. mōt; pres. pl. mōton; pret. 1 sg. mōste. may; be allowed to. Conventional citation label; no usable infinitive is supplied. See the finite forms in Unit 6. No past-participle citation is given in the checked source paradigm.\par
% glossary-block:B01022:end
% glossary-block:B01023:begin
\needspace{3\baselineskip}
\textbf{mugan} (v. preterite-present); pres. 1/3 sg. mæg; pres. pl. magon; pret. 1 sg. meahte. can; be able to. be able; could. Supplied modal form; ordinary capability, not permission. Conventional citation label; no usable infinitive is supplied. See the finite forms in Unit 6. No past-participle citation is given in the checked source paradigm.\par
% glossary-block:B01023:end
% glossary-block:B01024:begin
\needspace{3\baselineskip}
\textbf{munan} (v.). intend. Reference in a source or explanatory note.\par
% glossary-block:B01024:end
% glossary-block:B01025:begin
\needspace{3\baselineskip}
\textbf{munt} (n. m.); gen. sg. muntes; nom. pl. muntas. mountain.\par
% glossary-block:B01025:end
% glossary-block:B01026:begin
\needspace{3\baselineskip}
\textbf{munuc} (n. m.); gen. sg. munuces; nom. pl. munucas. monk.\par
% glossary-block:B01026:end
% glossary-block:B01027:begin
\needspace{3\baselineskip}
\textbf{mūð} (n. m.); gen. sg. mūðes; nom. pl. mūðas. mouth.\par
% glossary-block:B01027:end
% glossary-block:B01028:begin
\needspace{3\baselineskip}
\textbf{myngian} (v.); past ptcp. myngod. admonish; prompt. Supplied participial source use.\par
% glossary-block:B01028:end
% glossary-block:B01029:begin
\needspace{3\baselineskip}
\textbf{mynster} (n. n.); gen. sg. mynstres; nom. pl. mynstru. monastery.\par
% glossary-block:B01029:end
% book-navigation: standalone bookmark replaced by volume hierarchy
% glossary-block:B01030:begin
\needspace{5\baselineskip}
\subsection*{N}\BookHeading{nav-glossary-old-english-0015}{subsubsection}{N}{0}
% glossary-block:B01030:end
% glossary-block:B01031:begin
\textbf{nā} (adv.). not. Negates a contrasted constituent, distinct from finite ne.\par
% glossary-block:B01031:end
% glossary-block:B01032:begin
\needspace{3\baselineskip}
\textbf{nabban} (v. weak); pret. 1 sg. næfde; past ptcp. (ge)næfd; pres. 3 sg. næfð. have not.\par
% glossary-block:B01032:end
% glossary-block:B01033:begin
\needspace{3\baselineskip}
\textbf{nacod} (adj.). naked. naked people. Substantival adjective.\par
% glossary-block:B01033:end
% glossary-block:B01034:begin
\needspace{3\baselineskip}
\textbf{nafogār} (n.). auger. Reference in a source or explanatory note.\par
% glossary-block:B01034:end
% glossary-block:B01035:begin
\needspace{3\baselineskip}
\textbf{nāgan} (v.); pres. 3 sg. nāh. not have; not possess. Negative possessive verb in its supplied forms.\par
% glossary-block:B01035:end
% glossary-block:B01036:begin
\needspace{3\baselineskip}
\textbf{nāh} — pres. ind. 3 sg.: see \textbf{nāgan}.\par
% glossary-block:B01036:end
% glossary-block:B01037:begin
\needspace{3\baselineskip}
\textbf{nāht} (substantive n.). nothing; naught.\par
% glossary-block:B01037:end
% glossary-block:B01038:begin
\needspace{3\baselineskip}
\textbf{nāhwæðer} (adverb or particle). neither.\par
% glossary-block:B01038:end
% glossary-block:B01039:begin
\needspace{3\baselineskip}
\textbf{nālæs} (adv.). not at all.\par
% glossary-block:B01039:end
% glossary-block:B01040:begin
\needspace{3\baselineskip}
\textbf{nama} (n. m.); gen. sg. naman; nom. pl. naman. name.\par
% glossary-block:B01040:end
% glossary-block:B01041:begin
\needspace{3\baselineskip}
\textbf{nān} (indefinite). no; no one.\par
% glossary-block:B01041:end
% glossary-block:B01042:begin
\needspace{3\baselineskip}
\textbf{nāst} — negative pres. ind. 2 sg.: see \textbf{nytan}.\par
% glossary-block:B01042:end
% glossary-block:B01043:begin
\needspace{3\baselineskip}
\textbf{nāt} — negative pres. ind. 1 sg.; supplied source form: see \textbf{nytan}.\par
% glossary-block:B01043:end
% glossary-block:B01044:begin
\needspace{3\baselineskip}
\textbf{nāwiht} (indefinite pronoun/substantive). nothing. nothing; naught.\par
% glossary-block:B01044:end
% glossary-block:B01045:begin
\needspace{3\baselineskip}
\textbf{næddre} (n. f.). snake; adder. Supplied weak feminine plural.\par
% glossary-block:B01045:end
% glossary-block:B01046:begin
\needspace{3\baselineskip}
\textbf{nǣdl} (n. f.); gen. sg. nǣdle; nom. pl. nǣdla. needle.\par
% glossary-block:B01046:end
% glossary-block:B01047:begin
\needspace{3\baselineskip}
\textbf{nædre} (n. f.); gen. sg. nædran; nom. pl. nædran. serpent.\par
% glossary-block:B01047:end
% glossary-block:B01048:begin
\needspace{3\baselineskip}
\textbf{næfde} — pret. sg.: see \textbf{nabban}.\par
% glossary-block:B01048:end
% glossary-block:B01049:begin
\needspace{3\baselineskip}
\textbf{nǣfre} (adv.). never.\par
% glossary-block:B01049:end
% glossary-block:B01050:begin
\needspace{3\baselineskip}
\textbf{næfð} — pres. ind. 3 sg.: see \textbf{nabban}.\par
% glossary-block:B01050:end
% glossary-block:B01051:begin
\needspace{3\baselineskip}
\textbf{nǣnig} (indefinite pronoun/adjective). no; not any; none. Reference in a source or explanatory note.\par
% glossary-block:B01051:end
% glossary-block:B01052:begin
\needspace{3\baselineskip}
\textbf{nǣron} — negative pret. ind. pl.: see \textbf{wesan}.\par
% glossary-block:B01052:end
% glossary-block:B01053:begin
\needspace{3\baselineskip}
\textbf{næs} — negative pret. ind. 1/3 sg.: see \textbf{wesan}.\par
% glossary-block:B01053:end
% glossary-block:B01054:begin
\needspace{3\baselineskip}
\textbf{ne} (particle). not; simple clause negation.\par
% glossary-block:B01054:end
% glossary-block:B01055:begin
\needspace{3\baselineskip}
\textbf{nēah} (adv.). near; at hand.\par
% glossary-block:B01055:end
% glossary-block:B01056:begin
\needspace{3\baselineskip}
\textbf{nēahgebūr} (n. m.); gen. sg. nēahgebūres; nom. pl. nēahgebūras. neighbour.\par
% glossary-block:B01056:end
% glossary-block:B01057:begin
\needspace{3\baselineskip}
\textbf{nēalǣcan} (v. weak); pret. 1 sg. nēalǣcte; past ptcp. nēalǣced. approach; draw near. Dative person or place approached when the goal is expressed.\par
% glossary-block:B01057:end
% glossary-block:B01058:begin
\needspace{3\baselineskip}
\textbf{nearu} (adj.). narrow.\par
% glossary-block:B01058:end
% glossary-block:B01059:begin
\needspace{3\baselineskip}
\textbf{nele} — negative pres. ind. 3 sg.: see \textbf{willan}.\par
% glossary-block:B01059:end
% glossary-block:B01060:begin
\needspace{3\baselineskip}
\textbf{nelle} — negative pres. ind. 1 sg.: see \textbf{willan}.\par
% glossary-block:B01060:end
% glossary-block:B01061:begin
\needspace{3\baselineskip}
\textbf{nelt} — negative pres. ind. 2 sg.: see \textbf{willan}.\par
% glossary-block:B01061:end
% glossary-block:B01062:begin
\needspace{3\baselineskip}
\textbf{nēosian} (v.). seek; look for (+ genitive). genitive object; reference comparison with genēosian.\par
% glossary-block:B01062:end
% glossary-block:B01063:begin
\needspace{3\baselineskip}
\textbf{nēowolnes} (n. f.); gen. sg. nēowolnesse; nom. pl. nēowolnessa. abyss; the deep.\par
% glossary-block:B01063:end
% glossary-block:B01064:begin
\needspace{3\baselineskip}
\textbf{nęrian} (v. weak); pret. 1 sg. nęrede; past ptcp. (ge)nęred. save.\par
% glossary-block:B01064:end
% glossary-block:B01065:begin
\needspace{3\baselineskip}
\textbf{nest} (n. n.); gen. sg. nestes; nom. pl. nest. nest.\par
% glossary-block:B01065:end
% glossary-block:B01066:begin
\needspace{3\baselineskip}
\textbf{nętt} (n. n.); gen. sg. nęttes; nom. pl. nętt. net.\par
% glossary-block:B01066:end
% glossary-block:B01067:begin
\needspace{3\baselineskip}
\textbf{nīede} (adverb or particle). necessarily; of necessity.\par
% glossary-block:B01067:end
% glossary-block:B01068:begin
\needspace{3\baselineskip}
\textbf{nīehst} (adv.). at length. Whole supplied æt nīehstan.\par
% glossary-block:B01068:end
% glossary-block:B01069:begin
\needspace{3\baselineskip}
\textbf{nīeten} (n. n.); gen. sg. nīetenes; nom. pl. nīetenu. animal; beast; cattle. Cattle is the selected plural use; retain the early task distinction from dēor.\par
% glossary-block:B01069:end
% glossary-block:B01070:begin
\needspace{3\baselineskip}
\textbf{nigon} (cardinal numeral). nine. Cardinal 9.\par
% glossary-block:B01070:end
% glossary-block:B01071:begin
\needspace{3\baselineskip}
\textbf{nigoða} (ordinal). ninth.\par
% glossary-block:B01071:end
% glossary-block:B01072:begin
\needspace{3\baselineskip}
\textbf{niht} (n. f.); gen. sg. niht; nom. pl. niht. night.\par
% glossary-block:B01072:end
% glossary-block:B01073:begin
\needspace{3\baselineskip}
\textbf{niman} (v. strong); pret. 1 sg. nōm; pret. pl. nōmon; past ptcp. numen. take. Accusative person or thing taken, including food.\par
% glossary-block:B01073:end
% glossary-block:B01074:begin
\needspace{3\baselineskip}
\textbf{nis} — negative pres. ind. 3 sg.: see \textbf{wesan}.\par
% glossary-block:B01074:end
% glossary-block:B01075:begin
\needspace{3\baselineskip}
\textbf{nīwe} (adj.). new.\par
% glossary-block:B01075:end
% glossary-block:B01076:begin
\needspace{3\baselineskip}
\textbf{nōht} — citation/spelling variant: see \textbf{nāwiht}.\par
% glossary-block:B01076:end
% glossary-block:B01077:begin
\needspace{3\baselineskip}
\textbf{nolde} — negative pret. sg.: see \textbf{willan}.\par
% glossary-block:B01077:end
% glossary-block:B01078:begin
\needspace{3\baselineskip}
\textbf{norðan} (adv.). from the north.\par
% glossary-block:B01078:end
% glossary-block:B01079:begin
\needspace{3\baselineskip}
\textbf{nū} (adv.; conj.). now.\par
% glossary-block:B01079:end
% glossary-block:B01080:begin
\needspace{3\baselineskip}
\textbf{nytan} (v.). not know. Supplied negative knowing form. not know.\par
% glossary-block:B01080:end
% book-navigation: standalone bookmark replaced by volume hierarchy
% glossary-block:B01081:begin
\needspace{5\baselineskip}
\subsection*{O}\BookHeading{nav-glossary-old-english-0016}{subsubsection}{O}{0}
% glossary-block:B01081:end
% glossary-block:B01082:begin
\textbf{of} (prep.). from; out of. With dative. of; from. Portion from a whole in the supplied partitive phrases; dative whole. sense for sense. Literal correspondence in the supplied translation expression. from. Starting condition of old age in the supplied rejuvenation comparison. from. Starting generation in the supplied temporal range. of; from. Origin of wisdom in the supplied begetting statement.\par
% glossary-block:B01082:end
% glossary-block:B01083:begin
\needspace{3\baselineskip}
\textbf{of} (adv.). off.\par
% glossary-block:B01083:end
% glossary-block:B01084:begin
\needspace{3\baselineskip}
\textbf{ofen} (n. m.); gen. sg. ofenes; nom. pl. ofenas. oven.\par
% glossary-block:B01084:end
% glossary-block:B01085:begin
\needspace{3\baselineskip}
\textbf{ofer} (prep.). over; across. Spatial frame; case varies with the specific construction. Do not impose a universal motion rule. over. Accusative dependent in the taught authority frame. above. Rank, with accusative. toward. In gōd ofer + accusative people. over; about. Object of rejoicing, with the assigned accusative frame. above. Rank, with explicitly supplied dative complement; retain this limit. above; more than. Preference in the supplied comparison.\par
% glossary-block:B01085:end
% glossary-block:B01086:begin
\needspace{3\baselineskip}
\textbf{oferwinnan} (v. strong); pret. 1 sg. oferwann; pret. pl. oferwunnon; past ptcp. oferwunnen. overcome; conquer.\par
% glossary-block:B01086:end
% glossary-block:B01087:begin
\needspace{3\baselineskip}
\textbf{oferworhte} — pret. sg.: see \textbf{oferwyrcean}.\par
% glossary-block:B01087:end
% glossary-block:B01088:begin
\needspace{3\baselineskip}
\textbf{oferwyrcean} (v. weak); pret. 1 sg. oferworhte; past ptcp. oferworht. overlay; cover. Accusative object, mid + dative material used to cover it.\par
% glossary-block:B01088:end
% glossary-block:B01089:begin
\needspace{3\baselineskip}
\textbf{ofslēan} (v.); past ptcp. ofslægen. slay; kill.\par
% glossary-block:B01089:end
% glossary-block:B01090:begin
\needspace{3\baselineskip}
\textbf{oft} (adv.). often.\par
% glossary-block:B01090:end
% glossary-block:B01091:begin
\needspace{3\baselineskip}
\textbf{ofðyncean} (v. impersonal); pres. 3 sg. ofðyncð; pret. 3 sg. ofðūhte; past ptcp. ofðūht. grieve; displease. Impersonal third singular with dative experiencer and a content clause. Impersonal use: the course supplies third-singular finite forms.\par
% glossary-block:B01091:end
% glossary-block:B01092:begin
\needspace{3\baselineskip}
\textbf{ofwundrian} (v.); past ptcp. ofwundrod. be astonished.\par
% glossary-block:B01092:end
% glossary-block:B01093:begin
\needspace{3\baselineskip}
\textbf{ōht} — citation/spelling variant: see \textbf{āwiht}.\par
% glossary-block:B01093:end
% glossary-block:B01094:begin
\needspace{3\baselineskip}
\textbf{olfend} (n. m.); gen. sg. olfendes; nom. pl. olfendas. camel.\par
% glossary-block:B01094:end
% glossary-block:B01095:begin
\needspace{3\baselineskip}
\textbf{on} (prep.). in; on. Dative for the taught ordinary location construction. in the beginning. Whole supplied on anginne. on. Time preposition in the stated-day expression. in. Supplied circumstance of the storm. into; in two. Supplied division expression. alive. Only in supplied on līfe. according to. Only the supplied classificatory expression with hīw. in one's control. on + dative possessed geweald. in. Supported expression in his majesty/glory. at one's side. Accusative hand in the bounded placement expression. on one's side. The bounded on … hand side expression uses accusative, including with standan. in one's name. on + dative naman, with genitive possessor. rich in. Only in supplied berende on + dative. into; onto. Accusative goal of placement in the taught frame. in; amid. Circumstances of the stated unrest. in; with respect to. Quality or domain of description. in. Accusative person believed/trusted in. during; in. Supplied lifetime expression. among; in. Dative people among whom friends are obtained. in worldly life. Supplied temporal/situational expression. during. Dative time expression. in. Circumstance of the day's heat in the supplied source. fishing. Supplied on fiscnoð after gān. in; wearing. Supplied on scēapa gegierelum, in sheep's garments. at work; engaged in. Only the exact supplied activity phrases. upon; into. Supplied chancing-upon idiom. in. Source/object of joy in the supplied frame. to. Succeeding generation in the supplied temporal range. upon. Accusative beloved in the supplied love-falls-upon idiom. in the morning. Whole supplied on morgen; accusative singular. into. Destination of sale in the supplied expression.\par
% glossary-block:B01095:end
% glossary-block:B01096:begin
\needspace{3\baselineskip}
\textbf{on} (adv.). on. Adverb in the cloak-wearing relative clause.\par
% glossary-block:B01096:end
% glossary-block:B01097:begin
\needspace{3\baselineskip}
\textbf{oncnāwan} (v. strong); pret. 1 sg. oncnēow; pret. pl. oncnēowon; past ptcp. oncnāwen; pres. 3 sg. oncnǣwð. recognize. recognize; perceive. Supplied past source form. recognize; acknowledge. recognize; know by recognition.\par
% glossary-block:B01097:end
% glossary-block:B01098:begin
\needspace{3\baselineskip}
\textbf{oncnǣwð} — pres. ind. 3 sg.: see \textbf{oncnāwan}.\par
% glossary-block:B01098:end
% glossary-block:B01099:begin
\needspace{3\baselineskip}
\textbf{ondrǣdan} (v. strong); pret. 1 sg. ondrēd; pret. pl. ondrēdon; past ptcp. ondrǣden. fear.\par
% glossary-block:B01099:end
% glossary-block:B01100:begin
\needspace{3\baselineskip}
\textbf{onfēhð} — pres. ind. 3 sg.: see \textbf{onfōn}.\par
% glossary-block:B01100:end
% glossary-block:B01101:begin
\needspace{3\baselineskip}
\textbf{onfōn} (v. strong); pret. 1 sg. onfēng; pret. pl. onfēngon; past ptcp. onfangen; pres. 3 sg. onfēhð. receive; accept. Dative thing received in this selected frame; the later accusative payment frame is separate. receive; accept. A person received into care in the supplied frame.\par
% glossary-block:B01101:end
% glossary-block:B01102:begin
\needspace{3\baselineskip}
\textbf{ongēan} (adv.; prep.; conj.). against; opposite. Accusative landmark or opponent. back; again. by the time that. Complete supplied ongēan ðætte clause.\par
% glossary-block:B01102:end
% glossary-block:B01103:begin
\needspace{3\baselineskip}
\textbf{ongeslēan} (v.). inflict slaughter; slay.\par
% glossary-block:B01103:end
% glossary-block:B01104:begin
\needspace{3\baselineskip}
\textbf{ongietan} (v. strong); pret. 1 sg. ongeat; pret. pl. ongēaton; past ptcp. ongieten. understand; perceive. Supplied infinitive after a modal. perceive; realize. understand; perceive.\par
% glossary-block:B01104:end
% glossary-block:B01105:begin
\needspace{3\baselineskip}
\textbf{onginnan} (v.); pret. pl. ongunnon. begin. Supplied preterite began. begin. begin. Supplied playing/praising action as complement.\par
% glossary-block:B01105:end
% glossary-block:B01106:begin
\needspace{3\baselineskip}
\textbf{onlīesan} (v. weak); pret. 1 sg. onlīesde; past ptcp. onlīesed. release. Accusative person or thing released.\par
% glossary-block:B01106:end
% glossary-block:B01107:begin
\needspace{3\baselineskip}
\textbf{onsęndan} (v.). send.\par
% glossary-block:B01107:end
% glossary-block:B01108:begin
\needspace{3\baselineskip}
\textbf{ontȳnan} (v.). open. Supplied preterite opened.\par
% glossary-block:B01108:end
% glossary-block:B01109:begin
\needspace{3\baselineskip}
\textbf{onwęg} (adv.). away.\par
% glossary-block:B01109:end
% glossary-block:B01110:begin
\needspace{3\baselineskip}
\textbf{onwinnan} (v.). assail; attack.\par
% glossary-block:B01110:end
% glossary-block:B01111:begin
\needspace{3\baselineskip}
\textbf{open} (adj.). open.\par
% glossary-block:B01111:end
% glossary-block:B01112:begin
\needspace{3\baselineskip}
\textbf{ōra} (n. m.). vein of metal; ore (sense uncertain). Supplied dative plural after on.\par
% glossary-block:B01112:end
% glossary-block:B01113:begin
\needspace{3\baselineskip}
\textbf{ōrmōdnes} (n. f.). despair; hopelessness.\par
% glossary-block:B01113:end
% glossary-block:B01114:begin
\needspace{3\baselineskip}
\textbf{Osburg} (proper name). Osburg.\par
% glossary-block:B01114:end
% glossary-block:B01115:begin
\needspace{3\baselineskip}
\textbf{ōð} (prep.). as far as; up to.\par
% glossary-block:B01115:end
% glossary-block:B01116:begin
\needspace{3\baselineskip}
\textbf{oð} (prep.). up to; until. With accusative.\par
% glossary-block:B01116:end
% glossary-block:B01117:begin
\needspace{3\baselineskip}
\textbf{ōðer} (adj.; ordinal). second. other; another. remaining; the rest. The rest of the forest/earth in the source.\par
% glossary-block:B01117:end
% glossary-block:B01118:begin
\needspace{3\baselineskip}
\textbf{oð-ðæt} (conj.). until.\par
% glossary-block:B01118:end
% glossary-block:B01119:begin
\needspace{3\baselineskip}
\textbf{oððe} (conj.). or.\par
% glossary-block:B01119:end
% book-navigation: standalone bookmark replaced by volume hierarchy
% glossary-block:B01120:begin
\needspace{5\baselineskip}
\subsection*{P}\BookHeading{nav-glossary-old-english-0017}{subsubsection}{P}{0}
% glossary-block:B01120:end
% glossary-block:B01121:begin
\textbf{palmtwīg} (n. n.); gen. sg. palmtwīges; nom. pl. palmtwīg. palm twig; palm branch.\par
% glossary-block:B01121:end
% glossary-block:B01122:begin
\needspace{3\baselineskip}
\textbf{pęning} (n. m.); gen. sg. pęninges; nom. pl. pęningas. penny; coin.\par
% glossary-block:B01122:end
% glossary-block:B01123:begin
\needspace{3\baselineskip}
\textbf{Pentapolis} (proper place name). Pentapolis. Cyrenian region in this story.\par
% glossary-block:B01123:end
% glossary-block:B01124:begin
\needspace{3\baselineskip}
\textbf{Piht} (ethnonym). Picts. Supplied people-name forms.\par
% glossary-block:B01124:end
% glossary-block:B01125:begin
\needspace{3\baselineskip}
\textbf{plega} (n. m.). play; game.\par
% glossary-block:B01125:end
% glossary-block:B01126:begin
\needspace{3\baselineskip}
\textbf{plegian} (v. weak); pret. 1 sg. plegode; past ptcp. plegod. play. Intransitive; the reading supplies local wið participant and mid instrument phrases. act; perform. Performance context; ordinary play is a separate assigned sense.\par
% glossary-block:B01126:end
% glossary-block:B01127:begin
\needspace{3\baselineskip}
\textbf{pliht} (n. m.); gen. sg. plihtes; nom. pl. plihtas. peril; risk.\par
% glossary-block:B01127:end
% glossary-block:B01128:begin
\needspace{3\baselineskip}
\textbf{pund} (n. n.); gen. sg. pundes; nom. pl. pund. pound. Weight measure in the supplied gold quantity.\par
% glossary-block:B01128:end
% book-navigation: standalone bookmark replaced by volume hierarchy
% glossary-block:B01129:begin
\needspace{5\baselineskip}
\subsection*{R}\BookHeading{nav-glossary-old-english-0018}{subsubsection}{R}{0}
% glossary-block:B01129:end
% glossary-block:B01130:begin
\textbf{rǣdan} (v. weak); pret. 1 sg. rǣdde; past ptcp. gerǣd. read. Accusative text, or absolute with no expressed object.\par
% glossary-block:B01130:end
% glossary-block:B01131:begin
\needspace{3\baselineskip}
\textbf{rēaf} (n. n.); gen. sg. rēafes; nom. pl. rēaf. raiment; clothing. clothing; garments.\par
% glossary-block:B01131:end
% glossary-block:B01132:begin
\needspace{3\baselineskip}
\textbf{rēcelēasian} (v.). neglect; despise.\par
% glossary-block:B01132:end
% glossary-block:B01133:begin
\needspace{3\baselineskip}
\textbf{regn} (n. m.); gen. sg. regnes; nom. pl. regnas. rain.\par
% glossary-block:B01133:end
% glossary-block:B01134:begin
\needspace{3\baselineskip}
\textbf{ręst} (n. f.); gen. sg. ręste; nom. pl. ręsta. rest; repose.\par
% glossary-block:B01134:end
% glossary-block:B01135:begin
\needspace{3\baselineskip}
\textbf{rēðnes} (n. f.). fury; fierceness.\par
% glossary-block:B01135:end
% glossary-block:B01136:begin
\needspace{3\baselineskip}
\textbf{rēwð} — pres. ind. 3 sg.: see \textbf{rōwan}.\par
% glossary-block:B01136:end
% glossary-block:B01137:begin
\needspace{3\baselineskip}
\textbf{rīce} (n. n.); gen. sg. rīces; nom. pl. rīcu. kingdom.\par
% glossary-block:B01137:end
% glossary-block:B01138:begin
\needspace{3\baselineskip}
\textbf{rīce} (adj.). rich; wealthy. Includes substantival rich people in the assigned frame.\par
% glossary-block:B01138:end
% glossary-block:B01139:begin
\needspace{3\baselineskip}
\textbf{rīcsian} (v. weak); pret. 1 sg. rīcsode; past ptcp. rīcsod. rule; have dominion. Used absolutely in the taught frame; no new object government.\par
% glossary-block:B01139:end
% glossary-block:B01140:begin
\needspace{3\baselineskip}
\textbf{rīdan} (v. strong); pret. 1 sg. rād; pret. pl. ridon; past ptcp. riden. ride. Intransitive in the introduced frame.\par
% glossary-block:B01140:end
% glossary-block:B01141:begin
\needspace{3\baselineskip}
\textbf{riht} (adj.). right; just. direct.\par
% glossary-block:B01141:end
% glossary-block:B01142:begin
\needspace{3\baselineskip}
\textbf{rihtlīce} (adv.); comp. rihtlicor. In comparative rihtlicor: more correctly.\par
% glossary-block:B01142:end
% glossary-block:B01143:begin
\needspace{3\baselineskip}
\textbf{rihtlicor} — comparative adverb; exact source spelling: see \textbf{rihtlīce}.\par
% glossary-block:B01143:end
% glossary-block:B01144:begin
\needspace{3\baselineskip}
\textbf{rihtwīs} (adj.). righteous. righteous people. Substantival adjective.\par
% glossary-block:B01144:end
% glossary-block:B01145:begin
\needspace{3\baselineskip}
\textbf{rihtwīsnes} (n. f.); gen. sg. rihtwīsnesse; nom. pl. rihtwīsnessa. righteousness.\par
% glossary-block:B01145:end
% glossary-block:B01146:begin
\needspace{3\baselineskip}
\textbf{rīnan} (v.). rain. Supplied impersonal source clause with past rīnde.\par
% glossary-block:B01146:end
% glossary-block:B01147:begin
\needspace{3\baselineskip}
\textbf{rīpan} (v.). reap.\par
% glossary-block:B01147:end
% glossary-block:B01148:begin
\needspace{3\baselineskip}
\textbf{rīpe} (adj.). ripe.\par
% glossary-block:B01148:end
% glossary-block:B01149:begin
\needspace{3\baselineskip}
\textbf{rīsan} (v. strong); pret. 1 sg. rās; pret. pl. rison; past ptcp. risen. rise. rise up. Intransitive rīsan, optionally with ūp; ongēan + accusative marks opposition.\par
% glossary-block:B01149:end
% glossary-block:B01150:begin
\needspace{3\baselineskip}
\textbf{rōd} (n. f.); gen. sg. rōde; nom. pl. rōda. cross.\par
% glossary-block:B01150:end
% glossary-block:B01151:begin
\needspace{3\baselineskip}
\textbf{Rōmane} (ethnonym). Romans.\par
% glossary-block:B01151:end
% glossary-block:B01152:begin
\needspace{3\baselineskip}
\textbf{Rōmanisc} (adj.). Roman.\par
% glossary-block:B01152:end
% glossary-block:B01153:begin
\needspace{3\baselineskip}
\textbf{Rome} (proper name). Rome.\par
% glossary-block:B01153:end
% glossary-block:B01154:begin
\needspace{3\baselineskip}
\textbf{rōwan} (v. strong); pret. 1 sg. rēow; pret. pl. rēowon; past ptcp. rōwen; pres. 3 sg. rēwð. row. Intransitive; the assigned route is expressed with ofer.\par
% glossary-block:B01154:end
% book-navigation: standalone bookmark replaced by volume hierarchy
% glossary-block:B01155:begin
\needspace{5\baselineskip}
\subsection*{S}\BookHeading{nav-glossary-old-english-0019}{subsubsection}{S}{0}
% glossary-block:B01155:end
% glossary-block:B01156:begin
\textbf{Saba} (proper name). Sheba.\par
% glossary-block:B01156:end
% glossary-block:B01157:begin
\needspace{3\baselineskip}
\textbf{sacc} (n. m.); gen. sg. sacces; nom. pl. saccas. sack.\par
% glossary-block:B01157:end
% glossary-block:B01158:begin
\needspace{3\baselineskip}
\textbf{sāde} — pret. sg., source variant: see \textbf{sęcgan}.\par
% glossary-block:B01158:end
% glossary-block:B01159:begin
\needspace{3\baselineskip}
\textbf{sagað} — pres. ind. 3 sg., variant: see \textbf{sęcgan}.\par
% glossary-block:B01159:end
% glossary-block:B01160:begin
\needspace{3\baselineskip}
\textbf{Salomon} (proper name); gen. sg. Salomones. Solomon.\par
% glossary-block:B01160:end
% glossary-block:B01161:begin
\needspace{3\baselineskip}
\textbf{samod} (adv.). together.\par
% glossary-block:B01161:end
% glossary-block:B01162:begin
\needspace{3\baselineskip}
\textbf{Sanctus} (foreign adjective). Saint; holy. Latin word in the source title.\par
% glossary-block:B01162:end
% glossary-block:B01163:begin
\needspace{3\baselineskip}
\textbf{sang} (n. m.); gen. sg. sanges; nom. pl. sangas. song.\par
% glossary-block:B01163:end
% glossary-block:B01164:begin
\needspace{3\baselineskip}
\textbf{sārnes} (n. f.); gen. sg. sārnesse. grief.\par
% glossary-block:B01164:end
% glossary-block:B01165:begin
\needspace{3\baselineskip}
\textbf{sāwol} (n. f.); gen. sg. sāwle; nom. pl. sāwla. soul.\par
% glossary-block:B01165:end
% glossary-block:B01166:begin
\needspace{3\baselineskip}
\textbf{sǣ} (n. m.); gen. sg. sǣs; nom. pl. sǣs; source f. gen. sg. sǣ. sea. Also feminine in the supplied source passage.\par
% glossary-block:B01166:end
% glossary-block:B01167:begin
\needspace{3\baselineskip}
\textbf{sǣd} (n. n.); gen. sg. sǣdes; nom. pl. sǣd. seed. seed; descendants. Collective descendants, distinct from literal seed.\par
% glossary-block:B01167:end
% glossary-block:B01168:begin
\needspace{3\baselineskip}
\textbf{sǣde} — dat. n. sg.: see \textbf{sǣd}; pret. sg., variant: see \textbf{sęcgan}.\par
% glossary-block:B01168:end
% glossary-block:B01169:begin
\needspace{3\baselineskip}
\textbf{sǣdtīma} (n. m.); gen. sg. sǣdtīman; nom. pl. sǣdtīman. seedtime; sowing season.\par
% glossary-block:B01169:end
% glossary-block:B01170:begin
\needspace{3\baselineskip}
\textbf{sægde} — pret. sg.: see \textbf{sęcgan}.\par
% glossary-block:B01170:end
% glossary-block:B01171:begin
\needspace{3\baselineskip}
\textbf{sægð} — pres. ind. 3 sg.: see \textbf{sęcgan}.\par
% glossary-block:B01171:end
% glossary-block:B01172:begin
\needspace{3\baselineskip}
\textbf{sǣlic} (adj.). marine; of the sea.\par
% glossary-block:B01172:end
% glossary-block:B01173:begin
\needspace{3\baselineskip}
\textbf{sceal} — pres. ind. 1/3 sg.: see \textbf{sculan}.\par
% glossary-block:B01173:end
% glossary-block:B01174:begin
\needspace{3\baselineskip}
\textbf{scēap} (n. n.); gen. sg. scēapes; nom. pl. scēap. sheep.\par
% glossary-block:B01174:end
% glossary-block:B01175:begin
\needspace{3\baselineskip}
\textbf{scēaphierde} (n. m.); gen. sg. scēaphierdes; nom. pl. scēaphierdas. shepherd.\par
% glossary-block:B01175:end
% glossary-block:B01176:begin
\needspace{3\baselineskip}
\textbf{scēawere} (n. m.); gen. sg. scēaweres; nom. pl. scēaweras. spy.\par
% glossary-block:B01176:end
% glossary-block:B01177:begin
\needspace{3\baselineskip}
\textbf{scēawian} (v. weak); pret. 1 sg. scēawode; past ptcp. scēawod. watch; observe. Accusative person or event watched, including a journey.\par
% glossary-block:B01177:end
% glossary-block:B01178:begin
\needspace{3\baselineskip}
\textbf{scęncan} (v. weak); pret. 1 sg. scęncte; past ptcp. scęnced. pour out; give to drink. Accusative drink, dative recipient.\par
% glossary-block:B01178:end
% glossary-block:B01179:begin
\needspace{3\baselineskip}
\textbf{sceolde} — pret. sg.: see \textbf{sculan}.\par
% glossary-block:B01179:end
% glossary-block:B01180:begin
\needspace{3\baselineskip}
\textbf{scēowyrhta} (n. m.); gen. sg. scēowyrhtan; nom. pl. scēowyrhtan. shoemaker.\par
% glossary-block:B01180:end
% glossary-block:B01181:begin
\needspace{3\baselineskip}
\textbf{scęððan} (v. weak); pret. 1 sg. scęðede; past ptcp. scęðed. hurt; injure. Dative person harmed.\par
% glossary-block:B01181:end
% glossary-block:B01182:begin
\needspace{3\baselineskip}
\textbf{sciccels} (n. m.). cloak.\par
% glossary-block:B01182:end
% glossary-block:B01183:begin
\needspace{3\baselineskip}
\textbf{scield} (n. m.); gen. sg. scieldes; nom. pl. scieldas. shield.\par
% glossary-block:B01183:end
% glossary-block:B01184:begin
\needspace{3\baselineskip}
\textbf{scieppan} (v. strong); pret. 1 sg. scēop; pret. pl. scēopon; past ptcp. sceapen. create; shape.\par
% glossary-block:B01184:end
% glossary-block:B01185:begin
\needspace{3\baselineskip}
\textbf{scieppend} (n. m.). Creator.\par
% glossary-block:B01185:end
% glossary-block:B01186:begin
\needspace{3\baselineskip}
\textbf{scīete} (n. f.). linen cloth.\par
% glossary-block:B01186:end
% glossary-block:B01187:begin
\needspace{3\baselineskip}
\textbf{scīnan} (v. strong); pret. 1 sg. scān; pret. pl. scinon; past ptcp. scinen. shine. shining.\par
% glossary-block:B01187:end
% glossary-block:B01188:begin
\needspace{3\baselineskip}
\textbf{scip} (n. n.); gen. sg. scipes; nom. pl. scipu. ship.\par
% glossary-block:B01188:end
% glossary-block:B01189:begin
\needspace{3\baselineskip}
\textbf{sciphęre} (n. m.); gen. sg. sciphęres; nom. pl. sciphęras. fleet; naval force.\par
% glossary-block:B01189:end
% glossary-block:B01190:begin
\needspace{3\baselineskip}
\textbf{scolde} — pret. sg., variant: see \textbf{sculan}.\par
% glossary-block:B01190:end
% glossary-block:B01191:begin
\needspace{3\baselineskip}
\textbf{Scottas} (ethnonym). Scots.\par
% glossary-block:B01191:end
% glossary-block:B01192:begin
\needspace{3\baselineskip}
\textbf{scrȳdan} (v. weak); pret. 1 sg. scrȳdde; past ptcp. scrȳded. clothe; dress. Accusative person clothed in this introduced frame. dress oneself. Supplied reflexive use.\par
% glossary-block:B01192:end
% glossary-block:B01193:begin
\needspace{3\baselineskip}
\textbf{scūfan} (v. strong); pret. 1 sg. scēaf; pret. pl. scufon; past ptcp. scofen. thrust.\par
% glossary-block:B01193:end
% glossary-block:B01194:begin
\needspace{3\baselineskip}
\textbf{sculan} (v. preterite-present); pres. 1/3 sg. sceal; pres. pl. sculon; pret. 1 sg. sceolde. must; be obliged to. should; be to. Supplied question about what events were to mean; retain its reported/interpretive force. shall; be yet to come. Supplied gīet sculon fīf, with be explicitly omitted. Conventional citation label; no usable infinitive is supplied. See the finite forms in Unit 6. No past-participle citation is given in the checked source paradigm.\par
% glossary-block:B01194:end
% glossary-block:B01195:begin
\needspace{3\baselineskip}
\textbf{scyldig} (adj.). guilty; culpable.\par
% glossary-block:B01195:end
% glossary-block:B01196:begin
\needspace{3\baselineskip}
\textbf{sē} — nom. m sg.: see \textbf{se, sēo, ðæt}.\par
% glossary-block:B01196:end
% glossary-block:B01197:begin
\needspace{3\baselineskip}
\textbf{sealde} — pret. sg.: see \textbf{sęllan}.\par
% glossary-block:B01197:end
% glossary-block:B01198:begin
\needspace{3\baselineskip}
\textbf{sealf} (n. f.); gen. sg. sealfe. ointment.\par
% glossary-block:B01198:end
% glossary-block:B01199:begin
\needspace{3\baselineskip}
\textbf{sēamere} (n. m.); gen. sg. sēameres; nom. pl. sēameras. tailor.\par
% glossary-block:B01199:end
% glossary-block:B01200:begin
\needspace{3\baselineskip}
\textbf{sēcan} (v. weak); pret. 1 sg. sōhte; past ptcp. (ge)sōht. seek. Accusative person or thing sought.\par
% glossary-block:B01200:end
% glossary-block:B01201:begin
\needspace{3\baselineskip}
\textbf{sęcgan} (v. weak); pret. 1 sg. sægde; past ptcp. (ge)sægd; pres. 3 sg. sægð. say. tell; say. Dative addressee and a reported statement or question as content. say. Supplied source preterite. say; tell. call; declare to be. Object with the stated agreeing predicate.\par
% glossary-block:B01201:end
% glossary-block:B01202:begin
\needspace{3\baselineskip}
\textbf{sęgð} — pres. ind. 3 sg., variant: see \textbf{sęcgan}.\par
% glossary-block:B01202:end
% glossary-block:B01203:begin
\needspace{3\baselineskip}
\textbf{sęlen} (n. f.). bounty; bestowal.\par
% glossary-block:B01203:end
% glossary-block:B01204:begin
\needspace{3\baselineskip}
\textbf{self} (adj.). self; the same person or thing. Agreeing emphatic/reflexive self; ordinary reflexive personal pronouns do not always require it.\par
% glossary-block:B01204:end
% glossary-block:B01205:begin
\needspace{3\baselineskip}
\textbf{sęllan} (v. weak); pret. 1 sg. sealde; past ptcp. (ge)seald. give. give to drink. Supplied dative person + plain drincan, without tō. sell. Accusative thing sold.\par
% glossary-block:B01205:end
% glossary-block:B01206:begin
\needspace{3\baselineskip}
\textbf{sellic} (adj.). strange; remarkable.\par
% glossary-block:B01206:end
% glossary-block:B01207:begin
\needspace{3\baselineskip}
\textbf{sęmninga} (adv.). suddenly.\par
% glossary-block:B01207:end
% glossary-block:B01208:begin
\needspace{3\baselineskip}
\textbf{sęndan} (v. weak); pret. 1 sg. sęnde; past ptcp. (ge)sęnd(ed). send. Accusative person or thing sent, dative recipient.\par
% glossary-block:B01208:end
% glossary-block:B01209:begin
\needspace{3\baselineskip}
\textbf{sengan} (v.). singe. Pronunciation example.\par
% glossary-block:B01209:end
% glossary-block:B01210:begin
\needspace{3\baselineskip}
\textbf{sēo} — nom. f sg.: see \textbf{se, sēo, ðæt}; pres. ind. 1 sg.; pres. subj. sg.: see \textbf{sēon}.\par
% glossary-block:B01210:end
% glossary-block:B01211:begin
\needspace{3\baselineskip}
\textbf{seofon} (cardinal numeral). seven. Cardinal 7.\par
% glossary-block:B01211:end
% glossary-block:B01212:begin
\needspace{3\baselineskip}
\textbf{seofoða} (ordinal). seventh.\par
% glossary-block:B01212:end
% glossary-block:B01213:begin
\needspace{3\baselineskip}
\textbf{seolfor} (n. n.); gen. sg. seolfres; nom. pl. seolfor. silver. Metal noun.\par
% glossary-block:B01213:end
% glossary-block:B01214:begin
\needspace{3\baselineskip}
\textbf{seolfren} (adj.). silver; made of silver. Material adjective, distinct from the noun seolfor.\par
% glossary-block:B01214:end
% glossary-block:B01215:begin
\needspace{3\baselineskip}
\textbf{seolh} (n. m.); gen. sg. sēoles; nom. pl. sēolas. seal.\par
% glossary-block:B01215:end
% glossary-block:B01216:begin
\needspace{3\baselineskip}
\textbf{sēon} (v. strong); pret. 1 sg. seah; pret. pl. sāwon; past ptcp. sewen; pres. 3 sg. siehð. see.\par
% glossary-block:B01216:end
% glossary-block:B01217:begin
\needspace{3\baselineskip}
\textbf{sēoðan} (v. strong); pret. 1 sg. sēað; pret. pl. sudon; past ptcp. soden; pres. 3 sg. sīeð. boil.\par
% glossary-block:B01217:end
% glossary-block:B01218:begin
\needspace{3\baselineskip}
\textbf{se, sēo, ðæt} (article; demonstrative; relative pron.). the. that; those. those. Supplied substantival dative plural referring to waters. what; that which. who; which; that. Inflected relative; case follows its clause role. so that. Only in supplied tō ðon ðætte result connection. after that. Only in temporal æfter ðon. so; to that degree. Only in supplied tō ðan … ðæt.\par
% glossary-block:B01218:end
% glossary-block:B01219:begin
\needspace{3\baselineskip}
\textbf{setl} (n. n.); gen. sg. setles; nom. pl. setl. seat; throne.\par
% glossary-block:B01219:end
% glossary-block:B01220:begin
\needspace{3\baselineskip}
\textbf{sęttan} (v. weak); pret. 1 sg. sętte; past ptcp. (ge)sętt. set; place. appoint; establish. With accusative entity/decree in the assigned frame.\par
% glossary-block:B01220:end
% glossary-block:B01221:begin
\needspace{3\baselineskip}
\textbf{se-ðe} (relative). which; that. Supplied se-ðe relative complex referring to the storm.\par
% glossary-block:B01221:end
% glossary-block:B01222:begin
\needspace{3\baselineskip}
\textbf{sibb} (n. f.); gen. sg. sibbe; nom. pl. sibba. peace.\par
% glossary-block:B01222:end
% glossary-block:B01223:begin
\needspace{3\baselineskip}
\textbf{sīde} (n. f.); gen. sg. sīdan; nom. pl. sīdan. silk.\par
% glossary-block:B01223:end
% glossary-block:B01224:begin
\needspace{3\baselineskip}
\textbf{sīe} — pres. subj. sg.: see \textbf{wesan}.\par
% glossary-block:B01224:end
% glossary-block:B01225:begin
\needspace{3\baselineskip}
\textbf{siehð} — pres. ind. 3 sg.: see \textbf{sēon}.\par
% glossary-block:B01225:end
% glossary-block:B01226:begin
\needspace{3\baselineskip}
\textbf{sīeð} — pres. ind. 3 sg.: see \textbf{sēoðan}.\par
% glossary-block:B01226:end
% glossary-block:B01227:begin
\needspace{3\baselineskip}
\textbf{siex} (cardinal numeral). six. Cardinal 6.\par
% glossary-block:B01227:end
% glossary-block:B01228:begin
\needspace{3\baselineskip}
\textbf{siexta} (ordinal numeral). sixth.\par
% glossary-block:B01228:end
% glossary-block:B01229:begin
\needspace{3\baselineskip}
\textbf{siextig} (cardinal numeral). sixty. Within the supplied 365-days expression.\par
% glossary-block:B01229:end
% glossary-block:B01230:begin
\needspace{3\baselineskip}
\textbf{sige} (n. m.). victory.\par
% glossary-block:B01230:end
% glossary-block:B01231:begin
\needspace{3\baselineskip}
\textbf{simle} (adv.). always.\par
% glossary-block:B01231:end
% glossary-block:B01232:begin
\needspace{3\baselineskip}
\textbf{sīn} (possessive adjective). his; her; its (own). reflexive possessive; mostly poetic in Wright; reference only.\par
% glossary-block:B01232:end
% glossary-block:B01233:begin
\needspace{3\baselineskip}
\textbf{sind} — pres. ind. pl.: see \textbf{wesan}.\par
% glossary-block:B01233:end
% glossary-block:B01234:begin
\needspace{3\baselineskip}
\textbf{sindon} — pres. ind. pl., variant: see \textbf{wesan}.\par
% glossary-block:B01234:end
% glossary-block:B01235:begin
\needspace{3\baselineskip}
\textbf{singāl} (adj.). continual; constant.\par
% glossary-block:B01235:end
% glossary-block:B01236:begin
\needspace{3\baselineskip}
\textbf{singan} (v. strong); pret. 1 sg. sang; pret. pl. sungon; past ptcp. sungen. sing. Intransitive in the introduced frame. sing. Source genitive plural present participle, of people singing.\par
% glossary-block:B01236:end
% glossary-block:B01237:begin
\needspace{3\baselineskip}
\textbf{sint} — pres. ind. pl., variant: see \textbf{wesan}.\par
% glossary-block:B01237:end
% glossary-block:B01238:begin
\needspace{3\baselineskip}
\textbf{sittan} (v. strong); pret. 1 sg. sæt; pret. pl. sǣton; past ptcp. seten. sit.\par
% glossary-block:B01238:end
% glossary-block:B01239:begin
\needspace{3\baselineskip}
\textbf{sīð} (n. m.); gen. sg. sīðes; nom. pl. sīðas. time; occasion.\par
% glossary-block:B01239:end
% glossary-block:B01240:begin
\needspace{3\baselineskip}
\textbf{sīðfæt} (n. m.); gen. sg. sīðfætes; nom. pl. sīðfatas. journey; passage.\par
% glossary-block:B01240:end
% glossary-block:B01241:begin
\needspace{3\baselineskip}
\textbf{siððan} (adv.; conj.). since; after; afterwards. Conjunction since/after versus adverb afterwards.\par
% glossary-block:B01241:end
% glossary-block:B01242:begin
\needspace{3\baselineskip}
\textbf{slēan} (v. strong); pret. 1 sg. slōg; pret. pl. slōgon; past ptcp. slægen; pres. 3 sg. sliehð. strike; smite. strike; smite. Objectless striking in the assigned hammer sentences. slay; kill. Accusative person killed. make a signal. Only in supplied slōg tācen. kill. Accusative person; the bounded negative command uses selected short-vowel imperative sleah. strike. Supplied past source form.\par
% glossary-block:B01242:end
% glossary-block:B01243:begin
\needspace{3\baselineskip}
\textbf{slęcg} (n. f.); gen. sg. slęcge; nom. pl. slęcga. hammer.\par
% glossary-block:B01243:end
% glossary-block:B01244:begin
\needspace{3\baselineskip}
\textbf{sliehð} — pres. ind. 3 sg.: see \textbf{slēan}.\par
% glossary-block:B01244:end
% glossary-block:B01245:begin
\needspace{3\baselineskip}
\textbf{smēagan} (v.). ponder; consider. Supplied preterite source form.\par
% glossary-block:B01245:end
% glossary-block:B01246:begin
\needspace{3\baselineskip}
\textbf{smierwan} (v. weak); pret. 1 sg. smierede; past ptcp. smiered. anoint. Accusative body part or person, mid + dative material.\par
% glossary-block:B01246:end
% glossary-block:B01247:begin
\needspace{3\baselineskip}
\textbf{smið} (n. m.); gen. sg. smiðes; nom. pl. smiðas. blacksmith.\par
% glossary-block:B01247:end
% glossary-block:B01248:begin
\needspace{3\baselineskip}
\textbf{smiððe} (n. f.); gen. sg. smiððan; nom. pl. smiððan. smithy.\par
% glossary-block:B01248:end
% glossary-block:B01249:begin
\needspace{3\baselineskip}
\textbf{smylte} (adj.); supplied superl. acc. m. sg. smyltestan. calm; smooth.\par
% glossary-block:B01249:end
% glossary-block:B01250:begin
\needspace{3\baselineskip}
\textbf{smyltest} (adj.). very calm; calmest. Supplied superlative adjective describing the sea. Superlative of smylte; source form smyltestan.\par
% glossary-block:B01250:end
% glossary-block:B01251:begin
\needspace{3\baselineskip}
\textbf{smyltnes} (n. f.); gen. sg. smyltnesse. calm; calmness.\par
% glossary-block:B01251:end
% glossary-block:B01252:begin
\needspace{3\baselineskip}
\textbf{snāw} (n. m.). snow.\par
% glossary-block:B01252:end
% glossary-block:B01253:begin
\needspace{3\baselineskip}
\textbf{snelnes} (n. f.). skill; agility.\par
% glossary-block:B01253:end
% glossary-block:B01254:begin
\needspace{3\baselineskip}
\textbf{snīðan} (v. strong); pret. 1 sg. snāð; pret. pl. snidon; past ptcp. sniden. cut.\par
% glossary-block:B01254:end
% glossary-block:B01255:begin
\needspace{3\baselineskip}
\textbf{sōhte} — pret. sg.: see \textbf{sēcan}.\par
% glossary-block:B01255:end
% glossary-block:B01256:begin
\needspace{3\baselineskip}
\textbf{sōna} (adv.). as soon as. In sōna swā. immediately; at once.\par
% glossary-block:B01256:end
% glossary-block:B01257:begin
\needspace{3\baselineskip}
\textbf{sorg} (n. f.). distress; anxiety; trouble; sorrow.\par
% glossary-block:B01257:end
% glossary-block:B01258:begin
\needspace{3\baselineskip}
\textbf{sōð} (adj.). true.\par
% glossary-block:B01258:end
% glossary-block:B01259:begin
\needspace{3\baselineskip}
\textbf{sōðfæstnes} (n. f.); gen. sg. sōðfæstnesse; nom. pl. sōðfæstnessa. truth.\par
% glossary-block:B01259:end
% glossary-block:B01260:begin
\needspace{3\baselineskip}
\textbf{sōðlice} (adv.). now; but. Discourse connective corresponding to source Latin autem, distinct from truly. truly; indeed.\par
% glossary-block:B01260:end
% glossary-block:B01261:begin
\needspace{3\baselineskip}
\textbf{spǣtan} (v. weak); pret. 1 sg. spǣtte; past ptcp. spǣtt. spit.\par
% glossary-block:B01261:end
% glossary-block:B01262:begin
\needspace{3\baselineskip}
\textbf{spell} (n. n.). account; discourse.\par
% glossary-block:B01262:end
% glossary-block:B01263:begin
\needspace{3\baselineskip}
\textbf{sprǣc} (n. f.); gen. sg. sprǣce; nom. pl. sprǣca. saying; speech.\par
% glossary-block:B01263:end
% glossary-block:B01264:begin
\needspace{3\baselineskip}
\textbf{sprecan} (v. strong); pret. 1 sg. spræc; pret. pl. sprǣcon; past ptcp. sprecen; pres. 3 sg. spricð. speak. speak with. mid + dative interlocutor in the taught production frame.\par
% glossary-block:B01264:end
% glossary-block:B01265:begin
\needspace{3\baselineskip}
\textbf{spricð} — pres. ind. 3 sg.: see \textbf{sprecan}.\par
% glossary-block:B01265:end
% glossary-block:B01266:begin
\needspace{3\baselineskip}
\textbf{stān} (n. m.); gen. sg. stānes; nom. pl. stānas. stone.\par
% glossary-block:B01266:end
% glossary-block:B01267:begin
\needspace{3\baselineskip}
\textbf{standan} (v. strong); pret. 1 sg. stōd; pret. pl. stōdon; past ptcp. standen; pres. 3 sg. stęnt. oppose; stand against. Supplied standan ongēan frame, distinct from simple stand. stand.\par
% glossary-block:B01267:end
% glossary-block:B01268:begin
\needspace{3\baselineskip}
\textbf{stānen} (adj.). stone; made of stone. Material adjective.\par
% glossary-block:B01268:end
% glossary-block:B01269:begin
\needspace{3\baselineskip}
\textbf{stānweall} (n. m.); gen. sg. stānwealles; nom. pl. stānweallas. stone wall.\par
% glossary-block:B01269:end
% glossary-block:B01270:begin
\needspace{3\baselineskip}
\textbf{stefn} (n. f.); gen. sg. stefne; nom. pl. stefna. voice.\par
% glossary-block:B01270:end
% glossary-block:B01271:begin
\needspace{3\baselineskip}
\textbf{stelan} (v.). steal.\par
% glossary-block:B01271:end
% glossary-block:B01272:begin
\needspace{3\baselineskip}
\textbf{stęnt} — pres. ind. 3 sg.: see \textbf{standan}.\par
% glossary-block:B01272:end
% glossary-block:B01273:begin
\needspace{3\baselineskip}
\textbf{steorra} (n. m.); gen. sg. steorran; nom. pl. steorran. star.\par
% glossary-block:B01273:end
% glossary-block:B01274:begin
\needspace{3\baselineskip}
\textbf{stille} (adj.). still; quiet.\par
% glossary-block:B01274:end
% glossary-block:B01275:begin
\needspace{3\baselineskip}
\textbf{stilnes} (n. f.). stillness; calm. Supplied genitive singular. stillness; quietness.\par
% glossary-block:B01275:end
% glossary-block:B01276:begin
\needspace{3\baselineskip}
\textbf{storm} (n. m.). storm.\par
% glossary-block:B01276:end
% glossary-block:B01277:begin
\needspace{3\baselineskip}
\textbf{strand} (n. m.); gen. sg. strandes; nom. pl. strandas. strand; sea-shore.\par
% glossary-block:B01277:end
% glossary-block:B01278:begin
\needspace{3\baselineskip}
\textbf{strang} (adj.). strong. Masculine singular predicate in the source. strong. Attributive adjective with wind in the source.\par
% glossary-block:B01278:end
% glossary-block:B01279:begin
\needspace{3\baselineskip}
\textbf{strǣt} (n. f.); gen. sg. strǣte; nom. pl. strǣta. street. public place. Street used as a gathering place in the assigned prompt.\par
% glossary-block:B01279:end
% glossary-block:B01280:begin
\needspace{3\baselineskip}
\textbf{strēaw} (n. n.). straw.\par
% glossary-block:B01280:end
% glossary-block:B01281:begin
\needspace{3\baselineskip}
\textbf{stręngðu} (n. f.). strength.\par
% glossary-block:B01281:end
% glossary-block:B01282:begin
\needspace{3\baselineskip}
\textbf{strengu} (n. f.); gen. sg. strenge; nom. pl. strenga. strength. Recognition class-reference noun; preserve the source spelling and do not infer a plural count meaning.\par
% glossary-block:B01282:end
% glossary-block:B01283:begin
\needspace{3\baselineskip}
\textbf{styccemǣlum} (adv.). little by little.\par
% glossary-block:B01283:end
% glossary-block:B01284:begin
\needspace{3\baselineskip}
\textbf{styriendlic} (adj.). moving.\par
% glossary-block:B01284:end
% glossary-block:B01285:begin
\needspace{3\baselineskip}
\textbf{sulhscear} (n. m.); gen. sg. sulhsceares; nom. pl. sulhscearas. plowshare.\par
% glossary-block:B01285:end
% glossary-block:B01286:begin
\needspace{3\baselineskip}
\textbf{sum} (indefinite adjective/pronoun). a certain; some. some. Members selected from a group; genitive whole.\par
% glossary-block:B01286:end
% glossary-block:B01287:begin
\needspace{3\baselineskip}
\textbf{sumer} (n. m.); gen. sg. sumeres; nom. pl. sumeras. summer.\par
% glossary-block:B01287:end
% glossary-block:B01288:begin
\needspace{3\baselineskip}
\textbf{sund} (n.). swimming.\par
% glossary-block:B01288:end
% glossary-block:B01289:begin
\needspace{3\baselineskip}
\textbf{sunne} (n. f.); gen. sg. sunnan; nom. pl. sunnan. sun.\par
% glossary-block:B01289:end
% glossary-block:B01290:begin
\needspace{3\baselineskip}
\textbf{sunu} (n. m.); gen. sg. suna; nom. pl. suna. son.\par
% glossary-block:B01290:end
% glossary-block:B01291:begin
\needspace{3\baselineskip}
\textbf{sūðan} (adv.). in the south; from the south. Retain the supplied spatial expression.\par
% glossary-block:B01291:end
% glossary-block:B01292:begin
\needspace{3\baselineskip}
\textbf{sūðerne} (adj.). southern.\par
% glossary-block:B01292:end
% glossary-block:B01293:begin
\needspace{3\baselineskip}
\textbf{sūðwesterne} (adj.). southwestern.\par
% glossary-block:B01293:end
% glossary-block:B01294:begin
\needspace{3\baselineskip}
\textbf{swā} (adverb; conjunction). so; thus. as. Only in sōna swā. as long as; inasmuch as. Whole swā lange swā phrase; the assigned biblical expression need not imply duration. as; in the same way as. as. First degree-marker in equal comparison. as. Introduces the comparison standard. as; like. Predicative comparison. the … the. Proportional comparison, not the definite article. whoever; whatever. swā hwā/hwæt/hwæðer swā; the internal pronoun's case and type must fit the clause.\par
% glossary-block:B01294:end
% glossary-block:B01295:begin
\needspace{3\baselineskip}
\textbf{swā-hwæt-swā} (pron.). whatsoever; whatever.\par
% glossary-block:B01295:end
% glossary-block:B01296:begin
\needspace{3\baselineskip}
\textbf{swā-swā} (conj.). just as. about; approximately. Approximation in the supplied source expression.\par
% glossary-block:B01296:end
% glossary-block:B01297:begin
\needspace{3\baselineskip}
\textbf{swēg} (n. m.). music; sound.\par
% glossary-block:B01297:end
% glossary-block:B01298:begin
\needspace{3\baselineskip}
\textbf{swēgcræft} (n. m.). musical skill.\par
% glossary-block:B01298:end
% glossary-block:B01299:begin
\needspace{3\baselineskip}
\textbf{sweostor} (n. f.); gen. sg. sweostor; nom. pl. sweostor. sister.\par
% glossary-block:B01299:end
% glossary-block:B01300:begin
\needspace{3\baselineskip}
\textbf{sweotole} (adv.); superl. sweotolost. In superlative sweotolost: most clearly.\par
% glossary-block:B01300:end
% glossary-block:B01301:begin
\needspace{3\baselineskip}
\textbf{sweotollice} (adv.). clearly.\par
% glossary-block:B01301:end
% glossary-block:B01302:begin
\needspace{3\baselineskip}
\textbf{sweotolost} — superlative adverb; exact source spelling: see \textbf{sweotole}.\par
% glossary-block:B01302:end
% glossary-block:B01303:begin
\needspace{3\baselineskip}
\textbf{swęrian} (v. strong); pret. 1 sg. swōr; pret. pl. swōron; past ptcp. sworen. swear.\par
% glossary-block:B01303:end
% glossary-block:B01304:begin
\needspace{3\baselineskip}
\textbf{swēte} (adj.); comp. swētra; superl. swētest. sweet.\par
% glossary-block:B01304:end
% glossary-block:B01305:begin
\needspace{3\baselineskip}
\textbf{swētest} (adj.). sweetest. Superlative of swēte.\par
% glossary-block:B01305:end
% glossary-block:B01306:begin
\needspace{3\baselineskip}
\textbf{swētnes} (n. f.). sweetness.\par
% glossary-block:B01306:end
% glossary-block:B01307:begin
\needspace{3\baselineskip}
\textbf{swētra} (adj.). sweeter. Comparative of swēte.\par
% glossary-block:B01307:end
% glossary-block:B01308:begin
\needspace{3\baselineskip}
\textbf{swift} (adj.). swift.\par
% glossary-block:B01308:end
% glossary-block:B01309:begin
\needspace{3\baselineskip}
\textbf{swiftnes} (n. f.). swiftness.\par
% glossary-block:B01309:end
% glossary-block:B01310:begin
\needspace{3\baselineskip}
\textbf{swīge} (n. f.). silence.\par
% glossary-block:B01310:end
% glossary-block:B01311:begin
\needspace{3\baselineskip}
\textbf{swigian} (v.). be silent. Supplied source sense.\par
% glossary-block:B01311:end
% glossary-block:B01312:begin
\needspace{3\baselineskip}
\textbf{swilc} (indefinite adjective). such.\par
% glossary-block:B01312:end
% glossary-block:B01313:begin
\needspace{3\baselineskip}
\textbf{swilce} (adv.). likewise. Adverb, distinct from inflected swilc such. about; approximately. as if. Supplied apparent-rejuvenation comparison.\par
% glossary-block:B01313:end
% glossary-block:B01314:begin
\needspace{3\baselineskip}
\textbf{swimman} (v. strong); pret. 1 sg. swamm; pret. pl. swummon; past ptcp. swummen. swim.\par
% glossary-block:B01314:end
% glossary-block:B01315:begin
\needspace{3\baselineskip}
\textbf{swingan} (v.). swing; brandish. Supplied past source form.\par
% glossary-block:B01315:end
% glossary-block:B01316:begin
\needspace{3\baselineskip}
\textbf{swīð} (adjective citation); comp. swīðra. strong.\par
% glossary-block:B01316:end
% glossary-block:B01317:begin
\needspace{3\baselineskip}
\textbf{swīðe} (adv.); comp. swīðor; superl. swīðost. very. rather; more. Preferential comparative swīðor. most; very greatly. Superlative swīðost; recognition reference, not independent production. more. Comparative degree of an action/feeling, not a count of objects.\par
% glossary-block:B01317:end
% glossary-block:B01318:begin
\needspace{3\baselineskip}
\textbf{swīðlice} (adv.); comp. swīðlicor. In comparative swīðlicor: more violently; more severely.\par
% glossary-block:B01318:end
% glossary-block:B01319:begin
\needspace{3\baselineskip}
\textbf{swīðlicor} — comparative adverb; exact source spelling: see \textbf{swīðlice}.\par
% glossary-block:B01319:end
% glossary-block:B01320:begin
\needspace{3\baselineskip}
\textbf{swīðor} — comparative adverb: see \textbf{swīðe}.\par
% glossary-block:B01320:end
% glossary-block:B01321:begin
\needspace{3\baselineskip}
\textbf{swīðost} — superlative adverb: see \textbf{swīðe}.\par
% glossary-block:B01321:end
% glossary-block:B01322:begin
\needspace{3\baselineskip}
\textbf{swīðra} (adj.). right. Right-hand/side expression, distinct from stronger. stronger; fiercer. Comparative degree, distinct from right-side sense. In stronger/fiercer, comparative of swīð; right-hand/side is separate.\par
% glossary-block:B01322:end
% glossary-block:B01323:begin
\needspace{3\baselineskip}
\textbf{synfull} (adj.). sinful.\par
% glossary-block:B01323:end
% glossary-block:B01324:begin
\needspace{3\baselineskip}
\textbf{synn} (n. f.); gen. sg. synne; nom. pl. synna. sin.\par
% glossary-block:B01324:end
% book-navigation: standalone bookmark replaced by volume hierarchy
% glossary-block:B01325:begin
\needspace{5\baselineskip}
\subsection*{T}\BookHeading{nav-glossary-old-english-0020}{subsubsection}{T}{0}
% glossary-block:B01325:end
% glossary-block:B01326:begin
\textbf{tabele} (n. f.); gen. sg. tabelan; nom. pl. tabelan. tablet; table. Tables of the law, not an unrestricted furniture meaning.\par
% glossary-block:B01326:end
% glossary-block:B01327:begin
\needspace{3\baselineskip}
\textbf{tācen} (n. n.); gen. sg. tācnes; nom. pl. tācnu. sign.\par
% glossary-block:B01327:end
% glossary-block:B01328:begin
\needspace{3\baselineskip}
\textbf{tǣcan} (v.). teach; show.\par
% glossary-block:B01328:end
% glossary-block:B01329:begin
\needspace{3\baselineskip}
\textbf{tǣlan} (v. weak); pret. 1 sg. tǣlde; past ptcp. tǣled. blame; censure.\par
% glossary-block:B01329:end
% glossary-block:B01330:begin
\needspace{3\baselineskip}
\textbf{tēar} (n. m.); gen. sg. tēares; nom. pl. tēaras. tear.\par
% glossary-block:B01330:end
% glossary-block:B01331:begin
\needspace{3\baselineskip}
\textbf{tempel} (n. n.); gen. sg. temples; nom. pl. templu. temple.\par
% glossary-block:B01331:end
% glossary-block:B01332:begin
\needspace{3\baselineskip}
\textbf{tēon} (v. strong); pret. 1 sg. tēah; pret. pl. tugon; past ptcp. togen; pres. 3 sg. tīehð. draw; pull. bring forth; produce. Supplied third-plural present subjunctive, with waters as subject. bring forth; produce. Accusative thing produced with particle forð; use the assigned tēon forms.\par
% glossary-block:B01332:end
% glossary-block:B01333:begin
\needspace{3\baselineskip}
\textbf{tēð} — dat. sg.; nom./acc. pl.: see \textbf{tōð}.\par
% glossary-block:B01333:end
% glossary-block:B01334:begin
\needspace{3\baselineskip}
\textbf{tīd} (n. f.); gen. sg. tīde; nom. pl. tīda. season; time. hour. Time measure in the supplied intervals.\par
% glossary-block:B01334:end
% glossary-block:B01335:begin
\needspace{3\baselineskip}
\textbf{tīehð} — pres. ind. 3 sg.: see \textbf{tēon}.\par
% glossary-block:B01335:end
% glossary-block:B01336:begin
\needspace{3\baselineskip}
\textbf{tīen} (cardinal numeral). ten. Cardinal 10.\par
% glossary-block:B01336:end
% glossary-block:B01337:begin
\needspace{3\baselineskip}
\textbf{tilian} (v.). obtain; provide. Genitive thing procured; dative reflexive beneficiary in the assigned frame.\par
% glossary-block:B01337:end
% glossary-block:B01338:begin
\needspace{3\baselineskip}
\textbf{tin} (n. n.); gen. sg. tines; nom. pl. tinu. tin.\par
% glossary-block:B01338:end
% glossary-block:B01339:begin
\needspace{3\baselineskip}
\textbf{tō} (prep.). to; in order to. With the inflected infinitive; ordinary modal complements use a plain infinitive. for; to serve as. tō + dative function; a separate dative beneficiary may accompany it. to. Endpoint/result death. to. Dative destination. at; in response to. Only the exact at-the-prayers expression. to; toward. Person to whom knees/prayer are directed in the supplied phrase. so that. Only in supplied tō ðon ðætte result connection. for aid; for protection. Whole supplied purpose expressions. to. Postposed in the supplied cwæð him tō speech phrase. at; to. Bounded supplied root expression with genitive dependent and dative root. to. The person addressed in speech. to; into. Result in supplied tō lore, tō sorge or tō nāhte. why; for what purpose. Only in tō hwon. into; to. Resulting language/literary form. toward. Conduct toward a person in the supplied expression. in accordance with. A will as the goal of guidance in the supplied prayer. for; to the benefit of. for ever. Whole fixed tō worulde. for. Care of a body part in the supplied expression. as wife. Only the supplied tō wīfe role phrase. to. Complement of belonging in the supplied frame. to. Resulting condition of youth in the supplied comparison. so; to that degree. Whole supplied tō ðan … ðæt result construction. toward. Target of a look in the supplied frame.\par
% glossary-block:B01339:end
% glossary-block:B01340:begin
\needspace{3\baselineskip}
\textbf{tōberstan} (v. strong); pret. 1 sg. tōbærst; pret. pl. tōburston; past ptcp. tōborsten; pres. 3 sg. tōbierst. break up; go to pieces.\par
% glossary-block:B01340:end
% glossary-block:B01341:begin
\needspace{3\baselineskip}
\textbf{tōbierst} — pres. ind. 3 sg.: see \textbf{tōberstan}.\par
% glossary-block:B01341:end
% glossary-block:B01342:begin
\needspace{3\baselineskip}
\textbf{tōcuman} (v.). come; arrive.\par
% glossary-block:B01342:end
% glossary-block:B01343:begin
\needspace{3\baselineskip}
\textbf{tō-dæg} (adv.). today.\par
% glossary-block:B01343:end
% glossary-block:B01344:begin
\needspace{3\baselineskip}
\textbf{tōdæg} (adv.). today.\par
% glossary-block:B01344:end
% glossary-block:B01345:begin
\needspace{3\baselineskip}
\textbf{tōdǣlen / tōdǣlden (source forms)} (verb; supplied inflected forms). separate; distinguish. Source forms tōdǣlen/tōdǣlden are retained; no silently invented citation spelling. Only the cited source forms are supplied here.\par
% glossary-block:B01345:end
% glossary-block:B01346:begin
\needspace{3\baselineskip}
\textbf{tōforan} (prep.). before; above all.\par
% glossary-block:B01346:end
% glossary-block:B01347:begin
\needspace{3\baselineskip}
\textbf{tōgædere} (adv.). together.\par
% glossary-block:B01347:end
% glossary-block:B01348:begin
\needspace{3\baselineskip}
\textbf{tōgēanes} (postposition). towards; to meet.\par
% glossary-block:B01348:end
% glossary-block:B01349:begin
\needspace{3\baselineskip}
\textbf{tōgēiecan} (v.). add.\par
% glossary-block:B01349:end
% glossary-block:B01350:begin
\needspace{3\baselineskip}
\textbf{tōgelǣdan} (v. weak); pret. 1 sg. tōgelǣdde; past ptcp. tōgelǣded. bring. Accusative thing brought, dative recipient.\par
% glossary-block:B01350:end
% glossary-block:B01351:begin
\needspace{3\baselineskip}
\textbf{tōhopa} (n. m.). hope.\par
% glossary-block:B01351:end
% glossary-block:B01352:begin
\needspace{3\baselineskip}
\textbf{tōmiddes} (prep.). in the midst of. With dative landmark.\par
% glossary-block:B01352:end
% glossary-block:B01353:begin
\needspace{3\baselineskip}
\textbf{tō-morgen} (adv.). tomorrow.\par
% glossary-block:B01353:end
% glossary-block:B01354:begin
\needspace{3\baselineskip}
\textbf{top} (n. m.). ball?. Tentative source interpretation; retain the question mark and uncertainty.\par
% glossary-block:B01354:end
% glossary-block:B01355:begin
\needspace{3\baselineskip}
\textbf{torr} (n. m.); gen. sg. torres; nom. pl. torras. tower. crag. Probable contextual sense in Cook's note, distinct from earlier tower.\par
% glossary-block:B01355:end
% glossary-block:B01356:begin
\needspace{3\baselineskip}
\textbf{tōscǣt} — pres. ind. 3 sg.: see \textbf{tōsceadan}.\par
% glossary-block:B01356:end
% glossary-block:B01357:begin
\needspace{3\baselineskip}
\textbf{tōsceadan} (v.); pres. 3 sg. tōscǣt. separate; divide.\par
% glossary-block:B01357:end
% glossary-block:B01358:begin
\needspace{3\baselineskip}
\textbf{tō-slītan} (v.). tear. Supplied past source form.\par
% glossary-block:B01358:end
% glossary-block:B01359:begin
\needspace{3\baselineskip}
\textbf{tōtwǣman} (v. weak); pret. 1 sg. tōtwǣmde; past ptcp. tōtwǣmed. divide; separate. Waters divided from waters in the supplied frame. separate. Accusative beings separated; fram + dative other party.\par
% glossary-block:B01359:end
% glossary-block:B01360:begin
\needspace{3\baselineskip}
\textbf{tōð} (n. m.); gen. sg. tōðes; nom. pl. tēð. tooth.\par
% glossary-block:B01360:end
% glossary-block:B01361:begin
\needspace{3\baselineskip}
\textbf{tōweorpan} (v. strong); pret. 1 sg. tōwearp; pret. pl. tōwurpon; past ptcp. tōworpen; pres. 3 sg. tōwierpð. break; overthrow. Break a commandment in the assigned use.\par
% glossary-block:B01361:end
% glossary-block:B01362:begin
\needspace{3\baselineskip}
\textbf{tōwierpð} — pres. ind. 3 sg.: see \textbf{tōweorpan}.\par
% glossary-block:B01362:end
% glossary-block:B01363:begin
\needspace{3\baselineskip}
\textbf{trēo} (n. n.); gen. sg. trēowes; nom. pl. trēowu. tree.\par
% glossary-block:B01363:end
% glossary-block:B01364:begin
\needspace{3\baselineskip}
\textbf{trēownes} (n. f.). trust; confidence.\par
% glossary-block:B01364:end
% glossary-block:B01365:begin
\needspace{3\baselineskip}
\textbf{trēowwyrhta} (n. m.); gen. sg. trēowwyrhtan; nom. pl. trēowwyrhtan. carpenter.\par
% glossary-block:B01365:end
% glossary-block:B01366:begin
\needspace{3\baselineskip}
\textbf{tū} — nom./acc. neut. pl.: see \textbf{twēgen}.\par
% glossary-block:B01366:end
% glossary-block:B01367:begin
\needspace{3\baselineskip}
\textbf{twā} — nom./acc. fem./neut. pl.: see \textbf{twēgen}.\par
% glossary-block:B01367:end
% glossary-block:B01368:begin
\needspace{3\baselineskip}
\textbf{twām} — dat. pl.: see \textbf{twēgen}.\par
% glossary-block:B01368:end
% glossary-block:B01369:begin
\needspace{3\baselineskip}
\textbf{twǣm} — dat. pl., variant: see \textbf{twēgen}.\par
% glossary-block:B01369:end
% glossary-block:B01370:begin
\needspace{3\baselineskip}
\textbf{twēgen} (cardinal numeral). two. two; in two. Two parts in the supplied division expression.\par
% glossary-block:B01370:end
% glossary-block:B01371:begin
\needspace{3\baselineskip}
\textbf{twęlf} (cardinal numeral). twelve. Cardinal 12.\par
% glossary-block:B01371:end
% glossary-block:B01372:begin
\needspace{3\baselineskip}
\textbf{twēntig} (numeral). twenty.\par
% glossary-block:B01372:end
% glossary-block:B01373:begin
\needspace{3\baselineskip}
\textbf{twīfeald} (adj.). double; twofold.\par
% glossary-block:B01373:end
% book-navigation: standalone bookmark replaced by volume hierarchy
% glossary-block:B01374:begin
\needspace{5\baselineskip}
\subsection*{Þ / Ð}\BookHeading{nav-glossary-old-english-0021}{subsubsection}{Þ / Ð}{0}
% glossary-block:B01374:end
% glossary-block:B01375:begin
\textbf{ðā} (adv.). then. Also acc. f sg.; nom. pl.; acc. pl.: see \textbf{se, sēo, ðæt}.\par
% glossary-block:B01375:end
% glossary-block:B01376:begin
\needspace{3\baselineskip}
\textbf{ðā} (conj.). when; factual past time. Also acc. f sg.; nom. pl.; acc. pl.: see \textbf{se, sēo, ðæt}.\par
% glossary-block:B01376:end
% glossary-block:B01377:begin
\needspace{3\baselineskip}
\textbf{ðā-hwæðre} (conj.). nevertheless; yet. Whole supplied connective.\par
% glossary-block:B01377:end
% glossary-block:B01378:begin
\needspace{3\baselineskip}
\textbf{ðā-hwīle-ðe} (conj.). while.\par
% glossary-block:B01378:end
% glossary-block:B01379:begin
\needspace{3\baselineskip}
\textbf{ðām} — dat. m./n. sg.; dat. pl., variant: see \textbf{se, sēo, ðæt}.\par
% glossary-block:B01379:end
% glossary-block:B01380:begin
\needspace{3\baselineskip}
\textbf{ðan} — in supplied tō ðan, to that degree: see \textbf{se, sēo, ðæt}.\par
% glossary-block:B01380:end
% glossary-block:B01381:begin
\needspace{3\baselineskip}
\textbf{ðanan} (adv.). thence; from there.\par
% glossary-block:B01381:end
% glossary-block:B01382:begin
\needspace{3\baselineskip}
\textbf{ðancung} (n. f.); gen. sg. ðancunge. thanksgiving.\par
% glossary-block:B01382:end
% glossary-block:B01383:begin
\needspace{3\baselineskip}
\textbf{ðāra} — gen. pl.: see \textbf{se, sēo, ðæt}.\par
% glossary-block:B01383:end
% glossary-block:B01384:begin
\needspace{3\baselineskip}
\textbf{ðās} — acc. f sg.; nom. pl.; acc. pl.: see \textbf{ðes, ðēos, ðis}.\par
% glossary-block:B01384:end
% glossary-block:B01385:begin
\needspace{3\baselineskip}
\textbf{ðæce} (n. f.). roof.\par
% glossary-block:B01385:end
% glossary-block:B01386:begin
\needspace{3\baselineskip}
\textbf{ðǣm} — dat. m sg.; dat. n sg.; dat. pl.: see \textbf{se, sēo, ðæt}.\par
% glossary-block:B01386:end
% glossary-block:B01387:begin
\needspace{3\baselineskip}
\textbf{ðǣr} (adv.). there.\par
% glossary-block:B01387:end
% glossary-block:B01388:begin
\needspace{3\baselineskip}
\textbf{ðǣra} — gen. pl., variant: see \textbf{se, sēo, ðæt}.\par
% glossary-block:B01388:end
% glossary-block:B01389:begin
\needspace{3\baselineskip}
\textbf{ðǣre} — gen. f sg.; dat. f sg.: see \textbf{se, sēo, ðæt}.\par
% glossary-block:B01389:end
% glossary-block:B01390:begin
\needspace{3\baselineskip}
\textbf{ðǣron} (adverb or particle). therein; in it.\par
% glossary-block:B01390:end
% glossary-block:B01391:begin
\needspace{3\baselineskip}
\textbf{ðǣr-tō-ēacan} (adverbial expression). in addition to that; besides.\par
% glossary-block:B01391:end
% glossary-block:B01392:begin
\needspace{3\baselineskip}
\textbf{ðæs} — gen. m sg.; gen. n sg.: see \textbf{se, sēo, ðæt}.\par
% glossary-block:B01392:end
% glossary-block:B01393:begin
\needspace{3\baselineskip}
\textbf{ðæs-ðe} (conj.). as soon as. Complete source construction with sōna. as. Only in supplied ðæs-ðe mē ðyncð.\par
% glossary-block:B01393:end
% glossary-block:B01394:begin
\needspace{3\baselineskip}
\textbf{ðæt} (conj.). that. Content conjunction, distinct from demonstrative/relative that. so that; in order that. Purpose conjunction with the assigned subjunctive construction. that; so that. Supplied expanded spelling; distinguish content from result by the source clause. that; so that. Asserted result, distinct from intended purpose. that. Cook leaves the source linkage uncertain; do not impose one content/result classification. Also nom. n sg.; acc. n sg.: see \textbf{se, sēo, ðæt}.\par
% glossary-block:B01394:end
% glossary-block:B01395:begin
\needspace{3\baselineskip}
\textbf{ðæt-ðæt} (pron.). that which; what.\par
% glossary-block:B01395:end
% glossary-block:B01396:begin
\needspace{3\baselineskip}
\textbf{ðe} (relative particle; clause marker). that; which; who. which; that. Subject-relative referring to the waters. after. Only as the clause marker in the whole supplied æfter ðǣm ðe expression. clause marker. Only in supplied ðēah ðe, although; ðe alone is not an equivalent of although.\par
% glossary-block:B01396:end
% glossary-block:B01397:begin
\needspace{3\baselineskip}
\textbf{ðē} — dat. sg.; acc. sg.: see \textbf{ðū}.\par
% glossary-block:B01397:end
% glossary-block:B01398:begin
\needspace{3\baselineskip}
\textbf{ðēah} (conj.; adv.). although; even though. nevertheless; yet. Independent adverb, distinct from concessive conjunction although.\par
% glossary-block:B01398:end
% glossary-block:B01399:begin
\needspace{3\baselineskip}
\textbf{ðearf} (n. f.); gen. sg. ðearfe. benefit; good; need. Source contextual nouns; retain the local expression. Also pres. ind. 1/3 sg.: see \textbf{ðurfan}.\par
% glossary-block:B01399:end
% glossary-block:B01400:begin
\needspace{3\baselineskip}
\textbf{ðearfa} (n. m.); gen. sg. ðearfan; nom. pl. ðearfan. poor person; needy person. Noun referring to a person.\par
% glossary-block:B01400:end
% glossary-block:B01401:begin
\needspace{3\baselineskip}
\textbf{ðearle} (adv.). greatly; exceedingly.\par
% glossary-block:B01401:end
% glossary-block:B01402:begin
\needspace{3\baselineskip}
\textbf{ðearlice} (adv.). violently; fiercely.\par
% glossary-block:B01402:end
% glossary-block:B01403:begin
\needspace{3\baselineskip}
\textbf{ðec} — acc. sg., variant: see \textbf{ðū}.\par
% glossary-block:B01403:end
% glossary-block:B01404:begin
\needspace{3\baselineskip}
\textbf{ðegn} (n. m.); gen. sg. ðegnes; nom. pl. ðegnas. vassal; retainer; thane.\par
% glossary-block:B01404:end
% glossary-block:B01405:begin
\needspace{3\baselineskip}
\textbf{ðęncan} (v. weak); pret. 1 sg. ðōhte; past ptcp. (ge)ðōht. think. Supplied source preterite. think.\par
% glossary-block:B01405:end
% glossary-block:B01406:begin
\needspace{3\baselineskip}
\textbf{ðēnian} (v.). serve; attend. Supplied preterite source form.\par
% glossary-block:B01406:end
% glossary-block:B01407:begin
\needspace{3\baselineskip}
\textbf{ðēnung} (n. f.); gen. sg. ðēnunge; nom. pl. ðēnunga. service; attendance.\par
% glossary-block:B01407:end
% glossary-block:B01408:begin
\needspace{3\baselineskip}
\textbf{ðēod} (n. f.); gen. sg. ðēode; nom. pl. ðēoda. people; nation. region; country.\par
% glossary-block:B01408:end
% glossary-block:B01409:begin
\needspace{3\baselineskip}
\textbf{ðēos} — nom. f sg.: see \textbf{ðes, ðēos, ðis}.\par
% glossary-block:B01409:end
% glossary-block:B01410:begin
\needspace{3\baselineskip}
\textbf{ðēow} (adj.). the unfree; enslaved people. Substantival adjective in the source comparison.\par
% glossary-block:B01410:end
% glossary-block:B01411:begin
\needspace{3\baselineskip}
\textbf{ðēow} (n. m.); gen. sg. ðēowes; nom. pl. ðēowas. servant; enslaved person.\par
% glossary-block:B01411:end
% glossary-block:B01412:begin
\needspace{3\baselineskip}
\textbf{ðēowa} (n. m.); gen. sg. ðēowan; nom. pl. ðēowan. servant.\par
% glossary-block:B01412:end
% glossary-block:B01413:begin
\needspace{3\baselineskip}
\textbf{ðēowa manna} (noun phrase). enslaved men; servants. Whole supplied phrase ðēowa manna.\par
% glossary-block:B01413:end
% glossary-block:B01414:begin
\needspace{3\baselineskip}
\textbf{ðēowdōm} (n. m.); gen. sg. ðēowdōmes; nom. pl. ðēowdōmas. service.\par
% glossary-block:B01414:end
% glossary-block:B01415:begin
\needspace{3\baselineskip}
\textbf{ðēowian} (v.). enslave; reduce to servitude.\par
% glossary-block:B01415:end
% glossary-block:B01416:begin
\needspace{3\baselineskip}
\textbf{ðes, ðēos, ðis} (demonstrative). this; these.\par
% glossary-block:B01416:end
% glossary-block:B01417:begin
\needspace{3\baselineskip}
\textbf{ðider} (adv.). thither; to that place.\par
% glossary-block:B01417:end
% glossary-block:B01418:begin
\needspace{3\baselineskip}
\textbf{ðīestru} (n. n.); nom. pl. ðīestru. darkness. Plural only.\par
% glossary-block:B01418:end
% glossary-block:B01419:begin
\needspace{3\baselineskip}
\textbf{ðīn} (possessive adjective). your; yours. One possessor; agrees with the possessed noun. Also gen. sg.: see \textbf{ðū}.\par
% glossary-block:B01419:end
% glossary-block:B01420:begin
\needspace{3\baselineskip}
\textbf{ðīnen} (n. f.); gen. sg. ðīnenne; nom. pl. ðīnenna. handmaiden.\par
% glossary-block:B01420:end
% glossary-block:B01421:begin
\needspace{3\baselineskip}
\textbf{ðing} (n. n.); gen. sg. ðinges; nom. pl. ðing. thing. Ordinary singular/plural noun. thing; fact; matter. Matter known/discussed, distinct from a generic physical object.\par
% glossary-block:B01421:end
% glossary-block:B01422:begin
\needspace{3\baselineskip}
\textbf{þing} — citation/spelling variant: see \textbf{ðing}.\par
% glossary-block:B01422:end
% glossary-block:B01423:begin
\needspace{3\baselineskip}
\textbf{ðis} — nom. n sg.; acc. n sg.: see \textbf{ðes, ðēos, ðis}.\par
% glossary-block:B01423:end
% glossary-block:B01424:begin
\needspace{3\baselineskip}
\textbf{ðisne} — acc. m sg.: see \textbf{ðes, ðēos, ðis}.\par
% glossary-block:B01424:end
% glossary-block:B01425:begin
\needspace{3\baselineskip}
\textbf{ðissa} — gen. pl.: see \textbf{ðes, ðēos, ðis}.\par
% glossary-block:B01425:end
% glossary-block:B01426:begin
\needspace{3\baselineskip}
\textbf{ðisse} — gen. f sg.; dat. f sg.: see \textbf{ðes, ðēos, ðis}.\par
% glossary-block:B01426:end
% glossary-block:B01427:begin
\needspace{3\baselineskip}
\textbf{ðisses} — gen. m sg.; gen. n sg.: see \textbf{ðes, ðēos, ðis}.\par
% glossary-block:B01427:end
% glossary-block:B01428:begin
\needspace{3\baselineskip}
\textbf{ðissum} — dat. m sg.; dat. n sg.; dat. pl.: see \textbf{ðes, ðēos, ðis}.\par
% glossary-block:B01428:end
% glossary-block:B01429:begin
\needspace{3\baselineskip}
\textbf{ðisum} — dat. n pl.; dat. n sg.: see \textbf{ðes, ðēos, ðis}.\par
% glossary-block:B01429:end
% glossary-block:B01430:begin
\needspace{3\baselineskip}
\textbf{ðōhte} — pret. sg.: see \textbf{ðęncan}.\par
% glossary-block:B01430:end
% glossary-block:B01431:begin
\needspace{3\baselineskip}
\textbf{ðolian} (v. weak); pret. 1 sg. ðolode; past ptcp. ðolod. suffer; endure. Accusative hardship endured.\par
% glossary-block:B01431:end
% glossary-block:B01432:begin
\needspace{3\baselineskip}
\textbf{ðon} — instr. m./n. sg.; also supplied dat. n. sg.: see \textbf{se, sēo, ðæt}.\par
% glossary-block:B01432:end
% glossary-block:B01433:begin
\needspace{3\baselineskip}
\textbf{ðone} — acc. m sg.: see \textbf{se, sēo, ðæt}.\par
% glossary-block:B01433:end
% glossary-block:B01434:begin
\needspace{3\baselineskip}
\textbf{ðonne} (conj.; adv.). than. then. Sequence of future judgment events. then; at that point.\par
% glossary-block:B01434:end
% glossary-block:B01435:begin
\needspace{3\baselineskip}
\textbf{ðonne} (conj.). when; recurring or future time.\par
% glossary-block:B01435:end
% glossary-block:B01436:begin
\needspace{3\baselineskip}
\textbf{þonne} — citation/spelling variant: see \textbf{ðonne}.\par
% glossary-block:B01436:end
% glossary-block:B01437:begin
\needspace{3\baselineskip}
\textbf{ðorfte} — pret. sg.: see \textbf{ðurfan}.\par
% glossary-block:B01437:end
% glossary-block:B01438:begin
\needspace{3\baselineskip}
\textbf{ðorn} (n. m.); gen. sg. ðornes; nom. pl. ðornas. thorn.\par
% glossary-block:B01438:end
% glossary-block:B01439:begin
\needspace{3\baselineskip}
\textbf{ðoðor} (n. m.); gen. sg. ðoðores; nom. pl. ðoðoras. ball.\par
% glossary-block:B01439:end
% glossary-block:B01440:begin
\needspace{3\baselineskip}
\textbf{ðrēan} (v.). chastise; strike.\par
% glossary-block:B01440:end
% glossary-block:B01441:begin
\needspace{3\baselineskip}
\textbf{þrēo} — nom./acc. fem./neut. pl., spelling variant: see \textbf{ðrīe}.\par
% glossary-block:B01441:end
% glossary-block:B01442:begin
\needspace{3\baselineskip}
\textbf{ðrēo} — nom./acc. fem./neut. pl.: see \textbf{ðrīe}.\par
% glossary-block:B01442:end
% glossary-block:B01443:begin
\needspace{3\baselineskip}
\textbf{ðridda} (ordinal numeral). third.\par
% glossary-block:B01443:end
% glossary-block:B01444:begin
\needspace{3\baselineskip}
\textbf{ðrīe} (cardinal numeral). three.\par
% glossary-block:B01444:end
% glossary-block:B01445:begin
\needspace{3\baselineskip}
\textbf{ðrītig} (numeral). thirty.\par
% glossary-block:B01445:end
% glossary-block:B01446:begin
\needspace{3\baselineskip}
\textbf{ðrīwa} (adv.). thrice; three times.\par
% glossary-block:B01446:end
% glossary-block:B01447:begin
\needspace{3\baselineskip}
\textbf{ðrymsetl} (n. n.); gen. sg. ðrymsetles; nom. pl. ðrymsetl. throne; throne of glory.\par
% glossary-block:B01447:end
% glossary-block:B01448:begin
\needspace{3\baselineskip}
\textbf{ðū} (personal pron.). you. Personal singular/dual/plural and governed cases; dual subjects take plural verb forms. of you. Personal genitive citation identifying singular/dual/plural; not by itself an attributive possessive lesson.\par
% glossary-block:B01448:end
% glossary-block:B01449:begin
\needspace{3\baselineskip}
\textbf{ðūhte} — pret. sg.: see \textbf{ðyncan}.\par
% glossary-block:B01449:end
% glossary-block:B01450:begin
\needspace{3\baselineskip}
\textbf{ðurfan} (v. preterite-present); pres. 1/3 sg. ðearf; pres. pl. ðurfon; pret. 1 sg. ðorfte. need to. No past-participle citation is given in the checked source paradigm.\par
% glossary-block:B01450:end
% glossary-block:B01451:begin
\needspace{3\baselineskip}
\textbf{þurh} — citation/spelling variant: see \textbf{ðurh, þurh}.\par
% glossary-block:B01451:end
% glossary-block:B01452:begin
\needspace{3\baselineskip}
\textbf{ðurh, þurh} (prep.). through. Accusative landmark or medium. through; by. Mediated action through a person in the supplied frame. by. Accusative oath standard. through; by means of. Supplied service/evidence expressions, including Ðurh hwæt.\par
% glossary-block:B01452:end
% glossary-block:B01453:begin
\needspace{3\baselineskip}
\textbf{ðurhwunian} (v. weak); pret. 1 sg. ðurhwunode; past ptcp. ðurhwunod. remain; endure; last.\par
% glossary-block:B01453:end
% glossary-block:B01454:begin
\needspace{3\baselineskip}
\textbf{ðurstig} (adj.). thirsty.\par
% glossary-block:B01454:end
% glossary-block:B01455:begin
\needspace{3\baselineskip}
\textbf{ðus} (adv.). thus.\par
% glossary-block:B01455:end
% glossary-block:B01456:begin
\needspace{3\baselineskip}
\textbf{ðūsend} (numeral). thousand. Within the complete supplied high-number expression.\par
% glossary-block:B01456:end
% glossary-block:B01457:begin
\needspace{3\baselineskip}
\textbf{þūsend} — citation/spelling variant: see \textbf{ðūsend}.\par
% glossary-block:B01457:end
% glossary-block:B01458:begin
\needspace{3\baselineskip}
\textbf{ðwēal} (n. n.). bath; bathing place.\par
% glossary-block:B01458:end
% glossary-block:B01459:begin
\needspace{3\baselineskip}
\textbf{ðwēan} (v. strong); pret. 1 sg. ðwōh; pret. pl. ðwōgon; past ptcp. ðwegen; pres. 3 sg. ðwiehð. wash. Accusative body part, mid + dative liquid.\par
% glossary-block:B01459:end
% glossary-block:B01460:begin
\needspace{3\baselineskip}
\textbf{ðwiehð} — pres. ind. 3 sg.: see \textbf{ðwēan}.\par
% glossary-block:B01460:end
% glossary-block:B01461:begin
\needspace{3\baselineskip}
\textbf{ðȳ} — instr. m sg.; instr. n sg.: see \textbf{se, sēo, ðæt}.\par
% glossary-block:B01461:end
% glossary-block:B01462:begin
\needspace{3\baselineskip}
\textbf{ðȳ-lǣs-ðe} (conj.). lest; so that not.\par
% glossary-block:B01462:end
% glossary-block:B01463:begin
\needspace{3\baselineskip}
\textbf{ðyncan} (v. impersonal); pres. 3 sg. ðyncð; pret. 3 sg. ðūhte; past ptcp. geðūht. seem. Dative experiencer; the thing that seems is nominative. seem. Nominative thing that seems, dative experiencer and agreeing predicate. seem. Supplied past source form. Impersonal use: the course supplies third-singular finite forms.\par
% glossary-block:B01463:end
% glossary-block:B01464:begin
\needspace{3\baselineskip}
\textbf{ðyrel} (n. n.); gen. sg. ðyreles; nom. pl. ðyrelu. hole.\par
% glossary-block:B01464:end
% glossary-block:B01465:begin
\needspace{3\baselineskip}
\textbf{ðȳs} — instr. m sg.; instr. n sg.: see \textbf{ðes, ðēos, ðis}.\par
% glossary-block:B01465:end
% book-navigation: standalone bookmark replaced by volume hierarchy
% glossary-block:B01466:begin
\needspace{5\baselineskip}
\subsection*{U}\BookHeading{nav-glossary-old-english-0022}{subsubsection}{U}{0}
% glossary-block:B01466:end
% glossary-block:B01467:begin
\textbf{unāsęcgende} (adj.). indescribable.\par
% glossary-block:B01467:end
% glossary-block:B01468:begin
\needspace{3\baselineskip}
\textbf{unæðele} (adj.). people not noble. Substantival adjective in the supplied contrast.\par
% glossary-block:B01468:end
% glossary-block:B01469:begin
\needspace{3\baselineskip}
\textbf{unc} — dat. dual; acc. dual: see \textbf{ic}.\par
% glossary-block:B01469:end
% glossary-block:B01470:begin
\needspace{3\baselineskip}
\textbf{uncer} (possessive adjective). our; ours. Two possessors; agreeing possessive. Also gen. dual: see \textbf{ic}.\par
% glossary-block:B01470:end
% glossary-block:B01471:begin
\needspace{3\baselineskip}
\textbf{under} (prep.). under. Dative landmark in the taught static frame. under. Directional accusative in the assigned entry-under-the-roof frame.\par
% glossary-block:B01471:end
% glossary-block:B01472:begin
\needspace{3\baselineskip}
\textbf{underfōn} (v.); past ptcp. underfangen. receive; take charge of.\par
% glossary-block:B01472:end
% glossary-block:B01473:begin
\needspace{3\baselineskip}
\textbf{underðēodan} (v.). add; subjoin. Supplied passive/participial addition to speech.\par
% glossary-block:B01473:end
% glossary-block:B01474:begin
\needspace{3\baselineskip}
\textbf{underðēodnes} (n. f.). subjection; submission.\par
% glossary-block:B01474:end
% glossary-block:B01475:begin
\needspace{3\baselineskip}
\textbf{undēð} — pres. ind. 3 sg.: see \textbf{undōn}.\par
% glossary-block:B01475:end
% glossary-block:B01476:begin
\needspace{3\baselineskip}
\textbf{undōn} (v. irregular); pres. 1 sg. undō; pres. 3 sg. undēð; pres. pl. undōð; pret. 1 sg. undyde; pret. pl. undydon; past ptcp. undōn. undo; open.\par
% glossary-block:B01476:end
% glossary-block:B01477:begin
\needspace{3\baselineskip}
\textbf{undyde} — pret. sg.: see \textbf{undōn}.\par
% glossary-block:B01477:end
% glossary-block:B01478:begin
\needspace{3\baselineskip}
\textbf{unēaðe} (adv.). with difficulty; barely.\par
% glossary-block:B01478:end
% glossary-block:B01479:begin
\needspace{3\baselineskip}
\textbf{ungecnāwen} (past participle; supplied predicate); past ptcp. ungecnāwen. unknown.\par
% glossary-block:B01479:end
% glossary-block:B01480:begin
\needspace{3\baselineskip}
\textbf{ungesewenlic} (adj.). invisible.\par
% glossary-block:B01480:end
% glossary-block:B01481:begin
\needspace{3\baselineskip}
\textbf{ungeðancfull} (adj.). unthankful.\par
% glossary-block:B01481:end
% glossary-block:B01482:begin
\needspace{3\baselineskip}
\textbf{ungewunelic} (adj.). unusual; unfamiliar.\par
% glossary-block:B01482:end
% glossary-block:B01483:begin
\needspace{3\baselineskip}
\textbf{unnan} (v. preterite-present); pres. 1/3 sg. ann; pres. pl. unnon; pret. 1 sg. ūðe; past ptcp. (ge)unnen. grant; allow. Dative person, genitive thing granted.\par
% glossary-block:B01483:end
% glossary-block:B01484:begin
\needspace{3\baselineskip}
\textbf{unrihtwīs} (adj.). unrighteous. Adjective, distinct from the substantive unrighteous people. unrighteous people. Substantival adjective.\par
% glossary-block:B01484:end
% glossary-block:B01485:begin
\needspace{3\baselineskip}
\textbf{unrihtwīsnes} (n. f.); gen. sg. unrihtwīsnesse; nom. pl. unrihtwīsnessa. unrighteousness.\par
% glossary-block:B01485:end
% glossary-block:B01486:begin
\needspace{3\baselineskip}
\textbf{unrīm} (n. n.). abundance; uncounted quantity. Supplied quantity with genitive of the substance.\par
% glossary-block:B01486:end
% glossary-block:B01487:begin
\needspace{3\baselineskip}
\textbf{unscrȳdan} (v.). undress; take off. Supplied source form; retain its object/reflexive roles.\par
% glossary-block:B01487:end
% glossary-block:B01488:begin
\needspace{3\baselineskip}
\textbf{unstilnes} (n. f.). unrest; disorder.\par
% glossary-block:B01488:end
% glossary-block:B01489:begin
\needspace{3\baselineskip}
\textbf{untrum} (adj.). sick. sick people. Substantival adjective.\par
% glossary-block:B01489:end
% glossary-block:B01490:begin
\needspace{3\baselineskip}
\textbf{untrymnes} (n. f.); gen. sg. untrymnesse; nom. pl. untrymnessa. infirmity.\par
% glossary-block:B01490:end
% glossary-block:B01491:begin
\needspace{3\baselineskip}
\textbf{ūp} (adv.). up.\par
% glossary-block:B01491:end
% glossary-block:B01492:begin
\needspace{3\baselineskip}
\textbf{ūp-ā-hębban} (v. strong); pret. 1 sg. ūpāhōf; pret. pl. ūpāhōfon; past ptcp. ūpāhafen; pres. 3 sg. ūpāhęfð. lift up. Supplied source form. lift up. Supplied first-plural preterite. raise; exalt.\par
% glossary-block:B01492:end
% glossary-block:B01493:begin
\needspace{3\baselineskip}
\textbf{ūpāhęfð} — pres. ind. 3 sg.: see \textbf{ūp-ā-hębban}.\par
% glossary-block:B01493:end
% glossary-block:B01494:begin
\needspace{3\baselineskip}
\textbf{ūpāstīgan} (v.). rise up; ascend.\par
% glossary-block:B01494:end
% glossary-block:B01495:begin
\needspace{3\baselineskip}
\textbf{ūpēode} — pret. sg.: see \textbf{ūpgān}.\par
% glossary-block:B01495:end
% glossary-block:B01496:begin
\needspace{3\baselineskip}
\textbf{ūpgān} (v.). go up; come up. From the water in the supplied expression.\par
% glossary-block:B01496:end
% glossary-block:B01497:begin
\needspace{3\baselineskip}
\textbf{ūpgang} (n. m.). rising; sunrise.\par
% glossary-block:B01497:end
% glossary-block:B01498:begin
\needspace{3\baselineskip}
\textbf{upon} (prep.). on; onto.\par
% glossary-block:B01498:end
% glossary-block:B01499:begin
\needspace{3\baselineskip}
\textbf{uppan} (prep.). upon; on top of.\par
% glossary-block:B01499:end
% glossary-block:B01500:begin
\needspace{3\baselineskip}
\textbf{ūre} (possessive adjective). our; ours. Plural possessors; agreeing possessive. Also gen. pl.: see \textbf{ic}.\par
% glossary-block:B01500:end
% glossary-block:B01501:begin
\needspace{3\baselineskip}
\textbf{urnon} — pret. ind. pl.: see \textbf{iernan}.\par
% glossary-block:B01501:end
% glossary-block:B01502:begin
\needspace{3\baselineskip}
\textbf{ūs} — dat. pl.; acc. pl.: see \textbf{ic}.\par
% glossary-block:B01502:end
% glossary-block:B01503:begin
\needspace{3\baselineskip}
\textbf{ūser} — gen. pl., variant: see \textbf{ic}.\par
% glossary-block:B01503:end
% glossary-block:B01504:begin
\needspace{3\baselineskip}
\textbf{ūsic} — acc. pl., variant: see \textbf{ic}.\par
% glossary-block:B01504:end
% glossary-block:B01505:begin
\needspace{3\baselineskip}
\textbf{ūt} (adv.). out.\par
% glossary-block:B01505:end
% glossary-block:B01506:begin
\needspace{3\baselineskip}
\textbf{ūte} (adv.). outside.\par
% glossary-block:B01506:end
% glossary-block:B01507:begin
\needspace{3\baselineskip}
\textbf{ūtēode} — pret. sg.: see \textbf{ūtgān}.\par
% glossary-block:B01507:end
% glossary-block:B01508:begin
\needspace{3\baselineskip}
\textbf{ūtgān} (v.). go out. Supplied past source form. go out.\par
% glossary-block:B01508:end
% glossary-block:B01509:begin
\needspace{3\baselineskip}
\textbf{ūtgangende} (v.). come out. Supplied present participle; source permits its interpretation with earlier action. Only the cited source forms are supplied here.\par
% glossary-block:B01509:end
% glossary-block:B01510:begin
\needspace{3\baselineskip}
\textbf{ūðe} — pret. sg.: see \textbf{unnan}.\par
% glossary-block:B01510:end
% book-navigation: standalone bookmark replaced by volume hierarchy
% glossary-block:B01511:begin
\needspace{5\baselineskip}
\subsection*{W}\BookHeading{nav-glossary-old-english-0023}{subsubsection}{W}{0}
% glossary-block:B01511:end
% glossary-block:B01512:begin
\textbf{wafian} (v.). be amazed. Supplied source form.\par
% glossary-block:B01512:end
% glossary-block:B01513:begin
\needspace{3\baselineskip}
\textbf{wāt} — pres. ind. 1/3 sg.: see \textbf{witan}.\par
% glossary-block:B01513:end
% glossary-block:B01514:begin
\needspace{3\baselineskip}
\textbf{wǣd} (n. f.); gen. sg. wǣde; nom. pl. wǣda. garment.\par
% glossary-block:B01514:end
% glossary-block:B01515:begin
\needspace{3\baselineskip}
\textbf{wǣdla} (n. m.); gen. sg. wǣdlan; nom. pl. wǣdlan. poor or destitute man.\par
% glossary-block:B01515:end
% glossary-block:B01516:begin
\needspace{3\baselineskip}
\textbf{wǣfels} (n. m.). cloak.\par
% glossary-block:B01516:end
% glossary-block:B01517:begin
\needspace{3\baselineskip}
\textbf{wæl} (n. n.). slaughter.\par
% glossary-block:B01517:end
% glossary-block:B01518:begin
\needspace{3\baselineskip}
\textbf{wǣpen} (n. n.); gen. sg. wǣpnes; nom. pl. wǣpnu. weapon.\par
% glossary-block:B01518:end
% glossary-block:B01519:begin
\needspace{3\baselineskip}
\textbf{wǣron} — pret. ind. pl.: see \textbf{wesan}.\par
% glossary-block:B01519:end
% glossary-block:B01520:begin
\needspace{3\baselineskip}
\textbf{wæs} — pret. ind. 1/3 sg.: see \textbf{wesan}.\par
% glossary-block:B01520:end
% glossary-block:B01521:begin
\needspace{3\baselineskip}
\textbf{wæstm} (n. m.); gen. sg. wæstmes; nom. pl. wæstmas. fruit; produce.\par
% glossary-block:B01521:end
% glossary-block:B01522:begin
\needspace{3\baselineskip}
\textbf{wæstmbǣre} (adj.). fruitful.\par
% glossary-block:B01522:end
% glossary-block:B01523:begin
\needspace{3\baselineskip}
\textbf{wæter} (n. n.); gen. sg. wæteres; nom. pl. wæteru. water.\par
% glossary-block:B01523:end
% glossary-block:B01524:begin
\needspace{3\baselineskip}
\textbf{wē} — nom. pl.: see \textbf{ic}.\par
% glossary-block:B01524:end
% glossary-block:B01525:begin
\needspace{3\baselineskip}
\textbf{wealdend} (n. m.). Ruler.\par
% glossary-block:B01525:end
% glossary-block:B01526:begin
\needspace{3\baselineskip}
\textbf{wealhstōd} (n. m.). interpreter; translator.\par
% glossary-block:B01526:end
% glossary-block:B01527:begin
\needspace{3\baselineskip}
\textbf{weall} (n. m.); gen. sg. wealles; nom. pl. weallas. wall.\par
% glossary-block:B01527:end
% glossary-block:B01528:begin
\needspace{3\baselineskip}
\textbf{weallan} (v.). surge; boil. Supplied present participle used adjectivally of the sea.\par
% glossary-block:B01528:end
% glossary-block:B01529:begin
\needspace{3\baselineskip}
\textbf{weax} (n.). wax. Pronunciation example.\par
% glossary-block:B01529:end
% glossary-block:B01530:begin
\needspace{3\baselineskip}
\textbf{weaxan} (v. strong); pret. 1 sg. wēox; pret. pl. wēoxon; past ptcp. weaxen; pres. 3 sg. wiext, wiexð. grow.\par
% glossary-block:B01530:end
% glossary-block:B01531:begin
\needspace{3\baselineskip}
\textbf{węcg} (n. m.). metal; ingot. Supplied genitive plural dependent on veins.\par
% glossary-block:B01531:end
% glossary-block:B01532:begin
\needspace{3\baselineskip}
\textbf{weg} (n. m.); gen. sg. weges; nom. pl. wegas. way; path. way; means of access. Figurative access to the Father in the declaration; distinct from physical road.\par
% glossary-block:B01532:end
% glossary-block:B01533:begin
\needspace{3\baselineskip}
\textbf{wel} (adv.); comp. bęt; superl. bętst. well. Comparative bęt: better. Superlative bętst: best.\par
% glossary-block:B01533:end
% glossary-block:B01534:begin
\needspace{3\baselineskip}
\textbf{wela} (n. m.). riches; wealth.\par
% glossary-block:B01534:end
% glossary-block:B01535:begin
\needspace{3\baselineskip}
\textbf{welig} (adj.). rich; wealthy. rich in; abounding in. on + dative thing abounded in.\par
% glossary-block:B01535:end
% glossary-block:B01536:begin
\needspace{3\baselineskip}
\textbf{wēn} (n. f.); gen. sg. wēne. expectation; chance. wēn is ðæt + subjunctive supplies the taught perhaps construction.\par
% glossary-block:B01536:end
% glossary-block:B01537:begin
\needspace{3\baselineskip}
\textbf{wēnan} (v. weak); pret. 1 sg. wēnde; past ptcp. wēned. suppose; expect; think. Genitive thing expected in the selected nominal frame; content-clause uses are separately taught.\par
% glossary-block:B01537:end
% glossary-block:B01538:begin
\needspace{3\baselineskip}
\textbf{węndan} (v.). translate. turn; return. Supplied source form. Supplied preterite węnde: went.\par
% glossary-block:B01538:end
% glossary-block:B01539:begin
\needspace{3\baselineskip}
\textbf{weorc} (n. n.); gen. sg. weorces; nom. pl. weorc. work.\par
% glossary-block:B01539:end
% glossary-block:B01540:begin
\needspace{3\baselineskip}
\textbf{weorod} (n. n.). host; army.\par
% glossary-block:B01540:end
% glossary-block:B01541:begin
\needspace{3\baselineskip}
\textbf{weorpan} (v. strong); pret. 1 sg. wearp; pret. pl. wurpon; past ptcp. worpen; pres. 3 sg. wierpð. throw. Accusative thing thrown.\par
% glossary-block:B01541:end
% glossary-block:B01542:begin
\needspace{3\baselineskip}
\textbf{weorð} (n. n.); gen. sg. weorðes. price; value.\par
% glossary-block:B01542:end
% glossary-block:B01543:begin
\needspace{3\baselineskip}
\textbf{weorðan} (v. strong); pret. 1 sg. wearð; pret. pl. wurdon; past ptcp. worden; pres. 3 sg. wierð. become; be. Auxiliary of the assigned passive event with a past participle. become. Nominative subject and agreeing nominative predicate. arise; come about. Intransitive with a nominative subject, including the arising of a calm. come about; result. Supplied resultative expression.\par
% glossary-block:B01543:end
% glossary-block:B01544:begin
\needspace{3\baselineskip}
\textbf{weorðscipe} (n. m.); gen. sg. weorðscipes. honour; dignity.\par
% glossary-block:B01544:end
% glossary-block:B01545:begin
\needspace{3\baselineskip}
\textbf{wēpan} (v. strong); pret. 1 sg. wēop; pret. pl. wēopon; past ptcp. wōpen. weep. weeping; tearful. Supplied participial/adjectival use.\par
% glossary-block:B01545:end
% glossary-block:B01546:begin
\needspace{3\baselineskip}
\textbf{wer} (n. m.); nom. pl. weras. man. Supplied genitive singular. husband.\par
% glossary-block:B01546:end
% glossary-block:B01547:begin
\needspace{3\baselineskip}
\textbf{wesan} (v. irregular); pres. 1 sg. eom; pres. 3 sg. is; pres. pl. sindon; pret. 1 sg. wæs; pret. pl. wǣron; bēon-series pres. 1 sg. bēo; bēon-series pres. 3 sg. bið; bēon-series pres. pl. bēoð; negative pret. 1 sg. næs; negative pret. pl. nǣron; negative pres. 3 sg. nis. be. Copular present/past identification or state; other auxiliary constructions are separate senses. be. Passive auxiliary with the assigned participle. be. Auxiliary with intransitive past participle in the specified compound constructions. be. Auxiliary with a present participle in the supplied construction; not a universal English-progressive equivalent. be. Supplied plural present copula. seem. Only in the supplied impersonal expression it seemed to the king. not be. Supplied negative nǣre in reported supposition. would have come. Supplied past-subjunctive wesan + cumen in the counterfactual result. The checked source paradigm gives no past participle.\par
% glossary-block:B01547:end
% glossary-block:B01548:begin
\needspace{3\baselineskip}
\textbf{westan} (adv.). from the west.\par
% glossary-block:B01548:end
% glossary-block:B01549:begin
\needspace{3\baselineskip}
\textbf{wīc} (n. n.); gen. sg. wīces; nom. pl. wīc. dwelling; abode.\par
% glossary-block:B01549:end
% glossary-block:B01550:begin
\needspace{3\baselineskip}
\textbf{wician} (v. weak); pret. 1 sg. wicode; past ptcp. wicod. sojourn; lodge. Intransitive; mid + dative person with whom one lodges.\par
% glossary-block:B01550:end
% glossary-block:B01551:begin
\needspace{3\baselineskip}
\textbf{wīd} (adjective citation). wide.\par
% glossary-block:B01551:end
% glossary-block:B01552:begin
\needspace{3\baselineskip}
\textbf{wīde} (adv.); comp. wīdor; superl. wīdost. widely; far and wide. Comparative wīdor: more widely. Superlative wīdost: most widely.\par
% glossary-block:B01552:end
% glossary-block:B01553:begin
\needspace{3\baselineskip}
\textbf{wīdor} — comparative adverb: see \textbf{wīde}.\par
% glossary-block:B01553:end
% glossary-block:B01554:begin
\needspace{3\baselineskip}
\textbf{wīdost} — superlative adverb: see \textbf{wīde}.\par
% glossary-block:B01554:end
% glossary-block:B01555:begin
\needspace{3\baselineskip}
\textbf{wierpð} — pres. ind. 3 sg.: see \textbf{weorpan}.\par
% glossary-block:B01555:end
% glossary-block:B01556:begin
\needspace{3\baselineskip}
\textbf{wierrest} (adj.). worst. Superlative of yfel.\par
% glossary-block:B01556:end
% glossary-block:B01557:begin
\needspace{3\baselineskip}
\textbf{wiersa} (adj.). worse. Comparative of yfel.\par
% glossary-block:B01557:end
% glossary-block:B01558:begin
\needspace{3\baselineskip}
\textbf{wierð} — pres. ind. 3 sg.: see \textbf{weorðan}.\par
% glossary-block:B01558:end
% glossary-block:B01559:begin
\needspace{3\baselineskip}
\textbf{wierðe} (adj.). worthy.\par
% glossary-block:B01559:end
% glossary-block:B01560:begin
\needspace{3\baselineskip}
\textbf{wiext} — pres. ind. 3 sg.: see \textbf{weaxan}.\par
% glossary-block:B01560:end
% glossary-block:B01561:begin
\needspace{3\baselineskip}
\textbf{wiexð} — pres. ind. 3 sg.: see \textbf{weaxan}.\par
% glossary-block:B01561:end
% glossary-block:B01562:begin
\needspace{3\baselineskip}
\textbf{wīf} (n. n.); gen. sg. wīfes; nom. pl. wīf. wife. woman.\par
% glossary-block:B01562:end
% glossary-block:B01563:begin
\needspace{3\baselineskip}
\textbf{wilde} (adj.). wild. Reference in a source or explanatory note.\par
% glossary-block:B01563:end
% glossary-block:B01564:begin
\needspace{3\baselineskip}
\textbf{willa} (n. m.); nom. pl. willan. will; disposition.\par
% glossary-block:B01564:end
% glossary-block:B01565:begin
\needspace{3\baselineskip}
\textbf{willan} (v. irregular); pres. 1/3 sg. wille; pres. pl. willað; pret. 1 sg. wolde; pret. pl. woldon; negative pres. 3 sg. nele; negative pret. 1 sg. nolde; negative pret. pl. noldon; negative pres. 1 sg. nelle; negative pres. pl. nellað. wish; want to. be unwilling; refuse. With ne or a licensed negative contraction; followed by plain infinitive. might. Only in the exact supplied purpose construction with watching/seeing infinitives. The checked source paradigm gives no past participle.\par
% glossary-block:B01565:end
% glossary-block:B01566:begin
\needspace{3\baselineskip}
\textbf{wīn} (n. n.); gen. sg. wīnes; nom. pl. wīn. wine.\par
% glossary-block:B01566:end
% glossary-block:B01567:begin
\needspace{3\baselineskip}
\textbf{wind} (n. m.); gen. sg. windes; nom. pl. windas. wind.\par
% glossary-block:B01567:end
% glossary-block:B01568:begin
\needspace{3\baselineskip}
\textbf{windan} (v. strong); pret. 1 sg. wand; pret. pl. wundon; past ptcp. wunden. wind; plait.\par
% glossary-block:B01568:end
% glossary-block:B01569:begin
\needspace{3\baselineskip}
\textbf{wīngeard} (n. m.); gen. sg. wīngeardes; nom. pl. wīngeardas. vineyard.\par
% glossary-block:B01569:end
% glossary-block:B01570:begin
\needspace{3\baselineskip}
\textbf{winnan} (v. strong); pret. 1 sg. wann; pret. pl. wunnon; past ptcp. wunnen. struggle; toil.\par
% glossary-block:B01570:end
% glossary-block:B01571:begin
\needspace{3\baselineskip}
\textbf{winstre} (adj.). left. Supplied side phrase.\par
% glossary-block:B01571:end
% glossary-block:B01572:begin
\needspace{3\baselineskip}
\textbf{winter} (n. m.); gen. sg. wintres; nom. pl. wintras. winter.\par
% glossary-block:B01572:end
% glossary-block:B01573:begin
\needspace{3\baselineskip}
\textbf{winterlic} (adj.). winter; of winter.\par
% glossary-block:B01573:end
% glossary-block:B01574:begin
\needspace{3\baselineskip}
\textbf{wīs} (adj.); comp. wīsra. wise. wise man; wise one. Substantival direct address.\par
% glossary-block:B01574:end
% glossary-block:B01575:begin
\needspace{3\baselineskip}
\textbf{wīsdōm} (n. m.); gen. sg. wīsdōmes. wisdom.\par
% glossary-block:B01575:end
% glossary-block:B01576:begin
\needspace{3\baselineskip}
\textbf{wīsra} (adj.). wiser. Comparative of wīs.\par
% glossary-block:B01576:end
% glossary-block:B01577:begin
\needspace{3\baselineskip}
\textbf{wisse} — pret. sg., variant: see \textbf{witan}.\par
% glossary-block:B01577:end
% glossary-block:B01578:begin
\needspace{3\baselineskip}
\textbf{wiste} — pret. sg.: see \textbf{witan}.\par
% glossary-block:B01578:end
% glossary-block:B01579:begin
\needspace{3\baselineskip}
\textbf{wit} — nom. dual: see \textbf{ic}.\par
% glossary-block:B01579:end
% glossary-block:B01580:begin
\needspace{3\baselineskip}
\textbf{witan} (v. preterite-present); pres. 1/3 sg. wāt; pres. pl. witon; pret. 1 sg. wiste; past ptcp. (ge)witen. know. know. A stated fact or the answer to a question as complement. know. Accusative person in this taught frame; cunnan has a separate acquaintance sense.\par
% glossary-block:B01580:end
% glossary-block:B01581:begin
\needspace{3\baselineskip}
\textbf{wītan} (v.). blame; censure.\par
% glossary-block:B01581:end
% glossary-block:B01582:begin
\needspace{3\baselineskip}
\textbf{wīte} (n. n.). punishment; penalty.\par
% glossary-block:B01582:end
% glossary-block:B01583:begin
\needspace{3\baselineskip}
\textbf{wītga} (n. m.); gen. sg. wītgan; nom. pl. wītgan. prophet.\par
% glossary-block:B01583:end
% glossary-block:B01584:begin
\needspace{3\baselineskip}
\textbf{witodlice} (adv.). now; and. Discourse connective.\par
% glossary-block:B01584:end
% glossary-block:B01585:begin
\needspace{3\baselineskip}
\textbf{wið} (prep.). against; opposed to. Accusative opponent. by; beside; at. Accusative landmark in the taught proximity phrase. in return for. With dative; distinct from opponent/hostility. with; against. Dative fellow players in the exact supplied playing frame. with. Supplied conversational counterpart; do not infer hostility. toward. Conduct toward a person in the supplied local frame.\par
% glossary-block:B01585:end
% glossary-block:B01586:begin
\needspace{3\baselineskip}
\textbf{wiðerwinna} (n. m.). adversaries.\par
% glossary-block:B01586:end
% glossary-block:B01587:begin
\needspace{3\baselineskip}
\textbf{wiðstandan} (v. strong); pret. 1 sg. wiðstōd; pret. pl. wiðstōdon; past ptcp. wiðstanden; pres. 3 sg. wiðstęnt. withstand; resist.\par
% glossary-block:B01587:end
% glossary-block:B01588:begin
\needspace{3\baselineskip}
\textbf{wiðstęnt} — pres. ind. 3 sg.: see \textbf{wiðstandan}.\par
% glossary-block:B01588:end
% glossary-block:B01589:begin
\needspace{3\baselineskip}
\textbf{wlanc} (adj.). proud; wealthy. Contextually the rich man addressed.\par
% glossary-block:B01589:end
% glossary-block:B01590:begin
\needspace{3\baselineskip}
\textbf{wolcen} (n. n.); gen. sg. wolcnes; nom. pl. wolcnu. cloud.\par
% glossary-block:B01590:end
% glossary-block:B01591:begin
\needspace{3\baselineskip}
\textbf{wolde} — pret. sg.: see \textbf{willan}.\par
% glossary-block:B01591:end
% glossary-block:B01592:begin
\needspace{3\baselineskip}
\textbf{word} (n. n.); gen. sg. wordes; nom. pl. word. word. speech; utterance.\par
% glossary-block:B01592:end
% glossary-block:B01593:begin
\needspace{3\baselineskip}
\textbf{worhte} — pret. sg.: see \textbf{wyrcean}.\par
% glossary-block:B01593:end
% glossary-block:B01594:begin
\needspace{3\baselineskip}
\textbf{woruld} (n. f.); gen. sg. worulde; nom. pl. worulda. world.\par
% glossary-block:B01594:end
% glossary-block:B01595:begin
\needspace{3\baselineskip}
\textbf{woruldbisgu} (n. f.). worldly occupations; cares.\par
% glossary-block:B01595:end
% glossary-block:B01596:begin
\needspace{3\baselineskip}
\textbf{woruldcræft} (n. m.). worldly occupation.\par
% glossary-block:B01596:end
% glossary-block:B01597:begin
\needspace{3\baselineskip}
\textbf{woruldgeðyngð} (n. f.). worldly honour; worldly dignity.\par
% glossary-block:B01597:end
% glossary-block:B01598:begin
\needspace{3\baselineskip}
\textbf{woruldlīf} (n. n.). worldly life.\par
% glossary-block:B01598:end
% glossary-block:B01599:begin
\needspace{3\baselineskip}
\textbf{wrēon} (v. strong); pret. 1 sg. wrāh; pret. pl. wrigon; past ptcp. wrigen; pres. 3 sg. wrīhð. cover; clothe.\par
% glossary-block:B01599:end
% glossary-block:B01600:begin
\needspace{3\baselineskip}
\textbf{wrīhð} — pres. ind. 3 sg.: see \textbf{wrēon}.\par
% glossary-block:B01600:end
% glossary-block:B01601:begin
\needspace{3\baselineskip}
\textbf{wrītan} (v. strong); pret. 1 sg. wrāt; pret. pl. writon; past ptcp. writen. write. Accusative text written.\par
% glossary-block:B01601:end
% glossary-block:B01602:begin
\needspace{3\baselineskip}
\textbf{wudu} (n. m.). wood; forest. Collective vegetation, not a new material sense.\par
% glossary-block:B01602:end
% glossary-block:B01603:begin
\needspace{3\baselineskip}
\textbf{wudubearu} (n. m.). forest; grove.\par
% glossary-block:B01603:end
% glossary-block:B01604:begin
\needspace{3\baselineskip}
\textbf{wuldor} (n. n.); gen. sg. wuldres; nom. pl. wuldru. glory.\par
% glossary-block:B01604:end
% glossary-block:B01605:begin
\needspace{3\baselineskip}
\textbf{wuldrian} (v. weak); pret. 1 sg. wuldrode; past ptcp. wuldrod. glorify. Accusative person; an accusative reflexive expresses glorifying oneself.\par
% glossary-block:B01605:end
% glossary-block:B01606:begin
\needspace{3\baselineskip}
\textbf{wulf} (n. m.); gen. sg. wulfes; nom. pl. wulfas. wolf.\par
% glossary-block:B01606:end
% glossary-block:B01607:begin
\needspace{3\baselineskip}
\textbf{wund} (n. f.); gen. sg. wunde; nom. pl. wunda. wound; sore.\par
% glossary-block:B01607:end
% glossary-block:B01608:begin
\needspace{3\baselineskip}
\textbf{wundor} (n. n.); gen. sg. wundres; nom. pl. wundor. wonder.\par
% glossary-block:B01608:end
% glossary-block:B01609:begin
\needspace{3\baselineskip}
\textbf{wundrian} (v.). wonder. Supplied source form.\par
% glossary-block:B01609:end
% glossary-block:B01610:begin
\needspace{3\baselineskip}
\textbf{wunian} (v. weak); pret. 1 sg. wunode; past ptcp. wunod. dwell; remain.\par
% glossary-block:B01610:end
% glossary-block:B01611:begin
\needspace{3\baselineskip}
\textbf{wynsum} (adj.). pleasant; winsome.\par
% glossary-block:B01611:end
% glossary-block:B01612:begin
\needspace{3\baselineskip}
\textbf{wyrcean} (v. weak); pret. 1 sg. worhte; past ptcp. worht. make; build. Supplied unprefixed form; preserve explicit task controls. yield; produce. Accusative fruit in the taught/source expression. work; perform. Accusative work or deed.\par
% glossary-block:B01612:end
% glossary-block:B01613:begin
\needspace{3\baselineskip}
\textbf{wyrhta} (n. m.); gen. sg. wyrhtan; nom. pl. wyrhtan. workman; craftsman. Maker; Creator. Divine title, distinct from ordinary workman.\par
% glossary-block:B01613:end
% glossary-block:B01614:begin
\needspace{3\baselineskip}
\textbf{wyrt} (n. f.); gen. sg. wyrte; nom. pl. wyrta. herb; plant.\par
% glossary-block:B01614:end
% glossary-block:B01615:begin
\needspace{3\baselineskip}
\textbf{wyrtgemang} (n. n.); gen. sg. wyrtgemanges; nom. pl. wyrtgemang. spice; mixture of spices.\par
% glossary-block:B01615:end
% glossary-block:B01616:begin
\needspace{3\baselineskip}
\textbf{wyrttruma} (n. m.); gen. sg. wyrttruman; nom. pl. wyrttruman. root.\par
% glossary-block:B01616:end
% book-navigation: standalone bookmark replaced by volume hierarchy
% glossary-block:B01617:begin
\needspace{5\baselineskip}
\subsection*{Y}\BookHeading{nav-glossary-old-english-0024}{subsubsection}{Y}{0}
% glossary-block:B01617:end
% glossary-block:B01618:begin
\textbf{yfel} (adj.); comp. wiersa; superl. wierrest. bad; evil. evil; evil people. Substantival adjective in the assigned use. evil. Substantive noun in fram yfele.\par
% glossary-block:B01618:end
% glossary-block:B01619:begin
\needspace{3\baselineskip}
\textbf{yfele} (adv.). badly; wrongly.\par
% glossary-block:B01619:end
% glossary-block:B01620:begin
\needspace{3\baselineskip}
\textbf{ymbe} (prep.). about. Approximate time, with accusative.\par
% glossary-block:B01620:end
% glossary-block:B01621:begin
\needspace{3\baselineskip}
\textbf{ymbhwyrft} (n. m.); gen. sg. ymbhwyrftes; nom. pl. ymbhwyrftas. circuit. Whole on ymbhwyrfte expression means all around; supplied dependent landmark.\par
% glossary-block:B01621:end
% glossary-block:B01622:begin
\needspace{3\baselineskip}
\textbf{ymbsealde} — pret. sg.: see \textbf{ymbsęllan}.\par
% glossary-block:B01622:end
% glossary-block:B01623:begin
\needspace{3\baselineskip}
\textbf{ymbsęllan} (v.); past ptcp. ymbseald. surround; envelop. Supplied source form. surround; beset. Supplied source form.\par
% glossary-block:B01623:end
% glossary-block:B01624:begin
\needspace{3\baselineskip}
\textbf{ymbsittan} (v.). sit around. Supplied participle referring to people sitting around.\par
% glossary-block:B01624:end
% glossary-block:B01625:begin
\needspace{3\baselineskip}
\textbf{ȳst} (n. f.); gen. sg. ȳste; nom. pl. ȳsta. storm.\par
% glossary-block:B01625:end
% glossary-block:B01626:begin
\needspace{3\baselineskip}
\textbf{ȳð} (n. f.). wave; billow. Supplied dative plural after fram.\par
% glossary-block:B01626:end
\end{multicols*}
\clearpage
% book-navigation: standalone bookmark replaced by volume hierarchy
% glossary-block:B01627:begin
\needspace{5\baselineskip}
\section*{Quoted errors and reconstructions}\BookHeading{nav-glossary-old-english-0025}{subsection}{Quoted errors and reconstructions}{1}
% glossary-block:B01627:end
% glossary-block:B01628:begin
These notices preserve what the book quotes. A spelling rejected in a particular exercise is not thereby declared impossible in every historical context.\par
% glossary-block:B01628:end
% glossary-block:B01629:begin
\textbf{āgyltan}. Weak verb: sin; fail in duty. In Cook 34.9 this is an erroneous hint for reward; see āgyldan (strong III), whose contracted third singular is āgylt. [C06].\par
% glossary-block:B01629:end
% glossary-block:B01630:begin
\textbf{dægas}. Deliberate learner error; requested plural is dagas.\par
% glossary-block:B01630:end
% glossary-block:B01631:begin
\textbf{fiscnað}. Unselected respelling of the source word; the course retains fiscnoð.\par
% glossary-block:B01631:end
% glossary-block:B01632:begin
\textbf{fōnende}. Rejected mechanical present participle; selected form is fōnde.\par
% glossary-block:B01632:end
% glossary-block:B01633:begin
\textbf{hwælas}. Deliberate learner error; requested plural is hwalas.\par
% glossary-block:B01633:end
% glossary-block:B01634:begin
\textbf{mynsteru}. Unselected answer in the course exercise; see the selected plural under mynster. This does not rule out every historical use of the spelling.\par
% glossary-block:B01634:end
% glossary-block:B01635:begin
\textbf{sǣwa}. Reconstructed genitive plural of sǣ, marked with an asterisk in the book; it is not tested as an attested form.\par
% glossary-block:B01635:end
% glossary-block:B01636:begin
\textbf{swētrra}. Rejected unsimplified comparative genitive plural; selected form is swētra.\par
% glossary-block:B01636:end
% glossary-block:B01637:begin
\textbf{unnað}. Rejected manufactured indicative plural; selected present plural is unnon.\par
% glossary-block:B01637:end
% glossary-block:B01638:begin
\textbf{wastmas}. Deliberate learner error; requested plural is wæstmas.\par
% glossary-block:B01638:end
% glossary-block:B01639:begin
\textbf{wǣstm}. Cook’s erroneous quantity mark in the cited hint; see wæstm, with short æ. [C01].\par
% glossary-block:B01639:end
% glossary-block:B01640:begin
\textbf{wīsrra}. Rejected unsimplified comparative genitive plural; selected form is wīsra.\par
% glossary-block:B01640:end
% glossary-block:B01641:begin
\textbf{wundas}. Rejected masculine-pattern plural of the feminine wund in this context; selected plural is wunda.\par
% glossary-block:B01641:end

% book-source-body:end
\par\label{end-glossary-old-english}
\clearpage\endgroup

\BookRecto
\begingroup
\renewcommand{\runningtitle}{English–Old English}
\raggedcolumns
\newcommand{\needspace}[1]{\par\begingroup\skip0=\lastskip\vskip-\skip0\dimen0=#1\relax\advance\dimen0 by\skip0\vskip0pt plus\dimen0\penalty-100\vskip0pt plus-\dimen0\vskip\dimen0\penalty9999\vskip-\dimen0\vskip\skip0\endgroup}
\setlength{\columnsep}{6mm}
\hyphenpenalty=10000\exhyphenpenalty=10000
\emergencystretch=2em
% book-source-body:begin

\addfontfeatures{Ligatures=TeXOff}
\raggedright
% book-navigation: standalone bookmark replaced by volume hierarchy
% glossary-block:B00001:begin
\booktitle{\BookComponent{glossary-english}{English–Old English exercise glossary}GLOSSARY}{English–Old English exercise glossary}
% glossary-block:B00001:end
% glossary-block:B00002:begin
This glossary supplies English access to the vocabulary and expressions needed for the exercises in all eighteen units: Cook's items, original practice and reading work. The pronunciation practice supplies its own word meanings. English used only to give an instruction, such as “parse” or “give the genitive,” is not vocabulary to translate.\par
% glossary-block:B00002:end
% glossary-block:B00003:begin
Look up the intended meaning or the whole phrase. Choices under one English word are not always interchangeable. Read the qualifying meaning and case or construction note, then use the forms in the companion Old English–English glossary and the relevant unit. Exercise targets and supplied wording still apply; this reference does not turn supplied material into an independent-recall requirement.\par
% glossary-block:B00003:end
% glossary-block:B00004:begin
English articles, auxiliaries and dummy subjects often require a construction rather than a matching Old English word. Entries such as \textbf{a}, \textbf{do}, \textbf{will}, \textbf{there is}, and phrase entries explain these differences. A case label describes the Old English dependent: “genitive” requires a genitive form, not an extra word meaning “of.”\par
% glossary-block:B00004:end
% glossary-block:B00005:begin
Old English citation forms retain the book's spelling and length marks. Some source expressions are supplied only in their stated context. Inflected or conventional citation labels are identified where a normal infinitive or full paradigm is not established. Use the separate key to check a completed answer; this glossary is not a replacement key.\par
% glossary-block:B00005:end
% book-navigation: standalone bookmark replaced by volume hierarchy
\clearpage
\fontsize{10.5}{13.2}\selectfont
\setlength{\parskip}{4pt}
\begin{multicols}{2}
% glossary-block:B00006:begin
\needspace{5\baselineskip}
\section*{Entries}\BookHeading{nav-glossary-english-0001}{subsection}{Entries}{1}
% glossary-block:B00006:end
% book-navigation: standalone bookmark replaced by volume hierarchy
% glossary-block:B00007:begin
\needspace{5\baselineskip}
\subsection*{A}\BookHeading{nav-glossary-english-0002}{subsubsection}{A}{0}
% glossary-block:B00007:end
% glossary-block:B00008:begin
\textbf{a} — Ordinary English a/an usually has no separate Old English article. Ān is used in its taught a-certain-individual sense, not inserted into every indefinite noun phrase. No obligatory indefinite article: bare predicates/nouns often translate a/an. ān is used for the taught introducing-a-certain-individual sense, not inserted everywhere.\par
% glossary-block:B00008:end
% glossary-block:B00009:begin
\needspace{3\baselineskip}
\textbf{abandon} — \textbf{forlǣtan} (v.). abandon; forsake. Accusative person or thing left. Accusative thing/place left. Distinguish leave to someone and take leave of.\par
% glossary-block:B00009:end
% glossary-block:B00010:begin
\needspace{3\baselineskip}
\textbf{ability} — \textbf{miht} (n.). ability; power.\par
% glossary-block:B00010:end
% glossary-block:B00011:begin
\needspace{3\baselineskip}
\textbf{abode} — \textbf{wīc} (n.). dwelling; abode.\par
% glossary-block:B00011:end
% glossary-block:B00012:begin
\needspace{3\baselineskip}
\textbf{abounding in} — \textbf{welig} (adj.). rich in; abounding in. on + dative thing abounded in.\par
% glossary-block:B00012:end
% glossary-block:B00013:begin
\needspace{3\baselineskip}
\textbf{about} — \textbf{ābūtan} (prep.). around; about. With accusative. \textbf{be} (prep.). concerning; about. With dative. \textbf{ofer} (prep.). over; about. Object of rejoicing, with the assigned accusative frame. \textbf{swā-swā} (conj.). about; approximately. Approximation in the supplied source expression. \textbf{swilce} (adv.). about; approximately. \textbf{ymbe} (prep.). Approximate time, with accusative.\par
% glossary-block:B00013:end
% glossary-block:B00014:begin
\needspace{3\baselineskip}
\textbf{above} — \textbf{bufan} (prep.). With dative landmark in the taught spatial frame. \textbf{ofer} (prep.). over; across. Spatial frame; case varies with the specific construction. Do not impose a universal motion rule. Rank, with accusative. Rank, with explicitly supplied dative complement; retain this limit. above; more than. Preference in the supplied comparison. Spatial ofer can mean over/across; its case depends on the assigned spatial frame. Ofer in the authority sense over takes an accusative dependent. Rank and preference have further separately qualified senses.\par
% glossary-block:B00014:end
% glossary-block:B00015:begin
\needspace{3\baselineskip}
\textbf{above all} — \textbf{tōforan} (prep.). before; above all.\par
% glossary-block:B00015:end
% glossary-block:B00016:begin
\needspace{3\baselineskip}
\textbf{Abraham} — \textbf{Abraham} (proper name). Biblical patriarch; supplied name forms.\par
% glossary-block:B00016:end
% glossary-block:B00017:begin
\needspace{3\baselineskip}
\textbf{abundance} — \textbf{unrīm} (n.). abundance; uncounted quantity. Supplied quantity with genitive of the substance.\par
% glossary-block:B00017:end
% glossary-block:B00018:begin
\needspace{3\baselineskip}
\textbf{abyss} — \textbf{nēowolnes} (n.). abyss; the deep.\par
% glossary-block:B00018:end
% glossary-block:B00019:begin
\needspace{3\baselineskip}
\textbf{accept} — \textbf{onfōn} (v.). receive; accept. Dative thing received in this selected frame; the later accusative payment frame is separate. receive; accept. A person received into care in the supplied frame.\par
% glossary-block:B00019:end
% glossary-block:B00020:begin
\needspace{3\baselineskip}
\textbf{acceptable} — \textbf{gecwēme} (adj.). pleasing; acceptable.\par
% glossary-block:B00020:end
% glossary-block:B00021:begin
\needspace{3\baselineskip}
\textbf{accompanying} — \textbf{mid} (prep.; adv. in fixed expressions). with; accompanying. Dative person or cargo in the ordinary accompaniment frame.\par
% glossary-block:B00021:end
% glossary-block:B00022:begin
\needspace{3\baselineskip}
\textbf{accomplish} — \textbf{gefręmman} (v.). perform; accomplish. A work is the thing accomplished. \textbf{gefyllan} (v.). finish; accomplish. A work is the thing completed; distinct from fill. Gefyllan means finish/accomplish when a work is completed; the fill sense is separate.\par
% glossary-block:B00022:end
% glossary-block:B00023:begin
\needspace{3\baselineskip}
\textbf{accomplishment} — \textbf{cræft} (n.). skill; accomplishment.\par
% glossary-block:B00023:end
% glossary-block:B00024:begin
\needspace{3\baselineskip}
\textbf{accordingly} — \textbf{eornostlice} (adv.). accordingly; thus. \textbf{for-ðon-ðe} (conj.; adv. in for-ðon), \textbf{eornostlice} (adv.). The supplied consequence use of for-ðon-ðe means therefore/for that reason; eornostlice can mean accordingly/thus. Do not substitute the causal because sense.\par
% glossary-block:B00024:end
% glossary-block:B00025:begin
\needspace{3\baselineskip}
\textbf{according to} — \textbf{æfter} (prep.; conj. in fixed expressions). With dative. \textbf{be} (prep.). With dative. The supplied on + hīw expression means according to/in the forms of. Retain the whole supplied hīw phrase rather than treating on alone as a general according-to expression.\par
% glossary-block:B00025:end
% glossary-block:B00026:begin
\needspace{3\baselineskip}
\textbf{account} — \textbf{spell} (n.). account; discourse.\par
% glossary-block:B00026:end
% glossary-block:B00027:begin
\needspace{3\baselineskip}
\textbf{account heard} — \textbf{hlīsa} (n.). report; account heard.\par
% glossary-block:B00027:end
% glossary-block:B00028:begin
\needspace{3\baselineskip}
\textbf{a certain} — \textbf{ān} (numeral; adj.). Introduces an unnamed singular individual; not an obligatory indefinite article. \textbf{sum} (indefinite adjective/pronoun). a certain; some. \textbf{sum} (indefinite adjective/pronoun), \textbf{ān} (numeral; adj.). Sum or the taught introducing-an-individual ān can mean a certain. Ordinary English a/an does not require either word.\par
% glossary-block:B00028:end
% glossary-block:B00029:begin
\needspace{3\baselineskip}
\textbf{acknowledge} — \textbf{oncnāwan} (v.). recognize; acknowledge.\par
% glossary-block:B00029:end
% glossary-block:B00030:begin
\needspace{3\baselineskip}
\textbf{acorn} — \textbf{æcern} (n.).\par
% glossary-block:B00030:end
% glossary-block:B00031:begin
\needspace{3\baselineskip}
\textbf{across} — \textbf{ofer} (prep.). over; across. Spatial frame; case varies with the specific construction. Do not impose a universal motion rule.\par
% glossary-block:B00031:end
% glossary-block:B00032:begin
\needspace{3\baselineskip}
\textbf{act} — \textbf{plegian} (v.). act; perform. Performance context; ordinary play is a separate assigned sense.\par
% glossary-block:B00032:end
% glossary-block:B00033:begin
\needspace{3\baselineskip}
\textbf{add} — \textbf{tōgēiecan} (v.).  \textbf{underðēodan} (v.). add; subjoin. Supplied passive/participial addition to speech.\par
% glossary-block:B00033:end
% glossary-block:B00034:begin
\needspace{3\baselineskip}
\textbf{adder} — \textbf{næddre} (n.). snake; adder. Supplied weak feminine plural.\par
% glossary-block:B00034:end
% glossary-block:B00035:begin
\needspace{3\baselineskip}
\textbf{admonish} — \textbf{myngian} (v.). admonish; prompt. Supplied participial source use.\par
% glossary-block:B00035:end
% glossary-block:B00036:begin
\needspace{3\baselineskip}
\textbf{adorn} — \textbf{geglęngan} (v.). The supplied passive adornment expression uses mid + dative material. \textbf{geweorðian} (v.). distinguish; adorn. Accusative person adorned; mid + dative equipment or adornment.\par
% glossary-block:B00036:end
% glossary-block:B00037:begin
\needspace{3\baselineskip}
\textbf{adorned with} — \textbf{mid} (prep.; adv. in fixed expressions). with. Material or clothing in the supplied adornment phrase.\par
% glossary-block:B00037:end
% glossary-block:B00038:begin
\needspace{3\baselineskip}
\textbf{adversaries} — \textbf{wiðerwinna} (n.).\par
% glossary-block:B00038:end
% glossary-block:B00039:begin
\needspace{3\baselineskip}
\textbf{advise} — \textbf{lǣran} (v.). advise; instruct.\par
% glossary-block:B00039:end
% glossary-block:B00040:begin
\needspace{3\baselineskip}
\textbf{afar off} — \textbf{feorran} (adv.). feorran; the source supplies the from-afar paraphrase.\par
% glossary-block:B00040:end
% glossary-block:B00041:begin
\needspace{3\baselineskip}
\textbf{afflict} — \textbf{geðrǣstan} (v.). afflict; beat. Supplied participial source use. \textbf{ierman} (v.).\par
% glossary-block:B00041:end
% glossary-block:B00042:begin
\needspace{3\baselineskip}
\textbf{after} — \textbf{æfter} (prep.; conj. in fixed expressions). Temporal preposition with dative. In the supplied finite-clause conjunction æfter ðǣm ðe. after; following. Dative person followed, including motion after that person. after; in search of. Person sought or inquired after; keep this lexical frame distinct from temporal after. Temporal æfter + dative means after a person/event. The finite-clause conjunction is the complete supplied æfter ðǣm ðe. \textbf{siððan} (adv.; conj.). since; after; afterwards. Conjunction since/after versus adverb afterwards. Æfter ðǣm ðe is the complete supplied finite-clause conjunction after; ðe alone does not mean after.\par
% glossary-block:B00042:end
% glossary-block:B00043:begin
\needspace{3\baselineskip}
\textbf{after that} — Æfter ðon is the complete supplied temporal phrase after that.\par
% glossary-block:B00043:end
% glossary-block:B00044:begin
\needspace{3\baselineskip}
\textbf{afterwards} — \textbf{eft, ęft} (adv.). again; afterwards. \textbf{siððan} (adv.; conj.). since; after; afterwards. Conjunction since/after versus adverb afterwards.\par
% glossary-block:B00044:end
% glossary-block:B00045:begin
\needspace{3\baselineskip}
\textbf{again} — \textbf{eft, ęft} (adv.). again; also. The source note also permits the additive sense. again; back. again; afterwards. eft is the repeated-occurrence adverb. \textbf{mā} (adv.). more; again. In the negative phrase never more/never again. \textbf{ongēan} (adv.; prep.; conj.). back; again.\par
% glossary-block:B00045:end
% glossary-block:B00046:begin
\needspace{3\baselineskip}
\textbf{against} — \textbf{ongēan} (adv.; prep.; conj.). against; opposite. Accusative landmark or opponent. \textbf{wið} (prep.). against; opposed to. Accusative opponent. with; against. Dative fellow players in the exact supplied playing frame.\par
% glossary-block:B00046:end
% glossary-block:B00047:begin
\needspace{3\baselineskip}
\textbf{against that} — \textbf{ongēan} (adv.; prep.; conj.). by the time that. Complete supplied ongēan ðætte clause. Supplied ongēan ðætte + prospective subjunctive clause.\par
% glossary-block:B00047:end
% glossary-block:B00048:begin
\needspace{3\baselineskip}
\textbf{age} — \textbf{ieldu} (n.). age; old age.\par
% glossary-block:B00048:end
% glossary-block:B00049:begin
\needspace{3\baselineskip}
\textbf{agility} — \textbf{snelnes} (n.). skill; agility.\par
% glossary-block:B00049:end
% glossary-block:B00050:begin
\needspace{3\baselineskip}
\textbf{agreeably} — \textbf{gecwēmlīce} (adv.). pleasantly; agreeably.\par
% glossary-block:B00050:end
% glossary-block:B00051:begin
\needspace{3\baselineskip}
\textbf{agriculture} — \textbf{eorðtilð} (n.).\par
% glossary-block:B00051:end
% glossary-block:B00052:begin
\needspace{3\baselineskip}
\textbf{a hair} — \textbf{hǣr} (n.). hair; a hair.\par
% glossary-block:B00052:end
% glossary-block:B00053:begin
\needspace{3\baselineskip}
\textbf{ahead of} — \textbf{be-foran} (prep.). before; ahead of. Dative person, including going or sending ahead of that person.\par
% glossary-block:B00053:end
% glossary-block:B00054:begin
\needspace{3\baselineskip}
\textbf{a hundredfold} — \textbf{hundtēontig} (numeral). A hundred in return, in the supplied recompense expression.\par
% glossary-block:B00054:end
% glossary-block:B00055:begin
\needspace{3\baselineskip}
\textbf{aid} — \textbf{fultum} (n.). help; aid.\par
% glossary-block:B00055:end
% glossary-block:B00056:begin
\needspace{3\baselineskip}
\textbf{air} — \textbf{lyft} (n.).\par
% glossary-block:B00056:end
% glossary-block:B00057:begin
\needspace{3\baselineskip}
\textbf{a large number of} — \textbf{fela} (adj.). many; a large number of. Genitive of the things counted in the taught frame.\par
% glossary-block:B00057:end
% glossary-block:B00058:begin
\needspace{3\baselineskip}
\textbf{Alfred} — \textbf{Ælfred} (proper name).\par
% glossary-block:B00058:end
% glossary-block:B00059:begin
\needspace{3\baselineskip}
\textbf{a light} — \textbf{lēoht} (n.). light; luminary. Individual countable light, distinct from mass light. Mass light and an individual countable light/luminary are separate uses of lēoht. Use the count sense for a light or several lights.\par
% glossary-block:B00059:end
% glossary-block:B00060:begin
\needspace{3\baselineskip}
\textbf{alive} — \textbf{cucu} (adj.). alive; living. Supplied accusative neuter plural adjective. \textbf{cwic} (adj.). alive; living. \textbf{libban} (v.). living. Supplied weak participial adjective. The bracketed liveth yet is rendered by finite libban + gīet; an added copula is not required. The supplied whole phrase on līfe means alive; on alone does not.\par
% glossary-block:B00060:end
% glossary-block:B00061:begin
\needspace{3\baselineskip}
\textbf{all} — \textbf{eall} (adj.; substantive; adv.). all; whole. all people; everything. Substantival: plural for people, neuter singular for everything. Eall agrees with a noun in all/whole expressions. Substantival plural refers to all people; neuter singular can mean everything.\par
% glossary-block:B00061:end
% glossary-block:B00062:begin
\needspace{3\baselineskip}
\textbf{all around} — \textbf{ymbhwyrft} (n.). circuit. Whole on ymbhwyrfte expression means all around; supplied dependent landmark. Supplied on ymbhwyrfte + genitive landmark in this exercise.\par
% glossary-block:B00062:end
% glossary-block:B00063:begin
\needspace{3\baselineskip}
\textbf{allow} — \textbf{ālīefan} (v.). permit; allow. Impersonal permission frame, with dative permitted person; consult the complete supplied expression. \textbf{gelīefst} (v.; supplied pres. ind. 2 sg.). permit; allow. \textbf{līefan} (v.). allow; permit. \textbf{unnan} (v.). grant; allow. Dative person, genitive thing granted. Dative recipient and genitive thing granted; preserve the supported subjunctive command unne.\par
% glossary-block:B00063:end
% glossary-block:B00064:begin
\needspace{3\baselineskip}
\textbf{all people} — \textbf{eall} (adj.; substantive; adv.). all people; everything. Substantival: plural for people, neuter singular for everything.\par
% glossary-block:B00064:end
% glossary-block:B00065:begin
\needspace{3\baselineskip}
\textbf{almighty} — \textbf{ælmihtig} (adj.).\par
% glossary-block:B00065:end
% glossary-block:B00066:begin
\needspace{3\baselineskip}
\textbf{alone} — \textbf{ān} (numeral; adj.). Cardinal one/single and the agreeing adjective alone are distinct uses of ān. Preserve the supplied form where the exercise provides one.\par
% glossary-block:B00066:end
% glossary-block:B00067:begin
\needspace{3\baselineskip}
\textbf{along} — \textbf{geond} (prep.). along; through. With accusative.\par
% glossary-block:B00067:end
% glossary-block:B00068:begin
\needspace{3\baselineskip}
\textbf{along with} — \textbf{mid} (prep.; adv. in fixed expressions). with; along with. Supplied mid wes, be with them.\par
% glossary-block:B00068:end
% glossary-block:B00069:begin
\needspace{3\baselineskip}
\textbf{also} — \textbf{ēac} (adv.). Also in the combination ēac swilce. \textbf{eft, ęft} (adv.). again; also. The source note also permits the additive sense.\par
% glossary-block:B00069:end
% glossary-block:B00070:begin
\needspace{3\baselineskip}
\textbf{although} — \textbf{ðēah} (conj.; adv.). although; even though. ðēah with the taught concessive clause; a subjunctive need not make an admitted fact unreal. \textbf{ðēah ðe}. Supplied whole concessive conjunction in EX-XII-XIII-43-10.\par
% glossary-block:B00070:end
% glossary-block:B00071:begin
\needspace{3\baselineskip}
\textbf{always} — \textbf{ā} (adv.). ever; always. \textbf{ǣfre} (adv.). ever; always. \textbf{simle} (adv.).\par
% glossary-block:B00071:end
% glossary-block:B00072:begin
\needspace{3\baselineskip}
\textbf{am} — \textbf{wesan} (v.). be. Copular present/past identification or state; other auxiliary constructions are separate senses.\par
% glossary-block:B00072:end
% glossary-block:B00073:begin
\needspace{3\baselineskip}
\textbf{Amen} — \textbf{Amen} (interj.).\par
% glossary-block:B00073:end
% glossary-block:B00074:begin
\needspace{3\baselineskip}
\textbf{amid} — \textbf{on} (prep.). in; amid. Circumstances of the stated unrest.\par
% glossary-block:B00074:end
% glossary-block:B00075:begin
\needspace{3\baselineskip}
\textbf{among} — \textbf{betweox} (prep.). Accusative in the exact supplied phrases; not a universal case rule. among; between. \textbf{gemang} (prep.). Accusative group in the selected sending-among frame. \textbf{on} (prep.). among; in. Dative people among whom friends are obtained.\par
% glossary-block:B00075:end
% glossary-block:B00076:begin
\needspace{3\baselineskip}
\textbf{amount} — \textbf{micel} (adj.). much; an amount. Quantity sense. \textbf{dǣl} (n.). dǣl + genitive of the substance, as in sumne dǣl mętes; the amount is accusative object.\par
% glossary-block:B00076:end
% glossary-block:B00077:begin
\needspace{3\baselineskip}
\textbf{an} — Ordinary English a/an usually has no separate Old English article. Ān is used in its taught a-certain-individual sense, not inserted into every indefinite noun phrase. No obligatory indefinite article: bare predicates/nouns often translate a/an. ān is used for the taught introducing-a-certain-individual sense, not inserted everywhere.\par
% glossary-block:B00077:end
% glossary-block:B00078:begin
\needspace{3\baselineskip}
\textbf{and} — \textbf{and} (conj.). Simple coordination. \textbf{witodlice} (adv.). now; and. Discourse connective.\par
% glossary-block:B00078:end
% glossary-block:B00079:begin
\needspace{3\baselineskip}
\textbf{angel} — \textbf{ęngel} (n.).\par
% glossary-block:B00079:end
% glossary-block:B00080:begin
\needspace{3\baselineskip}
\textbf{animal} — \textbf{dēor} (n.). animal; beast. \textbf{nīeten} (n.). animal; beast; cattle. Cattle is the selected plural use; retain the early task distinction from dēor. \textbf{dēor} (n.), \textbf{nīeten} (n.). Dēor is the Unit 1 animal/beast target. Nīeten can also denote an animal in later contexts; use the noun specified by the exercise.\par
% glossary-block:B00080:end
% glossary-block:B00081:begin
\needspace{3\baselineskip}
\textbf{annihilate} — \textbf{fordilgian} (v.). destroy; annihilate.\par
% glossary-block:B00081:end
% glossary-block:B00082:begin
\needspace{3\baselineskip}
\textbf{announce} — \textbf{cȳðan} (v.). announce; make known. Dative person informed; accusative thing announced or a ðæt content clause.\par
% glossary-block:B00082:end
% glossary-block:B00083:begin
\needspace{3\baselineskip}
\textbf{anoint} — \textbf{gesmierwan} (v.). Supplied source form anointed. \textbf{smierwan} (v.). Accusative body part or person, mid + dative material.\par
% glossary-block:B00083:end
% glossary-block:B00084:begin
\needspace{3\baselineskip}
\textbf{another} — \textbf{ōðer} (adj.; ordinal). other; another.\par
% glossary-block:B00084:end
% glossary-block:B00085:begin
\needspace{3\baselineskip}
\textbf{answer} — \textbf{andswarian} (v.). answer; reply. Dative addressee; the reply may be given as direct speech.\par
% glossary-block:B00085:end
% glossary-block:B00086:begin
\needspace{3\baselineskip}
\textbf{anxiety} — \textbf{sorg} (n.). distress; anxiety; trouble; sorrow.\par
% glossary-block:B00086:end
% glossary-block:B00087:begin
\needspace{3\baselineskip}
\textbf{any} — \textbf{ǣnig} (indefinite adjective).  \textbf{hwilc} (interrogative/indefinite adj. or pron.). any; some. Indefinite adjective, distinct from an interrogative. \textbf{ǣnig} (indefinite adjective), \textbf{hwilc} (interrogative/indefinite adj. or pron.). Ǣnig means any; indefinite hwilc has a bounded any/some use. Hwā meaning anyone is a separate personal indefinite, not an adjective to insert before every noun.\par
% glossary-block:B00087:end
% glossary-block:B00088:begin
\needspace{3\baselineskip}
\textbf{any more} — \textbf{furðor} (adv.). further; any more. Adverb in the supplied no-spirit-more expression.\par
% glossary-block:B00088:end
% glossary-block:B00089:begin
\needspace{3\baselineskip}
\textbf{anyone} — \textbf{hwā, hwæt} (interrogative/indefinite pron.).\par
% glossary-block:B00089:end
% glossary-block:B00090:begin
\needspace{3\baselineskip}
\textbf{anything} — \textbf{āwiht} (indefinite substantive).\par
% glossary-block:B00090:end
% glossary-block:B00091:begin
\needspace{3\baselineskip}
\textbf{apart from} — \textbf{būtan} (prep.). besides; apart from. Dative excluded person or thing.\par
% glossary-block:B00091:end
% glossary-block:B00092:begin
\needspace{3\baselineskip}
\textbf{Apollo} — \textbf{Apollines} (proper name; source nom. sg.). The god, distinct from Apollonius.\par
% glossary-block:B00092:end
% glossary-block:B00093:begin
\needspace{3\baselineskip}
\textbf{Apollonius} — \textbf{Apollonius} (proper name). Supplied dative person in the assigned source expression.\par
% glossary-block:B00093:end
% glossary-block:B00094:begin
\needspace{3\baselineskip}
\textbf{appeal} — \textbf{bēn} (n.). petition; appeal.\par
% glossary-block:B00094:end
% glossary-block:B00095:begin
\needspace{3\baselineskip}
\textbf{appear} — \textbf{ætēowian} (v.).\par
% glossary-block:B00095:end
% glossary-block:B00096:begin
\needspace{3\baselineskip}
\textbf{appearance} — \textbf{ansīen} (n.).\par
% glossary-block:B00096:end
% glossary-block:B00097:begin
\needspace{3\baselineskip}
\textbf{appoint} — \textbf{gesęttan} (v.). appoint as lord. Supplied frame with person appointed, tō + dative role, and dative people governed. \textbf{sęttan} (v.). appoint; establish. With accusative entity/decree in the assigned frame.\par
% glossary-block:B00097:end
% glossary-block:B00098:begin
\needspace{3\baselineskip}
\textbf{appoint as} — \textbf{gesęttan} (v.). appoint as lord. Supplied frame with person appointed, tō + dative role, and dative people governed. Supplied appointment frame: accusative person, tō + dative role, and dative people governed.\par
% glossary-block:B00098:end
% glossary-block:B00099:begin
\needspace{3\baselineskip}
\textbf{approach} — \textbf{genēalǣcan} (v.). approach; draw near. Dative person, place or event approached. Supplied source form approached. \textbf{nēalǣcan} (v.). approach; draw near. Dative person or place approached when the goal is expressed. \textbf{nēalǣcan} (v.), \textbf{genēalǣcan} (v.). nēalǣcan and genēalǣcan mean approach/draw near; an expressed goal is dative. Preserve the verb selected by the exercise.\par
% glossary-block:B00099:end
% glossary-block:B00100:begin
\needspace{3\baselineskip}
\textbf{approximately} — \textbf{swā-swā} (conj.). about; approximately. Approximation in the supplied source expression. \textbf{swilce} (adv.). about; approximately.\par
% glossary-block:B00100:end
% glossary-block:B00101:begin
\needspace{3\baselineskip}
\textbf{Arcestrates} — \textbf{Arcestrates} (proper name). King of the region in the story; supplied name forms. The king; supplied name forms.\par
% glossary-block:B00101:end
% glossary-block:B00102:begin
\needspace{3\baselineskip}
\textbf{are} — \textbf{wesan} (v.). be. Copular present/past identification or state; other auxiliary constructions are separate senses. be. Supplied plural present copula.\par
% glossary-block:B00102:end
% glossary-block:B00103:begin
\needspace{3\baselineskip}
\textbf{arise} — \textbf{ārīsan} (v.). Intransitive arise; of + dative source in the bounded rise-from-death expression. \textbf{weorðan} (v.). arise; come about. Intransitive with a nominative subject, including the arising of a calm. \textbf{rīsan} (v.). Earlier taught rīsan supplies this command; prefixed ārīsan is separately used later.\par
% glossary-block:B00103:end
% glossary-block:B00104:begin
\needspace{3\baselineskip}
\textbf{arithmetic} — \textbf{gerīmcræft} (n.). arithmetic; computation.\par
% glossary-block:B00104:end
% glossary-block:B00105:begin
\needspace{3\baselineskip}
\textbf{army} — \textbf{hęre} (n.).  \textbf{weorod} (n.). host; army.\par
% glossary-block:B00105:end
% glossary-block:B00106:begin
\needspace{3\baselineskip}
\textbf{around} — \textbf{ābūtan} (prep.). around; about. With accusative.\par
% glossary-block:B00106:end
% glossary-block:B00107:begin
\needspace{3\baselineskip}
\textbf{arrange} — \textbf{gedōn} (v.). make; arrange. The present form/arrangement of the translation in the source.\par
% glossary-block:B00107:end
% glossary-block:B00108:begin
\needspace{3\baselineskip}
\textbf{array} — \textbf{frætwung} (n.).\par
% glossary-block:B00108:end
% glossary-block:B00109:begin
\needspace{3\baselineskip}
\textbf{arrive} — \textbf{becuman} (v.). come; arrive. \textbf{tōcuman} (v.). come; arrive.\par
% glossary-block:B00109:end
% glossary-block:B00110:begin
\needspace{3\baselineskip}
\textbf{as} — \textbf{swā} (adverb; conjunction). Only in sōna swā. as; in the same way as. as; like. Predicative comparison. Swā + adjective/adverb + swā frames equal degree, as … as: the first swā marks degree and the second introduces the standard. \textbf{ðæs-ðe} (conj.). Only in supplied ðæs-ðe mē ðyncð. Swā can introduce a predicative comparison, as/like; swā-swā means just as in its supplied comparison. Swilce means as if in the supplied apparent-rejuvenation comparison. Equal degree uses the separate paired swā … swā construction. \textbf{ðā} (conj.). Temporal as means when in the expressly supplied factual-past clause with ðā; comparison as uses a different construction.\par
% glossary-block:B00110:end
% glossary-block:B00111:begin
\needspace{3\baselineskip}
\textbf{as … as} — Swā + adjective/adverb + swā frames equal degree, as … as: the first swā marks degree and the second introduces the standard. Swā + adjective/adverb + swā frames equal degree, as … as. The first swā marks degree; the second introduces the standard. Comparative the … the uses a separate construction.\par
% glossary-block:B00111:end
% glossary-block:B00112:begin
\needspace{3\baselineskip}
\textbf{ascend} — \textbf{ūpāstīgan} (v.). rise up; ascend.\par
% glossary-block:B00112:end
% glossary-block:B00113:begin
\needspace{3\baselineskip}
\textbf{as far as} — \textbf{ōð} (prep.). as far as; up to.\par
% glossary-block:B00113:end
% glossary-block:B00114:begin
\needspace{3\baselineskip}
\textbf{as if} — \textbf{swilce} (adv.). Supplied apparent-rejuvenation comparison.\par
% glossary-block:B00114:end
% glossary-block:B00115:begin
\needspace{3\baselineskip}
\textbf{as it seems} — The supplied expression ðæs-ðe mē ðyncð means as it seems to me; ðæs-ðe alone introduces the as clause.\par
% glossary-block:B00115:end
% glossary-block:B00116:begin
\needspace{3\baselineskip}
\textbf{ask} — \textbf{āscian} (v.). ask about. Accusative person asked, genitive matter asked about; an indirect question may express the matter. \textbf{biddan} (v.). ask; request. Genitive thing requested; an accusative person may be added.\par
% glossary-block:B00116:end
% glossary-block:B00117:begin
\needspace{3\baselineskip}
\textbf{ask about} — \textbf{āscian} (v.). Accusative person asked, genitive matter asked about; an indirect question may express the matter. Supplied source form asked about.\par
% glossary-block:B00117:end
% glossary-block:B00118:begin
\needspace{3\baselineskip}
\textbf{ask for} — \textbf{giernan} (v.). desire; request; woo. Genitive desired thing/person in the selected frames, including the maiden wooed.\par
% glossary-block:B00118:end
% glossary-block:B00119:begin
\needspace{3\baselineskip}
\textbf{as long as} — The whole phrase swā lange swā means as long as and has inasmuch-as force in the assigned biblical context. Lange alone means long. Swā lange swā is the complete supplied phrase as long as; in the assigned biblical context it has inasmuch-as force. Swā lange swā is the complete supplied phrase as long as; in the assigned biblical context it has inasmuch-as force. Lange alone means long.\par
% glossary-block:B00119:end
% glossary-block:B00120:begin
\needspace{3\baselineskip}
\textbf{as lord} — \textbf{gesęttan} (v.). Supplied appointment frame: accusative person, tō + dative role, and dative people governed.\par
% glossary-block:B00120:end
% glossary-block:B00121:begin
\needspace{3\baselineskip}
\textbf{a smaller amount} — \textbf{lǣs} (degree word). less; a smaller amount.\par
% glossary-block:B00121:end
% glossary-block:B00122:begin
\needspace{3\baselineskip}
\textbf{ass} — \textbf{assa} (n.). ass; donkey.\par
% glossary-block:B00122:end
% glossary-block:B00123:begin
\needspace{3\baselineskip}
\textbf{assail} — \textbf{gestandan} (v.). assail; come upon. Supplied source sense. \textbf{onwinnan} (v.). assail; attack.\par
% glossary-block:B00123:end
% glossary-block:B00124:begin
\needspace{3\baselineskip}
\textbf{Asser} — \textbf{Asser} (proper name).\par
% glossary-block:B00124:end
% glossary-block:B00125:begin
\needspace{3\baselineskip}
\textbf{assist} — \textbf{gefultumian} (v.). help; assist.\par
% glossary-block:B00125:end
% glossary-block:B00126:begin
\needspace{3\baselineskip}
\textbf{as soon as} — \textbf{sōna} (adv.). In sōna swā. Sōna swā is the complete supplied expression as soon as; swā alone does not include the soon meaning. The complete supplied construction uses sōna with ðæs-ðe to mean as soon as. Retain both parts.\par
% glossary-block:B00126:end
% glossary-block:B00127:begin
\needspace{3\baselineskip}
\textbf{as wife} — Tō wīfe is the supplied role phrase as wife/to wife.\par
% glossary-block:B00127:end
% glossary-block:B00128:begin
\needspace{3\baselineskip}
\textbf{at} — \textbf{æt} (prep.).  \textbf{tō} (prep.). at; in response to. Only the exact at-the-prayers expression. at; to. Bounded supplied root expression with genitive dependent and dative root. \textbf{wið} (prep.). by; beside; at. Accusative landmark in the taught proximity phrase.\par
% glossary-block:B00128:end
% glossary-block:B00129:begin
\needspace{3\baselineskip}
\textbf{at first} — \textbf{ǣrest} (adv.). first; at first.\par
% glossary-block:B00129:end
% glossary-block:B00130:begin
\needspace{3\baselineskip}
\textbf{at hand} — \textbf{nēah} (adv.). near; at hand. nēah is the predicate in nigh at hand.\par
% glossary-block:B00130:end
% glossary-block:B00131:begin
\needspace{3\baselineskip}
\textbf{at length} — The supplied whole phrase æt nīehstan means at length. \textbf{nīehst} (adv.). Whole supplied æt nīehstan.\par
% glossary-block:B00131:end
% glossary-block:B00132:begin
\needspace{3\baselineskip}
\textbf{at noon} — The supplied whole phrase tō middes dæges means at noon/midday; middes and dæges are genitives. \textbf{middes} (adverb or particle). Whole supplied phrase tō middes dæges; preserve its genitives instead of applying ordinary tō + dative mechanically.\par
% glossary-block:B00132:end
% glossary-block:B00133:begin
\needspace{3\baselineskip}
\textbf{at once} — \textbf{ætgædere} (adv.). together; at once. Paired with samod in the supplied source. \textbf{sōna} (adv.). immediately; at once. sōna carries soon in this future-reference prompt.\par
% glossary-block:B00133:end
% glossary-block:B00134:begin
\needspace{3\baselineskip}
\textbf{at one's side} — \textbf{on} (prep.). Accusative hand in the bounded placement expression.\par
% glossary-block:B00134:end
% glossary-block:B00135:begin
\needspace{3\baselineskip}
\textbf{attack} — \textbf{onwinnan} (v.). assail; attack.\par
% glossary-block:B00135:end
% glossary-block:B00136:begin
\needspace{3\baselineskip}
\textbf{attend} — \textbf{ðēnian} (v.). serve; attend. Supplied preterite source form.\par
% glossary-block:B00136:end
% glossary-block:B00137:begin
\needspace{3\baselineskip}
\textbf{attendance} — \textbf{ðēnung} (n.). service; attendance.\par
% glossary-block:B00137:end
% glossary-block:B00138:begin
\needspace{3\baselineskip}
\textbf{attention} — \textbf{mōd} (n.). mind; attention.\par
% glossary-block:B00138:end
% glossary-block:B00139:begin
\needspace{3\baselineskip}
\textbf{attentive} — \textbf{behygdig} (adj.). attentive; thoughtful.\par
% glossary-block:B00139:end
% glossary-block:B00140:begin
\needspace{3\baselineskip}
\textbf{attentively} — \textbf{geornlice} (adv.). earnestly; attentively.\par
% glossary-block:B00140:end
% glossary-block:B00141:begin
\needspace{3\baselineskip}
\textbf{at that point} — \textbf{ðonne} (conj.; adv.). then; at that point.\par
% glossary-block:B00141:end
% glossary-block:B00142:begin
\needspace{3\baselineskip}
\textbf{at the left side} — \textbf{on} (prep.). at one's side. Accusative hand in the bounded placement expression.\par
% glossary-block:B00142:end
% glossary-block:B00143:begin
\needspace{3\baselineskip}
\textbf{at the outset} — \textbf{fyrmest} (adv.). first of all; at the outset.\par
% glossary-block:B00143:end
% glossary-block:B00144:begin
\needspace{3\baselineskip}
\textbf{at the right side} — \textbf{on} (prep.). at one's side. Accusative hand in the bounded placement expression.\par
% glossary-block:B00144:end
% glossary-block:B00145:begin
\needspace{3\baselineskip}
\textbf{at the root of} — The supplied phrase uses tō + dative root with a genitive dependent naming what has the root. Retain that full phrase; tō alone means at/to.\par
% glossary-block:B00145:end
% glossary-block:B00146:begin
\needspace{3\baselineskip}
\textbf{at work} — The supplied on + activity-noun phrases express at work/engaged in that activity. Use the whole activity phrase and its supplied form; on alone does not name the activity.\par
% glossary-block:B00146:end
% glossary-block:B00147:begin
\needspace{3\baselineskip}
\textbf{available time} — \textbf{ǣmetta} (n.). leisure; available time.\par
% glossary-block:B00147:end
% glossary-block:B00148:begin
\needspace{3\baselineskip}
\textbf{avoid} — \textbf{flēon} (v.). flee; flee from; avoid. Intransitive or with accusative person/thing avoided.\par
% glossary-block:B00148:end
% glossary-block:B00149:begin
\needspace{3\baselineskip}
\textbf{away} — \textbf{onwęg} (adv.).\par
% glossary-block:B00149:end
% glossary-block:B00150:begin
\needspace{3\baselineskip}
\textbf{a while ago} — The supplied whole phrase hwīlum ǣr means a while ago/formerly. Hwīlum alone has the separate earlier meaning sometimes. Hwīlum ǣr is the supplied whole expression a while ago. Hwīlum alone has the separate earlier meaning sometimes.\par
% glossary-block:B00150:end
% glossary-block:B00151:begin
\needspace{3\baselineskip}
\textbf{awl} — \textbf{ǣl} (n.).\par
% glossary-block:B00151:end
% glossary-block:B00152:begin
\needspace{3\baselineskip}
\textbf{ax} — \textbf{æx} (n.).\par
% glossary-block:B00152:end
% book-navigation: standalone bookmark replaced by volume hierarchy
% glossary-block:B00153:begin
\needspace{5\baselineskip}
\subsection*{B}\BookHeading{nav-glossary-english-0003}{subsubsection}{B}{0}
% glossary-block:B00153:end
% glossary-block:B00154:begin
\textbf{back} — \textbf{eft, ęft} (adv.). again; back. \textbf{ongēan} (adv.; prep.; conj.). back; again.\par
% glossary-block:B00154:end
% glossary-block:B00155:begin
\needspace{3\baselineskip}
\textbf{bad} — \textbf{yfel} (adj.). bad; evil.\par
% glossary-block:B00155:end
% glossary-block:B00156:begin
\needspace{3\baselineskip}
\textbf{badly} — \textbf{yfele} (adv.). badly; wrongly. Adverb modifying doing; distinguish the adjective yfel.\par
% glossary-block:B00156:end
% glossary-block:B00157:begin
\needspace{3\baselineskip}
\textbf{ball} — \textbf{ðoðor} (n.).  \textbf{top} (n.). ball?. Tentative source interpretation; retain the question mark and uncertainty.\par
% glossary-block:B00157:end
% glossary-block:B00158:begin
\needspace{3\baselineskip}
\textbf{banquet} — \textbf{bēorscipe} (n.). feast; banquet. \textbf{gebēorscipe} (n.). feast; banquet. \textbf{gereord} (n.). meal; banquet.\par
% glossary-block:B00158:end
% glossary-block:B00159:begin
\needspace{3\baselineskip}
\textbf{Baptist} — \textbf{fuluhtere} (n.). baptizer; Baptist. Baptist as John's title in the assigned passage.\par
% glossary-block:B00159:end
% glossary-block:B00160:begin
\needspace{3\baselineskip}
\textbf{baptizer} — \textbf{fuluhtere} (n.). baptizer; Baptist. Baptist as John's title in the assigned passage.\par
% glossary-block:B00160:end
% glossary-block:B00161:begin
\needspace{3\baselineskip}
\textbf{barely} — \textbf{unēaðe} (adv.). with difficulty; barely.\par
% glossary-block:B00161:end
% glossary-block:B00162:begin
\needspace{3\baselineskip}
\textbf{bath} — \textbf{bæð} (n.).  \textbf{ðwēal} (n.). bath; bathing place.\par
% glossary-block:B00162:end
% glossary-block:B00163:begin
\needspace{3\baselineskip}
\textbf{bathe} — \textbf{baðian} (v.). Supplied source form bathed.\par
% glossary-block:B00163:end
% glossary-block:B00164:begin
\needspace{3\baselineskip}
\textbf{bathing place} — \textbf{bæðstęde} (n.). bathing place; gymnasium. Gymnasium is the contextual interpretation in the supplied reading. \textbf{ðwēal} (n.). bath; bathing place.\par
% glossary-block:B00164:end
% glossary-block:B00165:begin
\needspace{3\baselineskip}
\textbf{be} — Passive weorðan combines with a past participle for the assigned event or becoming. Lexical weorðan become and copular wesan be have separate uses. \textbf{wesan} (v.). Copular present/past identification or state; other auxiliary constructions are separate senses. Supplied plural present copula. Copular eom/wæs and bēo indicative series; later passive and compound uses require their own construction. Wesan + the assigned intransitive past participle supplies the taught compound construction. With cumen come it can render have/had come; English have is not an extra word to add to that construction. Passive wesan combines with the assigned past participle; be done includes the participial meaning as well as the finite be auxiliary. The supplied progressive construction is wesan + present participle. Ordinary English be doing can also be rendered by a simple finite Old English verb; this compound is not the universal default.\par
% glossary-block:B00165:end
% glossary-block:B00166:begin
\needspace{3\baselineskip}
\textbf{be able to} — \textbf{cunnan} (v.). know how; be able. Learned skill; plain infinitive complement. Mugan expresses capability, can/be able to, with a plain infinitive. Learned ability can instead use cunnan; neither construction takes an inserted tō before that infinitive. The supplied mugan form expresses capability, can/could/might be able to, with a plain infinitive. It is not the permission sense of English may/might. \textbf{mugan} (v.), \textbf{cunnan} (v.). Mugan expresses ability or capability. Cunnan + plain infinitive emphasizes learned ability. A modal complement is a plain infinitive, without tō.\par
% glossary-block:B00166:end
% glossary-block:B00167:begin
\needspace{3\baselineskip}
\textbf{be accustomed} — \textbf{gewunian} (v.). be accustomed; be wont.\par
% glossary-block:B00167:end
% glossary-block:B00168:begin
\needspace{3\baselineskip}
\textbf{be accustomed to} — \textbf{gewunian} (v.). gewunian + plain infinitive; no separate be auxiliary in the model.\par
% glossary-block:B00168:end
% glossary-block:B00169:begin
\needspace{3\baselineskip}
\textbf{be acquainted with} — \textbf{cunnan} (v.). know; be acquainted with. Accusative person known; distinct from learned ability.\par
% glossary-block:B00169:end
% glossary-block:B00170:begin
\needspace{3\baselineskip}
\textbf{be allowed to} — Mōtan expresses permission, may/be allowed to, with a plain infinitive. A wish beginning with English may instead uses the taught independent subjunctive. \textbf{mōtan} (v.). mōtan + plain infinitive; permission/prohibition differs from ability and independent wishes.\par
% glossary-block:B00170:end
% glossary-block:B00171:begin
\needspace{3\baselineskip}
\textbf{be amazed} — \textbf{wafian} (v.). Supplied source form.\par
% glossary-block:B00171:end
% glossary-block:B00172:begin
\needspace{3\baselineskip}
\textbf{bear} — \textbf{āberan} (v.). carry. Supplied carried form; retain the ā-prefix. \textbf{beran} (v.). carry. Accusative person or thing carried. bear; produce. Accusative produce; distinct from physical carrying. Accusative produce; distinguish physical carrying. \textbf{gefęrian} (v.). carry; bear. Supplied passive movement; was carried/borne.\par
% glossary-block:B00172:end
% glossary-block:B00173:begin
\needspace{3\baselineskip}
\textbf{bear fruit} — \textbf{beran} (v.). bear; produce. Accusative produce; distinct from physical carrying. Accusative produce; distinguish physical carrying.\par
% glossary-block:B00173:end
% glossary-block:B00174:begin
\needspace{3\baselineskip}
\textbf{bear in mind} — \textbf{geðęncean} (v.). remember; bear in mind.\par
% glossary-block:B00174:end
% glossary-block:B00175:begin
\needspace{3\baselineskip}
\textbf{bear rule} — \textbf{rīcsian} (v.). rule; have dominion. Used absolutely in the taught frame; no new object government. rīcsian can be used absolutely; English have/bear need no independent translation.\par
% glossary-block:B00175:end
% glossary-block:B00176:begin
\needspace{3\baselineskip}
\textbf{bear seed} — \textbf{beran} (v.). bear; produce. Accusative produce; distinct from physical carrying.\par
% glossary-block:B00176:end
% glossary-block:B00177:begin
\needspace{3\baselineskip}
\textbf{beast} — \textbf{dēor} (n.). animal; beast. \textbf{nīeten} (n.). animal; beast; cattle. Cattle is the selected plural use; retain the early task distinction from dēor. \textbf{dēor} (n.), \textbf{nīeten} (n.). Dēor is the Unit 1 animal/beast target. Nīeten can also denote an animal in later contexts; use the noun specified by the exercise.\par
% glossary-block:B00177:end
% glossary-block:B00178:begin
\needspace{3\baselineskip}
\textbf{be astonished} — \textbf{ofwundrian} (v.). Subject plus wesan and supplied participle, with degree/result tō ðan … ðæt.\par
% glossary-block:B00178:end
% glossary-block:B00179:begin
\needspace{3\baselineskip}
\textbf{beat} — \textbf{bēatan} (v.). beat; strike. Accusative person or thing struck. \textbf{geðrǣstan} (v.). afflict; beat. Supplied participial source use.\par
% glossary-block:B00179:end
% glossary-block:B00180:begin
\needspace{3\baselineskip}
\textbf{be at rest} — \textbf{licgan} (v.). lie; be at rest.\par
% glossary-block:B00180:end
% glossary-block:B00181:begin
\needspace{3\baselineskip}
\textbf{beautiful} — \textbf{fæger} (adj.). beautiful; fair.\par
% glossary-block:B00181:end
% glossary-block:B00182:begin
\needspace{3\baselineskip}
\textbf{be born} — \textbf{ācęnnan} (v.). Source supplies ācęnned in a passive expression; distinguish active beget/produce.\par
% glossary-block:B00182:end
% glossary-block:B00183:begin
\needspace{3\baselineskip}
\textbf{be called} — \textbf{hātan} (v.). be called; be named. Special hātte form in the assigned naming expression. The special form hātte means is/was called or named, according to context; do not substitute an ordinary active hātan form.\par
% glossary-block:B00183:end
% glossary-block:B00184:begin
\needspace{3\baselineskip}
\textbf{because} — \textbf{for-ðām} (conj.).  \textbf{for-ðām-ðe} (conj.).  \textbf{for-ðon-ðe} (conj.; adv. in for-ðon). Causal conjunction. \textbf{for-ðām} (conj.), \textbf{for-ðon-ðe} (conj.; adv. in for-ðon). For-ðām and for-ðon-ðe are causal conjunctions meaning because/for. They are distinct from prepositions meaning because of or for someone’s benefit.\par
% glossary-block:B00184:end
% glossary-block:B00185:begin
\needspace{3\baselineskip}
\textbf{because of} — \textbf{for} (prep.). Dative hindering circumstance. because of; on account of. for + dative cause in this supplied money expression; distinguish causal conjunction because.\par
% glossary-block:B00185:end
% glossary-block:B00186:begin
\needspace{3\baselineskip}
\textbf{become} — \textbf{geweorðan} (v.). come to be; become. become; come about. Passive weorðan combines with a past participle for the assigned event or becoming. Lexical weorðan become and copular wesan be have separate uses. \textbf{weorðan} (v.). Nominative subject and agreeing nominative predicate.\par
% glossary-block:B00186:end
% glossary-block:B00187:begin
\needspace{3\baselineskip}
\textbf{become calm} — \textbf{gesmyltan} (v.). Supplied preterite of the sea becoming calm.\par
% glossary-block:B00187:end
% glossary-block:B00188:begin
\needspace{3\baselineskip}
\textbf{become (passive)} — Passive wesan + past participle describes the assigned state/passive construction; passive weorðan + past participle describes becoming or the event. Simple lexical become is separate.\par
% glossary-block:B00188:end
% glossary-block:B00189:begin
\needspace{3\baselineskip}
\textbf{become still} — \textbf{gestillan} (v.). cease; become still. Intransitive in the source storm relative.\par
% glossary-block:B00189:end
% glossary-block:B00190:begin
\needspace{3\baselineskip}
\textbf{be doing} — The supplied progressive construction is wesan + present participle. Ordinary English be doing can also be rendered by a simple finite Old English verb; this compound is not the universal default.\par
% glossary-block:B00190:end
% glossary-block:B00191:begin
\needspace{3\baselineskip}
\textbf{be done} — Passive wesan combines with the assigned past participle; be done includes the participial meaning as well as the finite be auxiliary.\par
% glossary-block:B00191:end
% glossary-block:B00192:begin
\needspace{3\baselineskip}
\textbf{befall} — \textbf{becuman} (v.). come upon; befall. Retain the source's affected participant. \textbf{gelimpan} (v.). happen; befall. Supplied preterite subjunctive in might befall us. \textbf{geweorðan} (v.). come to pass; befall.\par
% glossary-block:B00192:end
% glossary-block:B00193:begin
\needspace{3\baselineskip}
\textbf{before} — \textbf{ǣr} (conj.). Ǣr introduces before; use the mood required by the taught factual or anticipated clause. \textbf{ǣr} (adv.). before; earlier. \textbf{ætforan} (prep.). before; in front of. Dative landmark/person in the supplied frame; retain ætforan rather than replacing it with beforan. \textbf{be-foran} (prep.). before; ahead of. Dative person, including going or sending ahead of that person. before; in front of. Dative landmark in the static frame. before; in front of. Accusative landmark in the positioning construal; static dative can be a different valid construal. \textbf{tōforan} (prep.). before; above all.\par
% glossary-block:B00193:end
% glossary-block:B00194:begin
\needspace{3\baselineskip}
\textbf{beget} — \textbf{ācęnnan} (v.). beget; bring forth.\par
% glossary-block:B00194:end
% glossary-block:B00195:begin
\needspace{3\baselineskip}
\textbf{begin} — \textbf{āginnan} (v.). Supplied source form began. \textbf{onginnan} (v.). Supplied preterite began. Supplied playing/praising action as complement.\par
% glossary-block:B00195:end
% glossary-block:B00196:begin
\needspace{3\baselineskip}
\textbf{beginning} — \textbf{anginn} (n.).\par
% glossary-block:B00196:end
% glossary-block:B00197:begin
\needspace{3\baselineskip}
\textbf{begotten of} — \textbf{of} (prep.). of; from. Origin of wisdom in the supplied begetting statement.\par
% glossary-block:B00197:end
% glossary-block:B00198:begin
\needspace{3\baselineskip}
\textbf{be gracious toward} — \textbf{glæd} (adj.). glæd in this contextual favour sense, wið + person; not merely emotionally glad.\par
% glossary-block:B00198:end
% glossary-block:B00199:begin
\needspace{3\baselineskip}
\textbf{be hale} — Gesund means well/in good health; hāl means hale/healthy. These adjectives describe the subject in be well/healthy; wel is the separate adverb for doing something well.\par
% glossary-block:B00199:end
% glossary-block:B00200:begin
\needspace{3\baselineskip}
\textbf{behave toward} — \textbf{dōn} (v.). deal with; behave toward. Assigned wið person frame; manner is separately expressed.\par
% glossary-block:B00200:end
% glossary-block:B00201:begin
\needspace{3\baselineskip}
\textbf{behead} — \textbf{behēafdian} (v.). Accusative person beheaded.\par
% glossary-block:B00201:end
% glossary-block:B00202:begin
\needspace{3\baselineskip}
\textbf{be healthy} — Gesund means well/in good health; hāl means hale/healthy. These adjectives describe the subject in be well/healthy; wel is the separate adverb for doing something well.\par
% glossary-block:B00202:end
% glossary-block:B00203:begin
\needspace{3\baselineskip}
\textbf{behold} — \textbf{behealdan} (v.). behold; look at. Accusative object seen. behealdan takes accusative objects; distinguish interjection efne behold. \textbf{efne} (interj.). behold!. Interjection, distinct from the verb behealdan look at.\par
% glossary-block:B00203:end
% glossary-block:B00204:begin
\needspace{3\baselineskip}
\textbf{be hungry} — \textbf{hyngran} (v.). Impersonal construction with accusative experiencer. Hyngran expresses be hungry impersonally, with an accusative experiencer. Alternatively, the adjective hungrig agrees with the subject of be. Use the construction named by the exercise.\par
% glossary-block:B00204:end
% glossary-block:B00205:begin
\needspace{3\baselineskip}
\textbf{be in astonishment} — \textbf{wafian} (v.). be amazed. Supplied source form.\par
% glossary-block:B00205:end
% glossary-block:B00206:begin
\needspace{3\baselineskip}
\textbf{be … -ing} — Simple OE finite forms can render English progressive aspect; wesan + present participle is only one bounded source construction, not a required default.\par
% glossary-block:B00206:end
% glossary-block:B00207:begin
\needspace{3\baselineskip}
\textbf{be it so} — The wish be it so uses independent third-singular subjunctive sīe with hit and swā. It does not use a permission verb.\par
% glossary-block:B00207:end
% glossary-block:B00208:begin
\needspace{3\baselineskip}
\textbf{believe} — \textbf{gelīefan} (v.).\par
% glossary-block:B00208:end
% glossary-block:B00209:begin
\needspace{3\baselineskip}
\textbf{believe in} — \textbf{on} (prep.). in. Accusative person believed/trusted in. \textbf{gelīefan} (v.). gelīefan on + accusative person in this assigned frame.\par
% glossary-block:B00209:end
% glossary-block:B00210:begin
\needspace{3\baselineskip}
\textbf{bellows} — \textbf{bęlg} (n.). Selected plural instrument use.\par
% glossary-block:B00210:end
% glossary-block:B00211:begin
\needspace{3\baselineskip}
\textbf{belong} — \textbf{belimpan} (v.). belong; relate. Supplied participial construction.\par
% glossary-block:B00211:end
% glossary-block:B00212:begin
\needspace{3\baselineskip}
\textbf{belong to} — Tō introduces the dependent in the supplied belonging construction; the verb supplies belong. Retain the complete supplied frame.\par
% glossary-block:B00212:end
% glossary-block:B00213:begin
\needspace{3\baselineskip}
\textbf{below} — \textbf{under} (prep.). under. Dative landmark in the taught static frame.\par
% glossary-block:B00213:end
% glossary-block:B00214:begin
\needspace{3\baselineskip}
\textbf{be lying} — \textbf{licgan} (v.). lie; be lying. Supplied participle describing the possessed persons, not a habban perfect.\par
% glossary-block:B00214:end
% glossary-block:B00215:begin
\needspace{3\baselineskip}
\textbf{be multiplied} — \textbf{gemanigfieldan} (v.). Assigned passive with agreeing gemanigfield/gemanigfielde; no independent active causative frame is assigned. Assigned passive use with wesan and agreeing gemanigfield/gemanigfielde; do not infer an unassigned active causative frame.\par
% glossary-block:B00215:end
% glossary-block:B00216:begin
\needspace{3\baselineskip}
\textbf{be named} — \textbf{hātan} (v.). be called; be named. Special hātte form in the assigned naming expression. The special form hātte means is/was called or named, according to context; do not substitute an ordinary active hātan form.\par
% glossary-block:B00216:end
% glossary-block:B00217:begin
\needspace{3\baselineskip}
\textbf{bend} — \textbf{bīegan} (v.). Supplied past bent; retain the source's unmarked ie in its actual form. \textbf{būgan} (v.).\par
% glossary-block:B00217:end
% glossary-block:B00218:begin
\needspace{3\baselineskip}
\textbf{benefit} — \textbf{fręmman} (v.). Dative beneficiary. Dative person benefited. \textbf{gōd} (n.). good; benefit; kindness. Nominal good in the repayment contrast. \textbf{ðearf} (n.). benefit; good; need. Source contextual nouns; retain the local expression.\par
% glossary-block:B00218:end
% glossary-block:B00219:begin
\needspace{3\baselineskip}
\textbf{benevolent} — \textbf{gōd} (adj.). kind; benevolent.\par
% glossary-block:B00219:end
% glossary-block:B00220:begin
\needspace{3\baselineskip}
\textbf{Benjamin} — \textbf{Beniamin} (proper name).\par
% glossary-block:B00220:end
% glossary-block:B00221:begin
\needspace{3\baselineskip}
\textbf{be obliged to} — Sculan expresses obligation, must/be obliged to, with a plain infinitive. Pure future reference does not automatically require sculan.\par
% glossary-block:B00221:end
% glossary-block:B00222:begin
\needspace{3\baselineskip}
\textbf{be (passive)} — Passive wesan + past participle describes the assigned state/passive construction; passive weorðan + past participle describes becoming or the event. Simple lexical become is separate.\par
% glossary-block:B00222:end
% glossary-block:B00223:begin
\needspace{3\baselineskip}
\textbf{be permitted} — \textbf{ālīefan} (v.). permit; allow. Impersonal permission frame, with dative permitted person; consult the complete supplied expression.\par
% glossary-block:B00223:end
% glossary-block:B00224:begin
\needspace{3\baselineskip}
\textbf{bereft of} — \textbf{berēafian} (v.). despoil; deprive of. Accusative person deprived, genitive thing lost; passive retains that genitive. Passive berēafod with genitive of what is lost; active person deprived is accusative.\par
% glossary-block:B00224:end
% glossary-block:B00225:begin
\needspace{3\baselineskip}
\textbf{beseech} — \textbf{biddan} (v.). beseech; entreat. Accusative person; the requested action can be a ðæt clause. Accusative person beseeched; ðæt + subjunctive gives the requested act. Prayer to God with gebiddan is a different entry.\par
% glossary-block:B00225:end
% glossary-block:B00226:begin
\needspace{3\baselineskip}
\textbf{beset} — \textbf{ymbsęllan} (v.). surround; beset. Supplied source form.\par
% glossary-block:B00226:end
% glossary-block:B00227:begin
\needspace{3\baselineskip}
\textbf{beside} — \textbf{wið} (prep.). by; beside; at. Accusative landmark in the taught proximity phrase.\par
% glossary-block:B00227:end
% glossary-block:B00228:begin
\needspace{3\baselineskip}
\textbf{besides} — \textbf{būtan} (prep.). besides; apart from. Dative excluded person or thing. \textbf{ðǣr-tō-ēacan} (adverbial expression). in addition to that; besides.\par
% glossary-block:B00228:end
% glossary-block:B00229:begin
\needspace{3\baselineskip}
\textbf{be silent} — \textbf{swigian} (v.). Supplied source sense.\par
% glossary-block:B00229:end
% glossary-block:B00230:begin
\needspace{3\baselineskip}
\textbf{best} — \textbf{bętst} (adj.).  Bętst is the superlative adverb best, the superlative of wel well. Its adjectival use is separately identified. For better, bętera is adjectival and bęt adverbial. For best, bętst is the corresponding adjectival or adverbial superlative in the taught uses. Choose degree and part of speech from the task.\par
% glossary-block:B00230:end
% glossary-block:B00231:begin
\needspace{3\baselineskip}
\textbf{bestowal} — \textbf{sęlen} (n.). bounty; bestowal.\par
% glossary-block:B00231:end
% glossary-block:B00232:begin
\needspace{3\baselineskip}
\textbf{be sufficient for} — \textbf{genyhtsumian} (v.). suffice; be sufficient for. Complete supplied construction with dative beneficiary. Supplied complete construction with dative beneficiary; plural books require plural subjunctive.\par
% glossary-block:B00232:end
% glossary-block:B00233:begin
\needspace{3\baselineskip}
\textbf{be to} — In the supplied interpretive question, sculan expresses what events should/were to mean. Retain that reported or prospective interpretation rather than substituting the ordinary obligation sense.\par
% glossary-block:B00233:end
% glossary-block:B00234:begin
\needspace{3\baselineskip}
\textbf{be to come} — \textbf{tō} (prep.). Whole supplied is tō cumanne; bounded be-to use, distinct from purpose after a verb of coming.\par
% glossary-block:B00234:end
% glossary-block:B00235:begin
\needspace{3\baselineskip}
\textbf{better} — \textbf{bętera} (adj.).  Bęt is the comparative adverb better, the comparative of wel well. Use bętera for an adjective describing a noun. For better, bętera is adjectival and bęt adverbial. For best, bętst is the corresponding adjectival or adverbial superlative in the taught uses. Choose degree and part of speech from the task.\par
% glossary-block:B00235:end
% glossary-block:B00236:begin
\needspace{3\baselineskip}
\textbf{between} — \textbf{betwēonan} (adverb or particle). between; to one another. Supplied reciprocal betwēonan him in 45.6. \textbf{betweox} (prep.). among; between.\par
% glossary-block:B00236:end
% glossary-block:B00237:begin
\needspace{3\baselineskip}
\textbf{be unwilling} — \textbf{willan} (v.). be unwilling; refuse. With ne or a licensed negative contraction; followed by plain infinitive. ne + willan, or the licensed contracted negative, with a plain infinitive.\par
% glossary-block:B00237:end
% glossary-block:B00238:begin
\needspace{3\baselineskip}
\textbf{be well} — Gesund means well/in good health; hāl means hale/healthy. These adjectives describe the subject in be well/healthy; wel is the separate adverb for doing something well.\par
% glossary-block:B00238:end
% glossary-block:B00239:begin
\needspace{3\baselineskip}
\textbf{be willing to} — \textbf{willan} (v.). willan + plain infinitive denotes volition; intention does not assert the action occurred.\par
% glossary-block:B00239:end
% glossary-block:B00240:begin
\needspace{3\baselineskip}
\textbf{be wondering} — \textbf{wundrian} (v.). wonder. Supplied source form.\par
% glossary-block:B00240:end
% glossary-block:B00241:begin
\needspace{3\baselineskip}
\textbf{be wont} — \textbf{gewunian} (v.). be accustomed; be wont.\par
% glossary-block:B00241:end
% glossary-block:B00242:begin
\needspace{3\baselineskip}
\textbf{be wont to} — \textbf{gewunian} (v.). gewunian + plain infinitive; no separate be auxiliary in the model.\par
% glossary-block:B00242:end
% glossary-block:B00243:begin
\needspace{3\baselineskip}
\textbf{be yet to come} — Supplied gīet sculon fīf means that five are yet to come, with be expressly omitted. This bounded ellipsis is not a general rule that English shall always takes sculan.\par
% glossary-block:B00243:end
% glossary-block:B00244:begin
\needspace{3\baselineskip}
\textbf{bid} — \textbf{bēodan} (v.). command; bid.\par
% glossary-block:B00244:end
% glossary-block:B00245:begin
\needspace{3\baselineskip}
\textbf{bid farewell} — \textbf{biddan} (v.). wish well; bid farewell. Supplied farewell formula. \textbf{grētan} (v.). greet; take leave of. Accusative person/place; greeting and farewell depend on context. grētan + accusative person/place; no separate take or bid is needed.\par
% glossary-block:B00245:end
% glossary-block:B00246:begin
\needspace{3\baselineskip}
\textbf{bind} — \textbf{bindan} (v.). bind; tie. Accusative person or thing bound.\par
% glossary-block:B00246:end
% glossary-block:B00247:begin
\needspace{3\baselineskip}
\textbf{bird} — \textbf{fugol} (n.). Fowl is collective/plural English here; use the required Old English plural.\par
% glossary-block:B00247:end
% glossary-block:B00248:begin
\needspace{3\baselineskip}
\textbf{birth} — \textbf{gebyrd} (n.).\par
% glossary-block:B00248:end
% glossary-block:B00249:begin
\needspace{3\baselineskip}
\textbf{bishop} — \textbf{bisceop} (n.).\par
% glossary-block:B00249:end
% glossary-block:B00250:begin
\needspace{3\baselineskip}
\textbf{black} — \textbf{blæc} (adj.).\par
% glossary-block:B00250:end
% glossary-block:B00251:begin
\needspace{3\baselineskip}
\textbf{blacksmith} — \textbf{smið} (n.).\par
% glossary-block:B00251:end
% glossary-block:B00252:begin
\needspace{3\baselineskip}
\textbf{blame} — \textbf{tǣlan} (v.). blame; censure. \textbf{wītan} (v.). blame; censure.\par
% glossary-block:B00252:end
% glossary-block:B00253:begin
\needspace{3\baselineskip}
\textbf{blameless} — \textbf{clǣne} (adj.). clean; pure; blameless. Source's blameless is contextual clǣne, not just physical clean.\par
% glossary-block:B00253:end
% glossary-block:B00254:begin
\needspace{3\baselineskip}
\textbf{bless} — \textbf{blētsian} (v.). Supplied third-singular preterite; God is the understood subject. \textbf{geblētsian} (v.).\par
% glossary-block:B00254:end
% glossary-block:B00255:begin
\needspace{3\baselineskip}
\textbf{blessed} — \textbf{ēadig} (adj.). happy; blessed.\par
% glossary-block:B00255:end
% glossary-block:B00256:begin
\needspace{3\baselineskip}
\textbf{blessing} — \textbf{blētsung} (n.).\par
% glossary-block:B00256:end
% glossary-block:B00257:begin
\needspace{3\baselineskip}
\textbf{bliss} — \textbf{drēam} (n.). joy; bliss.\par
% glossary-block:B00257:end
% glossary-block:B00258:begin
\needspace{3\baselineskip}
\textbf{bloom} — \textbf{blōstmian} (v.). blossom; bloom.\par
% glossary-block:B00258:end
% glossary-block:B00259:begin
\needspace{3\baselineskip}
\textbf{blossom} — \textbf{blōstma} (n.).  \textbf{blōstmian} (v.). blossom; bloom.\par
% glossary-block:B00259:end
% glossary-block:B00260:begin
\needspace{3\baselineskip}
\textbf{blow} — \textbf{blāwan} (v.). Intransitive in the assigned bellows sentence.\par
% glossary-block:B00260:end
% glossary-block:B00261:begin
\needspace{3\baselineskip}
\textbf{board} — \textbf{āstīgan} (v.). go aboard; board. Accusative ship boarded in the assigned frame. āstīgan takes the ship as direct accusative in these exercises.\par
% glossary-block:B00261:end
% glossary-block:B00262:begin
\needspace{3\baselineskip}
\textbf{bodily} — \textbf{līchamlic} (adj.).\par
% glossary-block:B00262:end
% glossary-block:B00263:begin
\needspace{3\baselineskip}
\textbf{body} — \textbf{līchama} (n.).\par
% glossary-block:B00263:end
% glossary-block:B00264:begin
\needspace{3\baselineskip}
\textbf{boil} — \textbf{sēoðan} (v.).  \textbf{weallan} (v.). surge; boil. Supplied present participle used adjectivally of the sea.\par
% glossary-block:B00264:end
% glossary-block:B00265:begin
\needspace{3\baselineskip}
\textbf{bold} — \textbf{beald} (adj.). bold; brave.\par
% glossary-block:B00265:end
% glossary-block:B00266:begin
\needspace{3\baselineskip}
\textbf{bone} — \textbf{bān} (n.). bone; ivory. Ivory is elpes bān, literally elephant's bone.\par
% glossary-block:B00266:end
% glossary-block:B00267:begin
\needspace{3\baselineskip}
\textbf{book} — \textbf{bōc} (n.).\par
% glossary-block:B00267:end
% glossary-block:B00268:begin
\needspace{3\baselineskip}
\textbf{book-Latin} — \textbf{Bōclæden} (n.). book-Latin; Latin.\par
% glossary-block:B00268:end
% glossary-block:B00269:begin
\needspace{3\baselineskip}
\textbf{border} — \textbf{gemǣre} (n.). boundary; border.\par
% glossary-block:B00269:end
% glossary-block:B00270:begin
\needspace{3\baselineskip}
\textbf{borders} — \textbf{landgemǣre} (n.). territory; borders.\par
% glossary-block:B00270:end
% glossary-block:B00271:begin
\needspace{3\baselineskip}
\textbf{both} — The assigned whole expression ǣgðer ge … ge means both … and, coordinating the two members. \textbf{bēgen} (numeral).\par
% glossary-block:B00271:end
% glossary-block:B00272:begin
\needspace{3\baselineskip}
\textbf{both ... and} — Ge … ge is the paired expression both … and; the two occurrences coordinate the two members.\par
% glossary-block:B00272:end
% glossary-block:B00273:begin
\needspace{3\baselineskip}
\textbf{both … and} — The assigned whole expression ǣgðer ge … ge means both … and, coordinating the two members. \textbf{ge} (conj.). Paired ge … ge. Ge … ge means both … and; ǣgðer ge … ge is the expanded paired expression. The adjective bēgen both is separate.\par
% glossary-block:B00273:end
% glossary-block:B00274:begin
\needspace{3\baselineskip}
\textbf{boundary} — \textbf{gemǣre} (n.). boundary; border.\par
% glossary-block:B00274:end
% glossary-block:B00275:begin
\needspace{3\baselineskip}
\textbf{bounty} — \textbf{sęlen} (n.). bounty; bestowal.\par
% glossary-block:B00275:end
% glossary-block:B00276:begin
\needspace{3\baselineskip}
\textbf{bow} — \textbf{ālūtan} (v.). bow; incline oneself. Reflexive meaning in the assigned use.\par
% glossary-block:B00276:end
% glossary-block:B00277:begin
\needspace{3\baselineskip}
\textbf{boy} — \textbf{cnapa} (n.). Masculine weak noun; the task's chosen word for lad.\par
% glossary-block:B00277:end
% glossary-block:B00278:begin
\needspace{3\baselineskip}
\textbf{brandish} — \textbf{bregdan} (v.). brandish; draw. \textbf{swingan} (v.). swing; brandish. Supplied past source form.\par
% glossary-block:B00278:end
% glossary-block:B00279:begin
\needspace{3\baselineskip}
\textbf{brave} — \textbf{beald} (adj.). bold; brave.\par
% glossary-block:B00279:end
% glossary-block:B00280:begin
\needspace{3\baselineskip}
\textbf{bread} — \textbf{hlāf} (n.). bread; loaf.\par
% glossary-block:B00280:end
% glossary-block:B00281:begin
\needspace{3\baselineskip}
\textbf{break} — \textbf{tōweorpan} (v.). break; overthrow. Break a commandment in the assigned use.\par
% glossary-block:B00281:end
% glossary-block:B00282:begin
\needspace{3\baselineskip}
\textbf{break up} — \textbf{tōberstan} (v.). break up; go to pieces.\par
% glossary-block:B00282:end
% glossary-block:B00283:begin
\needspace{3\baselineskip}
\textbf{brethren} — \textbf{brōðor} (n.). Brethren is an English plural alias, not another Old English lexeme.\par
% glossary-block:B00283:end
% glossary-block:B00284:begin
\needspace{3\baselineskip}
\textbf{bride} — \textbf{brȳd} (n.).\par
% glossary-block:B00284:end
% glossary-block:B00285:begin
\needspace{3\baselineskip}
\textbf{bridge} — \textbf{brycg} (n.).\par
% glossary-block:B00285:end
% glossary-block:B00286:begin
\needspace{3\baselineskip}
\textbf{brief} — \textbf{medmicel} (adj.). short; brief.\par
% glossary-block:B00286:end
% glossary-block:B00287:begin
\needspace{3\baselineskip}
\textbf{brightness} — \textbf{beorhtnes} (n.).\par
% glossary-block:B00287:end
% glossary-block:B00288:begin
\needspace{3\baselineskip}
\textbf{bring} — \textbf{bringan} (v.). Accusative thing or person brought; tō + dative destination/person, or the assigned bare dative recipient. Cook's explicit carry hint selects bringan with accusative presents and dative recipient. \textbf{gebringan} (v.). Supplied preterite with plural winds as subject; retain the prefix. \textbf{lǣdan} (v.). lead; bring; take. A person is the object in the assigned frame. \textbf{tōgelǣdan} (v.). Accusative thing brought, dative recipient.\par
% glossary-block:B00288:end
% glossary-block:B00289:begin
\needspace{3\baselineskip}
\textbf{bring about} — \textbf{dōn} (v.). cause; bring about. With a finite ðæt clause; its actor is nominative within that clause.\par
% glossary-block:B00289:end
% glossary-block:B00290:begin
\needspace{3\baselineskip}
\textbf{bring forth} — \textbf{ācęnnan} (v.). produce; bring forth. beget; bring forth. \textbf{forð} (adv.). forth. Adverbial particle with the bringing/producing verb. \textbf{tēon} (v.). bring forth; produce. Accusative thing produced with particle forð; use the assigned tēon forms. bring forth; produce. Supplied third-plural present subjunctive, with waters as subject.\par
% glossary-block:B00290:end
% glossary-block:B00291:begin
\needspace{3\baselineskip}
\textbf{Britain} — \textbf{Breoton} (proper name).\par
% glossary-block:B00291:end
% glossary-block:B00292:begin
\needspace{3\baselineskip}
\textbf{Britons} — \textbf{Bryttas} (ethnonym).\par
% glossary-block:B00292:end
% glossary-block:B00293:begin
\needspace{3\baselineskip}
\textbf{broad} — \textbf{brād} (adj.).\par
% glossary-block:B00293:end
% glossary-block:B00294:begin
\needspace{3\baselineskip}
\textbf{brother} — \textbf{brōðor} (n.). Brethren is an English plural alias, not another Old English lexeme.\par
% glossary-block:B00294:end
% glossary-block:B00295:begin
\needspace{3\baselineskip}
\textbf{buffet} — \textbf{gewǣgan} (v.). buffet; shake. Supplied participial source use.\par
% glossary-block:B00295:end
% glossary-block:B00296:begin
\needspace{3\baselineskip}
\textbf{build} — \textbf{gewyrcean} (v.). make; build. \textbf{wyrcean} (v.). make; build. Supplied unprefixed form; preserve explicit task controls. \textbf{wyrcean} (v.), \textbf{gewyrcean} (v.). Wyrcean and gewyrcean mean make/build in these senses. Preserve the prefixed or unprefixed verb selected by the exercise.\par
% glossary-block:B00296:end
% glossary-block:B00297:begin
\needspace{3\baselineskip}
\textbf{building} — \textbf{getimbrung} (n.). building; structure built.\par
% glossary-block:B00297:end
% glossary-block:B00298:begin
\needspace{3\baselineskip}
\textbf{burden} — \textbf{hęfignes} (n.).\par
% glossary-block:B00298:end
% glossary-block:B00299:begin
\needspace{3\baselineskip}
\textbf{burst} — \textbf{berstan} (v.).\par
% glossary-block:B00299:end
% glossary-block:B00300:begin
\needspace{3\baselineskip}
\textbf{but} — \textbf{ac} (conj.).  \textbf{sōðlice} (adv.). now; but. Discourse connective corresponding to source Latin autem, distinct from truly.\par
% glossary-block:B00300:end
% glossary-block:B00301:begin
\needspace{3\baselineskip}
\textbf{buy} — \textbf{bycgan} (v.). Accusative goods bought. \textbf{cēapian} (v.). trade; traffic. Traffic means trade; retain the supplied trading frame.\par
% glossary-block:B00301:end
% glossary-block:B00302:begin
\needspace{3\baselineskip}
\textbf{by} — \textbf{be} (prep.). Body part by which a person is led; supplied grasping phrase. \textbf{for} (prep.). by; for the sake of. Grounds of entreaty in the prayer; keep the supplied phrase. \textbf{ðurh, þurh} (prep.). through; by. Mediated action through a person in the supplied frame. Accusative oath standard. \textbf{wið} (prep.). by; beside; at. Accusative landmark in the taught proximity phrase.\par
% glossary-block:B00302:end
% glossary-block:B00303:begin
\needspace{3\baselineskip}
\textbf{by means of} — \textbf{mid} (prep.; adv. in fixed expressions). with; by means of. Dative instrument or means. with; by means of. \textbf{ðurh, þurh} (prep.). through; by means of. Supplied service/evidence expressions, including Ðurh hwæt.\par
% glossary-block:B00303:end
% glossary-block:B00304:begin
\needspace{3\baselineskip}
\textbf{by night} — On niht is the whole supplied temporal expression by night. Retain the exact phrase.\par
% glossary-block:B00304:end
% glossary-block:B00305:begin
\needspace{3\baselineskip}
\textbf{by permission} — \textbf{be} (prep.). Supplied permission expression.\par
% glossary-block:B00305:end
% glossary-block:B00306:begin
\needspace{3\baselineskip}
\textbf{by that} — \textbf{be} (prep.). by that; from that. Refers to the preceding proposition; supplied expression. by that; on that evidence. Evidence/inference frame, not an unrestricted source preposition.\par
% glossary-block:B00306:end
% glossary-block:B00307:begin
\needspace{3\baselineskip}
\textbf{by the same way} — Ðȳ ilcan wege is the supplied instrumental expression by the same way; no additional preposition is required under the exercise’s omission control.\par
% glossary-block:B00307:end
% glossary-block:B00308:begin
\needspace{3\baselineskip}
\textbf{by the time that} — \textbf{ongēan} (adv.; prep.; conj.). Complete supplied ongēan ðætte clause.\par
% glossary-block:B00308:end
% glossary-block:B00309:begin
\needspace{3\baselineskip}
\textbf{by what} — \textbf{ðurh, þurh} (prep.). through; by means of. Supplied service/evidence expressions, including Ðurh hwæt.\par
% glossary-block:B00309:end
% book-navigation: standalone bookmark replaced by volume hierarchy
% glossary-block:B00310:begin
\needspace{5\baselineskip}
\subsection*{C}\BookHeading{nav-glossary-english-0004}{subsubsection}{C}{0}
% glossary-block:B00310:end
% glossary-block:B00311:begin
\textbf{call} — \textbf{cleopian} (v.). call; summon. Accusative person summoned; distinguish intransitive cry aloud. \textbf{hātan} (v.). call; name. \textbf{sęcgan} (v.). call; declare to be. Object with the stated agreeing predicate.\par
% glossary-block:B00311:end
% glossary-block:B00312:begin
\needspace{3\baselineskip}
\textbf{calm} — \textbf{smylte} (adj.). calm; smooth. \textbf{smyltnes} (n.). calm; calmness. \textbf{stilnes} (n.). stillness; calm. Supplied genitive singular.\par
% glossary-block:B00312:end
% glossary-block:B00313:begin
\needspace{3\baselineskip}
\textbf{calmest} — \textbf{smyltest} (adj.). very calm; calmest. Supplied superlative adjective describing the sea.\par
% glossary-block:B00313:end
% glossary-block:B00314:begin
\needspace{3\baselineskip}
\textbf{calmness} — \textbf{smyltnes} (n.). calm; calmness.\par
% glossary-block:B00314:end
% glossary-block:B00315:begin
\needspace{3\baselineskip}
\textbf{camel} — \textbf{olfend} (n.).\par
% glossary-block:B00315:end
% glossary-block:B00316:begin
\needspace{3\baselineskip}
\textbf{can} — Mugan expresses capability, can/be able to, with a plain infinitive. Learned ability can instead use cunnan; neither construction takes an inserted tō before that infinitive. The supplied mugan form expresses capability, can/could/might be able to, with a plain infinitive. It is not the permission sense of English may/might. \textbf{mugan} (v.), \textbf{cunnan} (v.). Mugan expresses ability or capability. Cunnan + plain infinitive emphasizes learned ability. A modal complement is a plain infinitive, without tō.\par
% glossary-block:B00316:end
% glossary-block:B00317:begin
\needspace{3\baselineskip}
\textbf{capacity} — \textbf{mǣð} (n.). measure; capacity.\par
% glossary-block:B00317:end
% glossary-block:B00318:begin
\needspace{3\baselineskip}
\textbf{cares} — \textbf{bisgu} (n.). occupations; cares. \textbf{woruldbisgu} (n.). worldly occupations; cares.\par
% glossary-block:B00318:end
% glossary-block:B00319:begin
\needspace{3\baselineskip}
\textbf{cargo} — \textbf{hlæst} (n.). wares; merchandise; cargo. Selected plural source use.\par
% glossary-block:B00319:end
% glossary-block:B00320:begin
\needspace{3\baselineskip}
\textbf{carpenter} — \textbf{trēowwyrhta} (n.).\par
% glossary-block:B00320:end
% glossary-block:B00321:begin
\needspace{3\baselineskip}
\textbf{carry} — \textbf{āberan} (v.). Supplied carried form; retain the ā-prefix. \textbf{beran} (v.). Accusative person or thing carried. \textbf{gefęrian} (v.). carry; bear. Supplied passive movement; was carried/borne. \textbf{lǣdan} (v.). carry; take. A person is conveyed in this supplied frame. \textbf{bringan} (v.). Cook's explicit carry hint selects bringan with accusative presents and dative recipient.\par
% glossary-block:B00321:end
% glossary-block:B00322:begin
\needspace{3\baselineskip}
\textbf{carry along} — \textbf{lǣdan} (v.). take; carry along.\par
% glossary-block:B00322:end
% glossary-block:B00323:begin
\needspace{3\baselineskip}
\textbf{carry out} — \textbf{dōn} (v.). fulfil; carry out. A word or command is fulfilled.\par
% glossary-block:B00323:end
% glossary-block:B00324:begin
\needspace{3\baselineskip}
\textbf{cast} — \textbf{weorpan} (v.). Accusative thing cast; a following destination has its own government.\par
% glossary-block:B00324:end
% glossary-block:B00325:begin
\needspace{3\baselineskip}
\textbf{cast a net} — \textbf{sęndan} (v.). sęndan takes the net as accusative object; on + accusative supplies the destination.\par
% glossary-block:B00325:end
% glossary-block:B00326:begin
\needspace{3\baselineskip}
\textbf{cast away} — \textbf{āweorpan} (v.). cast down; cast away.\par
% glossary-block:B00326:end
% glossary-block:B00327:begin
\needspace{3\baselineskip}
\textbf{cast down} — \textbf{āweorpan} (v.). cast down; cast away. Accusative people put down, of + dative point from which.\par
% glossary-block:B00327:end
% glossary-block:B00328:begin
\needspace{3\baselineskip}
\textbf{catch} — \textbf{fōn} (v.). seize; catch. Accusative patient in the selected frame. Accusative person overtaken in these tasks; the source's bracketed caught is the route. \textbf{gelǣccan} (v.). Supplied caught form.\par
% glossary-block:B00328:end
% glossary-block:B00329:begin
\needspace{3\baselineskip}
\textbf{cattle} — \textbf{nīeten} (n.). animal; beast; cattle. Cattle is the selected plural use; retain the early task distinction from dēor. The selected nīeten sense is plural cattle; dēor is the separate early animal/beast target.\par
% glossary-block:B00329:end
% glossary-block:B00330:begin
\needspace{3\baselineskip}
\textbf{cause} — \textbf{dōn} (v.). cause; bring about. With a finite ðæt clause; its actor is nominative within that clause. \textbf{intinga} (n.). cause; sake (in for … intingan). Within supplied for intingan expressions.\par
% glossary-block:B00330:end
% glossary-block:B00331:begin
\needspace{3\baselineskip}
\textbf{cause to be} — \textbf{gedōn} (v.). make; cause to be.\par
% glossary-block:B00331:end
% glossary-block:B00332:begin
\needspace{3\baselineskip}
\textbf{cause to stumble} — \textbf{ǣswician} (v.). offend; cause to stumble. Accusative person; ethical/metaphorical offence in the assigned use.\par
% glossary-block:B00332:end
% glossary-block:B00333:begin
\needspace{3\baselineskip}
\textbf{cease} — \textbf{blinnan} (v.). Supplied wæs blinnende periphrasis. \textbf{gestillan} (v.). cease; become still. Intransitive in the source storm relative. \textbf{geswīcan} (v.). Genitive thing ceased from. fail; cease. Intransitive without a genitive complement in the selected frame. Absolute command to the sea; distinct from cessation with a genitive complement.\par
% glossary-block:B00333:end
% glossary-block:B00334:begin
\needspace{3\baselineskip}
\textbf{cease from} — \textbf{geswīcan} (v.). cease. Genitive thing ceased from.\par
% glossary-block:B00334:end
% glossary-block:B00335:begin
\needspace{3\baselineskip}
\textbf{cease to hold} — \textbf{gewītan} (v.). pass away; cease to hold. Of a decree.\par
% glossary-block:B00335:end
% glossary-block:B00336:begin
\needspace{3\baselineskip}
\textbf{cedar} — \textbf{cēderbēam} (n.). cedar; cedar tree.\par
% glossary-block:B00336:end
% glossary-block:B00337:begin
\needspace{3\baselineskip}
\textbf{cedar tree} — \textbf{cēderbēam} (n.). cedar; cedar tree.\par
% glossary-block:B00337:end
% glossary-block:B00338:begin
\needspace{3\baselineskip}
\textbf{censure} — \textbf{tǣlan} (v.). blame; censure. \textbf{wītan} (v.). blame; censure.\par
% glossary-block:B00338:end
% glossary-block:B00339:begin
\needspace{3\baselineskip}
\textbf{chaff} — \textbf{ceaf} (n.).\par
% glossary-block:B00339:end
% glossary-block:B00340:begin
\needspace{3\baselineskip}
\textbf{chamber} — \textbf{būr} (n.). chamber; room.\par
% glossary-block:B00340:end
% glossary-block:B00341:begin
\needspace{3\baselineskip}
\textbf{chance} — \textbf{wēn} (n.). expectation; chance. wēn is ðæt + subjunctive supplies the taught perhaps construction.\par
% glossary-block:B00341:end
% glossary-block:B00342:begin
\needspace{3\baselineskip}
\textbf{chance upon} — \textbf{gefeallan} (v.).  \textbf{on} (prep.). upon; into. Supplied chancing-upon idiom.\par
% glossary-block:B00342:end
% glossary-block:B00343:begin
\needspace{3\baselineskip}
\textbf{change back} — \textbf{gewęndan} (v.). turn; change back. Supplied source form.\par
% glossary-block:B00343:end
% glossary-block:B00344:begin
\needspace{3\baselineskip}
\textbf{chastised} — \textbf{ðrēan} (v.). chastised; struck.\par
% glossary-block:B00344:end
% glossary-block:B00345:begin
\needspace{3\baselineskip}
\textbf{chief} — \textbf{ealdor} (n.). chief; king. This noun is distinct from the homonymous noun meaning life.\par
% glossary-block:B00345:end
% glossary-block:B00346:begin
\needspace{3\baselineskip}
\textbf{chief place} — \textbf{ealdordōm} (n.). chief place; primacy. \textbf{ealdorscipe} (n.). chief place; primacy.\par
% glossary-block:B00346:end
% glossary-block:B00347:begin
\needspace{3\baselineskip}
\textbf{child} — \textbf{bearn} (n.). son; child.\par
% glossary-block:B00347:end
% glossary-block:B00348:begin
\needspace{3\baselineskip}
\textbf{choose} — \textbf{cēosan} (v.). Accusative person or thing chosen. Accusative person selected.\par
% glossary-block:B00348:end
% glossary-block:B00349:begin
\needspace{3\baselineskip}
\textbf{Christ} — \textbf{Crīst} (n.).\par
% glossary-block:B00349:end
% glossary-block:B00350:begin
\needspace{3\baselineskip}
\textbf{church} — \textbf{cirice} (n.).\par
% glossary-block:B00350:end
% glossary-block:B00351:begin
\needspace{3\baselineskip}
\textbf{circuit} — \textbf{ymbhwyrft} (n.). Whole on ymbhwyrfte expression means all around; supplied dependent landmark.\par
% glossary-block:B00351:end
% glossary-block:B00352:begin
\needspace{3\baselineskip}
\textbf{citizens} — \textbf{ceasterwaran} (n.). townspeople; citizens.\par
% glossary-block:B00352:end
% glossary-block:B00353:begin
\needspace{3\baselineskip}
\textbf{city} — \textbf{burg} (n.). town; city. \textbf{ceaster} (n.). town; city.\par
% glossary-block:B00353:end
% glossary-block:B00354:begin
\needspace{3\baselineskip}
\textbf{clarion} — \textbf{biema} (n.). trumpet; clarion.\par
% glossary-block:B00354:end
% glossary-block:B00355:begin
\needspace{3\baselineskip}
\textbf{clean} — \textbf{clǣne} (adj.). pure; clean. clean; pure; blameless.\par
% glossary-block:B00355:end
% glossary-block:B00356:begin
\needspace{3\baselineskip}
\textbf{clearly} — Supplied superlative sweotolost means most clearly; do not use the positive sweotole for this degree. \textbf{sweotollice} (adv.).\par
% glossary-block:B00356:end
% glossary-block:B00357:begin
\needspace{3\baselineskip}
\textbf{cliff} — \textbf{clif} (n.).\par
% glossary-block:B00357:end
% glossary-block:B00358:begin
\needspace{3\baselineskip}
\textbf{cloak} — \textbf{sciccels} (n.).  \textbf{wǣfels} (n.).\par
% glossary-block:B00358:end
% glossary-block:B00359:begin
\needspace{3\baselineskip}
\textbf{close} — \textbf{geęndung} (n.). end; close.\par
% glossary-block:B00359:end
% glossary-block:B00360:begin
\needspace{3\baselineskip}
\textbf{clothe} — \textbf{begyrdan} (v.). gird; clothe. Supplied participial clothing expression. \textbf{bewǣfan} (v.). wrap; clothe. Supplied participial clothing expression. \textbf{gescrȳdan} (v.).  \textbf{scrȳdan} (v.). clothe; dress. Accusative person clothed in this introduced frame. \textbf{wrēon} (v.). cover; clothe.\par
% glossary-block:B00360:end
% glossary-block:B00361:begin
\needspace{3\baselineskip}
\textbf{clothed with} — \textbf{mid} (prep.; adv. in fixed expressions). with. Dative garment in the assigned clothing frame. with. Exact cloth/cloak complements of begyrd and bewǣfed.\par
% glossary-block:B00361:end
% glossary-block:B00362:begin
\needspace{3\baselineskip}
\textbf{clothing} — \textbf{gegierela} (n.). garment; clothing. \textbf{rēaf} (n.). raiment; clothing. clothing; garments. \textbf{rēaf} (n.), \textbf{gegierela} (n.). Wǣd denotes a garment, rēaf raiment/clothing, and gegierela a garment or clothing. Preserve any noun expressly supplied by the exercise.\par
% glossary-block:B00362:end
% glossary-block:B00363:begin
\needspace{3\baselineskip}
\textbf{cloud} — \textbf{wolcen} (n.).\par
% glossary-block:B00363:end
% glossary-block:B00364:begin
\needspace{3\baselineskip}
\textbf{coin} — \textbf{pęning} (n.). penny; coin.\par
% glossary-block:B00364:end
% glossary-block:B00365:begin
\needspace{3\baselineskip}
\textbf{collection} — \textbf{gegaderung} (n.). gathering; collection; congregation. Congregations of waters in the source.\par
% glossary-block:B00365:end
% glossary-block:B00366:begin
\needspace{3\baselineskip}
\textbf{come} — \textbf{becuman} (v.). come; arrive. \textbf{cuman} (v.).  \textbf{tōcuman} (v.). come; arrive.\par
% glossary-block:B00366:end
% glossary-block:B00367:begin
\needspace{3\baselineskip}
\textbf{come about} — \textbf{becuman} (v.). occur; come about. Of the calm in the source. \textbf{gelimpan} (v.). happen; come about. Supplied preterite source form. happen; come about. \textbf{geweorðan} (v.). become; come about. \textbf{weorðan} (v.). arise; come about. Intransitive with a nominative subject, including the arising of a calm. come about; result. Supplied resultative expression.\par
% glossary-block:B00367:end
% glossary-block:B00368:begin
\needspace{3\baselineskip}
\textbf{come nigh} — \textbf{nēalǣcan} (v.), \textbf{genēalǣcan} (v.). nēalǣcan and genēalǣcan mean approach/draw near; an expressed goal is dative. Preserve the verb selected by the exercise.\par
% glossary-block:B00368:end
% glossary-block:B00369:begin
\needspace{3\baselineskip}
\textbf{come out} — \textbf{ūtgangende} (v.). Supplied present participle; source permits its interpretation with earlier action.\par
% glossary-block:B00369:end
% glossary-block:B00370:begin
\needspace{3\baselineskip}
\textbf{come to an end} — \textbf{geęndian} (v.). end; come to an end.\par
% glossary-block:B00370:end
% glossary-block:B00371:begin
\needspace{3\baselineskip}
\textbf{come to be} — \textbf{geweorðan} (v.). come to be; become.\par
% glossary-block:B00371:end
% glossary-block:B00372:begin
\needspace{3\baselineskip}
\textbf{come to naught} — Weorðan tō nāhte is the supplied whole expression come to nothing/naught. Nāhte is the form in this result phrase; nāht alone means nothing.\par
% glossary-block:B00372:end
% glossary-block:B00373:begin
\needspace{3\baselineskip}
\textbf{come to nothing} — Weorðan tō nāhte is the supplied whole result expression come to nothing. Tō alone supplies only the result relation. Weorðan tō nāhte is the supplied whole expression come to nothing/naught. Nāhte is the form in this result phrase; nāht alone means nothing.\par
% glossary-block:B00373:end
% glossary-block:B00374:begin
\needspace{3\baselineskip}
\textbf{come to pass} — \textbf{geweorðan} (v.). come to pass; befall.\par
% glossary-block:B00374:end
% glossary-block:B00375:begin
\needspace{3\baselineskip}
\textbf{come up} — \textbf{ūpgān} (v.). go up; come up. From the water in the supplied expression.\par
% glossary-block:B00375:end
% glossary-block:B00376:begin
\needspace{3\baselineskip}
\textbf{come upon} — \textbf{becuman} (v.). come upon; befall. Retain the source's affected participant. \textbf{gestandan} (v.). assail; come upon. Supplied source sense.\par
% glossary-block:B00376:end
% glossary-block:B00377:begin
\needspace{3\baselineskip}
\textbf{comfort} — \textbf{frōfor} (n.). comfort; consolation.\par
% glossary-block:B00377:end
% glossary-block:B00378:begin
\needspace{3\baselineskip}
\textbf{command} — \textbf{bebēodan} (v.). command; instruct. Dative addressee and a ðæt clause with subjunctive for the commanded action. \textbf{bēodan} (v.). command; bid. \textbf{hǣs} (n.). command; order. \textbf{hātan} (v.). command; order. Accusative actor + plain infinitive for the commanded action; distinct from call/name.\par
% glossary-block:B00378:end
% glossary-block:B00379:begin
\needspace{3\baselineskip}
\textbf{commandment} — \textbf{bebod} (n.).\par
% glossary-block:B00379:end
% glossary-block:B00380:begin
\needspace{3\baselineskip}
\textbf{common} — \textbf{gemǣne} (adj.). common; shared.\par
% glossary-block:B00380:end
% glossary-block:B00381:begin
\needspace{3\baselineskip}
\textbf{companion} — \textbf{gefēra} (n.). companion; fellow.\par
% glossary-block:B00381:end
% glossary-block:B00382:begin
\needspace{3\baselineskip}
\textbf{companionship} — \textbf{gefērscipe} (n.).\par
% glossary-block:B00382:end
% glossary-block:B00383:begin
\needspace{3\baselineskip}
\textbf{company} — \textbf{flocc} (n.). group; company. \textbf{getæl} (n.). number; company. Singular collective getæl with the supplied dependent expression; retain source participial genitives. \textbf{męnigu} (n.). company; crowd. męnigu also denotes a multitude of things in the locally supplied spice expression.\par
% glossary-block:B00383:end
% glossary-block:B00384:begin
\needspace{3\baselineskip}
\textbf{compassion} — \textbf{mildheortnes} (n.).\par
% glossary-block:B00384:end
% glossary-block:B00385:begin
\needspace{3\baselineskip}
\textbf{compassionate} — \textbf{mildheort} (adj.). merciful; compassionate.\par
% glossary-block:B00385:end
% glossary-block:B00386:begin
\needspace{3\baselineskip}
\textbf{complete} — \textbf{fullfręmman} (v.). finish; complete. Supplied passive use.\par
% glossary-block:B00386:end
% glossary-block:B00387:begin
\needspace{3\baselineskip}
\textbf{completely} — Mid ealle is the fixed phrase completely; the individual adjective eall means all/whole. Mid ealle is the fixed phrase completely; mid alone does not have that meaning.\par
% glossary-block:B00387:end
% glossary-block:B00388:begin
\needspace{3\baselineskip}
\textbf{computation} — \textbf{gerīmcræft} (n.). arithmetic; computation.\par
% glossary-block:B00388:end
% glossary-block:B00389:begin
\needspace{3\baselineskip}
\textbf{conceal} — \textbf{diernan} (v.). conceal; hide. \textbf{forswīgian} (v.). keep silent about; conceal. \textbf{hȳdan} (v.). hide; conceal. Accusative thing hidden; wesan forms the assigned passive.\par
% glossary-block:B00389:end
% glossary-block:B00390:begin
\needspace{3\baselineskip}
\textbf{conceive} — \textbf{āsmēan} (v.). Exact supplied infinitive āsmēan in the heart-conceiving expression; not an independently assigned full paradigm. Supplied infinitive āsmēan; the heart is its subject.\par
% glossary-block:B00390:end
% glossary-block:B00391:begin
\needspace{3\baselineskip}
\textbf{concerning} — \textbf{be} (prep.). concerning; about. With dative.\par
% glossary-block:B00391:end
% glossary-block:B00392:begin
\needspace{3\baselineskip}
\textbf{condemn} — \textbf{dēman} (v.). Accusative person condemned.\par
% glossary-block:B00392:end
% glossary-block:B00393:begin
\needspace{3\baselineskip}
\textbf{confidence} — \textbf{trēownes} (n.). trust; confidence.\par
% glossary-block:B00393:end
% glossary-block:B00394:begin
\needspace{3\baselineskip}
\textbf{congregation} — \textbf{gegaderung} (n.). gathering; collection; congregation. Congregations of waters in the source.\par
% glossary-block:B00394:end
% glossary-block:B00395:begin
\needspace{3\baselineskip}
\textbf{conjure} — \textbf{hālsian} (v.). conjure; solemnly charge. Accusative person; ðurh + accusative oath standard; a requested ðæt clause may follow.\par
% glossary-block:B00395:end
% glossary-block:B00396:begin
\needspace{3\baselineskip}
\textbf{conjure by} — \textbf{ðurh, þurh} (prep.). by. Accusative oath standard.\par
% glossary-block:B00396:end
% glossary-block:B00397:begin
\needspace{3\baselineskip}
\textbf{conquer} — \textbf{oferwinnan} (v.). overcome; conquer.\par
% glossary-block:B00397:end
% glossary-block:B00398:begin
\needspace{3\baselineskip}
\textbf{consecrate} — \textbf{gehālgian} (v.). hallow; consecrate.\par
% glossary-block:B00398:end
% glossary-block:B00399:begin
\needspace{3\baselineskip}
\textbf{consider} — \textbf{smēagan} (v.). ponder; consider. Supplied preterite source form.\par
% glossary-block:B00399:end
% glossary-block:B00400:begin
\needspace{3\baselineskip}
\textbf{consolation} — \textbf{frōfor} (n.). comfort; consolation.\par
% glossary-block:B00400:end
% glossary-block:B00401:begin
\needspace{3\baselineskip}
\textbf{constant} — \textbf{singāl} (adj.). continual; constant.\par
% glossary-block:B00401:end
% glossary-block:B00402:begin
\needspace{3\baselineskip}
\textbf{container} — \textbf{fæt} (n.). vessel; container.\par
% glossary-block:B00402:end
% glossary-block:B00403:begin
\needspace{3\baselineskip}
\textbf{contend} — \textbf{feohtan} (v.). fight; contend. Intransitive; wið + accusative opponent.\par
% glossary-block:B00403:end
% glossary-block:B00404:begin
\needspace{3\baselineskip}
\textbf{continual} — \textbf{singāl} (adj.). continual; constant.\par
% glossary-block:B00404:end
% glossary-block:B00405:begin
\needspace{3\baselineskip}
\textbf{continue} — \textbf{ðurhwunian} (v.). Intransitive continuing of a building; distinguish the following dative-governing wiðstandan.\par
% glossary-block:B00405:end
% glossary-block:B00406:begin
\needspace{3\baselineskip}
\textbf{control} — \textbf{geweald} (n.). control; rule; dominion. Ordinary singular; on gewealde is the taught control phrase.\par
% glossary-block:B00406:end
% glossary-block:B00407:begin
\needspace{3\baselineskip}
\textbf{converse with} — \textbf{sprecan} (v.). speak with. mid + dative interlocutor in the taught production frame.\par
% glossary-block:B00407:end
% glossary-block:B00408:begin
\needspace{3\baselineskip}
\textbf{copper} — \textbf{ǣr} (n.).\par
% glossary-block:B00408:end
% glossary-block:B00409:begin
\needspace{3\baselineskip}
\textbf{correctly} — Supplied comparative rihtlicor means more correctly. Retain its source spelling without an inferred length mark.\par
% glossary-block:B00409:end
% glossary-block:B00410:begin
\needspace{3\baselineskip}
\textbf{costly} — \textbf{dēorwurð} (adj.). precious; costly.\par
% glossary-block:B00410:end
% glossary-block:B00411:begin
\needspace{3\baselineskip}
\textbf{could} — The supplied mugan form expresses capability, can/could/might be able to, with a plain infinitive. It is not the permission sense of English may/might. \textbf{mugan} (v.), \textbf{cunnan} (v.). Mugan expresses ability or capability. Cunnan + plain infinitive emphasizes learned ability. A modal complement is a plain infinitive, without tō.\par
% glossary-block:B00411:end
% glossary-block:B00412:begin
\needspace{3\baselineskip}
\textbf{counsellor} — \textbf{geðeahtend} (n.).\par
% glossary-block:B00412:end
% glossary-block:B00413:begin
\needspace{3\baselineskip}
\textbf{count} — \textbf{gerīman} (v.). number; count. Exact supplied infinitive gerīman; not an independently assigned full paradigm. Supplied infinitive gerīman, count; number.\par
% glossary-block:B00413:end
% glossary-block:B00414:begin
\needspace{3\baselineskip}
\textbf{countenance} — \textbf{andwlita} (n.). face; countenance.\par
% glossary-block:B00414:end
% glossary-block:B00415:begin
\needspace{3\baselineskip}
\textbf{country} — \textbf{land} (n.). land; country. \textbf{ðēod} (n.). region; country. \textbf{ēðel} (n.). ēðel in these home-country references; land remains a separate ordinary country word.\par
% glossary-block:B00415:end
% glossary-block:B00416:begin
\needspace{3\baselineskip}
\textbf{cover} — \textbf{oferwyrcean} (v.). overlay; cover. Accusative object, mid + dative material used to cover it. \textbf{wrēon} (v.). cover; clothe.\par
% glossary-block:B00416:end
% glossary-block:B00417:begin
\needspace{3\baselineskip}
\textbf{cowardly} — \textbf{earg} (adj.).\par
% glossary-block:B00417:end
% glossary-block:B00418:begin
\needspace{3\baselineskip}
\textbf{craft} — \textbf{cræft} (n.). trade; occupation; craft.\par
% glossary-block:B00418:end
% glossary-block:B00419:begin
\needspace{3\baselineskip}
\textbf{craftsman} — \textbf{wyrhta} (n.). workman; craftsman.\par
% glossary-block:B00419:end
% glossary-block:B00420:begin
\needspace{3\baselineskip}
\textbf{crag} — \textbf{torr} (n.). Probable contextual sense in Cook's note, distinct from earlier tower.\par
% glossary-block:B00420:end
% glossary-block:B00421:begin
\needspace{3\baselineskip}
\textbf{crash} — \textbf{gebrec} (n.). crash; noise; din. Supplied accusative singular.\par
% glossary-block:B00421:end
% glossary-block:B00422:begin
\needspace{3\baselineskip}
\textbf{create} — \textbf{gescieppan} (v.).  \textbf{scieppan} (v.). create; shape.\par
% glossary-block:B00422:end
% glossary-block:B00423:begin
\needspace{3\baselineskip}
\textbf{created being or thing} — \textbf{gesceaft} (n.). creature; created being or thing.\par
% glossary-block:B00423:end
% glossary-block:B00424:begin
\needspace{3\baselineskip}
\textbf{Creator} — \textbf{scieppend} (n.).  \textbf{wyrhta} (n.). Maker; Creator. Divine title, distinct from ordinary workman.\par
% glossary-block:B00424:end
% glossary-block:B00425:begin
\needspace{3\baselineskip}
\textbf{creature} — \textbf{gesceaft} (n.). creature; created being or thing.\par
% glossary-block:B00425:end
% glossary-block:B00426:begin
\needspace{3\baselineskip}
\textbf{crops} — \textbf{ierð} (n.). crops; produce of the earth.\par
% glossary-block:B00426:end
% glossary-block:B00427:begin
\needspace{3\baselineskip}
\textbf{cross} — \textbf{rōd} (n.).\par
% glossary-block:B00427:end
% glossary-block:B00428:begin
\needspace{3\baselineskip}
\textbf{crowd} — \textbf{męnigu} (n.). company; crowd. męnigu also denotes a multitude of things in the locally supplied spice expression.\par
% glossary-block:B00428:end
% glossary-block:B00429:begin
\needspace{3\baselineskip}
\textbf{crown} — \textbf{cynehelm} (n.).\par
% glossary-block:B00429:end
% glossary-block:B00430:begin
\needspace{3\baselineskip}
\textbf{cruelly} — \textbf{earmlice} (adverb or particle). cruelly; miserably.\par
% glossary-block:B00430:end
% glossary-block:B00431:begin
\needspace{3\baselineskip}
\textbf{cry} — \textbf{cleopian} (v.). cry aloud. Intransitive, with the words cried given as direct speech.\par
% glossary-block:B00431:end
% glossary-block:B00432:begin
\needspace{3\baselineskip}
\textbf{cry aloud} — \textbf{cleopian} (v.). Intransitive, with the words cried given as direct speech.\par
% glossary-block:B00432:end
% glossary-block:B00433:begin
\needspace{3\baselineskip}
\textbf{culpable} — \textbf{scyldig} (adj.). guilty; culpable.\par
% glossary-block:B00433:end
% glossary-block:B00434:begin
\needspace{3\baselineskip}
\textbf{cup} — \textbf{læfel} (n.).\par
% glossary-block:B00434:end
% glossary-block:B00435:begin
\needspace{3\baselineskip}
\textbf{custom} — \textbf{gewuna} (n.). custom; habit.\par
% glossary-block:B00435:end
% glossary-block:B00436:begin
\needspace{3\baselineskip}
\textbf{cut} — \textbf{snīðan} (v.).\par
% glossary-block:B00436:end
% glossary-block:B00437:begin
\needspace{3\baselineskip}
\textbf{cut down} — \textbf{forceorfan} (v.). cut down; cut off. Accusative tree or thing cut.\par
% glossary-block:B00437:end
% glossary-block:B00438:begin
\needspace{3\baselineskip}
\textbf{cut off} — \textbf{āceorfan} (v.). Accusative thing cut; of in the assigned expression is a separate off particle. \textbf{forceorfan} (v.). cut down; cut off. Accusative tree or thing cut.\par
% glossary-block:B00438:end
% glossary-block:B00439:begin
\needspace{3\baselineskip}
\textbf{Cyrenian} — \textbf{Cyrenisc} (adj.). Cyrenian; of Cyrene. Geographical adjective.\par
% glossary-block:B00439:end
% book-navigation: standalone bookmark replaced by volume hierarchy
% glossary-block:B00440:begin
\needspace{5\baselineskip}
\subsection*{D}\BookHeading{nav-glossary-english-0005}{subsubsection}{D}{0}
% glossary-block:B00440:end
% glossary-block:B00441:begin
\textbf{daily} — \textbf{dæghwāmlīc} (adj.).\par
% glossary-block:B00441:end
% glossary-block:B00442:begin
\needspace{3\baselineskip}
\textbf{Danes} — \textbf{Dęne} (n.). Plural people-name.\par
% glossary-block:B00442:end
% glossary-block:B00443:begin
\needspace{3\baselineskip}
\textbf{danger} — \textbf{frēcnes} (n.). danger; peril.\par
% glossary-block:B00443:end
% glossary-block:B00444:begin
\needspace{3\baselineskip}
\textbf{dare to} — \textbf{durran} (v.). durran + plain infinitive; ne dearr denies daring.\par
% glossary-block:B00444:end
% glossary-block:B00445:begin
\needspace{3\baselineskip}
\textbf{darkness} — \textbf{ðīestru} (n.). Grammatically plural ðīestru; no singular is assigned.\par
% glossary-block:B00445:end
% glossary-block:B00446:begin
\needspace{3\baselineskip}
\textbf{daughter} — \textbf{dohtor} (n.).\par
% glossary-block:B00446:end
% glossary-block:B00447:begin
\needspace{3\baselineskip}
\textbf{day} — \textbf{dæg} (n.).\par
% glossary-block:B00447:end
% glossary-block:B00448:begin
\needspace{3\baselineskip}
\textbf{Day of Judgment} — \textbf{dōm} (n.). judgment. dōmes dæg means Day of Judgment.\par
% glossary-block:B00448:end
% glossary-block:B00449:begin
\needspace{3\baselineskip}
\textbf{dead} — \textbf{dēad} (adj.).\par
% glossary-block:B00449:end
% glossary-block:B00450:begin
\needspace{3\baselineskip}
\textbf{dead people} — \textbf{dēad} (adj.). Substantival adjective.\par
% glossary-block:B00450:end
% glossary-block:B00451:begin
\needspace{3\baselineskip}
\textbf{deal ill with} — Use the whole expression dōn earmlice wið + person, deal cruelly/ill with that person. Earmlice expresses the manner; it is not a second verb meaning deal.\par
% glossary-block:B00451:end
% glossary-block:B00452:begin
\needspace{3\baselineskip}
\textbf{deal with} — \textbf{dōn} (v.). deal with; behave toward. Assigned wið person frame; manner is separately expressed.\par
% glossary-block:B00452:end
% glossary-block:B00453:begin
\needspace{3\baselineskip}
\textbf{dear} — \textbf{lēof} (adj.).\par
% glossary-block:B00453:end
% glossary-block:B00454:begin
\needspace{3\baselineskip}
\textbf{dearer} — \textbf{lēofra} (adj.).\par
% glossary-block:B00454:end
% glossary-block:B00455:begin
\needspace{3\baselineskip}
\textbf{dearest} — \textbf{lēofost} (adj.).\par
% glossary-block:B00455:end
% glossary-block:B00456:begin
\needspace{3\baselineskip}
\textbf{death} — \textbf{dēað} (n.).\par
% glossary-block:B00456:end
% glossary-block:B00457:begin
\needspace{3\baselineskip}
\textbf{debt} — \textbf{gylt} (n.). Includes the prayer's moral debt.\par
% glossary-block:B00457:end
% glossary-block:B00458:begin
\needspace{3\baselineskip}
\textbf{debtor} — \textbf{gyltend} (n.).\par
% glossary-block:B00458:end
% glossary-block:B00459:begin
\needspace{3\baselineskip}
\textbf{declare to be} — \textbf{sęcgan} (v.). call; declare to be. Object with the stated agreeing predicate.\par
% glossary-block:B00459:end
% glossary-block:B00460:begin
\needspace{3\baselineskip}
\textbf{deed} — \textbf{dǣd} (n.).  \textbf{weorc} (n.). weorc includes deed/action in these prompts; dǣd is a separate lexical option elsewhere.\par
% glossary-block:B00460:end
% glossary-block:B00461:begin
\needspace{3\baselineskip}
\textbf{deeds} — \textbf{weorc} (n.). weorc includes deed/action in these prompts; dǣd is a separate lexical option elsewhere.\par
% glossary-block:B00461:end
% glossary-block:B00462:begin
\needspace{3\baselineskip}
\textbf{defence} — \textbf{gebeorg} (n.). protection; defence.\par
% glossary-block:B00462:end
% glossary-block:B00463:begin
\needspace{3\baselineskip}
\textbf{deity} — \textbf{God} (n.). god; deity. Common deity sense, distinct from the proper divine name.\par
% glossary-block:B00463:end
% glossary-block:B00464:begin
\needspace{3\baselineskip}
\textbf{delay} — \textbf{ęlcung} (n.).\par
% glossary-block:B00464:end
% glossary-block:B00465:begin
\needspace{3\baselineskip}
\textbf{deliver} — \textbf{ālīesan} (v.). deliver; release. Accusative person delivered or released; fram + dative marks separation.\par
% glossary-block:B00465:end
% glossary-block:B00466:begin
\needspace{3\baselineskip}
\textbf{demon} — \textbf{dēofol} (n.). devil; demon.\par
% glossary-block:B00466:end
% glossary-block:B00467:begin
\needspace{3\baselineskip}
\textbf{deny} — \textbf{forwiernan} (v.). deny; refuse. Dative person denied, genitive thing denied. Dative person denied, genitive thing denied.\par
% glossary-block:B00467:end
% glossary-block:B00468:begin
\needspace{3\baselineskip}
\textbf{depart} — \textbf{āgān} (v.). go away. Supplied gone-away form; distinct from āgan possess. \textbf{gewītan} (v.). depart; go away. Intransitive; fram + dative departure point is optional. The assigned frame requires no reflexive.\par
% glossary-block:B00468:end
% glossary-block:B00469:begin
\needspace{3\baselineskip}
\textbf{departure} — \textbf{forðfōr} (n.).\par
% glossary-block:B00469:end
% glossary-block:B00470:begin
\needspace{3\baselineskip}
\textbf{deprived of} — \textbf{berēafian} (v.). Passive berēafod with genitive of what is lost; active person deprived is accusative.\par
% glossary-block:B00470:end
% glossary-block:B00471:begin
\needspace{3\baselineskip}
\textbf{deprive of} — \textbf{berēafian} (v.). despoil; deprive of. Accusative person deprived, genitive thing lost; passive retains that genitive.\par
% glossary-block:B00471:end
% glossary-block:B00472:begin
\needspace{3\baselineskip}
\textbf{descend} — \textbf{āstīgan} (v.). Intransitive descend; of + dative departure point.\par
% glossary-block:B00472:end
% glossary-block:B00473:begin
\needspace{3\baselineskip}
\textbf{descendants} — \textbf{sǣd} (n.). seed; descendants. Collective descendants, distinct from literal seed.\par
% glossary-block:B00473:end
% glossary-block:B00474:begin
\needspace{3\baselineskip}
\textbf{deserve} — \textbf{āwyrcean} (v.). do; deserve. In the source, deserving through one's conduct; retain the complete local frame.\par
% glossary-block:B00474:end
% glossary-block:B00475:begin
\needspace{3\baselineskip}
\textbf{desire} — \textbf{gewilnian} (v.). desire; wish. \textbf{giernan} (v.). desire; request; woo. Genitive desired thing/person in the selected frames, including the maiden wooed. \textbf{lystan} (v.). please; be desired. Impersonal source expression; keep the complete supplied frame.\par
% glossary-block:B00475:end
% glossary-block:B00476:begin
\needspace{3\baselineskip}
\textbf{despair} — \textbf{ōrmōdnes} (n.). despair; hopelessness.\par
% glossary-block:B00476:end
% glossary-block:B00477:begin
\needspace{3\baselineskip}
\textbf{despise} — \textbf{rēcelēasian} (v.). neglect; despise.\par
% glossary-block:B00477:end
% glossary-block:B00478:begin
\needspace{3\baselineskip}
\textbf{despoil} — \textbf{berēafian} (v.). despoil; deprive of. Accusative person deprived, genitive thing lost; passive retains that genitive.\par
% glossary-block:B00478:end
% glossary-block:B00479:begin
\needspace{3\baselineskip}
\textbf{despoiled of} — \textbf{berēafian} (v.). Passive berēafod with genitive of what is lost; active person deprived is accusative.\par
% glossary-block:B00479:end
% glossary-block:B00480:begin
\needspace{3\baselineskip}
\textbf{destroy} — \textbf{fordilgian} (v.). destroy; annihilate. \textbf{fordōn} (v.).  \textbf{forniman} (v.). destroy; lay waste.\par
% glossary-block:B00480:end
% glossary-block:B00481:begin
\needspace{3\baselineskip}
\textbf{destruction} — \textbf{lor} (n.). destruction; ruin.\par
% glossary-block:B00481:end
% glossary-block:B00482:begin
\needspace{3\baselineskip}
\textbf{devil} — \textbf{dēofol} (n.). devil; demon.\par
% glossary-block:B00482:end
% glossary-block:B00483:begin
\needspace{3\baselineskip}
\textbf{did} — Dōn can replace an action already mentioned, as do/did in a comparison. This substitute verb is distinct from English auxiliary do in questions and negatives.\par
% glossary-block:B00483:end
% glossary-block:B00484:begin
\needspace{3\baselineskip}
\textbf{did (question or negative helper)} — No dummy dōn: inflect the lexical verb, with ne for negation and the taught question order. Lexical do/perform and substitute dōn are separate.\par
% glossary-block:B00484:end
% glossary-block:B00485:begin
\needspace{3\baselineskip}
\textbf{did (repeat an action)} — \textbf{dōn} (v.). dōn can replace the action already mentioned in a comparison; distinguish modern English dummy do.\par
% glossary-block:B00485:end
% glossary-block:B00486:begin
\needspace{3\baselineskip}
\textbf{dignity} — \textbf{horn} (n.). horn; strength; dignity. Strength/dignity is the figurative force in 37.14, not a replacement for the literal source horn. \textbf{weorðscipe} (n.). honour; dignity.\par
% glossary-block:B00486:end
% glossary-block:B00487:begin
\needspace{3\baselineskip}
\textbf{direct} — \textbf{āðęnian} (v.). direct; stretch toward. Directing the mind toward what is heard; supplied figurative phrase. \textbf{getǣcan} (v.). show; direct. Supplied source form. \textbf{gewissian} (v.). guide; direct.\par
% glossary-block:B00487:end
% glossary-block:B00488:begin
\needspace{3\baselineskip}
\textbf{dirty} — \textbf{horg} (adj.). dirty; filthy.\par
% glossary-block:B00488:end
% glossary-block:B00489:begin
\needspace{3\baselineskip}
\textbf{disciple} — \textbf{leornungcniht} (n.).\par
% glossary-block:B00489:end
% glossary-block:B00490:begin
\needspace{3\baselineskip}
\textbf{discourse} — \textbf{spell} (n.). account; discourse.\par
% glossary-block:B00490:end
% glossary-block:B00491:begin
\needspace{3\baselineskip}
\textbf{dish} — \textbf{disc} (n.).\par
% glossary-block:B00491:end
% glossary-block:B00492:begin
\needspace{3\baselineskip}
\textbf{dismiss} — \textbf{forlǣtan} (v.). let go; dismiss.\par
% glossary-block:B00492:end
% glossary-block:B00493:begin
\needspace{3\baselineskip}
\textbf{disorder} — \textbf{unstilnes} (n.). unrest; disorder.\par
% glossary-block:B00493:end
% glossary-block:B00494:begin
\needspace{3\baselineskip}
\textbf{display} — \textbf{ætīewan} (v.). show; display; reveal.\par
% glossary-block:B00494:end
% glossary-block:B00495:begin
\needspace{3\baselineskip}
\textbf{displease} — \textbf{ofðyncean} (v.). grieve; displease. Impersonal third singular with dative experiencer and a content clause.\par
% glossary-block:B00495:end
% glossary-block:B00496:begin
\needspace{3\baselineskip}
\textbf{disposition} — \textbf{willa} (n.). will; disposition.\par
% glossary-block:B00496:end
% glossary-block:B00497:begin
\needspace{3\baselineskip}
\textbf{dispute} — \textbf{geflīt} (n.). dispute; strife.\par
% glossary-block:B00497:end
% glossary-block:B00498:begin
\needspace{3\baselineskip}
\textbf{distinguish} — \textbf{geweorðian} (v.). distinguish; adorn. Accusative person adorned; mid + dative equipment or adornment. \textbf{tōdǣlen / tōdǣlden (source forms)} (verb; supplied inflected forms). separate; distinguish. Source forms tōdǣlen/tōdǣlden are retained; no silently invented citation spelling.\par
% glossary-block:B00498:end
% glossary-block:B00499:begin
\needspace{3\baselineskip}
\textbf{distress} — \textbf{gewinn} (n.). distress; struggle. \textbf{sorg} (n.). distress; anxiety; trouble; sorrow.\par
% glossary-block:B00499:end
% glossary-block:B00500:begin
\needspace{3\baselineskip}
\textbf{divide} — \textbf{gedǣlan} (v.). divide; separate. Preserve the assigned separation frame. \textbf{tōsceadan} (v.). separate; divide. \textbf{tōtwǣman} (v.). divide; separate. Waters divided from waters in the supplied frame.\par
% glossary-block:B00500:end
% glossary-block:B00501:begin
\needspace{3\baselineskip}
\textbf{divine favour} — \textbf{giefu} (n.). grace; divine favour.\par
% glossary-block:B00501:end
% glossary-block:B00502:begin
\needspace{3\baselineskip}
\textbf{do} — \textbf{āwyrcean} (v.). do; deserve. In the source, deserving through one's conduct; retain the complete local frame. Lexical dōn means do/perform. English auxiliary do/does/did in questions or negatives has no matching dōn: inflect the lexical verb and use the taught question order or ne. \textbf{dōn} (v.). do for; do to. Dative beneficiary/affected person; thing done is accusative when expressed. dōn can replace the action already mentioned in a comparison; distinguish modern English dummy do. Dōn can replace an action already mentioned, as do/did in a comparison. This substitute verb is distinct from English auxiliary do in questions and negatives. \textbf{gedōn} (v.). do; make. Supplied passive fulfilment formula. do; perform. No dummy dōn: inflect the lexical verb, with ne for negation and the taught question order. Lexical do/perform and substitute dōn are separate.\par
% glossary-block:B00502:end
% glossary-block:B00503:begin
\needspace{3\baselineskip}
\textbf{does} — No dummy dōn: inflect the lexical verb, with ne for negation and the taught question order. Lexical do/perform and substitute dōn are separate.\par
% glossary-block:B00503:end
% glossary-block:B00504:begin
\needspace{3\baselineskip}
\textbf{do for} — \textbf{dōn} (v.). do for; do to. Dative beneficiary/affected person; thing done is accusative when expressed.\par
% glossary-block:B00504:end
% glossary-block:B00505:begin
\needspace{3\baselineskip}
\textbf{dog} — \textbf{hund} (n.).\par
% glossary-block:B00505:end
% glossary-block:B00506:begin
\needspace{3\baselineskip}
\textbf{do good to} — \textbf{fręmman} (v.). benefit. Dative beneficiary.\par
% glossary-block:B00506:end
% glossary-block:B00507:begin
\needspace{3\baselineskip}
\textbf{dolphin} — \textbf{męreswīn} (n.).\par
% glossary-block:B00507:end
% glossary-block:B00508:begin
\needspace{3\baselineskip}
\textbf{dominion} — \textbf{geweald} (n.). control; rule; dominion. Ordinary singular; on gewealde is the taught control phrase. \textbf{rīce} (n.). Genitive rīces after geswīcan in ceased to bear rule. Do not insert a literal bear verb.\par
% glossary-block:B00508:end
% glossary-block:B00509:begin
\needspace{3\baselineskip}
\textbf{donkey} — \textbf{assa} (n.). ass; donkey.\par
% glossary-block:B00509:end
% glossary-block:B00510:begin
\needspace{3\baselineskip}
\textbf{doom} — \textbf{dōm} (n.). doom; fate.\par
% glossary-block:B00510:end
% glossary-block:B00511:begin
\needspace{3\baselineskip}
\textbf{do to} — \textbf{dōn} (v.). do for; do to. Dative beneficiary/affected person; thing done is accusative when expressed.\par
% glossary-block:B00511:end
% glossary-block:B00512:begin
\needspace{3\baselineskip}
\textbf{double} — \textbf{twīfeald} (adj.). double; twofold.\par
% glossary-block:B00512:end
% glossary-block:B00513:begin
\needspace{3\baselineskip}
\textbf{dragon} — \textbf{draca} (n.).\par
% glossary-block:B00513:end
% glossary-block:B00514:begin
\needspace{3\baselineskip}
\textbf{draw} — \textbf{bregdan} (v.). brandish; draw. \textbf{tēon} (v.). draw; pull.\par
% glossary-block:B00514:end
% glossary-block:B00515:begin
\needspace{3\baselineskip}
\textbf{draw near} — \textbf{genēalǣcan} (v.). approach; draw near. Dative person, place or event approached. approach. Supplied source form approached. \textbf{nēalǣcan} (v.). approach; draw near. Dative person or place approached when the goal is expressed. \textbf{nēalǣcan} (v.), \textbf{genēalǣcan} (v.). nēalǣcan and genēalǣcan mean approach/draw near; an expressed goal is dative. Preserve the verb selected by the exercise.\par
% glossary-block:B00515:end
% glossary-block:B00516:begin
\needspace{3\baselineskip}
\textbf{draw nigh} — \textbf{nēalǣcan} (v.), \textbf{genēalǣcan} (v.). nēalǣcan and genēalǣcan mean approach/draw near; an expressed goal is dative. Preserve the verb selected by the exercise.\par
% glossary-block:B00516:end
% glossary-block:B00517:begin
\needspace{3\baselineskip}
\textbf{dreadful} — \textbf{ęgeslic} (adj.). dreadful; terrible.\par
% glossary-block:B00517:end
% glossary-block:B00518:begin
\needspace{3\baselineskip}
\textbf{dress} — \textbf{scrȳdan} (v.). clothe; dress. Accusative person clothed in this introduced frame.\par
% glossary-block:B00518:end
% glossary-block:B00519:begin
\needspace{3\baselineskip}
\textbf{dress oneself} — \textbf{scrȳdan} (v.). Supplied reflexive use.\par
% glossary-block:B00519:end
% glossary-block:B00520:begin
\needspace{3\baselineskip}
\textbf{drink} — \textbf{dręnc} (n.).  \textbf{drincan} (v.). Accusative liquid drunk.\par
% glossary-block:B00520:end
% glossary-block:B00521:begin
\needspace{3\baselineskip}
\textbf{drive} — \textbf{drīfan} (v.). Accusative person or thing driven.\par
% glossary-block:B00521:end
% glossary-block:B00522:begin
\needspace{3\baselineskip}
\textbf{drive out} — \textbf{ādrīfan} (v.). drive out; expel. Accusative person expelled; fram + dative may mark the point of separation.\par
% glossary-block:B00522:end
% glossary-block:B00523:begin
\needspace{3\baselineskip}
\textbf{dry land} — \textbf{drȳgnes} (n.).\par
% glossary-block:B00523:end
% glossary-block:B00524:begin
\needspace{3\baselineskip}
\textbf{during} — \textbf{on} (prep.). during; in. Supplied lifetime expression. Dative time expression.\par
% glossary-block:B00524:end
% glossary-block:B00525:begin
\needspace{3\baselineskip}
\textbf{dust} — \textbf{dūst} (n.).\par
% glossary-block:B00525:end
% glossary-block:B00526:begin
\needspace{3\baselineskip}
\textbf{dwell} — \textbf{wunian} (v.). dwell; remain.\par
% glossary-block:B00526:end
% glossary-block:B00527:begin
\needspace{3\baselineskip}
\textbf{dwelling} — \textbf{hām} (n.). home; dwelling. \textbf{hof} (n.).  \textbf{wīc} (n.). dwelling; abode.\par
% glossary-block:B00527:end
% book-navigation: standalone bookmark replaced by volume hierarchy
% glossary-block:B00528:begin
\needspace{5\baselineskip}
\subsection*{E}\BookHeading{nav-glossary-english-0006}{subsubsection}{E}{0}
% glossary-block:B00528:end
% glossary-block:B00529:begin
\textbf{each} — \textbf{ǣghwilc} (indefinite adjective). every; each. each; every. \textbf{ǣlc} (adj.). each; every. With singular noun; substantival each follows the case required by its clause. \textbf{gehwilc} (indefinite adjective). each; every. Ǣlc and gehwilc mean each/every with a singular noun. Personal gehwā means each person/everyone; neuter gehwæt refers to things. Select case from the clause role.\par
% glossary-block:B00529:end
% glossary-block:B00530:begin
\needspace{3\baselineskip}
\textbf{each one} — \textbf{gehwā / gehwæt} (pron.). each one; everyone; everything. Personal gehwā and neuter gehwæt; the early reference status does not by itself establish production.\par
% glossary-block:B00530:end
% glossary-block:B00531:begin
\needspace{3\baselineskip}
\textbf{ear} — \textbf{ēare} (n.).\par
% glossary-block:B00531:end
% glossary-block:B00532:begin
\needspace{3\baselineskip}
\textbf{earlier} — \textbf{ǣr} (adv.). before; earlier. \textbf{ǣrra} (adj.). former; earlier.\par
% glossary-block:B00532:end
% glossary-block:B00533:begin
\needspace{3\baselineskip}
\textbf{earnestly} — \textbf{geornlice} (adv.). earnestly; attentively.\par
% glossary-block:B00533:end
% glossary-block:B00534:begin
\needspace{3\baselineskip}
\textbf{earth} — \textbf{eorðe} (n.).  \textbf{middangeard} (n.). world; earth.\par
% glossary-block:B00534:end
% glossary-block:B00535:begin
\needspace{3\baselineskip}
\textbf{earthly} — \textbf{eorðlic} (adj.).\par
% glossary-block:B00535:end
% glossary-block:B00536:begin
\needspace{3\baselineskip}
\textbf{East Anglians} — \textbf{Eastęngle} (proper name).\par
% glossary-block:B00536:end
% glossary-block:B00537:begin
\needspace{3\baselineskip}
\textbf{eat} — \textbf{etan} (v.).\par
% glossary-block:B00537:end
% glossary-block:B00538:begin
\needspace{3\baselineskip}
\textbf{educate} — \textbf{getēon} (v.). train; educate.\par
% glossary-block:B00538:end
% glossary-block:B00539:begin
\needspace{3\baselineskip}
\textbf{educated} — \textbf{gelǣred} (adj.). educated; trained; skilled.\par
% glossary-block:B00539:end
% glossary-block:B00540:begin
\needspace{3\baselineskip}
\textbf{Egypt} — \textbf{Egypta} (proper name). Egyptians. Egypta land denotes Egypt, literally land of the Egyptians; supplied name forms.\par
% glossary-block:B00540:end
% glossary-block:B00541:begin
\needspace{3\baselineskip}
\textbf{Egyptians} — \textbf{Egypta} (proper name). Egypta land denotes Egypt, literally land of the Egyptians; supplied name forms. \textbf{Egypte} (ethnonym).\par
% glossary-block:B00541:end
% glossary-block:B00542:begin
\needspace{3\baselineskip}
\textbf{eight} — \textbf{eahta} (cardinal numeral). Cardinal 8.\par
% glossary-block:B00542:end
% glossary-block:B00543:begin
\needspace{3\baselineskip}
\textbf{eldest} — \textbf{ieldest} (adj.). eldest; oldest. Substantival person/eldest brother in this source context.\par
% glossary-block:B00543:end
% glossary-block:B00544:begin
\needspace{3\baselineskip}
\textbf{elephant} — \textbf{elp} (n.).\par
% glossary-block:B00544:end
% glossary-block:B00545:begin
\needspace{3\baselineskip}
\textbf{eleven} — \textbf{endlefan} (cardinal numeral). Cardinal 11. \textbf{endleofan} (numeral).\par
% glossary-block:B00545:end
% glossary-block:B00546:begin
\needspace{3\baselineskip}
\textbf{eleventh} — \textbf{endlefta} (ordinal).\par
% glossary-block:B00546:end
% glossary-block:B00547:begin
\needspace{3\baselineskip}
\textbf{empty} — \textbf{ǣmtig} (adj.). empty; void. \textbf{īdel} (adj.). empty; void.\par
% glossary-block:B00547:end
% glossary-block:B00548:begin
\needspace{3\baselineskip}
\textbf{end} — \textbf{ęnde} (n.).  \textbf{geęndian} (v.). end; come to an end. \textbf{geęndung} (n.). end; close.\par
% glossary-block:B00548:end
% glossary-block:B00549:begin
\needspace{3\baselineskip}
\textbf{endure} — \textbf{drēogan} (v.).  \textbf{ðolian} (v.). suffer; endure. Accusative hardship endured. \textbf{ðurhwunian} (v.). remain; endure; last. Intransitive continuing of a building; distinguish the following dative-governing wiðstandan.\par
% glossary-block:B00549:end
% glossary-block:B00550:begin
\needspace{3\baselineskip}
\textbf{enemy} — \textbf{fēond} (n.). enemy; foe. \textbf{gewinna} (n.). enemy; opponent. Supplied plural source forms.\par
% glossary-block:B00550:end
% glossary-block:B00551:begin
\needspace{3\baselineskip}
\textbf{engaged in} — The supplied on + activity-noun phrases express at work/engaged in that activity. Use the whole activity phrase and its supplied form; on alone does not name the activity.\par
% glossary-block:B00551:end
% glossary-block:B00552:begin
\needspace{3\baselineskip}
\textbf{English} — \textbf{Englisc} (adjective used substantivally). English, the language. Language noun. \textbf{Englisc} (adj.).\par
% glossary-block:B00552:end
% glossary-block:B00553:begin
\needspace{3\baselineskip}
\textbf{English people} — \textbf{Engle} (n.). the English. People, distinct from the language/adjective.\par
% glossary-block:B00553:end
% glossary-block:B00554:begin
\needspace{3\baselineskip}
\textbf{enslave} — \textbf{ðēowian} (v.). enslave; reduce to servitude.\par
% glossary-block:B00554:end
% glossary-block:B00555:begin
\needspace{3\baselineskip}
\textbf{enslaved men} — \textbf{ðēowa manna} (noun phrase). enslaved men; servants. Whole supplied phrase ðēowa manna.\par
% glossary-block:B00555:end
% glossary-block:B00556:begin
\needspace{3\baselineskip}
\textbf{enslaved people} — \textbf{ðēow} (adj.). the unfree; enslaved people. Substantival adjective in the source comparison.\par
% glossary-block:B00556:end
% glossary-block:B00557:begin
\needspace{3\baselineskip}
\textbf{enslaved person} — \textbf{ðēow} (n.). servant; enslaved person.\par
% glossary-block:B00557:end
% glossary-block:B00558:begin
\needspace{3\baselineskip}
\textbf{enter} — \textbf{gangan} (v.). go; enter. Supplied circumstantial participial phrase. \textbf{ingān} (v.). enter; go in. enter; go in. Supplied past source form. \textbf{ingangan} (v.). enter; go in. Supplied present-participle phrase.\par
% glossary-block:B00558:end
% glossary-block:B00559:begin
\needspace{3\baselineskip}
\textbf{entreat} — \textbf{biddan} (v.). beseech; entreat. Accusative person; the requested action can be a ðæt clause. \textbf{hālsian} (v.). entreat; implore.\par
% glossary-block:B00559:end
% glossary-block:B00560:begin
\needspace{3\baselineskip}
\textbf{envelop} — \textbf{ymbsęllan} (v.). surround; envelop. Supplied source form.\par
% glossary-block:B00560:end
% glossary-block:B00561:begin
\needspace{3\baselineskip}
\textbf{envoys} — \textbf{ǣrendwreca} (n.). messengers; envoys.\par
% glossary-block:B00561:end
% glossary-block:B00562:begin
\needspace{3\baselineskip}
\textbf{equal} — \textbf{gelīc} (adj.). equal; like. Dative standard of comparison. \textbf{gelīca} (n.). equal; peer.\par
% glossary-block:B00562:end
% glossary-block:B00563:begin
\needspace{3\baselineskip}
\textbf{escape} — \textbf{ætberstan} (v.).\par
% glossary-block:B00563:end
% glossary-block:B00564:begin
\needspace{3\baselineskip}
\textbf{establish} — \textbf{gesęttan} (v.). set; place; establish. \textbf{gestaðelian} (v.). establish; make steadfast. \textbf{sęttan} (v.). appoint; establish. With accusative entity/decree in the assigned frame.\par
% glossary-block:B00564:end
% glossary-block:B00565:begin
\needspace{3\baselineskip}
\textbf{eternal} — \textbf{ēce} (adj.). Adjective ēce, agreeing with the noun it qualifies.\par
% glossary-block:B00565:end
% glossary-block:B00566:begin
\needspace{3\baselineskip}
\textbf{Etheldred} — \textbf{Æðelðrȳð} (proper name). Etheldred; Æthelthryth. Proper name.\par
% glossary-block:B00566:end
% glossary-block:B00567:begin
\needspace{3\baselineskip}
\textbf{even} — \textbf{efne} (interj.). just; precisely; even.\par
% glossary-block:B00567:end
% glossary-block:B00568:begin
\needspace{3\baselineskip}
\textbf{evening} — \textbf{ǣfen} (n.).\par
% glossary-block:B00568:end
% glossary-block:B00569:begin
\needspace{3\baselineskip}
\textbf{even though} — \textbf{ðēah} (conj.; adv.). although; even though. ðēah with the taught concessive clause; a subjunctive need not make an admitted fact unreal. \textbf{ðēah ðe}. Supplied whole concessive conjunction in EX-XII-XIII-43-10.\par
% glossary-block:B00569:end
% glossary-block:B00570:begin
\needspace{3\baselineskip}
\textbf{ever} — \textbf{ā} (adv.). ever; always. \textbf{ǣfre} (adv.). ever; always.\par
% glossary-block:B00570:end
% glossary-block:B00571:begin
\needspace{3\baselineskip}
\textbf{everlasting} — \textbf{ēce} (adj.). Adjective ēce, agreeing with the noun it qualifies.\par
% glossary-block:B00571:end
% glossary-block:B00572:begin
\needspace{3\baselineskip}
\textbf{every} — \textbf{ǣghwilc} (indefinite adjective). every; each. each; every. \textbf{ǣlc} (adj.). each; every. With singular noun; substantival each follows the case required by its clause. \textbf{gehwilc} (indefinite adjective). each; every. Ǣlc and gehwilc mean each/every with a singular noun. Personal gehwā means each person/everyone; neuter gehwæt refers to things. Select case from the clause role.\par
% glossary-block:B00572:end
% glossary-block:B00573:begin
\needspace{3\baselineskip}
\textbf{everyone} — \textbf{gehwā / gehwæt} (pron.). each one; everyone; everything. Personal gehwā and neuter gehwæt; the early reference status does not by itself establish production. Ǣlc and gehwilc mean each/every with a singular noun. Personal gehwā means each person/everyone; neuter gehwæt refers to things. Select case from the clause role.\par
% glossary-block:B00573:end
% glossary-block:B00574:begin
\needspace{3\baselineskip}
\textbf{everything} — \textbf{eall} (adj.; substantive; adv.). all people; everything. Substantival: plural for people, neuter singular for everything. Eall agrees with a noun in all/whole expressions. Substantival plural refers to all people; neuter singular can mean everything. \textbf{gehwā / gehwæt} (pron.). each one; everyone; everything. Personal gehwā and neuter gehwæt; the early reference status does not by itself establish production.\par
% glossary-block:B00574:end
% glossary-block:B00575:begin
\needspace{3\baselineskip}
\textbf{evil} — \textbf{yfel} (adj.). bad; evil. Substantive noun in fram yfele. evil; evil people. Substantival adjective in the assigned use.\par
% glossary-block:B00575:end
% glossary-block:B00576:begin
\needspace{3\baselineskip}
\textbf{evil people} — \textbf{yfel} (adj.). evil; evil people. Substantival adjective in the assigned use.\par
% glossary-block:B00576:end
% glossary-block:B00577:begin
\needspace{3\baselineskip}
\textbf{exactly} — \textbf{eall} (adj.; substantive; adv.). just; exactly. Degree emphasis in eall swā.\par
% glossary-block:B00577:end
% glossary-block:B00578:begin
\needspace{3\baselineskip}
\textbf{exalt} — \textbf{ūp-ā-hębban} (v.). raise; exalt. Accusative people raised; the prefixed verb supplies up.\par
% glossary-block:B00578:end
% glossary-block:B00579:begin
\needspace{3\baselineskip}
\textbf{exalted} — \textbf{hēah} (adj.). high in rank; exalted. Rank/status sense in the supplied source.\par
% glossary-block:B00579:end
% glossary-block:B00580:begin
\needspace{3\baselineskip}
\textbf{example} — \textbf{bȳsen} (n.). example; illustration.\par
% glossary-block:B00580:end
% glossary-block:B00581:begin
\needspace{3\baselineskip}
\textbf{exceedingly} — \textbf{ðearle} (adv.). greatly; exceedingly.\par
% glossary-block:B00581:end
% glossary-block:B00582:begin
\needspace{3\baselineskip}
\textbf{except} — \textbf{būtan} (conj.). Būtan introducing the supplied conditional exception; distinguish a prepositional complement.\par
% glossary-block:B00582:end
% glossary-block:B00583:begin
\needspace{3\baselineskip}
\textbf{except if} — \textbf{būtan} (conj.). unless; except if.\par
% glossary-block:B00583:end
% glossary-block:B00584:begin
\needspace{3\baselineskip}
\textbf{expect} — \textbf{wēnan} (v.). suppose; expect; think. Genitive thing expected in the selected nominal frame; content-clause uses are separately taught.\par
% glossary-block:B00584:end
% glossary-block:B00585:begin
\needspace{3\baselineskip}
\textbf{expectation} — \textbf{wēn} (n.). expectation; chance. wēn is ðæt + subjunctive supplies the taught perhaps construction.\par
% glossary-block:B00585:end
% glossary-block:B00586:begin
\needspace{3\baselineskip}
\textbf{expel} — \textbf{ādrīfan} (v.). drive out; expel. Accusative person expelled; fram + dative may mark the point of separation. \textbf{āflīeman} (v.). put to flight; expel.\par
% glossary-block:B00586:end
% glossary-block:B00587:begin
\needspace{3\baselineskip}
\textbf{explain} — \textbf{geręccan} (v.). explain; render.\par
% glossary-block:B00587:end
% glossary-block:B00588:begin
\needspace{3\baselineskip}
\textbf{extend} — \textbf{brǣdan} (v.). spread; extend. Supplied source form spreads.\par
% glossary-block:B00588:end
% glossary-block:B00589:begin
\needspace{3\baselineskip}
\textbf{eye} — \textbf{ēage} (n.).\par
% glossary-block:B00589:end
% book-navigation: standalone bookmark replaced by volume hierarchy
% glossary-block:B00590:begin
\needspace{5\baselineskip}
\subsection*{F}\BookHeading{nav-glossary-english-0007}{subsubsection}{F}{0}
% glossary-block:B00590:end
% glossary-block:B00591:begin
\textbf{face} — \textbf{andwlita} (n.). face; countenance. \textbf{brādnes} (n.). surface; face. Of the deep in the assigned expression.\par
% glossary-block:B00591:end
% glossary-block:B00592:begin
\needspace{3\baselineskip}
\textbf{fact} — \textbf{ðing} (n.). thing; fact; matter. Matter known/discussed, distinct from a generic physical object.\par
% glossary-block:B00592:end
% glossary-block:B00593:begin
\needspace{3\baselineskip}
\textbf{fail} — \textbf{geswīcan} (v.). fail; cease. Intransitive without a genitive complement in the selected frame.\par
% glossary-block:B00593:end
% glossary-block:B00594:begin
\needspace{3\baselineskip}
\textbf{fair} — \textbf{fæger} (adj.). beautiful; fair. \textbf{gesyndig} (adj.). favourable; fair. Of winds.\par
% glossary-block:B00594:end
% glossary-block:B00595:begin
\needspace{3\baselineskip}
\textbf{fairer} — \textbf{fægerra} (adj.). more beautiful; fairer.\par
% glossary-block:B00595:end
% glossary-block:B00596:begin
\needspace{3\baselineskip}
\textbf{faith} — \textbf{gelēafa} (n.).\par
% glossary-block:B00596:end
% glossary-block:B00597:begin
\needspace{3\baselineskip}
\textbf{fall} — \textbf{āhrēosan} (v.). fall; fall upon. \textbf{feallan} (v.).  \textbf{hrēosan} (v.). Supplied present participle falling.\par
% glossary-block:B00597:end
% glossary-block:B00598:begin
\needspace{3\baselineskip}
\textbf{fall in love} — \textbf{gefeallan} (v.). Figurative falling of the heart; supplied idiom differs in prefix from befeallan in 42.1.\par
% glossary-block:B00598:end
% glossary-block:B00599:begin
\needspace{3\baselineskip}
\textbf{fall in love with} — \textbf{befeallan} (v.). fall in love. Supplied idiom makes lufu love the subject, its possessor the lover, and on + accusative the beloved. \textbf{on} (prep.). upon. Accusative beloved in the supplied love-falls-upon idiom. In the supplied idiom, the lover’s lufu love is the subject of befeallan, and on + accusative identifies the beloved. The whole expression means fall in love with; the person is not the literal subject of fall.\par
% glossary-block:B00599:end
% glossary-block:B00600:begin
\needspace{3\baselineskip}
\textbf{fall upon} — \textbf{āhrēosan} (v.). fall; fall upon.\par
% glossary-block:B00600:end
% glossary-block:B00601:begin
\needspace{3\baselineskip}
\textbf{fame} — \textbf{hlīsa} (n.).\par
% glossary-block:B00601:end
% glossary-block:B00602:begin
\needspace{3\baselineskip}
\textbf{famine} — \textbf{hunger} (n.). famine; hunger.\par
% glossary-block:B00602:end
% glossary-block:B00603:begin
\needspace{3\baselineskip}
\textbf{far} — \textbf{feorr} (adv.). far; far off.\par
% glossary-block:B00603:end
% glossary-block:B00604:begin
\needspace{3\baselineskip}
\textbf{far and wide} — \textbf{wīde} (adv.). widely; far and wide.\par
% glossary-block:B00604:end
% glossary-block:B00605:begin
\needspace{3\baselineskip}
\textbf{farmer} — \textbf{ierðling} (n.). farmer; plowman.\par
% glossary-block:B00605:end
% glossary-block:B00606:begin
\needspace{3\baselineskip}
\textbf{Farne} — \textbf{Farne} (proper name). Island; supplied proper-name form after in.\par
% glossary-block:B00606:end
% glossary-block:B00607:begin
\needspace{3\baselineskip}
\textbf{far off} — \textbf{feorr} (adv.). far; far off.\par
% glossary-block:B00607:end
% glossary-block:B00608:begin
\needspace{3\baselineskip}
\textbf{fate} — \textbf{dōm} (n.). doom; fate.\par
% glossary-block:B00608:end
% glossary-block:B00609:begin
\needspace{3\baselineskip}
\textbf{father} — \textbf{fæder} (n.). Includes the spiritual father relationship in the assigned prompts.\par
% glossary-block:B00609:end
% glossary-block:B00610:begin
\needspace{3\baselineskip}
\textbf{favorable} — \textbf{gesyndig} (adj.). favourable; fair. Of winds.\par
% glossary-block:B00610:end
% glossary-block:B00611:begin
\needspace{3\baselineskip}
\textbf{favorably disposed} — \textbf{glæd} (adj.). gracious; favourably disposed. With wið + person in the supplied gracious-toward frame.\par
% glossary-block:B00611:end
% glossary-block:B00612:begin
\needspace{3\baselineskip}
\textbf{favourable} — \textbf{gesyndig} (adj.). favourable; fair. Of winds.\par
% glossary-block:B00612:end
% glossary-block:B00613:begin
\needspace{3\baselineskip}
\textbf{favourably disposed} — \textbf{glæd} (adj.). gracious; favourably disposed. With wið + person in the supplied gracious-toward frame.\par
% glossary-block:B00613:end
% glossary-block:B00614:begin
\needspace{3\baselineskip}
\textbf{fear} — \textbf{ondrǣdan} (v.). Accusative person feared in 24.6; the supplied reflexive construction in 48.6 has dative self and genitive thing feared. Keep the frames separate.\par
% glossary-block:B00614:end
% glossary-block:B00615:begin
\needspace{3\baselineskip}
\textbf{feast} — \textbf{bēorscipe} (n.). feast; banquet. \textbf{gebēorscipe} (n.). feast; banquet.\par
% glossary-block:B00615:end
% glossary-block:B00616:begin
\needspace{3\baselineskip}
\textbf{feathered} — \textbf{gefiðred} (adj.). Adjective agreeing with the fowl.\par
% glossary-block:B00616:end
% glossary-block:B00617:begin
\needspace{3\baselineskip}
\textbf{feed} — \textbf{fēdan} (v.). feed; nourish. Accusative being fed; food may be mid + dative or a bare instrumental in the assigned expressions.\par
% glossary-block:B00617:end
% glossary-block:B00618:begin
\needspace{3\baselineskip}
\textbf{feet} — \textbf{fōt} (n.). foot. Supplied body-part form feet; distinct from the length-measure sense. foot. Unit of length, distinct from the body-part sense.\par
% glossary-block:B00618:end
% glossary-block:B00619:begin
\needspace{3\baselineskip}
\textbf{fellow} — \textbf{gefēra} (n.). companion; fellow.\par
% glossary-block:B00619:end
% glossary-block:B00620:begin
\needspace{3\baselineskip}
\textbf{fervently} — \textbf{inweardlīce} (adv.). fervently; wholeheartedly.\par
% glossary-block:B00620:end
% glossary-block:B00621:begin
\needspace{3\baselineskip}
\textbf{fetch} — \textbf{feccean} (v.).\par
% glossary-block:B00621:end
% glossary-block:B00622:begin
\needspace{3\baselineskip}
\textbf{few} — \textbf{fēawe} (adj.). few people. Substantival use in the assigned expression.\par
% glossary-block:B00622:end
% glossary-block:B00623:begin
\needspace{3\baselineskip}
\textbf{few people} — \textbf{fēawe} (adj.). Substantival use in the assigned expression.\par
% glossary-block:B00623:end
% glossary-block:B00624:begin
\needspace{3\baselineskip}
\textbf{field} — \textbf{æcer} (n.).\par
% glossary-block:B00624:end
% glossary-block:B00625:begin
\needspace{3\baselineskip}
\textbf{fierce} — \textbf{angrīslic} (adj.). fierce; raging.\par
% glossary-block:B00625:end
% glossary-block:B00626:begin
\needspace{3\baselineskip}
\textbf{fiercely} — \textbf{ðearlice} (adv.). violently; fiercely.\par
% glossary-block:B00626:end
% glossary-block:B00627:begin
\needspace{3\baselineskip}
\textbf{fierceness} — \textbf{rēðnes} (n.). fury; fierceness.\par
% glossary-block:B00627:end
% glossary-block:B00628:begin
\needspace{3\baselineskip}
\textbf{fiercer} — \textbf{swīðra} (adj.). stronger; fiercer. Comparative degree, distinct from right-side sense.\par
% glossary-block:B00628:end
% glossary-block:B00629:begin
\needspace{3\baselineskip}
\textbf{fifth} — \textbf{fīfta} (ordinal numeral).\par
% glossary-block:B00629:end
% glossary-block:B00630:begin
\needspace{3\baselineskip}
\textbf{fifty} — \textbf{fīftig} (numeral).\par
% glossary-block:B00630:end
% glossary-block:B00631:begin
\needspace{3\baselineskip}
\textbf{fight} — \textbf{campian} (v.).  \textbf{feohtan} (v.). fight; contend. Intransitive; wið + accusative opponent.\par
% glossary-block:B00631:end
% glossary-block:B00632:begin
\needspace{3\baselineskip}
\textbf{fill} — \textbf{āfyllan} (v.).  \textbf{gefyllan} (v.). Gefyllan means fill when something is filled; the finish/accomplish sense is separate.\par
% glossary-block:B00632:end
% glossary-block:B00633:begin
\needspace{3\baselineskip}
\textbf{filthy} — \textbf{horg} (adj.). dirty; filthy.\par
% glossary-block:B00633:end
% glossary-block:B00634:begin
\needspace{3\baselineskip}
\textbf{find} — \textbf{findan} (v.). Accusative person or thing found. \textbf{gemētan} (v.). find; meet. Accusative person or thing found/met. Supplied past participle gemētt in a passive. \textbf{mētan} (v.). meet; find.\par
% glossary-block:B00634:end
% glossary-block:B00635:begin
\needspace{3\baselineskip}
\textbf{find out} — \textbf{gewitan} (v.). find out; learn. Knowledge verb, distinct from gewītan depart.\par
% glossary-block:B00635:end
% glossary-block:B00636:begin
\needspace{3\baselineskip}
\textbf{finger} — \textbf{finger} (n.).\par
% glossary-block:B00636:end
% glossary-block:B00637:begin
\needspace{3\baselineskip}
\textbf{finish} — \textbf{fullfręmman} (v.). finish; complete. Supplied passive use. \textbf{gefyllan} (v.). finish; accomplish. A work is the thing completed; distinct from fill. Gefyllan means finish/accomplish when a work is completed; the fill sense is separate.\par
% glossary-block:B00637:end
% glossary-block:B00638:begin
\needspace{3\baselineskip}
\textbf{fire} — \textbf{fȳr} (n.).\par
% glossary-block:B00638:end
% glossary-block:B00639:begin
\needspace{3\baselineskip}
\textbf{fire from heaven} — \textbf{heofonfȳr} (n.). fire from heaven; heavenly fire.\par
% glossary-block:B00639:end
% glossary-block:B00640:begin
\needspace{3\baselineskip}
\textbf{firmament} — \textbf{fæstnes} (n.).\par
% glossary-block:B00640:end
% glossary-block:B00641:begin
\needspace{3\baselineskip}
\textbf{first} — \textbf{ǣrest} (adv.). first; at first. \textbf{ǣresta} (ordinal numeral).  \textbf{fyrmest} (adv.).\par
% glossary-block:B00641:end
% glossary-block:B00642:begin
\needspace{3\baselineskip}
\textbf{first of all} — \textbf{fyrmest} (adv.). first of all; at the outset.\par
% glossary-block:B00642:end
% glossary-block:B00643:begin
\needspace{3\baselineskip}
\textbf{fish} — \textbf{fisc} (n.).\par
% glossary-block:B00643:end
% glossary-block:B00644:begin
\needspace{3\baselineskip}
\textbf{fisher} — \textbf{fiscere} (n.). Masculine noun; fisherman is only one English surface form of this lookup.\par
% glossary-block:B00644:end
% glossary-block:B00645:begin
\needspace{3\baselineskip}
\textbf{fisherman} — \textbf{fiscere} (n.). Masculine noun; fisherman is only one English surface form of this lookup.\par
% glossary-block:B00645:end
% glossary-block:B00646:begin
\needspace{3\baselineskip}
\textbf{fishermen} — \textbf{fiscere} (n.). Masculine noun; fisherman is only one English surface form of this lookup.\par
% glossary-block:B00646:end
% glossary-block:B00647:begin
\needspace{3\baselineskip}
\textbf{fishers} — \textbf{fiscere} (n.). Masculine noun; fisherman is only one English surface form of this lookup.\par
% glossary-block:B00647:end
% glossary-block:B00648:begin
\needspace{3\baselineskip}
\textbf{fish-hook} — \textbf{angel} (n.). Compound English spelling and optional fish- lead to the same fishing noun; distinguish ęngel angel.\par
% glossary-block:B00648:end
% glossary-block:B00649:begin
\needspace{3\baselineskip}
\textbf{fishhook} — \textbf{angel} (n.). Compound English spelling and optional fish- lead to the same fishing noun; distinguish ęngel angel.\par
% glossary-block:B00649:end
% glossary-block:B00650:begin
\needspace{3\baselineskip}
\textbf{fishing} — \textbf{fiscnoð} (n.).  Gān on fiscnoð means go fishing; fiscnoð is the noun fishing and on is part of this supplied phrase.\par
% glossary-block:B00650:end
% glossary-block:B00651:begin
\needspace{3\baselineskip}
\textbf{fish kind} — \textbf{fisccynn} (n.). kinds of fish. Supplied accusative neuter plural.\par
% glossary-block:B00651:end
% glossary-block:B00652:begin
\needspace{3\baselineskip}
\textbf{five} — \textbf{fīf} (cardinal numeral). Cardinal 5.\par
% glossary-block:B00652:end
% glossary-block:B00653:begin
\needspace{3\baselineskip}
\textbf{flee} — \textbf{flēon} (v.). flee; flee from; avoid. Intransitive or with accusative person/thing avoided. Supplied third-plural present; distinct from flēogan fly.\par
% glossary-block:B00653:end
% glossary-block:B00654:begin
\needspace{3\baselineskip}
\textbf{flee from} — \textbf{flēon} (v.). flee; flee from; avoid. Intransitive or with accusative person/thing avoided.\par
% glossary-block:B00654:end
% glossary-block:B00655:begin
\needspace{3\baselineskip}
\textbf{fleet} — \textbf{sciphęre} (n.). fleet; naval force.\par
% glossary-block:B00655:end
% glossary-block:B00656:begin
\needspace{3\baselineskip}
\textbf{flesh} — \textbf{flǣsc} (n.).\par
% glossary-block:B00656:end
% glossary-block:B00657:begin
\needspace{3\baselineskip}
\textbf{flood} — \textbf{flōd} (n.). flood; stream of water.\par
% glossary-block:B00657:end
% glossary-block:B00658:begin
\needspace{3\baselineskip}
\textbf{flow} — \textbf{flōwan} (v.).\par
% glossary-block:B00658:end
% glossary-block:B00659:begin
\needspace{3\baselineskip}
\textbf{fly} — \textbf{flēogan} (v.).\par
% glossary-block:B00659:end
% glossary-block:B00660:begin
\needspace{3\baselineskip}
\textbf{fodder} — \textbf{fōdor} (n.). Provender is routed by Cook's bracketed fodder to fōdor.\par
% glossary-block:B00660:end
% glossary-block:B00661:begin
\needspace{3\baselineskip}
\textbf{foe} — \textbf{fēond} (n.). enemy; foe.\par
% glossary-block:B00661:end
% glossary-block:B00662:begin
\needspace{3\baselineskip}
\textbf{follow} — \textbf{æfter} (prep.; conj. in fixed expressions). after; following. Dative person followed, including motion after that person.\par
% glossary-block:B00662:end
% glossary-block:B00663:begin
\needspace{3\baselineskip}
\textbf{food} — \textbf{ēst} (n.). provisions; food. \textbf{męte} (n.).\par
% glossary-block:B00663:end
% glossary-block:B00664:begin
\needspace{3\baselineskip}
\textbf{foot} — \textbf{fōt} (n.). Supplied body-part form feet; distinct from the length-measure sense. Unit of length, distinct from the body-part sense.\par
% glossary-block:B00664:end
% glossary-block:B00665:begin
\needspace{3\baselineskip}
\textbf{footstool} — \textbf{fōtsceamul} (n.).\par
% glossary-block:B00665:end
% glossary-block:B00666:begin
\needspace{3\baselineskip}
\textbf{for} — \textbf{for} (prep.). for; on behalf of. Dative beneficiary. for; for the sake of. Case follows each complete supplied accusative/dative expression. in exchange for. With dative. for; to the benefit of. \textbf{tō} (prep.). for; to serve as. tō + dative function; a separate dative beneficiary may accompany it. for; to the benefit of. Care of a body part in the supplied expression. \textbf{wið} (prep.). in return for. With dative; distinct from opponent/hostility.\par
% glossary-block:B00666:end
% glossary-block:B00667:begin
\needspace{3\baselineskip}
\textbf{for aid} — Tō + the supplied aid/protection noun forms the whole purpose expression for aid/for protection; the noun supplies the aid/protection meaning.\par
% glossary-block:B00667:end
% glossary-block:B00668:begin
\needspace{3\baselineskip}
\textbf{for a sign} — Tō + a dative noun names the assigned function, for food or for a sign. The noun supplies food/sign; a separate dative beneficiary may also occur. \textbf{tō} (prep.). tō + dative function, with an optional separate dative beneficiary.\par
% glossary-block:B00668:end
% glossary-block:B00669:begin
\needspace{3\baselineskip}
\textbf{for (causal)} — \textbf{for-ðām} (conj.), \textbf{for-ðon-ðe} (conj.; adv. in for-ðon). For-ðām and for-ðon-ðe are causal conjunctions meaning because/for. They are distinct from prepositions meaning because of or for someone’s benefit.\par
% glossary-block:B00669:end
% glossary-block:B00670:begin
\needspace{3\baselineskip}
\textbf{foreign} — \textbf{ęlðēodig} (adj.).\par
% glossary-block:B00670:end
% glossary-block:B00671:begin
\needspace{3\baselineskip}
\textbf{foreigners} — \textbf{ęlðēodig} (adj.). foreigners; strangers. Substantival adjective.\par
% glossary-block:B00671:end
% glossary-block:B00672:begin
\needspace{3\baselineskip}
\textbf{foremost} — \textbf{furðra} (adj.). foremost; prior.\par
% glossary-block:B00672:end
% glossary-block:B00673:begin
\needspace{3\baselineskip}
\textbf{forest} — \textbf{wudu} (n.). wood; forest. Collective vegetation, not a new material sense. \textbf{wudubearu} (n.). forest; grove.\par
% glossary-block:B00673:end
% glossary-block:B00674:begin
\needspace{3\baselineskip}
\textbf{for ever} — Tō worulde is the whole fixed phrase for ever/forever.\par
% glossary-block:B00674:end
% glossary-block:B00675:begin
\needspace{3\baselineskip}
\textbf{forever} — Tō worulde is the whole fixed phrase for ever/forever.\par
% glossary-block:B00675:end
% glossary-block:B00676:begin
\needspace{3\baselineskip}
\textbf{for ever and ever} — \textbf{ā} (adv.). Preserve the source-supplied repetition ā ā ā where given; this is not three separately chosen words.\par
% glossary-block:B00676:end
% glossary-block:B00677:begin
\needspace{3\baselineskip}
\textbf{for food} — Tō + a dative noun names the assigned function, for food or for a sign. The noun supplies food/sign; a separate dative beneficiary may also occur. \textbf{tō} (prep.). tō + dative function, with an optional separate dative beneficiary.\par
% glossary-block:B00677:end
% glossary-block:B00678:begin
\needspace{3\baselineskip}
\textbf{forgive} — \textbf{forgiefan} (v.). Accusative sin, dative person forgiven; in a passive the sin becomes nominative and the person remains dative.\par
% glossary-block:B00678:end
% glossary-block:B00679:begin
\needspace{3\baselineskip}
\textbf{for its part} — \textbf{sōðlice} (adv.). now; but. Discourse connective corresponding to source Latin autem, distinct from truly.\par
% glossary-block:B00679:end
% glossary-block:B00680:begin
\needspace{3\baselineskip}
\textbf{form} — \textbf{hīw} (n.). form; kind. Supplied dative neuter plural.\par
% glossary-block:B00680:end
% glossary-block:B00681:begin
\needspace{3\baselineskip}
\textbf{former} — \textbf{ǣrra} (adj.). former; earlier.\par
% glossary-block:B00681:end
% glossary-block:B00682:begin
\needspace{3\baselineskip}
\textbf{formerly} — The supplied whole phrase hwīlum ǣr means a while ago/formerly. Hwīlum alone has the separate earlier meaning sometimes.\par
% glossary-block:B00682:end
% glossary-block:B00683:begin
\needspace{3\baselineskip}
\textbf{for protection} — Tō + the supplied aid/protection noun forms the whole purpose expression for aid/for protection; the noun supplies the aid/protection meaning.\par
% glossary-block:B00683:end
% glossary-block:B00684:begin
\needspace{3\baselineskip}
\textbf{for (recipient or beneficiary)} — A dative noun/pronoun often expresses the recipient/beneficiary without an added preposition. Distinguish destinations and lexical government.\par
% glossary-block:B00684:end
% glossary-block:B00685:begin
\needspace{3\baselineskip}
\textbf{forsake} — \textbf{forlǣtan} (v.). abandon; forsake. Accusative person or thing left.\par
% glossary-block:B00685:end
% glossary-block:B00686:begin
\needspace{3\baselineskip}
\textbf{forth} — \textbf{forð} (adv.). Adverbial particle with the bringing/producing verb.\par
% glossary-block:B00686:end
% glossary-block:B00687:begin
\needspace{3\baselineskip}
\textbf{for that reason} — \textbf{for-ðon-ðe} (conj.; adv. in for-ðon). therefore; for that reason. Adverbial consequence, distinct from because.\par
% glossary-block:B00687:end
% glossary-block:B00688:begin
\needspace{3\baselineskip}
\textbf{for the benefit of} — \textbf{for} (prep.). Supplied for + dative benefit phrase; its possessed noun controls possessive agreement.\par
% glossary-block:B00688:end
% glossary-block:B00689:begin
\needspace{3\baselineskip}
\textbf{for the sake of} — \textbf{for} (prep.). for; for the sake of. Case follows each complete supplied accusative/dative expression. by; for the sake of. Grounds of entreaty in the prayer; keep the supplied phrase. Intinga means cause/sake. The supplied for intingan expression renders for the sake of; retain its full supplied dependent phrase.\par
% glossary-block:B00689:end
% glossary-block:B00690:begin
\needspace{3\baselineskip}
\textbf{fortification} — \textbf{fæsten} (n.). fortification; stronghold.\par
% glossary-block:B00690:end
% glossary-block:B00691:begin
\needspace{3\baselineskip}
\textbf{forty} — \textbf{fēowertig} (numeral).\par
% glossary-block:B00691:end
% glossary-block:B00692:begin
\needspace{3\baselineskip}
\textbf{for what purpose} — Tō hwon is the whole expression why/for what purpose.\par
% glossary-block:B00692:end
% glossary-block:B00693:begin
\needspace{3\baselineskip}
\textbf{for what reason} — \textbf{for} (prep.). why; for what reason. In for hwȳ or for hwon.\par
% glossary-block:B00693:end
% glossary-block:B00694:begin
\needspace{3\baselineskip}
\textbf{foul} — \textbf{fūl} (adj.). foul; vile.\par
% glossary-block:B00694:end
% glossary-block:B00695:begin
\needspace{3\baselineskip}
\textbf{foundation} — \textbf{grundweall} (n.).\par
% glossary-block:B00695:end
% glossary-block:B00696:begin
\needspace{3\baselineskip}
\textbf{four} — \textbf{fēower} (cardinal numeral). Cardinal 4.\par
% glossary-block:B00696:end
% glossary-block:B00697:begin
\needspace{3\baselineskip}
\textbf{fourth} — \textbf{fēorða} (ordinal numeral).\par
% glossary-block:B00697:end
% glossary-block:B00698:begin
\needspace{3\baselineskip}
\textbf{fowl} — \textbf{fugol} (n.). Fowl is collective/plural English here; use the required Old English plural.\par
% glossary-block:B00698:end
% glossary-block:B00699:begin
\needspace{3\baselineskip}
\textbf{fox} — \textbf{fox} (n.).\par
% glossary-block:B00699:end
% glossary-block:B00700:begin
\needspace{3\baselineskip}
\textbf{free} — \textbf{frēo} (adj.).  Frēo is the adjective free. In go free it describes the subject with appropriate agreement; it is not itself a verb meaning go.\par
% glossary-block:B00700:end
% glossary-block:B00701:begin
\needspace{3\baselineskip}
\textbf{frequently} — \textbf{gelōmlice} (adv.). frequently; regularly. \textbf{genēahhe} (adv.). often; frequently.\par
% glossary-block:B00701:end
% glossary-block:B00702:begin
\needspace{3\baselineskip}
\textbf{friend} — \textbf{frēond} (n.).\par
% glossary-block:B00702:end
% glossary-block:B00703:begin
\needspace{3\baselineskip}
\textbf{frighten} — \textbf{āfǣran} (v.).\par
% glossary-block:B00703:end
% glossary-block:B00704:begin
\needspace{3\baselineskip}
\textbf{from} — \textbf{æt} (prep.). æt + dative person with learn. from; out of. Proceeds of work, in the assigned expression. \textbf{fram} (prep.). Separation, with dative in the taught frame. Only the supplied rest-from-work construction. \textbf{of} (prep.). from; out of. With dative. of; from. Portion from a whole in the supplied partitive phrases; dative whole. Starting condition of old age in the supplied rejuvenation comparison. of; from. Origin of wisdom in the supplied begetting statement. Starting generation in the supplied temporal range.\par
% glossary-block:B00704:end
% glossary-block:B00705:begin
\needspace{3\baselineskip}
\textbf{from afar} — \textbf{feorran} (adv.). feorran; the source supplies the from-afar paraphrase.\par
% glossary-block:B00705:end
% glossary-block:B00706:begin
\needspace{3\baselineskip}
\textbf{from generation to generation} — \textbf{of} (prep.). from. Starting generation in the supplied temporal range. On introduces the succeeding generation in the supplied from-generation-to-generation range. Retain the full supplied temporal phrase; on alone does not mean from generation to generation.\par
% glossary-block:B00706:end
% glossary-block:B00707:begin
\needspace{3\baselineskip}
\textbf{from heaven} — \textbf{heofonlice} (adv.). from heaven; heavenly. Source expression at Cook 143.17.\par
% glossary-block:B00707:end
% glossary-block:B00708:begin
\needspace{3\baselineskip}
\textbf{from here} — \textbf{heonan} (adv.). hence; from here.\par
% glossary-block:B00708:end
% glossary-block:B00709:begin
\needspace{3\baselineskip}
\textbf{from that} — \textbf{be} (prep.). by that; from that. Refers to the preceding proposition; supplied expression. by that; on that evidence. Evidence/inference frame, not an unrestricted source preposition.\par
% glossary-block:B00709:end
% glossary-block:B00710:begin
\needspace{3\baselineskip}
\textbf{from the north} — \textbf{norðan} (adv.).\par
% glossary-block:B00710:end
% glossary-block:B00711:begin
\needspace{3\baselineskip}
\textbf{from there} — \textbf{ðanan} (adv.). thence; from there.\par
% glossary-block:B00711:end
% glossary-block:B00712:begin
\needspace{3\baselineskip}
\textbf{from the south} — \textbf{sūðan} (adv.). in the south; from the south. Retain the supplied spatial expression.\par
% glossary-block:B00712:end
% glossary-block:B00713:begin
\needspace{3\baselineskip}
\textbf{from the west} — \textbf{westan} (adv.).\par
% glossary-block:B00713:end
% glossary-block:B00714:begin
\needspace{3\baselineskip}
\textbf{from where} — \textbf{hwanon} (adv.). whence; from where.\par
% glossary-block:B00714:end
% glossary-block:B00715:begin
\needspace{3\baselineskip}
\textbf{fruit} — \textbf{wæstm} (n.). fruit; produce. Corrected short æ in wæstm under C01, despite the retained erroneous source hint wǣstm.\par
% glossary-block:B00715:end
% glossary-block:B00716:begin
\needspace{3\baselineskip}
\textbf{fruit-bearing} — \textbf{æppelbǣre} (adj.).\par
% glossary-block:B00716:end
% glossary-block:B00717:begin
\needspace{3\baselineskip}
\textbf{fruitful} — \textbf{wæstmbǣre} (adj.).\par
% glossary-block:B00717:end
% glossary-block:B00718:begin
\needspace{3\baselineskip}
\textbf{fulfil} — \textbf{dōn} (v.). fulfil; carry out. A word or command is fulfilled. \textbf{gefyllan} (v.). Accusative prediction fulfilled; passive subjunctive in the assigned purpose construction.\par
% glossary-block:B00718:end
% glossary-block:B00719:begin
\needspace{3\baselineskip}
\textbf{fulfill} — \textbf{dōn} (v.). fulfil; carry out. A word or command is fulfilled. \textbf{gefyllan} (v.). fulfil. Accusative prediction fulfilled; passive subjunctive in the assigned purpose construction.\par
% glossary-block:B00719:end
% glossary-block:B00720:begin
\needspace{3\baselineskip}
\textbf{furlong} — \textbf{furlang} (n.).\par
% glossary-block:B00720:end
% glossary-block:B00721:begin
\needspace{3\baselineskip}
\textbf{further} — \textbf{furðor} (adv.). further; any more. Adverb in the supplied no-spirit-more expression. \textbf{gīet} (adv.). still; further.\par
% glossary-block:B00721:end
% glossary-block:B00722:begin
\needspace{3\baselineskip}
\textbf{fury} — \textbf{rēðnes} (n.). fury; fierceness.\par
% glossary-block:B00722:end
% book-navigation: standalone bookmark replaced by volume hierarchy
% glossary-block:B00723:begin
\needspace{5\baselineskip}
\subsection*{G}\BookHeading{nav-glossary-english-0008}{subsubsection}{G}{0}
% glossary-block:B00723:end
% glossary-block:B00724:begin
\textbf{gain} — \textbf{gestrēon} (n.). gain; profit.\par
% glossary-block:B00724:end
% glossary-block:B00725:begin
\needspace{3\baselineskip}
\textbf{game} — \textbf{plega} (n.). play; game.\par
% glossary-block:B00725:end
% glossary-block:B00726:begin
\needspace{3\baselineskip}
\textbf{garment} — \textbf{gegierela} (n.). garment; clothing. \textbf{wǣd} (n.).  \textbf{wǣd} (n.), \textbf{gegierela} (n.). Wǣd denotes a garment, rēaf raiment/clothing, and gegierela a garment or clothing. Preserve any noun expressly supplied by the exercise.\par
% glossary-block:B00726:end
% glossary-block:B00727:begin
\needspace{3\baselineskip}
\textbf{garments} — \textbf{rēaf} (n.). clothing; garments.\par
% glossary-block:B00727:end
% glossary-block:B00728:begin
\needspace{3\baselineskip}
\textbf{gate} — \textbf{geat} (n.).\par
% glossary-block:B00728:end
% glossary-block:B00729:begin
\needspace{3\baselineskip}
\textbf{gather} — \textbf{gegaderian} (v.). Accusative people or things gathered.\par
% glossary-block:B00729:end
% glossary-block:B00730:begin
\needspace{3\baselineskip}
\textbf{gathering} — \textbf{gegaderung} (n.). gathering; collection; congregation. Congregations of waters in the source.\par
% glossary-block:B00730:end
% glossary-block:B00731:begin
\needspace{3\baselineskip}
\textbf{gaze} — \textbf{gesihð} (n.). sight; gaze. Supplied accusative singular.\par
% glossary-block:B00731:end
% glossary-block:B00732:begin
\needspace{3\baselineskip}
\textbf{gem} — \textbf{gimm} (n.). gem; precious stone.\par
% glossary-block:B00732:end
% glossary-block:B00733:begin
\needspace{3\baselineskip}
\textbf{generation} — \textbf{cnēoris} (n.).\par
% glossary-block:B00733:end
% glossary-block:B00734:begin
\needspace{3\baselineskip}
\textbf{gift} — \textbf{giefu} (n.). Supplied plural gifts. Divine grace in Cook's bracketed gift route; not a license to merge all secular present/gift senses. \textbf{lāc} (n.). gift; present.\par
% glossary-block:B00734:end
% glossary-block:B00735:begin
\needspace{3\baselineskip}
\textbf{gird} — \textbf{begyrdan} (v.). gird; clothe. Supplied participial clothing expression.\par
% glossary-block:B00735:end
% glossary-block:B00736:begin
\needspace{3\baselineskip}
\textbf{girl} — \textbf{mǣden} (n.). girl; maiden.\par
% glossary-block:B00736:end
% glossary-block:B00737:begin
\needspace{3\baselineskip}
\textbf{give} — \textbf{forgiefan} (v.). give; grant. Supplied granted sense at Cook 143.18; not forgive here. give; grant. Accusative thing given and dative recipient; the thing may be understood from context. \textbf{gesęllan} (v.).  \textbf{giefan} (v.).  \textbf{sęllan} (v.). Sęllan means give in the giving sense; an expressed recipient is dative. Its sell sense is separate.\par
% glossary-block:B00737:end
% glossary-block:B00738:begin
\needspace{3\baselineskip}
\textbf{give back} — \textbf{āgiefan} (v.). give back; restore.\par
% glossary-block:B00738:end
% glossary-block:B00739:begin
\needspace{3\baselineskip}
\textbf{give light to} — \textbf{ālīehtan} (v.). illuminate; give light to. Accusative person or thing illuminated. ālīehtan takes an accusative object, not a dative inferred from English to.\par
% glossary-block:B00739:end
% glossary-block:B00740:begin
\needspace{3\baselineskip}
\textbf{give to drink} — \textbf{scęncan} (v.). pour out; give to drink. Accusative drink, dative recipient. \textbf{sęllan} (v.). Supplied dative person + plain drincan, without tō. Supplied sęllan + dative person + plain drincan; no added tō.\par
% glossary-block:B00740:end
% glossary-block:B00741:begin
\needspace{3\baselineskip}
\textbf{glad} — \textbf{blīðe} (adj.). merry; glad. \textbf{glæd} (adj.).\par
% glossary-block:B00741:end
% glossary-block:B00742:begin
\needspace{3\baselineskip}
\textbf{gladly} — \textbf{blīðe} (adv.). joyfully; gladly.\par
% glossary-block:B00742:end
% glossary-block:B00743:begin
\needspace{3\baselineskip}
\textbf{glorify} — \textbf{mǣrsian} (v.). make great; magnify; glorify. \textbf{wuldrian} (v.). Accusative person; an accusative reflexive expresses glorifying oneself.\par
% glossary-block:B00743:end
% glossary-block:B00744:begin
\needspace{3\baselineskip}
\textbf{glory} — \textbf{mægenðrymm} (n.). majesty; glory. \textbf{mægenðrymnes} (n.). glory; majesty. \textbf{wuldor} (n.).\par
% glossary-block:B00744:end
% glossary-block:B00745:begin
\needspace{3\baselineskip}
\textbf{go} — \textbf{faran} (v.). go; travel. \textbf{gān} (v.). go; walk. Supplied imperative gang. \textbf{gangan} (v.). go; enter. Supplied circumstantial participial phrase.\par
% glossary-block:B00745:end
% glossary-block:B00746:begin
\needspace{3\baselineskip}
\textbf{go aboard} — \textbf{āstīgan} (v.). go aboard; board. Accusative ship boarded in the assigned frame. āstīgan takes the ship as direct accusative in these exercises.\par
% glossary-block:B00746:end
% glossary-block:B00747:begin
\needspace{3\baselineskip}
\textbf{goat} — \textbf{gāt} (n.).\par
% glossary-block:B00747:end
% glossary-block:B00748:begin
\needspace{3\baselineskip}
\textbf{go away} — \textbf{āgān} (v.). Supplied gone-away form; distinct from āgan possess. \textbf{gewītan} (v.). depart; go away. Intransitive; fram + dative departure point is optional. The assigned frame requires no reflexive.\par
% glossary-block:B00748:end
% glossary-block:B00749:begin
\needspace{3\baselineskip}
\textbf{god} — \textbf{God} (n.). god; deity. Common deity sense, distinct from the proper divine name. A god or deity; this common-noun sense is distinct from the proper name God.\par
% glossary-block:B00749:end
% glossary-block:B00750:begin
\needspace{3\baselineskip}
\textbf{God} — \textbf{God} (n.). Proper divine name; no productive plural sense is assigned. The proper divine name God; no productive plural is assigned.\par
% glossary-block:B00750:end
% glossary-block:B00751:begin
\needspace{3\baselineskip}
\textbf{go fishing} — Gān on fiscnoð means go fishing; fiscnoð is the noun fishing and on is part of this supplied phrase. Gān on fiscnoð is the supplied whole expression go fishing. Fiscnoð is the noun fishing; on is part of this phrase, not a verb meaning go.\par
% glossary-block:B00751:end
% glossary-block:B00752:begin
\needspace{3\baselineskip}
\textbf{go free} — Frēo is the adjective free. In go free it describes the subject with appropriate agreement; it is not itself a verb meaning go.\par
% glossary-block:B00752:end
% glossary-block:B00753:begin
\needspace{3\baselineskip}
\textbf{go in} — \textbf{ingān} (v.). enter; go in. enter; go in. Supplied past source form. \textbf{ingangan} (v.). enter; go in. Supplied present-participle phrase.\par
% glossary-block:B00753:end
% glossary-block:B00754:begin
\needspace{3\baselineskip}
\textbf{gold} — \textbf{gold} (n.).\par
% glossary-block:B00754:end
% glossary-block:B00755:begin
\needspace{3\baselineskip}
\textbf{good} — \textbf{gōd} (adj.). Ordinary gōd good. The person-directed kind/benevolent use is a separate sense. \textbf{gōd} (n.). good; benefit; kindness. Nominal good in the repayment contrast. \textbf{ðearf} (n.). benefit; good; need. Source contextual nouns; retain the local expression.\par
% glossary-block:B00755:end
% glossary-block:B00756:begin
\needspace{3\baselineskip}
\textbf{goods} — \textbf{ǣht} (n.). possessions; goods. \textbf{gōd} (n.). goods; property. Plural lexical noun.\par
% glossary-block:B00756:end
% glossary-block:B00757:begin
\needspace{3\baselineskip}
\textbf{good service} — \textbf{earnung} (n.). merit; good service. Source plural merits/services.\par
% glossary-block:B00757:end
% glossary-block:B00758:begin
\needspace{3\baselineskip}
\textbf{good things} — \textbf{gōd} (n.). goods; property. Plural lexical noun.\par
% glossary-block:B00758:end
% glossary-block:B00759:begin
\needspace{3\baselineskip}
\textbf{go out} — \textbf{ūtgān} (v.). Supplied past source form.\par
% glossary-block:B00759:end
% glossary-block:B00760:begin
\needspace{3\baselineskip}
\textbf{Gospel} — \textbf{godspell} (n.).\par
% glossary-block:B00760:end
% glossary-block:B00761:begin
\needspace{3\baselineskip}
\textbf{go to pieces} — \textbf{tōberstan} (v.). break up; go to pieces.\par
% glossary-block:B00761:end
% glossary-block:B00762:begin
\needspace{3\baselineskip}
\textbf{go up} — \textbf{ūpgān} (v.). go up; come up. From the water in the supplied expression.\par
% glossary-block:B00762:end
% glossary-block:B00763:begin
\needspace{3\baselineskip}
\textbf{grace} — \textbf{giefu} (n.). grace; divine favour. Divine grace in Cook's bracketed gift route; not a license to merge all secular present/gift senses.\par
% glossary-block:B00763:end
% glossary-block:B00764:begin
\needspace{3\baselineskip}
\textbf{gracious} — \textbf{glæd} (adj.). gracious; favourably disposed. With wið + person in the supplied gracious-toward frame.\par
% glossary-block:B00764:end
% glossary-block:B00765:begin
\needspace{3\baselineskip}
\textbf{grant} — \textbf{forgiefan} (v.). give; grant. Supplied granted sense at Cook 143.18; not forgive here. give; grant. Accusative thing given and dative recipient; the thing may be understood from context. \textbf{unnan} (v.). grant; allow. Dative person, genitive thing granted. Dative recipient and genitive thing granted; preserve the supported subjunctive command unne.\par
% glossary-block:B00765:end
% glossary-block:B00766:begin
\needspace{3\baselineskip}
\textbf{grass} — \textbf{gærs} (n.).\par
% glossary-block:B00766:end
% glossary-block:B00767:begin
\needspace{3\baselineskip}
\textbf{great} — \textbf{mǣre} (adj.). great; renowned. Distinct from simple physical size where the task requires the renowned sense. \textbf{micel} (adj.). great; large. great; loud. Loud is the contextual force with a voice.\par
% glossary-block:B00767:end
% glossary-block:B00768:begin
\needspace{3\baselineskip}
\textbf{greater} — \textbf{māra} (adj.). greater; larger. Comparative adjective. greater; more important. Value/importance in the life/raiment comparison. Mā expresses more in amount or number in its separate uses; māra expresses a greater amount. Comparative swīðor expresses more intensely or preferential rather. Use an agreeing comparative adjective for greater quality/value; these uses are not interchangeable.\par
% glossary-block:B00768:end
% glossary-block:B00769:begin
\needspace{3\baselineskip}
\textbf{greater amount} — \textbf{māra} (adj.). more; a greater amount. Nominal/attributive quantity, distinct from more in number.\par
% glossary-block:B00769:end
% glossary-block:B00770:begin
\needspace{3\baselineskip}
\textbf{greatest} — \textbf{mǣst} (adj.). greatest; largest. Superlative adjective.\par
% glossary-block:B00770:end
% glossary-block:B00771:begin
\needspace{3\baselineskip}
\textbf{greatly} — \textbf{ðearle} (adv.). greatly; exceedingly. \textbf{swīðe} (adv.). Degree adverb swīðe; retain the separate great/large adjective micel.\par
% glossary-block:B00771:end
% glossary-block:B00772:begin
\needspace{3\baselineskip}
\textbf{greet} — \textbf{grētan} (v.). greet; take leave of. Accusative person/place; greeting and farewell depend on context. Greeting with accusative person; the same verb also has farewell/leave-taking uses.\par
% glossary-block:B00772:end
% glossary-block:B00773:begin
\needspace{3\baselineskip}
\textbf{grief} — \textbf{sārnes} (n.).\par
% glossary-block:B00773:end
% glossary-block:B00774:begin
\needspace{3\baselineskip}
\textbf{grieve} — \textbf{ofðyncean} (v.). grieve; displease. Impersonal third singular with dative experiencer and a content clause. Impersonal third-singular predicate, dative experiencer, and ðæt clause for the cause.\par
% glossary-block:B00774:end
% glossary-block:B00775:begin
\needspace{3\baselineskip}
\textbf{grievous} — \textbf{hęfig} (adj.).\par
% glossary-block:B00775:end
% glossary-block:B00776:begin
\needspace{3\baselineskip}
\textbf{group} — \textbf{flocc} (n.). group; company.\par
% glossary-block:B00776:end
% glossary-block:B00777:begin
\needspace{3\baselineskip}
\textbf{grove} — \textbf{wudubearu} (n.). forest; grove.\par
% glossary-block:B00777:end
% glossary-block:B00778:begin
\needspace{3\baselineskip}
\textbf{grow} — \textbf{weaxan} (v.). Intransitive increase in these commands/wishes; not the transitive īecan increase something.\par
% glossary-block:B00778:end
% glossary-block:B00779:begin
\needspace{3\baselineskip}
\textbf{guest} — \textbf{cuma} (n.). stranger; guest.\par
% glossary-block:B00779:end
% glossary-block:B00780:begin
\needspace{3\baselineskip}
\textbf{guest-house} — \textbf{giesthūs} (n.). inn; guest-house.\par
% glossary-block:B00780:end
% glossary-block:B00781:begin
\needspace{3\baselineskip}
\textbf{guide} — \textbf{gewissian} (v.). guide; direct.\par
% glossary-block:B00781:end
% glossary-block:B00782:begin
\needspace{3\baselineskip}
\textbf{guilty} — \textbf{scyldig} (adj.). guilty; culpable.\par
% glossary-block:B00782:end
% glossary-block:B00783:begin
\needspace{3\baselineskip}
\textbf{gymnasium} — \textbf{bæðstęde} (n.). bathing place; gymnasium. Gymnasium is the contextual interpretation in the supplied reading.\par
% glossary-block:B00783:end
% book-navigation: standalone bookmark replaced by volume hierarchy
% glossary-block:B00784:begin
\needspace{5\baselineskip}
\subsection*{H}\BookHeading{nav-glossary-english-0009}{subsubsection}{H}{0}
% glossary-block:B00784:end
% glossary-block:B00785:begin
\textbf{habit} — \textbf{gewuna} (n.). custom; habit.\par
% glossary-block:B00785:end
% glossary-block:B00786:begin
\needspace{3\baselineskip}
\textbf{had} — Auxiliary habban combines with a transitive past participle in the taught completed-action/result construction. English have/has/had does not require habban in every perfect; a simple Old English tense may express the intended time relation. \textbf{habban} (v.). Possessive habban takes an accusative possession. Auxiliary habban combines with a transitive past participle in the taught perfect construction; it is not an automatic translation of every English have.\par
% glossary-block:B00786:end
% glossary-block:B00787:begin
\needspace{3\baselineskip}
\textbf{had …} — Preserve the action's time relation. Simple preterite, sometimes with ǣr, often renders English perfect/pluperfect; use an analytic construction only where taught/requested.\par
% glossary-block:B00787:end
% glossary-block:B00788:begin
\needspace{3\baselineskip}
\textbf{had come} — Wesan + the assigned intransitive past participle supplies the taught compound construction. With cumen come it can render have/had come; English have is not an extra word to add to that construction.\par
% glossary-block:B00788:end
% glossary-block:B00789:begin
\needspace{3\baselineskip}
\textbf{had rather} — \textbf{lēofra} (adj.). preferable; dearer. Dative experiencer + is lēofre + infinitive in the taught preference frame. The preference construction is dative experiencer + is lēofre + plain infinitive, is preferable to someone. It renders had/would rather without requiring past time. Wician in the supplied example is the action preferred, sojourn/live with, not a word meaning prefer.\par
% glossary-block:B00789:end
% glossary-block:B00790:begin
\needspace{3\baselineskip}
\textbf{had to} — \textbf{sculan} (v.). sculan + plain infinitive for obligation; English have to is not possessive habban. Supplied have something to offer in Unit 6 is a different construction.\par
% glossary-block:B00790:end
% glossary-block:B00791:begin
\needspace{3\baselineskip}
\textbf{hail} — \textbf{hægel} (n.).\par
% glossary-block:B00791:end
% glossary-block:B00792:begin
\needspace{3\baselineskip}
\textbf{hair} — \textbf{hǣr} (n.). hair; a hair.\par
% glossary-block:B00792:end
% glossary-block:B00793:begin
\needspace{3\baselineskip}
\textbf{hale} — \textbf{hāl} (adj.). hale; well. Supplied greeting. hale; healthy.\par
% glossary-block:B00793:end
% glossary-block:B00794:begin
\needspace{3\baselineskip}
\textbf{half} — \textbf{healf} (adj.).\par
% glossary-block:B00794:end
% glossary-block:B00795:begin
\needspace{3\baselineskip}
\textbf{hall} — \textbf{heall} (n.).\par
% glossary-block:B00795:end
% glossary-block:B00796:begin
\needspace{3\baselineskip}
\textbf{hallow} — \textbf{gehālgian} (v.). hallow; consecrate.\par
% glossary-block:B00796:end
% glossary-block:B00797:begin
\needspace{3\baselineskip}
\textbf{halt} — \textbf{healt} (adj.). lame; halt. Adjective describing impaired walking; not the command halt/stop.\par
% glossary-block:B00797:end
% glossary-block:B00798:begin
\needspace{3\baselineskip}
\textbf{hammer} — \textbf{slęcg} (n.).\par
% glossary-block:B00798:end
% glossary-block:B00799:begin
\needspace{3\baselineskip}
\textbf{hand} — \textbf{hand} (n.). hand; side. Side in right-/left-hand expressions.\par
% glossary-block:B00799:end
% glossary-block:B00800:begin
\needspace{3\baselineskip}
\textbf{handmaiden} — \textbf{ðīnen} (n.).\par
% glossary-block:B00800:end
% glossary-block:B00801:begin
\needspace{3\baselineskip}
\textbf{happen} — \textbf{gelimpan} (v.). happen; befall. Supplied preterite subjunctive in might befall us. happen; come about. Supplied preterite source form. happen; come about.\par
% glossary-block:B00801:end
% glossary-block:B00802:begin
\needspace{3\baselineskip}
\textbf{happy} — \textbf{ēadig} (adj.). happy; blessed.\par
% glossary-block:B00802:end
% glossary-block:B00803:begin
\needspace{3\baselineskip}
\textbf{hard to count} — \textbf{earfoðrīme} (adj.). innumerable; hard to count. Assigned source adjective.\par
% glossary-block:B00803:end
% glossary-block:B00804:begin
\needspace{3\baselineskip}
\textbf{harp} — \textbf{hearpe} (n.).  \textbf{hearpian} (v.). Supplied genitive harper/harping construction, with instrument; do not reconstruct it from the infinitive alone. Preserve the supplied genitive harper/harping expression; headword alone does not reconstruct its syntax.\par
% glossary-block:B00804:end
% glossary-block:B00805:begin
\needspace{3\baselineskip}
\textbf{harper} — \textbf{hearpere} (n.).\par
% glossary-block:B00805:end
% glossary-block:B00806:begin
\needspace{3\baselineskip}
\textbf{harp pick} — \textbf{hearpenægl} (n.). plectrum; harp pick.\par
% glossary-block:B00806:end
% glossary-block:B00807:begin
\needspace{3\baselineskip}
\textbf{harp strings} — \textbf{hearpestręng} (n.).\par
% glossary-block:B00807:end
% glossary-block:B00808:begin
\needspace{3\baselineskip}
\textbf{harvest} — \textbf{hærfest} (n.). Masculine noun; rich in uses welig on + dative fruits.\par
% glossary-block:B00808:end
% glossary-block:B00809:begin
\needspace{3\baselineskip}
\textbf{has} — Auxiliary habban combines with a transitive past participle in the taught completed-action/result construction. English have/has/had does not require habban in every perfect; a simple Old English tense may express the intended time relation. \textbf{habban} (v.). Possessive habban takes an accusative possession. Auxiliary habban combines with a transitive past participle in the taught perfect construction; it is not an automatic translation of every English have.\par
% glossary-block:B00809:end
% glossary-block:B00810:begin
\needspace{3\baselineskip}
\textbf{hatchet} — \textbf{adesa} (n.).\par
% glossary-block:B00810:end
% glossary-block:B00811:begin
\needspace{3\baselineskip}
\textbf{hate} — \textbf{hatian} (v.).\par
% glossary-block:B00811:end
% glossary-block:B00812:begin
\needspace{3\baselineskip}
\textbf{have} — \textbf{habban} (v.). have; possess. Accusative possession; distinguish the perfect auxiliary. Possessive habban takes an accusative possession. Auxiliary habban combines with a transitive past participle in the taught perfect construction; it is not an automatic translation of every English have. Auxiliary habban combines with a transitive past participle in the taught completed-action/result construction. English have/has/had does not require habban in every perfect; a simple Old English tense may express the intended time relation. \textbf{āgan} (v.). have; possess. Accusative possession; positive present and preterite indicative in the taught scope.\par
% glossary-block:B00812:end
% glossary-block:B00813:begin
\needspace{3\baselineskip}
\textbf{have as wife} — \textbf{tō} (prep.). habban + accusative woman + tō wīfe in the supplied role frame.\par
% glossary-block:B00813:end
% glossary-block:B00814:begin
\needspace{3\baselineskip}
\textbf{have come} — Wesan + the assigned intransitive past participle supplies the taught compound construction. With cumen come it can render have/had come; English have is not an extra word to add to that construction.\par
% glossary-block:B00814:end
% glossary-block:B00815:begin
\needspace{3\baselineskip}
\textbf{have compassion on} — \textbf{gemiltsian} (v.). pity; have compassion on. Dative person. gemiltsian + dative person. A wish uses the appropriate subjunctive, not a separate have verb.\par
% glossary-block:B00815:end
% glossary-block:B00816:begin
\needspace{3\baselineskip}
\textbf{have … (completed action)} — Preserve the action's time relation. Simple preterite, sometimes with ǣr, often renders English perfect/pluperfect; use an analytic construction only where taught/requested.\par
% glossary-block:B00816:end
% glossary-block:B00817:begin
\needspace{3\baselineskip}
\textbf{have dominion} — \textbf{rīcsian} (v.). rule; have dominion. Used absolutely in the taught frame; no new object government. rīcsian can be used absolutely; English have/bear need no independent translation.\par
% glossary-block:B00817:end
% glossary-block:B00818:begin
\needspace{3\baselineskip}
\textbf{have dominion over} — \textbf{ofer} (prep.). over. Accusative dependent in the taught authority frame.\par
% glossary-block:B00818:end
% glossary-block:B00819:begin
\needspace{3\baselineskip}
\textbf{have mercy on} — \textbf{gemiltsian} (v.). pity; have compassion on. Dative person. gemiltsian + dative person. A wish uses the appropriate subjunctive, not a separate have verb.\par
% glossary-block:B00819:end
% glossary-block:B00820:begin
\needspace{3\baselineskip}
\textbf{have not} — \textbf{nabban} (v.).\par
% glossary-block:B00820:end
% glossary-block:B00821:begin
\needspace{3\baselineskip}
\textbf{have on} — \textbf{on} (adv.). on. Adverb in the cloak-wearing relative clause.\par
% glossary-block:B00821:end
% glossary-block:B00822:begin
\needspace{3\baselineskip}
\textbf{have speech with} — \textbf{wið} (prep.). with. Supplied conversational counterpart; do not infer hostility.\par
% glossary-block:B00822:end
% glossary-block:B00823:begin
\needspace{3\baselineskip}
\textbf{have to} — \textbf{sculan} (v.). sculan + plain infinitive for obligation; English have to is not possessive habban. Supplied have something to offer in Unit 6 is a different construction.\par
% glossary-block:B00823:end
% glossary-block:B00824:begin
\needspace{3\baselineskip}
\textbf{he} — Hē is masculine singular nominative. Third-person pronouns follow the grammatical gender of their referent. Grammatical gender and antecedent determine reference; it may be represented by masculine/feminine forms for a gendered inanimate noun.\par
% glossary-block:B00824:end
% glossary-block:B00825:begin
\needspace{3\baselineskip}
\textbf{head} — \textbf{hēafod} (n.).\par
% glossary-block:B00825:end
% glossary-block:B00826:begin
\needspace{3\baselineskip}
\textbf{healthy} — \textbf{hāl} (adj.). hale; healthy.\par
% glossary-block:B00826:end
% glossary-block:B00827:begin
\needspace{3\baselineskip}
\textbf{hear} — \textbf{gehīeran} (v.). Supplied prefixed third-singular preterite. hear; listen. \textbf{hīeran} (v.). hear an action. Accusative perceived actor plus plain infinitive in the taught perception construction. \textbf{hlystan} (v.). listen to; heed. Genitive person or thing listened to; hear is the contextual prompt wording. hlystan governs genitive of the speech heard; do not copy hīeran's ordinary accusative frame.\par
% glossary-block:B00827:end
% glossary-block:B00828:begin
\needspace{3\baselineskip}
\textbf{hear someone do something} — \textbf{hīeran} (v.). hear an action. Accusative perceived actor plus plain infinitive in the taught perception construction.\par
% glossary-block:B00828:end
% glossary-block:B00829:begin
\needspace{3\baselineskip}
\textbf{heart} — \textbf{heorte} (n.).  \textbf{mōd} (n.). heart; mind.\par
% glossary-block:B00829:end
% glossary-block:B00830:begin
\needspace{3\baselineskip}
\textbf{heat} — \textbf{hǣtu} (n.).\par
% glossary-block:B00830:end
% glossary-block:B00831:begin
\needspace{3\baselineskip}
\textbf{heathen people} — \textbf{hǣðen} (adj.).\par
% glossary-block:B00831:end
% glossary-block:B00832:begin
\needspace{3\baselineskip}
\textbf{heaven} — \textbf{heofon} (n.). heaven; sky. \textbf{heofone} (n.). heaven; sky. Weak feminine lexical alternative in the source.\par
% glossary-block:B00832:end
% glossary-block:B00833:begin
\needspace{3\baselineskip}
\textbf{heavenly} — \textbf{heofonlic} (adj.).  \textbf{heofonlice} (adv.). from heaven; heavenly. Source expression at Cook 143.17.\par
% glossary-block:B00833:end
% glossary-block:B00834:begin
\needspace{3\baselineskip}
\textbf{heavenly fire} — \textbf{heofonfȳr} (n.). fire from heaven; heavenly fire.\par
% glossary-block:B00834:end
% glossary-block:B00835:begin
\needspace{3\baselineskip}
\textbf{heavenly powers} — \textbf{mægen} (n.). power; strength. Heavenly powers in 37.2 is the contextual plural use.\par
% glossary-block:B00835:end
% glossary-block:B00836:begin
\needspace{3\baselineskip}
\textbf{heed} — \textbf{hlystan} (v.). listen to; heed. Genitive person or thing listened to; hear is the contextual prompt wording. hlystan governs genitive of the speech heard; do not copy hīeran's ordinary accusative frame.\par
% glossary-block:B00836:end
% glossary-block:B00837:begin
\needspace{3\baselineskip}
\textbf{height} — \textbf{hēanes} (n.). height; highest point.\par
% glossary-block:B00837:end
% glossary-block:B00838:begin
\needspace{3\baselineskip}
\textbf{heights} — \textbf{hēanes} (n.). heights; highest places.\par
% glossary-block:B00838:end
% glossary-block:B00839:begin
\needspace{3\baselineskip}
\textbf{hell} — \textbf{hęll} (n.).\par
% glossary-block:B00839:end
% glossary-block:B00840:begin
\needspace{3\baselineskip}
\textbf{help} — \textbf{fultum} (n.). help; aid. help; support. \textbf{gefultumian} (v.). help; assist. \textbf{helpan} (v.). Dative person helped.\par
% glossary-block:B00840:end
% glossary-block:B00841:begin
\needspace{3\baselineskip}
\textbf{hence} — \textbf{heonan} (adv.). hence; from here.\par
% glossary-block:B00841:end
% glossary-block:B00842:begin
\needspace{3\baselineskip}
\textbf{hence forth} — \textbf{heonanforð} (adv.). heonanforð as a whole adverb.\par
% glossary-block:B00842:end
% glossary-block:B00843:begin
\needspace{3\baselineskip}
\textbf{henceforth} — \textbf{heonanforð} (adv.). heonanforð as a whole adverb.\par
% glossary-block:B00843:end
% glossary-block:B00844:begin
\needspace{3\baselineskip}
\textbf{her} — Hīe is feminine singular accusative; hiere is dative or genitive. Possessive hiere remains invariant before the possessed noun. \textbf{hē} (personal pron.). his; her; its; their. Invariant genitive possessors, unlike agreeing possessive adjectives. Invariant personal genitive possessor; distinguish the same surface used as a governed personal case. Grammatical gender and antecedent determine reference; it may be represented by masculine/feminine forms for a gendered inanimate noun.\par
% glossary-block:B00844:end
% glossary-block:B00845:begin
\needspace{3\baselineskip}
\textbf{herb} — \textbf{wyrt} (n.). herb; plant. Cook specifically routes these spices to wyrt; wyrtgemang is the separate earlier spice word.\par
% glossary-block:B00845:end
% glossary-block:B00846:begin
\needspace{3\baselineskip}
\textbf{here} — \textbf{hēr} (adv.).  \textbf{hider} (adv.). hither; here. Motion toward the speaker, distinct from location hēr.\par
% glossary-block:B00846:end
% glossary-block:B00847:begin
\needspace{3\baselineskip}
\textbf{herself} — An ordinary reflexive uses the personal pronoun in the case governed by its verb or preposition. Agreeing self adds the taught emphasis; English -self does not automatically require a separate self. Ordinary reflexive personal pronouns take the governed case; agreeing self adds the taught emphasis. English -self does not automatically require self.\par
% glossary-block:B00847:end
% glossary-block:B00848:begin
\needspace{3\baselineskip}
\textbf{hew} — \textbf{hēawan} (v.).\par
% glossary-block:B00848:end
% glossary-block:B00849:begin
\needspace{3\baselineskip}
\textbf{hide} — \textbf{diernan} (v.). conceal; hide. \textbf{hȳdan} (v.). hide; conceal. Accusative thing hidden; wesan forms the assigned passive.\par
% glossary-block:B00849:end
% glossary-block:B00850:begin
\needspace{3\baselineskip}
\textbf{Hierusalem} — \textbf{Hierusalem} (proper name).\par
% glossary-block:B00850:end
% glossary-block:B00851:begin
\needspace{3\baselineskip}
\textbf{high} — \textbf{hēah} (adj.). high in rank; exalted. Rank/status sense in the supplied source. \textbf{hēah} (adv.). Adverbial height, distinct from an agreeing adjective.\par
% glossary-block:B00851:end
% glossary-block:B00852:begin
\needspace{3\baselineskip}
\textbf{high-born} — \textbf{æðelboren} (adj.).\par
% glossary-block:B00852:end
% glossary-block:B00853:begin
\needspace{3\baselineskip}
\textbf{high cliff} — \textbf{hēaclif} (n.).\par
% glossary-block:B00853:end
% glossary-block:B00854:begin
\needspace{3\baselineskip}
\textbf{Highest} — \textbf{hēah} (adj.). the Highest. Divine title.\par
% glossary-block:B00854:end
% glossary-block:B00855:begin
\needspace{3\baselineskip}
\textbf{highest places} — \textbf{hēanes} (n.). heights; highest places.\par
% glossary-block:B00855:end
% glossary-block:B00856:begin
\needspace{3\baselineskip}
\textbf{highest point} — \textbf{hēanes} (n.). height; highest point.\par
% glossary-block:B00856:end
% glossary-block:B00857:begin
\needspace{3\baselineskip}
\textbf{high office} — \textbf{ār} (n.). honour; high position. Contextual rank or office, not a new general meaning of every ār form.\par
% glossary-block:B00857:end
% glossary-block:B00858:begin
\needspace{3\baselineskip}
\textbf{high position} — \textbf{ār} (n.). honour; high position. Contextual rank or office, not a new general meaning of every ār form.\par
% glossary-block:B00858:end
% glossary-block:B00859:begin
\needspace{3\baselineskip}
\textbf{hight} — \textbf{hātan} (v.). be called; be named. Special hātte form in the assigned naming expression. The special form hātte means is/was called or named, according to context; do not substitute an ordinary active hātan form.\par
% glossary-block:B00859:end
% glossary-block:B00860:begin
\needspace{3\baselineskip}
\textbf{hill} — \textbf{dūn} (n.). mountain; hill. \textbf{hyll} (n.).\par
% glossary-block:B00860:end
% glossary-block:B00861:begin
\needspace{3\baselineskip}
\textbf{him} — Hine is masculine singular accusative; him is masculine/neuter singular dative. Choose the governed case. Grammatical gender and antecedent determine reference; it may be represented by masculine/feminine forms for a gendered inanimate noun.\par
% glossary-block:B00861:end
% glossary-block:B00862:begin
\needspace{3\baselineskip}
\textbf{himself} — An ordinary reflexive uses the personal pronoun in the case governed by its verb or preposition. Agreeing self adds the taught emphasis; English -self does not automatically require a separate self. Ordinary reflexive personal pronouns take the governed case; agreeing self adds the taught emphasis. English -self does not automatically require self.\par
% glossary-block:B00862:end
% glossary-block:B00863:begin
\needspace{3\baselineskip}
\textbf{hire} — \textbf{hȳrian} (v.).  \textbf{mēd} (n.). hire; wages paid for work.\par
% glossary-block:B00863:end
% glossary-block:B00864:begin
\needspace{3\baselineskip}
\textbf{his} — \textbf{hē} (personal pron.). his; her; its; their. Invariant genitive possessors, unlike agreeing possessive adjectives. Invariant personal genitive possessor; distinguish the same surface used as a governed personal case.\par
% glossary-block:B00864:end
% glossary-block:B00865:begin
\needspace{3\baselineskip}
\textbf{hither} — \textbf{hider} (adv.). hither; here. Motion toward the speaker, distinct from location hēr.\par
% glossary-block:B00865:end
% glossary-block:B00866:begin
\needspace{3\baselineskip}
\textbf{hoard} — \textbf{hord} (n.). hoard; treasure.\par
% glossary-block:B00866:end
% glossary-block:B00867:begin
\needspace{3\baselineskip}
\textbf{hoard of gold} — \textbf{goldhord} (n.). treasure; hoard of gold.\par
% glossary-block:B00867:end
% glossary-block:B00868:begin
\needspace{3\baselineskip}
\textbf{hold} — \textbf{healdan} (v.).\par
% glossary-block:B00868:end
% glossary-block:B00869:begin
\needspace{3\baselineskip}
\textbf{hold chief place} — Healdan ealdorscipe means hold chief place or primacy: healdan takes accusative ealdorscipe. Retain the exercise’s omission of the article.\par
% glossary-block:B00869:end
% glossary-block:B00870:begin
\needspace{3\baselineskip}
\textbf{hole} — \textbf{hol} (n.).  \textbf{ðyrel} (n.).\par
% glossary-block:B00870:end
% glossary-block:B00871:begin
\needspace{3\baselineskip}
\textbf{holy} — \textbf{hālig} (adj.).  \textbf{Sanctus} (foreign adjective). Saint; holy. Latin word in the source title.\par
% glossary-block:B00871:end
% glossary-block:B00872:begin
\needspace{3\baselineskip}
\textbf{holy one} — \textbf{hālga} (n.). saint; holy one.\par
% glossary-block:B00872:end
% glossary-block:B00873:begin
\needspace{3\baselineskip}
\textbf{holy ones} — \textbf{hālig} (adj.). holy ones; saints. Substantival adjective in the source.\par
% glossary-block:B00873:end
% glossary-block:B00874:begin
\needspace{3\baselineskip}
\textbf{home} — \textbf{hām} (adv.). home; homeward. Directional adverb; distinct from the noun home. \textbf{hām} (n.). home; dwelling.\par
% glossary-block:B00874:end
% glossary-block:B00875:begin
\needspace{3\baselineskip}
\textbf{homeland} — \textbf{ēðel} (n.). native land; homeland. ēðel in these home-country references; land remains a separate ordinary country word.\par
% glossary-block:B00875:end
% glossary-block:B00876:begin
\needspace{3\baselineskip}
\textbf{homeward} — \textbf{hām} (adv.). home; homeward. Directional adverb; distinct from the noun home. \textbf{hāmweard} (adverb or particle). Supplied hāmweard with wesan; distinguish noun hām and directional adverb hām.\par
% glossary-block:B00876:end
% glossary-block:B00877:begin
\needspace{3\baselineskip}
\textbf{homewards} — \textbf{hām} (adv.). home; homeward. Directional adverb; distinct from the noun home.\par
% glossary-block:B00877:end
% glossary-block:B00878:begin
\needspace{3\baselineskip}
\textbf{honey} — \textbf{hunig} (n.).\par
% glossary-block:B00878:end
% glossary-block:B00879:begin
\needspace{3\baselineskip}
\textbf{honor} — \textbf{ār} (n.). honour; high position. Contextual rank or office, not a new general meaning of every ār form.\par
% glossary-block:B00879:end
% glossary-block:B00880:begin
\needspace{3\baselineskip}
\textbf{honour} — \textbf{ār} (n.). honour; high position. Contextual rank or office, not a new general meaning of every ār form. \textbf{weorðscipe} (n.). honour; dignity.\par
% glossary-block:B00880:end
% glossary-block:B00881:begin
\needspace{3\baselineskip}
\textbf{hook} — \textbf{angel} (n.). Compound English spelling and optional fish- lead to the same fishing noun; distinguish ęngel angel.\par
% glossary-block:B00881:end
% glossary-block:B00882:begin
\needspace{3\baselineskip}
\textbf{hope} — \textbf{tōhopa} (n.).\par
% glossary-block:B00882:end
% glossary-block:B00883:begin
\needspace{3\baselineskip}
\textbf{hopelessness} — \textbf{ōrmōdnes} (n.). despair; hopelessness.\par
% glossary-block:B00883:end
% glossary-block:B00884:begin
\needspace{3\baselineskip}
\textbf{horn} — \textbf{horn} (n.). horn; strength; dignity. Strength/dignity is the figurative force in 37.14, not a replacement for the literal source horn.\par
% glossary-block:B00884:end
% glossary-block:B00885:begin
\needspace{3\baselineskip}
\textbf{horse} — \textbf{hors} (n.).\par
% glossary-block:B00885:end
% glossary-block:B00886:begin
\needspace{3\baselineskip}
\textbf{host} — \textbf{weorod} (n.). host; army.\par
% glossary-block:B00886:end
% glossary-block:B00887:begin
\needspace{3\baselineskip}
\textbf{hot} — \textbf{hāt} (adj.).\par
% glossary-block:B00887:end
% glossary-block:B00888:begin
\needspace{3\baselineskip}
\textbf{hour} — \textbf{tīd} (n.). Time measure in the supplied intervals.\par
% glossary-block:B00888:end
% glossary-block:B00889:begin
\needspace{3\baselineskip}
\textbf{house} — \textbf{hūs} (n.).\par
% glossary-block:B00889:end
% glossary-block:B00890:begin
\needspace{3\baselineskip}
\textbf{how} — \textbf{hū} (adv.). how?.\par
% glossary-block:B00890:end
% glossary-block:B00891:begin
\needspace{3\baselineskip}
\textbf{human being} — \textbf{mann} (n.; impersonal pronoun). man; human being.\par
% glossary-block:B00891:end
% glossary-block:B00892:begin
\needspace{3\baselineskip}
\textbf{humble} — \textbf{ēaðmōd} (adj.). lowly; humble. \textbf{hīenan} (v.). oppress; humble.\par
% glossary-block:B00892:end
% glossary-block:B00893:begin
\needspace{3\baselineskip}
\textbf{humble people} — \textbf{ēaðmōd} (adj.). humble people; the lowly. Substantival adjective.\par
% glossary-block:B00893:end
% glossary-block:B00894:begin
\needspace{3\baselineskip}
\textbf{humility} — \textbf{ēaðmōdnes} (n.).\par
% glossary-block:B00894:end
% glossary-block:B00895:begin
\needspace{3\baselineskip}
\textbf{hundred} — \textbf{hund} (cardinal). Only the complete supplied hundred-sheep expression at its first appearance; later numerical frames remain locally specified. \textbf{hundtēontig} (numeral). Within the complete supplied 144,000 expression; do not infer an independent high-number system. High quantities retain the complete local numerical expressions; a glossary headword does not replace the supplied genitive/counting construction.\par
% glossary-block:B00895:end
% glossary-block:B00896:begin
\needspace{3\baselineskip}
\textbf{hundredfold} — \textbf{hundtēontig} (numeral). a hundredfold. A hundred in return, in the supplied recompense expression.\par
% glossary-block:B00896:end
% glossary-block:B00897:begin
\needspace{3\baselineskip}
\textbf{hunger} — \textbf{hunger} (n.). famine; hunger. \textbf{hyngran} (v.). be hungry. Impersonal construction with accusative experiencer.\par
% glossary-block:B00897:end
% glossary-block:B00898:begin
\needspace{3\baselineskip}
\textbf{hungering} — \textbf{hyngran} (v.). hungry; hungering. Present participle used substantivally in the assigned frame.\par
% glossary-block:B00898:end
% glossary-block:B00899:begin
\needspace{3\baselineskip}
\textbf{hungry} — \textbf{hungrig} (adj.).  \textbf{hyngran} (v.). hungry; hungering. Present participle used substantivally in the assigned frame.\par
% glossary-block:B00899:end
% glossary-block:B00900:begin
\needspace{3\baselineskip}
\textbf{hungry people} — \textbf{hungrig} (adj.). Substantival adjective.\par
% glossary-block:B00900:end
% glossary-block:B00901:begin
\needspace{3\baselineskip}
\textbf{hurt} — \textbf{scęððan} (v.). hurt; injure. Dative person harmed.\par
% glossary-block:B00901:end
% glossary-block:B00902:begin
\needspace{3\baselineskip}
\textbf{husband} — \textbf{wer} (n.).\par
% glossary-block:B00902:end
% book-navigation: standalone bookmark replaced by volume hierarchy
% glossary-block:B00903:begin
\needspace{5\baselineskip}
\subsection*{I}\BookHeading{nav-glossary-english-0010}{subsubsection}{I}{0}
% glossary-block:B00903:end
% glossary-block:B00904:begin
\textbf{I} — Ic is the nominative singular subject I. The governed object form me is mē in both dative and accusative. Select case from the verb/preposition and preserve dual versus plural; dual subjects use plural verb forms.\par
% glossary-block:B00904:end
% glossary-block:B00905:begin
\needspace{3\baselineskip}
\textbf{ice} — \textbf{īs} (n.).\par
% glossary-block:B00905:end
% glossary-block:B00906:begin
\needspace{3\baselineskip}
\textbf{idle} — \textbf{īdel} (adj.). Distinct from empty/void.\par
% glossary-block:B00906:end
% glossary-block:B00907:begin
\needspace{3\baselineskip}
\textbf{idleness} — \textbf{īdelnes} (n.).\par
% glossary-block:B00907:end
% glossary-block:B00908:begin
\needspace{3\baselineskip}
\textbf{if} — \textbf{gif} (conj.). Mood is determined by the open, tentative or counterfactual condition, not by gif alone. gif does not by itself choose a mood; open, tentative and counterfactual conditions have distinct taught frames.\par
% glossary-block:B00908:end
% glossary-block:B00909:begin
\needspace{3\baselineskip}
\textbf{ill} — \textbf{yfele} (adv.). Adverb modifying doing; distinguish the adjective yfel.\par
% glossary-block:B00909:end
% glossary-block:B00910:begin
\needspace{3\baselineskip}
\textbf{illuminate} — \textbf{ālīehtan} (v.). illuminate; give light to. Accusative person or thing illuminated. ālīehtan takes an accusative object, not a dative inferred from English to.\par
% glossary-block:B00910:end
% glossary-block:B00911:begin
\needspace{3\baselineskip}
\textbf{illumination} — \textbf{līehting} (n.).\par
% glossary-block:B00911:end
% glossary-block:B00912:begin
\needspace{3\baselineskip}
\textbf{illustration} — \textbf{bȳsen} (n.). example; illustration.\par
% glossary-block:B00912:end
% glossary-block:B00913:begin
\needspace{3\baselineskip}
\textbf{illustrious} — \textbf{æðele} (adj.). noble; illustrious.\par
% glossary-block:B00913:end
% glossary-block:B00914:begin
\needspace{3\baselineskip}
\textbf{image} — \textbf{andlīcnes} (n.).\par
% glossary-block:B00914:end
% glossary-block:B00915:begin
\needspace{3\baselineskip}
\textbf{imagine} — \textbf{āsmēan} (v.). conceive. Exact supplied infinitive āsmēan in the heart-conceiving expression; not an independently assigned full paradigm. Supplied infinitive āsmēan; the heart is its subject.\par
% glossary-block:B00915:end
% glossary-block:B00916:begin
\needspace{3\baselineskip}
\textbf{immediately} — \textbf{sōna} (adv.). immediately; at once. sōna carries soon in this future-reference prompt.\par
% glossary-block:B00916:end
% glossary-block:B00917:begin
\needspace{3\baselineskip}
\textbf{implore} — \textbf{hālsian} (v.). entreat; implore.\par
% glossary-block:B00917:end
% glossary-block:B00918:begin
\needspace{3\baselineskip}
\textbf{in} — \textbf{in} (prep.). in; on. With dative in the exact locations/circumstances. in; inward. Adverbial direction. \textbf{on} (prep.). in; on. Dative for the taught ordinary location construction. Supported expression in his majesty/glory. in; with respect to. Quality or domain of description. Supplied circumstance of the storm. Accusative person believed/trusted in. among; in. Dative people among whom friends are obtained. Circumstance of the day's heat in the supplied source. during. Dative time expression. in; wearing. Supplied on scēapa gegierelum, in sheep's garments. Source/object of joy in the supplied frame. in; amid. Circumstances of the stated unrest.\par
% glossary-block:B00918:end
% glossary-block:B00919:begin
\needspace{3\baselineskip}
\textbf{in accordance with} — \textbf{tō} (prep.). A will as the goal of guidance in the supplied prayer.\par
% glossary-block:B00919:end
% glossary-block:B00920:begin
\needspace{3\baselineskip}
\textbf{in addition to that} — \textbf{ðǣr-tō-ēacan} (adverbial expression). in addition to that; besides.\par
% glossary-block:B00920:end
% glossary-block:B00921:begin
\needspace{3\baselineskip}
\textbf{inasmuch as} — Swā lange swā is the complete supplied phrase as long as; in the assigned biblical context it has inasmuch-as force. Swā lange swā is the complete supplied phrase as long as; in the assigned biblical context it has inasmuch-as force. Lange alone means long.\par
% glossary-block:B00921:end
% glossary-block:B00922:begin
\needspace{3\baselineskip}
\textbf{in a storm} — \textbf{on} (prep.). in. Supplied circumstance of the storm.\par
% glossary-block:B00922:end
% glossary-block:B00923:begin
\needspace{3\baselineskip}
\textbf{in circuit of} — \textbf{ymbhwyrft} (n.). Supplied on ymbhwyrfte + genitive landmark in this exercise.\par
% glossary-block:B00923:end
% glossary-block:B00924:begin
\needspace{3\baselineskip}
\textbf{incline} — \textbf{āhyldan} (v.). lay; incline.\par
% glossary-block:B00924:end
% glossary-block:B00925:begin
\needspace{3\baselineskip}
\textbf{incline oneself} — \textbf{ālūtan} (v.). bow; incline oneself. Reflexive meaning in the assigned use.\par
% glossary-block:B00925:end
% glossary-block:B00926:begin
\needspace{3\baselineskip}
\textbf{in control} — \textbf{on} (prep.). in one's control. on + dative possessed geweald.\par
% glossary-block:B00926:end
% glossary-block:B00927:begin
\needspace{3\baselineskip}
\textbf{in control of} — On + dative possessed geweald expresses in/under someone’s control. In the assigned habban construction, the things controlled are its accusative object; the controller is expressed within the geweald phrase.\par
% glossary-block:B00927:end
% glossary-block:B00928:begin
\needspace{3\baselineskip}
\textbf{increase} — \textbf{īecan} (v.).  \textbf{weaxan} (v.). Intransitive increase in these commands/wishes; not the transitive īecan increase something.\par
% glossary-block:B00928:end
% glossary-block:B00929:begin
\needspace{3\baselineskip}
\textbf{indeed} — \textbf{sōðlice} (adv.). truly; indeed. Adverbial confirmation; distinct from discourse connective sōðlice where the key uses but/now.\par
% glossary-block:B00929:end
% glossary-block:B00930:begin
\needspace{3\baselineskip}
\textbf{indescribable} — \textbf{unāsęcgende} (adj.).\par
% glossary-block:B00930:end
% glossary-block:B00931:begin
\needspace{3\baselineskip}
\textbf{in exchange for} — \textbf{for} (prep.). With dative. \textbf{wið} (prep.). in return for. With dative; distinct from opponent/hostility. wið + dative exchange, distinct from wið + accusative opponent.\par
% glossary-block:B00931:end
% glossary-block:B00932:begin
\needspace{3\baselineskip}
\textbf{in fact} — \textbf{sōðlice} (adv.). Adverbial confirmation; distinct from discourse connective sōðlice where the key uses but/now.\par
% glossary-block:B00932:end
% glossary-block:B00933:begin
\needspace{3\baselineskip}
\textbf{infirmity} — \textbf{untrymnes} (n.).\par
% glossary-block:B00933:end
% glossary-block:B00934:begin
\needspace{3\baselineskip}
\textbf{inflict slaughter} — \textbf{ongeslēan} (v.). inflict slaughter; slay.\par
% glossary-block:B00934:end
% glossary-block:B00935:begin
\needspace{3\baselineskip}
\textbf{in forms of} — The supplied on + hīw expression means according to/in the forms of. Retain the whole supplied hīw phrase rather than treating on alone as a general according-to expression.\par
% glossary-block:B00935:end
% glossary-block:B00936:begin
\needspace{3\baselineskip}
\textbf{in front of} — \textbf{ætforan} (prep.). before; in front of. Dative landmark/person in the supplied frame; retain ætforan rather than replacing it with beforan. \textbf{be-foran} (prep.). before; in front of. Dative landmark in the static frame. before; in front of. Accusative landmark in the positioning construal; static dative can be a different valid construal.\par
% glossary-block:B00936:end
% glossary-block:B00937:begin
\needspace{3\baselineskip}
\textbf{in glory} — \textbf{on} (prep.). in. Supported expression in his majesty/glory.\par
% glossary-block:B00937:end
% glossary-block:B00938:begin
\needspace{3\baselineskip}
\textbf{in good health} — \textbf{gesund} (adj.). well; in good health.\par
% glossary-block:B00938:end
% glossary-block:B00939:begin
\needspace{3\baselineskip}
\textbf{ingot} — \textbf{węcg} (n.). metal; ingot. Supplied genitive plural dependent on veins.\par
% glossary-block:B00939:end
% glossary-block:B00940:begin
\needspace{3\baselineskip}
\textbf{in it} — \textbf{ðǣron} (adverb or particle). therein; in it. Supplied ðǣron points back to the sacks' location.\par
% glossary-block:B00940:end
% glossary-block:B00941:begin
\needspace{3\baselineskip}
\textbf{injure} — \textbf{scęððan} (v.). hurt; injure. Dative person harmed.\par
% glossary-block:B00941:end
% glossary-block:B00942:begin
\needspace{3\baselineskip}
\textbf{inn} — \textbf{giesthūs} (n.). inn; guest-house.\par
% glossary-block:B00942:end
% glossary-block:B00943:begin
\needspace{3\baselineskip}
\textbf{innocent} — \textbf{clǣne} (adj.). Source's blameless is contextual clǣne, not just physical clean.\par
% glossary-block:B00943:end
% glossary-block:B00944:begin
\needspace{3\baselineskip}
\textbf{innumerable} — \textbf{earfoðrīme} (adj.). innumerable; hard to count. Assigned source adjective.\par
% glossary-block:B00944:end
% glossary-block:B00945:begin
\needspace{3\baselineskip}
\textbf{in one's life} — \textbf{on} (prep.). during; in. Supplied lifetime expression.\par
% glossary-block:B00945:end
% glossary-block:B00946:begin
\needspace{3\baselineskip}
\textbf{in one's name} — On + dative naman with a genitive possessor means in someone’s name. The whole name phrase, not on alone, supplies this meaning.\par
% glossary-block:B00946:end
% glossary-block:B00947:begin
\needspace{3\baselineskip}
\textbf{in order that} — \textbf{ðæt} (conj.). so that; in order that. Purpose conjunction with the assigned subjunctive construction.\par
% glossary-block:B00947:end
% glossary-block:B00948:begin
\needspace{3\baselineskip}
\textbf{in order to} — Tō + inflected infinitive expresses the taught purpose, to/in order to. Plain infinitives after modals do not take an inserted tō; recipient to may instead be expressed by a dative.\par
% glossary-block:B00948:end
% glossary-block:B00949:begin
\needspace{3\baselineskip}
\textbf{in place of} — \textbf{for} (prep.). instead of; in place of. Accusative in the supplied substitution frame. for + accusative in this supplied substitution frame.\par
% glossary-block:B00949:end
% glossary-block:B00950:begin
\needspace{3\baselineskip}
\textbf{inquire after} — \textbf{æfter} (prep.; conj. in fixed expressions). after; in search of. Person sought or inquired after; keep this lexical frame distinct from temporal after.\par
% glossary-block:B00950:end
% glossary-block:B00951:begin
\needspace{3\baselineskip}
\textbf{in readiness for} — \textbf{ongēan} (adv.; prep.; conj.). by the time that. Complete supplied ongēan ðætte clause. Supplied ongēan ðætte + prospective subjunctive clause.\par
% glossary-block:B00951:end
% glossary-block:B00952:begin
\needspace{3\baselineskip}
\textbf{in response to} — \textbf{tō} (prep.). at; in response to. Only the exact at-the-prayers expression.\par
% glossary-block:B00952:end
% glossary-block:B00953:begin
\needspace{3\baselineskip}
\textbf{in return for} — \textbf{wið} (prep.). With dative; distinct from opponent/hostility. wið + dative exchange, distinct from wið + accusative opponent.\par
% glossary-block:B00953:end
% glossary-block:B00954:begin
\needspace{3\baselineskip}
\textbf{in search of} — \textbf{æfter} (prep.; conj. in fixed expressions). after; in search of. Person sought or inquired after; keep this lexical frame distinct from temporal after.\par
% glossary-block:B00954:end
% glossary-block:B00955:begin
\needspace{3\baselineskip}
\textbf{inside} — \textbf{innan} (adv.). within; inside.\par
% glossary-block:B00955:end
% glossary-block:B00956:begin
\needspace{3\baselineskip}
\textbf{instead of} — \textbf{for} (prep.). instead of; in place of. Accusative in the supplied substitution frame. for + accusative in this supplied substitution frame.\par
% glossary-block:B00956:end
% glossary-block:B00957:begin
\needspace{3\baselineskip}
\textbf{instruct} — \textbf{bebēodan} (v.). command; instruct. Dative addressee and a ðæt clause with subjunctive for the commanded action. \textbf{lǣran} (v.). advise; instruct. teach; instruct.\par
% glossary-block:B00957:end
% glossary-block:B00958:begin
\needspace{3\baselineskip}
\textbf{intelligibly} — Supplied superlative andgietfullicost means most intelligibly; do not use the positive andgietfullice for this degree.\par
% glossary-block:B00958:end
% glossary-block:B00959:begin
\needspace{3\baselineskip}
\textbf{intend} — \textbf{willan} (v.). Willan + plain infinitive expresses the reported intention in the assigned task. Distinguish intention from a mechanically inserted future auxiliary.\par
% glossary-block:B00959:end
% glossary-block:B00960:begin
\needspace{3\baselineskip}
\textbf{intend to} — \textbf{willan} (v.). Willan + plain infinitive expresses the reported intention in the assigned task. Distinguish intention from a mechanically inserted future auxiliary.\par
% glossary-block:B00960:end
% glossary-block:B00961:begin
\needspace{3\baselineskip}
\textbf{intent} — \textbf{geornfull} (adj.). occupied; intent. Source mental-attention context.\par
% glossary-block:B00961:end
% glossary-block:B00962:begin
\needspace{3\baselineskip}
\textbf{interpreter} — \textbf{wealhstōd} (n.). interpreter; translator.\par
% glossary-block:B00962:end
% glossary-block:B00963:begin
\needspace{3\baselineskip}
\textbf{interval} — \textbf{fæc} (n.). interval; period of time; while. Noun in the supplied time expression. \textbf{fierst} (n.). period; interval. Masculine noun in the source time expression.\par
% glossary-block:B00963:end
% glossary-block:B00964:begin
\needspace{3\baselineskip}
\textbf{in the beginning} — The supplied whole phrase on anginne means in the beginning.\par
% glossary-block:B00964:end
% glossary-block:B00965:begin
\needspace{3\baselineskip}
\textbf{in the heat} — \textbf{on} (prep.). in. Circumstance of the day's heat in the supplied source.\par
% glossary-block:B00965:end
% glossary-block:B00966:begin
\needspace{3\baselineskip}
\textbf{in the midst of} — \textbf{tōmiddes} (prep.). With dative landmark. tōmiddes + dative landmark; English of does not require genitive here.\par
% glossary-block:B00966:end
% glossary-block:B00967:begin
\needspace{3\baselineskip}
\textbf{in the morning} — The supplied whole phrase on morgen means in the morning; morgen is accusative singular here. On morgen is the supplied expression in the morning; morgen is accusative singular in this expression.\par
% glossary-block:B00967:end
% glossary-block:B00968:begin
\needspace{3\baselineskip}
\textbf{in the name of} — On + dative naman with a genitive possessor means in someone’s name. The whole name phrase, not on alone, supplies this meaning.\par
% glossary-block:B00968:end
% glossary-block:B00969:begin
\needspace{3\baselineskip}
\textbf{in the same way as} — \textbf{swā} (adverb; conjunction). as; in the same way as.\par
% glossary-block:B00969:end
% glossary-block:B00970:begin
\needspace{3\baselineskip}
\textbf{in the south} — The supplied whole phrase be sūðan means in the south. Be alone does not mean south. \textbf{sūðan} (adv.). in the south; from the south. Retain the supplied spatial expression.\par
% glossary-block:B00970:end
% glossary-block:B00971:begin
\needspace{3\baselineskip}
\textbf{into} — \textbf{in-tō} (prep.). into; destination expressed with dative after intō. \textbf{on} (prep.). into; in two. Supplied division expression. into; onto. Accusative goal of placement in the taught frame. upon; into. Supplied chancing-upon idiom. Destination of sale in the supplied expression. \textbf{tō} (prep.). to; into. Result in supplied tō lore, tō sorge or tō nāhte. into; to. Resulting language/literary form.\par
% glossary-block:B00971:end
% glossary-block:B00972:begin
\needspace{3\baselineskip}
\textbf{in two} — \textbf{on} (prep.). into; in two. Supplied division expression. \textbf{twēgen} (cardinal numeral). two; in two. Two parts in the supplied division expression.\par
% glossary-block:B00972:end
% glossary-block:B00973:begin
\needspace{3\baselineskip}
\textbf{in vain} — \textbf{hōlunga} (adv.). hōlunga; not at all in vain uses nālæs hōlunga with the intended negative scope.\par
% glossary-block:B00973:end
% glossary-block:B00974:begin
\needspace{3\baselineskip}
\textbf{invisible} — \textbf{ungesewenlic} (adj.).\par
% glossary-block:B00974:end
% glossary-block:B00975:begin
\needspace{3\baselineskip}
\textbf{inward} — \textbf{in} (prep.). in; inward. Adverbial direction.\par
% glossary-block:B00975:end
% glossary-block:B00976:begin
\needspace{3\baselineskip}
\textbf{in worldly life} — \textbf{on} (prep.). Supplied temporal/situational expression.\par
% glossary-block:B00976:end
% glossary-block:B00977:begin
\needspace{3\baselineskip}
\textbf{iron} — \textbf{īsern} (n.). Metal noun. \textbf{īsern} (adj.). iron; made of iron. Adjective of material.\par
% glossary-block:B00977:end
% glossary-block:B00978:begin
\needspace{3\baselineskip}
\textbf{is} — \textbf{wesan} (v.). be. Copular present/past identification or state; other auxiliary constructions are separate senses. Copular eom/wæs and bēo indicative series; later passive and compound uses require their own construction.\par
% glossary-block:B00978:end
% glossary-block:B00979:begin
\needspace{3\baselineskip}
\textbf{island} — \textbf{īgland} (n.).\par
% glossary-block:B00979:end
% glossary-block:B00980:begin
\needspace{3\baselineskip}
\textbf{Israel} — \textbf{Israel} (proper name).\par
% glossary-block:B00980:end
% glossary-block:B00981:begin
\needspace{3\baselineskip}
\textbf{Israelites} — \textbf{Israela} (proper name). Supplied genitive plural name form.\par
% glossary-block:B00981:end
% glossary-block:B00982:begin
\needspace{3\baselineskip}
\textbf{is to come} — \textbf{tō} (prep.). Whole supplied is tō cumanne; bounded be-to use, distinct from purpose after a verb of coming.\par
% glossary-block:B00982:end
% glossary-block:B00983:begin
\needspace{3\baselineskip}
\textbf{it} — Hit is neuter singular nominative/accusative; dative is him and genitive his. An inanimate referent of masculine or feminine grammatical gender takes the corresponding gendered pronoun instead. Grammatical gender and antecedent determine reference; it may be represented by masculine/feminine forms for a gendered inanimate noun.\par
% glossary-block:B00983:end
% glossary-block:B00984:begin
\needspace{3\baselineskip}
\textbf{it (anticipatory or impersonal)} — Hit occurs in some taught impersonal/anticipatory frames; other supplied passives expressly omit it. Follow that whole frame.\par
% glossary-block:B00984:end
% glossary-block:B00985:begin
\needspace{3\baselineskip}
\textbf{its} — \textbf{hē} (personal pron.). his; her; its; their. Invariant genitive possessors, unlike agreeing possessive adjectives. Invariant personal genitive possessor; distinguish the same surface used as a governed personal case.\par
% glossary-block:B00985:end
% glossary-block:B00986:begin
\needspace{3\baselineskip}
\textbf{itself} — Emphatic self agrees with the person or thing emphasized: supplied dēaðes selfes means death itself. This is emphasis on an inanimate thing, not a reflexive object.\par
% glossary-block:B00986:end
% glossary-block:B00987:begin
\needspace{3\baselineskip}
\textbf{ivory} — \textbf{bān} (n.). bone; ivory. Ivory is elpes bān, literally elephant's bone. \textbf{elpenbǣnen} (adj.). ivory; made of ivory. Adjective of material.\par
% glossary-block:B00987:end
% book-navigation: standalone bookmark replaced by volume hierarchy
% glossary-block:B00988:begin
\needspace{5\baselineskip}
\subsection*{J}\BookHeading{nav-glossary-english-0011}{subsubsection}{J}{0}
% glossary-block:B00988:end
% glossary-block:B00989:begin
\textbf{Jacob} — \textbf{Iacob} (proper name).\par
% glossary-block:B00989:end
% glossary-block:B00990:begin
\needspace{3\baselineskip}
\textbf{Jerusalem} — \textbf{Hierusalem} (proper name). English spelling Jerusalem points to the retained Old English name Hierusalem.\par
% glossary-block:B00990:end
% glossary-block:B00991:begin
\needspace{3\baselineskip}
\textbf{Jesus} — \textbf{Hǣlend} (n.). Saviour; Jesus. Proper title.\par
% glossary-block:B00991:end
% glossary-block:B00992:begin
\needspace{3\baselineskip}
\textbf{jet} — \textbf{gagates} (n.). jet (Latin loan-name). The stone; Latin loan-name in the source.\par
% glossary-block:B00992:end
% glossary-block:B00993:begin
\needspace{3\baselineskip}
\textbf{Jews} — \textbf{Iudea} (proper name). Supplied genitive plural name form.\par
% glossary-block:B00993:end
% glossary-block:B00994:begin
\needspace{3\baselineskip}
\textbf{John} — \textbf{Johannes} (proper name).\par
% glossary-block:B00994:end
% glossary-block:B00995:begin
\needspace{3\baselineskip}
\textbf{join} — \textbf{gemęngan} (v.). join; mingle. Supplied source form.\par
% glossary-block:B00995:end
% glossary-block:B00996:begin
\needspace{3\baselineskip}
\textbf{Joseph} — \textbf{Iosep} (proper name).\par
% glossary-block:B00996:end
% glossary-block:B00997:begin
\needspace{3\baselineskip}
\textbf{journey} — \textbf{sīðfæt} (n.). journey; passage. \textbf{faran} (v.). Verb faran; distinct from noun journey sīðfæt.\par
% glossary-block:B00997:end
% glossary-block:B00998:begin
\needspace{3\baselineskip}
\textbf{joy} — \textbf{drēam} (n.). joy; bliss.\par
% glossary-block:B00998:end
% glossary-block:B00999:begin
\needspace{3\baselineskip}
\textbf{joyfully} — \textbf{blīðe} (adv.). joyfully; gladly.\par
% glossary-block:B00999:end
% glossary-block:B01000:begin
\needspace{3\baselineskip}
\textbf{Judah} — \textbf{Iudas} (proper name).\par
% glossary-block:B01000:end
% glossary-block:B01001:begin
\needspace{3\baselineskip}
\textbf{judge} — \textbf{dēma} (n.).  \textbf{dēman} (v.).\par
% glossary-block:B01001:end
% glossary-block:B01002:begin
\needspace{3\baselineskip}
\textbf{judgment} — \textbf{dōm} (n.). dōmes dæg means Day of Judgment. A judgment made in judging, distinct from doom/fate.\par
% glossary-block:B01002:end
% glossary-block:B01003:begin
\needspace{3\baselineskip}
\textbf{just} — \textbf{eall} (adj.; substantive; adv.). just; exactly. Degree emphasis in eall swā. \textbf{efne} (interj.). just; precisely; even. \textbf{riht} (adj.). right; just.\par
% glossary-block:B01003:end
% glossary-block:B01004:begin
\needspace{3\baselineskip}
\textbf{just as} — \textbf{swā-swā} (conj.).\par
% glossary-block:B01004:end
% glossary-block:B01005:begin
\needspace{3\baselineskip}
\textbf{just as … as} — Swā + adjective/adverb + swā frames equal degree, as … as. The first swā marks degree; the second introduces the standard. Comparative the … the uses a separate construction.\par
% glossary-block:B01005:end
% book-navigation: standalone bookmark replaced by volume hierarchy
% glossary-block:B01006:begin
\needspace{5\baselineskip}
\subsection*{K}\BookHeading{nav-glossary-english-0012}{subsubsection}{K}{0}
% glossary-block:B01006:end
% glossary-block:B01007:begin
\textbf{keep} — \textbf{healdan} (v.). keep; observe. Accusative saying or command observed.\par
% glossary-block:B01007:end
% glossary-block:B01008:begin
\needspace{3\baselineskip}
\textbf{keep silent} — \textbf{swigian} (v.). be silent. Supplied source sense.\par
% glossary-block:B01008:end
% glossary-block:B01009:begin
\needspace{3\baselineskip}
\textbf{keep silent about} — \textbf{forswīgian} (v.). keep silent about; conceal.\par
% glossary-block:B01009:end
% glossary-block:B01010:begin
\needspace{3\baselineskip}
\textbf{kill} — \textbf{cwellan} (v.).  \textbf{ofslēan} (v.). slay; kill. \textbf{slēan} (v.). slay; kill. Accusative person killed. Accusative person; the bounded negative command uses selected short-vowel imperative sleah.\par
% glossary-block:B01010:end
% glossary-block:B01011:begin
\needspace{3\baselineskip}
\textbf{kind} — \textbf{cynn} (n.).  \textbf{gōd} (adj.). kind; benevolent. Gōd means kind or benevolent in the taught person-directed frame. \textbf{hīw} (n.). form; kind. Supplied dative neuter plural.\par
% glossary-block:B01011:end
% glossary-block:B01012:begin
\needspace{3\baselineskip}
\textbf{kindness} — \textbf{ār-fæstnes} (n.).  \textbf{gōd} (n.). good; benefit; kindness. Nominal good in the repayment contrast.\par
% glossary-block:B01012:end
% glossary-block:B01013:begin
\needspace{3\baselineskip}
\textbf{kinds of fish} — \textbf{fisccynn} (n.). Supplied accusative neuter plural.\par
% glossary-block:B01013:end
% glossary-block:B01014:begin
\needspace{3\baselineskip}
\textbf{kind toward} — \textbf{ofer} (prep.). toward. In gōd ofer + accusative people.\par
% glossary-block:B01014:end
% glossary-block:B01015:begin
\needspace{3\baselineskip}
\textbf{king} — \textbf{cyning} (n.).  \textbf{ealdor} (n.). chief; king. This noun is distinct from the homonymous noun meaning life.\par
% glossary-block:B01015:end
% glossary-block:B01016:begin
\needspace{3\baselineskip}
\textbf{kingdom} — \textbf{rīce} (n.).\par
% glossary-block:B01016:end
% glossary-block:B01017:begin
\needspace{3\baselineskip}
\textbf{kiss} — \textbf{coss} (n.).  \textbf{cyssan} (v.).\par
% glossary-block:B01017:end
% glossary-block:B01018:begin
\needspace{3\baselineskip}
\textbf{knee} — \textbf{cnēo} (n.). Supplied accusative plural knees.\par
% glossary-block:B01018:end
% glossary-block:B01019:begin
\needspace{3\baselineskip}
\textbf{knock} — \textbf{cnucian} (v.).\par
% glossary-block:B01019:end
% glossary-block:B01020:begin
\needspace{3\baselineskip}
\textbf{know} — \textbf{cnāwan} (v.).  \textbf{cunnan} (v.). know; be acquainted with. Accusative person known; distinct from learned ability. \textbf{witan} (v.). A stated fact or the answer to a question as complement. Accusative person in this taught frame; cunnan has a separate acquaintance sense. \textbf{cunnan} (v.), \textbf{witan} (v.). Cunnan + accusative person expresses acquaintance. Witan can take an accusative person in its taught personal sense, or a fact/question as content. Learned ability is the separate cunnan + plain infinitive construction.\par
% glossary-block:B01020:end
% glossary-block:B01021:begin
\needspace{3\baselineskip}
\textbf{know a person} — \textbf{witan} (v.). know. Accusative person in this taught frame; cunnan has a separate acquaintance sense.\par
% glossary-block:B01021:end
% glossary-block:B01022:begin
\needspace{3\baselineskip}
\textbf{know by recognition} — \textbf{oncnāwan} (v.). recognize; know by recognition.\par
% glossary-block:B01022:end
% glossary-block:B01023:begin
\needspace{3\baselineskip}
\textbf{know how} — \textbf{cunnan} (v.). know how; be able. Learned skill; plain infinitive complement.\par
% glossary-block:B01023:end
% glossary-block:B01024:begin
\needspace{3\baselineskip}
\textbf{knowledge} — \textbf{andgiet} (n.). andgiet denotes the knowledge possessed and desired in this exercise.\par
% glossary-block:B01024:end
% glossary-block:B01025:begin
\needspace{3\baselineskip}
\textbf{know that} — \textbf{witan} (v.). know. A stated fact or the answer to a question as complement.\par
% glossary-block:B01025:end
% glossary-block:B01026:begin
\needspace{3\baselineskip}
\textbf{know whether} — \textbf{witan} (v.). know. A stated fact or the answer to a question as complement.\par
% glossary-block:B01026:end
% book-navigation: standalone bookmark replaced by volume hierarchy
% glossary-block:B01027:begin
\needspace{5\baselineskip}
\subsection*{L}\BookHeading{nav-glossary-english-0013}{subsubsection}{L}{0}
% glossary-block:B01027:end
% glossary-block:B01028:begin
\textbf{labor} — \textbf{gedeorf} (n.). labor; toil. \textbf{gewinn} (n.). labour; work. Noun, distinct from labour as a verb. Noun in the expression from his labour. \textbf{winnan} (v.). Verb; the subject labours with the hands.\par
% glossary-block:B01028:end
% glossary-block:B01029:begin
\needspace{3\baselineskip}
\textbf{labour} — \textbf{gewinn} (n.). labour; work. Noun, distinct from labour as a verb. Noun in the expression from his labour. \textbf{winnan} (v.). Verb; the subject labours with the hands.\par
% glossary-block:B01029:end
% glossary-block:B01030:begin
\needspace{3\baselineskip}
\textbf{lad} — \textbf{cnapa} (n.). Masculine weak noun; the task's chosen word for lad.\par
% glossary-block:B01030:end
% glossary-block:B01031:begin
\needspace{3\baselineskip}
\textbf{lady} — \textbf{hlǣfdige} (n.).\par
% glossary-block:B01031:end
% glossary-block:B01032:begin
\needspace{3\baselineskip}
\textbf{lame} — \textbf{healt} (adj.). lame; halt. Adjective describing impaired walking; not the command halt/stop.\par
% glossary-block:B01032:end
% glossary-block:B01033:begin
\needspace{3\baselineskip}
\textbf{land} — \textbf{land} (n.). land; country.\par
% glossary-block:B01033:end
% glossary-block:B01034:begin
\needspace{3\baselineskip}
\textbf{large} — \textbf{micel} (adj.). great; large.\par
% glossary-block:B01034:end
% glossary-block:B01035:begin
\needspace{3\baselineskip}
\textbf{larger} — \textbf{māra} (adj.). greater; larger. Comparative adjective.\par
% glossary-block:B01035:end
% glossary-block:B01036:begin
\needspace{3\baselineskip}
\textbf{largest} — \textbf{mǣst} (adj.). greatest; largest. Superlative adjective.\par
% glossary-block:B01036:end
% glossary-block:B01037:begin
\needspace{3\baselineskip}
\textbf{last} — \textbf{ðurhwunian} (v.). remain; endure; last.\par
% glossary-block:B01037:end
% glossary-block:B01038:begin
\needspace{3\baselineskip}
\textbf{Latin} — \textbf{Bōclæden} (n.). book-Latin; Latin. \textbf{Læden} (n.). The language.\par
% glossary-block:B01038:end
% glossary-block:B01039:begin
\needspace{3\baselineskip}
\textbf{law} — \textbf{ǣ} (n.).\par
% glossary-block:B01039:end
% glossary-block:B01040:begin
\needspace{3\baselineskip}
\textbf{lay} — \textbf{āhyldan} (v.). lay; incline. \textbf{lęcgean} (v.). lay; place. Supplied source form. lay; place.\par
% glossary-block:B01040:end
% glossary-block:B01041:begin
\needspace{3\baselineskip}
\textbf{layman} — \textbf{ceorl} (n.).\par
% glossary-block:B01041:end
% glossary-block:B01042:begin
\needspace{3\baselineskip}
\textbf{lay out} — \textbf{lęcgean} (v.). lęcgean + accusative thing + forð.\par
% glossary-block:B01042:end
% glossary-block:B01043:begin
\needspace{3\baselineskip}
\textbf{lay waste} — \textbf{forniman} (v.). destroy; lay waste.\par
% glossary-block:B01043:end
% glossary-block:B01044:begin
\needspace{3\baselineskip}
\textbf{Lazarus} — \textbf{Lazarus} (proper name). The parable's poor man; supplied plural denotes people like him.\par
% glossary-block:B01044:end
% glossary-block:B01045:begin
\needspace{3\baselineskip}
\textbf{Lazaruses} — \textbf{Lazarus} (proper name). Lazarus. The parable's poor man; supplied plural denotes people like him.\par
% glossary-block:B01045:end
% glossary-block:B01046:begin
\needspace{3\baselineskip}
\textbf{lead} — \textbf{lǣdan} (v.). lead; bring; take. A person is the object in the assigned frame. \textbf{lēad} (n.). Metal noun, not the verb lead.\par
% glossary-block:B01046:end
% glossary-block:B01047:begin
\needspace{3\baselineskip}
\textbf{leader} — \textbf{ealdormann} (n.). prince; leader; noble.\par
% glossary-block:B01047:end
% glossary-block:B01048:begin
\needspace{3\baselineskip}
\textbf{learn} — \textbf{geleornian} (v.). learn; study. learn; master. Learn a skill or performance in the source. \textbf{gewitan} (v.). find out; learn. Knowledge verb, distinct from gewītan depart. \textbf{leornian} (v.). Accusative thing learned; æt + dative person/source.\par
% glossary-block:B01048:end
% glossary-block:B01049:begin
\needspace{3\baselineskip}
\textbf{learn from} — \textbf{æt} (prep.). from. æt + dative person with learn. leornian æt + dative person in the assigned learning frame.\par
% glossary-block:B01049:end
% glossary-block:B01050:begin
\needspace{3\baselineskip}
\textbf{learning} — \textbf{lār} (n.). learning; study. \textbf{leornung} (n.). study; learning.\par
% glossary-block:B01050:end
% glossary-block:B01051:begin
\needspace{3\baselineskip}
\textbf{learn of} — \textbf{æt} (prep.). from. æt + dative person with learn. leornian æt + dative person in the assigned learning frame.\par
% glossary-block:B01051:end
% glossary-block:B01052:begin
\needspace{3\baselineskip}
\textbf{least} — \textbf{lǣst} (adj.). In rank or importance.\par
% glossary-block:B01052:end
% glossary-block:B01053:begin
\needspace{3\baselineskip}
\textbf{leave} — \textbf{forlǣtan} (v.). leave to. Thing left and recipient retain their taught roles. Accusative thing/place left. Distinguish leave to someone and take leave of.\par
% glossary-block:B01053:end
% glossary-block:B01054:begin
\needspace{3\baselineskip}
\textbf{leave to} — \textbf{forlǣtan} (v.). Thing left and recipient retain their taught roles.\par
% glossary-block:B01054:end
% glossary-block:B01055:begin
\needspace{3\baselineskip}
\textbf{left} — \textbf{winstre} (adj.). Supplied side phrase.\par
% glossary-block:B01055:end
% glossary-block:B01056:begin
\needspace{3\baselineskip}
\textbf{left-hand side} — \textbf{winstre} (adj.). left. Supplied side phrase.\par
% glossary-block:B01056:end
% glossary-block:B01057:begin
\needspace{3\baselineskip}
\textbf{leisure} — \textbf{ǣmetta} (n.). leisure; available time.\par
% glossary-block:B01057:end
% glossary-block:B01058:begin
\needspace{3\baselineskip}
\textbf{less} — \textbf{lǣs} (degree word). less; a smaller amount.\par
% glossary-block:B01058:end
% glossary-block:B01059:begin
\needspace{3\baselineskip}
\textbf{lest} — \textbf{ðȳ-lǣs-ðe} (conj.). lest; so that not. ðȳ-lǣs-ðe with the assigned subjunctive negative-purpose clause.\par
% glossary-block:B01059:end
% glossary-block:B01060:begin
\needspace{3\baselineskip}
\textbf{let} — Lǣtan is lexical let/allow. Wishes introduced by English let use the taught independent subjunctive; they do not require an inserted lǣtan. An independent wish/hortatory uses a subjunctive; let does not automatically require lǣtan. Literal permission to go uses lǣtan + accusative actor + infinitive.\par
% glossary-block:B01060:end
% glossary-block:B01061:begin
\needspace{3\baselineskip}
\textbf{let go} — \textbf{forlǣtan} (v.). let go; dismiss.\par
% glossary-block:B01061:end
% glossary-block:B01062:begin
\needspace{3\baselineskip}
\textbf{letter} — \textbf{gewrit} (n.). letter; written document.\par
% glossary-block:B01062:end
% glossary-block:B01063:begin
\needspace{3\baselineskip}
\textbf{levin} — \textbf{līeg} (n.). thunderbolt; lightning. Source levin sense.\par
% glossary-block:B01063:end
% glossary-block:B01064:begin
\needspace{3\baselineskip}
\textbf{lick} — \textbf{liccian} (v.).\par
% glossary-block:B01064:end
% glossary-block:B01065:begin
\needspace{3\baselineskip}
\textbf{lie} — \textbf{licgan} (v.). lie; be at rest. lie; be lying. Supplied participle describing the possessed persons, not a habban perfect.\par
% glossary-block:B01065:end
% glossary-block:B01066:begin
\needspace{3\baselineskip}
\textbf{life} — \textbf{līf} (n.). An individual's lived life, including cognate accusative with libban.\par
% glossary-block:B01066:end
% glossary-block:B01067:begin
\needspace{3\baselineskip}
\textbf{lift} — \textbf{āhębban} (v.). raise; lift. \textbf{hębban} (v.). raise; lift.\par
% glossary-block:B01067:end
% glossary-block:B01068:begin
\needspace{3\baselineskip}
\textbf{lift up} — \textbf{ūp-ā-hębban} (v.). Supplied source form. Supplied first-plural preterite. Accusative people raised; the prefixed verb supplies up.\par
% glossary-block:B01068:end
% glossary-block:B01069:begin
\needspace{3\baselineskip}
\textbf{light} — \textbf{lēoht} (n.). light; luminary. Individual countable light, distinct from mass light. Mass light and an individual countable light/luminary are separate uses of lēoht. Use the count sense for a light or several lights.\par
% glossary-block:B01069:end
% glossary-block:B01070:begin
\needspace{3\baselineskip}
\textbf{lightning} — \textbf{līeg} (n.). thunderbolt; lightning. Source levin sense.\par
% glossary-block:B01070:end
% glossary-block:B01071:begin
\needspace{3\baselineskip}
\textbf{like} — \textbf{gelīc} (adj.). equal; like. Dative standard of comparison. \textbf{swā} (adverb; conjunction). as; like. Predicative comparison. Swā can introduce a predicative comparison, as/like; swā-swā means just as in its supplied comparison. Swilce means as if in the supplied apparent-rejuvenation comparison. Equal degree uses the separate paired swā … swā construction.\par
% glossary-block:B01071:end
% glossary-block:B01072:begin
\needspace{3\baselineskip}
\textbf{likeness} — \textbf{gelīcnes} (n.).\par
% glossary-block:B01072:end
% glossary-block:B01073:begin
\needspace{3\baselineskip}
\textbf{likewise} — \textbf{gelīce} (adv.). likewise; similarly. \textbf{swilce} (adv.). Adverb, distinct from inflected swilc such.\par
% glossary-block:B01073:end
% glossary-block:B01074:begin
\needspace{3\baselineskip}
\textbf{linen cloth} — \textbf{scīete} (n.).\par
% glossary-block:B01074:end
% glossary-block:B01075:begin
\needspace{3\baselineskip}
\textbf{lion} — \textbf{lēo} (n.). lion (Latin loan-form).\par
% glossary-block:B01075:end
% glossary-block:B01076:begin
\needspace{3\baselineskip}
\textbf{listen} — \textbf{gehīeran} (v.). hear; listen.\par
% glossary-block:B01076:end
% glossary-block:B01077:begin
\needspace{3\baselineskip}
\textbf{listen to} — \textbf{hlystan} (v.). listen to; heed. Genitive person or thing listened to; hear is the contextual prompt wording. hlystan governs genitive of the speech heard; do not copy hīeran's ordinary accusative frame.\par
% glossary-block:B01077:end
% glossary-block:B01078:begin
\needspace{3\baselineskip}
\textbf{little} — \textbf{lȳtel} (adj.). little; small.\par
% glossary-block:B01078:end
% glossary-block:B01079:begin
\needspace{3\baselineskip}
\textbf{little by little} — \textbf{styccemǣlum} (adv.). Single adverb styccemǣlum.\par
% glossary-block:B01079:end
% glossary-block:B01080:begin
\needspace{3\baselineskip}
\textbf{live} — \textbf{libban} (v.). The bracketed liveth yet is rendered by finite libban + gīet; an added copula is not required.\par
% glossary-block:B01080:end
% glossary-block:B01081:begin
\needspace{3\baselineskip}
\textbf{lives} — \textbf{līf} (n.). life. An individual's lived life, including cognate accusative with libban.\par
% glossary-block:B01081:end
% glossary-block:B01082:begin
\needspace{3\baselineskip}
\textbf{living} — \textbf{cucu} (adj.). alive; living. Supplied accusative neuter plural adjective. \textbf{cwic} (adj.). alive; living. \textbf{libban} (v.). Supplied weak participial adjective.\par
% glossary-block:B01082:end
% glossary-block:B01083:begin
\needspace{3\baselineskip}
\textbf{living people} — \textbf{cwic} (adj.). Substantival adjective.\par
% glossary-block:B01083:end
% glossary-block:B01084:begin
\needspace{3\baselineskip}
\textbf{lodge} — \textbf{wician} (v.). sojourn; lodge. Intransitive; mid + dative person with whom one lodges.\par
% glossary-block:B01084:end
% glossary-block:B01085:begin
\needspace{3\baselineskip}
\textbf{long} — \textbf{lang} (adj.). Adjective agreeing with the time interval fæc. \textbf{lange} (adv.). In supplied swā lange swā; its biblical inasmuch-as force is explained locally.\par
% glossary-block:B01085:end
% glossary-block:B01086:begin
\needspace{3\baselineskip}
\textbf{longer} — \textbf{lęng} (adverb or particle). lęng as the comparative adverb; clause negation supplies no longer.\par
% glossary-block:B01086:end
% glossary-block:B01087:begin
\needspace{3\baselineskip}
\textbf{look} — \textbf{besēon} (v.). look; turn one's gaze. Supplied source form and frame.\par
% glossary-block:B01087:end
% glossary-block:B01088:begin
\needspace{3\baselineskip}
\textbf{look at} — \textbf{behealdan} (v.). behold; look at. Accusative object seen. observe. Supplied source form observed. behealdan takes accusative objects; distinguish interjection efne behold.\par
% glossary-block:B01088:end
% glossary-block:B01089:begin
\needspace{3\baselineskip}
\textbf{look for} — \textbf{sēcan} (v.). sēcan with accusative sought thing, no extra word for English for/out.\par
% glossary-block:B01089:end
% glossary-block:B01090:begin
\needspace{3\baselineskip}
\textbf{look toward} — \textbf{besēon} (v.). Target of the look follows the supplied phrase.\par
% glossary-block:B01090:end
% glossary-block:B01091:begin
\needspace{3\baselineskip}
\textbf{lord} — \textbf{dryhten} (n.). lord; the Lord. \textbf{hlāford} (n.). lord; master.\par
% glossary-block:B01091:end
% glossary-block:B01092:begin
\needspace{3\baselineskip}
\textbf{lose} — \textbf{forlēosan} (v.). Accusative thing lost.\par
% glossary-block:B01092:end
% glossary-block:B01093:begin
\needspace{3\baselineskip}
\textbf{loss} — \textbf{lyre} (n.).\par
% glossary-block:B01093:end
% glossary-block:B01094:begin
\needspace{3\baselineskip}
\textbf{loud} — \textbf{hlūd} (adj.).  \textbf{micel} (adj.). great; loud. Loud is the contextual force with a voice.\par
% glossary-block:B01094:end
% glossary-block:B01095:begin
\needspace{3\baselineskip}
\textbf{love} — \textbf{lufe} (n.).  \textbf{lufian} (v.).  \textbf{lufu} (n.).\par
% glossary-block:B01095:end
% glossary-block:B01096:begin
\needspace{3\baselineskip}
\textbf{lover} — \textbf{lufiend} (n.).\par
% glossary-block:B01096:end
% glossary-block:B01097:begin
\needspace{3\baselineskip}
\textbf{lowly} — \textbf{ēaðmōd} (adj.). lowly; humble.\par
% glossary-block:B01097:end
% glossary-block:B01098:begin
\needspace{3\baselineskip}
\textbf{luminary} — \textbf{lēoht} (n.). light; luminary. Individual countable light, distinct from mass light. Mass light and an individual countable light/luminary are separate uses of lēoht. Use the count sense for a light or several lights.\par
% glossary-block:B01098:end
% glossary-block:B01099:begin
\needspace{3\baselineskip}
\textbf{lust} — \textbf{gālnes} (n.). Assigned source sense.\par
% glossary-block:B01099:end
% book-navigation: standalone bookmark replaced by volume hierarchy
% glossary-block:B01100:begin
\needspace{5\baselineskip}
\subsection*{M}\BookHeading{nav-glossary-english-0014}{subsubsection}{M}{0}
% glossary-block:B01100:end
% glossary-block:B01101:begin
\textbf{made of iron} — \textbf{īsern} (adj.). iron; made of iron. Adjective of material.\par
% glossary-block:B01101:end
% glossary-block:B01102:begin
\needspace{3\baselineskip}
\textbf{made of ivory} — \textbf{elpenbǣnen} (adj.). ivory; made of ivory. Adjective of material.\par
% glossary-block:B01102:end
% glossary-block:B01103:begin
\needspace{3\baselineskip}
\textbf{made of silver} — \textbf{seolfren} (adj.). silver; made of silver. Material adjective, distinct from the noun seolfor.\par
% glossary-block:B01103:end
% glossary-block:B01104:begin
\needspace{3\baselineskip}
\textbf{made of stone} — \textbf{stānen} (adj.). stone; made of stone. Material adjective.\par
% glossary-block:B01104:end
% glossary-block:B01105:begin
\needspace{3\baselineskip}
\textbf{magnify} — \textbf{mǣrsian} (v.). make great; magnify; glorify.\par
% glossary-block:B01105:end
% glossary-block:B01106:begin
\needspace{3\baselineskip}
\textbf{maiden} — \textbf{fǣmne} (n.). maiden; young woman. \textbf{mǣden} (n.). girl; maiden.\par
% glossary-block:B01106:end
% glossary-block:B01107:begin
\needspace{3\baselineskip}
\textbf{maidenhood} — \textbf{mægðhād} (n.). maidenhood; virginity.\par
% glossary-block:B01107:end
% glossary-block:B01108:begin
\needspace{3\baselineskip}
\textbf{maintain} — \textbf{healdan} (v.). maintain; sustain.\par
% glossary-block:B01108:end
% glossary-block:B01109:begin
\needspace{3\baselineskip}
\textbf{majesty} — \textbf{mægenðrymm} (n.). majesty; glory. \textbf{mægenðrymnes} (n.). glory; majesty.\par
% glossary-block:B01109:end
% glossary-block:B01110:begin
\needspace{3\baselineskip}
\textbf{make} — \textbf{dōn} (v.). make; render. Accusative object with agreeing object predicate. \textbf{gedōn} (v.). do; make. Supplied passive fulfilment formula. make; arrange. The present form/arrangement of the translation in the source. make; cause to be. \textbf{gewyrcean} (v.). make; build. make; render. Accusative object with an agreeing predicate noun/adjective. \textbf{wyrcean} (v.). make; build. Supplied unprefixed form; preserve explicit task controls. \textbf{wyrcean} (v.), \textbf{gewyrcean} (v.). Wyrcean and gewyrcean mean make/build in these senses. Preserve the prefixed or unprefixed verb selected by the exercise.\par
% glossary-block:B01110:end
% glossary-block:B01111:begin
\needspace{3\baselineskip}
\textbf{make a decree} — \textbf{sęttan} (v.). appoint; establish. With accusative entity/decree in the assigned frame.\par
% glossary-block:B01111:end
% glossary-block:B01112:begin
\needspace{3\baselineskip}
\textbf{make a signal} — The supplied expression slōg tācen means made a signal; tācen is its object. Slēan alone has the separate strike sense.\par
% glossary-block:B01112:end
% glossary-block:B01113:begin
\needspace{3\baselineskip}
\textbf{make clear} — \textbf{gesweotolian} (v.). reveal; make clear. Supplied source form.\par
% glossary-block:B01113:end
% glossary-block:B01114:begin
\needspace{3\baselineskip}
\textbf{make great} — \textbf{mǣrsian} (v.). make great; magnify; glorify.\par
% glossary-block:B01114:end
% glossary-block:B01115:begin
\needspace{3\baselineskip}
\textbf{make known} — \textbf{cȳðan} (v.). announce; make known. Dative person informed; accusative thing announced or a ðæt content clause.\par
% glossary-block:B01115:end
% glossary-block:B01116:begin
\needspace{3\baselineskip}
\textbf{make manifest} — \textbf{gesweotolian} (v.). make manifest; reveal.\par
% glossary-block:B01116:end
% glossary-block:B01117:begin
\needspace{3\baselineskip}
\textbf{Maker} — \textbf{wyrhta} (n.). Maker; Creator. Divine title, distinct from ordinary workman.\par
% glossary-block:B01117:end
% glossary-block:B01118:begin
\needspace{3\baselineskip}
\textbf{make steadfast} — \textbf{gestaðelian} (v.). establish; make steadfast.\par
% glossary-block:B01118:end
% glossary-block:B01119:begin
\needspace{3\baselineskip}
\textbf{man} — \textbf{mann} (n.; impersonal pronoun). man; human being. \textbf{wer} (n.). Supplied genitive singular.\par
% glossary-block:B01119:end
% glossary-block:B01120:begin
\needspace{3\baselineskip}
\textbf{manifold} — \textbf{manigfeald} (adj.). manifold; many.\par
% glossary-block:B01120:end
% glossary-block:B01121:begin
\needspace{3\baselineskip}
\textbf{many} — \textbf{fela} (adj.). many; a large number of. Genitive of the things counted in the taught frame. \textbf{manig} (adj.).  \textbf{manigfeald} (adj.). manifold; many.\par
% glossary-block:B01121:end
% glossary-block:B01122:begin
\needspace{3\baselineskip}
\textbf{marine} — \textbf{sǣlic} (adj.). marine; of the sea.\par
% glossary-block:B01122:end
% glossary-block:B01123:begin
\needspace{3\baselineskip}
\textbf{Mary} — \textbf{Maria} (proper name).\par
% glossary-block:B01123:end
% glossary-block:B01124:begin
\needspace{3\baselineskip}
\textbf{master} — \textbf{geleornian} (v.). learn; master. Learn a skill or performance in the source. \textbf{hlāford} (n.). lord; master. \textbf{lārēow} (n.). master; teacher.\par
% glossary-block:B01124:end
% glossary-block:B01125:begin
\needspace{3\baselineskip}
\textbf{matter} — \textbf{ðing} (n.). thing; fact; matter. Matter known/discussed, distinct from a generic physical object.\par
% glossary-block:B01125:end
% glossary-block:B01126:begin
\needspace{3\baselineskip}
\textbf{may} — Mōtan expresses permission, may/be allowed to, with a plain infinitive. A wish beginning with English may instead uses the taught independent subjunctive.\par
% glossary-block:B01126:end
% glossary-block:B01127:begin
\needspace{3\baselineskip}
\textbf{may (permission)} — \textbf{mōtan} (v.). mōtan + plain infinitive; permission/prohibition differs from ability and independent wishes.\par
% glossary-block:B01127:end
% glossary-block:B01128:begin
\needspace{3\baselineskip}
\textbf{may (wish)} — An independent wish/hortatory uses a subjunctive; let does not automatically require lǣtan. Literal permission to go uses lǣtan + accusative actor + infinitive.\par
% glossary-block:B01128:end
% glossary-block:B01129:begin
\needspace{3\baselineskip}
\textbf{me} — Mē is first-person singular dative or accusative me; choose its case from the verb or preposition. Select case from the verb/preposition and preserve dual versus plural; dual subjects use plural verb forms.\par
% glossary-block:B01129:end
% glossary-block:B01130:begin
\needspace{3\baselineskip}
\textbf{meal} — \textbf{gereord} (n.). meal; banquet.\par
% glossary-block:B01130:end
% glossary-block:B01131:begin
\needspace{3\baselineskip}
\textbf{meaning} — \textbf{andgiet} (n.). meaning; sense.\par
% glossary-block:B01131:end
% glossary-block:B01132:begin
\needspace{3\baselineskip}
\textbf{means of access} — \textbf{weg} (n.). way; means of access. Figurative access to the Father in the declaration; distinct from physical road.\par
% glossary-block:B01132:end
% glossary-block:B01133:begin
\needspace{3\baselineskip}
\textbf{measure} — \textbf{mǣð} (n.). measure; capacity.\par
% glossary-block:B01133:end
% glossary-block:B01134:begin
\needspace{3\baselineskip}
\textbf{meet} — \textbf{gemētan} (v.). find; meet. Accusative person or thing found/met. \textbf{mētan} (v.). meet; find.\par
% glossary-block:B01134:end
% glossary-block:B01135:begin
\needspace{3\baselineskip}
\textbf{mention before} — \textbf{foresęcgean} (v.). mention before; tell of beforehand.\par
% glossary-block:B01135:end
% glossary-block:B01136:begin
\needspace{3\baselineskip}
\textbf{merchandise} — \textbf{hlæst} (n.). wares; merchandise; cargo. Selected plural source use.\par
% glossary-block:B01136:end
% glossary-block:B01137:begin
\needspace{3\baselineskip}
\textbf{merchant} — \textbf{mangere} (n.).\par
% glossary-block:B01137:end
% glossary-block:B01138:begin
\needspace{3\baselineskip}
\textbf{mercies} — \textbf{mildheortnes} (n.). mercy; mercies.\par
% glossary-block:B01138:end
% glossary-block:B01139:begin
\needspace{3\baselineskip}
\textbf{merciful} — \textbf{mildheort} (adj.). merciful; compassionate.\par
% glossary-block:B01139:end
% glossary-block:B01140:begin
\needspace{3\baselineskip}
\textbf{mercy} — \textbf{mildheortnes} (n.). mercy; mercies.\par
% glossary-block:B01140:end
% glossary-block:B01141:begin
\needspace{3\baselineskip}
\textbf{merit} — \textbf{earnung} (n.). merit; good service. Source plural merits/services.\par
% glossary-block:B01141:end
% glossary-block:B01142:begin
\needspace{3\baselineskip}
\textbf{merry} — \textbf{blīðe} (adj.). merry; glad.\par
% glossary-block:B01142:end
% glossary-block:B01143:begin
\needspace{3\baselineskip}
\textbf{messengers} — \textbf{ǣrendwreca} (n.). messengers; envoys.\par
% glossary-block:B01143:end
% glossary-block:B01144:begin
\needspace{3\baselineskip}
\textbf{metal} — \textbf{węcg} (n.). metal; ingot. Supplied genitive plural dependent on veins.\par
% glossary-block:B01144:end
% glossary-block:B01145:begin
\needspace{3\baselineskip}
\textbf{Michael} — \textbf{Michael} (proper name).\par
% glossary-block:B01145:end
% glossary-block:B01146:begin
\needspace{3\baselineskip}
\textbf{midday} — The supplied whole phrase tō middes dæges means at noon/midday; middes and dæges are genitives.\par
% glossary-block:B01146:end
% glossary-block:B01147:begin
\needspace{3\baselineskip}
\textbf{might} — The supplied mugan form expresses capability, can/could/might be able to, with a plain infinitive. It is not the permission sense of English may/might. In the exact supplied purpose construction, willan accompanies the watching/seeing infinitives and is rendered might. This does not make willan a general equivalent of English might or pure future will.\par
% glossary-block:B01147:end
% glossary-block:B01148:begin
\needspace{3\baselineskip}
\textbf{mighty} — \textbf{mihtig} (adj.). mighty; powerful.\par
% glossary-block:B01148:end
% glossary-block:B01149:begin
\needspace{3\baselineskip}
\textbf{mind} — \textbf{geðanc} (n.). thought; mind. \textbf{mōd} (n.). mind; attention. heart; mind.\par
% glossary-block:B01149:end
% glossary-block:B01150:begin
\needspace{3\baselineskip}
\textbf{mine} — \textbf{mīn} (possessive adjective). my; mine. Mīn my/mine: one first-person possessor; the possessive agrees with the possessed noun.\par
% glossary-block:B01150:end
% glossary-block:B01151:begin
\needspace{3\baselineskip}
\textbf{mingle} — \textbf{gemęngan} (v.). join; mingle. Supplied source form. mix. Mix sound with song in the supplied performance expression.\par
% glossary-block:B01151:end
% glossary-block:B01152:begin
\needspace{3\baselineskip}
\textbf{miserably} — \textbf{earmlice} (adverb or particle). cruelly; miserably.\par
% glossary-block:B01152:end
% glossary-block:B01153:begin
\needspace{3\baselineskip}
\textbf{mix} — \textbf{gemęngan} (v.). Mix sound with song in the supplied performance expression.\par
% glossary-block:B01153:end
% glossary-block:B01154:begin
\needspace{3\baselineskip}
\textbf{mixture of spices} — \textbf{wyrtgemang} (n.). spice; mixture of spices.\par
% glossary-block:B01154:end
% glossary-block:B01155:begin
\needspace{3\baselineskip}
\textbf{monastery} — \textbf{mynster} (n.).\par
% glossary-block:B01155:end
% glossary-block:B01156:begin
\needspace{3\baselineskip}
\textbf{money} — \textbf{feoh} (n.).\par
% glossary-block:B01156:end
% glossary-block:B01157:begin
\needspace{3\baselineskip}
\textbf{monk} — \textbf{munuc} (n.).\par
% glossary-block:B01157:end
% glossary-block:B01158:begin
\needspace{3\baselineskip}
\textbf{month} — \textbf{mōnað} (n.). Source plural months.\par
% glossary-block:B01158:end
% glossary-block:B01159:begin
\needspace{3\baselineskip}
\textbf{moon} — \textbf{mōna} (n.).\par
% glossary-block:B01159:end
% glossary-block:B01160:begin
\needspace{3\baselineskip}
\textbf{more} — \textbf{mā} (adv.). Adverb of amount. more; again. In the negative phrase never more/never again. More in number, distinct from degree or greater value. \textbf{māra} (adj.). greater; more important. Value/importance in the life/raiment comparison. more; a greater amount. Nominal/attributive quantity, distinct from more in number. Comparative swīðor expresses rather/more in the preference construction; the positive adverb swīðe means greatly/very. Comparative swīðor expresses more in the degree of an action or feeling, not a count of objects. Mā expresses more in amount or number in its separate uses; māra expresses a greater amount. Comparative swīðor expresses more intensely or preferential rather. Use an agreeing comparative adjective for greater quality/value; these uses are not interchangeable.\par
% glossary-block:B01160:end
% glossary-block:B01161:begin
\needspace{3\baselineskip}
\textbf{more beautiful} — \textbf{fægerra} (adj.). more beautiful; fairer.\par
% glossary-block:B01161:end
% glossary-block:B01162:begin
\needspace{3\baselineskip}
\textbf{more correctly} — Supplied comparative rihtlicor means more correctly. Retain its source spelling without an inferred length mark.\par
% glossary-block:B01162:end
% glossary-block:B01163:begin
\needspace{3\baselineskip}
\textbf{more important} — \textbf{māra} (adj.). greater; more important. Value/importance in the life/raiment comparison.\par
% glossary-block:B01163:end
% glossary-block:B01164:begin
\needspace{3\baselineskip}
\textbf{more precious} — \textbf{dēorwierðra} (adj.). Comparative adjective; distinct from an adverb modifying an action.\par
% glossary-block:B01164:end
% glossary-block:B01165:begin
\needspace{3\baselineskip}
\textbf{more severely} — Supplied comparative swīðlicor means more violently/more severely; it is not the positive swīðlice.\par
% glossary-block:B01165:end
% glossary-block:B01166:begin
\needspace{3\baselineskip}
\textbf{more than} — \textbf{ofer} (prep.). above; more than. Preference in the supplied comparison.\par
% glossary-block:B01166:end
% glossary-block:B01167:begin
\needspace{3\baselineskip}
\textbf{more violently} — Supplied comparative swīðlicor means more violently/more severely; it is not the positive swīðlice.\par
% glossary-block:B01167:end
% glossary-block:B01168:begin
\needspace{3\baselineskip}
\textbf{more widely} — Wīdor is the comparative adverb more widely; wīde is the positive widely.\par
% glossary-block:B01168:end
% glossary-block:B01169:begin
\needspace{3\baselineskip}
\textbf{morning} — \textbf{morgen} (n.).\par
% glossary-block:B01169:end
% glossary-block:B01170:begin
\needspace{3\baselineskip}
\textbf{most} — Swīðost means most/very greatly. It is the recognition-only superlative reference form of swīðe.\par
% glossary-block:B01170:end
% glossary-block:B01171:begin
\needspace{3\baselineskip}
\textbf{most beautiful} — \textbf{fægerest} (adj.).\par
% glossary-block:B01171:end
% glossary-block:B01172:begin
\needspace{3\baselineskip}
\textbf{most clearly} — Supplied superlative sweotolost means most clearly; do not use the positive sweotole for this degree.\par
% glossary-block:B01172:end
% glossary-block:B01173:begin
\needspace{3\baselineskip}
\textbf{most intelligibly} — Supplied superlative andgietfullicost means most intelligibly; do not use the positive andgietfullice for this degree.\par
% glossary-block:B01173:end
% glossary-block:B01174:begin
\needspace{3\baselineskip}
\textbf{most widely} — Wīdost is the superlative adverb most widely; wīde is the positive widely.\par
% glossary-block:B01174:end
% glossary-block:B01175:begin
\needspace{3\baselineskip}
\textbf{mother} — \textbf{mōdor} (n.).\par
% glossary-block:B01175:end
% glossary-block:B01176:begin
\needspace{3\baselineskip}
\textbf{mountain} — \textbf{dūn} (n.). mountain; hill. \textbf{munt} (n.).\par
% glossary-block:B01176:end
% glossary-block:B01177:begin
\needspace{3\baselineskip}
\textbf{mouth} — \textbf{mūð} (n.).\par
% glossary-block:B01177:end
% glossary-block:B01178:begin
\needspace{3\baselineskip}
\textbf{move} — \textbf{astyrian} (v.). move; set in motion. Supplied source sense.\par
% glossary-block:B01178:end
% glossary-block:B01179:begin
\needspace{3\baselineskip}
\textbf{moving} — \textbf{styriendlic} (adj.).\par
% glossary-block:B01179:end
% glossary-block:B01180:begin
\needspace{3\baselineskip}
\textbf{much} — \textbf{micel} (adj.). much; an amount. Quantity sense. Swīðe expresses much as degree, as in love very much; swā swīðe expresses so much. This differs from micel expressing quantity.\par
% glossary-block:B01180:end
% glossary-block:B01181:begin
\needspace{3\baselineskip}
\textbf{multiply} — \textbf{gemanigfieldan} (v.). be multiplied. Assigned passive with agreeing gemanigfield/gemanigfielde; no independent active causative frame is assigned. Assigned passive use with wesan and agreeing gemanigfield/gemanigfielde; do not infer an unassigned active causative frame.\par
% glossary-block:B01181:end
% glossary-block:B01182:begin
\needspace{3\baselineskip}
\textbf{multitude} — \textbf{getæl} (n.). Singular collective getæl with the supplied dependent expression; retain source participial genitives. \textbf{męnigu} (n.). męnigu also denotes a multitude of things in the locally supplied spice expression.\par
% glossary-block:B01182:end
% glossary-block:B01183:begin
\needspace{3\baselineskip}
\textbf{music} — \textbf{swēg} (n.). music; sound.\par
% glossary-block:B01183:end
% glossary-block:B01184:begin
\needspace{3\baselineskip}
\textbf{musical skill} — \textbf{swēgcræft} (n.).\par
% glossary-block:B01184:end
% glossary-block:B01185:begin
\needspace{3\baselineskip}
\textbf{must} — Sculan expresses obligation, must/be obliged to, with a plain infinitive. Pure future reference does not automatically require sculan. \textbf{sculan} (v.). sculan + plain infinitive for obligation; English have to is not possessive habban. Supplied have something to offer in Unit 6 is a different construction.\par
% glossary-block:B01185:end
% glossary-block:B01186:begin
\needspace{3\baselineskip}
\textbf{my} — \textbf{mīn} (possessive adjective). my; mine. Mīn my/mine: one first-person possessor; the possessive agrees with the possessed noun.\par
% glossary-block:B01186:end
% glossary-block:B01187:begin
\needspace{3\baselineskip}
\textbf{myself} — An ordinary reflexive uses the personal pronoun in the case governed by its verb or preposition. Agreeing self adds the taught emphasis; English -self does not automatically require a separate self. Ordinary reflexive personal pronouns take the governed case; agreeing self adds the taught emphasis. English -self does not automatically require self.\par
% glossary-block:B01187:end
% book-navigation: standalone bookmark replaced by volume hierarchy
% glossary-block:B01188:begin
\needspace{5\baselineskip}
\subsection*{N}\BookHeading{nav-glossary-english-0015}{subsubsection}{N}{0}
% glossary-block:B01188:end
% glossary-block:B01189:begin
\textbf{naked} — \textbf{nacod} (adj.).\par
% glossary-block:B01189:end
% glossary-block:B01190:begin
\needspace{3\baselineskip}
\textbf{naked people} — \textbf{nacod} (adj.). Substantival adjective.\par
% glossary-block:B01190:end
% glossary-block:B01191:begin
\needspace{3\baselineskip}
\textbf{name} — \textbf{hātan} (v.). call; name. Active hātan call/name; the supplied naming expression may have an accusative object and accusative name. \textbf{nama} (n.).\par
% glossary-block:B01191:end
% glossary-block:B01192:begin
\needspace{3\baselineskip}
\textbf{narrow} — \textbf{nearu} (adj.).\par
% glossary-block:B01192:end
% glossary-block:B01193:begin
\needspace{3\baselineskip}
\textbf{nation} — \textbf{mǣgð} (n.). tribe; nation. \textbf{ðēod} (n.). people; nation.\par
% glossary-block:B01193:end
% glossary-block:B01194:begin
\needspace{3\baselineskip}
\textbf{nations} — \textbf{folc} (n.). peoples. Distinct peoples or national groups; plural sense.\par
% glossary-block:B01194:end
% glossary-block:B01195:begin
\needspace{3\baselineskip}
\textbf{native land} — \textbf{ēðel} (n.). native land; homeland. ēðel in these home-country references; land remains a separate ordinary country word.\par
% glossary-block:B01195:end
% glossary-block:B01196:begin
\needspace{3\baselineskip}
\textbf{naught} — \textbf{nāht} (substantive). nothing; naught. \textbf{nāwiht} (indefinite pronoun/substantive). nothing; naught.\par
% glossary-block:B01196:end
% glossary-block:B01197:begin
\needspace{3\baselineskip}
\textbf{naval force} — \textbf{sciphęre} (n.). fleet; naval force.\par
% glossary-block:B01197:end
% glossary-block:B01198:begin
\needspace{3\baselineskip}
\textbf{near} — \textbf{nēah} (adv.). near; at hand. nēah is the predicate in nigh at hand.\par
% glossary-block:B01198:end
% glossary-block:B01199:begin
\needspace{3\baselineskip}
\textbf{necessarily} — \textbf{nīede} (adverb or particle). necessarily; of necessity. Archaic English needs is an adverb here; do not use the verb need.\par
% glossary-block:B01199:end
% glossary-block:B01200:begin
\needspace{3\baselineskip}
\textbf{need} — \textbf{beðurfan} (v.). need; require. Genitive thing needed. Genitive thing needed; a relative particle may represent that governed relation. \textbf{ðearf} (n.). benefit; good; need. Source contextual nouns; retain the local expression.\par
% glossary-block:B01200:end
% glossary-block:B01201:begin
\needspace{3\baselineskip}
\textbf{needle} — \textbf{nǣdl} (n.).\par
% glossary-block:B01201:end
% glossary-block:B01202:begin
\needspace{3\baselineskip}
\textbf{needs} — \textbf{nīede} (adverb or particle). Archaic English needs is an adverb here; do not use the verb need.\par
% glossary-block:B01202:end
% glossary-block:B01203:begin
\needspace{3\baselineskip}
\textbf{need to} — \textbf{ðurfan} (v.). ðurfan + plain infinitive; distinguish beðurfan + genitive thing needed.\par
% glossary-block:B01203:end
% glossary-block:B01204:begin
\needspace{3\baselineskip}
\textbf{needy person} — \textbf{ðearfa} (n.). poor person; needy person. Noun referring to a person.\par
% glossary-block:B01204:end
% glossary-block:B01205:begin
\needspace{3\baselineskip}
\textbf{neglect} — \textbf{rēcelēasian} (v.). neglect; despise.\par
% glossary-block:B01205:end
% glossary-block:B01206:begin
\needspace{3\baselineskip}
\textbf{neighbor} — \textbf{nēahgebūr} (n.).\par
% glossary-block:B01206:end
% glossary-block:B01207:begin
\needspace{3\baselineskip}
\textbf{neighbour} — \textbf{nēahgebūr} (n.).\par
% glossary-block:B01207:end
% glossary-block:B01208:begin
\needspace{3\baselineskip}
\textbf{neither} — \textbf{nāhwæðer} (adverb or particle).\par
% glossary-block:B01208:end
% glossary-block:B01209:begin
\needspace{3\baselineskip}
\textbf{neither … nor} — \textbf{nāhwæðer} (adverb or particle). Whole supplied nāhwæðer ne … ne with finite negative scope; neither excludes both alternatives.\par
% glossary-block:B01209:end
% glossary-block:B01210:begin
\needspace{3\baselineskip}
\textbf{nest} — \textbf{nest} (n.).\par
% glossary-block:B01210:end
% glossary-block:B01211:begin
\needspace{3\baselineskip}
\textbf{net} — \textbf{nętt} (n.).\par
% glossary-block:B01211:end
% glossary-block:B01212:begin
\needspace{3\baselineskip}
\textbf{never} — \textbf{nǣfre} (adv.). Nǣfre means never; retain finite ne or the licensed negative contraction in negative concord.\par
% glossary-block:B01212:end
% glossary-block:B01213:begin
\needspace{3\baselineskip}
\textbf{never more} — Mā means more/again within the negative expression never more/never again; the clause must supply the negation.\par
% glossary-block:B01213:end
% glossary-block:B01214:begin
\needspace{3\baselineskip}
\textbf{nevertheless} — \textbf{ðā-hwæðre} (conj.). nevertheless; yet. Whole supplied connective. \textbf{ðēah} (conj.; adv.). nevertheless; yet. Independent adverb, distinct from concessive conjunction although.\par
% glossary-block:B01214:end
% glossary-block:B01215:begin
\needspace{3\baselineskip}
\textbf{new} — \textbf{nīwe} (adj.).\par
% glossary-block:B01215:end
% glossary-block:B01216:begin
\needspace{3\baselineskip}
\textbf{nigh} — \textbf{nēah} (adv.). nēah is the predicate in nigh at hand.\par
% glossary-block:B01216:end
% glossary-block:B01217:begin
\needspace{3\baselineskip}
\textbf{night} — \textbf{niht} (n.).\par
% glossary-block:B01217:end
% glossary-block:B01218:begin
\needspace{3\baselineskip}
\textbf{nine} — \textbf{nigon} (cardinal numeral). Cardinal 9.\par
% glossary-block:B01218:end
% glossary-block:B01219:begin
\needspace{3\baselineskip}
\textbf{ninety} — \textbf{hundnigontig} (numeral).\par
% glossary-block:B01219:end
% glossary-block:B01220:begin
\needspace{3\baselineskip}
\textbf{ninth} — \textbf{nigoða} (ordinal).\par
% glossary-block:B01220:end
% glossary-block:B01221:begin
\needspace{3\baselineskip}
\textbf{no} — \textbf{nān} (indefinite). no; no one.\par
% glossary-block:B01221:end
% glossary-block:B01222:begin
\needspace{3\baselineskip}
\textbf{noble} — \textbf{æðele} (adj.). noble; illustrious. \textbf{ealdormann} (n.). prince; leader; noble.\par
% glossary-block:B01222:end
% glossary-block:B01223:begin
\needspace{3\baselineskip}
\textbf{noble people} — \textbf{æðele} (adj.). Substantival adjective in the source.\par
% glossary-block:B01223:end
% glossary-block:B01224:begin
\needspace{3\baselineskip}
\textbf{nobody} — \textbf{nān} (indefinite). Substantival nān means no one; retain finite ne or the licensed negative contraction in negative concord.\par
% glossary-block:B01224:end
% glossary-block:B01225:begin
\needspace{3\baselineskip}
\textbf{noise} — \textbf{gebrec} (n.). crash; noise; din. Supplied accusative singular.\par
% glossary-block:B01225:end
% glossary-block:B01226:begin
\needspace{3\baselineskip}
\textbf{no longer} — \textbf{lęng} (adverb or particle). lęng as the comparative adverb; clause negation supplies no longer.\par
% glossary-block:B01226:end
% glossary-block:B01227:begin
\needspace{3\baselineskip}
\textbf{no man} — \textbf{nān} (indefinite). Substantival nān can render no man/no one where the task omits a separate man; retain negative concord.\par
% glossary-block:B01227:end
% glossary-block:B01228:begin
\needspace{3\baselineskip}
\textbf{noon} — The supplied whole phrase tō middes dæges means at noon/midday; middes and dæges are genitives. \textbf{middes} (adverb or particle). Whole supplied phrase tō middes dæges; preserve its genitives instead of applying ordinary tō + dative mechanically.\par
% glossary-block:B01228:end
% glossary-block:B01229:begin
\needspace{3\baselineskip}
\textbf{no one} — \textbf{nān} (indefinite). no; no one. Substantival nān means no one; retain finite ne or the licensed negative contraction in negative concord.\par
% glossary-block:B01229:end
% glossary-block:B01230:begin
\needspace{3\baselineskip}
\textbf{northeasterly} — \textbf{ēastnorðerne} (adj.).\par
% glossary-block:B01230:end
% glossary-block:B01231:begin
\needspace{3\baselineskip}
\textbf{not} — \textbf{nā} (adv.). Negates a contrasted constituent, distinct from finite ne. Ne is the clause negative not, placed with the finite verb in the taught construction. English do not/did not requires no extra Old English auxiliary do.\par
% glossary-block:B01231:end
% glossary-block:B01232:begin
\needspace{3\baselineskip}
\textbf{not at all} — \textbf{nālæs} (adv.).\par
% glossary-block:B01232:end
% glossary-block:B01233:begin
\needspace{3\baselineskip}
\textbf{not be} — \textbf{wesan} (v.). Supplied negative nǣre in reported supposition.\par
% glossary-block:B01233:end
% glossary-block:B01234:begin
\needspace{3\baselineskip}
\textbf{not have} — \textbf{nāgan} (v.). not have; not possess. Negative possessive verb in its supplied forms.\par
% glossary-block:B01234:end
% glossary-block:B01235:begin
\needspace{3\baselineskip}
\textbf{nothing} — \textbf{nāht} (substantive). nothing; naught. \textbf{nāwiht} (indefinite pronoun/substantive). nothing; naught. Nāwiht means nothing in the bounded accusative construction; retain finite ne or the licensed negative contraction.\par
% glossary-block:B01235:end
% glossary-block:B01236:begin
\needspace{3\baselineskip}
\textbf{not know} — \textbf{nytan} (v.). Supplied negative knowing form.\par
% glossary-block:B01236:end
% glossary-block:B01237:begin
\needspace{3\baselineskip}
\textbf{not possess} — \textbf{nāgan} (v.). not have; not possess. Negative possessive verb in its supplied forms.\par
% glossary-block:B01237:end
% glossary-block:B01238:begin
\needspace{3\baselineskip}
\textbf{not working} — \textbf{īdel} (adj.). idle. Distinct from empty/void.\par
% glossary-block:B01238:end
% glossary-block:B01239:begin
\needspace{3\baselineskip}
\textbf{nourish} — \textbf{āfēdan} (v.). support; nourish. Accusative person supported. \textbf{fēdan} (v.). feed; nourish. Accusative being fed; food may be mid + dative or a bare instrumental in the assigned expressions.\par
% glossary-block:B01239:end
% glossary-block:B01240:begin
\needspace{3\baselineskip}
\textbf{now} — \textbf{nū} (adv.; conj.).  \textbf{sōðlice} (adv.). now; but. Discourse connective corresponding to source Latin autem, distinct from truly. \textbf{witodlice} (adv.). now; and. Discourse connective.\par
% glossary-block:B01240:end
% glossary-block:B01241:begin
\needspace{3\baselineskip}
\textbf{number} — \textbf{gerīman} (v.). number; count. Exact supplied infinitive gerīman; not an independently assigned full paradigm. Supplied infinitive gerīman, count; number. \textbf{getæl} (n.). number; company. Singular collective getæl with the supplied dependent expression; retain source participial genitives.\par
% glossary-block:B01241:end
% book-navigation: standalone bookmark replaced by volume hierarchy
% glossary-block:B01242:begin
\needspace{5\baselineskip}
\subsection*{O}\BookHeading{nav-glossary-english-0016}{subsubsection}{O}{0}
% glossary-block:B01242:end
% glossary-block:B01243:begin
\textbf{O} — \textbf{ēalā} (interj.). O!. Interjection in address, exclamation or wish.\par
% glossary-block:B01243:end
% glossary-block:B01244:begin
\needspace{3\baselineskip}
\textbf{oath} — \textbf{āð} (n.).\par
% glossary-block:B01244:end
% glossary-block:B01245:begin
\needspace{3\baselineskip}
\textbf{obedience} — \textbf{gehīersumnes} (n.).  \textbf{hīernes} (n.).\par
% glossary-block:B01245:end
% glossary-block:B01246:begin
\needspace{3\baselineskip}
\textbf{obey} — \textbf{gehīersumian} (v.). serve; obey. Supplied source form and government. \textbf{hīeran} (v.). Dative person obeyed. Dative person obeyed; distinct from hearing a sound/object.\par
% glossary-block:B01246:end
% glossary-block:B01247:begin
\needspace{3\baselineskip}
\textbf{observe} — \textbf{behealdan} (v.). Supplied source form observed. \textbf{healdan} (v.). keep; observe. Accusative saying or command observed. \textbf{scēawian} (v.). watch; observe. Accusative person or event watched, including a journey.\par
% glossary-block:B01247:end
% glossary-block:B01248:begin
\needspace{3\baselineskip}
\textbf{obtain} — \textbf{tilian} (v.). obtain; provide. Genitive thing procured; dative reflexive beneficiary in the assigned frame.\par
% glossary-block:B01248:end
% glossary-block:B01249:begin
\needspace{3\baselineskip}
\textbf{obtain for oneself} — \textbf{tilian} (v.). tilian + dative reflexive beneficiary + genitive thing obtained.\par
% glossary-block:B01249:end
% glossary-block:B01250:begin
\needspace{3\baselineskip}
\textbf{occasion} — \textbf{sīð} (n.). time; occasion.\par
% glossary-block:B01250:end
% glossary-block:B01251:begin
\needspace{3\baselineskip}
\textbf{occupation} — \textbf{cræft} (n.). trade; occupation; craft.\par
% glossary-block:B01251:end
% glossary-block:B01252:begin
\needspace{3\baselineskip}
\textbf{occupations} — \textbf{bisgu} (n.). occupations; cares.\par
% glossary-block:B01252:end
% glossary-block:B01253:begin
\needspace{3\baselineskip}
\textbf{occupied} — \textbf{geornfull} (adj.). occupied; intent. Source mental-attention context.\par
% glossary-block:B01253:end
% glossary-block:B01254:begin
\needspace{3\baselineskip}
\textbf{occupy} — \textbf{bisgian} (v.). occupy; trouble.\par
% glossary-block:B01254:end
% glossary-block:B01255:begin
\needspace{3\baselineskip}
\textbf{occur} — \textbf{becuman} (v.). occur; come about. Of the calm in the source.\par
% glossary-block:B01255:end
% glossary-block:B01256:begin
\needspace{3\baselineskip}
\textbf{of} — \textbf{of} (prep.). of; from. Portion from a whole in the supplied partitive phrases; dative whole. of; from. Origin of wisdom in the supplied begetting statement. Possessive/relational English of or 's often becomes a dependent genitive, not preposition of. Partitive, source and lexical government require their own frames.\par
% glossary-block:B01256:end
% glossary-block:B01257:begin
\needspace{3\baselineskip}
\textbf{of Cyrene} — \textbf{Cyrenisc} (adj.). Cyrenian; of Cyrene. Geographical adjective.\par
% glossary-block:B01257:end
% glossary-block:B01258:begin
\needspace{3\baselineskip}
\textbf{off} — \textbf{of} (adv.).\par
% glossary-block:B01258:end
% glossary-block:B01259:begin
\needspace{3\baselineskip}
\textbf{offend} — \textbf{ǣswician} (v.). offend; cause to stumble. Accusative person; ethical/metaphorical offence in the assigned use.\par
% glossary-block:B01259:end
% glossary-block:B01260:begin
\needspace{3\baselineskip}
\textbf{offer} — \textbf{bēodan} (v.).\par
% glossary-block:B01260:end
% glossary-block:B01261:begin
\needspace{3\baselineskip}
\textbf{offspring} — \textbf{sǣd} (n.). seed; descendants. Collective descendants, distinct from literal seed.\par
% glossary-block:B01261:end
% glossary-block:B01262:begin
\needspace{3\baselineskip}
\textbf{of me} — Mīn is the first-person singular personal genitive, of me. Do not add the preposition of; attributive my is a separately taught agreeing possessive.\par
% glossary-block:B01262:end
% glossary-block:B01263:begin
\needspace{3\baselineskip}
\textbf{of necessity} — \textbf{nīede} (adverb or particle). necessarily; of necessity. Archaic English needs is an adverb here; do not use the verb need.\par
% glossary-block:B01263:end
% glossary-block:B01264:begin
\needspace{3\baselineskip}
\textbf{often} — \textbf{genēahhe} (adv.). often; frequently. \textbf{oft} (adv.).\par
% glossary-block:B01264:end
% glossary-block:B01265:begin
\needspace{3\baselineskip}
\textbf{of the sea} — \textbf{sǣlic} (adj.). marine; of the sea.\par
% glossary-block:B01265:end
% glossary-block:B01266:begin
\needspace{3\baselineskip}
\textbf{of us} — Uncer is the personal genitive of us two; ūre is plural of us. Do not add the preposition of; attributive our uses separately taught agreeing possessives.\par
% glossary-block:B01266:end
% glossary-block:B01267:begin
\needspace{3\baselineskip}
\textbf{of us two} — \textbf{ic} (personal pron.). Personal genitive dual after a selection word.\par
% glossary-block:B01267:end
% glossary-block:B01268:begin
\needspace{3\baselineskip}
\textbf{of what sort} — \textbf{hūlic} (interrogative adjective). of what sort?.\par
% glossary-block:B01268:end
% glossary-block:B01269:begin
\needspace{3\baselineskip}
\textbf{of winter} — \textbf{winterlic} (adj.). winter; of winter.\par
% glossary-block:B01269:end
% glossary-block:B01270:begin
\needspace{3\baselineskip}
\textbf{of you} — Personal genitive ðīn refers to one addressee, incer to two, ēower to plural addressees. English of is expressed by the genitive, not an added preposition; attributive your is separately taught.\par
% glossary-block:B01270:end
% glossary-block:B01271:begin
\needspace{3\baselineskip}
\textbf{oil} — \textbf{ęle} (n.).\par
% glossary-block:B01271:end
% glossary-block:B01272:begin
\needspace{3\baselineskip}
\textbf{ointment} — \textbf{sealf} (n.).\par
% glossary-block:B01272:end
% glossary-block:B01273:begin
\needspace{3\baselineskip}
\textbf{old} — \textbf{eald} (adj.).\par
% glossary-block:B01273:end
% glossary-block:B01274:begin
\needspace{3\baselineskip}
\textbf{old age} — \textbf{ieldu} (n.). age; old age.\par
% glossary-block:B01274:end
% glossary-block:B01275:begin
\needspace{3\baselineskip}
\textbf{older} — \textbf{ieldra} (adj.).\par
% glossary-block:B01275:end
% glossary-block:B01276:begin
\needspace{3\baselineskip}
\textbf{oldest} — \textbf{ieldest} (adj.). eldest; oldest. Substantival person/eldest brother in this source context.\par
% glossary-block:B01276:end
% glossary-block:B01277:begin
\needspace{3\baselineskip}
\textbf{old people} — \textbf{eald} (adj.). old people; the elderly. Substantival adjective.\par
% glossary-block:B01277:end
% glossary-block:B01278:begin
\needspace{3\baselineskip}
\textbf{on} — \textbf{in} (prep.). in; on. With dative in the exact locations/circumstances. \textbf{on} (prep.). in; on. Dative for the taught ordinary location construction. Time preposition in the stated-day expression. \textbf{on} (adv.). Adverb in the cloak-wearing relative clause. \textbf{upon} (prep.). on; onto.\par
% glossary-block:B01278:end
% glossary-block:B01279:begin
\needspace{3\baselineskip}
\textbf{on account of} — \textbf{for} (prep.). because of; on account of. for + dative cause in this supplied money expression; distinguish causal conjunction because.\par
% glossary-block:B01279:end
% glossary-block:B01280:begin
\needspace{3\baselineskip}
\textbf{on behalf of} — \textbf{for} (prep.). for; on behalf of. Dative beneficiary.\par
% glossary-block:B01280:end
% glossary-block:B01281:begin
\needspace{3\baselineskip}
\textbf{once more} — \textbf{eft, ęft} (adv.). eft is the repeated-occurrence adverb.\par
% glossary-block:B01281:end
% glossary-block:B01282:begin
\needspace{3\baselineskip}
\textbf{one} — \textbf{ān} (numeral; adj.). one; single. Cardinal one/single and the agreeing adjective alone are distinct uses of ān. Preserve the supplied form where the exercise provides one. \textbf{mann} (n.; impersonal pronoun). one; people in general. Impersonal man with singular finite verb.\par
% glossary-block:B01282:end
% glossary-block:B01283:begin
\needspace{3\baselineskip}
\textbf{one of} — The supplied partitive expression combines the required word for one with of + dative whole. Of alone means of/from, not one of.\par
% glossary-block:B01283:end
% glossary-block:B01284:begin
\needspace{3\baselineskip}
\textbf{one (people in general)} — Impersonal man has a singular finite verb; distinguish numeral ān and one selected member of a group.\par
% glossary-block:B01284:end
% glossary-block:B01285:begin
\needspace{3\baselineskip}
\textbf{one who delivers} — \textbf{ālīesend} (n.). Redeemer; one who delivers.\par
% glossary-block:B01285:end
% glossary-block:B01286:begin
\needspace{3\baselineskip}
\textbf{one who shows mercy} — \textbf{gemiltsiend} (n.). one who shows mercy; pitier.\par
% glossary-block:B01286:end
% glossary-block:B01287:begin
\needspace{3\baselineskip}
\textbf{only-begotten} — \textbf{āncęnned} (adj.).\par
% glossary-block:B01287:end
% glossary-block:B01288:begin
\needspace{3\baselineskip}
\textbf{on one's side} — The bounded on … hand phrase means on someone’s side; hand is accusative even with standan. Retain the supplied possessor and side expression.\par
% glossary-block:B01288:end
% glossary-block:B01289:begin
\needspace{3\baselineskip}
\textbf{on that evidence} — \textbf{be} (prep.). by that; from that. Refers to the preceding proposition; supplied expression. by that; on that evidence. Evidence/inference frame, not an unrestricted source preposition.\par
% glossary-block:B01289:end
% glossary-block:B01290:begin
\needspace{3\baselineskip}
\textbf{on the side of} — The bounded on … hand phrase means on someone’s side; hand is accusative even with standan. Retain the supplied possessor and side expression.\par
% glossary-block:B01290:end
% glossary-block:B01291:begin
\needspace{3\baselineskip}
\textbf{on the way home} — \textbf{hāmweard} (adverb or particle). Supplied hāmweard with wesan; distinguish noun hām and directional adverb hām.\par
% glossary-block:B01291:end
% glossary-block:B01292:begin
\needspace{3\baselineskip}
\textbf{onto} — \textbf{in} (prep.). toward; onto. Supplied phrase directing the mind toward these things. \textbf{on} (prep.). into; onto. Accusative goal of placement in the taught frame. \textbf{upon} (prep.). on; onto.\par
% glossary-block:B01292:end
% glossary-block:B01293:begin
\needspace{3\baselineskip}
\textbf{on top of} — \textbf{uppan} (prep.). upon; on top of.\par
% glossary-block:B01293:end
% glossary-block:B01294:begin
\needspace{3\baselineskip}
\textbf{open} — \textbf{geopenian} (v.). Accusative thing opened, including mouth or eyes. \textbf{ontȳnan} (v.). Supplied preterite opened. \textbf{open} (adj.).  \textbf{undōn} (v.). undo; open.\par
% glossary-block:B01294:end
% glossary-block:B01295:begin
\needspace{3\baselineskip}
\textbf{opponent} — \textbf{gewinna} (n.). enemy; opponent. Supplied plural source forms.\par
% glossary-block:B01295:end
% glossary-block:B01296:begin
\needspace{3\baselineskip}
\textbf{oppose} — \textbf{standan} (v.). oppose; stand against. Supplied standan ongēan frame, distinct from simple stand.\par
% glossary-block:B01296:end
% glossary-block:B01297:begin
\needspace{3\baselineskip}
\textbf{opposed to} — \textbf{wið} (prep.). against; opposed to. Accusative opponent.\par
% glossary-block:B01297:end
% glossary-block:B01298:begin
\needspace{3\baselineskip}
\textbf{opposite} — \textbf{ongēan} (adv.; prep.; conj.). against; opposite. Accusative landmark or opponent.\par
% glossary-block:B01298:end
% glossary-block:B01299:begin
\needspace{3\baselineskip}
\textbf{oppress} — \textbf{hīenan} (v.). oppress; humble.\par
% glossary-block:B01299:end
% glossary-block:B01300:begin
\needspace{3\baselineskip}
\textbf{or} — \textbf{oððe} (conj.).\par
% glossary-block:B01300:end
% glossary-block:B01301:begin
\needspace{3\baselineskip}
\textbf{oratory} — \textbf{cirice} (n.). Chad's separate place of prayer; distinct from the ordinary church sense.\par
% glossary-block:B01301:end
% glossary-block:B01302:begin
\needspace{3\baselineskip}
\textbf{order} — \textbf{hǣs} (n.). command; order. \textbf{hātan} (v.). command; order. Accusative actor + plain infinitive for the commanded action; distinct from call/name.\par
% glossary-block:B01302:end
% glossary-block:B01303:begin
\needspace{3\baselineskip}
\textbf{order someone to} — \textbf{hātan} (v.). hātan + accusative actor + plain infinitive.\par
% glossary-block:B01303:end
% glossary-block:B01304:begin
\needspace{3\baselineskip}
\textbf{ore} — \textbf{ōra} (n.). vein of metal; ore (sense uncertain). Supplied dative plural after on.\par
% glossary-block:B01304:end
% glossary-block:B01305:begin
\needspace{3\baselineskip}
\textbf{Osburg} — \textbf{Osburg} (proper name).\par
% glossary-block:B01305:end
% glossary-block:B01306:begin
\needspace{3\baselineskip}
\textbf{other} — \textbf{ōðer} (adj.; ordinal). other; another.\par
% glossary-block:B01306:end
% glossary-block:B01307:begin
\needspace{3\baselineskip}
\textbf{ought to} — \textbf{sculan} (v.). sculan + plain infinitive for obligation; English have to is not possessive habban. Supplied have something to offer in Unit 6 is a different construction.\par
% glossary-block:B01307:end
% glossary-block:B01308:begin
\needspace{3\baselineskip}
\textbf{our} — \textbf{ic} (personal pron.). Personal genitive plural used attributively in the exact ūre hǣle expression. \textbf{uncer} (possessive adjective). our; ours. Two possessors; agreeing possessive. \textbf{ūre} (possessive adjective). our; ours. Plural possessors; agreeing possessive. \textbf{uncer} (possessive adjective), \textbf{ūre} (possessive adjective). Uncer means belonging to us two; ūre belonging to us in the plural. The possessive agrees with the possessed noun.\par
% glossary-block:B01308:end
% glossary-block:B01309:begin
\needspace{3\baselineskip}
\textbf{ours} — \textbf{uncer} (possessive adjective). our; ours. Two possessors; agreeing possessive. \textbf{ūre} (possessive adjective). our; ours. Plural possessors; agreeing possessive. \textbf{uncer} (possessive adjective), \textbf{ūre} (possessive adjective). Uncer means belonging to us two; ūre belonging to us in the plural. The possessive agrees with the possessed noun.\par
% glossary-block:B01309:end
% glossary-block:B01310:begin
\needspace{3\baselineskip}
\textbf{ourselves} — An ordinary reflexive uses the personal pronoun in the case governed by its verb or preposition. Agreeing self adds the taught emphasis; English -self does not automatically require a separate self. Ordinary reflexive personal pronouns take the governed case; agreeing self adds the taught emphasis. English -self does not automatically require self.\par
% glossary-block:B01310:end
% glossary-block:B01311:begin
\needspace{3\baselineskip}
\textbf{out} — \textbf{ūt} (adv.).\par
% glossary-block:B01311:end
% glossary-block:B01312:begin
\needspace{3\baselineskip}
\textbf{out of} — \textbf{æt} (prep.). from; out of. Proceeds of work, in the assigned expression. \textbf{of} (prep.). from; out of. With dative.\par
% glossary-block:B01312:end
% glossary-block:B01313:begin
\needspace{3\baselineskip}
\textbf{outside} — \textbf{būtan} (prep.). Būtan + dative means without a companion or outside a city in the respective sense. Conditional unless is separate. \textbf{ūte} (adv.).\par
% glossary-block:B01313:end
% glossary-block:B01314:begin
\needspace{3\baselineskip}
\textbf{oven} — \textbf{ofen} (n.).\par
% glossary-block:B01314:end
% glossary-block:B01315:begin
\needspace{3\baselineskip}
\textbf{over} — \textbf{ofer} (prep.). over; across. Spatial frame; case varies with the specific construction. Do not impose a universal motion rule. Accusative dependent in the taught authority frame. over; about. Object of rejoicing, with the assigned accusative frame. Spatial ofer can mean over/across; its case depends on the assigned spatial frame. Ofer in the authority sense over takes an accusative dependent. Rank and preference have further separately qualified senses.\par
% glossary-block:B01315:end
% glossary-block:B01316:begin
\needspace{3\baselineskip}
\textbf{overcome} — \textbf{oferwinnan} (v.). overcome; conquer.\par
% glossary-block:B01316:end
% glossary-block:B01317:begin
\needspace{3\baselineskip}
\textbf{overlay} — \textbf{oferwyrcean} (v.). overlay; cover. Accusative object, mid + dative material used to cover it.\par
% glossary-block:B01317:end
% glossary-block:B01318:begin
\needspace{3\baselineskip}
\textbf{overtake} — \textbf{fōn} (v.). Accusative person overtaken in these tasks; the source's bracketed caught is the route.\par
% glossary-block:B01318:end
% glossary-block:B01319:begin
\needspace{3\baselineskip}
\textbf{overthrow} — \textbf{tōweorpan} (v.). break; overthrow. Break a commandment in the assigned use.\par
% glossary-block:B01319:end
% glossary-block:B01320:begin
\needspace{3\baselineskip}
\textbf{own} — \textbf{āgen} (adj.).  \textbf{āgan} (v.). have; possess. Accusative possession; positive present and preterite indicative in the taught scope.\par
% glossary-block:B01320:end
% book-navigation: standalone bookmark replaced by volume hierarchy
% glossary-block:B01321:begin
\needspace{5\baselineskip}
\subsection*{P}\BookHeading{nav-glossary-english-0017}{subsubsection}{P}{0}
% glossary-block:B01321:end
% glossary-block:B01322:begin
\textbf{palm branch} — \textbf{palmtwīg} (n.). palm twig; palm branch.\par
% glossary-block:B01322:end
% glossary-block:B01323:begin
\needspace{3\baselineskip}
\textbf{palm twig} — \textbf{palmtwīg} (n.). palm twig; palm branch.\par
% glossary-block:B01323:end
% glossary-block:B01324:begin
\needspace{3\baselineskip}
\textbf{parch} — \textbf{forbærnan} (v.). parch; scorch.\par
% glossary-block:B01324:end
% glossary-block:B01325:begin
\needspace{3\baselineskip}
\textbf{part} — \textbf{dǣl} (n.). dǣl + genitive of the substance, as in sumne dǣl mętes; the amount is accusative object.\par
% glossary-block:B01325:end
% glossary-block:B01326:begin
\needspace{3\baselineskip}
\textbf{pass} — \textbf{gewītan} (v.). pass away; cease to hold. Of a decree.\par
% glossary-block:B01326:end
% glossary-block:B01327:begin
\needspace{3\baselineskip}
\textbf{passage} — \textbf{sīðfæt} (n.). journey; passage.\par
% glossary-block:B01327:end
% glossary-block:B01328:begin
\needspace{3\baselineskip}
\textbf{pass away} — \textbf{gewītan} (v.). pass away; cease to hold. Of a decree.\par
% glossary-block:B01328:end
% glossary-block:B01329:begin
\needspace{3\baselineskip}
\textbf{path} — \textbf{weg} (n.). way; path. Weg means road, path or way.\par
% glossary-block:B01329:end
% glossary-block:B01330:begin
\needspace{3\baselineskip}
\textbf{pay} — \textbf{forgieldan} (v.). pay; repay; requite. Dative recipient, accusative thing repaid in the selected frame.\par
% glossary-block:B01330:end
% glossary-block:B01331:begin
\needspace{3\baselineskip}
\textbf{peace} — \textbf{sibb} (n.).\par
% glossary-block:B01331:end
% glossary-block:B01332:begin
\needspace{3\baselineskip}
\textbf{pearl} — \textbf{męregrēote} (n.).\par
% glossary-block:B01332:end
% glossary-block:B01333:begin
\needspace{3\baselineskip}
\textbf{peer} — \textbf{gelīca} (n.). equal; peer.\par
% glossary-block:B01333:end
% glossary-block:B01334:begin
\needspace{3\baselineskip}
\textbf{penalty} — \textbf{wīte} (n.). punishment; penalty.\par
% glossary-block:B01334:end
% glossary-block:B01335:begin
\needspace{3\baselineskip}
\textbf{penetrate} — \textbf{fēolan} (v.). Preserve the supplied/taught lexical frame; not a synonym for every kind of entering.\par
% glossary-block:B01335:end
% glossary-block:B01336:begin
\needspace{3\baselineskip}
\textbf{penny} — \textbf{pęning} (n.). penny; coin.\par
% glossary-block:B01336:end
% glossary-block:B01337:begin
\needspace{3\baselineskip}
\textbf{Pentapolis} — \textbf{Pentapolis} (proper place name). Cyrenian region in this story.\par
% glossary-block:B01337:end
% glossary-block:B01338:begin
\needspace{3\baselineskip}
\textbf{people} — \textbf{folc} (n.). Collective singular. \textbf{ðēod} (n.). people; nation.\par
% glossary-block:B01338:end
% glossary-block:B01339:begin
\needspace{3\baselineskip}
\textbf{people in general} — \textbf{mann} (n.; impersonal pronoun). one; people in general. Impersonal man with singular finite verb.\par
% glossary-block:B01339:end
% glossary-block:B01340:begin
\needspace{3\baselineskip}
\textbf{people not noble} — \textbf{unæðele} (adj.). Substantival adjective in the supplied contrast.\par
% glossary-block:B01340:end
% glossary-block:B01341:begin
\needspace{3\baselineskip}
\textbf{peoples} — \textbf{folc} (n.). Distinct peoples or national groups; plural sense.\par
% glossary-block:B01341:end
% glossary-block:B01342:begin
\needspace{3\baselineskip}
\textbf{perceive} — \textbf{oncnāwan} (v.). recognize; perceive. Supplied past source form. \textbf{ongietan} (v.). understand; perceive. Supplied infinitive after a modal. perceive; realize. understand; perceive.\par
% glossary-block:B01342:end
% glossary-block:B01343:begin
\needspace{3\baselineskip}
\textbf{perform} — \textbf{āwyrcean} (v.). do; deserve. In the source, deserving through one's conduct; retain the complete local frame. \textbf{dōn} (v.). do; perform. \textbf{ęfnan} (v.).  \textbf{forðtēon} (v.). present; perform. Source performance sense; retain the actual form's spelling. \textbf{gedōn} (v.). do; perform. \textbf{gefręmman} (v.). perform; accomplish. A work is the thing accomplished. \textbf{plegian} (v.). act; perform. Performance context; ordinary play is a separate assigned sense. \textbf{wyrcean} (v.). work; perform. Accusative work or deed.\par
% glossary-block:B01343:end
% glossary-block:B01344:begin
\needspace{3\baselineskip}
\textbf{perhaps} — \textbf{wēn} (n.). expectation; chance. wēn is ðæt + subjunctive supplies the taught perhaps construction. wēn is ðæt + present subjunctive in the taught possibility frame.\par
% glossary-block:B01344:end
% glossary-block:B01345:begin
\needspace{3\baselineskip}
\textbf{peril} — \textbf{frēcnes} (n.). danger; peril. \textbf{pliht} (n.). peril; risk.\par
% glossary-block:B01345:end
% glossary-block:B01346:begin
\needspace{3\baselineskip}
\textbf{period} — \textbf{fierst} (n.). period; interval. Masculine noun in the source time expression.\par
% glossary-block:B01346:end
% glossary-block:B01347:begin
\needspace{3\baselineskip}
\textbf{perish} — \textbf{forweorðan} (v.).\par
% glossary-block:B01347:end
% glossary-block:B01348:begin
\needspace{3\baselineskip}
\textbf{permission} — \textbf{lēaf} (n.).\par
% glossary-block:B01348:end
% glossary-block:B01349:begin
\needspace{3\baselineskip}
\textbf{permit} — \textbf{ālīefan} (v.). permit; allow. Impersonal permission frame, with dative permitted person; consult the complete supplied expression. \textbf{gelīefst} (v.; supplied pres. ind. 2 sg.). permit; allow. \textbf{līefan} (v.). allow; permit.\par
% glossary-block:B01349:end
% glossary-block:B01350:begin
\needspace{3\baselineskip}
\textbf{persecute} — \textbf{ēhtan} (v.). pursue; persecute. Accusative person in the selected production frame; a separately supported genitive alternative is not the selected target.\par
% glossary-block:B01350:end
% glossary-block:B01351:begin
\needspace{3\baselineskip}
\textbf{petition} — \textbf{bēn} (n.). prayer; petition. Source dative plural after tō. petition; appeal.\par
% glossary-block:B01351:end
% glossary-block:B01352:begin
\needspace{3\baselineskip}
\textbf{Picts} — \textbf{Piht} (ethnonym). Supplied people-name forms.\par
% glossary-block:B01352:end
% glossary-block:B01353:begin
\needspace{3\baselineskip}
\textbf{pitier} — \textbf{gemiltsiend} (n.). one who shows mercy; pitier.\par
% glossary-block:B01353:end
% glossary-block:B01354:begin
\needspace{3\baselineskip}
\textbf{pity} — \textbf{gemiltsian} (v.). pity; have compassion on. Dative person. gemiltsian + dative person. A wish uses the appropriate subjunctive, not a separate have verb.\par
% glossary-block:B01354:end
% glossary-block:B01355:begin
\needspace{3\baselineskip}
\textbf{place} — \textbf{āsęttan} (v.). set; place. \textbf{dōn} (v.). put; place. Accusative object; on + accusative goal in the assigned frame. \textbf{gelōgian} (v.). place; set. \textbf{gesęttan} (v.). set; place. Supplied gesętte. set; place; establish. \textbf{lęcgean} (v.). lay; place. Supplied source form. lay; place. \textbf{sęttan} (v.). set; place. Place an accusative object; destination is a separate phrase.\par
% glossary-block:B01355:end
% glossary-block:B01356:begin
\needspace{3\baselineskip}
\textbf{placed} — \textbf{gesęttan} (v.). placed; situated. Supplied participial complement of understood us; not the identical-looking finite preterite.\par
% glossary-block:B01356:end
% glossary-block:B01357:begin
\needspace{3\baselineskip}
\textbf{place of prayer} — \textbf{cirice} (n.). oratory. Chad's separate place of prayer; distinct from the ordinary church sense.\par
% glossary-block:B01357:end
% glossary-block:B01358:begin
\needspace{3\baselineskip}
\textbf{plait} — \textbf{windan} (v.). wind; plait. Accusative crown; material follows of. Cook's platted means plaited.\par
% glossary-block:B01358:end
% glossary-block:B01359:begin
\needspace{3\baselineskip}
\textbf{plant} — \textbf{wyrt} (n.). herb; plant.\par
% glossary-block:B01359:end
% glossary-block:B01360:begin
\needspace{3\baselineskip}
\textbf{plat} — \textbf{windan} (v.). Accusative crown; material follows of. Cook's platted means plaited.\par
% glossary-block:B01360:end
% glossary-block:B01361:begin
\needspace{3\baselineskip}
\textbf{platted} — \textbf{windan} (v.). Accusative crown; material follows of. Cook's platted means plaited.\par
% glossary-block:B01361:end
% glossary-block:B01362:begin
\needspace{3\baselineskip}
\textbf{play} — \textbf{plega} (n.). play; game. \textbf{plegian} (v.). Intransitive; the reading supplies local wið participant and mid instrument phrases.\par
% glossary-block:B01362:end
% glossary-block:B01363:begin
\needspace{3\baselineskip}
\textbf{play harp} — \textbf{hearpian} (v.).\par
% glossary-block:B01363:end
% glossary-block:B01364:begin
\needspace{3\baselineskip}
\textbf{play the harp} — \textbf{hearpian} (v.). harp. Supplied genitive harper/harping construction, with instrument; do not reconstruct it from the infinitive alone. Preserve the supplied genitive harper/harping expression; headword alone does not reconstruct its syntax.\par
% glossary-block:B01364:end
% glossary-block:B01365:begin
\needspace{3\baselineskip}
\textbf{play with} — \textbf{wið} (prep.). with; against. Dative fellow players in the exact supplied playing frame.\par
% glossary-block:B01365:end
% glossary-block:B01366:begin
\needspace{3\baselineskip}
\textbf{pleasant} — \textbf{wynsum} (adj.). pleasant; winsome.\par
% glossary-block:B01366:end
% glossary-block:B01367:begin
\needspace{3\baselineskip}
\textbf{pleasantly} — \textbf{gecwēmlīce} (adv.). pleasantly; agreeably.\par
% glossary-block:B01367:end
% glossary-block:B01368:begin
\needspace{3\baselineskip}
\textbf{please} — \textbf{lician} (v.). The thing pleasing is the subject; the persons pleased are dative in the supplied reading. \textbf{lystan} (v.). please; be desired. Impersonal source expression; keep the complete supplied frame.\par
% glossary-block:B01368:end
% glossary-block:B01369:begin
\needspace{3\baselineskip}
\textbf{pleasing} — \textbf{gecwēme} (adj.). pleasing; acceptable.\par
% glossary-block:B01369:end
% glossary-block:B01370:begin
\needspace{3\baselineskip}
\textbf{plectrum} — \textbf{hearpenægl} (n.). plectrum; harp pick.\par
% glossary-block:B01370:end
% glossary-block:B01371:begin
\needspace{3\baselineskip}
\textbf{plow} — \textbf{ęrian} (v.).\par
% glossary-block:B01371:end
% glossary-block:B01372:begin
\needspace{3\baselineskip}
\textbf{plowman} — \textbf{ierðling} (n.). farmer; plowman.\par
% glossary-block:B01372:end
% glossary-block:B01373:begin
\needspace{3\baselineskip}
\textbf{plowshare} — \textbf{sulhscear} (n.).\par
% glossary-block:B01373:end
% glossary-block:B01374:begin
\needspace{3\baselineskip}
\textbf{plunder} — \textbf{hęrgian} (v.). plunder; ravage.\par
% glossary-block:B01374:end
% glossary-block:B01375:begin
\needspace{3\baselineskip}
\textbf{ply} — \textbf{begān} (v.). ply; practise. Accusative trade or occupation.\par
% glossary-block:B01375:end
% glossary-block:B01376:begin
\needspace{3\baselineskip}
\textbf{poem} — \textbf{lēoð} (n.). poem; poetry; verse.\par
% glossary-block:B01376:end
% glossary-block:B01377:begin
\needspace{3\baselineskip}
\textbf{poet} — \textbf{lēoðwyrhta} (n.).\par
% glossary-block:B01377:end
% glossary-block:B01378:begin
\needspace{3\baselineskip}
\textbf{poetry} — \textbf{lēoð} (n.). poem; poetry; verse.\par
% glossary-block:B01378:end
% glossary-block:B01379:begin
\needspace{3\baselineskip}
\textbf{ponder} — \textbf{smēagan} (v.). ponder; consider. Supplied preterite source form.\par
% glossary-block:B01379:end
% glossary-block:B01380:begin
\needspace{3\baselineskip}
\textbf{poor man} — \textbf{wǣdla} (n.), \textbf{ðearfa} (n.). Wǣdla denotes a poor or destitute man; ðearfa a poor or needy person. Use the noun specified by the exercise.\par
% glossary-block:B01380:end
% glossary-block:B01381:begin
\needspace{3\baselineskip}
\textbf{poor or destitute man} — \textbf{wǣdla} (n.).\par
% glossary-block:B01381:end
% glossary-block:B01382:begin
\needspace{3\baselineskip}
\textbf{poor person} — \textbf{ðearfa} (n.). poor person; needy person. Noun referring to a person. \textbf{wǣdla} (n.), \textbf{ðearfa} (n.). Wǣdla denotes a poor or destitute man; ðearfa a poor or needy person. Use the noun specified by the exercise.\par
% glossary-block:B01382:end
% glossary-block:B01383:begin
\needspace{3\baselineskip}
\textbf{possess} — \textbf{habban} (v.). have; possess. Accusative possession; distinguish the perfect auxiliary. \textbf{āgan} (v.). have; possess. Accusative possession; positive present and preterite indicative in the taught scope.\par
% glossary-block:B01383:end
% glossary-block:B01384:begin
\needspace{3\baselineskip}
\textbf{possessions} — \textbf{ǣht} (n.). possessions; goods.\par
% glossary-block:B01384:end
% glossary-block:B01385:begin
\needspace{3\baselineskip}
\textbf{pound} — \textbf{pund} (n.). Weight measure in the supplied gold quantity.\par
% glossary-block:B01385:end
% glossary-block:B01386:begin
\needspace{3\baselineskip}
\textbf{pour} — \textbf{scęncan} (v.). pour out; give to drink. Accusative drink, dative recipient.\par
% glossary-block:B01386:end
% glossary-block:B01387:begin
\needspace{3\baselineskip}
\textbf{pour out} — \textbf{scęncan} (v.). pour out; give to drink. Accusative drink, dative recipient.\par
% glossary-block:B01387:end
% glossary-block:B01388:begin
\needspace{3\baselineskip}
\textbf{power} — \textbf{mægen} (n.). power; strength. Heavenly powers in 37.2 is the contextual plural use. \textbf{miht} (n.). ability; power.\par
% glossary-block:B01388:end
% glossary-block:B01389:begin
\needspace{3\baselineskip}
\textbf{powerful} — \textbf{mihtig} (adj.). mighty; powerful.\par
% glossary-block:B01389:end
% glossary-block:B01390:begin
\needspace{3\baselineskip}
\textbf{powerful people} — \textbf{miht} (n.). powers; powerful people. Those holding high station in the source's concluding human application.\par
% glossary-block:B01390:end
% glossary-block:B01391:begin
\needspace{3\baselineskip}
\textbf{powers} — \textbf{miht} (n.). powers; powerful people. Those holding high station in the source's concluding human application.\par
% glossary-block:B01391:end
% glossary-block:B01392:begin
\needspace{3\baselineskip}
\textbf{practice} — \textbf{begān} (v.). ply; practise. Accusative trade or occupation. pursue; practise. Study is the thing pursued.\par
% glossary-block:B01392:end
% glossary-block:B01393:begin
\needspace{3\baselineskip}
\textbf{practise} — \textbf{begān} (v.). ply; practise. Accusative trade or occupation. pursue; practise. Study is the thing pursued.\par
% glossary-block:B01393:end
% glossary-block:B01394:begin
\needspace{3\baselineskip}
\textbf{praise} — \textbf{hęrian} (v.). Genitive plural present-participle use, of people praising. Supplied source form praised. \textbf{lof} (n.).\par
% glossary-block:B01394:end
% glossary-block:B01395:begin
\needspace{3\baselineskip}
\textbf{pray} — \textbf{biddan} (v.). beseech; entreat. Accusative person; the requested action can be a ðæt clause. Accusative person beseeched; ðæt + subjunctive gives the requested act. Prayer to God with gebiddan is a different entry. \textbf{gebiddan} (v.). Supplied wæs gebiddende, rendered prayed in this source. With an accusative reflexive pronoun in the assigned prayer frame. For prayer, gebiddan takes the assigned accusative reflexive pronoun. Older English pray meaning beseech/entreat a person uses biddan + accusative person; a requested action can be a ðæt clause.\par
% glossary-block:B01395:end
% glossary-block:B01396:begin
\needspace{3\baselineskip}
\textbf{prayer} — \textbf{bēn} (n.). prayer; petition. Source dative plural after tō. \textbf{gebed} (n.).\par
% glossary-block:B01396:end
% glossary-block:B01397:begin
\needspace{3\baselineskip}
\textbf{precious} — \textbf{dēorwierðe} (adj.). precious; valuable. \textbf{dēorwurð} (adj.). precious; costly.\par
% glossary-block:B01397:end
% glossary-block:B01398:begin
\needspace{3\baselineskip}
\textbf{precious stone} — \textbf{gimm} (n.). gem; precious stone.\par
% glossary-block:B01398:end
% glossary-block:B01399:begin
\needspace{3\baselineskip}
\textbf{precisely} — \textbf{efne} (interj.). just; precisely; even.\par
% glossary-block:B01399:end
% glossary-block:B01400:begin
\needspace{3\baselineskip}
\textbf{prefer} — \textbf{lēofra} (adj.). preferable; dearer. Dative experiencer + is lēofre + infinitive in the taught preference frame. The preference construction is dative experiencer + is lēofre + plain infinitive, is preferable to someone. It renders had/would rather without requiring past time. Wician in the supplied example is the action preferred, sojourn/live with, not a word meaning prefer.\par
% glossary-block:B01400:end
% glossary-block:B01401:begin
\needspace{3\baselineskip}
\textbf{preferable} — \textbf{lēofra} (adj.). preferable; dearer. Dative experiencer + is lēofre + infinitive in the taught preference frame.\par
% glossary-block:B01401:end
% glossary-block:B01402:begin
\needspace{3\baselineskip}
\textbf{prepare} — \textbf{gearcian} (v.).  \textbf{gierwan} (v.). Accusative thing prepared.\par
% glossary-block:B01402:end
% glossary-block:B01403:begin
\needspace{3\baselineskip}
\textbf{present} — \textbf{andweard} (adj.).  \textbf{forðtēon} (v.). present; perform. Source performance sense; retain the actual form's spelling. \textbf{giefu} (n.). gift. Supplied plural gifts. \textbf{lāc} (n.). gift; present.\par
% glossary-block:B01403:end
% glossary-block:B01404:begin
\needspace{3\baselineskip}
\textbf{price} — \textbf{weorð} (n.). price; value.\par
% glossary-block:B01404:end
% glossary-block:B01405:begin
\needspace{3\baselineskip}
\textbf{priest} — \textbf{mæsseprēost} (n.).\par
% glossary-block:B01405:end
% glossary-block:B01406:begin
\needspace{3\baselineskip}
\textbf{primacy} — \textbf{ealdordōm} (n.). chief place; primacy. \textbf{ealdorscipe} (n.). chief place; primacy.\par
% glossary-block:B01406:end
% glossary-block:B01407:begin
\needspace{3\baselineskip}
\textbf{prince} — \textbf{ealdormann} (n.). prince; leader; noble.\par
% glossary-block:B01407:end
% glossary-block:B01408:begin
\needspace{3\baselineskip}
\textbf{princess} — \textbf{cwēn} (n.). The context distinguishes princess from queen. The selected accusative singular cwēn is endingless.\par
% glossary-block:B01408:end
% glossary-block:B01409:begin
\needspace{3\baselineskip}
\textbf{prior} — \textbf{furðra} (adj.). foremost; prior.\par
% glossary-block:B01409:end
% glossary-block:B01410:begin
\needspace{3\baselineskip}
\textbf{prison} — \textbf{cweartern} (n.).\par
% glossary-block:B01410:end
% glossary-block:B01411:begin
\needspace{3\baselineskip}
\textbf{produce} — \textbf{ācęnnan} (v.). produce; bring forth. \textbf{beran} (v.). bear; produce. Accusative produce; distinct from physical carrying. \textbf{tēon} (v.). bring forth; produce. Accusative thing produced with particle forð; use the assigned tēon forms. bring forth; produce. Supplied third-plural present subjunctive, with waters as subject. \textbf{wæstm} (n.). fruit; produce.\par
% glossary-block:B01411:end
% glossary-block:B01412:begin
\needspace{3\baselineskip}
\textbf{produce fruit} — \textbf{wyrcean} (v.). yield; produce. Accusative fruit in the taught/source expression. Accusative fruit; birð wæstm is a recognized lexical alternative in the key.\par
% glossary-block:B01412:end
% glossary-block:B01413:begin
\needspace{3\baselineskip}
\textbf{produce of the earth} — \textbf{ierð} (n.). crops; produce of the earth.\par
% glossary-block:B01413:end
% glossary-block:B01414:begin
\needspace{3\baselineskip}
\textbf{produce seed} — \textbf{beran} (v.). Accusative produce; distinguish physical carrying.\par
% glossary-block:B01414:end
% glossary-block:B01415:begin
\needspace{3\baselineskip}
\textbf{productive of} — \textbf{beran} (v.). rich in; productive of. Supplied berende on + dative ore expression; participial adjective.\par
% glossary-block:B01415:end
% glossary-block:B01416:begin
\needspace{3\baselineskip}
\textbf{profit} — \textbf{gestrēon} (n.). gain; profit.\par
% glossary-block:B01416:end
% glossary-block:B01417:begin
\needspace{3\baselineskip}
\textbf{promise} — \textbf{behātan} (v.).  \textbf{gehātan} (v.). Accusative thing promised; a dative recipient may be added.\par
% glossary-block:B01417:end
% glossary-block:B01418:begin
\needspace{3\baselineskip}
\textbf{prompt} — \textbf{myngian} (v.). admonish; prompt. Supplied participial source use.\par
% glossary-block:B01418:end
% glossary-block:B01419:begin
\needspace{3\baselineskip}
\textbf{property} — \textbf{gōd} (n.). goods; property. Plural lexical noun.\par
% glossary-block:B01419:end
% glossary-block:B01420:begin
\needspace{3\baselineskip}
\textbf{prophet} — \textbf{wītga} (n.).\par
% glossary-block:B01420:end
% glossary-block:B01421:begin
\needspace{3\baselineskip}
\textbf{prostrate oneself} — \textbf{āstręccean} (v.). Reflexive accusative; the before phrase retains its own case/construal.\par
% glossary-block:B01421:end
% glossary-block:B01422:begin
\needspace{3\baselineskip}
\textbf{protect} — \textbf{gescieldan} (v.). shield; protect.\par
% glossary-block:B01422:end
% glossary-block:B01423:begin
\needspace{3\baselineskip}
\textbf{protection} — \textbf{gebeorg} (n.). protection; defence.\par
% glossary-block:B01423:end
% glossary-block:B01424:begin
\needspace{3\baselineskip}
\textbf{proud} — \textbf{wlanc} (adj.). proud; wealthy. Contextually the rich man addressed.\par
% glossary-block:B01424:end
% glossary-block:B01425:begin
\needspace{3\baselineskip}
\textbf{provender} — \textbf{fōdor} (n.). Provender is routed by Cook's bracketed fodder to fōdor.\par
% glossary-block:B01425:end
% glossary-block:B01426:begin
\needspace{3\baselineskip}
\textbf{provide} — \textbf{tilian} (v.). obtain; provide. Genitive thing procured; dative reflexive beneficiary in the assigned frame.\par
% glossary-block:B01426:end
% glossary-block:B01427:begin
\needspace{3\baselineskip}
\textbf{provide for oneself} — \textbf{tilian} (v.). obtain; provide. Genitive thing procured; dative reflexive beneficiary in the assigned frame. tilian + dative reflexive beneficiary + genitive thing obtained.\par
% glossary-block:B01427:end
% glossary-block:B01428:begin
\needspace{3\baselineskip}
\textbf{provisions} — \textbf{ēst} (n.). provisions; food.\par
% glossary-block:B01428:end
% glossary-block:B01429:begin
\needspace{3\baselineskip}
\textbf{public place} — \textbf{strǣt} (n.). Street used as a gathering place in the assigned prompt.\par
% glossary-block:B01429:end
% glossary-block:B01430:begin
\needspace{3\baselineskip}
\textbf{pull} — \textbf{tēon} (v.). draw; pull.\par
% glossary-block:B01430:end
% glossary-block:B01431:begin
\needspace{3\baselineskip}
\textbf{punishment} — \textbf{wīte} (n.). punishment; penalty.\par
% glossary-block:B01431:end
% glossary-block:B01432:begin
\needspace{3\baselineskip}
\textbf{pure} — \textbf{clǣne} (adj.). pure; clean. clean; pure; blameless.\par
% glossary-block:B01432:end
% glossary-block:B01433:begin
\needspace{3\baselineskip}
\textbf{purple cloth} — \textbf{gōdwębb} (n.).\par
% glossary-block:B01433:end
% glossary-block:B01434:begin
\needspace{3\baselineskip}
\textbf{pursue} — \textbf{begān} (v.). pursue; practise. Study is the thing pursued. \textbf{ēhtan} (v.). pursue; persecute. Accusative person in the selected production frame; a separately supported genitive alternative is not the selected target.\par
% glossary-block:B01434:end
% glossary-block:B01435:begin
\needspace{3\baselineskip}
\textbf{put} — \textbf{dōn} (v.). put; place. Accusative object; on + accusative goal in the assigned frame. \textbf{sęttan} (v.). Place an accusative object; destination is a separate phrase.\par
% glossary-block:B01435:end
% glossary-block:B01436:begin
\needspace{3\baselineskip}
\textbf{put down} — \textbf{āweorpan} (v.). Accusative people put down, of + dative point from which. \textbf{sęttan} (v.). Sęttan means put down/place an accusative thing, such as seed. This placement sense is distinct from āweorpan cast down/depose.\par
% glossary-block:B01436:end
% glossary-block:B01437:begin
\needspace{3\baselineskip}
\textbf{put to flight} — \textbf{āflīeman} (v.). put to flight; expel.\par
% glossary-block:B01437:end
% book-navigation: standalone bookmark replaced by volume hierarchy
% glossary-block:B01438:begin
\needspace{5\baselineskip}
\subsection*{Q}\BookHeading{nav-glossary-english-0018}{subsubsection}{Q}{0}
% glossary-block:B01438:end
% glossary-block:B01439:begin
\textbf{quantity} — \textbf{unrīm} (n.). abundance; uncounted quantity. Supplied quantity with genitive of the substance.\par
% glossary-block:B01439:end
% glossary-block:B01440:begin
\needspace{3\baselineskip}
\textbf{queen} — \textbf{cwēn} (n.). The context distinguishes princess from queen. The selected accusative singular cwēn is endingless.\par
% glossary-block:B01440:end
% glossary-block:B01441:begin
\needspace{3\baselineskip}
\textbf{quell} — \textbf{gestillan} (v.). quell; settle. Accusative disturbance settled; distinguish intransitive become still.\par
% glossary-block:B01441:end
% glossary-block:B01442:begin
\needspace{3\baselineskip}
\textbf{quickly} — \textbf{hrædlīce} (adv.).  \textbf{hraðe} (adv.).\par
% glossary-block:B01442:end
% glossary-block:B01443:begin
\needspace{3\baselineskip}
\textbf{quickness} — \textbf{hrædnes} (n.). speed; quickness.\par
% glossary-block:B01443:end
% glossary-block:B01444:begin
\needspace{3\baselineskip}
\textbf{quiet} — \textbf{stille} (adj.). still; quiet.\par
% glossary-block:B01444:end
% glossary-block:B01445:begin
\needspace{3\baselineskip}
\textbf{quietness} — \textbf{stilnes} (n.). stillness; quietness.\par
% glossary-block:B01445:end
% book-navigation: standalone bookmark replaced by volume hierarchy
% glossary-block:B01446:begin
\needspace{5\baselineskip}
\subsection*{R}\BookHeading{nav-glossary-english-0019}{subsubsection}{R}{0}
% glossary-block:B01446:end
% glossary-block:B01447:begin
\textbf{raging} — \textbf{angrīslic} (adj.). fierce; raging.\par
% glossary-block:B01447:end
% glossary-block:B01448:begin
\needspace{3\baselineskip}
\textbf{raiment} — \textbf{rēaf} (n.). raiment; clothing. Wǣd denotes a garment, rēaf raiment/clothing, and gegierela a garment or clothing. Preserve any noun expressly supplied by the exercise.\par
% glossary-block:B01448:end
% glossary-block:B01449:begin
\needspace{3\baselineskip}
\textbf{rain} — \textbf{regn} (n.).  \textbf{rīnan} (v.). Supplied impersonal source clause with past rīnde.\par
% glossary-block:B01449:end
% glossary-block:B01450:begin
\needspace{3\baselineskip}
\textbf{raise} — \textbf{āhębban} (v.). raise; lift. \textbf{hębban} (v.). raise; lift. \textbf{ūp-ā-hębban} (v.). lift up. Supplied source form. lift up. Supplied first-plural preterite. raise; exalt. Accusative people raised; the prefixed verb supplies up.\par
% glossary-block:B01450:end
% glossary-block:B01451:begin
\needspace{3\baselineskip}
\textbf{rather} — Comparative swīðor expresses rather/more in the preference construction; the positive adverb swīðe means greatly/very.\par
% glossary-block:B01451:end
% glossary-block:B01452:begin
\needspace{3\baselineskip}
\textbf{ravage} — \textbf{hęrgian} (v.). plunder; ravage.\par
% glossary-block:B01452:end
% glossary-block:B01453:begin
\needspace{3\baselineskip}
\textbf{read} — \textbf{rǣdan} (v.). Accusative text, or absolute with no expressed object.\par
% glossary-block:B01453:end
% glossary-block:B01454:begin
\needspace{3\baselineskip}
\textbf{reading books} — \textbf{bēcrǣding} (n.).\par
% glossary-block:B01454:end
% glossary-block:B01455:begin
\needspace{3\baselineskip}
\textbf{ready} — \textbf{gearu} (adj.).\par
% glossary-block:B01455:end
% glossary-block:B01456:begin
\needspace{3\baselineskip}
\textbf{realize} — \textbf{ongietan} (v.). perceive; realize.\par
% glossary-block:B01456:end
% glossary-block:B01457:begin
\needspace{3\baselineskip}
\textbf{reap} — \textbf{rīpan} (v.).\par
% glossary-block:B01457:end
% glossary-block:B01458:begin
\needspace{3\baselineskip}
\textbf{receive} — \textbf{onfōn} (v.). receive; accept. Dative thing received in this selected frame; the later accusative payment frame is separate. receive; accept. A person received into care in the supplied frame. Selected payment frame has accusative pęning in 28.11, but selected reward frame dative edlēane in 39.6; show both with their task-specific limits. \textbf{underfōn} (v.). receive; take charge of.\par
% glossary-block:B01458:end
% glossary-block:B01459:begin
\needspace{3\baselineskip}
\textbf{recognize} — \textbf{oncnāwan} (v.). recognize; perceive. Supplied past source form. recognize; acknowledge. recognize; know by recognition.\par
% glossary-block:B01459:end
% glossary-block:B01460:begin
\needspace{3\baselineskip}
\textbf{recompense} — \textbf{edlēan} (n.). reward; recompense. \textbf{forgieldan} (v.). pay; repay; requite. Dative recipient, accusative thing repaid in the selected frame.\par
% glossary-block:B01460:end
% glossary-block:B01461:begin
\needspace{3\baselineskip}
\textbf{Redeemer} — \textbf{ālīesend} (n.). Redeemer; one who delivers.\par
% glossary-block:B01461:end
% glossary-block:B01462:begin
\needspace{3\baselineskip}
\textbf{reduce to misery} — \textbf{forierman} (v.). ruin; reduce to misery.\par
% glossary-block:B01462:end
% glossary-block:B01463:begin
\needspace{3\baselineskip}
\textbf{reduce to servitude} — \textbf{ðēowian} (v.). enslave; reduce to servitude.\par
% glossary-block:B01463:end
% glossary-block:B01464:begin
\needspace{3\baselineskip}
\textbf{reed} — \textbf{hrēod} (n.).\par
% glossary-block:B01464:end
% glossary-block:B01465:begin
\needspace{3\baselineskip}
\textbf{reeve} — \textbf{gerēfa} (n.). steward; reeve.\par
% glossary-block:B01465:end
% glossary-block:B01466:begin
\needspace{3\baselineskip}
\textbf{refuse} — \textbf{forwiernan} (v.). deny; refuse. Dative person denied, genitive thing denied. Dative person denied, genitive thing denied. \textbf{willan} (v.). be unwilling; refuse. With ne or a licensed negative contraction; followed by plain infinitive.\par
% glossary-block:B01466:end
% glossary-block:B01467:begin
\needspace{3\baselineskip}
\textbf{refuse to} — \textbf{willan} (v.). ne + willan, or the licensed contracted negative, with a plain infinitive.\par
% glossary-block:B01467:end
% glossary-block:B01468:begin
\needspace{3\baselineskip}
\textbf{regard} — \textbf{gesēon} (v.). regard; take notice of.\par
% glossary-block:B01468:end
% glossary-block:B01469:begin
\needspace{3\baselineskip}
\textbf{region} — \textbf{ðēod} (n.). region; country.\par
% glossary-block:B01469:end
% glossary-block:B01470:begin
\needspace{3\baselineskip}
\textbf{regularly} — \textbf{gelōmlice} (adv.). frequently; regularly.\par
% glossary-block:B01470:end
% glossary-block:B01471:begin
\needspace{3\baselineskip}
\textbf{rejoice} — \textbf{blissian} (v.). Source genitive plural present-participle use, of people rejoicing. Intransitive; mid + dative companion or ofer + accusative cause of joy in the assigned frames.\par
% glossary-block:B01471:end
% glossary-block:B01472:begin
\needspace{3\baselineskip}
\textbf{rejoice in} — \textbf{on} (prep.). in. Source/object of joy in the supplied frame.\par
% glossary-block:B01472:end
% glossary-block:B01473:begin
\needspace{3\baselineskip}
\textbf{rejoice over} — \textbf{ofer} (prep.). over; about. Object of rejoicing, with the assigned accusative frame.\par
% glossary-block:B01473:end
% glossary-block:B01474:begin
\needspace{3\baselineskip}
\textbf{relate to} — \textbf{belimpan} (v.). belong; relate. Supplied participial construction.\par
% glossary-block:B01474:end
% glossary-block:B01475:begin
\needspace{3\baselineskip}
\textbf{release} — \textbf{ālīesan} (v.). deliver; release. Accusative person delivered or released; fram + dative marks separation. \textbf{onlīesan} (v.). Accusative person or thing released.\par
% glossary-block:B01475:end
% glossary-block:B01476:begin
\needspace{3\baselineskip}
\textbf{remain} — \textbf{ðurhwunian} (v.). remain; endure; last. Intransitive continuing of a building; distinguish the following dative-governing wiðstandan. \textbf{wunian} (v.). dwell; remain.\par
% glossary-block:B01476:end
% glossary-block:B01477:begin
\needspace{3\baselineskip}
\textbf{remaining} — \textbf{ōðer} (adj.; ordinal). remaining; the rest. The rest of the forest/earth in the source.\par
% glossary-block:B01477:end
% glossary-block:B01478:begin
\needspace{3\baselineskip}
\textbf{remarkable} — \textbf{sellic} (adj.). strange; remarkable.\par
% glossary-block:B01478:end
% glossary-block:B01479:begin
\needspace{3\baselineskip}
\textbf{remember} — \textbf{gemunan} (v.).  \textbf{geðęncean} (v.). remember; bear in mind.\par
% glossary-block:B01479:end
% glossary-block:B01480:begin
\needspace{3\baselineskip}
\textbf{remove} — \textbf{āfierran} (v.). remove; take away.\par
% glossary-block:B01480:end
% glossary-block:B01481:begin
\needspace{3\baselineskip}
\textbf{render} — \textbf{dōn} (v.). make; render. Accusative object with agreeing object predicate. \textbf{geręccan} (v.). explain; render. \textbf{gewyrcean} (v.). make; render. Accusative object with an agreeing predicate noun/adjective.\par
% glossary-block:B01481:end
% glossary-block:B01482:begin
\needspace{3\baselineskip}
\textbf{renowned} — \textbf{mǣre} (adj.). great; renowned. Distinct from simple physical size where the task requires the renowned sense.\par
% glossary-block:B01482:end
% glossary-block:B01483:begin
\needspace{3\baselineskip}
\textbf{repay} — \textbf{āgyldan} (v.). reward; repay. C06 corrects Cook's āgyltan to strong III āgyldan; supplied contracted present āgylt and dative recipient. \textbf{forgieldan} (v.). pay; repay; requite. Dative recipient, accusative thing repaid in the selected frame.\par
% glossary-block:B01483:end
% glossary-block:B01484:begin
\needspace{3\baselineskip}
\textbf{reply} — \textbf{andswarian} (v.). answer; reply. Dative addressee; the reply may be given as direct speech.\par
% glossary-block:B01484:end
% glossary-block:B01485:begin
\needspace{3\baselineskip}
\textbf{report} — \textbf{hlīsa} (n.). report; account heard.\par
% glossary-block:B01485:end
% glossary-block:B01486:begin
\needspace{3\baselineskip}
\textbf{repose} — \textbf{geręstan} (v.). rest; repose. \textbf{ręst} (n.). rest; repose.\par
% glossary-block:B01486:end
% glossary-block:B01487:begin
\needspace{3\baselineskip}
\textbf{request} — \textbf{biddan} (v.). ask; request. Genitive thing requested; an accusative person may be added. Accusative person beseeched; ðæt + subjunctive gives the requested act. Prayer to God with gebiddan is a different entry. \textbf{giernan} (v.). desire; request; woo. Genitive desired thing/person in the selected frames, including the maiden wooed.\par
% glossary-block:B01487:end
% glossary-block:B01488:begin
\needspace{3\baselineskip}
\textbf{require} — \textbf{beðurfan} (v.). need; require. Genitive thing needed. Genitive thing needed; a relative particle may represent that governed relation.\par
% glossary-block:B01488:end
% glossary-block:B01489:begin
\needspace{3\baselineskip}
\textbf{requite} — \textbf{forgieldan} (v.). pay; repay; requite. Dative recipient, accusative thing repaid in the selected frame.\par
% glossary-block:B01489:end
% glossary-block:B01490:begin
\needspace{3\baselineskip}
\textbf{rescue} — \textbf{hǣlu} (n.). safety; rescue.\par
% glossary-block:B01490:end
% glossary-block:B01491:begin
\needspace{3\baselineskip}
\textbf{resist} — \textbf{wiðstandan} (v.). withstand; resist. Dative of what is resisted.\par
% glossary-block:B01491:end
% glossary-block:B01492:begin
\needspace{3\baselineskip}
\textbf{rest} — \textbf{geręstan} (v.). rest; repose. \textbf{ręst} (n.). rest; repose.\par
% glossary-block:B01492:end
% glossary-block:B01493:begin
\needspace{3\baselineskip}
\textbf{rest from} — Fram + dative expresses from the work or activity in the supplied rest-from construction; the verb of resting must also be supplied.\par
% glossary-block:B01493:end
% glossary-block:B01494:begin
\needspace{3\baselineskip}
\textbf{restore} — \textbf{āgiefan} (v.). give back; restore.\par
% glossary-block:B01494:end
% glossary-block:B01495:begin
\needspace{3\baselineskip}
\textbf{result} — \textbf{weorðan} (v.). come about; result. Supplied resultative expression.\par
% glossary-block:B01495:end
% glossary-block:B01496:begin
\needspace{3\baselineskip}
\textbf{retainer} — \textbf{ðegn} (n.). vassal; retainer; thane.\par
% glossary-block:B01496:end
% glossary-block:B01497:begin
\needspace{3\baselineskip}
\textbf{retreat} — \textbf{dīegelnes} (n.). retreat; seclusion. Supplied plural after of.\par
% glossary-block:B01497:end
% glossary-block:B01498:begin
\needspace{3\baselineskip}
\textbf{return} — \textbf{ęfthweorfan} (v.). return; turn back. \textbf{gewęndan} (v.). return; turn. \textbf{hweorfan} (v.). return; turn back. Supplied source preterite. \textbf{węndan} (v.). turn; return. Supplied source form.\par
% glossary-block:B01498:end
% glossary-block:B01499:begin
\needspace{3\baselineskip}
\textbf{reveal} — \textbf{ætīewan} (v.). show; display; reveal. \textbf{gesweotolian} (v.). reveal; make clear. Supplied source form. make manifest; reveal.\par
% glossary-block:B01499:end
% glossary-block:B01500:begin
\needspace{3\baselineskip}
\textbf{reward} — \textbf{āgyldan} (v.). reward; repay. C06 corrects Cook's āgyltan to strong III āgyldan; supplied contracted present āgylt and dative recipient. \textbf{edlēan} (n.). reward; recompense. \textbf{mēd} (n.).\par
% glossary-block:B01500:end
% glossary-block:B01501:begin
\needspace{3\baselineskip}
\textbf{rich} — \textbf{rīce} (adj.). rich; wealthy. Includes substantival rich people in the assigned frame. \textbf{welig} (adj.). rich; wealthy. \textbf{wlanc} (adj.). proud; wealthy. Contextually the rich man addressed. \textbf{welig} (adj.), \textbf{rīce} (adj.). Welig and rīce can mean rich/wealthy; rīce also occurs substantivally for rich people. Inflect the selected adjective for its role.\par
% glossary-block:B01501:end
% glossary-block:B01502:begin
\needspace{3\baselineskip}
\textbf{riches} — \textbf{wela} (n.). riches; wealth.\par
% glossary-block:B01502:end
% glossary-block:B01503:begin
\needspace{3\baselineskip}
\textbf{rich in} — \textbf{beran} (v.). rich in; productive of. Supplied berende on + dative ore expression; participial adjective. Berende on + dative is the supplied rich-in/productive-of expression. The meaning belongs to the participial expression, not to on alone. \textbf{welig} (adj.). rich in; abounding in. on + dative thing abounded in. welig on + dative of what is abundant.\par
% glossary-block:B01503:end
% glossary-block:B01504:begin
\needspace{3\baselineskip}
\textbf{rich people} — \textbf{rīce} (adj.). rich; wealthy. Includes substantival rich people in the assigned frame.\par
% glossary-block:B01504:end
% glossary-block:B01505:begin
\needspace{3\baselineskip}
\textbf{ride} — \textbf{rīdan} (v.). Intransitive in the introduced frame.\par
% glossary-block:B01505:end
% glossary-block:B01506:begin
\needspace{3\baselineskip}
\textbf{right} — \textbf{riht} (adj.). right; just. \textbf{swīðra} (adj.). Right-hand/side expression, distinct from stronger.\par
% glossary-block:B01506:end
% glossary-block:B01507:begin
\needspace{3\baselineskip}
\textbf{righteous} — \textbf{rihtwīs} (adj.).\par
% glossary-block:B01507:end
% glossary-block:B01508:begin
\needspace{3\baselineskip}
\textbf{righteousness} — \textbf{rihtwīsnes} (n.).\par
% glossary-block:B01508:end
% glossary-block:B01509:begin
\needspace{3\baselineskip}
\textbf{righteous people} — \textbf{rihtwīs} (adj.). Substantival adjective.\par
% glossary-block:B01509:end
% glossary-block:B01510:begin
\needspace{3\baselineskip}
\textbf{right-hand side} — \textbf{swīðra} (adj.). right. Right-hand/side expression, distinct from stronger.\par
% glossary-block:B01510:end
% glossary-block:B01511:begin
\needspace{3\baselineskip}
\textbf{ripe} — \textbf{rīpe} (adj.).\par
% glossary-block:B01511:end
% glossary-block:B01512:begin
\needspace{3\baselineskip}
\textbf{rise} — \textbf{rīsan} (v.). Earlier taught rīsan supplies this command; prefixed ārīsan is separately used later.\par
% glossary-block:B01512:end
% glossary-block:B01513:begin
\needspace{3\baselineskip}
\textbf{rise up} — \textbf{rīsan} (v.). Intransitive rīsan, optionally with ūp; ongēan + accusative marks opposition. \textbf{ūpāstīgan} (v.). rise up; ascend.\par
% glossary-block:B01513:end
% glossary-block:B01514:begin
\needspace{3\baselineskip}
\textbf{rising} — \textbf{ūpgang} (n.). rising; sunrise.\par
% glossary-block:B01514:end
% glossary-block:B01515:begin
\needspace{3\baselineskip}
\textbf{risk} — \textbf{pliht} (n.). peril; risk.\par
% glossary-block:B01515:end
% glossary-block:B01516:begin
\needspace{3\baselineskip}
\textbf{road} — \textbf{weg} (n.). Weg means road, path or way.\par
% glossary-block:B01516:end
% glossary-block:B01517:begin
\needspace{3\baselineskip}
\textbf{Roman} — \textbf{Rōmanisc} (adj.).\par
% glossary-block:B01517:end
% glossary-block:B01518:begin
\needspace{3\baselineskip}
\textbf{Romans} — \textbf{Rōmane} (ethnonym).\par
% glossary-block:B01518:end
% glossary-block:B01519:begin
\needspace{3\baselineskip}
\textbf{Rome} — \textbf{Rome} (proper name).\par
% glossary-block:B01519:end
% glossary-block:B01520:begin
\needspace{3\baselineskip}
\textbf{roof} — \textbf{hrōf} (n.).  \textbf{ðæce} (n.).\par
% glossary-block:B01520:end
% glossary-block:B01521:begin
\needspace{3\baselineskip}
\textbf{room} — \textbf{būr} (n.). chamber; room.\par
% glossary-block:B01521:end
% glossary-block:B01522:begin
\needspace{3\baselineskip}
\textbf{root} — \textbf{wyrttruma} (n.).\par
% glossary-block:B01522:end
% glossary-block:B01523:begin
\needspace{3\baselineskip}
\textbf{round about} — \textbf{ymbhwyrft} (n.). circuit. Whole on ymbhwyrfte expression means all around; supplied dependent landmark. Supplied on ymbhwyrfte + genitive landmark in this exercise.\par
% glossary-block:B01523:end
% glossary-block:B01524:begin
\needspace{3\baselineskip}
\textbf{row} — \textbf{rōwan} (v.). Intransitive; the assigned route is expressed with ofer.\par
% glossary-block:B01524:end
% glossary-block:B01525:begin
\needspace{3\baselineskip}
\textbf{royal seat} — \textbf{cynesętl} (n.). royal seat; throne.\par
% glossary-block:B01525:end
% glossary-block:B01526:begin
\needspace{3\baselineskip}
\textbf{ruin} — \textbf{forierman} (v.). ruin; reduce to misery. \textbf{lor} (n.). destruction; ruin.\par
% glossary-block:B01526:end
% glossary-block:B01527:begin
\needspace{3\baselineskip}
\textbf{rule} — \textbf{geweald} (n.). control; rule; dominion. Ordinary singular; on gewealde is the taught control phrase. \textbf{rīcsian} (v.). rule; have dominion. Used absolutely in the taught frame; no new object government. rīcsian can be used absolutely; English have/bear need no independent translation. \textbf{rīce} (n.). Genitive rīces after geswīcan in ceased to bear rule. Do not insert a literal bear verb.\par
% glossary-block:B01527:end
% glossary-block:B01528:begin
\needspace{3\baselineskip}
\textbf{rule over} — \textbf{gīeman} (v.). Genitive of what is ruled; no additional over preposition. gīeman governs genitive of what is ruled; English over supplies no additional preposition.\par
% glossary-block:B01528:end
% glossary-block:B01529:begin
\needspace{3\baselineskip}
\textbf{Ruler} — \textbf{wealdend} (n.).\par
% glossary-block:B01529:end
% glossary-block:B01530:begin
\needspace{3\baselineskip}
\textbf{run} — \textbf{iernan} (v.). Intransitive; geond + accusative path. A perception construction may make its actor accusative.\par
% glossary-block:B01530:end
% book-navigation: standalone bookmark replaced by volume hierarchy
% glossary-block:B01531:begin
\needspace{5\baselineskip}
\subsection*{S}\BookHeading{nav-glossary-english-0020}{subsubsection}{S}{0}
% glossary-block:B01531:end
% glossary-block:B01532:begin
\textbf{sack} — \textbf{sacc} (n.).\par
% glossary-block:B01532:end
% glossary-block:B01533:begin
\needspace{3\baselineskip}
\textbf{safety} — \textbf{hǣlu} (n.). safety; rescue.\par
% glossary-block:B01533:end
% glossary-block:B01534:begin
\needspace{3\baselineskip}
\textbf{Saint} — \textbf{Sanctus} (foreign adjective). Saint; holy. Latin word in the source title.\par
% glossary-block:B01534:end
% glossary-block:B01535:begin
\needspace{3\baselineskip}
\textbf{saint} — \textbf{hālga} (n.). saint; holy one.\par
% glossary-block:B01535:end
% glossary-block:B01536:begin
\needspace{3\baselineskip}
\textbf{saints} — \textbf{hālig} (adj.). holy ones; saints. Substantival adjective in the source.\par
% glossary-block:B01536:end
% glossary-block:B01537:begin
\needspace{3\baselineskip}
\textbf{sake} — \textbf{intinga} (n.). cause; sake (in for … intingan). Within supplied for intingan expressions.\par
% glossary-block:B01537:end
% glossary-block:B01538:begin
\needspace{3\baselineskip}
\textbf{salute} — \textbf{grētan} (v.). greet; take leave of. Accusative person/place; greeting and farewell depend on context. Greeting with accusative person; the same verb also has farewell/leave-taking uses.\par
% glossary-block:B01538:end
% glossary-block:B01539:begin
\needspace{3\baselineskip}
\textbf{salvation} — \textbf{hǣlu} (n.).\par
% glossary-block:B01539:end
% glossary-block:B01540:begin
\needspace{3\baselineskip}
\textbf{same} — \textbf{ilca} (demonstrative adjective).\par
% glossary-block:B01540:end
% glossary-block:B01541:begin
\needspace{3\baselineskip}
\textbf{save} — \textbf{hǣlan} (v.).  \textbf{nęrian} (v.).\par
% glossary-block:B01541:end
% glossary-block:B01542:begin
\needspace{3\baselineskip}
\textbf{Savior} — \textbf{Hǣlend} (n.). Saviour; Jesus. Proper title.\par
% glossary-block:B01542:end
% glossary-block:B01543:begin
\needspace{3\baselineskip}
\textbf{Saviour} — \textbf{Hǣlend} (n.). Saviour; Jesus. Proper title.\par
% glossary-block:B01543:end
% glossary-block:B01544:begin
\needspace{3\baselineskip}
\textbf{say} — \textbf{cweðan} (v.). Source genitive plural present participle, of people saying. Supplied preterite said. \textbf{gecweðan} (v.). speak; say. \textbf{sęcgan} (v.). tell; say. Dative addressee and a reported statement or question as content. Supplied source preterite. say; tell. Person addressed is dative; direct or reported content is separate. \textbf{sprecan} (v.). This task explicitly selects sprecan for say; tō marks the addressee.\par
% glossary-block:B01544:end
% glossary-block:B01545:begin
\needspace{3\baselineskip}
\textbf{saying} — \textbf{sprǣc} (n.). saying; speech.\par
% glossary-block:B01545:end
% glossary-block:B01546:begin
\needspace{3\baselineskip}
\textbf{say to} — In supplied cwæð him tō, cwæð expresses said and postposed tō marks the addressee. Tō alone does not express say. Tō marks the dative addressee with the supplied verb of speaking; the verb supplies say, not tō alone.\par
% glossary-block:B01546:end
% glossary-block:B01547:begin
\needspace{3\baselineskip}
\textbf{scorch} — \textbf{forbærnan} (v.). parch; scorch.\par
% glossary-block:B01547:end
% glossary-block:B01548:begin
\needspace{3\baselineskip}
\textbf{Scots} — \textbf{Scottas} (ethnonym).\par
% glossary-block:B01548:end
% glossary-block:B01549:begin
\needspace{3\baselineskip}
\textbf{sea} — \textbf{sǣ} (n.).\par
% glossary-block:B01549:end
% glossary-block:B01550:begin
\needspace{3\baselineskip}
\textbf{seal} — \textbf{seolh} (n.).\par
% glossary-block:B01550:end
% glossary-block:B01551:begin
\needspace{3\baselineskip}
\textbf{sea-shore} — \textbf{strand} (n.). strand; sea-shore.\par
% glossary-block:B01551:end
% glossary-block:B01552:begin
\needspace{3\baselineskip}
\textbf{season} — \textbf{tīd} (n.). season; time.\par
% glossary-block:B01552:end
% glossary-block:B01553:begin
\needspace{3\baselineskip}
\textbf{seat} — \textbf{setl} (n.). seat; throne.\par
% glossary-block:B01553:end
% glossary-block:B01554:begin
\needspace{3\baselineskip}
\textbf{seclusion} — \textbf{dīegelnes} (n.). retreat; seclusion. Supplied plural after of.\par
% glossary-block:B01554:end
% glossary-block:B01555:begin
\needspace{3\baselineskip}
\textbf{second} — \textbf{ōðer} (adj.; ordinal).\par
% glossary-block:B01555:end
% glossary-block:B01556:begin
\needspace{3\baselineskip}
\textbf{second time} — Ōðre sīðe is the supplied instrumental singular phrase a second time, without a separate indefinite article.\par
% glossary-block:B01556:end
% glossary-block:B01557:begin
\needspace{3\baselineskip}
\textbf{see} — \textbf{gesēon} (v.). Absolute use, with no expressed object. Accusative perceived actor and plain infinitive. With a content clause stating the circumstance seen. With the assigned clause describing how a state of affairs stands. \textbf{sēon} (v.).\par
% glossary-block:B01557:end
% glossary-block:B01558:begin
\needspace{3\baselineskip}
\textbf{seed} — \textbf{sǣd} (n.). seed; descendants. Collective descendants, distinct from literal seed.\par
% glossary-block:B01558:end
% glossary-block:B01559:begin
\needspace{3\baselineskip}
\textbf{seedtime} — \textbf{sǣdtīma} (n.). seedtime; sowing season.\par
% glossary-block:B01559:end
% glossary-block:B01560:begin
\needspace{3\baselineskip}
\textbf{see how} — \textbf{gesēon} (v.). see. With the assigned clause describing how a state of affairs stands.\par
% glossary-block:B01560:end
% glossary-block:B01561:begin
\needspace{3\baselineskip}
\textbf{seek} — \textbf{sēcan} (v.). Accusative person or thing sought. sēcan with accusative sought thing, no extra word for English for/out.\par
% glossary-block:B01561:end
% glossary-block:B01562:begin
\needspace{3\baselineskip}
\textbf{seek in marriage} — \textbf{giernan} (v.). desire; request; woo. Genitive desired thing/person in the selected frames, including the maiden wooed. Genitive person wooed; keep the marriage-seeking sense distinct from a general want.\par
% glossary-block:B01562:end
% glossary-block:B01563:begin
\needspace{3\baselineskip}
\textbf{seek out} — \textbf{sēcan} (v.). sēcan with accusative sought thing, no extra word for English for/out.\par
% glossary-block:B01563:end
% glossary-block:B01564:begin
\needspace{3\baselineskip}
\textbf{seem} — \textbf{ðyncan} (v.). Nominative thing that seems, dative experiencer and agreeing predicate. Dative experiencer; the thing that seems is nominative. Supplied past source form. Wesan is rendered seem only within the supplied impersonal expression it seemed to the king. Retain its full dative-experiencer construction; ordinary wesan means be.\par
% glossary-block:B01564:end
% glossary-block:B01565:begin
\needspace{3\baselineskip}
\textbf{seem good to} — \textbf{ðyncan} (v.). seem. Nominative thing that seems, dative experiencer and agreeing predicate. Subject thing nominative, dative experiencer, agreeing predicate gōd.\par
% glossary-block:B01565:end
% glossary-block:B01566:begin
\needspace{3\baselineskip}
\textbf{see someone do something} — \textbf{gesēon} (v.). see. Accusative perceived actor and plain infinitive.\par
% glossary-block:B01566:end
% glossary-block:B01567:begin
\needspace{3\baselineskip}
\textbf{see that} — \textbf{gesēon} (v.). see. With a content clause stating the circumstance seen.\par
% glossary-block:B01567:end
% glossary-block:B01568:begin
\needspace{3\baselineskip}
\textbf{seize} — \textbf{fōn} (v.). seize; catch. Accusative patient in the selected frame. Accusative person overtaken in these tasks; the source's bracketed caught is the route.\par
% glossary-block:B01568:end
% glossary-block:B01569:begin
\needspace{3\baselineskip}
\textbf{select} — \textbf{cēosan} (v.). Accusative person selected.\par
% glossary-block:B01569:end
% glossary-block:B01570:begin
\needspace{3\baselineskip}
\textbf{self} — \textbf{self} (adj.). self; the same person or thing. Agreeing emphatic/reflexive self; ordinary reflexive personal pronouns do not always require it.\par
% glossary-block:B01570:end
% glossary-block:B01571:begin
\needspace{3\baselineskip}
\textbf{sell} — \textbf{cīepan} (v.). Accusative goods sold. \textbf{sęllan} (v.). Accusative thing sold. Sęllan means sell in the selected trading contexts; the thing sold is accusative.\par
% glossary-block:B01571:end
% glossary-block:B01572:begin
\needspace{3\baselineskip}
\textbf{sell into} — \textbf{on} (prep.). into. Destination of sale in the supplied expression.\par
% glossary-block:B01572:end
% glossary-block:B01573:begin
\needspace{3\baselineskip}
\textbf{send} — \textbf{āsęndan} (v.).  \textbf{onsęndan} (v.).  \textbf{sęndan} (v.). Accusative person or thing sent, dative recipient.\par
% glossary-block:B01573:end
% glossary-block:B01574:begin
\needspace{3\baselineskip}
\textbf{sense} — \textbf{andgiet} (n.). meaning; sense.\par
% glossary-block:B01574:end
% glossary-block:B01575:begin
\needspace{3\baselineskip}
\textbf{sense for sense} — \textbf{of} (prep.). Literal correspondence in the supplied translation expression.\par
% glossary-block:B01575:end
% glossary-block:B01576:begin
\needspace{3\baselineskip}
\textbf{sent back} — \textbf{gesęndan} (v.).\par
% glossary-block:B01576:end
% glossary-block:B01577:begin
\needspace{3\baselineskip}
\textbf{separate} — \textbf{gedǣlan} (v.). divide; separate. Preserve the assigned separation frame. \textbf{tōdǣlen / tōdǣlden (source forms)} (verb; supplied inflected forms). separate; distinguish. Source forms tōdǣlen/tōdǣlden are retained; no silently invented citation spelling. \textbf{tōsceadan} (v.). separate; divide. \textbf{tōtwǣman} (v.). divide; separate. Waters divided from waters in the supplied frame. Accusative beings separated; fram + dative other party.\par
% glossary-block:B01577:end
% glossary-block:B01578:begin
\needspace{3\baselineskip}
\textbf{serpent} — \textbf{nædre} (n.).\par
% glossary-block:B01578:end
% glossary-block:B01579:begin
\needspace{3\baselineskip}
\textbf{servant} — \textbf{cniht} (n.).  \textbf{ðēow} (n.). servant; enslaved person. \textbf{ðēowa} (n.).  \textbf{cniht} (n.), \textbf{ðēowa} (n.), \textbf{ðēow} (n.). Cniht, weak ðēowa and strong ðēow can denote a servant. Use the particular noun required by the exercise and its own declension.\par
% glossary-block:B01579:end
% glossary-block:B01580:begin
\needspace{3\baselineskip}
\textbf{servants} — \textbf{ðēowa manna} (noun phrase). enslaved men; servants. Whole supplied phrase ðēowa manna.\par
% glossary-block:B01580:end
% glossary-block:B01581:begin
\needspace{3\baselineskip}
\textbf{serve} — \textbf{gehīersumian} (v.). serve; obey. Supplied source form and government. \textbf{ðēnian} (v.). serve; attend. Supplied preterite source form.\par
% glossary-block:B01581:end
% glossary-block:B01582:begin
\needspace{3\baselineskip}
\textbf{service} — \textbf{ðēnung} (n.). service; attendance. \textbf{ðēowdōm} (n.).\par
% glossary-block:B01582:end
% glossary-block:B01583:begin
\needspace{3\baselineskip}
\textbf{set} — \textbf{āsęttan} (v.). set; place. \textbf{gelōgian} (v.). place; set. \textbf{gesęttan} (v.). set; place. Supplied gesętte. set; place; establish. \textbf{lęcgean} (v.). lay; place. Supplied source form. \textbf{sęttan} (v.). set; place. Place an accusative object; destination is a separate phrase.\par
% glossary-block:B01583:end
% glossary-block:B01584:begin
\needspace{3\baselineskip}
\textbf{set aside} — \textbf{forlǣtan} (v.). Supplied source sense.\par
% glossary-block:B01584:end
% glossary-block:B01585:begin
\needspace{3\baselineskip}
\textbf{set as lord} — \textbf{gesęttan} (v.). appoint as lord. Supplied frame with person appointed, tō + dative role, and dative people governed.\par
% glossary-block:B01585:end
% glossary-block:B01586:begin
\needspace{3\baselineskip}
\textbf{set forth} — \textbf{lęcgean} (v.). lęcgean + accusative thing + forð.\par
% glossary-block:B01586:end
% glossary-block:B01587:begin
\needspace{3\baselineskip}
\textbf{set in motion} — \textbf{astyrian} (v.). move; set in motion. Supplied source sense.\par
% glossary-block:B01587:end
% glossary-block:B01588:begin
\needspace{3\baselineskip}
\textbf{settle} — \textbf{gestillan} (v.). quell; settle. Accusative disturbance settled; distinguish intransitive become still.\par
% glossary-block:B01588:end
% glossary-block:B01589:begin
\needspace{3\baselineskip}
\textbf{seven} — \textbf{seofon} (cardinal numeral). Cardinal 7.\par
% glossary-block:B01589:end
% glossary-block:B01590:begin
\needspace{3\baselineskip}
\textbf{seventh} — \textbf{seofoða} (ordinal).\par
% glossary-block:B01590:end
% glossary-block:B01591:begin
\needspace{3\baselineskip}
\textbf{shake} — \textbf{geswęngan} (v.). toss; shake. Supplied participial source use. \textbf{gewǣgan} (v.). buffet; shake. Supplied participial source use.\par
% glossary-block:B01591:end
% glossary-block:B01592:begin
\needspace{3\baselineskip}
\textbf{shall} — Supplied gīet sculon fīf means that five are yet to come, with be expressly omitted. This bounded ellipsis is not a general rule that English shall always takes sculan.\par
% glossary-block:B01592:end
% glossary-block:B01593:begin
\needspace{3\baselineskip}
\textbf{shall (future reference)} — Present morphology can carry future reference; bēo is used in its bounded expected-state frame. Do not add willan unless volition or a stated source target requires it.\par
% glossary-block:B01593:end
% glossary-block:B01594:begin
\needspace{3\baselineskip}
\textbf{shape} — \textbf{scieppan} (v.). create; shape.\par
% glossary-block:B01594:end
% glossary-block:B01595:begin
\needspace{3\baselineskip}
\textbf{shared} — \textbf{gemǣne} (adj.). common; shared.\par
% glossary-block:B01595:end
% glossary-block:B01596:begin
\needspace{3\baselineskip}
\textbf{she} — Hēo is feminine singular nominative. Third-person pronouns follow the grammatical gender of their referent. Grammatical gender and antecedent determine reference; it may be represented by masculine/feminine forms for a gendered inanimate noun.\par
% glossary-block:B01596:end
% glossary-block:B01597:begin
\needspace{3\baselineskip}
\textbf{Sheba} — \textbf{Saba} (proper name).\par
% glossary-block:B01597:end
% glossary-block:B01598:begin
\needspace{3\baselineskip}
\textbf{sheep} — \textbf{scēap} (n.).\par
% glossary-block:B01598:end
% glossary-block:B01599:begin
\needspace{3\baselineskip}
\textbf{shepherd} — \textbf{scēaphierde} (n.).\par
% glossary-block:B01599:end
% glossary-block:B01600:begin
\needspace{3\baselineskip}
\textbf{shield} — \textbf{gescieldan} (v.). shield; protect.\par
% glossary-block:B01600:end
% glossary-block:B01601:begin
\needspace{3\baselineskip}
\textbf{shine} — \textbf{scīnan} (v.).\par
% glossary-block:B01601:end
% glossary-block:B01602:begin
\needspace{3\baselineskip}
\textbf{shining} — \textbf{scīnan} (v.).\par
% glossary-block:B01602:end
% glossary-block:B01603:begin
\needspace{3\baselineskip}
\textbf{ship} — \textbf{scip} (n.).\par
% glossary-block:B01603:end
% glossary-block:B01604:begin
\needspace{3\baselineskip}
\textbf{shipwreck} — \textbf{forlidennes} (n.).\par
% glossary-block:B01604:end
% glossary-block:B01605:begin
\needspace{3\baselineskip}
\textbf{shipwrecked} — \textbf{forliden} (adj.).\par
% glossary-block:B01605:end
% glossary-block:B01606:begin
\needspace{3\baselineskip}
\textbf{shoemaker} — \textbf{scēowyrhta} (n.).\par
% glossary-block:B01606:end
% glossary-block:B01607:begin
\needspace{3\baselineskip}
\textbf{shore} — \textbf{strand} (n.). Strand means shore or sea-shore in the assigned noun and translation tasks.\par
% glossary-block:B01607:end
% glossary-block:B01608:begin
\needspace{3\baselineskip}
\textbf{short} — \textbf{medmicel} (adj.). short; brief.\par
% glossary-block:B01608:end
% glossary-block:B01609:begin
\needspace{3\baselineskip}
\textbf{should} — In the supplied interpretive question, sculan expresses what events should/were to mean. Retain that reported or prospective interpretation rather than substituting the ordinary obligation sense. Select intention, obligation, reported prospect or counterfactual context; English spelling alone does not determine tense, mood or an OE modal.\par
% glossary-block:B01609:end
% glossary-block:B01610:begin
\needspace{3\baselineskip}
\textbf{show} — \textbf{ætīewan} (v.). show; display; reveal. \textbf{ge-cȳðan} (v.).  \textbf{getǣcan} (v.). show; direct. Supplied source form. \textbf{tǣcan} (v.). teach; show.\par
% glossary-block:B01610:end
% glossary-block:B01611:begin
\needspace{3\baselineskip}
\textbf{sick} — \textbf{untrum} (adj.).\par
% glossary-block:B01611:end
% glossary-block:B01612:begin
\needspace{3\baselineskip}
\textbf{sick people} — \textbf{untrum} (adj.). Substantival adjective.\par
% glossary-block:B01612:end
% glossary-block:B01613:begin
\needspace{3\baselineskip}
\textbf{side} — \textbf{hand} (n.). hand; side. Side in right-/left-hand expressions.\par
% glossary-block:B01613:end
% glossary-block:B01614:begin
\needspace{3\baselineskip}
\textbf{sight} — \textbf{gesihð} (n.). sight; gaze. Supplied accusative singular.\par
% glossary-block:B01614:end
% glossary-block:B01615:begin
\needspace{3\baselineskip}
\textbf{sign} — \textbf{tācen} (n.).\par
% glossary-block:B01615:end
% glossary-block:B01616:begin
\needspace{3\baselineskip}
\textbf{signal} — The supplied expression slōg tācen means made a signal; tācen is its object. Slēan alone has the separate strike sense.\par
% glossary-block:B01616:end
% glossary-block:B01617:begin
\needspace{3\baselineskip}
\textbf{silence} — \textbf{swīge} (n.).\par
% glossary-block:B01617:end
% glossary-block:B01618:begin
\needspace{3\baselineskip}
\textbf{silk} — \textbf{sīde} (n.).\par
% glossary-block:B01618:end
% glossary-block:B01619:begin
\needspace{3\baselineskip}
\textbf{silver} — \textbf{seolfor} (n.). Metal noun. \textbf{seolfren} (adj.). silver; made of silver. Material adjective, distinct from the noun seolfor.\par
% glossary-block:B01619:end
% glossary-block:B01620:begin
\needspace{3\baselineskip}
\textbf{similarly} — \textbf{gelīce} (adv.). likewise; similarly.\par
% glossary-block:B01620:end
% glossary-block:B01621:begin
\needspace{3\baselineskip}
\textbf{sin} — \textbf{synn} (n.).\par
% glossary-block:B01621:end
% glossary-block:B01622:begin
\needspace{3\baselineskip}
\textbf{since} — \textbf{siððan} (adv.; conj.). since; after; afterwards. Conjunction since/after versus adverb afterwards.\par
% glossary-block:B01622:end
% glossary-block:B01623:begin
\needspace{3\baselineskip}
\textbf{sinful} — \textbf{synfull} (adj.).\par
% glossary-block:B01623:end
% glossary-block:B01624:begin
\needspace{3\baselineskip}
\textbf{sing} — \textbf{singan} (v.). Intransitive in the introduced frame. Source genitive plural present participle, of people singing.\par
% glossary-block:B01624:end
% glossary-block:B01625:begin
\needspace{3\baselineskip}
\textbf{single} — \textbf{ān} (numeral; adj.). one; single. Cardinal one/single and the agreeing adjective alone are distinct uses of ān. Preserve the supplied form where the exercise provides one.\par
% glossary-block:B01625:end
% glossary-block:B01626:begin
\needspace{3\baselineskip}
\textbf{sir} — \textbf{lēof} (adj.). Supplied respectful address lēof. Supplied respectful lēof address; the older gloss dear alone does not give this lookup.\par
% glossary-block:B01626:end
% glossary-block:B01627:begin
\needspace{3\baselineskip}
\textbf{sister} — \textbf{sweostor} (n.).\par
% glossary-block:B01627:end
% glossary-block:B01628:begin
\needspace{3\baselineskip}
\textbf{sit} — \textbf{sittan} (v.).\par
% glossary-block:B01628:end
% glossary-block:B01629:begin
\needspace{3\baselineskip}
\textbf{sit around} — \textbf{ymbsittan} (v.). Supplied participle referring to people sitting around.\par
% glossary-block:B01629:end
% glossary-block:B01630:begin
\needspace{3\baselineskip}
\textbf{situated} — \textbf{gesęttan} (v.). placed; situated. Supplied participial complement of understood us; not the identical-looking finite preterite.\par
% glossary-block:B01630:end
% glossary-block:B01631:begin
\needspace{3\baselineskip}
\textbf{six} — \textbf{siex} (cardinal numeral). Cardinal 6.\par
% glossary-block:B01631:end
% glossary-block:B01632:begin
\needspace{3\baselineskip}
\textbf{sixth} — \textbf{siexta} (ordinal numeral).\par
% glossary-block:B01632:end
% glossary-block:B01633:begin
\needspace{3\baselineskip}
\textbf{sixty} — \textbf{siextig} (cardinal numeral). Within the supplied 365-days expression.\par
% glossary-block:B01633:end
% glossary-block:B01634:begin
\needspace{3\baselineskip}
\textbf{skill} — \textbf{cræft} (n.). skill; accomplishment. \textbf{snelnes} (n.). skill; agility.\par
% glossary-block:B01634:end
% glossary-block:B01635:begin
\needspace{3\baselineskip}
\textbf{skilled} — \textbf{gelǣred} (adj.). educated; trained; skilled.\par
% glossary-block:B01635:end
% glossary-block:B01636:begin
\needspace{3\baselineskip}
\textbf{sky} — \textbf{heofon} (n.). heaven; sky. \textbf{heofone} (n.). heaven; sky. Weak feminine lexical alternative in the source.\par
% glossary-block:B01636:end
% glossary-block:B01637:begin
\needspace{3\baselineskip}
\textbf{slaughter} — \textbf{wæl} (n.).\par
% glossary-block:B01637:end
% glossary-block:B01638:begin
\needspace{3\baselineskip}
\textbf{slay} — \textbf{ofslēan} (v.). slay; kill. \textbf{ongeslēan} (v.). inflict slaughter; slay. \textbf{slēan} (v.). slay; kill. Accusative person killed.\par
% glossary-block:B01638:end
% glossary-block:B01639:begin
\needspace{3\baselineskip}
\textbf{small} — \textbf{lȳtel} (adj.). little; small.\par
% glossary-block:B01639:end
% glossary-block:B01640:begin
\needspace{3\baselineskip}
\textbf{smaller} — \textbf{lǣssa} (adj.).\par
% glossary-block:B01640:end
% glossary-block:B01641:begin
\needspace{3\baselineskip}
\textbf{smallest} — \textbf{lǣst} (adj.).\par
% glossary-block:B01641:end
% glossary-block:B01642:begin
\needspace{3\baselineskip}
\textbf{smite} — \textbf{geslēan} (v.). strike; smite. Supplied past/participial source form. \textbf{slēan} (v.). strike; smite. strike; smite. Objectless striking in the assigned hammer sentences. \textbf{slēan} (v.), \textbf{bēatan} (v.). Slēan means strike/smite and can be objectless in the hammer sentences. Bēatan means beat/strike and takes an accusative object; the explicit bēatan target must be retained.\par
% glossary-block:B01642:end
% glossary-block:B01643:begin
\needspace{3\baselineskip}
\textbf{smithy} — \textbf{smiððe} (n.).\par
% glossary-block:B01643:end
% glossary-block:B01644:begin
\needspace{3\baselineskip}
\textbf{smooth} — \textbf{smylte} (adj.). calm; smooth.\par
% glossary-block:B01644:end
% glossary-block:B01645:begin
\needspace{3\baselineskip}
\textbf{snake} — \textbf{næddre} (n.). snake; adder. Supplied weak feminine plural.\par
% glossary-block:B01645:end
% glossary-block:B01646:begin
\needspace{3\baselineskip}
\textbf{snow} — \textbf{snāw} (n.).\par
% glossary-block:B01646:end
% glossary-block:B01647:begin
\needspace{3\baselineskip}
\textbf{so} — Tō ðan … ðæt is the supplied degree-and-result construction so/to that degree … that. Keep the whole construction. \textbf{swā} (adverb; conjunction). so; thus.\par
% glossary-block:B01647:end
% glossary-block:B01648:begin
\needspace{3\baselineskip}
\textbf{sojourn} — \textbf{wician} (v.). sojourn; lodge. Intransitive; mid + dative person with whom one lodges.\par
% glossary-block:B01648:end
% glossary-block:B01649:begin
\needspace{3\baselineskip}
\textbf{soldier} — \textbf{cęmpa} (n.).\par
% glossary-block:B01649:end
% glossary-block:B01650:begin
\needspace{3\baselineskip}
\textbf{solemnly charge} — \textbf{hālsian} (v.). conjure; solemnly charge. Accusative person; ðurh + accusative oath standard; a requested ðæt clause may follow.\par
% glossary-block:B01650:end
% glossary-block:B01651:begin
\needspace{3\baselineskip}
\textbf{Solomon} — \textbf{Salomon} (proper name).\par
% glossary-block:B01651:end
% glossary-block:B01652:begin
\needspace{3\baselineskip}
\textbf{some} — \textbf{hwilc} (interrogative/indefinite adj. or pron.). any; some. Indefinite adjective, distinct from an interrogative. \textbf{sum} (indefinite adjective/pronoun). Members selected from a group; genitive whole. a certain; some. \textbf{sum} (indefinite adjective/pronoun), \textbf{hwilc} (interrogative/indefinite adj. or pron.). Sum means some/a certain; partitive sum selects members from a genitive whole. Indefinite hwilc has the separate supplied some/any use.\par
% glossary-block:B01652:end
% glossary-block:B01653:begin
\needspace{3\baselineskip}
\textbf{some of} — Partitive sum selects some members from a genitive whole; English of is expressed by that genitive, without an extra preposition.\par
% glossary-block:B01653:end
% glossary-block:B01654:begin
\needspace{3\baselineskip}
\textbf{something} — \textbf{hwæthwega} (indefinite pronoun/substantive). Supplied local indefinite form; no unrestricted new paradigm. \textbf{hwæthwugu} (indefinite). Supplied hwæthwugu swilces.\par
% glossary-block:B01654:end
% glossary-block:B01655:begin
\needspace{3\baselineskip}
\textbf{something of that sort} — \textbf{hwæthwugu} (indefinite). something. Supplied hwæthwugu swilces.\par
% glossary-block:B01655:end
% glossary-block:B01656:begin
\needspace{3\baselineskip}
\textbf{something of this sort} — Hwæthwega swilces and the supplied variant hwæthwugu swilces mean something of this sort. Swilces is the genitive expression of quality; preserve the supplied indefinite variant.\par
% glossary-block:B01656:end
% glossary-block:B01657:begin
\needspace{3\baselineskip}
\textbf{sometimes} — \textbf{hwīlum} (adv.).\par
% glossary-block:B01657:end
% glossary-block:B01658:begin
\needspace{3\baselineskip}
\textbf{so much} — Swīðe expresses much as degree, as in love very much; swā swīðe expresses so much. This differs from micel expressing quantity.\par
% glossary-block:B01658:end
% glossary-block:B01659:begin
\needspace{3\baselineskip}
\textbf{son} — \textbf{bearn} (n.). son; child. \textbf{sunu} (n.).\par
% glossary-block:B01659:end
% glossary-block:B01660:begin
\needspace{3\baselineskip}
\textbf{song} — \textbf{sang} (n.).\par
% glossary-block:B01660:end
% glossary-block:B01661:begin
\needspace{3\baselineskip}
\textbf{song of joy} — \textbf{blissesang} (n.).\par
% glossary-block:B01661:end
% glossary-block:B01662:begin
\needspace{3\baselineskip}
\textbf{soon} — \textbf{sōna} (adv.). sōna carries soon in this future-reference prompt.\par
% glossary-block:B01662:end
% glossary-block:B01663:begin
\needspace{3\baselineskip}
\textbf{sore} — \textbf{wund} (n.). wound; sore.\par
% glossary-block:B01663:end
% glossary-block:B01664:begin
\needspace{3\baselineskip}
\textbf{sorrow} — \textbf{sorg} (n.). distress; anxiety; trouble; sorrow.\par
% glossary-block:B01664:end
% glossary-block:B01665:begin
\needspace{3\baselineskip}
\textbf{so that} — Tō ðon ðætte is the complete supplied result connection so that. The demonstrative component alone is not a conjunction meaning so that. \textbf{ðæt} (conj.). that; so that. Supplied expanded spelling; distinguish content from result by the source clause. so that; in order that. Purpose conjunction with the assigned subjunctive construction. that; so that. Asserted result, distinct from intended purpose. Purpose normally uses ðæt + subjunctive in the assigned frames. Do not add a permission modal merely for English may; retain mugan when the prompt explicitly requests it. Tō ðon ðætte is the complete supplied result connection so that; tō alone does not supply that connection.\par
% glossary-block:B01665:end
% glossary-block:B01666:begin
\needspace{3\baselineskip}
\textbf{so that not} — \textbf{ðȳ-lǣs-ðe} (conj.). lest; so that not. ðȳ-lǣs-ðe with the assigned subjunctive negative-purpose clause.\par
% glossary-block:B01666:end
% glossary-block:B01667:begin
\needspace{3\baselineskip}
\textbf{soul} — \textbf{sāwol} (n.).\par
% glossary-block:B01667:end
% glossary-block:B01668:begin
\needspace{3\baselineskip}
\textbf{sound} — \textbf{swēg} (n.). music; sound.\par
% glossary-block:B01668:end
% glossary-block:B01669:begin
\needspace{3\baselineskip}
\textbf{south} — The supplied whole phrase be sūðan means in the south. Be alone does not mean south. \textbf{sūðan} (adv.). in the south; from the south. Retain the supplied spatial expression.\par
% glossary-block:B01669:end
% glossary-block:B01670:begin
\needspace{3\baselineskip}
\textbf{southeastern part} — \textbf{ēastsūðdǣl} (n.).\par
% glossary-block:B01670:end
% glossary-block:B01671:begin
\needspace{3\baselineskip}
\textbf{southern} — \textbf{sūðerne} (adj.).\par
% glossary-block:B01671:end
% glossary-block:B01672:begin
\needspace{3\baselineskip}
\textbf{southwestern} — \textbf{sūðwesterne} (adj.).\par
% glossary-block:B01672:end
% glossary-block:B01673:begin
\needspace{3\baselineskip}
\textbf{sowing season} — \textbf{sǣdtīma} (n.). seedtime; sowing season.\par
% glossary-block:B01673:end
% glossary-block:B01674:begin
\needspace{3\baselineskip}
\textbf{spark} — \textbf{fȳrspearca} (n.).\par
% glossary-block:B01674:end
% glossary-block:B01675:begin
\needspace{3\baselineskip}
\textbf{speak} — \textbf{gecweðan} (v.). speak; say. \textbf{sprecan} (v.). speak with. mid + dative interlocutor in the taught production frame. This task explicitly selects sprecan for say; tō marks the addressee.\par
% glossary-block:B01675:end
% glossary-block:B01676:begin
\needspace{3\baselineskip}
\textbf{speak with} — \textbf{mid} (prep.; adv. in fixed expressions). with. Dative conversational partner. \textbf{sprecan} (v.). mid + dative interlocutor in the taught production frame.\par
% glossary-block:B01676:end
% glossary-block:B01677:begin
\needspace{3\baselineskip}
\textbf{speech} — \textbf{sprǣc} (n.). saying; speech. \textbf{word} (n.). speech; utterance.\par
% glossary-block:B01677:end
% glossary-block:B01678:begin
\needspace{3\baselineskip}
\textbf{speed} — \textbf{hrædnes} (n.). speed; quickness.\par
% glossary-block:B01678:end
% glossary-block:B01679:begin
\needspace{3\baselineskip}
\textbf{spice} — \textbf{wyrtgemang} (n.). spice; mixture of spices. \textbf{wyrt} (n.). Cook specifically routes these spices to wyrt; wyrtgemang is the separate earlier spice word.\par
% glossary-block:B01679:end
% glossary-block:B01680:begin
\needspace{3\baselineskip}
\textbf{spices} — \textbf{wyrt} (n.). Cook specifically routes these spices to wyrt; wyrtgemang is the separate earlier spice word.\par
% glossary-block:B01680:end
% glossary-block:B01681:begin
\needspace{3\baselineskip}
\textbf{spirit} — \textbf{gāst} (n.).\par
% glossary-block:B01681:end
% glossary-block:B01682:begin
\needspace{3\baselineskip}
\textbf{spit} — \textbf{spǣtan} (v.).\par
% glossary-block:B01682:end
% glossary-block:B01683:begin
\needspace{3\baselineskip}
\textbf{'s (possession)} — Possessive/relational English of or 's often becomes a dependent genitive, not preposition of. Partitive, source and lexical government require their own frames.\par
% glossary-block:B01683:end
% glossary-block:B01684:begin
\needspace{3\baselineskip}
\textbf{spread} — \textbf{brǣdan} (v.). spread; extend. Supplied source form spreads.\par
% glossary-block:B01684:end
% glossary-block:B01685:begin
\needspace{3\baselineskip}
\textbf{spy} — \textbf{scēawere} (n.).\par
% glossary-block:B01685:end
% glossary-block:B01686:begin
\needspace{3\baselineskip}
\textbf{stand} — \textbf{standan} (v.).\par
% glossary-block:B01686:end
% glossary-block:B01687:begin
\needspace{3\baselineskip}
\textbf{stand against} — \textbf{standan} (v.). oppose; stand against. Supplied standan ongēan frame, distinct from simple stand.\par
% glossary-block:B01687:end
% glossary-block:B01688:begin
\needspace{3\baselineskip}
\textbf{star} — \textbf{steorra} (n.).\par
% glossary-block:B01688:end
% glossary-block:B01689:begin
\needspace{3\baselineskip}
\textbf{steal} — \textbf{forstelan} (v.). steal away; steal. \textbf{stelan} (v.).\par
% glossary-block:B01689:end
% glossary-block:B01690:begin
\needspace{3\baselineskip}
\textbf{steal away} — \textbf{forstelan} (v.). steal away; steal.\par
% glossary-block:B01690:end
% glossary-block:B01691:begin
\needspace{3\baselineskip}
\textbf{steward} — \textbf{gerēfa} (n.). steward; reeve.\par
% glossary-block:B01691:end
% glossary-block:B01692:begin
\needspace{3\baselineskip}
\textbf{still} — \textbf{gīet} (adv.). still; further. \textbf{stille} (adj.). still; quiet.\par
% glossary-block:B01692:end
% glossary-block:B01693:begin
\needspace{3\baselineskip}
\textbf{stillness} — \textbf{stilnes} (n.). stillness; calm. Supplied genitive singular. stillness; quietness.\par
% glossary-block:B01693:end
% glossary-block:B01694:begin
\needspace{3\baselineskip}
\textbf{stone} — \textbf{stān} (n.).  \textbf{stānen} (adj.). stone; made of stone. Material adjective.\par
% glossary-block:B01694:end
% glossary-block:B01695:begin
\needspace{3\baselineskip}
\textbf{stone wall} — \textbf{stānweall} (n.).\par
% glossary-block:B01695:end
% glossary-block:B01696:begin
\needspace{3\baselineskip}
\textbf{storm} — \textbf{hrēohnes} (n.). tempest; storm. \textbf{storm} (n.).  \textbf{ȳst} (n.).\par
% glossary-block:B01696:end
% glossary-block:B01697:begin
\needspace{3\baselineskip}
\textbf{strand} — \textbf{strand} (n.). strand; sea-shore.\par
% glossary-block:B01697:end
% glossary-block:B01698:begin
\needspace{3\baselineskip}
\textbf{strange} — \textbf{sellic} (adj.). strange; remarkable.\par
% glossary-block:B01698:end
% glossary-block:B01699:begin
\needspace{3\baselineskip}
\textbf{stranger} — \textbf{cuma} (n.). stranger; guest.\par
% glossary-block:B01699:end
% glossary-block:B01700:begin
\needspace{3\baselineskip}
\textbf{strangers} — \textbf{ęlðēodig} (adj.). foreigners; strangers. Substantival adjective.\par
% glossary-block:B01700:end
% glossary-block:B01701:begin
\needspace{3\baselineskip}
\textbf{straw} — \textbf{strēaw} (n.).\par
% glossary-block:B01701:end
% glossary-block:B01702:begin
\needspace{3\baselineskip}
\textbf{stream of water} — \textbf{flōd} (n.). flood; stream of water.\par
% glossary-block:B01702:end
% glossary-block:B01703:begin
\needspace{3\baselineskip}
\textbf{street} — \textbf{strǣt} (n.).\par
% glossary-block:B01703:end
% glossary-block:B01704:begin
\needspace{3\baselineskip}
\textbf{strength} — \textbf{horn} (n.). horn; strength; dignity. Strength/dignity is the figurative force in 37.14, not a replacement for the literal source horn. \textbf{mægen} (n.). power; strength. Heavenly powers in 37.2 is the contextual plural use. \textbf{stręngðu} (n.).  \textbf{strengu} (n.). Recognition class-reference noun; preserve the source spelling and do not infer a plural count meaning.\par
% glossary-block:B01704:end
% glossary-block:B01705:begin
\needspace{3\baselineskip}
\textbf{strengthen} — \textbf{gestrangian} (v.).\par
% glossary-block:B01705:end
% glossary-block:B01706:begin
\needspace{3\baselineskip}
\textbf{strife} — \textbf{geflīt} (n.). dispute; strife.\par
% glossary-block:B01706:end
% glossary-block:B01707:begin
\needspace{3\baselineskip}
\textbf{strike} — \textbf{bēatan} (v.). beat; strike. Accusative person or thing struck. \textbf{geslēan} (v.). strike; smite. Supplied past/participial source form. Supplied source form struck. \textbf{slēan} (v.). strike; smite. strike; smite. Objectless striking in the assigned hammer sentences. Supplied past source form. \textbf{slēan} (v.), \textbf{bēatan} (v.). Slēan means strike/smite and can be objectless in the hammer sentences. Bēatan means beat/strike and takes an accusative object; the explicit bēatan target must be retained.\par
% glossary-block:B01707:end
% glossary-block:B01708:begin
\needspace{3\baselineskip}
\textbf{strong} — \textbf{strang} (adj.). Masculine singular predicate in the source. Attributive adjective with wind in the source. \textbf{swīð} (adjective citation).\par
% glossary-block:B01708:end
% glossary-block:B01709:begin
\needspace{3\baselineskip}
\textbf{stronger} — \textbf{swīðra} (adj.). stronger; fiercer. Comparative degree, distinct from right-side sense.\par
% glossary-block:B01709:end
% glossary-block:B01710:begin
\needspace{3\baselineskip}
\textbf{stronghold} — \textbf{fæsten} (n.). fortification; stronghold.\par
% glossary-block:B01710:end
% glossary-block:B01711:begin
\needspace{3\baselineskip}
\textbf{struck} — \textbf{ðrēan} (v.). chastised; struck.\par
% glossary-block:B01711:end
% glossary-block:B01712:begin
\needspace{3\baselineskip}
\textbf{structure built} — \textbf{getimbrung} (n.). building; structure built.\par
% glossary-block:B01712:end
% glossary-block:B01713:begin
\needspace{3\baselineskip}
\textbf{struggle} — \textbf{gewinn} (n.). distress; struggle. \textbf{winnan} (v.). struggle; toil.\par
% glossary-block:B01713:end
% glossary-block:B01714:begin
\needspace{3\baselineskip}
\textbf{study} — \textbf{geleornian} (v.). learn; study. \textbf{lār} (n.). learning; study. \textbf{leornung} (n.). study; learning.\par
% glossary-block:B01714:end
% glossary-block:B01715:begin
\needspace{3\baselineskip}
\textbf{subjection} — \textbf{underðēodnes} (n.). subjection; submission.\par
% glossary-block:B01715:end
% glossary-block:B01716:begin
\needspace{3\baselineskip}
\textbf{subjoin} — \textbf{underðēodan} (v.). add; subjoin. Supplied passive/participial addition to speech.\par
% glossary-block:B01716:end
% glossary-block:B01717:begin
\needspace{3\baselineskip}
\textbf{submission} — \textbf{underðēodnes} (n.). subjection; submission.\par
% glossary-block:B01717:end
% glossary-block:B01718:begin
\needspace{3\baselineskip}
\textbf{subside} — \textbf{gesmyltan} (v.). become calm. Supplied preterite of the sea becoming calm. \textbf{geswīcan} (v.). Absolute command to the sea; distinct from cessation with a genitive complement.\par
% glossary-block:B01718:end
% glossary-block:B01719:begin
\needspace{3\baselineskip}
\textbf{such} — \textbf{swilc} (indefinite adjective).\par
% glossary-block:B01719:end
% glossary-block:B01720:begin
\needspace{3\baselineskip}
\textbf{suddenly} — \textbf{fǣringa} (adv.).  \textbf{sęmninga} (adv.).\par
% glossary-block:B01720:end
% glossary-block:B01721:begin
\needspace{3\baselineskip}
\textbf{suffer} — \textbf{ðolian} (v.). suffer; endure. Accusative hardship endured.\par
% glossary-block:B01721:end
% glossary-block:B01722:begin
\needspace{3\baselineskip}
\textbf{suffice} — \textbf{genyhtsumian} (v.). suffice; be sufficient for. Complete supplied construction with dative beneficiary. Supplied complete construction with dative beneficiary; plural books require plural subjunctive.\par
% glossary-block:B01722:end
% glossary-block:B01723:begin
\needspace{3\baselineskip}
\textbf{summer} — \textbf{sumer} (n.).\par
% glossary-block:B01723:end
% glossary-block:B01724:begin
\needspace{3\baselineskip}
\textbf{summon} — \textbf{cleopian} (v.). call; summon. Accusative person summoned; distinguish intransitive cry aloud.\par
% glossary-block:B01724:end
% glossary-block:B01725:begin
\needspace{3\baselineskip}
\textbf{sun} — \textbf{sunne} (n.).\par
% glossary-block:B01725:end
% glossary-block:B01726:begin
\needspace{3\baselineskip}
\textbf{sunrise} — \textbf{ūpgang} (n.). rising; sunrise.\par
% glossary-block:B01726:end
% glossary-block:B01727:begin
\needspace{3\baselineskip}
\textbf{support} — \textbf{āfēdan} (v.). support; nourish. Accusative person supported. \textbf{fultum} (n.). help; support.\par
% glossary-block:B01727:end
% glossary-block:B01728:begin
\needspace{3\baselineskip}
\textbf{suppose} — \textbf{wēnan} (v.). suppose; expect; think. Genitive thing expected in the selected nominal frame; content-clause uses are separately taught.\par
% glossary-block:B01728:end
% glossary-block:B01729:begin
\needspace{3\baselineskip}
\textbf{surface} — \textbf{brādnes} (n.). surface; face. Of the deep in the assigned expression.\par
% glossary-block:B01729:end
% glossary-block:B01730:begin
\needspace{3\baselineskip}
\textbf{surge} — \textbf{weallan} (v.). surge; boil. Supplied present participle used adjectivally of the sea.\par
% glossary-block:B01730:end
% glossary-block:B01731:begin
\needspace{3\baselineskip}
\textbf{surround} — \textbf{ymbsęllan} (v.). surround; envelop. Supplied source form. surround; beset. Supplied source form.\par
% glossary-block:B01731:end
% glossary-block:B01732:begin
\needspace{3\baselineskip}
\textbf{sustain} — \textbf{healdan} (v.). maintain; sustain.\par
% glossary-block:B01732:end
% glossary-block:B01733:begin
\needspace{3\baselineskip}
\textbf{swear} — \textbf{swęrian} (v.).\par
% glossary-block:B01733:end
% glossary-block:B01734:begin
\needspace{3\baselineskip}
\textbf{swear by} — \textbf{ðurh, þurh} (prep.). by. Accusative oath standard.\par
% glossary-block:B01734:end
% glossary-block:B01735:begin
\needspace{3\baselineskip}
\textbf{sweet} — \textbf{swēte} (adj.).\par
% glossary-block:B01735:end
% glossary-block:B01736:begin
\needspace{3\baselineskip}
\textbf{sweeter} — \textbf{swētra} (adj.).\par
% glossary-block:B01736:end
% glossary-block:B01737:begin
\needspace{3\baselineskip}
\textbf{sweetest} — \textbf{swētest} (adj.).\par
% glossary-block:B01737:end
% glossary-block:B01738:begin
\needspace{3\baselineskip}
\textbf{sweetness} — \textbf{swētnes} (n.).\par
% glossary-block:B01738:end
% glossary-block:B01739:begin
\needspace{3\baselineskip}
\textbf{swift} — \textbf{swift} (adj.).\par
% glossary-block:B01739:end
% glossary-block:B01740:begin
\needspace{3\baselineskip}
\textbf{swiftness} — \textbf{swiftnes} (n.).\par
% glossary-block:B01740:end
% glossary-block:B01741:begin
\needspace{3\baselineskip}
\textbf{swim} — \textbf{swimman} (v.).\par
% glossary-block:B01741:end
% glossary-block:B01742:begin
\needspace{3\baselineskip}
\textbf{swimming} — \textbf{sund} (n.).\par
% glossary-block:B01742:end
% glossary-block:B01743:begin
\needspace{3\baselineskip}
\textbf{swing} — \textbf{swingan} (v.). swing; brandish. Supplied past source form.\par
% glossary-block:B01743:end
% glossary-block:B01744:begin
\needspace{3\baselineskip}
\textbf{swollen} — \textbf{āðindan} (v.). Supplied weak participial adjective describing the sea.\par
% glossary-block:B01744:end
% book-navigation: standalone bookmark replaced by volume hierarchy
% glossary-block:B01745:begin
\needspace{5\baselineskip}
\subsection*{T}\BookHeading{nav-glossary-english-0021}{subsubsection}{T}{0}
% glossary-block:B01745:end
% glossary-block:B01746:begin
\textbf{table of the law} — \textbf{tabele} (n.). tablet; table. Tables of the law, not an unrestricted furniture meaning.\par
% glossary-block:B01746:end
% glossary-block:B01747:begin
\needspace{3\baselineskip}
\textbf{tablet} — \textbf{tabele} (n.). tablet; table. Tables of the law, not an unrestricted furniture meaning.\par
% glossary-block:B01747:end
% glossary-block:B01748:begin
\needspace{3\baselineskip}
\textbf{tailor} — \textbf{sēamere} (n.).\par
% glossary-block:B01748:end
% glossary-block:B01749:begin
\needspace{3\baselineskip}
\textbf{take} — \textbf{geniman} (v.).  \textbf{lǣdan} (v.). lead; bring; take. A person is the object in the assigned frame. carry; take. A person is conveyed in this supplied frame. take; carry along. \textbf{niman} (v.). Accusative person or thing taken, including food.\par
% glossary-block:B01749:end
% glossary-block:B01750:begin
\needspace{3\baselineskip}
\textbf{take away} — \textbf{āfierran} (v.). remove; take away.\par
% glossary-block:B01750:end
% glossary-block:B01751:begin
\needspace{3\baselineskip}
\textbf{take charge of} — \textbf{underfōn} (v.). receive; take charge of.\par
% glossary-block:B01751:end
% glossary-block:B01752:begin
\needspace{3\baselineskip}
\textbf{take leave of} — \textbf{grētan} (v.). greet; take leave of. Accusative person/place; greeting and farewell depend on context. grētan + accusative person/place; no separate take or bid is needed.\par
% glossary-block:B01752:end
% glossary-block:B01753:begin
\needspace{3\baselineskip}
\textbf{take notice of} — \textbf{gesēon} (v.). regard; take notice of.\par
% glossary-block:B01753:end
% glossary-block:B01754:begin
\needspace{3\baselineskip}
\textbf{take off} — \textbf{unscrȳdan} (v.). undress; take off. Supplied source form; retain its object/reflexive roles.\par
% glossary-block:B01754:end
% glossary-block:B01755:begin
\needspace{3\baselineskip}
\textbf{teach} — \textbf{lǣran} (v.). teach; instruct. \textbf{tǣcan} (v.). teach; show.\par
% glossary-block:B01755:end
% glossary-block:B01756:begin
\needspace{3\baselineskip}
\textbf{teacher} — \textbf{lārēow} (n.). master; teacher.\par
% glossary-block:B01756:end
% glossary-block:B01757:begin
\needspace{3\baselineskip}
\textbf{tear} — \textbf{tēar} (n.).  \textbf{tō-slītan} (v.). Supplied past source form.\par
% glossary-block:B01757:end
% glossary-block:B01758:begin
\needspace{3\baselineskip}
\textbf{tear apart} — \textbf{tō-slītan} (v.). tear. Supplied past source form.\par
% glossary-block:B01758:end
% glossary-block:B01759:begin
\needspace{3\baselineskip}
\textbf{tearful} — \textbf{wēpan} (v.). weeping; tearful. Supplied participial/adjectival use.\par
% glossary-block:B01759:end
% glossary-block:B01760:begin
\needspace{3\baselineskip}
\textbf{tell} — \textbf{sęcgan} (v.). tell; say. Dative addressee and a reported statement or question as content. say; tell. Person addressed is dative; direct or reported content is separate.\par
% glossary-block:B01760:end
% glossary-block:B01761:begin
\needspace{3\baselineskip}
\textbf{tell of beforehand} — \textbf{foresęcgean} (v.). mention before; tell of beforehand.\par
% glossary-block:B01761:end
% glossary-block:B01762:begin
\needspace{3\baselineskip}
\textbf{tell someone to} — \textbf{hātan} (v.). hātan + accusative actor + plain infinitive.\par
% glossary-block:B01762:end
% glossary-block:B01763:begin
\needspace{3\baselineskip}
\textbf{tempest} — \textbf{hrēohnes} (n.). tempest; storm.\par
% glossary-block:B01763:end
% glossary-block:B01764:begin
\needspace{3\baselineskip}
\textbf{temple} — \textbf{tempel} (n.).\par
% glossary-block:B01764:end
% glossary-block:B01765:begin
\needspace{3\baselineskip}
\textbf{temptation} — \textbf{costnung} (n.).\par
% glossary-block:B01765:end
% glossary-block:B01766:begin
\needspace{3\baselineskip}
\textbf{ten} — \textbf{tīen} (cardinal numeral). Cardinal 10.\par
% glossary-block:B01766:end
% glossary-block:B01767:begin
\needspace{3\baselineskip}
\textbf{tenacious} — \textbf{fæsthafol} (adj.). tenacious; withholding. Genitive of what is held back in the source.\par
% glossary-block:B01767:end
% glossary-block:B01768:begin
\needspace{3\baselineskip}
\textbf{terrible} — \textbf{ęgeslic} (adj.). dreadful; terrible.\par
% glossary-block:B01768:end
% glossary-block:B01769:begin
\needspace{3\baselineskip}
\textbf{territory} — \textbf{landgemǣre} (n.). territory; borders.\par
% glossary-block:B01769:end
% glossary-block:B01770:begin
\needspace{3\baselineskip}
\textbf{than} — \textbf{ðonne} (conj.; adv.). Comparison standard follows ðonne; do not confuse temporal ðonne then/when.\par
% glossary-block:B01770:end
% glossary-block:B01771:begin
\needspace{3\baselineskip}
\textbf{thane} — \textbf{ðegn} (n.). vassal; retainer; thane.\par
% glossary-block:B01771:end
% glossary-block:B01772:begin
\needspace{3\baselineskip}
\textbf{thanksgiving} — \textbf{ðancung} (n.).\par
% glossary-block:B01772:end
% glossary-block:B01773:begin
\needspace{3\baselineskip}
\textbf{that} — \textbf{se, sēo, ðæt} (article; demonstrative; relative pron.). who; which; that. Inflected relative; case follows its clause role. that; those. Ðes supplies this/these; se supplies demonstrative that/those. Select the form by case, number and gender. Conjunction and relative that have separate entries. \textbf{se-ðe} (relative). which; that. Supplied se-ðe relative complex referring to the storm. \textbf{ðæt} (conj.). Content conjunction, distinct from demonstrative/relative that. that; so that. Supplied expanded spelling; distinguish content from result by the source clause. that; so that. Asserted result, distinct from intended purpose. Cook leaves the source linkage uncertain; do not impose one content/result classification. \textbf{ðe} (relative particle; clause marker). which; that. Subject-relative referring to the waters. that; which; who. \textbf{ðe} (relative particle; clause marker), \textbf{se, sēo, ðæt} (article; demonstrative; relative pron.). Relative ðe is invariant; the inflected relative/demonstrative se takes the form required by its role in the relative clause. The antecedent’s case does not automatically determine that form.\par
% glossary-block:B01773:end
% glossary-block:B01774:begin
\needspace{3\baselineskip}
\textbf{that … may} — \textbf{ðæt} (conj.). Purpose normally uses ðæt + subjunctive in the assigned frames. Do not add a permission modal merely for English may; retain mugan when the prompt explicitly requests it.\par
% glossary-block:B01774:end
% glossary-block:B01775:begin
\needspace{3\baselineskip}
\textbf{that … might} — \textbf{ðæt} (conj.). Purpose normally uses ðæt + subjunctive in the assigned frames. Do not add a permission modal merely for English may; retain mugan when the prompt explicitly requests it.\par
% glossary-block:B01775:end
% glossary-block:B01776:begin
\needspace{3\baselineskip}
\textbf{that (reported content)} — \textbf{ðæt} (conj.). Content ðæt is separate from demonstrative and relative that; choose mood/time from the reporting construction.\par
% glossary-block:B01776:end
% glossary-block:B01777:begin
\needspace{3\baselineskip}
\textbf{that (result)} — \textbf{ðæt} (conj.). Result clause, often with degree swā/tō ðan; distinguish intended purpose.\par
% glossary-block:B01777:end
% glossary-block:B01778:begin
\needspace{3\baselineskip}
\textbf{that which} — \textbf{se, sēo, ðæt} (article; demonstrative; relative pron.). what; that which. \textbf{ðæt-ðæt} (pron.). that which; what.\par
% glossary-block:B01778:end
% glossary-block:B01779:begin
\needspace{3\baselineskip}
\textbf{the} — The definite article/demonstrative has se, sēo, ðæt as its nominative singular forms; choose its case, gender and number from the noun phrase. Preserve the exercises’ explicit article omissions. Select article agreement only where required; preserve all explicit Cook omission controls and bare noun predicates.\par
% glossary-block:B01779:end
% glossary-block:B01780:begin
\needspace{3\baselineskip}
\textbf{the Danes} — \textbf{Dęne} (n.). Danes. Plural people-name.\par
% glossary-block:B01780:end
% glossary-block:B01781:begin
\needspace{3\baselineskip}
\textbf{the dead} — \textbf{dēad} (adj.). dead people. Substantival adjective. Substantival adjectives take their own agreement/case; do not add man/people/ones where Cook omits it.\par
% glossary-block:B01781:end
% glossary-block:B01782:begin
\needspace{3\baselineskip}
\textbf{the deep} — \textbf{nēowolnes} (n.). abyss; the deep.\par
% glossary-block:B01782:end
% glossary-block:B01783:begin
\needspace{3\baselineskip}
\textbf{thee} — Ðē is singular dative or accusative thee/you; the governed case follows the verb or preposition. Singular ðū, dual git and plural gē systems; thou/thee mark singular, ye plural. Unspecified you allows appropriate choices unless the task specifies number.\par
% glossary-block:B01783:end
% glossary-block:B01784:begin
\needspace{3\baselineskip}
\textbf{the elderly} — \textbf{eald} (adj.). old people; the elderly. Substantival adjective.\par
% glossary-block:B01784:end
% glossary-block:B01785:begin
\needspace{3\baselineskip}
\textbf{the English} — \textbf{Engle} (n.). People, distinct from the language/adjective.\par
% glossary-block:B01785:end
% glossary-block:B01786:begin
\needspace{3\baselineskip}
\textbf{the evil} — \textbf{yfel} (adj.). evil; evil people. Substantival adjective in the assigned use.\par
% glossary-block:B01786:end
% glossary-block:B01787:begin
\needspace{3\baselineskip}
\textbf{the Highest} — \textbf{hēah} (adj.). Divine title.\par
% glossary-block:B01787:end
% glossary-block:B01788:begin
\needspace{3\baselineskip}
\textbf{the hungry} — \textbf{hungrig} (adj.). hungry people. Substantival adjective. Substantival adjectives take their own agreement/case; do not add man/people/ones where Cook omits it.\par
% glossary-block:B01788:end
% glossary-block:B01789:begin
\needspace{3\baselineskip}
\textbf{their} — \textbf{hē} (personal pron.). his; her; its; their. Invariant genitive possessors, unlike agreeing possessive adjectives. Source genitive possessor. Invariant personal genitive possessor; distinguish the same surface used as a governed personal case.\par
% glossary-block:B01789:end
% glossary-block:B01790:begin
\needspace{3\baselineskip}
\textbf{the living} — \textbf{cwic} (adj.). living people. Substantival adjective. Substantival adjectives take their own agreement/case; do not add man/people/ones where Cook omits it.\par
% glossary-block:B01790:end
% glossary-block:B01791:begin
\needspace{3\baselineskip}
\textbf{the Lord} — \textbf{dryhten} (n.). lord; the Lord.\par
% glossary-block:B01791:end
% glossary-block:B01792:begin
\needspace{3\baselineskip}
\textbf{the lowly} — \textbf{ēaðmōd} (adj.). humble people; the lowly. Substantival adjective.\par
% glossary-block:B01792:end
% glossary-block:B01793:begin
\needspace{3\baselineskip}
\textbf{them} — Hīe is plural accusative; him is plural dative. These plural forms serve all genders. Grammatical gender and antecedent determine reference; it may be represented by masculine/feminine forms for a gendered inanimate noun.\par
% glossary-block:B01793:end
% glossary-block:B01794:begin
\needspace{3\baselineskip}
\textbf{the mighty} — \textbf{mihtig} (adj.). the powerful; the mighty. Substantival adjective.\par
% glossary-block:B01794:end
% glossary-block:B01795:begin
\needspace{3\baselineskip}
\textbf{the more … the more} — \textbf{swā} (adverb; conjunction). the … the. Proportional comparison, not the definite article.\par
% glossary-block:B01795:end
% glossary-block:B01796:begin
\needspace{3\baselineskip}
\textbf{themselves} — An ordinary reflexive uses the personal pronoun in the case governed by its verb or preposition. Agreeing self adds the taught emphasis; English -self does not automatically require a separate self. Ordinary reflexive personal pronouns take the governed case; agreeing self adds the taught emphasis. English -self does not automatically require self.\par
% glossary-block:B01796:end
% glossary-block:B01797:begin
\needspace{3\baselineskip}
\textbf{then} — \textbf{ðā} (adv.).  \textbf{ðonne} (conj.; adv.). Sequence of future judgment events. then; at that point.\par
% glossary-block:B01797:end
% glossary-block:B01798:begin
\needspace{3\baselineskip}
\textbf{the naked} — \textbf{nacod} (adj.). naked people. Substantival adjective.\par
% glossary-block:B01798:end
% glossary-block:B01799:begin
\needspace{3\baselineskip}
\textbf{thence} — \textbf{ðanan} (adv.). thence; from there.\par
% glossary-block:B01799:end
% glossary-block:B01800:begin
\needspace{3\baselineskip}
\textbf{the noble} — \textbf{æðele} (adj.). noble people. Substantival adjective in the source.\par
% glossary-block:B01800:end
% glossary-block:B01801:begin
\needspace{3\baselineskip}
\textbf{the poor} — \textbf{earm} (adj.). the poor; the unfortunate. Substantival adjective.\par
% glossary-block:B01801:end
% glossary-block:B01802:begin
\needspace{3\baselineskip}
\textbf{the powerful} — \textbf{mihtig} (adj.). the powerful; the mighty. Substantival adjective.\par
% glossary-block:B01802:end
% glossary-block:B01803:begin
\needspace{3\baselineskip}
\textbf{there} — Ðǣr is locative there, at that place. Presentational there in there is/there was has no obligatory matching word and follows the assigned whole construction.\par
% glossary-block:B01803:end
% glossary-block:B01804:begin
\needspace{3\baselineskip}
\textbf{therefore} — \textbf{for-ðon-ðe} (conj.; adv. in for-ðon). therefore; for that reason. Adverbial consequence, distinct from because. \textbf{for-ðon-ðe} (conj.; adv. in for-ðon), \textbf{eornostlice} (adv.). The supplied consequence use of for-ðon-ðe means therefore/for that reason; eornostlice can mean accordingly/thus. Do not substitute the causal because sense.\par
% glossary-block:B01804:end
% glossary-block:B01805:begin
\needspace{3\baselineskip}
\textbf{therein} — \textbf{ðǣron} (adverb or particle). therein; in it. Supplied ðǣron points back to the sacks' location.\par
% glossary-block:B01805:end
% glossary-block:B01806:begin
\needspace{3\baselineskip}
\textbf{there is} — Presentational there has no obligatory OE word. Cook may supply ðā, use weorðan, or explicitly omit there; locative ðǣr is a distinct retained word.\par
% glossary-block:B01806:end
% glossary-block:B01807:begin
\needspace{3\baselineskip}
\textbf{the rest} — \textbf{ōðer} (adj.; ordinal). remaining; the rest. The rest of the forest/earth in the source.\par
% glossary-block:B01807:end
% glossary-block:B01808:begin
\needspace{3\baselineskip}
\textbf{there was} — Presentational there has no obligatory OE word. Cook may supply ðā, use weorðan, or explicitly omit there; locative ðǣr is a distinct retained word.\par
% glossary-block:B01808:end
% glossary-block:B01809:begin
\needspace{3\baselineskip}
\textbf{the rich} — \textbf{rīce} (adj.). rich; wealthy. Includes substantival rich people in the assigned frame.\par
% glossary-block:B01809:end
% glossary-block:B01810:begin
\needspace{3\baselineskip}
\textbf{the righteous} — \textbf{rihtwīs} (adj.). righteous people. Substantival adjective. Substantival adjectives take their own agreement/case; do not add man/people/ones where Cook omits it.\par
% glossary-block:B01810:end
% glossary-block:B01811:begin
\needspace{3\baselineskip}
\textbf{these} — \textbf{ðes, ðēos, ðis} (demonstrative). this; these. Ðes supplies this/these; se supplies demonstrative that/those. Select the form by case, number and gender. Conjunction and relative that have separate entries.\par
% glossary-block:B01811:end
% glossary-block:B01812:begin
\needspace{3\baselineskip}
\textbf{the sick} — \textbf{untrum} (adj.). sick people. Substantival adjective. Substantival adjectives take their own agreement/case; do not add man/people/ones where Cook omits it.\par
% glossary-block:B01812:end
% glossary-block:B01813:begin
\needspace{3\baselineskip}
\textbf{the … the} — \textbf{swā} (adverb; conjunction). Proportional comparison, not the definite article.\par
% glossary-block:B01813:end
% glossary-block:B01814:begin
\needspace{3\baselineskip}
\textbf{the … the (proportional comparison)} — \textbf{swā} (adverb; conjunction). swā … swā with comparatives; English the is not the definite article here.\par
% glossary-block:B01814:end
% glossary-block:B01815:begin
\needspace{3\baselineskip}
\textbf{the unnoble} — \textbf{unæðele} (adj.). people not noble. Substantival adjective in the supplied contrast.\par
% glossary-block:B01815:end
% glossary-block:B01816:begin
\needspace{3\baselineskip}
\textbf{the unrighteous} — \textbf{unrihtwīs} (adj.). unrighteous people. Substantival adjective.\par
% glossary-block:B01816:end
% glossary-block:B01817:begin
\needspace{3\baselineskip}
\textbf{they} — Hīe is plural nominative they for all genders. The plural dative is him, accusative hīe, genitive hiera. Grammatical gender and antecedent determine reference; it may be represented by masculine/feminine forms for a gendered inanimate noun.\par
% glossary-block:B01817:end
% glossary-block:B01818:begin
\needspace{3\baselineskip}
\textbf{thine} — \textbf{ðīn} (possessive adjective). your; yours. One possessor; agrees with the possessed noun. ðīn identifies one possessor and agrees with the possessed noun. ēower and incer mark plural and dual possessors separately.\par
% glossary-block:B01818:end
% glossary-block:B01819:begin
\needspace{3\baselineskip}
\textbf{thing} — \textbf{ðing} (n.). Ordinary singular/plural noun. thing; fact; matter. Matter known/discussed, distinct from a generic physical object.\par
% glossary-block:B01819:end
% glossary-block:B01820:begin
\needspace{3\baselineskip}
\textbf{things} — \textbf{ðing} (n.). thing. Ordinary singular/plural noun.\par
% glossary-block:B01820:end
% glossary-block:B01821:begin
\needspace{3\baselineskip}
\textbf{think} — \textbf{hycgan} (v.).  \textbf{ðęncan} (v.). Supplied source preterite. \textbf{wēnan} (v.). suppose; expect; think. Genitive thing expected in the selected nominal frame; content-clause uses are separately taught.\par
% glossary-block:B01821:end
% glossary-block:B01822:begin
\needspace{3\baselineskip}
\textbf{third} — \textbf{ðridda} (ordinal numeral).\par
% glossary-block:B01822:end
% glossary-block:B01823:begin
\needspace{3\baselineskip}
\textbf{thirsty} — \textbf{ðurstig} (adj.).\par
% glossary-block:B01823:end
% glossary-block:B01824:begin
\needspace{3\baselineskip}
\textbf{thirty} — \textbf{ðrītig} (numeral).\par
% glossary-block:B01824:end
% glossary-block:B01825:begin
\needspace{3\baselineskip}
\textbf{this} — \textbf{ðes, ðēos, ðis} (demonstrative). this; these. Ðes supplies this/these; se supplies demonstrative that/those. Select the form by case, number and gender. Conjunction and relative that have separate entries.\par
% glossary-block:B01825:end
% glossary-block:B01826:begin
\needspace{3\baselineskip}
\textbf{thither} — \textbf{ðider} (adv.). thither; to that place.\par
% glossary-block:B01826:end
% glossary-block:B01827:begin
\needspace{3\baselineskip}
\textbf{thorn} — \textbf{ðorn} (n.).\par
% glossary-block:B01827:end
% glossary-block:B01828:begin
\needspace{3\baselineskip}
\textbf{those} — \textbf{se, sēo, ðæt} (article; demonstrative; relative pron.). that; those. Supplied substantival dative plural referring to waters. Ðes supplies this/these; se supplies demonstrative that/those. Select the form by case, number and gender. Conjunction and relative that have separate entries.\par
% glossary-block:B01828:end
% glossary-block:B01829:begin
\needspace{3\baselineskip}
\textbf{thou} — Ðū is the singular nominative subject thou/you. Singular ðū, dual git and plural gē systems; thou/thee mark singular, ye plural. Unspecified you allows appropriate choices unless the task specifies number.\par
% glossary-block:B01829:end
% glossary-block:B01830:begin
\needspace{3\baselineskip}
\textbf{though} — \textbf{ðēah ðe}. Supplied whole concessive conjunction in EX-XII-XIII-43-10. \textbf{ðēah} (conj.; adv.). ðēah with the taught concessive clause; a subjunctive need not make an admitted fact unreal.\par
% glossary-block:B01830:end
% glossary-block:B01831:begin
\needspace{3\baselineskip}
\textbf{thought} — \textbf{geðanc} (n.). thought; mind.\par
% glossary-block:B01831:end
% glossary-block:B01832:begin
\needspace{3\baselineskip}
\textbf{thoughtful} — \textbf{behygdig} (adj.). attentive; thoughtful.\par
% glossary-block:B01832:end
% glossary-block:B01833:begin
\needspace{3\baselineskip}
\textbf{thousand} — \textbf{ðūsend} (numeral). Within the complete supplied high-number expression. High quantities retain the complete local numerical expressions; a glossary headword does not replace the supplied genitive/counting construction.\par
% glossary-block:B01833:end
% glossary-block:B01834:begin
\needspace{3\baselineskip}
\textbf{three} — \textbf{ðrīe} (cardinal numeral).\par
% glossary-block:B01834:end
% glossary-block:B01835:begin
\needspace{3\baselineskip}
\textbf{three times} — \textbf{ðrīwa} (adv.). thrice; three times. ðrīwa is an adverb; not a cardinal with a silently invented time noun.\par
% glossary-block:B01835:end
% glossary-block:B01836:begin
\needspace{3\baselineskip}
\textbf{thrice} — \textbf{ðrīwa} (adv.). thrice; three times. ðrīwa is an adverb; not a cardinal with a silently invented time noun.\par
% glossary-block:B01836:end
% glossary-block:B01837:begin
\needspace{3\baselineskip}
\textbf{throne} — \textbf{cynesętl} (n.). royal seat; throne. \textbf{setl} (n.). seat; throne. \textbf{ðrymsetl} (n.). throne; throne of glory.\par
% glossary-block:B01837:end
% glossary-block:B01838:begin
\needspace{3\baselineskip}
\textbf{throne of glory} — \textbf{ðrymsetl} (n.). throne; throne of glory.\par
% glossary-block:B01838:end
% glossary-block:B01839:begin
\needspace{3\baselineskip}
\textbf{through} — \textbf{geond} (prep.). along; through. With accusative. \textbf{ðurh, þurh} (prep.). Accusative landmark or medium. through; by. Mediated action through a person in the supplied frame. through; by means of. Supplied service/evidence expressions, including Ðurh hwæt.\par
% glossary-block:B01839:end
% glossary-block:B01840:begin
\needspace{3\baselineskip}
\textbf{throw} — \textbf{weorpan} (v.). Accusative thing thrown. Accusative thing cast; a following destination has its own government.\par
% glossary-block:B01840:end
% glossary-block:B01841:begin
\needspace{3\baselineskip}
\textbf{thunderbolt} — \textbf{līeg} (n.). thunderbolt; lightning. Source levin sense.\par
% glossary-block:B01841:end
% glossary-block:B01842:begin
\needspace{3\baselineskip}
\textbf{thus} — \textbf{eornostlice} (adv.). accordingly; thus. \textbf{swā} (adverb; conjunction). so; thus. \textbf{ðus} (adv.).\par
% glossary-block:B01842:end
% glossary-block:B01843:begin
\needspace{3\baselineskip}
\textbf{thy} — \textbf{ðīn} (possessive adjective). your; yours. One possessor; agrees with the possessed noun. ðīn identifies one possessor and agrees with the possessed noun. ēower and incer mark plural and dual possessors separately.\par
% glossary-block:B01843:end
% glossary-block:B01844:begin
\needspace{3\baselineskip}
\textbf{tie} — \textbf{bindan} (v.). bind; tie. Accusative person or thing bound.\par
% glossary-block:B01844:end
% glossary-block:B01845:begin
\needspace{3\baselineskip}
\textbf{time} — \textbf{sīð} (n.). time; occasion. \textbf{tīd} (n.). season; time.\par
% glossary-block:B01845:end
% glossary-block:B01846:begin
\needspace{3\baselineskip}
\textbf{tin} — \textbf{tin} (n.).\par
% glossary-block:B01846:end
% glossary-block:B01847:begin
\needspace{3\baselineskip}
\textbf{to} — On introduces the succeeding generation in the supplied from-generation-to-generation range. Retain the full supplied temporal phrase; on alone does not mean from generation to generation. \textbf{tō} (prep.). Endpoint/result death. Dative destination. to; toward. Person to whom knees/prayer are directed in the supplied phrase. Postposed in the supplied cwæð him tō speech phrase. at; to. Bounded supplied root expression with genitive dependent and dative root. to; into. Result in supplied tō lore, tō sorge or tō nāhte. The person addressed in speech. into; to. Resulting language/literary form. Complement of belonging in the supplied frame. Resulting condition of youth in the supplied comparison. Tō + inflected infinitive expresses the taught purpose, to/in order to. Plain infinitives after modals do not take an inserted tō; recipient to may instead be expressed by a dative. A dative noun/pronoun often expresses the recipient/beneficiary without an added preposition. Distinguish destinations and lexical government.\par
% glossary-block:B01847:end
% glossary-block:B01848:begin
\needspace{3\baselineskip}
\textbf{to-day} — \textbf{tō-dæg} (adv.), \textbf{tōdæg} (adv.). Tō-dæg and tōdæg are the retained source/course spellings of today.\par
% glossary-block:B01848:end
% glossary-block:B01849:begin
\needspace{3\baselineskip}
\textbf{today} — \textbf{tō-dæg} (adv.).  \textbf{tōdæg} (adv.).  \textbf{tō-dæg} (adv.), \textbf{tōdæg} (adv.). Tō-dæg and tōdæg are the retained source/course spellings of today.\par
% glossary-block:B01849:end
% glossary-block:B01850:begin
\needspace{3\baselineskip}
\textbf{to death} — Tō introduces death as the endpoint or result; use the dative form of the death noun in the complete supplied expression.\par
% glossary-block:B01850:end
% glossary-block:B01851:begin
\needspace{3\baselineskip}
\textbf{together} — \textbf{ætgædere} (adv.). together; at once. Paired with samod in the supplied source. \textbf{samod} (adv.).  \textbf{tōgædere} (adv.).\par
% glossary-block:B01851:end
% glossary-block:B01852:begin
\needspace{3\baselineskip}
\textbf{toil} — \textbf{gedeorf} (n.). labor; toil. \textbf{winnan} (v.). struggle; toil. Verb; the subject labours with the hands.\par
% glossary-block:B01852:end
% glossary-block:B01853:begin
\needspace{3\baselineskip}
\textbf{to meet} — \textbf{tōgēanes} (postposition). towards; to meet.\par
% glossary-block:B01853:end
% glossary-block:B01854:begin
\needspace{3\baselineskip}
\textbf{to-morrow} — \textbf{tō-morgen} (adv.). Adverb tō-morgen, not separately to plus morning.\par
% glossary-block:B01854:end
% glossary-block:B01855:begin
\needspace{3\baselineskip}
\textbf{tomorrow} — \textbf{tō-morgen} (adv.). Adverb tō-morgen, not separately to plus morning.\par
% glossary-block:B01855:end
% glossary-block:B01856:begin
\needspace{3\baselineskip}
\textbf{to one another} — \textbf{betwēonan} (adverb or particle). between; to one another. Supplied reciprocal betwēonan him in 45.6. Supplied reciprocal betwēonan him; do not translate one and another independently.\par
% glossary-block:B01856:end
% glossary-block:B01857:begin
\needspace{3\baselineskip}
\textbf{tooth} — \textbf{tōð} (n.).\par
% glossary-block:B01857:end
% glossary-block:B01858:begin
\needspace{3\baselineskip}
\textbf{to (purpose)} — \textbf{tō} (prep.). tō + inflected infinitive expresses the assigned elementary purpose; modal complements use plain infinitives.\par
% glossary-block:B01858:end
% glossary-block:B01859:begin
\needspace{3\baselineskip}
\textbf{toss} — \textbf{geswęngan} (v.). toss; shake. Supplied participial source use.\par
% glossary-block:B01859:end
% glossary-block:B01860:begin
\needspace{3\baselineskip}
\textbf{to that degree} — Tō ðan … ðæt is the supplied degree-and-result construction so/to that degree … that. Keep the whole construction.\par
% glossary-block:B01860:end
% glossary-block:B01861:begin
\needspace{3\baselineskip}
\textbf{to that place} — \textbf{ðider} (adv.). thither; to that place.\par
% glossary-block:B01861:end
% glossary-block:B01862:begin
\needspace{3\baselineskip}
\textbf{to the benefit of} — \textbf{for} (prep.). for; to the benefit of. \textbf{tō} (prep.). for; to the benefit of.\par
% glossary-block:B01862:end
% glossary-block:B01863:begin
\needspace{3\baselineskip}
\textbf{toward} — \textbf{in} (prep.). toward; onto. Supplied phrase directing the mind toward these things. \textbf{ofer} (prep.). In gōd ofer + accusative people. \textbf{tō} (prep.). to; toward. Person to whom knees/prayer are directed in the supplied phrase. Conduct toward a person in the supplied expression. Target of a look in the supplied frame. \textbf{wið} (prep.). Conduct toward a person in the supplied local frame.\par
% glossary-block:B01863:end
% glossary-block:B01864:begin
\needspace{3\baselineskip}
\textbf{towards} — \textbf{in} (prep.). toward; onto. Supplied phrase directing the mind toward these things. \textbf{tō} (prep.). toward. Conduct toward a person in the supplied expression. toward. Target of a look in the supplied frame. \textbf{tōgēanes} (postposition). towards; to meet. \textbf{wið} (prep.). toward. Conduct toward a person in the supplied local frame.\par
% glossary-block:B01864:end
% glossary-block:B01865:begin
\needspace{3\baselineskip}
\textbf{tower} — \textbf{hlīfian} (v.). Verb meaning stand/tower high in the source; distinct from noun tower. \textbf{torr} (n.).\par
% glossary-block:B01865:end
% glossary-block:B01866:begin
\needspace{3\baselineskip}
\textbf{to where} — \textbf{hwider} (adv.). whither; to where.\par
% glossary-block:B01866:end
% glossary-block:B01867:begin
\needspace{3\baselineskip}
\textbf{to wife} — Tō wīfe is the supplied role phrase as wife/to wife.\par
% glossary-block:B01867:end
% glossary-block:B01868:begin
\needspace{3\baselineskip}
\textbf{town} — \textbf{burg} (n.). town; city. \textbf{ceaster} (n.). town; city.\par
% glossary-block:B01868:end
% glossary-block:B01869:begin
\needspace{3\baselineskip}
\textbf{townspeople} — \textbf{ceasterwaran} (n.). townspeople; citizens.\par
% glossary-block:B01869:end
% glossary-block:B01870:begin
\needspace{3\baselineskip}
\textbf{to youth} — Tō introduces the resulting youthful condition in the supplied comparison; the accompanying youth noun supplies that meaning.\par
% glossary-block:B01870:end
% glossary-block:B01871:begin
\needspace{3\baselineskip}
\textbf{trade} — \textbf{cēapian} (v.). trade; traffic. Traffic means trade; retain the supplied trading frame. Cook's traffic means trade/buy through trade, not modern road traffic. \textbf{cræft} (n.). trade; occupation; craft.\par
% glossary-block:B01871:end
% glossary-block:B01872:begin
\needspace{3\baselineskip}
\textbf{traffic} — \textbf{cēapian} (v.). trade; traffic. Traffic means trade; retain the supplied trading frame. Cook's traffic means trade/buy through trade, not modern road traffic.\par
% glossary-block:B01872:end
% glossary-block:B01873:begin
\needspace{3\baselineskip}
\textbf{train} — \textbf{getēon} (v.). train; educate.\par
% glossary-block:B01873:end
% glossary-block:B01874:begin
\needspace{3\baselineskip}
\textbf{trained} — \textbf{gelǣred} (adj.). educated; trained; skilled.\par
% glossary-block:B01874:end
% glossary-block:B01875:begin
\needspace{3\baselineskip}
\textbf{trample down} — \textbf{fortrēdan} (v.).\par
% glossary-block:B01875:end
% glossary-block:B01876:begin
\needspace{3\baselineskip}
\textbf{translate} — \textbf{gewęndan} (v.). translate; turn from one language into another. \textbf{węndan} (v.).\par
% glossary-block:B01876:end
% glossary-block:B01877:begin
\needspace{3\baselineskip}
\textbf{translate into} — Tō introduces the resulting language or literary form with the supplied translation verb; tō alone does not express translate.\par
% glossary-block:B01877:end
% glossary-block:B01878:begin
\needspace{3\baselineskip}
\textbf{translation} — \textbf{āwęndednes} (n.).\par
% glossary-block:B01878:end
% glossary-block:B01879:begin
\needspace{3\baselineskip}
\textbf{translator} — \textbf{wealhstōd} (n.). interpreter; translator.\par
% glossary-block:B01879:end
% glossary-block:B01880:begin
\needspace{3\baselineskip}
\textbf{travel} — \textbf{faran} (v.). go; travel. Verb faran; distinct from noun journey sīðfæt.\par
% glossary-block:B01880:end
% glossary-block:B01881:begin
\needspace{3\baselineskip}
\textbf{treasure} — \textbf{goldhord} (n.). treasure; hoard of gold. \textbf{hord} (n.). hoard; treasure.\par
% glossary-block:B01881:end
% glossary-block:B01882:begin
\needspace{3\baselineskip}
\textbf{tree} — \textbf{trēo} (n.).\par
% glossary-block:B01882:end
% glossary-block:B01883:begin
\needspace{3\baselineskip}
\textbf{tribe} — \textbf{mǣgð} (n.). tribe; nation.\par
% glossary-block:B01883:end
% glossary-block:B01884:begin
\needspace{3\baselineskip}
\textbf{trouble} — \textbf{bisgian} (v.). occupy; trouble. \textbf{sorg} (n.). distress; anxiety; trouble; sorrow.\par
% glossary-block:B01884:end
% glossary-block:B01885:begin
\needspace{3\baselineskip}
\textbf{true} — \textbf{sōð} (adj.).\par
% glossary-block:B01885:end
% glossary-block:B01886:begin
\needspace{3\baselineskip}
\textbf{truly} — \textbf{sōðlice} (adv.). truly; indeed. Adverbial confirmation; distinct from discourse connective sōðlice where the key uses but/now.\par
% glossary-block:B01886:end
% glossary-block:B01887:begin
\needspace{3\baselineskip}
\textbf{trumpet} — \textbf{biema} (n.). trumpet; clarion.\par
% glossary-block:B01887:end
% glossary-block:B01888:begin
\needspace{3\baselineskip}
\textbf{trust} — \textbf{trēownes} (n.). trust; confidence.\par
% glossary-block:B01888:end
% glossary-block:B01889:begin
\needspace{3\baselineskip}
\textbf{trust in} — \textbf{on} (prep.). in. Accusative person believed/trusted in. \textbf{gelīefan} (v.). gelīefan on + accusative person in this assigned frame.\par
% glossary-block:B01889:end
% glossary-block:B01890:begin
\needspace{3\baselineskip}
\textbf{truth} — \textbf{sōðfæstnes} (n.).\par
% glossary-block:B01890:end
% glossary-block:B01891:begin
\needspace{3\baselineskip}
\textbf{turn} — \textbf{bewęndan} (v.). The assigned reflexive frame uses dative him and tō + dative goal. \textbf{gewęndan} (v.). return; turn. turn; change back. Supplied source form. \textbf{węndan} (v.). turn; return. Supplied source form.\par
% glossary-block:B01891:end
% glossary-block:B01892:begin
\needspace{3\baselineskip}
\textbf{turn back} — \textbf{ęfthweorfan} (v.). return; turn back. \textbf{hweorfan} (v.). return; turn back. Supplied source preterite.\par
% glossary-block:B01892:end
% glossary-block:B01893:begin
\needspace{3\baselineskip}
\textbf{turn from one language into another} — \textbf{gewęndan} (v.). translate; turn from one language into another.\par
% glossary-block:B01893:end
% glossary-block:B01894:begin
\needspace{3\baselineskip}
\textbf{turn one's attention to} — \textbf{āðęnian} (v.). direct; stretch toward. Directing the mind toward what is heard; supplied figurative phrase.\par
% glossary-block:B01894:end
% glossary-block:B01895:begin
\needspace{3\baselineskip}
\textbf{turn one's gaze} — \textbf{besēon} (v.). look; turn one's gaze. Supplied source form and frame.\par
% glossary-block:B01895:end
% glossary-block:B01896:begin
\needspace{3\baselineskip}
\textbf{twelve} — \textbf{twęlf} (cardinal numeral). Cardinal 12.\par
% glossary-block:B01896:end
% glossary-block:B01897:begin
\needspace{3\baselineskip}
\textbf{twenty} — \textbf{twēntig} (numeral).\par
% glossary-block:B01897:end
% glossary-block:B01898:begin
\needspace{3\baselineskip}
\textbf{two} — \textbf{twēgen} (cardinal numeral). two; in two. Two parts in the supplied division expression.\par
% glossary-block:B01898:end
% glossary-block:B01899:begin
\needspace{3\baselineskip}
\textbf{twofold} — \textbf{twīfeald} (adj.). double; twofold.\par
% glossary-block:B01899:end
% book-navigation: standalone bookmark replaced by volume hierarchy
% glossary-block:B01900:begin
\needspace{5\baselineskip}
\subsection*{U}\BookHeading{nav-glossary-english-0022}{subsubsection}{U}{0}
% glossary-block:B01900:end
% glossary-block:B01901:begin
\textbf{uncounted quantity} — \textbf{unrīm} (n.). abundance; uncounted quantity. Supplied quantity with genitive of the substance.\par
% glossary-block:B01901:end
% glossary-block:B01902:begin
\needspace{3\baselineskip}
\textbf{under} — \textbf{under} (prep.). Dative landmark in the taught static frame. Directional accusative in the assigned entry-under-the-roof frame.\par
% glossary-block:B01902:end
% glossary-block:B01903:begin
\needspace{3\baselineskip}
\textbf{under one's control} — \textbf{on} (prep.). in one's control. on + dative possessed geweald. On + dative possessed geweald expresses in/under someone’s control. In the assigned habban construction, the things controlled are its accusative object; the controller is expressed within the geweald phrase.\par
% glossary-block:B01903:end
% glossary-block:B01904:begin
\needspace{3\baselineskip}
\textbf{understand} — \textbf{ongietan} (v.). understand; perceive. Supplied infinitive after a modal. understand; perceive.\par
% glossary-block:B01904:end
% glossary-block:B01905:begin
\needspace{3\baselineskip}
\textbf{understanding} — \textbf{andgiet} (n.). andgiet denotes the knowledge possessed and desired in this exercise.\par
% glossary-block:B01905:end
% glossary-block:B01906:begin
\needspace{3\baselineskip}
\textbf{undo} — \textbf{undōn} (v.). undo; open.\par
% glossary-block:B01906:end
% glossary-block:B01907:begin
\needspace{3\baselineskip}
\textbf{undress} — \textbf{unscrȳdan} (v.). undress; take off. Supplied source form; retain its object/reflexive roles.\par
% glossary-block:B01907:end
% glossary-block:B01908:begin
\needspace{3\baselineskip}
\textbf{unfamiliar} — \textbf{ungewunelic} (adj.). unusual; unfamiliar.\par
% glossary-block:B01908:end
% glossary-block:B01909:begin
\needspace{3\baselineskip}
\textbf{unfortunate people} — \textbf{earm} (adj.). the poor; the unfortunate. Substantival adjective.\par
% glossary-block:B01909:end
% glossary-block:B01910:begin
\needspace{3\baselineskip}
\textbf{unfree} — \textbf{ðēow} (adj.). the unfree; enslaved people. Substantival adjective in the source comparison.\par
% glossary-block:B01910:end
% glossary-block:B01911:begin
\needspace{3\baselineskip}
\textbf{unknown} — \textbf{ungecnāwen} (past participle; supplied predicate).\par
% glossary-block:B01911:end
% glossary-block:B01912:begin
\needspace{3\baselineskip}
\textbf{unless} — \textbf{būtan} (conj.). unless; except if. Būtan introducing the supplied conditional exception; distinguish a prepositional complement.\par
% glossary-block:B01912:end
% glossary-block:B01913:begin
\needspace{3\baselineskip}
\textbf{unrest} — \textbf{unstilnes} (n.). unrest; disorder.\par
% glossary-block:B01913:end
% glossary-block:B01914:begin
\needspace{3\baselineskip}
\textbf{unrighteous} — \textbf{unrihtwīs} (adj.). Adjective, distinct from the substantive unrighteous people.\par
% glossary-block:B01914:end
% glossary-block:B01915:begin
\needspace{3\baselineskip}
\textbf{unrighteousness} — \textbf{unrihtwīsnes} (n.).\par
% glossary-block:B01915:end
% glossary-block:B01916:begin
\needspace{3\baselineskip}
\textbf{unrighteous people} — \textbf{unrihtwīs} (adj.). Substantival adjective.\par
% glossary-block:B01916:end
% glossary-block:B01917:begin
\needspace{3\baselineskip}
\textbf{unthankful} — \textbf{ungeðancfull} (adj.).\par
% glossary-block:B01917:end
% glossary-block:B01918:begin
\needspace{3\baselineskip}
\textbf{until} — \textbf{oð-ðæt} (conj.). Oð-ðæt introduces until; an anticipated endpoint takes the taught subjunctive construction. \textbf{oð} (prep.). up to; until. With accusative.\par
% glossary-block:B01918:end
% glossary-block:B01919:begin
\needspace{3\baselineskip}
\textbf{unusual} — \textbf{ungewunelic} (adj.). unusual; unfamiliar.\par
% glossary-block:B01919:end
% glossary-block:B01920:begin
\needspace{3\baselineskip}
\textbf{up} — \textbf{ūp} (adv.).\par
% glossary-block:B01920:end
% glossary-block:B01921:begin
\needspace{3\baselineskip}
\textbf{upon} — \textbf{on} (prep.). upon; into. Supplied chancing-upon idiom. Accusative beloved in the supplied love-falls-upon idiom. \textbf{uppan} (prep.). upon; on top of.\par
% glossary-block:B01921:end
% glossary-block:B01922:begin
\needspace{3\baselineskip}
\textbf{up to} — \textbf{oð} (prep.). up to; until. With accusative. \textbf{ōð} (prep.). as far as; up to.\par
% glossary-block:B01922:end
% glossary-block:B01923:begin
\needspace{3\baselineskip}
\textbf{us} — Ūs is first-person plural dative or accusative us; unc is the explicitly dual us two. Select case from the verb/preposition and preserve dual versus plural; dual subjects use plural verb forms.\par
% glossary-block:B01923:end
% glossary-block:B01924:begin
\needspace{3\baselineskip}
\textbf{useful} — \textbf{behēfe} (adj.). Dative person to whom something is useful.\par
% glossary-block:B01924:end
% glossary-block:B01925:begin
\needspace{3\baselineskip}
\textbf{us two} — Unc is first-person dual dative or accusative us two. Select case from the verb/preposition and preserve dual versus plural; dual subjects use plural verb forms.\par
% glossary-block:B01925:end
% glossary-block:B01926:begin
\needspace{3\baselineskip}
\textbf{utterance} — \textbf{word} (n.). speech; utterance.\par
% glossary-block:B01926:end
% book-navigation: standalone bookmark replaced by volume hierarchy
% glossary-block:B01927:begin
\needspace{5\baselineskip}
\subsection*{V}\BookHeading{nav-glossary-english-0023}{subsubsection}{V}{0}
% glossary-block:B01927:end
% glossary-block:B01928:begin
\textbf{valuable} — \textbf{dēorwierðe} (adj.). precious; valuable.\par
% glossary-block:B01928:end
% glossary-block:B01929:begin
\needspace{3\baselineskip}
\textbf{value} — \textbf{weorð} (n.). price; value.\par
% glossary-block:B01929:end
% glossary-block:B01930:begin
\needspace{3\baselineskip}
\textbf{various} — \textbf{mislic} (adj.).\par
% glossary-block:B01930:end
% glossary-block:B01931:begin
\needspace{3\baselineskip}
\textbf{vassal} — \textbf{ðegn} (n.). vassal; retainer; thane.\par
% glossary-block:B01931:end
% glossary-block:B01932:begin
\needspace{3\baselineskip}
\textbf{vein of metal} — \textbf{ōra} (n.). vein of metal; ore (sense uncertain). Supplied dative plural after on.\par
% glossary-block:B01932:end
% glossary-block:B01933:begin
\needspace{3\baselineskip}
\textbf{verily} — \textbf{sōðlice} (adv.). Truly; indeed.\par
% glossary-block:B01933:end
% glossary-block:B01934:begin
\needspace{3\baselineskip}
\textbf{verse} — \textbf{lēoð} (n.). poem; poetry; verse.\par
% glossary-block:B01934:end
% glossary-block:B01935:begin
\needspace{3\baselineskip}
\textbf{very} — \textbf{full} (adv.).  \textbf{swīðe} (adv.). Degree adverb swīðe; retain the separate great/large adjective micel.\par
% glossary-block:B01935:end
% glossary-block:B01936:begin
\needspace{3\baselineskip}
\textbf{very calm} — \textbf{smyltest} (adj.). very calm; calmest. Supplied superlative adjective describing the sea.\par
% glossary-block:B01936:end
% glossary-block:B01937:begin
\needspace{3\baselineskip}
\textbf{very greatly} — Swīðost means most/very greatly. It is the recognition-only superlative reference form of swīðe.\par
% glossary-block:B01937:end
% glossary-block:B01938:begin
\needspace{3\baselineskip}
\textbf{vessel} — \textbf{fæt} (n.). vessel; container.\par
% glossary-block:B01938:end
% glossary-block:B01939:begin
\needspace{3\baselineskip}
\textbf{victory} — \textbf{sige} (n.).\par
% glossary-block:B01939:end
% glossary-block:B01940:begin
\needspace{3\baselineskip}
\textbf{vineyard} — \textbf{wīngeard} (n.).\par
% glossary-block:B01940:end
% glossary-block:B01941:begin
\needspace{3\baselineskip}
\textbf{violently} — \textbf{ðearlice} (adv.). violently; fiercely.\par
% glossary-block:B01941:end
% glossary-block:B01942:begin
\needspace{3\baselineskip}
\textbf{virginity} — \textbf{mægðhād} (n.). maidenhood; virginity.\par
% glossary-block:B01942:end
% glossary-block:B01943:begin
\needspace{3\baselineskip}
\textbf{visible} — \textbf{gesewenlic} (adj.).\par
% glossary-block:B01943:end
% glossary-block:B01944:begin
\needspace{3\baselineskip}
\textbf{visit} — \textbf{genēosian} (v.). Accusative person in the selected prefixed-verb frame.\par
% glossary-block:B01944:end
% glossary-block:B01945:begin
\needspace{3\baselineskip}
\textbf{voice} — \textbf{stefn} (n.).\par
% glossary-block:B01945:end
% glossary-block:B01946:begin
\needspace{3\baselineskip}
\textbf{void} — \textbf{ǣmtig} (adj.). empty; void. \textbf{īdel} (adj.). empty; void.\par
% glossary-block:B01946:end
% book-navigation: standalone bookmark replaced by volume hierarchy
% glossary-block:B01947:begin
\needspace{5\baselineskip}
\subsection*{W}\BookHeading{nav-glossary-english-0024}{subsubsection}{W}{0}
% glossary-block:B01947:end
% glossary-block:B01948:begin
\textbf{wages paid for work} — \textbf{mēd} (n.). hire; wages paid for work.\par
% glossary-block:B01948:end
% glossary-block:B01949:begin
\needspace{3\baselineskip}
\textbf{walk} — \textbf{gān} (v.). go; walk.\par
% glossary-block:B01949:end
% glossary-block:B01950:begin
\needspace{3\baselineskip}
\textbf{wall} — \textbf{weall} (n.).\par
% glossary-block:B01950:end
% glossary-block:B01951:begin
\needspace{3\baselineskip}
\textbf{want to} — \textbf{willan} (v.). wish; want to. willan + plain infinitive denotes volition; intention does not assert the action occurred.\par
% glossary-block:B01951:end
% glossary-block:B01952:begin
\needspace{3\baselineskip}
\textbf{wares} — \textbf{hlæst} (n.). wares; merchandise; cargo. Selected plural source use.\par
% glossary-block:B01952:end
% glossary-block:B01953:begin
\needspace{3\baselineskip}
\textbf{warrior} — \textbf{beorn} (n.).\par
% glossary-block:B01953:end
% glossary-block:B01954:begin
\needspace{3\baselineskip}
\textbf{was} — \textbf{wesan} (v.). be. Copular present/past identification or state; other auxiliary constructions are separate senses. Copular eom/wæs and bēo indicative series; later passive and compound uses require their own construction.\par
% glossary-block:B01954:end
% glossary-block:B01955:begin
\needspace{3\baselineskip}
\textbf{wash} — \textbf{ðwēan} (v.). Accusative body part, mid + dative liquid.\par
% glossary-block:B01955:end
% glossary-block:B01956:begin
\needspace{3\baselineskip}
\textbf{was … -ing} — Simple OE finite forms can render English progressive aspect; wesan + present participle is only one bounded source construction, not a required default.\par
% glossary-block:B01956:end
% glossary-block:B01957:begin
\needspace{3\baselineskip}
\textbf{watch} — \textbf{scēawian} (v.). watch; observe. Accusative person or event watched, including a journey.\par
% glossary-block:B01957:end
% glossary-block:B01958:begin
\needspace{3\baselineskip}
\textbf{water} — \textbf{wæter} (n.).\par
% glossary-block:B01958:end
% glossary-block:B01959:begin
\needspace{3\baselineskip}
\textbf{water for the feet} — The supplied phrase combines water with tō and the dative feet expression; tō supplies the relation for, not the meanings water or feet.\par
% glossary-block:B01959:end
% glossary-block:B01960:begin
\needspace{3\baselineskip}
\textbf{wave} — \textbf{ȳð} (n.). wave; billow. Supplied dative plural after fram.\par
% glossary-block:B01960:end
% glossary-block:B01961:begin
\needspace{3\baselineskip}
\textbf{way} — \textbf{weg} (n.). way; path. way; means of access. Figurative access to the Father in the declaration; distinct from physical road. Weg means road, path or way.\par
% glossary-block:B01961:end
% glossary-block:B01962:begin
\needspace{3\baselineskip}
\textbf{we} — Wē is nominative plural we; wit is the explicitly dual we two. Both take a plural verb form. Select case from the verb/preposition and preserve dual versus plural; dual subjects use plural verb forms.\par
% glossary-block:B01962:end
% glossary-block:B01963:begin
\needspace{3\baselineskip}
\textbf{wealth} — \textbf{wela} (n.). riches; wealth.\par
% glossary-block:B01963:end
% glossary-block:B01964:begin
\needspace{3\baselineskip}
\textbf{wealthy} — \textbf{rīce} (adj.). rich; wealthy. Includes substantival rich people in the assigned frame. \textbf{welig} (adj.). rich; wealthy. \textbf{wlanc} (adj.). proud; wealthy. Contextually the rich man addressed. \textbf{welig} (adj.), \textbf{rīce} (adj.). Welig and rīce can mean rich/wealthy; rīce also occurs substantivally for rich people. Inflect the selected adjective for its role.\par
% glossary-block:B01964:end
% glossary-block:B01965:begin
\needspace{3\baselineskip}
\textbf{weapon} — \textbf{wǣpen} (n.).\par
% glossary-block:B01965:end
% glossary-block:B01966:begin
\needspace{3\baselineskip}
\textbf{wearing} — \textbf{on} (prep.). in; wearing. Supplied on scēapa gegierelum, in sheep's garments.\par
% glossary-block:B01966:end
% glossary-block:B01967:begin
\needspace{3\baselineskip}
\textbf{weep} — \textbf{wēpan} (v.).\par
% glossary-block:B01967:end
% glossary-block:B01968:begin
\needspace{3\baselineskip}
\textbf{weeping} — \textbf{wēpan} (v.). weeping; tearful. Supplied participial/adjectival use.\par
% glossary-block:B01968:end
% glossary-block:B01969:begin
\needspace{3\baselineskip}
\textbf{weight} — \textbf{gewiht} (n.).\par
% glossary-block:B01969:end
% glossary-block:B01970:begin
\needspace{3\baselineskip}
\textbf{well} — \textbf{gesund} (adj.). well; in good health. \textbf{hāl} (adj.). hale; well. Supplied greeting. \textbf{wel} (adv.).\par
% glossary-block:B01970:end
% glossary-block:B01971:begin
\needspace{3\baselineskip}
\textbf{were} — \textbf{wesan} (v.). be. Copular present/past identification or state; other auxiliary constructions are separate senses. Copular eom/wæs and bēo indicative series; later passive and compound uses require their own construction.\par
% glossary-block:B01971:end
% glossary-block:B01972:begin
\needspace{3\baselineskip}
\textbf{were it not for} — \textbf{wesan} (v.). Whole supplied past-subjunctive conditional and wesan + cumen result; preserve its counterfactual anterior sense.\par
% glossary-block:B01972:end
% glossary-block:B01973:begin
\needspace{3\baselineskip}
\textbf{we two} — Wit is nominative dual we two and takes a plural verb form. Select case from the verb/preposition and preserve dual versus plural; dual subjects use plural verb forms.\par
% glossary-block:B01973:end
% glossary-block:B01974:begin
\needspace{3\baselineskip}
\textbf{whale} — \textbf{hwæl} (n.).\par
% glossary-block:B01974:end
% glossary-block:B01975:begin
\needspace{3\baselineskip}
\textbf{what} — \textbf{hwā, hwæt} (interrogative/indefinite pron.). who? what?. Case follows role/government; distinguish personal and neuter reference. Inflect the interrogative by its role/government; dependent questions and relatives are separate constructions. \textbf{se, sēo, ðæt} (article; demonstrative; relative pron.). what; that which. \textbf{ðæt-ðæt} (pron.). that which; what.\par
% glossary-block:B01975:end
% glossary-block:B01976:begin
\needspace{3\baselineskip}
\textbf{whatever} — \textbf{swā-hwæt-swā} (pron.). whatsoever; whatever. Swā hwā swā means whoever, swā hwæt swā whatever, and swā hwæðer swā whichever of two. Inflect the internal pronoun for its clause role; swā alone does not carry the whole free-choice meaning. Swā hwā swā means whoever; swā hwæt swā means whatever. Inflect the internal pronoun for its clause role, so whomsoever requires the appropriate oblique form. Retain any explicitly supplied attraction.\par
% glossary-block:B01976:end
% glossary-block:B01977:begin
\needspace{3\baselineskip}
\textbf{what reason} — \textbf{hwā, hwæt} (interrogative/indefinite pron.). Instrumental hwon in the supplied why expressions.\par
% glossary-block:B01977:end
% glossary-block:B01978:begin
\needspace{3\baselineskip}
\textbf{whatsoever} — \textbf{swā-hwæt-swā} (pron.). whatsoever; whatever. Swā hwā swā means whoever, swā hwæt swā whatever, and swā hwæðer swā whichever of two. Inflect the internal pronoun for its clause role; swā alone does not carry the whole free-choice meaning. Swā hwā swā means whoever; swā hwæt swā means whatever. Inflect the internal pronoun for its clause role, so whomsoever requires the appropriate oblique form. Retain any explicitly supplied attraction.\par
% glossary-block:B01978:end
% glossary-block:B01979:begin
\needspace{3\baselineskip}
\textbf{what sort of} — \textbf{hūlic} (interrogative adjective). Interrogative hūlic asks what sort/kind of; inflect it for the noun phrase required by the task.\par
% glossary-block:B01979:end
% glossary-block:B01980:begin
\needspace{3\baselineskip}
\textbf{wheat} — \textbf{hwǣte} (n.).\par
% glossary-block:B01980:end
% glossary-block:B01981:begin
\needspace{3\baselineskip}
\textbf{when} — \textbf{hwonne} (adv.). when?. Direct interrogative. \textbf{mid-ðȳ} (conj.). when; while. \textbf{ðā} (conj.). when; factual past time. \textbf{ðonne} (conj.). when; recurring or future time. \textbf{ðā} (conj.), \textbf{ðonne} (conj.), \textbf{hwonne} (adv.). Ðā introduces factual past when; ðonne recurring/future when; hwonne asks when? Clause mood follows the particular temporal construction.\par
% glossary-block:B01981:end
% glossary-block:B01982:begin
\needspace{3\baselineskip}
\textbf{whence} — \textbf{hwanon} (adv.). whence; from where.\par
% glossary-block:B01982:end
% glossary-block:B01983:begin
\needspace{3\baselineskip}
\textbf{where} — \textbf{hwǣr} (adv.).\par
% glossary-block:B01983:end
% glossary-block:B01984:begin
\needspace{3\baselineskip}
\textbf{whether} — \textbf{hwæðer} (conj.).\par
% glossary-block:B01984:end
% glossary-block:B01985:begin
\needspace{3\baselineskip}
\textbf{which} — \textbf{hwilc} (interrogative/indefinite adj. or pron.). Selection from a group expressed in genitive. \textbf{se, sēo, ðæt} (article; demonstrative; relative pron.). who; which; that. Inflected relative; case follows its clause role. \textbf{se-ðe} (relative). which; that. Supplied se-ðe relative complex referring to the storm. \textbf{ðe} (relative particle; clause marker). which; that. Subject-relative referring to the waters. that; which; who.\par
% glossary-block:B01985:end
% glossary-block:B01986:begin
\needspace{3\baselineskip}
\textbf{whichever} — Swā hwā swā means whoever, swā hwæt swā whatever, and swā hwæðer swā whichever of two. Inflect the internal pronoun for its clause role; swā alone does not carry the whole free-choice meaning.\par
% glossary-block:B01986:end
% glossary-block:B01987:begin
\needspace{3\baselineskip}
\textbf{whichever of two} — Swā hwæðer swā is the whole free-choice expression whichever of two; the supplied attraction example has a special case arrangement. \textbf{hwæðer} (interrogative pronoun/adjective). Use the whole free-choice expression swā hwæðer swā, with the pronoun inflected for its role. The supplied attraction example has a special case arrangement.\par
% glossary-block:B01987:end
% glossary-block:B01988:begin
\needspace{3\baselineskip}
\textbf{which of} — \textbf{hwilc} (interrogative/indefinite adj. or pron.). which. Selection from a group expressed in genitive.\par
% glossary-block:B01988:end
% glossary-block:B01989:begin
\needspace{3\baselineskip}
\textbf{which of two} — \textbf{hwæðer} (interrogative pronoun/adjective). which of two?. Interrogative hwæðer asks which of two; its case follows its role.\par
% glossary-block:B01989:end
% glossary-block:B01990:begin
\needspace{3\baselineskip}
\textbf{which? of what kind} — \textbf{hwilc} (interrogative/indefinite adj. or pron.). which? of what kind?.\par
% glossary-block:B01990:end
% glossary-block:B01991:begin
\needspace{3\baselineskip}
\textbf{which (relative)} — \textbf{ðe} (relative particle; clause marker), \textbf{se, sēo, ðæt} (article; demonstrative; relative pron.). Relative ðe is invariant; the inflected relative/demonstrative se takes the form required by its role in the relative clause. The antecedent’s case does not automatically determine that form.\par
% glossary-block:B01991:end
% glossary-block:B01992:begin
\needspace{3\baselineskip}
\textbf{while} — \textbf{fæc} (n.). interval; period of time; while. Noun in the supplied time expression. \textbf{mid-ðȳ} (conj.). when; while. \textbf{ðā-hwīle-ðe} (conj.).  \textbf{mid-ðȳ} (conj.), \textbf{ðā-hwīle-ðe} (conj.). Mid-ðȳ and ðā-hwīle-ðe are complete conjunctions meaning while in their assigned temporal frames.\par
% glossary-block:B01992:end
% glossary-block:B01993:begin
\needspace{3\baselineskip}
\textbf{white} — \textbf{hwīt} (adj.).\par
% glossary-block:B01993:end
% glossary-block:B01994:begin
\needspace{3\baselineskip}
\textbf{whither} — \textbf{hwider} (adv.). whither; to where.\par
% glossary-block:B01994:end
% glossary-block:B01995:begin
\needspace{3\baselineskip}
\textbf{who} — \textbf{hwā, hwæt} (interrogative/indefinite pron.). who? what?. Case follows role/government; distinguish personal and neuter reference. Inflect the interrogative by its role/government; dependent questions and relatives are separate constructions. \textbf{se, sēo, ðæt} (article; demonstrative; relative pron.). who; which; that. Inflected relative; case follows its clause role. \textbf{ðe} (relative particle; clause marker). that; which; who. \textbf{ðe} (relative particle; clause marker), \textbf{se, sēo, ðæt} (article; demonstrative; relative pron.). Relative ðe is invariant; the inflected relative/demonstrative se takes the form required by its role in the relative clause. The antecedent’s case does not automatically determine that form.\par
% glossary-block:B01995:end
% glossary-block:B01996:begin
\needspace{3\baselineskip}
\textbf{whoever} — Swā hwā swā means whoever, swā hwæt swā whatever, and swā hwæðer swā whichever of two. Inflect the internal pronoun for its clause role; swā alone does not carry the whole free-choice meaning. Swā hwā swā means whoever; swā hwæt swā means whatever. Inflect the internal pronoun for its clause role, so whomsoever requires the appropriate oblique form. Retain any explicitly supplied attraction.\par
% glossary-block:B01996:end
% glossary-block:B01997:begin
\needspace{3\baselineskip}
\textbf{whole} — \textbf{eall} (adj.; substantive; adv.). all; whole. Eall agrees with a noun in all/whole expressions. Substantival plural refers to all people; neuter singular can mean everything.\par
% glossary-block:B01997:end
% glossary-block:B01998:begin
\needspace{3\baselineskip}
\textbf{wholeheartedly} — \textbf{inweardlīce} (adv.). fervently; wholeheartedly.\par
% glossary-block:B01998:end
% glossary-block:B01999:begin
\needspace{3\baselineskip}
\textbf{whom} — \textbf{hwā, hwæt} (interrogative/indefinite pron.). who? what?. Case follows role/government; distinguish personal and neuter reference. Inflect the interrogative by its role/government; dependent questions and relatives are separate constructions. \textbf{se, sēo, ðæt} (article; demonstrative; relative pron.). who; which; that. Inflected relative; case follows its clause role.\par
% glossary-block:B01999:end
% glossary-block:B02000:begin
\needspace{3\baselineskip}
\textbf{whomsoever} — Swā hwā swā means whoever, swā hwæt swā whatever, and swā hwæðer swā whichever of two. Inflect the internal pronoun for its clause role; swā alone does not carry the whole free-choice meaning. Swā hwā swā means whoever; swā hwæt swā means whatever. Inflect the internal pronoun for its clause role, so whomsoever requires the appropriate oblique form. Retain any explicitly supplied attraction.\par
% glossary-block:B02000:end
% glossary-block:B02001:begin
\needspace{3\baselineskip}
\textbf{whose} — \textbf{hwā, hwæt} (interrogative/indefinite pron.). who? what?. Case follows role/government; distinguish personal and neuter reference. Inflect the interrogative by its role/government; dependent questions and relatives are separate constructions. \textbf{se, sēo, ðæt} (article; demonstrative; relative pron.). who; which; that. Inflected relative; case follows its clause role.\par
% glossary-block:B02001:end
% glossary-block:B02002:begin
\needspace{3\baselineskip}
\textbf{whosoever} — Swā hwā swā means whoever, swā hwæt swā whatever, and swā hwæðer swā whichever of two. Inflect the internal pronoun for its clause role; swā alone does not carry the whole free-choice meaning. Swā hwā swā means whoever; swā hwæt swā means whatever. Inflect the internal pronoun for its clause role, so whomsoever requires the appropriate oblique form. Retain any explicitly supplied attraction.\par
% glossary-block:B02002:end
% glossary-block:B02003:begin
\needspace{3\baselineskip}
\textbf{why} — \textbf{for} (prep.). why; for what reason. In for hwȳ or for hwon. \textbf{hwā, hwæt} (interrogative/indefinite pron.). what reason. Instrumental hwon in the supplied why expressions. \textbf{hwȳ} (adv.). why?. Tō hwon is the whole expression why/for what purpose.\par
% glossary-block:B02003:end
% glossary-block:B02004:begin
\needspace{3\baselineskip}
\textbf{wide} — \textbf{wīd} (adjective citation).\par
% glossary-block:B02004:end
% glossary-block:B02005:begin
\needspace{3\baselineskip}
\textbf{widely} — \textbf{wīde} (adv.). widely; far and wide.\par
% glossary-block:B02005:end
% glossary-block:B02006:begin
\needspace{3\baselineskip}
\textbf{wife} — \textbf{wīf} (n.). Wīf is grammatically neuter. Choose the woman or wife sense from the context.\par
% glossary-block:B02006:end
% glossary-block:B02007:begin
\needspace{3\baselineskip}
\textbf{will} — \textbf{willa} (n.). will; disposition. Present morphology can carry future reference; bēo is used in its bounded expected-state frame. Do not add willan unless volition or a stated source target requires it.\par
% glossary-block:B02007:end
% glossary-block:B02008:begin
\needspace{3\baselineskip}
\textbf{wind} — \textbf{wind} (n.).  \textbf{windan} (v.). wind; plait.\par
% glossary-block:B02008:end
% glossary-block:B02009:begin
\needspace{3\baselineskip}
\textbf{window} — \textbf{ēagðȳrel} (n.).\par
% glossary-block:B02009:end
% glossary-block:B02010:begin
\needspace{3\baselineskip}
\textbf{wine} — \textbf{wīn} (n.).\par
% glossary-block:B02010:end
% glossary-block:B02011:begin
\needspace{3\baselineskip}
\textbf{winged} — \textbf{gefiðred} (adj.). Adjective agreeing with the fowl.\par
% glossary-block:B02011:end
% glossary-block:B02012:begin
\needspace{3\baselineskip}
\textbf{winsome} — \textbf{wynsum} (adj.). pleasant; winsome.\par
% glossary-block:B02012:end
% glossary-block:B02013:begin
\needspace{3\baselineskip}
\textbf{winter} — \textbf{winter} (n.).  \textbf{winterlic} (adj.). winter; of winter.\par
% glossary-block:B02013:end
% glossary-block:B02014:begin
\needspace{3\baselineskip}
\textbf{wisdom} — \textbf{wīsdōm} (n.).\par
% glossary-block:B02014:end
% glossary-block:B02015:begin
\needspace{3\baselineskip}
\textbf{wise} — \textbf{wīs} (adj.).\par
% glossary-block:B02015:end
% glossary-block:B02016:begin
\needspace{3\baselineskip}
\textbf{wise man} — \textbf{wīs} (adj.). wise man; wise one. Substantival direct address.\par
% glossary-block:B02016:end
% glossary-block:B02017:begin
\needspace{3\baselineskip}
\textbf{wise one} — \textbf{wīs} (adj.). wise man; wise one. Substantival direct address.\par
% glossary-block:B02017:end
% glossary-block:B02018:begin
\needspace{3\baselineskip}
\textbf{wiser} — \textbf{wīsra} (adj.).\par
% glossary-block:B02018:end
% glossary-block:B02019:begin
\needspace{3\baselineskip}
\textbf{wish} — \textbf{gewilnian} (v.). desire; wish. \textbf{willan} (v.). wish; want to.\par
% glossary-block:B02019:end
% glossary-block:B02020:begin
\needspace{3\baselineskip}
\textbf{wish to} — \textbf{willan} (v.). willan + plain infinitive denotes volition; intention does not assert the action occurred.\par
% glossary-block:B02020:end
% glossary-block:B02021:begin
\needspace{3\baselineskip}
\textbf{wish well} — \textbf{biddan} (v.). wish well; bid farewell. Supplied farewell formula.\par
% glossary-block:B02021:end
% glossary-block:B02022:begin
\needspace{3\baselineskip}
\textbf{with} — \textbf{mid} (prep.; adv. in fixed expressions). with; accompanying. Dative person or cargo in the ordinary accompaniment frame. with; by means of. Dative instrument or means. Circumstances of risk or loss in the supplied phrase. Manner in the supplied compassion phrase. with; along with. Supplied mid wes, be with them. Dative garment in the assigned clothing frame. Dative conversational partner. Accompanying quality or mental disposition in the source. Material or clothing in the supplied adornment phrase. Exact cloth/cloak complements of begyrd and bewǣfed. with; by means of. Manner/circumstance of victory. \textbf{wið} (prep.). with; against. Dative fellow players in the exact supplied playing frame. Supplied conversational counterpart; do not infer hostility.\par
% glossary-block:B02022:end
% glossary-block:B02023:begin
\needspace{3\baselineskip}
\textbf{with compassion} — Mid + the supplied dative noun of compassion forms the whole manner expression with compassion; mid alone means with.\par
% glossary-block:B02023:end
% glossary-block:B02024:begin
\needspace{3\baselineskip}
\textbf{with difficulty} — \textbf{unēaðe} (adv.). with difficulty; barely.\par
% glossary-block:B02024:end
% glossary-block:B02025:begin
\needspace{3\baselineskip}
\textbf{withholding} — \textbf{fæsthafol} (adj.). tenacious; withholding. Genitive of what is held back in the source.\par
% glossary-block:B02025:end
% glossary-block:B02026:begin
\needspace{3\baselineskip}
\textbf{within} — \textbf{binnan} (prep.).  \textbf{innan} (adv.). within; inside.\par
% glossary-block:B02026:end
% glossary-block:B02027:begin
\needspace{3\baselineskip}
\textbf{without} — \textbf{būtan} (prep.). With dative. Būtan + dative means without a companion or outside a city in the respective sense. Conditional unless is separate.\par
% glossary-block:B02027:end
% glossary-block:B02028:begin
\needspace{3\baselineskip}
\textbf{with permission} — \textbf{be} (prep.). by permission. Supplied permission expression.\par
% glossary-block:B02028:end
% glossary-block:B02029:begin
\needspace{3\baselineskip}
\textbf{with respect to} — \textbf{on} (prep.). in; with respect to. Quality or domain of description.\par
% glossary-block:B02029:end
% glossary-block:B02030:begin
\needspace{3\baselineskip}
\textbf{withstand} — \textbf{wiðstandan} (v.). withstand; resist. Dative of what is resisted.\par
% glossary-block:B02030:end
% glossary-block:B02031:begin
\needspace{3\baselineskip}
\textbf{with victory} — \textbf{mid} (prep.; adv. in fixed expressions). with. Manner/circumstance of victory.\par
% glossary-block:B02031:end
% glossary-block:B02032:begin
\needspace{3\baselineskip}
\textbf{wolf} — \textbf{wulf} (n.).\par
% glossary-block:B02032:end
% glossary-block:B02033:begin
\needspace{3\baselineskip}
\textbf{woman} — \textbf{wīf} (n.). Wīf is grammatically neuter. Choose the woman or wife sense from the context.\par
% glossary-block:B02033:end
% glossary-block:B02034:begin
\needspace{3\baselineskip}
\textbf{wonder} — \textbf{wundor} (n.).  \textbf{wundrian} (v.). Supplied source form.\par
% glossary-block:B02034:end
% glossary-block:B02035:begin
\needspace{3\baselineskip}
\textbf{woo} — \textbf{giernan} (v.). desire; request; woo. Genitive desired thing/person in the selected frames, including the maiden wooed. Genitive person wooed; keep the marriage-seeking sense distinct from a general want.\par
% glossary-block:B02035:end
% glossary-block:B02036:begin
\needspace{3\baselineskip}
\textbf{wood} — \textbf{wudu} (n.). wood; forest. Collective vegetation, not a new material sense.\par
% glossary-block:B02036:end
% glossary-block:B02037:begin
\needspace{3\baselineskip}
\textbf{word} — \textbf{word} (n.).\par
% glossary-block:B02037:end
% glossary-block:B02038:begin
\needspace{3\baselineskip}
\textbf{word by word} — \textbf{be} (prep.). In word be worde; literal correspondence.\par
% glossary-block:B02038:end
% glossary-block:B02039:begin
\needspace{3\baselineskip}
\textbf{work} — \textbf{gewinn} (n.). labour; work. Noun, distinct from labour as a verb. Noun in the expression from his labour. \textbf{weorc} (n.). weorc includes deed/action in these prompts; dǣd is a separate lexical option elsewhere. \textbf{wyrcean} (v.). work; perform. Accusative work or deed.\par
% glossary-block:B02039:end
% glossary-block:B02040:begin
\needspace{3\baselineskip}
\textbf{workman} — \textbf{wyrhta} (n.). workman; craftsman.\par
% glossary-block:B02040:end
% glossary-block:B02041:begin
\needspace{3\baselineskip}
\textbf{world} — \textbf{middangeard} (n.). world; earth. \textbf{woruld} (n.).\par
% glossary-block:B02041:end
% glossary-block:B02042:begin
\needspace{3\baselineskip}
\textbf{worldly dignity} — \textbf{woruldgeðyngð} (n.). worldly honour; worldly dignity.\par
% glossary-block:B02042:end
% glossary-block:B02043:begin
\needspace{3\baselineskip}
\textbf{worldly honour} — \textbf{woruldgeðyngð} (n.). worldly honour; worldly dignity.\par
% glossary-block:B02043:end
% glossary-block:B02044:begin
\needspace{3\baselineskip}
\textbf{worldly life} — \textbf{woruldlīf} (n.).\par
% glossary-block:B02044:end
% glossary-block:B02045:begin
\needspace{3\baselineskip}
\textbf{worldly occupation} — \textbf{woruldcræft} (n.).\par
% glossary-block:B02045:end
% glossary-block:B02046:begin
\needspace{3\baselineskip}
\textbf{worldly occupations} — \textbf{woruldbisgu} (n.). worldly occupations; cares.\par
% glossary-block:B02046:end
% glossary-block:B02047:begin
\needspace{3\baselineskip}
\textbf{worse} — \textbf{wiersa} (adj.).\par
% glossary-block:B02047:end
% glossary-block:B02048:begin
\needspace{3\baselineskip}
\textbf{worst} — \textbf{wierrest} (adj.).\par
% glossary-block:B02048:end
% glossary-block:B02049:begin
\needspace{3\baselineskip}
\textbf{worthy} — \textbf{wierðe} (adj.).\par
% glossary-block:B02049:end
% glossary-block:B02050:begin
\needspace{3\baselineskip}
\textbf{would} — Select intention, obligation, reported prospect or counterfactual context; English spelling alone does not determine tense, mood or an OE modal.\par
% glossary-block:B02050:end
% glossary-block:B02051:begin
\needspace{3\baselineskip}
\textbf{would have come} — The supplied past-subjunctive wesan + cumen construction expresses would have come in the counterfactual result. It does not require separate Old English words for would and have. \textbf{wesan} (v.). Whole supplied past-subjunctive conditional and wesan + cumen result; preserve its counterfactual anterior sense.\par
% glossary-block:B02051:end
% glossary-block:B02052:begin
\needspace{3\baselineskip}
\textbf{would rather} — \textbf{lēofra} (adj.). preferable; dearer. Dative experiencer + is lēofre + infinitive in the taught preference frame. The preference construction is dative experiencer + is lēofre + plain infinitive, is preferable to someone. It renders had/would rather without requiring past time. Wician in the supplied example is the action preferred, sojourn/live with, not a word meaning prefer.\par
% glossary-block:B02052:end
% glossary-block:B02053:begin
\needspace{3\baselineskip}
\textbf{would that} — A present-unreal wish introduced by English would that uses the taught preterite subjunctive. The lexical verb expresses the desired state; no mandatory willan translates would.\par
% glossary-block:B02053:end
% glossary-block:B02054:begin
\needspace{3\baselineskip}
\textbf{wound} — \textbf{wund} (n.). wound; sore.\par
% glossary-block:B02054:end
% glossary-block:B02055:begin
\needspace{3\baselineskip}
\textbf{wrap} — \textbf{bewǣfan} (v.). wrap; clothe. Supplied participial clothing expression.\par
% glossary-block:B02055:end
% glossary-block:B02056:begin
\needspace{3\baselineskip}
\textbf{wrapped with} — \textbf{mid} (prep.; adv. in fixed expressions). with. Exact cloth/cloak complements of begyrd and bewǣfed.\par
% glossary-block:B02056:end
% glossary-block:B02057:begin
\needspace{3\baselineskip}
\textbf{write} — \textbf{āwrītan} (v.).  \textbf{wrītan} (v.). Accusative text written.\par
% glossary-block:B02057:end
% glossary-block:B02058:begin
\needspace{3\baselineskip}
\textbf{written document} — \textbf{gewrit} (n.). letter; written document.\par
% glossary-block:B02058:end
% glossary-block:B02059:begin
\needspace{3\baselineskip}
\textbf{wrongly} — \textbf{yfele} (adv.). badly; wrongly.\par
% glossary-block:B02059:end
% book-navigation: standalone bookmark replaced by volume hierarchy
% glossary-block:B02060:begin
\needspace{5\baselineskip}
\subsection*{Y}\BookHeading{nav-glossary-english-0025}{subsubsection}{Y}{0}
% glossary-block:B02060:end
% glossary-block:B02061:begin
\textbf{ye} — Gē is plural nominative ye/you; dative/accusative plural is ēow. Singular ðū, dual git and plural gē systems; thou/thee mark singular, ye plural. Unspecified you allows appropriate choices unless the task specifies number.\par
% glossary-block:B02061:end
% glossary-block:B02062:begin
\needspace{3\baselineskip}
\textbf{year} — \textbf{gēar} (n.).  \textbf{winter} (n.). Supplied age expression uses wintra; gēar is the ordinary year lemma elsewhere.\par
% glossary-block:B02062:end
% glossary-block:B02063:begin
\needspace{3\baselineskip}
\textbf{years of age} — \textbf{winter} (n.). Supplied age expression uses wintra; gēar is the ordinary year lemma elsewhere.\par
% glossary-block:B02063:end
% glossary-block:B02064:begin
\needspace{3\baselineskip}
\textbf{yet} — \textbf{ðā-hwæðre} (conj.). nevertheless; yet. Whole supplied connective. \textbf{ðēah} (conj.; adv.). nevertheless; yet. Independent adverb, distinct from concessive conjunction although. \textbf{gīet} (adv.). Still; continuing until now, distinct from adversative nevertheless. Retain exact nū gīet where it is supplied.\par
% glossary-block:B02064:end
% glossary-block:B02065:begin
\needspace{3\baselineskip}
\textbf{yield} — \textbf{gieldan} (v.).  \textbf{wyrcean} (v.). yield; produce. Accusative fruit in the taught/source expression. Accusative fruit; birð wæstm is a recognized lexical alternative in the key.\par
% glossary-block:B02065:end
% glossary-block:B02066:begin
\needspace{3\baselineskip}
\textbf{you} — Singular: nominative ðū, dative/accusative ðē; dual: nominative git, dative/accusative inc; plural: nominative gē, dative/accusative ēow. Dual subjects take plural verb forms. Singular ðū, dual git and plural gē systems; thou/thee mark singular, ye plural. Unspecified you allows appropriate choices unless the task specifies number.\par
% glossary-block:B02066:end
% glossary-block:B02067:begin
\needspace{3\baselineskip}
\textbf{young} — \textbf{geong} (adj.).\par
% glossary-block:B02067:end
% glossary-block:B02068:begin
\needspace{3\baselineskip}
\textbf{youngest} — Supplied weak superlative gingstan means youngest in the assigned accusative/genitive masculine singular contexts; positive geong means young.\par
% glossary-block:B02068:end
% glossary-block:B02069:begin
\needspace{3\baselineskip}
\textbf{young man} — \textbf{cniht} (n.). young man; youth.\par
% glossary-block:B02069:end
% glossary-block:B02070:begin
\needspace{3\baselineskip}
\textbf{young woman} — \textbf{fǣmne} (n.). maiden; young woman.\par
% glossary-block:B02070:end
% glossary-block:B02071:begin
\needspace{3\baselineskip}
\textbf{your} — \textbf{ēower} (possessive adjective). your; yours. Plural possessors; inflection agrees with the possessed noun. \textbf{incer} (possessive adjective). your; yours. Two possessors; agreeing possessive. \textbf{ðīn} (possessive adjective). your; yours. One possessor; agrees with the possessed noun. ðīn identifies one possessor and agrees with the possessed noun. ēower and incer mark plural and dual possessors separately. \textbf{ðīn} (possessive adjective), \textbf{incer} (possessive adjective), \textbf{ēower} (possessive adjective). Ðīn is for one second-person possessor, incer for two, ēower for plural possessors. The possessive agrees with the possessed noun.\par
% glossary-block:B02071:end
% glossary-block:B02072:begin
\needspace{3\baselineskip}
\textbf{yours} — \textbf{ēower} (possessive adjective). your; yours. Plural possessors; inflection agrees with the possessed noun. \textbf{incer} (possessive adjective). your; yours. Two possessors; agreeing possessive. \textbf{ðīn} (possessive adjective). your; yours. One possessor; agrees with the possessed noun. ðīn identifies one possessor and agrees with the possessed noun. ēower and incer mark plural and dual possessors separately. \textbf{ðīn} (possessive adjective), \textbf{incer} (possessive adjective), \textbf{ēower} (possessive adjective). Ðīn is for one second-person possessor, incer for two, ēower for plural possessors. The possessive agrees with the possessed noun.\par
% glossary-block:B02072:end
% glossary-block:B02073:begin
\needspace{3\baselineskip}
\textbf{yourself} — An ordinary reflexive uses the personal pronoun in the case governed by its verb or preposition. Agreeing self adds the taught emphasis; English -self does not automatically require a separate self. Ordinary reflexive personal pronouns take the governed case; agreeing self adds the taught emphasis. English -self does not automatically require self.\par
% glossary-block:B02073:end
% glossary-block:B02074:begin
\needspace{3\baselineskip}
\textbf{yourselves} — An ordinary reflexive uses the personal pronoun in the case governed by its verb or preposition. Agreeing self adds the taught emphasis; English -self does not automatically require a separate self. Ordinary reflexive personal pronouns take the governed case; agreeing self adds the taught emphasis. English -self does not automatically require self.\par
% glossary-block:B02074:end
% glossary-block:B02075:begin
\needspace{3\baselineskip}
\textbf{youth} — \textbf{cniht} (n.). young man; youth. \textbf{geoguð} (n.).\par
% glossary-block:B02075:end
% glossary-block:B02076:begin
\needspace{3\baselineskip}
\textbf{you two} — Git is nominative dual you two; inc is dative/accusative dual. A dual subject takes a plural verb form.\par
% glossary-block:B02076:end
% book-navigation: standalone bookmark replaced by volume hierarchy
% glossary-block:B02077:begin
\needspace{5\baselineskip}
\subsection*{Æ}\BookHeading{nav-glossary-english-0026}{subsubsection}{Æ}{0}
% glossary-block:B02077:end
% glossary-block:B02078:begin
\textbf{Æthelthryth} — \textbf{Æðelðrȳð} (proper name). Etheldred; Æthelthryth. Proper name.\par
% glossary-block:B02078:end
% glossary-block:B02079:begin
\needspace{3\baselineskip}
\textbf{Æthelwald} — \textbf{Æðelwald} (proper name). Proper name in apposition to father.\par
% glossary-block:B02079:end
% glossary-block:B02080:begin
\needspace{3\baselineskip}
\textbf{Æthelwulf} — \textbf{Æthelwulf} (proper name).\par
% glossary-block:B02080:end
\end{multicols}

% book-source-body:end
\par\label{end-glossary-english}
\clearpage\endgroup

\BookRecto\begingroup\renewcommand{\runningtitle}{Grammatical index}
\phantomsection\label{grammatical-index}\addcontentsline{toc}{section}{Grammatical index}
\booktitle{REFERENCE}{Grammatical index}
\fontsize{10.5}{13.2}\selectfont\setlength{\columnsep}{6mm}\raggedcolumns
\begin{multicols}{2}
\BookNeedspace{5\baselineskip}
\par\addvspace{6pt}{\large\bfseries A\par}
\BookNeedspace{3\baselineskip}
\par\addvspace{3pt}\noindent\hypertarget{index-0d583a3d28c1}{}\textbf{A-restoration}. \emph{See} \hyperlink{index-ff2eeaa96a81}{Sound changes, fronting and restoration}.\par
\BookNeedspace{3\baselineskip}
\par\addvspace{3pt}\noindent\hypertarget{index-9c40ae32bd6b}{}\textbf{Ability}. \emph{See} \hyperlink{index-09efa301eda1}{Modal verbs, learned skill and capability}.\par
\BookNeedspace{3\baselineskip}
\par\addvspace{3pt}\noindent\hypertarget{index-8ebb9081cbac}{}\textbf{Ablaut}\par\nopagebreak[4]
\par\noindent\hangindent=1em\hangafter=1\hspace*{1em}inherited vowel alternation versus umlaut, \hyperref[bk-sound-history-learner-7a97cb9f6faa]{781}\par
\BookNeedspace{3\baselineskip}
\par\addvspace{3pt}\noindent\hypertarget{index-bd9e40c9e577}{}\textbf{Accusative}\par\nopagebreak[4]
\par\noindent\hangindent=1em\hangafter=1\hspace*{1em}actor of an infinitive, \hyperref[bk-ch09-learner-637a34313b91]{235}, \hyperref[bk-ch09-learner-22cf9e7f4d0c]{236}\par
\par\noindent\hangindent=1em\hangafter=1\hspace*{1em}agreeing object predicate, \hyperref[bk-ch09-learner-33765ec333c9]{236}\par
\par\noindent\hangindent=1em\hangafter=1\hspace*{1em}cognate object, \hyperref[bk-ch08-learner-f078b5ed9d15]{218}\par
\par\noindent\hangindent=1em\hangafter=1\hspace*{1em}direct object, \hyperref[bk-ch01-learner-422453d6be7d]{20}, \hyperref[bk-ch04-learner-0e4e42b78c4b]{99}\par
\par\noindent\hangindent=1em\hangafter=1\hspace*{1em}duration, \hyperref[bk-ch12-learner-c457eea8440b]{382}\par
\par\noindent\hangindent=1em\hangafter=1\hspace*{1em}experiencer with impersonal hyngran, \hyperref[bk-ch10-learner-f11ca067063f]{290}\par
\par\noindent\hangindent=1em\hangafter=1\hspace*{1em}object and predicate in naming, \hyperref[bk-ch07-learner-645dc7063f24]{189}\par
\BookNeedspace{3\baselineskip}
\par\addvspace{3pt}\noindent\hypertarget{index-64f3006be77a}{}\textbf{Action substitution}\par\nopagebreak[4]
\par\noindent\hangindent=1em\hangafter=1\hspace*{1em}dōn repeats an earlier action, \hyperref[bk-ch16-learner-6a25e97d3147]{573}\par
\BookNeedspace{3\baselineskip}
\par\addvspace{3pt}\noindent\hypertarget{index-a47426d0da6f}{}\textbf{Adjectives}\par\nopagebreak[4]
\par\noindent\hangindent=1em\hangafter=1\hspace*{1em}accusative object predicates, \hyperref[bk-ch09-learner-33765ec333c9]{236}\par
\par\noindent\hangindent=1em\hangafter=1\hspace*{1em}after possessives, \hyperref[bk-ch03-learner-65f86edd6e8d]{58}\par
\par\noindent\hangindent=1em\hangafter=1\hspace*{1em}attributive and predicative uses, \hyperref[bk-ch03-learner-dbf98476e10a]{58}\par
\par\noindent\hangindent=1em\hangafter=1\hspace*{1em}bare strong forms naming people, \hyperref[bk-ch11-learner-deff8a6444cf]{334}\par
\par\noindent\hangindent=1em\hangafter=1\hspace*{1em}choice of strong or weak declension, \hyperref[bk-ch03-learner-dbf98476e10a]{58}\par
\par\noindent\hangindent=1em\hangafter=1\hspace*{1em}direct address, \hyperref[bk-ch09-learner-9466d7df4e42]{253}, \hyperref[bk-ch11-learner-af84902d6468]{335}\par
\par\noindent\hangindent=1em\hangafter=1\hspace*{1em}doubled consonants before endings, \hyperref[bk-ch10-learner-0c3379bd0a71]{285}\par
\par\noindent\hangindent=1em\hangafter=1\hspace*{1em}h-loss and contraction, \hyperref[bk-ch03-learner-480a776eec52]{62}\par
\par\noindent\hangindent=1em\hangafter=1\hspace*{1em}lexical final -e, \hyperref[bk-ch10-learner-85a5a29685fa]{283}\par
\par\noindent\hangindent=1em\hangafter=1\hspace*{1em}mixed-gender predicate agreement, \hyperref[bk-ch09-learner-db7c4490e70c]{253}\par
\par\noindent\hangindent=1em\hangafter=1\hspace*{1em}predicate agreement with become, \hyperref[bk-ch04-learner-bfb873b12018]{100}\par
\par\noindent\hangindent=1em\hangafter=1\hspace*{1em}predicate and attributive agreement, \hyperref[bk-ch10-learner-f9f1c5d1eb57]{281}\par
\par\noindent\hangindent=1em\hangafter=1\hspace*{1em}short stems and -u, \hyperref[bk-ch03-learner-fa69c0c29f5c]{60}\par
\par\noindent\hangindent=1em\hangafter=1\hspace*{1em}stem variants before vowel and consonant endings, \hyperref[bk-ch03-learner-ccc3cc60a047]{61}\par
\par\noindent\hangindent=1em\hangafter=1\hspace*{1em}strong declension, \hyperref[bk-ch03-learner-7fee105ca648]{59}\par
\par\noindent\hangindent=1em\hangafter=1\hspace*{1em}substantival use versus a separate noun, \hyperref[bk-ch14-learner-ecac515ca5c6]{457}\par
\par\noindent\hangindent=1em\hangafter=1\hspace*{1em}syncope and exceptions, \hyperref[bk-ch03-learner-1b4dba13d70a]{61}, \hyperref[bk-ch08-learner-56e2a1cf9d42]{209}\par
\par\noindent\hangindent=1em\hangafter=1\hspace*{1em}used as nouns for people, \hyperref[bk-ch10-learner-f4665e71eb65]{287}\par
\par\noindent\hangindent=1em\hangafter=1\hspace*{1em}weak declension and genitive plural, \hyperref[bk-ch03-learner-8c480e64ad30]{59}\par
\par\noindent\hangindent=1em\hangafter=1\hspace*{1em}\emph{See also} \hyperlink{index-2dff10f6a8b7}{Agreement}; \hyperlink{index-571527339d0e}{Comparison}; \hyperlink{index-8692d6f37e2d}{Participles}; \hyperlink{index-c0cf5a1df0a5}{Possessives}; \hyperlink{index-8065348567ea}{Pronominal adjectives}.\par
\BookNeedspace{3\baselineskip}
\par\addvspace{3pt}\noindent\hypertarget{index-cd019f642e7a}{}\textbf{Adverbial clauses}. \emph{See} \hyperlink{index-67f5ffbe4c23}{Clauses, adverbial relations and mood}.\par
\BookNeedspace{3\baselineskip}
\par\addvspace{3pt}\noindent\hypertarget{index-39effafd6188}{}\textbf{Adverbs}\par\nopagebreak[4]
\par\noindent\hangindent=1em\hangafter=1\hspace*{1em}comparison, \hyperref[bk-ch03-learner-b9fb825d30c1]{71}\par
\par\noindent\hangindent=1em\hangafter=1\hspace*{1em}comparison of swīðe, \hyperref[bk-ch15-learner-24b5b5f6352f]{502}\par
\par\noindent\hangindent=1em\hangafter=1\hspace*{1em}direction, time and relation, \hyperref[bk-ch16-learner-5b59859f9b28]{575}\par
\par\noindent\hangindent=1em\hangafter=1\hspace*{1em}formation in -e and -lice, \hyperref[bk-ch03-learner-b9fb825d30c1]{71}\par
\par\noindent\hangindent=1em\hangafter=1\hspace*{1em}frequency, \hyperref[bk-ch09-learner-a6e022a0fd7e]{261}\par
\par\noindent\hangindent=1em\hangafter=1\hspace*{1em}function and absence of agreement, \hyperref[bk-ch03-learner-b99b7da0a77f]{71}\par
\par\noindent\hangindent=1em\hangafter=1\hspace*{1em}intensifying adjectives, \hyperref[bk-ch03-learner-b9fb825d30c1]{71}\par
\par\noindent\hangindent=1em\hangafter=1\hspace*{1em}interrogative manner and reason, \hyperref[bk-ch08-learner-343e3ece044b]{205}\par
\BookNeedspace{3\baselineskip}
\par\addvspace{3pt}\noindent\hypertarget{index-2dff10f6a8b7}{}\textbf{Agreement}\par\nopagebreak[4]
\par\noindent\hangindent=1em\hangafter=1\hspace*{1em}adjective and noun, \hyperref[bk-ch03-learner-dbf98476e10a]{58}\par
\par\noindent\hangindent=1em\hangafter=1\hspace*{1em}article and noun, \hyperref[bk-ch01-learner-f6bfdc1fc647]{13}\par
\par\noindent\hangindent=1em\hangafter=1\hspace*{1em}dual subjects and plural verbs, \hyperref[bk-ch01-learner-f73641fc4bb6]{19}, \hyperref[bk-ch02-learner-b542d326d059]{37}\par
\par\noindent\hangindent=1em\hangafter=1\hspace*{1em}mixed-gender subjects and neuter predicates, \hyperref[bk-ch09-learner-db7c4490e70c]{253}\par
\par\noindent\hangindent=1em\hangafter=1\hspace*{1em}possessive adjective and possessed noun, \hyperref[bk-ch03-learner-0a3e4174931b]{65}\par
\par\noindent\hangindent=1em\hangafter=1\hspace*{1em}predicate noun may differ in number, \hyperref[bk-ch07-learner-73c8a76ced55]{189}\par
\par\noindent\hangindent=1em\hangafter=1\hspace*{1em}relative subject and personal antecedent, \hyperref[bk-ch08-learner-a5e8eb6b52df]{216}\par
\par\noindent\hangindent=1em\hangafter=1\hspace*{1em}subject and finite verb, \hyperref[bk-ch01-learner-814e0a394df4]{19}, \hyperref[bk-ch02-learner-d8738489e241]{44}\par
\par\noindent\hangindent=1em\hangafter=1\hspace*{1em}\emph{See also} \hyperlink{index-dda9f9e22ce1}{Object predicates}; \hyperlink{index-49a25b90c87f}{Passive voice}; \hyperlink{index-b4bb0a2bf943}{Relative clauses}.\par
\BookNeedspace{3\baselineskip}
\par\addvspace{3pt}\noindent\hypertarget{index-0bb3b18d9013}{}\textbf{Alphabet and spelling}\par\nopagebreak[4]
\par\noindent\hangindent=1em\hangafter=1\hspace*{1em}ash, thorn and eth, \hyperref[bk-pronunciation-learner-156b3d5ac4b3]{1}\par
\par\noindent\hangindent=1em\hangafter=1\hspace*{1em}Cook’s hooked vowels, \hyperref[bk-pronunciation-learner-156b3d5ac4b3]{1}\par
\BookNeedspace{3\baselineskip}
\par\addvspace{3pt}\noindent\hypertarget{index-52fd2e54d3bb}{}\textbf{Analogy}\par\nopagebreak[4]
\par\noindent\hangindent=1em\hangafter=1\hspace*{1em}restoration and reshaping, \hyperref[bk-sound-history-learner-1ac916116017]{779}\par
\BookNeedspace{3\baselineskip}
\par\addvspace{3pt}\noindent\hypertarget{index-f2071aa933f8}{}\textbf{Anglo-Frisian brightening}. \emph{See} \hyperlink{index-ff2eeaa96a81}{Sound changes, fronting and restoration}.\par
\BookNeedspace{3\baselineskip}
\par\addvspace{3pt}\noindent\hypertarget{index-c3366c831602}{}\textbf{Apocope}\par\nopagebreak[4]
\par\noindent\hangindent=1em\hangafter=1\hspace*{1em}final high-vowel loss, \hyperref[bk-sound-history-learner-f2eedb8e42ac]{780}\par
\BookNeedspace{3\baselineskip}
\par\addvspace{3pt}\noindent\hypertarget{index-1bf276624b72}{}\textbf{Apposition}\par\nopagebreak[4]
\par\noindent\hangindent=1em\hangafter=1\hspace*{1em}case agreement of names and titles, \hyperref[bk-ch11-learner-ef323531af96]{326}\par
\par\noindent\hangindent=1em\hangafter=1\hspace*{1em}some of them with nominative forms, \hyperref[bk-ch09-learner-07e15514cef1]{238}\par
\BookNeedspace{3\baselineskip}
\par\addvspace{3pt}\noindent\hypertarget{index-29d94922512b}{}\textbf{Article}\par\nopagebreak[4]
\par\noindent\hangindent=1em\hangafter=1\hspace*{1em}complete forms and agreement, \hyperref[bk-ch01-learner-f6bfdc1fc647]{13}\par
\par\noindent\hangindent=1em\hangafter=1\hspace*{1em}demonstrative use, \hyperref[bk-ch03-learner-c532ec2b751d]{77}\par
\par\noindent\hangindent=1em\hangafter=1\hspace*{1em}instrumental reference forms, \hyperref[bk-ch01-learner-f6bfdc1fc647]{13}\par
\par\noindent\hangindent=1em\hangafter=1\hspace*{1em}omission in noun predicates and genitive phrases, \hyperref[bk-ch01-learner-086b075156aa]{19}\par
\BookNeedspace{3\baselineskip}
\par\addvspace{3pt}\noindent\hypertarget{index-bda94d3dce78}{}\textbf{Ash}. \emph{See} \hyperlink{index-0bb3b18d9013}{Alphabet and spelling, ash, thorn and eth}.\par
\BookNeedspace{3\baselineskip}
\par\addvspace{3pt}\noindent\hypertarget{index-5695b23fe738}{}\textbf{Attributive adjectives}. \emph{See} \hyperlink{index-a47426d0da6f}{Adjectives, attributive and predicative uses}.\par
\BookNeedspace{5\baselineskip}
\par\addvspace{6pt}{\large\bfseries B\par}
\BookNeedspace{3\baselineskip}
\par\addvspace{3pt}\noindent\hypertarget{index-28dce19e9a43}{}\textbf{Back mutation}\par\nopagebreak[4]
\par\noindent\hangindent=1em\hangafter=1\hspace*{1em}\hyperref[bk-sound-history-learner-bcd0bead7661]{778}\par
\BookNeedspace{3\baselineskip}
\par\addvspace{3pt}\noindent\hypertarget{index-77ae9160b827}{}\textbf{Backshift}. \emph{See} \hyperlink{index-9f7aadf368e6}{Reported speech, simultaneous, anterior and intended-later events}.\par
\BookNeedspace{3\baselineskip}
\par\addvspace{3pt}\noindent\hypertarget{index-8ceb54144db0}{}\textbf{Bare infinitive}. \emph{See} \hyperlink{index-10a0d424f76e}{Infinitives, plain and inflected forms}.\par
\BookNeedspace{3\baselineskip}
\par\addvspace{3pt}\noindent\hypertarget{index-ef014e50b98c}{}\textbf{Be}. \emph{See} \hyperlink{index-0e45f24c3e51}{Copula}.\par
\BookNeedspace{3\baselineskip}
\par\addvspace{3pt}\noindent\hypertarget{index-d4af3c32590c}{}\textbf{Breaking}\par\nopagebreak[4]
\par\noindent\hangindent=1em\hangafter=1\hspace*{1em}\hyperref[bk-sound-history-learner-bf4f0271a73a]{774}\par
\BookNeedspace{3\baselineskip}
\par\addvspace{3pt}\noindent\hypertarget{index-1511d6a077bb}{}\textbf{Brightening}. \emph{See} \hyperlink{index-ff2eeaa96a81}{Sound changes, fronting and restoration}.\par
\BookNeedspace{5\baselineskip}
\par\addvspace{6pt}{\large\bfseries C\par}
\BookNeedspace{3\baselineskip}
\par\addvspace{3pt}\noindent\hypertarget{index-9b18fd3b7908}{}\textbf{Cardinal numbers}. \emph{See} \hyperlink{index-bd34477bd6c9}{Numerals, cardinals one through twelve}.\par
\BookNeedspace{3\baselineskip}
\par\addvspace{3pt}\noindent\hypertarget{index-15f86e24ed25}{}\textbf{Case}\par\nopagebreak[4]
\par\noindent\hangindent=1em\hangafter=1\hspace*{1em}case of a phrase and its internal dependents, \hyperref[bk-ch10-learner-9d51aaaeaec2]{275}\par
\par\noindent\hangindent=1em\hangafter=1\hspace*{1em}four-case system and ordinary functions, \hyperref[bk-ch01-learner-624dd0eb97a3]{13}\par
\par\noindent\hangindent=1em\hangafter=1\hspace*{1em}identical forms with different functions, \hyperref[bk-ch01-learner-624dd0eb97a3]{13}\par
\par\noindent\hangindent=1em\hangafter=1\hspace*{1em}\emph{See also} \hyperlink{index-bd9e40c9e577}{Accusative}; \hyperlink{index-b3d9506dd56e}{Dative}; \hyperlink{index-b09d245b9fd1}{Genitive}; \hyperlink{index-05082dc1436b}{Instrumental}; \hyperlink{index-b310f849cbca}{Nominative}; \hyperlink{index-f6bcda8b4f94}{Prepositions}.\par
\BookNeedspace{3\baselineskip}
\par\addvspace{3pt}\noindent\hypertarget{index-bae9e1053f31}{}\textbf{Case attraction}. \emph{See} \hyperlink{index-b4bb0a2bf943}{Relative clauses, case attraction: recognition}.\par
\BookNeedspace{3\baselineskip}
\par\addvspace{3pt}\noindent\hypertarget{index-2249bd503ec4}{}\textbf{Cases}. \emph{See} \hyperlink{index-15f86e24ed25}{Case}.\par
\BookNeedspace{3\baselineskip}
\par\addvspace{3pt}\noindent\hypertarget{index-383ebeeaa84c}{}\textbf{Causation}\par\nopagebreak[4]
\par\noindent\hangindent=1em\hangafter=1\hspace*{1em}dōn with a finite ðæt clause, \hyperref[bk-ch16-learner-b481bfab43f6]{573}\par
\BookNeedspace{3\baselineskip}
\par\addvspace{3pt}\noindent\hypertarget{index-82820a58e9e1}{}\textbf{Cause clauses}\par\nopagebreak[4]
\par\noindent\hangindent=1em\hangafter=1\hspace*{1em}factual because with indicative, \hyperref[bk-ch09-learner-a6e022a0fd7e]{261}\par
\par\noindent\hangindent=1em\hangafter=1\hspace*{1em}factual reasons, \hyperref[bk-ch11-learner-640cebcb16ea]{317}\par
\BookNeedspace{3\baselineskip}
\par\addvspace{3pt}\noindent\hypertarget{index-67f5ffbe4c23}{}\textbf{Clauses}\par\nopagebreak[4]
\par\noindent\hangindent=1em\hangafter=1\hspace*{1em}adverbial relations and mood, \hyperref[bk-ch11-learner-d6160974f921]{315}\par
\par\noindent\hangindent=1em\hangafter=1\hspace*{1em}boundaries in long sentences, \hyperref[bk-ch13-learner-ce6c8914b5b0]{408}\par
\BookNeedspace{3\baselineskip}
\par\addvspace{3pt}\noindent\hypertarget{index-4b69fa7c8fb9}{}\textbf{Cognate accusative}. \emph{See} \hyperlink{index-bd9e40c9e577}{Accusative, cognate object}.\par
\BookNeedspace{3\baselineskip}
\par\addvspace{3pt}\noindent\hypertarget{index-b269dc4e81a5}{}\textbf{Commands}\par\nopagebreak[4]
\par\noindent\hangindent=1em\hangafter=1\hspace*{1em}accusative actor and infinitive versus naming, \hyperref[bk-ch09-learner-637a34313b91]{235}\par
\par\noindent\hangindent=1em\hangafter=1\hspace*{1em}dative addressee with bēodan, \hyperref[bk-ch18-learner-a0509728a1ca]{738}\par
\par\noindent\hangindent=1em\hangafter=1\hspace*{1em}\emph{See also} \hyperlink{index-94de9d0bb1ea}{Imperative}; \hyperlink{index-10a0d424f76e}{Infinitives}; \hyperlink{index-49a25b90c87f}{Passive voice}; \hyperlink{index-05ba186b1cc6}{Subjunctive}.\par
\BookNeedspace{3\baselineskip}
\par\addvspace{3pt}\noindent\hypertarget{index-ac57521735f6}{}\textbf{Comparatives}\par\nopagebreak[4]
\par\noindent\hangindent=1em\hangafter=1\hspace*{1em}weak inflection, including predicates, \hyperref[bk-ch11-learner-d7ca3a0f46ee]{333}\par
\BookNeedspace{3\baselineskip}
\par\addvspace{3pt}\noindent\hypertarget{index-571527339d0e}{}\textbf{Comparison}\par\nopagebreak[4]
\par\noindent\hangindent=1em\hangafter=1\hspace*{1em}adverbial degrees, \hyperref[bk-ch03-learner-b9fb825d30c1]{71}\par
\par\noindent\hangindent=1em\hangafter=1\hspace*{1em}comparative genitive plural exceptions, \hyperref[bk-ch03-learner-c8e0934fdf36]{68}\par
\par\noindent\hangindent=1em\hangafter=1\hspace*{1em}comparatives always weak, \hyperref[bk-ch03-learner-c8e0934fdf36]{68}\par
\par\noindent\hangindent=1em\hangafter=1\hspace*{1em}equality: swā ... swā, \hyperref[bk-ch12-learner-514aac1c9ee6]{382}\par
\par\noindent\hangindent=1em\hangafter=1\hspace*{1em}exact equality of amount, \hyperref[bk-ch15-learner-019bcd600b0c]{503}\par
\par\noindent\hangindent=1em\hangafter=1\hspace*{1em}finite swā clauses, \hyperref[bk-ch11-learner-af27acafa7ea]{319}\par
\par\noindent\hangindent=1em\hangafter=1\hspace*{1em}mā as more in number, \hyperref[bk-ch12-learner-d84bed8d05f0]{382}\par
\par\noindent\hangindent=1em\hangafter=1\hspace*{1em}mā modifying an action, \hyperref[bk-ch11-learner-d6ee9c5b27b0]{319}\par
\par\noindent\hangindent=1em\hangafter=1\hspace*{1em}nominal swā phrase, \hyperref[bk-ch12-learner-161b6102df9e]{382}\par
\par\noindent\hangindent=1em\hangafter=1\hspace*{1em}nominative standard with ðonne, \hyperref[bk-ch03-learner-91a7506b6176]{70}\par
\par\noindent\hangindent=1em\hangafter=1\hspace*{1em}preference with parallel prepositional phrases, \hyperref[bk-ch09-learner-e3fb5ca9be3f]{237}\par
\par\noindent\hangindent=1em\hangafter=1\hspace*{1em}preference with swīðor ... ðonne, \hyperref[bk-ch15-learner-66e05fdac64f]{503}\par
\par\noindent\hangindent=1em\hangafter=1\hspace*{1em}proportional: the more ... the more, \hyperref[bk-ch15-learner-eacf155021af]{502}\par
\par\noindent\hangindent=1em\hangafter=1\hspace*{1em}regular, umlaut and irregular adjective series, \hyperref[bk-ch03-learner-3909b65cce85]{68}\par
\par\noindent\hangindent=1em\hangafter=1\hspace*{1em}superlative inflection and use, \hyperref[bk-ch03-learner-2b75c3552b41]{69}\par
\par\noindent\hangindent=1em\hangafter=1\hspace*{1em}\emph{See also} \hyperlink{index-64f3006be77a}{Action substitution}; \hyperlink{index-39effafd6188}{Adverbs}; \hyperlink{index-a0dc191bc1d4}{Preference}; \hyperlink{index-822bab8d41bc}{Quantity}.\par
\BookNeedspace{3\baselineskip}
\par\addvspace{3pt}\noindent\hypertarget{index-c0e8495eb498}{}\textbf{Compensatory lengthening}\par\nopagebreak[4]
\par\noindent\hangindent=1em\hangafter=1\hspace*{1em}\hyperref[bk-sound-history-learner-d2a70dff307f]{781}\par
\BookNeedspace{3\baselineskip}
\par\addvspace{3pt}\noindent\hypertarget{index-52fc470b96aa}{}\textbf{Complement clauses}. \emph{See} \hyperlink{index-6291aef59438}{Content clauses}.\par
\BookNeedspace{3\baselineskip}
\par\addvspace{3pt}\noindent\hypertarget{index-38d8bb29e9f3}{}\textbf{Compounds}. \emph{See} \hyperlink{index-de7476cd8f8e}{Word formation, compounds and the final head}.\par
\BookNeedspace{3\baselineskip}
\par\addvspace{3pt}\noindent\hypertarget{index-2cbbfe890bb5}{}\textbf{Concession}\par\nopagebreak[4]
\par\noindent\hangindent=1em\hangafter=1\hspace*{1em}ðēah and the subjunctive, \hyperref[bk-ch11-learner-3669ec8fac8a]{318}\par
\BookNeedspace{3\baselineskip}
\par\addvspace{3pt}\noindent\hypertarget{index-126995365015}{}\textbf{Conditionals}. \emph{See} \hyperlink{index-97d4be896069}{Conditions}.\par
\BookNeedspace{3\baselineskip}
\par\addvspace{3pt}\noindent\hypertarget{index-97d4be896069}{}\textbf{Conditions}\par\nopagebreak[4]
\par\noindent\hangindent=1em\hangafter=1\hspace*{1em}main clause has its own mood, \hyperref[bk-ch07-learner-49a97724b59e]{175}\par
\par\noindent\hangindent=1em\hangafter=1\hspace*{1em}open and contrary-present, \hyperref[bk-ch11-learner-32cd9ba34c77]{318}\par
\par\noindent\hangindent=1em\hangafter=1\hspace*{1em}open indicative and tentative subjunctive, \hyperref[bk-ch07-learner-49a97724b59e]{175}\par
\BookNeedspace{3\baselineskip}
\par\addvspace{3pt}\noindent\hypertarget{index-2c1e80268124}{}\textbf{Conjugation}. \emph{See} \hyperlink{index-4e436562432f}{Verbs}.\par
\BookNeedspace{3\baselineskip}
\par\addvspace{3pt}\noindent\hypertarget{index-8291944bd54e}{}\textbf{Consonant alternations}\par\nopagebreak[4]
\par\noindent\hangindent=1em\hangafter=1\hspace*{1em}j-present geminates, \hyperref[bk-ch05-learner-67a692c44be4]{115}, \hyperref[bk-ch05-learner-69fa4f57c0e3]{117}\par
\par\noindent\hangindent=1em\hangafter=1\hspace*{1em}present dental combinations and simplification, \hyperref[bk-ch02-learner-1310ddcd10cb]{40}, \hyperref[bk-ch04-learner-3d8a96a1ca23]{90}\par
\par\noindent\hangindent=1em\hangafter=1\hspace*{1em}strong principal-part stems, \hyperref[bk-ch04-learner-3d8a96a1ca23]{90}, \hyperref[bk-ch05-learner-7c5f2c14a3b1]{115}\par
\par\noindent\hangindent=1em\hangafter=1\hspace*{1em}\emph{See also} \hyperlink{index-626458dd9e7f}{Consonant length}; \hyperlink{index-c7020f8dc02f}{Contracted verbs}; \hyperlink{index-a74b1465a508}{Strong verbs}; \hyperlink{index-57ac9d925438}{Weak verbs}.\par
\BookNeedspace{3\baselineskip}
\par\addvspace{3pt}\noindent\hypertarget{index-626458dd9e7f}{}\textbf{Consonant length}\par\nopagebreak[4]
\par\noindent\hangindent=1em\hangafter=1\hspace*{1em}medial doubling and final duration, \hyperref[bk-pronunciation-learner-a18005f2eb82]{7}\par
\BookNeedspace{3\baselineskip}
\par\addvspace{3pt}\noindent\hypertarget{index-182d0cc13b07}{}\textbf{Consonants}\par\nopagebreak[4]
\par\noindent\hangindent=1em\hangafter=1\hspace*{1em}c and g; palatal and hard values, \hyperref[bk-pronunciation-learner-16b43725bb81]{3}\par
\par\noindent\hangindent=1em\hangafter=1\hspace*{1em}continuous g, \hyperref[bk-pronunciation-learner-b426b0db6cb3]{5}\par
\par\noindent\hangindent=1em\hangafter=1\hspace*{1em}h after front and back vowels, \hyperref[bk-pronunciation-learner-fce112ac8644]{5}\par
\par\noindent\hangindent=1em\hangafter=1\hspace*{1em}sounded initial clusters, \hyperref[bk-pronunciation-learner-fce112ac8644]{5}\par
\par\noindent\hangindent=1em\hangafter=1\hspace*{1em}voicing of f, s and þ/ð, \hyperref[bk-pronunciation-learner-74b76bb98627]{6}\par
\BookNeedspace{3\baselineskip}
\par\addvspace{3pt}\noindent\hypertarget{index-6291aef59438}{}\textbf{Content clauses}\par\nopagebreak[4]
\par\noindent\hangindent=1em\hangafter=1\hspace*{1em}anticipatory genitive ðæs with giernan, \hyperref[bk-ch18-learner-0bf6031e8afb]{737}\par
\par\noindent\hangindent=1em\hangafter=1\hspace*{1em}conjunction versus demonstrative and relative ðæt, \hyperref[bk-ch09-learner-a7e68e159a4d]{233}\par
\par\noindent\hangindent=1em\hangafter=1\hspace*{1em}finite ðæt clauses after knowing or saying, \hyperref[bk-ch09-learner-a7e68e159a4d]{233}\par
\par\noindent\hangindent=1em\hangafter=1\hspace*{1em}indicative fact and subjunctive report, \hyperref[bk-ch09-learner-a7e68e159a4d]{233}\par
\par\noindent\hangindent=1em\hangafter=1\hspace*{1em}\emph{See also} \hyperlink{index-9a72221a2747}{Questions}; \hyperlink{index-b4bb0a2bf943}{Relative clauses}; \hyperlink{index-9f7aadf368e6}{Reported speech}.\par
\BookNeedspace{3\baselineskip}
\par\addvspace{3pt}\noindent\hypertarget{index-c7020f8dc02f}{}\textbf{Contracted verbs}\par\nopagebreak[4]
\par\noindent\hangindent=1em\hangafter=1\hspace*{1em}Class I and II indicatives, \hyperref[bk-ch04-learner-04534b37414d]{92}, \hyperref[bk-ch04-learner-cde3298d0700]{94}\par
\par\noindent\hangindent=1em\hangafter=1\hspace*{1em}Class V and VI indicatives, \hyperref[bk-ch05-learner-852de88e21a6]{115}, \hyperref[bk-ch05-learner-02721c417980]{117}\par
\par\noindent\hangindent=1em\hangafter=1\hspace*{1em}Class VII indicative, \hyperref[bk-ch05-learner-be6ba8199b8f]{120}\par
\BookNeedspace{3\baselineskip}
\par\addvspace{3pt}\noindent\hypertarget{index-b993fdb83928}{}\textbf{Contraction}\par\nopagebreak[4]
\par\noindent\hangindent=1em\hangafter=1\hspace*{1em}loss of h and adjacent vowels, \hyperref[bk-sound-history-learner-fb537a53e749]{781}\par
\BookNeedspace{3\baselineskip}
\par\addvspace{3pt}\noindent\hypertarget{index-45c5b055409e}{}\textbf{Control expression}\par\nopagebreak[4]
\par\noindent\hangindent=1em\hangafter=1\hspace*{1em}have in one’s control, \hyperref[bk-ch08-learner-0244beef0db6]{219}\par
\BookNeedspace{3\baselineskip}
\par\addvspace{3pt}\noindent\hypertarget{index-9c0f7ed16a6e}{}\textbf{Coordination}\par\nopagebreak[4]
\par\noindent\hangindent=1em\hangafter=1\hspace*{1em}both ... and: ǣgðer ge ... ge, \hyperref[bk-ch15-learner-b36c32144c55]{502}\par
\par\noindent\hangindent=1em\hangafter=1\hspace*{1em}noun phrases and their articles, \hyperref[bk-ch01-learner-086b075156aa]{19}\par
\par\noindent\hangindent=1em\hangafter=1\hspace*{1em}shared subject with coordinated verbs, \hyperref[bk-ch02-learner-4f6bce65b786]{45}\par
\par\noindent\hangindent=1em\hangafter=1\hspace*{1em}shared subjects and omitted predicates, \hyperref[bk-ch13-learner-4edce9a3cc31]{407}\par
\BookNeedspace{3\baselineskip}
\par\addvspace{3pt}\noindent\hypertarget{index-0e45f24c3e51}{}\textbf{Copula}\par\nopagebreak[4]
\par\noindent\hangindent=1em\hangafter=1\hspace*{1em}alternative bēo present, \hyperref[bk-ch06-learner-666a2007807d]{141}\par
\par\noindent\hangindent=1em\hangafter=1\hspace*{1em}expected, habitual and general states, \hyperref[bk-ch06-learner-666a2007807d]{141}\par
\par\noindent\hangindent=1em\hangafter=1\hspace*{1em}present eom and preterite wæs, \hyperref[bk-ch01-learner-814e0a394df4]{19}\par
\par\noindent\hangindent=1em\hangafter=1\hspace*{1em}\emph{See also} \hyperlink{index-dda9f9e22ce1}{Object predicates}; \hyperlink{index-49a25b90c87f}{Passive voice}.\par
\BookNeedspace{3\baselineskip}
\par\addvspace{3pt}\noindent\hypertarget{index-8e8a3c2f8870}{}\textbf{Counterfactual conditions}. \emph{See} \hyperlink{index-97d4be896069}{Conditions, open and contrary-present}.\par
\BookNeedspace{5\baselineskip}
\par\addvspace{6pt}{\large\bfseries D\par}
\BookNeedspace{3\baselineskip}
\par\addvspace{3pt}\noindent\hypertarget{index-b3d9506dd56e}{}\textbf{Dative}\par\nopagebreak[4]
\par\noindent\hangindent=1em\hangafter=1\hspace*{1em}addressee of saying or telling, \hyperref[bk-ch09-learner-6a9c4e61c4a7]{234}\par
\par\noindent\hangindent=1em\hangafter=1\hspace*{1em}beneficiary and assigned function, \hyperref[bk-ch07-learner-7c9ec2ca6444]{189}\par
\par\noindent\hangindent=1em\hangafter=1\hspace*{1em}experiencer in preference expressions, \hyperref[bk-ch09-learner-57a9fc624ee1]{237}\par
\par\noindent\hangindent=1em\hangafter=1\hspace*{1em}experiencer with ðyncan, \hyperref[bk-ch10-learner-5ca7ed8d1974]{289}\par
\par\noindent\hangindent=1em\hangafter=1\hspace*{1em}object of onfōn and gemiltsian, \hyperref[bk-ch14-learner-48f75fc95810]{470}\par
\par\noindent\hangindent=1em\hangafter=1\hspace*{1em}person helped or benefited, \hyperref[bk-ch04-learner-97ceb9686099]{99}, \hyperref[bk-ch09-learner-2a6c766e81e1]{255}\par
\par\noindent\hangindent=1em\hangafter=1\hspace*{1em}recipient of giving, \hyperref[bk-ch01-learner-d64391595e92]{20}, \hyperref[bk-ch05-learner-95f29519f213]{122}\par
\par\noindent\hangindent=1em\hangafter=1\hspace*{1em}recipient, beneficiary and governed person, \hyperref[bk-ch10-learner-f12b5c0271f8]{289}\par
\par\noindent\hangindent=1em\hangafter=1\hspace*{1em}usefulness, \hyperref[bk-ch09-learner-c78da7bcc201]{253}\par
\par\noindent\hangindent=1em\hangafter=1\hspace*{1em}\emph{See also} \hyperlink{index-a0dc191bc1d4}{Preference}; \hyperlink{index-f6bcda8b4f94}{Prepositions}; \hyperlink{index-fed697438fa8}{Reflexive constructions}.\par
\BookNeedspace{3\baselineskip}
\par\addvspace{3pt}\noindent\hypertarget{index-c2c5121e5a51}{}\textbf{Defective verbs}. \emph{See} \hyperlink{index-67aacf1d4a8e}{Preterite-present verbs, missing forms and citation labels}.\par
\BookNeedspace{3\baselineskip}
\par\addvspace{3pt}\noindent\hypertarget{index-606f5a2a4f8f}{}\textbf{Definite article}. \emph{See} \hyperlink{index-29d94922512b}{Article}.\par
\BookNeedspace{3\baselineskip}
\par\addvspace{3pt}\noindent\hypertarget{index-c190a5ef3668}{}\textbf{Demonstratives}\par\nopagebreak[4]
\par\noindent\hangindent=1em\hangafter=1\hspace*{1em}that and this: agreement and use, \hyperref[bk-ch03-learner-c532ec2b751d]{77}\par
\par\noindent\hangindent=1em\hangafter=1\hspace*{1em}that: complete forms, \hyperref[bk-ch03-learner-59601dfdb337]{77}\par
\par\noindent\hangindent=1em\hangafter=1\hspace*{1em}this: complete forms, \hyperref[bk-ch03-learner-78aeff8dce51]{77}\par
\par\noindent\hangindent=1em\hangafter=1\hspace*{1em}\emph{See also} \hyperlink{index-29d94922512b}{Article}; \hyperlink{index-b4bb0a2bf943}{Relative clauses}.\par
\BookNeedspace{3\baselineskip}
\par\addvspace{3pt}\noindent\hypertarget{index-6879056b12bc}{}\textbf{Dental preterite}. \emph{See} \hyperlink{index-57ac9d925438}{Weak verbs, dental preterite and basic endings}.\par
\BookNeedspace{3\baselineskip}
\par\addvspace{3pt}\noindent\hypertarget{index-aa7412d997f2}{}\textbf{Derivation}. \emph{See} \hyperlink{index-de7476cd8f8e}{Word formation, derivation versus inflection}.\par
\BookNeedspace{3\baselineskip}
\par\addvspace{3pt}\noindent\hypertarget{index-5b88c0af13ba}{}\textbf{Diphthongs}\par\nopagebreak[4]
\par\noindent\hangindent=1em\hangafter=1\hspace*{1em}pronunciation and one-syllable status, \hyperref[bk-pronunciation-learner-1212ca9f676b]{2}\par
\par\noindent\hangindent=1em\hangafter=1\hspace*{1em}short and long, \hyperref[bk-pronunciation-learner-1212ca9f676b]{2}\par
\BookNeedspace{3\baselineskip}
\par\addvspace{3pt}\noindent\hypertarget{index-bfb78853cd67}{}\textbf{Direct address}\par\nopagebreak[4]
\par\noindent\hangindent=1em\hangafter=1\hspace*{1em}nominative noun and weak adjective, \hyperref[bk-ch09-learner-9466d7df4e42]{253}\par
\BookNeedspace{3\baselineskip}
\par\addvspace{3pt}\noindent\hypertarget{index-9e0a49d8d5b1}{}\textbf{Direct speech}\par\nopagebreak[4]
\par\noindent\hangindent=1em\hangafter=1\hspace*{1em}independent quotation and personal reference, \hyperref[bk-ch02-learner-693038b042d5]{45}, \hyperref[bk-ch05-learner-a3edb3da0846]{123}\par
\par\noindent\hangindent=1em\hangafter=1\hspace*{1em}quoted commands, \hyperref[bk-ch06-learner-c13dd1a2bbd6]{150}\par
\par\noindent\hangindent=1em\hangafter=1\hspace*{1em}speaker and addressee reference, \hyperref[bk-ch13-learner-d1987196da60]{407}\par
\BookNeedspace{3\baselineskip}
\par\addvspace{3pt}\noindent\hypertarget{index-9cef35e3a57a}{}\textbf{Distributives}. \emph{See} \hyperlink{index-28104b8f1375}{Pronouns, distributive gehwā and gehwæt}; \hyperlink{index-822bab8d41bc}{Quantity, singular distribution and a plural whole}.\par
\BookNeedspace{3\baselineskip}
\par\addvspace{3pt}\noindent\hypertarget{index-c7125c8226cf}{}\textbf{Do-support}. \emph{See} \hyperlink{index-9a72221a2747}{Questions, direct yes/no questions}.\par
\BookNeedspace{3\baselineskip}
\par\addvspace{3pt}\noindent\hypertarget{index-6c25dc08784f}{}\textbf{Double negatives}. \emph{See} \hyperlink{index-849bc3ea2c32}{Negation, negative concord}.\par
\BookNeedspace{3\baselineskip}
\par\addvspace{3pt}\noindent\hypertarget{index-66ab09544685}{}\textbf{Dual}. \emph{See} \hyperlink{index-bd82cf16699b}{Number, dual pronouns and plural verb agreement}.\par
\BookNeedspace{5\baselineskip}
\par\addvspace{6pt}{\large\bfseries E\par}
\BookNeedspace{3\baselineskip}
\par\addvspace{3pt}\noindent\hypertarget{index-4637e99b28a1}{}\textbf{Ellipsis}\par\nopagebreak[4]
\par\noindent\hangindent=1em\hangafter=1\hspace*{1em}recovering a predicate and its agreement, \hyperref[bk-ch13-learner-9ed625d4233e]{407}\par
\BookNeedspace{3\baselineskip}
\par\addvspace{3pt}\noindent\hypertarget{index-0ab4b2c43447}{}\textbf{Emphasis}\par\nopagebreak[4]
\par\noindent\hangindent=1em\hangafter=1\hspace*{1em}agreeing self, \hyperref[bk-ch08-learner-d0b92b93014c]{212}\par
\BookNeedspace{3\baselineskip}
\par\addvspace{3pt}\noindent\hypertarget{index-1b6043dad3a1}{}\textbf{Equative comparison}. \emph{See} \hyperlink{index-571527339d0e}{Comparison, equality: swā ... swā}.\par
\BookNeedspace{3\baselineskip}
\par\addvspace{3pt}\noindent\hypertarget{index-270d79aad669}{}\textbf{Exhortations}. \emph{See} \hyperlink{index-05ba186b1cc6}{Subjunctive, wishes, exhortations and commands}.\par
\BookNeedspace{5\baselineskip}
\par\addvspace{6pt}{\large\bfseries F\par}
\BookNeedspace{3\baselineskip}
\par\addvspace{3pt}\noindent\hypertarget{index-44709c37368b}{}\textbf{Figurative language}\par\nopagebreak[4]
\par\noindent\hangindent=1em\hangafter=1\hspace*{1em}grammar and contextual interpretation, \hyperref[bk-ch17-learner-4a295f1a1bed]{617}\par
\BookNeedspace{3\baselineskip}
\par\addvspace{3pt}\noindent\hypertarget{index-a75824a81d23}{}\textbf{Free relatives}. \emph{See} \hyperlink{index-b4bb0a2bf943}{Relative clauses, headless what in predicate use}.\par
\BookNeedspace{3\baselineskip}
\par\addvspace{3pt}\noindent\hypertarget{index-e1ed51bbecae}{}\textbf{Free-choice clauses}. \emph{See} \hyperlink{index-b4bb0a2bf943}{Relative clauses, free choice: whoever, whatever, whichever of two}.\par
\BookNeedspace{3\baselineskip}
\par\addvspace{3pt}\noindent\hypertarget{index-45b9dcc617bb}{}\textbf{Front mutation}. \emph{See} \hyperlink{index-f67ed2e92001}{Umlaut}.\par
\BookNeedspace{3\baselineskip}
\par\addvspace{3pt}\noindent\hypertarget{index-301b7d8d2511}{}\textbf{Fronting (sound change)}. \emph{See} \hyperlink{index-ff2eeaa96a81}{Sound changes, fronting and restoration}.\par
\BookNeedspace{3\baselineskip}
\par\addvspace{3pt}\noindent\hypertarget{index-b3f5018ed29d}{}\textbf{Fronting (word order)}. \emph{See} \hyperlink{index-6cab2d8cd0db}{Word order, constituents and fronting}.\par
\BookNeedspace{5\baselineskip}
\par\addvspace{6pt}{\large\bfseries G\par}
\BookNeedspace{3\baselineskip}
\par\addvspace{3pt}\noindent\hypertarget{index-92efc33be7aa}{}\textbf{Geminate consonants}. \emph{See} \hyperlink{index-626458dd9e7f}{Consonant length}.\par
\BookNeedspace{3\baselineskip}
\par\addvspace{3pt}\noindent\hypertarget{index-16cbf1cf52bb}{}\textbf{Gemination}\par\nopagebreak[4]
\par\noindent\hangindent=1em\hangafter=1\hspace*{1em}\hyperref[bk-sound-history-learner-8fefc2cf70de]{780}\par
\BookNeedspace{3\baselineskip}
\par\addvspace{3pt}\noindent\hypertarget{index-a04630ef8bc3}{}\textbf{Gender}\par\nopagebreak[4]
\par\noindent\hangindent=1em\hangafter=1\hspace*{1em}grammatical gender, \hyperref[bk-ch01-learner-624dd0eb97a3]{13}\par
\par\noindent\hangindent=1em\hangafter=1\hspace*{1em}pronoun reference to things, \hyperref[bk-ch01-learner-8f8ad9d5a04c]{19}\par
\BookNeedspace{3\baselineskip}
\par\addvspace{3pt}\noindent\hypertarget{index-b09d245b9fd1}{}\textbf{Genitive}\par\nopagebreak[4]
\par\noindent\hangindent=1em\hangafter=1\hspace*{1em}dependent and possessive, \hyperref[bk-ch10-learner-5d1ddbd64d12]{275}\par
\par\noindent\hangindent=1em\hangafter=1\hspace*{1em}governed by verbs, \hyperref[bk-ch10-learner-4bed4c080848]{276}\par
\par\noindent\hangindent=1em\hangafter=1\hspace*{1em}grounds for a reward, \hyperref[bk-ch17-learner-af5777028dbd]{669}\par
\par\noindent\hangindent=1em\hangafter=1\hspace*{1em}inside a prepositional phrase, \hyperref[bk-ch01-learner-0b109d29614d]{20}\par
\par\noindent\hangindent=1em\hangafter=1\hspace*{1em}measurement of breadth or height, \hyperref[bk-ch12-learner-f5ceef68b943]{382}\par
\par\noindent\hangindent=1em\hangafter=1\hspace*{1em}objective relationship: service, \hyperref[bk-ch09-learner-a6e022a0fd7e]{261}\par
\par\noindent\hangindent=1em\hangafter=1\hspace*{1em}partitive personal pronouns, \hyperref[bk-ch09-learner-555d1541a4a7]{237}\par
\par\noindent\hangindent=1em\hangafter=1\hspace*{1em}partitive whole after a selection word, \hyperref[bk-ch10-learner-2df8a95ea983]{288}\par
\par\noindent\hangindent=1em\hangafter=1\hspace*{1em}personal partitive uncer versus possessive uncer, \hyperref[bk-ch15-learner-8156c28183ed]{502}\par
\par\noindent\hangindent=1em\hangafter=1\hspace*{1em}possession and other relationships, \hyperref[bk-ch01-learner-086b075156aa]{19}\par
\par\noindent\hangindent=1em\hangafter=1\hspace*{1em}requesting and expecting a thing, \hyperref[bk-ch16-learner-9b2cc4c38e6b]{574}\par
\par\noindent\hangindent=1em\hangafter=1\hspace*{1em}subjective and objective, \hyperref[bk-ch10-learner-606f8f7495ea]{276}\par
\par\noindent\hangindent=1em\hangafter=1\hspace*{1em}thing with a dative person, \hyperref[bk-ch10-learner-3d7f287af6c0]{276}\par
\par\noindent\hangindent=1em\hangafter=1\hspace*{1em}with verbs: ceasing and ruling, \hyperref[bk-ch05-learner-e8f7cdddfe3f]{123}, \hyperref[bk-ch07-learner-4a4ffd61ee9a]{170}\par
\par\noindent\hangindent=1em\hangafter=1\hspace*{1em}\emph{See also} \hyperlink{index-c71e90d2f497}{Government}; \hyperlink{index-74d59939e1d0}{Partitive expressions}; \hyperlink{index-18e556f7feed}{Possession}; \hyperlink{index-c0cf5a1df0a5}{Possessives}; \hyperlink{index-b4bb0a2bf943}{Relative clauses}.\par
\BookNeedspace{3\baselineskip}
\par\addvspace{3pt}\noindent\hypertarget{index-c71e90d2f497}{}\textbf{Government}\par\nopagebreak[4]
\par\noindent\hangindent=1em\hangafter=1\hspace*{1em}a prefix does not preserve the base verb’s frame, \hyperref[bk-ch14-learner-ba4182dd4151]{469}\par
\par\noindent\hangindent=1em\hangafter=1\hspace*{1em}\emph{See also} \hyperlink{index-58e39e480c32}{Transitivity and government}.\par
\BookNeedspace{5\baselineskip}
\par\addvspace{6pt}{\large\bfseries H\par}
\BookNeedspace{3\baselineskip}
\par\addvspace{3pt}\noindent\hypertarget{index-5e55784cff7c}{}\textbf{Headless relatives}. \emph{See} \hyperlink{index-b4bb0a2bf943}{Relative clauses, headless what in predicate use}.\par
\BookNeedspace{5\baselineskip}
\par\addvspace{6pt}{\large\bfseries I\par}
\BookNeedspace{3\baselineskip}
\par\addvspace{3pt}\noindent\hypertarget{index-0a79c7efa0b7}{}\textbf{i-umlaut}. \emph{See} \hyperlink{index-f67ed2e92001}{Umlaut}.\par
\BookNeedspace{3\baselineskip}
\par\addvspace{3pt}\noindent\hypertarget{index-94de9d0bb1ea}{}\textbf{Imperative}\par\nopagebreak[4]
\par\noindent\hangindent=1em\hangafter=1\hspace*{1em}addressee, number and omitted subject, \hyperref[bk-ch06-learner-a6c353738baa]{146}\par
\par\noindent\hangindent=1em\hangafter=1\hspace*{1em}contracted verbs and vowel length, \hyperref[bk-ch06-learner-634c423f6df0]{147}, \hyperref[bk-ch06-learner-75a0f72a393d]{148}\par
\par\noindent\hangindent=1em\hangafter=1\hspace*{1em}irregular verbs and be, \hyperref[bk-ch06-learner-1a5c7d53adeb]{149}\par
\par\noindent\hangindent=1em\hangafter=1\hspace*{1em}j-present singulars, \hyperref[bk-ch06-learner-eb4f99a5a301]{148}\par
\par\noindent\hangindent=1em\hangafter=1\hspace*{1em}negative commands and government, \hyperref[bk-ch06-learner-c13dd1a2bbd6]{150}\par
\par\noindent\hangindent=1em\hangafter=1\hspace*{1em}ordinary strong verbs, \hyperref[bk-ch06-learner-634c423f6df0]{147}\par
\par\noindent\hangindent=1em\hangafter=1\hspace*{1em}shortened plural before a pronoun: reference, \hyperref[bk-ch06-learner-6a76a82f4e05]{150}\par
\par\noindent\hangindent=1em\hangafter=1\hspace*{1em}weak Class I models, \hyperref[bk-ch06-learner-ba352ea21d51]{146}\par
\par\noindent\hangindent=1em\hangafter=1\hspace*{1em}weak Classes II and III, \hyperref[bk-ch06-learner-8eef50183c9c]{147}\par
\BookNeedspace{3\baselineskip}
\par\addvspace{3pt}\noindent\hypertarget{index-a09295d4c986}{}\textbf{Impersonal constructions}\par\nopagebreak[4]
\par\noindent\hangindent=1em\hangafter=1\hspace*{1em}hyngran: accusative experiencer, \hyperref[bk-ch10-learner-f11ca067063f]{290}\par
\par\noindent\hangindent=1em\hangafter=1\hspace*{1em}man with a singular verb, \hyperref[bk-ch15-learner-7c10bcfe8c52]{501}\par
\par\noindent\hangindent=1em\hangafter=1\hspace*{1em}ofðyncean: dative experiencer and content clause, \hyperref[bk-ch18-learner-8cf9d181123a]{735}\par
\BookNeedspace{3\baselineskip}
\par\addvspace{3pt}\noindent\hypertarget{index-a3b050b8ae5a}{}\textbf{Indefinite pronouns}\par\nopagebreak[4]
\par\noindent\hangindent=1em\hangafter=1\hspace*{1em}ge-compound reference paradigm, \hyperref[bk-ch08-learner-14917e578635]{213}\par
\BookNeedspace{3\baselineskip}
\par\addvspace{3pt}\noindent\hypertarget{index-7b9279740289}{}\textbf{Indefinite reference}\par\nopagebreak[4]
\par\noindent\hangindent=1em\hangafter=1\hspace*{1em}ān: a certain individual, \hyperref[bk-ch15-learner-3439a8c6c2aa]{501}\par
\par\noindent\hangindent=1em\hangafter=1\hspace*{1em}hwā and hwæt, \hyperref[bk-ch15-learner-bfbf45918c8f]{501}\par
\BookNeedspace{3\baselineskip}
\par\addvspace{3pt}\noindent\hypertarget{index-747f1b861135}{}\textbf{Indicative}\par\nopagebreak[4]
\par\noindent\hangindent=1em\hangafter=1\hspace*{1em}present and preterite meaning, \hyperref[bk-ch02-learner-925bbc29948d]{39}, \hyperref[bk-ch04-learner-28ffaf280840]{89}\par
\BookNeedspace{3\baselineskip}
\par\addvspace{3pt}\noindent\hypertarget{index-39cbc15b8fca}{}\textbf{Indirect questions}. \emph{See} \hyperlink{index-9a72221a2747}{Questions, dependent wh-questions}.\par
\BookNeedspace{3\baselineskip}
\par\addvspace{3pt}\noindent\hypertarget{index-c482110e72bc}{}\textbf{Indirect speech}. \emph{See} \hyperlink{index-9f7aadf368e6}{Reported speech}.\par
\BookNeedspace{3\baselineskip}
\par\addvspace{3pt}\noindent\hypertarget{index-99e56f02d1ac}{}\textbf{Indirect statements}. \emph{See} \hyperlink{index-6291aef59438}{Content clauses}.\par
\BookNeedspace{3\baselineskip}
\par\addvspace{3pt}\noindent\hypertarget{index-10a0d424f76e}{}\textbf{Infinitives}\par\nopagebreak[4]
\par\noindent\hangindent=1em\hangafter=1\hspace*{1em}command with an accusative actor, \hyperref[bk-ch09-learner-637a34313b91]{235}\par
\par\noindent\hangindent=1em\hangafter=1\hspace*{1em}perception with an accusative actor, \hyperref[bk-ch09-learner-22cf9e7f4d0c]{236}\par
\par\noindent\hangindent=1em\hangafter=1\hspace*{1em}plain and inflected forms, \hyperref[bk-ch06-learner-530cdd5762f1]{143}\par
\par\noindent\hangindent=1em\hangafter=1\hspace*{1em}preference expressions, \hyperref[bk-ch09-learner-57a9fc624ee1]{237}\par
\par\noindent\hangindent=1em\hangafter=1\hspace*{1em}purpose after coming, \hyperref[bk-ch06-learner-c3a1991273ab]{145}\par
\par\noindent\hangindent=1em\hangafter=1\hspace*{1em}retained object government, \hyperref[bk-ch06-learner-c3a1991273ab]{145}\par
\par\noindent\hangindent=1em\hangafter=1\hspace*{1em}same-subject modal complement, \hyperref[bk-ch06-learner-d0a4c47b0324]{143}\par
\par\noindent\hangindent=1em\hangafter=1\hspace*{1em}selected -anne and variant -enne, \hyperref[bk-ch06-learner-530cdd5762f1]{143}\par
\par\noindent\hangindent=1em\hangafter=1\hspace*{1em}\emph{See also} \hyperlink{index-b269dc4e81a5}{Commands}; \hyperlink{index-09efa301eda1}{Modal verbs}; \hyperlink{index-bd18e320a6bb}{Perception}; \hyperlink{index-a0dc191bc1d4}{Preference}.\par
\BookNeedspace{3\baselineskip}
\par\addvspace{3pt}\noindent\hypertarget{index-c687f7654360}{}\textbf{Inflected infinitive}. \emph{See} \hyperlink{index-10a0d424f76e}{Infinitives, plain and inflected forms}.\par
\BookNeedspace{3\baselineskip}
\par\addvspace{3pt}\noindent\hypertarget{index-05082dc1436b}{}\textbf{Instrumental}\par\nopagebreak[4]
\par\noindent\hangindent=1em\hangafter=1\hspace*{1em}adjective and demonstrative reference forms, \hyperref[bk-ch03-learner-8c480e64ad30]{59}, \hyperref[bk-ch03-learner-59601dfdb337]{77}\par
\par\noindent\hangindent=1em\hangafter=1\hspace*{1em}bare means and mid with dative, \hyperref[bk-ch10-learner-accee2a2f1b3]{291}\par
\par\noindent\hangindent=1em\hangafter=1\hspace*{1em}distinct forms and shared dative shapes, \hyperref[bk-ch10-learner-c80743ab2698]{291}\par
\par\noindent\hangindent=1em\hangafter=1\hspace*{1em}pronoun reference forms, \hyperref[bk-ch08-learner-14917e578635]{213}\par
\BookNeedspace{3\baselineskip}
\par\addvspace{3pt}\noindent\hypertarget{index-94d858470fde}{}\textbf{Irregular verbs}\par\nopagebreak[4]
\par\noindent\hangindent=1em\hangafter=1\hspace*{1em}do, go and want: complete indicatives, \hyperref[bk-ch06-learner-ee684d5ec1f6]{140}\par
\BookNeedspace{5\baselineskip}
\par\addvspace{6pt}{\large\bfseries J\par}
\BookNeedspace{3\baselineskip}
\par\addvspace{3pt}\noindent\hypertarget{index-b363fe83e3a7}{}\textbf{Jussive}. \emph{See} \hyperlink{index-05ba186b1cc6}{Subjunctive, wishes, exhortations and commands}.\par
\BookNeedspace{5\baselineskip}
\par\addvspace{6pt}{\large\bfseries K\par}
\BookNeedspace{3\baselineskip}
\par\addvspace{3pt}\noindent\hypertarget{index-dcb3e1c00e79}{}\textbf{Knowledge}\par\nopagebreak[4]
\par\noindent\hangindent=1em\hangafter=1\hspace*{1em}knowing a matter versus knowing how, \hyperref[bk-ch06-learner-d0a4c47b0324]{143}, \hyperref[bk-ch06-learner-9fbb6a52092f]{144}\par
\BookNeedspace{5\baselineskip}
\par\addvspace{6pt}{\large\bfseries M\par}
\BookNeedspace{3\baselineskip}
\par\addvspace{3pt}\noindent\hypertarget{index-34368f93eb0d}{}\textbf{Macrons}. \emph{See} \hyperlink{index-59be352618d8}{Vowel length, macrons and duration}.\par
\BookNeedspace{3\baselineskip}
\par\addvspace{3pt}\noindent\hypertarget{index-0e9d5ac7f7bd}{}\textbf{Metathesis}\par\nopagebreak[4]
\par\noindent\hangindent=1em\hangafter=1\hspace*{1em}Class III stem pattern, \hyperref[bk-ch04-learner-d75e819526a2]{94}\par
\BookNeedspace{3\baselineskip}
\par\addvspace{3pt}\noindent\hypertarget{index-09efa301eda1}{}\textbf{Modal verbs}\par\nopagebreak[4]
\par\noindent\hangindent=1em\hangafter=1\hspace*{1em}desire is not future tense, \hyperref[bk-ch06-learner-f11fea9bc36d]{144}\par
\par\noindent\hangindent=1em\hangafter=1\hspace*{1em}finite verb and plain infinitive, \hyperref[bk-ch06-learner-d0a4c47b0324]{143}\par
\par\noindent\hangindent=1em\hangafter=1\hspace*{1em}learned skill and capability, \hyperref[bk-ch06-learner-d0a4c47b0324]{143}\par
\par\noindent\hangindent=1em\hangafter=1\hspace*{1em}meaning distinct from subjunctive mood, \hyperref[bk-ch07-learner-bd7d043cab91]{163}\par
\par\noindent\hangindent=1em\hangafter=1\hspace*{1em}need, daring and obligation, \hyperref[bk-ch06-learner-1041b72208e0]{144}\par
\par\noindent\hangindent=1em\hangafter=1\hspace*{1em}permission and past permission, \hyperref[bk-ch06-learner-95f9e5a43d91]{144}\par
\par\noindent\hangindent=1em\hangafter=1\hspace*{1em}\emph{See also} \hyperlink{index-67aacf1d4a8e}{Preterite-present verbs}; \hyperlink{index-05ba186b1cc6}{Subjunctive}.\par
\BookNeedspace{3\baselineskip}
\par\addvspace{3pt}\noindent\hypertarget{index-780fe43f6e8b}{}\textbf{Mood}\par\nopagebreak[4]
\par\noindent\hangindent=1em\hangafter=1\hspace*{1em}statement, command and wish in praise, \hyperref[bk-ch17-learner-f0d1bccd5033]{667}\par
\BookNeedspace{5\baselineskip}
\par\addvspace{6pt}{\large\bfseries N\par}
\BookNeedspace{3\baselineskip}
\par\addvspace{3pt}\noindent\hypertarget{index-849bc3ea2c32}{}\textbf{Negation}\par\nopagebreak[4]
\par\noindent\hangindent=1em\hangafter=1\hspace*{1em}constituent contrast with nā, \hyperref[bk-ch15-learner-89bd250f0c28]{501}\par
\par\noindent\hangindent=1em\hangafter=1\hspace*{1em}contracted have not: recognition, \hyperref[bk-ch02-learner-79fe9bfe07cd]{43}\par
\par\noindent\hangindent=1em\hangafter=1\hspace*{1em}finite ne and negative be-series, \hyperref[bk-ch15-learner-435744864c61]{499}\par
\par\noindent\hangindent=1em\hangafter=1\hspace*{1em}nabban: negative possession, \hyperref[bk-ch15-learner-7af15a67a1b7]{499}\par
\par\noindent\hangindent=1em\hangafter=1\hspace*{1em}nān: strong declension, \hyperref[bk-ch15-learner-1ec05076f1c3]{500}\par
\par\noindent\hangindent=1em\hangafter=1\hspace*{1em}ne before a finite verb, \hyperref[bk-ch02-learner-833c9271fac6]{45}\par
\par\noindent\hangindent=1em\hangafter=1\hspace*{1em}negative concord, \hyperref[bk-ch15-learner-73465f380192]{500}\par
\par\noindent\hangindent=1em\hangafter=1\hspace*{1em}negative imperatives, \hyperref[bk-ch06-learner-c13dd1a2bbd6]{150}\par
\par\noindent\hangindent=1em\hangafter=1\hspace*{1em}negative willan: unwillingness and refusal, \hyperref[bk-ch15-learner-ffdec0152920]{499}\par
\par\noindent\hangindent=1em\hangafter=1\hspace*{1em}unwillingness and refusal with willan, \hyperref[bk-ch09-learner-f4a46510a7af]{237}\par
\par\noindent\hangindent=1em\hangafter=1\hspace*{1em}\emph{See also} \hyperlink{index-94de9d0bb1ea}{Imperative}; \hyperlink{index-09efa301eda1}{Modal verbs}.\par
\BookNeedspace{3\baselineskip}
\par\addvspace{3pt}\noindent\hypertarget{index-b310f849cbca}{}\textbf{Nominative}\par\nopagebreak[4]
\par\noindent\hangindent=1em\hangafter=1\hspace*{1em}subjects and noun predicates, \hyperref[bk-ch01-learner-086b075156aa]{19}\par
\par\noindent\hangindent=1em\hangafter=1\hspace*{1em}personal-pronoun subjects, \hyperref[bk-ch01-learner-f4792afbd8d6]{19}\par
\BookNeedspace{3\baselineskip}
\par\addvspace{3pt}\noindent\hypertarget{index-d85ab1133877}{}\textbf{Nonfinite forms}\par\nopagebreak[4]
\par\noindent\hangindent=1em\hangafter=1\hspace*{1em}infinitive and participle versus finite verb, \hyperref[bk-ch06-learner-530cdd5762f1]{143}, \hyperref[bk-ch07-learner-ae50c7e839f0]{176}\par
\BookNeedspace{3\baselineskip}
\par\addvspace{3pt}\noindent\hypertarget{index-38dee9ba5543}{}\textbf{Noun declension}. \emph{See} \hyperlink{index-888b2617c1be}{Nouns}.\par
\BookNeedspace{3\baselineskip}
\par\addvspace{3pt}\noindent\hypertarget{index-888b2617c1be}{}\textbf{Nouns}\par\nopagebreak[4]
\par\noindent\hangindent=1em\hangafter=1\hspace*{1em}-nes and doubled s, \hyperref[bk-ch14-learner-2f64e0f5ae82]{456}\par
\par\noindent\hangindent=1em\hangafter=1\hspace*{1em}abstracts in -u, \hyperref[bk-ch15-learner-c5bb0898ad26]{496}\par
\par\noindent\hangindent=1em\hangafter=1\hspace*{1em}agent nouns in -end, \hyperref[bk-ch09-learner-a5032b3d8e24]{247}\par
\par\noindent\hangindent=1em\hangafter=1\hspace*{1em}count and mass senses, \hyperref[bk-ch03-learner-aad11cd34c34]{63}, \hyperref[bk-ch08-learner-f078b5ed9d15]{218}\par
\par\noindent\hangindent=1em\hangafter=1\hspace*{1em}feminines in -nes and -ung, \hyperref[bk-ch01-learner-71357d7f7267]{16}\par
\par\noindent\hangindent=1em\hangafter=1\hspace*{1em}frēond and its alternating stem, \hyperref[bk-ch13-learner-890b82e6326f]{413}\par
\par\noindent\hangindent=1em\hangafter=1\hspace*{1em}h-loss and compensatory lengthening: seolh, \hyperref[bk-ch11-learner-8c7893ec15d2]{323}\par
\par\noindent\hangindent=1em\hangafter=1\hspace*{1em}hand and woruld: endingless accusative, \hyperref[bk-ch10-learner-df6021e47304]{279}\par
\par\noindent\hangindent=1em\hangafter=1\hspace*{1em}irregular case and number patterns, \hyperref[bk-ch01-learner-3d6b15e7b6bf]{17}\par
\par\noindent\hangindent=1em\hangafter=1\hspace*{1em}kinship: fæder, \hyperref[bk-ch10-learner-53a33b32add5]{278}\par
\par\noindent\hangindent=1em\hangafter=1\hspace*{1em}kinship: mōdor, dohtor and brōðor, \hyperref[bk-ch13-learner-19b96c59b02d]{410}, \hyperref[bk-ch13-learner-912457bc89a0]{413}\par
\par\noindent\hangindent=1em\hangafter=1\hspace*{1em}kinship: sweostor, \hyperref[bk-ch15-learner-bff20a9dbee5]{496}\par
\par\noindent\hangindent=1em\hangafter=1\hspace*{1em}masculines in -e, \hyperref[bk-ch01-learner-246e855423d2]{14}\par
\par\noindent\hangindent=1em\hangafter=1\hspace*{1em}neuter nominative–accusative identity, \hyperref[bk-ch01-learner-737b70e12605]{15}\par
\par\noindent\hangindent=1em\hangafter=1\hspace*{1em}neuter w-stems and exceptional plural, \hyperref[bk-ch01-learner-737b70e12605]{15}\par
\par\noindent\hangindent=1em\hangafter=1\hspace*{1em}neuters in -e, \hyperref[bk-ch09-learner-47d8491e816c]{247}\par
\par\noindent\hangindent=1em\hangafter=1\hspace*{1em}plural-only nouns, \hyperref[bk-ch01-learner-3d6b15e7b6bf]{17}\par
\par\noindent\hangindent=1em\hangafter=1\hspace*{1em}root: bōc and tōð, \hyperref[bk-ch15-learner-2acd7f38336d]{495}\par
\par\noindent\hangindent=1em\hangafter=1\hspace*{1em}root: burg, \hyperref[bk-ch13-learner-317b0ddeb02e]{413}\par
\par\noindent\hangindent=1em\hangafter=1\hspace*{1em}root: gāt and its vowel alternation, \hyperref[bk-ch10-learner-55e62ee30134]{279}\par
\par\noindent\hangindent=1em\hangafter=1\hspace*{1em}shortened stems before endings, \hyperref[bk-ch10-learner-4c244822ba6e]{277}\par
\par\noindent\hangindent=1em\hangafter=1\hspace*{1em}strong feminine declension, \hyperref[bk-ch01-learner-71357d7f7267]{16}\par
\par\noindent\hangindent=1em\hangafter=1\hspace*{1em}strong masculine declension, \hyperref[bk-ch01-learner-246e855423d2]{14}\par
\par\noindent\hangindent=1em\hangafter=1\hspace*{1em}strong neuter declension and plural -u, \hyperref[bk-ch01-learner-737b70e12605]{15}\par
\par\noindent\hangindent=1em\hangafter=1\hspace*{1em}syncope and retained internal vowels, \hyperref[bk-ch01-learner-246e855423d2]{14}, \hyperref[bk-ch01-learner-737b70e12605]{15}\par
\par\noindent\hangindent=1em\hangafter=1\hspace*{1em}syncope: wuldor, \hyperref[bk-ch14-learner-fdc0479f56b0]{460}\par
\par\noindent\hangindent=1em\hangafter=1\hspace*{1em}u-stem declension, \hyperref[bk-ch09-learner-6f36190f5839]{246}\par
\par\noindent\hangindent=1em\hangafter=1\hspace*{1em}vowel alternation by number: bæð, geat, sīðfæt, \hyperref[bk-ch11-learner-7e1199ab8347]{321}, \hyperref[bk-ch11-learner-20cd80791fd4]{322}\par
\par\noindent\hangindent=1em\hangafter=1\hspace*{1em}weak declension, \hyperref[bk-ch01-learner-9c5fd7a11f03]{17}\par
\par\noindent\hangindent=1em\hangafter=1\hspace*{1em}æ/a stem alternation, \hyperref[bk-ch01-learner-246e855423d2]{14}\par
\par\noindent\hangindent=1em\hangafter=1\hspace*{1em}\emph{See also} \hyperlink{index-15f86e24ed25}{Case}; \hyperlink{index-a04630ef8bc3}{Gender}; \hyperlink{index-bd82cf16699b}{Number}; \hyperlink{index-aea538abc20e}{Syllable weight}; \hyperlink{index-0c2d377d343c}{Syncope}; \hyperlink{index-f67ed2e92001}{Umlaut}; \hyperlink{index-de7476cd8f8e}{Word formation}.\par
\BookNeedspace{3\baselineskip}
\par\addvspace{3pt}\noindent\hypertarget{index-bd82cf16699b}{}\textbf{Number}\par\nopagebreak[4]
\par\noindent\hangindent=1em\hangafter=1\hspace*{1em}count nouns and numeral agreement, \hyperref[bk-ch03-learner-8a4d537a7515]{72}\par
\par\noindent\hangindent=1em\hangafter=1\hspace*{1em}dual pronouns and plural verb agreement, \hyperref[bk-ch01-learner-f73641fc4bb6]{19}, \hyperref[bk-ch02-learner-b542d326d059]{37}\par
\BookNeedspace{3\baselineskip}
\par\addvspace{3pt}\noindent\hypertarget{index-bd34477bd6c9}{}\textbf{Numerals}\par\nopagebreak[4]
\par\noindent\hangindent=1em\hangafter=1\hspace*{1em}cardinals one through twelve, \hyperref[bk-ch03-learner-8a4d537a7515]{72}\par
\par\noindent\hangindent=1em\hangafter=1\hspace*{1em}four through twelve: invariant before nouns, \hyperref[bk-ch03-learner-2aa77dbd105d]{74}\par
\par\noindent\hangindent=1em\hangafter=1\hspace*{1em}one: singular cardinal use and reference plural, \hyperref[bk-ch03-learner-c138fbf9cb87]{72}\par
\par\noindent\hangindent=1em\hangafter=1\hspace*{1em}ordinals first through seventh, \hyperref[bk-ch03-learner-482b1a79fb62]{75}\par
\par\noindent\hangindent=1em\hangafter=1\hspace*{1em}ordinary weak ordinal pattern, \hyperref[bk-ch03-learner-f4e2c8b803ee]{75}\par
\par\noindent\hangindent=1em\hangafter=1\hspace*{1em}second remains strong, \hyperref[bk-ch03-learner-d9d0ac45cb4c]{76}\par
\par\noindent\hangindent=1em\hangafter=1\hspace*{1em}two and three: gender and case, \hyperref[bk-ch03-learner-02a6416f7b36]{73}\par
\BookNeedspace{5\baselineskip}
\par\addvspace{6pt}{\large\bfseries O\par}
\BookNeedspace{3\baselineskip}
\par\addvspace{3pt}\noindent\hypertarget{index-dda9f9e22ce1}{}\textbf{Object predicates}\par\nopagebreak[4]
\par\noindent\hangindent=1em\hangafter=1\hspace*{1em}adjective after perception, \hyperref[bk-ch10-learner-d733552c7139]{297}\par
\par\noindent\hangindent=1em\hangafter=1\hspace*{1em}calling someone blessed, \hyperref[bk-ch17-learner-cbd56e38041e]{668}\par
\par\noindent\hangindent=1em\hangafter=1\hspace*{1em}making equal and a dative standard, \hyperref[bk-ch15-learner-249aa2db805e]{517}\par
\par\noindent\hangindent=1em\hangafter=1\hspace*{1em}naming with two accusatives, \hyperref[bk-ch07-learner-645dc7063f24]{189}\par
\par\noindent\hangindent=1em\hangafter=1\hspace*{1em}strong adjective agreement, \hyperref[bk-ch09-learner-33765ec333c9]{236}\par
\BookNeedspace{3\baselineskip}
\par\addvspace{3pt}\noindent\hypertarget{index-7e0d7142ece5}{}\textbf{Objective genitive}. \emph{See} \hyperlink{index-b09d245b9fd1}{Genitive, subjective and objective}.\par
\BookNeedspace{3\baselineskip}
\par\addvspace{3pt}\noindent\hypertarget{index-fb99971e0209}{}\textbf{Obligation}. \emph{See} \hyperlink{index-09efa301eda1}{Modal verbs, need, daring and obligation}.\par
\BookNeedspace{3\baselineskip}
\par\addvspace{3pt}\noindent\hypertarget{index-dc160f109bd5}{}\textbf{Optative}. \emph{See} \hyperlink{index-05ba186b1cc6}{Subjunctive}.\par
\BookNeedspace{3\baselineskip}
\par\addvspace{3pt}\noindent\hypertarget{index-0eec9803ce7a}{}\textbf{Ordinal numbers}. \emph{See} \hyperlink{index-bd34477bd6c9}{Numerals, ordinals first through seventh}.\par
\BookNeedspace{3\baselineskip}
\par\addvspace{3pt}\noindent\hypertarget{index-9d43c453e386}{}\textbf{Orthography}. \emph{See} \hyperlink{index-0bb3b18d9013}{Alphabet and spelling}.\par
\BookNeedspace{5\baselineskip}
\par\addvspace{6pt}{\large\bfseries P\par}
\BookNeedspace{3\baselineskip}
\par\addvspace{3pt}\noindent\hypertarget{index-b2bdccc52d3d}{}\textbf{Palatal diphthongization}\par\nopagebreak[4]
\par\noindent\hangindent=1em\hangafter=1\hspace*{1em}\hyperref[bk-sound-history-learner-43e30165bf50]{778}\par
\BookNeedspace{3\baselineskip}
\par\addvspace{3pt}\noindent\hypertarget{index-ab24b5f8ac59}{}\textbf{Palatalization}\par\nopagebreak[4]
\par\noindent\hangindent=1em\hangafter=1\hspace*{1em}consonants and change order, \hyperref[bk-sound-history-learner-5d52be86de50]{779}\par
\BookNeedspace{3\baselineskip}
\par\addvspace{3pt}\noindent\hypertarget{index-8692d6f37e2d}{}\textbf{Participles}\par\nopagebreak[4]
\par\noindent\hangindent=1em\hangafter=1\hspace*{1em}agreement with the object in habban constructions, \hyperref[bk-ch12-learner-e89b8e1c8555]{359}\par
\par\noindent\hangindent=1em\hangafter=1\hspace*{1em}attributive strong and weak agreement, \hyperref[bk-ch07-learner-5cc5caec500a]{177}\par
\par\noindent\hangindent=1em\hangafter=1\hspace*{1em}bare and agreeing plural predicates, \hyperref[bk-ch17-learner-5930a132dc85]{667}\par
\par\noindent\hangindent=1em\hangafter=1\hspace*{1em}base forms and passive predicates, \hyperref[bk-ch11-learner-be45aa398fa7]{338}\par
\par\noindent\hangindent=1em\hangafter=1\hspace*{1em}Class II and monosyllabic dental models, \hyperref[bk-ch07-learner-3ddd770922b3]{181}, \hyperref[bk-ch07-learner-de8dd5a7122c]{183}\par
\par\noindent\hangindent=1em\hangafter=1\hspace*{1em}distinct from agent nouns in -end, \hyperref[bk-ch07-learner-5cc5caec500a]{177}, \hyperref[bk-ch09-learner-a5032b3d8e24]{247}\par
\par\noindent\hangindent=1em\hangafter=1\hspace*{1em}naming people while retaining verbal complements, \hyperref[bk-ch17-learner-07c1bbec4fe0]{667}\par
\par\noindent\hangindent=1em\hangafter=1\hspace*{1em}nonfinite status and formation, \hyperref[bk-ch07-learner-ae50c7e839f0]{176}\par
\par\noindent\hangindent=1em\hangafter=1\hspace*{1em}present versus past stems, \hyperref[bk-ch07-learner-3fa3aaeef76b]{177}\par
\par\noindent\hangindent=1em\hangafter=1\hspace*{1em}present-participle adjective model, \hyperref[bk-ch07-learner-22bdebc13439]{178}\par
\par\noindent\hangindent=1em\hangafter=1\hspace*{1em}strong past-participle model and syncope, \hyperref[bk-ch07-learner-ea64958329b3]{179}\par
\par\noindent\hangindent=1em\hangafter=1\hspace*{1em}verbal predicate versus adjective agreement, \hyperref[bk-ch07-learner-2de790126d8b]{184}\par
\par\noindent\hangindent=1em\hangafter=1\hspace*{1em}weak dental model and syncope, \hyperref[bk-ch07-learner-49e61c42cd78]{180}\par
\par\noindent\hangindent=1em\hangafter=1\hspace*{1em}wesan with intransitive past participle (recognition), \hyperref[bk-ch12-learner-a01000ad790b]{360}\par
\par\noindent\hangindent=1em\hangafter=1\hspace*{1em}wesan with present participle (recognition), \hyperref[bk-ch12-learner-d03b0ae4da78]{360}\par
\par\noindent\hangindent=1em\hangafter=1\hspace*{1em}\emph{See also} \hyperlink{index-a47426d0da6f}{Adjectives}; \hyperlink{index-888b2617c1be}{Nouns}; \hyperlink{index-49a25b90c87f}{Passive voice}; \hyperlink{index-0c2d377d343c}{Syncope}; \hyperlink{index-2cacf6c14ec8}{Tense}.\par
\BookNeedspace{3\baselineskip}
\par\addvspace{3pt}\noindent\hypertarget{index-74d59939e1d0}{}\textbf{Partitive expressions}\par\nopagebreak[4]
\par\noindent\hangindent=1em\hangafter=1\hspace*{1em}personal genitive, \hyperref[bk-ch09-learner-555d1541a4a7]{237}\par
\BookNeedspace{3\baselineskip}
\par\addvspace{3pt}\noindent\hypertarget{index-2c0cff128b81}{}\textbf{Partitive genitive}. \emph{See} \hyperlink{index-b09d245b9fd1}{Genitive, partitive whole after a selection word}.\par
\BookNeedspace{3\baselineskip}
\par\addvspace{3pt}\noindent\hypertarget{index-78c04330e989}{}\textbf{Passive}. \emph{See} \hyperlink{index-49a25b90c87f}{Passive voice}.\par
\BookNeedspace{3\baselineskip}
\par\addvspace{3pt}\noindent\hypertarget{index-49a25b90c87f}{}\textbf{Passive voice}\par\nopagebreak[4]
\par\noindent\hangindent=1em\hangafter=1\hspace*{1em}agentless clauses and the affected subject, \hyperref[bk-ch12-learner-313b0bb61c54]{360}\par
\par\noindent\hangindent=1em\hangafter=1\hspace*{1em}auxiliary choice and state or event, \hyperref[bk-ch07-learner-8d010318c225]{184}\par
\par\noindent\hangindent=1em\hangafter=1\hspace*{1em}commands and wishes, \hyperref[bk-ch07-learner-162da4968438]{185}\par
\par\noindent\hangindent=1em\hangafter=1\hspace*{1em}event and state readings, \hyperref[bk-ch12-learner-20f46332400d]{360}\par
\par\noindent\hangindent=1em\hangafter=1\hspace*{1em}formation and subject, \hyperref[bk-ch07-learner-8d010318c225]{184}\par
\par\noindent\hangindent=1em\hangafter=1\hspace*{1em}participle inflection and noninflection, \hyperref[bk-ch07-learner-2de790126d8b]{184}\par
\par\noindent\hangindent=1em\hangafter=1\hspace*{1em}retained genitive of possessions with berēafian, \hyperref[bk-ch18-learner-83e42b478fb6]{737}\par
\BookNeedspace{3\baselineskip}
\par\addvspace{3pt}\noindent\hypertarget{index-bd18e320a6bb}{}\textbf{Perception}\par\nopagebreak[4]
\par\noindent\hangindent=1em\hangafter=1\hspace*{1em}accusative actor and infinitive, \hyperref[bk-ch09-learner-22cf9e7f4d0c]{236}\par
\par\noindent\hangindent=1em\hangafter=1\hspace*{1em}accusative actor with infinitive, \hyperref[bk-ch10-learner-c4f790ec0bcb]{298}\par
\par\noindent\hangindent=1em\hangafter=1\hspace*{1em}infinitive actor versus finite-clause subject, \hyperref[bk-ch16-learner-88ba24c6485d]{573}\par
\BookNeedspace{3\baselineskip}
\par\addvspace{3pt}\noindent\hypertarget{index-f002f4dfc0aa}{}\textbf{Perfect constructions}. \emph{See} \hyperlink{index-8692d6f37e2d}{Participles, wesan with intransitive past participle (recognition)}; \hyperlink{index-2cacf6c14ec8}{Tense, completed action with habban}.\par
\BookNeedspace{3\baselineskip}
\par\addvspace{3pt}\noindent\hypertarget{index-229efc8f5263}{}\textbf{Permission}. \emph{See} \hyperlink{index-09efa301eda1}{Modal verbs, permission and past permission}.\par
\BookNeedspace{3\baselineskip}
\par\addvspace{3pt}\noindent\hypertarget{index-b527f140565b}{}\textbf{Personal pronouns}\par\nopagebreak[4]
\par\noindent\hangindent=1em\hangafter=1\hspace*{1em}case determined by clause function, \hyperref[bk-ch02-learner-8f105697c218]{38}\par
\par\noindent\hangindent=1em\hangafter=1\hspace*{1em}complete singular, dual and plural forms, \hyperref[bk-ch02-learner-b542d326d059]{37}\par
\par\noindent\hangindent=1em\hangafter=1\hspace*{1em}first- and second-person genitives as reference forms, \hyperref[bk-ch02-learner-b542d326d059]{37}\par
\par\noindent\hangindent=1em\hangafter=1\hspace*{1em}grammatical gender and reference, \hyperref[bk-ch01-learner-8f8ad9d5a04c]{19}, \hyperref[bk-ch02-learner-fc8f17ff0711]{38}\par
\par\noindent\hangindent=1em\hangafter=1\hspace*{1em}invariant third-person possessors, \hyperref[bk-ch02-learner-62b2d81d3d34]{38}\par
\par\noindent\hangindent=1em\hangafter=1\hspace*{1em}reflexive use, \hyperref[bk-ch02-learner-8fc55662d2c6]{44}\par
\par\noindent\hangindent=1em\hangafter=1\hspace*{1em}nominative subject forms (singular, dual and plural), \hyperref[bk-ch01-learner-f4792afbd8d6]{19}\par
\par\noindent\hangindent=1em\hangafter=1\hspace*{1em}\emph{See also} \hyperlink{index-2dff10f6a8b7}{Agreement}; \hyperlink{index-c0cf5a1df0a5}{Possessives}; \hyperlink{index-fed697438fa8}{Reflexive constructions}.\par
\BookNeedspace{3\baselineskip}
\par\addvspace{3pt}\noindent\hypertarget{index-219196e8bf0a}{}\textbf{Pluperfect meaning}. \emph{See} \hyperlink{index-2cacf6c14ec8}{Tense, completed action with habban: past viewpoint}; \hyperlink{index-2cacf6c14ec8}{Tense, preterite with anterior meaning}.\par
\BookNeedspace{3\baselineskip}
\par\addvspace{3pt}\noindent\hypertarget{index-18e556f7feed}{}\textbf{Possession}\par\nopagebreak[4]
\par\noindent\hangindent=1em\hangafter=1\hspace*{1em}predicative genitive and agreeing possessive, \hyperref[bk-ch12-learner-d4a523783302]{382}\par
\BookNeedspace{3\baselineskip}
\par\addvspace{3pt}\noindent\hypertarget{index-c0cf5a1df0a5}{}\textbf{Possessives}\par\nopagebreak[4]
\par\noindent\hangindent=1em\hangafter=1\hspace*{1em}dual models and syncope, \hyperref[bk-ch03-learner-c24379f75ac8]{66}\par
\par\noindent\hangindent=1em\hangafter=1\hspace*{1em}first- and second-person singular models, \hyperref[bk-ch03-learner-51b68bb635a4]{65}\par
\par\noindent\hangindent=1em\hangafter=1\hspace*{1em}invariant third-person genitives, \hyperref[bk-ch02-learner-62b2d81d3d34]{38}\par
\par\noindent\hangindent=1em\hangafter=1\hspace*{1em}plural models and syncope, \hyperref[bk-ch03-learner-ed122a88dd3d]{67}\par
\par\noindent\hangindent=1em\hangafter=1\hspace*{1em}possessor stem versus agreement ending, \hyperref[bk-ch03-learner-0a3e4174931b]{65}\par
\BookNeedspace{3\baselineskip}
\par\addvspace{3pt}\noindent\hypertarget{index-62b69d723b9f}{}\textbf{Predicative adjectives}. \emph{See} \hyperlink{index-a47426d0da6f}{Adjectives, attributive and predicative uses}.\par
\BookNeedspace{3\baselineskip}
\par\addvspace{3pt}\noindent\hypertarget{index-a0dc191bc1d4}{}\textbf{Preference}\par\nopagebreak[4]
\par\noindent\hangindent=1em\hangafter=1\hspace*{1em}dative experiencer and infinitive, \hyperref[bk-ch09-learner-57a9fc624ee1]{237}\par
\BookNeedspace{3\baselineskip}
\par\addvspace{3pt}\noindent\hypertarget{index-6f20e73b9302}{}\textbf{Prefixes}\par\nopagebreak[4]
\par\noindent\hangindent=1em\hangafter=1\hspace*{1em}lexical identity and limits of automatic formation, \hyperref[bk-ch05-learner-4385a8053534]{111}, \hyperref[bk-ch07-learner-ae50c7e839f0]{176}\par
\par\noindent\hangindent=1em\hangafter=1\hspace*{1em}\emph{See also} \hyperlink{index-de7476cd8f8e}{Word formation, prefixes and their lexical limits}.\par
\BookNeedspace{3\baselineskip}
\par\addvspace{3pt}\noindent\hypertarget{index-f6bcda8b4f94}{}\textbf{Prepositions}\par\nopagebreak[4]
\par\noindent\hangindent=1em\hangafter=1\hspace*{1em}above, below and in the midst, \hyperref[bk-ch07-learner-bae927b6adfb]{188}\par
\par\noindent\hangindent=1em\hangafter=1\hspace*{1em}accompaniment: mid with dative, \hyperref[bk-ch09-learner-e3fb5ca9be3f]{237}\par
\par\noindent\hangindent=1em\hangafter=1\hspace*{1em}assigned function: tō with dative, \hyperref[bk-ch07-learner-7c9ec2ca6444]{189}\par
\par\noindent\hangindent=1em\hangafter=1\hspace*{1em}case of a phrase containing a genitive, \hyperref[bk-ch10-learner-98bdf06c8402]{299}\par
\par\noindent\hangindent=1em\hangafter=1\hspace*{1em}location: on with dative, \hyperref[bk-ch01-learner-086b075156aa]{19}\par
\par\noindent\hangindent=1em\hangafter=1\hspace*{1em}numbered day: on with accusative, \hyperref[bk-ch09-learner-a6e022a0fd7e]{261}\par
\par\noindent\hangindent=1em\hangafter=1\hspace*{1em}on: the right-side idiom, \hyperref[bk-ch10-learner-902af315de33]{298}\par
\par\noindent\hangindent=1em\hangafter=1\hspace*{1em}separation: fram with dative, \hyperref[bk-ch01-learner-2b2114170c4e]{20}\par
\par\noindent\hangindent=1em\hangafter=1\hspace*{1em}spatial over and authority over, \hyperref[bk-ch07-learner-844cd74c2c5d]{188}\par
\BookNeedspace{3\baselineskip}
\par\addvspace{3pt}\noindent\hypertarget{index-9f89be8d7ed9}{}\textbf{Present stem changes}\par\nopagebreak[4]
\par\noindent\hangindent=1em\hangafter=1\hspace*{1em}second and third singular vowels, \hyperref[bk-ch04-learner-d95b114af330]{92}, \hyperref[bk-ch05-learner-4385a8053534]{111}\par
\par\noindent\hangindent=1em\hangafter=1\hspace*{1em}\emph{See also} \hyperlink{index-f67ed2e92001}{Umlaut}.\par
\BookNeedspace{3\baselineskip}
\par\addvspace{3pt}\noindent\hypertarget{index-67aacf1d4a8e}{}\textbf{Preterite-present verbs}\par\nopagebreak[4]
\par\noindent\hangindent=1em\hangafter=1\hspace*{1em}missing forms and citation labels, \hyperref[bk-ch06-learner-494df9f87625]{139}\par
\par\noindent\hangindent=1em\hangafter=1\hspace*{1em}present meaning and dental preterites, \hyperref[bk-ch06-learner-f7a979796f68]{138}\par
\BookNeedspace{3\baselineskip}
\par\addvspace{3pt}\noindent\hypertarget{index-3f182014da0b}{}\textbf{Principal parts}\par\nopagebreak[4]
\par\noindent\hangindent=1em\hangafter=1\hspace*{1em}citations versus finite forms, \hyperref[bk-ch04-learner-28ffaf280840]{89}\par
\par\noindent\hangindent=1em\hangafter=1\hspace*{1em}strong verbs and stem positions, \hyperref[bk-ch04-learner-28ffaf280840]{89}\par
\BookNeedspace{3\baselineskip}
\par\addvspace{3pt}\noindent\hypertarget{index-3619034d11d5}{}\textbf{Pro-verb dōn}. \emph{See} \hyperlink{index-64f3006be77a}{Action substitution, dōn repeats an earlier action}.\par
\BookNeedspace{3\baselineskip}
\par\addvspace{3pt}\noindent\hypertarget{index-dbde49fc47f5}{}\textbf{Progressive translation}. \emph{See} \hyperlink{index-8692d6f37e2d}{Participles, wesan with present participle (recognition)}.\par
\BookNeedspace{3\baselineskip}
\par\addvspace{3pt}\noindent\hypertarget{index-d65c09e289ea}{}\textbf{Prohibitions}. \emph{See} \hyperlink{index-94de9d0bb1ea}{Imperative, negative commands and government}.\par
\BookNeedspace{3\baselineskip}
\par\addvspace{3pt}\noindent\hypertarget{index-8065348567ea}{}\textbf{Pronominal adjectives}\par\nopagebreak[4]
\par\noindent\hangindent=1em\hangafter=1\hspace*{1em}other or another: strong declension, \hyperref[bk-ch08-learner-a575cb1bd9c9]{210}\par
\par\noindent\hangindent=1em\hangafter=1\hspace*{1em}same: weak declension only, \hyperref[bk-ch08-learner-788b7927080a]{211}\par
\par\noindent\hangindent=1em\hangafter=1\hspace*{1em}self: strong and weak agreement, \hyperref[bk-ch08-learner-d0b92b93014c]{212}\par
\BookNeedspace{3\baselineskip}
\par\addvspace{3pt}\noindent\hypertarget{index-28104b8f1375}{}\textbf{Pronouns}\par\nopagebreak[4]
\par\noindent\hangindent=1em\hangafter=1\hspace*{1em}bēgen: both, with plural inflection, \hyperref[bk-ch16-learner-02c84df7b74b]{563}\par
\par\noindent\hangindent=1em\hangafter=1\hspace*{1em}distributive gehwā and gehwæt, \hyperref[bk-ch15-learner-f424e7a8d086]{501}\par
\par\noindent\hangindent=1em\hangafter=1\hspace*{1em}\emph{See also} \hyperlink{index-a3b050b8ae5a}{Indefinite pronouns}; \hyperlink{index-7b9279740289}{Indefinite reference}; \hyperlink{index-b527f140565b}{Personal pronouns}; \hyperlink{index-b4bb0a2bf943}{Relative clauses}.\par
\BookNeedspace{3\baselineskip}
\par\addvspace{3pt}\noindent\hypertarget{index-2fdc7b9bbf5c}{}\textbf{Purpose clauses}\par\nopagebreak[4]
\par\noindent\hangindent=1em\hangafter=1\hspace*{1em}finite ðæt clauses and their subjects, \hyperref[bk-ch11-learner-0621543f1a64]{317}\par
\par\noindent\hangindent=1em\hangafter=1\hspace*{1em}negative: lest, \hyperref[bk-ch11-learner-5c0d35483e2b]{317}\par
\BookNeedspace{5\baselineskip}
\par\addvspace{6pt}{\large\bfseries Q\par}
\BookNeedspace{3\baselineskip}
\par\addvspace{3pt}\noindent\hypertarget{index-c0ed786005d1}{}\textbf{Quantifiers}. \emph{See} \hyperlink{index-822bab8d41bc}{Quantity}.\par
\BookNeedspace{3\baselineskip}
\par\addvspace{3pt}\noindent\hypertarget{index-822bab8d41bc}{}\textbf{Quantity}\par\nopagebreak[4]
\par\noindent\hangindent=1em\hangafter=1\hspace*{1em}a certain, some, any, each and such, \hyperref[bk-ch08-learner-c37835ca5690]{207}\par
\par\noindent\hangindent=1em\hangafter=1\hspace*{1em}all and each: strong agreement, \hyperref[bk-ch07-learner-6ca488fddd27]{186}\par
\par\noindent\hangindent=1em\hangafter=1\hspace*{1em}fela with genitive plural, \hyperref[bk-ch10-learner-92c9166ec67f]{288}\par
\par\noindent\hangindent=1em\hangafter=1\hspace*{1em}many with plural nouns, \hyperref[bk-ch09-learner-23e51d68091c]{251}\par
\par\noindent\hangindent=1em\hangafter=1\hspace*{1em}singular distribution and a plural whole, \hyperref[bk-ch15-learner-ca8a7ca630d6]{502}\par
\par\noindent\hangindent=1em\hangafter=1\hspace*{1em}substantival micel and māre, \hyperref[bk-ch15-learner-6ec683c21516]{503}\par
\BookNeedspace{3\baselineskip}
\par\addvspace{3pt}\noindent\hypertarget{index-d3075d033caf}{}\textbf{Quantity marks}. \emph{See} \hyperlink{index-59be352618d8}{Vowel length, macrons and duration}.\par
\BookNeedspace{3\baselineskip}
\par\addvspace{3pt}\noindent\hypertarget{index-9a72221a2747}{}\textbf{Questions}\par\nopagebreak[4]
\par\noindent\hangindent=1em\hangafter=1\hspace*{1em}case determined inside the question, \hyperref[bk-ch08-learner-f4ef5d866f80]{204}, \hyperref[bk-ch09-learner-6a9c4e61c4a7]{234}\par
\par\noindent\hangindent=1em\hangafter=1\hspace*{1em}dependent wh-questions, \hyperref[bk-ch09-learner-6a9c4e61c4a7]{234}\par
\par\noindent\hangindent=1em\hangafter=1\hspace*{1em}dependent whether versus inflected which of two, \hyperref[bk-ch09-learner-6a9c4e61c4a7]{234}\par
\par\noindent\hangindent=1em\hangafter=1\hspace*{1em}direct yes/no questions, \hyperref[bk-ch08-learner-da4d22fbf997]{203}\par
\par\noindent\hangindent=1em\hangafter=1\hspace*{1em}how versus why, \hyperref[bk-ch08-learner-343e3ece044b]{205}\par
\par\noindent\hangindent=1em\hangafter=1\hspace*{1em}interrogative phrases and constituent order, \hyperref[bk-ch08-learner-da4d22fbf997]{203}\par
\par\noindent\hangindent=1em\hangafter=1\hspace*{1em}mood in dependent questions, \hyperref[bk-ch09-learner-6a9c4e61c4a7]{234}\par
\par\noindent\hangindent=1em\hangafter=1\hspace*{1em}which, which of two and what sort, \hyperref[bk-ch08-learner-e70ad5862678]{205}\par
\par\noindent\hangindent=1em\hangafter=1\hspace*{1em}who and what: complete singular paradigm, \hyperref[bk-ch08-learner-f4ef5d866f80]{204}\par
\par\noindent\hangindent=1em\hangafter=1\hspace*{1em}\emph{See also} \hyperlink{index-8065348567ea}{Pronominal adjectives}; \hyperlink{index-b4bb0a2bf943}{Relative clauses}; \hyperlink{index-6cab2d8cd0db}{Word order}.\par
\BookNeedspace{5\baselineskip}
\par\addvspace{6pt}{\large\bfseries R\par}
\BookNeedspace{3\baselineskip}
\par\addvspace{3pt}\noindent\hypertarget{index-08d96720e75f}{}\textbf{Raising}\par\nopagebreak[4]
\par\noindent\hangindent=1em\hangafter=1\hspace*{1em}older e-to-i change versus later umlaut, \hyperref[bk-sound-history-learner-02dd4c54579e]{777}\par
\BookNeedspace{3\baselineskip}
\par\addvspace{3pt}\noindent\hypertarget{index-ba80f1f434d8}{}\textbf{Reason clauses}. \emph{See} \hyperlink{index-82820a58e9e1}{Cause clauses}.\par
\BookNeedspace{3\baselineskip}
\par\addvspace{3pt}\noindent\hypertarget{index-14947a56d0bb}{}\textbf{Reconstruction}\par\nopagebreak[4]
\par\noindent\hangindent=1em\hangafter=1\hspace*{1em}\hyperref[bk-sound-history-learner-fd30ba546e22]{773}\par
\BookNeedspace{3\baselineskip}
\par\addvspace{3pt}\noindent\hypertarget{index-34f460cee99d}{}\textbf{Reduplicating verbs}. \emph{See} \hyperlink{index-a74b1465a508}{Strong verbs, Class VII: pasts in ēo and ē}.\par
\BookNeedspace{3\baselineskip}
\par\addvspace{3pt}\noindent\hypertarget{index-71bf90935fc2}{}\textbf{Reference}\par\nopagebreak[4]
\par\noindent\hangindent=1em\hangafter=1\hspace*{1em}ambiguous antecedents, \hyperref[bk-ch13-learner-848f09dddf39]{406}\par
\par\noindent\hangindent=1em\hangafter=1\hspace*{1em}gender, number and case of pronouns, \hyperref[bk-ch13-learner-1c8bf8495a8f]{406}\par
\par\noindent\hangindent=1em\hangafter=1\hspace*{1em}\emph{See also} \hyperlink{index-9e0a49d8d5b1}{Direct speech}; \hyperlink{index-28104b8f1375}{Pronouns}; \hyperlink{index-9f7aadf368e6}{Reported speech}.\par
\BookNeedspace{3\baselineskip}
\par\addvspace{3pt}\noindent\hypertarget{index-fed697438fa8}{}\textbf{Reflexive constructions}\par\nopagebreak[4]
\par\noindent\hangindent=1em\hangafter=1\hspace*{1em}accusative object: conceal oneself, \hyperref[bk-ch18-learner-dbac285fef20]{739}\par
\par\noindent\hangindent=1em\hangafter=1\hspace*{1em}accusative with gebiddan, \hyperref[bk-ch14-learner-61ffa9daf947]{470}\par
\par\noindent\hangindent=1em\hangafter=1\hspace*{1em}beneficiary with tilian, \hyperref[bk-ch14-learner-1ec6cf84d79b]{470}\par
\par\noindent\hangindent=1em\hangafter=1\hspace*{1em}dative with a genitive feared thing, \hyperref[bk-ch18-learner-a19c79a380b1]{738}\par
\par\noindent\hangindent=1em\hangafter=1\hspace*{1em}dative with bewęndan, \hyperref[bk-ch16-learner-9a119a3340b1]{574}\par
\par\noindent\hangindent=1em\hangafter=1\hspace*{1em}ordinary personal pronouns, \hyperref[bk-ch02-learner-8fc55662d2c6]{44}\par
\par\noindent\hangindent=1em\hangafter=1\hspace*{1em}reinforcement with agreeing self, \hyperref[bk-ch08-learner-cdb9f6787322]{213}\par
\BookNeedspace{3\baselineskip}
\par\addvspace{3pt}\noindent\hypertarget{index-f4e88fdae1d6}{}\textbf{Reflexive reference}. \emph{See} \hyperlink{index-fed697438fa8}{Reflexive constructions}.\par
\BookNeedspace{3\baselineskip}
\par\addvspace{3pt}\noindent\hypertarget{index-b4bb0a2bf943}{}\textbf{Relative clauses}\par\nopagebreak[4]
\par\noindent\hangindent=1em\hangafter=1\hspace*{1em}antecedent and relative strategies, \hyperref[bk-ch08-learner-287d30ac3868]{214}\par
\par\noindent\hangindent=1em\hangafter=1\hspace*{1em}antecedent gender and number; internal case, \hyperref[bk-ch08-learner-748837dad0a8]{215}\par
\par\noindent\hangindent=1em\hangafter=1\hspace*{1em}case attraction: recognition, \hyperref[bk-ch08-learner-4171ac2e2f18]{216}\par
\par\noindent\hangindent=1em\hangafter=1\hspace*{1em}free choice: whoever, whatever, whichever of two, \hyperref[bk-ch15-learner-81d02f820a47]{502}\par
\par\noindent\hangindent=1em\hangafter=1\hspace*{1em}headless what in predicate use, \hyperref[bk-ch09-learner-5ae9fa426c67]{237}\par
\par\noindent\hangindent=1em\hangafter=1\hspace*{1em}indefinite or possible antecedents and the subjunctive, \hyperref[bk-ch17-learner-273409eecb31]{668}\par
\par\noindent\hangindent=1em\hangafter=1\hspace*{1em}inflected relative with or without ðe, \hyperref[bk-ch08-learner-748837dad0a8]{215}\par
\par\noindent\hangindent=1em\hangafter=1\hspace*{1em}internal genitive despite an accusative antecedent, \hyperref[bk-ch18-learner-49add5c35eec]{736}\par
\par\noindent\hangindent=1em\hangafter=1\hspace*{1em}invariant ðe, \hyperref[bk-ch08-learner-149b83a2f42c]{215}\par
\par\noindent\hangindent=1em\hangafter=1\hspace*{1em}person agreement with personal antecedent, \hyperref[bk-ch08-learner-a5e8eb6b52df]{216}\par
\par\noindent\hangindent=1em\hangafter=1\hspace*{1em}possessive relative, \hyperref[bk-ch08-learner-9b001c83b737]{216}\par
\par\noindent\hangindent=1em\hangafter=1\hspace*{1em}subjunctive in an optative setting, \hyperref[bk-ch18-learner-bbf20820beab]{739}\par
\par\noindent\hangindent=1em\hangafter=1\hspace*{1em}temporal: the day when, \hyperref[bk-ch11-learner-c34cc0292c16]{319}\par
\par\noindent\hangindent=1em\hangafter=1\hspace*{1em}ðe with a personal pronoun: reference, \hyperref[bk-ch08-learner-4171ac2e2f18]{216}\par
\par\noindent\hangindent=1em\hangafter=1\hspace*{1em}\emph{See also} \hyperlink{index-2dff10f6a8b7}{Agreement}; \hyperlink{index-6291aef59438}{Content clauses}; \hyperlink{index-9a72221a2747}{Questions}; \hyperlink{index-71bf90935fc2}{Reference}; \hyperlink{index-05ba186b1cc6}{Subjunctive}.\par
\BookNeedspace{3\baselineskip}
\par\addvspace{3pt}\noindent\hypertarget{index-9f7aadf368e6}{}\textbf{Reported speech}\par\nopagebreak[4]
\par\noindent\hangindent=1em\hangafter=1\hspace*{1em}anterior habban construction, \hyperref[bk-ch12-learner-234095e10d2a]{361}\par
\par\noindent\hangindent=1em\hangafter=1\hspace*{1em}person and place viewpoint, \hyperref[bk-ch16-learner-531548ae175f]{572}\par
\par\noindent\hangindent=1em\hangafter=1\hspace*{1em}person and simultaneous time reference, \hyperref[bk-ch09-learner-dc85c3b339be]{234}\par
\par\noindent\hangindent=1em\hangafter=1\hspace*{1em}simultaneous, anterior and intended-later events, \hyperref[bk-ch12-learner-f19bc233592b]{361}\par
\par\noindent\hangindent=1em\hangafter=1\hspace*{1em}willan and intended action, \hyperref[bk-ch12-learner-339d6a4d03a5]{362}\par
\BookNeedspace{3\baselineskip}
\par\addvspace{3pt}\noindent\hypertarget{index-487af01af2dd}{}\textbf{Result clauses}\par\nopagebreak[4]
\par\noindent\hangindent=1em\hangafter=1\hspace*{1em}asserted result versus intended purpose, \hyperref[bk-ch11-learner-678b89a87f59]{318}\par
\BookNeedspace{5\baselineskip}
\par\addvspace{6pt}{\large\bfseries S\par}
\BookNeedspace{3\baselineskip}
\par\addvspace{3pt}\noindent\hypertarget{index-fcae13f0e231}{}\textbf{Smoothing}\par\nopagebreak[4]
\par\noindent\hangindent=1em\hangafter=1\hspace*{1em}Late West Saxon and Anglian, \hyperref[bk-sound-history-learner-ffea67e44b42]{779}\par
\BookNeedspace{3\baselineskip}
\par\addvspace{3pt}\noindent\hypertarget{index-ff2eeaa96a81}{}\textbf{Sound changes}\par\nopagebreak[4]
\par\noindent\hangindent=1em\hangafter=1\hspace*{1em}\hyperref[bk-sound-history-learner-5a5f7fd73e0c]{773}\par
\par\noindent\hangindent=1em\hangafter=1\hspace*{1em}fronting and restoration, \hyperref[bk-sound-history-learner-af4fa27ce6c1]{774}\par
\BookNeedspace{3\baselineskip}
\par\addvspace{3pt}\noindent\hypertarget{index-2319cd90fbab}{}\textbf{Stress}\par\nopagebreak[4]
\par\noindent\hangindent=1em\hangafter=1\hspace*{1em}word and prefix stress, \hyperref[bk-pronunciation-learner-a18005f2eb82]{7}\par
\BookNeedspace{3\baselineskip}
\par\addvspace{3pt}\noindent\hypertarget{index-a74b1465a508}{}\textbf{Strong verbs}\par\nopagebreak[4]
\par\noindent\hangindent=1em\hangafter=1\hspace*{1em}Class I, \hyperref[bk-ch04-learner-70278eafbcc8]{91}\par
\par\noindent\hangindent=1em\hangafter=1\hspace*{1em}Class II, \hyperref[bk-ch04-learner-d95b114af330]{92}\par
\par\noindent\hangindent=1em\hangafter=1\hspace*{1em}Class III and short-vowel subpatterns, \hyperref[bk-ch04-learner-d75e819526a2]{94}\par
\par\noindent\hangindent=1em\hangafter=1\hspace*{1em}Class IV, \hyperref[bk-ch05-learner-8f7dcc852b82]{112}\par
\par\noindent\hangindent=1em\hangafter=1\hspace*{1em}Class IV nasal-stem exceptions, \hyperref[bk-ch05-learner-a2089fd14fa6]{113}\par
\par\noindent\hangindent=1em\hangafter=1\hspace*{1em}Class V, \hyperref[bk-ch05-learner-db2ab453e454]{113}\par
\par\noindent\hangindent=1em\hangafter=1\hspace*{1em}Class V palatal series, \hyperref[bk-ch05-learner-9cc23aad7df2]{115}\par
\par\noindent\hangindent=1em\hangafter=1\hspace*{1em}Class VI, \hyperref[bk-ch05-learner-92ffb8f8186c]{116}\par
\par\noindent\hangindent=1em\hangafter=1\hspace*{1em}Class VII: pasts in ēo and ē, \hyperref[bk-ch05-learner-55e28af3ed17]{118}, \hyperref[bk-ch05-learner-748db6b8492f]{120}\par
\par\noindent\hangindent=1em\hangafter=1\hspace*{1em}exceptional third singular preterite, \hyperref[bk-ch04-learner-b7272f944b5d]{96}\par
\par\noindent\hangindent=1em\hangafter=1\hspace*{1em}j-present stems, \hyperref[bk-ch05-learner-67a692c44be4]{115}, \hyperref[bk-ch05-learner-69fa4f57c0e3]{117}\par
\par\noindent\hangindent=1em\hangafter=1\hspace*{1em}present and preterite endings, \hyperref[bk-ch04-learner-3d8a96a1ca23]{90}\par
\par\noindent\hangindent=1em\hangafter=1\hspace*{1em}recognition-only exceptional systems, \hyperref[bk-ch04-learner-01b371fdc98a]{97}\par
\par\noindent\hangindent=1em\hangafter=1\hspace*{1em}second singular preterite uses plural stem, \hyperref[bk-ch04-learner-3d8a96a1ca23]{90}, \hyperref[bk-ch05-learner-4385a8053534]{111}\par
\par\noindent\hangindent=1em\hangafter=1\hspace*{1em}stem-vowel alternation and principal parts, \hyperref[bk-ch04-learner-28ffaf280840]{89}\par
\par\noindent\hangindent=1em\hangafter=1\hspace*{1em}\emph{See also} \hyperlink{index-8291944bd54e}{Consonant alternations}; \hyperlink{index-c7020f8dc02f}{Contracted verbs}; \hyperlink{index-94de9d0bb1ea}{Imperative}; \hyperlink{index-3f182014da0b}{Principal parts}; \hyperlink{index-05ba186b1cc6}{Subjunctive}.\par
\BookNeedspace{3\baselineskip}
\par\addvspace{3pt}\noindent\hypertarget{index-15b897b5a172}{}\textbf{Subject pronouns}. \emph{See} \hyperlink{index-b527f140565b}{Personal pronouns, nominative subject forms (singular, dual and plural)}.\par
\BookNeedspace{3\baselineskip}
\par\addvspace{3pt}\noindent\hypertarget{index-cbd03dfde58a}{}\textbf{Subjective genitive}. \emph{See} \hyperlink{index-b09d245b9fd1}{Genitive, subjective and objective}.\par
\BookNeedspace{3\baselineskip}
\par\addvspace{3pt}\noindent\hypertarget{index-05ba186b1cc6}{}\textbf{Subjunctive}\par\nopagebreak[4]
\par\noindent\hangindent=1em\hangafter=1\hspace*{1em}anticipated before and until, \hyperref[bk-ch11-learner-a018097f3125]{316}\par
\par\noindent\hangindent=1em\hangafter=1\hspace*{1em}before an event anticipated from a past viewpoint, \hyperref[bk-ch17-learner-41a025a3a7a7]{669}\par
\par\noindent\hangindent=1em\hangafter=1\hspace*{1em}commands with unnan, \hyperref[bk-ch10-learner-4160a02bc5fd]{297}\par
\par\noindent\hangindent=1em\hangafter=1\hspace*{1em}contrary-present conditions, \hyperref[bk-ch11-learner-126b97da7db2]{319}\par
\par\noindent\hangindent=1em\hangafter=1\hspace*{1em}dependent questions, \hyperref[bk-ch09-learner-6a9c4e61c4a7]{234}\par
\par\noindent\hangindent=1em\hangafter=1\hspace*{1em}expectation and a possible provider, \hyperref[bk-ch16-learner-71e2b7958144]{570}\par
\par\noindent\hangindent=1em\hangafter=1\hspace*{1em}nonpast preterite reference in a supplied condition, \hyperref[bk-ch07-learner-19b1b69abd23]{175}\par
\par\noindent\hangindent=1em\hangafter=1\hspace*{1em}present and preterite endings and stems, \hyperref[bk-ch07-learner-813bcaedf4f5]{164}\par
\par\noindent\hangindent=1em\hangafter=1\hspace*{1em}preterite-presents and irregular verbs, \hyperref[bk-ch07-learner-a651adcc0fcb]{169}\par
\par\noindent\hangindent=1em\hangafter=1\hspace*{1em}reported statements, \hyperref[bk-ch09-learner-a7e68e159a4d]{233}\par
\par\noindent\hangindent=1em\hangafter=1\hspace*{1em}strong-verb systems, \hyperref[bk-ch07-learner-2bd176407f8c]{166}, \hyperref[bk-ch07-learner-2c917f3ea037]{167}\par
\par\noindent\hangindent=1em\hangafter=1\hspace*{1em}unreal present wishes, \hyperref[bk-ch09-learner-84672e8deefc]{236}\par
\par\noindent\hangindent=1em\hangafter=1\hspace*{1em}weak-verb systems, \hyperref[bk-ch07-learner-407c09325356]{164}, \hyperref[bk-ch07-learner-eb3e06d33ada]{165}\par
\par\noindent\hangindent=1em\hangafter=1\hspace*{1em}wishes, exhortations and commands, \hyperref[bk-ch07-learner-bd7d043cab91]{163}\par
\par\noindent\hangindent=1em\hangafter=1\hspace*{1em}\emph{See also} \hyperlink{index-2cbbfe890bb5}{Concession}; \hyperlink{index-97d4be896069}{Conditions}; \hyperlink{index-6291aef59438}{Content clauses}; \hyperlink{index-09efa301eda1}{Modal verbs}; \hyperlink{index-2fdc7b9bbf5c}{Purpose clauses}; \hyperlink{index-9a72221a2747}{Questions}; \hyperlink{index-9f7aadf368e6}{Reported speech}; \hyperlink{index-403cc21502eb}{Wishes}.\par
\BookNeedspace{3\baselineskip}
\par\addvspace{3pt}\noindent\hypertarget{index-464d963bc323}{}\textbf{Substantival adjectives}. \emph{See} \hyperlink{index-a47426d0da6f}{Adjectives, bare strong forms naming people}; \hyperlink{index-a47426d0da6f}{Adjectives, used as nouns for people}.\par
\BookNeedspace{3\baselineskip}
\par\addvspace{3pt}\noindent\hypertarget{index-7fb23f65d163}{}\textbf{Substantival participles}. \emph{See} \hyperlink{index-8692d6f37e2d}{Participles, naming people while retaining verbal complements}.\par
\BookNeedspace{3\baselineskip}
\par\addvspace{3pt}\noindent\hypertarget{index-95ef5f9316e2}{}\textbf{Suffixes}. \emph{See} \hyperlink{index-de7476cd8f8e}{Word formation, adjectives in -lic}; \hyperlink{index-de7476cd8f8e}{Word formation, adverbs in -e and -lice}; \hyperlink{index-de7476cd8f8e}{Word formation, noun suffixes: -ung, -scipe, -nes, -dōm}.\par
\BookNeedspace{3\baselineskip}
\par\addvspace{3pt}\noindent\hypertarget{index-559abc673ae8}{}\textbf{Suppletion}\par\nopagebreak[4]
\par\noindent\hangindent=1em\hangafter=1\hspace*{1em}past of go, \hyperref[bk-ch06-learner-ee684d5ec1f6]{140}\par
\BookNeedspace{3\baselineskip}
\par\addvspace{3pt}\noindent\hypertarget{index-aea538abc20e}{}\textbf{Syllable weight}\par\nopagebreak[4]
\par\noindent\hangindent=1em\hangafter=1\hspace*{1em}light and heavy stems, \hyperref[bk-ch01-learner-737b70e12605]{15}, \hyperref[bk-ch04-learner-c034d18fad71]{98}, \hyperref[bk-sound-history-learner-771e0549912c]{780}\par
\BookNeedspace{3\baselineskip}
\par\addvspace{3pt}\noindent\hypertarget{index-4275dee070b1}{}\textbf{Syllables}\par\nopagebreak[4]
\par\noindent\hangindent=1em\hangafter=1\hspace*{1em}vowel count and consonant doubling, \hyperref[bk-pronunciation-learner-a18005f2eb82]{7}\par
\BookNeedspace{3\baselineskip}
\par\addvspace{3pt}\noindent\hypertarget{index-0c2d377d343c}{}\textbf{Syncope}\par\nopagebreak[4]
\par\noindent\hangindent=1em\hangafter=1\hspace*{1em}adjective and possessive stems, \hyperref[bk-ch03-learner-1b4dba13d70a]{61}, \hyperref[bk-ch03-learner-c24379f75ac8]{66}\par
\par\noindent\hangindent=1em\hangafter=1\hspace*{1em}internal vowel loss and analogy, \hyperref[bk-sound-history-learner-00c58def41fa]{781}\par
\par\noindent\hangindent=1em\hangafter=1\hspace*{1em}noun stems and exceptions, \hyperref[bk-ch01-learner-246e855423d2]{14}, \hyperref[bk-ch01-learner-737b70e12605]{15}\par
\par\noindent\hangindent=1em\hangafter=1\hspace*{1em}participle stems, \hyperref[bk-ch07-learner-ea64958329b3]{179}, \hyperref[bk-ch07-learner-49e61c42cd78]{180}\par
\BookNeedspace{3\baselineskip}
\par\addvspace{3pt}\noindent\hypertarget{index-afc1766c0274}{}\textbf{Syncretism}. \emph{See} \hyperlink{index-15f86e24ed25}{Case, identical forms with different functions}.\par
\BookNeedspace{5\baselineskip}
\par\addvspace{6pt}{\large\bfseries T\par}
\BookNeedspace{3\baselineskip}
\par\addvspace{3pt}\noindent\hypertarget{index-2cacf6c14ec8}{}\textbf{Tense}\par\nopagebreak[4]
\par\noindent\hangindent=1em\hangafter=1\hspace*{1em}completed action with habban, \hyperref[bk-ch12-learner-7f0ea6e7cad6]{358}\par
\par\noindent\hangindent=1em\hangafter=1\hspace*{1em}completed action with habban: past viewpoint, \hyperref[bk-ch12-learner-6f6cc62dab0b]{358}\par
\par\noindent\hangindent=1em\hangafter=1\hspace*{1em}narrative succession, \hyperref[bk-ch12-learner-24ebe58b96d5]{357}\par
\par\noindent\hangindent=1em\hangafter=1\hspace*{1em}present with future reference, \hyperref[bk-ch10-learner-5a53bc6a06e5]{298}\par
\par\noindent\hangindent=1em\hangafter=1\hspace*{1em}preterite with anterior meaning, \hyperref[bk-ch12-learner-e075379e48d6]{358}, \hyperref[bk-ch18-learner-8fe39e530b47]{739}\par
\BookNeedspace{3\baselineskip}
\par\addvspace{3pt}\noindent\hypertarget{index-0000f38892e0}{}\textbf{Thorn and eth}. \emph{See} \hyperlink{index-0bb3b18d9013}{Alphabet and spelling, ash, thorn and eth}.\par
\BookNeedspace{3\baselineskip}
\par\addvspace{3pt}\noindent\hypertarget{index-ae1903628a34}{}\textbf{Time clauses}\par\nopagebreak[4]
\par\noindent\hangindent=1em\hangafter=1\hspace*{1em}anticipated before and until, \hyperref[bk-ch11-learner-8e868f4d3742]{316}\par
\par\noindent\hangindent=1em\hangafter=1\hspace*{1em}mid-ðȳ, \hyperref[bk-ch11-learner-82c1c5623ba7]{316}\par
\par\noindent\hangindent=1em\hangafter=1\hspace*{1em}past when ... then: ðā, \hyperref[bk-ch11-learner-a850a4fba7f6]{315}\par
\par\noindent\hangindent=1em\hangafter=1\hspace*{1em}prospective when with subjunctive, \hyperref[bk-ch18-learner-0a7d2c4cd048]{739}\par
\par\noindent\hangindent=1em\hangafter=1\hspace*{1em}recurring or prospective when: ðonne, \hyperref[bk-ch11-learner-336ad2620b98]{315}\par
\par\noindent\hangindent=1em\hangafter=1\hspace*{1em}while and simultaneous events, \hyperref[bk-ch11-learner-8e5721384100]{316}\par
\BookNeedspace{3\baselineskip}
\par\addvspace{3pt}\noindent\hypertarget{index-a0d846a0e9c6}{}\textbf{Time expressions}\par\nopagebreak[4]
\par\noindent\hangindent=1em\hangafter=1\hspace*{1em}on a numbered day, \hyperref[bk-ch09-learner-a6e022a0fd7e]{261}\par
\BookNeedspace{3\baselineskip}
\par\addvspace{3pt}\noindent\hypertarget{index-58e39e480c32}{}\textbf{Transitivity and government}\par\nopagebreak[4]
\par\noindent\hangindent=1em\hangafter=1\hspace*{1em}conjugation does not determine complements, \hyperref[bk-ch02-learner-d8738489e241]{44}, \hyperref[bk-ch04-learner-0e4e42b78c4b]{99}\par
\par\noindent\hangindent=1em\hangafter=1\hspace*{1em}intransitive uses, \hyperref[bk-ch04-learner-0e4e42b78c4b]{99}, \hyperref[bk-ch05-learner-95f29519f213]{122}, \hyperref[bk-ch09-learner-a6e022a0fd7e]{261}\par
\BookNeedspace{5\baselineskip}
\par\addvspace{6pt}{\large\bfseries U\par}
\BookNeedspace{3\baselineskip}
\par\addvspace{3pt}\noindent\hypertarget{index-f67ed2e92001}{}\textbf{Umlaut}\par\nopagebreak[4]
\par\noindent\hangindent=1em\hangafter=1\hspace*{1em}historical i/j condition, \hyperref[bk-sound-history-learner-b884d4929d35]{776}\par
\par\noindent\hangindent=1em\hangafter=1\hspace*{1em}noun and derivational alternations, \hyperref[bk-sound-history-learner-a00d2110e9a4]{776}, \hyperref[bk-sound-history-learner-64f61747fcac]{777}\par
\par\noindent\hangindent=1em\hangafter=1\hspace*{1em}\emph{See also} \hyperlink{index-9f89be8d7ed9}{Present stem changes}.\par
\BookNeedspace{5\baselineskip}
\par\addvspace{6pt}{\large\bfseries V\par}
\BookNeedspace{3\baselineskip}
\par\addvspace{3pt}\noindent\hypertarget{index-40939b310e14}{}\textbf{Verb-second order}. \emph{See} \hyperlink{index-6cab2d8cd0db}{Word order, finite-second tendency in main clauses}.\par
\BookNeedspace{3\baselineskip}
\par\addvspace{3pt}\noindent\hypertarget{index-4e436562432f}{}\textbf{Verbs}\par\nopagebreak[4]
\par\noindent\hangindent=1em\hangafter=1\hspace*{1em}preterite-present unnan; evidence gap and command substitute, \hyperref[bk-ch10-learner-538269668304]{296}, \hyperref[bk-ch10-learner-4160a02bc5fd]{297}\par
\par\noindent\hangindent=1em\hangafter=1\hspace*{1em}strong and weak formation, \hyperref[bk-ch02-learner-925bbc29948d]{39}, \hyperref[bk-ch04-learner-28ffaf280840]{89}\par
\par\noindent\hangindent=1em\hangafter=1\hspace*{1em}\emph{See also} \hyperlink{index-8ebb9081cbac}{Ablaut}; \hyperlink{index-94de9d0bb1ea}{Imperative}; \hyperlink{index-747f1b861135}{Indicative}; \hyperlink{index-10a0d424f76e}{Infinitives}; \hyperlink{index-94d858470fde}{Irregular verbs}; \hyperlink{index-849bc3ea2c32}{Negation}; \hyperlink{index-8692d6f37e2d}{Participles}; \hyperlink{index-67aacf1d4a8e}{Preterite-present verbs}; \hyperlink{index-a74b1465a508}{Strong verbs}; \hyperlink{index-05ba186b1cc6}{Subjunctive}; \hyperlink{index-2cacf6c14ec8}{Tense}; \hyperlink{index-f67ed2e92001}{Umlaut}; \hyperlink{index-57ac9d925438}{Weak verbs}.\par
\BookNeedspace{3\baselineskip}
\par\addvspace{3pt}\noindent\hypertarget{index-8731b0b78a21}{}\textbf{Verner’s law}\par\nopagebreak[4]
\par\noindent\hangindent=1em\hangafter=1\hspace*{1em}inherited consonant alternation, \hyperref[bk-sound-history-learner-880e35b3cd92]{782}\par
\BookNeedspace{3\baselineskip}
\par\addvspace{3pt}\noindent\hypertarget{index-a2d2c96213de}{}\textbf{Vocative}. \emph{See} \hyperlink{index-bfb78853cd67}{Direct address}.\par
\BookNeedspace{3\baselineskip}
\par\addvspace{3pt}\noindent\hypertarget{index-df9654479f77}{}\textbf{Voicing}. \emph{See} \hyperlink{index-182d0cc13b07}{Consonants, voicing of f, s and þ/ð}.\par
\BookNeedspace{3\baselineskip}
\par\addvspace{3pt}\noindent\hypertarget{index-bb013e68ab47}{}\textbf{Vowel gradation}. \emph{See} \hyperlink{index-8ebb9081cbac}{Ablaut}.\par
\BookNeedspace{3\baselineskip}
\par\addvspace{3pt}\noindent\hypertarget{index-59be352618d8}{}\textbf{Vowel length}\par\nopagebreak[4]
\par\noindent\hangindent=1em\hangafter=1\hspace*{1em}distinct from stress and syllable weight, \hyperref[bk-pronunciation-learner-156b3d5ac4b3]{1}, \hyperref[bk-ch01-learner-737b70e12605]{15}\par
\par\noindent\hangindent=1em\hangafter=1\hspace*{1em}macrons and duration, \hyperref[bk-pronunciation-learner-156b3d5ac4b3]{1}\par
\par\noindent\hangindent=1em\hangafter=1\hspace*{1em}\emph{See also} \hyperlink{index-5b88c0af13ba}{Diphthongs}; \hyperlink{index-2319cd90fbab}{Stress}; \hyperlink{index-aea538abc20e}{Syllable weight}.\par
\BookNeedspace{3\baselineskip}
\par\addvspace{3pt}\noindent\hypertarget{index-1e2e2b6b6efe}{}\textbf{Vowel loss}. \emph{See} \hyperlink{index-c3366c831602}{Apocope}; \hyperlink{index-0c2d377d343c}{Syncope}.\par
\BookNeedspace{3\baselineskip}
\par\addvspace{3pt}\noindent\hypertarget{index-8d1ecf74ebb0}{}\textbf{Vowels}\par\nopagebreak[4]
\par\noindent\hangindent=1em\hangafter=1\hspace*{1em}rounded front vowels, \hyperref[bk-pronunciation-learner-c95193d444e8]{1}\par
\par\noindent\hangindent=1em\hangafter=1\hspace*{1em}simple vowel values, \hyperref[bk-pronunciation-learner-c95193d444e8]{1}\par
\par\noindent\hangindent=1em\hangafter=1\hspace*{1em}stressed a before nasal consonants, \hyperref[bk-pronunciation-learner-c95193d444e8]{1}\par
\par\noindent\hangindent=1em\hangafter=1\hspace*{1em}unstressed vowels remain audible, \hyperref[bk-pronunciation-learner-a18005f2eb82]{7}\par
\par\noindent\hangindent=1em\hangafter=1\hspace*{1em}\emph{See also} \hyperlink{index-5b88c0af13ba}{Diphthongs}; \hyperlink{index-2319cd90fbab}{Stress}; \hyperlink{index-59be352618d8}{Vowel length}.\par
\BookNeedspace{5\baselineskip}
\par\addvspace{6pt}{\large\bfseries W\par}
\BookNeedspace{3\baselineskip}
\par\addvspace{3pt}\noindent\hypertarget{index-57ac9d925438}{}\textbf{Weak verbs}\par\nopagebreak[4]
\par\noindent\hangindent=1em\hangafter=1\hspace*{1em}Class I: consonant clusters before dental endings, \hyperref[bk-ch02-learner-1310ddcd10cb]{40}\par
\par\noindent\hangindent=1em\hangafter=1\hspace*{1em}Class I: irregular preterite stems, \hyperref[bk-ch02-learner-759d9baec3e2]{41}\par
\par\noindent\hangindent=1em\hangafter=1\hspace*{1em}Class I: long stems and dental changes, \hyperref[bk-ch02-learner-fee1ccca58cc]{39}\par
\par\noindent\hangindent=1em\hangafter=1\hspace*{1em}Class I: short stems and linking e, \hyperref[bk-ch02-learner-759d9baec3e2]{41}\par
\par\noindent\hangindent=1em\hangafter=1\hspace*{1em}Class I: stems in rw/lw, \hyperref[bk-ch02-learner-1310ddcd10cb]{40}\par
\par\noindent\hangindent=1em\hangafter=1\hspace*{1em}Class II, \hyperref[bk-ch02-learner-cef52d0e20f6]{42}\par
\par\noindent\hangindent=1em\hangafter=1\hspace*{1em}Class III, \hyperref[bk-ch02-learner-79fe9bfe07cd]{43}\par
\par\noindent\hangindent=1em\hangafter=1\hspace*{1em}dental preterite and basic endings, \hyperref[bk-ch02-learner-925bbc29948d]{39}\par
\par\noindent\hangindent=1em\hangafter=1\hspace*{1em}infinitive ending does not determine class, \hyperref[bk-ch02-learner-cef52d0e20f6]{42}, \hyperref[bk-ch05-learner-aece1b132820]{117}\par
\par\noindent\hangindent=1em\hangafter=1\hspace*{1em}\emph{See also} \hyperlink{index-94de9d0bb1ea}{Imperative}; \hyperlink{index-8692d6f37e2d}{Participles}; \hyperlink{index-05ba186b1cc6}{Subjunctive}.\par
\BookNeedspace{3\baselineskip}
\par\addvspace{3pt}\noindent\hypertarget{index-403cc21502eb}{}\textbf{Wishes}\par\nopagebreak[4]
\par\noindent\hangindent=1em\hangafter=1\hspace*{1em}present subjunctive and exhortation, \hyperref[bk-ch07-learner-bd7d043cab91]{163}\par
\par\noindent\hangindent=1em\hangafter=1\hspace*{1em}unreal present wish with preterite subjunctive, \hyperref[bk-ch09-learner-84672e8deefc]{236}\par
\BookNeedspace{3\baselineskip}
\par\addvspace{3pt}\noindent\hypertarget{index-de7476cd8f8e}{}\textbf{Word formation}\par\nopagebreak[4]
\par\noindent\hangindent=1em\hangafter=1\hspace*{1em}adjectives in -lic, \hyperref[bk-ch14-learner-92d5683c753b]{456}\par
\par\noindent\hangindent=1em\hangafter=1\hspace*{1em}adverbs in -e and -lice, \hyperref[bk-ch14-learner-85c8a677bf69]{456}\par
\par\noindent\hangindent=1em\hangafter=1\hspace*{1em}agent nouns in -end, \hyperref[bk-ch14-learner-70727fa3d47b]{456}\par
\par\noindent\hangindent=1em\hangafter=1\hspace*{1em}compounds and the final head, \hyperref[bk-ch14-learner-79faadb7257c]{455}\par
\par\noindent\hangindent=1em\hangafter=1\hspace*{1em}derivation versus inflection, \hyperref[bk-ch14-learner-7f05ca79b8a2]{455}\par
\par\noindent\hangindent=1em\hangafter=1\hspace*{1em}ge- in lexical verbs and participles, \hyperref[bk-ch14-learner-22f05dfaf32e]{457}\par
\par\noindent\hangindent=1em\hangafter=1\hspace*{1em}noun suffixes: -ung, -scipe, -nes, -dōm, \hyperref[bk-ch14-learner-70727fa3d47b]{456}\par
\par\noindent\hangindent=1em\hangafter=1\hspace*{1em}prefixes and their lexical limits, \hyperref[bk-ch14-learner-dce2679d3490]{457}\par
\par\noindent\hangindent=1em\hangafter=1\hspace*{1em}whole-word meanings and limits of analysis, \hyperref[bk-ch14-learner-719f663802e5]{455}\par
\par\noindent\hangindent=1em\hangafter=1\hspace*{1em}\emph{See also} \hyperlink{index-a47426d0da6f}{Adjectives}; \hyperlink{index-888b2617c1be}{Nouns}.\par
\BookNeedspace{3\baselineskip}
\par\addvspace{3pt}\noindent\hypertarget{index-6cab2d8cd0db}{}\textbf{Word order}\par\nopagebreak[4]
\par\noindent\hangindent=1em\hangafter=1\hspace*{1em}constituents and fronting, \hyperref[bk-ch13-learner-54a05581dff2]{404}\par
\par\noindent\hangindent=1em\hangafter=1\hspace*{1em}dependent content clauses, \hyperref[bk-ch09-learner-8c696664c000]{234}\par
\par\noindent\hangindent=1em\hangafter=1\hspace*{1em}direct questions and clause types, \hyperref[bk-ch08-learner-da4d22fbf997]{203}\par
\par\noindent\hangindent=1em\hangafter=1\hspace*{1em}finite-second tendency in main clauses, \hyperref[bk-ch13-learner-30eeeceef5e9]{404}\par
\par\noindent\hangindent=1em\hangafter=1\hspace*{1em}fronting across a clause boundary (supplied analysis), \hyperref[bk-ch13-learner-52907367739b]{408}\par
\par\noindent\hangindent=1em\hangafter=1\hspace*{1em}late verbs in dependent clauses, \hyperref[bk-ch13-learner-d792ea8f2502]{404}\par
\par\noindent\hangindent=1em\hangafter=1\hspace*{1em}limits on verb-final rules, \hyperref[bk-ch13-learner-672058608ca5]{406}\par
\par\noindent\hangindent=1em\hangafter=1\hspace*{1em}ordinary active clauses, \hyperref[bk-ch02-learner-d8738489e241]{44}\par
\par\noindent\hangindent=1em\hangafter=1\hspace*{1em}subject pronouns after an initial adverbial, \hyperref[bk-ch13-learner-254b66d59004]{405}\par
\end{multicols}

\clearpage\endgroup
\end{document}
